\documentclass[11pt]{amsart}
\usepackage[margin=1.1in]{geometry}
\usepackage{amsmath,amssymb,mathrsfs}
\numberwithin{equation}{section}
\newtheorem{theorem}{Theorem}[section]
\newtheorem{proposition}[theorem]{Proposition}
\newtheorem{lemma}[theorem]{Lemma}
\newtheorem{corollary}[theorem]{Corollary}
\theoremstyle{definition}
\newtheorem{definition}[theorem]{Definition}
\newtheorem{example}[theorem]{Example}
\theoremstyle{remark}
\newtheorem{remark}[theorem]{Remark}
\newcommand{\R}{\mathbb{R}}
\newcommand{\Z}{\mathbb{Z}}
\newcommand{\C}{\mathbb{C}}
\newcommand{\N}{\mathbb{N}}
\newcommand{\Ocub}{O_{\mathrm{cub}}}
\newcommand{\Conf}{\operatorname{Conf}}
\newcommand{\Adm}{\operatorname{Adm}}
\newcommand{\sys}{\operatorname{sys}}
\newcommand{\supp}{\operatorname{supp}}
\newcommand{\Tr}{\operatorname{Tr}}
\newcommand{\idx}{\operatorname{idx}}
\newcommand{\Ins}{\operatorname{Ins}}
\newcommand{\hull}{\operatorname{hull}}
\newcommand{\Gr}{\operatorname{Gr}}
\newcommand{\red}{\operatorname{red}}
\newcommand{\fold}{\operatorname{fold}}
\newcommand{\Fold}{\operatorname{Fold}}
\newcommand{\diag}{\operatorname{diag}}
\newcommand{\End}{\operatorname{End}}
\renewcommand{\Re}{\operatorname{Re}}
\renewcommand{\Im}{\operatorname{Im}}
\newcommand{\Q}{\mathbb{Q}}
\renewcommand{\theenumi}{\alph{enumi}}
\renewcommand{\labelenumi}{(\theenumi)}
\renewcommand{\theenumii}{\roman{enumii}}
\renewcommand{\labelenumii}{(\theenumii)}
\makeatletter
\renewcommand{\p@enumii}{\theenumi.}
\newcommand{\itemmark}[1]{\protected@edef\@currentlabel{#1}}
\makeatother
\begin{document}
\title{Noncompact lattice QED in three dimensions: chamber fermions, reflection positivity at all mirrors and rotation invariance}
\author{}
\date{}

\begin{abstract}
The subject is noncompact lattice quantum electrodynamics in three dimensions with the chamber fermion,
whose sites are the points of one orbit of the symmetry group of $\Z^3$, one site in each chamber cut out by the mirror planes of $\Z^3$,
and whose hopping matrices contain the Wilson parameter $\rho$.
The mirror planes form nine families, $\{x_\mu=c\}$ and $\{x_\mu\pm x_\nu=c\}$.
For $-1\le\rho\le1$ and each of these families, the weight of photon and fermion together satisfies reflection positivity at every mirror of the family,
on every torus adapted to it.
For $0<\rho<1$, every real mass and arbitrary sides tending to infinity, the laws of the curvature on cubic tori form a tight sequence,
and each weak limit point obeys Gaussian domination with the constant one,
has the symmetries of the lattice and satisfies reflection positivity at the mirrors of all nine families.
In the continuum, for $0<\rho<1$, at every charge and for every choice of the lattice masses,
let the volume tend to infinity first and the spacing to zero afterwards.
Each limit point obtained in this way, and almost each of its ergodic components,
is left invariant by all Euclidean motions of $\R^3$, reflections included, and is reflection positive with respect to every plane;
the joint spectrum of energy and momentum lies in the round forward cone;
the two-point function has the K\"all\'en--Lehmann form, with a spectral measure of total mass at most one;
and the Schwinger functions satisfy the Osterwalder--Schrader axioms of linear growth, Euclidean invariance, positivity and symmetry,
together with clustering, at a rate for tests of compact support, in almost every ergodic component.
The results contain no lower bound: the point mass at the zero field has all of the properties stated.
\end{abstract}
\maketitle

\section{Introduction}\label{sec:introduction}

This paper concerns lattice quantum electrodynamics in three dimensions,
with a noncompact photon field on the links of $\Z^3$ and all flux sectors summed,
and the Euclidean symmetries of its limits.
The isometries of $\R^3$ that map $\Z^3$ onto itself form the group $W=\Z^3\rtimes\Ocub$, where $\Ocub$ is the cubic group of order $48$.
The group $W$ is generated by the reflections that it contains,
the reflections in the planes $\{x_\mu=c\}$, $c\in\frac12\Z$, and $\{x_\mu\pm x_\nu=c\}$, $c\in\Z$;
these planes are the mirrors, and they fall into three axis families and six diagonal families,
whose members are the axis mirrors and the diagonal mirrors.
Reflection positivity is the property from which a Hilbert space with a transfer operator, on the lattice,
and a Hilbert space with a non-negative Hamiltonian, in the continuum, are reconstructed.
Two questions are treated:
whether the infinite volume limits of the lattice theory on cubic tori are reflection positive at the mirrors of all nine families,
and whether the continuum limit points have the full Euclidean symmetry of $\R^3$, with reflection positivity at every plane.

The fermion is the chamber fermion.
It has one site in each chamber of the arrangement of the mirrors, so that no site lies on a mirror,
and each bond joins a site to its mirror image in a wall of its chamber.
At every mirror of the nine families the bonds that cross the mirror are orthogonal to it,
and the crossing part of the action is a sum of terms of positive type,
with coefficients that are non-negative for $-1\le\rho\le1$, $\rho$ the Wilson parameter (Proposition~\ref{prop:crossing}).
On a torus adapted to a family, a notion defined below, two mirrors of this family cut the torus into two halves that the reflection exchanges;
the cubic torus is adapted to the three axis families and to no diagonal family (Proposition~\ref{prop:untwistings}).

On every torus adapted to a family, the weight of photon and fermion together satisfies reflection positivity at each mirror of that family
(Theorem~\ref{thm:adapted-positivity}).
When the sides of cubic tori tend to infinity along an arbitrary sequence,
each weak limit point of the laws of the curvature has reflection positivity at every mirror of all nine families (Theorem~\ref{thm:cubic-limits}).
In the continuum, where the volume is sent to infinity before the spacing is sent to zero,
the limit points are reflection positive at the planes of the nine families and obey the cone inequality with the constant one (Theorem~\ref{thm:continuum-limits});
they are invariant under all Euclidean motions of $\R^3$ and reflection positive at every plane,
and the Schwinger functions of the curvature satisfy the axioms of Osterwalder and Schrader, the cluster property in almost every ergodic component
(Theorem~\ref{thm:euclidean}).

\subsection{The lattice theory}\label{sec:intro-lattice}

\emph{The photon field.}
The photon lattice is $\Z^3$, with its cells: sites, links, plaquettes and cubes.
The space of the real $k$-cochains is $C^k$, $d$ is the coboundary and $d^{*}$ its adjoint on finitely supported cochains;
an isometry $g$ of $\R^3$ with $g\Z^3=\Z^3$ acts on cochains by the pull-back $g^{*}$,
and $\tau_v$ is the translation by $v\in\Z^3$ (Definition~\ref{def:cochains}).
Let $q>0$ be the lattice charge.
A period lattice is a subgroup $\Gamma$ of $\Z^3$ of rank three with $\sys(\Gamma)=\min\{|g|_\infty:\ g\in\Gamma,\ g\neq0\}\ge4$.
A configuration on $\Z^3/\Gamma$ is a class of pairs $(A,z)$ of a real one-cochain $A$ and a closed integer two-cochain $z$, both $\Gamma$-periodic,
modulo $(A,z)\sim\left(A+\frac{2\pi}{q}\cdot k,\ z-dk\right)$ for integer $\Gamma$-periodic one-cochains $k$;
the configurations form the space $\Conf(\Gamma)$,
with the measure $\mathrm{d}(A,z)$ for which $\int_{\Conf(\Gamma)}\mathrm{d}(A,z)\,\Phi(A,z)$ is the sum over the closed $z$
of the integrals of $\Phi(A,z)$ over the values $A(b)\in[0,2\pi/q)$ at the links $b$ of $\Z^3/\Gamma$ (Definition~\ref{def:configurations}).
The curvature
\[
F:=dA+\frac{2\pi}{q}\cdot z
\]
is a closed real two-cochain.
The photon field is noncompact, and the sum over the closed $z$ includes every flux sector.
The link phases are the functions $\exp(iqA(b))$, $b$ a link,
and the periodic extension $\widetilde F$ of $F$ lies in the Polish space $\Omega$ of the closed real two-cochains on $\Z^3$.
A law on $\Omega$, that is, a Borel probability measure $\mu$ on $\Omega$, is said to satisfy
Gaussian domination with the constant one, \eqref{cond:gaussian-domination}, when
$\int\mathrm{d}\mu(F)\,\exp\langle\beta,F\rangle\le\exp\left(\frac12\|\beta\|^2\right)$ for every finitely supported real two-cochain $\beta$ (Definition~\ref{def:characters}).

\emph{Mirrors and chambers.}
The nine families are the linear forms $\ell=x_\mu$, the axis families, and $\ell=x_\mu\pm x_\nu$, $\mu<\nu$, the diagonal families;
they form the set $\mathcal{L}$.
For $\ell\in\mathcal{L}$ let $n_\ell\in\Z^3$ be the vector with $\ell(x)=n_\ell\cdot x$, and $\hat n_\ell:=n_\ell/|n_\ell|$ the upward unit normal.
A mirror is a pair $\Pi=(\ell,c)$ with $c\in\Adm(\ell)$,
where $\Adm(\ell)=\frac12\Z$ for an axis family, $c\in\Z$ giving a site plane and $c\in\Z+\frac12$ a link plane,
and $\Adm(\ell)=\Z$ for a diagonal family;
the reflection in the plane $\{\ell=c\}$ is $\theta_\Pi=\theta_{\ell,c}$.
These reflections generate $W=\Z^3\rtimes\Ocub$, with $\Ocub=\mathrm{O}(3)\cap\mathrm{GL}(3,\Z)$ (Lemma~\ref{lem:chambers});
an element $g\in W$ is a map $x\mapsto k+\vec g\,x$ with $k\in\Z^3$ and the linear part $\vec g\in\Ocub$ (Definition~\ref{def:mirrors}).
Two diagonal families have names, $\ell_\pm:=x_1\pm x_2$, and $e_\pm:=e_1\pm e_2$.
The complement of the union of the mirrors consists of open simplices, the chambers, and $W$ permutes the chambers simply transitively;
the fundamental chamber is $C_0=\left\{0<x_3<x_2<x_1<\frac12\right\}$, with the walls $\Pi_0,\dots,\Pi_3$ (Definition~\ref{def:chambers}).

\emph{The chamber fermion.}
Fix $v_0\in C_0$.
The fermion sites form the orbit $\mathbb{X}:=Wv_0$:
one site in each chamber, $48$ sites in each unit cube, none on a mirror.
A site $x=gv_0$, $g\in W$, has the neighbours $x^{(i)}:=gs_iv_0$, $i=0,\dots,3$, where $s_i$ is the reflection in $\Pi_i$;
so $x^{(i)}$ is the reflection of $x$ in the wall $g\Pi_i$ of the chamber $gC_0$ that contains $x$.
The bond $\{x,x^{(i)}\}$ has the type $i$, and $\pi(x)\in\Z^3$ is the nearest photon site (Definition~\ref{def:fermion-sites}).
A bond joins two different photon sites exactly when its type is $3$, that is, when it crosses a link plane,
and exactly eight bonds lie over each photon link (Proposition~\ref{prop:geometry}).
Let $\gamma_\mu=\diag(\sigma_\mu,-\sigma_\mu)$, $\sigma_\mu$ the Pauli matrices, be the Hermitian matrices of Definition~\ref{def:spin-matrices},
and put $\gamma_v:=\gamma\cdot v:=\sum_\mu v_\mu\gamma_\mu$ for $v\in\R^3$.
With hopping weights $w_0,\dots,w_3>0$, a real mass parameter $M$ and the Wilson parameter $\rho$,
the chamber operator acts on functions $\psi\colon\mathbb{X}\to\C^4$ by
\[
(D_A\psi)(x)=M\psi(x)-\sum_{i=0}^{3}w(x,x^{(i)})\cdot U_A(x,x^{(i)})\cdot\psi(x^{(i)}),\qquad
w(x,x'):=\frac{w_i}{2}\left(\rho-\gamma\cdot\frac{x'-x}{|x'-x|}\right),
\]
where $U_A(x,x'):=\exp\left(iqA(\pi(x),\pi(x'))\right)$ and $i$ is the type of the bond $\{x,x'\}$ (Definition~\ref{def:chamber-operator}).
The hopping matrix $w(x,x')$ has the eigenvalues $\frac{w_i}{2}(\rho-1)$ and $\frac{w_i}{2}(\rho+1)$, each twice;
it is invertible for $\rho^2\neq1$,
and the crossing coefficients $\frac{w_i}{2}(1\pm\rho)$ are non-negative for $|\rho|\le1$ and positive for $|\rho|<1$
(Lemma~\ref{lem:hopping}, Proposition~\ref{prop:crossing}).
The lattice mass is $m:=M-\frac{\rho}{2}\sum_iw_i$,
and the standard weights are $w_i=48\cdot|\Pi_i\cap\overline{C_0}|$, forty-eight times the areas of the walls of $C_0$.

\emph{Grassmann variables and tori.}
For $Y\subset\mathbb{X}$, the complex Grassmann algebra $\Gr(Y)$ has the generators $\psi_\alpha(x)$, $\bar\psi_\alpha(x)$, $x\in Y$, $\alpha=1,\dots,4$;
for finite $Y$ the Berezin integral $\int_Y$ is normalised so that $\int_Y\exp(-\bar\psi D\psi)=\det D$ \eqref{eq:setting-berezin}.
A torus is a pair $\mathbb{T}=(\Gamma,\chi)$ of a period lattice $\Gamma$ and a character $\chi\colon\Gamma\to\{1,-1\}$,
the boundary condition of the fermion:
on $\mathbb{T}$ the fields obey $\psi(x+g)=\chi(g)\cdot\psi(x)$ and $\bar\psi(x+g)=\chi(g)\cdot\bar\psi(x)$, $g\in\Gamma$,
and they generate the Grassmann algebra $\Gr_{\mathbb{T}}$,
in which the generators at a set of representatives of $\mathbb{X}/\Gamma$ are free generators.
For a real $\Gamma$-periodic one-cochain $A$, the action $S_{\mathbb{T}}(A)$ is $\sum_x\bar\psi(x)\,(D_A\psi)(x)$, the sum over such a set of representatives;
$\int_{\mathbb{T}}$ is the Berezin integral in the free generators at these representatives;
and $D^{\mathbb{T}}_A$ is the operator $D_A$ on the spinors of $\mathbb{T}$, the functions $\psi\colon\mathbb{X}\to\C^4$ with $\psi(x+g)=\chi(g)\cdot\psi(x)$ for $g\in\Gamma$.
Then $\int_{\mathbb{T}}\exp\left(-S_{\mathbb{T}}(A)\right)=\det D^{\mathbb{T}}_A$ (Definition~\ref{def:tori}).
The image of $\mathbb{T}$ under $g\in W$ is the torus $g\mathbb{T}:=\left(\vec g\,\Gamma,\chi\circ\vec g^{-1}\right)$.
Three kinds of tori have names:
the coordinate torus $C(n_1,n_2,n_3)$, with the periods $n_\mu e_\mu$, $n_\mu\ge4$, and $\chi=-1$ on them,
which is cubic if $n_1=n_2=n_3=n$ and is then written $C(n)$;
the torus $D(a,b,c)$, with the periods $a\,e_-$, $b\,e_+$, $c\,e_3$, $\min(a,b,c)\ge4$, and $\chi=-1$ on them;
and the torus $D'(a,b,c)$, with the same periods, $\chi(a\,e_-)=\chi(c\,e_3)=-1$ and $\chi(b\,e_+)=1$.

\emph{Reflections.}
Let $\Pi=(\ell,c)$ be a mirror, $\theta:=\theta_\Pi$ and $n:=\hat n_\ell$.
The reflection $\Theta_\Pi$ of Grassmann variables is antilinear, reverses the order of products, and
\[
\Theta_\Pi\psi(x)=-\bar\psi(\theta x)\,\gamma_n,\qquad
\Theta_\Pi\bar\psi(x)=-\gamma_n\,\psi(\theta x)
\]
(Definition~\ref{def:grassmann-reflections});
on coefficient functions it acts by complex conjugation composed with $\theta^{*}$ \eqref{eq:setting-reflection-coefficients}.
On a torus it maps $\Gr_{\mathbb{T}}$ onto $\Gr_{\theta\mathbb{T}}$ (Lemma~\ref{lem:reflected-algebras}).

\emph{The functional and the law.}
Let $\mathcal{A}(\R^3)$ be the class of the finite sums of products of a bounded Borel function of finitely many values $F_\sigma$, $\sigma$ a plaquette,
with even, gauge invariant, local integrands that have Gaussian growth of a constant below $\frac12$ in pairwise disjoint boxes;
a box is a closed cube $y+[-r,r]^3$, with the centre $y\in\R^3$ and the radius $r>0$ (Definition~\ref{def:integrands}).
For $G\in\mathcal{A}(\R^3)$ that embeds in the torus $\mathbb{T}=(\Gamma,\chi)$,
\[
Z_{\mathbb{T}}[G]:=\int_{\Conf(\Gamma)}\mathrm{d}(A,z)\,\exp\left(-\frac12\|F\|^2\right)\cdot\int_{\mathbb{T}}\exp\left(-S_{\mathbb{T}}(A)\right)\cdot G(A,z)
\]
(Definition~\ref{def:functional}):
the photon weight is the Maxwell weight, and integration over the fermion gives the determinant $\det D^{\mathbb{T}}_A$,
which is non-negative for every real one-cochain $A$ (Proposition~\ref{prop:accretivity}).
For every torus, every real $\rho$ and $M$, all weights and every charge, $0<Z_{\mathbb{T}}[1]<\infty$ (Theorem~\ref{thm:partition-positivity}).
So $\mathbb{E}_{\mathbb{T}}:=Z_{\mathbb{T}}[\,\cdot\,]/Z_{\mathbb{T}}[1]$ is defined,
and the law $\mu_{\mathbb{T}}$ of $\widetilde F$ under $\mathbb{E}_{\mathbb{T}}$ is a law on $\Omega$.
The integral that defines $Z_{\mathbb{T}}$ converges absolutely also on the algebra $\mathcal{B}(\mathbb{T})$ of the finite sums $\sum_j\Phi_j\xi_j$
with bounded Borel functions $\Phi_j$ on $\Conf(\Gamma)$ and $\xi_j\in\Gr_{\mathbb{T}}$ (Lemma~\ref{lem:bounded-integrands}).

\emph{Reflection positivity of laws.}
Let $(\ell,c)$ be a mirror.
A law $\mu$ on $\Omega$ is invariant at $(\ell,c)$ if the map $F\mapsto\theta_{\ell,c}^{*}F$ preserves it,
and it is reflection positive at $(\ell,c)$ if
\[
\int\mathrm{d}\mu(F)\,\overline{\Phi\left(\theta_{\ell,c}^{*}F\right)}\cdot\Phi(F)\ge0
\]
for every bounded Borel function $\Phi$ of finitely many values $F_\sigma$, $\sigma$ ranging over the plaquettes all of whose corners satisfy $\ell>c$
(Definition~\ref{def:reflection-positivity}).
Only the open side is read; nothing that touches the mirror enters.

\emph{Adapted tori.}
Let $\ell\in\mathcal{L}$ and let $\mathbb{T}=(\Gamma,\chi)$ be a torus.
A closing vector of $\mathbb{T}$ for $\ell$ is an element $g\in\Gamma$ with $\Gamma=(\Gamma\cap\{\ell=0\})+\Z g$ and $\ell(g)>0$;
the number $p=p(\mathbb{T},\ell):=\ell(g)$ does not depend on $g$ and is the height period.
The torus $\mathbb{T}$ is adapted to $\ell$ if the vector $v:=\frac{p}{|n_\ell|^2}\cdot n_\ell$ is integral,
$\Gamma=(\Gamma\cap\{\ell=0\})\oplus\Z v$, and $\chi(v)=-1$ (Definition~\ref{def:slabs});
then $v$ is a closing vector.
For such a torus and $c\in\Adm(\ell)$ the position $c+\frac p2$ is admissible (Definition~\ref{def:split-tori}\eqref{it:split-height}).
The linear part of $\theta_{\ell,c}$ fixes the plane $\{\ell=0\}$ pointwise and maps $v$ to $-v$;
hence it maps $\Gamma$ onto itself and preserves $\chi$, so that $\theta_{\ell,c}\mathbb{T}=\mathbb{T}$ and $\Theta_{\ell,c}$ maps $\Gr_{\mathbb{T}}$ onto itself.
Moreover, the reflections $\theta_{\ell,c}$ and $\theta_{\ell,c+p/2}=\tau_v\theta_{\ell,c}$ induce one and the same involution of $\R^3/\Gamma$,
which exchanges the images of the open slabs $\left\{c<\ell<c+\frac p2\right\}$ and $\left\{c+\frac p2<\ell<c+p\right\}$.
A coordinate torus $C(n_1,n_2,n_3)$ is adapted to $x_1$, $x_2$ and $x_3$, and $D(a,b,c)$ to $\ell_-$, $\ell_+$ and $x_3$ (Definition~\ref{def:cell-structures});
applying the elements of $\Ocub$ to these tori gives, for each of the nine families, tori adapted to it.

\subsection{The main theorems}\label{sec:intro-theorems}

\emph{Integrands of one half of an adapted torus.}
Let $\ell\in\mathcal{L}$, let $\mathbb{T}=(\Gamma,\chi)$ be a torus adapted to $\ell$ with the height period $p$,
let $c\in\Adm(\ell)$ and $c':=c+\frac p2$.
Let $\Gr_{\mathbb{T}}^{+}$ be the subalgebra of $\Gr_{\mathbb{T}}$ generated by the fields at the sites in the open slab $\{c<\ell<c'\}$.
Then $\mathcal{B}^{+}=\mathcal{B}^{+}(\mathbb{T},\ell,c)$ is the linear span of the products $\Phi\,\xi$ with $\xi\in\Gr_{\mathbb{T}}^{+}$
and $\Phi$ a bounded Borel function of the values $F_\sigma$ and of the link phases $\exp(iqA(b))$
of the plaquettes $\sigma$ and the links $b$ that are contained, modulo $\Gamma$, in the closed slab $\{c\le\ell\le c'\}$;
and $\mathcal{G}^{+}=\mathcal{G}^{+}(\mathbb{T},\ell,c)$ is the linear span of the products $\Phi\,\xi$ with $\xi\in\Gr_{\mathbb{T}}^{+}$
and $\Phi$ a Borel function of the same variables for which there are $C<\infty$ and $\varepsilon\in\left(0,\frac12\right]$ with
\[
|\Phi|\le C\cdot\exp\left(\left(\frac12-\varepsilon\right)\sum_\sigma F_\sigma^2\right),
\]
the sum being taken over the plaquettes whose relative interior lies, modulo $\Gamma$, in the open slab
(Definition~\ref{def:upper-integrands}).
So $\Phi$ may grow like a Gaussian with a constant below $\frac12$ in the values at these plaquettes,
but not at all in the values at the plaquettes contained in the two mirrors, which do not occur in the bound.
Thus $\mathcal{B}^{+}\subset\mathcal{G}^{+}$.
Finally, for the cross-section $\Xi=\Xi(\mathbb{T},\ell)$ of $\mathbb{T}$, formed with the lattice $\Gamma\cap\{\ell=0\}$ and the restriction of $\chi$ to it,
$\mathcal{A}_\Xi[c,c']$ is the linear span of the products of $\mathcal{A}(\{c<\ell<c'\})$ that embed in $\Gamma\cap\{\ell=0\}$ (Definition~\ref{def:slabs}):
their boxes and the plaquettes on which their bounded factors depend lie in the open slab.

\begin{theorem}[Reflection positivity on adapted tori]\label{thm:adapted-positivity}
Let $-1\le\rho\le1$, $M\in\R$, $w_0,\dots,w_3>0$ and $q>0$.
Let $\ell\in\mathcal{L}$, let $\mathbb{T}$ be a torus adapted to $\ell$ with the height period $p$, and let $c\in\Adm(\ell)$;
then $c':=c+\frac p2$ lies in $\Adm(\ell)$.
\begin{enumerate}
\item\label{it:adapted-invariance}
$Z_{\mathbb{T}}[\Theta_{\ell,c}G'']=\overline{Z_{\mathbb{T}}[G'']}$ for all $G''\in\mathcal{B}(\mathbb{T})$.
\item\label{it:adapted-representation}
There are a finite positive measure space $(\mathcal{Z},\mathfrak{m})$
and a linear map $G\mapsto\hat G$ from $\mathcal{B}^{+}$ to the bounded measurable functions on $\mathcal{Z}$ with
\[
Z_{\mathbb{T}}\left[\Theta_{\ell,c}(G')\,G\right]=\int_{\mathcal{Z}}\mathrm{d}\mathfrak{m}\,\overline{\widehat{G'}}\cdot\hat G\qquad\left(G,G'\in\mathcal{B}^{+}\right),
\]
whether $G$ and $G'$ are even, odd or of mixed parity.
\item\label{it:adapted-form}
Hence $(G',G)\mapsto Z_{\mathbb{T}}[\Theta_{\ell,c}(G')\,G]$ is a positive semidefinite Hermitian form on $\mathcal{B}^{+}$;
in particular $Z_{\mathbb{T}}[\Theta_{\ell,c}(G)\,G]\ge0$ for $G\in\mathcal{B}^{+}$.
\item\label{it:adapted-growth}
For $G,G'\in\mathcal{G}^{+}$ the integral that defines $Z_{\mathbb{T}}[\Theta_{\ell,c}(G')\,G]$ converges absolutely, coefficient by coefficient;
\eqref{it:adapted-representation} and \eqref{it:adapted-form} hold on $\mathcal{G}^{+}$,
with the same $(\mathcal{Z},\mathfrak{m})$ and the same formula \eqref{eq:chamber-transform} for $\hat G$, and $\hat G$ is bounded.
The class $\mathcal{A}_\Xi[c,c']$, $\Xi=\Xi(\mathbb{T},\ell)$, is contained in $\mathcal{G}^{+}$;
in particular $(G',G)\mapsto Z_{\mathbb{T}}[\Theta_{\ell,c}(G')\,G]$ is a positive semidefinite Hermitian form on $\mathcal{A}_\Xi[c,c']$.
\end{enumerate}
In particular, every finite sum $G$ of products of a bounded Borel function of $F$ and of the link phases, read in the open slab $\{c<\ell<c'\}$,
with an element, even or odd, of the Grassmann algebra of the sites in this slab lies in $\mathcal{B}^{+}$,
and $Z_{\mathbb{T}}[\Theta_{\ell,c}(G)\,G]\ge0$.
\end{theorem}

On the class $\mathcal{A}_\Xi[c,c']$ the form has a second representation, as a trace,
\[
Z_{\mathbb{T}}\left[\Theta_{\ell,c}(G')\,G\right]=\Tr\left(\mathfrak{S}[G']^{*}\,\mathfrak{S}[G]\right),
\]
with Hilbert--Schmidt slab operators $\mathfrak{S}[\,\cdot\,]$ (Corollary~\ref{cor:trace-forms}).
Every family has adapted tori among the images under $\Ocub$ of $C(n_1,n_2,n_3)$ and of $D(a,b,c)$,
and Theorem~\ref{thm:adapted-positivity} applies on each torus adapted to it.
The normals of two different families to which one torus is adapted are orthogonal (Proposition~\ref{prop:untwistings}),
whereas the unit normals of an axis family $x_\mu$ and of a diagonal family in which $x_\mu$ occurs have a scalar product of modulus $2^{-\frac12}$:
their mirrors meet at the angle of $45$ degrees.
Hence the cubic torus $C(n)$, which is adapted to the three axis families, is adapted to no diagonal family,
and Theorem~\ref{thm:adapted-positivity} says nothing about the diagonal mirrors on $C(n)$.
The laws of the cubic tori reach the diagonal mirrors through a chain of four tori of equal volume whose pressures agree in the limit
(Theorem~\ref{thm:four-tori}).

\begin{theorem}[Volume limits of cubic tori]\label{thm:cubic-limits}
Let $0<\rho<1$, $w_0,\dots,w_3>0$, $M\in\R$ and $q>0$.
Let $\mathfrak{n}=(n_i)$ be an arbitrary sequence of integers $n_i\ge4$ with $n_i\to\infty$,
and let $\mu_i:=\mu_{C(n_i)}$ be the law of $\widetilde F$ on the cubic torus $C(n_i)=C(n_i,n_i,n_i)$.
Then the sequence $(\mu_i)$ is tight on $\Omega$,
and every weak limit point $\nu$ satisfies \eqref{cond:gaussian-domination} and $\nu(dF=0)=1$,
is invariant under $W$ and under $F\mapsto-F$,
and is reflection positive (Definition~\ref{def:reflection-positivity}) at every mirror of the nine families.
\end{theorem}

The theorem holds for every real mass parameter $M$ (Theorem~\ref{thm:partition-positivity})
and along every sequence of sides, even or odd, with or without common divisors.

\emph{The continuum.}
Let $L\ge3$ be an odd integer and $0<\bar m\le1$.
For a cutoff, that is, a pair of integers $N,T\ge0$, put
\[
\eta:=L^{-N},\qquad n:=8\cdot L^{\lfloor\log_L\bar m^{-1}\rfloor}\cdot L^{T+N},\qquad R:=n\cdot\eta ;
\]
thus $n$ is even.
Let $e>0$ be the charge and $q:=e\cdot\eta^{\frac12}$ the lattice charge.
A mass assignment is an arbitrary sequence $N\mapsto m_{0,N}$ of real numbers;
with $m_N:=\bar m\cdot\eta+m_{0,N}$, the chamber operator at the cutoff $(N,T)$ has the mass parameter $M=m_N+\frac{\rho}{2}\sum_iw_i$,
so that its lattice mass is $m_N$.
With $\rho\in(0,1)$, a point $v_0$ of the fundamental chamber and the standard weights fixed, let $F^{\mathrm{u}}$ have the law $\mu_{C(n)}$.
The law $P^{\mathrm{ch}}_{N,T}$ is the law of the step two-form $F^\sharp$ of $\eta^{-\frac32}\cdot F^{\mathrm{u}}$:
for every plaquette $\sigma=(x;\mu\nu)$ of $\Z^3$, the $\mu\nu$-component of $F^\sharp$ equals $\eta^{-\frac32}\cdot F^{\mathrm{u}}(\sigma)$
on the cube $\eta\cdot\left(x+\frac12(e_\mu+e_\nu)+\left[-\frac12,\frac12\right)^3\right)$ (Definition~\ref{def:cutoffs}).
It is a law on the space $\mathcal{S}'=\mathcal{S}'(\R^3;\Lambda^2)$ of the distributions with values in two-forms,
carried by the Hermite--Sobolev space $H_{-3}\subset H_{-4}$ of \cite[Sections~2.4 and~7.1.5]{NoncompactQED}, with the norms $\|\cdot\|_s$.
A limit point of the family $(P^{\mathrm{ch}}_{N,T})$ is a weak limit on $H_{-4}$ along a sequence of cutoffs $(N_i,T_i)$ with $N_i\to\infty$ and $T_i\to\infty$.
For fixed $N$, the weak limits on $H_{-4}$ of $P^{\mathrm{ch}}_{N,T_i}$ along sequences $T_i\to\infty$ form the set $\mathcal{V}_N$ of the volume limits at the spacing $L^{-N}$,
and the weak limits on $H_{-4}$ of sequences $\mu^{(N_i)}\in\mathcal{V}_{N_i}$ with $N_i\to\infty$ are the volume-first limit points:
for them the volume is sent to infinity at a fixed spacing before the spacing is sent to zero (Definition~\ref{def:limit-points}).
They are the iterated limits taken with the volume first.
An iterated limit taken in the opposite order is obtained from a sequence $T_i\to\infty$:
for each $i$ one takes a weak limit on $H_{-4}$ of $P^{\mathrm{ch}}_{N,T_i}$ along a sequence of spacing indices $N\to\infty$,
and then a weak limit on $H_{-4}$ of these laws as $i\to\infty$.

\emph{Laws on $\mathcal{S}'$.}
A law is a probability measure $\mu$ on $\mathcal{S}'$ with the $\sigma$-algebra $\Sigma$ generated by the functions $F\mapsto F(f)$, $f$ a real test form;
for an open set $U\subset\R^3$, $\Sigma_U$ is generated by the $F(f)$ with $f$ real, smooth and of compact support in $U$;
$\mu(\Phi):=\int\mathrm{d}\mu(F)\,\Phi(F)$;
on $L^2(\mu)$, $\langle\Phi,\Phi'\rangle_\mu:=\mu\left(\overline{\Phi}\cdot\Phi'\right)$ and $\|\Phi\|_2:=\langle\Phi,\Phi\rangle_\mu^{\frac12}$.
A Euclidean map $\mathfrak{g}\colon x\mapsto gx+a$, $g\in\mathrm{O}(3)$, acts on two-forms by $(\mathfrak{g}\cdot f)(x)=(\wedge^2g)\,f\left(\mathfrak{g}^{-1}x\right)$,
on $\mathcal{S}'$ by $(\mathfrak{g}F)(f)=F\left(\mathfrak{g}^{-1}\cdot f\right)$,
on functions by $(\hat{\mathfrak{g}}\Phi)(F)=\Phi\left(\mathfrak{g}^{-1}F\right)$, and on laws by $\mathfrak{g}_*$;
$U(a):=\hat\tau_a$ for the translation $\tau_a$ by $a\in\R^3$,
$\theta_j$ is the reflection in the plane $\{x_j=0\}$ and $\Theta_j\Phi:=\overline{\hat\theta_j\Phi}$,
and $r_{n,c}$ is the reflection in the plane $\{n\cdot x=c\}$, $n$ a unit vector (Definition~\ref{def:laws}, \cite[Section~8.1.2]{NoncompactQED});
for $\ell\in\mathcal{L}$ and $c\in\R$ the reflection in the plane $\{\ell=c\}$ is $r_{\ell,c}:=r_{\hat n_\ell,c/|n_\ell|}$.
The conditions on laws are those of \cite[Sections~8.1.3 and~8.1.5]{NoncompactQED}:
Gaussian domination \textup{(G)}, $\int\mathrm{d}\mu(F)\,\exp F(f)\le\exp\left(\frac12\|f\|_2^2\right)$ for real test forms $f$;
the invariance \eqref{cond:translation-invariance} under the translations along the directions in $J\subset\{1,2,3\}$;
the invariance \eqref{cond:reflection-invariance} under $\theta_j$;
the reflection positivity \eqref{cond:reflection-positivity} at $\{x_j=0\}$;
and the hypothesis \eqref{cond:soft-hypothesis}: \textup{(G)}, invariance under the translations of $\R^3$ and under $\Ocub$, and $(RP_j)$ for $j=1,2,3$.
A law $\mu$ is reflection positive at the plane $\{n\cdot x=c\}$ if
\[
\sum_{\alpha,\alpha'}\overline{z_\alpha}\cdot z_{\alpha'}\cdot\mu\left(\exp\left(iF(f_{\alpha'})-iF(r_{n,c}\cdot f_\alpha)\right)\right)\ge0
\]
for finitely many real smooth test forms $f_\alpha$ with compact supports in $\{n\cdot x>c\}$ and all $z_\alpha\in\C$.
The planes $\{\ell=c\}$, $\ell\in\mathcal{L}$, $c\in\R$, form the nine families of planes, the continuum counterparts of the nine families of mirrors.
The class $\mathcal{C}_9$ consists of the laws that are invariant under the translations of $\R^3$
and, at every plane of the nine families, invariant under the reflection in this plane and reflection positive at it (Definition~\ref{def:laws}).
A law satisfies \eqref{cond:closedness} if it is carried by the closed distributions, $\mu(dF=0)=1$:
$F(f)=0$ whenever the dual vector field $\ast f:=(f_{23},-f_{13},f_{12})$ of $f$ is the gradient of a Schwartz function \cite[Section~8.1.6]{NoncompactQED}.

\emph{Reconstruction.}
Let $\mu$ satisfy $(T_{123})$, $(R_j)$ and $(RP_j)$, and put $\Sigma_j^+:=\Sigma_{\{x_j>0\}}$.
The reflected form $(\Phi,\Phi')\mapsto\mu\left((\Theta_j\Phi)\cdot\Phi'\right)$ is positive semidefinite on $L^2(\Sigma_j^+,\mu)$;
the completion of the quotient by its null space is the Hilbert space $\mathcal{H}_j$, with the class map $[\,\cdot\,]_j$ and $\Omega_j:=[1]_j$.
There are a self-adjoint operator $\mathsf{H}_j\ge0$ with $\exp(-t\cdot\mathsf{H}_j)[\Phi]_j=[U(t\cdot e_j)\Phi]_j$, $t\ge0$,
and a strongly continuous unitary representation $V_j$ of $e_j^\perp$ with $V_j(y)[\Phi]_j=[U(y)\Phi]_j$
\cite[Proposition~8.19]{NoncompactQED};
for $i\neq j$, $\mathsf{P}_i$ is the generator of $u\mapsto V_j(u\cdot e_i)$.
This is the reconstruction of $\mu$ in the direction $j$ (Definition~\ref{def:laws}), which is also called the picture $e_j$.
The joint spectral measure $\mathfrak{E}_j$ of $(\mathsf{H}_j,(\mathsf{P}_i)_{i\neq j})$ lives on $[0,\infty)\times e_j^\perp$, with the variables $(\lambda,k)$, $k=(k_i)_{i\neq j}$,
and
\[
\exp(-t\cdot\mathsf{H}_j)\,V_j(y)=\int\mathrm{d}\mathfrak{E}_j(\lambda,k)\,\exp(-t\cdot\lambda+i\,k\cdot y)\qquad\left(t\ge0,\ y\in e_j^\perp\right);
\]
its image under $(\lambda,k)\mapsto(\lambda,k_i)$ is the joint spectral measure of $(\mathsf{H}_j,\mathsf{P}_i)$ (Lemma~\ref{lem:joint-spectrum}).
For $J\subset\{1,2,3\}$, $P^J_\mu$ is the orthogonal projection of $L^2(\mu)$ onto the vectors that are invariant under the translations along the directions in $J$.
A law $\mu$ with \eqref{cond:soft-hypothesis} has the $\R^3$-ergodic decomposition $\mu=\int\mathrm{d}\nu(\xi)\,\mu_\xi$
(Definition~\ref{def:symmetry-conditions}, \cite[Theorem~8.1]{NoncompactQED}).

\emph{Moments and tensor tests.}
Let $\mu$ be a law with \textup{(G)}.
The function $F\mapsto F(f)$ extends complex linearly to all $f$ in the space $L^2(\R^3;\Lambda^2)$ of the square integrable two-forms with complex coefficients,
with values in every $L^{s}(\mu)$, $s<\infty$ (Lemma~\ref{lem:domination-consequences}),
and the moments $S_n(f_1,\dots,f_n):=\mu\left(F(f_1)\cdots F(f_n)\right)$ are defined for $f_1,\dots,f_n\in L^2(\R^3;\Lambda^2)$.
For $n\ge1$ let $\mathscr{S}_n$ be the space of the functions on $\R^{3n}$ with values in the complexification of $(\Lambda^2)^{\otimes n}$ whose components are Schwartz functions,
with the seminorms $q_m$ of \cite[Section~2.5]{NoncompactQED}, and let $\mathscr{S}_0:=\C$.
There is exactly one continuous linear map $F^{(n)}\colon\mathscr{S}_n\to L^2(\mu)$ with $F^{(n)}(f_1\otimes\cdots\otimes f_n)=F(f_1)\cdots F(f_n)$ for test forms $f_1,\dots,f_n$
(Lemma~\ref{lem:tensor-tests}), and $F^{(0)}(f):=f$;
the Schwinger functions of $\mu$ are $S_n(f):=\int\mathrm{d}\mu\,F^{(n)}(f)$ for $f\in\mathscr{S}_n$, $n\ge0$, so that $S_0(f)=f$
(Definition~\ref{def:tensor-schwinger}).
A Euclidean map $\mathfrak{g}$ acts on $\mathscr{S}_n$ argument by argument, by $(\mathfrak{g}\cdot)^{\otimes n}$, and $\tau_a^{\otimes n}:=(\tau_a\cdot)^{\otimes n}$;
for a permutation $\mathfrak{s}$ of $\{1,\dots,n\}$, $f^{\mathfrak{s}}$ arises from $f\in\mathscr{S}_n$ by permuting the arguments together with their tensor indices;
and for a unit vector $n\in\R^3$, an integer $m\ge0$ and $f\in\mathscr{S}_m$, $\Theta^n_mf:=\overline{(r_{n,0}\cdot)^{\otimes m}f}$.
The axioms of Osterwalder and Schrader are numbered as in \cite{OsterwalderSchrader1975}.
For a vector field $b$ on $\R^3$, $\hat b(p):=\int\mathrm{d}x\,\exp(-ip\cdot x)\,b(x)$ is its Fourier transform and $\hat b(p)^{\dagger}$ the conjugate transpose of $\hat b(p)$;
$P^\perp(p):=1-p\,p^{T}/|p|^2$ is the multiplier of the transverse projection (Definition~\ref{def:free-maxwell}).

\begin{theorem}[Continuum limit points]\label{thm:continuum-limits}
Let $\rho\in(0,1)$, let $v_0$ be a point of the fundamental chamber, let the weights be standard,
let $L\ge3$ be odd, $0<\bar m\le1$ and $e>0$, and let $N\mapsto m_{0,N}$ be an arbitrary sequence of real numbers.
\begin{enumerate}
\item\label{it:continuum-soft}
$\sup_{N,T}\int\mathrm{d}P^{\mathrm{ch}}_{N,T}(F)\,\|F\|_{-3}^2\le61$;
the family $(P^{\mathrm{ch}}_{N,T})$ is tight on $H_{-4}$, and on $\mathcal{S}'$ for the strong and for the weak-$*$ topology;
every limit point has the eight properties of Part~I\,(b) and the property~(c) of the main theorem of the companion paper,
\cite[Theorem~3.1, Part~I\,(b), (c)]{NoncompactQED};
in particular it satisfies \textup{(G)} and \eqref{cond:closedness}, is invariant under the translations of $\R^3$ and under $\Ocub$,
and is reflection positive at every plane $\{x_j=c\}$, $j=1,2,3$, $c\in\R$;
every iterated limit, in either order, is a limit point.
\item\label{it:continuum-volume-first}
Volume-first limit points exist:
$\mathcal{V}_N$ is not empty for every $N$,
and every sequence $\mu^{(N_i)}\in\mathcal{V}_{N_i}$, $N_i\to\infty$, has a subsequence that converges weakly on $H_{-4}$.
Every volume-first limit point $\mu$ has the properties of \eqref{it:continuum-soft} and, in addition, the following ones.
\begin{enumerate}
\item\label{it:continuum-nine-planes}
$\mu\in\mathcal{C}_9$:
for every $\ell\in\mathcal{L}$ and every $c\in\R$, $\mu$ is invariant under the reflection $r_{\ell,c}$ in the plane $\{\ell=c\}$,
and $\langle\widehat{r_{\ell,c}}\,\Phi,\Phi\rangle_\mu\ge0$ for all $\Phi$ in $L^2(\Sigma_{\{\ell>c\}},\mu)$, all of $L^2$ of the open side.
\item\label{it:continuum-cone}
For all $j\neq i$:
the joint spectral measure of $(\mathsf{H}_j,\mathsf{P}_i)$ is carried by $\{(\lambda,k_i):\ \lambda\ge|k_i|\}$, that is, $\mathsf{H}_j\ge|\mathsf{P}_i|$ with the constant one;
$\left(1-\exp(-s\cdot\mathsf{H}_j)\right)^2\ge\exp(-s\cdot\mathsf{H}_j)\left(2-V_j(s\cdot e_i)-V_j(-s\cdot e_i)\right)$ for all $s>0$;
and the five-point inequality
\[
u(0)-4\,u(t\cdot e_j)+u(2t\cdot e_j)+u(t\cdot e_j+t\cdot e_i)+u(t\cdot e_j-t\cdot e_i)\ge0,\qquad u(v):=\mu\left((\Theta_j\Phi)\cdot U(v)\Phi\right),
\]
holds for all $t>0$ and all $\Phi\in L^2(\Sigma_j^+,\mu)$,
in particular for all polynomials and exponentials in the $F(f)$, $f\in L^2(\R^3;\Lambda^2)$ vanishing on $\{x_j<0\}$,
and in the form of Corollary~\ref{cor:unbounded-observables} for the Schwinger functions.
\item\label{it:continuum-vacuum}
$\ker\mathsf{H}_j\subset\ker\mathsf{P}_i$ for $j\neq i$;
and $P^j_\mu=P^{123}_\mu$ for $j=1,2,3$, which is the property \eqref{it:projections-coincide} of Definition~\ref{def:symmetry-conditions}.
\item\label{it:continuum-rate}
For every $r>0$, all $\Phi,\Phi'\in L^2(\Sigma_{(-r,r)^3},\mu)$ and all $a\in\R^3$ with $|a|_\infty\ge C_1\cdot r$,
\[
\left|\langle\Phi,U(a)\Phi'\rangle_\mu-\langle\Phi,P^{123}_\mu\Phi'\rangle_\mu\right|\le C_0\cdot\|\Phi\|_2\cdot\|\Phi'\|_2\cdot\frac{r}{|a|_\infty},
\]
where $C_0$ and $C_1$ are the absolute constants of Theorem~\ref{thm:clustering-rate}.
\item\label{it:continuum-components}
Let $\mu=\int\mathrm{d}\nu(\xi)\,\mu_\xi$ be the $\R^3$-ergodic decomposition of $\mu$.
For $\nu$-almost every $\xi$:
$\mu_\xi\in\mathcal{C}_9$;
$\mu_\xi$ is invariant under $\R^3\rtimes\Ocub$, so that $\mu$ has the property \eqref{it:components-cubic} of Definition~\ref{def:symmetry-conditions}, and ergodic under $\R^3$;
it satisfies \textup{(G)} with the same right side;
$\mu_\xi(dF=0)=1$;
$\ker\mathsf{H}_j(\mu_\xi)=\C\,\Omega_j$, a unique vacuum, and $\mathsf{H}_j(\mu_\xi)\ge|\mathsf{P}_i|$ for $j\neq i$;
for every $r>0$, all $\Phi,\Phi'\in L^2(\Sigma_{(-r,r)^3},\mu_\xi)$ and all $a\in\R^3$ with $|a|_\infty\ge C_1\cdot r$,
\[
\left|\mu_\xi\left(\overline{\Phi}\cdot U(a)\Phi'\right)-\overline{\mu_\xi(\Phi)}\cdot\mu_\xi(\Phi')\right|\le C_0\cdot\|\Phi\|_2\cdot\|\Phi'\|_2\cdot\frac{r}{|a|_\infty};
\]
it clusters as $|a|\to\infty$ in any manner on all of $L^2(\mu_\xi)$;
and it has the properties 1--6 of \cite[Corollary~8.9]{NoncompactQED}.
\end{enumerate}
\end{enumerate}
\end{theorem}

\begin{theorem}[Euclidean invariance]\label{thm:euclidean}
Let $\rho$, $v_0$, the weights, $L$, $\bar m$, $e$ and the mass assignment be as in Theorem~\ref{thm:continuum-limits}.
Let $\mu$ be a volume-first limit point of $(P^{\mathrm{ch}}_{N,T})$, and let $\mu=\int\mathrm{d}\nu(\xi)\,\mu_\xi$ be its $\R^3$-ergodic decomposition.

\emph{Part~I.}
\begin{enumerate}
\item\label{it:limit-euclidean}
$\mu$ is invariant under the full Euclidean group $\R^3\rtimes\mathrm{O}(3)$:
$\mathfrak{g}_*\mu=\mu$, and $S_n(\mathfrak{g}\cdot f_1,\dots,\mathfrak{g}\cdot f_n)=S_n(f_1,\dots,f_n)$
for all $n$, all Euclidean maps $\mathfrak{g}$ and all $f_1,\dots,f_n\in L^2(\R^3;\Lambda^2)$, coinciding supports included.
\item\label{it:limit-all-planes}
$\mu$ is reflection positive at every plane:
for every unit vector $n$ and every $c\in\R$, $(r_{n,c})_*\mu=\mu$ and $\langle\widehat{r_{n,c}}\,\Phi,\Phi\rangle_\mu\ge0$
for all $\Phi\in L^2(\Sigma_{\{n\cdot x>c\}},\mu)$.
\item\label{it:limit-spectrum}
In each picture $e_j$, $j=1,2,3$:
the joint spectral measure $\mathfrak{E}_j$ of $(\mathsf{H}_j,(\mathsf{P}_i)_{i\neq j})$ is carried by the round forward cone $\{(\lambda,k):\ \lambda\ge|k|_2\}$,
$|k|_2$ being the Euclidean norm of $k\in e_j^\perp$;
the orthogonal maps $g$ with $ge_j=e_j$ act on $\mathcal{H}_j$ by a strongly continuous unitary representation $\mathcal{U}_j$
with $\mathcal{U}_j(g)\Omega_j=\Omega_j$, which commutes with $\mathsf{H}_j$ and satisfies $\mathcal{U}_j(g)\,V_j(y)\,\mathcal{U}_j(g)^{-1}=V_j(gy)$, $y\in e_j^\perp$;
and, for every $i\neq j$, the closed linear span of the plate vectors of the plane $(j,i)$ carries a unitary group of boosts,
the group of Proposition~\ref{prop:rotations-boosts} for the plane $(1,2)$ and its transport by a permutation matrix for the other planes.
\item\label{it:limit-kallen-lehmann}
There is exactly one positive measure $\mathfrak{K}_\mu$ on $[0,\infty)$, of mass at most $1$, with
\[
S_2(f,f')=(2\pi)^{-3}\int\mathrm{d}p\,\hat b(p)^{\dagger}P^\perp(p)\,\hat b'(p)\cdot\int\mathrm{d}\mathfrak{K}_\mu(m)\,\frac{|p|^2}{|p|^2+m^2}
\]
for all real $f,f'\in L^2(\R^3;\Lambda^2)$, where $b=\ast f$ and $b'=\ast f'$.
\item\label{it:limit-components}
For $\nu$-almost every $\xi$ the component $\mu_\xi$ has the properties \eqref{it:limit-euclidean}--\eqref{it:limit-kallen-lehmann}, with a measure $\mathfrak{K}_\xi$,
in addition to those of Theorem~\ref{thm:continuum-limits}\eqref{it:continuum-components}:
it is ergodic under $\R^3$, satisfies \textup{(G)} with the same right side, $\mu_\xi(dF=0)=1$, $\ker\mathsf{H}_j(\mu_\xi)=\C\,\Omega_j$, and
\[
\left|\mu_\xi\left(\overline{\Phi}\cdot U(a)\Phi'\right)-\overline{\mu_\xi(\Phi)}\cdot\mu_\xi(\Phi')\right|\le C_0\cdot\|\Phi\|_2\cdot\|\Phi'\|_2\cdot\frac{r}{|a|_\infty}
\]
for $r>0$, $\Phi,\Phi'\in L^2(\Sigma_{(-r,r)^3},\mu_\xi)$ and $|a|_\infty\ge C_1\cdot r$.
\item\label{it:limit-general-laws}
The assertions \eqref{it:limit-euclidean}--\eqref{it:limit-components} hold for every law $\mu$ on $(\mathcal{S}',\Sigma)$
that satisfies \eqref{cond:soft-hypothesis} and \eqref{cond:closedness} and lies in $\mathcal{C}_9$,
with its own $\R^3$-ergodic decomposition $\mu=\int\mathrm{d}\nu(\xi)\,\mu_\xi$.
\end{enumerate}

\emph{Part~II.}
Let $\tilde\mu$ be a law on $(\mathcal{S}',\Sigma)$ that satisfies \eqref{cond:soft-hypothesis} and \eqref{cond:closedness} and lies in $\mathcal{C}_9$,
for instance $\mu$ or, for $\nu$-almost every $\xi$, $\mu_\xi$.
Its Schwinger functions $S_n$, which are defined on all of $\mathscr{S}_n$, coinciding arguments included, have the following properties.
\begin{enumerate}
\item[\textup{(E0$'$)}]\itemmark{E0$'$}\label{cond:os-growth}
Linear growth:
$|S_n(f)|\le2\cdot66^n\cdot(n!)^{\frac12}\cdot q_{6n}(f)$ for $n\ge1$ and $f\in\mathscr{S}_n$.
\item[\textup{(E1)}]\itemmark{E1}\label{cond:os-invariance}
Euclidean invariance:
$S_n\left((\mathfrak{g}\cdot)^{\otimes n}f\right)=S_n(f)$ for every Euclidean map $\mathfrak{g}$ of $\R^3$, reflections included, and every $f\in\mathscr{S}_n$.
\item[\textup{(E2)}]\itemmark{E2}\label{cond:os-positivity}
Positivity:
for every unit vector $n$, every integer $N\ge0$, every $f_0\in\C$ and all $f_k\in\mathscr{S}_k$, $1\le k\le N$,
with the support of $f_k$ in $\left\{y\in\R^{3k}:\ n\cdot y_i\ge0\text{ for }1\le i\le k\right\}$,
\[
\sum_{k,l=0}^{N}S_{k+l}\left(\Theta^n_kf_k\otimes f_l\right)\ge0.
\]
\item[\textup{(E3)}]\itemmark{E3}\label{cond:os-symmetry}
Symmetry:
$S_n\left(f^{\mathfrak{s}}\right)=S_n(f)$ for every permutation $\mathfrak{s}$ of $\{1,\dots,n\}$ and every $f\in\mathscr{S}_n$.
\item[\textup{(E4)}]\itemmark{E4}\label{cond:os-cluster}
Cluster property:
if $\tilde\mu$ is ergodic under the translations of $\R^3$, then for $f\in\mathscr{S}_m$ and $f'\in\mathscr{S}_n$
\[
S_{m+n}\left(f\otimes\tau_a^{\otimes n}f'\right)\to S_m(f)\cdot S_n(f')\qquad\text{as }|a|\to\infty\text{ in any manner};
\]
if moreover $m,n\ge1$, $r>0$, and $f$ and $f'$ have compact supports in the open sets $\left((-r,r)^3\right)^m$ and $\left((-r,r)^3\right)^n$, then
\[
\left|S_{m+n}\left(f\otimes\tau_a^{\otimes n}f'\right)-S_m(f)\cdot S_n(f')\right|\le C_0\cdot\left\|F^{(m)}(f)\right\|_2\cdot\left\|F^{(n)}(f')\right\|_2\cdot\frac{r}{|a|_\infty}
\qquad\left(|a|_\infty\ge C_1\cdot r\right).
\]
If $\tilde\mu$ is not ergodic, the limit of $S_{m+n}\left(f\otimes\tau_a^{\otimes n}f'\right)$ as $|a|_\infty\to\infty$ is
$\left\langle\overline{F^{(m)}(f)},P^{123}_{\tilde\mu}F^{(n)}(f')\right\rangle_{\tilde\mu}$.
\end{enumerate}
Moreover:
\begin{enumerate}
\item[\textup{(E5)}]\itemmark{E5}\label{it:schwinger-reality}
Reality:
$\overline{S_n(f)}=S_n(\overline f)$ for $f\in\mathscr{S}_n$.
\item[\textup{(E6)}]\itemmark{E6}\label{it:schwinger-bianchi}
Bianchi identity:
$S_n(d^{*}u\otimes f')=0$ for $n\ge1$, every Schwartz function $u$ on $\R^3$ and every $f'\in\mathscr{S}_{n-1}$,
where $d^{*}u$ is the two-form whose dual vector field is $-\nabla u$.
\item[\textup{(E7)}]\itemmark{E7}\label{it:schwinger-determinacy}
Determinacy:
$\tilde\mu$ is the only law on $(\mathcal{S}',\Sigma)$ with the Schwinger functions $S_n$.
\item[\textup{(E8)}]\itemmark{E8}\label{it:schwinger-spectrum}
Spectrum and vacuum:
$\ker\mathsf{H}_j=\C\,\Omega_j$, $j=1,2,3$, if $\tilde\mu$ is ergodic under the translations of $\R^3$.
The round forward cone of \eqref{it:limit-spectrum}, which holds for $\tilde\mu$ by \eqref{it:limit-general-laws}, and this vacuum
hold also in the reconstruction at every plane $\{n\cdot x=c\}$, which is unitarily equivalent to the picture $e_1$.
This reconstruction is defined as follows:
by \eqref{it:limit-all-planes} for $\tilde\mu$, $\widehat{r_{n,c}}$ is a unitary involution of $L^2(\tilde\mu)$,
and $(\Phi,\Phi')\mapsto\langle\widehat{r_{n,c}}\,\Phi,\Phi'\rangle_{\tilde\mu}$ is positive semidefinite on $L^2(\Sigma_{\{n\cdot x>c\}},\tilde\mu)$;
$\mathcal{H}_{n,c}$ is the completion of this space modulo the null vectors of the form, with the canonical map $[\,\cdot\,]_{n,c}$,
and for $n=e_j$, $c=0$ it is $\mathcal{H}_j$.
There are a self-adjoint operator $\mathsf{H}_n\ge0$ and a unitary representation $V_n$ of $n^\perp$ on $\mathcal{H}_{n,c}$ with
\[
\exp(-t\cdot\mathsf{H}_n)\,V_n(y')\,[\Phi]_{n,c}=\left[U(t\cdot n+y')\Phi\right]_{n,c}\qquad
\left(t\ge0,\ y'\in n^\perp,\ \Phi\in L^2(\Sigma_{\{n\cdot x>c\}},\tilde\mu)\right);
\]
their joint spectral measure is carried by the round forward cone $\{(\lambda,k'):\ \lambda\ge|k'|_2\}$, $k'\in n^\perp$,
and $\ker\mathsf{H}_n=\C\,[1]_{n,c}$ if $\tilde\mu$ is ergodic.
\end{enumerate}
So for $\mu$ itself \eqref{cond:os-growth}, \eqref{cond:os-invariance}, \eqref{cond:os-positivity} and \eqref{cond:os-symmetry} hold,
and for $\nu$-almost every $\xi$ the ergodic component $\mu_\xi$ satisfies all five axioms \eqref{cond:os-growth}--\eqref{cond:os-cluster}.
\end{theorem}

\subsection{What the theorems do not contain}\label{sec:intro-no-lower-bound}

Theorems~\ref{thm:continuum-limits} and~\ref{thm:euclidean} contain no lower bound.
The point mass $\mu_0$ at the zero field $F=0$ satisfies \eqref{cond:soft-hypothesis} and \eqref{cond:closedness},
is invariant under $F\mapsto-F$ and lies in $\mathcal{C}_9$;
it has every property that Theorem~\ref{thm:continuum-limits}\eqref{it:continuum-soft} and~\eqref{it:continuum-volume-first} asserts of a volume-first limit point $\mu$ itself,
and every property that Theorem~\ref{thm:euclidean} asserts of $\mu$ and of $\mu_\xi$, with the K\"all\'en--Lehmann measure $0$
(Examples~\ref{ex:zero-field} and~\ref{ex:zero-field-axioms}).
Gaussian domination, the cone, the rate of clustering and the inequality $\mathfrak{K}_\mu([0,\infty))\le1$ are bounds from above,
and neither $\mu\neq\mu_0$ nor $\mathfrak{K}_\mu\neq0$ follows from them.

\subsection{The argument}\label{sec:intro-argument}

\emph{The chamber operator.}
Since no fermion site lies on a mirror, a bond that crosses a mirror $\Pi$ joins a site $x$ to $\theta_\Pi x$ and is orthogonal to $\Pi$ (Proposition~\ref{prop:geometry}).
For every mirror, with the upward unit normal $n$ and the reflection $\theta$, the hopping matrices obey $\gamma_n\,w(x,x')\,\gamma_n=w(\theta x',\theta x)$ (Lemma~\ref{lem:hopping});
consequently $\Theta_\Pi$ fixes the fermion action $S_F$,
and $W$ acts on the even integrands so that $Z_{g\mathbb{T}}[G^g]$ equals $Z_{\mathbb{T}}[G]$ or its complex conjugate (Proposition~\ref{prop:symmetries}).
On every torus $\det D^{\mathbb{T}}_A\ge0$ (Proposition~\ref{prop:accretivity}).
At a mirror $\Pi$, the part $S^{\times}_\Pi$ of the action that couples the two sides has the form
\[
-S^{\times}_\Pi=\sum_i\omega_i\cdot u_i\cdot\Theta_\Pi(\zeta_i)\,\zeta_i,
\]
where the $\zeta_i$ are the modes of the bonds across $\Pi$, the $u_i$ their phases,
and $\omega_i=\frac{w_j}{2}(1\pm\rho)$, $j$ the type of the bond, are the crossing coefficients;
all of them are non-negative for $-1\le\rho\le1$ (Proposition~\ref{prop:crossing}).
This form of the crossing part expresses the reflection positivity of the fermion weight at every mirror of the nine families.
The sites of the cubic Wilson fermion, by contrast, lie on the diagonal mirrors.

The partition function is positive on every torus and for every real mass, $Z_{\mathbb{T}}[1]>0$ (Theorem~\ref{thm:partition-positivity}).
Let $\rho^2\neq1$, and expand $\det D^{\mathbb{T}}_A$ in the link phases $u_b=\exp(iqA(b))$:
the monomial $\prod_{b\parallel e_\mu}u_b^{32}$ has for its coefficient a product of determinants of hopping matrices,
one factor for each fermion site of the torus, and this coefficient equals $C_*^{|\mathbb{T}|}$ with
\[
C_*:=\left(\frac{w_0}{2}\right)^{32}\cdot\left(\frac{w_1}{2}\right)^{64}\cdot\left(\frac{w_2}{2}\right)^{64}\cdot\left(\frac{w_3}{2}\right)^{32}\cdot\left(1-\rho^2\right)^{96}>0
\]
(Proposition~\ref{prop:extreme-coefficient});
here $|\mathbb{T}|$ is the number of photon sites of $\mathbb{T}$.
For $-1\le\rho\le1$ and an antiperiodic coordinate torus $\mathbb{T}$, the free determinant $\det D^{\mathbb{T}}_0$ is at least a power of the determinant of the hexagon operator,
which is known in closed form (Lemma~\ref{lem:hexagon-determinant}, Proposition~\ref{prop:free-determinants});
changing the character amounts to translating $A$ by a constant one-cochain (Lemma~\ref{lem:characters-coverings}, Theorem~\ref{thm:closed-interval}).

\emph{Adapted tori and slab operators.}
Let $-1\le\rho\le1$, and let $\mathbb{T}$ be adapted to $\ell$.
Its weight is a product of three kinds of factors: the weight of the upper half, the reflection of that weight,
and, at each of the two mirrors that bound the half, crossing factors that are of positive type;
at the second mirror the right sign is produced by $\chi(v)=-1$, $v$ the closing vector (Lemma~\ref{lem:second-mirror}).
For the photon, the crossing weight is of positive type, and Gauss' law holds at the links that the mirrors cut (Lemmas~\ref{lem:positive-type} and~\ref{lem:gauss-law});
together these facts give Theorem~\ref{thm:adapted-positivity}.
For integrands of the class $\mathcal{A}_\Xi[c,c']$ the same factorisation produces Hilbert--Schmidt slab operators $\mathfrak{S}[G]$ (Theorem~\ref{thm:slab-operators}).
They compose, they are covariant under translations, the reflection of $G$ replaces $\mathfrak{S}[G]$ by its adjoint up to a translation,
and every torus whose cross-section is the given one has a trace formula that contains the unitary twist $V_g^{-1}U_{\mathrm{f}}^{\frac{1+\chi(g)}{2}}$,
$g$ a closing vector and $U_{\mathrm{f}}$ the fermion parity.
At the coordinate link planes, the word formula writes the functional with the slice operator $\mathfrak{M}$
and the crossing operator $K$, which is a Gaussian in the electric field operators multiplied by the projection onto the gauge invariant vectors
(Proposition~\ref{prop:word-formula}).

\emph{Tori.}
Let $-1\le\rho\le1$.
In Schatten classes, H\"older's inequality is insensitive to the unitary twist:
for the contents of $2k$ consecutive slabs, the modulus of the functional of their product is bounded by the geometric mean of the functionals of their folds
on the straightened torus $\mathbb{T}^0$ of Definition~\ref{def:slabs} (Proposition~\ref{prop:twisted-holder}).
A chain of three untwistings joins the cubic torus to the adapted torus, and the adapted torus to an elongated coordinate torus:
\[
C(n)\to D'(n)\to D(n)\to C\left(\frac n2,2n,n\right),\qquad D(n)=D\left(\frac n2,n,n\right),\quad D'(n)=D'\left(\frac n2,n,n\right),
\]
for even $n\ge8$ (Definition~\ref{def:sides}, Proposition~\ref{prop:untwistings}), and chessboard estimates hold on these tori and between them (Theorem~\ref{thm:chessboard}).
On coordinate tori the pressure $f(n_1,n_2,n_3)=(n_1n_2n_3)^{-1}\cdot\log Z_{C(n_1,n_2,n_3)}[1]$ is monotone in each side, $\log Z_{C(n_1,n_2,n_3)}[1]$ is convex in each side,
and $f$ has a limit when the sides tend to infinity in an arbitrary way (Proposition~\ref{prop:coordinate-pressure}).
Hence the partition functions increase along the chain, and for every sequence of even sides $n_i\to\infty$
\[
\Delta_i:=f\left(\frac{n_i}{2},2n_i,n_i\right)-f(n_i,n_i,n_i)\ge0,\qquad\Delta_i\to0
\]
(Definition~\ref{def:sides}, Theorem~\ref{thm:four-tori}).
The four tori have the same volume $n_i^3$ and, in the limit, the same pressure:
the logarithm of the ratio of the partition functions of any two of them lies between $-n_i^3\cdot\Delta_i$ and $n_i^3\cdot\Delta_i$,
and the argument requires no bound on these ratios that is uniform in $i$.
On a coordinate torus $\mathbb{T}$ with the period lattice $\Gamma$, the expectation of $\exp\langle\beta,F\rangle$ is at most $\exp\left(\frac12\|\beta\|^2\right)$
for every real two-cochain $\beta$ of $\Z^3/\Gamma$ (Theorem~\ref{thm:coordinate-domination}).

\emph{Sequences of sides, data and comparison states.}
A trigonometric polynomial of the curvature is a finite sum $O=\sum_ka_k\cdot\varphi_{\beta_k}$, with complex coefficients,
of characters $\varphi_\beta(F)=\exp(i\langle\beta,F\rangle)$ indexed by finitely supported real two-cochains $\beta$;
one puts $|O|_1:=\sum_k|a_k|$, the support of $O$ is the union of the plaquettes at which some $\beta_k$ is not zero,
$r_O$ is the least integer $r\ge1$ such that the support lies in a cube $y+[-r,r]^3$ for some $y\in\R^3$,
and the real trigonometric polynomials form the set $\mathcal{T}$ (Definition~\ref{def:characters}).
For a sequence $\mathfrak{n}=(n_i)$ of sides, an integer belongs to $\mathcal{D}(\mathfrak{n})$ if it divides $n_i$ for all but finitely many $i$,
and \eqref{cond:divisibility} requires that $8\in\mathcal{D}(\mathfrak{n})$ and that $\mathcal{D}(\mathfrak{n})$ be infinite (Definition~\ref{def:sides}).
Let $\mathfrak{n}$ satisfy \eqref{cond:divisibility}.
Data $\mathfrak{d}=(c_0,c_+,c_3;h;O)$ consist of a cell structure $\mathcal{Q}_{\mathfrak{d}}$ of the diagonal frame $(\ell_-,\ell_+,x_3)$,
with the walls $(\ell_-,c_0+jh)$, $(\ell_+,c_++jh)$, $(x_3,c_3+jh)$, $j\in\Z$, where $c_0,c_+\in\Z$, $c_3\in\Z+\frac12$, and $h$ is even with $4h\in\mathcal{D}(\mathfrak{n})$,
and of an element $O\in\mathcal{T}$ whose support lies in the open cell $Q_0$ of $\mathcal{Q}_{\mathfrak{d}}$, with $h>8r_O+8$;
the volume of an open cell is $|Q_0|=h^3/2$ (Definitions~\ref{def:cell-structures} and~\ref{def:data}).
The folded array $O_{\mathfrak{d}}$ places a mirror copy of $O$ in every cell, and the perturbed functional is
\[
\mu^t_i(G)=\frac{Z_{D(n_i)}\left[\exp(t\cdot O_{\mathfrak{d}})\cdot G\right]}{Z_{D(n_i)}\left[\exp(t\cdot O_{\mathfrak{d}})\right]}\qquad(t\ge0).
\]
A comparison state is a law $\nu_0$ on $\Omega$ of the form $\nu_0=\lim_k\lim_m\mu^{t_k}_{i_m}$ with $t_k\to0+$, weak limits taken volume first and strength second;
it is subordinate to a subsequence of the indices if $(i_m)$ is a subsequence of it (Definition~\ref{def:comparison-states}).
Condition \eqref{cond:divisibility} is imposed on one sequence only, the sequence used for the adapted tori, for the data and for the comparison states:
cell structures and coarser structures of unbounded width have to fit into these tori.
The cubic tori of Theorem~\ref{thm:cubic-limits} are not subject to it.
To an arbitrary sequence $\mathfrak{n}$ of integers $n_i\ge4$ tending to infinity belongs a companion $\mathfrak{n}'=(n'_i)$:
a sequence with \eqref{cond:divisibility} such that $n'_i\le n_i$ for all but finitely many $i$ and $n'_i/n_i\to1$
(Definition~\ref{def:companions}, Lemma~\ref{lem:companions-exist}).
The proof of Theorem~\ref{thm:cubic-limits} builds everything that needs \eqref{cond:divisibility} on a companion,
while the cubic tori $C(n_i)$ and their laws $\mu_i$ keep the sides of $\mathfrak{n}$.

\emph{Three points of departure.}
The proof of Theorem~\ref{thm:cubic-limits} follows the corresponding argument for Yang--Mills theory in four dimensions
and departs from it at three points, each of which is a theorem in its own right;
the absence of layered states is their conclusion.
That argument uses pointwise positivity of integrands, which is not available here.
Its place is taken by H\"older's inequality in Schatten classes (Proposition~\ref{prop:twisted-holder});
by Jensen's inequality, which is applied to bosonic observables only, on the cubic torus and on the adapted torus (Proposition~\ref{prop:fold-finite});
by an amplification without a logarithm (Theorem~\ref{thm:amplification});
and by an elimination of layers without a division (Theorem~\ref{thm:two-layer-reduction}).

\emph{The fold inequality.}
Let $-1\le\rho\le1$, let $\mathfrak{n}$ satisfy \eqref{cond:divisibility}, and let $\mathfrak{d}$ be data.
For all $i\ge i_0(\mathfrak{d})$ (Definition~\ref{def:threshold}) and all $t>0$,
\[
\mu_i(O)\le\mu^t_i(O)+\frac{|Q_0|\cdot\Delta_i}{t}
\]
(Proposition~\ref{prop:fold-finite}), and $\Delta_i\ge0$, $\Delta_i\to0$ (Theorem~\ref{thm:four-tori});
because of the error term, the volume limit has to be taken before the strength tends to zero.
For each $t\ge0$ the laws $(\mu^t_i)_i$ form a tight family, and each subsequence of the indices carries comparison states subordinate to it;
each comparison state satisfies \textup{(G)}, is invariant under the group generated by the reflections in the walls of $\mathcal{Q}_{\mathfrak{d}}$,
and is reflection positive at every wall, the diagonal mirror $(\ell_-,c_0)$ among them;
finally, if $\mu_i\to\nu$ weakly along a subsequence, then $\nu(O)\le\nu_0(O)$ for every comparison state $\nu_0$ subordinate to it (Theorem~\ref{thm:comparison-states}).
Tightness and \textup{(G)} are consequences of the amplification (Proposition~\ref{prop:domination-tightness}).
The inequality involves only observables read in a single open cell, hence on one side of the diagonal mirror,
and on its own it gives no information at the mirror.
It rests on Jensen's inequality for the probability measures $\mu_i$, on the cubic torus, and $\mu^t_i$, on the adapted torus, applied to bosonic observables,
and on H\"older's inequality, with the unitary twist, along the four tori.
Now let $\mathfrak{n}$ be an arbitrary sequence of integers $n_i\ge4$ with $n_i\to\infty$, let $\mathfrak{n}'$ be a companion, and let the data be formed with $\mathfrak{n}'$.
Then the inequality compares the cubic tori of the given sides with the adapted tori of the companion:
for all large $i$, namely $i\ge i_0(\mathfrak{d})$ with $n'_i\le n_i$, and all $t>0$,
\[
\mu_{C(n_i)}(O)\le\mu^t_i(O)+\frac{|Q_0|\cdot\Delta'_i}{t},
\]
where $\Delta'_i$, defined for all large $i$, tends to $0$;
moreover $\nu(O)\le\nu_0(O)$ whenever $\nu$ is the weak limit of the laws of $C(n_i)$ along a subsequence and $\nu_0$ is a comparison state subordinate to it
(Proposition~\ref{prop:fold-any-side}, Theorem~\ref{thm:comparison-any-side}).
The proof uses H\"older's inequality for slabs of unequal heights (Proposition~\ref{prop:unequal-slabs}),
leaves one wider slab empty in each direction (Lemma~\ref{lem:one-step}),
and estimates the exceptional tori that occur by $Z_{\mathbb{T}}[1]\le\exp\left(\bar f\cdot|\mathbb{T}|\right)$,
a bound valid on every torus with a single constant $\bar f$ (Lemma~\ref{lem:uniform-ceiling}).
It is this form of the inequality that the proof of Theorem~\ref{thm:cubic-limits} uses.

\emph{Amplification and the two-layer principle.}
Let $\mathfrak{n}$ satisfy \eqref{cond:divisibility}, and let $\mathfrak{d}$ be data;
$k(\mathfrak{d},r)$ and $i_0(\mathfrak{d},r)$ denote the integers of Lemma~\ref{lem:coarse-structures}.
For $-1\le\rho\le1$:
if an integrand $G$ is bounded by $C$ with a box of radius $r$ (Definition~\ref{def:bounded-elements}), then
\[
\left|\mu^t_i(G)\right|\le C\cdot\gamma_{t,i}(k),\qquad
\gamma_{t,i}(k)=\exp\left(2t\cdot|O|_1\cdot k^3+|Q_0|\cdot k^3\cdot\Delta_i\right),\qquad k=k(\mathfrak{d},r),
\]
for all $t\ge0$ and all $i\ge i_0(\mathfrak{d},r)$ (Theorem~\ref{thm:amplification});
for fixed $k$, $\lim_{t\to0+}\limsup_i\gamma_{t,i}(k)=1$ (Definition~\ref{def:amplification-factor}).
Let now $0<\rho<1$.
For every coordinate link plane $(x_\mu,c)$, every monomial $Y\neq1$ of $\mathcal{P}_1(c)$, with $Y^{g_\mu}\in\mathcal{P}_1(\mu,c)$ (Definition~\ref{def:other-directions}),
every $\varepsilon>0$, every $t\ge0$ and every $i\ge i_0(\mathfrak{d},2r_0)$,
with $k=k(\mathfrak{d},2r_0)$, $r_0$ the radius of the box $B(c,S,\Lambda_Y)$, and $S$, $\Lambda_Y$ as in Definition~\ref{def:two-layer},
\[
\left|\mu^t_i\left(\Theta_{\mu,c}Y^{g_\mu}\right)-\overline{\mu^t_i(Y^{g_\mu})}\right|
\le2\cdot2^{\frac{\#|S|}{2}}\cdot\xi_Y(\varepsilon)\cdot\gamma_{t,i}(k)+4\cdot\left(\gamma_{t,i}(k)^2-1\right)^{\frac12}\cdot\|X_{Y,\varepsilon}\|_\#
\]
(Theorem~\ref{thm:two-layer-principle});
$k$ and $i_0(\mathfrak{d},2r_0)$ do not depend on $\varepsilon$, $\xi_Y(\varepsilon)\to0$ as $\varepsilon\to0$, and $\|X_{Y,\varepsilon}\|_\#$ is finite for every $\varepsilon>0$.
Hence
\[
\lim_{t\to0+}\,\limsup_{i\to\infty}\left|\mu^t_i\left(\Theta_{\mu,c}Y'\right)-\overline{\mu^t_i(Y')}\right|=0
\]
for all $Y'\in\mathcal{P}_1(\mu,c)$, elements with fields included;
the limits are taken in the order $i$, then $t$ at fixed $\varepsilon$, then $\varepsilon$.
The mechanism is an insertion:
an operator $X$ of the mid-plane algebra is placed between the two factors of $K=K^{\frac12}K^{\frac12}$, $K$ the crossing operator,
and since $K^{\frac12}$ factorises over links and sites, the result is a local integrand (Proposition~\ref{prop:insertions}).
By H\"older's inequality with insertions its size is controlled by $\|X\|_\#$ (Propositions~\ref{prop:holder-insertions} and~\ref{prop:bounded-elements}),
and chessboard estimates in two frames transfer this control to a single copy under $\mu^t_i$, at the cost of the factor $\gamma_{t,i}(k)$
(Lemma~\ref{lem:coarse-structures}, Theorem~\ref{thm:amplification}).
On a C$^{*}$-algebra, a unital functional whose norm does not exceed $(1+\delta)^{\frac12}$ is Hermitian up to an error $4\cdot\delta^{\frac12}$ (Theorem~\ref{thm:two-layer-principle}),
and each two-layer monomial differs from such an insertion by an element with a small bound (Proposition~\ref{prop:two-layer-insertions}).

\emph{Exact elimination of the outermost layers.}
Let $0<\rho<1$, let $\mathfrak{n}$ satisfy \eqref{cond:divisibility}, and let $c=j_0+\frac12$ with $j_0\in\Z$;
$\Theta_c:=\Theta_{x_1,c}$ is the reflection at the mirror $(x_1,c)$ (Definition~\ref{def:two-layer}).
For every $P\in\mathcal{P}$ (Definition~\ref{def:polynomials}), in particular for every character of the curvature, there is $Y_P\in\mathcal{P}_1(c)$,
defined on $\Z^3$ and independent of the data, of $i$, of $t$ and of the sequence of sides,
and for all data $\mathfrak{d}$ there are $C_{P,\mathfrak{d}}<\infty$ and $i_0$ with
\[
\left|\mu^t_i(P)-\mu^t_i(Y_P)\right|\le C_{P,\mathfrak{d}}\cdot t,\qquad
\left|\mu^t_i(\Theta_cP)-\mu^t_i(\Theta_cY_P)\right|\le C_{P,\mathfrak{d}}\cdot t\qquad(0\le t\le1,\ i\ge i_0)
\]
(Theorem~\ref{thm:two-layer-reduction});
the same holds at the link planes of the directions $e_2$ and $e_3$, with an element $Y^\mu_P$ of $\mathcal{P}_1(\mu,c)$, the reflection $\Theta_{\mu,c}$,
and constants that depend on $P$, $\mu$, $c$ and $\mathfrak{d}$.
Since the hopping matrices are invertible for $\rho^2\neq1$, Schwinger--Dyson identities trade the fields of the lowest thin layer for fields of higher layers
(Lemma~\ref{lem:layer-bonds}, Proposition~\ref{prop:schwinger-dyson});
an imaginary translation of the integration contour of the link variables in the next plane then removes the exponent on the lowest plane,
and under this translation the Maxwell weight yields a character and the fermion weight, its variables being nilpotent, a polynomial
(Propositions~\ref{prop:contour-shift} and~\ref{prop:outer-plane}).
These identities are valid on every torus in which the box in question embeds and for all entire bosonic factors, and therefore so are their mirror images;
the perturbation $\exp(t\cdot O_{\mathfrak{d}})$ contributes only errors of order $t$, which the a priori bound of Proposition~\ref{prop:a-priori} controls.
Through this bound the amplification, and with it the condition $0<\rho<1$, enters the elimination.

\emph{No layered states.}
Let $0<\rho<1$, and let $\mathfrak{n}$ satisfy \eqref{cond:divisibility}.
Whatever the data, a comparison state is invariant under the reflections $\theta_{x_\mu,c}^{*}$ in the link planes, $\mu=1,2,3$, $c\in\Z+\frac12$,
and consequently under the translations $\tau_v$, $v\in2\Z^3$ (Theorem~\ref{thm:no-layers});
the absence of layered states means precisely this invariance under the link reflections and under $2\Z^3$.
For its proof one applies the elimination and the two-layer principle to the characters $P=\varphi_\beta$, letting first $i\to\infty$, then $t\to0+$, then $\varepsilon\to0$;
this gives $\nu_0(\Theta_c\varphi_\beta)=\overline{\nu_0(\varphi_\beta)}$, and a law is determined by its values on the characters.
Of this invariance, the proof of Theorem~\ref{thm:cubic-limits} uses only the even translations along $e_1$.

\emph{The diagonal mirrors.}
Let $\nu$ be a weak limit point of $(\mu_i)$ and $c_0\in\Z$.
The laws on $\Omega$ that satisfy \textup{(G)} and are invariant and reflection positive at $(\ell_-,c_0)$ form a convex and weakly compact set $\mathcal{R}(c_0)$
(Lemma~\ref{lem:diagonal-laws}).
Both $\nu$ and the comparison states are invariant under the even translations along $e_1$, and such a translation carries an arbitrary $O\in\mathcal{T}$
into a single open cell on one side of $(\ell_-,c_0)$;
this produces $\nu_0\in\mathcal{R}(c_0)$ with $\nu(O)\le\nu_0(O)$ (Lemma~\ref{lem:diagonal-domination}).
By Lemma~\ref{lem:separation}, a law that no element of $\mathcal{T}$ separates from $\mathcal{R}(c_0)$ lies in $\mathcal{R}(c_0)$;
so $\nu\in\mathcal{R}(c_0)$, and the action of $\Ocub$ gives the other diagonal families.
The remaining assertions of Theorem~\ref{thm:cubic-limits} are obtained on the cubic tori themselves, for every integer side $n_i\ge4$.
Each law $\mu_i$ is invariant under $W$ and under $F\mapsto-F$.
Under $\mu_i$, Gaussian domination holds for the cochains $\beta$ whose support lies in a box that embeds in the torus (Theorem~\ref{thm:coordinate-domination}),
and reflection positivity at a mirror $(x_\mu,c)$ holds for the bounded continuous functions of finitely many $F_\sigma$ that are read in a slab $\{c<x_\mu<c+h\}$ with $h<n_i/2$
(Corollary~\ref{cor:trace-forms}).
The cochains supported on one plaquette fit into each of these tori, which yields tightness (Lemma~\ref{lem:domination-compactness});
a fixed cochain or test function fits as soon as $i$ is large, and the two inequalities, like the invariances, are preserved under weak limits.

\emph{The continuum limit.}
At each cutoff the law of the curvature on the torus has the symmetries of the lattice, obeys the Gaussian bound with the constant one for every real two-cochain of the torus,
and is reflection positive at the coordinate planes through the origin (Proposition~\ref{prop:cutoff-law});
hence the family $(P^{\mathrm{ch}}_{N,T})$ and its limit points have the properties of Part~I of the main theorem of the companion paper, \cite[Theorem~3.1]{NoncompactQED}
(Proposition~\ref{prop:chamber-limits}).
At a fixed spacing, a volume limit is the image of a weak limit on $\Omega$ of laws of cubic tori under an explicit Borel map,
and Theorem~\ref{thm:cubic-limits} applies to that weak limit (Proposition~\ref{prop:volume-limits}).
Since reflection positivity at a plane survives weak limits, the volume-first limit points belong to $\mathcal{C}_9$ (Proposition~\ref{prop:nine-families}).
For a law of $\mathcal{C}_9$, the reconstructions at $\{x_j=0\}$ and at $\{x_j+x_i=0\}$, compared through a two-constants estimate on the closed unit disc
(Lemma~\ref{lem:two-constants}), yield the cone with the constant one (Theorem~\ref{thm:cone});
it follows that zero energy forces zero momentum and $P^j_\mu=P^{123}_\mu$ (Proposition~\ref{prop:zero-momentum}),
and the correlations cluster at the rate of Theorem~\ref{thm:clustering-rate}.
Let $\mu$ be a law of $\mathcal{C}_9$ that satisfies \eqref{cond:soft-hypothesis}.
For almost every $\xi$, its ergodic component $\mu_\xi$ lies in $\mathcal{C}_9$, is ergodic under $\R^3$ and invariant under $\R^3\rtimes\Ocub$,
satisfies \textup{(G)} with the same right side, has a unique vacuum and the cone $\mathsf{H}_j\ge|\mathsf{P}_i|$, and clusters;
and $\mu_\xi$ is carried by the closed distributions if $\mu$ is (Proposition~\ref{prop:ergodic-components}).
Theorem~\ref{thm:clustering-rate}, applied to these components, gives them the rate, and \cite[Corollary~8.9]{NoncompactQED}, applied to $\mu$, the properties 1--6 listed there.
A volume-first limit point satisfies \eqref{cond:soft-hypothesis} by Proposition~\ref{prop:chamber-limits}
and \eqref{cond:closedness} by Theorem~\ref{thm:continuum-limits}\eqref{it:continuum-soft};
so all of this applies to it, and gives Theorem~\ref{thm:continuum-limits}\eqref{it:continuum-components}.
What the sections on the fold inequality, the amplification, the elimination and the layered states prove
on perturbed functionals and comparison states enters Theorem~\ref{thm:continuum-limits}\eqref{it:continuum-volume-first}
only through Theorem~\ref{thm:cubic-limits}.
Besides, its proof takes two auxiliary lemmas from the section on layered states, neither of which depends on the rest of these sections:
the stability of the Gaussian bound under weak limits of laws on $\Omega$ (Lemma~\ref{lem:domination-compactness}),
which enters the proof of Proposition~\ref{prop:volume-limits},
and the two-constants estimate for finite measures on the closed unit disc (Lemma~\ref{lem:two-constants}), on which the cone rests.
Part~\eqref{it:continuum-soft} rests on Proposition~\ref{prop:cutoff-law}, that is, on the symmetries of the functional, the sign of the determinant,
the convergence of the functional, the trace forms, the chessboard estimates, Gaussian domination on coordinate tori and the positivity of the partition function
(Propositions~\ref{prop:symmetries}, \ref{prop:accretivity}, \ref{prop:berezin}, Corollary~\ref{cor:trace-forms},
Theorems~\ref{thm:chessboard}, \ref{thm:coordinate-domination} and~\ref{thm:partition-positivity}).
Of the statements of the companion paper on single laws,
Theorem~\ref{thm:continuum-limits}\eqref{it:continuum-soft} uses \cite[Proposition~7.19]{NoncompactQED} and \cite[Lemma~21.36]{NoncompactQED},
and Theorem~\ref{thm:continuum-limits}\eqref{it:continuum-components} uses \cite[Theorem~8.8\,(v)]{NoncompactQED} besides the corollary cited above.

\emph{Rotations.}
Behind Theorem~\ref{thm:euclidean} stands a theorem on single laws (Theorem~\ref{thm:orthogonal-group}):
if a law on $\mathcal{S}'$ satisfies \eqref{cond:soft-hypothesis} and \eqref{cond:closedness},
and the cone $\mathsf{H}_j\ge|\mathsf{P}_i|$ with the constant one holds for a single ordered pair $j\neq i$,
then all maps of $\R^3\rtimes\mathrm{O}(3)$ leave the law invariant, and it is reflection positive at all planes.
Its proof treats one plane at a time:
\textup{(G)}, \eqref{cond:closedness}, invariance under the translations, under $\theta_j$ and under the quarter turn of the plane $(j,i)$,
together with $(RP_j)$ and the cone, yield invariance under every rotation of that plane (Theorem~\ref{thm:one-plane}).
The steps are these.
Gaussian domination and multiple reflections, without the cone, control a single field smeared with an $L^2$ function supported in a time slab,
inserted among any number of other fields, by a bound that does not depend on that number (Proposition~\ref{prop:insertion-bound}).
If the field is smeared in the single transverse direction only, the cone turns it into an operator valued function, holomorphic on the forward tube,
whose norm is at most a constant divided by the distance to the boundary (Propositions~\ref{prop:point-plates} and~\ref{prop:kernels}).
For two such plates, the two-point function extends to an entire function of the rotation angle $t$, of size at most a constant times $\exp\left(2\,|\Im t|\right)$;
so its expansion in $t$ contains no frequency of modulus above $2$,
and the Bianchi identity $dF=0$ removes the single term that is not invariant under the rotations (Theorem~\ref{thm:two-point}).
For $N$ points the anomaly index $\mathfrak{a}(N)$ does not exceed $N/2$; reflection positivity makes it convex, and regular polygons make it divisible;
these three properties give $\mathfrak{a}(2a)=a\cdot\mathfrak{a}(2)$ and $\mathfrak{a}(N)\le\frac N2\cdot\mathfrak{a}(2)$ (Proposition~\ref{prop:anomaly-vanishes}).
The value $\mathfrak{a}(2)=0$ is supplied by the theorem on the two-point function (Theorem~\ref{thm:two-point});
so closedness enters the argument for $N$ points a second time, through $\mathfrak{a}(2)=0$,
and without closedness the conclusion is false (Proposition~\ref{prop:sharpness}).
Finally, a bounded multilinear form on $L^2$ vanishes identically as soon as it vanishes on arguments with pairwise disjoint projections to the plane
and on arguments one of which is a gradient (Proposition~\ref{prop:no-contact-terms});
this carries the invariance from separated plates over to the law.
For laws with \eqref{cond:soft-hypothesis} and \eqref{cond:closedness}, the following properties are equivalent:
the cone, the five-point inequalities, invariance under $\mathrm{O}(3)$, reflection positivity at one diagonal plane, membership in $\mathcal{C}_9$,
and reflection positivity at every plane (Theorem~\ref{thm:cone-equivalences}).
Thus a single diagonal mirror of the limit law already yields reflection positivity at every plane:
the lattice has only nine families of mirrors, and the continuum supplies all other planes through Gaussian domination, closedness and the cone.
From the lattice, Theorem~\ref{thm:euclidean} uses only Theorem~\ref{thm:continuum-limits}\eqref{it:continuum-soft} and~\eqref{it:continuum-nine-planes}.
In Part~II, linear growth, positivity at the coordinate planes, symmetry, the Bianchi identity and determinacy
are taken from \cite[Theorems~7.15 and~7.17]{NoncompactQED},
which apply to every law with \eqref{cond:soft-hypothesis} and \eqref{cond:closedness} by \cite[Proposition~7.19]{NoncompactQED};
the cluster property rests on Proposition~\ref{prop:ergodic-components}, Theorem~\ref{thm:clustering-rate} and the tensor tests of Lemma~\ref{lem:tensor-tests}.

\subsection{The order of the limits}\label{sec:intro-limits}

No two limits are interchanged.
On the lattice the order is: volume ($i\to\infty$), then strength ($t\to0+$), then the regularisation ($\varepsilon\to0$).
The error term $|Q_0|\cdot\Delta_i/t$ of the fold inequality, and $|Q_0|\cdot\Delta'_i/t$ for arbitrary sides, forces the volume before the strength;
the norm $\|X_{Y,\varepsilon}\|_\#$, which is in general unbounded as $\varepsilon\to0$, forces the regularisation last.
In the continuum the order is: volume ($T\to\infty$) at a fixed spacing, then spacing ($N\to\infty$);
this is why Theorem~\ref{thm:continuum-limits}\eqref{it:continuum-volume-first} and Theorem~\ref{thm:euclidean} speak of volume-first limit points.
Limit points along joint sequences of cutoffs that approach the volume limits are volume-first limit points (Theorem~\ref{thm:joint-limits}).

\subsection{The Wilson parameter}\label{sec:intro-wilson}

The condition $-1\le\rho\le1$, that is, $\omega_i=\frac{w_j}{2}(1\pm\rho)\ge0$ (Proposition~\ref{prop:crossing}), is what every positivity needs:
Theorem~\ref{thm:adapted-positivity}, the slab operators, the comparison of tori, the fold inequality,
H\"older's inequality with insertions, the amplification and the domination of the perturbed laws
(Proposition~\ref{prop:holder-insertions}, Theorem~\ref{thm:amplification}, Proposition~\ref{prop:domination-tightness}),
and the fold inequality and the comparison states for arbitrary sides
(Propositions~\ref{prop:unequal-slabs} and~\ref{prop:fold-any-side}, Lemma~\ref{lem:one-step}, Theorem~\ref{thm:comparison-any-side});
the bound on all partition functions holds for every real $\rho$ (Lemma~\ref{lem:uniform-ceiling}).
At $\rho=1$ the coefficients $\frac{w_j}{2}(1-\rho)$ vanish, at $\rho=-1$ the coefficients $\frac{w_j}{2}(1+\rho)$;
so all $\omega_i$ are positive, and the diagonal matrix of the crossing coefficients is invertible, exactly when $-1<\rho<1$.
This enters with the two-layer monomials that contain fields (Definition~\ref{def:two-layer}, Proposition~\ref{prop:two-layer-insertions}, Theorem~\ref{thm:two-layer-principle}),
and, through Proposition~\ref{prop:two-layer-insertions}, the a priori bound of Proposition~\ref{prop:a-priori}.
The Schwinger--Dyson identities hold for every real $\rho$;
the invertibility of the hopping matrices, $\rho^2\neq1$, is what the replacement rules built on them need
(Proposition~\ref{prop:schwinger-dyson}\eqref{it:equation-of-motion} and~\eqref{it:replacement-rules}).
The positivity of the partition function needs no condition on $\rho$ (Theorems~\ref{thm:floor} and~\ref{thm:closed-interval}).
In Lemma~\ref{lem:diagonal-domination} and in Theorem~\ref{thm:cubic-limits} the condition $0<\rho<1$ enters through Theorem~\ref{thm:no-layers} only.

\subsection{Further results}\label{sec:intro-further}

\emph{The averaged law.}
For the harmonic-averaged law, the law of the curvature on a coordinate torus conditioned on the flux zero (Definition~\ref{def:averaged-law}),
the following holds for $0<\rho<1$ along every sequence of integers $n_i\ge4$ with $n_i\to\infty$:
the laws on the cubic tori $C(n_i)$ form a tight sequence, and each weak limit point satisfies \textup{(G)} and closedness,
is invariant under the group $W$ and the map $F\mapsto-F$, and is reflection positive at all mirrors of the six diagonal families,
and at all mirrors of the nine families when every $n_i$ is even (Theorem~\ref{thm:averaged-any-side}).
Theorem~\ref{thm:continuum-limits}\eqref{it:continuum-soft} and~\eqref{it:continuum-volume-first} hold with the averaged family in the place of $(P^{\mathrm{ch}}_{N,T})$,
whose sides are even (Corollary~\ref{cor:averaged-admissible}, Theorem~\ref{thm:averaged-limit-points});
its volume-first limit points satisfy \eqref{cond:soft-hypothesis} and \eqref{cond:closedness} and lie in $\mathcal{C}_9$ (Proposition~\ref{prop:soft-limits}),
so that Theorem~\ref{thm:euclidean}\eqref{it:limit-general-laws} and Part~II of Theorem~\ref{thm:euclidean} apply to them (Remark~\ref{rem:euclidean-soft-limits}).

\emph{Volume limits on the lattice.}
Let $\nu$ be a law on $\Omega$, invariant under $W$ and reflection positive at the mirrors of all nine families;
every $\nu$ of Theorem~\ref{thm:cubic-limits} is such a law.
In the reconstruction of $\nu$ in the direction $e_\mu$, the transfer operator $0\le T_\mu\le1$ of the translation by $e_\mu$
and the unitary $U_{\mu,\mu'}$ of the translation by $e_{\mu'}$, $\mu'\neq\mu$, commute and satisfy
$(1-T_\mu)^2\ge T_\mu\left(2-U_{\mu,\mu'}-U_{\mu,\mu'}^{*}\right)$;
at the points of their joint spectrum with $T_\mu=\exp(-\lambda)>0$ and $U_{\mu,\mu'}=\exp(ik)$ one has $\cosh\lambda+\cos k\ge2$;
functions in $L^2(\nu)$ that are invariant under $\tau_{e_\mu}$ are invariant under $\Z^3$;
and almost every ergodic component of $\nu$ is invariant under $W$ (Theorem~\ref{thm:lattice-cone}).
For $0<\rho<1$, the family $(\mu_{C(n)})_{n\ge4}$ of all cubic tori is tight,
and every weak limit along a sequence $n_i\to\infty$ has the properties of Theorem~\ref{thm:cubic-limits}
(Corollary~\ref{cor:all-cubic-tori}).

\emph{The functional and the partition function.}
The functional $Z_{\mathbb{T}}[G]$ does not depend on the character $\chi$ of the torus, for every $G\in\mathcal{A}(\R^3)$ that embeds in its period lattice;
in particular $Z_{\mathbb{T}}[1]$ and $\mu_{\mathbb{T}}$ do not depend on $\chi$ (Proposition~\ref{prop:symmetries}\eqref{it:sign-character}).
For $\rho^2\neq1$ there is the explicit floor
\[
|\mathbb{T}|^{-1}\cdot\log Z_{\mathbb{T}}[1]\ge\log C_*+3\log\frac{2\pi}{q}-6\left(\frac{2\pi}{q}\right)^2,
\]
with the number $C_*$ above, on every torus (Theorem~\ref{thm:floor});
for every real $\rho$ and $M$, the law $\mu_{\mathbb{T}}$ and the law of $\widetilde F$ under the normalised Maxwell weight, all flux sectors summed,
are mutually absolutely continuous (Theorems~\ref{thm:floor} and~\ref{thm:closed-interval});
and for $-1\le\rho\le1$ and every real $M$ the free determinant is positive on every antiperiodic coordinate torus (Proposition~\ref{prop:free-determinants}).

\emph{Joint limits.}
At fixed $N$, the Prokhorov distance $d_{\mathrm{P}}$ of $P^{\mathrm{ch}}_{N,T}$ from $\mathcal{V}_N$ tends to $0$ as $T\to\infty$;
the limit points along joint sequences $(N_i,T_i)$ with $d_{\mathrm{P}}(P^{\mathrm{ch}}_{N_i,T_i},\mathcal{V}_{N_i})\to0$ are volume-first limit points,
and for every sequence $N_i\to\infty$ there are such $T_i$ (Theorem~\ref{thm:joint-limits}).

\emph{Arbitrary sides.}
Let $-1\le\rho\le1$.
H\"older's inequality with a unitary twist holds for slabs of unequal heights (Proposition~\ref{prop:unequal-slabs}),
and the fold inequality leads from the cubic torus of an arbitrary integer side $n\ge n'$ to the adapted torus of the side $n'$ (Proposition~\ref{prop:fold-any-side}).
A second form of it holds for $n'\le n\le2n'$ (Propositions~\ref{prop:comparable-even} and~\ref{prop:comparable-odd}).
In this form no layer of cells is left out, and the error term $|Q_0|\cdot(n/n')^3\cdot\left(f(n'/2,n'/2,n'/2)-f_\infty\right)/t$,
$f_\infty$ the infimum of the pressure, does not contain the constant $\bar f$;
it tends to $0$ whenever $n'_i\le n_i\le2n'_i$.

\emph{Single laws on $\mathcal{S}'$.}
The cone with the constant one and the five-point inequality hold for every law invariant under the translations of $\R^3$
that satisfies $(R_j)$ and $(RP_j)$ and has, at the plane $\{x_j+x_i=0\}$, both invariance and reflection positivity (Theorem~\ref{thm:cone}),
and reflection positivity at one diagonal plane gives the cone also by a direct argument in the continuum (Proposition~\ref{prop:diagonal-mirror}).
Zero energy implies zero momentum, and the clustering rate holds, for every law of $\mathcal{C}_9$
(Proposition~\ref{prop:zero-momentum}, Theorem~\ref{thm:clustering-rate});
the statement on the ergodic components holds for every law of $\mathcal{C}_9$ with \eqref{cond:soft-hypothesis} (Proposition~\ref{prop:ergodic-components}).
The insertion bound holds without the cone, for every law with \textup{(G)}, $(T_1)$, $(R_1)$ and $(RP_1)$ (Proposition~\ref{prop:insertion-bound}),
and with the cone the plate bound has the exponent one (Proposition~\ref{prop:point-plates}).
A two-point function that is invariant under $\mathrm{O}(3)$ has the K\"all\'en--Lehmann form (Theorem~\ref{thm:two-point}).
For a law invariant under the rotations about a coordinate axis these rotations act by unitary operators on the reconstruction space;
with the cone the spectrum lies in the round forward cone;
and boosts act unitarily on the closed linear span of the plate vectors (Proposition~\ref{prop:rotations-boosts}).
Explicit laws show that none of the following can be omitted from Theorem~\ref{thm:one-plane}:
Gaussian domination, closedness, the quarter turn, the cone, and the constant one in the cone (Proposition~\ref{prop:sharpness}).

\subsection{Plan of the paper}\label{sec:intro-plan}

Section~\ref{sec:setting} fixes the notation and the definitions.
Section~\ref{sec:chamber} studies the chamber operator and proves Theorem~\ref{thm:adapted-positivity},
Section~\ref{sec:positivity} proves that the partition function is positive,
Section~\ref{sec:slabs} constructs the slab operators,
and Section~\ref{sec:tori} compares the partition functions of different tori.
Sections~\ref{sec:fold}, \ref{sec:amplification} and~\ref{sec:elimination} prove the fold inequality,
the amplification with the two-layer principle and the exact elimination of the outermost layers;
Section~\ref{sec:layered} proves the absence of layered states and its consequences on the lattice,
and Section~\ref{sec:sides} treats arbitrary sides and proves Theorem~\ref{thm:cubic-limits}.
Section~\ref{sec:continuum} proves Theorem~\ref{thm:continuum-limits},
Section~\ref{sec:averaged} treats the averaged law and joint limits,
and Section~\ref{sec:rotations} proves Theorem~\ref{thm:euclidean}.

\section{Setting and notation}\label{sec:setting}

\subsection{Mirrors, chambers and the photon field}\label{sec:mirrors-chambers}

\begin{definition}[Lattice and cochains]\label{def:cochains}
Points of $\R^3$ are written $x=(x_1,x_2,x_3)$, the unit vectors are $e_1,e_2,e_3$, and $|x|_\infty:=\max_\mu|x_\mu|$.
The \emph{photon lattice} is $\Z^3$, and its points are the photon sites.
Its \emph{cells} are closed subsets of $\R^3$, each of them carrying an orientation:
the sites $x\in\Z^3$, which are positively oriented;
the \emph{links} $(x,\mu):=[x,x+e_\mu]$, oriented from $x$ to $x+e_\mu$;
the \emph{plaquettes} $(x;\mu\nu)$, $\mu<\nu$, that is, the squares with the corners $x$, $x+e_\mu$, $x+e_\mu+e_\nu$, $x+e_\nu$, oriented by $e_\mu\wedge e_\nu$;
and the \emph{cubes} $x+[0,1]^3$, oriented by $e_1\wedge e_2\wedge e_3$.
A link is also denoted by $b$; plaquettes, and cells in general, are denoted by $\sigma$.

For $k=0,1,2,3$ let $C^k$ be the space of the real \emph{$k$-cochains}, the real functions on the $k$-cells;
$C^k_{\mathrm{c}}\subset C^k$ is the subspace of the finitely supported cochains, and $C^k(\Z)\subset C^k$ the group of the integer-valued ones.
The pairing
\[
\langle u,v\rangle:=\sum_{\sigma}u(\sigma)\cdot v(\sigma),
\]
the sum over all $k$-cells $\sigma$, is defined whenever the sum converges absolutely, in particular if $u$ or $v$ is finitely supported;
moreover $\|u\|^2:=\langle u,u\rangle\in[0,\infty]$ for every $u\in C^k$.
The values of $\lambda\in C^0$, $A\in C^1$ and $F\in C^2$ at $x$, at $(x,\mu)$ and at $(x;\mu\nu)$ are written $\lambda(x)$, $A_\mu(x)$ and $F_{\mu\nu}(x)$;
the value of $F$ at a plaquette $\sigma$ is also written $F_\sigma$, and the value of a three-cochain $u$ at the cube $x+[0,1]^3$ is written $u(x)$.

The \emph{coboundary} $d\colon C^k\to C^{k+1}$, $k=0,1,2$, is the one of \cite[Sections~2.1 and~10.1.1]{NoncompactQED}:
\[
\begin{aligned}
(d\lambda)_\mu(x)&=\lambda(x+e_\mu)-\lambda(x),\\
(dA)_{\mu\nu}(x)&=A_\mu(x)+A_\nu(x+e_\mu)-A_\mu(x+e_\nu)-A_\nu(x)\qquad(\mu<\nu),\\
(dF)(x)&=[F_{23}(x+e_1)-F_{23}(x)]-[F_{13}(x+e_2)-F_{13}(x)]+[F_{12}(x+e_3)-F_{12}(x)];
\end{aligned}
\]
there the same formulas are written on the tori $\Lambda_n$, and the value of a three-cochain at a cube carries the index $123$.
A cochain $u$ with $du=0$ is called \emph{closed}.
The adjoint of $d$ on finitely supported cochains is the map $d^{*}\colon C^{k+1}_{\mathrm{c}}\to C^k_{\mathrm{c}}$ with
$\langle d^{*}u,v\rangle=\langle u,dv\rangle$ for $u\in C^{k+1}_{\mathrm{c}}$ and $v\in C^k_{\mathrm{c}}$.

A one-cochain $A$ is also read as a function on the ordered pairs of nearest neighbours of $\Z^3$ that changes its sign when the two points are exchanged,
extended by zero to the pairs of equal points:
\[
A(x,x+e_\mu):=A_\mu(x),\qquad A(x+e_\mu,x):=-A_\mu(x),\qquad A(x,x):=0 .
\]
The same convention applies to complex one-cochains, the elements of $C^1+iC^1$.

Let $g$ be an isometry of $\R^3$ with $g\Z^3=\Z^3$.
It maps every $k$-cell onto a $k$-cell, and its \emph{pull-back} $g^{*}\colon C^k\to C^k$ is
\[
(g^{*}u)(\sigma):=\pm u(g\sigma),
\]
with the sign $+$ if $g$ carries the orientation of $\sigma$ to the orientation of the cell $g\sigma$, and $-$ otherwise.
For a one-cochain this means $(g^{*}A)(x,x')=A(gx,gx')$ for all nearest neighbours $x,x'\in\Z^3$, and also for $x=x'$.
The map $g^{*}$ is orthogonal for the pairing and commutes with $d$, and $(gg')^{*}=g'^{*}g^{*}$ for two such isometries $g$ and $g'$.
Indeed, $g$ permutes the cells of each dimension, the orientation signs of a composition multiply,
and the formulas for $d$ express $(du)(\sigma)$ as the sum of the values of $u$ at the faces of $\sigma$ of one dimension less,
each counted with the sign $+1$ or $-1$ according as its orientation agrees or not with the orientation induced on it from $\sigma$ (outward normal first);
an isometry respects induced orientations.
For $v\in\Z^3$, $\tau_v$ is the translation $x\mapsto x+v$;
it preserves orientations, so that $(\tau_v^{*}u)(\sigma)=u(\sigma+v)$.
\end{definition}

\begin{definition}[Mirror families]\label{def:mirrors}
The set $\mathcal{L}$ of the \emph{families} consists of the nine linear forms
\[
x_\mu\quad(\mu=1,2,3),\qquad x_\mu-x_\nu,\quad x_\mu+x_\nu\quad(1\le\mu<\nu\le3)
\]
on $\R^3$: the three \emph{axis families} and the six \emph{diagonal families}.
For $\ell\in\mathcal{L}$, $n_\ell\in\Z^3$ denotes the vector with $\ell(x)=n_\ell\cdot x$ for all $x$;
thus $|n_\ell|^2=1$ for an axis family and $|n_\ell|^2=2$ for a diagonal family.
The vector $\hat n_\ell:=n_\ell/|n_\ell|$ is the \emph{upward unit normal} of $\ell$, and the number $\ell(x)$ is the \emph{height} of the point $x$.
The \emph{admissible positions} are
\[
\Adm(\ell):=\tfrac12\Z\quad\text{for an axis family},\qquad\Adm(\ell):=\Z\quad\text{for a diagonal family}.
\]
A \emph{mirror} is a pair $\Pi=(\ell,c)$ with $\ell\in\mathcal{L}$ and $c\in\Adm(\ell)$;
the same letter denotes the plane $\{\ell=c\}=\{x\in\R^3:\ \ell(x)=c\}$.
A mirror $(x_\mu,c)$ of an axis family is a \emph{site plane} if $c\in\Z$, and then it contains photon sites;
it is a \emph{link plane} if $c\in\Z+\tfrac12$, and then it cuts the links $(x,\mu)$ with $x_\mu=c-\tfrac12$ at their midpoints.
The \emph{reflection} in the mirror $\Pi=(\ell,c)$ is the map
\[
\theta_\Pi(x)=\theta_{\ell,c}(x):=x-\frac{2}{|n_\ell|^2}\cdot(\ell(x)-c)\cdot n_\ell ,
\]
also written $\theta_c$ when the family is fixed.
As $\ell(n_\ell)=|n_\ell|^2$, the formula gives $\ell(\theta_c x)=2c-\ell(x)$ for every $x\in\R^3$.
The following identities hold for $c,c'\in\Adm(\ell)$ and $v\in\Z^3$.
\begin{enumerate}
\item\label{it:mirror-involution}
$\theta_c^2=\mathrm{id}$.
\item\label{it:mirror-height}
$\ell\circ\theta_c=2c-\ell$.
\item\label{it:mirror-product}
$\theta_{c'}\theta_c$ is the translation by the vector $\frac{2}{|n_\ell|^2}\cdot(c'-c)\cdot n_\ell$, which lies in $\Z^3$ and has the height $2(c'-c)$.
\item\label{it:mirror-conjugation}
$\tau_v\theta_c\tau_{-v}=\theta_{c+\ell(v)}$.
\end{enumerate}
Item~\eqref{it:mirror-height} is the relation just obtained.
Items~\eqref{it:mirror-involution} and~\eqref{it:mirror-product} follow
when the formula for $\theta_c$, respectively for $\theta_{c'}$, is evaluated at the point $\theta_cx$ and this relation is inserted;
the height of the vector in~\eqref{it:mirror-product} is $\frac{2}{|n_\ell|^2}\cdot(c'-c)\cdot\ell(n_\ell)=2(c'-c)$.
Item~\eqref{it:mirror-conjugation} comes from the formula for $\theta_c$ at the point $x-v$ together with the linearity of $\ell$, $\ell(x-v)=\ell(x)-\ell(v)$:
\[
\tau_v\theta_c\tau_{-v}(x)=\theta_c(x-v)+v=x-\frac{2}{|n_\ell|^2}\cdot(\ell(x)-\ell(v)-c)\cdot n_\ell=\theta_{c+\ell(v)}(x).
\]
In~\eqref{it:mirror-product} the vector is integral because $2(c'-c)\in\Z$ and $n_\ell=e_\mu$ for an axis family $x_\mu$,
while $c'-c\in\Z$ and $|n_\ell|^2=2$ for a diagonal family;
in~\eqref{it:mirror-conjugation} the position $c+\ell(v)$ is admissible because $\ell(v)\in\Z$.

Let $W$ be the group generated by all reflections $\theta_\Pi$, $\Pi$ a mirror.
For the family $x_\mu$ and the positions $c=0$ and $c'=\tfrac12$, item~\eqref{it:mirror-product} shows that $\theta_{c'}\theta_c=\tau_{e_\mu}$;
hence $W$ contains every translation $\tau_v$, $v\in\Z^3$.
The cubic group $\Ocub:=\mathrm{O}(3)\cap\mathrm{GL}(3,\Z)$ of \cite[Section~2.5]{NoncompactQED} consists of the $48$ signed permutation matrices,
since the columns of an orthogonal matrix with integer entries are vectors $\pm e_\mu$;
its diagonal matrices form the subgroup $\{\pm1\}^3$.
The elements of $W$ are maps $x\mapsto k+\vec g\,x$ with $k\in\Z^3$ and $\vec g\in\Ocub$, and every such map belongs to $W$
(Lemma~\ref{lem:chambers}\eqref{it:chambers-group}).
Thus $W=\Z^3\rtimes\Ocub$; it is the affine Weyl group of type $\widetilde C_3$.
For $g\in W$ the matrix $\vec g\in\Ocub$ is the \emph{linear part} of $g$,
and $g$ is called \emph{even} if $\det\vec g=1$ and \emph{odd} if $\det\vec g=-1$.
The maps $x\mapsto k+\vec g\,x$ with diagonal $\vec g$ form the subgroup
\[
W_{\mathrm{c}}:=\Z^3\rtimes\{\pm1\}^3 ,
\]
which is normal in $W$, because $\Ocub$ normalises $\{\pm1\}^3$.
It is the subgroup generated by the reflections in the axis families:
these reflections $x\mapsto x-2(x_\mu-c)\cdot e_\mu$ lie in $W_{\mathrm{c}}$,
and the subgroup generated by them contains the changes of sign $\theta_{x_\mu,0}$ of one coordinate
and the translations $\tau_{e_\mu}$, hence all of $W_{\mathrm{c}}$.
Finally, $e_\pm:=e_1\pm e_2$ and $\ell_\pm:=x_1\pm x_2$;
thus $\ell_-$ and $\ell_+$ are diagonal families, with $n_{\ell_\pm}=e_\pm$.
\end{definition}

\begin{definition}[Chambers]\label{def:chambers}
The mirror planes form a locally finite family, so that the complement of their union is open in $\R^3$;
its connected components are the \emph{chambers}.
The chambers are open simplices, and $W$ acts simply transitively on the set of the chambers
(Lemma~\ref{lem:chambers}\eqref{it:chambers-transitive}; see also \cite[Ch.~V, \S3]{Bourbaki1968}).
The \emph{fundamental chamber} is
\[
C_0:=\left\{x\in\R^3:\ 0<x_3<x_2<x_1<\frac{1}{2}\right\}.
\]
Its \emph{walls} are the mirrors
\[
\Pi_0=\{x_3=0\},\qquad\Pi_1=\{x_2=x_3\},\qquad\Pi_2=\{x_1=x_2\},\qquad\Pi_3=\left\{x_1=\frac{1}{2}\right\},
\]
that is, $\Pi_0=(x_3,0)$, $\Pi_1=(x_2-x_3,0)$, $\Pi_2=(x_1-x_2,0)$ and $\Pi_3=\left(x_1,\tfrac12\right)$;
for $i=0,1,2,3$, $s_i\in W$ denotes the reflection in $\Pi_i$.
The inward unit normals of the walls are
\[
\hat n_0=e_3,\qquad\hat n_1=\frac{e_2-e_3}{\sqrt{2}},\qquad\hat n_2=\frac{e_1-e_2}{\sqrt{2}},\qquad\hat n_3=-e_1 .
\]
The closure $\overline{C_0}$ is the tetrahedron with the vertices $0$, $\tfrac12e_1$, $\tfrac12(e_1+e_2)$ and $\tfrac12(e_1+e_2+e_3)$;
the areas of its faces $\Pi_i\cap\overline{C_0}$ and its volume are
\[
|\Pi_0\cap\overline{C_0}|=|\Pi_3\cap\overline{C_0}|=\frac{1}{8},\qquad
|\Pi_1\cap\overline{C_0}|=|\Pi_2\cap\overline{C_0}|=\frac{\sqrt{2}}{8},\qquad
|C_0|=\frac{1}{48}.
\]
Moreover, for every $x\in C_0$,
\[
\sum_{i=0}^{3}|\Pi_i\cap\overline{C_0}|\cdot\hat n_i=0,\qquad
\frac{1}{3}\sum_{i=0}^{3}|\Pi_i\cap\overline{C_0}|\cdot\operatorname{dist}(x,\Pi_i)=|C_0| .
\]
The first identity follows from the values above;
the second expresses the decomposition of $C_0$ into the four pyramids with the apex $x$ over its faces.
\end{definition}

\begin{definition}[Fermion sites and bonds]\label{def:fermion-sites}
A base point $v_0=(v_1,v_2,v_3)\in C_0$ is fixed; thus $0<v_3<v_2<v_1<\tfrac12$.
The \emph{fermion sites} are the points of the orbit
\[
\mathbb{X}:=Wv_0=\{gv_0:\ g\in W\}.
\]
The map $g\mapsto gv_0$ is a bijection of $W$ onto $\mathbb{X}$ (Lemma~\ref{lem:chambers}\eqref{it:chambers-transitive});
together with Lemma~\ref{lem:chambers}\eqref{it:chambers-group} it follows that $\mathbb{X}$ is the disjoint union
of the $48$ translates $\Z^3+\vec g\,v_0$, $\vec g\in\Ocub$, of $\Z^3$,
so that every cube $k+[0,1)^3$, $k\in\Z^3$, contains exactly $48$ fermion sites.

Let $x\in\mathbb{X}$, and let $g\in W$ be the element with $x=gv_0$.
For $i\in\{0,1,2,3\}$ put
\[
x^{(i)}:=gs_iv_0 .
\]
Since $gs_ig^{-1}$ is the reflection in the wall $g\Pi_i$ of the chamber $gC_0$, which contains $x$,
the site $x^{(i)}$ is the mirror image of $x$ in this wall;
it is called the \emph{neighbour of $x$ across the wall of type $i$}.
As $g\mapsto gv_0$ is injective and $s_0,\dots,s_3$ are four different reflections,
the five sites $x,x^{(0)},\dots,x^{(3)}$ are pairwise different;
moreover $\left(x^{(i)}\right)^{(i)}=x$, because $x^{(i)}=(gs_i)v_0$ and $s_i^2=\mathrm{id}$.
The \emph{bonds} are the unordered pairs $\{x,x^{(i)}\}$ with $x\in\mathbb{X}$ and $i\in\{0,1,2,3\}$.
By the two facts just stated, the pair determines the index $i$, which is called the \emph{type} of the bond.
The unit vector from $x$ to $x^{(i)}$ is
\[
\hat d_i(x):=\frac{x^{(i)}-x}{\left|x^{(i)}-x\right|}.
\]
Every $x\in\mathbb{X}$ has a unique representation $x=k+\vec g\,v_0$ with $k\in\Z^3$ and $\vec g\in\Ocub$
(Lemma~\ref{lem:chambers}\eqref{it:chambers-group} and~\eqref{it:chambers-transitive}).
The map
\[
\pi\colon\mathbb{X}\to\Z^3,\qquad\pi(k+\vec g\,v_0):=k ,
\]
assigns to $x$ its \emph{nearest photon site}:
all coordinates of $\vec g\,v_0$ lie in $\left(-\tfrac12,\tfrac12\right)$, so that $\pi(x)$ is the point of $\Z^3$ nearest to $x$.
Finally, for $x=k+\vec g\,v_0$,
\[
y(x):=x-\pi(x)=\vec g\,v_0 ,
\]
and the moduli of the coordinates of $y(x)$ are $v_1,v_2,v_3$ in some order.
\end{definition}

\begin{definition}[Closed two-cochains and characters]\label{def:characters}
The closed two-cochains form the space
\[
\Omega:=\left\{F\in C^2:\ dF=0\right\}
\]
with the product topology, that is, the topology of pointwise convergence on the plaquettes;
as a closed subspace of a countable product of copies of $\R$, it is a Polish space.
\emph{Local} functions on $\Omega$ are those that depend only on the values $F_\sigma$ at the plaquettes $\sigma$ of some finite set.
For $\beta\in C^2_{\mathrm{c}}$ the function
\[
\varphi_\beta(F):=\exp(i\langle\beta,F\rangle),\qquad F\in\Omega,
\]
is a continuous character of the additive group $\Omega$.
A \emph{trigonometric polynomial} is a finite sum
\[
O=\sum_k a_k\cdot\varphi_{\beta_k}
\]
with complex coefficients $a_k$ and pairwise different $\beta_k\in C^2_{\mathrm{c}}$.
For it one puts $|O|_1:=\sum_k|a_k|$;
$\supp O$ is the union of the plaquettes $\sigma$ with $\beta_k(\sigma)\neq0$ for some $k$;
and $r_O$ is the least integer $r\ge1$ such that $\supp O\subset y+[-r,r]^3$ for some $y\in\R^3$.
The real-valued trigonometric polynomials form the set $\mathcal{T}$.

Every $g\in W$ is an isometry of $\R^3$ with $g\Z^3=\Z^3$ (Definition~\ref{def:mirrors}),
and it maps $\Omega$ onto itself through $g^{*}$, which commutes with $d$ and has the inverse $(g^{-1})^{*}$ (Definition~\ref{def:cochains}).
The group $W$ acts on trigonometric polynomials by
\[
(gO)(F):=O(g^{*}F),\qquad F\in\Omega ;
\]
as $g^{*}$ is orthogonal, $\langle\beta,g^{*}F\rangle=\langle(g^{-1})^{*}\beta,F\rangle$, so that $g\varphi_\beta=\varphi_{(g^{-1})^{*}\beta}$,
and $gO$ is again a trigonometric polynomial, real-valued if $O$ is.
For a translation this reads $(\tau_vO)(F)=O(\tau_v^{*}F)$ with $(\tau_v^{*}F)_\sigma=F_{\sigma+v}$, $v\in\Z^3$;
so $\tau_vO$ depends only on the values of $F$ at the plaquettes of $\supp O+v$.

A Borel probability measure on $\Omega$ is called a \emph{law} on $\Omega$.
A law $\mu$ on $\Omega$ satisfies \emph{Gaussian domination with the constant one}, condition~\textup{(G)}, if for every $\beta\in C^2_{\mathrm{c}}$
\begin{equation}\tag{G}\label{cond:gaussian-domination}
\int\mathrm{d}\mu(F)\,\exp(\langle\beta,F\rangle)\le\exp\left(\frac{1}{2}\,\|\beta\|^2\right).
\end{equation}
\end{definition}

\begin{definition}[Configurations on a torus]\label{def:configurations}
The \emph{lattice charge} is a number $q>0$, and $2\pi/q$ is the period of the link variables.
A \emph{period lattice} is a subgroup $\Gamma$ of $\Z^3$ of rank three whose \emph{systole} satisfies
\[
\sys(\Gamma):=\min\{|g|_\infty:\ g\in\Gamma,\ g\neq0\}\ge4 ;
\]
its index in $\Z^3$ is $[\Z^3:\Gamma]$.

Fix a period lattice $\Gamma$.
Translations by elements of $\Gamma$ map cells onto cells and preserve orientations;
the $k$-cells of the torus $\Z^3/\Gamma$ are the $\Gamma$-orbits of $k$-cells of $\Z^3$.
Let $C^k(\Gamma)$ be the space of the real functions on the $k$-cells of $\Z^3/\Gamma$, and $C^k(\Gamma;\Z)$ the group of the integer-valued ones.
Composition with the quotient map identifies an element of $C^k(\Gamma)$ with a $\Gamma$-periodic element of $C^k$, its \emph{periodic extension}.
Through this identification $d$ maps $C^k(\Gamma)$ into $C^{k+1}(\Gamma)$ and $C^k(\Gamma;\Z)$ into $C^{k+1}(\Gamma;\Z)$,
as it commutes with the translations by elements of $\Gamma$,
and the elements of $C^1(\Gamma)$ are read as functions of ordered pairs of nearest neighbours as in Definition~\ref{def:cochains}.
On $C^k(\Gamma)$ the pairing and the norm are
\[
\langle u,v\rangle:=\sum_{\sigma}u(\sigma)\cdot v(\sigma),\qquad\|u\|^2:=\langle u,u\rangle,
\]
the sum ranging over the $k$-cells $\sigma$ of $\Z^3/\Gamma$.

The \emph{configuration space} of $\Gamma$ is the quotient
\[
\Conf(\Gamma):=\{(A,z):\ A\in C^1(\Gamma),\ z\in C^2(\Gamma;\Z),\ dz=0\}\big/{\sim}
\]
by the equivalence relation
\[
(A,z)\sim\left(A+\frac{2\pi}{q}\cdot k,\ z-dk\right),\qquad k\in C^1(\Gamma;\Z).
\]
For a function $\Phi$ on $\Conf(\Gamma)$, $\Phi(A,z)$ denotes its value at the class of $(A,z)$.
The values $A(b)$ at the $3\cdot[\Z^3:\Gamma]$ links $b$ of $\Z^3/\Gamma$ are coordinates which identify $C^1(\Gamma)$ with $\R^{3\cdot[\Z^3:\Gamma]}$,
and $\mathrm{d}A$ denotes the Lebesgue measure in these coordinates.
The closed elements $z$ of $C^2(\Gamma;\Z)$ form a countable set,
and a set of pairs $(A,z)$ is called Borel if, for each closed $z$, the $A\in C^1(\Gamma)$ for which $(A,z)$ belongs to the set form a Borel subset of $C^1(\Gamma)$.
A subset of $\Conf(\Gamma)$, or a function on $\Conf(\Gamma)$, is \emph{Borel} if its pull-back under the quotient map to the pairs $(A,z)$ is Borel.

The pairs $(A,z)$ with $A$ in the box $[0,2\pi/q)^{3\cdot[\Z^3:\Gamma]}$ form a fundamental domain of $\sim$, the fundamental domain given by the box:
every class contains exactly one of them, $k$ being determined by the condition that $A+\frac{2\pi}{q}\cdot k$ lie in the box.
The measure $\mathrm{d}(A,z)$ on the Borel sets of $\Conf(\Gamma)$ is defined as the image under the quotient map
of the product of the Lebesgue measure on the box with the counting measure on the closed $z$, a measure on this fundamental domain.
Thus
\[
\int_{\Conf(\Gamma)}\mathrm{d}(A,z)\,\Phi(A,z)=\sum_{z}\ \int_{[0,2\pi/q)^{3\cdot[\Z^3:\Gamma]}}\mathrm{d}A\,\Phi(A,z)
\]
for every Borel function $\Phi\ge0$ on $\Conf(\Gamma)$, the sum ranging over the closed elements $z$ of $C^2(\Gamma;\Z)$.
A complex Borel function $\Phi$ on $\Conf(\Gamma)$ is \emph{integrable} if this integral is finite for $|\Phi|$;
then the sum and the integrals on the right converge absolutely, and their value is the integral of $\Phi$ with respect to $\mathrm{d}(A,z)$.
The integral of $\Phi$ with respect to $\mathrm{d}(A,z)$ is also written $\int_{\Conf(\Gamma)}\Phi$.

On the set of all pairs, the product of $\mathrm{d}A$ with the counting measure on the closed $z$ is invariant
under the maps $(A,z)\mapsto\left(A+\frac{2\pi}{q}\cdot k,\ z-dk\right)$, $k\in C^1(\Gamma;\Z)$,
and the quotient map takes the same value at a pair and at its images.
Consequently every Borel set of pairs that meets every class in exactly one pair
can take the place of the fundamental domain given by the box in the definition of $\mathrm{d}(A,z)$:
such a set is the disjoint union of its intersections with the images of that fundamental domain under these maps,
and the inverse maps carry these intersections onto pairwise disjoint subsets of that fundamental domain whose union is the whole of it.

The \emph{curvature}
\[
F:=dA+\frac{2\pi}{q}\cdot z ,
\]
a closed element of $C^2(\Gamma)$, and the \emph{link phases} $\exp(iqA(b))$, $b$ a link, are functions on $\Conf(\Gamma)$:
they do not change when $(A,z)$ is replaced by an equivalent pair.
So are the functions $G(\hat A)$ of the local lifts $\hat A$, for the integrands $G$ of Definition~\ref{def:integrands}.
The periodic extension of $F$ is denoted by $\widetilde F\in\Omega$.

Let $g$ be an isometry of $\R^3$ with $g\Z^3=\Z^3$.
Its linear part $\vec g$ is orthogonal and maps $\Z^3$ onto $\Z^3$, so that $\vec g\in\Ocub$,
and $g\Gamma:=\vec g\,\Gamma$ is again a period lattice, as $\vec g$ preserves $|\cdot|_\infty$.
The pull-back $(A,z)\mapsto(g^{*}A,g^{*}z)$ induces a bijection of $\Conf(g\Gamma)$ onto $\Conf(\Gamma)$
which maps the Borel sets of $\Conf(g\Gamma)$ onto those of $\Conf(\Gamma)$ and the measure $\mathrm{d}(A,z)$ of $\Conf(g\Gamma)$ onto that of $\Conf(\Gamma)$:
\[
\int_{\Conf(g\Gamma)}\mathrm{d}(A,z)\,\Phi(g^{*}A,g^{*}z)=\int_{\Conf(\Gamma)}\mathrm{d}(A,z)\,\Phi(A,z)
\]
for every Borel function $\Phi\ge0$ on $\Conf(\Gamma)$.
Indeed, $g^{*}$ maps $g\Gamma$-periodic cochains to $\Gamma$-periodic ones and integer-valued cochains to integer-valued ones, with the inverse $(g^{-1})^{*}$;
it commutes with $d$, carries the equivalence relation of $g\Gamma$ to that of $\Gamma$, and permutes the coordinates $A(b)$ up to signs.
So it maps the product of Lebesgue measure and counting measure on the pairs for $g\Gamma$ onto the corresponding product for $\Gamma$,
and the fundamental domain given by the box for $g\Gamma$ onto a Borel set of pairs for $\Gamma$ which meets every class in exactly one pair
and can therefore take the place of the fundamental domain given by the box for $\Gamma$.
\end{definition}

For the period lattice $\Gamma=n_1\Z e_1+n_2\Z e_2+n_3\Z e_3$, $\min_\mu n_\mu\ge4$,
the integral over $\Conf(\Gamma)$ is the sum over all flux sectors of the integrals over the box
with which the partition functions of \cite[Definition~10.1]{NoncompactQED} are formed;
there the torus is written $\Lambda_n$ and the closed integer two-cochain $\nu$.

\subsection{The chamber fermion}\label{sec:chamber-fermion}

\begin{definition}[Spin matrices]\label{def:spin-matrices}
Let $\sigma_1,\sigma_2,\sigma_3$ be the Pauli matrices.
With respect to the decomposition $\C^4=\C^2\oplus\C^2$ put
\[
\gamma_\mu:=\diag(\sigma_\mu,-\sigma_\mu),\qquad \mu=1,2,3 ;
\]
these are Hermitian $4\times4$ matrices, and $\gamma_\mu\gamma_\nu+\gamma_\nu\gamma_\mu=2\delta_{\mu\nu}$.
For $v\in\R^3$ write $\gamma\cdot v=\gamma_v:=\sum_{\mu}v_\mu\gamma_\mu$.
If $n\in\R^3$ is a unit vector, then $\gamma_n$ is Hermitian with $\gamma_n^2=1$ and trace zero,
so that the matrices
\[
P^{\pm}_n:=\frac{1}{2}\,(1\pm\gamma_n)
\]
are the orthogonal projections onto the eigenspaces of $\gamma_n$ for the eigenvalues $\pm1$;
both have rank two, and $P^{\pm}_{-n}=P^{\mp}_n$.
For such $n$ and every $v\in\R^3$ one has
\[
\gamma_n\gamma_v\gamma_n=-\gamma_{v-2(v\cdot n)n},
\]
since $\gamma_v\gamma_n=2\,(v\cdot n)-\gamma_n\gamma_v$ and $\gamma_n^2=1$.
The matrix $\gamma_5:=\sigma_1\otimes 1_2$, with $1_2$ the unit matrix of size two,
exchanges the two summands of $\C^2\oplus\C^2$;
it anticommutes with $\gamma_1$, $\gamma_2$ and $\gamma_3$.
The charge conjugation matrix $\sigma_2\oplus\sigma_2=\diag(\sigma_2,\sigma_2)$ satisfies
\[
(\sigma_2\oplus\sigma_2)\,\gamma_\mu^{\top}\,(\sigma_2\oplus\sigma_2)^{-1}=-\gamma_\mu,\qquad \mu=1,2,3,
\]
where $\gamma_\mu^{\top}$ is the transposed matrix of $\gamma_\mu$;
this follows from $\sigma_2\,\sigma_\mu^{\top}\,\sigma_2^{-1}=-\sigma_\mu$.
The nine upward unit normals $\hat n_\ell$, $\ell\in\mathcal{L}$ (Definition~\ref{def:mirrors}), and their negatives are eighteen distinct unit vectors.
To each of these vectors $n$ belong two orthonormal bases, $(e^+_{n,1},e^+_{n,2})$ of $\operatorname{ran}P^+_n$ and $(e^-_{n,1},e^-_{n,2})$ of $\operatorname{ran}P^-_n$;
for $n=\hat n_\ell$ they are fixed arbitrarily,
and for $-n$ they are given by $e^{\pm}_{-n,k}:=e^{\mp}_{n,k}$, $k=1,2$, which is legitimate as $\operatorname{ran}P^{\pm}_{-n}=\operatorname{ran}P^{\mp}_n$.
Consequently $e^{\pm}_{-n,k}=e^{\mp}_{n,k}$ for all eighteen vectors $n$ and for $k=1,2$.
\end{definition}

\begin{definition}[Grassmann algebra]\label{def:grassmann}
Let $Y\subset\mathbb{X}$.
To every site $x\in Y$ and every index $\alpha\in\{1,2,3,4\}$ belong two symbols $\psi_\alpha(x)$ and $\bar\psi_\alpha(x)$.
The Grassmann algebra $\Gr(Y)$ over $\C$ is the exterior algebra of the complex vector space which has these symbols as a basis;
the symbols are its generators, so that any two of them anticommute and each of them has square zero.
The generators $\psi_1(x),\dots,\psi_4(x)$ at a site $x$ are arranged as a column $\psi(x)$, and $\bar\psi_1(x),\dots,\bar\psi_4(x)$ as a row $\bar\psi(x)$.
A product $\Psi$ of pairwise different generators is a monomial, and the number of its factors is its degree $\deg\Psi$;
a non-zero linear combination $G$ of monomials of one degree $k$ is homogeneous, of degree $\deg G:=k$.
The even part $\Gr^{\mathrm{ev}}(Y)$ is the span of the monomials of even degree,
and $(-1)^{\deg}$ is the automorphism of $\Gr(Y)$ that multiplies every monomial $\Psi$ by $(-1)^{\deg\Psi}$.
Derivatives with respect to the generators are left derivatives.

Now let $Y$ be finite.
The Berezin integral is the linear form $\int_Y$ on $\Gr(Y)$ written as
\[
\int_Y=\int\prod_{x\in Y}\prod_{\alpha=1}^{4}(\mathrm{d}\bar\psi_\alpha(x)\,\mathrm{d}\psi_\alpha(x))
\]
and normalised pair by pair by the first identity in \eqref{eq:setting-berezin} below.
Every pair $\mathrm{d}\bar\psi_\alpha(x)\,\mathrm{d}\psi_\alpha(x)$ is even, so that the order of the pairs is immaterial;
explicitly, $\int_Y G$ is the coefficient of the product of all pairs $\psi_\alpha(x)\,\bar\psi_\alpha(x)$, $x\in Y$, $\alpha=1,\dots,4$, in $G$.
For a complex matrix $D$ whose rows and columns are indexed by $Y\times\{1,\dots,4\}$ write
\[
\bar\psi D\psi:=\sum_{x,y\in Y}\ \sum_{\alpha,\beta=1}^{4}\bar\psi_\alpha(x)\,D_{(x,\alpha),(y,\beta)}\,\psi_\beta(y).
\]
The normalisation and its consequence for such quadratic elements read
\begin{equation}\label{eq:setting-berezin}
\int\mathrm{d}\bar\psi_\alpha(x)\,\mathrm{d}\psi_\alpha(x)\;\psi_\alpha(x)\,\bar\psi_\alpha(x)=1,
\qquad
\int_Y\exp\left(-\bar\psi D\psi\right)=\det D .
\end{equation}
The second identity holds because the component of top degree of $\exp(-\bar\psi D\psi)$ is $\det D$ times the product of all pairs $\psi_\alpha(x)\,\bar\psi_\alpha(x)$.
\end{definition}

\begin{definition}[Reflections of Grassmann variables]\label{def:grassmann-reflections}
Let $\Pi=(\ell,c)$ be a mirror, $\theta:=\theta_\Pi$ the reflection in $\Pi$ and $n:=\hat n_\ell$ the upward unit normal (Definition~\ref{def:mirrors}).
As an element of $W$, the reflection $\theta$ maps $\mathbb{X}=Wv_0$ onto itself.
The reflection $\Theta_\Pi$ of Grassmann variables is the antilinear map $\Gr(\mathbb{X})\to\Gr(\mathbb{X})$
that reverses the order of products, $\Theta_\Pi(GG')=\Theta_\Pi(G')\,\Theta_\Pi(G)$,
and acts on the generators by
\begin{equation}\label{eq:setting-reflection-fields}
\Theta_\Pi\psi(x)=-\bar\psi(\theta x)\,\gamma_n,\qquad
\Theta_\Pi\bar\psi(x)=-\gamma_n\,\psi(\theta x),\qquad x\in\mathbb{X};
\end{equation}
in components, $\Theta_\Pi\psi_\alpha(x)=-\sum_\beta\bar\psi_\beta(\theta x)\,(\gamma_n)_{\beta\alpha}$
and $\Theta_\Pi\bar\psi_\alpha(x)=-\sum_\beta(\gamma_n)_{\alpha\beta}\,\psi_\beta(\theta x)$.
There is exactly one such map, and $\Theta_\Pi^2=1$ (Lemma~\ref{lem:sign-rules}).
\end{definition}

\begin{definition}[The chamber operator]\label{def:chamber-operator}
The chamber operator depends on the Wilson parameter $\rho\in\R$, on the hopping weights $w_0,w_1,w_2,w_3>0$ and on the mass parameter $M\in\R$.
Where no other range of $\rho$ is stated, $0<\rho<1$.
The lattice mass is
\[
m:=M-\frac{\rho}{2}\sum_{i=0}^{3}w_i ,
\]
and the standard weights are $w_i:=48\cdot|\Pi_i\cap\overline{C_0}|$, $i=0,\dots,3$,
where $|\Pi_i\cap\overline{C_0}|$ is the area of the wall $\Pi_i$ of $C_0$ (Definition~\ref{def:chambers});
explicitly, $(w_0,w_1,w_2,w_3)=(6,6\sqrt{2},6\sqrt{2},6)$.

The photon sites of the two end points of a bond are equal or nearest neighbours: for every $x\in\mathbb{X}$
\[
\pi\left(x^{(i)}\right)=\pi(x)\quad(i=0,1,2),\qquad
\pi\left(x^{(3)}\right)-\pi(x)\in\{\pm e_1,\pm e_2,\pm e_3\}.
\]
To see this, let $g\in W$ be the element with $x=gv_0$, and write $g=\tau_k\circ\vec g$ with $k\in\Z^3$ and $\vec g\in\Ocub$ (Definition~\ref{def:mirrors});
then $\pi(x)=k$ and $x^{(i)}=gs_iv_0=k+\vec g\,s_iv_0$ (Definition~\ref{def:fermion-sites}).
The walls $\Pi_0$, $\Pi_1$, $\Pi_2$ contain the origin, so the reflections $s_0$, $s_1$, $s_2$ in them are linear and lie in $\Ocub$;
hence $\vec g\,s_i\in\Ocub$ and $\pi\left(x^{(i)}\right)=k$ for $i=0,1,2$.
The reflection in the wall $\Pi_3$ is $s_3=\tau_{e_1}\circ\diag(-1,1,1)$,
so that $x^{(3)}=(k+\vec g\,e_1)+\vec g\,\diag(-1,1,1)\,v_0$ with $\vec g\,\diag(-1,1,1)\in\Ocub$, and $\pi\left(x^{(3)}\right)=k+\vec g\,e_1$;
finally, $\vec g\,e_1$ is one of the vectors $\pm e_\mu$, because $\vec g$ is a signed permutation matrix.

Let $A$ be a one-cochain on $\Z^3$ with real or complex values, read as a function of pairs of points as in Definition~\ref{def:cochains},
and let $q>0$ be the lattice charge (Definition~\ref{def:configurations}).
For an ordered pair $(x,x')$ such that $\{x,x'\}$ is a bond of type $i$ (Definition~\ref{def:fermion-sites}), the value $A(\pi(x),\pi(x'))$ is therefore defined,
and the \emph{parallel transport} along the bond and the \emph{hopping matrix} of the bond are
\[
U_A(x,x'):=\exp\left(iqA(\pi(x),\pi(x'))\right),\qquad
w(x,x'):=\frac{w_i}{2}\,(\rho-\gamma\cdot\hat d),\qquad
\hat d:=\frac{x'-x}{|x'-x|} .
\]
The chamber operator $D_A$ acts on functions $\psi\colon\mathbb{X}\to\C^4$ by
\[
(D_A\psi)(x):=M\psi(x)-\sum_{i=0}^{3}w(x,x^{(i)})\cdot U_A(x,x^{(i)})\cdot\psi(x^{(i)}),\qquad x\in\mathbb{X}.
\]
The same expression in the generators of $\Gr(\mathbb{X})$ gives the fermion action
\[
S_F(A):=\bar\psi D_A\psi:=\sum_{x\in\mathbb{X}}\bar\psi(x)\,(D_A\psi)(x),
\]
a formal sum, $\mathbb{X}$ being infinite.

For a function $\lambda\colon\Z^3\to\R$ the gauge transformation by $\lambda$ is
\begin{equation}\label{eq:setting-gauge}
A\mapsto A-d\lambda,\qquad
\psi(x)\mapsto\exp\left(iq\lambda(\pi(x))\right)\cdot\psi(x),\qquad
\bar\psi(x)\mapsto\exp\left(-iq\lambda(\pi(x))\right)\cdot\bar\psi(x).
\end{equation}
It leaves $S_F$ invariant, term by term.
For points $y,y'\in\Z^3$ that are nearest neighbours or equal one has $(d\lambda)(y,y')=\lambda(y')-\lambda(y)$ (Definition~\ref{def:cochains});
for a bond $\{x,x'\}$, with $y=\pi(x)$ and $y'=\pi(x')$, this gives
\[
U_{A-d\lambda}(x,x')=\exp\left(iq\lambda(\pi(x))\right)\cdot U_A(x,x')\cdot\exp\left(-iq\lambda(\pi(x'))\right),
\]
and the two outer factors cancel against the factors acquired by $\bar\psi(x)$ and $\psi(x')$.
The charge density of a monomial $\Psi$ in the generators is the function $\mathfrak{q}(\Psi)\colon\Z^3\to\Z$ given by
\[
\begin{aligned}
\mathfrak{q}(\Psi)(y):={}&\#|\{(x,\alpha):\ \pi(x)=y,\ \psi_\alpha(x)\text{ is a factor of }\Psi\}|\\
&-\#|\{(x,\alpha):\ \pi(x)=y,\ \bar\psi_\alpha(x)\text{ is a factor of }\Psi\}| ;
\end{aligned}
\]
the gauge transformation \eqref{eq:setting-gauge} multiplies $\Psi$ by $\exp\left(iq\sum_{y\in\Z^3}\lambda(y)\cdot\mathfrak{q}(\Psi)(y)\right)$.
\end{definition}

\begin{definition}[Tori]\label{def:tori}
A torus is a pair $\mathbb{T}=(\Gamma,\chi)$ consisting of a period lattice $\Gamma$ (Definition~\ref{def:configurations})
and a character $\chi\colon\Gamma\to\{+1,-1\}$.
Its number of photon sites is $|\mathbb{T}|:=[\Z^3:\Gamma]$.
Translations by vectors of $\Z^3$ belong to $W$ (Definition~\ref{def:mirrors}\eqref{it:mirror-product}),
so that $\Gamma$ acts on $\mathbb{X}=Wv_0$ by translations;
$\mathbb{X}_{\mathbb{T}}:=\mathbb{X}/\Gamma$ is the set of the orbits.
Fields on $\mathbb{T}$ obey
\[
\psi(x+g)=\chi(g)\cdot\psi(x),\qquad
\bar\psi(x+g)=\chi(g)\cdot\bar\psi(x),\qquad
x\in\mathbb{X},\ g\in\Gamma .
\]
Accordingly, $\Gr_{\mathbb{T}}$ is the complex Grassmann algebra generated by elements $\psi_\alpha(x)$, $\bar\psi_\alpha(x)$, $x\in\mathbb{X}$, $\alpha=1,\dots,4$, subject to these relations;
the generators at the points of any set of representatives of $\mathbb{X}_{\mathbb{T}}$ are free generators of $\Gr_{\mathbb{T}}$.
The quotient homomorphism $r_{\mathbb{T}}\colon\Gr(\mathbb{X})\to\Gr_{\mathbb{T}}$ is the homomorphism of algebras
that sends every generator $\psi_\alpha(x)$, $\bar\psi_\alpha(x)$ of $\Gr(\mathbb{X})$ to the generator of $\Gr_{\mathbb{T}}$ with the same name.
Let $Y\subset\mathbb{X}$ be a finite set on which the projection $\mathbb{X}\to\mathbb{X}_{\mathbb{T}}$ is injective.
An even element $G$ of $\Gr(Y)$ defines the element $r_{\mathbb{T}}(G)$ of $\Gr_{\mathbb{T}}$.
This element does not change when $Y$ is translated as a whole by some $g\in\Gamma$,
that is, when $G$ is replaced by the element of $\Gr^{\mathrm{ev}}(Y+g)$ obtained by substituting $\psi_\alpha(x+g)$, $\bar\psi_\alpha(x+g)$ for $\psi_\alpha(x)$, $\bar\psi_\alpha(x)$:
the substitution multiplies the image under $r_{\mathbb{T}}$ of every monomial $\Psi$ that occurs in $G$ by $\chi(g)^{\deg\Psi}$, which is $1$ because $\deg\Psi$ is even.

Let $A$ be a $\Gamma$-periodic one-cochain.
Put
\[
S_{\mathbb{T}}(A):=\sum_{x}\bar\psi(x)\,(D_A\psi)(x)\in\Gr_{\mathbb{T}},\qquad
\int_{\mathbb{T}}:=\int_{\mathbb{X}_{\mathbb{T}}},
\]
where $x$ ranges over a set of representatives of $\mathbb{X}_{\mathbb{T}}$,
and $\int_{\mathbb{X}_{\mathbb{T}}}$ is the Berezin integral of Definition~\ref{def:grassmann} in the free generators at these representatives.
Neither depends on the choice of the representatives.
Indeed, the maps $x\mapsto x^{(i)}$ and $\pi$ commute with the translations by $\Gamma$, the hopping matrices are invariant under them, and $A$ is $\Gamma$-periodic;
hence $(D_A\psi)(x+g)=\chi(g)\cdot(D_A\psi)(x)$ in $\Gr_{\mathbb{T}}$ for $g\in\Gamma$,
the summands at $x$ and at $x+g$ coincide,
and so do the pairs $\psi_\alpha(x)\,\bar\psi_\alpha(x)$ and $\psi_\alpha(x+g)\,\bar\psi_\alpha(x+g)$, because $\chi(g)^2=1$.
The spinors of $\mathbb{T}$ are the functions $\psi\colon\mathbb{X}\to\C^4$ with $\psi(x+g)=\chi(g)\cdot\psi(x)$ for all $g\in\Gamma$;
by the same computation $D_A$ maps spinors of $\mathbb{T}$ to spinors of $\mathbb{T}$,
and $D^{\mathbb{T}}_A$ denotes the operator $D_A$ on this space.
Applied to the matrix of $D^{\mathbb{T}}_A$ in the coordinates given by the values at the representatives,
\eqref{eq:setting-berezin} gives
\[
\int_{\mathbb{T}}\exp\left(-S_{\mathbb{T}}(A)\right)=\det D^{\mathbb{T}}_A .
\]
Let $g\in W$, with linear part $\vec g\in\Ocub$.
By Definition~\ref{def:configurations}, $g\Gamma=\vec g\,\Gamma$ is a period lattice,
and $\chi\circ\vec g^{-1}$ is a character of $g\Gamma$ with values $\pm1$;
the torus
\[
g\mathbb{T}:=\left(g\Gamma,\chi\circ\vec g^{-1}\right)
\]
is the image of $\mathbb{T}$ under $g$.

Three kinds of tori carry names;
their images under elements $g\in\Ocub$ are written $g\,C(n_1,n_2,n_3)$, $g\,D(a,b,c)$, $g\,D'(a,b,c)$ as above.
Here $e_\pm=e_1\pm e_2$.
\begin{itemize}
\item For integers $n_1,n_2,n_3\ge4$, the coordinate torus $C(n_1,n_2,n_3)$ has the period lattice $\Gamma=n_1\Z e_1+n_2\Z e_2+n_3\Z e_3$
and the character with $\chi(n_\mu e_\mu)=-1$ for $\mu=1,2,3$.
It is cubic if $n_1=n_2=n_3$, and $|C(n_1,n_2,n_3)|=n_1n_2n_3$.
\item For integers $a,b,c$ with $\min(a,b,c)\ge4$, the adapted torus $D(a,b,c)$ has the period lattice $\Gamma=\Z a\,e_-+\Z b\,e_++\Z c\,e_3$
and the character with $\chi(a\,e_-)=\chi(b\,e_+)=\chi(c\,e_3)=-1$.
\item For the same $a,b,c$, the torus $D'(a,b,c)$ has the same period lattice
and the character with $\chi(a\,e_-)=\chi(c\,e_3)=-1$ and $\chi(b\,e_+)=+1$.
\end{itemize}
One has $|D(a,b,c)|=|D'(a,b,c)|=2abc$.
Under the stated conditions on the sides these lattices are period lattices:
the systole of the lattice of $C(n_1,n_2,n_3)$ is $\min_\mu n_\mu$,
and that of the lattice of $D(a,b,c)$ and $D'(a,b,c)$ is $\min(a,b,c)$ (Lemma~\ref{lem:systoles}).
For $n=(n_1,n_2,n_3)$ one also writes $C(n):=C(n_1,n_2,n_3)$ and $|n|:=n_1n_2n_3$, so that $|C(n)|=|n|$;
for an integer $n\ge4$, $C(n):=C(n,n,n)$ is the cubic torus of side $n$.
\end{definition}

\subsection{Integrands and the joint law}\label{sec:joint-law}

\begin{definition}[Boxes and integrands]\label{def:integrands}
\emph{Boxes.}
A \emph{box} is a closed cube $B=y+[-r,r]^3$ with $y\in\R^3$ and $r>0$.
Let $B^{(1)}$ and $B^{(2)}$ be the sets of the links and of the plaquettes of $\Z^3$ that are contained in $B$,
and let $\mathbb{X}(B)$ be the set of the fermion sites $x$ whose cube $\pi(x)+\left[-\frac12,\frac12\right]^3$ is contained in $B$;
for $x\in\mathbb{X}(B)$ the photon site $\pi(x)$ lies in $B$.
If $B$ lies in one of the two open half-spaces bounded by a mirror,
then no site $x$ whose cube $\pi(x)+\left[-\frac12,\frac12\right]^3$ meets the mirror belongs to $\mathbb{X}(B)$.
The cells of $\Z^3$ contained in $B$ are the cells of the product of the integer intervals $[\lceil y_\mu-r\rceil,\lfloor y_\mu+r\rfloor]$, $\mu=1,2,3$,
and there is no such cell if one of these intervals is empty.
This cubical complex is contractible;
hence each of its closed cochains of degree one or two is the coboundary of a cochain with values in the same group,
for real and for integer values alike.

\emph{Embedding of boxes.}
Let $\Gamma'$ be a subgroup of $\Z^3$ of rank three or two.
A box $B$ \emph{embeds} in $\Gamma'$ if the sets $B+[-1,1]^3+g$, $g\in\Gamma'$, are pairwise disjoint;
it embeds in a torus $\mathbb{T}=(\Gamma,\chi)$ if it embeds in $\Gamma$.
Boxes $B_1,\dots,B_r$ \emph{embed disjointly} in $\Gamma'$ if each of them embeds in $\Gamma'$ and
\[
(B_j+g)\cap(B_{j'}+g')=\emptyset\qquad\text{for }1\le j,j'\le r\text{ and }g,g'\in\Gamma'\text{ with }(j,g)\neq(j',g').
\]

\emph{Integrands.}
Let $B$ be a box, and fix a monomial basis of $\Gr(\mathbb{X}(B))$,
consisting of the products of distinct generators in a fixed order, the empty product $1$ included.
A Borel map
\[
G\colon\R^{B^{(1)}}\to\Gr(\mathbb{X}(B)),\qquad\hat A\mapsto G(\hat A)=\sum_\Psi a_\Psi(\hat A)\cdot\Psi ,
\]
the sum being over the monomials $\Psi$ of the basis, is called an \emph{integrand that reads in} $B$ if it is gauge invariant in the following sense.
For every monomial $\Psi$ of the basis the charge density $\mathfrak{q}(\Psi)$ of Definition~\ref{def:chamber-operator} vanishes off the sites of $\Z^3$ contained in $B$,
and gauge invariance means that
\[
a_\Psi(\hat A-d\lambda)\cdot\exp\left(iq\langle\lambda,\mathfrak{q}(\Psi)\rangle\right)=a_\Psi(\hat A)
\]
for every real function $\lambda$ on the sites of $\Z^3$ contained in $B$, every $\Psi$ and every $\hat A\in\R^{B^{(1)}}$,
the pairing being the sum over these sites.
By Definition~\ref{def:chamber-operator}, this says that $G$ is unchanged when the gauge transformation \eqref{eq:setting-gauge} by $\lambda$ is applied to $\hat A$ and to the generators.
If $d\hat A'=d\hat A$ on $B^{(2)}$, then $\hat A'=\hat A-d\lambda$ for a real function $\lambda$ on the sites of $B$, by the contractibility of the complex of $B$;
hence $|a_\Psi(\hat A)|$ depends on $d\hat A$ only,
and the coefficient $a_1(\hat A)$ of the unit, whose charge density vanishes, depends on $d\hat A$ only.
The set $\mathcal{G}(B)$ consists of the integrands $G$ that read in $B$ for which there are $C<\infty$ and $\varepsilon_G>0$ with
\[
|a_\Psi(\hat A)|\le C\cdot\exp\left(\left(\frac12-\varepsilon_G\right)\sum_{\sigma\in B^{(2)}}(d\hat A)_\sigma^2\right)
\]
for every $\Psi$ and every $\hat A\in\R^{B^{(1)}}$.
The even elements of $\mathcal{G}(B)$, those with values in $\Gr^{\mathrm{ev}}(\mathbb{X}(B))$, form the set $\mathcal{G}^{\mathrm{ev}}(B)$.
An integrand is \emph{bosonic} if all its values are multiples of the unit;
a bosonic integrand is then a function of $(F_\sigma)_{\sigma\in B^{(2)}}$, $F=d\hat A$.

\emph{Local lifts.}
Suppose now that $B$ embeds in a period lattice $\Gamma$, and fix a representative $(A,z)$ of a point of $\Conf(\Gamma)$;
its two components and the curvature $F=dA+\frac{2\pi}{q}\cdot z$ are regarded as $\Gamma$-periodic cochains on $\Z^3$ (Definition~\ref{def:configurations}).
The \emph{local lifts} on $B$ are the real cochains $\hat A$ on $B^{(1)}$ that satisfy
\[
d\hat A=F\quad\text{on }B^{(2)},\qquad\exp\left(iq\hat A(b)\right)=\exp\left(iqA(b)\right)\quad\text{for }b\in B^{(1)}.
\]
Local lifts exist:
the integer two-cochain $z$ is closed, so on the cells of $B$ it is the coboundary $dk$ of an integer one-cochain $k$,
and $\hat A:=A+(2\pi/q)\cdot k$ is a local lift.
If $\hat A$ and $\hat A'$ are local lifts, then $(q/2\pi)\cdot(\hat A'-\hat A)$ is a closed integer one-cochain on the cells of $B$;
so $\hat A'=\hat A+(2\pi/q)\cdot d\lambda$ with an integer valued function $\lambda$ on the sites of $B$
(Lemma~\ref{lem:product-integrands}\eqref{it:product-torus}).
For an integrand $G$ that reads in $B$, gauge invariance with the function $-(2\pi/q)\cdot\lambda$ gives $a_\Psi(\hat A')=a_\Psi(\hat A)$,
because $\exp\left(-2\pi i\langle\lambda,\mathfrak{q}(\Psi)\rangle\right)=1$ for integer valued $\lambda$.
A second representative $\left(A+(2\pi/q)\cdot k',\,z-dk'\right)$, $k'\in C^1(\Gamma;\Z)$, of the same point has the same curvature and the same link phases, hence the same local lifts.
Therefore $G(\hat A)$ depends neither on the local lift nor on the representative.
Let $\mathbb{T}=(\Gamma,\chi)$ be a torus, with the quotient homomorphism $r_{\mathbb{T}}$ of Definition~\ref{def:tori}.
For an even integrand $G$ that reads in $B$, the element $r_{\mathbb{T}}\left(G(\hat A)\right)$ of $\Gr_{\mathbb{T}}$ is thus a function of the point of $\Conf(\Gamma)$.
If moreover $G\in\mathcal{G}^{\mathrm{ev}}(B)$, this function is Borel:
by Lemma~\ref{lem:borel-lifts}\eqref{it:lift-coefficients} if $B$ contains a site of $\Z^3$,
and trivially otherwise, since then $B^{(1)}$ and $\mathbb{X}(B)$ are empty and $G$ is a constant.

\emph{Reading in a set.}
An integrand $G$ is said to \emph{read in} a set $U\subset\R^3$ if $U$ contains a box $B$ with $G\in\mathcal{G}(B)$.

\emph{Admissible products.}
Let $C^1$ be the space of the real one-cochains on $\Z^3$ (Definition~\ref{def:cochains}),
and let $\Gr(\mathbb{X})^{C^1}$ be the complex vector space of all maps from $C^1$ to $\Gr(\mathbb{X})$, an algebra under the pointwise product.
An integrand $G$ that reads in a box $B$ is regarded as the element $\hat A\mapsto G\left(\hat A|_{B^{(1)}}\right)$ of $\Gr(\mathbb{X})^{C^1}$,
and a bounded Borel function $\Phi$ of finitely many $F_\sigma$ as the element $\hat A\mapsto\Phi(d\hat A)$.
For an open set $U\subset\R^3$, the space $\mathcal{A}(U)$ is the linear span in $\Gr(\mathbb{X})^{C^1}$ of the products
\[
\Phi\cdot G_1\cdots G_r\colon\ \hat A\mapsto\Phi(d\hat A)\cdot G_1\left(\hat A|_{B_1^{(1)}}\right)\cdots G_r\left(\hat A|_{B_r^{(1)}}\right),\qquad r\ge0,
\]
in which $B_1,\dots,B_r$ are pairwise disjoint boxes contained in $U$, $G_j\in\mathcal{G}^{\mathrm{ev}}(B_j)$ for each $j$,
and $\Phi$ is a bounded Borel function of finitely many $F_\sigma$ with $\sigma\subset U$.
As the boxes are pairwise disjoint, no plaquette lies in two of them,
so that the value of $d\hat A$ at a given plaquette enters the growth bound of at most one of the factors $G_j$.
The elements of $\mathcal{A}(U)$ take their values in $\Gr^{\mathrm{ev}}(\mathbb{X})$.

\emph{Embedding of products.}
Whether a product $\Phi\cdot G_1\cdots G_r$ as above \emph{embeds} in $\Gamma'$ is decided by its boxes alone:
by definition, it does so when the family $B_1,\dots,B_r$ embeds disjointly in $\Gamma'$;
for $r=0$ this condition is void, and a product without integrands embeds in every $\Gamma'$.
The plaquettes on which $\Phi$ depends are not restricted, since on a torus $\Phi$ is applied to the periodically extended curvature $\widetilde F$.
For an element of $\mathcal{A}(U)$, to \emph{embed} in $\Gamma'$ means to be the sum of finitely many products each of which embeds in $\Gamma'$.

\emph{Elements of $\mathcal{A}(\R^3)$ on a torus.}
Let $G\in\mathcal{A}(\R^3)$ embed in a torus $\mathbb{T}=(\Gamma,\chi)$, and fix a representative $(A,z)$ of a point of $\Conf(\Gamma)$.
There is a real one-cochain $\hat A$ on $\Z^3$ with $d\hat A=\widetilde F$ and $\exp\left(iq\hat A(b)\right)=\exp\left(iqA(b)\right)$ for every link $b$ of $\Z^3$
(Lemma~\ref{lem:global-lift}\eqref{it:lift-torus});
its restriction to $B^{(1)}$ is a local lift on every box $B$ that embeds in $\Gamma$.
Put
\[
G(A,z):=r_{\mathbb{T}}\left(G(\hat A)\right)\in\Gr_{\mathbb{T}} .
\]
For a product $G=\Phi\cdot G_1\cdots G_r$ that embeds in $\mathbb{T}$ one has
\[
G(A,z)=\Phi(\widetilde F)\cdot r_{\mathbb{T}}\left(G_1(\hat A^{(1)})\right)\cdots r_{\mathbb{T}}\left(G_r(\hat A^{(r)})\right)
\]
for arbitrary local lifts $\hat A^{(j)}$ on $B_j$, and the right side does not depend on these lifts.
A general $G$ is a finite sum of such products, while the value $G(\hat A)$ is that of the map $G$ and involves no choice of such a sum;
so $G(A,z)$ depends only on $G$ and on the point of $\Conf(\Gamma)$, and it is linear in $G$
(Lemma~\ref{lem:product-integrands}\eqref{it:sum-torus}).
In this way every element of $\mathcal{A}(\R^3)$ that embeds in $\mathbb{T}$ is a Borel function on $\Conf(\Gamma)$ with values in the even part of $\Gr_{\mathbb{T}}$.
Likewise, for the open slab $U$ between two mirrors of one family,
the products of $\mathcal{A}(U)$ that embed in the lattice of a cross-section for this family (Definition~\ref{def:slabs}), and their finite sums,
assign to every configuration of the closed slab an element of the Grassmann algebra of the slab
(Lemma~\ref{lem:product-integrands}\eqref{it:product-slab} and~\eqref{it:sum-slab}).
\end{definition}

\paragraph{Reflections of integrands.}
Let $\Pi$ be a mirror, $\theta_\Pi$ its reflection and $\Theta_\Pi$ the reflection of Grassmann variables of Definition~\ref{def:grassmann-reflections}.
Let $G$ be either an integrand that reads in a box $B$ or an element of $\Gr(\mathbb{X})^{C^1}$,
and write $G(\hat A)=\sum_\Psi a_\Psi(\hat A)\cdot\Psi$ in a monomial basis.
The reflected map $\Theta_\Pi G$ is given by
\begin{equation}\label{eq:setting-reflection-coefficients}
(\Theta_\Pi G)(\hat A):=\sum_\Psi\overline{a_\Psi(\theta_\Pi^{*}\hat A)}\cdot\Theta_\Pi(\Psi),
\end{equation}
where $\hat A$ ranges over $\R^{(\theta_\Pi B)^{(1)}}$ in the first case, so that $\theta_\Pi^{*}\hat A\in\R^{B^{(1)}}$, and over $C^1$ in the second.
As $\Theta_\Pi$ is antilinear, $(\Theta_\Pi G)(\hat A)=\Theta_\Pi\left(G(\theta_\Pi^{*}\hat A)\right)$;
thus $\Theta_\Pi G$ does not depend on the monomial basis,
and the two readings agree when an integrand is regarded as an element of $\Gr(\mathbb{X})^{C^1}$.
The image $\theta_\Pi B$ of a box is a box, because the linear part of $\theta_\Pi$ lies in $\Ocub$.
By Lemma~\ref{lem:reflection-classes}, $\Theta_\Pi$ maps the integrands that read in $B$ onto those that read in $\theta_\Pi B$,
the sets $\mathcal{G}(B)$ and $\mathcal{G}^{\mathrm{ev}}(B)$ onto $\mathcal{G}(\theta_\Pi B)$ and $\mathcal{G}^{\mathrm{ev}}(\theta_\Pi B)$,
and $\mathcal{A}(U)$ onto $\mathcal{A}(\theta_\Pi U)$ for every open set $U\subset\R^3$.

\paragraph{The action of the group $W$.}
Let $G$ be an even element of $\Gr(\mathbb{X})^{C^1}$, for instance an element of $\mathcal{G}^{\mathrm{ev}}(B)$ or of $\mathcal{A}(U)$, and let $g\in W$.
The group $W$ is generated by the reflections in the mirrors (Definition~\ref{def:mirrors});
for a word $g=\theta_{\Pi^{(1)}}\cdots\theta_{\Pi^{(r)}}$ in such reflections put
\begin{equation}\label{eq:setting-group-action}
G^g:=\Theta_{\Pi^{(1)}}\cdots\Theta_{\Pi^{(r)}}G .
\end{equation}
By Proposition~\ref{prop:symmetries}\eqref{it:reflection-words}, $G^g$ does not depend on the word that represents $g$.
Juxtaposition of words gives $(G^g)^{g'}=G^{g'g}$ for $g,g'\in W$.
Each $\Theta_{\Pi^{(j)}}$ is an antilinear anti-automorphism of the algebra $\Gr(\mathbb{X})^{C^1}$,
and the determinant of the linear part of $g$ is $(-1)^r$;
hence $G\mapsto G^g$ is a linear automorphism of the algebra of the even elements of $\Gr(\mathbb{X})^{C^1}$ if $g$ is even,
and an antilinear anti-automorphism of it if $g$ is odd.
On bosonic $G$, regarded as functions of $F$, the action reads
\[
G^g(F)=G(g^{*}F)\quad\text{if $g$ is even},\qquad G^g(F)=\overline{G(g^{*}F)}\quad\text{if $g$ is odd};
\]
indeed, by \eqref{eq:setting-reflection-coefficients} each reflection conjugates the value and composes with $\theta_\Pi^{*}$,
the pull-backs commute with $d$,
and $(\theta_{\Pi^{(1)}}\cdots\theta_{\Pi^{(r)}})^{*}=\theta_{\Pi^{(r)}}^{*}\cdots\theta_{\Pi^{(1)}}^{*}$ (Definition~\ref{def:cochains}).
By Lemma~\ref{lem:reflection-classes}, $G^g\in\mathcal{G}^{\mathrm{ev}}(gB)$ if $G\in\mathcal{G}^{\mathrm{ev}}(B)$,
and $G^g\in\mathcal{A}(gU)$ if $G\in\mathcal{A}(U)$.

\begin{definition}[The functional and the law]\label{def:functional}
Let $\rho$ and $M$ be real, $w_0,\dots,w_3>0$ and $q>0$, and let $\mathbb{T}=(\Gamma,\chi)$ be a torus.
For $G\in\mathcal{A}(\R^3)$ that embeds in $\mathbb{T}$ put
\[
Z_{\mathbb{T}}[G]:=\int_{\Conf(\Gamma)}\mathrm{d}(A,z)\,\exp\left(-\frac12\|F\|^2\right)\cdot\int_{\mathbb{T}}\exp\left(-S_{\mathbb{T}}(A)\right)\cdot G(A,z),
\]
where $\mathrm{d}(A,z)$ is the measure of Definition~\ref{def:configurations},
$F=dA+\frac{2\pi}{q}\cdot z$ and $\|F\|^2$ is its norm in $C^2(\Gamma)$,
the action $S_{\mathbb{T}}(A)$ and the Berezin integral $\int_{\mathbb{T}}$ are those of Definition~\ref{def:tori},
and $G(A,z)\in\Gr_{\mathbb{T}}$ is the function of Definition~\ref{def:integrands}.
The action depends on $A$ only through the link phases $\exp\left(iqA(b)\right)$, so that the integrand is a function on $\Conf(\Gamma)$.
Every coefficient of $\exp\left(-\frac12\|F\|^2\right)\cdot\exp\left(-S_{\mathbb{T}}(A)\right)\cdot G(A,z)$ in a monomial basis of $\Gr_{\mathbb{T}}$ is absolutely integrable over $\Conf(\Gamma)$
(Lemma~\ref{lem:gaussian-growth}\eqref{it:chamber-convergence}),
and $Z_{\mathbb{T}}$ is linear on the space of the elements of $\mathcal{A}(\R^3)$ that embed in $\mathbb{T}$, since $G(A,z)$ is linear in $G$.

For $-1\le\rho\le1$ the same formula defines $Z_{\mathbb{T}}[G]$ on two further classes, both introduced in Definition~\ref{def:torus-integrands}:
\begin{enumerate}
\item\label{it:read-on-torus}
the space $\mathcal{A}(\mathbb{T})$ spanned by the products that are read on $\mathbb{T}$,
on which $Z_{\mathbb{T}}$ is well defined and linear (Lemma~\ref{lem:product-integrands}\eqref{it:sum-torus})
and coincides with the functional above at the elements that embed in $\mathbb{T}$;
\item
the class $\mathcal{I}(\mathbb{T})$ of that definition.
\end{enumerate}
For these values of $\rho$ the integral over $\Conf(\Gamma)$ converges absolutely on both classes as well, coefficient by coefficient (Proposition~\ref{prop:berezin}\eqref{it:torus-integrability}).

For a Borel function $\Phi\ge0$ on $\Conf(\Gamma)$, for instance a non-negative Borel function of $F$, put
\[
Z_{\mathbb{T}}[\Phi]:=\int_{\Conf(\Gamma)}\mathrm{d}(A,z)\,\exp\left(-\frac12\|F\|^2\right)\cdot\det D^{\mathbb{T}}_A\cdot\Phi(A,z)\in[0,\infty];
\]
the integrand is non-negative because $\det D^{\mathbb{T}}_A\ge0$ (Proposition~\ref{prop:accretivity}).
If $\Phi$ is at the same time the function $G(A,z)$ of an element $G$ for which $Z_{\mathbb{T}}[G]$ has been defined above, the two values are equal,
since a scalar factor can be taken out of the Berezin integral and $\int_{\mathbb{T}}\exp\left(-S_{\mathbb{T}}(A)\right)=\det D^{\mathbb{T}}_A$.

For all real $\rho$ and $M$, all $w_0,\dots,w_3>0$ and every $q>0$ one has
\[
0<Z_{\mathbb{T}}[1]<\infty
\]
on every torus $\mathbb{T}$, whatever its character, the character $\chi=1$ included (Theorem~\ref{thm:partition-positivity}).
The \emph{expectation} $\mathbb{E}_{\mathbb{T}}[G]:=Z_{\mathbb{T}}[G]/Z_{\mathbb{T}}[1]$ is defined for every $G$ for which $Z_{\mathbb{T}}[G]$ is defined.
The function
\[
Z_{\mathbb{T}}[1]^{-1}\cdot\exp\left(-\frac12\|F\|^2\right)\cdot\det D^{\mathbb{T}}_A
\]
is non-negative, again by Proposition~\ref{prop:accretivity}, and its integral with respect to $\mathrm{d}(A,z)$ is one;
it is therefore the density of a probability measure on $\Conf(\Gamma)$.
The \emph{law} $\mu_{\mathbb{T}}$ of $\mathbb{T}$ is the image of this probability measure under the Borel map $(A,z)\mapsto\widetilde F$ from $\Conf(\Gamma)$ to $\Omega$,
that is, the law of $\widetilde F$ under $\mathbb{E}_{\mathbb{T}}$;
it is a law on $\Omega$.
Let $\Phi$ be a bounded Borel function on $\Omega$.
The $F\in C^2(\Gamma)$ with $dF=0$ form a linear subspace of the finite-dimensional space $C^2(\Gamma)$, hence a closed set,
and $F\mapsto\widetilde F$ maps this subspace into $\Omega$, continuously because each value of $\widetilde F$ is a value of $F$.
So the function on $C^2(\Gamma)$ that equals $\Phi(\widetilde F)$ where $dF=0$ and vanishes elsewhere is bounded and Borel.
Choose one plaquette of $\Z^3$ in each of the finitely many classes of plaquettes modulo $\Gamma$;
the values at the chosen plaquettes are linear coordinates on $C^2(\Gamma)$,
in which that function becomes a bounded Borel function of finitely many $F_\sigma$,
that is, a product without integrands in $\mathcal{A}(\R^3)$.
This product, again denoted by $\Phi(\widetilde F)$, embeds in $\mathbb{T}$,
and its value at a point of $\Conf(\Gamma)$ (Definition~\ref{def:integrands}) is $\Phi(\widetilde F)$, as the curvature $F$ satisfies $dF=0$.
Since $\int_{\mathbb{T}}\exp\left(-S_{\mathbb{T}}(A)\right)\cdot\Phi(\widetilde F)=\Phi(\widetilde F)\cdot\det D^{\mathbb{T}}_A$, the definition of $\mu_{\mathbb{T}}$ yields, for all real $\rho$ and $M$,
\[
\int\mathrm{d}\mu_{\mathbb{T}}(F)\,\Phi(F)=\mathbb{E}_{\mathbb{T}}\left[\Phi(\widetilde F)\right].
\]
\end{definition}

Different formal sums $\sum_ka_k\cdot\varphi_{\beta_k}$ may define the same function on $\Omega$, as part~(a) of the following lemma shows.
A trigonometric polynomial $O$ is therefore understood as the formal sum,
that is, as the finite family of its coefficients $a_k$ together with the pairwise different cochains $\beta_k$;
the number $|O|_1$, the set $\supp O$ and the integer $r_O$ of Definition~\ref{def:characters} are formed with this representation.

\begin{lemma}[Characters of the curvature]\label{lem:curvature-characters}
Let $\beta\in C^2_{\mathrm{c}}$.
\begin{enumerate}
\item
If $\beta'\in C^2_{\mathrm{c}}$ and $\beta-\beta'\in d^{*}C^3_{\mathrm{c}}$, then $\varphi_\beta=\varphi_{\beta'}$ on $\Omega$.
\item
Let $\Gamma$ be a period lattice, and let $\supp\beta$ be contained in a box $B$ that embeds in $\Gamma$.
Let $\beta_\Gamma\in C^2(\Gamma)$ be the periodisation of $\beta$,
given by $\beta_\Gamma(\sigma+g):=\beta(\sigma)$ for the plaquettes $\sigma\subset B$ and the elements $g\in\Gamma$,
and by $\beta_\Gamma:=0$ on the plaquettes that are not of this form.
Then, for every point $(A,z)$ of $\Conf(\Gamma)$ and $F=dA+\frac{2\pi}{q}\cdot z$,
\[
\varphi_\beta(\widetilde F)=\exp\left(i\langle\beta_\Gamma,F\rangle\right),
\]
the pairing on the right being that of $C^2(\Gamma)$.
\end{enumerate}
\end{lemma}

\begin{proof}
(a)
Write $\beta-\beta'=d^{*}u$ with $u\in C^3_{\mathrm{c}}$, and let $F\in\Omega$.
The cochain $d^{*}u$ is supported on the faces of the cubes at which $u$ does not vanish.
Let $F'\in C^2_{\mathrm{c}}$ agree with $F$ on these finitely many faces and vanish elsewhere.
Then $dF'=dF=0$ on every cube at which $u$ does not vanish, and
\[
\langle\beta-\beta',F\rangle=\langle d^{*}u,F'\rangle=\langle u,dF'\rangle=0 ;
\]
hence $\exp\left(i\langle\beta,F\rangle\right)=\exp\left(i\langle\beta',F\rangle\right)$.

(b)
The sets $B+[-1,1]^3+g$, $g\in\Gamma$, are pairwise disjoint, and so are the sets $B+g$.
If $\sigma+g=\sigma'+g'$ with plaquettes $\sigma,\sigma'\subset B$ and $g,g'\in\Gamma$,
then $\sigma'$ lies in $B\cap(B+g-g')$, which forces $g=g'$ and $\sigma=\sigma'$.
Thus $\beta_\Gamma$ is well defined, and distinct plaquettes contained in $B$ represent distinct plaquettes of $\Z^3/\Gamma$.
As $\beta$ vanishes on the plaquettes not contained in $B$, and $\widetilde F$ is the periodic extension of $F$,
\[
\langle\beta,\widetilde F\rangle=\sum_{\sigma\subset B}\beta(\sigma)\cdot\widetilde F_\sigma=\sum_{\sigma\subset B}\beta_\Gamma(\sigma)\cdot F_\sigma=\langle\beta_\Gamma,F\rangle ,
\]
the plaquettes of $\Z^3/\Gamma$ without a representative in $B$ contributing zero to the last pairing.
\end{proof}

\begin{definition}[Polynomial classes]\label{def:polynomials}
The class $\mathcal{P}$ is the linear span in $\Gr(\mathbb{X})^{C^1}$ (Definition~\ref{def:integrands}) of the maps
\[
\hat A\mapsto\exp\left(i\langle\delta,\hat A\rangle\right)\cdot\Psi ,
\]
in which $\delta\in C^1_{\mathrm{c}}$ is real, $\Psi$ is a monomial in the generators $\psi_\alpha(x)$, $\bar\psi_\alpha(x)$, $x\in\mathbb{X}$, and
\[
d^{*}\delta=q\cdot\mathfrak{q}(\Psi).
\]
The gauge transformation \eqref{eq:setting-gauge} by a real function $\lambda$ on $\Z^3$ multiplies $\exp\left(i\langle\delta,\hat A\rangle\right)\cdot\Psi$
by $\exp\left(i\langle q\cdot\mathfrak{q}(\Psi)-d^{*}\delta,\lambda\rangle\right)$ (Definition~\ref{def:chamber-operator});
so the condition on $\delta$ says that these maps are gauge invariant.

The elements of $\mathcal{P}$ are even.
Indeed, $\sum_y(d^{*}\delta)(y)=0$ for finitely supported $\delta$, each link contributing its value at one end point and its negative at the other;
so $\sum_y\mathfrak{q}(\Psi)(y)=0$, and $\Psi$ has as many factors $\psi_\alpha(x)$ as factors $\bar\psi_\alpha(x)$.
Let $P=\sum_\alpha a_\alpha\cdot\exp\left(i\langle\delta_\alpha,\hat A\rangle\right)\cdot\Psi_\alpha$ be an element of $\mathcal{P}$,
and let $B$ be a box such that every $\delta_\alpha$ is supported on $B^{(1)}$ and every $\Psi_\alpha$ is a monomial in generators at sites of $\mathbb{X}(B)$.
Then $P\in\mathcal{G}(B)$, its coefficients having modulus at most $\sum_\alpha|a_\alpha|$;
so every element of $\mathcal{P}$ belongs to $\mathcal{G}(B)$ for a box $B$ that contains it.
The coefficients of $P$ extend to entire functions of the complex variables $\hat A(b)$.
An element of $\mathcal{P}$ is bosonic if and only if it is a finite linear combination of characters
\[
\hat A\mapsto\varphi_\beta(d\hat A)=\exp\left(i\langle d^{*}\beta,\hat A\rangle\right),\qquad\beta\in C^2_{\mathrm{c}},
\]
the characters of Definition~\ref{def:characters} being read as functions of $\hat A$.
Indeed, the maps $\hat A\mapsto\exp\left(i\langle\delta,\hat A\rangle\right)$ with pairwise different $\delta$ are linearly independent, being distinct characters of the additive group $C^1$;
so a bosonic element of $\mathcal{P}$ is a linear combination of such maps with $d^{*}\delta=q\cdot\mathfrak{q}(1)=0$,
and every finitely supported real one-cochain $\delta$ with $d^{*}\delta=0$ is of the form $d^{*}\beta$ with $\beta\in C^2_{\mathrm{c}}$
(Lemma~\ref{lem:lines-fillings}\eqref{it:integer-filling}).
Conversely, each such character is the map $\hat A\mapsto\exp\left(i\langle\delta,\hat A\rangle\right)\cdot1$ with $\delta=d^{*}\beta$, and $d^{*}\delta=d^{*}d^{*}\beta=0=q\cdot\mathfrak{q}(1)$.
The class $\mathcal{P}$ is stable under $G\mapsto G^g$ for every $g\in W$
(Proposition~\ref{prop:class-structure}\eqref{it:class-stable}; Lemma~\ref{lem:reflection-classes}).

Let $\mathbb{T}=(\Gamma,\chi)$ be a torus.
A function $\Phi$ on $C^2(\Gamma)$ is an \emph{entire bosonic factor} on $\mathbb{T}$ if it is the restriction of an entire function on $C^2(\Gamma)\otimes\C$
whose modulus is bounded on each of the tubes
\[
\left\{F\in C^2(\Gamma)\otimes\C:\ |\Im F_\sigma|\le r\ \text{for every plaquette }\sigma\right\},\qquad r>0 .
\]
For instance, $\exp(t\cdot V)$, with $t\in\R$ and $V$ a finite sum of trigonometric polynomials evaluated at $\widetilde F$, is an entire bosonic factor on every torus:
a character $F\mapsto\varphi_\beta(\widetilde F)=\exp\left(i\langle\beta,\widetilde F\rangle\right)$ is entire in $F$,
and on the set above its modulus is at most $\exp\left(r\cdot\sum_\sigma|\beta(\sigma)|\right)$.
\end{definition}

\begin{definition}[Reflection positivity of laws]\label{def:reflection-positivity}
Let $(\ell,c)$ be a mirror.
The set $\Sigma^{+}(\ell,c)$ consists of the plaquettes $\sigma$ with $\ell(x)>c$ at each corner $x$ of $\sigma$,
and $\mathcal{O}^{+}(\ell,c)$ is formed by the bounded Borel functions on $\Omega$ that depend only on the values $F_\sigma$ at the plaquettes of some finite subset of $\Sigma^{+}(\ell,c)$.
The pull-back $\theta_{\ell,c}^{*}$ commutes with $d$ and therefore maps $\Omega$ onto itself.
A law $\mu$ on $\Omega$ is said to be \emph{invariant at} $(\ell,c)$ if the map $F\mapsto\theta_{\ell,c}^{*}F$ carries $\mu$ to itself,
and \emph{reflection positive at} $(\ell,c)$ if
\[
\int\mathrm{d}\mu(F)\,\overline{\Phi(\theta_{\ell,c}^{*}F)}\cdot\Phi(F)\ge0\qquad\text{for every }\Phi\in\mathcal{O}^{+}(\ell,c).
\]
\end{definition}

\subsection{Cells, folds and comparison states}\label{sec:cells-folds}

\begin{definition}[Slabs and adapted tori]\label{def:slabs}
Fix a family $\ell\in\mathcal{L}$.
A \emph{cross-section} for $\ell$ is a pair $\Xi=(\Gamma_\parallel,\chi_\parallel)$
in which $\Gamma_\parallel$ is a subgroup of rank two of $\Z^3\cap\{\ell=0\}$ with $\sys(\Gamma_\parallel)\ge4$
and $\chi_\parallel\colon\Gamma_\parallel\to\{1,-1\}$ is a character.
For $c<c'$ in $\Adm(\ell)$ the open set
\[
\{c<\ell<c'\}:=\{x\in\R^3:\ c<\ell(x)<c'\}
\]
is the \emph{open slab} between the mirrors $(\ell,c)$ and $(\ell,c')$,
and $\mathcal{A}_\Xi[c,c']$ denotes the linear span of those products of $\mathcal{A}(\{c<\ell<c'\})$
that embed in $\Gamma_\parallel$ in the sense of Definition~\ref{def:integrands}.
The boxes of these products and the plaquettes on which their bounded factors depend lie in the open slab,
so that nothing that meets one of the planes $\{\ell=c\}$, $\{\ell=c'\}$ is read.

Let $\mathbb{T}=(\Gamma,\chi)$ be a torus.
Its cross-section for $\ell$ is
\[
\Xi(\mathbb{T},\ell):=\left(\Gamma\cap\{\ell=0\},\ \chi|_{\Gamma\cap\{\ell=0\}}\right);
\]
by Lemma~\ref{lem:closing-vectors} it is indeed a cross-section for $\ell$.
A \emph{closing vector} of $\mathbb{T}$ for $\ell$ is an element $g\in\Gamma$
with $\Gamma=(\Gamma\cap\{\ell=0\})+\Z g$ and $\ell(g)>0$;
such vectors exist (Lemma~\ref{lem:closing-vectors}).
For a closing vector $g$ one has $\ell(\Gamma)=\ell(g)\Z$,
so that the number $p(\mathbb{T},\ell):=\ell(g)$ does not depend on $g$;
it is the \emph{height period} of $\mathbb{T}$ for $\ell$.

Write $p:=p(\mathbb{T},\ell)$ and
\[
v:=\frac{p}{|n_\ell|^2}\cdot n_\ell\in\R^3 .
\]
Suppose that $v\in\Z^3$ and that the group
\[
\Gamma^0:=(\Gamma\cap\{\ell=0\})\oplus\Z v
\]
satisfies $\sys(\Gamma^0)\ge4$;
the sum is direct, since $\ell$ vanishes on the first summand and $\ell(v)=p\neq0$.
Then $\Gamma^0$ is a period lattice,
and the \emph{straightened torus} is
\[
\mathbb{T}^0=\mathbb{T}^0(\ell):=(\Gamma^0,\chi^0),
\]
where $\chi^0$ is the character of $\Gamma^0$ that agrees with $\chi$ on $\Gamma\cap\{\ell=0\}$
and takes the value $-1$ at $v$.
In this case $\mathbb{T}^0(\ell)$ is said to be defined, and $|\mathbb{T}^0|=|\mathbb{T}|$ (Lemma~\ref{lem:closing-vectors}).
If one of the two conditions fails, $\Gamma^0$ is not a period lattice, and $\mathbb{T}^0(\ell)$ is not defined.

The torus $\mathbb{T}$ is \emph{adapted to} $\ell$ if $\mathbb{T}^0(\ell)$ is defined and equal to $\mathbb{T}$.
Let $\mathbb{T}$ be adapted to $\ell$, and let $c\in\Adm(\ell)$ with $c+\frac p2\in\Adm(\ell)$.
Then $\Gamma=\Gamma^0=(\Gamma\cap\{\ell=0\})\oplus\Z v$, in particular $v\in\Gamma$;
the linear part of $\theta_{\ell,c}$ fixes $\{\ell=0\}$ pointwise and reverses $n_\ell$,
hence maps $\Gamma$ onto itself, and $\theta_{\ell,c}$ acts on $\R^3/\Gamma$.
By Definition~\ref{def:mirrors}\eqref{it:mirror-involution} and~\eqref{it:mirror-product},
\[
\theta_{\ell,c+p/2}=\tau_v\,\theta_{\ell,c},
\]
so that $\theta_{\ell,c}$ and $\theta_{\ell,c+p/2}$ induce the same map of $\R^3/\Gamma$.
The two mirrors $(\ell,c)$ and $\left(\ell,c+\frac p2\right)$ cut the torus into two halves,
the images in $\R^3/\Gamma$ of the open slabs $\left\{c<\ell<c+\frac p2\right\}$ and $\left\{c+\frac p2<\ell<c+p\right\}$:
as $\ell(\Gamma)=p\Z$, these images are disjoint, and their union is the complement of the images of the two planes.
By Definition~\ref{def:mirrors}\eqref{it:mirror-height}, $\theta_{\ell,c}$ maps the first slab onto $\left\{c-\frac p2<\ell<c\right\}$,
the translate of the second one by $-v$;
thus the map of $\R^3/\Gamma$ induced by $\theta_{\ell,c}$ exchanges the two halves.
\end{definition}

\begin{definition}[Frames and cell structures]\label{def:cell-structures}
A \emph{frame} is a triple $\mathfrak{F}=(\ell^1,\ell^2,\ell^3)$ of families
whose normals $n_a:=n_{\ell^a}$ are pairwise orthogonal.
The \emph{coordinate frame} $\mathfrak{F}_{\mathrm{c}}:=(x_1,x_2,x_3)$
and the \emph{diagonal frame} $\mathfrak{F}_{\mathrm{d}}:=(\ell_-,\ell_+,x_3)$ are frames,
and up to the order and the signs of its families every frame is the image of one of these two under an element of $\Ocub$.

Let $\mathfrak{F}$ be a frame,
and let $p_1,p_2,p_3$ be positive integers with
\[
\frac{p_a}{|n_a|^2}\cdot n_a\in\Z^3\qquad\text{and}\qquad p_a\ge4\,|n_a|^2\qquad(a=1,2,3).
\]
The \emph{rectangular torus} $\mathbb{T}(\mathfrak{F};p_1,p_2,p_3)$ has the lattice
\[
\Gamma=\Z\,\frac{p_1}{|n_1|^2}\cdot n_1\oplus\Z\,\frac{p_2}{|n_2|^2}\cdot n_2\oplus\Z\,\frac{p_3}{|n_3|^2}\cdot n_3
\]
and the character $\chi$ that takes the value $-1$ at each of the three generators.
This lattice has systole at least four (Lemma~\ref{lem:systoles}), so that $\mathbb{T}(\mathfrak{F};p_1,p_2,p_3)$ is a torus.
Since $\ell^a(n_b)=0$ for $b\neq a$, the normals being orthogonal, one has $\ell^a(\Gamma)=p_a\Z$:
the integer $p_a$ is the height period of $\mathbb{T}(\mathfrak{F};p_1,p_2,p_3)$ for $\ell^a$.
The coordinate tori and the adapted tori of Definition~\ref{def:tori} are rectangular:
\[
\mathbb{T}(\mathfrak{F}_{\mathrm{c}};p_1,p_2,p_3)=C(p_1,p_2,p_3),\qquad
\mathbb{T}(\mathfrak{F}_{\mathrm{d}};p_1,p_2,p_3)=D(p_1/2,p_2/2,p_3);
\]
for the diagonal frame one has $|n_1|^2=|n_2|^2=2$ and $|n_3|^2=1$, and the integrality condition says that $p_1$ and $p_2$ are even.
For these two frames the conditions $p_a\ge4\,|n_a|^2$ read $\min(p_1,p_2,p_3)\ge4$ and $\min(p_1/2,p_2/2,p_3)\ge4$, respectively,
which are the conditions on the sides in Definition~\ref{def:tori}.
A rectangular torus is adapted to each of its three families:
$\frac{p_a}{|n_a|^2}\cdot n_a$ is a closing vector for $\ell^a$,
$\Gamma\cap\{\ell^a=0\}$ is spanned by the two other generators,
and the straightened torus $\mathbb{T}^0(\ell^a)$ of Definition~\ref{def:slabs} has the lattice $\Gamma$ and the character $\chi$.

A \emph{cell structure} $\mathcal{Q}=(\mathfrak{F};o_1,o_2,o_3;h)$ consists of a frame $\mathfrak{F}=(\ell^1,\ell^2,\ell^3)$,
positions $o_a\in\Adm(\ell^a)$ and a positive integer $h$, the \emph{width}.
Its \emph{walls} are the mirrors $(\ell^a,o_a+j\cdot h)$, $a=1,2,3$, $j\in\Z$
(note that $o_a+j\cdot h\in\Adm(\ell^a)$, $h$ being an integer),
and its \emph{open cells} are the sets
\[
Q_\kappa:=\left\{x\in\R^3:\ o_a+\kappa_a\cdot h<\ell^a(x)<o_a+(\kappa_a+1)\cdot h,\ a=1,2,3\right\},\qquad\kappa\in\Z^3.
\]
The number
\[
|Q|:=\frac{h^3}{|n_1|\cdot|n_2|\cdot|n_3|}
\]
is the volume of an open cell;
it equals $h^3$ for $\mathfrak{F}_{\mathrm{c}}$ and $h^3/2$ for $\mathfrak{F}_{\mathrm{d}}$.
The group $W(\mathcal{Q})$ is the group generated by the reflections in the walls.
Reflections in walls of different families commute, because $\ell^b\circ\theta_{\ell^a,c}=\ell^b$ for $b\neq a$;
so $W(\mathcal{Q})$ is the direct product of three infinite dihedral groups, one for each family.
It acts simply transitively on the open cells (Definition~\ref{def:cell-families}).

Let $\mathbb{T}=(\Gamma,\chi)$ be a torus.
The structure $\mathcal{Q}$ \emph{cuts} $\mathbb{T}$ if $\ell^a(g)\in h\Z$ for all $g\in\Gamma$ and $a=1,2,3$.
In that case every translation $\tau_g$, $g\in\Gamma$, permutes the open cells:
it maps $Q_\kappa$ onto $Q_{\kappa'}$ with $\kappa'_a=\kappa_a+\ell^a(g)/h$.
The \emph{cells of} $\mathbb{T}$ are the $\Gamma$-orbits of open cells.
Their number is $|\mathbb{T}|/|Q|$:
the map $g\mapsto\left(\ell^1(g)/h,\ell^2(g)/h,\ell^3(g)/h\right)$ is an injective homomorphism of $\Gamma$ into $\Z^3$,
the orbits of cells correspond to the cosets of its image,
and this image has the index $[\Z^3:\Gamma]\cdot|n_1|\cdot|n_2|\cdot|n_3|\cdot h^{-3}=|\mathbb{T}|/|Q|$ in $\Z^3$.
In this sense $|Q|$ is the number of photon sites per cell.

The torus $\mathbb{T}$ \emph{carries} $\mathcal{Q}$ if $\mathbb{T}=\mathbb{T}(\mathfrak{F};p_1,p_2,p_3)$ with $p_a\in2h\Z$ for $a=1,2,3$,
that is, if in each of the three directions one period of $\mathbb{T}$ consists of an even number of slabs between consecutive walls.
Then $\mathcal{Q}$ cuts $\mathbb{T}$,
and the translations $\tau_g$, $g\in\Gamma$, belong to $W(\mathcal{Q})$:
by Definition~\ref{def:mirrors}\eqref{it:mirror-product}, $\theta_{\ell^a,o_a+h}\,\theta_{\ell^a,o_a}$ is the translation by $\frac{2h}{|n_a|^2}\cdot n_a$,
and the generators $\frac{p_a}{|n_a|^2}\cdot n_a$ of $\Gamma$ are integer multiples of these vectors.
The quotient $W(\mathcal{Q})/\Gamma$ acts simply transitively on the cells of $\mathbb{T}$ (Definition~\ref{def:cell-families}).

Let $\mathcal{Q}'=(\mathfrak{F};o'_1,o'_2,o'_3;h')$ be another cell structure with the same frame $\mathfrak{F}$.
When its width $h'$ and the three differences $o'_a-o_a$ of positions all lie in $h\Z$, one calls $\mathcal{Q}'$ \emph{coarser} than $\mathcal{Q}$.
For such a structure each wall of $\mathcal{Q}'$ is a wall of $\mathcal{Q}$,
and each open cell of $\mathcal{Q}'$ contains exactly $\left(\frac{h'}{h}\right)^3$ open cells of $\mathcal{Q}$, whose closures cover it.
\end{definition}

\begin{definition}[Mirror copies and folds]\label{def:folds}
Let $\mathcal{Q}$ be a cell structure, $\kappa\in\Z^3$, and $G\in\mathcal{A}(Q_\kappa)$ (Definition~\ref{def:integrands});
like every element of $\mathcal{A}(U)$, the integrand $G$ is even.
For $\kappa'\in\Z^3$ let $g_{\kappa'}$ be the element of $W(\mathcal{Q})$ with $g_{\kappa'}Q_\kappa=Q_{\kappa'}$.
The \emph{mirror copy} of $G$ in the cell $Q_{\kappa'}$ is
\[
G^{[\kappa']}:=G^{g_{\kappa'}},
\]
with the action of $W$ on even integrands given by \eqref{eq:setting-group-action}.
Elements of $\mathcal{A}(Q_\kappa)$ and of $\mathcal{A}(Q_{\kappa'})$ commute, both being even.

Let now $\mathbb{T}=(\Gamma,\chi)$ be a torus that carries $\mathcal{Q}$.
The \emph{fold} of $G$ on $\mathbb{T}$ is
\[
\Fold_{\mathbb{T}}G:=\prod_{\kappa'}G^{[\kappa']},
\]
where $\kappa'$ ranges over a set of representatives of the cells of $\mathbb{T}$.
Let $g\in\Gamma$ and $\tau_gQ_{\kappa'}=Q_{\kappa''}$.
Since $\tau_g$ lies in $W(\mathcal{Q})$ (Definition~\ref{def:cell-structures}), the element $\tau_g\,g_{\kappa'}$ of $W(\mathcal{Q})$ maps $Q_\kappa$ onto $Q_{\kappa''}$;
hence $\tau_g\,g_{\kappa'}=g_{\kappa''}$, and $(G^{[\kappa']})^{\tau_g}=G^{\tau_g\,g_{\kappa'}}=G^{[\kappa'']}$ by \eqref{eq:setting-group-action}.
With this relation, Definition~\ref{def:cell-families} shows that the product $\Fold_{\mathbb{T}}G$ embeds in $\mathbb{T}$
and that the function on $\Conf(\Gamma)$ which it defines does not depend on the choice of the representatives.
One has $Z_{\mathbb{T}}[\Fold_{\mathbb{T}}G]\ge0$ (Theorem~\ref{thm:chessboard}(a)),
and the \emph{chessboard seminorm} of $G$ on $\mathbb{T}$ is
\[
\|G\|_{\mathbb{T}}:=Z_{\mathbb{T}}\left[\Fold_{\mathbb{T}}G\right]^{\frac{|Q|}{|\mathbb{T}|}},
\]
the exponent being the inverse of the number of cells of $\mathbb{T}$.
\end{definition}

\begin{definition}[Arrays]\label{def:arrays}
Consider a cell structure $\mathcal{Q}=(\mathfrak{F}_{\mathrm{d}};c_0,c_+,c_3;h)$ of the diagonal frame with an even width $h$;
thus $c_0,c_+\in\Z$ and $c_3\in\frac12\Z$.
Let $Q_0$ be its open cell with $\kappa=0$,
and let $O\in\mathcal{T}$ be a real trigonometric polynomial with $\supp O\subset Q_0$ (Definition~\ref{def:characters}).
Put
\[
\Lambda_h:=\frac h2\cdot(\Z e_-+\Z e_+)+h\Z e_3 .
\]
The linear bijection $v\mapsto\left(\ell_-(v),\ell_+(v),v_3\right)$ of $\R^3$ takes $e_-$, $e_+$, $e_3$ to $(2,0,0)$, $(0,2,0)$, $(0,0,1)$;
so a vector $v\in\R^3$ has $\ell_-(v),\ell_+(v),v_3\in h\Z$ if and only if $v\in\Lambda_h$,
and $\Lambda_h\subset\Z^3$ because $h$ is even.
Hence
\[
\Lambda_h=\left\{v\in\Z^3:\ \ell_-(v),\ \ell_+(v),\ v_3\in h\Z\right\},
\]
and $v\mapsto\kappa_v:=\frac1h\cdot\left(\ell_-(v),\ell_+(v),v_3\right)$ is a bijection of $\Lambda_h$ onto $\Z^3$.
For $v,\kappa\in\Z^3$ the set $\tau_vQ_\kappa$ is described by the inequalities that define $Q_\kappa$,
with the bounds for $\ell_-$, $\ell_+$ and $x_3$ increased by $\ell_-(v)$, $\ell_+(v)$ and $v_3$;
it is an open cell if and only if these three numbers lie in $h\Z$, that is, if $v\in\Lambda_h$, and then $\tau_vQ_\kappa=Q_{\kappa+\kappa_v}$.
Consequently $\Lambda_h$ consists of the vectors $v\in\Z^3$ for which the translation $\tau_v$ maps open cells onto open cells,
and it acts simply transitively on the open cells.
The polynomial $O$ is a bounded Borel function of finitely many $F_\sigma$, $\sigma\subset Q_0$;
hence $O\in\mathcal{A}(Q_0)$, its mirror copies $O^{[\kappa]}$ are defined (Definition~\ref{def:folds}),
and, $O$ being real, $O^{[\kappa]}(F)=O(g_\kappa^*F)$ for the element $g_\kappa\in W(\mathcal{Q})$ with $g_\kappa Q_0=Q_\kappa$.
The \emph{unfolded array} and the \emph{folded array} of $O$ are the formal sums
\[
\widetilde{O}:=\sum_{v\in\Lambda_h}\tau_vO
\qquad\text{and}\qquad
\sum_{\kappa\in\Z^3}O^{[\kappa]}.
\]
Each of them has one term in each open cell:
$\tau_vO$ reads the plaquettes of $\supp O+v$, which lie in $Q_0+v$,
and $O^{[\kappa]}$ reads plaquettes contained in $Q_\kappa$.

Let $\mathbb{T}=(\Gamma,\chi)$ be a torus.
If $\mathcal{Q}$ cuts $\mathbb{T}$, then $\Gamma\subset\Lambda_h$,
and on $\mathbb{T}$ the unfolded array is the function $\sum_v(\tau_vO)(\widetilde{F})$ on $\Conf(\Gamma)$,
where $v$ ranges over a set of representatives of $\Lambda_h/\Gamma$.
If $\mathbb{T}$ carries $\mathcal{Q}$, then on $\mathbb{T}$ the folded array is the function $\sum_\kappa O^{[\kappa]}(\widetilde{F})$,
where $\kappa$ ranges over a set of representatives of the cells of $\mathbb{T}$.
These values do not depend on the representatives, since $\widetilde{F}$ is $\Gamma$-periodic
and, in the second case, $\tau_g\in W(\mathcal{Q})$ for $g\in\Gamma$.
In both cases the sum has $|\mathbb{T}|/|Q|$ terms, one for each cell of $\mathbb{T}$,
and each term has modulus at most $|O|_1$;
so both arrays are then bounded functions on $\Conf(\Gamma)$ of modulus at most $\frac{|\mathbb{T}|}{|Q|}\cdot|O|_1$.
\end{definition}

\begin{definition}[Sequences of sides]\label{def:sides}
Let $\mathfrak{n}=(n_i)_{i\ge1}$ be a sequence of positive integers with $n_i\to\infty$, and let
\[
\mathcal{D}(\mathfrak{n}):=\left\{k\in\N:\ k\text{ divides }n_i\text{ for all but finitely many }i\right\}.
\]
This set contains every divisor of each of its elements and the least common multiple of any two of them.
The sequence $\mathfrak{n}$ satisfies the \emph{divisibility condition} if
\begin{equation}\tag{Div}\label{cond:divisibility}
8\in\mathcal{D}(\mathfrak{n}),\qquad\#|\mathcal{D}(\mathfrak{n})|=\infty.
\end{equation}
For an even integer $n\ge8$ put
\[
D(n):=D(n/2,n,n),\qquad D'(n):=D'(n/2,n,n);
\]
the cubic torus is $C(n)=C(n,n,n)$, as in Definition~\ref{def:tori}.
For every index $i$ with $n_i$ even and $n_i\ge8$ the four tori
\[
C(n_i)\ \text{(cubic)},\qquad D'(n_i),\qquad D(n_i)\ \text{(adapted)},\qquad C(n_i/2,2n_i,n_i)\ \text{(elongated)}
\]
are defined, their lattices having systole at least four (Lemma~\ref{lem:systoles}),
and each of them has $n_i^3$ photon sites.
For integers $a_1,a_2,a_3\ge4$ the \emph{pressure} of the coordinate torus $C(a_1,a_2,a_3)$ is
\[
f(a_1,a_2,a_3):=\frac{1}{a_1\cdot a_2\cdot a_3}\cdot\log Z_{C(a_1,a_2,a_3)}[1],
\]
the partition function being positive by Theorem~\ref{thm:partition-positivity}.
For the indices $i$ above put
\[
\Delta_i:=f(n_i/2,2n_i,n_i)-f(n_i,n_i,n_i),\qquad\mu_i:=\mu_{C(n_i)};
\]
thus $\mu_i$ is the law of the cubic torus $C(n_i)$ (Definition~\ref{def:functional}).
\end{definition}

\begin{remark}
If every positive integer divides $n_i$ for all large $i$, then $\mathcal{D}(\mathfrak{n})=\N$, and \eqref{cond:divisibility} holds.
It holds as well for $n_i=8\cdot L^{\lfloor\log_L\bar m^{-1}\rfloor}\cdot L^{k_i}$ with an odd integer $L\ge3$, a number $0<\bar m\le1$ and integers $k_i\to\infty$,
which is the form of the sides in Definition~\ref{def:cutoffs}:
here $8L^j\in\mathcal{D}(\mathfrak{n})$ for every $j\ge0$.
For $n_i=2p_i$ with primes $p_i$ one has $\mathcal{D}(\mathfrak{n})=\{1,2\}$, and \eqref{cond:divisibility} fails.
Condition~\eqref{cond:divisibility} is imposed on a single sequence of sides,
from which the four tori above, the data of Definition~\ref{def:data} and the comparison states of Definition~\ref{def:comparison-states} are all built.
In Theorem~\ref{thm:cubic-limits}, by contrast, no such condition is placed on the sides of the cubic tori:
any sequence of integers $n_i\ge4$ with $n_i\to\infty$ possesses a companion sequence satisfying \eqref{cond:divisibility}
(Definition~\ref{def:companions} and Lemma~\ref{lem:companions-exist}).
\end{remark}

\begin{definition}[Data and perturbed functionals]\label{def:data}
Suppose that $\mathfrak{n}$ satisfies \eqref{cond:divisibility}.
\emph{Data} are $\mathfrak{d}=(c_0,c_+,c_3;h;O)$,
where $c_0,c_+\in\Z$, $c_3\in\Z+\frac12$, $h$ is an even positive integer with $4h\in\mathcal{D}(\mathfrak{n})$,
and $O\in\mathcal{T}$ satisfies $\supp O\subset Q_0$ for the open cell $Q_0$ with $\kappa=0$ of the cell structure
\[
\mathcal{Q}_{\mathfrak{d}}:=(\mathfrak{F}_{\mathrm{d}};c_0,c_+,c_3;h),
\]
and where moreover
\[
h>8r_O+8.
\]
Thus the walls of $\mathcal{Q}_{\mathfrak{d}}$ are the mirrors $(\ell_-,c_0+j\cdot h)$, $(\ell_+,c_++j\cdot h)$ and $(x_3,c_3+j\cdot h)$, $j\in\Z$,
and $|Q_0|=h^3/2$.
The \emph{folded array of the data} is the folded array of $O$ with respect to $\mathcal{Q}_{\mathfrak{d}}$ (Definition~\ref{def:arrays}),
\[
O_{\mathfrak{d}}:=\sum_{\kappa\in\Z^3}O^{[\kappa]},
\]
and $\widetilde{O}$ denotes the unfolded array of $O$ with respect to $\mathcal{Q}_{\mathfrak{d}}$.

For every index $i$ with $n_i\in2h\Z$ and $n_i\ge8$,
the structure $\mathcal{Q}_{\mathfrak{d}}$ cuts $C(n_i)$ and $D'(n_i)$, and $D(n_i)$ carries it (Theorem~\ref{thm:chessboard}(b));
the torus $D(n_i)$ then has $n_i^3/|Q_0|$ cells.
Every $i\ge i_0(\mathfrak{d})$ is such an index, because then $4h$ divides $n_i$ and $n_i\ge72$ (Definition~\ref{def:threshold});
so these assertions hold for all large $i$.
Let $i$ be an index for which $D(n_i)$ carries $\mathcal{Q}_{\mathfrak{d}}$, write $\mathbb{T}=(\Gamma,\chi):=D(n_i)$,
and let $t\ge0$.
For $G\in\mathcal{A}(\R^3)$ that embeds in $\mathbb{T}$ the \emph{perturbed functional} is
\[
\mu^t_i(G):=\frac{Z_{\mathbb{T}}\left[\exp(t\cdot O_{\mathfrak{d}})\cdot G\right]}{Z_{\mathbb{T}}\left[\exp(t\cdot O_{\mathfrak{d}})\right]}.
\]
Its denominator is positive:
on $\mathbb{T}$ one has $|O_{\mathfrak{d}}|\le\frac{n_i^3}{|Q_0|}\cdot|O|_1$ (Definition~\ref{def:arrays}),
and Lemma~\ref{lem:perturbed-denominators} together with Theorem~\ref{thm:partition-positivity} gives
\[
Z_{\mathbb{T}}\left[\exp(t\cdot O_{\mathfrak{d}})\right]\ge\exp\left(-t\cdot|O|_1\cdot\frac{n_i^3}{|Q_0|}\right)\cdot Z_{\mathbb{T}}[1]>0.
\]
On bosonic $G$ the functional $\mu^t_i$ is an expectation.
Indeed, the determinant $\det D^{\mathbb{T}}_A$ is non-negative (Proposition~\ref{prop:accretivity}),
so that
\[
\frac{\exp\left(t\cdot O_{\mathfrak{d}}(\widetilde{F})\right)\cdot\exp\left(-\frac12\|F\|^2\right)\cdot\det D^{\mathbb{T}}_A}{Z_{\mathbb{T}}\left[\exp(t\cdot O_{\mathfrak{d}})\right]}
\]
is the density of a probability measure on $\Conf(\Gamma)$ with respect to the measure of Definition~\ref{def:configurations};
the law of $\widetilde{F}$ under this probability measure is a Borel probability measure on $\Omega$, again denoted $\mu^t_i$,
and $\mu^t_i(G)=\int\mathrm{d}\mu^t_i(F)\,G(F)$ for every bosonic $G$ as above.

The data of this definition are called \emph{standard data}.
An element $g\in\Ocub$ permutes the coordinate axes and changes their signs;
the images under $g$ of the objects above make up the \emph{image data} $g\mathfrak{d}$ of Definition~\ref{def:image-data},
and they are called by the same names.
For $i$, $t$ and $G$ as above, $\mu^t_{i,g\mathfrak{d}}\left(G^g\right)$ is either $\mu^t_i(G)$ or its complex conjugate (Lemma~\ref{lem:data-transport}(iii)).
\end{definition}

\begin{definition}[Comparison states]\label{def:comparison-states}
Let $\mathfrak{n}$ satisfy \eqref{cond:divisibility}, and let $\mathfrak{d}$ be data.
Consider a subsequence $(i_m)_{m\ge1}$ of the indices and a sequence $(t_k)_{k\ge1}$ of positive numbers tending to $0$;
for every $k$ the law $\mu^{t_k}_{i_m}$ of Definition~\ref{def:data} is defined once $m$ is large.
If Borel probability measures $\nu_{t_k}$, $k\ge1$, and $\nu_0$ on $\Omega$ satisfy
\[
\mu^{t_k}_{i_m}\to\nu_{t_k}\ \text{weakly as }m\to\infty\ \text{for every }k,
\qquad
\nu_{t_k}\to\nu_0\ \text{weakly as }k\to\infty,
\]
then $\nu_0$ is called a \emph{comparison state} for $\mathfrak{d}$.
Given a further subsequence $(i'_m)_{m\ge1}$ of the indices, a comparison state $\nu_0$ for $\mathfrak{d}$ is \emph{subordinate} to $(i'_m)$
if it is obtained in this way from some subsequence $(i_m)$ of $(i'_m)$, with suitable numbers $t_k$ and measures $\nu_{t_k}$.
In every limit of this kind the volume limit $i\to\infty$ is taken first and the limit $t\to0+$ in the strength second;
where a regularisation parameter $\varepsilon$ occurs as well, the limit $\varepsilon\to0$ is taken last.
\end{definition}

\begin{definition}[Amplification factor]\label{def:amplification-factor}
For a sequence $\mathfrak{n}$ with \eqref{cond:divisibility} and data $\mathfrak{d}=(c_0,c_+,c_3;h;O)$,
for $k\in\N$, $t\ge0$ and every index $i$ with $n_i$ even and $n_i\ge8$, so that $\Delta_i$ is defined, put
\[
\gamma_{t,i}(k):=\exp\left(2t\cdot|O|_1\cdot k^3+|Q_0|\cdot k^3\cdot\Delta_i\right).
\]
By \eqref{cond:divisibility}, $8$ divides $n_i$ for all large $i$, and $n_i\to\infty$;
so these indices include all large $i$,
and $\Delta_i\to0$ as $i\to\infty$ (Theorem~\ref{thm:four-tori}).
Hence, for every fixed $k$,
\[
\lim_{t\to0+}\,\limsup_{i\to\infty}\,\gamma_{t,i}(k)=1.
\]
\end{definition}

\subsection{Slices and Fock spaces}\label{sec:slices-fock}

In this subsection the coordinate $x_1$ is singled out.
A point of $\Z^3$ is written $(j,y)$ with $j\in\Z$ and $y\in\Z^2$, where $y=(y_2,y_3)$ collects the coordinates $x_2,x_3$,
and $e_2,e_3$ also denote the unit vectors of $\Z^2$.

\begin{definition}[Slices]\label{def:slices}
\emph{The slice lattice.}
The slice lattice is $\Z^2$, with the coordinates $(x_2,x_3)$.
Its links are the segments $b=(y,\nu)=[y,y+e_\nu]$ with $y\in\Z^2$ and $\nu\in\{2,3\}$,
and its plaquettes are the closed unit squares $\sigma'=y+[0,1]^2$.
A function on the plaquettes is written as a function of the lower corner $y$ of $\sigma'$,
and a one-cochain of the slice is also read as an odd function of oriented pairs of neighbours, as in Definition~\ref{def:cochains}.
The coboundary $d_{\mathrm{s}}$ of the slice lattice acts on a function $a$ on the sites and on a one-cochain $A$ by
\[
(d_{\mathrm{s}}a)(y,\nu)=a(y+e_\nu)-a(y),
\qquad
(d_{\mathrm{s}}A)(y)=A(y,2)+A(y+e_2,3)-A(y+e_3,2)-A(y,3).
\]

\emph{Slice variables.}
Let $A$ be a real one-cochain on $\Z^3$ and $j\in\Z$.
The links of the plane $\{x_1=j\}$ carry the one-cochain $A_j$ of the slice,
and the links that join the plane $\{x_1=j\}$ to the plane $\{x_1=j+1\}$ carry the function $a_j$ on $\Z^2$:
\[
A_j(y,\nu):=A_\nu(j,y)\quad(\nu\in\{2,3\}),
\qquad
a_j(y):=A_1(j,y).
\]
For a real two-cochain $F$ on $\Z^3$, the electric part $F^{\mathrm{el}}_j$, a one-cochain of the slice,
and the magnetic part $F^{\mathrm{mag}}_j$, a function on the plaquettes of the slice, are
\[
F^{\mathrm{el}}_j(y,\nu):=F_{1\nu}(j,y)\quad(\nu\in\{2,3\}),
\qquad
F^{\mathrm{mag}}_j(y):=F_{23}(j,y);
\]
where the two-cochain has to be named, they are written $F^{\mathrm{el}}_j(F)$ and $F^{\mathrm{mag}}_j(F)$.
The formulas for $d$ in Definition~\ref{def:cochains} show that $F=dA$ has the parts
\[
F^{\mathrm{el}}_j=A_{j+1}-A_j-d_{\mathrm{s}}a_j,
\qquad
F^{\mathrm{mag}}_j=d_{\mathrm{s}}A_j,
\]
and that every two-cochain $F$ satisfies
$(dF)(j,y)=F^{\mathrm{mag}}_{j+1}(y)-F^{\mathrm{mag}}_j(y)-(d_{\mathrm{s}}F^{\mathrm{el}}_j)(y)$.
Thus $dF=0$ if and only if $F^{\mathrm{mag}}_{j+1}-F^{\mathrm{mag}}_j=d_{\mathrm{s}}F^{\mathrm{el}}_j$ for every $j\in\Z$.

\emph{Reflections of the slice variables.}
Let $j_0\in\Z$ and $c:=j_0+\frac12$, so that $(x_1,c)$ is a link plane.
Let further $\nu\in\{2,3\}$, $c'\in\frac12\Z$, and let $\theta_{\mathrm{s}}$ be the reflection of $\R^2$ in the line $\{x_\nu=c'\}$;
it permutes the sites, the links and the plaquettes of $\Z^2$.
For a one-cochain $u$ of the slice, $\theta_{\mathrm{s}}^*u$ is the pull-back,
$(\theta_{\mathrm{s}}^*u)(y,y')=u(\theta_{\mathrm{s}}y,\theta_{\mathrm{s}}y')$,
and a function on the sites or on the plaquettes is composed with $\theta_{\mathrm{s}}$ through the action of $\theta_{\mathrm{s}}$ on sites or plaquettes.
For an isometry $\theta$ of $\R^3$ with $\theta\Z^3=\Z^3$,
$(\theta^*A)_j$ and $(\theta^*a)_j$ denote the slice variables of the one-cochain $\theta^*A$,
and $F^{\mathrm{el}}_j(\theta^*F)$, $F^{\mathrm{mag}}_j(\theta^*F)$ the parts of the two-cochain $\theta^*F$ (pull-backs as in Definition~\ref{def:cochains}).
The reflection $\theta_{x_1,c}$ maps $(j,y)$ to $(2j_0+1-j,y)$;
it reverses the orientation of the links parallel to $e_1$ and of the plaquettes $(x;1\mu)$, $\mu\in\{2,3\}$, and preserves that of all other links and plaquettes.
The reflection $\theta_{x_\nu,c'}$ maps $(j,y)$ to $(j,\theta_{\mathrm{s}}y)$;
it reverses the orientation of the links parallel to $e_\nu$ and of the plaquettes $(x;1\nu)$ and $(x;23)$, and preserves that of all other links and plaquettes.
It follows that, for every $j\in\Z$,
\begin{equation}\label{eq:setting-slice-reflections}
\begin{aligned}
(\theta_{x_1,c}^*A)_j&=A_{2j_0+1-j},
&\qquad
(\theta_{x_1,c}^*a)_j&=-a_{2j_0-j},\\
F^{\mathrm{el}}_j(\theta_{x_1,c}^*F)&=-F^{\mathrm{el}}_{2j_0-j}(F),
&
F^{\mathrm{mag}}_j(\theta_{x_1,c}^*F)&=F^{\mathrm{mag}}_{2j_0+1-j}(F),\\
(\theta_{x_\nu,c'}^*A)_j&=\theta_{\mathrm{s}}^*A_j,
&
(\theta_{x_\nu,c'}^*a)_j&=a_j\circ\theta_{\mathrm{s}},\\
F^{\mathrm{el}}_j(\theta_{x_\nu,c'}^*F)&=\theta_{\mathrm{s}}^*F^{\mathrm{el}}_j(F),
&
F^{\mathrm{mag}}_j(\theta_{x_\nu,c'}^*F)&=-F^{\mathrm{mag}}_j(F)\circ\theta_{\mathrm{s}}.
\end{aligned}
\end{equation}

\emph{Rectangular cross-sections.}
Let $n_2,n_3\ge4$ be integers.
The rectangular cross-section with the sides $n_2,n_3$ is the cross-section
$\Xi=(n_2\Z e_2+n_3\Z e_3;\chi_\parallel)$ for the family $x_1$ (Definition~\ref{def:slabs})
with $\chi_\parallel(n_2e_2)=\chi_\parallel(n_3e_3)=-1$.
Its slice torus is $\Lambda_{\mathrm{s}}:=\Z^2/(n_2\Z\times n_3\Z)$.
It has $|\Lambda_{\mathrm{s}}|=n_2n_3$ sites,
the set $\Lambda^{(1)}_{\mathrm{s}}$ of its $2|\Lambda_{\mathrm{s}}|$ links
and the set $\Lambda^{(2)}_{\mathrm{s}}$ of its $|\Lambda_{\mathrm{s}}|$ plaquettes,
and $d_{\mathrm{s}}$ acts on its cochains by the formulas above.
The subgroup
\[
\Gamma_{\mathrm{s}}:=\left\{(2\pi/q)\cdot k:\ k\in\Z^{\Lambda^{(1)}_{\mathrm{s}}},\ d_{\mathrm{s}}k=0\right\}
\]
of $\R^{\Lambda^{(1)}_{\mathrm{s}}}$ consists of $2\pi/q$ times the closed integer one-cochains of $\Lambda_{\mathrm{s}}$,
and the quotient $\R^{\Lambda^{(1)}_{\mathrm{s}}}/\Gamma_{\mathrm{s}}$ carries the measure induced by Lebesgue measure.
A plaquette $\sigma_0\in\Lambda^{(2)}_{\mathrm{s}}$ is fixed;
it is the position of the flux string: in the slice operator of Definition~\ref{def:slice-operators} the flux $(2\pi/q)\cdot k$ of the sector $k$ is placed on $\sigma_0$.
Finally, for $a>0$ and $s\in\R$,
\[
g_a(s):=(2\pi a)^{-\frac12}\cdot\exp\left(-\frac{s^2}{2a}\right).
\]

\emph{Hulls.}
For a non-empty finite set $\Lambda$ of sites of $\Z^2$ let
$y^-_\nu:=\min_{y\in\Lambda}y_\nu$ and $y^+_\nu:=\max_{y\in\Lambda}y_\nu$ for $\nu=2,3$.
The hull
\[
\hull(\Lambda):=[y^-_2,y^+_2]\times[y^-_3,y^+_3]
\]
is the smallest product of two closed intervals with integer end points that contains $\Lambda$.
For a finite set $S$ of links of $\Z^2$ whose end points lie in $\Lambda$, $\hull(S,\Lambda):=\hull(\Lambda)$.
The hull \emph{embeds} in the rectangular cross-section $\Xi$ if the sets
$\hull(\Lambda)+[-1,1]^2+g$, $g\in n_2\Z\times n_3\Z$, are pairwise disjoint,
and it \emph{contains} the plaquette $\sigma_0$ of $\Lambda_{\mathrm{s}}$ if a representative of $\sigma_0$ in $\Z^2$,
which is a closed unit square, is a subset of $\hull(\Lambda)$.
\end{definition}

\begin{definition}[Modes and Fock spaces]\label{def:modes}
Let $\rho\in\R$ and $w_0,\dots,w_3>0$ be the Wilson parameter and the hopping weights of the chamber operator (Definition~\ref{def:chamber-operator}).
Bonds and their types are those of Definition~\ref{def:fermion-sites},
and $A$ is a real one-cochain on $\Z^3$.

\emph{Modes at a mirror.}
Let $\Pi=(\ell,c)$ be a mirror, $n:=\hat n_\ell$ and $\theta:=\theta_\Pi$.
Just as $\theta_\Pi$ is also written $\theta_{\ell,c}$, the reflection $\Theta_\Pi$ of Definition~\ref{def:grassmann-reflections} is also written $\Theta_{\ell,c}$.
A bond \emph{crosses} $\Pi$ if its end points lie on different sides of $\Pi$;
by Proposition~\ref{prop:geometry}\eqref{it:geometry-mirrors} it is then of the form $\{x,\theta x\}$,
and the labels are chosen so that $\ell(x)<c$: the site $x$ is the lower end of the bond and $\theta x$ its upper end.
Let $\{x,\theta x\}$ be a bond that crosses $\Pi$, let $j\in\{0,1,2,3\}$ be its type and $k\in\{1,2\}$.
The bond carries four groups of modes.
For each group and each $k$, the following list gives an element $\zeta$ of the Grassmann algebra $\Gr(\mathbb{X})$ (Definition~\ref{def:grassmann}),
a coefficient $\omega$ and a phase $u$; the vectors $e^\pm_{n,k}$ are those of Definition~\ref{def:spin-matrices}:
\begin{align*}
&\text{group 1:} &\zeta&=(e^-_{n,k})^*\psi(\theta x), &\omega&=\tfrac{w_j}{2}(1+\rho), &u&=U_A(x,\theta x),\\
&\text{group 2:} &\zeta&=(e^+_{n,k})^*\psi(\theta x), &\omega&=\tfrac{w_j}{2}(1-\rho), &u&=U_A(x,\theta x),\\
&\text{group 3:} &\zeta&=\bar\psi(\theta x)\,e^+_{n,k}, &\omega&=\tfrac{w_j}{2}(1+\rho), &u&=\overline{U_A(x,\theta x)},\\
&\text{group 4:} &\zeta&=\bar\psi(\theta x)\,e^-_{n,k}, &\omega&=\tfrac{w_j}{2}(1-\rho), &u&=\overline{U_A(x,\theta x)}.
\end{align*}
The set $\idx(\Pi)$ of the modes at $\Pi$ consists of the triples formed by a bond that crosses $\Pi$, a group and an index $k\in\{1,2\}$;
there are eight modes for each such bond.
For a mode $i\in\idx(\Pi)$, its element, its coefficient and its phase, as given by the list, are written $\zeta_i$, $\omega_i$ and $u_i$.
The weights being positive, $\omega_i\ge0$ for every $i\in\idx(\Pi)$ if $-1\le\rho\le1$,
and $\omega_i>0$ for every $i\in\idx(\Pi)$ if $-1<\rho<1$, in particular if $0<\rho<1$.
The four vectors $e^\pm_{n,k}$ form an orthonormal basis of $\C^4$, and a site is an end of at most one bond that crosses $\Pi$.
Therefore the $\zeta_i$, $i\in\idx(\Pi)$, form a basis of the linear span of the generators $\psi_\alpha(x')$, $\bar\psi_\alpha(x')$
at the upper ends $x'$ of the bonds that cross $\Pi$;
since $\Theta_\Pi$ maps the span of the generators at $\theta x$ antilinearly and bijectively onto the span of those at $x$
(Definition~\ref{def:grassmann-reflections}),
the $\Theta_\Pi(\zeta_i)$ form a basis of the span of the generators at the lower ends.

\emph{Translations and orders.}
For $v\in\Z^3$ and a mode $i$ at $\Pi$, the triple $i+v$ formed by the translated bond, the same group and the same $k$ is a mode at $(\ell,c+\ell(v))$.
The type of a bond and the vectors $e^\pm_{n,k}$ do not change under translations,
and $\theta_{\ell,c+\ell(v)}(x+v)=\theta_\Pi(x)+v$ by Definition~\ref{def:mirrors}\eqref{it:mirror-conjugation};
hence $\zeta_{i+v}$ arises from $\zeta_i$ by the translation of the fields by $v$, and $\omega_{i+v}=\omega_i$.
Each set $\idx(\Pi)$ carries a total order, fixed once and for all, and the orders are compatible with the normal translations:
$i\mapsto i+n_\ell$ is an order preserving bijection of $\idx((\ell,c))$ onto $\idx((\ell,c+|n_\ell|^2))$.
(It suffices to choose the orders for the positions $c$ in a set of representatives of $\Adm(\ell)$ modulo $|n_\ell|^2\Z$ and to transport them.)

\emph{Cross-sections.}
Let $\Xi=(\Gamma_\parallel,\chi_\parallel)$ be a cross-section for $\ell$ (Definition~\ref{def:slabs}).
The Grassmann algebra of $\Xi$ is generated by the fields at the sites taken modulo $\Gamma_\parallel$, subject to
$\psi(x+g)=\chi_\parallel(g)\psi(x)$ and $\bar\psi(x+g)=\chi_\parallel(g)\bar\psi(x)$ for $g\in\Gamma_\parallel$, as for tori in Definition~\ref{def:tori};
in it, $\zeta_{i+g}=\chi_\parallel(g)\cdot\zeta_i$ for $i\in\idx(\Pi)$ and $g\in\Gamma_\parallel$.
Since $\theta_\Pi(x+g)=\theta_\Pi(x)+g$ for $x\in\R^3$ and $g\in\Gamma_\parallel$ (Definition~\ref{def:mirrors}\eqref{it:mirror-conjugation}, as $\ell(g)=0$) and $\chi_\parallel$ is real,
the formulas \eqref{eq:setting-reflection-fields} show that $\Theta_\Pi$ maps these relations to linear combinations of them;
so $\Theta_\Pi$ acts on the Grassmann algebra of $\Xi$, again antilinearly and reversing the order of products.
The group $\Gamma_\parallel$, which lies in $\{\ell=0\}$, acts on $\idx(\Pi)$ by $i\mapsto i+g$; let
\[
\idx_\Xi(\Pi):=\idx(\Pi)/\Gamma_\parallel ,
\]
a finite set: the bonds that cross $\Pi$ fall into finitely many classes under the translations by $\Z^3\cap\{\ell=0\}$,
and $\Gamma_\parallel$ has finite index in this group of rank two.
A set of representatives of its elements is fixed, and $\idx_\Xi(\Pi)$ is ordered as the representatives are ordered in $\idx(\Pi)$.
The element and the coefficient of a class are those of its representative, and so is its phase,
which for $\Gamma_\parallel$-periodic $A$ does not depend on the choice of the representative.
The representatives for $(\ell,c+|n_\ell|^2)$ are the translates by $n_\ell$ of those for $(\ell,c)$;
so the orders of the sets $\idx_\Xi(\Pi)$ are again compatible with $\tau_{n_\ell}$.
Translations along the mirror behave differently:
for $v\in\Z^3\cap\{\ell=0\}$ with $v\notin\Gamma_\parallel$, the map $i\mapsto i+v$ permutes the finite set $\idx_\Xi(\Pi)$ without fixed points,
and therefore it preserves no total order of $\idx_\Xi(\Pi)$.
Accordingly, let $v\in\Z^3$ with $\ell(v)=0$.
The unitary $V_v$ of Corollary~\ref{cor:slab-admissible} maps the space $\mathcal{H}(c)$ of that corollary onto itself;
it acts on the bosonic mirror variables through the pull-back $\tau_v^{*}$,
and it moves the value at a set $I=\{i_1<\dots<i_d\}\subset\idx_\Xi(\Pi)$ of modes to the set $I+v=\{i_1+v,\dots,i_d+v\}$,
multiplied by the sign of the permutation that brings $i_1+v,\dots,i_d+v$ into increasing order
and by the product of the numbers $\chi_\parallel(g)$, where $g\in\Gamma_\parallel$ ranges over the $d$ vectors
by which the translates by $v$ of the representatives of $i_1,\dots,i_d$ differ from the representatives of $i_1+v,\dots,i_d+v$.

\emph{Fock spaces.}
For a finite set $\mathcal{J}$, the Fock space $\mathcal{F}(\mathcal{J})$ is the Hilbert space with the orthonormal basis $(e_I)_{I\subset\mathcal{J}}$.
The fermion parity is the unitary involution $U_{\mathrm{f}}$ with $U_{\mathrm{f}}e_I:=(-1)^{\#|I|}e_I$;
an operator $\Upsilon\in\End\mathcal{F}(\mathcal{J})$ is \emph{even} if it commutes with $U_{\mathrm{f}}$,
and $\Upsilon_{IJ}:=\langle e_I,\Upsilon e_J\rangle$ are its matrix entries.

\emph{Ordered monomials; matrices and elements.}
Let $\mathcal{J}$ be a finite subset of $\idx(\Pi)$ or of $\idx_\Xi(\Pi)$, with the induced order,
and let $I=\{i_1<\dots<i_d\}\subset\mathcal{J}$.
Put $\omega_I:=\prod_{i\in I}\omega_i$ and $u_I:=\prod_{i\in I}u_i$,
and let $\omega:=\diag(\omega_I)_{I\subset\mathcal{J}}$ be the diagonal operator of the coefficients on $\mathcal{F}(\mathcal{J})$.
For $\Upsilon\in\End\mathcal{F}(\mathcal{J})$,
\begin{equation}\label{eq:setting-ordered-monomials}
\begin{aligned}
\zeta_I&:=\zeta_{i_1}\zeta_{i_2}\cdots\zeta_{i_d},
\qquad
\Theta_\Pi(\zeta_I)=\Theta_\Pi(\zeta_{i_d})\cdots\Theta_\Pi(\zeta_{i_2})\,\Theta_\Pi(\zeta_{i_1}),\\
\Ins_\Pi[\Upsilon]&:=\sum_{I,J\subset\mathcal{J}}\Upsilon_{IJ}\cdot\Theta_\Pi(\zeta_J)\,\zeta_I ,
\end{aligned}
\end{equation}
with $\zeta_\emptyset=\omega_\emptyset=u_\emptyset=1$;
the second formula holds because $\Theta_\Pi$ reverses products.
In $\Ins_\Pi[\Upsilon]$ the row index $I$ belongs to the upper side of $\Pi$ and the column index $J$ to the lower side.
The elements $\Theta_\Pi(\zeta_i)$ are odd and anticommute; reversing the order of the factors gives
$\Theta_\Pi(\zeta_{i_1})\cdots\Theta_\Pi(\zeta_{i_d})=(-1)^{\binom{\#|I|}{2}}\cdot\Theta_\Pi(\zeta_I)$,
so that $(-1)^{\binom{\#|I|}{2}}$ is the sign of the reversal of $\#|I|$ odd factors.

\emph{Localisation.}
Let $\mathcal{J}$ again be a finite subset of $\idx(\Pi)$ or of $\idx_\Xi(\Pi)$ with the induced order,
and let $\mathcal{J}'\subset\mathcal{J}$, ordered as a subset of $\mathcal{J}$.
Every subset of $\mathcal{J}$ is, in exactly one way, of the form $I\cup R$ with $I\subset\mathcal{J}'$ and $R\subset\mathcal{J}\setminus\mathcal{J}'$.
For $\Upsilon\in\End\mathcal{F}(\mathcal{J}')$, the operator $\iota(\Upsilon)\in\End\mathcal{F}(\mathcal{J})$ is defined by its matrix entries
\[
\iota(\Upsilon)_{I\cup R,\,J\cup R'}
:=\delta_{RR'}\cdot(-1)^{\#\{(i,r)\in I\times R:\ i>r\}}\cdot(-1)^{\#\{(i,r)\in J\times R:\ i>r\}}\cdot\Upsilon_{IJ}
\]
for $I,J\subset\mathcal{J}'$ and $R,R'\subset\mathcal{J}\setminus\mathcal{J}'$.
By Proposition~\ref{prop:localisation}\eqref{it:localisation-homomorphism},
$\iota:\End\mathcal{F}(\mathcal{J}')\to\End\mathcal{F}(\mathcal{J})$ is a unital injective $*$-homomorphism, hence isometric.

\emph{Coordinate planes.}
Let $j\in\Z$.
The bonds that cross the link plane $(x_1,j+\frac12)$ are of type $3$ (Proposition~\ref{prop:geometry}\eqref{it:geometry-projection}),
and there are eight of them over each $y\in\Z^2$, that is, over the link from $(j,y)$ to $(j+1,y)$.
The lower end $x$ of such a bond has $\pi(x)=(j,y)$ and its upper end has $\pi(\theta_{x_1,j+1/2}x)=(j+1,y)$,
so that $U_A(x,\theta_{x_1,j+1/2}x)=\exp(iq\,a_j(y))$.
For a finite set $\Lambda\subset\Z^2$, $\mathcal{J}^{\mathrm{l}}(\Lambda)$ is the set of the $64\cdot\#|\Lambda|$ modes of the bonds over the sites of $\Lambda$,
ordered as a subset of $\idx((x_1,j+\frac12))$.
A mode $i$ of a bond over $y\in\Lambda$ has the charge $\mathfrak{q}(i):=\delta_y$ if it belongs to the group 1 or 2,
and $\mathfrak{q}(i):=-\delta_y$ if it belongs to the group 3 or 4, where $\delta_y\in\Z^\Lambda$ is the indicator of $y$;
for $I\subset\mathcal{J}^{\mathrm{l}}(\Lambda)$ let $\mathfrak{q}(I):=\sum_{i\in I}\mathfrak{q}(i)\in\Z^\Lambda$.
Then $u_I=\exp(iq\langle a_j,\mathfrak{q}(I)\rangle)$, and for $a\in\R^\Lambda$ the diagonal unitary
\[
u(a):=\diag\left(\exp(iq\langle a,\mathfrak{q}(I)\rangle)\right)_{I\subset\mathcal{J}^{\mathrm{l}}(\Lambda)}
\quad\text{on }\mathcal{F}(\mathcal{J}^{\mathrm{l}}(\Lambda))
\]
satisfies $\diag(u_I)_{I\subset\mathcal{J}^{\mathrm{l}}(\Lambda)}=u(a_j)$, with $a_j$ restricted to $\Lambda$.
The bonds that cross the site plane $(x_1,j)$ are of type $0$ (Proposition~\ref{prop:geometry}\eqref{it:geometry-projection}),
and they carry no phase, since $\pi(x)=\pi(\theta_{x_1,j}x)$ for such a bond $\{x,\theta_{x_1,j}x\}$, $x_1<j$.
This photon site is $(j,y)$ for some $y\in\Z^2$,
because $x_1$ and $2j-x_1$ both differ from the integer $\pi(x)_1$ by less than $\frac12$ (Definition~\ref{def:fermion-sites});
the bond is then said to lie over $y$.
There are eight such bonds over each $y\in\Z^2$:
by Proposition~\ref{prop:geometry}\eqref{it:geometry-midpoint}, the bond of type $0$ at the site $(j,y)+\vec g\,v_0$, $\vec g\in\Ocub$,
crosses $(x_1,j)$ if and only if the wall $(j,y)+\vec g\,\Pi_0$ of the chamber of this site is the plane $\{x_1=j\}$, that is, $\vec g\,e_3=\pm e_1$;
this holds for $16$ of the $48$ matrices $\vec g$, and these $16$ sites are the end points of eight bonds.
For a finite set $\Lambda\subset\Z^2$, $\mathcal{J}^{\mathrm{s}}(\Lambda)$ is the set of the $64\cdot\#|\Lambda|$ modes over $\Lambda$ at $(x_1,j)$, ordered as a subset of $\idx((x_1,j))$.

For a rectangular cross-section $\Xi$ (Definition~\ref{def:slices}) the sets of the modes over $\Lambda_{\mathrm{s}}$ are
\[
\mathcal{J}^{\mathrm{l}}(\Lambda_{\mathrm{s}}):=\idx_\Xi((x_1,j+\tfrac12)),
\qquad
\mathcal{J}^{\mathrm{s}}(\Lambda_{\mathrm{s}}):=\idx_\Xi((x_1,j)),
\]
the charges $\mathfrak{q}(i)$ and $\mathfrak{q}(I)$ are defined as above, with values in $\Z^{\Lambda_{\mathrm{s}}}$, and
\[
\mathcal{F}:=\mathcal{F}(\mathcal{J}^{\mathrm{l}}(\Lambda_{\mathrm{s}})),
\qquad
\mathcal{F}':=\mathcal{F}(\mathcal{J}^{\mathrm{s}}(\Lambda_{\mathrm{s}})).
\]
The diagonal operators of the coefficients on $\mathcal{F}$ and on $\mathcal{F}'$ are $\omega$ and $\omega'$,
and for $a\in\R^{\Lambda_{\mathrm{s}}}$ the diagonal unitary $u(a)$ on $\mathcal{F}$ is given by the same formula as above, with $\Lambda_{\mathrm{s}}$ in the place of $\Lambda$.
If $A$ is $\Gamma_\parallel$-periodic, then $a_j$ is a function on $\Lambda_{\mathrm{s}}$,
and, the phase of a class being that of its representative, $u_I=\exp\left(iq\langle a_j,\mathfrak{q}(I)\rangle\right)$ for $I\subset\mathcal{J}^{\mathrm{l}}(\Lambda_{\mathrm{s}})$, that is, $\diag(u_I)_{I\subset\mathcal{J}^{\mathrm{l}}(\Lambda_{\mathrm{s}})}=u(a_j)$.

\emph{Signed relabelling.}
Let $\Xi$ be a rectangular cross-section and let $\Lambda\subset\Z^2$ be a finite set on which the reduction map $\Z^2\to\Lambda_{\mathrm{s}}$ is injective;
this holds if $\hull(\Lambda)$ embeds in $\Xi$.
Every mode $i\in\mathcal{J}^{\mathrm{l}}(\Lambda)$ is $i=r(i)+g_i$ with a unique representative $r(i)$ of a class of $\mathcal{J}^{\mathrm{l}}(\Lambda_{\mathrm{s}})$
and a unique $g_i\in\Gamma_\parallel$, and $r$ is injective on $\mathcal{J}^{\mathrm{l}}(\Lambda)$.
For $I=\{i_1<\dots<i_d\}\subset\mathcal{J}^{\mathrm{l}}(\Lambda)$ let
\[
s(I):=\chi_\parallel(g_{i_1})\cdots\chi_\parallel(g_{i_d})\cdot\operatorname{sgn}\left(\text{the permutation that sorts }(r(i_1),\dots,r(i_d))\right),
\]
the sorting being done in the order of $\mathcal{J}^{\mathrm{l}}(\Lambda_{\mathrm{s}})$,
and let $U_\Lambda:\mathcal{F}(\mathcal{J}^{\mathrm{l}}(\Lambda))\to\mathcal{F}(r(\mathcal{J}^{\mathrm{l}}(\Lambda)))$ be the unitary with $U_\Lambda e_I:=s(I)\cdot e_{r(I)}$.
In the Grassmann algebra of $\Xi$ one has $\zeta_i=\chi_\parallel(g_i)\cdot\zeta_{r(i)}$, hence $\zeta_I=s(I)\cdot\zeta_{r(I)}$;
thus $U_\Lambda$ is the unitary of Proposition~\ref{prop:localisation}\eqref{it:mode-relabelling} for the relabelling $r$.
On $\Xi$, the localisation of $\Upsilon\in\End\mathcal{F}(\mathcal{J}^{\mathrm{l}}(\Lambda))$ is the operator
\[
\iota(\Upsilon):=\iota\left(U_\Lambda\Upsilon U_\Lambda^*\right)\in\End\mathcal{F},
\]
where on the right $\iota$ is formed with the subset $r(\mathcal{J}^{\mathrm{l}}(\Lambda))$ of $\mathcal{J}^{\mathrm{l}}(\Lambda_{\mathrm{s}})$.
If $\mathcal{J}^{\mathrm{l}}(\Lambda)$ consists of representatives, then $U_\Lambda$ is the identity.
Conjugation by $U_\Lambda$ preserves norms, adjoints and products.
Moreover $\omega_{r(i)}=\omega_i$, and $\mathfrak{q}(r(i))$ is the image of $\mathfrak{q}(i)$ under the reduction map.
Write $[y]\in\Lambda_{\mathrm{s}}$ for the class of $y\in\Z^2$; then, for every $a\in\R^{\Lambda_{\mathrm{s}}}$,
\[
U_\Lambda\,\omega\,U_\Lambda^*=\omega,
\qquad
U_\Lambda\,u\left(a([\,\cdot\,])\right)U_\Lambda^*=u(a)
\quad\text{on }\mathcal{F}(r(\mathcal{J}^{\mathrm{l}}(\Lambda))),
\]
where on the left $\omega$ and $u(a([\,\cdot\,]))$ act on $\mathcal{F}(\mathcal{J}^{\mathrm{l}}(\Lambda))$, $a([\,\cdot\,])$ being the function $y\mapsto a([y])$ on $\Lambda$,
and on the right $\omega$ and $u(a)$ are restricted to the span of the $e_I$, $I\subset r(\mathcal{J}^{\mathrm{l}}(\Lambda))$.

\emph{Half slab.}
Let $\Xi$ be a rectangular cross-section, let $j\in\Z$, and let $A$ be $\Gamma_\parallel$-periodic, so that $A_j$ is a one-cochain of $\Lambda_{\mathrm{s}}$.
The half slab $\mathbb{X}_{\mathrm{h}}$ is the set of the fermion sites $x$ with $j<x_1<j+\frac12$, taken modulo $\Gamma_\parallel=n_2\Z e_2+n_3\Z e_3$;
its fields satisfy the relations of the Grassmann algebra of $\Xi$,
and $\int_{\mathbb{X}_{\mathrm{h}}}$ is the Berezin integral of Definition~\ref{def:grassmann} over one representative of each class.
Let $S_{\mathrm{h}}$ be the sum of the terms of $\bar\psi D_A\psi$ (Definition~\ref{def:chamber-operator}) whose sites all lie in $\mathbb{X}_{\mathrm{h}}$.
A bond with both ends in $\mathbb{X}_{\mathrm{h}}$ that carries a phase is of type $3$
and lies over a link of the plane $\{x_1=j\}$ (Proposition~\ref{prop:geometry}\eqref{it:geometry-projection});
so $S_{\mathrm{h}}$ depends on $A$ only through the phases $\exp(iqA_j(b))$, $b\in\Lambda^{(1)}_{\mathrm{s}}$,
and it is written $S_{\mathrm{h}}(A_j)$, or $S_{\mathrm{h}}(A)$ with an arbitrary $A\in\R^{\Lambda^{(1)}_{\mathrm{s}}}$ in the place of $A_j$.
Let $\zeta'_{i'}$, $i'\in\mathcal{J}^{\mathrm{s}}(\Lambda_{\mathrm{s}})$, be the modes of the site plane $(x_1,j)$,
$\zeta'_{I'}$ their ordered products as in \eqref{eq:setting-ordered-monomials},
and $\zeta_i$, $i\in\mathcal{J}^{\mathrm{l}}(\Lambda_{\mathrm{s}})$, the modes of the link plane $(x_1,j+\frac12)$;
the elements $\zeta'_{i'}$ and $\Theta_{x_1,j+1/2}(\zeta_i)$ belong to the Grassmann algebra of $\mathbb{X}_{\mathrm{h}}$.
For $A\in\R^{\Lambda^{(1)}_{\mathrm{s}}}$, the half-slab matrix $\mathfrak{N}(A):\mathcal{F}'\to\mathcal{F}$ has the matrix entries
\[
\mathfrak{N}(A)_{I,I'}:=\int_{\mathbb{X}_{\mathrm{h}}}\zeta'_{I'}\exp\left(-S_{\mathrm{h}}(A)\right)\Theta_{x_1,j+1/2}(\zeta_I),
\qquad I\subset\mathcal{J}^{\mathrm{l}}(\Lambda_{\mathrm{s}}),\ I'\subset\mathcal{J}^{\mathrm{s}}(\Lambda_{\mathrm{s}}).
\]
If $-1\le\rho\le1$, the diagonal operators $\omega$ and $\omega'$ have non-negative entries;
let $\omega^{\frac12}$ and $\omega'^{\frac12}$ be the diagonal operators with the non-negative square roots of these entries, and put
\[
\widehat{\mathfrak{N}}(A):=\omega^{\frac12}\,\mathfrak{N}(A)\,\omega'^{\frac12}:\ \mathcal{F}'\to\mathcal{F}.
\]

\emph{Independence of $j$.}
For every $y\in\Z^2$, the translation $\tau_{e_1}$ carries the bonds over $y$ that cross $(x_1,j)$ and $(x_1,j+\frac12)$
onto those over $y$ that cross $(x_1,j+1)$ and $(x_1,j+\frac32)$.
So $i\mapsto i+e_1$ maps $\mathcal{J}^{\mathrm{s}}(\Lambda)$ and $\mathcal{J}^{\mathrm{l}}(\Lambda)$, formed at $j$, onto the same sets formed at $j+1$;
it preserves the orders (they are compatible with $\tau_{n_\ell}$, and $n_{x_1}=e_1$), the groups, the coefficients and the charges.
As the representatives for $(\ell,c+|n_\ell|^2)$ are the translates by $n_\ell$ of those for $(\ell,c)$,
the representatives of the classes at $(x_1,j+1)$ and at $(x_1,j+\frac32)$ are the translates by $e_1$ of those at $(x_1,j)$ and at $(x_1,j+\frac12)$.
Hence $i\mapsto i+e_1$ also maps $\mathcal{J}^{\mathrm{s}}(\Lambda_{\mathrm{s}})$ and $\mathcal{J}^{\mathrm{l}}(\Lambda_{\mathrm{s}})$ onto the same sets formed at $j+1$, with the same properties,
and in the signed relabelling $r(i+e_1)=r(i)+e_1$, $g_{i+e_1}=g_i$ and $s(I+e_1)=s(I)$.
The translation $\tau_{e_1}$ carries the half slab of $j$ onto that of $j+1$, and $S_{\mathrm{h}}(A)$, formed at $j$, to $S_{\mathrm{h}}(A)$, formed at $j+1$, for every $A\in\R^{\Lambda^{(1)}_{\mathrm{s}}}$,
because the hopping matrices $w(x,x')$ of Definition~\ref{def:chamber-operator} do not change under translations.
It carries $\zeta'_{I'}$ to $\zeta'_{I'+e_1}$ and, by Definition~\ref{def:mirrors}\eqref{it:mirror-conjugation}, $\Theta_{x_1,j+1/2}(\zeta_I)$ to $\Theta_{x_1,j+3/2}(\zeta_{I+e_1})$,
and the Berezin integrals over the two half slabs correspond to each other under it.
Identifying $e_I$ with $e_{I+e_1}$, therefore, the Fock spaces $\mathcal{F}(\mathcal{J}^{\mathrm{l}}(\Lambda))$, $\mathcal{F}(\mathcal{J}^{\mathrm{s}}(\Lambda))$, $\mathcal{F}$ and $\mathcal{F}'$,
the operators $\omega$, $\omega'$ and $u(a)$ on them, the charges, the unitary $U_\Lambda$, the localisation $\iota$ on $\Xi$,
and the operators $\mathfrak{N}(A)$ and $\widehat{\mathfrak{N}}(A)$ do not depend on $j$.
\end{definition}

\begin{definition}[Operators of a cross-section]\label{def:slice-operators}
Let $\Xi$ be a rectangular cross-section,
with $\Lambda_{\mathrm{s}}$, $\Gamma_{\mathrm{s}}$, $\sigma_0$ and $g_a$ as in Definition~\ref{def:slices}
and $\mathcal{F}$, $\mathcal{F}'$ and $u(a)$ as in Definition~\ref{def:modes}.
The Hilbert spaces of the link plane and of the site plane are
\[
\mathcal{H}^{\mathrm{l}}:=\bigoplus_{k\in\Z}L^2\left(\R^{\Lambda^{(1)}_{\mathrm{s}}}/\Gamma_{\mathrm{s}};\mathcal{F}\right),
\qquad
\mathcal{H}^{\mathrm{s}}:=\bigoplus_{k\in\Z}L^2\left(\R^{\Lambda^{(1)}_{\mathrm{s}}}/\Gamma_{\mathrm{s}};\mathcal{F}'\right),
\]
the $L^2$ spaces being formed with the measure of Definition~\ref{def:slices}.
The summation index $k$ is the flux index,
and a vector of a summand is written as a $\Gamma_{\mathrm{s}}$-periodic function $\Phi(A)$ of $A\in\R^{\Lambda^{(1)}_{\mathrm{s}}}$.
All operators below are diagonal in $k$.
An operator on $L^2(\R^{\Lambda^{(1)}_{\mathrm{s}}}/\Gamma_{\mathrm{s}})$ acts on
$L^2(\R^{\Lambda^{(1)}_{\mathrm{s}}}/\Gamma_{\mathrm{s}};\mathcal{F})=L^2(\R^{\Lambda^{(1)}_{\mathrm{s}}}/\Gamma_{\mathrm{s}})\otimes\mathcal{F}$
through the first factor, and in the same way on every summand.

\emph{The electric field.}
For $b\in\Lambda^{(1)}_{\mathrm{s}}$ let $E_b:=-i\,\partial/\partial A(b)$,
the self-adjoint generator of the translations of the coordinate $A(b)$ on $L^2(\R^{\Lambda^{(1)}_{\mathrm{s}}}/\Gamma_{\mathrm{s}})$.
These operators commute.
Let $E:=(E_b)_{b\in\Lambda^{(1)}_{\mathrm{s}}}$;
for a bounded Borel function $g$ on $\R^{\Lambda^{(1)}_{\mathrm{s}}}$ the operator $g(E)$ is defined by the joint spectral calculus of the commuting self-adjoint operators $E_b$.
With $|x|^2:=\sum_bx(b)^2$, one has for $x\in\R^{\Lambda^{(1)}_{\mathrm{s}}}$ and $a>0$
\begin{align*}
\left(\exp(-i\langle x,E\rangle)\Phi\right)(A)&=\Phi(A-x),\\
\left(\exp\left(-\tfrac{a}{2}|E|^2\right)\Phi\right)(A)
&=\int_{\R^{\Lambda^{(1)}_{\mathrm{s}}}/\Gamma_{\mathrm{s}}}\mathrm{d}A'\,\sum_{\gamma\in\Gamma_{\mathrm{s}}}\ \prod_{b\in\Lambda^{(1)}_{\mathrm{s}}}g_a\left(A(b)-A'(b)+\gamma(b)\right)\cdot\Phi(A'):
\end{align*}
$\exp\left(-\frac{a}{2}|E|^2\right)$ is the convolution with the function $\sum_{\gamma\in\Gamma_{\mathrm{s}}}\prod_bg_a(A(b)+\gamma(b))$
(Lemma~\ref{lem:slice-kernels}\eqref{it:gaussian-kernels}).

\emph{The gauge group.}
For $a\in\R^{\Lambda_{\mathrm{s}}}$ let $U_{\mathrm{G}}(a)$ be the operator on $\mathcal{H}^{\mathrm{l}}$ given by
\[
\left(U_{\mathrm{G}}(a)\Phi\right)(A):=u(a)\,\Phi(A-d_{\mathrm{s}}a).
\]
For $a\in(2\pi/q)\cdot\Z^{\Lambda_{\mathrm{s}}}$ one has $u(a)=1$, the charges being integers, and $d_{\mathrm{s}}a\in\Gamma_{\mathrm{s}}$;
hence $U_{\mathrm{G}}$ is a unitary representation of the compact group $(\R/(2\pi/q)\Z)^{\Lambda_{\mathrm{s}}}$.
Let $P_{\mathrm{G}}$ be the orthogonal projection onto its invariant vectors.
Then
\[
\int_{[0,2\pi/q)^{\Lambda_{\mathrm{s}}}}\mathrm{d}a\,U_{\mathrm{G}}(a)=(2\pi/q)^{|\Lambda_{\mathrm{s}}|}\cdot P_{\mathrm{G}},
\]
and $P_{\mathrm{G}}$ commutes with every $g(E)$ (Lemma~\ref{lem:slice-kernels}\eqref{it:gauge-average}).

\emph{Crossing operators.}
For $\beta\in\R^{\Lambda^{(1)}_{\mathrm{s}}}$ and a subset $S$ of $\Lambda^{(1)}_{\mathrm{s}}$ let $|x|_S^2:=\sum_{b\in S}x(b)^2$, and
\begin{align*}
K&:=(2\pi/q)^{|\Lambda_{\mathrm{s}}|}\cdot\exp\left(-\tfrac12|E|^2\right)P_{\mathrm{G}},
&
K^{1/2}&:=(2\pi/q)^{\frac{|\Lambda_{\mathrm{s}}|}{2}}\cdot\exp\left(-\tfrac14|E|^2\right)P_{\mathrm{G}},\\
K_\beta&:=(2\pi/q)^{|\Lambda_{\mathrm{s}}|}\cdot\exp\left(-\tfrac12|E-\beta|^2\right)P_{\mathrm{G}},
&
K_{\beta,S}&:=K_\beta\exp\left(\tfrac14|E-\beta|_S^2\right).
\end{align*}
Since $P_{\mathrm{G}}$ commutes with the functions of $E$, the operator $K^{1/2}$ is non-negative and $(K^{1/2})^2=K$.
Likewise
$K_{\beta,S}=(2\pi/q)^{|\Lambda_{\mathrm{s}}|}\cdot\exp\left(-\frac12|E-\beta|^2+\frac14|E-\beta|_S^2\right)P_{\mathrm{G}}$,
a bounded operator, the exponent being at most $-\frac14|E-\beta|^2$.

\emph{The slice operator.}
Let $-1\le\rho\le1$, so that the operators $\widehat{\mathfrak{N}}(A)$ of Definition~\ref{def:modes} are defined.
For $k\in\Z$ and $A\in\R^{\Lambda^{(1)}_{\mathrm{s}}}$ let
\[
\mathfrak{M}_k(A):=\exp\left(-\tfrac14\left\|d_{\mathrm{s}}A+(2\pi/q)\cdot k\cdot1_{\sigma_0}\right\|^2\right)\cdot\widehat{\mathfrak{N}}(A):\ \mathcal{F}'\to\mathcal{F},
\]
where $1_{\sigma_0}$ is the indicator of $\sigma_0$ and $\|\cdot\|$ the Euclidean norm of $\R^{\Lambda^{(2)}_{\mathrm{s}}}$.
The function $A\mapsto\mathfrak{M}_k(A)$ is $\Gamma_{\mathrm{s}}$-periodic, because $d_{\mathrm{s}}\gamma=0$ for $\gamma\in\Gamma_{\mathrm{s}}$
and $\widehat{\mathfrak{N}}(A)$ depends on $A$ only through the phases $\exp(iqA(b))$;
it is bounded uniformly in $k$, since the entries of $\mathfrak{N}(A)$ are polynomials in the phases $\exp(\pm iqA(b))$.
The slice operator
\[
\mathfrak{M}:=\bigoplus_{k\in\Z}\mathfrak{M}_k:\ \mathcal{H}^{\mathrm{s}}\to\mathcal{H}^{\mathrm{l}}
\]
acts on the summand $k$ as the operator of multiplication by $\mathfrak{M}_k(A)$; it is bounded, and it depends on the choice of $\sigma_0$.
The operators
\begin{align*}
&\mathfrak{M}\mathfrak{M}^*\quad\text{and}\quad K^{1/2}\mathfrak{M}\mathfrak{M}^*K^{1/2}=\left(K^{1/2}\mathfrak{M}\right)\left(K^{1/2}\mathfrak{M}\right)^* &&\text{on }\mathcal{H}^{\mathrm{l}},\\
&\mathfrak{M}^*K\mathfrak{M}=\left(K^{1/2}\mathfrak{M}\right)^*\left(K^{1/2}\mathfrak{M}\right) &&\text{on }\mathcal{H}^{\mathrm{s}}
\end{align*}
are non-negative; through $\mathfrak{M}$ they depend on $\sigma_0$ as well.

\emph{A commutation relation.}
Let $\delta\in\R^{\Lambda^{(1)}_{\mathrm{s}}}$ satisfy $\langle\delta,\gamma\rangle\in2\pi\Z$ for every $\gamma\in\Gamma_{\mathrm{s}}$.
Then the function $A\mapsto\exp(i\langle\delta,A\rangle)$ is $\Gamma_{\mathrm{s}}$-periodic,
and the operator of multiplication by it, written $\exp(i\langle\delta,A\rangle)$, is unitary.
Since $\exp(-i\langle\delta,A\rangle)\,E_b\,\exp(i\langle\delta,A\rangle)=E_b+\delta(b)$ for every link $b$,
every bounded Borel function $g$ on $\R^{\Lambda^{(1)}_{\mathrm{s}}}$ satisfies
\begin{equation}\label{eq:setting-weyl}
g(E)\exp(i\langle\delta,A\rangle)=\exp(i\langle\delta,A\rangle)\,g(E+\delta).
\end{equation}
For a finite set $S$ of links of $\Z^2$ the same formula $E_b=-i\,\partial/\partial A(b)$, $b\in S$,
defines commuting self-adjoint operators on $L^2(\R^S)$, and $E_S:=(E_b)_{b\in S}$;
on $L^2(\R^S)$ the identity \eqref{eq:setting-weyl} holds with $E_S$ in the place of $E$,
for every $\delta\in\R^S$ and every bounded Borel function $g$ on $\R^S$.
\end{definition}

\subsection{Insertions and two-layer elements}\label{sec:insertions}

\begin{definition}[The mid-plane algebra]\label{def:mid-plane}
Let $S$ be a finite set of links of the slice lattice $\Z^2$ (Definition~\ref{def:slices}),
and let $\Lambda$ be a non-empty finite set of sites of $\Z^2$ that contains the end points of the links of $S$.
Put
\[
\mathcal{H}_S:=L^2(\R^S)\otimes\mathcal{F}\left(\mathcal{J}^{\mathrm{l}}(\Lambda)\right),
\]
where $\mathcal{F}(\mathcal{J}^{\mathrm{l}}(\Lambda))$ is the Fock space of the modes over $\Lambda$ at a link plane (Definition~\ref{def:modes}).
The variable of $L^2(\R^S)$ is $A\in\R^S$,
and $E_S=(E_b)_{b\in S}$ is the family of commuting self-adjoint operators on $L^2(\R^S)$ of Definition~\ref{def:slice-operators};
for $S'\subseteq S$ and a bounded Borel function $g$ on $\R^{S'}$, the operator $g(E_{S'})$ is formed with the joint spectral calculus of $E_{S'}=(E_b)_{b\in S'}$.
Functions of $A$ and of $E_S$ act on the first factor of $\mathcal{H}_S$, operators on the Fock space on the second.
For $a\in\R^\Lambda$, $u(a)$ is the diagonal unitary on $\mathcal{F}(\mathcal{J}^{\mathrm{l}}(\Lambda))$ of Definition~\ref{def:modes}.
An operator $\Upsilon$ on this Fock space is \emph{homogeneous of charge} $\mathfrak{q}(\Upsilon)\in\Z^\Lambda$ if
\[
u(a)\,\Upsilon\,u(a)^{-1}=\exp\left(iq\langle a,\mathfrak{q}(\Upsilon)\rangle\right)\cdot\Upsilon\qquad\text{for all }a\in\R^\Lambda.
\]
The \emph{mid-plane algebra} $\mathfrak{W}(S,\Lambda)$ is the unital $*$-algebra of operators on $\mathcal{H}_S$ generated by
\begin{enumerate}
\item\label{it:gaussian-generators}
the operators $g(E_{S'})$, $S'\subseteq S$, where $g$ lies in the linear span of the constant function $1$ and of the functions
$x\mapsto\exp\left(-\mathfrak{Q}(x)+\langle\lambda,x\rangle\right)$ on $\R^{S'}$,
with $\mathfrak{Q}$ a positive definite real quadratic form on $\R^{S'}$ and $\lambda\in\C^{S'}$;
\item\label{it:exponential-generators}
the operators $\exp(i\langle\delta,A\rangle)\cdot\Upsilon$, where $\delta\in\R^S$,
$\Upsilon\in\End\mathcal{F}(\mathcal{J}^{\mathrm{l}}(\Lambda))$ is homogeneous of a charge $\mathfrak{q}(\Upsilon)$,
and $d_{\mathrm{s}}^{*}\delta=q\cdot\mathfrak{q}(\Upsilon)$ on $\Z^2$;
here $d_{\mathrm{s}}^{*}$ is the adjoint of $d_{\mathrm{s}}$ on finitely supported cochains,
and $\delta$ and $\mathfrak{q}(\Upsilon)$ are extended by zero to all links and all sites of $\Z^2$.
\end{enumerate}
Every element of $\mathfrak{W}(S,\Lambda)$ is even, that is, it commutes with the parity $U_{\mathrm{f}}$ of the Fock space,
and it is gauge invariant: it commutes with the operators $\Phi\mapsto u(a)\,\Phi\left(\,\cdot-(d_{\mathrm{s}}a)|_S\right)$, $a\in\R^\Lambda$
(Lemma~\ref{lem:normal-form}).

For $\beta\in\R^S$, $\alpha_\beta$ denotes the $*$-automorphism of $\mathfrak{W}(S,\Lambda)$ that maps $g(E_{S'})$ to $g(E_{S'}-\beta|_{S'})$
and fixes the generators~\eqref{it:exponential-generators};
it is the conjugation by the unitary operator $\exp(i\langle\beta,A\rangle)$ (Lemma~\ref{lem:normal-form}).

Let $\Xi$ be a rectangular cross-section such that $\hull(S,\Lambda)$ embeds in $\Xi$ and does not contain the plaquette $\sigma_0$ (Definition~\ref{def:slices}).
The realisation $X\mapsto X_\Xi$ of $\mathfrak{W}(S,\Lambda)$ on the Hilbert space $\mathcal{H}^{\mathrm{l}}$ of Definition~\ref{def:slice-operators}
acts in the same way on every summand of $\mathcal{H}^{\mathrm{l}}$:
the links of $S$ are read as links of $\Lambda_{\mathrm{s}}$,
the operators $g(E_{S'})$ and $\exp(i\langle\delta,A\rangle)$ are given by the same formulas on $L^2\left(\R^{\Lambda^{(1)}_{\mathrm{s}}}/\Gamma_{\mathrm{s}}\right)$,
and $\Upsilon$ is replaced by $\iota(\Upsilon)$, the localisation with the signed relabelling of Definition~\ref{def:modes}.
By Lemma~\ref{lem:normal-form} this is a well defined unital $*$-representation.
For $X\in\mathfrak{W}(S,\Lambda)$ put
\[
\|X\|_\#:=\sup_{\Xi,\,\beta}\left\|(\alpha_\beta X)_\Xi\right\|,
\]
the supremum over all rectangular cross-sections $\Xi$ as above and all $\beta\in\R^S$.

Let $S\subseteq S_1$ and $\Lambda\subseteq\Lambda_1$, where $S_1$ is a finite set of links and $\Lambda_1$ a finite set of sites of $\Z^2$ that contains the end points of the links of $S_1$.
Then $\mathfrak{W}(S,\Lambda)$ is a subalgebra of $\mathfrak{W}(S_1,\Lambda_1)$:
an element of $\mathfrak{W}(S,\Lambda)$ is identified with its tensor product with the identity of $L^2(\R^{S_1\setminus S})$,
its Fock part being carried to $\mathcal{F}(\mathcal{J}^{\mathrm{l}}(\Lambda_1))$ by the localisation $\iota$.
The union of all the algebras $\mathfrak{W}(S,\Lambda)$ is a $*$-algebra, denoted by $\mathfrak{W}$.
By Lemma~\ref{lem:sharp-norm}, for each single cross-section $\Xi$ as above the number $\sup_{\beta\in\R^S}\|(\alpha_\beta X)_\Xi\|$ is the operator norm of $X$ on $\mathcal{H}_S$.
Hence $\|\cdot\|_\#$ is the operator norm on $\mathcal{H}_S$:
a C$^{*}$-norm on $\mathfrak{W}(S,\Lambda)$, invariant under every $\alpha_\beta$, with $\|1\|_\#=1$;
it does not change under the inclusions just described, and so it is a C$^{*}$-norm on $\mathfrak{W}$.
One and the same algebra $\mathfrak{W}(S,\Lambda)$ is associated with each of the coordinate link planes $(x_\mu,c)$, $\mu\in\{1,2,3\}$, $c\in\Z+\tfrac12$.
\end{definition}

\begin{definition}[Insertions]\label{def:insertions}
For a link $b$ and a site $y$ of $\Z^2$ and for $j'\in\Z$, let $(j',b)$ be the link of $\Z^3$ above $b$ in the plane $\{x_1=j'\}$
and $(j',y)$ the site above $y$ in this plane.
Let $c=j+\tfrac12$ with $j\in\Z$, let $S$ and $\Lambda$ be as in Definition~\ref{def:mid-plane}, and let $X\in\mathfrak{W}(S,\Lambda)$.
Let $k_X(A',A)\in\End\mathcal{F}(\mathcal{J}^{\mathrm{l}}(\Lambda))$, $A',A\in\R^S$, be the integral kernel of the operator
$\exp\left(-\tfrac14|E_S|^2\right)X\exp\left(-\tfrac14|E_S|^2\right)$ on $\mathcal{H}_S$.
Let $S^{\times}_c(\Lambda)$ be the sum of the terms of $\bar\psi D_A\psi$ that belong to the bonds over $\Lambda$ crossing the plane $\{x_1=c\}$, that is,
\[
S^{\times}_c(\Lambda):=-\sum_{(x,x')}\bar\psi(x)\,w(x,x')\,U_A(x,x')\,\psi(x'),
\]
with the hopping matrices $w(x,x')$ and the parallel transports $U_A(x,x')$ of Definition~\ref{def:chamber-operator},
the sum being taken over the ordered pairs $(x,x')$ of fermion sites such that $\{x,x'\}$ is a bond that crosses $\{x_1=c\}$
and $\{\pi(x),\pi(x')\}=\{(j,y),(j+1,y)\}$ for some $y\in\Lambda$.
The \emph{insertion} of $X$ at $c$ is the integrand
\[
G_X=G^{(c)}_X:=\Ins_c\left[\omega^{\frac12}\,k_X\!\left(\hat A_{j+1}|_S,\,(\hat A_j+d_{\mathrm{s}}\hat a_j)|_S\right)u(\hat a_j)\,\omega^{\frac12}\right]\cdot\prod_{b\in S}g_1\!\left(F^{\mathrm{el}}_j(b)\right)^{-1}\cdot\exp\left(S^{\times}_c(\Lambda)\right).
\]
Here $\hat A$ is a local lift and $F=d\hat A$ (Definition~\ref{def:integrands});
$F^{\mathrm{el}}_j$ is the electric part of $F$, and $\hat A_j$, $\hat A_{j+1}$ and $\hat a_j$ are the slice variables of $\hat A$,
that is, those of Definition~\ref{def:slices} with $\hat A$ in the place of $A$;
$\Ins_c$ is formed at the mirror $(x_1,c)$ with the modes of $\mathcal{J}^{\mathrm{l}}(\Lambda)$, so that its argument is a matrix indexed by the subsets of $\mathcal{J}^{\mathrm{l}}(\Lambda)$;
$\omega$ is the diagonal operator of the crossing coefficients on $\mathcal{F}(\mathcal{J}^{\mathrm{l}}(\Lambda))$,
and $u(\hat a_j)$ is formed with the restriction of $\hat a_j$ to $\Lambda$;
$g_1$ is the Gaussian density of Definition~\ref{def:slices}.
In $G_X$ the crossing terms $S^{\times}_c(\Lambda)$ are formed with $\hat A$ in the place of $A$:
the term of an ordered pair $(x,x')$ with $\pi(x)=(j,y)$ and $\pi(x')=(j+1,y)$ carries the phase $U_{\hat A}(x,x')=\exp\left(iq\hat a_j(y)\right)$,
and the term of the pair $(x',x)$ the phase $U_{\hat A}(x',x)=\exp\left(-iq\hat a_j(y)\right)$.
The last factor, $\exp(S^{\times}_c(\Lambda))$, is the inverse of $\exp(-S^{\times}_c(\Lambda))$;
both exponentials are even, since $S^{\times}_c(\Lambda)$ is.

Let $\mathcal{K}(c,S,\Lambda)$ be the union of the links $(j,b)$ and $(j+1,b)$, $b\in S$,
of the links from $(j,y)$ to $(j+1,y)$, $y\in\Lambda$,
and of the cubes $(j,y)+\left[-\tfrac12,\tfrac12\right]^3$ and $(j+1,y)+\left[-\tfrac12,\tfrac12\right]^3$, $y\in\Lambda$.
A box (Definition~\ref{def:integrands}) is determined by its \emph{centre}, the centre of symmetry of the cube,
and by its \emph{radius}, half the length of its edges.
The box $B(c,S,\Lambda)$ is the box of least radius that contains $\mathcal{K}(c,S,\Lambda)$ and, among the boxes of this radius that contain it, has the lexicographically least centre.
(Boxes are cubes, so that several boxes of the least radius may contain $\mathcal{K}(c,S,\Lambda)$; the condition on the centre selects one of them.)
Let $\prod_{\mu}[x^-_\mu,x^+_\mu]$ be the smallest product of intervals that contains $\mathcal{K}(c,S,\Lambda)$, and put
\[
x^0:=\left(\tfrac12(x^-_\mu+x^+_\mu)\right)_{\mu=1,2,3},\qquad r_0:=\max_{\mu}\tfrac12(x^+_\mu-x^-_\mu),\qquad B^{\#}(c,S,\Lambda):=x^0+[-2r_0,2r_0]^3;
\]
the number $r_0$ is the radius of $B(c,S,\Lambda)$.
An identity between two integrands that read in different boxes is understood as an equality in the sense of Lemma~\ref{lem:equal-integrands}.
\end{definition}

\begin{remark}
Over $(S,\Lambda)$, the insertion $G_X$ puts the kernel of $K^{\frac12}XK^{\frac12}$ in the place of the crossing weight $\prod_{b\in S}g_1(F^{\mathrm{el}}_j(b))\cdot\exp(-S^{\times}_c(\Lambda))$.
The box $B(c,S,\Lambda)$ is neither covariant under reflections nor monotone in $(S,\Lambda)$.
Indeed, the reflection $\theta_c$ in the plane $\{x_1=c\}$ leaves $c$, $S$, $\Lambda$ and $\mathcal{K}(c,S,\Lambda)$ unchanged,
but, since $x^-_1=j-\tfrac12$ and $x^+_1=j+\tfrac32$, the first coordinate of the centre of $B(c,S,\Lambda)$ is $j+\tfrac32-r_0$, which differs from $c$ as soon as $r_0>1$;
then $\theta_cB(c,S,\Lambda)\neq B(c,S,\Lambda)$.
For $c=\tfrac12$, the box
$B(c,\emptyset,\{(0,0)\})=\left[-\tfrac12,\tfrac32\right]\times\left[-\tfrac32,\tfrac12\right]^2$
is not contained in
$B(c,\emptyset,\{(0,0),(1,0)\})=\left[-\tfrac12,\tfrac32\right]^2\times\left[-\tfrac32,\tfrac12\right]$.
The box $B^{\#}(c,S,\Lambda)$, on the other hand, is covariant under $W_{\mathrm{c}}$
and contains every box of radius $r_0$ that contains $\mathcal{K}(c,S,\Lambda)$, in particular $B(c,S,\Lambda)$ (Lemma~\ref{lem:covariant-box}).
\end{remark}

\begin{definition}[Bounded elements]\label{def:bounded-elements}
Let $B$ be a box and $\mathbb{T}=(\Gamma,\chi)=C(n_1,n_2,n_3)$ a coordinate torus (Definition~\ref{def:tori}).
A \emph{placement} of $B$ in $\mathbb{T}$ is a finite family $g_1,\dots,g_r$ of elements of $W_{\mathrm{c}}$ (Definition~\ref{def:mirrors}), $r\geq1$, with the following two properties.
\begin{enumerate}
\item\label{it:placement-disjoint}
The boxes $g_1B,\dots,g_rB$ embed disjointly in $\Gamma$ (Definition~\ref{def:integrands}).
\item\label{it:placement-columns}
(Free columns.) For each $\mu\in\{1,2,3\}$ there is a plaquette $\sigma$ orthogonal to $e_\mu$ such that none of the sets $g_jB+v$, $1\leq j\leq r$, $v\in\Gamma$,
contains one of the plaquettes $\sigma+k\,e_\mu$, $k\in\Z$.
\end{enumerate}
Let $G\in\mathcal{G}^{\mathrm{ev}}(B)$ and $C\geq0$.
The integrand $G$ is \emph{bounded by $C$ with the box $B$} if for every coordinate torus $\mathbb{T}=(\Gamma,\chi)$, every placement $g_1,\dots,g_r$ of $B$ in $\mathbb{T}$ and every $\beta\in C^2(\Gamma)$
\[
\left|Z_{\mathbb{T}}\!\left[\varphi_\beta\prod_{j=1}^{r}G^{g_j}\right]\right|\leq C^r\cdot Z_{\mathbb{T}}[1],
\]
where $\varphi_\beta(F):=\exp(i\langle\beta,F\rangle)$ with the pairing $\langle\cdot,\cdot\rangle$ of $C^2(\Gamma)$,
$G^{g_j}$ is the image of $G$ under $g_j$ \eqref{eq:setting-group-action},
and $Z_{\mathbb{T}}$ is the functional of Definition~\ref{def:functional}.
\end{definition}

\begin{lemma}[Placements]\label{lem:placements}
Let $B$ be a box and $\mathbb{T}=(\Gamma,\chi)=C(n_1,n_2,n_3)$ a coordinate torus.
\begin{enumerate}
\item\label{it:placements-conjugation}
If $g_1,\dots,g_r$ is a placement of $B$ in $\mathbb{T}$ and $g\in W$,
then $g\mathbb{T}$ is a coordinate torus and $gg_1g^{-1},\dots,gg_rg^{-1}$ is a placement of the box $gB$ in $g\mathbb{T}$.
\item\label{it:placements-cells}
Let $\mathcal{Q}=(\mathfrak{F}_{\mathrm{c}};o_1,o_2,o_3;h)$ be a coordinate cell structure with $o_\mu\in\Z+\tfrac12$, $\mu=1,2,3$, that cuts $\mathbb{T}$ (Definition~\ref{def:cell-structures}),
and let $g_1,\dots,g_r\in W_{\mathrm{c}}$ be such that each of the boxes $g_1B,\dots,g_rB$ lies in an open cell of $\mathcal{Q}$.
Then $g_1,\dots,g_r$ have the property~\eqref{it:placement-columns} of Definition~\ref{def:bounded-elements}.
\end{enumerate}
\end{lemma}

\begin{proof}
Consider first~\eqref{it:placements-conjugation}.
Let $\vec g\in\Ocub$ be the linear part of $g$, so that $g(x+v)=g(x)+\vec g\,v$ for $x,v\in\R^3$.
The matrix $\vec g$ is a signed permutation matrix: $\vec g\,e_\mu=\pm e_{\mu'}$, where $\mu\mapsto\mu'$ is a permutation of $\{1,2,3\}$.
Hence $\vec g\,\Gamma=\sum_\mu n_\mu\Z e_{\mu'}$,
and the character $\chi\circ\vec g^{-1}$ of $g\mathbb{T}=(\vec g\,\Gamma,\chi\circ\vec g^{-1})$ takes the value $-1$ on the generators $n_\mu e_{\mu'}$,
because $\chi(\pm n_\mu e_\mu)=-1$;
so $g\mathbb{T}$ is a coordinate torus.
The subgroup $W_{\mathrm{c}}$ is normal in $W$, so the elements $gg_jg^{-1}$ lie in $W_{\mathrm{c}}$; moreover $(gg_jg^{-1})(gB)=g(g_jB)$.
The set $gB$ is again a box: it has the radius of $B$, and its centre is the image under $g$ of the centre of $B$,
because the signed permutation matrix $\vec g$ maps every box centred at the origin onto itself.
For $v\in\Gamma$ one has
\[
g\left(g_jB+[-1,1]^3+v\right)=g(g_jB)+[-1,1]^3+\vec g\,v,\qquad g\left(g_jB+v\right)=g(g_jB)+\vec g\,v.
\]
As $g$ is a bijection of $\R^3$ and $v\mapsto\vec g\,v$ a bijection of $\Gamma$ onto $\vec g\,\Gamma$,
the sets on the left are pairwise disjoint if and only if the corresponding sets on the right are;
hence the boxes $g(g_jB)$ embed disjointly in $\vec g\,\Gamma$.
Let $\mu\in\{1,2,3\}$ and let $\sigma$ be a plaquette orthogonal to $e_\mu$ as in property~\eqref{it:placement-columns}.
Since $g$ maps $\Z^3$ onto itself, $g\sigma$ is a plaquette orthogonal to $e_{\mu'}$, and $g(\sigma+k\,e_\mu)=g\sigma\pm k\,e_{\mu'}$.
A set $g(g_jB)+\vec g\,v=g(g_jB+v)$ contains $g(\sigma+k\,e_\mu)$ only if $g_jB+v$ contains $\sigma+k\,e_\mu$, which is excluded.
So the plaquette $g\sigma$ has the property required for the direction $\mu'$,
and since $\mu\mapsto\mu'$ is a bijection, every direction is provided for.

For~\eqref{it:placements-cells}, note that the coordinates of every $v\in\Gamma$ lie in $h\Z$, because $\mathcal{Q}$ cuts $\mathbb{T}$;
the translation by $v$ therefore permutes the open cells of $\mathcal{Q}$, and every set $g_jB+v$ lies in an open cell.
Fix $\mu$ and choose $\nu\neq\mu$.
As $o_\nu-\tfrac12$ is an integer, there is a plaquette $\sigma$ orthogonal to $e_\mu$ whose projection to the $x_\nu$-axis is the interval $\left[o_\nu-\tfrac12,o_\nu+\tfrac12\right]$;
it meets the wall $\{x_\nu=o_\nu\}$ of $\mathcal{Q}$.
The translation by $k\,e_\mu$ does not change the coordinate $x_\nu$, so every plaquette $\sigma+k\,e_\mu$, $k\in\Z$, meets the same wall.
An open cell of $\mathcal{Q}$ lies strictly between two consecutive walls orthogonal to $e_\nu$ and does not meet $\{x_\nu=o_\nu\}$.
Hence no plaquette $\sigma+k\,e_\mu$ is contained in an open cell, and a fortiori none is contained in one of the sets $g_jB+v$.
\end{proof}

\begin{definition}[Two-layer elements]\label{def:two-layer}
For $j\in\Z$ the \emph{thick layer} $j$ is the open slab $\left\{|x_1-j|<\tfrac12\right\}$.
The photon links that it contains are those of the plane $\{x_1=j\}$, which carry the variables $A_j$,
and its fermion sites are the $x\in\mathbb{X}$ with $x_1-j\in\{-v_1,-v_2,-v_3,v_3,v_2,v_1\}$ (Definition~\ref{def:fermion-sites}).
The sets $\mathbb{X}_{j,1},\dots,\mathbb{X}_{j,6}$ of the fermion sites with $x_1-j=-v_1,-v_2,-v_3,v_3,v_2,v_1$, respectively,
are the \emph{thin layers} of the thick layer $j$; they are ordered by $x_1$.
The successor $\mathbb{X}^{+}_{j,k}$ of $\mathbb{X}_{j,k}$ is $\mathbb{X}_{j,k+1}$ if $k\leq5$ and $\mathbb{X}_{j+1,1}$ if $k=6$;
$\mathbb{X}^{++}_{j,k}$ is the successor of $\mathbb{X}^{+}_{j,k}$,
and the predecessor $\mathbb{X}^{-}_{j,k}$ is the thin layer whose successor is $\mathbb{X}_{j,k}$.

Every $P\in\mathcal{P}$ (Definition~\ref{def:polynomials}) has a unique representation
\[
P=\sum_{\alpha}a_\alpha\cdot\exp\left(i\langle\delta_\alpha,\hat A\rangle\right)\cdot\Psi_\alpha
\]
with finitely many terms, coefficients $a_\alpha\neq0$ and pairwise different pairs $(\delta_\alpha,\Psi_\alpha)$,
where each $\Psi_\alpha$ is a monomial, written as a product of components of the fields in a fixed total order
(the monomials are linearly independent, and so are the functions $\hat A\mapsto\exp(i\langle\delta,\hat A\rangle)$ for pairwise different $\delta$).
The element $P$ \emph{reads} the link $b$ if $\delta_\alpha(b)\neq0$ for some $\alpha$,
and it reads the site $x$ if some $\Psi_\alpha$ contains a field at $x$.
As an element of $\mathcal{P}$, a character $\varphi_\beta$, $\beta\in C^2_{\mathrm{c}}$, equals $\exp(i\langle d^{*}\beta,\hat A\rangle)$, and the links that it reads are those of $\supp d^{*}\beta$.
The least (largest) integer $j$ for which $P$ reads a site of the thick layer $j$ or a link meeting the plane $\{x_1=j\}$ is the \emph{lowest} (\emph{highest}) \emph{plane} of $P$.

Let $c=j_0+\tfrac12$ with $j_0\in\Z$.
A \emph{two-layer monomial} at $c$ is a gauge invariant element $\exp(i\langle\delta,\hat A\rangle)\Psi$ of $\mathcal{P}$
in which $\delta$ is supported on the links of the planes $\{x_1=j_0\}$, $\{x_1=j_0+1\}$ and on the links between these two planes,
and $\Psi$ is a monomial in the fields of $\mathbb{X}_{j_0,6}\cup\mathbb{X}_{j_0+1,1}$, the two ends of the bonds that cross the plane $\{x_1=c\}$.
The linear span $\mathcal{P}_1(c)$ of the two-layer monomials at $c$ is the space of \emph{two-layer elements} at $c$.
For such a $\delta$ and for a link $b$ and a site $y$ of $\Z^2$,
let $\delta_0(b)$ and $\delta_1(b)$ be the values of $\delta$ on the links $(j_0,b)$ and $(j_0+1,b)$,
and $\delta_\perp(y)$ its value on the link from $(j_0,y)$ to $(j_0+1,y)$;
the triple $(\delta_0,\delta_1,\delta_\perp)$ is the \emph{exponent} of the monomial.
Two-layer monomials with $\delta\in q\,C^1(\Z)$ are called \emph{integer}.

Let $Y=\exp(i\langle\delta,\hat A\rangle)\Psi$ be a two-layer monomial at $c$ with $Y\neq1$.
Let $S$ be the set of the links $b$ of $\Z^2$ with $\delta_0(b)\neq0$ or $\delta_1(b)\neq0$,
and let $\Lambda_Y$, the set of the \emph{sites in play}, be a finite set of sites of $\Z^2$ that contains
the end points of the links of $S$, the sites $y$ with $\delta_\perp(y)\neq0$,
and every $y$ for which $\Psi$ contains a field at a site $x$ with $\pi(x)\in\{(j_0,y),(j_0+1,y)\}$.
Since $Y\neq1$, the set $\Lambda_Y$ is not empty.
As $|\rho|<1$, the crossing coefficients are positive and $\omega$ is invertible;
$\Upsilon_Y\in\End\mathcal{F}(\mathcal{J}^{\mathrm{l}}(\Lambda_Y))$ is the operator with
\[
\left.\left(\Psi\cdot\exp\left(-S^{\times}_c(\Lambda_Y)\right)\right)\right|_{\hat a_{j_0}=0}=\Ins_c\left[\omega^{\frac12}\,\Upsilon_Y\,\omega^{\frac12}\right],
\]
here $S^{\times}_c(\Lambda_Y)$ is formed with $\hat A$ as in Definition~\ref{def:insertions},
and the evaluation at $\hat a_{j_0}=0$ makes all phases $U_{\hat A}(x,x')=\exp\left(\pm iq\hat a_{j_0}(y)\right)$, $y\in\Lambda_Y$, of the crossing bonds equal to $1$.
The operator $\Upsilon_Y$ exists and is unique by Proposition~\ref{prop:two-layer-insertions}(a).

Let $\Xi$ be a rectangular cross-section as in Definition~\ref{def:mid-plane}, with $\Lambda_Y$ in the place of $\Lambda$.
On $\mathcal{H}_S$ and on $\mathcal{H}^{\mathrm{l}}$ put
\[
k^0_Y:=\exp\left(i\langle\delta_0+\delta_1,A\rangle\right)\cdot\exp\left(-\tfrac12|E+\delta_0|^2\right)\cdot\Upsilon_Y,
\]
where $E=E_S$ on $\mathcal{H}_S$, while on $\mathcal{H}^{\mathrm{l}}$ the vector $E=(E_b)_{b\in\Lambda^{(1)}_{\mathrm{s}}}$ is the one of Definition~\ref{def:slice-operators},
$\delta_0$ and $\delta_1$ are read as cochains of $\Lambda_{\mathrm{s}}$, and $\Upsilon_Y$ acts through $\iota$;
on $\mathcal{H}^{\mathrm{l}}$ put moreover
\[
k_Y:=(2\pi/q)^{|\Lambda_{\mathrm{s}}|}\cdot k^0_Y\cdot P_{\mathrm{G}}.
\]
On $\mathcal{H}^{\mathrm{l}}$ the factor $\exp(i\langle\delta_0+\delta_1,A\rangle)$ is a function on $\R^{\Lambda^{(1)}_{\mathrm{s}}}/\Gamma_{\mathrm{s}}$,
because $d_{\mathrm{s}}^{*}(\delta_0+\delta_1)\in q\Z^{\Z^2}$
(Proposition~\ref{prop:two-layer-insertions}(a) and Lemma~\ref{lem:normal-form}).
On $\mathcal{H}_S$ the relation~\eqref{eq:setting-weyl}, which holds on $L^2(\R^S)$ as well, gives
$k^0_Y=\exp(i\langle\delta_1,A\rangle)\cdot\exp\left(-\tfrac12|E_S|^2\right)\cdot\Upsilon_Y\cdot\exp(i\langle\delta_0,A\rangle)$;
on $\mathcal{H}^{\mathrm{l}}$ the two factors $\exp(i\langle\delta_0,A\rangle)$ and $\exp(i\langle\delta_1,A\rangle)$ are in general not defined,
since in general only $\delta_0+\delta_1$ has its divergence in $q\Z^{\Z^2}$.

With $\widetilde E:=E_S+\delta_0$ and $\varepsilon>0$, let
\[
\begin{aligned}
X'_Y&:=\exp\left(i\langle\delta_0+\delta_1,A\rangle\right)\cdot\exp\left(-\tfrac14|\widetilde E|^2+\tfrac14\langle\widetilde E,\delta_1-\delta_0\rangle+\tfrac18\left(|\delta_0|^2+|\delta_1|^2\right)\right)\cdot\Upsilon_Y,\\
X_{Y,\varepsilon}&:=\exp\left(i\langle\delta_0+\delta_1,A\rangle\right)\cdot\exp\left(-\varepsilon|\widetilde E+\delta_1|^2-\varepsilon|\widetilde E-\delta_0|^2+\tfrac12\langle\widetilde E,\delta_1-\delta_0\rangle+\tfrac14\left(|\delta_0|^2+|\delta_1|^2\right)\right)\cdot\Upsilon_Y,\\
P_\varepsilon&:=\exp\left(-\varepsilon|E_S|^2\right),
\end{aligned}
\]
operators on $\mathcal{H}_S$.
They are bounded and lie in $\mathfrak{W}(S,\Lambda_Y)$, and on $\mathcal{H}^{\mathrm{l}}$
\begin{equation}\label{eq:setting-kernel-relations}
k_Y=K_{0,S}^{\frac12}\,(X'_Y)_\Xi\,K_{0,S}^{\frac12},\qquad K^{\frac12}\,(X_{Y,\varepsilon})_\Xi\,K^{\frac12}=(P_\varepsilon)_\Xi\,k_Y\,(P_\varepsilon)_\Xi
\end{equation}
(Proposition~\ref{prop:two-layer-insertions}(b) and~(c)).
Thus the bounded operators $X'_Y$ and $X_{Y,\varepsilon}$ take the places of the formal expressions
$K_{0,S}^{-\frac12}k_YK_{0,S}^{-\frac12}$ and $P_\varepsilon K^{-\frac12}k_YK^{-\frac12}P_\varepsilon$.
Put
\[
\xi_Y(\varepsilon):=\left\|X'_Y-P_\varepsilon X'_YP_\varepsilon\right\|_\#.
\]

The monomial $1$ is given no set of sites, and hence no box $B^{\#}(c,S,\Lambda_Y)$; for it, $\delta=0$ and $S=\emptyset$.
One puts $\Upsilon_1:=1$, the identity of the Fock spaces $\mathcal{F}(\mathcal{J}^{\mathrm{l}}(\Lambda))$,
and $X'_1:=X_{1,\varepsilon}:=1$, the unit of $\mathfrak{W}$.
On $\mathcal{H}_S$ one puts $k^0_1:=1$ as well, whereas on $\mathcal{H}^{\mathrm{l}}$
\[
k^0_1:=\exp\left(-\tfrac12|E|^2\right),
\qquad
k_1:=(2\pi/q)^{|\Lambda_{\mathrm{s}}|}\cdot k^0_1\cdot P_{\mathrm{G}}=K .
\]
These are the values of the formulas above for $\delta=0$ and $\Upsilon_1=1$:
for $S=\emptyset$ the vector $E_S$ has no components, while on $\mathcal{H}^{\mathrm{l}}$ the vector $E$ has one component for each link of $\Lambda_{\mathrm{s}}$.
Moreover, for $S=\emptyset$ one has $P_\varepsilon=1$ and $K_{0,\emptyset}=K$;
as the realisation is unital, the two relations~\eqref{eq:setting-kernel-relations} hold for $Y=1$ as well.
The maps $Y\mapsto\Upsilon_Y$, $Y\mapsto k^0_Y$, $Y\mapsto k_Y$, $Y\mapsto X'_Y$ and $Y\mapsto X_{Y,\varepsilon}$ are extended linearly to $\mathcal{P}_1(c)$,
operators attached to different monomials being compared through the inclusions of Definition~\ref{def:mid-plane}.
After this extension the relations~\eqref{eq:setting-kernel-relations} hold monomial by monomial:
for each two-layer monomial, the operators $K_{0,S}$ and $P_\varepsilon$ in~\eqref{eq:setting-kernel-relations} are built from the set $S$ of that monomial.
The map $Y\mapsto\xi_Y(\varepsilon)$ is not linear;
its formula is applied, unchanged, to the finite combinations $Y$ of two-layer monomials that all have the same exponent,
since such monomials share one set $S$ and one operator $P_\varepsilon$.
Let $\Theta_c:=\Theta_{(x_1,c)}$ be the reflection of Definition~\ref{def:grassmann-reflections} at the mirror $(x_1,c)$,
acting on $\mathcal{P}$ by~\eqref{eq:setting-reflection-coefficients}.
If $Y$ is a two-layer monomial with the exponent $(\delta_0,\delta_1,\delta_\perp)$,
then $\Theta_cY$ is a finite combination of two-layer monomials with the exponent $(-\delta_1,-\delta_0,\delta_\perp)$, hence with the same set $S$
(Proposition~\ref{prop:two-layer-insertions}(d)).
Indeed, $\Theta_cY=\exp(-i\langle\delta,\theta_c^{*}\hat A\rangle)\cdot\Theta_c(\Psi)$ by~\eqref{eq:setting-reflection-coefficients},
where $\langle\delta,\theta_c^{*}\hat A\rangle=\langle\delta_0,\hat A_{j_0+1}\rangle+\langle\delta_1,\hat A_{j_0}\rangle-\langle\delta_\perp,\hat a_{j_0}\rangle$ by~\eqref{eq:setting-slice-reflections},
and $\Theta_c(\Psi)$ is a combination of monomials in the fields of the two thin layers $\mathbb{X}_{j_0,6}$ and $\mathbb{X}_{j_0+1,1}$, which $\theta_c$ exchanges.
\end{definition}

\begin{definition}[Other coordinate directions]\label{def:other-directions}
The objects of Subsections~\ref{sec:slices-fock} and~\ref{sec:insertions} are attached to the direction $e_1$, the time direction.
For $\mu\in\{1,2,3\}$ let $g_\mu\in\Ocub$ be the cyclic permutation matrix with $g_\mu e_1=e_\mu$;
thus $g_1=1$, $g_2$ maps $e_1\mapsto e_2\mapsto e_3\mapsto e_1$, and $g_3=g_2^{-1}$.
Each $g_\mu$ is an even element of $W$, so that $G\mapsto G^{g_\mu}$ is linear \eqref{eq:setting-group-action}.
The objects attached to a coordinate link plane $(x_\mu,c)$, $c\in\Z+\tfrac12$, are the images under $g_\mu$ of the corresponding objects at $(x_1,c)$;
for instance, the insertion of $X\in\mathfrak{W}(S,\Lambda)$ at $(x_\mu,c)$ is $(G^{(c)}_X)^{g_\mu}$.
Moreover,
\[
\Theta_{\mu,c}:=\Theta_{(x_\mu,c)},\qquad\mathcal{P}_1(\mu,c):=\left\{Y^{g_\mu}:Y\in\mathcal{P}_1(c)\right\}.
\]
The space $\mathcal{P}_1(\mu,c)$ does not depend on the choice of the matrix:
$\mathcal{P}_1(\mu,c)=\{Y^{g}:Y\in\mathcal{P}_1(c)\}$ for every $g\in\Ocub$ with $g\,e_1=e_\mu$ (Lemma~\ref{lem:data-transport}(v)(a)).
\end{definition}

\subsection{Continuum notation}\label{sec:continuum-notation}

\begin{definition}[Cutoffs and continuum laws]\label{def:cutoffs}
Let $L\ge3$ be an odd integer and let $0<\bar m\le1$.
The integer $L^{\lfloor\log_L\bar m^{-1}\rfloor}\ge1$ is the scale of \cite[Section~2.1]{NoncompactQED}, where it is denoted $h$.
A pair $(N,T)$ of integers $N,T\ge0$ is called a \emph{cutoff}.
At the cutoff $(N,T)$ put
\[
\begin{aligned}
\eta&:=L^{-N},
&\qquad
n&:=8\cdot L^{\lfloor\log_L\bar m^{-1}\rfloor}\cdot L^{T+N},\\
R&:=n\cdot\eta=8\cdot L^{\lfloor\log_L\bar m^{-1}\rfloor}\cdot L^{T},
&\qquad
\Lambda_{\eta,R}&:=\eta\Z^3/R\Z^3 ;
\end{aligned}
\]
thus $n$ is an even integer with $n\ge8$, and $R\ge8$.

The plaquettes of $\Lambda_{\eta,R}$ are the pairs $\sigma=(x;\mu\nu)$ with $x\in\Lambda_{\eta,R}$ and $1\le\mu<\nu\le3$,
the plaquette $(x;\mu\nu)$ being spanned at $x$ by the vectors $\eta\cdot e_\mu$ and $\eta\cdot e_\nu$.
A two-cochain on $\Lambda_{\eta,R}$ is a real function on its plaquettes.
For complex functions $u,u'$ on the plaquettes put
\[
\langle u,u'\rangle_\eta:=\eta^3\cdot\sum_{\sigma}u(\sigma)\cdot u'(\sigma),
\]
the sum being taken over all plaquettes, and $|u|_\eta:=\langle u,u\rangle_\eta^{\frac12}$ for a two-cochain $u$.
The bijection $x\mapsto\eta\cdot x$ of $\Z^3/n\Z^3$ onto $\Lambda_{\eta,R}$ maps each plaquette $(x;\mu\nu)$ onto the plaquette $(\eta\cdot x;\mu\nu)$;
through it the $n\Z^3$-periodic two-cochains on $\Z^3$, that is, the elements of $C^2(n\Z^3)$ (Definition~\ref{def:configurations}),
are identified with the two-cochains on $\Lambda_{\eta,R}$.

Let $\mathcal S=\mathcal S(\R^3;\Lambda^2)$ be the real vector space of the two-forms $f=(f_{\mu\nu})_{\mu<\nu}$ on $\R^3$ whose coefficients are real Schwartz functions;
its elements are the test forms.
A complex test form is a two-form $f=f'+if''$ with $f',f''\in\mathcal S$;
its coefficients are $f_{\mu\nu}=f'_{\mu\nu}+if''_{\mu\nu}$.
The space $\Lambda^2$ carries the scalar product for which the forms $\mathrm{d}x_\mu\wedge\mathrm{d}x_\nu$, $1\le\mu<\nu\le3$, are orthonormal;
accordingly, as in \cite[Sections~2.4 and~8.1.1]{NoncompactQED}, the $L^2$ norm of a test form $f$, real or complex, is
\[
\|f\|_2:=\left(\sum_{\mu<\nu}\int_{\R^3}\mathrm{d}x\,\left|f_{\mu\nu}(x)\right|^2\right)^{\frac12}.
\]
Let $\mathcal S'=\mathcal S'(\R^3;\Lambda^2)$ be the space of the continuous linear forms on $\mathcal S$.
For $F\in\mathcal S'$ and $f\in\mathcal S$ the value of $F$ at $f$ is written $F(f)$,
and $F\mapsto F(f)$ is regarded as a function on $\mathcal S'$;
for a complex test form $f=f'+if''$ one puts $F(f):=F(f')+iF(f'')$.
The $\sigma$-algebra on $\mathcal S'$ generated by the functions $F\mapsto F(f)$, $f\in\mathcal S$, is denoted $\Sigma$.
For $s\in\R$ let $H_s\subset\mathcal S'$ be the Hermite--Sobolev space of \cite[Sections~2.4 and~7.1.5]{NoncompactQED},
a separable Hilbert space with the norm $\|\cdot\|_s$;
one has $H_{-3}\subset H_{-4}$, and $\|F\|_{-4}\le\|F\|_{-3}$ for $F\in H_{-3}$.

For a test form $f$, real or complex, and a plaquette $\sigma=(x;\mu\nu)$ of $\Lambda_{\eta,R}$ put
\[
f_\eta(\sigma):=\eta^{-3}\cdot\int_{Q(\sigma)}\mathrm{d}y\,\sum_{v\in\Z^3}f_{\mu\nu}(y+R\cdot v),\qquad
Q(\sigma):=x+\frac{\eta}{2}\cdot(e_\mu+e_\nu)+\left[-\frac{\eta}{2},\frac{\eta}{2}\right)^3,
\]
where $x\in\eta\Z^3$ is a representative of the corner of $\sigma$;
the integral does not depend on this choice, its integrand being $R$-periodic.
For fixed $\mu<\nu$ the sets $Q(x;\mu\nu)+R\Z^3$, $x\in\Lambda_{\eta,R}$, form a partition of $\R^3$.
For a two-cochain $F$ on $\Lambda_{\eta,R}$, the \emph{step two-form} of $F$ is the two-form $F^\sharp=(F^\sharp_{\mu\nu})_{\mu<\nu}$ on $\R^3$ with the coefficients
\[
F^\sharp_{\mu\nu}(y):=F(\sigma)\qquad\text{for }y\in Q(\sigma)+R\Z^3,\ \sigma=(x;\mu\nu).
\]
These coefficients are bounded and $R$-periodic; hence $F^\sharp\in\mathcal S'$, and
\[
F^\sharp(f)=\int\mathrm{d}y\,\sum_{\mu<\nu}f_{\mu\nu}(y)\cdot F^\sharp_{\mu\nu}(y)=\langle f_\eta,F\rangle_\eta\qquad(f\in\mathcal S).
\]
The map $f\mapsto f_\eta$ is the test map of \cite[Sections~2.4 and~7.1.3]{NoncompactQED}, where $F^\sharp(f)=\langle f_\eta,F\rangle_\eta$ is denoted $X_{\eta,R}(f)$.
The test map and $F^\sharp$ are defined, and the last identity holds, on every torus $\eta\Z^3/R\Z^3$ with $R\ge1$ and $R/\eta$ an integer \cite[Section~7.1.3]{NoncompactQED}.
By \cite[Lemma~7.5]{NoncompactQED}, which holds on each of these tori,
the argument at the end of the proof of \cite[Theorem~7.12(1)]{NoncompactQED} gives
\[
\|F^\sharp\|_{-3}\le8\cdot|F|_\eta
\]
for every two-cochain $F$ on $\Lambda_{\eta,R}$.
The translations $\tau_b$, $b\in\R^3$, the maps $U_r$, $r\in\Ocub$, and the coordinate reflections $\mathsf R_\nu$, $\nu=1,2,3$,
act on test forms, and their transposes $\tau_b^*$, $U_r^*$, $\mathsf R_\nu^*$ act on $\mathcal S'$, as in \cite[Section~2.5]{NoncompactQED}.

Let $e>0$ be the charge; at the cutoff $(N,T)$ the lattice charge is $q:=e\cdot\eta^{\frac12}$.
The Wilson parameter $\rho\in(0,1)$ and the point $v_0\in C_0$ (Definition~\ref{def:chambers}) do not depend on the cutoff,
and the weights $w_0,\dots,w_3$ are the standard weights of Definition~\ref{def:chamber-operator}.
A \emph{mass assignment} is an arbitrary sequence $N\mapsto m_{0,N}$, $N\ge0$, of real numbers; it does not depend on $T$.
Given a mass assignment, put
\[
m_N:=\bar m\cdot\eta+m_{0,N},\qquad M:=m_N+\frac{\rho}{2}\cdot\sum_{i=0}^{3}w_i\qquad(N\ge0);
\]
at every cutoff $(N,T)$ the chamber operator is formed with the mass parameter $M$,
so that its lattice mass $M-\frac{\rho}{2}\cdot\sum_{i=0}^{3}w_i$ (Definition~\ref{def:chamber-operator}) equals $m_N$.

Let a mass assignment be given, and let $(N,T)$ be a cutoff.
On the cubic torus $C(n,n,n)$ of Definition~\ref{def:tori}, with the side $n$ of the cutoff,
let the functional of Definition~\ref{def:functional} be formed with the lattice charge $q=e\cdot\eta^{\frac12}$,
the Wilson parameter $\rho$, the point $v_0$, the standard weights and the mass parameter $M=m_N+\frac{\rho}{2}\cdot\sum_{i=0}^{3}w_i$.
By Theorem~\ref{thm:partition-positivity}, $0<Z_{C(n,n,n)}[1]<\infty$;
hence the law $\mu_{C(n,n,n)}$ on $\Omega$ of Definition~\ref{def:functional} is defined.
It is the law of the periodic extension of the curvature;
so it is carried by the $n\Z^3$-periodic elements of $\Omega$, which are two-cochains on $\Lambda_{\eta,R}$ through the identification above.
Let $F^{\mathrm{u}}$ have the law $\mu_{C(n,n,n)}$.
Read on $\Lambda_{\eta,R}$, $F^{\mathrm{u}}$ is the curvature in lattice units, and $F:=\eta^{-\frac32}\cdot F^{\mathrm{u}}$ is the physical curvature.
The law $P^{\mathrm{ch}}_{N,T}$ is the law of the step two-form $F^\sharp$ of the physical curvature $F$.
The map $F\mapsto F^\sharp$ is linear on the finite-dimensional space of the two-cochains on $\Lambda_{\eta,R}$,
with values in $H_{-3}$ by the bound above;
it is therefore continuous into $H_{-3}$ and into $H_{-4}$.
So $P^{\mathrm{ch}}_{N,T}$ is a Borel probability measure on $H_{-3}$ and on $H_{-4}$, and a probability measure on $(\mathcal S',\Sigma)$.
\end{definition}

\begin{definition}[Limit points]\label{def:limit-points}
Fix $L$, $\bar m$ and $e$ as in Definition~\ref{def:cutoffs}, together with a mass assignment.
A \emph{limit point} of the family $\left(P^{\mathrm{ch}}_{N,T}\right)$ is a Borel probability measure $\mu$ on $H_{-4}$
such that $P^{\mathrm{ch}}_{N_i,T_i}\to\mu$ weakly on $H_{-4}$ for some sequence of cutoffs $(N_i,T_i)$ with $N_i\to\infty$ and $T_i\to\infty$.
For an integer $N\ge0$ let $\mathcal V_N$ be the set of the Borel probability measures on $H_{-4}$
to which $P^{\mathrm{ch}}_{N,T_i}$ converges weakly on $H_{-4}$ along some sequence of integers $T_i\ge0$ with $T_i\to\infty$;
the elements of $\mathcal V_N$ are the \emph{volume limits} at the spacing $L^{-N}$.
A \emph{volume-first limit point} is a Borel probability measure $\mu$ on $H_{-4}$
for which there are integers $N_i\ge0$ with $N_i\to\infty$ and measures $\mu^{(N_i)}\in\mathcal V_{N_i}$
such that $\mu^{(N_i)}\to\mu$ weakly on $H_{-4}$:
the volume index $T$ tends to infinity first, at a fixed spacing, and the spacing index $N$ afterwards.
\end{definition}

\begin{definition}[Laws on distributions]\label{def:laws}
A \emph{law on $\mathcal S'$} is a probability measure on $(\mathcal S',\Sigma)$;
where no law on $\Omega$ (Definition~\ref{def:characters}) is in question, it is simply called a \emph{law}.
In \cite[Section~8.1]{NoncompactQED} the space $\mathcal S'$ is denoted $Q$, and the function $F\mapsto F(f)$ is denoted $X(f)$.
For a law $\mu$ and a $\Sigma$-measurable function $\Phi$ on $\mathcal S'$ that is $\mu$-integrable or takes its values in $[0,\infty]$,
put $\mu(\Phi):=\int\mathrm{d}\mu(F)\,\Phi(F)$.
Let $L^2(\mu)$ be the complex Hilbert space of the $\mu$-square-integrable functions on $\mathcal S'$, taken modulo $\mu$-null functions,
with the scalar product $\langle\Phi,\Phi'\rangle_\mu:=\mu\left(\overline{\Phi}\cdot\Phi'\right)$, which is antilinear in the first entry.
For an open set $U\subset\R^3$ let $\Sigma_U\subset\Sigma$ be the $\sigma$-algebra generated by the functions $F\mapsto F(f)$
with $f\in C^\infty_c(\R^3;\Lambda^2)$ real and $\supp f\subset U$,
and let $L^2(\Sigma_U,\mu)$ be the closed subspace of $L^2(\mu)$ formed by the classes that contain a $\Sigma_U$-measurable function.
The set $H_{-4}$ belongs to $\Sigma$, and the trace of $\Sigma$ on $H_{-4}$ is the Borel $\sigma$-algebra of $H_{-4}$ \cite[Section~7.1.5]{NoncompactQED};
accordingly a Borel probability measure $\mu$ on $H_{-4}$, for instance a limit point, a volume limit or a volume-first limit point,
is identified with the law $B\mapsto\mu(B\cap H_{-4})$.

A Euclidean map is a map $\mathfrak g\colon x\mapsto gx+a$ of $\R^3$ onto itself with $g\in\mathrm{O}(3)$ and $a\in\R^3$.
It acts on test forms $f$, on elements $F$ of $\mathcal S'$, on functions $\Phi$ on $\mathcal S'$ and on laws $\mu$ by
\[
\begin{aligned}
(\mathfrak g\cdot f)(x)&:=(\wedge^2g)\,f\left(\mathfrak g^{-1}x\right),
&\qquad
(\mathfrak gF)(f)&:=F\left(\mathfrak g^{-1}\cdot f\right),\\
(\hat{\mathfrak g}\Phi)(F)&:=\Phi\left(\mathfrak g^{-1}F\right),
&\qquad
(\mathfrak g_*\mu)(B)&:=\mu\left(\mathfrak g^{-1}B\right),
\end{aligned}
\]
where $\wedge^2g$ is the linear map induced by $g$ on $\Lambda^2$, $B\in\Sigma$ and $\mathfrak g^{-1}B:=\{\mathfrak g^{-1}F:F\in B\}$ \cite[Section~8.1.2]{NoncompactQED}.
Thus $\hat{\mathfrak g}$ maps the function $F\mapsto F(f)$ onto the function $F\mapsto F(\mathfrak g\cdot f)$;
if $\mathfrak g_*\mu=\mu$, then $\hat{\mathfrak g}$ is a unitary operator on $L^2(\mu)$ \cite[Section~8.1.2]{NoncompactQED}.
An element $g$ of $\Ocub$ is identified with the Euclidean map $x\mapsto gx$.
For $a\in\R^3$ let $\tau_a$ be the translation $x\mapsto x+a$, and put $U(a):=\hat\tau_a$;
thus $(\tau_a\cdot f)(x)=f(x-a)$, and $U(a)$ maps $F\mapsto F(f)$ onto $F\mapsto F(\tau_a\cdot f)$.
For $j\in\{1,2,3\}$ let $\theta_j$ be the reflection of $\R^3$ in the plane $\{x_j=0\}$.
For $b\in\R^3$, $g\in\Ocub$ and $\nu\in\{1,2,3\}$ the test forms $\tau_b\cdot f$, $g\cdot f$ and $\theta_\nu\cdot f$ are the test forms $\tau_bf$, $U_gf$ and $\mathsf R_\nu f$ of Definition~\ref{def:cutoffs}.
The antilinear map $\Theta_j$ on functions on $\mathcal S'$ is $\Theta_j\Phi:=\overline{\hat\theta_j\Phi}$, that is,
\[
(\Theta_j\Phi)(F)=\overline{\Phi(\theta_jF)} .
\]
For a law $\mu$ with $(\theta_j)_*\mu=\mu$, the \emph{reflected form} of $\mu$ in the direction $j$ is the sesquilinear form
\[
(\Phi,\Phi')\mapsto\mu\left((\Theta_j\Phi)\cdot\Phi'\right)=\langle\hat\theta_j\Phi,\Phi'\rangle_\mu\qquad\left(\Phi,\Phi'\in L^2(\mu)\right).
\]
In \cite[Section~8.1.2]{NoncompactQED} the complex linear operator $\hat\theta_j$ is denoted $\Theta_j$,
and the antilinear operator $\Theta_j$ is the operator $\Theta_{j,0}$ of \cite[Section~2.5]{NoncompactQED}.

As for laws on $\Omega$ (Definition~\ref{def:characters}), a law $\mu$ on $\mathcal S'$ satisfies \emph{Gaussian domination with the constant one}, condition~\textup{(G)}, if
\[
\int\mathrm{d}\mu(F)\,\exp F(f)\le\exp\left(\frac12\cdot\|f\|_2^2\right)\qquad(f\in\mathcal S).
\]
For $J\subset\{1,2,3\}$ the law $\mu$ satisfies $(T_J)$ if
\begin{equation}\tag{$T_J$}\label{cond:translation-invariance}
(\tau_a)_*\mu=\mu\qquad\left(a\in\textstyle\sum_{l\in J}\R\cdot e_l\right);
\end{equation}
in indices, $123$ stands for $\{1,2,3\}$ and $j$ for $\{j\}$,
so that $(T_{123})$ is the invariance under all translations of $\R^3$ and $(T_j)$ the invariance under the translations along $e_j$.
For $j\in\{1,2,3\}$ the law $\mu$ satisfies $(R_j)$ if
\begin{equation}\tag{$R_j$}\label{cond:reflection-invariance}
(\theta_j)_*\mu=\mu ,
\end{equation}
and it satisfies $(RP_j)$ if
\begin{equation}\tag{$RP_j$}\label{cond:reflection-positivity}
\sum_{\alpha,\alpha'=1}^{k}\overline{z_\alpha}\cdot z_{\alpha'}\cdot\int\mathrm{d}\mu(F)\,\exp\left(i\left(F(f_{\alpha'})-F(\theta_j\cdot f_\alpha)\right)\right)\ge0
\end{equation}
for every integer $k\ge1$, all $z_1,\dots,z_k\in\C$ and all real $f_1,\dots,f_k\in C^\infty_c(\R^3;\Lambda^2)$ with $\supp f_\alpha\subset\{x_j>0\}$.
The map $\Theta_j$ takes the function $F\mapsto\exp\left(iF(f)\right)$, $f$ real, to the function $F\mapsto\exp\left(-iF(\theta_j\cdot f)\right)$;
hence the left side of \eqref{cond:reflection-positivity} equals $\mu\left((\Theta_j\Phi)\cdot\Phi\right)$ for $\Phi:=\sum_{\alpha=1}^{k}z_\alpha\cdot\exp\left(iF(f_\alpha)\right)$.

For a unit vector $n\in\R^3$ and $c\in\R$, the reflection in the plane $\{n\cdot x=c\}$ is the Euclidean map
\[
r_{n,c}\colon x\mapsto x-2\cdot(n\cdot x-c)\cdot n ;
\]
one has $r_{e_j,0}=\theta_j$.
The law $\mu$ is \emph{reflection positive at the plane} $\{n\cdot x=c\}$ if the inequality \eqref{cond:reflection-positivity} holds with $r_{n,c}\cdot f_\alpha$ in place of $\theta_j\cdot f_\alpha$
for every integer $k\ge1$, all $z_1,\dots,z_k\in\C$ and all real $f_1,\dots,f_k\in C^\infty_c(\R^3;\Lambda^2)$ with $\supp f_\alpha\subset\{n\cdot x>c\}$ \cite[Section~8.1.3]{NoncompactQED};
for $n=e_j$ and $c=0$ this is $(RP_j)$.
If $(r_{n,c})_*\mu=\mu$, reflection positivity at $\{n\cdot x=c\}$ is equivalent to reflection positivity at $\{(-n)\cdot x=-c\}$,
which is the same plane with the opposite half space \cite[Section~8.1.3]{NoncompactQED}.

The class $\mathcal C$ is the set of the laws that satisfy $(T_{123})$ and, for $j=1,2,3$, $(R_j)$ and $(RP_j)$ \cite[Section~8.1.3]{NoncompactQED}.
A law $\mu$ satisfies the hypothesis \eqref{cond:soft-hypothesis} if
\begin{equation}\tag{H}\label{cond:soft-hypothesis}
\mu\in\mathcal C,\qquad g_*\mu=\mu\quad(g\in\Ocub),\qquad\int\mathrm{d}\mu(F)\,\exp F(f)\le\exp\left(\frac12\cdot\|f\|_2^2\right)\quad(f\in\mathcal S).
\end{equation}
Since $\theta_1,\theta_2,\theta_3\in\Ocub$, a law satisfies \eqref{cond:soft-hypothesis} if and only if it satisfies \textup{(G)} and $(T_{123})$, is invariant under $\Ocub$ and satisfies $(RP_j)$ for $j=1,2,3$;
this is the hypothesis \textup{(H)} of \cite[Sections~2.5 and~8.1.5]{NoncompactQED}.

Let $J\subset\{1,2,3\}$, and let $\mu$ satisfy $(T_J)$.
The operators $U(a)$, $a\in\sum_{l\in J}\R\cdot e_l$, are then unitary on $L^2(\mu)$.
Let $\mathcal E^J(\mu)$ be the closed subspace of the $\Phi\in L^2(\mu)$ with $U(a)\Phi=\Phi$ for all these $a$,
and let $P^J_\mu$ be the orthogonal projection of $L^2(\mu)$ onto $\mathcal E^J(\mu)$.
The law $\mu$ is \emph{ergodic} under the translations $\tau_a$, $a\in\sum_{l\in J}\R\cdot e_l$, if $\mathcal E^J(\mu)=\C\cdot1$ \cite[Section~8.1.4]{NoncompactQED}.

Let $j\in\{1,2,3\}$, and let $\mu$ satisfy $(T_j)$, $(R_j)$ and $(RP_j)$.
On $L^2(\Sigma_{\{x_j>0\}},\mu)$ the reflected form is positive semidefinite,
and $\mu\left((\Theta_j\Phi)\cdot\Phi\right)\le\langle\Phi,\Phi\rangle_\mu$ for $\Phi\in L^2(\Sigma_{\{x_j>0\}},\mu)$ \cite[Proposition~8.19(i)]{NoncompactQED}.
Let $\mathcal H_j$ be the completion of the quotient of $L^2(\Sigma_{\{x_j>0\}},\mu)$ by the null space $\left\{\Phi:\mu\left((\Theta_j\Phi)\cdot\Phi\right)=0\right\}$,
with the scalar product induced by the reflected form,
and let $[\,\cdot\,]_j\colon L^2(\Sigma_{\{x_j>0\}},\mu)\to\mathcal H_j$ be the canonical map, so that
\[
\left\langle[\Phi]_j,[\Phi']_j\right\rangle=\mu\left((\Theta_j\Phi)\cdot\Phi'\right)\qquad\left(\Phi,\Phi'\in L^2(\Sigma_{\{x_j>0\}},\mu)\right);
\]
put $\Omega_j:=[1]_j$.
There is exactly one self-adjoint operator $\mathsf H_j\ge0$ on $\mathcal H_j$ with
$\exp\left(-t\cdot\mathsf H_j\right)[\Phi]_j=[U(t\cdot e_j)\Phi]_j$ for all $t\ge0$ and $\Phi\in L^2(\Sigma_{\{x_j>0\}},\mu)$ \cite[Proposition~8.19(ii)]{NoncompactQED}.
If $\mu$ satisfies moreover $(T_{123})$, then for every $a$ in $e_j^\perp:=\{a\in\R^3:a_j=0\}$ there is exactly one unitary operator $V_j(a)$ on $\mathcal H_j$
with $V_j(a)[\Phi]_j=[U(a)\Phi]_j$ for all $\Phi\in L^2(\Sigma_{\{x_j>0\}},\mu)$,
and $a\mapsto V_j(a)$ is a strongly continuous unitary representation of $e_j^\perp$ that commutes with $\exp\left(-t\cdot\mathsf H_j\right)$, $t\ge0$ \cite[Proposition~8.19(iii)]{NoncompactQED};
for $i\ne j$, $\mathsf P_i=\mathsf P_i^{(j)}$ is the self-adjoint operator on $\mathcal H_j$ with $V_j(s\cdot e_i)=\exp\left(is\cdot\mathsf P_i\right)$ for all $s\in\R$,
which exists and is unique by Stone's theorem.
The tuple $\left(\mathcal H_j,[\,\cdot\,]_j,\mathsf H_j,V_j,\Omega_j\right)$ is the \emph{reconstruction} of $\mu$ in the direction $j$.

The \emph{nine families of planes} are the families of the planes $\{\ell=c\}$, $c\in\R$, for the nine linear forms $\ell\in\mathcal L$ of Definition~\ref{def:mirrors}:
the planes $\{x_i=c\}$, $i=1,2,3$, and $\{x_i-x_l=c\}$, $\{x_i+x_l=c\}$, $1\le i<l\le3$, with $c\in\R$.
As $\ell(x)=n_\ell\cdot x$, the plane $\{\ell=c\}$ is the plane $\{\hat n_\ell\cdot x=c/|n_\ell|\}$;
the reflection in it is $r_{\hat n_\ell,c/|n_\ell|}$,
and reflection positivity at $\{\ell=c\}$ is reflection positivity at $\{\hat n_\ell\cdot x=c/|n_\ell|\}$, the test forms being supported in $\{\ell>c\}$.
The class $\mathcal C_9$ is the set of the laws $\mu$ that satisfy $(T_{123})$ and,
for every $\ell\in\mathcal L$ and every $c\in\R$, are invariant under the reflection in $\{\ell=c\}$ and reflection positive at $\{\ell=c\}$.
The planes $\{x_j=0\}$ belong to the nine families; hence $\mathcal C_9\subset\mathcal C$.
\end{definition}

\begin{definition}[Two symmetry conditions]\label{def:symmetry-conditions}
Let $\mu$ be a law that satisfies \eqref{cond:soft-hypothesis}.
By \cite[Theorem~8.1(a)]{NoncompactQED}, applied with $J=\{1,2,3\}$,
there is exactly one probability measure $\nu$ on the set of the laws, endowed with the $\sigma$-algebra generated by the maps $\xi\mapsto\xi(B)$, $B\in\Sigma$,
such that $\nu$-almost every law $\xi$ satisfies $(T_{123})$ and is ergodic under the translations of $\R^3$, and
\[
\mu(B)=\int\mathrm{d}\nu(\xi)\,\mu_\xi(B)\qquad(B\in\Sigma),\qquad\mu_\xi:=\xi ;
\]
the laws $\mu_\xi$ are the $\R^3$-ergodic components of $\mu$ \cite[Remark~8.2]{NoncompactQED}.
The two symmetry conditions on $\mu$ are the following.
\begin{enumerate}
\item\label{it:projections-coincide}
$P^j_\mu=P^{123}_\mu$ for $j=1,2,3$.
\item\label{it:components-cubic}
For $\nu$-almost every $\xi$ the component $\mu_\xi$ is invariant under $\Ocub$, that is, $g_*\mu_\xi=\mu_\xi$ for every $g\in\Ocub$.
\end{enumerate}
\end{definition}

\begin{remark}
Both symmetry conditions are identities between limits of correlations under large translations.
Let $\mu$ satisfy \eqref{cond:soft-hypothesis}, so that $\mu\in\mathcal C$,
and let $\Phi,\Phi'$ range over the countable set of \cite[Lemma~8.12]{NoncompactQED}:
the finite linear combinations, with coefficients in $\mathbb{Q}+i\mathbb{Q}$, of the functions $F\mapsto\exp\left(iF(f)\right)$,
$f$ in the countable space of test forms of \cite[Lemma~8.11(i)]{NoncompactQED}.
For $a\in\R^3$ put $\min_l|a_l|:=\min\{|a_1|,|a_2|,|a_3|\}$.
The limits below exist by \cite[Theorem~8.3(i)]{NoncompactQED}.
By the equivalence of the conditions (2) and (3) of \cite[Theorem~8.8(i)]{NoncompactQED}, applied for $j=1,2,3$,
condition \eqref{it:projections-coincide} of Definition~\ref{def:symmetry-conditions} holds if and only if
\[
\lim_{s\to\infty}\mu\left(\overline{\Phi}\cdot U(s\cdot e_j+v)\Phi'\right)=\lim_{s\to\infty}\mu\left(\overline{\Phi}\cdot U(s\cdot e_j)\Phi'\right)
\]
for $j=1,2,3$, every $v\in\mathbb{Q}^3$ and all $\Phi,\Phi'$ as above, $s$ ranging over $\R$.
By the equivalence of the conditions (1) and (3) of \cite[Theorem~8.5(ii)]{NoncompactQED}, applied to every $g\in\Ocub$,
condition \eqref{it:components-cubic} of Definition~\ref{def:symmetry-conditions} holds if and only if
\[
\lim_{\min_l|a_l|\to\infty}\mu\left(\overline{\Phi}\cdot U(a)\hat g\Phi'\right)=\lim_{\min_l|a_l|\to\infty}\mu\left(\overline{\Phi}\cdot U(a)\Phi'\right)
\]
for every $g\in\Ocub$ and all $\Phi,\Phi'$ as above.
\end{remark}

\section{The chamber operator}\label{sec:chamber}

In this section $\rho\in\R$, $w_0,\dots,w_3>0$, $M\in\R$ and $q>0$ are arbitrary, except where a statement restricts~$\rho$.

\subsection{The fermion graph}\label{sec:fermion-graph}

\begin{remark}[Wall reflections in coordinates]\label{rem:wall-reflections}
For a unit vector $n\in\R^3$ let $r_n$ be the linear reflection $v\mapsto v-2\cdot(v\cdot n)\cdot n$ of $\R^3$ in the plane $n^{\perp}$ orthogonal to $n$;
for instance, $r_{e_1}=\diag(-1,1,1)$.
In coordinates, the reflections $s_0,\dots,s_3$ in the walls $\Pi_0,\dots,\Pi_3$ of $C_0$ (Definition~\ref{def:chambers}) are
\[
\begin{aligned}
s_0(x)&=(x_1,x_2,-x_3), & s_1(x)&=(x_1,x_3,x_2),\\
s_2(x)&=(x_2,x_1,x_3), & s_3(x)&=(1-x_1,x_2,x_3)=e_1+r_{e_1}x .
\end{aligned}
\]
\end{remark}

The affine maps $x\mapsto k+\vec g\,x$ with $k\in\Z^3$ and $\vec g\in\Ocub$ form the group $\Z^3\rtimes\Ocub$, since $\Ocub$ preserves $\Z^3$.
Such a map is written $(k,\vec g)$, and $\vec g$ is its linear part;
thus $(k,\vec g)(k',\vec{g'})=(k+\vec g k',\vec g\vec{g'})$, and the translation $\tau_v$ is $(v,1)$.
Recall that $\Ocub$ consists of the signed permutation matrices (Definition~\ref{def:mirrors}).

\begin{lemma}[The mirror group and its chambers]\label{lem:chambers}
\begin{enumerate}
\item\label{it:chambers-group}
For every $\ell\in\mathcal{L}$ and every $c\in\Adm(\ell)$ the reflection $\theta_{\ell,c}$ maps $\Z^3$ onto $\Z^3$, and $W=\Z^3\rtimes\Ocub$.
\item\label{it:chambers-transitive}
For each $(k,\vec g)\in\Z^3\rtimes\Ocub$ the set $k+\vec g\,C_0$ is open and convex, and it meets no mirror.
Two such sets with different pairs $(k,\vec g)$ are disjoint,
and each point of $\R^3$ that lies on no mirror belongs to one of them.
Consequently the chambers are precisely the sets $k+\vec g\,C_0$,
and the action of $W$ on the set of the chambers is simply transitive.
Moreover, the map $g\mapsto gv_0$ from $W$ to $\mathbb{X}$ is bijective, no mirror contains a fermion site,
and each fermion site has exactly one representation $k+\vec g\,v_0$ with $k\in\Z^3$ and $\vec g\in\Ocub$;
so the maps $\pi$ and $y$ of Definition~\ref{def:fermion-sites} are well defined.
\item\label{it:chambers-generators}
The group $W$ is generated by $s_0$, $s_1$, $s_2$ and $s_3$.
\item\label{it:chambers-walls}
Every mirror is of the form $g\Pi_i$ with $g\in W$ and $i\in\{0,1,2,3\}$.
\end{enumerate}
\end{lemma}

\begin{proof}
\eqref{it:chambers-group}
Let $\ell\in\mathcal{L}$, $c\in\Adm(\ell)$ and $x\in\Z^3$.
The number $\frac{2}{|n_\ell|^2}\cdot(\ell(x)-c)$ is an integer:
for an axis family $\ell=x_\mu$ it equals $2(x_\mu-c)$ with $2c\in\Z$,
and for a diagonal family it equals $\ell(x)-c$ with $c\in\Z$.
As $n_\ell\in\Z^3$, the point $\theta_{\ell,c}(x)=x-\frac{2}{|n_\ell|^2}\cdot(\ell(x)-c)\cdot n_\ell$ lies in $\Z^3$.
Hence $\theta_{\ell,c}(\Z^3)\subset\Z^3$, and equality holds because $\theta_{\ell,c}^2=\mathrm{id}$ (Definition~\ref{def:mirrors}\eqref{it:mirror-involution}).
For every $x\in\R^3$ one has
\[
\theta_{\ell,c}(x)=r_{\hat n_\ell}x+\frac{2c}{|n_\ell|^2}\cdot n_\ell ,
\]
so the linear part of $\theta_{\ell,c}$ is the orthogonal matrix $r_{\hat n_\ell}$.
It maps $\Z^3$ into itself, since $r_{\hat n_\ell}x=\theta_{\ell,c}(x)-\theta_{\ell,c}(0)$ for $x\in\Z^3$, and it is its own inverse;
thus $r_{\hat n_\ell}\in\Ocub$, and $\theta_{\ell,c}$ belongs to $\Z^3\rtimes\Ocub$.
The group $W$, generated by these reflections, is therefore contained in $\Z^3\rtimes\Ocub$.
Conversely, $\theta_{x_\mu,1/2}\,\theta_{x_\mu,0}=\tau_{e_\mu}$ for $\mu=1,2,3$ by Definition~\ref{def:mirrors}\eqref{it:mirror-product},
so that $W$ contains every translation $\tau_k$, $k\in\Z^3$.
For $c=0$ the reflection $\theta_{\ell,0}=r_{\hat n_\ell}$ is linear:
for $\ell=x_\mu$ it changes the sign of the coordinate $x_\mu$,
for $\ell=x_\mu-x_\nu$ it exchanges $x_\mu$ and $x_\nu$,
and for $\ell=x_\mu+x_\nu$ it replaces $(x_\mu,x_\nu)$ by $(-x_\nu,-x_\mu)$.
Every signed permutation matrix is the product of a diagonal matrix with entries $\pm1$ and a permutation matrix, and every permutation is a product of transpositions;
so the group generated by the sign changes of single coordinates and the transpositions of two coordinates is $\Ocub$.
Hence $\Ocub\subset W$, and $W=\Z^3\rtimes\Ocub$.

\eqref{it:chambers-transitive}
Let $x\in\R^3$ lie on no mirror.
The planes $\{x_\mu=c\}$, $c\in\frac12\Z$, are mirrors, so $x_\mu\notin\frac12\Z$ for $\mu=1,2,3$.
Hence there is exactly one $k\in\Z^3$ with $y:=x-k\in\left(-\frac12,\frac12\right)^3$, and $y_\mu\neq0$ for every $\mu$.
For $\mu\neq\nu$ the planes $\{x_\mu-x_\nu=c\}$ and $\{x_\mu+x_\nu=c\}$, $c\in\Z$, are mirrors as well;
so $x_\mu-x_\nu$ and $x_\mu+x_\nu$ are not integers, and neither are $y_\mu-y_\nu$ and $y_\mu+y_\nu$, which differ from them by integers.
In particular $|y_\mu|\neq|y_\nu|$, and the moduli $|y_1|$, $|y_2|$, $|y_3|$ are three different numbers in $\left(0,\frac12\right)$.
Let $\mu_1,\mu_2,\mu_3$ be the indices with $|y_{\mu_1}|>|y_{\mu_2}|>|y_{\mu_3}|$.
For a signed permutation matrix $\vec g$ with $\vec g e_j=\pm e_{\nu_j}$, $j=1,2,3$,
the coordinates of $\vec g^{-1}y$ are $\left(\vec g^{-1}y\right)_j=(\vec g e_j)\cdot y=\pm y_{\nu_j}$, because $\vec g^{-1}$ is the transpose of $\vec g$.
They are positive and decrease with $j$ if and only if $\nu_j=\mu_j$ and $\vec g e_j=\operatorname{sgn}y_{\mu_j}\cdot e_{\mu_j}$ for $j=1,2,3$;
as their moduli are smaller than $\frac12$, this is the condition $\vec g^{-1}y\in C_0$.
So exactly one $\vec g\in\Ocub$ satisfies $\vec g^{-1}y\in C_0$, and $x\in k+\vec g\,C_0$.

The sets $k+\vec g\,C_0$ are pairwise disjoint.
Indeed, a point $x'$ of $k'+\vec{g'}C_0$ has the form $x'=k'+\vec{g'}x''$ with $x''\in C_0$,
and the coordinates of $\vec{g'}x''$ are the numbers $\pm x''_j$, which are non-zero, of pairwise different moduli, and of moduli smaller than $\frac12$.
Two points of $\Z^3$ whose differences with $x'$ both lie in $\left(-\frac12,\frac12\right)^3$ coincide;
so $k'$ is determined by $x'$, and by the preceding paragraph $\vec{g'}$ is the only signed permutation matrix with $\vec{g'}^{-1}(x'-k')\in C_0$.

No mirror meets $C_0$:
for $x\in C_0$ and $\mu\neq\nu$ one has $x_\mu\in\left(0,\frac12\right)$, $x_\mu-x_\nu\in\left(-\frac12,\frac12\right)\setminus\{0\}$ and $x_\mu+x_\nu\in(0,1)$,
and these three sets contain no point of $\frac12\Z$, of $\Z$ and of $\Z$, respectively.

The set of all mirrors is invariant under $W$.
For $v\in\Z^3$, Definition~\ref{def:mirrors}\eqref{it:mirror-conjugation} gives $\tau_v\theta_{\ell,c}\tau_{-v}=\theta_{\ell,c+\ell(v)}$;
comparing the planes of fixed points, one sees that $\tau_v$ maps the mirror $(\ell,c)$ onto the mirror $(\ell,c+\ell(v))$.
For $\vec g\in\Ocub$, the image of the plane $\{\ell=c\}$ under $\vec g$ is the plane $\{\ell\circ\vec g^{-1}=c\}$,
and $\ell\circ\vec g^{-1}$ is the linear form given by the vector $\vec g n_\ell$, as $\vec g^{-1}$ is the transpose of $\vec g$.
A signed permutation matrix maps every vector $e_\mu$ to some $\pm e_{\mu'}$, and every vector $e_\mu\pm e_\nu$, $\mu\neq\nu$, to some $\pm e_{\mu'}\pm e_{\nu'}$ with $\mu'\neq\nu'$;
so $\vec g n_\ell=\pm n_{\ell'}$ for a family $\ell'\in\mathcal{L}$ of the same kind as $\ell$,
and the image is the plane $\{\ell'=\pm c\}$, which is a mirror, since $\Adm(\ell')=\Adm(\ell)=-\Adm(\ell)$.
By \eqref{it:chambers-group} every element of $W$ is the product of a translation $\tau_k$ and an element of $\Ocub$, and the invariance follows.

Consequently no set $gC_0$, $g\in W$, meets a mirror $\Pi$, since otherwise the mirror $g^{-1}\Pi$ would meet $C_0$.
The sets $k+\vec g\,C_0=(k,\vec g)\,C_0$ therefore lie in the complement of the union of the mirrors, and by the first paragraph they cover it.
They are open and convex, hence connected, and they are pairwise disjoint;
so each of them is a connected component of the complement, that is, a chamber (Definition~\ref{def:chambers}), and every chamber is one of them.
For $g_1,g_2\in W$ the element $g_2g_1^{-1}$ maps the chamber $g_1C_0$ onto $g_2C_0$.
If also $g\,g_1C_0=g_2C_0$ for some $g\in W$, then $gg_1=g_2$, because distinct elements of $W$ map $C_0$ onto disjoint sets.
So $W$ acts simply transitively on the chambers.
Finally, $v_0\in C_0$ gives $gv_0\in gC_0$ for every $g\in W$.
If $gv_0=g'v_0$, the chambers $gC_0$ and $g'C_0$ meet, and $g=g'$;
thus $g\mapsto gv_0$ is a bijection of $W$ onto $\mathbb{X}=Wv_0$, and no fermion site lies on a mirror.
In particular, a site $x$ determines the pair $(k,\vec g)$ with $x=k+\vec g\,v_0$,
so that $\pi(x)=k$ and $y(x)=x-\pi(x)=\vec g\,v_0$ are well defined.

\eqref{it:chambers-generators}
By Remark~\ref{rem:wall-reflections}, $s_1$ and $s_2$ are the transpositions of the coordinates $x_2,x_3$ and $x_1,x_2$;
the group generated by them is the group of all permutation matrices.
For a permutation matrix $\vec{g'}$ one has $\vec{g'}s_0\vec{g'}^{-1}=\vec{g'}r_{e_3}\vec{g'}^{-1}=r_{\vec{g'}e_3}$,
the sign change of the coordinate $\mu$ for which $\vec{g'}e_3=e_\mu$, and every $\mu\in\{1,2,3\}$ occurs.
Hence the group generated by $s_0,s_1,s_2$ contains all permutation matrices and all diagonal matrices with entries $\pm1$, and it equals $\Ocub$.
By Remark~\ref{rem:wall-reflections} and $r_{e_1}^2=1$ one has $s_3r_{e_1}=\tau_{e_1}$, where $r_{e_1}\in\Ocub$;
moreover $\vec g\,\tau_{e_1}\vec g^{-1}=\tau_{\vec g e_1}$ for $\vec g\in\Ocub$, and $\vec g e_1$ takes each of the values $\pm e_1,\pm e_2,\pm e_3$.
So the group generated by $s_0,\dots,s_3$ contains $\Ocub$ and every translation $\tau_k$, $k\in\Z^3$, hence $W$ by~\eqref{it:chambers-group};
it is contained in $W$, as $s_i=\theta_{\Pi_i}$.

\eqref{it:chambers-walls}
Let $\Pi=(\ell,c)$ be a mirror.
A bounded subset of $\R^3$ meets only finitely many mirrors,
because each of the nine families is bounded on it and the positions $\Adm(\ell')\subset\frac12\Z$ of a family $\ell'$ form a discrete set.
A mirror other than $\Pi$ meets the plane $\Pi$ in a line or not at all.
A closed disc of positive radius in $\Pi$ meets only finitely many of these lines, and finitely many lines do not cover it;
so some point $x\in\Pi$ lies on no other mirror.
The finitely many mirrors other than $\Pi$ that meet the closed unit ball around $x$ are closed and do not contain $x$;
hence there is an open ball $U$ centred at $x$ that meets no mirror other than $\Pi$.
The open half balls $U\cap\{\ell>c\}$ and $U\cap\{\ell<c\}$ are convex and meet no mirror, so each of them lies in a single chamber.
By~\eqref{it:chambers-transitive} the chamber that contains $U\cap\{\ell>c\}$ is $gC_0$ for some $g\in W$.
Then $x$ lies in the closure of $gC_0$, but not in $gC_0$, because chambers meet no mirror;
so $x$ lies in the boundary of $gC_0$, which is the image under $g$ of the boundary of $C_0$.
The closure of $C_0$ is contained in the closed set $\left\{0\le x_3\le x_2\le x_1\le\frac12\right\}$,
and a point of this set that does not lie in $C_0$ satisfies one of the four inequalities with equality;
hence the boundary of $C_0$ lies in $\Pi_0\cup\Pi_1\cup\Pi_2\cup\Pi_3$, and $x\in g\Pi_i$ for some $i$.
The plane $g\Pi_i$ is a mirror, by the invariance shown in the proof of~\eqref{it:chambers-transitive}, and it contains $x$;
as $\Pi$ is the only mirror through $x$, one has $g\Pi_i=\Pi$.
\end{proof}

\begin{lemma}[The four neighbours]\label{lem:neighbours}
Let $x\in\mathbb{X}$, let $g=(k,\vec g)$ be the element of $W$ with $x=gv_0$ (Lemma~\ref{lem:chambers}\eqref{it:chambers-transitive}),
and put $y:=y(x)$, so that $\pi(x)=k$ and $y=\vec g\,v_0$.
For $j=1,2,3$ let $\mu_j$ be the index with $\vec g e_j=\pm e_{\mu_j}$.
Then $j\mapsto\mu_j$ is a permutation of $\{1,2,3\}$, and
\[
\vec g e_j=\operatorname{sgn}y_{\mu_j}\cdot e_{\mu_j},\qquad |y_{\mu_j}|=v_j\qquad(j=1,2,3);
\]
thus $\mu_1$, $\mu_2$ and $\mu_3$ index the coordinates of $y$ whose moduli are $v_1$, $v_2$ and $v_3$, respectively.
The neighbours $x^{(i)}=gs_iv_0$ are obtained as follows:
$x^{(0)}-k$ arises from $y$ by reversing the sign of $y_{\mu_3}$;
$x^{(1)}-k$ by exchanging the moduli $v_2$, $v_3$ of $y_{\mu_2}$, $y_{\mu_3}$, the signs being kept;
$x^{(2)}-k$ by exchanging the moduli $v_1$, $v_2$ of $y_{\mu_1}$, $y_{\mu_2}$, the signs being kept;
and $x^{(3)}-k$ by replacing $y_{\mu_1}=\operatorname{sgn}y_{\mu_1}\cdot v_1$ with $\operatorname{sgn}y_{\mu_1}\cdot(1-v_1)$.
The displacements $x^{(i)}-x$, the nearest photon sites $\pi\left(x^{(i)}\right)$ and the walls $g\Pi_i$ of the chamber $gC_0$ are
\begin{equation}\label{eq:chamber-neighbours}
\begin{aligned}
x^{(0)}-x&=-2v_3\cdot\operatorname{sgn}y_{\mu_3}\cdot e_{\mu_3}, &\pi\left(x^{(0)}\right)&=k,\\
x^{(1)}-x&=-(v_2-v_3)\cdot\left(\operatorname{sgn}y_{\mu_2}\cdot e_{\mu_2}-\operatorname{sgn}y_{\mu_3}\cdot e_{\mu_3}\right), &\pi\left(x^{(1)}\right)&=k,\\
x^{(2)}-x&=-(v_1-v_2)\cdot\left(\operatorname{sgn}y_{\mu_1}\cdot e_{\mu_1}-\operatorname{sgn}y_{\mu_2}\cdot e_{\mu_2}\right), &\pi\left(x^{(2)}\right)&=k,\\
x^{(3)}-x&=(1-2v_1)\cdot\operatorname{sgn}y_{\mu_1}\cdot e_{\mu_1}, &\pi\left(x^{(3)}\right)&=k+\operatorname{sgn}y_{\mu_1}\cdot e_{\mu_1},\\
g\Pi_0&=k+e_{\mu_3}^{\perp},\\
g\Pi_1&=k+\left(\operatorname{sgn}y_{\mu_2}\cdot e_{\mu_2}-\operatorname{sgn}y_{\mu_3}\cdot e_{\mu_3}\right)^{\perp},\\
g\Pi_2&=k+\left(\operatorname{sgn}y_{\mu_1}\cdot e_{\mu_1}-\operatorname{sgn}y_{\mu_2}\cdot e_{\mu_2}\right)^{\perp},\\
g\Pi_3&=k+\tfrac12\operatorname{sgn}y_{\mu_1}\cdot e_{\mu_1}+e_{\mu_1}^{\perp}.
\end{aligned}
\end{equation}
Hence $g\Pi_0$ is the site plane $(x_{\mu_3},k_{\mu_3})$ of the axis family $x_{\mu_3}$,
$g\Pi_3$ is the link plane $\left(x_{\mu_1},k_{\mu_1}+\tfrac12\operatorname{sgn}y_{\mu_1}\right)$ of the axis family $x_{\mu_1}$,
$g\Pi_1$ and $g\Pi_2$ are mirrors of diagonal families,
and $y\left(x^{(3)}\right)=y-2y_{\mu_1}\cdot e_{\mu_1}$.
With $d_i:=\operatorname{dist}(v_0,\Pi_i)$, that is,
\[
d_0=v_3,\qquad d_1=\frac{v_2-v_3}{\sqrt{2}},\qquad d_2=\frac{v_1-v_2}{\sqrt{2}},\qquad d_3=\frac12-v_1,
\]
one has, for $i=0,1,2,3$,
\[
x^{(i)}-x=-2d_i\cdot\vec g\hat n_i,\qquad \hat d_i(x)=-\vec g\hat n_i,\qquad \left|x^{(i)}-x\right|=2d_i<1 .
\]
Finally, for $i=0,1,2,3$ the vectors $|\Pi_i\cap\overline{C_0}|\cdot\hat n_i$ are $\frac18e_3$, $\frac18(e_2-e_3)$, $\frac18(e_1-e_2)$, $-\frac18e_1$,
and the numbers $|\Pi_i\cap\overline{C_0}|\cdot d_i$ are $\frac18v_3$, $\frac18(v_2-v_3)$, $\frac18(v_1-v_2)$, $\frac18\left(\frac12-v_1\right)$;
hence
\[
\begin{aligned}
\sum_{i=0}^{3}|\Pi_i\cap\overline{C_0}|\cdot\hat n_i&=0,\\
\frac13\sum_{i=0}^{3}|\Pi_i\cap\overline{C_0}|\cdot d_i&=\frac{1}{24}\left(v_3+(v_2-v_3)+(v_1-v_2)+\frac12-v_1\right)=\frac{1}{48}=|C_0| .
\end{aligned}
\]
\end{lemma}

\begin{proof}
Since $\vec g$ is a signed permutation matrix, $\vec g e_j=\pm e_{\mu_j}$ with a permutation $j\mapsto\mu_j$, and
\[
y=\vec g\,v_0=v_1\cdot\vec g e_1+v_2\cdot\vec g e_2+v_3\cdot\vec g e_3
\]
shows that $y_{\mu_j}=\pm v_j$, with the sign that occurs in $\vec g e_j=\pm e_{\mu_j}$.
As $v_j>0$, this sign is $\operatorname{sgn}y_{\mu_j}$, and $|y_{\mu_j}|=v_j$.

For $i=0,1,2$ the reflection $s_i$ is linear (Remark~\ref{rem:wall-reflections}), so that $gs_i=(k,\vec g s_i)$.
Hence $\pi\left(x^{(i)}\right)=k$ and $x^{(i)}-k=\vec g s_iv_0$, where
\[
\begin{aligned}
\vec g s_0v_0&=v_1\cdot\vec g e_1+v_2\cdot\vec g e_2-v_3\cdot\vec g e_3,\\
\vec g s_1v_0&=v_1\cdot\vec g e_1+v_3\cdot\vec g e_2+v_2\cdot\vec g e_3,\\
\vec g s_2v_0&=v_2\cdot\vec g e_1+v_1\cdot\vec g e_2+v_3\cdot\vec g e_3 .
\end{aligned}
\]
Compared with the expansion of $y$ above, these are the descriptions of $x^{(0)}$, $x^{(1)}$ and $x^{(2)}$ in the statement.
For $i=3$, Remark~\ref{rem:wall-reflections} gives $s_3=(e_1,r_{e_1})$, hence $gs_3=(k+\vec g e_1,\vec g r_{e_1})$, and Lemma~\ref{lem:chambers}\eqref{it:chambers-transitive} yields
\[
\pi\left(x^{(3)}\right)=k+\vec g e_1=k+\operatorname{sgn}y_{\mu_1}\cdot e_{\mu_1},\qquad
y\left(x^{(3)}\right)=\vec g r_{e_1}v_0=y-2v_1\cdot\vec g e_1=y-2y_{\mu_1}\cdot e_{\mu_1}.
\]
So $x^{(3)}-k=\vec g e_1+y-2v_1\cdot\vec g e_1=(1-v_1)\cdot\vec g e_1+v_2\cdot\vec g e_2+v_3\cdot\vec g e_3$, which is the description of $x^{(3)}$.

The coordinates of Remark~\ref{rem:wall-reflections} give
\[
\begin{aligned}
s_0v_0&=v_0-2v_3\cdot e_3, & s_1v_0&=v_0-(v_2-v_3)\cdot(e_2-e_3),\\
s_2v_0&=v_0-(v_1-v_2)\cdot(e_1-e_2), & s_3v_0&=v_0+(1-2v_1)\cdot e_1 ,
\end{aligned}
\]
that is, $s_iv_0=v_0-2d_i\cdot\hat n_i$ with the numbers $d_i$ of the statement and the inward unit normals $\hat n_0=e_3$, $\hat n_1=(e_2-e_3)/\sqrt{2}$, $\hat n_2=(e_1-e_2)/\sqrt{2}$, $\hat n_3=-e_1$ of Definition~\ref{def:chambers}.
The reflection in a plane moves a point by twice its distance from the plane along a unit normal;
since $d_i>0$ by $0<v_3<v_2<v_1<\frac12$, it follows that $d_i=\operatorname{dist}(v_0,\Pi_i)$.
As $g$ is affine with the linear part $\vec g$,
\[
x^{(i)}-x=gs_iv_0-gv_0=\vec g\,(s_iv_0-v_0)=-2d_i\cdot\vec g\hat n_i ,
\]
and as $\vec g$ is orthogonal, $\left|x^{(i)}-x\right|=2d_i$ and $\hat d_i(x)=-\vec g\hat n_i$.
Inserting
\[
\begin{aligned}
\vec g\hat n_0&=\operatorname{sgn}y_{\mu_3}\cdot e_{\mu_3}, & \sqrt{2}\,\vec g\hat n_1&=\operatorname{sgn}y_{\mu_2}\cdot e_{\mu_2}-\operatorname{sgn}y_{\mu_3}\cdot e_{\mu_3},\\
\sqrt{2}\,\vec g\hat n_2&=\operatorname{sgn}y_{\mu_1}\cdot e_{\mu_1}-\operatorname{sgn}y_{\mu_2}\cdot e_{\mu_2}, & \vec g\hat n_3&=-\operatorname{sgn}y_{\mu_1}\cdot e_{\mu_1}
\end{aligned}
\]
one obtains the displacements in the first four rows of~\eqref{eq:chamber-neighbours}; the nearest photon sites in these rows were found above.
From $0<v_3<v_2<v_1<\frac12$ one gets $2d_0=2v_3<1$, $2d_1=\sqrt{2}\,(v_2-v_3)<\sqrt{2}\,v_2<1$, $2d_2=\sqrt{2}\,(v_1-v_2)<\sqrt{2}\,v_1<1$ and $2d_3=1-2v_1<1$.

The walls $\Pi_0$, $\Pi_1$, $\Pi_2$ are the planes $\hat n_0^{\perp}$, $\hat n_1^{\perp}$, $\hat n_2^{\perp}$, and $\Pi_3=\frac12e_1+e_1^{\perp}$.
As $\vec g$ is orthogonal, $\vec g\left(\hat n_i^{\perp}\right)=(\vec g\hat n_i)^{\perp}$, so that
\[
g\Pi_i=k+(\vec g\hat n_i)^{\perp}\quad(i=0,1,2),\qquad g\Pi_3=k+\tfrac12\vec g e_1+(\vec g e_1)^{\perp},
\]
and the values of $\vec g\hat n_i$ and of $\vec g e_1=\operatorname{sgn}y_{\mu_1}\cdot e_{\mu_1}$ give the last four rows of~\eqref{eq:chamber-neighbours}.
Thus $g\Pi_0$ consists of the points whose coordinate $\mu_3$ equals the integer $k_{\mu_3}$,
and $g\Pi_3$ of the points whose coordinate $\mu_1$ equals $k_{\mu_1}+\frac12\operatorname{sgn}y_{\mu_1}\in\Z+\frac12$;
these are the site plane $(x_{\mu_3},k_{\mu_3})$ and the link plane $\left(x_{\mu_1},k_{\mu_1}+\frac12\operatorname{sgn}y_{\mu_1}\right)$.
The planes $g\Pi_1$ and $g\Pi_2$ contain $k$ and are orthogonal to vectors of the form $\pm e_\mu\pm e_\nu$, $\mu\neq\nu$, that is, to $n_{\ell'}$ for diagonal families $\ell'$;
so they are the mirrors $(\ell',\ell'(k))$, as $\ell'(k)\in\Z=\Adm(\ell')$.

By Definition~\ref{def:chambers}, the areas $|\Pi_i\cap\overline{C_0}|$ are $\frac18$, $\frac{\sqrt{2}}{8}$, $\frac{\sqrt{2}}{8}$, $\frac18$ for $i=0,1,2,3$, and $|C_0|=\frac{1}{48}$.
Their products with the normals $\hat n_i$ listed above and with the numbers $d_i$ are the stated vectors and numbers.
The four vectors add up to $\frac18\left(e_3+(e_2-e_3)+(e_1-e_2)-e_1\right)=0$,
and the four numbers telescope to $\frac18\cdot\frac12$, one third of which is $\frac{1}{48}$.
\end{proof}

\begin{remark}[Link phases]\label{rem:link-phases}
Let $A$ be a real or complex one-cochain on $\Z^3$, and let $x$, $k$, $y$ and $\mu_1$ be as in Lemma~\ref{lem:neighbours}.
For a bond $\{x,x'\}$ one has $U_A(x,x')=\exp\left(iqA(\pi(x),\pi(x'))\right)$ (Definition~\ref{def:chamber-operator}).
By~\eqref{eq:chamber-neighbours} the bonds of the types $0$, $1$ and $2$ join $x$ to sites with the same nearest photon site $k$,
and $A(k,k)=0$ (Definition~\ref{def:cochains});
so $U_A\left(x,x^{(0)}\right)=U_A\left(x,x^{(1)}\right)=U_A\left(x,x^{(2)}\right)=1$.
On the bond of type $3$ one has $\pi\left(x^{(3)}\right)=k+\operatorname{sgn}y_{\mu_1}\cdot e_{\mu_1}$, and therefore
\[
U_A\left(x,x^{(3)}\right)=
\begin{cases}
\exp\left(iqA_{\mu_1}(k)\right) & \text{if } y_{\mu_1}>0,\\
\exp\left(-iqA_{\mu_1}(k-e_{\mu_1})\right) & \text{if } y_{\mu_1}<0 .
\end{cases}
\]
If $y_{\mu_1}>0$, this is the phase $\exp(iqA(b))$ of the link $b=(k,\mu_1)$;
if $y_{\mu_1}<0$, it is the inverse of the phase of the link $b=(k-e_{\mu_1},\mu_1)$.
In both cases $\pi(x)$ and $\pi\left(x^{(3)}\right)$ are the two end points of $b$,
and $U_A\left(x,x^{(3)}\right)$ is the phase of $b$ taken from $\pi(x)$ to $\pi\left(x^{(3)}\right)$.
This orientation agrees with that of the hop from $x$ to $x^{(3)}$,
since $x^{(3)}-x=(1-2v_1)\cdot\left(\pi\left(x^{(3)}\right)-\pi(x)\right)$ by~\eqref{eq:chamber-neighbours} and $1-2v_1>0$.
By Proposition~\ref{prop:geometry}\eqref{it:geometry-projection}, exactly eight bonds have end points that $\pi$ maps onto the two end points of $b$;
all of them are of type~$3$, and $\left\{x,x^{(3)}\right\}$ is one of them.
Consequently $D_A$ depends on $A$ only through the link phases $\exp(iqA(b))$, $b$ a link of $\Z^3$.
\end{remark}

\begin{remark}[Independence of the base point]\label{rem:base-point}
By Lemma~\ref{lem:chambers}\eqref{it:chambers-transitive} the map $g\mapsto gv_0$ is a bijection of $W$ onto $\mathbb{X}$,
and the site $x=gv_0$ may be labelled by $g=(k,\vec g)$.
For $i=0,1,2,3$ let $\vec s_i:=1-2\hat n_i\hat n_i^{\top}=r_{\hat n_i}$ be the linear part of $s_i$;
thus $\vec s_i=s_i$ for $i=0,1,2$, and $\vec s_3=r_{e_1}$.
In these labels the fermion graph and the operator are described as follows.
\begin{itemize}
\item
The neighbours of $g$ are the elements $gs_i$, that is, $(k,\vec g\vec s_i)$ for $i=0,1,2$ and $(k+\vec g e_1,\vec g\vec s_3)$ for $i=3$.
\item
One has $\pi(gv_0)=k$, and $\hat d_i(x)=-\vec g\hat n_i$ for $x=gv_0$ by Lemma~\ref{lem:neighbours};
so the hopping matrix of the bond of type $i$ at $x$ is $w\left(x,x^{(i)}\right)=\frac{w_i}{2}\left(\rho+\gamma\cdot(\vec g\hat n_i)\right)$.
\item
For a mirror $\Pi$, the reflection $\theta_\Pi$ sends the site with the label $g$ to the site with the label $\theta_\Pi g$, as $\theta_\Pi(gv_0)=(\theta_\Pi g)v_0$.
\item
Let $\Pi=(\ell,c)$ be a mirror.
The affine function $\ell-c$ has no zero on the connected set $gC_0$, which meets no mirror;
so its sign is constant on $gC_0$, and its value at the point $gv_0$ of $gC_0$ has this sign.
\item
A period lattice $\Gamma$ acts by the left multiplications $g\mapsto\tau_{g'}g$, $g'\in\Gamma$,
so that the sites $\mathbb{X}_{\mathbb{T}}=\mathbb{X}/\Gamma$ of a torus $\mathbb{T}=(\Gamma,\chi)$ correspond to the cosets $\Gamma\backslash W$.
\item
The site with the label $(k,\vec g)$ belongs to the set $\mathbb{X}(B)$ of a box $B$ (Definition~\ref{def:integrands}) if and only if $k+\left[-\frac12,\frac12\right]^3\subset B$.
\end{itemize}
Written in the labels, these descriptions are free of $v_0$.
Consequently the following objects are the same for every choice of $v_0$ in $C_0$:
the operators $D_A$ and $\nabla_A^{*}\nabla_A$ (Definitions~\ref{def:chamber-operator} and~\ref{def:covariant-laplacian}), read as operators on functions on $W$;
the actions $S_{\mathbb{T}}$ and the algebras $\Gr_{\mathbb{T}}$ (Definition~\ref{def:tori});
the reflections $\Theta_\Pi$ (Definition~\ref{def:grassmann-reflections});
the classes $\mathcal{G}(B)$ and $\mathcal{A}(U)$, $U\subset\R^3$ open (Definition~\ref{def:integrands}), and $\mathcal{P}$ (Definition~\ref{def:polynomials});
the algebras $\mathcal{B}(\mathbb{T})$ of Lemma~\ref{lem:bounded-integrands};
the functionals $Z_{\mathbb{T}}$ and $\mathbb{E}_{\mathbb{T}}$ and the laws $\mu_{\mathbb{T}}$ (Definition~\ref{def:functional}).
\end{remark}

The complement of a mirror $\Pi=(\ell,c)$ in $\R^3$ consists of the two open half-spaces $\{\ell>c\}$ and $\{\ell<c\}$;
they are the two \emph{sides} of $\Pi$.
A mirror of a diagonal family is called a \emph{diagonal mirror}.
For distinct points $x,x'\in\R^3$ let $[x,x']$ be the closed segment joining them, and let $(x,x'):=[x,x']\setminus\{x,x'\}$ be the open one.
Geometric statements about a bond $\{x,x'\}$ refer to the segment $[x,x']$:
the bond lies on a side of $\Pi$, meets $\Pi$ or is parallel to a vector if $[x,x']$ does,
and a function takes on the bond the values that it takes on $[x,x']$.
Since half-spaces are convex, a bond lies on a side of $\Pi$ if and only if its two end points lie on that side.
As in Definition~\ref{def:modes}, a bond crosses $\Pi$ if its end points lie on different sides of $\Pi$.
A bond $\{x,x'\}$ lies \emph{over} a link $b$ of $\Z^3$ if $\pi(x)$ and $\pi(x')$ are the two end points of $b$.
The open link $\mathring b$ is the link $b$ without its end points,
and a mirror \emph{cuts} $b$ if it meets $\mathring b$ in exactly one point.

\begin{proposition}[Geometry of the fermion graph]\label{prop:geometry}
\begin{enumerate}
\item\label{it:geometry-midpoint}
Let $g\in W$, $x=gv_0$ and $i\in\{0,1,2,3\}$.
The open segment $\left(x,x^{(i)}\right)$ meets exactly one mirror, namely $g\Pi_i$;
it meets $g\Pi_i$ orthogonally and at its midpoint, and $x^{(i)}=\theta_{g\Pi_i}x$.
\item\label{it:geometry-mirrors}
No mirror contains a fermion site.
Let $\Pi=(\ell,c)$ be a mirror and let $\{x,x'\}$ be a bond that crosses $\Pi$;
then $x'=\theta_\Pi x$, and the bond is parallel to $n_\ell$.
\item\label{it:geometry-projection}
The map $\pi$ is $W$-equivariant: $\pi(gx)=g\pi(x)$ for all $g\in W$ and $x\in\mathbb{X}$.
For every $x\in\mathbb{X}$ one has $\pi\left(x^{(i)}\right)=\pi(x)$ for $i=0,1,2$, and
\[
\pi\left(x^{(3)}\right)=\pi(x)+\operatorname{sgn}\left(y_\mu(x)\right)\cdot e_\mu ,
\]
where $\mu$ is the index with $|y_\mu(x)|=v_1$.
Consequently, if a bond crosses a mirror $\Pi$, its type is $0$ when $\Pi$ is a site plane, $3$ when $\Pi$ is a link plane,
and $1$ or $2$ when $\Pi$ is a diagonal mirror.
The bonds of the types $0$, $1$ and $2$ carry no phase:
$U_A\left(x,x^{(i)}\right)=1$ for all $x\in\mathbb{X}$, all $i\in\{0,1,2\}$ and every real or complex one-cochain $A$ on $\Z^3$.
Let $b$ be a link of $\Z^3$.
Exactly eight bonds lie over $b$, and all of them are of type~$3$.
Let $\Pi$ be a mirror, and assume that $\Pi$ does not cut $b$.
If $\mathring b$ lies on one side of $\Pi$, then every bond over $b$ lies on this side;
if $b$ lies in $\Pi$, which happens only for site planes and diagonal mirrors, then four of the eight bonds over $b$ lie on each side of $\Pi$.
\item\label{it:geometry-graph}
The graph with the vertex set $\mathbb{X}$ and the bonds as edges is connected and bipartite,
and every fermion site is an end point of exactly four bonds, one of each type.
\end{enumerate}
\end{proposition}

\begin{proof}
\emph{Mirrors under $W$.}
By Lemma~\ref{lem:chambers}\eqref{it:chambers-group}, every $g\in W$ is of the form $(k,\vec g)$ with $k\in\Z^3$ and $\vec g\in\Ocub$.
Being orthogonal with integer entries, the matrix $\vec g$ maps $e_1,e_2,e_3$ to the vectors $\pm e_1,\pm e_2,\pm e_3$, in some order and with some signs.
Hence for every $\ell\in\mathcal{L}$ there is $\ell'\in\mathcal{L}$ with $\vec g\, n_\ell=\pm n_{\ell'}$,
and $\ell'$ is an axis family if and only if $\ell$ is one.
Then $n_{\ell'}=\pm\vec g\, n_\ell$ and $\vec g^{\top}\vec g=1$ give $\vec g^{\top}n_{\ell'}=\pm n_\ell$,
hence $\ell'(k+\vec g x)=\ell'(k)\pm\ell(x)$ for $x\in\R^3$, that is,
\begin{equation}\label{eq:chamber-mirror-image}
\ell'\circ g=\pm\ell+\ell'(k),
\end{equation}
with the sign for which $\vec g\, n_\ell=\pm n_{\ell'}$.
So $g$ maps the plane $\{\ell=c\}$ onto the plane $\{\ell'=\pm c+\ell'(k)\}$, and the two sides of the first onto the two sides of the second.
If $(\ell,c)$ is a mirror, then $\pm c+\ell'(k)\in\Adm(\ell')$,
because $\ell'(k)$ is an integer and $\Adm(\ell')=\Adm(\ell)$ is invariant under $c\mapsto-c$ and under the addition of integers;
for an axis family, moreover, $\pm c+\ell'(k)$ is an integer if and only if $c$ is one.
Thus $g$ maps every mirror $\Pi$ onto a mirror $g\Pi$:
site planes onto site planes, link planes onto link planes, diagonal mirrors onto diagonal mirrors.
The same holds for $g^{-1}$, so that $\Pi\mapsto g\Pi$ is a bijection of the set of mirrors.
Finally $g$ is an isometry of $\R^3$, and $g\theta_\Pi g^{-1}=\theta_{g\Pi}$ for every mirror $\Pi$,
since $g\theta_\Pi g^{-1}$ is an isometry other than the identity that fixes every point of the plane $g\Pi$.

\emph{The segment of a bond.}
Item~\eqref{it:geometry-midpoint} is first proved for $g=1$, that is, for $x=v_0$ and $x^{(i)}=s_iv_0$.
The reflection $s_i$ in the wall $\Pi_i$ (Definition~\ref{def:chambers}) carries $v_0$ to its mirror image $s_iv_0$.
Hence the vector $s_iv_0-v_0$ is normal to $\Pi_i$, and the point $\tfrac12(v_0+s_iv_0)$ lies on $\Pi_i$;
this point is thus the foot $f_i$ of the perpendicular from $v_0$ to $\Pi_i$.
Explicitly,
\[
f_0=(v_1,v_2,0),\quad
f_1=\left(v_1,\tfrac{v_2+v_3}{2},\tfrac{v_2+v_3}{2}\right),\quad
f_2=\left(\tfrac{v_1+v_2}{2},\tfrac{v_1+v_2}{2},v_3\right),\quad
f_3=\left(\tfrac12,v_2,v_3\right).
\]
The chamber $C_0$ is the set on which the four affine functions $x_3$, $x_2-x_3$, $x_1-x_2$ and $\tfrac12-x_1$ are positive.
They are positive at $v_0$, and, as $0<v_3<v_2<v_1<\tfrac12$, they are non-negative at $f_i$;
an affine function with these two properties is positive at every point of $[v_0,f_i]$ other than $f_i$.
Hence $[v_0,f_i]\setminus\{f_i\}$ lies in $C_0$,
and its image $[s_iv_0,f_i]\setminus\{f_i\}$ under $s_i$ lies in $s_iC_0$.
By Lemma~\ref{lem:chambers}\eqref{it:chambers-transitive}, $C_0$ and $s_iC_0$ are chambers, and no chamber meets a mirror;
so $f_i$ is the only point of $(v_0,s_iv_0)$ that can lie on a mirror.
A point lies on a mirror of the family $\ell$ if and only if the value of $\ell$ at this point is admissible,
that is, lies in $\tfrac12\Z$ for an axis family and in $\Z$ for a diagonal family (Definition~\ref{def:mirrors}).
At the four points the nine forms take the following values.
\begin{itemize}
\item
At $f_0$: $x_1=v_1$ and $x_2=v_2$ lie in $\left(0,\tfrac12\right)$, and $x_3=0$;
$x_1-x_2\in\left(0,\tfrac12\right)$ and $x_1+x_2\in(0,1)$;
$x_1\pm x_3=v_1$ and $x_2\pm x_3=v_2$ lie in $\left(0,\tfrac12\right)$.
\item
At $f_1$: $x_1=v_1$ and $x_2=x_3=\tfrac12(v_2+v_3)$ lie in $\left(0,\tfrac12\right)$, where $\tfrac12(v_2+v_3)<v_2<v_1$;
$x_2-x_3=0$ and $x_2+x_3=v_2+v_3\in(0,1)$;
$x_1-x_2=x_1-x_3\in\left(0,\tfrac12\right)$ and $x_1+x_2=x_1+x_3\in(0,1)$.
\item
At $f_2$: $x_1=x_2=\tfrac12(v_1+v_2)$ and $x_3=v_3$ lie in $\left(0,\tfrac12\right)$, where $v_3<v_2<\tfrac12(v_1+v_2)$;
$x_1-x_2=0$ and $x_1+x_2=v_1+v_2\in(0,1)$;
$x_1-x_3=x_2-x_3\in\left(0,\tfrac12\right)$ and $x_1+x_3=x_2+x_3\in(0,1)$.
\item
At $f_3$: $x_1=\tfrac12$, while $x_2=v_2$ and $x_3=v_3$ lie in $\left(0,\tfrac12\right)$;
$x_1-x_2=\tfrac12-v_2$ and $x_1-x_3=\tfrac12-v_3$ lie in $\left(0,\tfrac12\right)$, and $x_1+x_2$, $x_1+x_3$ lie in $\left(\tfrac12,1\right)$;
$x_2-x_3\in\left(0,\tfrac12\right)$ and $x_2+x_3\in(0,1)$.
\end{itemize}
In each case exactly one of the nine values is admissible,
namely the value of $x_3$, of $x_2-x_3$, of $x_1-x_2$ and of $x_1$ at $f_0$, $f_1$, $f_2$ and $f_3$ respectively.
Thus $f_i$ lies on $\Pi_i$ and on no other mirror, and \eqref{it:geometry-midpoint} holds for $g=1$.

For an arbitrary $g\in W$ the open segment $\left(x,x^{(i)}\right)=(gv_0,gs_iv_0)$ is the image of $(v_0,s_iv_0)$ under $g$.
As $\Pi\mapsto g\Pi$ is a bijection of the set of mirrors,
a mirror $\Pi$ meets $\left(x,x^{(i)}\right)$ if and only if the mirror $g^{-1}\Pi$ meets $(v_0,s_iv_0)$, that is, if and only if $\Pi=g\Pi_i$.
The isometry $g$ carries $[v_0,s_iv_0]$, which is orthogonal to $\Pi_i$ and has its midpoint $f_i$ on $\Pi_i$, onto $\left[x,x^{(i)}\right]$;
so $\left[x,x^{(i)}\right]$ is orthogonal to $g\Pi_i$ and has its midpoint $gf_i$ on $g\Pi_i$.
Finally $x^{(i)}=(gs_ig^{-1})(gv_0)=\theta_{g\Pi_i}x$.

\emph{Bonds that cross a mirror.}
Each fermion site $gv_0$, $g\in W$, lies in the chamber $gC_0$, because $v_0\in C_0$,
and chambers meet no mirror by Lemma~\ref{lem:chambers}\eqref{it:chambers-transitive};
so no mirror contains a fermion site.
Let $\Pi=(\ell,c)$ be a mirror.
For every $x\in\mathbb{X}$ and every $i\in\{0,1,2,3\}$, the bond $\left\{x,x^{(i)}\right\}$ crosses $\Pi$ if and only if the open segment $\left(x,x^{(i)}\right)$ meets $\Pi$.
Indeed, if the affine function $\ell-c$ has opposite signs at $x$ and at $x^{(i)}$, then it vanishes at a point of $\left(x,x^{(i)}\right)$.
Conversely, if $\left(x,x^{(i)}\right)$ meets $\Pi$, then $x^{(i)}=\theta_\Pi x$ by \eqref{it:geometry-midpoint},
and $\ell\left(x^{(i)}\right)-c=c-\ell(x)$ by Definition~\ref{def:mirrors}\eqref{it:mirror-height};
this number is not $0$, since $x$ does not lie on $\Pi$, so that $\ell-c$ takes values of opposite signs at $x$ and at $x^{(i)}$.

Let $\{x,x'\}$ be a bond that crosses $\Pi$, let $i$ be its type, and let $g\in W$ be the element with $x=gv_0$.
Then $x'=x^{(i)}$, because $\left(x^{(i)}\right)^{(i)}=x$ (Definition~\ref{def:fermion-sites}).
By the equivalence just proved, $(x,x')$ meets $\Pi$,
and \eqref{it:geometry-midpoint} yields $\Pi=g\Pi_i$ and $x'=\theta_\Pi x$.
Moreover $[x,x']$ is orthogonal to $\Pi$, so that $x'-x$ is a multiple of $n_\ell$ and the bond is parallel to $n_\ell$.
This proves \eqref{it:geometry-mirrors}.

\emph{Equivariance of $\pi$.}
Let $g'=(a,\vec{g'})\in W$, and let $x\in\mathbb{X}$.
By Lemma~\ref{lem:chambers}\eqref{it:chambers-transitive} and Definition~\ref{def:fermion-sites},
$x=k+\vec g v_0$ with $k=\pi(x)\in\Z^3$ and $\vec g\in\Ocub$.
Then $g'x=(a+\vec{g'}k)+\vec{g'}\vec g v_0$ with $a+\vec{g'}k\in\Z^3$ and $\vec{g'}\vec g\in\Ocub$,
so that $\pi(g'x)=a+\vec{g'}k=g'\pi(x)$.

\emph{The photon sites of the neighbours.}
Let $x=gv_0$ with $g=(k,\vec g)$; then $\pi(x)=k$ and
$y(x)=\vec g v_0=v_1\cdot\vec g e_1+v_2\cdot\vec g e_2+v_3\cdot\vec g e_3$.
For $j=1,2,3$ let $\mu_j$ be the index with $|y_{\mu_j}(x)|=v_j$.
Since $\vec g e_j\in\{\pm e_1,\pm e_2,\pm e_3\}$ and $v_1>v_2>v_3>0$,
\[
\vec g e_j=\operatorname{sgn}\left(y_{\mu_j}(x)\right)\cdot e_{\mu_j},\qquad j=1,2,3.
\]
The values of $\pi\left(x^{(i)}\right)$ are those of the table~\eqref{eq:chamber-neighbours} in Lemma~\ref{lem:neighbours}; they are obtained as follows.
For $i=0,1,2$ the reflection $s_i$ is linear and belongs to $\Ocub$ (Remark~\ref{rem:wall-reflections}),
so $x^{(i)}=k+\vec g s_iv_0$ and $\pi\left(x^{(i)}\right)=k$.
For $i=3$ one has $s_3=\tau_{e_1}r_{e_1}$, hence $x^{(3)}=(k+\vec g e_1)+\vec g r_{e_1}v_0$ with $\vec g r_{e_1}\in\Ocub$, and
\[
\pi\left(x^{(3)}\right)=k+\vec g e_1=\pi(x)+\operatorname{sgn}\left(y_{\mu_1}(x)\right)\cdot e_{\mu_1}.
\]
This is the asserted formula, since $\mu_1$ is the index $\mu$ with $|y_\mu(x)|=v_1$.

\emph{Types of bonds and kinds of mirrors.}
Apply \eqref{eq:chamber-mirror-image} with the values of $\vec g e_3$ and $\vec g e_1$ just found.
The element $g$ maps the wall $\Pi_0=(x_3,0)$ onto the mirror of the axis family $x_{\mu_3}$ at the position $k_{\mu_3}$, which is a site plane,
and the wall $\Pi_3=\left(x_1,\tfrac12\right)$ onto the mirror of the axis family $x_{\mu_1}$ at the position $k_{\mu_1}+\tfrac12\operatorname{sgn}\left(y_{\mu_1}(x)\right)$, which is a link plane;
the walls $\Pi_1=(x_2-x_3,0)$ and $\Pi_2=(x_1-x_2,0)$ are diagonal mirrors, and so are their images $g\Pi_1$ and $g\Pi_2$.
Let a bond cross a mirror $\Pi$; if $x=gv_0$ is one of its end points and $i$ its type, the bond is $\left\{x,x^{(i)}\right\}$.
By the equivalence shown in the proof of \eqref{it:geometry-mirrors}, the open segment $\left(x,x^{(i)}\right)$ meets $\Pi$, and \eqref{it:geometry-midpoint} gives $\Pi=g\Pi_i$.
Among the four mirrors $g\Pi_0,\dots,g\Pi_3$ only $g\Pi_0$ is a site plane and only $g\Pi_3$ is a link plane, while $g\Pi_1$ and $g\Pi_2$ are diagonal mirrors.
Hence $i=0$ if $\Pi$ is a site plane, $i=3$ if $\Pi$ is a link plane, and $i\in\{1,2\}$ if $\Pi$ is a diagonal mirror.

\emph{Parallel transport.}
Since $\pi\left(x^{(i)}\right)=\pi(x)$ for $i=0,1,2$, the two end points of a bond of type $0$, $1$ or $2$ have the same image under $\pi$.
Hence on such a bond $\{x,x'\}$ one has $U_A(x,x')=\exp\left(iqA(\pi(x),\pi(x))\right)=1$ for every real or complex one-cochain $A$ (Definition~\ref{def:chamber-operator}),
because $A(k,k)=0$ for every $k\in\Z^3$ (Definition~\ref{def:cochains}).
For a bond $\left\{x,x^{(3)}\right\}$ of type~$3$ the formula for $\pi\left(x^{(3)}\right)$ shows that $\pi(x)$ and $\pi\left(x^{(3)}\right)$ are nearest neighbours in $\Z^3$;
let $b$ be the link with these two end points.
Reading $A$ as a function on pairs of neighbours, as in Definition~\ref{def:cochains}, one finds
\[
U_A\left(x,x^{(3)}\right)=\exp\left(iqA\left(\pi(x),\pi\left(x^{(3)}\right)\right)\right)=\exp\left(\pm iqA(b)\right),
\]
with $+$ if $b$ is oriented from $\pi(x)$ to $\pi\left(x^{(3)}\right)$ and $-$ otherwise;
so the factor is the phase $\exp\left(iqA(b)\right)$ of $b$ or its inverse (Remark~\ref{rem:link-phases}).

\emph{Bonds over a link.}
Let $b$ be the link of $\Z^3$ with the end points $k$ and $k+e_\mu$, and let $\nu,\nu'$ be the two indices other than $\mu$.
The end points of a bond over $b$ have different images under $\pi$, so the bond is of type~$3$.
Let $x$ be its end point with $\pi(x)=k$; then the bond is $\left\{x,x^{(3)}\right\}$, and $\pi\left(x^{(3)}\right)=k+e_\mu$.
By the formula for $\pi\left(x^{(3)}\right)$ this means that $\mu_1=\mu$ and $y_\mu(x)=v_1$.
Conversely, let $x=k+y$ with $y_\mu=v_1$ and $\{|y_\nu|,|y_{\nu'}|\}=\{v_2,v_3\}$.
Then $y=\vec g v_0$ for the matrix $\vec g$ with $\vec g e_j=\operatorname{sgn}(y_{\mu_j})\cdot e_{\mu_j}$, $j=1,2,3$, where $|y_{\mu_j}|=v_j$;
it is orthogonal with integer entries, so $\vec g\in\Ocub$ and $(k,\vec g)\in W$ by Lemma~\ref{lem:chambers}\eqref{it:chambers-group}.
Hence $x=(k,\vec g)v_0$ is a fermion site with $\pi(x)=k$ and $\mu_1=\mu$,
so that $\pi\left(x^{(3)}\right)=k+e_\mu$ by the formula for $\pi\left(x^{(3)}\right)$, and $\left\{x,x^{(3)}\right\}$ lies over $b$.
There are $2\cdot2\cdot2=8$ such points $x$: two orders of the moduli $v_2,v_3$, and two signs for each of $y_\nu$ and $y_{\nu'}$.
Distinct points give distinct bonds, because a bond over $b$ has exactly one end point that $\pi$ maps to $k$.
So exactly eight bonds lie over $b$.

Fix one of them, $\left\{x,x^{(3)}\right\}$ with $\pi(x)=k$, write $x=gv_0$ with $g=(k,\vec g)$, and put $y:=y(x)$ and $y_\perp:=y-v_1\cdot e_\mu$.
As $\vec g e_1=e_\mu$, one has $\vec g r_{e_1}v_0=\vec g\left(v_0-2v_1\cdot e_1\right)=y-2v_1\cdot e_\mu$, so that $x^{(3)}=k+e_\mu+y-2v_1\cdot e_\mu$.
Hence
\[
\left[x,x^{(3)}\right]=\left\{k+t\cdot e_\mu+y_\perp:\ v_1\le t\le 1-v_1\right\},
\qquad
\mathring b=\left\{k+t\cdot e_\mu:\ 0<t<1\right\}.
\]

\emph{The bonds over $b$ and the mirrors.}
The link plane $\left(x_\mu,k_\mu+\tfrac12\right)$ meets $\mathring b$ in the midpoint of $b$ and nowhere else, so it cuts $b$.
Let $\Pi=(\ell,c)$ be a mirror different from this link plane.
The plane $\{\ell=c\}$ and its pair of sides do not change when $\ell$ and $c$ are replaced by $-\ell$ and $-c$;
after this replacement, if necessary, $\ell$ is one of the forms $x_\mu$, $x_\nu$, $x_\mu\pm x_\nu$, $x_\nu\pm x_{\nu'}$, where $\nu$ and $\nu'$ may be interchanged.
The moduli $|y_\nu|$ and $|y_{\nu'}|$ are $v_2$ and $v_3$ in some order, and $0<v_3<v_2<v_1<\tfrac12$.
\begin{itemize}
\item
$\ell=x_\mu$.
Here $c\in\tfrac12\Z$ and $c\neq k_\mu+\tfrac12$, so $c\le k_\mu$ or $c\ge k_\mu+1$.
On $\mathring b$ the form $x_\mu$ takes its values in $(k_\mu,k_\mu+1)$,
and on the bond in $[k_\mu+v_1,k_\mu+1-v_1]\subset(k_\mu,k_\mu+1)$, an interval that does not contain $c$.
So $\mathring b$ lies on one side of $\Pi$, and the bond lies on the same side.
\item
$\ell=x_\nu$.
The form $x_\nu$ equals $k_\nu$ on $b$ and $k_\nu+y_\nu$ on the bond, where $0<|y_\nu|\le v_2<\tfrac12$.
If $c\neq k_\nu$, then $|c-k_\nu|\ge\tfrac12>|y_\nu|$, so that $\ell-c$ has the sign of $k_\nu-c$ on $b$ and on the bond:
$\mathring b$ lies on one side of $\Pi$, and the bond lies on the same side.
In the remaining case $c=k_\nu$ the plane $\Pi$ contains $b$, whereas $\ell-c$ is constant on the bond, with the value $y_\nu\neq0$;
the bond is therefore contained in $\{\ell>c\}$ if $y_\nu>0$ and in $\{\ell<c\}$ if $y_\nu<0$.
\item
$\ell=x_\mu\pm x_\nu$.
Here $c\in\Z$.
On $\mathring b$ the form takes the values $k_\mu\pm k_\nu+t$, $0<t<1$, so $c\le k_\mu\pm k_\nu$ or $c\ge k_\mu\pm k_\nu+1$.
On the bond it takes the values $k_\mu\pm k_\nu+(t\pm y_\nu)$ with $v_1\le t\le1-v_1$,
and $0<t\pm y_\nu<1$ because $|y_\nu|<v_1$.
Both sets of values lie in the open interval between the integers $k_\mu\pm k_\nu$ and $k_\mu\pm k_\nu+1$, which does not contain $c$;
so $\mathring b$ lies on one side of $\Pi$, and the bond lies on the same side.
\item
$\ell=x_\nu\pm x_{\nu'}$.
Here $c\in\Z$.
The form equals $k_\nu\pm k_{\nu'}$ on $b$ and $k_\nu\pm k_{\nu'}+(y_\nu\pm y_{\nu'})$ on the bond,
and $0<|y_\nu\pm y_{\nu'}|\le v_2+v_3<1$, the first inequality because $|y_\nu|\neq|y_{\nu'}|$.
If $c\neq k_\nu\pm k_{\nu'}$, then $|c-(k_\nu\pm k_{\nu'})|\ge1$, so that $\mathring b$ lies on one side of $\Pi$ and the bond on the same side.
For $c=k_\nu\pm k_{\nu'}$, finally, $b$ is contained in the plane $\Pi$, and on the bond $\ell-c$ takes the single value $y_\nu\pm y_{\nu'}$, which is not $0$;
its sign decides which of the two sides of $\Pi$ contains the bond.
\end{itemize}
Every mirror other than $\left(x_\mu,k_\mu+\tfrac12\right)$ therefore either has $\mathring b$ on one of its sides or contains $b$,
and in neither case does it meet $\mathring b$ in exactly one point.
So $\left(x_\mu,k_\mu+\tfrac12\right)$ is the only mirror that cuts $b$,
and the mirrors that do not cut $b$ are precisely those of the four cases.
For such a mirror with $\mathring b$ on one side, each of the eight bonds over $b$ lies on that side,
since the bond treated in the four cases is an arbitrary one of them.
Moreover $b$ lies in $\Pi$ only if $\Pi$ is the site plane $\{x_\nu=k_\nu\}$ or the diagonal mirror $\{x_\nu\pm x_{\nu'}=k_\nu\pm k_{\nu'}\}$.
In these two cases no bond over $b$ meets $\Pi$, since $y_\nu\neq0$ and $y_\nu\pm y_{\nu'}\neq0$.
The eight bonds over $b$ correspond to the eight vectors $y_\perp$,
and $y_\perp\mapsto-y_\perp$ is a permutation of this set of vectors, because it keeps $y_\mu=v_1$ and the moduli of $y_\nu$ and $y_{\nu'}$.
It reverses the signs of $y_\nu$ and of $y_\nu\pm y_{\nu'}$, hence it exchanges the bonds on the two sides of $\Pi$;
so four of the eight bonds lie on each side.
This completes the proof of \eqref{it:geometry-projection}.

\emph{The graph.}
Let $x=gv_0\in\mathbb{X}$.
The elements $g,gs_0,\dots,gs_3$ of $W$ are pairwise distinct, as $s_0,\dots,s_3$ are the reflections in four distinct planes,
and $g'\mapsto g'v_0$ is injective on $W$ by Lemma~\ref{lem:chambers}\eqref{it:chambers-transitive};
so the four fermion sites $x^{(i)}=gs_iv_0$ are pairwise distinct and distinct from $x$.
Since $\left(x^{(i)}\right)^{(i)}=gs_is_iv_0=x$, a bond of type $i$ that has $x$ as an end point is $\left\{x,x^{(i)}\right\}$.
Therefore $x$ is an end point of exactly the four bonds $\left\{x,x^{(i)}\right\}$, $i=0,1,2,3$, one of each type.
By Lemma~\ref{lem:chambers}\eqref{it:chambers-generators}, and as $s_i^2=1$,
there are $r\ge0$ and $i_1,\dots,i_r\in\{0,1,2,3\}$ with $g=s_{i_1}s_{i_2}\cdots s_{i_r}$.
In the sequence $v_0,\ s_{i_1}v_0,\ s_{i_1}s_{i_2}v_0,\ \dots,\ gv_0=x$ two consecutive terms have the form $g'v_0$ and $g's_iv_0=(g'v_0)^{(i)}$,
so they are joined by a bond; every point of $\mathbb{X}$ is thus joined to $v_0$ by a path, and the graph is connected.
For bipartiteness, attach to $x=gv_0\in\mathbb{X}$ the number $\det\vec g$, where $\vec g$ is the linear part of $g$;
it depends only on $x$, because $g'\mapsto g'v_0$ is injective.
The linear part of $gs_i$ is $\vec g\,\vec s_i$, and $\det\vec s_i=-1$, as $\vec s_i$ is the linear part of a reflection.
Consequently the numbers attached to the end points $x=gv_0$ and $x^{(i)}=(gs_i)v_0$ of a bond are opposite,
and every bond joins a point of $\mathbb{X}$ with the number $1$ to a point with the number $-1$.
\end{proof}

\subsection{Hopping matrices and reflections}\label{sec:hopping-reflections}

Let $\gamma_\mu$, $\gamma_v=\gamma\cdot v$ and $P^{\pm}_n$ be the spin matrices of Definition~\ref{def:spin-matrices}.

\begin{lemma}[Identities of the spin matrices]\label{lem:spin-identities}
Let $u,v\in\R^3$, let $n\in\R^3$ be a unit vector, and write $r_nv=v-2(v\cdot n)\,n$ for the reflection of $v$ in the plane orthogonal to $n$.
\begin{enumerate}
\item[(a)] $\gamma_u\gamma_v+\gamma_v\gamma_u=2\,(u\cdot v)$.
\item[(b)] $\gamma_n\gamma_v\gamma_n=-\gamma_{r_nv}$.
\item[(c)] The matrix $\gamma_n$ is Hermitian, $\operatorname{tr}\gamma_n=0$ and $\gamma_n^2=1$;
each of the eigenvalues $1$ and $-1$ of $\gamma_n$ has multiplicity two, and $\det\gamma_n=1$.
\end{enumerate}
\end{lemma}

\begin{proof}
Expanding $\gamma_u=\sum_\mu u_\mu\gamma_\mu$ and $\gamma_v=\sum_\nu v_\nu\gamma_\nu$ and using $\gamma_\mu\gamma_\nu+\gamma_\nu\gamma_\mu=2\delta_{\mu\nu}$ (Definition~\ref{def:spin-matrices}) one obtains (a).
By (a), $\gamma_n^2=n\cdot n=1$ and $\gamma_v\gamma_n=2(v\cdot n)-\gamma_n\gamma_v$.
Multiplying the latter identity from the left by $\gamma_n$ gives
\[
\gamma_n\gamma_v\gamma_n=2(v\cdot n)\,\gamma_n-\gamma_v=-\gamma_{v-2(v\cdot n)n},
\]
since $v\mapsto\gamma_v$ is linear; this is (b).
The matrix $\gamma_n=\sum_\mu n_\mu\gamma_\mu$ is a real combination of the Hermitian matrices $\gamma_\mu$, hence Hermitian, and $\operatorname{tr}\gamma_\mu=\operatorname{tr}\sigma_\mu-\operatorname{tr}\sigma_\mu=0$ gives $\operatorname{tr}\gamma_n=0$.
A Hermitian matrix with square $1$ has no eigenvalues other than $1$ and $-1$.
The two multiplicities add up to $4$ and their difference is $\operatorname{tr}\gamma_n=0$, so each of them is $2$, and $\det\gamma_n=(-1)^2=1$.
\end{proof}

\begin{lemma}[Hopping matrices]\label{lem:hopping}
Let $\rho$ be real, let $\{x,x'\}$ be a bond of type $i$, and let $\hat d:=(x'-x)/|x'-x|$, so that $w(x,x')=\frac{w_i}{2}\left(\rho-\gamma\cdot\hat d\right)$ (Definition~\ref{def:chamber-operator}).
\begin{enumerate}
\item[(a)] $w(x,x')^*=w(x,x')$, and $w(x,x')$ has the eigenvalues $\frac{w_i}{2}(\rho-1)$ and $\frac{w_i}{2}(\rho+1)$, each of multiplicity two.
\item[(b)] If $\rho^2\neq1$, then
\[
w(x,x')^{-1}=\frac{2}{w_i\,(\rho^2-1)}\left(\rho+\gamma\cdot\hat d\right).
\]
\item[(c)] $w(x',x)=\frac{w_i}{2}\left(\rho+\gamma\cdot\hat d\right)$.
\item[(d)] For every mirror $\Pi$, with reflection $\theta=\theta_\Pi$ and a unit normal $n$,
\[
\gamma_n\,w(x,x')\,\gamma_n=w(\theta x',\theta x).
\]
\item[(e)] Let $\Pi$ be a mirror such that $x$ and $x'$ lie on different sides of $\Pi$.
Then $x'=\theta_\Pi x$, and for the unit normal $n$ of $\Pi$ that points from $x$ to $\theta_\Pi x$
\[
w(x,\theta_\Pi x)=\frac{w_i}{2}\left((1+\rho)P^-_n-(1-\rho)P^+_n\right),\qquad
w(\theta_\Pi x,x)=\frac{w_i}{2}\left((1+\rho)P^+_n-(1-\rho)P^-_n\right).
\]
\end{enumerate}
\end{lemma}

\begin{proof}
By Lemma~\ref{lem:spin-identities}(c), the matrix $\gamma\cdot\hat d$ is Hermitian, with the eigenvalues $1$ and $-1$, each of multiplicity two.
As $\rho$ is real, $w(x,x')$ is Hermitian, and on the eigenspaces of $\gamma\cdot\hat d$ for $1$ and $-1$ it acts as $\frac{w_i}{2}(\rho-1)$ and $\frac{w_i}{2}(\rho+1)$; this is (a).
From $\left(\rho-\gamma\cdot\hat d\right)\left(\rho+\gamma\cdot\hat d\right)=\rho^2-\left(\gamma\cdot\hat d\right)^2=\rho^2-1$ one reads off (b).
The unit vector from $x'$ to $x$ is $-\hat d$, and the bond has the type $i$ also when it is seen from $x'$, because $x=\left(x^{(i)}\right)^{(i)}$ (Definition~\ref{def:fermion-sites}); this gives (c).

For (d), the linear part $\vec\theta$ of $\theta$ is the reflection $r_n$.
Lemma~\ref{lem:spin-identities}(b),(c) yields
\[
\gamma_n\,w(x,x')\,\gamma_n=\frac{w_i}{2}\left(\rho\cdot\gamma_n^2-\gamma_n\gamma_{\hat d}\gamma_n\right)=\frac{w_i}{2}\left(\rho+\gamma_{r_n\hat d}\right).
\]
Write $x=gv_0$ with $g\in W$; then $x'=x^{(i)}=gs_iv_0$.
Since $\theta\in W$, one has $\theta x=(\theta g)v_0$ and $\theta x'=(\theta g)s_iv_0=(\theta x)^{(i)}$, so that $\{\theta x,\theta x'\}$ is a bond of type $i$.
The unit vector from $\theta x'$ to $\theta x$ is $\vec\theta(x-x')/|x-x'|=-r_n\hat d$.
Hence $w(\theta x',\theta x)=\frac{w_i}{2}\left(\rho+\gamma_{r_n\hat d}\right)$, and (d) follows.

In (e) the end points of the bond lie on different sides of $\Pi$; by Proposition~\ref{prop:geometry}\eqref{it:geometry-mirrors} the bond is $\{x,\theta_\Pi x\}$ and is parallel to the normal of $\Pi$.
With $n$ pointing from $x$ to $\theta_\Pi x$ one has $\hat d=n$, and (c) gives $w(x,\theta_\Pi x)=\frac{w_i}{2}(\rho-\gamma_n)$ and $w(\theta_\Pi x,x)=\frac{w_i}{2}(\rho+\gamma_n)$.
By the definition of $P^\pm_n$, $\rho=\rho\cdot\left(P^+_n+P^-_n\right)$ and $\gamma_n=P^+_n-P^-_n$, so that
\[
\rho-\gamma_n=(1+\rho)P^-_n-(1-\rho)P^+_n,\qquad
\rho+\gamma_n=(1+\rho)P^+_n-(1-\rho)P^-_n.
\]
\end{proof}

\begin{definition}[Spinors on a torus]\label{def:torus-spinors}
Let $\mathbb{T}=(\Gamma,\chi)$ be a torus.
A \emph{spinor of $\mathbb{T}$} is a map $\psi\colon\mathbb{X}\to\C^4$ with $\psi(x+g)=\chi(g)\cdot\psi(x)$ for all $x\in\mathbb{X}$ and $g\in\Gamma$.
Let $\mathcal{R}\subset\mathbb{X}$ contain exactly one element of each class of $\mathbb{X}_{\mathbb{T}}=\mathbb{X}/\Gamma$.
The scalar product of two spinors $\psi,\psi'$ of $\mathbb{T}$ is
\[
\langle\psi,\psi'\rangle:=\sum_{x\in\mathcal{R}}\psi(x)^*\,\psi'(x).
\]
For a $\Gamma$-periodic one-cochain $A$, the operator $D^{\mathbb{T}}_A$ is the restriction of $D_A$ to the spinors of $\mathbb{T}$.
\end{definition}

For $x\in\mathbb{X}$ and $g\in\Gamma$ the point $x+g$ lies in $\mathbb{X}$, because $\mathbb{X}=Wv_0$ and the translation $\tau_g$ belongs to $W$ (Lemma~\ref{lem:chambers}\eqref{it:chambers-group}).
By Lemma~\ref{lem:chambers}\eqref{it:chambers-group} and \eqref{it:chambers-transitive}, every site is $x=g'v_0$ for exactly one $g'\in W$,
and $g'=\tau_k\circ\vec{g'}$ with $k\in\Z^3$ and the linear part $\vec{g'}\in\Ocub$ of $g'$;
so $x=k+\vec{g'}v_0$, the pair $(k,\vec{g'})$ being determined by $x$.
As $x+g=\left(\tau_{k+g}\circ\vec{g'}\right)v_0$, the translation by $g\in\Gamma$ replaces $k$ by $k+g$ and leaves $\vec{g'}$ unchanged.
Hence $\Gamma$ acts freely on $\mathbb{X}$, every $x\in\mathbb{X}$ is $x_0+g$ with unique $x_0\in\mathcal{R}$ and $g\in\Gamma$, and $\mathbb{X}_{\mathbb{T}}$ is in bijection with $(\Z^3/\Gamma)\times\Ocub$; in particular
\[
\#|\mathbb{X}_{\mathbb{T}}|=48\,|\mathbb{T}|.
\]
A spinor of $\mathbb{T}$ is therefore determined by its values on $\mathcal{R}$, and these values can be prescribed arbitrarily.
Thus the spinors of $\mathbb{T}$ form a complex vector space of dimension $4\cdot\#|\mathbb{X}_{\mathbb{T}}|$, which $\psi\mapsto(\psi(x))_{x\in\mathcal{R}}$ identifies with $(\C^4)^{\mathcal{R}}$.
In this identification $\langle\cdot,\cdot\rangle$ is the standard scalar product, and it does not depend on $\mathcal{R}$, since $\chi(g)\in\{1,-1\}$ and
$\psi(x+g)^*\psi'(x+g)=\chi(g)^2\cdot\psi(x)^*\psi'(x)=\psi(x)^*\psi'(x)$.

Let $A$ be $\Gamma$-periodic, $x\in\mathbb{X}$ and $g\in\Gamma$, and write $x=g'v_0$ with $g'\in W$.
Then $x+g=\tau_gg'v_0$, and the neighbour of type $i$ of $x+g$ is $\tau_gg's_iv_0=x^{(i)}+g$.
The hopping matrix of a bond depends only on its type and its direction, so $w\left(x+g,x^{(i)}+g\right)=w\left(x,x^{(i)}\right)$.
Moreover $\pi(x+g)=\pi(x)+g$ (Definition~\ref{def:fermion-sites}), and the periodicity of $A$ gives $U_A\left(x+g,x^{(i)}+g\right)=U_A\left(x,x^{(i)}\right)$.
It follows that $(D_A\psi)(x+g)=\chi(g)\cdot(D_A\psi)(x)$ for every spinor $\psi$ of $\mathbb{T}$, so that $D^{\mathbb{T}}_A$ is an operator on the spinors of $\mathbb{T}$.

The same computation holds in the Grassmann algebra $\Gr_{\mathbb{T}}$ of Definition~\ref{def:tori}, whose generators satisfy $\psi(x+g)=\chi(g)\cdot\psi(x)$ and $\bar\psi(x+g)=\chi(g)\cdot\bar\psi(x)$.
There $(D_A\psi)(x+g)=\chi(g)\cdot(D_A\psi)(x)$, and as $\chi(g)^2=1$, each summand $\bar\psi(x)(D_A\psi)(x)$ of
\[
S_{\mathbb{T}}(A)=\sum_{x\in\mathcal{R}}\bar\psi(x)\,(D_A\psi)(x)
\]
is unchanged when $x$ is replaced by $x+g$.
So $S_{\mathbb{T}}(A)$ does not depend on $\mathcal{R}$.
In the free generators $\psi(x),\bar\psi(x)$, $x\in\mathcal{R}$, of $\Gr_{\mathbb{T}}$ it is the form $\bar\psi D\psi$ whose matrix $D$ is the matrix of $D^{\mathbb{T}}_A$ in the identification of the spinors of $\mathbb{T}$ with $(\C^4)^{\mathcal{R}}$;
hence $\int_{\mathbb{T}}\exp\left(-S_{\mathbb{T}}(A)\right)=\det D^{\mathbb{T}}_A$ by \eqref{eq:setting-berezin}.
For the named tori of Definition~\ref{def:tori}, a spinor of $C(n_1,n_2,n_3)$ changes its sign under each of the translations by $n_1e_1$, $n_2e_2$ and $n_3e_3$,
and a spinor of $D(a,b,c)$ under each of the translations by $a\,e_-$, $b\,e_+$ and $c\,e_3$;
a spinor of $D'(a,b,c)$ changes its sign under the translations by $a\,e_-$ and $c\,e_3$ and is invariant under the translation by $b\,e_+$.

\begin{lemma}[Sign rules]\label{lem:sign-rules}
Let $\Pi=(\ell,c)$ be a mirror, $\theta:=\theta_\Pi$ and $n:=\hat n_\ell$.
There is exactly one antilinear map $\Theta_\Pi\colon\Gr(\mathbb{X})\to\Gr(\mathbb{X})$ with $\Theta_\Pi(GG')=\Theta_\Pi(G')\,\Theta_\Pi(G)$ for all $G,G'\in\Gr(\mathbb{X})$ that acts on the generators by \eqref{eq:setting-reflection-fields};
it is the reflection of Definition~\ref{def:grassmann-reflections}.
\begin{enumerate}
\item[(a)] $\Theta_\Pi^2=1$.
\item[(b)] If $n$ is replaced by $-n$ in \eqref{eq:setting-reflection-fields}, the map so obtained is $\Theta_\Pi\circ(-1)^{\deg}=(-1)^{\deg}\circ\Theta_\Pi$,
where $(-1)^{\deg}$ is the automorphism of $\Gr(\mathbb{X})$ of Definition~\ref{def:grassmann}.
On $\Gr^{\mathrm{ev}}(\mathbb{X})$ it coincides with $\Theta_\Pi$; there $\Theta_\Pi$ does not depend on the orientation of the normal.
\item[(c)] For every complex $4\times4$ matrix $Y$ and all $x,x'\in\mathbb{X}$,
\begin{equation}\label{eq:setting-reflection-bilinear}
\Theta_\Pi\left(\bar\psi(x)\,Y\psi(x')\right)=\bar\psi(\theta x')\,\gamma_nY^*\gamma_n\,\psi(\theta x).
\end{equation}
\item[(d)] Let $\Pi'$ be any mirror, $\theta':=\theta_{\Pi'}$, and let $n'$ be the upward unit normal of $\Pi'$.
Then $\Theta_\Pi\Theta_{\Pi'}$ is a linear automorphism of $\Gr(\mathbb{X})$, and for $x\in\mathbb{X}$
\begin{equation}\label{eq:setting-two-reflections}
\Theta_\Pi\Theta_{\Pi'}\psi(x)=\gamma_{n'}\gamma_n\,\psi(\theta\theta'x),\qquad
\Theta_\Pi\Theta_{\Pi'}\bar\psi(x)=\bar\psi(\theta\theta'x)\,\gamma_n\gamma_{n'}.
\end{equation}
If $\Pi'=(\ell,c')$ belongs to the family of $\Pi$, then $\Theta_\Pi\Theta_{\Pi'}$ is the translation $\psi(x)\mapsto\psi(x+v)$, $\bar\psi(x)\mapsto\bar\psi(x+v)$ by the vector
\[
v=\frac{2}{|n_\ell|^2}\,(c-c')\,n_\ell .
\]
If $n\cdot n'=0$, then $\Theta_\Pi\Theta_{\Pi'}=(-1)^{\deg}\,\Theta_{\Pi'}\Theta_\Pi$.
\item[(e)] Let $k\geq1$, let $\Pi_1,\dots,\Pi_{2k}$ be mirrors, $\theta_j:=\theta_{\Pi_j}$, and let $n_j$ be the upward unit normal of $\Pi_j$.
Put $U:=\gamma_{n_{2k}}\cdots\gamma_{n_2}\gamma_{n_1}$ and $g:=\theta_1\cdots\theta_{2k}\in W$.
Then for $x\in\mathbb{X}$
\begin{equation}\label{eq:chamber-iteration}
\Theta_{\Pi_1}\cdots\Theta_{\Pi_{2k}}\psi(x)=U\psi(gx),\qquad
\Theta_{\Pi_1}\cdots\Theta_{\Pi_{2k}}\bar\psi(x)=\bar\psi(gx)\,U^{-1}.
\end{equation}
\end{enumerate}
\end{lemma}

\begin{proof}
The algebra $\Gr(\mathbb{X})$ of Definition~\ref{def:grassmann} is the exterior algebra of the complex vector space spanned by the generators $\psi_\alpha(x),\bar\psi_\alpha(x)$, $x\in\mathbb{X}$, $\alpha=1,\dots,4$.
Odd elements of $\Gr(\mathbb{X})$ anticommute with each other, and therefore every odd element has square zero.
By the universal property of the exterior algebra, every assignment of odd elements to the generators extends to exactly one algebra endomorphism of $\Gr(\mathbb{X})$.
The same property, applied to $\Gr(\mathbb{X})$ with the opposite product and the complex conjugate multiplication by scalars, in which odd elements still have square zero, shows that such an assignment also extends to exactly one antilinear map that reverses products.
Since $\theta$ maps $\mathbb{X}=Wv_0$ onto itself, the right sides of \eqref{eq:setting-reflection-fields} are odd elements of $\Gr(\mathbb{X})$; so $\Theta_\Pi$ exists and is unique.
As the generators generate $\Gr(\mathbb{X})$, two algebra endomorphisms, or two antilinear maps that reverse products, are equal as soon as they agree on the generators.
The product of two antilinear maps that reverse products is linear and multiplicative, that is, an algebra endomorphism.
In components, with $\alpha,\alpha'\in\{1,\dots,4\}$,
\[
\Theta_\Pi\psi_\alpha(x)=-\sum_{\alpha'}\bar\psi_{\alpha'}(\theta x)\cdot(\gamma_n)_{\alpha'\alpha},\qquad
\Theta_\Pi\bar\psi_\alpha(x)=-\sum_{\alpha'}(\gamma_n)_{\alpha\alpha'}\cdot\psi_{\alpha'}(\theta x).
\]

\emph{(c).}
As $\Theta_\Pi$ is antilinear and reverses products,
\[
\Theta_\Pi\left(\bar\psi_\alpha(x)\,Y_{\alpha\alpha'}\,\psi_{\alpha'}(x')\right)
=\overline{Y_{\alpha\alpha'}}\cdot\Theta_\Pi\psi_{\alpha'}(x')\;\Theta_\Pi\bar\psi_\alpha(x)
=\overline{Y_{\alpha\alpha'}}\cdot\left(\bar\psi(\theta x')\gamma_n\right)_{\alpha'}\left(\gamma_n\psi(\theta x)\right)_{\alpha},
\]
the two minus signs cancelling each other.
Summation over $\alpha$ and $\alpha'$, with $\overline{Y_{\alpha\alpha'}}=(Y^*)_{\alpha'\alpha}$, gives \eqref{eq:setting-reflection-bilinear}.

\emph{(d).}
The map $\Theta_\Pi\Theta_{\Pi'}$ is an algebra endomorphism.
Apply the component formulas first for $\Pi'$ and then for $\Pi$, and use that $\gamma_{n'}$ is Hermitian (Lemma~\ref{lem:spin-identities}(c)):
\[
\begin{aligned}
\Theta_\Pi\Theta_{\Pi'}\psi_\alpha(x)
&=-\sum_{\alpha'}\overline{(\gamma_{n'})_{\alpha'\alpha}}\cdot\Theta_\Pi\bar\psi_{\alpha'}(\theta'x)
=\sum_{\alpha',\alpha''}(\gamma_{n'})_{\alpha\alpha'}\cdot(\gamma_n)_{\alpha'\alpha''}\cdot\psi_{\alpha''}(\theta\theta'x),\\
\Theta_\Pi\Theta_{\Pi'}\bar\psi_\alpha(x)
&=-\sum_{\alpha'}\overline{(\gamma_{n'})_{\alpha\alpha'}}\cdot\Theta_\Pi\psi_{\alpha'}(\theta'x)
=\sum_{\alpha',\alpha''}\bar\psi_{\alpha''}(\theta\theta'x)\cdot(\gamma_n)_{\alpha''\alpha'}\cdot(\gamma_{n'})_{\alpha'\alpha}.
\end{aligned}
\]
These are the components of \eqref{eq:setting-two-reflections}.

\emph{(a).}
For $\Pi'=\Pi$ one has $\theta\theta=\mathrm{id}$ (Definition~\ref{def:mirrors}\eqref{it:mirror-involution}) and $\gamma_n^2=1$ (Lemma~\ref{lem:spin-identities}(c)).
By \eqref{eq:setting-two-reflections} the algebra endomorphism $\Theta_\Pi^2$ then fixes every generator, so $\Theta_\Pi^2=1$.
In particular $\Theta_\Pi$ is bijective, and $\Theta_\Pi\Theta_{\Pi'}$ is an automorphism, with inverse $\Theta_{\Pi'}\Theta_\Pi$, as asserted in (d).

\emph{The two special cases of (d).}
If $\Pi'=(\ell,c')$, then $n'=n$, so $\gamma_{n'}\gamma_n=\gamma_n^2=1$, while by Definition~\ref{def:mirrors}\eqref{it:mirror-product}, with the roles of $c$ and $c'$ exchanged, $\theta\theta'=\theta_{\ell,c}\theta_{\ell,c'}$ is the translation $\tau_v$ by $v=\frac{2}{|n_\ell|^2}(c-c')\,n_\ell$.
Thus \eqref{eq:setting-two-reflections} reads $\psi(x)\mapsto\psi(x+v)$, $\bar\psi(x)\mapsto\bar\psi(x+v)$.
Let now $n\cdot n'=0$ and $\Pi'=(\ell',c')$.
Then $\gamma_{n'}\gamma_n=-\gamma_n\gamma_{n'}$ by Lemma~\ref{lem:spin-identities}(a).
Moreover $n_\ell\cdot n_{\ell'}=0$ gives $\ell(\theta'x)=\ell(x)$ and $\ell'(\theta x)=\ell'(x)$ for $x\in\R^3$, and the formula for the reflections in Definition~\ref{def:mirrors} yields
\[
\theta\theta'x=x-\frac{2}{|n_{\ell'}|^2}\left(\ell'(x)-c'\right)n_{\ell'}-\frac{2}{|n_\ell|^2}\left(\ell(x)-c\right)n_\ell=\theta'\theta x .
\]
By \eqref{eq:setting-two-reflections}, applied to the pairs $(\Pi,\Pi')$ and $(\Pi',\Pi)$, the automorphisms $\Theta_\Pi\Theta_{\Pi'}$ and $\Theta_{\Pi'}\Theta_\Pi$ differ by the factor $-1$ on every generator.
Since $\Theta_{\Pi'}\Theta_\Pi$ maps generators to odd elements, the algebra endomorphisms $\Theta_\Pi\Theta_{\Pi'}$ and $(-1)^{\deg}\,\Theta_{\Pi'}\Theta_\Pi$ agree on the generators, and so they coincide.

\emph{(b).}
As $v\mapsto\gamma_v$ is linear, $\gamma_{-n}=-\gamma_n$;
so the map formed with $-n$ is the antilinear map that reverses products and sends every generator $\xi$ to $-\Theta_\Pi(\xi)$.
The maps $\Theta_\Pi\circ(-1)^{\deg}$ and $(-1)^{\deg}\circ\Theta_\Pi$ are antilinear, reverse products, and send $\xi$ to $-\Theta_\Pi(\xi)$ as well:
the first because $(-1)^{\deg}\xi=-\xi$, the second because $\Theta_\Pi(\xi)$ is odd.
By uniqueness, the three maps coincide.
On even elements $(-1)^{\deg}$ is the identity.

\emph{(e).}
For $k=1$ the assertion is \eqref{eq:setting-two-reflections} for $\Pi=\Pi_1$, $\Pi'=\Pi_2$.
Let $k\geq2$, and assume \eqref{eq:chamber-iteration} for the $2k-2$ mirrors $\Pi_3,\dots,\Pi_{2k}$, with $U':=\gamma_{n_{2k}}\cdots\gamma_{n_3}$ and $g':=\theta_3\cdots\theta_{2k}$.
The algebra endomorphism $\Theta_{\Pi_1}\Theta_{\Pi_2}$ is linear.
It therefore maps the components of $U'\psi(x')$, $x'\in\mathbb{X}$, to the components of $U'$ times the column of the images of $\psi_1(x'),\dots,\psi_4(x')$, and the components of $\bar\psi(x')\,U'^{-1}$ to those of the row of the images of $\bar\psi_1(x'),\dots,\bar\psi_4(x')$ times $U'^{-1}$.
Hence, by the induction hypothesis and by \eqref{eq:setting-two-reflections} for $\Pi_1,\Pi_2$,
\[
\Theta_{\Pi_1}\cdots\Theta_{\Pi_{2k}}\psi(x)=\Theta_{\Pi_1}\Theta_{\Pi_2}\left(U'\psi(g'x)\right)=U'\gamma_{n_2}\gamma_{n_1}\,\psi(\theta_1\theta_2g'x)=U\psi(gx),
\]
and likewise
$\Theta_{\Pi_1}\cdots\Theta_{\Pi_{2k}}\bar\psi(x)=\bar\psi(\theta_1\theta_2g'x)\,\gamma_{n_1}\gamma_{n_2}U'^{-1}=\bar\psi(gx)\,U^{-1}$, because $\gamma_{n_j}^{-1}=\gamma_{n_j}$.
Finally $g$ belongs to $W$, the group generated by the reflections in the mirrors.
\end{proof}

\begin{lemma}[Reflection of integrands]\label{lem:reflection-classes}
Let $\Pi$ be a mirror, $\theta:=\theta_\Pi$, and let $\vec\theta$ be the linear part of $\theta$.
\begin{enumerate}
\item[(a)] Let $\Gamma$ be a period lattice, and let $\theta^*\colon\Conf(\vec\theta\Gamma)\to\Conf(\Gamma)$, $(A,z)\mapsto(\theta^*A,\theta^*z)$, be the measure preserving bijection of Definition~\ref{def:configurations}.
For a function $f\colon\Conf(\Gamma)\to\C$ let $\Theta_\Pi f$ be the function on $\Conf(\vec\theta\Gamma)$ given by
\[
(\Theta_\Pi f)(A,z):=\overline{f(\theta^*A,\theta^*z)};
\]
for a function $f=\sum_\Psi f_\Psi\cdot\Psi$ on $\Conf(\Gamma)$ with values in $\Gr(\mathbb{X})$, where $\Psi$ ranges over the monomials and the $f_\Psi$ are complex functions, let $\Theta_\Pi f:=\sum_\Psi(\Theta_\Pi f_\Psi)\cdot\Theta_\Pi(\Psi)$, that is, $(\Theta_\Pi f)(A,z)=\Theta_\Pi\left(f(\theta^*A,\theta^*z)\right)$.
Let $B$ be a box that embeds in $\Gamma$.
Then $\theta B$ is a box that embeds in $\vec\theta\Gamma$, and if $\hat A'$ is a local lift on $\theta B$ at a point $(A,z)$ of $\Conf(\vec\theta\Gamma)$, then $\theta^*\hat A'$ is a local lift on $B$ at the point $(\theta^*A,\theta^*z)$ of $\Conf(\Gamma)$.
Consequently, if $G$ is an integrand that reads in $B$, then $(\Theta_\Pi G)(\hat A')=\Theta_\Pi\left(G(\theta^*\hat A')\right)$ for every such $\hat A'$, where $\Theta_\Pi G$ is the integrand of \eqref{eq:setting-reflection-coefficients}, which reads in $\theta B$ by (b):
the function on $\Conf(\vec\theta\Gamma)$ defined by $\Theta_\Pi G$ is $\Theta_\Pi$ applied to the function on $\Conf(\Gamma)$ defined by $G$.
\item[(b)] Let $B$ be a box.
Then $\Theta_\Pi$ maps the integrands that read in $B$ onto the integrands that read in $\theta B$, $\mathcal{G}(B)$ onto $\mathcal{G}(\theta B)$, and $\mathcal{G}^{\mathrm{ev}}(B)$ onto $\mathcal{G}^{\mathrm{ev}}(\theta B)$.
If $G\in\mathcal{G}(B)$ satisfies the bound of Definition~\ref{def:integrands} with the constants $C$ and $\varepsilon_G$,
then $\Theta_\Pi G$ satisfies it with the constants $C\cdot\sum_\Psi4^{\deg\Psi}$ and $\varepsilon_G$, the sum being over the monomials $\Psi$ of $\Gr(\mathbb{X}(B))$.
For every open set $U\subset\R^3$ it maps $\mathcal{A}(U)$ onto $\mathcal{A}(\theta U)$, and it maps $\mathcal{P}$ onto $\mathcal{P}$.
\end{enumerate}
\end{lemma}

\begin{proof}
\emph{Boxes, cells and sites.}
Let $B=y+[-r,r]^3$.
The linear part $\vec\theta$ lies in $\Ocub$ (Lemma~\ref{lem:chambers}\eqref{it:chambers-group}), so it permutes the coordinates up to signs and maps each of the cubes $[-r,r]^3$, $[-1,1]^3$ and $\left[-\tfrac12,\tfrac12\right]^3$ onto itself; thus $\theta B=\theta y+[-r,r]^3$ is a box.
As $\theta\Z^3=\Z^3$, the map $\theta$ permutes the cells of $\Z^3$ (Definition~\ref{def:cochains}), and a cell lies in $B$ if and only if its image lies in $\theta B$.
Thus $\theta$ maps $B^{(1)}$ and $B^{(2)}$ bijectively onto $(\theta B)^{(1)}$ and $(\theta B)^{(2)}$.
Moreover $\mathbb{X}(\theta B)=\theta\,\mathbb{X}(B)$: for $x\in\mathbb{X}$, since $\pi\theta=\theta\pi$ (Proposition~\ref{prop:geometry}\eqref{it:geometry-projection}),
\[
\pi(\theta x)+\left[-\tfrac12,\tfrac12\right]^3=\theta\left(\pi(x)+\left[-\tfrac12,\tfrac12\right]^3\right),
\]
and the left side is contained in $\theta B$ if and only if $\pi(x)+\left[-\tfrac12,\tfrac12\right]^3$ is contained in $B$.

\emph{Integrands.}
Let $G(\hat A)=\sum_\Psi a_\Psi(\hat A)\cdot\Psi$ be an integrand that reads in $B$, the sum being over the monomials $\Psi$ of $\Gr(\mathbb{X}(B))$.
By \eqref{eq:setting-reflection-coefficients} and the antilinearity of $\Theta_\Pi$, for $\hat A'\in\R^{(\theta B)^{(1)}}$,
\[
(\Theta_\Pi G)(\hat A')=\sum_\Psi\overline{a_\Psi(\theta^*\hat A')}\cdot\Theta_\Pi(\Psi)=\Theta_\Pi\left(G(\theta^*\hat A')\right).
\]
Each $\Theta_\Pi(\Psi)$ lies in $\Gr(\theta\,\mathbb{X}(B))=\Gr(\mathbb{X}(\theta B))$, and $\hat A'\mapsto\theta^*\hat A'$ is linear; so $\Theta_\Pi G$ is a Borel map from $\R^{(\theta B)^{(1)}}$ to $\Gr(\mathbb{X}(\theta B))$.
It remains to check gauge invariance.
Write $\Pi=(\ell,c)$, and let $n:=\hat n_\ell$ be the upward unit normal of $\Pi$, with which $\Theta_\Pi$ acts on the generators by \eqref{eq:setting-reflection-fields}.
The involution $\theta$ maps $\Z^3$ onto itself and exchanges $B$ and $\theta B$;
hence the real functions on the sites of $\Z^3$ contained in $\theta B$ are the functions $\lambda\circ\theta$, with $\lambda$ a real function on the sites contained in $B$.
Fix such a $\lambda$.
For $x\in\mathbb{X}(B)$ the components of $\Theta_\Pi\psi(x)=-\bar\psi(\theta x)\gamma_n$ are linear combinations of $\bar\psi_1(\theta x),\dots,\bar\psi_4(\theta x)$,
and those of $\Theta_\Pi\bar\psi(x)=-\gamma_n\psi(\theta x)$ are linear combinations of $\psi_1(\theta x),\dots,\psi_4(\theta x)$.
Since $\pi(\theta x)=\theta\pi(x)$ (Proposition~\ref{prop:geometry}\eqref{it:geometry-projection}) and $\theta\theta=\mathrm{id}$,
the gauge transformation \eqref{eq:setting-gauge} by $\lambda\circ\theta$ multiplies the former by $\exp\left(-iq\lambda(\pi(x))\right)$ and the latter by $\exp\left(iq\lambda(\pi(x))\right)$;
these are the complex conjugates of the factors that the gauge transformation by $\lambda$ gives to $\psi(x)$ and to $\bar\psi(x)$.
For a monomial $\Psi$, the element $\Theta_\Pi(\Psi)$ is the product of the images of the factors of $\Psi$, taken in the reverse order;
so the transformation by $\lambda\circ\theta$ multiplies $\Theta_\Pi(\Psi)$ by $\overline{\exp\bigl(iq\langle\lambda,\mathfrak{q}(\Psi)\rangle\bigr)}$,
the conjugate of the factor of $\Psi$ under the transformation by $\lambda$.
On the cochains one has $\theta^*\left(\hat A'-d(\lambda\circ\theta)\right)=\theta^*\hat A'-d\lambda$, because $\theta^*$ commutes with $d$ and $\theta^*(\lambda\circ\theta)=\lambda\circ\theta\circ\theta=\lambda$.
The gauge invariance of $G$ is the identity $a_\Psi(\hat A-d\lambda)\cdot\exp\bigl(iq\langle\lambda,\mathfrak{q}(\Psi)\rangle\bigr)=a_\Psi(\hat A)$ for all $\hat A$ and $\Psi$;
its complex conjugate at $\hat A=\theta^*\hat A'$ shows that the transformation by $\lambda\circ\theta$, applied to $\hat A'$ and to the generators, turns $(\Theta_\Pi G)(\hat A')$ into
\[
\sum_\Psi\overline{a_\Psi(\theta^*\hat A'-d\lambda)}\cdot\overline{\exp\bigl(iq\langle\lambda,\mathfrak{q}(\Psi)\rangle\bigr)}\cdot\Theta_\Pi(\Psi)
=\sum_\Psi\overline{a_\Psi(\theta^*\hat A')}\cdot\Theta_\Pi(\Psi)=(\Theta_\Pi G)(\hat A').
\]
So $\Theta_\Pi G$ is an integrand that reads in $\theta B$.
By Lemma~\ref{lem:sign-rules}(a) and $\theta^*\theta^*=(\theta\theta)^*=\mathrm{id}$,
\[
\left(\Theta_\Pi\Theta_\Pi G\right)(\hat A)=\Theta_\Pi\left(\Theta_\Pi\left(G(\theta^*\theta^*\hat A)\right)\right)=G(\hat A),
\]
and $\theta\theta B=B$.
Applied to $\theta B$ in place of $B$, the argument shows that $\Theta_\Pi G'$ reads in $B$ for every integrand $G'$ that reads in $\theta B$, and $G'=\Theta_\Pi(\Theta_\Pi G')$; so every integrand that reads in $\theta B$ is the image of one that reads in $B$.

\emph{Growth.}
Let $G\in\mathcal{G}(B)$, with constants $C$ and $\varepsilon_G$ as in Definition~\ref{def:integrands}.
By Definition~\ref{def:spin-matrices}, $\gamma_n=\diag\left(\sum_\mu n_\mu\sigma_\mu,\,-\sum_\mu n_\mu\sigma_\mu\right)$, and the entries of $\sum_\mu n_\mu\sigma_\mu$ are $\pm n_3$ and $n_1\pm in_2$;
so every entry of $\gamma_n$ is $0$, $\pm n_3$ or $\pm(n_1\pm in_2)$, of modulus at most $|n|=1$.
By \eqref{eq:setting-reflection-fields}, every generator is therefore mapped to a sum of four generators with coefficients of modulus at most one, and $\Theta_\Pi(\Psi)$ is a sum of at most $4^{\deg\Psi}$ terms, each of which is a monomial multiplied by a complex number of modulus at most one.
Hence every coefficient of $(\Theta_\Pi G)(\hat A')$ has modulus at most $\sum_\Psi4^{\deg\Psi}\cdot|a_\Psi(\theta^*\hat A')|$, the sum being over the finitely many monomials of $\Gr(\mathbb{X}(B))$.
Since $\left(d\theta^*\hat A'\right)_\sigma=\left(\theta^*d\hat A'\right)_\sigma=\pm\left(d\hat A'\right)_{\theta\sigma}$ and $\sigma\mapsto\theta\sigma$ is a bijection of $B^{(2)}$ onto $(\theta B)^{(2)}$,
\[
\sum_{\sigma\in B^{(2)}}\left(d\theta^*\hat A'\right)_\sigma^2=\sum_{\sigma\in(\theta B)^{(2)}}\left(d\hat A'\right)_\sigma^2 .
\]
So the coefficients of $\Theta_\Pi G$ satisfy the bound of Definition~\ref{def:integrands} with the same $\varepsilon_G$ and with the constant $C\cdot\sum_\Psi4^{\deg\Psi}$, and $\Theta_\Pi G\in\mathcal{G}(\theta B)$.
The map $\Theta_\Pi$ preserves degrees, so it maps $\mathcal{G}^{\mathrm{ev}}(B)$ into $\mathcal{G}^{\mathrm{ev}}(\theta B)$.
With $\theta B$ in place of $B$ the same holds for the maps from $\mathcal{G}(\theta B)$ and $\mathcal{G}^{\mathrm{ev}}(\theta B)$ into $\mathcal{G}(B)$ and $\mathcal{G}^{\mathrm{ev}}(B)$, and since $\Theta_\Pi\Theta_\Pi$ is the identity, all these maps are onto.

\emph{Local lifts.}
Let $B$ embed in $\Gamma$.
The sets $\theta B+[-1,1]^3+\vec\theta g=\theta\left(B+[-1,1]^3+g\right)$, $g\in\Gamma$, are pairwise disjoint, $\theta$ being bijective; so $\theta B$ embeds in $\vec\theta\Gamma$.
Let $\hat A'$ be a local lift on $\theta B$ at $(A,z)$, and let $F':=dA+\frac{2\pi}{q}\,z$.
The real cochain $\theta^*\hat A'$ is defined on $B^{(1)}$, and $d\theta^*\hat A'=\theta^*d\hat A'=\theta^*F'$ on $B^{(2)}$, because $\theta^*$ commutes with $d$ (Definition~\ref{def:cochains});
here $\theta^*F'=d\theta^*A+\frac{2\pi}{q}\,\theta^*z$ is the curvature of $(\theta^*A,\theta^*z)$.
For a link $b$ of $B$ the image $\theta b$ is a link of $\theta B$, and $\left(\theta^*\hat A'\right)(b)=\pm\hat A'(\theta b)$, $(\theta^*A)(b)=\pm A(\theta b)$ with one and the same sign (Definition~\ref{def:cochains}).
Since $\exp\left(iq\hat A'(\theta b)\right)=\exp\left(iqA(\theta b)\right)$, it follows for either sign that
\[
\exp\left(iq\left(\theta^*\hat A'\right)(b)\right)=\exp\left(iq(\theta^*A)(b)\right),
\]
and $\theta^*\hat A'$ is a local lift on $B$ at $(\theta^*A,\theta^*z)$.
For an integrand $G$ that reads in $B$ the value of $G$ at $(\theta^*A,\theta^*z)$ is therefore $G(\theta^*\hat A')$ (Definition~\ref{def:integrands}), and $(\Theta_\Pi G)(\hat A')=\Theta_\Pi\left(G(\theta^*\hat A')\right)$ was shown above.
This proves (a).

\emph{Admissible products.}
Let $\Phi G_1\cdots G_r$ be one of the products that span $\mathcal{A}(U)$ (Definition~\ref{def:integrands}): $G_j\in\mathcal{G}^{\mathrm{ev}}(B_j)$ with pairwise disjoint boxes $B_j\subset U$, and $\Phi$ a bounded Borel function of finitely many $F_\sigma$, $\sigma\subset U$.
By \eqref{eq:setting-reflection-coefficients}, because $\Theta_\Pi$ reverses products and even elements commute,
\[
\Theta_\Pi\left(\Phi G_1\cdots G_r\right)=(\Theta_\Pi\Phi)\,(\Theta_\Pi G_1)\cdots(\Theta_\Pi G_r),\qquad (\Theta_\Pi\Phi)(F)=\overline{\Phi(\theta^*F)}.
\]
Here $\Theta_\Pi G_j\in\mathcal{G}^{\mathrm{ev}}(\theta B_j)$, the boxes $\theta B_j$ are pairwise disjoint and contained in the open set $\theta U$, and $\Theta_\Pi\Phi$ is a bounded Borel function of the finitely many $F_{\theta\sigma}$, $\theta\sigma\subset\theta U$, since $(\theta^*F)_\sigma=\pm F_{\theta\sigma}$.
As $\Theta_\Pi$ is antilinear, it maps $\mathcal{A}(U)$ into $\mathcal{A}(\theta U)$; it maps onto $\mathcal{A}(\theta U)$, because $\Theta_\Pi\Theta_\Pi$ is the identity and $\theta\theta U=U$.

\emph{The class $\mathcal{P}$.}
Let $\delta\in C^1_{\mathrm{c}}$ be real and $\Psi$ a monomial with $d^*\delta=q\,\mathfrak{q}(\Psi)$, so that $\exp\left(i\langle\delta,\hat A\rangle\right)\cdot\Psi$ is one of the gauge invariant elements that span $\mathcal{P}$ (Definition~\ref{def:polynomials}).
The map $\theta^*$ is orthogonal and $\theta^*\theta^*=\mathrm{id}$, so $\langle\delta,\theta^*\hat A\rangle=\langle\theta^*\delta,\hat A\rangle$, and \eqref{eq:setting-reflection-coefficients} gives
\[
\Theta_\Pi\left(\exp\left(i\langle\delta,\hat A\rangle\right)\cdot\Psi\right)=\exp\left(-i\langle\theta^*\delta,\hat A\rangle\right)\cdot\Theta_\Pi(\Psi).
\]
The element $\Theta_\Pi(\Psi)$ is a linear combination of monomials $\Psi'$ that arise from $\Psi$ by replacing every factor $\psi_\alpha(x)$ by some $\bar\psi_{\alpha'}(\theta x)$ and every factor $\bar\psi_\alpha(x)$ by some $\psi_{\alpha'}(\theta x)$.
As $\pi\theta=\theta\pi$ and $\theta$ is an involution, each such $\Psi'$ has the charge density $\mathfrak{q}(\Psi')=-\mathfrak{q}(\Psi)\circ\theta$.
Being an orthogonal involution that commutes with $d$, the map $\theta^*$ commutes with $d^*$, and on functions on $\Z^3$ it is $f\mapsto f\circ\theta$; hence
\[
d^*\left(-\theta^*\delta\right)=-\theta^*\left(q\,\mathfrak{q}(\Psi)\right)=-q\,\mathfrak{q}(\Psi)\circ\theta=q\,\mathfrak{q}(\Psi').
\]
So every $\exp\left(-i\langle\theta^*\delta,\hat A\rangle\right)\cdot\Psi'$ is a gauge invariant element of the kind that spans $\mathcal{P}$, and $\Theta_\Pi$ maps $\mathcal{P}$ into $\mathcal{P}$; onto, since $\Theta_\Pi\Theta_\Pi$ is the identity.
\end{proof}

\begin{lemma}[Grassmann algebras of reflected tori]\label{lem:reflected-algebras}
Let $\mathbb{T}=(\Gamma,\chi)$ be a torus, let $\Pi$ be a mirror with reflection $\theta:=\theta_\Pi$ and upward unit normal $n$, and let $\theta\mathbb{T}=(\vec\theta\Gamma,\chi\circ\vec\theta^{-1})$ be the image of $\mathbb{T}$ (Definition~\ref{def:tori}).
Let $[\,\cdot\,]_{\mathbb{T}}\colon\Gr(\mathbb{X})\to\Gr_{\mathbb{T}}$ and $[\,\cdot\,]_{\theta\mathbb{T}}\colon\Gr(\mathbb{X})\to\Gr_{\theta\mathbb{T}}$ be the homomorphisms of algebras that send every generator to its class.
There is exactly one antilinear map $\Theta_\Pi\colon\Gr_{\mathbb{T}}\to\Gr_{\theta\mathbb{T}}$ with
\[
\Theta_\Pi[G]_{\mathbb{T}}=\left[\Theta_\Pi G\right]_{\theta\mathbb{T}}\qquad\left(G\in\Gr(\mathbb{X})\right).
\]
It reverses products, acts on the classes of the generators by \eqref{eq:setting-reflection-fields}, and is bijective; its inverse is the map $\Gr_{\theta\mathbb{T}}\to\Gr_{\mathbb{T}}$ obtained in the same way from the torus $\theta\mathbb{T}$, whose image under $\theta$ is $\mathbb{T}$.
In particular, if $\theta\mathbb{T}=\mathbb{T}$, then $\Theta_\Pi$ is an antilinear anti-automorphism of $\Gr_{\mathbb{T}}$ with $\Theta_\Pi^2=1$.
Accordingly, for a function $f$ on $\Conf(\Gamma)$ with values in $\Gr_{\mathbb{T}}$,
\[
(\Theta_\Pi f)(A,z):=\Theta_\Pi\left(f(\theta^*A,\theta^*z)\right)
\]
defines a function on $\Conf(\vec\theta\Gamma)$ with values in $\Gr_{\theta\mathbb{T}}$.
Let $G$ be an integrand that reads in a box $B$ that embeds in $\Gamma$.
As $G(\hat A)$ depends neither on the local lift $\hat A$ on $B$ at a point $(A,z)$ of $\Conf(\Gamma)$ nor on the representative $(A,z)$ of the point (Definition~\ref{def:integrands}),
the class $\left[G(\hat A)\right]_{\mathbb{T}}$ is a function of the point $(A,z)$ of $\Conf(\Gamma)$, with values in $\Gr_{\mathbb{T}}$.
Likewise, $\Theta_\Pi G$ is an integrand that reads in the box $\theta B$, which embeds in $\vec\theta\Gamma$ (Lemma~\ref{lem:reflection-classes}),
and the class $\left[(\Theta_\Pi G)(\hat A')\right]_{\theta\mathbb{T}}$, with $\hat A'$ a local lift on $\theta B$ at a point $(A,z)$ of $\Conf(\vec\theta\Gamma)$, is a function of that point, with values in $\Gr_{\theta\mathbb{T}}$.
The second of these two functions is $\Theta_\Pi$ applied to the first.
\end{lemma}

\begin{proof}
By Definition~\ref{def:tori}, $[\,\cdot\,]_{\mathbb{T}}$ is surjective, and its kernel is the two-sided ideal generated by the components of $\psi(x+g)-\chi(g)\cdot\psi(x)$ and of $\bar\psi(x+g)-\chi(g)\cdot\bar\psi(x)$, $x\in\mathbb{X}$, $g\in\Gamma$; the same holds for $\theta\mathbb{T}$.
Since $\theta(x+g)=\theta x+\vec\theta g$ and $\chi(g)$ is real, \eqref{eq:setting-reflection-fields} gives
\[
\Theta_\Pi\left(\psi(x+g)-\chi(g)\cdot\psi(x)\right)=-\left(\bar\psi(\theta x+\vec\theta g)-\chi(g)\cdot\bar\psi(\theta x)\right)\gamma_n,
\]
and similarly $\Theta_\Pi\left(\bar\psi(x+g)-\chi(g)\cdot\bar\psi(x)\right)=-\gamma_n\left(\psi(\theta x+\vec\theta g)-\chi(g)\cdot\psi(\theta x)\right)$.
Here $\vec\theta g\in\vec\theta\Gamma$ and $(\chi\circ\vec\theta^{-1})(\vec\theta g)=\chi(g)$, so the components of the right sides lie in the kernel of $[\,\cdot\,]_{\theta\mathbb{T}}$.
As $\Theta_\Pi$ is antilinear and $\Theta_\Pi(G'\xi G'')=\Theta_\Pi(G'')\,\Theta_\Pi(\xi)\,\Theta_\Pi(G')$ for $G',\xi,G''\in\Gr(\mathbb{X})$, it maps the two-sided ideal generated by a set of elements $\xi$ into the two-sided ideal generated by their images.
Hence $\Theta_\Pi$ sends the kernel of $[\,\cdot\,]_{\mathbb{T}}$ into the kernel of $[\,\cdot\,]_{\theta\mathbb{T}}$.
Therefore $[\Theta_\Pi G]_{\theta\mathbb{T}}$ depends only on $[G]_{\mathbb{T}}$, and this defines $\Theta_\Pi$ on $\Gr_{\mathbb{T}}$.
The map is unique because $[\,\cdot\,]_{\mathbb{T}}$ is surjective, and it is antilinear, reverses products and acts on the classes of the generators as stated, because $\Theta_\Pi$ has these properties on $\Gr(\mathbb{X})$.
As $\vec\theta^2=1$, the image of $\theta\mathbb{T}$ under $\theta$ is $(\Gamma,\chi)=\mathbb{T}$; let $\Theta'_\Pi\colon\Gr_{\theta\mathbb{T}}\to\Gr_{\mathbb{T}}$ be the map obtained from $\theta\mathbb{T}$.
Lemma~\ref{lem:sign-rules}(a) gives $\Theta'_\Pi\Theta_\Pi[G]_{\mathbb{T}}=[\Theta_\Pi^2G]_{\mathbb{T}}=[G]_{\mathbb{T}}$ and $\Theta_\Pi\Theta'_\Pi[G]_{\theta\mathbb{T}}=[G]_{\theta\mathbb{T}}$, so $\Theta'_\Pi$ is the inverse of $\Theta_\Pi$.
If $\theta\mathbb{T}=\mathbb{T}$, then $\Theta'_\Pi=\Theta_\Pi$, whence $\Theta_\Pi^2=1$.

For the last assertion let $(A,z)$ represent a point of $\Conf(\vec\theta\Gamma)$ and let $\hat A'$ be a local lift on $\theta B$ at $(A,z)$.
By Lemma~\ref{lem:reflection-classes}(a), $\theta^*\hat A'$ is a local lift on $B$ at $(\theta^*A,\theta^*z)$, and $(\Theta_\Pi G)(\hat A')=\Theta_\Pi\left(G(\theta^*\hat A')\right)$.
So the first function of the statement takes the value $\left[G(\theta^*\hat A')\right]_{\mathbb{T}}$ at $(\theta^*A,\theta^*z)$,
and $\Theta_\Pi$ applied to it takes the value $\Theta_\Pi\left[G(\theta^*\hat A')\right]_{\mathbb{T}}$ at $(A,z)$.
The value of the second function at $(A,z)$ is $\left[(\Theta_\Pi G)(\hat A')\right]_{\theta\mathbb{T}}$, and the first part of the lemma gives
\[
\left[(\Theta_\Pi G)(\hat A')\right]_{\theta\mathbb{T}}=\left[\Theta_\Pi\left(G(\theta^*\hat A')\right)\right]_{\theta\mathbb{T}}=\Theta_\Pi\left[G(\theta^*\hat A')\right]_{\mathbb{T}}.
\]
\end{proof}

\subsection{Convergence of the functional}\label{sec:functional-convergence}

For every subgroup $\Gamma$ of $\Z^3$ of rank three, the objects of Definition~\ref{def:configurations} are given by the same formulas as for a period lattice.
The $k$-cells of $\Z^3/\Gamma$ are the $\Gamma$-orbits of $k$-cells of $\Z^3$, and $C^k(\Gamma)$, $C^k(\Gamma;\Z)$ are the real and the integer-valued functions on them;
the norm of $F\in C^2(\Gamma)$ is given by $\|F\|^2=\sum_\sigma F_\sigma^2$, the sum being taken over the plaquettes $\sigma$ of $\Z^3/\Gamma$;
the configuration space $\Conf(\Gamma)$, its measure $\mathrm{d}(A,z)$ and the curvature $F=dA+\frac{2\pi}{q}\cdot z$ are defined as in Definition~\ref{def:configurations}.
Two links $(x,\mu)$ and $(x',\mu)$ of $\Z^3$ lie in one $\Gamma$-orbit if and only if $x-x'\in\Gamma$, and the same holds for two plaquettes $(x;\mu\nu)$ and $(x';\mu\nu)$;
hence $\Z^3/\Gamma$ has $3\cdot[\Z^3:\Gamma]$ links and $3\cdot[\Z^3:\Gamma]$ plaquettes.
For a Borel function $\Phi\ge0$ on $\Conf(\Gamma)$,
\[
\int_{\Conf(\Gamma)}\mathrm{d}(A,z)\,\Phi(A,z)=\sum_{z}\ \int_{[0,2\pi/q)^{3\cdot[\Z^3:\Gamma]}}\mathrm{d}A\,\Phi(A,z),
\]
where $z$ ranges over the closed elements of $C^2(\Gamma;\Z)$ and $\mathrm{d}A$ is the Lebesgue measure in the coordinates $A(b)$, $b$ a link of $\Z^3/\Gamma$.

\begin{lemma}[Bounded integrands]\label{lem:bounded-integrands}
Let $q>0$.
The number
\[
C_q:=\sup_{y\in\R}\ \sum_{k\in\Z}\exp\left(-\frac12\left(y+\frac{2\pi k}{q}\right)^2\right)
\]
is finite, and the following holds.
\begin{enumerate}
\item\label{it:chamber-maxwell}
For every subgroup $\Gamma$ of $\Z^3$ of rank three,
\[
\begin{aligned}
0&<\int_{\Conf(\Gamma)}\mathrm{d}(A,z)\,\exp\left(-\tfrac12\|F\|^2\right)\\
&\le\int_{[0,2\pi/q)^{3\cdot[\Z^3:\Gamma]}}\mathrm{d}A\,\prod_{\sigma}\ \sum_{z_\sigma\in\Z}\exp\left(-\frac12\left((dA)_\sigma+\frac{2\pi}{q}\cdot z_\sigma\right)^2\right)
\le\left(\frac{2\pi}{q}\right)^{3\cdot[\Z^3:\Gamma]}\cdot C_q^{3\cdot[\Z^3:\Gamma]},
\end{aligned}
\]
the product being taken over the plaquettes $\sigma$ of $\Z^3/\Gamma$.
\item\label{it:chamber-bounded}
Let $\rho$ and $M$ be real, $w_0,\dots,w_3>0$, and let $\mathbb{T}=(\Gamma,\chi)$ be a torus.
In a monomial basis of $\Gr_{\mathbb{T}}$, the coefficients of $\exp\left(-S_{\mathbb{T}}(A)\right)$ are polynomials in the link phases $u_b=\exp\left(iqA(b)\right)$ and in their complex conjugates $\overline{u_b}$;
in particular they are bounded Borel functions on $\Conf(\Gamma)$.
Let $\mathcal{B}(\mathbb{T})$ be the algebra of the bounded integrands, that is, of the finite sums $G=\sum_j\Phi_j\cdot\xi_j$ with bounded Borel functions $\Phi_j$ on $\Conf(\Gamma)$ and elements $\xi_j$ of $\Gr_{\mathbb{T}}$.
For every $G\in\mathcal{B}(\mathbb{T})$ the integral
\[
Z_{\mathbb{T}}[G]=\int_{\Conf(\Gamma)}\mathrm{d}(A,z)\,\exp\left(-\tfrac12\|F\|^2\right)\int_{\mathbb{T}}\exp\left(-S_{\mathbb{T}}(A)\right)\cdot G(A,z)
\]
of Definition~\ref{def:functional} converges absolutely, coefficient by coefficient:
each coefficient of $\exp\left(-S_{\mathbb{T}}(A)\right)\cdot G(A,z)$ in a monomial basis of $\Gr_{\mathbb{T}}$, multiplied by $\exp\left(-\frac12\|F\|^2\right)$, is an absolutely integrable function on $\Conf(\Gamma)$.
\end{enumerate}
\end{lemma}

\begin{proof}
\emph{Gaussian sums.}
Let $\varepsilon>0$ and $y\in\R$, and let $k_0$ be the least integer with $y+2\pi k_0/q\ge0$.
In the sum $\sum_{k\in\Z}\exp\left(-\varepsilon\left(y+2\pi k/q\right)^2\right)$ the two terms with $k=k_0$ and $k=k_0-1$ are at most $1$.
For $k>k_0$ let $J_k$ be the interval from $y+2\pi(k-1)/q$ to $y+2\pi k/q$;
it lies in $[0,\infty)$, where $s\mapsto\exp\left(-\varepsilon s^2\right)$ is decreasing.
For $k<k_0-1$ let $J_k$ be the interval from $y+2\pi k/q$ to $y+2\pi(k+1)/q$;
it lies in $(-\infty,0)$, where this function is increasing.
In both cases the term of index $k$ is the minimum of $\exp\left(-\varepsilon s^2\right)$ on $J_k$, an interval of length $2\pi/q$, so that
\[
\exp\left(-\varepsilon\left(y+\frac{2\pi k}{q}\right)^2\right)\le\frac{q}{2\pi}\int_{J_k}\mathrm{d}s\,\exp\left(-\varepsilon s^2\right).
\]
The intervals $J_k$ with $k>k_0$ have pairwise disjoint interiors and lie in $[0,\infty)$, and those with $k<k_0-1$ have pairwise disjoint interiors and lie in $(-\infty,0]$.
Summing over $k$, one obtains for every $\varepsilon>0$ and every $y\in\R$
\begin{equation}\label{eq:chamber-gaussian-sums}
\sum_{k\in\Z}\exp\left(-\varepsilon\left(y+\frac{2\pi k}{q}\right)^2\right)\le2+\frac{q}{2\pi}\int_{\R}\mathrm{d}s\,\exp\left(-\varepsilon s^2\right)=2+\frac{q}{2\pi}\cdot\sqrt{\frac{\pi}{\varepsilon}}.
\end{equation}
For $\varepsilon=\frac12$ the right side is $2+q/\sqrt{2\pi}$; so $C_q\le2+q/\sqrt{2\pi}$ is finite.

\emph{The Maxwell integral.}
By the formula for the measure,
\[
\int_{\Conf(\Gamma)}\mathrm{d}(A,z)\,\exp\left(-\tfrac12\|F\|^2\right)=\sum_{z}\ \int_{[0,2\pi/q)^{3\cdot[\Z^3:\Gamma]}}\mathrm{d}A\,\prod_{\sigma}\exp\left(-\frac12\left((dA)_\sigma+\frac{2\pi}{q}\cdot z_\sigma\right)^2\right),
\]
with $z$ ranging over the closed elements of $C^2(\Gamma;\Z)$ and $\sigma$ over the plaquettes of $\Z^3/\Gamma$.
All terms are non-negative, and the term $z=0$ is the integral of the positive continuous function $\exp\left(-\frac12\|dA\|^2\right)$ over a box of positive volume;
so the left side is positive.
If the condition $dz=0$ is omitted, $z$ ranges over all of $C^2(\Gamma;\Z)=\Z^{3\cdot[\Z^3:\Gamma]}$, and the sum can only grow.
By Tonelli's theorem this larger sum may be taken under the integral,
and at fixed $A$ the sum over $z=(z_\sigma)_\sigma$ of the products over $\sigma$ equals the product over $\sigma$ of the sums over $z_\sigma\in\Z$, all terms being non-negative;
this is the second inequality.
Each of the $3\cdot[\Z^3:\Gamma]$ sums over $z_\sigma$ is at most $C_q$, by the definition of $C_q$ with $y=(dA)_\sigma$,
and the box $[0,2\pi/q)^{3\cdot[\Z^3:\Gamma]}$ has the volume $(2\pi/q)^{3\cdot[\Z^3:\Gamma]}$;
this is the third inequality.

\emph{Bounded integrands.}
Fix a set of representatives of $\mathbb{X}_{\mathbb{T}}$.
The monomials in the generators at the representatives form a basis of $\Gr_{\mathbb{T}}$,
and a generator at a site $x+g$, $g\in\Gamma$, is $\chi(g)=\pm1$ times the generator of the same name at $x$ (Definition~\ref{def:tori}).
By Definitions~\ref{def:chamber-operator} and~\ref{def:tori}, $S_{\mathbb{T}}(A)$ is a linear combination of products $\bar\psi_\alpha(x)\,\psi_\beta(x')$ of such generators,
whose coefficients are, up to signs and values of $\chi$, the parameter $M$ and the entries of the matrices $w(x,x')\cdot U_A(x,x')$ of the bonds;
and $U_A(x,x')$ is $1$, $u_b$ or $\overline{u_b}$ for a link $b$ of $\Z^3/\Gamma$ (Remark~\ref{rem:link-phases}).
Since $S_{\mathbb{T}}(A)$ is a sum of products of two generators and $\Gr_{\mathbb{T}}$ has finitely many generators, the powers $\left(-S_{\mathbb{T}}(A)\right)^n$ vanish for large $n$;
so $\exp\left(-S_{\mathbb{T}}(A)\right)$ is the finite sum $\sum_{n\ge0}\left(-S_{\mathbb{T}}(A)\right)^n/n!$, and its coefficients are polynomials in the $u_b$ and $\overline{u_b}$.
As $|u_b|=1$, they are bounded.
They are continuous in $A$, do not depend on $z$, and do not change under $A\mapsto A+\frac{2\pi}{q}\cdot k$, $k\in C^1(\Gamma;\Z)$;
so they are bounded Borel functions on $\Conf(\Gamma)$.
Let $G=\sum_j\Phi_j\cdot\xi_j\in\mathcal{B}(\mathbb{T})$.
The product of two monomials of $\Gr_{\mathbb{T}}$ is $0$ or, up to a sign, a monomial.
Hence each coefficient of $\exp\left(-S_{\mathbb{T}}(A)\right)\cdot G$ is a finite sum of terms, each of which is, up to a sign, the product of some $\Phi_j$, a coefficient of $\exp\left(-S_{\mathbb{T}}(A)\right)$ and a coefficient of $\xi_j$;
it is a bounded Borel function on $\Conf(\Gamma)$.
By \eqref{it:chamber-maxwell}, the function $\exp\left(-\frac12\|F\|^2\right)$ is integrable on $\Conf(\Gamma)$, and \eqref{it:chamber-bounded} follows.
\end{proof}

\begin{lemma}[Integrands with Gaussian growth]\label{lem:gaussian-growth}
Let $\rho$ and $M$ be real, $w_0,\dots,w_3>0$ and $q>0$.
For $\varepsilon>0$ put
\[
C_q(\varepsilon):=\sup_{y\in\R}\ \sum_{k\in\Z}\exp\left(-\varepsilon\left(y+\frac{2\pi k}{q}\right)^2\right).
\]
Then $C_q(\varepsilon)\le2+\frac{q}{2\pi}\cdot\sqrt{\pi/\varepsilon}$, and $C_q=C_q\left(\frac12\right)$.
Let $\mathbb{T}=(\Gamma,\chi)$ be a torus.
For every $\varepsilon>0$,
\begin{equation}\label{eq:chamber-gaussian-integral}
\int_{\Conf(\Gamma)}\mathrm{d}(A,z)\,\prod_{\sigma}\exp\left(-\varepsilon F_\sigma^2\right)\le\left(\frac{2\pi}{q}\cdot C_q(\varepsilon)\right)^{3\cdot[\Z^3:\Gamma]},
\end{equation}
the product being taken over the plaquettes $\sigma$ of $\Z^3/\Gamma$.
Let $G=\Phi\cdot G_1\cdots G_r$ be a product as in Definition~\ref{def:integrands} that embeds in $\Gamma$, with $G_j\in\mathcal{G}^{\mathrm{ev}}(B_j)$,
let $C_j<\infty$ and $\varepsilon_j>0$ be constants for $G_j$ as in that definition,
and put $\varepsilon_G:=\min\left(\varepsilon_1,\dots,\varepsilon_r,\frac12\right)$.
\begin{enumerate}
\item\label{it:chamber-representatives}
For every plaquette of $\Z^3/\Gamma$ there is at most one pair $(j,\sigma)$ with $1\le j\le r$ and $\sigma\in B_j^{(2)}$ such that $\sigma$ represents that plaquette.
\item\label{it:chamber-majorant}
Write $\exp\left(-S_{\mathbb{T}}(A)\right)\cdot G=\sum_\Psi a_\Psi\cdot\Psi$, the sum being taken over a monomial basis of $\Gr_{\mathbb{T}}$.
There is a constant $C$ such that, at every point of $\Conf(\Gamma)$ and for every $\Psi$,
\[
\begin{aligned}
\left|a_\Psi\right|&\le C\cdot\exp\left(\sum_{j=1}^{r}\left(\frac12-\varepsilon_j\right)\sum_{\sigma\in B_j^{(2)}}F_\sigma^2\right),\\
\exp\left(-\tfrac12\|F\|^2\right)\cdot\left|a_\Psi\right|&\le C\cdot\prod_{\sigma}\exp\left(-\varepsilon_GF_\sigma^2\right),
\end{aligned}
\]
the last product being taken over the plaquettes $\sigma$ of $\Z^3/\Gamma$.
\item\label{it:chamber-convergence}
For every $G'\in\mathcal{A}(\R^3)$ that embeds in $\Gamma$, the integral of Definition~\ref{def:functional} defining $Z_{\mathbb{T}}[G']$ converges absolutely, coefficient by coefficient.
\end{enumerate}
\end{lemma}

\begin{proof}
The inequality for $C_q(\varepsilon)$ is \eqref{eq:chamber-gaussian-sums}, which holds for every $y\in\R$,
and $C_q=C_q\left(\frac12\right)$ by the definitions of the two numbers.

By the formula for the measure, the left side of \eqref{eq:chamber-gaussian-integral} equals
\[
\sum_{z}\ \int_{[0,2\pi/q)^{3\cdot[\Z^3:\Gamma]}}\mathrm{d}A\,\prod_{\sigma}\exp\left(-\varepsilon\left((dA)_\sigma+\frac{2\pi}{q}\cdot z_\sigma\right)^2\right),
\]
the sum being taken over the closed elements $z$ of $C^2(\Gamma;\Z)$.
Omitting the condition $dz=0$ can only increase this sum of non-negative terms.
By Tonelli's theorem the sum over all of $C^2(\Gamma;\Z)$ may then be taken under the integral,
and at fixed $A$ it is the product over $\sigma$ of the sums $\sum_{z_\sigma\in\Z}\exp\left(-\varepsilon\left((dA)_\sigma+2\pi z_\sigma/q\right)^2\right)$.
Each of these $3\cdot[\Z^3:\Gamma]$ sums is at most $C_q(\varepsilon)$, and the box has the volume $(2\pi/q)^{3\cdot[\Z^3:\Gamma]}$;
this proves \eqref{eq:chamber-gaussian-integral}.

To prove \eqref{it:chamber-representatives}, consider two such pairs for one plaquette of $\Z^3/\Gamma$.
Two representatives of a plaquette of $\Z^3/\Gamma$ differ by a translation by an element of $\Gamma$,
so the pairs can be written $(j,\sigma)$ and $(k,\sigma+g)$ with $g\in\Gamma$;
the inclusions $\sigma\subset B_j$ and $\sigma+g\subset B_k$ give $\sigma+g\subset(B_j+g)\cap B_k$.
As $G$ embeds in $\Gamma$, the sets $B_i+g'$ and $B_{i'}+g''$, where $1\le i,i'\le r$ and $g',g''\in\Gamma$, are disjoint whenever $(i,g')\neq(i',g'')$ (Definition~\ref{def:integrands});
since $\sigma+g$ is not empty, $(j,g)=(k,0)$, and the two pairs coincide.

For \eqref{it:chamber-majorant}, let $(A,z)$ represent a point of $\Conf(\Gamma)$, and for each $j$ let $\hat A^{(j)}$ be a local lift on $B_j$ (Definition~\ref{def:integrands}).
Since $d\hat A^{(j)}=F$ on $B_j^{(2)}$, the growth condition in the definition of $\mathcal{G}(B_j)$ bounds every coefficient of $G_j\left(\hat A^{(j)}\right)$ in modulus by
\[
C_j\cdot\exp\left(\left(\frac12-\varepsilon_j\right)\sum_{\sigma\in B_j^{(2)}}F_\sigma^2\right).
\]
The factor $\Phi$, evaluated at $\widetilde F$, is bounded, and the coefficients of $\exp\left(-S_{\mathbb{T}}(A)\right)$ are bounded by Lemma~\ref{lem:bounded-integrands}\eqref{it:chamber-bounded}.
By Definition~\ref{def:integrands}, the value of $G$ at the point represented by $(A,z)$ is
\[
G(A,z)=\Phi(\widetilde F)\cdot r_{\mathbb{T}}\left(G_1\left(\hat A^{(1)}\right)\right)\cdots r_{\mathbb{T}}\left(G_r\left(\hat A^{(r)}\right)\right),
\]
with the quotient homomorphism $r_{\mathbb{T}}$ of Definition~\ref{def:tori}.
It maps every monomial of $\Gr\left(\mathbb{X}(B_j)\right)$ to $0$ or, up to a sign, to a monomial of $\Gr_{\mathbb{T}}$:
each generator goes to $\pm1$ times a generator at a representative,
and a product of generators at representatives vanishes if two of its factors coincide and is otherwise, up to a sign, a monomial of the basis.
In the monomial basis of $\Gr_{\mathbb{T}}$, a coefficient of $\exp\left(-S_{\mathbb{T}}(A)\right)\cdot G(A,z)$ is therefore a sum of terms,
each of which is, up to a sign, the product of $\Phi(\widetilde F)$, of a coefficient of $\exp\left(-S_{\mathbb{T}}(A)\right)$ and of one coefficient of each $G_j\left(\hat A^{(j)}\right)$;
the number of terms depends only on $\mathbb{T}$ and $G$.
This gives the first inequality.
By \eqref{it:chamber-representatives}, distinct pairs $(j,\sigma)$ with $\sigma\in B_j^{(2)}$ represent distinct plaquettes of $\Z^3/\Gamma$.
Therefore
\[
-\frac12\|F\|^2+\sum_{j=1}^{r}\left(\frac12-\varepsilon_j\right)\sum_{\sigma\in B_j^{(2)}}F_\sigma^2
=-\sum_{j=1}^{r}\varepsilon_j\sum_{\sigma\in B_j^{(2)}}F_\sigma^2-\frac12\sum_{\sigma'}F_{\sigma'}^2,
\]
where $\sigma'$ ranges over the plaquettes of $\Z^3/\Gamma$ that have no representative in $B_1\cup\dots\cup B_r$.
Since $\varepsilon_G\le\varepsilon_j$ for every $j$ and $\varepsilon_G\le\frac12$, the right side is at most $-\varepsilon_G\sum_\sigma F_\sigma^2$, the sum being taken over all plaquettes of $\Z^3/\Gamma$;
together with the first inequality this gives the second.

Finally, let $G'\in\mathcal{A}(\R^3)$ embed in $\Gamma$.
Then $G'$ is a finite sum of products that embed in $\Gamma$ (Definition~\ref{def:integrands}),
and the coefficients of $\exp\left(-S_{\mathbb{T}}(A)\right)\cdot G'$ are the sums of the corresponding coefficients for these products.
These coefficients are Borel functions on $\Conf(\Gamma)$,
being finite sums of products, up to signs, of coefficients of $\exp\left(-S_{\mathbb{T}}(A)\right)$, which are Borel by Lemma~\ref{lem:bounded-integrands}\eqref{it:chamber-bounded},
and of coefficients of $G'$, which is a Borel function on $\Conf(\Gamma)$ with values in $\Gr_{\mathbb{T}}$ (Definition~\ref{def:integrands}; Lemma~\ref{lem:borel-lifts}\eqref{it:lift-coefficients}).
The factor $\exp\left(-\frac12\|F\|^2\right)$ is continuous in $A$ for each $z$, hence Borel as well.
For each of the products, \eqref{it:chamber-majorant} and \eqref{eq:chamber-gaussian-integral}, the latter with $\varepsilon$ equal to the number $\varepsilon_G$ of that product, show that each such coefficient, multiplied by $\exp\left(-\frac12\|F\|^2\right)$, is bounded in modulus by an integrable function on $\Conf(\Gamma)$.
Hence every coefficient of $\exp\left(-S_{\mathbb{T}}(A)\right)\cdot G'$, multiplied by $\exp\left(-\frac12\|F\|^2\right)$, is a Borel function on $\Conf(\Gamma)$ dominated by an integrable one, and so it is absolutely integrable on $\Conf(\Gamma)$.
\end{proof}

Let $B=y+[-r,r]^3$ be a box, and put $a_\mu:=\lceil y_\mu-r\rceil$ and $b_\mu:=\lfloor y_\mu+r\rfloor$ for $\mu=1,2,3$.
A cell of $\Z^3$ is the convex hull of its corners, which lie in $\Z^3$;
it is therefore contained in the convex set $B$ if and only if every corner $x$ satisfies $a_\mu\le x_\mu\le b_\mu$ for $\mu=1,2,3$.
So the cells of $\Z^3$ contained in $B$ are, as stated in Definition~\ref{def:integrands}, the cells of $\Z^3$ contained in the integer box
\[
B_{\Z}:=\prod_{\mu=1}^{3}[a_\mu,b_\mu],
\]
which are called the cells of $B_{\Z}$;
in particular, $B^{(1)}$ and $B^{(2)}$ are the sets of the links and of the plaquettes of $B_{\Z}$.
If $B$ contains no site of $\Z^3$, then $a_\mu>b_\mu$ for some $\mu$, and $B^{(1)}$, $B^{(2)}$ and $\mathbb{X}(B)$ are empty.

Let now $B$ contain a site of $\Z^3$, so that $a_\mu\le b_\mu$ for $\mu=1,2,3$.
Let $\mathfrak{T}$ be the set of the links $(x,3)$ of $B_{\Z}$, of the links $(x,2)$ of $B_{\Z}$ with $x_3=a_3$, and of the links $(x,1)$ of $B_{\Z}$ with $x_2=a_2$ and $x_3=a_3$.
A site $x$ of $B_{\Z}$ other than the corner $(a_1,a_2,a_3)$ is the upper end point of exactly one link of $\mathfrak{T}$:
of $(x-e_3,3)$ if $x_3>a_3$, of $(x-e_2,2)$ if $x_3=a_3$ and $x_2>a_2$, and of $(x-e_1,1)$ if $x_3=a_3$, $x_2=a_2$ and $x_1>a_1$;
the corner is the upper end point of no link of $\mathfrak{T}$.
Descending along these links from a site $x$, one arrives at the corner;
the links met on the way form the path of $\mathfrak{T}$ from the corner to $x$.
Thus $\mathfrak{T}$ connects all sites of $B_{\Z}$ and has one link fewer than $B_{\Z}$ has sites:
it is a spanning tree of the graph formed by the sites and the links of $B_{\Z}$.

For $F'\in\R^{B^{(2)}}$ and $t\in\R^{\mathfrak{T}}$ let $\hat A(F',t)$ be the real one-cochain $\hat A$ on $B^{(1)}$ that coincides with $t$ on $\mathfrak{T}$ and satisfies
\[
\begin{aligned}
\hat A_1(x+e_2)&=\hat A_1(x)+\hat A_2(x+e_1)-\hat A_2(x)-F'_{12}(x)&&\text{for }(x;12)\in B^{(2)}\text{ with }x_3=a_3,\\
\hat A_1(x+e_3)&=\hat A_1(x)+\hat A_3(x+e_1)-\hat A_3(x)-F'_{13}(x)&&\text{for }(x;13)\in B^{(2)},\\
\hat A_2(x+e_3)&=\hat A_2(x)+\hat A_3(x+e_2)-\hat A_3(x)-F'_{23}(x)&&\text{for }(x;23)\in B^{(2)},
\end{aligned}
\]
the values at links and plaquettes being written as in Definition~\ref{def:cochains}.
Each line is the equation $(d\hat A)_{\mu\nu}(x)=F'_{\mu\nu}(x)$, solved for the value at one link.
These conditions determine $\hat A$.
The first line, applied to the plaquettes with $x_2=a_2,a_2+1,\dots,b_2-1$ in this order, yields $\hat A_1$ on the links in the plane $\{x_3=a_3\}$
from its values on the line $\{x_2=a_2,\ x_3=a_3\}$ and from the values of $\hat A_2$ in that plane, all of which are values on $\mathfrak{T}$.
The second and the third line, applied to the plaquettes with $x_3=a_3,a_3+1,\dots,b_3-1$ in this order, then yield $\hat A_1$ and $\hat A_2$ at the height $x_3+1$
from their values at the height $x_3$ and from values of $\hat A_3$, which are values on $\mathfrak{T}$.
Every link of $B^{(1)}$ that does not belong to $\mathfrak{T}$ is reached in exactly one of these steps.

For a complex number $u$ with $|u|=1$ let $\arg_qu$ be the number in $[0,2\pi/q)$ with $\exp\left(iq\arg_qu\right)=u$;
the function $\arg_q$ on the unit circle is continuous outside the point $1$, hence Borel.
For a real one-cochain $A$ on $B^{(1)}$ and an integer two-cochain $z$ on $B^{(2)}$ let $F:=dA+\frac{2\pi}{q}\cdot z$ on $B^{(2)}$ and $u_b:=\exp\left(iqA(b)\right)$ for $b\in B^{(1)}$, and put
\[
\hat A^{B}:=\hat A\left(\left(F_\sigma\right)_{\sigma\in B^{(2)}},\left(\arg_qu_b\right)_{b\in\mathfrak{T}}\right).
\]

For an integer $t$ and a function $f$ on $\Z$, the \emph{oriented sum} $\sum_{s=0}^{t-1}f(s)$ is the usual sum if $t\ge0$, the empty sum being $0$, and
\[
\sum_{s=0}^{t-1}f(s):=-\sum_{s=t}^{-1}f(s)\qquad\text{if }t<0 .
\]
It is linear in $f$, and
\[
\sum_{s=0}^{t}f(s)=\sum_{s=0}^{t-1}f(s)+f(t)\qquad\text{for every }t\in\Z ;
\]
indeed, for $t\ge0$ both sides are usual sums, and for $t<0$ the two oriented sums are $-\sum_{s=t+1}^{-1}f(s)$ and $-\sum_{s=t}^{-1}f(s)$, the first being $0$ if $t=-1$.
Both sides of
\[
\sum_{s=0}^{t-1}\left(f(s+1)-f(s)\right)=f(t)-f(0)
\]
vanish for $t=0$, and by the preceding rule both change by $f(t+1)-f(t)$ when $t$ is replaced by $t+1$;
so this identity holds for every $t\in\Z$.

\begin{lemma}[A Borel local lift]\label{lem:borel-lifts}
Let $q>0$.
In \eqref{it:chamber-recursion}, \eqref{it:lift-explicit} and \eqref{it:lift-coefficients}, $B$ is a box that contains a site of $\Z^3$, and $B_{\Z}$, $\mathfrak{T}$, $\hat A(F',t)$ and $\hat A^{B}$ are as above.
\begin{enumerate}
\item\label{it:chamber-recursion}
The map $(F',t)\mapsto\hat A(F',t)$ is linear, and $d\hat A(F',t)=F'$ holds on the plaquettes $(x;13)$ and $(x;23)$ of $B_{\Z}$ and on the plaquettes $(x;12)$ of $B_{\Z}$ with $x_3=a_3$.
If $dF'=0$ on the cubes of $B_{\Z}$, then $d\hat A(F',t)=F'$ on $B^{(2)}$.
Every real one-cochain $\hat A$ on $B^{(1)}$ with $d\hat A=F'$ on $B^{(2)}$ and $\hat A=t$ on $\mathfrak{T}$ is equal to $\hat A(F',t)$.
\item\label{it:lift-explicit}
Let $A$ be a real one-cochain on $B^{(1)}$ and $z$ an integer two-cochain on $B^{(2)}$ with $(dz)(Q)=0$ for every cube $Q$ of $B_{\Z}$, and let $F:=dA+\frac{2\pi}{q}\cdot z$.
Then $\hat A^{B}$ is a local lift on $B$, that is,
\[
d\hat A^{B}=F\ \text{ on }B^{(2)}\qquad\text{and}\qquad\exp\left(iq\hat A^{B}(b)\right)=\exp\left(iqA(b)\right)\ \text{ for every }b\in B^{(1)}.
\]
\item\label{it:lift-coefficients}
Let $G=\sum_\Psi a_\Psi(\hat A)\cdot\Psi$ be an element of $\mathcal{G}(B)$, and let $C$ and $\varepsilon_G$ be constants for $G$ as in Definition~\ref{def:integrands}.
For every integer-valued function $\lambda$ on the sites of $B_{\Z}$, every real one-cochain $\hat A$ on $B^{(1)}$ and every monomial $\Psi$,
\[
a_\Psi\left(\hat A+\frac{2\pi}{q}\cdot d\lambda\right)=a_\Psi(\hat A).
\]
Let $A$ be a real one-cochain on $B^{(1)}$ and $z$ an integer two-cochain on $B^{(2)}$, and let $F:=dA+\frac{2\pi}{q}\cdot z$.
Then $a_\Psi\left(\hat A^{B}\right)$ is the value, at the family $\left((F_\sigma)_{\sigma\in B^{(2)}},(u_b)_{b\in\mathfrak{T}}\right)$, of a Borel function that is the same for all $A$ and $z$.
If $z$ is closed on the cubes of $B_{\Z}$, then $a_\Psi\left(\hat A^{B}\right)=a_\Psi(\hat A)$ for every local lift $\hat A$ on $B$,
that is, for every real one-cochain $\hat A$ on $B^{(1)}$ with $d\hat A=F$ on $B^{(2)}$ and $\exp\left(iq\hat A(b)\right)=\exp\left(iqA(b)\right)$ for $b\in B^{(1)}$;
and
\[
\left|a_\Psi\left(\hat A^{B}\right)\right|\le C\cdot\exp\left(\left(\frac12-\varepsilon_G\right)\sum_{\sigma\in B^{(2)}}F_\sigma^2\right).
\]
\item\label{it:chamber-sums}
Let $\rho$ and $M$ be real, $w_0,\dots,w_3>0$, and let $\mathbb{T}=(\Gamma,\chi)$ be a torus.
Let $(A,z)$ represent a point of $\Conf(\Gamma)$, and let $\widetilde A$ and $\widetilde z$ be the periodic extensions of $A$ and $z$ to $\Z^3$ (Definition~\ref{def:configurations}).
Then $\widetilde z$ takes integer values, and $d\widetilde z=0$.
Let $k$ be the one-cochain on $\Z^3$ whose components are given, with oriented sums, by
\[
\begin{aligned}
k_1(x)&:=0,\\
k_2(x)&:=\sum_{s=0}^{x_1-1}\widetilde z_{12}(s,x_2,x_3),\\
k_3(x)&:=\sum_{s=0}^{x_1-1}\widetilde z_{13}(s,x_2,x_3)+\sum_{s=0}^{x_2-1}\widetilde z_{23}(0,s,x_3),
\end{aligned}
\qquad x\in\Z^3 .
\]
Its values are integers, and $dk=\widetilde z$.
The real one-cochain $\hat A:=\widetilde A+\frac{2\pi}{q}\cdot k$ on $\Z^3$ satisfies
\[
d\hat A=\widetilde F\qquad\text{and}\qquad\exp\left(iq\hat A(b)\right)=\exp\left(iq\widetilde A(b)\right)\ \text{ for every link }b\text{ of }\Z^3,
\]
and its restriction to the links of any box that embeds in $\Gamma$ is a local lift on that box (Definition~\ref{def:integrands}).
Let $H=\Phi\cdot G_1\cdots G_r$ be a product as in Definition~\ref{def:integrands} that embeds in $\Gamma$, with $G_j\in\mathcal{G}^{\mathrm{ev}}(B_j)$.
Then, with the quotient homomorphism $r_{\mathbb{T}}$ of Definition~\ref{def:tori} and for arbitrary local lifts $\hat A^{(j)}$ on the boxes $B_j$,
\[
r_{\mathbb{T}}\left(H(\hat A)\right)=\Phi(\widetilde F)\cdot r_{\mathbb{T}}\left(G_1\left(\hat A^{(1)}\right)\right)\cdots r_{\mathbb{T}}\left(G_r\left(\hat A^{(r)}\right)\right),
\]
and the right side is the element $H(A,z)$ of $\Gr_{\mathbb{T}}$ that Definition~\ref{def:integrands} attaches to $H$ at the point represented by $(A,z)$.
If $G\in\mathcal{A}(\R^3)$ embeds in $\Gamma$ and $G=H_1+\dots+H_n$ with products $H_1,\dots,H_n$ that embed in $\Gamma$, then
\[
H_1(A,z)+\dots+H_n(A,z)=r_{\mathbb{T}}\left(G(\hat A)\right).
\]
Thus the function of $G$ on $\Conf(\Gamma)$, formed as the sum of the functions of the products of such a representation, depends only on $G$, and it is linear in $G$.
For every $G\in\mathcal{A}(\R^3)$ embedding in $\Gamma$, the number $Z_{\mathbb{T}}[G]$ is determined by $G$,
and on the space of all such $G$ the map $G\mapsto Z_{\mathbb{T}}[G]$ is linear.
\end{enumerate}
\end{lemma}

\begin{proof}
\emph{The recursion.}
The three lines of the recursion are linear in $\hat A$ and $F'$, and the initial values are given by $t$;
so $\hat A(F',t)$ depends linearly on $(F',t)$.
Each line states $(d\hat A)_{\mu\nu}(x)=F'_{\mu\nu}(x)$ for the plaquette at which it is imposed, and these are the plaquettes named in \eqref{it:chamber-recursion}.
Let now $dF'=0$ on the cubes of $B_{\Z}$, and write $\hat A:=\hat A(F',t)$.
It is shown by induction on $x_3$ that $(d\hat A)_{12}(x)=F'_{12}(x)$ for every plaquette $(x;12)$ of $B_{\Z}$;
for $x_3=a_3$ this has been seen.
Let $(x;12)$ be a plaquette of $B_{\Z}$ with $x_3<b_3$ at which the equality holds.
Then $x+[0,1]^3$ is a cube of $B_{\Z}$, and $d\hat A$ and $F'$ agree on five of its faces, namely on $(x;23)$, $(x+e_1;23)$, $(x;13)$, $(x+e_2;13)$ and $(x;12)$.
The values of $d(d\hat A)$ and of $dF'$ at this cube vanish, the first because $dd=0$;
both are given by the formula for $dF$ in Definition~\ref{def:cochains}, in which the sixth face $(x+e_3;12)$ enters with the coefficient $+1$.
Hence $d\hat A$ and $F'$ agree at $(x+e_3;12)$ as well.
Finally, let $\hat A'$ be a real one-cochain on $B^{(1)}$ with $d\hat A'=F'$ on $B^{(2)}$ and $\hat A'=t$ on $\mathfrak{T}$.
The three lines of the recursion are among the equations $(d\hat A')_{\mu\nu}(x)=F'_{\mu\nu}(x)$, solved for one link, so $\hat A'$ satisfies them;
as the recursion and the values on $\mathfrak{T}$ determine the cochain, $\hat A'=\hat A(F',t)$.

\emph{The lift.}
On every cube of $B_{\Z}$ one has $dF=d\,dA+\frac{2\pi}{q}\cdot dz=0$, and \eqref{it:chamber-recursion} gives $d\hat A^{B}=F$ on $B^{(2)}$.
For $b\in\mathfrak{T}$ the number $\hat A^{B}(b)-A(b)=\arg_qu_b-A(b)$ lies in $\frac{2\pi}{q}\Z$, by the definition of $\arg_q$.
The cochain $\hat A^{B}$ satisfies the recursion with $F'=dA+\frac{2\pi}{q}\cdot z$, and the formula for $dA$ in Definition~\ref{def:cochains} says that $A$ satisfies the same recursion with $F'=dA$.
Subtracting, one finds
\[
\begin{aligned}
\left(\hat A^{B}-A\right)_1(x+e_2)&=\left(\hat A^{B}-A\right)_1(x)+\left(\hat A^{B}-A\right)_2(x+e_1)-\left(\hat A^{B}-A\right)_2(x)-\frac{2\pi}{q}\cdot z_{12}(x),\\
\left(\hat A^{B}-A\right)_1(x+e_3)&=\left(\hat A^{B}-A\right)_1(x)+\left(\hat A^{B}-A\right)_3(x+e_1)-\left(\hat A^{B}-A\right)_3(x)-\frac{2\pi}{q}\cdot z_{13}(x),\\
\left(\hat A^{B}-A\right)_2(x+e_3)&=\left(\hat A^{B}-A\right)_2(x)+\left(\hat A^{B}-A\right)_3(x+e_2)-\left(\hat A^{B}-A\right)_3(x)-\frac{2\pi}{q}\cdot z_{23}(x),
\end{aligned}
\]
for the same plaquettes as in the recursion.
In the identity that belongs to one step of the recursion at a plaquette $(x;\mu\nu)$, the left side is the value of $\hat A^{B}-A$ at the link determined by that step,
and the right side is the sum, taken with signs, of the values of $\hat A^{B}-A$ at three links, each in $\mathfrak{T}$ or determined by an earlier step, and of the number $-\frac{2\pi}{q}\cdot z_{\mu\nu}(x)\in\frac{2\pi}{q}\Z$.
Since $\frac{2\pi}{q}\Z$ is a group, induction along the order of the recursion, starting from the values on $\mathfrak{T}$, shows that $\hat A^{B}-A$ takes its values in $\frac{2\pi}{q}\Z$ on all of $B^{(1)}$;
this means $\exp\left(iq\hat A^{B}(b)\right)=\exp\left(iqA(b)\right)$ for every $b\in B^{(1)}$.

\emph{The coefficients.}
Let $\lambda$ be an integer-valued function on the sites of $B_{\Z}$, let $\hat A$ be a real one-cochain on $B^{(1)}$, and let $\Psi$ be a monomial.
The gauge invariance of $G$ in Definition~\ref{def:integrands}, applied with the real function $-\frac{2\pi}{q}\cdot\lambda$ on the sites of $B_{\Z}$, gives
\[
a_\Psi\left(\hat A+\frac{2\pi}{q}\cdot d\lambda\right)\cdot\exp\left(-2\pi i\langle\lambda,\mathfrak{q}(\Psi)\rangle\right)=a_\Psi(\hat A),
\]
where $\langle\lambda,\mathfrak{q}(\Psi)\rangle=\sum_y\lambda(y)\cdot\mathfrak{q}(\Psi)(y)$, the sum being taken over the sites $y$ of $B_{\Z}$.
The charge density $\mathfrak{q}(\Psi)$ of Definition~\ref{def:chamber-operator} takes integer values, and so does $\lambda$;
hence the exponential factor equals $1$, and $a_\Psi\left(\hat A+\frac{2\pi}{q}\cdot d\lambda\right)=a_\Psi(\hat A)$.

Let $A$, $z$ and $F$ be as in the second part of \eqref{it:lift-coefficients}.
The map $(A,z)\mapsto\left((F_\sigma)_{\sigma\in B^{(2)}},(u_b)_{b\in\mathfrak{T}}\right)$ is continuous, $\arg_q$ is Borel, and $(F',t)\mapsto\hat A(F',t)$ is linear;
moreover $a_\Psi$ is a Borel function on $\R^{B^{(1)}}$, since $G$ is a Borel map into the finite-dimensional space $\Gr(\mathbb{X}(B))$.
By the definition of $\hat A^{B}$, the number $a_\Psi\left(\hat A^{B}\right)$ results from $\left((F_\sigma)_\sigma,(u_b)_b\right)$ by the map $\left((F_\sigma)_\sigma,(u_b)_b\right)\mapsto\hat A\left((F_\sigma)_\sigma,(\arg_qu_b)_b\right)$, followed by $a_\Psi$.
Neither of these two maps depends on $A$ or $z$, and both are Borel; hence so is their composite.

Let now $z$ be closed on the cubes of $B_{\Z}$, and let $\hat A$ be a local lift on $B$.
By \eqref{it:lift-explicit}, $d\left(\hat A-\hat A^{B}\right)=F-F=0$ on $B^{(2)}$ and $\exp\left(iq\left(\hat A-\hat A^{B}\right)(b)\right)=1$ for $b\in B^{(1)}$;
so $k:=\frac{q}{2\pi}\cdot\left(\hat A-\hat A^{B}\right)$ is an integer-valued one-cochain on $B^{(1)}$ with $dk=0$ on $B^{(2)}$.
For a site $x$ of $B_{\Z}$ let $\lambda(x)$ be the sum of the values of $k$ on the links of the path of $\mathfrak{T}$ from the corner to $x$.
If $(x',\mu)\in\mathfrak{T}$, the path to $x'+e_\mu$ consists of the path to $x'$ and the link $(x',\mu)$, the only link of $\mathfrak{T}$ with the upper end point $x'+e_\mu$;
hence $(d\lambda)_\mu(x')=\lambda(x'+e_\mu)-\lambda(x')=k_\mu(x')$, that is, $d\lambda=k$ on $\mathfrak{T}$.
The cochain $k-d\lambda$ vanishes on $\mathfrak{T}$ and satisfies $d(k-d\lambda)=dk=0$ on $B^{(2)}$;
by \eqref{it:chamber-recursion} it equals $\hat A(0,0)=0$.
Thus $\hat A=\hat A^{B}+\frac{2\pi}{q}\cdot d\lambda$ with an integer-valued $\lambda$, and $a_\Psi(\hat A)=a_\Psi\left(\hat A^{B}\right)$ by the first assertion of \eqref{it:lift-coefficients}.
Finally, $d\hat A^{B}=F$ on $B^{(2)}$ by \eqref{it:lift-explicit}, and the growth condition in the definition of $\mathcal{G}(B)$, applied at $\hat A^{B}$, gives
\[
\left|a_\Psi\left(\hat A^{B}\right)\right|\le C\cdot\exp\left(\left(\frac12-\varepsilon_G\right)\sum_{\sigma\in B^{(2)}}\left(d\hat A^{B}\right)_\sigma^2\right)=C\cdot\exp\left(\left(\frac12-\varepsilon_G\right)\sum_{\sigma\in B^{(2)}}F_\sigma^2\right).
\]

\emph{Sums of products.}
Let $\mathbb{T}$, $(A,z)$, $\widetilde A$, $\widetilde z$, $k$ and $\hat A$ be as in \eqref{it:chamber-sums}.
In Definition~\ref{def:configurations} the coboundary on $\Z^3/\Gamma$ is that of the periodic extensions;
since $dz=0$, the integer two-cochain $\widetilde z$ on $\Z^3$ is closed, and $\widetilde F=d\widetilde A+\frac{2\pi}{q}\cdot\widetilde z$.
The cochain $k$ is integer valued, as $\widetilde z$ is.
Since $k_1=0$, the formula for $dA$ in Definition~\ref{def:cochains} gives
\[
\begin{aligned}
(dk)_{12}(x)&=k_2(x+e_1)-k_2(x),\\
(dk)_{13}(x)&=k_3(x+e_1)-k_3(x),\\
(dk)_{23}(x)&=k_2(x)-k_2(x+e_3)+k_3(x+e_2)-k_3(x).
\end{aligned}
\]
The second sum in $k_3$ does not depend on $x_1$;
so the rule for oriented sums, applied with $t=x_1$, yields $(dk)_{12}(x)=\widetilde z_{12}(x)$ and $(dk)_{13}(x)=\widetilde z_{13}(x)$.
The same rule, applied with $t=x_2$ to the second sum in $k_3$, and the linearity of oriented sums give
\[
\begin{aligned}
(dk)_{23}(x)={}&\widetilde z_{23}(0,x_2,x_3)\\
&+\sum_{s=0}^{x_1-1}\left(\widetilde z_{12}(s,x_2,x_3)-\widetilde z_{12}(s,x_2,x_3+1)+\widetilde z_{13}(s,x_2+1,x_3)-\widetilde z_{13}(s,x_2,x_3)\right).
\end{aligned}
\]
The coboundary $d\widetilde z$ vanishes at the cube with the lowest corner $(s,x_2,x_3)$;
by the formula for $dF$ in Definition~\ref{def:cochains}, this says that the term of index $s$ equals $\widetilde z_{23}(s+1,x_2,x_3)-\widetilde z_{23}(s,x_2,x_3)$.
The sum therefore telescopes to $\widetilde z_{23}(x)-\widetilde z_{23}(0,x_2,x_3)$, and $(dk)_{23}(x)=\widetilde z_{23}(x)$.
Thus $dk=\widetilde z$, so that $d\hat A=d\widetilde A+\frac{2\pi}{q}\cdot\widetilde z=\widetilde F$, and
\[
\exp\left(iq\hat A(b)\right)=\exp\left(iq\widetilde A(b)\right)\cdot\exp\left(2\pi i\cdot k(b)\right)=\exp\left(iq\widetilde A(b)\right)
\]
for every link $b$ of $\Z^3$, the values of $k$ being integers.
For a box that embeds in $\Gamma$, the restriction of $\hat A$ to its links therefore has the coboundary $\widetilde F$ on its plaquettes and the phases of $\widetilde A$ on its links;
as $A$, $z$ and $F$ are read as $\Gamma$-periodic cochains in the definition of local lifts (Definition~\ref{def:integrands}), this restriction is a local lift on the box.

Let $H=\Phi\cdot G_1\cdots G_r$ be a product that embeds in $\Gamma$, with $G_j\in\mathcal{G}^{\mathrm{ev}}(B_j)$, and let $\hat A^{(j)}$ be a local lift on $B_j$ for each $j$.
The boxes of $H$ embed disjointly in $\Gamma$; in particular each $B_j$ embeds in $\Gamma$, and $\hat A|_{B_j^{(1)}}$ is a local lift on $B_j$.
If $B_j$ contains a site of $\Z^3$, then \eqref{it:lift-coefficients}, applied to the box $B_j$, to the restrictions of $\widetilde A$ and $\widetilde z$ to $B_j^{(1)}$ and $B_j^{(2)}$, the latter being closed on the cubes of the integer box of $B_j$ because $d\widetilde z=0$, and to the two local lifts $\hat A|_{B_j^{(1)}}$ and $\hat A^{(j)}$,
shows that every coefficient of $G_j$ takes the same value at these two cochains.
If $B_j$ contains no site of $\Z^3$, then $B_j^{(1)}$ is empty and the two cochains coincide.
In both cases $G_j\left(\hat A|_{B_j^{(1)}}\right)=G_j\left(\hat A^{(j)}\right)$.
As an element of $\Gr(\mathbb{X})^{C^1}$ (Definition~\ref{def:integrands}), $H$ takes at $\hat A$ the value
\[
H(\hat A)=\Phi(d\hat A)\cdot G_1\left(\hat A|_{B_1^{(1)}}\right)\cdots G_r\left(\hat A|_{B_r^{(1)}}\right),
\]
in which $\Phi(d\hat A)=\Phi(\widetilde F)$.
The map $r_{\mathbb{T}}$ is a homomorphism of algebras (Definition~\ref{def:tori});
applying it and replacing each restriction $\hat A|_{B_j^{(1)}}$ by $\hat A^{(j)}$, one obtains the formula for $r_{\mathbb{T}}\left(H(\hat A)\right)$ stated in \eqref{it:chamber-sums}.
Its right side is the element that Definition~\ref{def:integrands} attaches to the product $H$ at the point represented by $(A,z)$, formed with one local lift on each box.

Now let $G\in\mathcal{A}(\R^3)$ embed in $\Gamma$;
by Definition~\ref{def:integrands}, $G=H_1+\dots+H_n$ with products $H_i$ that embed in $\Gamma$.
Sums in $\Gr(\mathbb{X})^{C^1}$ are taken pointwise, so that $G(\hat A)=H_1(\hat A)+\dots+H_n(\hat A)$;
the linearity of $r_{\mathbb{T}}$ and the formula for products give
\[
r_{\mathbb{T}}\left(G(\hat A)\right)=\sum_{i=1}^{n}r_{\mathbb{T}}\left(H_i(\hat A)\right)=\sum_{i=1}^{n}H_i(A,z).
\]
The left side involves only the map $G$ and the cochain $\hat A$, which is determined by $(A,z)$;
the right side involves only the products $H_i$ and the point of $\Conf(\Gamma)$ represented by $(A,z)$.
Hence the common value depends only on $G$ and on that point.
The elements of $\mathcal{A}(\R^3)$ that embed in $\Gamma$ form a linear subspace:
a complex multiple of a product that embeds in $\Gamma$ is again such a product, with the same boxes and with $\Phi$ multiplied by the constant,
and the sum of two finite sums of products that embed in $\Gamma$ is again such a sum.
On this subspace the map $G\mapsto r_{\mathbb{T}}\left(G(\hat A)\right)$ is linear, because evaluation at $\hat A$ and the homomorphism $r_{\mathbb{T}}$ are linear.
Finally, $Z_{\mathbb{T}}[G]$ is the integral over $\Conf(\Gamma)$ of $\exp\left(-\frac12\|F\|^2\right)$ times the Berezin integral of $\exp\left(-S_{\mathbb{T}}(A)\right)\cdot G(A,z)$ (Definition~\ref{def:functional}),
where $G(A,z)$ is the function of $G$ just described;
so it is determined by $G$.
By Lemma~\ref{lem:gaussian-growth}\eqref{it:chamber-convergence}, every coefficient of $\exp\left(-\frac12\|F\|^2\right)\cdot\exp\left(-S_{\mathbb{T}}(A)\right)\cdot G(A,z)$ is absolutely integrable over $\Conf(\Gamma)$, for every element $G$ of the subspace.
The Berezin integral is a linear functional on $\Gr_{\mathbb{T}}$, the integral over $\Conf(\Gamma)$ is linear on integrable functions, and $G(A,z)$ is linear in $G$;
consequently $G\mapsto Z_{\mathbb{T}}[G]$ is linear on the subspace.
\end{proof}

\subsection{Symmetries}\label{sec:symmetries}

\begin{proposition}[Symmetries]\label{prop:symmetries}
Let $\rho$ and $M$ be real, $w_0,\dots,w_3>0$ and $q>0$,
and let $A$ be an arbitrary real one-cochain on $\Z^3$,
which is $\Gamma$-periodic wherever $D^{\mathbb{T}}_A$ or $S_{\mathbb{T}}(A)$ is formed for a torus $\mathbb{T}=(\Gamma,\chi)$;
the action $S_F$ is read as the formal sum of its terms.
\begin{enumerate}
\item\label{it:gauge-covariance}
For every function $\lambda\colon\Z^3\to\R$,
\[
D_{A-d\lambda}=\exp(iq\,\lambda\circ\pi)\,D_A\,\exp(-iq\,\lambda\circ\pi),
\]
where $\exp(iq\,\lambda\circ\pi)$ multiplies a spinor field at every site $x$ by $\exp(iq\lambda(\pi(x)))$;
equivalently, $S_F$ is invariant under the gauge transformation~\eqref{eq:setting-gauge}.
If $\mathbb{T}=(\Gamma,\chi)$ is a torus and $A$ and $\lambda$ are $\Gamma$-periodic,
the multiplication by $\exp(iq\,\lambda\circ\pi)$ maps the spinors of $\mathbb{T}$ onto themselves, and
$D^{\mathbb{T}}_{A-d\lambda}=\exp(iq\,\lambda\circ\pi)\,D^{\mathbb{T}}_A\,\exp(-iq\,\lambda\circ\pi)$.
Moreover, $D_{A+(2\pi/q)k}=D_A$ for every $k\in C^1(\Z)$.
\item\label{it:lattice-covariance}
For every unit vector $n$ the matrix $\gamma_n\gamma_5$ satisfies
\[
(\gamma_n\gamma_5)\,\gamma_v\,(\gamma_n\gamma_5)^{-1}=\gamma_{r_nv},\qquad r_nv=v-2(v\cdot n)n,\quad v\in\R^3 .
\]
Let $g\in W$, let $n_1,\dots,n_k$ be unit vectors with $\vec g=r_{n_1}\cdots r_{n_k}$, and put
\[
\tilde g:=(\gamma_{n_1}\gamma_5)\cdots(\gamma_{n_k}\gamma_5),\qquad
(U_g\psi)(x):=\tilde g\cdot\psi(g^{-1}x),
\]
for $\psi\colon\mathbb{X}\to\C^4$ and $x\in\mathbb{X}$.
Then
\[
U_g\,D_A\,U_g^{-1}=D_{(g^{-1})^{*}A}.
\]
If $\mathbb{T}=(\Gamma,\chi)$ is a torus and $A$ is $\Gamma$-periodic,
then $U_g$ is a linear bijection from the spinors of $\mathbb{T}$ onto those of $g\mathbb{T}$,
$U_g\,D^{\mathbb{T}}_A\,U_g^{-1}=D^{g\mathbb{T}}_{(g^{-1})^{*}A}$,
and hence $\det D^{g\mathbb{T}}_{(g^{-1})^{*}A}=\det D^{\mathbb{T}}_A$.
\item\label{it:charge-conjugation}
Let $(\psi',\psi):=\sum_x\psi'(x)^{\top}\psi(x)$,
the sum extending over $\mathbb{X}$ for finitely supported fields $\psi,\psi'\colon\mathbb{X}\to\C^4$
and over a set of representatives of $\mathbb{X}_{\mathbb{T}}$ for spinors of a torus $\mathbb{T}$;
let $(D_A)^{\top}$ and $(D^{\mathbb{T}}_A)^{\top}$ be the transposes of $D_A$ and of $D^{\mathbb{T}}_A$ with respect to this bilinear form,
and let $\sigma_2\oplus\sigma_2$ act at every site.
Then
\[
(\sigma_2\oplus\sigma_2)\,(D_A)^{\top}\,(\sigma_2\oplus\sigma_2)^{-1}=D_{-A},
\]
and on every torus $(\sigma_2\oplus\sigma_2)\,(D^{\mathbb{T}}_A)^{\top}\,(\sigma_2\oplus\sigma_2)^{-1}=D^{\mathbb{T}}_{-A}$.
Hence $\det D^{\mathbb{T}}_{-A}=\det D^{\mathbb{T}}_A$,
and the law $\mu_{\mathbb{T}}$ is invariant under $F\mapsto-F$.
\item\label{it:action-invariance}
For every mirror $\Pi$ one has $\Theta_\Pi(S_F)=S_F$,
where $\Theta_\Pi$ acts on the coefficients as in~\eqref{eq:setting-reflection-coefficients}:
it maps the term $-\bar\psi(x)\cdot w(x,x')\cdot U_A(x,x')\cdot\psi(x')$ of the oriented bond $(x,x')$
to the term of the oriented bond $(\theta_\Pi x',\theta_\Pi x)$.
On a torus $\mathbb{T}=(\Gamma,\chi)$, with $\Theta_\Pi$ acting as in Lemmas~\ref{lem:reflection-classes} and~\ref{lem:reflected-algebras},
$(\Theta_\Pi S_{\mathbb{T}})(A)=S_{\theta_\Pi\mathbb{T}}(A)$ whenever $A$ is $\vec\theta_\Pi\Gamma$-periodic,
$\vec\theta_\Pi$ being the linear part of $\theta_\Pi$.
\item\label{it:berezin-conjugate}
For every mirror $\Pi$, every finite set $Y\subset\mathbb{X}$ and every $G\in\Gr(Y)$,
\[
\int_{\theta_\Pi Y}\Theta_\Pi(G)=\overline{\int_YG},
\]
the Berezin integrals being normalised as in~\eqref{eq:setting-berezin}.
\item\label{it:coxeter-relations}
For $i\neq j$ in $\{0,1,2,3\}$ let $k_{ij}$ be the order of $s_is_j$;
thus $k_{01}=k_{23}=4$, $k_{12}=3$, and $k_{ij}=2$ for the remaining pairs.
Then
\[
\left(\Theta_{\Pi_i}\Theta_{\Pi_j}\right)^{k_{ij}}=(-1)^{(k_{ij}+1)\deg}.
\]
\item[\textup{(f$'$)}]\itemmark{f$'$}\label{it:conjugate-reflections}
Let $\Pi$ and $\Pi'$ be mirrors with upward unit normals $n$ and $n'$, and let $\Pi'':=\theta_{\Pi'}(\Pi)$.
Then
\[
\Theta_{\Pi'}\Theta_\Pi\Theta_{\Pi'}=(-1)^{\deg}\,\Theta_{\Pi''},
\]
where $\Theta_{\Pi''}$ is formed by~\eqref{eq:setting-reflection-fields} with the unit normal $n'':=n-2(n\cdot n')n'$ of $\Pi''$.
\item[\textup{(f$''$)}]\itemmark{f$''$}\label{it:reflection-words}
Let $\Pi^{(1)},\dots,\Pi^{(r)}$ and $\Pi'^{(1)},\dots,\Pi'^{(s)}$ be mirrors with
$\theta_{\Pi^{(1)}}\cdots\theta_{\Pi^{(r)}}=\theta_{\Pi'^{(1)}}\cdots\theta_{\Pi'^{(s)}}$.
Then, on $\Gr(\mathbb{X})$ and on integrands,
$\Theta_{\Pi^{(1)}}\cdots\Theta_{\Pi^{(r)}}$ is equal either to $\Theta_{\Pi'^{(1)}}\cdots\Theta_{\Pi'^{(s)}}$
or to $(-1)^{\deg}\,\Theta_{\Pi'^{(1)}}\cdots\Theta_{\Pi'^{(s)}}$.
Consequently, for $g\in W$ and an even integrand $G$,
the element $G^g$ of~\eqref{eq:setting-group-action} does not depend on the word $g=\theta_{\Pi^{(1)}}\cdots\theta_{\Pi^{(r)}}$ with which it is formed,
and $G\mapsto G^g$ is an action of $W$ on the even integrands:
$(G^g)^{g'}=G^{g'g}$ for $g,g'\in W$;
the map $G\mapsto G^g$ is a linear automorphism if $g$ is even and an antilinear anti-automorphism if $g$ is odd;
and for a bosonic integrand $G$, regarded as a function of $F$,
\[
G^g(F)=G(g^{*}F)\quad\text{if $g$ is even},\qquad
G^g(F)=\overline{G(g^{*}F)}\quad\text{if $g$ is odd}.
\]
\item\label{it:functional-covariance}
Let $\mathbb{T}$ be a torus.
For every $G\in\mathcal{A}(\R^3)$ that embeds in $\mathbb{T}$ and every $g\in W$,
\[
Z_{g\mathbb{T}}[G^g]=
\begin{cases}
Z_{\mathbb{T}}[G] & \text{if $g$ is even},\\
\overline{Z_{\mathbb{T}}[G]} & \text{if $g$ is odd};
\end{cases}
\]
and for every mirror $\Pi$ and every $G\in\mathcal{B}(\mathbb{T})$ (Lemma~\ref{lem:bounded-integrands}), even or odd,
\[
Z_{\theta_\Pi\mathbb{T}}[\Theta_\Pi G]=\overline{Z_{\mathbb{T}}[G]},
\]
the element $\Theta_\Pi G$ of $\mathcal{B}(\theta_\Pi\mathbb{T})$ being formed as in Lemmas~\ref{lem:reflection-classes} and~\ref{lem:reflected-algebras}.
In particular, $\mu_{g\mathbb{T}}$ is the image of $\mu_{\mathbb{T}}$ under the map $F\mapsto(g^{-1})^{*}F$ of $\Omega$;
the law $\mu_{\mathbb{T}}$ is invariant under the translations $\tau_v$, $v\in\Z^3$;
and for every cubic torus $C(n)$ the law $\mu_{C(n)}$ is invariant under $W$.
\item\label{it:sign-character}
Let $\chi,\chi'\colon\Gamma\to\{+1,-1\}$ be two characters of a period lattice $\Gamma$,
and let $\mathbb{T}:=(\Gamma,\chi)$ and $\mathbb{T}':=(\Gamma,\chi')$.
Then $Z_{\mathbb{T}}[G]=Z_{\mathbb{T}'}[G]$ for every $G\in\mathcal{A}(\R^3)$ that embeds in $\Gamma$.
In particular, $Z_{\mathbb{T}}[1]$, the law $\mu_{\mathbb{T}}$,
and the values $\mathbb{E}_{\mathbb{T}}[G]$ at the elements $G$ of $\mathcal{A}(\R^3)$ that embed in $\Gamma$
do not depend on $\chi$.
\end{enumerate}
\end{proposition}

\begin{proof}
\emph{(a)}
Let $\lambda\colon\Z^3\to\R$, and let $\{x,x'\}$ be a bond.
By Proposition~\ref{prop:geometry}\eqref{it:geometry-projection} the photon sites $\pi(x)$ and $\pi(x')$ either coincide or are nearest neighbours;
in both cases Definition~\ref{def:cochains} gives $(d\lambda)(\pi(x),\pi(x'))=\lambda(\pi(x'))-\lambda(\pi(x))$, and therefore
\[
U_{A-d\lambda}(x,x')=\exp(iq\lambda(\pi(x)))\cdot U_A(x,x')\cdot\exp(-iq\lambda(\pi(x'))).
\]
These scalar factors commute with the hopping matrix $w(x,x')$.
Let $\psi\colon\mathbb{X}\to\C^4$ and $x\in\mathbb{X}$.
The two exponential factors at $x$ cancel in the mass term,
so that the value of $\exp(iq\,\lambda\circ\pi)\,D_A\,\exp(-iq\,\lambda\circ\pi)\,\psi$ at $x$ is $M\psi(x)$
minus the sum over $i=0,\dots,3$ of the left sides of the identities
\[
\begin{aligned}
&\exp\left(iq\lambda(\pi(x))\right)\cdot w(x,x^{(i)})\cdot U_A(x,x^{(i)})\cdot\exp\left(-iq\lambda(\pi(x^{(i)}))\right)\cdot\psi(x^{(i)})\\
&\qquad=w(x,x^{(i)})\cdot U_{A-d\lambda}(x,x^{(i)})\cdot\psi(x^{(i)}).
\end{aligned}
\]
The same expression formed with the right sides is $\left(D_{A-d\lambda}\psi\right)(x)$ (Definition~\ref{def:chamber-operator});
hence $\exp(iq\,\lambda\circ\pi)\,D_A\,\exp(-iq\,\lambda\circ\pi)=D_{A-d\lambda}$.
Since $S_F(A)=\bar\psi D_A\psi$, the substitution~\eqref{eq:setting-gauge} turns $S_F(A)$ into
$\bar\psi\,\exp(-iq\,\lambda\circ\pi)\,D_{A-d\lambda}\,\exp(iq\,\lambda\circ\pi)\,\psi$,
which by the identity just proved equals $\bar\psi D_A\psi$;
this is the invariance of $S_F$.

Now let $\mathbb{T}=(\Gamma,\chi)$ be a torus, and let $A$ and $\lambda$ be $\Gamma$-periodic.
The translations by the elements of $\Gamma$ belong to $W$, and $\pi$ is $W$-equivariant (Proposition~\ref{prop:geometry}\eqref{it:geometry-projection});
so $\pi(x+v)=\pi(x)+v$ and $\lambda(\pi(x+v))=\lambda(\pi(x))$ for $v\in\Gamma$,
and the multiplication by $\exp(iq\,\lambda\circ\pi)$ preserves the relations $\psi(x+v)=\chi(v)\cdot\psi(x)$ that define the spinors of $\mathbb{T}$.
As $A-d\lambda$ is $\Gamma$-periodic as well, the identity restricts to the spinors of $\mathbb{T}$.

For $k\in C^1(\Z)$ the number $k(\pi(x),\pi(x'))$ is an integer, so that
\[
U_{A+(2\pi/q)k}(x,x')=U_A(x,x')\cdot\exp\left(2\pi i\,k(\pi(x),\pi(x'))\right)=U_A(x,x'),
\]
and $D_{A+(2\pi/q)k}=D_A$.
Two representatives of a point of $\Conf(\Gamma)$ differ in their first components by $(2\pi/q)k$ with $k\in C^1(\Gamma;\Z)$ (Definition~\ref{def:configurations}),
and $D^{\mathbb{T}}_A$ and $\exp(-S_{\mathbb{T}}(A))$ depend on $A$ only through the link phases $U_A(x,x^{(i)})$.
Hence $\det D^{\mathbb{T}}_A$ and $\exp(-S_{\mathbb{T}}(A))$ are functions on $\Conf(\Gamma)$.
Finally, $\det D^{\mathbb{T}}_{A-d\lambda}=\det D^{\mathbb{T}}_A$ for $\Gamma$-periodic $\lambda$, the two operators being conjugate.

\emph{(b)}
The matrix $\gamma_5$, which exchanges the two summands of $\C^2\oplus\C^2$, is Hermitian with $\gamma_5^2=1$,
and it anticommutes with $\gamma_1$, $\gamma_2$, $\gamma_3$, hence with every $\gamma_v$ (Definition~\ref{def:spin-matrices}).
For a unit vector $n$ the matrix $\gamma_n$ is Hermitian with $\gamma_n^2=1$.
Therefore $\gamma_n\gamma_5$ is unitary, $(\gamma_n\gamma_5)^{-1}=\gamma_5\gamma_n$, and by Lemma~\ref{lem:spin-identities}
\[
(\gamma_n\gamma_5)\,\gamma_v\,(\gamma_n\gamma_5)^{-1}=\gamma_n\gamma_5\gamma_v\gamma_5\gamma_n=-\gamma_n\gamma_v\gamma_n=\gamma_{r_nv},\qquad v\in\R^3 .
\]
Factorisations $\vec g=r_{n_1}\cdots r_{n_k}$ exist:
by Lemma~\ref{lem:chambers}\eqref{it:chambers-generators} the element $g$ is a product of the wall reflections $s_0,\dots,s_3$,
and the linear part $\vec s_i$ of $s_i$ is the reflection $r_{\hat n_i}$ attached to the unit normal $\hat n_i$ of $\Pi_i$.
Applying the last identity $k$ times one obtains
$\tilde g\,\gamma_v\,\tilde g^{-1}=\gamma_{r_{n_1}\cdots r_{n_k}v}=\gamma_{\vec gv}$ for every $v\in\R^3$.

Let $x\in\mathbb{X}$, and let $g'\in W$ be the element with $x=g'v_0$ (Lemma~\ref{lem:chambers}\eqref{it:chambers-transitive}).
By Definition~\ref{def:fermion-sites}, both $(g^{-1}x)^{(i)}$ and $g^{-1}x^{(i)}$ are equal to $g^{-1}g's_iv_0$;
thus $(g^{-1}x)^{(i)}=g^{-1}x^{(i)}$ for $i=0,\dots,3$.
As $U_g^{-1}\psi=\tilde g^{-1}\cdot(\psi\circ g)$, this gives
\[
\begin{aligned}
&\left(U_g\,D_A\,U_g^{-1}\psi\right)(x)\\
&\qquad=M\psi(x)-\sum_{i=0}^{3}\tilde g\cdot w(g^{-1}x,g^{-1}x^{(i)})\cdot\tilde g^{-1}\cdot U_A(g^{-1}x,g^{-1}x^{(i)})\cdot\psi(x^{(i)}).
\end{aligned}
\]
The bond $\{g^{-1}x,g^{-1}x^{(i)}\}$ has the type $i$,
and $g^{-1}$ is an isometry with linear part $\vec g^{-1}$, so that the unit vector from $g^{-1}x$ to $g^{-1}x^{(i)}$ is $\vec g^{-1}\hat d_i(x)$.
Hence
\[
\tilde g\cdot w(g^{-1}x,g^{-1}x^{(i)})\cdot\tilde g^{-1}
=\frac{w_i}{2}\left(\rho-\tilde g\,\gamma_{\vec g^{-1}\hat d_i(x)}\,\tilde g^{-1}\right)
=\frac{w_i}{2}\left(\rho-\gamma_{\hat d_i(x)}\right)
=w(x,x^{(i)}).
\]
Further, $\pi(g^{-1}x)=g^{-1}\pi(x)$ for $x\in\mathbb{X}$ (Proposition~\ref{prop:geometry}\eqref{it:geometry-projection}),
and $\left((g^{-1})^{*}A\right)(y,y')=A(g^{-1}y,g^{-1}y')$ for photon sites $y,y'$ (Definition~\ref{def:cochains}); whence
\[
U_A(g^{-1}x,g^{-1}x^{(i)})=\exp\left(iqA\left(g^{-1}\pi(x),g^{-1}\pi(x^{(i)})\right)\right)=U_{(g^{-1})^{*}A}(x,x^{(i)}).
\]
Inserting the last two identities one finds $U_g\,D_A\,U_g^{-1}=D_{(g^{-1})^{*}A}$.

Let now $\mathbb{T}=(\Gamma,\chi)$ be a torus, $\psi$ a spinor of $\mathbb{T}$ (Definition~\ref{def:torus-spinors}) and $v\in\Gamma$.
Since $g^{-1}(x+\vec gv)=g^{-1}x+v$,
\[
(U_g\psi)(x+\vec gv)=\tilde g\cdot\psi(g^{-1}x+v)=\chi(v)\cdot(U_g\psi)(x)=\left(\chi\circ\vec g^{-1}\right)(\vec gv)\cdot(U_g\psi)(x);
\]
so $U_g\psi$ is a spinor of $g\mathbb{T}=(\vec g\,\Gamma,\chi\circ\vec g^{-1})$ (Definition~\ref{def:tori}).
In the same way $U_g^{-1}$ maps the spinors of $g\mathbb{T}$ to those of $\mathbb{T}$,
and $U_g$ is a linear bijection between the two spaces.
If $A$ is $\Gamma$-periodic, then $(g^{-1})^{*}A$ is $\vec g\,\Gamma$-periodic,
because $\left((g^{-1})^{*}A\right)(y+\vec gv,y'+\vec gv)=A(g^{-1}y+v,g^{-1}y'+v)=A(g^{-1}y,g^{-1}y')$.
The identity of operators therefore restricts to $U_g\,D^{\mathbb{T}}_A\,U_g^{-1}=D^{g\mathbb{T}}_{(g^{-1})^{*}A}$,
and two operators that are conjugate by a linear bijection between spaces of finite dimension have the same determinant.

The matrix $\tilde g$ is not determined by $g$ alone.
Indeed $(\gamma_n\gamma_5)^2=-\gamma_n^2\gamma_5^2=-1$,
so that for $\vec g=1$ the empty factorisation and the factorisation $r_nr_n$ give $\tilde g=1$ and $\tilde g=-1$.
It is, however, determined up to a sign.
Let $\vec g=r_{n'_1}\cdots r_{n'_{k'}}$ be a second factorisation of the linear part of $g$, and let
\[
V:=(\gamma_{n'_1}\gamma_5)\cdots(\gamma_{n'_{k'}}\gamma_5)
\]
be the product formed with it.
Both $(-1)^k$ and $(-1)^{k'}$ equal $\det\vec g$, so $k+k'$ is even.
As $(\gamma_n\gamma_5)^{-1}=-\gamma_n\gamma_5$, the matrix $\tilde g\,V^{-1}$ is, up to a sign,
a product of $k+k'$ matrices $\gamma_n\gamma_5$ with unit vectors $n$.
For $v\in\R^3$ write $\sigma_v:=\sum_\mu v_\mu\sigma_\mu$, so that $\gamma_v=\diag(\sigma_v,-\sigma_v)$.
For unit vectors $n$ and $n'$,
\[
(\gamma_n\gamma_5)(\gamma_{n'}\gamma_5)=-\gamma_n\gamma_{n'}\gamma_5^2=-\diag(\sigma_n\sigma_{n'},\sigma_n\sigma_{n'}),
\]
and grouping the factors of $\tilde g\,V^{-1}$ in consecutive pairs one obtains $\tilde g\,V^{-1}=\pm\diag(u,u)$,
where $u$ is a product of the even number $k+k'$ of matrices $\sigma_n$ with unit vectors $n$.
Since $V^{-1}\gamma_vV=\gamma_{\vec g^{-1}v}$ and $\tilde g\,\gamma_{\vec g^{-1}v}\,\tilde g^{-1}=\gamma_v$,
the matrix $\tilde g\,V^{-1}$, and with it $\diag(u,u)$, commutes with every $\gamma_v$;
so $u$ commutes with $\sigma_1$, $\sigma_2$ and $\sigma_3$.
By the argument that leads to~\eqref{eq:chamber-central-block} in the proof of~\eqref{it:reflection-words},
a product of an even number of matrices $\sigma_n$ with unit vectors $n$ that commutes with $\sigma_1$, $\sigma_2$ and $\sigma_3$
is a scalar matrix of determinant $1$, hence equal to $1$ or $-1$;
there the unit vectors are normals of mirrors, but the argument holds for arbitrary unit vectors.
So $u=\pm1$, and $V=\pm\tilde g$.
Replacing $\tilde g$ by $-\tilde g$ replaces $U_g$ by $-U_g$ and leaves $U_g\,D_A\,U_g^{-1}$ unchanged;
the identities above hold for every factorisation of $\vec g$.

\emph{(c)}
Let $\mathbb{T}=(\Gamma,\chi)$ be a torus, and let $\psi'$, $\psi$ be spinors of $\mathbb{T}$.
Replacing $x$ by $x+v$, $v\in\Gamma$, multiplies $\psi'(x)^{\top}\psi(x)$ by $\chi(v)^2=1$,
so $(\psi',\psi)$ does not depend on the set of representatives of $\mathbb{X}_{\mathbb{T}}$.
In the coordinates given by the values at the representatives (Definition~\ref{def:torus-spinors}) it is the standard symmetric bilinear form;
the transpose $(D^{\mathbb{T}}_A)^{\top}$ is therefore represented by the transposed matrix, and $\det\left(D^{\mathbb{T}}_A\right)^{\top}=\det D^{\mathbb{T}}_A$.
For each $i$ the map $x\mapsto x^{(i)}$ is an involution of $\mathbb{X}$ that commutes with the translations by $\Gamma$ (Definition~\ref{def:tori}),
and so it induces a bijection of $\mathbb{X}_{\mathbb{T}}$.
For $\Gamma$-periodic $A$ the expression $\psi'(x)^{\top}\cdot w(x,x^{(i)})\cdot U_A(x,x^{(i)})\cdot\psi(x^{(i)})$ is $\Gamma$-invariant in $x$,
and re-indexing its sum by $x\mapsto x^{(i)}$ yields
\[
(\psi',D^{\mathbb{T}}_A\psi)=\sum_xM\psi'(x)^{\top}\psi(x)-\sum_x\sum_{i=0}^{3}\psi'(x^{(i)})^{\top}\cdot w(x^{(i)},x)\cdot U_A(x^{(i)},x)\cdot\psi(x),
\]
the sums being taken over a set of representatives of $\mathbb{X}_{\mathbb{T}}$.
The cochain $A$ is odd on oriented pairs, so $U_A(x^{(i)},x)=U_{-A}(x,x^{(i)})$, and the last identity says
\[
\left((D^{\mathbb{T}}_A)^{\top}\psi'\right)(x)=M\psi'(x)-\sum_{i=0}^{3}w(x^{(i)},x)^{\top}\cdot U_{-A}(x,x^{(i)})\cdot\psi'(x^{(i)}).
\]
The unit vector from $x^{(i)}$ to $x$ is $-\hat d_i(x)$, whence $w(x^{(i)},x)=\frac{w_i}{2}\left(\rho+\gamma\cdot\hat d_i(x)\right)$ (Lemma~\ref{lem:hopping});
moreover $(\sigma_2\oplus\sigma_2)\,(\gamma\cdot v)^{\top}\,(\sigma_2\oplus\sigma_2)^{-1}=-\gamma\cdot v$ for $v\in\R^3$ (Definition~\ref{def:spin-matrices}).
Consequently
\[
(\sigma_2\oplus\sigma_2)\,w(x^{(i)},x)^{\top}\,(\sigma_2\oplus\sigma_2)^{-1}=\frac{w_i}{2}\left(\rho-\gamma\cdot\hat d_i(x)\right)=w(x,x^{(i)}),
\]
and $(\sigma_2\oplus\sigma_2)\,(D^{\mathbb{T}}_A)^{\top}\,(\sigma_2\oplus\sigma_2)^{-1}=D^{\mathbb{T}}_{-A}$.
The same computation, with sums over $\mathbb{X}$ and finitely supported fields, proves the identity for $D_A$.
Taking determinants on the torus, $\det D^{\mathbb{T}}_{-A}=\det\left(D^{\mathbb{T}}_A\right)^{\top}=\det D^{\mathbb{T}}_A$.

The map $(A,z)\mapsto(-A,-z)$ sends $(A+(2\pi/q)k,\,z-dk)$ to $(-A+(2\pi/q)(-k),\,-z-d(-k))$;
it therefore respects the equivalence of Definition~\ref{def:configurations} and induces a bijection of $\Conf(\Gamma)$,
which maps the curvature $F$ to $-F$.
This bijection preserves the measure of $\Conf(\Gamma)$.
Indeed, the product of the Lebesgue measure $\mathrm{d}A$ on $C^1(\Gamma)$ with the counting measure on the closed elements of $C^2(\Gamma;\Z)$
is invariant under $(A,z)\mapsto(-A,-z)$,
as $A\mapsto-A$ is a linear map of determinant $\pm1$ and $z\mapsto-z$ permutes the closed integer two-cochains;
and $(A,z)\mapsto(-A,-z)$ carries the fundamental domain formed by the box $[0,2\pi/q)^{3\cdot[\Z^3:\Gamma]}$ and the closed $z$
onto the fundamental domain formed by $(-2\pi/q,0]^{3\cdot[\Z^3:\Gamma]}$ and the same cochains.
So the integral of a Borel function $\ge0$ on $\Conf(\Gamma)$ composed with the bijection, computed on the first of these domains,
equals the integral of the function itself computed on the second,
and every measurable fundamental domain gives the same value (Definition~\ref{def:configurations}).
Let $\Phi$ be a bounded Borel function of finitely many $F_\sigma$.
By $\int_{\mathbb{T}}\exp(-S_{\mathbb{T}}(A))=\det D^{\mathbb{T}}_A$ (Definition~\ref{def:tori}),
\[
\begin{aligned}
Z_{\mathbb{T}}[\Phi(-\widetilde F)]
&=\int_{\Conf(\Gamma)}\exp\left(-\frac{1}{2}\|F\|^2\right)\cdot\det D^{\mathbb{T}}_A\cdot\Phi(-\widetilde F)\\
&=\int_{\Conf(\Gamma)}\exp\left(-\frac{1}{2}\|F\|^2\right)\cdot\det D^{\mathbb{T}}_{-A}\cdot\Phi(\widetilde F)\\
&=Z_{\mathbb{T}}[\Phi(\widetilde F)];
\end{aligned}
\]
here the second equality is the substitution $(A,z)\mapsto(-A,-z)$, together with $\|-F\|=\|F\|$,
the third one is $\det D^{\mathbb{T}}_{-A}=\det D^{\mathbb{T}}_A$,
and the integrands are integrable by Lemma~\ref{lem:bounded-integrands}.
By Theorem~\ref{thm:partition-positivity}, $0<Z_{\mathbb{T}}[1]<\infty$,
and division by $Z_{\mathbb{T}}[1]$ gives, for the law $\mu_{\mathbb{T}}$ of $\widetilde F$ under $\mathbb{E}_{\mathbb{T}}$ (Definition~\ref{def:functional}),
\[
\int\mathrm{d}\mu_{\mathbb{T}}(F)\,\Phi(-F)=\int\mathrm{d}\mu_{\mathbb{T}}(F)\,\Phi(F).
\]
The Borel $\sigma$-algebra of $\Omega$ is generated by the cylinder sets,
so $\mu_{\mathbb{T}}$ is invariant under $F\mapsto-F$.

\emph{(d)}
Let $\Pi$ be a mirror, $\theta:=\theta_\Pi$, and let $n$ be the upward unit normal of $\Pi$, with which $\Theta_\Pi$ is formed in~\eqref{eq:setting-reflection-fields}.
The term of an oriented bond $(x,x')$ in $S_F$ is $-\bar\psi(x)\,D\,\psi(x')$ with the $4\times4$ matrix $D:=w(x,x')\cdot U_A(x,x')$, whose entries are functions of $A$.
The rule~\eqref{eq:setting-reflection-bilinear} for reflected bilinear elements,
together with the action~\eqref{eq:setting-reflection-coefficients} of $\Theta_\Pi$ on the coefficients,
maps this term to
\[
-\bar\psi(\theta x')\,\gamma_n\,w(x,x')^{*}\,\gamma_n\cdot\overline{U_{\theta^{*}A}(x,x')}\cdot\psi(\theta x).
\]
Since $\rho$ is real, Lemma~\ref{lem:hopping} gives $w(x,x')^{*}=w(x,x')$ and $\gamma_n\,w(x,x')\,\gamma_n=w(\theta x',\theta x)$.
The cochain $A$ is real and odd on oriented pairs,
$(\theta^{*}A)(y,y')=A(\theta y,\theta y')$ for photon sites $y,y'$ that are equal or nearest neighbours (Definition~\ref{def:cochains}),
and $\pi\circ\theta=\theta\circ\pi$ on $\mathbb{X}$ (Proposition~\ref{prop:geometry}\eqref{it:geometry-projection}); hence
\[
\overline{U_{\theta^{*}A}(x,x')}
=\exp\left(-iqA\left(\theta\pi(x),\theta\pi(x')\right)\right)
=\exp\left(iqA\left(\pi(\theta x'),\pi(\theta x)\right)\right)
=U_A(\theta x',\theta x).
\]
The image is therefore $-\bar\psi(\theta x')\cdot w(\theta x',\theta x)\cdot U_A(\theta x',\theta x)\cdot\psi(\theta x)$, the term of the oriented bond $(\theta x',\theta x)$.
The mass term $M\bar\psi(x)\psi(x)$ is mapped to $M\bar\psi(\theta x)\gamma_n\gamma_n\psi(\theta x)=M\bar\psi(\theta x)\psi(\theta x)$,
because $M$ is real and $\gamma_n^2=1$.
If $x=gv_0$ with $g\in W$, then $\theta x=(\theta g)v_0$, and Definition~\ref{def:fermion-sites} gives $(\theta x)^{(i)}=\theta gs_iv_0=\theta(x^{(i)})$.
So $(x,x')\mapsto(\theta x',\theta x)$ maps oriented bonds to oriented bonds;
it is an involution, hence a bijection of the set of oriented bonds, and $x\mapsto\theta x$ is a bijection of $\mathbb{X}$.
Consequently $\Theta_\Pi(S_F)=S_F$ as formal sums, the terms being permuted.

On a torus $\mathbb{T}=(\Gamma,\chi)$ the action $S_{\mathbb{T}}$ is the sum of the same terms over representatives of the classes of sites and of oriented bonds modulo~$\Gamma$.
As $\theta(x+g)=\theta x+\vec\theta g$ for $g\in\Gamma$, the map $x\mapsto\theta x$ induces bijections of these classes onto the classes modulo $\theta\Gamma=\vec\theta\,\Gamma$.
Let $A$ be a $\theta\Gamma$-periodic one-cochain; then $\theta^{*}A$ is $\Gamma$-periodic.
Read in $\Gr_{\mathbb{T}}$ and $\Gr_{\theta\mathbb{T}}$ through Lemmas~\ref{lem:reflection-classes} and~\ref{lem:reflected-algebras},
the computation above shows that $\Theta_\Pi$ maps $S_{\mathbb{T}}(\theta^{*}A)$ to $S_{\theta\mathbb{T}}(A)$,
that is, $(\Theta_\Pi S_{\mathbb{T}})(A)=S_{\theta\mathbb{T}}(A)$.

\emph{(e)}
With $\Pi$, $\theta$ and $n$ as above, let $Y\subset\mathbb{X}$ be finite, and let
\[
\Psi_Y:=\prod_{x\in Y}\prod_{\alpha=1}^{4}\psi_\alpha(x)\,\bar\psi_\alpha(x)
\]
be the product of all pairs.
Its factors are even, so that their order is immaterial,
and $\int_YG$ is the coefficient of $\Psi_Y$ in $G\in\Gr(Y)$ (Definition~\ref{def:grassmann}).
By~\eqref{eq:setting-reflection-fields}, $\Theta_\Pi$ sends every generator at a site $x$ to a linear combination of generators at $\theta x$;
a monomial of degree $k$ in the generators at the sites of $Y$ therefore goes to a linear combination of monomials of degree $k$ in the generators at the sites of $\theta Y$.
In particular $\Theta_\Pi(\Gr(Y))\subset\Gr(\theta Y)$, with equality because $\Theta_\Pi^2=1$ (Lemma~\ref{lem:sign-rules}(a)) and $\theta(\theta Y)=Y$.
The components of $G$ of degree less than $8\cdot\#|Y|$ are thus mapped to combinations of monomials of degree less than $8\cdot\#|\theta Y|$, whose integral vanishes;
as $\Theta_\Pi$ is antilinear,
\[
\int_{\theta Y}\Theta_\Pi(G)=\overline{\int_YG}\cdot\int_{\theta Y}\Theta_\Pi(\Psi_Y).
\]
By~\eqref{eq:setting-reflection-fields}, with $\psi':=\gamma_n\psi(\theta x)$ and $\bar\psi':=\bar\psi(\theta x)\,\gamma_n$,
\[
\Theta_\Pi(\psi_\alpha(x)\,\bar\psi_\alpha(x))=\Theta_\Pi(\bar\psi_\alpha(x))\,\Theta_\Pi(\psi_\alpha(x))=\psi'_\alpha\,\bar\psi'_\alpha ,
\]
the two signs $-1$ cancelling.
Reordering eight odd factors one has
\[
\prod_{\alpha=1}^{4}\psi'_\alpha\,\bar\psi'_\alpha=\pm\,(\psi'_1\psi'_2\psi'_3\psi'_4)\,(\bar\psi'_1\bar\psi'_2\bar\psi'_3\bar\psi'_4),
\]
where the sign is that of the reordering permutation alone, and so is the same for the unprimed fields at $\theta x$.
Since $\psi'_\alpha=\sum_\beta(\gamma_n)_{\alpha\beta}\,\psi_\beta(\theta x)$ and $\bar\psi'_\alpha=\sum_\beta\bar\psi_\beta(\theta x)\,(\gamma_n)_{\beta\alpha}$ are linear combinations of anticommuting elements,
\[
\psi'_1\psi'_2\psi'_3\psi'_4=\det(\gamma_n)\cdot\psi_1(\theta x)\cdots\psi_4(\theta x),\qquad
\bar\psi'_1\bar\psi'_2\bar\psi'_3\bar\psi'_4=\det(\gamma_n)\cdot\bar\psi_1(\theta x)\cdots\bar\psi_4(\theta x).
\]
The even pairs commute with each other, so the reversal of products by $\Theta_\Pi$ does not affect $\Theta_\Pi(\Psi_Y)$.
The matrix $\gamma_n$ has the eigenvalues $1$ and $-1$, each twice, so that $\det\gamma_n=1$ (Lemma~\ref{lem:spin-identities}), and
\[
\Theta_\Pi(\Psi_Y)=\det(\gamma_n)^{2\cdot\#|Y|}\cdot\Psi_{\theta Y}=\Psi_{\theta Y},\qquad\int_{\theta Y}\Psi_{\theta Y}=1 .
\]
Hence $\int_{\theta Y}\Theta_\Pi(G)=\overline{\int_YG}$ for every $G\in\Gr(Y)$.

The identity extends to functions on configurations.
Let $\Gamma$ be a period lattice, and let $G$ be a function on $\Conf(\Gamma)$ with values in $\Gr(Y)$.
By Lemma~\ref{lem:reflection-classes}(a), $\Theta_\Pi G$ is the function on $\Conf(\theta\Gamma)$ given by
$(\Theta_\Pi G)(A,z)=\Theta_\Pi\left(G(\theta^{*}(A,z))\right)$, where $\theta^{*}(A,z):=(\theta^{*}A,\theta^{*}z)$;
its values lie in $\Gr(\theta Y)$.
Let $\int_YG$ denote the complex function on $\Conf(\Gamma)$ whose value at a point is the Berezin integral of the value of $G$ there.
Applied to the element $G(\theta^{*}(A,z))$ of $\Gr(Y)$, the identity just proved gives
\[
\int_{\theta Y}(\Theta_\Pi G)(A,z)=\overline{\left(\int_YG\right)\left(\theta^{*}(A,z)\right)}
\]
at every point $(A,z)$ of $\Conf(\theta\Gamma)$.

\emph{(f)}
Let $i\neq j$, $k:=k_{ij}$, and let $n,n'$ be the upward unit normals of $\Pi_i,\Pi_j$, with which $\Theta_{\Pi_i},\Theta_{\Pi_j}$ are formed.
The iteration formula~\eqref{eq:chamber-iteration} for the word $\Pi_i,\Pi_j,\dots,\Pi_i,\Pi_j$ of length $2k$
has the matrix $(\gamma_{n'}\gamma_n)^k$ and the element $(s_is_j)^k=\mathrm{id}$ of $W$;
with $U:=\gamma_{n'}\gamma_n$ it gives
\[
(\Theta_{\Pi_i}\Theta_{\Pi_j})^k\psi(x)=U^k\psi(x),\qquad
(\Theta_{\Pi_i}\Theta_{\Pi_j})^k\bar\psi(x)=\bar\psi(x)\,U^{-k}.
\]
The power is first computed with the inward unit normals $\hat n_i,\hat n_j$ of Definition~\ref{def:chambers} in the place of $n,n'$.
One has $\hat n_0\cdot\hat n_1=\hat n_2\cdot\hat n_3=-1/\sqrt{2}$, $\hat n_1\cdot\hat n_2=-1/2$, and $\hat n_i\cdot\hat n_j=0$ for the three remaining pairs;
thus $\hat n_i\cdot\hat n_j=-\cos(\pi/k)$, and the angle between $\hat n_j$ and $\hat n_i$ is $t:=\pi-\pi/k\in(0,\pi)$.
Write $\hat n_i=\cos t\cdot\hat n_j+\sin t\cdot n_\perp$ with a unit vector $n_\perp$ orthogonal to $\hat n_j$.
By the anticommutation relations of Lemma~\ref{lem:spin-identities},
$\gamma_{\hat n_j}^2=1$, $\gamma_{n_\perp}\gamma_{\hat n_j}=-\gamma_{\hat n_j}\gamma_{n_\perp}$ and $(\gamma_{\hat n_j}\gamma_{n_\perp})^2=-1$, and
\[
\gamma_{\hat n_j}\gamma_{\hat n_i}=\cos t+\sin t\cdot\gamma_{\hat n_j}\gamma_{n_\perp}.
\]
The real span of $1$ and $\gamma_{\hat n_j}\gamma_{n_\perp}$ is thus a commutative algebra isomorphic to $\C$, the element $\gamma_{\hat n_j}\gamma_{n_\perp}$ corresponding to the imaginary unit,
and de Moivre's formula gives
\[
(\gamma_{\hat n_j}\gamma_{\hat n_i})^k=\cos(kt)+\sin(kt)\cdot\gamma_{\hat n_j}\gamma_{n_\perp}=(-1)^{k-1},
\]
because $kt=(k-1)\pi$.
For the walls $\Pi_0,\Pi_1,\Pi_2$ the upward unit normal is the inward one, $\hat n_0$, $\hat n_1$, $\hat n_2$ respectively,
whereas the upward unit normal of $\Pi_3=(x_1,\tfrac12)$ is $e_1=-\hat n_3$.
Changing the sign of one of the two normals multiplies $U$ by $-1$ and $U^k$ by $(-1)^k$,
and $k_{03}=k_{13}=2$ and $k_{23}=4$ are even.
So $U^k=(-1)^{k-1}=(-1)^{k+1}$ in every case, and $(\Theta_{\Pi_i}\Theta_{\Pi_j})^k$ multiplies each generator by $(-1)^{k+1}$.
A product of an even number of antilinear anti-automorphisms is a linear automorphism of $\Gr(\mathbb{X})$,
and an automorphism of $\Gr(\mathbb{X})$ is determined by its values on the generators;
thus $(\Theta_{\Pi_i}\Theta_{\Pi_j})^k=(-1)^{(k+1)\deg}$ on $\Gr(\mathbb{X})$.
For an integrand $G$, the definition~\eqref{eq:setting-reflection-coefficients}, in the form $(\Theta_\Pi G)(\hat A)=\Theta_\Pi\left(G(\theta_\Pi^{*}\hat A)\right)$,
and the rule for the pull-back by a composition (Definition~\ref{def:cochains}) give
\[
\left((\Theta_{\Pi_i}\Theta_{\Pi_j})^kG\right)(\hat A)=(\Theta_{\Pi_i}\Theta_{\Pi_j})^k\left(G\left(((s_is_j)^k)^{*}\hat A\right)\right)=(-1)^{(k+1)\deg}\,G(\hat A),
\]
since $(s_is_j)^k=\mathrm{id}$; the relation therefore holds on integrands as well.

\emph{(f$'$)}
Let $\theta:=\theta_\Pi$ and $\theta':=\theta_{\Pi'}$, and note that $n''=n-2(n\cdot n')\,n'=r_{n'}n$.
The reflection in $\Pi''=\theta'(\Pi)$ is $\theta'\theta\theta'$,
and $n''$ is a unit normal of $\Pi''$, since $r_{n'}$ is the linear part of $\theta'$.
By~\eqref{eq:setting-two-reflections}, $\Theta_\Pi\Theta_{\Pi'}\psi(x)=\gamma_{n'}\gamma_n\,\psi(\theta\theta'x)$.
The matrices $\gamma_n$ and $\gamma_{n'}$ are Hermitian, so that $(\gamma_{n'}\gamma_n)^{*}=\gamma_n\gamma_{n'}$, and $\gamma_{n'}\gamma_n\gamma_{n'}=-\gamma_{r_{n'}n}$ (Lemma~\ref{lem:spin-identities}).
With these two identities, \eqref{eq:setting-reflection-fields} shows that the antilinear map $\Theta_{\Pi'}$ acts on the component $\alpha$ of this column as follows:
\[
\begin{aligned}
\Theta_{\Pi'}\left((\gamma_{n'}\gamma_n\psi(\theta\theta'x))_\alpha\right)
&=-\sum_\beta\overline{(\gamma_{n'}\gamma_n)_{\alpha\beta}}\cdot(\bar\psi(\theta'\theta\theta'x)\,\gamma_{n'})_\beta\\
&=-(\bar\psi(\theta'\theta\theta'x)\,\gamma_{n'}\gamma_n\gamma_{n'})_\alpha
=(\bar\psi(\theta'\theta\theta'x)\,\gamma_{n''})_\alpha .
\end{aligned}
\]
As $\Theta_{\Pi''}$ is formed with the normal $n''$, this says $\Theta_{\Pi'}\Theta_\Pi\Theta_{\Pi'}\psi(x)=-\Theta_{\Pi''}\psi(x)$.
For the barred field, \eqref{eq:setting-two-reflections} gives $\Theta_\Pi\Theta_{\Pi'}\bar\psi(x)=\bar\psi(\theta\theta'x)\,\gamma_n\gamma_{n'}$, and in the same way
\[
\begin{aligned}
\Theta_{\Pi'}\left((\bar\psi(\theta\theta'x)\,\gamma_n\gamma_{n'})_\alpha\right)
&=-\sum_\beta\overline{(\gamma_n\gamma_{n'})_{\beta\alpha}}\cdot(\gamma_{n'}\psi(\theta'\theta\theta'x))_\beta\\
&=-(\gamma_{n'}\gamma_n\gamma_{n'}\psi(\theta'\theta\theta'x))_\alpha
=(\gamma_{n''}\psi(\theta'\theta\theta'x))_\alpha
=-\Theta_{\Pi''}\bar\psi_\alpha(x).
\end{aligned}
\]
Both $\Theta_{\Pi'}\Theta_\Pi\Theta_{\Pi'}$, a product of three antilinear anti-automorphisms, and $(-1)^{\deg}\,\Theta_{\Pi''}$ are antilinear anti-automorphisms of $\Gr(\mathbb{X})$.
They agree on the generators of the free Grassmann algebra $\Gr(\mathbb{X})$ and therefore coincide.

\emph{(f$''$)}
Let $\Pi^{(1)},\dots,\Pi^{(r)}$ and $\Pi'^{(1)},\dots,\Pi'^{(s)}$ be mirrors with $\theta_{\Pi^{(1)}}\cdots\theta_{\Pi^{(r)}}=\theta_{\Pi'^{(1)}}\cdots\theta_{\Pi'^{(s)}}$.
The linear part of a reflection has the determinant $-1$; so $(-1)^r=(-1)^s$, and $r+s$ is even.
As $\Theta_{\Pi'^{(j)}}^2=1$ for every $j$ (Lemma~\ref{lem:sign-rules}(a)),
\[
\Theta_{\Pi^{(1)}}\cdots\Theta_{\Pi^{(r)}}\left(\Theta_{\Pi'^{(1)}}\cdots\Theta_{\Pi'^{(s)}}\right)^{-1}
=\Theta_{\Pi^{(1)}}\cdots\Theta_{\Pi^{(r)}}\,\Theta_{\Pi'^{(s)}}\cdots\Theta_{\Pi'^{(1)}}.
\]
This is a product of an even number, $r+s$, of reflections of Grassmann variables, to which the iteration formula~\eqref{eq:chamber-iteration} applies;
its element of $W$ is
\[
g=\theta_{\Pi^{(1)}}\cdots\theta_{\Pi^{(r)}}\,\theta_{\Pi'^{(s)}}\cdots\theta_{\Pi'^{(1)}}=\mathrm{id}.
\]
Let $n_1,\dots,n_{r+s}$ be the upward unit normals of the mirrors $\Pi^{(1)},\dots,\Pi^{(r)},\Pi'^{(s)},\dots,\Pi'^{(1)}$, in this order,
and let $U=\gamma_{n_{r+s}}\cdots\gamma_{n_1}$ be the matrix of~\eqref{eq:chamber-iteration}.
Applying $\gamma_n\gamma_v\gamma_n=-\gamma_{r_nv}$ (Lemma~\ref{lem:spin-identities}) $r+s$ times one obtains, for $v\in\R^3$,
\[
U\gamma_vU^{-1}=(-1)^{r+s}\,\gamma_{r_{n_{r+s}}\cdots r_{n_1}v}=\gamma_{\vec g^{-1}v}=\gamma_v ,
\]
since $r_{n_{r+s}}\cdots r_{n_1}$ is the inverse of the linear part $\vec g=r_{n_1}\cdots r_{n_{r+s}}$ of $g=\mathrm{id}$.
With $\sigma_v=\sum_\mu v_\mu\sigma_\mu$ as in the proof of~\eqref{it:lattice-covariance} one has $\gamma_n=\diag(\sigma_n,-\sigma_n)$;
as $r+s$ is even, the signs in the lower block cancel, and $U=\diag(u,u)$ with $u:=\sigma_{n_{r+s}}\cdots\sigma_{n_1}$.
The relation $U\gamma_vU^{-1}=\gamma_v$, $v\in\R^3$, says that $u$ commutes with $\sigma_1$, $\sigma_2$ and $\sigma_3$;
together with $1$ these matrices span all complex $2\times2$ matrices, so $u$ is a multiple of $1$.
Moreover $\det u=\prod_{j=1}^{r+s}\det\sigma_{n_j}=(-1)^{r+s}=1$, because $\sigma_n$ has the eigenvalues $1$ and $-1$ for every unit vector $n$.
Therefore
\begin{equation}\label{eq:chamber-central-block}
U=\diag(u,u),\qquad u\in\{1,-1\}.
\end{equation}
By~\eqref{eq:chamber-iteration} the linear automorphism $\Theta_{\Pi^{(1)}}\cdots\Theta_{\Pi^{(r)}}\,\Theta_{\Pi'^{(s)}}\cdots\Theta_{\Pi'^{(1)}}$ multiplies $\psi(x)$ by $u$ and $\bar\psi(x)$ by $u^{-1}=u$, for every $x\in\mathbb{X}$;
it is the automorphism $u^{\deg}$.
Multiplication from the right by $\Theta_{\Pi'^{(1)}}\cdots\Theta_{\Pi'^{(s)}}$ gives
\[
\Theta_{\Pi^{(1)}}\cdots\Theta_{\Pi^{(r)}}=u^{\deg}\,\Theta_{\Pi'^{(1)}}\cdots\Theta_{\Pi'^{(s)}}
\]
on $\Gr(\mathbb{X})$, with the one sign $u$.
On coefficient functions the $r+s$ maps together act through $g^{*}=\mathrm{id}$, the complex conjugations cancelling in pairs;
so the identity holds on integrands as well.

Every $g\in W$ is a product of reflections (Definition~\ref{def:mirrors}),
and on an even integrand $G$ the automorphism $u^{\deg}$ is the identity.
Hence the element $G^g$ of~\eqref{eq:setting-group-action} does not depend on the word with which $g$ is written.
If $g$ and $g'$ are written as words in reflections, the word of $g'$ followed by the word of $g$ is a word for $g'g$;
applying~\eqref{eq:setting-group-action} first with the word of $g$ and then with the word of $g'$ gives $(G^g)^{g'}=G^{g'g}$.
For a word of length $r$ the map $G\mapsto G^g$ is a composition of $r$ antilinear anti-automorphisms, and $(-1)^r=\det\vec g$;
it is therefore a linear automorphism if $g$ is even and an antilinear anti-automorphism if $g$ is odd.
Let $G$ be bosonic, regarded as a function of $F=d\hat A$.
Then~\eqref{eq:setting-reflection-coefficients} gives $(\Theta_\Pi G)(F)=\overline{G(\theta_\Pi^{*}F)}$, because $\theta_\Pi^{*}$ commutes with $d$.
For two mirrors $\Pi$ and $\Pi'$, with $\theta:=\theta_\Pi$ and $\theta':=\theta_{\Pi'}$,
\[
(\Theta_\Pi\Theta_{\Pi'}G)(F)=\overline{(\Theta_{\Pi'}G)(\theta^{*}F)}=G(\theta'^{*}\theta^{*}F)=G\left((\theta\theta')^{*}F\right),
\]
by $(\theta\theta')^{*}=\theta'^{*}\theta^{*}$ (Definition~\ref{def:cochains}).
Induction on the length of the word gives $G^g(F)=G(g^{*}F)$ for even $g$ and $G^g(F)=\overline{G(g^{*}F)}$ for odd $g$.

\emph{(g)}
It suffices to show, for every torus $\mathbb{T}=(\Gamma,\chi)$, every mirror $\Pi$ and every $G\in\mathcal{A}(\R^3)$ that embeds in $\mathbb{T}$,
that $\Theta_\Pi G$ is an element of $\mathcal{A}(\R^3)$ that embeds in $\theta_\Pi\mathbb{T}$, and that
\[
Z_{\theta_\Pi\mathbb{T}}[\Theta_\Pi G]=\overline{Z_{\mathbb{T}}[G]} .
\]
Indeed, let $g=\theta_\Pi g'$, where $g'$ is written with a word shorter by one.
Then $G^g=\Theta_\Pi(G^{g'})$ and $g\mathbb{T}=\theta_\Pi(g'\mathbb{T})$, so that
\[
Z_{g\mathbb{T}}[G^g]=Z_{\theta_\Pi(g'\mathbb{T})}[\Theta_\Pi(G^{g'})]=\overline{Z_{g'\mathbb{T}}[G^{g'}]},
\]
and induction on the length of the word gives the first assertion of~\eqref{it:functional-covariance}, the parities of $g$ and $g'$ being opposite.

Let $\theta:=\theta_\Pi$, with linear part $\vec\theta$, so that $\theta\Gamma=\vec\theta\,\Gamma$, and let $G\in\mathcal{A}(\R^3)$ embed in $\Gamma$.
Consider first a product $G=\Phi\cdot G_1\cdots G_r$ that embeds in $\Gamma$, with pairwise disjoint boxes $B_1,\dots,B_r$ and $G_j\in\mathcal{G}^{\mathrm{ev}}(B_j)$ (Definition~\ref{def:integrands}).
The definition~\eqref{eq:setting-reflection-coefficients}, in the form $(\Theta_\Pi G)(\hat A)=\Theta_\Pi\left(G(\theta^{*}\hat A)\right)$, gives
\[
\Theta_\Pi G=(\Theta_\Pi\Phi)\cdot(\Theta_\Pi G_1)\cdots(\Theta_\Pi G_r),\qquad(\Theta_\Pi\Phi)(F)=\overline{\Phi(\theta^{*}F)},
\]
because $\Theta_\Pi$ is antilinear and reverses products, the values of the $G_j$ are even and so commute with each other, and $d\theta^{*}=\theta^{*}d$;
here $\Theta_\Pi G_j\in\mathcal{G}^{\mathrm{ev}}(\theta B_j)$ by Lemma~\ref{lem:reflection-classes}(b).
The reflection $\theta$ maps every box $B$ onto the box $\theta B$, and
\[
\theta\left(B+[-1,1]^3+g\right)=\theta B+[-1,1]^3+\vec\theta g,\qquad g\in\Gamma,
\]
because the signed permutation matrix $\vec\theta$ maps the cube $[-1,1]^3$ onto itself (Definition~\ref{def:mirrors}).
As $\theta$ is bijective, the conditions of disjointness that define the embedding of $B_1,\dots,B_r$ in $\Gamma$ hold for $\theta B_1,\dots,\theta B_r$ in $\theta\Gamma$,
and $\Theta_\Pi G$ is a product that embeds in $\theta\Gamma$.
A general $G$ that embeds in $\Gamma$ is a finite sum of such products;
by antilinearity, $\Theta_\Pi G$ is then an element of $\mathcal{A}(\R^3)$ that embeds in $\theta\Gamma$.

Let $\Theta_\Pi$ act on functions on $\Conf(\Gamma)$ with values in $\Gr_{\mathbb{T}}$ as in Lemma~\ref{lem:reflected-algebras},
$(\Theta_\Pi f)(A,z)=\Theta_\Pi\left(f(\theta^{*}(A,z))\right)$ for $(A,z)\in\Conf(\theta\Gamma)$, with $\theta^{*}(A,z)=(\theta^{*}A,\theta^{*}z)$ as above.
The function on $\Conf(\theta\Gamma)$ defined by the element $\Theta_\Pi G$ (Definition~\ref{def:integrands}) is $\Theta_\Pi$ applied to the function on $\Conf(\Gamma)$ defined by $G$, that is,
\begin{equation}\label{eq:chamber-reflected-values}
(\Theta_\Pi G)(A,z)=\Theta_\Pi\left(G\left(\theta^{*}(A,z)\right)\right),\qquad(A,z)\in\Conf(\theta\Gamma).
\end{equation}
Both sides are antilinear in $G$, since the value at a point depends linearly on the element (Definition~\ref{def:integrands});
so it suffices to take for $G$ a product $\Phi\cdot G_1\cdots G_r$ as above.
Let $F$ be the curvature of $(A,z)$.
As $\theta^{*}$ commutes with $d$, the periodic extension of the curvature of $\theta^{*}(A,z)$ is $\theta^{*}\widetilde F$,
and the factor $\Theta_\Pi\Phi$ takes at $(A,z)$ the value $\overline{\Phi(\theta^{*}\widetilde F)}$, the complex conjugate of the value of $\Phi$ at $\theta^{*}(A,z)$.
By the last assertion of Lemma~\ref{lem:reflected-algebras}, the value of $\Theta_\Pi G_j$ at $(A,z)$ is $\Theta_\Pi\xi_j$, where $\xi_j\in\Gr_{\mathbb{T}}$ denotes the value of $G_j$ at $\theta^{*}(A,z)$.
The elements $\xi_j$ are even, and so are the $\Theta_\Pi\xi_j$;
as $\Theta_\Pi$ is antilinear and reverses products,
\[
\begin{aligned}
\Theta_\Pi\left(\Phi(\theta^{*}\widetilde F)\cdot\xi_1\cdots\xi_r\right)
&=\overline{\Phi(\theta^{*}\widetilde F)}\cdot\Theta_\Pi\xi_r\cdots\Theta_\Pi\xi_1\\
&=\overline{\Phi(\theta^{*}\widetilde F)}\cdot\Theta_\Pi\xi_1\cdots\Theta_\Pi\xi_r .
\end{aligned}
\]
By the product formula of Definition~\ref{def:integrands}, the left side is the image under $\Theta_\Pi$ of the value of $G$ at $\theta^{*}(A,z)$,
and the right side is the value of $\Theta_\Pi G$ at $(A,z)$.

The function $\exp(-S_{\mathbb{T}})$ on $\Conf(\Gamma)$ takes even values, which are central in $\Gr_{\mathbb{T}}$,
and $\Theta_\Pi(\exp(-S_{\mathbb{T}}))=\exp(-\Theta_\Pi S_{\mathbb{T}})=\exp(-S_{\theta\mathbb{T}})$ by~\eqref{it:action-invariance}, the exponential series having real coefficients.
As $\Theta_\Pi$ reverses products, \eqref{eq:chamber-reflected-values} gives
\[
\Theta_\Pi\left(\exp(-S_{\mathbb{T}})\cdot G\right)=\Theta_\Pi(G)\cdot\Theta_\Pi\left(\exp(-S_{\mathbb{T}})\right)=\exp(-S_{\theta\mathbb{T}})\cdot\Theta_\Pi G ,
\]
an identity between functions on $\Conf(\theta\Gamma)$ with values in $\Gr_{\theta\mathbb{T}}$.
Let $Y$ be a set of representatives of $\mathbb{X}_{\mathbb{T}}$; then $\theta Y$ is a set of representatives of $\mathbb{X}_{\theta\mathbb{T}}$.
The homomorphisms $[\,\cdot\,]_{\mathbb{T}}$ and $[\,\cdot\,]_{\theta\mathbb{T}}$ of Lemma~\ref{lem:reflected-algebras} restrict to isomorphisms $\Gr(Y)\to\Gr_{\mathbb{T}}$ and $\Gr(\theta Y)\to\Gr_{\theta\mathbb{T}}$,
which carry $\int_Y$ to $\int_{\mathbb{T}}$ and $\int_{\theta Y}$ to $\int_{\theta\mathbb{T}}$ (Definition~\ref{def:tori});
by the same lemma they carry $\Theta_\Pi\colon\Gr(Y)\to\Gr(\theta Y)$ to $\Theta_\Pi\colon\Gr_{\mathbb{T}}\to\Gr_{\theta\mathbb{T}}$.
Through these isomorphisms, the identity for functions established in the proof of~\eqref{it:berezin-conjugate} applies to $\exp(-S_{\mathbb{T}})\cdot G$ and gives
\[
\int_{\theta\mathbb{T}}\exp\left(-S_{\theta\mathbb{T}}(A)\right)\cdot(\Theta_\Pi G)(A,z)
=\overline{\left(\int_{\mathbb{T}}\exp(-S_{\mathbb{T}})\cdot G\right)\left(\theta^{*}(A,z)\right)} .
\]
Further, $F(\theta^{*}(A,z))=\theta^{*}F(A,z)$, as $\theta^{*}$ commutes with $d$;
$\|\theta^{*}F\|=\|F\|$, as $\theta$ induces a bijection of the plaquettes modulo $\Gamma$ onto the plaquettes modulo $\theta\Gamma$;
and $\theta^{*}\colon\Conf(\theta\Gamma)\to\Conf(\Gamma)$ preserves the measure (Definition~\ref{def:configurations}).
The integral that defines $Z_{\mathbb{T}}[G]$ converges absolutely (Lemma~\ref{lem:gaussian-growth}),
and the substitution $(A,z)\mapsto\theta^{*}(A,z)$ gives
\[
\begin{aligned}
Z_{\theta\mathbb{T}}[\Theta_\Pi G]
&=\int_{\Conf(\theta\Gamma)}\mathrm{d}(A,z)\,\exp\left(-\frac{1}{2}\left\|F\left(\theta^{*}(A,z)\right)\right\|^2\right)\\
&\qquad\cdot\overline{\left(\int_{\mathbb{T}}\exp(-S_{\mathbb{T}})\cdot G\right)\left(\theta^{*}(A,z)\right)}\\
&=\overline{Z_{\mathbb{T}}[G]} .
\end{aligned}
\]
Let now $G\in\mathcal{B}(\mathbb{T})$, even or odd.
The function $\Theta_\Pi G$ of Lemmas~\ref{lem:reflection-classes}(a) and~\ref{lem:reflected-algebras} belongs to $\mathcal{B}(\theta\mathbb{T})$,
and for it \eqref{eq:chamber-reflected-values} holds by definition;
the central element $\exp(-S_{\mathbb{T}})$ commutes with $G$,
and the integrals converge absolutely by Lemma~\ref{lem:bounded-integrands}.
So the computation above gives $Z_{\theta\mathbb{T}}[\Theta_\Pi G]=\overline{Z_{\mathbb{T}}[G]}$.

Let $\Phi$ be a real bounded Borel function of finitely many $F_\sigma$, and let $g\in W$.
As a product without boxes, $\Phi$ is an element of $\mathcal{A}(\R^3)$ that embeds in every period lattice;
and $\Phi^g=\Phi\circ g^{*}$ for both parities of $g$, by~\eqref{it:reflection-words} and because $\Phi$ is real.
By the first assertion of~\eqref{it:functional-covariance}, applied to $\Phi$ and to $1$,
the number $\mathbb{E}_{g\mathbb{T}}[\Phi\circ g^{*}]$ equals $\mathbb{E}_{\mathbb{T}}[\Phi]$ if $g$ is even and $\overline{\mathbb{E}_{\mathbb{T}}[\Phi]}$ if $g$ is odd.
These are the same real number, since $\mathbb{E}_{\mathbb{T}}[\Phi]=\int\mathrm{d}\mu_{\mathbb{T}}(F)\,\Phi(F)$ with the law $\mu_{\mathbb{T}}$ of Definition~\ref{def:functional}.
Hence
\[
\int\mathrm{d}\mu_{g\mathbb{T}}(F)\,\Phi(g^{*}F)=\int\mathrm{d}\mu_{\mathbb{T}}(F)\,\Phi(F).
\]
The map $g^{*}$ is a homeomorphism of $\Omega$ with the inverse $(g^{-1})^{*}$,
and as $(g^{*}F)_\sigma=\pm F_{g\sigma}$ (Definition~\ref{def:cochains}), composition with $g^{*}$ or with $(g^{-1})^{*}$ turns a function of finitely many $F_\sigma$ into a function of finitely many $F_\sigma$.
So every real bounded Borel function $\Phi'$ of finitely many $F_\sigma$ is of the form $\Phi\circ g^{*}$ with $\Phi:=\Phi'\circ(g^{-1})^{*}$ of the same kind,
and the last identity says that $\mu_{g\mathbb{T}}$ and the image of $\mu_{\mathbb{T}}$ under the map $F\mapsto(g^{-1})^{*}F$ agree on all such $\Phi'$.
The cylinder sets are closed under finite intersections and generate the Borel $\sigma$-algebra of $\Omega$;
hence $\mu_{g\mathbb{T}}$ is the image of $\mu_{\mathbb{T}}$ under $F\mapsto(g^{-1})^{*}F$.
For $v\in\Z^3$ the translation $\tau_v$ has the linear part $1$, so that $\tau_v\mathbb{T}=\mathbb{T}$, and $\mu_{\mathbb{T}}$ is invariant under $\tau_v$.
For the cubic torus $C(n)$ one has $\Gamma=n\Z^3$ and $\chi(nv)=(-1)^{v_1+v_2+v_3}$, $v\in\Z^3$.
Every signed permutation matrix $\vec g\in\Ocub$ maps $n\Z^3$ onto itself and preserves the parity of $v_1+v_2+v_3$;
hence $gC(n)=C(n)$ for all $g\in W$, and $\mu_{C(n)}$ is invariant under $W$.

\emph{(h)}
The product $\chi\chi'$ is again a character of $\Gamma$ with the values $\pm1$.
Let $g_1,g_2,g_3$ be a basis of $\Gamma$, and let $k\colon\Gamma\to\Z$ be the homomorphism with $k(g_j)=0$ if $\chi\chi'(g_j)=1$ and $k(g_j)=1$ if $\chi\chi'(g_j)=-1$.
Then $(-1)^{k(g)}=\chi(g)\cdot\chi'(g)$ for all $g\in\Gamma$, since both sides are characters of $\Gamma$ that agree on the basis.
Choose a set $\mathcal{R}$ of representatives of $\Z^3/\Gamma$.
Every $x\in\Z^3$ is uniquely of the form $x=y+g$ with $y\in\mathcal{R}$ and $g\in\Gamma$; put $\lambda(x):=k(g)$.
This defines an integer-valued function $\lambda$ on $\Z^3$ with
\[
\lambda(x+g)=\lambda(x)+k(g),\qquad x\in\Z^3,\ g\in\Gamma,
\]
so that $d\lambda$ is an integer-valued, closed, $\Gamma$-periodic one-cochain.

The map
\[
\mathcal{U}\colon(A,z)\mapsto\left(A+\frac{\pi}{q}\,d\lambda,\ z\right)
\]
adds the same cochain to the first component of every representative of a configuration;
so it respects the equivalence of Definition~\ref{def:configurations} and is a bijection of $\Conf(\Gamma)$, with the inverse $(A,z)\mapsto\left(A-\tfrac{\pi}{q}d\lambda,z\right)$.
It maps the fundamental domain formed by the box $[0,2\pi/q)^{3\cdot[\Z^3:\Gamma]}$ and the closed $z$ onto its translate by $\tfrac{\pi}{q}d\lambda$, which is again a measurable fundamental domain;
as Lebesgue measure is invariant under translations, $\mathcal{U}$ preserves the measure of $\Conf(\Gamma)$ (Definition~\ref{def:configurations}).
Since $d(d\lambda)=0$, one has $F\circ\mathcal{U}=F$; in particular every bounded function $\Phi$ of $\widetilde F$ is invariant under $\mathcal{U}$.

Let $\alpha_\lambda$ be the automorphism of $\Gr(\mathbb{X})$ with
\[
\alpha_\lambda\psi(x)=(-1)^{\lambda(\pi(x))}\psi(x),\qquad
\alpha_\lambda\bar\psi(x)=(-1)^{\lambda(\pi(x))}\bar\psi(x),\qquad x\in\mathbb{X}.
\]
For $g\in\Gamma$ one has $\pi(x+g)=\pi(x)+g$ (Proposition~\ref{prop:geometry}\eqref{it:geometry-projection}).
Hence $\alpha_\lambda$ maps the relation $\psi(x+g)=\chi(g)\cdot\psi(x)$ to
$(-1)^{\lambda(\pi(x))+k(g)}\psi(x+g)=\chi(g)\cdot(-1)^{\lambda(\pi(x))}\psi(x)$,
that is, to $\psi(x+g)=\chi'(g)\cdot\psi(x)$; the same holds for $\bar\psi$.
The relations of $\Gr_{\mathbb{T}}$ (Definition~\ref{def:tori}) are thus mapped to those of $\Gr_{\mathbb{T}'}$.
Let $[\,\cdot\,]_{\mathbb{T}}\colon\Gr(\mathbb{X})\to\Gr_{\mathbb{T}}$ and $[\,\cdot\,]_{\mathbb{T}'}\colon\Gr(\mathbb{X})\to\Gr_{\mathbb{T}'}$ be the homomorphisms of algebras that send every generator to its class;
their kernels are the two-sided ideals generated by the respective relations, and the automorphism $\alpha_\lambda$ maps the first kernel into the second.
Hence $\alpha_\lambda$ induces a homomorphism of algebras $\Gr_{\mathbb{T}}\to\Gr_{\mathbb{T}'}$, again denoted $\alpha_\lambda$, with $\alpha_\lambda[\xi]_{\mathbb{T}}=[\alpha_\lambda\xi]_{\mathbb{T}'}$ for $\xi\in\Gr(\mathbb{X})$;
it is an isomorphism, its inverse being induced in the same way, with $\chi$ and $\chi'$ exchanged.
It preserves the degree and fixes every pair $\psi_\alpha(x)\,\bar\psi_\alpha(x)$,
hence the product of the pairs at a set of representatives of $\mathbb{X}_{\mathbb{T}}=\mathbb{X}_{\mathbb{T}'}$.
Therefore
\[
\int_{\mathbb{T}'}\alpha_\lambda(\xi)=\int_{\mathbb{T}}\xi,\qquad\xi\in\Gr_{\mathbb{T}}.
\]

Let $A$ be a $\Gamma$-periodic one-cochain.
In $\alpha_\lambda(S_{\mathbb{T}}(A))$ the mass terms are unchanged,
and the term of an oriented bond $(x,x')$ is multiplied by
\[
(-1)^{\lambda(\pi(x))+\lambda(\pi(x'))}=(-1)^{\lambda(\pi(x'))-\lambda(\pi(x))}=\exp\left(iq\cdot\frac{\pi}{q}\cdot(d\lambda)\left(\pi(x),\pi(x')\right)\right),
\]
where $(d\lambda)(\pi(x),\pi(x'))=\lambda(\pi(x'))-\lambda(\pi(x))$ as in the proof of~\eqref{it:gauge-covariance}.
This is the factor by which $U_A(x,x')$ changes when $A$ is replaced by $A+\tfrac{\pi}{q}d\lambda$.
Hence
\[
\alpha_\lambda\left(S_{\mathbb{T}}(A)\right)=S_{\mathbb{T}'}\left(A+\frac{\pi}{q}\,d\lambda\right).
\]

Let $B_j$ be a box that embeds in $\Gamma$, let $G_j\in\mathcal{G}^{\mathrm{ev}}(B_j)$, with $G_j(\hat A)=\sum_\Psi a_\Psi(\hat A)\cdot\Psi$ in a monomial basis of $\Gr(\mathbb{X}(B_j))$,
and let $\hat A$ be a local lift on $B_j$ at $(A,z)$ (Definition~\ref{def:integrands}).
Put $\lambda':=-\tfrac{\pi}{q}\lambda$.
Then $\hat A-d\lambda'=\hat A+\tfrac{\pi}{q}d\lambda$ is a local lift on $B_j$ at $\mathcal{U}(A,z)$:
its coboundary is $d\hat A=F$, and its link phases are those of $A+\tfrac{\pi}{q}d\lambda$.
Moreover $\exp(\pm iq\lambda'(y))=(-1)^{\lambda(y)}$ at every photon site $y$, as $\lambda$ is integer-valued.
The gauge invariance of $G_j$ (Definition~\ref{def:integrands}), applied with $\lambda'$ on the sites of $\Z^3$ contained in $B_j$, therefore gives
\[
a_\Psi(\hat A-d\lambda')=\exp\left(-iq\langle\lambda',\mathfrak{q}(\Psi)\rangle\right)\cdot a_\Psi(\hat A)=(-1)^{\langle\lambda,\mathfrak{q}(\Psi)\rangle}\cdot a_\Psi(\hat A)
\]
for every monomial $\Psi$ of the basis, $\mathfrak{q}(\Psi)$ being integer-valued.
The sign on the right is the factor by which $\alpha_\lambda$ multiplies $\Psi$.
Indeed, $\alpha_\lambda$ contributes a factor $(-1)^{\lambda(y)}$ for each factor $\psi_\alpha(x)$ or $\bar\psi_\alpha(x)$ of $\Psi$ with $\pi(x)=y$,
and the number of these factors differs from $\mathfrak{q}(\Psi)(y)$ by twice the number of the barred ones among them.
Consequently
\[
G_j(\hat A-d\lambda')=\alpha_\lambda\left(G_j(\hat A)\right).
\]
By Definition~\ref{def:integrands}, the value of $G_j$ on $\mathbb{T}'$ at $\mathcal{U}(A,z)$ is $[G_j(\hat A-d\lambda')]_{\mathbb{T}'}=[\alpha_\lambda(G_j(\hat A))]_{\mathbb{T}'}=\alpha_\lambda[G_j(\hat A)]_{\mathbb{T}}$,
the image under $\alpha_\lambda$ of the value $[G_j(\hat A)]_{\mathbb{T}}$ of $G_j$ on $\mathbb{T}$ at $(A,z)$.
The same holds for the translate of $G_j$ to the box $B_j+g$, $g\in\Gamma$, which defines the same elements of $\Gr_{\mathbb{T}}$ and $\Gr_{\mathbb{T}'}$:
on the sites of $B_j+g$ the function $\lambda$ exceeds its values at the corresponding sites of $B_j$ by the constant $k(g)$,
and a constant sign on all generators acts trivially on even elements.

Now let $G=\Phi\cdot G_1\cdots G_r$ be a product that embeds in $\Gamma$, with $G_j\in\mathcal{G}^{\mathrm{ev}}(B_j)$ and $\Phi$ a bounded Borel function of finitely many $F_\sigma$.
Since $\alpha_\lambda$ is a homomorphism of algebras, the identities for $\alpha_\lambda(S_{\mathbb{T}}(A))$ and for the values of the factors $G_j$,
together with the invariance of $\Phi$, show that $\exp(-S_{\mathbb{T}'})\cdot G$ at the point $\mathcal{U}(A,z)$ is the image under $\alpha_\lambda$ of $\exp(-S_{\mathbb{T}})\cdot G$ at $(A,z)$.
The invariance of the Berezin integral under $\alpha_\lambda$ then gives
\[
\left(\int_{\mathbb{T}}\exp(-S_{\mathbb{T}})\cdot G\right)(A,z)=\left(\int_{\mathbb{T}'}\exp(-S_{\mathbb{T}'})\cdot G\right)\left(\mathcal{U}(A,z)\right).
\]
Multiply by $\exp\left(-\frac{1}{2}\|F\|^2\right)$, which is invariant under $\mathcal{U}$, and integrate over $\Conf(\Gamma)$;
the integrals converge absolutely (Lemma~\ref{lem:gaussian-growth}), and $\mathcal{U}$ preserves the measure.
The result is $Z_{\mathbb{T}}[G]=Z_{\mathbb{T}'}[G]$, and by linearity this holds for every $G\in\mathcal{A}(\R^3)$ that embeds in $\Gamma$.
In particular $Z_{\mathbb{T}}[1]=Z_{\mathbb{T}'}[1]$, and $\mathbb{E}_{\mathbb{T}}[G]=\mathbb{E}_{\mathbb{T}'}[G]$ for all these $G$.
The bounded Borel functions of finitely many $F_\sigma$ are among them,
so the laws $\mu_{\mathbb{T}}$ and $\mu_{\mathbb{T}'}$ agree on the cylinder sets and coincide.
\end{proof}

\subsection{Accretivity}\label{sec:accretivity}

For $n\in\{2,4\}$ the scalar product on $\ell^2(\mathbb{X};\C^n)$ is
$\langle\psi',\psi\rangle=\sum_{x\in\mathbb{X}}\psi'(x)^{*}\psi(x)$;
like all scalar products below, it is antilinear in the first argument.
For a subset $Y$ of $\mathbb{X}$, let $\ell^2(Y;\C^n)$ be the subspace of the fields in $\ell^2(\mathbb{X};\C^n)$ that vanish outside $Y$, and let $P_Y$, for both values of $n$, be the orthogonal projection onto it.
An operator on $\ell^2(\mathbb{X};\C^n)$ is compressed to $Y$ by multiplying it by $P_Y$ on both sides and restricting the product to $\ell^2(Y;\C^n)$.

Let $\Gamma$ be a subgroup of $\Z^3$ of rank three, let $\chi$ be a homomorphism from $\Gamma$ to the unit circle, and put $\mathbb{T}=(\Gamma,\chi)$ and $\mathbb{X}_{\mathbb{T}}:=\mathbb{X}/\Gamma$.
A spinor of $\mathbb{T}$ is a map $\psi\colon\mathbb{X}\to\C^4$ with $\psi(x+v)=\chi(v)\,\psi(x)$ for all $x\in\mathbb{X}$ and $v\in\Gamma$.
For two spinors $\psi$, $\psi'$ of $\mathbb{T}$ the function $x\mapsto\psi'(x)^{*}\psi(x)$ is invariant under the translations by $\Gamma$, because $\overline{\chi(v)}\cdot\chi(v)=1$.
Hence
\[
\langle\psi',\psi\rangle:=\sum_{x\in\mathbb{X}_{\mathbb{T}}}\psi'(x)^{*}\psi(x),
\]
the sum being taken over a set of representatives of the classes, is a scalar product, and the spinors of $\mathbb{T}$ form a Hilbert space of dimension $4\cdot\#|\mathbb{X}_{\mathbb{T}}|$.
This dimension is finite: $\mathbb{X}$ has $48$ sites in each unit cell (Definition~\ref{def:fermion-sites}), and $\Gamma$ has finite index in $\Z^3$.
Let $A$ be a real $\Gamma$-periodic one-cochain on $\Z^3$, and let $v\in\Gamma$.
If $x=gv_0$ with $g\in W$, then $x+v=(\tau_vg)v_0$, and the definition of the neighbours gives $(x+v)^{(i)}=(\tau_vgs_i)v_0=x^{(i)}+v$;
moreover $\pi(x+v)=\pi(x)+v$ (Definition~\ref{def:fermion-sites}).
The bond $\left\{x+v,(x+v)^{(i)}\right\}$ is thus the translate of $\left\{x,x^{(i)}\right\}$, with the same type and the same direction, and
\[
A\left(\pi(x+v),\pi\left((x+v)^{(i)}\right)\right)=A\left(\pi(x),\pi\left(x^{(i)}\right)\right),
\]
$A$ being $\Gamma$-periodic.
It follows that the coefficients $w\left(x,x^{(i)}\right)U_A\left(x,x^{(i)}\right)$ of $D_A$ are invariant under the translations by $\Gamma$, and that $D_A$ maps spinors of $\mathbb{T}$ to spinors of $\mathbb{T}$.
The operator so obtained is denoted $D^{\mathbb{T}}_A$.
When $\Gamma$ is a period lattice and $\chi$ takes the values $\pm1$, the pair $\mathbb{T}$ is a torus, and its spinors and the operator $D^{\mathbb{T}}_A$ are those of Definitions~\ref{def:tori} and~\ref{def:torus-spinors}.

\begin{definition}[The weighted covariant Laplacian]\label{def:covariant-laplacian}
Let the weights $w_0,\dots,w_3>0$ and the charge $q>0$ be as in Definition~\ref{def:chamber-operator}, and let $A$ be a real one-cochain on $\Z^3$.
The weighted covariant difference $\nabla_A$ assigns to a field $\psi\colon\mathbb{X}\to\C^4$ the function on $\mathbb{X}\times\{0,1,2,3\}$ with values in $\C^4$ given by
\[
(\nabla_A\psi)(x,i):=\sqrt{\frac{w_i}{2}}\,\left(\psi(x)-U_A\left(x,x^{(i)}\right)\psi\left(x^{(i)}\right)\right).
\]
On $\ell^2(\mathbb{X};\C^4)$ it is a bounded operator into $\ell^2(\mathbb{X}\times\{0,1,2,3\};\C^4)$, of norm at most $\left(2\sum_{i}w_i\right)^{\frac12}$, since each map $x\mapsto x^{(i)}$ is a bijection of $\mathbb{X}$.
The operator $\nabla_A^{*}\nabla_A$ on $\ell^2(\mathbb{X};\C^4)$ is the \emph{weighted covariant Laplacian}.
If $\mathbb{T}=(\Gamma,\chi)$ is a pair as above and $A$ is $\Gamma$-periodic, then $\nabla_A$ maps the spinors of $\mathbb{T}$ to the functions on $\mathbb{X}\times\{0,1,2,3\}$ that are multiplied by $\chi(v)$ when $x$ is replaced by $x+v$, $v\in\Gamma$;
these functions form a Hilbert space for the norm whose square is the sum of the squared moduli over $\mathbb{X}_{\mathbb{T}}\times\{0,1,2,3\}$, and the operator $\nabla_A^{*}\nabla_A$ on the spinors of $\mathbb{T}$ is the weighted covariant Laplacian of $\mathbb{T}$.
In both cases
\begin{equation}\label{eq:chamber-laplacian}
\left\langle\psi,\nabla_A^{*}\nabla_A\psi\right\rangle=\frac12\sum_{x}\sum_{i=0}^{3}w_i\left|\psi(x)-U_A\left(x,x^{(i)}\right)\psi\left(x^{(i)}\right)\right|^2,
\end{equation}
where $x$ ranges over $\mathbb{X}$, respectively over $\mathbb{X}_{\mathbb{T}}$.
On $\ell^2(\mathbb{X};\C^4)$ the right-hand side equals the sum of $w_i\left|\psi(x)-U_A\left(x,x^{(i)}\right)\psi\left(x^{(i)}\right)\right|^2$ over all bonds $\left\{x,x^{(i)}\right\}$, $i$ being the type of the bond:
the pairs $(x,i)$ and $\left(x^{(i)},i\right)$ belong to the same bond and contribute the same term, because $U_A\left(x^{(i)},x\right)=\overline{U_A\left(x,x^{(i)}\right)}$ has modulus one.
For $A=0$ the operator $\nabla_0^{*}\nabla_0$ is the weighted Laplacian of the fermion graph.
\end{definition}

\begin{proposition}[Accretivity and the sign of the determinant]\label{prop:accretivity}
Let $\rho$ and $M$ be real, $w_0,\dots,w_3>0$, $q>0$, and $m=M-\frac{\rho}{2}\sum_{i=0}^{3}w_i$.
Let $A$ be a real one-cochain on $\Z^3$.
Let $\mathcal{H}$ be either $\ell^2(\mathbb{X};\C^4)$, or the space of the spinors of a pair $\mathbb{T}=(\Gamma,\chi)$ formed by a subgroup $\Gamma$ of $\Z^3$ of rank three and a homomorphism $\chi$ from $\Gamma$ to the unit circle; in the second case $A$ is assumed to be $\Gamma$-periodic.
Define $H_A$ and $J_A$ by
\[
(H_A\psi)(x):=\sum_{i=0}^{3}w_i\,U_A\left(x,x^{(i)}\right)\psi\left(x^{(i)}\right),
\qquad
(J_A\psi)(x):=\sum_{i=0}^{3}w_i\,\gamma\cdot\hat d_i(x)\,U_A\left(x,x^{(i)}\right)\psi\left(x^{(i)}\right).
\]
\begin{enumerate}
\item
$H_A$ and $J_A$ are bounded operators on $\mathcal{H}$, $H_A$ is Hermitian, $J_A$ is anti-Hermitian, and
\[
D_A=M-\frac{\rho}{2}H_A+\frac12J_A.
\]
For every subset $Y$ of $\mathbb{X}$, the compressions $P_YH_AP_Y$ and $P_YJ_AP_Y$ of the operators on $\ell^2(\mathbb{X};\C^4)$ are Hermitian and anti-Hermitian, respectively, and $P_YD_AP_Y=M-\frac{\rho}{2}P_YH_AP_Y+\frac12P_YJ_AP_Y$ on $\ell^2(Y;\C^4)$.
\item
The real part $\Re D_A:=\frac12\left(D_A+D_A^{*}\right)$ satisfies
\[
\Re D_A=M-\frac{\rho}{2}H_A=m+\frac{\rho}{2}\nabla_A^{*}\nabla_A,
\]
and $\Re\left(P_YD_AP_Y\right)=m+\frac{\rho}{2}P_Y\nabla_A^{*}\nabla_AP_Y$ on $\ell^2(Y;\C^4)$ for every subset $Y$ of $\mathbb{X}$.
If $\rho\ge0$, then $\Re D_A\ge m$ and $\Re\left(P_YD_AP_Y\right)\ge m$; that is, $D_A-m$ and $P_YD_AP_Y-m$ are accretive.
\item\label{it:operator-norm}
The norm of $D_A$ on $\mathcal{H}$ satisfies
\[
\|D_A\|\le|M|+\frac{1+|\rho|}{2}\sum_{i=0}^{3}w_i,
\]
and, for every subset $Y$ of $\mathbb{X}$, the same bound holds for the norm on $\ell^2(Y;\C^4)$ of the compression $P_YD_AP_Y$ of the operator on $\ell^2(\mathbb{X};\C^4)$.
\item\label{it:determinant-sign}
For $v\in\R^3$ let $\sigma\cdot v:=\sum_{\mu}v_\mu\sigma_\mu$.
Let $\check H_A$ and $\check J_A$ be given by the formulas for $H_A$ and $J_A$, with $\sigma\cdot\hat d_i(x)$ in place of $\gamma\cdot\hat d_i(x)$,
on the fields with values in $\C^2$:
on $\ell^2(\mathbb{X};\C^2)$, respectively on the maps $\psi\colon\mathbb{X}\to\C^2$ with $\psi(x+v)=\chi(v)\,\psi(x)$ for $v\in\Gamma$.
Put $\check D_A:=M-\frac{\rho}{2}\check H_A+\frac12\check J_A$;
on the second space $\check D_A$ is written $\check D^{\mathbb{T}}_A$.
In the decomposition $\C^4=\C^2\oplus\C^2$ in which $\gamma_\mu=\diag(\sigma_\mu,-\sigma_\mu)$,
\[
D_A=\diag\left(\check D_A,\check D_A^{*}\right).
\]
Consequently $\det D^{\mathbb{T}}_A=\left|\det\check D^{\mathbb{T}}_A\right|^2\ge0$ for every pair $\mathbb{T}$ as above, in particular for every torus.
For every subset $Y$ of $\mathbb{X}$, the compressions to $Y$ of $D_A$ on $\ell^2(\mathbb{X};\C^4)$ and of $\check D_A$ on $\ell^2(\mathbb{X};\C^2)$ are related by
\[
P_YD_AP_Y=\diag\left(P_Y\check D_AP_Y,\left(P_Y\check D_AP_Y\right)^{*}\right)
\]
on $\ell^2(Y;\C^4)=\ell^2(Y;\C^2)\oplus\ell^2(Y;\C^2)$;
if $Y$ is finite, then $\det\left(P_YD_AP_Y\right)=\left|\det\left(P_Y\check D_AP_Y\right)\right|^2\ge0$.
\end{enumerate}
\end{proposition}

\begin{proof}
Write $\sum_x$ for the sum over $\mathbb{X}$ if $\mathcal{H}=\ell^2(\mathbb{X};\C^4)$, and for the sum over $\mathbb{X}_{\mathbb{T}}$ if $\mathcal{H}$ is the space of the spinors of $\mathbb{T}$.
For each type $i$ the map $x\mapsto x^{(i)}$ is an involution of $\mathbb{X}$, and a bond has the same type seen from either of its ends (Definition~\ref{def:fermion-sites}).
As $(x+v)^{(i)}=x^{(i)}+v$ for $v\in\Gamma$, the map also induces an involution of $\mathbb{X}_{\mathbb{T}}$.
The unit vector from $x^{(i)}$ to $x$ is $\hat d_i\left(x^{(i)}\right)=-\hat d_i(x)$, and, $A$ being real and odd under the exchange of its two arguments,
\[
U_A\left(x^{(i)},x\right)=\overline{U_A\left(x,x^{(i)}\right)},\qquad\left|U_A\left(x,x^{(i)}\right)\right|=1.
\]
When $A$ is $\Gamma$-periodic, the coefficients of $H_A$ and $J_A$, like those of $D_A$, are invariant under the translations by $\Gamma$, so that both operators map spinors of $\mathbb{T}$ to spinors of $\mathbb{T}$.

\emph{Hermitian and anti-Hermitian part.}
In the $i$-th term of $H_A$ the value $\psi\left(x^{(i)}\right)$ is multiplied by $w_i\cdot U_A\left(x,x^{(i)}\right)$, in the $i$-th term of $J_A$ by $w_i\,\gamma\cdot\hat d_i(x)\,U_A\left(x,x^{(i)}\right)$.
Both are matrices of norm $w_i$: by Definition~\ref{def:spin-matrices}, $\gamma\cdot\hat d_i(x)$ is Hermitian with square $1$, hence unitary.
Since $x\mapsto x^{(i)}$ is a bijection, every term maps $\mathcal{H}$ to $\mathcal{H}$ with norm at most $w_i$, and $\|H_A\|$, $\|J_A\|$ are at most $\sum_iw_i$.
Let $\psi,\psi'\in\mathcal{H}$.
On $\ell^2(\mathbb{X};\C^4)$ the sums below converge absolutely, by the Cauchy--Schwarz inequality.
On the spinors of $\mathbb{T}$ each of their terms is invariant under the replacement of $x$ by $x+v$, $v\in\Gamma$, because the factors $\overline{\chi(v)}$ and $\chi(v)$ cancel; so the sums over $\mathbb{X}_{\mathbb{T}}$ are well defined.
In both cases $x^{(i)}$ may be substituted for $x$ in the $i$-th sum.
This substitution and the relations above give
\[
\begin{aligned}
\left\langle\psi',H_A\psi\right\rangle
&=\sum_x\sum_{i}w_i\,\psi'(x)^{*}U_A\left(x,x^{(i)}\right)\psi\left(x^{(i)}\right)\\
&=\sum_x\sum_{i}w_i\,\psi'\left(x^{(i)}\right)^{*}\overline{U_A\left(x,x^{(i)}\right)}\,\psi(x)
=\left\langle H_A\psi',\psi\right\rangle,
\end{aligned}
\]
and, with $\gamma\cdot\hat d_i\left(x^{(i)}\right)=-\gamma\cdot\hat d_i(x)$ and $\left(\gamma\cdot\hat d_i(x)\right)^{*}=\gamma\cdot\hat d_i(x)$,
\[
\begin{aligned}
\left\langle\psi',J_A\psi\right\rangle
&=\sum_x\sum_{i}w_i\,\psi'(x)^{*}\gamma\cdot\hat d_i(x)\,U_A\left(x,x^{(i)}\right)\psi\left(x^{(i)}\right)\\
&=-\sum_x\sum_{i}w_i\,\psi'\left(x^{(i)}\right)^{*}\gamma\cdot\hat d_i(x)\,\overline{U_A\left(x,x^{(i)}\right)}\,\psi(x)
=-\left\langle J_A\psi',\psi\right\rangle.
\end{aligned}
\]
Therefore $H_A^{*}=H_A$ and $J_A^{*}=-J_A$.
For a subset $Y$ of $\mathbb{X}$ it follows that $\left(P_YH_AP_Y\right)^{*}=P_YH_A^{*}P_Y=P_YH_AP_Y$ and $\left(P_YJ_AP_Y\right)^{*}=-P_YJ_AP_Y$.
By Definition~\ref{def:chamber-operator} the hopping matrix of a bond of type $i$ is $w\left(x,x^{(i)}\right)=\frac{w_i}{2}\left(\rho-\gamma\cdot\hat d_i(x)\right)$, so that
\[
(D_A\psi)(x)=M\psi(x)-\frac{\rho}{2}\sum_{i}w_i\,U_A\left(x,x^{(i)}\right)\psi\left(x^{(i)}\right)+\frac12\sum_{i}w_i\,\gamma\cdot\hat d_i(x)\,U_A\left(x,x^{(i)}\right)\psi\left(x^{(i)}\right),
\]
which is $D_A=M-\frac{\rho}{2}H_A+\frac12J_A$.
Multiplication by $P_Y$ on both sides gives the formula for $P_YD_AP_Y$, as $P_Y$ is the identity of $\ell^2(Y;\C^4)$.

\emph{The real part.}
The operator $M-\frac{\rho}{2}H_A$ is Hermitian and $\frac12J_A$ is anti-Hermitian, $M$ and $\rho$ being real.
So $D_A^{*}=M-\frac{\rho}{2}H_A-\frac12J_A$, and $\Re D_A=M-\frac{\rho}{2}H_A$.
In \eqref{eq:chamber-laplacian} one has
\[
\left|\psi(x)-U_A\left(x,x^{(i)}\right)\psi\left(x^{(i)}\right)\right|^2
=\left|\psi(x)\right|^2+\left|\psi\left(x^{(i)}\right)\right|^2-2\Re\left(\psi(x)^{*}U_A\left(x,x^{(i)}\right)\psi\left(x^{(i)}\right)\right).
\]
Every site lies in exactly one bond of each type (Proposition~\ref{prop:geometry}\eqref{it:geometry-graph}), so that $\sum_x\left|\psi\left(x^{(i)}\right)\right|^2=\|\psi\|^2$ for each $i$, on $\ell^2(\mathbb{X};\C^4)$ as well as on the spinors of $\mathbb{T}$.
Therefore
\[
\left\langle\psi,\nabla_A^{*}\nabla_A\psi\right\rangle
=\left(\sum_{i}w_i\right)\|\psi\|^2-\Re\left\langle\psi,H_A\psi\right\rangle
=\left(\sum_{i}w_i\right)\|\psi\|^2-\left\langle\psi,H_A\psi\right\rangle,
\]
the number $\langle\psi,H_A\psi\rangle$ being real.
A bounded operator on a complex Hilbert space is determined by its quadratic form.
Thus $\nabla_A^{*}\nabla_A=\sum_iw_i-H_A$, and
\[
\Re D_A=M-\frac{\rho}{2}\sum_{i}w_i+\frac{\rho}{2}\nabla_A^{*}\nabla_A=m+\frac{\rho}{2}\nabla_A^{*}\nabla_A.
\]
For a subset $Y$ of $\mathbb{X}$, $\left(P_YD_AP_Y\right)^{*}=P_YD_A^{*}P_Y$ gives $\Re\left(P_YD_AP_Y\right)=P_Y\left(\Re D_A\right)P_Y=m+\frac{\rho}{2}P_Y\nabla_A^{*}\nabla_AP_Y$ on $\ell^2(Y;\C^4)$.
Both $\nabla_A^{*}\nabla_A$ and $P_Y\nabla_A^{*}\nabla_AP_Y$ are non-negative, their quadratic forms being $\|\nabla_A\psi\|^2$ and $\|\nabla_AP_Y\psi\|^2$; for $\rho\ge0$ this yields $\Re D_A\ge m$ and $\Re\left(P_YD_AP_Y\right)\ge m$.

\emph{The norm.}
For $i=0,\dots,3$ let $D^{(i)}_A$ be the hop operator
\[
\left(D^{(i)}_A\psi\right)(x):=w\left(x,x^{(i)}\right)U_A\left(x,x^{(i)}\right)\psi\left(x^{(i)}\right),
\]
so that $D_A=M-\sum_iD^{(i)}_A$.
By Lemma~\ref{lem:hopping}, $w\left(x,x^{(i)}\right)$ is a Hermitian matrix with the eigenvalues $\frac{w_i}{2}(\rho-1)$ and $\frac{w_i}{2}(\rho+1)$; its norm is $\frac{w_i}{2}(1+|\rho|)$.
As $\left|U_A\left(x,x^{(i)}\right)\right|=1$ and $x\mapsto x^{(i)}$ is a bijection,
\[
\left\|D^{(i)}_A\psi\right\|^2=\sum_x\left|w\left(x,x^{(i)}\right)U_A\left(x,x^{(i)}\right)\psi\left(x^{(i)}\right)\right|^2\le\frac{w_i^2}{4}\cdot(1+|\rho|)^2\cdot\|\psi\|^2.
\]
Hence $\|D_A\|\le|M|+\sum_i\left\|D^{(i)}_A\right\|\le|M|+\frac{1+|\rho|}{2}\sum_iw_i$.
If $Y$ is a subset of $\mathbb{X}$, the norm of $P_YD_AP_Y$ on $\ell^2(Y;\C^4)$ is at most $\|P_Y\|^2\cdot\|D_A\|$, the norms of $P_Y$ and $D_A$ being taken on $\ell^2(\mathbb{X};\C^4)$;
an orthogonal projection has norm at most $1$, so the compression obeys the bound just proved for $\mathcal{H}=\ell^2(\mathbb{X};\C^4)$.

\emph{Blocks.}
By Definition~\ref{def:spin-matrices}, $\gamma\cdot v=\diag(\sigma\cdot v,-\sigma\cdot v)$ in the decomposition $\C^4=\C^2\oplus\C^2$, while the coefficients of $H_A$ are multiples of the unit matrix.
Every field splits into its two halves with values in $\C^2$, and on the spinors of $\mathbb{T}$ each half again satisfies $\psi(x+v)=\chi(v)\,\psi(x)$.
In this splitting
\[
H_A=\diag\left(\check H_A,\check H_A\right),\qquad J_A=\diag\left(\check J_A,-\check J_A\right).
\]
A diagonal block of a Hermitian operator is Hermitian, and a diagonal block of an anti-Hermitian operator is anti-Hermitian; so $\check H_A^{*}=\check H_A$ and $\check J_A^{*}=-\check J_A$.
Since $M$ and $\rho$ are real,
\[
D_A=\diag\left(M-\frac{\rho}{2}\check H_A+\frac12\check J_A,\;M-\frac{\rho}{2}\check H_A-\frac12\check J_A\right)=\diag\left(\check D_A,\check D_A^{*}\right).
\]
On the spinors of $\mathbb{T}$, a space of finite dimension, the determinant of this block diagonal operator is the product of the determinants of the blocks, and the determinant of an adjoint is the complex conjugate of the determinant.
Therefore
\[
\det D^{\mathbb{T}}_A=\det\check D^{\mathbb{T}}_A\cdot\overline{\det\check D^{\mathbb{T}}_A}=\left|\det\check D^{\mathbb{T}}_A\right|^2\ge0.
\]
Let now $Y$ be a subset of $\mathbb{X}$, and let $D_A$ and $\check D_A$ be the operators on $\ell^2(\mathbb{X};\C^4)$ and on $\ell^2(\mathbb{X};\C^2)$, respectively.
The projection $P_Y$ of $\ell^2(\mathbb{X};\C^4)$ acts on the site variable only: in the splitting it is $\diag\left(P_Y,P_Y\right)$, with the projection $P_Y$ of $\ell^2(\mathbb{X};\C^2)$ in both places.
Hence
\[
P_YD_AP_Y=\diag\left(P_Y\check D_AP_Y,P_Y\check D_A^{*}P_Y\right)=\diag\left(P_Y\check D_AP_Y,\left(P_Y\check D_AP_Y\right)^{*}\right)
\]
on $\ell^2(Y;\C^4)=\ell^2(Y;\C^2)\oplus\ell^2(Y;\C^2)$, the adjoint of a compression being the compression of the adjoint.
If $Y$ is finite, this space has finite dimension, and
\[
\det\left(P_YD_AP_Y\right)=\det\left(P_Y\check D_AP_Y\right)\cdot\overline{\det\left(P_Y\check D_AP_Y\right)}=\left|\det\left(P_Y\check D_AP_Y\right)\right|^2\ge0.
\]
\end{proof}

\begin{corollary}\label{cor:modulus-bound}
Let $\rho$ and $M$ be real, $w_0,\dots,w_3>0$ and $q>0$, and let $\mathbb{T}=(\Gamma,\chi)$ be a torus.
Let $G$ be scalar, that is, a complex-valued function on $\Conf(\Gamma)$, regarded as a multiple of the unit of $\Gr_{\mathbb{T}}$, for which $Z_{\mathbb{T}}[G]$ is defined.
This is the case, for instance, for a bosonic element of $\mathcal{A}(\R^3)$ that embeds in $\mathbb{T}$ (Definitions~\ref{def:integrands} and~\ref{def:functional}), and for a bounded Borel function of $F$ and of the link phases (Lemma~\ref{lem:bounded-integrands}).
Then $\det D^{\mathbb{T}}_A\ge0$ for every real $\Gamma$-periodic one-cochain $A$, and
\[
Z_{\mathbb{T}}[G]=\int_{\Conf(\Gamma)}\mathrm{d}(A,z)\,\exp\left(-\tfrac12\|F\|^2\right)\cdot\det D^{\mathbb{T}}_A\cdot G(A,z),
\]
where $\mathrm{d}(A,z)$ is the measure on $\Conf(\Gamma)$ of Definition~\ref{def:configurations}.
Consequently
\begin{equation}\label{eq:setting-modulus-bound}
\left|Z_{\mathbb{T}}[G]\right|\le Z_{\mathbb{T}}\left[\,|G|\,\right]
\end{equation}
whenever the right-hand side is finite, in particular for the two kinds of $G$ just named, and $\left|Z_{\mathbb{T}}[G]\right|\le Z_{\mathbb{T}}[G']$ for every scalar $G'$ with $|G|\le G'$.
Moreover, $\mu_{\mathbb{T}}$ is a Borel probability measure on $\Omega$.
\end{corollary}

\begin{proof}
The Berezin integral is linear, and $G$ is a multiple of the unit.
So $\int_{\mathbb{T}}\exp\left(-S_{\mathbb{T}}(A)\right)G=G\cdot\int_{\mathbb{T}}\exp\left(-S_{\mathbb{T}}(A)\right)=G\cdot\det D^{\mathbb{T}}_A$ by Definition~\ref{def:tori}, and the formula of Definition~\ref{def:functional} takes the form displayed in the statement.
A period lattice is a subgroup of $\Z^3$ of rank three, a sign character takes its values in the unit circle, and $A$ is real and $\Gamma$-periodic; Proposition~\ref{prop:accretivity}\eqref{it:determinant-sign} therefore gives $\det D^{\mathbb{T}}_A\ge0$.
The measure with the density $\exp\left(-\tfrac12\|F\|^2\right)\cdot\det D^{\mathbb{T}}_A$ with respect to $\mathrm{d}(A,z)$ is thus non-negative, and $Z_{\mathbb{T}}$ is the integral against it, on scalar functions as well as on non-negative Borel functions on $\Conf(\Gamma)$.
If $Z_{\mathbb{T}}[|G|]<\infty$, the integral of $G$ against this measure converges absolutely, and its modulus is at most $Z_{\mathbb{T}}[|G|]$; this is \eqref{eq:setting-modulus-bound}.
For a bounded Borel function of $F$ and of the link phases, and for a bosonic element of $\mathcal{A}(\R^3)$ that embeds in $\mathbb{T}$, the integral of Definition~\ref{def:functional} converges absolutely, coefficient by coefficient (Lemmas~\ref{lem:bounded-integrands} and~\ref{lem:gaussian-growth});
the Berezin integral of $\exp\left(-S_{\mathbb{T}}(A)\right)G$ is one of these coefficients, equal to $G\cdot\det D^{\mathbb{T}}_A$ for scalar $G$, and so $Z_{\mathbb{T}}[|G|]<\infty$.
Let $G'$ be scalar with $|G|\le G'$.
Then $Z_{\mathbb{T}}[|G|]\le Z_{\mathbb{T}}[G']$, the integrands being non-negative.
If $Z_{\mathbb{T}}[|G|]$ is finite, \eqref{eq:setting-modulus-bound} gives $\left|Z_{\mathbb{T}}[G]\right|\le Z_{\mathbb{T}}[G']$; if it is infinite, so is $Z_{\mathbb{T}}[G']$, and there is nothing to prove.

By Theorem~\ref{thm:partition-positivity}, $0<Z_{\mathbb{T}}[1]<\infty$.
Hence
\[
Z_{\mathbb{T}}[1]^{-1}\cdot\exp\left(-\tfrac12\|F\|^2\right)\cdot\det D^{\mathbb{T}}_A\,\mathrm{d}(A,z)
\]
is a probability measure on $\Conf(\Gamma)$, and on bounded Borel functions the normalised functional $\mathbb{E}_{\mathbb{T}}=Z_{\mathbb{T}}[\,\cdot\,]/Z_{\mathbb{T}}[1]$ is the integral against it.
The map $(A,z)\mapsto\widetilde F$ from $\Conf(\Gamma)$ to $\Omega$ is Borel: each coordinate $\widetilde F_\sigma$ is a Borel function on $\Conf(\Gamma)$, and the coordinates generate the Borel sets of $\Omega$.
The law $\mu_{\mathbb{T}}$ of $\widetilde F$ under $\mathbb{E}_{\mathbb{T}}$ is the image of this probability measure under that map, and so it is a Borel probability measure on $\Omega$.
\end{proof}

\subsection{Crossing a mirror}\label{sec:crossing}

For a mirror $\Pi=(\ell,c)$, a point $x\in\R^3$ lies above $\Pi$ if $\ell(x)>c$ and below $\Pi$ if $\ell(x)<c$;
the closed half-space $\{x\in\R^3:\ell(x)\ge c\}$ is the closed upper side of $\Pi$.
For every $x\in\R^3$ the formula for $\theta_\Pi$ in Definition~\ref{def:mirrors} gives
\begin{equation}\label{eq:chamber-displacement}
\theta_\Pi(x)-x=\frac{2}{|n_\ell|^2}\cdot\left(c-\ell(x)\right)\cdot n_\ell ,
\end{equation}
and $\ell\left(\theta_\Pi(x)\right)=2c-\ell(x)$ by Definition~\ref{def:mirrors}\eqref{it:mirror-height};
so $\theta_\Pi$ exchanges the two sides of $\Pi$.
By Proposition~\ref{prop:geometry}\eqref{it:geometry-mirrors} no site lies on $\Pi$,
and a bond whose end points lie on different sides of $\Pi$ (a bond that crosses $\Pi$) is of the form $\{x,\theta_\Pi x\}$, with $x$ below $\Pi$ and $\theta_\Pi x$ above $\Pi$.

Given a real one-cochain $A$ on $\Z^3$, the action $S_F$ of Definition~\ref{def:chamber-operator} is the formal sum of the mass terms $M\cdot\bar\psi(x)\psi(x)$, $x\in\mathbb{X}$, and of the terms
\[
-U_A(x,x')\cdot\bar\psi(x)\,w(x,x')\,\psi(x')
\]
of the oriented bonds $(x,x')$, that is, of the pairs $(x,x^{(j)})$ with $x\in\mathbb{X}$ and $j\in\{0,1,2,3\}$.
The site of a mass term is $x$, the sites of the term of $(x,x')$ are $x$ and $x'$,
and a bond $\{x,x'\}$ carries the two terms of $(x,x')$ and of $(x',x)$.
The modes $\zeta_i$ and their phases $u_i$, $i\in\idx(\Pi)$, are those of Definition~\ref{def:modes}; they do not involve $\rho$.
The coefficients $\omega_i$ are given, for every real $\rho$, by the formulas of that definition:
if $i$ is a mode of a bond of type $j$, then $\omega_i=\tfrac{w_j}{2}(1+\rho)$ when $i$ belongs to the group $1$ or $3$,
and $\omega_i=\tfrac{w_j}{2}(1-\rho)$ when $i$ belongs to the group $2$ or $4$.
For a finite set $I\subset\idx(\Pi)$ the element $\zeta_I$ and the numbers $\omega_I$, $u_I$ are formed as in \eqref{eq:setting-ordered-monomials}.

\begin{proposition}[The crossing expansion]\label{prop:crossing}
Let $\Pi=(\ell,c)$ be a mirror, let $\rho$ and $M$ be real, $w_0,\dots,w_3>0$ and $q>0$, and let $A$ be a real one-cochain on $\Z^3$.
Let $S_+$ be the sum of the terms of $S_F$ all of whose sites lie above $\Pi$,
and let $S^{\times}_\Pi$ be the sum of the terms of the bonds that cross $\Pi$.
\begin{enumerate}
\item\label{it:crossing-split}
$S_F=S_++\Theta_\Pi(S_+)+S^{\times}_\Pi$, where $\Theta_\Pi$ acts on the coefficient functions as in \eqref{eq:setting-reflection-coefficients}.
\item\label{it:crossing-modes}
One has
\[
-S^{\times}_\Pi=\sum_{i\in\idx(\Pi)}\omega_i\cdot u_i\cdot\Theta_\Pi(\zeta_i)\,\zeta_i ;
\]
here $\Theta_\Pi$ is applied to the elements $\zeta_i$, which do not depend on $A$, and the phases $u_i$ are scalar factors in front.
If $\Pi$ is not a link plane, then $u_i=1$ for every $i\in\idx(\Pi)$.
\item\label{it:crossing-signs}
One has $\omega_i\ge 0$ for every $i\in\idx(\Pi)$ if and only if $-1\le\rho\le 1$,
and $\omega_i>0$ for every $i\in\idx(\Pi)$ if and only if $-1<\rho<1$.
\item\label{it:crossing-exponential}
Let $\mathcal{J}\subset\idx(\Pi)$ be the set of the modes of finitely many bonds that cross $\Pi$,
and let $S^{\times}_\Pi(\mathcal{J})$ be the sum of the terms of these bonds.
For every $I\subset\mathcal{J}$
\[
\prod_{i\in I}\Theta_\Pi(\zeta_i)\,\zeta_i=\Theta_\Pi(\zeta_I)\,\zeta_I ,
\]
the factors of the product being even and commuting with each other, and
\[
\exp\left(-S^{\times}_\Pi(\mathcal{J})\right)=\sum_{I\subset\mathcal{J}}\omega_I\cdot u_I\cdot\Theta_\Pi(\zeta_I)\,\zeta_I .
\]
If $-1\le\rho\le 1$, this element equals $\Ins_\Pi\left[\omega^{\frac{1}{2}}\diag(u_I)\,\omega^{\frac{1}{2}}\right]$,
where $\omega$ and $\diag(u_I)$ are the diagonal operators on $\mathcal{F}(\mathcal{J})$ with the entries $\omega_I$ and $u_I$,
and $\omega^{\frac{1}{2}}$ is the diagonal operator with the entries $\omega_I^{\frac{1}{2}}$.
\item\label{it:crossing-upper-phases}
$S_+$ depends on $A$ only through the phases $\exp\left(iqA(b)\right)$ of the links $b$ contained in the closed upper side of $\Pi$.
\end{enumerate}
\end{proposition}

\begin{proof}
Write $\theta:=\theta_\Pi$ and $n:=\hat n_\ell$.

\emph{The splitting.}
As no site lies on $\Pi$, every term of $S_F$ has all its sites above $\Pi$, or all its sites below $\Pi$,
or it is one of the two terms of a bond that crosses $\Pi$; these three cases exclude each other.
By Proposition~\ref{prop:symmetries}\eqref{it:action-invariance}, $\Theta_\Pi$ maps the term of an oriented bond $(x,x')$ to the term of the oriented bond $(\theta x',\theta x)$.
Applied to the scalar matrix $M$, the rule \eqref{eq:setting-reflection-bilinear} shows that it maps the mass term at $x$ to
$M\cdot\bar\psi(\theta x)\gamma_n\gamma_n\psi(\theta x)=M\cdot\bar\psi(\theta x)\psi(\theta x)$, the mass term at $\theta x$, because $M$ is real and $\gamma_n^2=1$.
The maps $x\mapsto\theta x$ and $(x,x')\mapsto(\theta x',\theta x)$ are involutions, and $\theta$ exchanges the two sides of $\Pi$.
Consequently $\Theta_\Pi$ maps the terms of $S_+$ bijectively onto the terms of $S_F$ all of whose sites lie below $\Pi$,
whose sum is therefore $\Theta_\Pi(S_+)$; this is \eqref{it:crossing-split}.

\emph{One crossing bond.}
Let $\{x,\theta x\}$ be a bond of type $j$ that crosses $\Pi$, with $x$ below $\Pi$.
By \eqref{eq:chamber-displacement}, $\theta x-x$ is a positive multiple of $n_\ell$, so that the unit vector of the bond from $x$ to $\theta x$ is $n$.
Put $u:=U_A(x,\theta x)$.
As $A$ is real and $A\left(\pi(\theta x),\pi(x)\right)=-A\left(\pi(x),\pi(\theta x)\right)$, one has $U_A(\theta x,x)=\overline{u}$.
The two terms of the bond are $-u\cdot\bar\psi(x)\,w(x,\theta x)\,\psi(\theta x)$ and $-\overline{u}\cdot\bar\psi(\theta x)\,w(\theta x,x)\,\psi(x)$,
and the formulas of Lemma~\ref{lem:hopping} for a bond that crosses a mirror give
\begin{equation}\label{eq:chamber-one-bond}
\begin{aligned}
&u\cdot\bar\psi(x)\,w(x,\theta x)\,\psi(\theta x)+\overline{u}\cdot\bar\psi(\theta x)\,w(\theta x,x)\,\psi(x)\\
&\qquad=\tfrac{w_j}{2}\cdot u\cdot\left((1+\rho)\,\bar\psi(x)P^-_n\psi(\theta x)-(1-\rho)\,\bar\psi(x)P^+_n\psi(\theta x)\right)\\
&\qquad\quad+\tfrac{w_j}{2}\cdot\overline{u}\cdot\left((1+\rho)\,\bar\psi(\theta x)P^+_n\psi(x)-(1-\rho)\,\bar\psi(\theta x)P^-_n\psi(x)\right).
\end{aligned}
\end{equation}
By Definition~\ref{def:modes} the eight modes of the bond are, for $k=1,2$,
the elements $(e^-_{n,k})^*\psi(\theta x)$ and $(e^+_{n,k})^*\psi(\theta x)$ of the first and the second group,
with the phase $u$ and the coefficients $\tfrac{w_j}{2}(1+\rho)$ and $\tfrac{w_j}{2}(1-\rho)$ respectively,
and the elements $\bar\psi(\theta x)\,e^+_{n,k}$ and $\bar\psi(\theta x)\,e^-_{n,k}$ of the third and the fourth group,
with the phase $\overline{u}$ and the coefficients $\tfrac{w_j}{2}(1+\rho)$ and $\tfrac{w_j}{2}(1-\rho)$ respectively.

Let $e\in\C^4$ have the components $e_\alpha$, $\alpha=1,\dots,4$.
Since $\Theta_\Pi$ is antilinear, \eqref{eq:setting-reflection-fields} and $\theta^2=\mathrm{id}$ give
\[
\begin{aligned}
\Theta_\Pi\left(e^*\psi(\theta x)\right)&=\sum_{\alpha=1}^{4}e_\alpha\cdot\Theta_\Pi\left(\psi_\alpha(\theta x)\right)=-\bar\psi(x)\,\gamma_n e,\\
\Theta_\Pi\left(\bar\psi(\theta x)\,e\right)&=\sum_{\alpha=1}^{4}\overline{e_\alpha}\cdot\Theta_\Pi\left(\bar\psi_\alpha(\theta x)\right)=-e^*\gamma_n\psi(x)=-(\gamma_n e)^*\psi(x),
\end{aligned}
\]
the last step because $\gamma_n$ is Hermitian.
The vector $e^{\pm}_{n,k}$ lies in the range of $P^{\pm}_n=\frac{1}{2}(1\pm\gamma_n)$, so that $\gamma_n e^{\pm}_{n,k}=\pm e^{\pm}_{n,k}$; hence, for $k=1,2$,
\begin{equation}\label{eq:chamber-reflected-modes}
\begin{aligned}
\Theta_\Pi\left((e^-_{n,k})^*\psi(\theta x)\right)&=\bar\psi(x)\,e^-_{n,k}, &\qquad
\Theta_\Pi\left((e^+_{n,k})^*\psi(\theta x)\right)&=-\bar\psi(x)\,e^+_{n,k},\\
\Theta_\Pi\left(\bar\psi(\theta x)\,e^+_{n,k}\right)&=-(e^+_{n,k})^*\psi(x), &
\Theta_\Pi\left(\bar\psi(\theta x)\,e^-_{n,k}\right)&=(e^-_{n,k})^*\psi(x).
\end{aligned}
\end{equation}
For $e\in\C^4$ one has $\left(\bar\psi(x)\,e\right)\left(e^*\psi(\theta x)\right)=\bar\psi(x)\,ee^*\psi(\theta x)$,
and, two odd elements being anticommuting, $\left(e^*\psi(x)\right)\left(\bar\psi(\theta x)\,e\right)=-\bar\psi(\theta x)\,ee^*\psi(x)$.
With \eqref{eq:chamber-reflected-modes}, the products $\Theta_\Pi(\zeta)\,\zeta$ for the modes $\zeta$ of the first, second, third and fourth group are
\[
\begin{aligned}
\bar\psi(x)\,e^-_{n,k}(e^-_{n,k})^*\psi(\theta x),&\qquad
-\bar\psi(x)\,e^+_{n,k}(e^+_{n,k})^*\psi(\theta x),\\
\bar\psi(\theta x)\,e^+_{n,k}(e^+_{n,k})^*\psi(x),&\qquad
-\bar\psi(\theta x)\,e^-_{n,k}(e^-_{n,k})^*\psi(x).
\end{aligned}
\]
The pairs $(e^{\pm}_{n,1},e^{\pm}_{n,2})$ are orthonormal bases of the ranges of the orthogonal projections $P^{\pm}_n$ (Definition~\ref{def:spin-matrices}),
so $P^{\pm}_n=\sum_{k=1}^{2}e^{\pm}_{n,k}(e^{\pm}_{n,k})^*$, and the sums over $k$ of the four products are
\[
\bar\psi(x)P^-_n\psi(\theta x),\qquad
-\bar\psi(x)P^+_n\psi(\theta x),\qquad
\bar\psi(\theta x)P^+_n\psi(x),\qquad
-\bar\psi(\theta x)P^-_n\psi(x).
\]
Multiplied by the coefficients and the phases of their groups and added, these four sums give the right side of \eqref{eq:chamber-one-bond}.
Thus the negative of the sum of the two terms of the bond equals the sum of $\omega_i\cdot u_i\cdot\Theta_\Pi(\zeta_i)\,\zeta_i$ over the eight modes $i$ of the bond.
A bond that crosses $\Pi$ has exactly one end below $\Pi$, and its eight modes are the elements of $\idx(\Pi)$ that belong to it;
the sum over all bonds that cross $\Pi$ is the identity of \eqref{it:crossing-modes}.
If $\Pi$ is a site plane or a diagonal mirror, the bonds that cross $\Pi$ are of type $0$, $1$ or $2$ and carry no phase, by Proposition~\ref{prop:geometry}\eqref{it:geometry-projection};
then $u=1$ for every such bond, and all $u_i$ equal $1$.

\emph{Signs of the coefficients.}
The coefficients of the modes of a crossing bond of type $j$ take both values $\tfrac{w_j}{2}(1+\rho)$ and $\tfrac{w_j}{2}(1-\rho)$.
As $w_j>0$, both are non-negative if and only if $-1\le\rho\le 1$, and both are positive if and only if $-1<\rho<1$.
Moreover $\idx(\Pi)$ is not empty:
by Lemma~\ref{lem:chambers}\eqref{it:chambers-walls} one has $\Pi=g\Pi_j$ for some $g\in W$ and some $j\in\{0,1,2,3\}$,
and by Proposition~\ref{prop:geometry}\eqref{it:geometry-midpoint} the end points of the bond $\{gv_0,(gv_0)^{(j)}\}$ are images of each other under $\theta_{g\Pi_j}=\theta$;
as no site lies on $\Pi$ and $\theta$ exchanges the two sides of $\Pi$, this bond crosses $\Pi$.
Both equivalences of \eqref{it:crossing-signs} follow.

\emph{Products and the exponential.}
By \eqref{eq:setting-reflection-fields}, $\Theta_\Pi$ maps every generator to a linear combination of generators;
since it reverses the order of products, it maps odd elements to odd elements.
So $\Theta_\Pi(\zeta_i)\,\zeta_i$ is even for every $i$, and even elements commute with all elements of the Grassmann algebra.
Let $I=\{i_1<\dots<i_d\}\subset\mathcal{J}$.
As $\Theta_\Pi$ reverses products,
\[
\Theta_\Pi(\zeta_I)\,\zeta_I
=\Theta_\Pi(\zeta_{i_d})\cdots\Theta_\Pi(\zeta_{i_2})\left(\Theta_\Pi(\zeta_{i_1})\,\zeta_{i_1}\right)\zeta_{i_2}\cdots\zeta_{i_d}
=\Theta_\Pi(\zeta_{i_1})\,\zeta_{i_1}\cdot\Theta_\Pi\left(\zeta_{I\setminus\{i_1\}}\right)\zeta_{I\setminus\{i_1\}} ,
\]
where the even element in the middle has been moved to the front.
Induction on $d$, starting from $\Theta_\Pi(\zeta_\emptyset)\,\zeta_\emptyset=1$, gives the product formula of \eqref{it:crossing-exponential}.
Each $\zeta_i$ is a linear combination of generators, hence $\zeta_i^2=0$, and
\[
\left(\Theta_\Pi(\zeta_i)\,\zeta_i\right)^2=-\Theta_\Pi(\zeta_i)^2\,\zeta_i^2=0 .
\]
By the computation for one crossing bond, $-S^{\times}_\Pi(\mathcal{J})=\sum_{i\in\mathcal{J}}\omega_i\cdot u_i\cdot\Theta_\Pi(\zeta_i)\,\zeta_i$,
a finite sum of commuting even elements of the Grassmann algebra of the end points of the bonds; for fixed $A$ the phases $u_i$ are complex numbers.
Hence
\[
\begin{aligned}
\exp\left(-S^{\times}_\Pi(\mathcal{J})\right)
&=\prod_{i\in\mathcal{J}}\exp\left(\omega_i\cdot u_i\cdot\Theta_\Pi(\zeta_i)\,\zeta_i\right)\\
&=\prod_{i\in\mathcal{J}}\left(1+\omega_i\cdot u_i\cdot\Theta_\Pi(\zeta_i)\,\zeta_i\right)\\
&=\sum_{I\subset\mathcal{J}}\omega_I\cdot u_I\cdot\prod_{i\in I}\Theta_\Pi(\zeta_i)\,\zeta_i ,
\end{aligned}
\]
and the product formula turns the last sum into $\sum_{I\subset\mathcal{J}}\omega_I\cdot u_I\cdot\Theta_\Pi(\zeta_I)\,\zeta_I$.
If $-1\le\rho\le 1$, then $\omega_I\ge 0$ for all $I$ by \eqref{it:crossing-signs}, and $\omega^{\frac{1}{2}}\diag(u_I)\,\omega^{\frac{1}{2}}=\diag(\omega_I\cdot u_I)$;
by the definition of $\Ins_\Pi$ (Definition~\ref{def:modes}), $\Ins_\Pi\left[\diag(\omega_I\cdot u_I)\right]=\sum_{I\subset\mathcal{J}}\omega_I\cdot u_I\cdot\Theta_\Pi(\zeta_I)\,\zeta_I$,
which completes \eqref{it:crossing-exponential}.

\emph{The phases in $S_+$.}
A mass term does not depend on $A$.
Let $(x,x')$ be an oriented bond with both ends above $\Pi$.
If it is of type $0$, $1$ or $2$, then $U_A(x,x')=1$ by Proposition~\ref{prop:geometry}\eqref{it:geometry-projection}, and its term does not depend on $A$.
Let it be of type $3$.
By the same item, $\pi(x)$ and $\pi(x')$ are the end points of a photon link $b$, parallel to the $\mu$-th coordinate axis for some $\mu\in\{1,2,3\}$,
and $U_A(x,x')=\exp\left(iqA(\pi(x),\pi(x'))\right)$ is the phase $\exp\left(iqA(b)\right)$ or its complex conjugate.
It remains to show that $b$ lies in the closed upper side of $\Pi$.
The plane $\{\ell=c\}$ contains the open link $b$, or does not meet it, or meets it in exactly one point.
In the first case $b$ lies in $\{\ell=c\}$.
In the second case the open link lies on one side of $\Pi$;
by Proposition~\ref{prop:geometry}\eqref{it:geometry-projection} the bond lies on the same side, which is therefore the upper one, and the closed link $b$ lies in $\{\ell\ge c\}$.
The third case does not occur.
Indeed, on the open link the form $x_\mu$ takes the values strictly between the integers $\pi(x)_\mu$ and $\pi(x')_\mu$,
the forms of $\mathcal{L}$ that do not involve the coordinate $x_\mu$ are constant,
and the diagonal forms that involve it take the values of an open interval of length one with integer end points.
As $\Adm(x_\mu)=\frac{1}{2}\Z$ and $\Adm(\ell')=\Z$ for every diagonal family $\ell'$,
the only mirror that meets the open link in exactly one point is $\left(x_\mu,\tfrac{1}{2}\left(\pi(x)_\mu+\pi(x')_\mu\right)\right)$.
The coordinates of $y(x)=x-\pi(x)$ and of $y(x')$ have moduli smaller than $\frac{1}{2}$ (Definition~\ref{def:fermion-sites}), and $|\pi(x)_\mu-\pi(x')_\mu|=1$;
so $x_\mu$ and $x'_\mu$ lie on different sides of $\frac{1}{2}\left(\pi(x)_\mu+\pi(x')_\mu\right)$,
and the two ends of the bond lie on different sides of this mirror.
Since both ends lie above $\Pi$, this mirror is not $\Pi$.
\end{proof}

\subsection{Positivity at the mirrors of adapted tori}\label{sec:adapted-tori}

\begin{remark}
The proof of Theorem~\ref{thm:adapted-positivity} follows the scheme of the corresponding argument for Yang--Mills theory in four dimensions:
the variables are split,
the crossing weight is expanded into terms of positive type,
and the form $(G',G)\mapsto Z_{\mathbb{T}}[\Theta_{\ell,c}(G')\,G]$ is represented as a scalar product in an $L^2$ space.
Three features are new in comparison with that argument.
The first is the Grassmann part of the weight,
with a sign at the second mirror which the antiperiodicity of the fermion compensates.
The second is the closedness of the integer fluxes at the cubes cut by a mirror;
it is enforced by multipliers and Poisson summation, for all kinds of mirrors at once.
The third is Gauss' law at link planes.
\end{remark}

\subsubsection{Two mirrors, two halves}\label{sec:two-halves}

\begin{definition}[Split tori]\label{def:split-tori}
Let $\ell\in\mathcal{L}$ and $n:=\hat n_\ell$.
Let $\mathbb{T}=(\Gamma,\chi)$ be a torus,
$p:=p(\mathbb{T},\ell)$ its height period and $\Gamma_\parallel:=\Gamma\cap\{\ell=0\}$ (Definition~\ref{def:slabs}).
The torus $\mathbb{T}$ is \emph{split for} $\ell$ if
\[
g_0:=\frac{p}{|n_\ell|^2}\,n_\ell\in\Z^3
\qquad\text{and}\qquad
\Gamma=\Gamma_\parallel+\Z g_0 .
\]
A torus is adapted to $\ell$ if and only if it is split for $\ell$ and $\chi(g_0)=-1$.
Indeed, by Definition~\ref{def:slabs} the straightened torus $\mathbb{T}^0(\ell)$, when it is defined,
has the lattice $\Gamma_\parallel+\Z g_0$ and the character that agrees with $\chi$ on $\Gamma_\parallel$ and takes the value $-1$ at $g_0$;
and if $\mathbb{T}$ is split for $\ell$, this lattice is the period lattice $\Gamma$, so that $\mathbb{T}^0(\ell)$ is defined.
For instance, the coordinate torus $C(n_1,n_2,n_3)$ is adapted to $x_1$, $x_2$ and $x_3$,
and $D(a_1,a_2,a_3)$ is adapted to $\ell_-$, $\ell_+$ and $x_3$,
because a rectangular torus is adapted to each of the three families of its frame (Definition~\ref{def:cell-structures}).
For $g\in\Ocub$ and a rectangular torus $\mathbb{T}(\mathfrak{F};p_1,p_2,p_3)$ with the frame $\mathfrak{F}=(\ell^1,\ell^2,\ell^3)$,
the torus $g\,\mathbb{T}(\mathfrak{F};p_1,p_2,p_3)$ of Definition~\ref{def:tori} is the rectangular torus with the same height periods
for the frame formed by the families $\pm\ell^a\circ g^{-1}$, the signs being chosen so that these lie in $\mathcal{L}$.
Consequently, whatever the family $\ell$, some image of $C(n_1,n_2,n_3)$ or of $D(a_1,a_2,a_3)$ under $\Ocub$ is adapted to $\ell$.
For an axis family $\ell$, the torus $C(n_1,n_2,n_3)$ itself is such an image.
For a diagonal family $\ell$, there is $g\in\Ocub$ with $g(e_1-e_2)=n_\ell$,
because the signed permutation matrices that form $\Ocub$ (Definition~\ref{def:mirrors}) permute the twelve vectors $\pm e_\mu\pm e_\nu$, $\mu\neq\nu$, transitively;
as $g^{-1}$ is the transpose of $g$, one has $\ell_-\circ g^{-1}=\ell$,
so $\ell$ is a family of the frame of $g\,D(a_1,a_2,a_3)$, and this torus is adapted to $\ell$.

Let $\mathbb{T}$ be split for $\ell$, and let $c\in\Adm(\ell)$.
Put
\[
c':=c+\frac{p}{2},\qquad
\Pi:=(\ell,c),\qquad
\Pi':=(\ell,c'),\qquad
\theta:=\theta_\Pi,\qquad
\theta':=\theta_{\Pi'},
\]
\[
U_+:=\{c<\ell<c'\},\qquad
U_-:=\left\{c-\frac{p}{2}<\ell<c\right\},
\]
and let $\mathcal{M}$ be the union of the planes $\left\{\ell=c+j\,\frac{p}{2}\right\}$, $j\in\Z$.
Then the following hold.
\begin{enumerate}
\item\label{it:split-height}
If $\ell$ is an axis family, then $p\ge4$;
if $\ell$ is a diagonal family, then $p$ is even and $p\ge8$.
In both cases $p\ge4\,|n_\ell|^2$ and $c'\in\Adm(\ell)$, so that $\Pi'$ is a mirror.
Indeed, $g_0$ is a non-zero element of $\Gamma$, whence $|g_0|_\infty\ge\sys(\Gamma)\ge4$ (Definition~\ref{def:configurations}).
For $\ell=x_\mu$ one has $g_0=p\,e_\mu$.
For $\ell=x_\mu\mp x_\nu$ one has $g_0=\frac{p}{2}\,(e_\mu\mp e_\nu)$,
which is an integer vector only if $p$ is even, and $|g_0|_\infty=\frac{p}{2}$.
So $\frac{p}{2}$ lies in $\frac12\Z$, respectively in $\Z$, and $c'=c+\frac{p}{2}\in\Adm(\ell)$.
\item\label{it:split-reflections}
One has $\theta'=\tau_{g_0}\theta$ and $\theta\mathbb{T}=\mathbb{T}$,
and $\theta$ and $\theta'$ induce one and the same involution of $\R^3/\Gamma$,
of the set of the cells of $\Z^3/\Gamma$, of $\mathbb{X}_{\mathbb{T}}$ and, through $\theta^{*}$, of $\Conf(\Gamma)$.
Indeed, $\theta'\theta=\tau_{g_0}$ by Definition~\ref{def:mirrors}\eqref{it:mirror-product},
since $\frac{2}{|n_\ell|^2}\,(c'-c)\,n_\ell=g_0$,
and $\theta^2=\mathrm{id}$ by Definition~\ref{def:mirrors}\eqref{it:mirror-involution};
so $\theta'=\theta'\theta\theta=\tau_{g_0}\theta$.
The linear part $\vec\theta$ of $\theta$ is the reflection in the plane $\{\ell=0\}$;
it fixes $\Gamma_\parallel$ pointwise and maps $g_0$ to $-g_0$.
Hence $\vec\theta\,\Gamma=\Gamma$,
and as $\chi(-g_0)=\chi(g_0)$, the torus $\theta\mathbb{T}=(\vec\theta\,\Gamma,\chi\circ\vec\theta^{-1})$ of Definition~\ref{def:tori} equals $\mathbb{T}$.
Because $g_0\in\Gamma$, the maps induced on $\R^3/\Gamma$ by $\theta$ and by $\theta'=\tau_{g_0}\theta$ coincide.
The same holds for the cells, which $\theta$ permutes, since it maps $\Z^3$ onto $\Z^3$ (Lemma~\ref{lem:chambers}\eqref{it:chambers-group}),
and for the sites, which $\theta$ permutes as an element of $W$.
On $\Conf(\Gamma)$ one has $\theta'^{*}=\theta^{*}\tau_{g_0}^{*}$ (Definition~\ref{def:cochains}),
and $\tau_{g_0}^{*}$ acts trivially on the $\Gamma$-periodic cochains (Definition~\ref{def:configurations}).
These maps are involutions because $\theta^2=\mathrm{id}$.
\item\label{it:split-slabs}
One has $U_-=\theta U_+$, since $\ell\circ\theta=2c-\ell$ (Definition~\ref{def:mirrors}\eqref{it:mirror-height}).
The planes that form $\mathcal{M}$ are the translates $\Pi+jg_0$ and $\Pi'+jg_0$, $j\in\Z$, because $\ell(g_0)=p$.
Their positions $c+j\,\frac{p}{2}$ lie in $\Adm(\ell)$, since $\Adm(\ell)$ is a group that contains $c$ and, by \eqref{it:split-height}, $\frac{p}{2}$;
so these planes are mirrors.
The slabs
\[
U_++jg_0=\left\{c+jp<\ell<c+jp+\frac{p}{2}\right\},\qquad
U_-+jg_0=\left\{c+jp-\frac{p}{2}<\ell<c+jp\right\},\qquad j\in\Z,
\]
consist of the points whose heights lie strictly between two consecutive elements of $c+\frac{p}{2}\Z$.
Hence $\R^3\setminus\mathcal{M}$ is the disjoint union of these slabs.
\end{enumerate}
\end{definition}

\begin{lemma}[Products of the two halves]\label{lem:reflected-embedding}
Let $\ell\in\mathcal{L}$, let $\mathbb{T}=(\Gamma,\chi)$ be a torus that is split for $\ell$, and let $c\in\Adm(\ell)$;
let $c'$, $\Pi$, $\theta$, $U_+$ be as in Definition~\ref{def:split-tori},
and $\Xi:=\Xi(\mathbb{T},\ell)=(\Gamma_\parallel,\chi|_{\Gamma_\parallel})$.
If $G$ and $G'$ are products of $\mathcal{A}(U_+)$ that embed in $\Gamma_\parallel$,
then $\Theta_\Pi(G')\,G$ is a product of $\mathcal{A}(\R^3)$ that embeds in $\Gamma$ (Definition~\ref{def:integrands}).
Consequently, for all $G,G'\in\mathcal{A}_\Xi[c,c']$ (Definition~\ref{def:slabs}),
the element $\Theta_\Pi(G')\,G$ belongs to $\mathcal{A}(\R^3)$ and embeds in $\Gamma$.
\end{lemma}

\begin{proof}
Write $G=\Phi\,G_1\cdots G_r$ and $G'=\Phi'\,G'_1\cdots G'_{r'}$
with $G_j\in\mathcal{G}^{\mathrm{ev}}(B_j)$, $G'_k\in\mathcal{G}^{\mathrm{ev}}(B'_k)$,
pairwise disjoint boxes $B_1,\dots,B_r$ and pairwise disjoint boxes $B'_1,\dots,B'_{r'}$ contained in $U_+$,
and bounded Borel functions $\Phi$, $\Phi'$ of finitely many $F_\sigma$, $\sigma\subset U_+$.
The image $\theta B'_k$ of a box is a box,
because the linear part of $\theta$ is a signed permutation matrix;
it is contained in $U_-$ by Definition~\ref{def:split-tori}\eqref{it:split-slabs},
and $\Theta_\Pi$ maps $\mathcal{G}^{\mathrm{ev}}(B'_k)$ onto $\mathcal{G}^{\mathrm{ev}}(\theta B'_k)$ (Lemma~\ref{lem:reflection-classes}).
Moreover $\Theta_\Pi\Phi'$ is the bounded Borel function $F\mapsto\overline{\Phi'(\theta^{*}F)}$ of finitely many $F_\sigma$, $\sigma\subset U_-$.
All factors $G_j$, $G'_k$ are even, so their order is immaterial, and
\[
\Theta_\Pi(G')\,G=\left(\Theta_\Pi\Phi'\right)\cdot\Phi\cdot\Theta_\Pi(G'_1)\cdots\Theta_\Pi(G'_{r'})\cdot G_1\cdots G_r
\]
is a product of $\mathcal{A}(\R^3)$ with the boxes $B_1,\dots,B_r,\theta B'_1,\dots,\theta B'_{r'}$;
these are pairwise disjoint, the first $r$ of them lying in $U_+$ and the others in $U_-$.
It remains to show that this family of boxes embeds disjointly in $\Gamma$.

\emph{Each box embeds in $\Gamma$.}
The entries of $n_\ell$ lie in $\{0,1,-1\}$, so $|\ell|\le|n_\ell|^2$ on the cube $[-1,1]^3$.
Every element of $\Gamma$ is of the form $g+jg_0$ with $g\in\Gamma_\parallel$ and $j\in\Z$,
and the translation by $g+jg_0$ adds $jp$ to all heights.
Let $B$ be one of the boxes $B_1,\dots,B_r$.
It is a compact subset of $U_+$,
so the heights of the points of $B+[-1,1]^3$ lie in the open interval $(c-|n_\ell|^2,c'+|n_\ell|^2)$,
whose length is $\frac{p}{2}+2\,|n_\ell|^2\le p$ by Definition~\ref{def:split-tori}\eqref{it:split-height}.
For $j\neq0$ the shifted interval does not meet this one, its length being at most $p$,
so the translate of $B+[-1,1]^3$ by $g+jg_0$ is disjoint from $B+[-1,1]^3$.
For $j=0$ and $g\neq0$ the two sets are disjoint because $B$ embeds in $\Gamma_\parallel$.
Hence the translates of $B+[-1,1]^3$ by the elements of $\Gamma$ are pairwise disjoint, that is, $B$ embeds in $\Gamma$.
Now let $B=\theta B'_k$.
As $\ell\circ\theta=2c-\ell$ and $B'_k$ is a compact subset of $U_+$,
the heights of the points of $B+[-1,1]^3$ lie in the open interval $\left(c-\frac{p}{2}-|n_\ell|^2,c+|n_\ell|^2\right)$,
whose length is again at most $p$;
so the translate of $B+[-1,1]^3$ by $g+jg_0$ is disjoint from $B+[-1,1]^3$ if $j\neq0$.
For $g\in\Gamma_\parallel$ one has $\theta B'_k+[-1,1]^3+g=\theta\left(B'_k+[-1,1]^3+g\right)$,
because the linear part of $\theta$ fixes $g$ and maps $[-1,1]^3$ onto itself;
as $B'_k$ embeds in $\Gamma_\parallel$ and $\theta$ is injective, these sets are pairwise disjoint for $g\in\Gamma_\parallel$.
Hence $\theta B'_k$ embeds in $\Gamma$.

\emph{The translates are pairwise disjoint.}
Let $g\in\Gamma_\parallel$ and $i\in\Z$.
As $U_\pm+g=U_\pm$, the set $B_j+g+ig_0$ lies in $U_++ig_0$, and $\theta B'_k+g+ig_0$ lies in $U_-+ig_0$.
Two such sets in different slabs are disjoint by Definition~\ref{def:split-tori}\eqref{it:split-slabs}.
Inside $U_++ig_0$ the sets $B_j+g+ig_0$ are the translates by $ig_0$ of the sets $B_j+g$,
which are pairwise disjoint because the boxes of $G$ embed disjointly in $\Gamma_\parallel$.
Inside $U_-+ig_0$ the sets $\theta B'_k+g+ig_0$ are the translates by $ig_0$ of the sets $\theta\left(B'_k+g\right)$,
which are pairwise disjoint because the boxes of $G'$ embed disjointly in $\Gamma_\parallel$.

Thus the boxes of $\Theta_\Pi(G')\,G$ embed disjointly in $\Gamma$, which is the first assertion.
The second follows from it, applied to each pair of summands, because $\Theta_\Pi$ is antilinear
and an element of $\mathcal{A}(\R^3)$ that is a finite sum of products embedding in $\Gamma$ embeds in $\Gamma$.
\end{proof}

\paragraph{The two halves of the Grassmann algebra.}
Let $\mathbb{T}=(\Gamma,\chi)$ be split for $\ell$, and let $c\in\Adm(\ell)$, with the notation of Definition~\ref{def:split-tori}.
Every plane of $\mathcal{M}$ is a mirror by Definition~\ref{def:split-tori}\eqref{it:split-slabs},
so no site lies on $\mathcal{M}$ (Proposition~\ref{prop:geometry}\eqref{it:geometry-mirrors}).
For a site $x$, $g\in\Gamma_\parallel$ and $j\in\Z$ one has $\ell(x+g+jg_0)=\ell(x)+jp$.
Since $\ell(x)\notin c+\frac{p}{2}\Z$, there is exactly one $j\in\Z$ with $c-\frac{p}{2}<\ell(x)+jp<c+\frac{p}{2}$, and $\ell(x)+jp\neq c$;
hence every $\Gamma$-orbit of sites meets $(\mathbb{X}\cap U_+)\cup(\mathbb{X}\cap U_-)$ in exactly one $\Gamma_\parallel$-orbit.
Choose a set $\mathbb{X}_+$ of representatives of $(\mathbb{X}\cap U_+)/\Gamma_\parallel$, and put $\mathbb{X}_-:=\theta\,\mathbb{X}_+$.
Since $\theta$ maps $\mathbb{X}\cap U_+$ onto $\mathbb{X}\cap U_-$ and $\theta(x+g)=\theta x+g$ for $g\in\Gamma_\parallel$,
the set $\mathbb{X}_-$ represents $(\mathbb{X}\cap U_-)/\Gamma_\parallel$.
So $\mathbb{X}_+\cup\mathbb{X}_-$ is a set of representatives of $\mathbb{X}_{\mathbb{T}}$,
and the generators at its points are free generators of $\Gr_{\mathbb{T}}$ (Definition~\ref{def:tori}).
Put
\[
\Gr_{\mathbb{T}}^{\pm}:=\Gr(\mathbb{X}_\pm)\subset\Gr_{\mathbb{T}} .
\]
In $\Gr_{\mathbb{T}}$ the generators at a site of $U_\pm+\Z g_0$ are, up to a sign, generators at a site of $\mathbb{X}_\pm$;
so $\Gr_{\mathbb{T}}^{\pm}$ is the subalgebra generated by the generators at all sites of $U_\pm+\Z g_0$,
and it does not depend on the choice of $\mathbb{X}_+$.
The number $8\,\#|\mathbb{X}_+|$ of the generators of $\Gr_{\mathbb{T}}^{+}$ is even.
For a finite set $Y$ of sites let
\[
\Psi_Y:=\prod_{x\in Y}\prod_{\alpha=1}^{4}\psi_\alpha(x)\,\bar\psi_\alpha(x),
\]
a product of even elements, whose order is therefore immaterial;
by Definition~\ref{def:grassmann}, $\int_Y\xi$ is the coefficient of $\Psi_Y$ in $\xi\in\Gr(Y)$.
Put $\Psi_\pm:=\Psi_{\mathbb{X}_\pm}$, and let $\Theta:=\Theta_\Pi$.
Since $\theta\mathbb{T}=\mathbb{T}$ by Definition~\ref{def:split-tori}\eqref{it:split-reflections},
$\Theta$ is an antilinear anti-automorphism of $\Gr_{\mathbb{T}}$ (Lemma~\ref{lem:reflected-algebras}),
and $\Theta^2=1$ (Lemma~\ref{lem:sign-rules}).
By \eqref{eq:setting-reflection-fields}, $\Theta$ maps the generators at a site $x$ to linear combinations of the generators at $\theta x$;
so $\Theta$ preserves the degree and maps $\Gr_{\mathbb{T}}^{+}$ onto $\Gr_{\mathbb{T}}^{-}$ and $\Gr_{\mathbb{T}}^{-}$ onto $\Gr_{\mathbb{T}}^{+}$.
In particular $\Theta(\Psi_+)$ is a multiple of $\Psi_-$, and in fact
\[
\Theta(\Psi_+)=\Psi_- ,
\]
because $\int_{\mathbb{X}_-}\Theta(\Psi_+)=\overline{\int_{\mathbb{X}_+}\Psi_+}=1$ by Proposition~\ref{prop:symmetries}\eqref{it:berezin-conjugate}.

\paragraph{The action at the two mirrors.}
Let in addition $\rho,M\in\R$, $w_0,\dots,w_3>0$ and $q>0$ (Definition~\ref{def:chamber-operator}),
and let $A$ be a real $\Gamma$-periodic one-cochain.
Put
\[
S_+:=\sum_{x\in\mathbb{X}_+}\left(M\cdot\bar\psi(x)\,\psi(x)-\sum_{j:\ x^{(j)}\in U_+}U_A(x,x^{(j)})\cdot\bar\psi(x)\,w(x,x^{(j)})\,\psi(x^{(j)})\right).
\]
Since $\mathbb{X}_+$ represents $(\mathbb{X}\cap U_+)/\Gamma_\parallel$,
each $\Gamma_\parallel$-orbit of sites in $U_+$, and each $\Gamma_\parallel$-orbit of oriented bonds contained in $U_+$, contributes exactly one term to $S_+$.
Let
\[
S^\times_\Pi:=-\sum_{x}\left(U_A(x,\theta x)\cdot\bar\psi(x)\,w(x,\theta x)\,\psi(\theta x)+U_A(\theta x,x)\cdot\bar\psi(\theta x)\,w(\theta x,x)\,\psi(x)\right),
\]
the sum over those $x\in\mathbb{X}_+$ for which $\theta x$ is a neighbour of $x$,
and let $S^\times_{\Pi'}$ be the same expression with $\theta'$ in the place of $\theta$,
summed over those $x\in\mathbb{X}_+$ for which $\theta'x$ is a neighbour of $x$.

By Proposition~\ref{prop:geometry}\eqref{it:geometry-midpoint} the open segment of a bond meets exactly one mirror.
Hence a bond that crosses $\Pi$ does not meet $\Pi'$, and its upper end point lies in $U_+$;
modulo $\Gamma_\parallel$ it is a bond $\{\theta x,x\}$ with $x\in\mathbb{X}_+$, for exactly one $x$.
Conversely, if $\theta x$ is a neighbour of $x\in\mathbb{X}_+$, the bond $\{\theta x,x\}$ crosses $\Pi$;
its lower end point is $\theta x$, its upper end point is $x$,
and its eight modes at $\Pi$ (Definition~\ref{def:modes}) are linear combinations of the generators at $x$.
Let $\idx_{\mathbb{T}}(\Pi)$ be the set of the modes at $\Pi$ of all these bonds, a finite set endowed with a total order;
it represents $\idx(\Pi)/\Gamma_\parallel$.
For $i\in\idx_{\mathbb{T}}(\Pi)$ let $\zeta_i$, $\omega_i$, $u_i$ be the mode, its coefficient and its phase,
and for finite $I\subset\idx_{\mathbb{T}}(\Pi)$ let $\zeta_I$, $\omega_I$, $u_I$ be as in \eqref{eq:setting-ordered-monomials}.
Thus $\zeta_i\in\Gr_{\mathbb{T}}^{+}$ and $\Theta(\zeta_i)\in\Gr_{\mathbb{T}}^{-}$.
(Another choice of $\mathbb{X}_+$ multiplies a mode by $\chi(g)=\pm1$ for some $g\in\Gamma_\parallel$,
and changes none of $\Theta(\zeta_i)\,\zeta_i$, $\omega_i$ and $u_i$.)

A bond that crosses $\Pi'$ does not meet $\Pi$, so its lower end point lies in $U_+$;
modulo $\Gamma_\parallel$ it is a bond $\{x,\theta'x\}$ with $x\in\mathbb{X}_+$, for exactly one $x$.
If $\theta'x$ is a neighbour of $x\in\mathbb{X}_+$, the bond $\{x,\theta'x\}$ crosses $\Pi'$;
its lower end point is $x$, and its upper end point is $\theta'x=\theta x+g_0$.
Let $\idx_{\mathbb{T}}(\Pi')$ be the set of the modes at $\Pi'$ of all these bonds (Definition~\ref{def:modes}, for the mirror $\Pi'$), a finite set endowed with a total order,
and let $\zeta_{i'}$, $\omega_{i'}$, $u_{i'}$ and $\omega_{I'}$, $u_{I'}$ be defined as before, for $i'\in\idx_{\mathbb{T}}(\Pi')$ and finite $I'\subset\idx_{\mathbb{T}}(\Pi')$.
The mode $\zeta_{i'}$ is linear in the generators at $\theta'x$.
Put
\[
\check\zeta_{i'}:=\Theta_{\Pi'}(\zeta_{i'})\in\Gr_{\mathbb{T}}^{+},
\]
an element that is linear in the generators at $x$:
by \eqref{eq:chamber-reflected-modes}, applied at the mirror $\Pi'$,
\[
\check\zeta_{i'}=\bar\psi(x)\,e^-_{n,k},\qquad
-\bar\psi(x)\,e^+_{n,k},\qquad
-(e^+_{n,k})^{*}\,\psi(x),\qquad
(e^-_{n,k})^{*}\,\psi(x)
\]
for the modes of the groups $1$, $2$, $3$, $4$ with the index $k\in\{1,2\}$.
For a finite set $I'=\{i'_1<\dots<i'_d\}\subset\idx_{\mathbb{T}}(\Pi')$ put $\check\zeta_{I'}:=\check\zeta_{i'_1}\cdots\check\zeta_{i'_d}$.

\begin{lemma}[The sign at the second mirror]\label{lem:second-mirror}
Let $\rho,M\in\R$, $w_0,\dots,w_3>0$ and $q>0$.
Let $\ell\in\mathcal{L}$, let $\mathbb{T}=(\Gamma,\chi)$ be a torus that is split for $\ell$, let $c\in\Adm(\ell)$,
and let $A$ be a real $\Gamma$-periodic one-cochain;
the notation is that of Definition~\ref{def:split-tori} and of the paragraphs on the two halves of the Grassmann algebra and on the action at the two mirrors, with $\Theta=\Theta_\Pi$.
Then
\begin{equation}\label{eq:chamber-second-mirror}
\Theta_{\Pi'}=\chi(g_0)^{\deg}\,\Theta_\Pi\qquad\text{on }\Gr_{\mathbb{T}},
\end{equation}
where $\chi(g_0)^{\deg}$ is the automorphism of $\Gr_{\mathbb{T}}$ that multiplies every monomial of degree $d$ by $\chi(g_0)^d$.
The element $S_+$ of $\Gr_{\mathbb{T}}^{+}$ is even;
its coefficients are polynomials in the phases $\exp\left(\pm iqA(b)\right)$ of links $b$ of $\Z^3$ that are contained in the closed slab $\{c\le\ell\le c'\}$,
and each such link either lies in $\Pi$ or in $\Pi'$,
or the link without its two end points lies in $U_+$.
Moreover
\[
S_{\mathbb{T}}(A)=S_++\Theta(S_+)+S^\times_\Pi+S^\times_{\Pi'},
\]
where $\Theta$ acts on the coefficients as in \eqref{eq:setting-reflection-coefficients}, and
\begin{equation}\label{eq:chamber-action-split}
\begin{split}
-S_{\mathbb{T}}(A)={}&-S_+-\Theta(S_+)
+\sum_{i\in\idx_{\mathbb{T}}(\Pi)}\omega_i\cdot u_i\cdot\Theta(\zeta_i)\,\zeta_i\\
&-\chi(g_0)\sum_{i'\in\idx_{\mathbb{T}}(\Pi')}\omega_{i'}\cdot u_{i'}\cdot\Theta(\check\zeta_{i'})\,\check\zeta_{i'},
\end{split}
\end{equation}
\begin{equation}\label{eq:chamber-weight-expansion}
\begin{split}
\exp\left(-S_{\mathbb{T}}(A)\right)={}&\exp(-S_+)\cdot\Theta\left(\exp(-S_+)\right)\\
&\cdot\sum_{I\subset\idx_{\mathbb{T}}(\Pi)}\ \sum_{I'\subset\idx_{\mathbb{T}}(\Pi')}
\left(-\chi(g_0)\right)^{\#|I'|}\cdot\omega_I\cdot\omega_{I'}\cdot u_I\cdot u_{I'}\cdot\Theta(\zeta_I)\,\zeta_I\cdot\Theta(\check\zeta_{I'})\,\check\zeta_{I'} .
\end{split}
\end{equation}
\end{lemma}

\begin{proof}
\emph{The relation between the two reflections.}
The reflection $\theta'=\tau_{g_0}\theta$ has the same linear part as $\theta$, so $\theta'\mathbb{T}=\theta\mathbb{T}=\mathbb{T}$ by Definition~\ref{def:split-tori}\eqref{it:split-reflections};
by Lemma~\ref{lem:reflected-algebras}, $\Theta_{\Pi'}$ and $\Theta_\Pi$ are antilinear anti-automorphisms of $\Gr_{\mathbb{T}}$,
and so is $\chi(g_0)^{\deg}\,\Theta_\Pi$, because $\chi(g_0)^{\deg}$ is an automorphism.
The mirrors $\Pi$ and $\Pi'$ belong to the family $\ell$ and have the same upward unit normal $n$.
For every site $x$ one has $\theta'x=\theta x+g_0$ with $g_0\in\Gamma$,
so \eqref{eq:setting-reflection-fields} and the relations of $\Gr_{\mathbb{T}}$ (Definition~\ref{def:tori}) give
\[
\Theta_{\Pi'}\psi(x)=-\bar\psi(\theta x+g_0)\,\gamma_n=\chi(g_0)\cdot\Theta_\Pi\psi(x),\qquad
\Theta_{\Pi'}\bar\psi(x)=-\gamma_n\,\psi(\theta x+g_0)=\chi(g_0)\cdot\Theta_\Pi\bar\psi(x).
\]
Thus the two sides of \eqref{eq:chamber-second-mirror} agree on the generators, and therefore on $\Gr_{\mathbb{T}}$.

\emph{The splitting of the action.}
Since $\mathbb{X}_+\cup\mathbb{X}_-$ represents $\mathbb{X}_{\mathbb{T}}$,
every $\Gamma$-orbit of oriented bonds contains exactly one pair $(x,x^{(j)})$ with $x\in\mathbb{X}_+\cup\mathbb{X}_-$;
so $S_{\mathbb{T}}(A)$ is the sum of the mass terms $M\cdot\bar\psi(x)\,\psi(x)$, $x\in\mathbb{X}_+\cup\mathbb{X}_-$,
and of the terms $-U_A(x,x^{(j)})\cdot\bar\psi(x)\,w(x,x^{(j)})\,\psi(x^{(j)})$ of these pairs.
Let $x\in\mathbb{X}_+$ and $j\in\{0,\dots,3\}$ with $x^{(j)}\notin U_+$.
The site $x^{(j)}$ lies on no plane of $\mathcal{M}$,
so either $\ell(x^{(j)})<c$ and $x$, $x^{(j)}$ lie on different sides of $\Pi$,
or $\ell(x^{(j)})>c'$ and they lie on different sides of $\Pi'$.
By Proposition~\ref{prop:geometry}\eqref{it:geometry-mirrors}, $x^{(j)}=\theta x$ in the first case and $x^{(j)}=\theta'x=\theta x+g_0$ in the second;
moreover $\theta x\in\theta\,\mathbb{X}_+=\mathbb{X}_-$.
The neighbours are equivariant under $W$ (Definition~\ref{def:fermion-sites}): $(\theta x)^{(j)}=\theta(x^{(j)})$.
As $\theta$ maps $\mathbb{X}_+$ onto $\mathbb{X}_-$, the pairs with first point in $\mathbb{X}_-$ are the pairs $(\theta x,\theta(x^{(j)}))$, $x\in\mathbb{X}_+$.
According to the three cases, $\theta(x^{(j)})$ lies in $U_-$, or equals $x$,
or equals $\theta(\theta x+g_0)=x-g_0$;
in the last case $(\theta x,x-g_0)$ lies in the $\Gamma$-orbit of $(\theta'x,x)$.
Collecting the terms:
the mass terms at $\mathbb{X}_+$ and the pairs with first point in $\mathbb{X}_+$ and second point in $U_+$ give $S_+$;
the mass terms at $\mathbb{X}_-$ and the pairs with first point in $\mathbb{X}_-$ and second point in $U_-$ give an element $S_-$;
the pairs $(x,\theta x)$ and $(\theta x,x)$ give $S^\times_\Pi$;
the pairs $(x,\theta'x)$ and, in the form $(\theta x,x-g_0)$, the pairs $(\theta'x,x)$ give $S^\times_{\Pi'}$.
Hence $S_{\mathbb{T}}(A)=S_++S_-+S^\times_\Pi+S^\times_{\Pi'}$.
Since $\Theta$ is the map of $\Gr_{\mathbb{T}}$ induced by $\Theta_\Pi$ on $\Gr(\mathbb{X})$ (Lemma~\ref{lem:reflected-algebras}),
Proposition~\ref{prop:symmetries}\eqref{it:action-invariance} shows that $\Theta$ maps the term of a pair $(x,x')$ to the term of $(\theta x',\theta x)$,
the coefficients being transformed as in \eqref{eq:setting-reflection-coefficients}.
The mass terms are covered by the rule \eqref{eq:setting-reflection-bilinear}, applied to the scalar matrix $M$:
as $M$ is real and $\gamma_n^2=1$, the image of $M\cdot\bar\psi(x)\,\psi(x)$ is $M\cdot\bar\psi(\theta x)\,\gamma_n\gamma_n\,\psi(\theta x)=M\cdot\bar\psi(\theta x)\,\psi(\theta x)$, the mass term at $\theta x$.
As $(x,x')\mapsto(\theta x',\theta x)$ maps the $\Gamma_\parallel$-orbits of pairs with both points in $U_+$ bijectively onto those of pairs with both points in $U_-$,
and translations by $\Gamma_\parallel$ do not change the terms in $\Gr_{\mathbb{T}}$,
$\Theta$ maps the terms of $S_+$ bijectively onto those of $S_-$, and $S_-=\Theta(S_+)$.
Each term of $S_+$ is a linear combination of products of two generators of $\Gr_{\mathbb{T}}^{+}$,
a generator at a site of $U_+$ being, up to a sign, a generator at a site of $\mathbb{X}_+$;
so $S_+$ is an even element of $\Gr_{\mathbb{T}}^{+}$.

\emph{The link phases of $S_+$.}
The coefficients of $S_+$ are the number $M$ and the entries of the matrices $w(x,x')$, multiplied by the phases $U_A(x,x')=\exp\left(iqA(\pi(x),\pi(x'))\right)$.
Let $(x,x')$ be a pair of neighbouring sites in $U_+$ with $\pi(x)\neq\pi(x')$,
and let $b$ be the link with the end points $\pi(x)$ and $\pi(x')$, so that $U_A(x,x')=\exp\left(\pm iqA(b)\right)$.
Both $x$ and $x'$ lie above $\Pi$;
by Proposition~\ref{prop:crossing}\eqref{it:crossing-upper-phases}, applied at the mirror $\Pi$, the link $b$ is contained in $\{\ell\ge c\}$.
Both $x$ and $x'$ lie below $\Pi'$, so $\theta'x$ and $\theta'x'$ are neighbouring sites above $\Pi'$.
Since $\pi\circ\theta'=\theta'\circ\pi$ (Proposition~\ref{prop:geometry}\eqref{it:geometry-projection}),
the phase of the pair $(\theta'x,\theta'x')$ is the phase of the link $\theta'b$,
and the same statement, applied at the mirror $\Pi'$, gives $\theta'b\subset\{\ell\ge c'\}$.
As $\ell\circ\theta'=2c'-\ell$, it follows that $b\subset\{\ell\le c'\}$.
So $b$ is contained in the closed slab $\{c\le\ell\le c'\}$.
Finally, $\ell$ is affine on $b$.
If $\ell$ takes the value $c$ or $c'$ at a point of $b$ other than its end points,
then $\ell$ attains its minimum or its maximum over $b$ at that point, hence is constant on $b$, and $b$ lies in $\Pi$ or in $\Pi'$;
otherwise the link $b$ without its end points lies in $U_+$.

\emph{The crossing terms.}
The expansion of the crossing part in Proposition~\ref{prop:crossing} holds for every real $\rho$ and bond by bond;
for a bond $\{\theta x,x\}$, $x\in\mathbb{X}_+$, that crosses $\Pi$ it reads
\[
-\left(\text{terms of }\{\theta x,x\}\right)=\sum_{i}\omega_i\cdot u_i\cdot\Theta(\zeta_i)\,\zeta_i ,
\]
the sum over the eight modes of the bond;
this is an identity between even elements of the Grassmann algebra of the two points $\theta x$, $x$,
which represent different points of $\mathbb{X}_{\mathbb{T}}$, so it holds in $\Gr_{\mathbb{T}}$ as well (Definition~\ref{def:tori}).
Summation over $x$ gives
\[
-S^\times_\Pi=\sum_{i\in\idx_{\mathbb{T}}(\Pi)}\omega_i\cdot u_i\cdot\Theta(\zeta_i)\,\zeta_i .
\]
Proposition~\ref{prop:crossing}, applied at the mirror $\Pi'$ to the bonds $\{x,\theta'x\}$, $x\in\mathbb{X}_+$,
whose end points $x$ and $\theta'x$ also represent different points of $\mathbb{X}_{\mathbb{T}}$, gives
\[
-S^\times_{\Pi'}=\sum_{i'\in\idx_{\mathbb{T}}(\Pi')}\omega_{i'}\cdot u_{i'}\cdot\Theta_{\Pi'}(\zeta_{i'})\,\zeta_{i'} .
\]
Since $\Theta_{\Pi'}^2=1$ (Lemma~\ref{lem:sign-rules}), \eqref{eq:chamber-second-mirror} gives $\zeta_{i'}=\Theta_{\Pi'}(\check\zeta_{i'})=\chi(g_0)\cdot\Theta(\check\zeta_{i'})$, the degree being one;
as $\check\zeta_{i'}$ and $\Theta(\check\zeta_{i'})$ are odd, they anticommute, and
\[
\Theta_{\Pi'}(\zeta_{i'})\,\zeta_{i'}=\chi(g_0)\cdot\check\zeta_{i'}\,\Theta(\check\zeta_{i'})=-\chi(g_0)\cdot\Theta(\check\zeta_{i'})\,\check\zeta_{i'} .
\]
Together with $S_-=\Theta(S_+)$ this proves \eqref{eq:chamber-action-split}.

\emph{The exponential.}
The four parts of $S_{\mathbb{T}}(A)$ are even and commute with one another;
hence $\exp\left(-S_{\mathbb{T}}(A)\right)$ is the product of their exponentials.
Moreover $\Theta\left(\exp(-S_+)\right)=\exp\left(-\Theta(S_+)\right)$:
$\Theta$ maps a power of the even element $S_+$ to the same power of $\Theta(S_+)$, and the exponential series has real coefficients.
For $i\in\idx_{\mathbb{T}}(\Pi)$ the element $\Theta(\zeta_i)\,\zeta_i$ is even,
and its square is $-\Theta(\zeta_i)^2\,\zeta_i^2=0$, odd elements having square zero.
So $\exp\left(\omega_i\cdot u_i\cdot\Theta(\zeta_i)\,\zeta_i\right)=1+\omega_i\cdot u_i\cdot\Theta(\zeta_i)\,\zeta_i$,
and since these even elements commute,
\[
\exp\left(-S^\times_\Pi\right)=\prod_{i\in\idx_{\mathbb{T}}(\Pi)}\left(1+\omega_i\cdot u_i\cdot\Theta(\zeta_i)\,\zeta_i\right)
=\sum_{I\subset\idx_{\mathbb{T}}(\Pi)}\omega_I\cdot u_I\cdot\prod_{i\in I}\Theta(\zeta_i)\,\zeta_i .
\]
For $I=\{i_1<\dots<i_d\}$ one has
$\Theta(\zeta_I)\,\zeta_I=\Theta(\zeta_{i_d})\cdots\Theta(\zeta_{i_2})\cdot\Theta(\zeta_{i_1})\,\zeta_{i_1}\cdot\zeta_{i_2}\cdots\zeta_{i_d}$;
the even element $\Theta(\zeta_{i_1})\,\zeta_{i_1}$ may be moved to the left,
and induction on $d$ gives $\Theta(\zeta_I)\,\zeta_I=\prod_{i\in I}\Theta(\zeta_i)\,\zeta_i$.
This computation uses only that the $\zeta_i$ are odd;
applied to the odd elements $\check\zeta_{i'}$, with the coefficients $-\chi(g_0)\cdot\omega_{i'}\cdot u_{i'}$, it gives
\[
\exp\left(-S^\times_{\Pi'}\right)=\sum_{I'\subset\idx_{\mathbb{T}}(\Pi')}\left(-\chi(g_0)\right)^{\#|I'|}\cdot\omega_{I'}\cdot u_{I'}\cdot\Theta(\check\zeta_{I'})\,\check\zeta_{I'} .
\]
The product of the four exponentials is \eqref{eq:chamber-weight-expansion}.
\end{proof}

On a torus adapted to $\ell$ one has $\chi(g_0)=-1$,
and the crossing terms at the two mirrors enter \eqref{eq:chamber-action-split} and \eqref{eq:chamber-weight-expansion} with the same sign.
The factor $-\chi(g_0)$ at $\Pi'$ reflects the orientation:
seen from $U_+$, the normal $n$ points into $U_+$ at $\Pi$ and out of $U_+$ at $\Pi'$.

\begin{lemma}[Product rule and factorisation]\label{lem:product-rule}
Let $\ell\in\mathcal{L}$, let $\mathbb{T}$ be a torus that is split for $\ell$, and let $c\in\Adm(\ell)$;
let $\mathbb{X}_\pm$, $\Gr_{\mathbb{T}}^{\pm}$ and $\Theta=\Theta_\Pi$ be as in the paragraph on the two halves of the Grassmann algebra.
For $\xi,\xi',\xi''\in\Gr_{\mathbb{T}}^{+}$ of definite parities $|\xi|,|\xi'|,|\xi''|\in\{0,1\}$ one has
\begin{equation}\label{eq:chamber-product-rule}
\Theta(\xi')\,\xi\cdot\Theta(\xi'')\,\xi''=(-1)^{(|\xi|+|\xi'|)\,|\xi''|}\cdot\Theta(\xi'\xi'')\,\xi\,\xi''
\end{equation}
and
\begin{equation}\label{eq:chamber-factorisation}
\int_{\mathbb{T}}\Theta(\xi')\,\xi=\overline{\int_{\mathbb{X}_+}\xi'}\cdot\int_{\mathbb{X}_+}\xi .
\end{equation}
\end{lemma}

\begin{proof}
The map $\Theta$ preserves parities,
and two elements of definite parities of a Grassmann algebra commute if one of them is even and anticommute if both are odd.
Hence $\xi\,\Theta(\xi'')=(-1)^{|\xi|\,|\xi''|}\,\Theta(\xi'')\,\xi$,
and, $\Theta$ being an antilinear anti-automorphism and the sign real,
$\Theta(\xi')\,\Theta(\xi'')=\Theta(\xi''\xi')=(-1)^{|\xi'|\,|\xi''|}\,\Theta(\xi'\xi'')$.
Therefore
\[
\Theta(\xi')\,\xi\cdot\Theta(\xi'')\,\xi''
=(-1)^{|\xi|\,|\xi''|}\,\Theta(\xi')\,\Theta(\xi'')\,\xi\,\xi''
=(-1)^{(|\xi|+|\xi'|)\,|\xi''|}\,\Theta(\xi'\xi'')\,\xi\,\xi'' ,
\]
which is \eqref{eq:chamber-product-rule}.

Computed with the representatives $\mathbb{X}_-\cup\mathbb{X}_+$ (Definition~\ref{def:tori}),
$\int_{\mathbb{T}}$ is the coefficient of $\Psi_{\mathbb{X}_-\cup\mathbb{X}_+}=\Psi_-\,\Psi_+$;
no sign occurs in this factorisation, the pairs $\psi_\alpha(x)\,\bar\psi_\alpha(x)$ being even.
The element $\Theta(\xi')$ lies in $\Gr(\mathbb{X}_-)$ and $\xi$ in $\Gr(\mathbb{X}_+)$.
The product of a monomial in the generators at $\mathbb{X}_-$ and a monomial in the generators at $\mathbb{X}_+$
is a multiple of $\Psi_-\,\Psi_+$ only if the first is a multiple of $\Psi_-$ and the second a multiple of $\Psi_+$.
So the coefficient of $\Psi_-\,\Psi_+$ in $\Theta(\xi')\,\xi$ is the product of the coefficient of $\Psi_-$ in $\Theta(\xi')$ and the coefficient of $\Psi_+$ in $\xi$.
The latter is $\int_{\mathbb{X}_+}\xi$,
and the former is $\int_{\mathbb{X}_-}\Theta(\xi')=\overline{\int_{\mathbb{X}_+}\xi'}$ by Proposition~\ref{prop:symmetries}\eqref{it:berezin-conjugate}, as $\theta\,\mathbb{X}_+=\mathbb{X}_-$.
This is \eqref{eq:chamber-factorisation}.
If $\xi$ or $\xi'$ is odd, both sides of \eqref{eq:chamber-factorisation} vanish, since $\Psi_+$ and $\Psi_-$ have the even degree $8\,\#|\mathbb{X}_+|$.
\end{proof}

\subsubsection{Kinds of cells}\label{sec:cell-kinds}

The cells of $\Z^3$ are closed subsets of $\R^3$ (Definition~\ref{def:cochains}).
For a cell $\sigma$ let $\operatorname{aff}\sigma$ be its affine hull,
and let $\mathring{\sigma}$ be its relative interior, that is, the interior of $\sigma$ in $\operatorname{aff}\sigma$.
Every cell can be written in exactly one way as
\[
\sigma=x+\sum_{\mu\in J}[0,1]\,e_\mu,\qquad x\in\Z^3,\quad J\subset\{1,2,3\};
\]
the point $x$ is its first corner, $J$ is its set of \emph{directions},
and $\sigma$ is a site, a link, a plaquette or a cube according as $\#|J|=0,1,2,3$.
Its relative interior is
\[
\mathring{\sigma}=\left\{y\in\R^3:\ x_\mu<y_\mu<x_\mu+1\ \text{for }\mu\in J,\quad y_\mu=x_\mu\ \text{for }\mu\notin J\right\}.
\]
Consequently a point $y\in\R^3$ lies in the relative interior of exactly one cell,
namely the cell with the directions $J=\{\mu:\ y_\mu\notin\Z\}$ and the first corner $x=\left(\lfloor y_1\rfloor,\lfloor y_2\rfloor,\lfloor y_3\rfloor\right)$.
Two distinct cells of $\Z^3$, whether sites, links, plaquettes or cubes, therefore have disjoint relative interiors.

For a plaquette $\sigma$ and a link $b$ of $\Z^3$, and for a cube $Q$ and a plaquette $\sigma$,
the \emph{incidence numbers} $[\sigma:b]$ and $[Q:\sigma]$ are the coefficients of the coboundary of Definition~\ref{def:cochains}:
\[
(dA)_\sigma=\sum_b[\sigma:b]\cdot A(b),
\qquad
(dF)(Q)=\sum_\sigma[Q:\sigma]\cdot F_\sigma
\qquad\left(A\in C^1,\ F\in C^2\right).
\]
Thus $[\sigma:b]=\pm1$ if $b$ is one of the four links of $\sigma$ and $[\sigma:b]=0$ otherwise,
and $[Q:\sigma]=\pm1$ if $\sigma$ is one of the six plaquettes of $Q$ and $[Q:\sigma]=0$ otherwise.
Since $d(dA)=0$ for every $A\in C^1$ (the difference operators in the three coordinate directions commute),
\[
\sum_\sigma[Q:\sigma]\cdot[\sigma:b]=0
\qquad\text{for every cube $Q$ and every link $b$.}
\]

Let $\Gamma$ be a period lattice, with the cells and cochains of $\Z^3/\Gamma$ of Definition~\ref{def:configurations}.
Two distinct cells that are contained in one cube $x+[0,1]^3$ and are translates of each other
differ by a non-zero vector of $\{-1,0,1\}^3$;
its maximum norm is $1<4\le\sys(\Gamma)$, so the two cells are not equivalent modulo $\Gamma$.
For a plaquette $\sigma$ and a link $b$ of $\Z^3/\Gamma$ one may therefore define $[\sigma:b]$
as the incidence number of a representative of $\sigma$ and of the representative of $b$ that is one of its links,
and as $0$ if no representative of $b$ is a link of it.
The incidence numbers of $\Z^3$ do not change when both cells are translated by one vector of $\Z^3$,
the coefficients in the formulas of Definition~\ref{def:cochains} being independent of the point;
so this number does not depend on the representative of $\sigma$.
The incidence number $[Q:\sigma]$ of a cube and a plaquette of $\Z^3/\Gamma$ is defined in the same way.
With these numbers the two formulas for $dA$ and $dF$ hold for $A\in C^1(\Gamma)$ and $F\in C^2(\Gamma)$,
and $\sum_\sigma[Q:\sigma]\cdot[\sigma:b]=0$ for every cube $Q$ and every link $b$ of $\Z^3/\Gamma$,
all sums being taken over the cells of $\Z^3/\Gamma$:
the cells that occur in one of these formulas are contained in one cube,
and distinct ones among them represent distinct cells of $\Z^3/\Gamma$.

Let now $\mathbb{T}=(\Gamma,\chi)$ be a torus that is split for a family $\ell\in\mathcal{L}$,
let $p$ be its height period for $\ell$, let $c\in\Adm(\ell)$ and $c'=c+\frac p2$,
and let $\Gamma_\parallel$, $g_0$, $\theta=\theta_{\ell,c}$, $U_+$, $U_-$ and $\mathcal{M}$ be as in Definition~\ref{def:split-tori}.
The set $\mathcal{M}$ is the union of the planes $\left\{\ell=c+j\frac p2\right\}$, $j\in\Z$, which are called the planes of $\mathcal{M}$.
As $\ell(g_0)=p$, the translates of $U_+$ and $U_-$ by $jg_0$, $j\in\Z$, are the open slabs
\[
U_{+}+jg_0=\left\{c+jp<\ell<c'+jp\right\},
\qquad
U_{-}+jg_0=\left\{c-\frac p2+jp<\ell<c+jp\right\},
\]
and $U_{+}+\Z g_0$ and $U_{-}+\Z g_0$ denote the unions of these slabs over $j\in\Z$.
By Definition~\ref{def:split-tori}\eqref{it:split-slabs}, the set $\R^3\setminus\mathcal{M}$ is the disjoint union of all these slabs.
A link, plaquette or cube $\sigma$ of $\Z^3$
\begin{itemize}
\item \emph{lies in} $\mathcal{M}$ if $\sigma\subset\mathcal{M}$;
\item is \emph{cut} if $\mathring{\sigma}$ meets $\mathcal{M}$ and $\sigma\not\subset\mathcal{M}$,
and it is \emph{cut by} a plane $H$ of $\mathcal{M}$ if in addition $\mathring{\sigma}$ meets $H$;
\item is \emph{upper} if $\mathring{\sigma}\subset U_{+}+\Z g_0$;
\item is \emph{lower} if $\mathring{\sigma}\subset U_{-}+\Z g_0$.
\end{itemize}
These are the four \emph{kinds} of cells; they depend on $\mathbb{T}$, $\ell$ and $c$.
The \emph{faces} of a plaquette are its four links, and the faces of a cube are its six plaquettes.

\begin{lemma}[Kinds of cells]\label{lem:cell-kinds}
Let $\mathbb{T}=(\Gamma,\chi)$ be a torus that is split for $\ell\in\mathcal{L}$, let $c\in\Adm(\ell)$,
and let $\theta=\theta_{\ell,c}$ and $\mathcal{M}$ be as in Definition~\ref{def:split-tori}.
Every link, plaquette and cube of $\Z^3$ is of exactly one of the four kinds:
it lies in $\mathcal{M}$, or it is cut, or it is upper, or it is lower.
For $g\in\Gamma$ the cells $\sigma$ and $\sigma+g$ are of the same kind,
so that every link, plaquette and cube of $\Z^3/\Gamma$ has a kind.
Moreover the following hold.
\begin{enumerate}
\item\label{it:kinds-one-plane}
A link, plaquette or cube meets at most one plane of $\mathcal{M}$.
\item\label{it:kinds-reflection}
The reflection $\theta$ maps every upper cell onto a lower cell and every lower cell onto an upper cell.
It maps every cell $\sigma$ of $\Z^3/\Gamma$ that lies in $\mathcal{M}$ onto itself and preserves its orientation,
and it maps every cut cell $\sigma$ of $\Z^3/\Gamma$ onto itself and reverses its orientation:
for every $u\in C^k(\Gamma)$, $k$ the dimension of $\sigma$,
\[
(\theta^*u)(\sigma)=u(\sigma)\quad\text{if $\sigma$ lies in $\mathcal{M}$},
\qquad
(\theta^*u)(\sigma)=-u(\sigma)\quad\text{if $\sigma$ is cut}.
\]
\item\label{it:kinds-upper-faces}
The faces of an upper cell are upper or lie in $\mathcal{M}$.
The links of a plaquette that lies in $\mathcal{M}$ lie in $\mathcal{M}$.
\item\label{it:kinds-cut-faces}
The faces of a cut cell are upper, lower or cut.
\item\label{it:kinds-cut-cofaces}
If a link of a plaquette $\sigma$ is cut by a plane $H$ of $\mathcal{M}$, then $\sigma$ is cut by $H$;
if a plaquette of a cube $Q$ is cut by a plane $H$ of $\mathcal{M}$, then $Q$ is cut by $H$.
\end{enumerate}
\end{lemma}

\begin{proof}
The relative interior of a cell is convex, hence connected.
If $\mathring{\sigma}$ does not meet $\mathcal{M}$,
it is a connected subset of the disjoint union of the open slabs $U_{+}+jg_0$, $U_{-}+jg_0$, $j\in\Z$,
and so it lies in one of them: $\sigma$ is upper or lower.
If $\mathring{\sigma}$ meets $\mathcal{M}$, then $\sigma$ lies in $\mathcal{M}$ or is cut.
No cell is of two kinds.
Indeed, the relative interior of an upper or of a lower cell does not meet $\mathcal{M}$,
whereas the non-empty relative interior of a cell that lies in $\mathcal{M}$ or is cut meets $\mathcal{M}$;
the sets $U_{+}+\Z g_0$ and $U_{-}+\Z g_0$ are disjoint, the slabs being pairwise disjoint;
and a cut cell does not lie in $\mathcal{M}$ by definition.

Let $g\in\Gamma$.
Since $\mathbb{T}$ is split, $g$ is the sum of an element of $\Gamma_\parallel$, on which $\ell$ vanishes,
and of an integer multiple of $g_0$; hence $\ell(g)\in p\Z$.
The translation by $g$ maps the plane $\left\{\ell=c+j\frac p2\right\}$ onto the plane $\left\{\ell=c+j\frac p2+\ell(g)\right\}$,
which is again a plane of $\mathcal{M}$,
and it maps the slab $U_\pm+jg_0$ onto the slab $U_\pm+j'g_0$ with $j'=j+\ell(g)/p\in\Z$.
Therefore $\mathcal{M}+g=\mathcal{M}$ and $(U_\pm+\Z g_0)+g=U_\pm+\Z g_0$,
and, the relative interior of $\sigma+g$ being $\mathring{\sigma}+g$, the cells $\sigma$ and $\sigma+g$ are of the same kind.

\emph{Proof of \eqref{it:kinds-one-plane}.}
A cell is contained in a cube $x+[0,1]^3$, so that $|y_\mu-y'_\mu|\le1$ for any two of its points $y,y'$ and $\mu=1,2,3$.
The coordinates of $n_\ell$ lie in $\{-1,0,1\}$; hence
\[
|\ell(y)-\ell(y')|=|n_\ell\cdot(y-y')|\le\sum_{\mu=1}^{3}|(n_\ell)_\mu|=|n_\ell|^2 .
\]
The heights of two distinct planes of $\mathcal{M}$ differ by at least $\frac p2$, and $\frac p2>|n_\ell|^2$:
for an axis family $|n_\ell|^2=1$ and $p\ge4$, for a diagonal family $|n_\ell|^2=2$ and $p\ge8$
(Definition~\ref{def:split-tori}\eqref{it:split-height}).
So no cell meets two planes of $\mathcal{M}$.

\emph{Proof of \eqref{it:kinds-reflection}.}
Let $H=\left\{\ell=c+j\frac p2\right\}$ be a plane of $\mathcal{M}$.
The numbers $c$ and $c'$ lie in $\Adm(\ell)$ (Definition~\ref{def:split-tori}\eqref{it:split-height}),
and $\Adm(\ell)$ is a subgroup of $\R$; hence $c+j\frac p2\in\Adm(\ell)$, and $\left(\ell,c+j\frac p2\right)$ is a mirror.
Let $\theta_H$ be its reflection.
Since $\frac{2}{|n_\ell|^2}\cdot j\frac p2\cdot n_\ell=jg_0$,
Definition~\ref{def:mirrors}\eqref{it:mirror-involution} and~\eqref{it:mirror-product} give
\[
\theta_H=\tau_{jg_0}\,\theta .
\]
The reflection $\theta$ maps $\Z^3$ onto itself (Lemma~\ref{lem:chambers}\eqref{it:chambers-group}), and $jg_0\in\Z^3$;
so $\theta_H$ maps $\Z^3$ onto itself and permutes the cells of $\Z^3$ (Definition~\ref{def:cochains}).
Since $jg_0\in\Gamma$, the maps $\theta_H$ and $\theta$ induce the same permutation of the cells of $\Z^3/\Gamma$,
and $\theta_H^*=\theta^*\,\tau_{jg_0}^*=\theta^*$ on $\Gamma$-periodic cochains,
because a translation preserves the orientations of the cells.

The linear part of $\theta$ maps $g_0$ to $-g_0$ (Definition~\ref{def:split-tori}\eqref{it:split-reflections}),
so that $\theta(y+jg_0)=\theta y-jg_0$ for $y\in\R^3$ and $j\in\Z$.
Together with $\theta U_+=U_-$ (Definition~\ref{def:split-tori}\eqref{it:split-slabs}) and $\theta^2=\mathrm{id}$ this gives
\[
\theta\left(U_{+}+\Z g_0\right)=U_{-}+\Z g_0,
\qquad
\theta\left(U_{-}+\Z g_0\right)=U_{+}+\Z g_0 .
\]
An affine bijection maps the relative interior of a cell onto the relative interior of the image cell;
hence $\theta$ maps upper cells onto lower cells and lower cells onto upper cells.

Now let $\sigma$ be a cell of $\Z^3$ that lies in $\mathcal{M}$ or is cut.
Its relative interior meets a plane $H$ of $\mathcal{M}$:
if $\sigma$ is cut, by definition;
if $\sigma$ lies in $\mathcal{M}$, because by \eqref{it:kinds-one-plane} it is then contained in a single plane $H$ of $\mathcal{M}$.
Let $x\in\mathring{\sigma}\cap H$.
The point $x=\theta_Hx$ lies in $\theta_H(\mathring{\sigma})$, which is the relative interior of the cell $\theta_H\sigma$;
the relative interiors of distinct cells being disjoint, $\theta_H\sigma=\sigma$.
If $\sigma\subset H$, then $\theta_H$ fixes $\sigma$ pointwise and preserves its orientation.
If $\sigma\not\subset H$, then $\theta_H$ maps $\operatorname{aff}\sigma$ onto itself
and fixes pointwise the set $H\cap\operatorname{aff}\sigma$,
which is an affine hyperplane of $\operatorname{aff}\sigma$, since it contains $x$ and $\operatorname{aff}\sigma\not\subset H$;
moreover $\theta_H$ is not the identity on $\operatorname{aff}\sigma$, its fixed points in $\R^3$ being the points of $H$.
An isometry of $\operatorname{aff}\sigma$ that fixes an affine hyperplane pointwise and is not the identity is the reflection in this hyperplane.
Hence $\theta_H$ restricted to $\operatorname{aff}\sigma$ is the reflection in $H\cap\operatorname{aff}\sigma$,
and it reverses the orientation of $\sigma$.
In the first case $(\theta_H^*u)(\sigma)=u(\sigma)$, in the second $(\theta_H^*u)(\sigma)=-u(\sigma)$,
for every cochain $u$ whose degree is the dimension of $\sigma$.
The kind of a cell of $\Z^3/\Gamma$ is the kind of each of its representatives;
so the above applies to every representative of a cell of $\Z^3/\Gamma$ that lies in $\mathcal{M}$ or is cut,
and, as $\theta_H$ and $\theta$ induce the same permutation of the cells of $\Z^3/\Gamma$ and $\theta_H^*=\theta^*$ on $C^k(\Gamma)$,
the remaining assertions of \eqref{it:kinds-reflection} follow.

\emph{Proof of \eqref{it:kinds-upper-faces}.}
Let $\sigma$ be upper.
Its relative interior is connected and therefore lies in one slab $U_{+}+jg_0=\{c+jp<\ell<c'+jp\}$.
The cell $\sigma$ is the closure of $\mathring{\sigma}$;
so $\sigma$ and all its faces lie in the closed slab $\{c+jp\le\ell\le c'+jp\}$.
Let $\sigma'$ be a face of $\sigma$,
and suppose that $\mathring{\sigma'}$ contains a point $y$ of one of the two planes $\{\ell=c+jp\}$, $\{\ell=c'+jp\}$.
Then the affine function $\ell$ attains its minimum or its maximum over $\sigma'$ at $y$.
For $y'\in\sigma'$ and $\varepsilon>0$ small enough the point $y+\varepsilon(y-y')$ lies in $\sigma'$, because $y$ is a relative interior point, and
\[
\ell\left(y+\varepsilon(y-y')\right)-\ell(y)=\varepsilon\left(\ell(y)-\ell(y')\right).
\]
At a minimum the left side is non-negative, so $\ell(y')\le\ell(y)\le\ell(y')$;
at a maximum it is non-positive, so $\ell(y')\ge\ell(y)\ge\ell(y')$.
In both cases $\ell(y')=\ell(y)$.
Thus $\ell$ is constant on $\sigma'$, and $\sigma'$ is contained in a plane of $\mathcal{M}$.
If, on the contrary, $\mathring{\sigma'}$ meets neither of the two planes,
then it lies in the open slab $U_{+}+jg_0$, and $\sigma'$ is upper.
Finally, a link of a plaquette that lies in $\mathcal{M}$ is a subset of the plaquette, hence of $\mathcal{M}$.

\emph{Proof of \eqref{it:kinds-cut-faces}.}
Let $\sigma$ be a cut plaquette or a cut cube, and let $H$ be a plane of $\mathcal{M}$ that meets $\mathring{\sigma}$.
Suppose that a face $\sigma'$ of $\sigma$ lies in $\mathcal{M}$.
By \eqref{it:kinds-one-plane} the face $\sigma'$ is contained in one plane of $\mathcal{M}$;
this plane meets $\sigma\supset\sigma'$, and so, by \eqref{it:kinds-one-plane} again, it is $H$.
As shown above, $\theta_H\sigma=\sigma$, and $\theta_H$ restricted to $\operatorname{aff}\sigma$
is the reflection of $\operatorname{aff}\sigma$ in its hyperplane $H\cap\operatorname{aff}\sigma$.
This hyperplane contains $\sigma'$, and the affine hull of $\sigma'$ is a hyperplane of $\operatorname{aff}\sigma$;
hence $H\cap\operatorname{aff}\sigma=\operatorname{aff}\sigma'$.
Write $\sigma=x+\sum_{\mu\in J}[0,1]\,e_\mu$.
The face has the form $\sigma'=x'+\sum_{\mu\in J\setminus\{\nu\}}[0,1]\,e_\mu$ with $\nu\in J$ and $x'\in\{x,x+e_\nu\}$.
The reflection of $\operatorname{aff}\sigma$ in $\operatorname{aff}\sigma'$ fixes the directions $e_\mu$, $\mu\in J\setminus\{\nu\}$, and reverses $e_\nu$;
it maps $\sigma$ onto $\sigma-e_\nu$ if $x'=x$, and onto $\sigma+e_\nu$ if $x'=x+e_\nu$.
Neither of these cells is $\sigma$, which contradicts $\theta_H\sigma=\sigma$.
So no face of $\sigma$ lies in $\mathcal{M}$, and every face of $\sigma$ is upper, lower or cut.

\emph{Proof of \eqref{it:kinds-cut-cofaces}.}
Let $\sigma'$ be a face of a plaquette or of a cube $\sigma$,
and let $\sigma'$ be cut by the plane $H=\left\{\ell=c+j\frac p2\right\}$ of $\mathcal{M}$.
The affine function $\ell-c-j\frac p2$ vanishes at a point $y\in\mathring{\sigma'}$,
and it does not vanish identically on $\sigma'$, because $\sigma'\not\subset\mathcal{M}$.
Choose $y'\in\sigma'$ at which it does not vanish.
For $\varepsilon>0$ small enough the point $y+\varepsilon(y-y')$ lies in $\sigma'$,
and there the function takes the value $-\varepsilon\left(\ell(y')-c-j\frac p2\right)$,
whose sign is opposite to its sign at $y'$.
So the function takes both signs on $\sigma'$.
Every point of $\sigma$ is a limit of points of $\mathring{\sigma}$;
by continuity the function takes both signs on $\mathring{\sigma}$ as well,
and since $\mathring{\sigma}$ is convex, it vanishes at some point of $\mathring{\sigma}$.
Thus $\mathring{\sigma}$ meets $H$, while $\sigma\not\subset\mathcal{M}$ because $\sigma'\not\subset\mathcal{M}$:
the cell $\sigma$ is cut by $H$.
\end{proof}

\begin{example}
Let $H$ be a plane of $\mathcal{M}$,
and let $\sigma$ be a link, plaquette or cube with the first corner $x$ and the set of directions $J$.
By Lemma~\ref{lem:cell-kinds}\eqref{it:kinds-one-plane},
a cell that meets $H$ lies in $\mathcal{M}$ if and only if it is contained in $H$,
and a cell is cut by $H$ if and only if its relative interior meets $H$ and it is not contained in $H$.
The description of $\mathring{\sigma}$ in coordinates then gives the following.
\begin{enumerate}
\item
Let $\ell=x_\mu$ and $H=\{x_\mu=t\}$ with $t\in\Z$, a site plane.
If $\mathring{\sigma}$ meets $H$, then $\mu\notin J$, since the interval $(x_\mu,x_\mu+1)$ contains no integer;
then $x_\mu=t$ and $\sigma\subset H$.
So the links and the plaquettes contained in $H$ lie in $\mathcal{M}$, and no cell is cut by $H$.
\item
Let $\ell=x_\mu$ and $H=\left\{x_\mu=t+\frac12\right\}$ with $t\in\Z$, a link plane.
No cell is contained in $H$, since every cell contains a point of $\Z^3$.
The relative interior $\mathring{\sigma}$ meets $H$ if and only if $\mu\in J$ and $x_\mu=t$.
So the cells cut by $H$ are the links $(x,\mu)$, the plaquettes with the direction $\mu$, and the cubes $x+[0,1]^3$,
in each case with $x_\mu=t$.
\item
Let $\ell=x_\mu\mp x_\nu$ with $\mu<\nu$, let $\nu'$ be the third direction, and let $H=\{x_\mu\mp x_\nu=t\}$ with $t\in\Z$.
On $\mathring{\sigma}$ the function $\ell$ is constant, equal to the integer $\ell(x)$, if $J\cap\{\mu,\nu\}$ is empty;
its values fill an open interval between two consecutive integers if $J\cap\{\mu,\nu\}$ has one element;
and they fill an open interval of length two with an integer midpoint if $\{\mu,\nu\}\subset J$,
the midpoint being the value of $\ell$ at two opposite corners of each plaquette with the directions $\mu$, $\nu$ that is contained in $\sigma$.
Hence the cells contained in $H$ are the links $(x,\nu')$ whose end points lie on $H$, and these lie in $\mathcal{M}$;
no plaquette is contained in $H$, and no link is cut by $H$.
The cells cut by $H$ are the plaquettes $(x;\mu\nu)$ two opposite corners of which lie on $H$,
and the cubes that contain such a plaquette.
\end{enumerate}
\end{example}

The links of $\Z^3/\Gamma$ that are upper, lower, lie in $\mathcal{M}$ or are cut form the sets
\[
\mathbb{T}^{(1)}_{+},\qquad\mathbb{T}^{(1)}_{-},\qquad\mathbb{T}^{(1)}_{0},\qquad\mathbb{T}^{(1)}_{\times},
\]
respectively, and the plaquettes of the four kinds form the sets
$\mathbb{T}^{(2)}_{+}$, $\mathbb{T}^{(2)}_{-}$, $\mathbb{T}^{(2)}_{0}$, $\mathbb{T}^{(2)}_{\times}$.
A cube is not contained in a plane, so no cube lies in $\mathcal{M}$;
the upper, lower and cut cubes of $\Z^3/\Gamma$ form the sets
$\mathbb{T}^{(3)}_{+}$, $\mathbb{T}^{(3)}_{-}$, $\mathbb{T}^{(3)}_{\times}$.

\begin{lemma}[Coboundaries at cut cells]\label{lem:cut-coboundaries}
Let $\mathbb{T}=(\Gamma,\chi)$ be a torus that is split for $\ell\in\mathcal{L}$, let $c\in\Adm(\ell)$ and $\theta=\theta_{\ell,c}$,
and let the sets of cells of the four kinds be formed with respect to $\mathbb{T}$, $\ell$ and $c$.
For $\sigma\in\mathbb{T}^{(2)}_{\times}$, $Q\in\mathbb{T}^{(3)}_{\times}$, $A\in C^1(\Gamma)$ and $z\in C^2(\Gamma)$ put
\[
(d_+A)_\sigma:=\sum_{b\in\mathbb{T}^{(1)}_{+}}[\sigma:b]\cdot A(b),
\qquad
(d_\times A)_\sigma:=\sum_{b\in\mathbb{T}^{(1)}_{\times}}[\sigma:b]\cdot A(b),
\]
\[
(d_+z)(Q):=\sum_{\sigma'\in\mathbb{T}^{(2)}_{+}}[Q:\sigma']\cdot z_{\sigma'},
\qquad
(d_\times z)(Q):=\sum_{\sigma'\in\mathbb{T}^{(2)}_{\times}}[Q:\sigma']\cdot z_{\sigma'}.
\]
Then for all $\sigma\in\mathbb{T}^{(2)}_{\times}$, $Q\in\mathbb{T}^{(3)}_{\times}$, $A\in C^1(\Gamma)$ and $z\in C^2(\Gamma)$
\[
(dA)_\sigma=(d_+A)_\sigma-(d_+\theta^*A)_\sigma+(d_\times A)_\sigma,
\qquad
(dz)(Q)=(d_+z)(Q)-(d_+\theta^*z)(Q)+(d_\times z)(Q).
\]
\end{lemma}

\begin{proof}
Fix $\sigma\in\mathbb{T}^{(2)}_{\times}$.
By Lemma~\ref{lem:cell-kinds}\eqref{it:kinds-reflection}, $(\theta^*F)_\sigma=-F_\sigma$ for every $F\in C^2(\Gamma)$.
As $\theta^*$ commutes with $d$ (Definition~\ref{def:cochains}), every $A\in C^1(\Gamma)$ satisfies
\[
\sum_b[\sigma:b]\cdot(\theta^*A)(b)=(d\theta^*A)_\sigma=(\theta^*dA)_\sigma=-(dA)_\sigma=-\sum_b[\sigma:b]\cdot A(b),
\]
the sums being taken over the links of $\Z^3/\Gamma$.
For each link $b$ one has $(\theta^*A)(b)=\pm A(\theta b)$, with a sign that depends on $b$ only (Definition~\ref{def:cochains}),
and $\theta$ permutes the links of $\Z^3/\Gamma$ (Definition~\ref{def:split-tori}\eqref{it:split-reflections}).
Both sides are linear forms in the values of $A$ at the links of $\Z^3/\Gamma$.
On the left the only term that contains $A(\theta b)$ is the term of $b$, and its coefficient is $\pm[\sigma:b]$, with the sign of $(\theta^*A)(b)=\pm A(\theta b)$;
on the right the coefficient of $A(\theta b)$ is $-[\sigma:\theta b]$.
Comparing the two coefficients one obtains
\[
[\sigma:b]\cdot(\theta^*A)(b)=-[\sigma:\theta b]\cdot A(\theta b)
\qquad\text{for every link $b$ of $\Z^3/\Gamma$ and every $A\in C^1(\Gamma)$.}
\]
By Lemma~\ref{lem:cell-kinds}\eqref{it:kinds-reflection}, the involution $\theta$ maps $\mathbb{T}^{(1)}_{+}$ onto $\mathbb{T}^{(1)}_{-}$.
Summing the last identity over $b\in\mathbb{T}^{(1)}_{+}$ therefore gives
\[
(d_+\theta^*A)_\sigma=-\sum_{b\in\mathbb{T}^{(1)}_{+}}[\sigma:\theta b]\cdot A(\theta b)=-\sum_{b\in\mathbb{T}^{(1)}_{-}}[\sigma:b]\cdot A(b) .
\]
The links $b$ with $[\sigma:b]\neq0$ are the links of $\sigma$,
and by Lemma~\ref{lem:cell-kinds}\eqref{it:kinds-cut-faces} each of them is upper, lower or cut.
Splitting the sum $(dA)_\sigma=\sum_b[\sigma:b]\cdot A(b)$ over the links of these three kinds yields
$(dA)_\sigma=(d_+A)_\sigma-(d_+\theta^*A)_\sigma+(d_\times A)_\sigma$.

For the cubes, fix $Q\in\mathbb{T}^{(3)}_{\times}$.
Lemma~\ref{lem:cell-kinds}\eqref{it:kinds-reflection} gives $(\theta^*u)(Q)=-u(Q)$ for every $u\in C^3(\Gamma)$; hence every $z\in C^2(\Gamma)$ satisfies
\[
\sum_{\sigma'}[Q:\sigma']\cdot(\theta^*z)_{\sigma'}=(d\theta^*z)(Q)=(\theta^*dz)(Q)=-(dz)(Q)=-\sum_{\sigma'}[Q:\sigma']\cdot z_{\sigma'},
\]
the sums being taken over the plaquettes of $\Z^3/\Gamma$.
For each plaquette $\sigma'$ one has $(\theta^*z)_{\sigma'}=\pm z_{\theta\sigma'}$, with a sign that depends on $\sigma'$ only,
and $\theta$ permutes the plaquettes of $\Z^3/\Gamma$.
The coefficient of $z_{\theta\sigma'}$ is therefore $\pm[Q:\sigma']$ on the left, with the sign of $(\theta^*z)_{\sigma'}=\pm z_{\theta\sigma'}$,
and $-[Q:\theta\sigma']$ on the right, so that
\[
[Q:\sigma']\cdot(\theta^*z)_{\sigma'}=-[Q:\theta\sigma']\cdot z_{\theta\sigma'}
\qquad\text{for every plaquette $\sigma'$ of $\Z^3/\Gamma$ and every $z\in C^2(\Gamma)$.}
\]
Since $\theta$ maps $\mathbb{T}^{(2)}_{+}$ onto $\mathbb{T}^{(2)}_{-}$ (Lemma~\ref{lem:cell-kinds}\eqref{it:kinds-reflection}), summation over $\sigma'\in\mathbb{T}^{(2)}_{+}$ gives
\[
(d_+\theta^*z)(Q)=-\sum_{\sigma'\in\mathbb{T}^{(2)}_{-}}[Q:\sigma']\cdot z_{\sigma'} .
\]
The plaquettes $\sigma'$ with $[Q:\sigma']\neq0$ are the six plaquettes of $Q$,
and each of them is upper, lower or cut (Lemma~\ref{lem:cell-kinds}\eqref{it:kinds-cut-faces}).
Splitting $(dz)(Q)=\sum_{\sigma'}[Q:\sigma']\cdot z_{\sigma'}$ over these three kinds of plaquettes yields
$(dz)(Q)=(d_+z)(Q)-(d_+\theta^*z)(Q)+(d_\times z)(Q)$.
\end{proof}

\subsubsection{Half weights and the crossing weight}\label{sec:half-weights}

\begin{definition}[Variables and half weights]\label{def:half-weights}
Let $\mathbb{T}=(\Gamma,\chi)$ be a torus that is split for $\ell\in\mathcal{L}$ and let $c\in\Adm(\ell)$.
Let $\Pi$, $\Pi'$, $\theta$, $\theta'$ and $\mathcal{M}$ be as in Definition~\ref{def:split-tori},
and let $\Theta=\Theta_\Pi$, the halves $\mathbb{X}_\pm$ of the sites and the algebras $\Gr_{\mathbb{T}}^{\pm}$ be as in Lemmas~\ref{lem:second-mirror} and~\ref{lem:product-rule}.
The links, plaquettes and cubes of $\Z^3/\Gamma$ carry the four kinds of Lemma~\ref{lem:cell-kinds},
and $\mathbb{T}^{(1)}$ is the set of all links of $\Z^3/\Gamma$.

\emph{Intervals.}
For $b\in\mathbb{T}^{(1)}_{+}\cup\mathbb{T}^{(1)}_{0}$ let $J_b:=[0,2\pi/q)$.
For $b\in\mathbb{T}^{(1)}_{+}$ let $J_{\theta b}:=[0,2\pi/q)$ if $(\theta^*A)(b)=A(\theta b)$ for all $A\in C^1(\Gamma)$,
and $J_{\theta b}:=(-2\pi/q,0]$ if $(\theta^*A)(b)=-A(\theta b)$ for all $A\in C^1(\Gamma)$;
as $\theta$ maps $\mathbb{T}^{(1)}_{+}$ onto $\mathbb{T}^{(1)}_{-}$ (Lemma~\ref{lem:cell-kinds}\eqref{it:kinds-reflection}),
this defines $J_b$ for every lower link $b$.
For $b\in\mathbb{T}^{(1)}_{\times}$ let $J_b:=[-\pi/q,\pi/q)$.
Every $J_b$ is thus a half-open interval of length $2\pi/q$.

\emph{Variables.}
Every $A\in C^1(\Gamma)$ and every $z\in C^2(\Gamma;\Z)$ is written as
\[
A=(A_+,A_-,A_0,a),\qquad z=(z_+,z_-,z_0,z_\times),
\]
the components being the restrictions of $A$ to the upper links, the lower links, the links in $\mathcal{M}$ and the cut links,
and of $z$ to the plaquettes of the same four kinds;
thus $a=(a_b)_{b\in\mathbb{T}^{(1)}_{\times}}$, with $a_b=A(b)$, collects the values of $A$ at the cut links.
These are the upper, lower, mirror and cut variables, and $y_+:=(A_+,z_+)$ collects the upper ones.
Let
\[
\mathcal{Y}_+:=[0,2\pi/q)^{\mathbb{T}^{(1)}_{+}}\times\Z^{\mathbb{T}^{(2)}_{+}},
\qquad
\mathcal{Y}_0:=[0,2\pi/q)^{\mathbb{T}^{(1)}_{0}}\times\Z^{\mathbb{T}^{(2)}_{0}},
\]
and let $\mathrm{d}y_+$ be the product of Lebesgue measure and counting measure on $\mathcal{Y}_+$.
As $J_b=[0,2\pi/q)$ for the upper links and for the links in $\mathcal{M}$,
the sets $\mathcal{Y}_+$ and $\mathcal{Y}_0$ are the ranges of the upper variables $y_+$ and of the mirror variables $(A_0,z_0)$
of the pairs $(A,z)$ with $A(b)\in J_b$ for every $b\in\mathbb{T}^{(1)}$.

\emph{Reflection of functions of representatives.}
For a function $f\colon C^1(\Gamma)\times C^2(\Gamma;\Z)\to\C$ let
\[
(\Theta f)(A,z):=\overline{f(\theta^*A,\theta^*z)},
\]
and for a function $f=\sum_\Psi f_\Psi\cdot\Psi$ with values in $\Gr_{\mathbb{T}}$, the sum being over the monomials $\Psi$, let $\Theta f:=\sum_\Psi(\Theta f_\Psi)\cdot\Theta(\Psi)$.
For functions on $\Conf(\Gamma)$ these are the reflections $\Theta_\Pi$ of Lemma~\ref{lem:reflection-classes}, for complex values,
and of Lemma~\ref{lem:reflected-algebras}, for values in $\Gr_{\mathbb{T}}$, the torus $\theta\mathbb{T}$ being $\mathbb{T}$ (Definition~\ref{def:split-tori}\eqref{it:split-reflections}).
A function of the upper and mirror variables is regarded as a function of $(A,z)$;
in particular a function on $\mathcal{Y}_+\times\mathcal{Y}_0$ is a function of those $(A,z)$ whose upper and mirror variables lie in $\mathcal{Y}_+\times\mathcal{Y}_0$.

\emph{Half weights.}
With $F=dA+(2\pi/q)z$, the half weights are
\[
\mathcal{W}_+:=\prod_{\sigma\in\mathbb{T}^{(2)}_{+}}\exp\left(-\frac{F_\sigma^2}{2}\right)\cdot\prod_{Q\in\mathbb{T}^{(3)}_{+}}1_{\{(dz)(Q)=0\}},
\qquad
\mathcal{W}_0:=\prod_{\sigma\in\mathbb{T}^{(2)}_{0}}\exp\left(-\frac{F_\sigma^2}{2}\right).
\]
\end{definition}

\paragraph{The fundamental domain.}
Each point of $\Conf(\Gamma)$ has exactly one representative $(A,z)$ with $A(b)\in J_b$ for every $b\in\mathbb{T}^{(1)}$.
Indeed, the equivalence of Definition~\ref{def:configurations} adds to $A$ the product of $2\pi/q$ with an arbitrary integer one-cochain and changes $z$ accordingly,
while each coset of $(2\pi/q)\Z$ in $\R$ meets the half-open interval $J_b$ in exactly one point.
The Borel set
\[
\prod_{b\in\mathbb{T}^{(1)}}J_b\times\{z\in C^2(\Gamma;\Z):\ dz=0\}
\]
is therefore a measurable fundamental domain,
and Definition~\ref{def:configurations} gives
\[
\int_{\Conf(\Gamma)}\Phi=\sum_{z\in C^2(\Gamma;\Z),\ dz=0}\ \int_{\prod_{b}J_b}\mathrm{d}A\,\Phi(A,z)
\]
for every Borel function $\Phi\ge0$ on $\Conf(\Gamma)$, and, by linearity, for every integrable one.

\paragraph{The two halves.}
By Lemma~\ref{lem:cell-kinds}\eqref{it:kinds-reflection}, $\theta$ maps every link and every plaquette in $\mathcal{M}$ onto itself with orientation sign $+1$,
every cut link onto itself with orientation sign $-1$,
and the upper cells bijectively onto the lower ones.
For all $A\in C^1(\Gamma)$ and $z\in C^2(\Gamma;\Z)$ it follows that
\[
(\theta^*A)_0=A_0,\qquad(\theta^*z)_0=z_0,\qquad(\theta^*A)(b)=-a_b\quad\left(b\in\mathbb{T}^{(1)}_{\times}\right),
\]
and that $(\theta^*A)_+$ and $(\theta^*z)_+$ arise from $A_-$ and $z_-$ by a permutation of the coordinates and multiplication by orientation signs.
For $b\in\mathbb{T}^{(1)}_{+}$ the choice of the intervals says that $(\theta^*A)(b)\in[0,2\pi/q)$ if and only if $A(\theta b)\in J_{\theta b}$.
Hence the map
\[
(A_-,z_-)\longmapsto\left((\theta^*A)_+,(\theta^*z)_+\right)
\]
is a bijection of $\prod_{b\in\mathbb{T}^{(1)}_{-}}J_b\times\Z^{\mathbb{T}^{(2)}_{-}}$ onto $\mathcal{Y}_+$,
and it carries the product of Lebesgue and counting measure to the measure $\mathrm{d}y_+$.
Consequently, if $f$ is a complex function of the upper and mirror variables, then
\[
(\Theta f)(A,z)=\overline{f\left((\theta^*A)_+,(\theta^*z)_+;A_0,z_0\right)},
\]
and if $f=\sum_\Psi f_\Psi\cdot\Psi$ takes its values in $\Gr_{\mathbb{T}}^{+}$, then
\[
(\Theta f)(A,z)=\sum_\Psi\overline{f_\Psi\left((\theta^*A)_+,(\theta^*z)_+;A_0,z_0\right)}\cdot\Theta(\Psi).
\]
In both cases $\Theta f$ is a function of the lower and mirror variables;
for $f$ on $\mathcal{Y}_+\times\mathcal{Y}_0$ it is defined whenever $(A_-,z_-)\in\prod_{b\in\mathbb{T}^{(1)}_{-}}J_b\times\Z^{\mathbb{T}^{(2)}_{-}}$ and $(A_0,z_0)\in\mathcal{Y}_0$.

The map $(A,z)\mapsto(\theta^*A,\theta^*z)$ sends the fundamental domain into itself,
except on the null set of the points with $a_b=-\pi/q$ for some cut link $b$.
To see this, recall that $\theta^*$ commutes with $d$ (Definition~\ref{def:cochains}), so that it preserves closedness.
Since $\theta^*\theta^*$ is the identity, the orientation signs of $\theta$ at $b$ and at $\theta b$ coincide;
for $b\in\mathbb{T}^{(1)}_{+}$ this gives $(\theta^*A)(\theta b)=\pm A(b)$, with the orientation sign of $\theta$ at $b$,
and this number lies in $J_{\theta b}$ whenever $A(b)\in[0,2\pi/q)$.
On the upper links $(\theta^*A)(b)\in[0,2\pi/q)$ by the choice of $J_{\theta b}$, on the links in $\mathcal{M}$ the value is unchanged,
and on a cut link $b$ the value $(\theta^*A)(b)=-a_b$ lies in $[-\pi/q,\pi/q)$ unless $a_b=-\pi/q$.

\paragraph{Integration over the two halves.}
Let $f$ and $f'$ be Borel functions on $\mathcal{Y}_+\times\mathcal{Y}_0$ with values in $\Gr_{\mathbb{T}}^{+}$
whose coefficient functions are integrable over $\mathcal{Y}_+$ for every value of the mirror variables,
and let $\int\mathrm{d}(A_-,z_-)$ be the integral over $\prod_{b\in\mathbb{T}^{(1)}_{-}}J_b\times\Z^{\mathbb{T}^{(2)}_{-}}$
with respect to the product of Lebesgue and counting measure.
Then, for every fixed value $(A_0,z_0)\in\mathcal{Y}_0$ of the mirror variables,
\begin{equation}\label{eq:chamber-half-integrals}
\int\mathrm{d}(A_-,z_-)\int_{\mathcal{Y}_+}\mathrm{d}y_+\int_{\mathbb{T}}\Theta(f')\,f
=\overline{\int_{\mathcal{Y}_+}\mathrm{d}y_+\int_{\mathbb{X}_+}f'(y_+;A_0,z_0)}\cdot\int_{\mathcal{Y}_+}\mathrm{d}y_+\int_{\mathbb{X}_+}f(y_+;A_0,z_0),
\end{equation}
where on the left $\Theta(f')\,f$ is taken at the point $(A,z)$ with upper variables $y_+$, lower variables $(A_-,z_-)$ and mirror variables $(A_0,z_0)$;
the cut variables do not enter.
For the proof put $y'_+:=\left((\theta^*A)_+,(\theta^*z)_+\right)$.
Since $\Theta$ is antilinear, the formula for $\Theta f$ obtained above for functions with values in $\Gr_{\mathbb{T}}^{+}$ shows that $(\Theta f')(A,z)$ is the image under $\Theta$ of the element $f'(y'_+;A_0,z_0)$ of $\Gr_{\mathbb{T}}^{+}$.
Both sides of the factorisation \eqref{eq:chamber-factorisation} of Lemma~\ref{lem:product-rule} are linear in $\xi$ and antilinear in $\xi'$.
Every element of $\Gr_{\mathbb{T}}^{+}$, a linear combination of monomials in the generators at the sites of $\mathbb{X}_+$, is the sum of an even and an odd element of $\Gr_{\mathbb{T}}^{+}$;
as the factorisation holds for elements of definite parities, it holds for all $\xi,\xi'\in\Gr_{\mathbb{T}}^{+}$.
Applied with $\xi'=f'(y'_+;A_0,z_0)$ and $\xi=f(y_+;A_0,z_0)$, it gives at every such point
\[
\int_{\mathbb{T}}\Theta(f')\,f=\overline{\int_{\mathbb{X}_+}f'(y'_+;A_0,z_0)}\cdot\int_{\mathbb{X}_+}f(y_+;A_0,z_0).
\]
The Berezin integral over $\mathbb{X}_+$ is a fixed multiple of one coefficient.
So the right side is the product of an integrable function of the lower variables and an integrable function of the upper variables,
and Fubini's theorem, together with the measure preserving bijection $(A_-,z_-)\mapsto y'_+$, gives \eqref{eq:chamber-half-integrals}.

\paragraph{Properties of the half weights.}
The faces of an upper cell are upper or lie in $\mathcal{M}$, and the links of a plaquette in $\mathcal{M}$ lie in $\mathcal{M}$ (Lemma~\ref{lem:cell-kinds}\eqref{it:kinds-upper-faces}).
Consequently $F_\sigma$ depends only on $A_+$, $A_0$ and $z_\sigma$ if $\sigma$ is upper, and only on $A_0$ and $z_\sigma$ if $\sigma$ lies in $\mathcal{M}$;
for an upper cube $Q$, $(dz)(Q)$ depends only on $z_+$ and $z_0$.
So $\mathcal{W}_+$ is a function of the upper and mirror variables, $\mathcal{W}_0$ a function of the mirror variables,
and $0\le\mathcal{W}_+\le1$, $0\le\mathcal{W}_0\le1$.
As $\theta^*$ commutes with $d$, one has $F(\theta^*A,\theta^*z)=\theta^*F$,
that is, $F_\sigma(\theta^*A,\theta^*z)=\pm F_{\theta\sigma}(A,z)$ and $(d\theta^*z)(Q)=\pm(dz)(\theta Q)$ with the orientation signs of $\theta$.
Since $\mathcal{W}_+$ is real and $\theta$ maps the upper plaquettes and cubes bijectively onto the lower ones,
\[
\Theta(\mathcal{W}_+)=\prod_{\sigma\in\mathbb{T}^{(2)}_{-}}\exp\left(-\frac{F_\sigma^2}{2}\right)\cdot\prod_{Q\in\mathbb{T}^{(3)}_{-}}1_{\{(dz)(Q)=0\}}.
\]
Every plaquette of $\Z^3/\Gamma$ is of exactly one of the four kinds,
every cube is upper, lower or cut (a cube does not lie in a plane),
and $z$ is closed if and only if $(dz)(Q)=0$ for every cube $Q$.
Therefore, for all $A\in C^1(\Gamma)$ and $z\in C^2(\Gamma;\Z)$,
\begin{equation}\label{eq:chamber-weight-factors}
1_{\{dz=0\}}\cdot\exp\left(-\frac{\|F\|^2}{2}\right)
=\mathcal{W}_0\cdot\mathcal{W}_+\cdot\Theta(\mathcal{W}_+)\cdot\prod_{\sigma\in\mathbb{T}^{(2)}_{\times}}\exp\left(-\frac{F_\sigma^2}{2}\right)\cdot\prod_{Q\in\mathbb{T}^{(3)}_{\times}}1_{\{(dz)(Q)=0\}}.
\end{equation}

Let
\[
C_q=\sup_{y\in\R}\sum_{j\in\Z}\exp\left(-\frac{1}{2}\left(y+\frac{2\pi j}{q}\right)^2\right)
\]
be the constant of Lemma~\ref{lem:bounded-integrands}.
For every value of the mirror variables,
\[
\int_{\mathcal{Y}_+}\mathrm{d}y_+\,\mathcal{W}_+\le(2\pi/q)^{\#|\mathbb{T}^{(1)}_{+}|}\cdot C_q^{\#|\mathbb{T}^{(2)}_{+}|},
\]
and
\[
\sum_{z_0\in\Z^{\mathbb{T}^{(2)}_{0}}}\ \int_{[0,2\pi/q)^{\mathbb{T}^{(1)}_{0}}}\mathrm{d}A_0\,\mathcal{W}_0\le(2\pi/q)^{\#|\mathbb{T}^{(1)}_{0}|}\cdot C_q^{\#|\mathbb{T}^{(2)}_{0}|}.
\]
For the first inequality, replace the indicators of the upper cubes by $1$.
At fixed $A_+$ and fixed mirror variables the sum over $z_+$ becomes the product, over $\sigma\in\mathbb{T}^{(2)}_{+}$, of the sums
\[
\sum_{z_\sigma\in\Z}\exp\left(-\frac{1}{2}\left((dA)_\sigma+\frac{2\pi z_\sigma}{q}\right)^2\right)\le C_q,
\]
and the integral over $A_+\in[0,2\pi/q)^{\mathbb{T}^{(1)}_{+}}$ contributes the factor $(2\pi/q)^{\#|\mathbb{T}^{(1)}_{+}|}$.
For the second, at fixed $A_0$ the sum over $z_0$ is the product, over $\sigma\in\mathbb{T}^{(2)}_{0}$, of sums of the same form, each at most $C_q$,
and the integral over $A_0\in[0,2\pi/q)^{\mathbb{T}^{(1)}_{0}}$ contributes the factor $(2\pi/q)^{\#|\mathbb{T}^{(1)}_{0}|}$.

\begin{lemma}[Poisson summation]\label{lem:poisson}
Let $q>0$. For all real $y$ and $\alpha$,
\[
\sum_{j\in\Z}\exp\left(-\frac{1}{2}\left(y+\frac{2\pi j}{q}\right)^2+ij\alpha\right)
=\frac{q}{\sqrt{2\pi}}\sum_{k\in\Z}\exp\left(-\frac{\beta_k^2}{2}-i\beta_k y\right),
\qquad
\beta_k:=\frac{q}{2\pi}\,(\alpha-2\pi k).
\]
\end{lemma}

\begin{proof}
The function
\[
f(s):=\exp\left(-\frac{1}{2}\left(y+\frac{2\pi s}{q}\right)^2+is\alpha\right),\qquad s\in\R,
\]
is a Schwartz function, and Poisson's summation formula gives $\sum_{j\in\Z}f(j)=\sum_{k\in\Z}\hat f(2\pi k)$,
where $\hat f(\xi):=\int_{\R}\mathrm{d}s\,f(s)\cdot\exp(-i\xi s)$.
In the substitution $t=y+2\pi s/q$ one has $\mathrm{d}s=(q/2\pi)\,\mathrm{d}t$ and $s(\alpha-2\pi k)=(t-y)\beta_k$; hence
\[
\hat f(2\pi k)=\frac{q}{2\pi}\int_{\R}\mathrm{d}t\,\exp\left(-\frac{t^2}{2}+i(t-y)\beta_k\right)
=\frac{q}{2\pi}\cdot\sqrt{2\pi}\cdot\exp\left(-\frac{\beta_k^2}{2}-i\beta_k y\right),
\]
by the Fourier transform of the Gaussian,
\[
\int_{\R}\mathrm{d}t\,\exp\left(-\frac{t^2}{2}+it\beta\right)=\sqrt{2\pi}\cdot\exp\left(-\frac{\beta^2}{2}\right)\qquad(\beta\in\R).
\]
\end{proof}

\begin{lemma}[Positive type of the photon crossing]\label{lem:positive-type}
Let $\mathbb{T}$, $c$ and the notation be as in Definition~\ref{def:half-weights},
and let $(d_+A)_\sigma$, $(d_\times A)_\sigma$ and $(d_+z)(Q)$, for cut plaquettes $\sigma$ and cut cubes $Q$, be as in Lemma~\ref{lem:cut-coboundaries}.
For $\alpha=(\alpha_Q)\in[0,2\pi)^{\mathbb{T}^{(3)}_{\times}}$ and $k=(k_\sigma)\in\Z^{\mathbb{T}^{(2)}_{\times}}$ put
\[
\alpha'_\sigma:=\sum_{Q\in\mathbb{T}^{(3)}_{\times}}[Q:\sigma]\cdot\alpha_Q,
\qquad
\beta_\sigma:=\frac{q}{2\pi}\,(\alpha'_\sigma-2\pi k_\sigma)
\qquad\left(\sigma\in\mathbb{T}^{(2)}_{\times}\right),
\]
and let $u_{\alpha,k}$ be the function of the upper variables, of modulus one, given by
\[
u_{\alpha,k}(A_+,z_+):=\exp\left(i\sum_{Q\in\mathbb{T}^{(3)}_{\times}}\alpha_Q\cdot(d_+z)(Q)-i\sum_{\sigma\in\mathbb{T}^{(2)}_{\times}}\beta_\sigma\cdot(d_+A)_\sigma\right).
\]
Then, for every $A\in C^1(\Gamma)$ and all integer cochains $z_+$, $z_-$, $z_0$ on the upper plaquettes, on the lower plaquettes and on the plaquettes in $\mathcal{M}$,
\begin{multline}\label{eq:chamber-positive-type}
\sum_{z_\times\in\Z^{\mathbb{T}^{(2)}_{\times}}}\ \prod_{\sigma\in\mathbb{T}^{(2)}_{\times}}\exp\left(-\frac{F_\sigma^2}{2}\right)\cdot\prod_{Q\in\mathbb{T}^{(3)}_{\times}}1_{\{(dz)(Q)=0\}}
\\
=\int_{[0,2\pi)^{\mathbb{T}^{(3)}_{\times}}}\prod_{Q\in\mathbb{T}^{(3)}_{\times}}\frac{\mathrm{d}\alpha_Q}{2\pi}\ \sum_{k\in\Z^{\mathbb{T}^{(2)}_{\times}}}\ \prod_{\sigma\in\mathbb{T}^{(2)}_{\times}}\frac{q}{\sqrt{2\pi}}\exp\left(-\frac{\beta_\sigma^2}{2}\right)
\\
\cdot u_{\alpha,k}\cdot\Theta(u_{\alpha,k})\cdot\exp\left(iq\sum_{b\in\mathbb{T}^{(1)}_{\times}}a_b\sum_{\sigma\in\mathbb{T}^{(2)}_{\times}}[\sigma:b]\cdot k_\sigma\right),
\end{multline}
where on the left $z=(z_+,z_-,z_0,z_\times)$ and $F=dA+(2\pi/q)z$, and on the right $u_{\alpha,k}$ and $\Theta(u_{\alpha,k})$ are taken at $(A,z)$.
Moreover, with the constant $C_q$ of Lemma~\ref{lem:bounded-integrands},
\[
\int_{[0,2\pi)^{\mathbb{T}^{(3)}_{\times}}}\prod_{Q\in\mathbb{T}^{(3)}_{\times}}\frac{\mathrm{d}\alpha_Q}{2\pi}\ \sum_{k\in\Z^{\mathbb{T}^{(2)}_{\times}}}\ \prod_{\sigma\in\mathbb{T}^{(2)}_{\times}}\frac{q}{\sqrt{2\pi}}\exp\left(-\frac{\beta_\sigma^2}{2}\right)\le C_q^{\#|\mathbb{T}^{(2)}_{\times}|}.
\]
\end{lemma}

\begin{proof}
Every $(dz)(Q)$ is an integer, so that
\[
1_{\{(dz)(Q)=0\}}=\int_0^{2\pi}\frac{\mathrm{d}\alpha_Q}{2\pi}\,\exp\left(i\alpha_Q\cdot(dz)(Q)\right)\qquad\left(Q\in\mathbb{T}^{(3)}_{\times}\right).
\]
Lemma~\ref{lem:cut-coboundaries} gives, with all sums over $Q$ taken over the cut cubes,
\[
\sum_{Q}\alpha_Q\cdot(dz)(Q)=\sum_{Q}\alpha_Q\cdot(d_+z)(Q)-\sum_{Q}\alpha_Q\cdot(d_+\theta^*z)(Q)+\sum_{\sigma\in\mathbb{T}^{(2)}_{\times}}\alpha'_\sigma\cdot z_\sigma;
\]
the last sum is $\sum_Q\alpha_Q\cdot(d_\times z)(Q)$ rearranged, and the first two sums do not involve $z_\times$.
Once these integrals are inserted, the term of $z_\times$ on the left of \eqref{eq:chamber-positive-type} becomes an integral over $\alpha$;
for every $\alpha$ its integrand has the modulus $\prod_{\sigma}\exp\left(-\frac{F_\sigma^2}{2}\right)$, the other factors being of modulus one, and
\[
\sum_{z_\times}\ \prod_{\sigma\in\mathbb{T}^{(2)}_{\times}}\exp\left(-\frac{F_\sigma^2}{2}\right)
=\prod_{\sigma\in\mathbb{T}^{(2)}_{\times}}\ \sum_{j\in\Z}\exp\left(-\frac{1}{2}\left((dA)_\sigma+\frac{2\pi j}{q}\right)^2\right)\le C_q^{\#|\mathbb{T}^{(2)}_{\times}|}.
\]
The series of the integrands thus converges absolutely, with a majorant that does not depend on $\alpha$.
As the measure $\prod_Q\mathrm{d}\alpha_Q/(2\pi)$ on $[0,2\pi)^{\mathbb{T}^{(3)}_{\times}}$ is finite, of total mass one,
Fubini's theorem allows the sum over $z_\times$ to be interchanged with the integral over $\alpha$.
After the interchange the sum over $z_\times$ splits into a product over the cut plaquettes.
With $y_\sigma:=(dA)_\sigma$, Lemma~\ref{lem:poisson}, applied with $y=y_\sigma$ and $\alpha=\alpha'_\sigma$, gives for each cut plaquette $\sigma$
\[
\sum_{j\in\Z}\exp\left(-\frac{1}{2}\left(y_\sigma+\frac{2\pi j}{q}\right)^2+ij\alpha'_\sigma\right)
=\frac{q}{\sqrt{2\pi}}\sum_{k_\sigma\in\Z}\exp\left(-\frac{\beta_\sigma^2}{2}\right)\cdot\exp\left(-i\beta_\sigma y_\sigma\right),
\]
and the product over $\sigma$ of these absolutely convergent series is the sum, over $k\in\Z^{\mathbb{T}^{(2)}_{\times}}$, of the products of their terms.
Lemma~\ref{lem:cut-coboundaries} also gives $y_\sigma=(d_+A)_\sigma-(d_+\theta^*A)_\sigma+(d_\times A)_\sigma$, whence
\[
\exp\left(-i\beta_\sigma y_\sigma\right)=\exp\left(-i\beta_\sigma\cdot(d_+A)_\sigma\right)\cdot\exp\left(i\beta_\sigma\cdot(d_+\theta^*A)_\sigma\right)\cdot\exp\left(-i\beta_\sigma\cdot(d_\times A)_\sigma\right).
\]
The numbers $\alpha_Q$ and $\beta_\sigma$ are real, and by Definition~\ref{def:half-weights}
\[
\begin{aligned}
\Theta\left(\exp\left(-i\beta_\sigma\cdot(d_+A)_\sigma\right)\right)&=\exp\left(i\beta_\sigma\cdot(d_+\theta^*A)_\sigma\right),\\
\Theta\left(\exp\left(i\alpha_Q\cdot(d_+z)(Q)\right)\right)&=\exp\left(-i\alpha_Q\cdot(d_+\theta^*z)(Q)\right);
\end{aligned}
\]
on scalar functions $\Theta$ is multiplicative.
The factors that contain $d_+A$, $d_+\theta^*A$, $d_+z$ and $d_+\theta^*z$ therefore combine to $u_{\alpha,k}\cdot\Theta(u_{\alpha,k})$,
and it remains to compute $\prod_\sigma\exp\left(-i\beta_\sigma\cdot(d_\times A)_\sigma\right)$.
Here $\sum_\sigma\beta_\sigma\cdot(d_\times A)_\sigma=\sum_{b\in\mathbb{T}^{(1)}_{\times}}a_b\sum_{\sigma\in\mathbb{T}^{(2)}_{\times}}[\sigma:b]\cdot\beta_\sigma$,
and for every cut link $b$
\[
\sum_{\sigma\in\mathbb{T}^{(2)}_{\times}}[\sigma:b]\cdot\alpha'_\sigma
=\sum_{Q\in\mathbb{T}^{(3)}_{\times}}\alpha_Q\sum_{\sigma\in\mathbb{T}^{(2)}_{\times}}[Q:\sigma]\cdot[\sigma:b]
=\sum_{Q\in\mathbb{T}^{(3)}_{\times}}\alpha_Q\sum_{\sigma}[Q:\sigma]\cdot[\sigma:b]=0.
\]
In the middle equality the restriction of the inner sum to cut plaquettes is immaterial:
$[\sigma:b]=0$ unless $b$ is a link of $\sigma$, and a plaquette one of whose links is cut is itself cut (Lemma~\ref{lem:cell-kinds}\eqref{it:kinds-cut-cofaces}).
The sum over all plaquettes vanishes because $d\circ d=0$, applied to the one-cochain that equals one at $b$ and zero at all other links.
As $\beta_\sigma=\frac{q}{2\pi}\alpha'_\sigma-qk_\sigma$, it follows that $\sum_\sigma[\sigma:b]\cdot\beta_\sigma=-q\sum_\sigma[\sigma:b]\cdot k_\sigma$,
and the remaining factor is $\exp\left(iq\sum_{b}a_b\sum_{\sigma}[\sigma:b]\cdot k_\sigma\right)$.
This proves \eqref{eq:chamber-positive-type}.

For the bound, Lemma~\ref{lem:poisson} with $y=0$, read from right to left, gives for every cut plaquette $\sigma$
\[
\frac{q}{\sqrt{2\pi}}\sum_{k_\sigma\in\Z}\exp\left(-\frac{\beta_\sigma^2}{2}\right)
=\sum_{j\in\Z}\exp\left(-\frac{1}{2}\left(\frac{2\pi j}{q}\right)^2+ij\alpha'_\sigma\right)
\le\sum_{j\in\Z}\exp\left(-\frac{1}{2}\left(\frac{2\pi j}{q}\right)^2\right)\le C_q.
\]
The sum over $k$ of the products over $\sigma$ is the product of these sums,
and the measure $\prod_Q\mathrm{d}\alpha_Q/(2\pi)$ on $[0,2\pi)^{\mathbb{T}^{(3)}_{\times}}$ has total mass one.
\end{proof}

\paragraph{Charges and Gauss' law.}
The charge of a mode $i$ at $\Pi$ or at $\Pi'$ (Definition~\ref{def:modes}) is the number
\[
\mathfrak{q}(i):=
\begin{cases}
+1&\text{if $i$ belongs to the group $1$ or $2$},\\
-1&\text{if $i$ belongs to the group $3$ or $4$}.
\end{cases}
\]
At a link plane $(x_1,j+\frac12)$, where Definition~\ref{def:modes} assigns to a mode of a bond over $y\in\Z^2$ the charge $\pm\delta_y$,
the number $\mathfrak{q}(i)$ is the non-zero value of this function.
A bond $\{x,x'\}$ lies over a link $b$ of $\Z^3/\Gamma$ if it lies over a representative of $b$,
that is, if $\pi(x)$ and $\pi(x')$ are the two end points of a link of $\Z^3$ in the class $b$.
A link of $\Z^3$ is determined by its end points;
so a bond lies over at most one link of $\Z^3$, and over at most one link of $\Z^3/\Gamma$.
As the kind of a link of $\Z^3/\Gamma$ is the kind of each of its representatives (Lemma~\ref{lem:cell-kinds}),
a bond lies over a link of $\mathbb{T}^{(1)}_{\times}$ if and only if it lies over a cut link of $\Z^3$.
For finite sets $I\subset\idx_{\mathbb{T}}(\Pi)$ and $I'\subset\idx_{\mathbb{T}}(\Pi')$ and for $b\in\mathbb{T}^{(1)}_{\times}$ let
\[
\mathfrak{q}(I\cup I')(b):=\sum_{i}\mathfrak{q}(i),
\]
the sum being over the modes $i\in I\cup I'$ whose bond lies over $b$.
An element $k\in\Z^{\mathbb{T}^{(2)}_{\times}}$ is extended by zero to all plaquettes of $\Z^3/\Gamma$;
the adjoint $d^*$ of $d$ then gives $(d^*k)(b)=\sum_{\sigma}[\sigma:b]\cdot k_\sigma$.
The triple $(k,I,I')$ satisfies Gauss' law if
\begin{equation}\label{eq:chamber-gauss}
d^{*}k+\mathfrak{q}(I\cup I')=0\quad\text{on }\mathbb{T}^{(1)}_{\times},
\end{equation}
that is, if $\sum_{\sigma\in\mathbb{T}^{(2)}_{\times}}[\sigma:b]\cdot k_\sigma+\mathfrak{q}(I\cup I')(b)=0$ for every $b\in\mathbb{T}^{(1)}_{\times}$.
If no link is cut, the condition is empty.

\begin{lemma}[Gauss' law]\label{lem:gauss-law}
Let $\mathbb{T}$, $c$ and the notation be as in Definition~\ref{def:half-weights}.
\begin{enumerate}
\item[(a)] Let $b$ be a cut link of $\Z^3$, of direction $\mu$, and let $H$ be the plane of $\mathcal{M}$ that meets its relative interior
(there is only one, by Lemma~\ref{lem:cell-kinds}\eqref{it:kinds-one-plane}).
Then $\ell=x_\mu$, the plane $H$ is a link plane through the midpoint of $b$, and $b$ is oriented along $n_\ell=e_\mu$.
A bond crossing a link plane of $\mathcal{M}$ lies over a cut link of $\Z^3$.
\item[(b)] Let $b\in\mathbb{T}^{(1)}_{\times}$, and let $i$ be a mode at $\Pi$ or at $\Pi'$ whose bond lies over $b$.
Then $u_i=\exp\left(iq\cdot\mathfrak{q}(i)\cdot a_b\right)$, where $\mathfrak{q}(i)\in\{+1,-1\}$ is the charge of $i$.
For finite sets $I\subset\idx_{\mathbb{T}}(\Pi)$ and $I'\subset\idx_{\mathbb{T}}(\Pi')$,
\[
u_I\cdot u_{I'}=\exp\left(iq\sum_{b\in\mathbb{T}^{(1)}_{\times}}\mathfrak{q}(I\cup I')(b)\cdot a_b\right).
\]
\item[(c)] For $k\in\Z^{\mathbb{T}^{(2)}_{\times}}$ and $I$, $I'$ as in \textup{(b)}, with $\mathrm{d}a$ the Lebesgue measure,
\begin{multline}\label{eq:chamber-cut-links}
\int_{\prod_{b\in\mathbb{T}^{(1)}_{\times}}J_b}\mathrm{d}a\,\exp\left(iq\sum_{b\in\mathbb{T}^{(1)}_{\times}}a_b\sum_{\sigma\in\mathbb{T}^{(2)}_{\times}}[\sigma:b]\cdot k_\sigma\right)\cdot u_I\cdot u_{I'}
\\
=(2\pi/q)^{\#|\mathbb{T}^{(1)}_{\times}|}\cdot1_{\{d^{*}k+\mathfrak{q}(I\cup I')=0\ \text{on}\ \mathbb{T}^{(1)}_{\times}\}}.
\end{multline}
\end{enumerate}
\end{lemma}

\begin{proof}
(a) The plane $H$ is the translate of $\Pi$ or of $\Pi'$ by an integer vector (Definition~\ref{def:split-tori}\eqref{it:split-slabs}),
so its reflection $\theta_H$ is the conjugate of $\theta$ or $\theta'$ by an integer translation;
it maps $\Z^3$ onto itself (Lemma~\ref{lem:chambers}\eqref{it:chambers-group}), hence links onto links, and it fixes every point of $H$.
A point of $H$ in the relative interior of $b$ is fixed by $\theta_H$, so it lies in the relative interior of $\theta_H b$ as well.
Distinct links have disjoint relative interiors; therefore $\theta_H b=b$,
and as $b$ does not lie in $H$, the reflection $\theta_H$ exchanges the two end points of $b$.
For every $x\in\R^3$ the difference $\theta_H x-x$ is a multiple of $n_\ell$, while at an end point of $b$ it equals $\pm e_\mu$.
Among the nine families, $x_\mu$ is the only one whose normal is a multiple of $e_\mu$; hence $\ell=x_\mu$ and $n_\ell=e_\mu$.
The midpoint of $b$ is fixed by $\theta_H$, so it lies on $H$;
since the end points of $b$ are integer points, $H=\{x_\mu=t+1/2\}$ for an integer $t$, which is a link plane.
Finally, every link of direction $\mu$ is oriented along $e_\mu=n_\ell$.

Now let a bond cross a link plane $H=\{x_\mu=t+1/2\}$ of $\mathcal{M}$.
By Proposition~\ref{prop:geometry}\eqref{it:geometry-mirrors} the bond is $\{x,\theta_H x\}$, where $x$ denotes its end below $H$,
and by Proposition~\ref{prop:geometry}\eqref{it:geometry-projection} it is of type $3$ and lies over the link $b$ of $\Z^3$ whose end points are $\pi(x)$ and $\pi(\theta_H x)=\theta_H\pi(x)$.
These two integer points are exchanged by $\theta_H$ and do not lie on $H$, so they lie on different sides of $H$:
the relative interior of $b$ meets $H$, $b$ does not lie in $H$, and $b$ is cut.
Moreover, the coordinates of $x-\pi(x)$ have the moduli $v_1$, $v_2$, $v_3$, all less than $1/2$ (Definition~\ref{def:fermion-sites});
as $x$ lies below $H$, the $\mu$-th coordinate of $\pi(x)$ is at most $t$.
Hence $\pi(x)$ is the end point of $b$ below $H$, $\theta_H\pi(x)=\pi(x)+e_\mu$, and $b$ is the link $(\pi(x),\mu)$.

(b) Let $H$ be the plane, $\Pi$ or $\Pi'$, that the bond of $i$ crosses, and let $x$ be the lower end of this bond.
By Definition~\ref{def:modes}, at $\Pi'$ in the form of Lemma~\ref{lem:second-mirror},
$u_i=U_A(x,\theta_H x)$ for the groups $1$ and $2$, and $u_i=\overline{U_A(x,\theta_H x)}$ for the groups $3$ and $4$.
The bond lies over a representative $b'$ of $b$;
thus $b'$ is a cut link of $\Z^3$, and its end points are $\pi(x)$ and $\pi(\theta_H x)=\theta_H\pi(x)$ (Proposition~\ref{prop:geometry}\eqref{it:geometry-projection}).
Being distinct and exchanged by $\theta_H$, the two end points have their midpoint on $H$;
so $H$ meets the relative interior of $b'$, and (a), applied to $b'$, shows that $H$ is a link plane.
If $\mu$ denotes the direction of $b'$, the second paragraph of the proof of (a) now gives $b'=(\pi(x),\mu)$.
The periodic extension of $A$ takes the value $A(b)=a_b$ at the representative $b'$ of $b$; hence
\[
U_A(x,\theta_H x)=\exp\left(iqA(\pi(x),\pi(x)+e_\mu)\right)=\exp(iqa_b).
\]
As $a_b$ is real, $\overline{U_A(x,\theta_H x)}=\exp(-iqa_b)$,
and in both cases $u_i=\exp\left(iq\cdot\mathfrak{q}(i)\cdot a_b\right)$ by the definition of the charge.
For the product formula, let $i\in I\cup I'$ be a mode whose bond lies over no link of $\mathbb{T}^{(1)}_{\times}$, hence over no cut link of $\Z^3$.
By (a), the mirror at which $i$ is a mode is then not a link plane;
so the bond of $i$ is of type $0$, $1$ or $2$, its two ends have the same nearest photon site (Proposition~\ref{prop:geometry}\eqref{it:geometry-projection}),
and $u_i=1$, since $A$ vanishes on pairs of equal points (Definition~\ref{def:cochains}).
Grouping the remaining modes of $I\cup I'$ according to the link of $\mathbb{T}^{(1)}_{\times}$ over which their bonds lie, which is unique, gives
$u_I\cdot u_{I'}=\prod_{b\in\mathbb{T}^{(1)}_{\times}}\exp\left(iq\cdot\mathfrak{q}(I\cup I')(b)\cdot a_b\right)$, which is the formula.

(c) By (b), the integrand on the left of \eqref{eq:chamber-cut-links} equals
\[
\prod_{b\in\mathbb{T}^{(1)}_{\times}}\exp\left(iq\cdot\left(d^*k+\mathfrak{q}(I\cup I')\right)(b)\cdot a_b\right),
\]
and every number $\left(d^*k+\mathfrak{q}(I\cup I')\right)(b)$ is an integer.
The integral is the product, over the cut links $b$, of the integrals of these factors over $a_b\in J_b=[-\pi/q,\pi/q)$.
The length $2\pi/q$ of $J_b$ is a period of each factor;
so the integral of the factor at $b$ equals $2\pi/q$ if $\left(d^*k+\mathfrak{q}(I\cup I')\right)(b)=0$, and $0$ otherwise.
The product is $(2\pi/q)^{\#|\mathbb{T}^{(1)}_{\times}|}$ if $(k,I,I')$ satisfies \eqref{eq:chamber-gauss}, and $0$ if it does not.
\end{proof}

\subsubsection{The reflected form}\label{sec:reflected-form}

Throughout, $\ell\in\mathcal{L}$, the torus $\mathbb{T}=(\Gamma,\chi)$ is split for $\ell$ with height period $p$, and $c\in\Adm(\ell)$;
the mirrors $\Pi=(\ell,c)$ and $\Pi'=(\ell,c')$, $c'=c+\frac p2$, the reflection $\theta=\theta_{\ell,c}$,
the open slab $U_+=\{c<\ell<c'\}$, the set $\mathcal{M}$ and the lattice $\Gamma_\parallel=\Gamma\cap\{\ell=0\}$ are those of Definition~\ref{def:split-tori},
the kinds of cells are those of Lemma~\ref{lem:cell-kinds},
and $\Theta:=\Theta_\Pi=\Theta_{\ell,c}$.

The values $F_\sigma$ of the curvature at the plaquettes $\sigma\in\mathbb{T}^{(2)}_{+}\cup\mathbb{T}^{(2)}_{0}$
and the link phases $\exp\left(iqA(b)\right)$, $b\in\mathbb{T}^{(1)}_{+}\cup\mathbb{T}^{(1)}_{0}$, are functions on $\Conf(\Gamma)$.
The links of an upper plaquette are upper or lie in $\mathcal{M}$,
and the links of a plaquette that lies in $\mathcal{M}$ lie in $\mathcal{M}$ (Lemma~\ref{lem:cell-kinds}\eqref{it:kinds-upper-faces}).
In the coordinates $A=(A_+,A_-,A_0,a)$ and $z=(z_+,z_-,z_0,z_\times)$ of the fundamental domain of Definition~\ref{def:half-weights},
a Borel function of these values and phases therefore depends only on the upper variables $y_+=(A_+,z_+)\in\mathcal{Y}_+$
and on the mirror variables $(A_0,z_0)\in\mathcal{Y}_0$,
and it is a Borel function on $\mathcal{Y}_+\times\mathcal{Y}_0$.

\begin{definition}[Integrands of the upper half]\label{def:upper-integrands}
Let $\ell\in\mathcal{L}$, let $\mathbb{T}=(\Gamma,\chi)$ be a torus that is split for $\ell$, and let $c\in\Adm(\ell)$,
with the notation of Definitions~\ref{def:split-tori} and~\ref{def:half-weights};
the cells of $\Z^3/\Gamma$ carry the kinds of Lemma~\ref{lem:cell-kinds}, taken with respect to $\mathbb{T}$, $\ell$ and $c$.
The class $\mathcal{B}^{+}=\mathcal{B}^{+}(\mathbb{T},\ell,c)$ is the linear span of the products $\Phi\,\xi$
in which $\xi$ lies in the Grassmann algebra $\Gr_{\mathbb{T}}^{+}=\Gr(\mathbb{X}_+)$ of the upper half (Lemma~\ref{lem:second-mirror})
and $\Phi$ is a bounded Borel function of the values $F_\sigma$, $\sigma\in\mathbb{T}^{(2)}_{+}\cup\mathbb{T}^{(2)}_{0}$,
and of the link phases $\exp\left(iqA(b)\right)$, $b\in\mathbb{T}^{(1)}_{+}\cup\mathbb{T}^{(1)}_{0}$.
The class $\mathcal{G}^{+}=\mathcal{G}^{+}(\mathbb{T},\ell,c)$ is the linear span of the products $\Phi\,\xi$
in which $\xi\in\Gr_{\mathbb{T}}^{+}$ and $\Phi$ is a Borel function of the same values and phases
for which there are constants $C<\infty$ and $\varepsilon\in\left(0,\frac12\right]$ such that
\[
\mathcal{W}_+\cdot|\Phi|\le C\cdot\prod_{\sigma\in\mathbb{T}^{(2)}_{+}}\exp\left(-\varepsilon F_\sigma^2\right)
\qquad\text{on }\mathcal{Y}_+\times\mathcal{Y}_0 ,
\]
where $\mathcal{W}_+$ is the upper half weight
and $\Phi$ is regarded, through these values and phases, as a function of the upper and mirror variables (Lemma~\ref{lem:cell-kinds}\eqref{it:kinds-upper-faces}).
\end{definition}

The elements of $\mathcal{B}^{+}$ and of $\mathcal{G}^{+}$ are functions on $\Conf(\Gamma)$ with values in $\Gr_{\mathbb{T}}^{+}$,
and they are read on $\mathcal{Y}_+\times\mathcal{Y}_0$ as above.
At the points where $\mathcal{W}_+\neq0$, that is, where $(dz)(Q)=0$ for all $Q\in\mathbb{T}^{(3)}_{+}$,
this reading does not depend on the way in which an element is written as a sum of products $\Phi\,\xi$:
each such point $(A_+,z_+;A_0,z_0)$ has the values $F_\sigma$, $\sigma\in\mathbb{T}^{(2)}_{+}\cup\mathbb{T}^{(2)}_{0}$, and the link phases at $\mathbb{T}^{(1)}_{+}\cup\mathbb{T}^{(1)}_{0}$ of a configuration.
To see this, let $z\in C^2(\Gamma;\Z)$ have the components $z_+$, $z_0$ and $z_\times:=0$, and the component $z_-$ for which $(\theta^*z)_+=z_+$,
which exists because $\theta$ maps the upper plaquettes bijectively onto the lower ones (Lemma~\ref{lem:cell-kinds}\eqref{it:kinds-reflection}).
By the same item $\theta^*z=z$ on the plaquettes in $\mathcal{M}$;
so $\theta^*z$ and $z$ agree on the faces of each upper cube $Q$ (Lemma~\ref{lem:cell-kinds}\eqref{it:kinds-upper-faces}),
and, as $\theta^*$ commutes with $d$, one has $(dz)(\theta Q)=\pm(d\theta^*z)(Q)=\pm(dz)(Q)=0$;
the cubes $\theta Q$, $Q\in\mathbb{T}^{(3)}_{+}$, are the lower cubes.
At a cut cube $Q$, Lemma~\ref{lem:cut-coboundaries} gives $(dz)(Q)=(d_+z)(Q)-(d_+\theta^*z)(Q)+(d_\times z)(Q)=0$, since $(\theta^*z)_+=z_+$ and $z_\times=0$.
Every cube being upper, lower or cut, $z$ is closed,
and for every $A\in C^1(\Gamma)$ with the components $A_+$ and $A_0$ the class of $(A,z)$ in $\Conf(\Gamma)$ is such a configuration.

One has $\mathcal{B}^{+}\subset\mathcal{G}^{+}$:
the indicators in $\mathcal{W}_+$ do not exceed $1$, so that
$\mathcal{W}_+\le\prod_{\sigma\in\mathbb{T}^{(2)}_{+}}\exp\left(-\frac12F_\sigma^2\right)$,
and a bounded $\Phi$ satisfies the inequality of the definition with $\varepsilon=\frac12$.
At a point of $\mathcal{Y}_+\times\mathcal{Y}_0$ with $(dz)(Q)\neq0$ for some $Q\in\mathbb{T}^{(3)}_{+}$ the weight $\mathcal{W}_+$ vanishes, and the inequality holds for every $\Phi$;
at a point with $(dz)(Q)=0$ for all $Q\in\mathbb{T}^{(3)}_{+}$ it reads
\[
|\Phi|\le C\cdot\exp\left(\left(\frac12-\varepsilon\right)\sum_{\sigma\in\mathbb{T}^{(2)}_{+}}F_\sigma^2\right).
\]
Thus the functions $\Phi$ of the class $\mathcal{G}^{+}$ have Gaussian growth, with a constant below $\frac12$, in the values at the upper plaquettes,
and their bound does not involve the values at the plaquettes that lie in $\mathcal{M}$.

Let $\Phi$ be a bounded Borel function of the values $F_\sigma$ and of the link phases $\exp\left(iqA(b)\right)$
of plaquettes $\sigma$ and links $b$ contained in $U_+$,
and let $\xi$ lie in the subalgebra of $\Gr_{\mathbb{T}}$ generated by the fields at the sites in $U_+$.
A cell contained in $U_+$ has its relative interior in $U_+$, so it is upper;
and the fields at the sites in $U_+$ generate $\Gr_{\mathbb{T}}^{+}$ (Lemma~\ref{lem:second-mirror}).
Hence $\Phi\,\xi\in\mathcal{B}^{+}$, and every finite sum of such products belongs to $\mathcal{B}^{+}$.

Let $B$ be a box contained in $U_+$ that embeds in $\Gamma_\parallel$ (Definition~\ref{def:integrands}).
By Lemma~\ref{lem:borel-lifts}, every element of $\mathcal{G}(B)$ lies in $\mathcal{G}^{+}$;
the estimate is made as for~\eqref{eq:chamber-slab-growth} in the proof of Theorem~\ref{thm:adapted-positivity}\eqref{it:adapted-growth} below.
An element of $\mathcal{G}(B)$ with bounded coefficients,
for instance an element of $\mathcal{P}$ that reads in $B$ (Definition~\ref{def:polynomials}), lies in $\mathcal{B}^{+}$.

\begin{proof}[Proof of Theorem~\ref{thm:adapted-positivity}]
Let $\rho$, $M$, $w_i$, $q$, $\ell$, $\mathbb{T}$, $p$ and $c$ be as in the theorem.
The torus $\mathbb{T}$, being adapted to $\ell$, is split for $\ell$,
and $\chi(g_0)=-1$ (Definition~\ref{def:split-tori}).
Since $-1\le\rho\le1$, the crossing coefficients $\omega_i$ of the modes at $\Pi$ and at $\Pi'$ are non-negative (Proposition~\ref{prop:crossing}).

The measure space, the measure and the linear map of item~\eqref{it:adapted-representation} are the following.
Let
\begin{equation}\label{eq:chamber-measure-space}
\mathcal{Z}:=\mathcal{Y}_0\times[0,2\pi)^{\mathbb{T}^{(3)}_{\times}}\times\Z^{\mathbb{T}^{(2)}_{\times}}
\times\left\{I:\ I\subset\idx_{\mathbb{T}}(\Pi)\right\}\times\left\{I':\ I'\subset\idx_{\mathbb{T}}(\Pi')\right\},
\end{equation}
with the points $s=(A_0,z_0,\alpha,k,I,I')$;
here $\idx_{\mathbb{T}}(\Pi)$ and $\idx_{\mathbb{T}}(\Pi')$ are the finite ordered sets of the modes at the two mirrors (Lemma~\ref{lem:second-mirror}).
On $\mathcal{Z}$ let $\mathfrak{m}$ be the measure
\begin{equation}\label{eq:chamber-measure}
\begin{aligned}
\mathrm{d}\mathfrak{m}(s):={}&\mathcal{W}_0(A_0,z_0)\,\mathrm{d}A_0
\cdot\prod_{Q\in\mathbb{T}^{(3)}_{\times}}\frac{\mathrm{d}\alpha_Q}{2\pi}
\cdot\prod_{\sigma\in\mathbb{T}^{(2)}_{\times}}\frac{q}{\sqrt{2\pi}}\exp\left(-\frac{\beta_\sigma^2}{2}\right)\\
&\cdot\omega_I\cdot\omega_{I'}\cdot\left(\frac{2\pi}{q}\right)^{\#|\mathbb{T}^{(1)}_{\times}|}
\cdot1_{\{d^{*}k+\mathfrak{q}(I\cup I')=0\ \text{on}\ \mathbb{T}^{(1)}_{\times}\}} .
\end{aligned}
\end{equation}
Here $\mathrm{d}A_0$ is Lebesgue measure on $[0,2\pi/q)^{\mathbb{T}^{(1)}_{0}}$ and $\mathrm{d}\alpha_Q$ is Lebesgue measure on $[0,2\pi)$;
the variables $z_0\in\Z^{\mathbb{T}^{(2)}_{0}}$, $k$, $I$ and $I'$ carry the counting measure;
$\beta_\sigma=\beta_\sigma(\alpha,k)$ is the real number attached to $\alpha$ and $k$ in Lemma~\ref{lem:positive-type};
$\omega_I$ and $\omega_{I'}$ are the products of the crossing coefficients over $I$ and over $I'$;
and the last factor is the indicator of Gauss' law~\eqref{eq:chamber-gauss}.
For $G\in\mathcal{B}^{+}$ and $s\in\mathcal{Z}$ let
\begin{equation}\label{eq:chamber-transform}
\hat G(s):=\int_{\mathcal{Y}_+}\mathrm{d}y_+\,\mathcal{W}_+(y_+,A_0,z_0)\cdot u_{\alpha,k}(y_+)
\int_{\mathbb{X}_+}G(y_+,A_0,z_0)\,\exp\left(-S_+(y_+,A_0,z_0)\right)\zeta_I\,\check\zeta_{I'} .
\end{equation}
In this formula $\int_{\mathbb{X}_+}$ is the Berezin integral of $\Gr(\mathbb{X}_+)=\Gr_{\mathbb{T}}^{+}$;
$S_+$ is the even element of $\Gr_{\mathbb{T}}^{+}$ in~\eqref{eq:chamber-action-split};
$\zeta_I$ and $\check\zeta_{I'}$ are the ordered products of the modes $\zeta_i$, $i\in I$, at $\Pi$
and of the elements $\check\zeta_{i'}$, $i'\in I'$, at $\Pi'$ (Lemma~\ref{lem:second-mirror});
and $u_{\alpha,k}$ is the unimodular function on $\mathcal{Y}_+$ of Lemma~\ref{lem:positive-type}.
By Lemma~\ref{lem:second-mirror}, the coefficients of $S_+$ are polynomials in the phases of links contained in the closed slab $\{c\le\ell\le c'\}$,
and each of these links lies in $\Pi$ or in $\Pi'$, or has its relative interior in $U_+$;
so these links are upper or lie in $\mathcal{M}$.
Hence $S_+$ and $\exp(-S_+)$ are functions on $\mathcal{Y}_+\times\mathcal{Y}_0$ with values in $\Gr_{\mathbb{T}}^{+}$,
and the coefficients of $\exp(-S_+)$, which are polynomials in link phases, are bounded.

\eqref{it:adapted-invariance}
Let $G''\in\mathcal{B}(\mathbb{T})$, the algebra of the finite sums $\sum_j\Phi_j\cdot\xi_j$
of bounded Borel functions $\Phi_j$ on $\Conf(\Gamma)$ times elements $\xi_j\in\Gr_{\mathbb{T}}$ (Lemma~\ref{lem:bounded-integrands}).
The reflection $\theta$ maps the torus onto itself, $\theta\mathbb{T}=\mathbb{T}$ (Definition~\ref{def:split-tori}\eqref{it:split-reflections}),
and Proposition~\ref{prop:symmetries}\eqref{it:functional-covariance}, applied to $G''$ and to the mirror $\Pi$, gives
\[
Z_{\mathbb{T}}[\Theta_{\ell,c}G'']=Z_{\theta\mathbb{T}}[\Theta_{\Pi}G'']=\overline{Z_{\mathbb{T}}[G'']} .
\]

\eqref{it:adapted-representation}
Let first $G$ and $G'$ be elements of $\mathcal{B}^{+}$ of definite parities $|G|,|G'|\in\{0,1\}$.
For $I\subset\idx_{\mathbb{T}}(\Pi)$ and $I'\subset\idx_{\mathbb{T}}(\Pi')$ put
\[
\xi_{I,I'}:=\exp(-S_+)\,\zeta_I\,\check\zeta_{I'}\in\Gr_{\mathbb{T}}^{+},
\]
an element whose coefficients are bounded functions on $\mathcal{Y}_+\times\mathcal{Y}_0$.
\begin{enumerate}
\item[(1)]\itemmark{1}\label{it:integrand-step}
\emph{The integrand.}
Since $\chi(g_0)=-1$, the factors $\left(-\chi(g_0)\right)^{\#|I'|}$ in~\eqref{eq:chamber-weight-expansion} are equal to $1$;
even elements commute with all elements of $\Gr_{\mathbb{T}}$, and so
\[
\exp(-S_{\mathbb{T}})=\sum_{I,I'}\omega_I\cdot\omega_{I'}\cdot u_I\cdot u_{I'}\cdot
\Theta\left(\exp(-S_+)\right)\exp(-S_+)\,\Theta(\zeta_I)\,\zeta_I\,\Theta(\check\zeta_{I'})\,\check\zeta_{I'},
\]
the sum being taken over all subsets $I\subset\idx_{\mathbb{T}}(\Pi)$ and $I'\subset\idx_{\mathbb{T}}(\Pi')$.
Multiply from the left by $\Theta(G')\,G$ and apply the product rule~\eqref{eq:chamber-product-rule} three times,
with $\exp(-S_+)$, $\zeta_I$ and $\check\zeta_{I'}$ in turn as the third factor.
The first application produces no sign, $\exp(-S_+)$ being even.
In the second and in the third application the first two factors are $G'\exp(-S_+)$ and $G\exp(-S_+)$,
respectively $G'\exp(-S_+)\,\zeta_I$ and $G\exp(-S_+)\,\zeta_I$;
their parities add up to $|G|+|G'|$ modulo $2$,
and the signs produced are $(-1)^{(|G|+|G'|)\#|I|}$ and $(-1)^{(|G|+|G'|)\#|I'|}$.
Hence
\[
\Theta(G')\,G\,\exp(-S_{\mathbb{T}})=\sum_{I,I'}\omega_I\cdot\omega_{I'}\cdot u_I\cdot u_{I'}\cdot(-1)^{(|G|+|G'|)(\#|I|+\#|I'|)}\cdot
\Theta\left(G'\,\xi_{I,I'}\right)\left(G\,\xi_{I,I'}\right).
\]
In this identity $\Theta$ is applied to functions on $\mathcal{Y}_+\times\mathcal{Y}_0$, namely to the coefficients of $G'$ and of $S_+$,
as prescribed in Definition~\ref{def:half-weights} and in the same manner as in the term $\Theta(S_+)$ of~\eqref{eq:chamber-action-split}.
The product rule~\eqref{eq:chamber-product-rule} follows from commutations of elements of definite parities alone,
so it remains valid when the third factor under $\Theta$ and the third factor outside $\Theta$ are different elements of the same parity.
Hence the identity holds at every point $(A,z)$ with $A\in C^1(\Gamma)$ and $z\in C^2(\Gamma;\Z)$,
the factors under $\Theta$ being taken at $(\theta^*A,\theta^*z)$ and the other factors at $(A,z)$.
If $|G|=|G'|$, the sign is $1$.
If $|G|\neq|G'|$, then in each term one of the elements $G'\,\xi_{I,I'}$, $G\,\xi_{I,I'}$ is odd;
an odd element has no component along the top monomial of $\Gr(\mathbb{X}_+)$, which is a product of $8\,\#|\mathbb{X}_+|$ generators,
so its Berezin integral over $\mathbb{X}_+$ vanishes, and by~\eqref{eq:chamber-factorisation} the term is annihilated by $\int_{\mathbb{T}}$.
As $\exp(-S_{\mathbb{T}})$ is even, it follows in both cases that
\[
\int_{\mathbb{T}}\exp(-S_{\mathbb{T}})\,\Theta(G')\,G
=\sum_{I,I'}\omega_I\cdot\omega_{I'}\cdot u_I\cdot u_{I'}\int_{\mathbb{T}}\Theta\left(G'\,\xi_{I,I'}\right)\left(G\,\xi_{I,I'}\right).
\]

\item[(2)]\itemmark{2}\label{it:weight-step}
\emph{The weight.}
By Lemma~\ref{lem:reflection-classes}, $\Theta(G')$ is a finite sum of products of functions $\overline{\Phi'\circ\theta^*}$ with elements of $\Gr_{\mathbb{T}}$,
where $\Phi'$ ranges over the bounded factors of $G'$ and $\theta^*$ maps $\Conf(\Gamma)$ onto itself, because $\theta\mathbb{T}=\mathbb{T}$.
So $\Theta(G')\,G\in\mathcal{B}(\mathbb{T})$,
and the integral of Definition~\ref{def:functional} that defines $Z_{\mathbb{T}}[\Theta(G')\,G]$ converges absolutely (Lemma~\ref{lem:bounded-integrands}).
It may be computed on the fundamental domain $\prod_bJ_b\times\{z:\ dz=0\}$ of Definition~\ref{def:half-weights}, as an iterated integral.
Writing the condition $dz=0$ as the product of the indicators $1_{\{(dz)(Q)=0\}}$ over all cubes $Q$ of $\Z^3/\Gamma$ and using~\eqref{eq:chamber-weight-factors}, one obtains
\[
\begin{aligned}
Z_{\mathbb{T}}[\Theta(G')\,G]={}&\sum_{z_0}\int\mathrm{d}A_0\,\mathcal{W}_0\int_{\mathcal{Y}_+}\mathrm{d}y_+\int\mathrm{d}(A_-,z_-)\,\mathcal{W}_+\cdot\Theta(\mathcal{W}_+)\\
&\cdot\int\mathrm{d}a\sum_{z_\times}\prod_{\sigma\in\mathbb{T}^{(2)}_{\times}}\exp\left(-\frac{F_\sigma^2}{2}\right)\prod_{Q\in\mathbb{T}^{(3)}_{\times}}1_{\{(dz)(Q)=0\}}
\int_{\mathbb{T}}\exp(-S_{\mathbb{T}})\,\Theta(G')\,G .
\end{aligned}
\]
In this formula $z_0$ and $z_\times$ range over $\Z^{\mathbb{T}^{(2)}_{0}}$ and $\Z^{\mathbb{T}^{(2)}_{\times}}$,
$\mathrm{d}A_0$ is Lebesgue measure on $[0,2\pi/q)^{\mathbb{T}^{(1)}_{0}}$,
$\mathrm{d}(A_-,z_-)$ is the product of Lebesgue measure on $\prod_{b\in\mathbb{T}^{(1)}_{-}}J_b$ and of counting measure on $\Z^{\mathbb{T}^{(2)}_{-}}$,
and $\mathrm{d}a$ is Lebesgue measure on $\prod_{b\in\mathbb{T}^{(1)}_{\times}}J_b$.
The functions $G$ and $G'$ and the coefficients of $S_+$ are functions on $\mathcal{Y}_+\times\mathcal{Y}_0$,
and $\theta$ maps upper cells to lower cells and every cell that lies in $\mathcal{M}$ to itself (Lemma~\ref{lem:cell-kinds}\eqref{it:kinds-reflection});
so the Berezin integrals $\int_{\mathbb{T}}\Theta(G'\,\xi_{I,I'})(G\,\xi_{I,I'})$ of Step~\ref{it:integrand-step} depend neither on $z_\times$ nor on $a$.
The product $u_I\cdot u_{I'}$ depends on $a$ only (Lemma~\ref{lem:gauss-law}).
Insert the result of Step~\ref{it:integrand-step}, and replace the sum over $z_\times$ by the right side of~\eqref{eq:chamber-positive-type}:
\begin{equation}\label{eq:chamber-iterated}
\begin{aligned}
Z_{\mathbb{T}}[\Theta(G')\,G]={}&\sum_{z_0}\int\mathrm{d}A_0\,\mathcal{W}_0\int_{\mathcal{Y}_+}\mathrm{d}y_+\int\mathrm{d}(A_-,z_-)\,\mathcal{W}_+\cdot\Theta(\mathcal{W}_+)
\int\mathrm{d}a\int\prod_{Q\in\mathbb{T}^{(3)}_{\times}}\frac{\mathrm{d}\alpha_Q}{2\pi}\\
&\cdot\sum_{k}\prod_{\sigma\in\mathbb{T}^{(2)}_{\times}}\frac{q}{\sqrt{2\pi}}\exp\left(-\frac{\beta_\sigma^2}{2}\right)
\cdot u_{\alpha,k}\cdot\Theta(u_{\alpha,k})\cdot\exp\left(iq\sum_{b\in\mathbb{T}^{(1)}_{\times}}a_b\sum_{\sigma\in\mathbb{T}^{(2)}_{\times}}[\sigma:b]\cdot k_\sigma\right)\\
&\cdot\sum_{I,I'}\omega_I\cdot\omega_{I'}\cdot u_I\cdot u_{I'}\int_{\mathbb{T}}\Theta\left(G'\,\xi_{I,I'}\right)\left(G\,\xi_{I,I'}\right),
\end{aligned}
\end{equation}
with $\alpha\in[0,2\pi)^{\mathbb{T}^{(3)}_{\times}}$ and $k\in\Z^{\mathbb{T}^{(2)}_{\times}}$.
The integral over $a$ is to be computed with~\eqref{eq:chamber-cut-links},
after an exchange of the order of the integrations which Step~\ref{it:fubini-step} justifies.

\item[(3)]\itemmark{3}\label{it:fubini-step}
\emph{Integrability; the order of the integrations.}
By~\eqref{eq:chamber-factorisation}, $\int_{\mathbb{T}}\Theta(G'\,\xi_{I,I'})(G\,\xi_{I,I'})$ is the product of the complex conjugate of the Berezin integral over $\mathbb{X}_+$ of $G'\,\xi_{I,I'}$,
taken at $(\theta^*A,\theta^*z)$, and of the Berezin integral over $\mathbb{X}_+$ of $G\,\xi_{I,I'}$.
These integrals are coefficients of the top monomial of $\Gr(\mathbb{X}_+)$;
they are bounded, because the functions $\Phi$ of $G$ and of $G'$ and the coefficients of $\exp(-S_+)$ are bounded.
Let $C$ bound their moduli on $\mathcal{Y}_+\times\mathcal{Y}_0$ for the finitely many pairs $(I,I')$;
then the modulus of $\int_{\mathbb{T}}\Theta(G'\,\xi_{I,I'})(G\,\xi_{I,I'})$, a product of two such integrals, is at most $C^2$.
The functions $u_{\alpha,k}$, $\Theta(u_{\alpha,k})$, $u_I$, $u_{I'}$ and the exponential that contains $a$ in~\eqref{eq:chamber-iterated} have modulus $1$.
So the modulus of the integrand of~\eqref{eq:chamber-iterated} is at most the non-negative function
\[
C^2\cdot\omega_I\cdot\omega_{I'}\cdot\mathcal{W}_0\cdot\mathcal{W}_+\cdot\Theta(\mathcal{W}_+)\cdot\prod_{\sigma\in\mathbb{T}^{(2)}_{\times}}\frac{q}{\sqrt{2\pi}}\exp\left(-\frac{\beta_\sigma^2}{2}\right).
\]
Integrate this majorant successively over the following variables.
Over the upper variables the integral is at most $(2\pi/q)^{\#|\mathbb{T}^{(1)}_{+}|}\cdot C_q^{\#|\mathbb{T}^{(2)}_{+}|}$, uniformly in the mirror variables:
drop the constraints at the cubes, sum over each $z_\sigma$ first, and use the constant $C_q$ of Lemma~\ref{lem:bounded-integrands}.
Over the lower variables the same bound holds:
$\Theta(\mathcal{W}_+)$ is the function $\mathcal{W}_+$ evaluated at $\left((\theta^*A)_+,(\theta^*z)_+;A_0,z_0\right)$,
and $(A_-,z_-)\mapsto\left((\theta^*A)_+,(\theta^*z)_+\right)$ is a measure preserving bijection onto $\mathcal{Y}_+$ (Definition~\ref{def:half-weights}).
Over the mirror variables the integral of $\mathcal{W}_0$ is at most $(2\pi/q)^{\#|\mathbb{T}^{(1)}_{0}|}\cdot C_q^{\#|\mathbb{T}^{(2)}_{0}|}$, in the same way.
Over $a$ one obtains the factor $(2\pi/q)^{\#|\mathbb{T}^{(1)}_{\times}|}$,
over $(\alpha,k)$ at most $C_q^{\#|\mathbb{T}^{(2)}_{\times}|}$ (Lemma~\ref{lem:positive-type}),
and over $(I,I')$
\[
\sum_{I,I'}\omega_I\cdot\omega_{I'}=\prod_{i\in\idx_{\mathbb{T}}(\Pi)}(1+\omega_i)\cdot\prod_{i'\in\idx_{\mathbb{T}}(\Pi')}(1+\omega_{i'}).
\]
The reflection $\theta$ exchanges the upper and the lower cells,
and $\Z^3/\Gamma$ has $3|\mathbb{T}|$ links and $3|\mathbb{T}|$ plaquettes, each of exactly one kind.
By Tonelli's theorem the integral of the majorant over all the variables is therefore at most
\[
C^2\cdot\left(\frac{2\pi}{q}\cdot C_q\right)^{3|\mathbb{T}|}\cdot\prod_{i\in\idx_{\mathbb{T}}(\Pi)}(1+\omega_i)\cdot\prod_{i'\in\idx_{\mathbb{T}}(\Pi')}(1+\omega_{i'})<\infty ,
\]
and by Fubini's theorem the integrations in~\eqref{eq:chamber-iterated} may be carried out in any order.

Integrate first over $a$.
Among the factors of the integrand only $\exp\left(iq\sum_{b}a_b\sum_{\sigma}[\sigma:b]\cdot k_\sigma\right)\cdot u_I\cdot u_{I'}$ depends on $a$,
and by~\eqref{eq:chamber-cut-links} its integral equals
$(2\pi/q)^{\#|\mathbb{T}^{(1)}_{\times}|}\cdot1_{\{d^{*}k+\mathfrak{q}(I\cup I')=0\ \text{on}\ \mathbb{T}^{(1)}_{\times}\}}$.
Fix now $s=(A_0,z_0,\alpha,k,I,I')\in\mathcal{Z}$.
The functions $\mathcal{W}_+$ and $u_{\alpha,k}$ are scalar, and $\mathcal{W}_+$ is real.
For the functions
\[
f:=\mathcal{W}_+\cdot u_{\alpha,k}\cdot G\,\xi_{I,I'},
\qquad
f':=\mathcal{W}_+\cdot u_{\alpha,k}\cdot G'\,\xi_{I,I'}
\]
on $\mathcal{Y}_+\times\mathcal{Y}_0$, which take values in $\Gr_{\mathbb{T}}^{+}$ and are integrable over $\mathcal{Y}_+$, one therefore has
\[
\mathcal{W}_+\cdot\Theta(\mathcal{W}_+)\cdot u_{\alpha,k}\cdot\Theta(u_{\alpha,k})\cdot\Theta\left(G'\,\xi_{I,I'}\right)\left(G\,\xi_{I,I'}\right)=\Theta(f')\,f .
\]
By~\eqref{eq:chamber-half-integrals} and the definition~\eqref{eq:chamber-transform},
the integral of $\int_{\mathbb{T}}\Theta(f')\,f$ over the lower and the upper variables is
\[
\overline{\int_{\mathcal{Y}_+}\mathrm{d}y_+\int_{\mathbb{X}_+}f'}\cdot\int_{\mathcal{Y}_+}\mathrm{d}y_+\int_{\mathbb{X}_+}f
=\overline{\widehat{G'}(s)}\cdot\hat G(s).
\]
The integrations that remain are those over $(A_0,z_0)$, $\alpha$, $k$ and $(I,I')$, with the weights collected in~\eqref{eq:chamber-measure}; hence
\[
Z_{\mathbb{T}}[\Theta(G')\,G]=\int_{\mathcal{Z}}\mathrm{d}\mathfrak{m}(s)\,\overline{\widehat{G'}(s)}\cdot\hat G(s).
\]

\item[(4)]\itemmark{4}\label{it:measure-step}
\emph{The measure and the map $G\mapsto\hat G$.}
The measure $\mathfrak{m}$ is positive:
$\omega_I\ge0$ and $\omega_{I'}\ge0$ because $|\rho|\le1$ (Proposition~\ref{prop:crossing}),
and the other factors in~\eqref{eq:chamber-measure} are non-negative.
It is finite: with the bounds of Step~\ref{it:fubini-step} for the mirror variables and for $(\alpha,k)$,
\[
\mathfrak{m}(\mathcal{Z})\le\left(\frac{2\pi}{q}\right)^{\#|\mathbb{T}^{(1)}_{0}|+\#|\mathbb{T}^{(1)}_{\times}|}
\cdot C_q^{\#|\mathbb{T}^{(2)}_{0}|+\#|\mathbb{T}^{(2)}_{\times}|}
\cdot\prod_{i\in\idx_{\mathbb{T}}(\Pi)}(1+\omega_i)\cdot\prod_{i'\in\idx_{\mathbb{T}}(\Pi')}(1+\omega_{i'}).
\]
For $G\in\mathcal{B}^{+}$ the function $\hat G$ is measurable on $\mathcal{Z}$ by Fubini's theorem.
Let $C$ be a bound for the moduli of the Berezin integrals $\int_{\mathbb{X}_+}G\,\xi_{I,I'}$ on $\mathcal{Y}_+\times\mathcal{Y}_0$, for all pairs $(I,I')$;
it exists, since $G$ and $\exp(-S_+)$ have bounded coefficients.
As $|u_{\alpha,k}|=1$, the bound of Step~\ref{it:fubini-step} for the upper variables then gives
\[
|\hat G(s)|\le C\cdot\left(\frac{2\pi}{q}\right)^{\#|\mathbb{T}^{(1)}_{+}|}\cdot C_q^{\#|\mathbb{T}^{(2)}_{+}|}\qquad(s\in\mathcal{Z}).
\]
The map $G\mapsto\hat G$ is linear.
Every element of $\mathcal{B}^{+}$ is the sum of an even and an odd element of $\mathcal{B}^{+}$,
and both sides of the identity obtained in Step~\ref{it:fubini-step} are linear in $G$ and antilinear in $G'$;
so the identity holds for all $G,G'\in\mathcal{B}^{+}$.
With the finite positive measure space $(\mathcal{Z},\mathfrak{m})$ of~\eqref{eq:chamber-measure-space} and~\eqref{eq:chamber-measure}
and with the linear map $G\mapsto\hat G$ of~\eqref{eq:chamber-transform} into the bounded measurable functions on $\mathcal{Z}$,
this is item~\eqref{it:adapted-representation}.
\end{enumerate}

\eqref{it:adapted-form}
By~\eqref{it:adapted-representation}, $Z_{\mathbb{T}}[\Theta_{\ell,c}(G')\,G]$ is the scalar product of $\widehat{G'}$ and $\hat G$ in $L^2(\mathcal{Z},\mathfrak{m})$,
antilinear in $G'$ and linear in $G$.
An exchange of $G$ and $G'$ replaces it by its complex conjugate, and
\[
Z_{\mathbb{T}}[\Theta_{\ell,c}(G)\,G]=\int_{\mathcal{Z}}\mathrm{d}\mathfrak{m}(s)\,|\hat G(s)|^2\ge0 .
\]
So $(G',G)\mapsto Z_{\mathbb{T}}[\Theta_{\ell,c}(G')\,G]$ is a positive semidefinite Hermitian form on $\mathcal{B}^{+}$.

\eqref{it:adapted-growth}
\emph{The class $\mathcal{A}_\Xi[c,c']$ lies in $\mathcal{G}^{+}$.}
Here $c'=c+p/2$, $U_+=\{c<\ell<c'\}$, and $\Xi=\Xi(\mathbb{T},\ell)=(\Gamma_\parallel,\chi_\parallel)$ is the cross-section of $\mathbb{T}$ for $\ell$.
Both classes are linear spans.
It therefore suffices to consider a product $G=\Phi_0G_1\cdots G_r$ of $\mathcal{A}(U_+)$ that embeds in $\Gamma_\parallel$.
In it, $\Phi_0$ is a bounded Borel function of finitely many $F_\sigma$ with $\sigma\subset U_+$,
and $G_j\in\mathcal{G}^{\mathrm{ev}}(B_j)$ for boxes $B_j\subset U_+$;
by Definition~\ref{def:integrands} there are constants $C_j<\infty$ and $\varepsilon_j>0$ such that every coefficient function of $G_j$ is bounded in modulus by
\[
C_j\cdot\exp\left(\left(\tfrac12-\varepsilon_j\right)\sum_{\sigma\in B_j^{(2)}}(d\hat A)_\sigma^2\right)
\]
at every real one-cochain $\hat A$ on $B_j^{(1)}$.

Each box $B_j$ is a compact subset of the open slab $U_+$.
A cell contained in $B_j$ thus has its relative interior in $U_+$, and it is upper (Lemma~\ref{lem:cell-kinds});
its faces are contained in $B_j$ as well, so they are upper too.
The same holds for the plaquettes read by $\Phi_0$.
A site $x\in\mathbb{X}(B_j)$ lies in its cube $\pi(x)+[-\tfrac12,\tfrac12]^3$, which is contained in $B_j$; hence $x\in U_+$.
Consequently $x=x'+g$ with $x'\in\mathbb{X}_+$ and $g\in\Gamma_\parallel$,
and in $\Gr_{\mathbb{T}}$ one has $\psi(x)=\chi_\parallel(g)\psi(x')$ and $\bar\psi(x)=\chi_\parallel(g)\bar\psi(x')$.
Read on $\mathbb{T}$, every monomial of $G_j$ is therefore, up to a sign, a monomial of $\Gr_{\mathbb{T}}^{+}$.

A point of $\mathcal{Y}_+\times\mathcal{Y}_0$, read at the classes in $\Z^3/\Gamma$ of the cells of $B_j$, which are upper,
gives a real one-cochain $A$ on $B_j^{(1)}$ and an integer two-cochain $z$ on $B_j^{(2)}$,
and $F=dA+\frac{2\pi}{q}\cdot z$ on $B_j^{(2)}$.
Suppose first that $B_j$ contains a site of $\Z^3$.
Then Lemma~\ref{lem:borel-lifts}\eqref{it:lift-coefficients} applies to $G_j$, to the box $B_j$ and to these two cochains.
At the Borel local lift $\hat A^{B_j}$ of that lemma, each coefficient function of $G_j$ is a Borel function of the values $F_\sigma$, $\sigma\in B_j^{(2)}$, and of the phases of the links in $B_j^{(1)}$,
hence a Borel function on $\mathcal{Y}_+\times\mathcal{Y}_0$;
at every point of $\mathcal{Y}_+\times\mathcal{Y}_0$ at which $(dz)(Q)=0$ for all cubes $Q$ of $B_j$, its modulus is at most
\[
C_j\cdot\exp\left(\left(\tfrac12-\varepsilon_j\right)\sum_{\sigma\in B_j^{(2)}}F_\sigma^2\right).
\]
If $B_j$ contains no site of $\Z^3$, then $B_j^{(1)}$, $B_j^{(2)}$ and $\mathbb{X}(B_j)$ are empty, $G_j$ is a constant of modulus at most $C_j$,
and both assertions hold trivially, the sum in the bound being empty.
The cubes of $B_j$ are upper, and $\mathcal{W}_+$ vanishes unless $(dz)(Q)=0$ for every upper cube $Q$;
this estimate is therefore valid wherever $\mathcal{W}_+\neq0$.
By Lemma~\ref{lem:reflected-embedding}, applied with $1$ as the reflected factor, the boxes $B_1,\dots,B_r$ embed disjointly in $\Gamma$,
and then a plaquette of $\Z^3/\Gamma$ has a representative in at most one of them (Lemma~\ref{lem:gaussian-growth}).
Since each $B_j$ embeds in $\Gamma$, the classes in $\Z^3/\Gamma$ of the plaquettes in $B_1^{(2)},\dots,B_r^{(2)}$ are pairwise different upper plaquettes.

Let $\varepsilon:=\min(\varepsilon_1,\dots,\varepsilon_r,\tfrac12)$,
and write $G$, read on $\mathbb{T}$, as $\sum_\Psi\Phi_\Psi\Psi$, the sum over the monomials $\Psi$ of $\Gr_{\mathbb{T}}^{+}$.
Up to signs, each $\Phi_\Psi$ is a finite sum of products of $\Phi_0$ with one coefficient function of each $G_j$;
it is a Borel function of the $F_\sigma$, $\sigma$ upper, and of the link phases of upper links.
One has $\mathcal{W}_+\le\prod_{\sigma\in\mathbb{T}^{(2)}_{+}}\exp\left(-\frac{F_\sigma^2}{2}\right)$, and $\tfrac12-\varepsilon_j\le\tfrac12-\varepsilon$ for every $j$, where $\tfrac12-\varepsilon\ge0$.
With the preceding facts this gives a constant $C$ such that
\begin{equation}\label{eq:chamber-slab-growth}
\mathcal{W}_+\cdot|\Phi_\Psi|\le C\cdot\prod_{\sigma\in\mathbb{T}^{(2)}_{+}}\exp\left(-\varepsilon F_\sigma^2\right)\qquad\text{on }\mathcal{Y}_+\times\mathcal{Y}_0
\end{equation}
for every $\Psi$.
Hence $G\in\mathcal{G}^{+}$ (Definition~\ref{def:upper-integrands}).

Let now $G,G'\in\mathcal{A}_\Xi[c,c']$.
The element $\Theta_{\ell,c}G'$ is formed with the action \eqref{eq:setting-reflection-coefficients} of the reflection on coefficient functions;
by Lemmas~\ref{lem:reflection-classes} and~\ref{lem:reflected-algebras} it is the image under the map $\Theta$ of Definition~\ref{def:half-weights} of $G'$, read as a function on $\mathcal{Y}_+\times\mathcal{Y}_0$.
By Lemma~\ref{lem:reflected-embedding}, $\Theta_{\ell,c}(G')G$ lies in $\mathcal{A}(\R^3)$ and embeds in $\Gamma$,
so that $Z_{\mathbb{T}}[\Theta_{\ell,c}(G')G]$ is the functional of Definition~\ref{def:functional}.
Once \eqref{it:adapted-representation} and \eqref{it:adapted-form} are established on $\mathcal{G}^{+}$,
the last assertion of \eqref{it:adapted-growth} follows by restriction to $\mathcal{A}_\Xi[c,c']$.

\emph{The representation on $\mathcal{G}^{+}$.}
Both sides of \eqref{it:adapted-representation} are sesquilinear in $(G',G)$, and $G\mapsto\hat G$ is linear.
It suffices to take $G=\Phi\xi$ and $G'=\Phi'\xi'$,
with $\xi,\xi'\in\Gr_{\mathbb{T}}^{+}$ of definite parities and with $\Phi$, $\Phi'$ subject to the bound of Definition~\ref{def:upper-integrands} with the constants $(C,\varepsilon)$ and $(C',\varepsilon')$.
Step~\ref{it:integrand-step} is an identity in $\Gr_{\mathbb{T}}$ at every point and remains valid,
and the identities \eqref{eq:chamber-weight-expansion}, \eqref{eq:chamber-weight-factors}, \eqref{eq:chamber-positive-type} and \eqref{eq:chamber-cut-links} do not involve $\Phi$ and $\Phi'$.
The coefficients of $\exp(-S_{\mathbb{T}})$ are polynomials in the link phases (Lemma~\ref{lem:bounded-integrands}), hence bounded.
Written with the indicator of $dz=0$, the weight of the integral of Definition~\ref{def:functional} for $Z_{\mathbb{T}}[\Theta(G')G]$ factorises as in \eqref{eq:chamber-weight-factors};
since the indicators at the cut cubes do not exceed $1$, each coefficient of the integrand has modulus at most a constant times
\[
\mathcal{W}_0\cdot\left(\mathcal{W}_+|\Phi|\right)\cdot\Theta\left(\mathcal{W}_+|\Phi'|\right)\cdot\prod_{\sigma\in\mathbb{T}^{(2)}_{\times}}\exp\left(-\frac{F_\sigma^2}{2}\right).
\]
Uniformly in the mirror variables,
\begin{equation}\label{eq:chamber-upper-integral}
\int_{\mathcal{Y}_+}\mathrm{d}y_+\prod_{\sigma\in\mathbb{T}^{(2)}_{+}}\exp\left(-\varepsilon F_\sigma^2\right)
\le(2\pi/q)^{\#|\mathbb{T}^{(1)}_{+}|}\cdot C_q(\varepsilon)^{\#|\mathbb{T}^{(2)}_{+}|}.
\end{equation}
Indeed, at fixed $A_+$ and $A_0$ the sum over $z_\sigma\in\Z$ of $\exp\left(-\varepsilon\left((dA)_\sigma+\frac{2\pi}{q}\cdot z_\sigma\right)^2\right)$ is at most $C_q(\varepsilon)$ for each upper plaquette $\sigma$ (Lemma~\ref{lem:gaussian-growth}),
and the integral over $A_+\in[0,2\pi/q)^{\mathbb{T}^{(1)}_{+}}$ contributes the first factor.
The map $(A_-,z_-)\mapsto\left((\theta^*A)_+,(\theta^*z)_+\right)$ of Definition~\ref{def:half-weights} is a measure preserving bijection of the range of the lower variables onto $\mathcal{Y}_+$,
and $\Theta\left(\mathcal{W}_+|\Phi'|\right)$ is the function $\mathcal{W}_+|\Phi'|$ taken at $\left((\theta^*A)_+,(\theta^*z)_+;A_0,z_0\right)$;
for the lower variables the same bound therefore holds, with $\varepsilon'$ in place of $\varepsilon$.
Summation over $z_\times$ comes first and gives at most $C_q^{\#|\mathbb{T}^{(2)}_{\times}|}$, uniformly in $A$;
the integrals over the upper and over the lower variables are then bounded as just shown,
the one over $(A_0,z_0)$ by the bound of Definition~\ref{def:half-weights} for $\mathcal{W}_0$,
and the cut-link variables range over intervals of length $2\pi/q$.
By Tonelli's theorem the integral of Definition~\ref{def:functional} converges absolutely, coefficient by coefficient,
and Step~\ref{it:weight-step} holds.

After the insertion of \eqref{eq:chamber-positive-type}, the majorant of Step~\ref{it:fubini-step} becomes a constant multiple of
\[
\omega_I\cdot\omega_{I'}\cdot\mathcal{W}_0\cdot\prod_{\sigma\in\mathbb{T}^{(2)}_{+}}\exp\left(-\varepsilon F_\sigma^2\right)\cdot\prod_{\sigma\in\mathbb{T}^{(2)}_{-}}\exp\left(-\varepsilon'F_\sigma^2\right)\cdot\prod_{\sigma\in\mathbb{T}^{(2)}_{\times}}\frac{q}{\sqrt{2\pi}}\,\exp\left(-\frac{\beta_\sigma^2}{2}\right).
\]
Its integral is estimated in the order of Step~\ref{it:fubini-step}:
the upper variables come first, then the lower variables, $(A_0,z_0)$, $a$ and $(\alpha,k)$, and the sum over $(I,I')$ comes last.
The result is at most a constant times
\[
(2\pi/q)^{3|\mathbb{T}|}\cdot C_q(\varepsilon)^{\#|\mathbb{T}^{(2)}_{+}|}\cdot C_q(\varepsilon')^{\#|\mathbb{T}^{(2)}_{-}|}\cdot C_q^{\#|\mathbb{T}^{(2)}_{0}|+\#|\mathbb{T}^{(2)}_{\times}|}\cdot\prod_{i\in\idx_{\mathbb{T}}(\Pi)}(1+\omega_i)\cdot\prod_{i'\in\idx_{\mathbb{T}}(\Pi')}(1+\omega_{i'})<\infty.
\]
Here $3|\mathbb{T}|$ is the number of links of $\Z^3/\Gamma$,
the factor $C_q^{\#|\mathbb{T}^{(2)}_{\times}|}$ comes from Lemma~\ref{lem:positive-type},
and the two products over the modes are $\sum_{I,I'}\omega_I\cdot\omega_{I'}$.
The integrations may therefore be carried out in any order.
The functions $f=\mathcal{W}_+\cdot u_{\alpha,k}\cdot G\xi_{I,I'}$ and $f'=\mathcal{W}_+\cdot u_{\alpha,k}\cdot G'\xi_{I,I'}$ of Step~\ref{it:fubini-step}
are integrable over $\mathcal{Y}_+$ by the bound of Definition~\ref{def:upper-integrands} and \eqref{eq:chamber-upper-integral};
hence \eqref{eq:chamber-half-integrals} applies, and the remainder of Step~\ref{it:fubini-step} is unchanged.
In Step~\ref{it:measure-step}, the bound \eqref{eq:chamber-upper-integral} shows that $|\hat G(s)|$ is bounded, uniformly in $s\in\mathcal{Z}$, by a constant multiple of
\[
(2\pi/q)^{\#|\mathbb{T}^{(1)}_{+}|}\cdot C_q(\varepsilon)^{\#|\mathbb{T}^{(2)}_{+}|},
\]
and $\mathfrak{m}$ is the measure \eqref{eq:chamber-measure}.
Thus \eqref{it:adapted-representation} holds on $\mathcal{G}^{+}$, with the same $(\mathcal{Z},\mathfrak{m})$, the same formula \eqref{eq:chamber-transform} and bounded $\hat G$;
and \eqref{it:adapted-form} on $\mathcal{G}^{+}$ follows from it as on $\mathcal{B}^{+}$.
\end{proof}

\begin{corollary}[Consequences]\label{cor:adapted-consequences}
Let $-1\le\rho\le1$, $M\in\R$, $w_0,\dots,w_3>0$ and $q>0$.
Let $\ell\in\mathcal{L}$, let $\mathbb{T}=(\Gamma,\chi)$ be a torus adapted to $\ell$ with height period $p$, and let $c\in\Adm(\ell)$.
Put $c':=c+p/2$, $\Pi:=(\ell,c)$, $\Pi':=(\ell,c')$, $\theta:=\theta_{\ell,c}$ and $\Theta:=\Theta_{\ell,c}=\Theta_\Pi$,
and let $\mathcal{B}^{+}=\mathcal{B}^{+}(\mathbb{T},\ell,c)$ and $\mathcal{G}^{+}=\mathcal{G}^{+}(\mathbb{T},\ell,c)$ be the classes of Definition~\ref{def:upper-integrands}.
\begin{enumerate}
\item\label{it:adapted-law}
\emph{Law.}
For every bounded Borel function $\Phi$ of the $F_\sigma$, $\sigma\in\mathbb{T}^{(2)}_{+}\cup\mathbb{T}^{(2)}_{0}$,
\[
\int\mathrm{d}\mu_{\mathbb{T}}(F)\,\overline{\Phi(\theta^{*}F)}\,\Phi(F)\ge0,
\]
where $\Phi(F)$ and $\Phi(\theta^{*}F)$ are evaluated on the $\Gamma$-periodic elements $F$ of $\Omega$, which carry $\mu_{\mathbb{T}}$,
by reading $F$ and $\theta^{*}F$ at arbitrary representatives in $\Z^3$ of the plaquettes $\sigma$.
\item\label{it:adapted-perturbed}
\emph{Perturbed weights.}
Let $\Phi_+\in\mathcal{B}^{+}$ have no Grassmann part,
that is, let $\Phi_+$ be a bounded Borel function, real or complex valued, of the $F_\sigma$, $\sigma\in\mathbb{T}^{(2)}_{+}\cup\mathbb{T}^{(2)}_{0}$, and of the link phases of the links in $\mathbb{T}^{(1)}_{+}\cup\mathbb{T}^{(1)}_{0}$.
Then $G\mapsto G\Phi_+$ maps $\mathcal{B}^{+}$ into $\mathcal{B}^{+}$ and $\mathcal{G}^{+}$ into $\mathcal{G}^{+}$,
the product $\Phi_+\Theta(\Phi_+)$ is a scalar function, and for every $G\in\mathcal{G}^{+}$
\[
Z_{\mathbb{T}}\left[\Phi_+\Theta(\Phi_+)\Theta(G)G\right]=Z_{\mathbb{T}}\left[\Theta(G\Phi_+)G\Phi_+\right]\ge0.
\]
\item\label{it:adapted-lower-side}
\emph{The lower side and the second mirror.}
Let $G\in\Theta(\mathcal{B}^{+})$ have definite parity $|G|$.
Then $(-1)^{|G|}Z_{\mathbb{T}}[\Theta(G)G]\ge0$:
the form $G\mapsto Z_{\mathbb{T}}[\Theta(G)G]$ is non-negative on the even and non-positive on the odd elements of $\Theta(\mathcal{B}^{+})$,
while the form of the reflection defined by \eqref{eq:setting-reflection-fields} with the normal $-\hat n_\ell$ in place of $\hat n_\ell$ is non-negative on both.
For $G\in\mathcal{B}^{+}$ of definite parity, $\Theta_{\Pi'}(G)=(-1)^{|G|}\Theta_\Pi(G)$ and
\[
Z_{\mathbb{T}}[\Theta_{\Pi'}(G)G]=(-1)^{|G|}Z_{\mathbb{T}}[\Theta_\Pi(G)G].
\]
\item\label{it:adapted-gauss-parity}
\emph{Gauss' law and parity.}
Let $(\mathcal{Z},\mathfrak{m})$ and $G\mapsto\hat G$ be as in \eqref{eq:chamber-measure-space}, \eqref{eq:chamber-measure} and \eqref{eq:chamber-transform}, with points $s=(A_0,z_0,\alpha,k,I,I')$.
If $\Pi'$ is a link plane, then the following holds for $\mathfrak{m}$-almost every $s\in\mathcal{Z}$.
The total charge through $\Pi'$ vanishes,
\[
\sum_b\mathfrak{q}(I\cup I')(b)=0,
\]
where $\mathfrak{q}(I\cup I')(b)$ is the charge at the link $b$ in Gauss' law \eqref{eq:chamber-gauss}
and the sum is taken over the links $b$ of $\Z^3/\Gamma$ that have a representative cut by $\Pi'$;
the set $I'$ contains as many modes of the groups $1$ and $2$ as modes of the groups $3$ and $4$ (Definition~\ref{def:modes});
and $\#|I'|$ is even.
If $\Pi$ is a link plane, the same holds with $\Pi$ and $I$ in place of $\Pi'$ and $I'$.
If $\Pi$ and $\Pi'$ are link planes, then $\hat G=0$ $\mathfrak{m}$-almost everywhere for every odd $G\in\mathcal{G}^{+}$,
and $Z_{\mathbb{T}}[\Theta(G')G]=0$ for every $G'\in\mathcal{G}^{+}$ and every odd $G\in\mathcal{G}^{+}$.
\end{enumerate}
\end{corollary}

\begin{proof}
\eqref{it:adapted-law}
By Theorem~\ref{thm:partition-positivity}, $0<Z_{\mathbb{T}}[1]<\infty$,
and $\mu_{\mathbb{T}}$ is the law of the periodic extension $\widetilde F$ under $\mathbb{E}_{\mathbb{T}}=Z_{\mathbb{T}}[\,\cdot\,]/Z_{\mathbb{T}}[1]$ (Definition~\ref{def:functional});
it is a probability measure on $\Omega$ by Corollary~\ref{cor:modulus-bound}.
It is carried by the $\Gamma$-periodic two-cochains, so $\Phi(F)$ may be evaluated at any representatives in $\Z^3$ of the plaquettes $\sigma$.
The reflection $\theta$ permutes the $\Gamma$-orbits (Definition~\ref{def:split-tori}\eqref{it:split-reflections}),
so $\theta^{*}\widetilde F$ is the periodic extension of $\theta^{*}F$;
moreover $F(\theta^{*}(A,z))=\theta^{*}F(A,z)$, because $\theta^{*}$ commutes with $d$.
By Lemma~\ref{lem:reflection-classes}, $(\Theta\Phi)(A,z)=\overline{\Phi(\theta^{*}F(A,z))}$.
Therefore
\[
\int\mathrm{d}\mu_{\mathbb{T}}(F)\,\overline{\Phi(\theta^{*}F)}\,\Phi(F)=\frac{Z_{\mathbb{T}}[\Theta(\Phi)\Phi]}{Z_{\mathbb{T}}[1]}.
\]
As a function on $\Conf(\Gamma)$, $\Phi$ is an element of $\mathcal{B}^{+}$ without Grassmann part,
and the numerator is non-negative by Theorem~\ref{thm:adapted-positivity}\eqref{it:adapted-form}.

\eqref{it:adapted-perturbed}
Let $G=\Phi\xi$ with $\xi\in\Gr_{\mathbb{T}}^{+}$.
Then $G\Phi_+=(\Phi\cdot\Phi_+)\xi$, and $\Phi\cdot\Phi_+$ is a Borel function of the same variables as $\Phi$.
It is bounded if $\Phi$ is;
and if $\Phi$ satisfies the bound of Definition~\ref{def:upper-integrands} with the constants $(C,\varepsilon)$, then $\Phi\cdot\Phi_+$ satisfies it with $(C\cdot\sup|\Phi_+|,\varepsilon)$.
By linearity, $G\mapsto G\Phi_+$ maps $\mathcal{B}^{+}$ into $\mathcal{B}^{+}$ and $\mathcal{G}^{+}$ into $\mathcal{G}^{+}$.
The element $\Theta(\Phi_+)$ is the scalar function $\overline{\Phi_+\circ\theta^{*}}$ (Definition~\ref{def:half-weights}),
and scalar functions commute with every element of $\Gr_{\mathbb{T}}$.
As $\Theta$ reverses the order of products,
\[
\Theta(G\Phi_+)G\Phi_+=\Theta(\Phi_+)\Theta(G)G\Phi_+=\Phi_+\Theta(\Phi_+)\Theta(G)G.
\]
For $G\in\mathcal{G}^{+}$ the product $G\Phi_+$ lies in $\mathcal{G}^{+}$,
and Theorem~\ref{thm:adapted-positivity}\eqref{it:adapted-growth} gives $Z_{\mathbb{T}}[\Theta(G\Phi_+)G\Phi_+]\ge0$.

\eqref{it:adapted-lower-side}
Write $G=\Theta(G_1)$ with $G_1\in\mathcal{B}^{+}$.
By Lemma~\ref{lem:sign-rules}, $\Theta^2=1$ on $\Gr_{\mathbb{T}}$;
on coefficient functions $\Theta^2$ is the identity as well, since $(\theta^{*})^2=\mathrm{id}$.
Hence $\Theta(G)=G_1$, and $G_1$ has the parity of $G$.
Exchanging two factors of definite parities multiplies their product by $-1$ raised to the product of the parities;
for the factors $G_1$ and $\Theta(G_1)$, both of parity $|G|$, the sign is $(-1)^{|G|^2}=(-1)^{|G|}$.
So
\[
\Theta(G)G=G_1\Theta(G_1)=(-1)^{|G|}\Theta(G_1)G_1,
\qquad
Z_{\mathbb{T}}[\Theta(G)G]=(-1)^{|G|}Z_{\mathbb{T}}[\Theta(G_1)G_1],
\]
and $Z_{\mathbb{T}}[\Theta(G_1)G_1]\ge0$ by Theorem~\ref{thm:adapted-positivity}\eqref{it:adapted-form}.
With $-\hat n_\ell$ in place of $\hat n_\ell$ in \eqref{eq:setting-reflection-fields}, the reflection changes by the factor $(-1)^{\deg}$ (Lemma~\ref{lem:sign-rules});
on $G$ this is the factor $(-1)^{|G|}$, and the form becomes $Z_{\mathbb{T}}[\Theta(G_1)G_1]\ge0$ for both parities.
Since $\mathbb{T}$ is adapted to $\ell$, $\chi(g_0)=-1$,
and \eqref{eq:chamber-second-mirror} reads $\Theta_{\Pi'}=(-1)^{\deg}\Theta_\Pi$ on $\Gr_{\mathbb{T}}$ (Lemma~\ref{lem:second-mirror}).
On coefficient functions $\Theta_{\Pi'}$ and $\Theta_\Pi$ coincide,
because $\theta'$ and $\theta$ induce the same involution of $\Conf(\Gamma)$ (Definition~\ref{def:split-tori}\eqref{it:split-reflections}).
So $\Theta_{\Pi'}(G)=(-1)^{|G|}\Theta_\Pi(G)$ for $G\in\mathcal{B}^{+}$ of definite parity, and the identity for $Z_{\mathbb{T}}$ follows.

\eqref{it:adapted-gauss-parity}
Let $\Pi'$ be a link plane, so that $\ell=x_\mu$ for some $\mu$ and $c'\in\Z+\tfrac12$.
A link or a plaquette of $\Z^3/\Gamma$ is said to be cut by $\Pi'$ if one of its representatives is cut by $\Pi'$.
A link or a plaquette of $\Z^3$ without the direction $\mu$ lies in a plane $\{x_\mu=j\}$ with $j\in\Z$, which does not meet $\Pi'$;
so $\Pi'$ cuts no such cell.
Consequently a plaquette $\sigma_0$ of $\Z^3$ cut by $\Pi'$ has the direction $\mu$ and a second direction $\nu$.
Its first corner $x$ satisfies $x_\mu=c'-\tfrac12$, since the relative interior of $\sigma_0$ meets $\Pi'$,
and $\Pi'$ cuts the two links $(x,\mu)$ and $(x+e_\nu,\mu)$ of $\sigma_0$.
In the formula for $dA$ of Definition~\ref{def:cochains} the values $A_\mu(x)$ and $A_\mu(x+e_\nu)$ enter $(dA)_{\sigma_0}$ with opposite signs, whether $\mu<\nu$ or $\nu<\mu$;
that is, $[\sigma_0:(x,\mu)]+[\sigma_0:(x+e_\nu,\mu)]=0$.
As the incidence numbers of $\Z^3/\Gamma$ are read on representatives,
$\sum_b[\sigma:b]=0$ for every plaquette $\sigma$ of $\Z^3/\Gamma$ cut by $\Pi'$, the sum being taken over the links $b$ of $\Z^3/\Gamma$ cut by $\Pi'$.
A plaquette that contains a link cut by $\Pi'$ is itself cut by $\Pi'$ (Lemma~\ref{lem:cell-kinds}\eqref{it:kinds-cut-cofaces}).
Let $s\in\mathcal{Z}$ be a point at which Gauss' law \eqref{eq:chamber-gauss} holds.
Summation of \eqref{eq:chamber-gauss} over the links $b\in\mathbb{T}^{(1)}_{\times}$ cut by $\Pi'$ gives,
with $b$ ranging over these links and $\sigma$ over the plaquettes cut by $\Pi'$,
\[
\sum_b\mathfrak{q}(I\cup I')(b)=-\sum_\sigma k_\sigma\sum_b[\sigma:b]=0.
\]
A bond meets exactly one mirror (Proposition~\ref{prop:geometry}\eqref{it:geometry-midpoint}),
and the bonds that lie over the links cut by $\Pi'$ are exactly the bonds that cross $\Pi'$
(Proposition~\ref{prop:geometry}\eqref{it:geometry-projection}, Lemma~\ref{lem:gauss-law}).
Hence no mode of $I$ lies over a link cut by $\Pi'$, while the bond of each mode of $I'$ lies over exactly one such link.
By its definition, $\mathfrak{q}(I\cup I')(b)$ is the sum of the charges of the modes in $I\cup I'$ whose bond lies over $b$,
the charge of a mode being $+1$ for the groups $1$ and $2$ and $-1$ for the groups $3$ and $4$.
The left side of the last display is therefore the number of modes of $I'$ in the groups $1$ and $2$ minus the number of modes of $I'$ in the groups $3$ and $4$.
These two numbers are thus equal, and $\#|I'|$, which is their sum, is even.
The density of $\mathfrak{m}$ in \eqref{eq:chamber-measure} contains the indicator of Gauss' law;
all three assertions therefore hold for $\mathfrak{m}$-almost every $s$.
The same argument, with $\Pi$ in place of $\Pi'$, applies to $I$ when $\Pi$ is a link plane.

Now let $G\in\mathcal{G}^{+}$ have definite parity $|G|$.
The element $G\exp(-S_+)\zeta_I\check\zeta_{I'}$ of $\Gr_{\mathbb{T}}^{+}$ has the parity of $|G|+\#|I|+\#|I'|$,
since $\exp(-S_+)$ is even and every $\zeta_i$ and every $\check\zeta_{i'}$ is odd;
the top monomial $\Psi_+$ of $\Gr_{\mathbb{T}}^{+}$ has the even degree $8\,\#|\mathbb{X}_+|$.
The Berezin integral of this element over $\mathbb{X}_+$ thus vanishes unless $|G|+\#|I|+\#|I'|$ is even, and by \eqref{eq:chamber-transform} so does $\hat G(s)$.
If $\Pi$ and $\Pi'$ are link planes and $G$ is odd, it follows that $\hat G=0$ $\mathfrak{m}$-almost everywhere,
and Theorem~\ref{thm:adapted-positivity}\eqref{it:adapted-representation}, \eqref{it:adapted-growth} gives
\[
Z_{\mathbb{T}}[\Theta(G')G]=\int_{\mathcal{Z}}\mathrm{d}\mathfrak{m}(s)\,\overline{\widehat{G'}(s)}\,\hat G(s)=0
\]
for every $G'\in\mathcal{G}^{+}$.
\end{proof}

\begin{remark}
The function $\Phi$ in Corollary~\ref{cor:adapted-consequences}\eqref{it:adapted-law} reads only plaquettes that are upper or lie in $\mathcal{M}$, that is, plaquettes of the closed upper half of the torus.
The class $\mathcal{O}^{+}(\ell,c)$ of Definition~\ref{def:reflection-positivity} is formed with the set $\Sigma^{+}(\ell,c)$ of the plaquettes of $\Z^3$ all of whose corners lie in $\{\ell>c\}$.
Every plaquette contained in the open slab $\{c'<\ell<c+p\}=U_{-}+g_0$ belongs to $\Sigma^{+}(\ell,c)$ and is lower;
a function in $\mathcal{O}^{+}(\ell,c)$ may therefore read plaquettes whose classes in $\Z^3/\Gamma$ are lower.
Each plaquette of $\Z^3/\Gamma$ that is upper or lies in $\mathcal{M}$, on the other hand, has a representative in $\Sigma^{+}(\ell,c)$,
since $\ell(g_0)=p>0$ and every plaquette of $\Z^3$ therefore has a translate by a multiple of $g_0\in\Gamma$ all of whose corners lie in $\{\ell>c\}$;
and on the $\Gamma$-periodic two-cochains, which carry $\mu_{\mathbb{T}}$, the values of $F$ and of $\theta^{*}F$ at a plaquette do not change when the plaquette is translated by an element of $\Gamma$
(Definition~\ref{def:split-tori}\eqref{it:split-reflections}).
For the law $\mu_{\mathbb{T}}$, Corollary~\ref{cor:adapted-consequences}\eqref{it:adapted-law} is therefore exactly the inequality of Definition~\ref{def:reflection-positivity}
for those $\Phi\in\mathcal{O}^{+}(\ell,c)$ that are functions of finitely many $F_\sigma$ with plaquettes $\sigma\in\Sigma^{+}(\ell,c)$ that are upper or lie in $\mathcal{M}$;
it does not cover the functions in $\mathcal{O}^{+}(\ell,c)$ that read a lower or a cut plaquette.
\end{remark}

\begin{remark}
The cases of Theorem~\ref{thm:adapted-positivity}\eqref{it:adapted-form} and of Corollary~\ref{cor:adapted-consequences}\eqref{it:adapted-lower-side} obey one rule for the normal $n$ in \eqref{eq:setting-reflection-fields}.
If $n$ points from the mirror into the half of the torus on which $G$ depends,
the reflected form at that mirror, $Z_{\mathbb{T}}[\Theta_\Pi(G)G]$ or $Z_{\mathbb{T}}[\Theta_{\Pi'}(G)G]$, is non-negative for even and for odd $G$;
if $n$ points out of that half, it is non-negative for even and non-positive for odd $G$.
At $\Pi$ the normal $\hat n_\ell$ points out of the lower half, and at $\Pi'$ out of the upper half;
this is the origin of the sign $(-1)^{|G|}$ in Corollary~\ref{cor:adapted-consequences}\eqref{it:adapted-lower-side}.
For even $G$ the orientation of the normal is immaterial.
\end{remark}

\section{Positivity of the partition function}\label{sec:positivity}

\subsection{The determinant as a Laurent polynomial}\label{sec:laurent-polynomial}

\begin{definition}[General tori]\label{def:general-tori}
Throughout this section $\rho$ is a real number, $w_0,\dots,w_3>0$, $q>0$,
and $v_0\in C_0$ is the base point of Definition~\ref{def:fermion-sites}.
The mass $M$ is real;
in this subsection and in Subsections~\ref{sec:torus-layers} and~\ref{sec:extreme-coefficient} it may be any complex number,
the operator $D_A$ being given by the formula of Definition~\ref{def:chamber-operator}.

A \emph{general torus} is a pair $\mathbb{T}=(\Gamma,\chi)$ of a subgroup $\Gamma$ of $\Z^3$ of rank three
and a character $\chi\colon\Gamma\to\{z\in\C:\ |z|=1\}$;
its number of photon sites is $|\mathbb{T}|:=[\Z^3:\Gamma]$.
The tori of Definition~\ref{def:tori} are the general tori with $\sys(\Gamma)\ge4$ and $\chi(\Gamma)\subset\{1,-1\}$.
In particular $C(n_1,n_2,n_3)$, $D(a,b,c)$, $D'(a,b,c)$
and the tori $g\,C(n_1,n_2,n_3)$, $g\,D(a,b,c)$, $g\,D'(a,b,c)$ with $g\in\Ocub$ are general tori.

Let $\mathbb{T}=(\Gamma,\chi)$ be a general torus.
The group $\Gamma$ acts on $\mathbb{X}$ by translations, and $\mathbb{X}_{\mathbb{T}}:=\mathbb{X}/\Gamma$ is the set of the orbits.
As $\mathbb{X}$ is the disjoint union of the $48$ sets $y+\Z^3$, $y\in\Ocub v_0$ (Lemma~\ref{lem:chambers}\eqref{it:chambers-transitive}),
each of which consists of $|\mathbb{T}|$ orbits,
one has $\#|\mathbb{X}_{\mathbb{T}}|=48|\mathbb{T}|$.
A \emph{spinor} of $\mathbb{T}$ is a map $\psi\colon\mathbb{X}\to\C^4$ with $\psi(x+g)=\chi(g)\cdot\psi(x)$ for all $x\in\mathbb{X}$ and $g\in\Gamma$.
For a real $\Gamma$-periodic one-cochain $A$ the operator $D_A$ maps spinors of $\mathbb{T}$ to spinors of $\mathbb{T}$,
because $x\mapsto x^{(i)}$ and $\pi$ commute with the translations by $\Gamma$, the hopping matrices are invariant under them, and $A$ is $\Gamma$-periodic;
$D^{\mathbb{T}}_A$ denotes the operator $D_A$ on the spinors of $\mathbb{T}$.
The configuration space $\Conf(\Gamma)$ with its measure and, for real $M$, the functional $Z_{\mathbb{T}}$ on Borel functions $\Phi\ge0$
are given by the formulas of Definitions~\ref{def:configurations} and~\ref{def:functional}.

\emph{Links and phases.}
Let $\mathbb{T}^{(1)}$ be the set of the $\Gamma$-orbits of links of $\Z^3$,
and for $\mu\in\{1,2,3\}$ let $\mathbb{T}^{(1)}_\mu\subset\mathbb{T}^{(1)}$ be the set of the orbits of the links $(v,\mu)$, $v\in\Z^3$.
Two links $(v,\mu)$ and $(v',\mu)$ lie in one orbit if and only if $v-v'\in\Gamma$;
hence $\#|\mathbb{T}^{(1)}_\mu|=|\mathbb{T}|$ and $\#|\mathbb{T}^{(1)}|=3|\mathbb{T}|$.
A real $\Gamma$-periodic one-cochain $A$ is a function $b\mapsto A(b)$ on $\mathbb{T}^{(1)}$,
and $A\mapsto\left(A(b)\right)_{b\in\mathbb{T}^{(1)}}$ identifies these cochains with $\R^{\mathbb{T}^{(1)}}$.
The \emph{phase} of the link $b\in\mathbb{T}^{(1)}$ is $u_b:=\exp\left(iqA(b)\right)$.
\end{definition}

For real $M$, Proposition~\ref{prop:accretivity} gives $\det D^{\mathbb{T}}_A\ge0$
for every general torus $\mathbb{T}=(\Gamma,\chi)$ and every real $\Gamma$-periodic one-cochain $A$.

Every site $x\in\mathbb{X}$ is $x=\pi(x)+y(x)$ with a unique $\pi(x)\in\Z^3$ and $y(x)\in\Ocub v_0$
(Definition~\ref{def:fermion-sites}, Lemma~\ref{lem:chambers}\eqref{it:chambers-transitive}).
The coordinates of $y(x)$ have the moduli $v_1>v_2>v_3>0$ in some order, and
\[
\pi(x+g)=\pi(x)+g,\qquad y(x+g)=y(x),\qquad g\in\Z^3 .
\]

\begin{lemma}\label{lem:incongruent-neighbours}
Let $x\in\mathbb{X}$.
The five vectors $y(x)$, $y\left(x^{(0)}\right)$, $y\left(x^{(1)}\right)$, $y\left(x^{(2)}\right)$, $y\left(x^{(3)}\right)$ are pairwise different.
Consequently, for every subgroup $\Gamma$ of $\Z^3$, the five sites $x,x^{(0)},\dots,x^{(3)}$ lie in five different classes modulo $\Gamma$.
\end{lemma}

\begin{proof}
Let $\nu$ be the coordinate in which $y(x)$ has the modulus $v_1$.
Lemma~\ref{lem:neighbours} describes the four neighbours of $x$ in terms of $y(x)$.
In $y\left(x^{(0)}\right)$ the coordinate of modulus $v_3$ has changed its sign, and no other coordinate has changed.
In $y\left(x^{(1)}\right)$ the moduli $v_2$ and $v_3$ have exchanged their coordinates,
in $y\left(x^{(2)}\right)$ the moduli $v_1$ and $v_2$,
and in both every sign has stayed in its coordinate.
In $y\left(x^{(3)}\right)$ the coordinate $\nu$ has changed its sign, and no other coordinate has changed:
if $y_\nu(x)=\pm v_1$, then $x^{(3)}-\pi(x)$ arises from $y(x)$ by replacing $\pm v_1$ with $\pm(1-v_1)$ in the coordinate $\nu$,
and $\pi\left(x^{(3)}\right)=\pi(x)\pm e_\nu$, so that $y_\nu\left(x^{(3)}\right)=\mp v_1$.

Consider first the vectors of moduli $(|y_1|,|y_2|,|y_3|)$.
For $y(x)$, $y\left(x^{(0)}\right)$ and $y\left(x^{(3)}\right)$ it is one and the same vector.
The vector of moduli of $y\left(x^{(2)}\right)$ has the entry $v_2$ in the coordinate $\nu$,
where those of $y(x)$ and of $y\left(x^{(1)}\right)$ have the entry $v_1$;
the vector of moduli of $y\left(x^{(1)}\right)$ has the entry $v_2$ in the coordinate in which that of $y(x)$ has the entry $v_3$.
Since $v_1>v_2>v_3$, it follows that $y\left(x^{(1)}\right)$ and $y\left(x^{(2)}\right)$ differ from each other and from the three other vectors.
The remaining vectors $y(x)$, $y\left(x^{(0)}\right)$ and $y\left(x^{(3)}\right)$ have the same moduli and pairwise different vectors of signs,
because $y\left(x^{(0)}\right)$ and $y\left(x^{(3)}\right)$ differ from $y(x)$ in the sign of exactly one coordinate each,
and these two coordinates are different.
This proves the first assertion.
Two sites whose difference lies in $\Gamma\subset\Z^3$ have the same image under $y$;
so no two of the sites $x,x^{(0)},\dots,x^{(3)}$ are congruent modulo $\Gamma$.
\end{proof}

\begin{definition}[The matrix of the operator]\label{def:operator-matrix}
Let $\rho$, $w_0,\dots,w_3$, $q$ and $v_0$ be as in Definition~\ref{def:general-tori}, and let $M\in\C$.
Let $\mathbb{T}=(\Gamma,\chi)$ be a general torus and $A$ a real $\Gamma$-periodic one-cochain.
Let $\mathbb{F}\subset\mathbb{X}$ be a set of representatives of the classes in $\mathbb{X}_{\mathbb{T}}$, one site in each class;
for $x'\in\mathbb{X}$ let $\widehat{x'}\in\mathbb{F}$ be the representative of the class of $x'$.
For $x\in\mathbb{F}$ and $i\in\{0,1,2,3\}$ let $g_i(x)\in\Gamma$ be the element with
\[
x^{(i)}=\widehat{x^{(i)}}+g_i(x).
\]
The set $\mathbb{I}:=\mathbb{F}\times\{1,2,3,4\}$, identified with $\mathbb{X}_{\mathbb{T}}\times\{1,2,3,4\}$,
indexes the rows $(x,\alpha)$ and the columns $(x',\alpha')$ of the matrices below;
a row or a column $(x,\alpha)$ \emph{stands at the site} $x$.
The matrix $D=\left(D[(x,\alpha),(x',\alpha')]\right)_{(x,\alpha),(x',\alpha')\in\mathbb{I}}$ is defined by its $4\times4$ blocks
$D[x,x']:=\left(D[(x,\alpha),(x',\alpha')]\right)_{\alpha,\alpha'=1,\dots,4}$, $x,x'\in\mathbb{F}$:
\begin{equation}\label{eq:positivity-blocks}
D[x,x']=
\begin{cases}
M\cdot1, & x'=x,\\
-\chi(g_i(x))\cdot U_A\left(x,x^{(i)}\right)\cdot w\left(x,x^{(i)}\right), & x'=\widehat{x^{(i)}},\ i\in\{0,1,2,3\},\\
0, & x'\notin\left\{x,\widehat{x^{(0)}},\dots,\widehat{x^{(3)}}\right\},
\end{cases}
\end{equation}
where $1$ is the unit matrix of size four.
The cases of \eqref{eq:positivity-blocks} do not overlap:
by Lemma~\ref{lem:incongruent-neighbours} the classes of $x,x^{(0)},\dots,x^{(3)}$ are pairwise different,
hence so are their representatives $x,\widehat{x^{(0)}},\dots,\widehat{x^{(3)}}$,
and no block receives two contributions.
The matrix $D$ depends on $A$; the notation does not show it.

Let $\mathbb{F}'$ be a second set of representatives of the classes in $\mathbb{X}_{\mathbb{T}}$,
for $x\in\mathbb{F}$ let $g_x\in\Gamma$ be the element with $x+g_x\in\mathbb{F}'$,
and let $\tilde D$ be the matrix defined by \eqref{eq:positivity-blocks} with $\mathbb{F}'$ in place of $\mathbb{F}$,
its row and its column $(x+g_x,\alpha)$ being identified with $(x,\alpha)\in\mathbb{I}$.
The neighbours of $x+g_x$ are the sites $x^{(i)}+g_x$,
so that $g_i(x)$ is replaced by $g_i(x)+g_x-g_{x'}$ with $x'=\widehat{x^{(i)}}$,
while $U_A$ and the hopping matrices take the same values at $\left(x+g_x,x^{(i)}+g_x\right)$ as at $\left(x,x^{(i)}\right)$
(Definition~\ref{def:general-tori}).
As $\chi$ is a homomorphism, it follows that
\[
\tilde D[(x,\alpha),(x',\alpha')]=\chi(g_x)\cdot D[(x,\alpha),(x',\alpha')]\cdot\chi(g_{x'})^{-1},\qquad (x,\alpha),(x',\alpha')\in\mathbb{I}.
\]
Thus $\tilde D$ arises from $D$ by conjugation with the diagonal matrix whose entry at $(x,\alpha)$ is $\chi(g_x)$;
this matrix is unitary, its entries are values of $\chi$, and it does not depend on $A$.
In particular $\det\tilde D=\det D$.
\end{definition}

With the notation of Definition~\ref{def:operator-matrix}, the matrix $D$ represents $D^{\mathbb{T}}_A$.
Every site is $x'=\widehat{x'}+g$ with a unique $g\in\Gamma$,
and a spinor $\psi$ of $\mathbb{T}$ satisfies $\psi(x')=\chi(g)\cdot\psi\left(\widehat{x'}\right)$;
thus a spinor is determined by its restriction to $\mathbb{F}$, and every map $\mathbb{F}\to\C^4$ is the restriction of a spinor.
For $x\in\mathbb{F}$ one has $\psi\left(x^{(i)}\right)=\chi(g_i(x))\cdot\psi\left(\widehat{x^{(i)}}\right)$, and therefore
\[
\left(D^{\mathbb{T}}_A\psi\right)(x)=M\cdot\psi(x)-\sum_{i=0}^{3}\chi(g_i(x))\cdot w\left(x,x^{(i)}\right)\cdot U_A\left(x,x^{(i)}\right)\cdot\psi\left(\widehat{x^{(i)}}\right),
\]
that is, $\left(D^{\mathbb{T}}_A\psi\right)_\alpha(x)=\sum_{(x',\alpha')\in\mathbb{I}}D[(x,\alpha),(x',\alpha')]\cdot\psi_{\alpha'}(x')$ for $(x,\alpha)\in\mathbb{I}$.
Hence $\det D=\det D^{\mathbb{T}}_A$.
In the Grassmann generators at the sites of $\mathbb{F}$,
the form $\sum_{x\in\mathbb{F}}\bar\psi(x)\cdot\left(D^{\mathbb{T}}_A\psi\right)(x)$,
which for the tori of Definition~\ref{def:tori} is the action $S_{\mathbb{T}}(A)$,
equals $\bar\psi D\psi$ in the notation of Definition~\ref{def:grassmann}:
the row $(x,\alpha)$ belongs to $\bar\psi_\alpha(x)$, the column $(x',\alpha')$ to $\psi_{\alpha'}(x')$.

\begin{definition}[The determinant polynomial]\label{def:determinant-polynomial}
Let $\rho$, $w_0,\dots,w_3$, $q$ and $v_0$ be as in Definition~\ref{def:general-tori}, and let $M\in\C$.
Let $\mathbb{T}=(\Gamma,\chi)$ be a general torus, $A$ a real $\Gamma$-periodic one-cochain,
and $\mathbb{F}$, $\mathbb{I}$ and $D$ a set of representatives, the index set and the matrix of Definition~\ref{def:operator-matrix}.

For $x\in\mathbb{X}$ let $\nu$ be the coordinate with $|y_\nu(x)|=v_1$.
By Proposition~\ref{prop:geometry}\eqref{it:geometry-projection} and Remark~\ref{rem:link-phases},
$U_A\left(x,x^{(i)}\right)=1$ for $i=0,1,2$,
and the bond $\left\{x,x^{(3)}\right\}$ lies over a link parallel to $e_\nu$:
\begin{itemize}
\item
if $y_\nu(x)=v_1$, then $\pi\left(x^{(3)}\right)=\pi(x)+e_\nu$, the bond lies over the link $(\pi(x),\nu)$,
and $U_A\left(x,x^{(3)}\right)=u_b$ for the class $b\in\mathbb{T}^{(1)}_\nu$ of this link;
the site $x$ is called the \emph{lower end} of a bond over $b$;
\item
if $y_\nu(x)=-v_1$, then $\pi\left(x^{(3)}\right)=\pi(x)-e_\nu$, the bond lies over the link $(\pi(x)-e_\nu,\nu)$,
and $U_A\left(x,x^{(3)}\right)=u_b^{-1}$ for the class $b\in\mathbb{T}^{(1)}_\nu$ of this link;
the site $x$ is called the \emph{upper end} of a bond over $b$.
\end{itemize}
Consequently every block of $D$ other than the blocks $D\left[x,\widehat{x^{(3)}}\right]$, $x\in\mathbb{F}$,
is independent of $A$,
and $D\left[x,\widehat{x^{(3)}}\right]$ is the matrix $-\chi(g_3(x))\cdot w\left(x,x^{(3)}\right)$, which is independent of $A$,
multiplied by $u_b$ if $x$ is the lower end and by $u_b^{-1}$ if $x$ is the upper end of a bond over $b$.

The \emph{determinant polynomial} of $\mathbb{T}$ is the Laurent polynomial with complex coefficients
\begin{equation}\label{eq:positivity-leibniz}
P_{\mathbb{T}}:=\sum_{\mathfrak{s}\in\operatorname{Sym}(\mathbb{I})}\operatorname{sgn}(\mathfrak{s})\cdot\prod_{(x,\alpha)\in\mathbb{I}}D[(x,\alpha),\mathfrak{s}(x,\alpha)]
\end{equation}
in the indeterminates $u_b$, $b\in\mathbb{T}^{(1)}$, which are denoted by the same letters as the phases.
In \eqref{eq:positivity-leibniz} the entries of $D$ are read as elements of the ring $\C\left[u_b,u_b^{-1}:\ b\in\mathbb{T}^{(1)}\right]$:
in a block $D\left[x,\widehat{x^{(3)}}\right]$ the factor $U_A\left(x,x^{(3)}\right)$ is replaced by the indeterminate $u_b$ if $x$ is the lower end,
and by $u_b^{-1}$ if $x$ is the upper end, of a bond over $b$;
the entries of all other blocks are the complex numbers given by \eqref{eq:positivity-blocks}.
For $\lambda=(\lambda_b)_{b\in\mathbb{T}^{(1)}}\in\Z^{\mathbb{T}^{(1)}}$ put $u^\lambda:=\prod_{b\in\mathbb{T}^{(1)}}u_b^{\lambda_b}$.

Read in this way, the entries of $D$, and with them $P_{\mathbb{T}}$, are the same for every real $\Gamma$-periodic one-cochain $A$.
For each such $A$, the substitution $u_b=\exp\left(iqA(b)\right)$, $b\in\mathbb{T}^{(1)}$,
returns the entries of the matrix $D$ formed with this $A$,
which is the matrix of $D^{\mathbb{T}}_A$ in the coordinates given by the values at the sites of $\mathbb{F}$.
The Leibniz formula for the determinant therefore yields
\[
\det D^{\mathbb{T}}_A=P_{\mathbb{T}}\left(\left(\exp\left(iqA(b)\right)\right)_{b\in\mathbb{T}^{(1)}}\right).
\]
Moreover, $P_{\mathbb{T}}$ is the same for every set of representatives $\mathbb{F}$.
To see this, let $\mathbb{F}'$, $g_x$ and $\tilde D$ be as in Definition~\ref{def:operator-matrix}.
Let $x\in\mathbb{F}$ and $b\in\mathbb{T}^{(1)}$.
Since $y(x+g_x)=y(x)$ and $\pi(x+g_x)=\pi(x)+g_x$,
the site $x+g_x$ is the lower end of a bond over $b$ exactly when $x$ is,
and the upper end of a bond over $b$ exactly when $x$ is;
so $U_A\left(x,x^{(3)}\right)$ and $U_A\left(x+g_x,x^{(3)}+g_x\right)$ are replaced by the same indeterminate,
and the relation between the entries of $\tilde D$ and of $D$ in Definition~\ref{def:operator-matrix} persists for the entries read in the ring.
Written for $\tilde D$, the term of a permutation $\mathfrak{s}$ in \eqref{eq:positivity-leibniz} is therefore the term written for $D$,
multiplied by the factors $\chi(g_x)$ of all rows $(x,\alpha)$ and $\chi(g_{x'})^{-1}$ of all columns $(x',\alpha')$;
these factors cancel, because $\mathfrak{s}$ is a bijection of $\mathbb{I}$.
Hence $\tilde D$ and $D$ give the same polynomial.
The coefficients of $P_{\mathbb{T}}$ are functions of $M$, $\rho$, $(w_0,\dots,w_3)$, $\chi$ and $\Gamma$ alone:
the charge $q$ enters $\det D^{\mathbb{T}}_A$ only through the substitution $u_b=\exp\left(iqA(b)\right)$,
and replacing $v_0$ by another point of $C_0$ leaves the coefficients unchanged (Remark~\ref{rem:base-point}).

A permutation $\mathfrak{s}\in\operatorname{Sym}(\mathbb{I})$ is called \emph{admissible} if it sends each row $(x,\alpha)$
to a column that stands at $x$ or at the representative $\widehat{x^{(i)}}$ of one of the four neighbours of $x$.
Let $\mathfrak{s}$ be admissible, $\nu\in\{1,2,3\}$ and $b\in\mathbb{T}^{(1)}_\nu$.
Then $n^{\uparrow}_b(\mathfrak{s})$ is the number of the rows $(x,\alpha)\in\mathbb{I}$
such that $x$ is the lower end of a bond over $b$ and $\mathfrak{s}(x,\alpha)$ stands at $\widehat{x^{(3)}}$,
and $n^{\downarrow}_b(\mathfrak{s})$ is the number of the rows $(x,\alpha)\in\mathbb{I}$
such that $x$ is the upper end of a bond over $b$ and $\mathfrak{s}(x,\alpha)$ stands at $\widehat{x^{(3)}}$.
The \emph{exponent} of $\mathfrak{s}$ is $\lambda(\mathfrak{s})\in\Z^{\mathbb{T}^{(1)}}$,
\begin{equation}\label{eq:positivity-exponent}
\lambda_b(\mathfrak{s}):=n^{\uparrow}_b(\mathfrak{s})-n^{\downarrow}_b(\mathfrak{s}),\qquad b\in\mathbb{T}^{(1)} .
\end{equation}

The term of a permutation $\mathfrak{s}$ in \eqref{eq:positivity-leibniz} is $\operatorname{sgn}(\mathfrak{s})\cdot\prod_{(x,\alpha)\in\mathbb{I}}D[(x,\alpha),\mathfrak{s}(x,\alpha)]$.
If $\mathfrak{s}$ is not admissible, one of its factors is an entry of a block that vanishes by \eqref{eq:positivity-blocks}, and the term is $0$.
If $\mathfrak{s}$ is admissible, its term is a complex number, possibly $0$, times the single monomial $u^{\lambda(\mathfrak{s})}$,
with $\lambda(\mathfrak{s})$ as in \eqref{eq:positivity-exponent}.
Indeed, every $x\in\mathbb{F}$ is the lower end or the upper end of a bond over exactly one $b\in\mathbb{T}^{(1)}$, namely of the bond $\left\{x,x^{(3)}\right\}$.
The factor $D[(x,\alpha),\mathfrak{s}(x,\alpha)]$ is a complex number if $\mathfrak{s}(x,\alpha)$ does not stand at $\widehat{x^{(3)}}$;
otherwise it is a complex number times $u_b$ if $x$ is the lower end, and a complex number times $u_b^{-1}$ if $x$ is the upper end, of a bond over $b$.
The product of the factors is therefore a complex number times
$\prod_{b\in\mathbb{T}^{(1)}}u_b^{n^{\uparrow}_b(\mathfrak{s})-n^{\downarrow}_b(\mathfrak{s})}=u^{\lambda(\mathfrak{s})}$.
\end{definition}

\subsection{Thin layers on a torus}\label{sec:torus-layers}

\begin{definition}[The layer index]\label{def:layer-index}
Let $\mu\in\{1,2,3\}$ and $x\in\mathbb{X}$.
The coordinates of $y(x)=x-\pi(x)$ have the moduli $v_1,v_2,v_3$ in some order
(Definition~\ref{def:fermion-sites}),
so that $y_\mu(x)$ is one of the six numbers
\[
-v_1<-v_2<-v_3<v_3<v_2<v_1 .
\]
The \emph{layer index} $k_\mu(x)\in\{1,\dots,6\}$ is the position of $y_\mu(x)$ in this increasing list.
For $g\in\Z^3$ one has $\pi(x+g)=\pi(x)+g$ and $y(x+g)=y(x)$
(Definition~\ref{def:fermion-sites}),
hence $k_\mu(x+g)=k_\mu(x)$.
Let $\mathbb{T}=(\Gamma,\chi)$ be a general torus (Definition~\ref{def:general-tori}).
Then $k_\mu$ is constant on the classes of $\mathbb{X}_{\mathbb{T}}=\mathbb{X}/\Gamma$
and defines a function $k_\mu\colon\mathbb{X}_{\mathbb{T}}\to\{1,\dots,6\}$.
The \emph{thin layers} of $\mathbb{T}$ in the direction $e_\mu$ are the sets
\[
\mathbb{X}^{\mu}_{\mathbb{T},k}:=\{\xi\in\mathbb{X}_{\mathbb{T}}:k_\mu(\xi)=k\},
\qquad k=1,\dots,6,
\]
and $\#\left|\mathbb{X}^{\mu}_{\mathbb{T},k}\right|=8|\mathbb{T}|$ for each $k$.
Layer indices are regarded as residues modulo~$6$, each written as its representative in $\{1,\dots,6\}$;
accordingly, for a layer index $k$ the indices $k+1$ and $k-1$ are again taken in $\{1,\dots,6\}$,
the successor of $6$ being $1$ and the predecessor of $1$ being $6$.
\end{definition}

The number of elements of a thin layer is found from the decomposition of the sites.
By Lemma~\ref{lem:chambers}\eqref{it:chambers-group} and~\eqref{it:chambers-transitive},
every site has the form $x=\pi(x)+\vec g v_0$ with a unique pair $(\pi(x),\vec g)\in\Z^3\times\Ocub$,
and $\vec g\mapsto\vec g v_0$ is injective on $\Ocub$,
since the coordinates of $v_0$ are positive and pairwise different.
So $x\mapsto(\pi(x),y(x))$ is a bijection of $\mathbb{X}$ onto $\Z^3\times\Ocub v_0$,
and it carries the translation by $g\in\Gamma$ into the map $(v,y)\mapsto(v+g,y)$.
Consequently, for each of the $48$ vectors $y\in\Ocub v_0$,
exactly $[\Z^3:\Gamma]=|\mathbb{T}|$ classes of $\mathbb{X}_{\mathbb{T}}$ consist of sites $x$ with $y(x)=y$.
Among these $48$ vectors, exactly eight have a prescribed $\mu$-th coordinate:
the two other coordinates carry the two remaining moduli in one of two orders,
with one of four choices of signs.
Each thin layer therefore consists of $8|\mathbb{T}|$ classes.

For $\mu=1$ the layer index describes the thin layers of Definition~\ref{def:two-layer}:
a site $x$ lies in $\mathbb{X}_{j,k}$ if and only if $\pi(x)_1=j$ and $k_1(x)=k$.
Indeed, $|x_1-\pi(x)_1|=|y_1(x)|<1/2$,
so $\pi(x)_1$ is the only integer $j$ with $|x_1-j|<1/2$,
and $x_1-\pi(x)_1=y_1(x)$.

\begin{lemma}[Neighbours across the layers]\label{lem:torus-layers}
Let $\mathbb{T}=(\Gamma,\chi)$ be a general torus (Definition~\ref{def:general-tori}),
let $A$ be a real $\Gamma$-periodic one-cochain,
and let $\mu\in\{1,2,3\}$, $x\in\mathbb{X}$ and $k:=k_\mu(x)$.
\begin{enumerate}
\item\label{it:up-down-neighbours}
Exactly one of the four neighbours of $x$, denoted $x^{\uparrow}$, has the layer index $k_\mu\left(x^{\uparrow}\right)=k+1$,
and exactly one, denoted $x^{\downarrow}$, has the layer index $k_\mu\left(x^{\downarrow}\right)=k-1$;
the two other neighbours have the layer index $k$.
Moreover $\left(x^{\uparrow}\right)_\mu>x_\mu>\left(x^{\downarrow}\right)_\mu$.
\item\label{it:up-bond-type}
The bond $\left\{x,x^{\uparrow}\right\}$ has the type $i_k$, where
\[
(i_1,i_2,i_3,i_4,i_5,i_6):=(2,1,0,1,2,3).
\]
\item\label{it:up-bijection}
One has $\left(x^{\uparrow}\right)^{\downarrow}=x=\left(x^{\downarrow}\right)^{\uparrow}$,
and $(x+g)^{\uparrow}=x^{\uparrow}+g$ for every $g\in\Z^3$.
Consequently $x\mapsto x^{\uparrow}$ induces a bijection
\[
\mathfrak{u}\colon\mathbb{X}_{\mathbb{T}}\to\mathbb{X}_{\mathbb{T}},\qquad \xi\mapsto\xi^{\uparrow},
\]
whose inverse $\xi\mapsto\xi^{\downarrow}$ is induced by $x\mapsto x^{\downarrow}$,
and $\mathfrak{u}\left(\mathbb{X}^{\mu}_{\mathbb{T},k'}\right)=\mathbb{X}^{\mu}_{\mathbb{T},k'+1}$ for $k'=1,\dots,6$.
\item\label{it:up-phases}
If $k\le5$, then $U_A\left(x,x^{\uparrow}\right)=1=U_A\left(x^{\uparrow},x\right)$.
If $k=6$, then $U_A\left(x,x^{\uparrow}\right)=u_b$ and $U_A\left(x^{\uparrow},x\right)=u_b^{-1}$,
where $b\in\mathbb{T}^{(1)}_\mu$ is the class of the link $(\pi(x),\mu)$.
Read the entries of the matrix $D$ of Definition~\ref{def:operator-matrix}
as Laurent polynomials in the variables $u_{b'}$, $b'\in\mathbb{T}^{(1)}$,
as in Definition~\ref{def:determinant-polynomial},
and let $D[\xi,\xi']$ denote the block of $D$ at the representatives in $\mathbb{F}$
of two classes $\xi,\xi'\in\mathbb{X}_{\mathbb{T}}$.
For every $b\in\mathbb{T}^{(1)}_\mu$,
the variable $u_b$ occurs in the blocks $D\left[\xi,\xi^{\uparrow}\right]$
and the variable $u_b^{-1}$ in the blocks $D\left[\xi^{\uparrow},\xi\right]$,
where $\xi$ ranges over the eight classes in $\mathbb{X}^{\mu}_{\mathbb{T},6}$
whose sites $x'$ satisfy $(\pi(x'),\mu)\in b$;
neither $u_b$ nor $u_b^{-1}$ occurs in any other block of $D$.
If a site $x'\in\mathbb{X}$ satisfies $k_\mu(x')=6$, then $x'$ is the lower end of a bond over some class in $\mathbb{T}^{(1)}_\mu$,
and if it satisfies $k_\mu(x')=1$, then $x'$ is the upper end of a bond over some class in $\mathbb{T}^{(1)}_\mu$;
in either case there is only one such bond and only one such class.
\end{enumerate}
\end{lemma}

\begin{proof}
Let $\nu$ be the coordinate with $|y_\nu(x)|=v_1$.
By Lemma~\ref{lem:neighbours}, the vector $x^{(0)}-\pi(x)$ arises from $y(x)$ by a change of sign in the coordinate of modulus $v_3$;
the vectors $x^{(1)}-\pi(x)$ and $x^{(2)}-\pi(x)$ arise from $y(x)$
when the moduli $v_2$ and $v_3$, respectively $v_1$ and $v_2$, change places,
every coordinate keeping its sign;
and $x^{(3)}-\pi(x)$ arises from $y(x)$ when the coordinate $y_\nu(x)=\pm v_1$ is replaced by $\pm(1-v_1)$.
According to~\eqref{eq:chamber-neighbours}, moreover,
$\pi\left(x^{(i)}\right)=\pi(x)$ for $i=0,1,2$ and $\pi\left(x^{(3)}\right)=\pi(x)\pm e_\nu$.
Hence $y\left(x^{(i)}\right)=x^{(i)}-\pi(x)$ for $i=0,1,2$,
whereas $y\left(x^{(3)}\right)$ agrees with $y(x)$ outside the coordinate $\nu$,
and $y_\nu\left(x^{(3)}\right)=\pm(1-v_1)\mp1=\mp v_1$.
If $y_\mu\left(x^{(i)}\right)=y_\mu(x)$ and $\pi\left(x^{(i)}\right)_\mu=\pi(x)_\mu$ for some $i$,
then $\left(x^{(i)}\right)_\mu=x_\mu$ and $k_\mu\left(x^{(i)}\right)=k$.

Write $y_\mu(x)=\eta\cdot v_a$ with $\eta\in\{1,-1\}$ and $a\in\{1,2,3\}$.
The three values of $a$ are treated in turn.
\begin{itemize}
\item
$a=3$, that is, $k=3$ if $\eta=-1$ and $k=4$ if $\eta=1$.
The neighbour of type $0$ has $y_\mu=-\eta\cdot v_3$, hence the index $4$ if $k=3$ and the index $3$ if $k=4$.
The neighbour of type $1$ has $y_\mu=\eta\cdot v_2$, hence the index $2$ if $k=3$ and the index $5$ if $k=4$.
For $i=2,3$ one has $y_\mu\left(x^{(i)}\right)=y_\mu(x)$ and $\pi\left(x^{(i)}\right)_\mu=\pi(x)_\mu$.
\item
$a=2$, that is, $k=2$ if $\eta=-1$ and $k=5$ if $\eta=1$.
The neighbour of type $1$ has $y_\mu=\eta\cdot v_3$, hence the index $3$ if $k=2$ and the index $4$ if $k=5$.
The neighbour of type $2$ has $y_\mu=\eta\cdot v_1$, hence the index $1$ if $k=2$ and the index $6$ if $k=5$.
For $i=0,3$ one has $y_\mu\left(x^{(i)}\right)=y_\mu(x)$ and $\pi\left(x^{(i)}\right)_\mu=\pi(x)_\mu$.
\item
$a=1$, that is, $k=1$ if $\eta=-1$ and $k=6$ if $\eta=1$; here $\nu=\mu$.
The neighbour of type $2$ has $y_\mu=\eta\cdot v_2$, hence the index $2$ if $k=1$ and the index $5$ if $k=6$.
The neighbour of type $3$ has the $\mu$-th coordinate $\pi(x)_\mu+\eta\cdot(1-v_1)=\pi(x)_\mu+\eta-\eta\cdot v_1$:
for $k=6$ it has the layer index $1$ and the photon site $\pi(x)+e_\mu$,
and for $k=1$ the layer index $6$ and the photon site $\pi(x)-e_\mu$.
For $i=0,1$ one has $y_\mu\left(x^{(i)}\right)=y_\mu(x)$ and $\pi\left(x^{(i)}\right)_\mu=\pi(x)_\mu$.
\end{itemize}
Thus, whatever the value of $k$, the value $k$ occurs twice among the layer indices of $x^{(0)},\dots,x^{(3)}$,
and each of the values $k+1$ and $k-1$ occurs once;
these two values are different, because $6$ does not divide $(k+1)-(k-1)=2$.
This is the first assertion of~\eqref{it:up-down-neighbours}.
The neighbour with the index $k+1$ has the type $2,1,0,1,2,3$ for $k=1,\dots,6$,
which is~\eqref{it:up-bond-type},
and the neighbour with the index $k-1$ has the type $3,2,1,0,1,2$ for $k=1,\dots,6$.
For $k\le5$ the neighbour $x^{\uparrow}$, and for $k\ge2$ the neighbour $x^{\downarrow}$,
is therefore of type $0$, $1$ or $2$:
its photon site is $\pi(x)$,
and its coordinate $y_\mu$ is the next larger, respectively the next smaller, of the six numbers.
This gives $\left(x^{\uparrow}\right)_\mu>x_\mu$ for $k\le5$ and $\left(x^{\downarrow}\right)_\mu<x_\mu$ for $k\ge2$.
In the two remaining cases the neighbour is of type $3$, and $v_1<1/2$ gives
\[
\begin{aligned}
\left(x^{\uparrow}\right)_\mu&=\pi(x)_\mu+1-v_1>\pi(x)_\mu+v_1=x_\mu&&\text{for }k=6,\\
\left(x^{\downarrow}\right)_\mu&=\pi(x)_\mu-1+v_1<\pi(x)_\mu-v_1=x_\mu&&\text{for }k=1.
\end{aligned}
\]
This completes the proof of~\eqref{it:up-down-neighbours} and~\eqref{it:up-bond-type}.

For~\eqref{it:up-bijection}, recall that the neighbour relation is symmetric:
$\left(x^{(i)}\right)^{(i)}=x$ for $i=0,\dots,3$ (Definition~\ref{def:fermion-sites}).
If $x^{\uparrow}=x^{(i)}$, then $x$ is the neighbour of type $i$ of $x^{\uparrow}$,
and its index $k$ is the index $k+1$ of $x^{\uparrow}$ lowered by one, modulo $6$;
the uniqueness in~\eqref{it:up-down-neighbours}, applied to the site $x^{\uparrow}$, gives $\left(x^{\uparrow}\right)^{\downarrow}=x$.
The identity $\left(x^{\downarrow}\right)^{\uparrow}=x$ follows in the same way.
Now let $g\in\Z^3$, and write $x=g'v_0$ with $g'\in W$
(Lemma~\ref{lem:chambers}\eqref{it:chambers-transitive}).
The translation $\tau_g$ belongs to $W$ (Lemma~\ref{lem:chambers}\eqref{it:chambers-group}),
and $x+g=(\tau_g g')v_0$;
by Definition~\ref{def:fermion-sites},
\[
(x+g)^{(i)}=\tau_g g' s_i v_0=x^{(i)}+g,\qquad i=0,\dots,3.
\]
As $k_\mu$ is invariant under $\Z^3$,
the neighbour $x^{\uparrow}+g$ of $x+g$ has the index $k_\mu(x+g)+1$,
so that $(x+g)^{\uparrow}=x^{\uparrow}+g$ by uniqueness;
likewise $(x+g)^{\downarrow}=x^{\downarrow}+g$.
For $g\in\Gamma$ this shows that $x\mapsto x^{\uparrow}$ and $x\mapsto x^{\downarrow}$ map each class of $\mathbb{X}_{\mathbb{T}}$ into a class,
so that they induce maps $\xi\mapsto\xi^{\uparrow}$ and $\xi\mapsto\xi^{\downarrow}$ of $\mathbb{X}_{\mathbb{T}}$ into itself,
and by the identities $\left(x^{\uparrow}\right)^{\downarrow}=x=\left(x^{\downarrow}\right)^{\uparrow}$ these two maps are inverse to each other.
Hence $\mathfrak{u}$ is a bijection.
It raises the layer index by one and its inverse lowers it by one,
which gives $\mathfrak{u}\left(\mathbb{X}^{\mu}_{\mathbb{T},k'}\right)=\mathbb{X}^{\mu}_{\mathbb{T},k'+1}$.

It remains to prove~\eqref{it:up-phases}.
By Definition~\ref{def:determinant-polynomial},
$U_A\left(x',x'^{(i)}\right)=1$ for every site $x'$ and $i\in\{0,1,2\}$,
whereas the bond $\left\{x',x'^{(3)}\right\}$ lies over a link parallel to $e_\nu$,
$\nu$ the coordinate with $|y_\nu(x')|=v_1$.
If $y_\nu(x')=v_1$, the site $x'$ is the lower end of this bond,
the link is $(\pi(x'),\nu)$, and $U_A\left(x',x'^{(3)}\right)=u_b$ for the class $b$ of this link.
If $y_\nu(x')=-v_1$, the site $x'$ is the upper end of this bond,
the link is $(\pi(x')-e_\nu,\nu)$, and $U_A\left(x',x'^{(3)}\right)=u_b^{-1}$ for the class $b$ of that link.

Let $k\le5$.
By~\eqref{it:up-bond-type} the bond $\left\{x,x^{\uparrow}\right\}$ has the type $i_k\in\{0,1,2\}$,
and $x=\left(x^{\uparrow}\right)^{(i_k)}$;
so $U_A\left(x,x^{\uparrow}\right)=1$ and $U_A\left(x^{\uparrow},x\right)=1$.
Let $k=6$.
Then $y_\mu(x)=v_1$, so $\nu=\mu$ and $x^{\uparrow}=x^{(3)}$,
and $U_A\left(x,x^{\uparrow}\right)=u_b$ with $b$ the class of $(\pi(x),\mu)$.
The site $x^{\uparrow}$ satisfies $y_\mu\left(x^{\uparrow}\right)=-v_1$ and $\pi\left(x^{\uparrow}\right)=\pi(x)+e_\mu$,
and $x$ is its neighbour of type $3$;
since the link $\left(\pi\left(x^{\uparrow}\right)-e_\mu,\mu\right)=(\pi(x),\mu)$ has the class $b$,
one obtains $U_A\left(x^{\uparrow},x\right)=u_b^{-1}$.

Turning to the matrix $D$,
recall from~\eqref{eq:positivity-blocks} and Lemma~\ref{lem:incongruent-neighbours}
that the variables $u_{b'}$ enter $D$ only through the factor $U_A\left(x',x'^{(i)}\right)$
of the block $D\left[x',\widehat{x'^{(i)}}\right]$, $x'\in\mathbb{F}$, $i\in\{0,\dots,3\}$,
and that each block of $D$ contains at most one such factor.
By the description of $U_A$ just given, this factor is a variable only for $i=3$,
and it is $u_b$ or $u_b^{-1}$ with $b\in\mathbb{T}^{(1)}_\mu$ exactly when $\nu=\mu$,
that is, when $k_\mu(x')\in\{1,6\}$.
If $k_\mu(x')=6$, then $x'^{(3)}=x'^{\uparrow}$,
and the block $D\left[x',\widehat{x'^{\uparrow}}\right]=D\left[\xi,\xi^{\uparrow}\right]$, $\xi$ the class of $x'$,
contains $u_b$ with $b$ the class of $(\pi(x'),\mu)$.
If $k_\mu(x')=1$, then $x'^{(3)}=x'^{\downarrow}$;
the class $\xi$ of $x'^{\downarrow}$ lies in $\mathbb{X}^{\mu}_{\mathbb{T},6}$,
the class of $x'$ is $\xi^{\uparrow}$ by~\eqref{it:up-bijection},
and the block $D\left[x',\widehat{x'^{\downarrow}}\right]=D\left[\xi^{\uparrow},\xi\right]$ contains
$U_A\left(x',x'^{\downarrow}\right)=u_b^{-1}$ with $b$ the class of $\left(\pi\left(x'^{\downarrow}\right),\mu\right)$,
by the case $k=6$ applied to the site $x'^{\downarrow}$, whose neighbour $\left(x'^{\downarrow}\right)^{\uparrow}$ is $x'$.
Since $\pi(x'+g)=\pi(x')+g$ for $g\in\Gamma$,
the class of the link $(\pi(x'),\mu)$ is the same for all sites $x'$ of one class of $\mathbb{X}_{\mathbb{T}}$.
Therefore $u_b$ occurs exactly in the blocks $D\left[\xi,\xi^{\uparrow}\right]$
and $u_b^{-1}$ exactly in the blocks $D\left[\xi^{\uparrow},\xi\right]$,
where $\xi\in\mathbb{X}^{\mu}_{\mathbb{T},6}$ and $(\pi(x'),\mu)\in b$ for the sites $x'$ of $\xi$.

These classes $\xi$ are eight in number.
Let $(v,\mu)$, $v\in\Z^3$, be a link of the class $b$.
Each class with the required property contains a site $x'$ with $\pi(x')=v$,
obtained from any of its sites by a translation by an element of $\Gamma$,
and a class in $\mathbb{X}^{\mu}_{\mathbb{T},6}$ that contains such a site has the required property.
By the bijection $x\mapsto(\pi(x),y(x))$ of $\mathbb{X}$ onto $\Z^3\times\Ocub v_0$
described after Definition~\ref{def:layer-index},
the sites $x'$ with $\pi(x')=v$ and layer index $6$ are the eight sites $v+y$ with $y\in\Ocub v_0$ and $y_\mu=v_1$.
Two of them differ by a non-zero vector whose coordinates lie in $(-1,1)$,
so their difference is not in $\Gamma\subset\Z^3$, and they lie in eight different classes.

Finally, let $x'\in\mathbb{X}$ with $k_\mu(x')\in\{1,6\}$.
Then $y_\mu(x')=v_1$ if $k_\mu(x')=6$ and $y_\mu(x')=-v_1$ if $k_\mu(x')=1$,
so that $\mu$ is the coordinate in which $y(x')$ has the modulus $v_1$.
By Definition~\ref{def:determinant-polynomial},
the bond $\left\{x',x'^{(3)}\right\}$ lies over the link $(\pi(x'),\mu)$ in the first case and over the link $(\pi(x')-e_\mu,\mu)$ in the second,
and $x'$ is the lower end, respectively the upper end, of this bond over the class of that link, which belongs to $\mathbb{T}^{(1)}_\mu$.
The three other bonds at $x'$, of the types $0$, $1$ and $2$, lie over no link,
because their two end points have the same nearest photon site;
and a bond lies over at most one link, since a link is determined by its two end points.
Hence $\left\{x',x'^{(3)}\right\}$ is the only bond over a class in $\mathbb{T}^{(1)}_\mu$ of which $x'$ is the lower end, respectively the upper end,
and the class of its link is the only class in $\mathbb{T}^{(1)}_\mu$ with this property.
\end{proof}

\subsection{Degrees and the extreme coefficient}\label{sec:extreme-coefficient}

Let $\mathbb{T}=(\Gamma,\chi)$ be a general torus,
and let $\mathbb{F}$, $\mathbb{I}$ and $D$ be a set of representatives, the index set and the matrix of Definition~\ref{def:operator-matrix},
formed with some $M\in\C$ and some real $\Gamma$-periodic one-cochain $A$.
Through the identification of $\mathbb{I}=\mathbb{F}\times\{1,2,3,4\}$ with $\mathbb{X}_{\mathbb{T}}\times\{1,2,3,4\}$,
a class $\xi\in\mathbb{X}_{\mathbb{T}}$ with representative $x\in\mathbb{F}$ carries the four elements $(\xi,\alpha)=(x,\alpha)$, $\alpha\in\{1,2,3,4\}$, of $\mathbb{I}$.
As indices of rows of $D$ they are called the \emph{rows at}~$\xi$, as indices of columns the \emph{columns at}~$\xi$;
in the terms of Definition~\ref{def:operator-matrix}, these are the rows and the columns that stand at the site $x$.
A permutation $\mathfrak{s}$ of $\mathbb{I}$ maps the row $(\xi,\alpha)$ to the column $\mathfrak{s}(\xi,\alpha)$,
as in the Leibniz formula~\eqref{eq:positivity-leibniz}.

\begin{lemma}[Degrees]\label{lem:link-degrees}
Let $\mathbb{T}=(\Gamma,\chi)$ be a general torus (Definition~\ref{def:general-tori}),
let $\mathbb{F}$ and $\mathbb{I}$ be a set of representatives and the index set of Definition~\ref{def:operator-matrix},
and let $b\in\mathbb{T}^{(1)}$.
Lower and upper ends of bonds, admissible permutations $\mathfrak{s}$ of $\mathbb{I}$ and their numbers
$n^{\uparrow}_b(\mathfrak{s})$, $n^{\downarrow}_b(\mathfrak{s})$ and $\lambda_b(\mathfrak{s})$ are those of Definition~\ref{def:determinant-polynomial}.
Exactly eight sites of $\mathbb{F}$ are lower ends of bonds over $b$, and exactly eight are upper ends of bonds over $b$.
Every admissible permutation $\mathfrak{s}$ of $\mathbb{I}$ satisfies
\[
0\le n^{\uparrow}_b(\mathfrak{s})\le32
\qquad\text{and}\qquad
0\le n^{\downarrow}_b(\mathfrak{s})\le32,
\]
so that $\left|\lambda_b(\mathfrak{s})\right|\le32$.
Moreover, $\lambda_b(\mathfrak{s})=32$ if and only if $\mathfrak{s}$ has the following two properties:
for each of the eight lower ends $x\in\mathbb{F}$, the four columns $\mathfrak{s}(x,\alpha)$, $\alpha\in\{1,2,3,4\}$, stand at $\widehat{x^{(3)}}$,
the representative of the upper end of the bond $\left\{x,x^{(3)}\right\}$;
for each of the eight upper ends $x\in\mathbb{F}$, none of the four columns $\mathfrak{s}(x,\alpha)$ stands at $\widehat{x^{(3)}}$,
the representative of the lower end of that bond.
\end{lemma}

\begin{proof}
Let $\nu\in\{1,2,3\}$ be the direction with $b\in\mathbb{T}^{(1)}_\nu$,
and let $(v,\nu)$, $v\in\Z^3$, be a link in the class $b$;
the links in $b$ are the translates $(v+g,\nu)$, $g\in\Gamma$.
By Definition~\ref{def:determinant-polynomial},
the lower ends of the bonds over $b$ are the sites $x$ with $y_\nu(x)=v_1$ and $\pi(x)\in v+\Gamma$,
and the upper ends are the sites $x$ with $y_\nu(x)=-v_1$ and $\pi(x)\in v+e_\nu+\Gamma$.
There are exactly eight sites $x$ with $\pi(x)=v$ and $y_\nu(x)=v_1$,
and exactly eight sites with $\pi(x)=v+e_\nu$ and $y_\nu(x)=-v_1$,
because in either case the two other coordinates of $y(x)$ carry the moduli $v_2$ and $v_3$, in either order and with either sign.
By Lemma~\ref{lem:neighbours}, a site $x$ with $\pi(x)=v$ and $y_\nu(x)=v_1$ has the neighbour of type~$3$
\[
x^{(3)}=x+\left(1-2v_1\right)\cdot e_\nu,
\qquad
\pi\left(x^{(3)}\right)=v+e_\nu,
\qquad
y_\nu\left(x^{(3)}\right)=-v_1 .
\]
The eight bonds $\left\{x,x^{(3)}\right\}$ so obtained are distinct and lie over the link $(v,\nu)$;
by Proposition~\ref{prop:geometry}\eqref{it:geometry-projection}, exactly eight bonds lie over this link,
so these are all of them.
Their lower ends are the eight sites $x$,
and their upper ends are the eight sites $x^{(3)}$,
which exhaust the sites with $\pi=v+e_\nu$ and $y_\nu=-v_1$.

For $g\in\Gamma$ one has $\pi(x+g)=\pi(x)+g$ and $y(x+g)=y(x)$;
in particular, two sites with the same image under $\pi$ are congruent modulo $\Gamma$ only if they coincide.
Hence every lower end of a bond over a link in $b$ is congruent modulo $\Gamma$ to exactly one of the eight lower ends over $(v,\nu)$,
and every upper end to exactly one of the eight upper ends over $(v,\nu)$.
The lower ends of the bonds over $b$ therefore form exactly eight classes of $\mathbb{X}_{\mathbb{T}}$,
and the upper ends exactly eight further classes:
no class contains both a lower and an upper end, as $y_\nu$ is constant on classes.
The two conditions that describe the lower and the upper ends are invariant under the translations by $\Gamma$,
so each of these sixteen classes consists of ends of one kind only.
Consequently the sites of $\mathbb{F}$ that are lower ends of bonds over $b$ are the representatives of the eight classes of lower ends,
and those that are upper ends are the representatives of the eight classes of upper ends;
this proves the first assertion.
Four rows stand at each of these sixteen sites of $\mathbb{F}$.

By Definition~\ref{def:determinant-polynomial},
$n^{\uparrow}_b(\mathfrak{s})$ counts the rows $(x,\alpha)\in\mathbb{I}$ for which $x$ is one of the eight lower ends in $\mathbb{F}$
and $\mathfrak{s}(x,\alpha)$ stands at $\widehat{x^{(3)}}$,
and $n^{\downarrow}_b(\mathfrak{s})$ is the analogous count with the eight upper ends in place of the eight lower ends.
Each of the two numbers is therefore at least $0$ and at most $8\cdot4=32$,
and $\left|\lambda_b(\mathfrak{s})\right|\le32$ follows from~\eqref{eq:positivity-exponent}.
Since $n^{\uparrow}_b(\mathfrak{s})\le32$ and $n^{\downarrow}_b(\mathfrak{s})\ge0$,
the difference $\lambda_b(\mathfrak{s})=n^{\uparrow}_b(\mathfrak{s})-n^{\downarrow}_b(\mathfrak{s})$ equals $32$
exactly when $n^{\uparrow}_b(\mathfrak{s})=32$ and $n^{\downarrow}_b(\mathfrak{s})=0$.
The first of these equalities means that all $32$ rows $(x,\alpha)$ with $x$ a lower end are counted, which is the first of the two properties;
the second means that none of the $32$ rows $(x,\alpha)$ with $x$ an upper end is counted, which is the second.
\end{proof}

By Lemma~\ref{lem:link-degrees}, the exponent $\lambda(\mathfrak{s})$ of an admissible permutation $\mathfrak{s}$
takes at every link a value between $-32$ and $32$.

\begin{proposition}[The extreme coefficient]\label{prop:extreme-coefficient}
Let $\rho\in\R$, $M\in\C$, $w_0,w_1,w_2,w_3>0$ and $q>0$,
let $\mathbb{T}=(\Gamma,\chi)$ be a general torus (Definition~\ref{def:general-tori}),
and let $A$ be a real $\Gamma$-periodic one-cochain.
Let $\mathbb{F}$, $\mathbb{I}$ and $D$ be a set of representatives, the index set and the matrix of Definition~\ref{def:operator-matrix}, formed with these data;
admissible permutations of $\mathbb{I}$, their exponents $\lambda(\mathfrak{s})$ and the determinant polynomial $P_{\mathbb{T}}$ are those of Definition~\ref{def:determinant-polynomial}.
For $\xi\in\mathbb{X}_{\mathbb{T}}$ let $x\in\mathbb{F}$ denote the representative of $\xi$.
Let $\mu\in\{1,2,3\}$,
and let $x^{\uparrow}$, $x^{\downarrow}$, $\xi^{\uparrow}$, $\xi^{\downarrow}$ and $i_1,\dots,i_6$ be those of Lemma~\ref{lem:torus-layers} for this direction $\mu$.
Let $\lambda^{(\mu)}\in\Z^{\mathbb{T}^{(1)}}$ be the exponent with
$\lambda^{(\mu)}_b=32$ for $b\in\mathbb{T}^{(1)}_\mu$
and $\lambda^{(\mu)}_b=0$ for $b\in\mathbb{T}^{(1)}\setminus\mathbb{T}^{(1)}_\mu$.
\begin{enumerate}
\item\label{it:extreme-permutations}
Let $\mathfrak{s}$ be an admissible permutation of $\mathbb{I}$ with $\lambda_b(\mathfrak{s})=32$ for every $b\in\mathbb{T}^{(1)}_\mu$.
Then $\mathfrak{s}$ is of the form
\[
\mathfrak{s}(\xi,\alpha)=\left(\xi^{\uparrow},\mathfrak{p}_\xi(\alpha)\right),
\qquad(\xi,\alpha)\in\mathbb{I},
\]
with a permutation $\mathfrak{p}_\xi$ of $\{1,2,3,4\}$ for each $\xi\in\mathbb{X}_{\mathbb{T}}$.
Conversely, for every choice of permutations $\mathfrak{p}_\xi$ of $\{1,2,3,4\}$, $\xi\in\mathbb{X}_{\mathbb{T}}$,
this formula defines an admissible permutation $\mathfrak{s}$ of $\mathbb{I}$,
and $\lambda(\mathfrak{s})=\lambda^{(\mu)}$.
Consequently, $\lambda^{(\mu)}$ belongs to the set of the exponents $\lambda(\mathfrak{s})$ of admissible permutations $\mathfrak{s}$,
whose terms in~\eqref{eq:positivity-leibniz} may vanish,
and no other element of this set takes the value $32$ at every $b\in\mathbb{T}^{(1)}_\mu$.
\item\label{it:extreme-value}
For $x\in\mathbb{F}$ let $g^{\uparrow}(x)\in\Gamma$ be defined by $x^{\uparrow}=\widehat{x^{\uparrow}}+g^{\uparrow}(x)$,
and put
\[
\chi_{\mathbb{T}}:=\prod_{x\in\mathbb{F}}\chi\left(g^{\uparrow}(x)\right)^{4},
\qquad
C_*:=\left(\frac{w_0}{2}\right)^{32}\cdot\left(\frac{w_1}{2}\right)^{64}\cdot\left(\frac{w_2}{2}\right)^{64}\cdot\left(\frac{w_3}{2}\right)^{32}\cdot\left(1-\rho^{2}\right)^{96}.
\]
Then the coefficient of $u^{\lambda^{(\mu)}}$ in $P_{\mathbb{T}}$ is
\[
\chi_{\mathbb{T}}\cdot\prod_{\xi\in\mathbb{X}_{\mathbb{T}}}\det w\left(x,x^{\uparrow}\right)
=\chi_{\mathbb{T}}\cdot C_*^{|\mathbb{T}|},
\qquad\text{where}\qquad
C_*^{|\mathbb{T}|}=\prod_{k=1}^{6}\left(\left(\frac{w_{i_k}}{2}\right)^{4}\cdot\left(1-\rho^{2}\right)^{2}\right)^{8|\mathbb{T}|}.
\]
One has $\left|\chi_{\mathbb{T}}\right|=1$.
If $\chi$ takes only the values $\pm1$, as it does on every torus of Definition~\ref{def:tori},
then $\chi_{\mathbb{T}}=1$ and the coefficient is $C_*^{|\mathbb{T}|}$;
none of $M$, $\chi$, $q$, $v_0$ and $\mu$ enters this number,
which vanishes if $\rho^2=1$ and is positive if $\rho^2\neq1$.
\item\label{it:extreme-opposite}
The exponent $-\lambda^{(\mu)}$ behaves in the same way,
with $\xi^{\downarrow}$, $x^{\downarrow}$ and $\overline{\chi_{\mathbb{T}}}$ taking the places of $\xi^{\uparrow}$, $x^{\uparrow}$ and $\chi_{\mathbb{T}}$,
and with the same number $C_*^{|\mathbb{T}|}$.
Explicitly: if $\mathfrak{s}$ is admissible and $\lambda_b(\mathfrak{s})=-32$ for every $b\in\mathbb{T}^{(1)}_\mu$,
then $\mathfrak{s}(\xi,\alpha)=\left(\xi^{\downarrow},\mathfrak{p}_\xi(\alpha)\right)$ with a permutation $\mathfrak{p}_\xi$ of $\{1,2,3,4\}$ for each $\xi$;
conversely, for every choice of permutations $\mathfrak{p}_\xi$ of $\{1,2,3,4\}$, $\xi\in\mathbb{X}_{\mathbb{T}}$,
this formula defines an admissible permutation $\mathfrak{s}$ with $\lambda(\mathfrak{s})=-\lambda^{(\mu)}$,
so that $-\lambda^{(\mu)}$ belongs to the set of the exponents of admissible permutations
and no other element of this set takes the value $-32$ at every $b\in\mathbb{T}^{(1)}_\mu$;
and the coefficient of $u^{-\lambda^{(\mu)}}$ in $P_{\mathbb{T}}$ is
\[
\overline{\chi_{\mathbb{T}}}\cdot\prod_{\xi\in\mathbb{X}_{\mathbb{T}}}\det w\left(x,x^{\downarrow}\right)
=\overline{\chi_{\mathbb{T}}}\cdot C_*^{|\mathbb{T}|},
\]
which equals $C_*^{|\mathbb{T}|}$ if $\chi$ takes only the values $\pm1$.
\end{enumerate}
\end{proposition}

\begin{proof}
Write $\mathbb{X}_k:=\mathbb{X}^{\mu}_{\mathbb{T},k}$ for the thin layers in the direction $\mu$,
with the index $k$ read modulo~$6$ as in Definition~\ref{def:layer-index};
thus $\mathbb{X}_0=\mathbb{X}_6$ and $\mathbb{X}_7=\mathbb{X}_1$.
First, three properties of the bonds.

\emph{Neighbours.}
Let $\xi\in\mathbb{X}_k$.
The classes of the four neighbours of the representative of $\xi$ are
$\xi^{\uparrow}\in\mathbb{X}_{k+1}$, $\xi^{\downarrow}\in\mathbb{X}_{k-1}$ and two classes of $\mathbb{X}_k$
(Lemma~\ref{lem:torus-layers}(a)).
The map $\mathfrak{u}\colon\xi\mapsto\xi^{\uparrow}$ is a bijection of $\mathbb{X}_{\mathbb{T}}$ with the inverse $\xi\mapsto\xi^{\downarrow}$,
and it maps $\mathbb{X}_k$ onto $\mathbb{X}_{k+1}$
(Lemma~\ref{lem:torus-layers}(c)).
Consequently an admissible permutation maps every row at $\xi$
to a column at a class of $\mathbb{X}_k$, at $\xi^{\downarrow}$, or at $\xi^{\uparrow}$.

\emph{Bonds over links parallel to $e_\mu$.}
A site $x$ has $k_\mu(x)=6$ if and only if $y_\mu(x)=v_1$,
and $k_\mu(x)=1$ if and only if $y_\mu(x)=-v_1$ (Definition~\ref{def:layer-index}).
So the lower ends of the bonds over the links in the classes $b\in\mathbb{T}^{(1)}_\mu$ are exactly the sites in the classes of $\mathbb{X}_6$,
and their upper ends are exactly the sites in the classes of $\mathbb{X}_1$
(Lemma~\ref{lem:torus-layers}(d)).
If $k_\mu(x)=6$, the bond $\left\{x,x^{\uparrow}\right\}$ has the type $i_6=3$ (Lemma~\ref{lem:torus-layers}(b)),
so that $x^{(3)}=x^{\uparrow}$;
if $k_\mu(x)=1$, then $x=\left(x^{\downarrow}\right)^{\uparrow}$ with $k_\mu\left(x^{\downarrow}\right)=6$,
so that the bond $\left\{x^{\downarrow},x\right\}$ has the type $3$ and $x^{(3)}=x^{\downarrow}$.
Thus the bond at a lower end $\xi\in\mathbb{X}_6$ joins $\xi$ to $\xi^{\uparrow}$,
and the bond at an upper end $\xi\in\mathbb{X}_1$ joins $\xi$ to $\xi^{\downarrow}$.

\emph{Bonds over the other links.}
Let $\nu\neq\mu$, and let $\left\{x,x^{(3)}\right\}$ be a bond over a link parallel to $e_\nu$.
By Proposition~\ref{prop:geometry}\eqref{it:geometry-projection},
$\nu$ is the coordinate in which $y(x)$ has the modulus $v_1$, and $\pi\left(x^{(3)}\right)-\pi(x)=\pm e_\nu$.
The neighbour $x^{(3)}$ differs from $x$ only in this coordinate (Lemma~\ref{lem:neighbours}).
Hence $y_\mu\left(x^{(3)}\right)=y_\mu(x)$,
and the two ends of the bond lie in the same thin layer.

\emph{Proof of~\eqref{it:extreme-permutations}.}
Let $\mathfrak{s}$ be an admissible permutation with $\lambda_b(\mathfrak{s})=32$ for every $b\in\mathbb{T}^{(1)}_\mu$.
For $k\in\{1,\dots,6\}$ consider the assertions
\begin{itemize}
\item[$(\mathrm{A}_k)$] every row at a class $\xi\in\mathbb{X}_k$ is mapped by $\mathfrak{s}$ to a column at $\xi^{\uparrow}$;
\item[$(\mathrm{C}_{k+1})$] every column at a class of $\mathbb{X}_{k+1}$ is the image under $\mathfrak{s}$ of a row at a class of $\mathbb{X}_k$,
\end{itemize}
where $(\mathrm{C}_7)$ means $(\mathrm{C}_1)$, and the assertion
\begin{itemize}
\item[$(\mathrm{B})$] no row at a class $\xi\in\mathbb{X}_1$ is mapped by $\mathfrak{s}$ to a column at $\xi^{\downarrow}$.
\end{itemize}
By the description of the bonds over links parallel to $e_\mu$,
a site $x\in\mathbb{F}$ is the lower end of a bond over some $b\in\mathbb{T}^{(1)}_\mu$ exactly when its class $\xi$ lies in $\mathbb{X}_6$,
and then $\widehat{x^{(3)}}$ is the representative of $\xi^{\uparrow}$;
it is the upper end of such a bond exactly when $\xi\in\mathbb{X}_1$,
and then $\widehat{x^{(3)}}$ is the representative of $\xi^{\downarrow}$.
Lemma~\ref{lem:link-degrees}, applied to every $b\in\mathbb{T}^{(1)}_\mu$,
therefore turns the hypothesis into the conjunction of $(\mathrm{A}_6)$ and $(\mathrm{B})$.

Each assertion $(\mathrm{A}_k)$ entails $(\mathrm{C}_{k+1})$.
Indeed, assume $(\mathrm{A}_k)$ and let $\xi\in\mathbb{X}_k$.
By injectivity, the images of the four rows at $\xi$ are four distinct columns at $\xi^{\uparrow}$,
and no other column stands at $\xi^{\uparrow}$;
each column at $\xi^{\uparrow}$ therefore occurs among these images.
As $\xi$ ranges over $\mathbb{X}_k$, the class $\xi^{\uparrow}$ ranges over all of $\mathbb{X}_{k+1}$.

The assertions $(\mathrm{A}_k)$ are now derived one at a time, in the order $k=6,1,2,3,4,5$.
The first of them, $(\mathrm{A}_6)$, is part of the hypothesis, and it yields $(\mathrm{C}_1)$.
Next let $\xi\in\mathbb{X}_1$, and consider a row at $\xi$.
Its image is a column at a class of $\mathbb{X}_1$, at $\xi^{\downarrow}$, or at $\xi^{\uparrow}$.
A column of the first kind is already, by $(\mathrm{C}_1)$, the image of a row at a class of $\mathbb{X}_6$;
since the row under consideration stands at $\xi\notin\mathbb{X}_6$, injectivity excludes this case.
The assertion $(\mathrm{B})$ excludes $\xi^{\downarrow}$.
The image is therefore a column at $\xi^{\uparrow}$, which proves $(\mathrm{A}_1)$ and hence $(\mathrm{C}_2)$.
For the remaining values $2\le k\le5$,
the two steps preceding the step $k$ in the chosen order have provided $(\mathrm{C}_{k-1})$ and $(\mathrm{C}_k)$.
Let $\xi\in\mathbb{X}_k$, and consider a row at $\xi$.
Its image is a column at a class of $\mathbb{X}_k$, at $\xi^{\downarrow}\in\mathbb{X}_{k-1}$, or at $\xi^{\uparrow}$.
A column of the first kind comes, by $(\mathrm{C}_k)$, from a row at a class of $\mathbb{X}_{k-1}$,
and a column of the second kind comes, by $(\mathrm{C}_{k-1})$, from a row at a class of $\mathbb{X}_{k-2}$, where $\mathbb{X}_0=\mathbb{X}_6$.
The layers $\mathbb{X}_{k-2}$, $\mathbb{X}_{k-1}$ and $\mathbb{X}_k$ are pairwise distinct and the given row stands at $\xi\in\mathbb{X}_k$,
so by injectivity its image is a column at $\xi^{\uparrow}$.
This is $(\mathrm{A}_k)$, and $(\mathrm{C}_{k+1})$ follows.
Once the step $k=5$ is completed, each of $(\mathrm{A}_1),\dots,(\mathrm{A}_6)$ is established.
For every $\xi\in\mathbb{X}_{\mathbb{T}}$ the images of the four rows at $\xi$ are then four distinct columns at $\xi^{\uparrow}$, hence all of them,
and the equation $\mathfrak{s}(\xi,\alpha)=\left(\xi^{\uparrow},\mathfrak{p}_\xi(\alpha)\right)$, $\alpha\in\{1,2,3,4\}$,
defines a permutation $\mathfrak{p}_\xi$ of $\{1,2,3,4\}$.

Conversely, let $(\mathfrak{p}_\xi)_{\xi\in\mathbb{X}_{\mathbb{T}}}$ be a family of permutations of $\{1,2,3,4\}$,
and let $\mathfrak{s}(\xi,\alpha):=\left(\xi^{\uparrow},\mathfrak{p}_\xi(\alpha)\right)$.
Since $\mathfrak{u}$ is a bijection of $\mathbb{X}_{\mathbb{T}}$, the map $\mathfrak{s}$ is a bijection of $\mathbb{I}$;
it is admissible, because $\xi^{\uparrow}$ is the class of a neighbour of the representative of $\xi$.
Let $b\in\mathbb{T}^{(1)}_\mu$.
The lower ends of the bonds over $b$ are classes $\xi\in\mathbb{X}_6$, and the bond at such a $\xi$ joins it to $\xi^{\uparrow}$;
as $\mathfrak{s}$ maps every row at $\xi$ to a column at $\xi^{\uparrow}$, one has $n^{\uparrow}_b(\mathfrak{s})=32$.
The upper ends are classes $\xi\in\mathbb{X}_1$, and the bond at such a $\xi$ joins it to $\xi^{\downarrow}\in\mathbb{X}_6$,
whereas $\mathfrak{s}$ maps the rows at $\xi$ to columns at $\xi^{\uparrow}\in\mathbb{X}_2$;
hence $n^{\downarrow}_b(\mathfrak{s})=0$.
For $b\in\mathbb{T}^{(1)}_\nu$ with $\nu\neq\mu$, the two ends of every bond over $b$ lie in one thin layer,
while $\mathfrak{s}$ maps every row at a class of $\mathbb{X}_k$ to a column at a class of $\mathbb{X}_{k+1}$;
hence $n^{\uparrow}_b(\mathfrak{s})=n^{\downarrow}_b(\mathfrak{s})=0$.
By~\eqref{eq:positivity-exponent}, $\lambda(\mathfrak{s})=\lambda^{(\mu)}$.

Finally, the converse just proved shows that $\lambda^{(\mu)}$ is the exponent of an admissible permutation.
If $\mathfrak{s}$ is admissible and $\lambda_b(\mathfrak{s})=32$ for every $b\in\mathbb{T}^{(1)}_\mu$,
then $\mathfrak{s}$ is of the form treated in the converse, by the description obtained from $(\mathrm{A}_1),\dots,(\mathrm{A}_6)$,
and so $\lambda(\mathfrak{s})=\lambda^{(\mu)}$.
Hence no exponent of an admissible permutation other than $\lambda^{(\mu)}$ takes the value $32$ at every $b\in\mathbb{T}^{(1)}_\mu$.

\emph{Proof of~\eqref{it:extreme-value}.}
In the sum~\eqref{eq:positivity-leibniz}, the term of a permutation that is not admissible vanishes,
and the term of an admissible permutation $\mathfrak{s}$ is a complex multiple of the monomial $u^{\lambda(\mathfrak{s})}$
(Definition~\ref{def:determinant-polynomial}).
Distinct monomials being linearly independent,
the coefficient of $u^{\lambda^{(\mu)}}$ in $P_{\mathbb{T}}$ is the sum of the coefficients of the terms
of the admissible permutations $\mathfrak{s}$ with $\lambda(\mathfrak{s})=\lambda^{(\mu)}$.
By~\eqref{it:extreme-permutations}, these are exactly the permutations
\[
\mathfrak{s}=\mathfrak{s}_0\circ\mathfrak{t},
\qquad
\mathfrak{s}_0(\xi,\alpha):=\left(\xi^{\uparrow},\alpha\right),
\qquad
\mathfrak{t}(\xi,\alpha):=\left(\xi,\mathfrak{p}_\xi(\alpha)\right),
\]
where $\mathfrak{s}_0$ and $\mathfrak{t}$ are permutations of $\mathbb{I}$
and $(\mathfrak{p}_\xi)_{\xi\in\mathbb{X}_{\mathbb{T}}}$ ranges over all families of permutations of $\{1,2,3,4\}$;
distinct families give distinct permutations.

\emph{The sign.}
For each $\xi\in\mathbb{X}_{\mathbb{T}}$, the permutation $\mathfrak{t}$ maps the set $\{\xi\}\times\{1,2,3,4\}$ onto itself and acts there as $\mathfrak{p}_\xi$.
These sets form a partition of $\mathbb{I}$,
and a permutation that maps every block of a partition onto itself has as its sign the product of the signs of its restrictions to the blocks;
hence $\operatorname{sgn}(\mathfrak{t})=\prod_{\xi}\operatorname{sgn}(\mathfrak{p}_\xi)$.
The four sets $\mathbb{X}_{\mathbb{T}}\times\{\alpha\}$, $\alpha\in\{1,2,3,4\}$, form another partition of $\mathbb{I}$;
the permutation $\mathfrak{s}_0$ maps each of them onto itself,
and under the bijection $\xi\mapsto(\xi,\alpha)$ its restriction to $\mathbb{X}_{\mathbb{T}}\times\{\alpha\}$ becomes $\mathfrak{u}$.
Therefore
\[
\operatorname{sgn}(\mathfrak{s}_0)=\operatorname{sgn}(\mathfrak{u})^{4}=1
\qquad\text{and}\qquad
\operatorname{sgn}(\mathfrak{s}_0\circ\mathfrak{t})=\prod_{\xi\in\mathbb{X}_{\mathbb{T}}}\operatorname{sgn}(\mathfrak{p}_\xi).
\]

\emph{The sum.}
Let $\xi\in\mathbb{X}_{\mathbb{T}}$ have the representative $x$.
The site $x^{\uparrow}$ is one of the four neighbours of $x$, and $x^{\uparrow}=\widehat{x^{\uparrow}}+g^{\uparrow}(x)$;
so~\eqref{eq:positivity-blocks} gives the block of $D$ at $\left(\xi,\xi^{\uparrow}\right)$:
\[
D\left[\xi,\xi^{\uparrow}\right]=-\chi\left(g^{\uparrow}(x)\right)\cdot U_A\left(x,x^{\uparrow}\right)\cdot w\left(x,x^{\uparrow}\right).
\]
Here $U_A\left(x,x^{\uparrow}\right)$ is read in $\C\left[u_b,u_b^{-1}:\ b\in\mathbb{T}^{(1)}\right]$ as in Definition~\ref{def:determinant-polynomial}:
by Lemma~\ref{lem:torus-layers}(d) it equals $1$ if $\xi\notin\mathbb{X}_6$,
and the indeterminate $u_b$, with $b$ the class of the link $(\pi(x),\mu)$, if $\xi\in\mathbb{X}_6$.
Let $D\left[\xi,\xi^{\uparrow}\right]_{\alpha,\alpha'}:=D\left[(\xi,\alpha),\left(\xi^{\uparrow},\alpha'\right)\right]$ denote the entries of this block.
By the formula for the sign, the term of $\mathfrak{s}_0\circ\mathfrak{t}$ in~\eqref{eq:positivity-leibniz} is
$\prod_{\xi}\left(\operatorname{sgn}(\mathfrak{p}_\xi)\cdot\prod_{\alpha=1}^{4}D\left[\xi,\xi^{\uparrow}\right]_{\alpha,\mathfrak{p}_\xi(\alpha)}\right)$.
The Leibniz formula for $4\times4$ determinants and the distributive law give
\begin{equation}\label{eq:positivity-up-blocks}
\sum_{(\mathfrak{p}_\xi)_{\xi\in\mathbb{X}_{\mathbb{T}}}}\ \prod_{\xi\in\mathbb{X}_{\mathbb{T}}}\left(\operatorname{sgn}(\mathfrak{p}_\xi)\cdot\prod_{\alpha=1}^{4}D\left[\xi,\xi^{\uparrow}\right]_{\alpha,\mathfrak{p}_\xi(\alpha)}\right)
=\prod_{\xi\in\mathbb{X}_{\mathbb{T}}}\det D\left[\xi,\xi^{\uparrow}\right],
\end{equation}
and
\[
\det D\left[\xi,\xi^{\uparrow}\right]
=(-1)^{4}\cdot\chi\left(g^{\uparrow}(x)\right)^{4}\cdot U_A\left(x,x^{\uparrow}\right)^{4}\cdot\det w\left(x,x^{\uparrow}\right).
\]
The product over $\xi$ of the factors $\chi\left(g^{\uparrow}(x)\right)^{4}$ is $\chi_{\mathbb{T}}$.
For each $b\in\mathbb{T}^{(1)}_\mu$, exactly eight classes $\xi\in\mathbb{X}_6$ have the link $(\pi(x),\mu)$ in $b$
(Lemma~\ref{lem:torus-layers}(d)); hence
\[
\prod_{\xi\in\mathbb{X}_{\mathbb{T}}}U_A\left(x,x^{\uparrow}\right)^{4}
=\prod_{b\in\mathbb{T}^{(1)}_\mu}u_b^{4\cdot8}
=u^{\lambda^{(\mu)}}.
\]
By Lemma~\ref{lem:hopping}, the matrix $w\left(x,x^{\uparrow}\right)$ has the eigenvalues
$\frac{w_i}{2}\cdot\left(\rho-1\right)$ and $\frac{w_i}{2}\cdot\left(\rho+1\right)$, each twice,
where $i=i_{k_\mu(\xi)}$ is the type of the bond $\left\{x,x^{\uparrow}\right\}$ (Lemma~\ref{lem:torus-layers}(b)).
Consequently
\[
\det w\left(x,x^{\uparrow}\right)
=\left(\left(\frac{w_i}{2}\right)^{2}\cdot\left(\rho^{2}-1\right)\right)^{2}
=\left(\frac{w_i}{2}\right)^{4}\cdot\left(1-\rho^{2}\right)^{2}.
\]
As $\#|\mathbb{X}_k|=8|\mathbb{T}|$ for every $k$ (Definition~\ref{def:layer-index}),
the product of these determinants over $\xi\in\mathbb{X}_{\mathbb{T}}$ is
$\prod_{k=1}^{6}\left(\left(\frac{w_{i_k}}{2}\right)^{4}\cdot\left(1-\rho^{2}\right)^{2}\right)^{8|\mathbb{T}|}$.
In the sequence $(i_1,\dots,i_6)=(2,1,0,1,2,3)$ the value $0$ occurs once, $1$ twice, $2$ twice and $3$ once.
The exponents of $\frac{w_0}{2}$, $\frac{w_1}{2}$, $\frac{w_2}{2}$, $\frac{w_3}{2}$ in the product are therefore
$4\cdot8\cdot|\mathbb{T}|$ times $1$, $2$, $2$, $1$, that is, $32|\mathbb{T}|$, $64|\mathbb{T}|$, $64|\mathbb{T}|$, $32|\mathbb{T}|$,
and the exponent of $1-\rho^{2}$ is $2\cdot\#|\mathbb{X}_{\mathbb{T}}|=2\cdot48|\mathbb{T}|=96|\mathbb{T}|$;
the product is $C_*^{|\mathbb{T}|}$.

Altogether, the terms of the admissible permutations $\mathfrak{s}$ with $\lambda(\mathfrak{s})=\lambda^{(\mu)}$ add up to
$\chi_{\mathbb{T}}\cdot C_*^{|\mathbb{T}|}\cdot u^{\lambda^{(\mu)}}$,
and the coefficient has the asserted value, in both of its forms.
Since $|\chi(g)|=1$ for $g\in\Gamma$, one has $\left|\chi_{\mathbb{T}}\right|=1$;
if $\chi$ takes only the values $\pm1$, every factor of $\chi_{\mathbb{T}}$ is $(\pm1)^{4}=1$.
The number $C_*^{|\mathbb{T}|}$ is formed from $w_0,\dots,w_3$, $\rho$ and $|\mathbb{T}|$ alone;
the powers of $\frac{w_0}{2},\dots,\frac{w_3}{2}$ in it are positive,
and the power $\left(1-\rho^{2}\right)^{96|\mathbb{T}|}$ vanishes for $\rho^2=1$ and is positive for $\rho^2\neq1$, its exponent being even.

\emph{Proof of~\eqref{it:extreme-opposite}.}
Now $\xi^{\downarrow}$ plays the role of $\xi^{\uparrow}$ and conversely,
and the layers are visited in the order $1,6,5,4,3,2$.
Let $\mathfrak{s}$ be admissible with $\lambda_b(\mathfrak{s})=-32$ for every $b\in\mathbb{T}^{(1)}_\mu$.
By the bounds of Lemma~\ref{lem:link-degrees}, $n^{\downarrow}_b(\mathfrak{s})=32$ and $n^{\uparrow}_b(\mathfrak{s})=0$ for each such $b$.
As $\mathbb{F}$ contains exactly eight upper ends $x$ of bonds over $b$ (Lemma~\ref{lem:link-degrees}),
the first equality says that each of the $32$ rows at these sites goes to a column at $\widehat{x^{(3)}}$, the representative of the lower end of the bond at $x$;
the second says that no row at a lower end $x\in\mathbb{F}$ of a bond over $b$ goes to a column at $\widehat{x^{(3)}}$.
By the description of the bonds over links parallel to $e_\mu$, this means that
every row at a class $\xi\in\mathbb{X}_1$ is mapped to a column at $\xi^{\downarrow}$,
and that no row at a class $\xi\in\mathbb{X}_6$ is mapped to a column at $\xi^{\uparrow}$.

Suppose that, for some $k$, the rows at every class $\xi\in\mathbb{X}_k$ are known to go to columns at $\xi^{\downarrow}$.
Then each column at a class of $\mathbb{X}_{k-1}$ comes from a row at a class of $\mathbb{X}_k$:
by injectivity, the rows at $\xi$ occupy all four columns at $\xi^{\downarrow}$,
and $\xi\mapsto\xi^{\downarrow}$ maps $\mathbb{X}_k$ onto $\mathbb{X}_{k-1}$.
For $k=1$ this shows that each column at a class of $\mathbb{X}_6$ comes from a row at a class of $\mathbb{X}_1$.
Now let $\xi\in\mathbb{X}_6$, and consider a row at $\xi$.
Its image is a column at a class of $\mathbb{X}_6$, at $\xi^{\uparrow}$, or at $\xi^{\downarrow}$;
injectivity excludes the first case, since $\xi\notin\mathbb{X}_1$, and the hypothesis excludes the second.
So the rows at the classes $\xi\in\mathbb{X}_6$ go to columns at $\xi^{\downarrow}$,
and each column at a class of $\mathbb{X}_5$ comes from a row at a class of $\mathbb{X}_6$.
The layers $\mathbb{X}_5$, $\mathbb{X}_4$, $\mathbb{X}_3$, $\mathbb{X}_2$ are treated next, in this order.
When $\mathbb{X}_k$ is reached, $2\le k\le5$, the two preceding steps have shown that
each column at a class of $\mathbb{X}_{k+1}$ comes from a row at a class of $\mathbb{X}_{k+2}$, where $\mathbb{X}_7=\mathbb{X}_1$,
and each column at a class of $\mathbb{X}_k$ from a row at a class of $\mathbb{X}_{k+1}$.
For $\xi\in\mathbb{X}_k$, the image of a row at $\xi$ is a column at a class of $\mathbb{X}_k$, at $\xi^{\uparrow}\in\mathbb{X}_{k+1}$, or at $\xi^{\downarrow}$;
as $\xi$ belongs neither to $\mathbb{X}_{k+1}$ nor to $\mathbb{X}_{k+2}$, injectivity leaves only $\xi^{\downarrow}$,
and then each column at a class of $\mathbb{X}_{k-1}$ comes from a row at a class of $\mathbb{X}_k$.
At the end, the rows at every class $\xi\in\mathbb{X}_{\mathbb{T}}$ go to columns at $\xi^{\downarrow}$;
by injectivity they fill the four columns at $\xi^{\downarrow}$,
and $\mathfrak{s}(\xi,\alpha)=\left(\xi^{\downarrow},\mathfrak{p}_\xi(\alpha)\right)$ with permutations $\mathfrak{p}_\xi$ of $\{1,2,3,4\}$.

For the converse, let $(\mathfrak{p}_\xi)_{\xi\in\mathbb{X}_{\mathbb{T}}}$ be a family of permutations of $\{1,2,3,4\}$,
and let $\mathfrak{s}(\xi,\alpha):=\left(\xi^{\downarrow},\mathfrak{p}_\xi(\alpha)\right)$.
This is a bijection of $\mathbb{I}$, because $\xi\mapsto\xi^{\downarrow}$ is a bijection of $\mathbb{X}_{\mathbb{T}}$,
and it is admissible, because $\xi^{\downarrow}$ is the class of a neighbour of the representative of $\xi$.
Let $b\in\mathbb{T}^{(1)}_\mu$.
Each row at an upper end $\xi\in\mathbb{X}_1$ of a bond over $b$ goes to a column at $\xi^{\downarrow}$, the lower end of that bond,
whereas each row at a lower end $\xi\in\mathbb{X}_6$ goes to a column at $\xi^{\downarrow}\in\mathbb{X}_5$ and not to one at the upper end $\xi^{\uparrow}\in\mathbb{X}_1$;
thus $n^{\downarrow}_b(\mathfrak{s})=32$ and $n^{\uparrow}_b(\mathfrak{s})=0$.
For $b\in\mathbb{T}^{(1)}_\nu$ with $\nu\neq\mu$ both numbers vanish,
since the bonds over $b$ join classes of one layer and $\mathfrak{s}$ maps the rows at the classes of each layer to columns in another layer.
So $\lambda(\mathfrak{s})=-\lambda^{(\mu)}$.
Combined with the description of the admissible permutations with $\lambda_b(\mathfrak{s})=-32$ on $\mathbb{T}^{(1)}_\mu$ obtained before,
this shows that $-\lambda^{(\mu)}$ is the exponent of an admissible permutation
and that no other exponent of an admissible permutation takes the value $-32$ at every $b\in\mathbb{T}^{(1)}_\mu$.

The coefficient of $u^{-\lambda^{(\mu)}}$ in $P_{\mathbb{T}}$ is therefore the sum of the coefficients of the terms of the permutations
$\mathfrak{s}_0^{-1}\circ\mathfrak{t}$, where $\mathfrak{s}_0^{-1}(\xi,\alpha)=\left(\xi^{\downarrow},\alpha\right)$
and $\mathfrak{t}$ ranges over the permutations considered in the proof of~\eqref{it:extreme-value}.
As $\operatorname{sgn}\left(\mathfrak{s}_0^{-1}\right)=\operatorname{sgn}(\mathfrak{s}_0)=1$,
these terms add up to $\prod_{\xi}\det D\left[\xi,\xi^{\downarrow}\right]$,
by the identity~\eqref{eq:positivity-up-blocks} for the blocks $D\left[\xi,\xi^{\downarrow}\right]$ in place of $D\left[\xi,\xi^{\uparrow}\right]$.
For $x\in\mathbb{F}$ let $g^{\downarrow}(x)\in\Gamma$ be defined by $x^{\downarrow}=\widehat{x^{\downarrow}}+g^{\downarrow}(x)$.
For $\xi\in\mathbb{X}_{\mathbb{T}}$ with representative $x$, the identity~\eqref{eq:positivity-blocks} gives
\[
\begin{aligned}
D\left[\xi,\xi^{\downarrow}\right]&=-\chi\left(g^{\downarrow}(x)\right)\cdot U_A\left(x,x^{\downarrow}\right)\cdot w\left(x,x^{\downarrow}\right),\\
\det D\left[\xi,\xi^{\downarrow}\right]&=\chi\left(g^{\downarrow}(x)\right)^{4}\cdot U_A\left(x,x^{\downarrow}\right)^{4}\cdot\det w\left(x,x^{\downarrow}\right).
\end{aligned}
\]
The three products over $\xi$ are computed as follows.

\emph{Phases.}
Let $\xi\in\mathbb{X}_k$.
Then $U_A\left(x,x^{\downarrow}\right)=U_A\left(\left(x^{\downarrow}\right)^{\uparrow},x^{\downarrow}\right)$ and $k_\mu\left(x^{\downarrow}\right)=k-1$.
Lemma~\ref{lem:torus-layers}(d), applied to the site $x^{\downarrow}$, shows that
this phase is $1$ if $k\neq1$,
and that it is $u_b^{-1}$, with $b$ the class of the link $\left(\pi\left(x^{\downarrow}\right),\mu\right)$, if $k=1$.
The same item shows that for each $b\in\mathbb{T}^{(1)}_\mu$ exactly eight classes of $\mathbb{X}_6$
consist of sites $x'$ with $(\pi(x'),\mu)\in b$;
as $\xi\mapsto\xi^{\downarrow}$ maps $\mathbb{X}_1$ bijectively onto $\mathbb{X}_6$,
exactly eight classes $\xi\in\mathbb{X}_1$ contribute the factor $u_b^{-4}$,
and $\prod_{\xi}U_A\left(x,x^{\downarrow}\right)^{4}=u^{-\lambda^{(\mu)}}$.

\emph{Hopping matrices.}
Let $\xi\in\mathbb{X}_k$.
Since $k_\mu\left(x^{\downarrow}\right)=k-1$ and $\left(x^{\downarrow}\right)^{\uparrow}=x$,
the bond $\left\{x^{\downarrow},x\right\}$ has the type $i_{k-1}$ (Lemma~\ref{lem:torus-layers}(b)),
and Lemma~\ref{lem:hopping} gives $\det w\left(x,x^{\downarrow}\right)=\left(\frac{w_{i_{k-1}}}{2}\right)^{4}\cdot\left(1-\rho^{2}\right)^{2}$.
As $k$ ranges over $\{1,\dots,6\}$, so does $k-1$ modulo $6$;
the product of these determinants over $\xi\in\mathbb{X}_{\mathbb{T}}$ is again $C_*^{|\mathbb{T}|}$.

\emph{Characters.}
The map $x\mapsto x^{\downarrow}$ commutes with the translations by elements of $\Z^3$:
by Lemma~\ref{lem:torus-layers}(c), $\left(x^{\downarrow}+v\right)^{\uparrow}=x+v$ for $v\in\Z^3$,
and applying $\downarrow$ gives $(x+v)^{\downarrow}=x^{\downarrow}+v$.
Let $x\in\mathbb{F}$, and put $\tilde{x}:=\widehat{x^{\uparrow}}$ and $g:=g^{\uparrow}(x)$, so that $x^{\uparrow}=\tilde{x}+g$.
Then $\tilde{x}^{\downarrow}=\left(x^{\uparrow}-g\right)^{\downarrow}=x-g$;
that is, $\widehat{\tilde{x}^{\downarrow}}=x$ and $g^{\downarrow}(\tilde{x})=-g^{\uparrow}(x)$.
The map $x\mapsto\widehat{x^{\uparrow}}$ is a bijection of $\mathbb{F}$, since $\mathfrak{u}$ is a bijection of $\mathbb{X}_{\mathbb{T}}$,
and $\chi(-g)=\overline{\chi(g)}$ because $\chi$ is unitary.
Therefore
\[
\prod_{\tilde{x}\in\mathbb{F}}\chi\left(g^{\downarrow}(\tilde{x})\right)^{4}
=\prod_{x\in\mathbb{F}}\overline{\chi\left(g^{\uparrow}(x)\right)}^{\,4}
=\overline{\chi_{\mathbb{T}}}.
\]

Collecting the three products, one obtains for the coefficient of $u^{-\lambda^{(\mu)}}$ the value
$\overline{\chi_{\mathbb{T}}}\cdot\prod_{\xi}\det w\left(x,x^{\downarrow}\right)=\overline{\chi_{\mathbb{T}}}\cdot C_*^{|\mathbb{T}|}$.
If $\chi$ takes only the values $\pm1$, then $\overline{\chi_{\mathbb{T}}}=\chi_{\mathbb{T}}=1$ by~\eqref{it:extreme-value},
and the coefficient is $C_*^{|\mathbb{T}|}$.
\end{proof}

\subsection{From the coefficient to the partition function}\label{sec:from-coefficient}

\begin{definition}[The torus of phases]\label{def:phase-torus}
Let $n$ be a positive integer.
The \emph{torus of phases} is the compact abelian group
\[
(S^1)^n:=\{u=(u_1,\dots,u_n)\in\C^n:\ |u_j|=1\ \text{for}\ j=1,\dots,n\},
\]
with coordinatewise multiplication.
For $t=(t_1,\dots,t_n)\in\R^n$ put $\exp(it):=(\exp(it_1),\dots,\exp(it_n))$, a point of $(S^1)^n$.
The normalised Haar measure $\mathrm{d}u$ of $(S^1)^n$ is given by
\[
\int\mathrm{d}u\,f(u)=(2\pi)^{-n}\int_{[0,2\pi)^n}\mathrm{d}t\,f(\exp(it))
\]
for every Borel function $f$ on $(S^1)^n$ that is bounded or non-negative.
For $\lambda=(\lambda_1,\dots,\lambda_n)\in\Z^n$ put $u^\lambda:=\prod_{j=1}^{n}u_j^{\lambda_j}$.
A \emph{Laurent polynomial in $n$ variables} is a finite linear combination
\[
P=\sum_{\lambda\in\Z^n}a_\lambda u^\lambda
\]
of these monomials with complex coefficients $a_\lambda$, all but finitely many of them zero,
that is, an element of the ring $\C[u_1,u_1^{-1},\dots,u_n,u_n^{-1}]$;
$a_\lambda$ is the coefficient of $u^\lambda$ in $P$, $a_0$ is the constant coefficient,
and $P\neq0$ means that some coefficient of $P$ is different from zero.
Evaluation at the points of $(S^1)^n$ makes $P$ a continuous function on $(S^1)^n$.
For $\lambda,\lambda'\in\Z^n$ one has
\[
\int\mathrm{d}u\,u^\lambda\cdot\overline{u^{\lambda'}}=
\begin{cases}
1 & \text{if } \lambda=\lambda',\\
0 & \text{if } \lambda\neq\lambda'.
\end{cases}
\]
Indeed, at $u=\exp(it)$ the integrand is the product of the functions $\exp(i(\lambda_j-\lambda'_j)t_j)$, $j=1,\dots,n$,
so that by Fubini's theorem the integral is the product of the $n$ integrals
$(2\pi)^{-1}\int_0^{2\pi}\mathrm{d}t_j\,\exp(i(\lambda_j-\lambda'_j)t_j)$,
each of which equals $1$ if $\lambda_j=\lambda'_j$ and $0$ otherwise.
Integrating the finite sum $P(u)\cdot\overline{u^\lambda}=\sum_{\lambda'}a_{\lambda'}u^{\lambda'}\cdot\overline{u^\lambda}$ term by term,
one obtains for every Laurent polynomial $P=\sum_{\lambda}a_\lambda u^\lambda$ in $n$ variables
\begin{equation}\label{eq:positivity-fourier}
a_\lambda=\int\mathrm{d}u\,P(u)\cdot\overline{u^\lambda}\qquad\text{for every }\lambda\in\Z^n .
\end{equation}
A Laurent polynomial is therefore determined by its values on $(S^1)^n$:
the functions $u\mapsto u^\lambda$, $\lambda\in\Z^n$, on $(S^1)^n$ are linearly independent,
and a Laurent polynomial whose values on $(S^1)^n$ all vanish is the zero polynomial.
For variables indexed by a finite set of $n$ elements other than $\{1,\dots,n\}$,
such as the link phases $u_b$, $b\in\mathbb{T}^{(1)}$, of a torus $\mathbb{T}$,
the same definitions apply after an enumeration of that set.
\end{definition}

\begin{lemma}[Non-negative trigonometric polynomials]\label{lem:nonnegative-trigonometric}
Let $n$ be a positive integer and $P=\sum_{\lambda}a_\lambda u^\lambda$ a Laurent polynomial in $n$ variables.
\begin{enumerate}
\item\label{it:fourier-bound}
If $P(u)\ge0$ for every $u\in(S^1)^n$, then
\[
|a_\lambda|\le a_0=\int\mathrm{d}u\,P(u)\qquad\text{for every }\lambda\in\Z^n .
\]
\item\label{it:null-zero-set}
If $P\neq0$, then the set
\[
\{t\in[0,2\pi)^n:\ P(\exp(it))=0\}
\]
is a Borel set of Lebesgue measure zero.
\end{enumerate}
\end{lemma}

\begin{proof}
\eqref{it:fourier-bound}
On $(S^1)^n$ one has $|u^\lambda|=1$ and, by hypothesis, $|P(u)|=P(u)$.
Formula \eqref{eq:positivity-fourier}, applied to $\lambda$ and to $0$, therefore gives
\[
|a_\lambda|=\left|\int\mathrm{d}u\,P(u)\cdot\overline{u^\lambda}\right|
\le\int\mathrm{d}u\,|P(u)|\cdot|u^\lambda|
=\int\mathrm{d}u\,P(u)=a_0 .
\]

\eqref{it:null-zero-set}
Denote the set in question by $\mathcal{N}$.
The function $t\mapsto P(\exp(it))$ is continuous on $\R^n$;
hence $\mathcal{N}$ is the intersection of a closed subset of $\R^n$ with $[0,2\pi)^n$, and in particular a Borel set.
That $\mathcal{N}$ has Lebesgue measure zero is proved by induction on $n$.

Let $n=1$.
The product of $P$ with a suitable power of $u_1$ is a polynomial in $u_1$ with the same coefficients as $P$, so it is not the zero polynomial,
and it has only finitely many zeros in $\C$.
Since $u_1\neq0$ on $S^1$ and $t_1\mapsto\exp(it_1)$ is injective on $[0,2\pi)$, the set $\mathcal{N}$ is finite.

Now let $n\ge2$, and assume \eqref{it:null-zero-set} for Laurent polynomials in $n-1$ variables.
Write $u=(u',u_n)$ with $u'=(u_1,\dots,u_{n-1})$, and accordingly $t=(t',t_n)$.
Collecting the monomials of $P$ according to their power of $u_n$, one has
\[
P=\sum_{j}P_j\cdot u_n^{j},
\]
a finite sum over $j\in\Z$, in which every $P_j$ is a Laurent polynomial in the $n-1$ variables $u'$.
As $P\neq0$, there is an index $j_*$ with $P_{j_*}\neq0$.
By the induction hypothesis the set
\[
\mathcal{N}':=\{t'\in[0,2\pi)^{n-1}:\ P_{j_*}(\exp(it'))=0\}
\]
is a Borel set of Lebesgue measure zero.
Fix $t'\in[0,2\pi)^{n-1}$ outside $\mathcal{N}'$.
The Laurent polynomial $\sum_jP_j(\exp(it'))\cdot u_n^{j}$ in the one variable $u_n$
takes at $\exp(it_n)$ the value $P(\exp(it))$, $t=(t',t_n)$,
and its coefficient of $u_n^{j_*}$ is $P_{j_*}(\exp(it'))\neq0$;
by the case $n=1$, the section $\{t_n\in[0,2\pi):\ (t',t_n)\in\mathcal{N}\}$ is therefore a finite set.
Tonelli's theorem, applied to the indicator function of $\mathcal{N}$, a non-negative Borel function on $[0,2\pi)^{n-1}\times[0,2\pi)$,
expresses the Lebesgue measure of $\mathcal{N}$ as the integral over $t'\in[0,2\pi)^{n-1}$ of the Lebesgue measure of the section at $t'$.
This measure is at most $2\pi$ for $t'\in\mathcal{N}'$ and zero for every other $t'$.
So the Lebesgue measure of $\mathcal{N}$ is at most $2\pi$ times that of $\mathcal{N}'$, which is zero.
\end{proof}

\begin{definition}[The functional of a general torus]\label{def:general-functional}
Let $\mathbb{T}=(\Gamma,\chi)$ be a torus as in Definition~\ref{def:general-tori}.
The space $\Conf(\Gamma)$ and its measure are given by the formulas of Definition~\ref{def:configurations},
which make sense for every subgroup $\Gamma$ of $\Z^3$ of rank three.
For a Borel function $\Phi\ge0$ on $\Conf(\Gamma)$, the number $Z_{\mathbb{T}}[\Phi]$ is defined by the formula of Definition~\ref{def:functional},
as the integral of $\exp\left(-\frac12\|F\|^2\right)\cdot\det D^{\mathbb{T}}_A\cdot\Phi$, with the curvature $F=dA+\frac{2\pi}{q}z$, against this measure.
Written out as a sum over the closed flux variables, it is
\[
Z_{\mathbb{T}}[\Phi]=\sum_{\substack{z\in C^2(\Gamma;\Z)\\ dz=0}}\ \int_{[0,2\pi/q)^{\mathbb{T}^{(1)}}}\mathrm{d}A\,
\exp\left(-\frac12\left\|dA+\frac{2\pi}{q}z\right\|^2\right)\cdot\det D^{\mathbb{T}}_A\cdot\Phi(A,z),
\]
where $\Phi(A,z)$ is the value of $\Phi$ at the class of $(A,z)$,
and where, for a real $\Gamma$-periodic two-cochain $F$, $\|F\|^2$ is the sum of the squares $F_\sigma^2$
over the $3|\mathbb{T}|$ plaquettes $\sigma$ of $\Z^3/\Gamma$ (three for each class of photon sites).
The term of a closed $z$ in this sum is the \emph{flux term} of $z$.
Each flux term is the integral of a non-negative Borel function of $A$:
$\Phi\ge0$, and $\det D^{\mathbb{T}}_A\ge0$ for every real $A$ by Proposition~\ref{prop:accretivity}, because $M$ is real.
So the sum has a value $Z_{\mathbb{T}}[\Phi]\in[0,\infty]$, which does not depend on the order of summation.
\end{definition}

With $\Phi=1$ and the factor $\det D^{\mathbb{T}}_A$ replaced by $1$, the same sum is the total mass of the Maxwell weight
$\exp\left(-\frac12\|F\|^2\right)$ on $\Conf(\Gamma)$.
Lemma~\ref{lem:bounded-integrands}, applied to the subgroup $\Gamma$, bounds it:
\[
0<\sum_{\substack{z\in C^2(\Gamma;\Z)\\ dz=0}}\ \int_{[0,2\pi/q)^{\mathbb{T}^{(1)}}}\mathrm{d}A\,
\exp\left(-\frac12\left\|dA+\frac{2\pi}{q}z\right\|^2\right)
\le\left(\frac{2\pi}{q}\cdot C_q\right)^{3|\mathbb{T}|},
\]
with the constant $C_q$ of that lemma.

\begin{lemma}[Positivity of the flux terms]\label{lem:flux-terms}
Let $\rho$ and $M$ be real, $w_0,\dots,w_3>0$ and $q>0$,
let $\Gamma$ be a subgroup of $\Z^3$ of rank three, $\chi$ a unitary character of $\Gamma$,
and let $\mathbb{T}=(\Gamma,\chi)$, with the operator $D^{\mathbb{T}}_A$ and the functional $Z_{\mathbb{T}}$ of Definitions~\ref{def:general-tori} and~\ref{def:general-functional}.
Assume that the function $A\mapsto\det D^{\mathbb{T}}_A$ on $\R^{\mathbb{T}^{(1)}}$ is not identically zero.
\begin{enumerate}
\item\label{it:flux-positive}
For every closed $z\in C^2(\Gamma;\Z)$ the flux term of $z$ is positive:
\[
\int_{[0,2\pi/q)^{\mathbb{T}^{(1)}}}\mathrm{d}A\,\exp\left(-\frac12\left\|dA+\frac{2\pi}{q}\cdot z\right\|^2\right)\cdot\det D^{\mathbb{T}}_A>0,
\]
where $\|\cdot\|^2$ is the sum of the squares over the $3|\mathbb{T}|$ plaquettes of $\Z^3/\Gamma$.
Since $Z_{\mathbb{T}}[1]$ is the sum of the flux terms over all closed $z$ (Definition~\ref{def:general-functional}),
it follows that $Z_{\mathbb{T}}[1]>0$.
\item\label{it:flux-null-set}
The set of the zeros of the determinant,
\[
\mathcal{N}_{\mathbb{T}}:=\left\{A\in\R^{\mathbb{T}^{(1)}}:\ \det D^{\mathbb{T}}_A=0\right\},
\]
is a closed Lebesgue null set in $\R^{\mathbb{T}^{(1)}}$.
\item\label{it:flux-equivalence}
Let $\Phi\ge0$ be a Borel function on $\Conf(\Gamma)$.
Then $Z_{\mathbb{T}}[\Phi]=0$ if and only if $\{\Phi\neq0\}$ is a null set for the measure of $\Conf(\Gamma)$ (Definition~\ref{def:configurations}).
Consequently the two measures on $\Conf(\Gamma)$ with the densities
$\exp\left(-\frac12\|F\|^2\right)\cdot\det D^{\mathbb{T}}_A$ and $\exp\left(-\frac12\|F\|^2\right)$,
the second one being the Maxwell measure, have the same null sets.
Moreover $Z_{\mathbb{T}}[1]<\infty$,
and the law $\mu_{\mathbb{T}}$ of $\widetilde F$ under the probability measure $Z_{\mathbb{T}}[\,\cdot\,]/Z_{\mathbb{T}}[1]$ (Definition~\ref{def:functional})
and the law of $\widetilde F$ under the normalised Maxwell measure are mutually absolutely continuous.
\end{enumerate}
\end{lemma}

\begin{proof}
By Definition~\ref{def:determinant-polynomial} there is a Laurent polynomial $P_{\mathbb{T}}=\sum_\lambda a_\lambda u^\lambda$
in the indeterminates $u_b$, $b\in\mathbb{T}^{(1)}$, such that
\[
\det D^{\mathbb{T}}_A=P_{\mathbb{T}}\left(\left(\exp(iqA(b))\right)_{b\in\mathbb{T}^{(1)}}\right)
\qquad\left(A\in\R^{\mathbb{T}^{(1)}}\right).
\]
Three properties of the determinant can be read off from this formula.
It is a continuous function of $A$.
It does not change when $A$ is replaced by $A+\frac{2\pi}{q}\cdot k$ with $k\in\Z^{\mathbb{T}^{(1)}}$.
It is bounded by $\sum_\lambda|a_\lambda|$, since every monomial $u^\lambda$ has modulus one on $(S^1)^{\mathbb{T}^{(1)}}$.
Besides, it is non-negative by Proposition~\ref{prop:accretivity}, because $M$, $\rho$ and $A$ are real and $\chi$ is unitary.
Since the determinant does not vanish identically, $P_{\mathbb{T}}\neq0$.

\emph{Positivity of the flux terms.}
Choose $A_0\in\R^{\mathbb{T}^{(1)}}$ with $\det D^{\mathbb{T}}_{A_0}>0$.
A translation by an element of $\frac{2\pi}{q}\Z^{\mathbb{T}^{(1)}}$ leaves the determinant unchanged and brings $A_0$ into $[0,2\pi/q)^{\mathbb{T}^{(1)}}$;
so one may assume that $A_0$ lies there.
By continuity there is $\varepsilon>0$ such that $\det D^{\mathbb{T}}_A>0$ whenever $|A(b)-A_0(b)|<\varepsilon$ for all $b\in\mathbb{T}^{(1)}$.
The set of the $A$ with
\[
A_0(b)<A(b)<\min\left(A_0(b)+\varepsilon,\frac{2\pi}{q}\right)\qquad\left(b\in\mathbb{T}^{(1)}\right)
\]
is a non-empty open box contained in $(0,2\pi/q)^{\mathbb{T}^{(1)}}$, on which the determinant is positive;
it has positive Lebesgue measure.
Let $z$ be closed.
The integrand of the flux term of $z$ is non-negative on $[0,2\pi/q)^{\mathbb{T}^{(1)}}$,
and on this open box it is positive and continuous, its exponential factor being positive everywhere.
So the flux term of $z$ is positive.
All terms of the sum of Definition~\ref{def:general-functional} for $Z_{\mathbb{T}}[1]$ are non-negative;
hence $Z_{\mathbb{T}}[1]$ is at least the flux term of $z=0$, and $Z_{\mathbb{T}}[1]>0$.

\emph{The zero set.}
The set $\mathcal{N}_{\mathbb{T}}$ is closed, being the zero set of a continuous function.
As $P_{\mathbb{T}}\neq0$, Lemma~\ref{lem:nonnegative-trigonometric}\eqref{it:null-zero-set} shows that
\[
\left\{t\in[0,2\pi)^{\mathbb{T}^{(1)}}:\ P_{\mathbb{T}}\left(\left(\exp(it_b)\right)_{b\in\mathbb{T}^{(1)}}\right)=0\right\}
\]
is a Borel set of Lebesgue measure zero.
The linear bijection $t\mapsto t/q$ maps it onto $\mathcal{N}_{\mathbb{T}}\cap[0,2\pi/q)^{\mathbb{T}^{(1)}}$, which is therefore a null set.
Now $\R^{\mathbb{T}^{(1)}}$ is the union of the countably many translates $[0,2\pi/q)^{\mathbb{T}^{(1)}}+\frac{2\pi}{q}\cdot k$, $k\in\Z^{\mathbb{T}^{(1)}}$,
and $\mathcal{N}_{\mathbb{T}}$ is invariant under the translations by $\frac{2\pi}{q}\cdot k$.
Thus $\mathcal{N}_{\mathbb{T}}$ is the union of countably many translates of a null set, and it is a null set.

\emph{Null sets and laws.}
Let $\Phi\ge0$ be a Borel function on $\Conf(\Gamma)$, and write $\Phi(A,z)$ for its value at the class of $(A,z)$.
Definition~\ref{def:general-functional} gives
\[
Z_{\mathbb{T}}[\Phi]=\sum_{z}\int_{[0,2\pi/q)^{\mathbb{T}^{(1)}}}\mathrm{d}A\,\exp\left(-\frac12\left\|dA+\frac{2\pi}{q}\cdot z\right\|^2\right)\cdot\det D^{\mathbb{T}}_A\cdot\Phi(A,z),
\]
the sum being taken over the closed $z\in C^2(\Gamma;\Z)$; each term is the integral of a non-negative Borel function.
Hence $Z_{\mathbb{T}}[\Phi]$ is zero if and only if each of these integrands vanishes almost everywhere on $[0,2\pi/q)^{\mathbb{T}^{(1)}}$.
In the integrand of a closed $z$ the exponential factor is positive at every point,
and the determinant, being non-negative, is positive off the null set $\mathcal{N}_{\mathbb{T}}$;
so that integrand vanishes almost everywhere on the box exactly when $\Phi(\,\cdot\,,z)$ does.
Thus $Z_{\mathbb{T}}[\Phi]=0$ if and only if, for each closed $z$, the set of the $A\in[0,2\pi/q)^{\mathbb{T}^{(1)}}$ with $\Phi(A,z)\neq0$ has Lebesgue measure zero.
By Definition~\ref{def:configurations}, the measure of the subset $\{\Phi\neq0\}$ of $\Conf(\Gamma)$ is the sum, over the closed $z$, of the Lebesgue measures of these sets;
the condition just obtained is therefore equivalent to $\{\Phi\neq0\}$ being a null set.
The same reasoning, with the determinant replaced by the constant $1$, shows that the integral of $\Phi$ against the Maxwell measure vanishes if and only if $\Phi=0$ almost everywhere.
Taking for $\Phi$ the indicator functions of the Borel subsets of $\Conf(\Gamma)$, one sees that the two measures have the same null sets.

By Lemma~\ref{lem:bounded-integrands} the total mass of the Maxwell measure is positive and at most $\left(\frac{2\pi}{q}\cdot C_q\right)^{3|\mathbb{T}|}$,
with the constant $C_q<\infty$ of that lemma.
As the determinant is bounded by $\sum_\lambda|a_\lambda|$,
\[
Z_{\mathbb{T}}[1]\le\sum_\lambda|a_\lambda|\cdot\left(\frac{2\pi}{q}\cdot C_q\right)^{3|\mathbb{T}|}<\infty ,
\]
while $Z_{\mathbb{T}}[1]>0$ by item~\eqref{it:flux-positive}.
Hence $Z_{\mathbb{T}}[\,\cdot\,]/Z_{\mathbb{T}}[1]$ and the normalised Maxwell measure are probability measures on $\Conf(\Gamma)$;
they are positive multiples of the two measures above, and so they have the same null sets.
The same is then true of their images under the Borel map $(A,z)\mapsto\widetilde F$ from $\Conf(\Gamma)$ into the space $\Omega$ of closed real two-cochains on $\Z^3$:
a Borel subset of $\Omega$ is a null set for one image if and only if its preimage is a null set for the corresponding measure on $\Conf(\Gamma)$.
These images are $\mu_{\mathbb{T}}$ and the law of $\widetilde F$ under the normalised Maxwell measure.
\end{proof}

\begin{theorem}[The partition function for $\rho^2\neq1$]\label{thm:floor}
Let $\rho$ be a real number with $\rho^2\neq1$, let $M$ be real, $w_0,\dots,w_3>0$ and $q>0$,
let $\Gamma$ be a subgroup of $\Z^3$ of rank three, $\chi$ a unitary character of $\Gamma$, and $\mathbb{T}=(\Gamma,\chi)$.
Let
\[
C_*=\left(\frac{w_0}{2}\right)^{32}\cdot\left(\frac{w_1}{2}\right)^{64}\cdot\left(\frac{w_2}{2}\right)^{64}\cdot\left(\frac{w_3}{2}\right)^{32}\cdot\left(1-\rho^2\right)^{96}
\]
be the number of Proposition~\ref{prop:extreme-coefficient}; it is positive.
Then the following hold.
\begin{enumerate}
\item\label{it:floor-polynomial}
The function $A\mapsto\det D^{\mathbb{T}}_A$ on $\R^{\mathbb{T}^{(1)}}$ is a trigonometric polynomial in the variables $qA(b)$, $b\in\mathbb{T}^{(1)}$;
it is non-negative and not identically zero.
\item\label{it:floor-mean}
The mean of the determinant over the box $[0,2\pi/q)^{\mathbb{T}^{(1)}}$ satisfies
\[
\left(\frac{2\pi}{q}\right)^{-3|\mathbb{T}|}\int_{[0,2\pi/q)^{\mathbb{T}^{(1)}}}\mathrm{d}A\,\det D^{\mathbb{T}}_A\ \ge\ C_*^{|\mathbb{T}|}\ >\ 0 .
\]
\item\label{it:floor-bounds}
The partition function satisfies the lower bound
\begin{equation}\label{eq:positivity-floor}
\left(\frac{2\pi}{q}\cdot\exp\left(-2\left(\frac{2\pi}{q}\right)^2\right)\right)^{3|\mathbb{T}|}\cdot C_*^{|\mathbb{T}|}\ \le\ Z_{\mathbb{T}}[1].
\end{equation}
For every real $\rho$, the values $\rho=\pm1$ included, and with $M$, $w_0,\dots,w_3$, $q$ and $\mathbb{T}$ as above, it satisfies the upper bound
\begin{equation}\label{eq:positivity-ceiling}
Z_{\mathbb{T}}[1]\ \le\ \left(\frac{2\pi}{q}\cdot C_q\right)^{3|\mathbb{T}|}\cdot\left(|M|+\frac{1+|\rho|}{2}\sum_{i=0}^{3}w_i\right)^{192|\mathbb{T}|},
\end{equation}
where
\[
C_q=\sup_{y\in\R}\ \sum_{k\in\Z}\exp\left(-\frac12\left(y+\frac{2\pi k}{q}\right)^2\right)<\infty
\]
is the constant of Lemma~\ref{lem:bounded-integrands}.
For $\rho^2\neq1$ the two bounds together give $0<Z_{\mathbb{T}}[1]<\infty$,
and the lower bound gives, for every $\mathbb{T}$ as above,
\[
\frac{1}{|\mathbb{T}|}\log Z_{\mathbb{T}}[1]\ \ge\ \log C_*+3\log\frac{2\pi}{q}-6\left(\frac{2\pi}{q}\right)^2 .
\]
\item\label{it:floor-flux-terms}
Every flux term of $Z_{\mathbb{T}}[1]$ is bounded below explicitly: for every closed $z\in C^2(\Gamma;\Z)$,
\begin{multline*}
\int_{[0,2\pi/q)^{\mathbb{T}^{(1)}}}\mathrm{d}A\,\exp\left(-\frac12\left\|dA+\frac{2\pi}{q}\cdot z\right\|^2\right)\cdot\det D^{\mathbb{T}}_A\\
\ge\left(\frac{2\pi}{q}\right)^{3|\mathbb{T}|}\cdot C_*^{|\mathbb{T}|}\cdot\exp\left(-\frac12\left(\sqrt{12\,|\mathbb{T}|}\cdot\frac{2\pi}{q}+\frac{2\pi}{q}\cdot\|z\|\right)^2\right)>0 .
\end{multline*}
\item\label{it:floor-equivalence}
The zero set $\mathcal{N}_{\mathbb{T}}=\left\{A\in\R^{\mathbb{T}^{(1)}}:\det D^{\mathbb{T}}_A=0\right\}$ of the determinant is a closed Lebesgue null set.
A Borel function $\Phi\ge0$ on $\Conf(\Gamma)$ satisfies $Z_{\mathbb{T}}[\Phi]=0$ precisely when it vanishes outside a null set of $\Conf(\Gamma)$.
Hence the measures on $\Conf(\Gamma)$ with the densities $\exp\left(-\frac12\|F\|^2\right)\cdot\det D^{\mathbb{T}}_A$ and $\exp\left(-\frac12\|F\|^2\right)$,
the latter being the Maxwell measure summed over all closed $z$, have the same null sets,
and $\mu_{\mathbb{T}}$ is mutually absolutely continuous with the law of $\widetilde F$ under the normalised Maxwell measure.
\end{enumerate}
\end{theorem}

\begin{proof}
\emph{Item~\eqref{it:floor-polynomial}.}
Write $P_{\mathbb{T}}=\sum_\lambda a_\lambda u^\lambda$ for the Laurent polynomial of Definition~\ref{def:determinant-polynomial}. Then
\[
\begin{aligned}
\det D^{\mathbb{T}}_A&=P_{\mathbb{T}}\left(\left(\exp(iqA(b))\right)_{b\in\mathbb{T}^{(1)}}\right)\\
&=\sum_\lambda a_\lambda\cdot\exp\left(i\sum_{b\in\mathbb{T}^{(1)}}\lambda_b\cdot qA(b)\right),
\end{aligned}
\]
a trigonometric polynomial in the variables $qA(b)$.
Its values are non-negative by Proposition~\ref{prop:accretivity}, as $M$, $\rho$ and $A$ are real and $\chi$ is unitary.
Proposition~\ref{prop:extreme-coefficient}\eqref{it:extreme-value}, applied in the direction $\mu=1$,
gives $a_{\lambda^{(1)}}=\chi_{\mathbb{T}}\cdot C_*^{|\mathbb{T}|}$ with $|\chi_{\mathbb{T}}|=1$; thus
\[
\left|a_{\lambda^{(1)}}\right|=C_*^{|\mathbb{T}|}>0,
\]
because $\rho^2\neq1$ and the weights are positive.
In particular $P_{\mathbb{T}}\neq0$.
A Laurent polynomial that vanishes on $(S^1)^{\mathbb{T}^{(1)}}$ is zero (Definition~\ref{def:phase-torus});
so $P_{\mathbb{T}}(u)\neq0$ for some $u\in(S^1)^{\mathbb{T}^{(1)}}$.
Since $q\cdot2\pi/q=2\pi$, the map $A\mapsto\left(\exp(iqA(b))\right)_{b\in\mathbb{T}^{(1)}}$ sends the box $[0,2\pi/q)^{\mathbb{T}^{(1)}}$ onto $(S^1)^{\mathbb{T}^{(1)}}$,
and the determinant does not vanish identically.

\emph{Item~\eqref{it:floor-mean}.}
The set $\mathbb{T}^{(1)}$ has $3|\mathbb{T}|$ elements (Definition~\ref{def:general-tori}).
The substitution $t_b=qA(b)$, $b\in\mathbb{T}^{(1)}$, maps the box $[0,2\pi/q)^{\mathbb{T}^{(1)}}$ onto $[0,2\pi)^{\mathbb{T}^{(1)}}$, and $\mathrm{d}A=q^{-3|\mathbb{T}|}\,\mathrm{d}t$.
With the normalised Haar measure $\mathrm{d}u$ of Definition~\ref{def:phase-torus} one obtains
\[
\begin{aligned}
&\left(\frac{2\pi}{q}\right)^{-3|\mathbb{T}|}\int_{[0,2\pi/q)^{\mathbb{T}^{(1)}}}\mathrm{d}A\,\det D^{\mathbb{T}}_A\\
&\qquad=(2\pi)^{-3|\mathbb{T}|}\int_{[0,2\pi)^{\mathbb{T}^{(1)}}}\mathrm{d}t\,P_{\mathbb{T}}\left(\left(\exp(it_b)\right)_{b\in\mathbb{T}^{(1)}}\right)\\
&\qquad=\int\mathrm{d}u\,P_{\mathbb{T}}(u)=a_0 .
\end{aligned}
\]
Every point of $(S^1)^{\mathbb{T}^{(1)}}$ has the form $\left(\exp(iqA(b))\right)_b$ with a real $A$,
so $P_{\mathbb{T}}\ge0$ on $(S^1)^{\mathbb{T}^{(1)}}$ by item~\eqref{it:floor-polynomial}.
Lemma~\ref{lem:nonnegative-trigonometric}\eqref{it:fourier-bound} and the value of $|a_{\lambda^{(1)}}|$ found above give
$a_0\ge\left|a_{\lambda^{(1)}}\right|=C_*^{|\mathbb{T}|}>0$.

\emph{Item~\eqref{it:floor-flux-terms} and the lower bound~\eqref{eq:positivity-floor}.}
Let $A\in[0,2\pi/q)^{\mathbb{T}^{(1)}}$, and let $\sigma$ be a plaquette of $\Z^3/\Gamma$, the class of a plaquette $(x;\mu\nu)$ of $\Z^3$.
Definition~\ref{def:cochains} gives
\[
(dA)_\sigma=\left(A_\mu(x)+A_\nu(x+e_\mu)\right)-\left(A_\mu(x+e_\nu)+A_\nu(x)\right),
\]
and each of the four values of $A$ on the right is a coordinate $A(b)$, $b\in\mathbb{T}^{(1)}$, so that it lies in $[0,2\pi/q)$.
Both brackets thus lie in $\left[0,2\cdot\frac{2\pi}{q}\right)$, and $|(dA)_\sigma|<2\cdot\frac{2\pi}{q}$.
Squaring and summing over the $3|\mathbb{T}|$ plaquettes gives $\|dA\|^2<3|\mathbb{T}|\cdot4\left(\frac{2\pi}{q}\right)^2=12\left(\frac{2\pi}{q}\right)^2\cdot|\mathbb{T}|$,
and for every $z\in C^2(\Gamma;\Z)$ the triangle inequality yields
\[
\left\|dA+\frac{2\pi}{q}\cdot z\right\|\le\sqrt{12\,|\mathbb{T}|}\cdot\frac{2\pi}{q}+\frac{2\pi}{q}\cdot\|z\|
\]
on the whole box.
As the determinant is non-negative, the flux term of a closed $z$ is at least
\[
\exp\left(-\frac12\left(\sqrt{12\,|\mathbb{T}|}\cdot\frac{2\pi}{q}+\frac{2\pi}{q}\cdot\|z\|\right)^2\right)\cdot\int_{[0,2\pi/q)^{\mathbb{T}^{(1)}}}\mathrm{d}A\,\det D^{\mathbb{T}}_A ,
\]
and item~\eqref{it:floor-mean} turns this into item~\eqref{it:floor-flux-terms}.
At $z=0$ the bound of item~\eqref{it:floor-flux-terms} equals
\[
\begin{aligned}
&\left(\frac{2\pi}{q}\right)^{3|\mathbb{T}|}\cdot C_*^{|\mathbb{T}|}\cdot\exp\left(-6\left(\frac{2\pi}{q}\right)^2\cdot|\mathbb{T}|\right)\\
&\qquad=\left(\frac{2\pi}{q}\cdot\exp\left(-2\left(\frac{2\pi}{q}\right)^2\right)\right)^{3|\mathbb{T}|}\cdot C_*^{|\mathbb{T}|}.
\end{aligned}
\]
By Definition~\ref{def:general-functional}, $Z_{\mathbb{T}}[1]$ is the sum of the flux terms over the closed $z$, all of them non-negative;
the term of $z=0$ alone gives~\eqref{eq:positivity-floor}.

\emph{The upper bound~\eqref{eq:positivity-ceiling}.}
In this paragraph $\rho$ is an arbitrary real number, and $M$, $w_0,\dots,w_3$, $q$, $\mathbb{T}$ are as in the statement.
Let $A\in\R^{\mathbb{T}^{(1)}}$.
The spinors of $\mathbb{T}$ form a Hilbert space of dimension $4\cdot\#|\mathbb{X}_{\mathbb{T}}|=192|\mathbb{T}|$ (Definition~\ref{def:general-tori}),
on which $D^{\mathbb{T}}_A$ is a linear operator.
As for every operator on a finite-dimensional Hilbert space, $\left|\det D^{\mathbb{T}}_A\right|$ is the product of all $192|\mathbb{T}|$ singular values of $D^{\mathbb{T}}_A$, counted with multiplicity,
and every singular value is bounded by the operator norm $\left\|D^{\mathbb{T}}_A\right\|$; so
\[
\left|\det D^{\mathbb{T}}_A\right|\le\left\|D^{\mathbb{T}}_A\right\|^{192|\mathbb{T}|}.
\]
Proposition~\ref{prop:accretivity} shows that $\det D^{\mathbb{T}}_A\ge0$, so that the determinant coincides with its modulus,
and that $\left\|D^{\mathbb{T}}_A\right\|\le|M|+\frac{1+|\rho|}{2}\sum_iw_i$.
Hence
\[
0\le\det D^{\mathbb{T}}_A\le\left(|M|+\frac{1+|\rho|}{2}\sum_{i=0}^{3}w_i\right)^{192|\mathbb{T}|}
\qquad\left(A\in\R^{\mathbb{T}^{(1)}}\right).
\]
Inserted into the formula of Definition~\ref{def:general-functional} for $Z_{\mathbb{T}}[1]$, this bound leaves the Maxwell integral, and Lemma~\ref{lem:bounded-integrands} gives
\[
\sum_z\int_{[0,2\pi/q)^{\mathbb{T}^{(1)}}}\mathrm{d}A\,\exp\left(-\frac12\left\|dA+\frac{2\pi}{q}\cdot z\right\|^2\right)\le\left(\frac{2\pi}{q}\cdot C_q\right)^{3|\mathbb{T}|},
\]
the sum being taken over the closed $z\in C^2(\Gamma;\Z)$.
The two inequalities together are~\eqref{eq:positivity-ceiling}.

\emph{The rest of item~\eqref{it:floor-bounds}.}
For $\rho^2\neq1$ the number $C_*$ is positive, so the left side of~\eqref{eq:positivity-floor} is positive,
while the right side of~\eqref{eq:positivity-ceiling} is finite; hence $0<Z_{\mathbb{T}}[1]<\infty$.
Taking logarithms in~\eqref{eq:positivity-floor} and dividing by $|\mathbb{T}|$ gives the lower bound for $|\mathbb{T}|^{-1}\log Z_{\mathbb{T}}[1]$.

\emph{Item~\eqref{it:floor-equivalence}.}
By item~\eqref{it:floor-polynomial} the function $A\mapsto\det D^{\mathbb{T}}_A$ is not identically zero.
Lemma~\ref{lem:flux-terms}\eqref{it:flux-null-set} and~\eqref{it:flux-equivalence} therefore apply and give the assertions of item~\eqref{it:floor-equivalence}.
\end{proof}

\begin{lemma}[Perturbed denominators]\label{lem:perturbed-denominators}
Let $\rho$ and $M$ be real, $w_0,\dots,w_3>0$ and $q>0$,
let $\Gamma$ be a subgroup of $\Z^3$ of rank three, $\chi$ a unitary character of $\Gamma$, and $\mathbb{T}=(\Gamma,\chi)$.
\begin{enumerate}
\item\label{it:perturbed-bounded}
For every bounded real Borel function $V$ on $\Conf(\Gamma)$, with $\sup|V|$ the supremum of $|V|$ over $\Conf(\Gamma)$,
\[
Z_{\mathbb{T}}\left[\exp(V)\right]\ \ge\ \exp\left(-\sup|V|\right)\cdot Z_{\mathbb{T}}[1].
\]
\item\label{it:perturbed-gaussian}
Assume that $\rho^2\neq1$, or more generally that the function $A\mapsto\det D^{\mathbb{T}}_A$ on $\R^{\mathbb{T}^{(1)}}$ is not identically zero.
Then for every $\beta\in C^2(\Gamma)$
\[
Z_{\mathbb{T}}\left[\exp\left(-\langle\beta,F\rangle-\frac12\|\beta\|^2\right)\right]>0,
\]
where $\langle\beta,F\rangle=\sum_\sigma\beta_\sigma\cdot F_\sigma$ and $\|\beta\|^2=\langle\beta,\beta\rangle$, both sums being taken over the $3|\mathbb{T}|$ plaquettes $\sigma$ of $\Z^3/\Gamma$.
\end{enumerate}
\end{lemma}

\begin{proof}
Pointwise on $\Conf(\Gamma)$ one has $\exp(V)\ge\exp\left(-\sup|V|\right)$.
By Definition~\ref{def:general-functional}, $Z_{\mathbb{T}}[\Phi]$ is the integral of $\Phi$ against the measure on $\Conf(\Gamma)$ with the density $\exp\left(-\frac12\|F\|^2\right)\cdot\det D^{\mathbb{T}}_A$,
which is non-negative by Proposition~\ref{prop:accretivity}.
So $\Phi\mapsto Z_{\mathbb{T}}[\Phi]$ is monotone on non-negative Borel functions, and
\[
Z_{\mathbb{T}}\left[\exp(V)\right]\ge Z_{\mathbb{T}}\left[\exp\left(-\sup|V|\right)\right]=\exp\left(-\sup|V|\right)\cdot Z_{\mathbb{T}}[1],
\]
which is item~\eqref{it:perturbed-bounded}.

For item~\eqref{it:perturbed-gaussian}, note first that for $\rho^2\neq1$ the determinant does not vanish identically, by Theorem~\ref{thm:floor}\eqref{it:floor-polynomial}.
The curvature $F=dA+\frac{2\pi}{q}\cdot z$ is a function on $\Conf(\Gamma)$ (Definition~\ref{def:configurations}),
and $\Phi:=\exp\left(-\langle\beta,F\rangle-\frac12\|\beta\|^2\right)$ is a Borel function on $\Conf(\Gamma)$ with positive values.
The measure of $\Conf(\Gamma)$ is not zero, as the classes of the $(A,0)$ with $A\in[0,2\pi/q)^{\mathbb{T}^{(1)}}$ already have the measure $(2\pi/q)^{3|\mathbb{T}|}$.
Hence $\Phi$ is not zero almost everywhere, and Lemma~\ref{lem:flux-terms}\eqref{it:flux-equivalence} gives $Z_{\mathbb{T}}[\Phi]\neq0$.
Being non-negative, $Z_{\mathbb{T}}[\Phi]$ is positive.
\end{proof}

\subsection{Reflection positivity of free determinants}\label{sec:determinant-positivity}

In this subsection $\rho$ and $M$ are real and $w_0,\dots,w_3>0$;
the condition $-1\le\rho\le1$ is imposed where it is stated.

For $-1\le\rho\le1$, the end points $\rho=\pm1$ included,
the positivity of $Z_{\mathbb{T}}[1]$ on a general torus $\mathbb{T}=(\Gamma,\chi)$
is deduced from the positivity of free determinants, the determinants of the chamber operator at $A=0$,
on antiperiodic coordinate tori.
For integers $n_1,n_2,n_3\ge1$, the antiperiodic coordinate torus $C(n_1,n_2,n_3)$ is the general torus
formed by the group $n_1\Z e_1+n_2\Z e_2+n_3\Z e_3$
and the character of this group with the value $-1$ at $n_1e_1$, $n_2e_2$ and $n_3e_3$.
By Lemma~\ref{lem:flux-terms}, $Z_{\mathbb{T}}[1]$ and each of its flux terms are positive
as soon as the function $A\mapsto\det D^{\mathbb{T}}_A$ is not identically zero;
the remaining argument concerns determinants at $A=0$ and at constant one-cochains $A$ only.
Were this function identically zero,
Lemma~\ref{lem:characters-coverings} would give $\det D^{C(n_1,n_2,n_3)}_0=0$
whenever $n_1e_1$, $n_2e_2$ and $n_3e_3$ lie in $\Gamma$;
and such integers exist, since $\Gamma$ has finite index in $\Z^3$.

Free determinants are controlled by reflection positivity in a form that involves matrices only.
Let a set of sites be invariant under the reflection $\theta_{x_\mu,c}$ in a coordinate plane,
and let $a$ be the compression of $D_0$ to those sites of the set that lie on one side of the plane.
When the spinors on the other side are identified with those on the first side by means of the reflection,
the compression of $D_0$ to the whole set becomes the block operator
\[
\begin{pmatrix} a & -Q^{+}\\ Q^{-} & a^{*}\end{pmatrix},
\]
in which $Q^{+}$ and $Q^{-}$ are Hermitian operators contributed by the bonds that cross the plane,
both non-negative because $|\rho|\le1$.
This structure comes from Lemma~\ref{lem:hopping}
and from the counterpart for matrices of the sign at the second mirror,
which Lemma~\ref{lem:second-mirror} establishes for elements of the Grassmann algebra.
By Lemma~\ref{lem:reflection-determinants}, whose proof is pure linear algebra,
the determinants of such block operators obey the inequalities known from reflection positivity:
they are real and not smaller than $|\det a|^2$,
they do not decrease when $Q^{+}$ or $Q^{-}$ is enlarged,
and they satisfy a Cauchy--Schwarz inequality.

Once all bonds that cross site planes or link planes are deleted,
the bonds that remain join the sites in disjoint hexagons of six sites each.
Applied at all these planes, first in the direction $e_1$, then in $e_2$ and in $e_3$,
the three inequalities bound $\det D^{C(n_1,n_2,n_3)}_0$ from below
by the determinant of the hexagon operator raised to the power $8n_1n_2n_3$
(Lemma~\ref{lem:half-slab} and Proposition~\ref{prop:free-determinants}).
That determinant has a closed form, computed in Lemma~\ref{lem:hexagon-determinant},
and it is positive for every real $M$.

\paragraph{Linear algebra.}
Matrices are complex, and a matrix with no rows and no columns has determinant $1$.
For a matrix $Y$, $Y^{*}$ is its conjugate transpose;
$1_n$ is the unit matrix of size $n$, and $1$ is a unit matrix of the size that the context requires.
For Hermitian matrices $Z$, $Z'$ of one size, $Z\le Z'$ means that $Z'-Z$ is positive semidefinite,
and $Z\ge0$ means $0\le Z$.
Let $n,r,s\ge0$ be integers with $s\ge r$.
For an $r\times s$ matrix $X$ and a subset $I\subset\{1,\dots,s\}$ with $\#|I|=r$,
let $X_I$ be the $r\times r$ matrix formed by the columns of $X$ whose indices lie in $I$.
In the first two of the following five relations,
$E_{11}$, $E_{12}$, $E_{21}$, $E_{22}$ are the blocks of a square matrix with square diagonal blocks $E_{11}$ and $E_{22}$,
and in the second relation the four blocks are square matrices of one size;
in the fourth relation $Z$ and $Z'$ are Hermitian matrices of one size;
in the fifth, the sum extends over all subsets $I\subset\{1,\dots,s\}$ with $\#|I|=r$.
\begin{equation}\label{eq:positivity-identities}
\begin{aligned}
\det\begin{pmatrix} E_{11} & E_{12}\\ E_{21} & E_{22}\end{pmatrix}
&=\det E_{11}\cdot\det\left(E_{22}-E_{21}E_{11}^{-1}E_{12}\right)
&&\qquad\left(\det E_{11}\neq0\right),\\
\det\begin{pmatrix} E_{11} & E_{12}\\ E_{21} & E_{22}\end{pmatrix}
&=\det\left(E_{11}E_{22}-E_{12}E_{21}\right)
&&\qquad\left(E_{21}E_{22}=E_{22}E_{21}\right),\\
\det\left(1_n+XY\right)
&=\det\left(1_r+YX\right)
&&\qquad\left(X\in\C^{n\times r},\ Y\in\C^{r\times n}\right),\\
\det\left(1+Z\right)
&\le\det\left(1+Z'\right)
&&\qquad\left(0\le Z\le Z'\right),\\
\det\left(XY^{*}\right)
&=\sum_{I}\det X_I\cdot\overline{\det Y_I}
&&\qquad\left(X,Y\in\C^{r\times s}\right).
\end{aligned}
\end{equation}
The first relation is the formula for the Schur complement,
the third is Sylvester's identity,
and the fifth is the formula of Cauchy and Binet.
If $E_{21}$ and $E_{22}$ commute, then
\[
\begin{pmatrix} E_{11} & E_{12}\\ E_{21} & E_{22}\end{pmatrix}
\begin{pmatrix} E_{22} & 0\\ -E_{21} & 1\end{pmatrix}
=\begin{pmatrix} E_{11}E_{22}-E_{12}E_{21} & E_{12}\\ 0 & E_{22}\end{pmatrix},
\]
and taking determinants gives the second relation multiplied by $\det E_{22}$.
This proves the second relation for invertible $E_{22}$;
in general it follows by letting $\varepsilon\to0$ in the relation for $E_{22}+\varepsilon\cdot1$,
which commutes with $E_{21}$ and is invertible for all sufficiently small $\varepsilon>0$,
since both sides are polynomials in the matrix entries.
For the fourth relation, list the eigenvalues of $Z$ and those of $Z'$, with multiplicity, in non-increasing order.
By the min-max principle each eigenvalue of $Z$ is at most the eigenvalue of $Z'$ at the same place of the list,
and all of them are non-negative;
moreover $\det(1+Z)$ is the product of the numbers obtained by adding $1$ to the eigenvalues of $Z$, and likewise for $Z'$.

\begin{lemma}[Determinants with a reflection]\label{lem:reflection-determinants}
Let $\mathcal{H}_a$, $\mathcal{H}_b$ and $\mathcal{K}$ be finite-dimensional complex Hilbert spaces;
the zero space is admitted, and the determinant of the operator on the zero space is $1$.
Let $a\in\End\mathcal{H}_a$ and $b\in\End\mathcal{H}_b$,
let $j_a\colon\mathcal{K}\to\mathcal{H}_a$ and $j_b\colon\mathcal{K}\to\mathcal{H}_b$ be linear maps,
and let $Q^{+},Q^{-}\in\End\mathcal{K}$ be Hermitian.
Put
\[
\mathfrak{D}(a,b):=\det\begin{pmatrix} a & -j_aQ^{+}j_b^{*}\\ j_bQ^{-}j_a^{*} & b^{*}\end{pmatrix},
\]
the determinant of an operator on $\mathcal{H}_a\oplus\mathcal{H}_b$,
and define $\mathfrak{D}(a,a)$ and $\mathfrak{D}(b,b)$ by the same formula
with $(a,j_a)$, respectively $(b,j_b)$, in both places:
\[
\mathfrak{D}(a,a):=\det\begin{pmatrix} a & -j_aQ^{+}j_a^{*}\\ j_aQ^{-}j_a^{*} & a^{*}\end{pmatrix},
\qquad
\mathfrak{D}(b,b):=\det\begin{pmatrix} b & -j_bQ^{+}j_b^{*}\\ j_bQ^{-}j_b^{*} & b^{*}\end{pmatrix}.
\]
If $Q^{+}\ge0$ and $Q^{-}\ge0$, then the following hold.
\begin{enumerate}
\item\label{it:determinant-lower}
$\mathfrak{D}(a,a)$ is real, and $\mathfrak{D}(a,a)\ge|\det a|^2$.
\item\label{it:determinant-monotone}
$\mathfrak{D}(a,a)$ does not decrease if $Q^{+}$ is replaced by a Hermitian operator $Q^{+\prime}\ge Q^{+}$,
or $Q^{-}$ by a Hermitian operator $Q^{-\prime}\ge Q^{-}$.
\item\label{it:determinant-schwarz}
$|\mathfrak{D}(a,b)|^2\le\mathfrak{D}(a,a)\cdot\mathfrak{D}(b,b)$.
\end{enumerate}
\end{lemma}

\begin{proof}
Throughout the proof $Q^{+}\ge0$ and $Q^{-}\ge0$,
and $\left(Q^{+}\right)^{\frac{1}{2}}$, $\left(Q^{-}\right)^{\frac{1}{2}}$ denote their Hermitian square roots $\ge0$.
Operators between the spaces of the lemma are identified with their matrices in fixed orthonormal bases;
adjoints then become conjugate transposes, and the relations~\eqref{eq:positivity-identities} apply to them.

If $\mathcal{K}$ is the zero space, the two off-diagonal blocks vanish, so that
$\mathfrak{D}(a,b)=\det a\cdot\overline{\det b}$, $\mathfrak{D}(a,a)=|\det a|^2$ and $\mathfrak{D}(b,b)=|\det b|^2$;
then \eqref{it:determinant-lower}--\eqref{it:determinant-schwarz} hold with equality.
From now on $r:=\dim\mathcal{K}\ge1$.

\emph{Invertible $a$ and $b$.}
Let $a$ and $b$ be invertible, and define two operators on $\mathcal{K}$ by
\[
\tilde a:=j_a^{*}a^{-1}j_a,\qquad\tilde b:=j_b^{*}b^{-1}j_b ;
\]
then $\tilde b^{*}=j_b^{*}\left(b^{*}\right)^{-1}j_b$.
The upper left block $a$ of the operator that defines $\mathfrak{D}(a,b)$ is invertible,
and the first relation of~\eqref{eq:positivity-identities} gives
\[
\mathfrak{D}(a,b)=\det a\cdot\det\left(b^{*}+j_bQ^{-}\tilde aQ^{+}j_b^{*}\right).
\]
Factor out $b^{*}$ on the left, whose determinant is $\overline{\det b}$,
and apply Sylvester's identity, the third relation of~\eqref{eq:positivity-identities},
first with $X=\left(b^{*}\right)^{-1}j_b$ and $Y=Q^{-}\tilde aQ^{+}j_b^{*}$,
then with $X=\left(Q^{-}\right)^{\frac{1}{2}}$ and $Y=\left(Q^{-}\right)^{\frac{1}{2}}\tilde aQ^{+}\tilde b^{*}$.
With the operators
\[
\hat a:=\left(Q^{-}\right)^{\frac{1}{2}}\tilde a\left(Q^{+}\right)^{\frac{1}{2}},
\qquad
\hat b:=\left(Q^{-}\right)^{\frac{1}{2}}\tilde b\left(Q^{+}\right)^{\frac{1}{2}}
\]
on $\mathcal{K}$, for which $\hat b^{*}=\left(Q^{+}\right)^{\frac{1}{2}}\tilde b^{*}\left(Q^{-}\right)^{\frac{1}{2}}$, this yields
\begin{equation}\label{eq:positivity-factorisation}
\begin{aligned}
\mathfrak{D}(a,b)
&=\det a\cdot\overline{\det b}\cdot\det_{\mathcal{H}_b}\left(1+\left(b^{*}\right)^{-1}j_bQ^{-}\tilde aQ^{+}j_b^{*}\right)\\
&=\det a\cdot\overline{\det b}\cdot\det_{\mathcal{K}}\left(1+Q^{-}\tilde aQ^{+}\tilde b^{*}\right)
=\det a\cdot\overline{\det b}\cdot\det_{\mathcal{K}}\left(1+\hat a\hat b^{*}\right).
\end{aligned}
\end{equation}
For the pair $(a,a)$, with $(a,j_a)$ in both places, the operators $\tilde b$ and $\hat b$ are replaced by $\tilde a$ and $\hat a$;
for the pair $(b,b)$, the operators $\tilde a$ and $\hat a$ are replaced by $\tilde b$ and $\hat b$.
So~\eqref{eq:positivity-factorisation} also gives
\[
\mathfrak{D}(a,a)=|\det a|^2\cdot\det_{\mathcal{K}}\left(1+\hat a\hat a^{*}\right),
\qquad
\mathfrak{D}(b,b)=|\det b|^2\cdot\det_{\mathcal{K}}\left(1+\hat b\hat b^{*}\right).
\]

The $r\times2r$ matrices $\begin{pmatrix} 1_r & \hat a\end{pmatrix}$ and $\begin{pmatrix} 1_r & \hat b\end{pmatrix}$ satisfy
\[
\begin{pmatrix} 1_r & \hat a\end{pmatrix}\begin{pmatrix} 1_r & \hat b\end{pmatrix}^{*}=1_r+\hat a\hat b^{*},
\]
and similarly with $\hat a$ in both places, or with $\hat b$ in both places.
The formula of Cauchy and Binet, the fifth relation of~\eqref{eq:positivity-identities} with $s=2r$, therefore gives
\[
\det\left(1_r+\hat a\hat b^{*}\right)
=\sum_{I}\det\begin{pmatrix} 1_r & \hat a\end{pmatrix}_I\cdot\overline{\det\begin{pmatrix} 1_r & \hat b\end{pmatrix}_I},
\qquad
\det\left(1_r+\hat a\hat a^{*}\right)
=\sum_{I}\left|\det\begin{pmatrix} 1_r & \hat a\end{pmatrix}_I\right|^2,
\]
and $\det\left(1_r+\hat b\hat b^{*}\right)$ is the corresponding sum for $\hat b$;
here $I$ ranges over the subsets of $\{1,\dots,2r\}$ with $\#|I|=r$.
By the Cauchy--Schwarz inequality for these finite sums,
\[
\left|\det\left(1_r+\hat a\hat b^{*}\right)\right|^2
\le\det\left(1_r+\hat a\hat a^{*}\right)\cdot\det\left(1_r+\hat b\hat b^{*}\right),
\]
and multiplication by $|\det a|^2\cdot|\det b|^2$ turns this inequality into~\eqref{it:determinant-schwarz}, by the three factorisations above.
In the sum for $\det\left(1_r+\hat a\hat a^{*}\right)$ the term of $I=\{1,\dots,r\}$ is $|\det 1_r|^2=1$, and no term is negative.
Hence $\det\left(1_r+\hat a\hat a^{*}\right)\ge1$,
and $\mathfrak{D}(a,a)=|\det a|^2\cdot\det\left(1_r+\hat a\hat a^{*}\right)$ is real and at least $|\det a|^2$:
this is~\eqref{it:determinant-lower}.

Turning to~\eqref{it:determinant-monotone}, let $Z:=\tilde aQ^{+}\tilde a^{*}$,
a Hermitian operator on $\mathcal{K}$ with $Z\ge0$, which does not depend on $Q^{-}$.
As $\hat a\hat a^{*}=\left(Q^{-}\right)^{\frac{1}{2}}Z\left(Q^{-}\right)^{\frac{1}{2}}$,
Sylvester's identity with $X=\left(Q^{-}\right)^{\frac{1}{2}}Z^{\frac{1}{2}}$ and $Y=Z^{\frac{1}{2}}\left(Q^{-}\right)^{\frac{1}{2}}$ gives
\[
\det\left(1+\hat a\hat a^{*}\right)
=\det\left(1+\left(Q^{-}\right)^{\frac{1}{2}}\tilde aQ^{+}\tilde a^{*}\left(Q^{-}\right)^{\frac{1}{2}}\right)
=\det\left(1+Z^{\frac{1}{2}}Q^{-}Z^{\frac{1}{2}}\right).
\]
Conjugation by a fixed operator preserves the order of Hermitian operators.
If $Q^{+}$ is replaced by $Q^{+\prime}\ge Q^{+}$, then $\tilde aQ^{+}\tilde a^{*}\le\tilde aQ^{+\prime}\tilde a^{*}$,
and conjugation by $\left(Q^{-}\right)^{\frac{1}{2}}$ preserves this inequality;
so the operator in the middle member, which is $\ge0$, does not decrease.
If $Q^{-}$ is replaced by $Q^{-\prime}\ge Q^{-}$, then conjugation by $Z^{\frac{1}{2}}$ gives
$0\le Z^{\frac{1}{2}}Q^{-}Z^{\frac{1}{2}}\le Z^{\frac{1}{2}}Q^{-\prime}Z^{\frac{1}{2}}$.
In both cases the fourth relation of~\eqref{eq:positivity-identities} shows that $\det\left(1+\hat a\hat a^{*}\right)$ does not decrease.
The factor $|\det a|^2$ depends on neither $Q^{+}$ nor $Q^{-}$,
and the factorisation $\mathfrak{D}(a,a)=|\det a|^2\cdot\det\left(1+\hat a\hat a^{*}\right)$
holds for every pair of Hermitian operators $Q^{+}\ge0$, $Q^{-}\ge0$;
consequently \eqref{it:determinant-monotone} holds.

\emph{General $a$ and $b$.}
For fixed $j_a$, $j_b$, $Q^{+}$ and $Q^{-}$, the numbers $\mathfrak{D}(a,b)$, $\mathfrak{D}(a,a)$, $\mathfrak{D}(b,b)$ and $|\det a|^2$,
as well as the values of $\mathfrak{D}(a,a)$ after the replacements in~\eqref{it:determinant-monotone},
are polynomials in the matrix entries of $a$ and $b$ and in their complex conjugates,
hence continuous functions of $(a,b)$.
Since $a$ and $b$ have finitely many eigenvalues,
the operators $a+\varepsilon\cdot1$ and $b+\varepsilon\cdot1$ are invertible for all sufficiently small $\varepsilon>0$,
and \eqref{it:determinant-lower}--\eqref{it:determinant-schwarz} hold for them by the case already treated.
The inequalities in \eqref{it:determinant-lower}--\eqref{it:determinant-schwarz} are not strict, and a limit of real numbers is real;
letting $\varepsilon\to0$ gives \eqref{it:determinant-lower}--\eqref{it:determinant-schwarz} for $a$ and $b$.
\end{proof}

\subsubsection{Pieces and coordinate mirrors}\label{sec:pieces}

Elements of $W$ map bonds to bonds of the same type.
Indeed, let $g\in W$ and $x\in\mathbb{X}$, and write $x=g'v_0$ with $g'\in W$;
then $gx=(gg')v_0$, and Definition~\ref{def:fermion-sites} gives
\[
(gx)^{(i)}=g\,x^{(i)},\qquad i=0,\dots,3 .
\]
The translations by vectors $v\in\Z^3$ belong to $W$, and they preserve, in addition, the direction of every bond.
The hopping matrix $w(x,x')$ of a bond $\{x,x'\}$ depends only on the type of the bond and on the unit vector from $x$ to $x'$
(Definition~\ref{def:chamber-operator}).
Therefore
\[
(x+v)^{(i)}=x^{(i)}+v,\qquad w(x+v,x^{(i)}+v)=w(x,x^{(i)}),\qquad x\in\mathbb{X},\ v\in\Z^3,\ i=0,\dots,3 .
\]

\begin{definition}[Pieces]\label{def:pieces}
Let $\rho$ and $M$ be real and $w_0,\dots,w_3>0$.
A \emph{piece} is a triple $\mathfrak{P}=(\mathbb{Y},\Gamma_0,\chi_0)$ in which
$\Gamma_0$ is a subgroup of $\Z^3$ of any rank, the trivial group $\{0\}$ included,
$\chi_0\colon\Gamma_0\to\{+1,-1\}$ is a character,
and $\mathbb{Y}\subset\mathbb{X}$ is a set of sites with $\mathbb{Y}+\Gamma_0=\mathbb{Y}$
such that the set $\mathbb{Y}/\Gamma_0$ of the $\Gamma_0$-orbits in $\mathbb{Y}$ is finite.
The Hilbert space of $\mathfrak{P}$ is
\[
\mathcal{H}(\mathfrak{P}):=\left\{\psi\colon\mathbb{Y}\to\C^4:\ \psi(x+g)=\chi_0(g)\cdot\psi(x)\ \text{for all}\ x\in\mathbb{Y},\ g\in\Gamma_0\right\},
\]
with the scalar product
\[
\langle\psi',\psi\rangle:=\sum_{x\in\mathbb{Y}/\Gamma_0}\psi'(x)^{*}\cdot\psi(x),
\]
in which $x$ ranges over a set of representatives of the orbits;
the summand does not depend on the representative, since $|\chi_0(g)|=1$,
and $\dim\mathcal{H}(\mathfrak{P})=4\cdot\#|\mathbb{Y}/\Gamma_0|$.
The operator $D_{\mathfrak{P}}\in\End\mathcal{H}(\mathfrak{P})$ is the compression to $\mathbb{Y}$
of the chamber operator $D_0$ at $A=0$ (Definition~\ref{def:chamber-operator}):
\[
(D_{\mathfrak{P}}\psi)(x):=M\psi(x)-\sum_{i:\ x^{(i)}\in\mathbb{Y}}w(x,x^{(i)})\cdot\psi(x^{(i)}),\qquad x\in\mathbb{Y},
\]
the sum being over the types $i\in\{0,1,2,3\}$ with $x^{(i)}\in\mathbb{Y}$.
It maps $\mathcal{H}(\mathfrak{P})$ into itself:
for $g\in\Gamma_0$ one has $\mathbb{Y}+g=\mathbb{Y}$, $(x+g)^{(i)}=x^{(i)}+g$ and $w(x+g,x^{(i)}+g)=w(x,x^{(i)})$.
When the triple is written out, $\mathcal{H}(\mathbb{Y},\Gamma_0,\chi_0):=\mathcal{H}(\mathfrak{P})$
and $D_{(\mathbb{Y},\Gamma_0,\chi_0)}:=D_{\mathfrak{P}}$.
If $\mathbb{Y}$ is empty, then $\mathcal{H}(\mathfrak{P})=\{0\}$, and $\det D_{\mathfrak{P}}:=1$.
Pieces have the following three properties.
\begin{enumerate}
\item\label{it:piece-torus}
Let $\mathbb{T}=(\Gamma,\chi)$ be a torus in the sense of Definition~\ref{def:general-tori}
whose character takes only the values $+1$ and $-1$.
Then $(\mathbb{X},\Gamma,\chi)$ is a piece, and $D_{(\mathbb{X},\Gamma,\chi)}=D^{\mathbb{T}}_0$.
(The set $\mathbb{X}/\Gamma$ has $48\cdot|\mathbb{T}|$ elements,
$\mathcal{H}(\mathbb{X},\Gamma,\chi)$ is the space of the spinors of $\mathbb{T}$,
and all four neighbours of a site lie in $\mathbb{X}$.)
\item\label{it:piece-translation}
Let $(\mathbb{Y},\Gamma_0,\chi_0)$ be a piece and $v\in\Z^3$.
Then $(\mathbb{Y}+v,\Gamma_0,\chi_0)$ is a piece, and
\[
\det D_{(\mathbb{Y}+v,\Gamma_0,\chi_0)}=\det D_{(\mathbb{Y},\Gamma_0,\chi_0)} .
\]
(One has $\mathbb{X}+v=\mathbb{X}$.
The map $\psi\mapsto\psi(\cdot-v)$ is a unitary of $\mathcal{H}(\mathbb{Y},\Gamma_0,\chi_0)$ onto $\mathcal{H}(\mathbb{Y}+v,\Gamma_0,\chi_0)$,
and it intertwines the two operators, because $(x+v)^{(i)}=x^{(i)}+v$ and $w(x+v,x^{(i)}+v)=w(x,x^{(i)})$.)
\item\label{it:piece-adjoint}
The adjoint of $D_{\mathfrak{P}}$ is given by
\[
(D_{\mathfrak{P}}^{*}\psi)(x)=M\psi(x)-\sum_{i:\ x^{(i)}\in\mathbb{Y}}w(x^{(i)},x)\cdot\psi(x^{(i)}),\qquad x\in\mathbb{Y}.
\]
\end{enumerate}
\end{definition}

\begin{proof}[Proof of \eqref{it:piece-adjoint}]
Fix $\psi,\psi'\in\mathcal{H}(\mathfrak{P})$, and for $i\in\{0,1,2,3\}$ let $\mathbb{Y}_i:=\{x\in\mathbb{Y}:\ x^{(i)}\in\mathbb{Y}\}$.
Then
\[
\langle\psi',D_{\mathfrak{P}}\psi\rangle
=\sum_{x\in\mathbb{Y}/\Gamma_0}M\psi'(x)^{*}\cdot\psi(x)
-\sum_{i=0}^{3}\ \sum_{x\in\mathbb{Y}_i/\Gamma_0}\psi'(x)^{*}\cdot w(x,x^{(i)})\cdot\psi(x^{(i)}) .
\]
Since $(x^{(i)})^{(i)}=x$ (Definition~\ref{def:fermion-sites}), the map $x\mapsto x^{(i)}$ is an involution of $\mathbb{Y}_i$;
it commutes with the translations by $\Gamma_0$ and so induces an involution of $\mathbb{Y}_i/\Gamma_0$.
The summands of the inner sum are invariant under $\Gamma_0$,
because $\chi_0(g)^2=1$ and $w(x+g,x^{(i)}+g)=w(x,x^{(i)})$ for $g\in\Gamma_0$;
so the inner sum may be re-indexed by this involution:
\[
\sum_{x\in\mathbb{Y}_i/\Gamma_0}\psi'(x)^{*}\cdot w(x,x^{(i)})\cdot\psi(x^{(i)})
=\sum_{x\in\mathbb{Y}_i/\Gamma_0}\psi'(x^{(i)})^{*}\cdot w(x^{(i)},x)\cdot\psi(x)
=\sum_{x\in\mathbb{Y}_i/\Gamma_0}\left(w(x^{(i)},x)\cdot\psi'(x^{(i)})\right)^{*}\cdot\psi(x);
\]
the last equality holds because the hopping matrices are Hermitian (Lemma~\ref{lem:hopping}).
As $M$ is real, it follows that $\langle\psi',D_{\mathfrak{P}}\psi\rangle=\langle\psi'',\psi\rangle$ with
\[
\psi''(x):=M\psi'(x)-\sum_{i:\ x^{(i)}\in\mathbb{Y}}w(x^{(i)},x)\cdot\psi'(x^{(i)}),\qquad x\in\mathbb{Y}.
\]
The function $\psi''$ belongs to $\mathcal{H}(\mathfrak{P})$ for the same reason as $D_{\mathfrak{P}}\psi'$,
and $\psi$ is arbitrary; hence $D_{\mathfrak{P}}^{*}\psi'=\psi''$.
\end{proof}

\begin{definition}[Crossing matrices]\label{def:crossing-matrices}
Let $\rho$ be real, $w_0,\dots,w_3>0$, and $\mu\in\{1,2,3\}$;
the objects below depend on $\mu$, which is not shown in the notation.
For $c\in\frac{1}{2}\Z$ let $\theta_c:=\theta_{x_\mu,c}$ be the reflection in the coordinate plane $\{x_\mu=c\}$,
\[
\theta_cx=x-2(x_\mu-c)\cdot e_\mu,\qquad x\in\R^3
\]
(Definition~\ref{def:mirrors});
the mirror $(x_\mu,c)$ is a site plane if $c\in\Z$ and a link plane if $c\in\Z+\frac{1}{2}$.
Put $w^{\times}_c:=w_0$ if $c\in\Z$ and $w^{\times}_c:=w_3$ if $c\in\Z+\frac{1}{2}$.
For a set $\mathbb{Y}\subset\mathbb{X}$ and real numbers $s<s'$ put
\[
\mathbb{Y}(s,s'):=\{x\in\mathbb{Y}:\ s<x_\mu<s'\}.
\]
Let $\mathbb{B}_c$ be the set of the sites $x\in\mathbb{X}$ with $x_\mu>c$ for which $\theta_cx$ is a neighbour of $x$,
that is, $\theta_cx=x^{(i)}$ for some $i\in\{0,1,2,3\}$.
The crossing matrices at $c$ are the $4\times4$ matrices
\[
Q^{+}_c:=\frac{w^{\times}_c}{2}\,(1+\rho\gamma_\mu),\qquad
Q^{-}_c:=\frac{w^{\times}_c}{2}\,(1-\rho\gamma_\mu).
\]
\end{definition}

The matrix $\gamma_\mu$ is Hermitian, and its eigenvalues $1$ and $-1$ both have multiplicity two (Definition~\ref{def:spin-matrices}).
So $Q^{+}_c$ and $Q^{-}_c$ are Hermitian, each of them has the eigenvalues
$\frac{w^{\times}_c}{2}\,(1+\rho)$ and $\frac{w^{\times}_c}{2}\,(1-\rho)$, both of multiplicity two,
and $Q^{+}_c$, $Q^{-}_c$ are positive semidefinite if and only if $-1\le\rho\le1$.

A bond whose end points lie on different sides of the mirror $(x_\mu,c)$ has the form $\{x,\theta_cx\}$
(Proposition~\ref{prop:geometry}\eqref{it:geometry-mirrors});
thus $\mathbb{B}_c$ consists of the upper end points of the bonds that cross this mirror.
For $x\in\mathbb{B}_c$ the plane $\{x_\mu=c\}$ meets the bond $\{x,\theta_cx\}$ orthogonally at its midpoint
(Proposition~\ref{prop:geometry}\eqref{it:geometry-midpoint}),
and the length $|\theta_cx-x|=2(x_\mu-c)$ of the bond is smaller than $1$ (Lemma~\ref{lem:neighbours}).
Therefore
\[
\mathbb{B}_c\subset\left\{x\in\mathbb{X}:\ c<x_\mu<c+\tfrac{1}{2}\right\}.
\]
If $g\in\Z^3$ and $g_\mu=0$, then $\theta_c(x+g)=\theta_cx+g$ and $(x+g)^{(i)}=x^{(i)}+g$ for $x\in\mathbb{X}$;
so $\mathbb{B}_c+g=\mathbb{B}_c$.

Now let $x\in\mathbb{B}_c$.
The bond $\{x,\theta_cx\}$ crosses the plane $\{x_\mu=c\}$;
by Proposition~\ref{prop:geometry}\eqref{it:geometry-projection} it has the type $0$ if $c\in\Z$ and the type $3$ if $c\in\Z+\frac{1}{2}$,
and its weight is $w^{\times}_c$ in both cases.
The unit vector from $x$ to $\theta_cx$ is $-e_\mu$, and the one from $\theta_cx$ to $x$ is $e_\mu$.
Definition~\ref{def:chamber-operator} gives
\[
w(x,\theta_cx)=\frac{w^{\times}_c}{2}\,(\rho+\gamma_\mu),\qquad
w(\theta_cx,x)=\frac{w^{\times}_c}{2}\,(\rho-\gamma_\mu),
\]
and with $\gamma_\mu^2=1$ one obtains
\begin{equation}\label{eq:positivity-crossing}
-w(x,\theta_cx)\cdot\gamma_\mu=-Q^{+}_c,\qquad
-\gamma_\mu\cdot w(\theta_cx,x)=Q^{-}_c,\qquad x\in\mathbb{B}_c .
\end{equation}

\begin{lemma}[Gluing at a mirror]\label{lem:gluing}
Let $\rho$ and $M$ be real, $w_0,\dots,w_3>0$, $\mu\in\{1,2,3\}$ and $c\in\frac{1}{2}\Z$,
and let $\theta_c=\theta_{x_\mu,c}$, $\mathbb{B}_c$, $Q^{+}_c$ and $Q^{-}_c$ be formed with this $\mu$ as in Definition~\ref{def:crossing-matrices}.
Let $\Gamma_0$ be a subgroup of $\{g\in\Z^3:\ g_\mu=0\}$, let $\chi_0\colon\Gamma_0\to\{+1,-1\}$ be a character,
and let $(\mathbb{U},\Gamma_0,\chi_0)$ and $(\mathbb{V},\Gamma_0,\chi_0)$ be pieces
with $\mathbb{U}\subset\{x_\mu>c\}$ and $\mathbb{V}\subset\{x_\mu>c\}$.
Put
\[
a:=D_{(\mathbb{U},\Gamma_0,\chi_0)},\qquad
b:=D_{(\mathbb{V},\Gamma_0,\chi_0)},\qquad
\mathcal{K}:=\mathcal{H}(\mathbb{U}\cap\mathbb{V}\cap\mathbb{B}_c,\Gamma_0,\chi_0),
\]
let $j_a\colon\mathcal{K}\to\mathcal{H}(\mathbb{U},\Gamma_0,\chi_0)$ and $j_b\colon\mathcal{K}\to\mathcal{H}(\mathbb{V},\Gamma_0,\chi_0)$ be the extensions by zero,
and let $Q^{+}$ and $Q^{-}$ be the operators on $\mathcal{K}$ of multiplication by $Q^{+}_c$ and by $Q^{-}_c$.
Then $Q^{+}$ and $Q^{-}$ are Hermitian, $(\mathbb{U}\cup\theta_c\mathbb{V},\Gamma_0,\chi_0)$ is a piece, and
\[
\det D_{(\mathbb{U}\cup\theta_c\mathbb{V},\Gamma_0,\chi_0)}
=\mathfrak{D}(a,b)
=\det\begin{pmatrix} a & -j_aQ^{+}j_b^{*}\\ j_bQ^{-}j_a^{*} & b^{*}\end{pmatrix},
\]
where $\mathfrak{D}(a,b)$ is formed as in Lemma~\ref{lem:reflection-determinants}
with $\mathcal{H}_a:=\mathcal{H}(\mathbb{U},\Gamma_0,\chi_0)$, $\mathcal{H}_b:=\mathcal{H}(\mathbb{V},\Gamma_0,\chi_0)$,
the space $\mathcal{K}$, the maps $j_a$, $j_b$ and the operators $Q^{+}$, $Q^{-}$.
If $-1\le\rho\le1$, then $Q^{+}\ge0$ and $Q^{-}\ge0$.
\end{lemma}

\begin{proof}
Throughout the proof $\Gamma_0$ and $\chi_0$ are fixed;
for a piece $(\mathbb{Y},\Gamma_0,\chi_0)$ write $\mathcal{H}(\mathbb{Y}):=\mathcal{H}(\mathbb{Y},\Gamma_0,\chi_0)$
and $D_{\mathbb{Y}}:=D_{(\mathbb{Y},\Gamma_0,\chi_0)}$.

\emph{The glued piece.}
Every $g\in\Gamma_0$ satisfies $g_\mu=0$, whence $\theta_c(x+g)=\theta_cx+g$ for $x\in\R^3$.
The reflection $\theta_c$ belongs to $W$ and maps $\mathbb{X}$ onto itself.
Hence $\theta_c\mathbb{V}$ is a $\Gamma_0$-invariant subset of $\mathbb{X}$, contained in $\{x_\mu<c\}$,
and $\theta_c$ induces a bijection of $\mathbb{V}/\Gamma_0$ onto $\theta_c\mathbb{V}/\Gamma_0$.
Consequently $(\theta_c\mathbb{V},\Gamma_0,\chi_0)$ and $(\mathbb{U}\cup\theta_c\mathbb{V},\Gamma_0,\chi_0)$ are pieces;
so is $(\mathbb{U}\cap\mathbb{V}\cap\mathbb{B}_c,\Gamma_0,\chi_0)$, since $\mathbb{B}_c+g=\mathbb{B}_c$ for $g\in\Gamma_0$.
The sets $\mathbb{U}\subset\{x_\mu>c\}$ and $\theta_c\mathbb{V}$ are disjoint,
and restriction to them identifies $\mathcal{H}(\mathbb{U}\cup\theta_c\mathbb{V})$ with the orthogonal sum $\mathcal{H}(\mathbb{U})\oplus\mathcal{H}(\theta_c\mathbb{V})$.
Multiplication by a constant $4\times4$ matrix commutes with multiplication by the numbers $\chi_0(g)$;
so $Q^{+}$ and $Q^{-}$ map $\mathcal{K}$ into itself, and they are Hermitian, the matrices $Q^{\pm}_c$ being Hermitian.
The adjoints $j_a^{*}$ and $j_b^{*}$ are the restrictions to $\mathbb{U}\cap\mathbb{V}\cap\mathbb{B}_c$.

\emph{A unitary.}
For $\psi'\in\mathcal{H}(\mathbb{V})$ put
\[
(U\psi')(x'):=\gamma_\mu\cdot\psi'(\theta_cx'),\qquad x'\in\theta_c\mathbb{V}.
\]
For $g\in\Gamma_0$ one has $(U\psi')(x'+g)=\gamma_\mu\cdot\psi'(\theta_cx'+g)=\chi_0(g)\cdot(U\psi')(x')$,
so that $U$ maps $\mathcal{H}(\mathbb{V})$ into $\mathcal{H}(\theta_c\mathbb{V})$.
As $\gamma_\mu$ is unitary and $\theta_c$ induces a bijection of the sets of orbits, $U$ is unitary;
from $\gamma_\mu^2=1$ and $\theta_c^2=\mathrm{id}$ its inverse is
\[
(U^{-1}\psi)(x)=\gamma_\mu\cdot\psi(\theta_cx),\qquad x\in\mathbb{V},\ \psi\in\mathcal{H}(\theta_c\mathbb{V}).
\]

\emph{Blocks.}
Let $D:=D_{\mathbb{U}\cup\theta_c\mathbb{V}}$,
and let $D_{11}$, $D_{12}$, $D_{21}$, $D_{22}$ be its blocks with respect to $\mathcal{H}(\mathbb{U})\oplus\mathcal{H}(\theta_c\mathbb{V})$;
for example, $D_{12}\colon\mathcal{H}(\theta_c\mathbb{V})\to\mathcal{H}(\mathbb{U})$ maps $\psi$
to the restriction to $\mathbb{U}$ of $D$ applied to the extension of $\psi$ by zero.
If a spinor vanishes outside one of the two sets, the values of $D$ applied to it at the points of that set
are the values of the compression to that set;
hence $D_{11}=D_{\mathbb{U}}=a$ and $D_{22}=D_{\theta_c\mathbb{V}}$.
The map $1\oplus U$ is a unitary of $\mathcal{H}(\mathbb{U})\oplus\mathcal{H}(\mathbb{V})$ onto $\mathcal{H}(\mathbb{U})\oplus\mathcal{H}(\theta_c\mathbb{V})$, and
\[
(1\oplus U)^{-1}D\,(1\oplus U)=\begin{pmatrix} D_{11} & D_{12}U\\ U^{-1}D_{21} & U^{-1}D_{22}U\end{pmatrix}.
\]
It will be shown that the four blocks are
\begin{equation}\label{eq:positivity-mirror-blocks}
\begin{aligned}
D_{11}&=a, &\qquad D_{12}U&=-j_aQ^{+}j_b^{*},\\
U^{-1}D_{21}&=j_bQ^{-}j_a^{*}, &\qquad U^{-1}D_{\theta_c\mathbb{V}}U&=b^{*}.
\end{aligned}
\end{equation}
The first of these identities has just been obtained.

\emph{The upper right block.}
Let $\psi'\in\mathcal{H}(\mathbb{V})$ and $x\in\mathbb{U}$.
The point $x$ does not lie in $\theta_c\mathbb{V}$, so the mass term does not contribute, and
\[
(D_{12}U\psi')(x)=-\sum_{i:\ x^{(i)}\in\theta_c\mathbb{V}}w(x,x^{(i)})\cdot(U\psi')(x^{(i)}).
\]
If $x^{(i)}\in\theta_c\mathbb{V}$, the end points $x\in\{x_\mu>c\}$ and $x^{(i)}\in\{x_\mu<c\}$ of the bond
lie on different sides of the mirror $(x_\mu,c)$,
and $x^{(i)}=\theta_cx$ by Proposition~\ref{prop:geometry}\eqref{it:geometry-mirrors}.
The sum is therefore empty unless $\theta_cx$ is a neighbour of $x$ that lies in $\theta_c\mathbb{V}$,
that is, unless $x\in\mathbb{B}_c\cap\mathbb{V}$.
For $x\in\mathbb{U}\cap\mathbb{V}\cap\mathbb{B}_c$ it has exactly one term,
the four neighbours of $x$ being distinct (Lemma~\ref{lem:incongruent-neighbours}),
and \eqref{eq:positivity-crossing} gives
\[
(D_{12}U\psi')(x)=-w(x,\theta_cx)\cdot(U\psi')(\theta_cx)=-w(x,\theta_cx)\cdot\gamma_\mu\cdot\psi'(x)=-Q^{+}_c\cdot\psi'(x).
\]
So $D_{12}U\psi'$ is the extension by zero of $-Q^{+}j_b^{*}\psi'$,
which is the second identity of \eqref{eq:positivity-mirror-blocks}.

\emph{The lower left block.}
Let $\psi\in\mathcal{H}(\mathbb{U})$ and $x\in\mathbb{V}$.
The point $\theta_cx$ lies in $\theta_c\mathbb{V}$, hence not in $\mathbb{U}$, and
\[
(U^{-1}D_{21}\psi)(x)=\gamma_\mu\cdot(D_{21}\psi)(\theta_cx)
=-\gamma_\mu\sum_{i:\ (\theta_cx)^{(i)}\in\mathbb{U}}w(\theta_cx,(\theta_cx)^{(i)})\cdot\psi((\theta_cx)^{(i)}).
\]
A neighbour of $\theta_cx\in\{x_\mu<c\}$ that lies in $\mathbb{U}\subset\{x_\mu>c\}$ is on the other side of the mirror,
so it equals $\theta_c\theta_cx=x$, again by Proposition~\ref{prop:geometry}\eqref{it:geometry-mirrors}.
The sum is empty unless $x$ is a neighbour of $\theta_cx$ and $x\in\mathbb{U}$.
As $x_\mu>c$, and as $x=(\theta_cx)^{(i)}$ means $\theta_cx=x^{(i)}$, this happens exactly for $x\in\mathbb{U}\cap\mathbb{V}\cap\mathbb{B}_c$;
then the sum has exactly one term, the four neighbours of $\theta_cx$ being distinct,
and \eqref{eq:positivity-crossing} gives
\[
(U^{-1}D_{21}\psi)(x)=-\gamma_\mu\cdot w(\theta_cx,x)\cdot\psi(x)=Q^{-}_c\cdot\psi(x).
\]
Thus $U^{-1}D_{21}=j_bQ^{-}j_a^{*}$.

\emph{The lower right block.}
Let $\psi'\in\mathcal{H}(\mathbb{V})$ and $x\in\mathbb{V}$.
Since $\theta_c\in W$, the neighbours of $\theta_cx$ are the points $(\theta_cx)^{(i)}=\theta_cx^{(i)}$;
so $(\theta_cx)^{(i)}\in\theta_c\mathbb{V}$ if and only if $x^{(i)}\in\mathbb{V}$,
and $(U\psi')(\theta_cx^{(i)})=\gamma_\mu\cdot\psi'(x^{(i)})$.
Therefore
\[
\begin{aligned}
(U^{-1}D_{\theta_c\mathbb{V}}U\psi')(x)
&=\gamma_\mu\left(M\gamma_\mu\cdot\psi'(x)-\sum_{i:\ x^{(i)}\in\mathbb{V}}w(\theta_cx,\theta_cx^{(i)})\cdot\gamma_\mu\cdot\psi'(x^{(i)})\right)\\
&=M\psi'(x)-\sum_{i:\ x^{(i)}\in\mathbb{V}}\gamma_\mu\cdot w(\theta_cx,\theta_cx^{(i)})\cdot\gamma_\mu\cdot\psi'(x^{(i)}).
\end{aligned}
\]
The upward unit normal of the mirror $(x_\mu,c)$ is $e_\mu$,
and Lemma~\ref{lem:hopping}, applied to this mirror and to the bond $\{\theta_cx,\theta_cx^{(i)}\}$, yields
\[
\gamma_\mu\cdot w(\theta_cx,\theta_cx^{(i)})\cdot\gamma_\mu=w(\theta_c\theta_cx^{(i)},\theta_c\theta_cx)=w(x^{(i)},x).
\]
Hence
\[
(U^{-1}D_{\theta_c\mathbb{V}}U\psi')(x)=M\psi'(x)-\sum_{i:\ x^{(i)}\in\mathbb{V}}w(x^{(i)},x)\cdot\psi'(x^{(i)})=(b^{*}\psi')(x)
\]
by Definition~\ref{def:pieces}\eqref{it:piece-adjoint}, which is the last identity of \eqref{eq:positivity-mirror-blocks}.

\emph{The determinant.}
By \eqref{eq:positivity-mirror-blocks},
\[
(1\oplus U)^{-1}D\,(1\oplus U)=\begin{pmatrix} a & -j_aQ^{+}j_b^{*}\\ j_bQ^{-}j_a^{*} & b^{*}\end{pmatrix}
\]
as operators on $\mathcal{H}(\mathbb{U})\oplus\mathcal{H}(\mathbb{V})$,
and conjugation by the unitary $1\oplus U$ leaves the determinant unchanged.
The operator on the right is the one whose determinant defines $\mathfrak{D}(a,b)$ in Lemma~\ref{lem:reflection-determinants},
for $\mathcal{H}_a=\mathcal{H}(\mathbb{U})$, $\mathcal{H}_b=\mathcal{H}(\mathbb{V})$ and the Hermitian operators $Q^{+}$, $Q^{-}$ on $\mathcal{K}$;
this proves the asserted identity.
For $-1\le\rho\le1$ the matrices $Q^{\pm}_c$ are positive semidefinite,
so that $\langle\psi,Q^{\pm}\psi\rangle=\sum_{x}\psi(x)^{*}\cdot Q^{\pm}_c\cdot\psi(x)\ge0$ for $\psi\in\mathcal{K}$,
the sum being over a set of representatives of $(\mathbb{U}\cap\mathbb{V}\cap\mathbb{B}_c)/\Gamma_0$.
\end{proof}

\subsubsection{Rings and half slabs}\label{sec:rings}

\begin{lemma}[Cutting a ring]\label{lem:ring-cut}
Let $-1\le\rho\le1$, $M\in\R$ and $w_0,\dots,w_3>0$.
Let $\mu\in\{1,2,3\}$, let $\theta_c=\theta_{x_\mu,c}$ for $c\in\tfrac12\Z$, and, for a set $\mathbb{Y}$ of sites and real numbers $s<s'$, let $\mathbb{Y}(s,s')=\{x\in\mathbb{Y}:s<x_\mu<s'\}$, as in Definition~\ref{def:crossing-matrices}.
Let $n\ge1$ be an integer, $\Gamma_0'$ a subgroup of $\{g\in\Z^3:g_\mu=0\}$, $\Gamma_0:=\Gamma_0'+\Z ne_\mu$, and $\chi_0\colon\Gamma_0\to\{1,-1\}$ a character with
\[
\chi_0(ne_\mu)=-1 ;
\]
its restriction to $\Gamma_0'$ is again denoted by $\chi_0$.
Let $(\mathbb{Y},\Gamma_0,\chi_0)$ be a piece (Definition~\ref{def:pieces}) with $\theta_0\mathbb{Y}=\mathbb{Y}$.
Then $(\mathbb{Y}(0,n),\Gamma_0',\chi_0)$ is a piece, the determinants $\det D_{(\mathbb{Y},\Gamma_0,\chi_0)}$ and $\det D_{(\mathbb{Y}(0,n),\Gamma_0',\chi_0)}$ are real, and
\[
\det D_{(\mathbb{Y},\Gamma_0,\chi_0)}\ \ge\ \det D_{(\mathbb{Y}(0,n),\Gamma_0',\chi_0)} .
\]
\end{lemma}

\begin{proof}
Since the $\mu$-th coordinate of every element of $\Gamma_0'$ vanishes, each $g\in\Gamma_0$ is $g'+kne_\mu$ for exactly one pair $(g',k)\in\Gamma_0'\times\Z$.
The planes $\{x_\mu=c\}$ with $c\in\tfrac12\Z$ are mirrors (Definition~\ref{def:mirrors}), and no site lies on a mirror (Proposition~\ref{prop:geometry}\eqref{it:geometry-mirrors}).
Consequently $x_\mu\notin\tfrac n2\Z$ for every site $x$.

\emph{The piece of one period.}
The elements of $\Gamma_0'$ do not change the coordinate $x_\mu$, so $\mathbb{Y}(0,n)$ is invariant under $\Gamma_0'$.
For $y\in\mathbb{Y}$ there is exactly one $k\in\Z$ with $y_\mu+kn\in(0,n)$, because $y_\mu\notin n\Z$.
The $\Gamma_0$-orbit of $y$ therefore meets $\mathbb{Y}(0,n)$ in exactly one $\Gamma_0'$-orbit, and the inclusion of $\mathbb{Y}(0,n)$ in $\mathbb{Y}$ induces a bijection
\[
\mathbb{Y}(0,n)/\Gamma_0'\longrightarrow\mathbb{Y}/\Gamma_0 .
\]
So the quotient on the left is finite, and $(\mathbb{Y}(0,n),\Gamma_0',\chi_0)$ is a piece.

\emph{Two halves.}
Put
\[
\mathbb{Y}_+:=\left\{x\in\mathbb{Y}:x_\mu\in\left(0,\frac n2\right)+n\Z\right\},
\qquad
\mathbb{Y}_-:=\left\{x\in\mathbb{Y}:x_\mu\in\left(-\frac n2,0\right)+n\Z\right\}.
\]
As no coordinate $x_\mu$ of a site lies in $\tfrac n2\Z$, the set $\mathbb{Y}$ is the disjoint union of $\mathbb{Y}_+$ and $\mathbb{Y}_-$.
Both sets are $\Gamma_0$-invariant, and $\theta_0\mathbb{Y}_+=\mathbb{Y}_-$, because $\theta_0\mathbb{Y}=\mathbb{Y}$ and $\theta_0$ changes the sign of $x_\mu$.
Hence $(\mathbb{Y}_+,\Gamma_0,\chi_0)$ and $(\mathbb{Y}_-,\Gamma_0,\chi_0)$ are pieces.
In this proof $\mathcal{H}(\mathbb{Y})$ and $\mathcal{H}(\mathbb{Y}_\pm)$ denote the spaces of the pieces $(\mathbb{Y},\Gamma_0,\chi_0)$ and $(\mathbb{Y}_\pm,\Gamma_0,\chi_0)$.
Since $\mathbb{Y}/\Gamma_0$ is the disjoint union of $\mathbb{Y}_+/\Gamma_0$ and $\mathbb{Y}_-/\Gamma_0$, restriction to the two halves identifies $\mathcal{H}(\mathbb{Y})$ with the orthogonal sum $\mathcal{H}(\mathbb{Y}_+)\oplus\mathcal{H}(\mathbb{Y}_-)$.

\emph{A unitary.}
For $\psi'\in\mathcal{H}(\mathbb{Y}_+)$ and $x'\in\mathbb{Y}_-$ put
\[
(U\psi')(x'):=\gamma_\mu\,\psi'(\theta_0x').
\]
For $g'\in\Gamma_0'$ and $k\in\Z$ one has $\theta_0(x'+g'+kne_\mu)=\theta_0x'+g'-kne_\mu$, hence
\[
(U\psi')(x'+g'+kne_\mu)=\chi_0(g'-kne_\mu)\cdot(U\psi')(x')=\chi_0(g'+kne_\mu)\cdot(U\psi')(x'),
\]
because $\chi_0$ takes the values $\pm1$.
Thus $U$ maps $\mathcal{H}(\mathbb{Y}_+)$ into $\mathcal{H}(\mathbb{Y}_-)$.
The linear part of $\theta_0$ maps $\Gamma_0$ onto itself, so $\theta_0$ induces a bijection of $\mathbb{Y}_-/\Gamma_0$ onto $\mathbb{Y}_+/\Gamma_0$; as $\gamma_\mu$ is unitary, $U$ is unitary, and
\[
(U^{-1}\psi)(x)=\gamma_\mu\,\psi(\theta_0x),\qquad \psi\in\mathcal{H}(\mathbb{Y}_-),\ x\in\mathbb{Y}_+ .
\]

\emph{Blocks.}
Let $D:=D_{(\mathbb{Y},\Gamma_0,\chi_0)}$ and $a:=D_{(\mathbb{Y}_+,\Gamma_0,\chi_0)}$.
With respect to $\mathcal{H}(\mathbb{Y}_+)\oplus\mathcal{H}(\mathbb{Y}_-)$ the operator $D$ has the diagonal blocks $D_{++}=a$ and $D_{--}=D_{(\mathbb{Y}_-,\Gamma_0,\chi_0)}$,
the compressions of $D$ to the two halves,
and the off-diagonal blocks $D_{+-}\colon\mathcal{H}(\mathbb{Y}_-)\to\mathcal{H}(\mathbb{Y}_+)$ and $D_{-+}\colon\mathcal{H}(\mathbb{Y}_+)\to\mathcal{H}(\mathbb{Y}_-)$,
\[
\begin{aligned}
(D_{+-}\psi)(x)&=-\sum_{i:\,x^{(i)}\in\mathbb{Y}_-}w(x,x^{(i)})\,\psi(x^{(i)}),\\
(D_{-+}\psi')(x')&=-\sum_{i:\,{x'}^{(i)}\in\mathbb{Y}_+}w(x',{x'}^{(i)})\,\psi'({x'}^{(i)}),
\end{aligned}
\]
for $x\in\mathbb{Y}_+$ and $x'\in\mathbb{Y}_-$.
The operator $(1\oplus U)^{-1}D(1\oplus U)$ on $\mathcal{H}(\mathbb{Y}_+)\oplus\mathcal{H}(\mathbb{Y}_+)$ has the same determinant as $D$ and the blocks $a$, $D_{+-}U$, $U^{-1}D_{-+}$ and $U^{-1}D_{--}U$.

\emph{The lower right block.}
The neighbour relation commutes with the elements of $W$: $(gx)^{(i)}=g\,x^{(i)}$ for $g\in W$ and every site $x$ (Definition~\ref{def:fermion-sites}).
In particular $(\theta_0x)^{(i)}=\theta_0x^{(i)}$, and this site lies in $\mathbb{Y}_-$ if and only if $x^{(i)}\in\mathbb{Y}_+$.
For $\psi'\in\mathcal{H}(\mathbb{Y}_+)$ and $x\in\mathbb{Y}_+$ one obtains, with $\gamma_\mu^2=1$,
\[
\begin{aligned}
(U^{-1}D_{--}U\psi')(x)
&=\gamma_\mu\left(M\gamma_\mu\psi'(x)-\sum_{i:\,x^{(i)}\in\mathbb{Y}_+}w(\theta_0x,\theta_0x^{(i)})\,\gamma_\mu\psi'(x^{(i)})\right)\\
&=M\psi'(x)-\sum_{i:\,x^{(i)}\in\mathbb{Y}_+}w(x^{(i)},x)\,\psi'(x^{(i)}).
\end{aligned}
\]
The second equality holds by Lemma~\ref{lem:hopping}(d), applied to the mirror $\{x_\mu=0\}$, its upward unit normal $e_\mu$ and the bond $\{\theta_0x,\theta_0x^{(i)}\}$:
since $\theta_0^2=\mathrm{id}$, that lemma gives $\gamma_\mu\,w(\theta_0x,\theta_0x^{(i)})\,\gamma_\mu=w(x^{(i)},x)$.
By Definition~\ref{def:pieces}\eqref{it:piece-adjoint} the right side is $(a^*\psi')(x)$.
So $U^{-1}D_{--}U=a^*$; this is the computation of the last of the blocks \eqref{eq:positivity-mirror-blocks}, with $\theta_0$ in the place of $\theta_c$.

\emph{Bonds between the halves.}
Let $x\in\mathbb{Y}_+$, and let $k\in\Z$ be the integer with $kn<x_\mu<kn+n/2$.
A neighbour $x^{(i)}$ of $x$ lies in $\mathbb{Y}_-$ if and only if $x^{(i)}=\theta_cx$ for some $c\in\{kn,\,kn+n/2\}$.
Indeed, if $x^{(i)}\in\mathbb{Y}_-$, then $(x^{(i)})_\mu$ lies outside the closed interval $[kn,\,kn+n/2]$.
So $x$ and $x^{(i)}$ lie on the two sides of the plane $\{x_\mu=c\}$ for $c=kn$ or for $c=kn+n/2$.
This plane is a mirror, and $x^{(i)}=\theta_cx$ by Proposition~\ref{prop:geometry}\eqref{it:geometry-mirrors}.
Conversely, let $c\in\{kn,\,kn+n/2\}$.
Definition~\ref{def:mirrors}\eqref{it:mirror-product} gives
\[
\theta_c=\tau_{2ce_\mu}\theta_0,\qquad \theta_0\theta_c=\tau_{-2ce_\mu},
\]
and $2ce_\mu\in n\Z e_\mu\subset\Gamma_0$; together with $\theta_0\mathbb{Y}=\mathbb{Y}$ this yields $\theta_c\mathbb{Y}=\mathbb{Y}$.
The coordinate $2c-x_\mu$ of $\theta_cx$ lies in $(kn-n/2,\,kn)$ if $c=kn$, and in $(kn+n/2,\,kn+n)$ if $c=kn+n/2$; in both cases $\theta_cx\in\mathbb{Y}_-$.
Moreover, if $\theta_cx$ is a neighbour of $x$ for some $c\in n\Z$, then $c=kn$, and if it is one for some $c\in\tfrac n2+n\Z$, then $c=kn+n/2$.
Indeed, the open segment from $x$ to $\theta_cx$ meets the mirror $\{x_\mu=c\}$.
If $c$ were not an end point of the interval $(kn,\,kn+n/2)$, the segment would also meet the mirror $\{x_\mu=c'\}$, where $c'$ is the end point that lies between $x_\mu$ and $c$;
but the open segment of a bond meets exactly one mirror (Proposition~\ref{prop:geometry}\eqref{it:geometry-midpoint}).

Now let $c\in\{kn,\,kn+n/2\}$ be such that $\theta_cx$ is a neighbour of $x$.
One has $\chi_0(2ce_\mu)=1$ if $c=kn$, since then $2ce_\mu=2k\cdot ne_\mu$, and $\chi_0(2ce_\mu)=-1$ if $c=kn+n/2$, since then $2ce_\mu=(2k+1)\cdot ne_\mu$.
As $\theta_0\theta_cx=x-2ce_\mu$, it follows that for $\psi',\psi\in\mathcal{H}(\mathbb{Y}_+)$
\[
(U\psi')(\theta_cx)=\chi_0(2ce_\mu)\cdot\gamma_\mu\psi'(x),
\qquad
\psi(\theta_0\theta_cx)=\chi_0(2ce_\mu)\cdot\psi(x).
\]
The bond $\{x,\theta_cx\}$ contributes the term $-\chi_0(2ce_\mu)\cdot w(x,\theta_cx)\,\gamma_\mu\psi'(x)$ to $(D_{+-}U\psi')(x)$.
Its image $\{\theta_0x,\theta_0\theta_cx\}$ under $\theta_0$ joins $\theta_0x$ to a site of $\mathbb{Y}_+$,
and it contributes the term $-\chi_0(2ce_\mu)\cdot\gamma_\mu\,w(\theta_0x,\theta_0\theta_cx)\,\psi(x)$ to $(U^{-1}D_{-+}\psi)(x)=\gamma_\mu(D_{-+}\psi)(\theta_0x)$.
By the preceding paragraph and the equivariance of the neighbours, $(D_{+-}U\psi')(x)$ and $(U^{-1}D_{-+}\psi)(x)$ are the sums of these terms over those $c\in\{kn,\,kn+n/2\}$ for which $\theta_cx$ is a neighbour of $x$.

\emph{The lower plane, $c=kn$.}
Here $x\in\mathbb{B}_c$, $\chi_0(2ce_\mu)=1$, and $Q^{\pm}_c=Q^{\pm}_0$ because $c\in\Z$.
By \eqref{eq:positivity-crossing}, $-w(x,\theta_cx)\gamma_\mu=-Q^{+}_0$.
Moreover $\theta_0x=\theta_cx-2ce_\mu$ and $\theta_0\theta_cx=x-2ce_\mu$.
A translation by a vector of $\Z^3$ preserves the type and the direction of a bond, so that $w(\theta_0x,\theta_0\theta_cx)=w(\theta_cx,x)$;
and $-\gamma_\mu w(\theta_cx,x)=Q^{-}_0$, again by \eqref{eq:positivity-crossing}.
The lower plane thus contributes $-Q^{+}_0\psi'(x)$ to $(D_{+-}U\psi')(x)$ and $Q^{-}_0\psi(x)$ to $(U^{-1}D_{-+}\psi)(x)$, exactly as in the off-diagonal blocks of \eqref{eq:positivity-mirror-blocks}.

\emph{The upper plane, $c=kn+n/2$.}
Here $\chi_0(2ce_\mu)=-1$, and $Q^{\pm}_c=Q^{\pm}_{n/2}$ because $c-n/2\in\Z$.
The bond $\{x,\theta_cx\}$ crosses the mirror $\{x_\mu=c\}$, which is a site plane if $c\in\Z$ and a link plane if $c\in\Z+\tfrac12$.
By Proposition~\ref{prop:geometry}\eqref{it:geometry-projection} the bond has the type~$0$ in the first case and the type~$3$ in the second,
so that in either case its weight is the number $w^{\times}_c$ of Definition~\ref{def:crossing-matrices}.
The direction from $x$ to $\theta_cx$ is $+e_\mu$, and Definition~\ref{def:chamber-operator} gives
\[
w(x,\theta_cx)=\frac{w^{\times}_c}{2}(\rho-\gamma_\mu).
\]
Its image under $\theta_0$ is a bond of the same type, and the direction from $\theta_0x$ to $\theta_0\theta_cx$ is $-e_\mu$; hence $w(\theta_0x,\theta_0\theta_cx)=\tfrac{w^{\times}_c}{2}(\rho+\gamma_\mu)$.
The two contributions are
\[
\begin{aligned}
w(x,\theta_cx)\,\gamma_\mu\psi'(x)&=-\frac{w^{\times}_c}{2}(1-\rho\gamma_\mu)\psi'(x)=-Q^{-}_c\psi'(x),\\
\gamma_\mu\,w(\theta_0x,\theta_0\theta_cx)\,\psi(x)&=\frac{w^{\times}_c}{2}(1+\rho\gamma_\mu)\psi(x)=Q^{+}_c\psi(x).
\end{aligned}
\]
Seen from $\mathbb{Y}_+$, the normal $e_\mu$ points inwards at the planes of $n\Z$ and outwards at those of $\tfrac n2+n\Z$.
At the upper plane this exchanges the roles of $Q^{+}$ and $Q^{-}$ and reverses both signs: with $1$ in the place of $\chi_0(2ce_\mu)$ the contributions would be $+Q^{-}_c\psi'(x)$ and $-Q^{+}_c\psi(x)$.
The value $\chi_0(ne_\mu)=-1$ restores the signs; this is the sign at the second mirror in \eqref{eq:chamber-action-split}, written for matrices.

\emph{The block form.}
For $\psi\in\mathcal{H}(\mathbb{Y}_+)$ and $x\in\mathbb{Y}_+$ let $(R_0\psi)(x):=Q^{+}_0\psi(x)$ and $(R'_0\psi)(x):=Q^{-}_0\psi(x)$ if $\theta_cx$ is a neighbour of $x$ for some $c\in n\Z$, and $(R_0\psi)(x):=(R'_0\psi)(x):=0$ otherwise.
Likewise let $(R_1\psi)(x):=Q^{-}_{n/2}\psi(x)$ and $(R'_1\psi)(x):=Q^{+}_{n/2}\psi(x)$ if $\theta_cx$ is a neighbour of $x$ for some $c\in\tfrac n2+n\Z$, and $(R_1\psi)(x):=(R'_1\psi)(x):=0$ otherwise.
Both conditions on $x$ are invariant under $\Gamma_0$:
for $g\in\Gamma_0$ one has $\theta_{c+g_\mu}(x+g)=\theta_cx+g$ by Definition~\ref{def:mirrors}\eqref{it:mirror-conjugation},
the site $\theta_cx+g$ is a neighbour of $x+g$ if $\theta_cx$ is one of $x$, and $g_\mu\in n\Z$.
So $R_0$, $R'_0$, $R_1$, $R'_1$ are operators on $\mathcal{H}(\mathbb{Y}_+)$.
By the computations at the lower and at the upper plane, $D_{+-}U=-(R_0+R_1)$ and $U^{-1}D_{-+}=R'_0+R'_1$.
Therefore
\begin{equation}\label{eq:positivity-ring-blocks}
\det D=\det\begin{pmatrix} a & -(R_0+R_1)\\ R'_0+R'_1 & a^{*}\end{pmatrix}.
\end{equation}
The matrix $\gamma_\mu$ is Hermitian with the eigenvalues $1$ and $-1$ (Lemma~\ref{lem:spin-identities}(c)).
So the crossing matrices $Q^{\pm}_c=\tfrac{w^{\times}_c}{2}(1\pm\rho\gamma_\mu)$ of Definition~\ref{def:crossing-matrices} are Hermitian with the eigenvalues $\tfrac{w^{\times}_c}{2}(1+\rho)$ and $\tfrac{w^{\times}_c}{2}(1-\rho)$,
which are non-negative because $|\rho|\le1$.
Hence $R_0$, $R'_0$, $R_1$ and $R'_1$ are Hermitian and non-negative.

\emph{Omitting the bonds across $n\Z$.}
Lemma~\ref{lem:reflection-determinants} applies with $\mathcal{H}_a=\mathcal{H}_b=\mathcal{K}:=\mathcal{H}(\mathbb{Y}_+)$ and $j_a=j_b:=1$: the right side of \eqref{eq:positivity-ring-blocks} is $\mathfrak{D}(a,a)$ for $Q^{+}:=R_0+R_1$ and $Q^{-}:=R'_0+R'_1$.
Lemma~\ref{lem:reflection-determinants}\eqref{it:determinant-monotone} is applied twice:
first to $Q^{+}=R_1$, $Q^{-}=R'_1$, with $R_1$ replaced by $R_0+R_1\ge R_1$;
then to $Q^{+}=R_0+R_1$, $Q^{-}=R'_1$, with $R'_1$ replaced by $R'_0+R'_1\ge R'_1$.
Neither replacement decreases the determinant, so that
\[
\det D\ \ge\ \det\begin{pmatrix} a & -R_1\\ R'_1 & a^{*}\end{pmatrix},
\]
and both sides are real by Lemma~\ref{lem:reflection-determinants}\eqref{it:determinant-lower}.

\emph{The cut operator.}
Let $D^{\mathrm{cut}}$ be the operator on $\mathcal{H}(\mathbb{Y})$ defined by
\[
(D^{\mathrm{cut}}\psi)(x):=M\psi(x)-\sum_{i}w(x,x^{(i)})\,\psi(x^{(i)}),
\]
where the sum is over those $i$ for which $x^{(i)}\in\mathbb{Y}$ and the bond $\{x,x^{(i)}\}$ meets no plane $\{x_\mu=c\}$ with $c\in n\Z$;
this condition is invariant under the translations by $\Gamma_0$, since $g_\mu\in n\Z$ for $g\in\Gamma_0$, so that $D^{\mathrm{cut}}$ maps $\mathcal{H}(\mathbb{Y})$ into itself.
A bond with both ends in $\mathbb{Y}_+$, or both in $\mathbb{Y}_-$, meets no plane $\{x_\mu=c\}$ with $c\in\tfrac n2\Z$.
Otherwise its ends would lie in two different slabs of width $n/2$ of the same half, separated by a slab of the other half,
and the open bond would meet at least two mirrors, which Proposition~\ref{prop:geometry}\eqref{it:geometry-midpoint} excludes.
So $D^{\mathrm{cut}}$ has the diagonal blocks $a$ and $D_{--}$ of $D$.
In its off-diagonal blocks exactly the bonds across the planes of $n\Z$ are missing:
at $x\in\mathbb{Y}_+$ these are the bonds that gave $R_0$,
and at $\theta_0x$ their images under $\theta_0$, which gave $R'_0$, because $\theta_0$ maps the planes of $n\Z$ onto themselves.
Hence $(1\oplus U)^{-1}D^{\mathrm{cut}}(1\oplus U)$ has the blocks $a$, $-R_1$, $R'_1$, $a^*$, and
\[
\det D^{\mathrm{cut}}=\det\begin{pmatrix} a & -R_1\\ R'_1 & a^{*}\end{pmatrix}.
\]

\emph{Restriction to one period.}
For $\psi\in\mathcal{H}(\mathbb{Y})$ the restriction $\psi|_{\mathbb{Y}(0,n)}$ lies in $\mathcal{H}(\mathbb{Y}(0,n),\Gamma_0',\chi_0)$,
and the map $\psi\mapsto\psi|_{\mathbb{Y}(0,n)}$ is isometric, by the bijection $\mathbb{Y}(0,n)/\Gamma_0'\to\mathbb{Y}/\Gamma_0$.
It is onto.
Indeed, every $y\in\mathbb{Y}$ is $x+kne_\mu$ with unique $x\in\mathbb{Y}(0,n)$ and $k\in\Z$.
For $\psi'\in\mathcal{H}(\mathbb{Y}(0,n),\Gamma_0',\chi_0)$ the spinor $\psi(x+kne_\mu):=(-1)^k\psi'(x)$ lies in $\mathcal{H}(\mathbb{Y})$,
because $\chi_0(g'+kne_\mu)=(-1)^k\chi_0(g')$ for $g'\in\Gamma_0'$, and its restriction is $\psi'$.
Let $x\in\mathbb{Y}(0,n)$.
A neighbour $x^{(i)}\in\mathbb{Y}$ lies in $\mathbb{Y}(0,n)$ if and only if the bond $\{x,x^{(i)}\}$ meets no plane $\{x_\mu=c\}$ with $c\in n\Z$, since the slab $\{0<x_\mu<n\}$ is convex and no site lies on its boundary planes.
Hence
\[
(D^{\mathrm{cut}}\psi)|_{\mathbb{Y}(0,n)}=D_{(\mathbb{Y}(0,n),\Gamma_0',\chi_0)}\left(\psi|_{\mathbb{Y}(0,n)}\right),\qquad \psi\in\mathcal{H}(\mathbb{Y}),
\]
so the restriction is a unitary map that intertwines $D^{\mathrm{cut}}$ with $D_{(\mathbb{Y}(0,n),\Gamma_0',\chi_0)}$, and the two operators have the same determinant.

Putting the three comparisons together,
\[
\det D\ \ge\ \det\begin{pmatrix} a & -R_1\\ R'_1 & a^{*}\end{pmatrix}=\det D^{\mathrm{cut}}=\det D_{(\mathbb{Y}(0,n),\Gamma_0',\chi_0)},
\]
and all these numbers are real.
\end{proof}

\begin{lemma}[Reduction to a half slab]\label{lem:half-slab}
Let $-1\le\rho\le1$, $M\in\R$ and $w_0,\dots,w_3>0$.
Let $\mu\in\{1,2,3\}$, let $\theta_c=\theta_{x_\mu,c}$ for $c\in\tfrac12\Z$, and, for a set $\mathbb{Y}$ of sites and real numbers $s<s'$, let $\mathbb{Y}(s,s')=\{x\in\mathbb{Y}:s<x_\mu<s'\}$, as in Definition~\ref{def:crossing-matrices}.
Let $n\ge1$ be an integer, $\Gamma_0'$ a subgroup of $\{g\in\Z^3:g_\mu=0\}$, $\Gamma_0:=\Gamma_0'+\Z ne_\mu$, and $\chi_0\colon\Gamma_0\to\{1,-1\}$ a character with $\chi_0(ne_\mu)=-1$; its restriction to $\Gamma_0'$ is again denoted by $\chi_0$.
Let $(\mathbb{Y},\Gamma_0,\chi_0)$ be a piece such that $\theta_c\mathbb{Y}=\mathbb{Y}$ for every $c\in\tfrac12\Z$.
Then $(\mathbb{Y}(0,1/2),\Gamma_0',\chi_0)$ is a piece, and
\[
\det D_{(\mathbb{Y},\Gamma_0,\chi_0)}\ \ge\ \left|\det D_{(\mathbb{Y}(0,1/2),\Gamma_0',\chi_0)}\right|^{2n}.
\]
\end{lemma}

\begin{proof}
By Definition~\ref{def:mirrors}\eqref{it:mirror-product}, $\theta_{1/2}\theta_0=\tau_{e_\mu}$; hence $\tau_{e_\mu}\mathbb{Y}=\mathbb{Y}$.
For real $s<s'$ the set $\mathbb{Y}(s,s')$ is $\Gamma_0'$-invariant.
Two of its sites that lie in one $\Gamma_0$-orbit differ by a vector $g'+kne_\mu$ with $g'\in\Gamma_0'$ and $|k|n<s'-s$.
So every $\Gamma_0$-orbit meets $\mathbb{Y}(s,s')$ in finitely many $\Gamma_0'$-orbits, and $\mathbb{Y}(s,s')/\Gamma_0'$ is finite, as $\mathbb{Y}/\Gamma_0$ is.
Thus $(\mathbb{Y}(s,s'),\Gamma_0',\chi_0)$ is a piece for all real $s<s'$; for $s=0$ and $s'=1/2$ this is the first assertion.
Put
\[
d_0:=1,\qquad d_k:=\det D_{(\mathbb{Y}(0,k),\Gamma_0',\chi_0)}\quad(k=1,2,\dots).
\]

\emph{Step 1.}
Since $\theta_0\mathbb{Y}=\mathbb{Y}$, Lemma~\ref{lem:ring-cut} gives
\[
\det D_{(\mathbb{Y},\Gamma_0,\chi_0)}\ \ge\ d_n .
\]

\emph{Step 2: every $d_k$ is real and non-negative, and $d_1\ge\left|\det D_{(\mathbb{Y}(0,1/2),\Gamma_0',\chi_0)}\right|^2$.}
For $k=0$ there is nothing to prove.
Let $k\ge1$, $c:=k/2$, and $\mathbb{U}=\mathbb{V}:=\mathbb{Y}(k/2,k)$, a subset of $\{x_\mu>c\}$.
Since $\theta_c\mathbb{Y}=\mathbb{Y}$ and $\theta_c$ maps $\{k/2<x_\mu<k\}$ onto $\{0<x_\mu<k/2\}$, one has $\theta_c\mathbb{V}=\mathbb{Y}(0,k/2)$; and $\mathbb{U}\cup\theta_c\mathbb{V}=\mathbb{Y}(0,k)$, because no site lies on the mirror $\{x_\mu=c\}$.
Lemma~\ref{lem:gluing}, for the group $\Gamma_0'$, gives $d_k=\mathfrak{D}(a,a)$ with $a:=D_{(\mathbb{Y}(k/2,k),\Gamma_0',\chi_0)}$ and with $Q^{\pm}$ the multiplications by $Q^{\pm}_c$, which are non-negative because $|\rho|\le1$.
By Lemma~\ref{lem:reflection-determinants}\eqref{it:determinant-lower}, $d_k$ is real and $d_k\ge|\det a|^2\ge0$.
Let now $k=1$, so that $c=1/2$ and $\mathbb{V}=\mathbb{Y}(1/2,1)$.
The operator $D_{(\mathbb{Y}(0,1/2),\Gamma_0',\chi_0)}=D_{(\theta_{1/2}\mathbb{V},\Gamma_0',\chi_0)}$ is unitarily equivalent to $a^*$, by the last of the blocks \eqref{eq:positivity-mirror-blocks}, in which here $b=a$.
Hence $|\det a|=\left|\det D_{(\mathbb{Y}(0,1/2),\Gamma_0',\chi_0)}\right|$, and $d_1\ge\left|\det D_{(\mathbb{Y}(0,1/2),\Gamma_0',\chi_0)}\right|^2$.

\emph{Step 3: for every $k\ge1$,}
\begin{equation}\label{eq:positivity-log-convex}
d_k^2\ \le\ d_{k-1}\cdot d_{k+1}.
\end{equation}
Let $c:=(k-1)/2$, $\mathbb{U}:=\mathbb{Y}(c,k)$ and $a:=D_{(\mathbb{U},\Gamma_0',\chi_0)}$.
The reflection $\theta_c$ maps $\{c<x_\mu<k\}$ onto $\{-1<x_\mu<c\}$; as no site lies on $\{x_\mu=c\}$ and $\tau_{e_\mu}\mathbb{Y}=\mathbb{Y}$,
\[
\mathbb{U}\cup\theta_c\mathbb{U}=\mathbb{Y}(-1,k)=\mathbb{Y}(0,k+1)-e_\mu .
\]
By Definition~\ref{def:pieces}\eqref{it:piece-translation} and Lemma~\ref{lem:gluing}, $d_{k+1}=\mathfrak{D}(a,a)$, formed with the space $\mathcal{K}=\mathcal{H}(\mathbb{U}\cap\mathbb{B}_c,\Gamma_0',\chi_0)$ and with the multiplications by $Q^{\pm}_c\ge0$.

If $k=1$, then $c=0$, $\mathbb{U}=\mathbb{Y}(0,1)$ and $d_1=\det a$.
Lemma~\ref{lem:reflection-determinants}\eqref{it:determinant-lower} gives $d_2=\mathfrak{D}(a,a)\ge|\det a|^2=d_1^2$, which is \eqref{eq:positivity-log-convex}, as $d_0=1$.

Let $k\ge2$, $\mathbb{V}:=\mathbb{Y}(c,k-1)$ and $b:=D_{(\mathbb{V},\Gamma_0',\chi_0)}$.
Then $\theta_c\mathbb{V}=\mathbb{Y}(0,c)$, so that $\mathbb{U}\cup\theta_c\mathbb{V}=\mathbb{Y}(0,k)$ and $\mathbb{V}\cup\theta_c\mathbb{V}=\mathbb{Y}(0,k-1)$.
Lemma~\ref{lem:gluing} gives $d_k=\mathfrak{D}(a,b)$, formed with $\mathcal{H}(\mathbb{U}\cap\mathbb{V}\cap\mathbb{B}_c,\Gamma_0',\chi_0)$, and $d_{k-1}=\mathfrak{D}(b,b)$, formed with $\mathcal{H}(\mathbb{V}\cap\mathbb{B}_c,\Gamma_0',\chi_0)$.
Every $x\in\mathbb{B}_c$ satisfies $x_\mu>c$, and its neighbour $\theta_cx=x-2(x_\mu-c)e_\mu$ lies at the distance $2(x_\mu-c)$ from it (Definition~\ref{def:crossing-matrices}),
a distance smaller than $1$ by Lemma~\ref{lem:neighbours}.
Thus $\mathbb{B}_c\subset\{c<x_\mu<c+1/2\}$; as $c+1/2=k/2\le k-1$ for $k\ge2$, one has $\mathbb{Y}\cap\mathbb{B}_c\subset\mathbb{Y}(c,c+1/2)\subset\mathbb{V}\subset\mathbb{U}$, and therefore
\[
\mathbb{U}\cap\mathbb{V}\cap\mathbb{B}_c=\mathbb{U}\cap\mathbb{B}_c=\mathbb{V}\cap\mathbb{B}_c=\mathbb{Y}\cap\mathbb{B}_c .
\]
The three determinants are thus $\mathfrak{D}(a,b)$, $\mathfrak{D}(a,a)$ and $\mathfrak{D}(b,b)$ in the sense of Lemma~\ref{lem:reflection-determinants},
formed with one space $\mathcal{K}:=\mathcal{H}(\mathbb{Y}\cap\mathbb{B}_c,\Gamma_0',\chi_0)$,
the extensions by zero $j_a\colon\mathcal{K}\to\mathcal{H}(\mathbb{U},\Gamma_0',\chi_0)$ and $j_b\colon\mathcal{K}\to\mathcal{H}(\mathbb{V},\Gamma_0',\chi_0)$,
and the same operators $Q^{\pm}\ge0$, the multiplications by $Q^{\pm}_c$.
Lemma~\ref{lem:reflection-determinants}\eqref{it:determinant-schwarz} yields $|d_k|^2\le d_{k+1}\cdot d_{k-1}$, and $d_k$ is real by Step~2.

\emph{Step 4: $d_k\ge d_1^k$ for every $k\ge0$.}
If $d_1=0$, this follows from Step~2 and $d_0=1$.
Let $d_1>0$.
Induction on $k\ge1$ shows that $d_k>0$ and $d_k\ge d_1\cdot d_{k-1}$.
For $k=1$ both statements hold, the second with equality, since $d_0=1$.
If they hold for $1,\dots,k$, then $d_{k-1}>0$ (for $k=1$ because $d_0=1$), and \eqref{eq:positivity-log-convex} gives
\[
d_{k+1}\ \ge\ \frac{d_k^2}{d_{k-1}}=d_k\cdot\frac{d_k}{d_{k-1}}\ \ge\ d_k\cdot d_1\ >\ 0 .
\]
Repeated use of $d_k\ge d_1\cdot d_{k-1}$, starting from $d_0=1$, yields $d_k\ge d_1^k$.

Steps~1, 4 and~2, in this order, give
\[
\det D_{(\mathbb{Y},\Gamma_0,\chi_0)}\ \ge\ d_n\ \ge\ d_1^n\ \ge\ \left|\det D_{(\mathbb{Y}(0,1/2),\Gamma_0',\chi_0)}\right|^{2n}.
\]
\end{proof}

\subsubsection{The hexagon}\label{sec:hexagon}

Let
\[
H:=\mathbb{X}\cap\left(0,\tfrac12\right)^3
\]
be the set of the sites in the open cube $\left(0,\tfrac12\right)^3$.
Every site has the form $x=\pi(x)+y(x)$ with $\pi(x)\in\Z^3$,
and the coordinates of $y(x)$ are $\pm v_1$, $\pm v_2$, $\pm v_3$ in some order (Definition~\ref{def:fermion-sites}).
Since $0<v_3<v_2<v_1<\frac12$, every coordinate of $y(x)$ has modulus less than $\frac12$;
hence the coordinate $x_\mu=\pi(x)_\mu+y_\mu(x)$ lies in $\left(0,\tfrac12\right)$ if and only if $\pi(x)_\mu=0$ and $y_\mu(x)>0$.
Conversely, every permutation matrix belongs to $\Ocub\subset W$ (Definition~\ref{def:mirrors})
and therefore maps $v_0$ to a site, as $\mathbb{X}=Wv_0$.
It follows that $H$ consists of the six points whose coordinates are the permutations of $(v_1,v_2,v_3)$.
They are numbered cyclically by $k\in\Z/6$:
\[
\begin{aligned}
&z_0=(v_1,v_2,v_3), && z_1=(v_1,v_3,v_2), && z_2=(v_2,v_3,v_1),\\
&z_3=(v_3,v_2,v_1), && z_4=(v_3,v_1,v_2), && z_5=(v_2,v_1,v_3).
\end{aligned}
\]

The neighbours of these points are read off from Lemma~\ref{lem:neighbours}.
For $x\in H$, so that $\pi(x)=0$ and all coordinates of $x$ are positive,
the neighbour $x^{(0)}$ is obtained by changing the sign of the coordinate of modulus $v_3$,
and $x^{(3)}$ by replacing the coordinate $v_1$ by $1-v_1$;
neither lies in the cube, since $x^{(0)}$ has the coordinate $-v_3<0$ and $x^{(3)}$ the coordinate $1-v_1>\frac12$.
The neighbour $x^{(1)}$, resp.\ $x^{(2)}$, is obtained by exchanging the places of the moduli $v_2$ and $v_3$, resp.\ of $v_1$ and $v_2$,
while the signs stay where they are.
Thus
\[
z_0^{(1)}=z_1,\quad z_2^{(1)}=z_3,\quad z_4^{(1)}=z_5,\qquad
z_1^{(2)}=z_2,\quad z_3^{(2)}=z_4,\quad z_5^{(2)}=z_0 ,
\]
and, as $(x^{(i)})^{(i)}=x$ for every site $x$ and every type $i$ (Definition~\ref{def:fermion-sites}),
the neighbours of $z_k$ that lie in $H$ are exactly $z_{k-1}$ and $z_{k+1}$.
The bonds with both ends in $H$ are therefore the six bonds $\{z_k,z_{k+1}\}$, $k\in\Z/6$.
They form a cycle of length six, the hexagon,
and the bond $\{z_k,z_{k+1}\}$ has the type $1$ if $k$ is even and the type $2$ if $k$ is odd
(the parity of $k\in\Z/6$ is well defined).

\begin{definition}[The hexagon operator]\label{def:hexagon}
For $k\in\Z/6$ the bond $\{z_k,z_{k+1}\}$ has the weight
\[
\tilde w_k:=\begin{cases} w_1, & k\ \text{even},\\ w_2, & k\ \text{odd}.\end{cases}
\]
Let $\hat d_k$ be the unit vector from $z_k$ to $z_{k+1}$, and let $\gamma_{(k)}:=\gamma\cdot\hat d_k$.
The set $H$ is finite;
so $(H,\{0\},1)$, formed with the trivial group $\{0\}$ and its trivial character $1$, is a piece (Definition~\ref{def:pieces}).
The hexagon operator is the operator of this piece,
\[
D_H:=D_{(H,\{0\},1)} ;
\]
it acts on the space $\mathcal{H}(H,\{0\},1)$ of all maps $\psi\colon H\to\C^4$.
\end{definition}

The differences of consecutive points are
\[
\begin{aligned}
z_1-z_0&=(v_2-v_3)\,(0,-1,1), & z_4-z_3&=(v_1-v_2)\,(0,1,-1),\\
z_2-z_1&=(v_1-v_2)\,(-1,0,1), & z_5-z_4&=(v_2-v_3)\,(1,0,-1),\\
z_3-z_2&=(v_2-v_3)\,(-1,1,0), & z_0-z_5&=(v_1-v_2)\,(1,-1,0),
\end{aligned}
\]
with the positive factors $v_1-v_2$ and $v_2-v_3$.
Consequently
\[
\sqrt{2}\,\hat d_0=(0,-1,1),\qquad
\sqrt{2}\,\hat d_1=(-1,0,1),\qquad
\sqrt{2}\,\hat d_2=(-1,1,0),\qquad
\hat d_{k+3}=-\hat d_k ,
\]
and
\[
\hat d_k\cdot\hat d_{k+1}=\frac12,\qquad
\hat d_{k+2}=\hat d_{k+1}-\hat d_k,\qquad
k\in\Z/6 .
\]
For $k=0,1,2$ both relations are verified on the three vectors displayed, together with $\hat d_3=-\hat d_0$ and $\hat d_4=-\hat d_1$;
replacing $k$ by $k+3$ changes the sign of every vector that occurs in them and leaves both relations intact.

By Definition~\ref{def:pieces}, $D_H$ is the compression to $H$ of the chamber operator $D_0$ at $A=0$:
\[
(D_H\psi)(x)=M\psi(x)-\sum_{i:\ x^{(i)}\in H}w(x,x^{(i)})\cdot\psi(x^{(i)}),\qquad x\in H .
\]
At $x=z_k$ the sum has two terms, one for the bond $\{z_k,z_{k+1}\}$, of the weight $\tilde w_k$,
and one for the bond $\{z_{k-1},z_k\}$, of the weight $\tilde w_{k-1}$;
the unit vector from $z_k$ to $z_{k-1}$ is $-\hat d_{k-1}$.
The hopping matrices of Definition~\ref{def:chamber-operator} therefore give, for $k\in\Z/6$,
\begin{equation}\label{eq:positivity-hexagon-operator}
\begin{split}
(D_H\psi)(z_k)={}&M\psi(z_k)-\frac{\tilde w_k}{2}\,(\rho-\gamma_{(k)})\,\psi(z_{k+1})\\
&-\frac{\tilde w_{k-1}}{2}\,(\rho+\gamma_{(k-1)})\,\psi(z_{k-1}) .
\end{split}
\end{equation}

\begin{lemma}[The determinant of the hexagon]\label{lem:hexagon-determinant}
Let $\rho$ and $M$ be real numbers, let $w_1,w_2>0$, let $D_H$ be the hexagon operator of Definition~\ref{def:hexagon}, and put
\[
s_H:=\frac{(1-\rho^2)(w_1^2+w_2^2)}{4},\qquad p_H:=\frac{w_1w_2}{2}.
\]
Then
\begin{equation}\label{eq:positivity-hexagon}
\begin{split}
\det D_H={}&\left(M^2+s_H+p_H\left(1+\frac{\rho^2}{2}\right)\right)^4\\
&\cdot\left(\left(M^2+s_H-\frac{p_H}{2}-\frac{p_H\rho^2}{4}\right)^2+\frac{9}{4}\,p_H^2\rho^2\left(1-\frac{\rho^2}{4}\right)\right)^4 .
\end{split}
\end{equation}
If $-1\le\rho\le1$, then $\det D_H>0$ for every real $M$.
If $\rho=\pm1$, then
\[
\det D_H=\left(M^6+\left(\frac{3w_1w_2}{4}\right)^3\right)^4 .
\]
\end{lemma}

\begin{proof}
All indices of $z_k$, $\tilde w_k$, $\hat d_k$ and $\gamma_{(k)}$ are read modulo $6$;
in this reading \eqref{eq:positivity-hexagon-operator} holds for every $k\in\Z$.

\emph{Spin algebra.}
Each $\hat d_k$ is a real unit vector.
By Definition~\ref{def:spin-matrices} the matrix $\gamma_{(k)}$ is therefore Hermitian,
and bilinearity turns the anticommutation relations of the $\gamma_\mu$ into
\[
\gamma_{(k)}^2=1,\qquad
\gamma_{(k)}\gamma_{(k+1)}+\gamma_{(k+1)}\gamma_{(k)}=2\,\hat d_k\cdot\hat d_{k+1}=1 .
\]
Let
\[
U:=\gamma_{(0)}\gamma_{(1)} .
\]
The relation $\hat d_{k+2}=\hat d_{k+1}-\hat d_k$ and the linearity of $v\mapsto\gamma\cdot v$ give $\gamma_{(k+2)}=\gamma_{(k+1)}-\gamma_{(k)}$, whence
\[
\gamma_{(k+1)}\gamma_{(k+2)}=\gamma_{(k+1)}\left(\gamma_{(k+1)}-\gamma_{(k)}\right)=1-\gamma_{(k+1)}\gamma_{(k)}=\gamma_{(k)}\gamma_{(k+1)} ;
\]
by induction, $\gamma_{(k)}\gamma_{(k+1)}=U$ for every $k$.
A Hermitian matrix with square $1$ is unitary.
So $U$ is unitary, $U^{-1}=\gamma_{(1)}\gamma_{(0)}$, and the anticommutation relation for $k=0$ reads $U+U^{-1}=1$, that is,
\begin{equation}\label{eq:positivity-spin-quadratic}
U^2-U+1=0 .
\end{equation}
Multiplication of \eqref{eq:positivity-spin-quadratic} by $U+1$ gives $U^3=-1$.
Multiplying $\gamma_{(k)}\gamma_{(k+1)}=U$ from the left by $\gamma_{(k)}$, resp.\ from the right by $\gamma_{(k+1)}$,
one obtains $\gamma_{(k+1)}=\gamma_{(k)}U$ and $\gamma_{(k)}=U\gamma_{(k+1)}$.
Consequently
\[
\gamma_{(k)}U=\gamma_{(k+1)}=U^{-1}\gamma_{(k)},
\]
and $\gamma_{(k)}=U^{-k}\gamma_{(0)}$ for every $k\ge0$, hence for every $k\in\Z$, as $U^6=1$.
Now put
\[
R:=\frac{1}{\sqrt{3}}\,(1+U).
\]
From $U^2=U-1$ and $U^*=U^{-1}$,
\[
R^2=\frac{1+2U+U^2}{3}=U,\qquad
RR^*=\frac{2+U+U^{-1}}{3}=1 .
\]
Thus $R$ is unitary, $R^{-1}=R^*=\frac{1}{\sqrt{3}}\,(1+U^{-1})$, and
\[
R+R^{-1}=\frac{2+U+U^{-1}}{\sqrt{3}}=\sqrt{3},\qquad
R^6=U^3=-1 .
\]
The relation $\gamma_{(0)}U=U^{-1}\gamma_{(0)}$ gives $\gamma_{(0)}(1+U)=(1+U^{-1})\gamma_{(0)}$, that is, $\gamma_{(0)}R=R^{-1}\gamma_{(0)}$.
Hence $\gamma_{(0)}R^k=R^{-k}\gamma_{(0)}$ for every $k\in\Z$, and
\begin{equation}\label{eq:positivity-spin-conjugation}
\gamma_{(k)}=U^{-k}\gamma_{(0)}=R^{-2k}\gamma_{(0)}=R^{-k}\gamma_{(0)}R^{k},\qquad k\in\Z .
\end{equation}

\emph{A rotating frame.}
Let $\widehat{\mathcal{H}}$ be the space of the sequences $\hat\psi\colon\Z\to\C^4$ such that $\hat\psi(k+6)=-\hat\psi(k)$ for every $k\in\Z$.
Such a sequence is determined by $\hat\psi(0),\dots,\hat\psi(5)$, and these six vectors can be prescribed freely;
so $\dim\widehat{\mathcal{H}}=24$.
To a map $\psi\colon H\to\C^4$ attach the sequence
\[
\hat\psi(k):=R^k\,\psi(z_k),\qquad k\in\Z .
\]
It belongs to $\widehat{\mathcal{H}}$ because $R^6=-1$.
The map $\psi\mapsto\hat\psi$ is linear and bijective onto $\widehat{\mathcal{H}}$;
its inverse sends $\hat\psi$ to the map $z_k\mapsto R^{-k}\,\hat\psi(k)$, which is well defined since $R^{-6}=-1$.
Let $\widehat D_H$ be the operator on $\widehat{\mathcal{H}}$ that sends the sequence attached to $\psi$ to the sequence attached to $D_H\psi$.
It is similar to $D_H$, and so
\[
\det\widehat D_H=\det D_H .
\]
Let
\[
B_+:=(\rho-\gamma_{(0)})R^{-1},\qquad
B_-:=R(\rho+\gamma_{(0)}).
\]
By \eqref{eq:positivity-spin-conjugation}, $R^k\gamma_{(k)}R^{-k-1}=\gamma_{(0)}R^{-1}$ and $R^k\gamma_{(k-1)}R^{-k+1}=R\gamma_{(0)}$, so that
\[
R^k(\rho-\gamma_{(k)})R^{-k-1}=B_+,\qquad
R^k(\rho+\gamma_{(k-1)})R^{-k+1}=B_-
\]
for every $k\in\Z$.
Multiplying \eqref{eq:positivity-hexagon-operator} by $R^k$ and inserting $\psi(z_{k\pm1})=R^{-k\mp1}\,\hat\psi(k\pm1)$, one finds
\[
(\widehat D_H\hat\psi)(k)=M\hat\psi(k)-\frac{\tilde w_k}{2}\,B_+\,\hat\psi(k+1)-\frac{\tilde w_{k-1}}{2}\,B_-\,\hat\psi(k-1),\qquad k\in\Z .
\]

\emph{Fourier decomposition.}
The shift $\hat\psi\mapsto\hat\psi(\,\cdot+2)$ maps $\widehat{\mathcal{H}}$ onto itself, and its third power is $-1$ on $\widehat{\mathcal{H}}$.
It commutes with $\widehat D_H$, because $B_+$ and $B_-$ do not depend on $k$ and $\tilde w_{k+2}=\tilde w_k$.
The equation $\exp(3it)=-1$ has the three solutions $t=\pi/3,\ \pi,\ -\pi/3$ in $(-\pi,\pi]$.
For each of them let $\widehat{\mathcal{H}}_t$ be the eigenspace of the shift for the eigenvalue $\exp(it)$,
the space of the $\hat\psi\in\widehat{\mathcal{H}}$ with $\hat\psi(k+2)=\exp(it)\cdot\hat\psi(k)$ for all $k\in\Z$.
Its elements are the sequences
\[
\hat\psi(2j)=\exp(itj)\cdot\hat\psi(0),\qquad
\hat\psi(2j+1)=\exp(itj)\cdot\hat\psi(1),\qquad j\in\Z,
\]
with arbitrary $\hat\psi(0),\hat\psi(1)\in\C^4$:
every sequence of this form satisfies $\hat\psi(k+6)=\exp(3it)\cdot\hat\psi(k)=-\hat\psi(k)$.
So $\dim\widehat{\mathcal{H}}_t=8$.
Eigenvectors for distinct eigenvalues are linearly independent, and $3\cdot8=24=\dim\widehat{\mathcal{H}}$;
therefore
\[
\widehat{\mathcal{H}}=\widehat{\mathcal{H}}_{\pi/3}\oplus\widehat{\mathcal{H}}_{\pi}\oplus\widehat{\mathcal{H}}_{-\pi/3},
\]
each summand is invariant under $\widehat D_H$, and $\det\widehat D_H$ is the product of the determinants of the three restrictions.

Fix $t$ and $\hat\psi\in\widehat{\mathcal{H}}_t$.
Since $\tilde w_{2j}=w_1$ and $\tilde w_{2j-1}=\tilde w_{2j+1}=w_2$, the formula for $\widehat D_H$ gives
\[
\begin{aligned}
(\widehat D_H\hat\psi)(2j)&=\exp(itj)\cdot\left(M\hat\psi(0)-X\,\hat\psi(1)\right),\\
(\widehat D_H\hat\psi)(2j+1)&=\exp(itj)\cdot\left(M\hat\psi(1)-Y\,\hat\psi(0)\right),
\end{aligned}
\qquad j\in\Z,
\]
with the $4\times4$ matrices
\[
X:=\frac{w_1}{2}\,B_++\frac{w_2}{2}\,\exp(-it)\cdot B_-,\qquad
Y:=\frac{w_2}{2}\,\exp(it)\cdot B_++\frac{w_1}{2}\,B_- ,
\]
which depend on $t$.
In the coordinates $(\hat\psi(0),\hat\psi(1))\in\C^4\oplus\C^4$ of $\widehat{\mathcal{H}}_t$ the operator $\widehat D_H$ thus acts by the block matrix
\[
\begin{pmatrix} M & -X\\ -Y & M\end{pmatrix}.
\]
Its lower right block is a multiple of the identity and commutes with the lower left block;
by the identity for such block matrices in \eqref{eq:positivity-identities}, its determinant is $\det_{\C^4}\left(M^2-XY\right)$.
Altogether,
\[
\det D_H=\prod_{t\in\{\pi/3,\,\pi,\,-\pi/3\}}\det_{\C^4}\left(M^2-XY\right).
\]

\emph{The product $XY$.}
As $\gamma_{(0)}^2=1$,
\[
B_+B_-=(\rho-\gamma_{(0)})R^{-1}R(\rho+\gamma_{(0)})=\rho^2-1,\qquad
B_-B_+=R(\rho+\gamma_{(0)})(\rho-\gamma_{(0)})R^{-1}=\rho^2-1 .
\]
For the squares one uses $R^{-1}\gamma_{(0)}=\gamma_{(0)}R$ and $R^2=U$:
\[
\begin{aligned}
B_+^2&=(\rho-\gamma_{(0)})\left(\rho R^{-1}-\gamma_{(0)}R\right)R^{-1}=(\rho-\gamma_{(0)})\left(\rho U^{-1}-\gamma_{(0)}\right)\\
&=1+\rho^2U^{-1}-\rho\gamma_{(0)}\left(1+U^{-1}\right),\\
B_-^2&=R\left(\rho R+R^{-1}\gamma_{(0)}\right)(\rho+\gamma_{(0)})=\left(\rho U+\gamma_{(0)}\right)(\rho+\gamma_{(0)})\\
&=1+\rho^2U+\rho\,(1+U)\gamma_{(0)} .
\end{aligned}
\]
Moreover $(1+U)\gamma_{(0)}=\gamma_{(0)}(1+U^{-1})=\sqrt{3}\,\gamma_{(0)}R^{-1}$.
Let
\[
J:=\gamma_{(0)}R^{-1}.
\]
This matrix is a Hermitian involution, and conjugation by it inverts $U$:
\[
J^*=R\gamma_{(0)}=J,\qquad
J^2=\gamma_{(0)}R^{-1}\gamma_{(0)}R^{-1}=\gamma_{(0)}^2=1,\qquad
JU=\gamma_{(0)}R=R^{-1}\gamma_{(0)}=U^{-1}J .
\]
Hence
\[
B_+^2=1+\rho^2U^{-1}-\sqrt{3}\,\rho J,\qquad
B_-^2=1+\rho^2U+\sqrt{3}\,\rho J .
\]
In the product $XY$ the terms $B_+B_-$ and $B_-B_+$ carry the coefficients $w_1^2/4$ and $w_2^2/4$,
and the squares $B_+^2$ and $B_-^2$ the coefficients $\frac{w_1w_2}{4}\,\exp(it)$ and $\frac{w_1w_2}{4}\,\exp(-it)$;
it follows that
\[
\begin{aligned}
XY&=\frac{(w_1^2+w_2^2)(\rho^2-1)}{4}+\frac{w_1w_2}{4}\left(\exp(it)\cdot B_+^2+\exp(-it)\cdot B_-^2\right)\\
&=-s_H+p_H\left(\cos t+\frac{\rho^2}{2}\left(\exp(it)\cdot U^{-1}+\exp(-it)\cdot U\right)\right)\\
&\qquad-i\sqrt{3}\,p_H\rho\cdot\sin t\cdot J .
\end{aligned}
\]

\emph{The determinant.}
The unitary matrix $U$ is diagonalisable, and by \eqref{eq:positivity-spin-quadratic} its eigenvalues lie among the two roots $\exp\left(\pm\frac{i\pi}{3}\right)$ of the polynomial occurring there.
Let $\C^4_\pm$ be the eigenspace of $U$ for $\exp\left(\pm\frac{i\pi}{3}\right)$;
then $\C^4=\C^4_+\oplus\C^4_-$.
From $JU=U^{-1}J$ one gets $UJ=JU^{-1}$, which equals $\exp\left(\mp\frac{i\pi}{3}\right)\cdot J$ on $\C^4_\pm$;
thus $J$ maps $\C^4_\pm$ into $\C^4_\mp$.
Being invertible, $J$ maps each of the two spaces injectively into the other, and both have dimension two.
Let $J_{+-}\colon\C^4_-\to\C^4_+$ and $J_{-+}\colon\C^4_+\to\C^4_-$ be the two restrictions of $J$;
the relation $J^2=1$ gives $J_{+-}J_{-+}=1$ on $\C^4_+$.
On $\C^4_\pm$,
\[
\begin{aligned}
\frac12\left(\exp(it)\cdot U^{-1}+\exp(-it)\cdot U\right)
&=\frac12\left(\exp\left(i\left(t\mp\frac{\pi}{3}\right)\right)+\exp\left(-i\left(t\mp\frac{\pi}{3}\right)\right)\right)\\
&=\cos\left(t\mp\frac{\pi}{3}\right).
\end{aligned}
\]
With respect to $\C^4=\C^4_+\oplus\C^4_-$ the formula for $XY$ therefore yields
\[
M^2-XY=\begin{pmatrix} r_+ & i\kappa J_{+-}\\ i\kappa J_{-+} & r_-\end{pmatrix},
\]
where the diagonal blocks are the multiples of the identity by the numbers
\[
r_\pm:=r-p_H\rho^2\cdot\cos\left(t\mp\frac{\pi}{3}\right),\qquad
r:=M^2+s_H-p_H\cdot\cos t,
\]
and where
\[
\kappa:=\sqrt{3}\,p_H\rho\cdot\sin t .
\]
In bases of the two-dimensional spaces $\C^4_+$ and $\C^4_-$ the four blocks are $2\times2$ matrices, the lower right one is scalar,
and \eqref{eq:positivity-identities} gives
\[
\det_{\C^4}\left(M^2-XY\right)=\det_{\C^4_+}\left(r_+r_-+\kappa^2J_{+-}J_{-+}\right)=\left(r_+r_-+3p_H^2\rho^2\cdot\sin^2t\right)^2 .
\]

It remains to insert the three values of $t$.
For $t=\pi$ one has $\cos\left(\pi\mp\frac{\pi}{3}\right)=-\frac12$ and $\sin t=0$;
so $r=M^2+s_H+p_H$, $r_+=r_-=r+\frac{p_H\rho^2}{2}$, and the determinant is
\[
\left(r+\frac{p_H\rho^2}{2}\right)^4=\left(M^2+s_H+p_H\left(1+\frac{\rho^2}{2}\right)\right)^4 .
\]
For $t=\pm\frac{\pi}{3}$ one has $\cos t=\frac12$, $\sin^2t=\frac34$ and $\left\{\cos\left(t-\frac{\pi}{3}\right),\cos\left(t+\frac{\pi}{3}\right)\right\}=\left\{1,-\frac12\right\}$;
so $r=M^2+s_H-\frac{p_H}{2}$, and
\[
\begin{aligned}
r_+r_-+\frac{9}{4}\,p_H^2\rho^2
&=\left(r-p_H\rho^2\right)\left(r+\frac{p_H\rho^2}{2}\right)+\frac{9}{4}\,p_H^2\rho^2\\
&=r^2-\frac{p_H\rho^2}{2}\,r-\frac{p_H^2\rho^4}{2}+\frac{9}{4}\,p_H^2\rho^2\\
&=\left(r-\frac{p_H\rho^2}{4}\right)^2+\frac{9}{4}\,p_H^2\rho^2\left(1-\frac{\rho^2}{4}\right);
\end{aligned}
\]
for each of the two values of $t$ the determinant is the square of this number.
Since $r-\frac{p_H\rho^2}{4}=M^2+s_H-\frac{p_H}{2}-\frac{p_H\rho^2}{4}$,
the product of the three determinants is the right-hand side of \eqref{eq:positivity-hexagon}.

\emph{Positivity.}
Let $-1\le\rho\le1$.
Then $s_H\ge0$ and $p_H>0$, so the first bracket in \eqref{eq:positivity-hexagon} satisfies
\[
M^2+s_H+p_H\left(1+\frac{\rho^2}{2}\right)\ge p_H>0 .
\]
If $\rho\neq0$, the second bracket is at least $\frac{9}{4}\,p_H^2\rho^2\cdot\frac34>0$, because $1-\frac{\rho^2}{4}\ge\frac34$.
If $\rho=0$, then $s_H=\frac{w_1^2+w_2^2}{4}$, and the second bracket equals
\[
\left(M^2+\frac{w_1^2+w_2^2-w_1w_2}{4}\right)^2\ge\left(\frac{w_1w_2}{4}\right)^2>0,
\]
since $w_1^2+w_2^2\ge2w_1w_2$.
In every case $\det D_H>0$.

\emph{The values $\rho=\pm1$.}
Here $s_H=0$, and the two brackets in \eqref{eq:positivity-hexagon} become $M^2+\frac32\,p_H$ and
\[
\left(M^2-\frac34\,p_H\right)^2+\frac34\left(\frac32\,p_H\right)^2=M^4-\frac32\,p_H\,M^2+\left(\frac32\,p_H\right)^2 .
\]
By the factorisation of a sum of two cubes their product is $M^6+\left(\frac32\,p_H\right)^3$,
and $\frac32\,p_H=\frac{3w_1w_2}{4}$;
hence $\det D_H=\left(M^6+\left(\frac{3w_1w_2}{4}\right)^3\right)^4$.
\end{proof}

\subsection{Positivity on every torus}\label{sec:every-torus}

\begin{proposition}[Free determinants on antiperiodic coordinate tori]\label{prop:free-determinants}
Let $-1\le\rho\le1$, $M\in\R$ and $w_0,\dots,w_3>0$.
Let $n_1,n_2,n_3\ge1$ be integers, put $\Gamma:=n_1\Z e_1+n_2\Z e_2+n_3\Z e_3$,
and let $\chi\colon\Gamma\to\{+1,-1\}$ be the character with $\chi(n_\mu e_\mu)=-1$ for $\mu=1,2,3$.
Let $C(n_1,n_2,n_3):=(\Gamma,\chi)$, a general torus in the sense of Definition~\ref{def:general-tori};
for $n_1,n_2,n_3\ge4$ it is the coordinate torus of Definition~\ref{def:tori}.
Then
\[
\det D^{C(n_1,n_2,n_3)}_0\ \ge\ \left(\det D_H\right)^{8n_1n_2n_3}\ >\ 0,
\]
where $D_H$ is the hexagon operator of Definition~\ref{def:hexagon}.
\end{proposition}

\begin{proof}
Restrictions of $\chi$ to subgroups of $\Gamma$ are again denoted by $\chi$.
Consider the four triples
\[
\begin{aligned}
\mathfrak{P}_0&:=\left(\mathbb{X},\ \Gamma,\ \chi\right),\\
\mathfrak{P}_1&:=\left(\left\{x\in\mathbb{X}:\ 0<x_1<\tfrac12\right\},\ n_2\Z e_2+n_3\Z e_3,\ \chi\right),\\
\mathfrak{P}_2&:=\left(\left\{x\in\mathbb{X}:\ 0<x_1<\tfrac12,\ 0<x_2<\tfrac12\right\},\ n_3\Z e_3,\ \chi\right),\\
\mathfrak{P}_3&:=\left(H,\ \{0\},\ 1\right).
\end{aligned}
\]
They are pieces in the sense of Definition~\ref{def:pieces}.
For $\mathfrak{P}_0$ this is Definition~\ref{def:pieces}\eqref{it:piece-torus}, which also gives $D_{\mathfrak{P}_0}=D^{C(n_1,n_2,n_3)}_0$;
for $\mathfrak{P}_3$ it is part of Definition~\ref{def:hexagon}, and $D_{\mathfrak{P}_3}=D_H$.
Translations by vectors of $\Z^3$ belong to $W$ (Definition~\ref{def:mirrors}\eqref{it:mirror-product}) and therefore map $\mathbb{X}=Wv_0$ onto itself;
the translations by vectors of $n_2\Z e_2+n_3\Z e_3$ leave the coordinate $x_1$ unchanged,
and those by vectors of $n_3\Z e_3$ leave $x_1$ and $x_2$ unchanged.
So the sets of sites of $\mathfrak{P}_1$ and $\mathfrak{P}_2$ are invariant under the respective groups.
Every orbit of these groups meets a bounded subset of $\R^3$,
and $\mathbb{X}=\Z^3+\Ocub v_0$ has only finitely many points in a bounded set (Definition~\ref{def:fermion-sites});
hence the quotients are finite.
All characters involved take the values $\pm1$.

Lemma~\ref{lem:half-slab} is now applied in the directions $\mu=1,2,3$ in turn.
For $c\in\tfrac12\Z$ the reflection $\theta_{x_\mu,c}$ belongs to $W$, so that it maps $\mathbb{X}$ onto itself,
and it changes only the coordinate $x_\mu$ of a point (Definition~\ref{def:mirrors}).
\begin{itemize}
\item
For $\mu=1$ let $n:=n_1$ and $\Gamma_0':=n_2\Z e_2+n_3\Z e_3$, a subgroup of $\{g\in\Z^3:\ g_1=0\}$.
The group of $\mathfrak{P}_0$ is $\Gamma=\Gamma_0'+\Z n_1e_1$, one has $\chi(n_1e_1)=-1$,
and $\theta_{x_1,c}\mathbb{X}=\mathbb{X}$ for every $c\in\tfrac12\Z$.
The sites $x\in\mathbb{X}$ with $0<x_1<\tfrac12$ form the set of sites of $\mathfrak{P}_1$, whose group is $\Gamma_0'$.
\item
For $\mu=2$ let $n:=n_2$ and $\Gamma_0':=n_3\Z e_3$, a subgroup of $\{g\in\Z^3:\ g_2=0\}$.
The group of $\mathfrak{P}_1$ is $\Gamma_0'+\Z n_2e_2$, one has $\chi(n_2e_2)=-1$,
and $\theta_{x_2,c}$ maps the set of sites of $\mathfrak{P}_1$ onto itself, since it leaves $x_1$ unchanged.
The sites of $\mathfrak{P}_1$ with $0<x_2<\tfrac12$ form the set of sites of $\mathfrak{P}_2$, whose group is $\Gamma_0'$.
\item
For $\mu=3$ let $n:=n_3$ and $\Gamma_0':=\{0\}$.
The group of $\mathfrak{P}_2$ is $\Gamma_0'+\Z n_3e_3$, one has $\chi(n_3e_3)=-1$,
and $\theta_{x_3,c}$ maps the set of sites of $\mathfrak{P}_2$ onto itself, since it leaves $x_1$ and $x_2$ unchanged.
The sites of $\mathfrak{P}_2$ with $0<x_3<\tfrac12$ form the set $\mathbb{X}\cap\left(0,\tfrac12\right)^3=H$,
and the restriction of $\chi$ to $\{0\}$ is the trivial character.
\end{itemize}
The three applications of Lemma~\ref{lem:half-slab} yield
\[
\det D_{\mathfrak{P}_0}\ge\left|\det D_{\mathfrak{P}_1}\right|^{2n_1},\qquad
\det D_{\mathfrak{P}_1}\ge\left|\det D_{\mathfrak{P}_2}\right|^{2n_2},\qquad
\det D_{\mathfrak{P}_2}\ge\left|\det D_H\right|^{2n_3};
\]
in particular $\det D_{\mathfrak{P}_1}$ and $\det D_{\mathfrak{P}_2}$ are real and non-negative.
Since $-1\le\rho\le1$ and $M$ is real, Lemma~\ref{lem:hexagon-determinant} gives $\det D_H>0$.
Powers with positive integer exponents are non-decreasing on $[0,\infty)$, and the three inequalities combine to
\[
\det D^{C(n_1,n_2,n_3)}_0=\det D_{\mathfrak{P}_0}
\ge\left(\det D_{\mathfrak{P}_1}\right)^{2n_1}
\ge\left(\det D_{\mathfrak{P}_2}\right)^{4n_1n_2}
\ge\left(\det D_H\right)^{8n_1n_2n_3}
>0 .
\]
\end{proof}

\begin{lemma}[Characters and coverings]\label{lem:characters-coverings}
Let $\rho,M\in\C$, $w_0,\dots,w_3>0$ and $q>0$, and let $\Gamma$ be a subgroup of $\Z^3$ of rank three.
For a unitary character $\chi$ of $\Gamma$ and a $\Gamma$-periodic one-cochain $A$,
$D^{(\Gamma,\chi)}_A$ denotes the operator on the spinors of the general torus $(\Gamma,\chi)$ (Definition~\ref{def:general-tori})
that is given by the formula for $D_A$ in Definition~\ref{def:chamber-operator}.
\begin{enumerate}
\item\label{it:character-shift}
Let $\chi_1$ and $\chi_2$ be unitary characters of $\Gamma$.
There is a vector $\xi\in\R^3$ such that
\[
\chi_2(g)=\exp\left(i\,\xi\cdot g\right)\cdot\chi_1(g)\qquad\text{for all }g\in\Gamma .
\]
For such a vector $\xi$ let $a_\xi$ be the constant one-cochain with the components $(a_\xi)_\mu=\xi_\mu/q$, $\mu=1,2,3$.
Then
\[
\det D^{(\Gamma,\chi_1)}_A=\det D^{(\Gamma,\chi_2)}_{A-a_\xi}
\]
for every $\Gamma$-periodic one-cochain $A$.
\item\label{it:covering-product}
Let $\Gamma'$ be a subgroup of finite index of $\Gamma$, let $\chi'$ be a unitary character of $\Gamma'$,
and let $A$ be a $\Gamma$-periodic one-cochain.
Then $\Gamma$ has exactly $[\Gamma:\Gamma']$ unitary characters $\tilde\chi$ that extend $\chi'$, and
\[
\det D^{(\Gamma',\chi')}_A=\prod_{\tilde\chi}\det D^{(\Gamma,\tilde\chi)}_A ,
\]
the product being taken over these $[\Gamma:\Gamma']$ characters.
\end{enumerate}
\end{lemma}

\begin{proof}
\emph{Item~\eqref{it:character-shift}.}
Being a subgroup of rank three of $\Z^3$, the group $\Gamma$ is free abelian of rank three; let $f_1,f_2,f_3$ be a basis of $\Gamma$.
As $\Gamma$ has finite index in $\Z^3$, it spans $\R^3$, so that $f_1,f_2,f_3$ are linearly independent over $\R$.
The numbers $\chi_2(f_j)/\chi_1(f_j)$ have modulus one;
choose $t_1,t_2,t_3\in\R$ with $\chi_2(f_j)=\exp(it_j)\cdot\chi_1(f_j)$,
and let $\xi\in\R^3$ be the solution of the linear system $\xi\cdot f_j=t_j$, $j=1,2,3$.
The maps $g\mapsto\chi_2(g)$ and $g\mapsto\exp(i\,\xi\cdot g)\cdot\chi_1(g)$ are homomorphisms of $\Gamma$
into the multiplicative group of the complex numbers of modulus one,
and they coincide on a basis of $\Gamma$; so they coincide on $\Gamma$.

From now on $\xi$ is any vector with this property.
For maps $\psi\colon\mathbb{X}\to\C^4$ put
\[
(U_\xi\psi)(x):=\exp\left(i\,\xi\cdot\pi(x)\right)\cdot\psi(x),\qquad x\in\mathbb{X}.
\]
If $\psi$ is a spinor of $(\Gamma,\chi_1)$, then, as $\pi(x+g)=\pi(x)+g$ for $g\in\Z^3$ (Definition~\ref{def:fermion-sites}),
\[
(U_\xi\psi)(x+g)=\exp\left(i\,\xi\cdot g\right)\cdot\chi_1(g)\cdot(U_\xi\psi)(x)=\chi_2(g)\cdot(U_\xi\psi)(x),\qquad g\in\Gamma .
\]
Thus $U_\xi$ maps the spinors of $(\Gamma,\chi_1)$ into those of $(\Gamma,\chi_2)$;
by the same computation $U_{-\xi}$ maps the spinors of $(\Gamma,\chi_2)$ into those of $(\Gamma,\chi_1)$,
and $U_{-\xi}$ is the inverse of $U_\xi$.
For a spinor $\psi$ of $(\Gamma,\chi_2)$ and $x\in\mathbb{X}$, Definition~\ref{def:chamber-operator} gives
\[
\left(U_\xi D_AU_{-\xi}\psi\right)(x)
=M\psi(x)-\sum_{i=0}^{3}w(x,x^{(i)})\cdot U_A(x,x^{(i)})\cdot\exp\left(-i\,\xi\cdot\left(\pi(x^{(i)})-\pi(x)\right)\right)\cdot\psi(x^{(i)}).
\]
By Proposition~\ref{prop:geometry}\eqref{it:geometry-projection} the vector $\pi(x^{(i)})-\pi(x)$ is $0$ or one of $\pm e_1,\pm e_2,\pm e_3$.
For $y,y'\in\Z^3$ with $y'-y\in\{0,\pm e_1,\pm e_2,\pm e_3\}$ the constant one-cochain satisfies $a_\xi(y,y')=\xi\cdot(y'-y)/q$,
since a one-cochain is read as an odd function of oriented pairs of neighbours that vanishes on the pairs $(y,y)$ (Definition~\ref{def:cochains}).
Therefore
\[
U_A(x,x^{(i)})\cdot\exp\left(-i\,\xi\cdot\left(\pi(x^{(i)})-\pi(x)\right)\right)
=\exp\left(iq\,(A-a_\xi)\left(\pi(x),\pi(x^{(i)})\right)\right)
=U_{A-a_\xi}(x,x^{(i)}),
\]
and $A-a_\xi$ is $\Gamma$-periodic together with $A$.
This proves $U_\xi D^{(\Gamma,\chi_1)}_AU_{-\xi}=D^{(\Gamma,\chi_2)}_{A-a_\xi}$,
and operators that are conjugate under a linear bijection have the same determinant.

\emph{Item~\eqref{it:covering-product}.}
Let $\mathcal{H}'$ be the space of the spinors of $(\Gamma',\chi')$, with the scalar product
\[
\langle\psi',\psi\rangle:=\sum_{x}\psi'(x)^{*}\cdot\psi(x),
\]
in which $x$ ranges over a set of representatives of $\mathbb{X}/\Gamma'$;
the summands do not depend on the representatives, since $|\chi'|=1$.
A spinor of a general torus $\mathbb{T}$ is determined by its values at a set of representatives of $\mathbb{X}_{\mathbb{T}}$,
and these values are arbitrary, as the period group acts freely on $\mathbb{X}$;
so the spinors of $\mathbb{T}$ form a space of dimension $4\cdot\#|\mathbb{X}_{\mathbb{T}}|=192\cdot|\mathbb{T}|$ (Definition~\ref{def:general-tori}).
Hence $\dim\mathcal{H}'=192\cdot[\Z^3:\Gamma']$,
and for every unitary character $\tilde\chi$ of $\Gamma$ the spinors of $(\Gamma,\tilde\chi)$ form a space of dimension $192\cdot[\Z^3:\Gamma]$.

For $g\in\Gamma$ and $\psi\in\mathcal{H}'$ put $(V_g\psi)(x):=\psi(x+g)$.
As $(V_g\psi)(x+g')=\chi'(g')\cdot(V_g\psi)(x)$ for $g'\in\Gamma'$, the operator $V_g$ maps $\mathcal{H}'$ into itself.
It is unitary, because translation by $g$ maps a set of representatives of $\mathbb{X}/\Gamma'$ onto another such set.
Moreover $V_gV_{g'}=V_{g+g'}$ for $g,g'\in\Gamma$, and $V_g=\chi'(g)\cdot1$ for $g\in\Gamma'$.
Each $V_g$ commutes with $D_A$.
Indeed, $(x+g)^{(i)}=x^{(i)}+g$ and $\pi(x+g)=\pi(x)+g$ (Definition~\ref{def:fermion-sites}),
so that the bond $\{x+g,x^{(i)}+g\}$ has the type and the direction of the bond $\{x,x^{(i)}\}$ and $w(x+g,x^{(i)}+g)=w(x,x^{(i)})$;
moreover $U_A(x+g,x^{(i)}+g)=U_A(x,x^{(i)})$, because $A$ is $\Gamma$-periodic.
Hence $(D_AV_g\psi)(x)=(D_A\psi)(x+g)=(V_gD_A\psi)(x)$.

The operators $V_g$, $g\in\Gamma$, are commuting unitary operators on the finite-dimensional Hilbert space $\mathcal{H}'$.
By the spectral theorem, $\mathcal{H}'$ is the orthogonal direct sum of the joint eigenspaces
\[
\mathcal{H}'_{\tilde\chi}:=\left\{\psi\in\mathcal{H}':\ V_g\psi=\tilde\chi(g)\cdot\psi\ \text{ for all }g\in\Gamma\right\},
\]
where $\tilde\chi$ ranges over the maps from $\Gamma$ to the complex numbers of modulus one with $\mathcal{H}'_{\tilde\chi}\neq\{0\}$.
Applied to a non-zero vector of $\mathcal{H}'_{\tilde\chi}$,
the relations $V_gV_{g'}=V_{g+g'}$ and $V_g=\chi'(g)\cdot1$ for $g\in\Gamma'$ show that every such $\tilde\chi$ is a character of $\Gamma$ that extends $\chi'$.
Conversely, let $\tilde\chi$ be a unitary character of $\Gamma$ that extends $\chi'$.
Every spinor of $(\Gamma,\tilde\chi)$ lies in $\mathcal{H}'$, and $\mathcal{H}'_{\tilde\chi}$ is exactly the space of these spinors;
its dimension $192\cdot[\Z^3:\Gamma]$ is positive.
As $[\Z^3:\Gamma']=[\Z^3:\Gamma]\cdot[\Gamma:\Gamma']$,
a comparison of the dimensions shows that $\Gamma$ has exactly $[\Gamma:\Gamma']$ unitary characters that extend $\chi'$.
Since $D_A$ commutes with every $V_g$, it maps each $\mathcal{H}'_{\tilde\chi}$ into itself, and there it acts as $D^{(\Gamma,\tilde\chi)}_A$.
The determinant of $D^{(\Gamma',\chi')}_A$, the operator $D_A$ on $\mathcal{H}'$, is therefore the product of the determinants of these restrictions.
\end{proof}

\begin{theorem}[The closed interval]\label{thm:closed-interval}
Let $-1\le\rho\le1$, $M\in\R$, $w_0,\dots,w_3>0$ and $q>0$.
Let $\Gamma$ be a subgroup of $\Z^3$ of rank three, $\chi$ a unitary character of $\Gamma$, and $\mathbb{T}:=(\Gamma,\chi)$.
\begin{enumerate}
\item\label{it:closed-interval-cochain}
There is a constant real one-cochain $a$ with $\det D^{\mathbb{T}}_a>0$.
Consequently the function $A\mapsto\det D^{\mathbb{T}}_A$ on the real $\Gamma$-periodic one-cochains
is a trigonometric polynomial in the variables $qA(b)$, $b\in\mathbb{T}^{(1)}$,
which is non-negative and not identically zero.
\item\label{it:closed-interval-flux}
One has $0<Z_{\mathbb{T}}[1]<\infty$, and every flux term of $Z_{\mathbb{T}}[1]$ (Definition~\ref{def:general-functional}) is positive:
\[
\int_{[0,2\pi/q)^{\mathbb{T}^{(1)}}}\mathrm{d}A\,\exp\left(-\frac12\left\|dA+\frac{2\pi}{q}\,z\right\|^2\right)\cdot\det D^{\mathbb{T}}_A>0
\qquad\text{for every closed }z\in C^2(\Gamma;\Z).
\]
\item\label{it:closed-interval-null}
The set $\mathcal{N}_{\mathbb{T}}:=\left\{A\in\R^{\mathbb{T}^{(1)}}:\ \det D^{\mathbb{T}}_A=0\right\}$ is closed and of Lebesgue measure zero.
For every Borel function $\Phi\ge0$ on $\Conf(\Gamma)$ one has $Z_{\mathbb{T}}[\Phi]=0$
if and only if $\Phi=0$ almost everywhere with respect to the measure of $\Conf(\Gamma)$ (Definition~\ref{def:configurations}).
Hence the two measures on $\Conf(\Gamma)$ whose densities with respect to that measure are
$\exp\left(-\frac12\|F\|^2\right)\cdot\det D^{\mathbb{T}}_A$ and $\exp\left(-\frac12\|F\|^2\right)$ have the same null sets,
and the law $\mu_{\mathbb{T}}$ of the periodic extension $\widetilde F$ of $F$ under $Z_{\mathbb{T}}[\,\cdot\,]/Z_{\mathbb{T}}[1]$
and the law of $\widetilde F$ under the normalised measure with the density $\exp\left(-\frac12\|F\|^2\right)$
are mutually absolutely continuous.
\end{enumerate}
\end{theorem}

\begin{proof}
Put $n:=[\Z^3:\Gamma]=|\mathbb{T}|$.
By Lagrange's theorem the order of every element of the finite group $\Z^3/\Gamma$ divides $n$;
hence $ny\in\Gamma$ for every $y\in\Z^3$, and $\Gamma':=n\Z e_1+n\Z e_2+n\Z e_3$ is a subgroup of $\Gamma$ of index $n^2$.
Let $\chi'$ be the character of $\Gamma'$ with $\chi'(ne_\mu)=-1$ for $\mu=1,2,3$;
then $(\Gamma',\chi')$ is the torus $C(n,n,n)$ of Proposition~\ref{prop:free-determinants}.
Lemma~\ref{lem:characters-coverings}\eqref{it:covering-product}, applied with $A=0$, and Proposition~\ref{prop:free-determinants} give
\[
\prod_{\tilde\chi}\det D^{(\Gamma,\tilde\chi)}_0=\det D^{C(n,n,n)}_0\ \ge\ \left(\det D_H\right)^{8n^3}\ >\ 0,
\]
the product being taken over the unitary characters $\tilde\chi$ of $\Gamma$ that extend $\chi'$.
Hence $\det D^{(\Gamma,\tilde\chi)}_0\neq0$ for each of these characters; fix one of them.
Lemma~\ref{lem:characters-coverings}\eqref{it:character-shift}, applied with $\chi_1:=\chi$ and $\chi_2:=\tilde\chi$,
provides $\xi\in\R^3$ with $\tilde\chi(g)=\exp(i\,\xi\cdot g)\cdot\chi(g)$ for $g\in\Gamma$,
and for the $\Gamma$-periodic one-cochain $A:=a_\xi$ it gives
\[
\det D^{\mathbb{T}}_{a_\xi}=\det D^{(\Gamma,\tilde\chi)}_0\neq0 .
\]
The one-cochain $a:=a_\xi$ is constant and real.
As $\rho$, $M$ and $a$ are real and $\chi$ is unitary, Proposition~\ref{prop:accretivity}\eqref{it:determinant-sign} gives $\det D^{\mathbb{T}}_a\ge0$;
so $\det D^{\mathbb{T}}_a>0$.
For real $\Gamma$-periodic $A$ one has
\[
\det D^{\mathbb{T}}_A=P_{\mathbb{T}}\left(\left(\exp(iqA(b))\right)_{b\in\mathbb{T}^{(1)}}\right)
\]
with the determinant polynomial $P_{\mathbb{T}}$ of Definition~\ref{def:determinant-polynomial},
a trigonometric polynomial in the variables $qA(b)$.
Its values are non-negative by Proposition~\ref{prop:accretivity}\eqref{it:determinant-sign}, and its value at $A=a$ is positive.
This proves \eqref{it:closed-interval-cochain}.

The function $A\mapsto\det D^{\mathbb{T}}_A$ on $\R^{\mathbb{T}^{(1)}}$ is therefore continuous and non-negative,
it is $(2\pi/q)$-periodic in each variable $A(b)$, as it depends on $A$ only through the phases $\exp(iqA(b))$,
and it is not identically zero.
So Lemma~\ref{lem:flux-terms} applies to $\mathbb{T}$.
It shows that every flux term of $Z_{\mathbb{T}}[1]$ is positive, and hence that $Z_{\mathbb{T}}[1]>0$,
the flux terms being non-negative and their sum being $Z_{\mathbb{T}}[1]$;
and it gives all assertions of \eqref{it:closed-interval-null}.
Finally, $Z_{\mathbb{T}}[1]<\infty$ by the upper bound \eqref{eq:positivity-ceiling}, which holds for every real $\rho$.
\end{proof}

\begin{theorem}[Positivity of the partition function]\label{thm:partition-positivity}
Let $\rho\in\R$, $M\in\R$, $w_0,\dots,w_3>0$, $q>0$ and $v_0\in C_0$,
and let $\mathbb{T}=(\Gamma,\chi)$ be a torus in the sense of Definition~\ref{def:tori},
with an arbitrary period lattice $\Gamma$ and an arbitrary character $\chi\colon\Gamma\to\{+1,-1\}$, the trivial character included.
\begin{enumerate}
\item\label{it:partition-finite}
One has $0<Z_{\mathbb{T}}[1]<\infty$.
\item\label{it:partition-flux}
Every flux term of $Z_{\mathbb{T}}[1]$ (Definition~\ref{def:general-functional}) is positive, in particular the flux term of $z=0$.
\item\label{it:partition-perturbed}
For every real bounded Borel function $V$ on $\Conf(\Gamma)$ and every $\beta\in C^2(\Gamma)$,
\[
Z_{\mathbb{T}}\left[\exp V\right]>0\qquad\text{and}\qquad Z_{\mathbb{T}}\left[\exp\left(-\langle\beta,F\rangle-\frac12\|\beta\|^2\right)\right]>0,
\]
the pairing and the norm being sums over the plaquettes of $\Z^3/\Gamma$.
\end{enumerate}
In particular, the normalised functional $\mathbb{E}_{\mathbb{T}}=Z_{\mathbb{T}}[\,\cdot\,]/Z_{\mathbb{T}}[1]$
and the law $\mu_{\mathbb{T}}$ of Definition~\ref{def:functional} are defined for all real $\rho$ and $M$.
\end{theorem}

\begin{proof}
Every torus of Definition~\ref{def:tori} is a general torus in the sense of Definition~\ref{def:general-tori}:
its period lattice is a subgroup of $\Z^3$ of rank three, and a character with the values $\pm1$ is unitary.
If $\rho^2\neq1$, Theorem~\ref{thm:floor}\eqref{it:floor-bounds} gives $0<Z_{\mathbb{T}}[1]<\infty$,
and Theorem~\ref{thm:floor}\eqref{it:floor-flux-terms} gives the positivity of every flux term.
If $\rho^2=1$, then $-1\le\rho\le1$, and both facts are contained in Theorem~\ref{thm:closed-interval}\eqref{it:closed-interval-flux}.
The cochain $z=0$ is closed.
This proves \eqref{it:partition-finite} and \eqref{it:partition-flux}.

Let $V$ be a real bounded Borel function on $\Conf(\Gamma)$.
Lemma~\ref{lem:perturbed-denominators}, which holds for every real $\rho$, gives
\[
Z_{\mathbb{T}}\left[\exp V\right]\ \ge\ \exp\left(-\sup|V|\right)\cdot Z_{\mathbb{T}}[1],
\]
and the right side is positive by \eqref{it:partition-finite}.
Let $\beta\in C^2(\Gamma)$.
The function $\exp\left(-\langle\beta,F\rangle-\frac12\|\beta\|^2\right)$ on $\Conf(\Gamma)$ is Borel and positive everywhere,
and the measure of $\Conf(\Gamma)$ is not zero, the classes of the pairs $(A,0)$ with $A\in[0,2\pi/q)^{\mathbb{T}^{(1)}}$ alone having the measure $(2\pi/q)^{3|\mathbb{T}|}$;
so this function does not vanish almost everywhere.
By the equivalence of Theorem~\ref{thm:floor}\eqref{it:floor-equivalence} if $\rho^2\neq1$,
and of Theorem~\ref{thm:closed-interval}\eqref{it:closed-interval-null} if $\rho^2=1$,
the number $Z_{\mathbb{T}}\left[\exp\left(-\langle\beta,F\rangle-\frac12\|\beta\|^2\right)\right]$ is not zero;
being non-negative, it is positive.
This proves \eqref{it:partition-perturbed}.
The last assertion holds because $\mathbb{E}_{\mathbb{T}}$ and $\mu_{\mathbb{T}}$ are defined as soon as $0<Z_{\mathbb{T}}[1]<\infty$.
\end{proof}

\section{Slab operators}\label{sec:slabs}

Throughout this section the lattice charge $q>0$, the base point $v_0\in C_0$, the hopping weights $w_0,\dots,w_3>0$ and the mass parameter $M\in\R$ are fixed
(Definitions~\ref{def:configurations}, \ref{def:fermion-sites} and~\ref{def:chamber-operator}),
and the Wilson parameter satisfies $-1\le\rho\le1$.

\subsection{Cells and slabs}\label{sec:cells-slabs}

\begin{definition}[Incidence numbers]\label{def:incidence}
Let $x\in\Z^3$ and $k\in\{0,1,2,3\}$, and let $\mu_1<\dots<\mu_k$ be elements of $\{1,2,3\}$.
The $k$-cell
\[
\sigma=(x;\mu_1\cdots\mu_k):=x+\sum_{j=1}^{k}[0,1]\,e_{\mu_j}
\]
is oriented by $e_{\mu_1}\wedge\dots\wedge e_{\mu_k}$,
and its direction space is $\vec\sigma:=\operatorname{span}\{e_{\mu_1},\dots,e_{\mu_k}\}$;
for $k=0$ the cell is the site $\{x\}$, with the orientation $1$ and $\vec\sigma=\{0\}$.
Thus $(x;\mu)$ is the link $(x,\mu)$, $(x;\mu\nu)$ is the plaquette of Definition~\ref{def:cochains},
and $(x;123)$ is the cube $x+[0,1]^3$.
For $k\ge1$ the faces of $\sigma$ are the $2k$ cells
\[
\left(x+e_{\mu_j};\mu_1\cdots\widehat{\mu_j}\cdots\mu_k\right)
\quad\text{and}\quad
\left(x;\mu_1\cdots\widehat{\mu_j}\cdots\mu_k\right),
\qquad 1\le j\le k,
\]
where the hat marks the omitted direction; a site has no faces.
For $1\le j\le k$ the incidence numbers of $\sigma$ and of its two faces without the direction $\mu_j$ are
\[
\begin{aligned}
\left[\sigma:\left(x+e_{\mu_j};\mu_1\cdots\widehat{\mu_j}\cdots\mu_k\right)\right]&:=(-1)^{j-1},\\
\left[\sigma:\left(x;\mu_1\cdots\widehat{\mu_j}\cdots\mu_k\right)\right]&:=-(-1)^{j-1},
\end{aligned}
\]
and $[\sigma:\sigma']:=0$ for every cell $\sigma'$ other than the faces of $\sigma$.
\end{definition}

With these numbers the coboundary of Definition~\ref{def:cochains} takes the form
\[
(du)(\sigma)=\sum_{\sigma'}[\sigma:\sigma']\cdot u(\sigma')
\]
for every $k$-cell $\sigma$ with $k\ge1$ and every cochain $u$ on the $(k-1)$-cells,
the sum extending over all $(k-1)$-cells $\sigma'$:
for $k=1,2,3$ these are the three formulas by which $d$ is defined there.
In particular, for a plaquette $\sigma$ and a link $\sigma'$, and for a cube $\sigma$ and a plaquette $\sigma'$,
the number $[\sigma:\sigma']$ is the coefficient of $u(\sigma')$ in $(du)(\sigma)$,
that is, the incidence number of these two cells introduced in Section~\ref{sec:cell-kinds}.

Let $g$ be an isometry of $\R^3$ with $g\Z^3=\Z^3$, and let $\vec g$ be its linear part.
Then $g$ maps every $k$-cell $\sigma=(x;\mu_1\cdots\mu_k)$ onto a $k$-cell $g\sigma$,
and $\vec g$ maps each $e_\mu$ to $\pm e_{\mu'}$ for some $\mu'$;
hence $\vec g e_{\mu_1}\wedge\dots\wedge\vec g e_{\mu_k}$ equals $\varepsilon(g,\sigma)$ times the orientation of $g\sigma$,
with a sign $\varepsilon(g,\sigma)\in\{1,-1\}$.
This is the sign of the pull-back of Definition~\ref{def:cochains}:
$(g^{*}u)(\sigma)=\varepsilon(g,\sigma)\cdot u(g\sigma)$.
For a translation $g$ one has $\varepsilon(g,\sigma)=1$.
The incidence numbers satisfy
\begin{equation}\label{eq:slab-incidence-square}
\sum_{\sigma'}[\sigma:\sigma']\cdot[\sigma':\sigma'']=0
\end{equation}
for all cells $\sigma$, $\sigma''$ of $\Z^3$, and
\begin{equation}\label{eq:slab-incidence-isometry}
\varepsilon(g,\sigma)\cdot[g\sigma:g\sigma']=[\sigma:\sigma']\cdot\varepsilon(g,\sigma')
\end{equation}
for all cells $\sigma$, $\sigma'$ and every isometry $g$ as above.
The first identity is $dd=0$.
A term of its sum can be non-zero only if $\sigma''$ arises from $\sigma=(x;\mu_1\cdots\mu_k)$
by omitting two directions $\mu_i$, $\mu_j$, $i<j$,
and moving the corner to $x+a_{\mu_i}\cdot e_{\mu_i}+a_{\mu_j}\cdot e_{\mu_j}$ with $a_{\mu_i},a_{\mu_j}\in\{0,1\}$;
then $\sigma''$ is a face of exactly two faces $\sigma'$ of $\sigma$,
and the two products of incidence numbers, $(-1)^{i+j-a_{\mu_i}-a_{\mu_j}}$ and $(-1)^{i+j-1-a_{\mu_i}-a_{\mu_j}}$, cancel.
The second identity is the relation $g^{*}d=dg^{*}$ of Definition~\ref{def:cochains},
applied to the indicator cochain $\mathbf{1}_{g\sigma'}$ of the cell $g\sigma'$ and evaluated at $\sigma$:
since $g^{*}\mathbf{1}_{g\sigma'}=\varepsilon(g,\sigma')\cdot\mathbf{1}_{\sigma'}$,
the two sides become $\varepsilon(g,\sigma)\cdot[g\sigma:g\sigma']$ and $\varepsilon(g,\sigma')\cdot[\sigma:\sigma']$.

\paragraph{The cylinder.}
Fix a family $\ell\in\mathcal{L}$ and a cross-section $\Xi=(\Gamma_\parallel,\chi_\parallel)$ for $\ell$
(Definitions~\ref{def:mirrors} and~\ref{def:slabs}).
The number $|n_\ell|^2/2$ equals $1/2$ for an axis family and $1$ for a diagonal family,
and $\Adm(\ell)=\frac{|n_\ell|^2}{2}\cdot\Z$.
For $c\in\Adm(\ell)$ put
\[
c^{\pm}:=c\pm|n_\ell|^2/2,
\qquad
\theta_c:=\theta_{\ell,c},
\qquad
\Theta_c:=\Theta_{\ell,c},
\]
so that $c^-<c<c^+$ are consecutive elements of $\Adm(\ell)$,
and $\Theta_c$ is the reflection of Definition~\ref{def:grassmann-reflections} at the mirror $(\ell,c)$.

The group $\Gamma_\parallel$ acts by translations on $\R^3$, on the cells of $\Z^3$ and on the set $\mathbb{X}$ of the fermion sites.
The cells of the cylinder $\Z^3/\Gamma_\parallel$ are the $\Gamma_\parallel$-orbits of cells of $\Z^3$,
and its fermion sites are the $\Gamma_\parallel$-orbits in $\mathbb{X}$.
For a family of cells of the cylinder, the cochains on it are the functions $u$,
defined on the cells of $\Z^3$ that belong to these orbits,
which satisfy $u(\sigma+g)=u(\sigma)$ for every such cell $\sigma$ and every $g\in\Gamma_\parallel$.
In Subsections~\ref{sec:cut-cells}, \ref{sec:slab-fermions}, \ref{sec:mirror-photon} and~\ref{sec:slab-transfer},
all cells, fermion sites and cochains are understood in this sense, that is, on the cylinder.
Since $\Gamma_\parallel$ lies in the plane $\{\ell=0\}$, the form $\ell$ is constant on orbits,
and $\theta_c(x+g)=\theta_c(x)+g$ for $x\in\R^3$ and $g\in\Gamma_\parallel$;
so $\theta_c$, like every translation $\tau_v$ with $v\in\Z^3$, maps orbits to orbits.

In the maximum norm $|\cdot|_\infty$ a cell of $\Z^3$ has diameter at most $1$.
Two cells that share a face have a point in common,
so a cell, or the union of two cells that share a face, has diameter at most $2$.
If $0\neq g\in\Gamma_\parallel$, then each point of a cell lies at the distance $|g|_\infty\ge\sys(\Gamma_\parallel)\ge4$
from its translate by $g$, which is a point of the translated cell.
Hence a set of diameter at most $2$ never contains a cell together with another cell of its orbit;
distinct cells contained in one cell of $\Z^3$, or in the union of two cells that share a face, belong to distinct orbits.
This applies to the faces of a cell and to the faces of these faces,
as well as to the faces of two cells that share a face.
The incidence number of two cells $\sigma$, $\sigma'$ of the cylinder is therefore well defined:
$[\sigma:\sigma']:=[\sigma_0:\sigma'_0]$ if $\sigma_0$ is a cell of the orbit $\sigma$
and $\sigma'_0$ is a cell of the orbit $\sigma'$ that is a face of $\sigma_0$ (there is at most one),
and $[\sigma:\sigma']:=0$ if no cell of $\sigma'$ is a face of $\sigma_0$;
by \eqref{eq:slab-incidence-isometry} for translations, the choice of $\sigma_0$ does not matter.
The coboundary of $\Gamma_\parallel$-periodic cochains is $(du)(\sigma)=\sum_{\sigma'}[\sigma:\sigma']\cdot u(\sigma')$,
the sum over the cells $\sigma'$ of the cylinder,
and \eqref{eq:slab-incidence-square} and \eqref{eq:slab-incidence-isometry} hold for the cells of the cylinder,
the latter for $g=\theta_c$, for $g=\tau_v$ and for compositions of such maps.
Both identities for the cylinder follow from those for $\Z^3$,
applied to a representative $\sigma_0$ of $\sigma$ and to the cells contained in $\sigma_0$:
these cells lie in different orbits;
each of the maps $g$ just named commutes with the translations by the elements of $\Gamma_\parallel$;
and $\varepsilon(g,\sigma_0)$ depends only on the directions of $\sigma_0$, so that it is the same for all cells of the orbit $\sigma$.

For a cell $\sigma$ let $\min_\sigma\ell$ and $\max_\sigma\ell$ be the least and the greatest height of its points.
They are attained at corners of $\sigma$, hence are integers ($n_\ell$ is an integer vector),
and they are the same for all cells of an orbit.
The cells of the cylinder all of whose points have heights in a given bounded interval form a finite set.
This holds because $\Gamma_\parallel$ has rank two and lies in the group $\Z^3\cap\{\ell=0\}$, which has rank two as well,
so that its index there is finite and the sites of $\Z^3$ of one height form finitely many orbits,
while a cell $(x;\mu_1\cdots\mu_k)$ is determined by its directions and by its corner $x$,
whose height lies between $\min_\sigma\ell$ and $\max_\sigma\ell$.

\begin{definition}[Slab configurations]\label{def:slab-configurations}
Let $c\in\Adm(\ell)$.
For $k=0,\dots,3$ the $k$-cells of the cylinder contained in the plane $\{\ell=c\}$ form the set $\mathcal{M}^k(c)$ of \emph{cells in the mirror},
and the $k$-cells $\sigma$ of the cylinder with $\min_\sigma\ell<c<\max_\sigma\ell$ form the set $\mathcal{X}^k(c)$ of \emph{cells cut by the mirror}.
For $c'\in\Adm(\ell)$ with $c<c'$, the $k$-cells of the cylinder that are contained in the closed slab $\{c\le\ell\le c'\}$
but in neither of the planes $\{\ell=c\}$ and $\{\ell=c'\}$ form the set $\mathcal{B}^k(c,c')$ of \emph{slab cells};
a slab cell may touch the two planes.
The unions over $k=0,\dots,3$ are denoted $\mathcal{M}(c)$, $\mathcal{X}(c)$ and $\mathcal{B}(c,c')$.
The links of the cylinder that are contained in the closed slab form the finite set
\[
[c,c']^{(1)}:=\mathcal{M}^1(c)\cup\mathcal{B}^1(c,c')\cup\mathcal{M}^1(c'),
\]
and the plaquettes of the cylinder lying in the closed slab form the finite set
\[
[c,c']^{(2)}:=\mathcal{M}^2(c)\cup\mathcal{B}^2(c,c')\cup\mathcal{M}^2(c').
\]

\emph{Configurations.}
Let $c<c'$ be elements of $\Adm(\ell)$.
A configuration of the closed slab $\{c\le\ell\le c'\}$ is a pair $\varpi=(A,z)\in\R^{[c,c']^{(1)}}\times\Z^{[c,c']^{(2)}}$.
The faces of a cell are subsets of it;
so the faces of a plaquette in $[c,c']^{(2)}$ are links in $[c,c']^{(1)}$,
and the faces of a cube of the cylinder lying in the closed slab are plaquettes in $[c,c']^{(2)}$.
Since a cube lies in no plane, these cubes are exactly the elements of $\mathcal{B}^3(c,c')$.
Hence the curvature $F(\varpi):=dA+\frac{2\pi}{q}\cdot z$ is defined on $[c,c']^{(2)}$,
and $dz$ is defined on $\mathcal{B}^3(c,c')$.
Put
\[
\mathfrak{C}[c,c']:=\left\{\varpi=(A,z):\ (dz)(\sigma)=0\ \text{for all}\ \sigma\in\mathcal{B}^3(c,c')\right\}.
\]

\emph{Integer shifts.}
For an integer one-cochain $k\in\Z^{[c,c']^{(1)}}$ the integer shift of $\varpi=(A,z)$ is
\[
\varpi+[k]:=\left(A+\frac{2\pi}{q}\cdot k,\ z-dk\right).
\]
It changes neither $F(\varpi)$, nor $dz$ on $\mathcal{B}^3(c,c')$, nor the link phases $\exp\left(iqA(b)\right)$, $b\in[c,c']^{(1)}$,
because $d$ applied twice vanishes by \eqref{eq:slab-incidence-square} and $\exp\left(2\pi i\cdot k(b)\right)=1$.

\emph{Parts of a configuration.}
Mirror data at $c$ are pairs $\mathfrak{m}\in\R^{\mathcal{M}^1(c)}\times\Z^{\mathcal{M}^2(c)}$.
For mirror data $\mathfrak{m}$ at $c$ and $\mathfrak{m}'$ at $c'$ and for $\mathring{\varpi}\in\R^{\mathcal{B}^1(c,c')}\times\Z^{\mathcal{B}^2(c,c')}$,
the configuration $\mathfrak{m}\vee\mathring{\varpi}\vee\mathfrak{m}'$ of the closed slab is the one that agrees with $\mathfrak{m}$ on the links and plaquettes of $\mathcal{M}(c)$,
with $\mathring{\varpi}$ on those of $\mathcal{B}(c,c')$, and with $\mathfrak{m}'$ on those of $\mathcal{M}(c')$.
Every configuration $\varpi$ of the closed slab is of this form, and $\mathring{\varpi}$ is called its interior part.

\emph{Pull-backs.}
Let $g$ be a reflection $\theta_{c''}$ with $c''\in\Adm(\ell)$, a translation $\tau_v$ with $v\in\Z^3$, or a composition of finitely many such maps.
Then $g$ maps the closed slab $\{c\le\ell\le c'\}$ onto a closed slab $\{c_1\le\ell\le c_1'\}$ with $c_1<c_1'$ in $\Adm(\ell)$;
for $g=\theta_{c''}$ one has $(c_1,c_1')=(2c''-c',2c''-c)$, and for $g=\tau_v$ one has $(c_1,c_1')=(c+\ell(v),c'+\ell(v))$.
For a configuration $\varpi=(A,z)$ of the latter slab, $g^{*}\varpi:=(g^{*}A,g^{*}z)$ is a configuration of the former.
By \eqref{eq:slab-incidence-isometry} the pull-back $g^{*}$ commutes with $d$;
hence $F(g^{*}\varpi)=g^{*}F(\varpi)$, and $g^{*}$ maps $\mathfrak{C}[c_1,c_1']$ onto $\mathfrak{C}[c,c']$.

\emph{Slab weight.}
For a configuration $\varpi=(A,z)$ of the closed slab,
\[
\mathcal{W}_{c,c'}(\varpi):=\prod_{\sigma\in\mathcal{B}^2(c,c')}\exp\left(-\frac{1}{2}F_\sigma(\varpi)^2\right)
\cdot\prod_{\sigma\in\mathcal{B}^3(c,c')}\mathbf{1}\left[(dz)(\sigma)=0\right],
\]
where $\mathbf{1}[\,\cdot\,]$ is the indicator of the condition in the brackets: it equals $1$ if the condition holds and $0$ otherwise.

\emph{Interior integral.}
Let $\mathfrak{m}$ and $\mathfrak{m}'$ be mirror data at $c$ and at $c'$,
and let $\Phi\colon\R^{[c,c']^{(1)}}\times\Z^{[c,c']^{(2)}}\to\C$ be a Borel function.
The interior integral of $\Phi$ at $(\mathfrak{m},\mathfrak{m}')$ is
\[
\int\mathrm{d}\mathring{\varpi}\,\Phi\left(\mathfrak{m}\vee\mathring{\varpi}\vee\mathfrak{m}'\right)
:=\sum_{z\in\Z^{\mathcal{B}^2(c,c')}}\ \int_{[0,2\pi/q)^{\mathcal{B}^1(c,c')}}\mathrm{d}A\,\Phi\left(\mathfrak{m}\vee(A,z)\vee\mathfrak{m}'\right),
\]
with the Lebesgue measure $\mathrm{d}A$ on the finite-dimensional cube $[0,2\pi/q)^{\mathcal{B}^1(c,c')}$.
For $\Phi\ge0$ it is an element of $[0,\infty]$;
for general $\Phi$ it is defined, as an absolutely convergent sum of integrals, whenever the interior integral of $|\Phi|$ is finite.
\end{definition}

\begin{lemma}\label{lem:fundamental-intervals}
Let $c<c'$ be elements of $\Adm(\ell)$, let $\mathfrak{m}$ and $\mathfrak{m}'$ be mirror data at $c$ and at $c'$,
and let $\Phi\colon\R^{[c,c']^{(1)}}\times\Z^{[c,c']^{(2)}}\to\C$ be a Borel function such that
\[
\int\mathrm{d}\mathring{\varpi}\,\left|\Phi\left(\mathfrak{m}\vee\mathring{\varpi}\vee\mathfrak{m}'\right)\right|<\infty
\]
and $\Phi(\varpi+[k])=\Phi(\varpi)$ for every configuration $\varpi$ of the closed slab and every $k\in\Z^{\mathcal{B}^1(c,c')}$, extended by $0$ to $[c,c']^{(1)}$.
For each $b\in\mathcal{B}^1(c,c')$ let $\mathbb{I}_b\subset\R$ be an interval of length $2\pi/q$, open, closed or half open.
Then the sum over $z\in\Z^{\mathcal{B}^2(c,c')}$ of the integrals of $\Phi\left(\mathfrak{m}\vee(A,z)\vee\mathfrak{m}'\right)$ over $A\in\prod_b\mathbb{I}_b$ converges absolutely, and
\[
\sum_{z\in\Z^{\mathcal{B}^2(c,c')}}\ \int_{\prod_b\mathbb{I}_b}\mathrm{d}A\,\Phi\left(\mathfrak{m}\vee(A,z)\vee\mathfrak{m}'\right)
=\int\mathrm{d}\mathring{\varpi}\,\Phi\left(\mathfrak{m}\vee\mathring{\varpi}\vee\mathfrak{m}'\right).
\]
\end{lemma}

\begin{proof}
An interval of length $2\pi/q$ differs from the half open interval of the same length with the same left end point by at most two points, a null set;
so every $\mathbb{I}_b$ may be assumed to be of the form $[s,s+2\pi/q)$.
The intervals $[0,2\pi/q)$ are replaced by the intervals $\mathbb{I}_b$ one link at a time.
Consider a step at which the interval of the link $b\in\mathcal{B}^1(c,c')$ is changed from $[0,2\pi/q)$ to $\mathbb{I}_b=[s,s+2\pi/q)$,
while every other link $b'$ carries an interval of length $2\pi/q$ which is not changed in this step,
and the sum of the integrals of $|\Phi|$ over the current product of intervals is finite;
at the first step this is the hypothesis.
Write $A=(A_b,A_{\neq b})$, where $A_{\neq b}=(A_{b'})_{b'\neq b}$ ranges over the product of the intervals of the links $b'\neq b$,
and define, for $A_b\in\R$,
\[
\begin{aligned}
f^{+}(A_b)&:=\sum_{z\in\Z^{\mathcal{B}^2(c,c')}}\int\mathrm{d}A_{\neq b}\,\left|\Phi\left(\mathfrak{m}\vee(A,z)\vee\mathfrak{m}'\right)\right|\in[0,\infty],\\
f(A_b)&:=\sum_{z\in\Z^{\mathcal{B}^2(c,c')}}\int\mathrm{d}A_{\neq b}\,\Phi\left(\mathfrak{m}\vee(A,z)\vee\mathfrak{m}'\right),
\end{aligned}
\]
the latter wherever $f^{+}(A_b)<\infty$.
By Tonelli's theorem $f^{+}$ is a Borel function, and $\int_0^{2\pi/q}\mathrm{d}A_b\,f^{+}(A_b)$ is the finite sum of integrals of $|\Phi|$ at the current step.

Let $\mathbf{1}_b\in\Z^{\mathcal{B}^1(c,c')}$ be the indicator cochain of $b$, extended by $0$ to $[c,c']^{(1)}$.
Its coboundary vanishes on $\mathcal{M}^2(c)\cup\mathcal{M}^2(c')$,
since the faces of a plaquette contained in one of the two planes lie in that plane and $b$ lies in neither.
Hence the integer shift by $\mathbf{1}_b$ does not change the mirror data:
for every interior part $(A,z)$,
\[
\left(\mathfrak{m}\vee(A,z+d\mathbf{1}_b)\vee\mathfrak{m}'\right)+[\mathbf{1}_b]=\mathfrak{m}\vee\left(A+\frac{2\pi}{q}\cdot\mathbf{1}_b,\ z\right)\vee\mathfrak{m}',
\]
where on the left $d\mathbf{1}_b$ stands for its restriction to $\mathcal{B}^2(c,c')$.
By the invariance of $\Phi$ under the integer shifts,
\[
\Phi\left(\mathfrak{m}\vee\left(A+\frac{2\pi}{q}\cdot\mathbf{1}_b,\ z\right)\vee\mathfrak{m}'\right)=\Phi\left(\mathfrak{m}\vee(A,z+d\mathbf{1}_b)\vee\mathfrak{m}'\right).
\]
As $z\mapsto z+d\mathbf{1}_b$ is a bijection of $\Z^{\mathcal{B}^2(c,c')}$,
this gives $f^{+}(A_b+2\pi/q)=f^{+}(A_b)$ for all $A_b\in\R$,
and $f(A_b+2\pi/q)=f(A_b)$ wherever $f^{+}(A_b)<\infty$, the sum over $z$ being absolutely convergent there and so open to reordering.
Thus $f^{+}$ is periodic with period $2\pi/q$ and integrable over $[0,2\pi/q)$;
it is finite almost everywhere and integrable over every bounded interval.
The function $f$ is defined almost everywhere, periodic with period $2\pi/q$, and $|f|\le f^{+}$;
by Fubini's theorem, $\int_a^{a'}\mathrm{d}A_b\,f(A_b)$ is the sum over $z$ of the integrals of $\Phi\left(\mathfrak{m}\vee(A,z)\vee\mathfrak{m}'\right)$
over $A_b\in[a,a')$ and $A_{\neq b}$ as above, for all real $a<a'$.

It remains to show that
\[
\int_s^{s+2\pi/q}\mathrm{d}A_b\,f(A_b)=\int_0^{2\pi/q}\mathrm{d}A_b\,f(A_b),
\]
and the same identity for $f^{+}$, which then shows that the sum of the integrals of $|\Phi|$ stays finite after the step.
For real $a<a'$ the substitution $A_b\mapsto A_b+2\pi/q$ and the periodicity give
\begin{equation}\label{eq:slab-period-shift}
\int_{a+2\pi/q}^{a'+2\pi/q}\mathrm{d}A_b\,f(A_b)=\int_a^{a'}\mathrm{d}A_b\,f(A_b).
\end{equation}
Applied to $a-2\pi/q$ and $a'-2\pi/q$, it gives the shift by $-2\pi/q$;
by iteration, the integral over $[a,a')$ does not change if $a$ and $a'$ are moved by the same integer multiple of $2\pi/q$.
Write $s=2\pi j/q+s_0$ with $j\in\Z$ and $0\le s_0<2\pi/q$.
Then the integral of $f$ over $[s,s+2\pi/q)$ equals its integral over $[s_0,s_0+2\pi/q)$.
If $s_0>0$, this interval is the disjoint union of $[s_0,2\pi/q)$ and $[2\pi/q,2\pi/q+s_0)$,
and \eqref{eq:slab-period-shift} with $a=0$ and $a'=s_0$ turns the integral over the second piece into the integral over $[0,s_0)$;
together with the integral over $[s_0,2\pi/q)$ this is the integral over $[0,2\pi/q)$.
The same computation, with values in $[0,\infty]$, applies to $f^{+}$.
Hence the step keeps the sum of the integrals of $|\Phi|$ finite and does not change the sum of the integrals of $\Phi$,
and after the last link of $\mathcal{B}^1(c,c')$ every link $b$ carries its interval $\mathbb{I}_b$.
\end{proof}

\subsection{Cells cut by a mirror}\label{sec:cut-cells}

Throughout this subsection a family $\ell\in\mathcal{L}$ and a cross-section $\Xi=(\Gamma_\parallel,\chi_\parallel)$ for $\ell$ are fixed,
and every cell is a cell of the cylinder $\Z^3/\Gamma_\parallel$.
For $c<c'$ in $\Adm(\ell)$ and $k=0,1,2,3$, the sets $\mathcal{M}^k(c)$, $\mathcal{X}^k(c)$ and $\mathcal{B}^k(c,c')$ of $k$-cells,
as well as their unions $\mathcal{M}(c)$, $\mathcal{X}(c)$ and $\mathcal{B}(c,c')$ over $k$, are those of Definition~\ref{def:slab-configurations}.
As in Subsection~\ref{sec:cells-slabs}, $\min_\sigma\ell$ and $\max_\sigma\ell$ are the least and the greatest height of the points of a cell $\sigma$;
they do not depend on the representative of the orbit $\sigma$, are attained at corners of $\sigma$ and are integers.
A cell $\sigma$ is \emph{cut by} $c$ if $\sigma\in\mathcal{X}(c)$, that is, if $\min_\sigma\ell<c<\max_\sigma\ell$.
The positions $c^{\pm}=c\pm|n_\ell|^2/2$, introduced in the same subsection, are the neighbours of $c$ in $\Adm(\ell)=\frac{|n_\ell|^2}{2}\cdot\Z$:
$c^+$ is the least element of $\Adm(\ell)$ larger than $c$, and $c^-$ is the largest element smaller than $c$.
For $c,c'\in\Adm(\ell)$ put
\[
v_{c,c'}:=\frac{2}{|n_\ell|^2}\cdot(c'-c)\cdot n_\ell .
\]
By Definition~\ref{def:mirrors}\eqref{it:mirror-product}, $v_{c,c'}$ is an integer vector of height $\ell(v_{c,c'})=2(c'-c)$,
and $\theta_{c'}\theta_c=\tau_{v_{c,c'}}$.

Let $\sigma=(x;\mu_1\cdots\mu_k)$, $\mu_1<\dots<\mu_k$, be a $k$-cell (Definition~\ref{def:incidence}).
The indices $\mu_1,\dots,\mu_k$ are the directions of $\sigma$,
and the direction space $\vec\sigma$ is the span of $e_{\mu_1},\dots,e_{\mu_k}$.
The points of $\sigma$ are the points $y=x+\sum_\mu a_\mu\cdot e_\mu$,
where $\mu$ ranges over the directions of $\sigma$ and $a_\mu\in[0,1]$;
since $\ell$ is linear,
\begin{equation}\label{eq:slab-cell-heights}
\ell(y)=\ell(x)+\sum_{\mu}a_\mu\cdot\ell(e_\mu).
\end{equation}

\begin{lemma}[Cells and parallel mirrors]\label{lem:cells-mirrors}
Let $\ell\in\mathcal{L}$, let $\Xi=(\Gamma_\parallel,\chi_\parallel)$ be a cross-section for $\ell$,
let $c\in\Adm(\ell)$, and put $\theta:=\theta_c$.
The cells of $\Z^3/\Gamma_\parallel$ have the following properties.
\begin{enumerate}
\item\label{it:cut-heights}
For every cell $\sigma$, the difference $\max_\sigma\ell-\min_\sigma\ell$ lies in $\{0,1\}$ if $\ell$ is an axis family
and in $\{0,1,2\}$ if $\ell$ is a diagonal family,
and at most one element of $\Adm(\ell)$ lies strictly between $\min_\sigma\ell$ and $\max_\sigma\ell$.
A cell $\sigma$ cut by $c$ has the least height $c^-$ and the greatest height $c^+$, and its direction space $\vec\sigma$ contains $n_\ell$.
In particular no cell is cut by two different elements of $\Adm(\ell)$,
and every cell cut by $c$ is contained in the closed slab $\{c^-\le\ell\le c^+\}$.
\item\label{it:cut-orientation}
If $\sigma$ is cut by $c$, then $\theta\sigma=\sigma$ and $\theta$ reverses the orientation of $\sigma$;
that is, $(\theta^{*}u)(\sigma)=-u(\sigma)$ for every cochain $u$.
If $\sigma$ lies in $\{\ell=c\}$, then $\theta$ fixes $\sigma$ pointwise and preserves its orientation.
More generally, if $\sigma$ lies in $\{\ell=c\}$ and $c'\in\Adm(\ell)$, then $\theta_{c'}\sigma=\sigma+v_{c,c'}$,
and $\theta_{c'}$ maps the orientation of $\sigma$ onto the orientation of $\sigma+v_{c,c'}$;
that is, $(\theta_{c'}^{*}u)(\sigma)=u(\sigma+v_{c,c'})$ for every cochain $u$.
\item\label{it:cut-faces}
Let $\sigma$ be cut by $c$ and of dimension $2$ or $3$.
Every face $\sigma'$ of $\sigma$ is of exactly one of the following three kinds:
$\sigma'$ is cut by $c$;
or $\sigma'\subset\{c\le\ell\le c^+\}$ and $\sigma'\not\subset\{\ell=c\}$, and then $\sigma'$ is called an \emph{upper} face of $\sigma$;
or $\sigma'\subset\{c^-\le\ell\le c\}$ and $\sigma'\not\subset\{\ell=c\}$, and then $\sigma'$ is called a \emph{lower} face of $\sigma$.
The reflection $\theta$ maps the set of the upper faces of $\sigma$ bijectively onto the set of the lower faces,
and for every face $\sigma'$ of $\sigma$ and every cochain $u$
\[
[\sigma:\theta\sigma']\cdot u(\theta\sigma')=-[\sigma:\sigma']\cdot(\theta^{*}u)(\sigma') .
\]
In other words, $[\sigma:\theta\sigma']=-[\sigma:\sigma']$ if $\theta$ maps the orientation of $\sigma'$ onto the orientation of $\theta\sigma'$,
and $[\sigma:\theta\sigma']=[\sigma:\sigma']$ otherwise.
\item\label{it:cell-partition}
Let $(c_j)_{j\in\Z}$ be a strictly increasing sequence in $\Adm(\ell)$ that is unbounded above and below.
The sets $\mathcal{M}(c_j)$, $\mathcal{X}(c_j)$ and $\mathcal{B}(c_{j-1},c_j)$, $j\in\Z$, are pairwise disjoint,
and each cell of the cylinder lies in one of them.
For $c',c''\in\Adm(\ell)$ with $c<c'<c''$ and for $k=0,1,2,3$,
\[
\mathcal{B}^k(c,c'')=\mathcal{B}^k(c,c')\cup\mathcal{M}^k(c')\cup\mathcal{X}^k(c')\cup\mathcal{B}^k(c',c''),
\]
and the four sets on the right are pairwise disjoint.
\item\label{it:upper-faces}
Let $\sigma$ be as in \eqref{it:cut-faces}.
An upper face of $\sigma$ belongs to $\mathcal{B}(c,c')\cup\mathcal{M}(c')$ for every $c'\in\Adm(\ell)$ with $c'>c$,
and a lower face of $\sigma$ belongs to $\mathcal{M}(c')\cup\mathcal{B}(c',c)$ for every $c'\in\Adm(\ell)$ with $c'<c$.
\end{enumerate}
\end{lemma}

\begin{proof}
Let $\sigma=(x;\mu_1\cdots\mu_k)$, and let $y$ and the numbers $a_\mu$ be as in \eqref{eq:slab-cell-heights}.

\eqref{it:cut-heights}
Let first $\ell=x_\mu$ be an axis family.
By \eqref{eq:slab-cell-heights}, the function $\ell-\ell(x)$ takes on $\sigma$ the values $a_\mu\in[0,1]$ if $\mu$ is a direction of $\sigma$,
and only the value $0$ otherwise.
In the second case the open interval $(\min_\sigma\ell,\max_\sigma\ell)$ is empty.
In the first case $\min_\sigma\ell=\ell(x)$ and $\max_\sigma\ell=\ell(x)+1$,
and the interval $\left(\ell(x),\ell(x)+1\right)$ meets $\Adm(\ell)=\frac12\Z$ only in the point $\ell(x)+\frac12$.
If this point is $c$, then $\min_\sigma\ell=c-\frac12=c^-$, $\max_\sigma\ell=c^+$, and $n_\ell=e_\mu\in\vec\sigma$.

Now let $\ell$ be a diagonal family, $\ell=x_\mu-x_\nu$ or $\ell=x_\mu+x_\nu$ with $\mu<\nu$.
Then $\ell(e_\mu)=1$, $\ell(e_\nu)=\pm1$, the third unit vector lies in the kernel of $\ell$,
and $n_\ell=e_\mu+\ell(e_\nu)\cdot e_\nu$.
By \eqref{eq:slab-cell-heights},
\[
\ell(y)-\ell(x)=a_\mu+\ell(e_\nu)\cdot a_\nu ,
\]
where $a_\mu$ (resp.\ $a_\nu$) is to be read as $0$ if $\mu$ (resp.\ $\nu$) is not a direction of $\sigma$.
If $\mu$ and $\nu$ are both directions of $\sigma$, the values of $\ell-\ell(x)$ on $\sigma$ fill the interval $[0,2]$ if $\ell(e_\nu)=1$,
and the interval $[-1,1]$ if $\ell(e_\nu)=-1$.
Then $(\min_\sigma\ell,\max_\sigma\ell)$ is an open interval of length $2$ with integer end points;
the only integer in it is its midpoint, and $n_\ell\in\vec\sigma$.
If exactly one of $\mu,\nu$ is a direction of $\sigma$, the interval has length $1$ and integer end points, so it contains no integer;
if neither is, it is empty.
As $\Adm(\ell)=\Z$ and $c^{\pm}=c\pm1$, the first three assertions follow.

Two different positions cutting $\sigma$ would be two points of $\Adm(\ell)$ in $(\min_\sigma\ell,\max_\sigma\ell)$.
If $\sigma$ is cut by $c$, then $c^-=\min_\sigma\ell\le\ell\le\max_\sigma\ell=c^+$ on $\sigma$.

\eqref{it:cut-orientation}
Let $\sigma$ be cut by $c$.
By Definition~\ref{def:mirrors},
\[
\theta(y)=y-\frac{2}{|n_\ell|^2}\cdot\left(\ell(y)-c\right)\cdot n_\ell .
\]
For an axis family $\ell=x_\mu$, part \eqref{it:cut-heights} shows that $\mu$ is a direction of $\sigma$ and that $x_\mu=\min_\sigma\ell=c-\frac12$.
Hence $\ell(y)-c=a_\mu-\frac12$ and $\theta(y)=y-(2a_\mu-1)\cdot e_\mu$:
the coordinate $a_\mu$ is replaced by $1-a_\mu$, and the others are unchanged.
For a diagonal family, $n_\ell\in\vec\sigma$ shows that $\mu$ and $\nu$ are both directions of $\sigma$, and $\theta(y)=y-\left(\ell(y)-c\right)\cdot n_\ell$.
If $\ell(e_\nu)=-1$, then $c=\ell(x)$ is the midpoint of $[\ell(x)-1,\ell(x)+1]$, $\ell(y)-c=a_\mu-a_\nu$, and
\[
\theta(y)=y-(a_\mu-a_\nu)\cdot(e_\mu-e_\nu):
\]
the coordinates $a_\mu$ and $a_\nu$ are exchanged.
If $\ell(e_\nu)=1$, then $c=\ell(x)+1$, $\ell(y)-c=a_\mu+a_\nu-1$, and
\[
\theta(y)=y-(a_\mu+a_\nu-1)\cdot(e_\mu+e_\nu):
\]
the pair $(a_\mu,a_\nu)$ is replaced by $(1-a_\nu,1-a_\mu)$.
In each of the three cases the coordinates are transformed by an involution of $[0,1]^k$, so $\theta\sigma=\sigma$.

The linear part of $\theta$ is the orthogonal reflection of $\R^3$ in the plane $n_\ell^{\perp}$.
As $n_\ell\in\vec\sigma$, this reflection maps $\vec\sigma$ onto itself,
and its restriction to $\vec\sigma$ is the reflection in $\vec\sigma\cap n_\ell^{\perp}$, of determinant $-1$.
So $\theta$ maps the orientation of $\sigma$ onto the opposite of the orientation of $\theta\sigma=\sigma$,
and Definition~\ref{def:cochains} gives $(\theta^{*}u)(\sigma)=-u(\theta\sigma)=-u(\sigma)$.

Let now $\sigma$ lie in $\{\ell=c\}$.
Then $\ell$ is constant on $\sigma$, the space $\vec\sigma$ is orthogonal to $n_\ell$,
and $\theta$ is the identity on the plane $\{\ell=c\}$; in particular it fixes $\sigma$ pointwise.
For $c'\in\Adm(\ell)$, Definition~\ref{def:mirrors}\eqref{it:mirror-involution} and \eqref{it:mirror-product} give
\[
\theta_{c'}=(\theta_{c'}\theta_c)\,\theta_c=\tau_{v_{c,c'}}\theta_c ,
\]
hence $\theta_{c'}\sigma=\sigma+v_{c,c'}$.
The linear part of $\theta_{c'}$ is again the reflection in $n_\ell^{\perp}$, which is the identity on $\vec\sigma$.
It therefore fixes $e_{\mu_1}\wedge\dots\wedge e_{\mu_k}$, which orients $\sigma$ as well as $\sigma+v_{c,c'}$ (Definition~\ref{def:incidence}).
This is the assertion on $\theta_{c'}$, and $c'=c$ gives the assertion on $\theta$.

\eqref{it:cut-faces}
Let $\sigma'$ be a face of $\sigma$.
By \eqref{it:cut-heights}, $\sigma'\subset\sigma\subset\{c^-\le\ell\le c^+\}$.
If $\sigma'$ is not cut by $c$, then $\max_{\sigma'}\ell\le c$ or $\min_{\sigma'}\ell\ge c$,
so that $\sigma'$ lies in $\{\ell\le c\}$ or in $\{\ell\ge c\}$.
Suppose that $\sigma'\subset\{\ell=c\}$; then $\vec{\sigma'}$ is orthogonal to $n_\ell$.
The directions of the face $\sigma'$ are those of $\sigma$ except one.
For an axis family $\ell=x_\mu$, the index $\mu$ is a direction of $\sigma$ by \eqref{it:cut-heights}, and as $\vec{\sigma'}$ is orthogonal to $e_\mu$ it is the missing one;
so $\sigma'$ is one of the two faces of $\sigma$ on which $a_\mu=0$, resp.\ $a_\mu=1$.
Thus $\sigma=\sigma'+[0,1]\cdot e_\mu$ or $\sigma=\sigma'-[0,1]\cdot e_\mu$,
and $\sigma$ lies in $\{c\le\ell\le c+1\}$ or in $\{c-1\le\ell\le c\}$;
then $\sigma$ is not cut by $c$, a contradiction.
For a diagonal family, $\mu$ and $\nu$ are directions of $\sigma$ by \eqref{it:cut-heights}, so at least one of them is a direction of $\sigma'$;
as $\ell(e_\mu)=1$ and $\ell(e_\nu)=\pm1$, the space $\vec{\sigma'}$ is not orthogonal to $n_\ell$, again a contradiction.
Hence every face that is not cut by $c$ is an upper or a lower face.
No face is of two kinds:
a face cut by $c$ lies neither in $\{\ell\le c\}$ nor in $\{\ell\ge c\}$,
and a face that is upper and lower at the same time lies in $\{\ell=c\}$.

By \eqref{it:cut-orientation}, $\theta\sigma=\sigma$.
The reflection $\theta$ is an isometry of $\R^3$ that maps $\Z^3$ onto itself,
and $\theta(y+g)=\theta(y)+g$ for $g\in\Gamma_\parallel$, because $\ell(g)=0$;
so $\theta$ permutes the cells of $\Z^3/\Gamma_\parallel$, and in particular the faces of $\sigma$.
By Definition~\ref{def:mirrors}\eqref{it:mirror-height}, $\ell\circ\theta=2c-\ell$.
So $\theta$ exchanges $\{\ell\le c\}$ and $\{\ell\ge c\}$, maps $\{\ell=c\}$ onto itself,
and maps $\{c\le\ell\le c^+\}$ onto $\{c^-\le\ell\le c\}$, since $2c-c^+=c^-$.
Therefore the image of an upper face is a lower face, and the image of a lower face is an upper face.
As $\theta^2$ is the identity, $\theta$ maps the upper faces bijectively onto the lower faces.

For the identity, recall that $\theta^{*}$ commutes with $d$ (Definition~\ref{def:cochains} and \eqref{eq:slab-incidence-isometry}).
Let $u$ be a cochain.
Since $\theta\sigma=\sigma$ and $\theta$ reverses the orientation of $\sigma$,
\[
(\theta^{*}du)(\sigma)=-(du)(\sigma)=-\sum_{\sigma''}[\sigma:\sigma'']\cdot u(\sigma'')=-\sum_{\sigma'}[\sigma:\theta\sigma']\cdot u(\theta\sigma'),
\]
where the last sum is the same sum written with $\sigma''=\theta\sigma'$, the map $\theta$ being a permutation of the faces of $\sigma$;
on the other hand
\[
(d\theta^{*}u)(\sigma)=\sum_{\sigma'}[\sigma:\sigma']\cdot(\theta^{*}u)(\sigma').
\]
Fix a face $\sigma'_0$ and apply the equality of the two right-hand sides to the cochain $u_0$
that equals $u(\theta\sigma'_0)$ at $\theta\sigma'_0$ and vanishes at every other cell.
As $(\theta^{*}u_0)(\sigma')=\pm u_0(\theta\sigma')$, on each side only the term of $\sigma'=\sigma'_0$ survives, and
\[
-[\sigma:\theta\sigma'_0]\cdot u(\theta\sigma'_0)=[\sigma:\sigma'_0]\cdot(\theta^{*}u_0)(\sigma'_0)=[\sigma:\sigma'_0]\cdot(\theta^{*}u)(\sigma'_0),
\]
because $(\theta^{*}u)(\sigma'_0)$ depends on $u$ only through $u(\theta\sigma'_0)$.
This is the asserted identity.
For $u$ the indicator of $\theta\sigma'_0$, the number $(\theta^{*}u)(\sigma'_0)$ is $1$ or $-1$
according as $\theta$ maps the orientation of $\sigma'_0$ onto the orientation of $\theta\sigma'_0$ or onto its opposite,
which gives the formulation with signs.

\eqref{it:cell-partition}
Let $\sigma$ be a cell.
Suppose first that $\min_\sigma\ell=\max_\sigma\ell=c_j$ for some $j$; then $\sigma\in\mathcal{M}(c_j)$.
The cell lies in no $\mathcal{M}(c_i)$ with $i\ne j$, since its height is $c_j$,
and in no $\mathcal{X}(c_i)$, since the open interval $(\min_\sigma\ell,\max_\sigma\ell)$ is empty.
If $\sigma\subset\{c_{i-1}\le\ell\le c_i\}$ for some $i$, then $c_{i-1}\le c_j\le c_i$,
so $j\in\{i-1,i\}$ because the sequence is strictly increasing;
thus $\sigma$ lies in one of the two planes $\{\ell=c_{i-1}\}$, $\{\ell=c_i\}$, and $\sigma\notin\mathcal{B}(c_{i-1},c_i)$.

Suppose next that this is not the case, but that some $c_j$ lies in $(\min_\sigma\ell,\max_\sigma\ell)$.
Then $\sigma\in\mathcal{X}(c_j)$, and $j$ is unique by \eqref{it:cut-heights}.
Since $\min_\sigma\ell<\max_\sigma\ell$, the cell lies in no plane $\{\ell=c_i\}$.
It is contained in no closed slab $\{c_{i-1}\le\ell\le c_i\}$,
for otherwise $c_{i-1}\le\min_\sigma\ell<c_j<\max_\sigma\ell\le c_i$, which is impossible for two consecutive terms of the sequence.

In the remaining case $\sigma$ lies in no plane $\{\ell=c_i\}$, since that would be the first case,
and no $c_i$ lies in $(\min_\sigma\ell,\max_\sigma\ell)$; so $\sigma$ belongs to no $\mathcal{M}(c_i)$ and to no $\mathcal{X}(c_i)$.
As the sequence is strictly increasing and unbounded above and below, there is exactly one $j$ with $c_{j-1}\le\min_\sigma\ell<c_j$.
Then $\max_\sigma\ell\le c_j$, because $c_j$ does not lie in $(\min_\sigma\ell,\max_\sigma\ell)$.
So $\sigma\subset\{c_{j-1}\le\ell\le c_j\}$, and $\sigma$ lies in neither of the two planes: $\sigma\in\mathcal{B}(c_{j-1},c_j)$.
If also $\sigma\in\mathcal{B}(c_{i-1},c_i)$, then $c_{i-1}\le\min_\sigma\ell$ and $\max_\sigma\ell\le c_i$;
moreover $\min_\sigma\ell<c_i$, since otherwise $\sigma\subset\{\ell=c_i\}$.
Hence $c_{i-1}\le\min_\sigma\ell<c_i$, and $i=j$.

For the second statement let $c',c''\in\Adm(\ell)$ with $c<c'<c''$.
The elements of $\Adm(\ell)$ outside the open interval $(c,c'')$, in increasing order, form a sequence as above,
in which $c$ and $c''$ are consecutive terms;
inserting $c'$ gives a second such sequence, in which $c$, $c'$, $c''$ are consecutive terms.
By the first statement each of the two sequences partitions the set of the $k$-cells.
The two partitions consist of the same sets, except that the set $\mathcal{B}^k(c,c'')$ of the first
is replaced in the second by the four sets $\mathcal{B}^k(c,c')$, $\mathcal{M}^k(c')$, $\mathcal{X}^k(c')$, $\mathcal{B}^k(c',c'')$.
Hence $\mathcal{B}^k(c,c'')$ is the union of these four sets, and they are pairwise disjoint.

\eqref{it:upper-faces}
Let $\sigma'$ be an upper face of $\sigma$, and let $c'\in\Adm(\ell)$, $c'>c$.
By definition $\sigma'\subset\{c\le\ell\le c^+\}$ and $\sigma'\not\subset\{\ell=c\}$;
and $c'\ge c^+$, as $c^+$ is the least element of $\Adm(\ell)$ larger than $c$.
So $\sigma'\subset\{c\le\ell\le c'\}$.
If $\sigma'\subset\{\ell=c'\}$, then $\sigma'\in\mathcal{M}(c')$;
otherwise $\sigma'$ lies in neither of the planes $\{\ell=c\}$, $\{\ell=c'\}$, and $\sigma'\in\mathcal{B}(c,c')$.
Let next $\sigma'$ be a lower face and $c'\in\Adm(\ell)$, $c'<c$.
Then $\sigma'\subset\{c^-\le\ell\le c\}$, $\sigma'\not\subset\{\ell=c\}$, and $c'\le c^-$, so that $\sigma'\subset\{c'\le\ell\le c\}$;
the face $\sigma'$ belongs to $\mathcal{M}(c')$ if it lies in $\{\ell=c'\}$, and to $\mathcal{B}(c',c)$ otherwise.
\end{proof}

\begin{lemma}[The coboundary at a mirror]\label{lem:coboundary-mirror}
Let $\ell\in\mathcal{L}$, let $\Xi=(\Gamma_\parallel,\chi_\parallel)$ be a cross-section for $\ell$, and let $c\in\Adm(\ell)$.
For $\sigma\in\mathcal{X}^2(c)\cup\mathcal{X}^3(c)$ and a real cochain $u$ on the faces of $\sigma$ put
\[
(d^{+}_cu)(\sigma):=\sum_{\substack{\sigma'\ \text{upper face}\\ \text{of}\ \sigma}}[\sigma:\sigma']\cdot u(\sigma'),
\qquad
(d_cu)(\sigma):=\sum_{\substack{\sigma'\ \text{face of}\ \sigma,\\ \text{cut by}\ c}}[\sigma:\sigma']\cdot u(\sigma'),
\]
the upper faces being those of Lemma~\ref{lem:cells-mirrors}\eqref{it:cut-faces}.
The sets $\mathcal{X}^1(c)$, $\mathcal{X}^2(c)$, $\mathcal{X}^3(c)$ are finite, since their cells have heights in $[c^-,c^+]$
and only finitely many cells of $\Z^3/\Gamma_\parallel$ have heights in a bounded interval, $\Gamma_\parallel$ being of rank two in the plane $\{\ell=0\}$.
The second formula defines linear maps
\[
d_c\colon\R^{\mathcal{X}^1(c)}\to\R^{\mathcal{X}^2(c)},\qquad d_c\colon\R^{\mathcal{X}^2(c)}\to\R^{\mathcal{X}^3(c)},
\]
and $d_c^{*}$ denotes the adjoint of $d_c$ for the counting inner products $\sum_{\sigma}u(\sigma)\cdot u'(\sigma)$ on these spaces.
\begin{enumerate}
\item\label{it:coboundary-split}
For $\sigma\in\mathcal{X}^2(c)\cup\mathcal{X}^3(c)$ and every real cochain $u$ on the faces of $\sigma$,
\[
(du)(\sigma)=(d^{+}_cu)(\sigma)-(d^{+}_c\theta_c^{*}u)(\sigma)+(d_cu)(\sigma).
\]
\item\label{it:cut-complex}
The composite $\R^{\mathcal{X}^1(c)}\to\R^{\mathcal{X}^2(c)}\to\R^{\mathcal{X}^3(c)}$ of the two maps $d_c$ vanishes: $d_cd_c=0$.
\item\label{it:cut-identity}
Let $k$ be a one-cochain with real values, for example with integer values,
such that $k(b)=0$ for every link $b$ that is cut by $c$ or contained in $\{\ell\le c\}$.
Every cube $\sigma\in\mathcal{X}^3(c)$ satisfies
\[
(d^{+}_c\,dk)(\sigma)+\sum_{\sigma'\in\mathcal{X}^2(c)}[\sigma:\sigma']\cdot(d^{+}_ck)(\sigma')=0 .
\]
\end{enumerate}
\end{lemma}

\begin{proof}
Write $\theta:=\theta_c$.

\eqref{it:coboundary-split}
By Lemma~\ref{lem:cells-mirrors}\eqref{it:cut-faces}, every face of $\sigma$ is cut by $c$, upper or lower,
and the sum $(du)(\sigma)=\sum_{\sigma'}[\sigma:\sigma']\cdot u(\sigma')$ splits accordingly into three sums.
The sum over the upper faces is $(d^{+}_cu)(\sigma)$, and the sum over the faces cut by $c$ is $(d_cu)(\sigma)$.
The lower faces are the cells $\theta\sigma''$ with $\sigma''$ an upper face, each of them obtained once;
by the identity of Lemma~\ref{lem:cells-mirrors}\eqref{it:cut-faces},
\[
\sum_{\sigma'\ \text{lower}}[\sigma:\sigma']\cdot u(\sigma')
=\sum_{\sigma''\ \text{upper}}[\sigma:\theta\sigma'']\cdot u(\theta\sigma'')
=-\sum_{\sigma''\ \text{upper}}[\sigma:\sigma'']\cdot(\theta^{*}u)(\sigma'')
=-(d^{+}_c\theta^{*}u)(\sigma).
\]

\eqref{it:cut-complex}
Let $\sigma\in\mathcal{X}^3(c)$ and $u\in\R^{\mathcal{X}^1(c)}$.
As the incidence number of two cells vanishes unless the second is a face of the first,
\[
(d_cd_cu)(\sigma)=\sum_{b\in\mathcal{X}^1(c)}\left(\sum_{\sigma'\in\mathcal{X}^2(c)}[\sigma:\sigma']\cdot[\sigma':b]\right)\cdot u(b).
\]
If $[\sigma':b]\neq0$, the plaquette $\sigma'$ contains the link $b$, and
$\min_{\sigma'}\ell\le\min_b\ell<c<\max_b\ell\le\max_{\sigma'}\ell$; so $\sigma'$ is cut by $c$.
The inner sum is therefore the sum over all plaquettes $\sigma'$, and it vanishes by \eqref{eq:slab-incidence-square}.

\eqref{it:cut-identity}
Let $\sigma\in\mathcal{X}^3(c)$.
By \eqref{eq:slab-incidence-square}, $(ddk)(\sigma)=\sum_{\sigma'}[\sigma:\sigma']\cdot(dk)(\sigma')$ vanishes,
and this sum splits into the sums over the upper, the lower and the cut faces $\sigma'$ of $\sigma$.
The first of them is $(d^{+}_c\,dk)(\sigma)$.
A lower face lies in $\{\ell\le c\}$, and so do its links; $k$ vanishes on them, and $(dk)(\sigma')=0$.
A face $\sigma'$ cut by $c$ is a plaquette of $\mathcal{X}^2(c)$, and Lemma~\ref{lem:cells-mirrors}\eqref{it:cut-faces}, applied to $\sigma'$,
shows that its links are cut by $c$, upper or lower;
$k$ vanishes on the cut links and on the lower ones, which lie in $\{\ell\le c\}$, so that $(dk)(\sigma')=(d^{+}_ck)(\sigma')$.
Hence
\[
0=(ddk)(\sigma)=(d^{+}_c\,dk)(\sigma)+\sum_{\sigma'\in\mathcal{X}^2(c)}[\sigma:\sigma']\cdot(d^{+}_ck)(\sigma').
\]
\end{proof}

At the plaquettes and the cubes cut by $c$, part \eqref{it:coboundary-split} of Lemma~\ref{lem:coboundary-mirror}
applies to the two components of a configuration $(A,z)$.
For a real one-cochain $A$ and an integer two-cochain $z$, given on the faces of the cells concerned, it reads
\[
\begin{aligned}
(dA)(\sigma)&=(d^{+}_cA)(\sigma)-(d^{+}_c\theta_c^{*}A)(\sigma)+(d_cA)(\sigma), &\qquad &\sigma\in\mathcal{X}^2(c),\\
(dz)(\sigma)&=(d^{+}_cz)(\sigma)-(d^{+}_c\theta_c^{*}z)(\sigma)+(d_cz)(\sigma), &\qquad &\sigma\in\mathcal{X}^3(c).
\end{aligned}
\]

\begin{example}[Four kinds of mirrors]\label{ex:four-mirrors}
For the axis family $x_1$ and the diagonal families $\ell_-=x_1-x_2$, $\ell_+=x_1+x_2$ (Definition~\ref{def:mirrors}),
the cells in a mirror, the cells cut by it and the maps $d^{+}_c$, $d_c$ of Lemma~\ref{lem:coboundary-mirror} are as follows.
The incidence numbers are read off from the coboundary of Definition~\ref{def:cochains}:
for the plaquette $(x;\mu\nu)$ and the cube $(x)=x+[0,1]^3$,
\[
\begin{aligned}
(dA)_{\mu\nu}(x)&=A_\mu(x)+A_\nu(x+e_\mu)-A_\mu(x+e_\nu)-A_\nu(x),\\
(dz)(x)&=z_{23}(x+e_1)-z_{23}(x)-z_{13}(x+e_2)+z_{13}(x)+z_{12}(x+e_3)-z_{12}(x).
\end{aligned}
\]
For $\ell=x_1$ the lattice $\Gamma_\parallel$ lies in $\{0\}\times\Z^2$ and is regarded as a subgroup of the slice lattice $\Z^2$ of Definition~\ref{def:slices}.
The sites $y$, the links $b=(y,\nu)$, $\nu\in\{2,3\}$, and the plaquettes $\sigma$ of the slice are taken modulo $\Gamma_\parallel$,
and a plaquette $\sigma$ of the slice is identified with its lowest corner $y$.
For $j\in\Z$, let $(j;b)$ be the link $((j,y),\nu)$ and $(j;\sigma)$ the plaquette $((j,y);23)$ of the plane $\{x_1=j\}$,
and let $(j,\sigma)$ be the cube $(j,y)+[0,1]^3$ between the planes $\{x_1=j\}$ and $\{x_1=j+1\}$;
between the same two planes, the plaquette over the link $b$ is $((j,y);1\nu)$.
In the notation of Definition~\ref{def:slices}, $A_j(b)$ is the value of $A$ at the link $(j;b)$,
and $a_j(y)$ is its value at the link from $(j,y)$ to $(j+1,y)$.
For an integer two-cochain $z$ put, as $F^{\mathrm{el}}_j$ and $F^{\mathrm{mag}}_j$ are formed from $F$ in Definition~\ref{def:slices},
\[
z^{\mathrm{el}}_j(y,\nu):=z_{1\nu}(j,y),\qquad z^{\mathrm{mag}}_j(\sigma):=z_{23}(j,y);
\]
these are the values of $z$ at the plaquettes $((j,y);1\nu)$ and $(j;\sigma)$.
\begin{itemize}
\item \emph{Site planes $\{x_1=j\}$, $j\in\Z$.}
The links and the plaquettes contained in the mirror are the links $(j;b)$ and the plaquettes $(j;\sigma)$.
No cell is cut: by Lemma~\ref{lem:cells-mirrors}\eqref{it:cut-heights}, the least height of a cell cut by $j$ would be $j-\frac12$, which is not an integer.
Thus $\mathcal{X}^k(j)$ is empty for $k=0,1,2,3$.
\item \emph{Link planes $\{x_1=j+\frac12\}$, $j\in\Z$.}
Put $c:=j+\frac12$.
No cell lies in the mirror, since the corners of a cell have integer coordinates.
By Lemma~\ref{lem:cells-mirrors}\eqref{it:cut-heights}, the cells cut by $c$ are the cells that have the direction $1$ and the least height $j$:
the links from $(j,y)$ to $(j+1,y)$, the plaquettes $((j,y);1\nu)$ and the cubes $(j,\sigma)$.
The only upper face of $((j,y);1\nu)$ is the link $(j+1;b)$, and the only upper face of $(j,\sigma)$ is the plaquette $(j+1;\sigma)$,
both with the incidence number $1$; so
\[
(d^{+}_cA)((j,y);1\nu)=A_{j+1}(b),\qquad (d^{+}_cz)(j,\sigma)=z^{\mathrm{mag}}_{j+1}(\sigma).
\]
Identify $\R^{\mathcal{X}^1(c)}$, $\R^{\mathcal{X}^2(c)}$ and $\R^{\mathcal{X}^3(c)}$ with the spaces of functions on the sites, the links and the plaquettes of the slice,
the link from $(j,y)$ to $(j+1,y)$ corresponding to $y$, the plaquette $((j,y);1\nu)$ to $b$, and the cube $(j,\sigma)$ to $\sigma$.
Then $d_c=-d_{\mathrm{s}}$ in both degrees, and the coboundary formulas above give, for every link $b=(y,\nu)$ and every plaquette $\sigma$ of the slice,
\[
\begin{aligned}
(dA)_{1\nu}(j,y)&=A_{j+1}(b)-A_j(b)-(d_{\mathrm{s}}a_j)(b),\\
(dz)(j,\sigma)&=z^{\mathrm{mag}}_{j+1}(\sigma)-z^{\mathrm{mag}}_j(\sigma)-(d_{\mathrm{s}}z^{\mathrm{el}}_j)(\sigma).
\end{aligned}
\]
This agrees with Lemma~\ref{lem:coboundary-mirror}\eqref{it:coboundary-split}:
$\theta_c$ maps $(j+1;b)$ onto $(j;b)$ and $(j+1;\sigma)$ onto $(j;\sigma)$, preserving the orientations,
so that $(d^{+}_c\theta_c^{*}A)((j,y);1\nu)=A_j(b)$ and $(d^{+}_c\theta_c^{*}z)(j,\sigma)=z^{\mathrm{mag}}_j(\sigma)$.
In particular $(dz)(j,\sigma)=0$ if and only if $z^{\mathrm{mag}}_{j+1}(\sigma)-z^{\mathrm{mag}}_j(\sigma)=(d_{\mathrm{s}}z^{\mathrm{el}}_j)(\sigma)$.
\item \emph{Diagonal mirrors $\{\ell_-=c\}$, $c\in\Z$.}
The links in the mirror are the links $(x,3)$ with $\ell_-(x)=c$, and no plaquette lies in it.
No link is cut, as $n_{\ell_-}=e_1-e_2$ lies in the direction space of no link;
the cut plaquettes are the plaquettes $(x;12)$ with $\ell_-(x)=c$, and the cut cubes are the cubes $(x)$ with $\ell_-(x)=c$.
The corners $x$, $x+e_1$, $x+e_1+e_2$, $x+e_2$ of $(x;12)$ have the heights $c$, $c+1$, $c$, $c-1$,
and the upper faces of $(x;12)$ are the links $(x,1)$ and $(x+e_1,2)$.
In the cube $(x)$ the plaquettes $(x+e_1;23)$ and $(x;13)$ are upper faces, $(x;23)$ and $(x+e_2;13)$ are lower faces,
and $(x;12)$ and $(x+e_3;12)$ are cut by $c$.
Hence
\[
(d^{+}_cA)(x;12)=A_1(x)+A_2(x+e_1),\qquad (d^{+}_cz)(x)=z_{23}(x+e_1)+z_{13}(x),
\]
and, if the cut plaquette $(x;12)$ and the cut cube $(x)$ are both identified with the corner $x$ that names them,
the corner with the least coordinates, then $(d_cu)(x)=u(x+e_3)-u(x)$ for $u\in\R^{\mathcal{X}^2(c)}$.
\item \emph{Diagonal mirrors $\{\ell_+=c\}$, $c\in\Z$.}
The links in the mirror are the links $(x,3)$ with $\ell_+(x)=c$, and no plaquette lies in it.
No link is cut, as $n_{\ell_+}=e_1+e_2$ lies in the direction space of no link;
the cut plaquettes are the plaquettes $(x;12)$ with $\ell_+(x)=c-1$, and the cut cubes are the cubes $(x)$ with $\ell_+(x)=c-1$.
The corners $x$, $x+e_1$, $x+e_1+e_2$, $x+e_2$ of $(x;12)$ have the heights $c-1$, $c$, $c+1$, $c$,
and the upper faces of $(x;12)$ are the links $(x+e_2,1)$ and $(x+e_1,2)$.
In the cube $(x)$ the plaquettes $(x+e_1;23)$ and $(x+e_2;13)$ are upper faces, $(x;23)$ and $(x;13)$ are lower faces,
and $(x;12)$ and $(x+e_3;12)$ are cut by $c$.
Hence
\[
(d^{+}_cA)(x;12)=A_2(x+e_1)-A_1(x+e_2),\qquad (d^{+}_cz)(x)=z_{23}(x+e_1)-z_{13}(x+e_2),
\]
and, with the same identification, $(d_cu)(x)=u(x+e_3)-u(x)$ for $u\in\R^{\mathcal{X}^2(c)}$.
\end{itemize}
For the other families the definitions of $d^{+}_c$ and $d_c$ give formulas of the same kind, with signs that depend on the family.
For instance, for the link plane $\{x_2=j+\frac12\}$ and a cube $(x)$ with $x_2=j$, one has $(d^{+}_cz)(x)=-z_{13}(x+e_2)$.
\end{example}

\subsection{Integrands on slabs and tori}\label{sec:slab-integrands}

Throughout this subsection $\ell\in\mathcal{L}$ is a fixed family and $\Xi=(\Gamma_\parallel,\chi_\parallel)$ a fixed cross-section for $\ell$ (Definition~\ref{def:slabs}).

\begin{definition}[Grassmann algebras of slabs]\label{def:slab-grassmann}
Let $\mathbb{I}\subset\R$ be an open interval.
The fermion sites with heights in $\mathbb{I}$, and their classes in the cylinder $\Z^3/\Gamma_\parallel$, are
\[
\mathbb{X}(\mathbb{I}):=\{x\in\mathbb{X}:\ \ell(x)\in\mathbb{I}\},
\qquad
\mathbb{X}_\Xi(\mathbb{I}):=\mathbb{X}(\mathbb{I})/\Gamma_\parallel .
\]
The algebra $\Gr_\Xi(\mathbb{I})$ is the complex Grassmann algebra generated by the elements $\psi_\alpha(x)$, $\bar\psi_\alpha(x)$ with $x\in\mathbb{X}(\mathbb{I})$ and $1\le\alpha\le4$, subject to the relations
\[
\psi(x+g)=\chi_\parallel(g)\cdot\psi(x),
\qquad
\bar\psi(x+g)=\chi_\parallel(g)\cdot\bar\psi(x)
\qquad
(x\in\mathbb{X}(\mathbb{I}),\ g\in\Gamma_\parallel).
\]
The generators at the points of a set of representatives of $\mathbb{X}_\Xi(\mathbb{I})$ are independent.
The products of distinct ones among them, taken in a fixed order, form together with $1$ a basis of $\Gr_\Xi(\mathbb{I})$, called a \emph{monomial basis}.
\end{definition}

Since $\ell(g)=0$ for $g\in\Gamma_\parallel$, the set $\mathbb{X}(\mathbb{I})$ is a union of orbits of $\Gamma_\parallel$.
The relations are consistent because $\chi_\parallel$ is a character with values $\pm1$, and every generator is $\pm1$ times the generator of the same name at the chosen representative of its class.
Another set of representatives, or another order, multiplies each element of a monomial basis by $\pm1$; consequently the moduli of the coefficients of an element of $\Gr_\Xi(\mathbb{I})$ in a monomial basis do not depend on the choice of the basis.
If $\mathbb{I}$ is bounded, then $\mathbb{X}_\Xi(\mathbb{I})$ is finite, because $\Gamma_\parallel$ has rank two in the plane $\{\ell=0\}$, and $\Gr_\Xi(\mathbb{I})$ has finite dimension.
For open intervals $\mathbb{I}\subset\mathbb{I}'$ the algebra $\Gr_\Xi(\mathbb{I})$ is identified with the subalgebra of $\Gr_\Xi(\mathbb{I}')$ generated by the fields at the points of $\mathbb{X}(\mathbb{I})$.

For $v\in\Z^3$ the translation $x\mapsto x+v$ maps $\mathbb{X}=Wv_0$ onto itself, since $W$ contains the translations by $\Z^3$ (Definitions~\ref{def:mirrors} and~\ref{def:fermion-sites}).
The automorphism of $\Gr(\mathbb{X})$ (Definition~\ref{def:grassmann}) with
\[
\tau_v\psi(x)=\psi(x+v),
\qquad
\tau_v\bar\psi(x)=\bar\psi(x+v)
\qquad
(x\in\mathbb{X})
\]
is again denoted $\tau_v$.
For $c\in\Adm(\ell)$, $\Theta_c$ is the antilinear map of $\Gr(\mathbb{X})$ that Definition~\ref{def:grassmann-reflections} attaches to the mirror $(\ell,c)$: it reverses the order of products, $\Theta_c^2=1$, and by \eqref{eq:setting-reflection-fields}
\[
\Theta_c\psi(x)=-\bar\psi(\theta_cx)\,\gamma_n,
\qquad
\Theta_c\bar\psi(x)=-\gamma_n\,\psi(\theta_cx),
\]
where $n=\hat n_\ell$ is the upward unit normal of the family $\ell$.

\begin{lemma}[Translations and reflections]\label{lem:slab-automorphisms}
Let $\mathbb{I}$ be an open interval, $v,v'\in\Z^3$ and $c,c'\in\Adm(\ell)$.
The automorphism $\tau_v$ maps $\Gr(\mathbb{X}(\mathbb{I}))$ onto $\Gr(\mathbb{X}(\mathbb{I}+\ell(v)))$ and induces an isomorphism of algebras $\Gr_\Xi(\mathbb{I})\to\Gr_\Xi(\mathbb{I}+\ell(v))$.
The map $\Theta_c$ maps $\Gr(\mathbb{X}(\mathbb{I}))$ onto $\Gr(\mathbb{X}(2c-\mathbb{I}))$ and induces an antilinear map $\Gr_\Xi(\mathbb{I})\to\Gr_\Xi(2c-\mathbb{I})$ with $\Theta_c(GG')=\Theta_c(G')\cdot\Theta_c(G)$.
The induced maps are again denoted $\tau_v$ and $\Theta_c$, and the following identities hold on $\Gr(\mathbb{X})$ and on the algebras $\Gr_\Xi(\mathbb{I})$.
\begin{enumerate}
\item\label{it:translation-law}
$\tau_v\tau_{v'}=\tau_{v+v'}$.
If $v\in\Gamma_\parallel$, then $\tau_vG=\chi_\parallel(v)^{\deg G}\cdot G$ for every homogeneous $G\in\Gr_\Xi(\mathbb{I})$.
\item\label{it:translated-reflection}
$\tau_v\Theta_c\tau_v^{-1}=\Theta_{c+\ell(v)}$.
\item\label{it:parallel-reflections}
$\Theta_c^2=1$, and $\Theta_{c'}\Theta_c=\tau_{v_{c,c'}}$, where
\[
v_{c,c'}:=\frac{2(c'-c)}{|n_\ell|^2}\,n_\ell
\]
is an integer vector of height $\ell(v_{c,c'})=2(c'-c)$.
With $c'':=2c'-c$ one has $\Theta_{c''}\tau_{v_{c,c'}}=\Theta_{c'}$.
\end{enumerate}
\end{lemma}

\begin{proof}
Translation by $v$ raises every height by $\ell(v)$.
So $\tau_v$ maps $\Gr(\mathbb{X}(\mathbb{I}))$ onto $\Gr(\mathbb{X}(\mathbb{I}+\ell(v)))$,
and it carries the relations $\psi(x+g)-\chi_\parallel(g)\cdot\psi(x)$ and $\bar\psi(x+g)-\chi_\parallel(g)\cdot\bar\psi(x)$ of $\Gr_\Xi(\mathbb{I})$ to the relations of $\Gr_\Xi(\mathbb{I}+\ell(v))$ at the point $x+v$.
Hence $\tau_v$ induces a homomorphism $\Gr_\Xi(\mathbb{I})\to\Gr_\Xi(\mathbb{I}+\ell(v))$, whose inverse is induced by $\tau_{-v}$.
By Definition~\ref{def:mirrors}\eqref{it:mirror-height}, $\ell(\theta_cx)=2c-\ell(x)$; so $\Theta_c$ maps $\Gr(\mathbb{X}(\mathbb{I}))$ onto $\Gr(\mathbb{X}(2c-\mathbb{I}))$.
For $g\in\Gamma_\parallel$ one has $\ell(g)=0$, hence $\theta_c(x+g)=\theta_cx+g$, and $\chi_\parallel(g)$ is real.
Therefore
\[
\Theta_c\left(\psi(x+g)-\chi_\parallel(g)\cdot\psi(x)\right)
=-\left(\bar\psi(\theta_cx+g)-\chi_\parallel(g)\cdot\bar\psi(\theta_cx)\right)\gamma_n,
\]
\[
\Theta_c\left(\bar\psi(x+g)-\chi_\parallel(g)\cdot\bar\psi(x)\right)
=-\gamma_n\left(\psi(\theta_cx+g)-\chi_\parallel(g)\cdot\psi(\theta_cx)\right),
\]
and the right-hand sides are linear combinations of relations of $\Gr_\Xi(2c-\mathbb{I})$.
An antilinear map that reverses the order of products sends a product of three factors to the product of their images in the opposite order; hence it carries the two-sided ideal generated by a set into the two-sided ideal generated by the image of that set.
Consequently $\Theta_c$ induces an antilinear map $\Gr_\Xi(\mathbb{I})\to\Gr_\Xi(2c-\mathbb{I})$ which reverses the order of products.

The maps on the algebras $\Gr_\Xi(\mathbb{I})$ are induced, through the surjective quotient maps, by the maps with the same names on the subalgebras $\Gr(\mathbb{X}(\mathbb{I}))$ of $\Gr(\mathbb{X})$.
Every identity between maps on $\Gr(\mathbb{X})$ therefore holds between the induced maps.
On $\Gr(\mathbb{X})$, two homomorphisms that agree on the generators coincide, and so do two antilinear maps that reverse the order of products and agree on the generators.

\eqref{it:translation-law}
Both $\tau_v\tau_{v'}$ and $\tau_{v+v'}$ send $\psi(x)$ to $\psi(x+v+v')$ and $\bar\psi(x)$ to $\bar\psi(x+v+v')$.
Let $v\in\Gamma_\parallel$.
In $\Gr_\Xi(\mathbb{I})$ the relations give $\tau_v\psi(x)=\psi(x+v)=\chi_\parallel(v)\cdot\psi(x)$ and, in the same way, $\tau_v\bar\psi(x)=\chi_\parallel(v)\cdot\bar\psi(x)$.
Thus $\tau_v$ multiplies a product of $k$ generators by $\chi_\parallel(v)^k$, a number that depends only on the parity of $k$.

\eqref{it:translated-reflection}
By \eqref{it:translation-law}, $\tau_v^{-1}=\tau_{-v}$, and $c+\ell(v)\in\Adm(\ell)$ because $\ell(v)\in\Z$.
The map $\tau_v\Theta_c\tau_{-v}$ is antilinear and reverses the order of products, and so does $\Theta_{c+\ell(v)}$.
Both mirrors belong to the family $\ell$ and have the upward unit normal $n$.
Since $\theta_c(x-v)+v=\theta_{c+\ell(v)}x$ by Definition~\ref{def:mirrors}\eqref{it:mirror-conjugation},
\[
\tau_v\Theta_c\tau_{-v}\psi(x)
=\tau_v\left(-\bar\psi(\theta_c(x-v))\,\gamma_n\right)
=-\bar\psi\left(\theta_c(x-v)+v\right)\gamma_n
=-\bar\psi(\theta_{c+\ell(v)}x)\,\gamma_n
=\Theta_{c+\ell(v)}\psi(x),
\]
and likewise $\tau_v\Theta_c\tau_{-v}\bar\psi(x)=\tau_v\left(-\gamma_n\psi(\theta_c(x-v))\right)=-\gamma_n\psi(\theta_{c+\ell(v)}x)=\Theta_{c+\ell(v)}\bar\psi(x)$.

\eqref{it:parallel-reflections}
The identity $\Theta_c^2=1$ is part of Definition~\ref{def:grassmann-reflections}.
The composition $\Theta_{c'}\Theta_c$ is linear and multiplicative.
The mirrors $(\ell,c')$ and $(\ell,c)$ have the same normal $n$, and $\gamma_n\gamma_n=1$;
by \eqref{eq:setting-two-reflections} the map $\Theta_{c'}\Theta_c$ therefore sends $\psi(x)$ to $\psi(\theta_{c'}\theta_cx)$ and $\bar\psi(x)$ to $\bar\psi(\theta_{c'}\theta_cx)$.
By Definition~\ref{def:mirrors}\eqref{it:mirror-product}, $\theta_{c'}\theta_c$ is the translation by $v_{c,c'}$.
Hence $\Theta_{c'}\Theta_c=\tau_{v_{c,c'}}$.
Since $c'-c\in\left(|n_\ell|^2/2\right)\Z$, the vector $v_{c,c'}$ is an integer multiple of $n_\ell$, and $\ell(v_{c,c'})=v_{c,c'}\cdot n_\ell=2(c'-c)$.
Write $v=v_{c,c'}$; then $c''=c+\ell(v)$, and $\tau_v^{-1}=(\Theta_{c'}\Theta_c)^{-1}=\Theta_c\Theta_{c'}$ because $\Theta_c^2=\Theta_{c'}^2=1$.
By \eqref{it:translated-reflection},
\[
\Theta_{c''}=\tau_v\Theta_c\tau_v^{-1}=(\Theta_{c'}\Theta_c)\,\Theta_c\,(\Theta_c\Theta_{c'})=\Theta_{c'}\Theta_c\Theta_{c'},
\]
and multiplying on the right by $\tau_v=\Theta_{c'}\Theta_c$ yields
\[
\Theta_{c''}\tau_v=(\Theta_{c'}\Theta_c\Theta_{c'})(\Theta_{c'}\Theta_c)=\Theta_{c'}.
\]
\end{proof}

\begin{definition}[Slab integrands]\label{def:slab-integrands}
Let $c<c'$ be elements of $\Adm(\ell)$,
and let $\mathfrak{C}[c,c']$ be the space of configurations of the closed slab between the mirrors at $c$ and $c'$, with the integer shifts $\varpi\mapsto\varpi+[k]$ and the curvature $F(\varpi)$ of Definition~\ref{def:slab-configurations}.
The class $\mathcal{I}_\Xi[c,c']$ consists of the Borel maps $G:\mathfrak{C}[c,c']\to\Gr_\Xi((c,c'))$ with the following two properties.
\begin{enumerate}
\item\label{it:integrand-shifts}
$G(\varpi+[k])=G(\varpi)$ for all $\varpi\in\mathfrak{C}[c,c']$ and all $k\in\Z^{[c,c']^{(1)}}$.
\item\label{it:integrand-growth}
There are $C<\infty$ and $\varepsilon_G\in\left(0,\tfrac12\right]$ such that, for every $\varpi\in\mathfrak{C}[c,c']$, each coefficient of $G(\varpi)$ in a monomial basis has modulus at most
\[
C\cdot\exp\left(\left(\tfrac12-\varepsilon_G\right)\sum_{\sigma\in\mathcal{B}^2(c,c')}F_\sigma(\varpi)^2\right).
\]
\end{enumerate}
The class $\mathcal{I}_\Xi[c,c']$ is a complex vector space, graded by the parity of the values.
Let $c<c'<c''$ be elements of $\Adm(\ell)$, let $G\in\mathcal{I}_\Xi[c,c']$, $G'\in\mathcal{I}_\Xi[c',c'']$ and $v\in\Z^3$.
The \emph{product} $GG'\in\mathcal{I}_\Xi[c,c'']$ is
\[
(GG')(\varpi):=G(\varpi|_{[c,c']})\cdot G'(\varpi|_{[c',c'']}),
\qquad
\varpi\in\mathfrak{C}[c,c''],
\]
where $\varpi|_{[c,c']}$ denotes the restriction of $\varpi$ to $[c,c']^{(1)}$ and $[c,c']^{(2)}$.
The \emph{translate} $\tau_vG\in\mathcal{I}_\Xi[c+\ell(v),c'+\ell(v)]$ is
\[
(\tau_vG)(\varpi'):=\tau_v\left(G(\tau_v^*\varpi')\right),
\qquad
\varpi'\in\mathfrak{C}[c+\ell(v),c'+\ell(v)].
\]
The \emph{reflection} $\Theta_{c'}G\in\mathcal{I}_\Xi[c',2c'-c]$ in the upper mirror is
\[
(\Theta_{c'}G)(\varpi'):=\Theta_{c'}\left(G(\theta_{c'}^*\varpi')\right),
\qquad
\varpi'\in\mathfrak{C}[c',2c'-c].
\]
On the right-hand sides, $\tau_v$ and $\Theta_{c'}$ are the maps of Lemma~\ref{lem:slab-automorphisms}, and $\tau_v^*\varpi'$, $\theta_{c'}^*\varpi'$ are the pull-backs of configurations (Definition~\ref{def:slab-configurations}).
If $G(\varpi)=\sum_\Psi a_\Psi(\varpi)\cdot\Psi$ in a monomial basis, then
\[
(\Theta_{c'}G)(\varpi')=\sum_\Psi\overline{a_\Psi(\theta_{c'}^*\varpi')}\cdot\Theta_{c'}(\Psi),
\]
which is the rule \eqref{eq:setting-reflection-coefficients}, read on configurations of slabs.
\end{definition}

\begin{remark}\label{rem:slab-growth}
Consider majorants $C\cdot\exp\left(\left(\tfrac12-\varepsilon\right)\sum_\sigma F_\sigma^2\right)$, the sum being over a finite set of plaquettes.
For $0<\varepsilon\le\tfrac12$ one has $\tfrac12-\varepsilon\ge0$, so the majorant does not decrease if the set of plaquettes is enlarged, or if $\varepsilon$ is replaced by a smaller positive number.
For $\varepsilon>\tfrac12$ the majorant is at most $C$, which is its value for $\varepsilon=\tfrac12$; so the range $\left(0,\tfrac12\right]$ for $\varepsilon_G$ in Definition~\ref{def:slab-integrands}\eqref{it:integrand-growth} is no restriction.
If the coefficients of two elements of a Grassmann algebra obey such bounds over the same set of plaquettes,
with the constants $(C_1,\varepsilon_1)$ and $(C_2,\varepsilon_2)$, where $\varepsilon_1,\varepsilon_2\in\left(0,\tfrac12\right]$,
then the coefficients of their sum obey such a bound with $(C_1+C_2,\min\{\varepsilon_1,\varepsilon_2\})$.
The coefficients of the even and of the odd part of an element are among the coefficients of the element.
Since property~\eqref{it:integrand-shifts} of Definition~\ref{def:slab-integrands} is preserved by sums and by taking the even and the odd parts of the values,
$\mathcal{I}_\Xi[c,c']$ is a complex vector space, the direct sum of the subspaces of its even and of its odd elements.
\end{remark}

The three maps of Definition~\ref{def:slab-integrands} lie in the classes indicated there, for the following reasons.

\emph{Products.}
Let $\varpi\in\mathfrak{C}[c,c'']$.
A face of a cell contained in a closed slab is contained in that slab, and $\mathcal{B}^3(c,c')\subset\mathcal{B}^3(c,c'')$ by Lemma~\ref{lem:cells-mirrors}\eqref{it:cell-partition}.
Hence $\varpi|_{[c,c']}\in\mathfrak{C}[c,c']$, the curvature $F(\varpi|_{[c,c']})$ is the restriction of $F(\varpi)$, and $(\varpi+[k])|_{[c,c']}=\varpi|_{[c,c']}+[k']$ for $k\in\Z^{[c,c'']^{(1)}}$, where $k'$ is the restriction of $k$ to $[c,c']^{(1)}$.
The same holds for $[c',c'']$, so $GG'$ is a Borel map with property \eqref{it:integrand-shifts}.
The values $G(\varpi|_{[c,c']})$ and $G'(\varpi|_{[c',c'']})$ lie in the subalgebras of $\Gr_\Xi((c,c''))$ generated by the fields at the sites with heights in $(c,c')$ and in $(c',c'')$, respectively, and these two sets of generators are disjoint.
Since the moduli of the coefficients do not depend on the monomial basis,
one may use a monomial basis of $\Gr_\Xi((c,c''))$ formed with a set of representatives that contains sets of representatives of $\mathbb{X}_\Xi((c,c'))$ and of $\mathbb{X}_\Xi((c',c''))$;
in it each coefficient of $(GG')(\varpi)$ is $0$ or $\pm$ the product of a coefficient of $G(\varpi|_{[c,c']})$ and a coefficient of $G'(\varpi|_{[c',c'']})$.
Let $(C_G,\varepsilon_G)$ and $(C_{G'},\varepsilon_{G'})$ be constants as in property~\eqref{it:integrand-growth} for $G$ and for $G'$.
The sets $\mathcal{B}^2(c,c')$ and $\mathcal{B}^2(c',c'')$ are disjoint subsets of $\mathcal{B}^2(c,c'')$ (Lemma~\ref{lem:cells-mirrors}\eqref{it:cell-partition});
by Remark~\ref{rem:slab-growth}, $GG'$ satisfies \eqref{it:integrand-growth} with the constants $C_G\cdot C_{G'}$ and $\min\{\varepsilon_G,\varepsilon_{G'}\}$.

\emph{Translates and reflections.}
The translation $\tau_v$ maps the closed slab $\{c\le\ell\le c'\}$ onto $\{c+\ell(v)\le\ell\le c'+\ell(v)\}$, and $\theta_{c'}$ maps it onto $\{c'\le\ell\le 2c'-c\}$, because $\ell\circ\theta_{c'}=2c'-\ell$ (Definition~\ref{def:mirrors}\eqref{it:mirror-height}).
Both maps commute with the translations by $\Gamma_\parallel$ and send cells of the cylinder to cells; a cell contained in one of the two bounding planes goes to a cell contained in a bounding plane, and a slab cell to a slab cell.
Their pull-backs commute with $d$ (Definition~\ref{def:cochains}) and preserve integer cochains.
Hence $\tau_v^*$ maps $\mathfrak{C}[c+\ell(v),c'+\ell(v)]$ onto $\mathfrak{C}[c,c']$, and $\theta_{c'}^*$ maps $\mathfrak{C}[c',2c'-c]$ onto $\mathfrak{C}[c,c']$.
The pull-back of $\varpi'+[k]$ is the pull-back of $\varpi'$ shifted by the pull-back of $k$, and the curvature of the pull-back of $\varpi'$ is the pull-back of $F(\varpi')$.
Since the orientation signs disappear in the squares, the sum of $F_\sigma(\tau_v^*\varpi')^2$, and likewise of $F_\sigma(\theta_{c'}^*\varpi')^2$, over $\sigma\in\mathcal{B}^2(c,c')$ equals the sum of $F_\sigma(\varpi')^2$ over the slab plaquettes of the image slab.
So $\tau_vG$ and $\Theta_{c'}G$ are Borel maps with property \eqref{it:integrand-shifts}, with values in $\Gr_\Xi((c+\ell(v),c'+\ell(v)))$ and in $\Gr_\Xi((c',2c'-c))$ by Lemma~\ref{lem:slab-automorphisms}.
The automorphism $\tau_v$ carries the monomial basis formed with a set of representatives onto the monomial basis formed with its translate by $v$.
Thus the coefficients of $(\tau_vG)(\varpi')$ are those of $G(\tau_v^*\varpi')$, and $\tau_vG$ satisfies \eqref{it:integrand-growth} with the constants of $G$.
The map $\Theta_{c'}:\Gr_\Xi((c,c'))\to\Gr_\Xi((c',2c'-c))$ is antilinear, and both spaces have finite dimension because the two intervals are bounded.
Thus the coefficients of $(\Theta_{c'}G)(\varpi')$ are bounded by a constant, depending only on $c$ and $c'$, times the largest modulus of the coefficients of $G(\theta_{c'}^*\varpi')$, and $\Theta_{c'}G$ satisfies \eqref{it:integrand-growth} with $\varepsilon_G$.

\begin{definition}[Torus integrands]\label{def:torus-integrands}
Let $\mathbb{T}=(\Gamma,\chi)$ be a torus (Definition~\ref{def:tori}), and put
\[
\mathfrak{C}(\Gamma):=\left\{(A,z)\in C^1(\Gamma)\times C^2(\Gamma;\Z):\ dz=0\right\},
\]
so that $\Conf(\Gamma)$ is the quotient of $\mathfrak{C}(\Gamma)$ by the integer shifts $(A,z)\mapsto\left(A+(2\pi/q)\,k,\ z-dk\right)$, $k\in C^1(\Gamma;\Z)$ (Definition~\ref{def:configurations}).
The class $\mathcal{I}(\mathbb{T})$ consists of the Borel maps $G:\mathfrak{C}(\Gamma)\to\Gr_{\mathbb{T}}$ with
\[
G\left(A+(2\pi/q)\,k,\ z-dk\right)=G(A,z)
\qquad
(k\in C^1(\Gamma;\Z))
\]
for which there are $C<\infty$ and $\varepsilon_G\in\left(0,\tfrac12\right]$ such that every coefficient of $G(A,z)$ in a monomial basis of $\Gr_{\mathbb{T}}$ has modulus at most $C\cdot\exp\left(\left(\tfrac12-\varepsilon_G\right)\|F\|^2\right)$.
Here $F=dA+(2\pi/q)\,z$, the number $\|F\|^2$ is the sum of $F_\sigma^2$ over the plaquettes $\sigma$ of $\Z^3/\Gamma$,
and a monomial basis of $\Gr_{\mathbb{T}}$ is formed, as in Definition~\ref{def:slab-grassmann}, with a set of representatives of $\mathbb{X}_{\mathbb{T}}$.
By the invariance, an element of $\mathcal{I}(\mathbb{T})$ is a function on $\Conf(\Gamma)$; by Remark~\ref{rem:slab-growth}, applied with the set of all plaquettes of $\Z^3/\Gamma$, the class $\mathcal{I}(\mathbb{T})$ is a complex vector space.
For $G\in\mathcal{I}(\mathbb{T})$,
\[
Z_{\mathbb{T}}[G]:=\sum_{\substack{z\in C^2(\Gamma;\Z)\\ dz=0}}\ \int_{[0,2\pi/q)^{\mathbb{T}^{(1)}}}\mathrm{d}A\ \exp\left(-\tfrac12\|F\|^2\right)\cdot\int_{\mathbb{T}}\exp\left(-S_{\mathbb{T}}(A)\right)\cdot G(A,z).
\]
This is the formula of Definition~\ref{def:functional}, with the measure of Definition~\ref{def:configurations} written out; the sum and the integral converge absolutely by Proposition~\ref{prop:berezin}\eqref{it:torus-integrability}.
An admissible product $\Phi G_1\cdots G_r$ of $\mathcal{A}(\R^3)$, with $G_j\in\mathcal{G}^{\mathrm{ev}}(B_j)$ (Definition~\ref{def:integrands}), is \emph{read on} $\mathbb{T}$ if the sets $B_j+g$, $1\le j\le r$, $g\in\Gamma$, are pairwise disjoint.
The linear span of the products that are read on $\mathbb{T}$ is denoted $\mathcal{A}(\mathbb{T})$.
\end{definition}

Thus $\mathcal{A}(\mathbb{T})$ is a subspace of $\mathcal{A}(\R^3)$.
It contains every element of $\mathcal{A}(\R^3)$ that embeds in $\mathbb{T}$.
Indeed, if a product embeds in $\Gamma$, its boxes embed disjointly in $\Gamma$, which includes that the sets $B_j+g$ ($1\le j\le r$, $g\in\Gamma$) are pairwise disjoint;
and an element that embeds in $\Gamma$ is a finite sum of products that embed (Definition~\ref{def:integrands}).

\begin{definition}[Products as maps of lifts]\label{def:product-maps}
Let $U\subset\R^3$ be open, and write $C^1$ for the space of real one-cochains on $\Z^3$ (Definition~\ref{def:cochains}).
By Definition~\ref{def:integrands}, $\mathcal{A}(U)$ is a linear subspace of $\Gr(\mathbb{X})^{C^1}$,
the algebra of the maps $C^1\to\Gr(\mathbb{X})$ under the pointwise product, and its elements are even.
The products that span it are maps of this kind:
for $r\ge0$, pairwise disjoint boxes $B_1,\dots,B_r$ contained in $U$, integrands $G_j\in\mathcal{G}^{\mathrm{ev}}(B_j)$
and a bounded Borel function $\Phi$ of finitely many $F_\sigma$ with $\sigma\subset U$,
\[
\Phi\cdot G_1\cdots G_r\colon\ C^1\to\Gr(\mathbb{X}),\qquad
\hat A\mapsto\Phi(d\hat A)\cdot G_1\left(\hat A|_{B_1^{(1)}}\right)\cdots G_r\left(\hat A|_{B_r^{(1)}}\right).
\]
An element $G$ of $\mathcal{A}(U)$ is thus a map $\hat A\mapsto G(\hat A)$ on $C^1$.
In general $G$ is a finite sum of products in more than one way, and different such sums may involve different boxes.
The reflections \eqref{eq:setting-reflection-coefficients} are maps of $\Gr(\mathbb{X})^{C^1}$ into itself,
and the action \eqref{eq:setting-group-action} of $W$ is defined on its even elements;
both apply to $G$.
For a mirror $\Pi$, an element $G$ of $\Gr(\mathbb{X})^{C^1}$ and $\hat A\in C^1$, the reflected map $\Theta_\Pi G$ takes at $\hat A$ the value
\[
(\Theta_\Pi G)(\hat A)=\Theta_\Pi\left(G(\theta_\Pi^{*}\hat A)\right),
\]
where $\Theta_\Pi$ on the right is the antilinear reflection of Grassmann variables of Definition~\ref{def:grassmann-reflections},
which acts on the generators by \eqref{eq:setting-reflection-fields}.
Indeed, once $G(\theta_\Pi^{*}\hat A)$ is expanded in a monomial basis, the antilinearity of $\Theta_\Pi$ makes the right side equal to the sum in \eqref{eq:setting-reflection-coefficients},
in which the coefficients of $G$ at $\theta_\Pi^{*}\hat A$ appear complex conjugated.

Let further $\ell\in\mathcal{L}$, let $\Xi=(\Gamma_\parallel,\chi_\parallel)$ be a cross-section for $\ell$ (Definition~\ref{def:slabs}),
and let $\mathbb{I}\subset\R$ be an open interval.
The \emph{quotient homomorphism}
\[
r_\Xi\colon\Gr\left(\mathbb{X}(\mathbb{I})\right)\to\Gr_\Xi(\mathbb{I})
\]
is the homomorphism of unital algebras under which every generator $\psi_\alpha(x)$, $\bar\psi_\alpha(x)$ with $x\in\mathbb{X}(\mathbb{I})$
goes to the generator of $\Gr_\Xi(\mathbb{I})$ of the same name;
in $\Gr_\Xi(\mathbb{I})$ the relations $\psi(x+g)=\chi_\parallel(g)\cdot\psi(x)$ and $\bar\psi(x+g)=\chi_\parallel(g)\cdot\bar\psi(x)$ hold for $g\in\Gamma_\parallel$
(Definition~\ref{def:slab-grassmann}).
For a torus $\mathbb{T}$ let $r_{\mathbb{T}}\colon\Gr(\mathbb{X})\to\Gr_{\mathbb{T}}$ be its quotient homomorphism (Definition~\ref{def:tori}).
Both $r_{\mathbb{T}}$ and $r_\Xi$ are linear, multiplicative and unital.
As they send generators to generators, they map a product of generators to a product of as many generators,
and hence even elements to even elements and odd elements to odd ones.

For $v\in\Z^3$ and $c\in\Adm(\ell)$ one has, moreover,
\[
r_\Xi\circ\tau_v=\tau_v\circ r_\Xi
\qquad\text{and}\qquad
r_\Xi\circ\Theta_c=\Theta_c\circ r_\Xi
\qquad\text{on }\Gr\left(\mathbb{X}(\mathbb{I})\right).
\]
On the left of these identities, $\tau_v$ is the automorphism $\psi(x)\mapsto\psi(x+v)$, $\bar\psi(x)\mapsto\bar\psi(x+v)$ of $\Gr(\mathbb{X})$,
and $\Theta_c=\Theta_{\ell,c}$ is the reflection of Grassmann variables at the mirror $(\ell,c)$ (Definition~\ref{def:grassmann-reflections});
by Lemma~\ref{lem:slab-automorphisms} they map $\Gr(\mathbb{X}(\mathbb{I}))$ onto $\Gr(\mathbb{X}(\mathbb{I}+\ell(v)))$ and onto $\Gr(\mathbb{X}(2c-\mathbb{I}))$, respectively,
and the homomorphism $r_\Xi$ that follows them is the one of the interval $\mathbb{I}+\ell(v)$, resp. $2c-\mathbb{I}$.
On the right, $\tau_v\colon\Gr_\Xi(\mathbb{I})\to\Gr_\Xi(\mathbb{I}+\ell(v))$ and $\Theta_c\colon\Gr_\Xi(\mathbb{I})\to\Gr_\Xi(2c-\mathbb{I})$ are the maps that they induce, as in the same lemma.
The two sides of the first identity are homomorphisms of algebras;
the two sides of the second are antilinear and reverse the order of products.
A map of either kind on $\Gr(\mathbb{X}(\mathbb{I}))$ is determined by its values on the generators,
and on each generator the two sides take the same value, since the maps on $\Gr_\Xi(\mathbb{I})$ are induced by those on $\Gr(\mathbb{X}(\mathbb{I}))$.
\end{definition}

\begin{lemma}[One lift for all boxes]\label{lem:global-lift}
For $\ell\in\mathcal{L}$ and $c<c'$ in $\Adm(\ell)$ denote by $\Z^3[c,c']$ the complex formed by the cells of $\Z^3$ lying in the closed slab $\{c\le\ell\le c'\}$,
and let $\mathcal{K}$ be either the complex of all cells of $\Z^3$ or one of the complexes $\Z^3[c,c']$.
If $u$ is an integer two-cochain on $\mathcal{K}$ and $du$ vanishes on every cube of $\mathcal{K}$,
then $u=dk$ for some integer one-cochain $k$ on $\mathcal{K}$.
As a consequence, the following two assertions hold.
\begin{enumerate}
\item\label{it:lift-torus}
Let $\Gamma$ be a period lattice and $(A,z)\in\mathfrak{C}(\Gamma)$ (Definition~\ref{def:torus-integrands}),
and let $\widetilde A$, $\widetilde z$ and $\widetilde F$ be the periodic extensions to $\Z^3$
of $A$, of $z$ and of the curvature $F=dA+\frac{2\pi}{q}\cdot z$.
There is a real one-cochain $\hat A$ on $\Z^3$ such that
\[
d\hat A=\widetilde F
\qquad\text{and}\qquad
\exp\left(iq\hat A(b)\right)=\exp\left(iq\widetilde A(b)\right)\quad\text{for all links }b\text{ of }\Z^3 .
\]
\item\label{it:lift-slab}
Let $\ell\in\mathcal{L}$ and $c<c'$ in $\Adm(\ell)$,
let $\Xi=(\Gamma_\parallel,\chi_\parallel)$ be a cross-section for $\ell$,
and let $\varpi=(A,z)\in\mathfrak{C}[c,c']$ (Definition~\ref{def:slab-configurations}).
Let $\widetilde A$ be the $\Gamma_\parallel$-periodic extension of $A$ to the links of $\Z^3[c,c']$,
and let $\widetilde z$ and $\widetilde{F(\varpi)}$ be those of $z$ and of $F(\varpi)$ to its plaquettes.
There is a real one-cochain $\hat A$ on $\Z^3[c,c']$ such that
$d\hat A=\widetilde{F(\varpi)}$ on the plaquettes of $\Z^3[c,c']$
and $\exp\left(iq\hat A(b)\right)=\exp\left(iq\widetilde A(b)\right)$ for all links $b$ of $\Z^3[c,c']$.
\end{enumerate}
\end{lemma}

\begin{proof}
Each of the two complexes contains, with every cell, all faces of that cell;
so $du$ is defined on the cubes of $\mathcal{K}$ and $dk$ on its plaquettes.
Cells carry the orientations of Definition~\ref{def:cochains}.

\emph{Product structure.}
If $\mathcal{K}$ is the complex of all cells, put $\mu_0:=3$.
If $\mathcal{K}=\Z^3[c,c']$, the vector $n_\ell$ has at most two non-zero components (Definition~\ref{def:mirrors}),
and $\mu_0\in\{1,2,3\}$ is chosen such that the component of index $\mu_0$ of $n_\ell$ vanishes.
Let $\mu<\nu$ be the two other indices,
and identify $\R^3$ with $\R^2\times\R$ through $x\mapsto\left((x_\mu,x_\nu),x_{\mu_0}\right)$.
Then every cell of $\Z^3$ is of the form $\sigma\times\{j\}$ or $\sigma\times[j,j+1]$,
where $\sigma$ is a cell of the lattice $\Z^2$ of the coordinates $(x_\mu,x_\nu)$ and $j\in\Z$.
If $\mathcal{K}$ is the complex of all cells, let $\mathcal{K}_2$ be the complex of all cells of $\Z^2$.
If $\mathcal{K}=\Z^3[c,c']$, the form $\ell$ does not depend on $x_{\mu_0}$;
it is a function of $(x_\mu,x_\nu)$,
and the closed slab $\{c\le\ell\le c'\}$ is the product of a strip of the plane of the coordinates $(x_\mu,x_\nu)$ with $\R$.
Let then $\mathcal{K}_2$ be the complex of the cells of $\Z^2$ contained in this strip.
A cell of $\Z^3$, resp. of $\Z^2$, lies in the closed slab, resp. in the strip, if and only if its corners do,
because the slab and the strip are convex and each cell is the smallest convex set that contains its corners;
the corners of $\sigma\times\{j\}$ and of $\sigma\times[j,j+1]$ are the points $(y,j)$, resp. $(y,j)$ and $(y,j+1)$, with $y$ a corner of $\sigma$.
Hence, in both cases, $\sigma\times\{j\}$ and $\sigma\times[j,j+1]$ belong to $\mathcal{K}$ if and only if $\sigma$ belongs to $\mathcal{K}_2$:
\[
\mathcal{K}=\mathcal{K}_2\times\Z .
\]
For $j\in\Z$ the cells $\sigma\times\{j\}$ with $\sigma\in\mathcal{K}_2$ form a subcomplex $\mathcal{K}_2\times\{j\}$ of $\mathcal{K}$, called the \emph{sheet} of index $j$.
For integers $a$ and $r$ let $(a,r)$ denote the site $x_\mu=a$, $x_\nu=r$ of $\Z^2$,
let $\ell(a,r)$ be the value of $\ell$ there,
and let $Q(a,r)$ be the square with the corners $(a,r)$, $(a+1,r)$, $(a+1,r+1)$, $(a,r+1)$.
The link from $(a,r)$ to $(a,r+1)$ is written $[(a,r),(a,r+1)]$.

\emph{Rows.}
For $r\in\Z$ let $\mathbb{I}_r$ be the set of the integers $a$ for which $[(a,r),(a,r+1)]$ belongs to $\mathcal{K}_2$.
Then $\mathbb{I}_r$ is an interval of $\Z$, possibly empty or unbounded, and for all $a,r\in\Z$
\begin{equation}\label{eq:slab-row-intervals}
Q(a,r)\in\mathcal{K}_2
\iff
a\in\mathbb{I}_r\ \text{and}\ a+1\in\mathbb{I}_r .
\end{equation}
If $\mathcal{K}$ is the complex of all cells, then $\mathbb{I}_r=\Z$ and \eqref{eq:slab-row-intervals} is evident.
Otherwise a link or a square belongs to $\mathcal{K}_2$ if and only if its corners lie in the strip,
again by convexity.
The function $a\mapsto\ell(a,r)$ is affine,
so each of the two conditions $c\le\ell(a,r)\le c'$ and $c\le\ell(a,r+1)\le c'$ defines an interval of integers $a$,
and $\mathbb{I}_r$ is the intersection of these two intervals.
Every corner of $Q(a,r)$ is an end point of $[(a,r),(a,r+1)]$ or of $[(a+1,r),(a+1,r+1)]$,
and every end point of these two links is a corner of $Q(a,r)$;
this gives \eqref{eq:slab-row-intervals}.

\emph{Oriented sums.}
For integers $a$, $a'$ and a function $f$ on $\Z$,
the sum $\sum_{a\le s<a'}f(s)$ has its usual meaning if $a\le a'$, and
\[
\sum_{a\le s<a'}f(s):=-\sum_{a'\le s<a}f(s)\qquad\text{if }a'<a .
\]
With this convention $\sum_{a\le s<a}f(s)=0$ and
\[
\sum_{a\le s<a'+1}f(s)-\sum_{a\le s<a'}f(s)=f(a')\qquad\text{for all }a'\in\Z .
\]

\emph{Construction of $k$.}
On the links $\{y\}\times[j,j+1]$, $y$ a site of $\mathcal{K}_2$ and $j\in\Z$,
which are the links of $\mathcal{K}$ of the direction $\mu_0$, put $k:=0$.
On the links of the direction $\mu$ in the sheet of index $0$ put $k:=0$ as well.
For every $r$ with $\mathbb{I}_r\neq\emptyset$ choose $a_r\in\mathbb{I}_r$,
and for $a\in\mathbb{I}_r$ put
\[
k\left([(a,r),(a,r+1)]\times\{0\}\right):=\sum_{a_r\le s<a}u\left(Q(s,r)\times\{0\}\right).
\]
The squares on the right belong to $\mathcal{K}_2$:
for $\min(a,a_r)\le s<\max(a,a_r)$ the integers $s$ and $s+1$ lie between $a_r$ and $a$,
hence in the interval $\mathbb{I}_r$, and \eqref{eq:slab-row-intervals} applies.
The links of $\mathcal{K}_2$ of the direction $\nu$ are exactly the links $[(a,r),(a,r+1)]$ with $a\in\mathbb{I}_r$;
so $k$ is now defined on all links of the sheet of index $0$.
Finally,
for a link $b$ of $\mathcal{K}_2$ of the direction $\nu'\in\{\mu,\nu\}$ and for $j\in\Z$, put
\[
k\left(b\times\{j\}\right):=k\left(b\times\{0\}\right)+s_{\nu'}\cdot\sum_{0\le j'<j}u\left(b\times[j',j'+1]\right),
\qquad
s_{\nu'}:=
\begin{cases}
-1&\text{if }\nu'<\mu_0,\\
+1&\text{if }\nu'>\mu_0.
\end{cases}
\]
For $j=0$ this repeats the values already given,
and so it defines an integer one-cochain $k$ on $\mathcal{K}$.

\emph{Verification.}
Let $b$ be a link of $\mathcal{K}_2$ of the direction $\nu'$, let $j\in\Z$,
and let $x$ be the lowest corner of the plaquette $b\times[j,j+1]$.
As $k$ vanishes on the links of the direction $\mu_0$,
the formula for the coboundary of a one-cochain in Definition~\ref{def:cochains} gives
$(dk)_{\nu'\mu_0}(x)=k_{\nu'}(x)-k_{\nu'}(x+e_{\mu_0})$ if $\nu'<\mu_0$,
and $(dk)_{\mu_0\nu'}(x)=k_{\nu'}(x+e_{\mu_0})-k_{\nu'}(x)$ if $\nu'>\mu_0$.
Together with the definition of $k$ and the rule for oriented sums this gives
\[
(dk)\left(b\times[j,j+1]\right)
=s_{\nu'}\cdot\left(k\left(b\times\{j+1\}\right)-k\left(b\times\{j\}\right)\right)
=s_{\nu'}^2\cdot u\left(b\times[j,j+1]\right)
=u\left(b\times[j,j+1]\right).
\]
In the sheet of index $0$ the cochain $k$ vanishes on the links of the direction $\mu$.
Let $Q(a,r)$ be a square of $\mathcal{K}_2$;
then $a,a+1\in\mathbb{I}_r$ by \eqref{eq:slab-row-intervals}.
As $\mu<\nu$, the square $Q(a,r)\times\{0\}$ is the plaquette spanned by $e_\mu$ and $e_\nu$ at its lowest corner,
and Definition~\ref{def:cochains}, the definition of $k$ on the sheet of index $0$ and the rule for oriented sums give
\[
(dk)\left(Q(a,r)\times\{0\}\right)
=k\left([(a+1,r),(a+1,r+1)]\times\{0\}\right)-k\left([(a,r),(a,r+1)]\times\{0\}\right)
=u\left(Q(a,r)\times\{0\}\right).
\]
Consider now the integer two-cochain $u':=u-dk$ on $\mathcal{K}$.
Because $d\,dk=0$ by \eqref{eq:slab-incidence-square}, its coboundary vanishes on each cube of $\mathcal{K}$, as that of $u$ does;
by the two computations above, $u'$ is zero on each plaquette $b\times[j,j+1]$ and on each square $Q\times\{0\}$ with $Q$ a square of $\mathcal{K}_2$.
A cube of $\mathcal{K}$ has the form $Q\times[j,j+1]$, with $Q$ a square of $\mathcal{K}_2$ and $j\in\Z$.
Its two-dimensional faces are the bottom $Q\times\{j\}$, the top $Q\times\{j+1\}$ and, for each of the four edges $b$ of $Q$, the lateral face $b\times[j,j+1]$.
Since $u'$ is zero on the lateral faces, the formula for the coboundary of a two-cochain in Definition~\ref{def:cochains} reduces to
\[
0=(du')\left(Q\times[j,j+1]\right)=\pm\left(u'\left(Q\times\{j+1\}\right)-u'\left(Q\times\{j\}\right)\right).
\]
By induction upwards and downwards from $j=0$,
$u'(Q\times\{j\})=u'(Q\times\{0\})=0$ for all $j\in\Z$.
Every plaquette of $\mathcal{K}$ is of the form $b\times[j,j+1]$ or $Q\times\{j\}$,
so $u'=0$, that is, $u=dk$.

\emph{The two consequences.}
For \eqref{it:lift-torus} let $\mathcal{K}$ be the complex of all cells.
The cochains of $\Z^3/\Gamma$ are the $\Gamma$-periodic cochains on $\Z^3$ (Definition~\ref{def:configurations}),
so $\widetilde z$ is an integer two-cochain on $\Z^3$ with $d\widetilde z=0$,
and $\widetilde F=d\widetilde A+\frac{2\pi}{q}\cdot\widetilde z$.
By the first part, $\widetilde z=dk$ with an integer one-cochain $k$ on $\Z^3$.
The cochain $\hat A:=\widetilde A+\frac{2\pi}{q}\cdot k$ satisfies
\[
d\hat A=d\widetilde A+\frac{2\pi}{q}\cdot\widetilde z=\widetilde F,
\qquad
\exp\left(iq\hat A(b)\right)=\exp\left(iq\widetilde A(b)\right)\cdot\exp\left(2\pi i\cdot k(b)\right)=\exp\left(iq\widetilde A(b)\right).
\]
For \eqref{it:lift-slab} let $\mathcal{K}=\Z^3[c,c']$.
Since a cube is not contained in any plane, it is contained neither in $\{\ell=c\}$ nor in $\{\ell=c'\}$;
hence the $\Gamma_\parallel$-orbits of the cubes of $\Z^3[c,c']$ are precisely the elements of $\mathcal{B}^3(c,c')$,
and $dz=0$ on them (Definition~\ref{def:slab-configurations}).
Since $\sys(\Gamma_\parallel)\ge4$, distinct faces of a cell lie in distinct $\Gamma_\parallel$-orbits,
and the coboundary on the cylinder $\Z^3/\Gamma_\parallel$ is computed on representatives.
Hence $d\widetilde z$ is the periodic extension of $dz$ and vanishes on the cubes of $\Z^3[c,c']$,
and on the plaquettes of $\Z^3[c,c']$ the cochain $d\widetilde A$ is the periodic extension of $dA$,
so that $\widetilde{F(\varpi)}=d\widetilde A+\frac{2\pi}{q}\cdot\widetilde z$ there.
By the first part, $\widetilde z=dk$ with an integer one-cochain $k$ on $\Z^3[c,c']$,
and $\hat A:=\widetilde A+\frac{2\pi}{q}\cdot k$ has the two properties of \eqref{it:lift-slab},
by the computation made for \eqref{it:lift-torus}.
\end{proof}

The cochain $\hat A$ of Lemma~\ref{lem:global-lift}\eqref{it:lift-torus} is in general not $\Gamma$-periodic.
Indeed, if a $\Gamma$-periodic real one-cochain $\hat A$ has the two properties stated there,
then $\frac{q}{2\pi}\cdot\left(\hat A-\widetilde A\right)$ is a $\Gamma$-periodic integer one-cochain with the coboundary $\widetilde z$,
so that $z=dk'$ with $k'\in C^1(\Gamma;\Z)$;
conversely, if $z=dk'$ with such a $k'$, then $\widetilde A+\frac{2\pi}{q}\cdot k'$ is $\Gamma$-periodic and has the two properties.
Let $B$ be a box, and in the situation of Lemma~\ref{lem:global-lift}\eqref{it:lift-slab} let $B$ be contained in $\{c\le\ell\le c'\}$.
The restriction to $B^{(1)}$ of the cochain $\hat A$ of Lemma~\ref{lem:global-lift}\eqref{it:lift-torus}, resp. \eqref{it:lift-slab},
has on $B^{(2)}$ the coboundary $\widetilde F$, resp. $\widetilde{F(\varpi)}$, and on $B^{(1)}$ the same phases as $\widetilde A$.
Thus one cochain on $\Z^3$, resp. on $\Z^3[c,c']$, yields cochains with these two properties on all boxes simultaneously,
and any two of them coincide on the links common to their boxes.
On a torus, or on the cylinder $\Z^3/\Gamma_\parallel$, cochains chosen on single boxes do not in general fit together in this way:
the image of a union of boxes in the quotient may fail to be contractible,
and two such cochains may differ on a link that lies in the images of both boxes.

\begin{lemma}[Admissible products as integrands]\label{lem:product-integrands}
Let $\mathbb{T}=(\Gamma,\chi)$ be a torus, let $\ell\in\mathcal{L}$, let $\Xi=(\Gamma_\parallel,\chi_\parallel)$ be a cross-section for $\ell$ (Definition~\ref{def:slabs}), and let $c<c'$ be elements of $\Adm(\ell)$.
Elements of $\mathcal{A}(\R^3)$ are regarded as maps $\hat A\mapsto G(\hat A)$ from the real one-cochains of $\Z^3$ to $\Gr(\mathbb{X})$, as in Definition~\ref{def:product-maps}.
\begin{enumerate}
\item\label{it:product-torus}
Let $G=\Phi\cdot G_1\cdots G_r$ be a product as in Definition~\ref{def:integrands}, with $G_j\in\mathcal{G}^{\mathrm{ev}}(B_j)$, that is read on $\mathbb{T}$, and let $(A,z)\in\mathfrak{C}(\Gamma)$.
Denote by $\widetilde A$ and $\widetilde F$ the $\Gamma$-periodic extensions to $\Z^3$ of $A$ and of $F=dA+(2\pi/q)z$.
For $j=1,\dots,r$ there is a real one-cochain $\hat A^{(j)}$ on $B_j^{(1)}$ with
\[
d\hat A^{(j)}=\widetilde F\ \text{ on } B_j^{(2)},
\qquad
\exp\left(iq\hat A^{(j)}(b)\right)=\exp\left(iq\widetilde A(b)\right)\ \text{ for } b\in B_j^{(1)},
\]
and the real one-cochains on $B_j^{(1)}$ with these two properties are exactly the cochains $\hat A^{(j)}+(2\pi/q)\,d\lambda$, $\lambda$ an integer valued function on the sites of $B_j$.
The element
\[
G(A,z):=\Phi(\widetilde F)\cdot\prod_{j=1}^{r}r_{\mathbb{T}}\left(G_j(\hat A^{(j)})\right)\in\Gr_{\mathbb{T}}
\]
does not depend on the choice of the cochains $\hat A^{(j)}$, and $(A,z)\mapsto G(A,z)$ is an even element of $\mathcal{I}(\mathbb{T})$.
If the product embeds in $\mathbb{T}$, then $G(A,z)$ is the element of $\Gr_{\mathbb{T}}$ that Definition~\ref{def:integrands} assigns to the product at the point of $\Conf(\Gamma)$ represented by $(A,z)$.
\item\label{it:product-slab}
Let $G=\Phi\cdot G_1\cdots G_r$, with $G_j\in\mathcal{G}^{\mathrm{ev}}(B_j)$, be a product of $\mathcal{A}(\{c<\ell<c'\})$ that embeds in $\Gamma_\parallel$; thus the boxes $B_j$ and the plaquettes read by $\Phi$ are contained in $\{c<\ell<c'\}$.
Let $\varpi=(A,z)\in\mathfrak{C}[c,c']$, and denote by $\widetilde A$ and $\widetilde F$ the $\Gamma_\parallel$-periodic extensions of $A$ and of $F(\varpi)$ to the cells of $\Z^3[c,c']$.
For $j=1,\dots,r$ there is a real one-cochain $\hat A^{(j)}$ on $B_j^{(1)}$ with the two properties displayed in \eqref{it:product-torus}, unique up to the addition of $(2\pi/q)\,d\lambda$, $\lambda$ an integer valued function on the sites of $B_j$.
The element
\[
G(\varpi):=\Phi(\widetilde F)\cdot\prod_{j=1}^{r}r_\Xi\left(G_j(\hat A^{(j)})\right)\in\Gr_\Xi((c,c'))
\]
does not depend on the choice of the cochains $\hat A^{(j)}$, and $\varpi\mapsto G(\varpi)$ is an even element of $\mathcal{I}_\Xi[c,c']$.
\item\label{it:sum-torus}
For $G\in\mathcal{A}(\mathbb{T})$, $(A,z)\in\mathfrak{C}(\Gamma)$ and a real one-cochain $\hat A$ on $\Z^3$ as in Lemma~\ref{lem:global-lift}\eqref{it:lift-torus}, put
\[
G(A,z):=r_{\mathbb{T}}\left(G(\hat A)\right).
\]
This element of $\Gr_{\mathbb{T}}$ does not depend on the choice of $\hat A$; for a product that is read on $\mathbb{T}$ it is the element of \eqref{it:product-torus}.
The map $G\mapsto G(\cdot,\cdot)$ is linear from $\mathcal{A}(\mathbb{T})$ into the even part of $\mathcal{I}(\mathbb{T})$, and $Z_{\mathbb{T}}[G]$, the value of the functional of Definition~\ref{def:torus-integrands} at the function $G(\cdot,\cdot)$, is linear in $G\in\mathcal{A}(\mathbb{T})$.
If $G\in\mathcal{A}(\R^3)$ embeds in $\mathbb{T}$, so that $G\in\mathcal{A}(\mathbb{T})$ (Definition~\ref{def:torus-integrands}), then $G(\cdot,\cdot)$ is the function that Definition~\ref{def:integrands} assigns to $G$, and $Z_{\mathbb{T}}[G]$ is the value at $G$ of the functional of Definition~\ref{def:functional}.
\item\label{it:sum-slab}
For $G\in\mathcal{A}_\Xi[c,c']$, $\varpi\in\mathfrak{C}[c,c']$ and a real cochain $\hat A$ as in Lemma~\ref{lem:global-lift}\eqref{it:lift-slab}, extended by $0$ to the links of $\Z^3$ that are not contained in $\{c\le\ell\le c'\}$, put
\[
G(\varpi):=r_\Xi\left(G(\hat A)\right).
\]
This element of $\Gr_\Xi((c,c'))$ does not depend on the choice of $\hat A$; for a product of $\mathcal{A}(\{c<\ell<c'\})$ that embeds in $\Gamma_\parallel$ it is the element of \eqref{it:product-slab}.
The map $G\mapsto G(\cdot)$ is linear from $\mathcal{A}_\Xi[c,c']$ into the even part of $\mathcal{I}_\Xi[c,c']$.
Let moreover $c''\in\Adm(\ell)$ with $c''>c'$, $G'\in\mathcal{A}_\Xi[c',c'']$ and $v\in\Z^3$.
Then
\begin{align*}
G\cdot G'&\in\mathcal{A}_\Xi[c,c''], & G^{\tau_v}&\in\mathcal{A}_\Xi[c+\ell(v),c'+\ell(v)],\\
\Theta_{c'}G&\in\mathcal{A}_\Xi[c',2c'-c], & \Theta_cG&\in\mathcal{A}_\Xi[2c-c',c],
\end{align*}
and
\begin{align*}
(G\cdot G')(\varpi)&=G(\varpi|_{[c,c']})\cdot G'(\varpi|_{[c',c'']}) &&\text{for }\varpi\in\mathfrak{C}[c,c''],\\
G^{\tau_v}(\varpi)&=\tau_v\left(G(\tau_v^*\varpi)\right) &&\text{for }\varpi\in\mathfrak{C}[c+\ell(v),c'+\ell(v)],\\
(\Theta_{c'}G)(\varpi)&=\Theta_{c'}\left(G(\theta_{c'}^*\varpi)\right) &&\text{for }\varpi\in\mathfrak{C}[c',2c'-c],\\
(\Theta_cG)(\varpi)&=\Theta_c\left(G(\theta_c^*\varpi)\right) &&\text{for }\varpi\in\mathfrak{C}[2c-c',c];
\end{align*}
on the right, $\tau_v$, $\Theta_{c'}$ and $\Theta_c$ act on the Grassmann algebras of the cylinder as in Lemma~\ref{lem:slab-automorphisms}.
In particular, the map $G\mapsto G(\cdot)$ is multiplicative for adjacent slabs, with the product of Definition~\ref{def:slab-integrands}, and it intertwines $G\mapsto G^{\tau_v}$ with the translation $\tau_v$ and $G\mapsto\Theta_{c'}G$ with the reflection $\Theta_{c'}$ of that definition.
\end{enumerate}
\end{lemma}

\begin{proof}
\emph{Integer cochains on a product of intervals.}
Let $a_\mu\le a'_\mu$ be integers, $\mu=1,2,3$, and let $\mathcal{K}$ be the complex of the cells of $\Z^3$ that are contained in $\prod_{\mu=1}^{3}[a_\mu,a'_\mu]$, that is, of the cells all of whose corners lie in this product.
Over the integers the Poincar\'e lemma holds on $\mathcal{K}$ in degrees one and two, with explicit primitives.
Let first $u$ be an integer two-cochain on $\mathcal{K}$ with $du=0$ on the cubes of $\mathcal{K}$, and define an integer one-cochain $k$ on the links of $\mathcal{K}$ by
\begin{align*}
k_1&:=0,\\
k_2(x)&:=\sum_{a_1\le s<x_1}u_{12}(s,x_2,x_3),\\
k_3(x)&:=\sum_{a_1\le s<x_1}u_{13}(s,x_2,x_3)+\sum_{a_2\le s<x_2}u_{23}(a_1,s,x_3).
\end{align*}
The formulas of Definition~\ref{def:cochains} give $(dk)_{12}=u_{12}$ and $(dk)_{13}=u_{13}$ at once, and
\begin{align*}
(dk)_{23}(x)={}&u_{23}(a_1,x_2,x_3)+\sum_{a_1\le s<x_1}\bigl(u_{12}(s,x_2,x_3)-u_{12}(s,x_2,x_3+1)\\
&+u_{13}(s,x_2+1,x_3)-u_{13}(s,x_2,x_3)\bigr).
\end{align*}
As $du$ vanishes on the cube with the lowest corner $(s,x_2,x_3)$, the term of index $s$ equals $u_{23}(s+1,x_2,x_3)-u_{23}(s,x_2,x_3)$.
The sum telescopes, $(dk)_{23}=u_{23}$, and therefore $u=dk$.

Let next $k'$ be an integer one-cochain on $\mathcal{K}$ with $dk'=0$ on the plaquettes of $\mathcal{K}$.
Let $\lambda(x)$ be the sum of $k'$ along the path in $\mathcal{K}$ from $(a_1,a_2,a_3)$ to the site $x$ whose steps go first in the direction $e_1$, then in the direction $e_2$, then in the direction $e_3$:
\[
\lambda(x):=\sum_{a_1\le s<x_1}k'_1(s,a_2,a_3)+\sum_{a_2\le s<x_2}k'_2(x_1,s,a_3)+\sum_{a_3\le s<x_3}k'_3(x_1,x_2,s).
\]
Then $\lambda(x+e_3)-\lambda(x)=k'_3(x)$ by definition.
Next,
\[
\lambda(x+e_2)-\lambda(x)=k'_2(x_1,x_2,a_3)+\sum_{a_3\le s<x_3}\left(k'_3(x_1,x_2+1,s)-k'_3(x_1,x_2,s)\right),
\]
and $dk'=0$ on the plaquette spanned by $e_2,e_3$ at the site $(x_1,x_2,s)$ turns the term of index $s$ into $k'_2(x_1,x_2,s+1)-k'_2(x_1,x_2,s)$; the right side telescopes to $k'_2(x)$.
Finally,
\begin{align*}
\lambda(x+e_1)-\lambda(x)={}&k'_1(x_1,a_2,a_3)\\
&+\sum_{a_2\le s<x_2}\left(k'_2(x_1+1,s,a_3)-k'_2(x_1,s,a_3)\right)\\
&+\sum_{a_3\le s<x_3}\left(k'_3(x_1+1,x_2,s)-k'_3(x_1,x_2,s)\right).
\end{align*}
Here $dk'=0$ on the plaquettes spanned by $e_1,e_2$ at the sites $(x_1,s,a_3)$ and on those spanned by $e_1,e_3$ at the sites $(x_1,x_2,s)$ turns the terms into $k'_1(x_1,s+1,a_3)-k'_1(x_1,s,a_3)$ and $k'_1(x_1,x_2,s+1)-k'_1(x_1,x_2,s)$.
The two sums telescope to $k'_1(x_1,x_2,a_3)-k'_1(x_1,a_2,a_3)$ and $k'_1(x)-k'_1(x_1,x_2,a_3)$, so that $\lambda(x+e_1)-\lambda(x)=k'_1(x)$.
Hence $k'=d\lambda$, with $\lambda$ integer valued.

Consider now two real one-cochains $\hat A$, $\hat A'$ on the links of $\mathcal{K}$ with $d\hat A=d\hat A'$ on the plaquettes of $\mathcal{K}$ and $\exp\left(iq\hat A(b)\right)=\exp\left(iq\hat A'(b)\right)$ for every link $b$ of $\mathcal{K}$.
The cochain $k':=(q/2\pi)(\hat A'-\hat A)$ is integer valued, and $dk'=0$ on the plaquettes of $\mathcal{K}$.
By the second fact there is an integer valued function $\lambda$ on the sites of $\mathcal{K}$ with
\begin{equation}\label{eq:slab-lift-difference}
\hat A'=\hat A+\frac{2\pi}{q}\,d\lambda .
\end{equation}
Conversely, for every integer valued $\lambda$ the cochain $\hat A+(2\pi/q)\,d\lambda$ has the same coboundary as $\hat A$, since $dd\lambda=0$, and the same phases, since $\exp\left(2\pi i(d\lambda)(b)\right)=1$.

\emph{Proof of \eqref{it:product-torus} and \eqref{it:product-slab}.}
The two items are proved together.
For \eqref{it:product-torus} put $\Gamma':=\Gamma$; the configuration is $(A,z)\in\mathfrak{C}(\Gamma)$, and $\widetilde z$ denotes the $\Gamma$-periodic extension of $z$ to $\Z^3$.
For \eqref{it:product-slab} put $\Gamma':=\Gamma_\parallel$; the configuration is $\varpi=(A,z)\in\mathfrak{C}[c,c']$, the letter $F$ stands for $F(\varpi)$, and $\widetilde z$ denotes the $\Gamma_\parallel$-periodic extension of $z$ to the cells of $\Z^3[c,c']$.
In both cases the sets $B_j+g$, for all $j$ and all $g\in\Gamma'$, are pairwise disjoint: for \eqref{it:product-torus} because the product is read on $\mathbb{T}$ (Definition~\ref{def:torus-integrands}), for \eqref{it:product-slab} because it embeds in $\Gamma_\parallel$ (Definition~\ref{def:integrands}).
If a cell $\sigma$ is contained in $B_j$ and $\sigma+g$ is contained in $B_k$ for some $g\in\Gamma'$, then $\sigma+g\subset(B_j+g)\cap B_k$, whence $j=k$ and $g=0$.
So the cells contained in $B_1\cup\dots\cup B_r$ lie in pairwise distinct $\Gamma'$-orbits.
The same holds for the sites of $\mathbb{X}(B_1)\cup\dots\cup\mathbb{X}(B_r)$, because $\mathbb{X}(B_j)\subset B_j$: every site $x$ lies in the cube $\pi(x)+[-1/2,1/2]^3$ (Definition~\ref{def:fermion-sites}).

Write $B_j=y_j+[-r_j,r_j]^3$ with $r_j>0$, and let $y_{j,1},y_{j,2},y_{j,3}$ be the coordinates of $y_j$.
A cell of $\Z^3$ is contained in the convex set $B_j$ if and only if all its corners are; these cells form the complex $\mathcal{K}_j$ of the cells contained in the product of the integer intervals $\left[\lceil y_{j,\mu}-r_j\rceil,\lfloor y_{j,\mu}+r_j\rfloor\right]$, $\mu=1,2,3$.
If one of these intervals is empty, $B_j$ contains no site, the complex $\mathcal{K}_j$ is empty, and so are $B_j^{(1)}$, $B_j^{(2)}$ and $\mathbb{X}(B_j)$.
In the case of \eqref{it:product-slab} the cells contained in $B_j$ lie in the open slab $\{c<\ell<c'\}$; they are cells of $\Z^3[c,c']$, contained in neither of the planes $\{\ell=c\}$, $\{\ell=c'\}$, and their $\Gamma_\parallel$-orbits belong to $\mathcal{B}(c,c')$.
For the same reason the plaquettes read by $\Phi$ have their orbits in $\mathcal{B}^2(c,c')$, so that $\Phi(\widetilde F)$ is defined.
The integer two-cochain $\widetilde z$ is closed on the cubes of $\mathcal{K}_j$: for \eqref{it:product-torus} because $dz=0$, for \eqref{it:product-slab} because the orbits of these cubes lie in $\mathcal{B}^3(c,c')$, where $dz=0$ by Definition~\ref{def:slab-configurations}.

By the first of the two facts above there is an integer one-cochain $k^{(j)}$ on $\mathcal{K}_j$ with $dk^{(j)}=\widetilde z$ on the plaquettes of $\mathcal{K}_j$.
The cochain
\[
\hat A^{(j)}:=\widetilde A+\frac{2\pi}{q}\,k^{(j)}\quad\text{on }B_j^{(1)}
\]
satisfies $d\hat A^{(j)}=d\widetilde A+(2\pi/q)\widetilde z=\widetilde F$ on $B_j^{(2)}$, and
\[
\exp\left(iq\hat A^{(j)}(b)\right)=\exp\left(iq\widetilde A(b)\right)\cdot\exp\left(2\pi ik^{(j)}(b)\right)=\exp\left(iq\widetilde A(b)\right),\qquad b\in B_j^{(1)}.
\]
By \eqref{eq:slab-lift-difference}, applied to $\mathcal{K}_j$, and by the converse stated after it, the real one-cochains on $B_j^{(1)}$ with these two properties are exactly the cochains $\hat A^{(j)}+(2\pi/q)\,d\lambda$ with $\lambda$ integer valued on the sites of $B_j$.

Let $\lambda$ be such a function.
Gauge invariance of $G_j$ (Definition~\ref{def:integrands}), applied with the real function $(2\pi/q)\lambda$, says that $G_j\left(\hat A-(2\pi/q)\,d\lambda\right)$, transformed by $\psi(x)\mapsto\exp\left(2\pi i\lambda(\pi(x))\right)\psi(x)$ and $\bar\psi(x)\mapsto\exp\left(-2\pi i\lambda(\pi(x))\right)\bar\psi(x)$, equals $G_j(\hat A)$.
As $\lambda$ is integer valued, this transformation is the identity, and
\[
G_j\left(\hat A-\frac{2\pi}{q}\,d\lambda\right)=G_j(\hat A)
\]
for every real one-cochain $\hat A$ on $B_j^{(1)}$.
Therefore $G_j(\hat A^{(j)})$, and with it $G(A,z)$, resp. $G(\varpi)$, does not depend on the choice of the cochains $\hat A^{(j)}$.

The integer shifts, $(A,z)\mapsto(A+(2\pi/q)k,z-dk)$ with $k\in C^1(\Gamma;\Z)$ (Definition~\ref{def:configurations}) in the case of \eqref{it:product-torus}, and $\varpi\mapsto\varpi+[k]$ with $k\in\Z^{[c,c']^{(1)}}$ in the case of \eqref{it:product-slab}, change neither $F$ nor the phases $\exp\left(iqA(b)\right)$.
They leave the set of the cochains $\hat A^{(j)}$ with the two properties unchanged, and $G$ is invariant under them.
For measurability, note that $z$ ranges over a countable set, the torus and the closed slab of the cylinder having finitely many plaquettes.
At fixed $z$ the cochains $k^{(j)}$ are fixed and $A\mapsto\hat A^{(j)}$ is affine; as the $G_j$ and $\Phi$ are Borel and $F$ is a continuous function of $A$ at fixed $z$, every coefficient of $G$ is a Borel function of the configuration.

Each $G_j(\hat A^{(j)})$ is an even element of $\Gr(\mathbb{X}(B_j))$.
The sites of $\mathbb{X}(B_1)\cup\dots\cup\mathbb{X}(B_r)$ lie in pairwise distinct $\Gamma'$-orbits; hence $r_{\mathbb{T}}$, resp. $r_\Xi$, maps the monomials in the generators over this union injectively, up to signs, onto monomials of a monomial basis of $\Gr_{\mathbb{T}}$, resp. of $\Gr_\Xi((c,c'))$ (Definitions~\ref{def:tori} and \ref{def:slab-grassmann}).
In the case of \eqref{it:product-torus} this shows that $r_{\mathbb{T}}\left(G_j(\hat A^{(j)})\right)$ is the element of $\Gr_{\mathbb{T}}$ that the even element $G_j(\hat A^{(j)})$ defines by Definition~\ref{def:tori}; in the case of \eqref{it:product-slab} the values lie in $\Gr_\Xi((c,c'))$, because $\mathbb{X}(B_j)\subset B_j\subset\{c<\ell<c'\}$.
A monomial over the union is, up to sign, the product of its parts over the single sets $\mathbb{X}(B_j)$, and these parts are determined by it.
Consequently every coefficient of $G(A,z)$, resp. of $G(\varpi)$, in such a basis is either $0$ or, up to sign, $\Phi(\widetilde F)$ times a product of coefficients, one of each $G_j(\hat A^{(j)})$; and the value is even.

For the growth let $C_j<\infty$ and $\varepsilon_j>0$ be constants of $G_j$ as in Definition~\ref{def:integrands}, and put
\[
\varepsilon:=\min\{1/2,\varepsilon_1,\dots,\varepsilon_r\},
\]
so that $\varepsilon=1/2$ if $r=0$.
Then $0<\varepsilon\le1/2$, and $1/2-\varepsilon_j\le1/2-\varepsilon$ for every $j$.
The plaquettes of $B_1^{(2)}\cup\dots\cup B_r^{(2)}$ lie in pairwise distinct $\Gamma'$-orbits, which are plaquettes of the torus, resp. elements of $\mathcal{B}^2(c,c')$.
As $1/2-\varepsilon\ge0$, one obtains
\[
\sum_{j=1}^{r}\left(\frac12-\varepsilon_j\right)\sum_{\sigma\in B_j^{(2)}}\widetilde F_\sigma^2
\le\left(\frac12-\varepsilon\right)\sum_{j=1}^{r}\sum_{\sigma\in B_j^{(2)}}\widetilde F_\sigma^2
\le\left(\frac12-\varepsilon\right)\sum_{\sigma}F_\sigma^2,
\]
where the last sum is taken over the plaquettes of the torus, so that it equals $\|F\|^2$, in the case of \eqref{it:product-torus}, and over $\mathcal{B}^2(c,c')$ in the case of \eqref{it:product-slab}.
Since $d\hat A^{(j)}=\widetilde F$ on $B_j^{(2)}$, each coefficient of $G_j(\hat A^{(j)})$ has modulus at most
\[
C_j\cdot\exp\left(\left(\frac12-\varepsilon_j\right)\sum_{\sigma\in B_j^{(2)}}\widetilde F_\sigma^2\right).
\]
So every coefficient of the value of $G$ is bounded in modulus by
\[
\sup|\Phi|\cdot C_1\cdots C_r\cdot\exp\left(\left(\frac12-\varepsilon\right)\sum_{\sigma}F_\sigma^2\right).
\]
Together with the invariance and the measurability shown above, this proves that $G$ is an even element of $\mathcal{I}(\mathbb{T})$, resp. of $\mathcal{I}_\Xi[c,c']$, with the constants $\sup|\Phi|\cdot C_1\cdots C_r$ and $\varepsilon$.

If the product of \eqref{it:product-torus} embeds in $\mathbb{T}$, each $B_j$ embeds in $\Gamma$, and the cochains $\hat A^{(j)}$ are precisely the local lifts on the boxes $B_j$ of Definition~\ref{def:integrands}.
The element $G(A,z)$ is then, by construction, the one that this definition assigns to the product at the class of $(A,z)$.

\emph{Proof of \eqref{it:sum-torus} and of the first assertions of \eqref{it:sum-slab}.}
Consider first a product $H=\Phi\cdot G_1\cdots G_r$ that is read on $\mathbb{T}$, resp. a product of $\mathcal{A}(\{c<\ell<c'\})$ that embeds in $\Gamma_\parallel$.
Let $\hat A$ be a cochain as in Lemma~\ref{lem:global-lift}\eqref{it:lift-torus} for $(A,z)$, resp. as in Lemma~\ref{lem:global-lift}\eqref{it:lift-slab} for $\varpi$, extended by $0$.
By Definition~\ref{def:product-maps},
\[
H(\hat A)=\Phi(d\hat A)\cdot G_1\left(\hat A|_{B_1^{(1)}}\right)\cdots G_r\left(\hat A|_{B_r^{(1)}}\right).
\]
In the second case the boxes $B_j$ and the plaquettes read by $\Phi$ are contained in $\{c<\ell<c'\}$.
So $H(\hat A)$ involves only values of $\hat A$ on links of $\Z^3[c,c']$, where the extension by $0$ plays no role, and $H(\hat A)\in\Gr(\mathbb{X}((c,c')))$.
In both cases $d\hat A=\widetilde F$ on the plaquettes read by $\Phi$ and on each $B_j^{(2)}$, and the restriction of $\hat A$ to $B_j^{(1)}$ has the two properties of the cochain $\hat A^{(j)}$ of \eqref{it:product-torus}, resp. \eqref{it:product-slab}.
Since $r_{\mathbb{T}}$ and $r_\Xi$ are multiplicative (Definition~\ref{def:product-maps}), $r_{\mathbb{T}}\left(H(\hat A)\right)$, resp. $r_\Xi\left(H(\hat A)\right)$, is the element $H(A,z)$ of \eqref{it:product-torus}, resp. $H(\varpi)$ of \eqref{it:product-slab}; by these items it does not depend on $\hat A$.

Let now $G$ be an arbitrary element of $\mathcal{A}(\mathbb{T})$, resp. of $\mathcal{A}_\Xi[c,c']$.
A complex multiple of a product is again a product, with the same boxes; hence $G=H_1+\dots+H_n$ with products $H_i$ of the kind just considered, and
\[
r_{\mathbb{T}}\left(G(\hat A)\right)=\sum_{i=1}^{n}r_{\mathbb{T}}\left(H_i(\hat A)\right)=\sum_{i=1}^{n}H_i(A,z),
\qquad\text{resp.}\qquad
r_\Xi\left(G(\hat A)\right)=\sum_{i=1}^{n}H_i(\varpi).
\]
The left sides are formed with the map $G$ alone, not with its decomposition into products.
The right sides do not involve $\hat A$; as functions of the configuration they belong to the linear space $\mathcal{I}(\mathbb{T})$, resp. $\mathcal{I}_\Xi[c,c']$, and they are even.
So $G(A,z)$, resp. $G(\varpi)$, is well defined and independent of $\hat A$, and $G(\cdot,\cdot)$, resp. $G(\cdot)$, is an even element of that space.
At a fixed configuration and a fixed $\hat A$ the maps $G\mapsto r_{\mathbb{T}}\left(G(\hat A)\right)$ and $G\mapsto r_\Xi\left(G(\hat A)\right)$ are linear; hence $G\mapsto G(\cdot,\cdot)$ and $G\mapsto G(\cdot)$ are linear.
The functional of Definition~\ref{def:torus-integrands} is linear on $\mathcal{I}(\mathbb{T})$, and therefore $Z_{\mathbb{T}}[G]$ is linear in $G\in\mathcal{A}(\mathbb{T})$.
Finally, let $G\in\mathcal{A}(\R^3)$ embed in $\mathbb{T}$.
At the point of $\Conf(\Gamma)$ represented by $(A,z)$, Definition~\ref{def:integrands} assigns to $G$ the element $r_{\mathbb{T}}\left(G(\hat A)\right)$, formed with a cochain $\hat A$ as in Lemma~\ref{lem:global-lift}\eqref{it:lift-torus}; this is $G(A,z)$.
Applied to this function, the formula of Definition~\ref{def:torus-integrands} is that of Definition~\ref{def:functional}, and therefore $Z_{\mathbb{T}}[G]$ is the value at $G$ of the functional of Definition~\ref{def:functional}.

\emph{Products of adjacent slabs.}
Let $G\in\mathcal{A}_\Xi[c,c']$ and $G'\in\mathcal{A}_\Xi[c',c'']$.
The product $(G,G')\mapsto G\cdot G'$ is bilinear; so, for the first inclusion, it suffices to consider products $H=\Phi\cdot G_1\cdots G_r$ of $\mathcal{A}(\{c<\ell<c'\})$ and $H'=\Phi'\cdot G'_1\cdots G'_s$ of $\mathcal{A}(\{c'<\ell<c''\})$ that embed in $\Gamma_\parallel$, with boxes $B_j$ and $B'_k$.
The factors $G_j$ and $G'_k$ are even and commute; thus $H\cdot H'=(\Phi\cdot\Phi')\cdot G_1\cdots G_r\cdot G'_1\cdots G'_s$.
The open slabs $\{c<\ell<c'\}$ and $\{c'<\ell<c''\}$ are disjoint, contained in $\{c<\ell<c''\}$, and invariant under the translations by $\Gamma_\parallel$, since $\ell(g)=0$ for $g\in\Gamma_\parallel$.
Hence the $r+s$ boxes are pairwise disjoint and contained in $\{c<\ell<c''\}$, each of them embeds in $\Gamma_\parallel$, and the sets $B_j+g$, $B'_k+g$, for all $j$, $k$ and all $g\in\Gamma_\parallel$, are pairwise disjoint.
Moreover $\Phi\cdot\Phi'$ is a bounded Borel function of finitely many $F_\sigma$ with $\sigma$ contained in $\{c<\ell<c''\}$.
So $H\cdot H'$ is a product of $\mathcal{A}(\{c<\ell<c''\})$ that embeds in $\Gamma_\parallel$, and $G\cdot G'\in\mathcal{A}_\Xi[c,c'']$.

Let $\varpi\in\mathfrak{C}[c,c'']$, and let $\hat A$ be a cochain as in Lemma~\ref{lem:global-lift}\eqref{it:lift-slab} for $\varpi$ and the slab $\{c\le\ell\le c''\}$, extended by $0$.
By Lemma~\ref{lem:cells-mirrors}\eqref{it:cell-partition}, $\mathcal{B}^3(c,c')$ and $\mathcal{B}^3(c',c'')$ are contained in $\mathcal{B}^3(c,c'')$; so the restrictions $\varpi|_{[c,c']}$ and $\varpi|_{[c',c'']}$ belong to $\mathfrak{C}[c,c']$ and $\mathfrak{C}[c',c'']$.
The plaquettes of $\Z^3[c,c']$ and their links are cells of $\Z^3[c,c'']$; hence the restriction of $\hat A$ to the links of $\Z^3[c,c']$, extended by $0$, is a cochain as in Lemma~\ref{lem:global-lift}\eqref{it:lift-slab} for $\varpi|_{[c,c']}$, and the restriction to the links of $\Z^3[c',c'']$, extended by $0$, is one for $\varpi|_{[c',c'']}$.
Now $G(\hat A)$ involves only links in $\{c<\ell<c'\}$ and $G'(\hat A)$ only links in $\{c'<\ell<c''\}$; so they are the values of $G$ and $G'$ at these two restricted cochains.
The quotient maps of the intervals $(c,c')$ and $(c',c'')$ are the restrictions of that of $(c,c'')$, all three sending each generator to the generator of the same name, and $r_\Xi$ is multiplicative; therefore
\[
(G\cdot G')(\varpi)=r_\Xi\left(G(\hat A)\cdot G'(\hat A)\right)=r_\Xi\left(G(\hat A)\right)\cdot r_\Xi\left(G'(\hat A)\right)=G(\varpi|_{[c,c']})\cdot G'(\varpi|_{[c',c'']}).
\]

\emph{Reflections and translations of maps.}
For a mirror $\Pi$, an element $G$ of $\mathcal{A}(\R^3)$ and a real one-cochain $\hat A$ on $\Z^3$, \eqref{eq:setting-reflection-coefficients} and Definition~\ref{def:product-maps} give $(\Theta_\Pi G)(\hat A)=\Theta_\Pi\left(G(\theta_\Pi^*\hat A)\right)$, with the antilinear map $\Theta_\Pi$ of \eqref{eq:setting-reflection-fields} on the right.
This applies in particular to each factor $G_j\in\mathcal{G}^{\mathrm{ev}}(B_j)$ of a product, read as the map $\hat A\mapsto G_j\left(\hat A|_{B_j^{(1)}}\right)$, which is itself an element of $\mathcal{A}(\R^3)$.
For $\mu\in\{1,2,3\}$, Definition~\ref{def:mirrors}\eqref{it:mirror-product} gives $\theta_{(x_\mu,1/2)}\theta_{(x_\mu,0)}=\tau_{e_\mu}$, and hence $\theta_{(x_\mu,0)}\theta_{(x_\mu,1/2)}=\tau_{-e_\mu}$.
For $v\in\Z^3$ one may therefore write $\tau_v=\theta_{\Pi^{(1)}}\cdots\theta_{\Pi^{(2n)}}$ with a word in which each pair $(\Pi^{(2i-1)},\Pi^{(2i)})$ is $\left((x_\mu,1/2),(x_\mu,0)\right)$ or $\left((x_\mu,0),(x_\mu,1/2)\right)$ for some $\mu$.
The two mirrors of such a pair have the same upward unit normal $e_\mu$, and $\gamma_{e_\mu}^2=1$.
By \eqref{eq:setting-two-reflections}, $\Theta_{\Pi^{(2i-1)}}\Theta_{\Pi^{(2i)}}$ maps $\psi(x)$ to $\psi(\theta_{\Pi^{(2i-1)}}\theta_{\Pi^{(2i)}}x)$ and $\bar\psi(x)$ to $\bar\psi(\theta_{\Pi^{(2i-1)}}\theta_{\Pi^{(2i)}}x)$.
Hence $\Theta_{\Pi^{(1)}}\cdots\Theta_{\Pi^{(2n)}}$ is the automorphism $\tau_v$ of $\Gr(\mathbb{X})$.
For even $G$, \eqref{eq:setting-group-action} gives $G^{\tau_v}=\Theta_{\Pi^{(1)}}\cdots\Theta_{\Pi^{(2n)}}G$.
Applying the formula for $\Theta_\Pi$ $2n$ times and using $(gg')^*=g'^*g^*$ (Definition~\ref{def:cochains}), one finds
\[
G^{\tau_v}(\hat A)=\Theta_{\Pi^{(1)}}\cdots\Theta_{\Pi^{(2n)}}\left(G(\theta_{\Pi^{(2n)}}^*\cdots\theta_{\Pi^{(1)}}^*\hat A)\right)=\tau_v\left(G(\tau_v^*\hat A)\right);
\]
the complex conjugations cancel in pairs, and $G^{\tau_v}(\hat A)=\sum_\Psi a_\Psi(\tau_v^*\hat A)\cdot\tau_v(\Psi)$ if $G(\hat A)=\sum_\Psi a_\Psi(\hat A)\cdot\Psi$.

\emph{Stability of the classes of slab products.}
The map $G\mapsto G^{\tau_v}$ is linear, and $G\mapsto\Theta_{c'}G$, $G\mapsto\Theta_cG$ are antilinear.
It therefore suffices to prove the three remaining inclusions for a product $H=\Phi\cdot G_1\cdots G_r$ of $\mathcal{A}(\{c<\ell<c'\})$ that embeds in $\Gamma_\parallel$, where $\Phi$ is a bounded Borel function of the $F_\sigma$ with $\sigma$ in a finite set $S$ of plaquettes contained in $\{c<\ell<c'\}$.

Let $\theta$ be $\theta_{c'}$ or $\theta_c$, let $\Theta$ be $\Theta_{c'}$, resp. $\Theta_c$, and let $(c_1,c_2)$ be $(c',2c'-c)$, resp. $(2c-c',c)$; these are elements of $\Adm(\ell)$, which is a group.
The map $\Theta$ is antilinear and reverses products, $\theta^*$ commutes with $d$, and the even elements $G_j(\cdot)$ commute with each other; so the formula for $\Theta$ gives
\[
\Theta H=\Phi^{\theta}\cdot(\Theta G_1)\cdots(\Theta G_r),\qquad\Phi^{\theta}(F):=\overline{\Phi(\theta^*F)}.
\]
By Lemma~\ref{lem:reflection-classes}, $\Theta G_j\in\mathcal{G}^{\mathrm{ev}}(\theta B_j)$.
Here $\theta B_j=\theta(y_j)+[-r_j,r_j]^3$ is again a box, because the linear part $\vec\theta$ of $\theta$ is a signed permutation matrix and maps every cube centred at $0$ with edges parallel to the axes onto itself.
Since $(\theta^*F)_\sigma=\pm F_{\theta\sigma}$, the function $\Phi^{\theta}$ is a bounded Borel function of the $F_{\theta\sigma}$, $\sigma\in S$.
By Definition~\ref{def:mirrors}\eqref{it:mirror-height}, $\ell\circ\theta_{c'}=2c'-\ell$ and $\ell\circ\theta_c=2c-\ell$; so $\theta$ maps $\{c<\ell<c'\}$ onto $\{c_1<\ell<c_2\}$.
The boxes $\theta B_j$ are therefore pairwise disjoint and contained in $\{c_1<\ell<c_2\}$, as are the plaquettes $\theta\sigma$, $\sigma\in S$, and $\Theta H\in\mathcal{A}(\{c_1<\ell<c_2\})$.
For $g\in\Gamma_\parallel$ one has $\ell(g)=0$, hence $\theta(x+g)=\theta(x)+g$ by the formula for $\theta_{\ell,c}$ in Definition~\ref{def:mirrors}, that is, $\vec\theta g=g$.
Together with $\vec\theta[-1,1]^3=[-1,1]^3$ this gives
\[
\theta\left(B_j+[-1,1]^3+g\right)=\theta B_j+[-1,1]^3+g,\qquad
\theta\left(B_j+g\right)=\theta B_j+g .
\]
As $\theta$ is a bijection of $\R^3$, the sets $\theta B_j+[-1,1]^3+g$, $g\in\Gamma_\parallel$, are pairwise disjoint for each $j$, and the sets $\theta B_j+g$, for all $j$ and all $g\in\Gamma_\parallel$, are pairwise disjoint.
So $\Theta H$ embeds in $\Gamma_\parallel$, and $\Theta H\in\mathcal{A}_\Xi[c_1,c_2]$.

For the translation, the formula for $G^{\tau_v}$, applied to $H$ and to each $G_j$, the relation $d\tau_v^*=\tau_v^*d$ and the multiplicativity of the automorphism $\tau_v$ give
\[
H^{\tau_v}=\Phi^{\tau_v}\cdot G_1^{\tau_v}\cdots G_r^{\tau_v},\qquad\Phi^{\tau_v}(F):=\Phi(\tau_v^*F),
\]
and $\Phi^{\tau_v}$ is a bounded Borel function of the $F_{\sigma+v}$, $\sigma\in S$.
Each letter $\Theta_{\Pi^{(i)}}$ of the word above maps $\mathcal{G}^{\mathrm{ev}}(B)$ onto $\mathcal{G}^{\mathrm{ev}}(\theta_{\Pi^{(i)}}B)$ for every box $B$, by Lemma~\ref{lem:reflection-classes}; hence $G_j^{\tau_v}\in\mathcal{G}^{\mathrm{ev}}(\theta_{\Pi^{(1)}}\cdots\theta_{\Pi^{(2n)}}B_j)=\mathcal{G}^{\mathrm{ev}}(B_j+v)$.
The translation by $v$ maps $\{c<\ell<c'\}$ onto $\{c+\ell(v)<\ell<c'+\ell(v)\}$, and
\[
(B_j+v)+[-1,1]^3+g=\left(B_j+[-1,1]^3+g\right)+v,\qquad(B_j+v)+g=(B_j+g)+v .
\]
So $H^{\tau_v}$ is a product of $\mathcal{A}(\{c+\ell(v)<\ell<c'+\ell(v)\})$ that embeds in $\Gamma_\parallel$; as $\ell(v)\in\Z\subset\Adm(\ell)$, this means $H^{\tau_v}\in\mathcal{A}_\Xi[c+\ell(v),c'+\ell(v)]$.

\emph{The identities for reflections and translations.}
Let $\theta$, $\Theta$, $c_1$, $c_2$ be as above, let $\varpi\in\mathfrak{C}[c_1,c_2]$, and let $\hat A$ be a cochain as in Lemma~\ref{lem:global-lift}\eqref{it:lift-slab} for $\varpi$ and the slab $\{c_1\le\ell\le c_2\}$, extended by $0$.
The reflection $\theta$ commutes with the translations by $\Gamma_\parallel$, maps the cells of $\Z^3[c,c']$ onto those of $\Z^3[c_1,c_2]$, and maps the planes $\{\ell=c\}$ and $\{\ell=c'\}$ onto $\{\ell=c_2\}$ and $\{\ell=c_1\}$.
Hence it maps $\mathcal{B}^3(c,c')$ onto $\mathcal{B}^3(c_1,c_2)$, and, $\theta^*$ commuting with $d$, one has $\theta^*\varpi\in\mathfrak{C}[c,c']$ and $F(\theta^*\varpi)=\theta^*F(\varpi)$.
The cochain $\theta^*\hat A$ vanishes on the links that are not contained in $\{c\le\ell\le c'\}$.
On $\Z^3[c,c']$ one has $d(\theta^*\hat A)=\theta^*d\hat A$, the periodic extension of $F(\theta^*\varpi)$, and, with the sign $\pm1$ of Definition~\ref{def:cochains},
\[
\exp\left(iq(\theta^*\hat A)(b)\right)=\exp\left(\pm iq\hat A(\theta b)\right)=\exp\left(\pm iq\widetilde A(\theta b)\right)=\exp\left(iq(\theta^*\widetilde A)(b)\right),
\]
where $\widetilde A$ is the periodic extension of the first component of $\varpi$, so that $\theta^*\widetilde A$ is the periodic extension of the first component of $\theta^*\varpi$.
So $\theta^*\hat A$ is a cochain as in Lemma~\ref{lem:global-lift}\eqref{it:lift-slab} for $\theta^*\varpi$, extended by $0$.
The function of $\Theta G\in\mathcal{A}_\Xi[c_1,c_2]$ at $\varpi$ may be computed with the cochain $\hat A$; the formula $(\Theta G)(\hat A)=\Theta\left(G(\theta^*\hat A)\right)$ and the relation $r_\Xi\Theta=\Theta r_\Xi$ (Definition~\ref{def:product-maps}), with $\Theta$ acting on the Grassmann algebras of the cylinder as in Lemma~\ref{lem:slab-automorphisms}, then give
\[
(\Theta G)(\varpi)=r_\Xi\left(\Theta\left(G(\theta^*\hat A)\right)\right)=\Theta\left(r_\Xi\left(G(\theta^*\hat A)\right)\right)=\Theta\left(G(\theta^*\varpi)\right),
\]
which is the third identity of \eqref{it:sum-slab} for $\theta=\theta_{c'}$ and the fourth for $\theta=\theta_c$.

For the second identity let $\varpi\in\mathfrak{C}[c+\ell(v),c'+\ell(v)]$, and let $\hat A$ be a cochain as in Lemma~\ref{lem:global-lift}\eqref{it:lift-slab} for $\varpi$, extended by $0$.
The translation $\tau_v$ commutes with the translations by $\Gamma_\parallel$ and maps $\Z^3[c,c']$ onto $\Z^3[c+\ell(v),c'+\ell(v)]$, each plane onto the corresponding plane; $\tau_v^*$ commutes with $d$, and $(\tau_v^*\hat A)(b)=\hat A(b+v)$.
Hence $\tau_v^*\varpi\in\mathfrak{C}[c,c']$, $F(\tau_v^*\varpi)=\tau_v^*F(\varpi)$, and $\tau_v^*\hat A$ is a cochain as in Lemma~\ref{lem:global-lift}\eqref{it:lift-slab} for $\tau_v^*\varpi$, extended by $0$: its coboundary is $\tau_v^*d\hat A$, and its phase at a link $b$ is the phase of $\hat A$ at $b+v$, which is that of the first component of $\varpi$ at $b+v$.
Computing the function of $G^{\tau_v}\in\mathcal{A}_\Xi[c+\ell(v),c'+\ell(v)]$ at $\varpi$ with the cochain $\hat A$, and using the formula for $G^{\tau_v}$ and the relation $r_\Xi\tau_v=\tau_vr_\Xi$ (Definition~\ref{def:product-maps}), one finds
\[
G^{\tau_v}(\varpi)=r_\Xi\left(\tau_v\left(G(\tau_v^*\hat A)\right)\right)=\tau_v\left(r_\Xi\left(G(\tau_v^*\hat A)\right)\right)=\tau_v\left(G(\tau_v^*\varpi)\right).
\]
In the sense of Definition~\ref{def:slab-integrands}, the right sides of the first three identities are the product of $G(\cdot)$ and $G'(\cdot)$, the translate of $G(\cdot)$ by $\tau_v$ and the reflection of $G(\cdot)$ by $\Theta_{c'}$; so the map $G\mapsto G(\cdot)$ carries $G\cdot G'$, $G^{\tau_v}$ and $\Theta_{c'}G$ to the product, the translate and the reflection of the images.
\end{proof}

\subsection{Berezin integrals}\label{sec:berezin-integrals}

\begin{lemma}[Poisson summation and kernels]\label{lem:poisson-kernels}
Let $q>0$ be the lattice charge (Definition~\ref{def:configurations}).
\begin{enumerate}
\item\label{it:poisson-summation}
For all real numbers $y$ and $s$,
\[
\sum_{k\in\Z}\exp\left(-\frac12\left(y+\frac{2\pi}{q}\,k\right)^2+iks\right)
=\frac{q}{\sqrt{2\pi}}\sum_{k'\in\Z}\exp\left(-\frac12\,\varkappa_{k'}^2-i\varkappa_{k'}\,y\right),
\qquad
\varkappa_{k'}:=\frac{q}{2\pi}\cdot\left(s-2\pi k'\right).
\]
\item\label{it:gaussian-sums}
For $\varepsilon>0$ the number
\[
C_{q,\varepsilon}:=\sup_{y\in\R}\ \sum_{k\in\Z}\exp\left(-\varepsilon\left(y+\frac{2\pi}{q}\,k\right)^2\right)
\]
satisfies
\[
C_{q,\varepsilon}\le2+\frac{q}{2\pi}\cdot\sqrt{\frac{\pi}{\varepsilon}} .
\]
For $\varepsilon=\tfrac12$ it is the constant $C_q$ of Lemma~\ref{lem:bounded-integrands}.
\item\label{it:hilbert-schmidt-kernels}
Let $(\mathcal{Y},\mu)$, $(\mathcal{Y}',\mu')$ and $(\mathcal{Y}'',\mu'')$ be $\sigma$-finite measure spaces whose spaces of square integrable functions are separable,
and write $\xi$, $\xi'$, $\xi''$ for their points and $L^2(\mu)$, $L^2(\mu')$, $L^2(\mu'')$ for these spaces.
Every $k\in L^2(\mathcal{Y}'\times\mathcal{Y},\mu'\otimes\mu)$ is the kernel of the Hilbert--Schmidt operator
\[
L^2(\mu)\to L^2(\mu'),\qquad f\mapsto\int\mathrm{d}\mu(\xi)\,k(\,\cdot\,,\xi)\,f(\xi),
\]
and its Hilbert--Schmidt norm is $\|k\|_2$.
The adjoint of this operator has the kernel $(\xi,\xi')\mapsto\overline{k(\xi',\xi)}$.
For $k_1\in L^2(\mathcal{Y}'\times\mathcal{Y},\mu'\otimes\mu)$ and $k_2\in L^2(\mathcal{Y}''\times\mathcal{Y}',\mu''\otimes\mu')$,
the operator with the kernel $k_2$, composed with the operator with the kernel $k_1$, has the kernel
\[
(\xi'',\xi)\mapsto\int\mathrm{d}\mu'(\xi')\,k_2(\xi'',\xi')\,k_1(\xi',\xi).
\]
Let $X\colon L^2(\mu)\to L^2(\mu')$ and $X'\colon L^2(\mu')\to L^2(\mu)$ be Hilbert--Schmidt operators with the kernels $k_X$ and $k_{X'}$.
Then $X'X$ is of trace class, and
\[
\Tr\left(X'X\right)=\left\langle X'^{*},X\right\rangle_{\mathrm{HS}}
=\int\mathrm{d}\mu(\xi)\int\mathrm{d}\mu'(\xi')\,k_{X'}(\xi,\xi')\,k_X(\xi',\xi),
\]
where the Hilbert--Schmidt inner product $\langle\,\cdot\,,\,\cdot\,\rangle_{\mathrm{HS}}$ is antilinear in its first argument;
the double integral converges absolutely.
The same assertions hold when the functions on $\mathcal{Y}$, $\mathcal{Y}'$, $\mathcal{Y}''$ take their values in finite-dimensional Hilbert spaces $\mathcal{H}$, $\mathcal{H}'$, $\mathcal{H}''$
and the kernels take their values in the spaces of linear maps between them:
then $\|k\|_2$ is formed with the Hilbert--Schmidt norm of the values,
the kernel of the adjoint is $(\xi,\xi')\mapsto k(\xi',\xi)^{*}$,
and
\[
\Tr\left(X'X\right)=\int\mathrm{d}\mu(\xi)\int\mathrm{d}\mu'(\xi')\,\operatorname{tr}\left(k_{X'}(\xi,\xi')\,k_X(\xi',\xi)\right).
\]
\end{enumerate}
\end{lemma}

\begin{proof}
\eqref{it:poisson-summation}
For fixed $y$ and $s$ the function
\[
f(x):=\exp\left(-\frac12\left(y+\frac{2\pi}{q}\,x\right)^2+isx\right),\qquad x\in\R,
\]
is a Schwartz function.
Poisson's summation formula gives
\[
\sum_{k\in\Z}f(k)=\sum_{k'\in\Z}\hat f(2\pi k'),\qquad
\hat f(2\pi k'):=\int_{\R}\mathrm{d}x\,f(x)\exp\left(-2\pi ik'x\right).
\]
In the integral substitute $u=y+\frac{2\pi}{q}\,x$.
Then $\mathrm{d}x=\frac{q}{2\pi}\,\mathrm{d}u$ and $(s-2\pi k')\,x=\varkappa_{k'}\cdot(u-y)$, so that
\[
\hat f(2\pi k')
=\frac{q}{2\pi}\cdot\exp\left(-i\varkappa_{k'}\,y\right)\int_{\R}\mathrm{d}u\,\exp\left(-\frac{u^2}{2}+i\varkappa_{k'}\,u\right)
=\frac{q}{2\pi}\cdot\sqrt{2\pi}\cdot\exp\left(-\frac12\,\varkappa_{k'}^2-i\varkappa_{k'}\,y\right),
\]
the last integral being the Fourier transform of the Gaussian.
Since $\frac{q}{2\pi}\cdot\sqrt{2\pi}=\frac{q}{\sqrt{2\pi}}$, the identity follows.

\eqref{it:gaussian-sums}
Fix $y\in\R$ and put $t_k:=y+\frac{2\pi}{q}\,k$ for $k\in\Z$.
Consecutive points $t_k$ are at the distance $\frac{2\pi}{q}$;
hence exactly one of them lies in $\left[0,\frac{2\pi}{q}\right)$ and exactly one in $\left[-\frac{2\pi}{q},0\right)$,
and the terms of these two points are at most $1$ each.
Every other point satisfies $t_k\ge\frac{2\pi}{q}$ or $t_k<-\frac{2\pi}{q}$.
The function $u\mapsto\exp\left(-\varepsilon u^2\right)$ decreases in $|u|$, so that
\[
\begin{aligned}
\exp\left(-\varepsilon t_k^2\right)&\le\frac{q}{2\pi}\int_{\left[t_k-\frac{2\pi}{q},\,t_k\right]}\mathrm{d}u\,\exp\left(-\varepsilon u^2\right)&&\text{if }t_k\ge\frac{2\pi}{q},\\
\exp\left(-\varepsilon t_k^2\right)&\le\frac{q}{2\pi}\int_{\left[t_k,\,t_k+\frac{2\pi}{q}\right]}\mathrm{d}u\,\exp\left(-\varepsilon u^2\right)&&\text{if }t_k<-\frac{2\pi}{q}.
\end{aligned}
\]
The intervals of integration lie in $[0,\infty)$, resp. in $(-\infty,0)$, and two of them have at most an end point in common.
So the terms of all other points add up to at most
\[
\frac{q}{2\pi}\int_{\R}\mathrm{d}u\,\exp\left(-\varepsilon u^2\right)=\frac{q}{2\pi}\cdot\sqrt{\frac{\pi}{\varepsilon}} .
\]
The bound does not depend on $y$, and the estimate for $C_{q,\varepsilon}$ follows.

\eqref{it:hilbert-schmidt-kernels}
The Hilbert--Schmidt operators between two separable Hilbert spaces form a Hilbert space for the inner product $\langle X_1,X_2\rangle_{\mathrm{HS}}=\Tr\left(X_1^{*}X_2\right)$;
the product of two Hilbert--Schmidt operators is of trace class, with $\Tr\left(X'X\right)=\left\langle X'^{*},X\right\rangle_{\mathrm{HS}}$;
and on a space of square integrable functions the integral operators with square integrable kernels are Hilbert--Schmidt operators,
the Hilbert--Schmidt norm being the $L^2$-norm of the kernel \cite[Theorems~VI.22 and~VI.23]{ReedSimon1980}.
For operators between $L^2(\mu)$ and $L^2(\mu')$ these facts are applied on the orthogonal sum of the two spaces,
which is the space of square integrable functions on the disjoint union of $\mathcal{Y}$ and $\mathcal{Y}'$.
By polarisation, the linear map from kernels to operators preserves inner products:
\[
\left\langle X_1,X_2\right\rangle_{\mathrm{HS}}=\int\mathrm{d}\mu(\xi)\int\mathrm{d}\mu'(\xi')\,\overline{k_{X_1}(\xi',\xi)}\,k_{X_2}(\xi',\xi)
\]
for operators $X_1,X_2\colon L^2(\mu)\to L^2(\mu')$ with the kernels $k_{X_1},k_{X_2}$.
For $f\in L^2(\mu)$ and $g\in L^2(\mu')$ the Cauchy--Schwarz inequality on $\mathcal{Y}'\times\mathcal{Y}$ gives
\[
\int\mathrm{d}\mu'(\xi')\int\mathrm{d}\mu(\xi)\,|g(\xi')|\cdot|k(\xi',\xi)|\cdot|f(\xi)|\le\|g\|_2\cdot\|k\|_2\cdot\|f\|_2 ;
\]
so Fubini's theorem applies to the scalar product of $g$ with the image of $f$,
and the exchange of the two integrations shows that the adjoint has the kernel $(\xi,\xi')\mapsto\overline{k(\xi',\xi)}$.
For the composition, the Cauchy--Schwarz inequality on $\mathcal{Y}'$ gives
\[
\int\mathrm{d}\mu'(\xi')\,|k_2(\xi'',\xi')|\cdot|k_1(\xi',\xi)|\le\|k_2(\xi'',\cdot\,)\|_2\cdot\|k_1(\,\cdot\,,\xi)\|_2,
\]
and the right-hand side is square integrable in $(\xi'',\xi)$, the integral of its square being $\|k_2\|_2^2\cdot\|k_1\|_2^2$.
Hence the integral that defines the kernel of the composition converges absolutely for almost every $(\xi'',\xi)$,
this kernel is square integrable,
and Fubini's theorem identifies the integral operator with this kernel with the composition of the two operators.
Inserting the kernel $(\xi',\xi)\mapsto\overline{k_{X'}(\xi,\xi')}$ of $X'^{*}$ into the formula for $\left\langle X'^{*},X\right\rangle_{\mathrm{HS}}$ gives the double integral for $\Tr\left(X'X\right)$;
it converges absolutely by the Cauchy--Schwarz inequality on $\mathcal{Y}\times\mathcal{Y}'$.
Let finally the functions take values in finite-dimensional Hilbert spaces, and choose orthonormal bases of these spaces.
Then the space of square integrable functions on $\mathcal{Y}$ with values in $\mathcal{H}$
becomes the space of square integrable scalar functions on $\mathcal{Y}\times\{1,\dots,\dim\mathcal{H}\}$
for the product of $\mu$ with the counting measure, and likewise for $\mathcal{Y}'$ and $\mathcal{Y}''$;
a kernel with values in linear maps becomes the scalar kernel formed by the matrix entries of its values.
The assertions for scalar kernels, applied to these kernels, are the assertions for kernels with values in linear maps:
the sum of the squared moduli of the entries of a matrix is the square of its Hilbert--Schmidt norm,
the conjugate transpose of a matrix represents its adjoint,
and the sum over the diagonal entries is the trace $\operatorname{tr}$.
\end{proof}

\begin{definition}[Left derivatives]\label{def:left-derivatives}
Let $Y\subset\mathbb{X}$.
An element of $\Gr(Y)$ \emph{does not contain} a generator if it lies in the subalgebra generated by the other generators.
Let $x\in Y$ and $\alpha\in\{1,2,3,4\}$.
Every $G\in\Gr(Y)$ can be written in exactly one way as $G=G_0+\psi_\alpha(x)\,G_1$, where $G_0$ and $G_1$ do not contain $\psi_\alpha(x)$.
The \emph{left derivative} with respect to $\psi_\alpha(x)$ is the linear map
\[
\frac{\partial}{\partial\psi_\alpha(x)}\colon\Gr(Y)\to\Gr(Y),\qquad G_0+\psi_\alpha(x)\,G_1\mapsto G_1 ;
\]
it annihilates the elements that do not contain $\psi_\alpha(x)$, and it maps $\psi_\alpha(x)\,G$ to $G$ for every such element $G$.
The left derivative $\partial/\partial\bar\psi_\alpha(x)$ is given by the same rule, with $\bar\psi_\alpha(x)$ in place of $\psi_\alpha(x)$.
Put
\[
\partial_{x,\alpha}:=\frac{\partial}{\partial\bar\psi_\alpha(x)}\,\frac{\partial}{\partial\psi_\alpha(x)},
\qquad
\Psi^{\mathrm{top}}(x):=\prod_{\alpha=1}^{4}\psi_\alpha(x)\,\bar\psi_\alpha(x),
\]
and, for finite sets $Y'\subset Y$,
\[
\partial_{Y'}:=\prod_{x\in Y'}\prod_{\alpha=1}^{4}\partial_{x,\alpha},
\qquad
\Psi^{\mathrm{top}}(Y'):=\prod_{x\in Y'}\Psi^{\mathrm{top}}(x),
\]
the factors of $\partial_{Y'}$ being taken in a fixed order.
The factors $\psi_\alpha(x)\,\bar\psi_\alpha(x)$ and $\Psi^{\mathrm{top}}(x)$ are even, so that the order of the factors in the products that define $\Psi^{\mathrm{top}}$ is immaterial.
If $Y$ is finite, the elements of $\Gr(Y)$ of the top degree $8\#|Y|$ are the multiples of $\Psi^{\mathrm{top}}(Y)$,
and the \emph{coefficient of $\Psi^{\mathrm{top}}(Y)$} in $G\in\Gr(Y)$ is the number $a$ for which $a\cdot\Psi^{\mathrm{top}}(Y)$ is the component of $G$ of degree $8\#|Y|$.

Let $\mathbb{T}$ be a torus, or let $\Xi$ be a cross-section for a family $\ell$ (Definition~\ref{def:slabs}) and $\mathbb{I}$ an open interval.
The generators at the points of a set of representatives of $\mathbb{X}_{\mathbb{T}}$, resp. of $\mathbb{X}_\Xi(\mathbb{I})$,
are free generators of $\Gr_{\mathbb{T}}$, resp. of $\Gr_\Xi(\mathbb{I})$ (Definitions~\ref{def:tori} and~\ref{def:slab-grassmann}),
and with respect to them left derivatives, the operators $\partial_{x,\alpha}$ and $\partial_{Y'}$, and the elements $\Psi^{\mathrm{top}}$ are defined as above.
For a set $Y$ of sites of $\mathbb{X}_{\mathbb{T}}$, resp. of $\mathbb{X}_\Xi(\mathbb{I})$,
$\Gr(Y)$ denotes the subalgebra generated by the generators at the points of the classes in $Y$.
For finite $Y$, the \emph{Berezin integral} $\int_YG$ of $G\in\Gr(Y)$ is the coefficient of $\Psi^{\mathrm{top}}(Y)$ in $G$,
as it is for finite subsets of $\mathbb{X}$ by Definition~\ref{def:grassmann}.
If $\mathbb{I}$ is bounded, then $\mathbb{X}_\Xi(\mathbb{I})$ is finite, and
\[
\int_{\mathbb{I}}:=\int_{\mathbb{X}_\Xi(\mathbb{I})} .
\]
\end{definition}

In each of the algebras $\Gr(Y)$, $Y\subset\mathbb{X}$, $\Gr_{\mathbb{T}}$ and $\Gr_\Xi(\mathbb{I})$, left derivatives obey the following rules.
Left derivatives with respect to two different generators anticommute.
Indeed, every element is a sum of four terms:
elements that contain neither of the two generators, multiplied from the left by $1$, by the first generator, by the second generator, and by the product of the first and the second generator, in this order;
both compositions of the two derivatives annihilate the first three terms,
and they map the fourth term to one and the same element with opposite signs.
The square of a left derivative vanishes, since the derivative of an element does not contain the generator.
Let $G$ be homogeneous and not contain $\psi_\alpha(x)$, and let $G'=G'_0+\psi_\alpha(x)\,G'_1$ as in Definition~\ref{def:left-derivatives}.
Then $GG'=GG'_0+(-1)^{\deg G}\,\psi_\alpha(x)\,GG'_1$, where $GG'_0$ and $GG'_1$ do not contain $\psi_\alpha(x)$, and therefore
\[
\frac{\partial}{\partial\psi_\alpha(x)}\left(GG'\right)=(-1)^{\deg G}\,G\,\frac{\partial}{\partial\psi_\alpha(x)}\,G' ;
\]
the same rule holds with $\bar\psi_\alpha(x)$ in place of $\psi_\alpha(x)$.
Each $\partial_{x,\alpha}$ is a product of two left derivatives, so it commutes with every left derivative;
consequently the operators $\partial_{x,\alpha}$ commute with each other, and the order of the factors of $\partial_{Y'}$ is immaterial.
If $G$ contains neither $\psi_\alpha(x)$ nor $\bar\psi_\alpha(x)$, then
\[
\partial_{x,\alpha}\left[\psi_\alpha(x)\,\bar\psi_\alpha(x)\,G\right]=G,
\qquad\text{in particular}\qquad
\partial_{x,\alpha}\left[\psi_\alpha(x)\,\bar\psi_\alpha(x)\right]=1,
\]
whereas $\partial_{x,\alpha}$ annihilates every monomial that lacks $\psi_\alpha(x)$ or $\bar\psi_\alpha(x)$.
Now let $Y$ be a finite set of sites.
A monomial in the generators at the points of $Y$ that is not $\pm\Psi^{\mathrm{top}}(Y)$ lacks one of these generators,
and the factor of $\partial_Y$ that belongs to this generator, applied first, annihilates it;
on the other hand $\partial_Y\Psi^{\mathrm{top}}(Y)=1$, the pairs $\psi_\alpha(x)\,\bar\psi_\alpha(x)$ being removed one after the other.
Hence $\partial_YG$ is the coefficient of $\Psi^{\mathrm{top}}(Y)$ in $G$, and
\[
\int_Y=\partial_Y=\prod_{x\in Y}\prod_{\alpha=1}^{4}\partial_{x,\alpha}\qquad\text{on }\Gr(Y);
\]
for $Y\subset\mathbb{X}$ the normalisation~\eqref{eq:setting-berezin} of the Berezin integral is the identity $\partial_{x,\alpha}\left[\psi_\alpha(x)\,\bar\psi_\alpha(x)\right]=1$.
In particular $\int_YG=0$ if $G$ is odd, or if all homogeneous components of $G$ have degrees less than $8\#|Y|$.

On a torus $\mathbb{T}=(\Gamma,\chi)$ the generators at $x+g$, $g\in\Gamma$, are $\chi(g)$ times those at $x$.
Let a representative $x$ be replaced by $x+g$.
Each of the eight generators $\psi_\alpha(x)$, $\bar\psi_\alpha(x)$, $\alpha=1,\dots,4$, is then multiplied by $\chi(g)$,
the generators at the other representatives stay as they are,
and for every generator the subalgebra generated by the remaining generators is the same as before.
If $G=G_0+\psi_\alpha(x)\,G_1$ as in Definition~\ref{def:left-derivatives}, then $G=G_0+\psi_\alpha(x+g)\,\chi(g)\,G_1$, since $\chi(g)^2=1$,
and $G_0$, $G_1$ do not contain $\psi_\alpha(x+g)$;
so the left derivative with respect to $\psi_\alpha(x+g)$ is $\chi(g)$ times the one with respect to $\psi_\alpha(x)$, and likewise for $\bar\psi_\alpha$.
Thus each of the eight left derivatives at $x$ is multiplied by $\chi(g)$, while the left derivatives at the other representatives do not change.
It follows that $\partial_{x,\alpha}$, a product of two left derivatives at $x$, does not change; moreover
\[
\Psi^{\mathrm{top}}(x+g)=\chi(g)^8\cdot\Psi^{\mathrm{top}}(x)=\Psi^{\mathrm{top}}(x).
\]
Hence $\partial_Y$, $\Psi^{\mathrm{top}}(Y)$ and $\int_Y$ depend only on the set $Y$ of classes,
and $\int_{\mathbb{X}_{\mathbb{T}}}$ is the integral $\int_{\mathbb{T}}$ of Definition~\ref{def:tori}.
The same holds on $\Gr_\Xi(\mathbb{I})$, with the character $\chi_\parallel$ of $\Xi$ in place of $\chi$;
in particular $\int_{\mathbb{I}}$ does not depend on the choice of the representatives.

\begin{proposition}[Berezin calculus and integrability]\label{prop:berezin}
\begin{enumerate}
\item\label{it:berezin-factorisation}
Let $Y_1,\dots,Y_r$ be pairwise disjoint finite sets of sites,
either all contained in $\mathbb{X}$, or all contained in $\mathbb{X}_{\mathbb{T}}$ for a torus $\mathbb{T}$,
or all contained in $\mathbb{X}_\Xi(\mathbb{I})$ for a cross-section $\Xi$ for a family $\ell$ and an open interval $\mathbb{I}$,
and let $G_k\in\Gr(Y_k)$, $k=1,\dots,r$, be arbitrary elements.
Then
\[
\int_{Y_1\cup\dots\cup Y_r}G_1\,G_2\cdots G_r=\prod_{k=1}^{r}\int_{Y_k}G_k .
\]
No sign occurs, whatever the parities of the $G_k$;
if one of the $G_k$ is odd, both sides vanish.
\item\label{it:reflected-integral}
Let $\Pi=(\ell,c)$ be a mirror, and $\theta:=\theta_\Pi$.
For every finite set $Y\subset\mathbb{X}$ and every $G\in\Gr(Y)$ one has $\Theta_\Pi(G)\in\Gr(\theta Y)$ and
\[
\int_{\theta Y}\Theta_\Pi(G)=\overline{\int_YG} .
\]
For every torus $\mathbb{T}$ and every $G\in\Gr_{\mathbb{T}}$,
\[
\int_{\theta\mathbb{T}}\Theta_\Pi(G)=\overline{\int_{\mathbb{T}}G},
\]
with the map $\Theta_\Pi\colon\Gr_{\mathbb{T}}\to\Gr_{\theta\mathbb{T}}$ of Lemma~\ref{lem:reflected-algebras}.
For every cross-section $\Xi$ for $\ell$, every bounded open interval $\mathbb{I}$ and every $G\in\Gr_\Xi(\mathbb{I})$,
\[
\int_{2c-\mathbb{I}}\Theta_c(G)=\overline{\int_{\mathbb{I}}G},
\]
with the map $\Theta_c\colon\Gr_\Xi(\mathbb{I})\to\Gr_\Xi(2c-\mathbb{I})$ that $\Theta_\Pi$ induces (Lemma~\ref{lem:slab-automorphisms});
here $2c-\mathbb{I}:=\{2c-t:\ t\in\mathbb{I}\}$.
\item\label{it:torus-integrability}
Let $w_0,\dots,w_3>0$, $M\in\R$, $q>0$ and $-1\le\rho\le1$.
Let $\mathbb{T}=(\Gamma,\chi)$ be a torus, let $\mathbb{T}^{(1)}$ be the set of its links, and fix a monomial basis of $\Gr_{\mathbb{T}}$.
There is a number $C_{\mathbb{T}}<\infty$, depending only on $\mathbb{T}$ and on $\rho$, $w_0,\dots,w_3$ and $M$, with the following property.
Let $G\in\mathcal{I}(\mathbb{T})$, with the constants $C$ and $\varepsilon_G$ of Definition~\ref{def:torus-integrands},
and for a basis monomial $\Psi$ let $a_\Psi(A,z)$ be the coefficient of $\Psi$ in
\[
\exp\left(-\tfrac12\|F\|^2\right)\cdot\exp\left(-S_{\mathbb{T}}(A)\right)\cdot G(A,z),\qquad(A,z)\in\mathfrak{C}(\Gamma).
\]
Then $a_\Psi$ is a Borel function on $\Conf(\Gamma)$, absolutely integrable over $\Conf(\Gamma)$, and
\[
\sum_{\substack{z\in C^2(\Gamma;\Z)\\ dz=0}}\ \int_{[0,2\pi/q)^{\mathbb{T}^{(1)}}}\mathrm{d}A\ \left|a_\Psi(A,z)\right|
\le C_{\mathbb{T}}\cdot C\cdot\left(\frac{2\pi}{q}\right)^{3|\mathbb{T}|}\cdot C_{q,\varepsilon_G}^{3|\mathbb{T}|} .
\]
In particular $Z_{\mathbb{T}}[G]$ is defined for every $G\in\mathcal{I}(\mathbb{T})$ and for every $G\in\mathcal{A}(\mathbb{T})$,
hence for every $G\in\mathcal{A}(\R^3)$ that embeds in $\mathbb{T}$.
Moreover, for every bounded Borel function $\Phi$ on $\Conf(\Gamma)$ and every $G\in\mathcal{I}(\mathbb{T})$ with the constants $C$ and $\varepsilon_G$,
\[
\left|Z_{\mathbb{T}}[\Phi G]\right|\le\sup|\Phi|\cdot C_{\mathbb{T}}\cdot C\cdot\left(\frac{2\pi}{q}\right)^{3|\mathbb{T}|}\cdot C_{q,\varepsilon_G}^{3|\mathbb{T}|},
\]
and the same holds for $G\in\mathcal{A}(\mathbb{T})$, with the constants of the element of $\mathcal{I}(\mathbb{T})$ that $G$ defines (Lemma~\ref{lem:product-integrands}\eqref{it:sum-torus});
so $\Phi\mapsto Z_{\mathbb{T}}[\Phi G]$ is continuous for the supremum norm.
\end{enumerate}
\end{proposition}

\begin{proof}
\eqref{it:berezin-factorisation}
Put $Y:=Y_1\cup\dots\cup Y_r$; all elements below lie in $\Gr(Y)$.
Let $x$ be a site of $Y$ and $\alpha\in\{1,2,3,4\}$, and let $G\in\Gr(Y)$ contain neither $\psi_\alpha(x)$ nor $\bar\psi_\alpha(x)$.
If $G$ is homogeneous, the rule for the derivative of a product, applied to both left derivatives in $\partial_{x,\alpha}$, gives
\[
\partial_{x,\alpha}\left(GG'\right)=(-1)^{2\deg G}\,G\,\partial_{x,\alpha}G'=G\,\partial_{x,\alpha}G'\qquad\left(G'\in\Gr(Y)\right),
\]
and by linearity this identity holds for every $G$ that contains neither generator, homogeneous or not.
The element $G_1\cdots G_{r-1}$ contains no generator at a point of $Y_r$, and $\partial_{Y_r}G_r=\int_{Y_r}G_r$ is a number;
hence
\[
\partial_{Y_r}\left(G_1\cdots G_{r-1}\,G_r\right)=G_1\cdots G_{r-1}\,\partial_{Y_r}G_r=\left(\int_{Y_r}G_r\right)\cdot G_1\cdots G_{r-1}.
\]
The operators $\partial_{x,\alpha}$ commute, so that $\int_Y=\partial_Y=\partial_{Y_1\cup\dots\cup Y_{r-1}}\,\partial_{Y_r}$, and therefore
\[
\int_YG_1\cdots G_r=\left(\int_{Y_r}G_r\right)\cdot\int_{Y_1\cup\dots\cup Y_{r-1}}G_1\cdots G_{r-1}.
\]
Induction on $r$ gives the identity.
The element $\Psi^{\mathrm{top}}(Y_k)$ has the even degree $8\#|Y_k|$, so that $\int_{Y_k}G_k=0$ if $G_k$ is odd;
then the right-hand side vanishes, and with it the left-hand side.

\eqref{it:reflected-integral}
Let $Y$ be a finite subset of $\mathbb{X}$, or a set of representatives of $\mathbb{X}_{\mathbb{T}}$, or a set of representatives of $\mathbb{X}_\Xi(\mathbb{I})$.
In each case the generators at the points of $Y$ are free generators of the algebra in question, $\Gr(Y)$, $\Gr_{\mathbb{T}}$ or $\Gr_\Xi(\mathbb{I})$,
and the integral in question, $\int_Y$, $\int_{\mathbb{T}}$ or $\int_{\mathbb{I}}$, is the coefficient of $\Psi^{\mathrm{top}}(Y)$ (Definition~\ref{def:left-derivatives}).
The image $\theta Y$ is a finite subset of $\mathbb{X}$, resp. a set of representatives of $\mathbb{X}_{\theta\mathbb{T}}$, resp. a set of representatives of $\mathbb{X}_\Xi(2c-\mathbb{I})$.
Indeed, $\theta$ is an element of $W$ and permutes $\mathbb{X}$.
For a torus $\mathbb{T}=(\Gamma,\chi)$ one has $\theta(x+g)=\theta x+\vec\theta g$ for $g\in\Gamma$, where $\vec\theta$ is the linear part of $\theta$,
and $\vec\theta\,\Gamma$ is the lattice of $\theta\mathbb{T}$ (Definition~\ref{def:tori});
so $\theta$ maps the $\Gamma$-orbits in $\mathbb{X}$ bijectively onto the $\vec\theta\,\Gamma$-orbits.
For $g\in\Gamma_\parallel$ one has $\ell(g)=0$, hence $\theta(x+g)=\theta x+g$,
and $\ell(\theta x)=2c-\ell(x)$ by Definition~\ref{def:mirrors}\eqref{it:mirror-height};
so $\theta$ maps $\mathbb{X}(\mathbb{I})$ onto $\mathbb{X}(2c-\mathbb{I})$ and its $\Gamma_\parallel$-orbits onto those of $\mathbb{X}(2c-\mathbb{I})$.
The map $\Theta_\Pi$, which is $\Theta_c$ in the third case, is antilinear from the algebra of $Y$ into $\Gr(\theta Y)$, $\Gr_{\theta\mathbb{T}}$ or $\Gr_\Xi(2c-\mathbb{I})$, respectively;
it reverses products and acts on the generators at the points of $Y$ by~\eqref{eq:setting-reflection-fields}
(Definition~\ref{def:grassmann-reflections}, Lemma~\ref{lem:reflected-algebras} and Lemma~\ref{lem:slab-automorphisms}).
In particular $\Theta_\Pi(G)\in\Gr(\theta Y)$ for $G\in\Gr(Y)$ and finite $Y\subset\mathbb{X}$.

Both sides of each asserted identity are antilinear in $G$.
By~\eqref{eq:setting-reflection-fields} every generator is mapped to a linear combination of generators,
so that homogeneous elements of degree $d$ are mapped to homogeneous elements of degree $d$.
Write $G=a\cdot\Psi^{\mathrm{top}}(Y)+R$, where $a$ is the coefficient of $\Psi^{\mathrm{top}}(Y)$ in $G$, that is, its integral,
and all homogeneous components of $R$ have degrees less than $8\#|Y|$.
Then $\Theta_\Pi(G)=\overline{a}\cdot\Theta_\Pi\left(\Psi^{\mathrm{top}}(Y)\right)+\Theta_\Pi(R)$,
and the integral of $\Theta_\Pi(R)$ over $\theta Y$ vanishes, because $\theta Y$ has $\#|Y|$ points.
It therefore suffices to show
\[
\Theta_\Pi\left(\Psi^{\mathrm{top}}(Y)\right)=\Psi^{\mathrm{top}}(\theta Y),
\]
whose integral over $\theta Y$ is $1$.
Let $n:=\hat n_\ell$ be the upward unit normal of the family $\ell$ of $\Pi$, as in~\eqref{eq:setting-reflection-fields}, and put
\[
\psi^\theta(x):=\gamma_n\,\psi(\theta x),\qquad\bar\psi^\theta(x):=\bar\psi(\theta x)\,\gamma_n\qquad(x\in Y),
\]
a column and a row whose entries are linear combinations of the generators at $\theta x$.
By~\eqref{eq:setting-reflection-fields}, $\Theta_\Pi\left(\bar\psi_\alpha(x)\right)=-\psi^\theta_\alpha(x)$ and $\Theta_\Pi\left(\psi_\alpha(x)\right)=-\bar\psi^\theta_\alpha(x)$.
The map $\Theta_\Pi$ reverses products, and the factors $\Psi^{\mathrm{top}}(x)$ of $\Psi^{\mathrm{top}}(Y)$, as well as the factors $\psi_\alpha(x)\,\bar\psi_\alpha(x)$ of $\Psi^{\mathrm{top}}(x)$, are even;
hence, the two signs cancelling in each pair,
\[
\Theta_\Pi\left(\Psi^{\mathrm{top}}(Y)\right)
=\prod_{x\in Y}\prod_{\alpha=1}^{4}\Theta_\Pi\left(\bar\psi_\alpha(x)\right)\,\Theta_\Pi\left(\psi_\alpha(x)\right)
=\prod_{x\in Y}\prod_{\alpha=1}^{4}\psi^\theta_\alpha(x)\,\bar\psi^\theta_\alpha(x).
\]
Fix $x\in Y$.
The eight elements $\psi^\theta_\alpha(x)$, $\bar\psi^\theta_\alpha(x)$, $\alpha=1,\dots,4$, have degree one and anticommute pairwise.
Reordering $\prod_\alpha\psi^\theta_\alpha(x)\,\bar\psi^\theta_\alpha(x)$ into $\left(\psi^\theta_1(x)\cdots\psi^\theta_4(x)\right)\left(\bar\psi^\theta_1(x)\cdots\bar\psi^\theta_4(x)\right)$
therefore produces a sign $s_4\in\{1,-1\}$ that depends only on the permutation of the eight factors;
the same permutation of the factors of $\Psi^{\mathrm{top}}(\theta x)=\prod_\alpha\psi_\alpha(\theta x)\,\bar\psi_\alpha(\theta x)$ gives
\[
\Psi^{\mathrm{top}}(\theta x)=s_4\cdot\left(\psi_1(\theta x)\cdots\psi_4(\theta x)\right)\left(\bar\psi_1(\theta x)\cdots\bar\psi_4(\theta x)\right).
\]
In the product of the four entries $\psi^\theta_\alpha(x)=\sum_\beta(\gamma_n)_{\alpha\beta}\,\psi_\beta(\theta x)$, expanded, only the terms with four different generators survive, since each generator has square zero;
reordering these terms gives
\[
\psi^\theta_1(x)\cdots\psi^\theta_4(x)=\det(\gamma_n)\cdot\psi_1(\theta x)\cdots\psi_4(\theta x).
\]
The entries $\bar\psi^\theta_\alpha(x)=\sum_\beta(\gamma_n^{\mathrm{T}})_{\alpha\beta}\,\bar\psi_\beta(\theta x)$ arise from the generators $\bar\psi_\beta(\theta x)$ through the transposed matrix,
and the same expansion gives $\bar\psi^\theta_1(x)\cdots\bar\psi^\theta_4(x)=\det(\gamma_n^{\mathrm{T}})\cdot\bar\psi_1(\theta x)\cdots\bar\psi_4(\theta x)$, with $\det(\gamma_n^{\mathrm{T}})=\det(\gamma_n)$.
The matrix $\gamma_n$ is Hermitian, with $\gamma_n^2=1$ and trace zero (Definition~\ref{def:spin-matrices});
its eigenvalues are $1,1,-1,-1$, and $\det(\gamma_n)=1$.
As $s_4^2=1$, it follows that
\[
\prod_{\alpha=1}^{4}\psi^\theta_\alpha(x)\,\bar\psi^\theta_\alpha(x)
=s_4\cdot\det(\gamma_n)^2\cdot\left(\psi_1(\theta x)\cdots\psi_4(\theta x)\right)\left(\bar\psi_1(\theta x)\cdots\bar\psi_4(\theta x)\right)
=\det(\gamma_n)^2\cdot\Psi^{\mathrm{top}}(\theta x)
=\Psi^{\mathrm{top}}(\theta x),
\]
and $\Theta_\Pi\left(\Psi^{\mathrm{top}}(Y)\right)=\prod_{x\in Y}\Psi^{\mathrm{top}}(\theta x)=\Psi^{\mathrm{top}}(\theta Y)$.

\eqref{it:torus-integrability}
Let $Y$ be the set of representatives of $\mathbb{X}_{\mathbb{T}}$ with which the monomial basis is formed;
the basis has $2^{8\#|\mathbb{X}_{\mathbb{T}}|}$ elements.
The action $S_{\mathbb{T}}(A)$ is even, of degree two, so that
\[
\exp\left(-S_{\mathbb{T}}(A)\right)=\sum_{j=0}^{4\#|\mathbb{X}_{\mathbb{T}}|}\frac{1}{j!}\left(-S_{\mathbb{T}}(A)\right)^j
\]
is a polynomial in the generators.
Its coefficients in the monomial basis are polynomials, with coefficients independent of $A$,
in $M$, in the entries of the hopping matrices $w(x,x')$ and in the phases $\exp\left(\pm iqA(b)\right)$ of the links $b$ of $\mathbb{T}$
(Definitions~\ref{def:chamber-operator} and~\ref{def:tori}).
The phases have modulus one for real $A$;
so these coefficients are bounded, uniformly in $A$, by a number $C'_{\mathbb{T}}$ that depends only on $\mathbb{T}$, $\rho$, $w_0,\dots,w_3$ and $M$.
Put $C_{\mathbb{T}}:=2^{8\#|\mathbb{X}_{\mathbb{T}}|}\cdot C'_{\mathbb{T}}$.
For a basis monomial $\Psi$ and each basis monomial $\Psi'$ there is at most one basis monomial $\Psi''$ with $\Psi'\,\Psi''=\pm\Psi$.
Hence the coefficient of $\Psi$ in $\exp\left(-S_{\mathbb{T}}(A)\right)\cdot G(A,z)$ is a sum of at most $2^{8\#|\mathbb{X}_{\mathbb{T}}|}$ terms,
each of which is $\pm$ the product of a coefficient of $\exp\left(-S_{\mathbb{T}}(A)\right)$ and a coefficient of $G(A,z)$.
With the bound of Definition~\ref{def:torus-integrands} for the coefficients of $G(A,z)$,
\[
\left|a_\Psi(A,z)\right|
\le\exp\left(-\tfrac12\|F\|^2\right)\cdot C_{\mathbb{T}}\cdot C\cdot\exp\left(\left(\tfrac12-\varepsilon_G\right)\|F\|^2\right)
=C_{\mathbb{T}}\cdot C\cdot\prod_{\sigma\in\mathbb{T}^{(2)}}\exp\left(-\varepsilon_G\,F_\sigma^2\right),
\]
where $\mathbb{T}^{(2)}$ is the set of the plaquettes of $\mathbb{T}$, that is, of the $\Gamma$-orbits of plaquettes.
The torus has $3|\mathbb{T}|$ links and $3|\mathbb{T}|$ plaquettes, three of each at every photon site of $\Z^3/\Gamma$.
Since $F_\sigma=(dA)_\sigma+\frac{2\pi}{q}\,z_\sigma$, dropping the condition $dz=0$ and applying Lemma~\ref{lem:poisson-kernels}\eqref{it:gaussian-sums} with $\varepsilon=\varepsilon_G$ to each plaquette gives, for every $A$,
\[
\sum_{\substack{z\in C^2(\Gamma;\Z)\\ dz=0}}\ \prod_{\sigma\in\mathbb{T}^{(2)}}\exp\left(-\varepsilon_G\,F_\sigma^2\right)
\le\prod_{\sigma\in\mathbb{T}^{(2)}}\ \sum_{z_\sigma\in\Z}\exp\left(-\varepsilon_G\left((dA)_\sigma+\frac{2\pi}{q}\,z_\sigma\right)^2\right)
\le C_{q,\varepsilon_G}^{3|\mathbb{T}|} .
\]
The box $[0,2\pi/q)^{\mathbb{T}^{(1)}}$ has the volume $(2\pi/q)^{3|\mathbb{T}|}$.
By Tonelli's theorem the sum over $z$ and the integral over $A$ of the non-negative function $|a_\Psi|$ may be interchanged, and the asserted bound follows.
The function $a_\Psi$ is Borel, because $G$ is a Borel map and the coefficients of $\exp\left(-S_{\mathbb{T}}(A)\right)$ are continuous in $A$.
It is invariant under the integer shifts $(A,z)\mapsto\left(A+(2\pi/q)\,k,\ z-dk\right)$, $k\in C^1(\Gamma;\Z)$:
the map $G$ is invariant by Definition~\ref{def:torus-integrands}, the curvature $F$ does not change,
and $S_{\mathbb{T}}(A)$ depends on $A$ only through the phases $\exp\left(iqA(b)\right)$, which do not change either.
So $a_\Psi$ is a Borel function on $\Conf(\Gamma)$, and the left-hand side of the bound is the integral of $|a_\Psi|$ over $\Conf(\Gamma)$ (Definition~\ref{def:configurations}).

The integral $\int_{\mathbb{T}}$ of an element of $\Gr_{\mathbb{T}}$ is its coefficient of $\Psi^{\mathrm{top}}(Y)$ (Definition~\ref{def:left-derivatives}),
and $\Psi^{\mathrm{top}}(Y)$ is $\pm$ a basis monomial.
Hence the integrand in the definition of $Z_{\mathbb{T}}[G]$ (Definition~\ref{def:torus-integrands}) is $\pm a_\Psi$ for this monomial $\Psi$;
the sum and the integral that define $Z_{\mathbb{T}}[G]$ converge absolutely, and $\left|Z_{\mathbb{T}}[G]\right|$ is at most the right-hand side of the bound.
An element of $\mathcal{A}(\mathbb{T})$ defines, linearly, an even element of $\mathcal{I}(\mathbb{T})$, and its value under $Z_{\mathbb{T}}$ is the value of this element
(Lemma~\ref{lem:product-integrands}\eqref{it:sum-torus});
for a product that is read on $\mathbb{T}$ this element is the function of Lemma~\ref{lem:product-integrands}\eqref{it:product-torus}.
Every element of $\mathcal{A}(\R^3)$ that embeds in $\mathbb{T}$ belongs to $\mathcal{A}(\mathbb{T})$ (Definition~\ref{def:torus-integrands}).
Let $\Phi$ be a bounded Borel function on $\Conf(\Gamma)$.
Then $\Phi G$ belongs to $\mathcal{I}(\mathbb{T})$, with the constants $\sup|\Phi|\cdot C$ and $\varepsilon_G$,
and the bound for $\left|Z_{\mathbb{T}}[\Phi G]\right|$ is the preceding one for $\Phi G$.
The functional $Z_{\mathbb{T}}$ is linear on $\mathcal{I}(\mathbb{T})$;
so for bounded Borel functions $\Phi$ and $\Phi'$ on $\Conf(\Gamma)$
\[
\left|Z_{\mathbb{T}}[\Phi G]-Z_{\mathbb{T}}[\Phi'G]\right|=\left|Z_{\mathbb{T}}[(\Phi-\Phi')\,G]\right|
\le\sup|\Phi-\Phi'|\cdot C_{\mathbb{T}}\cdot C\cdot\left(\frac{2\pi}{q}\right)^{3|\mathbb{T}|}\cdot C_{q,\varepsilon_G}^{3|\mathbb{T}|},
\]
which is the continuity for the supremum norm.
\end{proof}

\subsection{Fermions between parallel mirrors}\label{sec:slab-fermions}

The family $\ell\in\mathcal{L}$ and the cross-section $\Xi=(\Gamma_\parallel,\chi_\parallel)$ for $\ell$ (Definition~\ref{def:slabs}) stay fixed,
and fermion sites, bonds and cells are taken on the cylinder $\Z^3/\Gamma_\parallel$, that is, modulo $\Gamma_\parallel$.
A real one-cochain on the cylinder is a $\Gamma_\parallel$-periodic real one-cochain $A$ on $\Z^3$;
it enters the fermion action through the factors $U_A(x,x')$ of Definition~\ref{def:chamber-operator}.

\begin{lemma}[Bonds across a mirror]\label{lem:slab-bonds}
Let $A$ be a real one-cochain on the cylinder.
\begin{enumerate}
\item\label{it:bonds-crossing}
No fermion site lies on a mirror.
Let $c\in\Adm(\ell)$.
If the end points of a bond lie on different sides of the mirror $(\ell,c)$, then the bond is $\{x,\theta_cx\}$ for either of its end points $x$,
its open segment meets no mirror other than $(\ell,c)$,
and both end points lie in $\{c^{-}<\ell<c^{+}\}$.
\item\label{it:bonds-slab-action}
Let $c<c'$ be elements of $\Adm(\ell)$, and put
\[
S_{(c,c')}(A):=\sum_{x}\left(M\cdot\bar\psi(x)\,\psi(x)-\sum_{i}U_A\left(x,x^{(i)}\right)\cdot\bar\psi(x)\,w\left(x,x^{(i)}\right)\psi\left(x^{(i)}\right)\right),
\]
where $x$ ranges over a set of representatives of $\mathbb{X}_\Xi((c,c'))$ (Definition~\ref{def:slab-grassmann})
and, for each such $x$, the type $i\in\{0,1,2,3\}$ ranges over the types with $x^{(i)}\in\mathbb{X}((c,c'))$.
Then $S_{(c,c')}(A)$ does not depend on the choice of the representatives,
it is an even element of $\Gr_\Xi((c,c'))$,
and each of its coefficients in a monomial basis is a polynomial in the phases $\exp\left(\pm iqA(b)\right)$ of the links $b\in[c,c']^{(1)}$,
that is, of the links of the cylinder contained in the closed slab $\{c\le\ell\le c'\}$ (Definition~\ref{def:slab-configurations}).
\item\label{it:bonds-action-split}
For $c\in\Adm(\ell)$ put
\[
S^{\times}_c(A):=-\sum_{(x,x')}U_A(x,x')\cdot\bar\psi(x)\,w(x,x')\,\psi(x'),
\]
where $(x,x')$ ranges over a set of representatives, modulo $\Gamma_\parallel$, of the ordered pairs of end points of the bonds that cross $(\ell,c)$.
Then $S^{\times}_c(A)$ is an even element of $\Gr_\Xi((c^{-},c^{+}))$.
It depends on $A$ only through the phases $\exp\left(\pm iqA(b)\right)$ of the links $b\in\mathcal{X}^1(c)$ cut by $(\ell,c)$,
and it does not depend on $A$ at all unless $(\ell,c)$ is a link plane.
For $c_0<c_1<\dots<c_r$ in $\Adm(\ell)$,
\[
S_{(c_0,c_r)}(A)=\sum_{k=1}^{r}S_{(c_{k-1},c_k)}(A)+\sum_{k=1}^{r-1}S^{\times}_{c_k}(A).
\]
\end{enumerate}
\end{lemma}

\begin{proof}
Proposition~\ref{prop:geometry}\eqref{it:geometry-mirrors} states that no site lies on a mirror
and that a bond whose end points lie on different sides of $(\ell,c)$ is of the form $\{x,\theta_cx\}$.
Let $x$ be the end point with $\ell(x)<c$; then $\ell(\theta_cx)=2c-\ell(x)>c$ by Definition~\ref{def:mirrors}\eqref{it:mirror-height}.
By Proposition~\ref{prop:geometry}\eqref{it:geometry-midpoint} the open segment meets exactly one mirror;
as its end points lie on different sides of $(\ell,c)$, this mirror is $(\ell,c)$.
The planes $\{\ell=c^{-}\}$ and $\{\ell=c^{+}\}$ are mirrors different from $(\ell,c)$.
If one had $\ell(x)\le c^{-}$, the closed segment, along which $\ell$ is an affine function with the values from $\ell(x)$ to $2c-\ell(x)>c^{-}$,
would meet the plane $\{\ell=c^{-}\}$, either at the site $x$ or at a point of the open segment; both are excluded.
Hence $c^{-}<\ell(x)<c$ and $c<\ell(\theta_cx)=2c-\ell(x)<2c-c^{-}=c^{+}$, which proves~\eqref{it:bonds-crossing}.

Bonds of type $3$ and links are related as follows.
Let $b$ be a link of direction $e_\mu$, and let $\{x,x'\}$ be a bond over $b$ (Proposition~\ref{prop:geometry}\eqref{it:geometry-projection}),
its end points named so that $\pi(x')=\pi(x)+e_\mu$.
On $b$ every form $\ell'\in\mathcal{L}$ takes the values between the integers $\ell'(\pi(x))$ and $\ell'(\pi(x))+(n_{\ell'})_\mu$, where $(n_{\ell'})_\mu\in\{-1,0,1\}$.
A mirror of the family $\ell'$ that cuts $b$ therefore has $(n_{\ell'})_\mu\neq0$ and a position strictly between two consecutive integers.
The positions of the diagonal mirrors are integers, so $\ell'$ is an axis family; as $(n_{\ell'})_\mu\neq0$, it is $x_\mu$, and the position is $\pi(x)_\mu+\frac{1}{2}$.
Thus $b$ is cut by exactly one mirror, the link plane $\left\{x_\mu=\pi(x)_\mu+\frac{1}{2}\right\}$.
The coordinates of $x-\pi(x)$ and of $x'-\pi(x')$ have the moduli $v_1,v_2,v_3$, and $v_1<\frac{1}{2}$ (Definition~\ref{def:fermion-sites}); therefore
\[
x_\mu\le\pi(x)_\mu+v_1<\pi(x)_\mu+\frac{1}{2}<\pi(x)_\mu+1-v_1\le x'_\mu .
\]
So every bond over $b$ crosses the link plane that cuts $b$,
and by Proposition~\ref{prop:geometry}\eqref{it:geometry-midpoint} it crosses no other mirror.

Consider now~\eqref{it:bonds-slab-action}, and let $g\in\Gamma_\parallel$.
Translations commute with the maps $x\mapsto x^{(i)}$ and preserve the types of bonds, so that $(x+g)^{(i)}=x^{(i)}+g$ and $w\left(x+g,x^{(i)}+g\right)=w\left(x,x^{(i)}\right)$;
moreover $U_A\left(x+g,x^{(i)}+g\right)=U_A\left(x,x^{(i)}\right)$, because $\pi(x+g)=\pi(x)+g$ and $A$ is $\Gamma_\parallel$-periodic;
and the two fields of a term acquire the factor $\chi_\parallel(g)^2=1$.
Hence $S_{(c,c')}(A)$ does not depend on the representatives.
Each term is a linear combination of products of two generators, so $S_{(c,c')}(A)$ is even,
and its coefficients are linear combinations of $1$ and of the factors $U_A\left(x,x^{(i)}\right)$ of the bonds with both end points in $\{c<\ell<c'\}$.
By Proposition~\ref{prop:geometry}\eqref{it:geometry-projection}, $U_A\left(x,x^{(i)}\right)=1$ for $i\neq3$,
while a bond of type $3$ lies over a link $b$ and its factor is $\exp\left(\pm iqA(b)\right)$.
Let such a bond have both end points in $\{c<\ell<c'\}$.
It crosses neither $(\ell,c)$ nor $(\ell,c')$, so by the preceding paragraph neither of these two mirrors cuts $b$.
If $b$ lies in the plane $\{\ell=c\}$, then $b\subset\{\ell\ge c\}$.
Otherwise the open link $\mathring b$ lies on one side of $(\ell,c)$, and by Proposition~\ref{prop:geometry}\eqref{it:geometry-projection} the bond lies on the same side,
which is $\{\ell>c\}$; again $b\subset\{\ell\ge c\}$.
The same argument at the mirror $(\ell,c')$, below which the bond lies, gives $b\subset\{\ell\le c'\}$.
Therefore $b$ is contained in the closed slab, $b\in[c,c']^{(1)}$;
in particular $S_{(c,c')}(A)$ is determined by the restriction of $A$ to $[c,c']^{(1)}$.

In~\eqref{it:bonds-action-split}, the end points of a bond that crosses $(\ell,c)$ lie in $\{c^{-}<\ell<c^{+}\}$ by~\eqref{it:bonds-crossing},
and the terms of $S^{\times}_c(A)$ do not depend on the representatives, by the argument just given;
so $S^{\times}_c(A)$ is an even element of $\Gr_\Xi((c^{-},c^{+}))$.
A bond that crosses $(\ell,c)$ carries a factor $U_A(x,x')\neq1$ only if it is of type $3$.
It then lies over a link $b$ and crosses no mirror other than the link plane that cuts $b$, as shown above;
this plane is therefore $(\ell,c)$.
Hence $(\ell,c)$ is a link plane, $b\in\mathcal{X}^1(c)$, and the factor is $\exp\left(\pm iqA(b)\right)$.

For the identity, note first that no site lies on a mirror;
so $\mathbb{X}((c_0,c_r))$ is the disjoint union of the sets $\mathbb{X}((c_{k-1},c_k))$, $1\le k\le r$,
and the diagonal terms $M\cdot\bar\psi(x)\,\psi(x)$ are distributed accordingly.
Let $(x,x')$ be an oriented bond with both end points in $\{c_0<\ell<c_r\}$.
If $x$ and $x'$ lie in one open slab $\{c_{k-1}<\ell<c_k\}$, the term of $(x,x')$ occurs in $S_{(c_{k-1},c_k)}(A)$ and in no other summand of the right-hand side.
Otherwise some $c_k$ with $1\le k\le r-1$ lies strictly between $\ell(x)$ and $\ell(x')$, and the bond crosses $(\ell,c_k)$;
by~\eqref{it:bonds-crossing} it crosses no other mirror, and its term occurs in $S^{\times}_{c_k}(A)$ and nowhere else on the right.
Conversely, let a bond cross $(\ell,c_k)$ with $1\le k\le r-1$.
Its end points lie in $\{c_k^{-}<\ell<c_k^{+}\}$,
and $c_0\le c_{k-1}\le c_k^{-}$, $c_k^{+}\le c_{k+1}\le c_r$, because two different points of $\Adm(\ell)$ are at least $\frac{|n_\ell|^2}{2}$ apart;
so both terms of the bond occur in $S_{(c_0,c_r)}(A)$.
\end{proof}

\begin{lemma}\label{lem:integral-covariance}
Let $\mathbb{I}$ be a bounded open interval, $G\in\Gr_\Xi(\mathbb{I})$, $v\in\Z^3$ and $c\in\Adm(\ell)$.
Then, with the Berezin integrals of Definition~\ref{def:left-derivatives}
over the sets $\mathbb{X}_\Xi(\mathbb{I})$, $\mathbb{X}_\Xi(\mathbb{I}+\ell(v))$ and $\mathbb{X}_\Xi(2c-\mathbb{I})$,
\begin{equation}\label{eq:slab-integral-covariance}
\int_{\mathbb{I}+\ell(v)}\tau_vG=\int_{\mathbb{I}}G,
\qquad
\int_{2c-\mathbb{I}}\Theta_cG=\overline{\int_{\mathbb{I}}G}.
\end{equation}
\end{lemma}

\begin{proof}
By Lemma~\ref{lem:slab-automorphisms}, $\tau_v$ is an isomorphism of algebras of $\Gr_\Xi(\mathbb{I})$ onto $\Gr_\Xi(\mathbb{I}+\ell(v))$
which preserves the degree and sends $\psi_\alpha(x)$, $\bar\psi_\alpha(x)$ to $\psi_\alpha(x+v)$, $\bar\psi_\alpha(x+v)$.
Let $Y$ be a set of representatives of $\mathbb{X}_\Xi(\mathbb{I})$.
Then $Y+v$ is a set of representatives of $\mathbb{X}_\Xi(\mathbb{I}+\ell(v))$, and $\tau_v\Psi^{\mathrm{top}}(Y)=\Psi^{\mathrm{top}}(Y+v)$.
The elements of top degree of $\Gr_\Xi(\mathbb{I})$ are the multiples of $\Psi^{\mathrm{top}}(Y)$,
and the component of top degree of $G$ is $\left(\int_{\mathbb{I}}G\right)\cdot\Psi^{\mathrm{top}}(Y)$ (Definition~\ref{def:left-derivatives}).
Hence the component of top degree of $\tau_vG$ is $\left(\int_{\mathbb{I}}G\right)\cdot\Psi^{\mathrm{top}}(Y+v)$, which is the first identity.

For the second, $\theta_c(x+g)=\theta_c(x)+g$ for $g\in\Gamma_\parallel$, and $\ell\circ\theta_c=2c-\ell$ (Definition~\ref{def:mirrors}\eqref{it:mirror-height});
so $\theta_c$ induces a bijection of $\mathbb{X}_\Xi(\mathbb{I})$ onto $\mathbb{X}_\Xi(2c-\mathbb{I})$,
and $\Theta_c$ maps $\Gr_\Xi(\mathbb{I})$ into $\Gr_\Xi(2c-\mathbb{I})$ (Lemma~\ref{lem:slab-automorphisms}).
The identity is the reflection rule of Proposition~\ref{prop:berezin}\eqref{it:reflected-integral} for the mirror $(\ell,c)$, the cross-section $\Xi$ and the interval $\mathbb{I}$.
\end{proof}

\begin{definition}[Translated modes]\label{def:mode-translations}
For $c\in\Adm(\ell)$, the modes at the mirror $(\ell,c)$ form the set $\idx(\ell,c)$ of Definition~\ref{def:modes};
for $i\in\idx(\ell,c)$, $\zeta_i$ is the mode, $\omega_i\ge0$ its crossing coefficient and $u_i$ its phase.
A mode $i$ consists of a bond $\{x,\theta_cx\}$ that crosses $(\ell,c)$, with $\ell(x)<c$, of one of the four groups of Definition~\ref{def:modes}, and of an index $k\in\{1,2\}$.

\emph{Translated modes.}
Let $v\in\Z^3$, and let $i\in\idx(\ell,c)$ be a mode with the bond $\{x,\theta_cx\}$, $\ell(x)<c$.
The bond $\{x+v,\theta_cx+v\}$ crosses the mirror $(\ell,c+\ell(v))$,
since $\theta_cx+v=\theta_{c+\ell(v)}(x+v)$ by Definition~\ref{def:mirrors}\eqref{it:mirror-conjugation} and $\ell(x+v)=\ell(x)+\ell(v)<c+\ell(v)$.
The mode of $(\ell,c+\ell(v))$ with this bond, the same group and the same index $k$ is denoted $i+v$.
One has
\[
\zeta_{i+v}=\tau_v\zeta_i,
\qquad
\omega_{i+v}=\omega_i ,
\]
because the bases $e^{\pm}_{\hat n_\ell,k}$ of Definition~\ref{def:spin-matrices} depend only on the unit normal $\hat n_\ell$, which all mirrors of the family share,
and because translations preserve the types of bonds.
If $g\in\Gamma_\parallel$, then $i+g$ is again a mode of $(\ell,c)$, and $\zeta_{i+g}=\chi_\parallel(g)\cdot\zeta_i$ in the algebras $\Gr_\Xi$.
The sets $\{i+g:\ g\in\Gamma_\parallel\}$ are the classes of modes modulo $\Gamma_\parallel$.

\emph{Representatives.}
For every $c\in\Adm(\ell)$, $\mathcal{J}(c)\subset\idx(\ell,c)$ is a set of representatives of the classes of modes modulo $\Gamma_\parallel$, endowed with a total order,
the sets being chosen compatibly with the translation by $n_\ell$, as in Definition~\ref{def:modes}:
\[
\mathcal{J}\left(c+|n_\ell|^2\right)=\mathcal{J}(c)+n_\ell,
\qquad
i+n_\ell<j+n_\ell\ \text{ in }\mathcal{J}\left(c+|n_\ell|^2\right)\iff i<j\ \text{ in }\mathcal{J}(c),
\]
for all $c\in\Adm(\ell)$ and $i,j\in\mathcal{J}(c)$.
Such sets exist: they may be chosen for $c\in\left\{0,\frac{1}{2}\right\}$ if $\ell$ is an axis family and for $c\in\{0,1\}$ if $\ell$ is a diagonal family,
and transported by the translations by $k\,n_\ell$, $k\in\Z$.
Each $\mathcal{J}(c)$ is finite:
a bond that crosses $(\ell,c)$ has its end points in $\{c^{-}<\ell<c^{+}\}$ (Lemma~\ref{lem:slab-bonds}\eqref{it:bonds-crossing}),
the set $\mathbb{X}_\Xi((c^{-},c^{+}))$ is finite since the interval is bounded,
and every bond carries eight modes.

\emph{Fock spaces.}
The space $\mathcal{F}(c):=\mathcal{F}(\mathcal{J}(c))$ is the Hilbert space of Definition~\ref{def:modes} with the orthonormal basis $e_I$, $I\subset\mathcal{J}(c)$.
On it act the parity, $U_{\mathrm{f}}e_I=(-1)^{\#|I|}\cdot e_I$,
and the diagonal operator $\omega_c:=\diag(\omega_I)$ with the entries $\omega_I:=\prod_{i\in I}\omega_i$.

\emph{Monomials.}
For $I\subset\mathcal{J}(c)$ let $\zeta^{(c)}_I$ be the product of the $\zeta_i$, $i\in I$, taken in increasing order,
that is, the ordered monomial of~\eqref{eq:setting-ordered-monomials} formed with the order of $\mathcal{J}(c)$, and let $\zeta^{(c)}_{\emptyset}:=1$;
the upper index names the mirror of the modes, and $\zeta^{(c)}_i:=\zeta_i$ for $i\in\mathcal{J}(c)$.
By Definition~\ref{def:modes} and the relation $\zeta_{i+g}=\chi_\parallel(g)\cdot\zeta_i$,
the $\zeta_i$ with $i\in\mathcal{J}(c)$ form a basis of the linear span of the generators of $\Gr_\Xi$ at the upper end points of the bonds that cross $(\ell,c)$,
and the $\Theta_c(\zeta_i)$ form a basis of the span of the generators at the lower end points;
by Lemma~\ref{lem:slab-bonds}\eqref{it:bonds-crossing}, $\zeta_i\in\Gr_\Xi((c,c^{+}))$ and $\Theta_c(\zeta_i)\in\Gr_\Xi((c^{-},c))$.
The linear map
\[
\zeta^{(c)}\colon\mathcal{F}(c)\to\Gr_\Xi\left((c,c^{+})\right),
\qquad
\zeta^{(c)}(e_I):=\zeta^{(c)}_I ,
\]
is injective, because the $\zeta_i$, $i\in\mathcal{J}(c)$, are linearly independent elements of degree one;
its image $\Gr^{\uparrow}(c)$ is the subalgebra generated by the fields at the upper end points of the bonds that cross $(\ell,c)$.

\emph{Translations.}
Let $v\in\Z^3$ and $i\in\mathcal{J}(c)$.
The class of $i+v$ has exactly one representative $r_v(i)\in\mathcal{J}(c+\ell(v))$,
and $g_v(i)\in\Gamma_\parallel$ is the element with $i+v=r_v(i)+g_v(i)$;
it is unique, since every non-zero element of $\Gamma_\parallel$ moves the lower end point of a bond and hence every mode.
For $I\subset\mathcal{J}(c)$ put $r_v(I):=\{r_v(i):\ i\in I\}$,
and let $s_v(I)\in\{1,-1\}$ be the product of the numbers $\chi_\parallel(g_v(i))$, $i\in I$,
and of the sign of the permutation that puts the family $(r_v(i))_{i\in I}$, listed in the increasing order of $i$, into increasing order.
Let
\[
U_v\colon\mathcal{F}(c)\to\mathcal{F}\left(c+\ell(v)\right),
\qquad
U_ve_I:=s_v(I)\cdot e_{r_v(I)} .
\]
These objects have the following properties, for all $c\in\Adm(\ell)$ and $v,v'\in\Z^3$.
\begin{enumerate}
\item\label{it:fock-translations}
The map $r_v$ is a bijection of $\mathcal{J}(c)$ onto $\mathcal{J}(c+\ell(v))$, one has $\omega_{r_v(I)}=\omega_I$, and
\[
\tau_v\zeta^{(c)}_I=s_v(I)\cdot\zeta^{(c+\ell(v))}_{r_v(I)}\quad\left(I\subset\mathcal{J}(c)\right),
\qquad\text{that is,}\qquad
\tau_v\circ\zeta^{(c)}=\zeta^{(c+\ell(v))}\circ U_v ;
\]
thus $U_v=\left(\zeta^{(c+\ell(v))}\right)^{-1}\circ\tau_v\circ\zeta^{(c)}$, the inverse being taken on $\Gr^{\uparrow}(c+\ell(v))$.
The map $U_v$ is unitary and its matrix entries are real.
One has $U_vU_{v'}=U_{v+v'}$ as maps from $\mathcal{F}(c)$ to $\mathcal{F}(c+\ell(v+v'))$,
and $U_ve_I=\chi_\parallel(v)^{\#|I|}\cdot e_I$ if $v\in\Gamma_\parallel$.
The map $U_v$ commutes with the parity, and $U_v\,\omega_c=\omega_{c+\ell(v)}\,U_v$.
If $v\in\Z n_\ell$, then $\mathcal{J}(c+\ell(v))=\mathcal{J}(c)+v$ with the transported order, and $U_ve_I=e_{I+v}$.
\item\label{it:reflected-modes}
Let $c'\in\Adm(\ell)$, put $c'':=2c'-c$, and let $v_{c,c'}=\frac{2(c'-c)}{|n_\ell|^2}\,n_\ell$ be the vector of Lemma~\ref{lem:slab-automorphisms}\eqref{it:parallel-reflections}.
Then $v_{c,c'}\in\Z n_\ell$, $c+\ell(v_{c,c'})=c''$, $\mathcal{J}(c'')=\mathcal{J}(c)+v_{c,c'}$, and
\[
\Theta_{c''}\left(\zeta^{(c'')}_{I+v_{c,c'}}\right)=\Theta_{c'}\left(\zeta^{(c)}_I\right)
\qquad\left(I\subset\mathcal{J}(c)\right).
\]
\end{enumerate}
\end{definition}

\begin{proof}[Proof of the properties in Definition~\ref{def:mode-translations}]
\eqref{it:fock-translations}
The map $i\mapsto i+v$ is a bijection of $\idx(\ell,c)$ onto $\idx(\ell,c+\ell(v))$, with the inverse $i\mapsto i+(-v)$,
and it commutes with the translations by $\Gamma_\parallel$.
It therefore induces a bijection of the classes modulo $\Gamma_\parallel$, and $r_v$ is this bijection, read on the representatives.
Since $r_v(i)=i+\left(v-g_v(i)\right)$ and the crossing coefficients do not change under translations,
$\omega_{r_v(i)}=\omega_i$, whence $\omega_{r_v(I)}=\omega_I$.

The translation $\tau_v$ is a homomorphism of algebras.
So $\tau_v\zeta^{(c)}_I$ is the product, in increasing order of $i\in I$, of the elements
\[
\tau_v\zeta_i=\zeta_{i+v}=\chi_\parallel\left(g_v(i)\right)\cdot\zeta_{r_v(i)} .
\]
These factors have degree one and anticommute;
putting them into the increasing order of the $r_v(i)$ multiplies the product by the sign of the sorting permutation.
Hence $\tau_v\zeta^{(c)}_I=s_v(I)\cdot\zeta^{(c+\ell(v))}_{r_v(I)}$,
which is the identity $\tau_v\circ\zeta^{(c)}=\zeta^{(c+\ell(v))}\circ U_v$ on the basis vectors $e_I$;
the formula for $U_v$ follows, $\zeta^{(c+\ell(v))}$ being injective.

As $r_v$ is a bijection, $I\mapsto r_v(I)$ is a bijection of the set of the subsets of $\mathcal{J}(c)$ onto the set of the subsets of $\mathcal{J}(c+\ell(v))$.
So $U_v$ maps the orthonormal basis $(e_I)$ onto the orthonormal basis $(e_{r_v(I)})$, up to the signs $s_v(I)$:
it is unitary, and its matrix entries lie in $\{0,1,-1\}$.

By Lemma~\ref{lem:slab-automorphisms}\eqref{it:translation-law}, $\tau_v\tau_{v'}=\tau_{v+v'}$.
With the identity just proved, applied at the positions $c+\ell(v')$, $c$ and $c$,
\[
\zeta^{(c+\ell(v+v'))}\circ U_vU_{v'}
=\tau_v\circ\zeta^{(c+\ell(v'))}\circ U_{v'}
=\tau_v\tau_{v'}\circ\zeta^{(c)}
=\tau_{v+v'}\circ\zeta^{(c)}
=\zeta^{(c+\ell(v+v'))}\circ U_{v+v'} ,
\]
and the injectivity of $\zeta^{(c+\ell(v+v'))}$ gives $U_vU_{v'}=U_{v+v'}$.

Let $v\in\Gamma_\parallel$, so that $\ell(v)=0$.
By the same item of Lemma~\ref{lem:slab-automorphisms}, $\tau_v$ multiplies the element $\zeta^{(c)}_I$, which is homogeneous of degree $\#|I|$, by $\chi_\parallel(v)^{\#|I|}$.
Thus $\zeta^{(c)}(U_ve_I)=\tau_v\zeta^{(c)}_I=\chi_\parallel(v)^{\#|I|}\cdot\zeta^{(c)}_I$, and $U_ve_I=\chi_\parallel(v)^{\#|I|}\cdot e_I$ by injectivity.

The vector $U_ve_I$ is a multiple of $e_{r_v(I)}$, and $r_v(I)$ has as many elements as $I$;
hence $U_{\mathrm{f}}U_ve_I=(-1)^{\#|I|}\cdot U_ve_I=U_vU_{\mathrm{f}}e_I$.
Together with $\omega_{r_v(I)}=\omega_I$ this also gives
$U_v\,\omega_ce_I=\omega_I\cdot s_v(I)\cdot e_{r_v(I)}=\omega_{c+\ell(v)}\,U_ve_I$.

Finally let $v=k\,n_\ell$ with $k\in\Z$.
Then $\ell(v)=k\,|n_\ell|^2$, and the compatibility of the representatives, applied $|k|$ times, gives $\mathcal{J}(c+\ell(v))=\mathcal{J}(c)+v$ with the transported order.
Consequently $r_v(i)=i+v$ and $g_v(i)=0$ for $i\in\mathcal{J}(c)$, the map $r_v$ preserves the order,
$s_v(I)=1$ and $r_v(I)=I+v$; so $U_ve_I=e_{I+v}$.

\eqref{it:reflected-modes}
Since $c'-c\in\frac{|n_\ell|^2}{2}\,\Z$, the vector $v_{c,c'}$ lies in $\Z n_\ell$,
and $\ell(v_{c,c'})=2(c'-c)$ gives $c+\ell(v_{c,c'})=c''$.
By the last assertion of~\eqref{it:fock-translations}, $\mathcal{J}(c'')=\mathcal{J}(c)+v_{c,c'}$ and
\[
\zeta^{(c'')}_{I+v_{c,c'}}=\zeta^{(c'')}\left(U_{v_{c,c'}}e_I\right)=\tau_{v_{c,c'}}\zeta^{(c)}_I .
\]
Lemma~\ref{lem:slab-automorphisms}\eqref{it:parallel-reflections} states that $\Theta_{c''}\tau_{v_{c,c'}}=\Theta_{c'}$;
applying $\Theta_{c''}$ to the last identity yields the assertion.
\end{proof}

\begin{definition}[Slab matrices]\label{def:slab-matrices}
Let $c\in\Adm(\ell)$, and let $\Upsilon$ be an operator on $\mathcal{F}(c)$,
with the matrix entries $\Upsilon_{IJ}:=\langle e_I,\Upsilon e_J\rangle$, $I,J\subset\mathcal{J}(c)$.
The \emph{insertion} of $\Upsilon$ at the mirror $(\ell,c)$ is
\[
\Ins_c[\Upsilon]:=\sum_{I,J\subset\mathcal{J}(c)}\Upsilon_{IJ}\cdot\Theta_c\left(\zeta^{(c)}_J\right)\zeta^{(c)}_I\ \in\ \Gr_\Xi\left((c^{-},c^{+})\right);
\]
as in~\eqref{eq:setting-ordered-monomials}, the row index $I$ belongs to the upper side of the mirror and the column index $J$ to the lower side.

Let $c<c'$ be elements of $\Adm(\ell)$, and let $G\in\Gr_\Xi((c,c'))$.
Since $c^{+}\le c'$ and $c\le{c'}^{-}$, the elements $\zeta^{(c)}_I\in\Gr_\Xi((c,c^{+}))$ and $\Theta_{c'}\left(\zeta^{(c')}_{I'}\right)\in\Gr_\Xi(({c'}^{-},c'))$
of Definition~\ref{def:mode-translations} belong to $\Gr_\Xi((c,c'))$.
The \emph{slab matrix} of $G$ is the linear map $\mathfrak{N}_{c,c'}[G]\colon\mathcal{F}(c)\to\mathcal{F}(c')$ with the matrix entries
\[
\mathfrak{N}_{c,c'}[G]_{I',I}:=\int_{(c,c')}\zeta^{(c)}_I\,G\,\Theta_{c'}\left(\zeta^{(c')}_{I'}\right),
\qquad I\subset\mathcal{J}(c),\ I'\subset\mathcal{J}(c'),
\]
that is, $\mathfrak{N}_{c,c'}[G]\,e_I=\sum_{I'\subset\mathcal{J}(c')}\mathfrak{N}_{c,c'}[G]_{I',I}\cdot e_{I'}$;
the column index $I$ refers to $\mathcal{F}(c)$ and the row index $I'$ to $\mathcal{F}(c')$.
\end{definition}

The crossing part of the action at a mirror is the insertion of a diagonal operator.
Let $c\in\Adm(\ell)$ and let $A$ be a real one-cochain on the cylinder.
For $I\subset\mathcal{J}(c)$ put $u_I(A):=\prod_{i\in I}u_i(A)$, where $u_i(A)$ is the phase of the mode $i$ (Definition~\ref{def:modes}),
and let $\diag\left(u_I(A)\right)$ be the diagonal operator on $\mathcal{F}(c)$ with these entries.
Then
\begin{equation}\label{eq:slab-crossing-factor}
\exp\left(-S^{\times}_c(A)\right)
=\sum_{I\subset\mathcal{J}(c)}\omega_I\cdot u_I(A)\cdot\Theta_c\left(\zeta^{(c)}_I\right)\zeta^{(c)}_I
=\Ins_c\left[\omega_c\diag\left(u_I(A)\right)\right].
\end{equation}
Indeed, by Proposition~\ref{prop:crossing} the two terms of a bond that crosses $(\ell,c)$ add up to
$-\sum_i\omega_i\cdot u_i(A)\cdot\Theta_c(\zeta_i)\,\zeta_i$, the sum extending over the eight modes $i$ of the bond.
For $g\in\Gamma_\parallel$ one has $\omega_{i+g}=\omega_i$, $u_{i+g}(A)=u_i(A)$ since $A$ is $\Gamma_\parallel$-periodic, and
\[
\Theta_c(\zeta_{i+g})\,\zeta_{i+g}=\chi_\parallel(g)^2\cdot\Theta_c(\zeta_i)\,\zeta_i=\Theta_c(\zeta_i)\,\zeta_i ,
\]
because $\Theta_c$ is antilinear and $\chi_\parallel(g)$ is real.
Summing over a set of representatives of the bonds that cross $(\ell,c)$,
and replacing every mode by the element of $\mathcal{J}(c)$ in its class, one therefore obtains
\[
-S^{\times}_c(A)=\sum_{i\in\mathcal{J}(c)}\omega_i\cdot u_i(A)\cdot\Theta_c(\zeta_i)\,\zeta_i .
\]
The elements $\Theta_c(\zeta_i)\,\zeta_i$ are even, so they commute with each other, and
$\left(\Theta_c(\zeta_i)\,\zeta_i\right)^2=-\Theta_c(\zeta_i)^2\,\zeta_i^2=0$, as $\zeta_i$ has degree one.
Let $I\subset\mathcal{J}(c)$ be non-empty.
Since $\Theta_c$ reverses the order of products, the two middle factors of the product $\Theta_c\left(\zeta^{(c)}_I\right)\zeta^{(c)}_I$ of elements of degree one
are $\Theta_c(\zeta_i)$ and $\zeta_i$, where $i$ is the least element of $I$;
this even pair may be moved to the far left, and induction on $\#|I|$ gives $\Theta_c\left(\zeta^{(c)}_I\right)\zeta^{(c)}_I=\prod_{j\in I}\Theta_c(\zeta_j)\,\zeta_j$,
as in Proposition~\ref{prop:crossing}.
The exponential of a sum of commuting elements is the product of their exponentials,
and the exponential of an element of square zero is the sum of $1$ and this element; hence
\[
\exp\left(-S^{\times}_c(A)\right)=\prod_{i\in\mathcal{J}(c)}\left(1+\omega_i\cdot u_i(A)\cdot\Theta_c(\zeta_i)\,\zeta_i\right).
\]
Expanding the product gives the first equality in~\eqref{eq:slab-crossing-factor};
the second is the definition of $\Ins_c$, applied to the diagonal operator $\omega_c\diag\left(u_I(A)\right)$, whose entries are the numbers $\omega_I\cdot u_I(A)$.

Every Grassmann algebra of Definitions~\ref{def:grassmann}, \ref{def:tori} and~\ref{def:slab-grassmann} is the direct sum of its even part,
spanned by the products of an even number of generators,
and of its odd part, spanned by the products of an odd number of generators.
An element $G$ of such an algebra which lies in one of the two parts, that is, which is even or odd, is also called \emph{homogeneous},
and its degree is then taken modulo~$2$:
$\deg G:=0\in\Z/2\Z$ if $G$ is even and $\deg G:=1$ if $G$ is odd, either value being admitted for $G=0$.
If $G$ is homogeneous in the sense of Definition~\ref{def:grassmann}, this is the class of the degree defined there.
Since an integral power of $1$ or of $-1$ depends only on the class of the exponent modulo~$2$,
the powers $(-1)^{\deg G}$ and $\chi(g)^{\deg G}$, for a character $\chi$ with values $\pm1$, are defined,
and so is the parity of a sum of integers and degrees.
An element of one of the classes of integrands of Definitions~\ref{def:slab-integrands} and~\ref{def:torus-integrands} is homogeneous
if all its values are homogeneous of one and the same degree, which is then its degree.
For homogeneous elements $G$ and $G'$ of one Grassmann algebra the product $GG'$ is homogeneous, of degree $\deg G+\deg G'$,
and $G'G=(-1)^{\deg G\cdot\deg G'}\cdot GG'$;
indeed, a product of $m$ generators times a product of $m'$ generators is a product of $m+m'$ generators,
and exchanging the two factors multiplies it by $(-1)^{mm'}$, since any two generators anticommute.

\begin{lemma}[Rules for slab matrices]\label{lem:matrix-rules}
Let $c<c'<c''$ be positions in $\Adm(\ell)$,
and let $G\in\Gr_\Xi((c,c'))$ and $G'\in\Gr_\Xi((c',c''))$ be arbitrary.
\begin{enumerate}
\item\label{it:matrix-parity}
\emph{Parity.}
If $G$ is homogeneous, then $\mathfrak{N}_{c,c'}[G]_{I',I}=0$
for all $I\subset\mathcal{J}(c)$ and $I'\subset\mathcal{J}(c')$ for which $\#|I|+\#|I'|+\deg G$ is odd.
Consequently
\[
U_{\mathrm{f}}\,\mathfrak{N}_{c,c'}[G]=(-1)^{\deg G}\cdot\mathfrak{N}_{c,c'}[G]\,U_{\mathrm{f}} .
\]
\item\label{it:matrix-composition}
\emph{Composition.}
For every linear operator $\Upsilon$ on $\mathcal{F}(c')$,
\[
\mathfrak{N}_{c,c''}\left[G\,\Ins_{c'}[\Upsilon]\,G'\right]=\mathfrak{N}_{c',c''}[G']\,\Upsilon\,\mathfrak{N}_{c,c'}[G] .
\]
\item\label{it:matrix-adjoint}
\emph{Adjoint.}
Let $c''=2c'-c$,
and let $v_{c,c'}=2(c'-c)\,n_\ell/|n_\ell|^2$ be the vector of Lemma~\ref{lem:slab-automorphisms}\eqref{it:parallel-reflections}.
Then
\[
\mathfrak{N}_{c',c''}[\Theta_{c'}G]_{I+v_{c,c'},\,I'}=\overline{\mathfrak{N}_{c,c'}[G]_{I',I}}
\qquad\left(I\subset\mathcal{J}(c),\ I'\subset\mathcal{J}(c')\right);
\]
equivalently, $\mathfrak{N}_{c',c''}[\Theta_{c'}G]=U_{v_{c,c'}}\,\mathfrak{N}_{c,c'}[G]^{*}$.
\item\label{it:matrix-translation}
\emph{Translation.}
For every $v\in\Z^3$,
\[
\mathfrak{N}_{c+\ell(v),\,c'+\ell(v)}[\tau_vG]=U_v\,\mathfrak{N}_{c,c'}[G]\,U_v^{-1} .
\]
\item\label{it:matrix-closure}
\emph{Closure.}
Let $\mathbb{T}=(\Gamma,\chi)$ be a torus with $\Xi(\mathbb{T},\ell)=\Xi$,
let $g$ be a closing vector of $\mathbb{T}$ for $\ell$ (Definition~\ref{def:slabs}),
put $p:=\ell(g)$,
and let $c_0<c_1<\dots<c_r=c_0+p$ be positions in $\Adm(\ell)$, where $r\ge1$.
The reading in $\Gr_{\mathbb{T}}$ is the homomorphism of algebras from $\Gr_\Xi((c_0-p,c_0+p))$ to $\Gr_{\mathbb{T}}$
that maps every generator $\psi_\alpha(x)$, $\bar\psi_\alpha(x)$ to the generator of $\Gr_{\mathbb{T}}$ with the same name;
it is well defined, since $\Gamma_\parallel\subset\Gamma$ and $\chi_\parallel$ is the restriction of $\chi$ to $\Gamma_\parallel$,
and an element of $\Gr_\Xi((c_0-p,c_0+p))$ is said to be read in $\Gr_{\mathbb{T}}$ as its image under the reading.
For $k=1,\dots,r$ let $G_k\in\Gr_\Xi((c_{k-1},c_k))$ be homogeneous, and put $\mathfrak{N}_k:=\mathfrak{N}_{c_{k-1},c_k}[G_k]$;
for $k=0,\dots,r-1$ let $\Upsilon_k$ be an even operator on $\mathcal{F}(c_k)$.
If $\sum_{k=1}^{r}\deg G_k$ is even, then
\[
\begin{aligned}
&\int_{\mathbb{T}}\Ins_{c_0}[\Upsilon_0]\,G_1\,\Ins_{c_1}[\Upsilon_1]\,G_2\cdots\Ins_{c_{r-1}}[\Upsilon_{r-1}]\,G_r\\
&\qquad=\operatorname{tr}_{\mathcal{F}(c_0)}\left(U_g^{-1}\,U_{\mathrm{f}}^{\frac{1+\chi(g)}{2}}\,\mathfrak{N}_r\Upsilon_{r-1}\mathfrak{N}_{r-1}\cdots\Upsilon_1\mathfrak{N}_1\Upsilon_0\right),
\end{aligned}
\]
where the product under the integral on the left, an element of $\Gr_\Xi((c_0-p,c_0+p))$, is read in $\Gr_{\mathbb{T}}$,
and where $U_{\mathrm{f}}^{\frac{1+\chi(g)}{2}}$ is the identity if $\chi(g)=-1$ and is $U_{\mathrm{f}}$ if $\chi(g)=+1$.
When $\sum_{k=1}^{r}\deg G_k$ is odd instead, the integral on the left of the display and the trace on its right are both equal to zero.
\end{enumerate}
\end{lemma}

\begin{proof}
\emph{Parity.}
Every mode is a linear combination of the generators $\psi_\alpha(x)$, $\bar\psi_\alpha(x)$ (Definition~\ref{def:modes}),
so it has degree one;
and $\Theta_{c'}$ preserves the degree,
since it maps each generator to a linear combination of generators (Definition~\ref{def:grassmann-reflections}).
For homogeneous $G$ the element $\zeta^{(c)}_I\,G\,\Theta_{c'}\left(\zeta^{(c')}_{I'}\right)$ is therefore homogeneous,
of the parity of $\#|I|+\deg G+\#|I'|$.
The Berezin integral of an odd element is zero (Proposition~\ref{prop:berezin}\eqref{it:berezin-factorisation}),
whence the first assertion.
For the second, note that the entries of $U_{\mathrm{f}}\,\mathfrak{N}_{c,c'}[G]$ and of $\mathfrak{N}_{c,c'}[G]\,U_{\mathrm{f}}$
at $(I',I)$ are $(-1)^{\#|I'|}\cdot\mathfrak{N}_{c,c'}[G]_{I',I}$ and $(-1)^{\#|I|}\cdot\mathfrak{N}_{c,c'}[G]_{I',I}$,
and that $(-1)^{\#|I'|}=(-1)^{\deg G}\cdot(-1)^{\#|I|}$ whenever $\mathfrak{N}_{c,c'}[G]_{I',I}\neq0$.

\emph{Composition.}
Fix $I\subset\mathcal{J}(c)$ and $I''\subset\mathcal{J}(c'')$.
The insertion at $c'$ is
$\Ins_{c'}[\Upsilon]=\sum_{I',J'}\Upsilon_{I'J'}\cdot\Theta_{c'}\left(\zeta^{(c')}_{J'}\right)\,\zeta^{(c')}_{I'}$,
the sum extending over all pairs of subsets $I',J'$ of $\mathcal{J}(c')$ (Definition~\ref{def:slab-matrices}).
Grouping the factors anew, without changing their order, one obtains
\[
\begin{aligned}
&\zeta^{(c)}_I\,G\,\Ins_{c'}[\Upsilon]\,G'\,\Theta_{c''}\left(\zeta^{(c'')}_{I''}\right)\\
&\qquad=\sum_{I',J'\subset\mathcal{J}(c')}\Upsilon_{I'J'}\cdot
\left[\zeta^{(c)}_I\,G\,\Theta_{c'}\left(\zeta^{(c')}_{J'}\right)\right]
\left[\zeta^{(c')}_{I'}\,G'\,\Theta_{c''}\left(\zeta^{(c'')}_{I''}\right)\right].
\end{aligned}
\]
Each factor in square brackets is an element of the algebra of a single slab:
the left one of $\Gr_\Xi((c,c'))$, the right one of $\Gr_\Xi((c',c''))$.
Indeed, $\zeta^{(c)}_I$ and $\zeta^{(c')}_{I'}$ lie in $\Gr_\Xi((c,c^{+}))$ and in $\Gr_\Xi((c',{c'}^{+}))$,
while $\Theta_{c'}\left(\zeta^{(c')}_{J'}\right)$ and $\Theta_{c''}\left(\zeta^{(c'')}_{I''}\right)$ lie in $\Gr_\Xi(({c'}^{-},c'))$ and in $\Gr_\Xi(({c''}^{-},c''))$
(Lemma~\ref{lem:slab-bonds}\eqref{it:bonds-crossing} and Definition~\ref{def:mode-translations});
moreover $c^{+}\le c'$, ${c'}^{-}\ge c$, ${c'}^{+}\le c''$ and ${c''}^{-}\ge c'$,
because two different points of $\Adm(\ell)$ are at least $|n_\ell|^2/2$ apart.
No site has the height $c'$ (Lemma~\ref{lem:slab-bonds}\eqref{it:bonds-crossing}),
so $\mathbb{X}_\Xi((c,c''))$ is the disjoint union of $\mathbb{X}_\Xi((c,c'))$ and $\mathbb{X}_\Xi((c',c''))$.
Proposition~\ref{prop:berezin}\eqref{it:berezin-factorisation}, which holds without any condition of homogeneity,
splits the Berezin integral of each term into the integral of the first bracket over $(c,c')$,
which is $\mathfrak{N}_{c,c'}[G]_{J',I}$,
times the integral of the second bracket over $(c',c'')$, which is $\mathfrak{N}_{c',c''}[G']_{I'',I'}$.
Hence
\[
\mathfrak{N}_{c,c''}\left[G\,\Ins_{c'}[\Upsilon]\,G'\right]_{I'',I}
=\sum_{I',J'\subset\mathcal{J}(c')}\mathfrak{N}_{c',c''}[G']_{I'',I'}\cdot\Upsilon_{I'J'}\cdot\mathfrak{N}_{c,c'}[G]_{J',I},
\]
and the right side is the entry of $\mathfrak{N}_{c',c''}[G']\,\Upsilon\,\mathfrak{N}_{c,c'}[G]$ at $(I'',I)$.

\emph{Adjoint.}
Let $I\subset\mathcal{J}(c)$ and $I'\subset\mathcal{J}(c')$.
By Definition~\ref{def:mode-translations}\eqref{it:reflected-modes},
$\Theta_{c''}\left(\zeta^{(c'')}_{I+v_{c,c'}}\right)=\Theta_{c'}\left(\zeta^{(c)}_I\right)$.
The map $\Theta_{c'}$ is antilinear and reverses the order of factors (Lemma~\ref{lem:slab-automorphisms}),
and $\Theta_{c'}^2=1$ (Lemma~\ref{lem:slab-automorphisms}\eqref{it:parallel-reflections});
it follows that
\[
\zeta^{(c')}_{I'}\,\Theta_{c'}G\,\Theta_{c''}\left(\zeta^{(c'')}_{I+v_{c,c'}}\right)
=\zeta^{(c')}_{I'}\,\Theta_{c'}G\,\Theta_{c'}\left(\zeta^{(c)}_I\right)
=\Theta_{c'}\left(\zeta^{(c)}_I\,G\,\Theta_{c'}\left(\zeta^{(c')}_{I'}\right)\right).
\]
The argument of $\Theta_{c'}$ on the right is an element of $\Gr_\Xi((c,c'))$,
and the reflection $\theta_{c'}$ maps $\mathbb{X}_\Xi((c,c'))$ onto $\mathbb{X}_\Xi((c',2c'-c))=\mathbb{X}_\Xi((c',c''))$.
The reflection identity of~\eqref{eq:slab-integral-covariance}, taken at the position $c'$, now gives
\[
\begin{aligned}
\mathfrak{N}_{c',c''}[\Theta_{c'}G]_{I+v_{c,c'},\,I'}
&=\int_{(c',c'')}\Theta_{c'}\left(\zeta^{(c)}_I\,G\,\Theta_{c'}\left(\zeta^{(c')}_{I'}\right)\right)\\
&=\overline{\int_{(c,c')}\zeta^{(c)}_I\,G\,\Theta_{c'}\left(\zeta^{(c')}_{I'}\right)}
=\overline{\mathfrak{N}_{c,c'}[G]_{I',I}} .
\end{aligned}
\]
As to the operator form, $v_{c,c'}$ is an integer multiple of $n_\ell$,
so that $U_{v_{c,c'}}e_I=e_{I+v_{c,c'}}$ (Definition~\ref{def:mode-translations}\eqref{it:fock-translations}).
The entry of $U_{v_{c,c'}}\,\mathfrak{N}_{c,c'}[G]^{*}$ at $(I+v_{c,c'},I')$ is therefore the entry of $\mathfrak{N}_{c,c'}[G]^{*}$ at $(I,I')$,
which is $\overline{\mathfrak{N}_{c,c'}[G]_{I',I}}$.
Since $\mathcal{J}(c'')=\mathcal{J}(c)+v_{c,c'}$ by the choice of the representatives (Definition~\ref{def:mode-translations}),
every subset of $\mathcal{J}(c'')$ is of the form $I+v_{c,c'}$ for exactly one $I\subset\mathcal{J}(c)$,
and the two operators have the same entries.

\emph{Translation.}
Let $v\in\Z^3$, $I\subset\mathcal{J}(c)$ and $I'\subset\mathcal{J}(c')$.
The automorphism $\tau_v$ carries $\Gr_\Xi((c,c'))$ onto $\Gr_\Xi((c+\ell(v),c'+\ell(v)))$ (Lemma~\ref{lem:slab-automorphisms}),
and $\tau_v\Theta_{c'}=\Theta_{c'+\ell(v)}\tau_v$ by Lemma~\ref{lem:slab-automorphisms}\eqref{it:translated-reflection}.
Moreover
\[
\tau_v\zeta^{(c)}_I=s_v(I)\cdot\zeta^{(c+\ell(v))}_{r_v(I)},\qquad
\tau_v\zeta^{(c')}_{I'}=s_v(I')\cdot\zeta^{(c'+\ell(v))}_{r_v(I')}
\]
(Definition~\ref{def:mode-translations}),
with signs $s_v(I),s_v(I')\in\{1,-1\}$, which the antilinear map $\Theta_{c'+\ell(v)}$ leaves unchanged.
With the translation identity of~\eqref{eq:slab-integral-covariance} one finds
\[
\begin{aligned}
\mathfrak{N}_{c,c'}[G]_{I',I}
&=\int_{(c+\ell(v),\,c'+\ell(v))}\tau_v\left(\zeta^{(c)}_I\,G\,\Theta_{c'}\left(\zeta^{(c')}_{I'}\right)\right)\\
&=\int_{(c+\ell(v),\,c'+\ell(v))}\left(\tau_v\zeta^{(c)}_I\right)\left(\tau_vG\right)\,\Theta_{c'+\ell(v)}\left(\tau_v\zeta^{(c')}_{I'}\right)\\
&=s_v(I)\cdot s_v(I')\cdot\mathfrak{N}_{c+\ell(v),\,c'+\ell(v)}[\tau_vG]_{r_v(I'),\,r_v(I)} .
\end{aligned}
\]
On the other hand $U_ve_I=s_v(I)\cdot e_{r_v(I)}$ and $U_ve_{I'}=s_v(I')\cdot e_{r_v(I')}$,
hence $U_v^{-1}e_{r_v(I)}=s_v(I)\cdot e_I$,
and the entry of $U_v\,\mathfrak{N}_{c,c'}[G]\,U_v^{-1}$ at $(r_v(I'),r_v(I))$ is $s_v(I)\cdot s_v(I')\cdot\mathfrak{N}_{c,c'}[G]_{I',I}$.
The squares of the signs being one, this is the entry of $\mathfrak{N}_{c+\ell(v),\,c'+\ell(v)}[\tau_vG]$ at the same place.
Every pair of subsets of $\mathcal{J}(c+\ell(v))$ and $\mathcal{J}(c'+\ell(v))$ arises as $(r_v(I),r_v(I'))$,
because $r_v$ maps $\mathcal{J}(c)$ bijectively onto $\mathcal{J}(c+\ell(v))$ and $\mathcal{J}(c')$ bijectively onto $\mathcal{J}(c'+\ell(v))$.

\emph{Closure.}
By Definition~\ref{def:slabs} and the hypothesis $\Xi(\mathbb{T},\ell)=\Xi$,
one has $\Gamma_\parallel=\Gamma\cap\{\ell=0\}$, $\Gamma=\Gamma_\parallel+\Z g$ and $p=\ell(g)>0$.
The number $p$ is an integer and $\Adm(\ell)+\Z=\Adm(\ell)$,
so every point of $c_0+p\Z$ is an admissible position;
since mirrors carry no sites (Lemma~\ref{lem:slab-bonds}\eqref{it:bonds-crossing}), the height of every site lies outside $c_0+p\Z$.
The heights of the sites of a class $x+\Gamma$ form the set $\ell(x)+p\Z$;
exactly one of these numbers lies in $(c_0,c_0+p)$,
and the sites of $x+\Gamma$ at this height form a single class modulo $\Gamma_\parallel$.
Thus every class of $\mathbb{X}_{\mathbb{T}}$ contains exactly one class of $\mathbb{X}_\Xi((c_0,c_r))$.
Consequently the reading maps $\Gr_\Xi((c_0,c_r))$ isomorphically onto $\Gr_{\mathbb{T}}$,
since it sends a set of free generators of the one onto a set of free generators of the other (Definition~\ref{def:tori});
and the integral over $\mathbb{T}$ of the reading of an element of $\Gr_\Xi((c_0,c_r))$ equals its integral over $(c_0,c_r)$,
both Berezin integrals being formed with the same representatives.
From here on an element of $\Gr_\Xi((c_0,c_r))$ is identified with its reading;
thus $\Gr_{\mathbb{T}}=\Gr_\Xi((c_0,c_r))$, and $\int_{\mathbb{T}}$ agrees with $\int_{(c_0,c_r)}$ on this algebra.

Below $c_0$ the reading takes the following form.
In $\Gr_{\mathbb{T}}$ one has $\psi_\alpha(x+g)=\chi(g)\cdot\psi_\alpha(x)$,
so $\psi_\alpha(x)=\chi(-g)\cdot\psi_\alpha(x+g)=\chi(g)\cdot\psi_\alpha(x+g)$,
and likewise $\bar\psi_\alpha(x)=\chi(g)\cdot\bar\psi_\alpha(x+g)$.
If $\ell(x)\in(c_0-p,c_0)$, the site $x+g$ has its height in $(c_0,c_r)$.
Applying this to every factor, one sees that a product of $m$ generators at sites with heights in $(c_0-p,c_0)$
is read as $\chi(g)^m$ times its image under $\tau_g$, an element of $\Gr_\Xi((c_0,c_r))$.
If $G\in\Gr_\Xi((c_0-p,c_0))$ is homogeneous, the degrees of all monomials that occur in $G$ lie in the class $\deg G$;
so $G$ is read as the element $\chi(g)^{\deg G}\cdot\tau_gG$ of $\Gr_\Xi((c_0,c_r))$.
Now let $J\subset\mathcal{J}(c_0)$.
The element $\Theta_{c_0}\left(\zeta^{(c_0)}_J\right)$ is homogeneous of degree $\#|J|$ and lies in $\Gr_\Xi((c_0^{-},c_0))$,
which is contained in $\Gr_\Xi((c_0-p,c_0))$ because $p\ge c_1-c_0\ge|n_\ell|^2/2$.
By Lemma~\ref{lem:slab-automorphisms}\eqref{it:translated-reflection}, $\tau_g\Theta_{c_0}=\Theta_{c_0+p}\tau_g=\Theta_{c_r}\tau_g$.
By Definition~\ref{def:mode-translations}\eqref{it:fock-translations},
\[
\tau_g\zeta^{(c_0)}_J=\sum_{J'\subset\mathcal{J}(c_r)}(U_g)_{J'J}\cdot\zeta^{(c_r)}_{J'},
\]
and the entries of $U_g$ are real, so that the antilinear map $\Theta_{c_r}$ does not conjugate them.
Therefore $\Theta_{c_0}\left(\zeta^{(c_0)}_J\right)$ is read as
\begin{equation}\label{eq:slab-closure-reading}
\chi(g)^{\#|J|}\cdot\Theta_{c_r}\left(\tau_g\zeta^{(c_0)}_J\right)
=\chi(g)^{\#|J|}\sum_{J'\subset\mathcal{J}(c_r)}(U_g)_{J'J}\cdot\Theta_{c_r}\left(\zeta^{(c_r)}_{J'}\right).
\end{equation}

For $I_0\subset\mathcal{J}(c_0)$ put
\[
G_{I_0}:=\zeta^{(c_0)}_{I_0}\,G_1\,\Ins_{c_1}[\Upsilon_1]\,G_2\cdots\Ins_{c_{r-1}}[\Upsilon_{r-1}]\,G_r ,
\]
an element of $\Gr_\Xi((c_0,c_r))$, for the reason given in the proof of part~\eqref{it:matrix-composition}.
Expanding the insertion at $c_0$ (Definition~\ref{def:slab-matrices}),
one writes the product under the integral in part~\eqref{it:matrix-closure} as the sum of the elements
$(\Upsilon_0)_{I_0J_0}\cdot\Theta_{c_0}\left(\zeta^{(c_0)}_{J_0}\right)G_{I_0}$, $I_0,J_0\subset\mathcal{J}(c_0)$, of $\Gr_\Xi((c_0-p,c_r))$,
and the left side of the identity of that part as
\[
\sum_{I_0,J_0\subset\mathcal{J}(c_0)}(\Upsilon_0)_{I_0J_0}\cdot\int_{\mathbb{T}}\Theta_{c_0}\left(\zeta^{(c_0)}_{J_0}\right)\,G_{I_0} .
\]
An even operator $\Upsilon$ commutes with $U_{\mathrm{f}}$ (Definition~\ref{def:modes}),
which means that $\Upsilon_{IJ}=0$ unless $\#|I|+\#|J|$ is even;
each insertion $\Ins_{c_k}[\Upsilon_k]$ is therefore an even element.

Assume now that $\sum_k\deg G_k$ is even.
Then $G_{I_0}$, a product of homogeneous factors in which every insertion is even,
is homogeneous, of the parity of $\#|I_0|+\sum_k\deg G_k$, that is, of the parity of $\#|I_0|$.
Where $(\Upsilon_0)_{I_0J_0}\neq0$, the numbers $\#|I_0|$ and $\#|J_0|$ have the same parity,
and the element $\Theta_{c_0}\left(\zeta^{(c_0)}_{J_0}\right)$, homogeneous of the parity of $\#|J_0|$,
moves from the left of $G_{I_0}$ to its right at the cost of the sign $(-1)^{\#|J_0|\cdot\#|J_0|}=(-1)^{\#|J_0|}$.
Inserting~\eqref{eq:slab-closure-reading} and replacing $\int_{\mathbb{T}}$ by $\int_{(c_0,c_r)}$,
one obtains for the left side the expression
\[
\begin{aligned}
&\sum_{I_0,J_0\subset\mathcal{J}(c_0)}\ \sum_{J'\subset\mathcal{J}(c_r)}
(\Upsilon_0)_{I_0J_0}\cdot\left(-\chi(g)\right)^{\#|J_0|}\cdot(U_g)_{J'J_0}\\
&\qquad\cdot\int_{(c_0,c_r)}\zeta^{(c_0)}_{I_0}\,G_1\,\Ins_{c_1}[\Upsilon_1]\,G_2\cdots\Ins_{c_{r-1}}[\Upsilon_{r-1}]\,G_r\,\Theta_{c_r}\left(\zeta^{(c_r)}_{J'}\right).
\end{aligned}
\]
By Definition~\ref{def:slab-matrices} the integral in this expression is the entry at $(J',I_0)$ of the slab matrix
$\mathfrak{N}_{c_0,c_r}\left[G_1\,\Ins_{c_1}[\Upsilon_1]\,G_2\cdots\Ins_{c_{r-1}}[\Upsilon_{r-1}]\,G_r\right]$.
For $k=1,\dots,r-1$,
part~\eqref{it:matrix-composition}, applied to the positions $c_0<c_k<c_{k+1}$,
to the element $G_1\,\Ins_{c_1}[\Upsilon_1]\cdots G_k$ of $\Gr_\Xi((c_0,c_k))$, to the operator $\Upsilon_k$ and to $G_{k+1}$, yields
\[
\mathfrak{N}_{c_0,c_{k+1}}\left[G_1\,\Ins_{c_1}[\Upsilon_1]\cdots G_k\,\Ins_{c_k}[\Upsilon_k]\,G_{k+1}\right]
=\mathfrak{N}_{k+1}\Upsilon_k\,\mathfrak{N}_{c_0,c_k}\left[G_1\,\Ins_{c_1}[\Upsilon_1]\cdots G_k\right].
\]
Starting from $\mathfrak{N}_{c_0,c_1}[G_1]=\mathfrak{N}_1$,
these $r-1$ identities show that the slab matrix in question equals $\mathfrak{N}_r\Upsilon_{r-1}\mathfrak{N}_{r-1}\cdots\Upsilon_1\mathfrak{N}_1$.
The unitary $U_g$ has real entries, so that $(U_g)_{J'J_0}=(U_g^{-1})_{J_0J'}$;
and $\left(-\chi(g)\right)^{\#|J_0|}$ is the diagonal entry at $J_0$ of the operator $U_{\mathrm{f}}^{\frac{1+\chi(g)}{2}}$ on $\mathcal{F}(c_0)$,
namely $1$ if $\chi(g)=-1$ and $(-1)^{\#|J_0|}$ if $\chi(g)=+1$.
The left side is thus equal to
\[
\begin{aligned}
&\sum_{I_0\subset\mathcal{J}(c_0)}\left(\Upsilon_0\,U_{\mathrm{f}}^{\frac{1+\chi(g)}{2}}\,U_g^{-1}\,\mathfrak{N}_r\Upsilon_{r-1}\mathfrak{N}_{r-1}\cdots\Upsilon_1\mathfrak{N}_1\right)_{I_0,I_0}\\
&\qquad=\operatorname{tr}_{\mathcal{F}(c_0)}\left(U_{\mathrm{f}}^{\frac{1+\chi(g)}{2}}\,U_g^{-1}\,\mathfrak{N}_r\Upsilon_{r-1}\mathfrak{N}_{r-1}\cdots\Upsilon_1\mathfrak{N}_1\Upsilon_0\right),
\end{aligned}
\]
by the invariance of the trace under cyclic permutations.
Finally, $U_g$ commutes with the parity (Definition~\ref{def:mode-translations}\eqref{it:fock-translations}),
so that $U_{\mathrm{f}}^{\frac{1+\chi(g)}{2}}U_g^{-1}=U_g^{-1}U_{\mathrm{f}}^{\frac{1+\chi(g)}{2}}$,
with the parity of $\mathcal{F}(c_r)$ on the right;
the last trace is the right side of the asserted identity.

Let, on the contrary, $\sum_k\deg G_k$ be odd.
The insertions being even, the product under the integral on the left is homogeneous and odd,
and so is its reading in $\Gr_{\mathbb{T}}$, which maps generators to generators;
its Berezin integral vanishes by Proposition~\ref{prop:berezin}\eqref{it:berezin-factorisation}.
On the right, the operators $\Upsilon_k$, $U_g^{-1}$ and $U_{\mathrm{f}}^{\frac{1+\chi(g)}{2}}$ commute with the parity,
whereas $U_{\mathrm{f}}\,\mathfrak{N}_k=(-1)^{\deg G_k}\cdot\mathfrak{N}_k\,U_{\mathrm{f}}$ by part~\eqref{it:matrix-parity}.
Carrying $U_{\mathrm{f}}$ from the left of the whole product inside the trace to its right
thus produces the sign $(-1)^{\sum_k\deg G_k}=-1$.
Conjugation by the involution $U_{\mathrm{f}}$ leaves the trace unchanged and turns this product into its negative;
its trace is therefore zero.
\end{proof}

\subsection{The photon at a mirror}\label{sec:mirror-photon}

As before, the family $\ell\in\mathcal{L}$ and the cross-section $\Xi=(\Gamma_\parallel,\chi_\parallel)$ for $\ell$ are fixed,
and cells and cochains are those of the cylinder $\Z^3/\Gamma_\parallel$.
For $c\in\Adm(\ell)$ the modes at $c$ are labelled by the ordered set $\mathcal{J}(c)$ of Definition~\ref{def:mode-translations},
and the sets $\mathcal{X}^1(c)$, $\mathcal{X}^2(c)$, $\mathcal{X}^3(c)$ of the links, plaquettes and cubes cut by $c$ are finite (Lemma~\ref{lem:coboundary-mirror}).
Upper and lower faces of a cell cut by $c$ are meant in the sense of Lemma~\ref{lem:cells-mirrors}\eqref{it:cut-faces},
and $d^{+}_c$, $d_c$ and the adjoint $d_c^{*}$ are the operators of Lemma~\ref{lem:coboundary-mirror}.
On each of the spaces $\R^{\mathcal{X}^1(c)}$, $\R^{\mathcal{X}^2(c)}$ and $\R^{\mathcal{X}^3(c)}$ the counting pairing is written $\langle\cdot,\cdot\rangle$.

\begin{definition}[Fourier variables at a mirror]\label{def:fourier-variables}
Let $c\in\Adm(\ell)$.
\begin{enumerate}
\item
A \emph{Fourier variable} at $c$ is a pair
\[
(\vartheta,k)\in[0,2\pi)^{\mathcal{X}^3(c)}\times\Z^{\mathcal{X}^2(c)},
\]
that is, an angle $\vartheta_{\sigma'}$ for every cube $\sigma'$ cut by $c$ and an integer $k_\sigma$ for every plaquette $\sigma$ cut by $c$.
For $\sigma\in\mathcal{X}^2(c)$ put
\[
\varkappa_\sigma(\vartheta,k):=\frac{q}{2\pi}\cdot\left((d_c^{*}\vartheta)(\sigma)-2\pi k_\sigma\right),
\qquad\text{where}\qquad
(d_c^{*}\vartheta)(\sigma)=\sum_{\sigma'\in\mathcal{X}^3(c)}[\sigma':\sigma]\cdot\vartheta_{\sigma'} .
\]
\item
The measure $\mu^{\times}_c$ on $[0,2\pi)^{\mathcal{X}^3(c)}\times\Z^{\mathcal{X}^2(c)}$ is
\[
\mathrm{d}\mu^{\times}_c(\vartheta,k):=\prod_{\sigma'\in\mathcal{X}^3(c)}\frac{\mathrm{d}\vartheta_{\sigma'}}{2\pi}
\cdot\prod_{\sigma\in\mathcal{X}^2(c)}\frac{q}{\sqrt{2\pi}}\exp\left(-\frac{\varkappa_\sigma(\vartheta,k)^2}{2}\right),
\]
the variable $k$ carrying the counting measure.
\item
Let $\varpi=(A,z)$ consist of a real one-cochain $A$, given at least on the upper faces of the plaquettes cut by $c$,
and of an integer two-cochain $z$, given at least on the upper faces of the cubes cut by $c$;
for instance, $\varpi$ may be a configuration of a closed slab $\{c'\le\ell\le c''\}$ with $c'\le c<c''$ (Lemma~\ref{lem:cells-mirrors}\eqref{it:upper-faces}).
The \emph{phase} of $\varpi$ at $(\vartheta,k)$ is
\[
\mathfrak{e}_c\left((\vartheta,k);\varpi\right):=
\exp\left(i\sum_{\sigma'\in\mathcal{X}^3(c)}\vartheta_{\sigma'}\cdot(d^{+}_cz)(\sigma')
-i\sum_{\sigma\in\mathcal{X}^2(c)}\varkappa_\sigma(\vartheta,k)\cdot(d^{+}_cA)(\sigma)\right).
\]
\item
A mode is said to lie \emph{over} a link $b$ if the nearest photon sites of the two ends of its bond are the two end points of $b$.
For $I\subset\mathcal{J}(c)$ and $b\in\mathcal{X}^1(c)$ put
\[
\begin{aligned}
\mathfrak{q}(I)(b):={}&\#|\{i\in I:\ i\ \text{lies over}\ b\ \text{and is of group 1 or 2}\}|\\
&{}-\#|\{i\in I:\ i\ \text{lies over}\ b\ \text{and is of group 3 or 4}\}|,
\end{aligned}
\]
the groups being those of Definition~\ref{def:modes}.
The \emph{Gauss factor} is
\[
\mathfrak{D}_c(k,I):=(2\pi/q)^{\#|\mathcal{X}^1(c)|}\prod_{b\in\mathcal{X}^1(c)}\mathbf{1}\left[(d_c^{*}k)(b)+\mathfrak{q}(I)(b)=0\right],
\]
where $(d_c^{*}k)(b)=\sum_{\sigma\in\mathcal{X}^2(c)}[\sigma:b]\cdot k_\sigma$ for $b\in\mathcal{X}^1(c)$;
in particular $\mathfrak{D}_c=1$ if no link is cut by $c$.
\end{enumerate}
\end{definition}

A cell $\sigma$ cut by $c$ has $n_\ell$ in its direction space, and $\min_\sigma\ell=c^-$ is an integer (Lemma~\ref{lem:cells-mirrors}\eqref{it:cut-heights}).
Links are therefore cut by $c$ only if $\ell=x_\mu$ is an axis family and $c\in\Z+\frac12$ is a link plane;
in that case every $b\in\mathcal{X}^1(c)$ is parallel to $e_\mu=n_\ell$, and its orientation $e_\mu$ points upwards.
Let $c$ be such a link plane, let $i$ be a mode at $c$, and let $\{x,\theta_cx\}$ be its bond, with lower end $x$.
The points $x$ and $\theta_cx$ differ only in the coordinate $\mu$, and $x_\mu<c<(\theta_cx)_\mu$.
The bond is of type $3$, so that $\pi(\theta_cx)-\pi(x)$ is a unit vector (Proposition~\ref{prop:geometry}\eqref{it:geometry-projection}).
For every site $x'$, each coordinate of the nearest photon site $\pi(x')$ is the integer nearest to the corresponding coordinate of $x'$ (Definition~\ref{def:fermion-sites});
as $c\in\Z+\frac12$, this gives $\pi(x)_\mu\le c-\frac12$ and $\pi(\theta_cx)_\mu\ge c+\frac12$, and therefore
$\pi(\theta_cx)=\pi(x)+e_\mu$ with $\pi(x)_\mu=c-\frac12$.
Hence $i$ lies over exactly one link, the link $b\in\mathcal{X}^1(c)$ that leads from $\pi(x)$ to $\pi(\theta_cx)$.
By Definition~\ref{def:modes} the phase of $i$ is $u_i=U_A(x,\theta_cx)=\exp\left(iqA(b)\right)$ if $i$ is of group 1 or 2,
and $u_i=\exp\left(-iqA(b)\right)$ if $i$ is of group 3 or 4.
For $I\subset\mathcal{J}(c)$ it follows that
\[
u_I(A)=\exp\left(iq\sum_{b\in\mathcal{X}^1(c)}\mathfrak{q}(I)(b)\cdot A(b)\right).
\]
Thus $u_I(A)$ depends on $A$ only through the restriction of $A$ to $\mathcal{X}^1(c)$,
and for $A^{\times}\in\R^{\mathcal{X}^1(c)}$ it is written $u_I(A^{\times})$.
If no link is cut by $c$, then $c$ is not a link plane, and $u_I=1$ (Proposition~\ref{prop:crossing}), in accordance with the formula.

At a site plane, that is, for an axis family $\ell$ and $c\in\Z$, no cell is cut by $c$, since $c^-=c-\frac12$ is not an integer.
There the Fourier variables reduce to a single point, $\mu^{\times}_c$ is the unit mass, and $\mathfrak{e}_c=1$.
In general, the phase $\mathfrak{e}_c((\vartheta,k);\varpi)$ depends on $\varpi$ only through the values of $A$ on the upper faces of the plaquettes cut by $c$
and of $z$ on the upper faces of the cubes cut by $c$.

The measure $\mu^{\times}_c$ is positive and finite.
To see this, fix $\vartheta$ and $\sigma\in\mathcal{X}^2(c)$.
The number $\varkappa_\sigma(\vartheta,k)$ depends on $k$ only through $k_\sigma$,
and Lemma~\ref{lem:poisson-kernels}\eqref{it:poisson-summation}, applied with $y=0$ and $s=(d_c^{*}\vartheta)(\sigma)$, gives
\begin{equation}\label{eq:slab-fourier-mass}
\frac{q}{\sqrt{2\pi}}\sum_{k_\sigma\in\Z}\exp\left(-\frac{\varkappa_\sigma(\vartheta,k)^2}{2}\right)
=\sum_{j\in\Z}\exp\left(-\frac{2\pi^2j^2}{q^2}+ij\cdot(d_c^{*}\vartheta)(\sigma)\right)
\le\sum_{k\in\Z}\exp(-2\pi^2k^2/q^2) .
\end{equation}
The density of $\mu^{\times}_c$ with respect to the product of the measures $\mathrm{d}\vartheta_{\sigma'}/(2\pi)$ and of the counting measure
is a product over $\sigma\in\mathcal{X}^2(c)$ of strictly positive factors, the factor of $\sigma$ depending on $k$ only through $k_\sigma$.
Summing over $k$ by means of \eqref{eq:slab-fourier-mass}, and then integrating over $\vartheta$ against a probability measure, one obtains
\[
\mu^{\times}_c\left([0,2\pi)^{\mathcal{X}^3(c)}\times\Z^{\mathcal{X}^2(c)}\right)
\le\left(\sum_{k\in\Z}\exp(-2\pi^2k^2/q^2)\right)^{\#|\mathcal{X}^2(c)|}.
\]

\begin{lemma}[Crossing identity for the photon]\label{lem:fourier-crossing}
Let $c\in\Adm(\ell)$ and $I\subset\mathcal{J}(c)$.
Let $\varpi=(A,z)$, where $A$ is a real function on the upper and the lower faces of the plaquettes cut by $c$,
and $z$ is an integer valued function on the upper and the lower faces of the cubes cut by $c$.
The reflection $\theta_c$ maps the upper faces of a cut cell onto its lower faces and conversely (Lemma~\ref{lem:cells-mirrors}\eqref{it:cut-faces}),
so that $\theta_c^{*}\varpi=(\theta_c^{*}A,\theta_c^{*}z)$ is given on the same faces.
For $A^{\times}\in[0,2\pi/q)^{\mathcal{X}^1(c)}$ and $z^{\times}\in\Z^{\mathcal{X}^2(c)}$,
form $F_\sigma:=(dA)(\sigma)+\frac{2\pi}{q}z^{\times}_\sigma$ for $\sigma\in\mathcal{X}^2(c)$ and $(dz)(\sigma')$ for $\sigma'\in\mathcal{X}^3(c)$
with the values of $\varpi$ on the upper and the lower faces, of $A^{\times}$ on the links cut by $c$ and of $z^{\times}$ on the plaquettes cut by $c$;
every face of a plaquette or a cube cut by $c$ is of one of these kinds (Lemma~\ref{lem:cells-mirrors}\eqref{it:cut-faces}).
Then
\[
\begin{aligned}
&\sum_{z^{\times}\in\Z^{\mathcal{X}^2(c)}}\ \int_{[0,2\pi/q)^{\mathcal{X}^1(c)}}\mathrm{d}A^{\times}\,
\prod_{\sigma\in\mathcal{X}^2(c)}\exp\left(-\frac{F_\sigma^2}{2}\right)
\cdot\prod_{\sigma'\in\mathcal{X}^3(c)}\mathbf{1}\left[(dz)(\sigma')=0\right]\cdot u_I(A^{\times})\\
&\qquad=\int\mathrm{d}\mu^{\times}_c(\vartheta,k)\,\mathfrak{D}_c(k,I)\cdot\mathfrak{e}_c\left((\vartheta,k);\varpi\right)
\cdot\overline{\mathfrak{e}_c\left((\vartheta,k);\theta_c^{*}\varpi\right)} ,
\end{aligned}
\]
and both sides converge absolutely, uniformly in $\varpi$.
\end{lemma}

\begin{proof}
In what follows $\sigma$ ranges over $\mathcal{X}^2(c)$ and $\sigma'$ over $\mathcal{X}^3(c)$.
By Lemma~\ref{lem:coboundary-mirror}\eqref{it:coboundary-split},
\[
F_\sigma=A^{\partial}_\sigma+(d_cA^{\times})(\sigma)+\frac{2\pi}{q}z^{\times}_\sigma,
\qquad
(dz)(\sigma')=z^{\partial}_{\sigma'}+(d_cz^{\times})(\sigma'),
\]
where the numbers
\[
A^{\partial}_\sigma:=(d^{+}_cA)(\sigma)-(d^{+}_c\theta_c^{*}A)(\sigma),
\qquad
z^{\partial}_{\sigma'}:=(d^{+}_cz)(\sigma')-(d^{+}_c\theta_c^{*}z)(\sigma')
\]
depend on $\varpi$ alone, and each $z^{\partial}_{\sigma'}$ is an integer.

\emph{The constraint at the cut cubes.}
For every integer $j$ one has $\mathbf{1}[j=0]=\int_0^{2\pi}\frac{\mathrm{d}\vartheta}{2\pi}\,\exp(i\vartheta j)$.
Used at every cut cube, and combined with $\langle\vartheta,d_cz^{\times}\rangle=\langle d_c^{*}\vartheta,z^{\times}\rangle$, this gives
\[
\prod_{\sigma'\in\mathcal{X}^3(c)}\mathbf{1}\left[(dz)(\sigma')=0\right]
=\int_{[0,2\pi)^{\mathcal{X}^3(c)}}\prod_{\sigma'\in\mathcal{X}^3(c)}\frac{\mathrm{d}\vartheta_{\sigma'}}{2\pi}\,
\exp\left(i\langle\vartheta,z^{\partial}\rangle+i\langle d_c^{*}\vartheta,z^{\times}\rangle\right).
\]
After this substitution the left side of the lemma is a sum over $z^{\times}$ of an integral over $(A^{\times},\vartheta)$,
and the modulus of the integrand is the product over $\sigma$ of the numbers $\exp\left(-\frac{F_\sigma^2}{2}\right)$.
For fixed $A^{\times}$, the sum of this modulus over $z^{\times}$ is a product over $\sigma$ of sums of the form
\[
\sum_{j\in\Z}\exp\left(-\frac12\left(y+\frac{2\pi}{q}\,j\right)^2\right)
\]
with real $y$, each of which is at most $C_{q,1/2}$ by Lemma~\ref{lem:poisson-kernels}\eqref{it:gaussian-sums}.
As $A^{\times}$ and $\vartheta$ range over sets of finite measure, the left side converges absolutely:
the sum over $z^{\times}$ of the integrals of the moduli is at most $(2\pi/q)^{\#|\mathcal{X}^1(c)|}\cdot C_{q,1/2}^{\#|\mathcal{X}^2(c)|}$ for every $\varpi$.
By Fubini's theorem the sum over $z^{\times}$ may be carried out first, at fixed $(A^{\times},\vartheta)$.

\emph{Poisson summation at the cut plaquettes.}
Let $(A^{\times},\vartheta)$ be fixed.
The factors $\exp\left(i\langle\vartheta,z^{\partial}\rangle\right)$ and $u_I(A^{\times})$ of the integrand do not depend on $z^{\times}$.
The remaining factor is
\[
\prod_{\sigma\in\mathcal{X}^2(c)}\exp\left(-\frac{F_\sigma^2}{2}+iz^{\times}_\sigma\cdot(d_c^{*}\vartheta)(\sigma)\right),
\]
and its sum over $z^{\times}$ is the product over $\sigma$ of the sums over $z^{\times}_\sigma\in\Z$,
since $F_\sigma$ depends on $z^{\times}$ only through $z^{\times}_\sigma$.
For each $\sigma$, Lemma~\ref{lem:poisson-kernels}\eqref{it:poisson-summation}, applied with $y=A^{\partial}_\sigma+(d_cA^{\times})(\sigma)$ and $s=(d_c^{*}\vartheta)(\sigma)$, gives
\[
\begin{aligned}
&\sum_{z^{\times}_\sigma\in\Z}\exp\left(-\frac{F_\sigma^2}{2}+iz^{\times}_\sigma\cdot(d_c^{*}\vartheta)(\sigma)\right)\\
&\qquad=\frac{q}{\sqrt{2\pi}}\sum_{k_\sigma\in\Z}\exp\left(-\frac{\varkappa_\sigma^2}{2}\right)
\cdot\exp\left(-i\varkappa_\sigma\cdot A^{\partial}_\sigma\right)\cdot\exp\left(-i\varkappa_\sigma\cdot(d_cA^{\times})(\sigma)\right),
\end{aligned}
\]
with $\varkappa_\sigma=\varkappa_\sigma(\vartheta,k)$.
The product over $\sigma$ of these series is a single sum over $k\in\Z^{\mathcal{X}^2(c)}$.
By \eqref{eq:slab-fourier-mass} it converges absolutely, uniformly in $(A^{\times},\vartheta)$,
since all factors other than the Gaussian ones have modulus one.
Fubini's theorem therefore allows the sum over $k$ to be taken outside the integrals over $A^{\times}$ and $\vartheta$.

\emph{The cut links.}
Let $(\vartheta,k)$ be fixed and put $\varkappa:=\left(\varkappa_\sigma(\vartheta,k)\right)_{\sigma\in\mathcal{X}^2(c)}$.
The integrand now depends on $A^{\times}$ only through the factor
\[
\exp\left(-i\langle\varkappa,d_cA^{\times}\rangle\right)\cdot u_I(A^{\times})
=\exp\left(-i\langle d_c^{*}\varkappa,A^{\times}\rangle+iq\langle\mathfrak{q}(I),A^{\times}\rangle\right).
\]
By Lemma~\ref{lem:coboundary-mirror}\eqref{it:cut-complex}, $d_c^{*}d_c^{*}=(d_cd_c)^{*}=0$, so that
\[
d_c^{*}\varkappa=\frac{q}{2\pi}\cdot\left(d_c^{*}d_c^{*}\vartheta-2\pi\,d_c^{*}k\right)=-q\,d_c^{*}k .
\]
The function $d_c^{*}k+\mathfrak{q}(I)$ on $\mathcal{X}^1(c)$ is integer valued, and the integral of $\exp\left(iqjA^{\times}(b)\right)$
over $A^{\times}(b)\in[0,2\pi/q)$ equals $(2\pi/q)\cdot\mathbf{1}[j=0]$ for every integer $j$, because $q\cdot2\pi/q=2\pi$.
Hence
\[
\int_{[0,2\pi/q)^{\mathcal{X}^1(c)}}\mathrm{d}A^{\times}\,\exp\left(iq\langle d_c^{*}k+\mathfrak{q}(I),A^{\times}\rangle\right)=\mathfrak{D}_c(k,I).
\]
If no link is cut by $c$, there is no variable $A^{\times}$, and both sides of this identity equal $1$.

\emph{The upper and the lower faces.}
Besides the density of $\mu^{\times}_c$ and the Gauss factor, the factors that remain are
\[
\exp\left(i\langle\vartheta,z^{\partial}\rangle\right)\cdot\prod_{\sigma\in\mathcal{X}^2(c)}\exp\left(-i\varkappa_\sigma\cdot A^{\partial}_\sigma\right)
=\mathfrak{e}_c\left((\vartheta,k);\varpi\right)\cdot\overline{\mathfrak{e}_c\left((\vartheta,k);\theta_c^{*}\varpi\right)} .
\]
Indeed, once the definitions of $z^{\partial}$ and $A^{\partial}$ are inserted,
the terms with $d^{+}_cz$ and $d^{+}_cA$ make up the exponent of $\mathfrak{e}_c((\vartheta,k);\varpi)$,
while the terms with $d^{+}_c\theta_c^{*}z$ and $d^{+}_c\theta_c^{*}A$ make up the exponent of the complex conjugate of $\mathfrak{e}_c((\vartheta,k);\theta_c^{*}\varpi)$,
the angles $\vartheta_{\sigma'}$ and the numbers $\varkappa_\sigma$ being real.

Combining the four computations, one obtains for the left side of the lemma the expression
\[
\begin{aligned}
&\sum_{k\in\Z^{\mathcal{X}^2(c)}}\ \int_{[0,2\pi)^{\mathcal{X}^3(c)}}\prod_{\sigma'\in\mathcal{X}^3(c)}\frac{\mathrm{d}\vartheta_{\sigma'}}{2\pi}\,
\prod_{\sigma\in\mathcal{X}^2(c)}\frac{q}{\sqrt{2\pi}}\exp\left(-\frac{\varkappa_\sigma(\vartheta,k)^2}{2}\right)\\
&\qquad\cdot\mathfrak{D}_c(k,I)\cdot\mathfrak{e}_c\left((\vartheta,k);\varpi\right)\cdot\overline{\mathfrak{e}_c\left((\vartheta,k);\theta_c^{*}\varpi\right)},
\end{aligned}
\]
which is the right side.
The right side converges absolutely, uniformly in $\varpi$:
the two phases have modulus one, their exponents being purely imaginary,
$0\le\mathfrak{D}_c(k,I)\le(2\pi/q)^{\#|\mathcal{X}^1(c)|}$, and $\mu^{\times}_c$ is a finite measure.
\end{proof}

\begin{lemma}\label{lem:phase-properties}
Let $c\in\Adm(\ell)$ and let $(\vartheta,k)$ be a Fourier variable at $c$.
\begin{enumerate}
\item\label{it:phase-modulus}
$\left|\mathfrak{e}_c((\vartheta,k);\varpi)\right|=1$ for every $\varpi$,
and for fixed $k$ and $z$ the phase $\mathfrak{e}_c((\vartheta,k);(A,z))$ is a continuous function of $(\vartheta,A)$.
\item\label{it:phase-shifts}
Let $c'<c''$ be elements of $\Adm(\ell)$, let $\varpi$ be a configuration of the closed slab $\{c'\le\ell\le c''\}$,
and let $n\in\Z^{[c',c'']^{(1)}}$ be an integer one-cochain.
If $c'\le c<c''$ and $n$ vanishes on the links contained in $\{\ell\le c\}$ and on the links cut by $c$, then
\[
\mathfrak{e}_c\left((\vartheta,k);\varpi+[n]\right)=\mathfrak{e}_c\left((\vartheta,k);\varpi\right).
\]
If $c'<c\le c''$ and $n$ vanishes on the links contained in $\{\ell\ge c\}$ and on the links cut by $c$, then
\[
\mathfrak{e}_c\left((\vartheta,k);\theta_c^{*}(\varpi+[n])\right)=\mathfrak{e}_c\left((\vartheta,k);\theta_c^{*}\varpi\right).
\]
\item\label{it:phase-translations}
Let $v\in\Z^3$.
The translation $\sigma\mapsto\sigma+v$ maps the cells cut by $c$ bijectively onto the cells cut by $c+\ell(v)$, preserving dimensions.
Accordingly, for a Fourier variable $(\vartheta',k')$ at $c+\ell(v)$,
the pair $\tau_v^{*}(\vartheta',k'):=(\tau_v^{*}\vartheta',\tau_v^{*}k')$, where $(\tau_v^{*}u)(\sigma):=u(\sigma+v)$, is a Fourier variable at $c$.
For every such $(\vartheta',k')$, every $\sigma\in\mathcal{X}^2(c)$, every $\varpi'=(A',z')$ given on the upper faces of the plaquettes and the cubes cut by $c+\ell(v)$,
with $\tau_v^{*}\varpi':=(\tau_v^{*}A',\tau_v^{*}z')$, and every $I\subset\mathcal{J}(c)$,
\[
\begin{aligned}
\varkappa_{\sigma+v}(\vartheta',k')&=\varkappa_\sigma\left(\tau_v^{*}(\vartheta',k')\right),\\
\mathfrak{e}_{c+\ell(v)}\left((\vartheta',k');\varpi'\right)&=\mathfrak{e}_c\left(\tau_v^{*}(\vartheta',k');\tau_v^{*}\varpi'\right),\\
\mathfrak{D}_{c+\ell(v)}\left(k',r_v(I)\right)&=\mathfrak{D}_c\left(\tau_v^{*}k',I\right).
\end{aligned}
\]
Moreover, the map $(\vartheta',k')\mapsto\tau_v^{*}(\vartheta',k')$ carries the measure $\mu^{\times}_{c+\ell(v)}$ to $\mu^{\times}_c$.
\end{enumerate}
\end{lemma}

\begin{proof}
\eqref{it:phase-modulus}
The angles $\vartheta_{\sigma'}$, the numbers $\varkappa_\sigma(\vartheta,k)$ and the values of $d^{+}_cz$ and of $d^{+}_cA$ are real,
so the exponent in the definition of $\mathfrak{e}_c$ is purely imaginary.
For fixed $k$ and $z$ this exponent is a polynomial of degree at most two in $(\vartheta,A)$,
since $\varkappa_\sigma(\vartheta,k)$ is an affine function of $\vartheta$ and $(d^{+}_cA)(\sigma)$ a linear function of $A$;
hence the phase is continuous in $(\vartheta,A)$.

\eqref{it:phase-shifts}
Let $c'\le c<c''$, and let $n$ vanish on the links contained in $\{\ell\le c\}$ and on the links cut by $c$.
The phase is defined at $\varpi$ and at $\varpi+[n]$, the upper faces of the cells cut by $c$ being cells of the closed slab
(Lemma~\ref{lem:cells-mirrors}\eqref{it:upper-faces}).
Extend $n$ by zero to the links of the cylinder that are not contained in the closed slab.
The extension still vanishes on the links contained in $\{\ell\le c\}$ and on the links cut by $c$,
and its coboundary agrees with $dn$ on the plaquettes of the closed slab, whose links belong to the slab.
For $\sigma\in\mathcal{X}^2(c)$ the number $(d^{+}_cn)(\sigma)$ is an integer, and
\[
\varkappa_\sigma(\vartheta,k)\cdot\frac{2\pi}{q}\cdot(d^{+}_cn)(\sigma)
=(d_c^{*}\vartheta)(\sigma)\cdot(d^{+}_cn)(\sigma)-2\pi k_\sigma\cdot(d^{+}_cn)(\sigma),
\]
where the last term is an integer multiple of $2\pi$.
For $\sigma'\in\mathcal{X}^3(c)$ one has $(d^{+}_c(z-dn))(\sigma')=(d^{+}_cz)(\sigma')-(d^{+}_cdn)(\sigma')$.
By the linearity of $d^{+}_c$, the phase $\mathfrak{e}_c((\vartheta,k);\varpi+[n])$ is therefore $\mathfrak{e}_c((\vartheta,k);\varpi)$ times
\[
\begin{aligned}
&\exp\left(-i\sum_{\sigma'\in\mathcal{X}^3(c)}\vartheta_{\sigma'}\cdot(d^{+}_cdn)(\sigma')
-i\sum_{\sigma\in\mathcal{X}^2(c)}(d_c^{*}\vartheta)(\sigma)\cdot(d^{+}_cn)(\sigma)\right)\\
&\qquad=\exp\left(-i\sum_{\sigma'\in\mathcal{X}^3(c)}\vartheta_{\sigma'}\cdot
\left((d^{+}_cdn)(\sigma')+\sum_{\sigma\in\mathcal{X}^2(c)}[\sigma':\sigma]\cdot(d^{+}_cn)(\sigma)\right)\right),
\end{aligned}
\]
the formula for $d_c^{*}$ having been inserted.
For every $\sigma'\in\mathcal{X}^3(c)$ the expression in the inner parentheses vanishes,
by Lemma~\ref{lem:coboundary-mirror}\eqref{it:cut-identity} applied to the extension of $n$.

Let now $c'<c\le c''$, and let $n$ vanish on the links contained in $\{\ell\ge c\}$ and on the links cut by $c$.
The pull-back $\theta_c^{*}$ commutes with $d$ by \eqref{eq:slab-incidence-isometry};
hence $\theta_c^{*}(\varpi+[n])=\theta_c^{*}\varpi+[\theta_c^{*}n]$,
a configuration of the closed slab $\{2c-c''\le\ell\le2c-c'\}$, and $2c-c''\le c<2c-c'$.
Since $\ell\circ\theta_c=2c-\ell$ (Definition~\ref{def:mirrors}\eqref{it:mirror-height}),
the reflection $\theta_c$ maps every link contained in $\{\ell\le c\}$ onto a link contained in $\{\ell\ge c\}$,
and it maps every link cut by $c$ onto itself (Lemma~\ref{lem:cells-mirrors}\eqref{it:cut-orientation}).
So $\theta_c^{*}n$ vanishes on the links contained in $\{\ell\le c\}$ and on the links cut by $c$,
and the first statement, applied to $\theta_c^{*}\varpi$ and $\theta_c^{*}n$, gives the second.

\eqref{it:phase-translations}
The number $\ell(v)$ is an integer, so $c+\ell(v)\in\Adm(\ell)$ and $(c+\ell(v))^{+}=c^{+}+\ell(v)$.
As $\ell(x+v)=\ell(x)+\ell(v)$ for every point $x$, the translation by $v$ maps
the cells contained in $\{\ell=c\}$ onto those contained in $\{\ell=c+\ell(v)\}$,
the cells cut by $c$ onto the cells cut by $c+\ell(v)$,
and the upper faces of a cell cut by $c$ onto the upper faces of its translate;
it commutes with $\Gamma_\parallel$ and therefore acts on the cylinder.
A translation preserves orientations, and $[\sigma+v:\sigma''+v]=[\sigma:\sigma'']$ for all cells $\sigma$, $\sigma''$ by \eqref{eq:slab-incidence-isometry}.
Consequently, for every cell $\sigma$ cut by $c$ and every function $u'$ on the appropriate cells at $c+\ell(v)$,
\[
\left(d_{c+\ell(v)}^{*}u'\right)(\sigma+v)=\left(d_c^{*}\tau_v^{*}u'\right)(\sigma),
\qquad
\left(d^{+}_{c+\ell(v)}u'\right)(\sigma+v)=\left(d^{+}_c\tau_v^{*}u'\right)(\sigma).
\]
The first identity for $u'=\vartheta'$, together with $k'_{\sigma+v}=(\tau_v^{*}k')_\sigma$, is the formula for $\varkappa$.
Inserting this formula and the second identity, for $u'=z'$ and for $u'=A'$, into the definition of $\mathfrak{e}_{c+\ell(v)}$,
and writing the summation variables as translates $\sigma+v$, $\sigma'+v$ of cells cut by $c$, one obtains the formula for the phase.
The map $\tau_v^{*}$ is a bijection from the Fourier variables at $c+\ell(v)$ onto those at $c$;
it carries the normalised Lebesgue measures of the angles and the counting measure of the integers at $c+\ell(v)$ to the corresponding measures at $c$,
and by the formula for $\varkappa$ it carries the density of $\mu^{\times}_{c+\ell(v)}$ to that of $\mu^{\times}_c$.
Finally, $\#|\mathcal{X}^1(c+\ell(v))|=\#|\mathcal{X}^1(c)|$, and for $b\in\mathcal{X}^1(c)$ one has
$\left(d_{c+\ell(v)}^{*}k'\right)(b+v)=\left(d_c^{*}\tau_v^{*}k'\right)(b)$ and
\[
\mathfrak{q}\left(r_v(I)\right)(b+v)=\mathfrak{q}(I)(b).
\]
For the last identity note that the shift $i\mapsto i+v$ of Definition~\ref{def:mode-translations} preserves the group of a mode
and carries a mode over $b$ to a mode over $b+v$,
and that $i+v$ and $r_v(i)$ differ by an element of $\Gamma_\parallel$, so that they lie over the same link of the cylinder;
moreover $r_v$ maps $\mathcal{J}(c)$ bijectively onto $\mathcal{J}(c+\ell(v))$.
With these three equalities the factors of $\mathfrak{D}_{c+\ell(v)}(k',r_v(I))$ and of $\mathfrak{D}_c(\tau_v^{*}k',I)$ agree link by link.
\end{proof}

\begin{remark}\label{rem:positivity-source}
The right side of Lemma~\ref{lem:fourier-crossing} is the integral, against the positive measure $\mathfrak{D}_c(k,I)\,\mathrm{d}\mu^{\times}_c(\vartheta,k)$,
of the function $\mathfrak{e}_c((\vartheta,k);\cdot)$ of the configuration on the upper faces
times the complex conjugate of the same function at the reflected configuration $\theta_c^{*}\varpi$.
Three facts produce this form.
The Fourier transform of the Maxwell weight $\exp\left(-\frac{s^2}{2}\right)$ is positive, a positive multiple of $\exp\left(-\frac{\varkappa^2}{2}\right)$;
the constraint $dz=0$ at the cut cubes is an average of characters with respect to the positive weight $\mathrm{d}\vartheta/(2\pi)$;
and the integral over the variable of a cut link is an indicator, the factor of Gauss' law.
The fermion contributes the coefficients $\omega_I$ of the crossing expansion of Proposition~\ref{prop:crossing},
and these are non-negative because $|\rho|\le1$.
\end{remark}

\subsection{Transfer across slabs}\label{sec:slab-transfer}

Throughout this subsection a family $\ell\in\mathcal{L}$ and a cross-section $\Xi=(\Gamma_\parallel,\chi_\parallel)$ for $\ell$ are fixed,
and cells, cochains, fermion sites and modes are those of the cylinder $\Z^3/\Gamma_\parallel$.
Since $-1\le\rho\le1$, every crossing coefficient $\omega_i$ is non-negative,
and so is every product $\omega_I$ (Definition~\ref{def:modes}).

\begin{definition}[Mirror spaces]\label{def:mirror-spaces}
Let $c\in\Adm(\ell)$.
The \emph{mirror variables} at $c$ are the quadruples $\xi=(\mathfrak{m},\vartheta,k,I)$
formed by mirror data $\mathfrak{m}=(A,z)\in[0,2\pi/q)^{\mathcal{M}^1(c)}\times\Z^{\mathcal{M}^2(c)}$ (Definition~\ref{def:slab-configurations}),
Fourier variables $(\vartheta,k)\in[0,2\pi)^{\mathcal{X}^3(c)}\times\Z^{\mathcal{X}^2(c)}$ (Definition~\ref{def:fourier-variables})
and a set $I\subset\mathcal{J}(c)$ of modes (Definition~\ref{def:mode-translations});
their set is denoted $\mathcal{Y}_c$.
On $\mathcal{Y}_c$ the measure $\mu_c$ is
\[
\mathrm{d}\mu_c(\xi):=\mathrm{d}A\cdot\prod_{\sigma\in\mathcal{M}^2(c)}\exp\left(-\frac{1}{2}F_\sigma(\mathfrak{m})^2\right)\cdot\mathrm{d}\mu^{\times}_c(\vartheta,k)\cdot\mathfrak{D}_c(k,I),
\]
where $\mathrm{d}A$ is Lebesgue measure on $[0,2\pi/q)^{\mathcal{M}^1(c)}$,
the variables $z$, $k$ and $I$ carry the counting measure,
$\mu^{\times}_c$ and the Gauss factor $\mathfrak{D}_c$ are those of Definition~\ref{def:fourier-variables},
and $F_\sigma(\mathfrak{m}):=(dA)(\sigma)+(2\pi/q)\,z(\sigma)$.
The \emph{mirror space} at $c$ is the Hilbert space
\[
\mathcal{H}(c):=L^2\left(\mathcal{Y}_c,\mu_c\right);
\]
it depends on the cross-section $\Xi$.
Its \emph{parity} is the operator $U_{\mathrm{f}}$ given by $(U_{\mathrm{f}}f)(\xi):=(-1)^{\#|I|}\cdot f(\xi)$.
For $v\in\Z^3$ the operator $V_v$ from $\mathcal{H}(c)$ to $\mathcal{H}(c+\ell(v))$ is defined by
\begin{equation}\label{eq:slab-translation-unitary}
(V_vf)\left(\mathfrak{m}',\vartheta',k',r_v(I)\right):=s_v(I)\cdot f\left(\tau_v^*\mathfrak{m}',\tau_v^*\vartheta',\tau_v^*k',I\right)
\end{equation}
for mirror data $\mathfrak{m}'$ and Fourier variables $(\vartheta',k')$ at $c+\ell(v)$ and for $I\subset\mathcal{J}(c)$;
here $r_v$ and $s_v$ are those of Definition~\ref{def:mode-translations},
and $(\tau_v^*u)(\sigma)=u(\sigma+v)$ for every cochain $u$.
\end{definition}

The plaquettes of $\mathcal{M}^2(c)$ have all their links in the plane $\{\ell=c\}$,
so $F_\sigma(\mathfrak{m})$ depends on the mirror data only.
The right side of \eqref{eq:slab-translation-unitary} makes sense.
Indeed, the translation $\tau_v$ maps the cells of $\mathcal{M}(c)$ and of $\mathcal{X}(c)$,
dimension by dimension, onto those of $\mathcal{M}(c+\ell(v))$ and of $\mathcal{X}(c+\ell(v))$,
and it preserves orientations;
so $\tau_v^*\mathfrak{m}'$ are mirror data at $c$, with link variables in $[0,2\pi/q)$,
and $(\tau_v^*\vartheta',\tau_v^*k')$ are Fourier variables at $c$, with angles in $[0,2\pi)$.
Moreover $r_v$ is a bijection of $\mathcal{J}(c)$ onto $\mathcal{J}(c+\ell(v))$,
whence every point of $\mathcal{Y}_{c+\ell(v)}$ is of the form $(\mathfrak{m}',\vartheta',k',r_v(I))$ for exactly one $I\subset\mathcal{J}(c)$.

The measure $\mu_c$ is a finite positive Borel measure.
For fixed $A$, the sum over $z$ of the product of the Gaussian factors of the mirror plaquettes
is the product, over $\sigma\in\mathcal{M}^2(c)$, of the sums over $z(\sigma)\in\Z$,
and each of these is at most $C_{q,1/2}$ by Lemma~\ref{lem:poisson-kernels}\eqref{it:gaussian-sums};
the integral over $A$ contributes the factor $(2\pi/q)^{\#|\mathcal{M}^1(c)|}$.
The total mass of $\mu^{\times}_c$ is at most $\left(\sum_{k\in\Z}\exp(-2\pi^2k^2/q^2)\right)^{\#|\mathcal{X}^2(c)|}$
(Definition~\ref{def:fourier-variables}),
the Gauss factor is at most $(2\pi/q)^{\#|\mathcal{X}^1(c)|}$,
and $\mathcal{J}(c)$ has $2^{\#|\mathcal{J}(c)|}$ subsets.
Therefore
\begin{equation}\label{eq:slab-mirror-mass}
\mu_c(\mathcal{Y}_c)\le(2\pi/q)^{\#|\mathcal{M}^1(c)|+\#|\mathcal{X}^1(c)|}\cdot C_{q,1/2}^{\#|\mathcal{M}^2(c)|}
\cdot\left(\sum_{k\in\Z}\exp(-2\pi^2k^2/q^2)\right)^{\#|\mathcal{X}^2(c)|}\cdot2^{\#|\mathcal{J}(c)|}<\infty .
\end{equation}
Fixing $(z,k,I)$ singles out a box $[0,2\pi/q)^{\mathcal{M}^1(c)}\times[0,2\pi)^{\mathcal{X}^3(c)}$,
and $\mathcal{Y}_c$ is the disjoint union of these countably many boxes.
Hence $\mathcal{H}(c)$, the $L^2$ space of a finite Borel measure on such a union, is separable.

The space $\mathcal{H}(c)$ also admits a description by vector-valued functions.
To a function $f$ on $\mathcal{Y}_c$ corresponds the $\mathcal{F}(c)$-valued function
\[
f(\mathfrak{m},\vartheta,k):=\sum_{I\subset\mathcal{J}(c)}f(\mathfrak{m},\vartheta,k,I)\,e_I
\]
of the bosonic mirror variables $(\mathfrak{m},\vartheta,k)$,
on which $U_{\mathrm{f}}$ acts, value by value, as the parity of the Fock space $\mathcal{F}(c)$ (Definition~\ref{def:modes}).
The Gauss factor takes only the values $0$ and $(2\pi/q)^{\#|\mathcal{X}^1(c)|}$,
and the values of $f$ at the points where it vanishes do not enter the norm of $f$.
Consequently the map that sends $f$ to
$\sum_{I:\ \mathfrak{D}_c(k,I)\neq0}f(\mathfrak{m},\vartheta,k,I)\,e_I$
identifies $\mathcal{H}(c)$, with its inner product multiplied by $(q/2\pi)^{\#|\mathcal{X}^1(c)|}$,
with the closed subspace of the square integrable $\mathcal{F}(c)$-valued functions of $(\mathfrak{m},\vartheta,k)$
for the measure $\mathrm{d}A\cdot\prod_{\sigma\in\mathcal{M}^2(c)}\exp\left(-\frac12F_\sigma(\mathfrak{m})^2\right)\cdot\mathrm{d}\mu^{\times}_c(\vartheta,k)$
that consists of the functions whose value at $(\mathfrak{m},\vartheta,k)$ lies in the span of the $e_I$ with $\mathfrak{D}_c(k,I)\neq0$.

\begin{definition}[Kernels of slab operators]\label{def:slab-kernels}
Let $c<c'$ be positions in $\Adm(\ell)$, and let $G\in\mathcal{I}_\Xi[c,c']$ (Definition~\ref{def:slab-integrands}).
For $\xi=(\mathfrak{m},\vartheta,k,I)\in\mathcal{Y}_c$ and $\xi'=(\mathfrak{m}',\vartheta',k',I')\in\mathcal{Y}_{c'}$
let $\varpi=(A,z):=\mathfrak{m}\vee\mathring{\varpi}\vee\mathfrak{m}'$ be the configuration of the closed slab
with the mirror data $\mathfrak{m}$, $\mathfrak{m}'$ and the interior part $\mathring{\varpi}$,
and put
\begin{multline}\label{eq:slab-operator-kernel}
\mathfrak{S}[G](\xi',\xi):=\left(\omega_{I'}\cdot\omega_I\right)^{\frac12}\int\mathrm{d}\mathring{\varpi}\,
\Bigl\{\mathcal{W}_{c,c'}(\varpi)\cdot\mathfrak{e}_c\left((\vartheta,k);\varpi\right)\cdot\overline{\mathfrak{e}_{c'}\left((\vartheta',k');\theta_{c'}^*\varpi\right)}\\
\cdot\mathfrak{N}_{c,c'}\left[G(\varpi)\exp\left(-S_{(c,c')}(A)\right)\right]_{I',I}\Bigr\}.
\end{multline}
Here the slab weight $\mathcal{W}_{c,c'}$ and the interior integral are those of Definition~\ref{def:slab-configurations},
$\mathfrak{e}_c$ and $\mathfrak{e}_{c'}$ are the phases of Definition~\ref{def:fourier-variables},
$\mathfrak{N}_{c,c'}$ is the slab matrix of Definition~\ref{def:slab-matrices},
and $S_{(c,c')}(A)$ is the action of the open slab (Lemma~\ref{lem:slab-bonds}\eqref{it:bonds-slab-action}).
Where $\mathcal{W}_{c,c'}(\varpi)=0$, that is, where $\varpi\notin\mathfrak{C}[c,c']$, the expression in braces is read as $0$;
elsewhere $G(\varpi)$ is defined.
If $\mathcal{B}^1(c,c')$ and $\mathcal{B}^2(c,c')$ are empty, the interior integral is the evaluation at $\varpi=\mathfrak{m}\vee\mathfrak{m}'$.
The \emph{slab operator} $\mathfrak{S}[G]$ from $\mathcal{H}(c)$ to $\mathcal{H}(c')$ is
\[
(\mathfrak{S}[G]f)(\xi'):=\int\mathrm{d}\mu_c(\xi)\,\mathfrak{S}[G](\xi',\xi)\cdot f(\xi),
\qquad f\in\mathcal{H}(c).
\]
For $G\in\mathcal{A}_\Xi[c,c']$, the operator $\mathfrak{S}[G]$ is the slab operator of the image of $G$ in $\mathcal{I}_\Xi[c,c']$
under the map of Lemma~\ref{lem:product-integrands}\eqref{it:sum-slab}.
\end{definition}

Both phases in \eqref{eq:slab-operator-kernel} read $\varpi$ on cells of the closed slab only.
The phase $\mathfrak{e}_c\left((\vartheta,k);\,\cdot\,\right)$ reads its argument on the upper faces of the cells cut by $c$,
and these lie in $\mathcal{B}(c,c')\cup\mathcal{M}(c')$ by Lemma~\ref{lem:cells-mirrors}\eqref{it:upper-faces}.
The factor $\mathfrak{e}_{c'}\left((\vartheta',k');\theta_{c'}^*\varpi\right)$ reads $\varpi$ on the images under $\theta_{c'}$
of the upper faces of the cells cut by $c'$;
these images are the lower faces of the same cells (Lemma~\ref{lem:cells-mirrors}\eqref{it:cut-faces}),
and they lie in $\mathcal{M}(c)\cup\mathcal{B}(c,c')$ by Lemma~\ref{lem:cells-mirrors}\eqref{it:upper-faces}.
The factor $\left(\omega_{I'}\cdot\omega_I\right)^{\frac12}$ in \eqref{eq:slab-operator-kernel} is the non-negative square root of a non-negative number.
The interior integral converges absolutely and the kernel is bounded on $\mathcal{Y}_{c'}\times\mathcal{Y}_c$
(Theorem~\ref{thm:slab-operators}\eqref{it:slab-hilbert-schmidt} below);
since $\mu_c$ is finite, the integral that defines $\mathfrak{S}[G]f$ therefore converges absolutely.

\begin{theorem}[Composition, reflection and trace]\label{thm:slab-operators}
Let $q>0$, $w_0,\dots,w_3>0$, $M\in\R$ and $-1\le\rho\le1$.
Let $\ell\in\mathcal{L}$, let $\Xi=(\Gamma_\parallel,\chi_\parallel)$ be a cross-section for $\ell$,
and let the spaces $\mathcal{H}(c)$, the operators $U_{\mathrm{f}}$ and $V_v$, and the slab operators be those of
Definitions~\ref{def:mirror-spaces} and~\ref{def:slab-kernels}.
Let $c<c'<c''$ be positions in $\Adm(\ell)$.
\begin{enumerate}
\item\label{it:slab-hilbert-schmidt}
The parity $U_{\mathrm{f}}$ is a unitary involution of $\mathcal{H}(c)$.
For $v,v'\in\Z^3$ the operator $V_v\colon\mathcal{H}(c)\to\mathcal{H}(c+\ell(v))$ is unitary,
$V_vV_{v'}=V_{v+v'}$ as operators from $\mathcal{H}(c)$ to $\mathcal{H}(c+\ell(v+v'))$,
and $V_vU_{\mathrm{f}}=U_{\mathrm{f}}V_v$.
If $v\in\Gamma_\parallel$, then $V_v$ is the identity on the even subspace $\{f:\ U_{\mathrm{f}}f=f\}$ of $\mathcal{H}(c)$
and the multiplication by $\chi_\parallel(v)$ on the odd subspace $\{f:\ U_{\mathrm{f}}f=-f\}$.
The map $G\mapsto\mathfrak{S}[G]$ is linear from $\mathcal{I}_\Xi[c,c']$ to the Hilbert--Schmidt operators from $\mathcal{H}(c)$ to $\mathcal{H}(c')$.
There is a number $C_{c,c'}<\infty$, depending only on $\ell$, $\Xi$, $c$, $c'$, on $q$, $w_0,\dots,w_3$, $M$, $\rho$
and on the base point $v_0$ of the fermion sites,
such that every $G\in\mathcal{I}_\Xi[c,c']$, with constants $C$ and $\varepsilon_G$ as in Definition~\ref{def:slab-integrands}\eqref{it:integrand-growth}, satisfies
\[
\left|\mathfrak{S}[G](\xi',\xi)\right|\le C_{c,c'}\cdot C\cdot C_{q,\varepsilon_G}^{\#|\mathcal{B}^2(c,c')|}
\qquad\text{for all }(\xi',\xi)\in\mathcal{Y}_{c'}\times\mathcal{Y}_c
\]
and
\[
\left\|\mathfrak{S}[G]\right\|_{\mathrm{HS}}\le\sqrt{\mu_c(\mathcal{Y}_c)\cdot\mu_{c'}(\mathcal{Y}_{c'})}\cdot C_{c,c'}\cdot C\cdot C_{q,\varepsilon_G}^{\#|\mathcal{B}^2(c,c')|}.
\]
If $G$ is homogeneous, then $U_{\mathrm{f}}\,\mathfrak{S}[G]=(-1)^{\deg G}\,\mathfrak{S}[G]\,U_{\mathrm{f}}$.
\item\label{it:slab-composition}
For $G\in\mathcal{I}_\Xi[c,c']$ and $G'\in\mathcal{I}_\Xi[c',c'']$, with the product $GG'\in\mathcal{I}_\Xi[c,c'']$ of Definition~\ref{def:slab-integrands},
\[
\mathfrak{S}[G']\,\mathfrak{S}[G]=\mathfrak{S}[GG'].
\]
\item\label{it:slab-covariance}
For $G\in\mathcal{I}_\Xi[c,c']$ and $v\in\Z^3$, with the translate $\tau_vG\in\mathcal{I}_\Xi[c+\ell(v),c'+\ell(v)]$ of Definition~\ref{def:slab-integrands},
\[
V_v\,\mathfrak{S}[G]\,V_v^{-1}=\mathfrak{S}[\tau_vG].
\]
\item\label{it:slab-reflection}
For $G\in\mathcal{I}_\Xi[c,c']$, with the reflection $\Theta_{c'}G\in\mathcal{I}_\Xi[c',2c'-c]$ of Definition~\ref{def:slab-integrands}
and with $v_{c,c'}:=\frac{2(c'-c)}{|n_\ell|^2}\,n_\ell$,
\[
\mathfrak{S}[\Theta_{c'}G]=V_{v_{c,c'}}\,\mathfrak{S}[G]^{*};
\]
both sides map $\mathcal{H}(c')$ to $\mathcal{H}(2c'-c)$.
\item\label{it:slab-trace}
Let $\mathbb{T}=(\Gamma,\chi)$ be a torus with $\Xi(\mathbb{T},\ell)=\Xi$ (Definitions~\ref{def:tori} and~\ref{def:slabs}),
let $g$ be a closing vector of $\mathbb{T}$ for $\ell$, and put $p:=\ell(g)$.
Let $r\ge2$, let $c_0<c_1<\dots<c_r=c_0+p$ be positions in $\Adm(\ell)$,
and let $G_j\in\mathcal{I}_\Xi[c_{j-1},c_j]$ be homogeneous for $j=1,\dots,r$.
For $(A,z)\in\mathfrak{C}(\Gamma)$, regarded as a pair of cochains of the cylinder (they are $\Gamma_\parallel$-periodic, since $\Gamma_\parallel\subset\Gamma$),
let $\varpi_j\in\mathfrak{C}[c_{j-1},c_j]$ be its restriction to the closed slab between the positions $c_{j-1}$ and $c_j$,
and read the product $G_1(\varpi_1)\cdots G_r(\varpi_r)\in\Gr_\Xi((c_0,c_r))$ in $\Gr_{\mathbb{T}}$ as in Lemma~\ref{lem:matrix-rules}\eqref{it:matrix-closure}.
Then $(A,z)\mapsto G_1(\varpi_1)\cdots G_r(\varpi_r)$ is an element $G_1\cdots G_r$ of $\mathcal{I}(\mathbb{T})$ (Definition~\ref{def:torus-integrands}), and
\[
Z_{\mathbb{T}}[G_1\cdots G_r]=\Tr_{\mathcal{H}(c_0)}\left(V_g^{-1}\,U_{\mathrm{f}}^{\frac{1+\chi(g)}{2}}\,\mathfrak{S}[G_r]\cdots\mathfrak{S}[G_1]\right),
\]
where $U_{\mathrm{f}}^{\frac{1+\chi(g)}{2}}$ is the identity if $\chi(g)=-1$ and is $U_{\mathrm{f}}$ if $\chi(g)=+1$.
\end{enumerate}
\end{theorem}

\begin{proof}
\eqref{it:slab-hilbert-schmidt}
\emph{The parity.}
The operator $U_{\mathrm{f}}$ is the multiplication by the function $\xi\mapsto(-1)^{\#|I|}$, whose values are $1$ and $-1$.
Hence $U_{\mathrm{f}}$ preserves norms and $U_{\mathrm{f}}^2=1$; it is a unitary involution.

\emph{The translations.}
Let $v\in\Z^3$.
The map
\[
\left(\mathfrak{m}',\vartheta',k',r_v(I)\right)\longmapsto\left(\tau_v^*\mathfrak{m}',\tau_v^*\vartheta',\tau_v^*k',I\right)
\]
from $\mathcal{Y}_{c+\ell(v)}$ to $\mathcal{Y}_c$ is a bijection:
it renames the cells by the translation $\tau_v$, which preserves orientations, and the sets of modes by the bijection $r_v$;
its inverse sends $(\mathfrak{m},\vartheta,k,I)$ to $(\tau_{-v}^*\mathfrak{m},\tau_{-v}^*\vartheta,\tau_{-v}^*k,r_v(I))$.
It carries $\mu_{c+\ell(v)}$ to $\mu_c$.
Indeed, Lebesgue measure and the counting measures go to Lebesgue measure and the counting measures,
since only coordinates are renamed and no sign occurs.
The pull-back $\tau_v^*$ commutes with the coboundary, so $F_\sigma(\tau_v^*\mathfrak{m}')=F_{\sigma+v}(\mathfrak{m}')$ for $\sigma\in\mathcal{M}^2(c)$,
and the Gaussian factors of the mirror plaquettes at $c+\ell(v)$ go to those at $c$.
Finally, $(\vartheta',k')\mapsto(\tau_v^*\vartheta',\tau_v^*k')$ carries $\mu^{\times}_{c+\ell(v)}$ to $\mu^{\times}_c$,
and $\mathfrak{D}_{c+\ell(v)}(k',r_v(I))=\mathfrak{D}_c(\tau_v^*k',I)$,
by Lemma~\ref{lem:phase-properties}\eqref{it:phase-translations}.
Equivalently, the inverse map, a bijection of $\mathcal{Y}_c$ onto $\mathcal{Y}_{c+\ell(v)}$, carries $\mu_c$ to $\mu_{c+\ell(v)}$.
In particular $V_v$ maps functions that agree almost everywhere to functions that agree almost everywhere.
As $s_v(I)\in\{1,-1\}$, formula \eqref{eq:slab-translation-unitary} shows that $V_v$ preserves norms.
It is onto: a function $h\in\mathcal{H}(c+\ell(v))$ is the image of the function
\[
(\mathfrak{m},\vartheta,k,I)\longmapsto s_v(I)\cdot h\left(\tau_{-v}^*\mathfrak{m},\tau_{-v}^*\vartheta,\tau_{-v}^*k,r_v(I)\right).
\]
So $V_v$ is unitary.

The remaining properties of $V_v$ are read off from the vector-valued functions introduced after Definition~\ref{def:mirror-spaces}.
Since $U_ve_I=s_v(I)\,e_{r_v(I)}$ (Definition~\ref{def:mode-translations}),
formula \eqref{eq:slab-translation-unitary} becomes
\[
(V_vf)(\mathfrak{m}',\vartheta',k')=U_v\,f\left(\tau_v^*\mathfrak{m}',\tau_v^*\vartheta',\tau_v^*k'\right).
\]
For $v,v'\in\Z^3$ one has $U_vU_{v'}=U_{v+v'}$ by Definition~\ref{def:mode-translations}\eqref{it:fock-translations},
and $\tau_{v'}^*\tau_v^*=(\tau_v\tau_{v'})^*=\tau_{v+v'}^*$;
applying the last formula twice gives $V_vV_{v'}f=V_{v+v'}f$.
The unitary $U_v$ commutes with the parity of the Fock spaces (same item),
and $U_{\mathrm{f}}$ acts value by value as this parity;
hence $V_vU_{\mathrm{f}}=U_{\mathrm{f}}V_v$.
Let now $v\in\Gamma_\parallel$.
Then $\ell(v)=0$;
every cochain of the cylinder is $\Gamma_\parallel$-periodic, so that $\tau_v^*$ leaves mirror data and Fourier variables unchanged;
and $U_ve_I=\chi_\parallel(v)^{\#|I|}\,e_I$ (same item).
Therefore $V_vf=f$ if $f$ vanishes at all points with $\#|I|$ odd,
and $V_vf=\chi_\parallel(v)\,f$ if $f$ vanishes at all points with $\#|I|$ even.

\emph{The kernel bound.}
Let $G\in\mathcal{I}_\Xi[c,c']$, with constants $C$ and $\varepsilon_G$ as in Definition~\ref{def:slab-integrands}\eqref{it:integrand-growth},
let $\xi\in\mathcal{Y}_c$ and $\xi'\in\mathcal{Y}_{c'}$,
and let $\varpi=(A,z)\in\mathfrak{C}[c,c']$ have the mirror data of $\xi$ and of $\xi'$.
The two phases in \eqref{eq:slab-operator-kernel} have modulus one (Lemma~\ref{lem:phase-properties}).
By Lemma~\ref{lem:slab-bonds}\eqref{it:bonds-slab-action}, $S_{(c,c')}(A)$ is an even element of $\Gr_\Xi((c,c'))$
whose coefficients are polynomials in the phases $\exp\left(\pm iqA(b)\right)$, $b\in[c,c']^{(1)}$, which have modulus one.
The algebra $\Gr_\Xi((c,c'))$ has finite dimension, and $S_{(c,c')}(A)$ is a linear combination of products of two generators;
so $\exp\left(-S_{(c,c')}(A)\right)$ is a finite sum of powers of $S_{(c,c')}(A)$,
and its coefficients are bounded uniformly in $A$.
The elements $\zeta^{(c)}_I$ and $\Theta_{c'}\left(\zeta^{(c')}_{I'}\right)$ of $\Gr_\Xi((c,c'))$ do not depend on $\varpi$.
By Definition~\ref{def:slab-matrices}, the entry $\mathfrak{N}_{c,c'}\left[G(\varpi)\exp\left(-S_{(c,c')}(A)\right)\right]_{I',I}$
is the Berezin integral over $(c,c')$ of the product
$\zeta^{(c)}_I\cdot G(\varpi)\cdot\exp\left(-S_{(c,c')}(A)\right)\cdot\Theta_{c'}\left(\zeta^{(c')}_{I'}\right)$,
that is, the coefficient of $\Psi^{\mathrm{top}}\left(\mathbb{X}_\Xi((c,c'))\right)$ in this product.
In a monomial basis this coefficient is a sum, with signs, of products of one coefficient of each of the four factors,
and the number of terms is bounded in terms of the number of generators of $\Gr_\Xi((c,c'))$.
Together with Definition~\ref{def:slab-integrands}\eqref{it:integrand-growth} this yields
\[
\left|\mathfrak{N}_{c,c'}\left[G(\varpi)\exp\left(-S_{(c,c')}(A)\right)\right]_{I',I}\right|
\le C'\cdot C\cdot\exp\left(\left(\frac{1}{2}-\varepsilon_G\right)\sum_{\sigma\in\mathcal{B}^2(c,c')}F_\sigma(\varpi)^2\right),
\]
where $C'$ depends only on $\ell$, $\Xi$, $c$, $c'$, $q$, $w_0,\dots,w_3$, $M$, $\rho$ and $v_0$.
The indicators in $\mathcal{W}_{c,c'}(\varpi)$ are at most $1$.
It follows that the expression in braces in \eqref{eq:slab-operator-kernel} has modulus at most
\[
C'\cdot C\cdot\prod_{\sigma\in\mathcal{B}^2(c,c')}\exp\left(-\varepsilon_G\,F_\sigma(\varpi)^2\right)
\]
for every interior part;
where $\mathcal{W}_{c,c'}(\varpi)=0$ the expression is $0$, and the bound holds trivially.
For $\sigma\in\mathcal{B}^2(c,c')$ one has $F_\sigma(\varpi)=(dA)(\sigma)+(2\pi/q)\,z(\sigma)$,
where the links of $\sigma$ lie in the closed slab and $z(\sigma)$ is a variable of summation of the interior integral.
Drop the conditions $dz=0$ on $\mathcal{B}^3(c,c')$ and sum the last product over all $z\in\Z^{\mathcal{B}^2(c,c')}$ at fixed $A$:
the sum is the product over $\sigma\in\mathcal{B}^2(c,c')$ of the sums over $z(\sigma)\in\Z$ of $\exp\left(-\varepsilon_G\,F_\sigma(\varpi)^2\right)$,
and each of these is at most $C_{q,\varepsilon_G}$, whatever the value of $(dA)(\sigma)$, by Lemma~\ref{lem:poisson-kernels}\eqref{it:gaussian-sums}.
The integral over the variables of the links of $\mathcal{B}^1(c,c')$ adds the factor $(2\pi/q)^{\#|\mathcal{B}^1(c,c')|}$.
So the interior integral in \eqref{eq:slab-operator-kernel} converges absolutely, and, for all values of the mirror variables,
\[
\left|\mathfrak{S}[G](\xi',\xi)\right|\le\left(\omega_{I'}\cdot\omega_I\right)^{\frac12}\cdot C'\cdot C\cdot(2\pi/q)^{\#|\mathcal{B}^1(c,c')|}\cdot C_{q,\varepsilon_G}^{\#|\mathcal{B}^2(c,c')|}.
\]
Finally, $\omega_I\le\max\left(1,\max_i\omega_i\right)^{\#|\mathcal{J}(c)|}$ and $\omega_{I'}\le\max\left(1,\max_i\omega_i\right)^{\#|\mathcal{J}(c')|}$,
where $i$ ranges over all modes.
The kernel bound of the theorem follows with
\[
C_{c,c'}:=C'\cdot(2\pi/q)^{\#|\mathcal{B}^1(c,c')|}\cdot\max\left(1,\max_i\omega_i\right)^{\frac{1}{2}\left(\#|\mathcal{J}(c)|+\#|\mathcal{J}(c')|\right)},
\]
a number that does not depend on $G$.

\emph{The Hilbert--Schmidt property.}
The expression in braces in \eqref{eq:slab-operator-kernel} is a Borel function of $(\xi',\xi,\mathring{\varpi})$:
the map $G$ is Borel,
the phases are continuous in the real variables for fixed integer variables (Lemma~\ref{lem:phase-properties}),
and the entry of the slab matrix is a finite sum of products of coefficients of $G(\varpi)$, of $\exp\left(-S_{(c,c')}(A)\right)$ and of the modes.
Its interior integral converges absolutely for every $(\xi',\xi)$, as shown above.
Tonelli's theorem, applied to the positive and negative parts of the real and imaginary parts of this expression,
shows that $\mathfrak{S}[G]$ is a Borel function on $\mathcal{Y}_{c'}\times\mathcal{Y}_c$; it is bounded.
The measures $\mu_c$ and $\mu_{c'}$ are finite by \eqref{eq:slab-mirror-mass}, hence
\[
\int\mathrm{d}\mu_{c'}(\xi')\int\mathrm{d}\mu_c(\xi)\,\left|\mathfrak{S}[G](\xi',\xi)\right|^2
\le\mu_c(\mathcal{Y}_c)\cdot\mu_{c'}(\mathcal{Y}_{c'})\cdot\left(C_{c,c'}\cdot C\cdot C_{q,\varepsilon_G}^{\#|\mathcal{B}^2(c,c')|}\right)^2 .
\]
The spaces $\mathcal{H}(c)$ and $\mathcal{H}(c')$ are separable, as shown after Definition~\ref{def:mirror-spaces}.
By Lemma~\ref{lem:poisson-kernels}\eqref{it:hilbert-schmidt-kernels}, applied to the finite measure spaces $(\mathcal{Y}_c,\mu_c)$ and $(\mathcal{Y}_{c'},\mu_{c'})$,
the kernel $\mathfrak{S}[G]$ defines a Hilbert--Schmidt operator from $\mathcal{H}(c)$ to $\mathcal{H}(c')$
whose Hilbert--Schmidt norm is the $L^2$ norm of the kernel.
This operator is the slab operator $\mathfrak{S}[G]$ of Definition~\ref{def:slab-kernels}:
every $f\in\mathcal{H}(c)$ is integrable for the finite measure $\mu_c$, and the kernel is bounded,
so the integral that defines $(\mathfrak{S}[G]f)(\xi')$ converges absolutely for every $\xi'\in\mathcal{Y}_{c'}$.
The bound on $\left\|\mathfrak{S}[G]\right\|_{\mathrm{HS}}$ is the square root of the last inequality.

\emph{Linearity.}
Let $G,G'\in\mathcal{I}_\Xi[c,c']$ and $a\in\C$; then $aG+G'\in\mathcal{I}_\Xi[c,c']$.
The slab matrix is linear in its argument, so the expression in braces in \eqref{eq:slab-operator-kernel} for $aG+G'$
is $a$ times the expression for $G$ plus the expression for $G'$.
All three interior integrals converge absolutely; hence $\mathfrak{S}[aG+G'](\xi',\xi)=a\,\mathfrak{S}[G](\xi',\xi)+\mathfrak{S}[G'](\xi',\xi)$,
and $G\mapsto\mathfrak{S}[G]$ is linear.

\emph{The parity rule.}
Let $G$ be homogeneous.
The element $\exp\left(-S_{(c,c')}(A)\right)$ is even, so $G(\varpi)\exp\left(-S_{(c,c')}(A)\right)$ is homogeneous of degree $\deg G$.
By Lemma~\ref{lem:matrix-rules}\eqref{it:matrix-parity}, its slab matrix has the entry $0$ at $(I',I)$ unless $\#|I|+\#|I'|+\deg G$ is even.
By \eqref{eq:slab-operator-kernel}, the same holds for $\mathfrak{S}[G](\xi',\xi)$, so that
\[
(-1)^{\#|I'|}\cdot\mathfrak{S}[G](\xi',\xi)=(-1)^{\deg G}\cdot\mathfrak{S}[G](\xi',\xi)\cdot(-1)^{\#|I|}
\]
for all $\xi=(\mathfrak{m},\vartheta,k,I)$ and $\xi'=(\mathfrak{m}',\vartheta',k',I')$.
Multiplying by $f(\xi)$, $f\in\mathcal{H}(c)$, and integrating against $\mu_c$ gives
$U_{\mathrm{f}}\,\mathfrak{S}[G]\,f=(-1)^{\deg G}\,\mathfrak{S}[G]\,U_{\mathrm{f}}\,f$.

\emph{The composition rule~\eqref{it:slab-composition}.}
Let $G\in\mathcal{I}_\Xi[c,c']$ and $G'\in\mathcal{I}_\Xi[c',c'']$;
their product $GG'\in\mathcal{I}_\Xi[c,c'']$ is the one of Definition~\ref{def:slab-integrands}.
For a configuration $\varpi$ of a closed slab, $A(\varpi)$ denotes its link part.
Every sum and every integral in the following computation converges absolutely, jointly in all its variables:
the measures $\mu_c$, $\mu_{c'}$, $\mu_{c''}$ are finite,
the kernels are bounded by item~\eqref{it:slab-hilbert-schmidt},
the interior integral in Definition~\ref{def:slab-kernels} converges absolutely, uniformly in the mirror variables, by the estimate in the proof of item~\eqref{it:slab-hilbert-schmidt},
and both sides of the identity of Lemma~\ref{lem:fourier-crossing} at the position $c'$ converge absolutely,
uniformly in the values given on the upper and the lower faces of the cells cut by $c'$.
Summations and integrations may therefore be interchanged freely, by Fubini's theorem.

By Lemma~\ref{lem:poisson-kernels}\eqref{it:hilbert-schmidt-kernels},
the operator $\mathfrak{S}[G']\mathfrak{S}[G]\colon\mathcal{H}(c)\to\mathcal{H}(c'')$ has the kernel
\[
(\xi'',\xi)\longmapsto\int\mathrm{d}\mu_{c'}(\xi')\,\mathfrak{S}[G'](\xi'',\xi')\cdot\mathfrak{S}[G](\xi',\xi),
\]
and it suffices to show that this function coincides with $\mathfrak{S}[GG']$ at every point of $\mathcal{Y}_{c''}\times\mathcal{Y}_c$.
Write $\xi=(\mathfrak{m},\vartheta,k,I)$, $\xi'=(\mathfrak{m}',\vartheta',k',I')$ and $\xi''=(\mathfrak{m}'',\vartheta'',k'',I'')$.
The two kernels are interior integrals over configurations
\[
\varpi_1=\mathfrak{m}\vee\mathring{\varpi}_1\vee\mathfrak{m}'
\qquad\text{and}\qquad
\varpi_2=\mathfrak{m}'\vee\mathring{\varpi}_2\vee\mathfrak{m}''
\]
of the closed slabs $[c,c']$ and $[c',c'']$.
Once both kernels are written out, the integrand depends on the Fourier variables $(\vartheta',k')$
only through the factor $\mathrm{d}\mu^{\times}_{c'}(\vartheta',k')\,\mathfrak{D}_{c'}(k',I')$ of $\mathrm{d}\mu_{c'}(\xi')$
and through the product
\[
\mathfrak{e}_{c'}((\vartheta',k');\varpi_2)\cdot\overline{\mathfrak{e}_{c'}((\vartheta',k');\theta_{c'}^{*}\varpi_1)}.
\]
Apart from the Gauss factor $\mathfrak{D}_{c'}(k',I')$, it depends on the index set $I'$ only through
\[
\omega_{I'}\cdot\mathfrak{N}_{c',c''}\left[G'(\varpi_2)\cdot\exp\left(-S_{(c',c'')}(A(\varpi_2))\right)\right]_{I'',I'}\cdot\mathfrak{N}_{c,c'}\left[G(\varpi_1)\cdot\exp\left(-S_{(c,c')}(A(\varpi_1))\right)\right]_{I',I},
\]
because the coefficients in front of the two interior integrals multiply to
\[
\sqrt{\omega_{I''}\cdot\omega_{I'}}\cdot\sqrt{\omega_{I'}\cdot\omega_I}=\sqrt{\omega_{I''}\cdot\omega_I}\cdot\omega_{I'}.
\]

The upper faces of the cells cut by $c'$ lie in $\mathcal{B}(c',c'')\cup\mathcal{M}(c'')$, where $\varpi_2$ is given,
and their lower faces lie in $\mathcal{M}(c)\cup\mathcal{B}(c,c')$, where $\varpi_1$ is given
(Lemma~\ref{lem:cells-mirrors}\eqref{it:upper-faces}).
For fixed $I'$, Lemma~\ref{lem:fourier-crossing} at the position $c'$,
with the upper faces read from $\varpi_2$ and the lower faces from $\varpi_1$, therefore yields
\begin{multline*}
\int\mathrm{d}\mu^{\times}_{c'}(\vartheta',k')\,\mathfrak{D}_{c'}(k',I')\cdot\mathfrak{e}_{c'}((\vartheta',k');\varpi_2)\cdot\overline{\mathfrak{e}_{c'}((\vartheta',k');\theta_{c'}^{*}\varpi_1)}\\
=\sum_{z^{\times}}\int\mathrm{d}A^{\times}\prod_{\sigma\in\mathcal{X}^2(c')}\exp\left(-\frac12F_\sigma^2\right)\cdot\prod_{\sigma\in\mathcal{X}^3(c')}\mathbf{1}[(dz)(\sigma)=0]\cdot u_{I'}(A^{\times}),
\end{multline*}
where $z^{\times}$ ranges over $\Z^{\mathcal{X}^2(c')}$ and $A^{\times}$ over $[0,2\pi/q)^{\mathcal{X}^1(c')}$,
and where $F_\sigma$ and $(dz)(\sigma)$ are formed with $(A^{\times},z^{\times})$ on the cut links and plaquettes
and with $\varpi_1$, $\varpi_2$ on the other faces.

Let $\varpi$ be the configuration of the closed slab $[c,c'']$ that coincides with $\varpi_1$ on the cells of $[c,c']$,
with $\varpi_2$ on those of $[c',c'']$, and with $(A^{\times},z^{\times})$ on $\mathcal{X}^1(c')\cup\mathcal{X}^2(c')$.
By the partition of Lemma~\ref{lem:cells-mirrors}\eqref{it:cell-partition},
$\mathcal{B}(c,c'')$ is the disjoint union of $\mathcal{B}(c,c')$, $\mathcal{M}(c')$, $\mathcal{X}(c')$ and $\mathcal{B}(c',c'')$.
Hence $(\mathring{\varpi}_1,\mathfrak{m}',A^{\times},z^{\times},\mathring{\varpi}_2)$ is precisely the interior part $\mathring{\varpi}$ of $\varpi$,
each link variable ranging over $[0,2\pi/q)$ and each integer variable over $\Z$, as in the interior integral of $[c,c'']$.
Consider first the values of these variables for which $\varpi\in\mathfrak{C}[c,c'']$;
these are the values at which $\mathcal{W}_{c,c''}(\varpi)\neq0$.
Every cube of $\mathcal{B}^3(c,c')$ belongs to $\mathcal{B}^3(c,c'')$ by the partition of the cells,
and $dz$ at such a cube is formed with the values of $z$ on plaquettes of the closed slab $[c,c']$, where $\varpi$ agrees with $\varpi_1$;
so $\varpi_1\in\mathfrak{C}[c,c']$, and likewise $\varpi_2\in\mathfrak{C}[c',c'']$.
Thus $G(\varpi_1)$, $G'(\varpi_2)$ and $(GG')(\varpi)$ are defined.
At such values,
the expansion \eqref{eq:slab-crossing-factor} of $\exp(-S^{\times}_{c'})$,
the composition rule of Lemma~\ref{lem:matrix-rules}\eqref{it:matrix-composition},
applied at $c'$ with the diagonal matrix whose entries are $\omega_{I'}\cdot u_{I'}(A^{\times})$,
and the decomposition $S_{(c,c'')}=S_{(c,c')}+S_{(c',c'')}+S^{\times}_{c'}$ of Lemma~\ref{lem:slab-bonds}\eqref{it:bonds-action-split} give
\begin{align*}
&\sum_{I'}\mathfrak{N}_{c',c''}\left[G'(\varpi_2)\cdot\exp\left(-S_{(c',c'')}\right)\right]_{I'',I'}\cdot\omega_{I'}\cdot u_{I'}(A^{\times})\cdot\mathfrak{N}_{c,c'}\left[G(\varpi_1)\cdot\exp\left(-S_{(c,c')}\right)\right]_{I',I}\\
&\qquad=\mathfrak{N}_{c,c''}\left[G(\varpi_1)\cdot\exp\left(-S_{(c,c')}\right)\cdot\exp\left(-S^{\times}_{c'}\right)\cdot G'(\varpi_2)\cdot\exp\left(-S_{(c',c'')}\right)\right]_{I'',I}\\
&\qquad=\mathfrak{N}_{c,c''}\left[(GG')(\varpi)\cdot\exp\left(-S_{(c,c'')}\right)\right]_{I'',I},
\end{align*}
all actions being evaluated at $A(\varpi)$.
This is legitimate because $S_{(c,c')}$ and $S_{(c',c'')}$ read only the links of their closed slabs, where $A(\varpi)$ agrees with $A(\varpi_1)$ and with $A(\varpi_2)$,
and $S^{\times}_{c'}$ reads only the cut links
(Lemma~\ref{lem:slab-bonds}\eqref{it:bonds-slab-action},~\eqref{it:bonds-action-split}).
In the last step the three exponentials, which are even and hence commute with every element of the Grassmann algebra,
are moved to the right and combined;
the order of $G(\varpi_1)$ and $G'(\varpi_2)$ is kept, and $G(\varpi_1)\cdot G'(\varpi_2)=(GG')(\varpi)$.

The same partition, together with the fact that no cube lies in the plane $\{\ell=c'\}$,
shows that the weights combine into the slab weight of $[c,c'']$:
\begin{multline*}
\mathcal{W}_{c,c'}(\varpi_1)\cdot\prod_{\sigma\in\mathcal{M}^2(c')}\exp\left(-\frac12F_\sigma^2\right)\cdot\prod_{\sigma\in\mathcal{X}^2(c')}\exp\left(-\frac12F_\sigma^2\right)\\
\cdot\prod_{\sigma\in\mathcal{X}^3(c')}\mathbf{1}[(dz)(\sigma)=0]\cdot\mathcal{W}_{c',c''}(\varpi_2)=\mathcal{W}_{c,c''}(\varpi),
\end{multline*}
with $F_\sigma=F_\sigma(\varpi)$;
the second factor on the left is the density of $\mu_{c'}$ in the mirror data $\mathfrak{m}'$,
and the third and the fourth come from Lemma~\ref{lem:fourier-crossing}.
At the remaining values of the variables, where $\varpi\notin\mathfrak{C}[c,c'']$, the right side of this identity is $0$;
as the Gaussian factors are positive, $\mathcal{W}_{c,c'}(\varpi_1)=0$, or $\mathcal{W}_{c',c''}(\varpi_2)=0$, or $(dz)(\sigma)\neq0$ for some $\sigma\in\mathcal{X}^3(c')$.
Accordingly, the expression in braces of $\mathfrak{S}[G]$ or of $\mathfrak{S}[G']$ is read as $0$ (Definition~\ref{def:slab-kernels}),
or an indicator at a cube cut by $c'$ vanishes;
in each case the integrand that results from the crossing identity, as a function of $(\mathring{\varpi}_1,\mathfrak{m}',A^{\times},z^{\times},\mathring{\varpi}_2)$, is $0$ at these values,
and the expression in braces of the kernel of $[c,c'']$ is read as $0$ as well.
Moreover
\[
\mathfrak{e}_c((\vartheta,k);\varpi_1)=\mathfrak{e}_c((\vartheta,k);\varpi),\qquad
\mathfrak{e}_{c''}((\vartheta'',k'');\theta_{c''}^{*}\varpi_2)=\mathfrak{e}_{c''}((\vartheta'',k'');\theta_{c''}^{*}\varpi),
\]
since in each equality both sides read the same cells:
the upper faces of the cells cut by $c$, which lie in $\mathcal{B}(c,c')\cup\mathcal{M}(c')$,
and the lower faces of the cells cut by $c''$, which lie in $\mathcal{M}(c')\cup\mathcal{B}(c',c'')$
(Lemma~\ref{lem:cells-mirrors}\eqref{it:upper-faces}).
Collecting all factors,
\begin{multline*}
\int\mathrm{d}\mu_{c'}(\xi')\,\mathfrak{S}[G'](\xi'',\xi')\cdot\mathfrak{S}[G](\xi',\xi)
=\sqrt{\omega_{I''}\cdot\omega_I}\int\mathrm{d}\mathring{\varpi}\,\Bigl\{\mathcal{W}_{c,c''}(\varpi)\cdot\mathfrak{e}_c((\vartheta,k);\varpi)\\
\cdot\overline{\mathfrak{e}_{c''}((\vartheta'',k'');\theta_{c''}^{*}\varpi)}\cdot\mathfrak{N}_{c,c''}\left[(GG')(\varpi)\cdot\exp\left(-S_{(c,c'')}(A(\varpi))\right)\right]_{I'',I}\Bigr\},
\end{multline*}
where $\varpi=\mathfrak{m}\vee\mathring{\varpi}\vee\mathfrak{m}''$ and the interior integral is that of $[c,c'']$.
By Definition~\ref{def:slab-kernels}, the right side is $\mathfrak{S}[GG'](\xi'',\xi)$.

\emph{Translation covariance~\eqref{it:slab-covariance}.}
Now let $G\in\mathcal{I}_\Xi[c,c']$ and $v\in\Z^3$;
then $\tau_vG\in\mathcal{I}_\Xi[c+\ell(v),c'+\ell(v)]$ (Definition~\ref{def:slab-integrands}).
For $\xi_1=(\mathfrak{m},\vartheta,k,I)\in\mathcal{Y}_c$ put
\[
\xi_1^v:=\left(\tau_{-v}^{*}\mathfrak{m},\tau_{-v}^{*}\vartheta,\tau_{-v}^{*}k,r_v(I)\right)\in\mathcal{Y}_{c+\ell(v)},
\]
and form $\xi_2^v\in\mathcal{Y}_{c'+\ell(v)}$ in the same way from $\xi_2=(\mathfrak{m}',\vartheta',k',I')\in\mathcal{Y}_{c'}$.
The kernels are related by
\[
\mathfrak{S}[\tau_vG](\xi_2^v,\xi_1^v)=s_v(I')\cdot s_v(I)\cdot\mathfrak{S}[G](\xi_2,\xi_1).
\]
To see this, write the configuration in the interior integral that defines the left side as $\tau_{-v}^{*}\varpi$,
where $\varpi=\mathfrak{m}\vee\mathring{\varpi}\vee\mathfrak{m}'$ is a configuration of $[c,c']$.
The substitution $\mathring{\varpi}\mapsto\tau_{-v}^{*}\mathring{\varpi}$ is a bijection of the interior variables of $[c,c']$
onto those of $[c+\ell(v),c'+\ell(v)]$;
since translations preserve the orientation of links, it carries the box $[0,2\pi/q)$ of each link to the box of the translated link,
and it preserves the counting measures.
Under this substitution the factors of the integrand behave as follows.
\begin{itemize}
\item The slab weight is translation invariant: $\mathcal{W}_{c+\ell(v),c'+\ell(v)}(\tau_{-v}^{*}\varpi)=\mathcal{W}_{c,c'}(\varpi)$.
\item By Lemma~\ref{lem:phase-properties}\eqref{it:phase-translations},
$\mathfrak{e}_{c+\ell(v)}((\tau_{-v}^{*}\vartheta,\tau_{-v}^{*}k);\tau_{-v}^{*}\varpi)=\mathfrak{e}_c((\vartheta,k);\varpi)$.
\item Definition~\ref{def:mirrors}\eqref{it:mirror-conjugation} gives $\tau_{-v}\theta_{c'+\ell(v)}\tau_v=\theta_{c'}$, hence
\[
\theta_{c'+\ell(v)}^{*}\tau_{-v}^{*}=\left(\tau_{-v}\theta_{c'+\ell(v)}\right)^{*}=\left(\theta_{c'}\tau_{-v}\right)^{*}=\tau_{-v}^{*}\theta_{c'}^{*},
\]
and Lemma~\ref{lem:phase-properties}\eqref{it:phase-translations}, applied once more, gives
\begin{multline*}
\mathfrak{e}_{c'+\ell(v)}\left((\tau_{-v}^{*}\vartheta',\tau_{-v}^{*}k');\theta_{c'+\ell(v)}^{*}\tau_{-v}^{*}\varpi\right)\\
=\mathfrak{e}_{c'+\ell(v)}\left((\tau_{-v}^{*}\vartheta',\tau_{-v}^{*}k');\tau_{-v}^{*}\theta_{c'}^{*}\varpi\right)
=\mathfrak{e}_{c'}((\vartheta',k');\theta_{c'}^{*}\varpi).
\end{multline*}
\item Let $\varpi\in\mathfrak{C}[c,c']$; then $\tau_{-v}^{*}\varpi\in\mathfrak{C}[c+\ell(v),c'+\ell(v)]$, since $\tau_{-v}^{*}$ commutes with $d$,
and by the definition of $\tau_v$ on slab integrands, $(\tau_vG)(\tau_{-v}^{*}\varpi)=\tau_v(G(\varpi))$.
Moreover $S_{(c+\ell(v),c'+\ell(v))}(\tau_{-v}^{*}A(\varpi))=\tau_vS_{(c,c')}(A(\varpi))$,
because the translation by $v$ carries the terms of the action with both sites in $(c,c')$, evaluated at $A(\varpi)$,
to the terms with both sites in $(c+\ell(v),c'+\ell(v))$, evaluated at $\tau_{-v}^{*}A(\varpi)$.
Lemma~\ref{lem:matrix-rules}\eqref{it:matrix-translation} therefore gives
\begin{multline*}
\mathfrak{N}_{c+\ell(v),c'+\ell(v)}\left[(\tau_vG)(\tau_{-v}^{*}\varpi)\cdot\exp\left(-S_{(c+\ell(v),c'+\ell(v))}(\tau_{-v}^{*}A(\varpi))\right)\right]_{r_v(I'),r_v(I)}\\
=s_v(I')\cdot s_v(I)\cdot\mathfrak{N}_{c,c'}\left[G(\varpi)\cdot\exp\left(-S_{(c,c')}(A(\varpi))\right)\right]_{I',I}.
\end{multline*}
If $\varpi\notin\mathfrak{C}[c,c']$, both slab weights in the first item are $0$, and both expressions in braces are read as $0$.
\item The crossing coefficients are invariant under translations (Definition~\ref{def:mode-translations}),
so $\omega_{r_v(I)}=\omega_I$ and $\omega_{r_v(I')}=\omega_{I'}$.
\end{itemize}
Inserting these relations into the interior integral gives the asserted identity of kernels.

By Definition~\ref{def:mirror-spaces}, $(V_vf)(\xi_1^v)=s_v(I)\cdot f(\xi_1)$ for $f\in\mathcal{H}(c)$ and $\xi_1\in\mathcal{Y}_c$,
and in the same way $(V_vf')(\xi_2^v)=s_v(I')\cdot f'(\xi_2)$ for $f'\in\mathcal{H}(c')$ and $\xi_2\in\mathcal{Y}_{c'}$.
Moreover, as shown in the proof of item~\eqref{it:slab-hilbert-schmidt},
$\xi_1\mapsto\xi_1^v$ is a bijection of $\mathcal{Y}_c$ onto $\mathcal{Y}_{c+\ell(v)}$ that carries $\mu_c$ to $\mu_{c+\ell(v)}$.
Consequently, for $f\in\mathcal{H}(c)$ and every $\xi_2\in\mathcal{Y}_{c'}$,
\begin{align*}
\left(\mathfrak{S}[\tau_vG]V_vf\right)(\xi_2^v)
&=\int\mathrm{d}\mu_c(\xi_1)\,\mathfrak{S}[\tau_vG](\xi_2^v,\xi_1^v)\cdot s_v(I)\cdot f(\xi_1)\\
&=s_v(I')\int\mathrm{d}\mu_c(\xi_1)\,\mathfrak{S}[G](\xi_2,\xi_1)\cdot f(\xi_1)
=\left(V_v\mathfrak{S}[G]f\right)(\xi_2^v),
\end{align*}
because $s_v(I)^2=1$.
Every point of $\mathcal{Y}_{c'+\ell(v)}$ is of the form $\xi_2^v$;
thus $\mathfrak{S}[\tau_vG]V_v=V_v\mathfrak{S}[G]$, and since $V_v$ is unitary by item~\eqref{it:slab-hilbert-schmidt},
$V_v\mathfrak{S}[G]V_v^{-1}=\mathfrak{S}[\tau_vG]$.

\emph{The reflection rule~\eqref{it:slab-reflection}.}
Let $G\in\mathcal{I}_\Xi[c,c']$, and take for $c''$ the position $2c'-c$ of $\Adm(\ell)$, which lies above $c'$;
then $c''=c+\ell(v_{c,c'})$, and $\Theta_{c'}G\in\mathcal{I}_\Xi[c',c'']$ (Definition~\ref{def:slab-integrands}).
As $v_{c,c'}$ lies in $\Z n_\ell$, Definition~\ref{def:mode-translations}\eqref{it:fock-translations} gives
$r_{v_{c,c'}}(I)=I+v_{c,c'}$ and $s_{v_{c,c'}}(I)=1$ for every $I\subset\mathcal{J}(c)$.
For $\xi=(\mathfrak{m},\vartheta,k,I)\in\mathcal{Y}_c$ let
\[
\xi^{v_{c,c'}}:=\left(\tau_{-v_{c,c'}}^{*}\mathfrak{m},\tau_{-v_{c,c'}}^{*}\vartheta,\tau_{-v_{c,c'}}^{*}k,I+v_{c,c'}\right)\in\mathcal{Y}_{c''}.
\]
The map $\xi\mapsto\xi^{v_{c,c'}}$ is a bijection of $\mathcal{Y}_c$ onto $\mathcal{Y}_{c''}$,
and the definition of $V_{v_{c,c'}}$ (Definition~\ref{def:mirror-spaces}) reads $(V_{v_{c,c'}}f)(\xi^{v_{c,c'}})=f(\xi)$ for every $f\in\mathcal{H}(c)$.
The assertion follows from the identity of kernels
\begin{equation}\label{eq:slab-reflected-kernel}
\mathfrak{S}[\Theta_{c'}G]\left(\xi^{v_{c,c'}},\xi'\right)=\overline{\mathfrak{S}[G](\xi',\xi)},
\end{equation}
valid for all $\xi\in\mathcal{Y}_c$ and $\xi'\in\mathcal{Y}_{c'}$.
Indeed, let $f\in\mathcal{H}(c')$.
By Lemma~\ref{lem:poisson-kernels}\eqref{it:hilbert-schmidt-kernels}, the adjoint $\mathfrak{S}[G]^{*}\colon\mathcal{H}(c')\to\mathcal{H}(c)$
has the kernel $(\xi,\xi')\mapsto\overline{\mathfrak{S}[G](\xi',\xi)}$, so that \eqref{eq:slab-reflected-kernel} gives
\begin{align*}
\left(V_{v_{c,c'}}\mathfrak{S}[G]^{*}f\right)\left(\xi^{v_{c,c'}}\right)
&=\left(\mathfrak{S}[G]^{*}f\right)(\xi)
=\int\mathrm{d}\mu_{c'}(\xi')\,\overline{\mathfrak{S}[G](\xi',\xi)}\cdot f(\xi')\\
&=\left(\mathfrak{S}[\Theta_{c'}G]f\right)\left(\xi^{v_{c,c'}}\right)
\end{align*}
for every $\xi\in\mathcal{Y}_c$.
Hence $\mathfrak{S}[\Theta_{c'}G]=V_{v_{c,c'}}\mathfrak{S}[G]^{*}$, both sides being operators from $\mathcal{H}(c')$ to $\mathcal{H}(c'')$.

It remains to prove \eqref{eq:slab-reflected-kernel}.
Fix $\xi=(\mathfrak{m},\vartheta,k,I)\in\mathcal{Y}_c$ and $\xi'=(\mathfrak{m}',\vartheta',k',I')\in\mathcal{Y}_{c'}$,
and write $z(\varpi)$ for the plaquette part of a configuration $\varpi$, whose link part is $A(\varpi)$.
By Definition~\ref{def:slab-kernels},
$\mathfrak{S}[G](\xi',\xi)=\sqrt{\omega_{I'}\cdot\omega_I}\int\mathrm{d}\mathring{\varpi}\,\Phi(\mathfrak{m}\vee\mathring{\varpi}\vee\mathfrak{m}')$, where
\[
\Phi(\varpi):=\mathcal{W}_{c,c'}(\varpi)\cdot\mathfrak{e}_c((\vartheta,k);\varpi)\cdot\overline{\mathfrak{e}_{c'}((\vartheta',k');\theta_{c'}^{*}\varpi)}\cdot\mathfrak{N}_{c,c'}\left[G(\varpi)\cdot\exp\left(-S_{(c,c')}(A(\varpi))\right)\right]_{I',I}
\]
for configurations $\varpi$ of $[c,c']$, read as $0$ where $\mathcal{W}_{c,c'}(\varpi)=0$; likewise
\begin{multline*}
\mathfrak{S}[\Theta_{c'}G]\left(\xi^{v_{c,c'}},\xi'\right)
=\sqrt{\omega_{I+v_{c,c'}}\cdot\omega_{I'}}\int\mathrm{d}\mathring{\check\varpi}\,\Bigl\{\mathcal{W}_{c',c''}(\check\varpi)\cdot\mathfrak{e}_{c'}((\vartheta',k');\check\varpi)\\
\cdot\overline{\mathfrak{e}_{c''}\left((\tau_{-v_{c,c'}}^{*}\vartheta,\tau_{-v_{c,c'}}^{*}k);\theta_{c''}^{*}\check\varpi\right)}\\
\cdot\mathfrak{N}_{c',c''}\left[(\Theta_{c'}G)(\check\varpi)\cdot\exp\left(-S_{(c',c'')}(A(\check\varpi))\right)\right]_{I+v_{c,c'},I'}\Bigr\},
\end{multline*}
where $\check\varpi=\mathfrak{m}'\vee\mathring{\check\varpi}\vee\tau_{-v_{c,c'}}^{*}\mathfrak{m}$ ranges over the configurations of the closed slab $[c',c'']$
with the mirror data of $\xi'$ and of $\xi^{v_{c,c'}}$.

To each such $\check\varpi$ assign the configuration $\varpi:=\theta_{c'}^{*}\check\varpi$ of $[c,c']$.
By Lemma~\ref{lem:cells-mirrors}\eqref{it:cut-orientation}, the reflection $\theta_{c'}$ fixes every cell of $\mathcal{M}(c')$ together with its orientation,
and it maps every cell $\sigma$ of $\mathcal{M}(c)$ onto $\sigma+v_{c,c'}$, again preserving the orientation.
Therefore $\varpi=\mathfrak{m}\vee\mathring{\varpi}\vee\mathfrak{m}'$ with $\mathring{\varpi}=\theta_{c'}^{*}\mathring{\check\varpi}$.
Since $\ell\circ\theta_{c'}=2c'-\ell$ (Definition~\ref{def:mirrors}\eqref{it:mirror-height}), $\theta_{c'}$ maps $\mathcal{B}(c,c')$ onto $\mathcal{B}(c',c'')$;
so the map $\mathring{\check\varpi}\mapsto\mathring{\varpi}$, given by $\mathring{\varpi}(\sigma)=\pm\mathring{\check\varpi}(\theta_{c'}\sigma)$ with the orientation sign of $\theta_{c'}$ at $\sigma$,
is a bijection of $\R^{\mathcal{B}^1(c',c'')}\times\Z^{\mathcal{B}^2(c',c'')}$ onto $\R^{\mathcal{B}^1(c,c')}\times\Z^{\mathcal{B}^2(c,c')}$.
It preserves the Lebesgue measure and the counting measure,
and it maps the box $[0,2\pi/q)^{\mathcal{B}^1(c',c'')}$ onto a product $\prod_{b\in\mathcal{B}^1(c,c')}\mathbb{I}_b$
in which each $\mathbb{I}_b$ is $[0,2\pi/q)$ or $(-2\pi/q,0]$.

The factors of the two integrands correspond to each other as follows.
\begin{enumerate}
\item[(i)] Since $A(\varpi)=\theta_{c'}^{*}A(\check\varpi)$, $z(\varpi)=\theta_{c'}^{*}z(\check\varpi)$ and pull-backs commute with $d$,
one has $F(\varpi)=\theta_{c'}^{*}F(\check\varpi)$ and $dz(\varpi)=\theta_{c'}^{*}dz(\check\varpi)$.
The orientation signs, whose squares are $1$, disappear from $F_\sigma^2$ and from the conditions $(dz)(\sigma)=0$;
hence $\mathcal{W}_{c',c''}(\check\varpi)=\mathcal{W}_{c,c'}(\varpi)$,
and in particular $\check\varpi\in\mathfrak{C}[c',c'']$ if and only if $\varpi\in\mathfrak{C}[c,c']$.
\item[(ii)] Since $\theta_{c'}^{*}\varpi=\check\varpi$ (Definition~\ref{def:mirrors}\eqref{it:mirror-involution}),
the factor $\mathfrak{e}_{c'}((\vartheta',k');\check\varpi)$ is the complex conjugate of the factor
$\overline{\mathfrak{e}_{c'}((\vartheta',k');\theta_{c'}^{*}\varpi)}$ of $\Phi(\varpi)$.
\item[(iii)] Lemma~\ref{lem:phase-properties}\eqref{it:phase-translations}, applied with the vector $v_{c,c'}$, gives
\[
\mathfrak{e}_{c''}\left((\tau_{-v_{c,c'}}^{*}\vartheta,\tau_{-v_{c,c'}}^{*}k);\theta_{c''}^{*}\check\varpi\right)=\mathfrak{e}_c\left((\vartheta,k);\tau_{v_{c,c'}}^{*}\theta_{c''}^{*}\check\varpi\right),
\]
and $\tau_{v_{c,c'}}^{*}\theta_{c''}^{*}=(\theta_{c''}\tau_{v_{c,c'}})^{*}=\theta_{c'}^{*}$,
because $\tau_{v_{c,c'}}=\theta_{c''}\theta_{c'}$ by Definition~\ref{def:mirrors}\eqref{it:mirror-product} and $\theta_{c''}^2=\mathrm{id}$.
So the factor that contains $\theta_{c''}^{*}$ in the formula for $\mathfrak{S}[\Theta_{c'}G](\xi^{v_{c,c'}},\xi')$ equals $\overline{\mathfrak{e}_c((\vartheta,k);\varpi)}$,
the complex conjugate of the factor $\mathfrak{e}_c((\vartheta,k);\varpi)$ of $\Phi(\varpi)$.
\item[(iv)] By Definition~\ref{def:slab-integrands}, $(\Theta_{c'}G)(\check\varpi)=\Theta_{c'}(G(\varpi))$.
The reflection $\theta_{c'}$ maps the sites of the open slab $(c,c')$ onto those of $(c',c'')$.
Applied term by term, Proposition~\ref{prop:symmetries}\eqref{it:action-invariance} carries the term of the oriented bond $(x,x')$, evaluated at $A(\varpi)$,
to the term of $(\theta_{c'}x',\theta_{c'}x)$, evaluated at $A(\check\varpi)$,
and \eqref{eq:setting-reflection-bilinear} carries $M\bar\psi(x)\psi(x)$ to $M\bar\psi(\theta_{c'}x)\psi(\theta_{c'}x)$; therefore
\[
\Theta_{c'}\left(S_{(c,c')}(A(\varpi))\right)=S_{(c',c'')}(A(\check\varpi)).
\]
The antilinear anti-automorphism $\Theta_{c'}$ maps $\exp(-S_{(c,c')}(A(\varpi)))$ to $\exp(-S_{(c',c'')}(A(\check\varpi)))$,
an even element, which commutes with $\Theta_{c'}(G(\varpi))$; hence
\[
(\Theta_{c'}G)(\check\varpi)\cdot\exp\left(-S_{(c',c'')}(A(\check\varpi))\right)=\Theta_{c'}\left(G(\varpi)\cdot\exp\left(-S_{(c,c')}(A(\varpi))\right)\right),
\]
and Lemma~\ref{lem:matrix-rules}\eqref{it:matrix-adjoint} gives
\begin{multline*}
\mathfrak{N}_{c',c''}\left[(\Theta_{c'}G)(\check\varpi)\cdot\exp\left(-S_{(c',c'')}(A(\check\varpi))\right)\right]_{I+v_{c,c'},I'}\\
=\overline{\mathfrak{N}_{c,c'}\left[G(\varpi)\cdot\exp\left(-S_{(c,c')}(A(\varpi))\right)\right]_{I',I}}.
\end{multline*}
\item[(v)] The crossing coefficients are invariant under translations (Definition~\ref{def:mode-translations}); thus $\omega_{I+v_{c,c'}}=\omega_I$.
\end{enumerate}
By (i)--(iv), and since the slab weight is real, the expression in braces in the formula for $\mathfrak{S}[\Theta_{c'}G](\xi^{v_{c,c'}},\xi')$,
taken at $\check\varpi$, equals $\overline{\Phi(\varpi)}$.
Together with (v) and the substitution $\mathring{\check\varpi}\mapsto\mathring{\varpi}$ this gives
\[
\mathfrak{S}[\Theta_{c'}G]\left(\xi^{v_{c,c'}},\xi'\right)
=\sqrt{\omega_I\cdot\omega_{I'}}\cdot\overline{\int_{\prod_b\mathbb{I}_b}\mathrm{d}\mathring{\varpi}\,\Phi(\mathfrak{m}\vee\mathring{\varpi}\vee\mathfrak{m}')},
\]
where $\int_{\prod_b\mathbb{I}_b}\mathrm{d}\mathring{\varpi}$ is the interior integral of $[c,c']$
in which the variable of each link $b\in\mathcal{B}^1(c,c')$ ranges over $\mathbb{I}_b$ instead of $[0,2\pi/q)$.
The function $\Phi$ does not change when $\varpi$ is replaced by its integer shift by an integer one-cochain that vanishes outside $\mathcal{B}^1(c,c')$:
the slab weight depends on $F$ and $dz$ only, and these are invariant under integer shifts;
the two phase factors are invariant by Lemma~\ref{lem:phase-properties}\eqref{it:phase-shifts},
because a link of $\mathcal{B}^1(c,c')$ is cut neither by $c$ nor by $c'$ and is contained neither in $\{\ell\le c\}$ nor in $\{\ell\ge c'\}$;
$G$ is invariant by Definition~\ref{def:slab-integrands}\eqref{it:integrand-shifts};
and $S_{(c,c')}$ depends on the link variables only through their phases (Lemma~\ref{lem:slab-bonds}\eqref{it:bonds-slab-action}), which are invariant as well.
The interior integral of $|\Phi|$ is finite by the estimate in the proof of item~\eqref{it:slab-hilbert-schmidt}.
Lemma~\ref{lem:fundamental-intervals} therefore allows each $\mathbb{I}_b$ to be replaced by $[0,2\pi/q)$, and the right side becomes
\[
\sqrt{\omega_I\cdot\omega_{I'}}\cdot\overline{\int\mathrm{d}\mathring{\varpi}\,\Phi(\mathfrak{m}\vee\mathring{\varpi}\vee\mathfrak{m}')}=\overline{\mathfrak{S}[G](\xi',\xi)},
\]
which is \eqref{eq:slab-reflected-kernel}.

\emph{Trace formula~\eqref{it:slab-trace}.}
Let $\mathbb{T}=(\Gamma,\chi)$, the closing vector $g$, the height period $p=\ell(g)$, the positions $c_0<c_1<\dots<c_r=c_0+p$ with $r\ge2$
and the homogeneous elements $G_j\in\mathcal{I}_\Xi[c_{j-1},c_j]$, $1\le j\le r$, be as in~\eqref{it:slab-trace};
for $(A,z)\in\mathfrak{C}(\Gamma)$ let $\varpi_j$ be, as there, the restriction of $(A,z)$ to the closed slab between $c_{j-1}$ and $c_j$.

\emph{Step 0: the left side.}
Since $dz=0$, the restriction $\varpi_j$ lies in $\mathfrak{C}[c_{j-1},c_j]$.
Continue the chain to a sequence $(c_j)_{j\in\Z}$ by the rule
\[
c_{j+r}:=c_j+p\qquad(j\in\Z).
\]
As $p=\ell(g)$ is an integer and $\Adm(\ell)+\Z=\Adm(\ell)$, all $c_j$ lie in $\Adm(\ell)$;
the sequence is strictly increasing and unbounded above and below, so that Lemma~\ref{lem:cells-mirrors}\eqref{it:cell-partition} applies to it.
The translation $\tau_g$ raises $\ell$ by $p$;
hence it maps $\mathcal{M}(c_j)$, $\mathcal{X}(c_j)$ and $\mathcal{B}(c_{j-1},c_j)$ onto $\mathcal{M}(c_{j+r})$, $\mathcal{X}(c_{j+r})$ and $\mathcal{B}(c_{j+r-1},c_{j+r})$, for every $j\in\Z$.
On the other hand $\Gamma=\Gamma_\parallel+\Z g$ (Definition~\ref{def:slabs}),
so a cell of $\mathbb{T}$ is the orbit of a cell of the cylinder $\Z^3/\Gamma_\parallel$ under the group generated by $\tau_g$.
Every cell of the cylinder lies in exactly one of the sets $\mathcal{M}(c_j)$, $\mathcal{X}(c_j)$, $\mathcal{B}(c_{j-1},c_j)$, $j\in\Z$,
and $\tau_g$ raises the index of this set by $r$.
Every integer is congruent modulo $r$ to exactly one of the numbers $0,\dots,r-1$, and to exactly one of the numbers $1,\dots,r$.
Consequently every cell of $\mathbb{T}$ is the orbit of exactly one cell of the cylinder that belongs to the set
\begin{equation}\label{eq:slab-cell-representatives}
\bigcup_{j=0}^{r-1}\left(\mathcal{M}(c_j)\cup\mathcal{X}(c_j)\right)\ \cup\ \bigcup_{j=1}^{r}\mathcal{B}(c_{j-1},c_j);
\end{equation}
in other words, the set~\eqref{eq:slab-cell-representatives} is a set of representatives of the cells of $\mathbb{T}$.

In particular the sets $\mathcal{B}^2(c_{j-1},c_j)$, $1\le j\le r$, are mapped injectively onto pairwise disjoint sets of plaquettes of $\mathbb{T}$,
and for $\sigma\in\mathcal{B}^2(c_{j-1},c_j)$ the number $F_\sigma(\varpi_j)$ is the value of $F=dA+(2\pi/q)\,z$ at the plaquette of $\mathbb{T}$ determined by $\sigma$.
The values $G_j(\varpi_j)\in\Gr_\Xi((c_{j-1},c_j))$ are read in $\Gr_{\mathbb{T}}$ through the reading of Lemma~\ref{lem:matrix-rules}\eqref{it:matrix-closure},
which maps $\Gr_\Xi((c_0,c_r))$ isomorphically onto $\Gr_{\mathbb{T}}$.
Indeed, since $\Gamma=\Gamma_\parallel+\Z g$ and $\ell(g)=p>0$, the heights of the sites of a class $x+\Gamma$ form the set $\ell(x)+p\Z$,
and the sites of the class at any one of these heights form a single class modulo $\Gamma_\parallel$.
The points of $c_0+p\Z$ are positions of the extended chain, and mirrors carry no sites (Lemma~\ref{lem:slab-bonds}\eqref{it:bonds-crossing});
so exactly one of these heights lies in $(c_0,c_r)$, and every class of $\mathbb{X}_{\mathbb{T}}$ contains exactly one class of $\mathbb{X}_\Xi((c_0,c_r))$.
A set of representatives of $\mathbb{X}_\Xi((c_0,c_r))$ is therefore one of $\mathbb{X}_{\mathbb{T}}$,
and the reading sends the generators at its points, which are free generators of $\Gr_\Xi((c_0,c_r))$ (Definition~\ref{def:slab-grassmann}),
to free generators of $\Gr_{\mathbb{T}}$ (Definition~\ref{def:tori}).
Let $C_j$ and $\varepsilon_{G_j}$ be constants for $G_j$ as in Definition~\ref{def:slab-integrands}\eqref{it:integrand-growth}.
In a monomial basis formed with representatives in the open slabs of the chain,
each coefficient of $G_1(\varpi_1)\cdots G_r(\varpi_r)$ is $0$ or $\pm$ a product of coefficients of the factors.
Since $\tfrac12-\varepsilon_{G_j}\ge0$ for every $j$, its modulus is therefore at most
\[
C_1\cdots C_r\cdot\exp\left(\sum_{j=1}^{r}\left(\tfrac12-\varepsilon_{G_j}\right)\sum_{\sigma\in\mathcal{B}^2(c_{j-1},c_j)}F_\sigma(\varpi_j)^2\right)
\le C_1\cdots C_r\cdot\exp\left(\left(\tfrac12-\min_j\varepsilon_{G_j}\right)\|F\|^2\right).
\]
The map $(A,z)\mapsto G_1(\varpi_1)\cdots G_r(\varpi_r)$ is Borel.
It is invariant under the integer shifts of $\mathfrak{C}(\Gamma)$:
the restriction of an integer shift of $(A,z)$ to a closed slab of the chain is an integer shift of the restriction,
under which $G_j$ is invariant by Definition~\ref{def:slab-integrands}\eqref{it:integrand-shifts}.
So this map belongs to $\mathcal{I}(\mathbb{T})$ (Definition~\ref{def:torus-integrands}), with the constant $\min_j\varepsilon_{G_j}\in(0,\tfrac12]$;
it is the element $G_1\cdots G_r$ to which $Z_{\mathbb{T}}$ is applied in~\eqref{it:slab-trace}.

\emph{Admissible products.}
Suppose now that, for $1\le j\le r$, a product $G_j\in\mathcal{A}_\Xi[c_{j-1},c_j]$ with the boxes $B_{j,i}$ is given (Definition~\ref{def:slabs}).
Let $\gamma\in\Gamma$.
Since $B_{j,i}\subset\{c_{j-1}<\ell<c_j\}$, the set $B_{j,i}+\gamma$ lies in the open slab $\{c_{j-1}+\ell(\gamma)<\ell<c_j+\ell(\gamma)\}$,
and $\ell(\gamma)\in\ell(\Gamma)=p\Z$.
These open slabs lie between consecutive terms of the extended chain;
two of them are disjoint unless they belong to the same $j\in\{1,\dots,r\}$ and to the same value of $\ell(\gamma)$.
If $B_{j,i}+\gamma$ and $B_{j,i'}+\gamma'$, with $(i,\gamma)\neq(i',\gamma')$, lie in the same open slab,
then $\gamma'-\gamma\in\Gamma\cap\{\ell=0\}=\Gamma_\parallel$,
and the two sets are disjoint because the boxes of $G_j$ embed disjointly in $\Gamma_\parallel$ (Definition~\ref{def:integrands}).
Thus the sets $B_{j,i}+\gamma$, for all $j$, $i$ and $\gamma\in\Gamma$, are pairwise disjoint,
and the product $G_1\cdots G_r$ is read on $\mathbb{T}$ (Definition~\ref{def:torus-integrands}).
A lift of $(A,z)$ on $B_{j,i}$ as in Lemma~\ref{lem:product-integrands}\eqref{it:product-torus} is a lift of $\varpi_j$ on $B_{j,i}$ as in Lemma~\ref{lem:product-integrands}\eqref{it:product-slab};
and the bounded factor of $G_j$ depends on plaquettes of the open slab between $c_{j-1}$ and $c_j$, at which $\varpi_j$ and $(A,z)$ have the same curvature.
Hence the function of $G_1\cdots G_r$ furnished by Lemma~\ref{lem:product-integrands}\eqref{it:product-torus} is $(A,z)\mapsto G_1(\varpi_1)\cdots G_r(\varpi_r)$,
where $G_j(\cdot)$ is the function of $G_j$ furnished by Lemma~\ref{lem:product-integrands}\eqref{it:product-slab}.
Let finally $G_j\in\mathcal{A}_\Xi[c_{j-1},c_j]$ be arbitrary, that is, finite sums of such products.
Then $G_1\cdots G_r$ is a finite sum of products that are read on $\mathbb{T}$.
For products, the functions of Lemma~\ref{lem:product-integrands}\eqref{it:sum-torus} and~\eqref{it:sum-slab} are those of parts~\eqref{it:product-torus} and~\eqref{it:product-slab} of that lemma;
both sides of the following identity are multilinear in $(G_1,\dots,G_r)$, and they coincide when every $G_j$ is a product.
Therefore
\begin{equation}\label{eq:slab-torus-reading}
G_1\cdots G_r\in\mathcal{A}(\mathbb{T}),\qquad
(G_1\cdots G_r)(A,z)=G_1(\varpi_1)\cdots G_r(\varpi_r)\quad\left((A,z)\in\mathfrak{C}(\Gamma)\right),
\end{equation}
where the left side of the identity is the function of Lemma~\ref{lem:product-integrands}\eqref{it:sum-torus}
and $G_j(\cdot)$ is the function of Lemma~\ref{lem:product-integrands}\eqref{it:sum-slab}.

\emph{Step 1: reduction to two slabs.}
Let $r\ge3$.
Applying~\eqref{it:slab-composition} $r-2$ times one obtains
\[
\mathfrak{S}[G_{r-1}]\cdots\mathfrak{S}[G_1]=\mathfrak{S}[G_1\cdots G_{r-1}],
\]
where $G_1\cdots G_{r-1}\in\mathcal{I}_\Xi[c_0,c_{r-1}]$ is the product of Definition~\ref{def:slab-integrands}.
This element is homogeneous,
and its value at the restriction of $(A,z)$ to the closed slab between $c_0$ and $c_{r-1}$ is $G_1(\varpi_1)\cdots G_{r-1}(\varpi_{r-1})$.
Neither side of the trace formula changes, therefore, when the positions are replaced by $c_0<c_{r-1}<c_r$ and $(G_1,\dots,G_r)$ by $(G_1\cdots G_{r-1},G_r)$.
It suffices to treat the case $r=2$, which is assumed from now on.

\emph{Step 2: the trace as an integral.}
The unitary $V_g$ maps $\mathcal{H}(c_0)$ onto $\mathcal{H}(c_0+\ell(g))=\mathcal{H}(c_2)$, and $\mathfrak{S}[G_2]$ maps $\mathcal{H}(c_1)$ to $\mathcal{H}(c_2)$.
So the operator $V_g^{-1}\,U_{\mathrm{f}}^{\frac{1+\chi(g)}{2}}\,\mathfrak{S}[G_2]$, with the parity $U_{\mathrm{f}}$ of $\mathcal{H}(c_2)$, maps $\mathcal{H}(c_1)$ to $\mathcal{H}(c_0)$,
and it is a Hilbert--Schmidt operator by~\eqref{it:slab-hilbert-schmidt}.
Its kernel is found as follows.
By~\eqref{it:slab-hilbert-schmidt} and Definition~\ref{def:mode-translations}, $V_g^{-1}=V_{-g}$ and $r_{-g}=r_g^{-1}$.
Let $(U_v)_{J,I}$ denote the entries of $U_v$ in the bases $(e_I)$, so that $U_ve_I=\sum_J(U_v)_{J,I}\,e_J$.
Since $U_g$ is unitary with real entries (Definition~\ref{def:mode-translations}\eqref{it:fock-translations}), $U_{-g}=U_g^{-1}$ is the transpose of $U_g$, and
\[
s_{-g}\left(r_g(I_0)\right)=(U_{-g})_{I_0,\,r_g(I_0)}=(U_g)_{r_g(I_0),\,I_0}=s_g(I_0)\qquad\left(I_0\subset\mathcal{J}(c_0)\right).
\]
For $\xi_0=(\mathfrak{m}_0,\vartheta_0,k_0,I_0)\in\mathcal{Y}_{c_0}$ put
\[
\xi_0^g:=\left(\tau_{-g}^*\mathfrak{m}_0,\ \tau_{-g}^*\vartheta_0,\ \tau_{-g}^*k_0,\ r_g(I_0)\right)\in\mathcal{Y}_{c_2}.
\]
The definitions of $V_{-g}$ and of $U_{\mathrm{f}}$ (Definition~\ref{def:mirror-spaces}), together with $\#|r_g(I_0)|=\#|I_0|$, give for $\xi_1\in\mathcal{Y}_{c_1}$
\[
\left(V_g^{-1}\,U_{\mathrm{f}}^{\frac{1+\chi(g)}{2}}\,\mathfrak{S}[G_2]\right)(\xi_0,\xi_1)=\left(-\chi(g)\right)^{\#|I_0|}\cdot s_g(I_0)\cdot\mathfrak{S}[G_2](\xi_0^g,\xi_1);
\]
here $\left(-\chi(g)\right)^{\#|I_0|}$ equals $1$ if $\chi(g)=-1$ and $(-1)^{\#|I_0|}$ if $\chi(g)=+1$.
By Lemma~\ref{lem:poisson-kernels}\eqref{it:hilbert-schmidt-kernels}, the operator $V_g^{-1}\,U_{\mathrm{f}}^{\frac{1+\chi(g)}{2}}\,\mathfrak{S}[G_2]\,\mathfrak{S}[G_1]$ on $\mathcal{H}(c_0)$ is of trace class, and
\[
\begin{aligned}
&\Tr_{\mathcal{H}(c_0)}\left(V_g^{-1}\,U_{\mathrm{f}}^{\frac{1+\chi(g)}{2}}\,\mathfrak{S}[G_2]\,\mathfrak{S}[G_1]\right)\\
&\qquad=\int\mathrm{d}\mu_{c_0}(\xi_0)\int\mathrm{d}\mu_{c_1}(\xi_1)\ \left(-\chi(g)\right)^{\#|I_0|}\cdot s_g(I_0)\cdot\mathfrak{S}[G_2](\xi_0^g,\xi_1)\cdot\mathfrak{S}[G_1](\xi_1,\xi_0),
\end{aligned}
\]
the integral being absolutely convergent.
It stays absolutely convergent when the two kernels are written out as interior integrals by Definition~\ref{def:slab-kernels}.
Indeed, by Definition~\ref{def:slab-integrands}\eqref{it:integrand-growth}, Lemma~\ref{lem:phase-properties} and Lemma~\ref{lem:slab-bonds}\eqref{it:bonds-slab-action},
the integrand of the interior integral that defines $\mathfrak{S}[G](\xi',\xi)$, $G\in\mathcal{I}_\Xi[c,c']$, is bounded in modulus, uniformly in the mirror variables,
by a constant times $\prod_{\sigma\in\mathcal{B}^2(c,c')}\exp\left(-\varepsilon_GF_\sigma^2\right)$;
the interior integral of this majorant is bounded uniformly by Lemma~\ref{lem:poisson-kernels}\eqref{it:gaussian-sums};
and $\mu_{c_0}$, $\mu_{c_1}$ are finite measures by~\eqref{eq:slab-mirror-mass}.
By Fubini's theorem, sums and integrals may be interchanged in all that follows.

\emph{Step 3: the photon at the two mirrors.}
Let $\xi_1=(\mathfrak{m}_1,\vartheta_1,k_1,I_1)$.
After the kernels have been written out, the configurations of the two slabs are
\[
\varpi_1=\mathfrak{m}_0\vee\mathring{\varpi}_1\vee\mathfrak{m}_1,\qquad
\varpi_2=\mathfrak{m}_1\vee\mathring{\varpi}_2\vee\tau_{-g}^*\mathfrak{m}_0,
\]
the interior parts $\mathring{\varpi}_1$ and $\mathring{\varpi}_2$ being the variables of the two interior integrals.
The integrand depends on the Fourier variables $(\vartheta_1,k_1)$ only through the factor $\mathrm{d}\mu^{\times}_{c_1}(\vartheta_1,k_1)\,\mathfrak{D}_{c_1}(k_1,I_1)$ of $\mathrm{d}\mu_{c_1}(\xi_1)$ and through the product
\[
\mathfrak{e}_{c_1}\left((\vartheta_1,k_1);\varpi_2\right)\cdot\overline{\mathfrak{e}_{c_1}\left((\vartheta_1,k_1);\theta_{c_1}^*\varpi_1\right)},
\]
and on $(\vartheta_0,k_0)$ only through the factor $\mathrm{d}\mu^{\times}_{c_0}(\vartheta_0,k_0)\,\mathfrak{D}_{c_0}(k_0,I_0)$ of $\mathrm{d}\mu_{c_0}(\xi_0)$ and through the product
\[
\mathfrak{e}_{c_0}\left((\vartheta_0,k_0);\varpi_1\right)\cdot\overline{\mathfrak{e}_{c_2}\left(\tau_{-g}^*(\vartheta_0,k_0);\theta_{c_2}^*\varpi_2\right)}.
\]
Lemma~\ref{lem:phase-properties}\eqref{it:phase-translations}, applied with $v=g$, gives
\[
\mathfrak{e}_{c_2}\left(\tau_{-g}^*(\vartheta_0,k_0);\theta_{c_2}^*\varpi_2\right)=\mathfrak{e}_{c_0}\left((\vartheta_0,k_0);\tau_g^*\theta_{c_2}^*\varpi_2\right),
\]
and $\tau_g^*\theta_{c_2}^*=(\theta_{c_2}\tau_g)^*=(\tau_g\theta_{c_0})^*=\theta_{c_0}^*\tau_g^*$,
because $\theta_{c_2}=\tau_g\theta_{c_0}\tau_{-g}$ by Definition~\ref{def:mirrors}\eqref{it:mirror-conjugation}.
The configuration $\tau_g^*\varpi_2$ is $\varpi_2$ moved down by one period:
it is a configuration of the closed slab between $c_1-p$ and $c_0$, and its mirror data at $c_0$ are $\mathfrak{m}_0$.
Both products are therefore of the form that appears on the right side of Lemma~\ref{lem:fourier-crossing}:
at $c_1$, with the upper faces of the cells cut by $c_1$ read from $\varpi_2$ and the lower faces read from $\varpi_1$;
at $c_0$, with the upper faces of the cells cut by $c_0$ read from $\varpi_1$ and the lower faces read from $\tau_g^*\varpi_2$.
By Lemma~\ref{lem:cells-mirrors}\eqref{it:upper-faces}, these faces lie in the closed slabs on which the respective configurations are given.
Fix $I_0$ and $I_1$, and read the identity of Lemma~\ref{lem:fourier-crossing} at $c_1$ and at $c_0$ from its right side to its left side.
The Fourier variables $(\vartheta_1,k_1)$ and $(\vartheta_0,k_0)$ then disappear from the integrand.
In their place come an integer variable $z(\sigma)$ for every plaquette $\sigma\in\mathcal{X}^2(c_0)\cup\mathcal{X}^2(c_1)$, summed over $\Z$,
and a real variable $A(b)$ for every link $b\in\mathcal{X}^1(c_0)\cup\mathcal{X}^1(c_1)$, integrated over $[0,2\pi/q)$;
the weight attached to these new variables is
\[
\prod_{\sigma\in\mathcal{X}^2(c_0)\cup\mathcal{X}^2(c_1)}\exp\left(-\tfrac12F_\sigma^2\right)\cdot\prod_{\sigma\in\mathcal{X}^3(c_0)\cup\mathcal{X}^3(c_1)}\mathbf{1}[(dz)(\sigma)=0]\cdot u_{I_1}(A)\cdot u_{I_0}(A).
\]
Here $F=dA+(2\pi/q)\,z$ and $dz$ are formed with the new variables on the cut cells and with the configurations given on their faces;
all variables together are the values of one pair $(A,z)\in C^1(\Gamma)\times C^2(\Gamma;\Z)$, as Step~5 shows.

\emph{Step 4: the fermions.}
For $j=1,2$ put
\[
\mathfrak{N}_j:=\mathfrak{N}_{c_{j-1},c_j}\left[G_j(\varpi_j)\cdot\exp\left(-S_{(c_{j-1},c_j)}(A)\right)\right],
\]
and for $j=0,1$ let $\Upsilon_j:=\omega_{c_j}\diag\left(u_I(A)\right)$ be the diagonal operator on $\mathcal{F}(c_j)$ with the entries $\omega_I\cdot u_I(A)$, $I\subset\mathcal{J}(c_j)$.
The kernels $\mathfrak{S}[G_1](\xi_1,\xi_0)$ and $\mathfrak{S}[G_2](\xi_0^g,\xi_1)$ carry the factors $\sqrt{\omega_{I_1}\cdot\omega_{I_0}}$ and $\sqrt{\omega_{r_g(I_0)}\cdot\omega_{I_1}}$,
and $\omega_{r_g(I_0)}=\omega_{I_0}$ by Definition~\ref{def:mode-translations}\eqref{it:fock-translations}.
For fixed values of all photon variables the terms that depend on $I_0$ and $I_1$ therefore add up to
\[
\begin{aligned}
&\sum_{I_0\subset\mathcal{J}(c_0)}\ \sum_{I_1\subset\mathcal{J}(c_1)}\left(-\chi(g)\right)^{\#|I_0|}\cdot s_g(I_0)\cdot(\mathfrak{N}_2)_{r_g(I_0),\,I_1}\cdot\omega_{I_1}\cdot u_{I_1}(A)\cdot(\mathfrak{N}_1)_{I_1,\,I_0}\cdot\omega_{I_0}\cdot u_{I_0}(A)\\
&\qquad=\operatorname{tr}_{\mathcal{F}(c_0)}\left(U_g^{-1}\,U_{\mathrm{f}}^{\frac{1+\chi(g)}{2}}\,\mathfrak{N}_2\,\Upsilon_1\,\mathfrak{N}_1\,\Upsilon_0\right);
\end{aligned}
\]
for the equality note that $U_g^{-1}e_{r_g(I_0)}=s_g(I_0)\cdot e_{I_0}$ by Step~2,
and that the parity of $\mathcal{F}(c_2)$ acts on $e_{r_g(I_0)}$ by the sign $(-1)^{\#|I_0|}$.
The operators $\Upsilon_0$, $\Upsilon_1$ are even, being diagonal,
and $\Ins_{c_j}[\Upsilon_j]=\exp\left(-S^{\times}_{c_j}(A)\right)$ by~\eqref{eq:slab-crossing-factor}.

The action of the torus splits along the chain:
\[
S_{\mathbb{T}}(A)=S_{(c_0,c_1)}(A)+S_{(c_1,c_2)}(A)+S^{\times}_{c_0}(A)+S^{\times}_{c_1}(A)\qquad\text{in }\Gr_{\mathbb{T}}.
\]
To see this, compute $S_{\mathbb{T}}(A)$ with representatives of $\mathbb{X}_{\mathbb{T}}$ whose heights lie in $(c_0,c_2)$, as Step~0 permits.
The height of such a representative $x$ lies in $(c_0,c_1)$ or in $(c_1,c_2)$, and the term $M\cdot\bar\psi(x)\,\psi(x)$ belongs to the action of that open slab.
Consider a bond with the end point $x$.
If its other end point has its height in the same open interval, the term of the bond at $x$ belongs to $S_{(c_0,c_1)}(A)$ or to $S_{(c_1,c_2)}(A)$.
Otherwise its two end points are separated by the mirror at $c_0$, at $c_1$ or at $c_2$;
by Lemma~\ref{lem:slab-bonds}\eqref{it:bonds-crossing}, applied at the positions of the extended chain, no other mirror of the chain meets the bond,
and its term at $x$ is a term of $S^{\times}_{c_0}(A)$ or of $S^{\times}_{c_1}(A)$, a bond across the mirror at $c_2=c_0+p$ being identified with its translate by $-g$.
Every term of the four elements on the right arises in this way exactly once.

If $\deg G_1+\deg G_2$ is even, Lemma~\ref{lem:matrix-rules}\eqref{it:matrix-closure} expresses the trace as the first of the following two integrals,
and the second integral is equal to it:
\[
\begin{aligned}
&\int_{\mathbb{T}}\exp\left(-S^{\times}_{c_0}(A)\right)\,G_1(\varpi_1)\,\exp\left(-S_{(c_0,c_1)}(A)\right)\,\exp\left(-S^{\times}_{c_1}(A)\right)\,G_2(\varpi_2)\,\exp\left(-S_{(c_1,c_2)}(A)\right)\\
&\qquad=\int_{\mathbb{T}}\exp\left(-S_{\mathbb{T}}(A)\right)\cdot G_1(\varpi_1)\,G_2(\varpi_2).
\end{aligned}
\]
Indeed, the four exponentials are even elements, which commute with every element of $\Gr_{\mathbb{T}}$;
they can thus be gathered into one factor while $G_1(\varpi_1)$ stays to the left of $G_2(\varpi_2)$,
and by the splitting of the action their product is $\exp\left(-S_{\mathbb{T}}(A)\right)$.
If $\deg G_1+\deg G_2$ is odd, the trace vanishes by the same lemma,
and so does $\int_{\mathbb{T}}\exp\left(-S_{\mathbb{T}}(A)\right)\cdot G_1(\varpi_1)\,G_2(\varpi_2)$, the Berezin integral of an odd element;
in this case both sides of the trace formula vanish, the operator under the trace anticommuting with $U_{\mathrm{f}}$ by the parity rule of~\eqref{it:slab-hilbert-schmidt}.

\emph{Step 5: collecting the variables.}
After Steps~3 and~4 the remaining variables of summation and integration are the following:
the mirror data $\mathfrak{m}_0$ and $\mathfrak{m}_1$, with the weights $\prod_{\sigma\in\mathcal{M}^2(c_j)}\exp\left(-\tfrac12F_\sigma^2\right)$, $j=0,1$, contained in $\mu_{c_0}$ and $\mu_{c_1}$;
the values of $A$ and $z$ on the cells cut by $c_0$ and by $c_1$, with the weight of Step~3;
and the interior parts $\mathring{\varpi}_1$ and $\mathring{\varpi}_2$, with the weight $\mathcal{W}_{c_0,c_1}(\varpi_1)\cdot\mathcal{W}_{c_1,c_2}(\varpi_2)$.
The set~\eqref{eq:slab-cell-representatives}, formed for the positions $c_0<c_1<c_2$, contains exactly one representative of each cell of $\mathbb{T}$.
So every link of $\mathbb{T}$ carries exactly one of the link variables, each of which ranges over $[0,2\pi/q)$,
and every plaquette of $\mathbb{T}$ exactly one of the integer variables, each of which ranges over $\Z$;
together these variables form a pair $(A,z)\in C^1(\Gamma)\times C^2(\Gamma;\Z)$, whose values on $\mathcal{M}(c_2)$ are those of $\tau_{-g}^*\mathfrak{m}_0$, in accordance with the periodicity.
No cube lies in a mirror, so the cubes of $\mathbb{T}$ are represented by $\mathcal{X}^3(c_0)$, $\mathcal{X}^3(c_1)$, $\mathcal{B}^3(c_0,c_1)$ and $\mathcal{B}^3(c_1,c_2)$,
and the product of all weights is $\exp\left(-\tfrac12\|F\|^2\right)$ times the indicator of $dz=0$ on every cube of $\mathbb{T}$.
Collecting Steps~2 to~5, one arrives at
\[
\begin{aligned}
&\Tr_{\mathcal{H}(c_0)}\left(V_g^{-1}\,U_{\mathrm{f}}^{\frac{1+\chi(g)}{2}}\,\mathfrak{S}[G_2]\,\mathfrak{S}[G_1]\right)\\
&\qquad=\sum_{\substack{z\in C^2(\Gamma;\Z)\\ dz=0}}\ \int_{[0,2\pi/q)^{\mathbb{T}^{(1)}}}\mathrm{d}A\ \exp\left(-\tfrac12\|F\|^2\right)\cdot\int_{\mathbb{T}}\exp\left(-S_{\mathbb{T}}(A)\right)\cdot G_1(\varpi_1)\,G_2(\varpi_2),
\end{aligned}
\]
and by Definition~\ref{def:torus-integrands} the right side is $Z_{\mathbb{T}}[G_1G_2]$.

The trace formula remains true for arbitrary entries.
Let $r\ge2$ and the positions $c_0<\dots<c_r$ be as in~\eqref{it:slab-trace}, and let $G_j\in\mathcal{I}_\Xi[c_{j-1},c_j]$, $1\le j\le r$, be arbitrary elements, homogeneous or not.
Step~0 applies to them without change.
By Remark~\ref{rem:slab-growth} each $G_j$ is the sum of an even and an odd element of $\mathcal{I}_\Xi[c_{j-1},c_j]$,
and both sides of the trace formula depend linearly on each $G_j$, through the linear functional $Z_{\mathbb{T}}$ on $\mathcal{I}(\mathbb{T})$ and the linear map $G\mapsto\mathfrak{S}[G]$ of~\eqref{it:slab-hilbert-schmidt};
expanding, one writes each side as the sum of the corresponding sides for homogeneous entries, which agree.
\end{proof}

\begin{corollary}[The case of admissible products]\label{cor:slab-admissible}
Let $q>0$, $w_0,\dots,w_3>0$, $M\in\R$ and $-1\le\rho\le1$.
Let $\ell\in\mathcal{L}$, and let $\Xi=(\Gamma_\parallel,\chi_\parallel)$ be a cross-section for $\ell$.
Let the spaces $\mathcal{H}(c)$, $c\in\Adm(\ell)$, the parity $U_{\mathrm{f}}$ and the operators $V_v$, $v\in\Z^3$, be those of Definition~\ref{def:mirror-spaces}.
For $c<c'$ in $\Adm(\ell)$ and $G\in\mathcal{A}_\Xi[c,c']$ let $\mathfrak{S}[G]$ be the slab operator of the image of $G$ in $\mathcal{I}_\Xi[c,c']$
under the map of Lemma~\ref{lem:product-integrands}\eqref{it:sum-slab}, as in Definition~\ref{def:slab-kernels}.
Then the following hold for all positions $c<c'<c''$ in $\Adm(\ell)$.
\begin{enumerate}
\item
For $G\in\mathcal{A}_\Xi[c,c']$ the operator $\mathfrak{S}[G]\colon\mathcal{H}(c)\to\mathcal{H}(c')$ is of Hilbert--Schmidt class, and it depends linearly on $G$.
\item
For $G\in\mathcal{A}_\Xi[c,c']$ and $G'\in\mathcal{A}_\Xi[c',c'']$ the product $GG'$ belongs to $\mathcal{A}_\Xi[c,c'']$,
and the composition rule of Theorem~\ref{thm:slab-operators}\eqref{it:slab-composition} holds for $G$, $G'$ and $GG'$.
\item
For $G\in\mathcal{A}_\Xi[c,c']$ and $v\in\Z^3$ the element $G^{\tau_v}$ of~\eqref{eq:setting-group-action} belongs to $\mathcal{A}_\Xi[c+\ell(v),c'+\ell(v)]$,
and the covariance of Theorem~\ref{thm:slab-operators}\eqref{it:slab-covariance} holds for $G$ and $v$, with $G^{\tau_v}$ in the place of the translate of $G$.
\item
For $G\in\mathcal{A}_\Xi[c,c']$ the element $\Theta_{\ell,c'}G$ belongs to $\mathcal{A}_\Xi[c',2c'-c]$,
and the reflection rule of Theorem~\ref{thm:slab-operators}\eqref{it:slab-reflection} holds for $G$, with $\Theta_{\ell,c'}G$ in the place of the reflection of $G$ in the upper mirror.
\item
Let $\mathbb{T}=(\Gamma,\chi)$ be a torus with $\Xi(\mathbb{T},\ell)=\Xi$, let $g$ be a closing vector of $\mathbb{T}$ for $\ell$,
and let $c_0<c_1<\dots<c_r=c_0+\ell(g)$ be positions in $\Adm(\ell)$ with $r\ge2$.
For all $G_j\in\mathcal{A}_\Xi[c_{j-1},c_j]$, $1\le j\le r$, the product $G_1\cdots G_r$ belongs to $\mathcal{A}(\mathbb{T})$,
and the trace formula of Theorem~\ref{thm:slab-operators}\eqref{it:slab-trace} holds for $G_1,\dots,G_r$,
with $Z_{\mathbb{T}}$ taken on $\mathcal{A}(\mathbb{T})$ as in Definition~\ref{def:torus-integrands} and Lemma~\ref{lem:product-integrands}\eqref{it:sum-torus}.
\item
$U_{\mathrm{f}}$ commutes with $\mathfrak{S}[G]$ for every $G\in\mathcal{A}_\Xi[c,c']$, and with $V_v$ for every $v\in\Z^3$.
\end{enumerate}
\end{corollary}

\begin{proof}
For $c<c'$ in $\Adm(\ell)$ denote by $G\mapsto G(\cdot)$ the map of Lemma~\ref{lem:product-integrands}\eqref{it:sum-slab} from $\mathcal{A}_\Xi[c,c']$ to $\mathcal{I}_\Xi[c,c']$;
thus $\mathfrak{S}[G]=\mathfrak{S}[G(\cdot)]$.
By that lemma the map is linear and its values are even;
moreover $GG'$, $G^{\tau_v}$ and $\Theta_{\ell,c'}G$ lie in the classes named in~(b), (c) and~(d),
and the map carries them to the product $G(\cdot)\,G'(\cdot)$, to the translate of $G(\cdot)$ by $\tau_v$ and to the reflection of $G(\cdot)$ in the mirror at $c'$,
product, translate and reflection being those of Definition~\ref{def:slab-integrands}.

In~(a), the map $G\mapsto\mathfrak{S}[G]$ is the composite of $G\mapsto G(\cdot)$ with the linear map $\mathfrak{S}$ of Theorem~\ref{thm:slab-operators}\eqref{it:slab-hilbert-schmidt},
whose values are Hilbert--Schmidt operators $\mathcal{H}(c)\to\mathcal{H}(c')$.
Since $(GG')(\cdot)=G(\cdot)\,G'(\cdot)$, assertion~(b) is Theorem~\ref{thm:slab-operators}\eqref{it:slab-composition} for $G(\cdot)$ and $G'(\cdot)$.
Likewise, the images of $G^{\tau_v}$ and of $\Theta_{\ell,c'}G$ being the translate and the reflection of $G(\cdot)$,
Theorem~\ref{thm:slab-operators}\eqref{it:slab-covariance} and~\eqref{it:slab-reflection}, applied to $G(\cdot)$, give~(c) and~(d).

For~(e), \eqref{eq:slab-torus-reading} shows that $G_1\cdots G_r$ lies in $\mathcal{A}(\mathbb{T})$,
and that the element of $\mathcal{I}(\mathbb{T})$ to which $Z_{\mathbb{T}}$ is applied (Lemma~\ref{lem:product-integrands}\eqref{it:sum-torus})
is the map $(A,z)\mapsto G_1(\varpi_1)\cdots G_r(\varpi_r)$ built from the functions $G_j(\cdot)$.
These functions are even, hence homogeneous,
and the trace formula of Theorem~\ref{thm:slab-operators}\eqref{it:slab-trace} for $G_1(\cdot),\dots,G_r(\cdot)$ is the identity asserted in~(e).
Assertion~(f) holds because $U_{\mathrm{f}}$ commutes with every $V_v$ by Theorem~\ref{thm:slab-operators}\eqref{it:slab-hilbert-schmidt},
and with $\mathfrak{S}[G]=\mathfrak{S}[G(\cdot)]$ by the parity rule of the same item, $G(\cdot)$ being even.
\end{proof}

\begin{remark}
In part~(e) of Corollary~\ref{cor:slab-admissible} the element $G_1\cdots G_r$ is a finite sum of products that are read on $\mathbb{T}$,
and of such a product it is only required that the translates of its boxes under $\Gamma$ be pairwise disjoint.
A product that embeds in $\mathbb{T}$ must satisfy more (Definition~\ref{def:integrands}):
the translates of each of its boxes, enlarged by $[-1,1]^3$, must be pairwise disjoint as well.
For the products that do embed in $\mathbb{T}$, the functional of Definition~\ref{def:torus-integrands} is that of Definition~\ref{def:functional},
by Lemma~\ref{lem:product-integrands}\eqref{it:product-torus}.
The products of part~(e) can fail to embed, as the following example shows.
Let $\ell=x_1$ and $\mathbb{T}=C(n_1,n_2,n_3)$ with $n_1\ge4$ and $n_2,n_3>n_1+1$;
then $g=n_1e_1$ is a closing vector, $p=n_1$, and $\Gamma_\parallel=n_2\Z e_2+n_3\Z e_3$.
Take $c_0=0$, $c_1=n_1-\tfrac12$, $c_2=n_1$, and the box
\[
B=y+\left[-\frac{n_1-1}{2},\frac{n_1-1}{2}\right]^3,\qquad y\in\R^3,\quad y_1=\frac{2n_1-1}{4}.
\]
It lies in the open slab $\{0<x_1<n_1-\tfrac12\}$ and embeds in $\Gamma_\parallel$.
Hence, for every $G_1\in\mathcal{G}^{\mathrm{ev}}(B)$, the product formed by $G_1$ with the single box $B$ belongs to $\mathcal{A}_\Xi[c_0,c_1]$,
and $G_2:=1$ belongs to $\mathcal{A}_\Xi[c_1,c_2]$.
The translates $B+\gamma$, $\gamma\in\Gamma$, are pairwise disjoint, so that the product $G_1G_2$, formed with the box $B$, is read on $\mathbb{T}$.
The box $B$, however, does not embed in $\Gamma$, since $B+[-1,1]^3$ meets $B+[-1,1]^3+g$;
this product is therefore read on $\mathbb{T}$ without being one of the products that embed in $\mathbb{T}$.
For a diagonal family $\ell$ one has $\Adm(\ell)=\Z$,
and for a torus with $p(\mathbb{T},\ell)=1$ there is no chain $c_0<c_1<\dots<c_r=c_0+1$ in $\Adm(\ell)$ with $r\ge2$,
so that part~(e) has no instance on such a torus.
\end{remark}

\subsection{Coordinate link planes}\label{sec:link-planes}

\begin{definition}[Slice notation]\label{def:slice-notation}
Let $n_1,n_2,n_3\ge4$ be integers and $\mathbb{T}:=C(n_1,n_2,n_3)=(\Gamma,\chi)$,
so that $\Gamma=n_1\Z e_1+n_2\Z e_2+n_3\Z e_3$.
The family is $\ell=x_1$, and $e_1$ is the direction of time.
The cross-section of $\mathbb{T}$ for $x_1$ is the rectangular cross-section of Definition~\ref{def:slices},
\[
\Xi=\Xi(\mathbb{T},x_1)=\left(n_2\Z e_2+n_3\Z e_3;\ \chi_\parallel\right),
\qquad
\chi_\parallel(n_2e_2)=\chi_\parallel(n_3e_3)=-1;
\]
the vector $g:=n_1e_1=n_1n_\ell$ is a closing vector of $\mathbb{T}$ for $x_1$,
with $\chi(g)=-1$ and the height period $p=n_1$.
For every open interval $\mathbb{I}\subset\R$, the \emph{reading} in $\Gr_{\mathbb{T}}$ is the homomorphism of algebras
from $\Gr_\Xi(\mathbb{I})$ (Definition~\ref{def:slab-grassmann}) to $\Gr_{\mathbb{T}}$ (Definition~\ref{def:tori})
that sends every generator $\psi_\alpha(x)$, $\bar\psi_\alpha(x)$, $x\in\mathbb{X}(\mathbb{I})$, to the generator of $\Gr_{\mathbb{T}}$ with the same name.
It is well defined:
$\Gamma_\parallel=n_2\Z e_2+n_3\Z e_3$ is contained in $\Gamma$, and $\chi_\parallel$ is the restriction of $\chi$ to $\Gamma_\parallel$,
so that the relations of $\Gr_\Xi(\mathbb{I})$ hold in $\Gr_{\mathbb{T}}$.
For $\mathbb{I}\subset\mathbb{I}'$ the reading of $\Gr_\Xi(\mathbb{I}')$ restricts to that of its subalgebra $\Gr_\Xi(\mathbb{I})$,
both maps sending each generator to the generator with the same name;
on $\Gr_\Xi((c-p,c+p))$, $c\in\frac12\Z$, the reading is that of Lemma~\ref{lem:matrix-rules}\eqref{it:matrix-closure}.
Under the integral $\int_{\mathbb{T}}$, and in identities between elements of $\Gr_{\mathbb{T}}$,
an element of one of the algebras $\Gr_\Xi(\mathbb{I})$ stands for its reading.
The slices, the modes, the Fock spaces and the operators of the cross-section are those of
Definitions~\ref{def:slices}, \ref{def:modes} and~\ref{def:slice-operators}.
They are tied to the objects of Definitions~\ref{def:mode-translations} and~\ref{def:slab-matrices} as follows.
\begin{enumerate}
\item
For $j\in\Z$ the ordered set of representatives of the modes at the site plane $\{x_1=j\}$
is $\mathcal{J}(j):=\mathcal{J}^{\mathrm{s}}(\Lambda_{\mathrm{s}})$, formed at this plane,
and that at the link plane $\{x_1=j+1/2\}$ is $\mathcal{J}(j+1/2):=\mathcal{J}^{\mathrm{l}}(\Lambda_{\mathrm{s}})$, formed at that plane.
The representatives and their orders are transported by $\tau_{e_1}$:
$\mathcal{J}(c+1)=\mathcal{J}(c)+e_1$, with the transported order, for every $c\in\frac12\Z$;
this is consistent with Definition~\ref{def:modes}, whose orders are compatible with $\tau_{e_1}$.
\item
For every $c\in\frac12\Z$ the space $\mathcal{F}(c+1)$ is identified with $\mathcal{F}(c)$ through $U_{e_1}$,
$U_{e_1}e_I=e_{I+e_1}$, which carries $\omega_c$ to $\omega_{c+1}$
(Definition~\ref{def:mode-translations}\eqref{it:fock-translations}).
In this identification $\mathcal{F}(j)=\mathcal{F}'$, $\mathcal{F}(j+1/2)=\mathcal{F}$,
$\omega_j=\omega'$ and $\omega_{j+1/2}=\omega$ for every $j\in\Z$.
\item
The half slab is $\mathbb{X}_{\mathrm{h}}=\mathbb{X}_\Xi((j,j+1/2))$,
with the Grassmann algebra $\Gr_\Xi((j,j+1/2))$ of Definition~\ref{def:slab-grassmann}.
For $A\in C^1(\Gamma)$ one has $S_{\mathrm{h}}(A_j)=S_{(j,j+1/2)}(A)$;
by Lemma~\ref{lem:slab-bonds}\eqref{it:bonds-slab-action} this element reads only the phases of the links of
$[j,j+1/2]^{(1)}=\mathcal{M}^1(j)$, that is, the phases of the slice $A_j$.
\item
For every $j\in\Z$ and every $A\in\R^{\Lambda^{(1)}_{\mathrm{s}}}$
the half-slab matrix is the slab matrix of the half slab $\mathbb{X}_\Xi((j,j+1/2))$,
\[
\mathfrak{N}(A)=\mathfrak{N}_{j,j+1/2}\left[\exp\left(-S_{\mathrm{h}}(A)\right)\right],
\qquad
\widehat{\mathfrak{N}}(A)=\omega^{\frac{1}{2}}\,\mathfrak{N}(A)\,\omega'^{\frac{1}{2}} .
\]
\item
At the link plane $\{x_1=j+1/2\}$ the products of the phases of the modes are
$u_I=\exp\left(iq\langle a_j,\mathfrak{q}(I)\rangle\right)$ (Definition~\ref{def:modes});
thus $\diag(u_I)=u(a_j)$.
\item
For $A\in\R^{\Lambda^{(1)}_{\mathrm{s}}}$ and $k\in\Z$ the magnetic curvature in the flux sector $k$ is
\[
F^{\mathrm{mag}}(A,k):=d_{\mathrm{s}}A+\frac{2\pi}{q}\,k\,1_{\sigma_0},
\]
and in this notation
$\mathfrak{M}_k(A)=\exp\left(-\frac{1}{4}\|F^{\mathrm{mag}}(A,k)\|^2\right)\cdot\widehat{\mathfrak{N}}(A)$
(Definition~\ref{def:slice-operators}).
\end{enumerate}
\end{definition}

\begin{proposition}[Slice factorisation]\label{prop:slice-factorisation}
Let $w_0,\dots,w_3>0$, $M\in\R$, $q>0$ and $-1\le\rho\le1$.
Let $n_1,n_2,n_3\ge4$ and $\mathbb{T}=C(n_1,n_2,n_3)$, with the notation of Definition~\ref{def:slice-notation}.
Let $\Upsilon_j\in\End\mathcal{F}$, $j\in\Z/n_1\Z$, be even operators, the class $j$ being represented by the integer $0\le j<n_1$,
and let $A\in C^1(\Gamma)$ be a real one-cochain on $\mathbb{T}$.
For $0\le j<n_1$ the elements $S^{\times}_{j+1/2}(A)$ and $\Ins_{j+1/2}\left[\omega^{\frac{1}{2}}\Upsilon_j\,\omega^{\frac{1}{2}}\right]$ of $\Gr_\Xi((j,j+1))$
are formed at the link plane $\{x_1=j+1/2\}$ (Lemma~\ref{lem:slab-bonds}\eqref{it:bonds-action-split} and Definition~\ref{def:slab-matrices}),
and on the left of the following identity they stand for their readings in $\Gr_{\mathbb{T}}$ (Definition~\ref{def:slice-notation}).
Then
\begin{equation}\label{eq:slab-slice-factorisation}
\begin{aligned}
&\int_{\mathbb{T}}\exp\left(-S_{\mathbb{T}}(A)\right)\prod_{j\in\Z/n_1\Z}\left(\exp\left(S^{\times}_{j+1/2}(A)\right)\,\Ins_{j+1/2}\left[\omega^{\frac{1}{2}}\Upsilon_j\,\omega^{\frac{1}{2}}\right]\right)\\
&\qquad=\operatorname{tr}_{\mathcal{F}'}\prod_{j=n_1-1,\dots,0}\left(\widehat{\mathfrak{N}}(A_{j+1})^{*}\,\Upsilon_j\,\widehat{\mathfrak{N}}(A_j)\right),
\qquad A_{n_1}=A_0 .
\end{aligned}
\end{equation}
The factors of the product on the left are even, so that their order is immaterial;
on the right, later times stand to the left.
If $\Upsilon_j=u(a_j)$ for every $j$, the left side equals $\det D^{\mathbb{T}}_A$.
\end{proposition}

\begin{proof}
\emph{The action along a chain of planes.}
Fix $A$, and let $c_0<c_1<\dots<c_r=c_0+n_1$ be points of $\frac12\Z$, where $r\ge1$.
Then, in $\Gr_{\mathbb{T}}$,
\begin{equation}\label{eq:slab-action-chain}
S_{\mathbb{T}}(A)=\sum_{k=1}^{r}S_{(c_{k-1},c_k)}(A)+\sum_{k=0}^{r-1}S^{\times}_{c_k}(A).
\end{equation}
To see this, extend the chain to a sequence $(c_k)_{k\in\Z}$ through the rule $c_{k+r}:=c_k+n_1$, $k\in\Z$;
the rule is compatible with the given points, since $c_r=c_0+n_1$.
The sequence so obtained is strictly increasing and unbounded in both directions,
and the translation by $g$ maps the open slab $\{c_{k-1}<x_1<c_k\}$ onto $\{c_{k+r-1}<x_1<c_{k+r}\}$
and the plane $\{x_1=c_k\}$ onto $\{x_1=c_{k+r}\}$.
The terms of $S_{\mathbb{T}}(A)$ are the diagonal terms $M\bar\psi(x)\psi(x)$ and the terms of the oriented bonds
(Definitions~\ref{def:chamber-operator} and~\ref{def:tori}).
No site lies on a mirror (Lemma~\ref{lem:slab-bonds}\eqref{it:bonds-crossing});
every site therefore lies in exactly one of the open slabs, and its diagonal term belongs to the action of that slab.
Let $\{x,x'\}$ be a bond.
If $x$ and $x'$ lie in one open slab, the terms of the bond belong to the action of this slab.
Otherwise $x$ and $x'$ lie on different sides of a plane $\{x_1=c_k\}$;
by Lemma~\ref{lem:slab-bonds}\eqref{it:bonds-crossing} the bond is $\{x,\theta_{c_k}x\}$
and its open segment meets no other mirror of the family,
so it crosses the plane $\{x_1=c_k\}$ and no other plane of the sequence, and its terms belong to $S^{\times}_{c_k}(A)$.
Read in $\Gr_{\mathbb{T}}$, a term does not change under the translation by $g$:
the hopping matrix of a bond depends only on its type and its direction,
its phase is unchanged because $A$ is $\Gamma$-periodic,
and the fields at the two ends of the bond acquire the same factor $\chi(g)$, whose square is $1$.
The slabs with $1\le k\le r$ and the planes with $0\le k\le r-1$ represent the slabs and the planes of the sequence
modulo this translation.
Hence every term of $S_{\mathbb{T}}(A)$ occurs exactly once on the right side of \eqref{eq:slab-action-chain},
and no other term occurs there;
the closing term $S^{\times}_{c_0}(A)$ collects the bonds that cross one of the planes $\{x_1=c\}$ with $c\in c_0+n_1\Z$.

\emph{The integrand.}
Apply \eqref{eq:slab-action-chain} to the chain with $r:=2n_1$,
\[
c_{2j}:=j,\qquad c_{2j+1}:=j+\frac{1}{2}\qquad(0\le j<n_1),\qquad c_r:=n_1 .
\]
All summands in \eqref{eq:slab-action-chain} are even, since every term of the action has degree two;
so they commute with each other, and $\exp\left(-S_{\mathbb{T}}(A)\right)$ is the product of their exponentials.
The factors $\exp\left(-S^{\times}_{j+1/2}(A)\right)$ cancel against the factors $\exp\left(S^{\times}_{j+1/2}(A)\right)$
on the left of \eqref{eq:slab-slice-factorisation}.
At the site plane $\{x_1=j\}$ both ends of a crossing bond have the same nearest photon site,
so that the bond carries no phase (Definition~\ref{def:modes}), and \eqref{eq:slab-crossing-factor} gives
$\exp\left(-S^{\times}_j(A)\right)=\Ins_j[\omega']$.
An even operator has the entry $0$ at $(I,J)$ unless $\#|I|+\#|J|$ is even (Definition~\ref{def:modes}),
and $\Theta_{j+1/2}(\zeta_J)\,\zeta_I$ has the parity of $\#|I|+\#|J|$;
so $\Ins_{j+1/2}\left[\omega^{\frac{1}{2}}\Upsilon_j\,\omega^{\frac{1}{2}}\right]$ is an even element.
The integrand on the left of \eqref{eq:slab-slice-factorisation} is thus the product of the mutually commuting even elements
\[
\Ins_j[\omega'],\qquad
\exp\left(-S_{(j,j+1/2)}(A)\right),\qquad
\Ins_{j+1/2}\left[\omega^{\frac{1}{2}}\Upsilon_j\,\omega^{\frac{1}{2}}\right],\qquad
\exp\left(-S_{(j+1/2,j+1)}(A)\right),
\]
$0\le j<n_1$.
Arranged as in Lemma~\ref{lem:matrix-rules}\eqref{it:matrix-closure}, this product reads
\[
\Ins_{c_0}\left[\widetilde{\Upsilon}_0\right]G_1\,\Ins_{c_1}\left[\widetilde{\Upsilon}_1\right]G_2\cdots\Ins_{c_{r-1}}\left[\widetilde{\Upsilon}_{r-1}\right]G_r,
\]
where, for $0\le j<n_1$,
\[
\begin{aligned}
\widetilde{\Upsilon}_{2j}&:=\omega',
&\qquad
\widetilde{\Upsilon}_{2j+1}&:=\omega^{\frac{1}{2}}\Upsilon_j\,\omega^{\frac{1}{2}},\\
G_{2j+1}&:=\exp\left(-S_{(j,j+1/2)}(A)\right),
&
G_{2j+2}&:=\exp\left(-S_{(j+1/2,j+1)}(A)\right).
\end{aligned}
\]

\emph{Closure.}
The hypotheses of Lemma~\ref{lem:matrix-rules}\eqref{it:matrix-closure} hold:
$\Xi(\mathbb{T},x_1)=\Xi$;
$g$ is a closing vector with $x_1(g)=n_1$;
$c_0<\dots<c_r=c_0+n_1$;
each $G_k\in\Gr_\Xi((c_{k-1},c_k))$ is even, so that the degrees add up to an even number;
each $\widetilde{\Upsilon}_k$ is an even operator on $\mathcal{F}(c_k)$;
and all factors lie in $\Gr_\Xi((-1/2,n_1))$, which is contained in $\Gr_\Xi((c_0-n_1,c_0+n_1))=\Gr_\Xi((-n_1,n_1))$ because $n_1\ge4$.
The reading of $\Gr_\Xi((-n_1,n_1))$ is a homomorphism of algebras,
and it restricts to the reading of every subalgebra $\Gr_\Xi(\mathbb{I})$ with $\mathbb{I}\subset(-n_1,n_1)$ (Definition~\ref{def:slice-notation});
hence the integrand on the left of \eqref{eq:slab-slice-factorisation}, the product of the readings of the factors,
is the reading of their product formed in $\Gr_\Xi((-n_1,n_1))$, as that lemma requires.
That lemma writes the integral as the trace over $\mathcal{F}(0)$ of $U_g^{-1}$,
times the parity factor $U_{\mathrm{f}}^{\frac{1+\chi(g)}{2}}$,
times the product of the slab matrices $\mathfrak{N}_{c_{k-1},c_k}[G_k]$ and the operators $\widetilde{\Upsilon}_k$.
The parity factor is the identity, since $\chi(g)=-1$.
Since $g\in\Z n_\ell$, one has $U_ge_I=e_{I+g}$ (Definition~\ref{def:mode-translations}\eqref{it:fock-translations}),
and $U_g$ is the identification of $\mathcal{F}(0)$ with $\mathcal{F}(n_1)$.
Therefore the left side of \eqref{eq:slab-slice-factorisation} equals
\[
\operatorname{tr}_{\mathcal{F}'}\left(\mathfrak{N}_{c_{r-1},c_r}[G_r]\,\widetilde{\Upsilon}_{r-1}\,\mathfrak{N}_{c_{r-2},c_{r-1}}[G_{r-1}]\cdots\widetilde{\Upsilon}_1\,\mathfrak{N}_{c_0,c_1}[G_1]\,\widetilde{\Upsilon}_0\right).
\]

\emph{The half slabs.}
By Definition~\ref{def:slice-notation}, $\mathfrak{N}_{j,j+1/2}[G_{2j+1}]=\mathfrak{N}(A_j)$.
For the half slab above the link plane, let $\theta:=\theta_{x_1,j+1/2}$;
it maps the open slab $\{j<x_1<j+1/2\}$ onto $\{j+1/2<x_1<j+1\}$.
By Proposition~\ref{prop:symmetries}\eqref{it:action-invariance},
the map $\Theta_{j+1/2}$ carries the term of an oriented bond $(x,x')$ at $\theta^{*}A$ to the term of $(\theta x',\theta x)$ at $A$.
It carries the diagonal term at $x$ to the diagonal term at $\theta x$:
the rule~\eqref{eq:setting-reflection-bilinear} for reflected bilinear elements (Lemma~\ref{lem:sign-rules}),
applied with the upward unit normal $n=e_1$ of the family $x_1$ and with $Y=M\cdot1$, where $1$ is the unit $4\times4$ matrix, gives
\[
\Theta_{j+1/2}\left(\bar\psi(x)\,(M\cdot1)\,\psi(x)\right)=\bar\psi(\theta x)\,\gamma_n\,(M\cdot1)^{*}\,\gamma_n\,\psi(\theta x)=M\bar\psi(\theta x)\,\psi(\theta x),
\]
the last equality because $(M\cdot1)^{*}=M\cdot1$, $M$ being real, and $\gamma_n^2=1$ (Lemma~\ref{lem:spin-identities}(c)).
Consequently $\Theta_{j+1/2}\left(S_{(j,j+1/2)}(\theta^{*}A)\right)=S_{(j+1/2,j+1)}(A)$.
By \eqref{eq:setting-slice-reflections}, $(\theta^{*}A)_j=A_{j+1}$, whence $S_{(j,j+1/2)}(\theta^{*}A)=S_{\mathrm{h}}(A_{j+1})$.
Every term of $S_{\mathrm{h}}(A_{j+1})$ has degree two, so this element is nilpotent,
and its exponential is a polynomial in it with real coefficients;
as $\Theta_{j+1/2}$ is antilinear and maps powers of an element to the powers of its image,
\[
G_{2j+2}=\Theta_{j+1/2}\left(\exp\left(-S_{\mathrm{h}}(A_{j+1})\right)\right).
\]
The element $\exp\left(-S_{\mathrm{h}}(A_{j+1})\right)$ of $\Gr_\Xi((j,j+1/2))$ has the slab matrix $\mathfrak{N}(A_{j+1})$
(Definition~\ref{def:slice-notation}),
and Lemma~\ref{lem:matrix-rules}\eqref{it:matrix-adjoint}, applied with $c=j$, $c'=j+1/2$, $c''=j+1$ and $v_{c,c'}=e_1$, gives
\[
\mathfrak{N}_{j+1/2,j+1}[G_{2j+2}]=U_{e_1}\,\mathfrak{N}(A_{j+1})^{*},
\]
which is $\mathfrak{N}(A_{j+1})^{*}$ in the identification of Definition~\ref{def:slice-notation}.

\emph{The trace.}
Since $A$ is $\Gamma$-periodic, $A_{n_1}=A_0$, and the left side of \eqref{eq:slab-slice-factorisation} is
\[
\operatorname{tr}_{\mathcal{F}'}\prod_{j=n_1-1,\dots,0}\left(\mathfrak{N}(A_{j+1})^{*}\,\omega^{\frac{1}{2}}\Upsilon_j\,\omega^{\frac{1}{2}}\,\mathfrak{N}(A_j)\,\omega'\right).
\]
By Definition~\ref{def:modes} the coefficient of a mode is $\frac{w_i}{2}(1+\rho)$ or $\frac{w_i}{2}(1-\rho)$, $i$ being the type of its bond;
as the weights are positive and $-1\le\rho\le1$, every coefficient is non-negative,
and so is every diagonal entry of $\omega$ and of $\omega'$, a product of coefficients.
Hence $\omega^{\frac{1}{2}}$ and $\omega'^{\frac{1}{2}}$ are diagonal operators with non-negative real entries;
in particular they are self-adjoint, and $\omega'=\omega'^{\frac{1}{2}}\omega'^{\frac{1}{2}}$.
Moving the last factor $\omega'^{\frac{1}{2}}$ to the front leaves the trace unchanged and yields
\[
\operatorname{tr}_{\mathcal{F}'}\prod_{j=n_1-1,\dots,0}\left(\omega'^{\frac{1}{2}}\,\mathfrak{N}(A_{j+1})^{*}\,\omega^{\frac{1}{2}}\;\Upsilon_j\;\omega^{\frac{1}{2}}\,\mathfrak{N}(A_j)\,\omega'^{\frac{1}{2}}\right).
\]
The self-adjointness of $\omega^{\frac{1}{2}}$ and $\omega'^{\frac{1}{2}}$ gives
\[
\omega'^{\frac{1}{2}}\,\mathfrak{N}(A_{j+1})^{*}\,\omega^{\frac{1}{2}}
=\left(\omega^{\frac{1}{2}}\,\mathfrak{N}(A_{j+1})\,\omega'^{\frac{1}{2}}\right)^{*}
=\widehat{\mathfrak{N}}(A_{j+1})^{*},
\]
while $\omega^{\frac{1}{2}}\,\mathfrak{N}(A_j)\,\omega'^{\frac{1}{2}}=\widehat{\mathfrak{N}}(A_j)$ by definition;
the last trace is therefore the right side of \eqref{eq:slab-slice-factorisation}.

\emph{The determinant.}
Let $\Upsilon_j=u(a_j)$ for every $j$.
Then $\omega^{\frac{1}{2}}u(a_j)\,\omega^{\frac{1}{2}}=\omega\,\diag(u_I)$,
with the phases $u_I$ of the link plane $\{x_1=j+1/2\}$ (Definition~\ref{def:slice-notation}),
and \eqref{eq:slab-crossing-factor} gives
$\Ins_{j+1/2}\left[\omega^{\frac{1}{2}}u(a_j)\,\omega^{\frac{1}{2}}\right]=\exp\left(-S^{\times}_{j+1/2}(A)\right)$.
The integrand on the left of \eqref{eq:slab-slice-factorisation} reduces to $\exp\left(-S_{\mathbb{T}}(A)\right)$,
and $\int_{\mathbb{T}}\exp\left(-S_{\mathbb{T}}(A)\right)=\det D^{\mathbb{T}}_A$ by Definition~\ref{def:tori}.
\end{proof}

\begin{remark}
In \eqref{eq:slab-slice-factorisation} the row index of $\Upsilon_j$,
which by \eqref{eq:setting-ordered-monomials} belongs to the upper side of the plane $\{x_1=j+1/2\}$,
is contracted with the factor $\widehat{\mathfrak{N}}(A_{j+1})^{*}$ of the slab above this plane;
this is the order of the composition rule, Lemma~\ref{lem:matrix-rules}\eqref{it:matrix-composition}.
For $\Upsilon_j=u(a_j)$ the identity writes $\det D^{\mathbb{T}}_A$ in the form of the factorisation of the determinant
over time slices in \cite[Section~10.2.1]{NoncompactQED}:
with $s=1$ there, $\widehat{\mathfrak{N}}$, $u$ and $n_1$ take the places of $Q$, $W$ and $n_0$,
and two Fock spaces take the place of one, the half-slab matrices mapping $\mathcal{F}'$ into $\mathcal{F}$
and the trace being taken over $\mathcal{F}'$.
\end{remark}

\begin{lemma}[Half-slab matrices]\label{lem:half-slab-matrices}
Let $w_0,\dots,w_3>0$, $M\in\R$, $q>0$ and $-1\le\rho\le1$.
Let $n_2,n_3\ge4$ be integers, let $\Xi$ be the rectangular cross-section with the sides $n_2,n_3$
and $\Lambda_{\mathrm{s}}$ its slice torus (Definition~\ref{def:slices}),
let $\mathcal{F}$, $\mathcal{F}'$, $\widehat{\mathfrak{N}}$ and $u$ be the objects of Definition~\ref{def:modes} for $\Xi$,
and let $\mathcal{H}^{\mathrm{l}}$, $U_{\mathrm{G}}$, $P_{\mathrm{G}}$ and $\mathfrak{M}$ be those of Definition~\ref{def:slice-operators}.
\begin{enumerate}
\item\label{it:half-slab-entries}
The entries of $\widehat{\mathfrak{N}}(A)$, $A\in\R^{\Lambda^{(1)}_{\mathrm{s}}}$,
are polynomials in the phases $\exp\left(\pm iqA(b)\right)$, $b\in\Lambda^{(1)}_{\mathrm{s}}$.
The map $A\mapsto\widehat{\mathfrak{N}}(A)$ is continuous and bounded, $\sup_A\|\widehat{\mathfrak{N}}(A)\|<\infty$,
and it is invariant under $A\mapsto A+\frac{2\pi}{q}k$ for every $k\in\Z^{\Lambda^{(1)}_{\mathrm{s}}}$.
The map $a\mapsto u(a)$ is a unitary representation of $\R^{\Lambda_{\mathrm{s}}}$ by diagonal matrices
with the integer charges $\mathfrak{q}(I)\in\Z^{\Lambda_{\mathrm{s}}}$,
and $u(a)=1$ for $a\in\frac{2\pi}{q}\Z^{\Lambda_{\mathrm{s}}}$.
\item\label{it:half-slab-gauge}
For a mode $\zeta'_{i'}$ of the site plane over $y\in\Lambda_{\mathrm{s}}$ put
$\mathfrak{q}'(i'):=\delta_y$ if it belongs to the group $1$ or $2$,
and $\mathfrak{q}'(i'):=-\delta_y$ if it belongs to the group $3$ or $4$;
let $\mathfrak{q}'(I'):=\sum_{i'\in I'}\mathfrak{q}'(i')$,
and let $u'(a)$ be the diagonal operator on $\mathcal{F}'$ with the entries $\exp\left(iq\langle a,\mathfrak{q}'(I')\rangle\right)$.
Then, for $A\in\R^{\Lambda^{(1)}_{\mathrm{s}}}$ and $\lambda\in\R^{\Lambda_{\mathrm{s}}}$,
\[
\widehat{\mathfrak{N}}(A-d_{\mathrm{s}}\lambda)=u(\lambda)^{-1}\,\widehat{\mathfrak{N}}(A)\,u'(\lambda).
\]
Consequently $U_{\mathrm{G}}(a)\,\mathfrak{M}\mathfrak{M}^{*}\,U_{\mathrm{G}}(a)^{-1}=\mathfrak{M}\mathfrak{M}^{*}$ for every $a\in\R^{\Lambda_{\mathrm{s}}}$:
the operator $\mathfrak{M}\mathfrak{M}^{*}$ on $\mathcal{H}^{\mathrm{l}}$ commutes with every $U_{\mathrm{G}}(a)$ and with $P_{\mathrm{G}}$.
\end{enumerate}
\end{lemma}

\begin{proof}
By Definition~\ref{def:slice-notation} the half-slab matrix may be computed in the half slab $\mathbb{X}_{\mathrm{h}}=\mathbb{X}_\Xi((j,j+1/2))$ for any fixed $j\in\Z$,
with the modes $\zeta'_{i'}$ of the site plane $\{x_1=j\}$ and the modes $\zeta_i$ of the link plane $\{x_1=j+1/2\}$:
\[
\mathfrak{N}(A)_{I,I'}=\int_{\mathbb{X}_{\mathrm{h}}}\zeta'_{I'}\,\exp\left(-S_{\mathrm{h}}(A)\right)\,\Theta_{j+1/2}(\zeta_I).
\]

Consider first \eqref{it:half-slab-entries}.
By Lemma~\ref{lem:slab-bonds}\eqref{it:bonds-slab-action} the element $S_{\mathrm{h}}(A)$ is even,
and its coefficients in a monomial basis are polynomials in the phases $\exp\left(\pm iqA(b)\right)$
of the links $b$ of $[j,j+1/2]^{(1)}=\mathcal{M}^1(j)$, which are the links of $\Lambda^{(1)}_{\mathrm{s}}$.
Each term of $S_{\mathrm{h}}(A)$ has degree two,
so every power of $S_{\mathrm{h}}(A)$ whose exponent is larger than half the number of generators of $\Gr_\Xi((j,j+1/2))$ vanishes;
hence $\exp\left(-S_{\mathrm{h}}(A)\right)$ is a polynomial in $S_{\mathrm{h}}(A)$ with rational coefficients,
and its coefficients are again polynomials in the phases.
The elements $\zeta'_{I'}$ and $\Theta_{j+1/2}(\zeta_I)$ do not depend on $A$,
and the Berezin integral over $\mathbb{X}_{\mathrm{h}}$ is the coefficient of $\Psi^{\mathrm{top}}(\mathbb{X}_{\mathrm{h}})$
(Definition~\ref{def:left-derivatives}).
Every entry of $\mathfrak{N}(A)$, and so every entry of $\widehat{\mathfrak{N}}(A)=\omega^{\frac{1}{2}}\,\mathfrak{N}(A)\,\omega'^{\frac{1}{2}}$,
is therefore a polynomial in the phases.
For real $A$ the phases have modulus one;
a polynomial in them is a continuous function of $A$, bounded by the sum of the moduli of its coefficients,
and it is invariant under $A\mapsto A+\frac{2\pi}{q}k$, because $\exp\left(iq\cdot\frac{2\pi}{q}k(b)\right)=1$.
The spaces $\mathcal{F}$ and $\mathcal{F}'$ have finite dimension, so the bound on the entries bounds the operator norm.
The operator $u(a)$ is diagonal with the entries $\exp\left(iq\langle a,\mathfrak{q}(I)\rangle\right)$, of modulus one;
$u(a+a')=u(a)\,u(a')$, $u(0)=1$, and $a\mapsto u(a)$ is continuous.
Each charge $\mathfrak{q}(i)$ is $\pm\delta_y$ for some $y$, so $\mathfrak{q}(I)\in\Z^{\Lambda_{\mathrm{s}}}$,
and for $a=\frac{2\pi}{q}k$ with $k\in\Z^{\Lambda_{\mathrm{s}}}$ every entry equals $\exp\left(2\pi i\langle k,\mathfrak{q}(I)\rangle\right)=1$.

Turning to \eqref{it:half-slab-gauge}, let $\lambda\in\R^{\Lambda_{\mathrm{s}}}$,
and let $G\mapsto G^{\lambda}$ be the automorphism of $\Gr_\Xi((j,j+1/2))$ with
\[
\psi(x)\mapsto\exp\left(iq\lambda(y)\right)\cdot\psi(x),
\qquad
\bar\psi(x)\mapsto\exp\left(-iq\lambda(y)\right)\cdot\bar\psi(x),
\qquad
(j,y)=\pi(x);
\]
it is compatible with the relations of the cross-section, $\lambda$ being a function on $\Lambda_{\mathrm{s}}$.
Each term of the action is invariant under the gauge transformation \eqref{eq:setting-gauge}.
Take there a gauge function equal to $\lambda(y)$ at the photon sites $(j,y)$;
on the links of the plane $\{x_1=j\}$ its coboundary is $d_{\mathrm{s}}\lambda$.
The terms of $S_{\mathrm{h}}$ join sites whose nearest photon sites lie in this plane, and they read only links of this plane; hence
\[
\left(S_{\mathrm{h}}(A-d_{\mathrm{s}}\lambda)\right)^{\lambda}=S_{\mathrm{h}}(A).
\]
The automorphism multiplies every monomial by a number and fixes each $\Psi^{\mathrm{top}}(x)=\prod_{\alpha=1}^4\psi_\alpha(x)\,\bar\psi_\alpha(x)$;
it does not change the coefficient of $\Psi^{\mathrm{top}}(\mathbb{X}_{\mathrm{h}})$, and so
$\int_{\mathbb{X}_{\mathrm{h}}}G^{\lambda}=\int_{\mathbb{X}_{\mathrm{h}}}G$ for every $G\in\Gr_\Xi((j,j+1/2))$.
By Definition~\ref{def:modes}, a mode of the group $1$ or $2$ is a linear combination of the components of $\psi$ at the upper end of its bond,
and a mode of the group $3$ or $4$ a linear combination of the components of $\bar\psi$ there;
by \eqref{eq:setting-reflection-fields}, its image under the reflection $\Theta$ of the mirror that the bond crosses
is a linear combination of the components of $\bar\psi$, respectively of $\psi$, at the lower end.
The modes $\zeta'_{i'}$ sit at upper ends and the elements $\Theta_{j+1/2}(\zeta_i)$ at lower ends,
all of them in $\mathbb{X}_{\mathrm{h}}$, at sites whose nearest photon site is $(j,y)$ when the bond lies over $y$.
It follows that
\begin{equation}\label{eq:slab-gauge-modes}
\left(\zeta'_{I'}\right)^{\lambda}=\exp\left(iq\langle\lambda,\mathfrak{q}'(I')\rangle\right)\cdot\zeta'_{I'},
\qquad
\left(\Theta_{j+1/2}(\zeta_I)\right)^{\lambda}=\exp\left(-iq\langle\lambda,\mathfrak{q}(I)\rangle\right)\cdot\Theta_{j+1/2}(\zeta_I).
\end{equation}
With the invariance of the Berezin integral and \eqref{eq:slab-gauge-modes},
\[
\begin{aligned}
\mathfrak{N}(A-d_{\mathrm{s}}\lambda)_{I,I'}
&=\int_{\mathbb{X}_{\mathrm{h}}}\left(\zeta'_{I'}\,\exp\left(-S_{\mathrm{h}}(A-d_{\mathrm{s}}\lambda)\right)\,\Theta_{j+1/2}(\zeta_I)\right)^{\lambda}\\
&=\exp\left(iq\langle\lambda,\mathfrak{q}'(I')\rangle-iq\langle\lambda,\mathfrak{q}(I)\rangle\right)\cdot\mathfrak{N}(A)_{I,I'},
\end{aligned}
\]
that is, $\mathfrak{N}(A-d_{\mathrm{s}}\lambda)=u(\lambda)^{-1}\,\mathfrak{N}(A)\,u'(\lambda)$.
The diagonal operators $u(\lambda)$ and $u'(\lambda)$ commute with $\omega$ and $\omega'$,
which gives the asserted identity for $\widehat{\mathfrak{N}}$.

In the flux sector $k$ the operator $\mathfrak{M}\mathfrak{M}^{*}$ acts as the multiplication by
\[
\mathfrak{M}_k(A)\,\mathfrak{M}_k(A)^{*}=\exp\left(-\frac{1}{2}\|F^{\mathrm{mag}}(A,k)\|^2\right)\cdot\widehat{\mathfrak{N}}(A)\,\widehat{\mathfrak{N}}(A)^{*}
\]
(Definitions~\ref{def:slice-operators} and~\ref{def:slice-notation}).
Let $a\in\R^{\Lambda_{\mathrm{s}}}$.
Since $d_{\mathrm{s}}d_{\mathrm{s}}a=0$, one has $F^{\mathrm{mag}}(A-d_{\mathrm{s}}a,k)=F^{\mathrm{mag}}(A,k)$;
together with the identity for $\widehat{\mathfrak{N}}$ and the unitarity of $u(a)$ and $u'(a)$ this gives
\[
\mathfrak{M}_k(A-d_{\mathrm{s}}a)\,\mathfrak{M}_k(A-d_{\mathrm{s}}a)^{*}=u(a)^{-1}\,\mathfrak{M}_k(A)\,\mathfrak{M}_k(A)^{*}\,u(a).
\]
For $\Phi=(\Phi_k)_{k\in\Z}\in\mathcal{H}^{\mathrm{l}}$ the definition of $U_{\mathrm{G}}(a)$ in Definition~\ref{def:slice-operators} gives
\[
\begin{aligned}
\left(U_{\mathrm{G}}(a)\,\mathfrak{M}\mathfrak{M}^{*}\,U_{\mathrm{G}}(a)^{-1}\Phi\right)_k(A)
&=u(a)\,\mathfrak{M}_k(A-d_{\mathrm{s}}a)\,\mathfrak{M}_k(A-d_{\mathrm{s}}a)^{*}\,u(a)^{-1}\,\Phi_k(A)\\
&=\mathfrak{M}_k(A)\,\mathfrak{M}_k(A)^{*}\,\Phi_k(A).
\end{aligned}
\]
So $\mathfrak{M}\mathfrak{M}^{*}$ commutes with every $U_{\mathrm{G}}(a)$,
and hence with
\[
P_{\mathrm{G}}=\left(\frac{q}{2\pi}\right)^{|\Lambda_{\mathrm{s}}|}\int_{[0,2\pi/q)^{\Lambda_{\mathrm{s}}}}\mathrm{d}a\,U_{\mathrm{G}}(a)
\]
(Definition~\ref{def:slice-operators}).
\end{proof}

\subsubsection{Local insertions}\label{sec:local-insertions}

\begin{definition}[Index sets of modes]\label{def:index-sets}
Let $\Pi=(\ell,c)$ be a mirror, and let $c^{\pm}=c\pm|n_\ell|^2/2$ be the neighbours of $c$ in $\Adm(\ell)$.
An \emph{index set of modes} at $\Pi$ is a finite subset $\mathcal{J}$ of $\idx(\Pi)$ together with a total order on $\mathcal{J}$;
examples are the order induced from $\idx(\Pi)$
and, for a cross-section $\Xi=(\Gamma_\parallel,\chi_\parallel)$ for $\ell$ (Definition~\ref{def:slabs})
and a set $\mathcal{J}(c)\supseteq\mathcal{J}$ of representatives of the modes at $\Pi$ modulo $\Gamma_\parallel$ (Definition~\ref{def:mode-translations}),
the order induced from $\mathcal{J}(c)$.
The modes $\zeta_i$, $i\in\mathcal{J}$, are those of Definition~\ref{def:modes};
by Lemma~\ref{lem:slab-bonds}\eqref{it:bonds-crossing} they lie in the subalgebra $\Gr(\mathbb{X}((c,c^{+})))$ of $\Gr(\mathbb{X})$,
and their reflections $\Theta_\Pi(\zeta_i)$ lie in $\Gr(\mathbb{X}((c^{-},c)))$,
where $\mathbb{X}(\mathbb{I})=\{x\in\mathbb{X}:\ \ell(x)\in\mathbb{I}\}$ as in Definition~\ref{def:slab-grassmann}.
The modes are read in one of two algebras.
The first is $\Gr(\mathbb{X})$, on which $\Theta_\Pi$ is the map of Definition~\ref{def:grassmann-reflections}.
The second is defined for a cross-section $\Xi=(\Gamma_\parallel,\chi_\parallel)$ for $\ell$ such that $i+g\notin\mathcal{J}$ for all $i\in\mathcal{J}$ and all $g\in\Gamma_\parallel\setminus\{0\}$:
it is the algebra $\Gr_\Xi((c^{-},c^{+}))$ of the cylinder $\Z^3/\Gamma_\parallel$ (Definition~\ref{def:slab-grassmann}).
In it, $\zeta_i$ and $\Theta_\Pi(\zeta_i)$ denote the images of these elements under the homomorphism
$[\,\cdot\,]_\Xi\colon\Gr(\mathbb{X}((c^{-},c^{+})))\to\Gr_\Xi((c^{-},c^{+}))$ that sends every generator to the generator with the same name,
and $\Theta_\Pi$ is the antilinear map of $\Gr_\Xi((c^{-},c^{+}))$ onto itself induced by the map $\Theta_\Pi$ of $\Gr(\mathbb{X}((c^{-},c^{+})))$, so that $\Theta_\Pi[G]_\Xi=[\Theta_\Pi G]_\Xi$
(Lemma~\ref{lem:slab-automorphisms}, applied to the interval $(c^{-},c^{+})=2c-(c^{-},c^{+})$).
In either algebra $\Theta_\Pi$ is antilinear and reverses the order of products, $\Theta_\Pi(GG')=\Theta_\Pi(G')\,\Theta_\Pi(G)$;
the elements and identities below are read in the algebra in which the modes are read.
For $I\subseteq\mathcal{J}$ the elements $\zeta_I$ and $\Theta_\Pi(\zeta_I)$ are the ordered products of \eqref{eq:setting-ordered-monomials}, formed with the order of $\mathcal{J}$,
and $\zeta_\emptyset=\Theta_\Pi(\zeta_\emptyset)=1$.

For $z=(z_i)_{i\in\mathcal{J}}\in\C^{\mathcal{J}}$ put $z_I:=\prod_{i\in I}z_i$, with $z_\emptyset:=1$,
and let $\diag(z_I)$ be the diagonal operator on $\mathcal{F}(\mathcal{J})$ with the entries $z_I$, $I\subseteq\mathcal{J}$.
The matrix entries $\Upsilon_{IJ}$ of an operator $\Upsilon$ on $\mathcal{F}(\mathcal{J})$ are defined by $\Upsilon e_J=\sum_{I\subseteq\mathcal{J}}\Upsilon_{IJ}\,e_I$, and
\[
\Ins^{\mathcal{J}}[\Upsilon]:=\sum_{I,J\subseteq\mathcal{J}}\Upsilon_{IJ}\cdot\Theta_\Pi(\zeta_J)\,\zeta_I ;
\]
this is the insertion of Definition~\ref{def:modes}, formed with the modes of $\mathcal{J}$ in the order of $\mathcal{J}$.
Let $\Gr[\mathcal{J}]$ be the subalgebra generated by the modes $\zeta_i$, $i\in\mathcal{J}$.
The \emph{canonical map} of $\mathcal{J}$ is the linear map $\zeta\colon\mathcal{F}(\mathcal{J})\to\Gr[\mathcal{J}]$ with $\zeta(e_I):=\zeta_I$ for $I\subseteq\mathcal{J}$.
For subsets $\mathcal{J}_0\subseteq\mathcal{J}_1$ of $\mathcal{J}$, both ordered as subsets of $\mathcal{J}$,
the localisation of Definition~\ref{def:modes} is denoted by $\iota_{\mathcal{J}_0\to\mathcal{J}_1}\colon\End\mathcal{F}(\mathcal{J}_0)\to\End\mathcal{F}(\mathcal{J}_1)$.

Let, in the second case, $\mathbb{T}=(\Gamma,\chi)$ be a torus with $\Xi(\mathbb{T},\ell)=\Xi$ (Definition~\ref{def:slabs});
then $\Gamma_\parallel\subseteq\Gamma$, and $\chi_\parallel$ is the restriction of $\chi$ to $\Gamma_\parallel$.
The \emph{reading} in $\Gr_{\mathbb{T}}$ is the homomorphism of algebras $\Gr_\Xi((c^{-},c^{+}))\to\Gr_{\mathbb{T}}$
that sends every generator to the generator of $\Gr_{\mathbb{T}}$ with the same name, as in Lemma~\ref{lem:matrix-rules}\eqref{it:matrix-closure};
it is well defined, the relations of $\Gr_\Xi((c^{-},c^{+}))$ being relations of $\Gr_{\mathbb{T}}$ (Definition~\ref{def:tori}).
An element of $\Gr_\Xi((c^{-},c^{+}))$, such as $\zeta_I$, $\Theta_\Pi(\zeta_I)$ or $\Ins^{\mathcal{J}}[\Upsilon]$, is \emph{read in} $\Gr_{\mathbb{T}}$ as its image under the reading.
\end{definition}

Let $\mathcal{J}$ be an index set of modes at a mirror $\Pi$, and let $I$ and $R$ be disjoint subsets of $\mathcal{J}$.
The product $\zeta_I\,\zeta_R$ is turned into $\zeta_{I\cup R}$ by moving each factor $\zeta_r$, $r\in R$, to the left past the factors $\zeta_i$ with $i\in I$ and $i>r$,
and every such step changes the sign, the modes being odd.
This is the first identity below.
The second holds because every pair $(i,r)\in I\times R$ satisfies exactly one of the inequalities $i>r$ and $r>i$:
\begin{equation}\label{eq:slab-sorting-sign}
\begin{aligned}
\zeta_{I\cup R}&=(-1)^{\#\{(i,r)\in I\times R:\ i>r\}}\cdot\zeta_I\,\zeta_R,\\
(-1)^{\#\{(i,r)\in I\times R:\ i>r\}}\cdot(-1)^{\#\{(r,i)\in R\times I:\ r>i\}}&=(-1)^{\#|I|\cdot\#|R|}.
\end{aligned}
\end{equation}
The sign in the first identity is the one that enters the localisation:
for $\mathcal{J}_0\subseteq\mathcal{J}_1\subseteq\mathcal{J}$ and $\Upsilon\in\End\mathcal{F}(\mathcal{J}_0)$, Definition~\ref{def:modes} gives
\begin{equation}\label{eq:slab-localised-entries}
\iota_{\mathcal{J}_0\to\mathcal{J}_1}(\Upsilon)_{I\cup R,\,J\cup R'}
=\delta_{RR'}\cdot(-1)^{\#\{(i,r)\in I\times R:\ i>r\}}\cdot(-1)^{\#\{(i,r)\in J\times R:\ i>r\}}\cdot\Upsilon_{IJ}
\end{equation}
for $I,J\subseteq\mathcal{J}_0$ and $R,R'\subseteq\mathcal{J}_1\setminus\mathcal{J}_0$.
An operator $\Upsilon$ on $\mathcal{F}(\mathcal{J})$ is even if it commutes with the parity $U_{\mathrm{f}}$, $U_{\mathrm{f}}e_I=(-1)^{\#|I|}e_I$ (Definition~\ref{def:modes});
equivalently, $\Upsilon_{IJ}=0$ unless $\#|I|-\#|J|$ is even.
Since $\#|I\cup R|-\#|J\cup R|=\#|I|-\#|J|$, the formula \eqref{eq:slab-localised-entries} shows that the localisation maps even operators to even operators.

The modes $\zeta_i$, $i\in\mathcal{J}$, are linearly independent elements of degree one.
In $\Gr(\mathbb{X})$ this is part of Definition~\ref{def:modes}.
In $\Gr_\Xi((c^{-},c^{+}))$ one has $\zeta_{i+g}=\chi_\parallel(g)\cdot\zeta_i$ for $g\in\Gamma_\parallel$, and the elements of $\mathcal{J}$ lie in distinct classes modulo $\Gamma_\parallel$;
so, up to signs, the $\zeta_i$, $i\in\mathcal{J}$, are distinct members of the family of the modes of a set $\mathcal{J}(c)$ of representatives,
which is linearly independent by Definition~\ref{def:mode-translations}.
Hence the ordered products $\zeta_I$, $I\subseteq\mathcal{J}$, form a basis of $\Gr[\mathcal{J}]$,
and the canonical map $\zeta$ is a linear bijection of $\mathcal{F}(\mathcal{J})$ onto $\Gr[\mathcal{J}]$.
For $\mathcal{J}_0\subseteq\mathcal{J}$, ordered as a subset of $\mathcal{J}$,
the space $\mathcal{F}(\mathcal{J}_0)$ is regarded as the subspace of $\mathcal{F}(\mathcal{J})$ spanned by the vectors $e_I$, $I\subseteq\mathcal{J}_0$;
the canonical map of $\mathcal{J}_0$ is then the restriction of $\zeta$ to $\mathcal{F}(\mathcal{J}_0)$,
and it maps $\mathcal{F}(\mathcal{J}_0)$ onto $\Gr[\mathcal{J}_0]$.

Suppose that the modes are read in $\Gr_\Xi((c^{-},c^{+}))$, and let $\mathbb{T}$ be a torus with $\Xi(\mathbb{T},\ell)=\Xi$.
The reading in $\Gr_{\mathbb{T}}$ is a homomorphism of algebras;
every identity between elements of $\Gr_\Xi((c^{-},c^{+}))$ proved below therefore holds between their images in $\Gr_{\mathbb{T}}$.
For $G\in\Gr(\mathbb{X}((c^{-},c^{+})))$ the reading of $[G]_\Xi$ is the class $[G]_{\mathbb{T}}$ of $G$ under the homomorphism $[\,\cdot\,]_{\mathbb{T}}\colon\Gr(\mathbb{X})\to\Gr_{\mathbb{T}}$ that sends every generator to the generator with the same name:
the map $G\mapsto[G]_{\mathbb{T}}$ and the composition of $[\,\cdot\,]_\Xi$ with the reading are homomorphisms of algebras that agree on the generators.
Also $\Theta_\Pi G\in\Gr(\mathbb{X}((c^{-},c^{+})))$ (Lemma~\ref{lem:slab-automorphisms}).
If moreover $\theta_\Pi\mathbb{T}=\mathbb{T}$, the reading of $\Theta_\Pi[G]_\Xi=[\Theta_\Pi G]_\Xi$ is therefore $[\Theta_\Pi G]_{\mathbb{T}}=\Theta_\Pi[G]_{\mathbb{T}}$,
where $\Theta_\Pi$ is the antilinear anti-automorphism of $\Gr_{\mathbb{T}}$ of Lemma~\ref{lem:reflected-algebras};
in this case the reading commutes with $\Theta_\Pi$.

\begin{proposition}[Localisation]\label{prop:localisation}
Let $\mathcal{J}$ be an index set of modes at a mirror $\Pi$,
let $\mathcal{J}'\subseteq\mathcal{J}$ and $\mathcal{J}'':=\mathcal{J}\setminus\mathcal{J}'$,
and put $\iota:=\iota_{\mathcal{J}'\to\mathcal{J}}$.
\begin{enumerate}
\item\label{it:insertion-split}
Let $I,J\subseteq\mathcal{J}'$ be such that $\#|I|-\#|J|$ is even, and let $R\subseteq\mathcal{J}''$. Then
\[
\Theta_\Pi(\zeta_{J\cup R})\,\zeta_{I\cup R}
=(-1)^{\#\{(i,r)\in I\times R:\ i>r\}}\cdot(-1)^{\#\{(i,r)\in J\times R:\ i>r\}}\cdot\Theta_\Pi(\zeta_J)\,\zeta_I\cdot\Theta_\Pi(\zeta_R)\,\zeta_R .
\]
\item\label{it:localisation-homomorphism}
The map $\iota\colon\End\mathcal{F}(\mathcal{J}')\to\End\mathcal{F}(\mathcal{J})$ is a unital injective $*$-homomorphism; in particular, it is isometric.
Let $\Upsilon\in\End\mathcal{F}(\mathcal{J}')$, and let $\Phi\in\mathcal{F}(\mathcal{J})$ be such that $\zeta(\Phi)=G_1\,G_2$ with $G_1\in\Gr[\mathcal{J}']$ and $G_2\in\Gr[\mathcal{J}'']$. Then
\[
\zeta\left(\iota(\Upsilon)\,\Phi\right)=\left(\Upsilon G_1\right)\,G_2,
\qquad\text{where}\quad
\Upsilon G_1:=\zeta\left(\Upsilon\,\zeta^{-1}(G_1)\right).
\]
Consequently the operator $\zeta\circ\iota(\Upsilon)\circ\zeta^{-1}$ of $\Gr[\mathcal{J}]$ is determined by the operator $\zeta\circ\Upsilon\circ\zeta^{-1}$ of $\Gr[\mathcal{J}']$, whatever the total order of $\mathcal{J}$.
Moreover, the localisation is transitive:
if $\mathcal{J}_0\subseteq\mathcal{J}'$, then $\iota_{\mathcal{J}'\to\mathcal{J}}\circ\iota_{\mathcal{J}_0\to\mathcal{J}'}=\iota_{\mathcal{J}_0\to\mathcal{J}}$.
\item\label{it:localisation-commute}
Let $\mathcal{J}_1$ and $\mathcal{J}_2$ be disjoint subsets of $\mathcal{J}$,
and let $\Upsilon_1\in\End\mathcal{F}(\mathcal{J}_1)$ and $\Upsilon_2\in\End\mathcal{F}(\mathcal{J}_2)$, at least one of them even.
Then $\iota_{\mathcal{J}_1\to\mathcal{J}}(\Upsilon_1)$ and $\iota_{\mathcal{J}_2\to\mathcal{J}}(\Upsilon_2)$ commute.
\item\label{it:localised-insertion}
Let $\Upsilon\in\End\mathcal{F}(\mathcal{J}')$ be even, and let $z\in\C^{\mathcal{J}}$. Then
\[
\Ins^{\mathcal{J}}\left[\iota(\Upsilon)\,\diag(z_I)\right]
=\Ins^{\mathcal{J}'}\left[\Upsilon\,\diag(z_I)_{I\subseteq\mathcal{J}'}\right]\cdot\prod_{i\in\mathcal{J}''}\left(1+z_i\cdot\Theta_\Pi(\zeta_i)\,\zeta_i\right).
\]
Let $\omega_i\geq0$ and $u_i$ be the crossing coefficients and the phases of the modes (Definition~\ref{def:modes}),
and let $\omega^{\frac12}$ and $\diag(u_I)$ be the diagonal operators with the entries $\omega_I^{\frac12}$ and $u_I$, on $\mathcal{F}(\mathcal{J})$ and on $\mathcal{F}(\mathcal{J}')$. Then
\[
\Ins^{\mathcal{J}}\left[\omega^{\frac12}\,\iota(\Upsilon)\,\diag(u_I)\,\omega^{\frac12}\right]
=\Ins^{\mathcal{J}'}\left[\omega^{\frac12}\,\Upsilon\,\diag(u_I)\,\omega^{\frac12}\right]\cdot\prod_{i\in\mathcal{J}''}\left(1+\omega_i\cdot u_i\cdot\Theta_\Pi(\zeta_i)\,\zeta_i\right).
\]
\item\label{it:several-insertions}
Let $s\geq1$, let $\mathcal{J}_1,\dots,\mathcal{J}_s$ be pairwise disjoint subsets of $\mathcal{J}$,
and let $\Upsilon_k\in\End\mathcal{F}(\mathcal{J}_k)$ be even, $k=1,\dots,s$.
Then, with $\iota_k:=\iota_{\mathcal{J}_k\to\mathcal{J}}$ and with $\omega^{\frac12}$ and $\diag(u_I)$ as in~\eqref{it:localised-insertion},
\[
\begin{aligned}
&\Ins^{\mathcal{J}}\left[\omega^{\frac12}\,\iota_1(\Upsilon_1)\cdots\iota_s(\Upsilon_s)\,\diag(u_I)\,\omega^{\frac12}\right]\\
&\qquad=\prod_{k=1}^{s}\Ins^{\mathcal{J}_k}\left[\omega^{\frac12}\,\Upsilon_k\,\diag(u_I)\,\omega^{\frac12}\right]\cdot\prod_{i\in\mathcal{J}\setminus(\mathcal{J}_1\cup\dots\cup\mathcal{J}_s)}\left(1+\omega_i\cdot u_i\cdot\Theta_\Pi(\zeta_i)\,\zeta_i\right).
\end{aligned}
\]
\item\label{it:mode-relabelling}
(Change of order and of representatives.)
Let $\tilde{\mathcal{J}}$ be a second index set of modes at $\Pi$, read in the same algebra as $\mathcal{J}$,
and let $r\colon\mathcal{J}\to\tilde{\mathcal{J}}$ be a bijection such that $\zeta_i=\pm\zeta_{r(i)}$ for every $i\in\mathcal{J}$.
Let $\tilde\zeta_{I'}$ and $\Theta_\Pi(\tilde\zeta_{I'})$, $I'\subseteq\tilde{\mathcal{J}}$, be the ordered products formed with the order of $\tilde{\mathcal{J}}$,
and let $\Ins^{\tilde{\mathcal{J}}}$, the canonical map $\tilde\zeta$ of $\tilde{\mathcal{J}}$ and the localisation $\tilde\iota:=\iota_{r(\mathcal{J}')\to\tilde{\mathcal{J}}}$ be formed with $\tilde{\mathcal{J}}$ as in Definition~\ref{def:index-sets}.
For $I\subseteq\mathcal{J}$ let $s(I)\in\{1,-1\}$ be the sign with $\zeta_I=s(I)\cdot\tilde\zeta_{r(I)}$,
let $U\colon\mathcal{F}(\mathcal{J})\to\mathcal{F}(\tilde{\mathcal{J}})$ be the linear map with $Ue_I:=s(I)\cdot e_{r(I)}$, $I\subseteq\mathcal{J}$,
and let $U'\colon\mathcal{F}(\mathcal{J}')\to\mathcal{F}(r(\mathcal{J}'))$ be the linear map with $U'e_I:=s(I)\cdot e_{r(I)}$, $I\subseteq\mathcal{J}'$.
Then $U$ is unitary, and
\[
\begin{aligned}
&\Ins^{\mathcal{J}}[\Upsilon]=\Ins^{\tilde{\mathcal{J}}}\left[U\Upsilon U^{*}\right],
\qquad
U\,\iota(\Upsilon')\,U^{*}=\tilde\iota\left(U'\Upsilon'U'^{*}\right),\\
&U\diag(z_I)\,U^{*}=\diag\left(z_{r^{-1}(I)}\right)_{I\subseteq\tilde{\mathcal{J}}}
\end{aligned}
\]
for all $\Upsilon\in\End\mathcal{F}(\mathcal{J})$, $\Upsilon'\in\End\mathcal{F}(\mathcal{J}')$ and $z\in\C^{\mathcal{J}}$.
\end{enumerate}
\end{proposition}

\begin{proof}
For disjoint subsets $I$ and $R$ of $\mathcal{J}$ write
\[
\#(I>R):=\#\{(i,r)\in I\times R:\ i>r\}.
\]
This number is additive in each of its two arguments,
and in this notation \eqref{eq:slab-sorting-sign} and \eqref{eq:slab-localised-entries} read
$\zeta_{I\cup R}=(-1)^{\#(I>R)}\cdot\zeta_I\,\zeta_R$,
$\#(I>R)+\#(R>I)\equiv\#|I|\cdot\#|R|\pmod 2$,
and $\iota_{\mathcal{J}_0\to\mathcal{J}_1}(\Upsilon)_{I\cup R,\,J\cup R'}=\delta_{RR'}\cdot(-1)^{\#(I>R)+\#(J>R)}\cdot\Upsilon_{IJ}$.

In the situation of~\eqref{it:insertion-split}, \eqref{eq:slab-sorting-sign} gives $\zeta_{J\cup R}=(-1)^{\#(J>R)}\cdot\zeta_J\,\zeta_R$ and $\zeta_{I\cup R}=(-1)^{\#(I>R)}\cdot\zeta_I\,\zeta_R$.
The reflection $\Theta_\Pi$ is antilinear and reverses the order of products (Definition~\ref{def:index-sets}), and the sign is real;
hence $\Theta_\Pi(\zeta_{J\cup R})=(-1)^{\#(J>R)}\cdot\Theta_\Pi(\zeta_R)\,\Theta_\Pi(\zeta_J)$, and
\[
\Theta_\Pi(\zeta_{J\cup R})\,\zeta_{I\cup R}=(-1)^{\#(I>R)+\#(J>R)}\cdot\Theta_\Pi(\zeta_R)\cdot\Theta_\Pi(\zeta_J)\,\zeta_I\cdot\zeta_R .
\]
The modes and their reflections are odd, so $\Theta_\Pi(\zeta_J)\,\zeta_I$ has the parity of $\#|I|+\#|J|$.
By hypothesis it is even; it therefore commutes with $\Theta_\Pi(\zeta_R)$ and may be moved to the left of it, which gives~\eqref{it:insertion-split}.

For~\eqref{it:localisation-homomorphism}, let $U_{\otimes}\colon\mathcal{F}(\mathcal{J}')\otimes\mathcal{F}(\mathcal{J}'')\to\mathcal{F}(\mathcal{J})$ be the linear map with
\[
U_{\otimes}(e_I\otimes e_R):=(-1)^{\#(I>R)}\cdot e_{I\cup R},\qquad I\subseteq\mathcal{J}',\ R\subseteq\mathcal{J}''.
\]
A subset of $\mathcal{J}$ is determined by its intersections with $\mathcal{J}'$ and with $\mathcal{J}''$, and every pair of subsets of $\mathcal{J}'$ and of $\mathcal{J}''$ arises in this way;
so $U_{\otimes}$ maps an orthonormal basis onto an orthonormal basis up to signs, and it is unitary, with $U_{\otimes}^{*}e_{I\cup R}=(-1)^{\#(I>R)}\cdot e_I\otimes e_R$.
For $\Upsilon\in\End\mathcal{F}(\mathcal{J}')$, $J\subseteq\mathcal{J}'$ and $R'\subseteq\mathcal{J}''$ it follows that
\[
U_{\otimes}(\Upsilon\otimes1)\,U_{\otimes}^{*}\,e_{J\cup R'}=(-1)^{\#(J>R')}\cdot U_{\otimes}\left(\Upsilon e_J\otimes e_{R'}\right)=\sum_{I\subseteq\mathcal{J}'}(-1)^{\#(I>R')+\#(J>R')}\cdot\Upsilon_{IJ}\cdot e_{I\cup R'} .
\]
These are the entries \eqref{eq:slab-localised-entries}; thus $\iota(\Upsilon)=U_{\otimes}(\Upsilon\otimes1)\,U_{\otimes}^{*}$.
The map $\Upsilon\mapsto\Upsilon\otimes1$ is a unital injective $*$-homomorphism with $\|\Upsilon\otimes1\|=\|\Upsilon\|$,
and conjugation by a unitary operator preserves these properties;
so $\iota$ is a unital injective $*$-homomorphism, and it is isometric.
Next, \eqref{eq:slab-sorting-sign} gives $\zeta\left(U_{\otimes}(e_I\otimes e_R)\right)=(-1)^{\#(I>R)}\cdot\zeta_{I\cup R}=\zeta_I\,\zeta_R$,
and by bilinearity $\zeta\left(U_{\otimes}(\Phi'\otimes\Phi'')\right)=\zeta(\Phi')\,\zeta(\Phi'')$ for $\Phi'\in\mathcal{F}(\mathcal{J}')$ and $\Phi''\in\mathcal{F}(\mathcal{J}'')$.
Let $\Phi$, $G_1$, $G_2$ be as in~\eqref{it:localisation-homomorphism}, and put $\Phi':=\zeta^{-1}(G_1)$ and $\Phi'':=\zeta^{-1}(G_2)$.
Then $\zeta\left(U_{\otimes}(\Phi'\otimes\Phi'')\right)=G_1\,G_2=\zeta(\Phi)$, so that $\Phi=U_{\otimes}(\Phi'\otimes\Phi'')$, $\zeta$ being injective, and
\[
\zeta\left(\iota(\Upsilon)\,\Phi\right)=\zeta\left(U_{\otimes}\left(\Upsilon\Phi'\otimes\Phi''\right)\right)=\zeta(\Upsilon\Phi')\,\zeta(\Phi'')=\left(\Upsilon G_1\right)\,G_2 .
\]
The products $G_1\,G_2$ span $\Gr[\mathcal{J}]$, since $\zeta_I\,\zeta_R=\pm\zeta_{I\cup R}$.
So this formula determines $\zeta\circ\iota(\Upsilon)\circ\zeta^{-1}$, and it involves only the operator $\zeta\circ\Upsilon\circ\zeta^{-1}$ and the multiplication of $\Gr[\mathcal{J}]$.
For the transitivity let $\mathcal{J}_0\subseteq\mathcal{J}'$ and $\Upsilon\in\End\mathcal{F}(\mathcal{J}_0)$,
and let $G_0\in\Gr[\mathcal{J}_0]$, $G_1\in\Gr[\mathcal{J}'\setminus\mathcal{J}_0]$ and $G_2\in\Gr[\mathcal{J}'']$.
Applied twice, the formula shows that $\zeta\circ\iota_{\mathcal{J}'\to\mathcal{J}}\left(\iota_{\mathcal{J}_0\to\mathcal{J}'}(\Upsilon)\right)\circ\zeta^{-1}$ maps $(G_0\,G_1)\,G_2$ to $\left((\Upsilon G_0)\,G_1\right)\,G_2$;
applied once, it shows that $\zeta\circ\iota_{\mathcal{J}_0\to\mathcal{J}}(\Upsilon)\circ\zeta^{-1}$ maps $G_0\,(G_1\,G_2)$ to $(\Upsilon G_0)\,(G_1\,G_2)$.
The two images coincide, and the products $G_0\,G_1\,G_2$ span $\Gr[\mathcal{J}]$.

The commutation~\eqref{it:localisation-commute} is a computation of signs.
Let $I_1,J_1\subseteq\mathcal{J}_1$, $I_2,J_2\subseteq\mathcal{J}_2$ and $R,R'\subseteq\mathcal{J}\setminus(\mathcal{J}_1\cup\mathcal{J}_2)$,
and write $\iota_k:=\iota_{\mathcal{J}_k\to\mathcal{J}}$.
By \eqref{eq:slab-localised-entries}, an entry of $\iota_1(\Upsilon_1)$ vanishes unless its row index and its column index have the same intersection with $\mathcal{J}\setminus\mathcal{J}_1$,
and an entry of $\iota_2(\Upsilon_2)$ vanishes unless they have the same intersection with $\mathcal{J}\setminus\mathcal{J}_2$.
Therefore the entry of $\iota_1(\Upsilon_1)\,\iota_2(\Upsilon_2)$ in the row $I_1\cup I_2\cup R$ and the column $J_1\cup J_2\cup R'$ vanishes unless $R=R'$,
and then only the intermediate index $J_1\cup I_2\cup R$ contributes to it; this entry is
\[
(-1)^{\#(I_1>I_2\cup R)+\#(J_1>I_2\cup R)+\#(I_2>J_1\cup R)+\#(J_2>J_1\cup R)}\cdot(\Upsilon_1)_{I_1J_1}\cdot(\Upsilon_2)_{I_2J_2}.
\]
In the same way, in $\iota_2(\Upsilon_2)\,\iota_1(\Upsilon_1)$ only the intermediate index $I_1\cup J_2\cup R$ contributes, and the entry is
\[
(-1)^{\#(I_2>I_1\cup R)+\#(J_2>I_1\cup R)+\#(I_1>J_2\cup R)+\#(J_1>J_2\cup R)}\cdot(\Upsilon_1)_{I_1J_1}\cdot(\Upsilon_2)_{I_2J_2}.
\]
Split every count against a union into the counts against its two parts.
The counts against $R$ are the same in both exponents, and by \eqref{eq:slab-sorting-sign} the sum of the two exponents is congruent modulo $2$ to
\[
\begin{aligned}
&\#(I_1>I_2)+\#(I_2>I_1)+\#(J_1>I_2)+\#(I_2>J_1)\\
&\qquad+\#(J_2>J_1)+\#(J_1>J_2)+\#(J_2>I_1)+\#(I_1>J_2)\\
&\qquad\equiv\#|I_1|\cdot\#|I_2|+\#|J_1|\cdot\#|I_2|+\#|J_1|\cdot\#|J_2|+\#|I_1|\cdot\#|J_2|\\
&\qquad=\left(\#|I_1|+\#|J_1|\right)\cdot\left(\#|I_2|+\#|J_2|\right).
\end{aligned}
\]
If $\Upsilon_1$ is even, then $(\Upsilon_1)_{I_1J_1}=0$ unless $\#|I_1|+\#|J_1|$ is even; if $\Upsilon_2$ is even, the same holds for $I_2$ and $J_2$.
In either case the two entries coincide.

Item~\eqref{it:localised-insertion} rests on~\eqref{it:insertion-split}.
Let $I,J\subseteq\mathcal{J}'$ and $R,R'\subseteq\mathcal{J}''$.
As $z_{J\cup R}=z_J\cdot z_R$, the entry of $\iota(\Upsilon)\diag(z_I)$ in the row $I\cup R$ and the column $J\cup R'$ is
$\delta_{RR'}\cdot(-1)^{\#(I>R)+\#(J>R)}\cdot\Upsilon_{IJ}\cdot z_J\cdot z_R$.
Since $\Upsilon$ is even, only terms with $\#|I|-\#|J|$ even occur, and~\eqref{it:insertion-split} applies to each of them;
the square of a sign being $1$,
\[
\Ins^{\mathcal{J}}\left[\iota(\Upsilon)\,\diag(z_I)\right]
=\sum_{I,J\subseteq\mathcal{J}'}\Upsilon_{IJ}\cdot z_J\cdot\Theta_\Pi(\zeta_J)\,\zeta_I\cdot\sum_{R\subseteq\mathcal{J}''}z_R\cdot\Theta_\Pi(\zeta_R)\,\zeta_R .
\]
The first factor is $\Ins^{\mathcal{J}'}\left[\Upsilon\,\diag(z_I)_{I\subseteq\mathcal{J}'}\right]$.
In the second, $\Theta_\Pi(\zeta_R)\,\zeta_R=\prod_{i\in R}\Theta_\Pi(\zeta_i)\,\zeta_i$, the factors being even and commuting with one another.
For the order induced from $\idx(\Pi)$ and in $\Gr(\mathbb{X})$ this is Proposition~\ref{prop:crossing}\eqref{it:crossing-exponential}, applied to the modes of the bonds that carry the elements of $R$.
Ordering $R$ as a subset of $\mathcal{J}$ instead multiplies $\zeta_R$ by the sign of a permutation and, $\Theta_\Pi$ being antilinear, $\Theta_\Pi(\zeta_R)$ by the same sign, which leaves the product unchanged;
and in $\Gr_\Xi((c^{-},c^{+}))$ the identity is the image under $[\,\cdot\,]_\Xi$ of the identity in $\Gr(\mathbb{X})$, because $[\,\cdot\,]_\Xi$ is a homomorphism with $\Theta_\Pi[G]_\Xi=[\Theta_\Pi G]_\Xi$ (Definition~\ref{def:index-sets}).
So the sum over $R$ is the expansion of $\prod_{i\in\mathcal{J}''}\left(1+z_i\cdot\Theta_\Pi(\zeta_i)\,\zeta_i\right)$.
This is the first identity.
The numbers $\omega_I^{\frac12}$ and $u_I$ are multiplicative in $I$ as well;
hence the entry of $\omega^{\frac12}\,\iota(\Upsilon)\,\diag(u_I)\,\omega^{\frac12}$ in the row $I\cup R$ and the column $J\cup R'$ is
\[
\delta_{RR'}\cdot(-1)^{\#(I>R)+\#(J>R)}\cdot\left(\omega_I^{\frac12}\cdot\Upsilon_{IJ}\cdot u_J\cdot\omega_J^{\frac12}\right)\cdot\omega_R\cdot u_R ,
\]
and by~\eqref{it:insertion-split} the left side of the second identity equals
\[
\sum_{I,J\subseteq\mathcal{J}'}\omega_I^{\frac12}\cdot\Upsilon_{IJ}\cdot u_J\cdot\omega_J^{\frac12}\cdot\Theta_\Pi(\zeta_J)\,\zeta_I\cdot\sum_{R\subseteq\mathcal{J}''}\omega_R\cdot u_R\cdot\Theta_\Pi(\zeta_R)\,\zeta_R .
\]
The same expansion of the product turns this into the right side.

Several insertions, as in~\eqref{it:several-insertions}, are first gathered into one.
Put $\mathcal{J}_0:=\mathcal{J}_1\cup\dots\cup\mathcal{J}_s$, an index set of modes with the order induced from $\mathcal{J}$.
By the transitivity in~\eqref{it:localisation-homomorphism}, $\iota_k=\iota_{\mathcal{J}_0\to\mathcal{J}}\circ\iota_{\mathcal{J}_k\to\mathcal{J}_0}$, and $\iota_{\mathcal{J}_0\to\mathcal{J}}$ is multiplicative; hence
\[
\iota_1(\Upsilon_1)\cdots\iota_s(\Upsilon_s)=\iota_{\mathcal{J}_0\to\mathcal{J}}(\Upsilon_0),
\qquad
\Upsilon_0:=\iota_{\mathcal{J}_1\to\mathcal{J}_0}(\Upsilon_1)\cdots\iota_{\mathcal{J}_s\to\mathcal{J}_0}(\Upsilon_s),
\]
and $\Upsilon_0$ is even, as a product of localisations of even operators.
The second identity of~\eqref{it:localised-insertion}, with $\mathcal{J}_0$ in place of $\mathcal{J}'$ and $\Upsilon_0$ in place of $\Upsilon$,
produces the factor $\prod_{i\in\mathcal{J}\setminus\mathcal{J}_0}\left(1+\omega_i\cdot u_i\cdot\Theta_\Pi(\zeta_i)\,\zeta_i\right)$
and leaves $\Ins^{\mathcal{J}_0}\left[\omega^{\frac12}\,\Upsilon_0\,\diag(u_I)\,\omega^{\frac12}\right]$.
It therefore suffices to prove~\eqref{it:several-insertions} when $\mathcal{J}$ is the union of the sets $\mathcal{J}_k$.
For $s=1$ there is then nothing to prove.
For $s\geq3$, replace $\mathcal{J}_2,\dots,\mathcal{J}_s$ by their union,
and $\Upsilon_2,\dots,\Upsilon_s$ by the product of their localisations into this union, an even operator;
by transitivity this does not change $\iota_2(\Upsilon_2)\cdots\iota_s(\Upsilon_s)$.
The case of two sets, followed by~\eqref{it:several-insertions} for the $s-1$ sets $\mathcal{J}_2,\dots,\mathcal{J}_s$ in their union, then gives the assertion by induction on $s$.

There remains the case $s=2$ with $\mathcal{J}=\mathcal{J}_1\cup\mathcal{J}_2$.
Let $I_1,J_1\subseteq\mathcal{J}_1$ and $I_2,J_2\subseteq\mathcal{J}_2$.
As in the proof of~\eqref{it:localisation-commute}, now without a third part, the entry of $\iota_1(\Upsilon_1)\,\iota_2(\Upsilon_2)$ in the row $I_1\cup I_2$ and the column $J_1\cup J_2$ is
\[
(-1)^{\#(I_1>I_2)+\#(J_1>I_2)+\#(I_2>J_1)+\#(J_2>J_1)}\cdot(\Upsilon_1)_{I_1J_1}\cdot(\Upsilon_2)_{I_2J_2}.
\]
If $(\Upsilon_1)_{I_1J_1}\neq0$, then $\#|I_1|+\#|J_1|$ is even and $\Theta_\Pi(\zeta_{J_1})\,\zeta_{I_1}$ is even.
By \eqref{eq:slab-sorting-sign} and the rules for $\Theta_\Pi$ (Definition~\ref{def:index-sets}),
\[
\begin{aligned}
\Theta_\Pi(\zeta_{J_1\cup J_2})\,\zeta_{I_1\cup I_2}
&=(-1)^{\#(J_1>J_2)+\#(I_1>I_2)}\cdot\Theta_\Pi(\zeta_{J_2})\cdot\Theta_\Pi(\zeta_{J_1})\,\zeta_{I_1}\cdot\zeta_{I_2}\\
&=(-1)^{\#(J_1>J_2)+\#(I_1>I_2)}\cdot\Theta_\Pi(\zeta_{J_1})\,\zeta_{I_1}\cdot\Theta_\Pi(\zeta_{J_2})\,\zeta_{I_2},
\end{aligned}
\]
the even middle factor commuting with $\Theta_\Pi(\zeta_{J_2})$.
The product of the two signs is
\[
(-1)^{\#(J_1>I_2)+\#(I_2>J_1)+\#(J_2>J_1)+\#(J_1>J_2)}=(-1)^{\#|J_1|\cdot\left(\#|I_2|+\#|J_2|\right)},
\]
and it equals $1$ whenever $(\Upsilon_2)_{I_2J_2}\neq0$, because $\Upsilon_2$ is even.
The diagonal factors split as well:
$\omega_{I_1\cup I_2}^{\frac12}=\omega_{I_1}^{\frac12}\cdot\omega_{I_2}^{\frac12}$,
$u_{J_1\cup J_2}=u_{J_1}\cdot u_{J_2}$
and $\omega_{J_1\cup J_2}^{\frac12}=\omega_{J_1}^{\frac12}\cdot\omega_{J_2}^{\frac12}$.
Summing over $I_1,J_1,I_2,J_2$ yields
\[
\begin{aligned}
&\Ins^{\mathcal{J}}\left[\omega^{\frac12}\,\iota_1(\Upsilon_1)\,\iota_2(\Upsilon_2)\,\diag(u_I)\,\omega^{\frac12}\right]\\
&\qquad=\Ins^{\mathcal{J}_1}\left[\omega^{\frac12}\,\Upsilon_1\,\diag(u_I)\,\omega^{\frac12}\right]\cdot\Ins^{\mathcal{J}_2}\left[\omega^{\frac12}\,\Upsilon_2\,\diag(u_I)\,\omega^{\frac12}\right],
\end{aligned}
\]
which is~\eqref{it:several-insertions} for $s=2$.
The insertions of even operators are even elements, so the order of the factors in~\eqref{it:several-insertions} is immaterial.

Turning to~\eqref{it:mode-relabelling}, let $I\subseteq\mathcal{J}$ have the elements $i_1<\dots<i_d$.
Then $\zeta_I=\zeta_{i_1}\cdots\zeta_{i_d}$ is, up to a sign, the product of $\zeta_{r(i_1)},\dots,\zeta_{r(i_d)}$ in this order,
and bringing these odd factors into the order of $\tilde{\mathcal{J}}$ changes the product by a further sign; thus $\zeta_I=\pm\tilde\zeta_{r(I)}$.
The modes $\zeta_j$, $j\in\tilde{\mathcal{J}}$, are linearly independent, $\tilde{\mathcal{J}}$ being an index set of modes read in the same algebra as $\mathcal{J}$;
so $\tilde\zeta_{r(I)}\neq0$ and the sign $s(I)$ is well defined.
Moreover $\Gr[\tilde{\mathcal{J}}]=\Gr[\mathcal{J}]$, every generator of either algebra being $\pm$ a generator of the other.
The bijection $r$ induces a bijection $I\mapsto r(I)$ of the subsets of $\mathcal{J}$ onto those of $\tilde{\mathcal{J}}$.
So $U$ maps the orthonormal basis $(e_I)$ onto the orthonormal basis of $\mathcal{F}(\tilde{\mathcal{J}})$ up to signs; it is unitary, with $U^{*}e_{r(I)}=s(I)\cdot e_I$, and $U'$ is unitary for the same reason.
As $s(J)$ is real and $\Theta_\Pi$ is antilinear, $\Theta_\Pi(\zeta_J)=s(J)\cdot\Theta_\Pi(\tilde\zeta_{r(J)})$, and therefore
\[
\Theta_\Pi(\zeta_J)\,\zeta_I=s(I)\cdot s(J)\cdot\Theta_\Pi(\tilde\zeta_{r(J)})\,\tilde\zeta_{r(I)},
\qquad
\left(U\Upsilon U^{*}\right)_{r(I)\,r(J)}=s(I)\cdot s(J)\cdot\Upsilon_{IJ}.
\]
In the sum that defines $\Ins^{\tilde{\mathcal{J}}}\left[U\Upsilon U^{*}\right]$, every pair of subsets of $\tilde{\mathcal{J}}$ is $(r(I),r(J))$ for exactly one pair $I,J\subseteq\mathcal{J}$.
Inserting the second identity for the entry and then the first for the element of the Grassmann algebra, one finds
\[
\begin{aligned}
\Ins^{\tilde{\mathcal{J}}}\left[U\Upsilon U^{*}\right]
&=\sum_{I,J\subseteq\mathcal{J}}\left(U\Upsilon U^{*}\right)_{r(I)\,r(J)}\cdot\Theta_\Pi(\tilde\zeta_{r(J)})\,\tilde\zeta_{r(I)}\\
&=\sum_{I,J\subseteq\mathcal{J}}\Upsilon_{IJ}\cdot s(I)\cdot s(J)\cdot\Theta_\Pi(\tilde\zeta_{r(J)})\,\tilde\zeta_{r(I)}\\
&=\sum_{I,J\subseteq\mathcal{J}}\Upsilon_{IJ}\cdot\Theta_\Pi(\zeta_J)\,\zeta_I
=\Ins^{\mathcal{J}}[\Upsilon].
\end{aligned}
\]
The canonical maps satisfy $\tilde\zeta\circ U=\zeta$, because $\tilde\zeta(Ue_I)=s(I)\cdot\tilde\zeta_{r(I)}=\zeta_I$; likewise $\tilde\zeta\circ U'=\zeta$ on $\mathcal{F}(\mathcal{J}')$.
Let $G_1\in\Gr[\mathcal{J}']=\Gr[r(\mathcal{J}')]$ and $G_2\in\Gr[\mathcal{J}'']=\Gr[r(\mathcal{J}'')]$, and note that $r(\mathcal{J}'')=\tilde{\mathcal{J}}\setminus r(\mathcal{J}')$.
By~\eqref{it:localisation-homomorphism} for the index set $\tilde{\mathcal{J}}$ and its subset $r(\mathcal{J}')$, and then for $\mathcal{J}$ and $\mathcal{J}'$,
\[
\begin{aligned}
\tilde\zeta\left(\tilde\iota\left(U'\Upsilon'U'^{*}\right)\tilde\zeta^{-1}(G_1\,G_2)\right)
&=\tilde\zeta\left(U'\Upsilon'U'^{*}\,\tilde\zeta^{-1}(G_1)\right)\,G_2
=\zeta\left(\Upsilon'\zeta^{-1}(G_1)\right)\,G_2\\
&=\zeta\left(\iota(\Upsilon')\,\zeta^{-1}(G_1\,G_2)\right)
=\tilde\zeta\left(U\iota(\Upsilon')\,U^{*}\,\tilde\zeta^{-1}(G_1\,G_2)\right),
\end{aligned}
\]
where $\tilde\zeta^{-1}=U'\zeta^{-1}$ on $\Gr[\mathcal{J}']$ and $\tilde\zeta^{-1}=U\zeta^{-1}$ on $\Gr[\mathcal{J}]$.
The products $G_1\,G_2$ span $\Gr[\mathcal{J}]$ and $\tilde\zeta$ is bijective; hence $\tilde\iota\left(U'\Upsilon'U'^{*}\right)=U\iota(\Upsilon')\,U^{*}$.
Finally, $U\diag(z_I)\,U^{*}e_{r(I)}=s(I)\cdot U\diag(z_I)\,e_I=s(I)^2\cdot z_I\cdot e_{r(I)}=z_I\cdot e_{r(I)}$ for every $I\subseteq\mathcal{J}$, which is the last identity.
\end{proof}

\begin{remark}[Localisation on a cross-section]\label{rem:cross-section-localisation}
Let $\Pi$ be the link plane $\{x_1=j+\tfrac12\}$, $j\in\Z$,
let $\Xi$ be a rectangular cross-section, with the lattice $\Gamma_\parallel=n_2\Z e_2+n_3\Z e_3$, the character $\chi_\parallel$ and the slice torus $\Lambda_{\mathrm{s}}$ (Definition~\ref{def:slices}),
and let $\Lambda\subseteq\Z^2$ be a finite set whose hull embeds in $\Xi$.
Here $c^{-}=j$ and $c^{+}=j+1$, so that the algebra of the cylinder in Definition~\ref{def:index-sets} is $\Gr_\Xi((j,j+1))$.
The modes over $\Lambda$ form the index set $\mathcal{J}^{\mathrm{l}}(\Lambda)\subseteq\idx(\Pi)$, ordered as a subset of $\idx(\Pi)$,
and the modes over $\Lambda_{\mathrm{s}}$ form the set $\mathcal{J}^{\mathrm{l}}(\Lambda_{\mathrm{s}})$ of representatives of the modes at $\Pi$ modulo $\Gamma_\parallel$, with its order (Definition~\ref{def:modes}).
Every $i\in\mathcal{J}^{\mathrm{l}}(\Lambda)$ is $i=r(i)+g_i$ with its representative $r(i)\in\mathcal{J}^{\mathrm{l}}(\Lambda_{\mathrm{s}})$ and with $g_i\in\Gamma_\parallel$.
No two sites of $\Lambda$ differ by a non-zero element of $n_2\Z\times n_3\Z$, because the hull of $\Lambda$ embeds in $\Xi$;
hence the map $r$ is injective, that is, $i+g\notin\mathcal{J}^{\mathrm{l}}(\Lambda)$ for $i\in\mathcal{J}^{\mathrm{l}}(\Lambda)$ and $g\in\Gamma_\parallel\setminus\{0\}$.
So $\mathcal{J}^{\mathrm{l}}(\Lambda)$, like every subset of $\mathcal{J}^{\mathrm{l}}(\Lambda_{\mathrm{s}})$, satisfies the condition of Definition~\ref{def:index-sets} for the algebra $\Gr_\Xi((j,j+1))$;
read there, its modes satisfy $\zeta_i=\chi_\parallel(g_i)\cdot\zeta_{r(i)}$ (Definition~\ref{def:mode-translations}).
Item~\eqref{it:mode-relabelling} of Proposition~\ref{prop:localisation} therefore applies in $\Gr_\Xi((j,j+1))$ to the index set $\mathcal{J}^{\mathrm{l}}(\Lambda)$,
to the index set $r(\mathcal{J}^{\mathrm{l}}(\Lambda))$, ordered as a subset of $\mathcal{J}^{\mathrm{l}}(\Lambda_{\mathrm{s}})$, and to the bijection $r$ of $\mathcal{J}^{\mathrm{l}}(\Lambda)$ onto $r(\mathcal{J}^{\mathrm{l}}(\Lambda))$.
Its sign $s(I)$ is the product of the numbers $\chi_\parallel(g_i)$, $i\in I$, and of the sign of the permutation that sorts $(r(i))_{i\in I}$, listed in the increasing order of $i$,
and its unitary $U$ is the signed relabelling $U_\Lambda$ of Definition~\ref{def:modes};
for $\Upsilon\in\End\mathcal{F}(\mathcal{J}^{\mathrm{l}}(\Lambda))$ the operator $\iota(\Upsilon)$ on $\mathcal{F}=\mathcal{F}(\mathcal{J}^{\mathrm{l}}(\Lambda_{\mathrm{s}}))$ is, as fixed there,
$\iota_{r(\mathcal{J}^{\mathrm{l}}(\Lambda))\to\mathcal{J}^{\mathrm{l}}(\Lambda_{\mathrm{s}})}\left(U_\Lambda\Upsilon U_\Lambda^{*}\right)$.
A mode and its representative have the same crossing coefficient and, through the quotient map $\Z^2\to\Lambda_{\mathrm{s}}$, the same charge.
By the last identity of Proposition~\ref{prop:localisation}\eqref{it:mode-relabelling}, therefore,
\[
U_\Lambda\,\omega\,U_\Lambda^{*}=\omega,
\qquad
U_\Lambda\,\omega^{\frac12}\,U_\Lambda^{*}=\omega^{\frac12},
\qquad
U_\Lambda\,u(a)\,U_\Lambda^{*}=u(a)\quad\text{for }a\in\R^{\Lambda_{\mathrm{s}}},
\]
where on $\mathcal{F}(\mathcal{J}^{\mathrm{l}}(\Lambda))$ the function $a$ is read on $\Lambda$ through the quotient map.
Moreover $U_\Lambda$ preserves the number of modes, so that $U_\Lambda\Upsilon U_\Lambda^{*}$ is even if $\Upsilon$ is even.

Let $\Upsilon$ be even, let $\mathbb{T}$ be a torus with $\Xi(\mathbb{T},x_1)=\Xi$,
and let $a_j$ be the link variables between the planes $\{x_1=j\}$ and $\{x_1=j+1\}$ of a configuration of $\mathbb{T}$, a function on $\Lambda_{\mathrm{s}}$;
by Definition~\ref{def:modes} the diagonal operator of the phases of the modes at $\Pi$ is $\diag(u_I)=u(a_j)$.
In $\Gr_\Xi((j,j+1))$ apply the second identity of Proposition~\ref{prop:localisation}\eqref{it:localised-insertion}
to the index set $\mathcal{J}^{\mathrm{l}}(\Lambda_{\mathrm{s}})$, its subset $r(\mathcal{J}^{\mathrm{l}}(\Lambda))$ and the even operator $U_\Lambda\Upsilon U_\Lambda^{*}$,
and then the first identity of Proposition~\ref{prop:localisation}\eqref{it:mode-relabelling} to the operator
$\omega^{\frac12}\,\Upsilon\,u(a_j)\,\omega^{\frac12}=U_\Lambda^{*}\,\omega^{\frac12}\,U_\Lambda\Upsilon U_\Lambda^{*}\,u(a_j)\,\omega^{\frac12}\,U_\Lambda$.
Read in $\Gr_{\mathbb{T}}$ (Definition~\ref{def:index-sets}), the resulting identity is
\[
\begin{aligned}
&\Ins^{\mathcal{J}^{\mathrm{l}}(\Lambda_{\mathrm{s}})}\left[\omega^{\frac12}\,\iota(\Upsilon)\,u(a_j)\,\omega^{\frac12}\right]\\
&\qquad=\left[\Ins^{\mathcal{J}^{\mathrm{l}}(\Lambda)}\left[\omega^{\frac12}\,\Upsilon\,u(a_j)\,\omega^{\frac12}\right]\right]_{\mathbb{T}}
\cdot\prod_{i\in\mathcal{J}^{\mathrm{l}}(\Lambda_{\mathrm{s}})\setminus r(\mathcal{J}^{\mathrm{l}}(\Lambda))}\left(1+\omega_i\cdot u_i\cdot\Theta_\Pi(\zeta_i)\,\zeta_i\right),
\end{aligned}
\]
where $[\,\cdot\,]_{\mathbb{T}}\colon\Gr(\mathbb{X})\to\Gr_{\mathbb{T}}$ is the class map, and the insertion inside it is formed in $\Gr(\mathbb{X})$.
Indeed, this insertion lies in $\Gr(\mathbb{X}((j,j+1)))$,
and its image under $[\,\cdot\,]_\Xi$ is the insertion formed in $\Gr_\Xi((j,j+1))$, since $[\,\cdot\,]_\Xi$ is a homomorphism with $\Theta_\Pi[G]_\Xi=[\Theta_\Pi G]_\Xi$;
the reading in $\Gr_{\mathbb{T}}$ of $[G]_\Xi$ is $[G]_{\mathbb{T}}$ for every $G\in\Gr(\mathbb{X}((j,j+1)))$, both maps being homomorphisms that send every generator to the generator of $\Gr_{\mathbb{T}}$ with the same name.

Let further $\Lambda_k\subseteq\Z^2$ be finitely many finite sets whose hulls embed in $\Xi$ and whose images in $\Lambda_{\mathrm{s}}$ are pairwise disjoint,
and let $\Upsilon_k\in\End\mathcal{F}(\mathcal{J}^{\mathrm{l}}(\Lambda_k))$ be even, the localisation $\iota(\Upsilon_k)$ being formed with $U_{\Lambda_k}$.
The sets $r(\mathcal{J}^{\mathrm{l}}(\Lambda_k))$ consist of the modes of $\mathcal{J}^{\mathrm{l}}(\Lambda_{\mathrm{s}})$ over the images of the sets $\Lambda_k$, and so they are pairwise disjoint;
the operators $\iota(\Upsilon_k)$ commute with one another by Proposition~\ref{prop:localisation}\eqref{it:localisation-commute}.
Proposition~\ref{prop:localisation}\eqref{it:several-insertions}, applied in $\Gr_\Xi((j,j+1))$ to the index set $\mathcal{J}^{\mathrm{l}}(\Lambda_{\mathrm{s}})$, its subsets $r(\mathcal{J}^{\mathrm{l}}(\Lambda_k))$ and the even operators $U_{\Lambda_k}\Upsilon_kU_{\Lambda_k}^{*}$,
followed by the first identity of Proposition~\ref{prop:localisation}\eqref{it:mode-relabelling} for each $k$ as in the case of a single set, gives, read in $\Gr_{\mathbb{T}}$,
\[
\begin{aligned}
&\Ins^{\mathcal{J}^{\mathrm{l}}(\Lambda_{\mathrm{s}})}\left[\omega^{\frac12}\,\prod_k\iota(\Upsilon_k)\,u(a_j)\,\omega^{\frac12}\right]\\
&\qquad=\prod_k\left[\Ins^{\mathcal{J}^{\mathrm{l}}(\Lambda_k)}\left[\omega^{\frac12}\,\Upsilon_k\,u(a_j)\,\omega^{\frac12}\right]\right]_{\mathbb{T}}
\cdot\prod_{i}\left(1+\omega_i\cdot u_i\cdot\Theta_\Pi(\zeta_i)\,\zeta_i\right),
\end{aligned}
\]
the last product extending over the modes $i\in\mathcal{J}^{\mathrm{l}}(\Lambda_{\mathrm{s}})$ outside the sets $r(\mathcal{J}^{\mathrm{l}}(\Lambda_k))$.
\end{remark}

\subsubsection{Paths and kernels}\label{sec:paths-kernels}

\begin{lemma}[Unfolding]\label{lem:unfolding}
Let $n_1,n_2,n_3\ge4$, and let $\Gamma:=n_1\Z e_1+n_2\Z e_2+n_3\Z e_3$ be the period lattice of $C(n)$;
the slice torus $\Lambda_{\mathrm{s}}$, the lattice $\Gamma_{\mathrm{s}}$ and the plaquette $\sigma_0$ are those of Definition~\ref{def:slices}.
A \emph{path} is a tuple
\[
\mathsf{p}=\left(A_0,\dots,A_{n_1};a_0,\dots,a_{n_1-1}\right)
\qquad\text{with}\qquad
A_j\in\R^{\Lambda^{(1)}_{\mathrm{s}}},\quad
A_{n_1}-A_0\in\Gamma_{\mathrm{s}},\quad
a_j\in[0,2\pi/q)^{\Lambda_{\mathrm{s}}}.
\]
For an integer two-cochain $z$ on $C(n)$ and $j\in\Z/n_1\Z$,
the magnetic part $z^{\mathrm{mag}}_j(y):=z_{23}(j,y)$ and the electric part $z^{\mathrm{el}}_j(y,\nu):=z_{1\nu}(j,y)$ are formed as $F^{\mathrm{mag}}_j$ and $F^{\mathrm{el}}_j$ in Definition~\ref{def:slices}.
For a path $\mathsf{p}$ and $k\in\Z$ let $\varpi(\mathsf{p},k):=(A,z)$,
where $A\in C^1(\Gamma)$ is the one-cochain whose slice variables, in the sense of Definition~\ref{def:slices},
are $A_j$ on the links of the plane $x_1=j$ and $a_j$ on the links from the plane $x_1=j$ to the plane $x_1=j+1$, for $0\le j<n_1$,
and $z\in C^2(\Gamma;\Z)$ is the two-cochain with the parts
\[
z^{\mathrm{mag}}_j=k\,1_{\sigma_0}\quad(0\le j<n_1),\qquad
z^{\mathrm{el}}_j=0\quad(0\le j<n_1-1),\qquad
z^{\mathrm{el}}_{n_1-1}=\frac{q}{2\pi}\left(A_{n_1}-A_0\right).
\]
Then $dz=0$, so that $\varpi(\mathsf{p},k)$ belongs to the set $\mathfrak{C}(\Gamma)$ of Definition~\ref{def:torus-integrands},
and the curvature $F=dA+(2\pi/q)z$ of $\varpi(\mathsf{p},k)$ has the parts
\[
F^{\mathrm{el}}_j=A_{j+1}-A_j-d_{\mathrm{s}}a_j,\qquad
F^{\mathrm{mag}}_j=F^{\mathrm{mag}}(A_j,k)=d_{\mathrm{s}}A_j+(2\pi/q)k\,1_{\sigma_0}\qquad(0\le j<n_1).
\]
For $x\in\R^{\Lambda^{(1)}_{\mathrm{s}}}$ let $[x]$ denote the class of $x$ in the quotient $\R^{\Lambda^{(1)}_{\mathrm{s}}}/\Gamma_{\mathrm{s}}$, and put
\[
\begin{aligned}
\int\mathrm{d}\mathsf{p}\,(\cdots)
&:=\int_{\R^{\Lambda^{(1)}_{\mathrm{s}}}/\Gamma_{\mathrm{s}}}\mathrm{d}[A_0]
\sum_{g\in\Gamma_{\mathrm{s}}}
\int_{\left(\R^{\Lambda^{(1)}_{\mathrm{s}}}\right)^{n_1-1}}\prod_{j=1}^{n_1-1}\mathrm{d}A_j\\
&\qquad\int_{\left([0,2\pi/q)^{\Lambda_{\mathrm{s}}}\right)^{n_1}}\prod_{j=0}^{n_1-1}\mathrm{d}a_j\,(\cdots),
\qquad A_{n_1}:=A_0+g,
\end{aligned}
\]
with Lebesgue measure in the variables $A_j$ and $a_j$, and with the measure of $\R^{\Lambda^{(1)}_{\mathrm{s}}}/\Gamma_{\mathrm{s}}$ induced by Lebesgue measure.
Let $\Phi\ge0$ be a Borel function on $\mathfrak{C}(\Gamma)$ such that
\[
\Phi\left(A+(2\pi/q)\mathfrak{k},\,z-d\mathfrak{k}\right)=\Phi(A,z)\qquad\text{for all }(A,z)\in\mathfrak{C}(\Gamma)\text{ and }\mathfrak{k}\in C^1(\Gamma;\Z).
\]
Then, for every $k\in\Z$, the function of $A_0$ under the first integral of $\int\mathrm{d}\mathsf{p}\,\Phi\left(\varpi(\mathsf{p},k)\right)$ is $\Gamma_{\mathrm{s}}$-periodic, and
\[
\int_{\Conf(\Gamma)}\Phi=\sum_{k\in\Z}\int\mathrm{d}\mathsf{p}\,\Phi\left(\varpi(\mathsf{p},k)\right).
\]
The same holds for every complex Borel function $\Phi$ on $\mathfrak{C}(\Gamma)$ with this invariance and with $\int_{\Conf(\Gamma)}|\Phi|<\infty$;
all sums and integrals on the right then converge absolutely.
\end{lemma}

\begin{proof}
The argument is that of \cite[Lemma~10.3]{NoncompactQED}.

\emph{Closedness and curvature.}
Every cube of $C(n)$ lies between two consecutive planes $x_1=j$ and $x_1=j+1$, $j\in\Z/n_1\Z$, and is the cube $(j,\sigma)$ over a plaquette $\sigma$ of $\Lambda_{\mathrm{s}}$.
By Example~\ref{ex:four-mirrors}, every integer two-cochain $z$ on $C(n)$ satisfies
\begin{equation}\label{eq:slab-closed-slices}
(dz)(j,\sigma)=z^{\mathrm{mag}}_{j+1}(\sigma)-z^{\mathrm{mag}}_j(\sigma)-\left(d_{\mathrm{s}}z^{\mathrm{el}}_j\right)(\sigma),\qquad j\in\Z/n_1\Z .
\end{equation}
For the cochain $z$ of $\varpi(\mathsf{p},k)$ the magnetic parts do not depend on $j$, the electric parts vanish for $j<n_1-1$,
and $d_{\mathrm{s}}z^{\mathrm{el}}_{n_1-1}=0$, because $\frac{q}{2\pi}\left(A_{n_1}-A_0\right)$ is an integer one-cochain of $\Lambda_{\mathrm{s}}$ with vanishing coboundary, by the definition of $\Gamma_{\mathrm{s}}$.
So the right side of \eqref{eq:slab-closed-slices} vanishes for every $j$, and $dz=0$.
On $C(n)$ the plane $x_1=n_1$ is the plane $x_1=0$, so that the slice of $A$ in the plane $x_1=n_1$ is $A_0$.
Definition~\ref{def:slices}, applied to $dA$ and to $z$, gives $F^{\mathrm{mag}}_j=d_{\mathrm{s}}A_j+(2\pi/q)k\,1_{\sigma_0}$ for all $j$,
$F^{\mathrm{el}}_j=A_{j+1}-A_j-d_{\mathrm{s}}a_j$ for $j<n_1-1$, and
\[
F^{\mathrm{el}}_{n_1-1}=A_0-A_{n_1-1}-d_{\mathrm{s}}a_{n_1-1}+\left(A_{n_1}-A_0\right)=A_{n_1}-A_{n_1-1}-d_{\mathrm{s}}a_{n_1-1}.
\]

\emph{The left side in slice variables.}
By Definition~\ref{def:configurations}, $\int_{\Conf(\Gamma)}\Phi$ is the sum, over the closed $z\in C^2(\Gamma;\Z)$, of the integrals of $\Phi(A,z)$ over the box $[0,2\pi/q)$ in the variable of every link of $C(n)$.
In slice variables this box is the product of the boxes $[0,2\pi/q)^{\Lambda^{(1)}_{\mathrm{s}}}$ for the $A_j$ and $[0,2\pi/q)^{\Lambda_{\mathrm{s}}}$ for the $a_j$, $0\le j<n_1$.

\emph{Closed cochains.}
For $z\in C^2(\Gamma;\Z)$ define $\mathfrak{k}_0:=0$ and $\mathfrak{k}_{j+1}:=\mathfrak{k}_j+z^{\mathrm{el}}_j$ for $0\le j<n_1$; these are elements of $\Z^{\Lambda^{(1)}_{\mathrm{s}}}$.
By \eqref{eq:slab-closed-slices}, the conditions $(dz)(i,\sigma)=0$ for $0\le i<n_1-1$ and all $\sigma$ are equivalent, by summation and by taking differences, to
\[
z^{\mathrm{mag}}_j=z^{\mathrm{mag}}_0+d_{\mathrm{s}}\mathfrak{k}_j\qquad(0\le j<n_1),
\]
and under these relations the conditions at $i=n_1-1$ read $d_{\mathrm{s}}\mathfrak{k}_{n_1}=0$.
Conversely, $z^{\mathrm{el}}_j=\mathfrak{k}_{j+1}-\mathfrak{k}_j$ and the display recover $z$ from $z^{\mathrm{mag}}_0$ and $\mathfrak{k}_1,\dots,\mathfrak{k}_{n_1}$.
Hence $z\mapsto\left(z^{\mathrm{mag}}_0;\mathfrak{k}_1,\dots,\mathfrak{k}_{n_1}\right)$ is a bijection of the set of the closed integer two-cochains on $C(n)$
onto the set of the tuples in $\Z^{\Lambda^{(2)}_{\mathrm{s}}}\times\left(\Z^{\Lambda^{(1)}_{\mathrm{s}}}\right)^{n_1}$ whose last entry satisfies $d_{\mathrm{s}}\mathfrak{k}_{n_1}=0$.

\emph{The electric integers.}
Let $z$ be closed, and let $\mathfrak{k}\in C^1(\Gamma;\Z)$ be the cochain whose slice in the plane $x_1=j$ is $\mathfrak{k}_j$ for $0\le j<n_1$ and which vanishes on the links between consecutive planes.
By Definition~\ref{def:slices}, applied to $d\mathfrak{k}$, the magnetic parts of $d\mathfrak{k}$ are $d_{\mathrm{s}}\mathfrak{k}_j$,
and its electric parts are $\mathfrak{k}_{j+1}-\mathfrak{k}_j=z^{\mathrm{el}}_j$ for $j<n_1-1$ and $\mathfrak{k}_0-\mathfrak{k}_{n_1-1}=z^{\mathrm{el}}_{n_1-1}-\mathfrak{k}_{n_1}$ for $j=n_1-1$.
Consequently the pair $\left(A+(2\pi/q)\mathfrak{k},\,z-d\mathfrak{k}\right)$, at which $\Phi$ takes the value $\Phi(A,z)$, has the slices
\[
A'_j:=A_j+(2\pi/q)\mathfrak{k}_j\quad\text{and}\quad a_j\qquad(0\le j<n_1),
\]
the magnetic part $z^{\mathrm{mag}}_0$ in every plane, and the electric parts $0$ for $j<n_1-1$ and $\mathfrak{k}_{n_1}$ for $j=n_1-1$; note that $A'_0=A_0$.
For $1\le j<n_1$ the map $(A_j,\mathfrak{k}_j)\mapsto A'_j$ is a bijection of $[0,2\pi/q)^{\Lambda^{(1)}_{\mathrm{s}}}\times\Z^{\Lambda^{(1)}_{\mathrm{s}}}$ onto $\R^{\Lambda^{(1)}_{\mathrm{s}}}$
which carries the product of Lebesgue measure and counting measure to Lebesgue measure;
and $\mathfrak{k}_{n_1}\mapsto g:=(2\pi/q)\mathfrak{k}_{n_1}$ is a bijection of the integer one-cochains of $\Lambda_{\mathrm{s}}$ with vanishing coboundary onto $\Gamma_{\mathrm{s}}$.
With the bijection of the preceding paragraph and Tonelli's theorem it follows that
\begin{equation}\label{eq:slab-electric-removed}
\int_{\Conf(\Gamma)}\Phi=
\sum_{z^{\mathrm{mag}}_0\in\Z^{\Lambda^{(2)}_{\mathrm{s}}}}
\int_{[0,2\pi/q)^{\Lambda^{(1)}_{\mathrm{s}}}}\mathrm{d}A_0
\sum_{g\in\Gamma_{\mathrm{s}}}
\int\prod_{j=1}^{n_1-1}\mathrm{d}A'_j
\int\prod_{j=0}^{n_1-1}\mathrm{d}a_j\,
\Phi\left(\varpi'\right),
\end{equation}
where the $A'_j$ range over $\R^{\Lambda^{(1)}_{\mathrm{s}}}$ and the $a_j$ over $[0,2\pi/q)^{\Lambda_{\mathrm{s}}}$,
and $\varpi'\in\mathfrak{C}(\Gamma)$ is the pair with the slices $A_0,A'_1,\dots,A'_{n_1-1}$ and $a_0,\dots,a_{n_1-1}$,
the magnetic part $z^{\mathrm{mag}}_0$ in every plane and the electric parts $0,\dots,0,\frac{q}{2\pi}g$.

\emph{The magnetic integers.}
Let $z^{\mathrm{mag}}_0\in\Z^{\Lambda^{(2)}_{\mathrm{s}}}$, and let $k\in\Z$ be the sum of its values over the plaquettes of $\Lambda_{\mathrm{s}}$.
The values of the integer two-cochain $z^{\mathrm{mag}}_0-k\,1_{\sigma_0}$ of $\Lambda_{\mathrm{s}}$ sum to zero;
therefore $z^{\mathrm{mag}}_0-k\,1_{\sigma_0}=d_{\mathrm{s}}\mathfrak{k}'$ with some $\mathfrak{k}'\in\Z^{\Lambda^{(1)}_{\mathrm{s}}}$ \cite[(10.2), last line]{NoncompactQED},
and $\mathfrak{k}'$ is unique up to an integer one-cochain with vanishing coboundary, that is, up to an element of $\frac{q}{2\pi}\Gamma_{\mathrm{s}}$.
Conversely, for $k\in\Z$ and $\mathfrak{k}'\in\Z^{\Lambda^{(1)}_{\mathrm{s}}}$ the values of $k\,1_{\sigma_0}+d_{\mathrm{s}}\mathfrak{k}'$ sum to $k$:
by the formula for $d_{\mathrm{s}}$ in Definition~\ref{def:slices}, every link of $\Lambda_{\mathrm{s}}$ occurs in the coboundary at exactly two plaquettes, with opposite signs,
so that the values of $d_{\mathrm{s}}\mathfrak{k}'$ sum to zero.
Let $\mathcal{N}\subset\Z^{\Lambda^{(1)}_{\mathrm{s}}}$ contain exactly one element of each coset of $\frac{q}{2\pi}\Gamma_{\mathrm{s}}$.
Then $(k,\mathfrak{k}')\mapsto k\,1_{\sigma_0}+d_{\mathrm{s}}\mathfrak{k}'$ is a bijection of $\Z\times\mathcal{N}$ onto $\Z^{\Lambda^{(2)}_{\mathrm{s}}}$.

Write $z^{\mathrm{mag}}_0=k\,1_{\sigma_0}+d_{\mathrm{s}}\mathfrak{k}'$ with $(k,\mathfrak{k}')\in\Z\times\mathcal{N}$,
and shift $\varpi'$ by the integer one-cochain that has the slice $\mathfrak{k}'$ in every plane and vanishes on the links between the planes.
The coboundary of this cochain has the electric parts $\mathfrak{k}'-\mathfrak{k}'=0$, also between the planes $x_1=n_1-1$ and $x_1=n_1$, and the magnetic parts $d_{\mathrm{s}}\mathfrak{k}'$.
The shifted pair therefore has the slices $A_0+(2\pi/q)\mathfrak{k}'$, $A'_j+(2\pi/q)\mathfrak{k}'$ for $1\le j<n_1$, and $a_j$,
the magnetic parts $k\,1_{\sigma_0}$ and the electric parts $0,\dots,0,\frac{q}{2\pi}g$;
it is $\varpi(\mathsf{p},k)$ for the path
\[
\begin{aligned}
\mathsf{p}={}&\left(A_0+(2\pi/q)\mathfrak{k}',\,A'_1+(2\pi/q)\mathfrak{k}',\,\dots,\,A'_{n_1-1}+(2\pi/q)\mathfrak{k}',\right.\\
&\left.\quad A_0+(2\pi/q)\mathfrak{k}'+g;\,a_0,\dots,a_{n_1-1}\right),
\end{aligned}
\]
and $\Phi(\varpi')=\Phi\left(\varpi(\mathsf{p},k)\right)$.
For $k\in\Z$ and $A_0\in\R^{\Lambda^{(1)}_{\mathrm{s}}}$ let
\[
\begin{aligned}
f_k(A_0)&:=\sum_{g\in\Gamma_{\mathrm{s}}}\int\prod_{j=1}^{n_1-1}\mathrm{d}A_j\int\prod_{j=0}^{n_1-1}\mathrm{d}a_j\,\Phi\left(\varpi(\mathsf{p},k)\right)\in[0,\infty],\\
\mathsf{p}&=\left(A_0,A_1,\dots,A_{n_1-1},A_0+g;a_0,\dots,a_{n_1-1}\right),
\end{aligned}
\]
with the same ranges of the variables; this is the function of $A_0$ under the first integral of $\int\mathrm{d}\mathsf{p}\,\Phi\left(\varpi(\mathsf{p},k)\right)$.
With the bijection of $\Z\times\mathcal{N}$ onto $\Z^{\Lambda^{(2)}_{\mathrm{s}}}$ in the sum over $z^{\mathrm{mag}}_0$,
and after the translation of the variables $A'_j$ by $(2\pi/q)\mathfrak{k}'$, which preserves Lebesgue measure, \eqref{eq:slab-electric-removed} becomes
\begin{equation}\label{eq:slab-magnetic-removed}
\int_{\Conf(\Gamma)}\Phi=\sum_{k\in\Z}\sum_{\mathfrak{k}'\in\mathcal{N}}\int_{[0,2\pi/q)^{\Lambda^{(1)}_{\mathrm{s}}}}\mathrm{d}A_0\,f_k\left(A_0+(2\pi/q)\mathfrak{k}'\right).
\end{equation}

\emph{Periodicity.}
Let $g'\in\Gamma_{\mathrm{s}}$.
The integer one-cochain with the slice $\frac{q}{2\pi}g'$ in every plane and the value $0$ on the links between the planes has vanishing coboundary:
its electric parts are $0$, and its magnetic parts are $\frac{q}{2\pi}d_{\mathrm{s}}g'=0$.
The shift by this cochain adds $g'$ to every slice $A_j$ and leaves the $a_j$ and $z$ unchanged;
it carries $\varpi(\mathsf{p},k)$ to $\varpi(\mathsf{p}',k)$, where $\mathsf{p}'$ arises from $\mathsf{p}$ when $g'$ is added to each of $A_0,\dots,A_{n_1}$.
Hence $\Phi\left(\varpi(\mathsf{p}',k)\right)=\Phi\left(\varpi(\mathsf{p},k)\right)$,
and the translation of the variables $A_1,\dots,A_{n_1-1}$ by $g'$ in the definition of $f_k$ shows that $f_k(A_0+g')=f_k(A_0)$.

\emph{A fundamental domain.}
The space $\R^{\Lambda^{(1)}_{\mathrm{s}}}$ is the disjoint union of the boxes $[0,2\pi/q)^{\Lambda^{(1)}_{\mathrm{s}}}+(2\pi/q)\mathfrak{k}''$, $\mathfrak{k}''\in\Z^{\Lambda^{(1)}_{\mathrm{s}}}$,
and every such $\mathfrak{k}''$ is $\mathfrak{k}'+\frac{q}{2\pi}g$ with a unique $\mathfrak{k}'\in\mathcal{N}$ and a unique $g\in\Gamma_{\mathrm{s}}$.
Hence the translates $\mathbb{D}+g$, $g\in\Gamma_{\mathrm{s}}$, of the set
\[
\mathbb{D}:=\bigcup_{\mathfrak{k}'\in\mathcal{N}}\left([0,2\pi/q)^{\Lambda^{(1)}_{\mathrm{s}}}+(2\pi/q)\mathfrak{k}'\right)
\]
are pairwise disjoint and cover $\R^{\Lambda^{(1)}_{\mathrm{s}}}$: $\mathbb{D}$ is a fundamental domain of $\Gamma_{\mathrm{s}}$.
As $f_k$ is $\Gamma_{\mathrm{s}}$-periodic, the term of index $k$ in \eqref{eq:slab-magnetic-removed} equals
\[
\int_{\mathbb{D}}\mathrm{d}A_0\,f_k(A_0)=\int_{\R^{\Lambda^{(1)}_{\mathrm{s}}}/\Gamma_{\mathrm{s}}}\mathrm{d}[A_0]\,f_k(A_0)=\int\mathrm{d}\mathsf{p}\,\Phi\left(\varpi(\mathsf{p},k)\right).
\]
For $\Phi\ge0$ every exchange of sums and integrals above is justified by Tonelli's theorem.
For a complex $\Phi$ with $\int_{\Conf(\Gamma)}|\Phi|<\infty$ the identity for $|\Phi|$ shows that the right side converges absolutely,
and the same computation, with Fubini's theorem in place of Tonelli's theorem, gives the identity for $\Phi$.
\end{proof}

\begin{lemma}[Convolutions and Gaussian kernels]\label{lem:slice-kernels}
Let $n_2,n_3\ge4$, let $\Lambda_{\mathrm{s}}$ be the slice torus with the lattice $\Gamma_{\mathrm{s}}$ and the plaquette $\sigma_0$ of Definition~\ref{def:slices},
and let the spaces $\mathcal{H}^{\mathrm{l}}$, $\mathcal{H}^{\mathrm{s}}$ and the operators $E$, $U_{\mathrm{G}}(a)$, $P_{\mathrm{G}}$, $K_{\beta,S}$ be those of Definition~\ref{def:slice-operators};
functions of $E$ act in every flux sector, and trivially on the Fock space $\mathcal{F}$.
For a function $f$ on $\R^{\Lambda^{(1)}_{\mathrm{s}}}$ let $f^{\mathrm{per}}([x]):=\sum_{g\in\Gamma_{\mathrm{s}}}f(x+g)$ at every $x$ at which the series converges absolutely;
this is a function on $\R^{\Lambda^{(1)}_{\mathrm{s}}}/\Gamma_{\mathrm{s}}$.
For $x\in\R^{\Lambda^{(1)}_{\mathrm{s}}}$ and $S\subset\Lambda^{(1)}_{\mathrm{s}}$ write $|x|_S^2:=\sum_{b\in S}x_b^2$ and $|x|^2:=\sum_b x_b^2$.
\begin{enumerate}
\item\label{it:electric-convolution}
For $f\in L^1\left(\R^{\Lambda^{(1)}_{\mathrm{s}}}\right)$,
\[
\int\mathrm{d}x\,f(x)\exp\left(-i\langle x,E\rangle\right)=\hat f(E),\qquad
\hat f(\xi):=\int\mathrm{d}x\,f(x)\exp\left(-i\langle\xi,x\rangle\right)\quad\left(\xi\in\R^{\Lambda^{(1)}_{\mathrm{s}}}\right),
\]
and this operator is the integral operator on $L^2\left(\R^{\Lambda^{(1)}_{\mathrm{s}}}/\Gamma_{\mathrm{s}}\right)$ with the kernel $f^{\mathrm{per}}([A'-A])$.
\item\label{it:gaussian-kernels}
Let $\beta\in\R^{\Lambda^{(1)}_{\mathrm{s}}}$, $S\subset\Lambda^{(1)}_{\mathrm{s}}$ and $a,a'>0$, and let
\[
f(x):=\exp\left(i\langle\beta,x\rangle\right)\prod_{b\notin S}g_a(x_b)\prod_{b\in S}g_{a'}(x_b).
\]
The integral operator with the kernel $f^{\mathrm{per}}([A'-A])$ is
\[
\exp\left(-\frac{a}{2}\left(|E-\beta|^2-|E-\beta|_S^2\right)-\frac{a'}{2}|E-\beta|_S^2\right).
\]
In particular, $\exp\left(-\frac12|E-\beta|^2\right)$ has the kernel $f^{\mathrm{per}}([A'-A])$ with $f(x)=\exp\left(i\langle\beta,x\rangle\right)\prod_b g_1(x_b)$;
and, since $\sqrt{2}\cdot\exp\left(-\frac{s^2}{2}\right)\cdot g_1(s)=g_{1/2}(s)$ for $s\in\R$,
the kernel $f^{\mathrm{per}}([A'-A])$ with
\[
f(x)=\exp\left(i\langle\beta,x\rangle\right)\prod_{b\in S}\left(\sqrt{2}\cdot\exp\left(-\frac{x_b^2}{2}\right)\right)\prod_b g_1(x_b)
\]
is the kernel of $\exp\left(-\frac12|E-\beta|^2+\frac14|E-\beta|_S^2\right)$.
\item\label{it:gauge-average}
The map $a\mapsto U_{\mathrm{G}}(a)$ is a strongly continuous unitary representation on $\mathcal{H}^{\mathrm{l}}$
of the compact group $\R^{\Lambda_{\mathrm{s}}}$ modulo its subgroup $(2\pi/q)\Z^{\Lambda_{\mathrm{s}}}$, and
\[
\int_{[0,2\pi/q)^{\Lambda_{\mathrm{s}}}}\mathrm{d}a\,U_{\mathrm{G}}(a)=(2\pi/q)^{|\Lambda_{\mathrm{s}}|}P_{\mathrm{G}}.
\]
The operators $f(E)$, $f$ a bounded Borel function on $\R^{\Lambda^{(1)}_{\mathrm{s}}}$,
the translations $\exp\left(-i\langle x,E\rangle\right)$, $x\in\R^{\Lambda^{(1)}_{\mathrm{s}}}$, the operators $U_{\mathrm{G}}(a)$ and $P_{\mathrm{G}}$ commute with each other.
For every $\beta\in\R^{\Lambda^{(1)}_{\mathrm{s}}}$ and every $S\subset\Lambda^{(1)}_{\mathrm{s}}$ one has $K_{\beta,S}\ge0$ and
\[
K_{\beta,S}^{1/2}=(2\pi/q)^{\frac{|\Lambda_{\mathrm{s}}|}{2}}\exp\left(-\frac14|E-\beta|^2+\frac18|E-\beta|_S^2\right)P_{\mathrm{G}}.
\]
\item\label{it:flux-summability}
For $\varepsilon>0$ and $\beta\in\R^{\Lambda^{(2)}_{\mathrm{s}}}$,
\[
\sum_{k\in\Z}\int_{\R^{\Lambda^{(1)}_{\mathrm{s}}}/\Gamma_{\mathrm{s}}}\mathrm{d}A\,\exp\left(-\varepsilon\left\|F^{\mathrm{mag}}(A,k)+\beta\right\|^2\right)<\infty,
\]
where $F^{\mathrm{mag}}(A,k)=d_{\mathrm{s}}A+(2\pi/q)k\,1_{\sigma_0}$ as in Definition~\ref{def:slice-notation}.
\item\label{it:kernel-hilbert-schmidt}
Let $k^+_0\ge0$ be a Borel function on $\R^{\Lambda^{(1)}_{\mathrm{s}}}$ with $k^+_0(x)\le C\cdot\exp\left(-\varepsilon|x|^2\right)$ for some $C<\infty$ and $\varepsilon>0$, and let $k^+:=(k^+_0)^{\mathrm{per}}$.
Then $k^+$ is bounded and integrable on $\R^{\Lambda^{(1)}_{\mathrm{s}}}/\Gamma_{\mathrm{s}}$, hence square integrable;
$\|k^+\|_2$ denotes its norm in $L^2\left(\R^{\Lambda^{(1)}_{\mathrm{s}}}/\Gamma_{\mathrm{s}}\right)$.
Let $K'$ be an operator on $\mathcal{H}^{\mathrm{l}}$, diagonal in the flux index $k$, which acts by
\[
\left(K'\Phi\right)_k(A')=\int\mathrm{d}[A]\,K'_k\left([A'],[A]\right)\Phi_k(A)
\]
with kernels $K'_k([A'],[A])\in\End\mathcal{F}$ such that $\left\|K'_k([A'],[A])\right\|\le k^+([A'-A])$.
Let $\mathfrak{T}=\bigoplus_k\mathfrak{T}_k(A)$ be a multiplication operator from $\mathcal{H}^{\mathrm{l}}$, or from $\mathcal{H}^{\mathrm{s}}$, into $\mathcal{H}^{\mathrm{l}}$ with
\[
\|\mathfrak{T}\|_2^2:=\sum_{k\in\Z}\int_{\R^{\Lambda^{(1)}_{\mathrm{s}}}/\Gamma_{\mathrm{s}}}\mathrm{d}A\,\left\|\mathfrak{T}_k(A)\right\|^2<\infty.
\]
Then $K'\mathfrak{T}$ and $K'P_{\mathrm{G}}\mathfrak{T}$ are Hilbert--Schmidt operators with Hilbert--Schmidt norms at most $\sqrt{\dim\mathcal{F}}\cdot\|k^+\|_2\cdot\|\mathfrak{T}\|_2$,
and $K'P_{\mathrm{G}}$ acts in the flux sector $k$ by the kernel
\[
(2\pi/q)^{-|\Lambda_{\mathrm{s}}|}\int_{[0,2\pi/q)^{\Lambda_{\mathrm{s}}}}\mathrm{d}a\,K'_k\left([A'],[A+d_{\mathrm{s}}a]\right)u(a).
\]
\end{enumerate}
\end{lemma}

\begin{proof}
\eqref{it:electric-convolution}
By Definition~\ref{def:slice-operators}, $\exp\left(-i\langle x,E\rangle\right)$ is the translation $\Phi\mapsto\Phi(\cdot-x)$ of $L^2\left(\R^{\Lambda^{(1)}_{\mathrm{s}}}/\Gamma_{\mathrm{s}}\right)$.
Let $\mathfrak{E}$ be the joint spectral measure of the commuting self-adjoint operators $E_b$, $b\in\Lambda^{(1)}_{\mathrm{s}}$.
For $\Phi,\Phi'$ in this space the complex measure $\left\langle\Phi',\mathfrak{E}(\cdot)\Phi\right\rangle$ is finite, and the spectral theorem gives
\[
\left\langle\Phi',\exp\left(-i\langle x,E\rangle\right)\Phi\right\rangle=\int\left\langle\Phi',\mathfrak{E}(\mathrm{d}\xi)\Phi\right\rangle\exp\left(-i\langle x,\xi\rangle\right).
\]
Multiplication by $f(x)$, integration over $x$ and Fubini's theorem give
\[
\left\langle\Phi',\int\mathrm{d}x\,f(x)\exp\left(-i\langle x,E\rangle\right)\Phi\right\rangle=\int\left\langle\Phi',\mathfrak{E}(\mathrm{d}\xi)\Phi\right\rangle\hat f(\xi)=\left\langle\Phi',\hat f(E)\Phi\right\rangle .
\]
On the other hand, $\Phi$ being $\Gamma_{\mathrm{s}}$-periodic,
\[
\begin{aligned}
\left(\int\mathrm{d}x\,f(x)\exp\left(-i\langle x,E\rangle\right)\Phi\right)(A')
&=\int_{\R^{\Lambda^{(1)}_{\mathrm{s}}}}\mathrm{d}A\,f(A'-A)\,\Phi(A)\\
&=\int_{\R^{\Lambda^{(1)}_{\mathrm{s}}}/\Gamma_{\mathrm{s}}}\mathrm{d}[A]\,f^{\mathrm{per}}([A'-A])\,\Phi(A);
\end{aligned}
\]
the second equality follows from the decomposition of $\R^{\Lambda^{(1)}_{\mathrm{s}}}$ into the translates of a fundamental domain of $\Gamma_{\mathrm{s}}$,
since $\sum_{g\in\Gamma_{\mathrm{s}}}f(A'-A-g)=f^{\mathrm{per}}([A'-A])$ (note $\Gamma_{\mathrm{s}}=-\Gamma_{\mathrm{s}}$).

\eqref{it:gaussian-kernels}
For $a>0$ and $\xi\in\R$ one has $\int\mathrm{d}s\,g_a(s)\exp(-i\xi s)=\exp\left(-\frac{a}{2}\xi^2\right)$,
and the factor $\exp\left(i\langle\beta,x\rangle\right)$ replaces $\xi$ by $\xi-\beta$ in a Fourier transform.
Hence
\[
\begin{aligned}
\hat f(\xi)&=\prod_{b\notin S}\exp\left(-\frac{a}{2}(\xi_b-\beta_b)^2\right)\prod_{b\in S}\exp\left(-\frac{a'}{2}(\xi_b-\beta_b)^2\right)\\
&=\exp\left(-\frac{a}{2}\left(|\xi-\beta|^2-|\xi-\beta|_S^2\right)-\frac{a'}{2}|\xi-\beta|_S^2\right),
\end{aligned}
\]
and \eqref{it:electric-convolution} gives the first assertion.
The second is the case $a=a'=1$.
Finally, $\sqrt{2}\cdot\exp\left(-\frac{s^2}{2}\right)\cdot g_1(s)=\pi^{-\frac12}\cdot\exp\left(-s^2\right)=g_{1/2}(s)$;
so the last function is $f$ with $a=1$ and $a'=\frac12$, and
\[
-\frac12\left(|\xi-\beta|^2-|\xi-\beta|_S^2\right)-\frac14|\xi-\beta|_S^2=-\frac12|\xi-\beta|^2+\frac14|\xi-\beta|_S^2 .
\]

\eqref{it:gauge-average}
For $a,\tilde a\in\R^{\Lambda_{\mathrm{s}}}$ one has $U_{\mathrm{G}}(a)U_{\mathrm{G}}(\tilde a)=U_{\mathrm{G}}(a+\tilde a)$,
since $u$ is a representation (Lemma~\ref{lem:half-slab-matrices}) and $d_{\mathrm{s}}$ is linear.
If $a\in(2\pi/q)\Z^{\Lambda_{\mathrm{s}}}$, then $u(a)=1$ by Lemma~\ref{lem:half-slab-matrices},
and $d_{\mathrm{s}}a\in\Gamma_{\mathrm{s}}$, because $\frac{q}{2\pi}d_{\mathrm{s}}a$ is an integer one-cochain whose coboundary vanishes;
as the vectors of $\mathcal{H}^{\mathrm{l}}$ are $\Gamma_{\mathrm{s}}$-periodic functions, $U_{\mathrm{G}}(a)=1$.
So $U_{\mathrm{G}}$ is a unitary representation of the compact group $\R^{\Lambda_{\mathrm{s}}}$ modulo $(2\pi/q)\Z^{\Lambda_{\mathrm{s}}}$;
it is strongly continuous, since translations act strongly continuously on $L^2$ and $a\mapsto u(a)$ is continuous.
The normalised Haar measure of this group is $(q/2\pi)^{|\Lambda_{\mathrm{s}}|}\mathrm{d}a$ on $[0,2\pi/q)^{\Lambda_{\mathrm{s}}}$.
The average $(q/2\pi)^{|\Lambda_{\mathrm{s}}|}\int\mathrm{d}a\,U_{\mathrm{G}}(a)$ is self-adjoint, as $a\mapsto-a$ preserves Haar measure;
its range consists of invariant vectors, by the invariance of Haar measure; and it fixes every invariant vector.
It is therefore the orthogonal projection onto the invariant vectors, which is $P_{\mathrm{G}}$ by Definition~\ref{def:slice-operators}.
Next, $U_{\mathrm{G}}(a)=u(a)\exp\left(-i\langle d_{\mathrm{s}}a,E\rangle\right)$, where the constant matrix $u(a)$ acts on $\mathcal{F}$ alone,
while the translations commute with each other and act trivially on $\mathcal{F}$.
Hence $U_{\mathrm{G}}(a)$ commutes with every translation, therefore with the joint spectral measure of $E$ and with every $f(E)$;
the same holds for $P_{\mathrm{G}}$, an integral of the $U_{\mathrm{G}}(a)$, which also commutes with every $U_{\mathrm{G}}(a)$.
By Definition~\ref{def:slice-operators}, $K_{\beta,S}=K_\beta\exp\left(\frac14|E-\beta|_S^2\right)$ with $K_\beta=(2\pi/q)^{|\Lambda_{\mathrm{s}}|}\exp\left(-\frac12|E-\beta|^2\right)P_{\mathrm{G}}$.
Since $P_{\mathrm{G}}$ commutes with the functions of $E$,
\[
K_{\beta,S}=(2\pi/q)^{|\Lambda_{\mathrm{s}}|}\exp\left(-\frac12|E-\beta|^2+\frac14|E-\beta|_S^2\right)P_{\mathrm{G}},
\]
a product of two commuting non-negative bounded operators, the function $\xi\mapsto\exp\left(-\frac12|\xi-\beta|^2+\frac14|\xi-\beta|_S^2\right)$ taking its values in $(0,1]$.
So $K_{\beta,S}\ge0$; the operator on the right of the asserted formula is non-negative, by the same argument, and its square is $K_{\beta,S}$;
by the uniqueness of non-negative square roots it is $K_{\beta,S}^{1/2}$.

\eqref{it:flux-summability}
Let $1_{\mathrm{s}}$ be the two-cochain of $\Lambda_{\mathrm{s}}$ with the value $1$ at every plaquette,
and read two-cochains of $\Lambda_{\mathrm{s}}$ as functions of the lowest corner $y$ of the plaquette, as in Definition~\ref{def:slices}.
For a real two-cochain $\beta'$ of $\Lambda_{\mathrm{s}}$, the formula for $d_{\mathrm{s}}$ in Definition~\ref{def:slices} gives the adjoint for the counting inner products:
\[
\left(d_{\mathrm{s}}^{*}\beta'\right)(y,2)=\beta'(y)-\beta'(y-e_3),\qquad
\left(d_{\mathrm{s}}^{*}\beta'\right)(y,3)=\beta'(y-e_2)-\beta'(y).
\]
So $d_{\mathrm{s}}^{*}\beta'=0$ if and only if $\beta'$ is invariant under the translations by $e_2$ and $e_3$, that is, constant on the torus $\Lambda_{\mathrm{s}}$:
$\ker d_{\mathrm{s}}^{*}=\R\,1_{\mathrm{s}}$.
Consequently the range of $d_{\mathrm{s}}$ is the orthogonal complement $1_{\mathrm{s}}^{\perp}$ of $1_{\mathrm{s}}$;
in particular $\left\langle d_{\mathrm{s}}A,1_{\mathrm{s}}\right\rangle=0$ for every $A\in\R^{\Lambda^{(1)}_{\mathrm{s}}}$.
Let $k\in\Z$.
The orthogonal projection of $F^{\mathrm{mag}}(A,k)+\beta$ onto $\R\,1_{\mathrm{s}}$ is
\[
\frac{1}{|\Lambda_{\mathrm{s}}|}\left((2\pi/q)k+\sum_{\sigma}\beta(\sigma)\right)1_{\mathrm{s}},
\]
the sum over the plaquettes $\sigma$ of $\Lambda_{\mathrm{s}}$; it does not depend on $A$.
The difference between $(2\pi/q)k\,1_{\sigma_0}+\beta$ and this projection lies in $1_{\mathrm{s}}^{\perp}$, the range of $d_{\mathrm{s}}$;
choose $A^{\perp}_k\in\R^{\Lambda^{(1)}_{\mathrm{s}}}$ such that $d_{\mathrm{s}}A^{\perp}_k$ is this difference.
Then $F^{\mathrm{mag}}(A,k)+\beta$ is the orthogonal sum of $d_{\mathrm{s}}\left(A+A^{\perp}_k\right)$ and of the projection, and $\|1_{\mathrm{s}}\|^2=|\Lambda_{\mathrm{s}}|$ gives
\[
\left\|F^{\mathrm{mag}}(A,k)+\beta\right\|^2=\left\|d_{\mathrm{s}}\left(A+A^{\perp}_k\right)\right\|^2+\frac{1}{|\Lambda_{\mathrm{s}}|}\left((2\pi/q)k+\sum_{\sigma}\beta(\sigma)\right)^2.
\]
The measure of $\R^{\Lambda^{(1)}_{\mathrm{s}}}/\Gamma_{\mathrm{s}}$ is invariant under translations, and $d_{\mathrm{s}}$ vanishes on $\Gamma_{\mathrm{s}}$; hence
\[
\int_{\R^{\Lambda^{(1)}_{\mathrm{s}}}/\Gamma_{\mathrm{s}}}\mathrm{d}A\,\exp\left(-\varepsilon\left\|F^{\mathrm{mag}}(A,k)+\beta\right\|^2\right)
=\exp\left(-\frac{\varepsilon}{|\Lambda_{\mathrm{s}}|}\left((2\pi/q)k+\sum_{\sigma}\beta(\sigma)\right)^2\right)\cdot C_{\mathrm{s}},
\]
where
\[
C_{\mathrm{s}}:=\int_{\R^{\Lambda^{(1)}_{\mathrm{s}}}/\Gamma_{\mathrm{s}}}\mathrm{d}A\,\exp\left(-\varepsilon\left\|d_{\mathrm{s}}A\right\|^2\right).
\]
The number $C_{\mathrm{s}}$ is finite.
Indeed, $\ker d_{\mathrm{s}}$ is the kernel of an integer matrix and is therefore spanned by integer vectors,
so that $\Gamma_{\mathrm{s}}=(2\pi/q)\left(\Z^{\Lambda^{(1)}_{\mathrm{s}}}\cap\ker d_{\mathrm{s}}\right)$ is a lattice of full rank in $\ker d_{\mathrm{s}}$.
Thus $\R^{\Lambda^{(1)}_{\mathrm{s}}}/\Gamma_{\mathrm{s}}$ is the product of the compact torus $\ker d_{\mathrm{s}}/\Gamma_{\mathrm{s}}$
and of the orthogonal complement $\left(\ker d_{\mathrm{s}}\right)^{\perp}$, with the product of the Lebesgue measures.
The integrand depends on the second factor only, and $d_{\mathrm{s}}$ is injective on it,
so that $\left\|d_{\mathrm{s}}x\right\|^2\ge C^{-1}|x|^2$ for $x\in\left(\ker d_{\mathrm{s}}\right)^{\perp}$ with some $C<\infty$;
the Gaussian integral over $\left(\ker d_{\mathrm{s}}\right)^{\perp}$ converges.
Finally, $\sum_{k\in\Z}\exp\left(-\frac{\varepsilon}{|\Lambda_{\mathrm{s}}|}\left((2\pi/q)k+\sum_{\sigma}\beta(\sigma)\right)^2\right)<\infty$.

\eqref{it:kernel-hilbert-schmidt}
For $x\in\R^{\Lambda^{(1)}_{\mathrm{s}}}$ let $x_0$ be its orthogonal projection onto $\ker d_{\mathrm{s}}$.
For $g\in\Gamma_{\mathrm{s}}\subset\ker d_{\mathrm{s}}$ one has $|x+g|^2=|x_0+g|^2+|x-x_0|^2\ge|x_0+g|^2$, hence
$k^+([x])\le C\cdot\sum_{g\in\Gamma_{\mathrm{s}}}\exp\left(-\varepsilon|x_0+g|^2\right)$.
The series on the right converges uniformly for $x_0$ in compact subsets of $\ker d_{\mathrm{s}}$,
by $|x_0+g|^2\ge\frac12|g|^2-|x_0|^2$ and the convergence of $\sum_{g\in\Gamma_{\mathrm{s}}}\exp\left(-\frac{\varepsilon}{2}|g|^2\right)$ over the lattice $\Gamma_{\mathrm{s}}$;
its sum is continuous and $\Gamma_{\mathrm{s}}$-periodic on $\ker d_{\mathrm{s}}$, hence bounded, as $\ker d_{\mathrm{s}}/\Gamma_{\mathrm{s}}$ is compact.
So $k^+$ is bounded.
Tonelli's theorem on a fundamental domain of $\Gamma_{\mathrm{s}}$ gives $\int\mathrm{d}[x]\,k^+([x])=\int\mathrm{d}x\,k^+_0(x)<\infty$,
and a bounded integrable function is square integrable.

For $\Phi\in\mathcal{H}^{\mathrm{l}}$ and $a\in\R^{\Lambda_{\mathrm{s}}}$ the substitution $A''=A+d_{\mathrm{s}}a$, which preserves the measure, gives
\[
\begin{aligned}
\left(K'U_{\mathrm{G}}(a)\Phi\right)_k(A')
&=\int\mathrm{d}[A'']\,K'_k\left([A'],[A'']\right)u(a)\,\Phi_k\left(A''-d_{\mathrm{s}}a\right)\\
&=\int\mathrm{d}[A]\,K'_k\left([A'],[A+d_{\mathrm{s}}a]\right)u(a)\,\Phi_k(A).
\end{aligned}
\]
Averaging over $a\in[0,2\pi/q)^{\Lambda_{\mathrm{s}}}$ with \eqref{it:gauge-average} and Fubini's theorem gives the asserted kernel of $K'P_{\mathrm{G}}$.
As $u(a)$ is unitary, the norm of this kernel is at most
\[
\tilde k^+([A'-A]):=(2\pi/q)^{-|\Lambda_{\mathrm{s}}|}\int_{[0,2\pi/q)^{\Lambda_{\mathrm{s}}}}\mathrm{d}a\,k^+\left([A'-A-d_{\mathrm{s}}a]\right),
\]
and Minkowski's integral inequality gives $\|\tilde k^+\|_2\le\|k^+\|_2$, every translate of $k^+$ having the norm $\|k^+\|_2$.

The operator $K'\mathfrak{T}$ acts in the flux sector $k$ by the kernel $K'_k([A'],[A])\,\mathfrak{T}_k(A)$,
and $K'P_{\mathrm{G}}\mathfrak{T}$ by the kernel of $K'P_{\mathrm{G}}$ multiplied on the right by $\mathfrak{T}_k(A)$;
the operator norms of these kernels are at most $k^+([A'-A])\cdot\left\|\mathfrak{T}_k(A)\right\|$ and $\tilde k^+([A'-A])\cdot\left\|\mathfrak{T}_k(A)\right\|$.
A linear map from $\mathcal{F}$ or from the Fock space $\mathcal{F}'$ of a site plane into $\mathcal{F}$ has rank at most $\dim\mathcal{F}$,
so that the square of its Hilbert--Schmidt norm is at most $\dim\mathcal{F}$ times the square of its operator norm.
By Lemma~\ref{lem:poisson-kernels}\eqref{it:hilbert-schmidt-kernels}, applied to the measure space $\Z\times\R^{\Lambda^{(1)}_{\mathrm{s}}}/\Gamma_{\mathrm{s}}$ with the counting measure in the flux index,
an operator with a square integrable kernel is Hilbert--Schmidt, its Hilbert--Schmidt norm being the norm of the kernel in $L^2$.
Therefore, by the translation invariance of the measure,
\[
\begin{aligned}
\left\|K'\mathfrak{T}\right\|_{\mathrm{HS}}^2
&\le\dim\mathcal{F}\cdot\sum_{k\in\Z}\int\mathrm{d}[A']\int\mathrm{d}[A]\,k^+([A'-A])^2\cdot\left\|\mathfrak{T}_k(A)\right\|^2\\
&=\dim\mathcal{F}\cdot\|k^+\|_2^2\cdot\|\mathfrak{T}\|_2^2,
\end{aligned}
\]
and the same computation with $\tilde k^+$ in place of $k^+$ bounds $\left\|K'P_{\mathrm{G}}\mathfrak{T}\right\|_{\mathrm{HS}}$.
\end{proof}

\begin{corollary}[Trace class]\label{cor:trace-class}
Let $w_0,\dots,w_3>0$, $M\in\R$, $q>0$ and $-1\le\rho\le1$.
Let $n_2,n_3\ge4$, let $\Lambda_{\mathrm{s}}$ be the slice torus of the rectangular cross-section with the sides $n_2,n_3$ (Definition~\ref{def:slices}),
and let $\mathcal{H}^{\mathrm{l}}$, $\mathcal{H}^{\mathrm{s}}$, $K$, $K_{\beta,S}$ and $\mathfrak{M}$ be the spaces and the operators of Definition~\ref{def:slice-operators} for this cross-section.
For every $\beta\in\R^{\Lambda^{(1)}_{\mathrm{s}}}$ and every subset $S$ of $\Lambda^{(1)}_{\mathrm{s}}$,
the operator $K_{\beta,S}^{1/2}\mathfrak{M}$ from $\mathcal{H}^{\mathrm{s}}$ to $\mathcal{H}^{\mathrm{l}}$ is Hilbert--Schmidt.
The operators
\[
\begin{aligned}
K^{1/2}\mathfrak{M}\mathfrak{M}^{*}K^{1/2}&=\left(K^{1/2}\mathfrak{M}\right)\left(K^{1/2}\mathfrak{M}\right)^{*}&&\text{on }\mathcal{H}^{\mathrm{l}},\\
\mathfrak{M}^{*}K\mathfrak{M}&=\left(K^{1/2}\mathfrak{M}\right)^{*}\left(K^{1/2}\mathfrak{M}\right)&&\text{on }\mathcal{H}^{\mathrm{s}}
\end{aligned}
\]
are non-negative and of trace class.
\end{corollary}

\begin{proof}
Fix $\beta\in\R^{\Lambda^{(1)}_{\mathrm{s}}}$ and $S\subset\Lambda^{(1)}_{\mathrm{s}}$.
By Lemma~\ref{lem:slice-kernels}\eqref{it:gauge-average},
\[
K_{\beta,S}^{1/2}=(2\pi/q)^{\frac{|\Lambda_{\mathrm{s}}|}{2}}K'P_{\mathrm{G}},\qquad
K':=\exp\left(-\frac14\left(|E-\beta|^2-|E-\beta|_S^2\right)-\frac18|E-\beta|_S^2\right),
\]
and by Lemma~\ref{lem:slice-kernels}\eqref{it:gaussian-kernels} with $a=\frac12$ and $a'=\frac14$,
the operator $K'$ acts in every flux sector, trivially on $\mathcal{F}$, by the kernel $f^{\mathrm{per}}([A'-A])$, where
\[
f(x):=\exp\left(i\langle\beta,x\rangle\right)\prod_{b\notin S}g_{1/2}(x_b)\prod_{b\in S}g_{1/4}(x_b).
\]
By the triangle inequality, $\left|f^{\mathrm{per}}\right|\le k^+:=(k^+_0)^{\mathrm{per}}$ with $k^+_0(x):=\prod_{b\notin S}g_{1/2}(x_b)\prod_{b\in S}g_{1/4}(x_b)$.
As $g_{1/2}(s)=\pi^{-\frac12}\cdot\exp\left(-s^2\right)$ and $g_{1/4}(s)=\sqrt{2}\cdot\pi^{-\frac12}\cdot\exp\left(-2s^2\right)\le\sqrt{2}\cdot\pi^{-\frac12}\cdot\exp\left(-s^2\right)$,
and as $\Lambda_{\mathrm{s}}$ has $2|\Lambda_{\mathrm{s}}|$ links,
\[
k^+_0(x)\le2^{\frac{\#|S|}{2}}\cdot\pi^{-|\Lambda_{\mathrm{s}}|}\cdot\exp\left(-|x|^2\right).
\]
By Definition~\ref{def:slice-operators}, $\mathfrak{M}$ is the multiplication operator from $\mathcal{H}^{\mathrm{s}}$ into $\mathcal{H}^{\mathrm{l}}$ by
\[
\mathfrak{M}_k(A)=\exp\left(-\frac14\left\|F^{\mathrm{mag}}(A,k)\right\|^2\right)\widehat{\mathfrak{N}}(A);
\]
hence
\[
\left\|\mathfrak{M}_k(A)\right\|^2\le\left(\sup_A\left\|\widehat{\mathfrak{N}}(A)\right\|\right)^2\cdot\exp\left(-\frac12\left\|F^{\mathrm{mag}}(A,k)\right\|^2\right),
\]
the supremum being finite by Lemma~\ref{lem:half-slab-matrices}.
In particular $\mathfrak{M}$ is bounded, and Lemma~\ref{lem:slice-kernels}\eqref{it:flux-summability} with $\varepsilon=\frac12$ and $\beta=0$ gives $\|\mathfrak{M}\|_2<\infty$.
Lemma~\ref{lem:slice-kernels}\eqref{it:kernel-hilbert-schmidt}, with $\mathfrak{T}=\mathfrak{M}$, now shows that $K'P_{\mathrm{G}}\mathfrak{M}$ is Hilbert--Schmidt, and so is $K_{\beta,S}^{1/2}\mathfrak{M}$.

For $\beta=0$ and $S=\emptyset$ one has $K_{0,\emptyset}=K$, and Lemma~\ref{lem:slice-kernels}\eqref{it:gauge-average} gives
\[
K_{0,\emptyset}^{1/2}=(2\pi/q)^{\frac{|\Lambda_{\mathrm{s}}|}{2}}\exp\left(-\frac14|E|^2\right)P_{\mathrm{G}}=K^{1/2}.
\]
Thus $K^{1/2}\mathfrak{M}$ is Hilbert--Schmidt, and the products of this operator with its adjoint, in either order, are non-negative and of trace class.
Since $K^{1/2}$ is self-adjoint with square $K$, these products are $K^{1/2}\mathfrak{M}\mathfrak{M}^{*}K^{1/2}$ and $\mathfrak{M}^{*}K\mathfrak{M}$.
\end{proof}

\subsubsection{Word formula and external field}\label{sec:words}

\begin{lemma}[A trace over paths]\label{lem:path-trace}
Let $w_0,\dots,w_3>0$, $M\in\R$, $q>0$ and $-1\le\rho\le1$.
Let $n_1\ge2$ and $n_2,n_3\ge4$ be integers.
For the rectangular cross-section with the sides $n_2,n_3$,
let $\Lambda_{\mathrm{s}}$, $\Gamma_{\mathrm{s}}$ and $\sigma_0$ be as in Definition~\ref{def:slices},
let $\mathcal{F}$, $\mathcal{F}'$, $u$ and $\widehat{\mathfrak{N}}$ be as in Definition~\ref{def:modes},
and let $\mathcal{H}^{\mathrm{l}}$ and $P_{\mathrm{G}}$ be as in Definition~\ref{def:slice-operators};
on $\R^{\Lambda^{(2)}_{\mathrm{s}}}$ the norm $\|\cdot\|$ is the Euclidean one,
and $F^{\mathrm{mag}}(A,k):=d_{\mathrm{s}}A+(2\pi/q)k\,1_{\sigma_0}$ for $A\in\R^{\Lambda^{(1)}_{\mathrm{s}}}$ and $k\in\Z$.
For $j\in\Z/n_1\Z$ let
\[
\mathcal{W}_j\colon\R^{\Lambda^{(2)}_{\mathrm{s}}}\to\C
\qquad\text{and}\qquad
\widetilde K_j\colon\R^{\Lambda^{(1)}_{\mathrm{s}}}\times\R^{\Lambda^{(1)}_{\mathrm{s}}}\to\End\mathcal{F}
\]
be Borel functions,
and let $C,\varepsilon>0$ be constants such that for every $j$, all $\beta\in\R^{\Lambda^{(2)}_{\mathrm{s}}}$, all $A,A'\in\R^{\Lambda^{(1)}_{\mathrm{s}}}$ and all $g\in\Gamma_{\mathrm{s}}$
\[
\left|\mathcal{W}_j(\beta)\right|\le C\cdot\exp\left(-\varepsilon\|\beta\|^2\right),
\qquad
\left\|\widetilde K_j(A',A)\right\|\le C\cdot\exp\left(-\varepsilon|A'-A|^2\right),
\qquad
\widetilde K_j(A'+g,A+g)=\widetilde K_j(A',A).
\]
Let $K_j$ be the operator on $\mathcal{H}^{\mathrm{l}}$ that acts in every flux sector by the same kernel
\[
K_j\left([A'],[A]\right):=\sum_{g\in\Gamma_{\mathrm{s}}}\widetilde K_j(A'+g,A),
\]
and let $\mathfrak{T}_j$ be the multiplication operator on $\mathcal{H}^{\mathrm{l}}$ that acts in the flux sector $k$ by the function
$A\mapsto\mathcal{W}_j\left(F^{\mathrm{mag}}(A,k)\right)\cdot\widehat{\mathfrak{N}}(A)\widehat{\mathfrak{N}}(A)^{*}$.
Then every operator $(2\pi/q)^{|\Lambda_{\mathrm{s}}|}K_jP_{\mathrm{G}}\mathfrak{T}_j$ is Hilbert--Schmidt,
the integral on the left of
\[
\begin{aligned}
&\sum_{k\in\Z}\int\mathrm{d}\mathsf{p}\,\prod_{j=0}^{n_1-1}\mathcal{W}_j\left(F^{\mathrm{mag}}(A_j,k)\right)\\
&\qquad\cdot\operatorname{tr}_{\mathcal{F}'}\prod_{j=n_1-1,\dots,0}\left(\widehat{\mathfrak{N}}(A_{j+1})^{*}\,\widetilde K_j\left(A_{j+1},A_j+d_{\mathrm{s}}a_j\right)u(a_j)\,\widehat{\mathfrak{N}}(A_j)\right)\\
&\quad=\Tr_{\mathcal{H}^{\mathrm{l}}}\prod_{j=n_1-1,\dots,0}\left((2\pi/q)^{|\Lambda_{\mathrm{s}}|}K_jP_{\mathrm{G}}\mathfrak{T}_j\right)
\end{aligned}
\]
converges absolutely, and the identity holds.
Here $\int\mathrm{d}\mathsf{p}$ is the integral over the paths $\mathsf{p}=(A_0,\dots,A_{n_1};a_0,\dots,a_{n_1-1})$ with $A_{n_1}-A_0\in\Gamma_{\mathrm{s}}$,
given by the formula of Lemma~\ref{lem:unfolding} with $n_1$ steps,
and in a product over $j=n_1-1,\dots,0$ the factors with larger $j$ stand to the left.
\end{lemma}

\begin{proof}
\emph{The factors.}
Let $k^{+}_0(x):=C\cdot\exp\left(-\varepsilon|x|^2\right)$ for $x\in\R^{\Lambda^{(1)}_{\mathrm{s}}}$, and $k^{+}:=(k^{+}_0)^{\mathrm{per}}$.
The series that defines $K_j([A'],[A])$ converges absolutely, and its norm is at most $k^{+}([A'-A])$.
Its sum depends only on the classes of $A'$ and $A$:
a translation of $A'$ by an element of $\Gamma_{\mathrm{s}}$ permutes the terms,
and for $g'\in\Gamma_{\mathrm{s}}$ the third hypothesis gives $\sum_{g}\widetilde K_j(A'+g,A+g')=\sum_{g}\widetilde K_j(A'+g-g',A)$,
which is again the same sum with permuted terms.
Since $k^{+}$ is integrable on $\R^{\Lambda^{(1)}_{\mathrm{s}}}/\Gamma_{\mathrm{s}}$ (Lemma~\ref{lem:slice-kernels}\eqref{it:kernel-hilbert-schmidt}),
Young's inequality shows that $K_j$ is a bounded operator.
By Lemma~\ref{lem:half-slab-matrices} the function $\widehat{\mathfrak{N}}$ is bounded, and for every flux $k$ and every $A$
\[
\left\|\mathcal{W}_j\left(F^{\mathrm{mag}}(A,k)\right)\cdot\widehat{\mathfrak{N}}(A)\widehat{\mathfrak{N}}(A)^{*}\right\|
\le C\cdot\sup_{A''}\left\|\widehat{\mathfrak{N}}(A'')\right\|^2\cdot\exp\left(-\varepsilon\left\|F^{\mathrm{mag}}(A,k)\right\|^2\right).
\]
The sum over $k$ of the integrals over $\R^{\Lambda^{(1)}_{\mathrm{s}}}/\Gamma_{\mathrm{s}}$ of the square of the right side is finite,
by Lemma~\ref{lem:slice-kernels}\eqref{it:flux-summability} with $2\varepsilon$ in place of $\varepsilon$ and $\beta=0$.
So Lemma~\ref{lem:slice-kernels}\eqref{it:kernel-hilbert-schmidt} applies to $K_j$ and $\mathfrak{T}_j$:
the operator $\mathcal{Z}_j:=(2\pi/q)^{|\Lambda_{\mathrm{s}}|}K_jP_{\mathrm{G}}\mathfrak{T}_j$ is Hilbert--Schmidt, and its kernel in the flux sector $k$ is
\begin{equation}\label{eq:slab-path-kernel}
\mathcal{Z}_j\left(k;[A'],[A]\right)
=\int_{[0,2\pi/q)^{\Lambda_{\mathrm{s}}}}\mathrm{d}a\,\sum_{g\in\Gamma_{\mathrm{s}}}\widetilde K_j\left(A'+g,A+d_{\mathrm{s}}a\right)u(a)\cdot\mathcal{W}_j\left(F^{\mathrm{mag}}(A,k)\right)\cdot\widehat{\mathfrak{N}}(A)\widehat{\mathfrak{N}}(A)^{*}.
\end{equation}

\emph{Majorants.}
Let $\mathcal{Z}^{+}_j$ be the operator on $\bigoplus_{k\in\Z}L^2\left(\R^{\Lambda^{(1)}_{\mathrm{s}}}/\Gamma_{\mathrm{s}}\right)$, diagonal in $k$, with the non-negative kernel
\[
\mathcal{Z}^{+}_j\left(k;[A'],[A]\right)
:=\int_{[0,2\pi/q)^{\Lambda_{\mathrm{s}}}}\mathrm{d}a\,k^{+}\left([A'-A-d_{\mathrm{s}}a]\right)\cdot C\cdot\sup_{A''}\left\|\widehat{\mathfrak{N}}(A'')\right\|^2\cdot\exp\left(-\varepsilon\left\|F^{\mathrm{mag}}(A,k)\right\|^2\right).
\]
By the hypotheses, the term with the indices $(a,g)$ in \eqref{eq:slab-path-kernel} has norm at most
\[
k^{+}_0\left(A'+g-A-d_{\mathrm{s}}a\right)\cdot C\cdot\sup_{A''}\left\|\widehat{\mathfrak{N}}(A'')\right\|^2\cdot\exp\left(-\varepsilon\left\|F^{\mathrm{mag}}(A,k)\right\|^2\right),
\]
and the sum over $g$ of these numbers is the integrand in the definition of $\mathcal{Z}^{+}_j$.
For fixed $k$ and $[A]$, Minkowski's inequality and the invariance of the measure of $\R^{\Lambda^{(1)}_{\mathrm{s}}}/\Gamma_{\mathrm{s}}$ under translations
bound the $L^2$ norm of $[A']\mapsto\int\mathrm{d}a\,k^{+}\left([A'-A-d_{\mathrm{s}}a]\right)$ by $(2\pi/q)^{|\Lambda_{\mathrm{s}}|}\cdot\|k^{+}\|_2$,
which is finite by Lemma~\ref{lem:slice-kernels}\eqref{it:kernel-hilbert-schmidt}.
Hence the square of the $L^2$ norm of the kernel of $\mathcal{Z}^{+}_j$ is at most
\[
(2\pi/q)^{2|\Lambda_{\mathrm{s}}|}\cdot\|k^{+}\|_2^2\cdot C^2\cdot\sup_{A''}\left\|\widehat{\mathfrak{N}}(A'')\right\|^4\cdot\sum_{k\in\Z}\int\mathrm{d}A\,\exp\left(-2\varepsilon\left\|F^{\mathrm{mag}}(A,k)\right\|^2\right),
\]
which is finite by Lemma~\ref{lem:slice-kernels}\eqref{it:flux-summability} once more;
by Lemma~\ref{lem:poisson-kernels}\eqref{it:hilbert-schmidt-kernels}, $\mathcal{Z}^{+}_j$ is a Hilbert--Schmidt operator, and its Hilbert--Schmidt norm is the $L^2$ norm of its kernel.

\emph{The trace as an integral of kernels.}
Since $n_1\ge2$, the product $\mathcal{Z}_{n_1-2}\cdots\mathcal{Z}_0$ is a Hilbert--Schmidt operator.
Lemma~\ref{lem:poisson-kernels}\eqref{it:hilbert-schmidt-kernels}, applied to $\mathcal{Z}_{n_1-1}$ and to this product, together with the rule given there for the kernel of a product, yields
\begin{equation}\label{eq:slab-path-trace}
\Tr\prod_{j=n_1-1,\dots,0}\mathcal{Z}_j
=\sum_{k\in\Z}\int\mathrm{d}[A_0]\cdots\mathrm{d}[A_{n_1-1}]\,\operatorname{tr}_{\mathcal{F}}\prod_{j=n_1-1,\dots,0}\mathcal{Z}_j\left(k;[A_{j+1}],[A_j]\right),
\qquad[A_{n_1}]:=[A_0].
\end{equation}
The same identity holds for the operators $\mathcal{Z}^{+}_j$, with non-negative integrands;
there the value is $\Tr\prod_j\mathcal{Z}^{+}_j\le\prod_j\|\mathcal{Z}^{+}_j\|_{\mathrm{HS}}<\infty$,
by the Cauchy--Schwarz inequality for the Hilbert--Schmidt inner product and the submultiplicativity of the Hilbert--Schmidt norm.
Insert \eqref{eq:slab-path-kernel} into \eqref{eq:slab-path-trace}:
the integrand becomes an integral over $(a_0,\dots,a_{n_1-1})$ of a sum over $(g'_0,\dots,g'_{n_1-1})\in\Gamma_{\mathrm{s}}^{n_1}$,
$a_j$ and $g'_j$ being the variables of the kernel of $\mathcal{Z}_j$.
The trace of an element of $\End\mathcal{F}$ is in modulus at most $\dim\mathcal{F}$ times its norm, and the norm of a product is at most the product of the norms;
so each term has modulus at most $\dim\mathcal{F}$ times the corresponding term for the operators $\mathcal{Z}^{+}_j$,
and by Tonelli's theorem the integral of the latter terms is $\Tr\prod_j\mathcal{Z}^{+}_j<\infty$.
Hence the integral in \eqref{eq:slab-path-trace}, with the kernels written out, converges absolutely,
and by Fubini's theorem its variables may be integrated in any order.

\emph{Removing the periodisations.}
For a path $\mathsf{p}$, a flux $k$ and $0\le j<n_1$ put
\[
\Phi_j(A',A):=\widetilde K_j\left(A',A+d_{\mathrm{s}}a_j\right)u(a_j)\cdot\mathcal{W}_j\left(F^{\mathrm{mag}}(A,k)\right)\cdot\widehat{\mathfrak{N}}(A)\widehat{\mathfrak{N}}(A)^{*}.
\]
Then $\Phi_j(A'+g,A+g)=\Phi_j(A',A)$ for $g\in\Gamma_{\mathrm{s}}$:
the third hypothesis applies;
$d_{\mathrm{s}}g=0$, so that $F^{\mathrm{mag}}(A+g,k)=F^{\mathrm{mag}}(A,k)$;
and $\widehat{\mathfrak{N}}$ is $(2\pi/q)\Z^{\Lambda^{(1)}_{\mathrm{s}}}$-periodic (Lemma~\ref{lem:half-slab-matrices}), while $\Gamma_{\mathrm{s}}\subset(2\pi/q)\Z^{\Lambda^{(1)}_{\mathrm{s}}}$.
In particular the integrand of the integral over paths below is invariant under the translation of all of $A_0,\dots,A_{n_1}$ by one element of $\Gamma_{\mathrm{s}}$,
so that, after the integration over $A_1,\dots,A_{n_1-1}$ and the summation over $A_{n_1}-A_0$, it is a function of $[A_0]$.
Fix a Borel fundamental domain of $\Gamma_{\mathrm{s}}$ in $\R^{\Lambda^{(1)}_{\mathrm{s}}}$.
For $1\le j<n_1$ write $A_j=A^{\flat}_j+g_j$ with $A^{\flat}_j$ in this domain and $g_j\in\Gamma_{\mathrm{s}}$;
let $A^{\flat}_0:=A_0$ be taken in the domain, and put $g_0:=0$, $g_{n_1}:=A_{n_1}-A_0$ and $A^{\flat}_{n_1}:=A^{\flat}_0$.
Then $\Phi_j(A_{j+1},A_j)=\Phi_j\left(A^{\flat}_{j+1}+g_{j+1}-g_j,A^{\flat}_j\right)$,
and the map $(g_j)_{1\le j\le n_1}\mapsto(g_{j+1}-g_j)_{0\le j<n_1}$ is a bijection of $\Gamma_{\mathrm{s}}^{n_1}$ onto itself.
Therefore
\[
\begin{aligned}
&\int\mathrm{d}\mathsf{p}\,\operatorname{tr}_{\mathcal{F}}\prod_{j=n_1-1,\dots,0}\Phi_j(A_{j+1},A_j)\\
&\quad=\int\mathrm{d}A^{\flat}_0\cdots\mathrm{d}A^{\flat}_{n_1-1}\int\mathrm{d}a_0\cdots\mathrm{d}a_{n_1-1}\,
\operatorname{tr}_{\mathcal{F}}\prod_{j=n_1-1,\dots,0}\sum_{g'_j\in\Gamma_{\mathrm{s}}}\Phi_j\left(A^{\flat}_{j+1}+g'_j,A^{\flat}_j\right),
\end{aligned}
\]
the $A^{\flat}_j$ ranging over the fundamental domain and the $a_j$ over $[0,2\pi/q)^{\Lambda_{\mathrm{s}}}$.
Summed over $k$, the right side is the right side of \eqref{eq:slab-path-trace} with the kernels \eqref{eq:slab-path-kernel} written out.
The same substitution, applied to the majorants of the previous step, shows that the left side, summed over $k$, converges absolutely.

\emph{Cyclicity.}
For $\Upsilon_0,\dots,\Upsilon_{n_1-1}\in\End\mathcal{F}$ one has
\[
\operatorname{tr}_{\mathcal{F}}\prod_{j=n_1-1,\dots,0}\left(\Upsilon_j\,\widehat{\mathfrak{N}}(A_j)\widehat{\mathfrak{N}}(A_j)^{*}\right)
=\operatorname{tr}_{\mathcal{F}'}\prod_{j=n_1-1,\dots,0}\left(\widehat{\mathfrak{N}}(A_{j+1})^{*}\,\Upsilon_j\,\widehat{\mathfrak{N}}(A_j)\right):
\]
the factor $\widehat{\mathfrak{N}}(A_0)^{*}$ at the right end may be moved to the left end,
and $\widehat{\mathfrak{N}}(A_{n_1})=\widehat{\mathfrak{N}}(A_0+g_{n_1})=\widehat{\mathfrak{N}}(A_0)$.
Put $\Upsilon_j:=\widetilde K_j(A_{j+1},A_j+d_{\mathrm{s}}a_j)u(a_j)$ and take the scalars $\mathcal{W}_j(F^{\mathrm{mag}}(A_j,k))$ out of the trace:
the integrand on the left of the asserted identity is $\operatorname{tr}_{\mathcal{F}}\prod_j\Phi_j(A_{j+1},A_j)$.
With the previous step and \eqref{eq:slab-path-trace}, the left side of the asserted identity is therefore equal to $\Tr\prod_j\mathcal{Z}_j$.
\end{proof}

\begin{proposition}[The word formula]\label{prop:word-formula}
Let $w_0,\dots,w_3>0$, $M\in\R$, $q>0$ and $-1\le\rho\le1$.
A \emph{local step datum} is a triple $(S,\Lambda,k)$ in which
$S$ is a finite set of links of the slice lattice $\Z^2$,
$\Lambda$ is a non-empty finite set of sites of $\Z^2$ that contains the end points of the links of $S$,
and $k=k(A',A)$ is a Borel function on $\R^S\times\R^S$ with values in the even part of $\End\mathcal{F}\left(\mathcal{J}^{\mathrm{l}}(\Lambda)\right)$,
the Fock space of the modes over $\Lambda$ at a link plane (Definition~\ref{def:modes}), such that
\[
u(a)^{-1}k\left(A'+d_{\mathrm{s}}a,A+d_{\mathrm{s}}a\right)u(a)=k(A',A)\qquad\text{for all }a\in\R^{\Lambda}\text{ and all }A',A\in\R^S,
\]
$d_{\mathrm{s}}a$ being read on $S$, and such that, for some constants $C<\infty$ and $\varepsilon_k>0$,
\[
\left\|k(A',A)\right\|\le C\cdot\prod_{b\in S}\exp\left(-\varepsilon_k\left(A'(b)-A(b)\right)^2\right)\qquad\text{for all }A',A\in\R^S.
\]
For $j\in\Z$ and a local step datum $(S,\Lambda,k)$, the integrand $G[k]$ at the plane $\{x_1=j+\frac12\}$ is
\[
G[k](\hat A):=\Ins_{j+\frac12}\left[\omega^{\frac12}\,k\left(\hat A_{j+1}|_S,\left(\hat A_j+d_{\mathrm{s}}\hat a_j\right)|_S\right)u(\hat a_j)\,\omega^{\frac12}\right]
\cdot\prod_{b\in S}g_1\left(F^{\mathrm{el}}_j(b)\right)^{-1}\cdot\exp\left(S^{\times}_{j+\frac12}(\Lambda)\right)
\]
for real one-cochains $\hat A$ on $\Z^3$.
Here $\hat A_j$, $\hat A_{j+1}$ and $\hat a_j$ are the slice variables of $\hat A$, $F^{\mathrm{el}}_j=\hat A_{j+1}-\hat A_j-d_{\mathrm{s}}\hat a_j$ is the electric part of $d\hat A$,
and $g_1$ is the Gaussian density (Definition~\ref{def:slices});
the insertion $\Ins_{j+\frac12}$, the operators $\omega$ and $u(\hat a_j)$ on $\mathcal{F}\left(\mathcal{J}^{\mathrm{l}}(\Lambda)\right)$ and the element $S^{\times}_{j+\frac12}(\Lambda)$
are formed as in Definition~\ref{def:insertions}, at the plane $\{x_1=j+\frac12\}$, with the modes $\mathcal{J}^{\mathrm{l}}(\Lambda)$ and with $\hat A$;
and $B\left(j+\frac12,S,\Lambda\right)$ is the box of that definition.
\begin{enumerate}
\item\label{it:step-class}
$G[k]$ belongs to $\mathcal{G}^{\mathrm{ev}}\left(B\left(j+\frac12,S,\Lambda\right)\right)$ (Definition~\ref{def:integrands}), with the constant $\varepsilon_G=\min\left(\varepsilon_k,\frac12\right)$;
its construction involves the lattice $\Z^3$ alone, and no torus enters it.
\item\label{it:word-identity}
Let $n_1,n_2,n_3\ge4$ and $n=(n_1,n_2,n_3)$, and let $\Gamma$ be the period lattice of $C(n)$.
For the cross-section of $C(n)$ for $x_1$, the slice torus $\Lambda_{\mathrm{s}}$, the lattice $\Gamma_{\mathrm{s}}$, the plaquette $\sigma_0$, the spaces $\mathcal{H}^{\mathrm{l}}$ and $\mathcal{H}^{\mathrm{s}}$,
the localisation $\iota$, the operators $P_{\mathrm{G}}$, $\mathfrak{M}$, $K$, $K_\beta$, $K_{\beta,S}$ and the magnetic curvature $F^{\mathrm{mag}}(A,k)$ are those of Definition~\ref{def:slice-notation};
the class $\mathcal{I}(C(n))$ and the functional $Z_{C(n)}$ are those of Definition~\ref{def:torus-integrands},
and $F^{\mathrm{mag}}_j$ and $F^{\mathrm{el}}_j$ are the parts of the curvature $F$ (Definition~\ref{def:slices}).
For $j\in\Z/n_1\Z$, represented by $0\le j<n_1$,
let $u_j$ be a bounded Borel function on $\R^{\Lambda^{(2)}_{\mathrm{s}}}$,
let $v_j$ be a Borel function on $\R^{\Lambda^{(1)}_{\mathrm{s}}}$ for which there are constants $C<\infty$ and $\varepsilon_v>0$ with
$|v_j(x)|\le C\cdot\exp\left(\left(\frac12-\varepsilon_v\right)|x|^2\right)$ for all $x\in\R^{\Lambda^{(1)}_{\mathrm{s}}}$,
and let $(S_{j,l},\Lambda_{j,l},k_{j,l})_l$ be a finite family of local step data, the integrands $G[k_{j,l}]$ being formed at the plane $\{x_1=j+\frac12\}$.
Assume:
\begin{enumerate}
\item\label{it:step-disjoint}
the boxes $B\left(j+\frac12,S_{j,l},\Lambda_{j,l}\right)$, for all $j$ and all $l$, embed disjointly in $C(n)$;
\item\label{it:step-string}
no translate of one of these boxes by an element of the period lattice $\Gamma$ contains a plaquette of $\Z^3$ that lies in a plane $\{x_1=j'\}$, $j'\in\Z$, and represents $\sigma_0$;
\item\label{it:step-free}
$v_j(x)$ does not depend on the coordinates $x_b$ with $b\in\bigcup_lS_{j,l}$.
\end{enumerate}
Then the integrand $\prod_j\left(u_j(F^{\mathrm{mag}}_j)\cdot v_j(F^{\mathrm{el}}_j)\cdot\prod_lG[k_{j,l}]\right)$ belongs to $\mathcal{I}(C(n))$, and
\[
Z_{C(n)}\left[\prod_{j}\left(u_j(F^{\mathrm{mag}}_j)\cdot v_j(F^{\mathrm{el}}_j)\cdot\prod_{l}G[k_{j,l}]\right)\right]
=(2\pi)^{|n|}\cdot\Tr_{\mathcal{H}^{\mathrm{l}}}\prod_{j=n_1-1,\dots,0}\left((2\pi/q)^{|\Lambda_{\mathrm{s}}|}K_jP_{\mathrm{G}}\,\mathfrak{M}u_j\mathfrak{M}^{*}\right),
\]
where $K_j$ is the operator on $\mathcal{H}^{\mathrm{l}}$ that acts in every flux sector by the periodisation over $\Gamma_{\mathrm{s}}$, as in Lemma~\ref{lem:path-trace}, of the kernel
\[
v_j(A'-A)\cdot\prod_{b\in\Lambda^{(1)}_{\mathrm{s}}\setminus\bigcup_lS_{j,l}}g_1\left(A'(b)-A(b)\right)\cdot\prod_l\iota\left(k_{j,l}(A',A)\right),
\]
and $u_j$ acts on $\mathcal{H}^{\mathrm{s}}$, in the flux sector $k$, as the multiplication by $A\mapsto u_j\left(F^{\mathrm{mag}}(A,k)\right)$.
Here the links of $S_{j,l}$ are identified with their images in $\Lambda^{(1)}_{\mathrm{s}}$,
and $k_{j,l}(A',A)$ is the value of $k_{j,l}$ at the restrictions of $A'$ and $A$ to $S_{j,l}$.
The $n_1$ operators $(2\pi/q)^{|\Lambda_{\mathrm{s}}|}K_jP_{\mathrm{G}}\,\mathfrak{M}u_j\mathfrak{M}^{*}$ whose product stands under the trace are Hilbert--Schmidt operators.
\item\label{it:word-cases}
In the situation of \eqref{it:word-identity}:
\begin{enumerate}
\item
if $v_j(x)=\exp\left(i\langle\beta_j,x\rangle\right)$ with $\beta_j\in\R^{\Lambda^{(1)}_{\mathrm{s}}}$, and the family of local step data at the step $j$ is empty,
then $(2\pi/q)^{|\Lambda_{\mathrm{s}}|}K_jP_{\mathrm{G}}=K_{\beta_j}$;
\item
if $v_j(x)=\exp\left(i\langle\beta,x\rangle\right)\cdot\prod_{b\in S}\left(2^{\frac12}\cdot\exp\left(-\frac{x_b^2}{2}\right)\right)$
with $\beta\in\R^{\Lambda^{(1)}_{\mathrm{s}}}$ and an arbitrary subset $S$ of $\Lambda^{(1)}_{\mathrm{s}}$,
and the family of local step data at the step $j$ is empty,
then $(2\pi/q)^{|\Lambda_{\mathrm{s}}|}K_jP_{\mathrm{G}}=K_{\beta,S}$;
\item
without local step data the conditions \eqref{it:step-disjoint}--\eqref{it:step-free} are void,
and the identity of \eqref{it:word-identity} holds for every choice of the plaquette $\sigma_0$;
\item
$Z_{C(n)}[1]
=(2\pi)^{|n|}\cdot\Tr\left(K\mathfrak{M}\mathfrak{M}^{*}\right)^{n_1}
=(2\pi)^{|n|}\cdot\Tr\left(K^{1/2}\mathfrak{M}\mathfrak{M}^{*}K^{1/2}\right)^{n_1}
=(2\pi)^{|n|}\cdot\Tr\left(\mathfrak{M}^{*}K\mathfrak{M}\right)^{n_1}$.
\end{enumerate}
\end{enumerate}
\end{proposition}

\begin{proof}
We begin with \eqref{it:step-class}.
Fix $j\in\Z$ and a local step datum $(S,\Lambda,k)$, and let $B:=B\left(j+\frac12,S,\Lambda\right)$.
The construction of $G[k]$ refers to the lattice $\Z^3$ only.

\emph{What $G[k]$ reads.}
The value $G[k](\hat A)$ depends on $\hat A$ only through its values on the copies of the links of $S$ in the planes $\{x_1=j\}$ and $\{x_1=j+1\}$
and on the links from $(j,y)$ to $(j+1,y)$, $y\in\Lambda$;
its Grassmann factors are formed with the fields at the end points of the bonds over $\Lambda$ that cross the plane $\{x_1=j+\frac12\}$.
By Proposition~\ref{prop:geometry}\eqref{it:geometry-projection} the map $\pi$ sends the lower end of such a bond over $y$ to $(j,y)$ and its upper end to $(j+1,y)$.
By Definition~\ref{def:insertions} the box $B$ contains all these links and the cubes $(j,y)+\left[-\frac12,\frac12\right]^3$ and $(j+1,y)+\left[-\frac12,\frac12\right]^3$, $y\in\Lambda$;
hence the end points belong to $\mathbb{X}(B)$ (Definition~\ref{def:integrands}).
So $G[k]$ is a map from $\R^{B^{(1)}}$ to $\Gr(\mathbb{X}(B))$;
it is a Borel map, since $k$ is Borel and the other factors are continuous.

\emph{Parity.}
The matrices $k(A',A)$, $u(a)$ and $\omega^{\frac12}$ are even, the last two being diagonal.
So the entry $(I,J)$ of the matrix in the insertion vanishes unless $\#|I|\equiv\#|J|\pmod 2$, and for such $I$, $J$ the element $\Theta_{j+\frac12}(\zeta_J)\zeta_I$ is even.
The element $S^{\times}_{j+\frac12}(\Lambda)$ is a sum of products of two generators, so its exponential is even.
Hence $G[k]$ is even.

\emph{The Gaussian bound.}
At $A'=\hat A_{j+1}|_S$ and $A=\left(\hat A_j+d_{\mathrm{s}}\hat a_j\right)|_S$ one has $A'-A=F^{\mathrm{el}}_j$ on $S$,
so that $\|k(A',A)\|\le C\cdot\prod_{b\in S}\exp\left(-\varepsilon_kF^{\mathrm{el}}_j(b)^2\right)$.
The entries of $u(\hat a_j)$ have modulus one and those of $\omega^{\frac12}$ are bounded;
the coefficients of $\exp\left(S^{\times}_{j+\frac12}(\Lambda)\right)$ are polynomials in the phases $\exp\left(\pm iq\hat a_j(y)\right)$, $y\in\Lambda$, with coefficients that depend on the parameters only;
and a coefficient of a product of two elements of the Grassmann algebra is a finite sum of products of coefficients of the factors.
Since $g_1(s)^{-1}=(2\pi)^{\frac12}\cdot\exp\left(\frac{s^2}{2}\right)$, every coefficient of $G[k](\hat A)$ is therefore bounded by
\[
C'\cdot C\cdot\prod_{b\in S}\left((2\pi)^{\frac12}\cdot\exp\left(-\varepsilon_kF^{\mathrm{el}}_j(b)^2+\frac{F^{\mathrm{el}}_j(b)^2}{2}\right)\right)
\le C''\cdot\exp\left(\left(\frac12-\varepsilon_G\right)\sum_{b\in S}F^{\mathrm{el}}_j(b)^2\right),
\qquad\varepsilon_G=\min\left(\varepsilon_k,\frac12\right),
\]
with constants $C',C''$ that depend on $(S,\Lambda,k)$ and on the parameters only.
For $b\in S$ the number $F^{\mathrm{el}}_j(b)$ is the value of $d\hat A$ on the plaquette bounded by the copies of $b$ in the planes $\{x_1=j\}$ and $\{x_1=j+1\}$
and by the two links between these planes over the end points of $b$;
these four links lie in $B$, and $B$ is convex, so this plaquette belongs to $B^{(2)}$.
As $\frac12-\varepsilon_G\ge0$, the sum over $b\in S$ may be replaced by the sum of $(d\hat A)_\sigma^2$ over $\sigma\in B^{(2)}$,
and the bound required in Definition~\ref{def:integrands} follows.

\emph{Gauge transformations.}
Let $\lambda$ be a real function on the sites of $B$, and let $\lambda_j:=\lambda(j,\cdot)$ and $\lambda_{j+1}:=\lambda(j+1,\cdot)$ on $\Lambda$.
Under the gauge transformation \eqref{eq:setting-gauge} with $\lambda$,
\[
\hat A_j\mapsto\hat A_j-d_{\mathrm{s}}\lambda_j,
\qquad
\hat A_{j+1}\mapsto\hat A_{j+1}-d_{\mathrm{s}}\lambda_{j+1},
\qquad
\hat a_j\mapsto\hat a_j-\lambda_{j+1}+\lambda_j,
\]
so that $\hat A_j+d_{\mathrm{s}}\hat a_j\mapsto\hat A_j+d_{\mathrm{s}}\hat a_j-d_{\mathrm{s}}\lambda_{j+1}$.
The covariance of $k$, with $a=-\lambda_{j+1}$, and the multiplicativity of $u$ turn the matrix $k\left(\hat A_{j+1},\hat A_j+d_{\mathrm{s}}\hat a_j\right)u(\hat a_j)$ (restrictions to $S$ understood) into
\[
u(\lambda_{j+1})^{-1}k\left(\hat A_{j+1},\hat A_j+d_{\mathrm{s}}\hat a_j\right)u(\lambda_{j+1})\cdot u(\hat a_j)\,u(\lambda_{j+1})^{-1}u(\lambda_j)
=u(\lambda_{j+1})^{-1}k\left(\hat A_{j+1},\hat A_j+d_{\mathrm{s}}\hat a_j\right)u(\hat a_j)\,u(\lambda_j),
\]
because diagonal matrices commute.
They commute with $\omega^{\frac12}$ as well,
so the entry $(I,J)$ of the matrix in the insertion is multiplied by $\exp\left(-iq\langle\lambda_{j+1},\mathfrak{q}(I)\rangle+iq\langle\lambda_j,\mathfrak{q}(J)\rangle\right)$.
On the other hand, by Definition~\ref{def:modes} a mode $\zeta_i$ over $y\in\Lambda$ is a combination of the components of $\psi$ at the upper end of its bond if $\mathfrak{q}(i)=\delta_y$,
and of those of $\bar\psi$ there if $\mathfrak{q}(i)=-\delta_y$;
by \eqref{eq:setting-reflection-fields}, the element $\Theta_{j+\frac12}(\zeta_i)$ is then a combination of the components of $\bar\psi$, respectively of $\psi$, at the lower end.
The upper end lies over $(j+1,y)$ and the lower end over $(j,y)$.
So the transformation multiplies $\zeta_I$ by $\exp\left(iq\langle\lambda_{j+1},\mathfrak{q}(I)\rangle\right)$,
and it multiplies $\Theta_{j+\frac12}(\zeta_J)$ by $\exp\left(-iq\langle\lambda_j,\mathfrak{q}(J)\rangle\right)$,
exactly as in the second identity of \eqref{eq:slab-gauge-modes} with $\lambda_j$ in the place of $\lambda$;
in $\sum_{I,J}(\cdots)_{IJ}\,\Theta_{j+\frac12}(\zeta_J)\zeta_I$ all phases cancel.
Each $F^{\mathrm{el}}_j(b)$ is a value of $d\hat A$, and each term of $S^{\times}_{j+\frac12}(\Lambda)$ is invariant by \eqref{eq:setting-gauge}.
So $G[k]$ is gauge invariant, and \eqref{it:step-class} is proved.

To prove~\eqref{it:word-identity}, let $\mathbb{T}:=C(n)$, with the period lattice $\Gamma$, and represent $j\in\Z/n_1\Z$ by $0\le j<n_1$.
For the local step data of~\eqref{it:word-identity} write
\[
B_{j,l}:=B\left(j+\tfrac12,S_{j,l},\Lambda_{j,l}\right),
\]
and let
\[
G:=\prod_j\left(u_j\left(F^{\mathrm{mag}}_j\right)\cdot v_j\left(F^{\mathrm{el}}_j\right)\cdot\prod_lG[k_{j,l}]\right)
\]
be the integrand on the left side of~\eqref{it:word-identity}.

\emph{Step~1: the integrand belongs to $\mathcal{I}(\mathbb{T})$.}
By~\eqref{it:step-disjoint} the boxes $B_{j,l}$ embed disjointly in $C(n)$;
so their translates by the elements of $\Gamma$, for all $j$ and all $l$, are pairwise disjoint, and the product of the $G[k_{j,l}]$ is read on $\mathbb{T}$.
Each $G[k_{j,l}]$ belongs to $\mathcal{G}^{\mathrm{ev}}(B_{j,l})$ by~\eqref{it:step-class}.
Lemma~\ref{lem:product-integrands}\eqref{it:product-torus} therefore makes the product of the $G[k_{j,l}]$ a Borel map on $\mathfrak{C}(\Gamma)$ with values in $\Gr_{\mathbb{T}}$,
invariant under the integer shifts of Definition~\ref{def:torus-integrands},
whose value at $(A,z)$ is obtained by evaluating every factor $G[k_{j,l}]$ at a local lift of $(A,z)$ on $B_{j,l}$;
the factors $u_j(F^{\mathrm{mag}}_j)$ and $v_j(F^{\mathrm{el}}_j)$ are Borel functions of the curvature $F=dA+(2\pi/q)z$, which the integer shifts do not change.
Hence $G$ is a Borel map on $\mathfrak{C}(\Gamma)$, invariant under the integer shifts.
As shown above in the proof of~\eqref{it:step-class}, every coefficient of $G[k_{j,l}]$, evaluated at a local lift on $B_{j,l}$, is bounded by
\[
C\cdot\exp\left(\left(\tfrac12-\min\left(\varepsilon_k,\tfrac12\right)\right)\sum_{b\in S_{j,l}}F^{\mathrm{el}}_j(b)^2\right),
\]
where $\varepsilon_k$ is the decay constant of $k_{j,l}$;
only the electric plaquettes over the links of $S_{j,l}$ occur in this bound.
For fixed $j$ the sets $S_{j,l}$ are pairwise disjoint subsets of $\Lambda^{(1)}_{\mathrm{s}}$:
the link $(j,b)$, $b\in S_{j,l}$, lies in $B_{j,l}$, and the boxes $B_{j,l}$ and their translates by $\Gamma$ are pairwise disjoint.
By~\eqref{it:step-free}, the value $v_j(x)$ does not change when the coordinates $x_b$, $b\in\bigcup_lS_{j,l}$, are replaced by $0$; hence
\[
\left|v_j(x)\right|\le C\cdot\exp\left(\left(\tfrac12-\varepsilon_v\right)\sum_{b\notin\bigcup_lS_{j,l}}x_b^2\right).
\]
Let $\varepsilon>0$ be the least of the numbers $\varepsilon_v$, $\frac12$ and the constants $\varepsilon_k$ of the finitely many local step data.
A coefficient of a product in a Grassmann algebra is a finite sum of products of coefficients of the factors, and the functions $u_j$ are bounded.
The numbers $F^{\mathrm{el}}_j(b)$, $j\in\Z/n_1\Z$, $b\in\Lambda^{(1)}_{\mathrm{s}}$, are the values of $F$ at distinct plaquettes of $\mathbb{T}$,
and in the resulting bound for a coefficient of $G(A,z)$ the square $F^{\mathrm{el}}_j(b)^2$ occurs with the factor $\frac12-\min\left(\varepsilon_k,\frac12\right)\le\frac12-\varepsilon$ if $b\in S_{j,l}$, $\varepsilon_k$ being the constant of $k_{j,l}$,
and with the factor $\frac12-\varepsilon_v\le\frac12-\varepsilon$ if $b\notin\bigcup_lS_{j,l}$.
As $\frac12-\varepsilon\ge0$, every coefficient of $G(A,z)$ is therefore bounded by
\[
C\cdot\exp\left(\left(\tfrac12-\varepsilon\right)\sum_{j,b}F^{\mathrm{el}}_j(b)^2\right)\le C\cdot\exp\left(\left(\tfrac12-\varepsilon\right)\|F\|^2\right).
\]
Thus $G\in\mathcal{I}(\mathbb{T})$, with the constant $\varepsilon_G:=\varepsilon\in\left(0,\frac12\right]$ of Definition~\ref{def:torus-integrands}.

For $(A,z)\in\mathfrak{C}(\Gamma)$ let
\[
\Phi_G(A,z):=\exp\left(-\tfrac12\|F\|^2\right)\cdot\int_{\mathbb{T}}\exp\left(-S_{\mathbb{T}}(A)\right)\cdot G(A,z).
\]
The integer shifts add to $A$ an element of $(2\pi/q)\,C^1(\Gamma;\Z)$; they change neither $F$ nor the phases $\exp\left(\pm iqA(b)\right)$,
and the action $S_{\mathbb{T}}(A)$ depends on $A$ only through these phases (Definitions~\ref{def:chamber-operator} and~\ref{def:tori}).
Together with the invariance of $G$, this makes $\Phi_G$ a Borel function on $\mathfrak{C}(\Gamma)$ with the invariance required in Lemma~\ref{lem:unfolding}.
The Berezin integral $\int_{\mathbb{T}}$ takes the coefficient of the single monomial $\Psi^{\mathrm{top}}(\mathbb{X}_{\mathbb{T}})$ (Definition~\ref{def:left-derivatives}),
which up to a sign belongs to every monomial basis of $\Gr_{\mathbb{T}}$;
so, up to a sign that does not depend on $(A,z)$, the number $\Phi_G(A,z)$ is a coefficient of $\exp\left(-\frac12\|F\|^2\right)\cdot\exp\left(-S_{\mathbb{T}}(A)\right)\cdot G(A,z)$ in a monomial basis,
and Proposition~\ref{prop:berezin}\eqref{it:torus-integrability} shows that $\Phi_G$ is absolutely integrable over $\Conf(\Gamma)$.
By Definition~\ref{def:torus-integrands}, $Z_{C(n)}[G]$ is the integral of $\Phi_G$ over $\Conf(\Gamma)$,
and Lemma~\ref{lem:unfolding}, in its clause on complex integrable functions, gives
\[
Z_{C(n)}[G]=\sum_{k\in\Z}\int\mathrm{d}\mathsf{p}\,\Phi_G\left(\varpi(\mathsf{p},k)\right),
\]
the right side converging absolutely.

\emph{Step~2: the factors $G[k_{j,l}]$ along a path.}
Fix a path $\mathsf{p}$, with $A_{n_1}=A_0+g$ and $g\in\Gamma_{\mathrm{s}}$, a flux $k\in\Z$ and a pair $(j,l)$,
and put $S:=S_{j,l}$, $\Lambda:=\Lambda_{j,l}$ and $B:=B_{j,l}$.
Let $\mathcal{K}'$ be the complex of the cells of $\Z^3$ contained in $[j,j+1]\times\hull(\Lambda)$.
By Definition~\ref{def:insertions} the cube $B$ contains the cubes $(j,y)+\left[-\frac12,\frac12\right]^3$ and $(j+1,y)+\left[-\frac12,\frac12\right]^3$, $y\in\Lambda$.
Being a product of three intervals, $B$ therefore contains
\[
\left[j-\tfrac12,j+\tfrac32\right]\times\left(\hull(\Lambda)+\left[-\tfrac12,\tfrac12\right]^2\right),
\]
and in particular every cell of $\mathcal{K}'$.
As $B$ embeds in $\Gamma$, the cells of $\mathcal{K}'$ lie in pairwise different orbits of $\Gamma$;
and since the translates of $B+[-1,1]^3$ by the elements of $\Gamma$ are pairwise disjoint, while $B+[-1,1]^3$ contains $\{j\}\times\left(\hull(\Lambda)+[-1,1]^2\right)$,
the hull $\hull(\Lambda)$ embeds in the cross-section $\Xi$ of $\mathbb{T}$ (Definitions~\ref{def:slices} and~\ref{def:slice-notation}).

Let $\hat A^{\mathsf{p}}$ be the real one-cochain on the links of $\mathcal{K}'$ that takes the values of the path:
$A_j(b)$ on the link $(j,b)$, $A_{j+1}(b)$ on the link $(j+1,b)$, and $a_j(y)$ on the link from $(j,y)$ to $(j+1,y)$,
where $b$ and $y$ are the links and the sites of $\Z^2$ contained in $\hull(\Lambda)$, read in $\Lambda_{\mathrm{s}}$.
This cochain is a lift of $\varpi(\mathsf{p},k)$ on $\mathcal{K}'$:
it has the phases of $\varpi(\mathsf{p},k)$, and its coboundary is $F$ on the plaquettes of $\mathcal{K}'$.
The phases agree because $g$ takes its values in $(2\pi/q)\Z$.
On the electric plaquette of $\mathcal{K}'$ bounded by $(j,b)$, $(j+1,b)$ and the two links between them, the coboundary is
$\left(A_{j+1}-A_j-d_{\mathrm{s}}a_j\right)(b)=F^{\mathrm{el}}_j(b)$ (Lemma~\ref{lem:unfolding}).
On a magnetic plaquette, the square $\sigma\subset\hull(\Lambda)$ in a plane $\{x_1=s\}$ with $s\in\{j,j+1\}$, the coboundary is $\left(d_{\mathrm{s}}A_s\right)(\sigma)$,
whereas $F$ takes there the value $\left(d_{\mathrm{s}}A_s\right)(\sigma)+(2\pi/q)k\cdot1_{\sigma_0}(\sigma)$, by Lemma~\ref{lem:unfolding} and $d_{\mathrm{s}}A_{n_1}=d_{\mathrm{s}}A_0$.
This plaquette lies in $B$, so by~\eqref{it:step-string} it does not represent $\sigma_0$, and the two values agree.

Let $\hat A$ be a local lift of $\varpi(\mathsf{p},k)$ on $B$ (Definition~\ref{def:integrands}).
On the links of $\mathcal{K}'$ it has the same coboundary and the same phases as $\hat A^{\mathsf{p}}$;
by~\eqref{eq:slab-lift-difference}, applied to the product of intervals $[j,j+1]\times\hull(\Lambda)$,
there is an integer valued function $\lambda$ on the sites of $\mathcal{K}'$ with $\hat A=\hat A^{\mathsf{p}}+(2\pi/q)\,d\lambda$ on the links of $\mathcal{K}'$.
Extend $\lambda$ by $0$ to the other sites of $B$.
Then $\hat A-(2\pi/q)\,d\lambda$ is again a local lift on $B$, with the same coboundary and the same phases,
and it coincides with $\hat A^{\mathsf{p}}$ on the links of $\mathcal{K}'$.
The integrand $G[k_{j,l}]$ is gauge invariant (Definition~\ref{def:integrands});
for the real function $(2\pi/q)\lambda$ the transformation of the fields is trivial, since $\exp\left(iq\cdot(2\pi/q)\lambda\right)=1$,
so that $G[k_{j,l}]$ takes the same value at the two lifts $\hat A$ and $\hat A-(2\pi/q)\,d\lambda$.
Moreover, $G[k_{j,l}]$ reads only the links $(j,b)$ and $(j+1,b)$, $b\in S$, and the links from $(j,y)$ to $(j+1,y)$, $y\in\Lambda$,
all of which belong to $\mathcal{K}'$, because $\Lambda$ contains the end points of the links of $S$.
Hence the value of $G[k_{j,l}]$ at $\varpi(\mathsf{p},k)$ is given by its defining formula with the slices of $\hat A^{\mathsf{p}}$:
in $\Gr_{\mathbb{T}}$,
\[
G[k_{j,l}]=\Ins^{\mathcal{J}^{\mathrm{l}}(\Lambda)}\left[\omega^{\frac12}\,\Upsilon_{j,l}\,u(a_j)\,\omega^{\frac12}\right]\cdot\prod_{b\in S}g_1\left(F^{\mathrm{el}}_j(b)\right)^{-1}\cdot\exp\left(S^{\times}_{j+\frac12}(\Lambda)\right),
\]
where the insertion is formed with the modes over $\Lambda$ at the plane $\{x_1=j+\frac12\}$,
$F^{\mathrm{el}}_j=A_{j+1}-A_j-d_{\mathrm{s}}a_j$, the operator $u(a_j)$ is formed with the restriction of $a_j$ to $\Lambda$, and
\[
\Upsilon_{j,l}:=k_{j,l}\left(A_{j+1},A_j+d_{\mathrm{s}}a_j\right)
\]
is an even operator on $\mathcal{F}\left(\mathcal{J}^{\mathrm{l}}(\Lambda)\right)$, the restrictions to $S$ being understood as in~\eqref{it:word-identity}.

\emph{Step~3: the fermions.}
Keep $\mathsf{p}$, $k$ and $j$ fixed, and let $A\in C^1(\Gamma)$ be the one-cochain of $\varpi(\mathsf{p},k)$;
by Lemma~\ref{lem:unfolding} its values on the links from the plane $\{x_1=j\}$ to the plane $\{x_1=j+1\}$ are those of $a_j$.
The images of the sets $\Lambda_{j,l}$ in $\Lambda_{\mathrm{s}}$ are pairwise disjoint,
since the cubes $(j,y)+\left[-\frac12,\frac12\right]^3$, $y\in\Lambda_{j,l}$, lie in $B_{j,l}$,
and the boxes $B_{j,l}$ and their translates by $\Gamma$ are pairwise disjoint;
the hulls of these sets embed in the cross-section, by Step~2.
Let
\[
\Upsilon_j:=\left(\prod_l\iota\left(\Upsilon_{j,l}\right)\right)u(a_j)\in\End\mathcal{F},
\]
the localisations being formed with the signed relabelling, as in Remark~\ref{rem:cross-section-localisation};
the factors $\iota(\Upsilon_{j,l})$ commute by Proposition~\ref{prop:localisation}\eqref{it:localisation-commute}, and $\Upsilon_j$ is even.
At a step without local step data the product over $l$ is empty, and $\Upsilon_j=u(a_j)$.
Let $\Ins_{j+\frac12}$ be the insertion at the plane $\{x_1=j+\frac12\}$ formed with $\mathcal{J}^{\mathrm{l}}(\Lambda_{\mathrm{s}})$,
and let $S^{\times}_{j+\frac12}(A)$ be the crossing part of the action at this plane, formed with the cochain $A$ (Lemma~\ref{lem:slab-bonds}\eqref{it:bonds-action-split});
as in Proposition~\ref{prop:slice-factorisation}, the insertions and $S^{\times}_{j+\frac12}(A)$ are read in $\Gr_{\mathbb{T}}$.
The phases $u_i=u_i(A)$ of the modes at this plane are therefore formed with $a_j$, and $\diag(u_I)=u(a_j)$ (Definition~\ref{def:slice-notation}).
By Proposition~\ref{prop:crossing} and~\eqref{eq:slab-crossing-factor},
$\exp\left(-S^{\times}_{j+\frac12}(A)\right)$ is the product of the elements $1+\omega_i\cdot u_i\cdot\Theta_{j+\frac12}(\zeta_i)\,\zeta_i$ over all $i\in\mathcal{J}^{\mathrm{l}}(\Lambda_{\mathrm{s}})$,
and $\exp\left(-S^{\times}_{j+\frac12}(\Lambda_{j,l})\right)$, formed with the cochain $\hat A^{\mathsf{p}}$ of Step~2 for the pair $(j,l)$, is the product of the same elements over the modes over $\Lambda_{j,l}$,
an element being unchanged when $\zeta_i$ is replaced by $-\zeta_i$.
These elements are even, commute with one another and are invertible.
In $\Gr_{\mathbb{T}}$ one has
\[
\begin{aligned}
\Ins_{j+\frac12}\left[\omega^{\frac12}\,\Upsilon_j\,\omega^{\frac12}\right]
&=\prod_l\Ins^{\mathcal{J}^{\mathrm{l}}(\Lambda_{j,l})}\left[\omega^{\frac12}\,\Upsilon_{j,l}\,u(a_j)\,\omega^{\frac12}\right]\\
&\qquad\cdot\prod_i\left(1+\omega_i\cdot u_i\cdot\Theta_{j+\frac12}(\zeta_i)\,\zeta_i\right),
\end{aligned}
\]
where $i$ ranges over the modes of $\mathcal{J}^{\mathrm{l}}(\Lambda_{\mathrm{s}})$ that do not represent a mode over one of the sets $\Lambda_{j,l}$.
If there are local step data at the step $j$, this is Proposition~\ref{prop:localisation}\eqref{it:several-insertions}, read on $\mathbb{T}$ as in Remark~\ref{rem:cross-section-localisation}.
If there are none, the first product on the right is empty and $i$ ranges over all modes;
then $\omega^{\frac12}\,u(a_j)\,\omega^{\frac12}=\omega\,u(a_j)$ is the diagonal operator with the entries $\omega_I\cdot u_I$ (Definition~\ref{def:slice-notation}),
and by~\eqref{eq:slab-crossing-factor} both sides equal $\exp\left(-S^{\times}_{j+\frac12}(A)\right)$.
In either case the last product in the display is $\exp\left(-S^{\times}_{j+\frac12}(A)\right)\prod_l\exp\left(S^{\times}_{j+\frac12}(\Lambda_{j,l})\right)$,
and Step~2 gives
\[
\prod_lG[k_{j,l}]=\exp\left(S^{\times}_{j+\frac12}(A)\right)\Ins_{j+\frac12}\left[\omega^{\frac12}\,\Upsilon_j\,\omega^{\frac12}\right]\cdot\prod_{b\in\bigcup_lS_{j,l}}g_1\left(F^{\mathrm{el}}_j(b)\right)^{-1}.
\]
Now let $j$ vary; the elements so obtained are even and commute.
Proposition~\ref{prop:slice-factorisation}, applied to $A$ and to the even operators $\Upsilon_j$,
together with the periodicity of $\widehat{\mathfrak{N}}$ under $(2\pi/q)\Z^{\Lambda^{(1)}_{\mathrm{s}}}\supset\Gamma_{\mathrm{s}}$ (Lemma~\ref{lem:half-slab-matrices}), which gives $\widehat{\mathfrak{N}}(A_0+g)=\widehat{\mathfrak{N}}(A_0)$, yields
\[
\int_{\mathbb{T}}\exp\left(-S_{\mathbb{T}}(A)\right)\prod_{j,l}G[k_{j,l}]
=\prod_j\prod_{b\in\bigcup_lS_{j,l}}g_1\left(F^{\mathrm{el}}_j(b)\right)^{-1}\cdot\operatorname{tr}_{\mathcal{F}'}\prod_{j=n_1-1,\dots,0}\left(\widehat{\mathfrak{N}}(A_{j+1})^{*}\,\Upsilon_j\,\widehat{\mathfrak{N}}(A_j)\right).
\]

\emph{Step~4: the Gaussian weight.}
The plaquettes of $\mathbb{T}$ are the magnetic ones, at which $F$ takes the values of the $F^{\mathrm{mag}}_j$,
and the electric ones, at which it takes the values $F^{\mathrm{el}}_j(b)$, $b\in\Lambda^{(1)}_{\mathrm{s}}$;
there are $n_1\cdot2|\Lambda_{\mathrm{s}}|=2|n|$ electric plaquettes.
Since $\exp\left(-\frac12s^2\right)=(2\pi)^{\frac12}g_1(s)$,
\[
\exp\left(-\tfrac12\|F\|^2\right)=(2\pi)^{|n|}\cdot\prod_j\exp\left(-\tfrac12\left\|F^{\mathrm{mag}}_j\right\|^2\right)\cdot\prod_{j,b}g_1\left(F^{\mathrm{el}}_j(b)\right).
\]
At $\varpi(\mathsf{p},k)$ one has $F^{\mathrm{mag}}_j=F^{\mathrm{mag}}(A_j,k)$ and $F^{\mathrm{el}}_j=A_{j+1}-A_j-d_{\mathrm{s}}a_j$ (Lemma~\ref{lem:unfolding}).
For $\beta\in\R^{\Lambda^{(2)}_{\mathrm{s}}}$ and $A',A\in\R^{\Lambda^{(1)}_{\mathrm{s}}}$ put
\[
\begin{aligned}
\mathcal{W}_j(\beta)&:=\exp\left(-\tfrac12\|\beta\|^2\right)\cdot u_j(\beta),\\
\widetilde K_j(A',A)&:=v_j(A'-A)\cdot\prod_{b\in\Lambda^{(1)}_{\mathrm{s}}\setminus\bigcup_lS_{j,l}}g_1\left(A'(b)-A(b)\right)\cdot\prod_l\iota\left(k_{j,l}(A',A)\right)\in\End\mathcal{F}.
\end{aligned}
\]
The factors $g_1\left(F^{\mathrm{el}}_j(b)\right)$, $b\in\bigcup_lS_{j,l}$, of the Gaussian weight cancel against those of Step~3, and
\[
v_j\left(F^{\mathrm{el}}_j\right)\cdot\prod_{b\in\Lambda^{(1)}_{\mathrm{s}}\setminus\bigcup_lS_{j,l}}g_1\left(F^{\mathrm{el}}_j(b)\right)\cdot\Upsilon_j=\widetilde K_j\left(A_{j+1},A_j+d_{\mathrm{s}}a_j\right)u(a_j),
\]
because $A_{j+1}-\left(A_j+d_{\mathrm{s}}a_j\right)=F^{\mathrm{el}}_j$.
Multiplying the result of Step~3 by the Gaussian weight and by the scalars $u_j\left(F^{\mathrm{mag}}_j\right)$ and $v_j\left(F^{\mathrm{el}}_j\right)$ therefore gives
\begin{equation}\label{eq:slab-path-integrand}
\begin{split}
\Phi_G\left(\varpi(\mathsf{p},k)\right)
&=(2\pi)^{|n|}\cdot\prod_{j=0}^{n_1-1}\mathcal{W}_j\left(F^{\mathrm{mag}}(A_j,k)\right)\\
&\qquad\cdot\operatorname{tr}_{\mathcal{F}'}\prod_{j=n_1-1,\dots,0}\left(\widehat{\mathfrak{N}}(A_{j+1})^{*}\,\widetilde K_j\left(A_{j+1},A_j+d_{\mathrm{s}}a_j\right)u(a_j)\,\widehat{\mathfrak{N}}(A_j)\right).
\end{split}
\end{equation}

\emph{Step~5: the trace over paths.}
The functions $\mathcal{W}_j$ and $\widetilde K_j$ satisfy the hypotheses of Lemma~\ref{lem:path-trace}, and $n_1\ge4$.
They are Borel, and $\left|\mathcal{W}_j(\beta)\right|\le\sup\left|u_j\right|\cdot\exp\left(-\frac12\|\beta\|^2\right)\le\sup\left|u_j\right|\cdot\exp\left(-\varepsilon\|\beta\|^2\right)$, with $\varepsilon\le\frac12$ as in Step~1.
With $x:=A'-A$, the bound on $v_j$ in Step~1 gives
\[
\left|v_j(x)\right|\cdot\prod_{b\in\Lambda^{(1)}_{\mathrm{s}}\setminus\bigcup_lS_{j,l}}g_1(x_b)\le C\cdot\exp\left(-\varepsilon_v\sum_{b\notin\bigcup_lS_{j,l}}x_b^2\right),
\]
and $\left\|\iota\left(k_{j,l}(A',A)\right)\right\|=\left\|k_{j,l}(A',A)\right\|\le C\cdot\prod_{b\in S_{j,l}}\exp\left(-\varepsilon_kx_b^2\right)$, the localisation being isometric (Proposition~\ref{prop:localisation}\eqref{it:localisation-homomorphism}).
Hence $\left\|\widetilde K_j(A',A)\right\|\le C\cdot\exp\left(-\varepsilon|A'-A|^2\right)$, with the same $\varepsilon$.

Let now $g\in\Gamma_{\mathrm{s}}$ be arbitrary.
Then $\frac{q}{2\pi}g$ is an integer one-cochain of $\Lambda_{\mathrm{s}}$ with $d_{\mathrm{s}}g=0$.
On the links of the complex of the cells contained in $\{j\}\times\hull(\Lambda_{j,l})$, read in $\Lambda_{\mathrm{s}}$,
the cochains $0$ and $g$ have the same coboundary, namely $0$, and the same phases;
by~\eqref{eq:slab-lift-difference} there is an integer valued function $a$ on the sites of $\hull(\Lambda_{j,l})$ with $g=(2\pi/q)\,d_{\mathrm{s}}a$ on the links of $\hull(\Lambda_{j,l})$.
The links of $S_{j,l}$ lie in $\hull(\Lambda_{j,l})$ and have their end points in $\Lambda_{j,l}$.
The gauge covariance of $k_{j,l}$, applied to the restriction of $(2\pi/q)a$ to $\Lambda_{j,l}$,
and $u\left((2\pi/q)a\right)=1$, the charges being integers (Lemma~\ref{lem:half-slab-matrices}), give
\[
k_{j,l}\left(A'+g,A+g\right)=k_{j,l}\left(A',A\right).
\]
The other factors of $\widetilde K_j$ depend on $A'-A$ only, so that $\widetilde K_j(A'+g,A+g)=\widetilde K_j(A',A)$.

The operator of Lemma~\ref{lem:path-trace} with the kernel $\sum_{g\in\Gamma_{\mathrm{s}}}\widetilde K_j(A'+g,A)$ is the operator $K_j$ of~\eqref{it:word-identity}.
The multiplication operator $\mathfrak{T}_j$ of that lemma acts in the flux sector $k$ as the multiplication by
$\mathcal{W}_j\left(F^{\mathrm{mag}}(A,k)\right)\cdot\widehat{\mathfrak{N}}(A)\widehat{\mathfrak{N}}(A)^{*}$.
Since $\mathfrak{M}$ acts in this sector as the multiplication by $\exp\left(-\frac14\left\|F^{\mathrm{mag}}(A,k)\right\|^2\right)\cdot\widehat{\mathfrak{N}}(A)$ (Definition~\ref{def:slice-notation})
and $u_j$ as the multiplication by $u_j\left(F^{\mathrm{mag}}(A,k)\right)$, one has $\mathfrak{T}_j=\mathfrak{M}u_j\mathfrak{M}^{*}$.
Step~1, \eqref{eq:slab-path-integrand} and Lemma~\ref{lem:path-trace} now give
\[
Z_{C(n)}[G]=\sum_{k\in\Z}\int\mathrm{d}\mathsf{p}\,\Phi_G\left(\varpi(\mathsf{p},k)\right)
=(2\pi)^{|n|}\cdot\Tr_{\mathcal{H}^{\mathrm{l}}}\prod_{j=n_1-1,\dots,0}\left((2\pi/q)^{|\Lambda_{\mathrm{s}}|}K_jP_{\mathrm{G}}\,\mathfrak{M}u_j\mathfrak{M}^{*}\right),
\]
and every factor under the trace is Hilbert--Schmidt.

\emph{The special cases.}
In the first special case of~\eqref{it:word-cases}, the kernel of $K_j$ is $f^{\mathrm{per}}([A'-A])$ with $f(x)=\exp\left(i\langle\beta_j,x\rangle\right)\cdot\prod_bg_1(x_b)$;
by Lemma~\ref{lem:slice-kernels}\eqref{it:gaussian-kernels}, $K_j=\exp\left(-\frac12|E-\beta_j|^2\right)$,
and $(2\pi/q)^{|\Lambda_{\mathrm{s}}|}K_jP_{\mathrm{G}}=K_{\beta_j}$ by Definition~\ref{def:slice-operators}.
In the second, the kernel of $K_j$ is $f^{\mathrm{per}}([A'-A])$ with $f(x)=\exp\left(i\langle\beta,x\rangle\right)\cdot\prod_{b\in S}\left(\sqrt{2}\cdot\exp\left(-\frac{x_b^2}{2}\right)\right)\cdot\prod_bg_1(x_b)$;
the same item gives $K_j=\exp\left(-\frac12|E-\beta|^2+\frac14|E-\beta|_S^2\right)$,
and since $P_{\mathrm{G}}$ commutes with the functions of $E$ (Lemma~\ref{lem:slice-kernels}\eqref{it:gauge-average}),
\[
(2\pi/q)^{|\Lambda_{\mathrm{s}}|}K_jP_{\mathrm{G}}=K_\beta\exp\left(\tfrac14|E-\beta|_S^2\right)=K_{\beta,S}.
\]
In the third, the conditions~\eqref{it:step-disjoint}--\eqref{it:step-free} are void;
as the plaquette $\sigma_0$ enters the hypotheses of~\eqref{it:word-identity} only through~\eqref{it:step-string}, the identity holds for every choice of $\sigma_0$.
For the fourth, apply~\eqref{it:word-identity} with $u_j=1$, $v_j=1$ and without local step data.
By the first case with $\beta_j=0$, one has $(2\pi/q)^{|\Lambda_{\mathrm{s}}|}K_jP_{\mathrm{G}}=K$, so that
\[
Z_{C(n)}[1]=(2\pi)^{|n|}\cdot\Tr\left(K\mathfrak{M}\mathfrak{M}^{*}\right)^{n_1}.
\]
By Corollary~\ref{cor:trace-class}, $K^{1/2}\mathfrak{M}$ is Hilbert--Schmidt and $K^{1/2}\mathfrak{M}\mathfrak{M}^{*}K^{1/2}$ is of trace class;
moreover $K=K^{1/2}K^{1/2}$, and $\Tr(XY)=\Tr(YX)$ whenever $X$ and $Y$ are Hilbert--Schmidt, or $X$ is of trace class and $Y$ is bounded.
In $\left(K\mathfrak{M}\mathfrak{M}^{*}\right)^{n_1}$, move the first factor $K^{1/2}$ to the end;
the product of the remaining factors begins with $K^{1/2}\mathfrak{M}\mathfrak{M}^{*}K^{1/2}$, as $n_1\ge2$, and is of trace class.
Then, in $\left(K^{1/2}\mathfrak{M}\mathfrak{M}^{*}K^{1/2}\right)^{n_1}$, move the first factor $K^{1/2}\mathfrak{M}$ to the end.
This gives
\[
\Tr\left(K\mathfrak{M}\mathfrak{M}^{*}\right)^{n_1}=\Tr\left(K^{1/2}\mathfrak{M}\mathfrak{M}^{*}K^{1/2}\right)^{n_1}=\Tr\left(\mathfrak{M}^{*}K\mathfrak{M}\right)^{n_1}.
\]
\end{proof}

\begin{proposition}[External field]\label{prop:external-field}
Let $w_0,\dots,w_3>0$, $M\in\R$, $q>0$ and $-1\le\rho\le1$.
Let $n_1,n_2,n_3\ge4$ and $n=(n_1,n_2,n_3)$, let $\Gamma$ be the period lattice of $C(n)$, and let $\beta\in C^2(\Gamma)$ be a real two-cochain on $C(n)$;
the pairing $\langle\beta,F\rangle$ and the norm $\|\beta\|^2=\langle\beta,\beta\rangle$ are formed over the plaquettes of $C(n)$.
The cross-section of $C(n)$ for $x_1$, with the slice torus $\Lambda_{\mathrm{s}}$, the spaces $\mathcal{H}^{\mathrm{l}}$ and $\mathcal{H}^{\mathrm{s}}$, the operators $E$, $K^{1/2}$ and $\mathfrak{M}$,
the matrices $\widehat{\mathfrak{N}}(A)$ and the magnetic curvature $F^{\mathrm{mag}}(A,k)$, is that of Definition~\ref{def:slice-notation}.
For $j\in\Z/n_1\Z$ let $\beta^{\mathrm{el}}_j:=F^{\mathrm{el}}_j(\beta)\in\R^{\Lambda^{(1)}_{\mathrm{s}}}$ and $\beta^{\mathrm{mag}}_j:=F^{\mathrm{mag}}_j(\beta)\in\R^{\Lambda^{(2)}_{\mathrm{s}}}$ be the electric and the magnetic part of $\beta$,
formed as in Definition~\ref{def:slices}:
\[
\beta^{\mathrm{el}}_j(y,\nu)=\beta_{1\nu}(j,y)\quad(\nu\in\{2,3\}),
\qquad
\beta^{\mathrm{mag}}_j(y)=\beta_{23}(j,y).
\]
For a real two-cochain $\beta'\in\R^{\Lambda^{(2)}_{\mathrm{s}}}$ of the slice torus,
let $\mathfrak{M}_{(\beta')}:\mathcal{H}^{\mathrm{s}}\to\mathcal{H}^{\mathrm{l}}$ be the operator that acts in the flux sector $k$ as the multiplication by
\[
\exp\left(-\tfrac14\left\|F^{\mathrm{mag}}(A,k)+\beta'\right\|^2\right)\cdot\widehat{\mathfrak{N}}(A);
\]
thus $\mathfrak{M}_{(0)}=\mathfrak{M}$.
Then
\[
\begin{aligned}
&Z_{C(n)}\left[\exp\left(-\langle\beta,F\rangle-\tfrac12\|\beta\|^2\right)\right]\\
&\qquad=(2\pi)^{|n|}\cdot\Tr_{\mathcal{H}^{\mathrm{s}}}\prod_{j=n_1-1,\dots,0}\left(\mathfrak{M}_{(\beta^{\mathrm{mag}}_{j+1})}^{*}\,K^{1/2}\exp\left(i\langle\beta^{\mathrm{el}}_j,E\rangle\right)K^{1/2}\,\mathfrak{M}_{(\beta^{\mathrm{mag}}_j)}\right),
\end{aligned}
\]
where the indices of the parts of $\beta$ are read modulo $n_1$, so that $\beta^{\mathrm{mag}}_{n_1}=\beta^{\mathrm{mag}}_0$, and the factor with $j=n_1-1$ stands on the left.
Every operator $K^{1/2}\mathfrak{M}_{(\beta^{\mathrm{mag}}_j)}:\mathcal{H}^{\mathrm{s}}\to\mathcal{H}^{\mathrm{l}}$ is Hilbert--Schmidt.
\end{proposition}

\begin{proof}
Let $\mathbb{T}:=C(n)$, and let $F=dA+(2\pi/q)z$ be the curvature of $(A,z)\in\mathfrak{C}(\Gamma)$.
For every $\varepsilon>0$ one has $\left|\langle\beta,F\rangle\right|\le\frac{1}{4\varepsilon}\|\beta\|^2+\varepsilon\|F\|^2$, and hence
\[
\exp\left(-\langle\beta,F\rangle\right)\le\exp\left(\tfrac{1}{4\varepsilon}\|\beta\|^2\right)\cdot\exp\left(\varepsilon\|F\|^2\right).
\]
Taking $0<\varepsilon<\frac12$, one sees that $G:=\exp\left(-\langle\beta,F\rangle-\frac12\|\beta\|^2\right)$,
a Borel function of $F$ and therefore invariant under the integer shifts, belongs to $\mathcal{I}(\mathbb{T})$, with $\varepsilon_G=\frac12-\varepsilon$.
For $(A,z)\in\mathfrak{C}(\Gamma)$ let
\[
\begin{aligned}
\Phi_G(A,z)&:=\exp\left(-\tfrac12\|F\|^2\right)\cdot\int_{\mathbb{T}}\exp\left(-S_{\mathbb{T}}(A)\right)\cdot G(A,z)\\
&=\exp\left(-\tfrac12\|F+\beta\|^2\right)\cdot\int_{\mathbb{T}}\exp\left(-S_{\mathbb{T}}(A)\right),
\end{aligned}
\]
the two expressions agreeing by completion of the square, $\|F\|^2+2\langle\beta,F\rangle+\|\beta\|^2=\|F+\beta\|^2$.
Under the integer shifts of Definition~\ref{def:torus-integrands} neither $F$ nor the phases $\exp\left(\pm iqA(b)\right)$ change,
and $S_{\mathbb{T}}(A)$ depends on $A$ only through these phases (Definitions~\ref{def:chamber-operator} and~\ref{def:tori});
so $\Phi_G$ is a Borel function on $\mathfrak{C}(\Gamma)$ with the invariance required in Lemma~\ref{lem:unfolding}.
Since $\int_{\mathbb{T}}$ is the coefficient of $\Psi^{\mathrm{top}}(\mathbb{X}_{\mathbb{T}})$ (Definition~\ref{def:left-derivatives}), which up to a sign is an element of a monomial basis of $\Gr_{\mathbb{T}}$,
the number $\Phi_G(A,z)$ is, up to a sign independent of $(A,z)$, a coefficient of $\exp\left(-\frac12\|F\|^2\right)\cdot\exp\left(-S_{\mathbb{T}}(A)\right)\cdot G(A,z)$ in that basis;
hence $\Phi_G$ is absolutely integrable over $\Conf(\Gamma)$ by Proposition~\ref{prop:berezin}\eqref{it:torus-integrability}.
As $Z_{C(n)}[G]$ is the integral of $\Phi_G$ over $\Conf(\Gamma)$ (Definition~\ref{def:torus-integrands}),
Lemma~\ref{lem:unfolding}, applied to the complex integrable function $\Phi_G$, yields
$Z_{C(n)}[G]=\sum_{k\in\Z}\int\mathrm{d}\mathsf{p}\,\Phi_G\left(\varpi(\mathsf{p},k)\right)$, the right side converging absolutely.

Fix a path $\mathsf{p}$, with $A_{n_1}=A_0+g$, and a flux $k$.
At $\varpi(\mathsf{p},k)$ the magnetic and the electric parts of $F+\beta$ are $F^{\mathrm{mag}}(A_j,k)+\beta^{\mathrm{mag}}_j$ and $A_{j+1}-A_j-d_{\mathrm{s}}a_j+\beta^{\mathrm{el}}_j$ (Lemma~\ref{lem:unfolding}).
The torus $\mathbb{T}$ has $2|n|$ electric plaquettes, and $\exp\left(-\frac12s^2\right)=(2\pi)^{\frac12}g_1(s)$; hence
\[
\begin{aligned}
\exp\left(-\tfrac12\|F+\beta\|^2\right)
&=(2\pi)^{|n|}\cdot\prod_j\exp\left(-\tfrac12\left\|F^{\mathrm{mag}}(A_j,k)+\beta^{\mathrm{mag}}_j\right\|^2\right)\\
&\qquad\cdot\prod_{j,b}g_1\left(A_{j+1}(b)-A_j(b)-\left(d_{\mathrm{s}}a_j\right)(b)+\beta^{\mathrm{el}}_j(b)\right).
\end{aligned}
\]
Proposition~\ref{prop:slice-factorisation} with $\Upsilon_j=u(a_j)$, for which its left side is $\int_{\mathbb{T}}\exp\left(-S_{\mathbb{T}}(A)\right)$ by~\eqref{eq:slab-crossing-factor},
together with the periodicity of $\widehat{\mathfrak{N}}$ (Lemma~\ref{lem:half-slab-matrices}), gives
\[
\int_{\mathbb{T}}\exp\left(-S_{\mathbb{T}}(A)\right)=\operatorname{tr}_{\mathcal{F}'}\prod_{j=n_1-1,\dots,0}\left(\widehat{\mathfrak{N}}(A_{j+1})^{*}\,u(a_j)\,\widehat{\mathfrak{N}}(A_j)\right).
\]
For $\beta'\in\R^{\Lambda^{(2)}_{\mathrm{s}}}$ and $A',A\in\R^{\Lambda^{(1)}_{\mathrm{s}}}$ put
\[
\mathcal{W}_j(\beta'):=\exp\left(-\tfrac12\left\|\beta'+\beta^{\mathrm{mag}}_j\right\|^2\right),
\qquad
\widetilde K_j(A',A):=\prod_{b\in\Lambda^{(1)}_{\mathrm{s}}}g_1\left(A'(b)-A(b)+\beta^{\mathrm{el}}_j(b)\right)\cdot1_{\mathcal{F}},
\]
with the identity $1_{\mathcal{F}}$ of $\mathcal{F}$.
Inserting the product formula for $\exp\left(-\frac12\|F+\beta\|^2\right)$ and the trace formula for $\int_{\mathbb{T}}\exp\left(-S_{\mathbb{T}}(A)\right)$ into the second expression for $\Phi_G$,
one finds that $\Phi_G\left(\varpi(\mathsf{p},k)\right)$ is given by the right side of~\eqref{eq:slab-path-integrand} with these functions $\mathcal{W}_j$ and $\widetilde K_j$.
They are Borel, and $\widetilde K_j$ depends on $A'-A$ only.
The inequality $|x|^2\le2|x+y|^2+2|y|^2$, valid for every Euclidean norm, gives $\exp\left(-\frac12|x+y|^2\right)\le\exp\left(\frac12|y|^2\right)\cdot\exp\left(-\frac14|x|^2\right)$.
Applied with the norm $\|\cdot\|$ of $\R^{\Lambda^{(2)}_{\mathrm{s}}}$, and with the norm of $\R^{\Lambda^{(1)}_{\mathrm{s}}}$ together with $\prod_bg_1(x_b)=(2\pi)^{-|\Lambda_{\mathrm{s}}|}\cdot\exp\left(-\frac12|x|^2\right)$ over the $2|\Lambda_{\mathrm{s}}|$ links of $\Lambda_{\mathrm{s}}$, it yields
\[
\begin{aligned}
\left|\mathcal{W}_j(\beta')\right|&\le\exp\left(\tfrac12\left\|\beta^{\mathrm{mag}}_j\right\|^2\right)\cdot\exp\left(-\tfrac14\|\beta'\|^2\right),\\
\left\|\widetilde K_j(A',A)\right\|&\le(2\pi)^{-|\Lambda_{\mathrm{s}}|}\cdot\exp\left(\tfrac12\left|\beta^{\mathrm{el}}_j\right|^2\right)\cdot\exp\left(-\tfrac14|A'-A|^2\right).
\end{aligned}
\]
As $n_1\ge4$, Lemma~\ref{lem:path-trace} applies with $\varepsilon=\frac14$ and gives
\[
Z_{C(n)}[G]=(2\pi)^{|n|}\cdot\Tr_{\mathcal{H}^{\mathrm{l}}}\prod_{j=n_1-1,\dots,0}\left((2\pi/q)^{|\Lambda_{\mathrm{s}}|}K_jP_{\mathrm{G}}\mathfrak{T}_j\right),
\]
where $K_j$ is the operator whose kernel is the periodisation of $\widetilde K_j$ and $\mathfrak{T}_j$ is the multiplication operator of that lemma.

The kernel of $K_j$ is $f^{\mathrm{per}}([A'-A])$ with $f(x):=\prod_bg_1\left(x_b+\beta^{\mathrm{el}}_j(b)\right)$.
Lemma~\ref{lem:slice-kernels}\eqref{it:electric-convolution}, the substitution $x\mapsto x-\beta^{\mathrm{el}}_j$ and Lemma~\ref{lem:slice-kernels}\eqref{it:gaussian-kernels} give
\[
K_j=\int\mathrm{d}x\,\prod_bg_1(x_b)\cdot\exp\left(-i\langle x-\beta^{\mathrm{el}}_j,E\rangle\right)=\exp\left(-\tfrac12|E|^2\right)\exp\left(i\langle\beta^{\mathrm{el}}_j,E\rangle\right).
\]
The functions of $E$ commute with each other and with $P_{\mathrm{G}}$ (Lemma~\ref{lem:slice-kernels}\eqref{it:gauge-average}), and $K=K^{1/2}K^{1/2}$; hence
\[
(2\pi/q)^{|\Lambda_{\mathrm{s}}|}K_jP_{\mathrm{G}}=K\exp\left(i\langle\beta^{\mathrm{el}}_j,E\rangle\right)=K^{1/2}\exp\left(i\langle\beta^{\mathrm{el}}_j,E\rangle\right)K^{1/2}.
\]
The operator $\mathfrak{T}_j$ acts in the flux sector $k$ as the multiplication by
$\exp\left(-\frac12\left\|F^{\mathrm{mag}}(A,k)+\beta^{\mathrm{mag}}_j\right\|^2\right)\cdot\widehat{\mathfrak{N}}(A)\widehat{\mathfrak{N}}(A)^{*}$,
that is, $\mathfrak{T}_j=\mathfrak{M}_{(\beta^{\mathrm{mag}}_j)}\mathfrak{M}_{(\beta^{\mathrm{mag}}_j)}^{*}$.

Each $K^{1/2}\mathfrak{M}_{(\beta^{\mathrm{mag}}_j)}$ is Hilbert--Schmidt.
Indeed, $K^{1/2}=(2\pi/q)^{\frac{|\Lambda_{\mathrm{s}}|}{2}}\exp\left(-\frac14|E|^2\right)P_{\mathrm{G}}$ (Definition~\ref{def:slice-operators}),
and by Lemma~\ref{lem:slice-kernels}\eqref{it:gaussian-kernels} the operator $\exp\left(-\frac14|E|^2\right)$ acts in every flux sector by the kernel $k^{+}([A'-A])$,
where $k^{+}:=(k^{+}_0)^{\mathrm{per}}$ and $k^{+}_0(x):=\prod_bg_{1/2}(x_b)=\pi^{-|\Lambda_{\mathrm{s}}|}\cdot\exp\left(-|x|^2\right)$.
With the integrals over $\R^{\Lambda^{(1)}_{\mathrm{s}}}/\Gamma_{\mathrm{s}}$,
\[
\begin{aligned}
&\sum_{k\in\Z}\int\mathrm{d}A\,\left\|\exp\left(-\tfrac14\left\|F^{\mathrm{mag}}(A,k)+\beta^{\mathrm{mag}}_j\right\|^2\right)\cdot\widehat{\mathfrak{N}}(A)\right\|^2\\
&\qquad\le\sup_A\left\|\widehat{\mathfrak{N}}(A)\right\|^2\cdot\sum_{k\in\Z}\int\mathrm{d}A\,\exp\left(-\tfrac12\left\|F^{\mathrm{mag}}(A,k)+\beta^{\mathrm{mag}}_j\right\|^2\right)<\infty,
\end{aligned}
\]
by Lemma~\ref{lem:half-slab-matrices} and Lemma~\ref{lem:slice-kernels}\eqref{it:flux-summability}.
Lemma~\ref{lem:slice-kernels}\eqref{it:kernel-hilbert-schmidt}, with $K'=\exp\left(-\frac14|E|^2\right)$ and $\mathfrak{T}=\mathfrak{M}_{(\beta^{\mathrm{mag}}_j)}$,
shows that $\exp\left(-\frac14|E|^2\right)P_{\mathrm{G}}\mathfrak{M}_{(\beta^{\mathrm{mag}}_j)}$ is Hilbert--Schmidt, and so is $K^{1/2}\mathfrak{M}_{(\beta^{\mathrm{mag}}_j)}$.

It remains to rearrange the trace.
In the product
\[
\prod_{j=n_1-1,\dots,0}\left(K^{1/2}\exp\left(i\langle\beta^{\mathrm{el}}_j,E\rangle\right)K^{1/2}\,\mathfrak{M}_{(\beta^{\mathrm{mag}}_j)}\mathfrak{M}_{(\beta^{\mathrm{mag}}_j)}^{*}\right)
\]
the last factor is the bounded operator $\mathfrak{M}_{(\beta^{\mathrm{mag}}_0)}^{*}:\mathcal{H}^{\mathrm{l}}\to\mathcal{H}^{\mathrm{s}}$.
The product of all the other factors is of trace class:
it contains the Hilbert--Schmidt operator $K^{1/2}\mathfrak{M}_{(\beta^{\mathrm{mag}}_{n_1-1})}$ followed by its adjoint $\mathfrak{M}_{(\beta^{\mathrm{mag}}_{n_1-1})}^{*}K^{1/2}$,
whose factor $K^{1/2}$ is the first factor of the term with $j=n_1-2\ge0$.
By the cyclicity of the trace, $\mathfrak{M}_{(\beta^{\mathrm{mag}}_0)}^{*}$ may be moved to the front, and the trace becomes the trace over $\mathcal{H}^{\mathrm{s}}$ of
\[
\begin{aligned}
&\mathfrak{M}_{(\beta^{\mathrm{mag}}_0)}^{*}\,K^{1/2}\exp\left(i\langle\beta^{\mathrm{el}}_{n_1-1},E\rangle\right)K^{1/2}\,\mathfrak{M}_{(\beta^{\mathrm{mag}}_{n_1-1})}\\
&\qquad\mathfrak{M}_{(\beta^{\mathrm{mag}}_{n_1-1})}^{*}\,K^{1/2}\exp\left(i\langle\beta^{\mathrm{el}}_{n_1-2},E\rangle\right)K^{1/2}\,\mathfrak{M}_{(\beta^{\mathrm{mag}}_{n_1-2})}\\
&\qquad\cdots\,\mathfrak{M}_{(\beta^{\mathrm{mag}}_1)}^{*}\,K^{1/2}\exp\left(i\langle\beta^{\mathrm{el}}_0,E\rangle\right)K^{1/2}\,\mathfrak{M}_{(\beta^{\mathrm{mag}}_0)}.
\end{aligned}
\]
Grouped as displayed, these are the factors $\mathfrak{M}_{(\beta^{\mathrm{mag}}_{j+1})}^{*}K^{1/2}\exp\left(i\langle\beta^{\mathrm{el}}_j,E\rangle\right)K^{1/2}\mathfrak{M}_{(\beta^{\mathrm{mag}}_j)}$
for $j=n_1-1,n_1-2,\dots,0$, with $\beta^{\mathrm{mag}}_{n_1}=\beta^{\mathrm{mag}}_0$, in the order of the right side of the proposition.
\end{proof}

\begin{remark}
For the Wilson fermion, whose determinant takes the place of $\det D^{\mathbb{T}}_A$, the identity corresponding to Proposition~\ref{prop:external-field}, with the same constant, is \cite[Proposition~10.5]{NoncompactQED};
there $h$ and $n_0$ stand for $\beta$ and $n_1$, the operators $M_{\beta'}$, $\beta'\in\R^{\Lambda^{(2)}_{\mathrm{s}}}$, for $\mathfrak{M}_{(\beta')}$,
and one Fock space takes the place of the two spaces $\mathcal{F}'$ and $\mathcal{F}$.
\end{remark}

\subsubsection{Other time directions}\label{sec:time-directions}

The reflection rule for the functional in Proposition~\ref{prop:symmetries}\eqref{it:functional-covariance}
extends to the class $\mathcal{I}(\mathbb{T})$ of Definition~\ref{def:torus-integrands}.

\begin{lemma}[Reflected tori]\label{lem:reflection-covariance}
Let $w_0,\dots,w_3>0$, $M\in\R$, $q>0$ and $-1\le\rho\le1$.
Let $\Pi=(\ell,c)$ be a mirror of any of the nine families, let $\mathbb{T}=(\Gamma,\chi)$ be a torus, and let $G\in\mathcal{I}(\mathbb{T})$.
For $(A',z')\in\mathfrak{C}(\theta_\Pi\Gamma)$ put
\[
(\Theta_\Pi G)(A',z'):=\Theta_\Pi\left(G\left(\theta_\Pi^{*}A',\theta_\Pi^{*}z'\right)\right),
\]
where $\Theta_\Pi$ on the right is the antilinear map $\Gr_{\mathbb{T}}\to\Gr_{\theta_\Pi\mathbb{T}}$ of Lemma~\ref{lem:reflected-algebras}.
Then $\Theta_\Pi G\in\mathcal{I}(\theta_\Pi\mathbb{T})$, and
\[
Z_{\theta_\Pi\mathbb{T}}[\Theta_\Pi G]=\overline{Z_{\mathbb{T}}[G]} .
\]
Consequently, if $\Pi'$ and $\Pi''$ are mirrors and $g:=\theta_{\Pi'}\theta_{\Pi''}$,
then $\Theta_{\Pi'}\Theta_{\Pi''}G\in\mathcal{I}(g\mathbb{T})$ and
\[
Z_{g\mathbb{T}}\left[\Theta_{\Pi'}\Theta_{\Pi''}G\right]=Z_{\mathbb{T}}[G] .
\]
\end{lemma}

\begin{proof}
Write $\theta:=\theta_\Pi$ and $n:=\hat n_\ell$,
and let $\vec\theta\in\Ocub$ be the linear part of $\theta$, so that $\theta(x+v)=\theta x+\vec\theta v$ for $x,v\in\R^3$;
then $\theta\mathbb{T}=\left(\vec\theta\,\Gamma,\chi\circ\vec\theta^{-1}\right)$ (Definition~\ref{def:tori}).

\emph{The Grassmann algebras.}
Let $v\in\Gamma$.
Since $\chi$ is real, the reflection \eqref{eq:setting-reflection-fields} carries the relations
$\psi(x+v)=\chi(v)\cdot\psi(x)$ and $\bar\psi(x+v)=\chi(v)\cdot\bar\psi(x)$ of $\Gr_{\mathbb{T}}$ to
\[
\bar\psi(\theta x+\vec\theta v)\,\gamma_n=\chi(v)\cdot\bar\psi(\theta x)\,\gamma_n,
\qquad
\gamma_n\,\psi(\theta x+\vec\theta v)=\chi(v)\cdot\gamma_n\,\psi(\theta x).
\]
As $\gamma_n$ is invertible and $\chi(v)=\left(\chi\circ\vec\theta^{-1}\right)(\vec\theta v)$, these are relations of $\Gr_{\theta\mathbb{T}}$.
Hence $\Theta_\Pi$ induces an antilinear map $\Gr_{\mathbb{T}}\to\Gr_{\theta\mathbb{T}}$ that reverses products;
it is bijective, because $\Theta_\Pi^2=1$ and $\theta(\theta\mathbb{T})=\mathbb{T}$.
Let $Y$ be a set of representatives of $\mathbb{X}_{\mathbb{T}}$.
Then $\theta Y$ is a set of representatives of $\mathbb{X}_{\theta\mathbb{T}}=\mathbb{X}/\vec\theta\,\Gamma$,
the generators at $Y$ and at $\theta Y$ are free generators of $\Gr_{\mathbb{T}}$ and of $\Gr_{\theta\mathbb{T}}$ (Definition~\ref{def:tori}),
and in these generators the induced map is the reflection $\Theta_\Pi\colon\Gr(Y)\to\Gr(\theta Y)$.

\emph{The class.}
By Definition~\ref{def:configurations}, the map
$(A',z')\mapsto\theta^{*}(A',z'):=(\theta^{*}A',\theta^{*}z')$
is a measure preserving bijection $\Conf(\theta\Gamma)\to\Conf(\Gamma)$.
The pull-back $\theta^{*}$ commutes with $d$ and carries integer cochains to integer cochains (Definition~\ref{def:cochains});
so it maps $\mathfrak{C}(\theta\Gamma)$ onto $\mathfrak{C}(\Gamma)$,
and the integer shift by $k'\in C^1(\theta\Gamma;\Z)$ to the integer shift by $\theta^{*}k'\in C^1(\Gamma;\Z)$.
Moreover
\[
F\left(\theta^{*}(A',z')\right)=d\,\theta^{*}A'+\frac{2\pi}{q}\,\theta^{*}z'=\theta^{*}F(A',z'),
\]
and $\left\|\theta^{*}F(A',z')\right\|=\|F(A',z')\|$, the pull-back being orthogonal.
By \eqref{eq:setting-reflection-fields}, $\Theta_\Pi$ maps a monomial $\Psi$ in the generators at $Y$
to a linear combination of at most $4^{\deg\Psi}$ monomials in the generators at $\theta Y$,
whose coefficients are products of entries of the unitary matrix $\gamma_n$ and have modulus at most one.
Therefore every coefficient of $(\Theta_\Pi G)(A',z')$ in the monomial basis at $\theta Y$
is a linear combination, with a bounded number of terms and coefficients of modulus at most one,
of complex conjugates of coefficients of $G\left(\theta^{*}(A',z')\right)$;
with the constants $C$ and $\varepsilon_G$ of $G$ it is bounded by a constant multiple of
$\exp\left(\left(\tfrac12-\varepsilon_G\right)\|F(A',z')\|^2\right)$.
The map $\Theta_\Pi G$ is Borel, being composed of Borel maps, and it is invariant under the integer shifts;
thus $\Theta_\Pi G\in\mathcal{I}(\theta\mathbb{T})$.

\emph{The functional.}
Proposition~\ref{prop:symmetries}\eqref{it:action-invariance}, read on the torus $\mathbb{T}$, gives
$\Theta_\Pi\left(S_{\mathbb{T}}(\theta^{*}A')\right)=S_{\theta\mathbb{T}}(A')$ for every $A'\in C^1(\theta\Gamma)$.
The element $S_{\mathbb{T}}(\theta^{*}A')$ is a sum of terms of degree two;
it is nilpotent, and its exponential is a polynomial in it with real coefficients.
As $\Theta_\Pi$ is antilinear and maps the powers of an element to the powers of its image,
\[
\Theta_\Pi\left(\exp\left(-S_{\mathbb{T}}(\theta^{*}A')\right)\right)=\exp\left(-S_{\theta\mathbb{T}}(A')\right).
\]
Let $(A,z):=\theta^{*}(A',z')$.
Since $\Theta_\Pi$ reverses products, and since the element $\exp\left(-S_{\mathbb{T}}(A)\right)$ is even and therefore commutes with $G(A,z)$,
\[
\exp\left(-S_{\theta\mathbb{T}}(A')\right)\cdot(\Theta_\Pi G)(A',z')
=\Theta_\Pi\left(G(A,z)\cdot\exp\left(-S_{\mathbb{T}}(A)\right)\right)
=\Theta_\Pi\left(\exp\left(-S_{\mathbb{T}}(A)\right)\cdot G(A,z)\right).
\]
Proposition~\ref{prop:berezin}\eqref{it:reflected-integral}, applied in the free generators at $Y$ and at $\theta Y$,
shows that the Berezin integral over $\theta\mathbb{T}$ of this element is the complex conjugate of
$\int_{\mathbb{T}}\exp\left(-S_{\mathbb{T}}(A)\right)\cdot G(A,z)$.
Put
\[
\Phi(A,z):=\exp\left(-\tfrac12\|F(A,z)\|^2\right)\cdot\int_{\mathbb{T}}\exp\left(-S_{\mathbb{T}}(A)\right)\cdot G(A,z)
\qquad\left((A,z)\in\mathfrak{C}(\Gamma)\right),
\]
and let $\Phi'$ be the function on $\mathfrak{C}(\theta\Gamma)$ formed in the same way with $\theta\mathbb{T}$ and $\Theta_\Pi G$.
Since $\|F(A',z')\|=\|F(\theta^{*}(A',z'))\|$, what has been shown reads
\[
\Phi'(A',z')=\overline{\Phi\left(\theta^{*}(A',z')\right)}\qquad\left((A',z')\in\mathfrak{C}(\theta\Gamma)\right).
\]
Both functions are invariant under the integer shifts,
the action depending on $A$ only through the phases (Proposition~\ref{prop:symmetries}\eqref{it:gauge-covariance}),
and they are integrable (Proposition~\ref{prop:berezin}\eqref{it:torus-integrability});
by Definition~\ref{def:torus-integrands}, their integrals over $\Conf(\Gamma)$ and over $\Conf(\theta\Gamma)$
are $Z_{\mathbb{T}}[G]$ and $Z_{\theta\mathbb{T}}[\Theta_\Pi G]$.
The measure preserving bijection $\theta^{*}$ now gives $Z_{\theta\mathbb{T}}[\Theta_\Pi G]=\overline{Z_{\mathbb{T}}[G]}$.

\emph{Two mirrors.}
Apply what has been proved to $\Pi''$ and $G$, and then to $\Pi'$ and to $\Theta_{\Pi''}G\in\mathcal{I}(\theta_{\Pi''}\mathbb{T})$.
Since the linear parts compose, $\theta_{\Pi'}(\theta_{\Pi''}\mathbb{T})=g\mathbb{T}$;
so $\Theta_{\Pi'}\Theta_{\Pi''}G\in\mathcal{I}(g\mathbb{T})$, and
\[
Z_{g\mathbb{T}}\left[\Theta_{\Pi'}\Theta_{\Pi''}G\right]
=\overline{Z_{\theta_{\Pi''}\mathbb{T}}\left[\Theta_{\Pi''}G\right]}
=\overline{\overline{Z_{\mathbb{T}}[G]}}
=Z_{\mathbb{T}}[G].
\]
Explicitly, $(\Theta_{\Pi'}\Theta_{\Pi''}G)(A',z')=\Theta_{\Pi'}\Theta_{\Pi''}\left(G(g^{*}A',g^{*}z')\right)$ for $(A',z')\in\mathfrak{C}(g\Gamma)$,
because $\theta_{\Pi''}^{*}\theta_{\Pi'}^{*}=(\theta_{\Pi'}\theta_{\Pi''})^{*}=g^{*}$ (Definition~\ref{def:cochains}).
\end{proof}

\emph{The cyclic permutation of the axes.}
For $\mu\in\{2,3\}$ let $g_\mu\in\Ocub$ be the cyclic permutation matrix with $g_\mu e_1=e_\mu$ (Definition~\ref{def:other-directions}).
It is the product of the reflections in two diagonal mirrors through the origin, and in particular it is even.
Namely, let
\begin{equation}\label{eq:slab-cyclic-mirrors}
(\Pi',\Pi''):=
\begin{cases}
\left((x_1-x_2,0),\,(x_2-x_3,0)\right) & \text{if }\mu=2,\\
\left((x_2-x_3,0),\,(x_1-x_2,0)\right) & \text{if }\mu=3.
\end{cases}
\end{equation}
By Definition~\ref{def:mirrors}, the reflection $\theta_{x_1-x_2,0}$ exchanges the coordinates $x_1$ and $x_2$,
and $\theta_{x_2-x_3,0}$ exchanges $x_2$ and $x_3$.
Hence $\theta_{\Pi'}\theta_{\Pi''}$ maps $(x_1,x_2,x_3)$ to $(x_3,x_1,x_2)$ if $\mu=2$ and to $(x_2,x_3,x_1)$ if $\mu=3$;
it is a cyclic permutation of the coordinates that maps $e_1$ to $e_\mu$, and so
\[
\theta_{\Pi'}\theta_{\Pi''}=g_\mu .
\]
The upward unit normals of $\Pi'$ and $\Pi''$ (Definition~\ref{def:mirrors}) are
$\hat n'=\tfrac{1}{\sqrt2}(e_1-e_2)$ and $\hat n''=\tfrac{1}{\sqrt2}(e_2-e_3)$ if $\mu=2$,
and $\hat n'=\tfrac{1}{\sqrt2}(e_2-e_3)$ and $\hat n''=\tfrac{1}{\sqrt2}(e_1-e_2)$ if $\mu=3$.
Put
\[
U_{g_\mu}:=\gamma_{\hat n''}\gamma_{\hat n'} .
\]
For every unit vector $n$ the matrix $\gamma_n$ is Hermitian with $\gamma_n^2=1$ (Definition~\ref{def:spin-matrices});
hence $U_{g_\mu}$ is unitary, and $U_{g_\mu}^{-1}=U_{g_\mu}^{*}=\gamma_{\hat n'}\gamma_{\hat n''}$.
The identity $\gamma_n\gamma_v\gamma_n=-\gamma_{v-2(v\cdot n)n}$ of Definition~\ref{def:spin-matrices},
applied with $n=\hat n''$ and then with $n=\hat n'$, gives
\begin{equation}\label{eq:slab-spin-rotation}
U_{g_\mu}^{-1}\,\gamma_v\,U_{g_\mu}
=\gamma_{\hat n'}\gamma_{\hat n''}\,\gamma_v\,\gamma_{\hat n''}\gamma_{\hat n'}
=\gamma_{g_\mu v},
\qquad v\in\R^3 ;
\end{equation}
indeed, $v\mapsto v-2(v\cdot n)n$ is the reflection in the plane through the origin orthogonal to $n$,
and for $n=\hat n''$ and $n=\hat n'$ these reflections are $\theta_{\Pi''}$ and $\theta_{\Pi'}$, whose product $\theta_{\Pi'}\theta_{\Pi''}$ is $g_\mu$.
If $\mathbb{T}=C(n_1,n_2,n_3)$ is a coordinate torus, then the torus $g_\mu\mathbb{T}$ of Definition~\ref{def:tori}
is the coordinate torus $C(n_3,n_1,n_2)$ if $\mu=2$ and $C(n_2,n_3,n_1)$ if $\mu=3$;
in both cases its side in the direction $e_\mu$ is $n_1$.

\begin{proposition}[Cyclic permutation of the axes]\label{prop:cyclic-axes}
Let $w_0,\dots,w_3>0$, $M\in\R$, $q>0$ and $-1\le\rho\le1$.
Let $\mu\in\{2,3\}$, let $\Pi'$ and $\Pi''$ be the mirrors \eqref{eq:slab-cyclic-mirrors},
so that $\theta_{\Pi'}\theta_{\Pi''}=g_\mu$ is the cyclic permutation matrix with $g_\mu e_1=e_\mu$,
and let $U_{g_\mu}=\gamma_{\hat n''}\gamma_{\hat n'}$ be the unitary matrix formed with their upward unit normals $\hat n'$, $\hat n''$.
Let $\mathbb{T}=(\Gamma,\chi)$ be a torus.
Then $\Theta_{\Pi'}\Theta_{\Pi''}$ is the linear automorphism of $\Gr(\mathbb{X})$ with
\[
\Theta_{\Pi'}\Theta_{\Pi''}\,\psi(x)=U_{g_\mu}\,\psi(g_\mu x),
\qquad
\Theta_{\Pi'}\Theta_{\Pi''}\,\bar\psi(x)=\bar\psi(g_\mu x)\,U_{g_\mu}^{-1}
\qquad(x\in\mathbb{X}),
\]
and it induces an isomorphism of algebras from $\Gr_{\mathbb{T}}$ onto $\Gr_{g_\mu\mathbb{T}}$.
If $\Xi=(\Gamma_\parallel,\chi_\parallel)$ is a cross-section for the family $x_1$ (Definition~\ref{def:slabs}),
then $g_\mu\Xi:=\left(g_\mu\Gamma_\parallel,\chi_\parallel\circ g_\mu^{-1}\right)$ is a cross-section for the family $x_\mu$,
and for every open interval $\mathbb{I}\subset\R$ the automorphism $\Theta_{\Pi'}\Theta_{\Pi''}$ induces an isomorphism of algebras
from $\Gr_\Xi(\mathbb{I})$, formed with the family $x_1$, onto $\Gr_{g_\mu\Xi}(\mathbb{I})$, formed with the family $x_\mu$ (Definition~\ref{def:slab-grassmann}).
Moreover, the following hold.
\begin{enumerate}
\item\label{it:transport-integral}
For every $G\in\Gr_{\mathbb{T}}$,
\[
\int_{g_\mu\mathbb{T}}\Theta_{\Pi'}\Theta_{\Pi''}(G)=\int_{\mathbb{T}}G .
\]
\item\label{it:transport-action}
Let $\hat A$ be a real one-cochain on $\Z^3$, and let $S_F(\hat A)$ be the fermion action of Definition~\ref{def:chamber-operator}, read as the formal sum of its terms.
For every oriented bond $(x,x')$ the automorphism $\Theta_{\Pi'}\Theta_{\Pi''}$ maps the term
$-\bar\psi(x)\cdot w(x,x')\cdot U_{\hat A}(x,x')\cdot\psi(x')$ of $(x,x')$ in $S_F(\hat A)$
to the term of the oriented bond $(g_\mu x,g_\mu x')$ in $S_F\left((g_\mu^{-1})^{*}\hat A\right)$,
and for every $x\in\mathbb{X}$ it maps the diagonal term $M\bar\psi(x)\psi(x)$ to the diagonal term $M\bar\psi(g_\mu x)\psi(g_\mu x)$.
Consequently, for every $A\in C^1(\Gamma)$,
\[
\Theta_{\Pi'}\Theta_{\Pi''}\left(S_{\mathbb{T}}(A)\right)=S_{g_\mu\mathbb{T}}\left((g_\mu^{-1})^{*}A\right).
\]
In particular, for $c\in\tfrac12\Z$ let $S^{\times}_{x_1,c}(A)$ be the sum of the terms of $S_{\mathbb{T}}(A)$ that belong to bonds crossing the plane $\{x_1=c\}$,
and let $S^{\times}_{x_\mu,c}$ be formed in the same way with $S_{g_\mu\mathbb{T}}$ and the plane $\{x_\mu=c\}$;
then
\[
\Theta_{\Pi'}\Theta_{\Pi''}\left(S^{\times}_{x_1,c}(A)\right)=S^{\times}_{x_\mu,c}\left((g_\mu^{-1})^{*}A\right).
\]
\item\label{it:transport-reflection}
For every $c\in\tfrac12\Z$,
\[
\Theta_{\Pi'}\Theta_{\Pi''}\,\Theta_{x_1,c}\,\Theta_{\Pi''}\Theta_{\Pi'}=\Theta_{x_\mu,c}
\]
on the even and on the odd elements of $\Gr(\mathbb{X})$;
here $\Theta_{\Pi''}\Theta_{\Pi'}$ is the inverse of $\Theta_{\Pi'}\Theta_{\Pi''}$,
and $\Theta_{x_\mu,c}=\Theta_{(x_\mu,c)}$ is formed with the upward unit normal $e_\mu$ of the family $x_\mu$.
\end{enumerate}
\end{proposition}

\begin{proof}
Write $g:=g_\mu$ and $U:=U_{g_\mu}$.
By \eqref{eq:setting-two-reflections}, with $\Pi'$ and $\Pi''$ in the places of the first and of the second mirror there,
$\Theta_{\Pi'}\Theta_{\Pi''}$ is a linear automorphism of $\Gr(\mathbb{X})$ and
\[
\Theta_{\Pi'}\Theta_{\Pi''}\,\psi(x)=\gamma_{\hat n''}\gamma_{\hat n'}\,\psi(\theta_{\Pi'}\theta_{\Pi''}x)=U\psi(gx),
\qquad
\Theta_{\Pi'}\Theta_{\Pi''}\,\bar\psi(x)=\bar\psi(gx)\,\gamma_{\hat n'}\gamma_{\hat n''}=\bar\psi(gx)\,U^{-1}.
\]
Since $\Theta_\Pi^2=1$ for every mirror $\Pi$ (Definition~\ref{def:grassmann-reflections}), the inverse of this automorphism is $\Theta_{\Pi''}\Theta_{\Pi'}$;
by \eqref{eq:setting-two-reflections}, now with $\Pi''$ in the first place, it maps $\psi(y)$ to $U^{-1}\psi(g^{-1}y)$ and $\bar\psi(y)$ to $\bar\psi(g^{-1}y)\,U$.
By Lemma~\ref{lem:reflected-algebras}, $\Theta_{\Pi''}$ induces an antilinear bijection $\Gr_{\mathbb{T}}\to\Gr_{\theta_{\Pi''}\mathbb{T}}$ that reverses products,
and $\Theta_{\Pi'}$ induces one from $\Gr_{\theta_{\Pi''}\mathbb{T}}$ onto $\Gr_{\theta_{\Pi'}\theta_{\Pi''}\mathbb{T}}=\Gr_{g\mathbb{T}}$;
their composition is an isomorphism of algebras $\Gr_{\mathbb{T}}\to\Gr_{g\mathbb{T}}$,
and it maps the class in $\Gr_{\mathbb{T}}$ of an element $G$ of $\Gr(\mathbb{X})$ to the class of $\Theta_{\Pi'}\Theta_{\Pi''}G$ in $\Gr_{g\mathbb{T}}$.

For $y\in\R^3$ one has $(gy)_\mu=gy\cdot e_\mu=y\cdot g^{-1}e_\mu=y_1$, as $g$ is orthogonal and $ge_1=e_\mu$;
hence $g$ maps every plane $\{x_1=t\}$, $t\in\R$, onto the plane $\{x_\mu=t\}$, and the two sides of the first onto the corresponding sides of the second.
Let $\Xi=(\Gamma_\parallel,\chi_\parallel)$ be a cross-section for $x_1$.
The permutation matrix $g$ maps $\Gamma_\parallel$ onto a subgroup $g\Gamma_\parallel$ of rank two of $\Z^3\cap\{x_\mu=0\}$ with $\sys(g\Gamma_\parallel)=\sys(\Gamma_\parallel)\ge4$,
and $\chi_\parallel\circ g^{-1}$ is a character of $g\Gamma_\parallel$ with values $\pm1$;
so $g\Xi$ is a cross-section for $x_\mu$ (Definition~\ref{def:slabs}).
Let $\mathbb{I}\subset\R$ be an open interval.
The map $g$ carries the fermion sites $x$ with $x_1\in\mathbb{I}$ onto those with $x_\mu\in\mathbb{I}$,
and $\Theta_{\Pi'}\Theta_{\Pi''}$ carries the defining relations $\psi(x+v)=\chi_\parallel(v)\cdot\psi(x)$ and $\bar\psi(x+v)=\chi_\parallel(v)\cdot\bar\psi(x)$, $v\in\Gamma_\parallel$,
of $\Gr_\Xi(\mathbb{I})$ (Definition~\ref{def:slab-grassmann}) to
\[
U\psi(gx+gv)=\chi_\parallel(v)\cdot U\psi(gx),
\qquad
\bar\psi(gx+gv)\,U^{-1}=\chi_\parallel(v)\cdot\bar\psi(gx)\,U^{-1}.
\]
As $U$ is invertible and $(\chi_\parallel\circ g^{-1})(gv)=\chi_\parallel(v)$, these follow from the defining relations of $\Gr_{g\Xi}(\mathbb{I})$;
so $\Theta_{\Pi'}\Theta_{\Pi''}$ induces a homomorphism of algebras $\Gr_\Xi(\mathbb{I})\to\Gr_{g\Xi}(\mathbb{I})$.
The same computation with $\Theta_{\Pi''}\Theta_{\Pi'}$ and $g^{-1}$, which maps $g\Xi$ back to $\Xi$, gives a homomorphism $\Gr_{g\Xi}(\mathbb{I})\to\Gr_\Xi(\mathbb{I})$;
the two homomorphisms are inverse to each other, being induced by mutually inverse automorphisms of $\Gr(\mathbb{X})$.

Consider~\eqref{it:transport-integral}.
Let $Y$ be a set of representatives of $\mathbb{X}_{\mathbb{T}}$;
then $\theta_{\Pi''}Y$ and $gY$ are sets of representatives of $\mathbb{X}_{\theta_{\Pi''}\mathbb{T}}$ and of $\mathbb{X}_{g\mathbb{T}}$.
Proposition~\ref{prop:berezin}\eqref{it:reflected-integral}, applied in the free generators at $Y$ with the mirror $\Pi''$,
and then in the free generators at $\theta_{\Pi''}Y$ with the mirror $\Pi'$, gives
\[
\int_{g\mathbb{T}}\Theta_{\Pi'}\Theta_{\Pi''}(G)
=\overline{\int_{\theta_{\Pi''}\mathbb{T}}\Theta_{\Pi''}(G)}
=\overline{\overline{\int_{\mathbb{T}}G}}
=\int_{\mathbb{T}}G .
\]

Turning to~\eqref{it:transport-action}, let $\Pi$ be a mirror, $\theta:=\theta_\Pi$, and let $n$ be its upward unit normal.
By Proposition~\ref{prop:symmetries}\eqref{it:action-invariance}, for every real one-cochain $\hat A'$ on $\Z^3$ the map $\Theta_\Pi$ carries the term of an oriented bond $(x,x')$ in $S_F(\theta^{*}\hat A')$
to the term of the oriented bond $(\theta x',\theta x)$ in $S_F(\hat A')$.
Take $\hat A':=\theta^{*}\hat A$; as $\theta^{*}\theta^{*}=(\theta\theta)^{*}$ is the identity (Definition~\ref{def:cochains}),
$\Theta_\Pi$ carries the term of $(x,x')$ in $S_F(\hat A)$ to the term of $(\theta x',\theta x)$ in $S_F(\theta^{*}\hat A)$.
The diagonal term at $x$ is mapped to the diagonal term at $\theta x$ by \eqref{eq:setting-reflection-bilinear},
since $M$ is real and $\gamma_n^2=1$.
Applied to $\Pi''$ and then to $\Pi'$, this shows that $\Theta_{\Pi'}\Theta_{\Pi''}$ maps the term of $(x,x')$ in $S_F(\hat A)$
first to the term of $(\theta_{\Pi''}x',\theta_{\Pi''}x)$ in $S_F(\theta_{\Pi''}^{*}\hat A)$
and then to the term of $(\theta_{\Pi'}\theta_{\Pi''}x,\theta_{\Pi'}\theta_{\Pi''}x')=(gx,gx')$ in $S_F(\theta_{\Pi'}^{*}\theta_{\Pi''}^{*}\hat A)$,
and the diagonal term at $x$ to the diagonal term at $gx$, where
\[
\theta_{\Pi'}^{*}\theta_{\Pi''}^{*}=(\theta_{\Pi''}\theta_{\Pi'})^{*}=(g^{-1})^{*},
\]
because the pull-back by a product of two isometries is the product of the pull-backs in the reverse order (Definition~\ref{def:cochains})
and $\theta_{\Pi''}\theta_{\Pi'}=(\theta_{\Pi'}\theta_{\Pi''})^{-1}$.
This is the assertion on the single terms.
Now let $A\in C^1(\Gamma)$, read as a $\Gamma$-periodic cochain on $\Z^3$, so that $(g^{-1})^{*}A$ is $g\Gamma$-periodic,
and let $Y$ be a set of representatives of $\mathbb{X}_{\mathbb{T}}$.
By Definitions~\ref{def:chamber-operator} and~\ref{def:tori}, $S_{\mathbb{T}}(A)$ is the class in $\Gr_{\mathbb{T}}$ of the sum of the diagonal terms at the sites $x\in Y$
and of the terms of the oriented bonds $(x,x^{(i)})$, $x\in Y$, $i\in\{0,1,2,3\}$, in $S_F(A)$.
The set $gY$ is a set of representatives of $\mathbb{X}_{g\mathbb{T}}$, since $g(x+v)=gx+gv$ for $v\in\Gamma$,
and $g$ maps the neighbour $x^{(i)}$ of $x$ to the neighbour $(gx)^{(i)}$ of $gx$ (Definition~\ref{def:fermion-sites}).
The assertion on the single terms and the behaviour of the induced isomorphism on classes therefore give
$\Theta_{\Pi'}\Theta_{\Pi''}\left(S_{\mathbb{T}}(A)\right)=S_{g\mathbb{T}}\left((g^{-1})^{*}A\right)$.
A bond $\{x,x'\}$ crosses the plane $\{x_1=c\}$ if and only if $\{gx,gx'\}$ crosses $\{x_\mu=c\}$,
and the identity for the crossing parts follows term by term.

For~\eqref{it:transport-reflection}, let $\theta:=\theta_{x_1,c}$.
The isometry $g$ maps the plane $\{x_1=c\}$ onto $\{x_\mu=c\}$, so that $g\theta g^{-1}=\theta_{x_\mu,c}$.
The map $\Theta_{x_1,c}$ is antilinear and is given by \eqref{eq:setting-reflection-fields} with the normal $e_1$;
for a constant $4\times4$ matrix $D$ it therefore maps $D\psi(y')$ to $-\bar\psi(\theta y')\,\gamma_{e_1}D^{*}$
and $\bar\psi(y')D$ to $-D^{*}\gamma_{e_1}\psi(\theta y')$.
With $(U^{-1})^{*}=U$, $U^{*}=U^{-1}$ and \eqref{eq:slab-spin-rotation} for $v=e_1$, one obtains for $y\in\mathbb{X}$
\[
\begin{aligned}
\Theta_{\Pi'}\Theta_{\Pi''}\,\Theta_{x_1,c}\left(U^{-1}\psi(g^{-1}y)\right)
&=\Theta_{\Pi'}\Theta_{\Pi''}\left(-\bar\psi(\theta g^{-1}y)\,\gamma_{e_1}U\right)\\
&=-\bar\psi(g\theta g^{-1}y)\,U^{-1}\gamma_{e_1}U
=-\bar\psi(\theta_{x_\mu,c}\,y)\,\gamma_{e_\mu},\\
\Theta_{\Pi'}\Theta_{\Pi''}\,\Theta_{x_1,c}\left(\bar\psi(g^{-1}y)\,U\right)
&=\Theta_{\Pi'}\Theta_{\Pi''}\left(-U^{-1}\gamma_{e_1}\psi(\theta g^{-1}y)\right)\\
&=-U^{-1}\gamma_{e_1}U\,\psi(g\theta g^{-1}y)
=-\gamma_{e_\mu}\psi(\theta_{x_\mu,c}\,y).
\end{aligned}
\]
The left sides are the images of $\psi(y)$ and of $\bar\psi(y)$ under the left side of~\eqref{it:transport-reflection},
and by \eqref{eq:setting-reflection-fields} with the normal $e_\mu$ the right sides are $\Theta_{x_\mu,c}\psi(y)$ and $\Theta_{x_\mu,c}\bar\psi(y)$.
Both sides of~\eqref{it:transport-reflection} are antilinear maps of $\Gr(\mathbb{X})$ that reverse products,
the left side being composed of two linear automorphisms and of one antilinear map that reverses products.
They agree on the generators, and $\Gr(\mathbb{X})$ is generated by them as an algebra;
hence they agree on all elements, even and odd.
\end{proof}

\emph{The objects of the time direction $e_\mu$.}
Let $\mu\in\{2,3\}$, and write $g:=g_\mu$ and $U:=U_{g_\mu}$.
The objects of Definitions~\ref{def:slices}, \ref{def:modes}, \ref{def:slice-operators}, \ref{def:mid-plane} and~\ref{def:insertions}
attached to the time direction $e_\mu$ are defined as the images under $g$ of the corresponding objects attached to $e_1$;
at the coordinate link planes this agrees with Definition~\ref{def:other-directions}.
For an even element $G$ the image is $G^{g}=\Theta_{\Pi'}\Theta_{\Pi''}G$ \eqref{eq:setting-group-action}, which does not depend on the word that represents $g$.
For an odd element, a mode for instance, the image under $g$ is defined as its image under the automorphism $\Theta_{\Pi'}\Theta_{\Pi''}$ of Proposition~\ref{prop:cyclic-axes},
that is, through the word $g=\theta_{\Pi'}\theta_{\Pi''}$ formed with the mirrors \eqref{eq:slab-cyclic-mirrors};
in the algebras $\Gr_{\mathbb{T}}$ and $\Gr_\Xi(\mathbb{I})$ the images are taken with the isomorphisms that this automorphism induces.
At the planes of the family $x_1$ the objects written $\Theta_c$, $\Ins_c$ and $S^{\times}_c$ in Subsection~\ref{sec:link-planes}
are written $\Theta_{x_1,c}$, $\Ins_{x_1,c}$ and $S^{\times}_{x_1,c}$ below.
Explicitly, the objects of the direction $e_\mu$ are the following.

The slice variables of a real one-cochain $A'$ in the direction $e_\mu$ are those of $g^{*}A'$ in the direction $e_1$ (Definition~\ref{def:slices}),
and the electric and magnetic parts of a real two-cochain $\beta'$ in the direction $e_\mu$ are those of $g^{*}\beta'$ in the direction $e_1$.
If $A'=(g^{-1})^{*}A$, then $g^{*}A'=A$;
so the slices of $A'$ in the direction $e_\mu$ are the slices $A_j$, $a_j$ of $A$ in the direction $e_1$.

In Definition~\ref{def:spin-matrices} the orthonormal bases at the unit vector $e_\mu$ are arbitrary.
They are normalised by
\begin{equation}\label{eq:slab-transported-bases}
e^{\pm}_{e_\mu,k}:=U^{-1}e^{\pm}_{e_1,k}\qquad(k=1,2),
\end{equation}
and then $e^{\pm}_{-e_\mu,k}=e^{\mp}_{e_\mu,k}=U^{-1}e^{\pm}_{-e_1,k}$, by the rule $e^{\pm}_{-n,k}=e^{\mp}_{n,k}$ of that definition.
The normalisation is admissible:
by \eqref{eq:slab-spin-rotation} with $v=e_1$, one has $\gamma_{e_\mu}U^{-1}=U^{-1}\gamma_{e_1}$,
so that the unitary matrix $U^{-1}$ maps $\operatorname{ran}P^{\pm}_{e_1}$ onto $\operatorname{ran}P^{\pm}_{e_\mu}$
and carries orthonormal bases of the first space to orthonormal bases of the second.

Let $c\in\tfrac12\Z$.
An index $i\in\idx(x_1,c)$ consists of a bond $\{x,\theta_{x_1,c}x\}$ that crosses $(x_1,c)$, with the lower end $x$,
of one of the four groups and of $k\in\{1,2\}$ (Definition~\ref{def:modes}).
Since $g$ is orthogonal and $ge_1=e_\mu$, one has $(gy)_\mu=gy\cdot e_\mu=y\cdot g^{-1}e_\mu=y_1$ for $y\in\R^3$;
so $g$ maps the plane $\{x_1=c\}$ onto $\{x_\mu=c\}$, one has $g\,\theta_{x_1,c}=\theta_{x_\mu,c}\,g$,
and $g\in W$ maps the bonds that cross $(x_1,c)$ bijectively onto the bonds that cross $(x_\mu,c)$, lower ends onto lower ends.
Let $gi\in\idx(x_\mu,c)$ consist of the bond $\{gx,\theta_{x_\mu,c}\,gx\}$, with the lower end $gx$, of the same group and of the same $k$.
The map $i\mapsto gi$ is a bijection of $\idx(x_1,c)$ onto $\idx(x_\mu,c)$, and $g(i+v)=gi+gv$ for $v\in\Z^3$;
for $I\subset\idx(x_1,c)$ let $gI:=\{gi:\ i\in I\}$.
Definition~\ref{def:modes} subjects the total orders of the sets $\idx(x_\mu,c)$, $c\in\tfrac12\Z$, only to the condition
that $i\mapsto i+e_\mu$ be an order preserving bijection of $\idx(x_\mu,c)$ onto $\idx(x_\mu,c+1)$.
These orders are normalised to be the images of the orders of the sets $\idx(x_1,c)$:
$gi<gi'$ in $\idx(x_\mu,c)$ if and only if $i<i'$ in $\idx(x_1,c)$.
The normalised orders satisfy the condition,
because $g(i+e_1)=gi+e_\mu$ and the orders of the sets $\idx(x_1,c)$ satisfy the corresponding condition with $e_1$.
Let further $\Xi=(\Gamma_\parallel,\chi_\parallel)$ be a cross-section for $x_1$, and let $g\Xi$ be the cross-section for $x_\mu$ of Proposition~\ref{prop:cyclic-axes}.
As $g(i+v)=gi+gv$, the map $i\mapsto gi$ carries the classes of the modes at $(x_1,c)$ modulo $\Gamma_\parallel$
onto the classes of the modes at $(x_\mu,c)$ modulo $g\Gamma_\parallel$.
The sets of representatives of the latter classes (Definitions~\ref{def:modes} and~\ref{def:mode-translations}) are normalised in the same way:
at $(x_\mu,c)$ the set of representatives for $g\Xi$ is $g\mathcal{J}(c)$, with the image of the order,
where $\mathcal{J}(c)\subset\idx(x_1,c)$ is the set of representatives at $(x_1,c)$ for $\Xi$.
Since $\mathcal{J}(c+1)=\mathcal{J}(c)+e_1$, one has $g\mathcal{J}(c+1)=g\mathcal{J}(c)+e_\mu$, as these definitions require.

With these bases and orders the modes of Definition~\ref{def:modes} at $(x_\mu,c)$ are the images of the modes at $(x_1,c)$.
Indeed, let $i\in\idx(x_1,c)$ have the lower end $x$, and let $\theta:=\theta_{x_1,c}$.
By Proposition~\ref{prop:cyclic-axes} and \eqref{eq:slab-transported-bases},
for the groups $1$ and $2$ (upper sign for the group $1$)
\[
\Theta_{\Pi'}\Theta_{\Pi''}\left(\left(e^{\mp}_{e_1,k}\right)^{*}\psi(\theta x)\right)
=\left(e^{\mp}_{e_1,k}\right)^{*}U\,\psi(g\theta x)
=\left(e^{\mp}_{e_\mu,k}\right)^{*}\psi(\theta_{x_\mu,c}\,gx),
\]
because $e^{*}U=(U^{*}e)^{*}=(U^{-1}e)^{*}$ for $e\in\C^4$,
and for the groups $3$ and $4$ (upper sign for the group $3$)
\[
\Theta_{\Pi'}\Theta_{\Pi''}\left(\bar\psi(\theta x)\,e^{\pm}_{e_1,k}\right)
=\bar\psi(g\theta x)\,U^{-1}e^{\pm}_{e_1,k}
=\bar\psi(\theta_{x_\mu,c}\,gx)\,e^{\pm}_{e_\mu,k}.
\]
Thus $\Theta_{\Pi'}\Theta_{\Pi''}(\zeta_i)$ is the mode $\zeta_{gi}$ of Definition~\ref{def:modes} at the mirror $(x_\mu,c)$.
The type of a bond is invariant under $W$,
the neighbour of $g'y$ across the wall of a given type being the image under $g'$ of the neighbour of $y$ across the wall of the same type, for $g'\in W$ and $y\in\mathbb{X}$ (Definition~\ref{def:fermion-sites});
so $\zeta_{gi}$ has the coefficient $\omega_i$.
Its phase for a real one-cochain $A'$ is $U_{A'}(gx,\theta_{x_\mu,c}\,gx)=U_{g^{*}A'}(x,\theta x)$ for the groups $1$ and $2$,
and the complex conjugate of this number for the groups $3$ and $4$,
since $\pi(gy)=g\,\pi(y)$ for $y\in\mathbb{X}$ (Proposition~\ref{prop:geometry}\eqref{it:geometry-projection}) and $(g^{*}A')(y,y')=A'(gy,gy')$ (Definition~\ref{def:cochains});
that is, it is the phase $u_i(g^{*}A')$ of $\zeta_i$ for the cochain $g^{*}A'$.
So the modes at $(x_\mu,c)$, with their order, their groups, their coefficients and their phases for $A'$,
are the elements $\Theta_{\Pi'}\Theta_{\Pi''}(\zeta_i)$, with the order, the groups and the coefficients of the modes $\zeta_i$ at $(x_1,c)$
and with their phases for $g^{*}A'$.

Accordingly, the Fock spaces $\mathcal{F}$ and $\mathcal{F}'$ with their bases $e_I$, the diagonal operators $\omega$, $\omega'$ and $u(a)$,
the localisation $\iota$ and the operators of Definition~\ref{def:slice-operators} are the same in the time direction $e_\mu$ as in the time direction $e_1$,
the sets of representatives $\mathcal{J}(c)$ and $g\mathcal{J}(c)$ corresponding to each other, with their orders, under $i\mapsto gi$.
So is the half-slab matrix $\widehat{\mathfrak{N}}$ as a function of the slice variable.
Let $\Xi=(\Gamma_\parallel,\chi_\parallel)$ be a rectangular cross-section and $j\in\Z$.
In the direction $e_\mu$ the half-slab matrix is formed in the algebra $\Gr_{g\Xi}((j,j+1/2))$ of the half slab $\{j<x_\mu<j+1/2\}$,
with the modes at the planes $\{x_\mu=j\}$ and $\{x_\mu=j+1/2\}$, with $\Theta_{x_\mu,j+1/2}$,
and with the sum of the terms of the action whose sites lie in this half slab.
The isomorphism of $\Gr_\Xi((j,j+1/2))$ onto $\Gr_{g\Xi}((j,j+1/2))$ induced by $\Theta_{\Pi'}\Theta_{\Pi''}$ (Proposition~\ref{prop:cyclic-axes})
maps the corresponding objects of the direction $e_1$ to these:
the modes by what has just been shown;
the elements $\Theta_{x_1,j+1/2}(\zeta_I)$ to the elements $\Theta_{x_\mu,j+1/2}(\zeta_{gI})$ by Proposition~\ref{prop:cyclic-axes}\eqref{it:transport-reflection};
and, by Proposition~\ref{prop:cyclic-axes}\eqref{it:transport-action} applied to the single terms, $g$ carrying the first half slab onto the second,
the terms at a real one-cochain $\hat A$ on $\Z^3$ whose sites lie in $\{j<x_1<j+1/2\}$ to the terms at $(g^{-1})^{*}\hat A$ whose sites lie in $\{j<x_\mu<j+1/2\}$,
where the slice of $(g^{-1})^{*}\hat A$ in the plane $\{x_\mu=j\}$ is the slice of $\hat A$ in the plane $\{x_1=j\}$.
The Berezin integral over the half slab is not changed:
if $Y$ is a finite set of representatives of the sites of $\{j<x_1<j+1/2\}$ modulo $\Gamma_\parallel$,
then $gY$ is a set of representatives of the sites of $\{j<x_\mu<j+1/2\}$ modulo $g\Gamma_\parallel$,
and Proposition~\ref{prop:berezin}\eqref{it:reflected-integral}, applied with the mirror $\Pi''$ to $Y$ and with the mirror $\Pi'$ to $\theta_{\Pi''}Y$,
gives $\int_{gY}\Theta_{\Pi'}\Theta_{\Pi''}(G)=\int_YG$ for $G\in\Gr(Y)$.

Let $\Upsilon$ be a matrix indexed by the finite subsets of a finite part of $\idx(x_1,c)$, read at $(x_\mu,c)$ through the bijection $i\mapsto gi$.
The automorphism $\Theta_{\Pi'}\Theta_{\Pi''}$ is multiplicative and maps $\zeta_I$ to $\zeta_{gI}$, the order being preserved;
by Proposition~\ref{prop:cyclic-axes}\eqref{it:transport-reflection},
$\Theta_{\Pi'}\Theta_{\Pi''}\left(\Theta_{x_1,c}(\zeta_J)\right)=\Theta_{x_\mu,c}\left(\Theta_{\Pi'}\Theta_{\Pi''}(\zeta_J)\right)$.
Hence
\begin{equation}\label{eq:slab-insertion-transport}
\Theta_{\Pi'}\Theta_{\Pi''}\left(\Ins_{x_1,c}[\Upsilon]\right)
=\sum_{I,J}\Upsilon_{IJ}\cdot\Theta_{x_\mu,c}\left(\Theta_{\Pi'}\Theta_{\Pi''}(\zeta_J)\right)\,\Theta_{\Pi'}\Theta_{\Pi''}(\zeta_I)
=\Ins_{x_\mu,c}[\Upsilon],
\end{equation}
the right side being formed with the modes at $(x_\mu,c)$.

\emph{The statements on coordinate link planes in the directions $e_2$ and $e_3$.}
Let $n_1,n_2,n_3\ge4$, let $\mathbb{T}=C(n_1,n_2,n_3)=(\Gamma,\chi)$, and let $\mu\in\{2,3\}$, $g:=g_\mu$.
The coordinate torus $g\mathbb{T}$ has the side $n_1$ in the direction $e_\mu$,
and its cross-section for the family $x_\mu$ is the image $g\Xi$ of the cross-section $\Xi=\Xi(\mathbb{T},x_1)$ of Definition~\ref{def:slice-notation}.

Let $A\in C^1(\Gamma)$, put $A':=(g^{-1})^{*}A\in C^1(g\Gamma)$,
and let $\Upsilon_j\in\End\mathcal{F}$, $j\in\Z/n_1\Z$, be even operators.
Since $\Theta_{\Pi'}\Theta_{\Pi''}$ is an automorphism of algebras,
Proposition~\ref{prop:cyclic-axes}\eqref{it:transport-action} and \eqref{eq:slab-insertion-transport} give
\[
\begin{aligned}
&\Theta_{\Pi'}\Theta_{\Pi''}\left(\exp\left(-S_{\mathbb{T}}(A)\right)\cdot\prod_{j\in\Z/n_1\Z}\left(\exp\left(S^{\times}_{x_1,j+1/2}(A)\right)\cdot\Ins_{x_1,j+1/2}\left[\omega^{\frac12}\Upsilon_j\,\omega^{\frac12}\right]\right)\right)\\
&\qquad=\exp\left(-S_{g\mathbb{T}}(A')\right)\cdot\prod_{j\in\Z/n_1\Z}\left(\exp\left(S^{\times}_{x_\mu,j+1/2}(A')\right)\cdot\Ins_{x_\mu,j+1/2}\left[\omega^{\frac12}\Upsilon_j\,\omega^{\frac12}\right]\right).
\end{aligned}
\]
On the left, $\Theta_{\Pi'}\Theta_{\Pi''}$ is applied to an element of $\Gr_{\mathbb{T}}$,
namely to the integrand of the left side of the identity of Proposition~\ref{prop:slice-factorisation} for $\mathbb{T}$ and $A$;
the right member, its image in $\Gr_{g\mathbb{T}}$,
is the same integrand formed on $g\mathbb{T}$ with the objects of the direction $e_\mu$ and the cochain $A'$.
By Proposition~\ref{prop:cyclic-axes}\eqref{it:transport-integral}, the Berezin integral of this image over $g\mathbb{T}$
equals the Berezin integral over $\mathbb{T}$ of the element to which $\Theta_{\Pi'}\Theta_{\Pi''}$ is applied;
the latter integral is the left side of that identity.
The right side of that identity is formed with the slices of $A$ in the direction $e_1$,
which are the slices of $A'$ in the direction $e_\mu$.
Every $A'\in C^1(g\Gamma)$ is of this form, with $A:=g^{*}A'$.
So the slice factorisation of Proposition~\ref{prop:slice-factorisation} holds on $g\mathbb{T}$ in the direction $e_\mu$, at every fixed cochain.

Let now $G\in\mathcal{I}(\mathbb{T})$ be the integrand of the left side of Proposition~\ref{prop:word-formula}\eqref{it:word-identity},
formed in the direction $e_1$ with bounded Borel functions $u_j$, electric factors $v_j$ and local step data $(S_{j,l},\Lambda_{j,l},k_{j,l})$.
The integrand formed on $g\mathbb{T}$ with the same data in the direction $e_\mu$ is the element $\Theta_{\Pi'}\Theta_{\Pi''}G$ of Lemma~\ref{lem:reflection-covariance}.
Indeed, let $(A',z')\in\mathfrak{C}(g\Gamma)$, the set of Definition~\ref{def:torus-integrands} for the period lattice $g\Gamma$, and let
\[
F':=dA'+\frac{2\pi}{q}\,z'
\]
be its curvature.
The pull-back $g^{*}$ maps $\mathfrak{C}(g\Gamma)$ onto $\mathfrak{C}(\Gamma)$ and commutes with $d$ (Definition~\ref{def:cochains}),
so that $g^{*}F'=d\,g^{*}A'+(2\pi/q)\,g^{*}z'$ is the curvature of $g^{*}(A',z'):=(g^{*}A',g^{*}z')$.
At $(A',z')$ the element $\Theta_{\Pi'}\Theta_{\Pi''}G$ takes the value $\Theta_{\Pi'}\Theta_{\Pi''}\left(G(g^{*}A',g^{*}z')\right)$.
The bosonic factors $u_j(F^{\mathrm{mag}}_j)$ and $v_j(F^{\mathrm{el}}_j)$ become the same functions of the magnetic and electric parts of $g^{*}F'$,
that is, of $F'$ in the direction $e_\mu$, the two complex conjugations of \eqref{eq:setting-reflection-coefficients} cancelling.
The integrand $G[k]$ of a local step datum becomes $G[k]^{g}=\Theta_{\Pi'}\Theta_{\Pi''}G[k]$ \eqref{eq:setting-group-action},
the integrand of the datum in the direction $e_\mu$;
by Proposition~\ref{prop:cyclic-axes}\eqref{it:transport-action}, applied to the single terms of the bonds that cross the plane $\{x_1=j+1/2\}$, and by \eqref{eq:slab-insertion-transport},
it is given by the formula of Proposition~\ref{prop:word-formula}
with the objects of the direction $e_\mu$, evaluated at the local lift $(g^{-1})^{*}\hat A$ of $(A',z')$, where $\hat A$ is a local lift of $g^{*}(A',z')$.
By Lemma~\ref{lem:reflection-covariance},
\[
\Theta_{\Pi'}\Theta_{\Pi''}G\in\mathcal{I}(g\mathbb{T}),
\qquad
Z_{g\mathbb{T}}\left[\Theta_{\Pi'}\Theta_{\Pi''}G\right]=Z_{\mathbb{T}}[G];
\]
for an integrand in $\mathcal{A}(\R^3)$ that embeds in $\mathbb{T}$ this is also contained in Proposition~\ref{prop:symmetries}\eqref{it:functional-covariance}.
In the same way, let $\beta'$ be a real two-cochain on $g\mathbb{T}$ and $\beta:=g^{*}\beta'$.
The bosonic integrand of Proposition~\ref{prop:external-field} on $g\mathbb{T}$ is the function
$(A',z')\mapsto\exp\left(-\langle\beta',F'\rangle-\tfrac12\|\beta'\|^2\right)$ on $\mathfrak{C}(g\Gamma)$, where $F'=dA'+(2\pi/q)\,z'$.
It is the image under $\Theta_{\Pi'}\Theta_{\Pi''}$ of the function $(A,z)\mapsto\exp\left(-\langle\beta,F\rangle-\tfrac12\|\beta\|^2\right)$ on $\mathfrak{C}(\Gamma)$,
where $F=dA+(2\pi/q)\,z$:
at $(A',z')$ this image takes the scalar value $\exp\left(-\langle\beta,g^{*}F'\rangle-\tfrac12\|\beta\|^2\right)$, the two complex conjugations cancelling,
and $\langle\beta,g^{*}F'\rangle=\langle\beta',F'\rangle$ and $\|\beta\|=\|\beta'\|$ because $g^{*}$ is orthogonal.
The electric and magnetic parts of $\beta'$ in the direction $e_\mu$ are those of $\beta$ in the direction $e_1$;
Lemma~\ref{lem:reflection-covariance} again gives equal left sides.
The right sides of Propositions~\ref{prop:word-formula} and~\ref{prop:external-field} are traces of products of the operators of Definition~\ref{def:slice-operators}
and of operators whose kernels are formed with the slice data;
they are the same in both directions.
The conditions \eqref{it:step-disjoint}, \eqref{it:step-string} and \eqref{it:step-free} of Proposition~\ref{prop:word-formula},
written for the data in the direction $e_\mu$, are the conditions for the data in the direction $e_1$
in which the boxes $B(j+1/2,S_{j,l},\Lambda_{j,l})$ and the lattice $\Gamma$ are replaced by their images $gB(j+1/2,S_{j,l},\Lambda_{j,l})$ and $g\Gamma$,
and the plaquettes of $\Z^3$ that lie in a plane $\{x_1=s\}$, $s\in\Z$, and represent $\sigma_0$ are replaced by their images under $g$, which lie in the planes $\{x_\mu=s\}$;
the condition \eqref{it:step-free} is unchanged.
Since $g$ is a linear isometry of $\R^3$ that maps $\Z^3$ and the cube $[-1,1]^3$ onto themselves and $\Gamma$ onto $g\Gamma$,
the two sets of conditions are equivalent.
Finally, for a local step datum $(S,\Lambda,k)$ at the plane $\{x_1=j+1/2\}$,
\[
G[k]^{g}\in\mathcal{G}^{\mathrm{ev}}\left(gB(j+1/2,S,\Lambda)\right),
\]
with the same constant $\varepsilon_G=\min\left(\varepsilon_k,\tfrac12\right)$ as in Proposition~\ref{prop:word-formula}\eqref{it:step-class}.
Indeed, for every mirror $\Pi$ and every box $B$ the map $\Theta_\Pi$ carries $\mathcal{G}^{\mathrm{ev}}(B)$ onto $\mathcal{G}^{\mathrm{ev}}(\theta_\Pi B)$ (Lemma~\ref{lem:reflection-classes}),
without change of the growth constant:
for a real cochain $\hat A$ on the links of $\theta_\Pi B$ one has
$\sum_{\sigma\in B^{(2)}}\left(d\theta_\Pi^{*}\hat A\right)_\sigma^2=\sum_{\sigma\in(\theta_\Pi B)^{(2)}}\left(d\hat A\right)_\sigma^2$,
and $\Theta_\Pi$ maps a monomial $\Psi$ to a linear combination of at most $4^{\deg\Psi}$ monomials with coefficients of modulus at most one.
Applying this with the mirror $\Pi''$ to the box $B(j+1/2,S,\Lambda)$, and then with the mirror $\Pi'$ to the box $\theta_{\Pi''}B(j+1/2,S,\Lambda)$,
whose image under $\theta_{\Pi'}$ is $gB(j+1/2,S,\Lambda)$,
one obtains the assertion, since $G[k]^{g}=\Theta_{\Pi'}\Theta_{\Pi''}G[k]$.
Like $G[k]$, the element $G[k]^{g}$ is defined on $\Z^3$ and does not depend on a torus.

Lemma~\ref{lem:half-slab-matrices}, Proposition~\ref{prop:localisation} and Corollary~\ref{cor:trace-class} do not involve the time direction.
The lemma concerns the half-slab matrices, the representation $u$ and the operators of Definition~\ref{def:slice-operators};
the proposition concerns the modes and the Fock spaces of an arbitrary mirror (Definition~\ref{def:index-sets});
the corollary concerns the operators $K_{\beta,S}^{1/2}\mathfrak{M}$, $K^{1/2}\mathfrak{M}\mathfrak{M}^{*}K^{1/2}$ and $\mathfrak{M}^{*}K\mathfrak{M}$ of a rectangular cross-section;
and none of these objects changes with the direction.
Hence Propositions~\ref{prop:slice-factorisation}, \ref{prop:localisation}, \ref{prop:word-formula} and~\ref{prop:external-field},
Lemma~\ref{lem:half-slab-matrices} and Corollary~\ref{cor:trace-class}
hold for the coordinate tori $g_\mu C(n_1,n_2,n_3)$ in the time directions $e_2$ and $e_3$,
with the objects of these directions described above.

\subsection{Positivity of trace forms}\label{sec:trace-positivity}

Let $\ell\in\mathcal{L}$, let $\Xi=(\Gamma_\parallel,\chi_\parallel)$ be a cross-section for $\ell$,
and let $c<c'$ lie in $\Adm(\ell)$;
then $2c-c'$ lies in $\Adm(\ell)$ as well, $\Adm(\ell)$ being a subgroup of $\R$.
Because $c\in\Adm(\ell)$, the reflection $\theta_c=\theta_{\ell,c}$ maps $\Z^3$ onto itself;
its linear part $\vec\theta_c$ lies in $\Ocub$, so that it maps cells onto cells of the same dimension.
The reflection fixes the plane $\{\ell=c\}$ pointwise,
and from $\ell\circ\theta_c=2c-\ell$ one reads off that it maps the closed slab $\{c\le\ell\le c'\}$ onto $\{2c-c'\le\ell\le c\}$
and the plane $\{\ell=c'\}$ onto $\{\ell=2c-c'\}$.
Moreover $\vec\theta_c$ fixes every vector of $\{\ell=0\}$,
so that $\theta_c(x+\gamma)=\theta_c(x)+\gamma$ for $x\in\R^3$ and $\gamma\in\Gamma_\parallel$.
Consequently $\theta_c$ acts on the cells of the cylinder $\Z^3/\Gamma_\parallel$;
preserving the dimension of cells, it maps $\mathcal{M}(c)$ onto itself,
$\mathcal{B}(c,c')$ onto $\mathcal{B}(2c-c',c)$
and $\mathcal{M}(c')$ onto $\mathcal{M}(2c-c')$.
The pull-back $\theta_c^{*}$ commutes with the coboundary, by~\eqref{eq:slab-incidence-isometry},
and $\left(\theta_c^{*}\right)^2=\left(\theta_c^2\right)^{*}$ is the identity.
For a configuration $\varpi'=(A',z')$ of the closed slab $\{2c-c'\le\ell\le c\}$
the condition $dz'=0$ on $\mathcal{B}^3(2c-c',c)$
is therefore equivalent to $d\left(\theta_c^{*}z'\right)=0$ on $\mathcal{B}^3(c,c')$,
and $\varpi'\mapsto\theta_c^{*}\varpi'=\left(\theta_c^{*}A',\theta_c^{*}z'\right)$ (Definition~\ref{def:slab-configurations})
is a bijection of $\mathfrak{C}[2c-c',c]$ onto $\mathfrak{C}[c,c']$.
For $\varpi'\in\mathfrak{C}[2c-c',c]$ and every integer one-cochain $k$ on $[2c-c',c]^{(1)}$,
\[
F\left(\theta_c^{*}\varpi'\right)=\theta_c^{*}F(\varpi'),
\qquad
\theta_c^{*}\left(\varpi'+[k]\right)=\theta_c^{*}\varpi'+\left[\theta_c^{*}k\right],
\]
and $\theta_c^{*}k$ is an integer one-cochain on $[c,c']^{(1)}$.

\begin{definition}[Reflection in the lower mirror]\label{def:lower-reflection}
Let $\ell\in\mathcal{L}$, let $\Xi$ be a cross-section for $\ell$, and let $c<c'$ lie in $\Adm(\ell)$.
The reflection in the lower mirror of $G\in\mathcal{I}_\Xi[c,c']$ is the element
$\Theta_cG\in\mathcal{I}_\Xi[2c-c',c]$ given by
\[
\left(\Theta_cG\right)(\varpi'):=\Theta_c\left(G\left(\theta_c^{*}\varpi'\right)\right),
\qquad \varpi'\in\mathfrak{C}[2c-c',c],
\]
where $\Theta_c$ on the right is the antilinear, product-reversing map
$\Gr_\Xi\left((c,c')\right)\to\Gr_\Xi\left((2c-c',c)\right)$ of Lemma~\ref{lem:slab-automorphisms}.
The map $G\mapsto\Theta_cG$ is antilinear and preserves the parity,
and it is inverse to the reflection $\Theta_c\colon\mathcal{I}_\Xi[2c-c',c]\to\mathcal{I}_\Xi[c,c']$
of the slab $[2c-c',c]$ in its upper mirror (Definition~\ref{def:slab-integrands}):
\[
\Theta_c\left(\Theta_cG\right)=G\quad\left(G\in\mathcal{I}_\Xi[c,c']\right),
\qquad
\Theta_c\left(\Theta_cH\right)=H\quad\left(H\in\mathcal{I}_\Xi[2c-c',c]\right).
\]
\end{definition}

The reflection in the lower mirror has the properties stated in Definition~\ref{def:lower-reflection}, for the following reasons.
Fix $G\in\mathcal{I}_\Xi[c,c']$,
and let $C$ and $\varepsilon_G$ be constants for $G$ as in Definition~\ref{def:slab-integrands}\eqref{it:integrand-growth}.
The map $\varpi'\mapsto\theta_c^{*}\varpi'$ is continuous in the real variables and bijective in the integer ones,
and the coefficients, in a monomial basis, of the image under $\Theta_c$ of an element of $\Gr_\Xi\left((c,c')\right)$
are linear combinations of the complex conjugates of the coefficients of that element;
hence $\Theta_cG$ is a Borel map.
For $\varpi'\in\mathfrak{C}[2c-c',c]$ and an integer one-cochain $k$ on $[2c-c',c]^{(1)}$,
the identity $\theta_c^{*}\left(\varpi'+[k]\right)=\theta_c^{*}\varpi'+\left[\theta_c^{*}k\right]$
and the invariance of $G$ under the integer shifts (Definition~\ref{def:slab-integrands}\eqref{it:integrand-shifts}) give
\[
\left(\Theta_cG\right)\left(\varpi'+[k]\right)
=\Theta_c\left(G\left(\theta_c^{*}\varpi'+\left[\theta_c^{*}k\right]\right)\right)
=\left(\Theta_cG\right)(\varpi').
\]
Every plaquette $\sigma\in[c,c']^{(2)}$ satisfies $F_\sigma\left(\theta_c^{*}\varpi'\right)=\pm F_{\theta_c\sigma}(\varpi')$,
and $\theta_c$ maps $\mathcal{B}^2(c,c')$ bijectively onto $\mathcal{B}^2(2c-c',c)$; therefore
\[
\sum_{\sigma\in\mathcal{B}^2(c,c')}F_\sigma\left(\theta_c^{*}\varpi'\right)^2
=\sum_{\sigma\in\mathcal{B}^2(2c-c',c)}F_\sigma(\varpi')^2 .
\]
The antilinear map $\Theta_c$ acts between the finite-dimensional spaces
$\Gr_\Xi\left((c,c')\right)$ and $\Gr_\Xi\left((2c-c',c)\right)$;
so there is a constant $C'$, depending on the slab only,
such that every coefficient of the image of an element under $\Theta_c$ has modulus at most $C'$
times the largest modulus of a coefficient of that element, both taken in monomial bases.
It follows that $\Theta_cG$ obeys the bound of Definition~\ref{def:slab-integrands}\eqref{it:integrand-growth},
for the slab $[2c-c',c]$, with the constants $C'\cdot C$ and $\varepsilon_G$,
and $\Theta_cG\in\mathcal{I}_\Xi[2c-c',c]$.
Antilinearity is inherited from $\Theta_c$ on the Grassmann algebras.
By~\eqref{eq:setting-reflection-fields} the image of a generator is a linear combination of generators,
and $\Theta_c$ reverses products;
so a product of generators goes to a linear combination of products of as many generators,
and the parity is preserved.
Finally, let $G\in\mathcal{I}_\Xi[c,c']$ and $\varpi\in\mathfrak{C}[c,c']$.
The reflection of Definition~\ref{def:slab-integrands} gives
\[
\left(\Theta_c\left(\Theta_cG\right)\right)(\varpi)
=\Theta_c\left(\left(\Theta_cG\right)\left(\theta_c^{*}\varpi\right)\right)
=\Theta_c\left(\Theta_c\left(G\left(\left(\theta_c^{*}\right)^2\varpi\right)\right)\right)
=G(\varpi),
\]
because $\left(\theta_c^{*}\right)^2$ is the identity
and $\Theta_c^2=1$ on the Grassmann algebras of the cylinder
(Lemma~\ref{lem:slab-automorphisms}\eqref{it:parallel-reflections}).
For $H\in\mathcal{I}_\Xi[2c-c',c]$ the same computation,
with the two reflections applied in the opposite order, gives $\Theta_c\left(\Theta_cH\right)=H$.

For $G\in\mathcal{A}_\Xi[c,c']$ let $G(\cdot)\in\mathcal{I}_\Xi[c,c']$ be the element attached to $G$
by Lemma~\ref{lem:product-integrands}\eqref{it:sum-slab}.
By that lemma $\Theta_{\ell,c}G$ belongs to $\mathcal{A}_\Xi[2c-c',c]$,
and its identity for the reflection in the mirror at $c$ says that the element of $\mathcal{I}_\Xi[2c-c',c]$
attached to $\Theta_{\ell,c}G$ is the reflection $\Theta_c\left(G(\cdot)\right)$ of Definition~\ref{def:lower-reflection}.

\begin{corollary}[Trace forms]\label{cor:trace-forms}
Let $-1\le\rho\le1$, $M\in\R$, $w_0,\dots,w_3>0$ and $q>0$.
\begin{enumerate}
\item\label{it:trace-form}
Let $\ell\in\mathcal{L}$, let $\mathbb{T}$ be a torus adapted to $\ell$ (Definition~\ref{def:slabs}),
let $p:=p(\mathbb{T},\ell)$ and $\Xi:=\Xi(\mathbb{T},\ell)$,
and let $c$ and $c+\tfrac{p}{2}$ lie in $\Adm(\ell)$.
For all $G,G'\in\mathcal{I}_\Xi\left[c,c+\tfrac{p}{2}\right]$, even, odd or neither,
\begin{equation}\label{eq:slab-trace-form}
Z_{\mathbb{T}}\left[\Theta_c(G')\cdot G\right]
=\Tr_{\mathcal{H}(c)}\left(\mathfrak{S}[G']^{*}\cdot\mathfrak{S}[G]\right)
=\left\langle\mathfrak{S}[G'],\mathfrak{S}[G]\right\rangle_{\mathrm{HS}} .
\end{equation}
Here $\Theta_c(G')\in\mathcal{I}_\Xi\left[c-\tfrac{p}{2},c\right]$ is the reflection of $G'$ in the lower mirror
(Definition~\ref{def:lower-reflection}),
the two factors stand in the order $\Theta_c(G')\cdot G$,
and their product is read on $\mathbb{T}$ as in Theorem~\ref{thm:slab-operators}\eqref{it:slab-trace}
for the positions $c-\tfrac{p}{2}<c<c+\tfrac{p}{2}$.
The map $(G',G)\mapsto Z_{\mathbb{T}}\left[\Theta_c(G')\cdot G\right]$
is a positive semidefinite Hermitian form on $\mathcal{I}_\Xi\left[c,c+\tfrac{p}{2}\right]$.
In particular, \eqref{eq:slab-trace-form} holds for $G,G'\in\mathcal{A}_\Xi\left[c,c+\tfrac{p}{2}\right]$,
with $\Theta_{\ell,c}G'$ in the place of $\Theta_c(G')$,
the left side read as in Corollary~\ref{cor:slab-admissible},
and the slab operators of the elements of $\mathcal{I}_\Xi\left[c,c+\tfrac{p}{2}\right]$ attached to $G$ and $G'$
(Definition~\ref{def:slab-kernels}).
\item\label{it:weighted-trace-form}
Let $\ell$, $\mathbb{T}$, $p$, $\Xi$, $c$, $G$ and $G'$ be as in~\eqref{it:trace-form},
and let $\Phi_+$ be a bosonic element of $\mathcal{I}_\Xi\left[c,c+\tfrac{p}{2}\right]$,
that is, one whose values are multiples of the unit,
such that the pointwise products $\Phi_+\cdot G$ and $\Phi_+\cdot G'$ lie in $\mathcal{I}_\Xi\left[c,c+\tfrac{p}{2}\right]$.
Then
\[
Z_{\mathbb{T}}\left[\Phi_+\cdot\Theta_c(\Phi_+)\cdot\Theta_c(G')\cdot G\right]
=\Tr_{\mathcal{H}(c)}\left(\mathfrak{S}\left[\Phi_+\cdot G'\right]^{*}\cdot\mathfrak{S}\left[\Phi_+\cdot G\right]\right),
\]
and the form of~\eqref{it:trace-form} with the weight multiplied by $\Phi_+\cdot\Theta_c(\Phi_+)$
is positive semidefinite and Hermitian.
The hypothesis on $\Phi_+$ holds for all $G,G'\in\mathcal{I}_\Xi\left[c,c+\tfrac{p}{2}\right]$
whenever $\Phi_+$ is a bounded Borel function of finitely many $F_\sigma$,
$\sigma$ a plaquette in the open slab $\left\{c<\ell<c+\tfrac{p}{2}\right\}$.
If the values of $\Phi_+$ are real multiples of the unit, then, for every $\varpi'\in\mathfrak{C}\left[c-\tfrac{p}{2},c\right]$,
the value $\left(\Theta_c(\Phi_+)\right)(\varpi')$ is the same multiple of the unit as $\Phi_+\left(\theta_c^{*}\varpi'\right)$;
in particular, if $\Phi_+$ is a function of the curvature,
that is, if $\Phi_+(\varpi)$ depends on $\varpi\in\mathfrak{C}\left[c,c+\tfrac{p}{2}\right]$ only through $F(\varpi)$,
then $\left(\Theta_c(\Phi_+)\right)(\varpi')$ is the same function of $\theta_c^{*}F(\varpi')$ as $\Phi_+(\varpi)$ is of $F(\varpi)$.
\item\label{it:unit-step}
Let $\ell=x_1$, let $n_2,n_3\ge4$ be integers,
let $\Xi:=\left(n_2\Z e_2+n_3\Z e_3,\chi_\parallel\right)$ with $\chi_\parallel(n_2e_2)=\chi_\parallel(n_3e_3)=-1$,
and let $c\in\tfrac12\Z$.
For positions $c'<c''$ in $\tfrac12\Z$ let $\mathfrak{S}[1]_{[c',c'']}\colon\mathcal{H}(c')\to\mathcal{H}(c'')$
be the slab operator of the constant integrand $1\in\mathcal{I}_\Xi[c',c'']$,
and let $\mathfrak{S}_{1/2}:=\mathfrak{S}[1]_{[c,c+1/2]}$ be the half step.
Then $\mathfrak{S}_{1/2}^{*}\mathfrak{S}_{1/2}$ is a non-negative operator of trace class on $\mathcal{H}(c)$.
It is formed with $\Xi$, $c$ and the numbers $\rho$, $M$, $w_0,\dots,w_3$, $q$ alone;
thus it depends only on $n_2$, $n_3$, $c$ and these numbers, and on $c$ only modulo $1$, in the sense that
\[
V_{e_1}\cdot\mathfrak{S}_{1/2}^{*}\mathfrak{S}_{1/2}\cdot V_{e_1}^{-1}
=\mathfrak{S}[1]_{[c+1,c+3/2]}^{*}\cdot\mathfrak{S}[1]_{[c+1,c+3/2]} .
\]
Moreover
\[
\mathfrak{S}[1]_{[c+1/2,c+1]}=V_{e_1}\cdot\mathfrak{S}_{1/2}^{*},
\qquad
\mathfrak{S}[1]_{[c,c+1]}=V_{e_1}\cdot\mathfrak{S}_{1/2}^{*}\mathfrak{S}_{1/2},
\]
and for every integer $n_1\ge4$
\[
Z_{C(n_1,n_2,n_3)}[1]=\Tr_{\mathcal{H}(c)}\left(\left(\mathfrak{S}_{1/2}^{*}\mathfrak{S}_{1/2}\right)^{n_1}\right),
\]
with one and the same operator $\mathfrak{S}_{1/2}^{*}\mathfrak{S}_{1/2}$ for all $n_1$.
The same holds for $\ell=x_2$ and for $\ell=x_3$,
with $e_2$ and $n_2$, respectively $e_3$ and $n_3$, in the places of $e_1$ and $n_1$,
and with the cross-section formed in the same way from the two other sides.
\end{enumerate}
\end{corollary}

\begin{proof}
\eqref{it:trace-form}
As $\mathbb{T}=(\Gamma,\chi)$ is adapted to $\ell$, it is its own straightened torus (Definition~\ref{def:slabs}):
the vector $g:=\frac{p}{|n_\ell|^2}\cdot n_\ell$ lies in $\Gamma$,
$\Gamma=\left(\Gamma\cap\{\ell=0\}\right)+\Z g$, and $\chi(g)=-1$.
Since $\ell(g)=p$, the vector $g$ is a closing vector of $\mathbb{T}$ for $\ell$.
The position $c-\tfrac{p}{2}=2c-\left(c+\tfrac{p}{2}\right)$ lies in $\Adm(\ell)$,
and by Definition~\ref{def:lower-reflection}
the element $\Theta_c(G')$ lies in $\mathcal{I}_\Xi\left[c-\tfrac{p}{2},c\right]$,
while the reflection of the slab $\left[c-\tfrac{p}{2},c\right]$ in its upper mirror maps it back to $G'$.
The vector of this slab in the reflection rule is
$v_{c-p/2,c}=\frac{p}{|n_\ell|^2}\cdot n_\ell=g$.
Theorem~\ref{thm:slab-operators}\eqref{it:slab-reflection}, applied to $\Theta_c(G')$,
therefore gives $\mathfrak{S}[G']=V_g\cdot\mathfrak{S}\left[\Theta_c(G')\right]^{*}$,
and, $V_g$ being unitary, taking adjoints yields
\[
\mathfrak{S}\left[\Theta_c(G')\right]=\mathfrak{S}[G']^{*}\cdot V_g .
\]
Suppose first that $G$ and $G'$ are homogeneous;
then so is $\Theta_c(G')$, by Definition~\ref{def:lower-reflection}.
The trace formula of Theorem~\ref{thm:slab-operators}\eqref{it:slab-trace}
applies to the closing vector $g$, the positions $c_0=c-\tfrac{p}{2}<c_1=c<c_2=c+\tfrac{p}{2}=c_0+p$
and the elements $G_1=\Theta_c(G')$, $G_2=G$;
its parity factor $U_{\mathrm{f}}^{\frac{1+\chi(g)}{2}}$ is the identity, as $\chi(g)=-1$.
The trace is invariant under conjugation by the unitary $V_g\colon\mathcal{H}(c-p/2)\to\mathcal{H}(c+p/2)$,
and the trace of the product of two Hilbert--Schmidt operators, one from a Hilbert space to another and one back,
does not depend on the order of the two factors
(Lemma~\ref{lem:poisson-kernels}\eqref{it:hilbert-schmidt-kernels}).
One obtains
\[
\begin{aligned}
Z_{\mathbb{T}}\left[\Theta_c(G')\cdot G\right]
&=\Tr_{\mathcal{H}(c-p/2)}\left(V_g^{-1}\cdot\mathfrak{S}[G]\cdot\mathfrak{S}\left[\Theta_c(G')\right]\right)
=\Tr_{\mathcal{H}(c-p/2)}\left(V_g^{-1}\cdot\mathfrak{S}[G]\cdot\mathfrak{S}[G']^{*}\cdot V_g\right)\\
&=\Tr_{\mathcal{H}(c+p/2)}\left(\mathfrak{S}[G]\cdot\mathfrak{S}[G']^{*}\right)
=\Tr_{\mathcal{H}(c)}\left(\mathfrak{S}[G']^{*}\cdot\mathfrak{S}[G]\right),
\end{aligned}
\]
and the last trace is the Hilbert--Schmidt inner product $\left\langle\mathfrak{S}[G'],\mathfrak{S}[G]\right\rangle_{\mathrm{HS}}$.
In general, both sides of~\eqref{eq:slab-trace-form} are linear in $G$ and antilinear in $G'$:
the left side because $\Theta_c$ is antilinear, the product read on $\mathbb{T}$ is bilinear and $Z_{\mathbb{T}}$ is linear,
the right side because $G\mapsto\mathfrak{S}[G]$ is linear
(Theorem~\ref{thm:slab-operators}\eqref{it:slab-hilbert-schmidt}).
The even and the odd part of an element of $\mathcal{I}_\Xi\left[c,c+\tfrac{p}{2}\right]$ belong to this space
(Definition~\ref{def:slab-integrands}),
so~\eqref{eq:slab-trace-form} follows from the homogeneous case by expanding both sides.
The Hilbert--Schmidt inner product is a positive semidefinite Hermitian form,
and composing it in both arguments with the linear map $G\mapsto\mathfrak{S}[G]$ keeps both properties.

Let now $G,G'\in\mathcal{A}_\Xi\left[c,c+\tfrac{p}{2}\right]$.
By Lemma~\ref{lem:product-integrands}\eqref{it:sum-slab}, applied to the slab $\left[c,c+\tfrac{p}{2}\right]$,
$\Theta_{\ell,c}G'$ lies in $\mathcal{A}_\Xi\left[c-\tfrac{p}{2},c\right]$,
and the element of $\mathcal{I}_\Xi\left[c-\tfrac{p}{2},c\right]$ attached to it
arises from the element attached to $G'$ by the reflection $\Theta_c$ of Definition~\ref{def:lower-reflection};
the slab operators of elements of $\mathcal{A}_\Xi$ are those of the attached elements (Definition~\ref{def:slab-kernels}).
The trace formula for admissible products, Corollary~\ref{cor:slab-admissible},
applied to the same closing vector and the same positions, gives
\[
Z_{\mathbb{T}}\left[\Theta_{\ell,c}(G')\cdot G\right]
=\Tr_{\mathcal{H}(c-p/2)}\left(V_g^{-1}\cdot\mathfrak{S}[G]\cdot\mathfrak{S}\left[\Theta_{\ell,c}G'\right]\right).
\]
The right side is the first trace of the chain above, formed with the attached elements,
and the remaining equalities of that chain lead to~\eqref{eq:slab-trace-form}.

\eqref{it:weighted-trace-form}
On the Grassmann algebras $\Theta_c$ is antilinear and reverses products,
and the values of $\Phi_+$ are multiples of the unit;
for $\varpi'\in\mathfrak{C}\left[c-\tfrac{p}{2},c\right]$ this gives
\[
\left(\Theta_c\left(\Phi_+\cdot G'\right)\right)(\varpi')
=\Theta_c\left(G'\left(\theta_c^{*}\varpi'\right)\right)\cdot\Theta_c\left(\Phi_+\left(\theta_c^{*}\varpi'\right)\right)
=\left(\Theta_c(\Phi_+)\right)(\varpi')\cdot\left(\Theta_c(G')\right)(\varpi').
\]
Multiples of the unit commute with every element of the Grassmann algebras,
so the product $\Phi_+\cdot\Theta_c(\Phi_+)\cdot\Theta_c(G')\cdot G$, read on $\mathbb{T}$,
equals $\Theta_c\left(\Phi_+\cdot G'\right)\cdot\left(\Phi_+\cdot G\right)$.
Part~\eqref{it:trace-form}, applied to $\Phi_+\cdot G$ and $\Phi_+\cdot G'$, gives the identity.
As $G\mapsto\Phi_+\cdot G$ is linear,
$(G',G)\mapsto\left\langle\mathfrak{S}\left[\Phi_+\cdot G'\right],\mathfrak{S}\left[\Phi_+\cdot G\right]\right\rangle_{\mathrm{HS}}$
is a positive semidefinite Hermitian form on the space of the $G\in\mathcal{I}_\Xi\left[c,c+\tfrac{p}{2}\right]$
with $\Phi_+\cdot G\in\mathcal{I}_\Xi\left[c,c+\tfrac{p}{2}\right]$.
Now let $\Phi_+$ be a bounded Borel function of finitely many $F_\sigma$ with $\sigma$ in the open slab.
It is a Borel map on $\mathfrak{C}\left[c,c+\tfrac{p}{2}\right]$,
invariant under the integer shifts because $F$ is,
and it obeys the bound of Definition~\ref{def:slab-integrands}\eqref{it:integrand-growth}
with the constants $\sup|\Phi_+|$ and $\tfrac12$;
so it is a bosonic element of $\mathcal{I}_\Xi\left[c,c+\tfrac{p}{2}\right]$.
For $G\in\mathcal{I}_\Xi\left[c,c+\tfrac{p}{2}\right]$ the pointwise product $\Phi_+\cdot G$ is a Borel map,
invariant under the integer shifts (Definition~\ref{def:slab-integrands}\eqref{it:integrand-shifts}) as both factors are;
and if $C$ and $\varepsilon_G$ are constants for $G$ as in Definition~\ref{def:slab-integrands}\eqref{it:integrand-growth},
then $\Phi_+\cdot G$ obeys the same bound with the constants $\sup|\Phi_+|\cdot C$ and $\varepsilon_G$.
Let, finally, the values of $\Phi_+$ be real multiples of the unit, and let $\varpi'\in\mathfrak{C}\left[c-\tfrac{p}{2},c\right]$.
Since $\Theta_c$ is antilinear and maps the unit to the unit,
$\left(\Theta_c(\Phi_+)\right)(\varpi')=\Theta_c\left(\Phi_+\left(\theta_c^{*}\varpi'\right)\right)$
is the unit multiplied by the complex conjugate of the real factor of $\Phi_+\left(\theta_c^{*}\varpi'\right)$, that is, by this factor itself.
If $\Phi_+(\varpi)$ depends on $\varpi\in\mathfrak{C}\left[c,c+\tfrac{p}{2}\right]$ only through $F(\varpi)$,
then $\Phi_+\left(\theta_c^{*}\varpi'\right)$ is the corresponding function of $F\left(\theta_c^{*}\varpi'\right)=\theta_c^{*}F(\varpi')$.

\eqref{it:unit-step}
The constant $1$ belongs to $\mathcal{A}_\Xi[c',c'']$, as the empty product, and to $\mathcal{I}_\Xi[c',c'']$,
for all $c'<c''$ in $\tfrac12\Z$.
The translations and the reflections of slab integrands fix it,
since $\tau_v$ and $\Theta_{c''}$ map the unit of the Grassmann algebras to the unit;
in particular $\Theta_{c+1/2}1=1$.
The vector of the slab $\left[c,c+\tfrac12\right]$ in the reflection rule is $v_{c,c+1/2}=e_1$,
and Theorem~\ref{thm:slab-operators}\eqref{it:slab-reflection} gives
\[
\mathfrak{S}[1]_{[c+1/2,c+1]}=\mathfrak{S}\left[\Theta_{c+1/2}1\right]=V_{e_1}\cdot\mathfrak{S}_{1/2}^{*}.
\]
The product of $1\in\mathcal{I}_\Xi\left[c,c+\tfrac12\right]$ and $1\in\mathcal{I}_\Xi\left[c+\tfrac12,c+1\right]$
is $1\in\mathcal{I}_\Xi[c,c+1]$,
and the composition rule of Theorem~\ref{thm:slab-operators}\eqref{it:slab-composition} gives
\[
\mathfrak{S}[1]_{[c,c+1]}=\mathfrak{S}[1]_{[c+1/2,c+1]}\cdot\mathfrak{S}_{1/2}=V_{e_1}\cdot\mathfrak{S}_{1/2}^{*}\mathfrak{S}_{1/2}.
\]
For $j\in\Z$, the covariance of Theorem~\ref{thm:slab-operators}\eqref{it:slab-covariance} with $v=je_1$ gives
\[
\mathfrak{S}[1]_{[c+j,c+j+1]}
=V_{je_1}\cdot\mathfrak{S}[1]_{[c,c+1]}\cdot V_{je_1}^{-1}
=V_{je_1}\cdot V_{e_1}\cdot\mathfrak{S}_{1/2}^{*}\mathfrak{S}_{1/2}\cdot V_{je_1}^{-1},
\]
and with $v=e_1$ and the slab $\left[c,c+\tfrac12\right]$ it gives
$\mathfrak{S}[1]_{[c+1,c+3/2]}=V_{e_1}\cdot\mathfrak{S}_{1/2}\cdot V_{e_1}^{-1}$.
The operators $V_{e_1}$ are unitary, so that
\[
\mathfrak{S}[1]_{[c+1,c+3/2]}^{*}\cdot\mathfrak{S}[1]_{[c+1,c+3/2]}
=V_{e_1}\cdot\mathfrak{S}_{1/2}^{*}\cdot V_{e_1}^{-1}\cdot V_{e_1}\cdot\mathfrak{S}_{1/2}\cdot V_{e_1}^{-1}
=V_{e_1}\cdot\mathfrak{S}_{1/2}^{*}\mathfrak{S}_{1/2}\cdot V_{e_1}^{-1}.
\]
Let $n_1\ge4$ be an integer and $\mathbb{T}:=C(n_1,n_2,n_3)$ (Definition~\ref{def:tori}).
Its cross-section for $x_1$ is $\Xi$;
the vector $g:=n_1e_1$ is a closing vector with $\ell(g)=n_1$ and $\chi(g)=-1$;
and the positions $c_j:=c+j$, $j=0,\dots,n_1$, satisfy $c_{n_1}=c_0+n_1$ and bound $n_1\ge2$ slabs.
In the order of the trace formula, the operators of these slabs multiply to
\[
\begin{aligned}
\mathfrak{S}[1]_{[c+n_1-1,c+n_1]}\cdots\mathfrak{S}[1]_{[c,c+1]}
&=\left(V_{(n_1-1)e_1}\cdot V_{e_1}\cdot\mathfrak{S}_{1/2}^{*}\mathfrak{S}_{1/2}\cdot V_{(n_1-1)e_1}^{-1}\right)\cdots
\left(V_{e_1}\cdot\mathfrak{S}_{1/2}^{*}\mathfrak{S}_{1/2}\right)\\
&=V_{n_1e_1}\cdot\left(\mathfrak{S}_{1/2}^{*}\mathfrak{S}_{1/2}\right)^{n_1},
\end{aligned}
\]
because between the factors with the indices $j$ and $j-1$ there stands $V_{je_1}^{-1}\cdot V_{(j-1)e_1}\cdot V_{e_1}=1$,
and $V_{(n_1-1)e_1}\cdot V_{e_1}=V_{n_1e_1}$,
by the group law of Theorem~\ref{thm:slab-operators}\eqref{it:slab-hilbert-schmidt}.
The product of the constant integrands, read on $\mathbb{T}$, is the constant $1$.
The trace formula of Theorem~\ref{thm:slab-operators}\eqref{it:slab-trace},
whose parity factor is the identity because $\chi(g)=-1$, now gives
\[
Z_{C(n_1,n_2,n_3)}[1]
=\Tr_{\mathcal{H}(c)}\left(V_{n_1e_1}^{-1}\cdot V_{n_1e_1}\cdot\left(\mathfrak{S}_{1/2}^{*}\mathfrak{S}_{1/2}\right)^{n_1}\right)
=\Tr_{\mathcal{H}(c)}\left(\left(\mathfrak{S}_{1/2}^{*}\mathfrak{S}_{1/2}\right)^{n_1}\right).
\]
The half step $\mathfrak{S}_{1/2}$ is a Hilbert--Schmidt operator
(Theorem~\ref{thm:slab-operators}\eqref{it:slab-hilbert-schmidt});
hence $\mathfrak{S}_{1/2}^{*}\mathfrak{S}_{1/2}$ is non-negative and of trace class.
The space $\mathcal{H}(c)$ and the operator $\mathfrak{S}_{1/2}$ are defined through $\ell$, $\Xi$, $c$ and the numbers $\rho$, $M$, $w_0,\dots,w_3$, $q$ alone
(Definitions~\ref{def:mirror-spaces} and~\ref{def:slab-kernels});
the side $n_1$ does not enter them, so the same operator serves for every $n_1\ge4$.
For $\ell=x_2$ the cross-section of $C(n_1,n_2,n_3)$ is $\left(n_1\Z e_1+n_3\Z e_3,\chi_\parallel\right)$
with $\chi_\parallel(n_1e_1)=\chi_\parallel(n_3e_3)=-1$,
a closing vector is $n_2e_2$, and the vector of the slab $\left[c,c+\tfrac12\right]$ in the reflection rule is $e_2$;
for $\ell=x_3$ these data are $\left(n_1\Z e_1+n_2\Z e_2,\chi_\parallel\right)$
with $\chi_\parallel(n_1e_1)=\chi_\parallel(n_2e_2)=-1$, the vector $n_3e_3$ and the vector $e_3$.
Theorem~\ref{thm:slab-operators} holds for every family,
and with these data each step above remains valid as written,
with $e_2,n_2$, respectively $e_3,n_3$, in the places of $e_1,n_1$.
\end{proof}

\section{Comparison of tori}\label{sec:tori}

Throughout this section the Wilson parameter satisfies $-1\le\rho\le1$,
the hopping weights $w_0,\dots,w_3$ are positive,
the mass parameter $M$ is real,
the lattice charge $q$ is positive,
and the base point $v_0\in C_0$ of the fermion sites is fixed.
For every torus $\mathbb{T}$ one has $0<Z_{\mathbb{T}}[1]<\infty$ by Theorem~\ref{thm:partition-positivity}.

\subsection{Schatten norms and admissible products}\label{sec:schatten-norms}

Let $\mathcal{H}$ be a separable Hilbert space, and let $1\le s<\infty$.
The Schatten class $\mathscr{J}_s=\mathscr{J}_s(\mathcal{H})$ consists of the compact operators $\mathfrak{Y}$ on $\mathcal{H}$ with
\[
\|\mathfrak{Y}\|_s:=\left(\Tr|\mathfrak{Y}|^s\right)^{\frac{1}{s}}<\infty,
\qquad
|\mathfrak{Y}|:=\left(\mathfrak{Y}^{*}\cdot\mathfrak{Y}\right)^{\frac12}.
\]
Thus $\mathscr{J}_1$ is the trace class, and $\mathscr{J}_2$ is the class of the Hilbert--Schmidt operators.
In \cite[Section~10.3.1]{NoncompactQED} the exponent is denoted $p$.

\begin{lemma}[{Schatten norms and traces \cite[Lemma~10.6]{NoncompactQED}}]\label{lem:schatten-traces}
Let $\mathcal{H}$ be a separable Hilbert space; all Schatten classes below are those of $\mathcal{H}$.
\begin{enumerate}
\item\label{it:schatten-holder}
Let $k\ge1$, let $s_1,\dots,s_k\in[1,\infty)$ with $\sum_{i=1}^{k}\frac{1}{s_i}=1$,
and let $\mathfrak{Y}_i\in\mathscr{J}_{s_i}$ for $i=1,\dots,k$.
Then $\mathfrak{Y}_1\cdots\mathfrak{Y}_k$ is of trace class, and
\[
\left|\Tr\left(\mathfrak{Y}_1\cdots\mathfrak{Y}_k\right)\right|\le\prod_{i=1}^{k}\|\mathfrak{Y}_i\|_{s_i}.
\]
\item\label{it:schatten-monotone}
Let $1\le s'\le s<\infty$.
Every $\mathfrak{Y}\in\mathscr{J}_{s'}$ belongs to $\mathscr{J}_s$, and $\|\mathfrak{Y}\|_s\le\|\mathfrak{Y}\|_{s'}$.
For $\mathfrak{Y}\in\mathscr{J}_s$ and every unitary operator $U$ on $\mathcal{H}$ one has
$\|U\cdot\mathfrak{Y}\|_s=\|\mathfrak{Y}\cdot U\|_s=\|\mathfrak{Y}\|_s$ and $\|\mathfrak{Y}^{*}\|_s=\|\mathfrak{Y}\|_s$;
if $s$ is an even integer, then $\|\mathfrak{Y}\|_s^s=\Tr\left(\mathfrak{Y}^{*}\cdot\mathfrak{Y}\right)^{\frac{s}{2}}$.
\item\label{it:schatten-convex}
Let $\mathfrak{Y}\ge0$ be of trace class.
Then $\Tr\mathfrak{Y}^k\le\left(\Tr\mathfrak{Y}\right)^k$ for every integer $k\ge1$.
If moreover $\mathfrak{Y}\neq0$, the function $s\mapsto\log\Tr\mathfrak{Y}^s$ is convex on $[1,\infty)$.
\end{enumerate}
\end{lemma}

\begin{definition}[Schatten norms between two spaces]\label{def:even-schatten}
Let $\mathcal{H}$ and $\mathcal{H}'$ be Hilbert spaces,
let $\mathfrak{Y}\colon\mathcal{H}\to\mathcal{H}'$ be a Hilbert--Schmidt operator,
and let $k\ge1$ be an integer.
Put
\[
\|\mathfrak{Y}\|_{2k}:=\left(\Tr_{\mathcal{H}}\left(\mathfrak{Y}^{*}\cdot\mathfrak{Y}\right)^k\right)^{\frac{1}{2k}}.
\]
\end{definition}

The trace on the right is a finite non-negative number.
Indeed, $\mathfrak{Y}^{*}\cdot\mathfrak{Y}$ is a non-negative operator of trace class on $\mathcal{H}$, since $\mathfrak{Y}$ is a Hilbert--Schmidt operator,
and so is its $k$-th power, the product of $\mathfrak{Y}^{*}\cdot\mathfrak{Y}$ with the bounded operator $\left(\mathfrak{Y}^{*}\cdot\mathfrak{Y}\right)^{k-1}$.
If $\mathcal{H}'=\mathcal{H}$ and $\mathcal{H}$ is separable, then $\mathfrak{Y}\in\mathscr{J}_2\subset\mathscr{J}_{2k}$,
and by Lemma~\ref{lem:schatten-traces}\eqref{it:schatten-monotone} the number $\|\mathfrak{Y}\|_{2k}$ is the norm of $\mathfrak{Y}$ in $\mathscr{J}_{2k}$.

\begin{lemma}[H\"older's inequality between different spaces]\label{lem:holder-spaces}
Let $r\ge2$, let $\mathcal{H}_0,\dots,\mathcal{H}_r$ be separable Hilbert spaces,
let $\mathfrak{Y}_j\colon\mathcal{H}_{j-1}\to\mathcal{H}_j$ be Hilbert--Schmidt operators for $j=1,\dots,r$,
and let $U\colon\mathcal{H}_r\to\mathcal{H}_0$ be unitary.
Then $U\cdot\mathfrak{Y}_r\cdots\mathfrak{Y}_1$ is an operator of trace class on $\mathcal{H}_0$.
If $r=2k$ is even, then
\[
\left|\Tr_{\mathcal{H}_0}\left(U\cdot\mathfrak{Y}_{2k}\cdots\mathfrak{Y}_1\right)\right|\le\prod_{j=1}^{2k}\|\mathfrak{Y}_j\|_{2k},
\]
the norms on the right being those of Definition~\ref{def:even-schatten}.
\end{lemma}

\begin{proof}
Consider the separable Hilbert space $\mathcal{H}:=\mathcal{H}_0\oplus\cdots\oplus\mathcal{H}_{r-1}$
and the canonical isometries $\mathfrak{i}_j\colon\mathcal{H}_j\to\mathcal{H}$, $0\le j\le r-1$;
thus $\mathfrak{i}_j^{*}\cdot\mathfrak{i}_j$ is the identity of $\mathcal{H}_j$.
Define operators on $\mathcal{H}$ by
\begin{equation}\label{eq:tori-direct-sum}
\widehat{\mathfrak{Y}}_j:=\mathfrak{i}_j\cdot\mathfrak{Y}_j\cdot\mathfrak{i}_{j-1}^{*}\quad(1\le j\le r-1),
\qquad
\widehat{\mathfrak{Y}}_r:=\mathfrak{i}_0\cdot U\cdot\mathfrak{Y}_r\cdot\mathfrak{i}_{r-1}^{*}.
\end{equation}
For $0\le j\le r-1$ and an operator $\mathfrak{Y}$ of trace class on $\mathcal{H}_j$, the operator $\mathfrak{i}_j\cdot\mathfrak{Y}\cdot\mathfrak{i}_j^{*}$ is of trace class on $\mathcal{H}$, and
\[
\Tr_{\mathcal{H}}\left(\mathfrak{i}_j\cdot\mathfrak{Y}\cdot\mathfrak{i}_j^{*}\right)=\Tr_{\mathcal{H}_j}\mathfrak{Y}:
\]
the isometry $\mathfrak{i}_j$ carries an orthonormal basis of $\mathcal{H}_j$ onto an orthonormal basis of its range,
and $\mathfrak{i}_j\cdot\mathfrak{Y}\cdot\mathfrak{i}_j^{*}$ vanishes on the orthogonal complement of that range.

Since $\mathfrak{i}_j^{*}\cdot\mathfrak{i}_j=1$ for $0\le j\le r-1$, the products of consecutive operators \eqref{eq:tori-direct-sum} telescope:
\[
\widehat{\mathfrak{Y}}_r\cdots\widehat{\mathfrak{Y}}_1=\mathfrak{i}_0\cdot\left(U\cdot\mathfrak{Y}_r\cdots\mathfrak{Y}_1\right)\cdot\mathfrak{i}_0^{*}.
\]
The operator $\mathfrak{Y}_2\cdot\mathfrak{Y}_1\colon\mathcal{H}_0\to\mathcal{H}_2$ is of trace class, being the composition of two Hilbert--Schmidt operators;
so is its composition with the bounded operator $U\cdot\mathfrak{Y}_r\cdots\mathfrak{Y}_3$ on the left, which is $U$ when $r=2$.
Thus $U\cdot\mathfrak{Y}_r\cdots\mathfrak{Y}_1$ is of trace class on $\mathcal{H}_0$, which is the first assertion,
and the first paragraph, applied for $j=0$ to this operator, together with the last display gives
\[
\Tr_{\mathcal{H}_0}\left(U\cdot\mathfrak{Y}_r\cdots\mathfrak{Y}_1\right)=\Tr_{\mathcal{H}}\left(\widehat{\mathfrak{Y}}_r\cdots\widehat{\mathfrak{Y}}_1\right).
\]

Now let $r=2k$.
For $1\le j\le r$ one has
\[
\widehat{\mathfrak{Y}}_j^{*}\cdot\widehat{\mathfrak{Y}}_j=\mathfrak{i}_{j-1}\cdot\mathfrak{Y}_j^{*}\cdot\mathfrak{Y}_j\cdot\mathfrak{i}_{j-1}^{*},
\]
because $\mathfrak{i}_j^{*}\cdot\mathfrak{i}_j=1$ for $j<r$, and because $\mathfrak{i}_0^{*}\cdot\mathfrak{i}_0=1$ and $U^{*}\cdot U=1$ for $j=r$.
As $\mathfrak{i}_{j-1}^{*}\cdot\mathfrak{i}_{j-1}=1$, it follows that
\[
\left(\widehat{\mathfrak{Y}}_j^{*}\cdot\widehat{\mathfrak{Y}}_j\right)^k=\mathfrak{i}_{j-1}\cdot\left(\mathfrak{Y}_j^{*}\cdot\mathfrak{Y}_j\right)^k\cdot\mathfrak{i}_{j-1}^{*}.
\]
The first paragraph, applied to the operators $\mathfrak{Y}_j^{*}\cdot\mathfrak{Y}_j$ and $\left(\mathfrak{Y}_j^{*}\cdot\mathfrak{Y}_j\right)^k$ of trace class on $\mathcal{H}_{j-1}$,
shows that $\widehat{\mathfrak{Y}}_j^{*}\cdot\widehat{\mathfrak{Y}}_j$ is of trace class,
so that $\widehat{\mathfrak{Y}}_j$ is a Hilbert--Schmidt operator on $\mathcal{H}$,
and that $\Tr_{\mathcal{H}}\left(\widehat{\mathfrak{Y}}_j^{*}\cdot\widehat{\mathfrak{Y}}_j\right)^k=\Tr_{\mathcal{H}_{j-1}}\left(\mathfrak{Y}_j^{*}\cdot\mathfrak{Y}_j\right)^k$.
By Lemma~\ref{lem:schatten-traces}\eqref{it:schatten-monotone}, $\widehat{\mathfrak{Y}}_j$ lies in $\mathscr{J}_2(\mathcal{H})\subset\mathscr{J}_{2k}(\mathcal{H})$, and
\[
\|\widehat{\mathfrak{Y}}_j\|_{2k}^{2k}=\Tr_{\mathcal{H}}\left(\widehat{\mathfrak{Y}}_j^{*}\cdot\widehat{\mathfrak{Y}}_j\right)^k=\|\mathfrak{Y}_j\|_{2k}^{2k},
\]
the norm on the left being that of $\mathscr{J}_{2k}(\mathcal{H})$ and the one on the right that of Definition~\ref{def:even-schatten}.
Lemma~\ref{lem:schatten-traces}\eqref{it:schatten-holder}, applied on $\mathcal{H}$ to the $2k$ factors $\widehat{\mathfrak{Y}}_{2k},\dots,\widehat{\mathfrak{Y}}_1$ with all exponents equal to $2k$, now gives
\[
\left|\Tr_{\mathcal{H}_0}\left(U\cdot\mathfrak{Y}_{2k}\cdots\mathfrak{Y}_1\right)\right|
=\left|\Tr_{\mathcal{H}}\left(\widehat{\mathfrak{Y}}_{2k}\cdots\widehat{\mathfrak{Y}}_1\right)\right|
\le\prod_{j=1}^{2k}\|\widehat{\mathfrak{Y}}_j\|_{2k}
=\prod_{j=1}^{2k}\|\mathfrak{Y}_j\|_{2k}.
\]
\end{proof}

By \eqref{eq:setting-group-action}, the group $W$ acts on the even elements of the algebra $\Gr(\mathbb{X})^{C^1}$ of Definition~\ref{def:integrands},
that is, on the maps from $C^1$ to $\Gr^{\mathrm{ev}}(\mathbb{X})$, each reflection acting through \eqref{eq:setting-reflection-coefficients}.

\begin{lemma}[Images of admissible products]\label{lem:admissible-images}
Let $g\in W$, and let $\vec g\in\Ocub$ be its linear part.
\begin{enumerate}
\item\label{it:images-products}
Let $G_1$, $G_2$ be even elements of $\Gr(\mathbb{X})^{C^1}$, and let $z\in\C$.
Then
\[
(G_1\cdot G_2)^g=G_1^g\cdot G_2^g,\qquad 1^g=1,
\]
and $(z\cdot G_1)^g=z\cdot G_1^g$ if $g$ is even, $(z\cdot G_1)^g=\overline{z}\cdot G_1^g$ if $g$ is odd.
The even elements of $\Gr(\mathbb{X})^{C^1}$ lie in the centre of this algebra;
hence a product of elements of $\mathcal{A}(\R^3)$ does not depend on the order of its factors.
\item\label{it:images-boxes}
Let $B=y+[-r,r]^3$ be a box.
Then $gB=gy+[-r,r]^3$ is a box, and
\[
(gB)^{(1)}=g\,B^{(1)},\qquad(gB)^{(2)}=g\,B^{(2)},\qquad\mathbb{X}(gB)=g\,\mathbb{X}(B).
\]
If $G\in\mathcal{G}^{\mathrm{ev}}(B)$, then $G^g\in\mathcal{G}^{\mathrm{ev}}(gB)$, with the same growth margin $\varepsilon_G$.
If $\Phi$ is a bounded Borel function of finitely many $F_\sigma$, $\sigma$ contained in an open set $U\subset\R^3$,
then $\Phi^g$ is a bounded Borel function of finitely many $F_\sigma$, $\sigma$ contained in $gU$, and $\sup|\Phi^g|=\sup|\Phi|$.
Consequently $G\in\mathcal{A}(U)$ implies $G^g\in\mathcal{A}(gU)$;
and if a product of $\mathcal{A}(\R^3)$ embeds in a subgroup $\Gamma'$ of $\Z^3$ of rank two or three, then its image under $g$ embeds in $\vec g\,\Gamma'$.
\item\label{it:images-periods}
Let $\mathbb{T}=(\Gamma,\chi)$ be a torus, let $G\in\mathcal{A}(\R^3)$ embed in $\mathbb{T}$, and let $v\in\Gamma$.
Then $G^{\tau_v}$ embeds in $\mathbb{T}$,
and $G^{\tau_v}$ and $G$ define the same $\Gr_{\mathbb{T}}$-valued function on $\Conf(\Gamma)$.
\end{enumerate}
\end{lemma}

\begin{proof}
\eqref{it:images-products}
Moving a generator of $\Gr(\mathbb{X})$ across a monomial of even degree changes the sign an even number of times;
so such a monomial commutes with every generator, hence with every element of $\Gr(\mathbb{X})$,
and by linearity every even element of $\Gr(\mathbb{X})$ does the same.
As the product of $\Gr(\mathbb{X})^{C^1}$ is taken pointwise, the even elements of $\Gr(\mathbb{X})^{C^1}$ lie in its centre;
among them are the elements of $\mathcal{A}(\R^3)$, whose values are even (Definition~\ref{def:integrands}).
This proves the last two assertions of \eqref{it:images-products}.

By \eqref{eq:setting-group-action}, the map $G\mapsto G^g$ is the composition of the maps $\Theta_\Pi$ along a word in reflections that represents $g$.
It is therefore enough to consider $g=\theta_\Pi$ for a mirror $\Pi$.
By \eqref{eq:setting-reflection-coefficients} and the antilinearity of $\Theta_\Pi$,
\[
(\Theta_\Pi G)(\hat A)=\Theta_\Pi\left(G(\theta_\Pi^{*}\hat A)\right),\qquad\hat A\in C^1,
\]
for every element $G$ of $\Gr(\mathbb{X})^{C^1}$;
since $\Theta_\Pi$ reverses the order of products (Definition~\ref{def:grassmann-reflections}), $\Theta_\Pi(G_1\cdot G_2)=\Theta_\Pi(G_2)\cdot\Theta_\Pi(G_1)$.
The images of the generators under $\Theta_\Pi$ are linear combinations of generators, by \eqref{eq:setting-reflection-fields};
consequently $\Theta_\Pi$ carries every monomial into a linear combination of monomials of the same degree,
and the display shows that $\Theta_\Pi(G_1)$ and $\Theta_\Pi(G_2)$ are again even.
By the first paragraph they commute, and $\Theta_\Pi(G_1\cdot G_2)=\Theta_\Pi(G_1)\cdot\Theta_\Pi(G_2)$.
Moreover $\Theta_\Pi$ fixes the unit of $\Gr(\mathbb{X})^{C^1}$, the constant map $1$:
by the display, its image is the constant map with the value $\Theta_\Pi(1)\in\Gr(\mathbb{X})$, and $\Theta_\Pi(1)=1$ in $\Gr(\mathbb{X})$.
Indeed, $\Theta_\Pi(1)\cdot\Theta_\Pi(G)=\Theta_\Pi(G\cdot1)=\Theta_\Pi(G)$ for every $G\in\Gr(\mathbb{X})$,
and $\Theta_\Pi$ is onto, its square being the identity (Definition~\ref{def:grassmann-reflections});
so $\Theta_\Pi(1)\cdot G'=G'$ for every $G'\in\Gr(\mathbb{X})$, and $G'=1$ gives the claim.
Each $\Theta_\Pi$ is antilinear,
so a composition of an even number of them is linear and a composition of an odd number is antilinear;
and $g$ is even exactly when the words in reflections that represent it have even length (Definition~\ref{def:mirrors}),
the linear part of every reflection having the determinant $-1$.

\eqref{it:images-boxes}
One has $W=\Z^3\rtimes\Ocub$ (Lemma~\ref{lem:chambers}\eqref{it:chambers-group}),
so that $g$ is the composition of its linear part $\vec g$ with a translation by a vector of $\Z^3$.
The elements of $\Ocub=\mathrm{O}(3)\cap\mathrm{GL}(3,\Z)$ are the signed permutation matrices,
their columns being integer vectors of length one;
they map $[-r,r]^3$ onto itself, whence $g(y+[-r,r]^3)=gy+[-r,r]^3$.
For the same reason $g$ maps $\Z^3$ onto itself and permutes the sites, links, plaquettes and cubes of $\Z^3$, regarded as subsets of $\R^3$;
as $g$ is a bijection of $\R^3$, a cell is contained in $B$ if and only if its image is contained in $gB$.
This proves $(gB)^{(1)}=g\,B^{(1)}$ and $(gB)^{(2)}=g\,B^{(2)}$.
The set $\mathbb{X}=Wv_0$ is invariant under $g$, and $\pi(gx)=g\,\pi(x)$ for $x\in\mathbb{X}$ (Proposition~\ref{prop:geometry}\eqref{it:geometry-projection});
moreover $g\left(\pi(x)+[-\frac12,\frac12]^3\right)=\pi(gx)+[-\frac12,\frac12]^3$.
Hence $\pi(x)+[-\frac12,\frac12]^3\subset B$ if and only if $\pi(gx)+[-\frac12,\frac12]^3\subset gB$,
that is, $\mathbb{X}(gB)=g\,\mathbb{X}(B)$.

For the assertion on $\mathcal{G}^{\mathrm{ev}}(B)$ it suffices, by \eqref{eq:setting-group-action} and what has just been shown, to take for $g$ a reflection $\theta:=\theta_\Pi$;
let $n$ be the upward unit normal of the mirror $\Pi$.
Let $G\in\mathcal{G}^{\mathrm{ev}}(B)$, $G(\hat A)=\sum_\Psi a_\Psi(\hat A)\cdot\Psi$,
where $\Psi$ ranges over the monomials in the generators at the sites of $\mathbb{X}(B)$.
By \eqref{eq:setting-reflection-coefficients},
\[
(\Theta_\Pi G)(\hat A)=\sum_\Psi\overline{a_\Psi(\theta^{*}\hat A)}\cdot\Theta_\Pi(\Psi),
\]
where $\hat A$ is a real cochain on $(\theta B)^{(1)}$ and $\theta^{*}\hat A$ is its pull-back (Definition~\ref{def:cochains}), a real cochain on $B^{(1)}$.
The map $\hat A\mapsto(\Theta_\Pi G)(\hat A)$ is Borel, since $\hat A\mapsto\theta^{*}\hat A$ is linear and the $a_\Psi$ are Borel;
its values are even, and they involve only generators at the sites of $\theta\,\mathbb{X}(B)=\mathbb{X}(\theta B)$, by \eqref{eq:setting-reflection-fields}.

\emph{Growth.}
By \eqref{eq:setting-reflection-fields}, the image under $\Theta_\Pi$ of a generator is a linear combination of generators with entries of $-\gamma_n$ as coefficients.
As $\Theta_\Pi$ reverses products, $\Theta_\Pi(\Psi)$ is a product of such combinations;
multiplying out, one writes it as a sum of monomials, each multiplied by a signed product of entries of $\gamma_n$, a number that does not depend on $\hat A$.
There are finitely many monomials in the generators at the sites of $\mathbb{X}(B)$;
hence there is a constant $C'$ such that every coefficient of $(\Theta_\Pi G)(\hat A)$ is bounded in modulus by $C'\cdot\max_\Psi\left|a_\Psi(\theta^{*}\hat A)\right|$.
As $G\in\mathcal{G}^{\mathrm{ev}}(B)$, there are constants $C<\infty$ and $\varepsilon_G>0$ such that for every $\Psi$
\[
\left|a_\Psi(\theta^{*}\hat A)\right|\le C\cdot\exp\left(\left(\frac12-\varepsilon_G\right)\sum_{\sigma\in B^{(2)}}\left((d\theta^{*}\hat A)_\sigma\right)^2\right).
\]
Here $d\theta^{*}\hat A=\theta^{*}d\hat A$, and $(\theta^{*}u)_\sigma=\pm u_{\theta\sigma}$ for every two-cochain $u$ (Definition~\ref{def:cochains});
as $\theta$ maps $B^{(2)}$ onto $(\theta B)^{(2)}$, the sum in the exponent equals $\sum_{\sigma\in(\theta B)^{(2)}}\left((d\hat A)_\sigma\right)^2$.
So the coefficients of $\Theta_\Pi G$ obey the bound of Definition~\ref{def:integrands} on the box $\theta B$,
with the constant $C'\cdot C$ and the same $\varepsilon_G$.

\emph{Gauge invariance.}
By Definition~\ref{def:chamber-operator}, the transformation \eqref{eq:setting-gauge} with a real function $\lambda$ on photon sites multiplies a monomial $\Psi$ by $\exp\left(iq\langle\lambda,\mathfrak{q}(\Psi)\rangle\right)$,
where $\langle\cdot,\cdot\rangle$ is the counting pairing of functions on $\Z^3$ (Definition~\ref{def:cochains}).
The gauge invariance of $G$ therefore says that
\[
a_\Psi(\hat A'-d\lambda')\cdot\exp\left(iq\langle\lambda',\mathfrak{q}(\Psi)\rangle\right)=a_\Psi(\hat A')
\]
for every monomial $\Psi$, every real cochain $\hat A'$ on $B^{(1)}$ and every real function $\lambda'$ on the photon sites of $B$.
By \eqref{eq:setting-reflection-fields}, a generator $\psi_\alpha(x)$ of $\Psi$ is replaced in $\Theta_\Pi(\Psi)$ by generators $\bar\psi_\beta(\theta x)$,
and a generator $\bar\psi_\alpha(x)$ by generators $\psi_\beta(\theta x)$;
since $\pi(\theta x)=\theta\,\pi(x)$, every monomial $\Psi'$ that occurs in $\Theta_\Pi(\Psi)$ has the charge density $\mathfrak{q}(\Psi')=-\mathfrak{q}(\Psi)\circ\theta$.
Now let $\lambda$ be a real function on the photon sites of $\theta B$, and put $\lambda':=\lambda\circ\theta$, a real function on the photon sites of $B$.
Then $\langle\lambda,\mathfrak{q}(\Psi')\rangle=-\langle\lambda',\mathfrak{q}(\Psi)\rangle$ for these $\Psi'$,
and $\theta^{*}(\hat A-d\lambda)=\theta^{*}\hat A-d\lambda'$,
because $\theta^{*}$ commutes with $d$ and acts on functions of sites by composition with $\theta$ (Definition~\ref{def:cochains}).
Hence the transformation \eqref{eq:setting-gauge} with $\lambda$ carries $(\Theta_\Pi G)(\hat A-d\lambda)$ into
\[
\sum_\Psi\overline{a_\Psi(\theta^{*}\hat A-d\lambda')\cdot\exp\left(iq\langle\lambda',\mathfrak{q}(\Psi)\rangle\right)}\cdot\Theta_\Pi(\Psi)
=\sum_\Psi\overline{a_\Psi(\theta^{*}\hat A)}\cdot\Theta_\Pi(\Psi)
=(\Theta_\Pi G)(\hat A),
\]
and $\Theta_\Pi G$ is gauge invariant.
Together with the growth bound this gives $\Theta_\Pi G\in\mathcal{G}^{\mathrm{ev}}(\theta B)$ with the growth margin $\varepsilon_G$.

\emph{Bounded factors.}
Let $\Phi$ depend only on the values $F_\sigma$ at the plaquettes $\sigma$ of a finite set of plaquettes contained in $U$.
By \eqref{eq:setting-group-action}, $\Phi^g(F)=\Phi(g^{*}F)$ if $g$ is even and $\Phi^g(F)=\overline{\Phi(g^{*}F)}$ if $g$ is odd,
and $(g^{*}F)_\sigma=\pm F_{g\sigma}$ (Definition~\ref{def:cochains}).
Hence $\Phi^g$ is a bounded Borel function of the values $F_{g\sigma}$ at the images $g\sigma$ of these plaquettes, which are contained in $gU$.
As $g^{*}$ is a bijection of the two-cochains onto themselves which maps closed cochains onto closed cochains, $\sup|\Phi^g|=\sup|\Phi|$.

\emph{Products.}
Let $\Phi\cdot G_1\cdots G_k$ be a product of $\mathcal{A}(U)$:
$G_j\in\mathcal{G}^{\mathrm{ev}}(B_j)$ for pairwise disjoint boxes $B_j\subset U$, and $\Phi$ as above.
By \eqref{it:images-products},
\[
\left(\Phi\cdot G_1\cdots G_k\right)^g=\Phi^g\cdot G_1^g\cdots G_k^g,
\]
where $G_j^g\in\mathcal{G}^{\mathrm{ev}}(gB_j)$, the boxes $gB_j$ are pairwise disjoint and contained in $gU$, $g$ being a bijection of $\R^3$,
and $\Phi^g$ is a bounded Borel function of finitely many $F_\sigma$, $\sigma\subset gU$.
So the image is a product of $\mathcal{A}(gU)$,
and, $G\mapsto G^g$ being linear or antilinear, $\mathcal{A}(U)$ is mapped into $\mathcal{A}(gU)$.
Let finally the product embed in a subgroup $\Gamma'$ of $\Z^3$ of rank two or three.
For $v\in\Gamma'$ one has
\[
g\left(B_j+[-1,1]^3+v\right)=gB_j+[-1,1]^3+\vec g\,v,
\qquad
g\left(B_j+v\right)=gB_j+\vec g\,v,
\]
because $\vec g$ maps $[-1,1]^3$ onto itself.
As $g$ is a bijection of $\R^3$ and $v\mapsto\vec g\,v$ is a bijection of $\Gamma'$ onto $\vec g\,\Gamma'$,
the disjointness conditions of Definition~\ref{def:integrands} for the boxes $B_j$ and the lattice $\Gamma'$
are equivalent to those for the boxes $gB_j$ and the lattice $\vec g\,\Gamma'$.
Hence the image embeds in $\vec g\,\Gamma'$.

\eqref{it:images-periods}
The linear part of $\tau_v$ is the identity;
so $G^{\tau_v}$ embeds in $\Gamma$, that is, in $\mathbb{T}$, by \eqref{it:images-boxes}, applied to the products that embed in $\Gamma$ and whose sum is $G$.

The translations act on integrands as follows.
For $\mu\in\{1,2,3\}$, a sign $\pm$ and $c\in\frac12\Z$, the mirrors $\Pi=(x_\mu,c\pm\frac12)$ and $\Pi'=(x_\mu,c)$ are parallel,
and $\theta_\Pi\,\theta_{\Pi'}=\tau_{\pm e_\mu}$ (Definition~\ref{def:mirrors}\eqref{it:mirror-product}).
Let $G'$ be an even element of $\Gr(\mathbb{X})^{C^1}$.
By \eqref{eq:setting-group-action}, $(G')^{\tau_{\pm e_\mu}}=\Theta_\Pi\Theta_{\Pi'}G'$.
Two applications of \eqref{eq:setting-reflection-coefficients}, together with $(\theta_\Pi\,\theta_{\Pi'})^{*}=\theta_{\Pi'}^{*}\theta_\Pi^{*}$ (Definition~\ref{def:cochains}), give
\[
(\Theta_\Pi\Theta_{\Pi'}G')(\hat A)=\Theta_\Pi\Theta_{\Pi'}\left(G'(\tau_{\pm e_\mu}^{*}\hat A)\right),
\]
and by \eqref{eq:setting-two-reflections}, the two mirrors having the same normal $n$ and $\gamma_n^2=1$,
the linear automorphism $\Theta_\Pi\Theta_{\Pi'}$ substitutes $\psi(x\pm e_\mu)$, $\bar\psi(x\pm e_\mu)$ for $\psi(x)$, $\bar\psi(x)$.
For a monomial $\Psi$ and $v\in\Z^3$ let $\tau_v\Psi$ be the monomial obtained from $\Psi$ by substituting $\psi(x+v)$, $\bar\psi(x+v)$ for $\psi(x)$, $\bar\psi(x)$.
Since translations commute, the composition rule of \eqref{eq:setting-group-action} gives, by induction, for every $v\in\Z^3$, every box $B$ and every $G'\in\mathcal{G}^{\mathrm{ev}}(B)$ with coefficients $a_\Psi$,
\[
(G')^{\tau_v}(\hat A)=\sum_\Psi a_\Psi(\tau_v^{*}\hat A)\cdot\tau_v\Psi,
\]
where $\hat A$ is a real cochain on $(B+v)^{(1)}$ and $\tau_v^{*}\hat A$ is the cochain $b\mapsto\hat A(b+v)$ on $B^{(1)}$.

Let now $v\in\Gamma$, and let $(A,z)$ represent a point of $\Conf(\Gamma)$.
In $\Gr_{\mathbb{T}}$ the substitution $\Psi\mapsto\tau_v\Psi$ does not change the image of an even element (Definition~\ref{def:tori}):
it multiplies the image of a monomial $\Psi$ of even degree by $\chi(v)^{\deg\Psi}=1$.
Let $B$ be a box that embeds in $\Gamma$, and let $G'\in\mathcal{G}^{\mathrm{ev}}(B)$;
the box $B+v$ embeds in $\Gamma$ as well.
Let $\hat A$ be a local lift of the point represented by $(A,z)$ on $B+v$ (Definition~\ref{def:integrands}).
Then $\tau_v^{*}\hat A$ is a local lift of the same point on $B$:
on $B^{(2)}$ one has $d\tau_v^{*}\hat A=\tau_v^{*}d\hat A=\tau_v^{*}F=F$,
and $\exp\left(iq\hat A(b+v)\right)=\exp\left(iqA(b+v)\right)=\exp\left(iqA(b)\right)$ for every link $b\in B^{(1)}$,
because $A$, $z$ and $F$ are $\Gamma$-periodic.
By the formula for $(G')^{\tau_v}$ and the invariance of $\Gr_{\mathbb{T}}$ under $\Psi\mapsto\tau_v\Psi$,
the value of $(G')^{\tau_v}$ at $(A,z)$, computed with the local lift $\hat A$, is the image in $\Gr_{\mathbb{T}}$ of $G'(\tau_v^{*}\hat A)$.
The value of an integrand at a point of $\Conf(\Gamma)$ does not depend on the local lift (Definition~\ref{def:integrands});
so this is the value of $G'$ at $(A,z)$.

The value at $(A,z)$ of a product $\Phi\cdot G_1\cdots G_k$ that embeds in $\Gamma$ is the number $\Phi(\widetilde F)$, formed with the periodic extension $\widetilde F$ of the curvature,
times the product of the values of $G_1,\dots,G_k$ at $(A,z)$ (Definition~\ref{def:integrands}).
As $\tau_v$ is even, \eqref{eq:setting-group-action} gives $\Phi^{\tau_v}(\widetilde F)=\Phi(\tau_v^{*}\widetilde F)$,
and $\tau_v^{*}\widetilde F=\widetilde F$ because $\widetilde F$ is $\Gamma$-periodic.
Finally, let $\Phi\cdot G_1\cdots G_k$ be one of the products that embed in $\Gamma$ and whose sum is $G$.
By \eqref{it:images-products}, its image under $\tau_v$ is $\Phi^{\tau_v}\cdot G_1^{\tau_v}\cdots G_k^{\tau_v}$, which embeds in $\Gamma$ by \eqref{it:images-boxes};
by the two facts just proved, this image and $\Phi\cdot G_1\cdots G_k$ take the same value at $(A,z)$.
The action of $\tau_v$ is linear by \eqref{it:images-products}, and so is the value at a point of $\Conf(\Gamma)$ (Definition~\ref{def:integrands});
hence $G^{\tau_v}$ and $G$ define the same $\Gr_{\mathbb{T}}$-valued function on $\Conf(\Gamma)$.
\end{proof}

\begin{lemma}[Embedding criteria]\label{lem:embedding-criteria}
\begin{enumerate}
\item\label{it:embedding-sup}
Let $B=y+[-r,r]^3$ be a box,
and let $\Gamma'$ be a subgroup of $\Z^3$ such that every $\gamma\in\Gamma'\setminus\{0\}$ satisfies $|\gamma|_\infty\ge4$ and $|\gamma|_\infty>4r$.
Then $B$ embeds in $\Gamma'$.
\item\label{it:embedding-height}
Let $B=y+[-r,r]^3$ be a box, let $\ell\in\mathcal{L}$, and let $c<c'$ in $\Adm(\ell)$.
If $B\subset\{c<\ell<c'\}$, then $2|n_\ell|^2\cdot r<c'-c$.
Moreover, $|\ell(v)|\le|n_\ell|^2\cdot|v|_\infty$ for every $v\in\Z^3$.
\item\label{it:embedding-slabs}
Let $\mathbb{T}=(\Gamma,\chi)$ be a torus and $\ell\in\mathcal{L}$,
with the cross-section $\Xi(\mathbb{T},\ell)=(\Gamma_\parallel,\chi_\parallel)$ and the height period $p=p(\mathbb{T},\ell)$.
Let $c_0<c_1<\dots<c_s=c_0+p$ be positions in $\Adm(\ell)$ with $c_k-c_{k-1}\le\frac{p}{2}$ for $k=1,\dots,s$,
and for each $k$ let $G_k$ be a product of $\mathcal{A}(\{c_{k-1}<\ell<c_k\})$ that embeds in $\Gamma_\parallel$.
Then $G_1\cdots G_s$ embeds in $\Gamma$.
\item\label{it:embedding-cells}
Let $\mathcal{Q}=(\mathfrak{F};o_1,o_2,o_3;h)$ be a cell structure with the frame $\mathfrak{F}=(\ell^1,\ell^2,\ell^3)$,
and let $\Gamma$ be a period lattice with $\ell^a(\Gamma)\subset2h\Z$ for $a=1,2,3$;
thus the translations by the elements of $\Gamma$ permute the open cells (Definition~\ref{def:cell-structures}).
Let $\kappa^1,\dots,\kappa^s\in\Z^3$ be such that the open cells $Q_{\kappa^1},\dots,Q_{\kappa^s}$ lie in pairwise different $\Gamma$-orbits,
and let $G_i$ be a product of $\mathcal{A}(Q_{\kappa^i})$ for $i=1,\dots,s$.
Then $G_1\cdots G_s$ embeds in $\Gamma$, hence in every subgroup of $\Gamma$.
\end{enumerate}
\end{lemma}

\begin{proof}
\eqref{it:embedding-sup}
The set $B+[-1,1]^3=y+[-r-1,r+1]^3$ meets its translate by $v\in\Z^3$ if and only if $|v|_\infty\le2r+2$.
Let $\gamma\in\Gamma'\setminus\{0\}$.
If $r<1$, then $2r+2<4\le|\gamma|_\infty$;
if $r\ge1$, then $2r+2\le4r<|\gamma|_\infty$.
Two of the sets $B+[-1,1]^3+\gamma$, $\gamma\in\Gamma'$, differ by a translation by a non-zero element of $\Gamma'$;
so they are pairwise disjoint, and $B$ embeds in $\Gamma'$.

\eqref{it:embedding-height}
The entries of $n_\ell$ lie in $\{0,1,-1\}$, and the sum of their moduli is $|n_\ell|^2$,
namely $1$ for an axis family and $2$ for a diagonal family.
Hence
\[
\max\left\{\ell(x):\ |x|_\infty\le r\right\}=r\cdot|n_\ell|^2,
\]
and $\ell(B)=\left[\ell(y)-|n_\ell|^2\cdot r,\ \ell(y)+|n_\ell|^2\cdot r\right]$, an interval of length $2|n_\ell|^2\cdot r$.
If $B\subset\{c<\ell<c'\}$, this closed interval lies in the open interval $(c,c')$, and $2|n_\ell|^2\cdot r<c'-c$.
For $v\in\Z^3$ one has likewise $|\ell(v)|=|n_\ell\cdot v|\le|n_\ell|^2\cdot|v|_\infty$.

\eqref{it:embedding-slabs}
By Definition~\ref{def:slabs}, $\Gamma_\parallel=\Gamma\cap\{\ell=0\}$ and $\ell(\Gamma)=p\Z$.
Let $B$ be a box of one of the products, say of $G_k$, let $r$ be its half side, and let $\gamma\in\Gamma\setminus\{0\}$.
If $\ell(\gamma)=0$, then $\gamma\in\Gamma_\parallel$, and the sets $B+[-1,1]^3$ and $B+[-1,1]^3+\gamma$ are disjoint because $G_k$ embeds in $\Gamma_\parallel$.
If $\ell(\gamma)\neq0$, then \eqref{it:embedding-height}, applied to $B\subset\{c_{k-1}<\ell<c_k\}$ and to $\gamma$, gives
\[
|n_\ell|^2\cdot|\gamma|_\infty\ge|\ell(\gamma)|\ge p\ge2\,(c_k-c_{k-1})>4|n_\ell|^2\cdot r,
\]
so that $|\gamma|_\infty>4r$;
moreover $|\gamma|_\infty\ge\sys(\Gamma)\ge4$.
These two bounds give $|\gamma|_\infty>2r+2$, since $2r+2<4$ if $r<1$ and $2r+2\le4r$ if $r\ge1$;
and the cube $B+[-1,1]^3$, of side $2r+2$, does not meet its translate by a vector whose maximum norm exceeds $2r+2$.
Translating, one finds that every box of $G_1\cdots G_s$ embeds in $\Gamma$.

Let now $B$ and $B'$ be boxes of $G_k$ and $G_{k'}$, and let $\gamma,\gamma'\in\Gamma$ with $(B+\gamma)\cap(B'+\gamma')\neq\emptyset$.
The open slabs
\[
\{c_{j-1}+i\cdot p<\ell<c_j+i\cdot p\},\qquad 1\le j\le s,\ i\in\Z,
\]
are pairwise disjoint, since $c_0<c_1<\dots<c_s=c_0+p$.
The set $B+\gamma$ lies in the one with $j=k$ and $i\cdot p=\ell(\gamma)$, and $B'+\gamma'$ in the one with $j=k'$ and $i\cdot p=\ell(\gamma')$.
Hence $k=k'$ and $\ell(\gamma)=\ell(\gamma')$, so that $\gamma-\gamma'\in\Gamma_\parallel$.
As $G_k$ embeds in $\Gamma_\parallel$, the sets $B''+\gamma''$, for the boxes $B''$ of $G_k$ and $\gamma''\in\Gamma_\parallel$, are pairwise disjoint;
since $\left(B+(\gamma-\gamma')\right)\cap B'\neq\emptyset$, it follows that $B=B'$ and $\gamma=\gamma'$.
In particular, for $\gamma=\gamma'=0$, the boxes of $G_1,\dots,G_s$ are pairwise disjoint;
the product of their bounded factors being again a bounded Borel function of finitely many $F_\sigma$,
$G_1\cdots G_s$ is a product of $\mathcal{A}(\R^3)$, and its boxes embed disjointly in $\Gamma$.

\eqref{it:embedding-cells}
Let $B$ be a box of $G_i$ and $r$ its half side.
For each $a$, the box $B$ lies in the open slab $\{o_a+\kappa^i_a\cdot h<\ell^a<o_a+(\kappa^i_a+1)\cdot h\}$, whose bounding positions lie in $\Adm(\ell^a)$;
by \eqref{it:embedding-height}, $2|n_a|^2\cdot r<h$.
Let $\gamma\in\Gamma\setminus\{0\}$.
The forms $\ell^1$, $\ell^2$, $\ell^3$ are linearly independent, their normals being pairwise orthogonal;
so $\ell^a(\gamma)\neq0$ for some $a$, and then $|\ell^a(\gamma)|\ge2h>4|n_a|^2\cdot r$, because $\ell^a(\Gamma)\subset2h\Z$.
By \eqref{it:embedding-height}, $|\gamma|_\infty\ge|\ell^a(\gamma)|/|n_a|^2>4r$;
moreover $|\gamma|_\infty\ge\sys(\Gamma)\ge4$.
By \eqref{it:embedding-sup}, $B$ embeds in $\Gamma$.

For $\gamma\in\Gamma$ the translation $\tau_\gamma$ maps $Q_\kappa$ onto $Q_{\kappa'}$ with $\kappa'_a=\kappa_a+\ell^a(\gamma)/h$ (Definition~\ref{def:cell-structures}),
and distinct open cells are disjoint.
Let $B$ and $B'$ be boxes of $G_i$ and $G_j$, so that $B\subset Q_{\kappa^i}$ and $B'\subset Q_{\kappa^j}$,
and let $\gamma,\gamma'\in\Gamma$ with $(B+\gamma)\cap(B'+\gamma')\neq\emptyset$.
Then the cells $\tau_\gamma Q_{\kappa^i}$ and $\tau_{\gamma'}Q_{\kappa^j}$ meet, hence coincide.
So $Q_{\kappa^i}$ and $Q_{\kappa^j}$ lie in one $\Gamma$-orbit, and $i=j$.
Then $\tau_{\gamma-\gamma'}$ maps the cell $Q_{\kappa^i}$ onto itself, so that $\ell^a(\gamma-\gamma')=0$ for $a=1,2,3$,
and $\gamma=\gamma'$ by the linear independence of the three forms.
Finally, $B$ and $B'$ are boxes of the one product $G_i$ that meet;
the boxes of a product being pairwise disjoint (Definition~\ref{def:integrands}), $B=B'$.
Thus the boxes of $G_1,\dots,G_s$ are pairwise disjoint and embed disjointly in $\Gamma$, and $G_1\cdots G_s$ is a product of $\mathcal{A}(\R^3)$ that embeds in $\Gamma$.
For a subgroup $\Gamma'$ of $\Gamma$, the disjointness conditions of Definition~\ref{def:integrands} for $\Gamma'$ are among those for $\Gamma$;
so $G_1\cdots G_s$ embeds in $\Gamma'$ as well.
\end{proof}

\subsection{Straightened tori}\label{sec:straightened-tori}

\begin{lemma}[Cross-sections and closing vectors]\label{lem:closing-vectors}
Let $\mathbb{T}=(\Gamma,\chi)$ be a torus and $\ell\in\mathcal{L}$ a family, and put $\Gamma_\parallel:=\Gamma\cap\{\ell=0\}$.
\begin{enumerate}
\item\label{it:closing-section}
The group $\Gamma_\parallel$ has rank two, and $\sys(\Gamma_\parallel)\ge4$.
Consequently $\Xi(\mathbb{T},\ell)=\left(\Gamma_\parallel,\chi|_{\Gamma_\parallel}\right)$ is a cross-section for $\ell$ in the sense of Definition~\ref{def:slabs}.
\item\label{it:closing-period}
One has $\ell(\Gamma)=p\Z$ with an integer $p\ge1$.
The closing vectors of $\mathbb{T}$ for $\ell$ are exactly the elements $g\in\Gamma$ with $\ell(g)=p$;
such elements exist, $\Gamma=\Gamma_\parallel\oplus\Z g$ for each of them, and $p(\mathbb{T},\ell)=p$.
\item\label{it:closing-index}
The number of photon sites of $\mathbb{T}$ is
\[
|\mathbb{T}|=p\cdot\left[\Z^3\cap\{\ell=0\}:\Gamma_\parallel\right].
\]
\item\label{it:closing-straightened}
Suppose that the straightened torus $\mathbb{T}^0=\mathbb{T}^0(\ell)=(\Gamma^0,\chi^0)$ of Definition~\ref{def:slabs} is defined, and put $g^0:=\frac{p}{|n_\ell|^2}\,n_\ell$.
Then its lattice is $\Gamma^0=\Gamma_\parallel\oplus\Z g^0$, and $\Gamma^0\cap\{\ell=0\}=\Gamma_\parallel$, so that $\Xi(\mathbb{T}^0,\ell)=\Xi(\mathbb{T},\ell)$.
Moreover $p(\mathbb{T}^0,\ell)=p$, the vector $g^0$ is a closing vector of $\mathbb{T}^0$ for $\ell$ with $\chi^0(g^0)=-1$,
the torus $\mathbb{T}^0$ is adapted to $\ell$, and $|\mathbb{T}^0|=|\mathbb{T}|$.
\end{enumerate}
\end{lemma}

\begin{proof}
The vector $n_\ell$ has an entry equal to $1$ or to $-1$ (Definition~\ref{def:mirrors}), so that $\ell$ maps $\Z^3$ onto $\Z$.
On $\Gamma$ the form $\ell$ does not vanish identically:
$\Gamma$ has rank three, whereas a subgroup of $\Z^3$ contained in the plane $\{\ell=0\}$ has rank at most two.
Hence $\ell(\Gamma)$ is a non-zero subgroup of $\Z$, that is, $\ell(\Gamma)=p\Z$ with an integer $p\ge1$.
The kernel of the homomorphism $\Gamma\to p\Z$ induced by $\ell$ is $\Gamma_\parallel$, and the quotient $\Gamma/\Gamma_\parallel\cong p\Z$ has rank one;
so $\Gamma_\parallel$ has rank two.
Its non-zero elements are non-zero elements of $\Gamma$, whence $\sys(\Gamma_\parallel)\ge\sys(\Gamma)\ge4$ (Definition~\ref{def:configurations}).
Since $\Gamma_\parallel$ is a subgroup of $\Z^3\cap\{\ell=0\}$ and $\chi|_{\Gamma_\parallel}$ is a character with the values $\pm1$, item~\eqref{it:closing-section} is proved.

Let $g\in\Gamma$ satisfy $\ell(g)=p$; such an element exists, as $\ell(\Gamma)=p\Z$.
For $\gamma\in\Gamma$ the number $\ell(\gamma)/p$ is an integer, and in the decomposition
\[
\gamma=\left(\gamma-\frac{\ell(\gamma)}{p}\,g\right)+\frac{\ell(\gamma)}{p}\,g
\]
the first term lies in $\Gamma$ and in $\{\ell=0\}$, hence in $\Gamma_\parallel$.
Therefore $\Gamma=\Gamma_\parallel+\Z g$.
The sum is direct, because $\ell$ vanishes on $\Gamma_\parallel$ while $\ell(i\cdot g)=i\cdot p\ne0$ for every integer $i\ne0$;
and $\ell(g)=p>0$, so that $g$ is a closing vector.
Conversely, a closing vector $g$ satisfies $p\Z=\ell(\Gamma)=\ell(\Gamma_\parallel+\Z g)=\ell(g)\Z$ and $\ell(g)>0$, whence $\ell(g)=p$.
This proves~\eqref{it:closing-period}; in particular $p(\mathbb{T},\ell)=p$.

For~\eqref{it:closing-index} consider the map $\Z^3/\Gamma\to\Z/p\Z$, $x+\Gamma\mapsto\ell(x)+p\Z$.
It is a well-defined homomorphism, since $\ell(\Gamma)=p\Z$, and it is onto, since $\ell(\Z^3)=\Z$.
A class $x+\Gamma$ lies in its kernel if and only if $\ell(x)=\ell(\gamma)$ for some $\gamma\in\Gamma$,
that is, if and only if $x\in\left(\Z^3\cap\{\ell=0\}\right)+\Gamma$.
The kernel is therefore $\left(\left(\Z^3\cap\{\ell=0\}\right)+\Gamma\right)/\Gamma$,
which is isomorphic to $\left(\Z^3\cap\{\ell=0\}\right)/\Gamma_\parallel$, because the intersection of $\Z^3\cap\{\ell=0\}$ with $\Gamma$ is $\Gamma_\parallel$.
As $\Gamma$ has rank three, the index $|\mathbb{T}|=[\Z^3:\Gamma]$ is finite (Definition~\ref{def:tori}),
and it equals $p$ times the order of this kernel, that is, $p\cdot\left[\Z^3\cap\{\ell=0\}:\Gamma_\parallel\right]$.

Now let $\mathbb{T}^0$ be defined.
By Definition~\ref{def:slabs} and by~\eqref{it:closing-period}, the vector $g^0$ has integer entries, the group $\Gamma^0=\Gamma_\parallel\oplus\Z g^0$ has systole at least four,
$\chi^0$ agrees with $\chi$ on $\Gamma_\parallel$, and $\chi^0(g^0)=-1$.
Since $\ell(n_\ell)=|n_\ell|^2$,
\[
\ell(g^0)=\frac{p}{|n_\ell|^2}\cdot|n_\ell|^2=p .
\]
An element $\gamma_\parallel+i\cdot g^0$ of $\Gamma^0$, with $\gamma_\parallel\in\Gamma_\parallel$ and $i\in\Z$, has the height $\ell\left(\gamma_\parallel+i\cdot g^0\right)=i\cdot p$,
which vanishes only for $i=0$.
Hence $\Gamma^0\cap\{\ell=0\}=\Gamma_\parallel$, and
\[
\Xi(\mathbb{T}^0,\ell)=\left(\Gamma_\parallel,\chi^0|_{\Gamma_\parallel}\right)=\left(\Gamma_\parallel,\chi|_{\Gamma_\parallel}\right)=\Xi(\mathbb{T},\ell).
\]
Moreover $\ell(\Gamma^0)=p\Z$,
and~\eqref{it:closing-period}, applied to $\mathbb{T}^0$, gives $p(\mathbb{T}^0,\ell)=p$ and shows that $g^0$ is a closing vector of $\mathbb{T}^0$.
The straightened torus of $\mathbb{T}^0$ for $\ell$ is built from the group $\Gamma^0\cap\{\ell=0\}=\Gamma_\parallel$, from the height period $p$
and from the restriction of $\chi^0$ to $\Gamma_\parallel$, which is $\chi|_{\Gamma_\parallel}$.
Its lattice is $\Gamma_\parallel\oplus\Z g^0=\Gamma^0$, whose systole is at least four,
and its character agrees with $\chi^0$ on $\Gamma_\parallel$ and at $g^0$.
So the straightened torus of $\mathbb{T}^0$ is defined and equals $\mathbb{T}^0$; that is, $\mathbb{T}^0$ is adapted to $\ell$.
Finally,~\eqref{it:closing-index}, applied to $\mathbb{T}^0$ and to $\mathbb{T}$, gives
$|\mathbb{T}^0|=p\cdot\left[\Z^3\cap\{\ell=0\}:\Gamma_\parallel\right]=|\mathbb{T}|$.
\end{proof}

\begin{definition}[Copies in slabs]\label{def:slab-copies}
Let $\ell\in\mathcal{L}$, $c_0\in\Adm(\ell)$ and $h'>0$ be such that $c_j:=c_0+j\cdot h'$ belongs to $\Adm(\ell)$ for every $j\in\Z$, and put
\[
v_{h'}:=\frac{2h'}{|n_\ell|^2}\,n_\ell .
\]
By Definition~\ref{def:mirrors}\eqref{it:mirror-product} for the positions $c_0$ and $c_1=c_0+h'$, the vector $v_{h'}$ has integer entries, and $\ell(v_{h'})=2h'$.
The same identity for the positions $c_j$ and $c_{j+s}$, together with $\theta_{\ell,c_j}^2=1$ (Definition~\ref{def:mirrors}\eqref{it:mirror-involution}), gives
\begin{equation}\label{eq:tori-reflection-shift}
\theta_{\ell,c_{j+s}}\,\theta_{\ell,c_j}=\tau_{s\cdot v_{h'}},\qquad
\theta_{\ell,c_{j+s}}=\tau_{s\cdot v_{h'}}\,\theta_{\ell,c_j}\qquad(j,s\in\Z).
\end{equation}
For $j\in\Z$ the open slab $\{c_{j-1}<\ell<c_j\}$ is called the \emph{slab $j$}.

The reflections $\theta_{\ell,c_j}$, $j\in\Z$, generate a subgroup $W(\ell;c_0,h')$ of $W$, and
\[
W(\ell;c_0,h')=\left\{\tau_{s\cdot v_{h'}}:\ s\in\Z\right\}\cup\left\{\theta_{\ell,c_j}:\ j\in\Z\right\}.
\]
The right side contains the generators,
and it is contained in $W(\ell;c_0,h')$, because $\tau_{s\cdot v_{h'}}=\theta_{\ell,c_s}\,\theta_{\ell,c_0}$ by~\eqref{eq:tori-reflection-shift};
it therefore suffices to show that it is a group.
It is closed under inverses, and it is closed under products by~\eqref{eq:tori-reflection-shift}, by $\tau_{s\cdot v_{h'}}\,\tau_{s'\cdot v_{h'}}=\tau_{(s+s')\cdot v_{h'}}$, and by the identity
\[
\theta_{\ell,c_j}\,\tau_{s\cdot v_{h'}}=\tau_{s\cdot v_{h'}}\,\theta_{\ell,c_j-2s\cdot h'}=\tau_{s\cdot v_{h'}}\,\theta_{\ell,c_{j-2s}}=\theta_{\ell,c_{j-s}},
\]
in which the first equality is Definition~\ref{def:mirrors}\eqref{it:mirror-conjugation} for the vector $-s\cdot v_{h'}$ and the last one is~\eqref{eq:tori-reflection-shift}.
By Definition~\ref{def:mirrors}\eqref{it:mirror-height}, $\ell\circ\theta_{\ell,c_{j'}}=2c_{j'}-\ell$, and $2c_{j'}-c_j=c_{2j'-j}$ for all $j,j'\in\Z$;
so the involution $\theta_{\ell,c_{j'}}$ maps the slab $j$ onto the slab $2j'-j+1$.
The translation $\tau_{s\cdot v_{h'}}$ maps the slab $j$ onto the slab $j+2s$, as $\ell\left(x+s\cdot v_{h'}\right)=\ell(x)+2s\cdot h'$ for $x\in\R^3$.
A reflection of $W(\ell;c_0,h')$ thus changes the parity of the number of a slab, and a translation preserves it.
Now let $i,j\in\Z$.
If $i-j$ is even, the translation $\tau_{s\cdot v_{h'}}$ with $2s=i-j$ is the only element of $W(\ell;c_0,h')$ that carries the slab $j$ onto the slab $i$;
if $i-j$ is odd, the reflection $\theta_{\ell,c_{j'}}$ with $2j'=i+j-1$ is the only one.
Writing $g_{i\leftarrow j}$ for this element, one has, by~\eqref{eq:tori-reflection-shift} in the odd case,
\[
g_{i\leftarrow j}=
\begin{cases}
\tau_{\frac{i-j}{2}\cdot v_{h'}} &\text{if }i-j\text{ is even},\\
\theta_{\ell,c_{(i+j-1)/2}}=\tau_{\frac{i-j-1}{2}\cdot v_{h'}}\,\theta_{\ell,c_j} &\text{if }i-j\text{ is odd}.
\end{cases}
\]
In particular $g_{j\leftarrow j}=1$ and $g_{j+1\leftarrow j}=\theta_{\ell,c_j}$;
and $g_{i+2s\leftarrow j}=\tau_{s\cdot v_{h'}}\,g_{i\leftarrow j}$ for $i,s\in\Z$, the right side being an element of $W(\ell;c_0,h')$ that maps the slab $j$ onto the slab $i+2s$.

Let $\Xi=(\Gamma_\parallel,\chi_\parallel)$ be a cross-section for $\ell$, let $j\in\Z$, and let $G\in\mathcal{A}_\Xi[c_{j-1},c_j]$;
the elements of the classes $\mathcal{A}(U)$ are even (Definition~\ref{def:integrands}), and so is $G$.
The \emph{copy} of $G$ in the slab $i\in\Z$ is
\[
G^{(i)}:=G^{g_{i\leftarrow j}},
\]
the action of $W$ on even integrands being that of~\eqref{eq:setting-group-action}.
The copy $G^{(i)}$ belongs to $\mathcal{A}_\Xi[c_{i-1},c_i]$ by Lemma~\ref{lem:admissible-images}\eqref{it:images-products} and~\eqref{it:images-boxes}:
the linear part of an element of $W(\ell;c_0,h')$ is either the identity or the linear reflection $x\mapsto x-\frac{2\ell(x)}{|n_\ell|^2}\,n_\ell$,
and in both cases it fixes the plane $\{\ell=0\}$, hence $\Gamma_\parallel$, pointwise.
By the relations between the elements $g_{i\leftarrow j}$ just stated and by~\eqref{eq:setting-group-action},
\[
G^{(j)}=G,\qquad G^{(j+1)}=\Theta_{\ell,c_j}G,\qquad G^{(i+2s)}=\left(G^{(i)}\right)^{\tau_{s\cdot v_{h'}}}\qquad(i,s\in\Z).
\]

Let finally $\mathbb{T}^0=(\Gamma^0,\chi^0)$ be a torus adapted to $\ell$ with $\Xi(\mathbb{T}^0,\ell)=\Xi$ and $p(\mathbb{T}^0,\ell)=2k\cdot h'$ for an integer $k\ge1$.
The \emph{fold} of $G$ in the direction $\ell$ is
\[
\fold_\ell G:=\prod_{i=1}^{2k}G^{(i)} .
\]
It embeds in $\mathbb{T}^0$.
Indeed, expanded into products, it is a sum of products $G'_1\cdots G'_{2k}$, where $G'_i$ is a product of $\mathcal{A}(\{c_{i-1}<\ell<c_i\})$ that embeds in $\Gamma_\parallel$;
and Lemma~\ref{lem:embedding-criteria}\eqref{it:embedding-slabs} applies to $\mathbb{T}^0$ and to the planes $c_0<c_1<\dots<c_{2k}=c_0+p(\mathbb{T}^0,\ell)$,
whose distances $h'$ do not exceed $p(\mathbb{T}^0,\ell)/2$.
For the same reason the product of the copies over any $2k$ consecutive slabs embeds in $\mathbb{T}^0$,
and so does every single copy $G^{(i)}$, the other slabs carrying the factor $1$.
As $\mathbb{T}^0$ is its own straightened torus (Definition~\ref{def:slabs}), its lattice $\Gamma^0$ contains the vector
\[
\frac{p(\mathbb{T}^0,\ell)}{|n_\ell|^2}\,n_\ell=k\cdot v_{h'} .
\]
Since $G^{(i+2k)}=\left(G^{(i)}\right)^{\tau_{k\cdot v_{h'}}}$, Lemma~\ref{lem:admissible-images}\eqref{it:images-periods} shows that $G^{(i+2k)}$ and $G^{(i)}$ define the same function on $\Conf(\Gamma^0)$, for every $i\in\Z$.
Now take $2k$ consecutive slabs.
Their numbers represent each residue class modulo $2k$ exactly once,
and the factors of the product of the copies of $G$ in these slabs commute with one another (Lemma~\ref{lem:admissible-images}\eqref{it:images-products}).
Replacing each factor $G^{(i)}$ by the copy of $G$ in the one slab among $1,\dots,2k$ whose number is congruent to $i$ modulo $2k$, and reordering the factors,
one finds that the product of the copies of $G$ in any $2k$ consecutive slabs defines the same function on $\Conf(\Gamma^0)$ as $\fold_\ell G$.
\end{definition}

\begin{proposition}[H\"older's inequality with a twist]\label{prop:twisted-holder}
Let $-1\le\rho\le1$, $w_0,\dots,w_3>0$, $M\in\R$ and $q>0$.
Let $\mathbb{T}=(\Gamma,\chi)$ be a torus, $\ell\in\mathcal{L}$ and $\Xi:=\Xi(\mathbb{T},\ell)$;
let $g$ be a closing vector of $\mathbb{T}$ for $\ell$, and $p:=\ell(g)$.
Let $k\ge1$ be an integer and $h'>0$ a number with $p=2k\cdot h'$,
let $c_0\in\Adm(\ell)$ be such that $c_j:=c_0+j\cdot h'\in\Adm(\ell)$ for all $j\in\Z$,
and suppose that the straightened torus $\mathbb{T}^0=\mathbb{T}^0(\ell)$ is defined (Definition~\ref{def:slabs}).
Let $\mathfrak{S}$ be the slab operators of the cross-section $\Xi$ (Corollary~\ref{cor:slab-admissible}),
and let $G_j\in\mathcal{A}_\Xi[c_{j-1},c_j]$ for $j=1,\dots,2k$.
Then the folds $\fold_\ell G_j$ on $\mathbb{T}^0$ (Definition~\ref{def:slab-copies}) satisfy
\begin{align}
Z_{\mathbb{T}^0}\left[\fold_\ell G_j\right]&=\left\|\mathfrak{S}[G_j]\right\|_{2k}^{2k}\ge0\qquad(j=1,\dots,2k),\label{eq:tori-fold-norm}\\
\left|Z_{\mathbb{T}}\left[G_1\cdots G_{2k}\right]\right|&\le\prod_{j=1}^{2k}Z_{\mathbb{T}^0}\left[\fold_\ell G_j\right]^{\frac{1}{2k}},\label{eq:tori-holder}
\end{align}
where $\|\cdot\|_{2k}$ is the norm of Definition~\ref{def:even-schatten}.
\end{proposition}

\begin{proof}
\emph{Step 1: preparations.}
By Lemma~\ref{lem:closing-vectors}, $\Xi$ is a cross-section for $\ell$ and $p(\mathbb{T},\ell)=p$.
With $\Gamma_\parallel:=\Gamma\cap\{\ell=0\}$ and $g^0:=\frac{p}{|n_\ell|^2}\,n_\ell$, the same lemma shows:
the lattice of $\mathbb{T}^0$ is $\Gamma^0=\Gamma_\parallel\oplus\Z g^0$;
one has $\Xi(\mathbb{T}^0,\ell)=\Xi$ and $p(\mathbb{T}^0,\ell)=p$;
the torus $\mathbb{T}^0$ is adapted to $\ell$;
and $g^0$ is a closing vector of $\mathbb{T}^0$ with $\chi^0(g^0)=-1$.
In the notation of Definition~\ref{def:slab-copies}, $g^0=k\cdot v_{h'}$.
So the folds $\fold_\ell G_j$ on $\mathbb{T}^0$ are defined, and they embed in $\mathbb{T}^0$ (Definition~\ref{def:slab-copies}).
Consecutive positions of the chain $c_0<c_1<\dots<c_{2k}=c_0+p$ differ by $h'=\frac{p}{2k}$, which does not exceed $\frac{p}{2}$;
so Lemma~\ref{lem:embedding-criteria}\eqref{it:embedding-slabs} applies to $\mathbb{T}$ and to this chain,
and every product in the expansion of $G_1\cdots G_{2k}$ embeds in $\Gamma$.
Thus both sides of~\eqref{eq:tori-fold-norm} and of~\eqref{eq:tori-holder} are defined (Definition~\ref{def:functional}).
By Corollary~\ref{cor:slab-admissible}, the assertions of Theorem~\ref{thm:slab-operators} hold for the elements of the classes $\mathcal{A}_\Xi[c,c']$.
The separable Hilbert spaces $\mathcal{H}(c)$, $c\in\Adm(\ell)$, the parity $U_{\mathrm{f}}$, the unitaries $V_v$ and the map $\mathfrak{S}$ of that theorem
are formed from $\ell$, $\Xi$ and the parameters alone (Definitions~\ref{def:mirror-spaces} and~\ref{def:slab-kernels});
no torus enters their construction.
As $\Xi(\mathbb{T},\ell)=\Xi(\mathbb{T}^0,\ell)=\Xi$, one system of spaces and operators therefore serves both tori,
and the trace formula of Theorem~\ref{thm:slab-operators}\eqref{it:slab-trace} holds, with these operators, on $\mathbb{T}$ as well as on $\mathbb{T}^0$.

\emph{Step 2: the identity~\eqref{eq:tori-fold-norm}.}
Fix $j\in\{1,\dots,2k\}$, write $G:=G_j$, and let $\mathfrak{Y}:=\mathfrak{S}[G]\colon\mathcal{H}(c_{j-1})\to\mathcal{H}(c_j)$, a Hilbert--Schmidt operator.
For the positions $c_{j-1}<c_j$ the vector of the reflection rule is $\frac{2(c_j-c_{j-1})}{|n_\ell|^2}\,n_\ell=v_{h'}$;
so Theorem~\ref{thm:slab-operators}\eqref{it:slab-reflection} gives
\[
\mathfrak{S}\left[\Theta_{\ell,c_j}G\right]=V_{v_{h'}}\cdot\mathfrak{Y}^*\colon\mathcal{H}(c_j)\to\mathcal{H}(c_{j+1}).
\]
The relations $G^{(j)}=G$, $G^{(j+1)}=\Theta_{\ell,c_j}G$ and $G^{(i+2s)}=\left(G^{(i)}\right)^{\tau_{s\cdot v_{h'}}}$ of Definition~\ref{def:slab-copies} give
\[
G^{(j+2s)}=G^{\tau_{s\cdot v_{h'}}},\qquad G^{(j+2s+1)}=\left(\Theta_{\ell,c_j}G\right)^{\tau_{s\cdot v_{h'}}}\qquad(s\in\Z).
\]
Theorem~\ref{thm:slab-operators}\eqref{it:slab-covariance}, together with the group law $V_v\cdot V_{v'}=V_{v+v'}$ of Theorem~\ref{thm:slab-operators}\eqref{it:slab-hilbert-schmidt}, therefore gives
\[
\begin{aligned}
\mathfrak{S}\left[G^{(j+2s)}\right]&=V_{s\cdot v_{h'}}\cdot\mathfrak{Y}\cdot V_{s\cdot v_{h'}}^{-1},\\
\mathfrak{S}\left[G^{(j+2s+1)}\right]&=V_{s\cdot v_{h'}}\cdot V_{v_{h'}}\cdot\mathfrak{Y}^*\cdot V_{s\cdot v_{h'}}^{-1}=V_{(s+1)\cdot v_{h'}}\cdot\mathfrak{Y}^*\cdot V_{s\cdot v_{h'}}^{-1},
\end{aligned}
\]
and hence
\[
\mathfrak{S}\left[G^{(j+2s+1)}\right]\cdot\mathfrak{S}\left[G^{(j+2s)}\right]=V_{(s+1)\cdot v_{h'}}\cdot\mathfrak{Y}^*\cdot\mathfrak{Y}\cdot V_{s\cdot v_{h'}}^{-1}.
\]
Now Theorem~\ref{thm:slab-operators}\eqref{it:slab-trace} applies to the torus $\mathbb{T}^0$,
with the planes $c_{j-1}<c_j<\dots<c_{j-1+2k}=c_{j-1}+p$, that is, with $2k\ge2$ slabs,
with the closing vector $g^0$, for which $U_{\mathrm{f}}^{\frac{1+\chi^0(g^0)}{2}}=1$,
and with the integrands $G^{(j)},\dots,G^{(j+2k-1)}$;
by Definition~\ref{def:slab-copies}, their product defines the same function on $\Conf(\Gamma^0)$ as $\fold_\ell G$.
Hence
\[
Z_{\mathbb{T}^0}\left[\fold_\ell G\right]=Z_{\mathbb{T}^0}\left[G^{(j)}\cdots G^{(j+2k-1)}\right]
=\Tr_{\mathcal{H}(c_{j-1})}\left(V_{g^0}^{-1}\cdot\mathfrak{S}\left[G^{(j+2k-1)}\right]\cdots\mathfrak{S}\left[G^{(j)}\right]\right).
\]
In the ordered product of the $2k$ slab operators the unitaries between consecutive brackets cancel, $V_{s\cdot v_{h'}}^{-1}\cdot V_{s\cdot v_{h'}}=1$,
and $V_0=1$, because $V_0$ is unitary and $V_0\cdot V_0=V_0$; thus
\[
\begin{aligned}
\mathfrak{S}\left[G^{(j+2k-1)}\right]\cdots\mathfrak{S}\left[G^{(j)}\right]
&=\left(V_{k\cdot v_{h'}}\cdot\mathfrak{Y}^*\cdot\mathfrak{Y}\cdot V_{(k-1)\cdot v_{h'}}^{-1}\right)\cdots\left(V_{v_{h'}}\cdot\mathfrak{Y}^*\cdot\mathfrak{Y}\cdot V_0^{-1}\right)\\
&=V_{k\cdot v_{h'}}\cdot\left(\mathfrak{Y}^*\cdot\mathfrak{Y}\right)^k .
\end{aligned}
\]
As $g^0=k\cdot v_{h'}$, it follows that
\[
Z_{\mathbb{T}^0}\left[\fold_\ell G\right]
=\Tr_{\mathcal{H}(c_{j-1})}\left(V_{k\cdot v_{h'}}^{-1}\cdot V_{k\cdot v_{h'}}\cdot\left(\mathfrak{Y}^*\cdot\mathfrak{Y}\right)^k\right)
=\Tr_{\mathcal{H}(c_{j-1})}\left(\left(\mathfrak{Y}^*\cdot\mathfrak{Y}\right)^k\right).
\]
The operator $\left(\mathfrak{Y}^*\cdot\mathfrak{Y}\right)^k$ is non-negative and of trace class, and its trace is $\|\mathfrak{Y}\|_{2k}^{2k}$ (Definition~\ref{def:even-schatten}).
This is~\eqref{eq:tori-fold-norm}.

\emph{Step 3: the inequality~\eqref{eq:tori-holder}.}
For $j=1,\dots,2k$ the operator $\mathfrak{Y}_j:=\mathfrak{S}[G_j]\colon\mathcal{H}(c_{j-1})\to\mathcal{H}(c_j)$ is Hilbert--Schmidt.
Theorem~\ref{thm:slab-operators}\eqref{it:slab-trace} for the torus $\mathbb{T}$, with the planes $c_0<c_1<\dots<c_{2k}=c_0+p$ and the closing vector $g$, reads
\[
Z_{\mathbb{T}}\left[G_1\cdots G_{2k}\right]=\Tr_{\mathcal{H}(c_0)}\left(U\cdot\mathfrak{Y}_{2k}\cdots\mathfrak{Y}_1\right),\qquad
U:=V_g^{-1}\cdot U_{\mathrm{f}}^{\frac{1+\chi(g)}{2}}\colon\mathcal{H}(c_0+p)\to\mathcal{H}(c_0),
\]
and $U$, a product of two unitaries, is unitary.
Lemma~\ref{lem:holder-spaces}, applied to the spaces $\mathcal{H}(c_0),\dots,\mathcal{H}(c_{2k})$, to the operators $\mathfrak{Y}_1,\dots,\mathfrak{Y}_{2k}$ and to $U$, together with~\eqref{eq:tori-fold-norm}, yields
\[
\left|Z_{\mathbb{T}}\left[G_1\cdots G_{2k}\right]\right|\le\prod_{j=1}^{2k}\|\mathfrak{Y}_j\|_{2k}=\prod_{j=1}^{2k}Z_{\mathbb{T}^0}\left[\fold_\ell G_j\right]^{\frac{1}{2k}} .
\]
\end{proof}

\begin{remark}\label{rem:twist}
Let the notation be that of Proposition~\ref{prop:twisted-holder} and of its proof.
\begin{enumerate}
\item\label{it:twist-unitary}
The bound of Lemma~\ref{lem:holder-spaces} holds for every unitary operator $\mathcal{H}(c_0+p)\to\mathcal{H}(c_0)$ in the place of $U$.
Since $\ell(g)=\ell(g^0)=p$, the vector $g^0-g$ lies in $\Z^3\cap\{\ell=0\}$, though in general not in $\Gamma_\parallel$;
and by the group law of the unitaries $V_v$,
\[
V_g^{-1}\cdot U_{\mathrm{f}}^{\frac{1+\chi(g)}{2}}=V_{g^0}^{-1}\cdot V_{g^0-g}\cdot U_{\mathrm{f}}^{\frac{1+\chi(g)}{2}},
\]
where $V_{g^0}^{-1}$ is the twist of the trace formula of $\mathbb{T}^0$ in Step~2 of the proof.
Thus the twist of $\mathbb{T}$ differs from that of $\mathbb{T}^0$ by the glide $V_{g^0-g}$ and, if $\chi(g)=+1$, by the parity $U_{\mathrm{f}}$;
these two unitaries disappear together in~\eqref{eq:tori-holder}.
\item\label{it:twist-even}
The number $2k$ of slabs in~\eqref{eq:tori-fold-norm} is even.
Let $r\ge3$ be odd, suppose that $c_0+i\cdot\frac{p}{r}\in\Adm(\ell)$ for all $i\in\Z$,
form the positions $c_i$, the slabs, the group $W(\ell;c_0,h')$, the elements $g_{i\leftarrow j}$ and the copies of Definition~\ref{def:slab-copies} with $h':=\frac{p}{r}$,
and let $j\in\Z$.
By that definition $g_{j+r\leftarrow j}=\theta_{\ell,c_{j+(r-1)/2}}$, a reflection, since $r$ is odd.
The translation $\tau_{g^0}$ by the closing vector $g^0$ of $\mathbb{T}^0$, on the other hand, does not belong to $W(\ell;c_0,h')$:
the translations in this group are the $\tau_{s\cdot v_{h'}}$, $s\in\Z$, whose heights $\ell\left(s\cdot v_{h'}\right)=2s\cdot h'$ lie in $2h'\Z$,
whereas $\ell(g^0)=p=r\cdot h'$ is an odd multiple of $h'$.
Hence, for a general $G\in\mathcal{A}_\Xi[c_{j-1},c_j]$, the copy $G^{(j+r)}$ of $G$ in the slab $j+r$ is not the translate $G^{\tau_{g^0}}$ of $G=G^{(j)}$ by the period $g^0$ of $\mathbb{T}^0$:
the copies of $G$ do not close up along one height period, and no fold of $G$ over $r$ slabs is defined.
The number $\Tr\left(\left(\mathfrak{S}[G]^*\cdot\mathfrak{S}[G]\right)^{\frac{r}{2}}\right)$ is finite, as $\mathfrak{S}[G]^*\cdot\mathfrak{S}[G]$ is non-negative and of trace class and $\frac{r}{2}\ge1$;
but for a general $G$ the identity~\eqref{eq:tori-fold-norm} has no counterpart for this number.
\item\label{it:twist-chessboard}
If $\mathbb{T}$ is adapted to $\ell$, then $\mathbb{T}^0=\mathbb{T}$ (Definition~\ref{def:slabs}), and~\eqref{eq:tori-holder} is the chessboard estimate in one direction,
\[
\left|Z_{\mathbb{T}}\left[G_1\cdots G_{2k}\right]\right|\le\prod_{j=1}^{2k}Z_{\mathbb{T}}\left[\fold_\ell G_j\right]^{\frac{1}{2k}} .
\]
For Yang--Mills theory in four dimensions the corresponding estimate is proved for positive functions, by repeated application of the Cauchy--Schwarz inequality;
the proof of Proposition~\ref{prop:twisted-holder} does not require the integrands $G_j$ to be positive.
\end{enumerate}
\end{remark}

\subsection{From cubic to adapted tori}\label{sec:cubic-to-adapted}

\begin{lemma}[Systoles]\label{lem:systoles}
Let $a,b,c$ be positive integers and $i_1,i_2,i_3$ integers.
The vector $v:=i_1\cdot a\cdot e_-+i_2\cdot b\cdot e_++i_3\cdot c\cdot e_3$ has the coordinates
$\left(i_1\cdot a+i_2\cdot b,\ i_2\cdot b-i_1\cdot a,\ i_3\cdot c\right)$, and its maximum norm is
\[
|v|_\infty=\max\left(|i_1|\cdot a+|i_2|\cdot b,\ |i_3|\cdot c\right).
\]
Consequently
\[
\sys\left(\Z a\,e_-+\Z b\,e_++\Z c\,e_3\right)=\min(a,b,c),
\]
and for positive integers $n_1,n_2,n_3$
\[
\sys\left(n_1\Z e_1+n_2\Z e_2+n_3\Z e_3\right)=\min_\mu n_\mu .
\]
More generally, let $\mathfrak{F}=(\ell^1,\ell^2,\ell^3)$ be a frame, with the normals $n_{\ell^1},n_{\ell^2},n_{\ell^3}$,
and let $p_1,p_2,p_3$ be positive integers with $\frac{p_a}{|n_{\ell^a}|^2}\cdot n_{\ell^a}\in\Z^3$ for $a=1,2,3$.
Then
\[
\sys\left(\Z\,\frac{p_1}{|n_{\ell^1}|^2}\cdot n_{\ell^1}+\Z\,\frac{p_2}{|n_{\ell^2}|^2}\cdot n_{\ell^2}+\Z\,\frac{p_3}{|n_{\ell^3}|^2}\cdot n_{\ell^3}\right)=\min_a\frac{p_a}{|n_{\ell^a}|^2}.
\]
Each of these lattices is a period lattice if and only if its systole is at least four.
For the last one this means $p_a\ge4\,|n_{\ell^a}|^2$ for $a=1,2,3$, and it is then the lattice of the rectangular torus $\mathbb{T}(\mathfrak{F};p_1,p_2,p_3)$.
\end{lemma}

\begin{proof}
The coordinates of $v$ follow from $e_\pm=e_1\pm e_2$.
Both $|i_1\cdot a+i_2\cdot b|$ and $|i_2\cdot b-i_1\cdot a|$ are at most $|i_1|\cdot a+|i_2|\cdot b$, since $a,b>0$;
the first equals $|i_1|\cdot a+|i_2|\cdot b$ if $i_1\cdot i_2\ge0$, and the second if $i_1\cdot i_2\le0$.
So the larger of the moduli of the first two coordinates of $v$ is $|i_1|\cdot a+|i_2|\cdot b$,
which gives the formula for $|v|_\infty$.

Let $v\neq0$.
If $(i_1,i_2)\neq(0,0)$, then $|v|_\infty\ge|i_1|\cdot a+|i_2|\cdot b\ge\min(a,b)$;
otherwise $i_3\neq0$ and $|v|_\infty=|i_3|\cdot c\ge c$.
The coefficients $(1,0,0)$, $(0,1,0)$ and $(0,0,1)$ give the values $a$, $b$ and $c$,
so the first lattice has the systole $\min(a,b,c)$.
A non-zero vector $(i_1\cdot n_1,i_2\cdot n_2,i_3\cdot n_3)$ of the second lattice has the maximum norm $\max_\mu|i_\mu|\cdot n_\mu\ge\min_\mu n_\mu$,
and the value $\min_\mu n_\mu$ is attained at $n_\mu\cdot e_\mu$ for an index $\mu$ of the smallest side.

Let now $\mathfrak{F}$ and $p_1,p_2,p_3$ be as in the statement.
Renumbering the families of $\mathfrak{F}$ together with the numbers $p_a$ changes neither the lattice nor $\min_a p_a/|n_{\ell^a}|^2$.
The normals of the coordinate frame $\mathfrak{F}_{\mathrm{c}}$ are $e_1,e_2,e_3$, and those of the diagonal frame $\mathfrak{F}_{\mathrm{d}}=(\ell_-,\ell_+,x_3)$ are $e_-,e_+,e_3$.
By Definition~\ref{def:cell-structures} one may therefore assume that there is $g\in\Ocub$ such that, with suitable signs,
either $n_{\ell^a}=\pm g\,e_a$ for $a=1,2,3$,
or $n_{\ell^1}=\pm g\,e_-$, $n_{\ell^2}=\pm g\,e_+$ and $n_{\ell^3}=\pm g\,e_3$.
Being a signed permutation matrix, $g$ preserves the Euclidean length, the maximum norm and the set $\Z^3$.
Hence $g^{-1}$ maps the lattice in question onto a lattice of the same systole,
and it maps the integer vectors $\frac{p_a}{|n_{\ell^a}|^2}\cdot n_{\ell^a}$ to integer vectors.
In the first case $|n_{\ell^a}|^2=1$ for every $a$, and the image lattice is $p_1\Z e_1+p_2\Z e_2+p_3\Z e_3$, of systole $\min_a p_a$.
In the second case $|n_{\ell^1}|^2=|n_{\ell^2}|^2=2$ and $|n_{\ell^3}|^2=1$, and the image lattice is $\Z\frac{p_1}{2}\,e_-+\Z\frac{p_2}{2}\,e_++\Z p_3\,e_3$,
where $\frac{p_1}{2}$ and $\frac{p_2}{2}$ are integers, the vectors $\pm\frac{p_1}{2}\cdot e_-$ and $\pm\frac{p_2}{2}\cdot e_+$ being integer vectors;
its systole is $\min\left(\frac{p_1}{2},\frac{p_2}{2},p_3\right)$ by what has been shown.
In both cases the systole is $\min_a\frac{p_a}{|n_{\ell^a}|^2}$.

All lattices of the lemma are subgroups of $\Z^3$ of rank three, since their three generators are linearly independent.
By Definition~\ref{def:configurations} such a subgroup is a period lattice exactly when its systole is at least four,
and $\min_a\frac{p_a}{|n_{\ell^a}|^2}\ge4$ means $p_a\ge4\,|n_{\ell^a}|^2$ for every $a$.
\end{proof}

\begin{proposition}[Three untwistings]\label{prop:untwistings}
For a torus $\mathbb{T}=(\Gamma,\chi)$ and a family $\ell$ let $\Xi(\mathbb{T},\ell)=(\Gamma_\parallel,\chi_\parallel)$ be the cross-section,
$p=p(\mathbb{T},\ell)$ the height period and $\mathbb{T}^0(\ell)$ the straightened torus, as in Definition~\ref{def:slabs}.
\begin{enumerate}
\item\label{it:untwist-cubic}
Let $n\ge8$ be even, $\mathbb{T}=C(n,n,n)$ and $\ell=\ell_-$.
Then $\Gamma_\parallel=n\Z e_++n\Z e_3$, with $\chi_\parallel(n\cdot e_+)=+1$ and $\chi_\parallel(n\cdot e_3)=-1$;
the vector $g=n\cdot e_1$ is a closing vector, and $\chi(g)=-1$;
the height period is $p=n$;
and $\mathbb{T}^0(\ell_-)$ is defined and equal to $D'(n/2,n,n)$.
\item\label{it:untwist-parity}
Let $n\ge8$ be even, $\mathbb{T}=D'(n/2,n,n)$ and $\ell=\ell_+$.
Then $\Gamma_\parallel=\frac{n}{2}\Z e_-+n\Z e_3$, and $\chi_\parallel$ takes the value $-1$ at $\tfrac{n}{2}\cdot e_-$ and at $n\cdot e_3$;
the vector $g=n\cdot e_+$ is a closing vector, and $\chi(g)=+1$;
the height period is $p=2n$;
and $\mathbb{T}^0(\ell_+)$ is defined and equal to $D(n/2,n,n)$.
\item\label{it:untwist-axis}
Let $a\ge4$ and $c\ge4$ be integers, $\mathbb{T}=D(a,2a,c)$ and $\ell=x_1$.
Then $\Gamma_\parallel=4a\Z e_2+c\Z e_3$, and $\chi_\parallel$ takes the value $-1$ at $4a\cdot e_2$ and at $c\cdot e_3$;
the vector $g=a\cdot e_-$ is a closing vector, and $\chi(g)=-1$;
the height period is $p=a$;
and $\mathbb{T}^0(x_1)$ is defined and equal to $C(a,4a,c)$.
\item\label{it:split-angles}
Let $\Gamma$ be a lattice in $\R^k$, that is, a discrete subgroup of $\R^k$ of rank $k$, and let $\hat n,\hat n'\in\R^k$ be unit vectors.
Suppose that $\Gamma$ is split by the reflections in $\hat n^\perp$ and in $(\hat n')^\perp$, that is,
\[
\Gamma=\left(\Gamma\cap\R\hat n\right)+\left(\Gamma\cap\hat n^\perp\right)
\qquad\text{and}\qquad
\Gamma=\left(\Gamma\cap\R\hat n'\right)+\left(\Gamma\cap(\hat n')^\perp\right).
\]
Then $\hat n\cdot\hat n'\in\{0,1,-1\}$.
\item\label{it:cut-in-two}
Let $\Gamma$ be a lattice in $\R^k$, let $\hat n\in\R^k$ be a unit vector and $c\in\R$, and define $\theta\colon\R^k\to\R^k$ by
\[
\theta(x):=x-2\cdot\left(\hat n\cdot x-c\right)\cdot\hat n .
\]
Assume that the linear part of $\theta$ maps $\Gamma$ onto $\Gamma$;
then $x+\Gamma\mapsto\theta(x)+\Gamma$ is a well-defined map of the flat torus $\R^k/\Gamma$ onto itself.
The complement in $\R^k/\Gamma$ of the fixed point set of this map has precisely two connected components, which the map interchanges,
when $\Gamma$ is split by the reflection in $\hat n^\perp$;
it is connected when $\Gamma$ is not split by the reflection in $\hat n^\perp$.
\end{enumerate}
Consequently, if a torus is adapted to two families $\ell$ and $\ell'$ (Definition~\ref{def:slabs}), then $\hat n_\ell\cdot\hat n_{\ell'}\in\{0,1,-1\}$.
If $\ell=x_\lambda$ is an axis family and $\ell'=x_\mu\pm x_\nu$ a diagonal family with $\lambda\in\{\mu,\nu\}$, then $|\hat n_\ell\cdot\hat n_{\ell'}|=1/\sqrt{2}$,
and the mirrors of the two families meet at an angle of $45^\circ$;
if $\ell$ and $\ell'$ are diagonal families with exactly one common index, then $|\hat n_\ell\cdot\hat n_{\ell'}|=1/2$, and their normals make an angle of $60^\circ$ or $120^\circ$.
Hence no pair of families of either of these two kinds admits a torus adapted to both of its members.
If a torus is adapted to two different families, their normals are orthogonal;
and the families to which a given torus is adapted are all among the three families of a single frame (Definition~\ref{def:cell-structures}).
A \emph{mirror} of a flat torus $\R^k/\Gamma$ is the fixed point set of the map of $\R^k/\Gamma$ induced by a map $\theta$ as in \eqref{it:cut-in-two};
its normal is $\hat n$, and it \emph{cuts $\R^k/\Gamma$ in two} if its complement in $\R^k/\Gamma$ consists of exactly two connected components.
Two mirrors of $\R^k/\Gamma$ whose normals $\hat n$ and $\hat n'$ satisfy $|\hat n\cdot\hat n'|=1/\sqrt{2}$, that is, two mirrors at an angle of $45^\circ$, do not both cut $\R^k/\Gamma$ in two.
\end{proposition}

\begin{proof}
Under the hypotheses of \eqref{it:untwist-cubic}--\eqref{it:untwist-axis} every lattice that occurs below has systole at least four, by Lemma~\ref{lem:systoles}:
the lattice of $C(n,n,n)$ has the systole $n$, the common lattice of $D'(n/2,n,n)$ and $D(n/2,n,n)$ the systole $n/2\ge4$,
and the lattices of $D(a,2a,c)$ and $C(a,4a,c)$ the systole $\min(a,c)\ge4$.
For a torus $\mathbb{T}=(\Gamma,\chi)$ and a family $\ell$ one has $\ell(\Gamma)=p\Z$ with $p=p(\mathbb{T},\ell)$,
and the closing vectors are the elements $g\in\Gamma$ with $\ell(g)=p$ (Lemma~\ref{lem:closing-vectors}).
Put
\[
g^0:=\frac{p}{|n_\ell|^2}\,n_\ell .
\]
By Definition~\ref{def:slabs}, $\mathbb{T}^0(\ell)$ is defined if $g^0\in\Z^3$ and $\sys(\Gamma_\parallel+\Z g^0)\ge4$;
its lattice is then $\Gamma_\parallel+\Z g^0$, and its character agrees with $\chi_\parallel$ on $\Gamma_\parallel$ and takes the value $-1$ at $g^0$.

\emph{Item \eqref{it:untwist-cubic}.}
Here $\Gamma=n\Z^3$ and $\chi(n\cdot e_\mu)=-1$ for $\mu=1,2,3$.
For integers $i_1,i_2,i_3$ one has $\ell_-\left(n\cdot(i_1,i_2,i_3)\right)=n\cdot(i_1-i_2)$, whence
\[
\Gamma\cap\{\ell_-=0\}=\left\{n\cdot(i_1,i_1,i_3):\ i_1,i_3\in\Z\right\}=n\Z e_++n\Z e_3,
\]
with $\chi(n\cdot e_+)=\chi(n\cdot e_1)\cdot\chi(n\cdot e_2)=+1$ and $\chi(n\cdot e_3)=-1$.
The same formula gives $\ell_-(\Gamma)=n\Z$, so that $p=n$;
the vector $g=n\cdot e_1$ lies in $\Gamma$, satisfies $\ell_-(g)=n$, and $\chi(g)=-1$.
Since $n_{\ell_-}=e_-$ and $|n_{\ell_-}|^2=2$, one has $g^0=\tfrac{n}{2}\cdot e_-$, which belongs to $\Z^3$ since $n$ is even.
The lattice $n\Z e_++n\Z e_3+\Z g^0$ is the lattice of $D'(n/2,n,n)$, of systole $n/2\ge4$;
thus $\mathbb{T}^0(\ell_-)$ is defined.
Its character takes the value $+1$ at $n\cdot e_+$ and the value $-1$ at $n\cdot e_3$ and at $\tfrac{n}{2}\cdot e_-$,
so $\mathbb{T}^0(\ell_-)=D'(n/2,n,n)$ by Definition~\ref{def:tori}.

\emph{Item \eqref{it:untwist-parity}.}
Now $\Gamma=\frac{n}{2}\Z e_-+n\Z e_++n\Z e_3$, and $\ell_+(e_-)=\ell_+(e_3)=0$, $\ell_+(e_+)=2$.
For integers $i_1,i_2,i_3$ therefore
\[
\ell_+\left(i_1\cdot\tfrac{n}{2}\cdot e_-+i_2\cdot n\cdot e_++i_3\cdot n\cdot e_3\right)=2n\cdot i_2 .
\]
It follows that $\Gamma\cap\{\ell_+=0\}=\frac{n}{2}\Z e_-+n\Z e_3$, at whose two generators $\chi$ takes the value $-1$ (Definition~\ref{def:tori}),
and that $\ell_+(\Gamma)=2n\Z$, so that $p=2n$;
the vector $g=n\cdot e_+$ satisfies $\ell_+(g)=2n$ and $\chi(g)=+1$.
Here $g^0=\frac{2n}{2}\cdot e_+=n\cdot e_+$ belongs to $\Gamma$, so $\Gamma_\parallel+\Z g^0=\Gamma$, of systole $n/2\ge4$, and $\mathbb{T}^0(\ell_+)$ is defined.
Its character agrees with $\chi$ on $\Gamma_\parallel$ and takes the value $-1$ at $n\cdot e_+$; hence $\mathbb{T}^0(\ell_+)=D(n/2,n,n)$.

\emph{Item \eqref{it:untwist-axis}.}
Here $\Gamma=a\Z e_-+2a\Z e_++c\Z e_3$, with $\chi=-1$ at the three generators.
For integers $i_1,i_2,i_3$,
\[
i_1\cdot a\cdot e_-+i_2\cdot 2a\cdot e_++i_3\cdot c\cdot e_3=\left((i_1+2i_2)\cdot a,\ (2i_2-i_1)\cdot a,\ i_3\cdot c\right).
\]
The first coordinate is zero exactly when $i_1=-2i_2$.
In that case the vector equals $i_2\cdot4a\cdot e_2+i_3\cdot c\cdot e_3$, and the value of $\chi$ at it is
\[
\chi(a\cdot e_-)^{-2i_2}\cdot\chi(2a\cdot e_+)^{i_2}\cdot\chi(c\cdot e_3)^{i_3}=(-1)^{i_2+i_3}.
\]
So $\Gamma\cap\{x_1=0\}=4a\Z e_2+c\Z e_3$, with the value $-1$ of $\chi$ at both generators.
As $i_1+2i_2$ takes every integer value, $x_1(\Gamma)=a\Z$ and $p=a$;
the vector $g=a\cdot e_-$ satisfies $x_1(g)=a$ and $\chi(g)=-1$.
With $n_{x_1}=e_1$ one finds $g^0=a\cdot e_1$,
and $4a\Z e_2+c\Z e_3+a\Z e_1$ is the lattice of $C(a,4a,c)$, of systole $\min(a,c)\ge4$; thus $\mathbb{T}^0(x_1)$ is defined.
Its character is $-1$ at $4a\cdot e_2$, at $c\cdot e_3$ and at $a\cdot e_1$, and $\mathbb{T}^0(x_1)=C(a,4a,c)$.

\emph{Item \eqref{it:split-angles}.}
Since $\R\hat n\cap\hat n^\perp=\{0\}$, the sum $\left(\Gamma\cap\R\hat n\right)+\left(\Gamma\cap\hat n^\perp\right)$ is direct.
A discrete subgroup of a real vector space has rank at most the dimension of the space,
so $\Gamma\cap\R\hat n$ has rank at most one and $\Gamma\cap\hat n^\perp$ rank at most $k-1$; hence
\[
k=\operatorname{rank}\Gamma=\operatorname{rank}\left(\Gamma\cap\R\hat n\right)+\operatorname{rank}\left(\Gamma\cap\hat n^\perp\right)\le1+(k-1),
\]
and equality holds in both places.
In particular $\Gamma\cap\R\hat n$ is a discrete subgroup of rank one of the line $\R\hat n$;
it is generated by $\alpha\cdot\hat n$ for some $\alpha>0$.
In the same way $\Gamma\cap\R\hat n'$ is generated by $\alpha'\cdot\hat n'$ with $\alpha'>0$.
Write $\gamma\in\Gamma$ as $\gamma=\gamma_1+\gamma_2$ with $\gamma_1\in\Gamma\cap\R\hat n$ and $\gamma_2\in\Gamma\cap\hat n^\perp$.
Then $\gamma_1=(\hat n\cdot\gamma)\cdot\hat n$ is the orthogonal projection of $\gamma$ to $\R\hat n$, and therefore $\hat n\cdot\gamma\in\alpha\Z$ for every $\gamma\in\Gamma$.
For $\gamma=\alpha'\cdot\hat n'$ this gives $\alpha'\cdot(\hat n\cdot\hat n')\in\alpha\Z$;
the splitting at $(\hat n')^\perp$, applied in the same way to $\gamma=\alpha\cdot\hat n$, gives $\alpha\cdot(\hat n\cdot\hat n')\in\alpha'\Z$.
Suppose that $\hat n\cdot\hat n'\neq0$.
Then the quotients $\alpha'\cdot(\hat n\cdot\hat n')/\alpha$ and $\alpha\cdot(\hat n\cdot\hat n')/\alpha'$ are non-zero integers,
and their product $(\hat n\cdot\hat n')^2$ is a positive integer.
As $|\hat n\cdot\hat n'|\le1$, it follows that $(\hat n\cdot\hat n')^2=1$.

\emph{Item \eqref{it:cut-in-two}.}
The linear part of $\theta$ is $\vec\theta(x)=x-2\cdot(\hat n\cdot x)\cdot\hat n$, and $\theta(x+\gamma)=\theta(x)+\vec\theta\gamma$;
as $\vec\theta\,\Gamma=\Gamma$, the induced map $x+\Gamma\mapsto\theta(x)+\Gamma$ is well defined.
Since $\Gamma$ has rank $k$, it is not contained in the hyperplane $\hat n^\perp$; choose $\gamma_0\in\Gamma$ with $\hat n\cdot\gamma_0\neq0$.
Then $\gamma_0-\vec\theta\gamma_0=2\cdot(\hat n\cdot\gamma_0)\cdot\hat n$ belongs to $\Gamma\cap\R\hat n$ and is different from $0$.
Consequently $\Gamma\cap\R\hat n$ is a non-zero discrete subgroup of the line $\R\hat n$, generated by $\alpha\cdot\hat n$ for some $\alpha>0$,
and for every $\gamma\in\Gamma$ the element $\gamma-\vec\theta\gamma=2\cdot(\hat n\cdot\gamma)\cdot\hat n$ of $\Gamma\cap\R\hat n$ shows that $2\cdot(\hat n\cdot\gamma)\in\alpha\Z$.
The group
\[
\hat n\cdot\Gamma:=\{\hat n\cdot\gamma:\ \gamma\in\Gamma\}
\]
therefore contains $\alpha\Z$, because $\alpha\cdot\hat n\in\Gamma$, and is contained in $\frac{\alpha}{2}\Z$.
As $\alpha\Z$ has index two in $\frac{\alpha}{2}\Z$, either $\hat n\cdot\Gamma=\alpha\Z$ or $\hat n\cdot\Gamma=\frac{\alpha}{2}\Z$.
If $\hat n\cdot\Gamma=\alpha\Z$, then $(\hat n\cdot\gamma)\cdot\hat n\in\Gamma$ and $\gamma-(\hat n\cdot\gamma)\cdot\hat n\in\Gamma\cap\hat n^\perp$ for every $\gamma\in\Gamma$,
so $\Gamma$ is split by the reflection in $\hat n^\perp$.
Conversely, if $\Gamma$ is split, then $(\hat n\cdot\gamma)\cdot\hat n\in\Gamma\cap\R\hat n$ for every $\gamma\in\Gamma$, as in the proof of \eqref{it:split-angles},
and $\hat n\cdot\Gamma=\alpha\Z$.
Thus $\Gamma$ is split by the reflection in $\hat n^\perp$ if and only if $\hat n\cdot\Gamma=\alpha\Z$.

A point $x+\Gamma$ is fixed by the induced map exactly when $\theta(x)-x=-2\cdot(\hat n\cdot x-c)\cdot\hat n$ belongs to $\Gamma$,
that is, to $\Gamma\cap\R\hat n$, that is, exactly when $\hat n\cdot x\in c+\frac{\alpha}{2}\Z$.
This condition is invariant under translations by $\Gamma$, since $\hat n\cdot\Gamma\subset\frac{\alpha}{2}\Z$.
With the open slabs
\[
U_i:=\left\{x\in\R^k:\ c+i\cdot\frac{\alpha}{2}<\hat n\cdot x<c+(i+1)\cdot\frac{\alpha}{2}\right\},\qquad i\in\Z,
\]
the preimage in $\R^k$ of the set of the non-fixed points is therefore $\bigcup_{i\in\Z}U_i$.
Each $U_i$ is open and convex, so its image in $\R^k/\Gamma$ is open and connected.
For $\gamma\in\Gamma$ one has $U_i+\gamma=U_{i'}$ with $(i'-i)\cdot\frac{\alpha}{2}=\hat n\cdot\gamma$.
Since the slabs $U_i$ are pairwise disjoint, the images of $U_i$ and $U_{i'}$ coincide if $(i'-i)\cdot\frac{\alpha}{2}\in\hat n\cdot\Gamma$, and they are disjoint otherwise.
If $\Gamma$ is split, then $\hat n\cdot\Gamma=\alpha\Z$:
the image of $U_i$ is the image of $U_0$ for even $i$ and the image of $U_1$ for odd $i$, and these two images are disjoint.
So the set of the points of $\R^k/\Gamma$ that are not fixed is the union of two disjoint, non-empty, open and connected sets, which are therefore its connected components.
Moreover $\hat n\cdot\theta(x)-c=-(\hat n\cdot x-c)$, so that $\theta$ maps $U_i$ onto $U_{-1-i}$;
as $i$ and $-1-i$ have different parities, the induced map exchanges the two components.
If $\Gamma$ is not split, then $\hat n\cdot\Gamma=\frac{\alpha}{2}\Z$, all the images coincide with the image of $U_0$, and the complement is connected.

\emph{The consequences.}
Let $\mathbb{T}=(\Gamma,\chi)$ be adapted to $\ell$, with the height period $p$.
Then $\mathbb{T}=\mathbb{T}^0(\ell)$, so that $\Gamma=\left(\Gamma\cap\{\ell=0\}\right)+\Z\,\frac{p}{|n_\ell|^2}\,n_\ell$ (Definition~\ref{def:slabs}).
Since $\{\ell=0\}=\hat n_\ell^\perp$ and $\frac{p}{|n_\ell|^2}\,n_\ell\in\Gamma\cap\R\hat n_\ell$, the lattice $\Gamma$ is split by the reflection in $\hat n_\ell^\perp$.
If $\mathbb{T}$ is adapted to $\ell$ and to $\ell'$, item \eqref{it:split-angles} with $k=3$ gives $\hat n_\ell\cdot\hat n_{\ell'}\in\{0,1,-1\}$.

The normal of an axis family $x_\lambda$ is $e_\lambda$, and that of a diagonal family $x_\mu\pm x_\nu$ is $(e_\mu\pm e_\nu)/\sqrt{2}$.
Hence $\hat n_{x_\lambda}\cdot\hat n_{x_\mu\pm x_\nu}$ equals $\pm1/\sqrt{2}$ if $\lambda\in\{\mu,\nu\}$ and $0$ otherwise;
two different axis families have orthogonal normals;
and for two different diagonal families the scalar product of the normals is $0$ if they have the same two indices,
that is, for $x_\mu-x_\nu$ and $x_\mu+x_\nu$, and $\pm1/2$ if they have exactly one common index, which is the only other possibility.
In particular the scalar product of the normals of two different families is never $\pm1$,
and if one torus is adapted to two different families, their normals are orthogonal.
The pairs of different families with orthogonal normals are therefore $\{x_\mu,x_\nu\}$, $\{x_\lambda,x_\mu\pm x_\nu\}$ with $\lambda\notin\{\mu,\nu\}$, and $\{x_\mu-x_\nu,x_\mu+x_\nu\}$.
A set of families with pairwise orthogonal normals that contains a diagonal family $x_\mu\pm x_\nu$ is consequently contained in $\{x_\mu-x_\nu,\ x_\mu+x_\nu,\ x_\lambda\}$, $\lambda\notin\{\mu,\nu\}$,
and one that contains no diagonal family is contained in $\{x_1,x_2,x_3\}$;
both triples have pairwise orthogonal normals and form frames (Definition~\ref{def:cell-structures}).

Finally, let two mirrors of a flat torus $\R^k/\Gamma$, with the normals $\hat n$ and $\hat n'$, both cut $\R^k/\Gamma$ in two.
By \eqref{it:cut-in-two}, the complement of the mirror with the normal $\hat n$ would be connected if $\Gamma$ were not split by the reflection in $\hat n^\perp$;
so $\Gamma$ is split by this reflection, and likewise by the reflection in $(\hat n')^\perp$.
Item \eqref{it:split-angles} now gives $\hat n\cdot\hat n'\in\{0,1,-1\}$, and in particular $|\hat n\cdot\hat n'|\neq1/\sqrt{2}$.
\end{proof}

\begin{remark}\label{rem:adapted-families}
The coordinate torus $C(n_1,n_2,n_3)$ is adapted to $x_1$, $x_2$ and $x_3$, and $D(a,b,c)$ is adapted to $\ell_-$, $\ell_+$ and $x_3$,
both being rectangular tori (Definition~\ref{def:cell-structures}).
The lattice $\Z a\,e_-+\Z b\,e_++\Z c\,e_3$, which $D'(a,b,c)$ shares with $D(a,b,c)$, is split by the reflections in $e_-^\perp$, $e_+^\perp$ and $e_3^\perp$,
but $D'(a,b,c)$ is adapted to $\ell_-$ and $x_3$ only, because $\chi(b\cdot e_+)=+1$.
At the mirrors of $\ell_+$ the form of $D'(a,b,c)$ contains the fermion parity $U_{\mathrm{f}}$:
let $c_0\in\Z$, let $\Xi:=\Xi(D'(a,b,c),\ell_+)$, whose lattice is $\Gamma_\parallel=\Z a\,e_-+\Z c\,e_3$, and let $G\in\mathcal{A}_\Xi[c_0,c_0+b]$; then
\[
Z_{D'(a,b,c)}\left[\Theta_{\ell_+,c_0}(G)\cdot G\right]=\Tr_{\mathcal{H}(c_0)}\left(U_{\mathrm{f}}\cdot\mathfrak{S}[G]^{*}\cdot\mathfrak{S}[G]\right).
\]
To see this, put $G':=\Theta_{\ell_+,c_0}(G)$.
It lies in $\mathcal{A}_\Xi[c_0-b,c_0]$ (Lemma~\ref{lem:admissible-images}\eqref{it:images-boxes}, Definition~\ref{def:mirrors}\eqref{it:mirror-height};
the linear part of $\theta_{\ell_+,c_0}$ fixes $\Gamma_\parallel$ pointwise), and $\Theta_{\ell_+,c_0}(G')=G$.
The reflection rule for admissible products, Corollary~\ref{cor:slab-admissible}(e), applied to $G'$ and to the positions $c_0-b<c_0$,
involves the vector $\frac{2b}{|n_{\ell_+}|^2}\cdot n_{\ell_+}=b\cdot e_+$;
it gives $\mathfrak{S}[G]=V_{b\cdot e_+}\cdot\mathfrak{S}[G']^{*}$, that is, $\mathfrak{S}[G']=\mathfrak{S}[G]^{*}\cdot V_{b\cdot e_+}$.
The trace formula of Corollary~\ref{cor:slab-admissible}(f), applied with the closing vector $b\cdot e_+$, at which the character equals $+1$,
and with the two slabs $\{c_0-b<\ell_+<c_0\}$ and $\{c_0<\ell_+<c_0+b\}$, gives
\[
Z_{D'(a,b,c)}\left[G'\cdot G\right]=\Tr_{\mathcal{H}(c_0-b)}\left(V_{b\cdot e_+}^{-1}\cdot U_{\mathrm{f}}\cdot\mathfrak{S}[G]\cdot\mathfrak{S}[G]^{*}\cdot V_{b\cdot e_+}\right)
=\Tr_{\mathcal{H}(c_0+b)}\left(U_{\mathrm{f}}\cdot\mathfrak{S}[G]\cdot\mathfrak{S}[G]^{*}\right),
\]
the second equality because $V_{b\cdot e_+}$ is a unitary operator from $\mathcal{H}(c_0-b)$ onto $\mathcal{H}(c_0+b)$ (Corollary~\ref{cor:slab-admissible}(a)).
The last trace equals $\Tr_{\mathcal{H}(c_0)}\left(\mathfrak{S}[G]^{*}\cdot U_{\mathrm{f}}\cdot\mathfrak{S}[G]\right)=\Tr_{\mathcal{H}(c_0)}\left(U_{\mathrm{f}}\cdot\mathfrak{S}[G]^{*}\cdot\mathfrak{S}[G]\right)$,
because the trace is cyclic on products of Hilbert--Schmidt operators and $U_{\mathrm{f}}$ commutes with $\mathfrak{S}[G]$ (Corollary~\ref{cor:slab-admissible}(d)).

For the lattice $n\Z^3$ of the cubic torus $C(n,n,n)$ one has $n\Z^3\cap\R e_-=\Z n\,e_-$, generated by $\alpha\cdot\hat n_{\ell_-}$ with $\alpha=n\sqrt{2}$,
whereas $\hat n_{\ell_-}\cdot n\Z^3=\frac{n}{\sqrt{2}}\Z=\frac{\alpha}{2}\Z$.
If $n\Z^3$ were split by the reflection in $e_-^\perp$, the component $(\hat n_{\ell_-}\cdot\gamma)\cdot\hat n_{\ell_-}$ of each $\gamma\in n\Z^3$ in $\R e_-$ would lie in $\Z n\,e_-$,
and $\hat n_{\ell_-}\cdot n\Z^3$ would be contained in $\alpha\Z$.
So the lattice of the cubic torus is not split by the reflection in $e_-^\perp$.
\end{remark}

\subsection{Chessboard estimates}\label{sec:chessboard}

\begin{definition}[Families of cell contents]\label{def:cell-families}
Let $\mathcal{Q}=(\mathfrak{F};o_1,o_2,o_3;h)$ be a cell structure (Definition~\ref{def:cell-structures}) with the frame $\mathfrak{F}=(\ell^1,\ell^2,\ell^3)$.
As in that definition, $n_a$ denotes the normal of $\ell^a$, and $Q_\kappa$, $\kappa\in\Z^3$, denote the open cells.
An element of $\mathcal{A}(Q_\kappa)$ is called a \emph{content} of the cell $Q_\kappa$.

\emph{Cell indices of translations.}
For $v\in\Z^3$ with $\ell^a(v)\in h\Z$ for $a=1,2,3$ put
\[
\kappa_v:=\left(\frac{\ell^1(v)}{h},\frac{\ell^2(v)}{h},\frac{\ell^3(v)}{h}\right)\in\Z^3 .
\]
Since $\ell^a(x+v)=\ell^a(x)+h\cdot(\kappa_v)_a$ for $x\in\R^3$, one has $\tau_vQ_\kappa=Q_{\kappa+\kappa_v}$ for every $\kappa\in\Z^3$.
The map $v\mapsto\kappa_v$ is additive.
It is injective, because the forms $\ell^1,\ell^2,\ell^3$ are linearly independent, their normals being pairwise orthogonal.
For a subgroup $\Gamma$ of $\Z^3$ with $\ell^a(\Gamma)\subset h\Z$, $a=1,2,3$, the image $\kappa_\Gamma:=\{\kappa_\gamma:\gamma\in\Gamma\}$ is a subgroup of $\Z^3$.

\emph{Even cuts.}
The cell structure $\mathcal{Q}$ \emph{cuts the torus $\mathbb{T}=(\Gamma,\chi)$ evenly} if $\ell^a(\Gamma)\subset 2h\Z$ for $a=1,2,3$;
then $\mathcal{Q}$ cuts $\mathbb{T}$.
For a torus $\mathbb{T}(\mathfrak{F};p_1,p_2,p_3)$ that carries $\mathcal{Q}$, the form $\ell^a$ takes at the generator $\left(p_b/|n_b|^2\right)n_b$ of its lattice the value $p_a\cdot\delta_{ab}$,
which lies in $2h\Z$ because $p_a\in2h\Z$ ($a,b=1,2,3$);
so $\mathcal{Q}$ cuts such a torus evenly.
If $\mathcal{Q}$ cuts $\mathbb{T}$, then $\tau_\gamma Q_\kappa=Q_{\kappa+\kappa_\gamma}$ for $\gamma\in\Gamma$.
Hence the cells of $\mathbb{T}$, that is, the $\Gamma$-orbits of cells, correspond to the cosets of $\kappa_\Gamma$ in $\Z^3$,
and a set of representatives of the cells of $\mathbb{T}$ is a subset of $\Z^3$ that meets each coset of $\kappa_\Gamma$ exactly once.

\emph{Families.}
Let $\mathcal{Q}$ cut $\mathbb{T}=(\Gamma,\chi)$.
A \emph{family on $(\mathbb{T},\mathcal{Q})$} is a collection $\mathfrak{G}=(G_\kappa)_{\kappa\in\Z^3}$ with $G_\kappa\in\mathcal{A}(Q_\kappa)$ for every $\kappa$ and
\[
G_{\kappa+\kappa_\gamma}=\left(G_\kappa\right)^{\tau_\gamma}\qquad\text{for all }\kappa\in\Z^3\text{ and }\gamma\in\Gamma ;
\]
the right side lies in $\mathcal{A}(Q_{\kappa+\kappa_\gamma})$ by Lemma~\ref{lem:admissible-images}\eqref{it:images-boxes}.
Let $\mathcal{K}$ be a set of representatives of the cells of $\mathbb{T}$, and let $G_\kappa\in\mathcal{A}(Q_\kappa)$ be given for $\kappa\in\mathcal{K}$.
Every point of $\Z^3$ is of the form $\kappa+\kappa_\gamma$ for exactly one pair $(\kappa,\gamma)\in\mathcal{K}\times\Gamma$,
since $\mathcal{K}$ meets each coset of $\kappa_\Gamma$ exactly once and $\gamma\mapsto\kappa_\gamma$ is injective.
A family with the given entries at the points of $\mathcal{K}$ must therefore satisfy $G_{\kappa+\kappa_\gamma}=(G_\kappa)^{\tau_\gamma}$;
conversely, this formula defines a family, because $(G^{\tau_v})^{\tau_{v'}}=G^{\tau_{v+v'}}$ by \eqref{eq:setting-group-action}.
Hence exactly one family on $(\mathbb{T},\mathcal{Q})$ has the given entries at the points of $\mathcal{K}$.

Let now $\mathcal{Q}$ cut $\mathbb{T}$ evenly, let $\mathfrak{G}$ be a family on $(\mathbb{T},\mathcal{Q})$, and let $\mathcal{K}$ be a set of representatives of the cells of $\mathbb{T}$.
Expanded into products, $\prod_{\kappa\in\mathcal{K}}G_\kappa$ is an element of $\mathcal{A}(\R^3)$ that embeds in $\mathbb{T}$ (Lemma~\ref{lem:embedding-criteria}\eqref{it:embedding-cells}),
so that $Z_{\mathbb{T}}\left[\prod_{\kappa\in\mathcal{K}}G_\kappa\right]$ is defined (Definition~\ref{def:functional}).
By Lemma~\ref{lem:admissible-images}\eqref{it:images-products}, the factors of this product may be taken in any order.
Any other set of representatives of the cells of $\mathbb{T}$ consists of points $\kappa+\kappa_\gamma$, one for each $\kappa\in\mathcal{K}$, with $\gamma\in\Gamma$ depending on $\kappa$;
in the product over this set the factor $G_\kappa$ is replaced by $G_{\kappa+\kappa_\gamma}=(G_\kappa)^{\tau_\gamma}$, and $(G_\kappa)^{\tau_\gamma}$ and $G_\kappa$ define the same function on $\Conf(\Gamma)$ (Lemma~\ref{lem:admissible-images}\eqref{it:images-periods}).
Therefore the function on $\Conf(\Gamma)$ given by the product, and with it the number $Z_{\mathbb{T}}\left[\prod_{\kappa\in\mathcal{K}}G_\kappa\right]$, is the same for every choice of $\mathcal{K}$.
This number is also written $Z_{\mathbb{T}}\left[\prod_\kappa G_\kappa\right]$, the product being taken over the cells $\kappa$ of $\mathbb{T}$.
Since $Q_\kappa\subset\R^3$, a content of the cell $Q_\kappa$ is an element of $\mathcal{A}(\R^3)$,
and Lemma~\ref{lem:embedding-criteria}\eqref{it:embedding-cells}, applied to the single cell $Q_\kappa$, shows that it embeds in each torus that $\mathcal{Q}$ cuts evenly.

\emph{The groups of the walls.}
For $a\in\{1,2,3\}$ let $W_a(\mathcal{Q})$ be the group generated by the reflections $\theta_{\ell^a,o_a+jh}$, $j\in\Z$.
The positions $o_a+jh$ lie in $\Adm(\ell^a)$ because $h$ is an integer, and $W_a(\mathcal{Q})$ is the group $W(\ell^a;o_a,h)$ of Definition~\ref{def:slab-copies}.
The set $\{o_a+jh<\ell^a<o_a+(j+1)h\}$ is called the \emph{layer $j$ of $\ell^a$}.
For $i,j\in\Z$ let $g^a_{i\leftarrow j}$ be the element of $W_a(\mathcal{Q})$ that maps the layer $j$ of $\ell^a$ onto the layer $i$;
these layers are the slabs $j+1$ and $i+1$ of Definition~\ref{def:slab-copies} for $c_0=o_a$ and $h'=h$,
so that $g^a_{i\leftarrow j}=g_{i+1\leftarrow j+1}$ exists and is unique.
For $\kappa\in\Z^3$, $a\in\{1,2,3\}$ and $i\in\Z$ let $\kappa[a{:}i]$ be the point $\kappa$ with its $a$-th entry replaced by $i$.

Let $b\neq a$.
Then $\ell^b(n_a)=n_b\cdot n_a=0$, hence $\ell^b\circ\theta_{\ell^a,c}=\ell^b$ for every $c$, and
\[
\theta_{\ell^a,c}\,\theta_{\ell^b,c'}\,x=x-\frac{2}{|n_a|^2}\left(\ell^a(x)-c\right)n_a-\frac{2}{|n_b|^2}\left(\ell^b(x)-c'\right)n_b=\theta_{\ell^b,c'}\,\theta_{\ell^a,c}\,x
\]
for all $x\in\R^3$ and all $c,c'$.
Consequently every element of $W_a(\mathcal{Q})$ leaves $\ell^b$ invariant for $b\neq a$,
the element $g^a_{i\leftarrow j}$ maps each cell $Q_\kappa$ with $\kappa_a=j$ onto $Q_{\kappa[a{:}i]}$,
and the groups $W_1(\mathcal{Q})$, $W_2(\mathcal{Q})$, $W_3(\mathcal{Q})$ commute with each other.
Their product $W_1(\mathcal{Q})\,W_2(\mathcal{Q})\,W_3(\mathcal{Q})$ is therefore a subgroup of $W(\mathcal{Q})$;
it contains the reflections in all walls, so that $W(\mathcal{Q})=W_1(\mathcal{Q})\,W_2(\mathcal{Q})\,W_3(\mathcal{Q})$.

The group $W(\mathcal{Q})$ acts simply transitively on the cells.
It acts transitively, since $g^3_{\kappa'_3\leftarrow\kappa_3}\,g^2_{\kappa'_2\leftarrow\kappa_2}\,g^1_{\kappa'_1\leftarrow\kappa_1}$ maps $Q_\kappa$ onto $Q_{\kappa'}$, each factor replacing one entry of the index.
Let, conversely, $g\in W(\mathcal{Q})$ satisfy $gQ_\kappa=Q_\kappa$; write $g=g^1g^2g^3$ with $g^b\in W_b(\mathcal{Q})$, and let $a\in\{1,2,3\}$.
The factors $g^b$ with $b\neq a$ commute with $g^a$ and leave $\ell^a$ invariant; therefore $\ell^a\circ g=\ell^a\circ g^a$,
and $g^a$ maps the points of $Q_\kappa$ into the layer $\kappa_a$ of $\ell^a$.
Now $g^a$ maps this layer onto a layer, which then meets the layer $\kappa_a$ and hence coincides with it;
so $g^a$ maps the layer $\kappa_a$ onto itself, and $g^a=1$ by the uniqueness in Definition~\ref{def:slab-copies}.
Thus $g=1$, and the element of $W(\mathcal{Q})$ that maps $Q_\kappa$ onto $Q_{\kappa'}$ is
\[
g^3_{\kappa'_3\leftarrow\kappa_3}\,g^2_{\kappa'_2\leftarrow\kappa_2}\,g^1_{\kappa'_1\leftarrow\kappa_1}.
\]

By Definition~\ref{def:mirrors}\eqref{it:mirror-product}, $\theta_{\ell^a,o_a+h}\,\theta_{\ell^a,o_a}$ is the translation by $\left(2h/|n_a|^2\right)n_a$, which therefore belongs to $W_a(\mathcal{Q})$.
If $v\in\Z^3$ satisfies $\ell^a(v)\in 2h\Z$ for $a=1,2,3$, then
\[
v=\sum_{a=1}^{3}\frac{\ell^a(v)}{|n_a|^2}\,n_a=\sum_{a=1}^{3}\frac{\ell^a(v)}{2h}\cdot\frac{2h}{|n_a|^2}\,n_a
\]
is an integer combination of these three vectors, and $\tau_v\in W(\mathcal{Q})$.
In particular $\tau_\gamma\in W(\mathcal{Q})$ for every $\gamma$ in the lattice of a torus that $\mathcal{Q}$ cuts evenly.
If $\mathbb{T}=(\Gamma,\chi)$ carries $\mathcal{Q}$, the translations $\tau_\gamma$, $\gamma\in\Gamma$, form a normal subgroup of $W(\mathcal{Q})$:
for a reflection $\theta$ in a wall of $\mathcal{Q}$ one has $\theta\tau_\gamma\theta^{-1}=\tau_{\vec\theta\gamma}$,
and the linear part $\vec\theta$ maps each generator $\left(p_b/|n_b|^2\right)n_b$ of $\Gamma$ to itself or to its negative.
The group $W(\mathcal{Q})/\Gamma$ then acts simply transitively on the cells of $\mathbb{T}$:
transitively because $W(\mathcal{Q})$ acts transitively on the cells,
and freely because $gQ_\kappa=\tau_\gamma Q_\kappa$ with $g\in W(\mathcal{Q})$ and $\gamma\in\Gamma$ implies $\tau_{-\gamma}g=1$.

\emph{Folds of a family.}
Let $\mathfrak{G}=(G_\kappa)$ be a family on $(\mathbb{T},\mathcal{Q})$, $a\in\{1,2,3\}$ and $j\in\Z$.
The \emph{fold of $\mathfrak{G}$ in the direction $a$ from the layer $j$} is the collection $\mathfrak{R}^j_a\mathfrak{G}$ with the entries
\[
\left(\mathfrak{R}^j_a\mathfrak{G}\right)_\kappa:=\left(G_{\kappa[a{:}j]}\right)^{g^a_{\kappa_a\leftarrow j}}\qquad(\kappa\in\Z^3).
\]
Since $g^a_{\kappa_a\leftarrow j}$ maps $Q_{\kappa[a{:}j]}$ onto $Q_\kappa$, the entry with the index $\kappa$ lies in $\mathcal{A}(Q_\kappa)$ (Lemma~\ref{lem:admissible-images}\eqref{it:images-boxes}).

\emph{Chessboard families.}
For $G\in\mathcal{A}(Q_\kappa)$ the \emph{chessboard family} of $G$ is $\mathfrak{C}(G):=\left(G^{[\kappa']}\right)_{\kappa'\in\Z^3}$,
where $G^{[\kappa']}$ is the mirror copy of $G$ in the cell $Q_{\kappa'}$ (Definition~\ref{def:folds});
it lies in $\mathcal{A}(Q_{\kappa'})$ by Lemma~\ref{lem:admissible-images}\eqref{it:images-boxes}.
\end{definition}

\begin{lemma}[Straightening in one direction]\label{lem:one-direction}
Let $\mathcal{Q}=(\mathfrak{F};o_1,o_2,o_3;h)$ be a cell structure that cuts the torus $\mathbb{T}=(\Gamma,\chi)$ evenly,
and let $a\in\{1,2,3\}$ be such that the straightened torus $\mathbb{T}^0=\mathbb{T}^0(\ell^a)$ is defined (Definition~\ref{def:slabs}).
Put $p:=p(\mathbb{T},\ell^a)$ and $\Gamma_\parallel:=\Gamma\cap\{\ell^a=0\}$, and let $\mathfrak{G}=(G_\kappa)$ be a family on $(\mathbb{T},\mathcal{Q})$.
\begin{enumerate}
\item\label{it:straight-even}
$p\in 2h\Z$, and $\mathcal{Q}$ cuts $\mathbb{T}^0$ evenly.
\item\label{it:straight-representatives}
Let $\mathcal{K}_\parallel$ be a set of representatives of $\{\kappa\in\Z^3:\kappa_a=0\}$ modulo $\kappa_{\Gamma_\parallel}$.
Then
\[
\mathcal{K}:=\left\{\kappa\in\Z^3:\ 0\le\kappa_a<\frac{p}{h},\ \kappa[a{:}0]\in\mathcal{K}_\parallel\right\}
\]
is a set of representatives of the cells of $\mathbb{T}$ and of the cells of $\mathbb{T}^0$.
\item\label{it:straight-family}
For every $j\in\Z$, $\mathfrak{R}^j_a\mathfrak{G}$ is a family on $(\mathbb{T}^0,\mathcal{Q})$.
\item\label{it:straight-inequality}
Let $\mathcal{K}$ be as in \eqref{it:straight-representatives}.
For every $j\in\Z$
\[
Z_{\mathbb{T}^0}\left[\prod_{\kappa\in\mathcal{K}}\left(\mathfrak{R}^j_a\mathfrak{G}\right)_\kappa\right]\ge 0,
\]
and
\[
\left|Z_{\mathbb{T}}\left[\prod_{\kappa\in\mathcal{K}}G_\kappa\right]\right|\le\prod_{0\le j<p/h}Z_{\mathbb{T}^0}\left[\prod_{\kappa\in\mathcal{K}}\left(\mathfrak{R}^j_a\mathfrak{G}\right)_\kappa\right]^{\frac{h}{p}} .
\]
\end{enumerate}
\end{lemma}

\begin{proof}
\emph{Item~\eqref{it:straight-even}.}
By Lemma~\ref{lem:closing-vectors}, $\ell^a(\Gamma)=p\Z$, and $\ell^a(\Gamma)\subset 2h\Z$ by hypothesis; hence $p\in 2h\Z$.
Write $p/h=2k$ with a positive integer $k$.
The lattice of $\mathbb{T}^0$ is $\Gamma^0=\Gamma_\parallel\oplus\Z g^0$ with $g^0:=\left(p/|n_a|^2\right)n_a$ (Definition~\ref{def:slabs}).
One has $\ell^a(\Gamma^0)=p\Z\subset 2h\Z$.
For $b\neq a$ one has $\ell^b(g^0)=0$, whence $\ell^b(\Gamma^0)=\ell^b(\Gamma_\parallel)\subset\ell^b(\Gamma)\subset 2h\Z$.

\emph{Item~\eqref{it:straight-representatives}.}
Fix $j_0\in\Z$, and let $\mathcal{K}'$ be the set of the $\kappa\in\Z^3$ with $j_0\le\kappa_a<j_0+2k$ and $\kappa[a{:}0]\in\mathcal{K}_\parallel$;
for $j_0=0$ this is $\mathcal{K}$.
It will be shown that $\mathcal{K}'$ represents the cells of $\mathbb{T}$ and those of $\mathbb{T}^0$.
Let $g$ be a closing vector of $\mathbb{T}$, so that $\Gamma=\Gamma_\parallel\oplus\Z g$ and $\ell^a(g)=p$ (Definition~\ref{def:slabs}, Lemma~\ref{lem:closing-vectors}).
The $a$-th entry of $\kappa_g$ is $2k$, and that of $\kappa_\gamma$ vanishes for $\gamma\in\Gamma_\parallel$.
Let $\kappa\in\Z^3$.
There is exactly one $i\in\Z$ with $j_0\le\kappa_a-2k\cdot i<j_0+2k$, and the point $\kappa':=\kappa-i\cdot\kappa_g$ has its $a$-th entry in the window $[j_0,j_0+2k)$.
There is exactly one $\gamma_\parallel\in\Gamma_\parallel$ with $\kappa'[a{:}0]+\kappa_{\gamma_\parallel}\in\mathcal{K}_\parallel$,
and since the $a$-th entry of $\kappa_{\gamma_\parallel}$ is $0$, the point $\kappa'+\kappa_{\gamma_\parallel}=\kappa+\kappa_{\gamma_\parallel-ig}$ belongs to $\mathcal{K}'$.
Let, conversely, $\kappa$ and $\kappa+\kappa_\gamma$ both belong to $\mathcal{K}'$, where $\gamma=\gamma_\parallel+ig$ with $\gamma_\parallel\in\Gamma_\parallel$ and $i\in\Z$.
Their $a$-th entries differ by $2k\cdot i$ and lie in an interval of length $2k$; so $i=0$.
Then $\kappa[a{:}0]$ and $\kappa[a{:}0]+\kappa_{\gamma_\parallel}$ both lie in $\mathcal{K}_\parallel$, so that $\kappa_{\gamma_\parallel}=0$, and $\gamma_\parallel=0$ by injectivity.
Thus $\mathcal{K}'$ meets each coset of $\kappa_\Gamma$ exactly once, that is, it represents the cells of $\mathbb{T}$ (Definition~\ref{def:cell-families}).
By \eqref{it:straight-even}, the cells of $\mathbb{T}^0$ correspond to the cosets of $\kappa_{\Gamma^0}$.
By Lemma~\ref{lem:closing-vectors}, $\Gamma^0\cap\{\ell^a=0\}=\Gamma_\parallel$, and $g^0$ is a closing vector of $\mathbb{T}^0$ with $\ell^a(g^0)=p$.
The same argument, with $g^0$ in place of $g$, shows that $\mathcal{K}'$ represents the cells of $\mathbb{T}^0$.

\emph{Item~\eqref{it:straight-family}.}
Let $j\in\Z$.
The entries of $\mathfrak{R}^j_a\mathfrak{G}$ satisfy $\left(\mathfrak{R}^j_a\mathfrak{G}\right)_\kappa\in\mathcal{A}(Q_\kappa)$, and $\mathcal{Q}$ cuts $\mathbb{T}^0$ by \eqref{it:straight-even}.
By \eqref{eq:setting-group-action}, the vectors $v\in\Gamma^0$ with $\left(\mathfrak{R}^j_a\mathfrak{G}\right)_{\kappa+\kappa_v}=\left(\left(\mathfrak{R}^j_a\mathfrak{G}\right)_\kappa\right)^{\tau_v}$ for all $\kappa\in\Z^3$ form a subgroup of $\Gamma^0$.
As $\Gamma^0$ is generated by $\Gamma_\parallel$ and $g^0$, it suffices to treat $v\in\Gamma_\parallel$ and $v=g^0$.

Let $\gamma\in\Gamma_\parallel$.
By Definition~\ref{def:mirrors}\eqref{it:mirror-conjugation}, $\tau_\gamma\theta_{\ell^a,c}\tau_{-\gamma}=\theta_{\ell^a,c+\ell^a(\gamma)}=\theta_{\ell^a,c}$;
so $\tau_\gamma$ commutes with the generators of $W_a(\mathcal{Q})$, hence with all its elements.
The $a$-th entry of $\kappa_\gamma$ vanishes, so that $(\kappa+\kappa_\gamma)[a{:}j]=\kappa[a{:}j]+\kappa_\gamma$ and $(\kappa+\kappa_\gamma)_a=\kappa_a$.
The periodicity of $\mathfrak{G}$ and the rule $(G^{g_1})^{g_2}=G^{g_2g_1}$ of \eqref{eq:setting-group-action} give
\begin{align*}
\left(\mathfrak{R}^j_a\mathfrak{G}\right)_{\kappa+\kappa_\gamma}
&=\left(G_{\kappa[a{:}j]+\kappa_\gamma}\right)^{g^a_{\kappa_a\leftarrow j}}
=\left(\left(G_{\kappa[a{:}j]}\right)^{\tau_\gamma}\right)^{g^a_{\kappa_a\leftarrow j}}
=\left(G_{\kappa[a{:}j]}\right)^{g^a_{\kappa_a\leftarrow j}\tau_\gamma}\\
&=\left(G_{\kappa[a{:}j]}\right)^{\tau_\gamma g^a_{\kappa_a\leftarrow j}}
=\left(\left(G_{\kappa[a{:}j]}\right)^{g^a_{\kappa_a\leftarrow j}}\right)^{\tau_\gamma}
=\left(\left(\mathfrak{R}^j_a\mathfrak{G}\right)_\kappa\right)^{\tau_\gamma}.
\end{align*}

Now let $v_h:=\left(2h/|n_a|^2\right)n_a$.
Then $g^0=k\cdot v_h$, and $\tau_{v_h}\in W_a(\mathcal{Q})$ (Definition~\ref{def:slab-copies}); hence $\tau_{g^0}\in W_a(\mathcal{Q})$.
Moreover $\kappa_{g^0}=2k\cdot e_a$, because $\ell^b(g^0)=p\cdot\delta_{ab}$.
Since $\ell^a(x+g^0)=\ell^a(x)+p$, the element $\tau_{g^0}g^a_{\kappa_a\leftarrow j}$ of $W_a(\mathcal{Q})$ maps the layer $j$ of $\ell^a$ onto the layer $\kappa_a+2k$;
by uniqueness, $\tau_{g^0}g^a_{\kappa_a\leftarrow j}=g^a_{\kappa_a+2k\leftarrow j}$.
As the $a$-th entry of $\kappa+\kappa_{g^0}$ is $\kappa_a+2k$ and $\left(\kappa+\kappa_{g^0}\right)[a{:}j]=\kappa[a{:}j]$, it follows that
\[
\left(\mathfrak{R}^j_a\mathfrak{G}\right)_{\kappa+\kappa_{g^0}}
=\left(G_{\kappa[a{:}j]}\right)^{g^a_{\kappa_a+2k\leftarrow j}}
=\left(G_{\kappa[a{:}j]}\right)^{\tau_{g^0}g^a_{\kappa_a\leftarrow j}}
=\left(\left(\mathfrak{R}^j_a\mathfrak{G}\right)_\kappa\right)^{\tau_{g^0}} .
\]

\emph{Item~\eqref{it:straight-inequality}.}
Let $j_0\in\Z$ and $\mathcal{K}'$ be as in the proof of \eqref{it:straight-representatives}, and put $c_i:=o_a+(j_0+i)h$ for $i\in\Z$;
these positions lie in $\Adm(\ell^a)$.
Let $\Xi:=\Xi(\mathbb{T},\ell^a)$ be the cross-section formed by $\Gamma_\parallel$ and the restriction of $\chi$.
For $i=1,\dots,2k$ put
\[
S_i:=\prod_{\kappa\in\mathcal{K}',\ \kappa_a=j_0+i-1}G_\kappa .
\]
The cells $Q_\kappa$ with $\kappa_a=j_0+i-1$ lie in the open slab $\{c_{i-1}<\ell^a<c_i\}$, and those with $\kappa\in\mathcal{K}'$ belong to pairwise different $\Gamma$-orbits.
Expanded into products, $S_i$ is therefore a sum of products of $\mathcal{A}(\{c_{i-1}<\ell^a<c_i\})$ that embed in $\Gamma$, hence in $\Gamma_\parallel$ (Lemma~\ref{lem:embedding-criteria}\eqref{it:embedding-cells}).
So $S_i\in\mathcal{A}_\Xi[c_{i-1},c_i]$ (Definition~\ref{def:slabs}), and $\prod_{\kappa\in\mathcal{K}'}G_\kappa=S_1\cdots S_{2k}$.

For the planes $c_i$ and $h'=h$, the group of Definition~\ref{def:slab-copies} is $W_a(\mathcal{Q})$,
and the element that maps the slab $\{c_{i-1}<\ell^a<c_i\}$ onto the slab $\{c_{i'-1}<\ell^a<c_{i'}\}$ is $g^a_{j_0+i'-1\leftarrow j_0+i-1}$.
By Lemma~\ref{lem:admissible-images}\eqref{it:images-products},
\[
\fold_{\ell^a}S_i=\prod_{i'=1}^{2k}\ \prod_{\kappa\in\mathcal{K}',\ \kappa_a=j_0+i-1}\left(G_\kappa\right)^{g^a_{j_0+i'-1\leftarrow j_0+i-1}} .
\]
For $\kappa\in\mathcal{K}'$ with $\kappa_a=j_0+i-1$ one has $\kappa=\kappa[a{:}j_0+i'-1][a{:}j_0+i-1]$;
so the factor with the index $\kappa$ is the entry of $\mathfrak{R}^{j_0+i-1}_a\mathfrak{G}$ with the index $\kappa[a{:}j_0+i'-1]$.
Membership in $\mathcal{K}'$ depends only on the $a$-th entry and on $\kappa[a{:}0]$;
hence $\kappa\mapsto\kappa[a{:}j_0+i'-1]$ maps $\{\kappa\in\mathcal{K}':\kappa_a=j_0+i-1\}$ bijectively onto $\{\kappa\in\mathcal{K}':\kappa_a=j_0+i'-1\}$.
Taking the product over $i'=1,\dots,2k$, one obtains
\begin{equation}\label{eq:tori-layer-fold}
\fold_{\ell^a}S_i=\prod_{\kappa\in\mathcal{K}'}\left(\mathfrak{R}^{j_0+i-1}_a\mathfrak{G}\right)_\kappa\qquad(i=1,\dots,2k),
\end{equation}
the product of the family $\mathfrak{R}^{j_0+i-1}_a\mathfrak{G}$ on $(\mathbb{T}^0,\mathcal{Q})$ over the set $\mathcal{K}'$ of representatives of the cells of $\mathbb{T}^0$.

Proposition~\ref{prop:twisted-holder} applies to the torus $\mathbb{T}$, the family $\ell^a$, a closing vector of $\mathbb{T}$, $h'=h$, the planes $c_i$ and the slab contents $S_1,\dots,S_{2k}$.
By \eqref{eq:tori-fold-norm} and \eqref{eq:tori-holder},
\[
Z_{\mathbb{T}^0}\left[\fold_{\ell^a}S_i\right]\ge0\quad(i=1,\dots,2k),\qquad
\left|Z_{\mathbb{T}}\left[S_1\cdots S_{2k}\right]\right|\le\prod_{i=1}^{2k}Z_{\mathbb{T}^0}\left[\fold_{\ell^a}S_i\right]^{\frac{1}{2k}} .
\]
For $i=1$, \eqref{eq:tori-layer-fold} turns the first of these relations into $Z_{\mathbb{T}^0}\left[\prod_{\kappa\in\mathcal{K}'}\left(\mathfrak{R}^{j_0}_a\mathfrak{G}\right)_\kappa\right]\ge0$.
Both $\mathcal{K}$ and $\mathcal{K}'$ represent the cells of $\mathbb{T}^0$, and $\mathfrak{R}^{j_0}_a\mathfrak{G}$ is a family on $(\mathbb{T}^0,\mathcal{Q})$;
so its products over $\mathcal{K}$ and over $\mathcal{K}'$ define the same function on $\Conf(\Gamma^0)$ (Definition~\ref{def:cell-families}).
As $j_0\in\Z$ is arbitrary, this is the first assertion of \eqref{it:straight-inequality}.
For $j_0=0$ one has $\mathcal{K}'=\mathcal{K}$ and $\prod_{\kappa\in\mathcal{K}}G_\kappa=S_1\cdots S_{2k}$;
by \eqref{eq:tori-layer-fold} with $j=i-1$, and since $\frac{1}{2k}=\frac{h}{p}$, the inequality above is the second assertion of \eqref{it:straight-inequality}.
\end{proof}

\begin{proposition}[Straightening in three directions]\label{prop:three-directions}
Let $\mathcal{Q}=(\mathfrak{F};o_1,o_2,o_3;h)$ be a cell structure that cuts the torus $\mathbb{T}_0$ evenly,
and assume that the straightened tori
\[
\mathbb{T}_1:=\mathbb{T}_0^0(\ell^1),\qquad\mathbb{T}_2:=\mathbb{T}_1^0(\ell^2),\qquad\mathbb{T}_3:=\mathbb{T}_2^0(\ell^3)
\]
are defined.
Put $p_1:=p(\mathbb{T}_0,\ell^1)$, $p_2:=p(\mathbb{T}_1,\ell^2)$ and $p_3:=p(\mathbb{T}_2,\ell^3)$.
Then $\mathbb{T}_3=\mathbb{T}(\mathfrak{F};p_1,p_2,p_3)$, and $\mathbb{T}_3$ carries $\mathcal{Q}$.
The brick
\[
\mathcal{K}_p:=\left\{\kappa\in\Z^3:\ 0\le\kappa_a<\frac{p_a}{h},\ a=1,2,3\right\}
\]
is a set of representatives of the cells of each of the tori $\mathbb{T}_0$, $\mathbb{T}_1$, $\mathbb{T}_2$, $\mathbb{T}_3$, and
\[
\#|\mathcal{K}_p|=\frac{p_1p_2p_3}{h^3}=\frac{|\mathbb{T}_i|}{|Q|}\qquad(i=0,1,2,3).
\]
For every family $\mathfrak{G}=(G_\kappa)$ on $(\mathbb{T}_0,\mathcal{Q})$
\begin{align*}
\left|Z_{\mathbb{T}_0}\left[\prod_{\kappa\in\mathcal{K}_p}G_\kappa\right]\right|
&\le\prod_{j_1}Z_{\mathbb{T}_1}\left[\prod_{\kappa\in\mathcal{K}_p}\left(\mathfrak{R}^{j_1}_1\mathfrak{G}\right)_\kappa\right]^{\frac{h}{p_1}}\\
&\le\prod_{j_1,j_2}Z_{\mathbb{T}_2}\left[\prod_{\kappa\in\mathcal{K}_p}\left(\mathfrak{R}^{j_2}_2\mathfrak{R}^{j_1}_1\mathfrak{G}\right)_\kappa\right]^{\frac{h^2}{p_1p_2}}\\
&\le\prod_{\kappa\in\mathcal{K}_p}\|G_\kappa\|_{\mathbb{T}_3},
\end{align*}
where $j_a$ ranges over the integers with $0\le j_a<p_a/h$; all the factors in this chain are non-negative.
The same assertions hold if the three directions are taken in another order,
the tori $\mathbb{T}_1,\mathbb{T}_2,\mathbb{T}_3$ being obtained by straightening in that order,
$p_a$ being the height period of $\ell^a$ on the torus that is straightened in the direction $\ell^a$,
and the folds $\mathfrak{R}^{j_a}_a$ being applied in that order.
\end{proposition}

\begin{proof}
\emph{Lattices.}
For $i=0,1,2,3$ let $\Gamma_i$ be the lattice and $\chi_i$ the character of $\mathbb{T}_i$,
and for $a=1,2,3$ let $\gamma_a:=\left(p_a/|n_a|^2\right)n_a$, so that $\ell^b(\gamma_a)=p_a\cdot\delta_{ab}$.
By Definition~\ref{def:slabs}, $\Gamma_1=(\Gamma_0\cap\{\ell^1=0\})\oplus\Z\gamma_1$.
Since $\ell^2(\gamma_1)=0$, this gives $\Gamma_1\cap\{\ell^2=0\}=(\Gamma_0\cap\{\ell^1=\ell^2=0\})\oplus\Z\gamma_1$.
The group $\Gamma_0\cap\{\ell^1=\ell^2=0\}$ lies on the line $\R n_3$; it is the kernel of the homomorphism $(\ell^1,\ell^2)\colon\Gamma_0\to\Z^2$.
As $\Gamma_0$ spans $\R^3$ and the forms $\ell^1$, $\ell^2$ are linearly independent, the image of this homomorphism has rank two, and the kernel has rank $3-2=1$.
So $\Gamma_0\cap\{\ell^1=\ell^2=0\}=\Z\gamma'$ with a non-zero multiple $\gamma'$ of $n_3$, whose sign is fixed by the requirement $\ell^3(\gamma')>0$.
Hence $\Gamma_2=(\Gamma_1\cap\{\ell^2=0\})\oplus\Z\gamma_2=\Z\gamma'\oplus\Z\gamma_1\oplus\Z\gamma_2$,
so that $\Gamma_2\cap\{\ell^3=0\}=\Z\gamma_1\oplus\Z\gamma_2$ and $\ell^3(\Gamma_2)=\ell^3(\gamma')\Z$.
By Lemma~\ref{lem:closing-vectors}, $p_3=\ell^3(\gamma')$, and since $\gamma'$ is a multiple of $n_3$, $\gamma'=\gamma_3$.
Consequently
\[
\Gamma_3=(\Gamma_2\cap\{\ell^3=0\})\oplus\Z\gamma_3=\Z\gamma_1\oplus\Z\gamma_2\oplus\Z\gamma_3=\Gamma_2 .
\]

\emph{Characters.}
By Definition~\ref{def:slabs}, $\chi_1(\gamma_1)=\chi_2(\gamma_2)=\chi_3(\gamma_3)=-1$,
while $\chi_2$ agrees with $\chi_1$ on $\Gamma_1\cap\{\ell^2=0\}$ and $\chi_3$ agrees with $\chi_2$ on $\Gamma_2\cap\{\ell^3=0\}=\Z\gamma_1\oplus\Z\gamma_2$.
The group $\Gamma_1\cap\{\ell^2=0\}$ contains $\gamma_1$, so that $\chi_2(\gamma_1)=\chi_1(\gamma_1)=-1$;
consequently $\chi_3(\gamma_1)=\chi_2(\gamma_1)=-1$ and $\chi_3(\gamma_2)=\chi_2(\gamma_2)=-1$.
Thus $\chi_3(\gamma_a)=-1$ for $a=1,2,3$.

Lemma~\ref{lem:one-direction}\eqref{it:straight-even}, applied to $(\mathbb{T}_0,\ell^1)$, $(\mathbb{T}_1,\ell^2)$ and $(\mathbb{T}_2,\ell^3)$ in turn,
shows that $\mathcal{Q}$ cuts $\mathbb{T}_1$, $\mathbb{T}_2$, $\mathbb{T}_3$ evenly and that $p_1,p_2,p_3\in 2h\Z$.
As height periods, $p_1,p_2,p_3$ are positive (Definition~\ref{def:slabs}); they are thus positive integers,
and they satisfy the two conditions of Definition~\ref{def:cell-structures}.
First, the vectors $\gamma_a=\left(p_a/|n_a|^2\right)n_a$ belong to $\Gamma_3\subset\Z^3$.
Secondly, $\mathbb{T}_3$ is defined, so that $\Gamma_3$ is a period lattice (Definition~\ref{def:slabs}),
and its non-zero elements $\gamma_a$ satisfy $|\gamma_a|_\infty\ge\sys(\Gamma_3)\ge4$ (Definition~\ref{def:configurations});
as the entries of $n_a$ lie in $\{0,1,-1\}$ and are not all zero (Definition~\ref{def:mirrors}), one has $|\gamma_a|_\infty=p_a/|n_a|^2$,
and hence $p_a\ge4\,|n_a|^2$ for $a=1,2,3$.
Consequently $\mathbb{T}_3=(\Gamma_3,\chi_3)$ is the rectangular torus $\mathbb{T}(\mathfrak{F};p_1,p_2,p_3)$,
its lattice being $\Z\gamma_1\oplus\Z\gamma_2\oplus\Z\gamma_3$ and its character taking the value $-1$ at the three generators;
and $\mathbb{T}_3$ carries $\mathcal{Q}$, because $p_1,p_2,p_3\in2h\Z$.

\emph{The brick.}
One has $\kappa_{\gamma_a}=\left(p_a/h\right)e_a$.
So $\kappa_{\Gamma_2}=\kappa_{\Gamma_3}$ consists of the $\kappa$ with $\kappa_a\in\left(p_a/h\right)\Z$ for $a=1,2,3$,
and $\mathcal{K}_p$ represents the cells of $\mathbb{T}_2$ and of $\mathbb{T}_3$.
As $\Gamma_2\cap\{\ell^3=0\}=\Z\gamma_1\oplus\Z\gamma_2$, the set $\mathcal{K}_p$ is also the set $\mathcal{K}$ of Lemma~\ref{lem:one-direction}\eqref{it:straight-representatives} for $(\mathbb{T}_2,\ell^3)$, formed with $\mathcal{K}_\parallel=\mathcal{K}_p\cap\{\kappa_3=0\}$.
Next consider $(\mathbb{T}_1,\ell^2)$.
Here $\Gamma_1\cap\{\ell^2=0\}=\Z\gamma_1\oplus\Z\gamma_3$, and $\mathcal{K}_p\cap\{\kappa_2=0\}$ represents $\{\kappa\in\Z^3:\kappa_2=0\}$ modulo $\kappa_{\Z\gamma_1\oplus\Z\gamma_3}$.
So $\mathcal{K}_p$ is the set $\mathcal{K}$ of Lemma~\ref{lem:one-direction}\eqref{it:straight-representatives} for $(\mathbb{T}_1,\ell^2)$, formed with $\mathcal{K}_\parallel=\mathcal{K}_p\cap\{\kappa_2=0\}$,
and it represents the cells of $\mathbb{T}_1$.
Finally consider $(\mathbb{T}_0,\ell^1)$, and let $\mathcal{K}_\parallel:=\mathcal{K}_p\cap\{\kappa_1=0\}$.
Two points of $\mathcal{K}_\parallel$ that differ by an element of $\kappa_{\Gamma_0\cap\{\ell^1=0\}}$ differ by an element of $\kappa_{\Gamma_1}$, and are therefore equal.
A point $\kappa$ with $\kappa_1=0$ can be written $\kappa=\kappa'-\kappa_\gamma-i\cdot\kappa_{\gamma_1}$ with $\kappa'\in\mathcal{K}_p$, $\gamma\in\Gamma_0\cap\{\ell^1=0\}$ and $i\in\Z$,
because $\mathcal{K}_p$ represents $\Z^3$ modulo $\kappa_{\Gamma_1}$ and $\Gamma_1=(\Gamma_0\cap\{\ell^1=0\})\oplus\Z\gamma_1$.
The first entries give $\kappa'_1=i\cdot p_1/h\in[0,p_1/h)$; hence $i=0$ and $\kappa'\in\mathcal{K}_\parallel$.
Thus $\mathcal{K}_\parallel$ represents $\{\kappa\in\Z^3:\kappa_1=0\}$ modulo $\kappa_{\Gamma_0\cap\{\ell^1=0\}}$,
$\mathcal{K}_p$ is the set $\mathcal{K}$ of Lemma~\ref{lem:one-direction}\eqref{it:straight-representatives} for $(\mathbb{T}_0,\ell^1)$,
and it represents the cells of $\mathbb{T}_0$.

\emph{Volumes.}
The normals $n_1,n_2,n_3$ being pairwise orthogonal,
\[
|\mathbb{T}_3|=\left|\det(\gamma_1,\gamma_2,\gamma_3)\right|=\frac{p_1p_2p_3}{|n_1|\cdot|n_2|\cdot|n_3|}=\#|\mathcal{K}_p|\cdot\frac{h^3}{|n_1|\cdot|n_2|\cdot|n_3|}=\#|\mathcal{K}_p|\cdot|Q|,
\]
with $|Q|$ as in Definition~\ref{def:cell-structures};
and $|\mathbb{T}_0|=|\mathbb{T}_1|=|\mathbb{T}_2|=|\mathbb{T}_3|$ by Lemma~\ref{lem:closing-vectors}.

\emph{Inequalities.}
Lemma~\ref{lem:one-direction}\eqref{it:straight-inequality} for $(\mathbb{T}_0,\ell^1,\mathfrak{G})$, with $\mathcal{K}=\mathcal{K}_p$, is the first inequality, and its right side is a product of non-negative factors.
For $0\le j_1<p_1/h$, the collection $\mathfrak{R}^{j_1}_1\mathfrak{G}$ is a family on $(\mathbb{T}_1,\mathcal{Q})$ by Lemma~\ref{lem:one-direction}\eqref{it:straight-family},
and the same lemma for $(\mathbb{T}_1,\ell^2,\mathfrak{R}^{j_1}_1\mathfrak{G})$ gives
\[
Z_{\mathbb{T}_1}\left[\prod_{\kappa\in\mathcal{K}_p}\left(\mathfrak{R}^{j_1}_1\mathfrak{G}\right)_\kappa\right]\le\prod_{j_2}Z_{\mathbb{T}_2}\left[\prod_{\kappa\in\mathcal{K}_p}\left(\mathfrak{R}^{j_2}_2\mathfrak{R}^{j_1}_1\mathfrak{G}\right)_\kappa\right]^{\frac{h}{p_2}},
\]
with non-negative factors on both sides.
Raising to the power $h/p_1$ and multiplying over $j_1$ yields the second inequality.
In the same way, the lemma for $(\mathbb{T}_2,\ell^3,\mathfrak{R}^{j_2}_2\mathfrak{R}^{j_1}_1\mathfrak{G})$ gives
\[
\prod_{j_1,j_2}Z_{\mathbb{T}_2}\left[\prod_{\kappa\in\mathcal{K}_p}\left(\mathfrak{R}^{j_2}_2\mathfrak{R}^{j_1}_1\mathfrak{G}\right)_\kappa\right]^{\frac{h^2}{p_1p_2}}
\le\prod_{j_1,j_2,j_3}Z_{\mathbb{T}_3}\left[\prod_{\kappa\in\mathcal{K}_p}\left(\mathfrak{R}^{j_3}_3\mathfrak{R}^{j_2}_2\mathfrak{R}^{j_1}_1\mathfrak{G}\right)_\kappa\right]^{\frac{h^3}{p_1p_2p_3}},
\]
where $0\le j_3<p_3/h$ and all factors on the right are non-negative.
For $\kappa\in\Z^3$ put $g^a:=g^a_{\kappa_a\leftarrow j_a}$, $a=1,2,3$.
By the definition of the folds and by \eqref{eq:setting-group-action},
\[
\left(\mathfrak{R}^{j_3}_3\mathfrak{R}^{j_2}_2\mathfrak{R}^{j_1}_1\mathfrak{G}\right)_\kappa
=\left(\left(\left(G_{(j_1,j_2,j_3)}\right)^{g^1}\right)^{g^2}\right)^{g^3}
=\left(G_{(j_1,j_2,j_3)}\right)^{g^3g^2g^1}
=\left(G_{(j_1,j_2,j_3)}\right)^{[\kappa]};
\]
the last equality is Definition~\ref{def:folds}, since $g^3g^2g^1$ is the element of $W(\mathcal{Q})$ that maps $Q_{(j_1,j_2,j_3)}$ onto $Q_\kappa$ (Definition~\ref{def:cell-families}).
Thus $\mathfrak{R}^{j_3}_3\mathfrak{R}^{j_2}_2\mathfrak{R}^{j_1}_1\mathfrak{G}=\mathfrak{C}(G_{(j_1,j_2,j_3)})$,
and its product over $\mathcal{K}_p$, a set of representatives of the cells of $\mathbb{T}_3$, is $\Fold_{\mathbb{T}_3}G_{(j_1,j_2,j_3)}$.
The value of $Z_{\mathbb{T}_3}$ on it is non-negative and equals $\|G_{(j_1,j_2,j_3)}\|_{\mathbb{T}_3}^{\#|\mathcal{K}_p|}$ (Definition~\ref{def:folds}),
while $\frac{h^3}{p_1p_2p_3}=\frac{1}{\#|\mathcal{K}_p|}$.
As $(j_1,j_2,j_3)$ ranges over $\mathcal{K}_p$, the right side of the last inequality is $\prod_{\kappa\in\mathcal{K}_p}\|G_\kappa\|_{\mathbb{T}_3}$; this is the third inequality.

\emph{Other orders.}
The argument depends on the frame only through the pairwise orthogonality of its normals,
and the cells, the group $W(\mathcal{Q})$, the mirror copies and the chessboard seminorms are defined without reference to an order of the directions.
With the indices permuted, the same argument therefore applies to every order of the three steps.
\end{proof}

The last torus $\mathbb{T}_3$ depends on the order of the steps.
Let $\mathfrak{F}=(\ell_-,\ell_+,x_3)$ be the diagonal frame and $\mathbb{T}_0=C(n,n,n)$ with an even integer $n\ge8$.
If $\ell_-$ is straightened before $\ell_+$, then $\mathbb{T}_3=D(n/2,n,n)$;
in the three other orders $\mathbb{T}_3=D(n,n/2,n)$, the image (Definition~\ref{def:tori}) of $D(n/2,n,n)$ under the reflection $x_2\mapsto-x_2$, an element of $\Ocub$.
Indeed, by Proposition~\ref{prop:untwistings}(a), (b), straightening $C(n,n,n)$ in the direction $\ell_-$ gives $D'(n/2,n,n)$, and straightening this torus in the direction $\ell_+$ gives $D(n/2,n,n)$;
the tori $C(n,n,n)$, $D'(n/2,n,n)$ and $D(n/2,n,n)$ are adapted to $x_3$ (Remark~\ref{rem:adapted-families}), so that a step in the direction $x_3$ leaves them unchanged;
and the reflection $x_2\mapsto-x_2$ maps $C(n,n,n)$ onto itself and exchanges $\ell_-$ and $\ell_+$.

\begin{theorem}[Chessboard estimates]\label{thm:chessboard}
Let $-1\le\rho\le1$, $w_0,\dots,w_3>0$, $M\in\R$ and $q>0$.
\begin{enumerate}
\item\label{it:chessboard-carried}
Let $\mathcal{Q}=(\mathfrak{F};o_1,o_2,o_3;h)$ be a cell structure, and let $\mathbb{T}$ be a torus that carries $\mathcal{Q}$.
Then
\[
Z_{\mathbb{T}}\left[\Fold_{\mathbb{T}}G\right]\ge0\qquad\text{for every }\kappa\in\Z^3\text{ and every }G\in\mathcal{A}(Q_\kappa).
\]
If $\mathcal{K}$ is a set of representatives of the cells of $\mathbb{T}$ and $G_\kappa\in\mathcal{A}(Q_\kappa)$ for $\kappa\in\mathcal{K}$, then
\[
\left|Z_{\mathbb{T}}\left[\prod_{\kappa\in\mathcal{K}}G_\kappa\right]\right|\le\prod_{\kappa\in\mathcal{K}}\|G_\kappa\|_{\mathbb{T}} .
\]
\item\label{it:chessboard-diagonal}
Let $\mathcal{Q}=(\mathfrak{F}_{\mathrm{d}};o_1,o_2,o_3;h)$ be a cell structure of the diagonal frame $\mathfrak{F}_{\mathrm{d}}=(\ell_-,\ell_+,x_3)$,
and let $n\in2h\Z$ satisfy $n\ge8$.
Then $\mathcal{Q}$ cuts the tori $C(n,n,n)$ and $D'(n/2,n,n)$ evenly, and $D(n/2,n,n)$ carries $\mathcal{Q}$.
The brick of Proposition~\ref{prop:three-directions} for $p=(n,2n,n)$,
\[
\mathcal{K}_p=\left\{\kappa\in\Z^3:\ 0\le\kappa_1<\frac{n}{h},\ 0\le\kappa_2<\frac{2n}{h},\ 0\le\kappa_3<\frac{n}{h}\right\},
\]
is a set of representatives of the cells of each of these three tori, and each of them has $\#|\mathcal{K}_p|=2n^3/h^3$ cells.
Let $G_\kappa\in\mathcal{A}(Q_\kappa)$ for $\kappa\in\mathcal{K}_p$.
For $j=1,\dots,n/h$ the product
\[
S_j:=\prod_{\kappa\in\mathcal{K}_p,\ \kappa_1=j-1}G_\kappa
\]
is the content of the $j$-th slab $\{o_1+(j-1)\cdot h<\ell_-<o_1+j\cdot h\}$ of $\ell_-$:
it belongs to $\mathcal{A}_\Xi\left[o_1+(j-1)\cdot h,\ o_1+j\cdot h\right]$, where $\Xi:=\Xi(C(n,n,n),\ell_-)$.
Let $\fold_{\ell_-}S_j$ be its fold on $D'(n/2,n,n)$ with respect to the planes $o_1+i\cdot h$, $i\in\Z$ (Definition~\ref{def:slab-copies}).
Then
\[
\left|Z_{C(n,n,n)}\left[\prod_{\kappa\in\mathcal{K}_p}G_\kappa\right]\right|
\le\prod_{j=1}^{n/h}Z_{D'(n/2,n,n)}\left[\fold_{\ell_-}S_j\right]^{\frac{h}{n}}
\le\prod_{\kappa\in\mathcal{K}_p}\|G_\kappa\|_{D(n/2,n,n)},
\]
and moreover
\[
\left|Z_{D'(n/2,n,n)}\left[\prod_{\kappa\in\mathcal{K}_p}G_\kappa\right]\right|
\le\prod_{\kappa\in\mathcal{K}_p}\|G_\kappa\|_{D(n/2,n,n)} .
\]
\item\label{it:chessboard-coordinate}
Let $\mathcal{Q}=(\mathfrak{F}_{\mathrm{c}};o_1,o_2,o_3;h'')$ be a cell structure of the coordinate frame $\mathfrak{F}_{\mathrm{c}}=(x_1,x_2,x_3)$,
and let $a,c\in2h''\Z$ satisfy $a\ge4$ and $c\ge4$.
Then $\mathcal{Q}$ cuts $D(a,2a,c)$ evenly, and $C(a,4a,c)$ carries $\mathcal{Q}$.
The brick of Proposition~\ref{prop:three-directions} for $p=(a,4a,c)$ and the width $h''$,
\[
\mathcal{K}_p=\left\{\kappa\in\Z^3:\ 0\le\kappa_1<\frac{a}{h''},\ 0\le\kappa_2<\frac{4a}{h''},\ 0\le\kappa_3<\frac{c}{h''}\right\},
\]
is a set of representatives of the cells of both tori, each of which has $\#|\mathcal{K}_p|=4a^2\cdot c/h''^3$ cells.
For $G_\kappa\in\mathcal{A}(Q_\kappa)$, $\kappa\in\mathcal{K}_p$, one has
\[
\left|Z_{D(a,2a,c)}\left[\prod_{\kappa\in\mathcal{K}_p}G_\kappa\right]\right|
\le\prod_{\kappa\in\mathcal{K}_p}\|G_\kappa\|_{C(a,4a,c)} .
\]
\end{enumerate}
Let $\mathbb{T}_0$ be the torus on the left-hand side of one of the inequalities in \eqref{it:chessboard-carried}--\eqref{it:chessboard-coordinate},
and let $\mathbb{T}_3$ be the torus whose chessboard seminorms occur on its right-hand side.
If $\gamma$ belongs to the period lattice of $\mathbb{T}_0$, then
\[
\mathfrak{C}\left(G^{\tau_\gamma}\right)=\mathfrak{C}(G)
\qquad\text{and}\qquad
\left\|G^{\tau_\gamma}\right\|_{\mathbb{T}_3}=\|G\|_{\mathbb{T}_3}
\]
for every $\kappa\in\Z^3$ and every $G\in\mathcal{A}(Q_\kappa)$.
Consequently, for a family $(G_\kappa)_{\kappa\in\Z^3}$ on $(\mathbb{T}_0,\mathcal{Q})$,
the number $\|G_\kappa\|_{\mathbb{T}_3}$ depends only on the cell of $\mathbb{T}_0$ to which $Q_\kappa$ belongs.
\end{theorem}

\begin{proof}
\emph{Chessboard families.}
Let $\mathcal{Q}=(\mathfrak{F};o_1,o_2,o_3;h)$ be a cell structure, let $\mathbb{T}=(\Gamma,\chi)$ be a torus that carries $\mathcal{Q}$,
and let $\kappa\in\Z^3$ and $G\in\mathcal{A}(Q_\kappa)$.
Recall that for $\kappa'\in\Z^3$ exactly one element $g$ of $W(\mathcal{Q})$ maps $Q_\kappa$ onto $Q_{\kappa'}$,
and that $G^{[\kappa']}=G^{g}$ for this element (Definitions~\ref{def:cell-families} and~\ref{def:folds}).
Let $\gamma\in\Gamma$ and $\kappa'\in\Z^3$, and let $g\in W(\mathcal{Q})$ map $Q_\kappa$ onto $Q_{\kappa'}$.
The translation $\tau_\gamma$ belongs to $W(\mathcal{Q})$ (Definition~\ref{def:cell-structures})
and maps $Q_{\kappa'}$ onto $Q_{\kappa'+\kappa_\gamma}$ (Definition~\ref{def:cell-families}).
So $\tau_\gamma\,g$ is the element of $W(\mathcal{Q})$ that maps $Q_\kappa$ onto $Q_{\kappa'+\kappa_\gamma}$, and \eqref{eq:setting-group-action} gives
\[
G^{[\kappa'+\kappa_\gamma]}=G^{\tau_\gamma g}=\left(G^{g}\right)^{\tau_\gamma}=\left(G^{[\kappa']}\right)^{\tau_\gamma}.
\]
Next let $a\in\{1,2,3\}$ and $j\in\Z$.
The element $g^a_{\kappa'_a\leftarrow j}$ of $W_a(\mathcal{Q})$ maps $Q_{\kappa'[a{:}j]}$ onto $Q_{\kappa'}$ (Definition~\ref{def:cell-families}).
If $g\in W(\mathcal{Q})$ maps $Q_\kappa$ onto $Q_{\kappa'[a{:}j]}$, then $g^a_{\kappa'_a\leftarrow j}\,g$ maps $Q_\kappa$ onto $Q_{\kappa'}$, and
\[
\left(G^{[\kappa'[a{:}j]]}\right)^{g^a_{\kappa'_a\leftarrow j}}=G^{g^a_{\kappa'_a\leftarrow j}\,g}=G^{[\kappa']}.
\]
For the chessboard family $\mathfrak{C}(G)=\left(G^{[\kappa']}\right)_{\kappa'\in\Z^3}$ these two computations state that
\begin{equation}\label{eq:tori-chessboard-family}
\mathfrak{C}(G)_{\kappa'+\kappa_\gamma}=\left(\mathfrak{C}(G)_{\kappa'}\right)^{\tau_\gamma}\quad(\kappa'\in\Z^3,\ \gamma\in\Gamma),
\qquad
\mathfrak{R}^j_a\,\mathfrak{C}(G)=\mathfrak{C}(G)\quad(a\in\{1,2,3\},\ j\in\Z).
\end{equation}
Since $G^{[\kappa']}\in\mathcal{A}(Q_{\kappa'})$ for every $\kappa'$, the first relation says that $\mathfrak{C}(G)$ is a family on $(\mathbb{T},\mathcal{Q})$;
its product over a set of representatives of the cells of $\mathbb{T}$ is $\Fold_{\mathbb{T}}G$ (Definition~\ref{def:folds}).

\emph{The sign of a fold.}
Being rectangular, $\mathbb{T}$ is adapted to $\ell^1$ (Definition~\ref{def:cell-structures}), so that $\mathbb{T}^0(\ell^1)=\mathbb{T}$;
and $\mathcal{Q}$ cuts $\mathbb{T}$ evenly (Definition~\ref{def:cell-families}).
Choose a set of representatives of $\{\kappa'\in\Z^3:\kappa'_1=0\}$ modulo $\kappa_{\Gamma\cap\{\ell^1=0\}}$,
and let $\mathcal{K}$ be the set of representatives of the cells of $\mathbb{T}$ formed from it in Lemma~\ref{lem:one-direction}\eqref{it:straight-representatives} for $a=1$.
Lemma~\ref{lem:one-direction}\eqref{it:straight-inequality}, applied to the family $\mathfrak{C}(G)$ with $a=1$ and $j=0$, together with \eqref{eq:tori-chessboard-family}, gives
\[
Z_{\mathbb{T}}\left[\Fold_{\mathbb{T}}G\right]=Z_{\mathbb{T}}\left[\prod_{\kappa'\in\mathcal{K}}\left(\mathfrak{R}^0_1\,\mathfrak{C}(G)\right)_{\kappa'}\right]\ge0 .
\]
This is the first assertion of \eqref{it:chessboard-carried}; in particular the seminorm $\|G\|_{\mathbb{T}}$ is defined.

\emph{Translations by periods.}
Let now $\mathbb{T}_0$ be a torus that $\mathcal{Q}$ cuts evenly, with the period lattice $\Gamma_0$,
and let $\mathbb{T}_3$ be a torus that carries $\mathcal{Q}$.
Let $\gamma\in\Gamma_0$, $\kappa\in\Z^3$ and $G\in\mathcal{A}(Q_\kappa)$.
Then $\tau_\gamma\in W(\mathcal{Q})$ and $\tau_\gamma Q_\kappa=Q_{\kappa+\kappa_\gamma}$ (Definition~\ref{def:cell-families}),
and $G^{\tau_\gamma}\in\mathcal{A}(Q_{\kappa+\kappa_\gamma})$ by Lemma~\ref{lem:admissible-images}\eqref{it:images-boxes}.
For $\kappa'\in\Z^3$ let $g\in W(\mathcal{Q})$ map $Q_{\kappa+\kappa_\gamma}$ onto $Q_{\kappa'}$.
Then $g\,\tau_\gamma\in W(\mathcal{Q})$ maps $Q_\kappa$ onto $Q_{\kappa'}$,
and by \eqref{eq:setting-group-action} the mirror copy of $G^{\tau_\gamma}$ in $Q_{\kappa'}$ is
\[
\left(G^{\tau_\gamma}\right)^{g}=G^{g\,\tau_\gamma}=G^{[\kappa']}.
\]
Hence $\mathfrak{C}(G^{\tau_\gamma})=\mathfrak{C}(G)$.
The folds of $G^{\tau_\gamma}$ and of $G$ on $\mathbb{T}_3$, formed with one set of representatives, are therefore the same integrand,
and $\|G^{\tau_\gamma}\|_{\mathbb{T}_3}=\|G\|_{\mathbb{T}_3}$, both seminorms being defined by the preceding paragraph.
If $(G_\kappa)$ is a family on $(\mathbb{T}_0,\mathcal{Q})$, then $G_{\kappa+\kappa_\gamma}=(G_\kappa)^{\tau_\gamma}$,
and so $\|G_{\kappa+\kappa_\gamma}\|_{\mathbb{T}_3}=\|G_\kappa\|_{\mathbb{T}_3}$ for all $\kappa\in\Z^3$ and $\gamma\in\Gamma_0$.

\emph{The inequality in \eqref{it:chessboard-carried}.}
Let $\mathbb{T}=(\Gamma,\chi)=\mathbb{T}(\mathfrak{F};p_1,p_2,p_3)$ carry $\mathcal{Q}$, so that $p_1,p_2,p_3\in2h\Z$.
This torus is adapted to each of its three families, with the height periods $p_1,p_2,p_3$ (Definition~\ref{def:cell-structures}),
and $\mathcal{Q}$ cuts it evenly.
Proposition~\ref{prop:three-directions} therefore applies with $\mathbb{T}_0=\mathbb{T}_1=\mathbb{T}_2=\mathbb{T}_3=\mathbb{T}$.
Let $\mathcal{K}$ and $G_\kappa$, $\kappa\in\mathcal{K}$, be as in \eqref{it:chessboard-carried},
and let $\mathfrak{G}=(G_\kappa)_{\kappa\in\Z^3}$ be the family on $(\mathbb{T},\mathcal{Q})$ that extends these contents (Definition~\ref{def:cell-families}).
The proposition gives
\[
\left|Z_{\mathbb{T}}\left[\prod_{\kappa\in\mathcal{K}_p}G_\kappa\right]\right|\le\prod_{\kappa\in\mathcal{K}_p}\|G_\kappa\|_{\mathbb{T}} .
\]
Every $\kappa\in\mathcal{K}$ can be written as $\kappa'+\kappa_\gamma$ with a unique $\kappa'\in\mathcal{K}_p$ and a unique $\gamma\in\Gamma$,
and $\kappa\mapsto\kappa'$ is a bijection of $\mathcal{K}$ onto $\mathcal{K}_p$, both sets being sets of representatives of the cells of $\mathbb{T}$.
The products of $\mathfrak{G}$ over $\mathcal{K}$ and over $\mathcal{K}_p$ define the same function on $\Conf(\Gamma)$ (Definition~\ref{def:cell-families}),
and $\|G_\kappa\|_{\mathbb{T}}=\|G_{\kappa'}\|_{\mathbb{T}}$ by the paragraph on translations, applied with $\mathbb{T}_0=\mathbb{T}_3=\mathbb{T}$.
So neither side of the last inequality changes when $\mathcal{K}_p$ is replaced by $\mathcal{K}$, and \eqref{it:chessboard-carried} is proved.

\emph{Item \eqref{it:chessboard-diagonal}.}
The torus $C(n,n,n)$ has the period lattice $n\Z^3$, which the forms $\ell_-$, $\ell_+$ and $x_3$ map onto $n\Z$.
As $n\Z\subset2h\Z$, the structure $\mathcal{Q}$ cuts $C(n,n,n)$ evenly.
Since $n$ is even and $n\ge8$, Proposition~\ref{prop:untwistings}(a) applies:
straightening $C(n,n,n)$ in the direction $\ell_-$ gives $D'(n/2,n,n)$, the height period is $p(C(n,n,n),\ell_-)=n$,
and the lattice of the cross-section $\Xi=\Xi(C(n,n,n),\ell_-)$ is
\[
\Gamma_\parallel:=n\Z^3\cap\{\ell_-=0\}=\Z\,n\cdot e_+\oplus\Z\,n\cdot e_3 .
\]
By Proposition~\ref{prop:untwistings}(b), straightening $D'(n/2,n,n)$ in the direction $\ell_+$ gives $D(n/2,n,n)$,
with the height period $p(D'(n/2,n,n),\ell_+)=2n$.
The torus $D(n/2,n,n)=\mathbb{T}(\mathfrak{F}_{\mathrm{d}};n,2n,n)$ is rectangular, hence adapted to $x_3$, with the height period $n$ (Definition~\ref{def:cell-structures}).
Proposition~\ref{prop:three-directions} thus applies to
\[
\mathbb{T}_0:=C(n,n,n),\qquad\mathbb{T}_1:=D'(n/2,n,n),\qquad\mathbb{T}_2=\mathbb{T}_3:=D(n/2,n,n),\qquad p=(n,2n,n).
\]
It shows that $D(n/2,n,n)$ carries $\mathcal{Q}$, that $\mathcal{K}_p$ represents the cells of the three tori,
and that $\#|\mathcal{K}_p|=\frac{n\cdot2n\cdot n}{h^3}=\frac{2n^3}{h^3}$.
Lemma~\ref{lem:one-direction}\eqref{it:straight-even}, applied to $C(n,n,n)$ and $\ell_-$, shows that $\mathcal{Q}$ cuts $D'(n/2,n,n)$ evenly.

Let $\mathfrak{G}=(G_\kappa)$ be the family on $(C(n,n,n),\mathcal{Q})$ that extends the contents $G_\kappa$, $\kappa\in\mathcal{K}_p$.
One has $\kappa_{n\cdot e_+}=(0,2n/h,0)$ and $\kappa_{n\cdot e_3}=(0,0,n/h)$, so that
\[
\kappa_{\Gamma_\parallel}=\{0\}\times\frac{2n}{h}\Z\times\frac{n}{h}\Z ,
\]
and $\mathcal{K}_\parallel:=\mathcal{K}_p\cap\{\kappa_1=0\}$ is a set of representatives of $\{\kappa\in\Z^3:\kappa_1=0\}$ modulo $\kappa_{\Gamma_\parallel}$.
Accordingly $\mathcal{K}_p=\left\{\kappa\in\Z^3:\ 0\le\kappa_1<n/h,\ \kappa[1{:}0]\in\mathcal{K}_\parallel\right\}$
is the set of Lemma~\ref{lem:one-direction}\eqref{it:straight-representatives} for $C(n,n,n)$ and $\ell^1=\ell_-$.
The cells $Q_\kappa$ with $\kappa_1=j-1$ lie in the open slab $\{o_1+(j-1)\cdot h<\ell_-<o_1+j\cdot h\}$,
and the cells $Q_\kappa$ with $\kappa\in\mathcal{K}_p$ lie in pairwise different orbits of $n\Z^3$.
By Lemma~\ref{lem:embedding-criteria}\eqref{it:embedding-cells}, $S_j$ is therefore, expanded into products, a sum of products of
$\mathcal{A}(\{o_1+(j-1)\cdot h<\ell_-<o_1+j\cdot h\})$ that embed in $n\Z^3$ and hence in $\Gamma_\parallel$;
that is, $S_j\in\mathcal{A}_\Xi\left[o_1+(j-1)\cdot h,\ o_1+j\cdot h\right]$ (Definition~\ref{def:slabs}).
The torus $D'(n/2,n,n)$ is adapted to $\ell_-$, with the cross-section $\Xi$ and the height period $n$ (Lemma~\ref{lem:closing-vectors}\eqref{it:closing-straightened}),
so that the fold of $S_j$ on it is the product of $n/h$ copies.
Identity~\eqref{eq:tori-layer-fold}, taken for the torus $C(n,n,n)$, the direction $\ell^1=\ell_-$, the family $\mathfrak{G}$ and the window of the layers $0\le\kappa_1<n/h$, reads
\[
\fold_{\ell_-}S_j=\prod_{\kappa\in\mathcal{K}_p}\left(\mathfrak{R}^{j-1}_1\mathfrak{G}\right)_\kappa\qquad\left(j=1,\dots,\frac{n}{h}\right).
\]
The second member of the chain of inequalities in Proposition~\ref{prop:three-directions} is therefore
\[
\prod_{0\le j_1<n/h}Z_{D'(n/2,n,n)}\left[\prod_{\kappa\in\mathcal{K}_p}\left(\mathfrak{R}^{j_1}_1\mathfrak{G}\right)_\kappa\right]^{\frac{h}{n}}
=\prod_{j=1}^{n/h}Z_{D'(n/2,n,n)}\left[\fold_{\ell_-}S_j\right]^{\frac{h}{n}} .
\]
The first member of that chain is $\left|Z_{C(n,n,n)}\left[\prod_{\kappa\in\mathcal{K}_p}G_\kappa\right]\right|$, its last member is $\prod_{\kappa\in\mathcal{K}_p}\|G_\kappa\|_{D(n/2,n,n)}$,
and the first two inequalities of \eqref{it:chessboard-diagonal} follow.
The exponents $\frac{h}{n}$, $\frac{h}{2n}$ and $\frac{h}{n}$ of the three steps multiply to $\frac{h^3}{2n^3}=\frac{1}{\#|\mathcal{K}_p|}$.

For the inequality on $D'(n/2,n,n)$, apply Proposition~\ref{prop:three-directions} to $\mathbb{T}_0:=D'(n/2,n,n)$, which $\mathcal{Q}$ cuts evenly,
with the three directions taken in the order $\ell_+$, $x_3$, $\ell_-$.
Straightening in the direction $\ell_+$ gives $D(n/2,n,n)$, with the height period $2n$ (Proposition~\ref{prop:untwistings}(b)).
This rectangular torus is adapted to $x_3$ and to $\ell_-$, with the height periods $n$ and $n$, so that the two further steps leave it unchanged.
The height periods of $\ell_-$, $\ell_+$, $x_3$ are thus again $n$, $2n$, $n$, the last torus is $D(n/2,n,n)$, and the brick is again $\mathcal{K}_p$.
The proposition, applied to the family on $(D'(n/2,n,n),\mathcal{Q})$ that extends the contents $G_\kappa$, $\kappa\in\mathcal{K}_p$,
gives the last inequality of \eqref{it:chessboard-diagonal}.

\emph{Item \eqref{it:chessboard-coordinate}.}
The period lattice $\Gamma$ of $D(a,2a,c)$ is generated by $a\cdot e_-$, $2a\cdot e_+$ and $c\cdot e_3$.
By Lemma~\ref{lem:systoles}, applied to the sides $a$, $2a$ and $c$, one has for integers $i_1,i_2,i_3$
\[
i_1\cdot a\cdot e_-+i_2\cdot2a\cdot e_++i_3\cdot c\cdot e_3=\left((i_1+2i_2)\cdot a,\ (2i_2-i_1)\cdot a,\ i_3\cdot c\right).
\]
As $i_1+2i_2$ and $2i_2-i_1$ take every integer value, $x_1(\Gamma)=x_2(\Gamma)=a\Z$ and $x_3(\Gamma)=c\Z$.
These groups lie in $2h''\Z$, so that $\mathcal{Q}$ cuts $D(a,2a,c)$ evenly.
Since $a\ge4$ and $c\ge4$, Proposition~\ref{prop:untwistings}(c) shows that straightening $D(a,2a,c)$ in the direction $x_1$ gives $C(a,4a,c)$,
with the height period $p(D(a,2a,c),x_1)=a$.
The torus $C(a,4a,c)=\mathbb{T}(\mathfrak{F}_{\mathrm{c}};a,4a,c)$ is adapted to $x_2$ and to $x_3$, with the height periods $4a$ and $c$ (Definition~\ref{def:cell-structures}).
Proposition~\ref{prop:three-directions} applies to
\[
\mathbb{T}_0:=D(a,2a,c),\qquad\mathbb{T}_1=\mathbb{T}_2=\mathbb{T}_3:=C(a,4a,c),\qquad p=(a,4a,c).
\]
It shows that $C(a,4a,c)$ carries $\mathcal{Q}$, that $\mathcal{K}_p$ represents the cells of both tori,
and that $\#|\mathcal{K}_p|=\frac{a\cdot4a\cdot c}{h''^3}=\frac{4a^2\cdot c}{h''^3}$;
for the family on $(D(a,2a,c),\mathcal{Q})$ that extends the contents $G_\kappa$, $\kappa\in\mathcal{K}_p$, its chain of inequalities gives \eqref{it:chessboard-coordinate}.
The exponents $\frac{h''}{a}$, $\frac{h''}{4a}$ and $\frac{h''}{c}$ of the three steps multiply to $\frac{1}{\#|\mathcal{K}_p|}$.

\emph{The last assertion.}
In \eqref{it:chessboard-carried} the torus on the left and the torus of the seminorms are both $\mathbb{T}$, which carries $\mathcal{Q}$ and is cut evenly by it.
In \eqref{it:chessboard-diagonal} the tori on the left are $C(n,n,n)$ and $D'(n/2,n,n)$, which $\mathcal{Q}$ cuts evenly,
and the seminorms are taken on $D(n/2,n,n)$, which carries $\mathcal{Q}$;
in \eqref{it:chessboard-coordinate} the same holds for $D(a,2a,c)$ and $C(a,4a,c)$.
The paragraph on translations therefore applies in every case.
\end{proof}

\begin{lemma}[Symmetries that respect the cells]\label{lem:cell-symmetries}
Let $-1\le\rho\le1$, $w_0,\dots,w_3>0$, $M\in\R$ and $q>0$.
Let $\mathcal{Q}=(\mathfrak{F};o_1,o_2,o_3;h)$ be a cell structure, and let $\mathbb{T}=(\Gamma,\chi)$ be a torus that carries $\mathcal{Q}$.
Let $g\in W$ satisfy $g\mathbb{T}=\mathbb{T}$ (Definition~\ref{def:tori}), and let $g$ map the set of the walls of $\mathcal{Q}$ onto itself.
Then, for every $\kappa\in\Z^3$ and every $G\in\mathcal{A}(Q_\kappa)$, the set $gQ_\kappa$ is an open cell of $\mathcal{Q}$, $G^g\in\mathcal{A}(gQ_\kappa)$, and
\[
\left\|G^g\right\|_{\mathbb{T}}=\|G\|_{\mathbb{T}} .
\]
These hypotheses on $g$ are satisfied by every translation $\tau_v$ with $v\in\Z^3$ and $\ell^a(v)\in h\Z$ for $a=1,2,3$.
\end{lemma}

\begin{proof}
For a wall $\Pi$ of $\mathcal{Q}$, the element $g\,\theta_\Pi\,g^{-1}$ of $W$ is the reflection in the plane $g\Pi$, which is again a wall.
As $g^{-1}$ also maps the set of the walls onto itself, one has $g\,W(\mathcal{Q})\,g^{-1}=W(\mathcal{Q})$.
The open cells of $\mathcal{Q}$ are convex, open and pairwise disjoint, and their union is the complement in $\R^3$ of the union of the walls;
so they are the connected components of this complement.
The homeomorphism $g$ of $\R^3$ maps the union of the walls onto itself and therefore permutes these components:
$gQ_{\kappa'}$ and $g^{-1}Q_{\kappa'}$ are open cells for every $\kappa'\in\Z^3$.
Moreover $G^g\in\mathcal{A}(gQ_\kappa)$ by Lemma~\ref{lem:admissible-images}\eqref{it:images-boxes}.

Let $\mathcal{K}$ be a set of representatives of the cells of $\mathbb{T}$.
For $\kappa'\in\mathcal{K}$ let $g_{\kappa'}$ be the element of $W(\mathcal{Q})$ that maps $gQ_\kappa$ onto $Q_{\kappa'}$, and put $g':=g^{-1}g_{\kappa'}\,g$.
Then $g'\in W(\mathcal{Q})$, and $g'$ maps $Q_\kappa$ onto the open cell $g^{-1}Q_{\kappa'}=:Q_{\kappa''}$.
By Definition~\ref{def:folds} and \eqref{eq:setting-group-action},
\[
\left(G^g\right)^{[\kappa']}=\left(G^g\right)^{g_{\kappa'}}=G^{g_{\kappa'}g}=G^{g\,g'}=\left(G^{g'}\right)^{g}=\left(G^{[\kappa'']}\right)^{g}.
\]
Let $\mathcal{K}''$ be the set of the indices $\kappa''$ obtained in this way from the points $\kappa'\in\mathcal{K}$; distinct points $\kappa'$ give distinct $\kappa''$.
Lemma~\ref{lem:admissible-images}\eqref{it:images-products} yields the identity of integrands on $\R^3$
\[
\Fold_{\mathbb{T}}\left(G^g\right)=\prod_{\kappa'\in\mathcal{K}}\left(G^g\right)^{[\kappa']}=\widetilde G^{\,g},
\qquad
\widetilde G:=\prod_{\kappa''\in\mathcal{K}''}G^{[\kappa'']} .
\]

Let $\vec g$ be the linear part of $g$.
By Definition~\ref{def:tori}, the equality $g\mathbb{T}=\mathbb{T}$ means $\vec g\,\Gamma=\Gamma$ and $\chi\circ\vec g^{-1}=\chi$.
For $\kappa'\in\Z^3$ and $\gamma\in\Gamma$ one has $g^{-1}(Q_{\kappa'}+\gamma)=g^{-1}Q_{\kappa'}+\vec g^{-1}\gamma$, and $\vec g^{-1}\gamma\in\Gamma$.
Hence $g^{-1}$ maps every $\Gamma$-orbit of open cells onto a $\Gamma$-orbit of open cells;
as $g$ does the same, $g^{-1}$ induces a bijection of the set of the cells of $\mathbb{T}$ onto itself.
The set $\mathcal{K}$ meets every cell of $\mathbb{T}$ exactly once, and therefore so does $\mathcal{K}''$:
it is again a set of representatives of the cells of $\mathbb{T}$.

By \eqref{eq:tori-chessboard-family}, $\mathfrak{C}(G)$ is a family on $(\mathbb{T},\mathcal{Q})$.
Its products over $\mathcal{K}''$ and over $\mathcal{K}$ are $\widetilde G$ and $\Fold_{\mathbb{T}}G$,
and they define the same function on $\Conf(\Gamma)$ (Definition~\ref{def:cell-families});
thus $Z_{\mathbb{T}}[\widetilde G]=Z_{\mathbb{T}}[\Fold_{\mathbb{T}}G]$.
Expanded into products, $\widetilde G$ embeds in $\mathbb{T}$ (Lemma~\ref{lem:embedding-criteria}\eqref{it:embedding-cells}).
Proposition~\ref{prop:symmetries}\eqref{it:functional-covariance} and $g\mathbb{T}=\mathbb{T}$ now give
\[
Z_{\mathbb{T}}\left[\Fold_{\mathbb{T}}\left(G^g\right)\right]=Z_{g\mathbb{T}}\left[\widetilde G^{\,g}\right]
=\begin{cases}
Z_{\mathbb{T}}[\widetilde G]&\text{if $g$ is even},\\
\overline{Z_{\mathbb{T}}[\widetilde G]}&\text{if $g$ is odd}.
\end{cases}
\]
The number $Z_{\mathbb{T}}[\widetilde G]=Z_{\mathbb{T}}[\Fold_{\mathbb{T}}G]$ is real, indeed non-negative, by Theorem~\ref{thm:chessboard}\eqref{it:chessboard-carried}.
Consequently $Z_{\mathbb{T}}\left[\Fold_{\mathbb{T}}\left(G^g\right)\right]=Z_{\mathbb{T}}\left[\Fold_{\mathbb{T}}G\right]$,
and the two seminorms, the powers of these numbers with the exponent $|Q|/|\mathbb{T}|$, coincide.

Finally let $v\in\Z^3$ satisfy $\ell^a(v)\in h\Z$ for $a=1,2,3$.
The linear part of $\tau_v$ is the identity, so that $\tau_v\mathbb{T}=\mathbb{T}$ (Definition~\ref{def:tori}).
Since $\ell^a(x+v)=\ell^a(x)+\ell^a(v)$ for $x\in\R^3$, the translation $\tau_v$ maps the wall $(\ell^a,o_a+j\cdot h)$ onto the plane $\{\ell^a=o_a+j\cdot h+\ell^a(v)\}$,
which is the wall $\left(\ell^a,o_a+(j+\ell^a(v)/h)\cdot h\right)$;
correspondingly $\tau_v\,\theta_{\ell^a,c}\,\tau_{-v}=\theta_{\ell^a,c+\ell^a(v)}$ (Definition~\ref{def:mirrors}\eqref{it:mirror-conjugation}).
The same holds for $\tau_{-v}$, and so $\tau_v$ maps the set of the walls of $\mathcal{Q}$ onto itself.
\end{proof}

\subsection{Pressure and shape independence}\label{sec:tori-pressure}

\begin{lemma}[Separately convex functions]\label{lem:separately-convex}
For $\mu=1,2,3$ let $\mathcal{N}_\mu\subset(0,\infty)$ be a set that is not bounded above,
and fix $a_\mu,b_\mu\in\mathcal{N}_\mu$ with $a_\mu<b_\mu$.
On the product $\mathcal{N}:=\mathcal{N}_1\times\mathcal{N}_2\times\mathcal{N}_3$ put $|n|:=n_1n_2n_3$ for $n=(n_1,n_2,n_3)$.
Let $u\colon\mathcal{N}\to\R$ be a function.
Given $\mu\in\{1,2,3\}$ and values of the two coordinates other than the $\mu$-th,
denote by $\tilde u\colon\mathcal{N}_\mu\to\R$ the map that sends $s$ to the value of $u$ at the point whose $\mu$-th coordinate is $s$ and whose two other coordinates are the given ones.
Suppose that every function $\tilde u$ obtained in this way satisfies
\begin{equation}\tag{Mon}\label{cond:monotone-ratio}
\frac{\tilde u(s')}{s'}\le\frac{\tilde u(s)}{s}\qquad\text{for }s<s'\text{ in }\mathcal{N}_\mu,
\end{equation}
and
\begin{equation}\tag{Cvx}\label{cond:convexity}
(s''-s)\cdot\tilde u(s')\le(s''-s')\cdot\tilde u(s)+(s'-s)\cdot\tilde u(s'')\qquad\text{for }s<s'<s''\text{ in }\mathcal{N}_\mu .
\end{equation}
Then the following hold.
\begin{enumerate}
\item\label{it:ratio-monotone}
The quotient $u(n)/|n|$ is a non-increasing function of each variable $n_\mu$.
\item\label{it:four-values}
Let $n\in\mathcal{N}$ satisfy $n_\mu\ge b_\mu$ for $\mu=1,2,3$,
and put $\alpha_\mu:=(n_\mu-a_\mu)/(b_\mu-a_\mu)$, so that $\alpha_\mu\ge1$.
Then
\[
\begin{aligned}
u(n)\ge{}&\alpha_1\alpha_2\alpha_3\cdot u(b_1,b_2,b_3)-\alpha_1\alpha_2(\alpha_3-1)\cdot u(b_1,b_2,a_3)\\
&-\alpha_1(\alpha_2-1)\cdot\frac{n_3}{a_3}\cdot u(b_1,a_2,a_3)-(\alpha_1-1)\cdot\frac{n_2n_3}{a_2a_3}\cdot u(a_1,a_2,a_3).
\end{aligned}
\]
\item\label{it:convex-floor}
For all $n\in\mathcal{N}$ one has $u(n)/|n|\ge-C_u$, where
\[
C_u:=\frac{\left|u(b_1,b_2,b_3)\right|+\left|u(b_1,b_2,a_3)\right|}{(b_1-a_1)(b_2-a_2)(b_3-a_3)}
+\frac{\left|u(b_1,a_2,a_3)\right|}{(b_1-a_1)(b_2-a_2)\cdot a_3}
+\frac{\left|u(a_1,a_2,a_3)\right|}{(b_1-a_1)\cdot a_2a_3}.
\]
\item\label{it:any-manner}
The number $u_\infty:=\inf\left\{u(n)/|n|:\ n\in\mathcal{N}\right\}$ is finite,
and $u(n)/|n|$ converges to $u_\infty$ as $\min_\mu n_\mu\to\infty$ in any manner:
for every $\varepsilon>0$ there is $n^0\in\mathcal{N}$ such that
\[
\left|\frac{u(n)}{|n|}-u_\infty\right|<\varepsilon
\qquad\text{for all }n\in\mathcal{N}\text{ with }n_\mu\ge n^0_\mu,\ \mu=1,2,3.
\]
\end{enumerate}
\end{lemma}

\begin{proof}
The argument is the one of \cite[Lemma~10.10]{NoncompactQED},
carried out here for three independent variables,
with an arbitrary pair $a_\mu<b_\mu$ in each variable in the place of the base points $1<j$ used there.

For \eqref{it:ratio-monotone}, fix $\mu$ together with the two coordinates other than the $\mu$-th,
and let $\tilde u$ be the corresponding function of the statement.
Multiplication of $\tilde u(n_\mu)/n_\mu$ by the reciprocal of the product of the two fixed coordinates,
a positive factor independent of $n_\mu$, yields $u(n)/|n|$.
By \eqref{cond:monotone-ratio} the quotient $\tilde u(n_\mu)/n_\mu$ does not increase with $n_\mu$, and therefore neither does $u(n)/|n|$.

\emph{Two inequalities in one variable.}
Again fix $\mu$ and the two other coordinates, and let $\tilde u$ be as in the statement.
For $s\in\mathcal{N}_\mu$ one has
\begin{align}
\tilde u(s)&\ge\alpha\cdot\tilde u(b_\mu)-(\alpha-1)\cdot\tilde u(a_\mu),\qquad\alpha:=\frac{s-a_\mu}{b_\mu-a_\mu},&&\text{if }s\ge b_\mu,\label{eq:tori-secant}\\
\tilde u(s)&\le\frac{s}{a_\mu}\cdot\tilde u(a_\mu)&&\text{if }s\ge a_\mu.\label{eq:tori-ratio}
\end{align}
For $s=b_\mu$ the number $\alpha$ equals $1$, and \eqref{eq:tori-secant} is an equality.
For $s>b_\mu$, the condition \eqref{cond:convexity} at the three points $a_\mu<b_\mu<s$ reads
\[
(s-a_\mu)\cdot\tilde u(b_\mu)\le(s-b_\mu)\cdot\tilde u(a_\mu)+(b_\mu-a_\mu)\cdot\tilde u(s);
\]
dividing by $b_\mu-a_\mu>0$ and observing that $(s-a_\mu)/(b_\mu-a_\mu)=\alpha$ and $(s-b_\mu)/(b_\mu-a_\mu)=\alpha-1$, one obtains \eqref{eq:tori-secant}.
The inequality \eqref{eq:tori-ratio} is an equality for $s=a_\mu$;
for $s>a_\mu$ it is \eqref{cond:monotone-ratio} for the pair $a_\mu<s$, multiplied by $s>0$.
At $s=n_\mu$ the number $\alpha$ is the coefficient $\alpha_\mu$ of \eqref{it:four-values}.

Now let $n$ be as in \eqref{it:four-values}.
Apply \eqref{eq:tori-secant} in the first variable with the other coordinates $n_2,n_3$,
in the second variable with the other coordinates $b_1,n_3$,
and in the third variable with the other coordinates $b_1,b_2$:
\[
\begin{aligned}
u(n_1,n_2,n_3)&\ge\alpha_1\cdot u(b_1,n_2,n_3)-(\alpha_1-1)\cdot u(a_1,n_2,n_3),\\
u(b_1,n_2,n_3)&\ge\alpha_2\cdot u(b_1,b_2,n_3)-(\alpha_2-1)\cdot u(b_1,a_2,n_3),\\
u(b_1,b_2,n_3)&\ge\alpha_3\cdot u(b_1,b_2,b_3)-(\alpha_3-1)\cdot u(b_1,b_2,a_3).
\end{aligned}
\]
Since $n_2>a_2$ and $n_3>a_3$, the inequality \eqref{eq:tori-ratio} in the second and in the third variable gives
\[
u(a_1,n_2,n_3)\le\frac{n_2}{a_2}\cdot u(a_1,a_2,n_3)\le\frac{n_2n_3}{a_2a_3}\cdot u(a_1,a_2,a_3),
\qquad
u(b_1,a_2,n_3)\le\frac{n_3}{a_3}\cdot u(b_1,a_2,a_3);
\]
in the first chain, the second inequality is \eqref{eq:tori-ratio} in the third variable, multiplied by $n_2/a_2>0$.
In the first two of the three inequalities obtained from \eqref{eq:tori-secant}, the values $u(a_1,n_2,n_3)$ and $u(b_1,a_2,n_3)$ carry the coefficients $-(\alpha_1-1)\le0$ and $-(\alpha_2-1)\le0$.
Multiplied by these coefficients, the upper bounds just obtained become lower bounds, whatever the signs of the values of $u$:
\[
\begin{aligned}
u(n)&\ge\alpha_1\cdot u(b_1,n_2,n_3)-(\alpha_1-1)\cdot\frac{n_2n_3}{a_2a_3}\cdot u(a_1,a_2,a_3),\\
u(b_1,n_2,n_3)&\ge\alpha_2\cdot u(b_1,b_2,n_3)-(\alpha_2-1)\cdot\frac{n_3}{a_3}\cdot u(b_1,a_2,a_3).
\end{aligned}
\]
Insert the lower bound for $u(b_1,b_2,n_3)$ into the second of these inequalities, whose coefficient $\alpha_2$ is non-negative,
and the resulting lower bound for $u(b_1,n_2,n_3)$ into the first one, whose coefficient $\alpha_1$ is non-negative.
The outcome is the inequality of \eqref{it:four-values}.

To prove \eqref{it:convex-floor}, assume first that $n_\mu\ge b_\mu$ for $\mu=1,2,3$.
Because $a_\mu>0$,
\[
0\le\alpha_\mu-1\le\alpha_\mu\le\frac{n_\mu}{b_\mu-a_\mu}\qquad(\mu=1,2,3).
\]
Each of the four terms on the right of the inequality in \eqref{it:four-values} is a value of $u$ multiplied by a real number,
and it is at least minus the modulus of this number times the modulus of the value.
The moduli of the four numbers are
$\alpha_1\alpha_2\alpha_3$, $\alpha_1\alpha_2(\alpha_3-1)$, $\alpha_1(\alpha_2-1)\cdot n_3/a_3$ and $(\alpha_1-1)\cdot n_2n_3/(a_2a_3)$,
and by the last display
\[
\begin{aligned}
&\frac{\alpha_1\alpha_2\alpha_3}{|n|}\le\frac{1}{(b_1-a_1)(b_2-a_2)(b_3-a_3)},&&
\frac{\alpha_1\alpha_2(\alpha_3-1)}{|n|}\le\frac{1}{(b_1-a_1)(b_2-a_2)(b_3-a_3)},\\
&\frac{\alpha_1(\alpha_2-1)}{|n|}\cdot\frac{n_3}{a_3}\le\frac{1}{(b_1-a_1)(b_2-a_2)\cdot a_3},&&
\frac{\alpha_1-1}{|n|}\cdot\frac{n_2n_3}{a_2a_3}\le\frac{1}{(b_1-a_1)\cdot a_2a_3}.
\end{aligned}
\]
Division of \eqref{it:four-values} by $|n|$ therefore gives $u(n)/|n|\ge-C_u$.
There remains the case of an arbitrary $n\in\mathcal{N}$.
For each $\mu$ the number $\max(n_\mu,b_\mu)$ equals $n_\mu$ or $b_\mu$ and hence lies in $\mathcal{N}_\mu$;
let $n'\in\mathcal{N}$ be the point with the coordinates $n'_\mu:=\max(n_\mu,b_\mu)$, $\mu=1,2,3$.
It satisfies $n'_\mu\ge b_\mu$ and $n'_\mu\ge n_\mu$ for each $\mu$.
Raising the coordinates of $n$ to those of $n'$ one at a time,
one obtains from part \eqref{it:ratio-monotone} at each of the three steps, and from the case already treated at the end, that
\[
\frac{u(n)}{|n|}\ge\frac{u(n')}{|n'|}\ge-C_u .
\]

Finally, $u_\infty\ge-C_u$ by \eqref{it:convex-floor}, so that the infimum in \eqref{it:any-manner} is finite.
Let $\varepsilon>0$ and choose $n^0\in\mathcal{N}$ with $u(n^0)/|n^0|<u_\infty+\varepsilon$.
If $n\in\mathcal{N}$ and $n_\mu\ge n^0_\mu$ for $\mu=1,2,3$,
then \eqref{it:ratio-monotone}, applied in the three variables one after the other, yields
\[
u_\infty\le\frac{u(n)}{|n|}\le\frac{u(n^0)}{|n^0|}<u_\infty+\varepsilon .
\]
\end{proof}

\begin{proposition}[The pressure of coordinate tori]\label{prop:coordinate-pressure}
Let $-1\le\rho\le1$, $w_0,\dots,w_3>0$, $M\in\R$ and $q>0$.
For $n=(n_1,n_2,n_3)\in\Z^3$ with $n_\mu\ge4$ for $\mu=1,2,3$,
let $f(n)=|n|^{-1}\log Z_{C(n)}[1]$ be the pressure of the coordinate torus $C(n)$ (Definitions~\ref{def:tori} and~\ref{def:sides}).
\begin{enumerate}
\item
The pressure $f$ is non-increasing in each side $n_\mu$.
\item
For each $\mu$ and fixed values of the two other sides,
the function $n_\mu\mapsto\log Z_{C(n)}[1]$ is the restriction to the integers $n_\mu\ge4$ of a convex function on $[1,\infty)$.
\item
The pressure $f$ is bounded below.
More precisely, for integers $4\le a_0<b_0$ one has $f(n)\ge-C_u$ for all $n$,
where $C_u$ is the number of Lemma~\ref{lem:separately-convex}\eqref{it:convex-floor}
for the function $u(n):=\log Z_{C(n)}[1]$ and the points $a_\mu:=a_0$, $b_\mu:=b_0$ ($\mu=1,2,3$);
it is formed from the four values $\log Z_{\mathbb{T}}[1]$ at $\mathbb{T}=C(a_0,a_0,a_0)$, $C(b_0,a_0,a_0)$, $C(b_0,b_0,a_0)$, $C(b_0,b_0,b_0)$:
\[
C_u=\frac{\left|\log Z_{C(b_0,b_0,b_0)}[1]\right|+\left|\log Z_{C(b_0,b_0,a_0)}[1]\right|}{(b_0-a_0)^3}
+\frac{\left|\log Z_{C(b_0,a_0,a_0)}[1]\right|}{(b_0-a_0)^2\cdot a_0}
+\frac{\left|\log Z_{C(a_0,a_0,a_0)}[1]\right|}{(b_0-a_0)\cdot a_0^2}.
\]
Moreover, let
\[
C_*=\left(\frac{w_0}{2}\right)^{32}\left(\frac{w_1}{2}\right)^{64}\left(\frac{w_2}{2}\right)^{64}\left(\frac{w_3}{2}\right)^{32}\left(1-\rho^2\right)^{96}
\]
be the constant of Theorem~\ref{thm:floor}.
If $\rho^2\neq1$, every torus $\mathbb{T}$ satisfies
\[
\frac{1}{|\mathbb{T}|}\log Z_{\mathbb{T}}[1]\ge\log C_*+3\log\frac{2\pi}{q}-6\left(\frac{2\pi}{q}\right)^2,
\]
so that $\inf f$ is at least the right side;
and every torus $\mathbb{T}$ satisfies
\[
\frac{1}{|\mathbb{T}|}\log Z_{\mathbb{T}}[1]\le192\log\left(|M|+\frac{1+|\rho|}{2}\sum_{i=0}^{3}w_i\right)+3\log\left(\frac{2\pi}{q}\cdot C_q\right),
\]
where $C_q=\sup_{y\in\R}\sum_{k\in\Z}\exp\left(-\frac{1}{2}\left(y+\frac{2\pi}{q}\cdot k\right)^{2}\right)$ is the constant of Lemma~\ref{lem:bounded-integrands}.
These two bounds on arbitrary tori hold for every real $\rho$, the lower one under the condition $\rho^2\neq1$.
\item
As $\min_\mu n_\mu\to\infty$ in any manner, in the sense of Lemma~\ref{lem:separately-convex}\eqref{it:any-manner},
$f(n)$ converges to $f_\infty:=\inf f$, and $f_\infty$ is a finite number.
\end{enumerate}
\end{proposition}

\begin{proof}
Fix integers $n_2,n_3\ge4$ and a position $c\in\frac{1}{2}\Z$ of the family $x_1$.
Corollary~\ref{cor:trace-forms}\eqref{it:unit-step},
applied to the family $x_1$ and to the cross-section $\Xi=(\Gamma_\parallel,\chi_\parallel)$ with $\Gamma_\parallel=n_2\Z e_2+n_3\Z e_3$ and $\chi_\parallel(n_2e_2)=\chi_\parallel(n_3e_3)=-1$,
gives the following for the half step $\mathfrak{S}_{1/2}=\mathfrak{S}[1]\colon\mathcal{H}(c)\to\mathcal{H}\left(c+\frac{1}{2}\right)$:
the unit step $\mathfrak{S}_{1/2}^{*}\mathfrak{S}_{1/2}$ is a non-negative operator of trace class on $\mathcal{H}(c)$, it does not depend on $n_1$, and
\begin{equation}\label{eq:tori-unit-step}
Z_{C(n)}[1]=\Tr\left(\mathfrak{S}_{1/2}^{*}\mathfrak{S}_{1/2}\right)^{n_1}\qquad\text{for every integer }n_1\ge4.
\end{equation}
The same corollary provides the analogous identities in the directions $e_2$ and $e_3$, with the roles of the sides interchanged.
Proposition~\ref{prop:twisted-holder} leads to \eqref{eq:tori-unit-step} as well.
By Definition~\ref{def:cell-structures}, the torus $C(n)$ is adapted to $x_1$ and has the height period $n_1$ for $x_1$;
its cross-section $\Xi(C(n),x_1)$ is $\Xi$.
In that proposition take $\mathbb{T}=C(n)$, whose straightened torus is then $\mathbb{T}^0=C(n)$,
together with $\ell=x_1$, $c_0=c$, $k=n_1$ and $h'=\frac{1}{2}$;
thus $p=n_1=2k\cdot h'$, and all positions $c+\frac{j}{2}$, $j\in\Z$, lie in $\Adm(x_1)=\frac{1}{2}\Z$.
With $G_1=\dots=G_{2k}=1$ one has $\fold_{x_1}1=1$,
and \eqref{eq:tori-fold-norm} for $j=1$ equates $Z_{C(n)}[1]$ with $\left\|\mathfrak{S}[1]\right\|_{2n_1}^{2n_1}$,
where $\mathfrak{S}[1]\colon\mathcal{H}(c)\to\mathcal{H}\left(c+\frac{1}{2}\right)$ is the half step $\mathfrak{S}_{1/2}$;
by Definition~\ref{def:even-schatten} this number is $\Tr\left(\mathfrak{S}_{1/2}^{*}\mathfrak{S}_{1/2}\right)^{n_1}$.

Theorem~\ref{thm:partition-positivity} gives $Z_{C(n)}[1]>0$.
Consequently $\mathfrak{S}_{1/2}^{*}\mathfrak{S}_{1/2}\neq0$, and $u(n):=\log Z_{C(n)}[1]$ is a real number for every $n$.

(a) For $s\ge1$, the Schatten norm of the non-negative operator $\mathfrak{S}_{1/2}^{*}\mathfrak{S}_{1/2}$ is
$\left\|\mathfrak{S}_{1/2}^{*}\mathfrak{S}_{1/2}\right\|_s=\left(\Tr\left(\mathfrak{S}_{1/2}^{*}\mathfrak{S}_{1/2}\right)^{s}\right)^{\frac{1}{s}}$,
so that \eqref{eq:tori-unit-step} can be written
\[
f(n)=\frac{1}{n_2n_3}\cdot\log\left\|\mathfrak{S}_{1/2}^{*}\mathfrak{S}_{1/2}\right\|_{n_1}.
\]
By Lemma~\ref{lem:schatten-traces}(b), $\left\|\mathfrak{S}_{1/2}^{*}\mathfrak{S}_{1/2}\right\|_s\le\left\|\mathfrak{S}_{1/2}^{*}\mathfrak{S}_{1/2}\right\|_{s'}$ whenever $s\ge s'\ge1$.
Hence $f$ is non-increasing in $n_1$;
the identities in the directions $e_2$ and $e_3$ give the same in $n_2$ and in $n_3$.

(b) By \eqref{eq:tori-unit-step}, $u(n)=\log\Tr\left(\mathfrak{S}_{1/2}^{*}\mathfrak{S}_{1/2}\right)^{n_1}$.
The operator $\mathfrak{S}_{1/2}^{*}\mathfrak{S}_{1/2}$ is non-negative, of trace class and different from zero;
so Lemma~\ref{lem:schatten-traces}(c) shows that $s\mapsto\log\Tr\left(\mathfrak{S}_{1/2}^{*}\mathfrak{S}_{1/2}\right)^{s}$ is convex on $[1,\infty)$,
and $n_1\mapsto u(n)$ is its restriction to the integers $n_1\ge4$.
The directions $e_2$ and $e_3$ are treated in the same way.

(c), (d) Put $\mathcal{N}_\mu:=\{4,5,6,\dots\}$ for $\mu=1,2,3$, a subset of $(0,\infty)$ that is unbounded above.
When two sides are fixed, $u(n)/n_\mu$ equals $f(n)$ multiplied by the product of these two sides;
so $u$ satisfies \eqref{cond:monotone-ratio} by (a).
It satisfies \eqref{cond:convexity} by (b):
for $s<s'<s''$ the point $s'$ is the convex combination
$\frac{s''-s'}{s''-s}\cdot s+\frac{s'-s}{s''-s}\cdot s''$ of $s$ and $s''$,
and a convex function on $[1,\infty)$ takes at $s'$ a value that is at most the same combination of its values at $s$ and $s''$;
multiplied by $s''-s>0$, this is \eqref{cond:convexity}.
Lemma~\ref{lem:separately-convex}, with $a_\mu:=a_0$ and $b_\mu:=b_0$ for $\mu=1,2,3$, therefore applies to $u$.
Since $f(n)=u(n)/|n|$, its part \eqref{it:convex-floor} is the bound $f\ge-C_u$ of (c),
and its part \eqref{it:any-manner} is (d), with $f_\infty=u_\infty$.

It remains to derive the two bounds of (c) on arbitrary tori.
For every torus $\mathbb{T}$, the lower bound \eqref{eq:positivity-floor} of Theorem~\ref{thm:floor}(c), valid for $\rho^2\neq1$, gives
\[
\begin{aligned}
\log Z_{\mathbb{T}}[1]
&\ge\log\left(\left(\frac{2\pi}{q}\cdot\exp\left(-2\left(\frac{2\pi}{q}\right)^2\right)\right)^{3|\mathbb{T}|}\cdot C_*^{|\mathbb{T}|}\right)\\
&=|\mathbb{T}|\cdot\left(\log C_*+3\log\frac{2\pi}{q}-6\left(\frac{2\pi}{q}\right)^2\right),
\end{aligned}
\]
and the upper bound \eqref{eq:positivity-ceiling} of the same theorem, valid for every real $\rho$, gives
\[
\begin{aligned}
\log Z_{\mathbb{T}}[1]
&\le\log\left(\left(\frac{2\pi}{q}\cdot C_q\right)^{3|\mathbb{T}|}\cdot\left(|M|+\frac{1+|\rho|}{2}\sum_{i=0}^{3}w_i\right)^{192|\mathbb{T}|}\right)\\
&=|\mathbb{T}|\cdot\left(3\log\left(\frac{2\pi}{q}\cdot C_q\right)+192\log\left(|M|+\frac{1+|\rho|}{2}\sum_{i=0}^{3}w_i\right)\right).
\end{aligned}
\]
Division by $|\mathbb{T}|$ gives the two bounds.
For $\mathbb{T}=C(n)$ one has $|\mathbb{T}|=|n|$, and the lower bound becomes $f(n)\ge\log C_*+3\log\frac{2\pi}{q}-6\left(\frac{2\pi}{q}\right)^2$ for every $n$;
the same number is then a lower bound for $\inf f$.
\end{proof}

\begin{theorem}[The chain of four tori]\label{thm:four-tori}
Let $-1\le\rho\le1$, $w_0,\dots,w_3>0$, $M\in\R$ and $q>0$,
and let $f(n_1,n_2,n_3)=\left(n_1n_2n_3\right)^{-1}\cdot\log Z_{C(n_1,n_2,n_3)}[1]$ be the pressure of Definition~\ref{def:sides}.
Let $n\ge8$ be an even integer, and consider the four tori
$C(n)=C(n,n,n)$, $D'(n)=D'(n/2,n,n)$, $D(n)=D(n/2,n,n)$ and $C(n/2,2n,n)$.
Then
\begin{equation}\label{eq:tori-chain-order}
Z_{C(n)}[1]\le Z_{D'(n)}[1]\le Z_{D(n)}[1]\le Z_{C(n/2,2n,n)}[1] .
\end{equation}
With
\[
\Delta(n):=f(n/2,2n,n)-f(n,n,n)
\]
one has
\[
0\le\log Z_{C(n/2,2n,n)}[1]-\log Z_{C(n)}[1]=n^3\cdot\Delta(n)
\]
and
\[
0\le\Delta(n)\le f(n/2,n/2,n/2)-f_\infty\longrightarrow0\qquad(n\to\infty),
\]
where $f_\infty=\inf f$ is the finite number of Proposition~\ref{prop:coordinate-pressure}(d).
Consequently, for any two of the four tori,
the logarithm of the ratio of their partition functions lies between $-n^3\cdot\Delta(n)$ and $n^3\cdot\Delta(n)$;
in particular it is $o(n^3)$ as $n\to\infty$.
Finally, let $(n_i)$ be a sequence of even integers with $n_i\to\infty$,
and let $\Delta_i=\Delta(n_i)$ for the indices $i$ with $n_i\ge8$ (Definition~\ref{def:sides}).
Then $\Delta_i\ge0$ for all these indices, $\Delta_i\to0$ as $i\to\infty$, and
\begin{equation}\label{eq:tori-chain}
0\le\log Z_{\mathbb{T}'}[1]-\log Z_{\mathbb{T}}[1]\le n_i^3\cdot\Delta_i
\end{equation}
whenever $n_i\ge8$ and $\mathbb{T}$ stands to the left of $\mathbb{T}'$ in the list
$C(n_i)$, $D'(n_i)$, $D(n_i)$, $C(n_i/2,2n_i,n_i)$.
\end{theorem}

\begin{proof}
By Theorem~\ref{thm:partition-positivity}, $0<Z_{\mathbb{T}}[1]<\infty$ for every torus $\mathbb{T}$,
so that all the logarithms that occur are real numbers.

\emph{One straightening.}
Let $\mathbb{T}$ be a torus, $\ell\in\mathcal{L}$ a family and $g$ a closing vector of $\mathbb{T}$ for $\ell$,
with the height period $p=\ell(g)$, and suppose that the straightened torus $\mathbb{T}^0=\mathbb{T}^0(\ell)$ is defined.
Let $c_0\in\Adm(\ell)$, let $k\ge1$ be an integer and $h'>0$ a number with $p=2k\cdot h'$ and $c_j:=c_0+j\cdot h'\in\Adm(\ell)$ for all $j\in\Z$.
The constant $1$ is an admissible product in the sense of Definition~\ref{def:integrands},
with no factor of Gaussian growth and with the constant bounded factor $1$;
it has no box, so it embeds in every lattice,
and it lies in $\mathcal{A}_\Xi[c,c']$ for every cross-section $\Xi$ for $\ell$ and all $c<c'$ in $\Adm(\ell)$ (Definition~\ref{def:slabs}).
Every element of $W$ maps $1$ to $1$ (Lemma~\ref{lem:admissible-images}\eqref{it:images-products}).
Consequently every copy of $1$ in a slab, in the sense of Definition~\ref{def:slab-copies}, is $1$,
and so is the fold $\fold_\ell1$ of that definition, the product of $2k$ such copies.
With $G_1=\dots=G_{2k}=1$, the inequality~\eqref{eq:tori-holder} of Proposition~\ref{prop:twisted-holder} therefore becomes
\begin{equation}\label{eq:tori-straighten}
Z_{\mathbb{T}}[1]=\left|Z_{\mathbb{T}}[1]\right|\le\prod_{j=1}^{2k}Z_{\mathbb{T}^0}[1]^{\frac{1}{2k}}=Z_{\mathbb{T}^0}[1] .
\end{equation}

\emph{Three straightenings.}
In each of the following three cases take $c_0=0$, $k=1$ and $h'=p/2$, that is, two slabs.
\begin{itemize}
\item
$\mathbb{T}=C(n)$ and $\ell=\ell_-$.
By Proposition~\ref{prop:untwistings}(a), $p=n$ and $\mathbb{T}^0=D'(n)$.
The number $h'=n/2$ is an integer because $n$ is even, and $c_j=j\cdot n/2\in\Z=\Adm(\ell_-)$.
\item
$\mathbb{T}=D'(n)$ and $\ell=\ell_+$.
By Proposition~\ref{prop:untwistings}(b), $p=2n$ and $\mathbb{T}^0=D(n)$;
here $h'=n$ and $c_j=j\cdot n\in\Z=\Adm(\ell_+)$.
The closing vector $g=n\cdot e_+$ has $\chi(g)=+1$,
and the twist of the trace formula then contains the fermion parity $U_{\mathrm{f}}$ (Remark~\ref{rem:twist}).
\item
$\mathbb{T}=D(n)$ and $\ell=x_1$.
The torus $D(n)=D(n/2,n,n)$ is the torus of Proposition~\ref{prop:untwistings}(c)
whose first side is $n/2\ge4$ and whose third side is $n\ge4$;
so $p=n/2$ and $\mathbb{T}^0=C(n/2,2n,n)$.
As $n$ is even, $2h'=n/2$ is an integer;
thus $h'=n/4\in\tfrac12\Z$, and $c_j=j\cdot n/4\in\tfrac12\Z=\Adm(x_1)$.
When $n\equiv2\pmod4$, mirrors of $x_1$ of both kinds bound the two slabs:
the positions $c_0=0$ and $c_2=n/2$ are integers, so that their mirrors are site planes,
while $c_1=n/4$ lies in $\Z+\tfrac12$, and its mirror is a link plane.
\end{itemize}
In the three cases, \eqref{eq:tori-straighten} yields the three inequalities of~\eqref{eq:tori-chain-order}.
Other admissible choices give the same inequalities,
for instance $h'=1,1,\tfrac12$ with $k=n/2,n,n/2$ in the three cases.

\emph{Pressures.}
By Definition~\ref{def:tori}, $|C(n_1,n_2,n_3)|=n_1n_2n_3$ and $|D(a,b,c)|=|D'(a,b,c)|=2abc$;
hence each of the four tori has $n^3$ photon sites.
As $\log Z_{C(n_1,n_2,n_3)}[1]=n_1n_2n_3\cdot f(n_1,n_2,n_3)$,
\[
\log Z_{C(n/2,2n,n)}[1]-\log Z_{C(n)}[1]=n^3\cdot\left(f(n/2,2n,n)-f(n,n,n)\right)=n^3\cdot\Delta(n),
\]
and this number is non-negative by~\eqref{eq:tori-chain-order}.
All sides of the tori $C(n/2,n/2,n/2)$, $C(n/2,n/2,n)$ and $C(n/2,2n,n)$ are at least $4$,
and $f$ is non-increasing in each side (Proposition~\ref{prop:coordinate-pressure}(a)); thus
\[
f(n/2,2n,n)\le f(n/2,n/2,n)\le f(n/2,n/2,n/2).
\]
Together with $f(n,n,n)\ge\inf f=f_\infty$ this gives $\Delta(n)\le f(n/2,n/2,n/2)-f_\infty$.
The smallest side $n/2$ of $C(n/2,n/2,n/2)$ tends to infinity with $n$,
and Proposition~\ref{prop:coordinate-pressure}(d) shows that the right side tends to $0$.

\emph{Ratios.}
By~\eqref{eq:tori-chain-order}, the logarithms of the partition functions of $C(n)$, $D'(n)$, $D(n)$ and $C(n/2,2n,n)$, taken in this order,
form a non-decreasing sequence, whose last term exceeds its first by $n^3\cdot\Delta(n)$.
Hence any two of these four logarithms differ by at most $n^3\cdot\Delta(n)$ in absolute value,
and their difference is $o(n^3)$ as $n\to\infty$, since $\Delta(n)\to0$.
For a sequence of even integers $n_i\to\infty$ and the indices $i$ with $n_i\ge8$,
these statements with $n=n_i$ give $\Delta_i\ge0$ and~\eqref{eq:tori-chain};
the limit $\Delta_i\to0$ is the limit $\Delta(n)\to0$ along the sequence $(n_i)$.
\end{proof}

\begin{remark}
The tori $D'(n)$ and $D(n)$ have the same lattice and differ only in their sign characters (Definition~\ref{def:tori}).
The constant $1$ lies in $\mathcal{A}(\R^3)$ and embeds in every lattice,
and by Proposition~\ref{prop:symmetries}\eqref{it:sign-character} the value of the functional on such an integrand
does not depend on the sign character of the torus.
Therefore the second inequality in~\eqref{eq:tori-chain-order} is an equality:
$Z_{D'(n)}[1]=Z_{D(n)}[1]$.
\end{remark}

\begin{definition}[Shapes]\label{def:shapes}
Let $\mathcal{Q}=(\mathfrak{F};o_1,o_2,o_3;h)$ be a cell structure (Definition~\ref{def:cell-structures}),
let $\ell^1,\ell^2,\ell^3$ be the families of the frame $\mathfrak{F}$,
with the normals $n_1,n_2,n_3$,
and let $Q_0$ be the cell of $\mathcal{Q}$ with the index $\kappa=0$.
A \emph{shape} for $\mathcal{Q}$ is a triple $\mathfrak{m}=(m_1,m_2,m_3)$ of even positive integers
for which the rectangular torus
\[
\mathbb{T}(\mathfrak{F};h\mathfrak{m}):=\mathbb{T}(\mathfrak{F};hm_1,hm_2,hm_3)
\]
is defined in the sense of Definition~\ref{def:cell-structures}.
The set of the shapes for $\mathcal{Q}$ is denoted by $\mathcal{M}(\mathcal{Q})$, and $|\mathfrak{m}|:=m_1m_2m_3$.
For $a\in\{1,2,3\}$, the set $\mathcal{N}_a$ consists of the even positive integers $s$ with
$hs\ge4$ if $\ell^a$ is an axis family, and with $hs\ge8$ if $\ell^a$ is a diagonal family.
\end{definition}

The set of shapes is a product:
\[
\mathcal{M}(\mathcal{Q})=\mathcal{N}_1\times\mathcal{N}_2\times\mathcal{N}_3 .
\]
To see this, let $m_1,m_2,m_3$ be even positive integers and put $p_a:=hm_a$, an even number.
By Definition~\ref{def:cell-structures}, the rectangular torus $\mathbb{T}(\mathfrak{F};p_1,p_2,p_3)$ is defined exactly when
$\left(p_a/|n_a|^2\right)n_a\in\Z^3$ and $p_a\ge4\,|n_a|^2$ for $a=1,2,3$.
The first condition holds, since $|n_a|^2\in\{1,2\}$ and $p_a$ is even.
Lemma~\ref{lem:systoles} shows that the lattice generated by the three integer vectors $\left(p_a/|n_a|^2\right)n_a$ has the systole $\min_ap_a/|n_a|^2$.
For the coordinate frame this lattice is $p_1\Z e_1+p_2\Z e_2+p_3\Z e_3$, with the systole $\min\{p_1,p_2,p_3\}$;
for the diagonal frame, whose normals are $e_-$, $e_+$ and $e_3$, it is $\Z\frac{p_1}{2}\,e_-+\Z\frac{p_2}{2}\,e_++\Z p_3\,e_3$, with the systole $\min\{p_1/2,p_2/2,p_3\}$.
The second condition is therefore the condition that the systole $\min_ap_a/|n_a|^2$ be at least $4$.
Since $|n_a|^2$ equals $1$ for an axis family and $2$ for a diagonal family,
it reads $hm_a\ge4$ for the axis families and $hm_a\ge8$ for the diagonal families among $\ell^1,\ell^2,\ell^3$;
that is, $m_a\in\mathcal{N}_a$ for $a=1,2,3$.

Let $\mathfrak{m}\in\mathcal{M}(\mathcal{Q})$ and $\mathbb{T}=\mathbb{T}(\mathfrak{F};h\mathfrak{m})$.
All the height periods $hm_a$ lie in $2h\Z$, and so $\mathbb{T}$ carries $\mathcal{Q}$ (Definition~\ref{def:cell-structures}).
The lattice of $\mathbb{T}$ has the pairwise orthogonal generators $\gamma_a:=\left(hm_a/|n_a|^2\right)n_a$, $a=1,2,3$,
of the lengths $hm_a/|n_a|$.
Its index $|\mathbb{T}|$ in $\Z^3$ is the volume of the box spanned by these generators,
and $|Q|=h^3\cdot\left(|n_1|^2\cdot|n_2|^2\cdot|n_3|^2\right)^{-\frac12}$ (Definition~\ref{def:cell-structures}); hence
\[
|\mathbb{T}|=\prod_{a=1}^{3}\frac{hm_a}{|n_a|}=\frac{h^3}{|n_1|\cdot|n_2|\cdot|n_3|}\cdot|\mathfrak{m}|=|Q|\cdot|\mathfrak{m}|,
\]
and $\mathbb{T}$ has $|\mathbb{T}|/|Q|=|\mathfrak{m}|$ cells.
For $G\in\mathcal{A}(Q_0)$, the \emph{folded partition function}
$Z_{\mathbb{T}(\mathfrak{F};h\mathfrak{m})}\left[\Fold_{\mathbb{T}(\mathfrak{F};h\mathfrak{m})}G\right]$
is a non-negative number by Theorem~\ref{thm:chessboard}(a).

\begin{proposition}[Mirror-periodic perturbations]\label{prop:periodic-perturbations}
Let $-1\le\rho\le1$, $w_0,\dots,w_3>0$, $M\in\R$ and $q>0$.
Let $\mathcal{Q}=(\mathfrak{F};o_1,o_2,o_3;h)$ be a cell structure and $G\in\mathcal{A}(Q_0)$.
\begin{enumerate}
\item\label{it:shape-trace}
Let $a\in\{1,2,3\}$, let $b$ and $c$ be the two other indices, and fix $m_b\in\mathcal{N}_b$ and $m_c\in\mathcal{N}_c$;
for $m_a\in\mathcal{N}_a$ let $\mathfrak{m}=(m_1,m_2,m_3)$.
There is a Hilbert--Schmidt operator $\mathfrak{Y}$ between two separable Hilbert spaces, not depending on $m_a$, such that
\[
Z_{\mathbb{T}(\mathfrak{F};h\mathfrak{m})}\left[\Fold_{\mathbb{T}(\mathfrak{F};h\mathfrak{m})}G\right]
=\Tr\left(\mathfrak{Y}^*\mathfrak{Y}\right)^{\frac{m_a}{2}}
\qquad\text{for all }m_a\in\mathcal{N}_a ;
\]
the operator $\mathfrak{Y}^*\mathfrak{Y}$ is non-negative and of trace class.
Hence the function
\[
m_a\longmapsto\left(Z_{\mathbb{T}(\mathfrak{F};h\mathfrak{m})}\left[\Fold_{\mathbb{T}(\mathfrak{F};h\mathfrak{m})}G\right]\right)^{\frac{1}{m_a}}
\]
is non-increasing on $\mathcal{N}_a$;
the folded partition function is either positive for all $m_a\in\mathcal{N}_a$ or zero for all $m_a\in\mathcal{N}_a$;
and in the first case the function
$m_a\mapsto\log Z_{\mathbb{T}(\mathfrak{F};h\mathfrak{m})}\left[\Fold_{\mathbb{T}(\mathfrak{F};h\mathfrak{m})}G\right]$
is convex on $\mathcal{N}_a$ in the sense of~\eqref{cond:convexity}.
\item\label{it:shape-limit}
If $Z_{\mathbb{T}(\mathfrak{F};h\mathfrak{m}^0)}\left[\Fold_{\mathbb{T}(\mathfrak{F};h\mathfrak{m}^0)}G\right]>0$ for one shape $\mathfrak{m}^0\in\mathcal{M}(\mathcal{Q})$,
then the folded partition function is positive for every $\mathfrak{m}\in\mathcal{M}(\mathcal{Q})$, and the function
\[
\mathfrak{m}\longmapsto|\mathfrak{m}|^{-1}\cdot\log Z_{\mathbb{T}(\mathfrak{F};h\mathfrak{m})}\left[\Fold_{\mathbb{T}(\mathfrak{F};h\mathfrak{m})}G\right]
\]
on $\mathcal{M}(\mathcal{Q})$ is non-increasing in each $m_a$, is bounded below,
and converges to its infimum, which is finite, as $\min_am_a\to\infty$ in any manner,
in the sense of Lemma~\ref{lem:separately-convex}\eqref{it:any-manner}.
\end{enumerate}
\end{proposition}

\begin{proof}
\eqref{it:shape-trace}
Fix $a$, $m_b$ and $m_c$ as in the statement, let $m_a\in\mathcal{N}_a$,
and abbreviate $\mathbb{T}:=\mathbb{T}(\mathfrak{F};h\mathfrak{m})$; let $\Gamma$ be its lattice.

\emph{The cross-section.}
The lattice $\Gamma$ is generated by the vectors $\gamma_{a'}=\left(hm_{a'}/|n_{a'}|^2\right)n_{a'}$, $a'=1,2,3$.
The normals of the frame are pairwise orthogonal,
so $\ell^a(\gamma_{a'})=n_a\cdot\gamma_{a'}$ equals $hm_a$ for $a'=a$ and vanishes for $a'\ne a$.
Hence $\ell^a(\Gamma)=hm_a\Z$, the vector $\gamma_a$ is a closing vector,
the height period is $p(\mathbb{T},\ell^a)=hm_a$, and $\Gamma_\parallel=\Gamma\cap\{\ell^a=0\}=\Z\gamma_b\oplus\Z\gamma_c$.
A rectangular torus is adapted to its three families (Definition~\ref{def:cell-structures});
thus $\mathbb{T}^0(\ell^a)=\mathbb{T}$ (Definition~\ref{def:slabs}), and the cross-section of $\mathbb{T}$ for $\ell^a$ is
\[
\Xi:=\Xi(\mathbb{T},\ell^a)=\left(\Z\gamma_b\oplus\Z\gamma_c;\,-1,-1\right),
\]
with the character $-1$ on $\gamma_b$ and on $\gamma_c$.
It depends on $m_b$ and $m_c$ only.

\emph{The content of one layer.}
Let $\mathcal{K}:=\{\kappa\in\Z^3:\ 0\le\kappa_{a'}<m_{a'}\ \text{for}\ a'=1,2,3\}$.
The cell indices of $\gamma_b$ and $\gamma_c$ (Definition~\ref{def:cell-families}) are $\kappa_{\gamma_b}=m_b\cdot e_b$ and $\kappa_{\gamma_c}=m_c\cdot e_c$,
with the unit vectors $e_b,e_c$ of $\Z^3$.
So the set of the $\kappa\in\mathcal{K}$ with $\kappa_a=0$ represents $\{\kappa\in\Z^3:\kappa_a=0\}$ modulo $\kappa_{\Gamma_\parallel}$,
and $\mathcal{K}$ is the set of representatives that Lemma~\ref{lem:one-direction}(b) forms from it, with $p/h=m_a$;
in particular $\mathcal{K}$ is a set of representatives of the cells of $\mathbb{T}$.
Put
\[
S:=\prod_{\kappa\in\mathcal{K},\ \kappa_a=0}G^{[\kappa]} ,
\]
the content of the layer $\kappa_a=0$; it depends on $m_b$ and $m_c$ only.
The cells $Q_\kappa$ with $\kappa_a=0$ lie in the open slab $\{o_a<\ell^a<o_a+h\}$,
the cells with indices in $\mathcal{K}$ lie in pairwise different $\Gamma$-orbits,
$G^{[\kappa]}\in\mathcal{A}(Q_\kappa)$ (Lemma~\ref{lem:admissible-images}\eqref{it:images-boxes}),
and $\ell^{a'}(\Gamma)=hm_{a'}\Z\subset2h\Z$ for $a'=1,2,3$.
By Lemma~\ref{lem:embedding-criteria}\eqref{it:embedding-cells},
every product in the expansion of $S$ embeds in $\Gamma$, hence in $\Gamma_\parallel$;
so $S\in\mathcal{A}_\Xi[o_a,o_a+h]$ (Definition~\ref{def:slabs}).

\emph{The fold.}
Since $\mathbb{T}$ carries $\mathcal{Q}$ and $G\in\mathcal{A}(Q_0)$, the relations~\eqref{eq:tori-chessboard-family} apply:
the chessboard family $\mathfrak{C}(G)=\left(G^{[\kappa]}\right)_{\kappa\in\Z^3}$ is a family on $(\mathbb{T},\mathcal{Q})$
in the sense of Definition~\ref{def:cell-families}, and $\mathfrak{R}^0_a\mathfrak{C}(G)=\mathfrak{C}(G)$.
Take the identity~\eqref{eq:tori-layer-fold} of the proof of Lemma~\ref{lem:one-direction}\eqref{it:straight-inequality}
for the family $\mathfrak{C}(G)$ on $(\mathbb{T},\mathcal{Q})$, with $\mathbb{T}^0=\mathbb{T}$, in the case $j_0=0$ and $i=1$;
then the set $\mathcal{K}'$ of that proof is $\mathcal{K}$, and $S_1=S$.
Together with Definition~\ref{def:folds}, the identity gives
\[
\fold_{\ell^a}S=\prod_{\kappa\in\mathcal{K}}\left(\mathfrak{R}^0_a\mathfrak{C}(G)\right)_\kappa=\prod_{\kappa\in\mathcal{K}}G^{[\kappa]}=\Fold_{\mathbb{T}}G .
\]

\emph{The trace.}
The planes $o_a+j\cdot h$, $j\in\Z$, lie in $\Adm(\ell^a)$, because $o_a\in\Adm(\ell^a)$ and $h$ is an integer;
moreover $p(\mathbb{T},\ell^a)=hm_a=2k\cdot h$ with $k:=m_a/2$.
Let
\[
\mathfrak{Y}:=\mathfrak{S}[S]:\ \mathcal{H}(o_a)\longrightarrow\mathcal{H}(o_a+h)
\]
be the slab operator of $S$ for the family $\ell^a$ and the cross-section $\Xi$,
a Hilbert--Schmidt operator between separable Hilbert spaces (Corollary~\ref{cor:slab-admissible}).
The slab spaces and the map $\mathfrak{S}$ are attached to $\ell^a$ and $\Xi$ alone;
as $\Xi$ and $S$ depend on $m_b$ and $m_c$ only, so does $\mathfrak{Y}$.
The identity~\eqref{eq:tori-fold-norm} of Proposition~\ref{prop:twisted-holder},
applied to the torus $\mathbb{T}=\mathbb{T}^0$, the family $\ell^a$, the closing vector $\gamma_a$,
$c_0=o_a$, $h'=h$, $2k=m_a$ and the content $G_1=S$ of the first slab,
together with Definition~\ref{def:even-schatten}, gives
\[
Z_{\mathbb{T}}\left[\Fold_{\mathbb{T}}G\right]
=Z_{\mathbb{T}}\left[\fold_{\ell^a}S\right]
=\|\mathfrak{Y}\|_{m_a}^{m_a}
=\Tr\left(\mathfrak{Y}^*\mathfrak{Y}\right)^{\frac{m_a}{2}} ,
\]
and $\mathfrak{Y}^*\mathfrak{Y}$ is non-negative and of trace class.
This is the identity asserted in~\eqref{it:shape-trace}.

\emph{Monotonicity, positivity, convexity.}
For $s\ge1$, the Schatten norm of exponent $s$ of the non-negative trace-class operator $\mathfrak{Y}^*\mathfrak{Y}$ is
$\left\|\mathfrak{Y}^*\mathfrak{Y}\right\|_s=\left(\Tr\left(\mathfrak{Y}^*\mathfrak{Y}\right)^s\right)^{\frac1s}$.
Therefore
\[
\left(Z_{\mathbb{T}}\left[\Fold_{\mathbb{T}}G\right]\right)^{\frac{1}{m_a}}
=\left\|\mathfrak{Y}^*\mathfrak{Y}\right\|_{\frac{m_a}{2}}^{\frac12},
\]
where $m_a/2\ge1$ for $m_a\in\mathcal{N}_a$.
By Lemma~\ref{lem:schatten-traces}(b), $s\mapsto\left\|\mathfrak{Y}^*\mathfrak{Y}\right\|_s$ is non-increasing on $[1,\infty)$,
and the first assertion follows.
If $\mathfrak{Y}^*\mathfrak{Y}=0$, the folded partition function vanishes for all $m_a\in\mathcal{N}_a$.
Otherwise $\mathfrak{Y}^*\mathfrak{Y}$ has a positive eigenvalue,
and for every $s\ge1$ the trace $\Tr\left(\mathfrak{Y}^*\mathfrak{Y}\right)^s$ is at least the $s$-th power of it;
so the folded partition function is positive for all $m_a\in\mathcal{N}_a$.
In that case
\[
\log Z_{\mathbb{T}}\left[\Fold_{\mathbb{T}}G\right]=\log\Tr\left(\mathfrak{Y}^*\mathfrak{Y}\right)^{\frac{m_a}{2}} ,
\]
and $s\mapsto\log\Tr\left(\mathfrak{Y}^*\mathfrak{Y}\right)^s$ is convex on $[1,\infty)$ by Lemma~\ref{lem:schatten-traces}(c).
Its composition with $m_a\mapsto m_a/2$ is convex on $[2,\infty)$, which contains $\mathcal{N}_a$,
and the restriction of a convex function to $\mathcal{N}_a$ satisfies~\eqref{cond:convexity}.

\eqref{it:shape-limit}
Let $\mathfrak{m}\in\mathcal{M}(\mathcal{Q})$.
The folded partition function is positive at $\mathfrak{m}^0$;
item~\eqref{it:shape-trace} in the direction $1$, with the entries $m^0_2,m^0_3$ fixed,
therefore shows that it is positive at $\left(m_1,m^0_2,m^0_3\right)$;
the same item in the direction $2$, with the entries $m_1,m^0_3$ fixed, shows that it is positive at $\left(m_1,m_2,m^0_3\right)$;
and in the direction $3$, with the entries $m_1,m_2$ fixed, at $\mathfrak{m}$.
All these triples are shapes, because $\mathcal{M}(\mathcal{Q})=\mathcal{N}_1\times\mathcal{N}_2\times\mathcal{N}_3$.
Consequently
\[
u(\mathfrak{m}):=\log Z_{\mathbb{T}(\mathfrak{F};h\mathfrak{m})}\left[\Fold_{\mathbb{T}(\mathfrak{F};h\mathfrak{m})}G\right]
\]
is a real function on $\mathcal{N}_1\times\mathcal{N}_2\times\mathcal{N}_3$,
and each $\mathcal{N}_a$ is a subset of $(0,\infty)$ that is unbounded above.
Fix $a$ and the two other entries.
As a function of $m_a\in\mathcal{N}_a$, the quotient $u/m_a$ is the logarithm of the positive function
that~\eqref{it:shape-trace} shows to be non-increasing on $\mathcal{N}_a$;
so $u/m_a$ is non-increasing, and $u$ satisfies~\eqref{cond:monotone-ratio}.
Moreover $u$ is convex on $\mathcal{N}_a$ in the sense of~\eqref{cond:convexity}, again by~\eqref{it:shape-trace}.
Lemma~\ref{lem:separately-convex}, applied to $u$ with two points chosen in each of the sets $\mathcal{N}_1,\mathcal{N}_2,\mathcal{N}_3$,
yields the remaining assertions:
$u/|\mathfrak{m}|$ is non-increasing in each $m_a$ by its item~\eqref{it:ratio-monotone},
it is bounded below by its item~\eqref{it:convex-floor},
and it converges to its finite infimum as $\min_am_a\to\infty$ in any manner by its item~\eqref{it:any-manner}.
\end{proof}

\subsection{Gaussian domination}\label{sec:domination-tori}

\begin{definition}[External fields]\label{def:external-fields}
Let $\mathbb{T}=(\Gamma,\chi)$ be a torus.
An external field on $\mathbb{T}$ is a real two-cochain $\beta\in C^2(\Gamma)$.
For such a field the function $F\mapsto\exp\left(-\langle\beta,F\rangle-\frac12\|\beta\|^2\right)$ of the curvature is Borel and non-negative,
and one puts
\[
\begin{aligned}
Z_{\mathbb{T}}(\beta)&:=Z_{\mathbb{T}}\left[\exp\left(-\langle\beta,F\rangle-\frac12\|\beta\|^2\right)\right]\\
&=\int_{\Conf(\Gamma)}\mathrm{d}(A,z)\,\exp\left(-\frac12\|F+\beta\|^2\right)\cdot\det D^{\mathbb{T}}_A\ \in[0,\infty],
\end{aligned}
\]
the functional being evaluated on non-negative Borel functions as in Definition~\ref{def:functional}.
The integrand on the right is non-negative by Proposition~\ref{prop:accretivity}\eqref{it:determinant-sign},
and the two expressions agree because $\|F\|^2+2\langle\beta,F\rangle+\|\beta\|^2=\|F+\beta\|^2$;
in particular $Z_{\mathbb{T}}(0)=Z_{\mathbb{T}}[1]$.

Let now $\mathbb{T}=C(n_1,n_2,n_3)$ be a coordinate torus,
with the direction $e_1$ as time and with the slice torus $\Lambda_{\mathrm{s}}=\Z^2/(n_2\Z\times n_3\Z)$ (Definition~\ref{def:slices}).
The electric and the magnetic part of an external field $\beta$ at the time $j\in\Z$ are formed as those of the curvature:
\[
\beta^{\mathrm{el}}_j(y,\nu):=\beta_{1\nu}(j,y)\quad(\nu=2,3),\qquad
\beta^{\mathrm{mag}}_j(y):=\beta_{23}(j,y),
\]
where $\beta_{\mu\nu}(x)$ is the value of $\beta$ at the plaquette $(x;\mu\nu)$
and $(j,y)\in\Z^3$ is the point with $x_1=j$ and $(x_2,x_3)=y$.
Thus $\beta^{\mathrm{el}}_j$ is a function on the links $(y,\nu)$ and $\beta^{\mathrm{mag}}_j$ a function on the plaquettes of $\Lambda_{\mathrm{s}}$,
and both depend on $j$ modulo $n_1$ only.
\end{definition}

\begin{lemma}[Hilbert--Schmidt property]\label{lem:shifted-hilbert-schmidt}
Let $-1\le\rho\le1$, $w_0,\dots,w_3>0$, $M\in\R$ and $q>0$.
Let $n_2,n_3\ge4$, and let $\mathcal{H}^{\mathrm{s}}$, $\mathcal{H}^{\mathrm{l}}$ and $K^{1/2}$ be the Hilbert spaces and the operator of the rectangular cross-section
with the sides $n_2,n_3$ and the plaquette $\sigma_0$ (Definitions~\ref{def:slices} and~\ref{def:slice-operators}).
For $\beta\in\R^{\Lambda^{(2)}_{\mathrm{s}}}$ let $\mathfrak{M}_{(\beta)}=\bigoplus_{k\in\Z}\mathfrak{M}_{k,(\beta)}\colon\mathcal{H}^{\mathrm{s}}\to\mathcal{H}^{\mathrm{l}}$
be the multiplication operator of Proposition~\ref{prop:external-field}, which acts in the flux sector $k$ by
\[
\mathfrak{M}_{k,(\beta)}(A)=\exp\left(-\frac14\left\|d_{\mathrm{s}}A+(2\pi/q)k\,1_{\sigma_0}+\beta\right\|^2\right)\cdot\widehat{\mathfrak{N}}(A).
\]
Then $X_\beta:=K^{1/2}\mathfrak{M}_{(\beta)}\colon\mathcal{H}^{\mathrm{s}}\to\mathcal{H}^{\mathrm{l}}$ is a Hilbert--Schmidt operator.
\end{lemma}

\begin{proof}
Unless another domain is indicated, integrals over $A$ and over $A'$ extend over the quotient $\R^{\Lambda^{(1)}_{\mathrm{s}}}/\Gamma_{\mathrm{s}}$ with the measure of Definition~\ref{def:slices},
and $\|\cdot\|_1$, $\|\cdot\|_2$, $\|\cdot\|_\infty$ are the norms of $L^1$, $L^2$ and $L^\infty$ of this quotient.
The argument is that of the proof of \cite[Lemma~10.4]{NoncompactQED}, which treats the weight given by the determinant of the Wilson operator;
there one Fock space and one Hilbert space take the place of $\mathcal{F}',\mathcal{F}$ and of $\mathcal{H}^{\mathrm{s}},\mathcal{H}^{\mathrm{l}}$.

For $a>0$ let $\tilde g_a$ be the periodised Gaussian
\[
\tilde g_a([A]):=\sum_{\gamma\in\Gamma_{\mathrm{s}}}\ \prod_{b\in\Lambda^{(1)}_{\mathrm{s}}}g_a\left(A(b)+\gamma(b)\right)
\qquad\left(A\in\R^{\Lambda^{(1)}_{\mathrm{s}}}\right),
\]
a function of the class $[A]$ of $A$ in the quotient.
By Definition~\ref{def:slice-operators}, the operator $\exp\left(-\frac{a}{2}|E|^2\right)$ acts in every flux sector as the convolution with $\tilde g_a$.
The function $\tilde g_{1/2}$ is non-negative, and its integral over the quotient equals $1$:
the integral of a periodisation over the quotient is the integral of the function itself over $\R^{\Lambda^{(1)}_{\mathrm{s}}}$,
and $g_{1/2}$ is a probability density on $\R$.
Since $\Gamma_{\mathrm{s}}\subset(2\pi/q)\Z^{\Lambda^{(1)}_{\mathrm{s}}}$, the series that defines $\tilde g_{1/2}$ is dominated by the sum over all of $(2\pi/q)\Z^{\Lambda^{(1)}_{\mathrm{s}}}$,
which factorises over the $2|\Lambda_{\mathrm{s}}|$ links; hence
\[
\sup\tilde g_{1/2}\le\left(\sup_{s\in\R}\sum_{i\in\Z}g_{1/2}\left(s+(2\pi/q)i\right)\right)^{2|\Lambda_{\mathrm{s}}|}<\infty,
\]
the inner supremum being finite by Lemma~\ref{lem:poisson-kernels}\eqref{it:gaussian-sums} with $\varepsilon=1$, as $g_{1/2}(s)=\pi^{-\frac12}\cdot\exp\left(-s^2\right)$.

The definition of $K^{1/2}$ and the formula for $P_{\mathrm{G}}$ in Definition~\ref{def:slice-operators} give
\[
\begin{aligned}
K^{1/2}&=(2\pi/q)^{\frac{|\Lambda_{\mathrm{s}}|}{2}}\cdot\exp\left(-\frac14|E|^2\right)P_{\mathrm{G}}\\
&=(2\pi/q)^{-\frac{|\Lambda_{\mathrm{s}}|}{2}}\int_{[0,2\pi/q)^{\Lambda_{\mathrm{s}}}}\mathrm{d}a\,\exp\left(-\frac14|E|^2\right)U_{\mathrm{G}}(a).
\end{aligned}
\]
Let $a\in[0,2\pi/q)^{\Lambda_{\mathrm{s}}}$, and let $\Phi$ be a vector of a flux sector of $\mathcal{H}^{\mathrm{l}}$.
By the convolution formula above, at the value $\frac12$ of the parameter of the periodised Gaussians, the operator $\exp\left(-\frac14|E|^2\right)$ is the convolution with $\tilde g_{1/2}$;
it acts on the functions only and therefore commutes with the constant matrix $u(a)$ on $\mathcal{F}$.
Together with $\left(U_{\mathrm{G}}(a)\Phi\right)(A)=u(a)\,\Phi(A-d_{\mathrm{s}}a)$ and the substitution $A\mapsto A+d_{\mathrm{s}}a$, which preserves the measure, this gives
\[
\left(\exp\left(-\frac14|E|^2\right)U_{\mathrm{G}}(a)\Phi\right)(A')=\int\mathrm{d}A\,\tilde g_{1/2}([A'-A-d_{\mathrm{s}}a])\cdot u(a)\,\Phi(A).
\]
Integrating over $a$, and exchanging the integrals over $a$ and over $A$ by Fubini's theorem
(the function $\tilde g_{1/2}$ is bounded and integrable, hence square integrable, like $\Phi$),
one finds that in every flux sector $K^{1/2}$ acts by integration against the kernel
\[
K^{1/2}(A',A):=(2\pi/q)^{-\frac{|\Lambda_{\mathrm{s}}|}{2}}\int_{[0,2\pi/q)^{\Lambda_{\mathrm{s}}}}\mathrm{d}a\,\tilde g_{1/2}([A'-A-d_{\mathrm{s}}a])\cdot u(a)\in\End\mathcal{F}.
\]
Since $\tilde g_{1/2}$ is non-negative and every $u(a)$ is unitary, $\left\|K^{1/2}(A',A)\right\|\le\hat g([A'-A])$ for the function
\[
\hat g([A]):=(2\pi/q)^{-\frac{|\Lambda_{\mathrm{s}}|}{2}}\int_{[0,2\pi/q)^{\Lambda_{\mathrm{s}}}}\mathrm{d}a\,\tilde g_{1/2}([A-d_{\mathrm{s}}a]).
\]
By Tonelli's theorem and the invariance of the measure under translations, the value $1$ of the integral of $\tilde g_{1/2}$ yields
$\|\hat g\|_1=(2\pi/q)^{-\frac{|\Lambda_{\mathrm{s}}|}{2}}\cdot(2\pi/q)^{|\Lambda_{\mathrm{s}}|}=(2\pi/q)^{\frac{|\Lambda_{\mathrm{s}}|}{2}}$.
Moreover $\hat g\le(2\pi/q)^{\frac{|\Lambda_{\mathrm{s}}|}{2}}\cdot\sup\tilde g_{1/2}$, so that $\hat g$ is bounded and
\[
\|\hat g\|_2^2\le\|\hat g\|_\infty\cdot\|\hat g\|_1<\infty .
\]

In the flux sector $k$, the operator $X_\beta$ acts by the kernel $K^{1/2}(A',A)\,\mathfrak{M}_{k,(\beta)}(A)$, a linear map from $\mathcal{F}'$ to $\mathcal{F}$.
The number
\[
C:=\sup_A\left\|\widehat{\mathfrak{N}}(A)\right\|
\]
is finite by Lemma~\ref{lem:half-slab-matrices}\eqref{it:half-slab-entries},
and $\left\|\mathfrak{M}_{k,(\beta)}(A)\right\|\le C\cdot\exp\left(-\frac14\left\|d_{\mathrm{s}}A+(2\pi/q)k\,1_{\sigma_0}+\beta\right\|^2\right)$.
Summing over an orthonormal basis of $\mathcal{F}'$ shows that the square of the Hilbert--Schmidt norm of a linear map from $\mathcal{F}'$ to $\mathcal{F}$
is at most $\dim\mathcal{F}'$ times the square of its operator norm.
Integrating first over $A'$, with the invariance of the measure under translations, one obtains
\[
\begin{aligned}
&\sum_{k\in\Z}\int\mathrm{d}A\int\mathrm{d}A'\,\left\|K^{1/2}(A',A)\,\mathfrak{M}_{k,(\beta)}(A)\right\|_{\mathrm{HS}}^2\\
&\qquad\le\dim\mathcal{F}'\cdot\|\hat g\|_2^2\cdot C^2\cdot\sum_{k\in\Z}\int\mathrm{d}A\,\exp\left(-\frac12\left\|d_{\mathrm{s}}A+(2\pi/q)k\,1_{\sigma_0}+\beta\right\|^2\right).
\end{aligned}
\]
The left side is the square of the $L^2$ norm of the kernel of $X_\beta$.
By Lemma~\ref{lem:poisson-kernels}\eqref{it:hilbert-schmidt-kernels}, applied to the measure spaces formed by the flux index, the quotient $\R^{\Lambda^{(1)}_{\mathrm{s}}}/\Gamma_{\mathrm{s}}$
and the index of the orthonormal basis $(e_I)$ of $\mathcal{F}'$, respectively of $\mathcal{F}$, with the counting measure on the discrete factors,
$X_\beta$ is a Hilbert--Schmidt operator as soon as the left side is finite, and $\|X_\beta\|_{\mathrm{HS}}^2$ is then equal to it.
It remains to show that the sum on the right is finite.

Being defined by linear equations with integer coefficients, the subspace $\ker d_{\mathrm{s}}$ of $\R^{\Lambda^{(1)}_{\mathrm{s}}}$ has a basis of integer vectors;
so $\Gamma_{\mathrm{s}}=(2\pi/q)\left(\Z^{\Lambda^{(1)}_{\mathrm{s}}}\cap\ker d_{\mathrm{s}}\right)$ is a lattice of full rank in $\ker d_{\mathrm{s}}$.
As $\Gamma_{\mathrm{s}}$ lies in $\ker d_{\mathrm{s}}$, the orthogonal decomposition $\R^{\Lambda^{(1)}_{\mathrm{s}}}=\ker d_{\mathrm{s}}\oplus(\ker d_{\mathrm{s}})^{\perp}$
identifies the quotient $\R^{\Lambda^{(1)}_{\mathrm{s}}}/\Gamma_{\mathrm{s}}$, with its measure,
with $(\ker d_{\mathrm{s}}/\Gamma_{\mathrm{s}})\times(\ker d_{\mathrm{s}})^{\perp}$, where the first factor is a compact torus and both factors carry their Lebesgue measures;
on the second factor $d_{\mathrm{s}}$ is injective.
Let $1_{\Lambda^{(2)}_{\mathrm{s}}}$ be the constant function $1$ on the plaquettes of $\Lambda_{\mathrm{s}}$, so that $|\Lambda_{\mathrm{s}}|^{-\frac12}\cdot1_{\Lambda^{(2)}_{\mathrm{s}}}$ is a unit vector.
Each link of $\Lambda_{\mathrm{s}}$ lies in exactly two plaquettes, whose boundaries traverse it in opposite directions;
in the formula for $d_{\mathrm{s}}$ of Definition~\ref{def:slices} it therefore occurs once with the sign $+$ and once with the sign $-$,
so that $\sum_\sigma(d_{\mathrm{s}}A)(\sigma)=0$, the sum extending over the plaquettes $\sigma$ of $\Lambda_{\mathrm{s}}$.
Thus the range $\operatorname{ran}d_{\mathrm{s}}$ is orthogonal to $1_{\Lambda^{(2)}_{\mathrm{s}}}$, and for all $k\in\Z$ and $A\in\R^{\Lambda^{(1)}_{\mathrm{s}}}$
\[
\left\langle d_{\mathrm{s}}A+(2\pi/q)k\,1_{\sigma_0}+\beta,\ |\Lambda_{\mathrm{s}}|^{-\frac12}\cdot1_{\Lambda^{(2)}_{\mathrm{s}}}\right\rangle
=|\Lambda_{\mathrm{s}}|^{-\frac12}\cdot\left((2\pi/q)k+\sum_\sigma\beta(\sigma)\right).
\]
Let $k\in\Z$.
Decompose $(2\pi/q)k\,1_{\sigma_0}+\beta$ into a vector of $\operatorname{ran}d_{\mathrm{s}}$ and a vector orthogonal to $\operatorname{ran}d_{\mathrm{s}}$;
as $d_{\mathrm{s}}$ maps $(\ker d_{\mathrm{s}})^{\perp}$ bijectively onto its range, the first vector is $d_{\mathrm{s}}\tilde A$ for exactly one $\tilde A\in(\ker d_{\mathrm{s}})^{\perp}$,
which depends on $k$ and $\beta$.
The orthogonal projection of $d_{\mathrm{s}}A+(2\pi/q)k\,1_{\sigma_0}+\beta$ onto $\operatorname{ran}d_{\mathrm{s}}$ is then $d_{\mathrm{s}}A+d_{\mathrm{s}}\tilde A$,
and since $\operatorname{ran}d_{\mathrm{s}}$ is orthogonal to the line spanned by $1_{\Lambda^{(2)}_{\mathrm{s}}}$, Pythagoras' theorem gives
\[
\left\|d_{\mathrm{s}}A+(2\pi/q)k\,1_{\sigma_0}+\beta\right\|^2
\ge\frac{1}{|\Lambda_{\mathrm{s}}|}\cdot\left((2\pi/q)k+\sum_\sigma\beta(\sigma)\right)^2+\left\|d_{\mathrm{s}}A+d_{\mathrm{s}}\tilde A\right\|^2 .
\]
In the last term $A$ enters only through $d_{\mathrm{s}}A$, which is determined by the component of $A$ in $(\ker d_{\mathrm{s}})^{\perp}$,
and the substitution $A\mapsto A-\tilde A$ in $(\ker d_{\mathrm{s}})^{\perp}$ preserves Lebesgue measure; hence
\[
\int_{(\ker d_{\mathrm{s}})^{\perp}}\mathrm{d}A\,\exp\left(-\frac12\left\|d_{\mathrm{s}}A+d_{\mathrm{s}}\tilde A\right\|^2\right)
=\int_{(\ker d_{\mathrm{s}})^{\perp}}\mathrm{d}A\,\exp\left(-\frac12\left\|d_{\mathrm{s}}A\right\|^2\right)=:C_1 .
\]
The number $C_1$ depends neither on $k$ nor on $\beta$,
and it is finite, $A\mapsto\|d_{\mathrm{s}}A\|^2$ being a positive definite quadratic form on $(\ker d_{\mathrm{s}})^{\perp}$.
Integrating over the product of $\ker d_{\mathrm{s}}/\Gamma_{\mathrm{s}}$ and $(\ker d_{\mathrm{s}})^{\perp}$ and summing over $k$, one arrives at
\[
\begin{aligned}
&\sum_{k\in\Z}\int\mathrm{d}A\,\exp\left(-\frac12\left\|d_{\mathrm{s}}A+(2\pi/q)k\,1_{\sigma_0}+\beta\right\|^2\right)\\
&\qquad\le\operatorname{vol}\left(\ker d_{\mathrm{s}}/\Gamma_{\mathrm{s}}\right)\cdot C_1\cdot\sum_{k\in\Z}\exp\left(-\frac{\left((2\pi/q)k+\sum_\sigma\beta(\sigma)\right)^2}{2|\Lambda_{\mathrm{s}}|}\right)<\infty,
\end{aligned}
\]
the Gaussian series in $k$ being convergent.
So the kernel of $X_\beta$ is square integrable, and $X_\beta$ is a Hilbert--Schmidt operator.
\end{proof}

\begin{lemma}[Invariance under lattice isometries]\label{lem:isometry-invariance}
Let $-1\le\rho\le1$, $w_0,\dots,w_3>0$, $M\in\R$ and $q>0$.
Let $g\in W$, let $\mathbb{T}=(\Gamma,\chi)$ be a torus, with the image torus $g\mathbb{T}$ and its period lattice $g\Gamma$ (Definition~\ref{def:tori}),
and let $G\ge0$ be a Borel function on $C^2(\Gamma)$.
Then $G\circ g^{*}$ is a non-negative Borel function on $C^2(g\Gamma)$, and
\[
Z_{g\mathbb{T}}\left[G\circ g^{*}\right]=Z_{\mathbb{T}}[G],
\]
both functionals being evaluated on non-negative Borel functions of the curvature (Definition~\ref{def:functional}).
In particular, for every $\beta\in C^2(\Gamma)$,
\[
\left\|(g^{-1})^{*}\beta\right\|=\|\beta\|
\qquad\text{and}\qquad
Z_{g\mathbb{T}}\left((g^{-1})^{*}\beta\right)=Z_{\mathbb{T}}(\beta).
\]
\end{lemma}

\begin{proof}
By Definition~\ref{def:configurations}, the pull-back $(A',z')\mapsto(g^{*}A',g^{*}z')$ is a measure preserving bijection from $\Conf(g\Gamma)$ onto $\Conf(\Gamma)$.
The map $g^{*}$ commutes with $d$ and is orthogonal (Definition~\ref{def:cochains}).
Hence the curvature of $(g^{*}A',g^{*}z')$ is $g^{*}F'$, where $F'=dA'+\frac{2\pi}{q}\cdot z'$ is the curvature of $(A',z')$, and $\|g^{*}F'\|=\|F'\|$.
Moreover $(g^{-1})^{*}g^{*}=(gg^{-1})^{*}$ is the identity (Definition~\ref{def:cochains}),
so that Proposition~\ref{prop:symmetries}\eqref{it:lattice-covariance}, applied to the $\Gamma$-periodic one-cochain $A=g^{*}A'$,
gives $\det D^{g\mathbb{T}}_{A'}=\det D^{\mathbb{T}}_{g^{*}A'}$.
The change of variables $(A,z)=(g^{*}A',g^{*}z')$ therefore yields the identity
\[
\begin{aligned}
Z_{g\mathbb{T}}\left[G\circ g^{*}\right]
&=\int_{\Conf(g\Gamma)}\mathrm{d}(A',z')\,\exp\left(-\frac12\|F'\|^2\right)\cdot\det D^{g\mathbb{T}}_{A'}\cdot G(g^{*}F')\\
&=\int_{\Conf(\Gamma)}\mathrm{d}(A,z)\,\exp\left(-\frac12\|F\|^2\right)\cdot\det D^{\mathbb{T}}_A\cdot G(F)\\
&=Z_{\mathbb{T}}[G]
\end{aligned}
\]
between numbers in $[0,\infty]$.
The adjoint of the orthogonal map $g^{*}\colon C^2(g\Gamma)\to C^2(\Gamma)$ is its inverse $(g^{-1})^{*}$.
So $(g^{-1})^{*}$ is orthogonal as well, $\|(g^{-1})^{*}\beta\|=\|\beta\|$, and
$\langle\beta,g^{*}F'\rangle=\langle(g^{-1})^{*}\beta,F'\rangle$ for all $F'\in C^2(g\Gamma)$.
For $G=\exp\left(-\langle\beta,\cdot\,\rangle-\frac12\|\beta\|^2\right)$ the function $G\circ g^{*}$ is therefore
$\exp\left(-\langle(g^{-1})^{*}\beta,\cdot\,\rangle-\frac12\|(g^{-1})^{*}\beta\|^2\right)$,
and the identity above becomes $Z_{g\mathbb{T}}\left((g^{-1})^{*}\beta\right)=Z_{\mathbb{T}}(\beta)$.
\end{proof}

\begin{theorem}[Gaussian domination on coordinate tori]\label{thm:coordinate-domination}
Let $-1\le\rho\le1$, $w_0,\dots,w_3>0$, $M\in\R$ and $q>0$.
Let $n_1,n_2,n_3\ge4$ be integers,
let $\mathbb{T}:=C(n_1,n_2,n_3)$ be the coordinate torus with the period lattice $\Gamma=n_1\Z e_1+n_2\Z e_2+n_3\Z e_3$,
and let $\mathbb{E}:=\mathbb{E}_{\mathbb{T}}$.
For every external field $\beta\in C^2(\Gamma)$,
\[
0<Z_{\mathbb{T}}(\beta)\le Z_{\mathbb{T}}[1].
\]
Consequently the following hold.
\begin{enumerate}
\item
For every $\beta\in C^2(\Gamma)$ one has $\mathbb{E}\exp\langle\beta,F\rangle\le\exp\left(\frac12\|\beta\|^2\right)$.
For $\beta_1,\dots,\beta_s\in C^2(\Gamma)$ and complex numbers $z_1,\dots,z_s$,
\[
\left|\mathbb{E}\exp\left(\sum_{i=1}^sz_i\cdot\langle\beta_i,F\rangle\right)\right|\le\exp\left(\frac12\left\|\sum_{i=1}^s(\Re z_i)\cdot\beta_i\right\|^2\right).
\]
\item
For every $\beta\in C^2(\Gamma)$ and every integer $r\ge1$,
\begin{gather*}
\mathbb{E}\langle\beta,F\rangle=0,\qquad
\mathbb{E}\langle\beta,F\rangle^2\le\|\beta\|^2,\\
\mathbb{E}\langle\beta,F\rangle^{2r}\le\left(2r\cdot\exp(1)\cdot\|\beta\|^2\right)^r,\qquad
\mathbb{E}\langle\beta,F\rangle^4\le\frac32\exp(2)\cdot\|\beta\|^4 .
\end{gather*}
For $\beta_1,\dots,\beta_4\in C^2(\Gamma)$ the joint fourth cumulant under $\mathbb{E}$ satisfies
\[
\left|\kappa_4\left(\langle\beta_1,F\rangle,\dots,\langle\beta_4,F\rangle\right)\right|\le15\prod_{i=1}^4\|\beta_i\| .
\]
\item
Let $\beta_1,\dots,\beta_s\in C^2(\Gamma)$ be non-zero with pairwise disjoint supports, and let $\xi_1,\dots,\xi_s\ge0$.
Then the event that $|\langle\beta_i,F\rangle|\ge\xi_i\cdot\|\beta_i\|$ for $i=1,\dots,s$ has probability at most
\[
\prod_{i=1}^s2\exp\left(-\frac12\xi_i^2\right)
\]
under $\mathbb{E}$.
\item
The law $\mu_{\mathbb{T}}$ satisfies the inequality of condition~\eqref{cond:gaussian-domination},
\[
\int\mathrm{d}\mu_{\mathbb{T}}(F)\,\exp\langle\beta,F\rangle\le\exp\left(\frac12\|\beta\|^2\right),
\]
for every $\beta\in C^2_{\mathrm{c}}$ whose support lies in a box that embeds in $\Gamma$ (Definition~\ref{def:integrands}).
\end{enumerate}
\end{theorem}

\begin{proof}
The main inequality and items (a)--(c) are proved as in \cite[Theorem~10.7 and Corollary~10.8]{NoncompactQED},
where the coordinate torus is written $\Lambda_n$, the external field $h$,
and the fermion weight is the determinant of the Wilson operator instead of $\det D^{\mathbb{T}}_A$.

\emph{Positivity.}
On $\Conf(\Gamma)$ the function $\exp\left(-\frac12\|F+\beta\|^2\right)$ is strictly positive,
and $\det D^{\mathbb{T}}_A\ge0$ by Proposition~\ref{prop:accretivity}\eqref{it:determinant-sign}.
By Theorem~\ref{thm:partition-positivity},
\[
Z_{\mathbb{T}}[1]=\int_{\Conf(\Gamma)}\mathrm{d}(A,z)\,\exp\left(-\frac12\|F\|^2\right)\cdot\det D^{\mathbb{T}}_A>0 ,
\]
so the set of the configurations at which $\det D^{\mathbb{T}}_A>0$ has positive measure.
On this set the integrand of $Z_{\mathbb{T}}(\beta)$ is strictly positive, and therefore $Z_{\mathbb{T}}(\beta)>0$.
(The zero set of $A\mapsto\det D^{\mathbb{T}}_A$ is even a null set:
Theorem~\ref{thm:floor}(e) for $\rho^2\neq1$, and Theorem~\ref{thm:closed-interval} for $-1\le\rho\le1$.)

\emph{Reduction to magnetic fields.}
Take $e_1$ as time, and let $|n|=n_1n_2n_3$ (Definition~\ref{def:tori}).
For $j\in\Z/n_1\Z$ consider the operators
\[
X_j:=K^{1/2}\mathfrak{M}_{(\beta^{\mathrm{mag}}_j)}\colon\mathcal{H}^{\mathrm{s}}\to\mathcal{H}^{\mathrm{l}},
\qquad
U_j:=\exp\left(i\langle\beta^{\mathrm{el}}_j,E\rangle\right)
\]
of the cross-section with the sides $n_2,n_3$ (Lemma~\ref{lem:shifted-hilbert-schmidt} and Definition~\ref{def:slice-operators});
$U_j$ is the unitary operator on $\mathcal{H}^{\mathrm{l}}$ that acts in every flux sector by $(U_j\Phi)(A)=\Phi(A+\beta^{\mathrm{el}}_j)$.
The left side of the identity of Proposition~\ref{prop:external-field} is $Z_{\mathbb{T}}(\beta)$ (Definitions~\ref{def:functional} and~\ref{def:external-fields}),
and that identity reads
\[
(2\pi)^{-|n|}\cdot Z_{\mathbb{T}}(\beta)
=\Tr_{\mathcal{H}^{\mathrm{s}}}\left(X_0^{*}U_{n_1-1}X_{n_1-1}\,X_{n_1-1}^{*}U_{n_1-2}X_{n_1-2}\cdots X_1^{*}U_0X_0\right).
\]
By Lemma~\ref{lem:shifted-hilbert-schmidt}, every $X_j$ is a Hilbert--Schmidt operator.
For Hilbert--Schmidt operators $\mathfrak{Y}\colon\mathcal{H}^{\mathrm{l}}\to\mathcal{H}^{\mathrm{s}}$ and $\mathfrak{Y}'\colon\mathcal{H}^{\mathrm{s}}\to\mathcal{H}^{\mathrm{l}}$ one has
\[
\Tr_{\mathcal{H}^{\mathrm{s}}}\left(\mathfrak{Y}\mathfrak{Y}'\right)=\Tr_{\mathcal{H}^{\mathrm{l}}}\left(\mathfrak{Y}'\mathfrak{Y}\right):
\]
in orthonormal bases of $\mathcal{H}^{\mathrm{s}}$ and of $\mathcal{H}^{\mathrm{l}}$ both traces are the sum of one double series,
whose terms are the products of a matrix element of $\mathfrak{Y}$ with one of $\mathfrak{Y}'$,
and this series converges absolutely by the Cauchy--Schwarz inequality,
since the sums of the squared moduli of the matrix elements are the squared Hilbert--Schmidt norms.
Take $\mathfrak{Y}:=X_0^{*}$, and let $\mathfrak{Y}'$ be the product of the remaining factors under the trace above;
it is a Hilbert--Schmidt operator, because its right-most factor $X_0$ is one.
This gives
\[
(2\pi)^{-|n|}\cdot Z_{\mathbb{T}}(\beta)
=\Tr_{\mathcal{H}^{\mathrm{l}}}\left(U_{n_1-1}X_{n_1-1}X_{n_1-1}^{*}\;U_{n_1-2}X_{n_1-2}X_{n_1-2}^{*}\cdots U_0X_0X_0^{*}\right),
\]
where every $X_jX_j^{*}$ is a non-negative operator of trace class on $\mathcal{H}^{\mathrm{l}}$;
as $Z_{\mathbb{T}}(\beta)\ge0$, the left side is the modulus of the trace on the right.
Operators of trace class belong to every Schatten class (Lemma~\ref{lem:schatten-traces}(b)).
H\"older's inequality for traces, Lemma~\ref{lem:schatten-traces}(a), applied to the $n_1$ operators $U_jX_jX_j^{*}$ with all exponents equal to $n_1$,
together with the invariance of the Schatten norms under unitary factors, Lemma~\ref{lem:schatten-traces}(b), therefore gives
\[
(2\pi)^{-|n|}\cdot Z_{\mathbb{T}}(\beta)
\le\prod_{j\in\Z/n_1\Z}\left\|U_jX_jX_j^{*}\right\|_{n_1}
=\prod_{j\in\Z/n_1\Z}\left\|X_jX_j^{*}\right\|_{n_1}
=\prod_{j\in\Z/n_1\Z}\left(\Tr_{\mathcal{H}^{\mathrm{s}}}\left(X_j^{*}X_j\right)^{n_1}\right)^{\frac{1}{n_1}} .
\]
In the last step, $\|X_jX_j^{*}\|_{n_1}^{n_1}=\Tr_{\mathcal{H}^{\mathrm{l}}}\left(X_jX_j^{*}\right)^{n_1}$ because $X_jX_j^{*}\ge0$,
and the operators $X_jX_j^{*}$ and $X_j^{*}X_j$ have the same non-zero eigenvalues, with the same multiplicities.

For $j\in\Z$ let $\beta^{(j)}\in C^2(\Gamma)$ be the external field whose electric parts vanish and whose magnetic part at every time is $\beta^{\mathrm{mag}}_j$:
\[
\beta^{(j)}\bigl((j',y);1\nu\bigr):=0,\qquad
\beta^{(j)}\bigl((j',y);23\bigr):=\beta\bigl((j,y);23\bigr)
\qquad\left(j'\in\Z,\ y\in\Z^2,\ \nu=2,3\right).
\]
It depends on $j$ modulo $n_1$ only, and it vanishes off the plaquettes of the form $(x;23)$.
The unitary operators that Proposition~\ref{prop:external-field} attaches to $\beta^{(j)}$ are all equal to $\exp\left(i\langle0,E\rangle\right)=1$,
and each of the operators $K^{1/2}\mathfrak{M}_{((\beta^{(j)})^{\mathrm{mag}}_{j'})}$ attached to it equals $X_j$.
That proposition therefore gives $(2\pi)^{-|n|}\cdot Z_{\mathbb{T}}\left(\beta^{(j)}\right)=\Tr_{\mathcal{H}^{\mathrm{s}}}\left(X_j^{*}X_j\right)^{n_1}$,
and the preceding inequality becomes
\begin{equation}\label{eq:tori-magnetic}
Z_{\mathbb{T}}(\beta)\le\prod_{j\in\Z/n_1\Z}Z_{\mathbb{T}}\left(\beta^{(j)}\right)^{\frac{1}{n_1}} .
\end{equation}
If all magnetic parts of $\beta$ vanish, then $\beta^{(j)}=0$ for every $j$, and \eqref{eq:tori-magnetic} reads $Z_{\mathbb{T}}(\beta)\le Z_{\mathbb{T}}(0)$.
Both inequalities hold on every coordinate torus whose three sides are at least $4$, with $e_1$ as time.

\emph{The second time direction.}
Let $\beta'\in C^2(\Gamma)$ vanish off the plaquettes of the form $(x;23)$,
and let $g\in\Ocub\subset W$ be the transposition of $e_1$ and $e_2$.
Its linear part maps $n_1e_1$ to $n_1e_2$ and $n_2e_2$ to $n_2e_1$, and it fixes $n_3e_3$.
So $g\Gamma=n_2\Z e_1+n_1\Z e_2+n_3\Z e_3$, the character of $g\mathbb{T}$ takes the value $-1$ at these three generators,
and $g\mathbb{T}=C(n_2,n_1,n_3)$ (Definition~\ref{def:tori}).
Since $g^{-1}=g$, the field $\beta'':=(g^{-1})^{*}\beta'=g^{*}\beta'\in C^2(g\Gamma)$ vanishes at every plaquette $\sigma$ for which $g\sigma$ is not of the form $(x;23)$,
that is, off the plaquettes of the form $(x;13)$.
With $e_1$ as time on $g\mathbb{T}$, all magnetic parts of $\beta''$ therefore vanish.
Lemma~\ref{lem:isometry-invariance}, the inequality $Z_{g\mathbb{T}}(\beta'')\le Z_{g\mathbb{T}}(0)$ of the preceding step on the coordinate torus $g\mathbb{T}$,
and Lemma~\ref{lem:isometry-invariance} once more, for the field $0$, give
\[
Z_{\mathbb{T}}(\beta')=Z_{g\mathbb{T}}(\beta'')\le Z_{g\mathbb{T}}(0)=Z_{\mathbb{T}}(0).
\]
Applied with $\beta'=\beta^{(j)}$ for every $j\in\Z/n_1\Z$, this turns \eqref{eq:tori-magnetic} into
\[
Z_{\mathbb{T}}(\beta)\le\prod_{j\in\Z/n_1\Z}Z_{\mathbb{T}}(0)^{\frac{1}{n_1}}=Z_{\mathbb{T}}(0)=Z_{\mathbb{T}}[1].
\]

\emph{Item (a).}
Replacing $\beta$ by $-\beta$ in Definition~\ref{def:external-fields} gives
$Z_{\mathbb{T}}(-\beta)=\exp\left(-\frac12\|\beta\|^2\right)\cdot Z_{\mathbb{T}}\left[\exp\langle\beta,F\rangle\right]$, and hence
\[
\mathbb{E}\exp\langle\beta,F\rangle=\exp\left(\frac12\|\beta\|^2\right)\cdot\frac{Z_{\mathbb{T}}(-\beta)}{Z_{\mathbb{T}}(0)}\le\exp\left(\frac12\|\beta\|^2\right).
\]
As $\mathbb{E}$ is the integral with respect to a probability measure and the numbers $\langle\beta_i,F\rangle$ are real,
\[
\left|\mathbb{E}\exp\left(\sum_{i=1}^sz_i\cdot\langle\beta_i,F\rangle\right)\right|
\le\mathbb{E}\exp\left(\sum_{i=1}^s(\Re z_i)\cdot\langle\beta_i,F\rangle\right)
=\mathbb{E}\exp\left\langle\sum_{i=1}^s(\Re z_i)\cdot\beta_i,F\right\rangle ,
\]
and the first assertion, applied to $\sum_i(\Re z_i)\cdot\beta_i$, bounds the right side.

\emph{Item (b).}
Since $\exp|\langle\beta,F\rangle|\le\exp\langle\beta,F\rangle+\exp\langle-\beta,F\rangle$, item (a), applied to $\beta$ and to $-\beta$, gives $\mathbb{E}\exp|\langle\beta,F\rangle|\le2\exp\left(\frac12\|\beta\|^2\right)$;
as $|\langle\beta,F\rangle|^{2r}\le(2r)!\cdot\exp|\langle\beta,F\rangle|$ for every integer $r\ge1$, all moments of $\langle\beta,F\rangle$ are finite.
The law $\mu_{\mathbb{T}}$ is invariant under $F\mapsto-F$ (Proposition~\ref{prop:symmetries}\eqref{it:charge-conjugation}).
Therefore $\mathbb{E}\langle\beta,F\rangle=0$, and (a), applied to $s\cdot\beta$, gives for every real $s$
\[
\mathbb{E}\cosh\left(s\cdot\langle\beta,F\rangle\right)=\mathbb{E}\exp\left(s\cdot\langle\beta,F\rangle\right)\le\exp\left(\frac12s^2\cdot\|\beta\|^2\right).
\]
From $\cosh\left(s\cdot\langle\beta,F\rangle\right)\ge1+\frac12s^2\cdot\langle\beta,F\rangle^2$ it follows that
$\mathbb{E}\langle\beta,F\rangle^2\le2s^{-2}\cdot\left(\exp\left(\frac12s^2\cdot\|\beta\|^2\right)-1\right)$ for every $s\neq0$,
and the right side tends to $\|\beta\|^2$ as $s\to0$.
From $s^{2r}\cdot\langle\beta,F\rangle^{2r}\le(2r)!\cdot\cosh\left(s\cdot\langle\beta,F\rangle\right)$ it follows that
$\mathbb{E}\langle\beta,F\rangle^{2r}\le(2r)!\cdot s^{-2r}\cdot\exp\left(\frac12s^2\cdot\|\beta\|^2\right)$ for every $s\neq0$.
For $\beta=0$ the bounds on the moments hold trivially.
For $\beta\neq0$, the choice $s^2=2r/\|\beta\|^2$ and the inequality $(2r)!\le(2r)^{2r}$ give
\[
\mathbb{E}\langle\beta,F\rangle^{2r}
\le(2r)!\cdot\left(\frac{\|\beta\|^2}{2r}\right)^r\cdot\exp(r)
\le(2r)^{2r}\cdot\left(\frac{\|\beta\|^2}{2r}\right)^r\cdot\exp(r)
=\left(2r\cdot\exp(1)\cdot\|\beta\|^2\right)^r ,
\]
and for $r=2$ the first bound of this chain, with $4!=24$, reads
$\mathbb{E}\langle\beta,F\rangle^4\le24\cdot\left(\frac14\|\beta\|^2\right)^2\cdot\exp(2)=\frac32\exp(2)\cdot\|\beta\|^4$.
The variables $\langle\beta_1,F\rangle,\dots,\langle\beta_4,F\rangle$ are centred, so that their joint fourth cumulant is
\begin{multline*}
\kappa_4\left(\langle\beta_1,F\rangle,\dots,\langle\beta_4,F\rangle\right)
=\mathbb{E}\left[\langle\beta_1,F\rangle\cdot\langle\beta_2,F\rangle\cdot\langle\beta_3,F\rangle\cdot\langle\beta_4,F\rangle\right]\\
-\mathbb{E}\left[\langle\beta_1,F\rangle\cdot\langle\beta_2,F\rangle\right]\cdot\mathbb{E}\left[\langle\beta_3,F\rangle\cdot\langle\beta_4,F\rangle\right]\\
-\mathbb{E}\left[\langle\beta_1,F\rangle\cdot\langle\beta_3,F\rangle\right]\cdot\mathbb{E}\left[\langle\beta_2,F\rangle\cdot\langle\beta_4,F\rangle\right]\\
-\mathbb{E}\left[\langle\beta_1,F\rangle\cdot\langle\beta_4,F\rangle\right]\cdot\mathbb{E}\left[\langle\beta_2,F\rangle\cdot\langle\beta_3,F\rangle\right].
\end{multline*}
By H\"older's inequality and the bound on the fourth moments, the first term has modulus at most $\frac32\exp(2)\cdot\prod_{i=1}^4\|\beta_i\|$;
by the Cauchy--Schwarz inequality and the bound on the second moments, each of the three products has modulus at most $\prod_{i=1}^4\|\beta_i\|$.
As $\exp(2)<8$, one has $\frac32\exp(2)+3<15$, which gives the bound on $\kappa_4$.

\emph{Item (c).}
Fix one of the $2^s$ choices of a sign for each $i\in\{1,\dots,s\}$, and put
\[
\beta:=\sum_{i=1}^s\pm\,\xi_i\cdot\frac{\beta_i}{\|\beta_i\|},
\]
the $i$-th term carrying the sign chosen for $i$.
As the supports of the $\beta_i$ are pairwise disjoint, $\|\beta\|^2=\sum_{i=1}^s\xi_i^2$.
On the event that $\pm\langle\beta_i,F\rangle\ge\xi_i\cdot\|\beta_i\|$ for $i=1,\dots,s$, with the same signs, one has $\langle\beta,F\rangle\ge\sum_{i=1}^s\xi_i^2$.
By Markov's inequality for $\exp\langle\beta,F\rangle$ and by (a), this event has probability at most
\[
\exp\left(-\sum_{i=1}^s\xi_i^2\right)\cdot\mathbb{E}\exp\langle\beta,F\rangle\le\exp\left(-\frac12\sum_{i=1}^s\xi_i^2\right).
\]
The event of (c) is the union of these $2^s$ events, and its probability is therefore at most
$2^s\cdot\exp\left(-\frac12\sum_{i=1}^s\xi_i^2\right)=\prod_{i=1}^s2\exp\left(-\frac12\xi_i^2\right)$.

\emph{Item (d).}
Let $\beta\in C^2_{\mathrm{c}}$ have its support in a box $B$ that embeds in $\Gamma$.
The sets $B+\gamma$, $\gamma\in\Gamma$, are then pairwise disjoint (Definition~\ref{def:integrands}),
so that two different plaquettes of the support of $\beta$ are never congruent modulo $\Gamma$.
Consequently the periodisation $\beta_\Gamma\in C^2(\Gamma)$,
\[
\beta_\Gamma(\sigma+\Gamma):=\sum_{\gamma\in\Gamma}\beta(\sigma+\gamma),
\]
has at most one non-zero term in each of these sums.
Hence $\|\beta_\Gamma\|=\|\beta\|$, and $\langle\beta_\Gamma,F\rangle=\langle\beta,\widetilde F\rangle$ for every $F\in C^2(\Gamma)$ with the periodic extension $\widetilde F$.
Since $\mu_{\mathbb{T}}$ is the law of $\widetilde F$ under $\mathbb{E}$ (Definition~\ref{def:functional}), item (a) yields
\[
\int\mathrm{d}\mu_{\mathbb{T}}(F')\,\exp\langle\beta,F'\rangle
=\mathbb{E}\exp\langle\beta_\Gamma,F\rangle
\le\exp\left(\frac12\|\beta_\Gamma\|^2\right)
=\exp\left(\frac12\|\beta\|^2\right).
\]
\end{proof}

\section{The fold inequality}\label{sec:fold}

Throughout this section the Wilson parameter satisfies $-1\le\rho\le1$,
while the hopping weights $w_0,\dots,w_3>0$, the mass parameter $M\in\R$ and the lattice charge $q>0$ are arbitrary.
The sequence $\mathfrak{n}=(n_i)$ of sides satisfies \eqref{cond:divisibility} (Definition~\ref{def:sides}),
and $\mathfrak{d}=(c_0,c_+,c_3;h;O)$ are data in the sense of Definition~\ref{def:data}.
Their cell structure is $\mathcal{Q}_{\mathfrak{d}}=(\mathfrak{F}_{\mathrm{d}};c_0,c_+,c_3;h)$;
its open cells are the sets $Q_\kappa$, $\kappa\in\Z^3$, and $Q_0$ is the base cell
(Definition~\ref{def:cell-structures}).

\subsection{Heights and walls}\label{sec:fold-heights}

The families of the diagonal frame, their normals and the positions of the walls of $\mathcal{Q}_{\mathfrak{d}}$
are numbered as follows:
\begin{equation}\label{eq:fold-frame}
\begin{aligned}
(\ell^1,\ell^2,\ell^3)&:=(\ell_-,\ell_+,x_3), &\qquad (n_1,n_2,n_3)&:=(e_-,e_+,e_3),\\
\left(|n_1|^2,|n_2|^2,|n_3|^2\right)&=(2,2,1), &\qquad (o_1,o_2,o_3)&:=(c_0,c_+,c_3).
\end{aligned}
\end{equation}
Here $n_a$ is the integer vector of the family $\ell^a$ in the sense of Definition~\ref{def:mirrors},
and the three normals are pairwise orthogonal.
The walls of $\mathcal{Q}_{\mathfrak{d}}$ are the mirrors $(\ell^a,o_a+j\cdot h)$ with $a\in\{1,2,3\}$ and $j\in\Z$,
and $W(\mathcal{Q}_{\mathfrak{d}})$ is the group generated by the reflections in them.
For $x\in\R^3$ the number $\ell^a(x)$ is the height of $x$ for the family $\ell^a$.
In a statement about one torus of the sequence, its side $n_i$ is written $n$.

\begin{lemma}[Height coordinates]\label{lem:heights}
Let $\operatorname{ht}\colon\R^3\to\R^3$ be the map $\operatorname{ht}(x):=\left(\ell^1(x),\ell^2(x),\ell^3(x)\right)$.
\begin{enumerate}
\item\label{it:heights-lattice}
The map $\operatorname{ht}$ is linear and invertible.
The image of the integer lattice under it is the subgroup
\[
\operatorname{ht}(\Z^3)=\{u\in\Z^3:\ u_1+u_2\in2\Z\},
\]
whose index in $\Z^3$ equals two.
\item\label{it:heights-reflection}
Let $a\in\{1,2,3\}$, and let $(\ell^a,c)$ be a mirror.
Then for every $x\in\R^3$
\[
\ell^a(\theta_{\ell^a,c}\,x)=2c-\ell^a(x),
\qquad
\ell^b(\theta_{\ell^a,c}\,x)=\ell^b(x)\quad(b\neq a).
\]
\item\label{it:heights-translation}
For $b\in\{1,2,3\}$, $v\in\Z^3$ and $x\in\R^3$ one has $\ell^b(\tau_vx)=\ell^b(x)+\ell^b(v)$.
\item\label{it:heights-cells}
For every $\kappa\in\Z^3$ the open cell $Q_\kappa$ of $\mathcal{Q}_{\mathfrak{d}}$ is
\[
Q_\kappa=\operatorname{ht}^{-1}\left(\prod_{a=1}^{3}\left(o_a+\kappa_a\cdot h,\ o_a+(\kappa_a+1)\cdot h\right)\right).
\]
\end{enumerate}
\end{lemma}

\begin{proof}
Each $\ell^a$ is a linear form, and $\operatorname{ht}(x)=(x_1-x_2,\ x_1+x_2,\ x_3)$.
For $u\in\R^3$ the equation $\operatorname{ht}(x)=u$ has exactly one solution, namely
\[
x_1=\frac{u_1+u_2}{2},\qquad x_2=\frac{u_2-u_1}{2},\qquad x_3=u_3.
\]
If $x\in\Z^3$, then $u=\operatorname{ht}(x)$ is integral and $u_1+u_2=2x_1$ is even.
Conversely, if $u\in\Z^3$ and $u_1+u_2$ is even, then $u_2-u_1=(u_1+u_2)-2u_1$ is even as well,
and the solution $x$ is integral.
The set $\{u\in\Z^3:\ u_1+u_2\in2\Z\}$ is the kernel of the surjective homomorphism
$u\mapsto u_1+u_2+2\Z$ of $\Z^3$ onto $\Z/2\Z$, so its index in $\Z^3$ is two.
This proves \eqref{it:heights-lattice}.

The reflection in a mirror $(\ell,c)$ is
$\theta_{\ell,c}(x)=x-\frac{2}{|n_\ell|^2}\cdot\left(\ell(x)-c\right)\cdot n_\ell$
(Definition~\ref{def:mirrors}).
As the normals of the frame are pairwise orthogonal, $\ell^b(n_a)=|n_a|^2\cdot\delta_{ab}$, and therefore
\[
\ell^b(\theta_{\ell^a,c}\,x)
=\ell^b(x)-\frac{2}{|n_a|^2}\cdot\left(\ell^a(x)-c\right)\cdot\ell^b(n_a)
=\ell^b(x)-2\cdot\left(\ell^a(x)-c\right)\cdot\delta_{ab}.
\]
For $b=a$ the right side equals $2c-\ell^a(x)$, and for $b\neq a$ it equals $\ell^b(x)$,
which is \eqref{it:heights-reflection}.
Item \eqref{it:heights-translation} holds because $\tau_vx=x+v$ (Definition~\ref{def:cochains}) and $\ell^b$ is linear.
Finally, by Definition~\ref{def:cell-structures} the cell $Q_\kappa$ consists of the points $x$ with
$o_a+\kappa_a\cdot h<\ell^a(x)<o_a+(\kappa_a+1)\cdot h$ for $a=1,2,3$,
that is, of the points whose three heights lie in the three open intervals of \eqref{it:heights-cells}.
\end{proof}

\begin{lemma}[The group of the walls]\label{lem:wall-group}
In the notation \eqref{eq:fold-frame}, for $\xi\in\{1,-1\}^3$ and $\alpha\in\Z^3$
let $g_{\xi,\alpha}$ be the bijection of $\R^3$ determined by
\[
\ell^a(g_{\xi,\alpha}\,x)=o_a+\xi_a\cdot\left(\ell^a(x)-o_a\right)+2h\cdot\alpha_a,
\qquad a=1,2,3,\quad x\in\R^3.
\]
\begin{enumerate}
\item\label{it:walls-elements}
One has $W(\mathcal{Q}_{\mathfrak{d}})=\{g_{\xi,\alpha}:\ \xi\in\{1,-1\}^3,\ \alpha\in\Z^3\}$.
\item\label{it:walls-cells}
The element $g_{\xi,\alpha}$ carries the base cell $Q_0$ onto the open cell $Q_\kappa$ whose index has the coordinates
\[
\kappa_a=\begin{cases}2\alpha_a&\text{if }\xi_a=1,\\2\alpha_a-1&\text{if }\xi_a=-1,\end{cases}\qquad a=1,2,3.
\]
Consequently every element of $W(\mathcal{Q}_{\mathfrak{d}})$ maps $Q_0$ onto an open cell,
and every open cell of $\mathcal{Q}_{\mathfrak{d}}$ is the image of $Q_0$ under exactly one element of $W(\mathcal{Q}_{\mathfrak{d}})$;
for the cell $Q_\kappa$, $\kappa\in\Z^3$, this element is denoted $g_\kappa$, so that $g_\kappa Q_0=Q_\kappa$.
Whenever $g\in W(\mathcal{Q}_{\mathfrak{d}})$ and $\kappa\in\Z^3$, the set $gQ_\kappa$ is an open cell $Q_{\kappa'}$, and $g\,g_\kappa=g_{\kappa'}$.
Finally, $g_\kappa$ and the integer $\kappa_1+\kappa_2+\kappa_3$ have the same parity.
\item\label{it:walls-translations}
For every $v$ in the lattice
\[
2\Lambda_h=h\Z e_-+h\Z e_++2h\Z e_3
\]
the translation $\tau_v$ belongs to $W(\mathcal{Q}_{\mathfrak{d}})$,
and every translation in $W(\mathcal{Q}_{\mathfrak{d}})$ is of this form.
\item\label{it:walls-lattice}
An integer vector belongs to $\Lambda_h$ if and only if its three heights are multiples of $h$,
and $\operatorname{ht}$ maps $\Lambda_h$ onto $h\Z^3$:
\[
\Lambda_h=\{v\in\Z^3:\ \operatorname{ht}(v)\in h\Z^3\},\qquad\operatorname{ht}(\Lambda_h)=h\Z^3.
\]
For $v\in\Lambda_h$ the vector $\kappa_v:=\operatorname{ht}(v)/h$ therefore lies in $\Z^3$,
and $\tau_v$ maps every open cell $Q_\kappa$, $\kappa\in\Z^3$, onto $Q_{\kappa+\kappa_v}$.
The resulting action of $\Lambda_h$ on the open cells of $\mathcal{Q}_{\mathfrak{d}}$ is simply transitive,
and the index of $\Lambda_h$ in $\Z^3$ is $[\Z^3:\Lambda_h]=h^3/2=|Q_0|$.
\item\label{it:walls-conjugation}
Let $v\in\Lambda_h$.
Conjugation by $\tau_v$ maps $W(\mathcal{Q}_{\mathfrak{d}})$ onto itself, $\tau_v\,W(\mathcal{Q}_{\mathfrak{d}})\,\tau_v^{-1}=W(\mathcal{Q}_{\mathfrak{d}})$.
Moreover, for every $\kappa\in\Z^3$ exactly one element of $W(\mathcal{Q}_{\mathfrak{d}})$ maps $Q_{\kappa_v}$ onto $Q_{\kappa+\kappa_v}$,
namely $\tau_v\,g_\kappa\,\tau_v^{-1}$.
\item\label{it:walls-tori}
Let $n$ be a positive multiple of $2h$.
Then $n/2\ge h\ge18$, as $h$ is even and $h>8r_O+8$ with $r_O\ge1$ (Definitions~\ref{def:characters} and~\ref{def:data});
so $C(n)$ and $D(n)=D(n/2,n,n)$ are tori (Definitions~\ref{def:tori} and~\ref{def:sides}).
Their period lattices are, respectively, $n\Z^3$ and the lattice $\Gamma$ generated by $\tfrac{n}{2}\cdot e_-$, $n\cdot e_+$ and $n\cdot e_3$.
The generators of $n\Z^3$ and of $\Gamma$ have the heights
\begin{gather}
\operatorname{ht}(n\cdot e_1)=(n,n,0),\qquad\operatorname{ht}(n\cdot e_2)=(-n,n,0),\qquad\operatorname{ht}(n\cdot e_3)=(0,0,n),\notag\\
\operatorname{ht}\left(\tfrac{n}{2}\cdot e_-\right)=(n,0,0),\qquad\operatorname{ht}(n\cdot e_+)=(0,2n,0),\qquad\operatorname{ht}(n\cdot e_3)=(0,0,n),\label{eq:fold-generators}
\end{gather}
so that $\operatorname{ht}(\Gamma)=n\Z\times2n\Z\times n\Z$.
Both lattices lie in $\Lambda_h$, with the same index:
\[
n\Z^3\subset\Lambda_h,\qquad\Gamma\subset2\Lambda_h\subset\Lambda_h,\qquad[\Lambda_h:n\Z^3]=[\Lambda_h:\Gamma]=\frac{n^3}{|Q_0|}.
\]
Translations by vectors of $\Z^3$ leave every torus unchanged, $\tau_v\mathbb{T}=\mathbb{T}$,
and every element $g$ of $W(\mathcal{Q}_{\mathfrak{d}})$ leaves $D(n)$ unchanged, $g\,D(n)=D(n)$.
\end{enumerate}
\end{lemma}

\begin{proof}
By Lemma~\ref{lem:heights}\eqref{it:heights-lattice} the defining identities determine an affine bijection $g_{\xi,\alpha}$ of $\R^3$.
In height coordinates it acts on each height separately:
the $a$-th height $u_a$ is mapped to $o_a+\xi_a\cdot(u_a-o_a)+2h\cdot\alpha_a$, and the three heights are not mixed.
Two such maps are therefore composed height by height:
\[
o_a+\xi_a\cdot\left(\left[o_a+\xi'_a\cdot(u_a-o_a)+2h\cdot\alpha'_a\right]-o_a\right)+2h\cdot\alpha_a
=o_a+\xi_a\cdot\xi'_a\cdot(u_a-o_a)+2h\cdot\left(\alpha_a+\xi_a\cdot\alpha'_a\right).
\]
This gives the composition rule
\[
g_{\xi,\alpha}\,g_{\xi',\alpha'}=g_{\xi'',\alpha''},
\qquad
\xi''_a=\xi_a\cdot\xi'_a,\quad\alpha''_a=\alpha_a+\xi_a\cdot\alpha'_a\quad(a=1,2,3).
\]
In particular $g_{(1,1,1),0}$ is the identity, and the inverse of $g_{\xi,\alpha}$ is $g_{\xi,\alpha'}$ with $\alpha'_a=-\xi_a\cdot\alpha_a$.
So the maps $g_{\xi,\alpha}$ form a group of bijections of $\R^3$.

\emph{Item \eqref{it:walls-elements}.}
By Lemma~\ref{lem:heights}\eqref{it:heights-reflection} the reflection in the wall $(\ell^a,o_a+j\cdot h)$
maps the $a$-th height $u_a$ to $2(o_a+j\cdot h)-u_a=o_a-(u_a-o_a)+2h\cdot j$ and fixes the two other heights;
it is the map $g_{\xi,\alpha}$ with $\xi_a=-1$, $\alpha_a=j$, and $\xi_b=1$, $\alpha_b=0$ for $b\neq a$.
The group of the maps $g_{\xi,\alpha}$ therefore contains the generators of $W(\mathcal{Q}_{\mathfrak{d}})$, hence the whole group.
Conversely, let $\xi$ and $\alpha$ be given.
For $a\in\{1,2,3\}$ consider the map that acts on the $a$-th height as $g_{\xi,\alpha}$ does and fixes the two other heights.
If $\xi_a=-1$, it is the reflection in the wall $(\ell^a,o_a+\alpha_a\cdot h)$.
If $\xi_a=1$, it is the product $\theta_{\ell^a,o_a+\alpha_a\cdot h}\,\theta_{\ell^a,o_a}$ of two wall reflections,
by the composition rule with $(-1)\cdot(-1)=1$ and $\alpha_a+(-1)\cdot0=\alpha_a$.
By the same rule $g_{\xi,\alpha}$ is the product of these three maps, and so it belongs to $W(\mathcal{Q}_{\mathfrak{d}})$.

\emph{Item \eqref{it:walls-cells}.}
In the $a$-th height the interval $(o_a,o_a+h)$ is mapped onto $(o_a+2h\cdot\alpha_a,\ o_a+2h\cdot\alpha_a+h)$ if $\xi_a=1$,
and onto $(o_a+2h\cdot\alpha_a-h,\ o_a+2h\cdot\alpha_a)$ if $\xi_a=-1$;
in both cases the image is $(o_a+\kappa_a\cdot h,\ o_a+(\kappa_a+1)\cdot h)$ with $\kappa_a$ as stated.
Since $g_{\xi,\alpha}$ acts on each height separately, Lemma~\ref{lem:heights}\eqref{it:heights-cells} gives $g_{\xi,\alpha}Q_0=Q_\kappa$.
The assignment $(\xi_a,\alpha_a)\mapsto\kappa_a$ is a bijection of $\{1,-1\}\times\Z$ onto $\Z$:
the value $\xi_a=1$ produces the even integers, the value $\xi_a=-1$ the odd ones.
Distinct indices label distinct cells, since by Lemma~\ref{lem:heights}\eqref{it:heights-cells}
the cells are the preimages under $\operatorname{ht}$ of pairwise disjoint non-empty products of open intervals.
Together with \eqref{it:walls-elements} it follows that $(\xi,\alpha)\mapsto g_{\xi,\alpha}$ is injective
and that $g\mapsto gQ_0$ is a bijection of $W(\mathcal{Q}_{\mathfrak{d}})$ onto the set of the cells.
For $g\in W(\mathcal{Q}_{\mathfrak{d}})$ the element $g\,g_\kappa$ of $W(\mathcal{Q}_{\mathfrak{d}})$ maps $Q_0$ onto $gQ_\kappa$;
hence $gQ_\kappa$ is a cell $Q_{\kappa'}$, and $g\,g_\kappa=g_{\kappa'}$ by the bijectivity just shown.

For the parity, let $\vec g$ be the linear part of $g=g_{\xi,\alpha}$.
Subtracting the defining identities at two points gives $\ell^a(\vec g\,x)=\xi_a\cdot\ell^a(x)$ for all $x\in\R^3$ and $a=1,2,3$.
Since $\ell^a(n_b)=|n_a|^2\cdot\delta_{ab}$ and $\operatorname{ht}$ is injective, it follows that
\[
\vec g\,n_b=\xi_b\cdot n_b\qquad(b=1,2,3).
\]
The normals $n_1,n_2,n_3$ form a basis of $\R^3$, so $\det\vec g=\xi_1\cdot\xi_2\cdot\xi_3$.
By the formula for $\kappa$, $\xi_a=-1$ if and only if $\kappa_a$ is odd.
An element of $W$ is odd when the determinant of its linear part is $-1$ (Definition~\ref{def:mirrors});
thus $g_\kappa$ is odd if and only if an odd number of the $\kappa_a$ are odd, that is, if and only if $\kappa_1+\kappa_2+\kappa_3$ is odd.

\emph{Item \eqref{it:walls-translations}.}
An element $g_{\xi,\alpha}$ is a translation if and only if its linear part is the identity,
that is, since its linear part maps $n_b$ to $\xi_b\cdot n_b$, if and only if $\xi=(1,1,1)$.
For $\xi=(1,1,1)$ all heights are shifted, $\operatorname{ht}(g_{\xi,\alpha}x)=\operatorname{ht}(x)+2h\cdot\alpha$.
From $\operatorname{ht}(h\cdot e_-)=(2h,0,0)$, $\operatorname{ht}(h\cdot e_+)=(0,2h,0)$ and $\operatorname{ht}(2h\cdot e_3)=(0,0,2h)$
one obtains $\operatorname{ht}(v)=2h\cdot\alpha$ for the integral vector
$v:=h\cdot\left(\alpha_1\cdot e_-+\alpha_2\cdot e_++2\alpha_3\cdot e_3\right)$;
as $\operatorname{ht}$ is linear and injective, $g_{\xi,\alpha}=\tau_v$.
When $\alpha$ ranges over $\Z^3$, the vector $v$ ranges over $h\Z e_-+h\Z e_++2h\Z e_3$,
and this lattice is $2\Lambda_h$, by the definition $\Lambda_h=\tfrac{h}{2}(\Z e_-+\Z e_+)+h\Z e_3$ (Definition~\ref{def:arrays}).

\emph{Item \eqref{it:walls-lattice}.}
The generators $\tfrac{h}{2}\cdot e_-$, $\tfrac{h}{2}\cdot e_+$ and $h\cdot e_3$ of $\Lambda_h$
are integral vectors, $h$ being even (Definition~\ref{def:data}), and their heights are
\[
\operatorname{ht}\left(\tfrac{h}{2}\cdot e_-\right)=(h,0,0),\qquad
\operatorname{ht}\left(\tfrac{h}{2}\cdot e_+\right)=(0,h,0),\qquad
\operatorname{ht}(h\cdot e_3)=(0,0,h).
\]
By linearity $\operatorname{ht}(\Lambda_h)=h\Z^3$.
As $\operatorname{ht}$ is a bijection, $\Lambda_h=\operatorname{ht}^{-1}(h\Z^3)$,
and since $\Lambda_h\subset\Z^3$, also $\Lambda_h=\{v\in\Z^3:\ \operatorname{ht}(v)\in h\Z^3\}$.
In particular $v\mapsto\kappa_v$ is a bijection of $\Lambda_h$ onto $\Z^3$.
For $v\in\Lambda_h$ and $x\in\R^3$, Lemma~\ref{lem:heights}\eqref{it:heights-translation} gives
$\ell^a(x+v)=\ell^a(x)+(\kappa_v)_a\cdot h$,
and Lemma~\ref{lem:heights}\eqref{it:heights-cells} shows that $x\in Q_\kappa$ if and only if $x+v\in Q_{\kappa+\kappa_v}$;
thus $\tau_vQ_\kappa=Q_{\kappa+\kappa_v}$.
Given two cells $Q_\kappa$ and $Q_{\kappa'}$, an element $v\in\Lambda_h$ satisfies $\tau_vQ_\kappa=Q_{\kappa'}$
if and only if $\kappa_v=\kappa'-\kappa$, and exactly one $v\in\Lambda_h$ has this property;
so the action of $\Lambda_h$ on the cells is simply transitive.
Finally, $h\Z^3\subset\operatorname{ht}(\Z^3)$ by Lemma~\ref{lem:heights}\eqref{it:heights-lattice}, $h$ being even,
and $\operatorname{ht}$ restricts to an isomorphism of $\Z^3$ onto $\operatorname{ht}(\Z^3)$ that maps $\Lambda_h$ onto $h\Z^3$.
Therefore
\[
[\Z^3:\Lambda_h]=[\operatorname{ht}(\Z^3):h\Z^3]=\frac{[\Z^3:h\Z^3]}{[\Z^3:\operatorname{ht}(\Z^3)]}=\frac{h^3}{2},
\]
and $h^3/2=|Q_0|$ is the volume of an open cell of $\mathcal{Q}_{\mathfrak{d}}$ (Definitions~\ref{def:cell-structures} and~\ref{def:data}).

\emph{Item \eqref{it:walls-conjugation}.}
For every mirror $(\ell,c)$ and every $v\in\Z^3$ one has $\tau_v\theta_{\ell,c}\tau_{-v}=\theta_{\ell,c+\ell(v)}$
(Definition~\ref{def:mirrors}\eqref{it:mirror-conjugation}).
Let $v\in\Lambda_h$.
For a wall $(\ell^a,o_a+j\cdot h)$ one has $\ell^a(v)=(\kappa_v)_a\cdot h$ by \eqref{it:walls-lattice},
so $\tau_v\theta_{\ell^a,o_a+j\cdot h}\tau_v^{-1}$ is the reflection in the wall $(\ell^a,o_a+(j+(\kappa_v)_a)\cdot h)$.
Conjugation by $\tau_v$ therefore maps the set of the wall reflections into itself,
and so does conjugation by $\tau_v^{-1}=\tau_{-v}$, since $-v\in\Lambda_h$.
Hence $\tau_v\,W(\mathcal{Q}_{\mathfrak{d}})\,\tau_v^{-1}=W(\mathcal{Q}_{\mathfrak{d}})$.
By \eqref{it:walls-lattice}, the element $\tau_v\,g_\kappa\,\tau_v^{-1}$ of this group maps $\tau_vQ_0=Q_{\kappa_v}$
onto $\tau_v\,g_\kappa Q_0=\tau_vQ_\kappa=Q_{\kappa+\kappa_v}$.
No other element of $W(\mathcal{Q}_{\mathfrak{d}})$ does so:
if $gQ_{\kappa_v}=Q_{\kappa+\kappa_v}$ for some $g\in W(\mathcal{Q}_{\mathfrak{d}})$, then $g\,g_{\kappa_v}=g_{\kappa+\kappa_v}$ by \eqref{it:walls-cells},
and hence $g=g_{\kappa+\kappa_v}\,g_{\kappa_v}^{-1}$.

\emph{Item \eqref{it:walls-tori}.}
The values
$\ell_-(e_1)=\ell_+(e_1)=\ell_+(e_2)=1$, $\ell_-(e_2)=-1$, $\ell_-(e_3)=\ell_+(e_3)=0$,
$\ell_-(e_-)=\ell_+(e_+)=2$ and $\ell_-(e_+)=\ell_+(e_-)=0$
give the heights listed in \eqref{eq:fold-generators} and in the line above it.
Since $n\in2h\Z$, the heights of $n\cdot e_1$, $n\cdot e_2$, $n\cdot e_3$ lie in $h\Z^3$,
and those of $\tfrac{n}{2}\cdot e_-$, $n\cdot e_+$, $n\cdot e_3$ lie in $2h\Z^3$.
By linearity $\operatorname{ht}(n\Z^3)\subset h\Z^3$ and $\operatorname{ht}(\Gamma)=n\Z\times2n\Z\times n\Z\subset2h\Z^3$.
With \eqref{it:walls-lattice} it follows that $n\Z^3\subset\operatorname{ht}^{-1}(h\Z^3)=\Lambda_h$
and $\Gamma\subset\operatorname{ht}^{-1}(2h\Z^3)=2\Lambda_h\subset\Lambda_h$.
The index of $n\Z^3$ in $\Z^3$ is $n^3$, and so is the index of $\Gamma$,
which is the number of photon sites $2\cdot\tfrac{n}{2}\cdot n\cdot n$ of $D(n/2,n,n)$ (Definition~\ref{def:tori}).
Consequently $[\Lambda_h:n\Z^3]=[\Lambda_h:\Gamma]=n^3/[\Z^3:\Lambda_h]=n^3/|Q_0|$.

By Definition~\ref{def:tori}, the image of a torus under an element of $W$ is formed with the image of its period lattice
under the linear part of that element and with its character composed with the inverse of the linear part.
A translation has the identity as linear part, so $\tau_v\mathbb{T}=\mathbb{T}$ for every torus $\mathbb{T}$ and every $v\in\Z^3$.
Now let $g=g_{\xi,\alpha}\in W(\mathcal{Q}_{\mathfrak{d}})$.
By the proof of \eqref{it:walls-cells}, $\vec g\,n_a=\xi_a\cdot n_a$ and hence $\vec g^{-1}n_a=\xi_a\cdot n_a$ for $a=1,2,3$.
The generators $\tfrac{n}{2}\cdot e_-=\tfrac{n}{2}\cdot n_1$, $n\cdot e_+=n\cdot n_2$ and $n\cdot e_3=n\cdot n_3$ of $\Gamma$
are therefore mapped by $\vec g$ and by $\vec g^{-1}$ to themselves or to their negatives, and $\vec g\,\Gamma=\Gamma$.
The character $\chi$ of $D(n)$ takes the values $\pm1$, so $\chi(-v)=\chi(v)$ for $v\in\Gamma$;
consequently $\chi(\vec g^{-1}v)=\chi(v)$ on the three generators, $\chi\circ\vec g^{-1}=\chi$, and $g\,D(n)=D(n)$.
\end{proof}

Let $n$ be a positive multiple of $2h$, and let $\mathbb{T}$ be one of the tori $C(n)$ and $D(n)$,
with period lattice $\Gamma$; thus $\Gamma=n\Z^3$ for $\mathbb{T}=C(n)$.
By Lemma~\ref{lem:wall-group}\eqref{it:walls-lattice}, \eqref{it:walls-tori},
$\operatorname{ht}(\Gamma)\subset\operatorname{ht}(\Lambda_h)=h\Z^3$, that is, $\ell^a(\Gamma)\subset h\Z$ for $a=1,2,3$;
so $\mathcal{Q}_{\mathfrak{d}}$ cuts $\mathbb{T}$, and the cells of $\mathbb{T}$ are the $\Gamma$-orbits of the cells of $\mathcal{Q}_{\mathfrak{d}}$
(Definition~\ref{def:cell-structures}).
The lattice $\Lambda_h$ contains $\Gamma$ and acts simply transitively on the cells of $\mathcal{Q}_{\mathfrak{d}}$
(Lemma~\ref{lem:wall-group}\eqref{it:walls-lattice}, \eqref{it:walls-tori}).
Every cell is therefore of the form $Q_0+v$ with $v\in\Lambda_h$,
and two cells $Q_0+v$ and $Q_0+v'$ lie in one $\Gamma$-orbit if and only if $v'-v\in\Gamma$.
It follows that the map
\begin{equation}\label{eq:fold-cells}
\Lambda_h/\Gamma\ni v+\Gamma\longmapsto\left\{Q_0+v+g:\ g\in\Gamma\right\}
\end{equation}
is a bijection of $\Lambda_h/\Gamma$ onto the set of the cells of $\mathbb{T}$,
and that $\mathbb{T}$ has exactly $[\Lambda_h:\Gamma]=n^3/|Q_0|$ cells.

\subsection{Two tori, two arrays}\label{sec:fold-arrays}

\begin{definition}[Threshold]\label{def:threshold}
Let the sequence $\mathfrak{n}=(n_i)$ satisfy \eqref{cond:divisibility},
and let $\mathfrak{d}=(c_0,c_+,c_3;h;O)$ be data (Definition~\ref{def:data}).
The \emph{threshold} $i_0(\mathfrak{d})$ is the smallest index $j$ for which all the sides $n_i$ with $i\ge j$ are multiples of $4h$.
It exists, because $4h\in\mathcal{D}(\mathfrak{n})$, that is, $4h$ divides all but finitely many of the $n_i$.
Since $r_O\ge1$, the inequality $h>8r_O+8$ gives $h>16$, and as $h$ is even, $h\ge18$.
Consequently, for every $i\ge i_0(\mathfrak{d})$ the side $n_i$ is even and $n_i\ge4h\ge72$,
and the four objects $C(n_i)$, $D'(n_i)$, $D(n_i)$ and $C(n_i/2,2n_i,n_i)$ of Definition~\ref{def:sides}
are tori in the sense of Definition~\ref{def:tori}:
all their sides are integers not smaller than $n_i/2\ge36$,
so that the systoles of their period lattices are at least four.
\end{definition}

\emph{Order of composition.}
Let $g,g'\in W$.
Pull-backs of cochains and the action of $W$ on even integrands obey the composition rules
\begin{equation}\label{eq:fold-composition}
\left(g\,g'\right)^{*}=g'^{*}\,g^{*},
\qquad
\left(G^{g}\right)^{g'}=G^{g'g}
\end{equation}
for every even integrand $G$ (Definition~\ref{def:cochains} and \eqref{eq:setting-group-action}).
On a function $G$ of the curvature the action reads $G^{g}(F)=G\left(g^{*}F\right)$ for even $g$,
while for odd $g$ the value $G^{g}(F)$ is the complex conjugate of $G\left(g^{*}F\right)$.
The two identities in \eqref{eq:fold-composition} are compatible on such functions.
For even $g$ and $g'$,
\[
\left(G^{g}\right)^{g'}(F)
=G^{g}\left(g'^{*}F\right)
=G\left(g^{*}g'^{*}F\right)
=G\left(\left(g'g\right)^{*}F\right)
=G^{g'g}(F);
\]
if $g$ or $g'$ is odd, the same chain holds with one complex conjugation for each odd factor,
and the number of these conjugations is odd exactly when $g'g$ is odd,
the parity of an element of $W$ being the sign of the determinant of its linear part.
Moreover, if $G$ reads in a set $U\subset\R^3$, then $G^{g}$ reads in $gU$
(Lemma~\ref{lem:admissible-images}\eqref{it:images-boxes}).

\begin{lemma}[Bosonic integrands]\label{lem:bosonic-measure}
Let $\mathbb{T}=(\Gamma,\chi)$ be a torus.
\begin{enumerate}
\item\label{it:bosonic-mass}
The measure on $\Conf(\Gamma)$ that has the density $\exp\left(-\frac{1}{2}\|F\|^2\right)\cdot\det D^{\mathbb{T}}_A$
with respect to the measure of Definition~\ref{def:configurations}
is a finite positive Borel measure.
Its total mass is $Z_{\mathbb{T}}[1]$, and $0<Z_{\mathbb{T}}[1]<\infty$.
\item\label{it:bosonic-integral}
Let $\Phi$ be a bounded Borel function of finitely many $F_\sigma$.
Then $Z_{\mathbb{T}}[\Phi]$ is the integral of $(A,z)\mapsto\Phi(\widetilde F)$ against the measure of \eqref{it:bosonic-mass}:
\[
Z_{\mathbb{T}}[\Phi]=\int_{\Conf(\Gamma)}\exp\left(-\frac{1}{2}\|F\|^2\right)\cdot\det D^{\mathbb{T}}_A\cdot\Phi(\widetilde F),
\]
the integral on the right being taken with respect to the measure of Definition~\ref{def:configurations}.
It follows that $\mu_{\mathbb{T}}$ is a probability measure on the Borel $\sigma$-algebra of $\Omega$
and that $\int\mathrm{d}\mu_{\mathbb{T}}(F)\,\Phi(F)=\mathbb{E}_{\mathbb{T}}[\Phi]$.
\item\label{it:bosonic-invariance}
Let $g\in W$ with $g\mathbb{T}=\mathbb{T}$.
The image of $\mu_{\mathbb{T}}$ under the map $\Omega\to\Omega$, $F\mapsto g^{*}F$, is again $\mu_{\mathbb{T}}$.
\end{enumerate}
\end{lemma}

\begin{proof}
\emph{Item \eqref{it:bosonic-mass}.}
The operator $D^{\mathbb{T}}_A$ depends only on the class of $(A,z)$ in $\Conf(\Gamma)$,
since $D_{A+(2\pi/q)k}=D_A$ for every integer one-cochain $k$
(Proposition~\ref{prop:symmetries}\eqref{it:gauge-covariance}).
Its determinant is a continuous function of $A$,
the entries of $D^{\mathbb{T}}_A$ being continuous functions of the link phases $\exp\left(iqA(b)\right)$,
and $\det D^{\mathbb{T}}_A\ge0$ by Proposition~\ref{prop:accretivity}.
So $\exp\left(-\frac{1}{2}\|F\|^2\right)\cdot\det D^{\mathbb{T}}_A$ is a non-negative Borel function on $\Conf(\Gamma)$,
and the measure with this density is positive.
As $\int_{\mathbb{T}}\exp\left(-S_{\mathbb{T}}(A)\right)=\det D^{\mathbb{T}}_A$ (Definition~\ref{def:tori}),
its total mass is
\[
\int_{\Conf(\Gamma)}\exp\left(-\frac{1}{2}\|F\|^2\right)\cdot\det D^{\mathbb{T}}_A
=\int_{\Conf(\Gamma)}\exp\left(-\frac{1}{2}\|F\|^2\right)\cdot\int_{\mathbb{T}}\exp\left(-S_{\mathbb{T}}(A)\right)
=Z_{\mathbb{T}}[1]
\]
(Definition~\ref{def:functional}),
and this integral converges by Proposition~\ref{prop:berezin}\eqref{it:torus-integrability}, applied to the integrand $1$.
Hence the measure is finite.
Its mass $Z_{\mathbb{T}}[1]$ is positive by Theorem~\ref{thm:partition-positivity}.

\emph{Item \eqref{it:bosonic-integral}.}
The function $\Phi$ is an admissible product without factors of Gaussian growth.
Such a product embeds in every lattice,
and its bounded factor is evaluated at the periodic extension $\widetilde F$ (Definition~\ref{def:integrands});
thus $Z_{\mathbb{T}}[\Phi]$ is given by Definition~\ref{def:functional},
with the integrand $\exp\left(-\frac{1}{2}\|F\|^2\right)\cdot\int_{\mathbb{T}}\exp\left(-S_{\mathbb{T}}(A)\right)\cdot\Phi(\widetilde F)$.
The Berezin integral is linear, and $\Phi(\widetilde F)$ is a scalar, so that
\[
\int_{\mathbb{T}}\exp\left(-S_{\mathbb{T}}(A)\right)\cdot\Phi(\widetilde F)=\Phi(\widetilde F)\cdot\det D^{\mathbb{T}}_A ;
\]
this is the asserted formula for $Z_{\mathbb{T}}[\Phi]$.
After division by $Z_{\mathbb{T}}[1]$, which is positive and finite by \eqref{it:bosonic-mass},
the measure of \eqref{it:bosonic-mass} becomes a probability measure on $\Conf(\Gamma)$.
The map $(A,z)\mapsto\widetilde F$ of $\Conf(\Gamma)$ into $\Omega$ is Borel,
because for fixed $z$ it is a continuous affine function of $A$.
The law $\mu_{\mathbb{T}}$ of $\widetilde F$ under $\mathbb{E}_{\mathbb{T}}$ (Definition~\ref{def:functional})
is the image of this probability measure under that map.
Therefore $\mu_{\mathbb{T}}$ is a Borel probability measure on $\Omega$, and
\[
\mathbb{E}_{\mathbb{T}}[\Phi]=\frac{Z_{\mathbb{T}}[\Phi]}{Z_{\mathbb{T}}[1]}=\int\mathrm{d}\mu_{\mathbb{T}}(F)\,\Phi(F).
\]

\emph{Item \eqref{it:bosonic-invariance}.}
Let $\Phi$ be a real bounded Borel function of finitely many $F_\sigma$.
By \eqref{eq:setting-group-action}, $\Phi^{g}(F)=\Phi\left(g^{*}F\right)$;
for odd $g$ the complex conjugation in that rule does not change the real number $\Phi\left(g^{*}F\right)$.
The function $\Phi^{g}$ is again a real bounded Borel function of finitely many $F_\sigma$,
namely of the $F_{g\sigma}$ with $\sigma$ among the plaquettes on which $\Phi$ depends.
By Proposition~\ref{prop:symmetries}\eqref{it:functional-covariance} and $g\mathbb{T}=\mathbb{T}$,
the number $Z_{\mathbb{T}}\left[\Phi^{g}\right]=Z_{g\mathbb{T}}\left[\Phi^{g}\right]$ equals $Z_{\mathbb{T}}[\Phi]$ if $g$ is even
and $\overline{Z_{\mathbb{T}}[\Phi]}$ if $g$ is odd.
By \eqref{it:bosonic-integral} the number $Z_{\mathbb{T}}[\Phi]$ is the integral of a real function against a positive measure,
hence real, so that $Z_{\mathbb{T}}\left[\Phi^{g}\right]=Z_{\mathbb{T}}[\Phi]$ in both cases.
Dividing by $Z_{\mathbb{T}}[1]$ and applying \eqref{it:bosonic-integral} to $\Phi$ and to $\Phi^{g}$, one obtains
\[
\int\mathrm{d}\mu_{\mathbb{T}}(F)\,\Phi\left(g^{*}F\right)=\int\mathrm{d}\mu_{\mathbb{T}}(F)\,\Phi(F).
\]
The map $F\mapsto g^{*}F$ is a homeomorphism of $\Omega$ onto itself:
it changes the coordinates $F_\sigma$ by a permutation and by signs, it commutes with $d$,
and its inverse is $F\mapsto\left(g^{-1}\right)^{*}F$ by \eqref{eq:fold-composition}.
The left side above is the integral of $\Phi$ against the image of $\mu_{\mathbb{T}}$ under this map.
Let $\Sigma_0$ be a finite set of plaquettes.
Taking for $\Phi$ the indicator functions of the preimages of Borel subsets of $\R^{\Sigma_0}$
under $F\mapsto(F_\sigma)_{\sigma\in\Sigma_0}$
shows that the image measure and $\mu_{\mathbb{T}}$ agree on these cylinder sets.
The cylinder sets, formed with all finite $\Sigma_0$, are stable under finite intersections.
The $\sigma$-algebra generated by them is the Borel $\sigma$-algebra of $\Omega$,
since $\Omega$ is a subspace of the countable product of copies of $\R$ indexed by the plaquettes (Definition~\ref{def:characters}).
By Dynkin's $\pi$--$\lambda$ theorem the image measure equals $\mu_{\mathbb{T}}$.
\end{proof}

\begin{lemma}[The arrays as functions of the curvature]\label{lem:two-arrays}
Let $i\ge i_0(\mathfrak{d})$, let $\Gamma_i$ be the period lattice of $D(n_i)$, and put
\[
\mathcal{K}_i:=\left\{\kappa\in\Z^3:\ 0\le\kappa_1<\frac{n_i}{h},\ 0\le\kappa_2<\frac{2n_i}{h},\ 0\le\kappa_3<\frac{n_i}{h}\right\}.
\]
For $F\in\Omega$ let
\[
\widetilde O_i(F):=\sum_{v\in\Lambda_h\cap[0,n_i)^3}O\left(\tau_v^{*}F\right),
\qquad
O_{\mathfrak{d},i}(F):=\sum_{\kappa\in\mathcal{K}_i}O\left(g_\kappa^{*}F\right)=\sum_{\kappa\in\mathcal{K}_i}O^{[\kappa]}(F).
\]
\begin{enumerate}
\item\label{it:arrays-bounds}
The functions $\widetilde O_i$ and $O_{\mathfrak{d},i}$ are real valued trigonometric polynomials,
each a sum of $n_i^3/|Q_0|$ terms, and for every $F\in\Omega$
\[
\left|\widetilde O_i(F)\right|\le\frac{n_i^3}{|Q_0|}\cdot|O|_1,
\qquad
\left|O_{\mathfrak{d},i}(F)\right|\le\frac{n_i^3}{|Q_0|}\cdot|O|_1 .
\]
For $v\in\Lambda_h$ the term $\tau_vO$, that is $F\mapsto O\left(\tau_v^{*}F\right)$, reads plaquettes contained in $Q_0+v$,
and for $\kappa\in\Z^3$ the term $O^{[\kappa]}$ reads plaquettes contained in $Q_\kappa$.
\item\label{it:arrays-representatives}
The set $\Lambda_h\cap[0,n_i)^3$ is a set of representatives of $\Lambda_h/n_i\Z^3$,
and the cells $Q_\kappa$, $\kappa\in\mathcal{K}_i$, represent the $\Gamma_i$-orbits of open cells, each orbit once.
If $F$ is $n_i\Z^3$-periodic, the value $\widetilde O_i(F)$ does not change
when $\Lambda_h\cap[0,n_i)^3$ is replaced by any other set of representatives of $\Lambda_h/n_i\Z^3$.
If $F$ is $\Gamma_i$-periodic, the value $O_{\mathfrak{d},i}(F)$ does not change
when $\mathcal{K}_i$ is replaced by the set of indices of any other system of representatives of the $\Gamma_i$-orbits of open cells.
Hence, for $(A,z)\in\Conf(n_i\Z^3)$, the number $\widetilde O_i(\widetilde F)$ is the value at $(A,z)$ of the unfolded array on $C(n_i)$,
and for $(A,z)\in\Conf(\Gamma_i)$ the number $O_{\mathfrak{d},i}(\widetilde F)$ is the value at $(A,z)$ of the folded array $O_{\mathfrak{d}}$ on $D(n_i)$
(Definition~\ref{def:arrays}).
\item\label{it:arrays-invariance}
If $g\in W(\mathcal{Q}_{\mathfrak{d}})$ and $F\in\Omega$ is $\Gamma_i$-periodic,
then $g^{*}F$ is $\Gamma_i$-periodic, and $O_{\mathfrak{d},i}\left(g^{*}F\right)=O_{\mathfrak{d},i}(F)$.
\end{enumerate}
\end{lemma}

\begin{proof}
As $i\ge i_0(\mathfrak{d})$, the side $n_i$ is a multiple of $4h$ (Definition~\ref{def:threshold}).
In particular $n_i/h$ is an integer, and $n_i\in2h\Z$,
so that Lemma~\ref{lem:wall-group}\eqref{it:walls-tori} applies with $n=n_i$.

\emph{Item \eqref{it:arrays-bounds}.}
Let $g$ be an isometry of $\R^3$ with $g\Z^3=\Z^3$.
For every plaquette $\sigma$ one has $\left(g^{*}F\right)_\sigma=\pm F_{g\sigma}$,
with the orientation sign of Definition~\ref{def:cochains}, which depends on $g$ and $\sigma$ only.
Write $O=\sum_k a_k\cdot\varphi_{\beta_k}$.
Then
\[
O\left(g^{*}F\right)
=\sum_k a_k\cdot\exp\left(i\sum_\sigma\beta_k(\sigma)\cdot\left(g^{*}F\right)_\sigma\right)
=\sum_k a_k\cdot\exp\left(i\sum_\sigma\pm\beta_k(\sigma)\cdot F_{g\sigma}\right),
\]
the sign in front of $\beta_k(\sigma)$ being the orientation sign of $g$ at $\sigma$.
This is a trigonometric polynomial with the same coefficients $a_k$;
its exponents are again pairwise different,
and it reads the plaquettes $g\sigma$ with $\sigma$ contained in $\supp O$, which lie in $gQ_0$.
In particular $\left|O\left(g^{*}F\right)\right|\le\sum_k|a_k|=|O|_1$ for $F\in\Omega$.
The map $g^{*}$ commutes with $d$ and therefore maps $\Omega$ into $\Omega$;
since $O$ is real on $\Omega$, so is $F\mapsto O\left(g^{*}F\right)$.
For $g=\tau_v$ with $v\in\Lambda_h$ one has $gQ_0=Q_0+v$.
For $g=g_\kappa$ one has $g_\kappa Q_0=Q_\kappa$,
and $O^{[\kappa]}=O^{g_\kappa}$ (Definition~\ref{def:folds}),
which is the function $F\mapsto O\left(g_\kappa^{*}F\right)$ by \eqref{eq:setting-group-action}, $O$ being real.

It remains to count the terms.
The lattice $n_i\Z^3$ is a subgroup of $\Lambda_h$ (Lemma~\ref{lem:wall-group}\eqref{it:walls-tori}),
and every coset $v+n_i\Z^3$ with $v\in\Z^3$ meets the cube $[0,n_i)^3$ in exactly one point;
for $v\in\Lambda_h$ the whole coset lies in $\Lambda_h$.
Hence $\Lambda_h\cap[0,n_i)^3$ is a set of representatives of $\Lambda_h/n_i\Z^3$, and
\[
\#\left|\Lambda_h\cap[0,n_i)^3\right|=\left[\Lambda_h:n_i\Z^3\right]=\frac{n_i^3}{|Q_0|}
\]
by Lemma~\ref{lem:wall-group}\eqref{it:walls-tori}.
Since $|Q_0|=h^3/2$,
\[
\#|\mathcal{K}_i|=\frac{n_i}{h}\cdot\frac{2n_i}{h}\cdot\frac{n_i}{h}=\frac{2n_i^3}{h^3}=\frac{n_i^3}{|Q_0|}.
\]
Each term has modulus at most $|O|_1$ on $\Omega$, which gives the two bounds.

\emph{Item \eqref{it:arrays-representatives}.}
Let $F$ be $n_i\Z^3$-periodic, that is, $\tau_y^{*}F=F$ for $y\in n_i\Z^3$;
translations preserve orientations, so that $\left(\tau_y^{*}F\right)_\sigma=F_{\sigma+y}$.
For $v\in\Lambda_h$ and $y\in n_i\Z^3$ one has $\tau_{v+y}=\tau_y\tau_v$, and \eqref{eq:fold-composition} gives
\[
O\left(\tau_{v+y}^{*}F\right)=O\left(\tau_v^{*}\tau_y^{*}F\right)=O\left(\tau_v^{*}F\right).
\]
The term of $v$ depends only on the coset $v+n_i\Z^3$,
and the sum over any set of representatives of $\Lambda_h/n_i\Z^3$ equals $\widetilde O_i(F)$.

Turning to the folded array,
Lemma~\ref{lem:wall-group}\eqref{it:walls-tori} and \eqref{eq:fold-generators} give
$\operatorname{ht}(\Gamma_i)=n_i\Z\times2n_i\Z\times n_i\Z$ and $\Gamma_i\subset2\Lambda_h\subset\Lambda_h$.
For $y\in\Gamma_i$ the index $\kappa_y=\operatorname{ht}(y)/h$ of Lemma~\ref{lem:wall-group}\eqref{it:walls-lattice} is defined,
and $\tau_yQ_\kappa=Q_{\kappa+\kappa_y}$ for every $\kappa\in\Z^3$.
As $y$ ranges over $\Gamma_i$, the index $\kappa_y$ ranges over the lattice $\frac{n_i}{h}\Z\times\frac{2n_i}{h}\Z\times\frac{n_i}{h}\Z$.
Two open cells $Q_\kappa$ and $Q_{\kappa'}$ therefore lie in one $\Gamma_i$-orbit
if and only if $\kappa'-\kappa$ belongs to this lattice,
and $\mathcal{K}_i$ contains exactly one index of each orbit.
Now let $F$ be $\Gamma_i$-periodic, and let $y\in\Gamma_i$.
The translation $\tau_y$ belongs to $W(\mathcal{Q}_{\mathfrak{d}})$ (Lemma~\ref{lem:wall-group}\eqref{it:walls-translations}),
and $\tau_yg_\kappa$ maps $Q_0$ onto $\tau_yQ_\kappa=Q_{\kappa+\kappa_y}$;
so $\tau_yg_\kappa=g_{\kappa+\kappa_y}$ by Lemma~\ref{lem:wall-group}\eqref{it:walls-cells}.
With \eqref{eq:fold-composition} it follows that
\[
O\left(g_{\kappa+\kappa_y}^{*}F\right)
=O\left(\left(\tau_yg_\kappa\right)^{*}F\right)
=O\left(g_\kappa^{*}\tau_y^{*}F\right)
=O\left(g_\kappa^{*}F\right).
\]
The term of $\kappa$ depends only on the $\Gamma_i$-orbit of $Q_\kappa$,
and the sum over the indices of any system of representatives of the orbits equals $O_{\mathfrak{d},i}(F)$.

For the last assertion recall that $\mathcal{Q}_{\mathfrak{d}}$ cuts $C(n_i)$ and that $D(n_i)$ carries $\mathcal{Q}_{\mathfrak{d}}$
(Definition~\ref{def:data}).
If $(A,z)\in\Conf(n_i\Z^3)$, then $\widetilde F$ is $n_i\Z^3$-periodic,
and by Definition~\ref{def:arrays} the unfolded array on $C(n_i)$ is the sum of the numbers $(\tau_vO)(\widetilde F)$
over a set of representatives of $\Lambda_h/n_i\Z^3$; by the above, this sum is $\widetilde O_i(\widetilde F)$.
If $(A,z)\in\Conf(\Gamma_i)$, then $\widetilde F$ is $\Gamma_i$-periodic,
and the folded array on $D(n_i)$ is the sum of the numbers $O^{[\kappa]}(\widetilde F)$
over a set of representatives of the cells of $D(n_i)$, which are the $\Gamma_i$-orbits of open cells;
this sum is $O_{\mathfrak{d},i}(\widetilde F)$.

\emph{Item \eqref{it:arrays-invariance}.}
Let $g\in W(\mathcal{Q}_{\mathfrak{d}})$ have the linear part $\vec g$, and let $F$ be $\Gamma_i$-periodic.
By Lemma~\ref{lem:wall-group}\eqref{it:walls-tori}, $g\,D(n_i)=D(n_i)$,
so that $\vec g\,\Gamma_i=\Gamma_i$ (Definition~\ref{def:tori}).
For $y\in\Gamma_i$ and $x\in\R^3$ one has $g(x+y)=g(x)+\vec g\,y$, that is, $g\,\tau_y=\tau_{\vec g\,y}\,g$.
Since $\vec g\,y\in\Gamma_i$, \eqref{eq:fold-composition} yields
\[
\tau_y^{*}g^{*}F=\left(g\,\tau_y\right)^{*}F=\left(\tau_{\vec g\,y}\,g\right)^{*}F=g^{*}\tau_{\vec g\,y}^{*}F=g^{*}F,
\]
and $g^{*}F$ is $\Gamma_i$-periodic.
Again by \eqref{eq:fold-composition}, and by Lemma~\ref{lem:wall-group}\eqref{it:walls-cells},
\[
O_{\mathfrak{d},i}\left(g^{*}F\right)
=\sum_{\kappa\in\mathcal{K}_i}O\left(g_\kappa^{*}g^{*}F\right)
=\sum_{\kappa\in\mathcal{K}_i}O\left(\left(g\,g_\kappa\right)^{*}F\right)
=\sum_{\kappa\in\mathcal{K}_i}O\left(g_{\kappa'}^{*}F\right),
\]
where $\kappa'$ is the index with $Q_{\kappa'}=gQ_\kappa$.
The bijection $Q_\kappa\mapsto gQ_\kappa$ of the set of open cells maps each $\Gamma_i$-orbit onto a $\Gamma_i$-orbit,
since $g\,\tau_yQ_\kappa=\tau_{\vec g\,y}\,gQ_\kappa$ and $\vec g\,\Gamma_i=\Gamma_i$;
the same holds for its inverse, which is induced by $g^{-1}$.
The cells $Q_{\kappa'}$, $\kappa\in\mathcal{K}_i$, are therefore again a system of representatives of the $\Gamma_i$-orbits,
and \eqref{it:arrays-representatives}, applied to the $\Gamma_i$-periodic $F$, gives $O_{\mathfrak{d},i}\left(g^{*}F\right)=O_{\mathfrak{d},i}(F)$.
\end{proof}

\subsection{The inequality in finite volume}\label{sec:fold-finite}

By Definition~\ref{def:data}, the observable $O$ of data $\mathfrak{d}$ is a real trigonometric polynomial of the curvature
whose support is confined to a single open cell, namely the base cell $Q_0$ of the cell structure $\mathcal{Q}_{\mathfrak{d}}$ of the diagonal frame.
In the corresponding argument for Yang--Mills theory in four dimensions the array is folded in one direction only,
across the parallel mirrors of a single family,
and accordingly the observable there is placed in one slab of that family.
In the present setting the folding is carried out in all three directions of the diagonal frame;
the reason is given in Remark~\ref{rem:three-directions}.
If the support of the observable is bounded, confinement to one cell is no stronger a requirement
than confinement to one slab of the family $\ell_-$, provided that the width exceeds $4r_O+1$:

\begin{lemma}[One slab, one cell]\label{lem:one-cell}
Let $O\in\mathcal{T}$, let $c_0\in\Z$, and let $h$ be an integer with $h>4r_O+1$.
Suppose that $\supp O$ is contained in the open slab $\{c_0<\ell_-<c_0+h\}$.
Then there are $c_+\in\Z$ and $c_3\in\Z+\frac12$ such that $\supp O$ is contained in the open cell $Q_0$
of the cell structure $(\mathfrak{F}_{\mathrm{d}};c_0,c_+,c_3;h)$.
\end{lemma}

\begin{proof}
Write $r:=r_O$, and let $y\in\R^3$ be a point with $\supp O\subset y+[-r,r]^3$ (Definition~\ref{def:characters}).
Every point $x$ of this cube satisfies $|x_\mu-y_\mu|\le r$ for $\mu=1,2,3$, and therefore
\[
\ell_+(y)-2r\le\ell_+(x)\le\ell_+(y)+2r,
\qquad
y_3-r\le x_3\le y_3+r .
\]
Since $h-4r>1$, the open interval $\left(\ell_+(y)+2r-h,\ \ell_+(y)-2r\right)$, whose length is $h-4r$, contains an integer $c_+$;
with this choice $c_+<\ell_+(x)<c_++h$ at every point $x$ of the cube.
In the third coordinate one uses $h-2r>2r+1>1$:
the open interval $\left(y_3+r-h,\ y_3-r\right)$, of length $h-2r$, contains a point $c_3$ of $\Z+\frac12$,
and then $c_3<x_3<c_3+h$ at every point $x$ of the cube.
The positions $c_0$, $c_+$ and $c_3$ are admissible for $\ell_-$, $\ell_+$ and $x_3$ respectively
(Definition~\ref{def:mirrors}), and $h$ is a positive integer, since $h>4r_O+1\ge5$;
so $(\mathfrak{F}_{\mathrm{d}};c_0,c_+,c_3;h)$ is a cell structure (Definition~\ref{def:cell-structures}), whose open cell $Q_0$ is
\[
Q_0=\left\{x\in\R^3:\ c_0<\ell_-(x)<c_0+h,\ \ c_+<\ell_+(x)<c_++h,\ \ c_3<x_3<c_3+h\right\}.
\]
The set $\supp O$ lies in the slab $\{c_0<\ell_-<c_0+h\}$ by hypothesis,
and in the cube, on which $c_+<\ell_+<c_++h$ and $c_3<x_3<c_3+h$;
it is therefore contained in $Q_0$.
\end{proof}

For a finite positive measure $\mu$ on a measurable space and a bounded measurable function $V$ on that space,
$\mu(V)$ denotes the integral $\int\mathrm{d}\mu(x)\,V(x)$;
in particular $\mu(1)$ is the mass of $\mu$.

\begin{lemma}[Two Jensen inequalities]\label{lem:two-jensen}
Let $\mu$ and $\mu'$ be finite positive measures, each on a measurable space of its own,
with masses $\mu(1)>0$ and $\mu'(1)>0$.
Let $V$ and $V'$ be bounded real measurable functions on the respective spaces, and let $t>0$.
If
\[
\mu\left(\exp(t\cdot V)\right)\le\mu'\left(\exp(t\cdot V')\right),
\]
then
\[
\frac{\mu(V)}{\mu(1)}
\le\frac{\mu'\left(V'\cdot\exp(t\cdot V')\right)}{\mu'\left(\exp(t\cdot V')\right)}
+\frac1t\cdot\log\frac{\mu'(1)}{\mu(1)} .
\]
\end{lemma}

\begin{proof}
From $\exp(t\cdot V)\ge\exp\left(-t\cdot\sup|V|\right)$ one gets $\mu\left(\exp(t\cdot V)\right)\ge\exp\left(-t\cdot\sup|V|\right)\cdot\mu(1)>0$,
and in the same way $\mu'\left(\exp(t\cdot V')\right)>0$.
The exponential function is convex and $\mu/\mu(1)$ is a probability measure,
so Jensen's inequality gives $\exp\left(t\cdot\mu(V)/\mu(1)\right)\le\mu\left(\exp(t\cdot V)\right)/\mu(1)$, that is,
\begin{equation}\label{eq:fold-lower-tangent}
\log\frac{\mu\left(\exp(t\cdot V)\right)}{\mu(1)}\ge t\cdot\frac{\mu(V)}{\mu(1)} .
\end{equation}
Let $\mu'_t$ be the probability measure with density $\exp(t\cdot V')/\mu'\left(\exp(t\cdot V')\right)$ with respect to $\mu'$;
thus $\mu'_t(H)=\mu'\left(H\cdot\exp(t\cdot V')\right)/\mu'\left(\exp(t\cdot V')\right)$ for bounded measurable $H$.
Jensen's inequality for $\mu'_t$, applied to the function $-t\cdot V'$, gives
\[
\exp\left(-t\cdot\mu'_t(V')\right)\le\mu'_t\left(\exp(-t\cdot V')\right)=\frac{\mu'(1)}{\mu'\left(\exp(t\cdot V')\right)},
\]
and taking logarithms,
\begin{equation}\label{eq:fold-upper-tangent}
\log\frac{\mu'\left(\exp(t\cdot V')\right)}{\mu'(1)}\le t\cdot\mu'_t(V') .
\end{equation}
Now \eqref{eq:fold-lower-tangent}, the hypothesis and \eqref{eq:fold-upper-tangent} give, in this order,
\[
\begin{aligned}
t\cdot\frac{\mu(V)}{\mu(1)}
&\le\log\mu\left(\exp(t\cdot V)\right)-\log\mu(1)
\le\log\mu'\left(\exp(t\cdot V')\right)-\log\mu(1)\\
&=\log\frac{\mu'\left(\exp(t\cdot V')\right)}{\mu'(1)}+\log\frac{\mu'(1)}{\mu(1)}
\le t\cdot\mu'_t(V')+\log\frac{\mu'(1)}{\mu(1)} .
\end{aligned}
\]
Since $\mu'_t(V')=\mu'\left(V'\cdot\exp(t\cdot V')\right)/\mu'\left(\exp(t\cdot V')\right)$,
the first and the last member of this chain, divided by $t>0$, are the two sides of the asserted inequality.
\end{proof}

\begin{proposition}[The fold inequality in finite volume]\label{prop:fold-finite}
Let $-1\le\rho\le1$, let $\mathfrak{n}=(n_i)$ be a sequence of sides that satisfies \eqref{cond:divisibility},
and let $\mathfrak{d}=(c_0,c_+,c_3;h;O)$ be data (Definition~\ref{def:data}).
Let $i\ge i_0(\mathfrak{d})$ (Definition~\ref{def:threshold}) and $t>0$.
The unfolded array $\widetilde O$ is read on the cubic torus $C(n_i)$ and the folded array $O_{\mathfrak{d}}$ on the adapted torus $D(n_i)$
(Definition~\ref{def:arrays}, Lemma~\ref{lem:two-arrays}),
and chessboard seminorms on $D(n_i)$ are formed with the cell structure $\mathcal{Q}_{\mathfrak{d}}$.
Then the following hold.
\begin{enumerate}
\item\label{it:finite-lower}
\[
\log\left(\frac{Z_{C(n_i)}\left[\exp(t\cdot\widetilde O)\right]}{Z_{C(n_i)}[1]}\right)\ge t\cdot\frac{n_i^3}{|Q_0|}\cdot\mu_i(O).
\]
\item\label{it:finite-change}
\[
Z_{C(n_i)}\left[\exp(t\cdot\widetilde O)\right]
\le\left\|\exp(t\cdot O)\right\|_{D(n_i)}^{\frac{n_i^3}{|Q_0|}}
=Z_{D(n_i)}\left[\exp(t\cdot O_{\mathfrak{d}})\right].
\]
\item\label{it:finite-upper}
\[
\log\left(\frac{Z_{D(n_i)}\left[\exp(t\cdot O_{\mathfrak{d}})\right]}{Z_{D(n_i)}[1]}\right)\le t\cdot\frac{n_i^3}{|Q_0|}\cdot\mu^t_i(O).
\]
\item\label{it:finite-fold}
\[
\mu_i(O)\le\mu^t_i(O)+\frac{|Q_0|\cdot\Delta_i}{t}.
\]
\end{enumerate}
\end{proposition}

\begin{proof}
Fix $i\ge i_0(\mathfrak{d})$ and $t>0$, and write $n:=n_i$.
By Definition~\ref{def:threshold}, $4h$ divides $n$ and $n\ge72$;
so $n$ is even, $n\in2h\Z$, and $C(n)$, $D(n)$ and $C(n/2,2n,n)$ are tori.
Lemma~\ref{lem:bosonic-measure} applies to each of them;
in particular $Z_{C(n)}[1]$ and $Z_{D(n)}[1]$ are positive (Lemma~\ref{lem:bosonic-measure}\eqref{it:bosonic-mass}).
By Theorem~\ref{thm:chessboard}(b), the structure $\mathcal{Q}_{\mathfrak{d}}$ cuts $C(n)$, the torus $D(n)$ carries $\mathcal{Q}_{\mathfrak{d}}$,
and the set
\[
\mathcal{K}_i=\left\{\kappa\in\Z^3:\ 0\le\kappa_1<\frac nh,\ \ 0\le\kappa_2<\frac{2n}{h},\ \ 0\le\kappa_3<\frac nh\right\}
\]
of Lemma~\ref{lem:two-arrays} represents the cells of $C(n)$ as well as those of $D(n)$;
it has $\#|\mathcal{K}_i|=\frac nh\cdot\frac{2n}{h}\cdot\frac nh=\frac{2n^3}{h^3}=\frac{n^3}{|Q_0|}$ elements.
Let $\Gamma_i$ be the period lattice of $D(n)$,
and let $\widetilde O_i$ and $O_{\mathfrak{d},i}$ be the functions on $\Omega$ of Lemma~\ref{lem:two-arrays}:
\[
\widetilde O_i(F)=\sum_{v\in\Lambda_h\cap[0,n)^3}O(\tau_v^*F),
\qquad
O_{\mathfrak{d},i}(F)=\sum_{\kappa\in\mathcal{K}_i}O(g_\kappa^*F)=\sum_{\kappa\in\mathcal{K}_i}O^{[\kappa]}(F).
\]
Evaluated at $\widetilde F$, they are the unfolded array read on $C(n)$ and the folded array read on $D(n)$
(Lemma~\ref{lem:two-arrays}\eqref{it:arrays-representatives}).
Both are real trigonometric polynomials with $n^3/|Q_0|$ terms (Lemma~\ref{lem:two-arrays}\eqref{it:arrays-bounds});
hence $\exp(t\cdot\widetilde O_i)$ and $\exp(t\cdot O_{\mathfrak{d},i})$ are bounded continuous functions of finitely many $F_\sigma$.

\eqref{it:finite-lower}
By Lemma~\ref{lem:bosonic-measure}\eqref{it:bosonic-integral},
\[
\frac{Z_{C(n)}\left[\exp(t\cdot\widetilde O)\right]}{Z_{C(n)}[1]}=\int\mathrm{d}\mu_i(F)\,\exp\left(t\cdot\widetilde O_i(F)\right).
\]
The inequality \eqref{eq:fold-lower-tangent}, for the probability measure $\mu_i$ and $V:=\widetilde O_i$,
bounds the logarithm of this number from below by
\[
t\cdot\int\mathrm{d}\mu_i(F)\,\widetilde O_i(F)=t\cdot\sum_{v\in\Lambda_h\cap[0,n)^3}\int\mathrm{d}\mu_i(F)\,O(\tau_v^*F).
\]
Every translation maps the torus $C(n)$ onto itself (Lemma~\ref{lem:wall-group}\eqref{it:walls-tori}),
so $\mu_i$ is invariant under $F\mapsto\tau_v^*F$ (Lemma~\ref{lem:bosonic-measure}\eqref{it:bosonic-invariance}).
Each of the $n^3/|Q_0|$ terms of the sum is therefore equal to $\mu_i(O)$, and \eqref{it:finite-lower} follows.

\eqref{it:finite-change}
\emph{The product.}
For $v\in\Lambda_h$ let $G_v$ be the function $F\mapsto\exp\left(t\cdot O(\tau_v^*F)\right)$ on $\Omega$, the cell content $\exp(t\cdot\tau_vO)$.
The polynomial $\tau_vO$ reads the plaquettes of $\supp O+v$, which are contained in $Q_0+v$ (Definition~\ref{def:arrays}),
and $Q_0+v=Q_{\kappa_v}$ (Lemma~\ref{lem:wall-group}\eqref{it:walls-lattice}).
So $G_v$ is a bounded continuous function of finitely many $F_\sigma$ with $\sigma\subset Q_{\kappa_v}$.
It is an element of $\mathcal{A}(Q_{\kappa_v})$, namely an admissible product $\Phi\cdot G_1\cdots G_r$ with $r=0$
(Definition~\ref{def:integrands});
such a product has no boxes, and it embeds in every torus.
The map $v\mapsto\kappa_v$ is a bijection of $\Lambda_h$ onto $\Z^3$ (Lemma~\ref{lem:wall-group}\eqref{it:walls-lattice}).
As $\mathcal{K}_i$ represents the cells of $C(n)$, the bijection \eqref{eq:fold-cells} shows that
the vectors $v\in\Lambda_h$ with $\kappa_v\in\mathcal{K}_i$ form a set of representatives of $\Lambda_h/n\Z^3$.
On $n\Z^3$-periodic $F$ the sum $\widetilde O_i(F)$ may be taken over these vectors (Lemma~\ref{lem:two-arrays}\eqref{it:arrays-representatives});
hence, as functions on $\Conf(n\Z^3)$,
\[
\exp\left(t\cdot\widetilde O_i(\widetilde F)\right)=\prod_{v\in\Lambda_h:\ \kappa_v\in\mathcal{K}_i}G_v(\widetilde F),
\]
a product with one content in each cell of $C(n)$.

\emph{The estimate.}
By Lemma~\ref{lem:bosonic-measure}\eqref{it:bosonic-integral},
$Z_{C(n)}\left[\exp(t\cdot\widetilde O)\right]=Z_{C(n)}[1]\cdot\mu_i\left(\exp(t\cdot\widetilde O_i)\right)$
is the integral of a positive function against a positive measure; so it is a positive number and equal to its modulus.
Since $n\in2h\Z$ and $n\ge8$, Theorem~\ref{thm:chessboard}(b),
applied to the cell structure $\mathcal{Q}_{\mathfrak{d}}$ of the diagonal frame and to the contents $G_v\in\mathcal{A}(Q_{\kappa_v})$, $\kappa_v\in\mathcal{K}_i$, gives
\[
Z_{C(n)}\left[\exp(t\cdot\widetilde O)\right]
=\left|Z_{C(n)}\left[\prod_{v\in\Lambda_h:\ \kappa_v\in\mathcal{K}_i}G_v\right]\right|
\le\prod_{v\in\Lambda_h:\ \kappa_v\in\mathcal{K}_i}\|G_v\|_{D(n)} .
\]

\emph{All factors are equal.}
Put $G_0:=\exp(t\cdot O)$, an element of $\mathcal{A}(Q_0)$, and let $v\in\Lambda_h$.
The translation $\tau_v$ is even, and \eqref{eq:setting-group-action} gives $G_0^{\tau_v}(F)=G_0(\tau_v^*F)=G_v(F)$,
that is, $G_v=G_0^{\tau_v}$.
Let $\kappa'\in\Z^3$, and let $g'$ be the element of $W(\mathcal{Q}_{\mathfrak{d}})$ with $g'Q_{\kappa_v}=Q_{\kappa'}$.
As $G_v$ belongs to $\mathcal{A}(Q_{\kappa_v})$, its mirror copy in the cell $Q_{\kappa'}$ is $G_v^{[\kappa']}=G_v^{g'}$ (Definition~\ref{def:folds}),
and Lemma~\ref{lem:wall-group}\eqref{it:walls-conjugation}, applied to the index $\kappa'-\kappa_v$, shows that $g'=\tau_v\,g_{\kappa'-\kappa_v}\,\tau_v^{-1}$.
Two applications of the composition rule for the action in \eqref{eq:fold-composition},
together with Definition~\ref{def:folds} for $G_0\in\mathcal{A}(Q_0)$, give
\[
G_v^{[\kappa']}
=\left(G_0^{\tau_v}\right)^{g'}
=G_0^{g'\,\tau_v}
=G_0^{\tau_v\,g_{\kappa'-\kappa_v}}
=\left(G_0^{[\kappa'-\kappa_v]}\right)^{\tau_v}.
\]
As a function of the curvature this mirror copy is,
by \eqref{eq:setting-group-action}, \eqref{eq:fold-composition} and Definition~\ref{def:arrays},
\[
G_v^{[\kappa']}(F)
=\exp\left(t\cdot O\left(g_{\kappa'-\kappa_v}^{*}\tau_v^{*}F\right)\right)
=\exp\left(t\cdot O^{[\kappa'-\kappa_v]}\left(\tau_v^{*}F\right)\right),
\qquad F\in\Omega;
\]
all these functions are real, so that the complex conjugation which \eqref{eq:setting-group-action} attaches to odd elements changes nothing.

The translation $\tau_{-v}$ maps $Q_{\kappa'}$ onto $Q_{\kappa'-\kappa_v}$ (Lemma~\ref{lem:wall-group}\eqref{it:walls-lattice})
and commutes with every translation $\tau_y$, $y\in\Gamma_i$.
Hence it carries each $\Gamma_i$-orbit of open cells onto a $\Gamma_i$-orbit, and, $\tau_v$ doing the same, it permutes the cells of $D(n)$;
as $\kappa'$ ranges over $\mathcal{K}_i$, the index $\kappa'-\kappa_v$ therefore ranges over a set of representatives of the cells of $D(n)$.
Now let $F\in\Omega$ be $\Gamma_i$-periodic.
For $y\in\Gamma_i$, \eqref{eq:fold-composition} gives
$\tau_y^{*}\tau_v^{*}F=\left(\tau_v\,\tau_y\right)^{*}F=\left(\tau_y\,\tau_v\right)^{*}F=\tau_v^{*}\tau_y^{*}F=\tau_v^{*}F$,
so that $\tau_v^{*}F$ is $\Gamma_i$-periodic as well.
Lemma~\ref{lem:two-arrays}\eqref{it:arrays-representatives}, applied to $\tau_v^{*}F$ and to the representatives $\kappa'-\kappa_v$, $\kappa'\in\mathcal{K}_i$,
and the formula for $G_v^{[\kappa']}(F)$, first as it stands and then for $v=0$, yield
\[
\prod_{\kappa'\in\mathcal{K}_i}G_v^{[\kappa']}(F)
=\exp\left(t\cdot\sum_{\kappa'\in\mathcal{K}_i}O^{[\kappa'-\kappa_v]}\left(\tau_v^{*}F\right)\right)
=\exp\left(t\cdot O_{\mathfrak{d},i}\left(\tau_v^{*}F\right)\right)
=\prod_{\kappa\in\mathcal{K}_i}G_0^{[\kappa]}\left(\tau_v^{*}F\right).
\]
Form the folds of Definition~\ref{def:folds} over the representatives $\mathcal{K}_i$.
As the even element $\tau_v$ acts on a function $H$ of the curvature by $H^{\tau_v}(F)=H(\tau_v^{*}F)$ (see \eqref{eq:setting-group-action}),
the last display, taken at $F=\widetilde F$ with $(A,z)\in\Conf(\Gamma_i)$, says that, as functions on $\Conf(\Gamma_i)$,
\[
\Fold_{D(n)}G_v
=\prod_{\kappa'\in\mathcal{K}_i}G_v^{[\kappa']}
=\left(\prod_{\kappa\in\mathcal{K}_i}G_0^{[\kappa]}\right)^{\tau_v}
=\left(\Fold_{D(n)}G_0\right)^{\tau_v}.
\]
The two ends of this chain are bounded continuous functions of finitely many $F_\sigma$,
and by Lemma~\ref{lem:bosonic-measure}\eqref{it:bosonic-integral} the value of $Z_{D(n)}$ at such a function $\Phi$
is the integral of $(A,z)\mapsto\Phi(\widetilde F)$ against a measure on $\Conf(\Gamma_i)$;
so the two ends have the same value under $Z_{D(n)}$.
Moreover $\Fold_{D(n)}G_0$ embeds in $D(n)$ (Definition~\ref{def:folds}), $\tau_v$ is even,
and $\tau_v\,D(n)=D(n)$ (Lemma~\ref{lem:wall-group}\eqref{it:walls-tori}).
Proposition~\ref{prop:symmetries}\eqref{it:functional-covariance} therefore gives
\[
Z_{D(n)}\left[\Fold_{D(n)}G_v\right]
=Z_{\tau_v\,D(n)}\left[\left(\Fold_{D(n)}G_0\right)^{\tau_v}\right]
=Z_{D(n)}\left[\Fold_{D(n)}G_0\right].
\]
The chessboard seminorms of $G_v$ and of $G_0$ on $D(n)$ are the powers of these two numbers with the exponent $|Q_0|/n^3$
(Definition~\ref{def:folds}), and $v\in\Lambda_h$ was arbitrary; consequently
\[
\|G_v\|_{D(n)}=\|G_0\|_{D(n)}\qquad(v\in\Lambda_h).
\]

\emph{The fold of $G_0$.}
For $\kappa\in\Z^3$ the mirror copy $G_0^{[\kappa]}=G_0^{g_\kappa}$ is the function $F\mapsto\exp\left(t\cdot O(g_\kappa^*F)\right)=\exp\left(t\cdot O^{[\kappa]}(F)\right)$,
by \eqref{eq:setting-group-action} and Definition~\ref{def:arrays};
for odd $g_\kappa$ the complex conjugation in \eqref{eq:setting-group-action} leaves this real function unchanged.
Since $\mathcal{K}_i$ represents the cells of $D(n)$, the fold of Definition~\ref{def:folds} may be formed over $\mathcal{K}_i$;
as a function on $\Conf(\Gamma_i)$ it is
\[
\left(\Fold_{D(n)}G_0\right)(\widetilde F)=\prod_{\kappa\in\mathcal{K}_i}G_0^{[\kappa]}(\widetilde F)
=\exp\left(t\cdot\sum_{\kappa\in\mathcal{K}_i}O^{[\kappa]}(\widetilde F)\right)
=\exp\left(t\cdot O_{\mathfrak{d},i}(\widetilde F)\right),
\]
and the last function is $\exp(t\cdot O_{\mathfrak{d}})$ read on $D(n)$ (Lemma~\ref{lem:two-arrays}\eqref{it:arrays-representatives}).
The torus $D(n)$ has $n^3/|Q_0|$ cells, and the definition of the chessboard seminorm (Definition~\ref{def:folds}) gives
\[
\|G_0\|_{D(n)}^{\frac{n^3}{|Q_0|}}=Z_{D(n)}\left[\Fold_{D(n)}G_0\right]=Z_{D(n)}\left[\exp(t\cdot O_{\mathfrak{d}})\right].
\]
The product in the estimate has $\#|\mathcal{K}_i|=n^3/|Q_0|$ factors, each equal to $\|G_0\|_{D(n)}=\|\exp(t\cdot O)\|_{D(n)}$;
together with the last display this is \eqref{it:finite-change}.

\eqref{it:finite-upper}
For every bounded Borel function $\Phi$ of finitely many $F_\sigma$,
Definition~\ref{def:data} and Lemma~\ref{lem:bosonic-measure}\eqref{it:bosonic-integral} give
\[
\mu^t_i(\Phi)
=\frac{Z_{D(n)}\left[\exp(t\cdot O_{\mathfrak{d}})\cdot\Phi\right]}{Z_{D(n)}\left[\exp(t\cdot O_{\mathfrak{d}})\right]}
=\frac{\int\mathrm{d}\mu_{D(n)}(F)\,\Phi(F)\cdot\exp\left(t\cdot O_{\mathfrak{d},i}(F)\right)}{\int\mathrm{d}\mu_{D(n)}(F)\,\exp\left(t\cdot O_{\mathfrak{d},i}(F)\right)} .
\]
On these functions $\mu^t_i$ therefore agrees with the tilted probability measure $\mu'_t$
that the proof of Lemma~\ref{lem:two-jensen} forms from $\mu':=\mu_{D(n)}$ and $V':=O_{\mathfrak{d},i}$.
For these choices $\mu'\left(\exp(t\cdot V')\right)/\mu'(1)=Z_{D(n)}\left[\exp(t\cdot O_{\mathfrak{d}})\right]/Z_{D(n)}[1]$,
again by Lemma~\ref{lem:bosonic-measure}\eqref{it:bosonic-integral},
and \eqref{eq:fold-upper-tangent} reads
\[
\log\left(\frac{Z_{D(n)}\left[\exp(t\cdot O_{\mathfrak{d}})\right]}{Z_{D(n)}[1]}\right)
\le t\cdot\mu^t_i\left(O_{\mathfrak{d},i}\right)
=t\cdot\sum_{\kappa\in\mathcal{K}_i}\mu^t_i\left(O^{[\kappa]}\right).
\]
Let $g\in W(\mathcal{Q}_{\mathfrak{d}})$.
Then $g\,D(n)=D(n)$ (Lemma~\ref{lem:wall-group}\eqref{it:walls-tori}),
and $\mu_{D(n)}$ is invariant under $F\mapsto g^*F$ (Lemma~\ref{lem:bosonic-measure}\eqref{it:bosonic-invariance}).
Being the law of the periodic extension $\widetilde F$ (Definition~\ref{def:functional}),
$\mu_{D(n)}$ is carried by the $\Gamma_i$-periodic elements of $\Omega$,
and on these $O_{\mathfrak{d},i}(g^*F)=O_{\mathfrak{d},i}(F)$ (Lemma~\ref{lem:two-arrays}\eqref{it:arrays-invariance}).
For every bounded Borel function $\Phi$ of finitely many $F_\sigma$ it follows that
\begin{equation}\label{eq:fold-invariance}
\begin{aligned}
\int\mathrm{d}\mu_{D(n)}(F)\,\Phi(g^*F)\cdot\exp\left(t\cdot O_{\mathfrak{d},i}(F)\right)
&=\int\mathrm{d}\mu_{D(n)}(F)\,\Phi(g^*F)\cdot\exp\left(t\cdot O_{\mathfrak{d},i}(g^*F)\right)\\
&=\int\mathrm{d}\mu_{D(n)}(F)\,\Phi(F)\cdot\exp\left(t\cdot O_{\mathfrak{d},i}(F)\right).
\end{aligned}
\end{equation}
For $\kappa\in\mathcal{K}_i$ and $g:=g_\kappa$ one has $O(g_\kappa^*F)=O^{[\kappa]}(F)$ (Definition~\ref{def:arrays});
so \eqref{eq:fold-invariance} with $\Phi:=O$, divided by $\int\mathrm{d}\mu_{D(n)}(F)\,\exp\left(t\cdot O_{\mathfrak{d},i}(F)\right)$,
gives $\mu^t_i\left(O^{[\kappa]}\right)=\mu^t_i(O)$.
The sum over $\kappa\in\mathcal{K}_i$ has $n^3/|Q_0|$ terms, and \eqref{it:finite-upper} follows.

\eqref{it:finite-fold}
Taking logarithms in \eqref{it:finite-change} and combining with \eqref{it:finite-lower} and \eqref{it:finite-upper}, one obtains
\[
\begin{aligned}
t\cdot\frac{n^3}{|Q_0|}\cdot\mu_i(O)
&\le\log Z_{C(n)}\left[\exp(t\cdot\widetilde O)\right]-\log Z_{C(n)}[1]\\
&\le\log Z_{D(n)}\left[\exp(t\cdot O_{\mathfrak{d}})\right]-\log Z_{C(n)}[1]\\
&\le t\cdot\frac{n^3}{|Q_0|}\cdot\mu^t_i(O)+\log Z_{D(n)}[1]-\log Z_{C(n)}[1].
\end{aligned}
\]
This is the argument of Lemma~\ref{lem:two-jensen} for the measures $Z_{C(n)}[1]\cdot\mu_i$ and $Z_{D(n)}[1]\cdot\mu_{D(n)}$ on $\Omega$
and the functions $\widetilde O_i$ and $O_{\mathfrak{d},i}$, the hypothesis of that lemma being \eqref{it:finite-change}.
Since $n$ is even and $n\ge8$, Theorem~\ref{thm:four-tori} gives $\log Z_{D(n)}[1]\le\log Z_{C(n/2,2n,n)}[1]$.
Both $C(n/2,2n,n)$ and $C(n)$ have $n^3$ photon sites, so Definition~\ref{def:sides} gives
\[
\log Z_{C(n/2,2n,n)}[1]-\log Z_{C(n)}[1]=n^3\cdot\left(f(n/2,2n,n)-f(n,n,n)\right)=n^3\cdot\Delta_i .
\]
The last line of the chain is thus at most $t\cdot\frac{n^3}{|Q_0|}\cdot\mu^t_i(O)+n^3\cdot\Delta_i$.
Division by $t\cdot n^3/|Q_0|>0$ gives \eqref{it:finite-fold}.
\end{proof}

\subsection{Mirrors of the perturbed functional}\label{sec:fold-mirrors}

Let $\mathfrak{d}$ be data and let $i\ge i_0(\mathfrak{d})$.
In the notation~\eqref{eq:fold-frame},
the torus $D(n_i)=D(n_i/2,n_i,n_i)$ is the rectangular torus of the diagonal frame
with the height periods $p_1=n_i$, $p_2=2n_i$, $p_3=n_i$
(Definition~\ref{def:cell-structures}).
Its period lattice $\Gamma_i$ is generated by the three vectors
\[
\frac{p_1}{|n_1|^2}\cdot n_1=\frac{n_i}{2}\cdot e_-,
\qquad
\frac{p_2}{|n_2|^2}\cdot n_2=n_i\cdot e_+,
\qquad
\frac{p_3}{|n_3|^2}\cdot n_3=n_i\cdot e_3,
\]
and its character $\chi$ takes the value $-1$ on each of them
(Definition~\ref{def:tori}).
The normals of a frame are pairwise orthogonal,
so that $\ell^a(n_b)=|n_a|^2\cdot\delta_{ab}$:
the generator $(p_b/|n_b|^2)\cdot n_b$ has the height $p_b$ for $\ell^b$ and the height $0$ for the two other families.
Moreover $4h$ divides $n_i$ (Definition~\ref{def:threshold}), so that the three height periods are multiples of $2h$.
Let $a\in\{1,2,3\}$.
The height for $\ell^a$ of an element of $\Gamma_i$ is $p_a$ times its coefficient at $(p_a/|n_a|^2)\cdot n_a$ in the basis formed by the three generators;
in particular it vanishes exactly when this coefficient does.
Hence $\Gamma_i$ is the direct sum of $\Gamma_i\cap\{\ell^a=0\}$, which is spanned by the two other generators, and $\Z\cdot(p_a/|n_a|^2)\cdot n_a$;
the vector $(p_a/|n_a|^2)\cdot n_a$ is a closing vector of $D(n_i)$ for $\ell^a$,
the height period of $D(n_i)$ for $\ell^a$ is $p_a$,
and $\chi$ takes the value $-1$ at this closing vector.
The straightened torus $D(n_i)^0(\ell^a)$ is therefore defined and equal to $D(n_i)$
(Definitions~\ref{def:slabs} and~\ref{def:cell-structures}).
For $a,b\in\{1,2,3\}$ this gives
\begin{equation}\label{eq:fold-adapted}
\begin{aligned}
&D(n_i)=\mathbb{T}(\mathfrak{F}_{\mathrm{d}};n_i,2n_i,n_i)=(\Gamma_i,\chi),
\qquad
(p_1,p_2,p_3)=(n_i,2n_i,n_i)\in(2h\Z)^3,
\\
&\ell^a\left(\frac{p_b}{|n_b|^2}\cdot n_b\right)=p_b\cdot\delta_{ab},
\qquad
\ell^a(\Gamma_i)=p_a\Z,
\qquad
\Gamma_i=\left(\Gamma_i\cap\{\ell^a=0\}\right)\oplus\Z\,\frac{p_a}{|n_a|^2}\cdot n_a,
\\
&\chi\left(\frac{p_a}{|n_a|^2}\cdot n_a\right)=-1,
\qquad
D(n_i)^0(\ell^a)=D(n_i),
\\
&\Xi_a:=\Xi\left(D(n_i),\ell^a\right)=\left(\Gamma_i\cap\{\ell^a=0\},\ \chi|_{\Gamma_i\cap\{\ell^a=0\}}\right).
\end{aligned}
\end{equation}
Thus $D(n_i)$ is adapted to each of the three families $\ell^1$, $\ell^2$, $\ell^3$,
with the height period $p_a$ and the cross-section $\Xi_a$ for $\ell^a$
(Definition~\ref{def:slabs}).

For $t\ge0$ the perturbed functional $\mu^t_i$ is defined
on the elements of $\mathcal{A}(\R^3)$ that embed in $D(n_i)$
(Definition~\ref{def:data}).
By Definition~\ref{def:integrands} these are the finite sums of products
$\Phi\cdot G_1\cdots G_r$ of $\mathcal{A}(\R^3)$, $G_j\in\mathcal{G}^{\mathrm{ev}}(B_j)$,
whose boxes satisfy two conditions:
each $B_j$ embeds in $\Gamma_i$,
that is, the sets $B_j+[-1,1]^3+g$, $g\in\Gamma_i$, are pairwise disjoint;
and the sets $B_j+g$, $1\le j\le r$, $g\in\Gamma_i$, are pairwise disjoint.
Let $a\in\{1,2,3\}$ and $c\in\Adm(\ell^a)$.
The products that span $\mathcal{A}_{\Xi_a}[c,c+p_a/2]$
are only required to embed in the lattice $\Gamma_i\cap\{\ell^a=0\}$, which has rank two.
For a product $\Theta_{\ell^a,c}(G')\cdot G$ of two of them both conditions hold nevertheless:
this is Lemma~\ref{lem:reflected-products}\eqref{it:reflected-halves} below.

Let now $c<c''<c+p_a$ be positions in $\Adm(\ell^a)$,
so that the open slabs $\{c<\ell^a<c''\}$ and $\{c''<\ell^a<c+p_a\}$ together fill one period,
and let $G$ and $G''$ be products of $\mathcal{A}(\{c<\ell^a<c''\})$ and of $\mathcal{A}(\{c''<\ell^a<c+p_a\})$, respectively,
that embed in $\Gamma_i\cap\{\ell^a=0\}$.
The factors being even, $G\cdot G''$ is a product of $\mathcal{A}(\R^3)$ whose boxes are those of $G$ and those of $G''$.
For such a product the second condition always holds,
whereas the first one, which involves the margin $[-1,1]^3$, may fail when one of the two slabs is higher than half a period.
As to the second condition, a translation by $g\in\Gamma_i$ adds $\ell^a(g)\in p_a\Z$ to every height, by~\eqref{eq:fold-adapted},
and the open slabs $\{c+k\cdot p_a<\ell^a<c''+k\cdot p_a\}$ and $\{c''+k\cdot p_a<\ell^a<c+(k+1)\cdot p_a\}$, $k\in\Z$, are pairwise disjoint.
Hence, if translates $B+g$ and $B'+g'$ of two boxes of $G\cdot G''$ by elements of $\Gamma_i$ meet,
then $B$ and $B'$ are boxes of the same factor and $\ell^a(g)=\ell^a(g')$;
so $g-g'\in\Gamma_i\cap\{\ell^a=0\}$ and $B+(g-g')$ meets $B'$,
and as the boxes of each factor embed disjointly in $\Gamma_i\cap\{\ell^a=0\}$, one gets $B=B'$ and $g=g'$.
That the first condition may fail is seen for $a=1$, where $p_1=n_i$, and $c''=c+n_i-1$;
the slab $\{c<\ell^1<c''\}$ then has the height $n_i-1>p_1/2$.
Consider the box
\[
B:=y+\left[-\frac{n_i-2}{4},\frac{n_i-2}{4}\right]^3,
\qquad
\ell^1(y)=c+\frac{n_i-1}{2}.
\]
The form $\ell^1=x_1-x_2$ maps $\left[-\frac{n_i-2}{4},\frac{n_i-2}{4}\right]^3$ onto $\left[-\frac{n_i-2}{2},\frac{n_i-2}{2}\right]$,
so the heights of the points of $B$ fill the interval $\left[c+\frac12,c''-\frac12\right]$, and $B$ lies in $\{c<\ell^1<c''\}$;
moreover $B+[-1,1]^3$ is a closed cube of side $\frac{n_i}{2}+1$.
Every non-zero element $g$ of $\Gamma_i\cap\{\ell^1=0\}=\Z\,n_i\cdot e_+\oplus\Z\,n_i\cdot e_3$ satisfies $|g|_\infty\ge n_i>\frac{n_i}{2}+1$,
so $B$ embeds in this lattice;
for $G\in\mathcal{G}^{\mathrm{ev}}(B)$ and $G'':=1$, the product $G\cdot G''$, whose only box is $B$, is therefore of the kind just considered.
The closing vector $\frac{p_1}{|n_1|^2}\cdot n_1=\frac{n_i}{2}\cdot e_-$, however, satisfies $\left|\frac{n_i}{2}\cdot e_-\right|_\infty=\frac{n_i}{2}<\frac{n_i}{2}+1$;
so $B+[-1,1]^3$ meets its translate by this vector, and $B$ does not embed in $\Gamma_i$.
Thus the first condition fails for the box of $G\cdot G''$,
and the product $G\cdot G''$ with the box $B$ does not embed in $\Gamma_i$ (Definition~\ref{def:integrands}).

\begin{lemma}[Reflected products]\label{lem:reflected-products}
Let $\Pi=(\ell,c)$ be a mirror,
$\theta:=\theta_\Pi$ the reflection in $\Pi$,
$\vec\theta$ the linear part of $\theta$,
and $\Theta:=\Theta_\Pi$ the reflection of Definition~\ref{def:grassmann-reflections}.
\begin{enumerate}
\item\label{it:reflected-boxes}
For a box $B=y+[-r,r]^3$, the set $\theta B$ is the box $\theta y+[-r,r]^3$,
and $\Theta$ maps $\mathcal{G}^{\mathrm{ev}}(B)$ into $\mathcal{G}^{\mathrm{ev}}(\theta B)$,
an element $G$ and its image $\Theta G$ having the same growth margin $\varepsilon_G$.
For an open set $U\subset\R^3$ and a bounded Borel function $\Phi$ of finitely many $F_\sigma$ with $\sigma\subset U$,
the function $\Theta\Phi=\overline{\Phi(\theta^*\,\cdot\,)}$ is a bounded Borel function of finitely many $F_\sigma$ with $\sigma\subset\theta U$.
Accordingly, $\Theta$ carries a product $\Phi\cdot G_1\cdots G_r$ of $\mathcal{A}(U)$ with the boxes $B_1,\dots,B_r$
to a product of $\mathcal{A}(\theta U)$ with the boxes $\theta B_1,\dots,\theta B_r$,
and this image embeds in $\vec\theta\,\Gamma'$ whenever the given product embeds in a subgroup $\Gamma'$ of $\Z^3$ of rank three or two.
\item\label{it:reflected-halves}
Let $\mathfrak{d}$ be data, $i\ge i_0(\mathfrak{d})$ and $a\in\{1,2,3\}$,
with $\Gamma_i$, $p_a$ and $\Xi_a$ as in~\eqref{eq:fold-adapted};
let the mirror be $\Pi=(\ell^a,c)$, $c\in\Adm(\ell^a)$,
and put $\Gamma_\parallel:=\Gamma_i\cap\{\ell^a=0\}$.
The reflected product $\Theta(G')\cdot G$ of two products $G$, $G'$ of $\mathcal{A}(\{c<\ell^a<c+p_a/2\})$ that embed in $\Gamma_\parallel$
is itself a product of $\mathcal{A}(\R^3)$, and it embeds in $\Gamma_i$.
The same holds for sums:
for all $G,G'\in\mathcal{A}_{\Xi_a}[c,c+p_a/2]$ the element $\Theta(G')\cdot G$ belongs to $\mathcal{A}(\R^3)$ and embeds in $D(n_i)$.
\end{enumerate}
\end{lemma}

\begin{proof}
Consider first item~\eqref{it:reflected-boxes}, and let $B=y+[-r,r]^3$ be a box.
The linear part $\vec\theta$ lies in $\Ocub$;
it is a signed permutation matrix and maps $[-r,r]^3$ onto itself,
so that $\theta B=\theta y+\vec\theta\,[-r,r]^3=\theta y+[-r,r]^3$.
The reflection $\theta$ permutes the cells of $\Z^3$,
and, being a bijection of $\R^3$,
it maps the cells contained in $B$ onto the cells contained in $\theta B$:
\[
(\theta B)^{(1)}=\theta\left(B^{(1)}\right),
\qquad
(\theta B)^{(2)}=\theta\left(B^{(2)}\right).
\]
Moreover $\theta$ maps $\mathbb{X}$ onto itself (Definition~\ref{def:fermion-sites}),
one has $\pi(\theta x)=\theta\pi(x)$ for $x\in\mathbb{X}$
(Proposition~\ref{prop:geometry}\eqref{it:geometry-projection}),
and $\vec\theta$ maps the cube $[-1/2,1/2]^3$ onto itself.
Hence $\theta\left(\pi(x)+[-1/2,1/2]^3\right)=\pi(\theta x)+[-1/2,1/2]^3$;
as $\mathbb{X}(B)$ consists of the sites $x$ with $\pi(x)+[-1/2,1/2]^3\subset B$ (Definition~\ref{def:integrands}),
it follows that $\mathbb{X}(\theta B)=\theta\,\mathbb{X}(B)$.

Let $G\in\mathcal{G}^{\mathrm{ev}}(B)$ be written as $G(\hat A)=\sum_\Psi a_\Psi(\hat A)\cdot\Psi$,
the sum extending over the monomials $\Psi$ in the fields $\psi_\alpha(x)$, $\bar\psi_\alpha(x)$, $x\in\mathbb{X}(B)$,
and put $n:=\hat n_\ell$.
By~\eqref{eq:setting-reflection-coefficients},
\[
(\Theta G)(\hat A)=\sum_\Psi\overline{a_\Psi(\theta^*\hat A)}\cdot\Theta(\Psi)
\]
for every real cochain $\hat A$ on $(\theta B)^{(1)}$,
where the pull-back $\theta^*\hat A$ (Definition~\ref{def:cochains}) is a cochain on $B^{(1)}$.
By~\eqref{eq:setting-reflection-fields}, $\Theta$ is antilinear,
reverses the order of the factors of a product,
maps each component of $\psi(x)$ to a linear combination of the components of $\bar\psi(\theta x)$,
and each component of $\bar\psi(x)$ to a linear combination of the components of $\psi(\theta x)$,
the coefficients being entries of $-\gamma_n$.
So $\Theta(\Psi)$ is a linear combination of monomials of the same degree as $\Psi$
in the fields at the sites $\theta x$, $x\in\mathbb{X}(B)$,
with coefficients that are, up to sign, products of entries of $\gamma_n$.
It follows that $\Theta G$ is a Borel map from $\R^{(\theta B)^{(1)}}$ to $\Gr^{\mathrm{ev}}(\mathbb{X}(\theta B))$.

Each coefficient of $(\Theta G)(\hat A)$ with respect to the monomials in the fields at the sites of $\mathbb{X}(\theta B)$
is a finite sum of numbers $\overline{a_\Psi(\theta^*\hat A)}$ multiplied by products of entries of $\gamma_n$;
it is therefore bounded by a constant times $\max_\Psi|a_\Psi(\theta^*\hat A)|$.
Let $C$ and $\varepsilon_G$ be the constants of $G$ in the growth condition of Definition~\ref{def:integrands}.
Then
\[
|a_\Psi(\theta^*\hat A)|
\le C\cdot\exp\left(\left(\frac12-\varepsilon_G\right)\sum_{\sigma\in B^{(2)}}(d\theta^*\hat A)_\sigma^2\right).
\]
The pull-back commutes with $d$,
and $(\theta^*u)_\sigma=\pm u_{\theta\sigma}$ for every two-cochain $u$ (Definition~\ref{def:cochains});
thus $(d\theta^*\hat A)_\sigma^2=(d\hat A)_{\theta\sigma}^2$,
and since $\theta$ maps $B^{(2)}$ onto $(\theta B)^{(2)}$,
\[
\sum_{\sigma\in B^{(2)}}(d\theta^*\hat A)_\sigma^2=\sum_{\sigma\in(\theta B)^{(2)}}(d\hat A)_\sigma^2.
\]
So the coefficients of $\Theta G$ satisfy the growth condition on $\theta B$ with the same $\varepsilon_G$.

For $\Theta G\in\mathcal{G}^{\mathrm{ev}}(\theta B)$ it remains to verify gauge invariance.
Fix a real function $\lambda$ on the sites of $\theta B$,
and put $\lambda':=\lambda\circ\theta$, which is a real function on the sites of $B$.
Since $\theta^*$ commutes with $d$, one has $\theta^*(\hat A-d\lambda)=\theta^*\hat A-d\lambda'$.
Under the transformation~\eqref{eq:setting-gauge} with $\lambda'$,
a factor $\psi_\alpha(x)$ of $\Psi$ is multiplied by $\exp(iq\cdot\lambda'(\pi x))$
and a factor $\bar\psi_\alpha(x)$ by $\exp(-iq\cdot\lambda'(\pi x))$;
so $\Psi$ is multiplied by $\exp\left(iq\cdot\langle\lambda',\mathfrak{q}(\Psi)\rangle\right)$,
where $\mathfrak{q}(\Psi)$ is the charge density of $\Psi$ (Definition~\ref{def:chamber-operator}):
its value at a photon site is the number of factors $\psi_\alpha(x)$ of $\Psi$ with $\pi(x)$ at this site
minus the number of factors $\bar\psi_\alpha(x)$ with $\pi(x)$ at this site.
Now apply the transformation with $\lambda$ to $\Theta(\Psi)$.
For each factor $\psi_\alpha(x)$ of $\Psi$, a monomial of $\Theta(\Psi)$ contains a factor $\bar\psi_{\alpha'}(\theta x)$,
which acquires the phase $\exp(-iq\cdot\lambda(\pi\theta x))=\exp(-iq\cdot\lambda'(\pi x))$, since $\pi\theta x=\theta\pi x$;
likewise, for each factor $\bar\psi_\alpha(x)$ of $\Psi$ it contains a factor $\psi_{\alpha'}(\theta x)$,
which acquires the phase $\exp(iq\cdot\lambda'(\pi x))$.
Multiplying these phases over all factors,
one finds that the transformation with $\lambda$ multiplies each monomial of $\Theta(\Psi)$ by the number
\[
\exp\left(-iq\cdot\langle\lambda',\mathfrak{q}(\Psi)\rangle\right)=\overline{\exp\left(iq\cdot\langle\lambda',\mathfrak{q}(\Psi)\rangle\right)}.
\]
The gauge invariance of $G$ at the cochain $\theta^*\hat A$, read coefficient by coefficient, is the identity
\[
a_\Psi(\theta^*\hat A-d\lambda')\cdot\exp\left(iq\cdot\langle\lambda',\mathfrak{q}(\Psi)\rangle\right)=a_\Psi(\theta^*\hat A),
\]
valid for every monomial $\Psi$.
Evaluating $\Theta G$ at $\hat A-d\lambda$ and transforming the fields with $\lambda$ therefore gives
\[
\sum_\Psi\overline{a_\Psi(\theta^*\hat A-d\lambda')\cdot\exp\left(iq\cdot\langle\lambda',\mathfrak{q}(\Psi)\rangle\right)}\cdot\Theta(\Psi)
=\sum_\Psi\overline{a_\Psi(\theta^*\hat A)}\cdot\Theta(\Psi)
=(\Theta G)(\hat A).
\]
Thus $\Theta G$ is gauge invariant, and $\Theta G\in\mathcal{G}^{\mathrm{ev}}(\theta B)$ with the same $\varepsilon_G$.

For a bounded factor $\Phi$,
\eqref{eq:setting-reflection-coefficients} gives $(\Theta\Phi)(F)=\overline{\Phi(\theta^*F)}$.
If $\Phi$ depends only on the $F_\sigma$ with $\sigma$ in a finite set of plaquettes contained in $U$,
then, as $(\theta^*F)_\sigma=\pm F_{\theta\sigma}$,
the function $\Theta\Phi$ is a bounded Borel function of the $F_{\theta\sigma}$ for the same plaquettes $\sigma$,
and the plaquettes $\theta\sigma$ are contained in $\theta U$.

Now let $\Phi\cdot G_1\cdots G_r$ be a product of $\mathcal{A}(U)$ with $G_j\in\mathcal{G}^{\mathrm{ev}}(B_j)$.
By~\eqref{eq:setting-reflection-fields} and~\eqref{eq:setting-reflection-coefficients},
\[
\Theta(\Phi\cdot G_1\cdots G_r)=\Theta(G_r)\cdots\Theta(G_1)\cdot\Theta(\Phi).
\]
Each factor on the right is even and thus commutes with every element of the Grassmann algebra;
the factors may therefore be rearranged without any change of sign,
and the right side equals $\Theta(\Phi)\cdot\Theta(G_1)\cdots\Theta(G_r)$.
The boxes $\theta B_1,\dots,\theta B_r$ lie in $\theta U$,
and no two of them meet, $\theta$ being injective.
By the preceding steps, $\Theta(\Phi\cdot G_1\cdots G_r)$ is a product of $\mathcal{A}(\theta U)$ with these boxes.
For $x\in\R^3$ and $g\in\Z^3$ one has $\theta(x+g)=\theta x+\vec\theta g$, and $\vec\theta\,[-1,1]^3=[-1,1]^3$;
in particular, for $g\in\Gamma'$,
\[
\theta\left(B_j+[-1,1]^3+g\right)=\theta B_j+[-1,1]^3+\vec\theta g,
\qquad
\theta(B_j+g)=\theta B_j+\vec\theta g.
\]
The bijection $\theta$ preserves disjointness, and so does its inverse.
Hence, if for each $j$ the sets $B_j+[-1,1]^3+g$, $g\in\Gamma'$, are pairwise disjoint,
and the sets $B_j+g$, $1\le j\le r$, $g\in\Gamma'$, are pairwise disjoint,
then the same holds for the boxes $\theta B_j$ and the lattice $\vec\theta\,\Gamma'$:
the reflected product embeds in $\vec\theta\,\Gamma'$ whenever the given one embeds in $\Gamma'$.

Turn to item~\eqref{it:reflected-halves}.
Here $\theta=\theta_{\ell^a,c}$;
let $\Phi$ and $\Phi'$ be the bounded factors of $G$ and $G'$.
By Lemma~\ref{lem:heights}\eqref{it:heights-reflection}, $\ell^a(\theta x)=2c-\ell^a(x)$,
so that $\theta$ maps the open slab $\{c<\ell^a<c+p_a/2\}$ onto $\{c-p_a/2<\ell^a<c\}$.
By Definition~\ref{def:mirrors},
$\vec\theta x=x-(2/|n_a|^2)\cdot\ell^a(x)\cdot n_a$;
thus $\vec\theta$ fixes the plane $\{\ell^a=0\}$ pointwise, and $\vec\theta\,\Gamma_\parallel=\Gamma_\parallel$.
By~\eqref{it:reflected-boxes},
$\Theta(G')$ is a product of $\mathcal{A}(\{c-p_a/2<\ell^a<c\})$ with the bounded factor $\Theta(\Phi')$,
its boxes are the images $\theta B$ of the boxes $B$ of $G'$,
and it embeds in $\Gamma_\parallel$.
The boxes of $G$ lie in $\{\ell^a>c\}$, those of $\Theta(G')$ in $\{\ell^a<c\}$,
and within each of the two products the boxes are pairwise disjoint;
consequently no two boxes of $\Theta(G')\cdot G$ meet.
All factors being even, $\Theta(G')\cdot G$ is a product of $\mathcal{A}(\R^3)$
with the bounded factor $\Theta(\Phi')\cdot\Phi$,
whose boxes are those of $G$ and those of $\Theta(G')$.
It remains to show that these boxes embed disjointly in $\Gamma_i$.

\emph{Each box embeds in $\Gamma_i$.}
Let $B=y+[-r,r]^3$ be a box of $G$ or of $\Theta(G')$,
and let $|n_a|_1$ denote the sum of the moduli of the entries of $n_a$.
For $x\in[-r,r]^3$ one has $|\ell^a(x)|=|n_a\cdot x|\le|n_a|_1\cdot r$, with equality at a suitable vertex of the cube.
The entries of $n_a$ lie in $\{0,1,-1\}$, so that $|n_a|_1=|n_a|^2$: both numbers equal $2$ for $a=1,2$ and $1$ for $a=3$.
So $\ell^a(B)$ is the closed interval of length $2|n_a|^2\cdot r$ centred at $\ell^a(y)$.
It is contained in an open interval of length $p_a/2$,
namely $(c,c+p_a/2)$ or $(c-p_a/2,c)$;
therefore $4|n_a|^2\cdot r<p_a$.
Let $0\neq g\in\Gamma_i$.
When $\ell^a(g)=0$, the vector $g$ lies in $\Gamma_\parallel$,
and the embedding of $B$ in $\Gamma_\parallel$ makes $B+[-1,1]^3$ and $B+[-1,1]^3+g$ disjoint.
When $\ell^a(g)\neq0$, one has $|\ell^a(g)|\ge p_a$ by~\eqref{eq:fold-adapted}, and
\[
p_a\le|\ell^a(g)|\le|n_a|_1\cdot|g|_\infty=|n_a|^2\cdot|g|_\infty;
\]
hence $|g|_\infty>4r$.
In addition $|g|_\infty\ge\sys(\Gamma_i)\ge4$,
since $D(n_i)$ is a torus (Definitions~\ref{def:threshold} and~\ref{def:configurations}).
It follows that $|g|_\infty>2r+2$:
for $r<1$ because $2r+2<4$, and for $r\ge1$ because $2r+2\le4r$.
Now $B+[-1,1]^3=y+[-r-1,r+1]^3$ is a closed cube of side $2r+2$,
and its translate by a vector of $\R^3$ misses it precisely when some coordinate of that vector has modulus larger than $2r+2$.
So $B+[-1,1]^3$ and $B+[-1,1]^3+g$ are disjoint in this case too.
Applied to the differences of elements of $\Gamma_i$, this shows that the sets $B+[-1,1]^3+g$, $g\in\Gamma_i$, are pairwise disjoint,
that is, $B$ embeds in $\Gamma_i$.

\emph{The translates of the boxes are pairwise disjoint.}
For $k\in\Z$ call $\{c+k\cdot p_a/2<\ell^a<c+(k+1)\cdot p_a/2\}$ the open slab with the index $k$;
slabs with different indices do not meet.
In this numbering the boxes of $G$ belong to the slab with the index $0$ and the boxes of $\Theta(G')$ to the slab with the index $-1$.
A translation by $g\in\Gamma_i$ adds $\ell^a(g)$ to every height,
and $2\ell^a(g)/p_a$ is an even integer by~\eqref{eq:fold-adapted};
so this translation carries the slab with the index $k$ onto the slab with the index $k+2\ell^a(g)/p_a$, which has the parity of $k$.
Consequently the translates by elements of $\Gamma_i$ of the boxes of $G$ occupy only slabs with even indices,
and those of the boxes of $\Theta(G')$ only slabs with odd indices;
more precisely, the translation by $g$ moves the former into the slab with the index $2\ell^a(g)/p_a$
and the latter into the slab with the index $2\ell^a(g)/p_a-1$.
Let $B$ and $B'$ be boxes of $\Theta(G')\cdot G$ and $g,g'\in\Gamma_i$ such that $B+g$ meets $B'+g'$.
Each of these two sets lies in a single slab, and since they meet, it is one and the same slab.
The parity of its index shows that $B$ and $B'$ belong to the same one of the two products $G$ and $\Theta(G')$;
the index itself then determines the height of the translation,
so that $\ell^a(g)=\ell^a(g')$ and $g-g'\in\Gamma_\parallel$.
The boxes of $G$ embed disjointly in $\Gamma_\parallel$, and so do the boxes of $\Theta(G')$;
since $B+(g-g')$ meets $B'$, it follows that $B=B'$ and $g=g'$.
Together with the preceding step,
the boxes of $\Theta(G')\cdot G$ embed disjointly in $\Gamma_i$,
which means that $\Theta(G')\cdot G$ embeds in $\Gamma_i$.

\emph{Sums.}
An element of $\mathcal{A}_{\Xi_a}[c,c+p_a/2]$ is a finite sum of products of $\mathcal{A}(\{c<\ell^a<c+p_a/2\})$ that embed in $\Gamma_\parallel$,
a complex multiple of such a product being again one.
If $G=\sum_jG_j$ and $G'=\sum_kG'_k$ are such sums,
then, $\Theta$ being additive and the product bilinear,
\[
\Theta(G')\cdot G=\sum_{j,k}\Theta(G'_k)\cdot G_j,
\]
a finite sum of products of $\mathcal{A}(\R^3)$ that embed in $\Gamma_i$.
According to Definition~\ref{def:integrands}, such a sum is an element of $\mathcal{A}(\R^3)$ that embeds in $D(n_i)$.
\end{proof}

Consequently, for $t\ge0$, $a\in\{1,2,3\}$, $c\in\Adm(\ell^a)$ and $G,G'\in\mathcal{A}_{\Xi_a}[c,c+p_a/2]$,
the number $\mu^t_i\left(\Theta_{\ell^a,c}(G')\cdot G\right)$ is defined by Definition~\ref{def:data}.
The functional $Z_{D(n_i)}$ of Definition~\ref{def:functional} is linear on the elements of $\mathcal{A}(\R^3)$ that embed in $D(n_i)$,
and its value at such an element depends only on the element,
not on its representation as a finite sum of products that embed
(Lemma~\ref{lem:product-integrands}\eqref{it:sum-torus}).
For a bounded Borel function $\Phi$ of finitely many $F_\sigma$ the same is plain from Lemma~\ref{lem:bosonic-measure}\eqref{it:bosonic-integral}:
$Z_{D(n_i)}[\Phi]$ is the integral of $\Phi(\widetilde F)$ against the finite positive Borel measure on $\Conf(\Gamma_i)$
of Lemma~\ref{lem:bosonic-measure}\eqref{it:bosonic-mass}.

\begin{proposition}[Positivity at the walls]\label{prop:wall-positivity}
Let $-1\le\rho\le1$, $w_0,\dots,w_3>0$, $M\in\R$ and $q>0$,
and let $\mathfrak d=(c_0,c_+,c_3;h;O)$ be data (Definition~\ref{def:data})
for a sequence of sides $\mathfrak n$ that satisfies \eqref{cond:divisibility}.
Fix $i\ge i_0(\mathfrak d)$ (Definition~\ref{def:threshold}), $t\ge0$, $a\in\{1,2,3\}$ and $j\in\Z$.
With $\ell^a$ and $o_a$ as in \eqref{eq:fold-frame},
the mirror $(\ell^a,c)$ at the position $c:=o_a+j\cdot h$ is a wall of $\mathcal Q_{\mathfrak d}$;
$\theta:=\theta_{\ell^a,c}$ and $\Theta:=\Theta_{\ell^a,c}$ denote the reflections in this wall
(Definitions~\ref{def:mirrors} and~\ref{def:grassmann-reflections}, and \eqref{eq:setting-reflection-coefficients}).
The torus $D(n_i)$, whose period lattice is $\Gamma_i$,
has for $\ell^a$ the height period $p_a$ and the cross-section $\Xi_a=\Xi(D(n_i),\ell^a)$ of \eqref{eq:fold-adapted}.
\begin{enumerate}
\item\label{it:wall-factor}
On $\Conf(\Gamma_i)$ the function $\exp(t\cdot O_{\mathfrak d})$ admits a factorisation
\[
\exp(t\cdot O_{\mathfrak d})=\Phi_+\cdot\Theta(\Phi_+)
\]
in which $\Phi_+\in\mathcal A_{\Xi_a}[c,c+p_a/2]$ is a real, bounded and continuous function of finitely many $F_\sigma$,
each plaquette $\sigma$ lying in the open slab $\{c<\ell^a<c+p_a/2\}$.
\item\label{it:wall-symmetry}
Let $G\in\mathcal A(\R^3)$ embed in $D(n_i)$.
Then $\Theta G$ is again an element of $\mathcal A(\R^3)$ that embeds in $D(n_i)$,
and the values of $\mu^t_i$ at $G$ and at $\Theta G$ are complex conjugate to each other:
\[
\mu^t_i(\Theta G)=\overline{\mu^t_i(G)} .
\]
\item\label{it:wall-form}
Whenever $G$ and $G'$ belong to $\mathcal A_{\Xi_a}[c,c+p_a/2]$,
the element $\Theta(G')\cdot G$ belongs to $\mathcal A(\R^3)$ and embeds in $D(n_i)$.
On $\mathcal A_{\Xi_a}[c,c+p_a/2]$ the assignment
\[
(G',G)\longmapsto\mu^t_i\left(\Theta(G')\cdot G\right)
\]
defines a Hermitian form, and this form is positive semidefinite;
its diagonal values $\mu^t_i(\Theta(G)\cdot G)$ are therefore real and non-negative.
\item\label{it:wall-invariance}
For each $g\in W(\mathcal Q_{\mathfrak d})$ the law $\mu^t_i$ on $\Omega$ coincides with its image under the map $F\mapsto g^*F$.
The three translations $\tau_{h\cdot e_-}$, $\tau_{h\cdot e_+}$ and $\tau_{2h\cdot e_3}$ are elements of $W(\mathcal Q_{\mathfrak d})$;
hence $\mu^t_i$ is invariant under each of them.
\end{enumerate}
\end{proposition}

\begin{proof}
Since $p_a\in2h\Z$ by \eqref{eq:fold-adapted}, both $c$ and $c+p_a/2$ belong to $o_a+h\Z$.
For $a=1,2$ the positions $o_1=c_0$ and $o_2=c_+$ are integers and $\Adm(\ell^1)=\Adm(\ell^2)=\Z$;
for $a=3$ one has $o_3=c_3\in\Z+\frac12$ and $\Adm(\ell^3)=\frac12\Z$.
Hence $c$ and $c+p_a/2$ lie in $\Adm(\ell^a)$, and the class $\mathcal A_{\Xi_a}[c,c+p_a/2]$ of Definition~\ref{def:slabs} is defined.
For $\kappa\in\Z^3$ and $F\in\Omega$ the mirror copy of $O$ in the cell $Q_\kappa$ is $O^{[\kappa]}(F)=O(g_\kappa^*F)$, the observable being real (Definitions~\ref{def:folds} and~\ref{def:arrays}).
By Lemma~\ref{lem:two-arrays}\eqref{it:arrays-representatives} the folded array $O_{\mathfrak d}$, read on $D(n_i)$, is the function $(A,z)\mapsto O_{\mathfrak d,i}(\widetilde F)$ on $\Conf(\Gamma_i)$,
where $O_{\mathfrak d,i}=\sum_{\kappa\in\mathcal K_i}O^{[\kappa]}$.
Accordingly, for every $G\in\mathcal A(\R^3)$ that embeds in $D(n_i)$,
\[
\mu^t_i(G)=\frac{Z_{D(n_i)}[\exp(t\cdot O_{\mathfrak d,i})\cdot G]}{Z_{D(n_i)}[\exp(t\cdot O_{\mathfrak d,i})]},
\]
the bounded function $\exp(t\cdot O_{\mathfrak d,i})$ being evaluated at $\widetilde F$ (Definition~\ref{def:integrands}).

\emph{Proof of~\eqref{it:wall-factor}.}
The integer $p_a/h$ is even; it counts the slabs of height $h$ of $\mathcal Q_{\mathfrak d}$ in one height period of $D(n_i)$ for $\ell^a$.
Consider the index sets
\[
\begin{aligned}
\mathcal K&:=\{\kappa\in\Z^3:\ j\le\kappa_a<j+p_a/h,\ \ 0\le\kappa_b<p_b/h\ \text{ for } b\ne a\},\\
\mathcal K_+&:=\{\kappa\in\mathcal K:\ \kappa_a<j+p_a/(2h)\},\qquad
\mathcal K_-:=\mathcal K\setminus\mathcal K_+ .
\end{aligned}
\]
Lemma~\ref{lem:wall-group}\eqref{it:walls-tori} applies with $n=n_i$, because $n_i=p_1\in2h\Z$.
By \eqref{eq:fold-generators} one has $\operatorname{ht}(\Gamma_i)=n_i\Z\times2n_i\Z\times n_i\Z=p_1\Z\times p_2\Z\times p_3\Z$,
and the same item gives $\Gamma_i\subset\Lambda_h$.
The translation by an element of $\Lambda_h$ shifts the index of every cell by the heights of that element divided by $h$ (Lemma~\ref{lem:wall-group}\eqref{it:walls-lattice}).
Two cells $Q_\kappa$ and $Q_{\kappa'}$ therefore lie in one $\Gamma_i$-orbit if and only if $\kappa'-\kappa\in(p_1/h)\Z\times(p_2/h)\Z\times(p_3/h)\Z$,
and $\mathcal K$ contains exactly one index of each $\Gamma_i$-orbit of cells.
If $\kappa\in\mathcal K_+$, then $j\le\kappa_a\le j+p_a/(2h)-1$;
so the interval $(o_a+\kappa_a\cdot h,\,o_a+(\kappa_a+1)\cdot h)$ of the values of $\ell^a$ on $Q_\kappa$ (Lemma~\ref{lem:heights}\eqref{it:heights-cells}) lies in $(c,\,c+p_a/2)$.
Thus the cells $Q_\kappa$, $\kappa\in\mathcal K_+$, are contained in the open slab $\{c<\ell^a<c+p_a/2\}$.

For $F\in\Omega$ put
\[
\Phi_+(F):=\exp\left(t\cdot\sum_{\kappa\in\mathcal K_+}O^{[\kappa]}(F)\right).
\]
Every mirror copy $O^{[\kappa]}=O\circ g_\kappa^*$, $\kappa\in\Z^3$, is a real continuous function on $\Omega$ of modulus at most $|O|_1$:
the polynomial $O$ has these three properties,
and $g_\kappa^*$ maps $\Omega$ continuously into itself, since it permutes the coordinates $F_\sigma$ up to signs and commutes with $d$ (Definition~\ref{def:cochains}).
Moreover, $O^{[\kappa]}$ reads only plaquettes contained in $Q_\kappa$ (Lemma~\ref{lem:two-arrays}\eqref{it:arrays-bounds}).
Consequently $\Phi_+$ is real, bounded and continuous on $\Omega$,
and it depends on finitely many $F_\sigma$, every such plaquette $\sigma$ being contained in one of the cells $Q_\kappa$, $\kappa\in\mathcal K_+$, and therefore in the open slab.
It is thus an admissible product of $\mathcal A(\{c<\ell^a<c+p_a/2\})$ that consists of a bounded factor alone, without boxes.
Such a product embeds in every lattice (Definition~\ref{def:integrands}), in particular in the lattice $\Gamma_i\cap\{\ell^a=0\}$ of $\Xi_a$;
so $\Phi_+\in\mathcal A_{\Xi_a}[c,c+p_a/2]$ by Definition~\ref{def:slabs}.

Since $\Phi_+$ is real, the rule \eqref{eq:setting-reflection-coefficients} and the identity $g^*\theta^*=(\theta\,g)^*$ of Definition~\ref{def:cochains} give
\[
(\Theta\Phi_+)(F)=\Phi_+(\theta^*F)=\exp\left(t\cdot\sum_{\kappa\in\mathcal K_+}O\left((\theta\,g_\kappa)^*F\right)\right).
\]
The reflection $\theta$ in a wall of $\mathcal Q_{\mathfrak d}$ belongs to $W(\mathcal Q_{\mathfrak d})$;
by Lemma~\ref{lem:wall-group}\eqref{it:walls-cells}, $\theta\,g_\kappa=g_{\kappa'}$ with $Q_{\kappa'}=\theta Q_\kappa$.
By Lemma~\ref{lem:heights}\eqref{it:heights-reflection}, $\theta$ replaces the height $\ell^a$ by $2c-\ell^a=2o_a+2j\cdot h-\ell^a$ and keeps the two other heights;
it maps the interval $(o_a+\kappa_a\cdot h,\,o_a+(\kappa_a+1)\cdot h)$ onto $(o_a+(2j-1-\kappa_a)\cdot h,\,o_a+(2j-\kappa_a)\cdot h)$.
By Lemma~\ref{lem:heights}\eqref{it:heights-cells} it follows that $\kappa'_a=2j-1-\kappa_a$ and $\kappa'_b=\kappa_b$ for $b\ne a$.
So $\kappa\mapsto\kappa'$ is a bijection of $\mathcal K_+$ onto the set of the indices with $j-p_a/(2h)\le\kappa'_a\le j-1$ and $0\le\kappa'_b<p_b/h$ for $b\ne a$.

This set is the translate of $\mathcal K_-$ by the cell index of a period of $D(n_i)$.
Let $v:=(p_a/|n_a|^2)\cdot n_a$ be the closing vector of $D(n_i)$ for $\ell^a$;
it lies in $\Gamma_i$ by \eqref{eq:fold-adapted},
and $\operatorname{ht}(v)=p_a\cdot(\delta_{ab})_{b=1,2,3}$, because $\ell^b(n_a)=|n_a|^2\cdot\delta_{ab}$ for the orthogonal normals of the frame.
Hence its cell index $\kappa_v=\operatorname{ht}(v)/h$ is $(p_a/h)\cdot(\delta_{ab})_{b=1,2,3}$, and the image of $\mathcal K_+$ under $\kappa\mapsto\kappa'$ is $\mathcal K_--\kappa_v$.
As $v\in\Gamma_i\subset2\Lambda_h$ (Lemma~\ref{lem:wall-group}\eqref{it:walls-tori}),
the translation $\tau_v$ belongs to $W(\mathcal Q_{\mathfrak d})$ by Lemma~\ref{lem:wall-group}\eqref{it:walls-translations};
it maps $Q_\kappa$ onto $Q_{\kappa+\kappa_v}$ by Lemma~\ref{lem:wall-group}\eqref{it:walls-lattice},
and therefore $g_{\kappa+\kappa_v}=\tau_v\,g_\kappa$ by Lemma~\ref{lem:wall-group}\eqref{it:walls-cells}.
If $F$ is $\Gamma_i$-periodic, then $\tau_v^*F=F$, and
\[
O^{[\kappa+\kappa_v]}(F)=O(g_\kappa^*\tau_v^*F)=O^{[\kappa]}(F)\qquad(\kappa\in\Z^3).
\]
It follows that for every $\Gamma_i$-periodic $F\in\Omega$
\[
\begin{aligned}
(\Theta\Phi_+)(F)&=\exp\left(t\cdot\sum_{\kappa\in\mathcal K_-}O^{[\kappa]}(F)\right),\\
\Phi_+(F)\cdot(\Theta\Phi_+)(F)&=\exp\left(t\cdot\sum_{\kappa\in\mathcal K}O^{[\kappa]}(F)\right)=\exp\left(t\cdot O_{\mathfrak d,i}(F)\right),
\end{aligned}
\]
where the last equality is Lemma~\ref{lem:two-arrays}\eqref{it:arrays-representatives}, $\mathcal K$ and $\mathcal K_i$ being two sets of representatives of the $\Gamma_i$-orbits of cells.
On $\Conf(\Gamma_i)$ all three functions are evaluated at the periodic extension $\widetilde F$, which is $\Gamma_i$-periodic.
This proves~\eqref{it:wall-factor}.

\emph{Proof of~\eqref{it:wall-symmetry}.}
Let $G\in\mathcal A(\R^3)$ embed in $D(n_i)$, that is, in $\Gamma_i$.
By Lemma~\ref{lem:reflected-products}\eqref{it:reflected-boxes}, $\Theta$ maps each admissible product of $\mathcal A(\R^3)$ that embeds in $\Gamma_i$
to an admissible product of $\mathcal A(\theta\R^3)=\mathcal A(\R^3)$ that embeds in $\vec\theta\,\Gamma_i$.
As $\theta\in W(\mathcal Q_{\mathfrak d})$, Lemma~\ref{lem:wall-group}\eqref{it:walls-tori} gives $\theta\,D(n_i)=D(n_i)$, and in particular $\vec\theta\,\Gamma_i=\Gamma_i$ (Definition~\ref{def:tori}).
The map $\Theta$ is antilinear, so $\Theta G$ is a finite sum of such products:
it lies in $\mathcal A(\R^3)$ and embeds in $D(n_i)$, and $\mu^t_i(G)$ and $\mu^t_i(\Theta G)$ are both defined (Definition~\ref{def:data}).

Multiplication by $\exp(t\cdot O_{\mathfrak d,i})$ keeps $G$ in the domain of $Z_{D(n_i)}$.
The function $\exp(t\cdot O_{\mathfrak d,i})$ is bounded and continuous and depends on finitely many $F_\sigma$ (Lemma~\ref{lem:two-arrays}\eqref{it:arrays-bounds});
write $G$ as a finite sum of products $\Phi\cdot G_1\cdots G_r$ of $\mathcal A(\R^3)$ that embed in $\Gamma_i$.
In each of these products the multiplication affects the bounded factor only:
$\Phi$ becomes the bounded Borel function $\exp(t\cdot O_{\mathfrak d,i})\cdot\Phi$ of finitely many $F_\sigma$,
while $G_1,\dots,G_r$, and with them the boxes, stay as they were.
Embedding is a condition on the boxes alone;
Definition~\ref{def:integrands} imposes nothing on the plaquettes on which a bounded factor depends, and these may be read by $G_1,\dots,G_r$ as well.
So $\exp(t\cdot O_{\mathfrak d,i})\cdot G$ is again a finite sum of products of $\mathcal A(\R^3)$ that embed in $\Gamma_i$,
that is, an element of $\mathcal A(\R^3)$ that embeds in $D(n_i)$.
By \eqref{eq:setting-reflection-fields} the map $\Theta$ reverses the order of products.
By \eqref{eq:setting-reflection-coefficients} it sends the real function $\exp(t\cdot O_{\mathfrak d,i})$ to the scalar function $F\mapsto\exp\left(t\cdot O_{\mathfrak d,i}(\theta^*F)\right)$, which commutes with $\Theta G$.
On $\Conf(\Gamma_i)$ the argument $\widetilde F$ is $\Gamma_i$-periodic, and $O_{\mathfrak d,i}(\theta^*\widetilde F)=O_{\mathfrak d,i}(\widetilde F)$ by Lemma~\ref{lem:two-arrays}\eqref{it:arrays-invariance}.
Hence, as functions on $\Conf(\Gamma_i)$,
\[
\Theta(\exp(t\cdot O_{\mathfrak d,i})\cdot G)=\Theta(G)\cdot\Theta(\exp(t\cdot O_{\mathfrak d,i}))=\exp(t\cdot O_{\mathfrak d,i})\cdot\Theta G .
\]
The reflection $\theta$ is odd, $\theta\,D(n_i)=D(n_i)$, and $(\exp(t\cdot O_{\mathfrak d,i})\cdot G)^\theta=\Theta(\exp(t\cdot O_{\mathfrak d,i})\cdot G)$ by \eqref{eq:setting-group-action}.
Together with Proposition~\ref{prop:symmetries}\eqref{it:functional-covariance}, applied to the odd element $\theta$ of $W$ and to $\exp(t\cdot O_{\mathfrak d,i})\cdot G$, this yields
\[
Z_{D(n_i)}[\exp(t\cdot O_{\mathfrak d,i})\cdot\Theta G]=Z_{D(n_i)}[(\exp(t\cdot O_{\mathfrak d,i})\cdot G)^\theta]=\overline{Z_{D(n_i)}[\exp(t\cdot O_{\mathfrak d,i})\cdot G]} .
\]
The denominator of $\mu^t_i$ is real and positive.
Indeed, $|t\cdot O_{\mathfrak d,i}|\le t\cdot|O|_1\cdot n_i^3/|Q_0|$ by Lemma~\ref{lem:two-arrays}\eqref{it:arrays-bounds}, so that Lemma~\ref{lem:perturbed-denominators} and Theorem~\ref{thm:partition-positivity} give
\[
Z_{D(n_i)}[\exp(t\cdot O_{\mathfrak d,i})]\ge\exp\left(-t\cdot|O|_1\cdot\frac{n_i^3}{|Q_0|}\right)\cdot Z_{D(n_i)}[1]>0 .
\]
Dividing the preceding identity by this number one obtains $\mu^t_i(\Theta G)=\overline{\mu^t_i(G)}$.

\emph{Proof of~\eqref{it:wall-form}.}
Let $G,G'\in\mathcal A_{\Xi_a}[c,c+p_a/2]$.
That $\Theta(G')\cdot G$ belongs to $\mathcal A(\R^3)$ and embeds in $D(n_i)$ is Lemma~\ref{lem:reflected-products}\eqref{it:reflected-halves}.
Definition~\ref{def:data} therefore assigns to $\Theta(G')\cdot G$ the number $\mu^t_i(\Theta(G')\cdot G)$,
whose numerator is the value of the functional $Z_{D(n_i)}$ of Definition~\ref{def:functional} at $\exp(t\cdot O_{\mathfrak d,i})\cdot\Theta(G')\cdot G$.
The elements $\Phi_+\cdot G$ and $\Phi_+\cdot G'$ again belong to $\mathcal A_{\Xi_a}[c,c+p_a/2]$:
multiplication by $\Phi_+$ replaces the bounded factor $\Phi$ of each product by $\Phi_+\cdot\Phi$,
a bounded Borel function of finitely many $F_\sigma$ with $\sigma$ in the open slab, and leaves the boxes unchanged.
The map $\Theta$ reverses products \eqref{eq:setting-reflection-fields},
and the scalar functions $\Phi_+$ and $\Theta(\Phi_+)$ commute with every element of the Grassmann algebra;
therefore
\[
\Theta(\Phi_+\cdot G')\cdot(\Phi_+\cdot G)=\Theta(G')\cdot\Theta(\Phi_+)\cdot\Phi_+\cdot G=\Phi_+\cdot\Theta(\Phi_+)\cdot\Theta(G')\cdot G ,
\]
and by~\eqref{it:wall-factor}
\[
Z_{D(n_i)}[\exp(t\cdot O_{\mathfrak d,i})\cdot\Theta(G')\cdot G]=Z_{D(n_i)}[\Theta(\Phi_+\cdot G')\cdot(\Phi_+\cdot G)] .
\]
By \eqref{eq:fold-adapted} the torus $D(n_i)$ is rectangular, hence adapted to $\ell^a$ (Definition~\ref{def:cell-structures}),
with height period $p_a$ and cross-section $\Xi_a$; moreover $c,\,c+p_a/2\in\Adm(\ell^a)$.
Corollary~\ref{cor:trace-forms}\eqref{it:trace-form}, applied to $\Phi_+\cdot G$ and $\Phi_+\cdot G'$, gives
\[
\mu^t_i(\Theta(G')\cdot G)=\frac{\Tr_{\mathcal H(c)}\left(\mathfrak S[\Phi_+\cdot G']^*\cdot\mathfrak S[\Phi_+\cdot G]\right)}{Z_{D(n_i)}[\exp(t\cdot O_{\mathfrak d,i})]} .
\]
The same corollary shows that the trace in the numerator is a positive semidefinite Hermitian form in the pair $(\Phi_+\cdot G',\Phi_+\cdot G)$;
as $G\mapsto\Phi_+\cdot G$ is linear, the trace is such a form in $(G',G)$ as well.
Division by the denominator, which is positive, as was shown in the proof of~\eqref{it:wall-symmetry}, preserves both properties.

\emph{Proof of~\eqref{it:wall-invariance}.}
Fix $g\in W(\mathcal Q_{\mathfrak d})$.
Whenever $\Phi$ is a bounded Borel function of finitely many $F_\sigma$, so is the composite $\Phi\circ g^*$,
because $(g^*F)_\sigma=\pm F_{g\sigma}$ for every plaquette $\sigma$.
For functions of this kind Definition~\ref{def:data} and Lemma~\ref{lem:bosonic-measure}\eqref{it:bosonic-integral} express $\mu^t_i$ through the law $\mu_{D(n_i)}$:
$\mu^t_i(\Phi)$ is the integral of $\Phi\cdot\exp(t\cdot O_{\mathfrak d,i})$ against $\mu_{D(n_i)}$, divided by the integral of $\exp(t\cdot O_{\mathfrak d,i})$,
and likewise for $\Phi\circ g^*$.
The two denominators are one and the same number,
while the two numerators agree by the identity \eqref{eq:fold-invariance} for the torus $D(n_i)$, which holds for every real $t$.
Hence $\mu^t_i(\Phi\circ g^*)=\mu^t_i(\Phi)$.
Now let $\Sigma_0$ be a finite set of plaquettes, and let $\Phi$ be the indicator function of a cylinder set,
the preimage of a Borel subset of $\R^{\Sigma_0}$ under $F\mapsto(F_\sigma)_{\sigma\in\Sigma_0}$.
As $\mu^t_i$ assigns to a bounded Borel function its integral against the law $\mu^t_i$ (Definition~\ref{def:data}),
the equality just proved says that the image of $\mu^t_i$ under $F\mapsto g^*F$ and $\mu^t_i$ itself give the same mass to this cylinder set.
The cylinder sets form a $\pi$-system,
and the Borel $\sigma$-algebra of $\Omega$, a subspace of the countable product of copies of $\R$ indexed by the plaquettes, is generated by them;
by Dynkin's $\pi$--$\lambda$ theorem the two laws coincide.
For the three translations, Definition~\ref{def:mirrors}\eqref{it:mirror-product} with $c'-c=h$ and $|n_1|^2=|n_2|^2=2$, $|n_3|^2=1$ gives
\[
\theta_{\ell_-,c_0+h}\,\theta_{\ell_-,c_0}=\tau_{h\cdot e_-},\qquad
\theta_{\ell_+,c_++h}\,\theta_{\ell_+,c_+}=\tau_{h\cdot e_+},\qquad
\theta_{x_3,c_3+h}\,\theta_{x_3,c_3}=\tau_{2h\cdot e_3},
\]
and the six reflections on the left are reflections in walls of $\mathcal Q_{\mathfrak d}$.
So $\tau_{h\cdot e_-}$, $\tau_{h\cdot e_+}$ and $\tau_{2h\cdot e_3}$ belong to $W(\mathcal Q_{\mathfrak d})$.
\end{proof}

\subsection{Volume first, then strength}\label{sec:fold-limits}

The space $\Omega$ is a closed subset of $C^2$, the space of all real two-cochains with the product topology:
the equation $dF=0$ is one closed condition at each unit cube,
and each of these conditions involves six coordinates $F_\sigma$.
As a closed subset of a product of countably many copies of $\R$, the space $\Omega$ is Polish.
A law on $\Omega$ is a Borel probability measure on $\Omega$.
For a law $\mu$ and a Borel function $H$ on $\Omega$ that is bounded or non-negative,
$\mu(H)$ denotes the integral $\int\mathrm{d}\mu(F)\,H(F)$.
Laws $\mu_m$ converge weakly to a law $\mu$ if $\mu_m(H)\to\mu(H)$ for every bounded continuous function $H$ on $\Omega$.
An element $g\in W$ acts on $\Omega$ by $F\mapsto g^{*}F$.
This map permutes the coordinates up to signs and commutes with $d$ (Definition~\ref{def:cochains});
it is therefore a homeomorphism of $\Omega$ onto itself, with inverse $F\mapsto(g^{-1})^{*}F$.

\begin{lemma}[Tightness and closed conditions]\label{lem:closed-conditions}
Let $\mu$ be a law on $\Omega$,
and let $(\mu_m)$ be a family of laws on $\Omega$, indexed by an arbitrary set.
\begin{enumerate}
\item\label{it:closed-tightness}
If there is a constant $C<\infty$ such that
\[
\mu_m\left(\exp(F_\sigma)\right)+\mu_m\left(\exp(-F_\sigma)\right)\le C
\]
for every member $\mu_m$ of the family and every plaquette $\sigma$,
then the family $(\mu_m)$ is tight.
\item\label{it:closed-invariance}
Let the family be a sequence $(\mu_m)_{m\in\N}$ that converges weakly to $\mu$, and let $g\in W$.
If every $\mu_m$ is invariant under $F\mapsto g^{*}F$, then $\mu$ is invariant under $F\mapsto g^{*}F$.
\item\label{it:closed-positivity}
Let the family be a sequence $(\mu_m)_{m\in\N}$ that converges weakly to $\mu$,
let $(\ell,c)$ be a mirror and $\theta:=\theta_{\ell,c}$.
Suppose that for every bounded continuous $\Phi\in\mathcal{O}^{+}(\ell,c)$ the inequality
\[
\int\mathrm{d}\mu_m(F)\,\overline{\Phi(\theta^{*}F)}\cdot\Phi(F)\ge0
\]
holds for all sufficiently large $m$, where how large $m$ has to be may depend on $\Phi$.
Then $\int\mathrm{d}\mu(F)\,\overline{\Phi(\theta^{*}F)}\cdot\Phi(F)\ge0$
for every bounded continuous $\Phi\in\mathcal{O}^{+}(\ell,c)$.
\item\label{it:closed-continuous}
Let $(\ell,c)$ be a mirror and $\theta:=\theta_{\ell,c}$.
Suppose that $\mu$ is invariant at $(\ell,c)$
and that $\int\mathrm{d}\mu(F)\,\overline{\Phi(\theta^{*}F)}\cdot\Phi(F)\ge0$
for every bounded continuous $\Phi\in\mathcal{O}^{+}(\ell,c)$.
Then $\mu$ is reflection positive at $(\ell,c)$ in the sense of Definition~\ref{def:reflection-positivity}:
the inequality holds for every $\Phi\in\mathcal{O}^{+}(\ell,c)$, continuous or not.
\item\label{it:closed-domination}
Let the family be a sequence $(\mu_m)_{m\in\N}$ that converges weakly to $\mu$,
let $\beta\in C^2_{\mathrm{c}}$ and $C<\infty$.
If $\mu_m(\exp\langle\beta,F\rangle)\le C$ for all $m$,
then $\mu(\exp\langle\beta,F\rangle)\le C$.
In particular, the weak limit of a sequence of laws that satisfy \eqref{cond:gaussian-domination}
satisfies \eqref{cond:gaussian-domination}.
\end{enumerate}
\end{lemma}

\begin{proof}
\eqref{it:closed-tightness}
For real numbers $x$ and $r$ with $|x|>r$ one has $\exp(r)<\exp(|x|)\le\exp(x)+\exp(-x)$.
Integrating this inequality over the set on which $|F_\sigma|>r$,
one finds for every plaquette $\sigma$, every real $r$ and every member of the family
\[
\mu_m\left(\{F\in\Omega:|F_\sigma|>r\}\right)
\le\exp(-r)\cdot\left(\mu_m\left(\exp(F_\sigma)\right)+\mu_m\left(\exp(-F_\sigma)\right)\right)
\le C\cdot\exp(-r).
\]
Let $\sigma_1,\sigma_2,\dots$ be an enumeration of the plaquettes and let $\varepsilon>0$.
Put
\[
r_k:=\log\left(\frac{2^k\cdot C}{\varepsilon}\right)\quad(k\ge1),
\qquad
\Omega_\varepsilon:=\left\{F\in\Omega:|F_{\sigma_k}|\le r_k\text{ for all }k\ge1\right\}.
\]
In the coordinates $(F_{\sigma_k})_{k\ge1}$ the set $\Omega_\varepsilon$ is the intersection of $\Omega$
with the product of the intervals $[-r_k,r_k]$.
This product is compact by Tychonoff's theorem,
and $\Omega_\varepsilon$ is a closed subset of it, since $\Omega$ is closed;
so $\Omega_\varepsilon$ is compact.
For every member of the family,
\[
\mu_m(\Omega\setminus\Omega_\varepsilon)
\le\sum_{k\ge1}\mu_m\left(\{F\in\Omega:|F_{\sigma_k}|>r_k\}\right)
\le\sum_{k\ge1}C\cdot\exp(-r_k)
=\sum_{k\ge1}2^{-k}\cdot\varepsilon
=\varepsilon .
\]

\eqref{it:closed-invariance}
Let $H$ be a bounded continuous function on $\Omega$.
Then $H\circ g^{*}$ is bounded and continuous as well, and
\[
\mu(H\circ g^{*})=\lim_{m\to\infty}\mu_m(H\circ g^{*})=\lim_{m\to\infty}\mu_m(H)=\mu(H).
\]
Thus the image of $\mu$ under $F\mapsto g^{*}F$ and $\mu$ itself have the same integrals of all bounded continuous functions.
On a metric space two Borel probability measures with this property coincide.

\eqref{it:closed-positivity}
Fix a bounded continuous $\Phi\in\mathcal{O}^{+}(\ell,c)$.
Since $\theta^{*}$ is continuous, the function $F\mapsto\overline{\Phi(\theta^{*}F)}\cdot\Phi(F)$ is bounded and continuous on $\Omega$.
Weak convergence, applied to its real and to its imaginary part, gives
\[
\int\mathrm{d}\mu(F)\,\overline{\Phi(\theta^{*}F)}\cdot\Phi(F)
=\lim_{m\to\infty}\int\mathrm{d}\mu_m(F)\,\overline{\Phi(\theta^{*}F)}\cdot\Phi(F),
\]
and the numbers on the right are $\ge0$ for all large $m$.

\eqref{it:closed-continuous}
Let $\Phi\in\mathcal{O}^{+}(\ell,c)$.
There is a finite set $\Sigma_0\subset\Sigma^{+}(\ell,c)$ of plaquettes
such that $\Phi$ is a bounded Borel function of $(F_\sigma)_{\sigma\in\Sigma_0}$;
a function on $\R^{\Sigma_0}$ and the function on $\Omega$ that it defines through $F\mapsto(F_\sigma)_{\sigma\in\Sigma_0}$
are denoted by the same letter.
The image of $\mu$ under the map $F\mapsto(F_\sigma)_{\sigma\in\Sigma_0}$ is a Borel probability measure on $\R^{\Sigma_0}$;
it is denoted $\mu_{\Sigma_0}$.
Continuous functions of compact support are dense in $L^2(\mu_{\Sigma_0})$ \cite[Theorem~3.14]{Rudin1987};
choose $\Phi_k\in C_{\mathrm{c}}(\R^{\Sigma_0})$, $k\ge1$, with $\|\Phi_k-\Phi\|_{L^2(\mu_{\Sigma_0})}\to0$.
Each $\Phi_k$ is a bounded continuous element of $\mathcal{O}^{+}(\ell,c)$,
and $\|\Phi_k-\Phi\|_{L^2(\mu)}=\|\Phi_k-\Phi\|_{L^2(\mu_{\Sigma_0})}$.
The invariance of $\mu$ under $\theta^{*}$ gives $\|H\circ\theta^{*}\|_{L^2(\mu)}=\|H\|_{L^2(\mu)}$
for every bounded Borel function $H$ on $\Omega$.
From the identity
\[
\overline{\Phi\circ\theta^{*}}\cdot\Phi-\overline{\Phi_k\circ\theta^{*}}\cdot\Phi_k
=\overline{(\Phi-\Phi_k)\circ\theta^{*}}\cdot\Phi+\overline{\Phi_k\circ\theta^{*}}\cdot(\Phi-\Phi_k)
\]
and the Cauchy--Schwarz inequality in $L^2(\mu)$, applied to each of the two terms, one obtains
\begin{equation}\label{eq:fold-approximation}
\left|\mu\left(\overline{\Phi\circ\theta^{*}}\cdot\Phi\right)-\mu\left(\overline{\Phi_k\circ\theta^{*}}\cdot\Phi_k\right)\right|
\le\|\Phi-\Phi_k\|_{L^2(\mu)}\cdot\left(\|\Phi\|_{L^2(\mu)}+\|\Phi_k\|_{L^2(\mu)}\right).
\end{equation}
The norms $\|\Phi_k\|_{L^2(\mu)}$ are bounded in $k$,
so the right side of \eqref{eq:fold-approximation} tends to $0$ as $k\to\infty$.
By hypothesis $\mu\left(\overline{\Phi_k\circ\theta^{*}}\cdot\Phi_k\right)\ge0$ for every $k$;
hence $\mu\left(\overline{\Phi\circ\theta^{*}}\cdot\Phi\right)\ge0$.

\eqref{it:closed-domination}
As $\beta$ is finitely supported, $\langle\beta,F\rangle$ is a finite linear combination of coordinates of $F$.
For $r>0$ the function $F\mapsto\min\left(\exp\langle\beta,F\rangle,r\right)$ is therefore bounded and continuous on $\Omega$, and
\[
\mu\left(\min\left(\exp\langle\beta,F\rangle,r\right)\right)
=\lim_{m\to\infty}\mu_m\left(\min\left(\exp\langle\beta,F\rangle,r\right)\right)\le C.
\]
As $r\to\infty$ these functions increase to $\exp\langle\beta,F\rangle$,
and the monotone convergence theorem gives $\mu(\exp\langle\beta,F\rangle)\le C$.
The last assertion is this bound for every $\beta\in C^2_{\mathrm{c}}$, with $C=\exp\left(\tfrac12\|\beta\|^2\right)$.
\end{proof}

\begin{theorem}[Comparison states]\label{thm:comparison-states}
Let $-1\le\rho\le1$, $w_0,\dots,w_3>0$, $M\in\R$ and $q>0$.
Let $\mathfrak{n}=(n_i)$ be a sequence of sides that satisfies \eqref{cond:divisibility},
and let $\mathfrak{d}=(c_0,c_+,c_3;h;O)$ be data (Definition~\ref{def:data}).
\begin{enumerate}
\item\label{it:comparison-existence}
For every $t\ge0$ the sequence $(\mu^t_i)_{i\ge i_0(\mathfrak{d})}$ of laws on $\Omega$ is tight,
$i_0(\mathfrak{d})$ being the threshold of Definition~\ref{def:threshold}.
Given any subsequence of the indices, some comparison state for $\mathfrak{d}$ is subordinate to it
(Definition~\ref{def:comparison-states}).
\item\label{it:comparison-properties}
Every comparison state $\nu_0$ for $\mathfrak{d}$ satisfies \eqref{cond:gaussian-domination}.
It is invariant under $F\mapsto g^{*}F$ for every $g\in W(\mathcal{Q}_{\mathfrak{d}})$;
in particular, it is invariant at every wall of $\mathcal{Q}_{\mathfrak{d}}$
and invariant under $\tau_{he_-}$, $\tau_{he_+}$ and $\tau_{2he_3}$.
It is reflection positive (Definition~\ref{def:reflection-positivity}) at every wall of $\mathcal{Q}_{\mathfrak{d}}$,
in particular at $(\ell_-,c_0)$.
\item\label{it:comparison-fold}
Let $(i'_m)_m$ be a subsequence of the indices,
and suppose that the laws $\mu_{i'_m}$ converge weakly on $\Omega$ to a law $\nu$ as $m\to\infty$.
Then each comparison state $\nu_0$ for $\mathfrak{d}$ that is subordinate to $(i'_m)_m$ satisfies
\[
\nu(O)\le\nu_0(O).
\]
\end{enumerate}
\end{theorem}

\begin{proof}
\eqref{it:comparison-existence}
Fix $t\ge0$.
By Proposition~\ref{prop:domination-tightness}, the laws $\mu^t_i$ with $i\ge i_0(\mathfrak{d},1)$ form a tight sequence,
$i_0(\mathfrak{d},1)$ being the threshold of Lemma~\ref{lem:coarse-structures} for boxes of radius $1$.
Only finitely many indices $i$ satisfy $i_0(\mathfrak{d})\le i<i_0(\mathfrak{d},1)$.
Let $\varepsilon>0$.
There is a compact subset of $\Omega$ whose complement has $\mu^t_i$-measure at most $\varepsilon$ for every $i\ge i_0(\mathfrak{d},1)$;
for each of the finitely many remaining indices $i$ there is such a compact set for $\mu^t_i$ alone,
because a Borel probability measure on a Polish space is inner regular with respect to compact sets.
The union of these finitely many compact sets is compact,
and its complement has $\mu^t_i$-measure at most $\varepsilon$ for every $i\ge i_0(\mathfrak{d})$.

Now let $(i'_m)_m$ be a subsequence of the indices.
Since $i'_m\to\infty$, only finitely many of its terms are smaller than $i_0(\mathfrak{d})$;
deleting them leaves a subsequence of $(i'_m)_m$ all of whose terms are at least $i_0(\mathfrak{d})$.
A subsequence of the remaining sequence is a subsequence of $(i'_m)_m$,
so a comparison state subordinate to the remaining sequence is subordinate to $(i'_m)_m$.
One may therefore assume that $i'_m\ge i_0(\mathfrak{d})$ for every $m$.
Let $t_k:=1/k$ for $k\ge1$.

\emph{Volume.}
For each $k$ the laws $\mu^{t_k}_i$ with $i\ge i_0(\mathfrak{d})$ form a tight family;
so Prokhorov's theorem, on the Polish space $\Omega$, extracts from every sequence of such laws a subsequence that converges weakly to a law on $\Omega$.
Applying this first to $(\mu^{t_1}_{i'_m})_m$ and then, for $k=1,2,\dots$ in turn, to $(\mu^{t_{k+1}}_{i^{(k)}_m})_m$,
one obtains subsequences $(i^{(k)}_m)_m$ of $(i'_m)_m$ and laws $\nu_{t_k}$ on $\Omega$
such that $(i^{(k+1)}_m)_m$ is a subsequence of $(i^{(k)}_m)_m$ and $\mu^{t_k}_{i^{(k)}_m}\to\nu_{t_k}$ weakly as $m\to\infty$, for every $k\ge1$.
The diagonal sequence $i_m:=i^{(m)}_m$ is a subsequence of $(i'_m)_m$,
and for every $k$ the sequence $(i_m)_{m\ge k}$ is a subsequence of $(i^{(k)}_m)_m$.
Consequently $\mu^{t_k}_{i_m}\to\nu_{t_k}$ weakly as $m\to\infty$, for every $k$.

\emph{Strength.}
For a plaquette $\sigma$ let $1_\sigma$ be the two-cochain with the value $1$ at $\sigma$ and $0$ at all other plaquettes,
so that $\langle\pm1_\sigma,F\rangle=\pm F_\sigma$ and $\|\pm1_\sigma\|^2=1$.
Every plaquette $\sigma$ lies in a box of radius $1$, and this box contains the support of $\pm1_\sigma$.
The number $k(\mathfrak{d},1)$ of Lemma~\ref{lem:coarse-structures} depends on $\mathfrak{n}$, $h$ and the radius only;
in particular it is the same for all plaquettes.
For each $k$ the law $\nu_{t_k}$ is a weak limit point of $(\mu^{t_k}_i)_i$, and $t_k\le1$.
Proposition~\ref{prop:domination-tightness}, applied with $\beta=\pm1_\sigma$ and $r=1$, therefore gives
\[
\nu_{t_k}\left(\exp(F_\sigma)\right)+\nu_{t_k}\left(\exp(-F_\sigma)\right)
\le2\exp\left(\tfrac12+2|O|_1\cdot k(\mathfrak{d},1)^3\right)
\]
for all $k$ and all plaquettes $\sigma$.
Lemma~\ref{lem:closed-conditions}\eqref{it:closed-tightness} now shows that the laws $\nu_{t_k}$, $k\ge1$, form a tight sequence,
and Prokhorov's theorem provides indices $k_1<k_2<\cdots$ and a law $\nu_0$ on $\Omega$ such that $\nu_{t_{k_j}}\to\nu_0$ weakly as $j\to\infty$.
The subsequence $(i_m)_m$ of $(i'_m)_m$, the numbers $t_{k_j}>0$, which tend to $0$, and the laws $\nu_{t_{k_j}}$
have the properties required in Definition~\ref{def:comparison-states}.
So $\nu_0$ is a comparison state for $\mathfrak{d}$ that is subordinate to $(i'_m)_m$.

\eqref{it:comparison-properties}
Let $\nu_0$ be a comparison state for $\mathfrak{d}$,
and let $(i_m)_m$, $t_k$ and $\nu_{t_k}$ be as in Definition~\ref{def:comparison-states}.
By Proposition~\ref{prop:domination-tightness}, $\nu_0$ satisfies \eqref{cond:gaussian-domination}.
Since $(i_m)_m$ is a subsequence of the indices, $i_m\ge i_0(\mathfrak{d})$ for all large $m$;
for each $k$ the laws $\mu^{t_k}_{i_m}$ with $i_m\ge i_0(\mathfrak{d})$ thus form a tail of $(\mu^{t_k}_{i_m})_m$,
and this tail converges weakly to $\nu_{t_k}$ as $m\to\infty$.

\emph{Invariance.}
Let $g\in W(\mathcal{Q}_{\mathfrak{d}})$.
Every law $\mu^{t_k}_{i_m}$ with $i_m\ge i_0(\mathfrak{d})$ is invariant under $F\mapsto g^{*}F$ (Proposition~\ref{prop:wall-positivity}\eqref{it:wall-invariance}).
Lemma~\ref{lem:closed-conditions}\eqref{it:closed-invariance}, applied at fixed $k$ to the tail of $(\mu^{t_k}_{i_m})_m$ formed by these laws,
shows that every $\nu_{t_k}$ is invariant;
applied once more, as $k\to\infty$, it shows that $\nu_0$ is invariant.
This applies in particular to the reflections in the walls of $\mathcal{Q}_{\mathfrak{d}}$, which generate $W(\mathcal{Q}_{\mathfrak{d}})$,
and to $\tau_{he_-}$, $\tau_{he_+}$ and $\tau_{2he_3}$, which belong to $W(\mathcal{Q}_{\mathfrak{d}})$
(Lemma~\ref{lem:wall-group}\eqref{it:walls-translations}).

\emph{Positivity.}
Let $a\in\{1,2,3\}$, let $(\ell^a,c)$ be a wall of $\mathcal{Q}_{\mathfrak{d}}$,
and put $\theta:=\theta_{\ell^a,c}$ and $\Theta:=\Theta_{\ell^a,c}$;
the wall $(\ell_-,c_0)$ is one of these.
Let $\Phi$ be a bounded continuous element of $\mathcal{O}^{+}(\ell^a,c)$;
it is a function of finitely many coordinates $F_\sigma$, each $\sigma$ a plaquette all of whose corners satisfy $\ell^a>c$.
Every point of a plaquette is a convex combination of its four corners,
and the linear function $\ell^a$ takes at such a point the same convex combination of its values at the corners;
hence $\ell^a>c$ at every point of these finitely many plaquettes.
Their union is bounded, so there is a real number $r$ with $\ell^a<c+r$ at every point of it.
Let $i\ge i_0(\mathfrak{d})$ with $n_i\ge2r$.
By \eqref{eq:fold-adapted} the height periods of $D(n_i)$ are $p_1=n_i$, $p_2=2n_i$, $p_3=n_i$, so that $p_a\ge n_i$.
Hence these plaquettes lie in the open slab $\{c<\ell^a<c+p_a/2\}$,
and $\Phi$ belongs to $\mathcal{A}_{\Xi_a}[c,c+p_a/2]$:
it is an admissible product without factors of Gaussian growth, and such a product embeds in every lattice
(Definitions~\ref{def:integrands} and~\ref{def:slabs}).
By \eqref{eq:setting-reflection-coefficients}, $(\Theta\Phi)(F)=\overline{\Phi(\theta^{*}F)}$.
Proposition~\ref{prop:wall-positivity}\eqref{it:wall-form} and Lemma~\ref{lem:bosonic-measure}\eqref{it:bosonic-integral} give
\[
\int\mathrm{d}\mu^{t_k}_i(F)\,\overline{\Phi(\theta^{*}F)}\cdot\Phi(F)=\mu^{t_k}_i\left(\Theta(\Phi)\cdot\Phi\right)\ge0
\]
for all these $i$ and all $k$;
the condition on $i$ depends on $\Phi$ and not on $k$, and it holds for all large $i$, since $n_i\to\infty$.
Lemma~\ref{lem:closed-conditions}\eqref{it:closed-positivity},
applied first at fixed $k$ to the same tail of $(\mu^{t_k}_{i_m})_m$ and then to the sequence $(\nu_{t_k})_k$,
yields the same inequality for every $\nu_{t_k}$ and for $\nu_0$,
for every bounded continuous $\Phi\in\mathcal{O}^{+}(\ell^a,c)$.
By the invariance shown above, $\nu_0$ is invariant at $(\ell^a,c)$,
and Lemma~\ref{lem:closed-conditions}\eqref{it:closed-continuous} shows that $\nu_0$ is reflection positive at $(\ell^a,c)$.

\eqref{it:comparison-fold}
Let $(i_m)_m$, $t_k$ and $\nu_{t_k}$ be as in Definition~\ref{def:comparison-states},
with $(i_m)_m$ a subsequence of $(i'_m)_m$.
The observable $O$ is a trigonometric polynomial, hence bounded and continuous on $\Omega$.
Fix $k$.
Proposition~\ref{prop:fold-finite}\eqref{it:finite-fold} gives
\[
\mu_{i_m}(O)\le\mu^{t_k}_{i_m}(O)+\frac{|Q_0|\cdot\Delta_{i_m}}{t_k}
\]
for every $m$ with $i_m\ge i_0(\mathfrak{d})$.
Let $m\to\infty$ at fixed $k$.
Weak convergence gives $\mu_{i_m}(O)\to\nu(O)$ and $\mu^{t_k}_{i_m}(O)\to\nu_{t_k}(O)$,
Theorem~\ref{thm:four-tori} gives $\Delta_{i_m}\to0$, and $t_k>0$ stays fixed.
So $\nu(O)\le\nu_{t_k}(O)$ for every $k$.
Letting now $k\to\infty$, one has $\nu_{t_k}(O)\to\nu_0(O)$, and the inequality $\nu(O)\le\nu_0(O)$ follows.
\end{proof}

\subsection{The Yang--Mills case}\label{sec:fold-yang-mills}

\begin{remark}[Correspondence of the steps]\label{rem:yang-mills-fold}
In the argument for Yang--Mills theory in four dimensions,
the fold inequality is derived from a chessboard estimate in one direction,
a surface bound, a lower tangent inequality and a statement on weak limits.
Its steps correspond to the arguments of this section as follows.
\begin{itemize}
\item
\emph{The arrays.}
In the Yang--Mills case the observable lies in one closed slab of the family of the mirrors,
and two arrays are formed:
the unfolded one is summed over one lattice,
the folded one over the lattice generated by the even translations parallel to the mirrors and by one translation across the slabs;
the reflections in the walls permute the terms of the folded array.
Here the observable $O$ lies in the open cell $Q_0$ of $\mathcal{Q}_{\mathfrak{d}}$ (Definition~\ref{def:data}),
each of the arrays $\widetilde O$ and $O_{\mathfrak{d}}$ of Definition~\ref{def:arrays} has one term in each open cell,
and $W(\mathcal{Q}_{\mathfrak{d}})$ permutes the terms of $O_{\mathfrak{d}}$
(Lemmas~\ref{lem:wall-group} and~\ref{lem:two-arrays}).
\item
\emph{The change of the array.}
In the Yang--Mills case the partition function with the unfolded array is bounded by the partition function with the folded array
on one and the same torus, which is adapted to the family of the mirrors.
The tool is the chessboard estimate in one direction, proved for positive functions,
and the chessboard norms of the contents of all slabs coincide by translation.
Here the corresponding inequality is Proposition~\ref{prop:fold-finite}\eqref{it:finite-change}, and it involves two tori:
$\widetilde O$ is read on $C(n_i)$ and $O_{\mathfrak{d}}$ on $D(n_i)$ (Lemma~\ref{lem:two-arrays}).
The tool is Theorem~\ref{thm:chessboard}(b).
It rests on H\"older's inequality in Schatten classes with unitary twists,
built from the translation unitaries and the fermion parity,
and it requires no positivity of the cell contents.
The cell contents $G_v=\exp(t\cdot\tau_vO)$, $v\in\Lambda_h$, all have the chessboard seminorm $\|G_0\|_{D(n_i)}$ of $G_0=\exp(t\cdot O)$,
again by translation.
\item
\emph{The lower tangent.}
In the Yang--Mills case this step treats an arbitrary Gibbs state $\nu$ that is periodic under the lattice of the unfolded array.
On the torus of the preceding step,
the logarithm of the ratio of the partition function with $t$ times the unfolded array to the partition function without it
is bounded below by $t$ times the density of the array, times the expectation of the observable under $\nu$, times the volume of the torus,
minus surface terms.
The proof of this inequality involves a box with a boundary condition:
it uses a positive local specification and Jensen's inequality for it,
and the energy of a plaquette lies between $0$ and $2$, which bounds the surface terms uniformly in the boundary condition.
Here the lower tangent is Proposition~\ref{prop:fold-finite}\eqref{it:finite-lower},
obtained from Jensen's inequality \eqref{eq:fold-lower-tangent} for the law $\mu_i$ on the cubic torus $C(n_i)$ itself.
Boxes, boundary conditions and specifications play no part in this argument,
and the law $\nu$ of Theorem~\ref{thm:comparison-states}\eqref{it:comparison-fold} enters only as a weak limit of the laws $\mu_i$.
\item
\emph{The error term.}
In the Yang--Mills case the error is of the order of the ratio of the surface to the volume;
this rests on the compactness of the fibre and on the finite range of the interaction.
Here the error term of Proposition~\ref{prop:fold-finite}\eqref{it:finite-fold} is $|Q_0|\cdot\Delta_i/t$, where
\[
\Delta_i=n_i^{-3}\cdot\left(\log Z_{C(n_i/2,2n_i,n_i)}[1]-\log Z_{C(n_i)}[1]\right)
\]
is the logarithm of the ratio of two partition functions without source, per photon site
(both tori have $n_i^3$ photon sites);
it tends to $0$ by Theorem~\ref{thm:four-tori}.
\item
\emph{The upper tangent.}
In the Yang--Mills case the logarithm of the normalised partition function with the folded array,
a convex function of $t$, is bounded by $t$ times its derivative;
this requires a positive measure.
Here the upper tangent is Proposition~\ref{prop:fold-finite}\eqref{it:finite-upper},
obtained from Jensen's inequality \eqref{eq:fold-upper-tangent} for the probability measure $\mu^t_i$:
the same inequality, without a derivative.
\item
\emph{Positivity of the perturbed measure.}
In the Yang--Mills case $t$ times the folded array is the sum of a function and of its composition with the reflection in a wall,
and the perturbed measure on the torus is reflection positive.
Here this is Proposition~\ref{prop:wall-positivity}\eqref{it:wall-factor} and~\eqref{it:wall-form},
the second of which holds for Grassmann-valued integrands as well;
that the reflected products to which the functional is applied belong to its domain is Lemma~\ref{lem:reflected-products}.
\item
\emph{The limits.}
In the Yang--Mills case the volume tends to infinity first and the strength $t$ to $0$ afterwards,
and the limit $\nu_0$ is a Gibbs state that is periodic under the lattice of the folded array;
this uses the compactness of the configuration space and the statement on weak limits.
Here Theorem~\ref{thm:comparison-states} takes the limits in the same order (Definition~\ref{def:comparison-states}).
Tightness comes from Proposition~\ref{prop:domination-tightness},
and the properties of the limits come from Lemma~\ref{lem:closed-conditions}.
In place of the Gibbs property, the comparison state $\nu_0$ has the properties listed in Theorem~\ref{thm:comparison-states}\eqref{it:comparison-properties}.
\item
\emph{Continuous test functions.}
In the Yang--Mills case continuous test functions suffice for reflection positivity, by the invariance at the mirror.
Here this is Lemma~\ref{lem:closed-conditions}\eqref{it:closed-continuous}, with the same approximation argument.
\end{itemize}

\emph{Obstructions to the lower tangent of the Yang--Mills case.}
The photon field admits no specification,
since any finite collection of links that includes every link at some site has a gauge direction that is not compact.
No bound holds for the energy $F_\sigma^2/2$ of the plaquettes on the boundary of a box.
The weight $\exp\left(-\tfrac12\|F\|^2\right)\cdot\det D^{\mathbb{T}}_A$ that remains after the Berezin integral is non-local,
whereas the weight that precedes this integral fails to be positive.
Finally, the flux sectors are of a global nature.
The proofs of Propositions~\ref{prop:fold-finite} and~\ref{prop:wall-positivity} and of Theorem~\ref{thm:comparison-states}
use no specification, no box with a boundary condition,
no locality of the weight $\exp\left(-\tfrac12\|F\|^2\right)\cdot\det D^{\mathbb{T}}_A$
and no decomposition into flux sectors.

\emph{What is lost.}
\begin{enumerate}
\item
The class of the laws $\nu$ that are compared:
in the Yang--Mills case it consists of all periodic Gibbs states;
here, in Theorem~\ref{thm:comparison-states}\eqref{it:comparison-fold}, of the weak limits of the laws $\mu_i$ of the cubic tori.
\item
In the Yang--Mills case the comparison state is a periodic Gibbs state,
and a general theorem on equilibrium states can be applied to it.
Here no class of equilibrium states within the reach of such a theorem is known to contain $\nu_0$;
for this reason the proof of Theorem~\ref{thm:no-layers} deals with the comparison states themselves,
by way of the functionals $\mu^t_i$.
\item
A second torus enters, and with it the term $\Delta_i$,
which is controlled by Proposition~\ref{prop:coordinate-pressure} and Theorem~\ref{thm:four-tori}.
\item
Lateral invariance.
In the Yang--Mills case the comparison state is periodic under the lattice generated by the even translations parallel to the mirrors
and one translation across the slabs;
since no folding takes place in the directions parallel to the mirrors,
every even translation parallel to them leaves the comparison state invariant.
The array $O_{\mathfrak{d}}$, on the other hand, is folded in all three directions (Remark~\ref{rem:three-directions}),
and what Theorem~\ref{thm:comparison-states}\eqref{it:comparison-properties} provides is the invariance of $\nu_0$
under the translations by the vectors of $2\Lambda_h$ (Lemma~\ref{lem:wall-group}\eqref{it:walls-translations});
the invariance under the translations by the vectors of $2\Z^3$ is a consequence of Theorem~\ref{thm:no-layers}.
\end{enumerate}

\emph{What is gained.}
The fold inequality of Proposition~\ref{prop:fold-finite} holds for every real $M$, whatever the sign of the lattice mass $m$,
and neither a bound on $D_A^{-1}$ nor any locality of the marginal law of the photon field enters.
Given Theorem~\ref{thm:chessboard}(b) and Theorem~\ref{thm:four-tori},
the proof of Proposition~\ref{prop:fold-finite} is an application of the two Jensen inequalities of Lemma~\ref{lem:two-jensen}.
\end{remark}

\begin{remark}[Three directions]\label{rem:three-directions}
The array $O_{\mathfrak{d}}$ is folded in all three directions of the diagonal frame,
and the perturbed functional $\mu^t_i$ is formed on $D(n_i)$.
A fold in one direction already gives an inequality of the same kind in finite volume;
the reason for the three directions lies in the limits.

Let $i\ge i_0(\mathfrak{d})$ (Definition~\ref{def:threshold}) and $t>0$.
The group $W_1(\mathcal{Q}_{\mathfrak{d}})$, the layers of $\ell_-$ and the elements $g^1_{\kappa_1\leftarrow0}$, $\kappa_1\in\Z$,
are those of Definition~\ref{def:cell-families} for the structure $\mathcal{Q}_{\mathfrak{d}}$, in which $\ell^1=\ell_-$ and $o_1=c_0$ by \eqref{eq:fold-frame}.
Thus $W_1(\mathcal{Q}_{\mathfrak{d}})$ is generated by the reflections in the walls $(\ell_-,c_0+j\cdot h)$, $j\in\Z$,
and $g^1_{\kappa_1\leftarrow0}$ is its only element that maps the layer $0$ of $\ell_-$ onto the layer $\kappa_1$.
The array folded in the direction $\ell_-$ only is the formal sum
\[
O'_{\mathfrak{d}}:=\sum_{\kappa\in\Z^3}O'_\kappa,
\qquad
O'_\kappa:=\left(\tau_vO\right)^{g^1_{\kappa_1\leftarrow0}},
\quad v\in\Lambda_h,\ \kappa_v=(0,\kappa_2,\kappa_3),
\]
with one term in each cell, $O'_\kappa$ reading in $Q_\kappa$.
On the layer $\kappa_1=0$ it coincides with the unfolded array $\widetilde O$,
and its other terms are the images of these under $W_1(\mathcal{Q}_{\mathfrak{d}})$.

Two kinds of maps permute the terms of $O'_{\mathfrak{d}}$.
First, let $v'\in\Lambda_h$ satisfy $\ell_-(v')=0$.
By Definition~\ref{def:mirrors}\eqref{it:mirror-conjugation}, $\tau_{v'}$ commutes with the reflections in the walls of $\ell_-$,
and hence with every element of $W_1(\mathcal{Q}_{\mathfrak{d}})$;
as $\left(\tau_vO\right)^{\tau_{v'}}=\tau_{v+v'}O$ and the first entry of $\kappa_{v'}$ vanishes,
\eqref{eq:fold-composition} gives $\left(O'_\kappa\right)^{\tau_{v'}}=O'_{\kappa+\kappa_{v'}}$.
Secondly, let $g\in W_1(\mathcal{Q}_{\mathfrak{d}})$.
It leaves $\ell_+$ and $x_3$ invariant and maps the layer $\kappa_1$ of $\ell_-$ onto a layer $\kappa'_1$
(Definitions~\ref{def:slab-copies} and~\ref{def:cell-families});
so $g\,g^1_{\kappa_1\leftarrow0}=g^1_{\kappa'_1\leftarrow0}$ by uniqueness, and $\left(O'_\kappa\right)^g=O'_{(\kappa'_1,\kappa_2,\kappa_3)}$.
As $O'_\kappa$ is the image of $O'_0=O$ under $g^1_{\kappa_1\leftarrow0}\,\tau_v$, maps of the two kinds carry $O$ to every term.
The period lattice $\Gamma_i$ of $D(n_i)$ (Lemma~\ref{lem:two-arrays}) is also that of $D'(n_i)$ (Definition~\ref{def:tori}).
It is generated by $\frac{n_i}{2}\cdot e_-$, $n_i\cdot e_+$ and $n_i\cdot e_3$,
whose indices $\kappa_\gamma$ are $(n_i/h,0,0)$, $(0,2n_i/h,0)$ and $(0,0,n_i/h)$
(Lemma~\ref{lem:wall-group}\eqref{it:walls-lattice}, \eqref{it:walls-tori}).
The last two generators are vectors $v'$ of the first kind.
By Definition~\ref{def:mirrors}\eqref{it:mirror-product}, $\theta_{\ell_-,c_0+h}\,\theta_{\ell_-,c_0}=\tau_{h\cdot e_-}$;
as $4h$ divides $n_i$ (Definition~\ref{def:threshold}), the translation by the first generator is
$\left(\theta_{\ell_-,c_0+h}\,\theta_{\ell_-,c_0}\right)^{\frac{n_i}{2h}}$,
an element of $W_1(\mathcal{Q}_{\mathfrak{d}})$ that maps the layer $\kappa_1$ onto the layer $\kappa_1+n_i/h$.
The vectors $\gamma\in\Lambda_h$ with $\left(O'_\kappa\right)^{\tau_\gamma}=O'_{\kappa+\kappa_\gamma}$ for all $\kappa\in\Z^3$ form a group,
by \eqref{eq:fold-composition} and the additivity of $\gamma\mapsto\kappa_\gamma$;
it contains the three generators, hence $\Gamma_i$.
For a $\Gamma_i$-periodic $F$ and for $\gamma\in\Gamma_i$ it follows that
\[
O'_{\kappa+\kappa_\gamma}(F)=O'_\kappa\left(\tau_\gamma^*F\right)=O'_\kappa(F)\qquad(\kappa\in\Z^3).
\]
On $D'(n_i)$ and on $D(n_i)$ the array $O'_{\mathfrak{d}}$ is read as the sum of its terms over a set of representatives of the cells of the torus;
by the last identity, this sum does not depend on the set of representatives.

The first of the three steps of Theorem~\ref{thm:chessboard}(b) leads from $C(n_i)$ to $D'(n_i)$.
Its intermediate inequality, applied to the cell contents $G_v=\exp(t\cdot\tau_vO)$, $v\in\Lambda_h$,
of the proof of Proposition~\ref{prop:fold-finite}\eqref{it:finite-change}, reads
\[
Z_{C(n_i)}\left[\exp(t\cdot\widetilde O)\right]\le\prod_{j=1}^{\frac{n_i}{h}}Z_{D'(n_i)}\left[\fold_{\ell_-}S_j\right]^{\frac{h}{n_i}},
\]
$S_j$ being the content of the $j$-th slab of $\ell_-$,
that is, the product of the $G_v$ with $\kappa_v$ in the brick $\mathcal{K}_i$ of Lemma~\ref{lem:two-arrays} and $(\kappa_v)_1=j-1$;
the left side is positive (Lemma~\ref{lem:bosonic-measure}), hence equal to its modulus.

The factors on the right do not depend on $j$.
Let $1\le j\le n_i/h$, and put $u:=\frac{(j-1)\cdot h}{2}\cdot e_-$.
Then $u\in\Lambda_h$ (Definition~\ref{def:arrays}) and $\kappa_u=(j-1,0,0)$ (Lemma~\ref{lem:wall-group}\eqref{it:walls-lattice}).
With $G_0=\exp(t\cdot O)$ one has $G_v=G_0^{\tau_v}$, so that $G_v^{\tau_u}=G_{u+v}$ by \eqref{eq:fold-composition}, and $\kappa_{u+v}=\kappa_v+\kappa_u$.
The shift $\kappa\mapsto\kappa+\kappa_u$ maps the points of $\mathcal{K}_i$ with first entry $0$ onto those with first entry $j-1$;
Lemma~\ref{lem:admissible-images}\eqref{it:images-products} therefore gives $S_j=S_1^{\tau_u}$.
For every position $c$ one has $\tau_u\,\theta_{\ell_-,c}\,\tau_u^{-1}=\theta_{\ell_-,c+(j-1)\cdot h}$ (Definition~\ref{def:mirrors}\eqref{it:mirror-conjugation});
so conjugation by $\tau_u$ maps $W_1(\mathcal{Q}_{\mathfrak{d}})$ onto itself.
With the slabs of $\ell_-$ numbered as in Definition~\ref{def:slab-copies} for the planes $c_0+j'\cdot h$, $j'\in\Z$,
the translation $\tau_u$ maps the slab $j'$ onto the slab $j'+j-1$, for every $j'\in\Z$.
If $g\in W_1(\mathcal{Q}_{\mathfrak{d}})$ maps the slab $1$ onto the slab $j'$,
then $\tau_u\,g\,\tau_u^{-1}$ is the element of $W_1(\mathcal{Q}_{\mathfrak{d}})$ that maps the slab $j$ onto the slab $j'+j-1$,
and, by \eqref{eq:fold-composition}, the copy of $S_j$ in the slab $j'+j-1$ is
\[
\left(S_1^{\tau_u}\right)^{\tau_u\,g\,\tau_u^{-1}}=S_1^{\tau_u\,g}=\left(S_1^{g}\right)^{\tau_u},
\]
the image under $\tau_u$ of the copy of $S_1$ in the slab $j'$.
By Definition~\ref{def:slab-copies}, $\fold_{\ell_-}S_j$ defines the same function on $\Conf(\Gamma_i)$
as the product of the copies of $S_j$ in the $n_i/h$ consecutive slabs $j,\dots,j+n_i/h-1$;
by the last display and Lemma~\ref{lem:admissible-images}\eqref{it:images-products}, that product is $\left(\fold_{\ell_-}S_1\right)^{\tau_u}$.
The translation $\tau_u$ is even, and $\tau_u\,D'(n_i)=D'(n_i)$ (Lemma~\ref{lem:wall-group}\eqref{it:walls-tori});
so Proposition~\ref{prop:symmetries}\eqref{it:functional-covariance} gives
\[
Z_{D'(n_i)}\left[\fold_{\ell_-}S_j\right]=Z_{D'(n_i)}\left[\left(\fold_{\ell_-}S_1\right)^{\tau_u}\right]=Z_{D'(n_i)}\left[\fold_{\ell_-}S_1\right].
\]
The copy of $S_1$ in the slab $j'$, which is the layer $j'-1$ of $\ell_-$, is $S_1^{g^1_{j'-1\leftarrow0}}$;
since $\tau_vO$ is real, it is the product of the contents $\exp(t\cdot O'_\kappa)$ over the $\kappa\in\mathcal{K}_i$ with $\kappa_1=j'-1$.
As $\mathcal{K}_i$ represents the cells of $D'(n_i)$ (Theorem~\ref{thm:chessboard}(b)),
$\fold_{\ell_-}S_1=\exp(t\cdot O'_{\mathfrak{d}})$ on $D'(n_i)$.
Each of the $n_i/h$ factors on the right of the first inequality thus equals $Z_{D'(n_i)}\left[\exp(t\cdot O'_{\mathfrak{d}})\right]^{\frac{h}{n_i}}$, and
\[
Z_{C(n_i)}\left[\exp(t\cdot\widetilde O)\right]\le Z_{D'(n_i)}\left[\exp(t\cdot O'_{\mathfrak{d}})\right].
\]
The perturbed functional on $D'(n_i)$ is the map
\[
G\longmapsto\frac{Z_{D'(n_i)}\left[\exp(t\cdot O'_{\mathfrak{d}})\cdot G\right]}{Z_{D'(n_i)}\left[\exp(t\cdot O'_{\mathfrak{d}})\right]}
\]
on the elements $G\in\mathcal{A}(\R^3)$ that embed in $D'(n_i)$; its denominator is positive by the last inequality.
Lemma~\ref{lem:two-jensen}, applied to that inequality as in the proof of Proposition~\ref{prop:fold-finite}\eqref{it:finite-fold},
gives
\[
\mu_i(O)\le\frac{Z_{D'(n_i)}\left[\exp(t\cdot O'_{\mathfrak{d}})\cdot O\right]}{Z_{D'(n_i)}\left[\exp(t\cdot O'_{\mathfrak{d}})\right]}+\frac{|Q_0|\cdot\Delta_i}{t}.
\]
Indeed, the maps of the two kinds considered above map $D'(n_i)$ onto itself,
since their linear parts fix $e_+$ and $e_3$ and map $e_-$ to $\pm e_-$ (Definition~\ref{def:tori}),
and they carry $O$ to every term of $O'_{\mathfrak{d}}$ while permuting these terms;
as in the computation \eqref{eq:fold-invariance}, every term of $O'_{\mathfrak{d}}$ therefore has the same expectation as $O$ under the perturbed functional on $D'(n_i)$;
and $\log Z_{D'(n_i)}[1]-\log Z_{C(n_i)}[1]\le n_i^3\cdot\Delta_i$ by \eqref{eq:tori-chain}.
Since $D'(n_i)$ is adapted to $\ell_-$ (Definition~\ref{def:slabs}),
the walls $(\ell_-,c_0+j\cdot h)$, $j\in\Z$, are mirrors of this functional,
in the sense in which Proposition~\ref{prop:wall-positivity} makes the walls of $\mathcal{Q}_{\mathfrak{d}}$ mirrors of $\mu^t_i$.
This is the inequality of Proposition~\ref{prop:fold-finite}\eqref{it:finite-fold} with one H\"older step in place of three.

This inequality is not enough for the limits.
The tightness of Proposition~\ref{prop:domination-tightness} and the amplification of Theorem~\ref{thm:amplification}
both come from Theorem~\ref{thm:chessboard}(a).
That theorem is stated for a torus that carries the cell structure,
a torus which is then adapted to all three families of the frame (Definition~\ref{def:cell-structures}),
and it is applied there with a positive density of cells.
Of the tori $D(n_i)$ and $D'(n_i)$, only $D(n_i)$ carries $\mathcal{Q}_{\mathfrak{d}}$; $D'(n_i)$ fails to be adapted to $\ell_+$ (Definition~\ref{def:slabs}).
Both arguments rest, moreover, on an array that is folded completely:
in step~\eqref{it:amplification-chessboard} of the proof of Theorem~\ref{thm:amplification} one cell is singled out,
and the content of each of the other cells is a mirror copy of a single content;
the chessboard seminorms of these other contents therefore multiply to a power of the normalisation $Z_{D(n_i)}\left[\exp(t\cdot O_{\mathfrak{d}})\right]$.

For the array $O'_{\mathfrak{d}}$, folded in the direction $\ell_-$ only, this argument does not give tightness.
Read $O'_{\mathfrak{d}}$ on $D(n_i)$,
and let $\beta\in C^2_{\mathrm{c}}$ be such that $\exp\langle\beta,F\rangle\in\mathcal{A}(Q_\kappa)$ for one open cell $Q_\kappa$.
The content $\exp(t\cdot O'_{\kappa'})$ of every other cell $Q_{\kappa'}$ is a mirror copy of one of the $G_v$ (Definition~\ref{def:folds}),
so that its chessboard seminorm equals that of $G_v$, which is $\|G_0\|_{D(n_i)}$ (Lemma~\ref{lem:cell-symmetries});
the power of $\|G_0\|_{D(n_i)}$ with the exponent $n_i^3/|Q_0|$, the number of cells, is $Z_{D(n_i)}\left[\exp(t\cdot O_{\mathfrak{d}})\right]$
(Proposition~\ref{prop:fold-finite}\eqref{it:finite-change}).
Theorem~\ref{thm:chessboard}(a) therefore gives
\[
Z_{D(n_i)}\left[\exp(t\cdot O'_{\mathfrak{d}})\cdot\exp\langle\beta,F\rangle\right]
\le Z_{D(n_i)}\left[\exp(t\cdot O_{\mathfrak{d}})\right]^{1-\frac{|Q_0|}{n_i^3}}
\cdot\left\|\exp(t\cdot O'_\kappa)\cdot\exp\langle\beta,F\rangle\right\|_{D(n_i)}.
\]
After division by $Z_{D(n_i)}\left[\exp(t\cdot O'_{\mathfrak{d}})\right]$ the ratio
\[
\frac{Z_{D(n_i)}\left[\exp(t\cdot O_{\mathfrak{d}})\right]}{Z_{D(n_i)}\left[\exp(t\cdot O'_{\mathfrak{d}})\right]}\ge1
\]
appears, the inequality being the case $\beta=0$ of the same estimate,
and the argument gives no bound on this ratio.
The obstacle lies in the ratio and not in the torus:
by Proposition~\ref{prop:symmetries}\eqref{it:sign-character},
$Z_{D'(n_i)}[G]=Z_{D(n_i)}[G]$ for every $G\in\mathcal{A}(\R^3)$ that embeds in the common period lattice of the two tori,
so that the chessboard estimate of Theorem~\ref{thm:chessboard}(a) holds for $Z_{D'(n_i)}$ as well.
\end{remark}

\section{Amplification and two layers}\label{sec:amplification}

Let $c$ be a link plane of the direction $e_1$, let $K$ be the crossing operator at $c$ (Definition~\ref{def:slice-operators}), and let $X$ be an element of the mid-plane algebra $\mathfrak{W}(S,\Lambda)$ (Definition~\ref{def:mid-plane}).
Put between the two factors of $K=K^{\frac12}K^{\frac12}$, the operator $X$ defines an integrand $G_X$, the insertion of $X$ (Definition~\ref{def:insertions}), and this integrand is \emph{local}: it reads in a bounded box.
On a coordinate torus $C(n)$, the functional $Z_{C(n)}$ of a character times a product of insertions $G_{X_1},\dots,G_{X_r}$ at link planes of $C(n)$, in suitably separated boxes, is bounded in modulus by $\|X_1\|_\#\cdots\|X_r\|_\#\cdot Z_{C(n)}[1]$
(Propositions~\ref{prop:holder-insertions}, \ref{prop:insertions} and~\ref{prop:bounded-elements}).

The perturbed functional $\mu^t_i(G)=Z_{D(n_i)}[\exp(t\cdot O_{\mathfrak{d}})\cdot G]/Z_{D(n_i)}[\exp(t\cdot O_{\mathfrak{d}})]$ of Definition~\ref{def:data} is formed on the adapted torus $D(n_i)$, and on this torus a single insertion is amplified to many.
By the chessboard estimate in the diagonal frame (Theorem~\ref{thm:chessboard}), $\left|\mu^t_i(G_X)\right|$ raised to the power $n_i^3/(|Q_0|\,k^3)$ is at most the value of $\mu^t_i$ at the array that carries one mirror copy of $G_X$ in each of the $n_i^3/(|Q_0|\,k^3)$ cells of a coarser diagonal cell structure of width $kh$;
here $k$ is an integer determined by the data and by the radius of the box of $G_X$.
The chessboard estimate between the diagonal and the coordinate frame carries this array to the elongated coordinate torus $C(n_i/2,2n_i,n_i)$, where it is bounded by $\|X\|_\#^{\frac{n_i^3}{|Q_0|\,k^3}}\cdot Z_{C(n_i/2,2n_i,n_i)}[1]$,
and so is its product with any character.
On $D(n_i)$ the folded array $O_{\mathfrak{d}}$ is the sum of $n_i^3/|Q_0|$ mirror copies of $O$, each a combination of characters whose coefficients have moduli adding up to $|O|_1$ (Definition~\ref{def:arrays});
in particular $\left|O_{\mathfrak{d}}\right|\le|O|_1\cdot n_i^3/|Q_0|$.
Accordingly the denominator of $\mu^t_i$ is at least $\exp\left(-t\cdot|O|_1\cdot n_i^3/|Q_0|\right)\cdot Z_{D(n_i)}[1]$ (Definition~\ref{def:data}),
and in its numerator the factor $\exp(t\cdot O_{\mathfrak{d}})$ is the sum of an absolutely convergent series of characters, the moduli of whose coefficients add up to $\exp\left(t\cdot|O|_1\cdot n_i^3/|Q_0|\right)$.
Numerator and denominator thus contribute at most a factor $\exp\left(t\cdot|O|_1\cdot n_i^3/|Q_0|\right)$ each,
and the root of order $n_i^3/(|Q_0|\,k^3)$ turns the product of these two factors into $\exp\left(2t\cdot|O|_1\cdot k^3\right)$.
The quotient $Z_{C(n_i/2,2n_i,n_i)}[1]/Z_{D(n_i)}[1]$ lies between $1$ and $\exp\left(n_i^3\cdot\Delta_i\right)$, where $\Delta_i\to0$ (Theorem~\ref{thm:four-tori});
so it is $\exp\left(o\left(n_i^3\right)\right)$, while its root of order $n_i^3/(|Q_0|\,k^3)$ is at most $\exp\left(|Q_0|\cdot k^3\cdot\Delta_i\right)=\exp(o(1))$.
It follows that the linear functional $X\mapsto\mu^t_i(G_X)$ is unital and, for all large $i$, has norm at most
\[
\gamma_{t,i}(k)=\exp\left(2t\cdot|O|_1\cdot k^3+|Q_0|\cdot k^3\cdot\Delta_i\right)
\]
with respect to $\|\cdot\|_\#$ (Theorem~\ref{thm:amplification});
moreover $\lim_{t\to0}\limsup_{i\to\infty}\gamma_{t,i}(k)=1$ (Definition~\ref{def:amplification-factor}).

H\"older's inequality with insertions, Proposition~\ref{prop:holder-insertions}, has its counterpart in the four-dimensional Yang--Mills argument.
The two arguments differ at the step from that inequality to a state.
In the Yang--Mills case this step rests on Jensen's inequality for the logarithm of a kernel and therefore requires the kernel to be strictly positive;
here amplification takes its place, and no logarithm occurs.

If a unital linear functional $\mathfrak{o}$ on a C$^*$-algebra has norm at most $\gamma_{t,i}(k)$, then
\[
\left|\mathfrak{o}(X^*)-\overline{\mathfrak{o}(X)}\right|\le4\left(\gamma_{t,i}(k)^2-1\right)^{\frac{1}{2}}\cdot\|X\|
\]
for every $X$.
Applied to $X\mapsto\mu^t_i(G_X)$ and combined with the identity $\Theta_cG_X=G_{X^*}$ (Proposition~\ref{prop:insertions}), this is the two-layer principle for insertions (Theorem~\ref{thm:two-layer-principle}), and it holds also for elements $X$ with a non-trivial Fock part.

A two-layer element $Y\in\mathcal{P}_1(c)$ (Definition~\ref{def:two-layer}) is in general not the insertion of a bounded operator into $K$.
It is, however, the insertion of a bounded operator into the comparison operator $K_{0,S}$, whose kernel is pointwise at most $2^{\frac{\#|S|}{2}}$ times that of $K$.
This yields an approximation of $Y$ by insertions $G_X$ (Proposition~\ref{prop:two-layer-insertions}) and the two-layer principle for $Y$, with the limits taken in the order $i\to\infty$, then $t\to0$ at fixed $\varepsilon$, then $\varepsilon\to0$ (Theorem~\ref{thm:two-layer-principle}).

\paragraph{Standing assumptions.}
Throughout this section the hopping weights $w_0,\dots,w_3>0$, the mass $M\in\R$ and the lattice charge $q>0$ are arbitrary.
The Wilson parameter satisfies $0<\rho<1$, except in statements that name another range of $\rho$ among their hypotheses;
several conclusions remain true for $-1\le\rho\le1$, among them everything that concerns the purely photonic elements of the mid-plane algebra, whose Fock part is the identity (Remark~\ref{rem:wilson-range}).
The sequence of sides $\mathfrak{n}=(n_i)$ satisfies \textup{(Div)}, and $\mathfrak{d}=(c_0,c_+,c_3;h;O)$ are data for $\mathfrak{n}$ (Definition~\ref{def:data}).
Condition \textup{(Div)} is imposed on the sequence with which the adapted tori $D(n_i)$, the data and the functionals $\mu^t_i$ are formed.
Let the sides of a sequence of cubic tori be arbitrary integers at least $4$ that tend to infinity.
A companion $\mathfrak{n}'=(n'_i)$ of this sequence of sides satisfies \textup{(Div)} (Definition~\ref{def:companions}),
so that the results of this section hold with $\mathfrak{n}'$ in the place of $\mathfrak{n}$ and with data for $\mathfrak{n}'$.
Every object of this section that carries the index $i$ is then formed with the side $n'_i$:
the adapted torus $D(n'_i)$, the elongated torus $C(n'_i/2,2n'_i,n'_i)$, the cubic torus $C(n'_i)$ and with these tori the number $\Delta_i$, the functional $\mu^t_i$ and the factor $\gamma_{t,i}(k)$.

\paragraph{Positivity of the measure.}
Let $\mathbb{T}=(\Gamma,\chi)$ be a torus.
For all real $M$ and $\rho$, Corollary~\ref{cor:modulus-bound} gives $\det D^{\mathbb{T}}_A\ge0$ for every real $\Gamma$-periodic one-cochain $A$;
moreover, for every bosonic $G$ for which $Z_{\mathbb{T}}[G]$ is defined, it represents $Z_{\mathbb{T}}[G]$ as the integral of $G$ over $\Conf(\Gamma)$ against the measure $\exp\left(-\tfrac12\|F\|^2\right)\cdot\det D^{\mathbb{T}}_A\,\mathrm{d}(A,z)$,
which is therefore non-negative; here $\mathrm{d}(A,z)$ is the measure of $\Conf(\Gamma)$.
Consequently $\left|Z_{\mathbb{T}}[G]\right|\le Z_{\mathbb{T}}\left[|G|\right]$ whenever the right-hand side is finite, and
\begin{equation}\label{eq:amplification-monotone}
\left|Z_{\mathbb{T}}[G]\right|\le Z_{\mathbb{T}}[G']
\end{equation}
for all bosonic $G$ and $G'$ with $|G|\le G'$ for which $Z_{\mathbb{T}}[G]$ and $Z_{\mathbb{T}}[G']$ are defined;
the right-hand side may be infinite.

\paragraph{Conventions.}
Unless another direction is named, time is the direction $e_1$, and a link plane is a plane $\{x_1=c\}$ with $c\in\Z+\frac{1}{2}$.
The two other coordinate directions are reached through the matrices $g_\mu\in\Ocub$ of Definition~\ref{def:other-directions}:
$g_1=1$, and for $\mu=2,3$ the matrix $g_\mu$ permutes the unit vectors $e_1,e_2,e_3$ cyclically and takes $e_1$ to $e_\mu$.
Being of determinant one, these matrices are even elements of $W$, so that $G\mapsto G^{g_\mu}$ is linear;
and Definition~\ref{def:other-directions} transports by $g_\mu$ whatever is attached to the plane $(x_1,c)$ to the plane $(x_\mu,c)$.
For the insertion of $X$ at $(x_\mu,c)$ the notation is
\[
G^{(\mu,c)}_X:=\left(G^{(c)}_X\right)^{g_\mu},\qquad\mu\in\{1,2,3\},
\]
with $G^{(c)}_X$ the insertion of Definition~\ref{def:insertions}; thus $G^{(1,c)}_X=G^{(c)}_X$.

\begin{definition}[Charged Fock operators]\label{def:charged-operators}
Let $\Lambda$ be a finite set of sites of the slice lattice $\Z^2$.
The Fock space of the modes over $\Lambda$ at a link plane is
\[
\mathcal{F}_\Lambda:=\mathcal{F}\left(\mathcal{J}^{\mathrm{l}}(\Lambda)\right),
\]
with the orthonormal basis $(e_I)$, $I\subset\mathcal{J}^{\mathrm{l}}(\Lambda)$ (Definition~\ref{def:modes}).
For $a\in\R^\Lambda$ let $u(a)$ be the diagonal unitary on $\mathcal{F}_\Lambda$ with $u(a)e_I=\exp(iq\langle a,\mathfrak{q}(I)\rangle)\cdot e_I$, where $\mathfrak{q}(I)\in\Z^\Lambda$ is the charge of $I$.
For $z\in\Z^\Lambda$, an operator $\Upsilon$ on $\mathcal{F}_\Lambda$ has \emph{charge} $z$ if
\[
u(a)\,\Upsilon\,u(a)^{-1}=\exp\left(iq\langle a,z\rangle\right)\cdot\Upsilon\qquad\text{for all }a\in\R^\Lambda,
\]
and $\End_z(\Lambda)$ denotes the space of the operators on $\mathcal{F}_\Lambda$ of charge $z$.
\end{definition}

\begin{definition}[Integral kernels]\label{def:integral-kernels}
Fix a finite set $S$ of links of $\Z^2$ and a finite set $\Lambda$ of sites of $\Z^2$ that contains the end points of the links of $S$,
and let $\mathcal{H}_S=L^2(\R^S)\otimes\mathcal{F}_\Lambda$ (Definition~\ref{def:mid-plane}); its elements are functions $\Phi:\R^S\to\mathcal{F}_\Lambda$.
Vectors $\delta\in\R^S$ are regarded as functions on the links of $\Z^2$, or on those of $\Lambda_{\mathrm{s}}$, vanishing outside $S$.
On $\mathcal{H}_S$ let $E_S=(E_b)_{b\in S}$, $E_b=-i\,\partial/\partial A(b)$.
The variable $\xi\in\R^S$ is the spectral variable of $E_S$:
after the Fourier transform $\Phi\mapsto\int_{\R^S}\mathrm{d}A\,\exp(-i\langle\xi,A\rangle)\cdot\Phi(A)$, the operator $E_b$ acts as the multiplication by $\xi_b$,
and for a bounded Borel function $f$ on $\R^S$ with values in $\C$ or in $\End\mathcal{F}_\Lambda$ the operator $f(E_S)$ acts as the multiplication by $f(\xi)$.
For $f\in L^1(\R^S)$, again with values in $\C$ or in $\End\mathcal{F}_\Lambda$, the inverse Fourier transform is
\[
f^{\vee}(x):=(2\pi)^{-\#|S|}\int_{\R^S}\mathrm{d}\xi\,f(\xi)\cdot\exp\left(i\langle\xi,x\rangle\right),\qquad x\in\R^S.
\]
An operator $T$ on $\mathcal{H}_S$ \emph{has the kernel} $k$, where $k:\R^S\times\R^S\to\End\mathcal{F}_\Lambda$ is measurable, if
\[
(T\Phi)(A')=\int_{\R^S}\mathrm{d}A\,k(A',A)\,\Phi(A)
\]
for every $\Phi\in\mathcal{H}_S$ and almost every $A'\in\R^S$;
a complex value of $k$ stands for that multiple of the identity of $\mathcal{F}_\Lambda$.
Since the Fourier transform turns convolution into multiplication, for a bounded $f\in L^1(\R^S)$ the operator $f(E_S)$ has the kernel $f^{\vee}(A'-A)$,
and for $\delta\in\R^S$ the operator $\exp(i\langle\delta,A\rangle)\,f(E_S)$ has the kernel $\exp(i\langle\delta,A'\rangle)\cdot f^{\vee}(A'-A)$.
For instance, for a real number $a>0$ the operator $\exp\left(-\tfrac{a}{2}|E_S|^2\right)$, where $|E_S|^2=\sum_{b\in S}E_b^2$, has the kernel
\[
\prod_{b\in S}g_a\left(A'(b)-A(b)\right),\qquad g_a(s)=(2\pi a)^{-\frac12}\cdot\exp\left(-\frac{s^2}{2a}\right);
\]
for $S=\{b\}$ this is the kernel $g_a(A'(b)-A(b))$ of $\exp\left(-\tfrac{a}{2}E_b^2\right)$.
\end{definition}

For an operator $\Upsilon$ on a Fock space $\mathcal{F}(\mathcal{J})$, $\overline{\Upsilon}$ denotes the operator whose matrix in the basis $(e_I)$ is the entrywise complex conjugate of the matrix of $\Upsilon$.

\paragraph{Localisation on a cross-section.}
Let $\Xi=(\Gamma_\parallel,\chi_\parallel)$ be a rectangular cross-section (Definition~\ref{def:slices}), $\Pi=(x_1,c)$ a link plane, and $\Lambda$ a finite set of sites of $\Z^2$ such that $\hull(\Lambda)$ embeds in $\Xi$.
The modes $\mathcal{J}^{\mathrm{l}}(\Lambda)$ are identified with a subset of $\mathcal{J}^{\mathrm{l}}(\Lambda_{\mathrm{s}})=\idx_\Xi(\Pi)$ as in the word formula, Proposition~\ref{prop:word-formula}.
Each mode $i$ over $\Lambda$ is the translate of the representative $r(i)$ of its class by a unique element of $\Gamma_\parallel$.
If the map $r$ is not the identity of $\mathcal{J}^{\mathrm{l}}(\Lambda)$, that is, if some mode over $\Lambda$ is not itself a representative, then this identification involves signs and a relabelling by $r$, which need not be monotone on $\mathcal{J}^{\mathrm{l}}(\Lambda)$.
It is given by the unitary $U_\Lambda$ from $\mathcal{F}_\Lambda$ onto the Fock space of $r(\mathcal{J}^{\mathrm{l}}(\Lambda))$ with $U_\Lambda e_I=s(I)\cdot e_{r(I)}$;
here the sign $s(I)$ is the product of the values of $\chi_\parallel$ at the translations that carry $r(i)$ to $i$, $i\in I$, and of the sign of the permutation that sorts $(r(i))_{i\in I}$ into the order of $\mathcal{J}^{\mathrm{l}}(\Lambda_{\mathrm{s}})$,
and $U_\Lambda$ is the identity if $r$ is (Definition~\ref{def:modes}).
The localisation $\iota$ in the realisation of the mid-plane algebra on $\Xi$ (Definition~\ref{def:mid-plane}) and in the word formula includes this unitary:
for an operator $\Upsilon$ on $\mathcal{F}_\Lambda$, $\iota(\Upsilon)$ stands for $\iota(U_\Lambda\Upsilon U_\Lambda^*)$,
where $\iota$ on the right is the localisation of Definition~\ref{def:modes} from the Fock space of $r(\mathcal{J}^{\mathrm{l}}(\Lambda))$ to $\mathcal{F}(\mathcal{J}^{\mathrm{l}}(\Lambda_{\mathrm{s}}))$.

By Proposition~\ref{prop:localisation}, $\iota$ is a unital injective $*$-homomorphism \eqref{it:localisation-homomorphism}, localised even operators on disjoint sets of modes commute \eqref{it:localisation-commute},
and the insertion of a localised even operator factorises into the insertion of the operator itself and the crossing factors of the remaining modes.
The localisation also respects charges.
Let $u(a)$, $a\in\R^{\Lambda_{\mathrm{s}}}$, be the diagonal unitary of Definition~\ref{def:modes} on $\mathcal{F}(\mathcal{J}^{\mathrm{l}}(\Lambda_{\mathrm{s}}))$, and regard $z\in\Z^\Lambda$ as an element of $\Z^{\Lambda_{\mathrm{s}}}$ through the injection of $\Lambda$ into $\Lambda_{\mathrm{s}}$, extended by zero.
Then
\begin{equation}\label{eq:amplification-iota-charge}
u(a)\,\iota(\Upsilon)\,u(a)^{-1}=\exp\left(iq\langle a,z\rangle\right)\cdot\iota(\Upsilon),\qquad\Upsilon\in\End_z(\Lambda),\ a\in\R^{\Lambda_{\mathrm{s}}}.
\end{equation}

To see this, let $\mathcal{J}':=r(\mathcal{J}^{\mathrm{l}}(\Lambda))$, let $\mathcal{J}''$ be its complement in $\mathcal{J}^{\mathrm{l}}(\Lambda_{\mathrm{s}})$, put $\Upsilon':=U_\Lambda\Upsilon U_\Lambda^*$, so that $\iota(\Upsilon)=\iota(\Upsilon')$,
and for $a\in\R^{\Lambda_{\mathrm{s}}}$ let $a_\Lambda\in\R^\Lambda$ be its pull-back under the injection of $\Lambda$ into $\Lambda_{\mathrm{s}}$.
A mode and its representative belong to the same group and lie over the same site of $\Lambda_{\mathrm{s}}$;
hence they have the same crossing coefficient, and $\mathfrak{q}(r(I))$ is the image of $\mathfrak{q}(I)$ in $\Z^{\Lambda_{\mathrm{s}}}$.
As $U_\Lambda$ maps $e_I$ to $\pm e_{r(I)}$, it follows that $U_\Lambda\,\omega\,U_\Lambda^*=\omega$ and that $U_\Lambda\,u(a_\Lambda)\,U_\Lambda^*$ is the diagonal unitary $u(a)$ on the Fock space of $\mathcal{J}'$
(Definition~\ref{def:modes}; Proposition~\ref{prop:localisation}\eqref{it:mode-relabelling}).
Therefore
\[
u(a)\,\Upsilon'\,u(a)^{-1}=U_\Lambda\,u(a_\Lambda)\,\Upsilon\,u(a_\Lambda)^{-1}\,U_\Lambda^*=\exp\left(iq\langle a,z\rangle\right)\cdot\Upsilon',
\]
that is, $\exp(iq\langle a,\mathfrak{q}(I)-\mathfrak{q}(J)\rangle)\cdot\Upsilon'_{IJ}=\exp(iq\langle a,z\rangle)\cdot\Upsilon'_{IJ}$ for $I,J\subset\mathcal{J}'$.
By Definition~\ref{def:modes}, for $I,J\subset\mathcal{J}'$ and $R,R'\subset\mathcal{J}''$,
\[
\iota(\Upsilon')_{I\cup R,\,J\cup R'}=\delta_{RR'}\cdot(-1)^{\#\{(i,r)\in I\times R:\ i>r\}}\cdot(-1)^{\#\{(j,r)\in J\times R:\ j>r\}}\cdot\Upsilon'_{IJ}.
\]
Since the charge is additive over disjoint sets of modes, $u(a)$ multiplies $e_{I\cup R}$ by $\exp(iq\langle a,\mathfrak{q}(I)+\mathfrak{q}(R)\rangle)$.
Hence the entry of $u(a)\,\iota(\Upsilon')\,u(a)^{-1}$ at $(I\cup R,J\cup R)$ equals $\exp(iq\langle a,\mathfrak{q}(I)-\mathfrak{q}(J)\rangle)\cdot\iota(\Upsilon')_{I\cup R,\,J\cup R}$,
which by the charge of $\Upsilon'$ is $\exp(iq\langle a,z\rangle)\cdot\iota(\Upsilon')_{I\cup R,\,J\cup R}$;
the entries at $(I\cup R,J\cup R')$ with $R\neq R'$ vanish on both sides.

\subsection{Signs at a link plane}\label{sec:link-plane-signs}

Throughout this subsection $\Pi=(x_1,c)$ is a coordinate link plane,
$c=j_0+\frac12$ with $j_0\in\Z$,
and $\theta_c:=\theta_\Pi$, $\Theta_c:=\Theta_\Pi$;
the upward unit normal of $\Pi$ is $n=e_1$.
A bond that crosses $\Pi$ is written $\{x,\theta_c x\}$, with $x$ its lower end.
For $k=1,2$ let $e^{\pm}_k:=e^{\pm}_{e_1,k}$,
so that $(e^{+}_1,e^{+}_2)$ and $(e^{-}_1,e^{-}_2)$ are the orthonormal bases of the ranges of $P^{+}_n$ and $P^{-}_n$ fixed in Definition~\ref{def:spin-matrices}.

Let $\Lambda$ be a finite set of sites of $\Z^2$.
The modes $\zeta_i$, $i\in\mathcal{J}^{\mathrm{l}}(\Lambda)$, at the plane $c$,
their coefficients $\omega_i$ and charges $\mathfrak{q}(i)\in\Z^\Lambda$,
the Fock space $\mathcal{F}_\Lambda$ with its basis $(e_I)$,
the operators $\omega$, $U_{\mathrm{f}}$ and $u(a)$, $a\in\R^\Lambda$,
and the elements
\[
\Ins_c[\Upsilon]=\sum_{I,J}\Upsilon_{IJ}\cdot\Theta_c(\zeta_J)\,\zeta_I
\qquad\left(\Upsilon\in\End\mathcal{F}_\Lambda\right)
\]
are those of Definition~\ref{def:modes} and~\eqref{eq:setting-ordered-monomials};
the row index $I$ belongs to the upper side, the column index $J$ to the lower side.
A mode over $y\in\Lambda$ belongs to a bond $\{x,\theta_c x\}$ with $\pi(x)=(j_0,y)$ and $\pi(\theta_c x)=(j_0+1,y)$.
For a monomial $\Psi$ in the fields $\psi_\alpha(x)$, $\bar\psi_\alpha(x)$ let $\mathfrak{q}(\Psi)$ be its charge density (Definition~\ref{def:chamber-operator}),
and let
\[
\mathfrak{q}_-(\Psi)(y):=\mathfrak{q}(\Psi)(j_0,y),
\qquad
\mathfrak{q}_+(\Psi)(y):=\mathfrak{q}(\Psi)(j_0+1,y)
\qquad(y\in\Lambda)
\]
be its charge densities on the lower plane $\{x_1=j_0\}$ and on the upper plane $\{x_1=j_0+1\}$.

Let $C^1$ be the space of the real one-cochains on $\Z^3$ (Definition~\ref{def:cochains}).
For a map $\Upsilon\colon C^1\to\End\mathcal{F}_\Lambda$, $\hat A\mapsto\Upsilon(\hat A)$, the element $\Ins_c[\Upsilon]$ is formed pointwise,
\[
\Ins_c[\Upsilon](\hat A):=\Ins_c\left[\Upsilon(\hat A)\right]\qquad\left(\hat A\in C^1\right);
\]
it belongs to the algebra $\Gr(\mathbb{X})^{C^1}$ of the maps from $C^1$ to $\Gr(\mathbb{X})$ (Definition~\ref{def:integrands}).
Hence, for every mirror $\Pi'$, the reflection $\Theta_{\Pi'}$ acts on it by~\eqref{eq:setting-reflection-coefficients},
and, if its values lie in $\Gr^{\mathrm{ev}}(\mathbb{X})$, every $g\in W$ acts on it by~\eqref{eq:setting-group-action}.
A map $\Upsilon$ defined on $\R^{B^{(1)}}$ for a box $B$ is regarded as the map $\hat A\mapsto\Upsilon\left(\hat A|_{B^{(1)}}\right)$ on $C^1$,
in the same way as an integrand that reads in $B$.
For $g\in W$, for instance $g=\theta_c$ or $g=\tau_v$ with $v\in\Z^3$, the map $\hat A\mapsto\Upsilon(g^{*}\hat A)$ then depends only on the restriction of $\hat A$ to $(gB)^{(1)}$, the set of the links contained in the box $gB$:
the isometry $g$ maps $B^{(1)}$ onto $(gB)^{(1)}$, and $(g^{*}\hat A)(b)=\pm\hat A(gb)$ for every link $b$.

\begin{lemma}[Reflection of contracted fields]\label{lem:contracted-fields}
Let $\Pi'$ be a mirror, $\theta_{\Pi'}$ its reflection and $n'$ its upward unit normal.
For $u\in\C^4$ and $x\in\mathbb{X}$ put
$u^*\psi(x):=\sum_\alpha\overline{u_\alpha}\,\psi_\alpha(x)$ and
$\bar\psi(x)u:=\sum_\alpha\bar\psi_\alpha(x)\,u_\alpha$.
Then
\[
\Theta_{\Pi'}\left(u^*\psi(x)\right)=-\bar\psi(\theta_{\Pi'}x)\,\gamma_{n'}u,
\qquad
\Theta_{\Pi'}\left(\bar\psi(x)u\right)=-u^*\gamma_{n'}\,\psi(\theta_{\Pi'}x)=-(\gamma_{n'}u)^*\,\psi(\theta_{\Pi'}x).
\]
\end{lemma}

\begin{proof}
Written in components, \eqref{eq:setting-reflection-fields} reads
\[
\Theta_{\Pi'}\psi_\alpha(x)=-\sum_\beta\bar\psi_\beta(\theta_{\Pi'}x)\cdot(\gamma_{n'})_{\beta\alpha},
\qquad
\Theta_{\Pi'}\bar\psi_\alpha(x)=-\sum_\beta(\gamma_{n'})_{\alpha\beta}\cdot\psi_\beta(\theta_{\Pi'}x).
\]
As $\Theta_{\Pi'}$ is antilinear, it maps $u^*\psi(x)$ to
\[
\sum_\alpha u_\alpha\,\Theta_{\Pi'}\psi_\alpha(x)=-\sum_{\alpha,\beta}\bar\psi_\beta(\theta_{\Pi'}x)\cdot(\gamma_{n'})_{\beta\alpha}\,u_\alpha=-\bar\psi(\theta_{\Pi'}x)\,\gamma_{n'}u,
\]
and it maps $\bar\psi(x)u$ to
\[
\sum_\alpha\overline{u_\alpha}\,\Theta_{\Pi'}\bar\psi_\alpha(x)
=-\sum_{\alpha,\beta}\overline{u_\alpha}\,(\gamma_{n'})_{\alpha\beta}\cdot\psi_\beta(\theta_{\Pi'}x)
=-u^*\gamma_{n'}\,\psi(\theta_{\Pi'}x).
\]
The matrix $\gamma_{n'}$ is Hermitian (Definition~\ref{def:spin-matrices}),
so that $u^*\gamma_{n'}=(\gamma_{n'}u)^*$.
\end{proof}

\begin{lemma}[Reflected modes and their charges]\label{lem:reflected-modes}
Let $\{x,\theta_c x\}$ be a bond that crosses $\Pi$, and let $k\in\{1,2\}$.
The modes of the four groups at this bond and their images under $\Theta_c$ are
\[
\begin{aligned}
\text{group }1:\quad\zeta&=(e^-_k)^*\psi(\theta_c x),&\qquad\Theta_c(\zeta)&=+\bar\psi(x)\,e^-_k,\\
\text{group }2:\quad\zeta&=(e^+_k)^*\psi(\theta_c x),&\qquad\Theta_c(\zeta)&=-\bar\psi(x)\,e^+_k,\\
\text{group }3:\quad\zeta&=\bar\psi(\theta_c x)\,e^+_k,&\qquad\Theta_c(\zeta)&=-(e^+_k)^*\psi(x),\\
\text{group }4:\quad\zeta&=\bar\psi(\theta_c x)\,e^-_k,&\qquad\Theta_c(\zeta)&=+(e^-_k)^*\psi(x).
\end{aligned}
\]
Let $I,J\subset\mathcal{J}^{\mathrm{l}}(\Lambda)$.
Then $\zeta_I$ is a linear combination of monomials whose charge density equals $\mathfrak{q}(I)$ on the plane $\{x_1=j_0+1\}$ and vanishes at all other sites,
and $\Theta_c(\zeta_J)$ is a linear combination of monomials whose charge density equals $-\mathfrak{q}(J)$ on the plane $\{x_1=j_0\}$ and vanishes at all other sites.
In particular every monomial $\Psi$ in the expansion of $\Theta_c(\zeta_J)\,\zeta_I$ satisfies
$\mathfrak{q}_+(\Psi)=\mathfrak{q}(I)$ and $\mathfrak{q}_-(\Psi)=-\mathfrak{q}(J)$.
\end{lemma}

\begin{proof}
Since $P^{\pm}_n=\frac12(1\pm\gamma_n)$, one has $\gamma_n e^{\pm}_k=\pm e^{\pm}_k$;
moreover $\theta_c^2=\mathrm{id}$ by Definition~\ref{def:mirrors}\eqref{it:mirror-involution}.
Lemma~\ref{lem:contracted-fields}, applied with $\Pi'=\Pi$ at the upper end $\theta_c x$, gives for every $u\in\C^4$
\[
\Theta_c\left(u^*\psi(\theta_c x)\right)=-\bar\psi(x)\,\gamma_n u,
\qquad
\Theta_c\left(\bar\psi(\theta_c x)\,u\right)=-(\gamma_n u)^*\psi(x).
\]
The choices $u=e^-_k$ and $u=e^+_k$ in the first identity give the lines of the groups $1$ and $2$;
the choices $u=e^+_k$ and $u=e^-_k$ in the second identity give those of the groups $3$ and $4$.

The charge density of a monomial counts, at each photon site, the factors $\psi_\alpha(x')$ whose nearest photon site it is,
minus the factors $\bar\psi_\alpha(x')$ whose nearest photon site it is;
it is additive under products.
A mode $\zeta_i$ over $y$ is a combination of the components of $\psi(\theta_c x)$ (groups $1$, $2$) or of $\bar\psi(\theta_c x)$ (groups $3$, $4$),
and $\pi(\theta_c x)=(j_0+1,y)$;
each of its monomials therefore has the charge density $\mathfrak{q}(i)=\pm\delta_y$ on the plane $\{x_1=j_0+1\}$ and no charge elsewhere.
By the table, $\Theta_c(\zeta_i)$ is a combination of the components of $\bar\psi(x)$ (groups $1$, $2$) or of $\psi(x)$ (groups $3$, $4$),
with $\pi(x)=(j_0,y)$,
and these carry the opposite charge on the plane $\{x_1=j_0\}$.
Additivity over the factors of $\zeta_I$ and of $\Theta_c(\zeta_J)$ gives the two assertions on these elements, and the last assertion follows from them.
\end{proof}

\begin{lemma}[Phases at the lower ends]\label{lem:lower-phases}
Let $a\in\R^\Lambda$, and let $\alpha^-_a$ be the automorphism of the Grassmann algebra $\Gr(\mathbb{X})$ (Definition~\ref{def:grassmann})
that multiplies the components of $\bar\psi(x)$ by $\exp\left(iq\,a(y)\right)$ and those of $\psi(x)$ by $\exp\left(-iq\,a(y)\right)$
whenever $x$ is the lower end of a bond that crosses $\Pi$ and $\pi(x)=(j_0,y)$ with $y\in\Lambda$,
and that fixes all other generators.
Then for every $J\subset\mathcal{J}^{\mathrm{l}}(\Lambda)$, in particular for every set $J$ consisting of a single mode,
\[
\alpha^-_a\Theta_c(\zeta_J)=\exp\left(iq\langle a,\mathfrak{q}(J)\rangle\right)\cdot\Theta_c(\zeta_J),
\]
and for every $\Upsilon\in\End\mathcal{F}_\Lambda$, as well as, pointwise in $\hat A$, for every map $\Upsilon\colon C^1\to\End\mathcal{F}_\Lambda$,
\[
\alpha^-_a\left(\Ins_c[\Upsilon]\right)=\Ins_c\left[\Upsilon\,u(a)\right].
\]
Moreover
$\alpha^-_a(\Psi)=\exp\left(-iq\langle\mathfrak{q}_-(\Psi),a\rangle\right)\cdot\Psi$
for every monomial $\Psi$ whose factors at the sites $x'$ with $\pi(x')\in\{j_0\}\times\Lambda$ are all fields at lower ends of bonds that cross $\Pi$.
\end{lemma}

\begin{proof}
Let $i$ be a mode over $y\in\Lambda$, at the bond $\{x,\theta_c x\}$, so that $\pi(x)=(j_0,y)$.
By Lemma~\ref{lem:reflected-modes}, $\Theta_c(\zeta_i)$ is a combination of the components of $\bar\psi(x)$ if $i$ belongs to the group $1$ or $2$,
and of the components of $\psi(x)$ if $i$ belongs to the group $3$ or $4$.
In the first case $\alpha^-_a$ multiplies it by $\exp\left(iq\,a(y)\right)$ and $\mathfrak{q}(i)=\delta_y$;
in the second case it multiplies it by $\exp\left(-iq\,a(y)\right)$ and $\mathfrak{q}(i)=-\delta_y$.
In both cases the factor is $\exp\left(iq\langle a,\mathfrak{q}(i)\rangle\right)$.
The formula for an arbitrary $J$ follows, since $\alpha^-_a$ is multiplicative, $\Theta_c(\zeta_J)$ is the product of the $\Theta_c(\zeta_i)$, $i\in J$, in decreasing order of $i$, and $\mathfrak{q}(J)=\sum_{i\in J}\mathfrak{q}(i)$.
The modes $\zeta_i$ are formed from fields at upper ends, which $\alpha^-_a$ fixes; hence
\[
\alpha^-_a\left(\Ins_c[\Upsilon]\right)
=\sum_{I,J}\Upsilon_{IJ}\cdot\exp\left(iq\langle a,\mathfrak{q}(J)\rangle\right)\cdot\Theta_c(\zeta_J)\,\zeta_I
=\Ins_c\left[\Upsilon\,u(a)\right],
\]
because the diagonal matrix $u(a)$, multiplied from the right, multiplies the column $J$, which belongs to the lower side, by $\exp\left(iq\langle a,\mathfrak{q}(J)\rangle\right)$.
Finally let $\Psi$ be a monomial as in the last assertion.
Each of its factors that is a component of $\bar\psi(x)$, with $x$ a lower end over $y\in\Lambda$, is multiplied by $\exp\left(iq\,a(y)\right)$,
each such component of $\psi(x)$ by $\exp\left(-iq\,a(y)\right)$,
and every other factor is fixed.
By the hypothesis on $\Psi$ and the definition of the charge density, the product of these numbers is
$\exp\left(-iq\sum_{y\in\Lambda}a(y)\,\mathfrak{q}_-(\Psi)(y)\right)$.
\end{proof}

\begin{lemma}[Reflection in the link plane]\label{lem:plane-reflection}
For $I,J\subset\mathcal{J}^{\mathrm{l}}(\Lambda)$ one has
$\Theta_c\left(\Theta_c(\zeta_J)\,\zeta_I\right)=\Theta_c(\zeta_I)\,\zeta_J$.
Let $\Upsilon\colon C^1\to\End\mathcal{F}_\Lambda$ be a map, and let $\Theta_c$ act on the element $\Ins_c[\Upsilon]$ of $\Gr(\mathbb{X})^{C^1}$ by~\eqref{eq:setting-reflection-coefficients}.
Then
\[
\Theta_c\left(\Ins_c[\Upsilon]\right)(\hat A)=\Ins_c\left[\Upsilon(\theta_c^*\hat A)^*\right].
\]
\end{lemma}

\begin{proof}
Since $\Theta_c$ reverses products and $\Theta_c^2=1$ (Lemma~\ref{lem:sign-rules}(a)),
\[
\Theta_c\left(\Theta_c(\zeta_J)\,\zeta_I\right)=\Theta_c(\zeta_I)\,\Theta_c^2(\zeta_J)=\Theta_c(\zeta_I)\,\zeta_J .
\]
The elements $\Theta_c(\zeta_J)\zeta_I$ do not depend on $\hat A$;
the rule~\eqref{eq:setting-reflection-coefficients} and the antilinearity of $\Theta_c$ give
\[
\Theta_c\left(\Ins_c[\Upsilon]\right)(\hat A)=\sum_{I,J}\overline{\Upsilon_{IJ}(\theta_c^*\hat A)}\cdot\Theta_c(\zeta_I)\,\zeta_J .
\]
In the term $\Theta_c(\zeta_I)\zeta_J$ the row index is $J$ and the column index is $I$,
and its coefficient $\overline{\Upsilon_{IJ}(\theta_c^*\hat A)}$ is the entry $(J,I)$ of the adjoint matrix $\Upsilon(\theta_c^*\hat A)^*$.
\end{proof}

\begin{lemma}[Translations]\label{lem:translated-insertions}
Let $v=(v_1,v_2,v_3)\in\Z^3$ and $v_{\mathrm{s}}:=(v_2,v_3)$.
Let $\alpha_v$ be the automorphism of $\Gr(\mathbb{X})$ with $\alpha_v\psi(x)=\psi(x+v)$ and $\alpha_v\bar\psi(x)=\bar\psi(x+v)$,
acting on the maps $G\colon C^1\to\Gr(\mathbb{X})$, written $G(\hat A)=\sum_\Psi a_\Psi(\hat A)\,\Psi$ in a monomial basis, by
$(\alpha_v G)(\hat A):=\sum_\Psi a_\Psi(\tau_v^*\hat A)\cdot\alpha_v(\Psi)$.
\begin{enumerate}
\item[(a)] $G^{\tau_v}=\alpha_v G$ for every such map $G$ with values in $\Gr^{\mathrm{ev}}(\mathbb{X})$, in particular for every even integrand $G$.
\item[(b)] $\alpha_v\Theta_c=\Theta_{c+v_1}\alpha_v$.
\item[(c)] For a mode $i$ of $\Pi$ let $i+v$ be the mode of the translated bond with the same group and the same index $k$.
The map $i\mapsto i+v$ is a bijection of $\idx\left((x_1,c)\right)$ onto $\idx\left((x_1,c+v_1)\right)$,
which maps the modes over $\Lambda$ at the plane $c$ onto the modes over $\Lambda+v_{\mathrm{s}}$ at the plane $c+v_1$, and
\[
\alpha_v\zeta_i=\zeta_{i+v},
\qquad
\omega_{i+v}=\omega_i,
\qquad
\mathfrak{q}(i+v)=\mathfrak{q}(i)(\,\cdot\,-v_{\mathrm{s}}).
\]
For $I=\{i_1<\dots<i_r\}$ let $\operatorname{sgn}_v(I)$ be the sign of the permutation that sorts $(i_1+v,\dots,i_r+v)$.
Then $\alpha_v\zeta_I=\operatorname{sgn}_v(I)\cdot\zeta_{I+v}$, and $\operatorname{sgn}_v(I)=1$ if $v$ is parallel to $e_1$.
\item[(d)] Let $U_v\colon\mathcal{F}_\Lambda\to\mathcal{F}_{\Lambda+v_{\mathrm{s}}}$ be the linear map with $U_v e_I:=\operatorname{sgn}_v(I)\cdot e_{I+v}$.
It is unitary, and for every map $\Upsilon\colon C^1\to\End\mathcal{F}_\Lambda$
\[
\alpha_v\left(\Ins_c[\Upsilon]\right)(\hat A)=\Ins_{c+v_1}\left[U_v\,\Upsilon(\tau_v^*\hat A)\,U_v^*\right];
\]
the left side equals $\left(\Ins_c[\Upsilon]\right)^{\tau_v}(\hat A)$ if the values of $\Upsilon$ are even.
Moreover $U_v\,\omega=\omega\,U_v$, and
$U_v\,u\left(a(\,\cdot\,+v_{\mathrm{s}})\right)\,U_v^*=u(a)$ for $a\in\R^{\Lambda+v_{\mathrm{s}}}$.
\end{enumerate}
\end{lemma}

\begin{proof}
(a) Let $G\colon C^1\to\Gr^{\mathrm{ev}}(\mathbb{X})$.
For $\mu=1,2,3$ Definition~\ref{def:mirrors}\eqref{it:mirror-product} gives
$\theta_{x_\mu,v_\mu/2}\,\theta_{x_\mu,0}=\tau_{v_\mu e_\mu}$,
so that $\tau_v$ is the product of the six reflections
$\theta_{x_1,v_1/2}\,\theta_{x_1,0}\,\theta_{x_2,v_2/2}\,\theta_{x_2,0}\,\theta_{x_3,v_3/2}\,\theta_{x_3,0}$,
and by~\eqref{eq:setting-group-action}
\[
G^{\tau_v}=\Theta_{(x_1,v_1/2)}\Theta_{(x_1,0)}\,\Theta_{(x_2,v_2/2)}\Theta_{(x_2,0)}\,\Theta_{(x_3,v_3/2)}\Theta_{(x_3,0)}\,G .
\]
The mirrors $(x_\mu,v_\mu/2)$ and $(x_\mu,0)$ have the same upward normal $e_\mu$, and $\gamma_{e_\mu}^2=1$.
By~\eqref{eq:setting-two-reflections} the linear automorphism $\Theta_{(x_\mu,v_\mu/2)}\Theta_{(x_\mu,0)}$ therefore maps $\psi(x)$ to $\psi(x+v_\mu e_\mu)$ and $\bar\psi(x)$ to $\bar\psi(x+v_\mu e_\mu)$.
Applied twice, the rule~\eqref{eq:setting-reflection-coefficients} shows that it maps a coefficient $a_\Psi(\hat A)$ to
$a_\Psi\left(\theta_{x_\mu,0}^*\,\theta_{x_\mu,v_\mu/2}^*\hat A\right)=a_\Psi\left(\tau_{v_\mu e_\mu}^*\hat A\right)$:
the two complex conjugations cancel, and $(gg')^*=g'^*g^*$ (Definition~\ref{def:cochains}).
Thus $\Theta_{(x_\mu,v_\mu/2)}\Theta_{(x_\mu,0)}=\alpha_{v_\mu e_\mu}$, and the composition of the three maps is $\alpha_v$.

(b) By Definition~\ref{def:mirrors}\eqref{it:mirror-conjugation}, $\tau_v\theta_c\tau_{-v}=\theta_{c+v_1}$;
equivalently $\tau_v\theta_c=\theta_{c+v_1}\tau_v$, that is, $\theta_{c+v_1}(x+v)=\theta_c(x)+v$.
On the generators this gives
\[
\begin{aligned}
\alpha_v\Theta_c\psi(x)&=-\bar\psi(\theta_c x+v)\,\gamma_n=-\bar\psi\left(\theta_{c+v_1}(x+v)\right)\gamma_n=\Theta_{c+v_1}\alpha_v\psi(x),\\
\alpha_v\Theta_c\bar\psi(x)&=-\gamma_n\,\psi(\theta_c x+v)=-\gamma_n\,\psi\left(\theta_{c+v_1}(x+v)\right)=\Theta_{c+v_1}\alpha_v\bar\psi(x).
\end{aligned}
\]
Both $\alpha_v\Theta_c$ and $\Theta_{c+v_1}\alpha_v$ are antilinear and reverse products, so they agree on $\Gr(\mathbb{X})$.
On a coefficient they produce $\overline{a_\Psi(\theta_c^*\tau_v^*\hat A)}$ and $\overline{a_\Psi(\tau_v^*\theta_{c+v_1}^*\hat A)}$,
which coincide because $\theta_c^*\tau_v^*=(\tau_v\theta_c)^*=(\theta_{c+v_1}\tau_v)^*=\tau_v^*\theta_{c+v_1}^*$.

(c) The translation $\tau_v$ belongs to $W$, which maps bonds onto bonds of the same type (Definition~\ref{def:fermion-sites}),
and it adds $v_1$ to the first coordinate.
The translate $\{x+v,\theta_c x+v\}=\{x+v,\theta_{c+v_1}(x+v)\}$ of a bond that crosses $\Pi$ is therefore a bond that crosses $(x_1,c+v_1)$, with the lower end $x+v$.
At both planes the vectors $e^{\pm}_k$ are those of the normal $e_1$;
hence $i\mapsto i+v$ is a bijection of $\idx\left((x_1,c)\right)$ onto $\idx\left((x_1,c+v_1)\right)$, with the inverse $i\mapsto i-v$, and $\alpha_v\zeta_i=\zeta_{i+v}$,
as for instance $\alpha_v\left((e^-_k)^*\psi(\theta_c x)\right)=(e^-_k)^*\psi\left(\theta_{c+v_1}(x+v)\right)$.
The coefficient $\omega_i$ depends only on the type of the bond and on the group, so $\omega_{i+v}=\omega_i$.
Since $\pi(x+v)=\pi(x)+v$ (Proposition~\ref{prop:geometry}\eqref{it:geometry-projection}), a bond over $y$ is carried to a bond over $y+v_{\mathrm{s}}$;
this gives the statement on the modes over $\Lambda$ and the rule for the charges.
The modes anticommute, so that $\alpha_v\zeta_I=\zeta_{i_1+v}\cdots\zeta_{i_r+v}=\operatorname{sgn}_v(I)\cdot\zeta_{I+v}$.
If $v=v_1e_1$, the translation preserves the orders of the index sets,
because these orders are compatible with the normal translations (Definition~\ref{def:modes});
then $\operatorname{sgn}_v(I)=1$.

(d) The map $U_v$ carries the orthonormal basis $(e_I)$ onto the orthonormal basis $(\pm e_{I+v})$; so it is unitary, and $U_v^*e_{I+v}=\operatorname{sgn}_v(I)\cdot e_I$.
By (b) and (c),
\[
\alpha_v\left(\Theta_c(\zeta_J)\,\zeta_I\right)=\Theta_{c+v_1}(\alpha_v\zeta_J)\,\alpha_v\zeta_I
=\operatorname{sgn}_v(I)\cdot\operatorname{sgn}_v(J)\cdot\Theta_{c+v_1}(\zeta_{J+v})\,\zeta_{I+v},
\]
while the coefficient $\Upsilon_{IJ}(\hat A)$ is replaced by $\Upsilon_{IJ}(\tau_v^*\hat A)$.
The entry $(I+v,J+v)$ of $U_v\,\Upsilon(\tau_v^*\hat A)\,U_v^*$ is $\operatorname{sgn}_v(I)\cdot\operatorname{sgn}_v(J)\cdot\Upsilon_{IJ}(\tau_v^*\hat A)$,
which proves the identity.
If the values of $\Upsilon$ are even, that is, commute with $U_{\mathrm{f}}$, then $\Upsilon_{IJ}=0$ unless $\#|I|+\#|J|$ is even;
as $\Theta_c(\zeta_J)\zeta_I$ has the degree $\#|I|+\#|J|$, the map $\Ins_c[\Upsilon]$ then takes its values in $\Gr^{\mathrm{ev}}(\mathbb{X})$,
and (a), applied to $G=\Ins_c[\Upsilon]$, identifies $\alpha_v\left(\Ins_c[\Upsilon]\right)$ with $\left(\Ins_c[\Upsilon]\right)^{\tau_v}$.
From $\omega_{I+v}=\omega_I$ one gets $\omega\,U_v e_I=\operatorname{sgn}_v(I)\cdot\omega_{I+v}\,e_{I+v}=U_v\,\omega\,e_I$.
Finally, for $a\in\R^{\Lambda+v_{\mathrm{s}}}$,
\[
\left\langle a(\,\cdot\,+v_{\mathrm{s}}),\mathfrak{q}(I)\right\rangle
=\left\langle a,\mathfrak{q}(I)(\,\cdot\,-v_{\mathrm{s}})\right\rangle
=\left\langle a,\mathfrak{q}(I+v)\right\rangle,
\]
so that $U_v\,u\left(a(\,\cdot\,+v_{\mathrm{s}})\right)\,U_v^*e_{I+v}=\exp\left(iq\langle a,\mathfrak{q}(I+v)\rangle\right)\cdot e_{I+v}=u(a)\,e_{I+v}$.
\end{proof}

\begin{lemma}[Spatial mirrors]\label{lem:spatial-mirror-signs}
Let $\nu\in\{2,3\}$ and $w\in\frac12\Z$,
let $\Pi_w:=(x_\nu,w)$, let $\theta_w$ be its reflection, $\Theta_w:=\Theta_{\Pi_w}$ and $\gamma_w:=\gamma_\nu$,
and let $\theta_{\mathrm{s}}$ be the reflection of $\Z^2$ with $\theta_w(x_1,y)=(x_1,\theta_{\mathrm{s}}y)$.
Then $\theta_w$ commutes with $\theta_c$, preserves the coordinate $x_1$,
and maps the bonds that cross $\Pi$ onto bonds that cross $\Pi$, lower ends onto lower ends.
For such a bond $\mathfrak{b}=\{x,\theta_c x\}$ put $\mathfrak{b}':=\theta_w\mathfrak{b}$,
and let $\zeta_{(\mathfrak{b},1,k)},\dots,\zeta_{(\mathfrak{b},4,k)}$ be its modes of the groups $1,\dots,4$ with the index $k$.
\begin{enumerate}
\item[(a)] One has
\[
\begin{aligned}
\Theta_w\left(\zeta_{(\mathfrak{b},1,k)}\right)&=-\bar\psi(\theta_w\theta_c x)\,\gamma_w e^-_k=\sum_{k'}\left[-(e^+_{k'})^*\gamma_w e^-_k\right]\cdot\zeta_{(\mathfrak{b}',3,k')},\\
\Theta_w\left(\zeta_{(\mathfrak{b},2,k)}\right)&=\sum_{k'}\left[-(e^-_{k'})^*\gamma_w e^+_k\right]\cdot\zeta_{(\mathfrak{b}',4,k')},\\
\Theta_w\left(\zeta_{(\mathfrak{b},3,k)}\right)&=-(\gamma_w e^+_k)^*\,\psi(\theta_w\theta_c x)=\sum_{k'}\left[-\overline{(e^-_{k'})^*\gamma_w e^+_k}\right]\cdot\zeta_{(\mathfrak{b}',1,k')},\\
\Theta_w\left(\zeta_{(\mathfrak{b},4,k)}\right)&=\sum_{k'}\left[-\overline{(e^+_{k'})^*\gamma_w e^-_k}\right]\cdot\zeta_{(\mathfrak{b}',2,k')}.
\end{aligned}
\]
The matrix $\mathfrak{u}=(\mathfrak{u}_{i'i})$, $i\in\mathcal{J}^{\mathrm{l}}(\Lambda)$, $i'\in\mathcal{J}^{\mathrm{l}}(\theta_{\mathrm{s}}\Lambda)$,
defined by $\Theta_w(\zeta_i)=\sum_{i'}\mathfrak{u}_{i'i}\,\zeta_{i'}$,
is unitary, with unitary $2\times2$ blocks,
and $\mathfrak{u}_{i'i}\neq0$ only if $\omega_{i'}=\omega_i$ and $\mathfrak{q}(i')=-\mathfrak{q}(i)\circ\theta_{\mathrm{s}}$.
\item[(b)] For finite $I\subset\mathcal{J}^{\mathrm{l}}(\Lambda)$ and $I'\subset\mathcal{J}^{\mathrm{l}}(\theta_{\mathrm{s}}\Lambda)$ let
$(\wedge\mathfrak{u})_{I'I}:=\det\left(\mathfrak{u}_{i'i}\right)_{i'\in I',\,i\in I}$, rows and columns in increasing order, if $\#|I'|=\#|I|$,
and $(\wedge\mathfrak{u})_{I'I}:=0$ otherwise;
for $I=\{i_1<\dots<i_r\}$ let $\zeta^w_I:=\Theta_w(\zeta_{i_1})\cdots\Theta_w(\zeta_{i_r})$.
Then
\[
\zeta^w_I=\sum_{I'}(\wedge\mathfrak{u})_{I'I}\cdot\zeta_{I'}.
\]
The operator $\wedge\mathfrak{u}\colon\mathcal{F}_\Lambda\to\mathcal{F}_{\theta_{\mathrm{s}}\Lambda}$ is unitary,
$(\wedge\mathfrak{u})_{I'I}\neq0$ only if $\omega_{I'}=\omega_I$ and $\mathfrak{q}(I')=-\mathfrak{q}(I)\circ\theta_{\mathrm{s}}$,
and $\wedge\mathfrak{u}$ commutes with $\omega$.
\item[(c)] For finite $I,J\subset\mathcal{J}^{\mathrm{l}}(\Lambda)$:
\begin{enumerate}
\item[(i)] $\Theta_w(\zeta_I)=(-1)^{\binom{\#|I|}{2}}\,\zeta^w_I$;
\item[(ii)] $\Theta_w\Theta_c G=(-1)^{\deg G}\,\Theta_c\Theta_w G$ for every homogeneous element $G$ of $\Gr(\mathbb{X})$;
\item[(iii)] if $\#|I|\equiv\#|J|$ modulo $2$, then
$\Theta_w\left(\Theta_c(\zeta_J)\,\zeta_I\right)=(-1)^{\binom{\#|I|}{2}+\binom{\#|J|}{2}}\,\Theta_c(\zeta^w_J)\,\zeta^w_I$;
\item[(iv)] $\Theta_c(\zeta^w_J)=\sum_{J'}\overline{(\wedge\mathfrak{u})_{J'J}}\cdot\Theta_c(\zeta_{J'})$.
\end{enumerate}
\item[(d)] Let $\mathfrak{r}$ be the diagonal operator on $\mathcal{F}_\Lambda$ with $\mathfrak{r}e_I:=(-1)^{\binom{\#|I|}{2}}e_I$, and $\mathfrak{U}:=(\wedge\mathfrak{u})\,\mathfrak{r}$.
Then $\mathfrak{U}\colon\mathcal{F}_\Lambda\to\mathcal{F}_{\theta_{\mathrm{s}}\Lambda}$ is unitary and commutes with $\omega$,
and for every map $\Upsilon\colon C^1\to\End\mathcal{F}_\Lambda$ with even values
\[
\Theta_w\left(\Ins_c[\Upsilon]\right)(\hat A)=\Ins_c\left[\mathfrak{U}\,\overline{\Upsilon(\theta_w^*\hat A)}\,\mathfrak{U}^*\right],
\qquad
\mathfrak{U}\,\overline{u(a\circ\theta_{\mathrm{s}})}\,\mathfrak{U}^*=u(a)\quad\left(a\in\R^{\theta_{\mathrm{s}}\Lambda}\right),
\]
where the bar denotes the entrywise complex conjugate in the bases $(e_I)$,
and $\Ins_c$ on the right is formed with the modes over $\theta_{\mathrm{s}}\Lambda$.
\end{enumerate}
\end{lemma}

\begin{proof}
The planes $\{x_\nu=w\}$ and $\{x_1=c\}$ are orthogonal;
so $\theta_w$ and $\theta_c$ commute, and $\theta_w$ changes only the coordinate $x_\nu$.
As an element of $W$, $\theta_w$ maps bonds onto bonds of the same type (Definition~\ref{def:fermion-sites});
it maps $\{x,\theta_c x\}$ onto $\{\theta_w x,\theta_c\theta_w x\}$, which crosses $\Pi$ and has the lower end $\theta_w x$.
By Proposition~\ref{prop:geometry}\eqref{it:geometry-projection}, $\pi(\theta_w x)=\theta_w\pi(x)$;
a bond over $y$ is therefore mapped onto a bond over $\theta_{\mathrm{s}}y$,
and the modes over $\Lambda$ correspond to the modes over $\theta_{\mathrm{s}}\Lambda$.

(a) By Definition~\ref{def:spin-matrices} the matrix $\gamma_w$ is Hermitian, $\gamma_w^2=1$, and $\gamma_w\gamma_n=-\gamma_n\gamma_w$ because $\nu\neq1$.
If $\gamma_n u=\pm u$, then $\gamma_n\gamma_w u=\mp\gamma_w u$:
the matrix $\gamma_w$ maps $\operatorname{ran}P^{\pm}_n$ isometrically into, hence onto, $\operatorname{ran}P^{\mp}_n$, both ranges having dimension $2$.
Consequently
\[
\gamma_w e^-_k=\sum_{k'}\left((e^+_{k'})^*\gamma_w e^-_k\right)\cdot e^+_{k'},
\qquad
\gamma_w e^+_k=\sum_{k'}\left((e^-_{k'})^*\gamma_w e^+_k\right)\cdot e^-_{k'}.
\]
The upward unit normal of $\Pi_w$ is $e_\nu$, and $\theta_w\theta_c x=\theta_c(\theta_w x)$ is the upper end of $\mathfrak{b}'$.
Lemma~\ref{lem:contracted-fields} and the two expansions give
\[
\begin{aligned}
\Theta_w\left((e^-_k)^*\psi(\theta_c x)\right)&=-\bar\psi(\theta_c\theta_w x)\,\gamma_w e^-_k=\sum_{k'}\left[-(e^+_{k'})^*\gamma_w e^-_k\right]\cdot\bar\psi(\theta_c\theta_w x)\,e^+_{k'},\\
\Theta_w\left((e^+_k)^*\psi(\theta_c x)\right)&=-\bar\psi(\theta_c\theta_w x)\,\gamma_w e^+_k=\sum_{k'}\left[-(e^-_{k'})^*\gamma_w e^+_k\right]\cdot\bar\psi(\theta_c\theta_w x)\,e^-_{k'},\\
\Theta_w\left(\bar\psi(\theta_c x)\,e^+_k\right)&=-(\gamma_w e^+_k)^*\,\psi(\theta_c\theta_w x)=\sum_{k'}\left[-\overline{(e^-_{k'})^*\gamma_w e^+_k}\right]\cdot(e^-_{k'})^*\psi(\theta_c\theta_w x),\\
\Theta_w\left(\bar\psi(\theta_c x)\,e^-_k\right)&=-(\gamma_w e^-_k)^*\,\psi(\theta_c\theta_w x)=\sum_{k'}\left[-\overline{(e^+_{k'})^*\gamma_w e^-_k}\right]\cdot(e^+_{k'})^*\psi(\theta_c\theta_w x),
\end{aligned}
\]
where the last two lines use $(\gamma_w e^{\pm}_k)^*=\sum_{k'}\overline{(e^{\mp}_{k'})^*\gamma_w e^{\pm}_k}\cdot(e^{\mp}_{k'})^*$.
By Definition~\ref{def:modes}, $\bar\psi(\theta_c\theta_w x)\,e^+_{k'}$, $\bar\psi(\theta_c\theta_w x)\,e^-_{k'}$, $(e^-_{k'})^*\psi(\theta_c\theta_w x)$ and $(e^+_{k'})^*\psi(\theta_c\theta_w x)$
are the modes $\zeta_{(\mathfrak{b}',3,k')}$, $\zeta_{(\mathfrak{b}',4,k')}$, $\zeta_{(\mathfrak{b}',1,k')}$ and $\zeta_{(\mathfrak{b}',2,k')}$;
this proves the four lines.

For a fixed bond and a fixed group the coefficients form, up to the factor $-1$ and possibly an entrywise complex conjugation, the matrix of the isometry $\gamma_w$ between two of these ranges in orthonormal bases;
this is a unitary $2\times2$ matrix.
The modes of the groups $1,2,3,4$ of $\mathfrak{b}$ are sent to combinations of those of the groups $3,4,1,2$ of $\mathfrak{b}'$,
and $\mathfrak{b}\mapsto\theta_w\mathfrak{b}$ is a bijection of the bonds over $\Lambda$ onto the bonds over $\theta_{\mathrm{s}}\Lambda$.
Up to the corresponding relabelling of the index sets, $\mathfrak{u}$ is therefore block diagonal with unitary $2\times2$ blocks, and it is unitary.
By Definition~\ref{def:modes} the coefficient of a mode depends only on the type of its bond and on its group,
and the groups $1$, $3$ have the same coefficient, as have the groups $2$, $4$.
A mode over $y$ of the group $1$ or $2$ has the charge $\delta_y$, and the modes over $\theta_{\mathrm{s}}y$ of the groups $3$, $4$ have the charge $-\delta_{\theta_{\mathrm{s}}y}=-\delta_y\circ\theta_{\mathrm{s}}$;
a mode over $y$ of the group $3$ or $4$ has the charge $-\delta_y$, and the modes over $\theta_{\mathrm{s}}y$ of the groups $1$, $2$ have the charge $\delta_{\theta_{\mathrm{s}}y}=\delta_y\circ\theta_{\mathrm{s}}$.
This proves the last assertion of (a).

(b) Each $\zeta_{i'}$ is odd; so the $\zeta_{i'}$ anticommute, $\zeta_{i'}^2=0$, and
\[
\zeta^w_I=\sum_{i'_1,\dots,i'_r}\mathfrak{u}_{i'_1i_1}\cdots\mathfrak{u}_{i'_ri_r}\cdot\zeta_{i'_1}\cdots\zeta_{i'_r}.
\]
The terms with two equal indices vanish;
in every other term $\zeta_{i'_1}\cdots\zeta_{i'_r}$ is $\zeta_{I'}$, $I'=\{i'_1,\dots,i'_r\}$, times the sign of the permutation that sorts $(i'_1,\dots,i'_r)$.
Collecting the terms with a given $I'$ gives the Leibniz expansion of the minor $(\wedge\mathfrak{u})_{I'I}$.
For any matrix $T$ let $\wedge T$ denote the matrix of its minors, formed as for $\mathfrak{u}$.
By the Cauchy--Binet formula $\wedge(TT')=(\wedge T)(\wedge T')$ for matrices $T$, $T'$ of compatible sizes;
moreover $\wedge(T^*)=(\wedge T)^*$, because the determinant of the conjugate transpose of a square matrix is the complex conjugate of its determinant.
Hence
\[
(\wedge\mathfrak{u})^*(\wedge\mathfrak{u})=\wedge(\mathfrak{u}^*\mathfrak{u})=\wedge1=1,
\]
and $\wedge\mathfrak{u}$ is unitary, $\mathcal{F}_\Lambda$ and $\mathcal{F}_{\theta_{\mathrm{s}}\Lambda}$ having the same finite dimension.
If $(\wedge\mathfrak{u})_{I'I}\neq0$, some term of its Leibniz expansion is not zero:
there is a bijection $i\mapsto i'$ of $I$ onto $I'$ with $\mathfrak{u}_{i'i}\neq0$ for all $i\in I$.
By (a), $\omega_{i'}=\omega_i$ and $\mathfrak{q}(i')=-\mathfrak{q}(i)\circ\theta_{\mathrm{s}}$ for every $i\in I$;
the product of the first relations and the sum of the second ones over $i\in I$ give $\omega_{I'}=\omega_I$ and $\mathfrak{q}(I')=-\mathfrak{q}(I)\circ\theta_{\mathrm{s}}$.
The entries $(I',I)$ of $(\wedge\mathfrak{u})\,\omega$ and of $\omega\,(\wedge\mathfrak{u})$ are $(\wedge\mathfrak{u})_{I'I}\cdot\omega_I$ and $\omega_{I'}\cdot(\wedge\mathfrak{u})_{I'I}$;
they coincide by what has just been shown, so that $(\wedge\mathfrak{u})\,\omega=\omega\,(\wedge\mathfrak{u})$.

(c) Reversing the order of $r$ anticommuting factors produces the sign $(-1)^{\binom{r}{2}}$.
As $\Theta_w$ reverses products,
$\Theta_w(\zeta_I)=\Theta_w(\zeta_{i_r})\cdots\Theta_w(\zeta_{i_1})=(-1)^{\binom{r}{2}}\,\zeta^w_I$, which is (i).
The normals $e_\nu$ and $e_1$ are orthogonal, so that $\Theta_w\Theta_c=(-1)^{\deg}\,\Theta_c\Theta_w$ by Lemma~\ref{lem:sign-rules}(d);
the automorphism $\Theta_c\Theta_w$ maps generators to linear combinations of generators,
hence homogeneous elements to homogeneous elements of the same degree, and (ii) follows.
For (iii) let $\#|I|\equiv\#|J|$ modulo $2$.
Since $\zeta_J$ is homogeneous of degree $\#|J|$, (i) and (ii) give
\[
\begin{aligned}
\Theta_w\left(\Theta_c(\zeta_J)\,\zeta_I\right)
&=\Theta_w(\zeta_I)\,\Theta_w\Theta_c(\zeta_J)
=(-1)^{\binom{\#|I|}{2}+\#|J|}\,\zeta^w_I\,\Theta_c\left((-1)^{\binom{\#|J|}{2}}\zeta^w_J\right)\\
&=(-1)^{\binom{\#|I|}{2}+\binom{\#|J|}{2}}\cdot(-1)^{\#|J|\,(1+\#|I|)}\,\Theta_c(\zeta^w_J)\,\zeta^w_I,
\end{aligned}
\]
where the last step moves $\Theta_c(\zeta^w_J)$, of degree $\#|J|$, to the left of $\zeta^w_I$, of degree $\#|I|$.
The exponent $\#|J|\,(1+\#|I|)$ is even, which gives (iii).
Finally (iv) follows from (b), since $\Theta_c$ is antilinear.

(d) The operator $\mathfrak{r}$ is diagonal, real and unitary;
hence $\mathfrak{U}$ is unitary, and it commutes with $\omega$ because $\wedge\mathfrak{u}$ and $\mathfrak{r}$ do.
Write $\mathfrak{r}_I:=(-1)^{\binom{\#|I|}{2}}$.
If the values of $\Upsilon$ are even, then $\Upsilon_{IJ}=0$ unless $\#|I|\equiv\#|J|$ modulo $2$, since $U_{\mathrm{f}}e_I=(-1)^{\#|I|}e_I$.
The rule~\eqref{eq:setting-reflection-coefficients} together with (iii), (iv) and (b) gives
\[
\begin{aligned}
\Theta_w\left(\Ins_c[\Upsilon]\right)(\hat A)
&=\sum_{I,J}\overline{\Upsilon_{IJ}(\theta_w^*\hat A)}\cdot\Theta_w\left(\Theta_c(\zeta_J)\,\zeta_I\right)
=\sum_{I,J}\mathfrak{r}_I\cdot\overline{\Upsilon_{IJ}(\theta_w^*\hat A)}\cdot\mathfrak{r}_J\cdot\Theta_c(\zeta^w_J)\,\zeta^w_I\\
&=\sum_{I',J'}\left[\sum_{I,J}(\wedge\mathfrak{u})_{I'I}\cdot\mathfrak{r}_I\cdot\overline{\Upsilon_{IJ}(\theta_w^*\hat A)}\cdot\mathfrak{r}_J\cdot\overline{(\wedge\mathfrak{u})_{J'J}}\right]\cdot\Theta_c(\zeta_{J'})\,\zeta_{I'}.
\end{aligned}
\]
The bracket is the entry $(I',J')$ of $\mathfrak{U}\,\overline{\Upsilon(\theta_w^*\hat A)}\,\mathfrak{U}^*$,
because $\mathfrak{U}_{I'I}=(\wedge\mathfrak{u})_{I'I}\cdot\mathfrak{r}_I$ and $(\mathfrak{U}^*)_{JJ'}=\overline{\mathfrak{U}_{J'J}}=\mathfrak{r}_J\cdot\overline{(\wedge\mathfrak{u})_{J'J}}$.
This is the first identity.
For the second let $a\in\R^{\theta_{\mathrm{s}}\Lambda}$.
The diagonal operators $\mathfrak{r}$ and $\overline{u(a\circ\theta_{\mathrm{s}})}$ commute and $\mathfrak{r}^2=1$,
so that $\mathfrak{U}\,\overline{u(a\circ\theta_{\mathrm{s}})}\,\mathfrak{U}^*=(\wedge\mathfrak{u})\,\overline{u(a\circ\theta_{\mathrm{s}})}\,(\wedge\mathfrak{u})^*$.
Whenever $(\wedge\mathfrak{u})_{I'I}\neq0$, (b) gives $\mathfrak{q}(I')=-\mathfrak{q}(I)\circ\theta_{\mathrm{s}}$,
hence $\langle a\circ\theta_{\mathrm{s}},\mathfrak{q}(I)\rangle=-\langle a,\mathfrak{q}(I')\rangle$, $\theta_{\mathrm{s}}$ being an involution of $\Z^2$, and
\[
\left(\mathfrak{U}\,\overline{u(a\circ\theta_{\mathrm{s}})}\,\mathfrak{U}^*\right)_{I'J'}
=\sum_I(\wedge\mathfrak{u})_{I'I}\cdot\exp\left(-iq\langle a\circ\theta_{\mathrm{s}},\mathfrak{q}(I)\rangle\right)\cdot\overline{(\wedge\mathfrak{u})_{J'I}}
=\exp\left(iq\langle a,\mathfrak{q}(I')\rangle\right)\cdot\left((\wedge\mathfrak{u})(\wedge\mathfrak{u})^*\right)_{I'J'}.
\]
By (b) the last matrix is the identity, and the right side is the entry $(I',J')$ of $u(a)$.
\end{proof}

\begin{lemma}[Parity]\label{lem:fock-parity}
On $\mathcal{F}_\Lambda$ one has $U_{\mathrm{f}}=u(a)$ for the constant function $a\equiv\pi/q$ on $\Lambda$.
Consequently every $\Upsilon\in\End_z(\Lambda)$, $z\in\Z^\Lambda$, satisfies
\[
U_{\mathrm{f}}\,\Upsilon\,U_{\mathrm{f}}=(-1)^{\sum_{y\in\Lambda}z(y)}\,\Upsilon;
\]
in particular the elements of total charge zero are even.
\end{lemma}

\begin{proof}
For $I\subset\mathcal{J}^{\mathrm{l}}(\Lambda)$ the charge $\mathfrak{q}(I)=\sum_{i\in I}\mathfrak{q}(i)$ is a sum of $\#|I|$ functions of the form $\pm\delta_y$;
therefore $\sum_y\mathfrak{q}(I)(y)\equiv\#|I|$ modulo $2$, and
\[
u(a)\,e_I=\exp\left(i\pi\sum_{y}\mathfrak{q}(I)(y)\right)\cdot e_I=(-1)^{\#|I|}\,e_I=U_{\mathrm{f}}\,e_I .
\]
For $\Upsilon\in\End_z(\Lambda)$ the defining relation of Definition~\ref{def:charged-operators}, taken with this $a$, reads
$U_{\mathrm{f}}\,\Upsilon\,U_{\mathrm{f}}^{-1}=\exp\left(i\pi\sum_yz(y)\right)\cdot\Upsilon$,
and $U_{\mathrm{f}}^{-1}=U_{\mathrm{f}}$.
\end{proof}

\begin{lemma}[Reflected integrands]\label{lem:reflected-integrands}
Let $\Pi'$ be a mirror with the reflection $\theta_{\Pi'}$, let $B$ be a box and $G\in\mathcal{G}^{\mathrm{ev}}(B)$.
Then $\Theta_{\Pi'}G\in\mathcal{G}^{\mathrm{ev}}(\theta_{\Pi'}B)$, with the same growth margin $\varepsilon_G$.
Hence $G^g\in\mathcal{G}^{\mathrm{ev}}(gB)$ for every $g\in W$.
\end{lemma}

\begin{proof}
Let $n'$ be the upward unit normal of $\Pi'$.
The linear part of $\theta_{\Pi'}$ maps every cube centred at the origin onto itself.
Hence $\theta_{\Pi'}B$ is a box,
$\theta_{\Pi'}$ maps $B^{(1)}$ and $B^{(2)}$ onto $(\theta_{\Pi'}B)^{(1)}$ and $(\theta_{\Pi'}B)^{(2)}$,
and $\mathbb{X}(\theta_{\Pi'}B)=\theta_{\Pi'}\mathbb{X}(B)$,
because $\pi\circ\theta_{\Pi'}=\theta_{\Pi'}\circ\pi$ (Proposition~\ref{prop:geometry}\eqref{it:geometry-projection}) and the cube $[-\frac12,\frac12]^3$ is mapped onto itself.

Write $G(\hat A)=\sum_\Psi a_\Psi(\hat A)\,\Psi$.
By~\eqref{eq:setting-reflection-coefficients},
$(\Theta_{\Pi'}G)(\hat A)=\sum_\Psi\overline{a_\Psi(\theta_{\Pi'}^*\hat A)}\cdot\Theta_{\Pi'}(\Psi)$ for $\hat A\in\R^{(\theta_{\Pi'}B)^{(1)}}$,
and by~\eqref{eq:setting-reflection-fields} each $\Theta_{\Pi'}(\Psi)$ is a linear combination of monomials of the same degree in the fields at the sites of $\theta_{\Pi'}\mathbb{X}(B)$,
with coefficients that are products of entries of $\gamma_{n'}$, up to sign.
So every coefficient of $\Theta_{\Pi'}G$ is a finite combination, with constant coefficients, of the functions $\overline{a_\Psi(\theta_{\Pi'}^*\hat A)}$.
The pull-back $\theta_{\Pi'}^*$ commutes with $d$ (Definition~\ref{def:cochains}) and carries the values on the plaquettes of $\theta_{\Pi'}B$ to the values on the plaquettes of $B$, up to signs; hence
\[
\sum_{\sigma\in B^{(2)}}\left(d\,\theta_{\Pi'}^*\hat A\right)_\sigma^2=\sum_{\sigma\in(\theta_{\Pi'}B)^{(2)}}\left(d\hat A\right)_\sigma^2 ,
\]
and the bound of Definition~\ref{def:integrands} holds for $\Theta_{\Pi'}G$ with the same $\varepsilon_G$ and a larger constant.

Gauge invariance remains.
For a real function $\lambda$ on the sites of a box let $U_\lambda$ be the gauge transformation~\eqref{eq:setting-gauge} acting on integrands:
$(U_\lambda G)(\hat A)$ arises from $G(\hat A-d\lambda)$ by the substitution
$\psi(x)\mapsto\exp\left(iq\lambda(\pi x)\right)\psi(x)$, $\bar\psi(x)\mapsto\exp\left(-iq\lambda(\pi x)\right)\bar\psi(x)$,
and the gauge invariance of $G$ states that $U_\lambda G=G$ for every real $\lambda$ on the sites of $B$.
As $\pi\circ\theta_{\Pi'}=\theta_{\Pi'}\circ\pi$ and $\theta_{\Pi'}^2=\mathrm{id}$, one has $\lambda(\pi x)=(\lambda\circ\theta_{\Pi'})(\pi\,\theta_{\Pi'}x)$,
and antilinearity gives on the generators
\[
\begin{aligned}
\Theta_{\Pi'}\left(\exp\left(iq\lambda(\pi x)\right)\psi(x)\right)
&=-\exp\left(-iq(\lambda\circ\theta_{\Pi'})(\pi\,\theta_{\Pi'}x)\right)\bar\psi(\theta_{\Pi'}x)\,\gamma_{n'}
=U_{\lambda\circ\theta_{\Pi'}}\left(\Theta_{\Pi'}\psi(x)\right),\\
\Theta_{\Pi'}\left(\exp\left(-iq\lambda(\pi x)\right)\bar\psi(x)\right)
&=-\exp\left(iq(\lambda\circ\theta_{\Pi'})(\pi\,\theta_{\Pi'}x)\right)\gamma_{n'}\,\psi(\theta_{\Pi'}x)
=U_{\lambda\circ\theta_{\Pi'}}\left(\Theta_{\Pi'}\bar\psi(x)\right).
\end{aligned}
\]
On the Grassmann algebra the maps $\Theta_{\Pi'}U_\lambda$ and $U_{\lambda\circ\theta_{\Pi'}}\Theta_{\Pi'}$ are antilinear and reverse products;
as they agree on the generators, they agree everywhere.
On the coefficients,
$\overline{a_\Psi\left(\theta_{\Pi'}^*\hat A-d\lambda\right)}=\overline{a_\Psi\left(\theta_{\Pi'}^*\left(\hat A-d(\lambda\circ\theta_{\Pi'})\right)\right)}$,
because $\theta_{\Pi'}^*$ commutes with $d$ and $\theta_{\Pi'}^*(\lambda\circ\theta_{\Pi'})=\lambda\circ\theta_{\Pi'}\circ\theta_{\Pi'}=\lambda$.
Therefore $\Theta_{\Pi'}U_\lambda=U_{\lambda\circ\theta_{\Pi'}}\Theta_{\Pi'}$.
Given a real function $\lambda'$ on the sites of $\theta_{\Pi'}B$, put $\lambda:=\lambda'\circ\theta_{\Pi'}$;
then $U_{\lambda'}\Theta_{\Pi'}G=\Theta_{\Pi'}U_\lambda G=\Theta_{\Pi'}G$.
Degrees are preserved, so $\Theta_{\Pi'}G$ is even, and $\Theta_{\Pi'}G\in\mathcal{G}^{\mathrm{ev}}(\theta_{\Pi'}B)$.

The group $W$ is generated by the reflections in mirrors (Definition~\ref{def:mirrors}),
and for a word $g=\theta_{\Pi_1}\cdots\theta_{\Pi_r}$ one has $G^g=\Theta_{\Pi_1}\cdots\Theta_{\Pi_r}G$ by~\eqref{eq:setting-group-action}.
Applying the first part $r$ times gives $G^g\in\mathcal{G}^{\mathrm{ev}}(gB)$.
\end{proof}

\subsection{The mid-plane algebra}\label{sec:mid-plane-algebra}

Throughout this subsection $S$ is a finite set of links of the slice lattice $\Z^2$,
and $\Lambda$ is a non-empty finite set of sites of $\Z^2$ that contains the end points of the links of $S$.
As in Definition~\ref{def:mid-plane}, $\mathcal{H}_S=L^2(\R^S)\otimes\mathcal{F}_\Lambda$,
where $\mathcal{F}_\Lambda$ is the Fock space of the modes over $\Lambda$.
For $z\in\Z^\Lambda$, $\End_z(\Lambda)$ is the space of the operators $\Upsilon$ on $\mathcal{F}_\Lambda$ of charge $z$,
that is, of those with $u(a)\Upsilon u(a)^{-1}=\exp(iq\langle a,z\rangle)\cdot\Upsilon$ for all $a\in\R^\Lambda$
(Definition~\ref{def:charged-operators}).
The variable of the first factor of $\mathcal{H}_S$ is $A\in\R^S$;
$E_S=(E_b)_{b\in S}$ is the family of the commuting self-adjoint operators $E_b=-i\,\partial/\partial A(b)$,
and $\xi\in\R^S$ is the spectral variable of $E_S$.
A vector $\delta\in\R^S$ is extended by zero to all links of $\Z^2$,
as in Definition~\ref{def:mid-plane}\eqref{it:exponential-generators}.

\begin{definition}[Gaussian symbols and exponents]\label{def:gaussian-symbols}
Let $S$ be a finite set of links in the slice lattice $\Z^2$.
For a subset $S'$ of $S$, a positive definite real quadratic form $\mathfrak{Q}$ on $\R^{S'}$ and $\lambda\in\C^{S'}$,
the function $f_{S',\mathfrak{Q},\lambda}:\R^S\to\C$ is
\[
f_{S',\mathfrak{Q},\lambda}(\xi):=\exp\left(-\mathfrak{Q}(\xi_{S'})+\langle\lambda,\xi_{S'}\rangle\right),
\]
where $\xi_{S'}$ is the restriction of $\xi$ to $S'$ and $\langle\lambda,\xi_{S'}\rangle=\sum_{b\in S'}\lambda_b\cdot\xi_b$;
for $S'=\emptyset$ it is the constant function $1$.
The linear span of all these functions is denoted by $\mathfrak{A}(S)$,
and its elements are called \emph{Gaussian symbols} on $\R^S$.
An \emph{exponent} is an element of
\[
\mathfrak{D}(S):=\left\{\delta\in\R^S:\ d_{\mathrm{s}}^*\delta\in q\Z^{\Z^2}\right\},
\]
where $\delta$ is extended by zero to all links of $\Z^2$ and $d_{\mathrm{s}}^*$ is the adjoint of $d_{\mathrm{s}}$ on finitely supported cochains;
the \emph{charge} of an exponent $\delta$ is $z(\delta):=q^{-1}d_{\mathrm{s}}^*\delta$.
\end{definition}

The charge of an exponent lies in $\Z^\Lambda$:
the end points of the links of $S$ belong to $\Lambda$, so that $d_{\mathrm{s}}^*\delta$ vanishes outside $\Lambda$.
Its total charge is zero, $\sum_yz(\delta)(y)=0$,
because $\langle d_{\mathrm{s}}^*\delta,a\rangle=\langle\delta,d_{\mathrm{s}}a\rangle=0$ for the constant function $a=1$;
in other words, $z(\delta)$ is the divergence of a cochain of finite support.
The exponents form a subgroup of $\R^S$, and $z(\delta+\delta')=z(\delta)+z(\delta')$.

The elements of $\mathfrak{A}(S)$ are bounded continuous functions on $\R^S$:
the real part of the exponent of $f_{S',\mathfrak{Q},\lambda}$ is $-\mathfrak{Q}(\xi_{S'})+\langle\Re\lambda,\xi_{S'}\rangle$,
which is bounded above because $\mathfrak{Q}$ is positive definite.
A \emph{symbol} is an element $\Upsilon$ of $\mathfrak{A}(S)\otimes\End(\mathcal{F}_\Lambda)$, that is, a finite sum
\[
\Upsilon=\sum_jf_j\otimes\Upsilon_j\qquad\text{with }f_j\in\mathfrak{A}(S)\text{ and }\Upsilon_j\in\End(\mathcal{F}_\Lambda).
\]
Expanding the operators $\Upsilon_j$ in a basis of the finite-dimensional space $\End(\mathcal{F}_\Lambda)$,
one sees that $\Upsilon$ is determined by the function
\[
\R^S\to\End(\mathcal{F}_\Lambda),\qquad\xi\mapsto\Upsilon(\xi):=\sum_jf_j(\xi)\cdot\Upsilon_j,
\]
and $\Upsilon$ is identified with this bounded continuous function.
Thus an operator $\Upsilon$ on $\mathcal{F}_\Lambda$ is the constant symbol $1\otimes\Upsilon$,
and for a symbol $\Upsilon$ the function $\Upsilon^*$ is given by $\Upsilon^*(\xi):=\Upsilon(\xi)^*$.
For a symbol $\Upsilon$, $\Upsilon(E_S)$ is the operator on $\mathcal{H}_S$ that Definition~\ref{def:integral-kernels} attaches to the bounded function $\Upsilon$:
after the Fourier transform it acts as the multiplication by $\Upsilon(\xi)$.
For $x\in\R^S$, the operator $\Upsilon(E_S+x)$ is attached in the same way to the function $\xi\mapsto\Upsilon(\xi+x)$.
Definition~\ref{def:integral-kernels}, applied to the scalar functions $f_j$, gives
\[
\Upsilon(E_S)=\sum_jf_j(E_S)\otimes\Upsilon_j,\qquad\Upsilon(E_S+x)=\sum_jf_j(E_S+x)\otimes\Upsilon_j,
\]
where $f_j(E_S)$ and $f_j(E_S+x)$ are the operators of the joint spectral calculus of $E_S$ on $L^2(\R^S)$ attached to $f_j$ and to $\xi\mapsto f_j(\xi+x)$.
With another family of commuting self-adjoint operators indexed by $S$ in the place of $E_S$,
the same two sums define operators on the tensor product of the space on which that family acts with $\mathcal{F}_\Lambda$;
expanding in a basis of $\End(\mathcal{F}_\Lambda)$ and using the linearity of the joint spectral calculus,
one sees that they depend only on the function $\Upsilon$, not on the chosen sum.

\begin{lemma}[Normal form]\label{lem:normal-form}
Let $S$ be a finite set of links and $\Lambda$ a non-empty finite set of sites of the slice lattice $\Z^2$,
such that $\Lambda$ contains the end points of the links of $S$.
\begin{enumerate}
\item[(a)]
The linear space $\mathfrak{A}(S)$ contains the constant function $1$,
and for $f,f'\in\mathfrak{A}(S)$ and $x\in\R^S$ it contains the product $f\cdot f'$, the complex conjugate $\overline{f}$ and the translate $f(\cdot+x)$.
Every element of $\mathfrak{A}(S)$ is a bounded, uniformly continuous function on $\R^S$.
\item[(b)]
$\mathfrak{W}(S,\Lambda)$ is the set of the finite sums
\[
X=\sum_{\delta\in\mathfrak{D}_X}\exp(i\langle\delta,A\rangle)\cdot\Upsilon_\delta(E_S),
\]
where $\mathfrak{D}_X$ is a finite subset of $\mathfrak{D}(S)$
and $\Upsilon_\delta\in\mathfrak{A}(S)\otimes\End_{z(\delta)}(\Lambda)$ for every $\delta\in\mathfrak{D}_X$.
Such a sum is called a \emph{normal form} of $X$, with the exponents $\delta\in\mathfrak{D}_X$ and the symbols $\Upsilon_\delta$.
For $\delta,\delta'\in\mathfrak{D}(S)$, $\Upsilon\in\mathfrak{A}(S)\otimes\End_{z(\delta)}(\Lambda)$
and $\Upsilon'\in\mathfrak{A}(S)\otimes\End_{z(\delta')}(\Lambda)$ one has
\begin{equation}\label{eq:amplification-product-rule}
\begin{split}
&\exp(i\langle\delta,A\rangle)\cdot\Upsilon(E_S)\cdot\exp(i\langle\delta',A\rangle)\cdot\Upsilon'(E_S)\\
&\qquad=\exp(i\langle\delta+\delta',A\rangle)\cdot\Upsilon(E_S+\delta')\cdot\Upsilon'(E_S),
\end{split}
\end{equation}
\begin{equation}\label{eq:amplification-adjoint-rule}
\left(\exp(i\langle\delta,A\rangle)\cdot\Upsilon(E_S)\right)^*=\exp(-i\langle\delta,A\rangle)\cdot\Upsilon^*(E_S-\delta).
\end{equation}
\item[(c)]
A sum as in~(b) is the zero operator on $\mathcal{H}_S$ if and only if $\Upsilon_\delta=0$ for every $\delta\in\mathfrak{D}_X$.
\item[(d)]
Every $X\in\mathfrak{W}(S,\Lambda)$ is even, that is, it commutes with the parity $U_{\mathrm{f}}$ of the factor $\mathcal{F}_\Lambda$;
and it commutes with the gauge transformations $U^{S}_{\mathrm{G}}(a)$, $a\in\R^\Lambda$, of $\mathcal{H}_S$,
\[
\left(U^{S}_{\mathrm{G}}(a)\Phi\right)(A):=u(a)\,\Phi(A-d_{\mathrm{s}}a),
\]
where $d_{\mathrm{s}}a$ is restricted to the links of $S$.
\item[(e)]
For $\beta\in\R^S$ and $X$ as in~(b), the map $\alpha_\beta$ of Definition~\ref{def:mid-plane} is given by
\[
\alpha_\beta X=\sum_{\delta\in\mathfrak{D}_X}\exp(i\langle\delta,A\rangle)\cdot\Upsilon_\delta(E_S-\beta)
=\exp(i\langle\beta,A\rangle)\,X\,\exp(-i\langle\beta,A\rangle),
\]
and it is a $*$-automorphism of $\mathfrak{W}(S,\Lambda)$.
\item[(f)]
Let $\Xi$ be a rectangular cross-section in which $\hull(S,\Lambda)$ embeds (Definition~\ref{def:slices}).
Through the embedding, the links of $S$ are links of $\Lambda_{\mathrm{s}}$ and the sites of $\Lambda$ are sites of $\Lambda_{\mathrm{s}}$;
a vector $\delta\in\R^S$ is extended by zero to $\Lambda^{(1)}_{\mathrm{s}}$.
For $\delta\in\mathfrak{D}(S)$ the function $A\mapsto\exp(i\langle\delta,A\rangle)$ on $\R^{\Lambda^{(1)}_{\mathrm{s}}}$ is $\Gamma_{\mathrm{s}}$-periodic.
For $X$ as in~(b), let $X_\Xi$ be the operator on $\mathcal{H}^{\mathrm{l}}$
that acts on each summand $L^2\left(\R^{\Lambda^{(1)}_{\mathrm{s}}}/\Gamma_{\mathrm{s}};\mathcal{F}\right)$ of the direct sum as
\[
\sum_{\delta\in\mathfrak{D}_X}\exp(i\langle\delta,A\rangle)\cdot\iota\left(\Upsilon_\delta(E_S)\right),
\]
where $\Upsilon_\delta(E_S)$ is formed with the operators $E_b$, $b\in S$, of Definition~\ref{def:slice-operators},
and the localisation $\iota$ of Definition~\ref{def:modes} is applied to the Fock factor.
Then $X_\Xi$ is well defined,
and $X\mapsto X_\Xi$ is a unital $*$-representation of $\mathfrak{W}(S,\Lambda)$ on $\mathcal{H}^{\mathrm{l}}$,
which on the generators is given by the formulas of Definition~\ref{def:mid-plane}.
Every $X_\Xi$ commutes with $U_{\mathrm{G}}(a)$ for all $a\in\R^{\Lambda_{\mathrm{s}}}$, hence with $P_{\mathrm{G}}$,
and with $f(E_b)$ for every link $b\in\Lambda^{(1)}_{\mathrm{s}}\setminus S$ and every bounded Borel function $f$ on $\R$.
The operator $X_\Xi$ does not depend on the plaquette $\sigma_0$.
\end{enumerate}
\end{lemma}

\begin{proof}
(a) Let $S''$ be a second subset of $S$, $\mathfrak{Q}''$ a positive definite real quadratic form on $\R^{S''}$, and $\lambda''\in\C^{S''}$.
Then
\[
f_{S',\mathfrak{Q},\lambda}\cdot f_{S'',\mathfrak{Q}'',\lambda''}=f_{S'\cup S'',\,\mathfrak{Q}+\mathfrak{Q}'',\,\lambda+\lambda''},
\]
where the quadratic form $\mathfrak{Q}+\mathfrak{Q}''$ on $\R^{S'\cup S''}$ takes the value $\mathfrak{Q}(\xi_{S'})+\mathfrak{Q}''(\xi_{S''})$ at $\xi_{S'\cup S''}$,
and $\lambda$, $\lambda''$ are extended by zero to $S'\cup S''$.
This form is positive definite:
it is non-negative, and $\mathfrak{Q}(\xi_{S'})+\mathfrak{Q}''(\xi_{S''})=0$ forces $\xi_{S'}=0$ and $\xi_{S''}=0$, that is, $\xi_{S'\cup S''}=0$.
Since $\mathfrak{Q}$ is real and $\xi$ is real, the complex conjugate of $f_{S',\mathfrak{Q},\lambda}$ is $f_{S',\mathfrak{Q},\overline{\lambda}}$.
Denote by $\mathfrak{Q}(\cdot,\cdot)$ the symmetric bilinear form with $\mathfrak{Q}(\xi_{S'},\xi_{S'})=\mathfrak{Q}(\xi_{S'})$.
For $x\in\R^S$ the polar identity
$\mathfrak{Q}(\xi_{S'}+x_{S'})=\mathfrak{Q}(\xi_{S'})+2\,\mathfrak{Q}(\xi_{S'},x_{S'})+\mathfrak{Q}(x_{S'})$ gives
\[
f_{S',\mathfrak{Q},\lambda}(\xi+x)=\exp\left(-\mathfrak{Q}(x_{S'})+\langle\lambda,x_{S'}\rangle\right)\cdot f_{S',\mathfrak{Q},\lambda'}(\xi),
\]
where $\lambda'\in\C^{S'}$ is determined by $\langle\lambda',\xi_{S'}\rangle=\langle\lambda,\xi_{S'}\rangle-2\,\mathfrak{Q}(\xi_{S'},x_{S'})$ for all $\xi$.
So the span $\mathfrak{A}(S)$, which contains $1$, is closed under products, under complex conjugation and under translations.
Finally,
$\left|f_{S',\mathfrak{Q},\lambda}(\xi)\right|=\exp\left(-\mathfrak{Q}(\xi_{S'})+\langle\Re\lambda,\xi_{S'}\rangle\right)$
with $\mathfrak{Q}$ positive definite;
hence $f_{S',\mathfrak{Q},\lambda}$ is a Schwartz function of $\xi_{S'}$ that does not depend on the other components of $\xi$.
Such a function is smooth and bounded together with all its derivatives, and so is every finite linear combination of such functions;
a bounded gradient makes it uniformly continuous.

(b) For $x,\delta\in\R^S$ the two operators
$\exp(-i\langle\delta,A\rangle)\exp(-i\langle x,E_S\rangle)\exp(i\langle\delta,A\rangle)$ and $\exp(-i\langle x,E_S+\delta\rangle)$
both map $\Phi\in\mathcal{H}_S$ to the function $A\mapsto\exp(-i\langle\delta,x\rangle)\cdot\Phi(A-x)$.
A unitary group $x\mapsto\exp(-i\langle x,\cdot\rangle)$ determines the joint spectral calculus of its generators;
therefore
\[
f(E_S)\cdot\exp(i\langle\delta,A\rangle)=\exp(i\langle\delta,A\rangle)\cdot f(E_S+\delta)
\]
for every bounded Borel function $f$ on $\R^S$:
this is the relation~\eqref{eq:setting-weyl}, read on $\mathcal{H}_S$.
The operators on the Fock factor commute with the multiplication operators $\exp(i\langle\delta,A\rangle)$.
It follows that
\[
\begin{aligned}
&\exp(i\langle\delta,A\rangle)\cdot\Upsilon(E_S)\cdot\exp(i\langle\delta',A\rangle)\cdot\Upsilon'(E_S)\\
&\qquad=\exp(i\langle\delta,A\rangle)\cdot\exp(i\langle\delta',A\rangle)\cdot\Upsilon(E_S+\delta')\cdot\Upsilon'(E_S),
\end{aligned}
\]
which is~\eqref{eq:amplification-product-rule}, and that
\[
\left(\exp(i\langle\delta,A\rangle)\cdot\Upsilon(E_S)\right)^*
=\Upsilon^*(E_S)\cdot\exp(-i\langle\delta,A\rangle)
=\exp(-i\langle\delta,A\rangle)\cdot\Upsilon^*(E_S-\delta),
\]
which is~\eqref{eq:amplification-adjoint-rule}.
After the Fourier transform, the operator $\Upsilon(E_S+\delta')\cdot\Upsilon'(E_S)$ acts as the multiplication by $\Upsilon(\xi+\delta')\cdot\Upsilon'(\xi)$.
By~(a), the functions $\xi\mapsto\Upsilon(\xi+\delta')\cdot\Upsilon'(\xi)$ and $\xi\mapsto\Upsilon^*(\xi-\delta)$ are symbols,
with Fock factors in $\End_{z(\delta)+z(\delta')}(\Lambda)=\End_{z(\delta+\delta')}(\Lambda)$ and in $\End_{-z(\delta)}(\Lambda)=\End_{z(-\delta)}(\Lambda)$ respectively;
indeed $\End_z(\Lambda)\cdot\End_{z'}(\Lambda)\subset\End_{z+z'}(\Lambda)$, and $\End_z(\Lambda)^*=\End_{-z}(\Lambda)$ because $u(a)$ is unitary.
Hence the set of the finite sums in~(b) is closed under products and adjoints.
It is a linear space and contains the identity, the sum with $\mathfrak{D}_X=\{0\}$ and $\Upsilon_0=1$;
so it is a unital $*$-algebra of operators on $\mathcal{H}_S$.
It contains the generators of $\mathfrak{W}(S,\Lambda)$.
A generator $f(E_{S'})$ as in Definition~\ref{def:mid-plane}\eqref{it:gaussian-generators},
with $f$ in the span of $1$ and of the functions $\exp\left(-\mathfrak{Q}(\cdot)+\langle\lambda,\cdot\rangle\right)$ on $\R^{S'}$,
is the operator of the element $\xi\mapsto f(\xi_{S'})$ of $\mathfrak{A}(S)$, with the exponent $0$.
A generator $\exp(i\langle\delta,A\rangle)\cdot\Upsilon$ as in Definition~\ref{def:mid-plane}\eqref{it:exponential-generators}
has $\delta\in\mathfrak{D}(S)$ and $\Upsilon\in\End_{z(\delta)}(\Lambda)$, because $d_{\mathrm{s}}^*\delta=q\cdot\mathfrak{q}(\Upsilon)$.
Conversely, each term of a sum in~(b) is a finite sum of operators $\exp(i\langle\delta,A\rangle)\cdot\Upsilon\cdot f(E_S)$
with $\Upsilon\in\End_{z(\delta)}(\Lambda)$ and $f\in\mathfrak{A}(S)$,
and each of these is the product of the generator $\exp(i\langle\delta,A\rangle)\cdot\Upsilon$
with $f(E_S)$, a linear combination of generators as in Definition~\ref{def:mid-plane}\eqref{it:gaussian-generators}.
Therefore the set of the sums in~(b) is the unital $*$-algebra generated by these generators, namely $\mathfrak{W}(S,\Lambda)$.

(c) Let $\Phi\mapsto\widehat{\Phi}$ be the unitary Fourier transform of $\mathcal{H}_S$ onto $L^2(\R^S,\mathrm{d}\xi;\mathcal{F}_\Lambda)$
under which $E_S$ becomes the multiplication by $\xi$.
It carries $\exp(i\langle\delta,A\rangle)$ into the translation by $\delta$, so that for $X$ as in~(b)
\[
\widehat{X\Phi}(\xi)=\sum_{\delta\in\mathfrak{D}_X}\Upsilon_\delta(\xi-\delta)\,\widehat{\Phi}(\xi-\delta).
\]
Let $r>0$ be smaller than half the distance between any two distinct points of $\mathfrak{D}_X$, let $\xi_0\in\R^S$,
and let $\widehat{\Phi}$ vanish outside the ball of radius $r$ about $\xi_0$.
The term of $\delta$ then vanishes outside the ball of radius $r$ about $\xi_0+\delta$, and these balls are pairwise disjoint.
So $X=0$ implies $\Upsilon_\delta(\xi)\,\widehat{\Phi}(\xi)=0$ for almost every $\xi$, for every $\delta\in\mathfrak{D}_X$ and every such $\Phi$.
Taking for $\widehat{\Phi}$, in turn, each vector of a basis of $\mathcal{F}_\Lambda$ on the ball and zero outside,
and using the continuity of the symbols~(a), one obtains $\Upsilon_\delta=0$ on the ball.
As $\xi_0$ is arbitrary, $\Upsilon_\delta=0$.
If, conversely, all the symbols vanish, the sum is zero.

(d) Let $a\in\R^\Lambda$, $\delta\in\mathfrak{D}(S)$ and $\Upsilon\in\End_{z(\delta)}(\Lambda)$.
Since $\langle\delta,d_{\mathrm{s}}a\rangle=\langle d_{\mathrm{s}}^*\delta,a\rangle=q\langle z(\delta),a\rangle$,
\[
\begin{aligned}
&U^{S}_{\mathrm{G}}(a)\,\exp(i\langle\delta,A\rangle)\cdot\Upsilon\,U^{S}_{\mathrm{G}}(a)^{-1}\\
&\qquad=\exp(i\langle\delta,A-d_{\mathrm{s}}a\rangle)\cdot u(a)\Upsilon u(a)^{-1}\\
&\qquad=\exp(i\langle\delta,A\rangle)\cdot\exp(-iq\langle z(\delta),a\rangle)\cdot\exp(iq\langle a,z(\delta)\rangle)\cdot\Upsilon,
\end{aligned}
\]
and the right side is $\exp(i\langle\delta,A\rangle)\cdot\Upsilon$.
The translations $\Phi\mapsto\Phi(\cdot-d_{\mathrm{s}}a)$ commute with all functions of $E_S$, and $u(a)$ acts on the second factor only.
Hence $U^{S}_{\mathrm{G}}(a)$ commutes with every operator $\exp(i\langle\delta,A\rangle)\cdot\Upsilon\cdot f(E_S)$ with $f\in\mathfrak{A}(S)$,
therefore with every term of a sum as in~(b), and with every $X\in\mathfrak{W}(S,\Lambda)$.
As to evenness, Lemma~\ref{lem:fock-parity} gives $U_{\mathrm{f}}\Upsilon U_{\mathrm{f}}=(-1)^{\sum_yz(y)}\cdot\Upsilon$ for $\Upsilon\in\End_z(\Lambda)$.
The total charge of $z(\delta)$ is zero; so the Fock factors of the symbols $\Upsilon_\delta$ commute with $U_{\mathrm{f}}$, and so does $X$.

(e) The multiplication by $\exp(i\langle\beta,A\rangle)$ is a unitary operator on $\mathcal{H}_S$.
By the relation obtained in the proof of~(b),
$\exp(i\langle\beta,A\rangle)\,f(E_S)\,\exp(-i\langle\beta,A\rangle)=f(E_S-\beta)$ for every bounded Borel function $f$ on $\R^S$,
while $\exp(i\langle\beta,A\rangle)$ commutes with $\exp(i\langle\delta,A\rangle)$ and with the operators on the Fock factor.
This gives the identity in~(e).
By~(a), the symbols $\xi\mapsto\Upsilon_\delta(\xi-\beta)$ lie again in $\mathfrak{A}(S)\otimes\End_{z(\delta)}(\Lambda)$;
so, by~(b), $\alpha_\beta$ maps $\mathfrak{W}(S,\Lambda)$ into itself, and so does its inverse $\alpha_{-\beta}$.
A conjugation by a unitary operator is a $*$-automorphism of the algebra of bounded operators.
Hence $\alpha_\beta$ is a $*$-automorphism of $\mathfrak{W}(S,\Lambda)$;
it maps a generator $f(E_{S'})$ to $f(E_{S'}-\beta_{S'})$ and fixes the generators $\exp(i\langle\delta,A\rangle)\cdot\Upsilon$,
as in Definition~\ref{def:mid-plane}.

(f) Let $\gamma\in\Gamma_{\mathrm{s}}$, so that $\gamma=(2\pi/q)\cdot k$ with a closed integer one-cochain $k$ of $\Lambda_{\mathrm{s}}$.
On the links of $\hull(S,\Lambda)$, a product of two integer intervals that embeds in $\Xi$, the cochain $k$ is closed;
as this product is contractible, $k=d_{\mathrm{s}}a'$ there, with an integer-valued function $a'$ on the sites of $\hull(S,\Lambda)$.
The cochain $\delta$ is supported in $S$ and $d_{\mathrm{s}}^*\delta$ in $\Lambda$, both inside the hull; hence
\[
\langle\delta,\gamma\rangle=\frac{2\pi}{q}\cdot\langle\delta,d_{\mathrm{s}}a'\rangle=\frac{2\pi}{q}\cdot\langle d_{\mathrm{s}}^*\delta,a'\rangle=2\pi\langle z(\delta),a'\rangle\in2\pi\Z,
\]
and $\exp(i\langle\delta,A\rangle)$ is $\Gamma_{\mathrm{s}}$-periodic.
It is therefore a unitary multiplication operator on $L^2\left(\R^{\Lambda^{(1)}_{\mathrm{s}}}/\Gamma_{\mathrm{s}}\right)$,
and there the relation~\eqref{eq:setting-weyl} holds (Definition~\ref{def:slice-operators}).
The localisation $\iota$ is a unital $*$-homomorphism (Proposition~\ref{prop:localisation}\eqref{it:localisation-homomorphism}).
Consequently the rules~\eqref{eq:amplification-product-rule} and~\eqref{eq:amplification-adjoint-rule} remain true
when every operator $\exp(i\langle\delta,A\rangle)\cdot\Upsilon(E_S)$ on $\mathcal{H}_S$ in them
is replaced by the operator $\exp(i\langle\delta,A\rangle)\cdot\iota\left(\Upsilon(E_S)\right)$ on $L^2\left(\R^{\Lambda^{(1)}_{\mathrm{s}}}/\Gamma_{\mathrm{s}}\right)\otimes\mathcal{F}$,
and every operator of a symbol on the right sides by its image under $\iota$.
By~(c), the symbols $\Upsilon_\delta$ are determined by $X$; hence $X_\Xi$ is well defined, and $X\mapsto X_\Xi$ is linear.
Comparing, by means of~(b), the normal forms of products and of adjoints,
one finds that this map is multiplicative and preserves adjoints; it is unital because $\iota$ is.
On a generator $f(E_{S'})$ it gives $f(E_{S'})$ on each summand, and on a generator $\exp(i\langle\delta,A\rangle)\cdot\Upsilon$ it gives $\exp(i\langle\delta,A\rangle)\cdot\iota(\Upsilon)$,
as in Definition~\ref{def:mid-plane}.

Let $a\in\R^{\Lambda_{\mathrm{s}}}$.
The computation in the proof of~(d),
with the charge rule~\eqref{eq:amplification-iota-charge}, $u(a)\iota(\Upsilon)u(a)^{-1}=\exp(iq\langle a,z\rangle)\cdot\iota(\Upsilon)$ for $\Upsilon\in\End_z(\Lambda)$,
in the place of the definition of $\End_z(\Lambda)$,
shows that $U_{\mathrm{G}}(a)$ commutes with every operator $\exp(i\langle\delta,A\rangle)\cdot\iota\left(\Upsilon(E_S)\right)$, hence with $X_\Xi$.
The projection $P_{\mathrm{G}}$ is an average of the operators $U_{\mathrm{G}}(a)$ (Definition~\ref{def:slice-operators}),
and so it commutes with $X_\Xi$ as well.
For a link $b\in\Lambda^{(1)}_{\mathrm{s}}\setminus S$ the translations of the variable $A(b)$ commute with the multiplication operators $\exp(i\langle\delta,A\rangle)$,
which do not depend on $A(b)$, with the functions of $E_S$ and with the operators on $\mathcal{F}$;
so $X_\Xi$ commutes with every bounded Borel function of $E_b$.
Neither the plaquette $\sigma_0$ nor the flux index of the direct sum $\mathcal{H}^{\mathrm{l}}$ enters the definition of $X_\Xi$.
\end{proof}

\begin{lemma}[The norm]\label{lem:sharp-norm}
Let $S$ be a finite set of links and $\Lambda$ a non-empty finite set of sites of the slice lattice $\Z^2$,
such that $\Lambda$ contains the end points of the links of $S$,
and let $\Xi$ be a rectangular cross-section in which $\hull(S,\Lambda)$ embeds.
For every $X\in\mathfrak{W}(S,\Lambda)$
\begin{equation}\label{eq:amplification-norm-supremum}
\sup_{\beta\in\R^S}\left\|(\alpha_\beta X)_\Xi\right\|=\|X\|,
\end{equation}
where $\|X\|$ is the operator norm of $X$ on $\mathcal{H}_S$.
Consequently $\|X\|_\#=\|X\|$, and $\|\cdot\|_\#$ is a $C^*$-norm on $\mathfrak{W}(S,\Lambda)$ with
\[
\|1\|_\#=1,\qquad\|X^*\|_\#=\|X\|_\#,\qquad\|\alpha_\beta X\|_\#=\|X\|_\#\quad(\beta\in\R^S).
\]
For every rectangular cross-section $\Xi'$ in which $\hull(S,\Lambda)$ embeds and every $\beta\in\R^S$,
$\left\|(\alpha_\beta X)_{\Xi'}\right\|\le\|X\|_\#$.
For $\delta\in\mathfrak{D}(S)$, $\Upsilon\in\End_{z(\delta)}(\Lambda)$ and $f\in\mathfrak{A}(S)$,
\[
\left\|\exp(i\langle\delta,A\rangle)\cdot\Upsilon\cdot f(E_S)\right\|_\#=\|\Upsilon\|\cdot\sup|f| .
\]
Finally, $\|X\|_\#$ does not change under the inclusions $\mathfrak{W}(S,\Lambda)\subset\mathfrak{W}(S_1,\Lambda_1)$, $S\subset S_1$, $\Lambda\subset\Lambda_1$,
of Definition~\ref{def:mid-plane}.
\end{lemma}

\begin{proof}
\emph{Step~1: removal of the localisation.}
For $X$ with a normal form as in Lemma~\ref{lem:normal-form}(b) let
\[
X^{0}_{\Xi}:=\sum_{\delta\in\mathfrak{D}_X}\exp(i\langle\delta,A\rangle)\cdot\Upsilon_\delta(E_S),
\]
an operator on $L^2\left(\R^{\Lambda^{(1)}_{\mathrm{s}}}/\Gamma_{\mathrm{s}}\right)\otimes\mathcal{F}_\Lambda$,
well defined by Lemma~\ref{lem:normal-form}(c) and~(f).
On each summand of $\mathcal{H}^{\mathrm{l}}$ one has $X_\Xi=(\mathrm{id}\otimes\iota)(X^{0}_{\Xi})$,
where $\mathrm{id}\otimes\iota$ applies $\iota$ to the Fock factor.
As $\mathcal{F}_\Lambda$ and $\mathcal{F}$ are finite-dimensional,
$\mathrm{id}\otimes\iota$ is defined on the $C^*$-algebra of all bounded operators on $L^2\left(\R^{\Lambda^{(1)}_{\mathrm{s}}}/\Gamma_{\mathrm{s}}\right)\otimes\mathcal{F}_\Lambda$,
and it is a $*$-homomorphism into the bounded operators on $L^2\left(\R^{\Lambda^{(1)}_{\mathrm{s}}}/\Gamma_{\mathrm{s}}\right)\otimes\mathcal{F}$;
it is injective because $\iota$ is (Proposition~\ref{prop:localisation}\eqref{it:localisation-homomorphism}).
An injective $*$-homomorphism of $C^*$-algebras is isometric.
Therefore $\|X_\Xi\|=\|X^{0}_{\Xi}\|$.

\emph{Step~2: a unitary equivalence.}
Let
\[
\widetilde{\mathcal{H}}:=L^2(\R^S,\mathrm{d}z)\otimes L^2\left(\R^{\Lambda^{(1)}_{\mathrm{s}}}/\Gamma_{\mathrm{s}}\right)\otimes\mathcal{F}_\Lambda,
\]
whose elements are written $\Phi'(z,[A])$, with $[A]$ the class of $A\in\R^{\Lambda^{(1)}_{\mathrm{s}}}$,
and let $E'$ be the family of the operators $-i\,\partial/\partial z_b$, $b\in S$, on the first factor.
The operator
\[
X(z,E'):=\sum_{\delta\in\mathfrak{D}_X}\exp(i\langle\delta,z\rangle)\cdot\Upsilon_\delta(E')
\]
is $X$ acting on the first and the third factor, with $z$ in the place of $A$;
it is unitarily equivalent to $X\otimes1$ on $\mathcal{H}_S\otimes L^2\left(\R^{\Lambda^{(1)}_{\mathrm{s}}}/\Gamma_{\mathrm{s}}\right)$, and so $\|X(z,E')\|=\|X\|$.
Choose a Borel section $[A]\mapsto A^{\flat}$ of the quotient map $\R^{\Lambda^{(1)}_{\mathrm{s}}}\to\R^{\Lambda^{(1)}_{\mathrm{s}}}/\Gamma_{\mathrm{s}}$, and put
\[
(U_1\Phi')(z,[A]):=\Phi'\left(z+A^{\flat}|_S,[A]\right),\qquad(U_2\Phi')(z,[A]):=\Phi'\left(z,[A-z]\right),
\]
where $z$ is extended by zero outside $S$.
For fixed $[A]$, resp.\ for fixed $z$, these are translations; hence $U_1$ and $U_2$ are unitary.
The operator $U_1$ commutes with the translations $\exp(-i\langle x,E'\rangle)$, $x\in\R^S$, of the variable $z$,
hence with every bounded Borel function of $E'$; and
\[
U_1\exp(i\langle\delta,z\rangle)\,U_1^{-1}=\exp\left(i\langle\delta,z+A^{\flat}|_S\rangle\right)=\exp(i\langle\delta,z\rangle)\cdot\exp(i\langle\delta,A\rangle),
\]
because $\delta$ is supported in $S$ and $\exp(i\langle\delta,\cdot\rangle)$ is a function on $\R^{\Lambda^{(1)}_{\mathrm{s}}}/\Gamma_{\mathrm{s}}$
(Lemma~\ref{lem:normal-form}(f)).
For $x\in\R^S$, next,
\[
\left(U_2\exp(-i\langle x,E'\rangle)\,U_2^{-1}\Phi'\right)(z,[A])=\left(U_2^{-1}\Phi'\right)(z-x,[A-z])=\Phi'(z-x,[A-x]),
\]
that is, $U_2\exp(-i\langle x,E'\rangle)\,U_2^{-1}=\exp(-i\langle x,E'+E_S\rangle)$, with $E_S$ acting on the second factor.
These unitary groups determine the joint spectral calculus of the commuting families $E'$ and $E'+E_S$;
so $U_2f(E')\,U_2^{-1}=f(E'+E_S)$ for every bounded Borel function $f$ on $\R^S$, in particular for $f\in\mathfrak{A}(S)$.
Moreover
\[
U_2\exp(i\langle\delta,z\rangle)\cdot\exp(i\langle\delta,A\rangle)\,U_2^{-1}=\exp(i\langle\delta,z\rangle)\cdot\exp(i\langle\delta,A-z\rangle)=\exp(i\langle\delta,A\rangle).
\]
Both unitaries commute with the operators on the third factor.
With $U:=U_2U_1$ it follows that
\begin{equation}\label{eq:amplification-equivalence}
U\,X(z,E')\,U^{-1}=\sum_{\delta\in\mathfrak{D}_X}\exp(i\langle\delta,A\rangle)\cdot\Upsilon_\delta(E_S+E').
\end{equation}

\emph{Step~3: decomposition.}
Under the Fourier transform in the variable $z$, the family $E'$ becomes the multiplication by $p\in\R^S$,
and $\widetilde{\mathcal{H}}$ becomes $L^2\left(\R^S,\mathrm{d}p;L^2\left(\R^{\Lambda^{(1)}_{\mathrm{s}}}/\Gamma_{\mathrm{s}}\right)\otimes\mathcal{F}_\Lambda\right)$.
The right side of~\eqref{eq:amplification-equivalence} becomes the operator $\Phi'\mapsto\left(T_p\Phi'(p)\right)_p$, where, by Lemma~\ref{lem:normal-form}(e),
\[
T_p:=\sum_{\delta\in\mathfrak{D}_X}\exp(i\langle\delta,A\rangle)\cdot\Upsilon_\delta(E_S+p)=(\alpha_{-p}X)^{0}_{\Xi}.
\]
The symbols are bounded and uniformly continuous (Lemma~\ref{lem:normal-form}(a)).
Hence the family $(T_p)_p$ is bounded, and it is norm continuous:
\[
\left\|T_p-T_{p'}\right\|\le\sum_{\delta\in\mathfrak{D}_X}\sup_{\xi\in\R^S}\left\|\Upsilon_\delta(\xi+p)-\Upsilon_\delta(\xi+p')\right\|,
\]
and the right side tends to zero as $p'\to p$.
For such a family the operator $\Phi'\mapsto\left(T_p\Phi'(p)\right)_p$ has the norm $\sup_p\|T_p\|$.
Its norm is at most this supremum.
Conversely, let $p_0\in\R^S$ and $\varepsilon>0$, let $v$ be a unit vector with $\|T_{p_0}v\|\ge\|T_{p_0}\|-\varepsilon$,
and let $\Phi'$ be the product of $v$ with the normalised indicator function of a ball about $p_0$ on which $\|T_p-T_{p_0}\|\le\varepsilon$.
On this ball $\|T_pv\|\ge\|T_{p_0}\|-2\varepsilon$, and so the image of $\Phi'$ has a norm at least $\|T_{p_0}\|-2\varepsilon$.
Collecting the three steps, with $U$ unitary,
\[
\|X\|=\left\|X(z,E')\right\|=\sup_{p\in\R^S}\left\|(\alpha_{-p}X)^{0}_{\Xi}\right\|=\sup_{p\in\R^S}\left\|(\alpha_{-p}X)_\Xi\right\|,
\]
where Step~1 has been applied to the elements $\alpha_{-p}X$ of $\mathfrak{W}(S,\Lambda)$ (Lemma~\ref{lem:normal-form}(e)).
This is~\eqref{eq:amplification-norm-supremum}.

\emph{Consequences.}
By Definition~\ref{def:mid-plane}, $\|X\|_\#$ is the supremum of $\left\|(\alpha_\beta X)_{\Xi'}\right\|$ over all $\beta\in\R^S$
and all rectangular cross-sections $\Xi'$ admitted there, in each of which $\hull(S,\Lambda)$ embeds.
Such cross-sections exist.
Write $\hull(\Lambda)=[y^-_2,y^+_2]\times[y^-_3,y^+_3]$ as in Definition~\ref{def:slices},
and let $n_2,n_3\ge4$ be integers with $n_\nu\ge y^+_\nu-y^-_\nu+3$ for $\nu=2,3$.
Then the translates of $\hull(\Lambda)+[-1,1]^2$ by the elements of $n_2\Z\times n_3\Z$ are pairwise disjoint,
so that $\hull(S,\Lambda)$ embeds in the rectangular cross-section with the sides $n_2,n_3$;
moreover, at most $(y^+_2-y^-_2)\cdot(y^+_3-y^-_3)<n_2\cdot n_3$ plaquettes of its slice torus are contained in $\hull(S,\Lambda)$,
and the plaquette $\sigma_0$ can be chosen among the others,
a choice on which the operators $(\alpha_\beta X)_{\Xi'}$ do not depend (Lemma~\ref{lem:normal-form}(f)).
By~\eqref{eq:amplification-norm-supremum}, applied to each admitted $\Xi'$, the supremum over $\beta$ equals $\|X\|$; hence $\|X\|_\#=\|X\|$.
Applied to an arbitrary rectangular cross-section $\Xi'$ in which $\hull(S,\Lambda)$ embeds,
\eqref{eq:amplification-norm-supremum} also gives $\left\|(\alpha_\beta X)_{\Xi'}\right\|\le\|X\|=\|X\|_\#$.
The operator norm, restricted to the unital $*$-algebra $\mathfrak{W}(S,\Lambda)$ of operators on $\mathcal{H}_S$,
is a $C^*$-norm with $\|1\|=1$ and $\|X^*\|=\|X\|$;
and $\|\alpha_\beta X\|=\|X\|$, since $\alpha_\beta$ is the conjugation by a unitary operator (Lemma~\ref{lem:normal-form}(e)).
For a single term, $\exp(i\langle\delta,A\rangle)$ is unitary,
and $\Upsilon\cdot f(E_S)=f(E_S)\otimes\Upsilon$ has the norm $\|\Upsilon\|\cdot\sup|f|$,
because $f$ is continuous and the spectrum of $E_S$ on $L^2(\R^S)$ is $\R^S$.
Let finally $S\subset S_1$ and $\Lambda\subset\Lambda_1$, where $\Lambda_1$ is finite and contains the end points of the links of $S_1$.
The inclusion of $\mathfrak{W}(S,\Lambda)$ into $\mathfrak{W}(S_1,\Lambda_1)$ is the map $X\mapsto X\otimes1$, with $1$ the identity of $L^2(\R^{S_1\setminus S})$,
followed by the application to the Fock factor of the localisation $\iota_{\Lambda\to\Lambda_1}:\End(\mathcal{F}_\Lambda)\to\End(\mathcal{F}_{\Lambda_1})$ of Definition~\ref{def:modes}.
Both maps preserve operator norms, the second as an injective $*$-homomorphism of $C^*$-algebras;
so $\|X\|_\#$ does not change.

These inclusions are compatible with one another:
for finite sets $\Lambda\subset\Lambda_1\subset\Lambda_2$ of sites, $\iota_{\Lambda_1\to\Lambda_2}\circ\iota_{\Lambda\to\Lambda_1}=\iota_{\Lambda\to\Lambda_2}$.
To see this, let $I$, $J$ be sets of modes over $\Lambda$, let $R$ be a set of modes over $\Lambda_1$ none of which lies over $\Lambda$,
and let $R'$ be a set of modes over $\Lambda_2$ none of which lies over $\Lambda_1$.
By Definition~\ref{def:modes}, the entries of $\iota_{\Lambda_1\to\Lambda_2}\left(\iota_{\Lambda\to\Lambda_1}(\Upsilon)\right)$ vanish
unless the row and the column index have the forms $I\cup R\cup R'$ and $J\cup R\cup R'$, with the same $R$ and the same $R'$.
The entry at such a place is $\Upsilon_{IJ}$ times the power of $-1$ whose exponent is the sum of four numbers,
namely the numbers of the pairs $(i,r)$ with $i>r$ in the sets $(I\cup R)\times R'$, $(J\cup R)\times R'$, $I\times R$ and $J\times R$.
After $I\cup R$ and $J\cup R$ have been split, the pairs in $R\times R'$ are counted twice,
and the remaining ones are the pairs with $i>r$ in $I\times(R\cup R')$ and in $J\times(R\cup R')$.
The entry of $\iota_{\Lambda\to\Lambda_2}(\Upsilon)$ at the same place is $\Upsilon_{IJ}$ times the power of $-1$ whose exponent is the number of these remaining pairs,
and the entries of $\iota_{\Lambda\to\Lambda_2}(\Upsilon)$ at all other places vanish.
\end{proof}

\begin{remark}\label{rem:mid-plane-completion}
By~\eqref{eq:amplification-norm-supremum}, one rectangular cross-section $\Xi$ in which $\hull(S,\Lambda)$ embeds already gives the value of $\|X\|_\#$:
the supremum over the cross-sections in Definition~\ref{def:mid-plane} may be omitted.
Since $\|\cdot\|_\#$ is a $C^*$-norm that does not change under the inclusions (Lemma~\ref{lem:sharp-norm}),
and the inclusions are compatible with one another,
the union $\mathfrak{W}$ of the algebras $\mathfrak{W}(S,\Lambda)$ is a normed $*$-algebra, and its completion is a $C^*$-algebra.
The supremum over $\beta$ behaves differently.
On $L^2\left(\R^{\Lambda^{(1)}_{\mathrm{s}}}/\Gamma_{\mathrm{s}}\right)$ the operator $\exp(-i\langle\gamma,E\rangle)$ is the translation by $\gamma$,
which is the identity for $\gamma\in\Gamma_{\mathrm{s}}$;
so the joint spectrum of $E$ lies in the set of the $\xi$ with $\langle\gamma,\xi\rangle\in2\pi\Z$ for all $\gamma\in\Gamma_{\mathrm{s}}$.
For $\gamma=(2\pi/q)\cdot d_{\mathrm{s}}a$, with $a$ the indicator function of a site $y$, this condition reads $(d_{\mathrm{s}}^*\xi)(y)\in q\Z$.
If $S$ contains the four links at $y$, the condition concerns the components of $\xi$ on $S$ only,
and the spectrum of $E_S$ on this space is a proper closed subset of $\R^S$.
For $f\in\mathfrak{A}(S)$ the norm of $\left(f(E_S)\right)_\Xi$ is then the supremum of $|f|$ over this subset,
whereas $\|f(E_S)\|_\#=\sup|f|$.
For real $\lambda$ the function $f_{S,\mathfrak{Q},\lambda}$ is positive, tends to zero at infinity and has a single maximum point,
which the choice of $\lambda$ places anywhere in $\R^S$;
when this point lies outside the subset, the first number is strictly smaller than the second.
Hence the supremum over $\beta$ cannot be omitted.
\end{remark}

\subsection{Step elements}\label{sec:step-elements}

\subsubsection{Kernels}\label{sec:step-kernels}

For a finite set $S$ of links of the slice lattice $\Z^2$
and a non-empty finite set $\Lambda$ of sites of $\Z^2$ that contains the end points of the links of $S$,
the Hilbert space $\mathcal{H}_S=L^2(\R^S)\otimes\mathcal{F}_\Lambda$
and the mid-plane algebra $\mathfrak{W}(S,\Lambda)$ of operators on $\mathcal{H}_S$ are those of Definition~\ref{def:mid-plane}.
The spectral variable of $E_S=(E_b)_{b\in S}$ is $\xi\in\R^S$;
for $x\in\R^S$ put $|x|^2:=\sum_{b\in S}x(b)^2$, and $|E_S|^2:=\sum_{b\in S}E_b^2$.
Kernels of operators on $\mathcal{H}_S$, and the inverse Fourier transform $f^{\vee}$ of an integrable function $f$ on $\R^S$,
are understood as in Definition~\ref{def:integral-kernels}.
For $a\in\R^\Lambda$, the coboundary $d_{\mathrm{s}}a$ is defined on the links of $S$, whose end points lie in $\Lambda$;
its restriction to $S$ is again denoted by $d_{\mathrm{s}}a\in\R^S$.

\begin{lemma}[Kernels of regularised elements]\label{lem:regularised-kernels}
Let $S$ and $\Lambda$ be as in Definition~\ref{def:mid-plane}, let $0<s\le\frac14$, and let
\[
X=\sum_{\delta\in\mathfrak{D}_X}\exp\bigl(i\langle\delta,A\rangle\bigr)\cdot\Upsilon_\delta(E_S)
\]
be an element of $\mathfrak{W}(S,\Lambda)$ in its normal form (Lemma~\ref{lem:normal-form}(b)).
For $\delta\in\mathfrak{D}_X$ let $\Upsilon^{(s)}_\delta\colon\R^S\to\End\mathcal{F}_\Lambda$ be the function
\[
\Upsilon^{(s)}_\delta(\xi):=\exp\left(-s|\xi+\delta|^2\right)\cdot\Upsilon_\delta(\xi)\cdot\exp\left(-s|\xi|^2\right).
\]
Then the operator $\exp\left(-s|E_S|^2\right)\cdot X\cdot\exp\left(-s|E_S|^2\right)$ on $\mathcal{H}_S$
has the continuous kernel $k^{(s)}_X\colon\R^S\times\R^S\to\End\mathcal{F}_\Lambda$,
\[
k^{(s)}_X(A',A)=\sum_{\delta\in\mathfrak{D}_X}\exp\bigl(i\langle\delta,A'\rangle\bigr)\cdot\left(\Upsilon^{(s)}_\delta\right)^{\vee}(A'-A),
\]
and the following hold.
\begin{enumerate}
\item[(a)]
The map $X\mapsto k^{(s)}_X$ is linear;
the values of $k^{(s)}_X$ are even elements of $\End\mathcal{F}_\Lambda$;
and $k^{(s)}_{X^*}(A',A)=k^{(s)}_X(A,A')^*$ for all $A',A\in\R^S$.
\item[(b)]
\emph{Covariance.}
For every $a\in\R^\Lambda$ and all $A',A\in\R^S$,
\[
u(a)^{-1}\cdot k^{(s)}_X\bigl(A'+d_{\mathrm{s}}a,\,A+d_{\mathrm{s}}a\bigr)\cdot u(a)=k^{(s)}_X(A',A).
\]
\item[(c)]
\emph{Explicit form; Gaussian bound.}
Write $\langle\lambda',\lambda''\rangle:=\sum_{b\in S}\lambda'(b)\cdot\lambda''(b)$ for the bilinear pairing of $\C^S$.
Let $\delta\in\mathfrak{D}_X$, and suppose that $\Upsilon_\delta=f_{S',\mathfrak{Q},\lambda}\cdot\Upsilon$,
with $f_{S',\mathfrak{Q},\lambda}$ as in Definition~\ref{def:gaussian-symbols} and $\Upsilon\in\End_{z(\delta)}(\Lambda)$.
Regard $\mathfrak{Q}$ as the symmetric matrix on $\R^S$ with $\mathfrak{Q}(\xi_{S'})=\langle\xi,\mathfrak{Q}\xi\rangle$, extended by $0$ on $S\setminus S'$;
extend $\lambda$ by $0$ to an element of $\C^S$; and put
\[
\mathfrak{Q}_s:=4s+2\mathfrak{Q},\qquad\lambda_s:=\lambda-2s\delta.
\]
Then for every $x\in\R^S$
\begin{align*}
\left(\Upsilon^{(s)}_\delta\right)^{\vee}(x)
&=(2\pi)^{-\frac{\#|S|}{2}}\cdot(\det\mathfrak{Q}_s)^{-\frac12}\\
&\qquad\cdot\exp\left(-s|\delta|^2+\tfrac12\left\langle\lambda_s+ix,\mathfrak{Q}_s^{-1}(\lambda_s+ix)\right\rangle\right)\cdot\Upsilon,\\
\left\|\left(\Upsilon^{(s)}_\delta\right)^{\vee}(x)\right\|
&=(2\pi)^{-\frac{\#|S|}{2}}\cdot(\det\mathfrak{Q}_s)^{-\frac12}\cdot\|\Upsilon\|\\
&\qquad\cdot\exp\left(-s|\delta|^2+\tfrac12\left\langle\Re\lambda_s,\mathfrak{Q}_s^{-1}\Re\lambda_s\right\rangle-\tfrac12\left\langle x+\Im\lambda,\mathfrak{Q}_s^{-1}(x+\Im\lambda)\right\rangle\right).
\end{align*}
Consequently there are constants $C<\infty$ and $\varepsilon_k>0$ such that for all $A',A\in\R^S$
\[
\left\|k^{(s)}_X(A',A)\right\|\le C\prod_{b\in S}\exp\left(-\varepsilon_k\bigl(A'(b)-A(b)\bigr)^2\right).
\]
\item[(d)]
\emph{Characters.}
For every $\beta\in\R^S$ and all $A',A\in\R^S$,
\[
\exp\bigl(i\langle\beta,A'-A\rangle\bigr)\cdot k^{(s)}_X(A',A)
=\sum_{\delta\in\mathfrak{D}_X}\exp\bigl(i\langle\delta,A'\rangle\bigr)\cdot\left(\Upsilon^{(s)}_\delta(\cdot-\beta)\right)^{\vee}(A'-A),
\]
and this is the kernel of $\exp\left(-s|E_S-\beta|^2\right)\cdot(\alpha_\beta X)\cdot\exp\left(-s|E_S-\beta|^2\right)$.
\item[(e)]
For every $\beta\in\R^S$ the triples $\left(S,\Lambda,k^{(s)}_X\right)$ and $\left(S,\Lambda,\exp(i\langle\beta,A'-A\rangle)\cdot k^{(s)}_X\right)$
are local step data in the sense of Proposition~\ref{prop:word-formula}.
\end{enumerate}
\end{lemma}

\begin{proof}
Fix $\delta\in\mathfrak{D}_X$.
On $L^2(\R^S)$ one has $\exp(-i\langle\delta,A\rangle)\cdot E_b\cdot\exp(i\langle\delta,A\rangle)=E_b+\delta(b)$ for $b\in S$,
so the relation \eqref{eq:setting-weyl}, $f(E_S)\cdot\exp(i\langle\delta,A\rangle)=\exp(i\langle\delta,A\rangle)\cdot f(E_S+\delta)$,
holds there for every bounded Borel function $f$ on $\R^S$.
Applied to the function $\xi\mapsto\exp\left(-s|\xi|^2\right)$ it gives
\[
\exp\left(-s|E_S|^2\right)\cdot\exp\bigl(i\langle\delta,A\rangle\bigr)\cdot\Upsilon_\delta(E_S)\cdot\exp\left(-s|E_S|^2\right)
=\exp\bigl(i\langle\delta,A\rangle\bigr)\cdot\Upsilon^{(s)}_\delta(E_S).
\]
The function $\Upsilon^{(s)}_\delta$ is a Schwartz function on $\R^S$ with values in $\End\mathcal{F}_\Lambda$.
Indeed, $\Upsilon_\delta$ is a finite sum of products of elements of $\mathfrak{A}(S)$ with constant operators,
and every element of $\mathfrak{A}(S)$ is a linear combination of $1$ and of functions $f_{S',\mathfrak{Q},\lambda}$;
as a function of $\xi_{S'}$, such an $f_{S',\mathfrak{Q},\lambda}$ is a Schwartz function (the form $\mathfrak{Q}$ is positive definite),
and it does not depend on the other coordinates of $\xi$.
So $\Upsilon_\delta$ is smooth and bounded together with all its derivatives (compare Lemma~\ref{lem:normal-form}(a)),
and its product with the Gaussian function $\exp\left(-s|\xi+\delta|^2-s|\xi|^2\right)$,
which decreases in every direction of $\R^S$, is a Schwartz function.
The inverse Fourier transform $\left(\Upsilon^{(s)}_\delta\right)^{\vee}$ is then a Schwartz function as well,
in particular continuous and integrable,
and by Definition~\ref{def:integral-kernels} the operator $\exp(i\langle\delta,A\rangle)\cdot\Upsilon^{(s)}_\delta(E_S)$ has the kernel
$\exp(i\langle\delta,A'\rangle)\cdot\left(\Upsilon^{(s)}_\delta\right)^{\vee}(A'-A)$.
Summation over the finite set $\mathfrak{D}_X$ gives the formula for $k^{(s)}_X$.
This kernel is continuous, and $\left\|k^{(s)}_X(A',A)\right\|$ is bounded by the integrable function
$\sum_{\delta\in\mathfrak{D}_X}\left\|\left(\Upsilon^{(s)}_\delta\right)^{\vee}(A'-A)\right\|$ of $A'-A$.

Write $T:=\exp\left(-s|E_S|^2\right)\cdot X\cdot\exp\left(-s|E_S|^2\right)$, a bounded operator on $\mathcal{H}_S$.
A continuous kernel of a bounded operator on $\mathcal{H}_S$ is determined by the operator at every point:
two such kernels agree almost everywhere on $\R^S\times\R^S$, and therefore everywhere.

(a)
The operator $T$ depends linearly on $X$, and therefore so does its continuous kernel.
By Lemma~\ref{lem:normal-form}(d) the operator $X$ is even, that is, it commutes with $1\otimes U_{\mathrm{f}}$;
so does $T$, because $\exp\left(-s|E_S|^2\right)$ acts on the factor $L^2(\R^S)$ only.
The operator $(1\otimes U_{\mathrm{f}})\cdot T\cdot(1\otimes U_{\mathrm{f}})$ has the continuous kernel $U_{\mathrm{f}}\cdot k^{(s)}_X(A',A)\cdot U_{\mathrm{f}}$;
hence $U_{\mathrm{f}}\cdot k^{(s)}_X(A',A)\cdot U_{\mathrm{f}}=k^{(s)}_X(A',A)$, and the values of $k^{(s)}_X$ are even.
The algebra $\mathfrak{W}(S,\Lambda)$ contains $X^*$,
and $\exp\left(-s|E_S|^2\right)\cdot X^*\cdot\exp\left(-s|E_S|^2\right)=T^*$, since $\exp\left(-s|E_S|^2\right)$ is self-adjoint.
If the norm of the kernel $k$ of an operator is bounded by an integrable function of $A'-A$,
Fubini's theorem shows that the adjoint operator has the kernel $(A',A)\mapsto k(A,A')^*$.
With the majorant found above this gives $k^{(s)}_{X^*}(A',A)=k^{(s)}_X(A,A')^*$.

(b)
Let $a\in\R^\Lambda$.
By Lemma~\ref{lem:normal-form}(d) the operator $X$ commutes with $U^{S}_{\mathrm{G}}(a)$,
where $\bigl(U^{S}_{\mathrm{G}}(a)\Phi\bigr)(A)=u(a)\cdot\Phi(A-d_{\mathrm{s}}a)$.
The operator $\exp\left(-s|E_S|^2\right)$ commutes with $U^{S}_{\mathrm{G}}(a)$ as well,
because functions of $E_S$ commute with the translations of the variable $A$, and $u(a)$ acts on $\mathcal{F}_\Lambda$ only.
Hence $T$ commutes with $U^{S}_{\mathrm{G}}(a)$.
For $\Phi\in\mathcal{H}_S$ and $k:=k^{(s)}_X$,
\begin{align*}
\bigl(U^{S}_{\mathrm{G}}(a)\cdot T\Phi\bigr)(A')&=\int\mathrm{d}A\,u(a)\cdot k(A'-d_{\mathrm{s}}a,A)\cdot\Phi(A),\\
\bigl(T\cdot U^{S}_{\mathrm{G}}(a)\Phi\bigr)(A')&=\int\mathrm{d}A\,k(A',A)\cdot u(a)\cdot\Phi(A-d_{\mathrm{s}}a)
=\int\mathrm{d}A\,k(A',A+d_{\mathrm{s}}a)\cdot u(a)\cdot\Phi(A).
\end{align*}
Both kernels are continuous, so $u(a)\cdot k(A'-d_{\mathrm{s}}a,A)=k(A',A+d_{\mathrm{s}}a)\cdot u(a)$ for all $A',A\in\R^S$;
replacing $A'$ by $A'+d_{\mathrm{s}}a$ gives (b).

(c)
Suppose first that $\Upsilon_\delta=f_{S',\mathfrak{Q},\lambda}\cdot\Upsilon$ as in (c).
With the extended $\mathfrak{Q}$ and $\lambda$ one has
$f_{S',\mathfrak{Q},\lambda}(\xi)=\exp\left(-\langle\xi,\mathfrak{Q}\xi\rangle+\langle\lambda,\xi\rangle\right)$ for $\xi\in\R^S$, and
\[
-s|\xi+\delta|^2-s|\xi|^2-\langle\xi,\mathfrak{Q}\xi\rangle+\langle\lambda,\xi\rangle
=-\tfrac12\langle\xi,\mathfrak{Q}_s\xi\rangle+\langle\lambda_s,\xi\rangle-s|\delta|^2;
\]
hence
\[
\Upsilon^{(s)}_\delta(\xi)=\exp\left(-\tfrac12\langle\xi,\mathfrak{Q}_s\xi\rangle+\langle\lambda_s,\xi\rangle-s|\delta|^2\right)\cdot\Upsilon.
\]
The matrix $\mathfrak{Q}_s$ is real, symmetric and positive definite, because $s>0$ and $\mathfrak{Q}\ge0$.
For every $\lambda'\in\C^S$ one has
\[
\int_{\R^S}\mathrm{d}\xi\,\exp\left(-\tfrac12\langle\xi,\mathfrak{Q}_s\xi\rangle+\langle\lambda',\xi\rangle\right)
=(2\pi)^{\frac{\#|S|}{2}}\cdot(\det\mathfrak{Q}_s)^{-\frac12}\cdot\exp\left(\tfrac12\left\langle\lambda',\mathfrak{Q}_s^{-1}\lambda'\right\rangle\right).
\]
For real $\lambda'$ this follows from the completed square
\[
-\tfrac12\langle\xi,\mathfrak{Q}_s\xi\rangle+\langle\lambda',\xi\rangle
=-\tfrac12\left\langle\xi-\mathfrak{Q}_s^{-1}\lambda',\mathfrak{Q}_s\left(\xi-\mathfrak{Q}_s^{-1}\lambda'\right)\right\rangle
+\tfrac12\left\langle\lambda',\mathfrak{Q}_s^{-1}\lambda'\right\rangle,
\]
from the translation $\xi\mapsto\xi+\mathfrak{Q}_s^{-1}\lambda'$,
and from $\int_{\R^S}\mathrm{d}\xi\,\exp\left(-\tfrac12\langle\xi,\mathfrak{Q}_s\xi\rangle\right)=(2\pi)^{\frac{\#|S|}{2}}\cdot(\det\mathfrak{Q}_s)^{-\frac12}$.
For complex $\lambda'$ put $|\lambda'|^2:=\sum_{b\in S}|\lambda'(b)|^2$.
The right side of the integral formula is a constant multiple of the exponential of a quadratic polynomial in the coordinates $\lambda'(b)$, $b\in S$;
as a function of any one of these coordinates, the others being fixed, it is therefore holomorphic on $\C$.
The left side has the same property.
To see this, let $r>0$ and $|\lambda'|\le r$.
The Cauchy--Schwarz inequality gives $\left|\exp\left(\langle\lambda',\xi\rangle\right)\right|\le\exp\left(r|\xi|\right)$,
and the derivative of the integrand with respect to $\lambda'(b)$ is the integrand multiplied by $\xi(b)$;
so the integrand and each of these derivatives have modulus at most
\[
(1+|\xi|)\cdot\exp\left(-\tfrac12\langle\xi,\mathfrak{Q}_s\xi\rangle+r|\xi|\right),
\]
a function of $\xi$ alone, integrable over $\R^S$ because $\langle\xi,\mathfrak{Q}_s\xi\rangle\ge4s|\xi|^2$.
With this majorant, the theorem on differentiation of integrals depending on a parameter
(dominated convergence, applied to difference quotients)
shows that the left side is complex differentiable with respect to each coordinate $\lambda'(b)$ at every point of $\{|\lambda'|<r\}$;
and $r>0$ is arbitrary.
On $\R^S$ the two sides agree.
Induction on the number of coordinates of $\lambda'$ that are not real extends the agreement to $\C^S$:
given $\lambda'$ with $\lambda'(b)\notin\R$ for some $b\in S$, keep the other coordinates of $\lambda'$ fixed;
as functions of $\lambda'(b)\in\C$, the two sides are holomorphic and, by the induction hypothesis, equal on $\R$,
hence equal on $\C$ by the identity theorem.
This proves the integral formula for all $\lambda'\in\C^S$.
By the definition of the inverse Fourier transform,
\[
\left(\Upsilon^{(s)}_\delta\right)^{\vee}(x)
=(2\pi)^{-\#|S|}\cdot\exp\left(-s|\delta|^2\right)\cdot\int_{\R^S}\mathrm{d}\xi\,\exp\left(-\tfrac12\langle\xi,\mathfrak{Q}_s\xi\rangle+\langle\lambda_s+ix,\xi\rangle\right)\cdot\Upsilon,
\]
and the integral formula with $\lambda'=\lambda_s+ix$ gives the first formula of (c).
The modulus of the exponential of a complex number is the exponential of its real part, and,
$\mathfrak{Q}_s^{-1}$ being real and symmetric,
\[
\Re\left\langle\lambda',\mathfrak{Q}_s^{-1}\lambda'\right\rangle
=\left\langle\Re\lambda',\mathfrak{Q}_s^{-1}\Re\lambda'\right\rangle-\left\langle\Im\lambda',\mathfrak{Q}_s^{-1}\Im\lambda'\right\rangle;
\]
for $\lambda'=\lambda_s+ix$ one has $\Re\lambda'=\Re\lambda_s$ and $\Im\lambda'=x+\Im\lambda$, since $\delta$ is real.
This is the second formula.

Let $\|\mathfrak{Q}\|$ be the largest eigenvalue of the extended matrix $\mathfrak{Q}$.
Then $\mathfrak{Q}_s\le(4s+2\|\mathfrak{Q}\|)\cdot1$, hence $\mathfrak{Q}_s^{-1}\ge(4s+2\|\mathfrak{Q}\|)^{-1}\cdot1>0$,
and $|x+\Im\lambda|^2\ge\frac12|x|^2-|\Im\lambda|^2$.
Therefore
\[
-\tfrac12\left\langle x+\Im\lambda,\mathfrak{Q}_s^{-1}(x+\Im\lambda)\right\rangle
\le-\frac{|x|^2}{16s+8\|\mathfrak{Q}\|}+\frac{|\Im\lambda|^2}{8s+4\|\mathfrak{Q}\|},
\]
and the second formula yields
\[
\left\|\left(\Upsilon^{(s)}_\delta\right)^{\vee}(x)\right\|\le C_1\prod_{b\in S}\exp\left(-\frac{x(b)^2}{16s+8\|\mathfrak{Q}\|}\right)
\]
with a constant $C_1$ that does not depend on $x$.
For a general element $X$, each symbol $\Upsilon_\delta$ is a finite sum of products $f_{S',\mathfrak{Q},\lambda}\cdot\Upsilon$ of the kind just treated
(the function $1$ is $f_{S',\mathfrak{Q},\lambda}$ with $S'=\emptyset$),
and $\Upsilon_\delta\mapsto\Upsilon^{(s)}_\delta$ is linear;
so the estimate holds for each of these products.
Since $\mathfrak{D}_X$ is finite and $|\exp(i\langle\delta,A'\rangle)|=1$,
the sum of the finitely many estimates, taken at $x=A'-A$, gives the Gaussian bound for $k^{(s)}_X$,
with $C$ the sum of the constants $C_1$ and $\varepsilon_k$ the least of the numbers $(16s+8\|\mathfrak{Q}\|)^{-1}$ that occur.

(d)
Given $\beta\in\R^S$ and an integrable function $f$ on $\R^S$,
the substitution $\xi\mapsto\xi+\beta$ in the integral that defines $\left(f(\cdot-\beta)\right)^{\vee}(x)$
gives $\left(f(\cdot-\beta)\right)^{\vee}(x)=\exp(i\langle\beta,x\rangle)\cdot f^{\vee}(x)$.
Take $f=\Upsilon^{(s)}_\delta$ and $x=A'-A$, multiply by $\exp(i\langle\delta,A'\rangle)$ and sum over $\delta\in\mathfrak{D}_X$:
this is the identity of (d).
By Lemma~\ref{lem:normal-form}(e), $\alpha_\beta X=\sum_{\delta\in\mathfrak{D}_X}\exp(i\langle\delta,A\rangle)\cdot\Upsilon_\delta(E_S-\beta)$.
The relation \eqref{eq:setting-weyl}, applied to the function $\xi\mapsto\exp\left(-s|\xi-\beta|^2\right)$, gives
\[
\exp\left(-s|E_S-\beta|^2\right)\cdot\exp\bigl(i\langle\delta,A\rangle\bigr)\cdot\Upsilon_\delta(E_S-\beta)\cdot\exp\left(-s|E_S-\beta|^2\right)
=\exp\bigl(i\langle\delta,A\rangle\bigr)\cdot\Upsilon^{(s)}_\delta(E_S-\beta),
\]
because $\Upsilon^{(s)}_\delta(\xi-\beta)=\exp\left(-s|\xi-\beta+\delta|^2\right)\cdot\Upsilon_\delta(\xi-\beta)\cdot\exp\left(-s|\xi-\beta|^2\right)$.
The function $\Upsilon^{(s)}_\delta(\cdot-\beta)$ is a Schwartz function,
and by Definition~\ref{def:integral-kernels} the right side has the kernel
$\exp(i\langle\delta,A'\rangle)\cdot\left(\Upsilon^{(s)}_\delta(\cdot-\beta)\right)^{\vee}(A'-A)$.
Summation over $\delta\in\mathfrak{D}_X$ gives the kernel of $\exp\left(-s|E_S-\beta|^2\right)\cdot(\alpha_\beta X)\cdot\exp\left(-s|E_S-\beta|^2\right)$.

(e)
For a local step datum $(S,\Lambda,k)$ in the sense of Proposition~\ref{prop:word-formula},
the function $k$ is required to be a Borel function on $\R^S\times\R^S$ with values in the even part of $\End\mathcal{F}_\Lambda$,
covariant as in (b), and bounded as in (c) with some $C<\infty$ and $\varepsilon_k>0$.
The kernel $k^{(s)}_X$ is continuous, hence Borel, and it has the three properties by (a), (b) and (c).
The factor $\exp(i\langle\beta,A'-A\rangle)$ is a continuous scalar function of modulus one,
and it does not change under $(A',A)\mapsto(A'+d_{\mathrm{s}}a,\,A+d_{\mathrm{s}}a)$;
so $\exp(i\langle\beta,A'-A\rangle)\cdot k^{(s)}_X$ has the same properties.
\end{proof}

\begin{definition}[Two kinds of insertions]\label{def:two-insertions}
Let $j\in\Z$ and $c=j+\frac12$.
For every local step datum $(S,\Lambda,k)$ the integrand $G[k]$ of Proposition~\ref{prop:word-formula},
formed at the link plane $c$, is denoted by $G^{(c)}[k]$;
it is called the \emph{step element} of the datum at $c$.
Let $S$ and $\Lambda$ be as in Definition~\ref{def:mid-plane}, and let $X\in\mathfrak{W}(S,\Lambda)$.
By Lemma~\ref{lem:regularised-kernels}(e), $\left(S,\Lambda,k^{(\frac14)}_X\right)$ and $\left(S,\Lambda,k^{(\frac18)}_X\right)$ are local step data.
The insertion $G_X=G^{(c)}_X$ of $X$ at $c$ (Definition~\ref{def:insertions}) is formed with the kernel $k_X$ of
$\exp\left(-\tfrac14|E_S|^2\right)\cdot X\cdot\exp\left(-\tfrac14|E_S|^2\right)$,
which is the continuous kernel $k^{(\frac14)}_X$ of Lemma~\ref{lem:regularised-kernels};
and the formula by which Definition~\ref{def:insertions} defines $G_X$ is the formula for $G[k]$ in Proposition~\ref{prop:word-formula},
formed at the link plane $c$, with $k_X$ in the place of $k$.
It follows that
\[
G_X=G^{(c)}\left[k^{(\frac14)}_X\right].
\]
Put
\[
G'_X=G^{\prime(c)}_X:=G^{(c)}\left[k^{(\frac18)}_X\right].
\]
Over $(S,\Lambda)$, the place of the crossing weight
$\prod_{b\in S}g_1\left(F^{\mathrm{el}}_j(b)\right)\cdot\exp\left(-S^{\times}_c(\Lambda)\right)$
is taken in $G'_X$ by the kernel of $K_{0,S}^{1/2}\cdot X\cdot K_{0,S}^{1/2}$ and in $G_X$ by the kernel of $K^{1/2}\cdot X\cdot K^{1/2}$,
where $K$ and $K_{0,S}$ are the operators of Definition~\ref{def:slice-operators}.
Both $G_X$ and $G'_X$ are linear in $X$,
because $X\mapsto k^{(s)}_X$ is linear (Lemma~\ref{lem:regularised-kernels}(a)) and $k\mapsto G^{(c)}[k]$ is linear;
and both lie in $\mathcal{G}^{\mathrm{ev}}\bigl(B(c,S,\Lambda)\bigr)$ by Proposition~\ref{prop:word-formula}\eqref{it:step-class}.
The operator $\exp\left(-\tfrac18|E_S|^2\right)$ belongs to $\mathfrak{W}(S,\Lambda)$:
it is a generator of the kind \eqref{it:gaussian-generators} in Definition~\ref{def:mid-plane},
with $S'=S$, the quadratic form $\frac18|\xi|^2$ and $\lambda=0$.
Hence
\[
\widetilde{X}:=\exp\left(-\tfrac18|E_S|^2\right)\cdot X\cdot\exp\left(-\tfrac18|E_S|^2\right)
\]
lies in $\mathfrak{W}(S,\Lambda)$.
Since $\exp\left(-\tfrac14|E_S|^2\right)\cdot X\cdot\exp\left(-\tfrac14|E_S|^2\right)
=\exp\left(-\tfrac18|E_S|^2\right)\cdot\widetilde{X}\cdot\exp\left(-\tfrac18|E_S|^2\right)$,
one has $k^{(\frac14)}_X=k^{(\frac18)}_{\widetilde{X}}$, and therefore
\[
G_X=G^{(c)}\left[k^{(\frac14)}_X\right]=G^{(c)}\left[k^{(\frac18)}_{\widetilde{X}}\right]=G'_{\widetilde{X}}.
\]
\end{definition}

\subsubsection{The word formula}\label{sec:mixed-words}

Let $\Xi$ be a rectangular cross-section,
with the slice torus $\Lambda_{\mathrm{s}}$, its set of links $\Lambda^{(1)}_{\mathrm{s}}$ and the group $\Gamma_{\mathrm{s}}$ of Definition~\ref{def:slices},
and let $V$ be a finite-dimensional Hilbert space.
For $A\in\R^{\Lambda^{(1)}_{\mathrm{s}}}$ write $[A]$ for its class in $\R^{\Lambda^{(1)}_{\mathrm{s}}}/\Gamma_{\mathrm{s}}$,
and let $\mathrm{d}[A]$ be the measure on this quotient that is induced by Lebesgue measure (Definition~\ref{def:slices}).
Let $k$ be a Borel function on $\R^{\Lambda^{(1)}_{\mathrm{s}}}\times\R^{\Lambda^{(1)}_{\mathrm{s}}}$ with values in $\End V$
such that $k(A'+\gamma,A+\gamma)=k(A',A)$ for all $\gamma\in\Gamma_{\mathrm{s}}$.
The operator with the kernel $k$ periodised over $\Gamma_{\mathrm{s}}$ is the operator $K[k]$ on $L^2(\R^{\Lambda^{(1)}_{\mathrm{s}}}/\Gamma_{\mathrm{s}};V)$ given by
\[
(K[k]\Phi)([A'])
:=\int_{\R^{\Lambda^{(1)}_{\mathrm{s}}}/\Gamma_{\mathrm{s}}}\mathrm{d}[A]\,\left(\sum_{\gamma\in\Gamma_{\mathrm{s}}}k(A',A+\gamma)\right)\Phi([A])
=\int_{\R^{\Lambda^{(1)}_{\mathrm{s}}}}\mathrm{d}A\,k(A',A)\,\Phi([A])
\]
on those $\Phi$ for which the integrals converge absolutely for almost every $[A']$.
The periodicity of $k$ makes the sum over $\gamma$ a function of the classes $[A']$ and $[A]$ alone.
The step operators of the word formula, Proposition~\ref{prop:word-formula}, are operators of this kind, acting in the same way on every summand of $\mathcal{H}^{\mathrm{l}}$.

\begin{lemma}[Words with both kinds of insertions]\label{lem:mixed-word-formula}
Let $\mathbb{T}=C(n_1,n_2,n_3)=(\Gamma,\chi)$ be a coordinate torus with time direction $e_1$,
and let $\Xi=\Xi(\mathbb{T},x_1)$ be its cross-section for the family $x_1$.
By Definition~\ref{def:slice-notation}, $\Xi$ is the rectangular cross-section with the sides $n_2,n_3$:
its lattice is $\Gamma_\parallel=n_2\Z e_2+n_3\Z e_3$, and its character takes the value $-1$ at $n_2e_2$ and at $n_3e_3$.
The slice torus $\Lambda_{\mathrm{s}}$, the sets $\Lambda^{(1)}_{\mathrm{s}}$ and $\Lambda^{(2)}_{\mathrm{s}}$, the group $\Gamma_{\mathrm{s}}$,
the spaces $\mathcal{H}^{\mathrm{l}}$, $\mathcal{H}^{\mathrm{s}}$ and the operators of Definition~\ref{def:slice-operators} are those of this cross-section.
Let $\sigma_0$ be a plaquette of $\Lambda_{\mathrm{s}}$.
A closed unit square of $\R^2$ that represents $\sigma_0$ is fixed and denoted again by $\sigma_0$;
the representatives of $\sigma_0$ are then the squares $\sigma_0+v_{\mathrm{s}}$, $v\in\Gamma_\parallel$, where $v_{\mathrm{s}}=(v_2,v_3)$.
\begin{enumerate}
\item[(o)]
Let $j\in\Z$, let $\Lambda$ be a non-empty finite set of sites of $\Z^2$,
and let $B$ be a box that contains the cubes $(j,y)+[-1/2,1/2]^3$, $y\in\Lambda$.
If $B$ embeds in $\Gamma$, then $\hull(\Lambda)$ embeds in $\Xi$.
If none of the sets $B+v$, $v\in\Gamma$, contains one of the plaquettes $\{j'\}\times\sigma_0$, $j'\in\Z$,
then $\hull(\Lambda)$ does not contain $\sigma_0$.
\item[(a)]
Let $V$ be a finite-dimensional Hilbert space,
let $f$ be a Schwartz function on $\R^{\Lambda^{(1)}_{\mathrm{s}}}$ with values in $\End V$,
and let $\delta\in\R^{\Lambda^{(1)}_{\mathrm{s}}}$ be such that $A\mapsto\exp(i\langle\delta,A\rangle)$ is $\Gamma_{\mathrm{s}}$-periodic.
Let $f^{\vee}$ be the inverse Fourier transform of $f$,
given by the formula of Definition~\ref{def:integral-kernels} with $\Lambda^{(1)}_{\mathrm{s}}$ in the place of $S$:
\[
f^{\vee}(x)=(2\pi)^{-\#|\Lambda^{(1)}_{\mathrm{s}}|}\int_{\R^{\Lambda^{(1)}_{\mathrm{s}}}}\mathrm{d}\xi\,f(\xi)\cdot\exp\left(i\langle\xi,x\rangle\right),
\qquad x\in\R^{\Lambda^{(1)}_{\mathrm{s}}},
\]
a function with values in $\End V$.
Then the operator $\exp(i\langle\delta,A\rangle)f(E)$ on $L^2(\R^{\Lambda^{(1)}_{\mathrm{s}}}/\Gamma_{\mathrm{s}};V)$,
with $f(E)$ given by the joint spectral calculus of $E$,
is the operator $K[k]$ with the kernel
\[
k(A',A)=\exp(i\langle\delta,A'\rangle)\cdot f^{\vee}(A'-A)
\]
periodised over $\Gamma_{\mathrm{s}}$.
\item[(b)]
Let $\mathfrak{M}$ be formed with $\sigma_0$ (Definition~\ref{def:slice-operators}).
For $a=1,\dots,r$ let $j_a\in\{0,1,\dots,n_1-1\}$ and $c_a:=j_a+1/2$;
let $S_a$ be a finite set of links of $\Z^2$
and $\Lambda_a$ a non-empty finite set of sites of $\Z^2$ that contains the end points of the links of $S_a$;
let $X_a\in\mathfrak{W}(S_a,\Lambda_a)$ and $s_a\in\{1/4,1/8\}$, and let
\[
G^{\circ}_a:=G^{(c_a)}_{X_a}\quad\text{if }s_a=1/4,
\qquad
G^{\circ}_a:=G^{\prime(c_a)}_{X_a}\quad\text{if }s_a=1/8
\]
be the insertions of Definition~\ref{def:two-insertions}.
Assume that the boxes $B_a:=B(c_a,S_a,\Lambda_a)$ of Definition~\ref{def:insertions} embed disjointly in $\Gamma$ (Definition~\ref{def:integrands}),
and that none of the sets $B_a+v$, $a=1,\dots,r$, $v\in\Gamma$, contains a plaquette $\{j'\}\times\sigma_0$, $j'\in\Z$;
these are the conditions \eqref{it:step-disjoint} and \eqref{it:step-string} of Proposition~\ref{prop:word-formula} for the boxes $B_a$.
Let $\beta\in C^2(\Gamma)$.
For $j\in\{0,\dots,n_1-1\}$ let $\beta_j\in\R^{\Lambda^{(1)}_{\mathrm{s}}}$ and $\beta^{\mathrm{mag}}_j\in\R^{\Lambda^{(2)}_{\mathrm{s}}}$ be its electric and magnetic parts in the plane $j$,
$\beta_j(y,\nu):=\beta_{1\nu}(j,y)$ and $\beta^{\mathrm{mag}}_j(y):=\beta_{23}(j,y)$;
let $u_j:=\exp(i\langle\beta^{\mathrm{mag}}_j,\cdot\,\rangle)$;
and let $S_j\subset\Lambda^{(1)}_{\mathrm{s}}$ be the union of the images of those $S_a$ for which $j_a=j$ and $s_a=1/8$.
Then, for every $a$, $\hull(\Lambda_a)$ embeds in $\Xi$ and does not contain $\sigma_0$;
hence the realisation $X\mapsto X_\Xi$ is defined on $\mathfrak{W}(S_a,\Lambda_a)$,
and $\exp(i\langle\delta,\cdot\,\rangle)$ is $\Gamma_{\mathrm{s}}$-periodic for every $\delta\in\mathfrak{D}(S_a)$.
For every $j$ the operators $(\alpha_{\beta_j}X_a)_\Xi$ with $j_a=j$ commute with each other,
$\alpha_{\beta_j}$ being formed with the restriction of $\beta_j$ to the image of $S_a$;
and
\[
Z_{\mathbb{T}}\left[\varphi_\beta\prod_{a=1}^{r}G^{\circ}_a\right]
=(2\pi)^{|n|}\cdot\Tr_{\mathcal{H}^{\mathrm{l}}}\prod_{j=n_1-1,\dots,0}\left(K^{1/2}_{\beta_j,S_j}\left[\prod_{a:\,j_a=j}(\alpha_{\beta_j}X_a)_\Xi\right]K^{1/2}_{\beta_j,S_j}\,\mathfrak{M}\,u_j\,\mathfrak{M}^{*}\right).
\]
Here $|n|=n_1n_2n_3$;
$\varphi_\beta(F)=\exp(i\langle\beta,F\rangle)$ with the pairing of $C^2(\Gamma)$;
the factors of the product over $j$ stand in the order $j=n_1-1,\dots,0$ from left to right;
a product over an empty set of indices $a$ is $1$;
and $u_j$ acts on $\mathcal{H}^{\mathrm{s}}$ as in Proposition~\ref{prop:word-formula}.
\end{enumerate}
\end{lemma}

\begin{proof}
\emph{(o).}
The box $B$ is a product of three closed intervals and contains the points $(j,y)$, $y\in\Lambda$.
Its first factor therefore contains $j$, and its second and third factors are intervals containing the second and third coordinates of all these points;
since $\hull(\Lambda)$ is the smallest product of two intervals with this property, $\{j\}\times\hull(\Lambda)\subset B$.
Every $v\in\Gamma_\parallel$ lies in $\Gamma$, and
\[
\{j\}\times\left(\hull(\Lambda)+[-1,1]^2+v_{\mathrm{s}}\right)\subset B+[-1,1]^3+v .
\]
If $B$ embeds in $\Gamma$, the sets on the right are pairwise disjoint for different $v\in\Gamma_\parallel$ (Definition~\ref{def:integrands}).
Then the sets $\hull(\Lambda)+[-1,1]^2+v_{\mathrm{s}}$, $v\in\Gamma_\parallel$, are pairwise disjoint as well;
that is, $\hull(\Lambda)$ embeds in $\Xi$ (Definition~\ref{def:slices}).
Suppose now that a representative $\sigma_0+v_{\mathrm{s}}$ of $\sigma_0$, $v\in\Gamma_\parallel$, is contained in $\hull(\Lambda)$.
Then $\{j\}\times(\sigma_0+v_{\mathrm{s}})\subset B$, so that the translate $B-v$, where $-v\in\Gamma$, contains the plaquette $\{j\}\times\sigma_0$.
Under the hypothesis of the second assertion this cannot happen.

\emph{(a).}
Expanding $f$ in a basis of $\End V$ reduces the assertion to scalar $f$,
with $f(E)$ acting as $f(E)\otimes 1$ on $L^2(\R^{\Lambda^{(1)}_{\mathrm{s}}}/\Gamma_{\mathrm{s}})\otimes V$;
let $f$ be scalar.
As $f$ is a Schwartz function, so is $f^{\vee}$; in particular $f^{\vee}$ is integrable, and by the Fourier inversion theorem
\[
f(\xi)=\int_{\R^{\Lambda^{(1)}_{\mathrm{s}}}}\mathrm{d}x\,f^{\vee}(x)\,\exp(-i\langle\xi,x\rangle),
\qquad \xi\in\R^{\Lambda^{(1)}_{\mathrm{s}}}.
\]
Let $\Phi,\Phi'\in L^2(\R^{\Lambda^{(1)}_{\mathrm{s}}}/\Gamma_{\mathrm{s}};V)$ and let $P_E$ be the joint spectral measure of $E$;
then $\langle\Phi',P_E(\cdot)\Phi\rangle$ is a finite complex Borel measure on $\R^{\Lambda^{(1)}_{\mathrm{s}}}$.
Since this measure is finite and $f^{\vee}$ is integrable, Fubini's theorem gives
\[
\begin{aligned}
\langle\Phi',f(E)\Phi\rangle
&=\int\mathrm{d}\langle\Phi',P_E(\xi)\Phi\rangle\,f(\xi)
=\int\mathrm{d}x\,f^{\vee}(x)\int\mathrm{d}\langle\Phi',P_E(\xi)\Phi\rangle\,\exp(-i\langle\xi,x\rangle)\\
&=\int\mathrm{d}x\,f^{\vee}(x)\,\langle\Phi',\exp(-i\langle x,E\rangle)\Phi\rangle .
\end{aligned}
\]
By Definition~\ref{def:slice-operators}, $(\exp(-i\langle x,E\rangle)\Phi)([A])=\Phi([A-x])$,
so the last integrand is $f^{\vee}(x)$ times $\int\mathrm{d}[A']\,\langle\Phi'([A']),\Phi([A'-x])\rangle_V$.
The bound
\[
\int\mathrm{d}x\,|f^{\vee}(x)|\int\mathrm{d}[A']\,\|\Phi'([A'])\|\,\|\Phi([A'-x])\|\le\|f^{\vee}\|_{L^1}\cdot\|\Phi'\|\cdot\|\Phi\|
\]
allows a second application of Fubini's theorem, which yields
\[
\langle\Phi',f(E)\Phi\rangle=\int\mathrm{d}[A']\,\left\langle\Phi'([A']),\int_{\R^{\Lambda^{(1)}_{\mathrm{s}}}}\mathrm{d}x\,f^{\vee}(x)\,\Phi([A'-x])\right\rangle_V .
\]
By Minkowski's integral inequality the function $[A']\mapsto\int\mathrm{d}x\,f^{\vee}(x)\,\Phi([A'-x])$ is square integrable,
with norm at most $\|f^{\vee}\|_{L^1}\cdot\|\Phi\|$.
As $\Phi'$ is arbitrary, $f(E)\Phi$ is this function; substituting $A=A'-x$,
\[
(f(E)\Phi)([A'])=\int_{\R^{\Lambda^{(1)}_{\mathrm{s}}}}\mathrm{d}A\,f^{\vee}(A'-A)\,\Phi([A])
\quad\text{for almost every }[A'] .
\]
Thus $f(E)=K[k_0]$ with $k_0(A',A):=f^{\vee}(A'-A)$.
Because $\exp(i\langle\delta,\cdot\,\rangle)$ is $\Gamma_{\mathrm{s}}$-periodic, it is a function on the quotient, and $\langle\delta,\gamma\rangle\in2\pi\Z$ for $\gamma\in\Gamma_{\mathrm{s}}$;
multiplying by $\exp(i\langle\delta,A'\rangle)$ turns $K[k_0]$ into $K[k]$, and $k(A'+\gamma,A+\gamma)=k(A',A)$.

\emph{(b).}
Fix $a$.
By Definition~\ref{def:insertions} the box $B_a$ contains the links $\{j_a\}\times b$, $b\in S_a$, of the plane $x_1=j_a$,
and the cubes $(j_a,y)+[-1/2,1/2]^3$, $y\in\Lambda_a$.
By hypothesis $B_a$ embeds in $\Gamma$, and none of its translates by elements of $\Gamma$ contains a plaquette $\{j'\}\times\sigma_0$.
By (o), $\hull(\Lambda_a)$ embeds in $\Xi$ and does not contain $\sigma_0$;
it is the hull of $(S_a,\Lambda_a)$, because $\Lambda_a$ contains the end points of the links of $S_a$ (Definition~\ref{def:slices}).
Definition~\ref{def:mid-plane} and Lemma~\ref{lem:normal-form}(f) now show that the realisation $X\mapsto X_\Xi$ is defined on $\mathfrak{W}(S_a,\Lambda_a)$,
that every $X_\Xi$ commutes with $P_{\mathrm{G}}$,
and that $\exp(i\langle\delta,\cdot\,\rangle)$ is $\Gamma_{\mathrm{s}}$-periodic for $\delta\in\mathfrak{D}(S_a)$.
The sets $B_a+v$, $a=1,\dots,r$, $v\in\Gamma$, are pairwise disjoint.
Let $j_a=j_{a'}=j$, let $b$ be a link of $\Z^2$ and $v\in\Gamma_\parallel$.
If $\{j\}\times b\subset B_a$ and $\{j\}\times b+v\subset B_{a'}$, then $B_a+v$ meets $B_{a'}$, whence $a=a'$ and $v=0$;
the same holds for the points $(j,y)$, $y\in\Lambda_a$.
So the sets $S_a$ with $j_a=j$ are mapped injectively and with pairwise disjoint images into $\Lambda^{(1)}_{\mathrm{s}}$,
and the sets $\Lambda_a$ with $j_a=j$ likewise into $\Lambda_{\mathrm{s}}$.
Through these injections a link of $S_a$ is at the same time a link of $\Lambda_{\mathrm{s}}$, namely its image,
and a site of $\Lambda_a$ is at the same time a site of $\Lambda_{\mathrm{s}}$,
as in Lemma~\ref{lem:normal-form}(f) and in Proposition~\ref{prop:word-formula}\eqref{it:word-identity}.
The sets $S_a$ and $\Lambda_a$ themselves remain sets of links and of sites of $\Z^2$;
for $A\in\R^{\Lambda^{(1)}_{\mathrm{s}}}$ the vector $A|_{S_a}\in\R^{S_a}$ consists of the components of $A$ at the images of the links of $S_a$,
and a vector of $\R^{S_a}$ is extended by zero to $\Lambda^{(1)}_{\mathrm{s}}$.
Let $S^{(j)}\subset\Lambda^{(1)}_{\mathrm{s}}$ be the union of the images of the sets $S_a$ with $j_a=j$.

\emph{Reduction to single terms.}
The maps $X\mapsto G^{(c)}_X$ and $X\mapsto G^{\prime(c)}_X$ are linear (Definition~\ref{def:two-insertions}),
and so are $X\mapsto X_\Xi$ and $\alpha_\beta$ (Lemma~\ref{lem:normal-form}(e), (f)).
The operator under the trace in (b) is of trace class for all $X_a$:
all its factors are bounded, and it contains the product
$\left(K^{1/2}_{\beta_{n_1-1},S_{n_1-1}}\mathfrak{M}\right)u_{n_1-1}\left(\mathfrak{M}^{*}K^{1/2}_{\beta_{n_1-2},S_{n_1-2}}\right)$
of two Hilbert--Schmidt operators and a bounded one (Corollary~\ref{cor:trace-class}).
Hence both sides of the asserted identity are linear in each $X_a$, and so are the commutators of the operators $(\alpha_{\beta_j}X_a)_\Xi$.
By Lemma~\ref{lem:normal-form}(b), $X_a$ is a finite sum of elements of $\mathfrak{W}(S_a,\Lambda_a)$ of the form $\exp(i\langle\delta,A\rangle)\,\Upsilon f(E_{S_a})$
with $\delta\in\mathfrak{D}(S_a)$, $\Upsilon\in\End_{z(\delta)}(\Lambda_a)$ and $f=f_{S',\mathfrak{Q},\lambda}$ a Gaussian symbol, $S'\subset S_a$ ($f=1$ if $S'$ is empty).
It suffices therefore to prove the assertions when, for every $a$,
\[
X_a=\exp(i\langle\delta_a,A\rangle)\,\Upsilon_a\,f_a(E_{S_a}),
\qquad
\delta_a\in\mathfrak{D}(S_a),\quad \Upsilon_a\in\End_{z(\delta_a)}(\Lambda_a),\quad f_a=f_{S'_a,\mathfrak{Q}_a,\lambda_a}.
\]
Each $\Upsilon_a$ is then even.
Indeed, $X_a$ is even (Lemma~\ref{lem:normal-form}(d)), so it equals $\exp(i\langle\delta_a,A\rangle)\,(U_{\mathrm{f}}\Upsilon_aU_{\mathrm{f}})\,f_a(E_{S_a})$, $U_{\mathrm{f}}$ the fermion parity;
by the uniqueness of the normal form (Lemma~\ref{lem:normal-form}(c)) the symbols agree, $U_{\mathrm{f}}\Upsilon_aU_{\mathrm{f}}\,f_a=\Upsilon_a f_a$, and $f_a$ has no zero.

\emph{Commutation.}
Let $j$ be given and let $j_a=j$.
By Lemma~\ref{lem:normal-form}(e), (f), on every summand of $\mathcal{H}^{\mathrm{l}}$
\[
(\alpha_{\beta_j}X_a)_\Xi=\exp(i\langle\delta_a,A\rangle)\,\iota(\Upsilon_a)\,f_a(E_{S_a}-\beta_j),
\]
where $\delta_a$ is extended by zero to $\Lambda^{(1)}_{\mathrm{s}}$ and $\beta_j$ is restricted to $S_a$.
Let $a'\neq a$ with $j_{a'}=j$.
The operators $\iota(\Upsilon_a)$ and $\iota(\Upsilon_{a'})$ commute, being even and attached to disjoint sets of modes (Proposition~\ref{prop:localisation}\eqref{it:localisation-commute});
they act on the Fock factor and therefore commute with multiplication operators and with functions of $E$.
Multiplication operators commute with each other, and so do functions of $E$.
Finally, $\delta_a$ vanishes on $S_{a'}$, and the relation \eqref{eq:setting-weyl} gives
\[
\begin{aligned}
f_{a'}(E_{S_{a'}}-\beta_j)\,\exp(i\langle\delta_a,A\rangle)
&=\exp(i\langle\delta_a,A\rangle)\,f_{a'}\left(E_{S_{a'}}+\delta_a|_{S_{a'}}-\beta_j\right)\\
&=\exp(i\langle\delta_a,A\rangle)\,f_{a'}(E_{S_{a'}}-\beta_j).
\end{aligned}
\]
So the two realisations commute.
With $\delta^{(j)}:=\sum_{a:\,j_a=j}\delta_a$, which vanishes outside $S^{(j)}$ and coincides with $\delta_a$ on $S_a$, it follows that
\begin{equation}\label{eq:amplification-realised-product}
\prod_{a:\,j_a=j}(\alpha_{\beta_j}X_a)_\Xi
=\exp(i\langle\delta^{(j)},A\rangle)\left[\prod_{a:\,j_a=j}\iota(\Upsilon_a)\right]\prod_{a:\,j_a=j}f_a(E_{S_a}-\beta_j).
\end{equation}

\emph{The integrand as a word.}
For $j\in\{0,\dots,n_1-1\}$ and $x\in\R^{\Lambda^{(1)}_{\mathrm{s}}}$ put
\[
v_j(x):=\exp\left(i\sum_{b\in\Lambda^{(1)}_{\mathrm{s}}\setminus S^{(j)}}\beta_j(b)\,x_b\right),
\]
and for every $a$ with $j_a=j$ let
\[
k_{j,a}(A',A):=\exp(i\langle\beta_j|_{S_a},A'-A\rangle)\cdot k^{(s_a)}_{X_a}(A',A),
\qquad A',A\in\R^{S_a}.
\]
By Lemma~\ref{lem:regularised-kernels}(e), $(S_a,\Lambda_a,k_{j,a})$ is a local step datum.
For any local step datum $(S_a,\Lambda_a,k)$ and any real one-cochain $\hat A$,
the value $G^{(c_a)}[k](\hat A)$ given by Proposition~\ref{prop:word-formula} involves $k$ only through the one matrix $k(A',A)$
with $A'=\hat A_{j+1}|_{S_a}$ and $A=(\hat A_j+d_{\mathrm{s}}\hat a_j)|_{S_a}$,
and it is linear in this matrix, since the insertion $\Ins_{c_a}$ is linear in its argument.
At these two arguments $A'-A$ is the restriction to $S_a$ of $F^{\mathrm{el}}_j=\hat A_{j+1}-\hat A_j-d_{\mathrm{s}}\hat a_j$,
the electric part of $d\hat A$ (Definition~\ref{def:slices});
for a local lift $\hat A$ on $B_a$ it is the electric part of $F$, since then $d\hat A=F$ on the plaquettes of $B_a$ (Definition~\ref{def:integrands}).
Multiplying $k$ by the scalar function $(A',A)\mapsto\exp(i\langle\beta_j|_{S_a},A'-A\rangle)$
therefore multiplies $G^{(c_a)}[k]$ by $\exp(i\langle\beta_j|_{S_a},F^{\mathrm{el}}_j|_{S_a}\rangle)$.
For $k=k^{(s_a)}_{X_a}$, with $G^{\circ}_a=G^{(c_a)}[k^{(s_a)}_{X_a}]$ (Definition~\ref{def:two-insertions}), this gives
\[
G^{(c_a)}[k_{j,a}]=\exp(i\langle\beta_j|_{S_a},F^{\mathrm{el}}_j|_{S_a}\rangle)\cdot G^{\circ}_a .
\]
Further,
\[
\langle\beta,F\rangle=\sum_{j=0}^{n_1-1}\left(\langle\beta_j,F^{\mathrm{el}}_j\rangle+\langle\beta^{\mathrm{mag}}_j,F^{\mathrm{mag}}_j\rangle\right),
\]
and $\Lambda^{(1)}_{\mathrm{s}}$ is the disjoint union of $\Lambda^{(1)}_{\mathrm{s}}\setminus S^{(j)}$ and of the sets $S_a$, $j_a=j$.
Every $a$ belongs to exactly one plane $j_a\in\{0,\dots,n_1-1\}$, and the even elements $G^{\circ}_a$ commute with each other.
Therefore
\[
\varphi_\beta\prod_{a=1}^{r}G^{\circ}_a
=\prod_{j=0}^{n_1-1}\left(u_j(F^{\mathrm{mag}}_j)\cdot v_j(F^{\mathrm{el}}_j)\prod_{a:\,j_a=j}G^{(c_a)}[k_{j,a}]\right).
\]

\emph{The word formula.}
The functions $u_j$ are Borel and unimodular;
the functions $v_j$ are unimodular, so that they satisfy the growth condition of Proposition~\ref{prop:word-formula},
and they do not depend on the variables $x_b$, $b\in S^{(j)}$, which is condition \eqref{it:step-free};
the boxes of the step data $(S_a,\Lambda_a,k_{j_a,a})$ are the boxes $B_a$, which satisfy \eqref{it:step-disjoint} and \eqref{it:step-string}.
Proposition~\ref{prop:word-formula}\eqref{it:word-identity} therefore gives
\begin{equation}\label{eq:amplification-word-trace}
Z_{\mathbb{T}}\left[\varphi_\beta\prod_{a=1}^{r}G^{\circ}_a\right]
=(2\pi)^{|n|}\cdot\Tr_{\mathcal{H}^{\mathrm{l}}}\prod_{j=n_1-1,\dots,0}\left((2\pi/q)^{|\Lambda_{\mathrm{s}}|}\cdot K_j\,P_{\mathrm{G}}\,\mathfrak{M}\,u_j\,\mathfrak{M}^{*}\right),
\end{equation}
where $K_j$ acts on every summand of $\mathcal{H}^{\mathrm{l}}$ as the operator with the kernel
\[
v_j(A'-A)\cdot\prod_{b\in\Lambda^{(1)}_{\mathrm{s}}\setminus S^{(j)}}g_1(A'(b)-A(b))\cdot\prod_{a:\,j_a=j}\iota\left(k_{j,a}(A',A)\right)
\]
periodised over $\Gamma_{\mathrm{s}}$.

\emph{The step operators.}
Consider $j\in\{0,\dots,n_1-1\}$ and an index $a$ with $j_a=j$.
Lemma~\ref{lem:regularised-kernels}(d), applied to the single term $X_a$ and to $\beta_j|_{S_a}$, gives
\[
k_{j,a}(A',A)=\exp(i\langle\delta_a,A'\rangle)\cdot\left(\Upsilon^{(s_a)}_{\delta_a}(\cdot-\beta_j|_{S_a})\right)^{\vee}(A'-A),
\]
where, for $\xi\in\R^{S_a}$,
\[
\Upsilon^{(s_a)}_{\delta_a}(\xi)=\exp\left(-s_a|\xi+\delta_a|^2\right)\cdot f_a(\xi)\cdot\exp\left(-s_a|\xi|^2\right)\cdot\Upsilon_a .
\]
For $b\in\Lambda^{(1)}_{\mathrm{s}}\setminus S^{(j)}$ and $x\in\R$,
\[
\exp(i\beta_j(b)x)\cdot g_1(x)
=\frac{1}{2\pi}\int_{\R}\mathrm{d}\xi\,\exp\left(-\frac{(\xi-\beta_j(b))^2}{2}+i\xi x\right),
\]
the inverse Fourier transform of $\xi\mapsto\exp\left(-(\xi-\beta_j(b))^2/2\right)$.
Insert these formulas into the kernel of $K_j$.
Since $\iota$ is linear and $\Upsilon_a$ is a constant factor of $\Upsilon^{(s_a)}_{\delta_a}$,
the kernel becomes the product of $\exp(i\langle\delta^{(j)},A'\rangle)\cdot\left[\prod_{a:\,j_a=j}\iota(\Upsilon_a)\right]$ with scalar inverse Fourier transforms:
for each link $b\notin S^{(j)}$ the transform of a function of the single variable $\xi_b$, taken at $A'(b)-A(b)$,
and for each $a$ with $j_a=j$ the transform of a function of $\xi|_{S_a}$, taken at $(A'-A)|_{S_a}$.
These groups of links partition $\Lambda^{(1)}_{\mathrm{s}}$.
Accordingly, the constant $(2\pi)^{-\#|\Lambda^{(1)}_{\mathrm{s}}|}$ is the product of the constants of the factors,
and Fubini's theorem writes the integral over $\R^{\Lambda^{(1)}_{\mathrm{s}}}$ of a product of integrable functions of the separate groups of variables
as the product of the separate integrals.
Hence this product of scalar transforms is the inverse Fourier transform on $\R^{\Lambda^{(1)}_{\mathrm{s}}}$, taken at $A'-A$,
of the product of the corresponding functions of $\xi$.
The kernel of $K_j$ is therefore
\[
\exp(i\langle\delta^{(j)},A'\rangle)\cdot\left[\prod_{a:\,j_a=j}\iota(\Upsilon_a)\right]\cdot\left(f^{(j)}\right)^{\vee}(A'-A),
\]
where, for $\xi\in\R^{\Lambda^{(1)}_{\mathrm{s}}}$ and with $|\cdot|_{S_a}$ the Euclidean norm of the components in $S_a$,
\[
\begin{aligned}
f^{(j)}(\xi):={}&\prod_{b\in\Lambda^{(1)}_{\mathrm{s}}\setminus S^{(j)}}\exp\left(-\frac{(\xi_b-\beta_j(b))^2}{2}\right)\\
&\cdot\prod_{a:\,j_a=j}\left(\exp\left(-s_a|\xi-\beta_j+\delta_a|_{S_a}^2\right)\cdot f_a\left((\xi-\beta_j)|_{S_a}\right)\cdot\exp\left(-s_a|\xi-\beta_j|_{S_a}^2\right)\right).
\end{aligned}
\]
The exponents of the factors of $f^{(j)}$ add up to a quadratic polynomial in $\xi$ with complex coefficients,
whose terms of degree two sum to
\[
-\sum_{b\in\Lambda^{(1)}_{\mathrm{s}}\setminus S^{(j)}}\frac{\xi_b^2}{2}-\sum_{a:\,j_a=j}\left(2s_a|\xi|_{S_a}^2+\mathfrak{Q}_a(\xi_{S'_a})\right).
\]
This quadratic form has real coefficients, and it is negative definite on $\R^{\Lambda^{(1)}_{\mathrm{s}}}$:
every link lies either outside $S^{(j)}$ or in exactly one $S_a$, each $s_a$ is positive, and each $\mathfrak{Q}_a$ is positive definite.
So $f^{(j)}$ is a Schwartz function.
The function $\exp(i\langle\delta^{(j)},\cdot\,\rangle)=\prod_{a:\,j_a=j}\exp(i\langle\delta_a,\cdot\,\rangle)$ is $\Gamma_{\mathrm{s}}$-periodic.
Part (a), applied with the Fock space $\mathcal{F}$ of the link plane (Definition~\ref{def:modes}) in the place of $V$ and with the $\End\mathcal{F}$-valued Schwartz function $\left[\prod_{a:\,j_a=j}\iota(\Upsilon_a)\right]f^{(j)}$, yields
\[
K_j=\exp(i\langle\delta^{(j)},A\rangle)\left[\prod_{a:\,j_a=j}\iota(\Upsilon_a)\right]f^{(j)}(E).
\]

\emph{The factors of the word.}
Let $s(b):=1/8$ for $b\in S_j$ and $s(b):=1/4$ for $b\in\Lambda^{(1)}_{\mathrm{s}}\setminus S_j$.
The projection $P_{\mathrm{G}}$ is the average of the gauge transformations of the slice, each of which is a translation of the variables composed with an operator on the Fock factor;
hence $P_{\mathrm{G}}$ commutes with all functions of $E$.
By Definition~\ref{def:slice-operators} it follows that
\[
K_{\beta_j,S_j}=(2\pi/q)^{|\Lambda_{\mathrm{s}}|}\cdot\exp\left(-2\sum_{b\in\Lambda^{(1)}_{\mathrm{s}}}s(b)\,(E_b-\beta_j(b))^2\right)P_{\mathrm{G}},
\]
a product of commuting non-negative operators, whose non-negative square root is
\[
K^{1/2}_{\beta_j,S_j}=(2\pi/q)^{\frac{|\Lambda_{\mathrm{s}}|}{2}}\cdot P_{\mathrm{G}}\exp\left(-\sum_{b\in\Lambda^{(1)}_{\mathrm{s}}}s(b)\,(E_b-\beta_j(b))^2\right).
\]
Every $(\alpha_{\beta_j}X_a)_\Xi$ commutes with $P_{\mathrm{G}}$, as was shown at the beginning of the proof of (b), and $P_{\mathrm{G}}^2=P_{\mathrm{G}}$.
Consequently
\[
\begin{aligned}
&K^{1/2}_{\beta_j,S_j}\left[\prod_{a:\,j_a=j}(\alpha_{\beta_j}X_a)_\Xi\right]K^{1/2}_{\beta_j,S_j}\\
&\qquad=(2\pi/q)^{|\Lambda_{\mathrm{s}}|}\cdot P_{\mathrm{G}}\exp\left(-\sum_{b}s(b)\,(E_b-\beta_j(b))^2\right)\\
&\qquad\qquad\cdot\left[\prod_{a:\,j_a=j}(\alpha_{\beta_j}X_a)_\Xi\right]\exp\left(-\sum_{b}s(b)\,(E_b-\beta_j(b))^2\right),
\end{aligned}
\]
and the product of the last three factors commutes with $P_{\mathrm{G}}$.
To compute this product, insert \eqref{eq:amplification-realised-product}.
By \eqref{eq:setting-weyl}, applied to $\delta^{(j)}$ and to the function $\xi\mapsto\exp\left(-\sum_bs(b)\,(\xi_b-\beta_j(b))^2\right)$,
\[
\begin{aligned}
&\exp\left(-\sum_{b}s(b)\,(E_b-\beta_j(b))^2\right)\exp(i\langle\delta^{(j)},A\rangle)\\
&\qquad=\exp(i\langle\delta^{(j)},A\rangle)\exp\left(-\sum_{b}s(b)\,(E_b+\delta^{(j)}(b)-\beta_j(b))^2\right),
\end{aligned}
\]
and the operators $\iota(\Upsilon_a)$ commute with functions of $E$; therefore
\[
\begin{aligned}
&\exp\left(-\sum_{b}s(b)\,(E_b-\beta_j(b))^2\right)\left[\prod_{a:\,j_a=j}(\alpha_{\beta_j}X_a)_\Xi\right]\exp\left(-\sum_{b}s(b)\,(E_b-\beta_j(b))^2\right)\\
&\qquad=\exp(i\langle\delta^{(j)},A\rangle)\left[\prod_{a:\,j_a=j}\iota(\Upsilon_a)\right]\exp\left(-\sum_{b}s(b)\,(E_b+\delta^{(j)}(b)-\beta_j(b))^2\right)\\
&\qquad\qquad\cdot\prod_{a:\,j_a=j}f_a(E_{S_a}-\beta_j)\cdot\exp\left(-\sum_{b}s(b)\,(E_b-\beta_j(b))^2\right).
\end{aligned}
\]
For $b\notin S^{(j)}$ one has $s(b)=1/4$ and $\delta^{(j)}(b)=0$.
For $b\in S_a$ with $j_a=j$ one has $\delta^{(j)}(b)=\delta_a(b)$ and $s(b)=s_a$:
if $s_a=1/8$, then $S_a\subset S_j$;
if $s_a=1/4$, then $S_a$ is disjoint from $S_j$, since the sets $S_{a'}$ with $j_{a'}=j$ are pairwise disjoint.
Comparing with the definition of $f^{(j)}$, one sees that the three functions of $E$ on the right multiply to $f^{(j)}(E)$,
so that both sides of the last identity are equal to $K_j$.
As $P_{\mathrm{G}}K_j=K_jP_{\mathrm{G}}$, it follows that
\begin{equation}\label{eq:amplification-step-operator}
K^{1/2}_{\beta_j,S_j}\left[\prod_{a:\,j_a=j}(\alpha_{\beta_j}X_a)_\Xi\right]K^{1/2}_{\beta_j,S_j}
=(2\pi/q)^{|\Lambda_{\mathrm{s}}|}\cdot K_j\,P_{\mathrm{G}} .
\end{equation}
Substituting \eqref{eq:amplification-step-operator} for every $j$ into \eqref{eq:amplification-word-trace} yields the identity of (b).
\end{proof}

\subsubsection{The covariant box}\label{sec:covariant-box}

Let $c=j+\frac12$ with $j\in\Z$, let $S$ be a finite set of links of the slice lattice $\Z^2$,
and let $\Lambda$ be a non-empty finite set of sites of $\Z^2$ that contains the end points of the links of $S$.
Definition~\ref{def:insertions} attaches to $(c,S,\Lambda)$ the set $\mathcal{K}(c,S,\Lambda)$,
the smallest product $\prod_\mu[x^-_\mu,x^+_\mu]$ of intervals that contains it,
the centre $x^0$ of this product,
the number $r_0=\max_\mu\frac12\left(x^+_\mu-x^-_\mu\right)$, which is the radius of the box $B(c,S,\Lambda)$,
and the box $B^{\#}(c,S,\Lambda)=x^0+[-2r_0,2r_0]^3$.
Every point of $\mathcal{K}(c,S,\Lambda)$ satisfies $j-\frac12\le x_1\le j+\frac32$,
and since $\Lambda$ is not empty, both bounds are attained on the cubes $(j,y)+\left[-\frac12,\frac12\right]^3$ and $(j+1,y)+\left[-\frac12,\frac12\right]^3$, $y\in\Lambda$;
hence $x^+_1-x^-_1=2$ and $r_0\ge1$.

The group $W_{\mathrm{c}}$ acts on such triples.
An element $g\in W_{\mathrm{c}}$ has the form $g(x)=\vec g\,x+v$, with a diagonal matrix $\vec g\in\{\pm1\}^3$ and $v\in\Z^3$ (Definition~\ref{def:mirrors}).
It maps every plane orthogonal to $e_1$ onto a plane orthogonal to $e_1$,
and there is an isometry $y\mapsto gy$ of $\R^2$ that maps $\Z^2$ onto itself
such that $g$ maps the point above $y$ of a plane orthogonal to $e_1$ to the point above $gy$ of the image plane.
Let $g(c)\in\Z+\frac12$ be defined by $g\{x_1=c\}=\{x_1=g(c)\}$,
and let $gS$ and $g\Lambda$ be the images of $S$ and $\Lambda$ under $y\mapsto gy$;
the set $g\Lambda$ is again non-empty and contains the end points of the links of $gS$.

\begin{lemma}[A covariant box]\label{lem:covariant-box}
Let $c=j+\frac12$ with $j\in\Z$, let $S$ be a finite set of links of $\Z^2$,
and let $\Lambda$ be a non-empty finite set of sites of $\Z^2$ that contains the end points of the links of $S$.
Let $\mathcal{K}(c,S,\Lambda)$, $r_0$, $B(c,S,\Lambda)$ and $B^{\#}(c,S,\Lambda)$ be as in Definition~\ref{def:insertions}.
\begin{enumerate}
\item[(a)]
For every $g\in W_{\mathrm{c}}$,
\[
g\,\mathcal{K}(c,S,\Lambda)=\mathcal{K}(g(c),gS,g\Lambda),
\qquad
g\,B^{\#}(c,S,\Lambda)=B^{\#}(g(c),gS,g\Lambda).
\]
\item[(b)]
Every box of radius $r_0$ that contains $\mathcal{K}(c,S,\Lambda)$ is contained in $B^{\#}(c,S,\Lambda)$;
in particular $B(c,S,\Lambda)\subseteq B^{\#}(c,S,\Lambda)$.
\item[(c)]
Let $\mathbb{T}=C(n_1,n_2,n_3)=(\Gamma,\chi)$ be a coordinate torus,
and let $g_1,\dots,g_r$ be a placement of $B^{\#}(c,S,\Lambda)$ in $\mathbb{T}$ (Definition~\ref{def:bounded-elements}).
Let $\sigma=\{j''\}\times\sigma'$, with $j''\in\Z$ and a plaquette $\sigma'$ of $\Z^2$,
be a plaquette orthogonal to $e_1$ as in condition~\eqref{it:placement-columns} of Definition~\ref{def:bounded-elements} for $\mu=1$,
and let $\sigma_0$ be the plaquette of the slice torus $\Lambda_{\mathrm{s}}=\Z^2/(n_2\Z\times n_3\Z)$ represented by $\sigma'$.
Then the boxes $B(g_a(c),g_aS,g_a\Lambda)$, $a=1,\dots,r$, satisfy the conditions~\eqref{it:step-disjoint} and~\eqref{it:step-string} of Proposition~\ref{prop:word-formula},
the latter with this plaquette $\sigma_0$;
and for $a=1,\dots,r$ the hull $\hull(g_aS,g_a\Lambda)$ embeds in the cross-section $\Xi(\mathbb{T},x_1)$ (Definitions~\ref{def:slabs} and~\ref{def:slices}).
\end{enumerate}
\end{lemma}

\begin{proof}
(a)
Let $g\in W_{\mathrm{c}}$, and write $g(c)=j'+\frac12$ with $j'\in\Z$.
Being an isometry that maps planes orthogonal to $e_1$ onto such planes,
$g$ maps the two planes $\{x_1=j\}$ and $\{x_1=j+1\}$, which lie at distance $\frac12$ from $\{x_1=c\}$,
onto the planes $\{x_1=j'\}$ and $\{x_1=j'+1\}$, in this order or in the opposite one.
Consequently $g$ maps the pair of links $(j,b)$, $(j+1,b)$, $b\in S$, onto the pair $(j',gb)$, $(j'+1,gb)$,
and the link from $(j,y)$ to $(j+1,y)$, $y\in\Lambda$, onto the link from $(j',gy)$ to $(j'+1,gy)$.
As $\vec g\,\left[-\frac12,\frac12\right]^3=\left[-\frac12,\frac12\right]^3$, the unit cube centred at a site goes onto the unit cube centred at the image site;
so the cubes $(j,y)+\left[-\frac12,\frac12\right]^3$ and $(j+1,y)+\left[-\frac12,\frac12\right]^3$, $y\in\Lambda$,
go onto the cubes $(j',gy)+\left[-\frac12,\frac12\right]^3$ and $(j'+1,gy)+\left[-\frac12,\frac12\right]^3$.
This is the first identity.
In each coordinate, $g$ is an affine bijection $x_\mu\mapsto\pm x_\mu+v_\mu$ of $\R$.
Hence it maps every product of intervals onto a product of intervals with sides of the same lengths,
the smallest product of intervals that contains a set onto the smallest product of intervals that contains the image of the set,
and the centre of the former onto the centre of the latter.
By the first identity, the number $r_0$ of $(g(c),gS,g\Lambda)$ is that of $(c,S,\Lambda)$,
and the centre attached to $(g(c),gS,g\Lambda)$ is $g\,x^0$.
Since $\vec g\,[-2r_0,2r_0]^3=[-2r_0,2r_0]^3$, it follows that
\[
g\,B^{\#}(c,S,\Lambda)=g\,x^0+[-2r_0,2r_0]^3=B^{\#}(g(c),gS,g\Lambda).
\]

(b)
A box $x+[-r_0,r_0]^3$, $x\in\R^3$, is itself a product of intervals.
It contains $\mathcal{K}(c,S,\Lambda)$ if and only if it contains $\prod_\mu[x^-_\mu,x^+_\mu]$, that is, if and only if
\[
x^+_\mu-r_0\le x_\mu\le x^-_\mu+r_0\qquad(\mu=1,2,3).
\]
In that case $|x_\mu-x^0_\mu|\le r_0-\frac12\left(x^+_\mu-x^-_\mu\right)\le r_0$ for every $\mu$,
and every point $x'$ of the box satisfies $|x'_\mu-x^0_\mu|\le|x'_\mu-x_\mu|+|x_\mu-x^0_\mu|\le2r_0$;
so the box lies in $B^{\#}(c,S,\Lambda)$.
The box $B(c,S,\Lambda)$ is a box of radius $r_0$ that contains $\mathcal{K}(c,S,\Lambda)$.

(c)
Let $a\in\{1,\dots,r\}$.
By (a), applied to $g_a$, and by (b), applied to the triple $(g_a(c),g_aS,g_a\Lambda)$,
\[
B(g_a(c),g_aS,g_a\Lambda)\subseteq B^{\#}(g_a(c),g_aS,g_a\Lambda)=g_a\,B^{\#}(c,S,\Lambda).
\]
The properties of boxes that enter here are inherited by smaller boxes.
A box embeds in $\Gamma$ if its translates by the elements of $\Gamma$, each enlarged by $[-1,1]^3$, are pairwise disjoint;
then every box contained in it embeds in $\Gamma$.
Boxes $B'_1,\dots,B'_r$ embed disjointly in $\Gamma$ if each of them embeds and the sets $B'_a+v$, $a=1,\dots,r$, $v\in\Gamma$, are pairwise disjoint (Definition~\ref{def:integrands});
then boxes $B_a\subseteq B'_a$ embed disjointly in $\Gamma$ as well.
And a set that contains none of the plaquettes of a given family has no subset that contains one of them.
By condition~\eqref{it:placement-disjoint} of Definition~\ref{def:bounded-elements} the boxes $g_aB^{\#}(c,S,\Lambda)$ embed disjointly in $\Gamma$;
hence the boxes $B(g_a(c),g_aS,g_a\Lambda)$ satisfy condition~\eqref{it:step-disjoint} of Proposition~\ref{prop:word-formula}.
The plaquettes $\{j'\}\times\sigma'$, $j'\in\Z$, are the translates $\sigma+(j'-j'')\,e_1$,
and by condition~\eqref{it:placement-columns} none of them is contained in a set $g_aB^{\#}(c,S,\Lambda)+v$, $v\in\Gamma$.
The representatives of $\sigma_0$ in $\Z^2$ are the squares $\sigma'+v'_{\mathrm{s}}$,
where $v'$ ranges over $\Gamma_\parallel=n_2\Z e_2+n_3\Z e_3$, a subgroup of $\Gamma$, and $v'_{\mathrm{s}}=(v'_2,v'_3)$;
the plaquette $\{j'\}\times(\sigma'+v'_{\mathrm{s}})$ is the translate of $\{j'\}\times\sigma'$ by $v'$.
So it is contained in $g_aB^{\#}(c,S,\Lambda)+v$ only if $\{j'\}\times\sigma'$ is contained in $g_aB^{\#}(c,S,\Lambda)+v-v'$, which is excluded.
Therefore no set $g_aB^{\#}(c,S,\Lambda)+v$, and no subset $B(g_a(c),g_aS,g_a\Lambda)+v$ of it, $v\in\Gamma$,
contains a plaquette of $\Z^3$ that lies in a plane $\{x_1=j'\}$, $j'\in\Z$, and represents $\sigma_0$:
this is condition~\eqref{it:step-string} of Proposition~\ref{prop:word-formula}.
Finally, write $g_a(c)=j_a+\frac12$.
The box $B(g_a(c),g_aS,g_a\Lambda)$ embeds in $\Gamma$, as shown above,
and it contains the cubes $(j_a,y)+\left[-\frac12,\frac12\right]^3$, $y\in g_a\Lambda$, which are part of $\mathcal{K}(g_a(c),g_aS,g_a\Lambda)$.
Lemma~\ref{lem:mixed-word-formula}(o), applied to this box, shows that $\hull(g_a\Lambda)$ embeds in $\Xi(\mathbb{T},x_1)$;
and $\hull(g_aS,g_a\Lambda)=\hull(g_a\Lambda)$, because $g_a\Lambda$ contains the end points of the links of $g_aS$ (Definition~\ref{def:slices}).
\end{proof}

\paragraph{Integrands read in different boxes.}
An integrand that reads in a box (Definition~\ref{def:integrands}) can be regarded as an integrand of any larger box,
and two integrands of different boxes can then be compared in a box that contains both.

\begin{lemma}[Equality of integrands]\label{lem:equal-integrands}
Let $B$, $B'$, $B_1$ and $B_2$ be boxes.
\begin{enumerate}
\item[(a)]
Let $B\subseteq B'$.
Then $B^{(1)}\subseteq B'^{(1)}$, $B^{(2)}\subseteq B'^{(2)}$ and $\mathbb{X}(B)\subseteq\mathbb{X}(B')$.
For $G\in\mathcal{G}^{\mathrm{ev}}(B)$, the function $\hat A\mapsto G\left(\hat A|_{B^{(1)}}\right)$ of $\hat A\in\R^{B'^{(1)}}$,
which does not depend on the values of $\hat A$ on the links of $B'^{(1)}\setminus B^{(1)}$, is called \emph{$G$ read in $B'$};
it is an element of $\mathcal{G}^{\mathrm{ev}}(B')$, with the same constant $C$ as $G$ and with $\min\left(\varepsilon_G,\frac12\right)$ in the place of $\varepsilon_G$.
\item[(b)]
Let $G_1\in\mathcal{G}^{\mathrm{ev}}(B_1)$ and $G_2\in\mathcal{G}^{\mathrm{ev}}(B_2)$.
If $G_1$ and $G_2$, read in $B'$, coincide for one box $B'$ that contains $B_1$ and $B_2$,
then $G_1$ and $G_2$, read in $B''$, coincide for every box $B''$ that contains $B_1$ and $B_2$.
In this case $G_1$ and $G_2$ are called \emph{equal}, and one writes $G_1=G_2$.
\item[(c)]
Let $\Gamma$ be a period lattice in which $B_1$ and $B_2$ embed,
and let $G_1\in\mathcal{G}^{\mathrm{ev}}(B_1)$ and $G_2\in\mathcal{G}^{\mathrm{ev}}(B_2)$ be equal.
Then, for every torus $\mathbb{T}=(\Gamma,\chi)$, the integrands $G_1$ and $G_2$ define the same function on $\Conf(\Gamma)$ with values in $\Gr_{\mathbb{T}}$;
in particular $Z_{\mathbb{T}}[G_1]=Z_{\mathbb{T}}[G_2]$.
\end{enumerate}
\end{lemma}

\begin{proof}
(a)
A link, a plaquette or a unit cube that is contained in $B$ is contained in $B'$; this gives the three inclusions.
Let $G\in\mathcal{G}^{\mathrm{ev}}(B)$, and let $G'$ be $G$ read in $B'$.
It is a Borel function on $\R^{B'^{(1)}}$ with values in the even part of $\Gr(\mathbb{X}(B))$, a subalgebra of $\Gr(\mathbb{X}(B'))$.
Let $\lambda$ be a real function on the sites of $B'$.
The gauge transformation by $\lambda$ changes the values of $\hat A$ on the links of $B$ by those of $d\lambda$,
which depend only on the values of $\lambda$ at the end points of these links,
and it multiplies the fields at a site $x\in\mathbb{X}(B)$ by $\exp\left(\pm iq\lambda(\pi(x))\right)$.
The end points of the links of $B$ and the sites $\pi(x)$, $x\in\mathbb{X}(B)$, are sites of $B$;
so only the restriction of $\lambda$ to the sites of $B$ enters, and the gauge invariance of $G$ gives that of $G'$.
Let $a_\Psi$ be the coefficient of a monomial $\Psi$ in $G$, so that
\[
\left|a_\Psi(\hat A)\right|\le C\cdot\exp\left(\left(\tfrac12-\varepsilon_G\right)\sum_{\sigma\in B^{(2)}}(d\hat A)_\sigma^2\right).
\]
If $\varepsilon_G\le\frac12$, the right side does not decrease when the sum is extended over $B'^{(2)}$;
if $\varepsilon_G>\frac12$, the right side is at most $C$.
In both cases
\[
\left|a_\Psi(\hat A)\right|\le C\cdot\exp\left(\left(\tfrac12-\min\left(\varepsilon_G,\tfrac12\right)\right)\sum_{\sigma\in B'^{(2)}}(d\hat A)_\sigma^2\right),
\]
and $\min\left(\varepsilon_G,\tfrac12\right)>0$.
(For $\varepsilon_G>\frac12$ the bound with $\varepsilon_G$ itself, summed over the larger set $B'^{(2)}$, would be stronger than the given one.)

(b)
Read in $B'$, the integrands $G_1$ and $G_2$ are the functions $\hat A\mapsto G_1\left(\hat A|_{B_1^{(1)}}\right)$ and $\hat A\mapsto G_2\left(\hat A|_{B_2^{(1)}}\right)$ of $\hat A\in\R^{B'^{(1)}}$,
with values in $\Gr(\mathbb{X}(B'))$.
Every real function on $B_1^{(1)}\cup B_2^{(1)}$ is the restriction of a real function on $B'^{(1)}$.
So the two coincide if and only if $G_1\left(\hat A|_{B_1^{(1)}}\right)=G_2\left(\hat A|_{B_2^{(1)}}\right)$ for every real function $\hat A$ on $B_1^{(1)}\cup B_2^{(1)}$,
an identity in the Grassmann algebra $\Gr(\mathbb{X})$ that does not involve $B'$.

(c)
Fix a torus $\mathbb{T}=(\Gamma,\chi)$.
Since $B_1$ and $B_2$ embed in $\Gamma$, the integrands $G_1$ and $G_2$ are $\Gr_{\mathbb{T}}$-valued functions on $\Conf(\Gamma)$:
their values at a point are $G_1(\hat A_1)$ and $G_2(\hat A_2)$, where $\hat A_1$ and $\hat A_2$ are local lifts of the point on $B_1$ and on $B_2$,
and these values do not depend on the choice of the lifts (Definition~\ref{def:integrands}).
For a box $B'$ that contains $B_1$ and $B_2$, let $G'\in\mathcal{G}^{\mathrm{ev}}(B')$ be the common value of $G_1$ and $G_2$ read in $B'$.

Suppose first that some box $B'$ that contains $B_1$ and $B_2$ embeds in $\Gamma$.
This is the case, for instance, for $B'=B^{\#}(c,S,\Lambda)$
when $B_1$ and $B_2$ are boxes of radius $r_0$ that contain $\mathcal{K}(c,S,\Lambda)$ and $B^{\#}(c,S,\Lambda)$ embeds in $\Gamma$ (Lemma~\ref{lem:covariant-box}(b)).
A local lift $\hat A'$ of a point of $\Conf(\Gamma)$ on $B'$ restricts to local lifts on $B_1$ and on $B_2$,
because $B_1^{(1)},B_2^{(1)}\subseteq B'^{(1)}$ and $B_1^{(2)},B_2^{(2)}\subseteq B'^{(2)}$;
hence $G_1\left(\hat A'|_{B_1^{(1)}}\right)=G'(\hat A')=G_2\left(\hat A'|_{B_2^{(1)}}\right)$.

In general, let $B'$ be any box that contains $B_1$ and $B_2$.
The integrand $G'$ does not depend on the values of $\hat A\in\R^{B'^{(1)}}$ on the links outside $B_1^{(1)}$, nor on those outside $B_2^{(1)}$.
A change of the values on the links outside $B_1^{(1)}\cap B_2^{(1)}$ is a change outside $B_1^{(1)}$ followed by a change outside $B_2^{(1)}$;
so $G'$ depends only on the values on $B_1^{(1)}\cap B_2^{(1)}$, which is the set of the links contained in $B_1\cap B_2$.
Its values lie in $\Gr(\mathbb{X}(B_1))\cap\Gr(\mathbb{X}(B_2))=\Gr\left(\mathbb{X}(B_1)\cap\mathbb{X}(B_2)\right)$, as the expansion in the basis of monomials shows.
Let $\hat A_1$ and $\hat A_2$ be local lifts of one point of $\Conf(\Gamma)$ on $B_1$ and on $B_2$.
By the definition of $G'$, one has $G_1(\hat A_1)=G'(\hat A)$ for every $\hat A\in\R^{B'^{(1)}}$ that extends $\hat A_1$,
and $G_2(\hat A_2)=G'(\hat A)$ for every $\hat A\in\R^{B'^{(1)}}$ that extends $\hat A_2$.
If $B_1\cap B_2$ contains no link, $G'$ is a constant element, and $G_1(\hat A_1)=G_2(\hat A_2)$.
Otherwise $B_1\cap B_2$ is a product of three closed intervals,
and the cells of $\Z^3$ contained in it are the cells of a product of three intervals with integer end points;
they form a contractible cell complex.
On the links of this complex the cochain $\hat A_1-\hat A_2$ takes values in $(2\pi/q)\Z$,
because $\exp\left(iq\hat A_1\right)=\exp(iqA)=\exp\left(iq\hat A_2\right)$ there,
and it is closed, because $d\hat A_1=F=d\hat A_2$ on the plaquettes of the complex.
Hence $\hat A_1-\hat A_2=(2\pi/q)\,d\lambda$ on the links contained in $B_1\cap B_2$, with an integer valued function $\lambda$ on the sites of $B_1\cap B_2$;
extend $\lambda$ by $0$ to the remaining sites of $B_1$.
The real function $(2\pi/q)\lambda$ satisfies $\exp\left(iq(2\pi/q)\lambda\right)=1$,
so the gauge transformation by it leaves the fields unchanged, and the gauge invariance of $G_1$ gives
\[
G_1(\hat A_1)=G_1\left(\hat A_1-(2\pi/q)\,d\lambda\right).
\]
On the links contained in $B_1\cap B_2$, whose end points are sites of $B_1\cap B_2$, the cochain $\hat A_1-(2\pi/q)\,d\lambda$ coincides with $\hat A_2$.
As $G'$ depends only on the values on these links, $G_1\left(\hat A_1-(2\pi/q)\,d\lambda\right)=G_2(\hat A_2)$.
Thus $G_1$ and $G_2$ take the same value at every point of $\Conf(\Gamma)$.
By Definition~\ref{def:functional}, $Z_{\mathbb{T}}[G_1]$ and $Z_{\mathbb{T}}[G_2]$ are the integrals of these equal functions
against the same weight $\exp\left(-\frac12\|F\|^2\right)$ and the same Berezin integral over the fields of $\mathbb{T}$;
so they are equal.
\end{proof}

Let now $(S,\Lambda,k)$ be a local step datum (Proposition~\ref{prop:word-formula}) with $\Lambda\neq\emptyset$,
and let $G^{(c)}[k]$ be its step element at $c=j+\frac12$ (Definition~\ref{def:two-insertions}).
It reads only the links of $\mathcal{K}(c,S,\Lambda)$ and the fields at the fermion sites $x$ with $\pi(x)+\left[-\frac12,\frac12\right]^3\subseteq\mathcal{K}(c,S,\Lambda)$:
the arguments $\hat A_{j+1}|_S$ and $\left(\hat A_j+d_{\mathrm{s}}\hat a_j\right)|_S$ of $k$, the operator $u(\hat a_j)$,
the electric variables $F^{\mathrm{el}}_j(b)$, $b\in S$, the modes over $\Lambda$ and the crossing terms $S^{\times}_c(\Lambda)$
involve only the links $(j,b)$ and $(j+1,b)$, $b\in S$, the links from $(j,y)$ to $(j+1,y)$, $y\in\Lambda$,
and the fields at the sites $x$ with $\pi(x)\in\{(j,y),(j+1,y)\}$, $y\in\Lambda$.
The set $\mathcal{K}(c,S,\Lambda)$ is covariant under $W_{\mathrm{c}}$ by Lemma~\ref{lem:covariant-box}(a), and it increases with $S$ and with $\Lambda$.
By Proposition~\ref{prop:word-formula}\eqref{it:step-class}, $G^{(c)}[k]$ is an element of $\mathcal{G}^{\mathrm{ev}}\left(B(c,S,\Lambda)\right)$;
read in $B^{\#}(c,S,\Lambda)$, which contains $B(c,S,\Lambda)$ by Lemma~\ref{lem:covariant-box}(b),
it is an element of $\mathcal{G}^{\mathrm{ev}}\left(B^{\#}(c,S,\Lambda)\right)$ by Lemma~\ref{lem:equal-integrands}(a).
Identities between step elements of different data, such as those of Lemma~\ref{lem:enlarged-datum} below, of Lemma~\ref{lem:coordinate-images},
of Proposition~\ref{prop:insertions} and of Proposition~\ref{prop:two-layer-insertions},
are equalities in the sense of Lemma~\ref{lem:equal-integrands}(b).
They relate integrands that are read in different boxes,
because the box $B(c,S,\Lambda)$ is neither covariant under reflections nor increasing in $(S,\Lambda)$ (see the remark following Definition~\ref{def:insertions}):
for $c=\frac12$, the box $B(c,\emptyset,\{(0,0)\})=\left[-\frac12,\frac32\right]\times\left[-\frac32,\frac12\right]^2$
is not contained in $B(c,\emptyset,\{(0,0),(1,0)\})=\left[-\frac12,\frac32\right]^2\times\left[-\frac32,\frac12\right]$.

\begin{remark}
(1)
In the word formulas of Lemma~\ref{lem:mixed-word-formula}(b) and of Proposition~\ref{prop:insertions}(d),
the conditions~\eqref{it:step-disjoint} and~\eqref{it:step-string} of Proposition~\ref{prop:word-formula} are imposed on the boxes $B(c_a,S_a,\Lambda_a)$ of the factors,
and each factor is read in its own box $B(c_a,S_a,\Lambda_a)$;
the larger boxes $B^{\#}(c_a,S_a,\Lambda_a)$ need not embed.
For instance, for $c=\frac12$, $S=\emptyset$ and $\Lambda=\{(0,0)\}$ one has $r_0=1$,
\[
B(c,S,\Lambda)=\left[-\tfrac12,\tfrac32\right]\times\left[-\tfrac32,\tfrac12\right]^2,
\qquad
B^{\#}(c,S,\Lambda)=\left(\tfrac12,0,0\right)+[-2,2]^3;
\]
enlarged by $[-1,1]^3$, these boxes have the sides $4$ and $6$,
so the first embeds in the period lattice of $C(n_1,n_2,n_3)$ if and only if $n_1,n_2,n_3\ge5$, the second if and only if $n_1,n_2,n_3\ge7$.
By Lemma~\ref{lem:equal-integrands}(c), two equal step elements define the same function on $\Conf(\Gamma)$, and have the same value under $Z_{\mathbb{T}}$, as soon as the two boxes in which they are read embed in $\Gamma$.
This holds in particular for two step elements read in the box $B^{\#}(c,S,\Lambda)$ of a datum, or in its image $g\,B^{\#}(c,S,\Lambda)$ under an element $g$ of $W$, if this box embeds in $\Gamma$.
If $g_1,\dots,g_r$ is a placement of $B^{\#}(c,S,\Lambda)$ in a coordinate torus,
the boxes $g_aB^{\#}(c,S,\Lambda)=B^{\#}(g_a(c),g_aS,g_a\Lambda)$ embed in its period lattice,
by condition~\eqref{it:placement-disjoint} of Definition~\ref{def:bounded-elements} and Lemma~\ref{lem:covariant-box}(a).

(2)
The box $B^{\#}(c,S,\Lambda)$ does not increase with $(S,\Lambda)$ either.
Let $c=\frac12$ and $S=\emptyset$.
For $\Lambda=\{(0,0),(9,0)\}$ one finds $r_0=5$, and the projection of $B^{\#}(c,S,\Lambda)$ to the $x_3$-axis is $[-10,10]$;
for $\Lambda_1:=\Lambda\cup\{(0,1)\}$ again $r_0=5$, while the projection of $B^{\#}(c,S,\Lambda_1)$ to the $x_3$-axis is $\left[-\frac{19}{2},\frac{21}{2}\right]$.
Thus, for $S\subseteq S_1$ and $\Lambda\subseteq\Lambda_1$, neither of the boxes $B^{\#}(c,S,\Lambda)$ and $B^{\#}(c,S_1,\Lambda_1)$ need contain the other,
and an equality between the step elements of the two data, as in Lemma~\ref{lem:enlarged-datum},
relates integrands read in two boxes of which neither contains the other.
Lemma~\ref{lem:equal-integrands}(c) applies to such an equality whenever the two boxes embed in $\Gamma$, whether or not some box that contains both of them embeds.
\end{remark}

\paragraph{Enlargement of a datum.}
A local step datum can be enlarged by further links and sites without changing its step element.
Recall the Fock spaces $\mathcal{F}_\Lambda=\mathcal{F}\left(\mathcal{J}^{\mathrm{l}}(\Lambda)\right)$ and the diagonal unitaries $u(a)$, $a\in\R^\Lambda$, of Definition~\ref{def:charged-operators}.

\begin{lemma}[Enlargement]\label{lem:enlarged-datum}
Let $c=j+\frac12$ with $j\in\Z$, and let $(S,\Lambda,k)$ be a local step datum (Proposition~\ref{prop:word-formula}) with $\Lambda\neq\emptyset$.
Let $S_1\supseteq S$ be a finite set of links of $\Z^2$ and $\Lambda_1\supseteq\Lambda$ a finite set of sites of $\Z^2$ that contains the end points of the links of $S_1$,
let $\iota\colon\End\mathcal{F}_\Lambda\to\End\mathcal{F}_{\Lambda_1}$ be the localisation (Definition~\ref{def:modes}), and put
\[
k_1(A',A):=\iota\left(k\left(A'|_S,A|_S\right)\right)\cdot\prod_{b\in S_1\setminus S}g_1\left(A'(b)-A(b)\right),
\qquad A',A\in\R^{S_1}.
\]
Then $(S_1,\Lambda_1,k_1)$ is a local step datum, and
\[
G^{(c)}[k_1]=G^{(c)}[k].
\]
If $X\in\mathfrak{W}(S,\Lambda)$ and $X_1$ is its image in $\mathfrak{W}(S_1,\Lambda_1)$ (Definition~\ref{def:mid-plane}), then $G^{(c)}_{X_1}=G^{(c)}_X$.
Both identities are equalities in the sense of Lemma~\ref{lem:equal-integrands}.
\end{lemma}

\begin{proof}
\emph{The datum.}
The function $k_1$ is Borel, and its values are even.
Indeed, for $\Upsilon\in\End\mathcal{F}_\Lambda$, $I,J\subseteq\mathcal{J}^{\mathrm{l}}(\Lambda)$ and $R,R'\subseteq\mathcal{J}^{\mathrm{l}}(\Lambda_1)\setminus\mathcal{J}^{\mathrm{l}}(\Lambda)$,
the entry of $\iota(\Upsilon)$ at $(I\cup R,J\cup R')$ vanishes for $R\neq R'$ and equals $\pm\Upsilon_{IJ}$ for $R=R'$ (Definition~\ref{def:modes});
as $\#|I\cup R|-\#|J\cup R|=\#|I|-\#|J|$, the operator $\iota(\Upsilon)$ is even when $\Upsilon$ is.
The localisation is isometric (Proposition~\ref{prop:localisation}\eqref{it:localisation-homomorphism}), and $g_1(s)\le\exp\left(-\frac12s^2\right)$ for real $s$.
So the Gaussian bound of $k$, with constants $C$ and $\varepsilon_k$, gives
\[
\left\|k_1(A',A)\right\|\le C\cdot\prod_{b\in S_1}\exp\left(-\min\left(\varepsilon_k,\tfrac12\right)\left(A'(b)-A(b)\right)^2\right).
\]
For the covariance let $a\in\R^{\Lambda_1}$.
For $b\in S$ the value $(d_{\mathrm{s}}a)(b)$ depends only on the values of $a$ at the end points of $b$, which lie in $\Lambda$;
so the restriction of $d_{\mathrm{s}}a$ to $S$ is $d_{\mathrm{s}}(a|_\Lambda)$, read on $S$.
The operator $u(a)$ on $\mathcal{F}_{\Lambda_1}$ is diagonal with the entries $\exp\left(iq\langle a,\mathfrak{q}(I')\rangle\right)$, $I'\subseteq\mathcal{J}^{\mathrm{l}}(\Lambda_1)$,
and $\mathfrak{q}(I\cup R)=\mathfrak{q}(I)+\mathfrak{q}(R)$ for disjoint $I$ and $R$.
Hence the entry of $u(a)^{-1}\,\iota(\Upsilon)\,u(a)$ at $(I\cup R,J\cup R)$ is that of $\iota(\Upsilon)$ multiplied by $\exp\left(-iq\langle a,\mathfrak{q}(I)-\mathfrak{q}(J)\rangle\right)$.
The charges $\mathfrak{q}(I)$ and $\mathfrak{q}(J)$ of sets of modes over $\Lambda$ are supported in $\Lambda$;
so this factor is $\exp\left(-iq\langle a|_\Lambda,\mathfrak{q}(I)-\mathfrak{q}(J)\rangle\right)$,
the factor by which $u(a|_\Lambda)^{-1}\,\Upsilon\,u(a|_\Lambda)$ multiplies the entry $\Upsilon_{IJ}$.
Since $\iota$ is defined entry by entry, it follows that
\[
u(a)^{-1}\,\iota(\Upsilon)\,u(a)=\iota\left(u(a|_\Lambda)^{-1}\,\Upsilon\,u(a|_\Lambda)\right),\qquad\Upsilon\in\End\mathcal{F}_\Lambda.
\]
The factors $g_1\left(A'(b)-A(b)\right)$ do not change when $A'$ and $A$ are translated by the same vector.
Together with the covariance of $k$ this gives, for $A',A\in\R^{S_1}$ and with $d_{\mathrm{s}}a$ read on $S_1$,
\begin{align*}
u(a)^{-1}\,k_1\left(A'+d_{\mathrm{s}}a,A+d_{\mathrm{s}}a\right)u(a)
&=\iota\left(u(a|_\Lambda)^{-1}\,k\left(A'|_S+d_{\mathrm{s}}(a|_\Lambda),A|_S+d_{\mathrm{s}}(a|_\Lambda)\right)u(a|_\Lambda)\right)\\
&\qquad\cdot\prod_{b\in S_1\setminus S}g_1\left(A'(b)-A(b)\right)\\
&=k_1(A',A).
\end{align*}
Thus $(S_1,\Lambda_1,k_1)$ is a local step datum.

\emph{The step elements.}
Let $\hat A$ be a real one-cochain on $\Z^3$, at which both step elements are evaluated as in Proposition~\ref{prop:word-formula}.
At the arguments $A'=\hat A_{j+1}|_{S_1}$ and $A=\left(\hat A_j+d_{\mathrm{s}}\hat a_j\right)|_{S_1}$ one has $A'(b)-A(b)=F^{\mathrm{el}}_j(b)$ for $b\in S_1$,
since $F^{\mathrm{el}}_j=\hat A_{j+1}-\hat A_j-d_{\mathrm{s}}\hat a_j$ is the electric part of $d\hat A$ (Proposition~\ref{prop:word-formula}).
Hence, with the even operator
\[
\Upsilon:=k\left(\hat A_{j+1}|_S,\,\left(\hat A_j+d_{\mathrm{s}}\hat a_j\right)|_S\right)\in\End\mathcal{F}_\Lambda,
\]
one has
\[
k_1\left(\hat A_{j+1}|_{S_1},\,\left(\hat A_j+d_{\mathrm{s}}\hat a_j\right)|_{S_1}\right)=\iota(\Upsilon)\cdot\prod_{b\in S_1\setminus S}g_1\left(F^{\mathrm{el}}_j(b)\right).
\]
On $\mathcal{F}_{\Lambda_1}$ the operator $u(\hat a_j)$ is the diagonal operator with the entries $u_{I'}=\prod_{i\in I'}u_i$,
where $u_i=\exp\left(iq\langle\hat a_j,\mathfrak{q}(i)\rangle\right)$ is the phase of the mode $i$, and on $\mathcal{F}_\Lambda$ it is the diagonal operator with the same entries for $I'\subseteq\mathcal{J}^{\mathrm{l}}(\Lambda)$.
The insertion of a localised even operator factorises (Proposition~\ref{prop:localisation}, for the index sets $\mathcal{J}^{\mathrm{l}}(\Lambda)\subseteq\mathcal{J}^{\mathrm{l}}(\Lambda_1)$):
\[
\Ins^{\mathcal{J}^{\mathrm{l}}(\Lambda_1)}_c\left[\omega^{\frac12}\,\iota(\Upsilon)\,u(\hat a_j)\,\omega^{\frac12}\right]
=\Ins^{\mathcal{J}^{\mathrm{l}}(\Lambda)}_c\left[\omega^{\frac12}\,\Upsilon\,u(\hat a_j)\,\omega^{\frac12}\right]\cdot\prod_{i}\left(1+\omega_i\cdot u_i\cdot\Theta_c(\zeta_i)\,\zeta_i\right).
\]
Here the insertions are formed at the plane $\{x_1=c\}$ with the modes indicated,
$\omega$ and $u(\hat a_j)$ act on $\mathcal{F}_{\Lambda_1}$ on the left side and on $\mathcal{F}_\Lambda$ on the right side,
and the product extends over the modes $i\in\mathcal{J}^{\mathrm{l}}(\Lambda_1)\setminus\mathcal{J}^{\mathrm{l}}(\Lambda)$.
These are the modes of the bonds over $\Lambda_1\setminus\Lambda$ that cross $\{x_1=c\}$.
The factors of the product are even and commute;
expanding the product and applying Proposition~\ref{prop:crossing} to these bonds,
one finds that it equals $\exp\left(-S^{\times}_c(\Lambda_1\setminus\Lambda)\right)$,
where $S^{\times}_c(\Lambda_1\setminus\Lambda)$ is the sum of the terms of these bonds, defined by the formula of Definition~\ref{def:insertions}.
Moreover $S^{\times}_c(\Lambda_1)=S^{\times}_c(\Lambda)+S^{\times}_c(\Lambda_1\setminus\Lambda)$, a sum of two even and therefore commuting elements,
so that $\exp\left(S^{\times}_c(\Lambda_1)\right)=\exp\left(S^{\times}_c(\Lambda)\right)\cdot\exp\left(S^{\times}_c(\Lambda_1\setminus\Lambda)\right)$.
In the definition of $G^{(c)}[k_1]$ the scalar factors $g_1\left(F^{\mathrm{el}}_j(b)\right)$, $b\in S_1\setminus S$, may be taken out of the insertion,
and they cancel against the corresponding factors of $\prod_{b\in S_1}g_1\left(F^{\mathrm{el}}_j(b)\right)^{-1}$.
Since all the factors are even, they commute, and
\begin{align*}
G^{(c)}[k_1]
&=\Ins^{\mathcal{J}^{\mathrm{l}}(\Lambda)}_c\left[\omega^{\frac12}\,\Upsilon\,u(\hat a_j)\,\omega^{\frac12}\right]\cdot\exp\left(-S^{\times}_c(\Lambda_1\setminus\Lambda)\right)\\
&\qquad\cdot\prod_{b\in S}g_1\left(F^{\mathrm{el}}_j(b)\right)^{-1}\cdot\exp\left(S^{\times}_c(\Lambda)\right)\cdot\exp\left(S^{\times}_c(\Lambda_1\setminus\Lambda)\right)\\
&=G^{(c)}[k].
\end{align*}
This holds for every real one-cochain $\hat A$ on $\Z^3$;
so $G^{(c)}[k_1]$ and $G^{(c)}[k]$ coincide when they are read in a box that contains $B(c,S_1,\Lambda_1)$ and $B(c,S,\Lambda)$,
and they are equal in the sense of Lemma~\ref{lem:equal-integrands}(b).

\emph{Insertions.}
Let $X\in\mathfrak{W}(S,\Lambda)$.
Its image $X_1$ is the tensor product of $X$, with the Fock part carried to $\mathcal{F}_{\Lambda_1}$ by $\iota$, with the identity of $L^2(\R^{S_1\setminus S})$ (Definition~\ref{def:mid-plane}).
The operator $\exp\left(-\frac14|E_{S_1}|^2\right)$ is the product of $\exp\left(-\frac14|E_S|^2\right)$ with the operators $\exp\left(-\frac14E_b^2\right)$, $b\in S_1\setminus S$,
which act on the factors of $L^2(\R^{S_1\setminus S})$.
Hence $\exp\left(-\frac14|E_{S_1}|^2\right)X_1\exp\left(-\frac14|E_{S_1}|^2\right)$ is the tensor product of the image of $\exp\left(-\frac14|E_S|^2\right)X\exp\left(-\frac14|E_S|^2\right)$
with the operators $\exp\left(-\frac14E_b^2\right)\cdot1\cdot\exp\left(-\frac14E_b^2\right)=\exp\left(-\frac12E_b^2\right)$, $b\in S_1\setminus S$;
and $\exp\left(-\frac12E_b^2\right)$ has the kernel $g_1\left(A'(b)-A(b)\right)$ (Definition~\ref{def:integral-kernels}).
Therefore
\[
k_{X_1}(A',A)=\iota\left(k_X\left(A'|_S,A|_S\right)\right)\cdot\prod_{b\in S_1\setminus S}g_1\left(A'(b)-A(b)\right),
\]
that is, $k_{X_1}$ is formed from $k_X$ as $k_1$ is formed from $k$.
By Lemma~\ref{lem:regularised-kernels}(e) the triple $(S,\Lambda,k_X)$ is a local step datum, and $G^{(c)}_X=G^{(c)}[k_X]$ by Definition~\ref{def:two-insertions}.
The first part of the lemma now gives $G^{(c)}_{X_1}=G^{(c)}[k_{X_1}]=G^{(c)}[k_X]=G^{(c)}_X$.
\end{proof}

\begin{remark}
For the insertions into the comparison operator, an enlargement of the set of links changes the integrand.
Let $X$, $X_1$, $S_1$ and $\Lambda_1$ be as in Lemma~\ref{lem:enlarged-datum}.
The operator $\exp\left(-\frac18E_b^2\right)\cdot1\cdot\exp\left(-\frac18E_b^2\right)=\exp\left(-\frac14E_b^2\right)$ has the kernel $g_{\frac12}\left(A'(b)-A(b)\right)$,
so the argument of the last part of that proof gives
\[
k^{(\frac18)}_{X_1}(A',A)=\iota\left(k^{(\frac18)}_X\left(A'|_S,A|_S\right)\right)\cdot\prod_{b\in S_1\setminus S}g_{\frac12}\left(A'(b)-A(b)\right).
\]
The computation of the step elements in the same proof, with $g_{\frac12}$ in the place of $g_1$ in the factors over $S_1\setminus S$,
leaves the factors $g_{\frac12}\left(F^{\mathrm{el}}_j(b)\right)\cdot g_1\left(F^{\mathrm{el}}_j(b)\right)^{-1}=2^{\frac12}\cdot\exp\left(-\frac12F^{\mathrm{el}}_j(b)^2\right)$, $b\in S_1\setminus S$; hence
\[
G^{\prime(c)}_{X_1}=G^{\prime(c)}_X\cdot\prod_{b\in S_1\setminus S}2^{\frac12}\cdot\exp\left(-\tfrac12F^{\mathrm{el}}_j(b)^2\right).
\]
Unlike $G_X$, the integrand $G'_X$ therefore depends on the set $S$ of the algebra $\mathfrak{W}(S,\Lambda)$ in which $X$ is taken, and not only on $X$ as an element of $\mathfrak{W}$.
\end{remark}

\subsubsection{Images under the coordinate group}\label{sec:coordinate-images}

For a link plane $\{x_1=c\}$, $c=j+\frac12$ with $j\in\Z$,
the reflections attached to the mirror $(x_1,c)$ are denoted by $\theta_c$ and $\Theta_c$.
For a local step datum $(S,\Lambda,k)$ (Proposition~\ref{prop:word-formula}) put
\[
\Upsilon^{(c)}[k](\hat A):=\omega^{\frac12}\,k\left(\hat A_{j+1}|_S,\,\left(\hat A_j+d_{\mathrm{s}}\hat a_j\right)|_S\right)u(\hat a_j)\,\omega^{\frac12},
\qquad
g_S:=\prod_{b\in S}g_1\left(F^{\mathrm{el}}_j(b)\right),
\]
the function $\hat a_j$ being restricted to $\Lambda$ in $u(\hat a_j)$.
The step element of the datum at the plane $c$ is
\[
G^{(c)}[k]=\Ins_c\left[\Upsilon^{(c)}[k]\right]\cdot g_S^{-1}\cdot\exp\left(S^{\times}_c(\Lambda)\right),
\]
with $\Ins_c$, $\omega$, $u$ and $S^{\times}_c(\Lambda)$ as in Definition~\ref{def:insertions};
at another link plane of the same direction the same formulas are used with the time coordinate of that plane in the place of $j$.

A \emph{spatial mirror} is a mirror $\Pi_w=(x_\nu,w)$ with $\nu\in\{2,3\}$ and $w\in\frac12\Z$.
Its reflection $\theta_w$ leaves the coordinate $x_1$ unchanged and commutes with $\theta_c$;
$\Theta_w$ is the antilinear reflection attached to $\Pi_w$,
and $\theta_{\mathrm{s}}$ is the reflection of the slice lattice $\Z^2$ in the line $\{x_\nu=w\}$, which $\theta_w$ induces on the last two coordinates.
For a link $b$ of $\Z^2$ let $\operatorname{sgn}(\theta_{\mathrm{s}},b):=-1$ if $b$ is parallel to $e_\nu$, and $\operatorname{sgn}(\theta_{\mathrm{s}},b):=1$ otherwise.
This is the factor by which $\theta_{\mathrm{s}}$ changes the orientation of $b$,
and $\operatorname{sgn}(\theta_{\mathrm{s}},\theta_{\mathrm{s}}b)=\operatorname{sgn}(\theta_{\mathrm{s}},b)$.
For a finite set $S$ of links of $\Z^2$ and $A\in\R^{\theta_{\mathrm{s}}S}$ the pull-back $\theta_{\mathrm{s}}^{*}A\in\R^S$ is
\[
\left(\theta_{\mathrm{s}}^{*}A\right)(b):=\operatorname{sgn}(\theta_{\mathrm{s}},b)\cdot A(\theta_{\mathrm{s}}b),\qquad b\in S,
\]
as in Definition~\ref{def:cochains},
and for a function $a$ on sites $\theta_{\mathrm{s}}^{*}a:=a\circ\theta_{\mathrm{s}}$.
The pull-back permutes the coordinates up to signs;
it is therefore orthogonal and preserves Lebesgue measure,
the pull-back from $\R^S$ to $\R^{\theta_{\mathrm{s}}S}$ is its inverse,
and $\theta_{\mathrm{s}}^{*}d_{\mathrm{s}}a=d_{\mathrm{s}}(a\circ\theta_{\mathrm{s}})$.

\begin{lemma}[Images of step elements]\label{lem:coordinate-images}
Let $c=j+\frac12$ with $j\in\Z$, and let $(S,\Lambda,k)$ be a local step datum with $\Lambda\neq\emptyset$,
with the constants $C$ and $\varepsilon_k$ of its bound
$\|k(A',A)\|\le C\cdot\prod_{b\in S}\exp\left(-\varepsilon_k\left(A'(b)-A(b)\right)^2\right)$.
\begin{enumerate}
\item\label{it:images-translation}
Let $v\in\Z^3$ and $v_{\mathrm{s}}:=(v_2,v_3)$.
For $A\in\R^{S+v_{\mathrm{s}}}$ let $\tau^{*}_{v_{\mathrm{s}}}A\in\R^S$ be the function $b\mapsto A(b+v_{\mathrm{s}})$,
let $U_v\colon\mathcal{F}_\Lambda\to\mathcal{F}_{\Lambda+v_{\mathrm{s}}}$ be the unitary of Lemma~\ref{lem:translated-insertions}, and put
\[
k^{\tau_v}(A',A):=U_v\,k\left(\tau^{*}_{v_{\mathrm{s}}}A',\tau^{*}_{v_{\mathrm{s}}}A\right)U_v^{*},\qquad A',A\in\R^{S+v_{\mathrm{s}}}.
\]
Then $\left(S+v_{\mathrm{s}},\Lambda+v_{\mathrm{s}},k^{\tau_v}\right)$ is a local step datum with the constants $C$ and $\varepsilon_k$, and
\[
\left(G^{(c)}[k]\right)^{\tau_v}=G^{(c+v_1)}\left[k^{\tau_v}\right].
\]
\item\label{it:images-time-reflection}
Put $k^{\theta_c}(A',A):=k(A,A')^{*}$ for $A',A\in\R^S$.
Then $\left(S,\Lambda,k^{\theta_c}\right)$ is a local step datum with the constants $C$ and $\varepsilon_k$, and
\[
\Theta_c\,G^{(c)}[k]=G^{(c)}\left[k^{\theta_c}\right].
\]
\item\label{it:images-spatial-mirror}
Let $\Pi_w$ be a spatial mirror, and let $\mathfrak{U}\colon\mathcal{F}_\Lambda\to\mathcal{F}_{\theta_{\mathrm{s}}\Lambda}$ be the unitary of Lemma~\ref{lem:spatial-mirror-signs}.
Put
\[
k^{\theta_w}(A',A):=\mathfrak{U}\,\overline{k\left(\theta_{\mathrm{s}}^{*}A',\theta_{\mathrm{s}}^{*}A\right)}\,\mathfrak{U}^{*},\qquad A',A\in\R^{\theta_{\mathrm{s}}S},
\]
where the bar is the entrywise complex conjugation in the basis $(e_I)$.
Then $\left(\theta_{\mathrm{s}}S,\theta_{\mathrm{s}}\Lambda,k^{\theta_w}\right)$ is a local step datum with the constants $C$ and $\varepsilon_k$, and
\[
\Theta_w\,G^{(c)}[k]=G^{(c)}\left[k^{\theta_w}\right].
\]
\item\label{it:images-elements}
Let $v$ and $\Pi_w$ be as in \eqref{it:images-translation} and \eqref{it:images-spatial-mirror}.
Let the unitary $\mathfrak{T}_v\colon\mathcal{H}_S\to\mathcal{H}_{S+v_{\mathrm{s}}}$ and the antiunitary $\mathfrak{J}\colon\mathcal{H}_S\to\mathcal{H}_{\theta_{\mathrm{s}}S}$ be given by
\[
\left(\mathfrak{T}_v\Phi\right)(A):=U_v\,\Phi\left(\tau^{*}_{v_{\mathrm{s}}}A\right),\qquad
\left(\mathfrak{J}\Phi\right)(A):=\mathfrak{U}\,\overline{\Phi\left(\theta_{\mathrm{s}}^{*}A\right)},
\]
and for $X\in\mathfrak{W}(S,\Lambda)$ put
\[
X^{\tau_v}:=\mathfrak{T}_v\,X\,\mathfrak{T}_v^{-1},\qquad X^{\theta_c}:=X^{*},\qquad X^{\theta_w}:=\mathfrak{J}\,X\,\mathfrak{J}^{-1}.
\]
The maps $X\mapsto X^{\tau_v}$, $X\mapsto X^{\theta_c}$ and $X\mapsto X^{\theta_w}$ are bijections of $\mathfrak{W}(S,\Lambda)$ onto
$\mathfrak{W}(S+v_{\mathrm{s}},\Lambda+v_{\mathrm{s}})$, $\mathfrak{W}(S,\Lambda)$ and $\mathfrak{W}(\theta_{\mathrm{s}}S,\theta_{\mathrm{s}}\Lambda)$, respectively;
the first is a linear $*$-isomorphism, the second is the adjoint, and the third is an antilinear $*$-isomorphism.
The three maps preserve $\|\cdot\|_\#$, and for $0<s\le\frac14$
\[
\left(k^{(s)}_X\right)^{\tau_v}=k^{(s)}_{X^{\tau_v}},\qquad
\left(k^{(s)}_X\right)^{\theta_c}=k^{(s)}_{X^{*}},\qquad
\left(k^{(s)}_X\right)^{\theta_w}=k^{(s)}_{X^{\theta_w}}.
\]
\item\label{it:images-group}
Every $g\in W_{\mathrm{c}}$ has exactly one \emph{normal word}
\[
g=\tau_v\,\theta_c^{p_1}\,\theta_{x_2,0}^{p_2}\,\theta_{x_3,0}^{p_3},\qquad p=(p_1,p_2,p_3)\in\{0,1\}^3,\quad v\in\Z^3,
\]
where $\theta_{x_2,0}$ and $\theta_{x_3,0}$ are the reflections of the spatial mirrors $(x_2,0)$ and $(x_3,0)$.
For $X\in\mathfrak{W}(S,\Lambda)$ put
\[
X^g:=\left(\left(\left(X^{\theta_{x_3,0}^{p_3}}\right)^{\theta_{x_2,0}^{p_2}}\right)^{\theta_c^{p_1}}\right)^{\tau_v},
\]
each transport being that of \eqref{it:images-elements} and an exponent $0$ standing for the identity.
If $g=\theta_w$ is the reflection of a spatial mirror $\Pi_w=(x_\nu,w)$, then its normal word is $\tau_{2w\,e_\nu}\,\theta_{x_\nu,0}$,
and $X^g$ coincides with the element $X^{\theta_w}=\mathfrak{J}X\mathfrak{J}^{-1}$ of \eqref{it:images-elements}.
For every $g\in W_{\mathrm{c}}$ one has $X^g\in\mathfrak{W}(gS,g\Lambda)$, where $gS$ and $g\Lambda$ are the images of $S$ and $\Lambda$ under the map of $\Z^2$ that $g$ induces on the last two coordinates;
the insertions of $X^g$ are formed at the plane $g(c):=c+v_1$, the image of the plane $\{x_1=c\}$ under $g$, and
\[
\left(G^{(c)}_X\right)^g=G^{(g(c))}_{X^g},\qquad\left(G'^{(c)}_X\right)^g=G'^{(g(c))}_{X^g},\qquad\left\|X^g\right\|_\#=\|X\|_\#.
\]
In particular, for a mirror $(x_1,c')$ with $c'\in\frac12\Z$ and for $v:=2(c'-c)\,e_1$,
\[
\Theta_{c'}\,G^{(c)}_X=G^{(2c'-c)}_{\left(X^{*}\right)^{\tau_v}}.
\]
\end{enumerate}
The identities between integrands in \eqref{it:images-translation}, \eqref{it:images-time-reflection}, \eqref{it:images-spatial-mirror} and \eqref{it:images-group}
are equalities in the sense of Lemma~\ref{lem:equal-integrands}.
Both sides of each of them read only links contained in $\mathcal{K}$ and fields at fermion sites $x$ with $\pi(x)+\left[-\frac12,\frac12\right]^3\subseteq\mathcal{K}$,
and both lie in $\mathcal{G}^{\mathrm{ev}}\left(B^{\#}\right)$,
where $\mathcal{K}$ and $B^{\#}$ are formed with the plane and the sets of the right-hand side:
with $(c+v_1,S+v_{\mathrm{s}},\Lambda+v_{\mathrm{s}})$ in \eqref{it:images-translation}, with $(c,S,\Lambda)$ in \eqref{it:images-time-reflection},
with $(c,\theta_{\mathrm{s}}S,\theta_{\mathrm{s}}\Lambda)$ in \eqref{it:images-spatial-mirror} and with $(g(c),gS,g\Lambda)$ in \eqref{it:images-group}.
\end{lemma}

\begin{proof}
The values of $k$ are even operators, and $\omega^{\frac12}$ and $u(\hat a_j)$ are diagonal;
so $\Upsilon^{(c)}[k]$ takes even values, and $\Ins_c\left[\Upsilon^{(c)}[k]\right]$ is an even element.
The factors $g_S^{-1}$ and $\exp\left(S^{\times}_c(\Lambda)\right)$ are even as well, and even elements commute with one another.
Consequently an automorphism maps $G^{(c)}[k]$ to the product of the images of its three factors,
and so does an antilinear anti-automorphism such as $\Theta_c$ or $\Theta_w$, since the images of the three factors are again even.
Recall that in the datum $(S,\Lambda,k)$ the set $S$ is a finite set of links of $\Z^2$,
the set $\Lambda$, which is not empty by hypothesis, is a finite set of sites that contains their end points,
and $k$ is a Borel function on $\R^S\times\R^S$ whose values are even operators on $\mathcal{F}_\Lambda$,
which obeys the Gaussian bound above and is covariant:
\[
u(a)^{-1}\,k\left(A'+d_{\mathrm{s}}a,\,A+d_{\mathrm{s}}a\right)u(a)=k(A',A),\qquad a\in\R^\Lambda .
\]
The maps of $\Z^2$ that $\tau_v$, $\theta_w$ and an element $g\in W_{\mathrm{c}}$ induce on the last two coordinates are isometries,
and they map the end points of a link onto the end points of its image;
so $\Lambda+v_{\mathrm{s}}$, $\theta_{\mathrm{s}}\Lambda$ and $g\Lambda$ are non-empty finite sets of sites that contain the end points of the links of $S+v_{\mathrm{s}}$, $\theta_{\mathrm{s}}S$ and $gS$, respectively.
The functions $k^{\tau_v}$, $k^{\theta_c}$ and $k^{\theta_w}$ are Borel functions.

\emph{Translations.}
The pull-back by $\tau_v$ carries no signs: $\left(\tau_v^{*}\hat A\right)(b')=\hat A(b'+v)$ for every link $b'$ of $\Z^3$.
In the slice variables this reads, for $j'\in\Z$ and for links $b$ and sites $y$ of $\Z^2$,
\[
\left(\tau_v^{*}\hat A\right)_{j'}(b)=\hat A_{j'+v_1}(b+v_{\mathrm{s}}),\qquad
\left(\tau_v^{*}\hat a\right)_{j'}(y)=\hat a_{j'+v_1}(y+v_{\mathrm{s}}),\qquad
F^{\mathrm{el}}_{j'}\left(\tau_v^{*}F\right)(b)=F^{\mathrm{el}}_{j'+v_1}(F)(b+v_{\mathrm{s}}),
\]
and $d_{\mathrm{s}}$ commutes with the translation by $v_{\mathrm{s}}$.
By Lemma~\ref{lem:translated-insertions}, $G^{\tau_v}=\alpha_vG$ for every even integrand $G$;
for every function $\hat A\mapsto\Upsilon(\hat A)$ whose values are even elements of $\End\mathcal{F}_\Lambda$
\[
\left(\Ins_c[\Upsilon]\right)^{\tau_v}(\hat A)=\Ins_{c+v_1}\left[U_v\,\Upsilon\left(\tau_v^{*}\hat A\right)U_v^{*}\right];
\]
moreover $U_v\,\omega=\omega\,U_v$, and $U_v\,u\left(a(\,\cdot+v_{\mathrm{s}})\right)U_v^{*}=u(a)$ for $a\in\R^{\Lambda+v_{\mathrm{s}}}$.
This applies to $\Upsilon=\Upsilon^{(c)}[k]$, whose values are even;
with the restrictions to $S+v_{\mathrm{s}}$ and to $\Lambda+v_{\mathrm{s}}$ understood, it follows that
\[
\begin{aligned}
U_v\,\Upsilon^{(c)}[k]\left(\tau_v^{*}\hat A\right)U_v^{*}
&=\omega^{\frac12}\,U_v\,k\left(\tau^{*}_{v_{\mathrm{s}}}\hat A_{j+v_1+1},\,\tau^{*}_{v_{\mathrm{s}}}\left(\hat A_{j+v_1}+d_{\mathrm{s}}\hat a_{j+v_1}\right)\right)U_v^{*}\,U_v\,u\left(\hat a_{j+v_1}(\,\cdot+v_{\mathrm{s}})\right)U_v^{*}\,\omega^{\frac12}\\
&=\omega^{\frac12}\,k^{\tau_v}\left(\hat A_{j+v_1+1},\,\hat A_{j+v_1}+d_{\mathrm{s}}\hat a_{j+v_1}\right)u\left(\hat a_{j+v_1}\right)\omega^{\frac12},
\end{aligned}
\]
which is $\Upsilon^{(c+v_1)}\left[k^{\tau_v}\right](\hat A)$, formed at the plane $c+v_1=(j+v_1)+\frac12$ with the datum $\left(S+v_{\mathrm{s}},\Lambda+v_{\mathrm{s}},k^{\tau_v}\right)$.
So $\Ins_c\left[\Upsilon^{(c)}[k]\right]$ is carried to $\Ins_{c+v_1}\left[\Upsilon^{(c+v_1)}\left[k^{\tau_v}\right]\right]$.
The factor $g_S^{-1}$ is carried to
\[
\prod_{b\in S}g_1\left(F^{\mathrm{el}}_{j+v_1}(b+v_{\mathrm{s}})\right)^{-1}=\prod_{b\in S+v_{\mathrm{s}}}g_1\left(F^{\mathrm{el}}_{j+v_1}(b)\right)^{-1},
\]
the corresponding factor of the datum $\left(S+v_{\mathrm{s}},\Lambda+v_{\mathrm{s}},k^{\tau_v}\right)$ at the plane $c+v_1$.
Finally, $\alpha_v$ maps the term $-\bar\psi(x)\,w(x,x')\,U_A(x,x')\,\psi(x')$ of an oriented bond $(x,x')$ of $S_F$ to the term of $(x+v,x'+v)$:
the type and the direction of a bond do not change under a translation by a vector of $\Z^3$, so that $w(x+v,x'+v)=w(x,x')$, and $\pi(x+v)=\pi(x)+v$.
The bonds over $\Lambda$ that cross $\{x_1=c\}$ are thus carried onto the bonds over $\Lambda+v_{\mathrm{s}}$ that cross $\{x_1=c+v_1\}$,
and $\left(\exp\left(S^{\times}_c(\Lambda)\right)\right)^{\tau_v}=\exp\left(S^{\times}_{c+v_1}(\Lambda+v_{\mathrm{s}})\right)$.
The product of the three images is $G^{(c+v_1)}\left[k^{\tau_v}\right]$.

The values of $k^{\tau_v}$ are even, because $U_v$ maps each $e_I$ to $\pm e_{I+v}$ and $\#|I+v|=\#|I|$.
As $U_v$ is unitary and $\tau^{*}_{v_{\mathrm{s}}}$ only renames the coordinates,
\[
\left\|k^{\tau_v}(A',A)\right\|\le C\cdot\prod_{b\in S}\exp\left(-\varepsilon_k\left(A'(b+v_{\mathrm{s}})-A(b+v_{\mathrm{s}})\right)^2\right)
=C\cdot\prod_{b\in S+v_{\mathrm{s}}}\exp\left(-\varepsilon_k\left(A'(b)-A(b)\right)^2\right).
\]
For the covariance let $a\in\R^{\Lambda+v_{\mathrm{s}}}$ and $a':=a(\,\cdot+v_{\mathrm{s}})\in\R^\Lambda$.
Then $\tau^{*}_{v_{\mathrm{s}}}d_{\mathrm{s}}a=d_{\mathrm{s}}a'$, and the relation $U_v\,u(a')\,U_v^{*}=u(a)$ gives
$u(a)^{-1}U_v=U_v\,u(a')^{-1}$ and $U_v^{*}u(a)=u(a')\,U_v^{*}$;
hence, by the covariance of $k$ for $a'$,
\[
u(a)^{-1}\,k^{\tau_v}\left(A'+d_{\mathrm{s}}a,\,A+d_{\mathrm{s}}a\right)u(a)
=U_v\,u(a')^{-1}\,k\left(\tau^{*}_{v_{\mathrm{s}}}A'+d_{\mathrm{s}}a',\,\tau^{*}_{v_{\mathrm{s}}}A+d_{\mathrm{s}}a'\right)u(a')\,U_v^{*}
=k^{\tau_v}(A',A).
\]

\emph{The link plane.}
By \eqref{eq:setting-slice-reflections}, the reflection $\theta_c$ acts on the slice variables by
\[
\left(\theta_c^{*}\hat A\right)_j=\hat A_{j+1},\qquad\left(\theta_c^{*}\hat A\right)_{j+1}=\hat A_j,\qquad\left(\theta_c^{*}\hat a\right)_j=-\hat a_j,\qquad F^{\mathrm{el}}_j\left(\theta_c^{*}F\right)=-F^{\mathrm{el}}_j(F).
\]
Hence $\Upsilon^{(c)}[k]\left(\theta_c^{*}\hat A\right)=\omega^{\frac12}\,k\left(\hat A_j,\,\hat A_{j+1}-d_{\mathrm{s}}\hat a_j\right)u(-\hat a_j)\,\omega^{\frac12}$,
the restrictions to $S$ and to $\Lambda$ being understood;
since $\omega$ is real and diagonal and $u(-\hat a_j)^{*}=u(\hat a_j)$,
\[
\Upsilon^{(c)}[k]\left(\theta_c^{*}\hat A\right)^{*}
=\omega^{\frac12}\,u(\hat a_j)\,k\left(\hat A_j,\,\hat A_{j+1}-d_{\mathrm{s}}\hat a_j\right)^{*}\omega^{\frac12}
=\omega^{\frac12}\,u(\hat a_j)\,k^{\theta_c}\left(\hat A_{j+1}-d_{\mathrm{s}}\hat a_j,\,\hat A_j\right)\omega^{\frac12}.
\]
Taking adjoints in the covariance of $k$ and exchanging the names of the two arguments one obtains the covariance of $k^{\theta_c}$,
$u(a)^{-1}\,k^{\theta_c}(A'+d_{\mathrm{s}}a,A+d_{\mathrm{s}}a)\,u(a)=k^{\theta_c}(A',A)$ for $a\in\R^\Lambda$.
Replacing $A'$ by $A'-d_{\mathrm{s}}a$ and multiplying by $u(a)$ from the left gives
\[
u(a)\,k^{\theta_c}\left(A'-d_{\mathrm{s}}a,\,A\right)=k^{\theta_c}\left(A',\,A+d_{\mathrm{s}}a\right)u(a),
\]
and with $a=\hat a_j$, $A'=\hat A_{j+1}$, $A=\hat A_j$ the preceding display becomes
\[
\Upsilon^{(c)}[k]\left(\theta_c^{*}\hat A\right)^{*}=\omega^{\frac12}\,k^{\theta_c}\left(\hat A_{j+1},\,\hat A_j+d_{\mathrm{s}}\hat a_j\right)u(\hat a_j)\,\omega^{\frac12}=\Upsilon^{(c)}\left[k^{\theta_c}\right](\hat A).
\]
By Lemma~\ref{lem:plane-reflection}, $\Theta_c\left(\Ins_c[\Upsilon]\right)(\hat A)=\Ins_c\left[\Upsilon\left(\theta_c^{*}\hat A\right)^{*}\right]$ for every matrix valued function $\Upsilon$;
so $\Theta_c\Ins_c\left[\Upsilon^{(c)}[k]\right]=\Ins_c\left[\Upsilon^{(c)}\left[k^{\theta_c}\right]\right]$.
The function $g_S^{-1}$ is real and, $g_1$ being even, an even function of $F^{\mathrm{el}}_j$;
by \eqref{eq:setting-reflection-coefficients} it is fixed by $\Theta_c$.
By Proposition~\ref{prop:symmetries}\eqref{it:action-invariance}, $\Theta_c$ maps the term of an oriented bond $(x,x')$ of $S_F$ to the term of $(\theta_cx',\theta_cx)$.
For a bond $\{x,\theta_cx\}$ that crosses $\{x_1=c\}$ the term of $(x,\theta_cx)$ therefore goes to the term of $(\theta_c\theta_cx,\theta_cx)=(x,\theta_cx)$,
and the term of $(\theta_cx,x)$ goes to itself as well.
Hence $\Theta_cS^{\times}_c(\Lambda)=S^{\times}_c(\Lambda)$,
and $\Theta_c\exp\left(S^{\times}_c(\Lambda)\right)=\exp\left(S^{\times}_c(\Lambda)\right)$, because $\Theta_c$ is an antilinear anti-automorphism and the exponential series has real coefficients.
The three images multiply to $G^{(c)}\left[k^{\theta_c}\right]$.
The values of $k^{\theta_c}$ are even, the adjoint of an even operator being even,
and $\|k(A,A')^{*}\|=\|k(A,A')\|$ obeys the Gaussian bound, which is symmetric in $A$ and $A'$.

\emph{Spatial mirrors.}
By \eqref{eq:setting-slice-reflections}, for $j'\in\Z$,
\[
\left(\theta_w^{*}\hat A\right)_{j'}=\theta_{\mathrm{s}}^{*}\hat A_{j'},\qquad\left(\theta_w^{*}\hat a\right)_{j'}=\hat a_{j'}\circ\theta_{\mathrm{s}},\qquad F^{\mathrm{el}}_j\left(\theta_w^{*}F\right)=\theta_{\mathrm{s}}^{*}F^{\mathrm{el}}_j(F),
\]
and $d_{\mathrm{s}}(\hat a_j\circ\theta_{\mathrm{s}})=\theta_{\mathrm{s}}^{*}d_{\mathrm{s}}\hat a_j$.
Hence
\[
\Upsilon^{(c)}[k]\left(\theta_w^{*}\hat A\right)=\omega^{\frac12}\,k\left(\theta_{\mathrm{s}}^{*}\hat A_{j+1},\,\theta_{\mathrm{s}}^{*}\left(\hat A_j+d_{\mathrm{s}}\hat a_j\right)\right)u\left(\hat a_j\circ\theta_{\mathrm{s}}\right)\omega^{\frac12},
\]
where $\hat A_{j+1}$ and $\hat A_j+d_{\mathrm{s}}\hat a_j$ are restricted to $\theta_{\mathrm{s}}S$, and $\hat a_j$ to $\theta_{\mathrm{s}}\Lambda$.
Lemma~\ref{lem:spatial-mirror-signs} gives, for every even matrix valued function $\Upsilon$,
\[
\Theta_w\left(\Ins_c[\Upsilon]\right)(\hat A)=\Ins_c\left[\mathfrak{U}\,\overline{\Upsilon\left(\theta_w^{*}\hat A\right)}\,\mathfrak{U}^{*}\right],
\]
together with $\mathfrak{U}\,\omega=\omega\,\mathfrak{U}$ and $\mathfrak{U}\,\overline{u(a\circ\theta_{\mathrm{s}})}\,\mathfrak{U}^{*}=u(a)$ for $a\in\R^{\theta_{\mathrm{s}}\Lambda}$.
Since $\omega$ is real, the insertion of $\mathfrak{U}^{*}\mathfrak{U}=1$ yields
\[
\begin{aligned}
\mathfrak{U}\,\overline{\Upsilon^{(c)}[k]\left(\theta_w^{*}\hat A\right)}\,\mathfrak{U}^{*}
&=\omega^{\frac12}\,\mathfrak{U}\,\overline{k\left(\theta_{\mathrm{s}}^{*}\hat A_{j+1},\,\theta_{\mathrm{s}}^{*}\left(\hat A_j+d_{\mathrm{s}}\hat a_j\right)\right)}\,\mathfrak{U}^{*}\,\mathfrak{U}\,\overline{u\left(\hat a_j\circ\theta_{\mathrm{s}}\right)}\,\mathfrak{U}^{*}\,\omega^{\frac12}\\
&=\omega^{\frac12}\,k^{\theta_w}\left(\hat A_{j+1},\,\hat A_j+d_{\mathrm{s}}\hat a_j\right)u(\hat a_j)\,\omega^{\frac12},
\end{aligned}
\]
and the right-hand side is $\Upsilon^{(c)}\left[k^{\theta_w}\right](\hat A)$ for the datum $\left(\theta_{\mathrm{s}}S,\theta_{\mathrm{s}}\Lambda,k^{\theta_w}\right)$.
Thus $\Theta_w\Ins_c\left[\Upsilon^{(c)}[k]\right]=\Ins_c\left[\Upsilon^{(c)}\left[k^{\theta_w}\right]\right]$.
Next, $g_1$ is real and even; by \eqref{eq:setting-reflection-coefficients},
\[
\Theta_w\,g_S^{-1}=\prod_{b\in S}g_1\left(\left(\theta_{\mathrm{s}}^{*}F^{\mathrm{el}}_j\right)(b)\right)^{-1}=\prod_{b\in S}g_1\left(F^{\mathrm{el}}_j(\theta_{\mathrm{s}}b)\right)^{-1}=\prod_{b\in\theta_{\mathrm{s}}S}g_1\left(F^{\mathrm{el}}_j(b)\right)^{-1}.
\]
By Proposition~\ref{prop:symmetries}\eqref{it:action-invariance}, $\Theta_w$ maps the term of an oriented bond $(x,x')$ to the term of $(\theta_wx',\theta_wx)$.
As $\theta_w$ commutes with $\theta_c$ and $\pi(\theta_wx)=\theta_w\pi(x)$ (Proposition~\ref{prop:geometry}\eqref{it:geometry-projection}),
the reflection $\theta_w$ maps the bonds $\{x,\theta_cx\}$ over a site $y$ that cross $\{x_1=c\}$ onto the bonds $\{\theta_wx,\theta_c\theta_wx\}$ over $\theta_{\mathrm{s}}y$.
Therefore $\Theta_wS^{\times}_c(\Lambda)=S^{\times}_c(\theta_{\mathrm{s}}\Lambda)$ and $\Theta_w\exp\left(S^{\times}_c(\Lambda)\right)=\exp\left(S^{\times}_c(\theta_{\mathrm{s}}\Lambda)\right)$,
and the product of the three images is $G^{(c)}\left[k^{\theta_w}\right]$, formed with $\left(\theta_{\mathrm{s}}S,\theta_{\mathrm{s}}\Lambda,k^{\theta_w}\right)$.

For every $d$ the unitary $\mathfrak{U}$ maps the span of the $e_I$ with $\#|I|=d$ into the span of the $e_{I'}$ with $\#|I'|=d$ (Lemma~\ref{lem:spatial-mirror-signs});
so the conjugation by $\mathfrak{U}$ preserves evenness, and so does the entrywise conjugation.
Both conjugations preserve the operator norm, and $\left(\theta_{\mathrm{s}}^{*}A'\right)(b)-\left(\theta_{\mathrm{s}}^{*}A\right)(b)=\pm\left(A'(\theta_{\mathrm{s}}b)-A(\theta_{\mathrm{s}}b)\right)$;
hence $\|k^{\theta_w}(A',A)\|\le C\cdot\prod_{b\in\theta_{\mathrm{s}}S}\exp\left(-\varepsilon_k\left(A'(b)-A(b)\right)^2\right)$.
For the covariance let $a\in\R^{\theta_{\mathrm{s}}\Lambda}$.
The relation $\mathfrak{U}\,\overline{u(a\circ\theta_{\mathrm{s}})}\,\mathfrak{U}^{*}=u(a)$ gives
$u(a)^{-1}\mathfrak{U}=\mathfrak{U}\,\overline{u(a\circ\theta_{\mathrm{s}})}^{\,-1}$ and $\mathfrak{U}^{*}u(a)=\overline{u(a\circ\theta_{\mathrm{s}})}\,\mathfrak{U}^{*}$.
Together with $\theta_{\mathrm{s}}^{*}d_{\mathrm{s}}a=d_{\mathrm{s}}(a\circ\theta_{\mathrm{s}})$ and the covariance of $k$ for $a\circ\theta_{\mathrm{s}}\in\R^\Lambda$ this yields
\[
\begin{aligned}
u(a)^{-1}\,k^{\theta_w}\left(A'+d_{\mathrm{s}}a,\,A+d_{\mathrm{s}}a\right)u(a)
&=\mathfrak{U}\,\overline{u(a\circ\theta_{\mathrm{s}})^{-1}\,k\left(\theta_{\mathrm{s}}^{*}A'+d_{\mathrm{s}}(a\circ\theta_{\mathrm{s}}),\,\theta_{\mathrm{s}}^{*}A+d_{\mathrm{s}}(a\circ\theta_{\mathrm{s}})\right)u(a\circ\theta_{\mathrm{s}})}\,\mathfrak{U}^{*}\\
&=k^{\theta_w}(A',A).
\end{aligned}
\]

\emph{Elements of the mid-plane algebra.}
An operator $T$ on $\mathcal{H}_S$ has the kernel $k$ if $(T\Phi)(A')=\int_{\R^S}\mathrm{d}A\,k(A',A)\,\Phi(A)$ (Definition~\ref{def:integral-kernels}).
Let $T$ be a bounded operator with the kernel $k$, and suppose that $\|k(A',A)\|\le m(A'-A)$ for an integrable function $m\ge0$ on $\R^S$.
For $\Phi,\Phi'\in\mathcal{H}_S$ the function $(A',A)\mapsto\|\Phi'(A')\|\cdot m(A'-A)\cdot\|\Phi(A)\|$ is then integrable on $\R^S\times\R^S$,
with integral at most $\|m\|_{L^1}\cdot\|\Phi'\|\cdot\|\Phi\|$
(Cauchy--Schwarz inequality for the factors $m(A'-A)^{\frac12}\cdot\|\Phi'(A')\|$ and $m(A'-A)^{\frac12}\cdot\|\Phi(A)\|$).
So Fubini's theorem allows the two integrations in $\langle\Phi',T\Phi\rangle$ to be interchanged,
and $T^{*}$ has the kernel $(A',A)\mapsto k(A,A')^{*}=k^{\theta_c}(A',A)$.
The operator $\mathfrak{T}_vT\mathfrak{T}_v^{-1}$ has the kernel $k^{\tau_v}$:
for $\Phi'\in\mathcal{H}_{S+v_{\mathrm{s}}}$ and $A\in\R^{S+v_{\mathrm{s}}}$ one has $\left(\mathfrak{T}_v^{-1}\Phi'\right)\left(\tau^{*}_{v_{\mathrm{s}}}A\right)=U_v^{*}\Phi'(A)$,
and the substitution $A''=\tau^{*}_{v_{\mathrm{s}}}A$, a measure preserving bijection of $\R^{S+v_{\mathrm{s}}}$ onto $\R^S$, gives
\[
\left(\mathfrak{T}_vT\mathfrak{T}_v^{-1}\Phi'\right)(A')
=U_v\int_{\R^S}\mathrm{d}A''\,k\left(\tau^{*}_{v_{\mathrm{s}}}A',A''\right)\left(\mathfrak{T}_v^{-1}\Phi'\right)(A'')
=\int_{\R^{S+v_{\mathrm{s}}}}\mathrm{d}A\,U_v\,k\left(\tau^{*}_{v_{\mathrm{s}}}A',\tau^{*}_{v_{\mathrm{s}}}A\right)U_v^{*}\,\Phi'(A).
\]
The operator $\mathfrak{J}T\mathfrak{J}^{-1}$ has the kernel $k^{\theta_w}$:
if $\Phi'=\mathfrak{J}\Phi$, then $\Phi\left(\theta_{\mathrm{s}}^{*}A\right)=\overline{\mathfrak{U}^{*}\Phi'(A)}$ for $A\in\R^{\theta_{\mathrm{s}}S}$,
and $\theta_{\mathrm{s}}^{*}$ is a measure preserving bijection of $\R^{\theta_{\mathrm{s}}S}$ onto $\R^S$; so
\[
\left(\mathfrak{J}T\mathfrak{J}^{-1}\Phi'\right)(A')
=\mathfrak{U}\,\overline{\int_{\R^{\theta_{\mathrm{s}}S}}\mathrm{d}A\,k\left(\theta_{\mathrm{s}}^{*}A',\theta_{\mathrm{s}}^{*}A\right)\overline{\mathfrak{U}^{*}\Phi'(A)}}
=\int_{\R^{\theta_{\mathrm{s}}S}}\mathrm{d}A\,\mathfrak{U}\,\overline{k\left(\theta_{\mathrm{s}}^{*}A',\theta_{\mathrm{s}}^{*}A\right)}\,\mathfrak{U}^{*}\,\Phi'(A).
\]
For $0<s\le\frac14$ the operator $\exp\left(-s|E_S|^2\right)$ has the kernel $\prod_{b\in S}g_{2s}\left(A'(b)-A(b)\right)$ times the identity of $\mathcal{F}_\Lambda$ (Definition~\ref{def:integral-kernels}),
a real function of the Euclidean norm $|A'-A|$.
The maps $\tau^{*}_{v_{\mathrm{s}}}$ and $\theta_{\mathrm{s}}^{*}$ preserve this norm, and $U_v$, $\mathfrak{U}$ commute with scalars;
the rules for $\mathfrak{T}_v$ and for $\mathfrak{J}$ therefore give
\begin{equation}\label{eq:amplification-gaussian-transport}
\mathfrak{T}_v\exp\left(-s|E_S|^2\right)\mathfrak{T}_v^{-1}=\exp\left(-s|E_{S+v_{\mathrm{s}}}|^2\right),\qquad
\mathfrak{J}\exp\left(-s|E_S|^2\right)\mathfrak{J}^{-1}=\exp\left(-s|E_{\theta_{\mathrm{s}}S}|^2\right).
\end{equation}

Consider now the generators of Definition~\ref{def:mid-plane}.
Since $\left(\mathfrak{T}_v^{-1}\Phi'\right)(A'')=U_v^{*}\Phi'\left(A''(\,\cdot-v_{\mathrm{s}})\right)$ for $A''\in\R^S$,
one has $\mathfrak{T}_vE_b\mathfrak{T}_v^{-1}=E_{b+v_{\mathrm{s}}}$ for $b\in S$.
So $\mathfrak{T}_v$ carries a function of $E_S$ to the same function of $(E_{b+v_{\mathrm{s}}})_{b\in S}$,
and the generators of type~\eqref{it:gaussian-generators} of $\mathfrak{W}(S,\Lambda)$ onto those of $\mathfrak{W}(S+v_{\mathrm{s}},\Lambda+v_{\mathrm{s}})$.
Moreover
\[
\mathfrak{T}_v\exp\left(i\langle\delta,A\rangle\right)\Upsilon\,\mathfrak{T}_v^{-1}=\exp\left(i\langle\delta(\,\cdot-v_{\mathrm{s}}),A\rangle\right)U_v\,\Upsilon\,U_v^{*};
\]
if $\Upsilon\in\End_z(\Lambda)$ (Definition~\ref{def:charged-operators}), then $U_v\Upsilon U_v^{*}\in\End_{z(\cdot-v_{\mathrm{s}})}(\Lambda+v_{\mathrm{s}})$ by the relation $U_v\,u\left(a(\,\cdot+v_{\mathrm{s}})\right)U_v^{*}=u(a)$,
and $d_{\mathrm{s}}^{*}\left(\delta(\,\cdot-v_{\mathrm{s}})\right)=\left(d_{\mathrm{s}}^{*}\delta\right)(\,\cdot-v_{\mathrm{s}})$.
So the generators of type~\eqref{it:exponential-generators} of $\mathfrak{W}(S,\Lambda)$ are carried onto those of $\mathfrak{W}(S+v_{\mathrm{s}},\Lambda+v_{\mathrm{s}})$.

For $\mathfrak{J}$, the chain rule applied to $A\mapsto\Phi\left(\theta_{\mathrm{s}}^{*}A\right)$ produces the factor $\operatorname{sgn}(\theta_{\mathrm{s}},b)$ in the derivative with respect to $A(\theta_{\mathrm{s}}b)$,
and the complex conjugation in $\mathfrak{J}$ changes the sign of $-i\,\partial$; hence
\[
\mathfrak{J}E_b\mathfrak{J}^{-1}=-\operatorname{sgn}(\theta_{\mathrm{s}},b)\cdot E_{\theta_{\mathrm{s}}b},\qquad b\in S.
\]
As $\mathfrak{J}$ is antilinear, it follows for $f\in\mathfrak{A}(S)$ (Definition~\ref{def:gaussian-symbols}) that
\[
\mathfrak{J}f(E_S)\mathfrak{J}^{-1}=f^{\theta_{\mathrm{s}}}\left(E_{\theta_{\mathrm{s}}S}\right),\qquad
f^{\theta_{\mathrm{s}}}(\xi):=\overline{f\left(\xi^{\theta_{\mathrm{s}}}\right)},\qquad
\left(\xi^{\theta_{\mathrm{s}}}\right)_b:=-\operatorname{sgn}(\theta_{\mathrm{s}},b)\cdot\xi_{\theta_{\mathrm{s}}b}\quad\left(b\in S,\ \xi\in\R^{\theta_{\mathrm{s}}S}\right).
\]
For a Gaussian symbol $f_{S',\mathfrak{Q},\lambda}$ with $\mathfrak{Q}(\xi)=\sum_{b,b'\in S'}\mathfrak{Q}_{bb'}\,\xi_b\,\xi_{b'}$ the reality of $\mathfrak{Q}$ gives
\[
\overline{f_{S',\mathfrak{Q},\lambda}\left(\xi^{\theta_{\mathrm{s}}}\right)}
=\exp\left(-\sum_{b,b'\in S'}\mathfrak{Q}_{bb'}\cdot\operatorname{sgn}(\theta_{\mathrm{s}},b)\cdot\operatorname{sgn}(\theta_{\mathrm{s}},b')\cdot\xi_{\theta_{\mathrm{s}}b}\,\xi_{\theta_{\mathrm{s}}b'}
-\sum_{b\in S'}\overline{\lambda_b}\cdot\operatorname{sgn}(\theta_{\mathrm{s}},b)\cdot\xi_{\theta_{\mathrm{s}}b}\right),
\]
hence
\[
\mathfrak{J}f_{S',\mathfrak{Q},\lambda}(E_S)\mathfrak{J}^{-1}=f_{\theta_{\mathrm{s}}S',\mathfrak{Q}',\lambda'}\left(E_{\theta_{\mathrm{s}}S}\right),\qquad
\mathfrak{Q}'_{\theta_{\mathrm{s}}b,\theta_{\mathrm{s}}b'}:=\operatorname{sgn}(\theta_{\mathrm{s}},b)\cdot\operatorname{sgn}(\theta_{\mathrm{s}},b')\cdot\mathfrak{Q}_{bb'},\qquad
\lambda'_{\theta_{\mathrm{s}}b}:=-\operatorname{sgn}(\theta_{\mathrm{s}},b)\cdot\overline{\lambda_b},
\]
and $\mathfrak{Q}'$ is again a real positive definite quadratic form.
Every Gaussian symbol on $\R^{\theta_{\mathrm{s}}S}$ arises in this way,
so the generators of type~\eqref{it:gaussian-generators} of $\mathfrak{W}(S,\Lambda)$ are carried onto those of $\mathfrak{W}(\theta_{\mathrm{s}}S,\theta_{\mathrm{s}}\Lambda)$.
Further, $\langle\delta,\theta_{\mathrm{s}}^{*}A\rangle=-\langle\delta^{\theta_{\mathrm{s}}},A\rangle$ for $\delta\in\R^S$ and $A\in\R^{\theta_{\mathrm{s}}S}$,
where $\delta^{\theta_{\mathrm{s}}}(\theta_{\mathrm{s}}b):=-\operatorname{sgn}(\theta_{\mathrm{s}},b)\cdot\delta(b)$; therefore
\[
\mathfrak{J}\exp\left(i\langle\delta,A\rangle\right)\Upsilon\,\mathfrak{J}^{-1}=\exp\left(i\langle\delta^{\theta_{\mathrm{s}}},A\rangle\right)\mathfrak{U}\,\overline{\Upsilon}\,\mathfrak{U}^{*}.
\]
If $\Upsilon\in\End_z(\Lambda)$ and $a\in\R^{\theta_{\mathrm{s}}\Lambda}$, the relation $\mathfrak{U}\,\overline{u(a\circ\theta_{\mathrm{s}})}\,\mathfrak{U}^{*}=u(a)$ gives
\[
u(a)\,\mathfrak{U}\,\overline{\Upsilon}\,\mathfrak{U}^{*}\,u(a)^{-1}
=\mathfrak{U}\,\overline{u(a\circ\theta_{\mathrm{s}})\,\Upsilon\,u(a\circ\theta_{\mathrm{s}})^{-1}}\,\mathfrak{U}^{*}
=\exp\left(iq\langle a,-z\circ\theta_{\mathrm{s}}\rangle\right)\cdot\mathfrak{U}\,\overline{\Upsilon}\,\mathfrak{U}^{*},
\]
so $\mathfrak{U}\overline{\Upsilon}\mathfrak{U}^{*}$ has the charge $-z\circ\theta_{\mathrm{s}}$.
On the other hand $\delta^{\theta_{\mathrm{s}}}=-\theta_{\mathrm{s}}^{*}\delta$, because $\operatorname{sgn}(\theta_{\mathrm{s}},\theta_{\mathrm{s}}b)=\operatorname{sgn}(\theta_{\mathrm{s}},b)$,
and $d_{\mathrm{s}}^{*}$ commutes with $\theta_{\mathrm{s}}^{*}$, since $d_{\mathrm{s}}$ does and $\theta_{\mathrm{s}}^{*}$ is orthogonal;
thus $d_{\mathrm{s}}^{*}\delta^{\theta_{\mathrm{s}}}=-\left(d_{\mathrm{s}}^{*}\delta\right)\circ\theta_{\mathrm{s}}$,
and $d_{\mathrm{s}}^{*}\delta=q\cdot z$ implies $d_{\mathrm{s}}^{*}\delta^{\theta_{\mathrm{s}}}=q\cdot(-z\circ\theta_{\mathrm{s}})$.
As the maps $\delta\mapsto\delta^{\theta_{\mathrm{s}}}$ and $\Upsilon\mapsto\mathfrak{U}\overline{\Upsilon}\mathfrak{U}^{*}$ are bijective,
the generators of type~\eqref{it:exponential-generators} of $\mathfrak{W}(S,\Lambda)$ are carried onto those of $\mathfrak{W}(\theta_{\mathrm{s}}S,\theta_{\mathrm{s}}\Lambda)$.

Conjugation by a unitary operator is an injective linear map that respects products and adjoints,
and conjugation by an antiunitary operator is an injective antilinear map with the same two properties.
Each of them maps the $*$-algebra generated by a set of operators onto the $*$-algebra generated by the image of this set.
Since the generators of $\mathfrak{W}(S,\Lambda)$ are carried onto those of $\mathfrak{W}(S+v_{\mathrm{s}},\Lambda+v_{\mathrm{s}})$ and of $\mathfrak{W}(\theta_{\mathrm{s}}S,\theta_{\mathrm{s}}\Lambda)$,
the maps $X\mapsto X^{\tau_v}$ and $X\mapsto X^{\theta_w}$ are bijections of the kinds stated;
the adjoint is a bijection of the $*$-algebra $\mathfrak{W}(S,\Lambda)$ onto itself.

Let $X\in\mathfrak{W}(S,\Lambda)$ and $0<s\le\frac14$.
By \eqref{eq:amplification-gaussian-transport},
\[
\mathfrak{T}_v\left(\exp\left(-s|E_S|^2\right)X\exp\left(-s|E_S|^2\right)\right)\mathfrak{T}_v^{-1}
=\exp\left(-s|E_{S+v_{\mathrm{s}}}|^2\right)X^{\tau_v}\exp\left(-s|E_{S+v_{\mathrm{s}}}|^2\right).
\]
The left side has the kernel $\left(k^{(s)}_X\right)^{\tau_v}$ by the rule for $\mathfrak{T}_v$,
and the right side has the kernel $k^{(s)}_{X^{\tau_v}}$ (Lemma~\ref{lem:regularised-kernels}, applied in $\mathfrak{W}(S+v_{\mathrm{s}},\Lambda+v_{\mathrm{s}})$).
Both kernels are continuous, and two continuous kernels of one operator coincide;
so $\left(k^{(s)}_X\right)^{\tau_v}=k^{(s)}_{X^{\tau_v}}$.
The adjoint of $\exp\left(-s|E_S|^2\right)\,X\exp\left(-s|E_S|^2\right)$ is $\exp\left(-s|E_S|^2\right)\,X^{*}\exp\left(-s|E_S|^2\right)$.
By the Gaussian bound of Lemma~\ref{lem:regularised-kernels}(c), the norm of $k^{(s)}_X(A',A)$ is bounded by an integrable function of $A'-A$;
so the rule for adjoints applies, and the adjoint has the continuous kernel $\left(k^{(s)}_X\right)^{\theta_c}$ besides the continuous kernel $k^{(s)}_{X^{*}}$.
Hence $\left(k^{(s)}_X\right)^{\theta_c}=k^{(s)}_{X^{*}}$.
By \eqref{eq:amplification-gaussian-transport} once more,
\[
\mathfrak{J}\left(\exp\left(-s|E_S|^2\right)X\exp\left(-s|E_S|^2\right)\right)\mathfrak{J}^{-1}
=\exp\left(-s|E_{\theta_{\mathrm{s}}S}|^2\right)X^{\theta_w}\exp\left(-s|E_{\theta_{\mathrm{s}}S}|^2\right),
\]
with the continuous kernels $\left(k^{(s)}_X\right)^{\theta_w}$ on the left and $k^{(s)}_{X^{\theta_w}}$ on the right; hence $\left(k^{(s)}_X\right)^{\theta_w}=k^{(s)}_{X^{\theta_w}}$.
Finally, $\|\cdot\|_\#$ is the operator norm on $\mathcal{H}_S$, on $\mathcal{H}_{S+v_{\mathrm{s}}}$ and on $\mathcal{H}_{\theta_{\mathrm{s}}S}$ (Lemma~\ref{lem:sharp-norm}),
and the operator norm is preserved by the conjugation with a unitary or an antiunitary operator and by the adjoint.

\emph{The group $W_{\mathrm{c}}$.}
An element $g\in W_{\mathrm{c}}$ acts by $x\mapsto\vec g\,x+g(0)$, with a diagonal matrix $\vec g$ whose entries are $\pm1$ and with $g(0)\in\Z^3$ (Definition~\ref{def:mirrors}).
The reflections $\theta_c$, $\theta_{x_2,0}$, $\theta_{x_3,0}$ belong to $W_{\mathrm{c}}$ and have the linear parts $\diag(-1,1,1)$, $\diag(1,-1,1)$, $\diag(1,1,-1)$;
here $\theta_c(x)=(2c-x_1,x_2,x_3)$, and $2c$ is an odd integer.
Let $p\in\{0,1\}^3$ be given by $\vec g=\diag\left((-1)^{p_1},(-1)^{p_2},(-1)^{p_3}\right)$.
Then $g\left(\theta_c^{p_1}\theta_{x_2,0}^{p_2}\theta_{x_3,0}^{p_3}\right)^{-1}$ is an element of $W_{\mathrm{c}}$ with the linear part $1$,
that is, a translation $\tau_v$ with $v\in\Z^3$, and $g=\tau_v\theta_c^{p_1}\theta_{x_2,0}^{p_2}\theta_{x_3,0}^{p_3}$.
Conversely, the linear part of $\tau_v\theta_c^{p_1}\theta_{x_2,0}^{p_2}\theta_{x_3,0}^{p_3}$ is $\diag\left((-1)^{p_1},(-1)^{p_2},(-1)^{p_3}\right)$;
so $p$ is determined by $\vec g$, and then $\tau_v=g\left(\theta_c^{p_1}\theta_{x_2,0}^{p_2}\theta_{x_3,0}^{p_3}\right)^{-1}$ is determined by $g$.
This proves the existence and the uniqueness of the normal word.

By \eqref{eq:setting-group-action}, the image of an even integrand under a product of elements of $W$ is obtained by applying the factors one after the other, the rightmost first:
\[
G^g=\left(\left(\left(G^{\theta_{x_3,0}^{p_3}}\right)^{\theta_{x_2,0}^{p_2}}\right)^{\theta_c^{p_1}}\right)^{\tau_v},
\]
and $G^{\theta_\Pi}=\Theta_\Pi G$ for every mirror $\Pi$.
The insertions are step elements: $G^{(c)}_X=G^{(c)}\left[k^{\left(\frac14\right)}_X\right]$ and $G'^{(c)}_X=G^{(c)}\left[k^{\left(\frac18\right)}_X\right]$ (Definition~\ref{def:two-insertions}),
and $\left(S,\Lambda,k^{(s)}_X\right)$ is a local step datum for $0<s\le\frac14$ (Lemma~\ref{lem:regularised-kernels}(e)).
Let $s\in\{\frac14,\frac18\}$ and put
\[
X_3:=X^{\theta_{x_3,0}^{p_3}},\qquad X_2:=X_3^{\theta_{x_2,0}^{p_2}},\qquad X_1:=X_2^{\theta_c^{p_1}},
\]
so that $X^g=X_1^{\tau_v}$.
If $p_3=1$, then \eqref{it:images-spatial-mirror} and \eqref{it:images-elements} for the mirror $(x_3,0)$ give
$\Theta_{x_3,0}G^{(c)}\left[k^{(s)}_X\right]=G^{(c)}\left[\left(k^{(s)}_X\right)^{\theta_w}\right]=G^{(c)}\left[k^{(s)}_{X_3}\right]$, at the same plane $c$.
If $p_2=1$, then \eqref{it:images-spatial-mirror} and \eqref{it:images-elements} for the mirror $(x_2,0)$ give
$\Theta_{x_2,0}G^{(c)}\left[k^{(s)}_{X_3}\right]=G^{(c)}\left[k^{(s)}_{X_2}\right]$, again at the plane $c$.
If $p_1=1$, then \eqref{it:images-time-reflection} and \eqref{it:images-elements} give
$\Theta_cG^{(c)}\left[k^{(s)}_{X_2}\right]=G^{(c)}\left[\left(k^{(s)}_{X_2}\right)^{\theta_c}\right]=G^{(c)}\left[k^{(s)}_{X_1}\right]$.
At last \eqref{it:images-translation} and \eqref{it:images-elements} give
$\left(G^{(c)}\left[k^{(s)}_{X_1}\right]\right)^{\tau_v}=G^{(c+v_1)}\left[k^{(s)}_{X^g}\right]$.
An exponent $0$ contributes the identity at its step.
Hence $\left(G^{(c)}\left[k^{(s)}_X\right]\right)^g=G^{(g(c))}\left[k^{(s)}_{X^g}\right]$ for $s=\frac14$ and for $s=\frac18$, which are the first two identities of \eqref{it:images-group};
and \eqref{it:images-elements}, applied at each step, gives $X^g\in\mathfrak{W}(gS,g\Lambda)$ and $\|X^g\|_\#=\|X\|_\#$.
For the mirror $(x_1,c')$, Definition~\ref{def:mirrors}\eqref{it:mirror-product} gives $\theta_{c'}\theta_c=\tau_v$ with $v=2(c'-c)\,e_1$;
so $\theta_{c'}=\tau_v\theta_c$ is the normal word of $\theta_{c'}$, with $p=(1,0,0)$.
Then $X^{\theta_{c'}}=\left(X^{*}\right)^{\tau_v}$, the plane is $\theta_{c'}(c)=c+2(c'-c)=2c'-c$, and $G^{\theta_{c'}}=\Theta_{c'}G$;
this is the last identity.

There remains the reflection $\theta_w$ of a spatial mirror $\Pi_w=(x_\nu,w)$.
Put $v:=2w\,e_\nu\in\Z^3$, so that $v_1=0$.
By Definition~\ref{def:mirrors}\eqref{it:mirror-product} and~\eqref{it:mirror-involution}, $\theta_w\theta_{x_\nu,0}=\tau_v$,
hence $\theta_w=\tau_v\theta_{x_\nu,0}$, and this is the normal word of $\theta_w$.
Let $\Theta_{x_\nu,0}$ be the antilinear reflection attached to the mirror $(x_\nu,0)$,
and let $\theta'_{\mathrm{s}}$, $\mathfrak{u}'$, $\mathfrak{U}'$ and $\mathfrak{J}'$ be the objects of Lemma~\ref{lem:spatial-mirror-signs} and of \eqref{it:images-elements} that belong to this mirror in the place of $\Pi_w$.
On the last two coordinates the identity $\theta_w=\tau_v\theta_{x_\nu,0}$ reads $\theta_{\mathrm{s}}y=\theta'_{\mathrm{s}}y+v_{\mathrm{s}}$ for the sites $y$ of $\Z^2$;
so $\theta_{\mathrm{s}}b=\theta'_{\mathrm{s}}b+v_{\mathrm{s}}$ and $\operatorname{sgn}(\theta_{\mathrm{s}},b)=\operatorname{sgn}(\theta'_{\mathrm{s}},b)$ for every link $b$,
and the definition of the pull-backs gives
\[
\theta_{\mathrm{s}}^{*}A=\left(\theta'_{\mathrm{s}}\right)^{*}\tau^{*}_{v_{\mathrm{s}}}A,\qquad A\in\R^{\theta_{\mathrm{s}}S},
\]
with $\tau^{*}_{v_{\mathrm{s}}}\colon\R^{\theta_{\mathrm{s}}S}\to\R^{\theta'_{\mathrm{s}}S}$ as in \eqref{it:images-translation}.
By Lemma~\ref{lem:sign-rules}(d) the linear automorphism $\Theta_w\Theta_{x_\nu,0}$ of $\Gr(\mathbb{X})$ maps $\psi(x)$ to $\psi(x+v)$ and $\bar\psi(x)$ to $\bar\psi(x+v)$;
it is therefore $\alpha_v$, and since $\Theta_{x_\nu,0}^2=1$ (Lemma~\ref{lem:sign-rules}(a)), $\Theta_w=\alpha_v\Theta_{x_\nu,0}$.
Let $I=\{i_1<\dots<i_r\}\subseteq\mathcal{J}^{\mathrm{l}}(\Lambda)$.
As $\alpha_v$ is multiplicative, Lemma~\ref{lem:spatial-mirror-signs}(b), applied to both mirrors, and Lemma~\ref{lem:translated-insertions}(c) give
\[
\sum_{I''}(\wedge\mathfrak{u})_{I''I}\cdot\zeta_{I''}
=\Theta_w(\zeta_{i_1})\cdots\Theta_w(\zeta_{i_r})
=\alpha_v\left(\Theta_{x_\nu,0}(\zeta_{i_1})\cdots\Theta_{x_\nu,0}(\zeta_{i_r})\right)
=\sum_{I'}(\wedge\mathfrak{u}')_{I'I}\cdot\operatorname{sgn}_v(I')\cdot\zeta_{I'+v},
\]
where $I''$ ranges over the subsets of $\mathcal{J}^{\mathrm{l}}(\theta_{\mathrm{s}}\Lambda)$ and $I'$ over those of $\mathcal{J}^{\mathrm{l}}(\theta'_{\mathrm{s}}\Lambda)$.
The map $I'\mapsto I'+v$ is a bijection of the second family onto the first,
and the products $\zeta_{I''}$ are linearly independent, the modes being a basis of the span of the generators at the upper ends (Definition~\ref{def:modes}).
Hence $(\wedge\mathfrak{u})_{I'+v,I}=\operatorname{sgn}_v(I')\cdot(\wedge\mathfrak{u}')_{I'I}$,
that is, $\wedge\mathfrak{u}=U_v\,(\wedge\mathfrak{u}')$ with the unitary $U_v\colon\mathcal{F}_{\theta'_{\mathrm{s}}\Lambda}\to\mathcal{F}_{\theta_{\mathrm{s}}\Lambda}$ of Lemma~\ref{lem:translated-insertions}(d).
The diagonal operator $\mathfrak{r}$ of Lemma~\ref{lem:spatial-mirror-signs}(d) acts on $\mathcal{F}_\Lambda$ and is the same for both mirrors;
multiplying by it from the right one obtains $\mathfrak{U}=U_v\,\mathfrak{U}'$.
For $\Phi\in\mathcal{H}_S$ and $A\in\R^{\theta_{\mathrm{s}}S}$ it follows that
\[
(\mathfrak{J}\Phi)(A)
=U_v\,\mathfrak{U}'\,\overline{\Phi\left(\left(\theta'_{\mathrm{s}}\right)^{*}\tau^{*}_{v_{\mathrm{s}}}A\right)}
=U_v\left(\mathfrak{J}'\Phi\right)\left(\tau^{*}_{v_{\mathrm{s}}}A\right)
=\left(\mathfrak{T}_v\mathfrak{J}'\Phi\right)(A),
\]
with $\mathfrak{T}_v\colon\mathcal{H}_{\theta'_{\mathrm{s}}S}\to\mathcal{H}_{\theta_{\mathrm{s}}S}$ as in \eqref{it:images-elements}.
Thus $\mathfrak{J}=\mathfrak{T}_v\mathfrak{J}'$, and
\[
\mathfrak{J}X\mathfrak{J}^{-1}=\mathfrak{T}_v\left(\mathfrak{J}'X\left(\mathfrak{J}'\right)^{-1}\right)\mathfrak{T}_v^{-1}=\left(X^{\theta_{x_\nu,0}}\right)^{\tau_v},
\]
and the right side is $X^g$ for $g=\theta_w$, by the definition of $X^g$ through the normal word $\tau_v\theta_{x_\nu,0}$.

\emph{The sense of the identities.}
By Proposition~\ref{prop:word-formula}\eqref{it:step-class} a step element $G^{(c)}[k]$ lies in $\mathcal{G}^{\mathrm{ev}}(B(c,S,\Lambda))$,
and $B(c,S,\Lambda)\subseteq B^{\#}(c,S,\Lambda)$ by Lemma~\ref{lem:covariant-box}(b);
so $G^{(c)}[k]$ is an element of $\mathcal{G}^{\mathrm{ev}}\left(B^{\#}(c,S,\Lambda)\right)$ (Lemma~\ref{lem:equal-integrands}).
By its definition it reads only links contained in $\mathcal{K}(c,S,\Lambda)$ and fields at fermion sites $x$ with $\pi(x)+\left[-\frac12,\frac12\right]^3\subseteq\mathcal{K}(c,S,\Lambda)$.
For $g\in W_{\mathrm{c}}$ the integrand $\left(G^{(c)}[k]\right)^g$ reads only the images under $g$ of these links and sites, by \eqref{eq:setting-reflection-fields} and \eqref{eq:setting-reflection-coefficients},
and it lies in $\mathcal{G}^{\mathrm{ev}}\left(gB^{\#}(c,S,\Lambda)\right)$ by Lemma~\ref{lem:reflected-integrands};
moreover $g\mathcal{K}(c,S,\Lambda)=\mathcal{K}(g(c),gS,g\Lambda)$ and $gB^{\#}(c,S,\Lambda)=B^{\#}(g(c),gS,g\Lambda)$ by Lemma~\ref{lem:covariant-box}(a).
In each identity of \eqref{it:images-translation}, \eqref{it:images-time-reflection}, \eqref{it:images-spatial-mirror} and \eqref{it:images-group}
the left side is the image of a step element under an element of $W_{\mathrm{c}}$, namely $\tau_v$, $\theta_c$, $\theta_w$ or $g$,
and the right side is a step element formed with the plane and the sets named in the statement.
Both sides therefore read only links and sites of the same set $\mathcal{K}$ and lie in $\mathcal{G}^{\mathrm{ev}}$ of the same box $B^{\#}$;
the computations above, made with the link variables of this box, show that they coincide as elements of $\mathcal{G}^{\mathrm{ev}}\left(B^{\#}\right)$,
which is equality in the sense of Lemma~\ref{lem:equal-integrands}.

The boxes of least radius of the two sides can be different.
The set $\mathcal{K}(c,S,\Lambda)$ extends from $j-\frac12$ to $j+\frac32$ in the direction $e_1$,
so the first coordinate of the centre of $B(c,S,\Lambda)$ is $j+\frac32-r_0$, which differs from $c$ when the radius $r_0$ of $B(c,S,\Lambda)$ exceeds $1$.
In \eqref{it:images-time-reflection} the left side lies in $\mathcal{G}^{\mathrm{ev}}(\theta_cB(c,S,\Lambda))$ by Lemma~\ref{lem:reflected-integrands} and the right side in $\mathcal{G}^{\mathrm{ev}}(B(c,S,\Lambda))$,
and $\theta_cB(c,S,\Lambda)\neq B(c,S,\Lambda)$ as soon as $r_0>1$.
In \eqref{it:images-translation}, on the other hand, the two boxes coincide:
a translation by a vector of $\Z^3$ preserves radii and the lexicographic order of centres, so that $B(c+v_1,S+v_{\mathrm{s}},\Lambda+v_{\mathrm{s}})=B(c,S,\Lambda)+v$.
\end{proof}

\begin{lemma}[Injectivity]\label{lem:injective-insertions}
Let $-1<\rho<1$, let $c=j+\frac12$ with $j\in\Z$, and let $(S,\Lambda,k)$ be a local step datum.
If $G^{(c)}[k]=0$, then $k=0$.
Consequently the maps $X\mapsto G^{(c)}_X$ and $X\mapsto G'^{(c)}_X$ are injective on $\mathfrak{W}(S,\Lambda)$.
\end{lemma}

\begin{proof}
The element $\exp\left(S^{\times}_c(\Lambda)\right)$ is invertible, and $g_S^{-1}$ is a positive function.
Hence $G^{(c)}[k]=\Ins_c\left[\Upsilon^{(c)}[k]\right]\cdot g_S^{-1}\cdot\exp\left(S^{\times}_c(\Lambda)\right)=0$ implies $\Ins_c\left[\Upsilon^{(c)}[k](\hat A)\right]=0$ for every $\hat A$.
The modes $\zeta_i$, $i\in\mathcal{J}^{\mathrm{l}}(\Lambda)$, form a basis of the generators at the upper ends of the bonds over $\Lambda$ that cross $\{x_1=c\}$,
and the $\Theta_c(\zeta_i)$ form a basis of the generators at their lower ends (Definition~\ref{def:modes}).
So the products $\Theta_c(\zeta_J)\,\zeta_I$, $I,J\subseteq\mathcal{J}^{\mathrm{l}}(\Lambda)$, are linearly independent, and $\Upsilon^{(c)}[k](\hat A)=0$.
For $-1<\rho<1$ the crossing coefficients $\omega_i$ are positive, so that $\omega$ is invertible, and $u(\hat a_j)$ is unitary.
It follows that $k\left(\hat A_{j+1}|_S,\,(\hat A_j+d_{\mathrm{s}}\hat a_j)|_S\right)=0$ for every $\hat A$.
With $\hat a_j=0$ and arbitrary values of $\hat A_{j+1}$ and $\hat A_j$ on $S$, these pairs of arguments exhaust $\R^S\times\R^S$; thus $k=0$.

Let now $X\in\mathfrak{W}(S,\Lambda)$ with $G^{(c)}_X=0$ or with $G'^{(c)}_X=0$.
These integrands are $G^{(c)}\left[k^{(s)}_X\right]$ with $s=\frac14$, respectively $s=\frac18$ (Definition~\ref{def:two-insertions}),
and $\left(S,\Lambda,k^{(s)}_X\right)$ is a local step datum (Lemma~\ref{lem:regularised-kernels}(e)).
By the first part $k^{(s)}_X=0$, that is, $\exp\left(-s|E_S|^2\right)X\exp\left(-s|E_S|^2\right)=0$.
The operator $\exp\left(-s|E_S|^2\right)$ is self-adjoint and injective, so that its range is dense.
For every $\Phi\in\mathcal{H}_S$ the vector $X\exp\left(-s|E_S|^2\right)\Phi$ lies in the kernel of $\exp\left(-s|E_S|^2\right)$ and therefore vanishes;
thus the bounded operator $X$ vanishes on a dense subspace, and $X=0$.
As both maps are linear (Lemma~\ref{lem:regularised-kernels}(a)), they are injective.
\end{proof}

\subsection{H\"older's inequality with insertions}\label{sec:holder-insertions}

Let $\mathcal{H}$ be a separable Hilbert space,
and for $1\le p<\infty$ let $\mathscr{J}_p(\mathcal{H})$ be its Schatten class of exponent $p$,
with the norm $\|T\|_p=\left(\Tr|T|^p\right)^{\frac{1}{p}}$,
as in Lemma~\ref{lem:schatten-traces}.
Let $R$ and $R'$ be bounded operators on $\mathcal{H}$ and let $T\in\mathscr{J}_p(\mathcal{H})$.
Arranged in decreasing order and counted with multiplicity,
the singular values of $RTR'$ are bounded, one by one,
by $\|R\|\cdot\|R'\|$ times the singular values of $T$ \cite[Theorem~1.6]{Simon2005}.
Raising them to the power $p$ and summing, one finds $RTR'\in\mathscr{J}_p(\mathcal{H})$ and
\begin{equation}\label{eq:amplification-ideal}
\|RTR'\|_p\le\|R\|\cdot\|T\|_p\cdot\|R'\|.
\end{equation}
In Lemma~\ref{lem:schatten-traces} and in \eqref{eq:amplification-ideal} every operator maps $\mathcal{H}$ into itself.
An operator from one Hilbert space into another is turned into an operator on the direct sum of the two spaces as follows.

\begin{lemma}[Operators between two spaces]\label{lem:two-spaces}
Let $\mathcal{H}_1$ and $\mathcal{H}_2$ be separable Hilbert spaces,
let $\mathcal{H}:=\mathcal{H}_1\oplus\mathcal{H}_2$,
and let $J_a\colon\mathcal{H}_a\to\mathcal{H}$, $a=1,2$, be the canonical isometries,
so that $J_a^*J_a=1$, and $J_a^*J_b=0$ for $a\neq b$.
For $a,b\in\{1,2\}$ and a bounded operator $T\colon\mathcal{H}_a\to\mathcal{H}_b$ put
\[
\widehat{T}:=J_b\,T\,J_a^*,
\]
a bounded operator on $\mathcal{H}$.
\begin{enumerate}
\item[(i)] If $T\colon\mathcal{H}_a\to\mathcal{H}_b$ and $T'\colon\mathcal{H}_b\to\mathcal{H}_{b'}$ are bounded, $a,b,b'\in\{1,2\}$,
then $\widehat{T'T}=\widehat{T'}\,\widehat{T}$.
Every bounded $T\colon\mathcal{H}_a\to\mathcal{H}_b$ satisfies $\widehat{T^*}=\widehat{T}^{\,*}$ and $\|\widehat{T}\|=\|T\|$.
\item[(ii)] For every bounded $T\colon\mathcal{H}_a\to\mathcal{H}_b$ and every integer $r\ge1$,
\[
\left(\widehat{T}^{\,*}\widehat{T}\right)^r=J_a\left(T^*T\right)^rJ_a^*
\qquad\text{and}\qquad
\Tr_{\mathcal{H}}\left(\widehat{T}^{\,*}\widehat{T}\right)^r=\Tr_{\mathcal{H}_a}\left(T^*T\right)^r,
\]
where both traces are numbers in $[0,\infty]$.
\item[(iii)] If $T\colon\mathcal{H}_a\to\mathcal{H}_a$ is of trace class,
then $\widehat{T}$ is of trace class and $\Tr_{\mathcal{H}}\widehat{T}=\Tr_{\mathcal{H}_a}T$.
\end{enumerate}
\end{lemma}

\begin{proof}
(i) The relation $J_b^*J_b=1$ gives
\[
\widehat{T'}\,\widehat{T}=J_{b'}\,T'\,J_b^*J_b\,T\,J_a^*=J_{b'}\,T'T\,J_a^*=\widehat{T'T},
\]
and the adjoint of $J_b\,T\,J_a^*$ is $J_a\,T^*J_b^*=\widehat{T^*}$.
Since $J_b$ is an isometry, $\|\widehat{T}v\|=\|T\,J_a^*v\|$ for every $v\in\mathcal{H}$.
The contraction $J_a^*$ maps the closed unit ball of $\mathcal{H}$ onto the closed unit ball of $\mathcal{H}_a$,
because $J_a^*J_aw=w$ and $\|J_aw\|=\|w\|$ for $w\in\mathcal{H}_a$.
Taking the supremum over the unit ball of $\mathcal{H}$ therefore yields $\|\widehat{T}\|=\|T\|$.

(ii) By (i), $\widehat{T}^{\,*}\widehat{T}=\widehat{T^*T}=J_a\,T^*T\,J_a^*$.
From $J_a^*J_a=1$ one gets $\left(J_a\,R\,J_a^*\right)^r=J_a\,R^r\,J_a^*$ for every bounded operator $R$ on $\mathcal{H}_a$,
and $R=T^*T$ gives the first identity.
For the second, fix orthonormal bases of $\mathcal{H}_1$ and of $\mathcal{H}_2$;
their images under $J_1$ and $J_2$ together form an orthonormal basis of $\mathcal{H}$.
Let $R$ be a bounded operator on $\mathcal{H}_a$, and let $a'$ be the element of $\{1,2\}$ different from $a$.
If $v$ belongs to the basis of $\mathcal{H}_{a'}$, then $J_a^*J_{a'}v=0$,
hence $\langle J_{a'}v,J_a\,R\,J_a^*J_{a'}v\rangle=0$;
if $v$ belongs to the basis of $\mathcal{H}_a$, then $\langle J_av,J_a\,R\,J_a^*J_av\rangle=\langle v,Rv\rangle$.
For $R\ge0$ summation over the basis of $\mathcal{H}$ gives
\[
\Tr_{\mathcal{H}}J_a\,R\,J_a^*=\Tr_{\mathcal{H}_a}R\in[0,\infty],
\]
and the choice $R=\left(T^*T\right)^r$ gives the second identity.

(iii) The operator $J_a|T|J_a^*$ is non-negative,
and by the computation in (ii) its square is $J_a|T|^2J_a^*=J_a\,T^*T\,J_a^*=\widehat{T}^{\,*}\widehat{T}$.
By the uniqueness of non-negative square roots, $|\widehat{T}|=J_a|T|J_a^*$,
and the trace identity of (ii) with $R=|T|$ gives $\Tr_{\mathcal{H}}|\widehat{T}|=\Tr_{\mathcal{H}_a}|T|<\infty$.
So $\widehat{T}$ is of trace class.
Its trace may be computed in the basis of (ii);
the diagonal entries found there, now for $R=T$, add up to $\Tr_{\mathcal{H}_a}T$.
\end{proof}

Let $T\colon\mathcal{H}_a\to\mathcal{H}_b$ be bounded.
The non-negative operator $J_a|T|J_a^*$ has the square $J_a\,T^*T\,J_a^*=\widehat{T}^{\,*}\widehat{T}$ (Lemma~\ref{lem:two-spaces}(ii)),
so that $|\widehat{T}|=J_a|T|J_a^*$.
An eigenvector of $|\widehat{T}|$ for a non-zero eigenvalue lies in the range of $|\widehat{T}|$, hence in the range of $J_a$;
it follows that $J_a$ maps each eigenspace of $|T|$ for a non-zero eigenvalue isometrically onto the eigenspace of $|\widehat{T}|$ for the same eigenvalue.
If $T$ is compact, then so is $\widehat{T}$,
and the non-zero singular values of the two operators coincide, counted with multiplicity.
Let $p\ge2$ be an even integer.
For a bounded operator $T\colon\mathcal{H}_a\to\mathcal{H}_b$ put
\[
\|T\|_p:=\left(\Tr_{\mathcal{H}_a}\left(T^*T\right)^{\frac{p}{2}}\right)^{\frac{1}{p}}\in[0,\infty].
\]
Since $|\widehat{T}|^p=\left(\widehat{T}^{\,*}\widehat{T}\right)^{\frac{p}{2}}$ for even $p$, as in Lemma~\ref{lem:schatten-traces}(b),
Lemma~\ref{lem:two-spaces}(ii) gives $\Tr_{\mathcal{H}}|\widehat{T}|^p=\|T\|_p^p$.
Thus $\widehat{T}$ lies in $\mathscr{J}_p(\mathcal{H})$ exactly when $\|T\|_p<\infty$,
and then $\|\widehat{T}\|_p=\|T\|_p$.
For a Hilbert--Schmidt operator $T$ the number $\|T\|_p$ is the norm of Definition~\ref{def:even-schatten}.

\begin{proposition}[H\"older's inequality with insertions]\label{prop:holder-insertions}
Let $-1\le\rho\le1$, $w_0,\dots,w_3>0$, $M\in\R$ and $q>0$.
Let $n=(n_1,n_2,n_3)$ be a triple of integers $n_\mu\ge4$,
and consider the coordinate torus $C(n)=C(n_1,n_2,n_3)$ with the time direction $e_1$;
its cross-section is rectangular, with the slice torus $\Lambda_{\mathrm{s}}=\Z^2/(n_2\Z\times n_3\Z)$.
Let $\sigma_0$ be any plaquette of $\Lambda_{\mathrm{s}}$,
and let $\mathcal{H}^{\mathrm{l}}$, $\mathcal{H}^{\mathrm{s}}$, $K_{\beta,S}$ and $\mathfrak{M}$ be the spaces and operators of Definition~\ref{def:slice-operators}
for this cross-section, the operator $\mathfrak{M}$ being formed with the plaquette $\sigma_0$.
For $j\in\Z/n_1\Z$ let $\beta_j\in\R^{\Lambda^{(1)}_{\mathrm{s}}}$,
let $S_j$ be a subset of $\Lambda^{(1)}_{\mathrm{s}}$,
let $u_j\colon\R^{\Lambda^{(2)}_{\mathrm{s}}}\to\C$ be a Borel function with $|u_j|\le1$,
acting on $\mathcal{H}^{\mathrm{s}}=\bigoplus_{k\in\Z}L^2\left(\R^{\Lambda^{(1)}_{\mathrm{s}}}/\Gamma_{\mathrm{s}};\mathcal{F}'\right)$
as the multiplication by $A\mapsto u_j\left(d_{\mathrm{s}}A+(2\pi/q)\cdot k\cdot1_{\sigma_0}\right)$ on the summand with index $k$,
and let $\mathcal{B}_j$ be a bounded operator on $\mathcal{H}^{\mathrm{l}}$.
Then the product
\[
\prod_{j=n_1-1,\dots,0}\left(K^{1/2}_{\beta_j,S_j}\,\mathcal{B}_j\,K^{1/2}_{\beta_j,S_j}\,\mathfrak{M}\,u_j\,\mathfrak{M}^*\right),
\]
whose factors stand from left to right in the order $j=n_1-1,\dots,1,0$,
is an operator of trace class on $\mathcal{H}^{\mathrm{l}}$, and
\[
\begin{aligned}
&\left|\Tr_{\mathcal{H}^{\mathrm{l}}}\prod_{j=n_1-1,\dots,0}\left(K^{1/2}_{\beta_j,S_j}\,\mathcal{B}_j\,K^{1/2}_{\beta_j,S_j}\,\mathfrak{M}\,u_j\,\mathfrak{M}^*\right)\right|\\
&\qquad\le\prod_{j\in\Z/n_1\Z}\left(\|\mathcal{B}_j\|\cdot2^{\frac{\#|S_j|}{2}}\right)\cdot(2\pi)^{-|n|}\cdot Z_{C(n)}[1].
\end{aligned}
\]
\end{proposition}

\begin{proof}
The projection $P_{\mathrm{G}}$ commutes with every function of $E$:
by Definition~\ref{def:slice-operators}, $(2\pi/q)^{|\Lambda_{\mathrm{s}}|}\cdot P_{\mathrm{G}}$ is the integral of the operators $U_{\mathrm{G}}(a)$ over $a\in[0,2\pi/q)^{\Lambda_{\mathrm{s}}}$,
and each $U_{\mathrm{G}}(a)$ commutes with every function of $E$, since it acts on the variable $A$ by a translation and on the Fock part by a constant matrix.
So, for $\beta\in\R^{\Lambda^{(1)}_{\mathrm{s}}}$ and $S\subset\Lambda^{(1)}_{\mathrm{s}}$,
Definition~\ref{def:slice-operators} gives
\[
K_{\beta,S}=(2\pi/q)^{|\Lambda_{\mathrm{s}}|}\cdot\exp\left(-\frac{1}{2}|E-\beta|^2+\frac{1}{4}|E-\beta|_S^2\right)P_{\mathrm{G}},
\]
where $|\cdot|_S$ is the norm of the components in $S$.
The exponent is at most zero, so the function of $E$ in this formula is bounded and non-negative.
Hence $K_{\beta,S}$ is a bounded non-negative operator, with the square root
\[
K^{1/2}_{\beta,S}=(2\pi/q)^{\frac{|\Lambda_{\mathrm{s}}|}{2}}\cdot\exp\left(-\frac{1}{4}|E-\beta|^2+\frac{1}{8}|E-\beta|_S^2\right)P_{\mathrm{G}}.
\]
The operator $\mathfrak{M}\colon\mathcal{H}^{\mathrm{s}}\to\mathcal{H}^{\mathrm{l}}$ is bounded,
because the half-slab matrices are bounded functions of $A$ (Lemma~\ref{lem:half-slab-matrices}).
By Corollary~\ref{cor:trace-class} the operators
\[
T_j:=K^{1/2}_{\beta_j,S_j}\,\mathfrak{M}\colon\mathcal{H}^{\mathrm{s}}\to\mathcal{H}^{\mathrm{l}},\qquad j\in\Z/n_1\Z,
\]
are Hilbert--Schmidt operators.
In the notation of Lemma~\ref{lem:two-spaces} let $\mathcal{H}_1:=\mathcal{H}^{\mathrm{l}}$ and $\mathcal{H}_2:=\mathcal{H}^{\mathrm{s}}$,
both of which are separable, and $\mathcal{H}:=\mathcal{H}^{\mathrm{l}}\oplus\mathcal{H}^{\mathrm{s}}$.

\emph{Traces of powers.}
Let $\beta\in\R^{\Lambda^{(1)}_{\mathrm{s}}}$, let $S\subset\Lambda^{(1)}_{\mathrm{s}}$, and put $T:=K^{1/2}_{\beta,S}\,\mathfrak{M}$,
a Hilbert--Schmidt operator from $\mathcal{H}^{\mathrm{s}}$ to $\mathcal{H}^{\mathrm{l}}$ by Corollary~\ref{cor:trace-class}.
Then $T^*T$ and $TT^*$, each a product of two Hilbert--Schmidt operators, are of trace class on $\mathcal{H}^{\mathrm{s}}$ and on $\mathcal{H}^{\mathrm{l}}$ respectively,
and so are all their powers with positive integer exponents.
With $F^{\mathrm{el}}_j$ the electric part of the curvature at the time $j$ (Definition~\ref{def:slices}), one has
\begin{equation}\label{eq:amplification-trace-power}
\begin{split}
\Tr_{\mathcal{H}^{\mathrm{s}}}\left(T^*T\right)^{n_1}
&=\Tr_{\mathcal{H}^{\mathrm{l}}}\left(K_{\beta,S}\,\mathfrak{M}\mathfrak{M}^*\right)^{n_1}\\
&=(2\pi)^{-|n|}\cdot Z_{C(n)}\left[\prod_{j\in\Z/n_1\Z}\left(\exp\left(i\langle\beta,F^{\mathrm{el}}_j\rangle\right)\cdot\prod_{b\in S}\left(\sqrt{2}\cdot\exp\left(-\frac{1}{2}F^{\mathrm{el}}_j(b)^2\right)\right)\right)\right]\\
&\le2^{\frac{\#|S|\cdot n_1}{2}}\cdot(2\pi)^{-|n|}\cdot Z_{C(n)}[1].
\end{split}
\end{equation}
The first equality combines two identities between traces.
On the one hand, $\Tr_{\mathcal{H}^{\mathrm{s}}}\left(T^*T\right)^{n_1}=\Tr_{\mathcal{H}^{\mathrm{l}}}\left(TT^*\right)^{n_1}$.
Indeed, $\widehat{T}$ is a Hilbert--Schmidt operator on $\mathcal{H}$,
since $\Tr_{\mathcal{H}}\widehat{T}^{\,*}\widehat{T}=\Tr_{\mathcal{H}^{\mathrm{s}}}T^*T<\infty$ by Lemma~\ref{lem:two-spaces}(ii) with $r=1$.
As $\left(T^*T\right)^{n_1}$ and $\left(TT^*\right)^{n_1}$ are of trace class,
Lemma~\ref{lem:two-spaces}(i) and (iii) identify their traces with the traces on $\mathcal{H}$ of
\[
\left(\widehat{T}^{\,*}\widehat{T}\right)^{n_1}=\widehat{T}^{\,*}\left(\widehat{T}\left(\widehat{T}^{\,*}\widehat{T}\right)^{n_1-1}\right)
\qquad\text{and}\qquad
\left(\widehat{T}\,\widehat{T}^{\,*}\right)^{n_1}=\left(\widehat{T}\left(\widehat{T}^{\,*}\widehat{T}\right)^{n_1-1}\right)\widehat{T}^{\,*}.
\]
Here $\widehat{T}^{\,*}$ and $\widehat{T}\left(\widehat{T}^{\,*}\widehat{T}\right)^{n_1-1}$ are Hilbert--Schmidt operators on $\mathcal{H}$, as $\widehat{T}$ is,
and for two Hilbert--Schmidt operators on one Hilbert space the trace of their product does not change when they are interchanged.
On the other hand, $\Tr_{\mathcal{H}^{\mathrm{l}}}\left(TT^*\right)^{n_1}=\Tr_{\mathcal{H}^{\mathrm{l}}}\left(K_{\beta,S}\,\mathfrak{M}\mathfrak{M}^*\right)^{n_1}$.
Since $K^{1/2}_{\beta,S}$ is self-adjoint, $T^*=\mathfrak{M}^*K^{1/2}_{\beta,S}$.
Put $R:=\left(TT^*\right)^{n_1-1}T\,\mathfrak{M}^*$;
it is an operator of trace class on $\mathcal{H}^{\mathrm{l}}$, because $\left(TT^*\right)^{n_1-1}$ is of trace class ($n_1-1\ge1$) and $T\,\mathfrak{M}^*$ is bounded.
One has
\[
\left(TT^*\right)^{n_1}=R\,K^{1/2}_{\beta,S},
\qquad
K^{1/2}_{\beta,S}\,R=\left(K_{\beta,S}\,\mathfrak{M}\mathfrak{M}^*\right)^{n_1},
\]
and the traces of $R\,K^{1/2}_{\beta,S}$ and of $K^{1/2}_{\beta,S}\,R$ are equal, $R$ being of trace class and $K^{1/2}_{\beta,S}$ bounded.

The second equality is the word formula, Proposition~\ref{prop:word-formula}\eqref{it:word-identity},
with all magnetic factors equal to $1$, without local step data,
and with the electric factor $\exp\left(i\langle\beta,F^{\mathrm{el}}_j\rangle\right)\cdot\prod_{b\in S}\left(\sqrt{2}\cdot\exp\left(-\frac{1}{2}F^{\mathrm{el}}_j(b)^2\right)\right)$ at every time $j$.
Without local step data the conditions on their boxes are void and the plaquette $\sigma_0$ is arbitrary,
and for this electric factor the operator of every step, multiplied by $(2\pi/q)^{|\Lambda_{\mathrm{s}}|}$ and by $P_{\mathrm{G}}$ from the right,
equals $K_{\beta,S}$ (Proposition~\ref{prop:word-formula}\eqref{it:word-cases}):
for $s\in\R$ one has $\sqrt{2}\cdot\exp\left(-\frac{s^2}{2}\right)\cdot g_1(s)=g_{1/2}(s)$,
and $g_{1/2}$ is the kernel of $\exp\left(-\frac{1}{4}E_b^2\right)$.

For the inequality, observe that the integrand in the middle member is a bounded Borel function of the curvature, hence bosonic,
and that its modulus is at most $2^{\frac{\#|S|\cdot n_1}{2}}$, each of its $n_1$ factors having modulus at most $2^{\frac{\#|S|}{2}}$.
By Proposition~\ref{prop:accretivity} and Corollary~\ref{cor:modulus-bound},
$\left|Z_{C(n)}[G]\right|\le Z_{C(n)}[G']$ for bosonic $G$ and $G'$ with $|G|\le G'$;
with the constant $G'=2^{\frac{\#|S|\cdot n_1}{2}}$ this bounds the modulus of the middle member by the right-hand side,
and the left-hand side is non-negative.

\emph{Reduction of the word.}
The factor with index $j$ of the product in the statement is
\[
K^{1/2}_{\beta_j,S_j}\,\mathcal{B}_j\,K^{1/2}_{\beta_j,S_j}\,\mathfrak{M}\,u_j\,\mathfrak{M}^*=K^{1/2}_{\beta_j,S_j}\,\mathcal{B}_j\,T_j\,u_j\,\mathfrak{M}^*,
\]
the Hilbert--Schmidt operator $T_j$ multiplied on both sides by bounded operators;
hence it is a Hilbert--Schmidt operator on $\mathcal{H}^{\mathrm{l}}$.
As $n_1\ge4$, there are at least two factors, and the product of the two leftmost ones is of trace class;
multiplication from the right by the remaining factors, which are bounded, keeps it in the trace class.
Since $K^{1/2}_{\beta_{j-1},S_{j-1}}$ is self-adjoint, $\mathfrak{M}^*K^{1/2}_{\beta_{j-1},S_{j-1}}=T_{j-1}^*$,
and therefore $T_j\,u_j\,\mathfrak{M}^*K^{1/2}_{\beta_{j-1},S_{j-1}}=T_j\,u_j\,T_{j-1}^*$, with indices modulo $n_1$.
The leftmost factor $K^{1/2}_{\beta_{n_1-1},S_{n_1-1}}$ is bounded, and the product of the remaining operators is of trace class;
moving it to the right end leaves the trace unchanged and gives
\[
\Tr_{\mathcal{H}^{\mathrm{l}}}\prod_{j=n_1-1,\dots,0}\left(K^{1/2}_{\beta_j,S_j}\,\mathcal{B}_j\,K^{1/2}_{\beta_j,S_j}\,\mathfrak{M}\,u_j\,\mathfrak{M}^*\right)
=\Tr_{\mathcal{H}^{\mathrm{l}}}\prod_{j=n_1-1,\dots,0}\left(\mathcal{B}_j\,T_j\,u_j\,T_{j-1}^*\right),
\]
where $T_{-1}=T_{n_1-1}$.

\emph{H\"older's inequality on one space.}
On the right-hand side of the last display,
operators $\mathcal{B}_j\,T_j\colon\mathcal{H}^{\mathrm{s}}\to\mathcal{H}^{\mathrm{l}}$ alternate with operators $u_j\,T_{j-1}^*\colon\mathcal{H}^{\mathrm{l}}\to\mathcal{H}^{\mathrm{s}}$,
$2n_1$ factors in all.
Their images $\widehat{\mathcal{B}_j\,T_j}=\widehat{\mathcal{B}_j}\,\widehat{T_j}$ and $\widehat{u_j\,T_{j-1}^*}=\widehat{u_j}\,\widehat{T_{j-1}}^{\,*}$ are operators on the single space $\mathcal{H}$,
and Lemma~\ref{lem:two-spaces}(i) and (iii) give
\[
\Tr_{\mathcal{H}^{\mathrm{l}}}\prod_{j=n_1-1,\dots,0}\left(\mathcal{B}_j\,T_j\,u_j\,T_{j-1}^*\right)
=\Tr_{\mathcal{H}}\prod_{j=n_1-1,\dots,0}\left(\left(\widehat{\mathcal{B}_j}\,\widehat{T_j}\right)\left(\widehat{u_j}\,\widehat{T_{j-1}}^{\,*}\right)\right).
\]
By Lemma~\ref{lem:two-spaces}(ii) with $r=1$ each $\widehat{T_j}$ is Hilbert--Schmidt on $\mathcal{H}$,
so it lies in $\mathscr{J}_{2n_1}(\mathcal{H})$,
because the Schatten norms decrease as the exponent increases (Lemma~\ref{lem:schatten-traces}(b)).
By \eqref{eq:amplification-ideal}, by Lemma~\ref{lem:two-spaces}(i), by the equality $\|\widehat{T_j}\|_{2n_1}=\|T_j\|_{2n_1}$
and by the invariance of the Schatten norms under adjoints (Lemma~\ref{lem:schatten-traces}(b)),
the $2n_1$ factors $\widehat{\mathcal{B}_j}\,\widehat{T_j}$ and $\widehat{u_j}\,\widehat{T_{j-1}}^{\,*}$ lie in $\mathscr{J}_{2n_1}(\mathcal{H})$, with
\[
\left\|\widehat{\mathcal{B}_j}\,\widehat{T_j}\right\|_{2n_1}\le\|\mathcal{B}_j\|\cdot\|T_j\|_{2n_1},
\qquad
\left\|\widehat{u_j}\,\widehat{T_{j-1}}^{\,*}\right\|_{2n_1}\le\left\|\widehat{T_{j-1}}^{\,*}\right\|_{2n_1}=\|T_{j-1}\|_{2n_1};
\]
in the second bound $\|\widehat{u_j}\|=\|u_j\|\le1$, since $|u_j|\le1$.
H\"older's inequality on the one space $\mathcal{H}$, Lemma~\ref{lem:schatten-traces}(a),
applied with all $2n_1$ exponents equal to $2n_1$, so that their reciprocals add up to one, gives
\[
\left|\Tr_{\mathcal{H}^{\mathrm{l}}}\prod_{j=n_1-1,\dots,0}\left(\mathcal{B}_j\,T_j\,u_j\,T_{j-1}^*\right)\right|
\le\prod_{j\in\Z/n_1\Z}\left(\|\mathcal{B}_j\|\cdot\|T_j\|_{2n_1}\cdot\|T_{j-1}\|_{2n_1}\right)
=\prod_{j\in\Z/n_1\Z}\left(\|\mathcal{B}_j\|\cdot\|T_j\|_{2n_1}^2\right).
\]
Finally, the definition of $\|T_j\|_{2n_1}$ and \eqref{eq:amplification-trace-power} with $\beta=\beta_j$ and $S=S_j$ give
\[
\|T_j\|_{2n_1}^2=\left(\Tr_{\mathcal{H}^{\mathrm{s}}}\left(T_j^*T_j\right)^{n_1}\right)^{\frac{1}{n_1}}
\le2^{\frac{\#|S_j|}{2}}\cdot\left((2\pi)^{-|n|}\cdot Z_{C(n)}[1]\right)^{\frac{1}{n_1}},
\]
and the product of these $n_1$ bounds, multiplied by $\prod_j\|\mathcal{B}_j\|$, is the right-hand side of the assertion.
\end{proof}

\begin{remark}\label{rem:holder-jensen}
An inequality of the same form as Proposition~\ref{prop:holder-insertions} occurs in the argument for four-dimensional Yang--Mills theory,
with an operator of that theory in the place of $K_{\beta_j,S_j}$ and with inserted operators of norm at most one.
The step that follows it there, from this trace inequality to a state by Jensen's inequality for the logarithm of a kernel, does not carry over.
Here its place is taken by the amplification of Theorem~\ref{thm:amplification}.
\end{remark}

\subsection{Properties of insertions}\label{sec:insertion-properties}

\begin{proposition}[Insertions]\label{prop:insertions}
Let $S$ be a finite set of links of the slice lattice $\Z^2$,
let $\Lambda$ be a finite non-empty set of sites of $\Z^2$ that contains the end points of the links of $S$,
let $j\in\Z$ and $c=j+\frac12$,
and let $X\in\mathfrak{W}(S,\Lambda)$.
The identities between integrands in (a), (b) and (c) compare integrands that read in different boxes;
two such integrands are equal when they coincide after both have been read in one box that contains the two boxes (Lemma~\ref{lem:equal-integrands}(b)).
\begin{enumerate}
\item[(a)]
The map $X\mapsto G_X$ is linear, and $G_1=1$.
The integrand $G_X$ is an element of $\mathcal{G}^{\mathrm{ev}}(B)$, $B=B(c,S,\Lambda)$,
that is, a map from $\R^{B^{(1)}}$ to $\Gr(\mathbb{X}(B))$ with the properties of Definition~\ref{def:integrands},
and its definition involves no torus.
Let $S_1\supset S$ be a finite set of links,
let $\Lambda_1\supset\Lambda$ be a finite set of sites that contains the end points of the links of $S_1$,
and let $X_1$ be the image of $X$ under the inclusion of $\mathfrak{W}(S,\Lambda)$ in $\mathfrak{W}(S_1,\Lambda_1)$ of Definition~\ref{def:mid-plane}.
Then $G_{X_1}=G_X$,
where $G_{X_1}\in\mathcal{G}^{\mathrm{ev}}\left(B^{\#}(c,S_1,\Lambda_1)\right)$ and $G_X\in\mathcal{G}^{\mathrm{ev}}\left(B^{\#}(c,S,\Lambda)\right)$.

Let $X$ be of the form $X=\exp\left(i\langle\delta,A\rangle\right)\Upsilon\,f_{S',\mathfrak{Q},\lambda}(E_S)$,
where $\delta\in\mathfrak{D}(S)$, $\Upsilon\in\End_{z(\delta)}(\Lambda)$, $S'\subset S$,
$\mathfrak{Q}$ is a positive definite real quadratic form on $\R^{S'}$, and $\lambda\in\C^{S'}$.
Denote by $\mathfrak{Q}$ also the real symmetric matrix on $\R^S$ whose rows and columns indexed by the links of $S\setminus S'$ vanish
and which satisfies $\langle\xi,\mathfrak{Q}\xi\rangle=\mathfrak{Q}(\xi_{S'})$ for $\xi\in\R^S$;
extend $\lambda$ by zero to an element of $\C^S$,
write $\langle\cdot,\cdot\rangle$ also for the bilinear extension of the scalar product of $\R^S$ to $\C^S$,
and put $x:=F^{\mathrm{el}}_j|_S\in\R^S$.
Then
\[
\begin{split}
G_X={}&\det(1+2\mathfrak{Q})^{-\frac12}\cdot\exp\left(-\frac{|\delta|^2}{4}+\frac{|x|^2}{2}\right)\\
&\cdot\exp\left(\frac12\left\langle\lambda-\frac{\delta}{2}+ix,\,(1+2\mathfrak{Q})^{-1}\left(\lambda-\frac{\delta}{2}+ix\right)\right\rangle\right)\\
&\cdot\exp\left(i\left\langle\delta,\hat A_{j+1}\right\rangle\right)\cdot\Ins_c\left[\omega^{\frac12}\,\Upsilon\,u(\hat a_j)\,\omega^{\frac12}\right]\cdot\exp\left(S^{\times}_c(\Lambda)\right),
\end{split}
\]
and the product of the four scalar factors in front of $\Ins_c$ has the modulus
\[
C\cdot\exp\left(\mathfrak{Q}'(x)-\left\langle\Im\lambda,(1+2\mathfrak{Q})^{-1}x\right\rangle\right),
\]
with a constant $C>0$ that does not depend on $\hat A$;
here $\mathfrak{Q}'(x):=\langle x,\mathfrak{Q}'x\rangle$ is the quadratic form of the real symmetric matrix
\[
\mathfrak{Q}':=\frac12\left(1-(1+2\mathfrak{Q})^{-1}\right)=\mathfrak{Q}\,(1+2\mathfrak{Q})^{-1}.
\]
The number $\mathfrak{Q}'(x)$ depends on $x_{S'}$ only,
and $0\le\mathfrak{Q}'<\frac12$ as symmetric matrices.
\item[(b)]
$\Theta_cG_X=G_{X^*}$ and $\Theta_cG'_X=G'_{X^*}$;
the four integrands lie in $\mathcal{G}^{\mathrm{ev}}\left(B^{\#}(c,S,\Lambda)\right)$.
\item[(c)]
Let $g\in W_{\mathrm{c}}$,
let $g(c)\in\Z+\frac12$ be given by $g\{x_1=c\}=\{x_1=g(c)\}$,
and let $gS$ and $g\Lambda$ be the images of $S$ and $\Lambda$ under the action of $g$ on the coordinates $(x_2,x_3)$.
The map $X\mapsto X^g$ of Lemma~\ref{lem:coordinate-images}\eqref{it:images-group}, formed with the normal word of $g$,
is a bijection of $\mathfrak{W}(S,\Lambda)$, attached to the plane $c$,
onto $\mathfrak{W}(gS,g\Lambda)$, attached to the plane $g(c)$,
and
\[
\left(G^{(c)}_X\right)^g=G^{(g(c))}_{X^g},\qquad
\left(G'^{(c)}_X\right)^g=G'^{(g(c))}_{X^g},\qquad
\|X^g\|_\#=\|X\|_\#.
\]
For a translation $g=\tau_v$, $v\in\Z^3$, one has $X^g=\mathfrak{T}_vX\mathfrak{T}_v^{-1}$,
with the unitary $\mathfrak{T}_v$ from $\mathcal{H}_S$ to $\mathcal{H}_{S+v_{\mathrm{s}}}$ of Lemma~\ref{lem:coordinate-images}.
For the reflection $g$ in a time mirror $(x_1,c')$, $c'\in\frac12\Z$,
one has $X^g=\mathfrak{T}_vX^*\mathfrak{T}_v^{-1}$ with $v=2(c'-c)e_1$.
For the reflection $g=\theta_w$ in a spatial mirror $\Pi_w=(x_\nu,w)$, $\nu\in\{2,3\}$, $w\in\frac12\Z$,
one has $X^g=\mathfrak{J}X\mathfrak{J}^{-1}$;
here $\theta_{\mathrm{s}}$ is the reflection of the slice induced by $\theta_w$,
and $\mathfrak{J}$ is the antiunitary map from $\mathcal{H}_S$ to $\mathcal{H}_{\theta_{\mathrm{s}}S}$ given by
$(\mathfrak{J}\Phi)(A)=\mathfrak{U}\,\overline{\Phi(\theta_{\mathrm{s}}^*A)}$,
with the unitary $\mathfrak{U}=(\wedge\mathfrak{u})\,\mathfrak{r}$ of Lemma~\ref{lem:spatial-mirror-signs}.
If $|\rho|<1$, then $X^g$ is the only element $X'$ of $\mathfrak{W}(gS,g\Lambda)$ with $G^{(g(c))}_{X'}=\left(G^{(c)}_X\right)^g$;
in particular, if $g$ is written in any way as a product of translations and reflections in coordinate planes
and $X$ is transported factor by factor with the three rules above, each read at the current plane, the result is $X^g$.
For every local step datum $(S,\Lambda,k)$ (Proposition~\ref{prop:word-formula}) one has likewise
$\left(G^{(c)}[k]\right)^g=G^{(g(c))}[k^g]$,
where $k^g$ is obtained from $k$ along the normal word of $g$ by the maps $k\mapsto k^{\tau_v}$, $k\mapsto k^{\theta_c}$, $k\mapsto k^{\theta_w}$ of Lemma~\ref{lem:coordinate-images}.
All integrands in (c) lie in $\mathcal{G}^{\mathrm{ev}}\left(B^{\#}(g(c),gS,g\Lambda)\right)$.
\item[(d)]
Let $C(n)=C(n_1,n_2,n_3)=(\Gamma,\chi)$ be a coordinate torus,
with the time direction $e_1$ and $|n|=n_1n_2n_3$,
let $\Xi=\Xi(C(n),x_1)$ be its cross-section for the family $x_1$,
let $\sigma_0$ be a plaquette of the slice torus $\Lambda_{\mathrm{s}}$ of $\Xi$,
and let $\mathfrak{M}$ be formed with $\sigma_0$ (Definition~\ref{def:slice-operators}).
For $a=1,\dots,r$ let $j_a\in\{0,\dots,n_1-1\}$ and $c_a=j_a+\frac12$,
let $S_a$, $\Lambda_a$ be sets as $S$, $\Lambda$ above,
and let $X_a\in\mathfrak{W}(S_a,\Lambda_a)$.
Assume that the boxes $B(c_a,S_a,\Lambda_a)$ satisfy the conditions \eqref{it:step-disjoint} and \eqref{it:step-string} of Proposition~\ref{prop:word-formula}:
they embed disjointly in $\Gamma$,
and no translate of one of them by an element of $\Gamma$ contains a plaquette of a plane $\{x_1=j'\}$, $j'\in\Z$, that represents $\sigma_0$.
Let $\beta\in C^2(\Gamma)$, and let $\varphi_\beta(F)=\exp\left(i\langle\beta,F\rangle\right)$ be formed with the pairing of $C^2(\Gamma)$.
For $j=0,\dots,n_1-1$ let $\beta_j\in\R^{\Lambda^{(1)}_{\mathrm{s}}}$, $\beta_j(y,\nu):=\beta_{1\nu}(j,y)$, be the electric part of $\beta$ at the step $j$,
and let $u_j:=\exp\left(i\left\langle\beta^{\mathrm{mag}}_j,\cdot\right\rangle\right)$, with the magnetic part $\beta^{\mathrm{mag}}_j(y):=\beta_{23}(j,y)$,
act on $\mathcal{H}^{\mathrm{s}}$ as in Proposition~\ref{prop:word-formula}.
Then $\hull(S_a,\Lambda_a)$ embeds in $\Xi$ and does not contain $\sigma_0$,
so that $\left(\alpha_{\beta_j}X_a\right)_\Xi$ is defined, $\alpha_{\beta_j}$ being formed with the restriction of $\beta_j$ to $S_a$,
and
\[
\begin{split}
&Z_{C(n)}\left[\varphi_\beta\cdot\prod_{a=1}^{r}G^{(c_a)}_{X_a}\right]\\
&\qquad=(2\pi)^{|n|}\cdot\Tr_{\mathcal{H}^{\mathrm{l}}}\prod_{j=n_1-1,\dots,0}\left(K_{\beta_j}^{\frac12}\left[\prod_{a:\,j_a=j}\left(\alpha_{\beta_j}X_a\right)_\Xi\right]K_{\beta_j}^{\frac12}\,\mathfrak{M}\,u_j\,\mathfrak{M}^{*}\right),
\end{split}
\]
the factors of the product over $j$ standing from left to right in the order $j=n_1-1,\dots,0$.
\end{enumerate}
\end{proposition}

\begin{proof}
(a)
By Lemma~\ref{lem:regularised-kernels}(a) the map $X\mapsto k_X$ is linear,
and $k\mapsto G^{(c)}[k]$ is linear because $\Ins_c$ is linear in its matrix;
so $X\mapsto G_X=G^{(c)}[k_X]$ is linear.
For $X=1$ the operator
$\exp\left(-\frac14|E_S|^2\right)\cdot1\cdot\exp\left(-\frac14|E_S|^2\right)=\exp\left(-\frac12|E_S|^2\right)$
has the kernel $\prod_{b\in S}g_1\left(A'(b)-A(b)\right)$ times the identity of $\mathcal{F}_\Lambda$ (Definition~\ref{def:integral-kernels}).
At $A'=\hat A_{j+1}|_S$ and $A=(\hat A_j+d_{\mathrm{s}}\hat a_j)|_S$ one has $A'-A=F^{\mathrm{el}}_j|_S$,
since $d\hat A=F$ on the plaquettes of $B(c,S,\Lambda)$ (Definition~\ref{def:integrands}),
and these include the plaquettes spanned by $e_1$ and the links $b\in S$ in the plane $\{x_1=j\}$.
So the kernel takes the value $g_S:=\prod_{b\in S}g_1\left(F^{\mathrm{el}}_j(b)\right)$,
the factor $g_S^{-1}$ of Definition~\ref{def:insertions} cancels it, and
\[
G_1=\Ins_c\left[\omega^{\frac12}\,u(\hat a_j)\,\omega^{\frac12}\right]\cdot\exp\left(S^{\times}_c(\Lambda)\right).
\]
The matrix $\omega^{\frac12}u(\hat a_j)\,\omega^{\frac12}$ is diagonal,
with the entries $\omega_Iu_I$, $u_I=\exp\left(iq\left\langle\hat a_j,\mathfrak{q}(I)\right\rangle\right)$, $I\subset\mathcal{J}^{\mathrm{l}}(\Lambda)$ (Definition~\ref{def:modes});
by Proposition~\ref{prop:crossing}, applied to the crossing bonds over $\Lambda$,
its image under $\Ins_c$ is $\exp\left(-S^{\times}_c(\Lambda)\right)$.
Hence $G_1=1$.

By Lemma~\ref{lem:regularised-kernels}(e) the triple $(S,\Lambda,k^{(s)}_X)$ is a local step datum for $s=\frac14$ and for $s=\frac18$,
and $G_X$, $G'_X$ are the integrands $G^{(c)}\left[k^{(s)}_X\right]$ of these data (Definition~\ref{def:two-insertions}).
For every local step datum $(S,\Lambda,k)$ the integrand $G^{(c)}[k]$ is an element of $\mathcal{G}^{\mathrm{ev}}(B(c,S,\Lambda))$ whose definition involves no torus
(Proposition~\ref{prop:word-formula}\eqref{it:step-class});
this is the second assertion of (a), and it holds for $G'_X$ as well.
The equality $G_{X_1}=G_X$ is Lemma~\ref{lem:enlarged-datum}.
By Lemma~\ref{lem:covariant-box}(b) one has $B(c,S,\Lambda)\subset B^{\#}(c,S,\Lambda)$ and $B(c,S_1,\Lambda_1)\subset B^{\#}(c,S_1,\Lambda_1)$,
and an element of $\mathcal{G}^{\mathrm{ev}}$ of a box is an element of $\mathcal{G}^{\mathrm{ev}}$ of every box that contains it (Lemma~\ref{lem:equal-integrands});
so $G_X$ and $G_{X_1}$ lie in the classes named in (a).

Let now $X=\exp\left(i\langle\delta,A\rangle\right)\Upsilon\,f_{S',\mathfrak{Q},\lambda}(E_S)$ be as in the second part of (a).
This is a normal form of $X$, with the single exponent $\delta$ and the symbol $f_{S',\mathfrak{Q},\lambda}\Upsilon$ (Lemma~\ref{lem:normal-form}(b)).
Lemma~\ref{lem:regularised-kernels}, with its part (c) for $s=\frac14$, $\mathfrak{Q}_{\frac14}=1+2\mathfrak{Q}$ and $\lambda_{\frac14}=\lambda-\frac{\delta}{2}$, gives
\[
\begin{split}
k_X(A',A)={}&(2\pi)^{-\frac{\#|S|}{2}}\cdot\det(1+2\mathfrak{Q})^{-\frac12}\cdot\exp\left(i\langle\delta,A'\rangle\right)\\
&\cdot\exp\left(-\frac{|\delta|^2}{4}+\frac12\left\langle\lambda-\frac{\delta}{2}+i(A'-A),\,(1+2\mathfrak{Q})^{-1}\left(\lambda-\frac{\delta}{2}+i(A'-A)\right)\right\rangle\right)\Upsilon.
\end{split}
\]
At $A'=\hat A_{j+1}|_S$, $A=(\hat A_j+d_{\mathrm{s}}\hat a_j)|_S$ one has $A'-A=x$, as shown above.
Scalar factors can be taken out of $\Ins_c$,
and $g_S=\prod_{b\in S}g_1\left(x(b)\right)=(2\pi)^{-\frac{\#|S|}{2}}\cdot\exp\left(-\frac12|x|^2\right)$.
Inserting the value of $k_X$ into
\[
G_X=\Ins_c\left[\omega^{\frac12}\,k_X\left(\hat A_{j+1}|_S,(\hat A_j+d_{\mathrm{s}}\hat a_j)|_S\right)u(\hat a_j)\,\omega^{\frac12}\right]\cdot g_S^{-1}\cdot\exp\left(S^{\times}_c(\Lambda)\right)
\]
(Definition~\ref{def:insertions}) one obtains the display of (a).

The matrix $(1+2\mathfrak{Q})^{-1}$ is real and symmetric, and $\delta$ and $x$ are real.
Therefore the real part of $\frac12\left\langle\lambda-\frac{\delta}{2}+ix,(1+2\mathfrak{Q})^{-1}\left(\lambda-\frac{\delta}{2}+ix\right)\right\rangle$ is
\[
\frac12\left\langle\Re\lambda-\frac{\delta}{2},(1+2\mathfrak{Q})^{-1}\left(\Re\lambda-\frac{\delta}{2}\right)\right\rangle
-\frac12\left\langle x+\Im\lambda,(1+2\mathfrak{Q})^{-1}(x+\Im\lambda)\right\rangle,
\]
while $\exp\left(i\langle\delta,\hat A_{j+1}\rangle\right)$ has modulus one.
Expanding the second term, adding $\frac12|x|^2$ and using $\frac12|x|^2-\frac12\left\langle x,(1+2\mathfrak{Q})^{-1}x\right\rangle=\mathfrak{Q}'(x)$,
one finds the modulus asserted in (a), with
\[
\begin{split}
C={}&\det(1+2\mathfrak{Q})^{-\frac12}\cdot\exp\left(-\frac{|\delta|^2}{4}-\frac12\left\langle\Im\lambda,(1+2\mathfrak{Q})^{-1}\Im\lambda\right\rangle\right)\\
&\cdot\exp\left(\frac12\left\langle\Re\lambda-\frac{\delta}{2},(1+2\mathfrak{Q})^{-1}\left(\Re\lambda-\frac{\delta}{2}\right)\right\rangle\right).
\end{split}
\]
The matrix $\mathfrak{Q}$ vanishes on the coordinates of $S\setminus S'$;
so $1+2\mathfrak{Q}$ and its inverse act as the identity on these coordinates,
$1-(1+2\mathfrak{Q})^{-1}$ vanishes there, and $\mathfrak{Q}'(x)$ depends on $x_{S'}$ only.
From $1\le1+2\mathfrak{Q}\le1+2\|\mathfrak{Q}\|$ one obtains $(1+2\|\mathfrak{Q}\|)^{-1}\le(1+2\mathfrak{Q})^{-1}\le1$, hence
\[
0\le\mathfrak{Q}'\le\frac12\left(1-(1+2\|\mathfrak{Q}\|)^{-1}\right)<\frac12 .
\]
Finally, $\frac12\left(1-(1+2\mathfrak{Q})^{-1}\right)=\frac12\left((1+2\mathfrak{Q})-1\right)(1+2\mathfrak{Q})^{-1}=\mathfrak{Q}\,(1+2\mathfrak{Q})^{-1}$,
a symmetric matrix, as $\mathfrak{Q}$ and $(1+2\mathfrak{Q})^{-1}$ commute.

The boxes in (b) and (c) are treated next.
Let $g\in W_{\mathrm{c}}$ and let $(S,\Lambda,k)$ be a local step datum.
By Lemma~\ref{lem:reflected-integrands}, $\left(G^{(c)}[k]\right)^g$ lies in $\mathcal{G}^{\mathrm{ev}}\left(gB(c,S,\Lambda)\right)$.
The box $gB(c,S,\Lambda)$ has the radius $r_0$ of $B(c,S,\Lambda)$,
and it contains $g\mathcal{K}(c,S,\Lambda)=\mathcal{K}(g(c),gS,g\Lambda)$ (Lemma~\ref{lem:covariant-box}(a)).
Since $gB^{\#}(c,S,\Lambda)=B^{\#}(g(c),gS,g\Lambda)$ by the same lemma, and $g$ does not change the radii of boxes,
$r_0$ is also the radius of $B(g(c),gS,g\Lambda)$.
Lemma~\ref{lem:covariant-box}(b), applied at the plane $g(c)$, now shows that $gB(c,S,\Lambda)$ lies in $B^{\#}(g(c),gS,g\Lambda)$;
so does $B(g(c),gS,g\Lambda)$.
The right-hand sides of the equalities in (b) and (c) are integrands of local step data with the sets $gS$, $g\Lambda$ at the plane $g(c)$
(Lemma~\ref{lem:regularised-kernels}(e), Lemma~\ref{lem:coordinate-images}),
hence elements of $\mathcal{G}^{\mathrm{ev}}(B(g(c),gS,g\Lambda))$ by Proposition~\ref{prop:word-formula}\eqref{it:step-class};
the left-hand sides are images of integrands of local step data at the plane $c$.
With Lemma~\ref{lem:equal-integrands} this places all integrands of (c) in $\mathcal{G}^{\mathrm{ev}}\left(B^{\#}(g(c),gS,g\Lambda)\right)$,
and, for $g=\theta_c$, where $g(c)=c$, $gS=S$ and $g\Lambda=\Lambda$, all integrands of (b) in $\mathcal{G}^{\mathrm{ev}}\left(B^{\#}(c,S,\Lambda)\right)$.

(b)
For every local step datum, Lemma~\ref{lem:coordinate-images}\eqref{it:images-time-reflection} gives $\Theta_cG^{(c)}[k]=G^{(c)}[k^{\theta_c}]$ with $k^{\theta_c}(A',A)=k(A,A')^*$.
For $k=k^{(s)}_X$ one has $\left(k^{(s)}_X\right)^{\theta_c}=k^{(s)}_{X^*}$ by Lemma~\ref{lem:regularised-kernels}(a).
With $s=\frac14$ this is $\Theta_cG_X=G_{X^*}$, and with $s=\frac18$ it is $\Theta_cG'_X=G'_{X^*}$.

(c)
By Lemma~\ref{lem:coordinate-images}\eqref{it:images-elements} the maps $X\mapsto\mathfrak{T}_vX\mathfrak{T}_v^{-1}$, $X\mapsto X^*$ and $X\mapsto\mathfrak{J}X\mathfrak{J}^{-1}$
are bijections of $\mathfrak{W}(S,\Lambda)$ onto $\mathfrak{W}(S+v_{\mathrm{s}},\Lambda+v_{\mathrm{s}})$, $\mathfrak{W}(S,\Lambda)$ and $\mathfrak{W}(\theta_{\mathrm{s}}S,\theta_{\mathrm{s}}\Lambda)$,
and they preserve $\|\cdot\|_\#$.
The map $X\mapsto X^g$ is the composition of such maps along the normal word of $g$;
it is therefore a bijection onto $\mathfrak{W}(gS,g\Lambda)$ with $\|X^g\|_\#=\|X\|_\#$,
and the two identities for the insertions are Lemma~\ref{lem:coordinate-images}\eqref{it:images-group}.

Let $|\rho|<1$.
By Lemma~\ref{lem:injective-insertions} the map $X'\mapsto G^{(g(c))}_{X'}$ is injective on $\mathfrak{W}(gS,g\Lambda)$,
so $X^g$ is the only element with $G^{(g(c))}_{X^g}=\left(G^{(c)}_X\right)^g$.
Write $g$ in any way as a product of translations and reflections in coordinate planes,
and follow this product factor by factor, beginning with the rightmost factor, each time with the rule of (c) at the current plane.
At each step an element $X_0$ is replaced by an element $X_2$ with $G_{X_2}=\left(G_{X_0}\right)^{g'}$, $g'$ the factor of that step:
for a translation by Lemma~\ref{lem:coordinate-images}\eqref{it:images-translation} and \eqref{it:images-elements},
for a reflection in a time mirror by the last assertion of Lemma~\ref{lem:coordinate-images}\eqref{it:images-group},
for a reflection in a spatial mirror by Lemma~\ref{lem:coordinate-images}\eqref{it:images-spatial-mirror} and \eqref{it:images-elements}.
By the composition rule of \eqref{eq:setting-group-action},
which holds for words in arbitrary mirrors by Proposition~\ref{prop:symmetries}\eqref{it:reflection-words},
the final element $X'$ satisfies $G_{X'}=\left(G_X\right)^g$; hence $X'=X^g$.

For $g=\tau_v$ the normal word is $\tau_v$, and $X^g=\mathfrak{T}_vX\mathfrak{T}_v^{-1}$ by definition.
For the reflection $\theta_{c'}$ in the time mirror $(x_1,c')$, Definition~\ref{def:mirrors}\eqref{it:mirror-involution} and~\eqref{it:mirror-product} give $\theta_{c'}=\tau_v\theta_c$ with $v=2(c'-c)e_1$;
this is the normal word, and $X^g=\left(X^*\right)^{\tau_v}=\mathfrak{T}_vX^*\mathfrak{T}_v^{-1}$.
For the reflection $g=\theta_w$ in a spatial mirror $\Pi_w=(x_\nu,w)$, whatever the value of $w\in\frac12\Z$,
Lemma~\ref{lem:coordinate-images}\eqref{it:images-group} gives the normal word $\tau_{2w\,e_\nu}\theta_{x_\nu,0}$ and the identity $X^g=\mathfrak{J}X\mathfrak{J}^{-1}$.

Let $(S,\Lambda,k)$ be a local step datum,
and let $g=\tau_v\theta_c^{p_1}\theta_{x_2,0}^{p_2}\theta_{x_3,0}^{p_3}$, $p\in\{0,1\}^3$, be the normal word of $g$.
By \eqref{eq:setting-group-action} and Proposition~\ref{prop:symmetries}\eqref{it:reflection-words},
\[
G^g=\left(\left(\left(G^{\theta_{x_3,0}^{p_3}}\right)^{\theta_{x_2,0}^{p_2}}\right)^{\theta_c^{p_1}}\right)^{\tau_v}
\]
for every even $G$, and $G^{\theta_\Pi}=\Theta_\Pi G$ for a mirror $\Pi$.
Lemma~\ref{lem:coordinate-images}\eqref{it:images-spatial-mirror}, applied to the two spatial reflections, which do not move the plane $c$,
then \eqref{it:images-time-reflection} and \eqref{it:images-translation} of the same lemma give $\left(G^{(c)}[k]\right)^g=G^{(g(c))}[k^g]$;
after each step the transformed kernel belongs to a local step datum (Lemma~\ref{lem:coordinate-images}), so that the next step applies.

(d)
The box $B(c_a,S_a,\Lambda_a)$ contains the cubes $(j_a,y)+\left[-\frac12,\frac12\right]^3$, $y\in\Lambda_a$.
By condition~\eqref{it:step-disjoint} of Proposition~\ref{prop:word-formula} it embeds in $\Gamma$,
and by condition~\eqref{it:step-string} of that proposition none of its translates by elements of $\Gamma$ contains a plaquette of a plane $\{x_1=j'\}$, $j'\in\Z$, that represents $\sigma_0$.
Lemma~\ref{lem:mixed-word-formula}(o) therefore shows that $\hull(\Lambda_a)$, which is $\hull(S_a,\Lambda_a)$, embeds in $\Xi$ and does not contain $\sigma_0$.
Consequently the realisation on $\Xi$ is defined on $\mathfrak{W}(S_a,\Lambda_a)$ (Lemma~\ref{lem:normal-form}(f)),
and it applies to $\alpha_{\beta_{j_a}}X_a$, which lies in $\mathfrak{W}(S_a,\Lambda_a)$ by Lemma~\ref{lem:normal-form}(e).
Now apply Lemma~\ref{lem:mixed-word-formula}(b) with $s_a=\frac14$ for all $a$.
Every insertion of that lemma is then $G^{(c_a)}_{X_a}$,
the set $S_j$ formed there with the insertions into the comparison operator is empty for every $j$,
and $K_{\beta_j,\emptyset}=K_{\beta_j}$.
The identity of that lemma becomes the display of (d), and the factors in the brackets commute by the same lemma.
\end{proof}

\begin{remark}\label{rem:comparison-transport}
Let $S$, $\Lambda$, $v$, $\Pi_w$, $\theta_{\mathrm{s}}$, $\mathfrak{T}_v$ and $\mathfrak{J}$ be as in Proposition~\ref{prop:insertions}(c), and let $\beta\in\R^S$ and $s>0$.
The unitary $\mathfrak{T}_v$ and the antiunitary $\mathfrak{J}$ act on $\mathcal{H}_S=L^2(\R^S)\otimes\mathcal{F}_\Lambda$ factor by factor.
The first is the tensor product of the composition with $\tau^{*}_{v_{\mathrm{s}}}$ on $L^2(\R^S)$ and of the unitary $U_v$ on the Fock factor.
In the formula $(\mathfrak{J}\Phi)(A)=\mathfrak{U}\,\overline{\Phi(\theta_{\mathrm{s}}^{*}A)}$ the bar is the complex conjugation of the components of $\Phi(\theta_{\mathrm{s}}^{*}A)$ in the basis $(e_I)$ of $\mathcal{F}_\Lambda$;
so $\mathfrak{J}$ is the tensor product of two antilinear maps,
the composition with $\theta_{\mathrm{s}}^{*}$ followed by the complex conjugation on $L^2(\R^S)$,
and the entrywise complex conjugation in the basis $(e_I)$ followed by the unitary $\mathfrak{U}=(\wedge\mathfrak{u})\,\mathfrak{r}$ on the Fock factor.

\emph{The Fock factor.}
The unitary $U_v$ maps each basis vector $e_I$ to $\pm e_{I+v}$.
The matrix $\mathfrak{u}$ of Lemma~\ref{lem:spatial-mirror-signs}, which expresses $\Theta_w$ on the modes, consists of unitary $2\times2$ blocks;
these blocks depend on the bases $e^{\pm}_{n,k}$ of Definition~\ref{def:spin-matrices}, and they need not be diagonal.
Accordingly $\Theta_w(\zeta_i)$ need not be a multiple of a single mode.
For $I\subset\mathcal{J}^{\mathrm{l}}(\Lambda)$ the product $\zeta^w_I$ of the images $\Theta_w(\zeta_i)$, $i\in I$, taken in the increasing order of $I$,
is the combination $\sum_{I'}(\wedge\mathfrak{u})_{I'I}\,\zeta_{I'}$ of products of modes, with the minors as coefficients (Lemma~\ref{lem:spatial-mirror-signs}(b)).
The image of the product $\zeta_I$ itself carries in addition the reordering sign, since $\Theta_w$ reverses the order of products:
\[
\Theta_w(\zeta_I)=(-1)^{\binom{\#|I|}{2}}\,\zeta^w_I=\sum_{I'}\mathfrak{U}_{I'I}\,\zeta_{I'}
\]
by Lemma~\ref{lem:spatial-mirror-signs}(c)(i), the operator $\mathfrak{r}$ in $\mathfrak{U}=(\wedge\mathfrak{u})\,\mathfrak{r}$ being diagonal with the entries $(-1)^{\binom{\#|I|}{2}}$.
The coefficients of $\Theta_w(\zeta_I)$ are thus the entries of the column of $\mathfrak{U}$ with the index $I$,
and $\mathfrak{U}$ need not map a basis vector $e_I$ to a multiple of a basis vector.

\emph{The comparison operators.}
Let $\Xi$ be a rectangular cross-section in which $\hull(S,\Lambda)$ embeds,
so that the links of $S$ are links of $\Lambda_{\mathrm{s}}$ (Lemma~\ref{lem:normal-form}(f)),
and let $\beta'\in\R^{\Lambda^{(1)}_{\mathrm{s}}}$ coincide with $\beta$ on $S$.
By Definition~\ref{def:slice-operators} and Lemma~\ref{lem:slice-kernels}\eqref{it:gauge-average}, the latter applied also with the empty set in the place of $S$,
the square root $K_{\beta'}^{\frac12}$ of the shifted crossing operator is $(2\pi/q)^{\frac{|\Lambda_{\mathrm{s}}|}{2}}\cdot P_{\mathrm{G}}$ times the product of the commuting operators $\exp\left(-\frac14\left(E_b-\beta'(b)\right)^2\right)$, $b\in\Lambda^{(1)}_{\mathrm{s}}$,
and the square root $K_{\beta',S}^{\frac12}$ of the shifted comparison operator arises from it when the factors of the links $b\in S$ are replaced by $\exp\left(-\frac18\left(E_b-\beta'(b)\right)^2\right)$.
The factors of the links of $S$ thus make up the operator $\exp\left(-\frac14|E-\beta'|_S^2\right)$ in $K_{\beta'}^{\frac12}$ and the operator $\exp\left(-\frac18|E-\beta'|_S^2\right)$ in $K_{\beta',S}^{\frac12}$.
On $\mathcal{H}_S$ the same functions of the components $E_b$, $b\in S$, are the operators $\exp\left(-s|E_S-\beta|^2\right)$ with $s=\frac14$ and $s=\frac18$.

\emph{Transport of the local factors.}
For $x\in\R^S$ the unitary $\exp\left(-i\langle x,E_S\rangle\right)$ is the translation $\Phi\mapsto\Phi(\cdot-x)$.
As $(\mathfrak{T}_v\Phi)(A)=U_v\Phi(\tau^*_{v_{\mathrm{s}}}A)$,
the conjugate $\mathfrak{T}_v\exp\left(-i\langle x,E_S\rangle\right)\mathfrak{T}_v^{-1}$ is the translation by the vector of $\R^{S+v_{\mathrm{s}}}$ with the component $x(b)$ at $b+v_{\mathrm{s}}$.
As $(\theta_{\mathrm{s}}^*A)(b)=\operatorname{sgn}(\theta_{\mathrm{s}},b)\,A(\theta_{\mathrm{s}}b)$,
the conjugate $\mathfrak{J}\exp\left(-i\langle x,E_S\rangle\right)\mathfrak{J}^{-1}$ is the translation by the vector of $\R^{\theta_{\mathrm{s}}S}$ with the component $\operatorname{sgn}(\theta_{\mathrm{s}},b)\,x(b)$ at $\theta_{\mathrm{s}}b$;
since $\mathfrak{J}$ is antilinear, the same conjugate equals $\exp\left(i\left\langle x,\mathfrak{J}E_S\mathfrak{J}^{-1}\right\rangle\right)$.
Comparing the generators of these groups of unitaries one finds, for $b\in S$,
\[
\mathfrak{T}_vE_b\mathfrak{T}_v^{-1}=E_{b+v_{\mathrm{s}}},\qquad
\mathfrak{J}E_b\mathfrak{J}^{-1}=-\operatorname{sgn}(\theta_{\mathrm{s}},b)\,E_{\theta_{\mathrm{s}}b}.
\]
A unitary conjugation commutes with the joint spectral calculus of the commuting family $E_S$,
and an antiunitary one does so on real-valued functions.
Put
\[
\beta^{\tau_v}(b+v_{\mathrm{s}}):=\beta(b),\qquad
\beta^{\theta_w}(\theta_{\mathrm{s}}b):=-\operatorname{sgn}(\theta_{\mathrm{s}},b)\,\beta(b)\qquad(b\in S).
\]
Since $\left(-\operatorname{sgn}(\theta_{\mathrm{s}},b)\,E_{\theta_{\mathrm{s}}b}-\beta(b)\right)^2=\left(E_{\theta_{\mathrm{s}}b}-\beta^{\theta_w}(\theta_{\mathrm{s}}b)\right)^2$,
it follows that
\[
\begin{split}
\mathfrak{T}_v\exp\left(-s|E_S-\beta|^2\right)\mathfrak{T}_v^{-1}&=\exp\left(-s\left|E_{S+v_{\mathrm{s}}}-\beta^{\tau_v}\right|^2\right),\\
\mathfrak{J}\exp\left(-s|E_S-\beta|^2\right)\mathfrak{J}^{-1}&=\exp\left(-s\left|E_{\theta_{\mathrm{s}}S}-\beta^{\theta_w}\right|^2\right).
\end{split}
\]
In the same way $\mathfrak{T}_v\exp\left(i\langle\beta,A\rangle\right)\mathfrak{T}_v^{-1}=\exp\left(i\langle\beta^{\tau_v},A\rangle\right)$,
and, $\mathfrak{J}$ being antilinear,
$\mathfrak{J}\exp\left(i\langle\beta,A\rangle\right)\mathfrak{J}^{-1}=\exp\left(-i\left\langle\beta,\theta_{\mathrm{s}}^*A\right\rangle\right)=\exp\left(i\left\langle\beta^{\theta_w},A\right\rangle\right)$.
Since $\alpha_\beta X=\exp\left(i\langle\beta,A\rangle\right)X\exp\left(-i\langle\beta,A\rangle\right)$ (Lemma~\ref{lem:normal-form}(e)), one obtains
\[
\alpha_{\beta^g}\left(X^g\right)=\left(\alpha_\beta X\right)^g\qquad\text{for } g=\tau_v \text{ and for } g=\theta_w .
\]
The rule for $\beta^{\theta_w}$ is the one by which $\Theta_w$ acts on the characters of the curvature.
To see this, let $\eta$ be a finitely supported real two-cochain on $\Z^3$,
with the character $\varphi_\eta(F)=\exp\left(i\langle\eta,F\rangle\right)$ and the electric part $\eta_j(y,\mu):=\eta_{1\mu}(j,y)$, $\mu\in\{2,3\}$, at the step $j$.
Since $\theta_w^*$ is an involutive isometry of the two-cochains, \eqref{eq:setting-reflection-coefficients} gives
$\Theta_w\varphi_\eta=\overline{\varphi_\eta\circ\theta_w^*}=\varphi_{-\theta_w^*\eta}$.
By \eqref{eq:setting-slice-reflections} the electric part of $-\theta_w^*\eta$ at the step $j$ is $-\theta_{\mathrm{s}}^*\eta_j$;
its value at the link $\theta_{\mathrm{s}}b$ is $-\operatorname{sgn}(\theta_{\mathrm{s}},b)\,\eta_j(b)$, because $\operatorname{sgn}(\theta_{\mathrm{s}},\theta_{\mathrm{s}}b)=\operatorname{sgn}(\theta_{\mathrm{s}},b)$,
and this is the rule that defines $\beta^{\theta_w}$, with $\eta_j$ in the place of $\beta$.
For the reflection $\theta_c$ in the plane $c$ itself, finally, $X^{\theta_c}=X^*$, the set $S$ is not moved,
$\exp\left(-s|E_S-\beta|^2\right)$ is self-adjoint,
and $\left(\alpha_\beta X\right)^*=\alpha_\beta\left(X^*\right)$ because $\alpha_\beta$ is a $*$-automorphism;
a time mirror $(x_1,c')$ combines this with a translation $\tau_v$ along $e_1$, for which $v_{\mathrm{s}}=0$ and $\beta^{\tau_v}=\beta$.
With $\beta^g:=\beta$ for the reflections in time mirrors,
each of the transports of Proposition~\ref{prop:insertions}(c) thus carries $\exp\left(-s|E_S-\beta|^2\right)$ to $\exp\left(-s|E_{gS}-\beta^g|^2\right)$ and $\alpha_\beta X$ to $\alpha_{\beta^g}\left(X^g\right)$.
\end{remark}

\subsection{Bounded elements}\label{sec:bounded-elements}

Recall from Definition~\ref{def:bounded-elements} the placements of a box in a coordinate torus
and the notion of a bounded element.
An even integrand $G\in\mathcal{G}^{\mathrm{ev}}(B)$ is bounded by a constant $C\ge0$ with the box $B$
if, for every coordinate torus $\mathbb{T}=(\Gamma,\chi)$,
every placement $(g_1,\dots,g_r)$ of $B$ in $\mathbb{T}$
and every $\beta\in C^2(\Gamma)$,
\[
\left|Z_{\mathbb{T}}\left[\varphi_\beta\prod_{j=1}^rG^{g_j}\right]\right|\le C^r\cdot Z_{\mathbb{T}}[1],
\]
where $\varphi_\beta(F)=\exp\left(i\langle\beta,F\rangle\right)$ is formed with the pairing of $C^2(\Gamma)$.
The insertions of an element $X\in\mathfrak{W}(S,\Lambda)$ at a coordinate link plane $\{x_\mu=c\}$,
$\mu\in\{1,2,3\}$, $c\in\Z+\frac12$, are those of Definition~\ref{def:other-directions}.
For $\mu=1$ they are the two insertions $G^{(1,c)}_X=G^{(c)}_X$ and $G'^{(1,c)}_X=G'^{(c)}_X$
of Definitions~\ref{def:insertions} and~\ref{def:two-insertions},
and in general
\[
G^{(\mu,c)}_X=\left(G^{(1,c)}_X\right)^{g_\mu},
\qquad
G'^{(\mu,c)}_X=\left(G'^{(1,c)}_X\right)^{g_\mu},
\]
with the cyclic permutation matrix $g_\mu\in\Ocub$ that maps $e_1$ to $e_\mu$;
in particular $g_1=1$.

\begin{proposition}\label{prop:bounded-elements}
Let $-1\le\rho\le1$, $w_i>0$ for $i=0,\dots,3$, $M\in\R$ and $q>0$,
and let boundedness refer to the functionals $Z_{\mathbb{T}}$ formed with these parameters.
\begin{enumerate}
\item[(a)]
Let $G$ be bounded by $C\ge0$ with a box $B$, and let $g\in W$.
Then $G^g$ is bounded by $C$ with the box $gB$.
\item[(b)]
Let $G$ be bounded by $C\ge0$ with a box $B$.
Then $G$ is bounded by $C$ with every box that contains $B$.
Moreover, for every $\beta\in C^2_{\mathrm{c}}$ the integrand $\varphi_\beta G$ is bounded by $C$
with every box that contains $B$ and $\supp\beta$,
the union of the plaquettes $\sigma$ with $\beta(\sigma)\neq0$.
\item[(c)]
Let $G\in\mathcal{G}(B)$ be bosonic with $\sup|G|<\infty$.
Then $G$ is bounded by $\sup|G|$ with the box $B$.
\item[(d)]
Let in addition $0<\rho<1$.
Let $\mu\in\{1,2,3\}$ and $c\in\Z+\frac12$,
let $S$ be a finite set of links of the slice lattice $\Z^2$,
let $\Lambda$ be a non-empty finite set of sites of $\Z^2$ that contains the end points of the links of $S$,
and let $X\in\mathfrak{W}(S,\Lambda)$.
Then $G^{(\mu,c)}_X$ is bounded by $\|X\|_\#$ with the box $g_\mu B^{\#}(c,S,\Lambda)$.
\item[(d$'$)]
Under the hypotheses of (d),
$G'^{(\mu,c)}_X$ is bounded by $2^{\frac{\#|S|}{2}}\cdot\|X\|_\#$ with the box $g_\mu B^{\#}(c,S,\Lambda)$.
\end{enumerate}
\end{proposition}

\begin{proof}
For $g\in W$ let $\vec g\in\Ocub$ be its linear part.
The element $g$ acts on tori by $g\mathbb{T}=\left(\vec g\,\Gamma,\chi\circ\vec g^{-1}\right)$ for $\mathbb{T}=(\Gamma,\chi)$
(Definition~\ref{def:tori}),
and it maps a box onto the box of the same radius centred at the image of the centre,
because $\vec g$ preserves every cube centred at $0$.

\emph{Conjugate placements.}
Let $g\in W$, let $\mathbb{T}=(\Gamma,\chi)$ be the coordinate torus $C(n_1,n_2,n_3)$,
and let $(g_1,\dots,g_r)$ be a placement of a box $B$ in $\mathbb{T}$.
Then $g\mathbb{T}$ is a coordinate torus,
and $\left(gg_1g^{-1},\dots,gg_rg^{-1}\right)$ is a placement of $gB$ in $g\mathbb{T}$.
Indeed, $\vec g$ is a signed permutation matrix:
there is a permutation $\mu\mapsto\mu'$ of $\{1,2,3\}$ with $\vec ge_\mu=\pm e_{\mu'}$.
Hence $\vec g\,\Gamma=\sum_\mu n_\mu\Z e_{\mu'}$,
and on the generator $n_\mu e_{\mu'}$ the character $\chi\circ\vec g^{-1}$ takes the value $\chi\left(\pm n_\mu e_\mu\right)=-1$;
so $g\mathbb{T}$ is the coordinate torus with the side $n_\mu$ in the direction $e_{\mu'}$.
The elements $gg_jg^{-1}$ belong to $W_{\mathrm{c}}$, a normal subgroup of $W$ (Definition~\ref{def:mirrors}).
Since $g$ is affine with linear part $\vec g$ and $\vec g[-1,1]^3=[-1,1]^3$, one has, for $v\in\Gamma$,
\[
g\left(g_jB+[-1,1]^3+v\right)=\left(gg_jg^{-1}\right)(gB)+[-1,1]^3+\vec gv,
\qquad
g\left(g_jB+v\right)=\left(gg_jg^{-1}\right)(gB)+\vec gv,
\]
and $\vec gv$ ranges over $\vec g\,\Gamma$ as $v$ ranges over $\Gamma$.
As $g$ is a bijection of $\R^3$, a family of sets on the left is pairwise disjoint
if and only if the corresponding family on the right is.
So condition~\ref{it:placement-disjoint} of Definition~\ref{def:bounded-elements} for $(g_j)$, $B$ and $\Gamma$
gives the same condition for $\left(gg_jg^{-1}\right)$, $gB$ and $\vec g\,\Gamma$.
For condition~\ref{it:placement-columns},
let $\mu\in\{1,2,3\}$ and let $\sigma$ be a plaquette orthogonal to $e_\mu$
such that no set $g_jB+v$, $v\in\Gamma$, contains one of the plaquettes $\sigma+se_\mu$, $s\in\Z$.
The map $g$ preserves $\Z^3$ and permutes the cells of $\Z^3$ (Definition~\ref{def:cochains});
so $g\sigma$ is a plaquette orthogonal to $e_{\mu'}$,
and $g\left(\sigma+se_\mu\right)=g\sigma\pm se_{\mu'}$:
the column of $\sigma$ in the direction $e_\mu$ is mapped onto the column of $g\sigma$ in the direction $e_{\mu'}$.
By the second identity above, no set $\left(gg_jg^{-1}\right)(gB)+v'$, $v'\in\vec g\,\Gamma$,
contains a plaquette $g\sigma+se_{\mu'}$, $s\in\Z$.
Since $\mu\mapsto\mu'$ is onto,
condition~\ref{it:placement-columns} holds for $\left(gg_jg^{-1}\right)$, $gB$ and $g\mathbb{T}$ in every direction.

\emph{Proof of (a).}
Let $G$ be bounded by $C$ with the box $B$, and let $g\in W$.
By Lemma~\ref{lem:reflected-integrands}, $G^g\in\mathcal{G}^{\mathrm{ev}}(gB)$.
Let $\mathbb{T}'=(\Gamma',\chi')$ be a coordinate torus,
$(g'_1,\dots,g'_r)$ a placement of $gB$ in $\mathbb{T}'$,
and $\beta'\in C^2(\Gamma')$.
The preceding paragraph, applied to the element $g^{-1}$ and to this placement of $gB$,
shows that $\mathbb{T}:=g^{-1}\mathbb{T}'=(\Gamma,\chi)$,
with $\Gamma=\vec g^{-1}\Gamma'$, is a coordinate torus
and that the elements $g_j:=g^{-1}g'_jg$, $j=1,\dots,r$, form a placement of $g^{-1}(gB)=B$ in $\mathbb{T}$;
moreover $g\mathbb{T}=\mathbb{T}'$.
The composition rule \eqref{eq:setting-group-action} gives
\[
\left(G^g\right)^{g'_j}=G^{g'_jg}=G^{gg_j}=\left(G^{g_j}\right)^g.
\]
On even elements the map $H\mapsto H^g$ is multiplicative if $g$ is even and reverses the order of products if $g$ is odd;
since the elements $\left(G^{g_j}\right)^g$ are even, they commute, and in both cases
\[
\prod_{j=1}^r\left(G^g\right)^{g'_j}=\left(\prod_{j=1}^rG^{g_j}\right)^g.
\]
Put $\beta:=g^*\beta'$ if $g$ is even and $\beta:=-g^*\beta'$ if $g$ is odd;
this is an element of $C^2(\Gamma)$, because $\Gamma=\vec g^{-1}\Gamma'$.
The map $\sigma\mapsto g\sigma$ is a bijection from the $\Gamma$-orbits of plaquettes onto the $\Gamma'$-orbits,
and $g^*$ changes signs only;
hence $\langle g^*\beta',g^*F\rangle=\langle\beta',F\rangle$ for $F\in C^2(\Gamma')$,
the pairings being those of $C^2(\Gamma)$ and of $C^2(\Gamma')$.
On functions of $F$ the action \eqref{eq:setting-group-action} reads $H^g(F)=H(g^*F)$ for even $g$
and $H^g(F)=\overline{H(g^*F)}$ for odd $g$; in both cases
\[
\left(\varphi_\beta\right)^g(F)=\exp\left(i\langle g^*\beta',g^*F\rangle\right)=\varphi_{\beta'}(F).
\]
As $\varphi_\beta$ is a scalar function, it follows that
\[
\varphi_{\beta'}\prod_{j=1}^r\left(G^g\right)^{g'_j}=\left(\varphi_\beta\prod_{j=1}^rG^{g_j}\right)^g.
\]
The integrand $H:=\varphi_\beta\prod_jG^{g_j}$ embeds in $\mathbb{T}$:
$G^{g_j}\in\mathcal{G}^{\mathrm{ev}}(g_jB)$ by Lemma~\ref{lem:reflected-integrands},
and the boxes $g_jB$ embed disjointly in $\Gamma$ by condition~\ref{it:placement-disjoint}.
By Proposition~\ref{prop:symmetries}\eqref{it:functional-covariance},
$Z_{g\mathbb{T}}\left[H^g\right]$ is $Z_{\mathbb{T}}[H]$ for even $g$ and $\overline{Z_{\mathbb{T}}[H]}$ for odd $g$;
in particular $Z_{\mathbb{T}'}[1]=Z_{\mathbb{T}}[1]$.
Therefore
\[
\left|Z_{\mathbb{T}'}\left[\varphi_{\beta'}\prod_{j=1}^r\left(G^g\right)^{g'_j}\right]\right|
=\left|Z_{\mathbb{T}}\left[\varphi_\beta\prod_{j=1}^rG^{g_j}\right]\right|
\le C^r\cdot Z_{\mathbb{T}}[1]
=C^r\cdot Z_{\mathbb{T}'}[1],
\]
where the inequality is the boundedness of $G$ with $B$, applied to the placement $(g_j)$ of $B$ in $\mathbb{T}$ and to $\beta$.

\emph{Proof of (b).}
Let $G$ be bounded by $C$ with $B$, and let $B'$ be a box that contains $B$.
By Lemma~\ref{lem:equal-integrands}, $G$ is an element of $\mathcal{G}^{\mathrm{ev}}(B')$.
Let $(g_1,\dots,g_r)$ be a placement of $B'$ in a coordinate torus $\mathbb{T}=(\Gamma,\chi)$;
it is a placement of $B$ in $\mathbb{T}$.
Indeed, the sets $g_jB+[-1,1]^3+v$ and $g_jB+v$ are subsets of $g_jB'+[-1,1]^3+v$ and of $g_jB'+v$,
so the disjointness required in condition~\ref{it:placement-disjoint} holds for $B$;
and a set $g_jB+v$ contains no plaquette that $g_jB'+v$ does not contain,
so the plaquettes of condition~\ref{it:placement-columns} for $B'$ serve for $B$ as well.
The boxes $g_jB\subset g_jB'$ embed in $\Gamma$,
so that $G^{g_j}$ defines the same function on $\Conf(\Gamma)$ whether $G$ is read in $B$ or in $B'$
(Lemma~\ref{lem:equal-integrands}).
The inequality of Definition~\ref{def:bounded-elements} for the box $B'$ is thus an instance of the inequality for $B$,
and $G$ is bounded by $C$ with $B'$.

Let now $\beta\in C^2_{\mathrm{c}}$, and let the box $B'$ contain $B$ and $\supp\beta$.
Then $\varphi_\beta$ is a bounded function of the $F_\sigma$ with $\sigma\subset B'$,
and $\varphi_\beta G\in\mathcal{G}^{\mathrm{ev}}(B')$.
Let $(g_1,\dots,g_r)$ be a placement of $B'$ in a coordinate torus $\mathbb{T}=(\Gamma,\chi)$, and let $\beta'\in C^2(\Gamma)$.
Since $\varphi_\beta$ is a scalar function, \eqref{eq:setting-group-action} gives
$\left(\varphi_\beta G\right)^{g_j}=\left(\varphi_\beta\right)^{g_j}G^{g_j}$.
The pull-back $g_j^*$ is orthogonal with inverse $\left(g_j^{-1}\right)^*$ (Definition~\ref{def:cochains});
so $\langle\beta,g_j^*F\rangle=\left\langle\left(g_j^{-1}\right)^*\beta,F\right\rangle$, and
\[
\left(\varphi_\beta\right)^{g_j}=\varphi_{\beta^{g_j}},
\qquad
\beta^{g_j}:=
\begin{cases}
\left(g_j^{-1}\right)^*\beta & \text{if $g_j$ is even},\\
-\left(g_j^{-1}\right)^*\beta & \text{if $g_j$ is odd}.
\end{cases}
\]
The cochain $\beta^{g_j}\in C^2_{\mathrm{c}}$ has its support in $g_j\supp\beta\subset g_jB'$, and the box $g_jB'$ embeds in $\Gamma$.
By Lemma~\ref{lem:curvature-characters}(b), applied to $\beta^{g_j}$ and to the box $g_jB'$,
the character $\varphi_{\beta^{g_j}}$, evaluated at the periodic extension of the curvature on $\mathbb{T}$,
is the character of $C^2(\Gamma)$ attached to the periodisation of $\beta^{g_j}$.
At a plaquette $\sigma$ this periodisation takes the value $\sum_{v\in\Gamma}\beta^{g_j}(\sigma+v)$,
since $\beta^{g_j}(\sigma+v)=0$ unless $\sigma+v\subset g_jB'$,
and this inclusion holds for at most one $v\in\Gamma$, the translates of $g_jB'$ by the elements of $\Gamma$ being pairwise disjoint.
Hence
\[
\varphi_{\beta'}\prod_{j=1}^r\left(\varphi_\beta G\right)^{g_j}=\varphi_{\beta''}\prod_{j=1}^rG^{g_j},
\qquad
\beta''(\sigma):=\beta'(\sigma)+\sum_{j=1}^r\sum_{v\in\Gamma}\beta^{g_j}(\sigma+v),
\]
where $\beta''\in C^2(\Gamma)$.
Since $(g_j)$ is also a placement of $B$, the boundedness of $G$ with $B$, applied with $\beta''$, gives
\[
\left|Z_{\mathbb{T}}\left[\varphi_{\beta'}\prod_{j=1}^r\left(\varphi_\beta G\right)^{g_j}\right]\right|
=\left|Z_{\mathbb{T}}\left[\varphi_{\beta''}\prod_{j=1}^rG^{g_j}\right]\right|
\le C^r\cdot Z_{\mathbb{T}}[1].
\]

\emph{Proof of (c).}
Let $G\in\mathcal{G}(B)$ be bosonic with $\sup|G|<\infty$.
Its values are scalars, so that $G\in\mathcal{G}^{\mathrm{ev}}(B)$.
Let $(g_1,\dots,g_r)$ be a placement of $B$ in a coordinate torus $\mathbb{T}=(\Gamma,\chi)$, and let $\beta\in C^2(\Gamma)$.
By \eqref{eq:setting-group-action}, $G^{g_j}(F)$ is $G\left(g_j^*F\right)$ or its complex conjugate,
so that $\left|G^{g_j}\right|\le\sup|G|$; and $|\varphi_\beta|=1$.
The bosonic integrand $\varphi_\beta\prod_jG^{g_j}$ therefore has modulus at most the constant $\left(\sup|G|\right)^r$.
The measure that defines $Z_{\mathbb{T}}$ on bosonic integrands is non-negative (Proposition~\ref{prop:accretivity}),
and the monotone form of the modulus bound in Corollary~\ref{cor:modulus-bound} yields
\[
\left|Z_{\mathbb{T}}\left[\varphi_\beta\prod_{j=1}^rG^{g_j}\right]\right|
\le Z_{\mathbb{T}}\left[\left(\sup|G|\right)^r\right]
=\left(\sup|G|\right)^r\cdot Z_{\mathbb{T}}[1].
\]

\emph{Proof of (d) and (d$'$).}
Consider first the direction $\mu=1$.
Fix $s\in\left\{\frac14,\frac18\right\}$,
and for a plane $c''\in\Z+\frac12$ and an element $X''$ of a mid-plane algebra
write $G^{\circ(c'')}_{X''}$ for $G^{(c'')}_{X''}$ if $s=\frac14$ and for $G'^{(c'')}_{X''}$ if $s=\frac18$;
thus $G^{\circ(c'')}_{X''}=G^{(c'')}\left[k^{(s)}_{X''}\right]$ (Definition~\ref{def:two-insertions}).
By Proposition~\ref{prop:insertions}(a) and Definition~\ref{def:two-insertions},
$G^{\circ(c)}_X$ lies in $\mathcal{G}^{\mathrm{ev}}\left(B(c,S,\Lambda)\right)$;
as $B(c,S,\Lambda)\subset B^{\#}(c,S,\Lambda)$ (Lemma~\ref{lem:covariant-box}(b)),
it is an element of $\mathcal{G}^{\mathrm{ev}}\left(B^{\#}(c,S,\Lambda)\right)$ (Lemma~\ref{lem:equal-integrands}).
Let $\mathbb{T}=(\Gamma,\chi)$ be the coordinate torus $C(n_1,n_2,n_3)$, $|n|=n_1n_2n_3$,
let $(g_1,\dots,g_r)$ be a placement of $B^{\#}(c,S,\Lambda)$ in $\mathbb{T}$, and let $\beta\in C^2(\Gamma)$.
Each $g_a\in W_{\mathrm{c}}$ maps the plane $\{x_1=c\}$ onto a plane $\{x_1=c_a\}$,
$c_a=g_a(c)=j_a+\frac12$ with $j_a\in\Z$.
Put $S_a:=g_aS$, $\Lambda_a:=g_a\Lambda$ and $X_a:=X^{g_a}\in\mathfrak{W}(S_a,\Lambda_a)$.

\emph{Images.}
By Lemma~\ref{lem:coordinate-images}\eqref{it:images-group},
\[
\left(G^{\circ(c)}_X\right)^{g_a}=G^{\circ(c_a)}_{X_a},
\qquad
\|X_a\|_\#=\|X\|_\#.
\]
The two sides of the first identity are equal integrands in the sense of Lemma~\ref{lem:equal-integrands},
both elements of $\mathcal{G}^{\mathrm{ev}}$ of the box $B^{\#}(c_a,S_a,\Lambda_a)=g_aB^{\#}(c,S,\Lambda)$
(Lemma~\ref{lem:covariant-box}(a)),
and this box embeds in $\Gamma$ by condition~\ref{it:placement-disjoint}.
They therefore define the same function on $\Conf(\Gamma)$, and
\[
Z_{\mathbb{T}}\left[\varphi_\beta\prod_{a=1}^r\left(G^{\circ(c)}_X\right)^{g_a}\right]
=Z_{\mathbb{T}}\left[\varphi_\beta\prod_{a=1}^rG^{\circ(c_a)}_{X_a}\right].
\]
By Lemma~\ref{lem:covariant-box}(c), the boxes $B(c_a,S_a,\Lambda_a)$, $a=1,\dots,r$,
satisfy conditions~\ref{it:step-disjoint} and~\ref{it:step-string} of Proposition~\ref{prop:word-formula},
the latter with the plaquette $\sigma_0$ of the slice torus $\Lambda_{\mathrm{s}}$
that is the projection to the slice of a plaquette orthogonal to $e_1$ as in condition~\ref{it:placement-columns};
and $\hull(S_a,\Lambda_a)$ embeds in the cross-section $\Xi:=\Xi(\mathbb{T},x_1)$.

\emph{Reduction to one period.}
Let $j'_a\in\{0,\dots,n_1-1\}$ be the residue of $j_a$ modulo $n_1$,
put $c'_a:=j'_a+\frac12$,
and let $v_a:=(j_a-j'_a)e_1$, an element of $n_1\Z e_1\subset\Gamma$.
First, the set of cells that $B(c_a,S_a,\Lambda_a)$ is required to contain in Definition~\ref{def:insertions}
is the translate by $v_a$ of the set for $B(c'_a,S_a,\Lambda_a)$,
and the rule that selects the box among the boxes of least radius containing this set commutes with translations;
hence $B(c_a,S_a,\Lambda_a)=B(c'_a,S_a,\Lambda_a)+v_a$.
Conditions~\ref{it:step-disjoint} and~\ref{it:step-string} depend on the boxes only through the family of their translates by elements of $\Gamma$,
and this family is the same for the planes $c_a$ and $c'_a$;
so the boxes $B(c'_a,S_a,\Lambda_a)$ satisfy both conditions, with the same $\sigma_0$,
and $\hull(S_a,\Lambda_a)$ is unchanged.
Second, Lemma~\ref{lem:coordinate-images}\eqref{it:images-translation}, applied with the vector $v_a$, gives
\[
\left(G^{\circ(c'_a)}_{X_a}\right)^{\tau_{v_a}}=G^{\circ(c_a)}_{X_a}.
\]
Here the slice part of $v_a$ is $0$,
so that the pull-back $\tau^*_{v_{\mathrm{s}}}$ is the identity,
and the unitary $U_{v_a}$ of Lemma~\ref{lem:translated-insertions} is the identity of $\mathcal{F}_{\Lambda_a}$,
the signs $\operatorname{sgn}_{v_a}(I)$ being $1$ for a vector parallel to $e_1$;
hence the kernel $k^{(s)}_{X_a}$ is carried to itself, and only the plane moves from $c'_a$ to $c'_a+j_a-j'_a=c_a$.
Third, the integrand $G^{\circ(c'_a)}_{X_a}$ embeds in $\mathbb{T}$,
since its box is the translate by $-v_a\in\Gamma$ of the box $B(c_a,S_a,\Lambda_a)$, which embeds;
as $v_a\in\Gamma$, it defines the same function on $\Conf(\Gamma)$ as its translate by $v_a$
(Lemma~\ref{lem:admissible-images}\eqref{it:images-periods}).
Together,
\[
Z_{\mathbb{T}}\left[\varphi_\beta\prod_{a=1}^rG^{\circ(c_a)}_{X_a}\right]
=Z_{\mathbb{T}}\left[\varphi_\beta\prod_{a=1}^rG^{\circ(c'_a)}_{X_a}\right].
\]

\emph{The word and the bound.}
The planes $c'_a=j'_a+\frac12$ satisfy $0\le j'_a<n_1$,
and their boxes satisfy conditions~\ref{it:step-disjoint} and~\ref{it:step-string}.
Lemma~\ref{lem:mixed-word-formula}(b), applied with the plaquette $\sigma_0$,
with the elements $X_a\in\mathfrak{W}(S_a,\Lambda_a)$ inserted at the planes $c'_a$, all of the kind $s$,
and with $\beta$, gives
\[
Z_{\mathbb{T}}\left[\varphi_\beta\prod_{a=1}^rG^{\circ(c'_a)}_{X_a}\right]
=(2\pi)^{|n|}\cdot\Tr_{\mathcal{H}^{\mathrm{l}}}\prod_{j=n_1-1,\dots,0}
\left(K^{\frac12}_{\beta_j,S_j}\,\mathcal{B}_j\,K^{\frac12}_{\beta_j,S_j}\,\mathfrak{M}\,u_j\,\mathfrak{M}^{*}\right).
\]
Here $\beta_j\in\R^{\Lambda^{(1)}_{\mathrm{s}}}$ is the electric part of $\beta$ at the step $j$,
$u_j=\exp\left(i\left\langle\beta^{\mathrm{mag}}_j,\cdot\right\rangle\right)$,
the slice operator $\mathfrak{M}$ is formed with $\sigma_0$,
\[
\mathcal{B}_j:=\prod_{a:\,j'_a=j}\left(\alpha_{\beta_j}X_a\right)_\Xi
\]
is a product of commuting factors, equal to $1$ if no $j'_a$ equals $j$,
in which $\alpha_{\beta_j}X_a$ is formed with the restriction of $\beta_j$ to $S_a$,
and $S_j:=\emptyset$ if $s=\frac14$, $S_j:=\bigcup_{a:\,j'_a=j}S_a$ if $s=\frac18$.
Proposition~\ref{prop:holder-insertions}, applied with these $\beta_j$ and $S_j$,
with the functions $u_j$, for which $|u_j|=1$,
and with the bounded operators $\mathcal{B}_j$ on $\mathcal{H}^{\mathrm{l}}$,
bounds the modulus of the trace by
\[
\prod_{j=0}^{n_1-1}\left(\|\mathcal{B}_j\|\cdot2^{\frac{\#|S_j|}{2}}\right)\cdot(2\pi)^{-|n|}\cdot Z_{\mathbb{T}}[1].
\]
Let $r_j:=\#\left|\left\{a:j'_a=j\right\}\right|$, so that $\sum_jr_j=r$.
By Lemma~\ref{lem:sharp-norm}, $\left\|\left(\alpha_{\beta_j}X_a\right)_\Xi\right\|\le\|X_a\|_\#=\|X\|_\#$,
and hence $\|\mathcal{B}_j\|\le\|X\|_\#^{r_j}$.
For $s=\frac18$, the boxes $B(c'_a,S_a,\Lambda_a)$ with $j'_a=j$ contain the links $(j,b)$, $b\in S_a$,
of the plane $\{x_1=j\}$, and they embed disjointly in $\Gamma$;
so these sets $S_a$ are mapped injectively onto pairwise disjoint subsets of $\Lambda^{(1)}_{\mathrm{s}}$,
each of cardinality $\#|S|$, and $\#|S_j|=r_j\cdot\#|S|$.
Collecting the three identities and the bound, one obtains for $s=\frac14$
\[
\left|Z_{\mathbb{T}}\left[\varphi_\beta\prod_{a=1}^r\left(G^{(1,c)}_X\right)^{g_a}\right]\right|
\le\|X\|_\#^r\cdot Z_{\mathbb{T}}[1],
\]
and for $s=\frac18$
\[
\left|Z_{\mathbb{T}}\left[\varphi_\beta\prod_{a=1}^r\left(G'^{(1,c)}_X\right)^{g_a}\right]\right|
\le\left(2^{\frac{\#|S|}{2}}\cdot\|X\|_\#\right)^r\cdot Z_{\mathbb{T}}[1].
\]
The torus $\mathbb{T}$, the placement and $\beta$ being arbitrary,
$G^{(1,c)}_X$ is bounded by $\|X\|_\#$ and $G'^{(1,c)}_X$ by $2^{\frac{\#|S|}{2}}\cdot\|X\|_\#$,
both with the box $B^{\#}(c,S,\Lambda)$.
For $\mu=2,3$, statement (a) with $g=g_\mu$ carries these two bounds
to $G^{(\mu,c)}_X=\left(G^{(1,c)}_X\right)^{g_\mu}$ and $G'^{(\mu,c)}_X=\left(G'^{(1,c)}_X\right)^{g_\mu}$,
with the box $g_\mu B^{\#}(c,S,\Lambda)$.
\end{proof}

\subsection{Coarse structures}\label{sec:coarse-structures}

\begin{lemma}\label{lem:coarse-structures}
Let $\mathfrak{n}=(n_i)$ be a sequence of positive integers with $n_i\to\infty$
that satisfies the divisibility condition \eqref{cond:divisibility} of Definition~\ref{def:sides}.
\begin{enumerate}
\item\label{it:coarse-widths}
The set of even integers $h$ with $4h\in\mathcal{D}(\mathfrak{n})$ meets every half-line $[h_0,\infty)$, $h_0\in\R$.
If $h$ belongs to this set and $r>0$,
then there are $h''\in h\Z$ and an integer $k$ with $kh\in h''\Z$ such that
\begin{equation*}
h''\ge 2r+6h,\qquad kh\ge 16h'',\qquad 4h''\in\mathcal{D}(\mathfrak{n}),\qquad 2kh\in\mathcal{D}(\mathfrak{n}).
\end{equation*}
For data $\mathfrak{d}$ of width $h$ (Definition~\ref{def:data}) and for $r>0$,
let $h''(r)$ be the least number $h''$ for which such an integer $k$ exists,
let $k(\mathfrak{d},r)$ be the least such $k$ for $h''=h''(r)$,
and let $i_0(\mathfrak{d},r)$ be the least index $j$ such that $n_i$ is divisible by $4h$, by $4h''(r)$ and by $2k(\mathfrak{d},r)\cdot h$ for every $i\ge j$.
These three numbers depend only on $\mathfrak{n}$, $h$ and $r$;
in particular they do not depend on the positions $c_0$, $c_+$, $c_3$ or on the observable $O$ of $\mathfrak{d}=(c_0,c_+,c_3;h;O)$,
on the position of a box of radius $r$, or on the strength $t$ of the functionals $\mu^t_i$.
Moreover $i_0(\mathfrak{d},r)\ge i_0(\mathfrak{d})$ (Definition~\ref{def:threshold}).
\end{enumerate}
Let $\mathfrak{d}=(c_0,c_+,c_3;h;O)$ be data, let $B=y+[-r,r]^3$ be a box with $y\in\R^3$ and $r>0$,
and put $h'':=h''(r)$ and $k:=k(\mathfrak{d},r)$.
There are a cell structure $\mathcal{Q}'=(\mathfrak{F}_{\mathrm{d}};c'_-,c'_+,c'_3;kh)$ of the diagonal frame
and a cell structure $\mathcal{Q}''=(\mathfrak{F}_{\mathrm{c}};o_1,o_2,o_3;h'')$ of the coordinate frame,
with $o_1,o_2,o_3\in\Z+\tfrac12$, such that the following hold.
\begin{enumerate}
\setcounter{enumi}{1}
\item\label{it:coarse-diagonal}
$\mathcal{Q}'$ is coarser than $\mathcal{Q}_{\mathfrak{d}}$, and $B$ lies in an open cell of $\mathcal{Q}'$.
\item\label{it:coarse-coordinate}
For every $g\in W(\mathcal{Q}')$ the box $gB$ lies in an open cell of $\mathcal{Q}''$,
and $gB$, $g'B$ lie in different cells whenever $g$ and $g'$ are different elements of $W(\mathcal{Q}')$.
\item\label{it:coarse-walls}
Each wall of $\mathcal{Q}''$ is disjoint from $g(\supp O)$ for every $g\in W(\mathcal{Q}_{\mathfrak{d}})$.
\item\label{it:coarse-tori}
Let $i\ge i_0(\mathfrak{d},r)$, and let $\Gamma$ be the period lattice of $D(n_i)$.
Then $D(n_i)$ carries $\mathcal{Q}_{\mathfrak{d}}$ and $\mathcal{Q}'$, and $\Gamma$ lies in the translation subgroup of $W(\mathcal{Q}')$;
$\mathcal{Q}''$ cuts $D(n_i)$, and $C(n_i/2,2n_i,n_i)$ carries $\mathcal{Q}''$;
every box of radius $r$, in particular $B$, embeds in $D(n_i)$ and in $C(n_i/2,2n_i,n_i)$.
If $g,g'\in W(\mathcal{Q}')$ and $g'\neq\tau_vg$ for every $v\in\Gamma$,
then the cells of $\mathcal{Q}''$ that contain $gB$ and $g'B$ do not differ by a translation by a vector of $\Gamma$.
The torus $D(n_i)$ has $n_i^3/(|Q_0|\,k^3)$ cells of $\mathcal{Q}'$,
and each of the tori $D(n_i)$ and $C(n_i/2,2n_i,n_i)$ has $n_i^3/h''^3$ cells of $\mathcal{Q}''$.
\item\label{it:coarse-placement}
Let $i\ge i_0(\mathfrak{d},r)$, let $\Gamma'$ be the period lattice of $C(n_i/2,2n_i,n_i)$, let $g\in W(\mathcal{Q}')$,
and let $(g''_j)_j$ be a family of elements of $W(\mathcal{Q}'')$ with exactly one member in each coset of $\Gamma'$ in $W(\mathcal{Q}'')$.
Then $(g''_j)_j$ is a placement of $gB$ in $C(n_i/2,2n_i,n_i)$.
\end{enumerate}
\end{lemma}

\begin{proof}
\emph{Item~\eqref{it:coarse-widths}.}
By Definition~\ref{def:sides} the set $\mathcal{D}(\mathfrak{n})$ is closed under divisors and under least common multiples;
by \eqref{cond:divisibility} it contains $8$ and is infinite, so that it contains integers as large as one wishes.
Given $h_0$, choose an element of $\mathcal{D}(\mathfrak{n})$ that is at least $4h_0$,
and define $h$ by the requirement that $4h$ be the least common multiple of $8$ and of this element.
Then $4h\in\mathcal{D}(\mathfrak{n})$;
$h$ is even, because $8$ divides $4h$;
and $h$ is at least one quarter of the chosen element, hence $h\ge h_0$.

Now let $h$ be even with $4h\in\mathcal{D}(\mathfrak{n})$, and let $r>0$.
Choose an element of $\mathcal{D}(\mathfrak{n})$ that is at least $4(2r+6h)$,
and let $4h''$ be the least common multiple of $4h$ and of this element.
Then $h''$ is a multiple of $h$, $4h''\in\mathcal{D}(\mathfrak{n})$,
and $h''$ is at least one quarter of the chosen element, so that $h''\ge 2r+6h$.
Choose next an element of $\mathcal{D}(\mathfrak{n})$ that is at least $32h''$,
and let $2kh$ be the least common multiple of $4h''$ and of this element;
$k$ is an integer, since $2h$ divides $4h''$.
Then $kh$ is a multiple of $2h''$, $2kh\in\mathcal{D}(\mathfrak{n})$,
and $kh$ is at least one half of the chosen element, so that $kh\ge 16h''$.

Consequently the numbers $h''$ with the required properties form a non-empty set of positive integers,
and $h''(r)$ is its least element.
For $h''=h''(r)$ the admissible integers $k$ form a non-empty set of positive integers (note that $kh\ge 16h''>0$),
and $k(\mathfrak{d},r)$ is its least element.
The numbers $4h$, $4h''(r)$ and $2k(\mathfrak{d},r)\cdot h$ lie in $\mathcal{D}(\mathfrak{n})$;
each of them divides $n_i$ for all but finitely many $i$, and so $i_0(\mathfrak{d},r)$ exists.
None of these three definitions involves anything but $\mathfrak{n}$, $h$ and $r$.
Finally, $4h$ divides $n_i$ for every $i\ge i_0(\mathfrak{d},r)$,
whereas $i_0(\mathfrak{d})$ is the least index $j$ such that $4h$ divides $n_i$ for every $i\ge j$ (Definition~\ref{def:threshold});
hence $i_0(\mathfrak{d})\le i_0(\mathfrak{d},r)$.

\emph{The structure $\mathcal{Q}'$, and item~\eqref{it:coarse-diagonal}.}
In the rest of the proof $h''=h''(r)$ and $k=k(\mathfrak{d},r)$.
Every coset of $h\Z$ in $\R$ meets every half-open interval of length $h$ in exactly one point.
Hence there are $c'_-\in c_0+h\Z$, $c'_+\in c_++h\Z$ and $c'_3\in c_3+h\Z$ with
\begin{align*}
u_-&:=\ell_-(y)-c'_-\in[2h'',2h''+h),\\
u_+&:=\ell_+(y)-c'_+\in[4h'',4h''+h),\\
u_3&:=y_3-c'_3\in[h'',h''+h).
\end{align*}
As $c_0,c_+\in\Z$ and $c_3\in\Z+\tfrac12$ (Definition~\ref{def:data}),
the positions $c'_\pm$ lie in $\Z=\Adm(\ell_\pm)$ and $c'_3$ lies in $\Z+\tfrac12\subset\Adm(x_3)$ (Definition~\ref{def:mirrors}).
Since moreover $kh\in h\Z$, the triple of positions and the width $kh$ define a cell structure
$\mathcal{Q}'=(\mathfrak{F}_{\mathrm{d}};c'_-,c'_+,c'_3;kh)$ that is coarser than $\mathcal{Q}_{\mathfrak{d}}$
(Definition~\ref{def:cell-structures}).

Let $x\in B$.
Then $|x_\mu-y_\mu|\le r$ for $\mu=1,2,3$, hence $|\ell_\pm(x)-\ell_\pm(y)|\le 2r$, and
\begin{align*}
\ell_-(x)-c'_-&\in[2h''-2r,\,2h''+h+2r],\\
\ell_+(x)-c'_+&\in[4h''-2r,\,4h''+h+2r],\\
x_3-c'_3&\in[h''-r,\,h''+h+r].
\end{align*}
From $h''\ge 2r+6h$ one obtains $2r<h''$ and $h+2r<2h''$.
So the three intervals are contained in $(0,6h'')$, and $6h''<16h''\le kh$.
Thus $B$ lies in the open cell
\begin{equation*}
Q'_0:=\{x\in\R^3:\ 0<\ell_-(x)-c'_-<kh,\ \ 0<\ell_+(x)-c'_+<kh,\ \ 0<x_3-c'_3<kh\}
\end{equation*}
of $\mathcal{Q}'$.

\emph{The group $W(\mathcal{Q}')$.}
For the three families of the diagonal frame, the reflections of Definition~\ref{def:mirrors} are
\begin{equation*}
\theta_{\ell_-,c}(x)=x-(\ell_-(x)-c)\,e_-,\qquad
\theta_{\ell_+,c}(x)=x-(\ell_+(x)-c)\,e_+,\qquad
\theta_{x_3,c}(x)=x-2(x_3-c)\,e_3,
\end{equation*}
because $|e_\pm|^2=2$ and $|e_3|^2=1$.
By Definition~\ref{def:mirrors}\eqref{it:mirror-product}, the reflections in the walls of $\mathcal{Q}'$ are the maps
\begin{equation*}
\tau_{jkh\,e_-}\,\theta_{\ell_-,c'_-},\qquad
\tau_{jkh\,e_+}\,\theta_{\ell_+,c'_+},\qquad
\tau_{2jkh\,e_3}\,\theta_{x_3,c'_3}\qquad(j\in\Z).
\end{equation*}
The vectors $e_-$, $e_+$, $e_3$ are pairwise orthogonal,
so that $\ell_-$, $\ell_+$, $x_3$ are linear coordinates on $\R^3$.
The involution $\theta_{\ell_-,c'_-}$ replaces $\ell_-$ by $2c'_--\ell_-$ and leaves $\ell_+$ and $x_3$ unchanged;
in the same way $\theta_{\ell_+,c'_+}$ and $\theta_{x_3,c'_3}$ change only $\ell_+$, respectively only $x_3$.
Hence the three involutions commute with each other,
and the linear part of each of them reverses one of the vectors $e_-$, $e_+$, $e_3$ and fixes the other two;
in particular these linear parts map the lattice $kh(\Z e_-+\Z e_++2\Z e_3)$ onto itself.
Moving translations to the left through the three involutions changes their vectors only by these linear parts;
it follows that the elements
\begin{equation}\label{eq:amplification-wall-group}
g=\tau_{kh(z_-e_-+z_+e_++2z_3e_3)}\,\theta_{\ell_-,c'_-}^{a_-}\,\theta_{\ell_+,c'_+}^{a_+}\,\theta_{x_3,c'_3}^{a_3},
\qquad (z_-,z_+,z_3)\in\Z^3,\quad a=(a_-,a_+,a_3)\in\{0,1\}^3,
\end{equation}
form a group.
This group contains the reflections in the walls of $\mathcal{Q}'$,
and each of its elements is a product of such reflections
(for instance $\tau_{kh\,e_-}=\theta_{\ell_-,c'_-+kh}\,\theta_{\ell_-,c'_-}$);
so it is $W(\mathcal{Q}')$.
An element $g\in W(\mathcal{Q}')$ has only one representation \eqref{eq:amplification-wall-group}:
its linear part determines $a$, and then its translation part determines $(z_-,z_+,z_3)$.

Write $u_-=2h''+s_-$, $u_+=4h''+s_+$ and $u_3=h''+s_3$ with $s_-,s_+,s_3\in[0,h)$.
For $g$ as in \eqref{eq:amplification-wall-group},
\begin{equation}\label{eq:amplification-image-centre}
gy=y-a_-u_-\,e_--a_+u_+\,e_+-2a_3u_3\,e_3+kh\,(z_-e_-+z_+e_++2z_3e_3)=y+h''\cdot\kappa(g)+s(a),
\end{equation}
where
\begin{align*}
\kappa(g)&=\kappa_0(a)+\frac{kh}{h''}\,(z_-+z_+,\ z_+-z_-,\ 2z_3),\qquad
\kappa_0(a)=(-2a_--4a_+,\ 2a_--4a_+,\ -2a_3),\\
s(a)&=(-a_-s_--a_+s_+,\ a_-s_--a_+s_+,\ -2a_3s_3).
\end{align*}
Since $kh\in h''\Z$, the vector $\kappa(g)$ lies in $\Z^3$;
and $|s(a)|_\infty<2h$, because $s_-$, $s_+$, $s_3$ lie in $[0,h)$.
The linear part of $g$ belongs to $\Ocub$ and maps the cube $[-r,r]^3$ onto itself, so that $gB=gy+[-r,r]^3$.

The map $g\mapsto\kappa(g)$ is injective on $W(\mathcal{Q}')$.
Indeed, the eight vectors $\kappa_0(a)$ are the points of $\{(0,0),(-2,2),(-4,-4),(-6,-2)\}\times\{0,-2\}$;
they are pairwise different, and any two of them differ by a vector of sup norm at most $6$.
Let $g,g'\in W(\mathcal{Q}')$ have the integers $(z_-,z_+,z_3)$, $(z'_-,z'_+,z'_3)$ and the exponents $a$, $a'$,
and suppose that $\kappa(g)=\kappa(g')$.
Then $\kappa_0(a')-\kappa_0(a)$ equals $kh/h''$ times the integer vector
\begin{equation*}
(z_-+z_+,\ z_+-z_-,\ 2z_3)-(z'_-+z'_+,\ z'_+-z'_-,\ 2z'_3).
\end{equation*}
As $kh/h''\ge16$ and the sup norm of $\kappa_0(a')-\kappa_0(a)$ is at most $6$, this integer vector vanishes,
which gives $z'_-=z_-$, $z'_+=z_+$ and $z'_3=z_3$.
Then $\kappa_0(a)=\kappa_0(a')$, so $a=a'$, and $g=g'$.

\emph{The structure $\mathcal{Q}''$.}
For $\mu=1,2,3$ put
\begin{equation*}
I_\mu:=(y_\mu+r+2h-h'',\ y_\mu-r-2h),
\end{equation*}
an open interval of length $h''-2r-4h\ge 2h$.
Let $y^{O}$ be the centre of a cube $y^{O}+[-r_O,r_O]^3$ that contains $\supp O$ (Definition~\ref{def:characters}), and put
\begin{equation*}
u^O_-:=\ell_-(y^O)-c_0,\qquad u^O_+:=\ell_+(y^O)-c_+,\qquad u^O_3:=y^O_3-c_3.
\end{equation*}
The walls of $\mathcal{Q}_{\mathfrak{d}}$ are the planes $\{\ell_-=c_0+jh\}$, $\{\ell_+=c_++jh\}$ and $\{x_3=c_3+jh\}$, $j\in\Z$.
The argument that led to \eqref{eq:amplification-wall-group} and to the first equality in \eqref{eq:amplification-image-centre},
carried out with $(c_0,c_+,c_3;h)$ in place of $(c'_-,c'_+,c'_3;kh)$,
shows that every $g\in W(\mathcal{Q}_{\mathfrak{d}})$ satisfies
\begin{equation*}
gy^O=y^O-a_-u^O_-\,e_--a_+u^O_+\,e_+-2a_3u^O_3\,e_3+h\,(z_-e_-+z_+e_++2z_3e_3)
\end{equation*}
for some $(z_-,z_+,z_3)\in\Z^3$ and some $a\in\{0,1\}^3$.
Modulo $h$, the three coordinates of $gy^O$ are therefore congruent to
\begin{equation*}
y^O_1-a_-u^O_--a_+u^O_+,\qquad y^O_2+a_-u^O_--a_+u^O_+,\qquad y^O_3-2a_3u^O_3;
\end{equation*}
as $g$ ranges over $W(\mathcal{Q}_{\mathfrak{d}})$, the first and the second coordinate of $gy^O$ take at most four values modulo $h$,
and the third at most two.
The image of the cube under $g$ is $gy^O+[-r_O,r_O]^3$, and it contains the image of $\supp O$.
Hence, for each $\mu$, the $\mu$-th coordinate of a point of an image of $\supp O$
lies within distance $r_O$ of a number that is congruent modulo $h$ to one of at most four given numbers.
Since $r_O$ is an integer, a closed interval of length $2r_O$ contains at most $2r_O+1$ points of $\Z+\tfrac12$.
Therefore at most $4(2r_O+1)=8r_O+4$ of the $h$ classes of $\Z+\tfrac12$ modulo $h$
contain the $\mu$-th coordinate of a point of an image of $\supp O$ under $W(\mathcal{Q}_{\mathfrak{d}})$;
call the other classes free for $\mu$.
Free classes exist, because $8r_O+4<h$ by Definition~\ref{def:data}.
Since the open interval $I_\mu$ has length at least $2h$, hence greater than $h$, it meets each class of $\Z+\tfrac12$ modulo $h$,
and $o_\mu$ is taken to be a point of $I_\mu$ that lies in a class free for $\mu$.
Let $\mathcal{Q}''=(\mathfrak{F}_{\mathrm{c}};o_1,o_2,o_3;h'')$;
it is a cell structure of the coordinate frame, as $o_\mu\in\Z+\tfrac12\subset\Adm(x_\mu)$.

\emph{Item~\eqref{it:coarse-coordinate}.}
Let $g\in W(\mathcal{Q}')$.
As $o_\mu\in I_\mu$, the interval $[y_\mu-r-2h,\,y_\mu+r+2h]$ is contained in $(o_\mu,\,o_\mu+h'')$.
Let $\kappa_\mu(g)$ denote the $\mu$-th coordinate of $\kappa(g)$.
By \eqref{eq:amplification-image-centre} and $|s(a)|_\infty<2h$,
the $\mu$-th coordinates of the points of $gB=gy+[-r,r]^3$ lie in
$[y_\mu-r-2h,\,y_\mu+r+2h]+h''\cdot\kappa_\mu(g)$,
hence in $(o_\mu+h''\cdot\kappa_\mu(g),\,o_\mu+h''\cdot(\kappa_\mu(g)+1))$.
So $gB$ lies in the open cell $Q''_{\kappa(g)}$ of $\mathcal{Q}''$ with index $\kappa(g)$,
and different elements of $W(\mathcal{Q}')$ give different cells, by the injectivity of $\kappa$.

\emph{Item~\eqref{it:coarse-walls}.}
The walls of $\mathcal{Q}''$ orthogonal to $e_\mu$ are the planes $\{x_\mu=o_\mu+jh''\}$, $j\in\Z$,
and $o_\mu+jh''$ lies in the class of $o_\mu$ modulo $h$, because $h''\in h\Z$.
A point of an image of $\supp O$ on one of these planes would have its $\mu$-th coordinate in a class that is free for $\mu$,
which is impossible.

\emph{Item~\eqref{it:coarse-tori}.}
Let $i\ge i_0(\mathfrak{d},r)$; then $4h$, $4h''$ and $2kh$ divide $n_i$.
By Definitions~\ref{def:tori}, \ref{def:sides} and~\ref{def:cell-structures},
$D(n_i)=\mathbb{T}(\mathfrak{F}_{\mathrm{d}};n_i,2n_i,n_i)$ and
$C(n_i/2,2n_i,n_i)=\mathbb{T}(\mathfrak{F}_{\mathrm{c}};n_i/2,2n_i,n_i)$, with the period lattices
\begin{equation*}
\Gamma=\Z\,(n_i/2)\,e_-+\Z\,n_ie_++\Z\,n_ie_3,\qquad
\Gamma'=\Z\,(n_i/2)\,e_1+\Z\,2n_ie_2+\Z\,n_ie_3.
\end{equation*}
The height periods $n_i$, $2n_i$, $n_i$ of $D(n_i)$ are multiples of $2h$ and of $2kh$.
Hence $D(n_i)$ carries $\mathcal{Q}_{\mathfrak{d}}$ and $\mathcal{Q}'$,
and $\Gamma$ lies in the translation subgroup of $W(\mathcal{Q}')$ (Definition~\ref{def:cell-structures}).
The height periods $n_i/2$, $2n_i$, $n_i$ of $C(n_i/2,2n_i,n_i)$ are multiples of $2h''$;
so $C(n_i/2,2n_i,n_i)$ carries $\mathcal{Q}''$.
Every coordinate of a vector of $\Gamma$ or of $\Gamma'$ is a multiple of $n_i/2$, and $n_i/2$ lies in $h''\Z$.
In particular $x_\mu(v)\in h''\Z$ for $\mu=1,2,3$ and for every $v\in\Gamma$ or $v\in\Gamma'$;
that is, $\mathcal{Q}''$ cuts $D(n_i)$, and it also cuts $C(n_i/2,2n_i,n_i)$.
Moreover every non-zero vector $v$ of $\Gamma$ or of $\Gamma'$ satisfies
\begin{equation*}
|v|_\infty\ \ge\ \frac{n_i}{2}\ \ge\ kh\ \ge\ 16h''\ \ge\ 16(2r+6h)\ >\ 2r+2,
\end{equation*}
where $n_i/2\ge kh$ because $2kh$ divides $n_i$.
For $y'\in\R^3$ the sets $y'+[-r,r]^3+[-1,1]^3+v=y'+v+[-r-1,r+1]^3$, with $v\in\Gamma$ (resp. $v\in\Gamma'$),
are cubes of radius $r+1$ whose centres differ by vectors of sup norm greater than $2r+2$;
they are pairwise disjoint.
So the box $y'+[-r,r]^3$ embeds in $D(n_i)$ and in $C(n_i/2,2n_i,n_i)$ (Definition~\ref{def:integrands}).

Let $g,g'\in W(\mathcal{Q}')$, let $Q''$ be the open cell of $\mathcal{Q}''$ that contains $gB$,
and suppose that $g'B$ lies in $Q''+v$ for some $v\in\Gamma$.
As $\mathcal{Q}''$ cuts $D(n_i)$, translations by vectors of $\Gamma$ permute the open cells of $\mathcal{Q}''$,
and $Q''+v$ is again an open cell (Definition~\ref{def:cell-structures}).
As $\Gamma$ lies in the translation subgroup of $W(\mathcal{Q}')$, the element $\tau_vg$ belongs to $W(\mathcal{Q}')$,
and $\tau_vgB=gB+v$ lies in $Q''+v$.
The boxes $\tau_vgB$ and $g'B$ thus lie in one cell, and item~\eqref{it:coarse-coordinate} gives $g'=\tau_vg$.
Hence, if $g'\neq\tau_vg$ for every $v\in\Gamma$, the cells of $\mathcal{Q}''$ that contain $gB$ and $g'B$
do not differ by a translation by a vector of $\Gamma$.

By Definition~\ref{def:tori}, $|D(n_i)|=2\cdot(n_i/2)\cdot n_i\cdot n_i=n_i^3$ and $|C(n_i/2,2n_i,n_i)|=n_i^3$.
A cell of $\mathcal{Q}'$ has $|Q'|=(kh)^3/2=k^3\cdot|Q_0|$ photon sites, and a cell of $\mathcal{Q}''$ has $|Q''|=h''^3$
(Definitions~\ref{def:cell-structures} and~\ref{def:data}).
The heights $\ell_-(v)$, $\ell_+(v)$, $x_3(v)$ of the vectors $v\in\Gamma$ are integer multiples of $n_i$,
hence lie in $kh\Z$; so $\mathcal{Q}'$ cuts $D(n_i)$, and $\mathcal{Q}''$ cuts $D(n_i)$ and $C(n_i/2,2n_i,n_i)$ by the above.
A torus that is cut by a cell structure has as many cells as the quotient of its number of photon sites by that of a cell
(Definition~\ref{def:cell-structures}); this gives the numbers $n_i^3/(|Q_0|\,k^3)$ and $n_i^3/h''^3$.

\emph{Item~\eqref{it:coarse-placement}.}
Let $i$, $\Gamma'$, $g$ and $(g''_j)_j$ be as in item~\eqref{it:coarse-placement}.
The walls of $\mathcal{Q}''$ are mirrors of the axis families; so $W(\mathcal{Q}'')\subset W_{\mathrm{c}}$,
and the linear parts of the elements of $W(\mathcal{Q}'')$ are diagonal matrices with entries $\pm1$ (Definition~\ref{def:mirrors}).
By item~\eqref{it:coarse-tori} the torus $C(n_i/2,2n_i,n_i)$ carries $\mathcal{Q}''$;
so the translations $\tau_v$, $v\in\Gamma'$, belong to $W(\mathcal{Q}'')$,
and $W(\mathcal{Q}'')/\Gamma'$ acts simply transitively on the $n_i^3/h''^3$ cells of $C(n_i/2,2n_i,n_i)$
(Definition~\ref{def:cell-structures}).
There are therefore exactly $n_i^3/h''^3$ cosets of $\Gamma'$ in $W(\mathcal{Q}'')$,
and $(g''_j)_j$ is a finite non-empty family in $W_{\mathrm{c}}$, with one member for each cell of $C(n_i/2,2n_i,n_i)$.
Since $\Gamma'$ is invariant under the diagonal matrices with entries $\pm1$,
the translations $\tau_v$, $v\in\Gamma'$, form a normal subgroup of $W(\mathcal{Q}'')$.
By normality, $\tau_vg''_j$ lies in the coset of $g''_j$ for every $v\in\Gamma'$;
hence the elements $\tau_vg''_j$, $v\in\Gamma'$, all $j$, are pairwise different.
By item~\eqref{it:coarse-coordinate}, $gB$ lies in an open cell $Q''$ of $\mathcal{Q}''$.
As $W(\mathcal{Q}'')$ acts simply transitively on the open cells of $\mathcal{Q}''$ (Definition~\ref{def:cell-structures}),
the open cells $\tau_vg''_jQ''=g''_jQ''+v$ are pairwise different, hence pairwise disjoint.
They contain the sets $g''_jgB+v$, which are therefore pairwise disjoint.
Each $g''_jgB$ is a box of radius $r$, since the linear part of $g''_jg$ lies in $\Ocub$;
by item~\eqref{it:coarse-tori} it embeds in $C(n_i/2,2n_i,n_i)$.
So the boxes $g''_jgB$ embed disjointly in $\Gamma'$ (Definition~\ref{def:integrands}),
and the family $(g''_j)_j$ satisfies condition~\eqref{it:placement-disjoint} of Definition~\ref{def:bounded-elements}.
These boxes lie in open cells of $\mathcal{Q}''$,
a cell structure of the coordinate frame with positions $o_\mu\in\Z+\tfrac12$
which cuts $C(n_i/2,2n_i,n_i)$ (the coordinates of the vectors of $\Gamma'$ lie in $h''\Z$).
By Lemma~\ref{lem:placements} the family therefore satisfies the condition~\eqref{it:placement-columns} on free columns.
Thus $(g''_j)_j$ is a placement of $gB$ in $C(n_i/2,2n_i,n_i)$.
\end{proof}

\subsection{Amplification}\label{sec:amplification-bound}

Let $k$ be a positive integer.
The amplification factor $\gamma_{t,i}(k)$ of Definition~\ref{def:amplification-factor} is defined for the indices $i$ with $n_i$ even and $n_i\ge8$.
Among them are all indices $i\ge i_0(\mathfrak{d},r)$, for every $r>0$:
such an $n_i$ is a positive multiple of $4h$ (Lemma~\ref{lem:coarse-structures}\eqref{it:coarse-widths}),
and $4h>32r_O+32\ge64$ (Definition~\ref{def:data}).
Suprema and upper limits over $i$ are taken over the indices for which the factor is defined.
It splits as the product
\[
\gamma_{t,i}(k)=\exp\left(2t\cdot|O|_1\cdot k^3\right)\cdot\exp\left(|Q_0|\cdot k^3\cdot\Delta_i\right)
\]
of a factor that does not contain $i$ and a factor that does not contain $t$.
The second factor is at least $1$ and tends to $1$ as $i\to\infty$, because $\Delta_i\ge0$ and $\Delta_i\to0$ (Theorem~\ref{thm:four-tori});
in particular it is bounded in $i$.
For fixed $k$ it follows that $\gamma_{t,i}(k)$ tends to $1$ as $t\to0$ and $i\to\infty$ in any manner,
so that $\lim_{t\to0}\limsup_i\gamma_{t,i}(k)=1$,
and that $\gamma_{t,i}(k)$ is bounded, uniformly in $i$, for $t$ in a bounded interval:
\begin{equation}\label{eq:amplification-factor-limit}
\lim_{t\to0,\ i\to\infty}\gamma_{t,i}(k)=1,
\qquad
\sup_{0\le t\le t_*}\ \sup_i\ \gamma_{t,i}(k)<\infty\qquad(0\le t_*<\infty).
\end{equation}

\begin{theorem}[Amplification]\label{thm:amplification}
Let $-1\le\rho\le1$, $w_0,\dots,w_3>0$, $M\in\R$ and $q>0$.
Let the sequence of sides $\mathfrak{n}=(n_i)$ satisfy \eqref{cond:divisibility},
and let $\mathfrak{d}=(c_0,c_+,c_3;h;O)$ be data (Definition~\ref{def:data}).
Let $G$ be bounded by $C$ with a box $B$ of radius $r$
(Definition~\ref{def:bounded-elements}, the functionals $Z_{\mathbb{T}}$ being formed with these parameters),
and let $k:=k(\mathfrak{d},r)$ (Lemma~\ref{lem:coarse-structures}).
Then, for all $t\ge0$ and all $i\ge i_0(\mathfrak{d},r)$,
\[
\left|\mu^t_i(G)\right|\le C\cdot\gamma_{t,i}(k).
\]
\end{theorem}

\begin{proof}
Fix $i\ge i_0(\mathfrak{d},r)$.
Let $\mathcal{Q}'=(\mathfrak{F}_{\mathrm{d}};c'_-,c'_+,c'_3;kh)$ and $\mathcal{Q}''=(\mathfrak{F}_{\mathrm{c}};o_1,o_2,o_3;h'')$, with $h''=h''(r)$,
be cell structures that Lemma~\ref{lem:coarse-structures} attaches to the box $B$,
and let $\Gamma$ and $\Gamma'$ be the period lattices of $D(n_i)$ and of $C(n_i/2,2n_i,n_i)$.
By item~\eqref{it:coarse-diagonal} of that lemma, $B$ lies in an open cell of $\mathcal{Q}'$.
Changing $c'_-$, $c'_+$ and $c'_3$ by integer multiples of $kh$ gives positions of walls of $\mathcal{Q}'$, which are admissible (Definition~\ref{def:cell-structures}),
and a cell structure with the same frame, width and walls, the same positions modulo $h$, the same open cells and the same group $W(\mathcal{Q}')$;
only the indices of the open cells are shifted.
The items of the lemma involve $\mathcal{Q}'$ through these data alone, and so they hold for the new positions as well.
The positions may therefore be taken such that $B$ lies in the open cell of index $0$,
\[
Q'_0=\left\{x\in\R^3:\ c'_-<\ell_-(x)<c'_-+kh,\ \ c'_+<\ell_+(x)<c'_++kh,\ \ c'_3<x_3<c'_3+kh\right\}.
\]
By Lemma~\ref{lem:coarse-structures}\eqref{it:coarse-tori},
the torus $D(n_i)$ carries $\mathcal{Q}_{\mathfrak{d}}$ and $\mathcal{Q}'$,
the lattice $\Gamma$ lies in the translation subgroup of $W(\mathcal{Q}')$,
the structure $\mathcal{Q}''$ cuts $D(n_i)$,
the torus $C(n_i/2,2n_i,n_i)$ carries $\mathcal{Q}''$,
and $B$ embeds in $D(n_i)$.
The torus $D(n_i)$ has $n_i^3/|Q_0|$ cells of $\mathcal{Q}_{\mathfrak{d}}$ (Definition~\ref{def:data}),
$n_i^3/(|Q_0|\,k^3)$ cells of $\mathcal{Q}'$ and $n_i^3/h''^3$ cells of $\mathcal{Q}''$,
and $C(n_i/2,2n_i,n_i)$ has $n_i^3/h''^3$ cells of $\mathcal{Q}''$ (Lemma~\ref{lem:coarse-structures}\eqref{it:coarse-tori}).
As $G\in\mathcal{G}^{\mathrm{ev}}(B)$ lies in $\mathcal{A}(\R^3)$ and embeds in $D(n_i)$, the number $\mu^t_i(G)$ is defined.

For an open cell $Q'$ of $\mathcal{Q}'$ let $O^{Q'}_{\mathfrak{d}}$ denote the sum of the mirror copies $O^{[\kappa]}$
over the $k^3$ open cells $Q_\kappa$ of $\mathcal{Q}_{\mathfrak{d}}$ contained in $Q'$ (Definition~\ref{def:cell-structures}).
For $\kappa'\in\Z^3$ let $g_{\kappa'}$ be the element of $W(\mathcal{Q}')$ with $g_{\kappa'}Q'_0=Q'_{\kappa'}$;
the mirror copy of $H\in\mathcal{A}(Q'_0)$ in $Q'_{\kappa'}$ with respect to $\mathcal{Q}'$ is then $H^{[\kappa']}=H^{g_{\kappa'}}$ (Definition~\ref{def:folds}).
Products over $\kappa'$ are taken over a set of representatives of the cells of $\mathcal{Q}'$ in $D(n_i)$ that contains $0$;
thus $Q'_0$ itself is one of the cells $Q'_{\kappa'}$.

\begin{enumerate}
\item\label{it:amplification-coarse-array}
\emph{The folded array as an array of $\mathcal{Q}'$.}
As $\mathcal{Q}'$ is coarser than $\mathcal{Q}_{\mathfrak{d}}$, every wall of $\mathcal{Q}'$ is a wall of $\mathcal{Q}_{\mathfrak{d}}$ (Definition~\ref{def:cell-structures}),
and $W(\mathcal{Q}')$ is a subgroup of $W(\mathcal{Q}_{\mathfrak{d}})$.
Let $g'\in W(\mathcal{Q}')$, and let $Q_\kappa\subset Q'_0$ be an open cell of $\mathcal{Q}_{\mathfrak{d}}$;
then $O^{[\kappa]}=O^{g_\kappa}$, where $g_\kappa$ is the element of $W(\mathcal{Q}_{\mathfrak{d}})$ with $g_\kappa Q_0=Q_\kappa$.
By \eqref{eq:setting-group-action}, $\left(O^{g_\kappa}\right)^{g'}=O^{g'g_\kappa}$,
and $g'g_\kappa$ is the element of $W(\mathcal{Q}_{\mathfrak{d}})$ that maps $Q_0$ onto the cell $g'Q_\kappa$,
since $W(\mathcal{Q}_{\mathfrak{d}})$ acts simply transitively on the open cells;
so the image of $O^{[\kappa]}$ under $g'$ is the mirror copy of $O$ in $g'Q_\kappa$.
The functions $O^{[\kappa]}$ are real and bounded, and each of them depends on finitely many $F_\sigma$ with $\sigma\subset Q_\kappa$ (Definition~\ref{def:arrays});
hence $\exp\bigl(t\,O^{Q'_0}_{\mathfrak{d}}\bigr)$ takes positive values, is bounded and,
being a Borel function of finitely many $F_\sigma$ with $\sigma\subset Q'_0$, belongs to $\mathcal{A}(Q'_0)$.
By \eqref{eq:setting-group-action} the image under $g'$ of a real function $V$ of $F$ is $F\mapsto V(g'^*F)$, also for odd $g'$, without complex conjugation,
and in particular $\exp(t\,V)^{g'}=\exp\left(t\,V^{g'}\right)$.
As $g'$ maps the open cells of $\mathcal{Q}_{\mathfrak{d}}$ contained in $Q'_0$ onto those contained in $g'Q'_0$,
the mirror copy of $\exp\bigl(t\,O^{Q'_0}_{\mathfrak{d}}\bigr)$ in the cell $Q'_{\kappa'}$ is
\[
\left(\exp\bigl(t\,O^{Q'_0}_{\mathfrak{d}}\bigr)\right)^{[\kappa']}=\exp\bigl(t\,O^{Q'_{\kappa'}}_{\mathfrak{d}}\bigr),
\]
the exponential of $t$ times the sum of the terms of $O_{\mathfrak{d}}$ in $Q'_{\kappa'}$.
On $D(n_i)$ the array $O_{\mathfrak{d}}$ is the sum of the $O^{[\kappa]}$ over representatives of the cells of $\mathcal{Q}_{\mathfrak{d}}$ in the torus (Definition~\ref{def:arrays}).
Every open cell of $\mathcal{Q}_{\mathfrak{d}}$ lies in exactly one open cell of $\mathcal{Q}'$,
and the translations by vectors of $\Gamma$ map open cells of either structure onto open cells of the same structure;
hence the open cells of $\mathcal{Q}_{\mathfrak{d}}$ contained in the representatives $Q'_{\kappa'}$ represent the cells of $\mathcal{Q}_{\mathfrak{d}}$ in $D(n_i)$.
Therefore
\[
\prod_{\kappa'}\exp\bigl(t\,O^{Q'_{\kappa'}}_{\mathfrak{d}}\bigr)=\exp\bigl(t\,O_{\mathfrak{d}}\bigr)\qquad\text{on }D(n_i).
\]

\item\label{it:amplification-chessboard}
\emph{The chessboard estimate in the perturbed functional.}
As the box $B$ lies in $Q'_0$, multiplying $G\in\mathcal{G}^{\mathrm{ev}}(B)$ by the bounded factor $\exp\bigl(t\,O^{Q'_0}_{\mathfrak{d}}\bigr)$
gives an element of $\mathcal{A}(Q'_0)$;
that the two factors read some plaquettes in common is admitted by Definition~\ref{def:integrands}.
Theorem~\ref{thm:chessboard}(a) applies on $D(n_i)$, which carries $\mathcal{Q}'$,
to the content $\exp\bigl(t\,O^{Q'_0}_{\mathfrak{d}}\bigr)\cdot G$ of the cell $Q'_0$
and the contents $\exp\bigl(t\,O^{Q'_{\kappa'}}_{\mathfrak{d}}\bigr)$ of the cells $Q'_{\kappa'}$ with $\kappa'\neq0$;
by step~\eqref{it:amplification-coarse-array}, the product of all these contents is $\exp\bigl(t\,O_{\mathfrak{d}}\bigr)\cdot G$.
For $H\in\mathcal{A}(Q'_0)$ and $\kappa_1,\kappa_2\in\Z^3$ one has $\left(H^{[\kappa_1]}\right)^{[\kappa_2]}=H^{[\kappa_2]}$,
by \eqref{eq:setting-group-action} and the simple transitivity of $W(\mathcal{Q}')$ on the open cells.
Together with step~\eqref{it:amplification-coarse-array} and Lemma~\ref{lem:admissible-images}\eqref{it:images-products},
this gives the folds on $D(n_i)$ with respect to $\mathcal{Q}'$:
\[
\Fold_{D(n_i)}\exp\bigl(t\,O^{Q'_{\kappa'}}_{\mathfrak{d}}\bigr)=\exp\bigl(t\,O_{\mathfrak{d}}\bigr),
\qquad
\Fold_{D(n_i)}\left(\exp\bigl(t\,O^{Q'_0}_{\mathfrak{d}}\bigr)\cdot G\right)=\exp\bigl(t\,O_{\mathfrak{d}}\bigr)\cdot\prod_{\kappa'}G^{[\kappa']}.
\]
Since $D(n_i)$ has $n_i^3/(|Q_0|\,k^3)$ cells of $\mathcal{Q}'$,
Theorem~\ref{thm:chessboard}(a) and Definition~\ref{def:folds} yield
\[
\left|Z_{D(n_i)}\left[\exp\bigl(t\,O_{\mathfrak{d}}\bigr)\cdot G\right]\right|
\le
Z_{D(n_i)}\left[\exp\bigl(t\,O_{\mathfrak{d}}\bigr)\cdot\prod_{\kappa'}G^{[\kappa']}\right]^{\frac{|Q_0|\,k^3}{n_i^3}}
\cdot
Z_{D(n_i)}\left[\exp\bigl(t\,O_{\mathfrak{d}}\bigr)\right]^{1-\frac{|Q_0|\,k^3}{n_i^3}},
\]
where both numbers under the roots are non-negative, by the same theorem.
The second of them is positive by step~\eqref{it:amplification-denominator},
and division by it gives
\begin{equation}\label{eq:amplification-power-bound}
\left|\mu^t_i(G)\right|^{\frac{n_i^3}{|Q_0|\,k^3}}
\le
\frac{Z_{D(n_i)}\left[\exp\bigl(t\,O_{\mathfrak{d}}\bigr)\cdot\prod_{\kappa'}G^{[\kappa']}\right]}{Z_{D(n_i)}\left[\exp\bigl(t\,O_{\mathfrak{d}}\bigr)\right]}.
\end{equation}

\item\label{it:amplification-denominator}
\emph{The denominator.}
On $D(n_i)$ the array $O_{\mathfrak{d}}$ is a real function with $\left|O_{\mathfrak{d}}\right|\le\frac{n_i^3}{|Q_0|}\cdot|O|_1$ (Definition~\ref{def:arrays}),
and on bosonic integrands $Z_{D(n_i)}$ is the integral against a non-negative measure (Proposition~\ref{prop:accretivity}\eqref{it:determinant-sign}).
Lemma~\ref{lem:perturbed-denominators}, applied to the function $t\,O_{\mathfrak{d}}$, and Theorem~\ref{thm:partition-positivity} give
\[
Z_{D(n_i)}\left[\exp\bigl(t\,O_{\mathfrak{d}}\bigr)\right]\ge\exp\left(-t\cdot|O|_1\cdot\frac{n_i^3}{|Q_0|}\right)\cdot Z_{D(n_i)}[1]>0.
\]

\item\label{it:amplification-numerator}
\emph{The numerator.}
Write $O=\sum_\beta a_\beta\cdot\varphi_\beta$, the sum over the finitely many pairwise different exponents $\beta$ of $O$,
so that $\sum_\beta|a_\beta|=|O|_1$ (Definition~\ref{def:characters}).
Let $\mathcal{N}$ be the set of the pairs $\lambda=(\kappa,\beta)$ formed by a representative $\kappa$ of a cell of $\mathcal{Q}_{\mathfrak{d}}$ in $D(n_i)$, as in Definition~\ref{def:arrays},
and an exponent $\beta$ of $O$.
For $\lambda=(\kappa,\beta)\in\mathcal{N}$ put
\[
a_\lambda:=a_\beta,
\qquad
\varphi_\lambda:=\varphi_{(g_\kappa^{-1})^*\beta},
\qquad\text{so that}\qquad
\varphi_\lambda(F)=\varphi_\beta\left(g_\kappa^*F\right);
\]
the support of $\varphi_\lambda$ is the image of $\supp\beta$ under $g_\kappa$,
and it lies in $g_\kappa(\supp O)$, an image of $\supp O$ under $W(\mathcal{Q}_{\mathfrak{d}})$.
As $O^{[\kappa]}(F)=O(g_\kappa^*F)$ (Definition~\ref{def:arrays}), one has on $D(n_i)$
\[
O_{\mathfrak{d}}=\sum_{\lambda\in\mathcal{N}}a_\lambda\cdot\varphi_\lambda,
\qquad
\sum_{\lambda\in\mathcal{N}}|a_\lambda|=\frac{n_i^3}{|Q_0|}\cdot|O|_1.
\]
For $\alpha=(j;\lambda_1,\dots,\lambda_j)$ with $j\ge0$ and $\lambda_1,\dots,\lambda_j\in\mathcal{N}$ put
\[
a_\alpha:=\frac{t^j}{j!}\cdot a_{\lambda_1}\cdots a_{\lambda_j},
\qquad
\varphi_\alpha:=\varphi_{\lambda_1}\cdots\varphi_{\lambda_j}.
\]
Then $\sum_\alpha|a_\alpha|=\exp\left(t\cdot|O|_1\cdot n_i^3/|Q_0|\right)$;
and since $|\varphi_\alpha|=1$, the partial sums $\sum_{j\le J}$ of the series $\sum_\alpha a_\alpha\cdot\varphi_\alpha$
converge to $\exp\bigl(t\,O_{\mathfrak{d}}\bigr)$ uniformly on the configurations of $D(n_i)$.

Put $H:=\prod_{\kappa'}G^{[\kappa']}$.
Its factors $G^{[\kappa']}=G^{g_{\kappa'}}$ lie in $\mathcal{G}^{\mathrm{ev}}(g_{\kappa'}B)$ (Lemma~\ref{lem:admissible-images}\eqref{it:images-boxes}),
with $g_{\kappa'}B\subset g_{\kappa'}Q'_0=Q'_{\kappa'}$.
These boxes lie in pairwise different open cells and are therefore pairwise disjoint; so $H\in\mathcal{A}(\R^3)$.
Moreover the cells $Q'_{\kappa'}$ lie in pairwise different $\Gamma$-orbits,
and $\ell_-(\Gamma)=n_i\Z$, $\ell_+(\Gamma)=2n_i\Z$ and $x_3(\Gamma)=n_i\Z$ are contained in $2kh\Z$, because $2kh$ divides $n_i$;
by Lemma~\ref{lem:embedding-criteria}\eqref{it:embedding-cells}, $H$ embeds in $D(n_i)$.
Proposition~\ref{prop:berezin}\eqref{it:torus-integrability} shows that $\Phi\mapsto Z_{D(n_i)}[\Phi\cdot H]$ is continuous for the supremum norm of bounded Borel functions $\Phi$;
consequently
\[
Z_{D(n_i)}\left[\exp\bigl(t\,O_{\mathfrak{d}}\bigr)\cdot H\right]=\sum_\alpha a_\alpha\cdot Z_{D(n_i)}\left[\varphi_\alpha\cdot H\right],
\]
the series being summed over $j$.

Let $\alpha=(j;\lambda_1,\dots,\lambda_j)$ be fixed,
and let $\beta_\alpha\in C^2_{\mathrm{c}}$ be the sum of the cochains of $\varphi_{\lambda_1},\dots,\varphi_{\lambda_j}$, so that $\varphi_\alpha=\varphi_{\beta_\alpha}$.
A plaquette $\sigma\subset\supp\beta_\alpha$ is contained in the support of one of the characters $\varphi_{\lambda_1},\dots,\varphi_{\lambda_j}$,
hence in an image of $\supp O$ under $W(\mathcal{Q}_{\mathfrak{d}})$;
as no wall of $\mathcal{Q}''$ meets this image (Lemma~\ref{lem:coarse-structures}\eqref{it:coarse-walls}), the plaquette $\sigma$ lies in an open cell of $\mathcal{Q}''$.
As a function of $F$, the character $\varphi_\alpha$ splits into characters of single plaquettes:
\[
\varphi_\alpha(F)=\prod_{\sigma\subset\supp\beta_\alpha}\exp\left(i\cdot\beta_\alpha(\sigma)\cdot F_\sigma\right).
\]
The remaining factors of $\varphi_\alpha\cdot H$ are the mirror copies $G^{[\kappa']}=G^{g_{\kappa'}}$, with $g_{\kappa'}\in W(\mathcal{Q}')$,
and the box $g_{\kappa'}B$ of each of them lies in an open cell of $\mathcal{Q}''$ (Lemma~\ref{lem:coarse-structures}\eqref{it:coarse-coordinate}).
Let $\kappa'_1\neq\kappa'_2$ be two of the representatives, and let $v\in\Gamma$.
The element $\tau_vg_{\kappa'_1}$ maps $Q'_0$ onto $Q'_{\kappa'_1}+v$,
which represents the same cell of $D(n_i)$ as $Q'_{\kappa'_1}$ and therefore differs from $Q'_{\kappa'_2}$;
so $g_{\kappa'_2}\neq\tau_vg_{\kappa'_1}$.
Lemma~\ref{lem:coarse-structures}\eqref{it:coarse-tori} now shows that the open cells of $\mathcal{Q}''$ that contain $g_{\kappa'_1}B$ and $g_{\kappa'_2}B$
are not translates of each other by vectors of $\Gamma$.

\emph{Contents of the cells of a brick.}
Theorem~\ref{thm:chessboard}(c), for the structure $\mathcal{Q}''$ and with $n_i/2$ and $n_i$ in the roles of its sides $a$ and $c$,
compares the tori $D(n_i/2,n_i,n_i)=D(n_i)$ and $C(n_i/2,2n_i,n_i)$.
Its hypotheses are satisfied:
$n_i/2$ and $n_i$ lie in $2h''\Z$, because $4h''$ divides $n_i$,
and $n_i\ge n_i/2\ge kh\ge16h''\ge4$, because $2kh$ divides $n_i$ (Lemma~\ref{lem:coarse-structures}\eqref{it:coarse-widths}).
By that theorem the brick
\[
\mathcal{K}'':=\left\{\kappa''\in\Z^3:\ 0\le\kappa''_1<\frac{n_i}{2h''},\ 0\le\kappa''_2<\frac{2n_i}{h''},\ 0\le\kappa''_3<\frac{n_i}{h''}\right\}
\]
is a set of representatives of the cells of $\mathcal{Q}''$ in $D(n_i)$, and of those in $C(n_i/2,2n_i,n_i)$; it has $n_i^3/h''^3$ elements.

Each factor of $\varphi_\alpha\cdot H$ belongs to $\mathcal{A}(Q'')$ for an open cell $Q''$ of $\mathcal{Q}''$:
a character of one plaquette $\sigma$ for the cell that contains $\sigma$, a mirror copy $G^{g_{\kappa'}}$ for the cell that contains $g_{\kappa'}B$.
Since $\mathcal{Q}''$ cuts $D(n_i)$, the translates of $Q''$ by the vectors of $\Gamma$ are open cells,
and exactly one of them is a cell of the brick:
$Q''-v=Q''_{\kappa''}$ with $\kappa''\in\mathcal{K}''$ for one and only one $v\in\Gamma$ (Definition~\ref{def:cell-families}).
The image of the factor under $\tau_{-v}$ belongs to $\mathcal{A}(Q''_{\kappa''})$ (Lemma~\ref{lem:admissible-images}\eqref{it:images-boxes}),
and in $\varphi_\alpha\cdot H$ the factor is exchanged for this image.
For a mirror copy $G^g$, $g=g_{\kappa'}$, the image is $\left(G^g\right)^{\tau_{-v}}=G^{\tau_{-v}g}$ by \eqref{eq:setting-group-action},
and $\tau_{-v}g$ lies again in $W(\mathcal{Q}')$, since $\Gamma$ lies in the translation subgroup of $W(\mathcal{Q}')$.
For $\kappa''\in\mathcal{K}''$ let $G_{\kappa''}$ be the product of the translated factors that belong to $\mathcal{A}(Q''_{\kappa''})$, the empty product being $1$.
The mirror copies come from open cells of $\mathcal{Q}''$ in pairwise different $\Gamma$-orbits,
and so they are moved into pairwise different cells of the brick.
Consequently each $G_{\kappa''}$ is of one of two kinds:
a character supported in $Q''_{\kappa''}$,
or the product $\varphi\cdot G^g$ of such a character $\varphi$ and of one element $G^g$ with $g\in W(\mathcal{Q}')$ and $gB\subset Q''_{\kappa''}$.
In either case $G_{\kappa''}\in\mathcal{A}(Q''_{\kappa''})$,
and the second kind occurs for exactly $n_i^3/(|Q_0|\,k^3)$ indices $\kappa''$, one for each mirror copy.

The product $\prod_{\kappa''\in\mathcal{K}''}G_{\kappa''}$ embeds in $D(n_i)$ by Lemma~\ref{lem:embedding-criteria}\eqref{it:embedding-cells},
since the cells $Q''_{\kappa''}$, $\kappa''\in\mathcal{K}''$, lie in pairwise different $\Gamma$-orbits
and the coordinates of the vectors of $\Gamma$ are multiples of $n_i/2\in2h''\Z$;
and $\varphi_\alpha\cdot H$ embeds in $D(n_i)$, its boxes being those of $H$.
Every factor of $\varphi_\alpha\cdot H$ embeds in $D(n_i)$ as well:
a character of one plaquette has no box, and the box $g_{\kappa'}B$ of a mirror copy has radius $r$ (Lemma~\ref{lem:coarse-structures}\eqref{it:coarse-tori}).
By Lemma~\ref{lem:admissible-images}\eqref{it:images-periods}, the image of such a factor under its translation $\tau_{-v}$, $v\in\Gamma$, defines the same function on $\Conf(\Gamma)$ as the factor itself;
for a character of one plaquette this is the $\Gamma$-periodicity of the field $\widetilde F$ at which it is evaluated.
The function on $\Conf(\Gamma)$ defined by a product that embeds in $D(n_i)$ is the product of the functions defined by its factors (Definition~\ref{def:integrands}).
Hence $\varphi_\alpha\cdot H$ and $\prod_{\kappa''\in\mathcal{K}''}G_{\kappa''}$ define the same function on $\Conf(\Gamma)$,
and Theorem~\ref{thm:chessboard}(c) gives
\begin{equation}\label{eq:amplification-cell-contents}
\left|Z_{D(n_i)}\left[\varphi_\alpha\cdot H\right]\right|\le\prod_{\kappa''\in\mathcal{K}''}\left\|G_{\kappa''}\right\|_{C(n_i/2,2n_i,n_i)},
\end{equation}
where, by Definition~\ref{def:folds},
\[
\left\|G_{\kappa''}\right\|_{C(n_i/2,2n_i,n_i)}^{\frac{n_i^3}{h''^3}}=Z_{C(n_i/2,2n_i,n_i)}\left[\Fold_{C(n_i/2,2n_i,n_i)}G_{\kappa''}\right].
\]
Any other set of representatives of the cells of $\mathcal{Q}''$ in $D(n_i)$ gives the same right side in \eqref{eq:amplification-cell-contents}.
Indeed, moving the factors into the cells of such a set instead replaces each $G_{\kappa''}$ by its image under a translation $\tau_v$ with $v\in\Gamma$.
The coordinates of $v$ lie in $2h''\Z$, so that $\tau_v\in W(\mathcal{Q}'')$ (Definition~\ref{def:cell-families});
by \eqref{eq:setting-group-action} and the simple transitivity of $W(\mathcal{Q}'')$ on the open cells,
$G_{\kappa''}$ and $G_{\kappa''}^{\tau_v}$ then have the same mirror copies with respect to $\mathcal{Q}''$, and
\[
\Fold_{C(n_i/2,2n_i,n_i)}\left(G_{\kappa''}^{\tau_v}\right)=\Fold_{C(n_i/2,2n_i,n_i)}G_{\kappa''}.
\]

If $G_{\kappa''}$ is a character, its fold on $C(n_i/2,2n_i,n_i)$ is a product of characters and has modulus one;
Proposition~\ref{prop:accretivity} and Corollary~\ref{cor:modulus-bound} give
\[
\left\|G_{\kappa''}\right\|_{C(n_i/2,2n_i,n_i)}^{\frac{n_i^3}{h''^3}}\le Z_{C(n_i/2,2n_i,n_i)}[1].
\]
Let now $G_{\kappa''}=\varphi\cdot G^g$ with a character $\varphi$ supported in $Q''_{\kappa''}$ and $g\in W(\mathcal{Q}')$.
By Definition~\ref{def:cell-structures} and Lemma~\ref{lem:admissible-images}\eqref{it:images-products},
\[
\Fold_{C(n_i/2,2n_i,n_i)}\left(\varphi\cdot G^g\right)=\left(\prod_{g''}\varphi^{g''}\right)\cdot\prod_{g''}\left(G^g\right)^{g''},
\]
where $g''$ ranges over a set of representatives in $W(\mathcal{Q}'')$ of $W(\mathcal{Q}'')/\Gamma'$.
Each $\varphi^{g''}$ is a character supported in an open cell of $\mathcal{Q}''$.
Such a cell lies in a box of radius $h''/2$,
and this box embeds in $\Gamma'$ by Lemma~\ref{lem:embedding-criteria}\eqref{it:embedding-sup}:
every non-zero vector of $\Gamma'$ has maximum norm at least $n_i/2\ge16h''$, a number that exceeds both $4$ and $2h''$.
By Lemma~\ref{lem:curvature-characters} the first product is therefore a character $\varphi_{\beta'}$ of $C(n_i/2,2n_i,n_i)$, $\beta'\in C^2(\Gamma')$.
The element $G^g$ is bounded by $C$ with the box $gB$ (Proposition~\ref{prop:bounded-elements}(a)).
As $g$ belongs to $W(\mathcal{Q}')$, Lemma~\ref{lem:coarse-structures}\eqref{it:coarse-placement} applies to it:
the representatives $g''$ form a placement of $gB$ in $C(n_i/2,2n_i,n_i)$ with $n_i^3/h''^3$ members.
Definition~\ref{def:bounded-elements} gives
\[
\left\|\varphi\cdot G^g\right\|_{C(n_i/2,2n_i,n_i)}^{\frac{n_i^3}{h''^3}}\le C^{\frac{n_i^3}{h''^3}}\cdot Z_{C(n_i/2,2n_i,n_i)}[1].
\]
Of the $n_i^3/h''^3$ contents, exactly $n_i^3/(|Q_0|\,k^3)$ are of the second kind.
Taking roots of order $n_i^3/h''^3$ in the last two bounds and inserting them into \eqref{eq:amplification-cell-contents}, one finds
\[
\left|Z_{D(n_i)}\left[\varphi_\alpha\cdot H\right]\right|\le C^{\frac{n_i^3}{|Q_0|\,k^3}}\cdot Z_{C(n_i/2,2n_i,n_i)}[1],
\]
uniformly in $\alpha$.
Summation over $\alpha$, with $\sum_\alpha|a_\alpha|=\exp\left(t\cdot|O|_1\cdot n_i^3/|Q_0|\right)$, yields
\[
\left|Z_{D(n_i)}\left[\exp\bigl(t\,O_{\mathfrak{d}}\bigr)\cdot H\right]\right|
\le\exp\left(t\cdot|O|_1\cdot\frac{n_i^3}{|Q_0|}\right)\cdot C^{\frac{n_i^3}{|Q_0|\,k^3}}\cdot Z_{C(n_i/2,2n_i,n_i)}[1].
\]

\item\label{it:amplification-ratio}
\emph{The ratio of the partition functions, and the root.}
The numerator in \eqref{eq:amplification-power-bound} is non-negative, hence equal to its modulus.
With the bounds of steps~\eqref{it:amplification-denominator} and~\eqref{it:amplification-numerator}, \eqref{eq:amplification-power-bound} becomes
\[
\left|\mu^t_i(G)\right|^{\frac{n_i^3}{|Q_0|\,k^3}}
\le\exp\left(2t\cdot|O|_1\cdot\frac{n_i^3}{|Q_0|}\right)\cdot C^{\frac{n_i^3}{|Q_0|\,k^3}}\cdot\frac{Z_{C(n_i/2,2n_i,n_i)}[1]}{Z_{D(n_i)}[1]}.
\]
The side $n_i$ is even and at least $8$.
Theorem~\ref{thm:four-tori} therefore gives $Z_{C(n_i)}[1]\le Z_{D(n_i)}[1]$
and $\log Z_{C(n_i/2,2n_i,n_i)}[1]-\log Z_{C(n_i)}[1]=n_i^3\cdot\Delta_i$,
the partition functions being positive (Theorem~\ref{thm:partition-positivity}).
Hence
\[
\frac{Z_{C(n_i/2,2n_i,n_i)}[1]}{Z_{D(n_i)}[1]}\le\frac{Z_{C(n_i/2,2n_i,n_i)}[1]}{Z_{C(n_i)}[1]}=\exp\left(n_i^3\cdot\Delta_i\right).
\]
Finally $\frac{n_i^3}{|Q_0|}=k^3\cdot\frac{n_i^3}{|Q_0|\,k^3}$ and $n_i^3=|Q_0|\cdot k^3\cdot\frac{n_i^3}{|Q_0|\,k^3}$.
Combining the last two inequalities and taking the root of order $n_i^3/(|Q_0|\,k^3)$, one obtains
\[
\left|\mu^t_i(G)\right|\le C\cdot\exp\left(2t\cdot|O|_1\cdot k^3+|Q_0|\cdot k^3\cdot\Delta_i\right)=C\cdot\gamma_{t,i}(k).
\]
\end{enumerate}
\end{proof}

\begin{remark}\label{rem:arrays-converge}
Let $\mathfrak{d}$ be data, $B$ a box of radius $r$, $G\in\mathcal{G}^{\mathrm{ev}}(B)$ and $i\ge i_0(\mathfrak{d},r)$,
and let $\mathcal{Q}'$, $\mathcal{Q}''$, $\Gamma$, $\Gamma'$ and $g_{\kappa'}$ be as in the proof of Theorem~\ref{thm:amplification}.
Consider three products of images of $G$, each on a torus:
on $D(n_i)$, the product $\prod_{\kappa'}G^{[\kappa']}$ of the mirror copies $G^{[\kappa']}=G^{g_{\kappa'}}$ with respect to $\mathcal{Q}'$,
$\kappa'$ ranging over representatives of the cells of $\mathcal{Q}'$ in $D(n_i)$;
on $C(n_i/2,2n_i,n_i)$ and for $g\in W(\mathcal{Q}')$, the product $\prod_{g''}\left(G^g\right)^{g''}$,
$g''$ ranging over representatives in $W(\mathcal{Q}'')$ of $W(\mathcal{Q}'')/\Gamma'$;
on a coordinate torus, the product $\prod_jG^{g_j}$ over a placement $(g_j)$ of $B$ in that torus.
By Lemma~\ref{lem:admissible-images}\eqref{it:images-boxes}, their factors lie in $\mathcal{G}^{\mathrm{ev}}(g_{\kappa'}B)$, $\mathcal{G}^{\mathrm{ev}}(g''gB)$ and $\mathcal{G}^{\mathrm{ev}}(g_jB)$, respectively.
For each of the three products, these boxes embed disjointly in the period lattice of its torus.
For the third product this is property~\eqref{it:placement-disjoint} of Definition~\ref{def:bounded-elements}.
The second is a product of the same kind, for the element $G^g$ of $\mathcal{G}^{\mathrm{ev}}(gB)$,
because the representatives $g''$ form a placement of $gB$ in $C(n_i/2,2n_i,n_i)$ (Lemma~\ref{lem:coarse-structures}\eqref{it:coarse-placement}).
In the first, the box $g_{\kappa'}B$ lies in the open cell $g_{\kappa'}Q'_0=Q'_{\kappa'}$,
these cells belong to pairwise different $\Gamma$-orbits,
and $\ell_-(\Gamma)$, $\ell_+(\Gamma)$ and $x_3(\Gamma)$ are contained in $2kh\Z$;
as in step~\eqref{it:amplification-numerator} of that proof, Lemma~\ref{lem:embedding-criteria}\eqref{it:embedding-cells} shows that this product embeds in $\Gamma$,
which by Definition~\ref{def:integrands} means that its boxes embed disjointly in $\Gamma$.
In particular the boxes of two different factors are disjoint (take the zero vector twice in the definition of a disjoint embedding),
so that each of the three products belongs to $\mathcal{A}(\R^3)$ and embeds in its torus.
As the translates of these boxes by the vectors of the lattice are pairwise disjoint,
a plaquette of the torus has at most one representative contained in one of the boxes, and this representative lies in a single box;
hence the curvature at a plaquette of the torus enters the growth bound of Definition~\ref{def:integrands} for one factor at most.
A bounded Borel function $\Phi$ of finitely many $F_\sigma$, for instance $\exp\bigl(t\,O_{\mathfrak{d}}\bigr)$ on $D(n_i)$ or a character,
may depend on plaquettes inside these boxes:
Definition~\ref{def:integrands} places no restriction on them,
and the product of $\Phi$ with any of the three products is again an element of $\mathcal{A}(\R^3)$ that embeds in the torus, its boxes being the same.
For all parameters $-1\le\rho\le1$, $w_0,\dots,w_3>0$, $M\in\R$ and $q>0$ of the functionals,
Proposition~\ref{prop:berezin}\eqref{it:torus-integrability} therefore shows that the integral defining $Z_{\mathbb{T}}$ of each of the three products on its torus $\mathbb{T}$,
with or without such a factor $\Phi$, converges absolutely.
\end{remark}

\begin{remark}\label{rem:amplification}
\begin{enumerate}
\item[(1)]
\emph{Amplification in place of Jensen's inequality.}
Let $\mathbb{E}$ be the expectation of a probability measure.
For real functions $f_v$ Jensen's inequality gives
\[
\mathbb{E}\exp\left(\sum_v f_v\right)\ge\exp\left(\sum_v\mathbb{E}\,f_v\right),
\]
and in order to apply it to an integrand $G$ one has to write $G=\exp(f)$ with a real function $f$;
this requires $G>0$.
In the proof of Theorem~\ref{thm:amplification} the inequality \eqref{eq:amplification-power-bound} takes the place of Jensen's inequality.
Written with the functional $\mu^t_i$ on both sides, it is
\[
\left|\mu^t_i(G)\right|^{\frac{n_i^3}{|Q_0|\,k^3}}\le\mu^t_i\left(\prod_{\kappa'}G^{[\kappa']}\right),
\]
the product being taken over the cells $\kappa'$ of $\mathcal{Q}'$ in $D(n_i)$,
and it holds for every element $G$ of the space $\mathcal{A}(Q'_0)$ (Definition~\ref{def:integrands}),
the linear span of the admissible products in the open cell $Q'_0$ of $\mathcal{Q}'$ that contains the box $B$:
$G$ may be complex valued, unbounded, or take values in the Grassmann algebra.
Its derivation involves no logarithm:
the inequality comes from reflection positivity at the mirrors of the three families of the diagonal frame,
through the chessboard estimate of Theorem~\ref{thm:chessboard}(a).
\item[(2)]
\emph{The ratio of two partition functions.}
Of the ratio $Z_{C(n_i/2,2n_i,n_i)}[1]/Z_{D(n_i)}[1]$ it is only known that it is of the form $\exp(o(n_i^3))$:
by \eqref{eq:tori-chain} it lies between $1$ and $\exp(n_i^3\cdot\Delta_i)$, and $\Delta_i\to0$ (Theorem~\ref{thm:four-tori}).
In step~\ref{it:amplification-ratio} of the proof of Theorem~\ref{thm:amplification} the ratio is raised to the power $\frac{|Q_0|\,k^3}{n_i^3}$.
This is the inverse of the number of the cells of $\mathcal{Q}'$ in $D(n_i)$,
and that number is the fixed fraction $\frac{1}{|Q_0|\,k^3}$ of the volume $n_i^3$ of $D(n_i)$.
Hence the ratio contributes at most the factor $\exp(|Q_0|\cdot k^3\cdot\Delta_i)$, which tends to $1$.
\item[(3)]
\emph{From the diagonal to the coordinate frame.}
The reflections in the walls of $\mathcal{Q}'$ that belong to the families $\ell_-$ and $\ell_+$ exchange the link planes of the direction $e_1$ with those of the direction $e_2$.
If $G=G_X$ is an insertion at a link plane of the direction $e_1$,
some of its mirror copies $G^{[\kappa']}$ are therefore insertions at link planes of the direction $e_2$,
and in the picture with time $e_1$ such an insertion is not given by an operator.
The chessboard estimate between the two frames, Theorem~\ref{thm:chessboard}(c), separates the copies.
After it, the fold of each cell content is a character times the images of a single element under elements of $W_{\mathrm{c}}$;
these elements map link planes of one direction to link planes of the same direction,
so that all insertions in one fold sit at link planes of one direction.
\end{enumerate}
\end{remark}

In the corresponding argument for Yang--Mills theory in four dimensions,
the counterpart of Proposition~\ref{prop:holder-insertions} is applied to the operators $\exp(sD)$ with $D$ real and $D^*=-D$,
and the transition to a state is made by Jensen's inequality for the logarithm of the kernel of $\exp(sD)$ placed between two square roots of the crossing operator
(compare Remark~\ref{rem:holder-jensen}).
This kernel therefore has to be strictly positive.
At a coordinate link plane, the Gaussian factor of the crossing operator over a finite set $S$ of links is $\exp\left(-\frac12|E_S|^2\right)$,
with the square root $\exp\left(-\frac14|E_S|^2\right)$ (Definition~\ref{def:slice-operators}).
For this factor, and for a suitable real operator $D$ with $D^*=-D$, the kernel in question changes sign.

\begin{proposition}[Sign change of a regularised kernel]\label{prop:sign-change}
Let $S$ be a finite set of links of the slice lattice,
let $E_b:=-i\,\partial/\partial A(b)$, $b\in S$, act on $L^2(\R^S)$, $E_S:=(E_b)_{b\in S}$,
let $x\in\R^S$, and let
\[
D:=\exp\left(-i\langle x,E_S\rangle\right)-\exp\left(i\langle x,E_S\rangle\right).
\]
For every real $s$ the operator $\exp\left(-\frac14|E_S|^2\right)\exp(sD)\exp\left(-\frac14|E_S|^2\right)$ on $L^2(\R^S)$ has, in the sense of Definition~\ref{def:integral-kernels}, the integral kernel
\[
(A',A)\mapsto R_s\left(\langle x,A'-A\rangle\right)\cdot\prod_{b\in S}g_1(A'(b)-A(b)),
\qquad
R_s(\xi):=\sum_{j\in\Z}J_j(2s)\cdot\exp\left(j\xi-\frac{j^2|x|^2}{2}\right),
\]
where $g_1$ is the centred Gaussian density of variance one (Definition~\ref{def:slices}),
$J_j$ is the Bessel function of the first kind of order $j$,
and the series converges absolutely for every real $\xi$.
If $|x|^2\ge4$ and $0<s\le\frac12$, then
\[
R_s(\xi_0)\le-\frac32
\qquad\text{at}\qquad
\xi_0:=-\log\frac4s-\frac{|x|^2}{2}.
\]
\end{proposition}

\begin{proof}
Write $U:=\exp\left(-i\langle x,E_S\rangle\right)$.
It is the translation $(U\Phi)(A)=\Phi(A-x)$ (Definition~\ref{def:slice-operators}),
a unitary operator that commutes with complex conjugation, with $U^{-1}=U^*=\exp\left(i\langle x,E_S\rangle\right)$;
its power $U^j$, $j\in\Z$, is the translation by $jx$.
So $D=U-U^{-1}$ is real and $D^*=-D$.

\emph{Bessel functions.}
For real $s$ the generating function of the Bessel functions gives
\[
\exp\left(s\left(z-z^{-1}\right)\right)=\sum_{j\in\Z}J_j(2s)\,z^j,
\qquad z\in\C\setminus\{0\},
\]
a Laurent series of a function holomorphic on $\C\setminus\{0\}$, hence absolutely convergent at every $z\neq0$.
Moreover $J_{-j}=(-1)^jJ_j$, and for $j\ge0$
\[
J_j(2s)=\sum_{\upsilon\ge0}\frac{(-1)^\upsilon\,s^{j+2\upsilon}}{\upsilon!\,(j+\upsilon)!}.
\]

\emph{The kernel.}
On the unit circle $\overline{z}=z^{-1}$, and the coefficients $J_j(2s)$ are absolutely summable.
The continuous functional calculus of the unitary $U$, applied to the function $z\mapsto\exp\left(s\left(z-\overline{z}\right)\right)$, gives
\[
\exp(sD)=\exp\left(s\left(U-U^*\right)\right)=\sum_{j\in\Z}J_j(2s)\cdot U^j,
\]
with convergence in operator norm.
Each $U^j$ is a function of $E_S$ and commutes with $\exp\left(-\frac14|E_S|^2\right)$.
Hence the operator $T:=\exp\left(-\frac14|E_S|^2\right)\exp(sD)\exp\left(-\frac14|E_S|^2\right)$ is the sum of the series
\[
T=\sum_{j\in\Z}J_j(2s)\cdot\exp\left(-\frac12|E_S|^2\right)U^j,
\]
which converges in operator norm.
The operator $\exp\left(-\frac12|E_S|^2\right)$ has the kernel $\prod_{b\in S}g_1(A'(b)-A(b))$ (Definition~\ref{def:integral-kernels}),
and $U^j$ shifts the variable by $jx$;
hence $\exp\left(-\frac12|E_S|^2\right)U^j$ has the kernel
\[
k_j(A',A):=\prod_{b\in S}g_1(A'(b)-A(b)-jx_b).
\]
These kernels are continuous and bounded by $(2\pi)^{-\frac{\#|S|}{2}}$, and $\sum_{j\in\Z}|J_j(2s)|<\infty$;
so the series $k:=\sum_{j\in\Z}J_j(2s)\cdot k_j$ converges absolutely and uniformly on $\R^S\times\R^S$, and $k$ is continuous.
Moreover $g_1^2=(4\pi)^{-\frac12}\cdot g_{\frac12}$, and $g_{\frac12}$ is a probability density on $\R$;
consequently
\[
\int_{\R^S}\mathrm{d}A\,k_j(A',A)^2=(4\pi)^{-\frac{\#|S|}{2}}
\qquad\text{for all }A'\in\R^S\text{ and }j\in\Z.
\]
Now let $\Phi\in L^2(\R^S)$, with the norm $\|\Phi\|$.
For every $A'\in\R^S$ the Cauchy--Schwarz inequality gives
\[
\sum_{j\in\Z}|J_j(2s)|\cdot\int_{\R^S}\mathrm{d}A\,k_j(A',A)\cdot|\Phi(A)|
\le(4\pi)^{-\frac{\#|S|}{4}}\cdot\|\Phi\|\cdot\sum_{j\in\Z}|J_j(2s)|<\infty.
\]
By Fubini's theorem for the sum over $j$ and the integral over $A$,
the integral $\int_{\R^S}\mathrm{d}A\,k(A',A)\,\Phi(A)$ therefore converges absolutely for every $A'$, and it is the sum of the series
\[
\sum_{j\in\Z}J_j(2s)\cdot\int_{\R^S}\mathrm{d}A\,k_j(A',A)\,\Phi(A),
\]
whose $j$-th term has modulus at most $(4\pi)^{-\frac{\#|S|}{4}}\cdot\|\Phi\|\cdot|J_j(2s)|$, whatever $A'$ is.
So the partial sums of this series, as functions of $A'$, converge uniformly on $\R^S$ to $A'\mapsto\int_{\R^S}\mathrm{d}A\,k(A',A)\,\Phi(A)$.
On the other hand, the $j$-th term, as a function of $A'$, represents the vector $J_j(2s)\cdot\exp\left(-\frac12|E_S|^2\right)U^j\Phi$ of $L^2(\R^S)$,
because $k_j$ is the kernel of $\exp\left(-\frac12|E_S|^2\right)U^j$;
by the convergence in operator norm the same partial sums converge to $T\Phi$ in $L^2(\R^S)$,
and a subsequence of them converges to $T\Phi$ almost everywhere.
Comparing the two limits, one obtains
\[
(T\Phi)(A')=\int_{\R^S}\mathrm{d}A\,k(A',A)\,\Phi(A)\qquad\text{for almost every }A'\in\R^S.
\]
As $\Phi\in L^2(\R^S)$ was arbitrary, $T$ has the kernel $k$ in the sense of Definition~\ref{def:integral-kernels}.
For every $b\in S$ one has
\[
g_1(A'(b)-A(b)-jx_b)=g_1(A'(b)-A(b))\cdot\exp\left(jx_b\left(A'(b)-A(b)\right)-\frac{j^2x_b^2}{2}\right);
\]
therefore
\[
\prod_{b\in S}g_1(A'(b)-A(b)-jx_b)
=\prod_{b\in S}g_1(A'(b)-A(b))\cdot\exp\left(j\langle x,A'-A\rangle-\frac{j^2|x|^2}{2}\right),
\]
and summation over $j$ with the coefficients $J_j(2s)$ shows that $k$ is the kernel of the statement.
The series $R_s(\xi)$ converges absolutely,
because $\left|\exp\left(j\xi-\frac12j^2|x|^2\right)\right|\le\exp(j\xi)+\exp(-j\xi)$
and the Laurent series above converges absolutely at $z=\exp(\xi)$ and at $z=\exp(-\xi)$.

\emph{Bounds for the coefficients.}
Let $0<s\le\frac12$ and $j\ge0$.
In the series for $J_j(2s)$ the ratio of the moduli of the terms of index $\upsilon+1$ and $\upsilon$ is
\[
\frac{s^2}{(\upsilon+1)(j+\upsilon+1)}\le s^2<1:
\]
the series alternates and its terms decrease in modulus.
So its sum lies between its first term and the sum of its first two terms, and
\[
0<J_j(2s)\le\frac{s^j}{j!}\quad(j\ge0),
\qquad
J_1(2s)\ge s-\frac{s^3}{2}.
\]
In particular $J_0(2s)\le1$ and $|J_j(2s)|\le s^{|j|}/|j|!$ for all $j\in\Z$.

\emph{The value at $\xi_0$.}
Let $|x|^2\ge4$, $0<s\le\frac12$, and put $\lambda:=\exp\left(-\frac12|x|^2\right)\le\exp(-2)$.
Then $\exp(-\xi_0)=4/(s\lambda)$, $\exp(\xi_0)=s\lambda/4$, and
\[
R_s(\xi_0)=\sum_{j\in\Z}J_j(2s)\cdot\exp(j\xi_0)\cdot\lambda^{j^2}.
\]
The term $j=0$ is $J_0(2s)\le1$.
The term $j=-1$ is
\[
-J_1(2s)\cdot\exp(-\xi_0)\cdot\lambda=-\frac{4J_1(2s)}{s}\le-4\left(1-\frac{s^2}{2}\right)\le-\frac72 .
\]
For $j\le-2$ the modulus of the term is at most
\[
\frac{s^{|j|}}{|j|!}\left(\frac{4}{s\lambda}\right)^{|j|}\lambda^{j^2}
=\frac{4^{|j|}}{|j|!}\,\lambda^{j^2-|j|}
\le\frac{4^{|j|}}{|j|!}\,\lambda^2,
\]
because $\lambda\le1$ and $j^2-|j|\ge2$;
the sum of these bounds over $j\le-2$ is $\left(\exp(4)-5\right)\lambda^2\le\left(\exp(4)-5\right)\exp(-4)<0.91$.
For $j\ge1$ the modulus of the term is at most
\[
\frac{s^j}{j!}\left(\frac{s\lambda}{4}\right)^j
=\frac{1}{j!}\left(\frac{s^2\lambda}{4}\right)^j,
\]
since $\lambda^{j^2}\le1$;
the sum of these bounds over $j\ge1$ is
\[
\exp\left(\frac{s^2\lambda}{4}\right)-1\le\exp\left(\frac{\exp(-2)}{16}\right)-1<0.01.
\]
Adding the four contributions,
\[
R_s(\xi_0)\le1-\frac72+0.91+0.01=-1.58<-\frac32.
\]
\end{proof}

Let $|x|^2\ge4$ and $0<s\le\frac12$.
On every interval $[-\xi',\xi']$, $\xi'>0$, the series that defines $R_s$ converges uniformly:
for $|\xi|\le\xi'$ the $j$-th term of $R_s(\xi)$ has modulus at most $|J_j(2s)|\cdot\left(\exp(j\xi')+\exp(-j\xi')\right)$,
and these bounds are summable over $j$, because the Laurent series of the generating function converges absolutely at $z=\exp(\xi')$ and at $z=\exp(-\xi')$.
So $R_s$ is continuous on $\R$, and since $R_s(\xi_0)\le-\frac32$, it is negative on an open interval around $\xi_0$.
At the point $0$ the terms of the orders $j$ and $-j$ cancel for odd $j$, because $J_{-j}=(-1)^jJ_j$, and what remains is
\[
R_s(0)=J_0(2s)+2\sum_{j\ge2\ \text{even}}J_j(2s)\cdot\exp\left(-\frac{j^2|x|^2}{2}\right)>0,
\]
since $J_j(2s)>0$ for $j\ge0$, the power series of $J_j(2s)$ being alternating with a positive first term and terms of decreasing modulus.
As $x\neq0$, it follows that the kernel $(A',A)\mapsto R_s\left(\langle x,A'-A\rangle\right)\cdot\prod_{b\in S}g_1(A'(b)-A(b))$ is positive on an open neighbourhood of the diagonal $\{A'=A\}$
and negative on the non-empty open set of the pairs $(A',A)$ for which $\langle x,A'-A\rangle$ lies in that interval around $\xi_0$.
Two kernels of the same operator agree almost everywhere on $\R^S\times\R^S$,
as one sees by testing them with the indicator functions of the boxes with rational corners.
Therefore every kernel of $\exp\left(-\frac14|E_S|^2\right)\exp(sD)\exp\left(-\frac14|E_S|^2\right)$ is negative on a set of positive measure,
and no kernel of this operator has a real logarithm.

When the kernel is inserted at a link plane $c=j+\frac12$, the difference $A'-A$ is read as the electric field $F^{\mathrm{el}}_j$ on $S$ (Definition~\ref{def:insertions}),
and the kernel is negative at the configurations for which $\langle x,F^{\mathrm{el}}_j\rangle$ lies in the interval around $\xi_0$.
The inequality \eqref{eq:amplification-power-bound}, which takes the place of Jensen's inequality in the proof of Theorem~\ref{thm:amplification},
holds for integrands of arbitrary sign (Remark~\ref{rem:amplification}(1));
accordingly, Theorem~\ref{thm:two-layer-principle}(a) holds for every $X$ in the mid-plane algebra $\mathfrak{W}$.

\begin{proposition}[Domination and tightness]\label{prop:domination-tightness}
Let $-1\le\rho\le1$, let the weights $w_0,\dots,w_3$ be positive, $M\in\R$ and $q>0$.
Let the sequence of sides $\mathfrak{n}=(n_i)$ satisfy \eqref{cond:divisibility}, and let $\mathfrak{d}$ be data (Definition~\ref{def:data}).
Let $\beta\in C^2_{\mathrm{c}}$ be supported in a box $B$ of radius $r$, that is, every plaquette of $\supp\beta$ is contained in $B$, and let $k=k(\mathfrak{d},r)$.
Then for all $t\ge0$ and all $i\ge i_0(\mathfrak{d},r)$
\[
\mu^t_i\left(\exp\langle\beta,F\rangle\right)\le\exp\left(\frac{\|\beta\|^2}{2}\right)\cdot\gamma_{t,i}(k),
\]
with the amplification factor $\gamma_{t,i}(k)$ of Definition~\ref{def:amplification-factor}.
Consequently, for every $t\ge0$ the family $(\mu^t_i)_i$ of laws on $\Omega$ is tight;
every weak limit point $\nu_t$ of $(\mu^t_i)_i$ satisfies
\[
\nu_t\left(\exp\langle\beta,F\rangle\right)\le\exp\left(\frac{\|\beta\|^2}{2}+2t\cdot|O|_1\cdot k^3\right)
\]
for all $\beta$, $r$ as above and $k=k(\mathfrak{d},r)$;
and every comparison state for $\mathfrak{d}$ (Definition~\ref{def:comparison-states}) satisfies \eqref{cond:gaussian-domination}.
\end{proposition}

\begin{proof}
\emph{The bound.}
Fix $t\ge0$ and $i\ge i_0(\mathfrak{d},r)$, and let $G:=\exp\langle\beta,F\rangle$.
Since $|\beta(\sigma)F_\sigma|\le\beta(\sigma)^2+\frac14F_\sigma^2$ for every plaquette $\sigma$,
\[
0<G\le\exp\left(\|\beta\|^2+\frac14\sum_{\sigma\in\supp\beta}F_\sigma^2\right);
\]
so $G$ is a positive bosonic element of $\mathcal{G}^{\mathrm{ev}}(B)$, with the growth margin $\frac14$ (Definition~\ref{def:integrands}).
Let $\mathcal{Q}'$ and $\mathcal{Q}''$ be the cell structures of Lemma~\ref{lem:coarse-structures} for the box $B$,
let $Q'_0$ be the open cell of $\mathcal{Q}'$ that contains $B$,
and for a cell $\kappa'$ of $\mathcal{Q}'$ let $g_{\kappa'}\in W(\mathcal{Q}')$ be the element with $g_{\kappa'}Q'_0=Q'_{\kappa'}$.

Steps~\ref{it:amplification-coarse-array}--\ref{it:amplification-denominator} of the proof of Theorem~\ref{thm:amplification} apply to every element of $\mathcal{G}^{\mathrm{ev}}(B)$.
For $G$ they give the inequality \eqref{eq:amplification-power-bound}, in which $\mu^t_i(G)>0$
($\mu^t_i$ is the expectation of a probability measure on bosonic integrands, Definition~\ref{def:data}),
together with
\[
Z_{D(n_i)}\left[\exp(tO_{\mathfrak{d}})\right]\ge\exp\left(-\frac{t\cdot|O|_1\cdot n_i^3}{|Q_0|}\right)\cdot Z_{D(n_i)}[1]>0
\]
(Lemma~\ref{lem:perturbed-denominators} and Theorem~\ref{thm:partition-positivity}).
Hence
\[
\mu^t_i(G)^{\frac{n_i^3}{|Q_0|\,k^3}}
\le\exp\left(\frac{t\cdot|O|_1\cdot n_i^3}{|Q_0|}\right)\cdot
\frac{Z_{D(n_i)}\left[\exp(tO_{\mathfrak{d}})\prod_{\kappa'}G^{[\kappa']}\right]}{Z_{D(n_i)}[1]},
\]
the product being taken over the $n_i^3/(|Q_0|\,k^3)$ cells $\kappa'$ of $\mathcal{Q}'$ in $D(n_i)$.

The function $G$ is real and bosonic;
so, by \eqref{eq:setting-group-action}, its mirror copies are $G^{[\kappa']}(F)=G(g_{\kappa'}^*F)$, without complex conjugation also for odd $g_{\kappa'}$.
The pull-back $g^*$ by an element $g$ of $W$ is orthogonal for the pairing, and $(g^*)^{-1}=(g^{-1})^*$; hence
\[
G^{[\kappa']}=\exp\langle\beta^{\kappa'},F\rangle,
\qquad
\beta^{\kappa'}:=\left(g_{\kappa'}^{-1}\right)^*\beta,
\qquad
\|\beta^{\kappa'}\|=\|\beta\|,
\]
and every plaquette of $\supp\beta^{\kappa'}$ lies in $g_{\kappa'}B$.
In particular $\prod_{\kappa'}G^{[\kappa']}>0$.
On $D(n_i)$ one has $|O_{\mathfrak{d}}|\le|O|_1\cdot n_i^3/|Q_0|$ (Definition~\ref{def:arrays}),
so that \eqref{eq:amplification-monotone} gives
\[
Z_{D(n_i)}\left[\exp(tO_{\mathfrak{d}})\prod_{\kappa'}G^{[\kappa']}\right]
\le\exp\left(\frac{t\cdot|O|_1\cdot n_i^3}{|Q_0|}\right)\cdot Z_{D(n_i)}\left[\prod_{\kappa'}G^{[\kappa']}\right].
\]

By Lemma~\ref{lem:coarse-structures}\eqref{it:coarse-coordinate}, in the form that \eqref{it:coarse-tori} of the same lemma gives on $D(n_i)$,
the boxes $g_{\kappa'}B$ lie in open cells of $\mathcal{Q}''$ that are pairwise different modulo the period lattice of $D(n_i)$.
Therefore $\prod_{\kappa'}G^{[\kappa']}=\prod_{\kappa''}G_{\kappa''}$, the product over the $n_i^3/h''^3$ cells $\kappa''$ of $\mathcal{Q}''$ in $D(n_i)$,
where the content $G_{\kappa''}$ is the copy $G^{[\kappa']}$ whose box lies in the cell $\kappa''$ if there is one, and $G_{\kappa''}=1$ otherwise;
exactly $n_i^3/(|Q_0|\,k^3)$ contents are of the first kind.
As for \eqref{eq:amplification-cell-contents}, each cell $\kappa''$ of $D(n_i)$ is represented by a cell $Q''$ of $\mathcal{Q}''$ in $\R^3$,
and a content $\exp\langle\beta^{\kappa'},F\rangle$ that lies in $Q''+v$, $v$ in the period lattice of $D(n_i)$, is replaced by its translate by $-v$;
the translate defines the same function on the configurations of $D(n_i)$, because the field $\widetilde F$ is periodic.
It is of the form $\exp\langle\beta',F\rangle$ with $\|\beta'\|=\|\beta\|$ and $\beta'$ supported in $Q''$.
Theorem~\ref{thm:chessboard}(c), applied to $D(n_i)=D(n_i/2,n_i,n_i)$ and $C(n_i/2,2n_i,n_i)$ as in step~\ref{it:amplification-numerator}, gives
\[
Z_{D(n_i)}\left[\prod_{\kappa'}G^{[\kappa']}\right]\le\prod_{\kappa''}\|G_{\kappa''}\|_{C(n_i/2,2n_i,n_i)},
\qquad
\|G_{\kappa''}\|_{C(n_i/2,2n_i,n_i)}^{\frac{n_i^3}{h''^3}}=Z_{C(n_i/2,2n_i,n_i)}\left[\Fold_{C(n_i/2,2n_i,n_i)}G_{\kappa''}\right].
\]
For a content $1$ the right side of the identity is $Z_{C(n_i/2,2n_i,n_i)}[1]$.
For a content $\exp\langle\beta',F\rangle$ the fold is $\exp\langle\beta'',F\rangle$,
where $\beta''$ is the sum of the $n_i^3/h''^3$ cochains $(g^{-1})^*\beta'$,
$g$ ranging over a set of representatives in $W(\mathcal{Q}'')$ of the quotient by the period lattice of $C(n_i/2,2n_i,n_i)$,
read as a cochain on this torus.
These images lie in pairwise different cells of the torus;
so they have pairwise disjoint supports, and $\|\beta''\|^2=\frac{n_i^3}{h''^3}\cdot\|\beta\|^2$.
The sides of $C(n_i/2,2n_i,n_i)$ are at least $4$ (step~\ref{it:amplification-numerator}),
and Theorem~\ref{thm:coordinate-domination}, applied to the cochain $\beta''$, gives
\[
\|\exp\langle\beta',F\rangle\|_{C(n_i/2,2n_i,n_i)}^{\frac{n_i^3}{h''^3}}
\le\exp\left(\frac{n_i^3\cdot\|\beta\|^2}{2h''^3}\right)\cdot Z_{C(n_i/2,2n_i,n_i)}[1].
\]
Since exactly $n_i^3/(|Q_0|\,k^3)$ of the $n_i^3/h''^3$ contents are of this kind, it follows that
\[
Z_{D(n_i)}\left[\prod_{\kappa'}G^{[\kappa']}\right]
\le\exp\left(\frac{n_i^3\cdot\|\beta\|^2}{2|Q_0|\,k^3}\right)\cdot Z_{C(n_i/2,2n_i,n_i)}[1].
\]
By \eqref{eq:tori-chain}, $Z_{C(n_i/2,2n_i,n_i)}[1]\le\exp(n_i^3\cdot\Delta_i)\cdot Z_{D(n_i)}[1]$.
Collecting the inequalities,
\[
\mu^t_i(G)^{\frac{n_i^3}{|Q_0|\,k^3}}
\le\exp\left(\frac{2t\cdot|O|_1\cdot n_i^3}{|Q_0|}+\frac{n_i^3\cdot\|\beta\|^2}{2|Q_0|\,k^3}+n_i^3\cdot\Delta_i\right),
\]
and the root of order $n_i^3/(|Q_0|\,k^3)$ gives
\[
\mu^t_i(G)\le\exp\left(2t\cdot|O|_1\cdot k^3+\frac{\|\beta\|^2}{2}+|Q_0|\cdot k^3\cdot\Delta_i\right)=\exp\left(\frac{\|\beta\|^2}{2}\right)\cdot\gamma_{t,i}(k).
\]

\emph{Tightness.}
The space $\Omega$ is a closed subset of the product of the real lines indexed by the plaquettes, with the product topology.
Let $\sigma$ be a plaquette, $s>0$, and $1_\sigma$ the two-cochain with the value $1$ at $\sigma$ and $0$ elsewhere.
The cochains $\pm s\cdot1_\sigma$ are supported in a box of radius $1$ and have the norm $s$.
By the bound just proved, with $r=1$, and by Markov's inequality,
\[
\mu^t_i\left(|F_\sigma|\ge R\right)\le\exp(-sR)\cdot\left(\mu^t_i\left(\exp(sF_\sigma)\right)+\mu^t_i\left(\exp(-sF_\sigma)\right)\right)
\le2\exp\left(-sR+\frac{s^2}{2}\right)\cdot\gamma_{t,i}(k(\mathfrak{d},1))
\]
for all $R>0$ and $i\ge i_0(\mathfrak{d},1)$.
By Lemma~\ref{lem:coarse-structures}, $k(\mathfrak{d},1)$ and $i_0(\mathfrak{d},1)$ do not depend on the position of the box, hence not on $\sigma$;
and $\gamma_{t,i}(k(\mathfrak{d},1))$ is bounded in $i$, because $\Delta_i\to0$ (Theorem~\ref{thm:four-tori}).
Let $\varepsilon>0$ and let $j$ be a bijection from the set of the plaquettes onto the positive integers.
There are numbers $R_\sigma>0$ with $\mu^t_i(|F_\sigma|>R_\sigma)\le\varepsilon\cdot2^{-j(\sigma)}$ for every plaquette $\sigma$ and every $i\ge i_0(\mathfrak{d},1)$.
The set $\Omega\cap\prod_\sigma[-R_\sigma,R_\sigma]$ is compact, as a closed subset of a product of compact intervals (Tychonoff's theorem),
and for $i\ge i_0(\mathfrak{d},1)$ its complement has $\mu^t_i$-measure at most $\sum_\sigma\varepsilon\cdot2^{-j(\sigma)}=\varepsilon$.
Each of the finitely many laws $\mu^t_i$ with $i<i_0(\mathfrak{d},1)$ is tight by itself, as every Borel probability measure on the Polish space $\Omega$.
So the family $(\mu^t_i)_i$ is tight.

\emph{Limits.}
The function $\exp\langle\beta,F\rangle$ is continuous and non-negative on $\Omega$.
If a sequence of laws on $\Omega$ converges weakly to a law, then for every $R>0$ the integral of the bounded continuous function $\min\left(\exp\langle\beta,F\rangle,R\right)$ under the limit is the limit of its integrals,
hence at most the lower limit of the integrals of $\exp\langle\beta,F\rangle$;
letting $R\to\infty$ (monotone convergence), the integral of $\exp\langle\beta,F\rangle$ under the limit law is at most the lower limit of its integrals.
Let $\nu_t$ be the weak limit of a subsequence $(\mu^t_{i_m})_m$.
As $\gamma_{t,i}(k)\to\exp(2t\cdot|O|_1\cdot k^3)$ when $i\to\infty$, since $\Delta_i\to0$, the bound gives
\[
\nu_t\left(\exp\langle\beta,F\rangle\right)\le\liminf_{m\to\infty}\mu^t_{i_m}\left(\exp\langle\beta,F\rangle\right)\le\exp\left(\frac{\|\beta\|^2}{2}+2t\cdot|O|_1\cdot k^3\right).
\]
A comparison state $\nu_0$ is the weak limit of laws $\nu_t$ of this kind, $t$ ranging over a sequence of strengths that tends to $0$ (Definition~\ref{def:comparison-states}).
By the same lower semicontinuity,
\[
\nu_0\left(\exp\langle\beta,F\rangle\right)\le\liminf_{t\to0}\exp\left(\frac{\|\beta\|^2}{2}+2t\cdot|O|_1\cdot k^3\right)=\exp\left(\frac{\|\beta\|^2}{2}\right).
\]
Every $\beta\in C^2_{\mathrm{c}}$ is supported in some box, since its support is finite;
so $\nu_0$ satisfies \eqref{cond:gaussian-domination}.
\end{proof}

\subsection{Two-layer elements}\label{sec:two-layer-elements}

Throughout this subsection the time direction is $e_1$,
$c=j_0+\frac12$ with $j_0\in\Z$,
and
\[
Y=\exp\left(i\langle\delta,\hat A\rangle\right)\,\Psi
\]
is a monomial of $\mathcal{P}_1(c)$ with $Y\neq1$ (Definition~\ref{def:two-layer}).
Its exponent $\delta$ is the sum of its restrictions $\delta_0$ to the links of the plane $\{x_1=j_0\}$,
$\delta_1$ to the links of the plane $\{x_1=j_0+1\}$,
and $\delta_\perp$ to the links between these two planes;
we write $\delta=(\delta_0,\delta_1,\delta_\perp)$.
Here $\delta_0$ and $\delta_1$ are read as one-cochains of the slice lattice $\Z^2$,
and $\delta_\perp$ as the function on the sites of $\Z^2$ whose value at $y$ is the value of $\delta$ on the link from $(j_0,y)$ to $(j_0+1,y)$.

As in Definition~\ref{def:two-layer},
$\widetilde E=E_S+\delta_0$,
$S$ is the set of the links $b$ of $\Z^2$ with $\delta_0(b)\neq0$ or $\delta_1(b)\neq0$,
so that $\delta_0$ and $\delta_1$ are also vectors of $\R^S$,
and $\Lambda_Y$ is a set of sites in play of $Y$:
a finite set of sites of $\Z^2$ that contains the end points of the links of $S$,
the support of $\delta_\perp$,
and every $y$ for which $\Psi$ contains a field at a site $x$ with $\pi(x)=(j_0,y)$ or $\pi(x)=(j_0+1,y)$.

The fields of $\Psi$ lie in the thin layers $\mathbb{X}_{j_0,6}$ and $\mathbb{X}_{j_0+1,1}$,
the lower and the upper ends of the bonds that cross the plane $\{x_1=c\}$,
and $\pi$ maps these two layers into the planes $\{x_1=j_0\}$ and $\{x_1=j_0+1\}$.
Hence the charge density $\mathfrak{q}(\Psi)$ of Definition~\ref{def:chamber-operator} vanishes outside the points $(j_0,y)$ and $(j_0+1,y)$, $y\in\Lambda_Y$.
The set $\Lambda_Y$ is not empty, because $Y\neq1$:
either $\Psi$ contains a field,
or $\delta\neq0$, and then $S$ or the support of $\delta_\perp$ is not empty.

Operators on $\mathcal{H}^{\mathrm{l}}$ refer to a rectangular cross-section $\Xi$ in which $\hull(S,\Lambda_Y)$ embeds.
An element $X$ of $\mathfrak{W}(S,\Lambda_Y)$ acts there as the operator $X_\Xi$ of Lemma~\ref{lem:normal-form}(f):
the realisation $X\mapsto X_\Xi$ is defined for every such cross-section, does not depend on the plaquette $\sigma_0$,
and is given on the generators by the formulas of Definition~\ref{def:mid-plane}.
The operator $k^0_Y$ of Definition~\ref{def:two-layer} is considered on $\mathcal{H}_S$, where $E$ stands for $E_S$,
and on $\mathcal{H}^{\mathrm{l}}$, where $\Upsilon_Y$ stands for $\iota(\Upsilon_Y)$;
the operator $k_Y$ acts on $\mathcal{H}^{\mathrm{l}}$.
Equalities between integrands are understood in the sense of Lemma~\ref{lem:equal-integrands}.

\begin{proposition}[Two-layer elements as insertions]\label{prop:two-layer-insertions}
Let $0<\rho<1$, $w_0,\dots,w_3>0$, $M\in\R$ and $q>0$.
Let $c=j_0+\frac12$ with $j_0\in\Z$,
and let $Y=\exp\left(i\langle\delta,\hat A\rangle\right)\,\Psi$ be a monomial of $\mathcal{P}_1(c)$ with $Y\neq1$ and with the exponent $\delta=(\delta_0,\delta_1,\delta_\perp)$.
Let $S$ be the set of the links $b$ of $\Z^2$ with $\delta_0(b)\neq0$ or $\delta_1(b)\neq0$,
and let $\Lambda_Y$ be a set of sites in play of $Y$ (Definition~\ref{def:two-layer}).
The charge densities of $\Psi$ on the planes $\{x_1=j_0\}$ and $\{x_1=j_0+1\}$ are the functions $\mathfrak{q}_-,\mathfrak{q}_+\in\Z^{\Lambda_Y}$ given by
\[
\mathfrak{q}_-(y):=\mathfrak{q}(\Psi)(j_0,y),\qquad
\mathfrak{q}_+(y):=\mathfrak{q}(\Psi)(j_0+1,y),\qquad
y\in\Lambda_Y,
\]
where $\mathfrak{q}(\Psi)$ is the charge density of Definition~\ref{def:chamber-operator};
both are extended by zero to $\Z^2$.
Let $\Xi$ be a rectangular cross-section in which $\hull(S,\Lambda_Y)$ embeds.
\begin{enumerate}
\item[(a)]
The exponent obeys Gauss' law on the two planes:
\begin{equation}\label{eq:amplification-gauss}
d_{\mathrm{s}}^{*}\delta_0-\delta_\perp=q\cdot\mathfrak{q}_-,
\qquad
d_{\mathrm{s}}^{*}\delta_1+\delta_\perp=q\cdot\mathfrak{q}_+ .
\end{equation}
The equation of Definition~\ref{def:two-layer} for $\Upsilon_Y$,
\[
\left.\Psi\exp\left(-S^{\times}_c(\Lambda_Y)\right)\right|_{\hat a_{j_0}=0}
=\Ins_c\left[\omega^{\frac12}\,\Upsilon\,\omega^{\frac12}\right],
\]
has exactly one solution $\Upsilon=\Upsilon_Y$ in $\End\left(\mathcal{F}_{\Lambda_Y}\right)$,
and $\Upsilon_Y\in\End_z(\Lambda_Y)$ with $z=\mathfrak{q}_-+\mathfrak{q}_+=z(\delta_0+\delta_1)$.
With
\[
k^0_Y(A',A):=\exp\left(i\langle\delta_1,A'\rangle\right)\prod_{b\in S}g_1\left(A'(b)-A(b)\right)\exp\left(i\langle\delta_0,A\rangle\right)\,\Upsilon_Y,
\qquad A',A\in\R^S,
\]
the triple $\left(S,\Lambda_Y,k^0_Y\right)$ is a local step datum,
$k^0_Y(A',A)$ is the integral kernel of the operator $k^0_Y$ on $\mathcal{H}_S$,
and $Y=G[k^0_Y]$, the step element being formed at the plane $c$.

The map $Y\mapsto\Upsilon_Y$ being extended linearly to $\mathcal{P}_1(c)$ as in Definition~\ref{def:two-layer},
let $\Psi$ be, instead of a monomial, a product of modes
$\Psi=\pm\Theta_c(\zeta_{J_0})\zeta_{I_0}$ with $I_0,J_0\subset\mathcal{J}^{\mathrm{l}}(\Lambda_Y)$,
such that $Y=\exp\left(i\langle\delta,\hat A\rangle\right)\,\Psi$ lies in $\mathcal{P}_1(c)$.
Then $\Upsilon_Y$ has at most one non-zero entry in each row and in each column,
all non-zero entries have the modulus $\left(\omega_{I_0}\cdot\omega_{J_0}\right)^{-\frac12}$,
and $\|\Upsilon_Y\|=\left(\omega_{I_0}\cdot\omega_{J_0}\right)^{-\frac12}$.
A monomial in the fields at the sites $x$ of $\mathbb{X}_{j_0,6}\cup\mathbb{X}_{j_0+1,1}$ with $\pi(x)\in\{j_0,j_0+1\}\times\Lambda_Y$
is a finite linear combination of products $\Theta_c(\zeta_J)\zeta_I$, $I,J\subset\mathcal{J}^{\mathrm{l}}(\Lambda_Y)$,
for which the charge density of the monomial equals $\mathfrak{q}(I)$ on the plane $\{x_1=j_0+1\}$ and $-\mathfrak{q}(J)$ on the plane $\{x_1=j_0\}$.
If, moreover, $\exp\left(i\langle\delta,\hat A\rangle\right)$ times the monomial is a monomial of $\mathcal{P}_1(c)$,
then $\exp\left(i\langle\delta,\hat A\rangle\right)\Theta_c(\zeta_J)\zeta_I$ lies in $\mathcal{P}_1(c)$ for each of these products.
\item[(b)]
The element $X'_Y$ lies in $\mathfrak{W}(S,\Lambda_Y)$, and $Y=G'_{X'_Y}$.
On $\mathcal{H}^{\mathrm{l}}$,
\[
k_Y=K_{0,S}^{\frac12}\,\left(X'_Y\right)_\Xi\,K_{0,S}^{\frac12};
\]
and for every $\beta\in\R^{\Lambda^{(1)}_{\mathrm{s}}}$ this identity remains true when,
on the left, the kernel with which $k^0_Y$ acts on $\mathcal{H}^{\mathrm{l}}$
(periodised over $\Gamma_{\mathrm{s}}$, as in Lemma~\ref{lem:mixed-word-formula})
is multiplied by $\exp\left(i\langle\beta,A'-A\rangle\right)$,
and, on the right, $K_{0,S}$ and $X'_Y$ are replaced by $K_{\beta,S}$ and $\alpha_\beta X'_Y$,
$\alpha_\beta$ being formed with the restriction of $\beta$ to $S$.
Moreover
\[
\|X'_Y\|_\#=\|\Upsilon_Y\|\cdot\exp\left(\tfrac1{16}|\delta_1-\delta_0|^2+\tfrac18\left(|\delta_0|^2+|\delta_1|^2\right)\right),
\]
and $Y$ is bounded by $2^{\frac{\#|S|}{2}}\cdot\|X'_Y\|_\#$ with the box $B^{\#}(c,S,\Lambda_Y)$ (Definition~\ref{def:bounded-elements}).
\item[(c)]
For every $\varepsilon>0$ the element $X_{Y,\varepsilon}$ lies in $\mathfrak{W}(S,\Lambda_Y)$,
\[
\exp\left(-\tfrac14|E_S|^2\right)X_{Y,\varepsilon}\exp\left(-\tfrac14|E_S|^2\right)=P_\varepsilon\,k^0_Y\,P_\varepsilon
\qquad\text{on }\mathcal{H}_S,
\]
and consequently $K^{\frac12}\,\left(X_{Y,\varepsilon}\right)_\Xi\,K^{\frac12}=P_\varepsilon\,k_Y\,P_\varepsilon$ on $\mathcal{H}^{\mathrm{l}}$.
The integrand $Y-G_{X_{Y,\varepsilon}}$ equals $G'_{X'_Y-P_\varepsilon X'_YP_\varepsilon}$
and is bounded by $2^{\frac{\#|S|}{2}}\cdot\xi_Y(\varepsilon)$ with the box $B^{\#}(c,S,\Lambda_Y)$, where
\[
\begin{aligned}
\xi_Y(\varepsilon)
={}&\|\Upsilon_Y\|\cdot\sup_{x\in\R^S}\Bigl(\exp\left(-\tfrac14|x|^2+\tfrac14\langle x,\delta_1-\delta_0\rangle+\tfrac18\left(|\delta_0|^2+|\delta_1|^2\right)\right)\\
&\qquad\times\left(1-\exp\left(-\varepsilon|x+\delta_1|^2-\varepsilon|x-\delta_0|^2\right)\right)\Bigr)
\;\le\;C_Y\cdot\varepsilon
\end{aligned}
\]
with a constant $C_Y<\infty$ that does not depend on $\varepsilon$.
Finally
\[
\|X_{Y,\varepsilon}\|_\#=\|\Upsilon_Y\|\cdot\exp\left(\frac{\left(\frac12-2\varepsilon\right)^2|\delta_1-\delta_0|^2}{8\varepsilon}+\left(\tfrac14-\varepsilon\right)\left(|\delta_0|^2+|\delta_1|^2\right)\right).
\]
\item[(d)]
The element $\Theta_cY$ lies in $\mathcal{P}_1(c)$:
it is a finite linear combination of monomials with the exponent $(-\delta_1,-\delta_0,\delta_\perp)$.
Each of them has the set of links $S$, and $\Lambda_Y$ is a set of sites in play for each of them,
because $\theta_c$ exchanges the sites over $(j_0,y)$ with those over $(j_0+1,y)$.
For the linear extensions of the maps of Definition~\ref{def:two-layer}, formed with $S$ and $\Lambda_Y$,
\[
\Upsilon_{\Theta_cY}=\Upsilon_Y^{*},\qquad
X'_{\Theta_cY}=\left(X'_Y\right)^{*},\qquad
X_{\Theta_cY,\varepsilon}=\left(X_{Y,\varepsilon}\right)^{*}\quad(\varepsilon>0);
\]
moreover $k^0_{\Theta_cY}=\left(k^0_Y\right)^{*}$ holds both on $\mathcal{H}_S$ and on $\mathcal{H}^{\mathrm{l}}$,
and $k_{\Theta_cY}=k_Y^{*}$ on $\mathcal{H}^{\mathrm{l}}$.
\end{enumerate}
\end{proposition}

\begin{proof}
(a) \emph{Gauss' law.}
The element $Y$ of $\mathcal{P}$ is gauge invariant,
which means $d^{*}\delta=q\cdot\mathfrak{q}(\Psi)$ (Definition~\ref{def:polynomials}).
Here $d^{*}$ is the adjoint of the coboundary $d$ with respect to the pairing of Definition~\ref{def:cochains}.
As $d$ maps a zero-cochain $\lambda$ to the one-cochain with the value $\lambda(x+e_\mu)-\lambda(x)$ at the link $(x,\mu)=[x,x+e_\mu]$, one has
\[
\left(d^{*}\delta\right)(x)=\sum_{\mu=1}^{3}\left(\delta(x-e_\mu,\mu)-\delta(x,\mu)\right),
\]
where $\delta(x,\mu)$ is the value of $\delta$ at the link $(x,\mu)$;
in the same way $\left(d_{\mathrm{s}}^{*}u\right)(y)=\sum_{\nu=2,3}\left(u(y-e_\nu,\nu)-u(y,\nu)\right)$
for a finitely supported one-cochain $u$ of the slice lattice,
$d_{\mathrm{s}}$ being its coboundary (Definition~\ref{def:slices}).
The cochain $\delta$ vanishes on the links from the plane $\{x_1=j_0-1\}$ to the plane $\{x_1=j_0\}$
and on those from $\{x_1=j_0+1\}$ to $\{x_1=j_0+2\}$.
At a site $(j_0,y)$ the terms with $\mu=1$ therefore contribute $-\delta_\perp(y)$,
and those with $\mu=2,3$ contribute $\left(d_{\mathrm{s}}^{*}\delta_0\right)(y)$;
at $(j_0+1,y)$ they contribute $\delta_\perp(y)$ and $\left(d_{\mathrm{s}}^{*}\delta_1\right)(y)$.
At these two sites $q\cdot\mathfrak{q}(\Psi)$ takes the values $q\cdot\mathfrak{q}_-(y)$ and $q\cdot\mathfrak{q}_+(y)$,
and \eqref{eq:amplification-gauss} follows.

\emph{Existence, uniqueness and charge of $\Upsilon_Y$.}
Let $\mathbb{X}'$ be the set of the sites of $\mathbb{X}_{j_0,6}$ and of $\mathbb{X}_{j_0+1,1}$ over $\Lambda_Y$,
that is, of the two ends of the bonds over $\Lambda_Y$ that cross the plane $\{x_1=c\}$.
By Definition~\ref{def:modes} the modes $\zeta_i$, $i\in\mathcal{J}^{\mathrm{l}}(\Lambda_Y)$,
form a basis of the span of the generators $\psi_\alpha(x)$, $\bar\psi_\alpha(x)$ at the upper ends,
and the $\Theta_c(\zeta_i)$ form a basis of the span of the generators at the lower ends;
together they form a basis of the span of all generators at the sites of $\mathbb{X}'$.
A product of distinct elements of a basis of the generators is, up to a sign, an ordered monomial in this basis.
So the elements $\Theta_c(\zeta_J)\zeta_I$, $I,J\subset\mathcal{J}^{\mathrm{l}}(\Lambda_Y)$, form a basis of $\Gr(\mathbb{X}')$,
and every element of $\Gr(\mathbb{X}')$ equals $\Ins_c[\Upsilon]$ for exactly one matrix $\Upsilon=\left(\Upsilon_{IJ}\right)$ indexed by these sets.
By the choice of $\Lambda_Y$ the monomial $\Psi$ lies in $\Gr(\mathbb{X}')$.
At the plane $c$ the crossing bonds are of type $3$, the coefficients of their modes are $\omega_i=\frac{w_3}{2}(1\pm\rho)$,
and their phases are $u_I=\exp\left(iq\langle\hat a_{j_0},\mathfrak{q}(I)\rangle\right)$ (Definition~\ref{def:modes});
the phases equal $1$ at $\hat a_{j_0}=0$.
Proposition~\ref{prop:crossing}, applied to the crossing bonds over $\Lambda_Y$, gives
\[
\left.\exp\left(-S^{\times}_c(\Lambda_Y)\right)\right|_{\hat a_{j_0}=0}
=\sum_{R\subset\mathcal{J}^{\mathrm{l}}(\Lambda_Y)}\omega_R\cdot\Theta_c(\zeta_R)\zeta_R
=\Ins_c[\omega].
\]
Since $|\rho|<1$ and $w_3>0$, all $\omega_i$ are positive, and the diagonal operator $\omega$ is invertible.
Let $\Upsilon'$ be the matrix with
$\left.\Psi\exp\left(-S^{\times}_c(\Lambda_Y)\right)\right|_{\hat a_{j_0}=0}=\Ins_c[\Upsilon']$.
The defining equation of $\Upsilon_Y$ then reads $\omega^{\frac12}\,\Upsilon_Y\,\omega^{\frac12}=\Upsilon'$,
and its only solution is $\Upsilon_Y=\omega^{-\frac12}\,\Upsilon'\,\omega^{-\frac12}$.

To determine the charge, call an element of $\Gr(\mathbb{X}')$ homogeneous
if it is a linear combination of monomials in the generators that all have the same charge density,
and call this common charge density the charge density of the element.
The monomials form a basis;
so $\Gr(\mathbb{X}')$ is the direct sum of the spaces of homogeneous elements of the various charge densities,
and a product of homogeneous elements is homogeneous, its charge density being the sum of theirs.
By Lemma~\ref{lem:reflected-modes} every basis element $\Theta_c(\zeta_J)\zeta_I$ is homogeneous,
with the charge density $\mathfrak{q}(I)$ on the plane $\{x_1=j_0+1\}$ and $-\mathfrak{q}(J)$ on the plane $\{x_1=j_0\}$.
Consequently, if a homogeneous element is written as $\Ins_c[\Upsilon]$,
then $\Upsilon_{IJ}=0$ unless $\Theta_c(\zeta_J)\zeta_I$ has the charge density of that element.
The monomial $\Psi$ is homogeneous, with the charge densities $\mathfrak{q}_+$ on the upper and $\mathfrak{q}_-$ on the lower plane,
and $\Theta_c(\zeta_R)\zeta_R$ is homogeneous, with $\mathfrak{q}(R)$ and $-\mathfrak{q}(R)$.
Write $\omega_R\cdot\Psi\,\Theta_c(\zeta_R)\zeta_R=\Ins_c[\Upsilon'_R]$ for $R\subset\mathcal{J}^{\mathrm{l}}(\Lambda_Y)$.
Then $\Upsilon'=\sum_R\Upsilon'_R$, and $\left(\Upsilon'_R\right)_{IJ}\neq0$ only if
\[
\mathfrak{q}(I)=\mathfrak{q}_++\mathfrak{q}(R),\qquad -\mathfrak{q}(J)=\mathfrak{q}_--\mathfrak{q}(R).
\]
Thus $\Upsilon'_{IJ}\neq0$ implies $\mathfrak{q}(I)-\mathfrak{q}(J)=\mathfrak{q}_++\mathfrak{q}_-=z$.
As $\omega$ is diagonal, the entries of $\Upsilon_Y$ vanish outside the same set of pairs $(I,J)$;
as $u(a)$ is diagonal with the entries $\exp\left(iq\langle a,\mathfrak{q}(I)\rangle\right)$, it follows that
\[
u(a)\,\Upsilon_Y\,u(a)^{-1}=\exp\left(iq\langle a,z\rangle\right)\,\Upsilon_Y\qquad\text{for all }a\in\R^{\Lambda_Y},
\]
that is, $\Upsilon_Y\in\End_z(\Lambda_Y)$ (Definition~\ref{def:charged-operators}).
The sum of the two identities \eqref{eq:amplification-gauss} is $q\cdot z=d_{\mathrm{s}}^{*}(\delta_0+\delta_1)$;
hence $\delta_0+\delta_1$ is an exponent in $\mathfrak{D}(S)$, and $z=z(\delta_0+\delta_1)$ (Definition~\ref{def:gaussian-symbols}).
Pairing $d_{\mathrm{s}}^{*}(\delta_0+\delta_1)$ with the constant function $1$, for which $d_{\mathrm{s}}1=0$,
shows that $\sum_y z(y)=0$;
by Lemma~\ref{lem:fock-parity} the operator $\Upsilon_Y$ is even.

\emph{Monomials and products of modes.}
Let a monomial in the generators at the sites $x$ of $\mathbb{X}_{j_0,6}\cup\mathbb{X}_{j_0+1,1}$ with $\pi(x)\in\{j_0,j_0+1\}\times\Lambda_Y$ be given.
It lies in $\Gr(\mathbb{X}')$ and is homogeneous,
and its charge density vanishes outside the points $(j_0,y)$ and $(j_0+1,y)$, $y\in\Lambda_Y$.
Hence it equals $\Ins_c[\Upsilon]=\sum_{I,J}\Upsilon_{IJ}\cdot\Theta_c(\zeta_J)\zeta_I$ for a matrix $\Upsilon$ with $\Upsilon_{IJ}=0$
unless the charge density of the monomial equals $\mathfrak{q}(I)$ on the plane $\{x_1=j_0+1\}$ and $-\mathfrak{q}(J)$ on the plane $\{x_1=j_0\}$.
Suppose now that $\exp\left(i\langle\delta,\hat A\rangle\right)$ times the monomial is a monomial of $\mathcal{P}_1(c)$;
then $d^{*}\delta$ is $q$ times the charge density of the monomial (Definition~\ref{def:polynomials}).
Let $I,J\subset\mathcal{J}^{\mathrm{l}}(\Lambda_Y)$ be such that $\Theta_c(\zeta_J)\zeta_I$ has this charge density.
Every monomial in the expansion of $\Theta_c(\zeta_J)\zeta_I$ involves only generators at sites of $\mathbb{X}_{j_0,6}\cup\mathbb{X}_{j_0+1,1}$ and has the same charge density;
multiplied by $\exp\left(i\langle\delta,\hat A\rangle\right)$, it is again gauge invariant, hence a monomial of $\mathcal{P}_1(c)$ (Definition~\ref{def:two-layer}).
So $\exp\left(i\langle\delta,\hat A\rangle\right)\Theta_c(\zeta_J)\zeta_I$ lies in $\mathcal{P}_1(c)$,
and the last two assertions of (a) are proved.

\emph{The step datum.}
Let $a\in\R^{\Lambda_Y}$, extended by zero to $\Z^2$.
The values of $d_{\mathrm{s}}a$ on the links of $S$ involve only the values of $a$ at the end points of these links, which lie in $\Lambda_Y$,
and by \eqref{eq:amplification-gauss}
\[
\langle\delta_0+\delta_1,d_{\mathrm{s}}a\rangle=\langle d_{\mathrm{s}}^{*}(\delta_0+\delta_1),a\rangle=q\cdot\langle z,a\rangle .
\]
Since $\Upsilon_Y\in\End_z(\Lambda_Y)$, one has $u(a)^{-1}\,\Upsilon_Y\,u(a)=\exp\left(-iq\langle a,z\rangle\right)\,\Upsilon_Y$;
and the factors $g_1\left(A'(b)-A(b)\right)$ do not change when $d_{\mathrm{s}}a$ is added to both $A'$ and $A$.
Therefore
\[
u(a)^{-1}\,k^0_Y\left(A'+(d_{\mathrm{s}}a)|_S,\,A+(d_{\mathrm{s}}a)|_S\right)u(a)
=\exp\left(i\langle\delta_0+\delta_1,d_{\mathrm{s}}a\rangle-iq\langle a,z\rangle\right)k^0_Y(A',A)
=k^0_Y(A',A),
\]
so that $k^0_Y$ is gauge covariant.
Its values are even, because $\Upsilon_Y$ is even, and
\[
\left\|k^0_Y(A',A)\right\|
=\|\Upsilon_Y\|\cdot\prod_{b\in S}g_1\left(A'(b)-A(b)\right)
=(2\pi)^{-\frac{\#|S|}{2}}\cdot\|\Upsilon_Y\|\cdot\prod_{b\in S}\exp\left(-\tfrac12\left(A'(b)-A(b)\right)^2\right),
\]
which is the Gaussian bound of Proposition~\ref{prop:word-formula} with the decay constant $\varepsilon_k=\frac12$.
The function $k^0_Y$ is continuous, hence Borel;
thus $\left(S,\Lambda_Y,k^0_Y\right)$ is a local step datum.

On $\mathcal{H}_S$ the operator of Definition~\ref{def:two-layer} reads
$k^0_Y=\exp\left(i\langle\delta_0+\delta_1,A\rangle\right)\,f(E_S)\,\Upsilon_Y$
with $f(\xi):=\exp\left(-\frac12|\xi+\delta_0|^2\right)$, $\xi\in\R^S$,
a bounded integrable function.
Its inverse Fourier transform (Definition~\ref{def:integral-kernels}) at a point $x\in\R^S$ is
\[
f^{\vee}(x)=(2\pi)^{-\#|S|}\int_{\R^S}\mathrm{d}\xi\,\exp\left(-\tfrac12|\xi+\delta_0|^2+i\langle\xi,x\rangle\right).
\]
The substitution $\xi\mapsto\xi-\delta_0$ takes the factor $\exp\left(-i\langle\delta_0,x\rangle\right)$ out of this integral,
and the Gaussian integral that remains factorises over the coordinates $\xi_b$, $b\in S$.
Since $(2\pi)^{-1}\int_{\R}\mathrm{d}\xi_b\,\exp\left(-\frac12\xi_b^2+i\xi_bx_b\right)=g_1(x_b)$ for every $b\in S$, it follows that
\[
f^{\vee}(x)=\exp\left(-i\langle\delta_0,x\rangle\right)\prod_{b\in S}g_1(x_b),
\qquad x\in\R^S.
\]
By Definition~\ref{def:integral-kernels}, $k^0_Y$ is the operator with the kernel
\[
\exp\left(i\langle\delta_0+\delta_1,A'\rangle\right)\,f^{\vee}(A'-A)\,\Upsilon_Y
=\exp\left(i\langle\delta_1,A'\rangle\right)\prod_{b\in S}g_1\left(A'(b)-A(b)\right)\exp\left(i\langle\delta_0,A\rangle\right)\,\Upsilon_Y
=k^0_Y(A',A).
\]

\emph{The identity $Y=G[k^0_Y]$.}
In the step element of the local step datum $\left(S,\Lambda_Y,k^0_Y\right)$ (Proposition~\ref{prop:word-formula}) the kernel is evaluated at
\[
A'=\hat A_{j_0+1}|_S,\qquad A=\left(\hat A_{j_0}+d_{\mathrm{s}}\hat a_{j_0}\right)|_S,
\]
and at these arguments $A'-A=F^{\mathrm{el}}_{j_0}|_S$ (Definition~\ref{def:slices}).
The factor $\prod_{b\in S}g_1\left(A'(b)-A(b)\right)$ of $k^0_Y(A',A)$ therefore cancels against the factor $\prod_{b\in S}g_1\left(F^{\mathrm{el}}_{j_0}(b)\right)^{-1}$ of the step element.
As $\delta_0$ vanishes outside $S$ and $\langle\delta_0,d_{\mathrm{s}}\hat a_{j_0}\rangle=\langle d_{\mathrm{s}}^{*}\delta_0,\hat a_{j_0}\rangle$, one obtains
\begin{multline*}
G\left[k^0_Y\right]=\exp\left(i\langle\delta_1,\hat A_{j_0+1}\rangle+i\langle\delta_0,\hat A_{j_0}\rangle+i\langle d_{\mathrm{s}}^{*}\delta_0,\hat a_{j_0}\rangle\right)\\
\cdot\Ins_c\left[\omega^{\frac12}\,\Upsilon_Y\,u(\hat a_{j_0})\,\omega^{\frac12}\right]\cdot\exp\left(S^{\times}_c(\Lambda_Y)\right),
\end{multline*}
where $u(\hat a_{j_0})$ is formed with the restriction of $\hat a_{j_0}$ to $\Lambda_Y$.

The crossing terms are expanded as in the construction of $\Upsilon_Y$, but now at an arbitrary value of $\hat a_{j_0}$.
The phases $u_I=\exp\left(iq\langle\hat a_{j_0},\mathfrak{q}(I)\rangle\right)$ of the crossing modes are the diagonal entries of $u(\hat a_{j_0})$,
and Proposition~\ref{prop:crossing} gives
\[
\exp\left(-S^{\times}_c(\Lambda_Y)\right)=\Ins_c\left[\omega^{\frac12}\,u(\hat a_{j_0})\,\omega^{\frac12}\right]=\Ins_c\left[\omega\,u(\hat a_{j_0})\right];
\]
at $\hat a_{j_0}=0$ it reduces to the expansion $\Ins_c[\omega]$ of that construction.
The identity that defines $\Upsilon_Y$ (Definition~\ref{def:two-layer}) thus reads
\[
\Psi\cdot\Ins_c[\omega]=\Ins_c\left[\omega^{\frac12}\,\Upsilon_Y\,\omega^{\frac12}\right].
\]
Let $a\in\R^{\Lambda_Y}$, and let $\alpha^{-}_a$ be the automorphism of the Grassmann algebra of Lemma~\ref{lem:lower-phases},
which multiplies the fields at the lower ends of the crossing bonds over $\Lambda_Y$ by phases.
By that lemma, $\alpha^{-}_a\left(\Ins_c[\Upsilon]\right)=\Ins_c\left[\Upsilon\,u(a)\right]$ for every matrix $\Upsilon$ indexed by the subsets of $\mathcal{J}^{\mathrm{l}}(\Lambda_Y)$,
and $\alpha^{-}_a(\Psi)=\exp\left(-iq\langle\mathfrak{q}_-,a\rangle\right)\cdot\Psi$, since $\mathfrak{q}_-$ is the charge density of $\Psi$ on the plane $\{x_1=j_0\}$.
Applying the automorphism $\alpha^{-}_a$ to both sides of the last identity, and using that the diagonal matrices $\omega^{\frac12}$ and $u(a)$ commute, one finds
\[
\begin{aligned}
\exp\left(-iq\langle\mathfrak{q}_-,a\rangle\right)\cdot\Psi\cdot\Ins_c\left[\omega\,u(a)\right]
&=\Ins_c\left[\omega^{\frac12}\,\Upsilon_Y\,\omega^{\frac12}\,u(a)\right]\\
&=\Ins_c\left[\omega^{\frac12}\,\Upsilon_Y\,u(a)\,\omega^{\frac12}\right].
\end{aligned}
\]
For $a$ the restriction of $\hat a_{j_0}$ to $\Lambda_Y$, the left side is $\exp\left(-iq\langle\mathfrak{q}_-,\hat a_{j_0}\rangle\right)\cdot\Psi\cdot\exp\left(-S^{\times}_c(\Lambda_Y)\right)$.
Substituting this into the formula for $G[k^0_Y]$, the factors $\exp\left(\mp S^{\times}_c(\Lambda_Y)\right)$ cancel, and
\[
G\left[k^0_Y\right]=\exp\left(i\langle\delta_1,\hat A_{j_0+1}\rangle+i\langle\delta_0,\hat A_{j_0}\rangle+i\langle d_{\mathrm{s}}^{*}\delta_0-q\cdot\mathfrak{q}_-,\hat a_{j_0}\rangle\right)\cdot\Psi .
\]
By the first identity of Gauss' law in (a), $d_{\mathrm{s}}^{*}\delta_0-q\cdot\mathfrak{q}_-=\delta_\perp$.
The exponent is therefore $i\langle\delta_0,\hat A_{j_0}\rangle+i\langle\delta_1,\hat A_{j_0+1}\rangle+i\langle\delta_\perp,\hat a_{j_0}\rangle=i\langle\delta,\hat A\rangle$,
and $G[k^0_Y]=Y$.

\emph{The entries of $\Upsilon_Y$.}
Let now $\Psi=\pm\Theta_c(\zeta_{J_0})\,\zeta_{I_0}$ be a product of modes, with $I_0,J_0\subset\mathcal{J}^{\mathrm{l}}(\Lambda_Y)$.
Such a product is a finite combination of monomials, so $Y$ is a finite combination of two-layer monomials with the exponent $(\delta_0,\delta_1,\delta_\perp)$,
and $\Upsilon_Y$ is the operator of the linear extension in Definition~\ref{def:two-layer};
by linearity, $\Psi\cdot\Ins_c[\omega]=\Ins_c\left[\omega^{\frac12}\,\Upsilon_Y\,\omega^{\frac12}\right]$ again.
In this identity $\Ins_c[\omega]$ is, as found in the construction of $\Upsilon_Y$, the sum of the elements $\omega_R\cdot\Theta_c(\zeta_R)\,\zeta_R$ over all $R\subset\mathcal{J}^{\mathrm{l}}(\Lambda_Y)$.
The modes and their images under $\Theta_c$ are odd elements with vanishing squares.
Moving $\Theta_c(\zeta_R)$ to the left of $\zeta_{I_0}$ gives
\[
\Theta_c(\zeta_{J_0})\,\zeta_{I_0}\,\Theta_c(\zeta_R)\,\zeta_R=\pm\Theta_c(\zeta_{J_0})\,\Theta_c(\zeta_R)\,\zeta_{I_0}\,\zeta_R ;
\]
the right side vanishes if $R$ meets $I_0$ or $J_0$, since a factor then occurs twice,
and it equals $\pm\Theta_c(\zeta_{J_0\cup R})\,\zeta_{I_0\cup R}$ otherwise.
Both sides of $\Psi\cdot\Ins_c[\omega]=\Ins_c\left[\omega^{\frac12}\,\Upsilon_Y\,\omega^{\frac12}\right]$ lie in $\Gr(\mathbb{X}')$,
in which the elements $\Theta_c(\zeta_J)\,\zeta_I$, $I,J\subset\mathcal{J}^{\mathrm{l}}(\Lambda_Y)$, form a basis;
comparing their coefficients, one finds
\[
\left(\omega^{\frac12}\,\Upsilon_Y\,\omega^{\frac12}\right)_{I_0\cup R,\,J_0\cup R}=\pm\omega_R\qquad\text{if }R\cap\left(I_0\cup J_0\right)=\emptyset,
\]
while all other entries vanish.
Since $\omega_{I\cup R}=\omega_I\cdot\omega_R$ for disjoint $I$ and $R$, it follows that
\[
\left(\Upsilon_Y\right)_{I_0\cup R,\,J_0\cup R}=\pm\omega_R\cdot\left(\omega_{I_0\cup R}\cdot\omega_{J_0\cup R}\right)^{-\frac12}=\pm\left(\omega_{I_0}\cdot\omega_{J_0}\right)^{-\frac12}.
\]
A row $I$ of $\Upsilon_Y$ contains a non-zero entry only if $I\supset I_0$, and then only in the column $J_0\cup(I\setminus I_0)$;
in the same way, a column $J$ contains a non-zero entry only in the row $I_0\cup(J\setminus J_0)$.
Thus $\Upsilon_Y$ has at most one non-zero entry in each row and in each column, and all these entries have the modulus $\left(\omega_{I_0}\cdot\omega_{J_0}\right)^{-\frac12}$.
The matrix $\left(\omega_{I_0}\cdot\omega_{J_0}\right)^{\frac12}\cdot\Upsilon_Y$ has the entries $0$ and $\pm1$, with at most one non-zero entry in each row and in each column;
its product with its adjoint on the left is a diagonal matrix with the entries $0$ and $1$.
It is therefore a partial isometry, and it is not zero, its entry in the row $I_0$ and the column $J_0$ being $\pm1$.
Whatever the signs, $\left\|\Upsilon_Y\right\|=\left(\omega_{I_0}\cdot\omega_{J_0}\right)^{-\frac12}$.

(b) \emph{The element $X'_Y$.}
Since $\widetilde E=E_S+\delta_0$ and $\Upsilon_Y$ commutes with the functions of $E_S$, the definition of $X'_Y$ reads
\[
X'_Y=\left(\exp\left(i\langle\delta_0+\delta_1,A\rangle\right)\cdot\Upsilon_Y\right)\cdot f'(E_S),
\]
where the function $f'$ on $\R^S$ is
\[
f'(\xi):=\exp\left(-\tfrac14|\xi+\delta_0|^2+\tfrac14\langle\xi+\delta_0,\delta_1-\delta_0\rangle+\tfrac18\left(|\delta_0|^2+|\delta_1|^2\right)\right).
\]
By (a), $\Upsilon_Y\in\End_z(\Lambda_Y)$ with $z=\mathfrak{q}_-+\mathfrak{q}_+=z(\delta_0+\delta_1)$:
the operator $\Upsilon_Y$ is homogeneous of the charge $z$, and $d_{\mathrm{s}}^{*}(\delta_0+\delta_1)=q\cdot z$.
So $\exp\left(i\langle\delta_0+\delta_1,A\rangle\right)\cdot\Upsilon_Y$ is a generator of type~\eqref{it:exponential-generators} of $\mathfrak{W}(S,\Lambda_Y)$ (Definition~\ref{def:mid-plane}).
Expanding the square, one finds that $f'$ is a positive multiple of the function $\xi\mapsto\exp\left(-\tfrac14|\xi|^2+\tfrac14\langle\delta_1-3\delta_0,\xi\rangle\right)$.
Hence $f'\in\mathfrak{A}(S)$ (Definition~\ref{def:gaussian-symbols}), with the quadratic form $\mathfrak{Q}=\tfrac14|\cdot|^2$, and $f'=1$ if $S=\emptyset$;
the operator $f'(E_S)$ is a multiple of a generator of type~\eqref{it:gaussian-generators}, and $X'_Y\in\mathfrak{W}(S,\Lambda_Y)$.

\emph{The kernel on $\mathcal{H}_S$.}
The relation~\eqref{eq:setting-weyl}, read on $\mathcal{H}_S$ with the vector $\delta_0+\delta_1$ and the function $\xi\mapsto\exp\left(-\tfrac18|\xi|^2\right)$, gives
\[
\exp\left(-\tfrac18|E_S|^2\right)\cdot\exp\left(i\langle\delta_0+\delta_1,A\rangle\right)=\exp\left(i\langle\delta_0+\delta_1,A\rangle\right)\cdot\exp\left(-\tfrac18|E_S+\delta_0+\delta_1|^2\right).
\]
As $E_S+\delta_0+\delta_1=\widetilde E+\delta_1$ and $E_S=\widetilde E-\delta_0$, it follows that
\begin{multline*}
\exp\left(-\tfrac18|E_S|^2\right)X'_Y\exp\left(-\tfrac18|E_S|^2\right)\\
=\exp\left(i\langle\delta_0+\delta_1,A\rangle\right)\cdot\exp\left(-\tfrac18|\widetilde E+\delta_1|^2-\tfrac18|\widetilde E-\delta_0|^2\right)\cdot f'(E_S)\cdot\Upsilon_Y ,
\end{multline*}
where
\[
f'(E_S)=\exp\left(-\tfrac14|\widetilde E|^2+\tfrac14\langle\widetilde E,\delta_1-\delta_0\rangle+\tfrac18\left(|\delta_0|^2+|\delta_1|^2\right)\right).
\]
For $x\in\R^S$ one has
\[
\tfrac18|x+\delta_1|^2+\tfrac18|x-\delta_0|^2=\tfrac14|x|^2+\tfrac14\langle x,\delta_1-\delta_0\rangle+\tfrac18\left(|\delta_0|^2+|\delta_1|^2\right).
\]
The linear and the constant terms cancel, and
\begin{equation}\label{eq:amplification-kernel-prime}
\begin{split}
\exp\left(-\tfrac18|E_S|^2\right)X'_Y\exp\left(-\tfrac18|E_S|^2\right)
&=\exp\left(i\langle\delta_0+\delta_1,A\rangle\right)\cdot\exp\left(-\tfrac12|\widetilde E|^2\right)\cdot\Upsilon_Y\\
&=k^0_Y .
\end{split}
\end{equation}
In particular $k^0_Y\in\mathfrak{W}(S,\Lambda_Y)$.
By (a), $k^0_Y(A',A)$ is the kernel of the operator $k^0_Y$ on $\mathcal{H}_S$ in the sense of Definition~\ref{def:integral-kernels}, and it is continuous;
the kernel $k^{(\frac18)}_{X'_Y}$ of the left side is continuous as well (Lemma~\ref{lem:regularised-kernels}).
Two continuous kernels of one operator coincide.
Thus $k^{(\frac18)}_{X'_Y}=k^0_Y$, and Definition~\ref{def:two-insertions} together with (a) gives
\[
G'_{X'_Y}=G\left[k^{(\frac18)}_{X'_Y}\right]=G\left[k^0_Y\right]=Y .
\]

\emph{The operator $k_Y$.}
Let $\Xi$ be the cross-section of the proposition.
On $\mathcal{H}^{\mathrm{l}}$ the vector $E=(E_b)_{b\in\Lambda^{(1)}_{\mathrm{s}}}$ has the components $E_b$ with $b\in S$, which form $E_S$, and those with $b\in\Lambda^{(1)}_{\mathrm{s}}\setminus S$;
here and below, a sum over $b\notin S$ extends over the links $b\in\Lambda^{(1)}_{\mathrm{s}}\setminus S$.
By Definition~\ref{def:slice-operators},
\[
K_{0,S}^{\frac12}=(2\pi/q)^{\frac{|\Lambda_{\mathrm{s}}|}{2}}\cdot P_{\mathrm{G}}\cdot\exp\left(-\tfrac14\sum_{b\notin S}E_b^2-\tfrac18|E_S|^2\right),
\]
and $P_{\mathrm{G}}$ commutes with the functions of $E$.
The realisation $X\mapsto X_\Xi$ is a $*$-representation;
it represents the functions of $E_S$ and the operator $\exp\left(i\langle\delta_0+\delta_1,A\rangle\right)$ by the same formulas and $\Upsilon_Y$ by $\iota(\Upsilon_Y)$,
and for every $X\in\mathfrak{W}(S,\Lambda_Y)$ the operator $X_\Xi$ commutes with $P_{\mathrm{G}}$ and with the functions of the $E_b$, $b\notin S$ (Lemma~\ref{lem:normal-form}).
With~\eqref{eq:amplification-kernel-prime} it follows that
\[
K_{0,S}^{\frac12}\left(X'_Y\right)_\Xi K_{0,S}^{\frac12}
=(2\pi/q)^{|\Lambda_{\mathrm{s}}|}\cdot P_{\mathrm{G}}\cdot\exp\left(-\tfrac12\sum_{b\notin S}E_b^2\right)\cdot\left(k^0_Y\right)_\Xi,
\]
where
\[
\left(k^0_Y\right)_\Xi=\exp\left(i\langle\delta_0+\delta_1,A\rangle\right)\cdot\exp\left(-\tfrac12|E_S+\delta_0|^2\right)\cdot\iota(\Upsilon_Y).
\]
Since $\delta_0$ vanishes outside $S$, the operator $k^0_Y$ on $\mathcal{H}^{\mathrm{l}}$ of Definition~\ref{def:two-layer} is
\begin{equation}\label{eq:amplification-kernel-link}
k^0_Y=\exp\left(-\tfrac12\sum_{b\notin S}E_b^2\right)\cdot\left(k^0_Y\right)_\Xi ,
\end{equation}
and it commutes with $P_{\mathrm{G}}$.
Consequently
\[
K_{0,S}^{\frac12}\left(X'_Y\right)_\Xi K_{0,S}^{\frac12}=(2\pi/q)^{|\Lambda_{\mathrm{s}}|}\cdot k^0_Y\cdot P_{\mathrm{G}}=k_Y .
\]

\emph{Shifted comparison operators.}
Let $\beta\in\R^{\Lambda^{(1)}_{\mathrm{s}}}$; the element $\alpha_\beta X'_Y$ is formed with the restriction $\beta|_S$.
By Definition~\ref{def:mid-plane},
\[
\alpha_\beta X'_Y=\left(\exp\left(i\langle\delta_0+\delta_1,A\rangle\right)\cdot\Upsilon_Y\right)\cdot f'\left(E_S-\beta|_S\right),
\]
and by Definition~\ref{def:slice-operators}
\begin{multline*}
K_{\beta,S}^{\frac12}=(2\pi/q)^{\frac{|\Lambda_{\mathrm{s}}|}{2}}\cdot P_{\mathrm{G}}\\
\cdot\exp\left(-\tfrac14\sum_{b\notin S}\left(E_b-\beta(b)\right)^2-\tfrac18\sum_{b\in S}\left(E_b-\beta(b)\right)^2\right).
\end{multline*}
The computation that leads to~\eqref{eq:amplification-kernel-prime}, carried out with $E_S-\beta|_S$ in the place of $E_S$, and the argument just given for $k_Y$ yield
\begin{multline*}
K_{\beta,S}^{\frac12}\left(\alpha_\beta X'_Y\right)_\Xi K_{\beta,S}^{\frac12}\\
=(2\pi/q)^{|\Lambda_{\mathrm{s}}|}\cdot P_{\mathrm{G}}\cdot\exp\left(i\langle\delta_0+\delta_1,A\rangle\right)\cdot f_\beta(E)\cdot\iota(\Upsilon_Y),
\end{multline*}
where the function $f_\beta$ on $\R^{\Lambda^{(1)}_{\mathrm{s}}}$ is
\[
f_\beta(\xi):=\exp\left(-\tfrac12\sum_{b\notin S}\left(\xi_b-\beta(b)\right)^2-\tfrac12\sum_{b\in S}\left(\xi_b-\beta(b)+\delta_0(b)\right)^2\right).
\]
It is a Schwartz function with the inverse Fourier transform
\[
f_\beta^{\vee}(x)=\exp\left(i\langle\beta-\delta_0,x\rangle\right)\cdot\prod_{b\in\Lambda^{(1)}_{\mathrm{s}}}g_1(x_b),
\]
and the function $\exp\left(i\langle\delta_0+\delta_1,\cdot\,\rangle\right)$ is $\Gamma_{\mathrm{s}}$-periodic (Definition~\ref{def:two-layer}).
By Lemma~\ref{lem:mixed-word-formula}(a), the operator $\exp\left(i\langle\delta_0+\delta_1,A\rangle\right)\cdot f_\beta(E)\cdot\iota(\Upsilon_Y)$ on $\mathcal{H}^{\mathrm{l}}$
is therefore the operator with the kernel
\[
\exp\left(i\langle\beta,A'-A\rangle\right)\cdot\iota\left(k^0_Y\left(A'|_S,A|_S\right)\right)\cdot\prod_{b\notin S}g_1\left(A'(b)-A(b)\right),
\]
$A',A\in\R^{\Lambda^{(1)}_{\mathrm{s}}}$, periodised over $\Gamma_{\mathrm{s}}$; for $\beta=0$ it is the operator $k^0_Y$ on $\mathcal{H}^{\mathrm{l}}$.
Like $k^0_Y$, this operator commutes with $P_{\mathrm{G}}$:
it is the product of $\exp\left(-\tfrac12\sum_{b\notin S}\left(E_b-\beta(b)\right)^2\right)$ with the realisation of the element
\[
\exp\left(-\tfrac18\sum_{b\in S}\left(E_b-\beta(b)\right)^2\right)\left(\alpha_\beta X'_Y\right)\exp\left(-\tfrac18\sum_{b\in S}\left(E_b-\beta(b)\right)^2\right)
\]
of $\mathfrak{W}(S,\Lambda_Y)$.
Hence $K_{\beta,S}^{\frac12}\left(\alpha_\beta X'_Y\right)_\Xi K_{\beta,S}^{\frac12}$ is $(2\pi/q)^{|\Lambda_{\mathrm{s}}|}$ times the product of this operator with $P_{\mathrm{G}}$:
the identity $k_Y=K_{0,S}^{\frac12}\left(X'_Y\right)_\Xi K_{0,S}^{\frac12}$ remains true when the kernel of $k^0_Y$ on $\mathcal{H}^{\mathrm{l}}$ is multiplied by $\exp\left(i\langle\beta,A'-A\rangle\right)$
and $K_{0,S}$, $X'_Y$ are replaced by $K_{\beta,S}$, $\alpha_\beta X'_Y$.
Consider, in Lemma~\ref{lem:mixed-word-formula}(b), a single insertion, $r=1$, with $S_1=S$, $\Lambda_1=\Lambda_Y$, $X_1=X'_Y$ and $s_1=\frac18$ in the plane $j=j_1$,
and a two-cochain whose electric part $\beta_j$ in this plane is $\beta$;
then $S^{(j)}=S_j=S$, and as $k^{(\frac18)}_{X'_Y}=k^0_Y$, the identity~\eqref{eq:amplification-step-operator} for this $j$ is the identity just obtained,
its operator $K_j$ being the operator with the kernel displayed above.
So this kernel is the kernel, in the sense of Proposition~\ref{prop:word-formula} and before the projection $P_{\mathrm{G}}$,
of the step operator that the word formula attaches to this insertion.

\emph{The norm.}
By (a), the exponent $\delta_0+\delta_1$ lies in $\mathfrak{D}(S)$ and $\Upsilon_Y$ in $\End_{z(\delta_0+\delta_1)}(\Lambda_Y)$, and $f'\in\mathfrak{A}(S)$;
so Lemma~\ref{lem:sharp-norm}, applied to $X'_Y=\exp\left(i\langle\delta_0+\delta_1,A\rangle\right)\cdot\Upsilon_Y\cdot f'(E_S)$, gives $\left\|X'_Y\right\|_\#=\left\|\Upsilon_Y\right\|\cdot\sup|f'|$.
As $\xi$ ranges over $\R^S$, so does $x:=\xi+\delta_0$, and
\[
\sup_{x\in\R^S}\left(-\tfrac14|x|^2+\tfrac14\langle x,\delta_1-\delta_0\rangle\right)=\tfrac{1}{16}|\delta_1-\delta_0|^2,
\]
the supremum being attained at $x=\frac12(\delta_1-\delta_0)$.
Hence
\[
\left\|X'_Y\right\|_\#=\left\|\Upsilon_Y\right\|\cdot\exp\left(\tfrac{1}{16}|\delta_1-\delta_0|^2+\tfrac18\left(|\delta_0|^2+|\delta_1|^2\right)\right).
\]

\emph{The bound.}
Since $Y=G'_{X'_Y}$, Proposition~\ref{prop:bounded-elements}(d') shows that $Y$ is bounded by $2^{\frac{\#|S|}{2}}\cdot\left\|X'_Y\right\|_\#$ with the box $B^{\#}(c,S,\Lambda_Y)$.

(c) Let $\varepsilon>0$.
As in (b), the definition of $X_{Y,\varepsilon}$ reads $X_{Y,\varepsilon}=\left(\exp\left(i\langle\delta_0+\delta_1,A\rangle\right)\cdot\Upsilon_Y\right)\cdot f_\varepsilon(E_S)$,
where the function $f_\varepsilon$ on $\R^S$ is
\[
f_\varepsilon(\xi):=\exp\left(-\varepsilon|\xi+\delta_0+\delta_1|^2-\varepsilon|\xi|^2+\tfrac12\langle\xi+\delta_0,\delta_1-\delta_0\rangle+\tfrac14\left(|\delta_0|^2+|\delta_1|^2\right)\right).
\]
The quadratic part of the exponent is $-2\varepsilon|\xi|^2$.
So $f_\varepsilon\in\mathfrak{A}(S)$, with the quadratic form $\mathfrak{Q}=2\varepsilon|\cdot|^2$ (and $f_\varepsilon=1$ if $S=\emptyset$),
and $X_{Y,\varepsilon}\in\mathfrak{W}(S,\Lambda_Y)$ for the same reasons as $X'_Y$.

\emph{The kernel.}
The relation~\eqref{eq:setting-weyl}, read as in (b) but with the function $\xi\mapsto\exp\left(-\tfrac14|\xi|^2\right)$, gives
\begin{multline*}
\exp\left(-\tfrac14|E_S|^2\right)X_{Y,\varepsilon}\exp\left(-\tfrac14|E_S|^2\right)\\
=\exp\left(i\langle\delta_0+\delta_1,A\rangle\right)\cdot\exp\left(-\tfrac14|\widetilde E+\delta_1|^2-\tfrac14|\widetilde E-\delta_0|^2\right)\cdot f_\varepsilon(E_S)\cdot\Upsilon_Y,
\end{multline*}
where
\[
f_\varepsilon(E_S)=\exp\left(-\varepsilon|\widetilde E+\delta_1|^2-\varepsilon|\widetilde E-\delta_0|^2+\tfrac12\langle\widetilde E,\delta_1-\delta_0\rangle+\tfrac14\left(|\delta_0|^2+|\delta_1|^2\right)\right).
\]
Since for $x\in\R^S$
\[
\tfrac14|x+\delta_1|^2+\tfrac14|x-\delta_0|^2=\tfrac12|x|^2+\tfrac12\langle x,\delta_1-\delta_0\rangle+\tfrac14\left(|\delta_0|^2+|\delta_1|^2\right),
\]
the right side equals
\[
\exp\left(i\langle\delta_0+\delta_1,A\rangle\right)\cdot\exp\left(-\tfrac12|\widetilde E|^2-\varepsilon|\widetilde E+\delta_1|^2-\varepsilon|\widetilde E-\delta_0|^2\right)\cdot\Upsilon_Y .
\]
On the other hand, $P_\varepsilon=\exp\left(-\varepsilon|\widetilde E-\delta_0|^2\right)$, and by the same relation
\[
P_\varepsilon\cdot\exp\left(i\langle\delta_0+\delta_1,A\rangle\right)=\exp\left(i\langle\delta_0+\delta_1,A\rangle\right)\cdot\exp\left(-\varepsilon|\widetilde E+\delta_1|^2\right).
\]
The functions of $E_S$ commute with one another and with $\Upsilon_Y$; hence
\begin{equation}\label{eq:amplification-kernel-epsilon}
\exp\left(-\tfrac14|E_S|^2\right)X_{Y,\varepsilon}\exp\left(-\tfrac14|E_S|^2\right)=P_\varepsilon\,k^0_Y\,P_\varepsilon .
\end{equation}
On $\mathcal{H}^{\mathrm{l}}$ one has, by Definition~\ref{def:slice-operators},
\[
K^{\frac12}=(2\pi/q)^{\frac{|\Lambda_{\mathrm{s}}|}{2}}\cdot P_{\mathrm{G}}\cdot\exp\left(-\tfrac14\sum_{b\notin S}E_b^2-\tfrac14|E_S|^2\right),
\]
and $(P_\varepsilon)_\Xi=\exp\left(-\varepsilon|E_S|^2\right)$ commutes with $P_{\mathrm{G}}$ and with the functions of the $E_b$, $b\notin S$.
The argument given in (b) for $k_Y$, with~\eqref{eq:amplification-kernel-epsilon} in the place of~\eqref{eq:amplification-kernel-prime}, and the identity~\eqref{eq:amplification-kernel-link} give
\[
\begin{aligned}
K^{\frac12}\left(X_{Y,\varepsilon}\right)_\Xi K^{\frac12}
&=(2\pi/q)^{|\Lambda_{\mathrm{s}}|}\cdot P_{\mathrm{G}}\cdot\exp\left(-\tfrac12\sum_{b\notin S}E_b^2\right)\cdot(P_\varepsilon)_\Xi\left(k^0_Y\right)_\Xi(P_\varepsilon)_\Xi\\
&=(P_\varepsilon)_\Xi\,k_Y\,(P_\varepsilon)_\Xi .
\end{aligned}
\]

\emph{The difference $Y-G_{X_{Y,\varepsilon}}$.}
The operator $P_\varepsilon$ commutes with $\exp\left(-\tfrac18|E_S|^2\right)$, so that~\eqref{eq:amplification-kernel-epsilon} and~\eqref{eq:amplification-kernel-prime} give
\[
\exp\left(-\tfrac14|E_S|^2\right)X_{Y,\varepsilon}\exp\left(-\tfrac14|E_S|^2\right)=\exp\left(-\tfrac18|E_S|^2\right)\left(P_\varepsilon X'_YP_\varepsilon\right)\exp\left(-\tfrac18|E_S|^2\right).
\]
As $P_\varepsilon$ is a generator of type~\eqref{it:gaussian-generators} of $\mathfrak{W}(S,\Lambda_Y)$, the element $P_\varepsilon X'_YP_\varepsilon$ lies in $\mathfrak{W}(S,\Lambda_Y)$,
and the continuous kernels of the two sides coincide: $k^{(\frac14)}_{X_{Y,\varepsilon}}=k^{(\frac18)}_{P_\varepsilon X'_YP_\varepsilon}$ (Lemma~\ref{lem:regularised-kernels}).
By Definition~\ref{def:two-insertions}, $G_{X_{Y,\varepsilon}}=G'_{P_\varepsilon X'_YP_\varepsilon}$.
The map $X\mapsto G'_X$ is linear (Definition~\ref{def:two-insertions}), and $Y=G'_{X'_Y}$ by (b); therefore
\[
Y-G_{X_{Y,\varepsilon}}=G'_{X'_Y-P_\varepsilon X'_YP_\varepsilon},
\]
and by Proposition~\ref{prop:bounded-elements}(d') this integrand is bounded by
$2^{\frac{\#|S|}{2}}\cdot\left\|X'_Y-P_\varepsilon X'_YP_\varepsilon\right\|_\#=2^{\frac{\#|S|}{2}}\cdot\xi_Y(\varepsilon)$ with the box $B^{\#}(c,S,\Lambda_Y)$.

\emph{The number $\xi_Y(\varepsilon)$.}
By~\eqref{eq:setting-weyl},
\[
P_\varepsilon X'_YP_\varepsilon=\left(\exp\left(i\langle\delta_0+\delta_1,A\rangle\right)\cdot\Upsilon_Y\right)\cdot f'(E_S)\cdot\exp\left(-\varepsilon|\widetilde E+\delta_1|^2-\varepsilon|\widetilde E-\delta_0|^2\right).
\]
The difference is therefore a single term,
\begin{multline*}
X'_Y-P_\varepsilon X'_YP_\varepsilon\\
=\left(\exp\left(i\langle\delta_0+\delta_1,A\rangle\right)\cdot\Upsilon_Y\right)\cdot f'(E_S)\cdot\left(1-\exp\left(-\varepsilon|\widetilde E+\delta_1|^2-\varepsilon|\widetilde E-\delta_0|^2\right)\right),
\end{multline*}
whose function of $E_S$ is non-negative.
With Lemma~\ref{lem:sharp-norm}, in the variable $x=\xi+\delta_0$, one obtains
\begin{multline*}
\xi_Y(\varepsilon)=\left\|\Upsilon_Y\right\|\cdot\sup_{x\in\R^S}\Bigl[\exp\left(-\tfrac14|x|^2+\tfrac14\langle x,\delta_1-\delta_0\rangle+\tfrac18\left(|\delta_0|^2+|\delta_1|^2\right)\right)\\
\cdot\left(1-\exp\left(-\varepsilon|x+\delta_1|^2-\varepsilon|x-\delta_0|^2\right)\right)\Bigr].
\end{multline*}
As $1-\exp(-u)\le u$ for $u\ge0$, one obtains $\xi_Y(\varepsilon)\le C_Y\cdot\varepsilon$ with
\begin{multline*}
C_Y:=\left\|\Upsilon_Y\right\|\cdot\sup_{x\in\R^S}\Bigl[\exp\left(-\tfrac14|x|^2+\tfrac14\langle x,\delta_1-\delta_0\rangle+\tfrac18\left(|\delta_0|^2+|\delta_1|^2\right)\right)\\
\cdot\left(|x+\delta_1|^2+|x-\delta_0|^2\right)\Bigr],
\end{multline*}
which is finite, a Gaussian times a polynomial being bounded.

\emph{The norm of $X_{Y,\varepsilon}$.}
In the variable $x=\xi+\delta_0$ the exponent of $f_\varepsilon$ is
\[
\begin{aligned}
&-\varepsilon|x+\delta_1|^2-\varepsilon|x-\delta_0|^2+\tfrac12\langle x,\delta_1-\delta_0\rangle+\tfrac14\left(|\delta_0|^2+|\delta_1|^2\right)\\
&\qquad=-2\varepsilon|x|^2+\left(\tfrac12-2\varepsilon\right)\langle x,\delta_1-\delta_0\rangle+\left(\tfrac14-\varepsilon\right)\left(|\delta_0|^2+|\delta_1|^2\right).
\end{aligned}
\]
Its supremum over $x\in\R^S$ is attained at $x=\frac{1}{4\varepsilon}\left(\frac12-2\varepsilon\right)(\delta_1-\delta_0)$, and Lemma~\ref{lem:sharp-norm} gives
\[
\left\|X_{Y,\varepsilon}\right\|_\#=\left\|\Upsilon_Y\right\|\cdot\exp\left(\frac{\left(\frac12-2\varepsilon\right)^2|\delta_1-\delta_0|^2}{8\varepsilon}+\left(\tfrac14-\varepsilon\right)\left(|\delta_0|^2+|\delta_1|^2\right)\right).
\]
If $S$ is empty, then $\delta_0=\delta_1=0$, $P_\varepsilon=1$, $\xi_Y(\varepsilon)=0$ and $X_{Y,\varepsilon}=X'_Y=\Upsilon_Y$.

(d) By the rule~\eqref{eq:setting-reflection-coefficients} for coefficient functions,
\[
\Theta_cY=\exp\left(-i\langle\delta,\theta_c^{*}\hat A\rangle\right)\cdot\Theta_c(\Psi).
\]
The reflection $\theta_c$ exchanges the links of the planes $\{x_1=j_0\}$ and $\{x_1=j_0+1\}$ and reverses the links between them;
by~\eqref{eq:setting-slice-reflections},
\[
\langle\delta,\theta_c^{*}\hat A\rangle=\langle\delta_0,\hat A_{j_0+1}\rangle+\langle\delta_1,\hat A_{j_0}\rangle-\langle\delta_\perp,\hat a_{j_0}\rangle .
\]
So the exponent of $\Theta_cY$ is $(-\delta_1,-\delta_0,\delta_\perp)$, and its set $S$ is that of $Y$.
By~\eqref{eq:setting-reflection-fields}, $\Theta_c$ is antilinear, reverses products, and maps each field at a site $x$ to a linear combination of fields at $\theta_cx$.
As $\theta_c$ exchanges the thin layers $\mathbb{X}_{j_0,6}$ and $\mathbb{X}_{j_0+1,1}$ and leaves the coordinates $x_2$, $x_3$ unchanged,
$\Theta_c(\Psi)$ is a finite combination of monomials in the fields of these two thin layers over the same sites of $\Z^2$ as $\Psi$.
Moreover $\Theta_cY=Y^{\theta_c}$ by~\eqref{eq:setting-group-action}, and $\mathcal{P}$ is stable under the action of $W$ (Definition~\ref{def:polynomials}).
Thus $\Theta_cY\in\mathcal{P}$, and by the uniqueness of the representation in Definition~\ref{def:two-layer}
it is a finite combination of two-layer monomials with the exponent $(-\delta_1,-\delta_0,\delta_\perp)$, for each of which $\Lambda_Y$ is a set of sites in play.
In what follows, $\Upsilon_{\Theta_cY}$, $k^0_{\Theta_cY}$, $k_{\Theta_cY}$, $X'_{\Theta_cY}$ and $X_{\Theta_cY,\varepsilon}$ are those of the linear extension in Definition~\ref{def:two-layer};
all the monomials of $\Theta_cY$ have the set $S$, and hence the operator $P_\varepsilon$.

\emph{The Fock operator.}
By Lemma~\ref{lem:plane-reflection}, $\Theta_c\left(\Ins_c[\Upsilon]\right)=\Ins_c\left[\Upsilon^{*}\right]$ for every constant matrix $\Upsilon$;
in particular $\Theta_c\left(\Ins_c[\omega]\right)=\Ins_c[\omega]$, as $\omega$ is real and diagonal.
The element $\Ins_c[\omega]$ is even and commutes with every element of the Grassmann algebra.
Applying $\Theta_c$, which reverses products, to the identity $\Psi\cdot\Ins_c[\omega]=\Ins_c\left[\omega^{\frac12}\,\Upsilon_Y\,\omega^{\frac12}\right]$, one finds
\[
\Theta_c(\Psi)\cdot\Ins_c[\omega]=\Ins_c[\omega]\cdot\Theta_c(\Psi)=\Theta_c\left(\Psi\cdot\Ins_c[\omega]\right)=\Ins_c\left[\omega^{\frac12}\,\Upsilon_Y^{*}\,\omega^{\frac12}\right].
\]
By linearity the left side is $\Ins_c\left[\omega^{\frac12}\,\Upsilon_{\Theta_cY}\,\omega^{\frac12}\right]$.
Both sides lie in $\Gr(\mathbb{X}')$, where an element $\Ins_c[\Upsilon]$ determines the matrix $\Upsilon$;
since $\omega$ is invertible, it follows that $\Upsilon_{\Theta_cY}=\Upsilon_Y^{*}$.

\emph{The operators $k^0$.}
By~\eqref{eq:setting-weyl},
\[
\begin{aligned}
\left(k^0_Y\right)^{*}&=\Upsilon_Y^{*}\cdot\exp\left(-\tfrac12|E_S+\delta_0|^2\right)\cdot\exp\left(-i\langle\delta_0+\delta_1,A\rangle\right)\\
&=\exp\left(-i\langle\delta_0+\delta_1,A\rangle\right)\cdot\exp\left(-\tfrac12|E_S-\delta_1|^2\right)\cdot\Upsilon_Y^{*}.
\end{aligned}
\]
This is the operator of Definition~\ref{def:two-layer} for the exponent $(-\delta_1,-\delta_0,\delta_\perp)$ and the Fock operator $\Upsilon_Y^{*}=\Upsilon_{\Theta_cY}$;
as that operator depends linearly on the Fock operator, $\left(k^0_Y\right)^{*}=k^0_{\Theta_cY}$ on $\mathcal{H}_S$.

\emph{The elements $X'$ and $X_\varepsilon$.}
Applied to the monomials of $\Theta_cY$ and summed, \eqref{eq:amplification-kernel-prime} and~\eqref{eq:amplification-kernel-epsilon} give
\[
\begin{aligned}
\exp\left(-\tfrac18|E_S|^2\right)X'_{\Theta_cY}\exp\left(-\tfrac18|E_S|^2\right)&=k^0_{\Theta_cY},\\
\exp\left(-\tfrac14|E_S|^2\right)X_{\Theta_cY,\varepsilon}\exp\left(-\tfrac14|E_S|^2\right)&=P_\varepsilon\,k^0_{\Theta_cY}\,P_\varepsilon .
\end{aligned}
\]
Taking adjoints in~\eqref{eq:amplification-kernel-prime} and~\eqref{eq:amplification-kernel-epsilon}, one sees that $\left(X'_Y\right)^{*}$ and $\left(X_{Y,\varepsilon}\right)^{*}$ satisfy the same two identities,
because $\left(k^0_Y\right)^{*}=k^0_{\Theta_cY}$ and $P_\varepsilon$ is self-adjoint.
For $s>0$ the operator $\exp\left(-s|E_S|^2\right)$ is bounded, self-adjoint and injective, and hence has a dense range;
a bounded operator that vanishes after multiplication by $\exp\left(-s|E_S|^2\right)$ on both sides is therefore zero.
It follows that $X'_{\Theta_cY}=\left(X'_Y\right)^{*}$ and $X_{\Theta_cY,\varepsilon}=\left(X_{Y,\varepsilon}\right)^{*}$.

\emph{The operators on $\mathcal{H}^{\mathrm{l}}$.}
By~\eqref{eq:amplification-kernel-link}, applied to $Y$ and to the monomials of $\Theta_cY$, and since $X\mapsto X_\Xi$ is a $*$-representation,
the adjoint of the operator $k^0_Y$ on $\mathcal{H}^{\mathrm{l}}$ is
$\exp\left(-\tfrac12\sum_{b\notin S}E_b^2\right)\cdot\left(k^0_{\Theta_cY}\right)_\Xi=k^0_{\Theta_cY}$;
so $k^0_{\Theta_cY}=\left(k^0_Y\right)^{*}$ holds on $\mathcal{H}^{\mathrm{l}}$ as well, and both operators commute with $P_{\mathrm{G}}$.
For the operators $k_Y$ this gives
\[
k_Y^{*}=(2\pi/q)^{|\Lambda_{\mathrm{s}}|}\cdot P_{\mathrm{G}}\cdot\left(k^0_Y\right)^{*}=(2\pi/q)^{|\Lambda_{\mathrm{s}}|}\cdot k^0_{\Theta_cY}\cdot P_{\mathrm{G}}=k_{\Theta_cY}.
\]
\end{proof}

\begin{remark}\label{rem:two-layer-insertions}
\begin{enumerate}
\item[(1)]
\emph{An unbounded operator.}
Formally, $K^{-\frac12}k_YK^{-\frac12}$ is the operator obtained by putting $\varepsilon=0$ in the definition of $X_{Y,\varepsilon}$,
\[
X_{Y,0}=\exp\left(i\langle\delta_0+\delta_1,A\rangle\right)\cdot\exp\left(\tfrac12\langle\widetilde E,\delta_1-\delta_0\rangle+\tfrac14\left(|\delta_0|^2+|\delta_1|^2\right)\right)\cdot\Upsilon_Y ;
\]
the computation in the proof of Proposition~\ref{prop:two-layer-insertions}(c) gives $\exp\left(-\tfrac14|E_S|^2\right)X_{Y,0}\exp\left(-\tfrac14|E_S|^2\right)=k^0_Y$.
The operator $X_{Y,0}$ contains the factor $\exp\left(\tfrac12\langle\widetilde E,\delta_1-\delta_0\rangle\right)$ and is unbounded unless $\delta_1=\delta_0$.
H\"older's inequality with insertions, Proposition~\ref{prop:holder-insertions}, applies to $k_Y$ all the same, through Proposition~\ref{prop:two-layer-insertions}(b):
there $k_Y=K_{0,S}^{\frac12}\left(X'_Y\right)_\Xi K_{0,S}^{\frac12}$ with $X'_Y$ bounded.
The kernel $g_{\frac12}$ of $\exp\left(-\tfrac14E_b^2\right)$ is at most $2^{\frac12}$ times the kernel $g_1$ of $\exp\left(-\tfrac12E_b^2\right)$,
so that the kernel of $K_{0,S}$ is pointwise at most $2^{\frac{\#|S|}{2}}$ times that of $K$;
as the measure is positive, the comparison operators $K_{0,S}$ cost at most the factor $2^{\frac{\#|S|}{2}}$ per step in~\eqref{eq:amplification-trace-power}.
\item[(2)]
\emph{The order of the limits.}
Since $\frac{1}{8\varepsilon}\left(\frac12-2\varepsilon\right)^2=\frac{1}{32\varepsilon}-\frac14+\frac{\varepsilon}{2}$,
the formula of Proposition~\ref{prop:two-layer-insertions}(c) expresses $\left\|X_{Y,\varepsilon}\right\|_\#$ as the product of $\exp\left(\frac{|\delta_1-\delta_0|^2}{32\varepsilon}\right)$
and a factor that tends to $\left\|\Upsilon_Y\right\|\cdot\exp\left(\frac12\langle\delta_0,\delta_1\rangle\right)$ as $\varepsilon\to0$.
Because of this growth, the limits in Theorem~\ref{thm:two-layer-principle}(b) are taken in a fixed order:
first $i\to\infty$, then $t\to0$ with $\varepsilon$ fixed, and finally $\varepsilon\to0$.
\item[(3)]
\emph{The condition $\rho<1$.}
At a coordinate link plane the crossing bonds are of type $3$, and the modes of the groups $2$ and $4$ have the coefficient $\frac{w_3}{2}(1-\rho)$ (Definition~\ref{def:modes}).
If $\Psi=\pm\Theta_c(\zeta_{J_0})\,\zeta_{I_0}$, the norm $\left\|\Upsilon_Y\right\|=\left(\omega_{I_0}\cdot\omega_{J_0}\right)^{-\frac12}$ of Proposition~\ref{prop:two-layer-insertions}(a)
therefore contains the factor $\left(\frac{w_3}{2}(1-\rho)\right)^{-\frac12}$ for every mode of these two groups in $I_0$ and in $J_0$;
these factors are finite because $\rho<1$.
\end{enumerate}
\end{remark}

\subsection{The two-layer principle}\label{sec:two-layer-principle}

Throughout this subsection $\mu\in\{1,2,3\}$ and $c\in\Z+\frac12$,
and $(\mu,c)$ denotes the coordinate link plane $\{x_\mu=c\}$.
As in Definition~\ref{def:other-directions},
$g_\mu\in\Ocub$ is the cyclic permutation matrix with $g_\mu e_1=e_\mu$, so that $g_1=1$;
it is an even element of $W$.
The same definition gives $\Theta_{\mu,c}=\Theta_{(x_\mu,c)}$ and $\mathcal{P}_1(\mu,c)=\left\{Y^{g_\mu}:Y\in\mathcal{P}_1(c)\right\}$,
and it takes the objects attached to the plane $(\mu,c)$ to be the images under $g_\mu$ of the corresponding objects at $(1,c)$.
Thus, for $X\in\mathfrak{W}(S,\Lambda)$ (Definition~\ref{def:mid-plane})
with insertion $G^{(1,c)}_X$ at the plane $\{x_1=c\}$ (Definition~\ref{def:insertions}),
the insertion of $X$ at $(\mu,c)$ is
\[
G_X:=\left(G^{(1,c)}_X\right)^{g_\mu}.
\]
The monomials of $\mathcal{P}_1(\mu,c)$ are the elements $Y^{g_\mu}$ with $Y$ a monomial of $\mathcal{P}_1(c)$,
and the sets $S$, $\Lambda_Y$, the elements $X_{Y,\varepsilon}$ and the numbers $\xi_Y(\varepsilon)$ attached to $Y^{g_\mu}$
are those of $Y$ (Definition~\ref{def:two-layer}).

The reflections in the planes $\{x_\mu=c\}$ and $\{x_1=c\}$ are conjugate,
\[
\theta_{x_\mu,c}=g_\mu\,\theta_{x_1,c}\,g_\mu^{-1},
\]
since both maps send $x$ to $x-2(x_\mu-c)e_\mu$.
Together with the rule \eqref{eq:setting-group-action} for the action of $W$ on even integrands this gives,
for every even integrand $H$,
\begin{equation}\label{eq:amplification-directions}
\Theta_{\mu,c}\left(H^{g_\mu}\right)
=H^{\theta_{x_\mu,c}g_\mu}
=H^{g_\mu\theta_{x_1,c}}
=\left(\Theta_{1,c}H\right)^{g_\mu}.
\end{equation}
As $g_\mu$ is even, the map $H\mapsto H^{g_\mu}$ is linear.

\begin{theorem}[The two-layer principle]\label{thm:two-layer-principle}
Let $0<\rho<1$, $w_0,\dots,w_3>0$, $M\in\R$ and $q>0$.
Let the sequence of sides $\mathfrak{n}$ satisfy \textup{(Div)}, and let $\mathfrak{d}$ be data for $\mathfrak{n}$.
Let $\mu^t_i$ be the perturbed functionals of Definition~\ref{def:data},
$\gamma_{t,i}(k)$ the amplification factors of Definition~\ref{def:amplification-factor},
and $k(\mathfrak{d},r)$, $i_0(\mathfrak{d},r)$ the integers of Lemma~\ref{lem:coarse-structures}.
Let $(\mu,c)$ be a coordinate link plane.
\begin{enumerate}
\item[(a)]
Let $X\in\mathfrak{W}(S,\Lambda)$, with $S$ and $\Lambda$ as in Definition~\ref{def:mid-plane},
and let $G_X=\left(G^{(1,c)}_X\right)^{g_\mu}$ be its insertion at $(\mu,c)$.
Let $r_0$ be the radius of the box $B(c,S,\Lambda)$ of Definition~\ref{def:insertions},
and put $k:=k(\mathfrak{d},2r_0)$.
For all $t\ge0$ and all $i\ge i_0(\mathfrak{d},2r_0)$,
\[
\left|\mu^t_i(G_X)\right|\le\gamma_{t,i}(k)\cdot\|X\|_\#
\]
and
\[
\left|\mu^t_i\left(\Theta_{\mu,c}G_X\right)-\overline{\mu^t_i(G_X)}\right|
\le4\left(\gamma_{t,i}(k)^2-1\right)^{\frac12}\cdot\|X\|_\#.
\]
\item[(b)]
Let $Y\neq1$ be a monomial of $\mathcal{P}_1(c)$,
with the set $S$, the sites in play $\Lambda_Y$, the elements $X_{Y,\varepsilon}$ and the numbers $\xi_Y(\varepsilon)$
of Definition~\ref{def:two-layer};
let $r_0$ be the radius of $B(c,S,\Lambda_Y)$, and put $k:=k(\mathfrak{d},2r_0)$.
These objects of $Y$ are also those of the monomial $Y^{g_\mu}$ of $\mathcal{P}_1(\mu,c)$.
For all $\varepsilon>0$, $t\ge0$ and $i\ge i_0(\mathfrak{d},2r_0)$,
\[
\left|\mu^t_i\left(\Theta_{\mu,c}Y^{g_\mu}\right)-\overline{\mu^t_i\left(Y^{g_\mu}\right)}\right|
\le2\cdot2^{\frac{\#|S|}{2}}\cdot\xi_Y(\varepsilon)\cdot\gamma_{t,i}(k)
+4\left(\gamma_{t,i}(k)^2-1\right)^{\frac12}\cdot\left\|X_{Y,\varepsilon}\right\|_\#.
\]
For $Y=1$ the left side vanishes.
Consequently, for every $Y'\in\mathcal{P}_1(\mu,c)$, a Grassmann-even element that may contain fermion fields,
\[
\lim_{t\to0}\,\limsup_{i\to\infty}\left|\mu^t_i\left(\Theta_{\mu,c}Y'\right)-\overline{\mu^t_i(Y')}\right|=0;
\]
in particular
$\lim_{t\to0}\limsup_{i\to\infty}\left|\Re\mu^t_i(Y')\right|=0$
if $\Theta_{\mu,c}Y'=-Y'$ (antisymmetric elements),
and
$\lim_{t\to0}\limsup_{i\to\infty}\left|\Im\mu^t_i(Y')\right|=0$
if $\Theta_{\mu,c}Y'=Y'$ (symmetric elements).
In the bound the limits are taken in the order
$i\to\infty$, then $t\to0$ at fixed $\varepsilon$, then $\varepsilon\to0$.
\end{enumerate}
\end{theorem}

\begin{proof}[Proof of Theorem~\ref{thm:two-layer-principle}(a)]
Fix $t\ge0$ and $i\ge i_0(\mathfrak{d},2r_0)$, and write $\gamma:=\gamma_{t,i}(k)$.
For every $X'\in\mathfrak{W}(S,\Lambda)$ the insertion $G_{X'}$ at $(\mu,c)$
is bounded by $\|X'\|_\#$ with the box $g_\mu B^{\#}(c,S,\Lambda)$
(Proposition~\ref{prop:bounded-elements}(d));
this box is the same for all $X'$, and its radius is $2r_0$.
Theorem~\ref{thm:amplification}, applied with $r=2r_0$, shows therefore that the functional
\[
\mathfrak{o}^t_i:\mathfrak{W}(S,\Lambda)\to\C,\qquad
\mathfrak{o}^t_i(X'):=\mu^t_i\left(G_{X'}\right)=\mu^t_i\left(\left(G^{(1,c)}_{X'}\right)^{g_\mu}\right),
\]
satisfies
\[
\left|\mathfrak{o}^t_i(X')\right|\le\gamma\cdot\|X'\|_\#
\qquad\text{for all }X'\in\mathfrak{W}(S,\Lambda).
\]
For $X'=X$ this is the first inequality.
The functional $\mathfrak{o}^t_i$ is linear,
because $X'\mapsto G^{(1,c)}_{X'}$ is linear (Proposition~\ref{prop:insertions}(a)),
$H\mapsto H^{g_\mu}$ is linear and $\mu^t_i$ is linear.
It is unital:
$G^{(1,c)}_1=1$ by the same proposition, so that $\mathfrak{o}^t_i(1)=\mu^t_i(1)=1$.
In particular $1=\mathfrak{o}^t_i(1)\le\gamma\cdot\|1\|_\#=\gamma$, and $\gamma^2-1\ge0$.

\emph{Self-adjoint elements.}
Let $X_0\in\mathfrak{W}(S,\Lambda)$ be self-adjoint with $\|X_0\|_\#\le1$.
Then
\[
\left|\Im\mathfrak{o}^t_i(X_0)\right|\le\left(\gamma^2-1\right)^{\frac12}.
\]
Indeed, $\|\cdot\|_\#$ is a $C^*$-norm (Lemma~\ref{lem:sharp-norm}).
For real $s$, the $C^*$-property and the triangle inequality give
\[
\left\|X_0+is\right\|_\#^2
=\left\|\left(X_0+is\right)^*\left(X_0+is\right)\right\|_\#
=\left\|X_0^2+s^2\right\|_\#
\le1+s^2.
\]
Write $\mathfrak{o}^t_i(X_0)=u+iv$ with real $u$ and $v$.
Since $\mathfrak{o}^t_i(X_0+is)=u+i(v+s)$,
\[
(v+s)^2\le\left|\mathfrak{o}^t_i(X_0+is)\right|^2\le\gamma^2\cdot\left(1+s^2\right)
\qquad\text{for all }s\in\R.
\]
If $\gamma>1$, choose $s:=v/\left(\gamma^2-1\right)$.
Then $v+s=\gamma^2\cdot v/\left(\gamma^2-1\right)$,
and after multiplication by $\left(\gamma^2-1\right)^2\cdot\gamma^{-2}$ the inequality becomes
\[
v^2\cdot\gamma^2\le\left(\gamma^2-1\right)^2+v^2,
\qquad\text{that is,}\qquad
v^2\cdot\left(\gamma^2-1\right)\le\left(\gamma^2-1\right)^2;
\]
hence $v^2\le\gamma^2-1$.
If $\gamma=1$, the inequality reads $2vs+v^2\le1$ for every real $s$, and this forces $v=0$.

\emph{Arbitrary elements.}
Let $X\in\mathfrak{W}(S,\Lambda)$.
If $\|X\|_\#=0$, then $X=0$, and both sides of the second inequality vanish.
Otherwise put
\[
X_1:=\frac12\left(X+X^*\right),\qquad X_2:=\frac1{2i}\left(X-X^*\right).
\]
These are self-adjoint elements of $\mathfrak{W}(S,\Lambda)$ with $\|X_j\|_\#\le\|X\|_\#$,
and $X=X_1+iX_2$, $X^*=X_1-iX_2$.
Consequently
\begin{align*}
\mathfrak{o}^t_i(X^*)-\overline{\mathfrak{o}^t_i(X)}
&=\left[\mathfrak{o}^t_i(X_1)-\overline{\mathfrak{o}^t_i(X_1)}\right]
-i\left[\mathfrak{o}^t_i(X_2)-\overline{\mathfrak{o}^t_i(X_2)}\right]\\
&=2i\Im\mathfrak{o}^t_i(X_1)+2\Im\mathfrak{o}^t_i(X_2).
\end{align*}
On the other hand, $G^{(1,c)}_{X^*}=\Theta_{1,c}G^{(1,c)}_X$ by Proposition~\ref{prop:insertions}(b).
Both sides of this identity are elements of $\mathcal{G}^{\mathrm{ev}}\left(B^{\#}(c,S,\Lambda)\right)$
that are equal in the sense of Lemma~\ref{lem:equal-integrands},
and the image of this box under $g_\mu$ embeds in $D(n_i)$ (Lemma~\ref{lem:coarse-structures}\eqref{it:coarse-tori});
so the images of the two sides under $g_\mu$ have the same value under $\mu^t_i$.
With \eqref{eq:amplification-directions} it follows that
\[
\mathfrak{o}^t_i(X^*)
=\mu^t_i\left(\left(\Theta_{1,c}G^{(1,c)}_X\right)^{g_\mu}\right)
=\mu^t_i\left(\Theta_{\mu,c}G_X\right).
\]
The left side of the second inequality of (a) is therefore
$\left|2i\Im\mathfrak{o}^t_i(X_1)+2\Im\mathfrak{o}^t_i(X_2)\right|$,
and the bound for self-adjoint elements, applied to $X_1/\|X\|_\#$ and to $X_2/\|X\|_\#$,
shows that it is at most $4\left(\gamma^2-1\right)^{\frac12}\cdot\|X\|_\#$.
\end{proof}

\begin{proof}[Proof of Theorem~\ref{thm:two-layer-principle}(b)]
If $Y=1$, then $\Theta_{\mu,c}Y^{g_\mu}=1$ and $\mu^t_i(1)=1$, so the left side vanishes.
Let $Y\neq1$, and fix $\varepsilon>0$, $t\ge0$ and $i\ge i_0(\mathfrak{d},2r_0)$.
Put $X:=X_{Y,\varepsilon}$, an element of $\mathfrak{W}(S,\Lambda_Y)$.
The difference $Y-G^{(1,c)}_X$ is an even integrand, and
\[
Y^{g_\mu}=\left(Y-G^{(1,c)}_X\right)^{g_\mu}+G_X.
\]
Since $\Theta_{\mu,c}$ is additive, \eqref{eq:amplification-directions} gives
$\Theta_{\mu,c}Y^{g_\mu}=\left(\Theta_{1,c}\left(Y-G^{(1,c)}_X\right)\right)^{g_\mu}+\Theta_{\mu,c}G_X$,
and the linearity of $\mu^t_i$ yields
\begin{align*}
\mu^t_i\left(\Theta_{\mu,c}Y^{g_\mu}\right)-\overline{\mu^t_i\left(Y^{g_\mu}\right)}
&=\mu^t_i\left(\left(\Theta_{1,c}\left(Y-G^{(1,c)}_X\right)\right)^{g_\mu}\right)
+\left[\mu^t_i\left(\Theta_{\mu,c}G_X\right)-\overline{\mu^t_i(G_X)}\right]\\
&\quad-\overline{\mu^t_i\left(\left(Y-G^{(1,c)}_X\right)^{g_\mu}\right)}.
\end{align*}
By Proposition~\ref{prop:two-layer-insertions}(c),
$Y-G^{(1,c)}_X$ is bounded by $2^{\frac{\#|S|}{2}}\cdot\xi_Y(\varepsilon)$ with the box $B^{\#}(c,S,\Lambda_Y)$,
whose radius is $2r_0$.
By \eqref{eq:setting-group-action},
$\left(\Theta_{1,c}\left(Y-G^{(1,c)}_X\right)\right)^{g_\mu}=\left(Y-G^{(1,c)}_X\right)^{g_\mu\theta_{x_1,c}}$.
Proposition~\ref{prop:bounded-elements}(a) shows therefore that
$\left(Y-G^{(1,c)}_X\right)^{g_\mu}$ and $\left(\Theta_{1,c}\left(Y-G^{(1,c)}_X\right)\right)^{g_\mu}$
are bounded by the same constant,
with the images of $B^{\#}(c,S,\Lambda_Y)$ under $g_\mu$ and under $g_\mu\theta_{x_1,c}$,
which are again boxes of radius $2r_0$.
Theorem~\ref{thm:amplification} with $r=2r_0$ bounds the first and the third term on the right
by $2^{\frac{\#|S|}{2}}\cdot\xi_Y(\varepsilon)\cdot\gamma_{t,i}(k)$ each,
and part (a), applied to $X\in\mathfrak{W}(S,\Lambda_Y)$,
bounds the term in brackets by $4\left(\gamma_{t,i}(k)^2-1\right)^{\frac12}\cdot\left\|X_{Y,\varepsilon}\right\|_\#$.
This is the inequality of (b).

\emph{Limits.}
Neither $k=k(\mathfrak{d},2r_0)$ nor $i_0(\mathfrak{d},2r_0)$ depends on $\varepsilon$:
for every $\varepsilon>0$ the element $X_{Y,\varepsilon}$ lies in $\mathfrak{W}(S,\Lambda_Y)$,
and the box $B^{\#}(c,S,\Lambda_Y)$ does not change with $\varepsilon$.
As $\xi_Y(\varepsilon)\le C_Y\cdot\varepsilon$ (Proposition~\ref{prop:two-layer-insertions}(c)),
for $0\le t\le1$ and $i\ge i_0(\mathfrak{d},2r_0)$ the first term of the bound is at most
\[
2\cdot2^{\frac{\#|S|}{2}}\cdot C_Y\cdot\varepsilon\cdot
\sup\left\{\gamma_{t',i'}(k):0\le t'\le1,\ i'\ge i_0(\mathfrak{d},2r_0)\right\},
\]
with one $k$ and one $i_0$ for all $\varepsilon$;
the supremum is finite by \eqref{eq:amplification-factor-limit}.
Now fix $\varepsilon>0$.
By \eqref{eq:amplification-factor-limit}, $\lim_{t\to0}\limsup_{i\to\infty}\gamma_{t,i}(k)=1$,
hence also $\lim_{t\to0}\limsup_{i\to\infty}\left(\gamma_{t,i}(k)^2-1\right)^{\frac12}=0$,
while $\left\|X_{Y,\varepsilon}\right\|_\#$ is a finite number that depends neither on $t$ nor on $i$
(Proposition~\ref{prop:two-layer-insertions}(c)).
Taking in the inequality of (b) first the upper limit as $i\to\infty$ and then the upper limit as $t\to0$,
one obtains
\[
\limsup_{t\to0}\,\limsup_{i\to\infty}
\left|\mu^t_i\left(\Theta_{\mu,c}Y^{g_\mu}\right)-\overline{\mu^t_i\left(Y^{g_\mu}\right)}\right|
\le2\cdot2^{\frac{\#|S|}{2}}\cdot\xi_Y(\varepsilon).
\]
The left side does not depend on $\varepsilon$, and $\xi_Y(\varepsilon)\le C_Y\cdot\varepsilon$.
Letting $\varepsilon\to0$, one finds that the limit as $t\to0$ exists and equals $0$.
The regularisation is removed last because,
by the formula for $\left\|X_{Y,\varepsilon}\right\|_\#$ in Proposition~\ref{prop:two-layer-insertions}(c),
this norm is unbounded as $\varepsilon\to0$ whenever $\delta_1\neq\delta_0$,
$(\delta_0,\delta_1,\delta_\perp)$ being the exponent of $Y$ (Definition~\ref{def:two-layer}).

\emph{Combinations of monomials.}
An element $Y'\in\mathcal{P}_1(\mu,c)$ is a finite sum $Y'=\sum_az_a\cdot Y_a^{g_\mu}$
with $z_a\in\C$ and monomials $Y_a$ of $\mathcal{P}_1(c)$.
As $\Theta_{\mu,c}$ is antilinear and $\mu^t_i$ is linear,
\[
\mu^t_i\left(\Theta_{\mu,c}Y'\right)-\overline{\mu^t_i(Y')}
=\sum_a\overline{z_a}\cdot\left[\mu^t_i\left(\Theta_{\mu,c}Y_a^{g_\mu}\right)-\overline{\mu^t_i\left(Y_a^{g_\mu}\right)}\right].
\]
The terms with $Y_a=1$ vanish,
and each of the other terms tends to $0$ when the upper limit as $i\to\infty$ is followed by the limit $t\to0$;
so does the finite sum.
For the last two assertions it suffices to note that the left side equals $-2\Re\mu^t_i(Y')$ if $\Theta_{\mu,c}Y'=-Y'$,
and $2i\Im\mu^t_i(Y')$ if $\Theta_{\mu,c}Y'=Y'$.
\end{proof}

\begin{remark}[The range of the Wilson parameter]\label{rem:wilson-range}
Several statements of this section hold on a larger range of $\rho$ than $0<\rho<1$.
Proposition~\ref{prop:holder-insertions}, Proposition~\ref{prop:bounded-elements}(a)--(c),
Theorem~\ref{thm:amplification} and Proposition~\ref{prop:domination-tightness}
are stated for $-1\le\rho\le1$,
the closed interval on which all crossing coefficients $\omega_i$ of Definition~\ref{def:modes} are non-negative
(Proposition~\ref{prop:crossing}).
The parameter $\rho$ does not occur in Lemma~\ref{lem:coarse-structures}.
When $X$ has no Fock part,
that is, when $X$ lies in the subalgebra of $\mathfrak{W}$ generated by the elements of type~\eqref{it:gaussian-generators}
and the elements $\exp\left(i\langle\delta,A\rangle\right)$ of type~\eqref{it:exponential-generators} with $\Upsilon=1$
(Definition~\ref{def:mid-plane}),
the statements of this section on its insertion $G_X$ remain true for $-1\le\rho\le1$, with unchanged proofs;
excepted are Lemma~\ref{lem:injective-insertions} and the assertion of Proposition~\ref{prop:insertions}(c) that is made for $|\rho|<1$.

For every $\rho$ with $-1\le\rho\le1$ the formula of Definition~\ref{def:insertions} yields an integrand $G_X$ for every $X\in\mathfrak{W}(S,\Lambda)$.
Proposition~\ref{prop:insertions}, Proposition~\ref{prop:bounded-elements}(d) and Theorem~\ref{thm:two-layer-principle}(a)
hold for all these $\rho$, with the same proofs,
apart from the assertion of Proposition~\ref{prop:insertions}(c) that is made for $|\rho|<1$, which holds for $-1<\rho<1$.
In the formula of Definition~\ref{def:insertions} the kernel $k_X$ enters only through the matrices $\omega^{\frac12}\,k_X(A',A)\,u\,\omega^{\frac12}$,
where $u$ is the diagonal unitary $u(\hat a_j)$ of that definition.
The diagonal operators $\omega^{\frac12}$ and $u\,\omega^{\frac12}$ vanish on every basis vector $e_I$ with $\omega_I=0$,
that is, on every $e_I$ for which $I$ contains a mode with vanishing coefficient.
These matrices are therefore unchanged when $k_X(A',A)$ is replaced by its compression
to the subspace spanned by the basis vectors $e_I$ with $\omega_I>0$,
and at $\rho=\pm1$ the integrand $G_X$ is determined by these compressions.

The operator $\Upsilon_Y$ of Definition~\ref{def:two-layer}, in contrast, need not exist at $\rho=\pm1$.
At a coordinate link plane the crossing bonds are of type $3$;
at $\rho=1$ the coefficients $\frac{w_3}{2}(1-\rho)$ of the modes of the groups $2$ and $4$ vanish,
and at $\rho=-1$ the coefficients $\frac{w_3}{2}(1+\rho)$ of the modes of the groups $1$ and $3$.
Let $c=j_0+\frac12$ with $j_0\in\Z$,
and let $Y=\exp\left(i\langle\delta,\hat A\rangle\right)\Psi\in\mathcal{P}_1(c)$ be as in Proposition~\ref{prop:two-layer-insertions}(a),
with a product of modes $\Psi=\pm\Theta_c(\zeta_{J_0})\zeta_{I_0}$ in the place of a monomial:
$\delta_0$, $\delta_1$ and $\delta_\perp$ are formed from $\delta$ as in Definition~\ref{def:two-layer},
$\Lambda_Y$ is a finite set of sites of $\Z^2$ that contains the support of $\delta_\perp$
and the end points of the links $b$ of $\Z^2$ with $\delta_0(b)\neq0$ or $\delta_1(b)\neq0$,
and $I_0,J_0\subset\mathcal{J}^{\mathrm{l}}(\Lambda_Y)$.
Assume that $I_0\cup J_0$ contains a mode with vanishing coefficient, so that $\omega_{I_0}\cdot\omega_{J_0}=0$.
Elements of this kind exist at $\rho=\pm1$.
Indeed, let $i_1$ be a mode with $\omega_{i_1}=0$ of a bond over a site $y$,
and let $\delta'$ be the real one-cochain that vanishes off the link from $(j_0,y)$ to $(j_0+1,y)$
and takes there the value $q$ if $i_1$ belongs to the group $1$ or $2$, and the value $-q$ if it belongs to the group $3$ or $4$.
Then $d^{*}\delta'=q\cdot\mathfrak{q}(\Psi')$ for every monomial $\Psi'$ in the expansion of $\Theta_c(\zeta_{i_1})\zeta_{i_1}$:
both sides vanish off the sites $(j_0,y)$ and $(j_0+1,y)$,
where they take the values $-q$ and $q$, respectively, in the first case, and $q$ and $-q$ in the second,
by Lemma~\ref{lem:reflected-modes} and the definition of $\mathfrak{q}(i_1)$ in Definition~\ref{def:modes}.
Each $\exp\left(i\langle\delta',\hat A\rangle\right)\Psi'$ is therefore a two-layer monomial at $c$
(Definitions~\ref{def:polynomials} and~\ref{def:two-layer}),
and $\exp\left(i\langle\delta',\hat A\rangle\right)\Theta_c(\zeta_{i_1})\zeta_{i_1}$ is an element of the kind just described,
with $I_0=J_0=\{i_1\}$ and with $\{y\}$ in the place of $\Lambda_Y$.
For $Y$ as above, the products $\Theta_c(\zeta_J)\zeta_I$, $I,J\subset\mathcal{J}^{\mathrm{l}}(\Lambda_Y)$,
form a basis of the Grassmann algebra of the fields at the ends of the bonds over $\Lambda_Y$ that cross the plane $\{x_1=c\}$,
because the $\zeta_i$ and the $\Theta_c(\zeta_i)$ form bases of the spans of the generators at the upper and at the lower ends
(Definition~\ref{def:modes}).
Proposition~\ref{prop:crossing}, applied to these bonds, whose phases equal $1$ at $\hat a_{j_0}=0$,
shows that the left side of the equation that defines $\Upsilon_Y$ is
\[
\sum_{R\subset\mathcal{J}^{\mathrm{l}}(\Lambda_Y)}\omega_R\cdot\Psi\,\Theta_c(\zeta_R)\,\zeta_R .
\]
The term of $R$ vanishes if $R$ meets $I_0\cup J_0$,
and otherwise it equals $\pm\omega_R\cdot\Theta_c(\zeta_{J_0\cup R})\,\zeta_{I_0\cup R}$;
distinct sets $R$ give distinct basis elements.
On the left side the basis element $\Theta_c(\zeta_{J_0})\zeta_{I_0}$ thus has the coefficient $\pm1$, coming from $R=\emptyset$,
whereas its coefficient in $\Ins_c\left[\omega^{\frac12}\,\Upsilon\,\omega^{\frac12}\right]$
is $\left(\omega_{I_0}\cdot\omega_{J_0}\right)^{\frac12}\cdot\Upsilon_{I_0J_0}=0$, whatever the operator $\Upsilon$.
Hence the equation has no solution.
Since $Y$ lies in $\mathcal{P}_1(c)$, the unique representation of Definition~\ref{def:two-layer}
writes it as a finite linear combination of two-layer monomials $Y_\alpha=\exp\left(i\langle\delta,\hat A\rangle\right)\Psi_\alpha$,
all with the exponent $(\delta_0,\delta_1,\delta_\perp)$;
$\Lambda_Y$ is a set of sites in play for each of them,
since the fields of $\Psi_\alpha$ sit at the ends of bonds over $\Lambda_Y$.
As the equation is linear in $\Psi$ and in $\Upsilon$,
the equation that defines $\Upsilon_{Y_\alpha}$, formed with the set $\Lambda_Y$, has no solution for at least one $\alpha$.
Proposition~\ref{prop:two-layer-insertions} and Theorem~\ref{thm:two-layer-principle}(b),
for monomials that contain fields,
therefore require the invertibility of $\omega$, that is $|\rho|<1$;
so does the injectivity of the maps $X\mapsto G_X$ and $k\mapsto G[k]$
(Lemma~\ref{lem:injective-insertions}).
\end{remark}

\section{Elimination of outer layers}\label{sec:elimination}

In this section $\rho\in\R$, $M\in\R$, $w_0,\dots,w_3>0$ and $q>0$,
except where a narrower range of $\rho$ is stated.

Along the coordinate direction $e_1$ the photon field and the Grassmann field form a chain of layers,
the thick layers of Definition~\ref{def:two-layer}:
the thick layer $j$ contains the photon links of the plane $\{x_1=j\}$ and six thin layers of fermion sites.
Consider an element $P$ of the class $\mathcal{P}$ of Definition~\ref{def:polynomials},
that is, a finite combination of products $\exp(i\langle\delta,\hat A\rangle)\cdot\Psi$
of an exponential of the local lift and a monomial in the fields, each product being gauge invariant.
Two exact identities move $P$ out of its lowest plane and into the next two planes.

\emph{The fermion.}
Along $e_1$ the fermion graph is a chain of thin layers,
in which consecutive layers are joined by perfect matchings (Lemma~\ref{lem:layer-bonds}).
Let $x$ be a site, let $x'$ be its neighbour in the next thin layer, and let $i$ be the type of the bond $\{x,x'\}$.
The hopping matrix $w(x,x')=\frac{w_i}{2}\,(\rho-\gamma\cdot\hat d)$ has the eigenvalues $\frac{w_i}{2}\,(\rho\mp1)$;
it is invertible exactly when $\rho^2\neq1$.
In that case the equation of motion at $x'$ can be solved for the field at $x$ by means of $w(x,x')^{-1}$,
and the outcome is a polynomial rule which expresses the field at $x$
through the fields at $x'$ and at the three other neighbours of $x'$
(Proposition~\ref{prop:schwinger-dyson}).

\emph{The photon.}
Fix a finite set of links, all of them contained in one coordinate plane;
the variable of each of these links ranges over the interval $[0,2\pi/q)$.
When the integration over these variables is combined with the summation over the fluxes $z+dk$,
$k$ an integer one-cochain on the fixed links,
each of these variables comes to range over all of $\R$; here one uses that the orbits $\{z+dk\}$ are free.
The construction is local, and it can be carried out on every torus.
As a function of the new variables the integrand extends to an entire function that decays like a Gaussian,
which allows the contour of integration to be moved by $-i\delta$.
The Maxwell weight then produces the factor $\exp\left(\frac12\|d\delta\|^2\right)\cdot\varphi_{d\delta}$,
and the fermion weight produces the polynomial $\mathfrak{V}_\delta$
(a polynomial because the Grassmann generators are nilpotent),
which reads a single thick layer (Proposition~\ref{prop:contour-shift}).
Let $\delta$ be the exponent of a term of $P$ on the plane $\{x_1=j\}$, placed on the plane $\{x_1=j+1\}$.
The exponent $d^*d\delta$ of $\varphi_{d\delta}$ then cancels everything that this term reads on the plane $\{x_1=j\}$
and on the links between the two planes (Proposition~\ref{prop:outer-plane}).

Under these operations the class $\mathcal{P}$ is preserved exactly (Proposition~\ref{prop:class-structure});
in contrast with the corresponding step for Yang--Mills theory in four dimensions,
the elimination involves no approximation, no division by a weight and no average over gauge transformations
(Remark~\ref{rem:yang-mills-elimination}).
Both identities hold on every torus in which a certain box determined by $P$ embeds,
and they remain true after multiplication of the integrand by an arbitrary factor $\Phi$,
that is, by a bounded Borel function of the curvature, which has to be entire wherever the contour is moved.
Owing to this universality, the elimination at a link plane $\{x_1=c\}$ for the reflected element $\Theta_cP$ needs no computation of its own:
the mirrored rule (Proposition~\ref{prop:schwinger-dyson}\eqref{it:mirrored-rule})
and the reflected identity (Proposition~\ref{prop:outer-plane}\eqref{it:reflected-elimination})
follow from it together with the covariance of the functional under reflections
(Proposition~\ref{prop:symmetries}\eqref{it:functional-covariance}).
By composing two reflections one transfers the identities in the same manner to the directions $e_2$ and $e_3$
(Lemma~\ref{lem:reflection-identities}\eqref{it:reflected-universal}).
Consequently the estimates for these directions, too, are carried out on the torus $D(n_i)$ itself (Theorem~\ref{thm:two-layer-reduction}),
even though not all permutations of the coordinate axes leave this torus invariant.

When the identities are applied to the perturbed functional $\mu^t_i$ of Definition~\ref{def:data},
there is exactly one source of error.
The factor $\exp(t\cdot O_{\mathfrak{d}})$ of $\mu^t_i$ takes part in the shift of the contour,
and for real $F$
\[
\exp\left(t\cdot O_{\mathfrak{d}}(F-i\,d\delta)\right)
=\exp\left(t\cdot O_{\mathfrak{d}}(F)\right)\cdot\left(1+\sum_\alpha b_\alpha(t)\cdot\varphi_{\beta_\alpha}(F)\right),
\qquad
\sum_\alpha|b_\alpha(t)|\le\exp(C\cdot t)-1,
\]
where $\alpha$ ranges over a countable set,
the cochains $\beta_\alpha$, together with a box that contains all their supports, can be chosen independently of $t$ and of the torus,
and the constant $C$ depends only on $\delta$ and on the data $\mathfrak{d}$
(Proposition~\ref{prop:contour-shift}\eqref{it:shift-arrays}).
The series of characters therefore converges absolutely,
and for $0\le t\le1$ the sum of the moduli of its coefficients is at most $(\exp(C)-1)\cdot t$.
The terms that it contributes are bounded by means of Proposition~\ref{prop:a-priori}, which requires $0<\rho<1$;
the amplification of Section~\ref{sec:amplification} enters the elimination through the proof of that proposition.

\subsection{Covariant classes and boxes}\label{sec:covariant-classes}

\begin{definition}[Covariant classes]\label{def:covariant-classes}
The set of the generators $\psi_\alpha(x)$, $\bar\psi_\alpha(x)$, $x\in\mathbb{X}$, $\alpha=1,\dots,4$,
of $\Gr(\mathbb{X})$ (Definition~\ref{def:grassmann}) carries a fixed total order,
the one in which the monomials of the representation in Definition~\ref{def:two-layer} are written.
An \emph{ordered monomial} is a product $\Psi=\xi_1\cdots\xi_n$ of generators $\xi_1<\dots<\xi_n$, where $n\geq0$;
for $n=0$ it is $\Psi=1$.
The ordered monomials form a basis of $\Gr(\mathbb{X})$,
and a product of generators taken in an arbitrary order is either $0$ or $\pm$ an ordered monomial.
In this section the word monomial always refers to an ordered monomial.

Let $y_0\in\Z^3$.
The \emph{covariant class} $\mathcal{P}^{+}(y_0)$, respectively $\mathcal{P}^{-}(y_0)$,
is the linear span of the elements $\exp(i\langle\delta,\hat A\rangle)\cdot\Psi$
in which $\delta\in C^1_{\mathrm{c}}$ is real, $\Psi$ is a monomial, and
\begin{equation}\label{eq:elimination-covariance}
d^*\delta=q\cdot\left(\mathfrak{q}(\Psi)\mp\delta_{y_0}\right),
\end{equation}
with the upper sign for $\mathcal{P}^{+}(y_0)$ and the lower sign for $\mathcal{P}^{-}(y_0)$;
here $\delta_{y_0}$ is the indicator function of the photon site $y_0$,
and $\mathfrak{q}(\Psi)$ is the charge density of $\Psi$ (Definition~\ref{def:chamber-operator}).
The elements $\exp(i\langle\delta,\hat A\rangle)\cdot\Psi$, with $(\delta,\Psi)$ ranging over the pairs subject to \eqref{eq:elimination-covariance},
form a basis of $\mathcal{P}^{\pm}(y_0)$,
and those with $d^*\delta=q\cdot\mathfrak{q}(\Psi)$ form a basis of $\mathcal{P}$;
indeed, the monomials are linearly independent,
and so are the functions $\hat A\mapsto\exp(i\langle\delta,\hat A\rangle)$ for pairwise different $\delta$,
which are pairwise different characters of the additive group of a real vector space.
Accordingly, every element $P$ of $\mathcal{P}$ or of $\mathcal{P}^{\pm}(y_0)$ has a unique representation
\[
P=\sum_a\lambda_a\cdot\exp(i\langle\delta_a,\hat A\rangle)\cdot\Psi_a
\]
by such basis elements, with finitely many terms, pairwise different pairs $(\delta_a,\Psi_a)$ and coefficients $\lambda_a\neq0$.
For the elements of $\mathcal{P}^{\pm}(y_0)$ the links and the sites that are read,
and the lowest and the highest plane, are defined through this representation,
in the same way as for the elements of $\mathcal{P}$ in Definition~\ref{def:two-layer}.
A constant element of $\mathcal{P}$ reads no link and no site;
its lowest plane is $+\infty$ and its highest plane is $-\infty$.
\end{definition}

\paragraph{Gauge behaviour and parity.}
The condition $d^*\delta=q\cdot\mathfrak{q}(\Psi)$ is the gauge invariance by which $\mathcal{P}$ is defined (Definition~\ref{def:polynomials}),
whereas \eqref{eq:elimination-covariance} prescribes a definite behaviour under gauge transformations.
Let $\lambda\colon\Z^3\to\R$.
As $\langle\delta,d\lambda\rangle=\langle d^*\delta,\lambda\rangle$ for $\delta$ of finite support,
the gauge transformation \eqref{eq:setting-gauge} by $\lambda$ carries $\exp(i\langle\delta,\hat A\rangle)\cdot\Psi$ into
\[
\exp(i\langle\delta,\hat A-d\lambda\rangle)\cdot\exp\left(iq\langle\mathfrak{q}(\Psi),\lambda\rangle\right)\cdot\Psi
=\exp\left(i\left\langle q\cdot\mathfrak{q}(\Psi)-d^*\delta,\lambda\right\rangle\right)\cdot\exp(i\langle\delta,\hat A\rangle)\cdot\Psi .
\]
Under \eqref{eq:elimination-covariance} the first factor on the right equals $\exp\left(\pm iq\lambda(y_0)\right)$.
So an element of $\mathcal{P}^{+}(y_0)$ is multiplied by $\exp(iq\lambda(y_0))$,
as is the field $\psi(x)$ at a site $x$ with $\pi(x)=y_0$,
an element of $\mathcal{P}^{-}(y_0)$ is multiplied by $\exp(-iq\lambda(y_0))$, as is $\bar\psi(x)$,
and the elements of $\mathcal{P}$ are invariant.

The adjoint of the coboundary on zero-cochains (Definition~\ref{def:cochains}) is given by
\[
(d^*u)(y)=\sum_{\mu=1}^{3}\left(u_\mu(y-e_\mu)-u_\mu(y)\right),\qquad y\in\Z^3,
\]
for every one-cochain $u$ of finite support;
in particular $\sum_{y\in\Z^3}(d^*u)(y)=0$.
Summing \eqref{eq:elimination-covariance} over $\Z^3$ therefore gives $\sum_{y}\mathfrak{q}(\Psi)(y)=\pm1$,
and summing the gauge condition of $\mathcal{P}$ gives $\sum_{y}\mathfrak{q}(\Psi)(y)=0$.
This sum is the number of the factors $\psi_\alpha(x)$ of $\Psi$ minus the number of its factors $\bar\psi_\alpha(x)$,
and it has the parity of the degree of $\Psi$.
Hence the elements of $\mathcal{P}$ are even, and the elements of $\mathcal{P}^{\pm}(y_0)$ are odd.

\paragraph{Cochains of oriented pairs.}
Let $C^1_{\mathrm{c}}(\Z):=C^1_{\mathrm{c}}\cap C^1(\Z)$ be the group of the finitely supported integer-valued one-cochains (Definition~\ref{def:cochains}).
Two photon sites $y,y'$ are nearest neighbours if $y'-y\in\{\pm e_1,\pm e_2,\pm e_3\}$.
For such an oriented pair $(y,y')$ let $1_{(y,y')}\in C^1_{\mathrm{c}}(\Z)$ be the one-cochain
with $\langle 1_{(y,y')},A\rangle=A(y,y')$ for every one-cochain $A$,
read as an odd function of oriented pairs (Definition~\ref{def:cochains}).
The support of $1_{(y,y')}$, that is, the set of the links at which it does not vanish,
consists of the single link that joins $y$ and $y'$.
Put $1_{(y,y)}:=0$; its support is empty.
For every bond $\{x,x'\}$ the vector $\pi(x')-\pi(x)$ is $0$ or one of the unit vectors $\pm e_\mu$
(Proposition~\ref{prop:geometry}\eqref{it:geometry-projection}),
so that the cochain $1_{(\pi(x),\pi(x'))}$ is defined.
For photon sites $y,y'$ that are nearest neighbours or equal, and for every bond $\{x,x'\}$,
the formula for $d^*$ and the definition of $U_A$ (Definition~\ref{def:chamber-operator}) give
\begin{equation}\label{eq:elimination-pair}
d^*1_{(y,y')}=\delta_{y'}-\delta_y,
\qquad
U_A(x,x')=\exp\left(iq\langle 1_{(\pi(x),\pi(x'))},A\rangle\right).
\end{equation}
The second identity remains true with a local lift $\hat A$ in the place of $A$,
because $\exp(iq\hat A(b))=\exp(iqA(b))$ for every link $b$ of the box of the lift (Definition~\ref{def:integrands}).

\paragraph{Rules.}
The covariant classes are tied to $\mathcal{P}$ by the following inclusions,
in which the derivatives are left derivatives (Definition~\ref{def:grassmann}).
For $x\in\mathbb{X}$ and $\alpha=1,\dots,4$ each of the four sets
\begin{equation}\label{eq:elimination-rules}
\bar\psi_\alpha(x)\cdot\mathcal{P}^{+}(\pi(x)),\qquad
\psi_\alpha(x)\cdot\mathcal{P}^{-}(\pi(x)),\qquad
\frac{\partial}{\partial\psi_\alpha(x)}\,\mathcal{P}^{+}(\pi(x)),\qquad
\frac{\partial}{\partial\bar\psi_\alpha(x)}\,\mathcal{P}^{-}(\pi(x))
\end{equation}
is contained in $\mathcal{P}$,
and for every bond $\{x,x'\}$
\begin{equation}\label{eq:elimination-phases}
U_A(x',x)\cdot\mathcal{P}^{+}(\pi(x))\subset\mathcal{P}^{+}(\pi(x')),
\qquad
U_A(x,x')\cdot\mathcal{P}^{-}(\pi(x))\subset\mathcal{P}^{-}(\pi(x')).
\end{equation}

For \eqref{eq:elimination-rules} it suffices to consider basis elements $\exp(i\langle\delta,\hat A\rangle)\cdot\Psi$.
The charge density of a product of generators is the sum of the charge densities of its factors,
and the generators $\psi_\alpha(x)$ and $\bar\psi_\alpha(x)$ have the charge densities $\delta_{\pi(x)}$ and $-\delta_{\pi(x)}$.
Left multiplication by $\bar\psi_\alpha(x)$, and the derivative with respect to $\psi_\alpha(x)$, which removes a factor $\psi_\alpha(x)$,
turn $\Psi$ into $0$ or into $\pm$ a monomial $\Psi'$ with $\mathfrak{q}(\Psi')=\mathfrak{q}(\Psi)-\delta_{\pi(x)}$.
For a basis element of $\mathcal{P}^{+}(\pi(x))$ this gives
$d^*\delta=q\cdot\left(\mathfrak{q}(\Psi)-\delta_{\pi(x)}\right)=q\cdot\mathfrak{q}(\Psi')$,
which is the gauge condition of $\mathcal{P}$.
Likewise, left multiplication by $\psi_\alpha(x)$ and the derivative with respect to $\bar\psi_\alpha(x)$ lead to
$\mathfrak{q}(\Psi')=\mathfrak{q}(\Psi)+\delta_{\pi(x)}$,
and for a basis element of $\mathcal{P}^{-}(\pi(x))$ to $d^*\delta=q\cdot\left(\mathfrak{q}(\Psi)+\delta_{\pi(x)}\right)=q\cdot\mathfrak{q}(\Psi')$.

For \eqref{eq:elimination-phases}, note that by \eqref{eq:elimination-pair} multiplication by $U_A(x',x)$ adds $q\cdot1_{(\pi(x'),\pi(x))}$
to the exponent of a basis element, and multiplication by $U_A(x,x')$ adds $q\cdot1_{(\pi(x),\pi(x'))}$;
the monomial does not change.
For a basis element of $\mathcal{P}^{+}(\pi(x))$ the new exponent $\delta'=\delta+q\cdot1_{(\pi(x'),\pi(x))}$ satisfies
\[
d^*\delta'=q\cdot\left(\mathfrak{q}(\Psi)-\delta_{\pi(x)}\right)+q\cdot\left(\delta_{\pi(x)}-\delta_{\pi(x')}\right)
=q\cdot\left(\mathfrak{q}(\Psi)-\delta_{\pi(x')}\right),
\]
and for a basis element of $\mathcal{P}^{-}(\pi(x))$ the new exponent $\delta'=\delta+q\cdot1_{(\pi(x),\pi(x'))}$ satisfies
\[
d^*\delta'=q\cdot\left(\mathfrak{q}(\Psi)+\delta_{\pi(x)}\right)+q\cdot\left(\delta_{\pi(x')}-\delta_{\pi(x)}\right)
=q\cdot\left(\mathfrak{q}(\Psi)+\delta_{\pi(x')}\right).
\]

\begin{definition}[Enlarged boxes]\label{def:enlarged-boxes}
Let $B=y+[-r,r]^3$ be a box (Definition~\ref{def:integrands}) and let $n\geq0$ be an integer.
The box $B$ \emph{enlarged by} $n$ is
\[
B^{+n}:=y+[-r-n,r+n]^3=B+[-n,n]^3 .
\]
An element of $\mathcal{P}$, or of a covariant class $\mathcal{P}^{\pm}(y_0)$, is \emph{read in} $B$
if the links that it reads lie in $B^{(1)}$ and the sites that it reads lie in $\mathbb{X}(B)$.
\end{definition}

Let $B$ be a box.
An element $P=\sum_a\lambda_a\cdot\exp(i\langle\delta_a,\hat A\rangle)\cdot\Psi_a$ of $\mathcal{P}$ that is read in $B$ belongs to $\mathcal{G}^{\mathrm{ev}}(B)$
(Definition~\ref{def:integrands}).
It is a continuous map from the real cochains on $B^{(1)}$ to $\Gr^{\mathrm{ev}}(\mathbb{X}(B))$,
it is gauge invariant term by term,
and for real $\hat A$ the coefficient of each monomial has modulus at most $\sum_a|\lambda_a|$,
so that the growth condition holds with the margin $\frac12$.
For $g\in W$ the image $gB$ of a box is again a box:
the linear part $\vec g$ lies in $\Ocub$ (Definition~\ref{def:mirrors}), whose elements are the signed permutation matrices,
so that $\vec g$ maps $[-r,r]^3$ onto itself and $g\left(y+[-r,r]^3\right)=gy+[-r,r]^3$.

For sites $x,x'\in\mathbb{X}$ write $x'\sim x$ if $x'$ is one of the four neighbours $x^{(0)},\dots,x^{(3)}$ of $x$
(Definition~\ref{def:fermion-sites});
as $(x^{(i)})^{(i)}=x$, the relation is symmetric, and $\{x,x'\}$ is then a bond.
Let $B$ be a box, let $x\in\mathbb{X}(B)$, let $x'\sim x$ and let $x''\sim x'$.
Then any two of the points $\pi(x)$, $\pi(x')$, $\pi(x'')$ are equal or nearest neighbours, and
\begin{equation}\label{eq:elimination-neighbours}
x',\,x''\in\mathbb{X}(B^{+1}),
\qquad
\supp 1_{(y,y')}\subset(B^{+1})^{(1)}\quad\text{for all }y,y'\in\{\pi(x),\pi(x'),\pi(x'')\}.
\end{equation}
To see this, recall from Proposition~\ref{prop:geometry}\eqref{it:geometry-projection}
that $\pi$ takes the same value at the two ends of a bond of type $0$, $1$ or $2$,
and that the two values differ by a unit vector $\pm e_\mu$ at the ends of a bond of type $3$.
If $x''=x$, the three points are $\pi(x)$ and $\pi(x')$, and $\pi(x')-\pi(x)$ is $0$ or a unit vector.
If $x''\neq x$, then $\{x,x'\}$ and $\{x',x''\}$ are two different bonds at $x'$;
the four bonds at $x'$ have four different types, so at most one of the two has type $3$,
and at most one of the differences $\pi(x')-\pi(x)$, $\pi(x'')-\pi(x')$ is different from zero.
In both cases the three points are $\pi(x)$ and at most one of the six points $\pi(x)\pm e_\mu$, $\mu=1,2,3$.
Since $x\in\mathbb{X}(B)$, the cube $\pi(x)+\left[-\frac12,\frac12\right]^3$ lies in $B$;
the cubes of side one centred at $\pi(x')$ and at $\pi(x'')$ are translates of it by vectors of $[-1,1]^3$,
and therefore lie in $B+[-1,1]^3=B^{+1}$.
This proves $x',x''\in\mathbb{X}(B^{+1})$.
A link that joins two different points among the three is the segment from $\pi(x)$ to one of the points $\pi(x)\pm e_\mu$;
since $\pi(x)\in B$, the segment lies in $B+[-1,1]^3=B^{+1}$.

\begin{lemma}[Elements of $\mathcal{P}$ on a torus]\label{lem:torus-elements}
Let $\rho\in\R$, $M\in\R$, $w_0,\dots,w_3>0$ and $q>0$.
Let $\mathbb{T}=(\Gamma,\chi)$ be a torus and let $B$ be a box that embeds in $\mathbb{T}$ (Definition~\ref{def:integrands}),
so that the sets $B^{+1}+g$, $g\in\Gamma$, are pairwise disjoint.
\begin{enumerate}
\item\label{it:torus-injective}
The projection modulo $\Gamma$ is injective on the set of the cells of $\Z^3$ (sites, links and plaquettes) that are contained in $B^{+1}$,
and on $\mathbb{X}(B^{+1})$.
\item\label{it:torus-derivatives}
Let $\mathbb{X}'\subset\mathbb{X}$ be a set of representatives of $\mathbb{X}_{\mathbb{T}}$ that contains $\mathbb{X}(B)$;
such sets exist by \eqref{it:torus-injective}.
The generators at the sites of $\mathbb{X}'$ are free generators of $\Gr_{\mathbb{T}}$,
and $\Gr(\mathbb{X}(B))$ is a subalgebra of $\Gr_{\mathbb{T}}$.
Let $x'\in\mathbb{X}$ be a site such that $x'$ and its four neighbours lie in $\mathbb{X}(B)$.
Then for $\alpha=1,\dots,4$ and every complex one-cochain $A\in C^1(\Gamma)\otimes\C$ one has, in $\Gr_{\mathbb{T}}$,
\[
\frac{\partial S_{\mathbb{T}}(A)}{\partial\psi_\alpha(x')}=-(\bar\psi D_A)_\alpha(x'),
\qquad
\frac{\partial S_{\mathbb{T}}(A)}{\partial\bar\psi_\alpha(x')}=(D_A\psi)_\alpha(x'),
\]
where
\[
\begin{aligned}
(\bar\psi D_A)(x')&:=M\bar\psi(x')-\sum_{x''\sim x'}\bar\psi(x'')\cdot w(x'',x')\cdot U_A(x'',x'),\\
(D_A\psi)(x')&:=M\psi(x')-\sum_{x''\sim x'}w(x',x'')\cdot U_A(x',x'')\cdot\psi(x'').
\end{aligned}
\]
In particular, the character $\chi$ does not enter these two formulas.
\item\label{it:torus-lift-phase}
Let $\delta\in C^1_{\mathrm{c}}$ be supported in $B^{(1)}$, with $d^*\delta\in q\cdot C^0(\Z)$.
For a point of $\Conf(\Gamma)$ with a representative $(A,z)$ and a local lift $\hat A$ on $B$ (Definition~\ref{def:integrands}),
the number $\exp(i\langle\delta,\hat A\rangle)$ depends neither on the lift nor on the representative.
As a function on $\Conf(\Gamma)$ it is Borel measurable and of modulus one.
\item\label{it:torus-convergence}
The integral $\int_{\Conf(\Gamma)}\exp\left(-\frac12\|F\|^2\right)$ is finite.
Let $P\in\mathcal{P}$ be read in $B$, and for a point of $\Conf(\Gamma)$ with a representative $(A,z)$ put
\[
J_{\mathbb{T}}[P](A,z):=\int_{\mathbb{T}}\exp\left(-S_{\mathbb{T}}(A)\right)\cdot P(\hat A),
\]
$\hat A$ a local lift on $B$.
Then $J_{\mathbb{T}}[P]$ is a bounded Borel function on $\Conf(\Gamma)$,
and for every bounded Borel function $\Phi$ on $C^2(\Gamma)$
\[
Z_{\mathbb{T}}[\Phi\cdot P]=\int_{\Conf(\Gamma)}\exp\left(-\frac12\|F\|^2\right)\cdot\Phi(F)\cdot J_{\mathbb{T}}[P],
\]
where the integral converges absolutely,
and so do the integrals of the terms that correspond to the terms of the representation of $P$.
Moreover
\[
\left|Z_{\mathbb{T}}[\Phi\cdot P]\right|\le\sup|\Phi|\cdot\sup|J_{\mathbb{T}}[P]|\cdot\int_{\Conf(\Gamma)}\exp\left(-\frac12\|F\|^2\right).
\]
In particular, $\Phi\mapsto Z_{\mathbb{T}}[\Phi\cdot P]$ is continuous for the supremum norm.
\end{enumerate}
\end{lemma}

\begin{proof}
\eqref{it:torus-injective}
Let $\sigma$ be a cell contained in $B^{+1}$ and let $g\in\Gamma\setminus\{0\}$.
Then $\sigma+g\subset B^{+1}+g$, and $B^{+1}+g$ does not meet $B^{+1}$;
so $\sigma+g$ is not contained in $B^{+1}$.
On the fermion sites, $\pi$ commutes with the translations by $\Gamma$ (Definition~\ref{def:tori}).
If $x$ and $x+g$ both lay in $\mathbb{X}(B^{+1})$,
the cubes $\pi(x)+\left[-\frac12,\frac12\right]^3$ and $\pi(x+g)+\left[-\frac12,\frac12\right]^3=\pi(x)+g+\left[-\frac12,\frac12\right]^3$ would both lie in $B^{+1}$,
and the point $\pi(x)+g$ would belong to $B^{+1}\cap(B^{+1}+g)$, which is empty.

\eqref{it:torus-derivatives}
The first two assertions are Definition~\ref{def:tori} together with $\mathbb{X}(B)\subset\mathbb{X}'$.
By the same definition,
\[
S_{\mathbb{T}}(A)=\sum_{x\in\mathbb{X}'}\bar\psi(x)\cdot(D_A\psi)(x),
\qquad
(D_A\psi)(x)=M\psi(x)-\sum_{i=0}^{3}w(x,x^{(i)})\cdot U_A(x,x^{(i)})\cdot\psi(x^{(i)}),
\]
where a column $\psi(x^{(i)})$ with $x^{(i)}\notin\mathbb{X}'$ stands for $\chi(g)\cdot\psi(x^{(i)}-g)$,
$g\in\Gamma$ being the vector with $x^{(i)}-g\in\mathbb{X}'$.
The generators $\psi_{\alpha'}(x')$, $\alpha'=1,\dots,4$, occur in the summand of $x=x'$ through the term $M\psi(x')$,
and in the term of index $i$ of the summand of $x\in\mathbb{X}'$ whenever $x^{(i)}-g=x'$ for some $g\in\Gamma$.
In the latter case $x-g$ is a neighbour of $x'$,
since translations by $\Gamma$ map bonds to bonds of the same type (Definition~\ref{def:tori});
hence $x-g\in\mathbb{X}(B)\subset\mathbb{X}'$.
As $x$ lies in $\mathbb{X}'$ as well, $g=0$;
so $x$ is a neighbour of $x'$, and the term is $-\bar\psi(x)\cdot w(x,x')\cdot U_A(x,x')\cdot\psi(x')$, without a factor $\chi(g)$.
Conversely, every neighbour $x''$ of $x'$ belongs to $\mathbb{X}'$ and contributes exactly one such term.
Therefore the terms of $S_{\mathbb{T}}(A)$ that contain one of the generators $\psi_1(x'),\dots,\psi_4(x')$ add up to
\[
\bar\psi(x')\cdot M\cdot\psi(x')-\sum_{x''\sim x'}\bar\psi(x'')\cdot w(x'',x')\cdot U_A(x'',x')\cdot\psi(x').
\]
The left derivative with respect to $\psi_\alpha(x')$ annihilates all other terms.
In these terms it has to be moved across the odd generator $\bar\psi_{\alpha'}(x'')$, $x''=x'$ or $x''\sim x'$,
that stands to the left of $\psi_\alpha(x')$, and this produces the sign $-1$:
\[
\frac{\partial S_{\mathbb{T}}(A)}{\partial\psi_\alpha(x')}
=-M\bar\psi_\alpha(x')+\sum_{x''\sim x'}\ \sum_{\alpha'=1}^{4}\bar\psi_{\alpha'}(x'')\cdot w(x'',x')_{\alpha'\alpha}\cdot U_A(x'',x')
=-(\bar\psi D_A)_\alpha(x').
\]
The generators $\bar\psi_{\alpha'}(x')$ occur only in the summand of $x=x'$, as its leftmost factor,
because the summands are indexed by the different sites of $\mathbb{X}'$.
The neighbours of $x'$ lie in $\mathbb{X}(B)\subset\mathbb{X}'$, so that no column of $(D_A\psi)(x')$ has to be rewritten,
and the derivative of this summand with respect to $\bar\psi_\alpha(x')$ is $(D_A\psi)_\alpha(x')$.

\eqref{it:torus-lift-phase}
The conditions that define a local lift on $B$, namely $d\hat A=F$ on $B^{(2)}$ and $\exp(iq\hat A(b))=\exp(iqA(b))$ for $b\in B^{(1)}$,
involve only $F$ and the link phases,
which are invariant under $(A,z)\mapsto\left(A+\frac{2\pi}{q}\cdot k,z-dk\right)$, $k\in C^1(\Gamma;\Z)$;
so the set of the local lifts on $B$ depends on the point of $\Conf(\Gamma)$ alone.
Two of these lifts differ by $\frac{2\pi}{q}\cdot d\lambda$,
where $\lambda$ is an integer valued function on the sites of $B$ (Definition~\ref{def:integrands}); extend $\lambda$ by zero.
Since $\delta$ is supported in $B^{(1)}$,
\[
\frac{2\pi}{q}\cdot\langle\delta,d\lambda\rangle=\frac{2\pi}{q}\cdot\langle d^*\delta,\lambda\rangle\in\frac{2\pi}{q}\cdot q\Z=2\pi\Z ,
\]
so that $\exp(i\langle\delta,\hat A\rangle)$ depends neither on the lift nor on the representative.
The modulus is one because $\delta$ and $\hat A$ are real.
For measurability, fix one of the countably many closed fluxes $z\in C^2(\Gamma;\Z)$,
and choose an integer one-cochain $k'$ on $B^{(1)}$ with $dk'=z$ on $B^{(2)}$ (Definition~\ref{def:integrands}).
Then $\hat A:=A|_{B^{(1)}}+\frac{2\pi}{q}\cdot k'$ is a local lift of the class of $(A,z)$ for every $A$,
and $\exp(i\langle\delta,\hat A\rangle)$ is a continuous function of $A$.

\eqref{it:torus-convergence}
Let $\mathbb{T}^{(1)}$ be the set of the links of $\Z^3/\Gamma$.
There are $3|\mathbb{T}|$ links and $3|\mathbb{T}|$ plaquettes $\sigma$ of $\Z^3/\Gamma$, three of each for every class of photon sites.
Put
\[
C_q:=\sup_{y\in\R}\,\sum_{k\in\Z}\exp\left(-\frac12\left(y+\frac{2\pi}{q}\cdot k\right)^2\right),
\]
which is finite, the sum being a continuous function of $y$ with the period $2\pi/q$.
Extending the sum over the closed fluxes to all of $C^2(\Gamma;\Z)$ can only increase it, all terms being positive,
and the extended sum factorises over the plaquettes;
hence, for every real $A\in C^1(\Gamma)$,
\[
\sum_{\substack{z\in C^2(\Gamma;\Z)\\ dz=0}}\exp\left(-\frac12\left\|dA+\frac{2\pi}{q}\cdot z\right\|^2\right)
\le\prod_{\sigma}\ \sum_{k\in\Z}\exp\left(-\frac12\left((dA)_\sigma+\frac{2\pi}{q}\cdot k\right)^2\right)
\le C_q^{3|\mathbb{T}|},
\]
the product being taken over the plaquettes of $\Z^3/\Gamma$.
Integrating over $A\in[0,2\pi/q)^{\mathbb{T}^{(1)}}$, a box of volume $(2\pi/q)^{3|\mathbb{T}|}$,
and recalling the measure of $\Conf(\Gamma)$ (Definition~\ref{def:configurations}), whose terms are all non-negative here, one obtains
\[
\int_{\Conf(\Gamma)}\exp\left(-\frac12\|F\|^2\right)
=\sum_{\substack{z\in C^2(\Gamma;\Z)\\ dz=0}}\ \int_{[0,2\pi/q)^{\mathbb{T}^{(1)}}}\mathrm{d}A\,\exp\left(-\frac12\left\|dA+\frac{2\pi}{q}\cdot z\right\|^2\right)
\le\left(\frac{2\pi}{q}\cdot C_q\right)^{3|\mathbb{T}|}.
\]

Let $P=\sum_a\lambda_a\cdot\exp(i\langle\delta_a,\hat A\rangle)\cdot\Psi_a$ be the representation of $P$ (Definition~\ref{def:covariant-classes}).
As $P$ is read in $B$, every $\delta_a$ is supported in $B^{(1)}$, with $d^*\delta_a=q\cdot\mathfrak{q}(\Psi_a)\in q\cdot C^0(\Z)$,
and every $\Psi_a$ lies in $\Gr(\mathbb{X}(B))$.
The projection is injective on $\mathbb{X}(B)$ by \eqref{it:torus-injective}, and $P(\hat A)$ is even;
so $P(\hat A)$ defines an element of $\Gr_{\mathbb{T}}$ (Definition~\ref{def:tori}), and
\[
J_{\mathbb{T}}[P](A,z)=\sum_a\lambda_a\cdot\exp(i\langle\delta_a,\hat A\rangle)\cdot\int_{\mathbb{T}}\exp\left(-S_{\mathbb{T}}(A)\right)\cdot\Psi_a .
\]
The element $\exp\left(-S_{\mathbb{T}}(A)\right)$ is a finite sum, the generators being nilpotent,
and its coefficients are polynomials in the link phases $\exp(\pm iqA(b))$, $b\in\mathbb{T}^{(1)}$, with coefficients that do not depend on $A$.
Hence each Berezin integral $\int_{\mathbb{T}}\exp\left(-S_{\mathbb{T}}(A)\right)\cdot\Psi_a$ is such a polynomial:
a continuous function of $A$, invariant under $A\mapsto A+\frac{2\pi}{q}\cdot k$, $k\in C^1(\Gamma;\Z)$,
and bounded by the sum of the moduli of its coefficients, the phases having modulus one.
By \eqref{it:torus-lift-phase} every factor $\exp(i\langle\delta_a,\hat A\rangle)$ is a Borel function of modulus one on $\Conf(\Gamma)$.
So every term of $J_{\mathbb{T}}[P]$, and $J_{\mathbb{T}}[P]$ itself, is a bounded Borel function on $\Conf(\Gamma)$.

Let $\Phi$ be a bounded Borel function on $C^2(\Gamma)$.
The space $C^2(\Gamma)$ has finite dimension,
so that $\Phi$ depends on finitely many values $F_\sigma$ of the curvature,
and the embedding of a product imposes no condition on a factor of this kind (Definition~\ref{def:integrands}).
Because $P\in\mathcal{G}^{\mathrm{ev}}(B)$ and $B$ embeds in $\mathbb{T}$,
the product $\Phi\cdot P$ embeds in $\mathbb{T}$ and lies in the domain of $Z_{\mathbb{T}}$.
Since $\Phi(F)$ is a scalar,
$\int_{\mathbb{T}}\exp\left(-S_{\mathbb{T}}(A)\right)\cdot\Phi(F)\cdot P(\hat A)=\Phi(F)\cdot J_{\mathbb{T}}[P](A,z)$,
and Definition~\ref{def:functional} gives the asserted formula for $Z_{\mathbb{T}}[\Phi\cdot P]$.
Its integrand, and the integrand of each of its terms, is bounded in modulus by a constant times $\exp\left(-\frac12\|F\|^2\right)$,
the constant being $\sup|\Phi|\cdot\sup|J_{\mathbb{T}}[P]|$ for the whole integrand;
by the first part of the proof these integrals converge absolutely, and the bound on $|Z_{\mathbb{T}}[\Phi\cdot P]|$ follows.
Applied to the difference of two bounded Borel functions, this bound gives the continuity of the linear map $\Phi\mapsto Z_{\mathbb{T}}[\Phi\cdot P]$ for the supremum norm.
\end{proof}

\subsection{Factors and universal identities}\label{sec:universal-identities}

\begin{definition}[Factors]\label{def:factors}
Let $\mathbb{T}=(\Gamma,\chi)$ be a torus.

A \emph{factor} on $\mathbb{T}$ is a bounded Borel function $\Phi\colon C^2(\Gamma)\to\C$.
Choose one plaquette of $\Z^3$ in each of the finitely many classes of plaquettes modulo $\Gamma$;
the values at the chosen plaquettes are linear coordinates on $C^2(\Gamma)$.
Written in these coordinates, $\Phi$ is a bounded Borel function of finitely many $F_\sigma$, hence a product without integrands in $\mathcal{A}(\R^3)$ (Definition~\ref{def:integrands}),
and at the periodic extension $\widetilde F$ of $F\in C^2(\Gamma)$ it takes the value $\Phi(F)$.
Whether a product embeds in a lattice depends on its boxes alone;
a factor has no box and therefore embeds in every lattice,
and multiplication by $\Phi$ leaves unchanged the boxes of the products of which an element of $\mathcal{A}(\R^3)$ is the sum.
Consequently, if $G\in\mathcal{A}(\R^3)$ embeds in $\mathbb{T}$, then so does $\Phi\cdot G\in\mathcal{A}(\R^3)$,
and in $Z_{\mathbb{T}}[\Phi\cdot G]$ (Definition~\ref{def:functional}) the factor is evaluated at the curvature $F=dA+(2\pi/q)\cdot z$ of the point $(A,z)$ of $\Conf(\Gamma)$.

An \emph{entire factor} on $\mathbb{T}$ is an entire bosonic factor on $\mathbb{T}$ in the sense of Definition~\ref{def:polynomials}.
Its extension to an entire function on $C^2(\Gamma)\otimes\C$ is unique:
two such extensions differ by an entire function that vanishes on $C^2(\Gamma)$,
and this function vanishes identically, its Taylor coefficients at the origin being derivatives along $C^2(\Gamma)$.
The extension is denoted by the same letter.
For an entire factor $\Phi$ on $\mathbb{T}$ and $r\geq0$ put
\[
\|\Phi\|_r:=\sup\left\{|\Phi(F)| : F\in C^2(\Gamma)\otimes\C,\ |\Im F_\sigma|\leq r\ \text{for every plaquette}\ \sigma\right\}.
\]
This number is finite by Definition~\ref{def:polynomials};
for $r=0$ the supremum is taken over a subset of the set that belongs to $r=1$.
An entire factor is a factor, being continuous on $C^2(\Gamma)$ and bounded there by $\|\Phi\|_0$.

The \emph{periodisation} of a cochain $u\in C^k_{\mathrm{c}}$ is the $\Gamma$-periodic cochain $\sigma\mapsto\sum_{v\in\Gamma}u(\sigma+v)$; at each cell the sum is finite.
If $B$ is a box that embeds in $\mathbb{T}$, a cochain $u\in C^k_{\mathrm{c}}$ whose support lies in $B^{+1}$ (Definition~\ref{def:enlarged-boxes}) is identified with its periodisation.
For two such cochains, the pairing over the cells of $\Z^3$ (Definition~\ref{def:cochains}) equals the pairing of their periodisations over the cells of $\Z^3/\Gamma$ (Definition~\ref{def:configurations}), and the norms agree likewise.
Indeed, by Lemma~\ref{lem:torus-elements}\eqref{it:torus-injective} each class of cells modulo $\Gamma$ contains at most one cell contained in $B^{+1}$;
at a class with such a cell the periodisation takes the value of the cochain at that cell, and at every other class it vanishes.

For an entire factor $\Phi$ on $\mathbb{T}$ and a real $\eta\in C^2(\Gamma)$ the \emph{shifted factor} is
\[
\Phi^{[\eta]}(F):=\Phi(F-i\eta),\qquad F\in C^2(\Gamma)\otimes\C .
\]
It is again an entire factor on $\mathbb{T}$:
it is entire, and with $\|\eta\|_\infty:=\max_\sigma|\eta_\sigma|$ the inequality $|\Im(F-i\eta)_\sigma|\leq|\Im F_\sigma|+\|\eta\|_\infty$ gives $\|\Phi^{[\eta]}\|_r\leq\|\Phi\|_{r+\|\eta\|_\infty}$.
For a real $\delta\in C^1_{\mathrm{c}}$ one writes $\Phi^{(\delta)}:=\Phi^{[d\delta]}$, where $d\delta$ stands for its periodisation.

Let $\Pi$ be a mirror, $\theta:=\theta_\Pi$ its reflection and $\vec\theta$ the linear part of $\theta$.
The \emph{reflected torus} is $\theta\mathbb{T}=(\theta\Gamma,\chi\circ\vec\theta)$, the image of $\mathbb{T}$ under $\theta$ as in Definition~\ref{def:tori} (note $\vec\theta^{-1}=\vec\theta$);
since $\vec\theta^{2}=1$, one has $\theta(\theta\mathbb{T})=\mathbb{T}$.
The pull-back $\theta^{*}$ maps $C^2(\theta\Gamma)$ onto $C^2(\Gamma)$ (Definition~\ref{def:configurations});
it is orthogonal and commutes with $d$ (Definition~\ref{def:cochains}), $(\theta^{*})^{2}=(\theta\theta)^{*}=1$, and its complex linear extension commutes with complex conjugation.
The \emph{reflected factor} of a factor $\Phi$ on $\mathbb{T}$ is
\[
\Phi^{\theta}(F):=\overline{\Phi\left(\theta^{*}\overline{F}\right)},
\]
for $F\in C^2(\theta\Gamma)$, and for $F\in C^2(\theta\Gamma)\otimes\C$ if $\Phi$ is entire.
On real $F$ it equals $\overline{\Phi(\theta^{*}F)}$, which is the image of $\Phi$ under the reflection $\Theta_\Pi$ of Definition~\ref{def:grassmann-reflections}, acting on coefficient functions as in \eqref{eq:setting-reflection-coefficients}.
The reflected factor $\Phi^{\theta}$ is a factor on $\theta\mathbb{T}$, and
\[
\left(\Phi^{\theta}\right)^{\theta}(F)=\overline{\Phi^{\theta}\left(\theta^{*}\overline{F}\right)}=\Phi\left(\theta^{*}\theta^{*}F\right)=\Phi(F),
\]
so that $(\Phi^{\theta})^{\theta}=\Phi$.
If $\Phi$ is entire, then so is $\Phi^{\theta}$, the complex conjugate of a holomorphic function of $\overline{F}$ being holomorphic in $F$;
moreover $\|\Phi^{\theta}\|_r=\|\Phi\|_r$, because $F\mapsto\theta^{*}\overline{F}$ maps the set of the $F\in C^2(\theta\Gamma)\otimes\C$ with $|\Im F_\sigma|\leq r$ for all $\sigma$ bijectively onto the corresponding subset of $C^2(\Gamma)\otimes\C$
(the map $\theta^{*}$ permutes the values at the plaquettes up to signs).

For $g\in W$ put
\[
g_{*}:=(g^{-1})^{*},\qquad \operatorname{sgn}(g):=\det\vec g ;
\]
thus $\operatorname{sgn}(g)=1$ for even and $\operatorname{sgn}(g)=-1$ for odd $g$ (Definition~\ref{def:mirrors}).
By Definition~\ref{def:cochains}, $(g^{-1})^{*}g^{*}=(gg^{-1})^{*}=1$ and $g^{*}(g^{-1})^{*}=(g^{-1}g)^{*}=1$, so that $g_{*}=(g^{*})^{-1}$.
As $g^{*}$ maps $C^k(g\Gamma)$ to $C^k(\Gamma)$ (Definition~\ref{def:configurations}), $g_{*}$ maps $C^k(\Gamma)$ to $C^k(g\Gamma)$; it also maps $C^k_{\mathrm{c}}$ to $C^k_{\mathrm{c}}$.
Like $g^{*}$, the map $g_{*}$ is orthogonal and commutes with $d$;
it commutes with $d^{*}$ as well, as one sees by taking adjoints in $g^{*}d=dg^{*}$, the adjoint of $g^{*}$ being $g_{*}$.
For $g,g'\in W$ the rule for the pull-back of a product in Definition~\ref{def:cochains} gives
\[
(gg')_{*}=\left(g'^{-1}g^{-1}\right)^{*}=(g^{-1})^{*}(g'^{-1})^{*}=g_{*}g'_{*}.
\]
For a cochain $u$ and a cell $\sigma$ one has $(g_{*}u)(g\sigma)=\pm u(\sigma)$ with the orientation sign of $g$ at $\sigma$,
which is also the sign with which $g^{-1}$ carries the orientation of $g\sigma$ to that of $\sigma$;
hence $\supp g_{*}u=g(\supp u)$.
The push-forward commutes with periodisation: for $u\in C^k_{\mathrm{c}}$,
\[
g_{*}\left(\sum_{v\in\Gamma}u(\cdot+v)\right)=\sum_{v'\in g\Gamma}(g_{*}u)(\cdot+v'),
\]
because the orientation sign of $g^{-1}$ at a cell does not change when the cell is translated, and $g^{-1}(\sigma+\vec g v)=g^{-1}\sigma+v$ for $v\in\Gamma$.
For the reflection $\theta$ one has $\theta^{-1}=\theta$, hence $\theta_{*}=\theta^{*}$.
\end{definition}

\begin{definition}[Universal identities]\label{def:universal-identity}
Let $\rho,M\in\R$, $w_0,\dots,w_3>0$ and $q>0$, and let the functionals $Z_{\mathbb{T}}$ be formed with these parameters (Definition~\ref{def:functional}).
A \emph{universal identity} is a triple
\[
\mathcal{U}=\left(P;(\eta_a,P_a)_{a\in I};B'\right)
\]
formed by a box $B'$, a finite set $I$, elements $P$ and $P_a$, $a\in I$, of $\mathcal{P}$ (Definition~\ref{def:polynomials}) that are read in $B'$ (Definition~\ref{def:enlarged-boxes}),
and real cochains $\eta_a\in C^2_{\mathrm{c}}$ supported on plaquettes contained in $B'$, such that
\begin{equation}\label{eq:elimination-universal}
Z_{\mathbb{T}}[\Phi\cdot P]=\sum_{a\in I}Z_{\mathbb{T}}\left[\Phi^{[\eta_a]}\cdot P_a\right]
\end{equation}
for every torus $\mathbb{T}$ in which $B'$ embeds and every entire factor $\Phi$ on $\mathbb{T}$;
in \eqref{eq:elimination-universal} each $\eta_a$ stands for its periodisation (Definition~\ref{def:factors}).
The universal identity $\mathcal{U}$ is \emph{of the first kind} if $\eta_a=0$ for every $a\in I$ and
\[
Z_{\mathbb{T}}[\Phi\cdot P]=\sum_{a\in I}Z_{\mathbb{T}}[\Phi\cdot P_a]
\]
holds for every torus $\mathbb{T}$ in which $B'$ embeds and every factor $\Phi$ on $\mathbb{T}$, entire or not.

For a mirror $\Pi$ with the reflection $\theta=\theta_\Pi$, and for $g\in W$, one sets
\[
\begin{aligned}
\Theta_\Pi\mathcal{U}&:=\left(\Theta_\Pi P;(-\theta^{*}\eta_a,\Theta_\Pi P_a)_{a\in I};\theta B'\right),\\
\mathcal{U}^{g}&:=\left(P^{g};(\operatorname{sgn}(g)\cdot g_{*}\eta_a,P_a^{g})_{a\in I};gB'\right),
\end{aligned}
\]
where $P^{g}$ and $P_a^{g}$ are the images of the even elements $P$ and $P_a$ under $g$, as in \eqref{eq:setting-group-action}.
\end{definition}

\begin{lemma}[Reflection of identities]\label{lem:reflection-identities}
Let $\rho,M\in\R$, $w_0,\dots,w_3>0$ and $q>0$; the functionals $Z_{\mathbb{T}}$ and the universal identities are those of these parameters.
Let $\Pi$ be a mirror and $\theta:=\theta_\Pi$.
\begin{enumerate}
\item\label{it:reflected-functional}
Let $\mathbb{T}=(\Gamma,\chi)$ be a torus, and let $G\in\mathcal{P}$ be read in a box $B$ such that $\theta B$ embeds in $\mathbb{T}$.
Then $\Theta_\Pi G$ lies in $\mathcal{P}$ and is read in $\theta B$ (Proposition~\ref{prop:class-structure}\eqref{it:class-stable}),
and every factor $\Phi$ on $\mathbb{T}$ satisfies
\[
Z_{\mathbb{T}}[\Phi\cdot\Theta_\Pi G]=\overline{Z_{\theta\mathbb{T}}\left[\Phi^{\theta}\cdot G\right]}.
\]
\item\label{it:reflected-shift}
Let $\Phi$ be an entire factor on a torus $\mathbb{T}=(\Gamma,\chi)$, and let $\eta\in C^2(\theta\Gamma)$ be real.
Then
\[
\left(\left(\Phi^{\theta}\right)^{[\eta]}\right)^{\theta}=\Phi^{[-\theta^{*}\eta]};
\]
explicitly, $\left((\Phi^{\theta})^{[\eta]}\right)^{\theta}(F)=\Phi(F+i\theta^{*}\eta)$ for $F\in C^2(\Gamma)\otimes\C$.
\item\label{it:reflected-universal}
Let $\mathcal{U}$ be a universal identity.
Then $\Theta_\Pi\mathcal{U}$ is a universal identity, and consequently $\mathcal{U}^{g}$ is a universal identity for every $g\in W$.
If $\mathcal{U}$ is of the first kind, then so are $\Theta_\Pi\mathcal{U}$ and $\mathcal{U}^{g}$.
\end{enumerate}
\end{lemma}

\begin{proof}
\eqref{it:reflected-functional}
By Definition~\ref{def:mirrors}, $\theta$ is an affine map of $\R^3$ whose linear part $\vec\theta$ is a signed permutation matrix;
thus $\theta(x+v)=\theta(x)+\vec\theta v$ for all $x,v\in\R^3$, and $\vec\theta$ maps the cube $[-1,1]^3$ onto itself.
In particular the image $\theta B$ of the box $B$ is a box (Lemma~\ref{lem:admissible-images}\eqref{it:images-boxes}), and
\[
\theta\left(B+[-1,1]^3+v\right)=\theta B+[-1,1]^3+\vec\theta v,\qquad v\in\theta\Gamma .
\]
The map $v\mapsto\vec\theta v$ is a bijection of $\theta\Gamma=\vec\theta\Gamma$ onto $\Gamma$, because $\vec\theta^{2}=1$;
as $\theta B$ embeds in $\mathbb{T}$, the sets on the right side, $v\in\theta\Gamma$, are pairwise disjoint.
The reflection $\theta$ is a bijection of $\R^3$, so the sets $B+[-1,1]^3+v$, $v\in\theta\Gamma$, are pairwise disjoint as well:
the box $B$ embeds in $\theta\Gamma$, that is, in the torus $\theta\mathbb{T}$ (Definition~\ref{def:integrands}).
The element $G$ belongs to $\mathcal{G}^{\mathrm{ev}}(B)$ (Definitions~\ref{def:enlarged-boxes} and~\ref{def:polynomials}), and $\Phi^{\theta}$ is a factor on $\theta\mathbb{T}$;
so $\Phi^{\theta}\cdot G$ is an element of $\mathcal{A}(\R^3)$ that embeds in $\theta\mathbb{T}$ (Definition~\ref{def:integrands}).
Apply Proposition~\ref{prop:symmetries}\eqref{it:functional-covariance} on the torus $\theta\mathbb{T}$ to the odd element $\theta$ of $W$ and to $\Phi^{\theta}\cdot G$.
Since $\theta(\theta\mathbb{T})=\mathbb{T}$, it gives
\[
Z_{\mathbb{T}}\left[\left(\Phi^{\theta}\cdot G\right)^{\theta}\right]=\overline{Z_{\theta\mathbb{T}}\left[\Phi^{\theta}\cdot G\right]}.
\]
By \eqref{eq:setting-group-action}, $(\Phi^{\theta}\cdot G)^{\theta}=\Theta_\Pi(\Phi^{\theta}\cdot G)$.
By \eqref{eq:setting-reflection-coefficients}, $\Theta_\Pi$ replaces the scalar factor $\Phi^{\theta}$ by the function
\[
F\mapsto\overline{\Phi^{\theta}(\theta^{*}F)}=\Phi(\theta^{*}\theta^{*}F)=\Phi(F)
\]
of the real curvature $F$.
Hence $(\Phi^{\theta}\cdot G)^{\theta}=\Phi\cdot\Theta_\Pi G$,
and the identity just obtained from the covariance of the functional takes the form
\[
Z_{\mathbb{T}}[\Phi\cdot\Theta_\Pi G]=\overline{Z_{\theta\mathbb{T}}\left[\Phi^{\theta}\cdot G\right]}
\]
asserted in \eqref{it:reflected-functional}.

\eqref{it:reflected-shift}
For $F\in C^2(\Gamma)\otimes\C$ the definitions of the reflected and of the shifted factor give
\[
\begin{aligned}
\left(\left(\Phi^{\theta}\right)^{[\eta]}\right)^{\theta}(F)
&=\overline{\left(\Phi^{\theta}\right)^{[\eta]}\left(\theta^{*}\overline{F}\right)}
=\overline{\Phi^{\theta}\left(\theta^{*}\overline{F}-i\eta\right)}\\
&=\Phi\left(\theta^{*}\,\overline{\theta^{*}\overline{F}-i\eta}\right)
=\Phi\left(\theta^{*}\left(\theta^{*}F+i\eta\right)\right)
=\Phi\left(F+i\theta^{*}\eta\right),
\end{aligned}
\]
because $\theta^{*}$ commutes with complex conjugation, $\eta$ is real and $(\theta^{*})^{2}=1$.
The right side is $\Phi^{[-\theta^{*}\eta]}(F)$.

\eqref{it:reflected-universal}
Write $\mathcal{U}=(P;(\eta_a,P_a)_{a\in I};B')$, and consider first one mirror.
By Proposition~\ref{prop:class-structure}\eqref{it:class-stable}, $\Theta_\Pi P$ and the elements $\Theta_\Pi P_a$ lie in $\mathcal{P}$ and are read in $\theta B'$.
The cochains $-\theta^{*}\eta_a=-\theta_{*}\eta_a$ lie in $C^2_{\mathrm{c}}$, are real, and are supported on $\theta(\supp\eta_a)$, that is, on plaquettes contained in $\theta B'$.
Let $\mathbb{T}=(\Gamma,\chi)$ be a torus in which $\theta B'$ embeds, and let $\Phi$ be an entire factor on $\mathbb{T}$.
The argument at the beginning of the proof of \eqref{it:reflected-functional}, with $B'$ in place of $B$, shows that $B'$ embeds in $\theta\mathbb{T}$;
and $\Phi^{\theta}$ is an entire factor on $\theta\mathbb{T}$.
For $a\in I$ let
\[
\Phi_a:=\left(\left(\Phi^{\theta}\right)^{[\eta_a]}\right)^{\theta},
\]
with $\eta_a$ periodised over $\theta\Gamma$.
It is an entire factor on $\theta(\theta\mathbb{T})=\mathbb{T}$, and $\Phi_a^{\theta}=(\Phi^{\theta})^{[\eta_a]}$
by the identity $(\Phi^{\theta})^{\theta}=\Phi$ of Definition~\ref{def:factors}, applied to the factor $(\Phi^{\theta})^{[\eta_a]}$ on $\theta\mathbb{T}$.
The map $\theta^{*}=\theta_{*}$ carries the periodisation of $\eta_a$ over $\theta\Gamma$ to the periodisation of $\theta^{*}\eta_a$ over $\Gamma$;
so \eqref{it:reflected-shift} yields $\Phi_a=\Phi^{[-\theta^{*}\eta_a]}$, with $\theta^{*}\eta_a$ periodised over $\Gamma$.
The three equalities of the next display are, in this order,
\eqref{it:reflected-functional} for $\Phi$ and $P$,
the identity \eqref{eq:elimination-universal} of $\mathcal{U}$ on the torus $\theta\mathbb{T}$ with the entire factor $\Phi^{\theta}$,
and \eqref{it:reflected-functional} for $\Phi_a$ and $P_a$, read from right to left:
\[
\begin{aligned}
Z_{\mathbb{T}}[\Phi\cdot\Theta_\Pi P]
&=\overline{Z_{\theta\mathbb{T}}\left[\Phi^{\theta}\cdot P\right]}
=\sum_{a\in I}\overline{Z_{\theta\mathbb{T}}\left[\left(\Phi^{\theta}\right)^{[\eta_a]}\cdot P_a\right]}\\
&=\sum_{a\in I}Z_{\mathbb{T}}[\Phi_a\cdot\Theta_\Pi P_a].
\end{aligned}
\]
As $\Phi_a=\Phi^{[-\theta^{*}\eta_a]}$, this is \eqref{eq:elimination-universal} for $\Theta_\Pi\mathcal{U}$, which is therefore a universal identity.
Suppose that $\mathcal{U}$ is of the first kind.
Then every $\eta_a$ vanishes, hence so does every $-\theta^{*}\eta_a$, and the identity of $\mathcal{U}$ holds on $\theta\mathbb{T}$ for every factor on $\theta\mathbb{T}$.
For an arbitrary factor $\Phi$ on $\mathbb{T}$ the same three equalities hold with $\Phi^{\theta}$ in place of $(\Phi^{\theta})^{[\eta_a]}$ and with $\Phi_a=(\Phi^{\theta})^{\theta}=\Phi$;
thus $\Theta_\Pi\mathcal{U}$ is of the first kind.

Let now $g\in W$.
The group $W$ is generated by the reflections, and each of them is an involution (Definition~\ref{def:mirrors});
hence $g=\theta_{\Pi_1}\cdots\theta_{\Pi_r}$ for some mirrors $\Pi_1,\dots,\Pi_r$.
By \eqref{eq:setting-group-action}, whose right side does not depend on the choice of the word (Proposition~\ref{prop:symmetries}\eqref{it:reflection-words}),
\[
P^{g}=\Theta_{\Pi_1}\cdots\Theta_{\Pi_r}P,\qquad P_a^{g}=\Theta_{\Pi_1}\cdots\Theta_{\Pi_r}P_a .
\]
Starting from $\mathcal{U}$, form successively the triples
\[
\Theta_{\Pi_r}\mathcal{U},\quad\Theta_{\Pi_{r-1}}\Theta_{\Pi_r}\mathcal{U},\quad\dots,\quad\Theta_{\Pi_1}\cdots\Theta_{\Pi_r}\mathcal{U},
\]
each from the preceding one as in Definition~\ref{def:universal-identity}.
By the case of one mirror, applied at each step, all of them are universal identities, and they are of the first kind if $\mathcal{U}$ is.
Each step replaces every cochain by the negative of its pull-back under the reflection in the mirror concerned;
so the last triple is
\[
\left(\Theta_{\Pi_1}\cdots\Theta_{\Pi_r}P;\left((-1)^{r}\cdot\theta_{\Pi_1}^{*}\cdots\theta_{\Pi_r}^{*}\eta_a,\Theta_{\Pi_1}\cdots\Theta_{\Pi_r}P_a\right)_{a\in I};\theta_{\Pi_1}\cdots\theta_{\Pi_r}B'\right),
\]
whose elements are $P^{g}$ and $P_a^{g}$ and whose box is $gB'$.
By Definition~\ref{def:cochains},
\[
\theta_{\Pi_1}^{*}\cdots\theta_{\Pi_r}^{*}=\left(\theta_{\Pi_r}\cdots\theta_{\Pi_1}\right)^{*}=(g^{-1})^{*}=g_{*},
\]
and $(-1)^{r}=\det\vec g=\operatorname{sgn}(g)$, the linear part of a reflection having determinant $-1$.
The universal identity so obtained is therefore $\mathcal{U}^{g}$.
\end{proof}

\subsection{Slices and thin layers}\label{sec:slices-layers}

From now on the coordinate direction $e_1$ plays a distinguished role.
Unless another direction is named, slices, thick and thin layers, and the lowest and the highest plane of an element of $\mathcal{P}$ are those of Definitions~\ref{def:slices} and~\ref{def:two-layer}, which refer to the coordinate $x_1$.
A point of $\Z^3$ is written $(j,y)$, with $j\in\Z$ and $y$ a site of the slice lattice $\Z^2$;
for a link $b=(y,\nu)$ of the slice lattice, $\nu\in\{2,3\}$, the link $((j,y),\nu)$ of the plane $\{x_1=j\}$ is denoted $(j,b)$.
Pairings and norms of cochains of the slice lattice are the counting ones, as on $\Z^3$ (Definition~\ref{def:cochains}).

\begin{definition}[Slice components]\label{def:slice-components}
Let $j\in\Z$.
For $u\in C^1$ the \emph{slice components} $u_j$ and $u^{\perp}_j$ are the one-cochain and the zero-cochain of the slice lattice given by
\[
u_j(y,\nu):=u((j,y),\nu),\qquad u^{\perp}_j(y):=u((j,y),1);
\]
thus $u_j$ records the values of $u$ on the links of the plane $\{x_1=j\}$, and $u^{\perp}_j$ those on the links from the plane $\{x_1=j\}$ to the plane $\{x_1=j+1\}$.
For $\beta\in C^2$ the \emph{electric component} $\beta^{\mathrm{el}}_j$, a one-cochain of the slice lattice, and the \emph{magnetic component} $\beta^{\mathrm{mag}}_j$, a two-cochain of the slice lattice, are given by
\[
\begin{aligned}
\beta^{\mathrm{el}}_j(y,\nu)&:=\beta((j,y);1\nu),\\
\beta^{\mathrm{mag}}_j(\sigma)&:=\beta((j,y);23)\quad\text{for the plaquette }\sigma=(y;23)\text{ of the slice lattice};
\end{aligned}
\]
the plaquette $((j,y);1\nu)$ lies between the planes $\{x_1=j\}$ and $\{x_1=j+1\}$, and the plaquette $((j,y);23)$ lies in the plane $\{x_1=j\}$.
For the field $A$ and its curvature $F$ these are the components $A_j$, $a_j$, $F^{\mathrm{el}}_j$ and $F^{\mathrm{mag}}_j$ of Definition~\ref{def:slices}.
Finally, $d_{\mathrm{s}}^{*}$ denotes the adjoint of the coboundary $d_{\mathrm{s}}$ of the slice lattice on finitely supported cochains.
Both adjoints are extended to one- and two-cochains of arbitrary support:
for $u\in C^1$, $\beta\in C^2$, a site $x$ and a link $b$ of $\Z^3$,
\[
(d^*u)(x):=\langle u,d(\delta_x)\rangle,\qquad(d^*\beta)(b):=\langle\beta,d(1_b)\rangle,
\]
where the indicators $\delta_x\in C^0_{\mathrm{c}}$ and $1_b\in C^1_{\mathrm{c}}$ take the value $1$ at $x$, respectively at $b$, and the value $0$ at every other site, respectively link;
on the slice lattice $d_{\mathrm{s}}^{*}$ is extended in the same way, with its sites and links and with $d_{\mathrm{s}}$ in the place of $d$.
These pairings are finite sums.
Indeed, $d(\delta_x)$ vanishes off the six links that end at $x$, and $d(1_b)$ off the four plaquettes that contain $b$,
each of these cells carrying the coefficient $1$ or $-1$ with which the value at $x$, respectively at $b$, enters the formula for the coboundary at that cell;
on the slice lattice four links end at a site, and two plaquettes contain a link.
On finitely supported cochains the extended maps are the adjoints,
because the adjoint $d^*$ of Definition~\ref{def:cochains} satisfies $\langle d^*u,\delta_x\rangle=\langle u,d(\delta_x)\rangle$ and $\langle d^*\beta,1_b\rangle=\langle\beta,d(1_b)\rangle$,
and the adjoint $d_{\mathrm{s}}^{*}$ satisfies the same identities on the slice lattice.
By the formulas for $d$ in Definition~\ref{def:cochains} and for $d_{\mathrm{s}}$ in Definition~\ref{def:slices}, for all $u\in C^1$ and $\beta\in C^2$, every site $x$ of $\Z^3$ and every site $y$ of the slice lattice
\[
\begin{aligned}
(d^*u)(x)&=\sum_{\mu=1}^{3}\left[u_\mu(x-e_\mu)-u_\mu(x)\right],\\
(d_{\mathrm{s}}^{*}u_j)(y)&=\sum_{\nu\in\{2,3\}}\left[u_j(y-e_\nu,\nu)-u_j(y,\nu)\right],\\
(d_{\mathrm{s}}^{*}\beta^{\mathrm{mag}}_j)(y,2)&=\beta^{\mathrm{mag}}_j(y;23)-\beta^{\mathrm{mag}}_j(y-e_3;23),\\
(d_{\mathrm{s}}^{*}\beta^{\mathrm{mag}}_j)(y,3)&=\beta^{\mathrm{mag}}_j(y-e_2;23)-\beta^{\mathrm{mag}}_j(y;23),
\end{aligned}
\]
where $\beta^{\mathrm{mag}}_j(y;23)$ is the value of $\beta^{\mathrm{mag}}_j$ at the plaquette $(y;23)$.
\end{definition}

\begin{lemma}\label{lem:slice-formulas}
Let $u\in C^1$, $\beta\in C^2$ and $j\in\Z$.
\begin{enumerate}
\item\label{it:slice-coboundary}
$(du)^{\mathrm{el}}_j=u_{j+1}-u_j-d_{\mathrm{s}}u^{\perp}_j$ and $(du)^{\mathrm{mag}}_j=d_{\mathrm{s}}u_j$.
\item\label{it:slice-adjoint-two}
$(d^*\beta)_j=\beta^{\mathrm{el}}_{j-1}-\beta^{\mathrm{el}}_j+d_{\mathrm{s}}^{*}\beta^{\mathrm{mag}}_j$ and $(d^*\beta)^{\perp}_j=-d_{\mathrm{s}}^{*}\beta^{\mathrm{el}}_j$.
\item\label{it:slice-adjoint-one}
$(d^*u)(j,y)=(d_{\mathrm{s}}^{*}u_j)(y)+u^{\perp}_{j-1}(y)-u^{\perp}_j(y)$ for every site $y$ of the slice lattice.
\item\label{it:slice-one-plane}
Let $\delta\in C^1_{\mathrm{c}}$ be supported on the links of the plane $\{x_1=j\}$.
Then
\[
(d\delta)^{\mathrm{el}}_{j-1}=\delta_j,\qquad(d\delta)^{\mathrm{el}}_j=-\delta_j,\qquad(d\delta)^{\mathrm{mag}}_j=d_{\mathrm{s}}\delta_j,
\]
all other electric and magnetic components of $d\delta$ vanish,
$\|d\delta\|^2=2\|\delta_j\|^2+\|d_{\mathrm{s}}\delta_j\|^2$,
and
\[
\begin{aligned}
(d^*d\delta)_{j-1}&=-\delta_j, &\qquad (d^*d\delta)^{\perp}_{j-1}&=-d_{\mathrm{s}}^{*}\delta_j,\\
(d^*d\delta)_j&=(2+d_{\mathrm{s}}^{*}d_{\mathrm{s}})\delta_j, &\qquad (d^*d\delta)^{\perp}_j&=d_{\mathrm{s}}^{*}\delta_j,\\
(d^*d\delta)_{j+1}&=-\delta_j,
\end{aligned}
\]
while all other slice components of $d^*d\delta$ vanish.
\end{enumerate}
\end{lemma}

\begin{proof}
\eqref{it:slice-coboundary}
Let $x=(j,y)$ and $\nu\in\{2,3\}$.
By Definition~\ref{def:cochains},
\[
\begin{aligned}
(du)_{1\nu}(x)&=u_1(x)+u_\nu(x+e_1)-u_1(x+e_\nu)-u_\nu(x)\\
&=u_{j+1}(y,\nu)-u_j(y,\nu)-\left[u^{\perp}_j(y+e_\nu)-u^{\perp}_j(y)\right],
\end{aligned}
\]
and the bracket is $(d_{\mathrm{s}}u^{\perp}_j)(y,\nu)$ by Definition~\ref{def:slices}.
Likewise
\[
\begin{aligned}
(du)_{23}(x)&=u_2(x)+u_3(x+e_2)-u_2(x+e_3)-u_3(x)\\
&=u_j(y,2)+u_j(y+e_2,3)-u_j(y+e_3,2)-u_j(y,3),
\end{aligned}
\]
which is the value of $d_{\mathrm{s}}u_j$ at the plaquette $(y;23)$.

\eqref{it:slice-adjoint-two}
By Definition~\ref{def:slice-components}, at a given link either side of each of the two identities depends only on the values of $\beta$ at the plaquettes that contain this link;
it therefore suffices to treat $\beta\in C^2_{\mathrm{c}}$.
Every link of $\Z^3$ is a link $(j',b)$ of a plane $\{x_1=j'\}$ or a link $((j',y),1)$ between two consecutive planes,
and every plaquette of $\Z^3$ is of the form $((j',y);1\nu)$ or $((j',y);23)$.
Hence the pairing of two cochains on $\Z^3$ is the sum over $j'\in\Z$ of the pairings of their slice components.
Let $u\in C^1_{\mathrm{c}}$.
Part~\eqref{it:slice-coboundary}, a shift of the index in the terms that contain $u_{j'+1}$, and the adjointness of $d_{\mathrm{s}}$ and $d_{\mathrm{s}}^{*}$ give
\[
\begin{aligned}
\langle\beta,du\rangle
&=\sum_{j'\in\Z}\left(\left\langle\beta^{\mathrm{el}}_{j'},u_{j'+1}-u_{j'}-d_{\mathrm{s}}u^{\perp}_{j'}\right\rangle+\left\langle\beta^{\mathrm{mag}}_{j'},d_{\mathrm{s}}u_{j'}\right\rangle\right)\\
&=\sum_{j'\in\Z}\left\langle\beta^{\mathrm{el}}_{j'-1}-\beta^{\mathrm{el}}_{j'}+d_{\mathrm{s}}^{*}\beta^{\mathrm{mag}}_{j'},u_{j'}\right\rangle-\sum_{j'\in\Z}\left\langle d_{\mathrm{s}}^{*}\beta^{\mathrm{el}}_{j'},u^{\perp}_{j'}\right\rangle,
\end{aligned}
\]
whereas
\[
\langle\beta,du\rangle=\langle d^*\beta,u\rangle=\sum_{j'\in\Z}\left(\left\langle(d^*\beta)_{j'},u_{j'}\right\rangle+\left\langle(d^*\beta)^{\perp}_{j'},u^{\perp}_{j'}\right\rangle\right).
\]
The slice components of $u\in C^1_{\mathrm{c}}$ may be prescribed as arbitrary finitely supported cochains of the slice lattice, all but finitely many of them zero.
Comparing the two expressions therefore gives both identities, for every $j$.

\eqref{it:slice-adjoint-one}
Take $x=(j,y)$ in the formula for $d^*u$ in Definition~\ref{def:slice-components}.
The term with $\mu=1$ is $u^{\perp}_{j-1}(y)-u^{\perp}_j(y)$, and the terms with $\mu=2,3$ add up to $(d_{\mathrm{s}}^{*}u_j)(y)$.

\eqref{it:slice-one-plane}
Here $\delta_{j'}=0$ for $j'\neq j$, and $\delta^{\perp}_{j'}=0$ for every $j'$.
Part~\eqref{it:slice-coboundary} gives $(d\delta)^{\mathrm{el}}_{j'}=\delta_{j'+1}-\delta_{j'}$ and $(d\delta)^{\mathrm{mag}}_{j'}=d_{\mathrm{s}}\delta_{j'}$ for all $j'$,
which are the stated components of $d\delta$;
splitting the norm into slice components yields $\|d\delta\|^2=\|\delta_j\|^2+\|-\delta_j\|^2+\|d_{\mathrm{s}}\delta_j\|^2$.
Part~\eqref{it:slice-adjoint-two}, applied to $d\delta$, gives
\[
(d^*d\delta)_{j'}=(d\delta)^{\mathrm{el}}_{j'-1}-(d\delta)^{\mathrm{el}}_{j'}+d_{\mathrm{s}}^{*}(d\delta)^{\mathrm{mag}}_{j'},\qquad
(d^*d\delta)^{\perp}_{j'}=-d_{\mathrm{s}}^{*}(d\delta)^{\mathrm{el}}_{j'}.
\]
The first expression equals $0-\delta_j$ for $j'=j-1$, $\delta_j+\delta_j+d_{\mathrm{s}}^{*}d_{\mathrm{s}}\delta_j$ for $j'=j$, $-\delta_j-0$ for $j'=j+1$, and $0$ for all other $j'$;
the second equals $-d_{\mathrm{s}}^{*}\delta_j$ for $j'=j-1$, $d_{\mathrm{s}}^{*}\delta_j$ for $j'=j$, and $0$ otherwise.
\end{proof}

\begin{definition}\label{def:integer-box}
An \emph{integer box} is a set
\[
\mathbb{B}=[a_1,a'_1]\times[a_2,a'_2]\times[a_3,a'_3]\subset\R^3
\]
with integers $a_\mu\le a'_\mu$, $\mu=1,2,3$.
\end{definition}

\begin{lemma}[Lines and fillings]\label{lem:lines-fillings}
\begin{enumerate}
\item\label{it:integer-lines}
\emph{Lines.}
Let $V$ be a finite set of sites of $\Z^3$ such that the graph whose vertices are the sites of $V$ and whose edges are the links with both end points in $V$ is connected.
Let $\mathfrak{q}:V\to\Z$ satisfy $\sum_{y\in V}\mathfrak{q}(y)=0$, and regard $\mathfrak{q}$ as a zero-cochain on $\Z^3$ that vanishes off $V$.
Then $\mathfrak{q}=d^*\mathfrak{j}$ for an integer one-cochain $\mathfrak{j}\in C^1_{\mathrm{c}}(\Z)$ that vanishes at every link which does not join two sites of $V$.
\item\label{it:integer-filling}
\emph{Filling.}
Let $u\in C^1$ be a real one-cochain with $d^*u=0$, and let $\mathbb{B}$ be an integer box that contains every link at which $u$ does not vanish.
Then $u=d^*\beta$ for a real two-cochain $\beta\in C^2$ that vanishes at every plaquette not contained in $\mathbb{B}$.
\end{enumerate}
\end{lemma}

\begin{proof}
\eqref{it:integer-lines}
Fix a site $y_*\in V$.
For $y\in V$ choose a path in the graph from $y_*$ to $y$, that is, a finite sequence of sites of $V$ that begins at $y_*$, ends at $y$, and in which any two consecutive sites $y'$, $y''$ are the end points of a link;
let $\mathfrak{j}_y$ be the sum of the cochains $1_{(y',y'')}$ of \eqref{eq:elimination-pair} over the consecutive pairs $(y',y'')$ of this sequence.
For $y=y_*$ the sequence consisting of $y_*$ alone is admitted, and $\mathfrak{j}_{y_*}=0$.
Each $\mathfrak{j}_y$ is an integer-valued cochain supported on links with both end points in $V$,
and by \eqref{eq:elimination-pair} the sum $d^*\mathfrak{j}_y$ of the differences $\delta_{y''}-\delta_{y'}$ telescopes:
$d^*\mathfrak{j}_y=\delta_y-\delta_{y_*}$.
Put $\mathfrak{j}:=\sum_{y\in V}\mathfrak{q}(y)\cdot\mathfrak{j}_y$.
Then
\[
d^*\mathfrak{j}=\sum_{y\in V}\mathfrak{q}(y)\cdot\left(\delta_y-\delta_{y_*}\right)=\mathfrak{q}-\left(\sum_{y\in V}\mathfrak{q}(y)\right)\cdot\delta_{y_*}=\mathfrak{q}.
\]

\eqref{it:integer-filling}
Write $\mathbb{B}=\prod_\mu[a_\mu,a'_\mu]$.
Every link of a plaquette contained in $\mathbb{B}$ is contained in $\mathbb{B}$;
so $d^*$ maps a two-cochain supported on plaquettes contained in $\mathbb{B}$ to a one-cochain supported on links contained in $\mathbb{B}$.
For $a_1\le j\le a'_1$ the link $(j,b)$ is contained in $\mathbb{B}$ if and only if the link $b$ of the slice lattice is contained in the rectangle $[a_2,a'_2]\times[a_3,a'_3]$.

\emph{First sweep.}
Let $\beta'\in C^2$ be the cochain whose electric components are
\[
\beta'^{\mathrm{el}}_j(b):=-\sum_{j'=a_1}^{j}u_{j'}(b)
\]
for $a_1\le j<a'_1$ and the links $b$ of the slice lattice contained in $[a_2,a'_2]\times[a_3,a'_3]$,
and all of whose other electric components and all of whose magnetic components vanish;
if $a_1=a'_1$, then $\beta'=0$.
The plaquettes $((j,y);1\nu)$ on which $\beta'$ may differ from zero are contained in $\mathbb{B}$.
By Lemma~\ref{lem:slice-formulas}\eqref{it:slice-adjoint-two}, $(d^*\beta')_j=\beta'^{\mathrm{el}}_{j-1}-\beta'^{\mathrm{el}}_j$ for every $j$.
For $a_1\le j<a'_1$ this equals $u_j$:
at a link $b$ contained in the rectangle the two sums differ by the single term $u_j(b)$,
and at the other links of the slice lattice both sides vanish, because $u$ is supported on links contained in $\mathbb{B}$.
For $j<a_1$ both $(d^*\beta')_j$ and $u_j$ vanish.
Consequently $u':=u-d^*\beta'$ is supported on the links contained in $\mathbb{B}$,
$d^*u'=d^*u-d^*d^*\beta'=0$,
and $u'_j=0$ for every $j<a'_1$.
For $j<a'_1$ and every site $y$ of the slice lattice, Lemma~\ref{lem:slice-formulas}\eqref{it:slice-adjoint-one} therefore gives
\[
0=(d^*u')(j,y)=u'^{\perp}_{j-1}(y)-u'^{\perp}_j(y).
\]
The links from the plane $\{x_1=a_1-1\}$ to the plane $\{x_1=a_1\}$ are not contained in $\mathbb{B}$, so $u'^{\perp}_{a_1-1}=0$,
and induction over $j=a_1,a_1+1,\dots,a'_1-1$ yields $u'^{\perp}_j=0$ for these $j$;
for the remaining $j$ the cochain $u'^{\perp}_j$ vanishes because $u'$ is supported on links contained in $\mathbb{B}$.
The only links at which $u'$ can differ from zero are therefore the links of the plane $\{x_1=a'_1\}$ that are contained in $\mathbb{B}$.

\emph{Second sweep.}
Let $\beta''\in C^2$ be the cochain with the values
\[
\beta''((a'_1,y);23):=-\sum_{y'_2=a_2}^{y_2}u'((a'_1,y'_2,y_3),3)
\]
at the plaquettes $((a'_1,y);23)$ with $y=(y_2,y_3)$, $a_2\le y_2<a'_2$ and $a_3\le y_3<a'_3$, and with the value $0$ at all other plaquettes.
These plaquettes are contained in $\mathbb{B}$.
By Lemma~\ref{lem:slice-formulas}\eqref{it:slice-adjoint-two}, $d^*\beta''$ is supported on the links of the plane $\{x_1=a'_1\}$, where its slice component is $d_{\mathrm{s}}^{*}\beta''^{\mathrm{mag}}_{a'_1}$.
Put $u'':=u'-d^*\beta''$.
It is supported on links of the plane $\{x_1=a'_1\}$ contained in $\mathbb{B}$, and $d^*u''=0$.
Let $y$ be a site of the slice lattice with $y_2<a'_2$.
The formula for $d_{\mathrm{s}}^{*}$ on magnetic components in Definition~\ref{def:slice-components} gives
\[
u''((a'_1,y),3)=u'((a'_1,y),3)-\beta''((a'_1,y-e_2);23)+\beta''((a'_1,y);23).
\]
If $a_2\le y_2$ and $a_3\le y_3<a'_3$, the last two terms add up to $-u'((a'_1,y),3)$, the value of $\beta''$ at the plaquette $((a'_1,y-e_2);23)$ being $0$ when $y_2=a_2$.
Otherwise the link $((a'_1,y),3)$ is not contained in $\mathbb{B}$, so that $u'$ vanishes on it, and $\beta''$ vanishes at both plaquettes.
In both cases $u''((a'_1,y),3)=0$:
the cochain $u''$ vanishes on the links of direction $e_3$ with $x_2<a'_2$.
Next, Lemma~\ref{lem:slice-formulas}\eqref{it:slice-adjoint-one} and $d^*u''=0$ give $d_{\mathrm{s}}^{*}u''_{a'_1}=0$;
at a site $y$ with $y_2<a'_2$ the two terms of direction $e_3$ in this equation vanish, and it reads
\[
0=u''((a'_1,y-e_2),2)-u''((a'_1,y),2).
\]
The links $((a'_1,y),2)$ with $y_2<a_2$ or $y_2\ge a'_2$ are not contained in $\mathbb{B}$, so $u''$ vanishes on them;
induction over $y_2=a_2,\dots,a'_2-1$ shows that $u''$ vanishes on all links of direction $e_2$.

\emph{The remaining segment.}
Consequently $u''$ is supported on the links $((a'_1,a'_2,x_3),3)$ with $a_3\le x_3<a'_3$.
Since $u''$ vanishes on all links of direction $e_2$, the equation $d_{\mathrm{s}}^{*}u''_{a'_1}=0$ at the site $(a'_2,x_3)$ of the slice lattice reads,
by the formula for $d_{\mathrm{s}}^{*}$ on one-cochains in Definition~\ref{def:slice-components},
\[
0=u''((a'_1,a'_2,x_3-1),3)-u''((a'_1,a'_2,x_3),3).
\]
For $x_3=a_3$ the first term vanishes, and induction upwards in $x_3$ gives $u''=0$.
Therefore $u=d^*\beta$ for $\beta:=\beta'+\beta''$, which is real and vanishes at every plaquette not contained in $\mathbb{B}$, as $\beta'$ and $\beta''$ do.
\end{proof}

\emph{Bosonic elements of $\mathcal{P}$.}
Let $\delta\in C^1_{\mathrm{c}}$ satisfy $d^*\delta=0$, let $B$ be a box that contains the links in $\supp\delta$, and let $\hat A$ be a local lift on $B$ (Definition~\ref{def:integrands}).
If $\delta=0$, put $\beta:=0$.
Otherwise let $\mathbb{B}$ be the smallest integer box that contains the links in $\supp\delta$;
by Lemma~\ref{lem:lines-fillings}\eqref{it:integer-filling} there is $\beta\in C^2_{\mathrm{c}}$, supported on the plaquettes contained in $\mathbb{B}$, with $d^*\beta=\delta$.
A box is a product of three closed intervals, so $\mathbb{B}\subset B$;
hence the plaquettes in $\supp\beta$ lie in $B$, and there $d\hat A=F$.
It follows that
\[
\exp\left(i\langle\delta,\hat A\rangle\right)=\exp\left(i\langle d^*\beta,\hat A\rangle\right)=\exp\left(i\langle\beta,d\hat A\rangle\right)=\varphi_\beta(F)
\]
(Definition~\ref{def:characters}).
Now let $P\in\mathcal{P}$ (Definition~\ref{def:polynomials}) be read in $B$ and bosonic, that is, all its values are multiples of the unit (Definition~\ref{def:integrands}),
and let $P=\sum_a\lambda_a\cdot\exp(i\langle\delta_a,\hat A\rangle)\cdot\Psi_a$ be its unique representation (Definition~\ref{def:two-layer}), so that every $\delta_a$ is supported on links contained in $B$.
As $P(\hat A)$ is a scalar for every $\hat A$, every $\Psi_a$ is the empty monomial;
hence $d^*\delta_a=q\cdot\mathfrak{q}(\Psi_a)=0$, and $P=\sum_a\lambda_a\cdot\varphi_{\beta_a}(F)$ with the cochains $\beta_a$ obtained from $\delta_a$ as above.
Conversely, if $\beta\in C^2_{\mathrm{c}}$ is supported on plaquettes contained in $B$, then $\varphi_\beta(F)=\exp(i\langle d^*\beta,\hat A\rangle)$ with $d^*(d^*\beta)=0$, an element of $\mathcal{P}$ without fields.
Thus the bosonic elements of $\mathcal{P}$ are exactly the finite linear combinations of characters $\varphi_\beta$, $\beta\in C^2_{\mathrm{c}}$.

\emph{Coordinates of the neighbours.}
Let $x\in\mathbb{X}$, let $g\in W$ be the element with $x=gv_0$ (Definition~\ref{def:fermion-sites}), and let $v\in\Z^3$ and $\vec g\in\Ocub$ be its translation part and its linear part, so that $x=v+\vec g v_0$.
Then $\pi(x)=v$ and $y(x)=\vec g v_0$.
The matrix $\vec g$ is a signed permutation matrix;
hence the coordinates of $y(x)$ are $\pm v_1$, $\pm v_2$, $\pm v_3$ in some order,
and each of the three moduli $v_1>v_2>v_3$ is carried by exactly one coordinate $y_\mu(x)$, with the sign $\operatorname{sgn}y_\mu(x)$.
Since $y(x)=v_1\cdot\vec g e_1+v_2\cdot\vec g e_2+v_3\cdot\vec g e_3$, the vectors $\vec g e_1$, $\vec g e_2$, $\vec g e_3$ are the vectors $\operatorname{sgn}y_\mu(x)\cdot e_\mu$ for the coordinates $\mu$ of modulus $v_1$, $v_2$, $v_3$, respectively.
By Definition~\ref{def:chambers},
\[
s_0v_0=(v_1,v_2,-v_3),\qquad s_1v_0=(v_1,v_3,v_2),\qquad s_2v_0=(v_2,v_1,v_3),\qquad s_3v_0=(1-v_1,v_2,v_3),
\]
and $x^{(i)}=gs_iv_0=v+\vec g s_iv_0$.
For $i=0,1,2$ the reflection $s_i$ is linear and belongs to $\Ocub$, so that $\pi(x^{(i)})=v$ and $y(x^{(i)})=\vec g s_iv_0$,
where $s_0v_0=v_0-2v_3\cdot e_3$, $s_1v_0=v_0-(v_2-v_3)\cdot(e_2-e_3)$ and $s_2v_0=v_0-(v_1-v_2)\cdot(e_1-e_2)$.
For $i=3$ one has $s_3v_0=e_1+\diag(-1,1,1)\,v_0$, hence
\[
x^{(3)}=\left(v+\vec g e_1\right)+\vec g\diag(-1,1,1)\,v_0,
\]
with $v+\vec g e_1\in\Z^3$ and $\vec g\diag(-1,1,1)\in\Ocub$;
therefore $\pi(x^{(3)})=v+\vec g e_1$ and $y(x^{(3)})=\vec g\diag(-1,1,1)\,v_0=y(x)-2v_1\cdot\vec g e_1$.
Altogether,
\begin{equation}\label{eq:elimination-coordinates}
\begin{aligned}
\pi\left(x^{(0)}\right)&=\pi(x), &\qquad y\left(x^{(0)}\right)&=y(x)-2v_3\cdot\vec g e_3,\\
\pi\left(x^{(1)}\right)&=\pi(x), &\qquad y\left(x^{(1)}\right)&=y(x)-(v_2-v_3)\cdot\vec g\left(e_2-e_3\right),\\
\pi\left(x^{(2)}\right)&=\pi(x), &\qquad y\left(x^{(2)}\right)&=y(x)-(v_1-v_2)\cdot\vec g\left(e_1-e_2\right),\\
\pi\left(x^{(3)}\right)&=\pi(x)+\vec g e_1, &\qquad y\left(x^{(3)}\right)&=y(x)-2v_1\cdot\vec g e_1,
\end{aligned}
\end{equation}
in accordance with the table~\eqref{eq:chamber-neighbours} of Lemma~\ref{lem:neighbours}.
Read coordinate by coordinate, \eqref{eq:elimination-coordinates} says the following.
Along the bonds of the types $0$, $1$, $2$ the nearest photon site does not change.
For $i=0$ the vector $y(x^{(0)})$ arises from $y(x)$ by negating its coordinate of modulus $v_3$.
For $i=1$ the moduli $v_2$ and $v_3$ trade places between the two coordinates of $y(x)$ that carry them, and each of these coordinates retains its sign;
for $i=2$ the moduli $v_1$ and $v_2$ trade places in the same manner.
The bond of type $3$ alters a single coordinate of $x$, namely the coordinate $\mu$ of modulus $v_1$:
$x_\mu=\pi(x)_\mu+\operatorname{sgn}y_\mu(x)\cdot v_1$ is replaced by $\pi(x)_\mu+\operatorname{sgn}y_\mu(x)\cdot(1-v_1)$,
so that $\pi(x^{(3)})=\pi(x)+\operatorname{sgn}y_\mu(x)\cdot e_\mu$ and $y_\mu(x^{(3)})=-\operatorname{sgn}y_\mu(x)\cdot v_1$.

\emph{Thin layers.}
By Definition~\ref{def:two-layer}, the thin layers $\mathbb{X}_{j,1},\dots,\mathbb{X}_{j,6}$ of the thick layer $j$ consist of the sites $x$ with $x_1-j=-v_1,-v_2,-v_3,v_3,v_2,v_1$, respectively;
in particular each of them lies in a plane orthogonal to $e_1$.
Since $\pi(x)_1$ is an integer and $|y_1(x)|\le v_1<\tfrac12$, the integer $\pi(x)_1$ is the one nearest to $x_1$.
Hence $\mathbb{X}_{j,k}$ is the set of the sites $x$ with $\pi(x)_1=j$ and $y_1(x)=-v_1,-v_2,-v_3,v_3,v_2,v_1$ for $k=1,\dots,6$, respectively.
The thin layer that follows $\mathbb{X}_{j,k}$ is $\mathbb{X}_{j,k}^{+}=\mathbb{X}_{j,k+1}$ for $k\le5$ and $\mathbb{X}_{j,6}^{+}=\mathbb{X}_{j+1,1}$;
the one that precedes it is $\mathbb{X}_{j,k}^{-}=\mathbb{X}_{j,k-1}$ for $k\ge2$ and $\mathbb{X}_{j,1}^{-}=\mathbb{X}_{j-1,6}$.

\begin{lemma}[Thin layers]\label{lem:layer-bonds}
Let $j\in\Z$, $k\in\{1,\dots,6\}$ and $x\in\mathbb{X}_{j,k}$.
\begin{enumerate}
\item\label{it:layer-neighbours}
Of the four neighbours of $x$, two lie in $\mathbb{X}_{j,k}$ itself, while one lies in the next thin layer $\mathbb{X}_{j,k}^{+}$ and one in the preceding thin layer $\mathbb{X}_{j,k}^{-}$.
\item\label{it:layer-types}
\emph{Types.}
The bond from $x$ to its neighbour in $\mathbb{X}_{j,k}^{+}$ has the type $2$, $1$, $0$, $1$, $2$, $3$ for $k=1,2,3,4,5,6$, respectively.
Thus the bonds between $\mathbb{X}_{j,1}$ and $\mathbb{X}_{j,2}$ and between $\mathbb{X}_{j,5}$ and $\mathbb{X}_{j,6}$ have the type $2$,
those between $\mathbb{X}_{j,2}$ and $\mathbb{X}_{j,3}$ and between $\mathbb{X}_{j,4}$ and $\mathbb{X}_{j,5}$ the type $1$,
those between $\mathbb{X}_{j,3}$ and $\mathbb{X}_{j,4}$ the type $0$,
and those between $\mathbb{X}_{j,6}$ and $\mathbb{X}_{j+1,1}$ the type $3$.
The two bonds from $x$ to its neighbours in $\mathbb{X}_{j,k}$ have the types $0$ and $1$ if $k\in\{1,6\}$, the types $0$ and $3$ if $k\in\{2,5\}$, and the types $2$ and $3$ if $k\in\{3,4\}$.
\item\label{it:layer-phases}
\emph{Phases.}
Let $q>0$, and let $A$ be a real or complex one-cochain on $\Z^3$.
If $x'$ is the neighbour of $x$ of type $0$, $1$ or $2$, then $U_A(x,x')=1$.
If $k=6$ and $x'\in\mathbb{X}_{j+1,1}$ is the neighbour of $x$ in the next thin layer, then
\[
U_A(x,x')=\exp\left(iq\,a_j\left(\pi(x)_2,\pi(x)_3\right)\right),\qquad
U_A(x',x)=\exp\left(-iq\,a_j\left(\pi(x)_2,\pi(x)_3\right)\right).
\]
If $k\in\{2,3,4,5\}$ and $x'$ is the neighbour of $x$ of type $3$, then $x'\in\mathbb{X}_{j,k}$, $\pi(x')=\pi(x)\pm e_\nu$ for some $\nu\in\{2,3\}$, and
\[
U_A(x,x')=\exp\left(\pm iq\,A_j(b)\right)
\]
with the same sign, where $b$ is the link of the slice lattice that joins $\left(\pi(x)_2,\pi(x)_3\right)$ and $\left(\pi(x')_2,\pi(x')_3\right)$.
As in Proposition~\ref{prop:geometry}\eqref{it:geometry-projection}, a bond $\{x',x''\}$ is said to lie over a link of $\Z^3$ if $\pi(x')$ and $\pi(x'')$ are the two end points of this link.
Over every link $(j,b)$ of the plane $\{x_1=j\}$ there are eight bonds;
each of them has both ends in one of the thin layers $\mathbb{X}_{j,2},\dots,\mathbb{X}_{j,5}$, and each of these four thin layers contains two of them.
Over every link that joins a site of the plane $\{x_1=j\}$ to its translate by $e_1$ there are eight bonds, and each of them joins a site of $\mathbb{X}_{j,6}$ to a site of $\mathbb{X}_{j+1,1}$.
\item\label{it:layer-mirrors}
\emph{Mirrors.}
Let $c\in\R$ be such that the plane $\{x_1=c\}$ separates $\mathbb{X}_{j,k}$ from $\mathbb{X}_{j,k}^{+}$, that is, $c$ lies strictly between the values of $x_1$ on these two thin layers.
Then $\{x_1=c\}$ is a mirror if and only if either $k=3$ and $c=j$, a site plane, or $k=6$ and $c=j+\tfrac12$, a link plane.
\item\label{it:layer-displacement}
\emph{Displacement.}
For every neighbour $x'$ of $x$ the difference $\pi(x')-\pi(x)$ is $0$ or one of the vectors $\pm e_1$, $\pm e_2$, $\pm e_3$.
It equals $\pm e_1$ only if the bond $\{x,x'\}$ joins a site of $\mathbb{X}_{j',6}$ to a site of $\mathbb{X}_{j'+1,1}$ for some $j'\in\Z$,
and it equals $\pm e_\nu$ with $\nu\in\{2,3\}$ only if $\{x,x'\}$ is a bond of type $3$ inside one of the thin layers $\mathbb{X}_{j,2},\dots,\mathbb{X}_{j,5}$.
\end{enumerate}
\end{lemma}

\begin{proof}
Write $x_1=\pi(x)_1+y_1(x)$, where $\pi(x)_1=j$ and $|y_1(x)|=v_1,v_2,v_3,v_3,v_2,v_1$ for $k=1,\dots,6$.
Read in the first coordinate, \eqref{eq:elimination-coordinates} shows how $x_1$ changes along each of the four bonds at $x$;
note that the bonds of the types $0$, $1$, $2$ leave $\pi$ unchanged, and that the bond of type $3$ changes $\pi$ only in the coordinate of modulus $v_1$.

\emph{The case $k\in\{3,4\}$, in which $|y_1(x)|=v_3$.}
The bond of type $0$ reverses the sign of $y_1$: it joins $\mathbb{X}_{j,3}$ and $\mathbb{X}_{j,4}$.
The bond of type $1$ replaces $|y_1|$ by $v_2$ and keeps the sign: it leads from $\mathbb{X}_{j,3}$ to $\mathbb{X}_{j,2}$ and from $\mathbb{X}_{j,4}$ to $\mathbb{X}_{j,5}$.
The bonds of the types $2$ and $3$ change neither $y_1$ nor the first coordinate of $\pi$, since the first coordinate of $y(x)$ carries neither $v_1$ nor $v_2$;
so these two neighbours lie in $\mathbb{X}_{j,k}$.

\emph{The case $k\in\{2,5\}$, in which $|y_1(x)|=v_2$.}
The bond of type $1$ replaces $|y_1|$ by $v_3$: it leads from $\mathbb{X}_{j,2}$ to $\mathbb{X}_{j,3}$ and from $\mathbb{X}_{j,5}$ to $\mathbb{X}_{j,4}$.
The bond of type $2$ replaces $|y_1|$ by $v_1$: it leads from $\mathbb{X}_{j,2}$ to $\mathbb{X}_{j,1}$ and from $\mathbb{X}_{j,5}$ to $\mathbb{X}_{j,6}$.
The bonds of the types $0$ and $3$ do not change the first coordinate, which carries neither $v_3$ nor $v_1$.

\emph{The case $k\in\{1,6\}$, in which $|y_1(x)|=v_1$.}
The bond of type $2$ replaces $|y_1|$ by $v_2$: it leads from $\mathbb{X}_{j,1}$ to $\mathbb{X}_{j,2}$ and from $\mathbb{X}_{j,6}$ to $\mathbb{X}_{j,5}$.
For the bond of type $3$ the coordinate of modulus $v_1$ is the first one.
If $k=6$, then $y_1(x)=v_1$, and $x^{(3)}_1=j+1-v_1$, $\pi(x^{(3)})_1=j+1$, $y_1(x^{(3)})=-v_1$, so that $x^{(3)}\in\mathbb{X}_{j+1,1}$.
If $k=1$, then $y_1(x)=-v_1$, and $x^{(3)}_1=j-1+v_1$, $\pi(x^{(3)})_1=j-1$, $y_1(x^{(3)})=v_1$, so that $x^{(3)}\in\mathbb{X}_{j-1,6}$.
The bonds of the types $0$ and $1$ do not change the first coordinate, which carries neither $v_3$ nor $v_2$.

In each case one neighbour lies in $\mathbb{X}_{j,k}^{+}$, one in $\mathbb{X}_{j,k}^{-}$ and two in $\mathbb{X}_{j,k}$, and the types are those listed in \eqref{it:layer-types}.
This proves \eqref{it:layer-neighbours} and \eqref{it:layer-types}.

By~\eqref{eq:elimination-coordinates}, $\pi(x')=\pi(x)$ for the neighbours $x'$ of the types $0$, $1$, $2$,
while $\pi(x^{(3)})-\pi(x)=\operatorname{sgn}y_\mu(x)\cdot e_\mu$, $\mu$ the coordinate of modulus $v_1$;
moreover $\mu=1$ if and only if $k\in\{1,6\}$.
So $\pi(x')-\pi(x)$ is $0$ or one of the vectors $\pm e_1$, $\pm e_2$, $\pm e_3$.
If it equals $\pm e_1$, then $x'=x^{(3)}$ and $k\in\{1,6\}$, and by the case $k\in\{1,6\}$ above the bond joins $\mathbb{X}_{j,6}$ to $\mathbb{X}_{j+1,1}$ if $k=6$, and $\mathbb{X}_{j-1,6}$ to $\mathbb{X}_{j,1}$ if $k=1$.
If it equals $\pm e_\nu$ with $\nu\in\{2,3\}$, then $x'=x^{(3)}$ and $k\in\{2,\dots,5\}$, and by the cases $k\in\{2,3,4,5\}$ above the bond lies inside $\mathbb{X}_{j,k}$.
This proves \eqref{it:layer-displacement}.

For \eqref{it:layer-phases} recall that $U_A(x,x')=\exp\left(iq\,A(\pi(x),\pi(x'))\right)$ (Definition~\ref{def:chamber-operator}),
and that the one-cochain $A$, read as a function of pairs of sites of $\Z^3$, satisfies $A(v,v)=0$ and $A(v,v+e_\mu)=A_\mu(v)=-A(v+e_\mu,v)$ for $v\in\Z^3$ and $\mu\in\{1,2,3\}$ (Definition~\ref{def:cochains}).
Hence $U_A(x,x')=1$ for the neighbours of the types $0$, $1$, $2$.
Let $k=6$ and $x'=x^{(3)}$.
Then $\pi(x')=\pi(x)+e_1$, and Definition~\ref{def:slices} gives
\[
U_A(x,x')=\exp\left(iq\,A_1(\pi(x))\right)=\exp\left(iq\,a_j\left(\pi(x)_2,\pi(x)_3\right)\right),
\qquad
U_A(x',x)=\exp\left(-iq\,A_1(\pi(x))\right).
\]
Let $k\in\{2,3,4,5\}$ and $x'=x^{(3)}$.
Then $x'\in\mathbb{X}_{j,k}$ by \eqref{it:layer-types}, and the coordinate $\nu$ of modulus $v_1$ lies in $\{2,3\}$.
If $y_\nu(x)=v_1$, then $\pi(x')=\pi(x)+e_\nu$ and $U_A(x,x')=\exp\left(iq\,A_\nu(\pi(x))\right)=\exp\left(iq\,A_j(b)\right)$,
where $b$ is the link of the slice lattice from $\left(\pi(x)_2,\pi(x)_3\right)$ to $\left(\pi(x')_2,\pi(x')_3\right)$.
If $y_\nu(x)=-v_1$, then $\pi(x')=\pi(x)-e_\nu$ and $U_A(x,x')=\exp\left(-iq\,A_\nu(\pi(x'))\right)=\exp\left(-iq\,A_j(b)\right)$,
where $b$ is the link of the slice lattice from $\left(\pi(x')_2,\pi(x')_3\right)$ to $\left(\pi(x)_2,\pi(x)_3\right)$.

To count the bonds over a link, let $v\in\Z^3$ and $\nu\in\{1,2,3\}$, and consider the link from $v$ to $v+e_\nu$.
By~\eqref{eq:elimination-coordinates}, only bonds of type $3$ join different photon sites,
and a bond lies over this link exactly when one of its ends, $x'$, satisfies $\pi(x')=v$ and $y_\nu(x')=v_1$;
its other end $x''=x'^{(3)}$ then satisfies $\pi(x'')=v+e_\nu$, so that the end $x'$ is determined by the bond.
The sites $x'$ with $\pi(x')=v$ are the points $v+\vec g v_0$ with $\vec g\in\Ocub$ (Definition~\ref{def:fermion-sites}),
and $\Ocub$ consists of all signed permutation matrices;
hence $y(x')=x'-v$ takes, each exactly once, all values whose coordinates are $\pm v_1$, $\pm v_2$, $\pm v_3$ in some order.
The values with $y_\nu(x')=v_1$ carry the moduli $v_2$, $v_3$ in the two other coordinates, in either order and with either sign:
there are eight of them, and so there are eight bonds over the link.
Let $v$ lie in the plane $\{x_1=j\}$.
If $\nu\in\{2,3\}$, the first coordinate $y_1(x')$ of these eight sites takes each of the values $-v_2$, $-v_3$, $v_3$, $v_2$ exactly twice,
and both ends of each of the eight bonds lie in the same thin layer, because the coordinate of modulus $v_1$ is not the first one;
so the eight bonds lie in $\mathbb{X}_{j,2}$, $\mathbb{X}_{j,3}$, $\mathbb{X}_{j,4}$, $\mathbb{X}_{j,5}$, two in each.
If $\nu=1$, the eight ends $x'$ with $\pi(x')=v$ satisfy $y_1(x')=v_1$, so they lie in $\mathbb{X}_{j,6}$, and their neighbours of type $3$ lie in $\mathbb{X}_{j+1,1}$ by the case $k=6$ above.
This completes the proof of \eqref{it:layer-phases}.

For \eqref{it:layer-mirrors} note that the thin layers $\mathbb{X}_{j,1},\dots,\mathbb{X}_{j,6}$, $\mathbb{X}_{j+1,1}$ lie, in this order, in the planes on which $x_1$ takes the values
\[
j-v_1<j-v_2<j-v_3<j+v_3<j+v_2<j+v_1<j+1-v_1;
\]
the inequalities hold because $0<v_3<v_2<v_1<\tfrac12$.
The plane $\{x_1=c\}$ is a mirror if and only if $c\in\Adm(x_1)=\tfrac12\Z$ (Definition~\ref{def:mirrors}).
For $k\in\{1,2\}$ the open interval between the two relevant numbers lies in $\left(j-\tfrac12,j\right)$, and for $k\in\{4,5\}$ it lies in $\left(j,j+\tfrac12\right)$; it contains no point of $\tfrac12\Z$.
For $k=3$ the interval $(j-v_3,j+v_3)$ lies in $\left(j-\tfrac12,j+\tfrac12\right)$ and contains exactly one point of $\tfrac12\Z$, namely $j$.
For $k=6$ the interval $(j+v_1,j+1-v_1)$ lies in $(j,j+1)$ and contains exactly one point of $\tfrac12\Z$, namely $j+\tfrac12$.
\end{proof}

\subsection{The fermion recursion}\label{sec:fermion-recursion}

As in Definition~\ref{def:two-layer}, the thin layers $\mathbb{X}_{j,k}$, $j\in\Z$, $k=1,\dots,6$, are ordered in the direction of $e_1$.
For a thin layer $\mathbb{L}$ let $\mathbb{L}^{+}$ and $\mathbb{L}^{++}$ be the next two thin layers, and $\mathbb{L}^{-}$ and $\mathbb{L}^{--}$ the two preceding ones;
thus $\mathbb{X}_{j,k}^{+}=\mathbb{X}_{j,k+1}$ for $k\leq5$, $\mathbb{X}_{j,6}^{+}=\mathbb{X}_{j+1,1}$ and $\mathbb{X}_{j,1}^{--}=\mathbb{X}_{j-1,5}$.
Every $P\in\mathcal{P}$ has a unique representation $P=\sum_a\lambda_a\cdot\exp(i\langle\delta_a,\hat A\rangle)\cdot\Psi_a$
with pairwise different pairs $(\delta_a,\Psi_a)$ and $\lambda_a\neq0$ (Definition~\ref{def:covariant-classes}).
The fields of $P$ are the generators $\psi_\alpha(x)$ and $\bar\psi_\alpha(x)$ that are factors of one of the monomials $\Psi_a$;
a field at the site $x$ lies in the thin layer that contains $x$.
One says that $P$ has no field below $\mathbb{L}$ (above $\mathbb{L}$) if no field of $P$ lies in a thin layer that precedes $\mathbb{L}$ (follows $\mathbb{L}$).

\begin{proposition}[Schwinger--Dyson identities]\label{prop:schwinger-dyson}
Let $M\in\R$, $w_0,\dots,w_3>0$, $q>0$ and $\rho\in\R$, and assume in addition $\rho^2\neq1$ in \eqref{it:replacement-rules}--\eqref{it:mirrored-rule}.
For a photon site $y_0\in\Z^3$ let $\mathcal{P}^{+}(y_0)$ and $\mathcal{P}^{-}(y_0)$ be the covariant classes of Definition~\ref{def:covariant-classes}:
the linear spans of the elements $\exp(i\langle\delta,\hat A\rangle)\cdot\Psi$, with $\delta\in C^1_{\mathrm{c}}$ real and $\Psi$ a monomial, for which
\[
d^{*}\delta=q\cdot\left(\mathfrak{q}(\Psi)-\delta_{y_0}\right),
\qquad\text{respectively}\qquad
d^{*}\delta=q\cdot\left(\mathfrak{q}(\Psi)+\delta_{y_0}\right)
\]
(see \eqref{eq:elimination-covariance}; the condition $d^{*}\delta=q\cdot\mathfrak{q}(\Psi)$ gives $\mathcal{P}$ itself, Definition~\ref{def:polynomials}).
Their elements are odd.
The gauge transformation \eqref{eq:setting-gauge} by a function $\lambda$ multiplies the elements of $\mathcal{P}^{+}(y_0)$ by $\exp(iq\lambda(y_0))$ and those of $\mathcal{P}^{-}(y_0)$ by $\exp(-iq\lambda(y_0))$;
so they transform like $\psi(x)$, respectively like $\bar\psi(x)$, at a site $x$ with $\pi(x)=y_0$.
\begin{enumerate}
\item\label{it:equation-of-motion}
\emph{Equations of motion.}
Let $\mathbb{T}=(\Gamma,\chi)$ be a torus, $B$ a box that embeds in $\mathbb{T}$ (Definition~\ref{def:integrands}),
$x'\in\mathbb{X}$ a site such that $x'$ and its four neighbours lie in $\mathbb{X}(B)$, $\alpha\in\{1,\dots,4\}$,
and $\Phi$ a factor on $\mathbb{T}$ (Definition~\ref{def:factors}).
Let $(\bar\psi D_A)(x')$ and $(D_A\psi)(x')$ be as in Lemma~\ref{lem:torus-elements}\eqref{it:torus-derivatives}.
If $G\in\mathcal{P}^{+}(\pi(x'))$ is read in $B$ (Definition~\ref{def:enlarged-boxes}),
then $(\bar\psi D_A)_\alpha(x')\cdot G$ and $\partial G/\partial\psi_\alpha(x')$ are elements of $\mathcal{P}$ that are read in $B$, and
\[
Z_{\mathbb{T}}\left[\Phi\cdot(\bar\psi D_A)_\alpha(x')\cdot G\right]=-Z_{\mathbb{T}}\left[\Phi\cdot\frac{\partial G}{\partial\psi_\alpha(x')}\right].
\]
If $G\in\mathcal{P}^{-}(\pi(x'))$ is read in $B$,
then $(D_A\psi)_\alpha(x')\cdot G$ and $\partial G/\partial\bar\psi_\alpha(x')$ are elements of $\mathcal{P}$ that are read in $B$, and
\[
Z_{\mathbb{T}}\left[\Phi\cdot(D_A\psi)_\alpha(x')\cdot G\right]=Z_{\mathbb{T}}\left[\Phi\cdot\frac{\partial G}{\partial\bar\psi_\alpha(x')}\right].
\]
At each point of $\Conf(\Gamma)$ (Definition~\ref{def:configurations}) the same relations hold between the Berezin integrals $J_{\mathbb{T}}$ of Lemma~\ref{lem:torus-elements}\eqref{it:torus-convergence}:
\[
J_{\mathbb{T}}\left[(\bar\psi D_A)_\alpha(x')\cdot G\right]=-J_{\mathbb{T}}\left[\frac{\partial G}{\partial\psi_\alpha(x')}\right],
\qquad
J_{\mathbb{T}}\left[(D_A\psi)_\alpha(x')\cdot G\right]=J_{\mathbb{T}}\left[\frac{\partial G}{\partial\bar\psi_\alpha(x')}\right]
\]
for $G\in\mathcal{P}^{+}(\pi(x'))$, respectively $G\in\mathcal{P}^{-}(\pi(x'))$, read in $B$.
In the Grassmann algebra $\Gr_{\mathbb{T}}$ the corresponding identities hold for all elements and for complex fields.
Namely, let the left derivatives on $\Gr_{\mathbb{T}}$ be taken with respect to the free generators at the sites of a set of representatives of $\mathbb{X}_{\mathbb{T}}$ that contains $\mathbb{X}(B)$ (Definition~\ref{def:tori});
then for every $A\in C^1(\Gamma)\otimes\C$ and every $G\in\Gr_{\mathbb{T}}$, of any parity,
\[
\begin{aligned}
\int_{\mathbb{T}}\exp\left(-S_{\mathbb{T}}(A)\right)\cdot(\bar\psi D_A)_\alpha(x')\cdot G&=-\int_{\mathbb{T}}\exp\left(-S_{\mathbb{T}}(A)\right)\cdot\frac{\partial G}{\partial\psi_\alpha(x')},\\
\int_{\mathbb{T}}\exp\left(-S_{\mathbb{T}}(A)\right)\cdot(D_A\psi)_\alpha(x')\cdot G&=\int_{\mathbb{T}}\exp\left(-S_{\mathbb{T}}(A)\right)\cdot\frac{\partial G}{\partial\bar\psi_\alpha(x')}.
\end{aligned}
\]
\item\label{it:replacement-rules}
\emph{Replacement rules.}
Let $\{x,x'\}$ be a bond (Definition~\ref{def:fermion-sites}) and $\alpha'\in\{1,\dots,4\}$;
the hopping matrices $w(x,x')$ and $w(x',x)$ are invertible (Lemma~\ref{lem:hopping}).
For $G\in\mathcal{P}^{+}(\pi(x))$ put
\[
\begin{aligned}
\mathfrak{R}[\bar\psi_{\alpha'}(x);x'](G):=\sum_{\alpha=1}^{4}\biggl(&M\bar\psi_\alpha(x')\cdot G\\
&-\sum_{x''\sim x',\,x''\neq x}\ \sum_{\alpha''=1}^{4}\bar\psi_{\alpha''}(x'')\cdot w(x'',x')_{\alpha''\alpha}\cdot U_A(x'',x')\cdot G\\
&+\frac{\partial G}{\partial\psi_\alpha(x')}\biggr)\cdot\left(w(x,x')^{-1}\right)_{\alpha\alpha'}\cdot U_A(x',x),
\end{aligned}
\]
and for $G\in\mathcal{P}^{-}(\pi(x))$ put
\[
\begin{aligned}
\mathfrak{R}[\psi_{\alpha'}(x);x'](G):=\sum_{\alpha=1}^{4}U_A(x,x')\cdot\left(w(x',x)^{-1}\right)_{\alpha'\alpha}\cdot\biggl(&M\psi_\alpha(x')\cdot G\\
&-\sum_{x''\sim x',\,x''\neq x}\ \sum_{\alpha''=1}^{4}w(x',x'')_{\alpha\alpha''}\cdot U_A(x',x'')\cdot\psi_{\alpha''}(x'')\cdot G\\
&-\frac{\partial G}{\partial\bar\psi_\alpha(x')}\biggr);
\end{aligned}
\]
here $x''$ ranges over the neighbours of $x'$ other than $x$, and the phases are read on the local lift as in \eqref{eq:elimination-pair}.
These are elements of $\mathcal{P}$.
Let $\mathbb{T}$ be a torus, $B$ a box that embeds in $\mathbb{T}$ such that $x'$ and its four neighbours lie in $\mathbb{X}(B)$, and $\Phi$ a factor on $\mathbb{T}$.
If $G\in\mathcal{P}^{+}(\pi(x))$ is read in $B$,
then $\bar\psi_{\alpha'}(x)\cdot G$ and $\mathfrak{R}[\bar\psi_{\alpha'}(x);x'](G)$ are elements of $\mathcal{P}$ that are read in $B$, and
\[
Z_{\mathbb{T}}\left[\Phi\cdot\bar\psi_{\alpha'}(x)\cdot G\right]=Z_{\mathbb{T}}\left[\Phi\cdot\mathfrak{R}[\bar\psi_{\alpha'}(x);x'](G)\right].
\]
If $G\in\mathcal{P}^{-}(\pi(x))$ is read in $B$,
then $\psi_{\alpha'}(x)\cdot G$ and $\mathfrak{R}[\psi_{\alpha'}(x);x'](G)$ are elements of $\mathcal{P}$ that are read in $B$, and
\[
Z_{\mathbb{T}}\left[\Phi\cdot\psi_{\alpha'}(x)\cdot G\right]=Z_{\mathbb{T}}\left[\Phi\cdot\mathfrak{R}[\psi_{\alpha'}(x);x'](G)\right].
\]
In matrix notation, with $\bar\psi$ a row and $\psi$ a column,
with $\partial/\partial\psi(x')$ the row of the derivatives $\partial/\partial\psi_\alpha(x')$ and $\partial/\partial\bar\psi(x')$ the column of the derivatives $\partial/\partial\bar\psi_\alpha(x')$,
and with the replaced field written at the far left of the product, so that the derivatives act on the factor to its right, the two rules read
\[
\begin{aligned}
\bar\psi(x)&\longmapsto\left(M\bar\psi(x')-\sum_{x''\sim x',\,x''\neq x}\bar\psi(x'')\cdot w(x'',x')\cdot U_A(x'',x')+\frac{\partial}{\partial\psi(x')}\right)\cdot w(x,x')^{-1}\cdot U_A(x',x),\\
\psi(x)&\longmapsto U_A(x,x')\cdot w(x',x)^{-1}\cdot\left(M\psi(x')-\sum_{x''\sim x',\,x''\neq x}w(x',x'')\cdot U_A(x',x'')\cdot\psi(x'')-\frac{\partial}{\partial\bar\psi(x')}\right).
\end{aligned}
\]
\item\label{it:replacement-maps}
\emph{Upward and downward maps.}
For every thin layer $\mathbb{L}$ there is a linear map $\mathfrak{R}^{+}_{\mathbb{L}}\colon\mathcal{P}\to\mathcal{P}$,
defined on $\Z^3$ without reference to a torus or a factor, such that every $P\in\mathcal{P}$ has the following four properties.
\begin{enumerate}
\item\label{it:replacement-fields}
No field of $\mathfrak{R}^{+}_{\mathbb{L}}(P)$ lies in $\mathbb{L}$,
and each field of $\mathfrak{R}^{+}_{\mathbb{L}}(P)$ which is not among the fields of $P$ lies in $\mathbb{L}^{+}$ or in $\mathbb{L}^{++}$.
When $P$ has no field below $\mathbb{L}$, the element $\mathfrak{R}^{+}_{\mathbb{L}}(P)$ has a field neither in $\mathbb{L}$ nor below $\mathbb{L}$.
\item\label{it:replacement-links}
Every link that $\mathfrak{R}^{+}_{\mathbb{L}}(P)$ reads and $P$ does not read (Definition~\ref{def:two-layer})
joins $\pi(x)$ to $\pi(x')$ for a bond $\{x,x'\}$ of type~$3$ one of whose ends lies in $\mathbb{L}^{+}$;
thus the new phases are phases of bonds at $\mathbb{L}^{+}$.
\item\label{it:replacement-identity}
If $P$ is read in a box $B$, then $\mathfrak{R}^{+}_{\mathbb{L}}(P)$ is read in $B^{+1}$ (Definition~\ref{def:enlarged-boxes}), and
\[
Z_{\mathbb{T}}[\Phi\cdot P]=Z_{\mathbb{T}}\left[\Phi\cdot\mathfrak{R}^{+}_{\mathbb{L}}(P)\right]
\]
for every torus $\mathbb{T}$ in which $B^{+1}$ embeds and every factor $\Phi$ on $\mathbb{T}$.
\item\label{it:replacement-characters}
For every $\beta\in C^2_{\mathrm{c}}$, the character $\varphi_\beta$ (Definition~\ref{def:characters}) being regarded as the element $\exp(i\langle d^{*}\beta,\hat A\rangle)$ of $\mathcal{P}$ (Definition~\ref{def:polynomials}),
one has $\mathfrak{R}^{+}_{\mathbb{L}}(\varphi_\beta\cdot P)=\varphi_\beta\cdot\mathfrak{R}^{+}_{\mathbb{L}}(P)$;
and $\mathfrak{R}^{+}_{\mathbb{L}}(P)=P$ if $P$ has no field in $\mathbb{L}$.
\end{enumerate}
Likewise, for every thin layer $\mathbb{L}$ there is a linear map $\mathfrak{R}^{-}_{\mathbb{L}}\colon\mathcal{P}\to\mathcal{P}$,
defined on $\Z^3$ without reference to a torus or a factor,
with the properties \eqref{it:replacement-fields}--\eqref{it:replacement-characters} in which $\mathbb{L}^{+}$, $\mathbb{L}^{++}$ and ``below'' are replaced by $\mathbb{L}^{-}$, $\mathbb{L}^{--}$ and ``above''.
\item\label{it:mirrored-rule}
\emph{Mirrored rule.}
Let $\Pi$ be a mirror (Definition~\ref{def:mirrors}) and $\theta:=\theta_\Pi$, let $\mathbb{L}$ be a thin layer, and let $P\in\mathcal{P}$ be read in a box $B$.
For every torus $\mathbb{T}$ in which $\theta(B^{+1})$ embeds and every factor $\Phi$ on $\mathbb{T}$,
\[
Z_{\mathbb{T}}\left[\Phi\cdot\Theta_\Pi\left(\mathfrak{R}^{+}_{\mathbb{L}}P\right)\right]=Z_{\mathbb{T}}\left[\Phi\cdot\Theta_\Pi P\right],
\qquad
Z_{\mathbb{T}}\left[\Phi\cdot\Theta_\Pi\left(\mathfrak{R}^{-}_{\mathbb{L}}P\right)\right]=Z_{\mathbb{T}}\left[\Phi\cdot\Theta_\Pi P\right],
\]
where $\Theta_\Pi$ acts on the monomials as the reflection of Definition~\ref{def:grassmann-reflections} and on the coefficient functions as in \eqref{eq:setting-reflection-coefficients};
by Lemma~\ref{lem:reflection-identities}, the elements $\Theta_\Pi P$ and $\Theta_\Pi\left(\mathfrak{R}^{\pm}_{\mathbb{L}}P\right)$ lie in $\mathcal{P}$ and are read in $\theta(B^{+1})$.
\end{enumerate}
\end{proposition}

\begin{proof}[Proof of~\eqref{it:equation-of-motion}]
Fix a set $\mathbb{X}'$ of representatives of $\mathbb{X}_{\mathbb{T}}$ that contains $\mathbb{X}(B)$, as in Lemma~\ref{lem:torus-elements}\eqref{it:torus-derivatives}.
The generators at the sites of $\mathbb{X}'$ are free generators of $\Gr_{\mathbb{T}}$ (Definition~\ref{def:tori}),
and $\Gr(\mathbb{X}(B))$ is a subalgebra of $\Gr_{\mathbb{T}}$.
All left derivatives below are taken with respect to these free generators;
among them are $\psi_\alpha(x')$ and $\bar\psi_\alpha(x')$, since $x'\in\mathbb{X}(B)$.

\emph{The identities in $\Gr_{\mathbb{T}}$.}
Let $A\in C^1(\Gamma)\otimes\C$ and $G\in\Gr_{\mathbb{T}}$.
Each term of $S_{\mathbb{T}}(A)$ is a multiple of a product of two generators;
so $S_{\mathbb{T}}(A)$ is even and nilpotent, and $\exp\left(-S_{\mathbb{T}}(A)\right)$ is a polynomial in $S_{\mathbb{T}}(A)$, again even, which commutes with every element of $\Gr_{\mathbb{T}}$.
For a free generator $\xi$ and $G_1,G_2\in\Gr_{\mathbb{T}}$ with $G_1$ even, the left derivative obeys the product rule
\[
\frac{\partial(G_1\cdot G_2)}{\partial\xi}=\frac{\partial G_1}{\partial\xi}\cdot G_2+G_1\cdot\frac{\partial G_2}{\partial\xi}.
\]
As $S_{\mathbb{T}}(A)$ is even and commutes with $\partial S_{\mathbb{T}}(A)/\partial\xi$,
induction on $k$ gives $\partial\left(S_{\mathbb{T}}(A)^k\right)/\partial\xi=k\cdot S_{\mathbb{T}}(A)^{k-1}\cdot\partial S_{\mathbb{T}}(A)/\partial\xi$ for $k\geq1$,
and summation of the exponential series yields
\[
\frac{\partial\exp\left(-S_{\mathbb{T}}(A)\right)}{\partial\xi}=-\frac{\partial S_{\mathbb{T}}(A)}{\partial\xi}\cdot\exp\left(-S_{\mathbb{T}}(A)\right).
\]
Lemma~\ref{lem:torus-elements}\eqref{it:torus-derivatives}, which is valid for complex $A$, gives
$\partial S_{\mathbb{T}}(A)/\partial\psi_\alpha(x')=-(\bar\psi D_A)_\alpha(x')$ and $\partial S_{\mathbb{T}}(A)/\partial\bar\psi_\alpha(x')=(D_A\psi)_\alpha(x')$.
Consequently
\[
\frac{\partial\exp\left(-S_{\mathbb{T}}(A)\right)}{\partial\psi_\alpha(x')}=(\bar\psi D_A)_\alpha(x')\cdot\exp\left(-S_{\mathbb{T}}(A)\right),
\qquad
\frac{\partial\exp\left(-S_{\mathbb{T}}(A)\right)}{\partial\bar\psi_\alpha(x')}=-(D_A\psi)_\alpha(x')\cdot\exp\left(-S_{\mathbb{T}}(A)\right).
\]
The Berezin integral annihilates derivatives.
Indeed, by Definition~\ref{def:grassmann} and \eqref{eq:setting-berezin}, $\int_{\mathbb{T}}G_1$ is, up to a sign, the coefficient in $G_1\in\Gr_{\mathbb{T}}$ of the product of all free generators,
and this product contains $\xi$;
the left derivative $\partial G_1/\partial\xi$ is a combination of monomials that do not contain $\xi$,
so that $\int_{\mathbb{T}}\partial G_1/\partial\xi=0$ for every $G_1\in\Gr_{\mathbb{T}}$.
Take $G_1:=\exp\left(-S_{\mathbb{T}}(A)\right)\cdot G$.
The product rule, the two formulas just obtained and the commutativity of $\exp\left(-S_{\mathbb{T}}(A)\right)$ with every element give
\[
\begin{aligned}
0&=\int_{\mathbb{T}}\frac{\partial\left(\exp\left(-S_{\mathbb{T}}(A)\right)\cdot G\right)}{\partial\psi_\alpha(x')}\\
&=\int_{\mathbb{T}}\exp\left(-S_{\mathbb{T}}(A)\right)\cdot(\bar\psi D_A)_\alpha(x')\cdot G+\int_{\mathbb{T}}\exp\left(-S_{\mathbb{T}}(A)\right)\cdot\frac{\partial G}{\partial\psi_\alpha(x')},\\
0&=\int_{\mathbb{T}}\frac{\partial\left(\exp\left(-S_{\mathbb{T}}(A)\right)\cdot G\right)}{\partial\bar\psi_\alpha(x')}\\
&=-\int_{\mathbb{T}}\exp\left(-S_{\mathbb{T}}(A)\right)\cdot(D_A\psi)_\alpha(x')\cdot G+\int_{\mathbb{T}}\exp\left(-S_{\mathbb{T}}(A)\right)\cdot\frac{\partial G}{\partial\bar\psi_\alpha(x')},
\end{aligned}
\]
which are the identities in $\Gr_{\mathbb{T}}$.

\emph{Elements of $\mathcal{P}$.}
Let $G\in\mathcal{P}^{+}(\pi(x'))$ be read in $B$.
By Lemma~\ref{lem:torus-elements}\eqref{it:torus-derivatives},
\[
(\bar\psi D_A)_\alpha(x')\cdot G=M\bar\psi_\alpha(x')\cdot G-\sum_{x''\sim x'}\ \sum_{\alpha''=1}^{4}w(x'',x')_{\alpha''\alpha}\cdot\bar\psi_{\alpha''}(x'')\cdot U_A(x'',x')\cdot G .
\]
The first term and $\partial G/\partial\psi_\alpha(x')$ lie in $\mathcal{P}$ by \eqref{eq:elimination-rules}.
For a neighbour $x''$ of $x'$, \eqref{eq:elimination-phases}, applied to the bond $\{x',x''\}$, shows that $U_A(x'',x')\cdot G$ lies in $\mathcal{P}^{+}(\pi(x''))$;
then $\bar\psi_{\alpha''}(x'')\cdot U_A(x'',x')\cdot G$ lies in $\mathcal{P}$ by \eqref{eq:elimination-rules}.
Each of the elements obtained in this way reads at most the sites that $G$ reads, the site $x'$ and its neighbours,
and, besides the links that $G$ reads, at most the links that join $\pi(x')$ to $\pi(x'')$ for neighbours $x''$ of $x'$.
These sites lie in $\mathbb{X}(B)$.
The links lie in $B$:
the box $B$ is convex and contains the cubes $\pi(x')+[-\frac12,\frac12]^3$ and $\pi(x'')+[-\frac12,\frac12]^3$ (Definition~\ref{def:integrands}),
hence the segment from $\pi(x')$ to $\pi(x'')$.
So $(\bar\psi D_A)_\alpha(x')\cdot G$ and $\partial G/\partial\psi_\alpha(x')$ are elements of $\mathcal{P}$ read in $B$.
Let next $G\in\mathcal{P}^{-}(\pi(x'))$ be read in $B$.
Here
\[
(D_A\psi)_\alpha(x')\cdot G=M\psi_\alpha(x')\cdot G-\sum_{x''\sim x'}\ \sum_{\alpha''=1}^{4}w(x',x'')_{\alpha\alpha''}\cdot\psi_{\alpha''}(x'')\cdot U_A(x',x'')\cdot G .
\]
The first term and $\partial G/\partial\bar\psi_\alpha(x')$ lie in $\mathcal{P}$ by \eqref{eq:elimination-rules};
for a neighbour $x''$ of $x'$, \eqref{eq:elimination-phases} gives $U_A(x',x'')\cdot G\in\mathcal{P}^{-}(\pi(x''))$,
and then $\psi_{\alpha''}(x'')\cdot U_A(x',x'')\cdot G\in\mathcal{P}$ by \eqref{eq:elimination-rules}.
These elements read the same sites and links as in the first case, so that $(D_A\psi)_\alpha(x')\cdot G$ and $\partial G/\partial\bar\psi_\alpha(x')$ are read in $B$.

\emph{Pointwise and integrated identities.}
Fix a point of $\Conf(\Gamma)$, a representative $(A,z)$ of it and a local lift $\hat A$ on $B$ (Definition~\ref{def:integrands}).
An element of $\mathcal{P}$, or of one of the classes $\mathcal{P}^{\pm}(y_0)$, that is read in $B$ takes at $\hat A$ a value in $\Gr(\mathbb{X}(B))\subset\Gr_{\mathbb{T}}$.
Since $\exp(iq\hat A(b))=\exp(iqA(b))$ for every link $b$ of $B$,
the value of $(\bar\psi D_A)_\alpha(x')\cdot G$ at $\hat A$ is the product of $(\bar\psi D_A)_\alpha(x')$, formed with the representative $A$, and of the value $G(\hat A)$;
and since the left derivatives act on the generators only, the value of $\partial G/\partial\psi_\alpha(x')$ at $\hat A$ is $\partial G(\hat A)/\partial\psi_\alpha(x')$.
The same holds for $(D_A\psi)_\alpha(x')\cdot G$ and $\partial G/\partial\bar\psi_\alpha(x')$.
The identities in $\Gr_{\mathbb{T}}$, applied to the real cochain $A$ and to $G(\hat A)$,
therefore give the two identities for $J_{\mathbb{T}}$ at the chosen point, and the point was arbitrary.
Finally, by Lemma~\ref{lem:torus-elements}\eqref{it:torus-convergence}, for every $P\in\mathcal{P}$ that is read in $B$,
\[
Z_{\mathbb{T}}[\Phi\cdot P]=\int_{\Conf(\Gamma)}\exp\left(-\frac{1}{2}\|F\|^2\right)\cdot\Phi(F)\cdot J_{\mathbb{T}}[P],
\]
the integral converging absolutely.
Multiplying the identities for $J_{\mathbb{T}}$ by $\exp\left(-\frac{1}{2}\|F\|^2\right)\cdot\Phi(F)$ and integrating over $\Conf(\Gamma)$ one arrives at the identities for $Z_{\mathbb{T}}$.

\emph{Proof of \eqref{it:replacement-rules}: the rule for $\bar\psi$.}
In the definition of $(\bar\psi D_A)(x')$ in Lemma~\ref{lem:torus-elements}\eqref{it:torus-derivatives} single out the term of the neighbour $x$ of $x'$:
\[
(\bar\psi D_A)(x')=M\bar\psi(x')-\bar\psi(x)\cdot w(x,x')\cdot U_A(x,x')-\sum_{x''\sim x',\,x''\neq x}\bar\psi(x'')\cdot w(x'',x')\cdot U_A(x'',x'),
\]
an identity between rows with four entries.
Multiply it from the right by the matrix $w(x,x')^{-1}$, which exists by Lemma~\ref{lem:hopping} because $\rho^2\neq1$,
and by the scalar $U_A(x',x)=U_A(x,x')^{-1}$.
Solving for $\bar\psi(x)$ and taking the entry with the index $\alpha'$, one obtains
\[
\begin{aligned}
\bar\psi_{\alpha'}(x)=\sum_{\alpha=1}^{4}\biggl(&M\bar\psi_\alpha(x')-\sum_{x''\sim x',\,x''\neq x}\left(\bar\psi(x'')\cdot w(x'',x')\right)_\alpha\cdot U_A(x'',x')\\
&-(\bar\psi D_A)_\alpha(x')\biggr)\cdot\left(w(x,x')^{-1}\right)_{\alpha\alpha'}\cdot U_A(x',x).
\end{aligned}
\]
Now let $G\in\mathcal{P}^{+}(\pi(x))$, and put
\[
G'_\alpha:=\left(w(x,x')^{-1}\right)_{\alpha\alpha'}\cdot U_A(x',x)\cdot G,\qquad\alpha=1,\dots,4.
\]
By \eqref{eq:elimination-phases}, applied to the bond $\{x,x'\}$, every $G'_\alpha$ lies in $\mathcal{P}^{+}(\pi(x'))$.
The scalar factor $\left(w(x,x')^{-1}\right)_{\alpha\alpha'}\cdot U_A(x',x)$ commutes with the generators and with the left derivatives;
hence the definition of the replacement operator can be written as
\[
\begin{aligned}
\mathfrak{R}[\bar\psi_{\alpha'}(x);x'](G)=\sum_{\alpha=1}^{4}\biggl(&M\bar\psi_\alpha(x')\cdot G'_\alpha+\frac{\partial G'_\alpha}{\partial\psi_\alpha(x')}\\
&-\sum_{x''\sim x',\,x''\neq x}\ \sum_{\alpha''=1}^{4}w(x'',x')_{\alpha''\alpha}\cdot\bar\psi_{\alpha''}(x'')\cdot U_A(x'',x')\cdot G'_\alpha\biggr).
\end{aligned}
\]
Here $\bar\psi_\alpha(x')\cdot G'_\alpha$ and $\partial G'_\alpha/\partial\psi_\alpha(x')$ lie in $\mathcal{P}$ by \eqref{eq:elimination-rules};
$U_A(x'',x')\cdot G'_\alpha$ lies in $\mathcal{P}^{+}(\pi(x''))$ by \eqref{eq:elimination-phases}, applied to the bond $\{x',x''\}$,
and then $\bar\psi_{\alpha''}(x'')\cdot U_A(x'',x')\cdot G'_\alpha$ lies in $\mathcal{P}$ by \eqref{eq:elimination-rules}.
So $\mathfrak{R}[\bar\psi_{\alpha'}(x);x'](G)\in\mathcal{P}$, and $\bar\psi_{\alpha'}(x)\cdot G\in\mathcal{P}$ by \eqref{eq:elimination-rules}.
Multiplying the identity for $\bar\psi_{\alpha'}(x)$ from the right by $G$ and comparing with the last expression, one finds
\[
\bar\psi_{\alpha'}(x)\cdot G=\mathfrak{R}[\bar\psi_{\alpha'}(x);x'](G)-\sum_{\alpha=1}^{4}\left((\bar\psi D_A)_\alpha(x')\cdot G'_\alpha+\frac{\partial G'_\alpha}{\partial\psi_\alpha(x')}\right).
\]
Let now $\mathbb{T}$, $B$ and $\Phi$ be as in \eqref{it:replacement-rules}, and let $G$ be read in $B$.
Besides the sites and links that $G$ reads, the elements $G'_\alpha$, $\bar\psi_{\alpha'}(x)\cdot G$, $\mathfrak{R}[\bar\psi_{\alpha'}(x);x'](G)$ and all terms of the last identity read only the sites $x$, $x'$ and the neighbours of $x'$, which lie in $\mathbb{X}(B)$,
and the links that join $\pi(x')$ to $\pi(x)$ and to $\pi(x'')$, $x''\sim x'$;
these links lie in $B$ by the convexity of $B$, as in the proof of \eqref{it:equation-of-motion}.
Therefore the elements $G'_\alpha$, which lie in $\mathcal{P}^{+}(\pi(x'))$, are read in $B$,
and so are the remaining elements of the list, which lie in $\mathcal{P}$.
Item~\eqref{it:equation-of-motion}, applied to the elements $G'_\alpha\in\mathcal{P}^{+}(\pi(x'))$, gives
\[
Z_{\mathbb{T}}\left[\Phi\cdot(\bar\psi D_A)_\alpha(x')\cdot G'_\alpha\right]=-Z_{\mathbb{T}}\left[\Phi\cdot\frac{\partial G'_\alpha}{\partial\psi_\alpha(x')}\right],\qquad\alpha=1,\dots,4.
\]
The map $P\mapsto Z_{\mathbb{T}}[\Phi\cdot P]$ is linear on the elements of $\mathcal{P}$ read in $B$ (Lemma~\ref{lem:torus-elements}\eqref{it:torus-convergence});
so the sum over $\alpha$ contributes nothing, and $Z_{\mathbb{T}}\left[\Phi\cdot\bar\psi_{\alpha'}(x)\cdot G\right]=Z_{\mathbb{T}}\left[\Phi\cdot\mathfrak{R}[\bar\psi_{\alpha'}(x);x'](G)\right]$.

\emph{The rule for $\psi$.}
The definition of $(D_A\psi)(x')$ in Lemma~\ref{lem:torus-elements}\eqref{it:torus-derivatives} gives, as an identity between columns,
\[
(D_A\psi)(x')=M\psi(x')-w(x',x)\cdot U_A(x',x)\cdot\psi(x)-\sum_{x''\sim x',\,x''\neq x}w(x',x'')\cdot U_A(x',x'')\cdot\psi(x'').
\]
The matrix $w(x',x)$ is invertible by Lemma~\ref{lem:hopping}, and $U_A(x,x')=U_A(x',x)^{-1}$.
Multiplying from the left by $U_A(x,x')\cdot w(x',x)^{-1}$ and taking the entry with the index $\alpha'$ yields
\[
\begin{aligned}
\psi_{\alpha'}(x)=\sum_{\alpha=1}^{4}U_A(x,x')\cdot\left(w(x',x)^{-1}\right)_{\alpha'\alpha}\cdot\biggl(&M\psi_\alpha(x')-\sum_{x''\sim x',\,x''\neq x}\left(w(x',x'')\cdot\psi(x'')\right)_\alpha\cdot U_A(x',x'')\\
&-(D_A\psi)_\alpha(x')\biggr).
\end{aligned}
\]
Let $G\in\mathcal{P}^{-}(\pi(x))$ and put $G':=U_A(x,x')\cdot G$;
by \eqref{eq:elimination-phases}, applied to the bond $\{x,x'\}$, $G'$ lies in $\mathcal{P}^{-}(\pi(x'))$.
Moving the scalar $U_A(x,x')$ to $G$, one rewrites the definition of the replacement operator as
\[
\begin{aligned}
\mathfrak{R}[\psi_{\alpha'}(x);x'](G)=\sum_{\alpha=1}^{4}\left(w(x',x)^{-1}\right)_{\alpha'\alpha}\cdot\biggl(&M\psi_\alpha(x')\cdot G'-\frac{\partial G'}{\partial\bar\psi_\alpha(x')}\\
&-\sum_{x''\sim x',\,x''\neq x}\ \sum_{\alpha''=1}^{4}w(x',x'')_{\alpha\alpha''}\cdot\psi_{\alpha''}(x'')\cdot U_A(x',x'')\cdot G'\biggr).
\end{aligned}
\]
By \eqref{eq:elimination-rules}, and by \eqref{eq:elimination-phases} for the bonds $\{x',x''\}$, which gives $U_A(x',x'')\cdot G'\in\mathcal{P}^{-}(\pi(x''))$,
every term lies in $\mathcal{P}$; so does $\psi_{\alpha'}(x)\cdot G$.
Multiplication of the identity for $\psi_{\alpha'}(x)$ from the right by $G$ gives
\[
\psi_{\alpha'}(x)\cdot G=\mathfrak{R}[\psi_{\alpha'}(x);x'](G)+\sum_{\alpha=1}^{4}\left(w(x',x)^{-1}\right)_{\alpha'\alpha}\cdot\left(\frac{\partial G'}{\partial\bar\psi_\alpha(x')}-(D_A\psi)_\alpha(x')\cdot G'\right).
\]
Let $\mathbb{T}$, $B$ and $\Phi$ be as in \eqref{it:replacement-rules}, and let $G$ be read in $B$.
The elements $G'$, $\psi_{\alpha'}(x)\cdot G$, $\mathfrak{R}[\psi_{\alpha'}(x);x'](G)$ and all terms of this identity read, besides the sites and links that $G$ reads, only the sites $x$, $x'$ and the neighbours of $x'$, which lie in $\mathbb{X}(B)$,
and the links that join $\pi(x')$ to $\pi(x)$ and to $\pi(x'')$, $x''\sim x'$, which lie in $B$;
so they are read in $B$.
Item~\eqref{it:equation-of-motion}, applied to $G'\in\mathcal{P}^{-}(\pi(x'))$, gives
$Z_{\mathbb{T}}\left[\Phi\cdot(D_A\psi)_\alpha(x')\cdot G'\right]=Z_{\mathbb{T}}\left[\Phi\cdot\partial G'/\partial\bar\psi_\alpha(x')\right]$ for every $\alpha$,
whence the sum over $\alpha$ vanishes under $P\mapsto Z_{\mathbb{T}}[\Phi\cdot P]$ and $Z_{\mathbb{T}}\left[\Phi\cdot\psi_{\alpha'}(x)\cdot G\right]=Z_{\mathbb{T}}\left[\Phi\cdot\mathfrak{R}[\psi_{\alpha'}(x);x'](G)\right]$.

\emph{Proof of \eqref{it:replacement-maps}: construction of the upward map.}
Fix a thin layer $\mathbb{L}$.
By Definition~\ref{def:covariant-classes} the elements $\exp(i\langle\delta,\hat A\rangle)\cdot\Psi$,
where $\Psi$ is an ordered monomial and $d^{*}\delta=q\cdot\mathfrak{q}(\Psi)$, form a basis of $\mathcal{P}$.
The map $\mathfrak{R}^{+}_{\mathbb{L}}$ is defined on this basis and then extended linearly.
For a basis element let $\deg_{\mathbb{L}}\Psi$ be the number of generators of $\Psi$ at sites of $\mathbb{L}$;
the definition proceeds by induction on this number.
If $\deg_{\mathbb{L}}\Psi=0$, put
\[
\mathfrak{R}^{+}_{\mathbb{L}}\left(\exp(i\langle\delta,\hat A\rangle)\cdot\Psi\right):=\exp(i\langle\delta,\hat A\rangle)\cdot\Psi .
\]
If $\deg_{\mathbb{L}}\Psi>0$, let $\xi$ be the first generator of $\Psi$,
in the order of the generators fixed in Definition~\ref{def:covariant-classes}, that sits at a site of $\mathbb{L}$;
let $x\in\mathbb{L}$ be this site, and let $k$ be the position of $\xi$ in $\Psi$, so that $\xi$ is the $k$-th generator of $\Psi$.
Deleting $\xi$ from $\Psi$ leaves an ordered monomial $\Psi'$, and
\[
\Psi=(-1)^{k-1}\cdot\xi\cdot\Psi' ,
\]
since $\xi$ is moved to the front across $k-1$ generators.
The charge density of $\Psi'$ is $\mathfrak{q}(\Psi)+\delta_{\pi(x)}$ if $\xi=\bar\psi_{\alpha'}(x)$,
and $\mathfrak{q}(\Psi)-\delta_{\pi(x)}$ if $\xi=\psi_{\alpha'}(x)$.
Hence $d^{*}\delta=q\cdot\left(\mathfrak{q}(\Psi')-\delta_{\pi(x)}\right)$ in the first case
and $d^{*}\delta=q\cdot\left(\mathfrak{q}(\Psi')+\delta_{\pi(x)}\right)$ in the second;
that is, $\exp(i\langle\delta,\hat A\rangle)\cdot\Psi'$ lies in $\mathcal{P}^{+}(\pi(x))$, respectively in $\mathcal{P}^{-}(\pi(x))$
(Definition~\ref{def:covariant-classes}).
Let $x^{\uparrow}$ be the neighbour of $x$ in $\mathbb{L}^{+}$ (Lemma~\ref{lem:layer-bonds}\eqref{it:layer-neighbours}),
and let $\mathfrak{R}[\xi;x^{\uparrow}]$ be the replacement operator of Proposition~\ref{prop:schwinger-dyson}\eqref{it:replacement-rules} for the bond $\{x,x^{\uparrow}\}$.
Put
\begin{equation}\label{eq:elimination-recursion}
\mathfrak{R}^{+}_{\mathbb{L}}\left(\exp(i\langle\delta,\hat A\rangle)\cdot\Psi\right)
:=(-1)^{k-1}\cdot\mathfrak{R}^{+}_{\mathbb{L}}\left(\mathfrak{R}[\xi;x^{\uparrow}]\left(\exp(i\langle\delta,\hat A\rangle)\cdot\Psi'\right)\right).
\end{equation}
The right side makes sense as soon as $\mathfrak{R}^{+}_{\mathbb{L}}$ is defined on the basis elements with $\deg_{\mathbb{L}}\Psi-1$ generators in $\mathbb{L}$.
Indeed, $\mathfrak{R}[\xi;x^{\uparrow}]\left(\exp(i\langle\delta,\hat A\rangle)\cdot\Psi'\right)$ lies in $\mathcal{P}$
by Proposition~\ref{prop:schwinger-dyson}\eqref{it:replacement-rules},
and it arises from $\exp(i\langle\delta,\hat A\rangle)\cdot\Psi'$ by left multiplication with generators at $x^{\uparrow}$ or at neighbours $x''\neq x$ of $x^{\uparrow}$,
by left derivatives with respect to generators at $x^{\uparrow}$,
and by multiplication with link phases and constant matrices.
By Lemma~\ref{lem:layer-bonds}\eqref{it:layer-neighbours}, applied to $x^{\uparrow}\in\mathbb{L}^{+}$,
exactly one neighbour of $x^{\uparrow}$ lies in $(\mathbb{L}^{+})^{-}=\mathbb{L}$, namely $x$,
two neighbours lie in $\mathbb{L}^{+}$ and one lies in $\mathbb{L}^{++}$.
Since neither $x^{\uparrow}$ nor any $x''$ lies in $\mathbb{L}$, none of the operations just listed creates or removes a generator at a site of $\mathbb{L}$;
so every basis element in the expansion of $\mathfrak{R}[\xi;x^{\uparrow}]\left(\exp(i\langle\delta,\hat A\rangle)\cdot\Psi'\right)$
has exactly $\deg_{\mathbb{L}}\Psi-1$ generators at sites of $\mathbb{L}$.
By induction, \eqref{eq:elimination-recursion} defines $\mathfrak{R}^{+}_{\mathbb{L}}$ on the whole basis.
The linear map $\mathfrak{R}^{+}_{\mathbb{L}}\colon\mathcal{P}\to\mathcal{P}$ obtained in this way is defined on $\Z^3$
through the basis of $\mathcal{P}$, the order of the generators and the replacement operators,
without reference to a torus or a factor.

\emph{Fields and exponents of the image.}
Of the operations of $\mathfrak{R}[\xi;x^{\uparrow}]$ listed above, only the left multiplications create generators,
namely at $x^{\uparrow}$ and at the sites $x''$, all of which lie in $\mathbb{L}^{+}\cup\mathbb{L}^{++}$;
the left derivatives remove generators, and the phases and the constant matrices leave the generators unchanged.
As the generators of $\Psi'$ are generators of $\Psi$,
every generator of the monomial of a basis element in the expansion of $\mathfrak{R}[\xi;x^{\uparrow}]\left(\exp(i\langle\delta,\hat A\rangle)\cdot\Psi'\right)$
is a generator of $\Psi$ or sits at a site of $\mathbb{L}^{+}\cup\mathbb{L}^{++}$.
The phases that occur in $\mathfrak{R}[\xi;x^{\uparrow}]$ are $U_A(x^{\uparrow},x)$ and $U_A(x'',x^{\uparrow})$ if $\xi=\bar\psi_{\alpha'}(x)$,
and $U_A(x,x^{\uparrow})$ and $U_A(x^{\uparrow},x'')$ if $\xi=\psi_{\alpha'}(x)$:
phases of bonds with the end $x^{\uparrow}\in\mathbb{L}^{+}$.
The phase of a bond of type $0$, $1$ or $2$ equals $1$ (Lemma~\ref{lem:layer-bonds}\eqref{it:layer-phases}).
For a bond $\{x',x''\}$ of type $3$ the phase is $U_A(x',x'')=\exp\left(iq\langle 1_{(\pi(x'),\pi(x''))},\hat A\rangle\right)$ by \eqref{eq:elimination-pair},
and multiplication by it adds $q\cdot1_{(\pi(x'),\pi(x''))}$ to the exponent.
Induction on $\deg_{\mathbb{L}}\Psi$ along \eqref{eq:elimination-recursion} therefore shows that
$\mathfrak{R}^{+}_{\mathbb{L}}\left(\exp(i\langle\delta,\hat A\rangle)\cdot\Psi\right)$ is a linear combination of basis elements with two properties.
First, their monomials have no generator at a site of $\mathbb{L}$,
and every generator of one of their monomials is a generator of $\Psi$ or sits at a site of $\mathbb{L}^{+}\cup\mathbb{L}^{++}$.
Second, their exponents differ from $\delta$ by $q$ times a sum of cochains $\pm1_{(\pi(x'),\pi(x''))}$,
where $\{x',x''\}$ ranges over bonds of type $3$ with $x'\in\mathbb{L}^{+}$.
Now let $P\in\mathcal{P}$, and let $P=\sum_a\lambda_a\cdot\exp(i\langle\delta_a,\hat A\rangle)\cdot\Psi_a$ be its representation in the basis,
with pairwise different pairs $(\delta_a,\Psi_a)$ and $\lambda_a\neq0$.
The sites read by $P$ are those that carry a generator of some $\Psi_a$,
and the links read by $P$ are those of the supports of the $\delta_a$ (Definition~\ref{def:two-layer}).

\emph{Property \eqref{it:replacement-fields}.}
By linearity, every basis element in the representation of $\mathfrak{R}^{+}_{\mathbb{L}}(P)$
occurs in the expansion of $\mathfrak{R}^{+}_{\mathbb{L}}\left(\exp(i\langle\delta_a,\hat A\rangle)\cdot\Psi_a\right)$ for some $a$.
Hence $\mathfrak{R}^{+}_{\mathbb{L}}(P)$ has no field in $\mathbb{L}$,
and a field of $\mathfrak{R}^{+}_{\mathbb{L}}(P)$ that is not a field of $P$, and so is a generator of none of the $\Psi_a$,
sits at a site of $\mathbb{L}^{+}\cup\mathbb{L}^{++}$; thus it lies in $\mathbb{L}^{+}$ or in $\mathbb{L}^{++}$.
Let now $P$ have no field below $\mathbb{L}$.
Then a field of $\mathfrak{R}^{+}_{\mathbb{L}}(P)$ is either a field of $P$ outside $\mathbb{L}$, which lies above $\mathbb{L}$,
or it lies in $\mathbb{L}^{+}$ or in $\mathbb{L}^{++}$, which follow $\mathbb{L}$;
so $\mathfrak{R}^{+}_{\mathbb{L}}(P)$ has no field in or below $\mathbb{L}$.

\emph{Property \eqref{it:replacement-links}.}
A link that $\mathfrak{R}^{+}_{\mathbb{L}}(P)$ reads and $P$ does not read lies in the support of one of the cochains $1_{(\pi(x'),\pi(x''))}$ described above.
So it joins $\pi(x')$ to $\pi(x'')$ for a bond $\{x',x''\}$ of type $3$ whose end $x'$ lies in $\mathbb{L}^{+}$.

\emph{Property \eqref{it:replacement-identity}.}
Let $P$ be read in a box $B$, let $\mathbb{T}$ be a torus in which $B^{+1}$ embeds, and let $\Phi$ be a factor on $\mathbb{T}$.
Induction on $\deg_{\mathbb{L}}\Psi$ proves the following claim.
If a basis element $\exp(i\langle\delta,\hat A\rangle)\cdot\Psi$ is read in $B^{+1}$
and all generators of $\Psi$ at sites of $\mathbb{L}$ sit at sites of $\mathbb{X}(B)$,
then $\mathfrak{R}^{+}_{\mathbb{L}}\left(\exp(i\langle\delta,\hat A\rangle)\cdot\Psi\right)$ is read in $B^{+1}$, and
\[
Z_{\mathbb{T}}\left[\Phi\cdot\exp(i\langle\delta,\hat A\rangle)\cdot\Psi\right]
=Z_{\mathbb{T}}\left[\Phi\cdot\mathfrak{R}^{+}_{\mathbb{L}}\left(\exp(i\langle\delta,\hat A\rangle)\cdot\Psi\right)\right].
\]
For $\deg_{\mathbb{L}}\Psi=0$ there is nothing to prove.
Let $\deg_{\mathbb{L}}\Psi>0$, and let $\xi$, $x$, $k$, $\Psi'$ and $x^{\uparrow}$ be as in \eqref{eq:elimination-recursion}.
The site $x$ lies in $\mathbb{X}(B)$.
By \eqref{eq:elimination-neighbours}, the site $x^{\uparrow}$ and its neighbours lie in $\mathbb{X}(B^{+1})$,
and the links that join $\pi(x^{\uparrow})$ to the projections of these neighbours, where the two points differ, belong to $(B^{+1})^{(1)}$.
The element $\exp(i\langle\delta,\hat A\rangle)\cdot\Psi'$ is read in $B^{+1}$,
because it reads no site and no link that $\exp(i\langle\delta,\hat A\rangle)\cdot\Psi$ does not read.
Proposition~\ref{prop:schwinger-dyson}\eqref{it:replacement-rules}, applied with the box $B^{+1}$ and the bond $\{x,x^{\uparrow}\}$, gives
\[
Z_{\mathbb{T}}\left[\Phi\cdot\exp(i\langle\delta,\hat A\rangle)\cdot\Psi\right]
=(-1)^{k-1}\cdot Z_{\mathbb{T}}\left[\Phi\cdot\mathfrak{R}[\xi;x^{\uparrow}]\left(\exp(i\langle\delta,\hat A\rangle)\cdot\Psi'\right)\right].
\]
Each basis element in the expansion of $\mathfrak{R}[\xi;x^{\uparrow}]\left(\exp(i\langle\delta,\hat A\rangle)\cdot\Psi'\right)$ is read in $B^{+1}$:
its generators sit at sites of $\Psi'$, at $x^{\uparrow}$ or at neighbours of $x^{\uparrow}$,
and its exponent is supported on links of $\supp\delta$ and on links that join $\pi(x^{\uparrow})$ to the projection of a neighbour of $x^{\uparrow}$.
Its generators at sites of $\mathbb{L}$ are generators of $\Psi$, so they sit at sites of $\mathbb{X}(B)$, and there are $\deg_{\mathbb{L}}\Psi-1$ of them.
The induction hypothesis applies to each of these basis elements;
together with \eqref{eq:elimination-recursion} and with the linearity of $G\mapsto Z_{\mathbb{T}}[\Phi\cdot G]$ on the elements of $\mathcal{P}$ that are read in $B^{+1}$,
it yields the claim.
Every basis element in the representation of $P$ is read in $B\subset B^{+1}$ and has all its generators at sites of $\mathbb{X}(B)$.
By linearity, $\mathfrak{R}^{+}_{\mathbb{L}}(P)$ is read in $B^{+1}$, and $Z_{\mathbb{T}}[\Phi\cdot P]=Z_{\mathbb{T}}\left[\Phi\cdot\mathfrak{R}^{+}_{\mathbb{L}}(P)\right]$.

\emph{Property \eqref{it:replacement-characters}.}
Let $\beta\in C^2_{\mathrm{c}}$.
As an element of $\mathcal{P}$, the character $\varphi_\beta$ equals $\exp(i\langle d^{*}\beta,\hat A\rangle)$,
because $\langle d^{*}\beta,\hat A\rangle=\langle\beta,d\hat A\rangle$ and $d\hat A=F$ for a local lift $\hat A$ on a box that contains the support of $\beta$
(Definition~\ref{def:integrands}).
Hence $\varphi_\beta\cdot\exp(i\langle\delta,\hat A\rangle)\cdot\Psi=\exp(i\langle\delta+d^{*}\beta,\hat A\rangle)\cdot\Psi$ is again a basis element,
with the same monomial, as $d^{*}d^{*}\beta=0$.
The data $\xi$, $x$, $k$, $\Psi'$, $x^{\uparrow}$ in \eqref{eq:elimination-recursion} depend on $\Psi$ alone,
and $\mathfrak{R}[\xi;x^{\uparrow}]$ commutes with the multiplication by the scalar function $\exp(i\langle d^{*}\beta,\hat A\rangle)$,
since it consists of left multiplications by Grassmann elements with scalar coefficients and of left derivatives with respect to generators.
Induction on $\deg_{\mathbb{L}}\Psi$ gives
\[
\mathfrak{R}^{+}_{\mathbb{L}}\left(\varphi_\beta\cdot\exp(i\langle\delta,\hat A\rangle)\cdot\Psi\right)
=\varphi_\beta\cdot\mathfrak{R}^{+}_{\mathbb{L}}\left(\exp(i\langle\delta,\hat A\rangle)\cdot\Psi\right),
\]
and linearity gives $\mathfrak{R}^{+}_{\mathbb{L}}(\varphi_\beta\cdot P)=\varphi_\beta\cdot\mathfrak{R}^{+}_{\mathbb{L}}(P)$.
If $P$ has no field in $\mathbb{L}$, then $\deg_{\mathbb{L}}\Psi_a=0$ for every $a$, and $\mathfrak{R}^{+}_{\mathbb{L}}(P)=P$.

\emph{The downward map.}
The replacement operators of Proposition~\ref{prop:schwinger-dyson}\eqref{it:replacement-rules} are defined for every bond.
Define $\mathfrak{R}^{-}_{\mathbb{L}}$ by the recursion \eqref{eq:elimination-recursion},
with $\mathfrak{R}^{-}_{\mathbb{L}}$ in place of $\mathfrak{R}^{+}_{\mathbb{L}}$ and with the neighbour $x^{\downarrow}$ of $x$ in $\mathbb{L}^{-}$ in place of $x^{\uparrow}$.
By Lemma~\ref{lem:layer-bonds}\eqref{it:layer-neighbours}, applied to $x^{\downarrow}\in\mathbb{L}^{-}$,
the site $x$ is the only neighbour of $x^{\downarrow}$ in $(\mathbb{L}^{-})^{+}=\mathbb{L}$,
and the other three neighbours of $x^{\downarrow}$ lie in $\mathbb{L}^{-}$ (two of them) and in $\mathbb{L}^{--}$ (one);
the derivatives remove generators at $x^{\downarrow}\in\mathbb{L}^{-}$.
So every step of the recursion again lowers the number of generators in $\mathbb{L}$ by one,
and $\mathfrak{R}^{-}_{\mathbb{L}}\colon\mathcal{P}\to\mathcal{P}$ is a linear map defined on $\Z^3$.
The arguments given for the upward map, with $\mathbb{L}^{-}$, $\mathbb{L}^{--}$ and ``above'' in place of $\mathbb{L}^{+}$, $\mathbb{L}^{++}$ and ``below'',
prove \eqref{it:replacement-fields} and \eqref{it:replacement-links} for $\mathfrak{R}^{-}_{\mathbb{L}}$:
the new fields lie in $\mathbb{L}^{-}\cup\mathbb{L}^{--}$,
and a new link joins the projections of the two ends of a bond of type $3$ with an end in $\mathbb{L}^{-}$.
In the proof of \eqref{it:replacement-identity}, a generator at a site $x\in\mathbb{L}\cap\mathbb{X}(B)$ is now replaced through the partner $x^{\downarrow}$;
by \eqref{eq:elimination-neighbours} the site $x^{\downarrow}$ and its neighbours lie in $\mathbb{X}(B^{+1})$,
and the links that join $\pi(x^{\downarrow})$ to their projections lie in $(B^{+1})^{(1)}$,
so that Proposition~\ref{prop:schwinger-dyson}\eqref{it:replacement-rules} applies to the bond $\{x,x^{\downarrow}\}$ with the box $B^{+1}$,
and the induction is the same as before.
Property \eqref{it:replacement-characters} follows as before,
since the data of the recursion for $\mathfrak{R}^{-}_{\mathbb{L}}$ again depend on the monomial alone
and $\mathfrak{R}[\xi;x^{\downarrow}]$ commutes with the multiplication by every scalar function $\exp(i\langle d^{*}\beta,\hat A\rangle)$, $\beta\in C^2_{\mathrm{c}}$.

\emph{Proof of \eqref{it:mirrored-rule}.}
Let $\Pi$ be a mirror and $\theta:=\theta_\Pi$, let $P\in\mathcal{P}$ be read in a box $B$,
let $\mathbb{T}=(\Gamma,\chi)$ be a torus in which $\theta(B^{+1})$ embeds, and let $\Phi$ be a factor on $\mathbb{T}$.
The box $B^{+1}$ embeds in the reflected torus $\theta\mathbb{T}$, whose period lattice is $\theta\Gamma$:
the linear part $\vec\theta$ is an involution that preserves the cube $[-1,1]^3$,
so the map $\theta$ carries the sets $B^{+1}+[-1,1]^3+\vec\theta g$, $g\in\Gamma$,
onto the sets $\theta(B^{+1})+[-1,1]^3+g$, which are pairwise disjoint.
By \eqref{it:replacement-identity} the element $\mathfrak{R}^{+}_{\mathbb{L}}P$ is read in $B^{+1}$, and so is $P$.
Lemma~\ref{lem:reflection-identities}\eqref{it:reflected-functional} for $G=\mathfrak{R}^{+}_{\mathbb{L}}P$,
then \eqref{it:replacement-identity} on the torus $\theta\mathbb{T}$ with the factor $\Phi^{\theta}$ (Definition~\ref{def:factors}),
and finally Lemma~\ref{lem:reflection-identities}\eqref{it:reflected-functional} for $G=P$ give
\[
Z_{\mathbb{T}}\left[\Phi\cdot\Theta_\Pi\left(\mathfrak{R}^{+}_{\mathbb{L}}P\right)\right]
=\overline{Z_{\theta\mathbb{T}}\left[\Phi^{\theta}\cdot\mathfrak{R}^{+}_{\mathbb{L}}P\right]}
=\overline{Z_{\theta\mathbb{T}}\left[\Phi^{\theta}\cdot P\right]}
=Z_{\mathbb{T}}\left[\Phi\cdot\Theta_\Pi P\right].
\]
With \eqref{it:replacement-identity} for $\mathfrak{R}^{-}_{\mathbb{L}}$, the same three equalities hold for the downward map.
\end{proof}

\paragraph{The normal form at a link plane.}
Applied from below and from above, the maps $\mathfrak{R}^{\pm}_{\mathbb{L}}$ bring every field of an element of $\mathcal{P}$
to the two thin layers adjacent to a link plane.
The exponent of each resulting term then splits into a character and an integer exponent
carried by the two neighbouring planes and by the links between them.

\begin{proposition}[Normal form]\label{prop:normal-form}
Let $\rho\in\R$ with $\rho^2\neq1$, let $M\in\R$, $w_0,\dots,w_3>0$ and $q>0$.
Let $c=j_0+\tfrac12$ with $j_0\in\Z$, and let $P\in\mathcal{P}$.
There are a finite index set $I$ and, for $a\in I$, complex numbers $\lambda_a$, two-cochains $\beta_a\in C^2_{\mathrm{c}}$
and integer two-layer monomials $Y_a\in\mathcal{P}_1(c)$
(elements $\exp(i\langle\delta,\hat A\rangle)\cdot\Psi$ of $\mathcal{P}_1(c)$ with $\delta\in q\cdot C^1(\Z)$ and $\Psi$ a monomial;
Definition~\ref{def:two-layer}),
and there is a box $B'=B'(P,c)$, all of them defined on $\Z^3$ without reference to a torus or a factor,
such that $P$ and all $\varphi_{\beta_a}\cdot Y_a$ are read in $B'$,
every $\beta_a$ is supported on plaquettes contained in $B'$,
and the following holds.
\begin{enumerate}
\item\label{it:normal-form-identity}
For every torus $\mathbb{T}$ in which $B'$ embeds and every factor $\Phi$ on $\mathbb{T}$,
\[
Z_{\mathbb{T}}[\Phi\cdot P]=\sum_{a\in I}\lambda_a\cdot Z_{\mathbb{T}}\left[\Phi\cdot\varphi_{\beta_a}\cdot Y_a\right].
\]
\item\label{it:normal-form-mirror}
For every mirror $\Pi$, every torus $\mathbb{T}$ in which $\theta_\Pi B'$ embeds and every factor $\Phi$ on $\mathbb{T}$,
\[
Z_{\mathbb{T}}[\Phi\cdot\Theta_\Pi P]=\sum_{a\in I}\overline{\lambda_a}\cdot Z_{\mathbb{T}}\left[\Phi\cdot\Theta_\Pi\left(\varphi_{\beta_a}\cdot Y_a\right)\right].
\]
\end{enumerate}
\end{proposition}

\begin{proof}
Let $B$ be a box in which $P$ is read.

\emph{Step~1: all fields are brought to the two thin layers at $c$.}
If $P$ has no field in or below $\mathbb{X}_{j_0,5}$, put $r_1:=0$ and $P':=P$.
Otherwise let $\mathbb{L}_1$ be the lowest thin layer that contains a field of $P$, put $\mathbb{L}_{s+1}:=\mathbb{L}_s^{+}$,
let $r_1\geq1$ be the index with $\mathbb{L}_{r_1}=\mathbb{X}_{j_0,5}$, and put
\[
P':=\mathfrak{R}^{+}_{\mathbb{L}_{r_1}}\cdots\mathfrak{R}^{+}_{\mathbb{L}_2}\mathfrak{R}^{+}_{\mathbb{L}_1}P .
\]
The element $P$ has no field below $\mathbb{L}_1$.
If $\mathfrak{R}^{+}_{\mathbb{L}_s}\cdots\mathfrak{R}^{+}_{\mathbb{L}_1}P$ has no field in or below $\mathbb{L}_s$, then it has no field below $\mathbb{L}_{s+1}$,
and by Proposition~\ref{prop:schwinger-dyson}\eqref{it:replacement-fields} its image under $\mathfrak{R}^{+}_{\mathbb{L}_{s+1}}$ has no field in or below $\mathbb{L}_{s+1}$.
By induction on $s$, the element $P'$ has no field in or below $\mathbb{X}_{j_0,5}$.
Next, if $P'$ has no field in or above $\mathbb{X}_{j_0+1,2}$, put $r_2:=0$ and $P'':=P'$.
Otherwise let $\mathbb{L}'_1$ be the highest thin layer that contains a field of $P'$, put $\mathbb{L}'_{s+1}:=(\mathbb{L}'_s)^{-}$,
let $r_2\geq1$ be the index with $\mathbb{L}'_{r_2}=\mathbb{X}_{j_0+1,2}$, and put
\[
P'':=\mathfrak{R}^{-}_{\mathbb{L}'_{r_2}}\cdots\mathfrak{R}^{-}_{\mathbb{L}'_2}\mathfrak{R}^{-}_{\mathbb{L}'_1}P' .
\]
In the same way, Proposition~\ref{prop:schwinger-dyson}\eqref{it:replacement-fields} for the downward maps shows that $P''$ has no field in or above $\mathbb{X}_{j_0+1,2}$.
A field of the image under $\mathfrak{R}^{-}_{\mathbb{L}'_s}$ that is not a field of the argument lies in $(\mathbb{L}'_s)^{-}$ or in $(\mathbb{L}'_s)^{--}$,
and these thin layers are not below $\mathbb{X}_{j_0+1,2}^{--}=\mathbb{X}_{j_0,6}$;
hence $P''$, like $P'$, has no field in or below $\mathbb{X}_{j_0,5}$.
All fields of $P''$ therefore lie in the thin layers $\mathbb{X}_{j_0,6}$ and $\mathbb{X}_{j_0+1,1}$,
which contain the two ends of the bonds that cross the plane $x_1=c$.
Put $B':=B^{+(r_1+r_2)}$ (Definition~\ref{def:enlarged-boxes}).
Applied $r_1+r_2$ times, Proposition~\ref{prop:schwinger-dyson}\eqref{it:replacement-identity} shows that $P''$ is read in $B'$ and that
\[
Z_{\mathbb{T}}[\Phi\cdot P]=Z_{\mathbb{T}}[\Phi\cdot P'']
\]
for every torus $\mathbb{T}$ in which $B'$ embeds and every factor $\Phi$ on $\mathbb{T}$;
the boxes $B^{+s}$, $1\leq s\leq r_1+r_2$, of the single steps are contained in $B'$ and therefore embed in $\mathbb{T}$ as well.

\emph{Step~2: splitting off the characters.}
Let $P''=\sum_{a\in I}\lambda_a\cdot\exp(i\langle\delta_a,\hat A\rangle)\cdot\Psi_a$ be the representation of $P''$ in the basis of Definition~\ref{def:covariant-classes},
with a finite index set $I$, pairwise different pairs $(\delta_a,\Psi_a)$ and $\lambda_a\neq0$.
Fix $a\in I$.
The charge density $\mathfrak{q}(\Psi_a)$ is integer valued and vanishes off the planes $x_1=j_0$ and $x_1=j_0+1$,
because $\pi$ maps $\mathbb{X}_{j_0,6}$ into the first of these planes and $\mathbb{X}_{j_0+1,1}$ into the second.
Its sum over $\Z^3$ is zero:
summing the gauge condition $d^{*}\delta_a=q\cdot\mathfrak{q}(\Psi_a)$ (Definition~\ref{def:polynomials}) over $\Z^3$ gives zero on the left, since $\delta_a$ has finite support.

If $\delta_a=0$, then also $\mathfrak{q}(\Psi_a)=q^{-1}\cdot d^{*}\delta_a=0$, and one puts $\beta_a:=0$ and $\mathfrak{j}_a:=0$.
Let now $\delta_a\neq0$, and let $\mathbb{B}_a$ be the smallest integer box (Definition~\ref{def:integer-box})
that contains the links of $\supp\delta_a$ and the sites of $\supp\mathfrak{q}(\Psi_a)$.
Since $P''$ is read in $B'$, the box $B'$ contains these links and the unit cubes centred at these sites;
being a product of intervals, it contains $\mathbb{B}_a$.
If $\mathfrak{q}(\Psi_a)=0$, put $\mathfrak{j}_a:=0$.
If $\mathfrak{q}(\Psi_a)\neq0$, let $V_a$ be the set of the sites of $\mathbb{B}_a$ in the planes $x_1=j_0$ and $x_1=j_0+1$.
It consists of the sites of one rectangle in one of the two planes,
or of the sites of the same rectangle in both planes, joined by the links between the planes;
so the graph formed by $V_a$ and the links with both ends in $V_a$ is connected, and $V_a$ contains the support of $\mathfrak{q}(\Psi_a)$.
Lemma~\ref{lem:lines-fillings}\eqref{it:integer-lines}, applied to $V_a$ and to the restriction of $\mathfrak{q}(\Psi_a)$ to $V_a$,
yields a finitely supported integer one-cochain $\mathfrak{j}_a$ with $d^{*}\mathfrak{j}_a=\mathfrak{q}(\Psi_a)$, supported on links with both ends in $V_a$;
these are links of the planes $x_1=j_0$, $x_1=j_0+1$ and links between them, all contained in $\mathbb{B}_a$.
In both cases the real one-cochain
\[
u:=\delta_a-q\cdot\mathfrak{j}_a
\]
is supported on links contained in $\mathbb{B}_a$, and $d^{*}u=q\cdot\mathfrak{q}(\Psi_a)-q\cdot\mathfrak{q}(\Psi_a)=0$.
Lemma~\ref{lem:lines-fillings}\eqref{it:integer-filling} yields a real two-cochain $\beta_a$,
supported on plaquettes contained in $\mathbb{B}_a$, hence in $B'$, with $d^{*}\beta_a=u$.

For every $a\in I$ one has $\delta_a=d^{*}\beta_a+q\cdot\mathfrak{j}_a$.
Let $\hat A$ be a local lift on $B'$ (Definition~\ref{def:integrands}).
As $\beta_a$ is supported on plaquettes of $B'$, where $d\hat A=F$,
\[
\begin{aligned}
\exp(i\langle\delta_a,\hat A\rangle)
&=\exp(i\langle d^{*}\beta_a,\hat A\rangle)\cdot\exp\left(iq\langle\mathfrak{j}_a,\hat A\rangle\right)
=\exp(i\langle\beta_a,d\hat A\rangle)\cdot\exp\left(iq\langle\mathfrak{j}_a,\hat A\rangle\right)\\
&=\varphi_{\beta_a}(F)\cdot\exp\left(iq\langle\mathfrak{j}_a,\hat A\rangle\right).
\end{aligned}
\]
Put
\[
Y_a:=\exp\left(iq\langle\mathfrak{j}_a,\hat A\rangle\right)\cdot\Psi_a .
\]
It is gauge invariant, since $d^{*}(q\cdot\mathfrak{j}_a)=q\cdot\mathfrak{q}(\Psi_a)$;
its exponent $q\cdot\mathfrak{j}_a$ lies in $q\cdot C^1(\Z)$ and is supported on links of the planes $x_1=j_0$, $x_1=j_0+1$ and on links between them;
its fields lie in $\mathbb{X}_{j_0,6}$ and $\mathbb{X}_{j_0+1,1}$.
Thus $Y_a$ is an integer two-layer monomial in $\mathcal{P}_1(c)$, and $\exp(i\langle\delta_a,\hat A\rangle)\cdot\Psi_a=\varphi_{\beta_a}\cdot Y_a$.
Consequently
\[
P''=\sum_{a\in I}\lambda_a\cdot\varphi_{\beta_a}\cdot Y_a ,
\]
every term $\varphi_{\beta_a}\cdot Y_a$ is read in $B'$, as $P''$ is, and $P$ is read in $B\subset B'$.
The numbers $\lambda_a$, the cochains $\beta_a$, the monomials $Y_a$ and the box $B'$ arise from $P$ and $c$ by constructions on $\Z^3$ alone.
Step~1 and the linearity of $G\mapsto Z_{\mathbb{T}}[\Phi\cdot G]$ now give \eqref{it:normal-form-identity}.

\emph{Proof of \eqref{it:normal-form-mirror}.}
Let $\Pi$ be a mirror and $\theta:=\theta_\Pi$, let $\mathbb{T}=(\Gamma,\chi)$ be a torus in which $\theta B'$ embeds, and let $\Phi$ be a factor on $\mathbb{T}$.
The box $B'$ embeds in $\theta\mathbb{T}$,
since $\theta$ carries the sets $B'+[-1,1]^3+\vec\theta g$, $g\in\Gamma$, onto the pairwise disjoint sets $\theta B'+[-1,1]^3+g$;
and $\Phi^{\theta}$ is a factor on $\theta\mathbb{T}$ (Definition~\ref{def:factors}).
Lemma~\ref{lem:reflection-identities}\eqref{it:reflected-functional} for $G=P$,
then \eqref{it:normal-form-identity} on the torus $\theta\mathbb{T}$ with the factor $\Phi^{\theta}$,
and then the same lemma for the elements $G=\varphi_{\beta_a}\cdot Y_a$, which are read in $B'$, give
\[
\begin{aligned}
Z_{\mathbb{T}}[\Phi\cdot\Theta_\Pi P]
&=\overline{Z_{\theta\mathbb{T}}\left[\Phi^{\theta}\cdot P\right]}
=\sum_{a\in I}\overline{\lambda_a}\cdot\overline{Z_{\theta\mathbb{T}}\left[\Phi^{\theta}\cdot\varphi_{\beta_a}\cdot Y_a\right]}\\
&=\sum_{a\in I}\overline{\lambda_a}\cdot Z_{\mathbb{T}}\left[\Phi\cdot\Theta_\Pi\left(\varphi_{\beta_a}\cdot Y_a\right)\right].
\end{aligned}
\]
\end{proof}

\begin{remark}\label{rem:one-link-plane}
In the normal form every charge is located at an end of a bond that crosses the single link plane $x_1=c$:
the fields of each $Y_a$ lie in $\mathbb{X}_{j_0,6}$ and $\mathbb{X}_{j_0+1,1}$,
and its exponent is carried by the planes $x_1=j_0$, $x_1=j_0+1$ and by the links between them.
Consequently, for $0<\rho<1$, the trace formula for the slab operators meets neither static charges
(time-like Wilson lines joining different planes)
nor odd Fock operators placed at different planes.
In this range each $Y_a$ is the integrand of a single local step datum (Proposition~\ref{prop:two-layer-insertions}(a)),
and the remaining character $\varphi_{\beta_a}$ is absorbed by Proposition~\ref{prop:bounded-elements}(b) into the bound of $Y_a$.
\end{remark}

\subsection{An a priori bound}\label{sec:a-priori-bound}

Throughout this subsection $c=\tfrac12$, so that $j_0=0$:
the two-layer elements of $\mathcal{P}_1(c)$ are formed with the planes $\{x_1=0\}$ and $\{x_1=1\}$
and with the fields of the thin layers $\mathbb{X}_{0,6}$ and $\mathbb{X}_{1,1}$ (Definition~\ref{def:two-layer}).
Any other coordinate link plane could take the place of this one.

\begin{definition}\label{def:element-constant}
Let $0<\rho<1$, $w_0,\dots,w_3>0$, $M\in\R$ and $q>0$, and let $P\in\mathcal{P}$ (Definition~\ref{def:polynomials}).
Let the finite index set $I$, the numbers $\lambda_a\in\C$, the two-cochains $\beta_a\in C^2_{\mathrm{c}}$,
the integer two-layer monomials $Y_a\in\mathcal{P}_1(c)$, $a\in I$, and the box $B'(P,c)$
constitute the normal form of $P$ at the plane $c=\tfrac12$ given by Proposition~\ref{prop:normal-form}.
For an index $a$ with $Y_a\neq1$ write $Y_a=\exp(i\langle\delta,\hat A\rangle)\cdot\Psi$, with the exponent $(\delta_0,\delta_1,\delta_\perp)$;
let $S_a$ be the set of the links $b$ of $\Z^2$ with $\delta_0(b)\neq0$ or $\delta_1(b)\neq0$,
and let $\Lambda_{Y_a}$ be the smallest set of sites that Definition~\ref{def:two-layer} admits as the set of the sites in play of $Y_a$,
namely the set that consists of the end points of the links of $S_a$, of the sites $y$ with $\delta_\perp(y)\neq0$,
and of the sites $y$ for which $\Psi$ contains a field at a site $x$ with $\pi(x)\in\{(0,y),(1,y)\}$.
For an index $a$ with $Y_a=1$ put $S_a:=\emptyset$;
no set of sites and no box is attached to $Y_a$ in this case, and $X'_{Y_a}=X'_1=1$ (Definition~\ref{def:two-layer}).
The constant of $P$ is
\[
C_P:=\sum_{a\in I}|\lambda_a|\cdot2^{\frac{\#|S_a|}{2}}\cdot\left\|X'_{Y_a}\right\|_\# .
\]
By Proposition~\ref{prop:two-layer-insertions}(b) each norm $\|X'_{Y_a}\|_\#$ with $Y_a\neq1$ is finite, and so $C_P<\infty$.
The number $C_P$ depends on $P$ alone, in the following sense.
It is determined by $P$,
by the total order of the generators fixed in Definition~\ref{def:covariant-classes},
by the choice of the integer one-cochains made in the construction of the normal form,
and by the parameters $\rho$, $w_0,\dots,w_3$, $M$, $q$ and $v_0$;
it does not depend on a two-cochain $\beta$, a box $B_0$, data $\mathfrak{d}$, a mirror $\Pi$,
the strength $t$, the index $i$ or the sequence of sides $\mathfrak{n}$.
\end{definition}

\begin{proposition}\label{prop:a-priori}
Let $0<\rho<1$, $w_0,\dots,w_3>0$, $M\in\R$ and $q>0$.
Let $\mathfrak{n}=(n_i)$ be a sequence of positive integers with $n_i\to\infty$ that satisfies~\eqref{cond:divisibility},
and let $\mathfrak{d}=(c_0,c_+,c_3;h;O)$ be data (Definition~\ref{def:data}),
with the perturbed functionals $\mu^t_i$ on the adapted tori $D(n_i)$
and the amplification factors $\gamma_{t,i}(k)$ (Definition~\ref{def:amplification-factor}).
Let $P\in\mathcal{P}$, let $C_P$ be its constant (Definition~\ref{def:element-constant}), and let $B_0$ be a box.
\begin{enumerate}
\item\label{it:a-priori-element}
There are $k_P=k_P(P,B_0,\mathfrak{d})\in\N$ and an index $i_0=i_0(P,B_0,\mathfrak{d})$ such that
\[
\left|\mu^t_i\left(\varphi_\beta\cdot P\right)\right|\le C_P\cdot\gamma_{t,i}(k_P)
\]
for every $\beta\in C^2_{\mathrm{c}}$ with $\supp\beta\subset B_0$, every $t\ge0$ and every $i\ge i_0$.
\item\label{it:a-priori-mirror}
For every mirror $\Pi$ there are $k_P=k_P(P,B_0,\mathfrak{d},\Pi)\in\N$ and an index $i_0=i_0(P,B_0,\mathfrak{d},\Pi)$ such that
\[
\left|\mu^t_i\left(\varphi_\beta\cdot\Theta_\Pi P\right)\right|\le C_P\cdot\gamma_{t,i}(k_P)
\]
for every $\beta\in C^2_{\mathrm{c}}$ with $\supp\beta\subset B_0$, every $t\ge0$ and every $i\ge i_0$;
the constant $C_P$ is the same as in~\eqref{it:a-priori-element}.
\end{enumerate}
\end{proposition}

\begin{proof}
Let $I$, $\lambda_a$, $\beta_a$, $Y_a$, $S_a$, $\Lambda_{Y_a}$ and $B'(P,c)$ be as in Definition~\ref{def:element-constant}, and put
\[
C_a:=2^{\frac{\#|S_a|}{2}}\cdot\left\|X'_{Y_a}\right\|_\#\qquad(a\in I).
\]
Then $C_P=\sum_{a\in I}|\lambda_a|\cdot C_a$, and $C_a=1$ for every $a$ with $Y_a=1$.
Each $\beta_a$ is supported on plaquettes contained in $B'(P,c)$ (Proposition~\ref{prop:normal-form}).

\emph{The denominator.}
Let $i$ be an index for which $D(n_i)$ carries $\mathcal{Q}_{\mathfrak{d}}$, and let $t\ge0$.
On $D(n_i)$ the folded array satisfies $|O_{\mathfrak{d}}|\le\frac{n_i^3}{|Q_0|}\cdot|O|_1$ (Definition~\ref{def:arrays}),
and Lemma~\ref{lem:perturbed-denominators} together with Theorem~\ref{thm:partition-positivity} yields
\begin{equation}\label{eq:elimination-denominator}
Z_{D(n_i)}\left[\exp(t\cdot O_{\mathfrak{d}})\right]\ge\exp\left(-t\cdot|O|_1\cdot\frac{n_i^3}{|Q_0|}\right)\cdot Z_{D(n_i)}[1]>0 .
\end{equation}

\emph{Proof of~\eqref{it:a-priori-element}.}
Let $B_P$ be a box that contains $B_0$, $B'(P,c)$
and the boxes $B^{\#}(c,S_a,\Lambda_{Y_a})$ of the indices $a$ with $Y_a\neq1$ (Definition~\ref{def:insertions});
let $r_P$ be its radius, and put
\[
k_P:=k(\mathfrak{d},r_P)
\]
(Lemma~\ref{lem:coarse-structures}\eqref{it:coarse-widths}).
Choose an index $i_0$ such that, for every $i\ge i_0$, the box $B_P$ embeds in $D(n_i)$,
the torus $D(n_i)$ carries $\mathcal{Q}_{\mathfrak{d}}$, and $i\ge i_0(\mathfrak{d},r_P)$;
by Lemma~\ref{lem:coarse-structures}\eqref{it:coarse-tori} the index $i_0(\mathfrak{d},r_P)$ itself has these properties.
The box $B_P$, the integer $k_P$ and the index $i_0$ are thus chosen in terms of $P$, $B_0$ and $\mathfrak{d}$ only.

Now fix $i\ge i_0$, $t\ge0$ and $\beta\in C^2_{\mathrm{c}}$ with $\supp\beta\subset B_0$.
Being contained in $B_P$, the box $B'(P,c)$ embeds in $D(n_i)$ (Definition~\ref{def:integrands}).
As $P$ is read in $B'(P,c)$, the product $\varphi_\beta\cdot P$ embeds in $D(n_i)$,
and $\mu^t_i(\varphi_\beta\cdot P)$ is defined (Definition~\ref{def:data}).
The function
\[
\Phi:=\exp(t\cdot O_{\mathfrak{d}})\cdot\varphi_\beta
\]
is a factor on $D(n_i)$ (Definition~\ref{def:factors}).
Indeed, $O_{\mathfrak{d}}$ is a finite sum of trigonometric polynomials on $D(n_i)$, hence a bounded Borel function there,
and $\varphi_\beta$, evaluated at the periodic extension $\widetilde F$, has modulus one;
since $\supp\beta$ lies in the box $B_P$, which embeds in $D(n_i)$,
the function $\varphi_\beta(\widetilde F)$ is the character of $D(n_i)$ given by the periodisation of $\beta$ (Lemma~\ref{lem:curvature-characters}(b)).
Proposition~\ref{prop:normal-form}\eqref{it:normal-form-identity}, applied to the torus $D(n_i)$, in which $B'(P,c)$ embeds, and to the factor $\Phi$, gives
\[
Z_{D(n_i)}\left[\exp(t\cdot O_{\mathfrak{d}})\cdot\varphi_\beta\cdot P\right]
=\sum_{a\in I}\lambda_a\cdot Z_{D(n_i)}\left[\exp(t\cdot O_{\mathfrak{d}})\cdot\varphi_\beta\cdot\varphi_{\beta_a}\cdot Y_a\right].
\]
In the element $\varphi_{\beta_a}\cdot Y_a$ of $\mathcal{P}$ the character $\varphi_{\beta_a}$ is the function $\varphi_{\beta_a}(\widetilde F)$,
because $\beta_a$ is supported on plaquettes of $B'(P,c)$, on which $d\hat A=F$ for a local lift $\hat A$ (Definition~\ref{def:integrands});
consequently $\varphi_\beta\cdot\varphi_{\beta_a}=\varphi_{\beta+\beta_a}$.
Division by $Z_{D(n_i)}[\exp(t\cdot O_{\mathfrak{d}})]$, which is positive by~\eqref{eq:elimination-denominator}, leads to
\[
\mu^t_i\left(\varphi_\beta\cdot P\right)=\sum_{a\in I}\lambda_a\cdot\mu^t_i\left(\varphi_{\beta+\beta_a}\cdot Y_a\right).
\]

Each term on the right is estimated by amplification.
Consider first an index $a$ with $Y_a\neq1$.
By Proposition~\ref{prop:two-layer-insertions}(b), the monomial $Y_a$ is bounded by $C_a$ with the box $B^{\#}(c,S_a,\Lambda_{Y_a})$
(Definition~\ref{def:bounded-elements}).
The box $B_P$ contains this box and the support of $\beta+\beta_a$,
which lies in $\supp\beta\cup\supp\beta_a\subset B_0\cup B'(P,c)$;
by Proposition~\ref{prop:bounded-elements}(b), therefore, $\varphi_{\beta+\beta_a}\cdot Y_a$ is bounded by $C_a$ with the box $B_P$.
For an index $a$ with $Y_a=1$ the element $\varphi_{\beta+\beta_a}\cdot Y_a=\varphi_{\beta+\beta_a}$ is a bosonic integrand in $\mathcal{G}(B_P)$,
a Borel function of modulus one of the $F_\sigma$ with $\sigma\in B_P^{(2)}$;
by Proposition~\ref{prop:bounded-elements}(c) it is bounded by $1=C_a$ with the box $B_P$.
In both cases Theorem~\ref{thm:amplification}, applied with the box $B_P$ of radius $r_P$, gives
\[
\left|\mu^t_i\left(\varphi_{\beta+\beta_a}\cdot Y_a\right)\right|\le C_a\cdot\gamma_{t,i}(k_P),
\]
because $k_P=k(\mathfrak{d},r_P)$ and $i\ge i_0(\mathfrak{d},r_P)$.
Summation over $a\in I$ yields
\[
\left|\mu^t_i\left(\varphi_\beta\cdot P\right)\right|
\le\sum_{a\in I}|\lambda_a|\cdot C_a\cdot\gamma_{t,i}(k_P)
=C_P\cdot\gamma_{t,i}(k_P).
\]

\emph{Proof of~\eqref{it:a-priori-mirror}.}
Let $\Pi$ be a mirror, and put $G_a:=\varphi_{\beta_a}\cdot Y_a$ for $a\in I$.
Let $B''$ be a box that contains $B'(P,c)$ and the boxes $B^{\#}(c,S_a,\Lambda_{Y_a})$ of the indices $a$ with $Y_a\neq1$.
Every $G_a$ is bounded by $C_a$ with the box $B''$.
For $Y_a\neq1$ this follows from Proposition~\ref{prop:two-layer-insertions}(b) and Proposition~\ref{prop:bounded-elements}(b),
since $B''$ contains $B^{\#}(c,S_a,\Lambda_{Y_a})$ and $\supp\beta_a\subset B'(P,c)$;
for $Y_a=1$ the element $G_a=\varphi_{\beta_a}$ is a bosonic integrand in $\mathcal{G}(B'')$ of modulus one,
and Proposition~\ref{prop:bounded-elements}(c) applies.
By~\eqref{eq:setting-group-action} one has $\Theta_\Pi G_a=G_a^{\theta_\Pi}$,
and Proposition~\ref{prop:bounded-elements}(a) shows that $\Theta_\Pi G_a$ is bounded by $C_a$ with the box $\theta_\Pi B''$.
Let $B_\Pi$ be a box that contains $\theta_\Pi B''$ and $B_0$, let $r_\Pi$ be its radius, and put
\[
k_P:=k(\mathfrak{d},r_\Pi).
\]
For every $\beta\in C^2_{\mathrm{c}}$ with $\supp\beta\subset B_0$ the element $\varphi_\beta\cdot\Theta_\Pi G_a$ is then bounded by $C_a$ with the box $B_\Pi$,
by Proposition~\ref{prop:bounded-elements}(b).
Let $i_0$ be an index such that every $i\ge i_0$ has three properties:
the box $B_\Pi$ embeds in $D(n_i)$, the torus $D(n_i)$ carries $\mathcal{Q}_{\mathfrak{d}}$, and $i\ge i_0(\mathfrak{d},r_\Pi)$;
the index $i_0(\mathfrak{d},r_\Pi)$ is one such choice, by Lemma~\ref{lem:coarse-structures}\eqref{it:coarse-tori}.
Here $B_\Pi$, $k_P$ and $i_0$ are chosen in terms of $P$, $B_0$, $\mathfrak{d}$ and $\Pi$.

Fix $i\ge i_0$, $t\ge0$ and $\beta\in C^2_{\mathrm{c}}$ with $\supp\beta\subset B_0$.
The box $\theta_\Pi B'(P,c)$ lies in $B_\Pi$ and therefore embeds in $D(n_i)$;
the element $\Theta_\Pi P$ belongs to $\mathcal{P}$ and is read in $\theta_\Pi B'(P,c)$ (Lemma~\ref{lem:reflection-identities}),
so that $\mu^t_i(\varphi_\beta\cdot\Theta_\Pi P)$ is defined.
The function $\Phi:=\exp(t\cdot O_{\mathfrak{d}})\cdot\varphi_\beta$ is again a bounded Borel function of $F$ on $D(n_i)$, that is, a factor on $D(n_i)$.
Proposition~\ref{prop:normal-form}\eqref{it:normal-form-mirror}, for the mirror $\Pi$, the torus $D(n_i)$ and this factor,
divided by the positive number $Z_{D(n_i)}[\exp(t\cdot O_{\mathfrak{d}})]$ of~\eqref{eq:elimination-denominator}, gives
\[
\mu^t_i\left(\varphi_\beta\cdot\Theta_\Pi P\right)=\sum_{a\in I}\overline{\lambda_a}\cdot\mu^t_i\left(\varphi_\beta\cdot\Theta_\Pi G_a\right).
\]
For each $a\in I$, Theorem~\ref{thm:amplification} with the box $B_\Pi$ of radius $r_\Pi$ bounds the modulus of the term $\mu^t_i(\varphi_\beta\cdot\Theta_\Pi G_a)$ by $C_a\cdot\gamma_{t,i}(k_P)$,
since $k_P=k(\mathfrak{d},r_\Pi)$ and $i\ge i_0(\mathfrak{d},r_\Pi)$.
As $|\overline{\lambda_a}|=|\lambda_a|$, summation over $a$ gives
\[
\left|\mu^t_i\left(\varphi_\beta\cdot\Theta_\Pi P\right)\right|
\le\sum_{a\in I}|\lambda_a|\cdot C_a\cdot\gamma_{t,i}(k_P)
=C_P\cdot\gamma_{t,i}(k_P),
\]
with the same constant $C_P$ as in~\eqref{it:a-priori-element}.
\end{proof}

\subsection{Shifting the contour}\label{sec:contour-shift}

\begin{definition}[Shifted elements]\label{def:shifted-elements}
Let $S$ be a finite set of links of $\Z^3$ and let $\delta\in\R^S$, with the value $\delta_b$ at the link $b\in S$.
Extended by zero, $\delta$ becomes a real one-cochain of finite support on $\Z^3$, which is again denoted by $\delta$.

For an entire factor $\Phi$ on a torus $\mathbb{T}$ (Definition~\ref{def:factors}) put
\[
\Phi^{(\delta)}:=\Phi^{[d\delta]},\qquad\text{that is,}\qquad\Phi^{(\delta)}(F)=\Phi(F-i\,d\delta),
\]
where on $\mathbb{T}$ the cochain $d\delta$ is identified with its periodisation (Definition~\ref{def:factors}).

Let $P\in\mathcal{P}$, and let $P=\sum_a\lambda_a\cdot\exp(i\langle\delta'_a,\hat A\rangle)\cdot\Psi_a$ be its representation in the basis of $\mathcal{P}$ (Definition~\ref{def:covariant-classes}).
Put
\[
P^{(\delta)}:=\sum_a\lambda_a\cdot\exp\left(\langle\delta'_a,\delta\rangle\right)\cdot\exp\left(i\langle\delta'_a,\hat A\rangle\right)\cdot\Psi_a .
\]
Since $\exp\left(i\langle\delta'_a,\hat A-i\delta\rangle\right)=\exp\left(\langle\delta'_a,\delta\rangle\right)\cdot\exp\left(i\langle\delta'_a,\hat A\rangle\right)$,
the element $P^{(\delta)}$ arises from $P$ when $\hat A$ is replaced by $\hat A-i\delta$.

Suppose, in addition, that all links of $S$ lie in one coordinate plane $\{x_\mu=j\}$, where $\mu\in\{1,2,3\}$ and $j\in\Z$.
For a link $b=(y,\nu)\in S$, write the eight bonds over $b$ (Proposition~\ref{prop:geometry}\eqref{it:geometry-projection}) as ordered pairs $(x,x')$ oriented like $b$,
that is, with $\pi(x)=y$ and $\pi(x')=y+e_\nu$.
With the phases $U_A(x,x')=\exp\left(iqA(\pi(x),\pi(x'))\right)$ and the hopping matrices $w(x,x')$ of Definition~\ref{def:chamber-operator} put
\[
\begin{aligned}
\mathfrak{Q}_\delta:=\sum_{b\in S}\ \sum_{(x,x')}[\,&\left(\exp(q\delta_b)-1\right)\cdot\bar\psi(x)\cdot w(x,x')\cdot U_A(x,x')\cdot\psi(x')\\
&+\left(\exp(-q\delta_b)-1\right)\cdot\bar\psi(x')\cdot w(x',x)\cdot U_A(x',x)\cdot\psi(x)\,],
\end{aligned}
\]
where, for each $b\in S$, the inner sum ranges over the eight bonds $(x,x')$ over $b$, oriented like $b$; and put
\[
\mathfrak{V}_\delta:=\exp(\mathfrak{Q}_\delta).
\]
\end{definition}

\begin{proposition}[Imaginary shift of the contour]\label{prop:contour-shift}
Let $\rho\in\R$, $M\in\R$, $w_0,\dots,w_3>0$ and $q>0$.
Let $\mathbb{T}=(\Gamma,\chi)$ be any torus (Definition~\ref{def:tori}),
let $S$ be a finite set of links that all lie in one plane $\{x_\mu=j\}$ with $\mu\in\{1,2,3\}$ and $j\in\Z$,
let $P\in\mathcal{P}$, and let $B$ be a box such that
\begin{equation}\label{eq:elimination-shift-box}
B\ \text{embeds in}\ \mathbb{T},\qquad
P\ \text{is read in}\ B,\qquad
b+[-1,1]^3\subset B\quad\text{for every }b\in S
\end{equation}
(Definitions~\ref{def:integrands} and~\ref{def:enlarged-boxes}).
In particular, for every $b\in S$ the box $B$ contains the four plaquettes that contain $b$
and the cubes $y+\left[-\frac12,\frac12\right]^3$ around the two end points $y$ of $b$.
Let $\delta\in\R^S$, and let $P^{(\delta)}$, $\mathfrak{Q}_\delta$ and $\mathfrak{V}_\delta$ be as in Definition~\ref{def:shifted-elements}.
\begin{enumerate}
\item\label{it:shift-polynomial}
\emph{The polynomial $\mathfrak{V}_\delta$.}
The element $\mathfrak{Q}_\delta$ is even and lies in $\mathcal{P}$, and $\mathfrak{Q}_\delta^r=0$ for every integer $r>64\cdot\#|S|$;
hence
\[
\mathfrak{V}_\delta=\sum_{r=0}^{64\cdot\#|S|}\frac{1}{r!}\,\mathfrak{Q}_\delta^r
\]
is a polynomial in $\mathfrak{Q}_\delta$, and $\mathfrak{V}_\delta\in\mathcal{P}$.
Every link read by $\mathfrak{V}_\delta$ is a link $b\in S$ with $\delta_b\neq0$,
and every site read by $\mathfrak{V}_\delta$ is an end point of a bond over such a link (Definition~\ref{def:two-layer}).
If $\mu=1$, these sites lie in the thin layers $\mathbb{X}_{j,2},\dots,\mathbb{X}_{j,5}$,
so that $\mathfrak{V}_\delta$ reads only links and sites of the thick layer $j$.
The element $\mathfrak{V}_\delta$ is read in $B$.
Finally, let $A\in C^1(\Gamma)\otimes\C$ be a complex $\Gamma$-periodic one-cochain,
let $\delta$ be identified with its periodisation,
and let $\mathfrak{Q}_\delta$ and $\mathfrak{V}_\delta$ be formed with the phases $U_A(x,x')=\exp\left(iqA(\pi(x),\pi(x'))\right)$ of this $A$.
Then, in $\Gr_{\mathbb{T}}$,
\[
S_{\mathbb{T}}(A-i\delta)-S_{\mathbb{T}}(A)=-\mathfrak{Q}_\delta,
\qquad
\exp\left(-S_{\mathbb{T}}(A-i\delta)\right)=\exp\left(-S_{\mathbb{T}}(A)\right)\cdot\mathfrak{V}_\delta .
\]
\item\label{it:shift-identity}
\emph{Shift of the contour.}
For every entire factor $\Phi$ on $\mathbb{T}$,
\[
Z_{\mathbb{T}}[\Phi\cdot P]=Z_{\mathbb{T}}\left[\Phi^{(\delta)}\cdot\exp\left(\frac12\|d\delta\|^2\right)\cdot\varphi_{d\delta}\cdot\mathfrak{V}_\delta\cdot P^{(\delta)}\right].
\]
The element $\exp\left(\frac12\|d\delta\|^2\right)\cdot\varphi_{d\delta}\cdot\mathfrak{V}_\delta\cdot P^{(\delta)}$ lies in $\mathcal{P}$ and is read in $B$;
here $\varphi_{d\delta}=\exp\left(i\langle d^*d\delta,\hat A\rangle\right)$ for a local lift $\hat A$ on $B$.
\item\label{it:shift-arrays}
\emph{Shifted arrays.}
Let $\mathcal{Q}=(\mathfrak{F};o_1,o_2,o_3;h)$ be a cell structure with an arbitrary frame $\mathfrak{F}$ and width $h$ (Definition~\ref{def:cell-structures}),
and let $Q_0$ be its open cell with $\kappa=0$.
For $\kappa\in\Z^3$ let $g_\kappa$ be the element of $W(\mathcal{Q})$ with $g_\kappa Q_0=Q_\kappa$, and put $g_{\kappa*}:=(g_\kappa^{-1})^*$ (Definition~\ref{def:factors}).
Let $O=\sum_ka_k\cdot\varphi_{\beta_k}\in\mathcal{T}$ with $\supp O\subset Q_0$ (Definition~\ref{def:characters}).
Its mirror copies $O^{[\kappa]}$ (Definition~\ref{def:folds}) are the functions $O^{[\kappa]}(F)=O(g_\kappa^*F)$, by \eqref{eq:setting-group-action} and because $O$ is real.
The folded array of $O$ is the formal sum $O_{\mathcal{Q}}:=\sum_{\kappa\in\Z^3}O^{[\kappa]}$, with one term in each open cell;
on a torus that carries $\mathcal{Q}$ it is the function $\sum_\kappa O^{[\kappa]}(\widetilde F)$, where $\kappa$ ranges over a set of representatives of the cells of the torus,
as in Definition~\ref{def:arrays}.
Let $\eta\in C^2_{\mathrm{c}}$ be real, with $\supp\eta$ contained in a box $B_\eta$, and put
\[
\begin{aligned}
O_{\mathcal{Q}}^{[\eta]}-O_{\mathcal{Q}}&:=\sum_{\kappa\in\Z^3}\sum_ka_k\cdot\left(\exp\left(\langle g_{\kappa*}\beta_k,\eta\rangle\right)-1\right)\cdot\varphi_{g_{\kappa*}\beta_k},\\
C_{\mathcal{Q}}(\eta)&:=\sum_{\kappa\in\Z^3}\sum_k|a_k|\cdot\left|\exp\left(\langle g_{\kappa*}\beta_k,\eta\rangle\right)-1\right|.
\end{aligned}
\]
In both sums the term of index $(\kappa,k)$ vanishes unless the open cell $Q_\kappa$ contains a plaquette of $\supp\eta$.
At most $\#|\supp\eta|$ open cells contain such a plaquette, so that both sums are finite,
and these cells are contained in $B_\eta^{+h}$ (Definition~\ref{def:enlarged-boxes}).
With $\|\beta\|_1:=\sum_\sigma|\beta(\sigma)|$ and $\|\eta\|_\infty:=\max_\sigma|\eta(\sigma)|$, where $\sigma$ ranges over the plaquettes,
\[
C_{\mathcal{Q}}(\eta)\le\#|\supp\eta|\cdot|O|_1\cdot\left(\exp\left(\max_k\|\beta_k\|_1\cdot\|\eta\|_\infty\right)-1\right).
\]
On every torus that carries $\mathcal{Q}$ and in which $B_\eta^{+h}$ embeds, $\eta$ being identified there with its periodisation (Definition~\ref{def:factors}),
one has $O_{\mathcal{Q}}(F-i\eta)-O_{\mathcal{Q}}(F)=\left(O_{\mathcal{Q}}^{[\eta]}-O_{\mathcal{Q}}\right)(F)$ for every real $F$, and for every $t\in\R$
\[
\left(\exp\left(tO_{\mathcal{Q}}\right)\right)^{[\eta]}=\exp\left(tO_{\mathcal{Q}}\right)\cdot\exp\left(t\left(O_{\mathcal{Q}}^{[\eta]}-O_{\mathcal{Q}}\right)\right),
\qquad
\exp\left(t\left(O_{\mathcal{Q}}^{[\eta]}-O_{\mathcal{Q}}\right)\right)=1+\sum_\alpha b_\alpha(t)\cdot\varphi_{\beta_\alpha},
\]
where $\alpha$ ranges over a countable index set and
\[
\sum_\alpha|b_\alpha(t)|\le\exp\left(C_{\mathcal{Q}}(\eta)\cdot|t|\right)-1 ;
\]
the series converges absolutely and, for real $F$, uniformly in $F$,
and the cochains $\beta_\alpha\in C^2_{\mathrm{c}}$ are supported in $B_\eta^{+h}$ and depend neither on $t$ nor on the torus.
\end{enumerate}
\end{proposition}

\begin{proof}[Proof of Proposition~\ref{prop:contour-shift}\eqref{it:shift-polynomial}]
\emph{The element $\mathfrak{Q}_\delta$.}
Let $b\in S$, let $(x,x')$ be a bond over $b$, oriented like $b$, and let $\alpha,\alpha'\in\{1,\dots,4\}$.
A local lift $\hat A$ on a box that contains $b$ satisfies $\exp(iq\hat A(b))=\exp(iqA(b))$ (Definition~\ref{def:integrands}),
so that \eqref{eq:elimination-pair} gives $U_A(x,x')=\exp\left(i\langle q\cdot1_{(\pi(x),\pi(x'))},\hat A\rangle\right)$ and
\[
\bar\psi_\alpha(x)\cdot w(x,x')_{\alpha\alpha'}\cdot U_A(x,x')\cdot\psi_{\alpha'}(x')
=w(x,x')_{\alpha\alpha'}\cdot\exp\left(i\langle q\cdot1_{(\pi(x),\pi(x'))},\hat A\rangle\right)\cdot\bar\psi_\alpha(x)\,\psi_{\alpha'}(x').
\]
The charge density of the monomial $\bar\psi_\alpha(x)\,\psi_{\alpha'}(x')$ is $\delta_{\pi(x')}-\delta_{\pi(x)}$ (Definition~\ref{def:chamber-operator}),
and by \eqref{eq:elimination-pair}
\[
d^*\left(q\cdot1_{(\pi(x),\pi(x'))}\right)=q\cdot\left(\delta_{\pi(x')}-\delta_{\pi(x)}\right).
\]
So the element $\exp\left(i\langle q\cdot1_{(\pi(x),\pi(x'))},\hat A\rangle\right)\cdot\bar\psi_\alpha(x)\,\psi_{\alpha'}(x')$ is gauge invariant
and belongs to $\mathcal{P}$ (Definition~\ref{def:polynomials}).
The terms with $x$ and $x'$ exchanged are treated in the same way, with the cochain $q\cdot1_{(\pi(x'),\pi(x))}=-q\cdot1_{(\pi(x),\pi(x'))}$.
Hence $\mathfrak{Q}_\delta$ lies in $\mathcal{P}$;
it is even, since each of its terms contains exactly two generators.

\emph{Nilpotency.}
By Proposition~\ref{prop:geometry}\eqref{it:geometry-projection} the bonds over a link are of type $3$, and there are exactly eight of them.
Every site has exactly one neighbour of each type (Definition~\ref{def:fermion-sites}), so two different bonds of type $3$ have no end point in common;
and a bond over a link determines this link, whose end points are the images under $\pi$ of the end points of the bond.
Therefore the $8\cdot\#|S|$ bonds over the links of $S$ are pairwise different, and their $16\cdot\#|S|$ end points are pairwise different sites,
which carry $128\cdot\#|S|$ generators $\psi_\alpha(x)$, $\bar\psi_\alpha(x)$.
Every term of $\mathfrak{Q}_\delta$ is a function of $\hat A$ times a product of two of these generators,
so that $\mathfrak{Q}_\delta^r$ is a sum of functions of $\hat A$ times products of $2r$ of them.
A product of more than $128\cdot\#|S|$ of these generators contains one of them twice and vanishes;
consequently $\mathfrak{Q}_\delta^r=0$ for $r>64\cdot\#|S|$.
In a Grassmann algebra the even elements commute with all elements;
the exponential series of an even nilpotent element is a finite sum,
and, by the binomial formula, the exponential of the sum of two even nilpotent elements is the product of their exponentials.
In particular $\mathfrak{V}_\delta=\sum_{r=0}^{64\cdot\#|S|}\frac{1}{r!}\,\mathfrak{Q}_\delta^r$.
The product of two basis elements of $\mathcal{P}$ is zero or, up to a sign, again a basis element:
the exponents add, the monomials multiply,
and the charge density of a non-zero product of two monomials is the sum of their charge densities,
so that the gauge condition of Definition~\ref{def:polynomials} is preserved.
Thus $\mathcal{P}$ is closed under multiplication, and $\mathfrak{V}_\delta\in\mathcal{P}$.

\emph{Links and sites read by $\mathfrak{V}_\delta$.}
The terms of $\mathfrak{Q}_\delta$ with $\delta_b=0$ vanish, because $\exp(0)-1=0$.
By the first step, each of the other terms is a multiple of a basis element
whose exponent is $\pm q\cdot1_{(\pi(x),\pi(x'))}$ for a bond $(x,x')$ over a link $b\in S$ with $\delta_b\neq0$,
and whose monomial is the product of a generator at $x$ and a generator at $x'$.
In a product of such basis elements the exponents add and the monomials multiply,
and collecting equal basis elements produces neither a new exponent nor a new monomial.
Hence, in the representation of $\mathfrak{V}_\delta=1+\sum_{r\ge1}\frac{1}{r!}\,\mathfrak{Q}_\delta^r$ in the basis of $\mathcal{P}$,
every exponent is supported on links $b\in S$ with $\delta_b\neq0$,
and every monomial is a product of generators at end points of bonds over such links;
by Definition~\ref{def:two-layer} this is the assertion on the links and the sites that $\mathfrak{V}_\delta$ reads.
If $\mu=1$, the end points of the eight bonds over each link of $S$ lie in $\mathbb{X}_{j,2},\dots,\mathbb{X}_{j,5}$
by Lemma~\ref{lem:layer-bonds}\eqref{it:layer-phases};
these four thin layers and the links of the plane $\{x_1=j\}$ belong to the thick layer $j$ (Definition~\ref{def:two-layer}).

\emph{Reading in $B$.}
Let $b\in S$, and let $x$ be an end point of a bond over $b$.
Then $\pi(x)$ is an end point of $b$, and
$\pi(x)+\left[-\frac12,\frac12\right]^3\subset b+[-1,1]^3\subset B$ by \eqref{eq:elimination-shift-box};
that is, $x\in\mathbb{X}(B)$ (Definition~\ref{def:integrands}).
Moreover $b\subset b+[-1,1]^3\subset B$, that is, $b\in B^{(1)}$.
Together with the previous step this shows that $\mathfrak{V}_\delta$ is read in $B$ (Definition~\ref{def:enlarged-boxes}).

\emph{The action on the torus.}
Let $A\in C^1(\Gamma)\otimes\C$.
The links of $S$ lie in $B^{(1)}$, and $B$ embeds in $\mathbb{T}$;
by Lemma~\ref{lem:torus-elements}\eqref{it:torus-injective} the links of $S$ are therefore pairwise different modulo $\Gamma$,
and the projection $\mathbb{X}\to\mathbb{X}_{\mathbb{T}}$ is injective on $\mathbb{X}(B)$.
In particular the periodisation of $\delta$, with which $\delta$ is identified (Definition~\ref{def:factors}),
takes the value $\delta_b$ on every link congruent modulo $\Gamma$ to a link $b\in S$ and vanishes on all other links.
Let $\mathbb{X}'$ be a set of representatives of $\mathbb{X}_{\mathbb{T}}$ that contains $\mathbb{X}(B)$.
The generators at the sites of $\mathbb{X}'$ are free generators of $\Gr_{\mathbb{T}}$, and by Definitions~\ref{def:chamber-operator} and~\ref{def:tori}
\[
S_{\mathbb{T}}(A)=\sum_{x\in\mathbb{X}'}M\cdot\bar\psi(x)\,\psi(x)-\sum_{x\in\mathbb{X}'}\sum_{i=0}^{3}\bar\psi(x)\cdot w(x,x^{(i)})\cdot U_A(x,x^{(i)})\cdot\psi(x^{(i)}),
\]
where $\psi(x^{(i)})$ stands for $\chi(g)\cdot\psi(x^{(i)}-g)$, $g$ being the element of $\Gamma$ with $x^{(i)}-g\in\mathbb{X}'$.
The pairs $(x,x^{(i)})$ with $x\in\mathbb{X}'$ and $i\in\{0,\dots,3\}$ are oriented bonds,
and every oriented bond is congruent modulo $\Gamma$ to exactly one of them, since the maps $x\mapsto x^{(i)}$ commute with the translations by $\Gamma$.
A component $A(b)$ with $b\in S$ enters only through the phases of the bonds over the links congruent to $b$,
that is, of the translates by elements of $\Gamma$ of the eight bonds over $b$;
bonds over different links of $S$ are different modulo $\Gamma$, as these links are.
So the summands of $S_{\mathbb{T}}(A)$ that contain a component $A(b)$, $b\in S$, are,
for every $b\in S$ and every bond $(x,x')$ over $b$ oriented like $b$,
the summand $-\bar\psi(x)\cdot w(x,x')\cdot U_A(x,x')\cdot\psi(x')$ belonging to $x$
and the summand $-\bar\psi(x')\cdot w(x',x)\cdot U_A(x',x)\cdot\psi(x)$ belonging to $x'$;
no factor $\chi(g)$ appears in them, because $x$ and $x'$ lie in $\mathbb{X}(B)\subset\mathbb{X}'$.
The cochains $A-i\delta$ and $A$ coincide on all links that are not congruent to a link of $S$,
and for a bond $(x,x')$ over $b\in S$, oriented like $b$,
\[
U_{A-i\delta}(x,x')=\exp\left(iq(A(b)-i\delta_b)\right)=\exp(q\delta_b)\cdot U_A(x,x'),
\qquad
U_{A-i\delta}(x',x)=\exp(-q\delta_b)\cdot U_A(x',x).
\]
Inserting this into the summands just listed gives $S_{\mathbb{T}}(A-i\delta)-S_{\mathbb{T}}(A)=-\mathfrak{Q}_\delta$, with $\mathfrak{Q}_\delta$ formed with the phases of $A$.
The algebra $\Gr_{\mathbb{T}}$ has finitely many generators;
hence its even elements without scalar part, among them $S_{\mathbb{T}}(A)$ and $\mathfrak{Q}_\delta$, are nilpotent, and they commute.
The binomial formula now yields
\[
\exp\left(-S_{\mathbb{T}}(A-i\delta)\right)=\exp\left(-S_{\mathbb{T}}(A)+\mathfrak{Q}_\delta\right)
=\exp\left(-S_{\mathbb{T}}(A)\right)\cdot\exp(\mathfrak{Q}_\delta)=\exp\left(-S_{\mathbb{T}}(A)\right)\cdot\mathfrak{V}_\delta .
\]
\end{proof}

\begin{proof}[Proof of Proposition~\ref{prop:contour-shift}\eqref{it:shift-identity}]
Let $\mu\in\{1,2,3\}$ and $j\in\Z$ be such that every link of $S$ lies in the plane $\{x_\mu=j\}$.
Assume first that $\mu=1$.
Denote by $\mathbb{T}^{(1)}$ the set of the links of $\mathbb{T}$, that is, of the classes modulo $\Gamma$ of the links of $\Z^3$.
By Lemma~\ref{lem:torus-elements}\eqref{it:torus-injective}, a cell of $\Z^3$ that is contained in $B$ is identified with its class;
every link of $S$ is contained in $B$ by \eqref{eq:elimination-shift-box}, so that $S$ becomes a subset of $\mathbb{T}^{(1)}$.
For a real function $A'$ on $\mathbb{T}^{(1)}\setminus S$ and for $\zeta\in\C^S$,
let $(A',\zeta)$ be the complex one-cochain on $\mathbb{T}$ that equals $A'$ off $S$ and $\zeta$ on $S$.

\emph{Step~1: free orbits.}
Let $b=(x,\nu)\in S$, so that $x_1=j$ and $\nu\in\{2,3\}$.
The plaquette over $b$ that leaves the plane is $\sigma(b):=(x;1\nu)$, and by Definition~\ref{def:cochains}
\[
(dA)_{\sigma(b)}=A_1(x)+A_\nu(x+e_1)-A_1(x+e_\nu)-A_\nu(x).
\]
The four links of $\sigma(b)$ are contained in $b+[-1,1]^3$, hence in $B$ by \eqref{eq:elimination-shift-box}.
Two of them, $(x,1)$ and $(x+e_\nu,1)$, have the direction $e_1$, and the link $(x+e_1,\nu)$ lies in the plane $\{x_1=j+1\}$.
None of these three links belongs to $S$;
as they and the links of $S$ lie in $B$, Lemma~\ref{lem:torus-elements}\eqref{it:torus-injective} shows that none of them is congruent modulo $\Gamma$ to a link of $S$.
The plaquette $(x;1\nu)$ determines $x$ and $\nu$;
therefore the plaquettes $\sigma(b)$, $b\in S$, are pairwise different, also modulo $\Gamma$.
Let $u$ be a one-cochain on $\mathbb{T}$ with values in $\R$, $\C$ or $\Z$ and with support in $S$.
Of the four links of $\sigma(b)$ only $b$ can carry a non-zero value of $u$, and the sign of $A_\nu(x)$ in the formula for $(dA)_{\sigma(b)}$ is negative; hence
\begin{equation}\label{eq:elimination-free-orbits}
(du)_{\sigma(b)}=-u_b,\qquad b\in S.
\end{equation}
An element $k\in\Z^S$ is read as the integer one-cochain on $\mathbb{T}$ that equals $k$ on $S$ and vanishes off $S$;
for such $k$, \eqref{eq:elimination-free-orbits} shows that $dk=0$ only if $k=0$.
The set
\[
Z^2(\Gamma;\Z):=\left\{z\in C^2(\Gamma;\Z):\ dz=0\right\}
\]
is countable.
Since $d(dk)=0$, the group $\Z^S$ acts on $Z^2(\Gamma;\Z)$ by $z\mapsto z+dk$,
and by \eqref{eq:elimination-free-orbits} this action is free.
Let $\mathcal{R}\subset Z^2(\Gamma;\Z)$ be a set of representatives of its orbits.

\emph{Step~2: unfolding.}
Let $H$ be a Borel function on $\Conf(\Gamma)$ that is either non-negative or integrable.
By Definition~\ref{def:configurations}, $\int_{\Conf(\Gamma)}H$ is the sum over $z\in Z^2(\Gamma;\Z)$ of the integrals of $H(\cdot,z)$ over $[0,2\pi/q)^{\mathbb{T}^{(1)}}$.
Each $z\in Z^2(\Gamma;\Z)$ can be written as $z=z_0+dk$ with a unique $z_0\in\mathcal{R}$ and a unique $k\in\Z^S$,
and $(A,z_0+dk)\sim\left(A+\tfrac{2\pi}{q}k,z_0\right)$ by the same definition.
The translates $[0,2\pi/q)^S+\tfrac{2\pi}{q}k$, $k\in\Z^S$, form a partition of $\R^S$.
By the theorems of Tonelli and Fubini, therefore,
\begin{equation}\label{eq:elimination-unfolding}
\begin{aligned}
\int_{\Conf(\Gamma)}H
&=\sum_{z_0\in\mathcal{R}}\ \sum_{k\in\Z^S}\ \int_{[0,2\pi/q)^{\mathbb{T}^{(1)}}}\mathrm{d}A\,H\left(A+\frac{2\pi}{q}k,\,z_0\right)\\
&=\sum_{z_0\in\mathcal{R}}\ \int_{[0,2\pi/q)^{\mathbb{T}^{(1)}\setminus S}}\mathrm{d}A'\int_{\R^S}\mathrm{d}\xi\,H\left((A',\xi),\,z_0\right).
\end{aligned}
\end{equation}

\emph{Step~3: the entire integrand.}
Fix $z_0\in\mathcal{R}$ and $A'\in[0,2\pi/q)^{\mathbb{T}^{(1)}\setminus S}$.
As $B$ embeds in $\mathbb{T}$, Definition~\ref{def:integrands} provides an integer one-cochain $k'$ on $B^{(1)}$ whose coboundary agrees with $z_0$ on $B^{(2)}$;
for each $z_0\in\mathcal{R}$ one such $k'$ is fixed, the same for all $A'$.
For $\zeta\in\C^S$ put
\[
A(\zeta):=(A',\zeta),\qquad
F(\zeta):=dA(\zeta)+\frac{2\pi}{q}z_0,\qquad
\hat A(\zeta):=A(\zeta)|_{B^{(1)}}+\frac{2\pi}{q}k'.
\]
Both $F(\zeta)$ and $\hat A(\zeta)$ are affine in $\zeta$, and $F(\zeta-i\delta)=F(\zeta)-i\,d\delta$.
Now let $\zeta$ be real.
On $B^{(2)}$ the coboundary of $\hat A(\zeta)$ is $dA(\zeta)+\frac{2\pi}{q}\,dk'=dA(\zeta)+\frac{2\pi}{q}z_0=F(\zeta)$,
and on a link $b\in B^{(1)}$ the integer $k'(b)$ drops out of the phase:
\[
\exp\left(iq\hat A(\zeta)(b)\right)=\exp\left(iqA(\zeta)(b)\right)\cdot\exp\left(2\pi i\,k'(b)\right)=\exp\left(iqA(\zeta)(b)\right).
\]
These are the two properties that define a local lift in Definition~\ref{def:integrands};
so $\hat A(\zeta)$ is a local lift on $B$ of the point of $\Conf(\Gamma)$ that $(A(\zeta),z_0)$ represents.
Write $P=\sum_a\lambda_a\cdot\exp(i\langle\delta'_a,\hat A\rangle)\cdot\Psi_a$ and put
\[
J_a(A):=\int_{\mathbb{T}}\exp\left(-S_{\mathbb{T}}(A)\right)\cdot\Psi_a .
\]
Expanding the exponential one sees that $J_a$ is a polynomial in the variables $\exp(\pm iqA(b))$, $b\in\mathbb{T}^{(1)}$,
with coefficients that do not depend on $A$; for complex $A$ it is given by the same polynomial.
Extend the pairing $\langle\cdot,\cdot\rangle$ bilinearly, without complex conjugation, to complex cochains, and let
\begin{equation}\label{eq:elimination-integrand}
\begin{aligned}
f(\zeta):={}&\exp\left(-\frac12\left\langle F(\zeta),F(\zeta)\right\rangle\right)\cdot\Phi\left(F(\zeta)\right)\\
&\cdot\sum_a\lambda_a\cdot\exp\left(i\left\langle\delta'_a,\hat A(\zeta)\right\rangle\right)\cdot J_a\left(A(\zeta)\right),
\qquad\zeta\in\C^S.
\end{aligned}
\end{equation}
The factor $\Phi$ extends to an entire function on $C^2(\Gamma)\otimes\C$ (Definition~\ref{def:factors}),
and $J_a(A(\zeta))$ is a polynomial in exponentials of affine functions of $\zeta$;
hence $f$ is an entire function on $\C^S$.
For real $\zeta$ one has $\langle F(\zeta),F(\zeta)\rangle=\|F(\zeta)\|^2$ and
\[
\int_{\mathbb{T}}\exp\left(-S_{\mathbb{T}}(A(\zeta))\right)\cdot P\left(\hat A(\zeta)\right)
=\sum_a\lambda_a\cdot\exp\left(i\left\langle\delta'_a,\hat A(\zeta)\right\rangle\right)\cdot J_a\left(A(\zeta)\right),
\]
so that $f(\zeta)=H_1\left((A',\zeta),z_0\right)$, where
\[
H_1:=\exp\left(-\frac12\|F\|^2\right)\cdot\Phi(F)\cdot J_{\mathbb{T}}[P]
\]
is a function on $\Conf(\Gamma)$ as in Lemma~\ref{lem:torus-elements}\eqref{it:torus-convergence}.

\emph{Step~4: a bound on a strip.}
Put $\|\delta\|_\infty:=\max_{b\in S}|\delta_b|$,
let $\zeta\in\C^S$ satisfy $|\Im\zeta_b|\le\|\delta\|_\infty$ for all $b\in S$, and write $\xi:=\Re\zeta$.
As $A'$ and $z_0$ are real, $\Im F(\zeta)=d(0,\Im\zeta)$.
This two-cochain can be non-zero only at a plaquette with a side in $S$, and a link is a side of four plaquettes;
so it has at most $4\cdot\#|S|$ non-zero values,
each of them a sum of at most four numbers $\pm\Im\zeta_b$ and therefore of modulus at most $4\|\delta\|_\infty$.
Summing the squares one obtains
\[
\|\Im F(\zeta)\|^2\le4\cdot\#|S|\cdot\left(4\|\delta\|_\infty\right)^2=64\cdot\|\delta\|_\infty^2\cdot\#|S|.
\]
For $b=(x_b,\nu)\in S$, with initial point $x_b$, \eqref{eq:elimination-free-orbits} yields $\Re F(\zeta)_{\sigma(b)}=\xi^0_b-\xi_b$, where
\[
\xi^0_b:=\left(d(A',0)+\frac{2\pi}{q}z_0\right)_{\sigma(b)}
=A'_1(x_b)+A'_\nu(x_b+e_1)-A'_1(x_b+e_\nu)+\frac{2\pi}{q}z_0\left(\sigma(b)\right)
\]
does not depend on $\zeta$.
As the plaquettes $\sigma(b)$ are pairwise different,
\[
\Re\left\langle F(\zeta),F(\zeta)\right\rangle=\|\Re F(\zeta)\|^2-\|\Im F(\zeta)\|^2
\ge\sum_{b\in S}\left(\xi_b-\xi^0_b\right)^2-64\cdot\|\delta\|_\infty^2\cdot\#|S|.
\]
The other factors of $f$ are bounded as follows.
First, $|\Phi(F(\zeta))|\le\|\Phi\|_{4\|\delta\|_\infty}$ (Definition~\ref{def:factors}).
Second, $\Im\hat A(\zeta)$ is the restriction of $(0,\Im\zeta)$ to $B^{(1)}$; with $\|\delta'_a\|_1:=\sum_b|\delta'_a(b)|$ this gives
$\left|\exp\left(i\langle\delta'_a,\hat A(\zeta)\rangle\right)\right|\le\exp\left(\|\delta\|_\infty\cdot\|\delta'_a\|_1\right)$.
Third, let $|J_a|_1$ be the sum of the moduli of the coefficients of the polynomial $J_a$,
and $\deg J_a$ its total degree in the variables $\exp(\pm iqA(b))$ with $b\in S$.
At $A(\zeta)$ the variables with $b\notin S$ have modulus one and those with $b\in S$ modulus at most $\exp(q\cdot\|\delta\|_\infty)$;
so $|J_a(A(\zeta))|\le|J_a|_1\cdot\exp\left(q\cdot\|\delta\|_\infty\cdot\deg J_a\right)$.
It follows that
\begin{equation}\label{eq:elimination-strip}
|f(\zeta)|\le C_f\cdot\prod_{b\in S}\exp\left(-\frac12\left(\Re\zeta_b-\xi^0_b\right)^2\right)
\qquad\text{for }\zeta\in\C^S\text{ with }\max_{b\in S}|\Im\zeta_b|\le\|\delta\|_\infty,
\end{equation}
where
\[
\begin{aligned}
C_f:={}&\exp\left(32\cdot\|\delta\|_\infty^2\cdot\#|S|\right)\cdot\|\Phi\|_{4\|\delta\|_\infty}\\
&\cdot\sum_a|\lambda_a|\cdot\exp\left(\|\delta\|_\infty\cdot\|\delta'_a\|_1\right)\cdot|J_a|_1\cdot\exp\left(q\cdot\|\delta\|_\infty\cdot\deg J_a\right)
\end{aligned}
\]
depends neither on $\zeta$ nor on $A'$ and $z_0$.

\emph{Step~5: shift by Cauchy's theorem.}
Enumerate $S=\{b_1,\dots,b_{\#|S|}\}$ and put
$\delta_{\le\ell}:=(\delta_{b_1},\dots,\delta_{b_\ell},0,\dots,0)\in\R^S$ for $0\le\ell\le\#|S|$.
By \eqref{eq:elimination-strip}, every function $\xi\mapsto f(\xi-i\delta_{\le\ell})$ is integrable on $\R^S$.
Let $1\le\ell\le\#|S|$, and fix real numbers $\xi_b$ for $b\neq b_\ell$.
With $\zeta_b:=\xi_b-i\delta_b$ for $b=b_1,\dots,b_{\ell-1}$ and $\zeta_b:=\xi_b$ for $b=b_{\ell+1},\dots,b_{\#|S|}$,
the function $\zeta_{b_\ell}\mapsto f(\zeta)$ is entire on $\C$.
Apply Cauchy's theorem to it on the rectangle with the corners $-R$, $R$, $R-i\delta_{b_\ell}$, $-R-i\delta_{b_\ell}$, where $R>0$.
On the two vertical sides $\zeta$ lies in the strip of \eqref{eq:elimination-strip}, and there the factors with $b\neq b_\ell$ are at most one;
so these sides contribute at most
\[
|\delta_{b_\ell}|\cdot C_f\cdot\left(\exp\left(-\frac12\left(R-\xi^0_{b_\ell}\right)^2\right)+\exp\left(-\frac12\left(R+\xi^0_{b_\ell}\right)^2\right)\right),
\]
which tends to $0$ as $R\to\infty$.
In the limit $R\to\infty$ the two horizontal sides give
\[
\int_{\R}\mathrm{d}\zeta_{b_\ell}\,f(\dots,\zeta_{b_\ell},\dots)=\int_{\R}\mathrm{d}\zeta_{b_\ell}\,f(\dots,\zeta_{b_\ell}-i\delta_{b_\ell},\dots),
\]
both integrals being taken over the real line.
Integration over the other real variables, which is allowed by Fubini's theorem and \eqref{eq:elimination-strip}, gives
\[
\int_{\R^S}\mathrm{d}\xi\,f(\xi-i\delta_{\le\ell-1})=\int_{\R^S}\mathrm{d}\xi\,f(\xi-i\delta_{\le\ell}).
\]
Since $\delta_{\le0}=0$ and $\delta_{\le\#|S|}=\delta$, induction on $\ell$ yields
\[
\int_{\R^S}\mathrm{d}\xi\,f(\xi)=\int_{\R^S}\mathrm{d}\xi\,f(\xi-i\delta).
\]

\emph{Step~6: identification and folding back.}
Let $\xi\in\R^S$ and abbreviate $F:=F(\xi)$, $A:=A(\xi)$, $\hat A:=\hat A(\xi)$.
Since $\delta$ vanishes off $S\subset B^{(1)}$, one has $A(\xi-i\delta)=A-i\delta$, $\hat A(\xi-i\delta)=\hat A-i\delta$ and $F(\xi-i\delta)=F-i\,d\delta$.
Hence
\[
\begin{aligned}
\left\langle F-i\,d\delta,F-i\,d\delta\right\rangle&=\|F\|^2-2i\langle d\delta,F\rangle-\|d\delta\|^2,\\
\exp\left(i\left\langle\delta'_a,\hat A-i\delta\right\rangle\right)&=\exp\left(\langle\delta'_a,\delta\rangle\right)\cdot\exp\left(i\langle\delta'_a,\hat A\rangle\right).
\end{aligned}
\]
Moreover $\exp(-S_{\mathbb{T}}(A-i\delta))=\exp(-S_{\mathbb{T}}(A))\cdot\mathfrak{V}_\delta$ in $\Gr_{\mathbb{T}}$ by Proposition~\ref{prop:contour-shift}\eqref{it:shift-polynomial},
so that $J_a(A-i\delta)=\int_{\mathbb{T}}\exp(-S_{\mathbb{T}}(A))\cdot\mathfrak{V}_\delta\cdot\Psi_a$.
Inserting these three facts into \eqref{eq:elimination-integrand} and recalling the definition of $P^{(\delta)}$ (Definition~\ref{def:shifted-elements}), one finds
\[
\begin{aligned}
f(\xi-i\delta)={}&\exp\left(-\frac12\|F\|^2\right)\cdot\Phi(F-i\,d\delta)\cdot\exp\left(\frac12\|d\delta\|^2\right)\cdot\exp\left(i\langle d\delta,F\rangle\right)\\
&\cdot\int_{\mathbb{T}}\exp\left(-S_{\mathbb{T}}(A)\right)\cdot\mathfrak{V}_\delta\cdot P^{(\delta)}(\hat A).
\end{aligned}
\]
The cochain $d\delta$ is supported on plaquettes that contain a link of $S$, and these lie in $B$;
hence $\|d\delta\|$ and $\langle d\delta,F\rangle$ take the same values on $\mathbb{T}$ as on $\Z^3$ (Definition~\ref{def:factors}).
As $d\hat A=F$ on $B^{(2)}$,
\[
\exp\left(i\langle d\delta,F\rangle\right)=\exp\left(i\langle d\delta,d\hat A\rangle\right)=\exp\left(i\langle d^*d\delta,\hat A\rangle\right),
\]
where $d^*d\delta$ is supported in $B^{(1)}$ and $d^*(d^*d\delta)=0$.
Thus $\varphi_{d\delta}=\exp(i\langle d^*d\delta,\hat A\rangle)$ is an element of $\mathcal{P}$ that is read in $B$ (Definition~\ref{def:polynomials}).
The same holds for $\mathfrak{V}_\delta$, by Proposition~\ref{prop:contour-shift}\eqref{it:shift-polynomial},
and for $P^{(\delta)}$, which has the same basis elements as $P$.
Products of elements of $\mathcal{P}$ that are read in $B$ are again such elements:
for two basis elements the exponents add, the monomials multiply, and the gauge condition of Definition~\ref{def:polynomials} is additive.
Therefore $\exp\left(\frac12\|d\delta\|^2\right)\cdot\varphi_{d\delta}\cdot\mathfrak{V}_\delta\cdot P^{(\delta)}$ lies in $\mathcal{P}$ and is read in $B$, and
\[
\begin{aligned}
f(\xi-i\delta)&=H_2\left((A',\xi),z_0\right),\\
H_2&:=\exp\left(-\frac12\|F\|^2\right)\cdot\Phi^{(\delta)}(F)\cdot J_{\mathbb{T}}\left[\exp\left(\frac12\|d\delta\|^2\right)\cdot\varphi_{d\delta}\cdot\mathfrak{V}_\delta\cdot P^{(\delta)}\right].
\end{aligned}
\]
The entire factors $\Phi$ and $\Phi^{(\delta)}=\Phi^{[d\delta]}$ are factors (Definition~\ref{def:factors}),
so $H_1$ and $H_2$ are integrable on $\Conf(\Gamma)$ by Lemma~\ref{lem:torus-elements}\eqref{it:torus-convergence}.
Let $f$ be formed as in Step~3 with $z_0\in\mathcal{R}$ and $A'\in[0,2\pi/q)^{\mathbb{T}^{(1)}\setminus S}$.
Then \eqref{eq:elimination-unfolding} for $H_1$, Step~5 for every pair $(z_0,A')$, and \eqref{eq:elimination-unfolding} for $H_2$ give
\[
\begin{aligned}
Z_{\mathbb{T}}[\Phi\cdot P]=\int_{\Conf(\Gamma)}H_1
&=\sum_{z_0\in\mathcal{R}}\ \int_{[0,2\pi/q)^{\mathbb{T}^{(1)}\setminus S}}\mathrm{d}A'\int_{\R^S}\mathrm{d}\xi\,f(\xi)\\
&=\sum_{z_0\in\mathcal{R}}\ \int_{[0,2\pi/q)^{\mathbb{T}^{(1)}\setminus S}}\mathrm{d}A'\int_{\R^S}\mathrm{d}\xi\,f(\xi-i\delta)\\
&=\int_{\Conf(\Gamma)}H_2
=Z_{\mathbb{T}}\left[\Phi^{(\delta)}\cdot\exp\left(\frac12\|d\delta\|^2\right)\cdot\varphi_{d\delta}\cdot\mathfrak{V}_\delta\cdot P^{(\delta)}\right],
\end{aligned}
\]
where the first and the last equality hold by Lemma~\ref{lem:torus-elements}\eqref{it:torus-convergence}.

There remain the planes $\{x_\mu=j\}$ with $\mu\in\{2,3\}$, for which the six steps are modified at two points.
In this case each link of $S$ has the form $b=(x,\nu)$ with $x_\mu=j$ and $\nu\neq\mu$.

The first modification concerns the plaquettes.
The symbol $\sigma(b)$ now stands for the plaquette spanned by $b$ and $e_\mu$, whose corners are $x$, $x+e_\mu$, $x+e_\nu$ and $x+e_\mu+e_\nu$;
in the notation of Definition~\ref{def:cochains} it is $(x;\mu\nu)$ when $\mu<\nu$ and $(x;\nu\mu)$ when $\nu<\mu$.
It lies in $b+[-1,1]^3$, and it determines $b$.
Besides $b$, its sides are $(x,\mu)$ and $(x+e_\nu,\mu)$, which are parallel to $e_\mu$, and $(x+e_\mu,\nu)$, which lies in the plane $\{x_\mu=j+1\}$;
as in Step~1, none of these three is congruent modulo $\Gamma$ to a link of $S$.
Expanding $(dA)_{\sigma(b)}$ by the formula of Definition~\ref{def:cochains} one meets the term $s(b)\cdot A_\nu(x)$, where
\[
s(b):=\begin{cases}-1&\text{if }\mu<\nu,\\+1&\text{if }\nu<\mu.\end{cases}
\]
The reasoning of Step~1 therefore leads to
\begin{equation}\label{eq:elimination-other-planes}
(du)_{\sigma(b)}=s(b)\cdot u_b,\qquad b\in S,
\end{equation}
for every one-cochain $u$ on $\mathbb{T}$ with values in $\R$, $\C$ or $\Z$ and with support in $S$,
and this identity takes over the role of \eqref{eq:elimination-free-orbits}.

The second modification concerns the centres of the Gaussian factors in Step~4.
With $\xi^0_b:=\left(d(A',0)+\frac{2\pi}{q}z_0\right)_{\sigma(b)}$ one has $\Re F(\zeta)_{\sigma(b)}=\xi^0_b+s(b)\cdot\xi_b$,
and because $s(b)^2=1$,
\[
\left(\xi^0_b+s(b)\cdot\xi_b\right)^2=\left(\xi_b+s(b)\cdot\xi^0_b\right)^2.
\]
Consequently \eqref{eq:elimination-strip}, as well as the bound for the vertical sides in Step~5, holds with $-s(b)\cdot\xi^0_b$ in the place of $\xi^0_b$.
With these two modifications Steps~1 to~6 establish the identity for the planes $\{x_\mu=j\}$, $\mu\in\{2,3\}$, too.
\end{proof}

\begin{proof}[Proof of Proposition~\ref{prop:contour-shift}\eqref{it:shift-arrays}]
For $g\in W$ and $\beta\in C^2_{\mathrm{c}}$ one has $g_*\beta=(g^{-1})^*\beta=(g^*)^{-1}\beta$, and $g^*$ is orthogonal (Definition~\ref{def:cochains}); hence
\[
\varphi_\beta(g^*F)=\exp\left(i\langle\beta,g^*F\rangle\right)=\exp\left(i\langle g_*\beta,F\rangle\right)=\varphi_{g_*\beta}(F),
\]
and $g_*\beta$ is supported on $g(\supp\beta)$.
Together with $O^{[\kappa]}(F)=O(g_\kappa^*F)$ (Definitions~\ref{def:folds} and~\ref{def:arrays}; $O$ is real) this shows that
$O^{[\kappa]}=\sum_ka_k\cdot\varphi_{g_{\kappa*}\beta_k}$,
where $g_{\kappa*}\beta_k$ is supported on $g_\kappa(\supp\beta_k)\subset g_\kappa Q_0=Q_\kappa$.
A plaquette at which both $g_{\kappa*}\beta_k$ and $\eta$ are non-zero is therefore a plaquette of $\supp\eta$ contained in $Q_\kappa$;
if the open cell $Q_\kappa$ contains no plaquette of $\supp\eta$, then $\langle g_{\kappa*}\beta_k,\eta\rangle=0$ for every $k$.
A cell that does contain a plaquette of $\supp\eta$ is called contributing, and so are the pairs $(\kappa,k)$ formed with it.
For every other pair the factor $\exp(\langle g_{\kappa*}\beta_k,\eta\rangle)-1$ vanishes,
and the pair enters neither $O_{\mathcal{Q}}^{[\eta]}-O_{\mathcal{Q}}$ nor $C_{\mathcal{Q}}(\eta)$.
Distinct open cells are disjoint, so a plaquette lies in at most one cell, while one cell may contain several plaquettes of $\supp\eta$:
there are at most $\#|\supp\eta|$ contributing cells.

Every cell has sup-diameter at most $h$.
For the coordinate frame the cells are open cubes of side $h$.
For the diagonal frame, $x_1=(\ell_++\ell_-)/2$ and $x_2=(\ell_+-\ell_-)/2$ range through open intervals of length $h$ on a cell, and so does $x_3$.
Each further frame of Definition~\ref{def:cell-structures} arises from one of these two when the coordinates are permuted and some of their signs reversed,
and a map of this kind is an isometry for $|\cdot|_\infty$.
A cell that meets $B_\eta$, in particular a contributing cell, is therefore contained in $B_\eta^{+h}$ (Definition~\ref{def:enlarged-boxes}).
The pairing $\langle g_{\kappa*}\beta_k,\eta\rangle$ is real, so that
\[
\left|\exp\left(\langle g_{\kappa*}\beta_k,\eta\rangle\right)-1\right|\le\exp\left(\left|\langle g_{\kappa*}\beta_k,\eta\rangle\right|\right)-1;
\]
moreover $|\langle g_{\kappa*}\beta_k,\eta\rangle|\le\|g_{\kappa*}\beta_k\|_1\cdot\|\eta\|_\infty=\|\beta_k\|_1\cdot\|\eta\|_\infty$,
because $g_{\kappa*}$ permutes the plaquettes up to signs.
Summation over $k$ and over the at most $\#|\supp\eta|$ contributing cells yields
\[
C_{\mathcal{Q}}(\eta)\le\#|\supp\eta|\cdot|O|_1\cdot\left(\exp\left(\max_k\|\beta_k\|_1\cdot\|\eta\|_\infty\right)-1\right).
\]

Now let $\mathbb{T}=(\Gamma,\chi)$ be a torus that carries $\mathcal{Q}$ and in which $B_\eta^{+h}$ embeds.
On $\mathbb{T}$, $O_{\mathcal{Q}}(F)$ is the sum of the terms $O(g_\kappa^*\widetilde F)$ over the cells of $\mathbb{T}$,
that is, over the classes modulo $\Gamma$ of the open cells of $\mathcal{Q}$,
one open cell $Q_\kappa$, $\kappa\in\Z^3$, being chosen in each class (Definitions~\ref{def:cell-structures} and~\ref{def:arrays});
and $\eta$ is identified with its periodisation, the sum of its translates by the elements of $\Gamma$ (Definition~\ref{def:factors}).
Since $B_\eta^{+h}$ embeds, the boxes $B_\eta^{+h}+g$, $g\in\Gamma$, are pairwise disjoint (Definition~\ref{def:integrands}).
If $Q_\kappa$ and $Q_\kappa+g$ both met $\supp\eta\subset B_\eta$, both would lie in $B_\eta^{+h}$;
then $B_\eta^{+h}$ would meet $B_\eta^{+h}+g$, and so $g=0$.
Thus each class of cells contains at most one cell that meets $\supp\eta$; where such a cell exists, it is chosen as the representative of its class.
For such a representative the support of $g_{\kappa*}\beta_k$ lies in $Q_\kappa\subset B_\eta^{+h}$,
whereas the translate of $\eta$ by $g\in\Gamma\setminus\{0\}$ is supported in $B_\eta+g\subset B_\eta^{+h}+g$;
hence the pairing of $g_{\kappa*}\beta_k$ with the periodisation of $\eta$ equals $\langle g_{\kappa*}\beta_k,\eta\rangle$.
For a class none of whose cells meets $\supp\eta$ this pairing vanishes:
the translate of $\eta$ by $g$ pairs non-trivially with $g_{\kappa*}\beta_k$ only if the cell $Q_\kappa-g$ of that class contains a plaquette of $\supp\eta$.
As $\varphi_\beta(F-i\eta)=\exp(\langle\beta,\eta\rangle)\cdot\varphi_\beta(F)$, it follows that on $\mathbb{T}$
\[
O_{\mathcal{Q}}(F-i\eta)-O_{\mathcal{Q}}(F)=\sum_\kappa\sum_ka_k\cdot\left(\exp\left(\langle g_{\kappa*}\beta_k,\eta\rangle\right)-1\right)\cdot\varphi_{g_{\kappa*}\beta_k}(F),
\]
the sum extending over the $\kappa\in\Z^3$ for which the open cell $Q_\kappa$ of $\mathcal{Q}$ is contributing;
the sum of the moduli of its coefficients is at most $C_{\mathcal{Q}}(\eta)$.
In particular, for every real $t$,
\[
\left(\exp\left(tO_{\mathcal{Q}}\right)\right)^{[\eta]}=\exp\left(tO_{\mathcal{Q}}\right)\cdot\exp\left(t\left(O_{\mathcal{Q}}^{[\eta]}-O_{\mathcal{Q}}\right)\right).
\]
By $\varphi_\beta\cdot\varphi_{\beta'}=\varphi_{\beta+\beta'}$, the $r$-th power of $O_{\mathcal{Q}}^{[\eta]}-O_{\mathcal{Q}}$ is a finite combination of characters
whose exponents are sums of $r$ of the cochains $g_{\kappa*}\beta_k$ of contributing pairs,
and the sum of the moduli of its coefficients, which is submultiplicative, is at most $C_{\mathcal{Q}}(\eta)^r$.
Therefore
\[
\exp\left(t\left(O_{\mathcal{Q}}^{[\eta]}-O_{\mathcal{Q}}\right)\right)-1=\sum_{r\ge1}\frac{t^r}{r!}\left(O_{\mathcal{Q}}^{[\eta]}-O_{\mathcal{Q}}\right)^r=\sum_\alpha b_\alpha(t)\cdot\varphi_{\beta_\alpha},
\]
where $\alpha$ ranges over the countable set of the finite non-empty sequences of contributing pairs $(\kappa,k)$;
for a sequence $\alpha$ of length $r$, $b_\alpha(t)$ is $t^r/r!$ times the product of the coefficients $a_k\cdot\left(\exp(\langle g_{\kappa*}\beta_k,\eta\rangle)-1\right)$ of its members,
and $\beta_\alpha$ is the sum of their cochains $g_{\kappa*}\beta_k$.
Consequently
\[
\sum_\alpha|b_\alpha(t)|\le\sum_{r\ge1}\frac{|t|^r}{r!}\cdot C_{\mathcal{Q}}(\eta)^r=\exp\left(C_{\mathcal{Q}}(\eta)\cdot|t|\right)-1,
\]
and, since $|\varphi_\beta(F)|=1$ for real $F$, the series converges absolutely and uniformly in $F$.
Each $\beta_\alpha\in C^2_{\mathrm{c}}$ is a sum of cochains supported in contributing cells, which are contained in $B_\eta^{+h}$;
so $\beta_\alpha$ is supported in $B_\eta^{+h}$, and it depends neither on $t$ nor on $\mathbb{T}$.
\end{proof}

\begin{remark}\label{rem:imaginary-translation}
The shift of the contour has the following meaning in terms of the operators of a rectangular cross-section (Definition~\ref{def:slice-operators}).
Let $\delta\in\R^{\Lambda^{(1)}_{\mathrm{s}}}$ be a real one-cochain on the links of the slice torus $\Lambda_{\mathrm{s}}$
with $\langle\delta,\gamma\rangle\in2\pi\Z$ for every $\gamma\in\Gamma_{\mathrm{s}}$.
For such $\delta$ the function $A\mapsto\exp(i\langle\delta,A\rangle)$ is $\Gamma_{\mathrm{s}}$-periodic,
so that multiplication by it is a unitary operator, denoted $\exp(i\langle\delta,A\rangle)$ as well,
and conjugation by this operator shifts each $E_b$ by $\delta_b$:
$\exp(-i\langle\delta,A\rangle)\cdot E_b\cdot\exp(i\langle\delta,A\rangle)=E_b+\delta_b$ for every $b\in\Lambda^{(1)}_{\mathrm{s}}$ (Definition~\ref{def:slice-operators}).
Through the joint spectral calculus this unitary equivalence yields the relation \eqref{eq:setting-weyl} for every Borel function $g$ on $\R^{\Lambda^{(1)}_{\mathrm{s}}}$,
as an identity of possibly unbounded operators, both sides being defined exactly on the domain of $g(E+\delta)$.
Write $\langle\delta,E\rangle:=\sum_b\delta_bE_b$, $|E|^2:=\sum_bE_b^2$ and $|\delta|^2:=\sum_b\delta_b^2$, where $b$ ranges over $\Lambda^{(1)}_{\mathrm{s}}$.
Take $g(x):=\exp\left(\frac12|x|^2\right)$ for $x\in\R^{\Lambda^{(1)}_{\mathrm{s}}}$,
so that $g(E)$ is the inverse of the Gaussian factor $\exp\left(-\frac12|E|^2\right)$ of the crossing operator $K$,
and multiply the relation from the right by this bounded factor.
As $\frac12|x+\delta|^2-\frac12|x|^2=\langle\delta,x\rangle+\frac12|\delta|^2$, one obtains
\[
\exp\left(\frac12|E|^2\right)\cdot\exp\left(i\langle\delta,A\rangle\right)\cdot\exp\left(-\frac12|E|^2\right)
=\exp\left(i\langle\delta,A\rangle\right)\cdot\exp\left(\langle\delta,E\rangle+\frac12|\delta|^2\right),
\]
where both sides are defined exactly on the domain of $\exp(\langle\delta,E\rangle)$.
Here $\exp(\langle\delta,E\rangle)=\exp(-i\langle i\delta,E\rangle)$ is the translation operator of Definition~\ref{def:slice-operators} at the imaginary vector $i\delta$:
it replaces the argument $A$ by $A-i\delta$.
This translation is an unbounded operator.
For this reason the shift is made in the integral over configurations, on entire integrands,
and not in the algebra $\mathfrak{W}$ of Definition~\ref{def:mid-plane}, the union of the mid-plane algebras.
\end{remark}

\subsection{One elimination step}\label{sec:elimination-step}

\begin{proposition}[Elimination of an outer plane]\label{prop:outer-plane}
Let $\rho\in\R$ with $\rho^2\neq1$, let $M\in\R$, $w_0,\dots,w_3>0$ and $q>0$.
Let $j\in\Z$ and let $B$ be a box.
\begin{enumerate}
\item\label{it:lowest-plane}
\emph{Lowest plane.}
Let $P\in\mathcal{P}$ be read in $B$, and let the lowest plane of $P$ be at least $j$ (Definition~\ref{def:two-layer}).
Put
\[
P^{\uparrow}:=\mathfrak{R}^{+}_{\mathbb{X}_{j,6}}\mathfrak{R}^{+}_{\mathbb{X}_{j,5}}\cdots\mathfrak{R}^{+}_{\mathbb{X}_{j,1}}P ,
\]
and let $P^{\uparrow}=\sum_{a\in I}\lambda_a\cdot\exp(i\langle\delta'_a,\hat A\rangle)\cdot\Psi_a$ be its unique representation in the basis of $\mathcal{P}$, with a finite index set $I$.
For $a\in I$ let $(\delta'_a)_j$ be the component of $\delta'_a$ on the plane $\{x_1=j\}$, a cochain on the links of the slice (Definition~\ref{def:slice-components});
let $\delta_a\in C^1_{\mathrm{c}}$ be the cochain that is supported on the links of the plane $\{x_1=j+1\}$ and has there the component $(\delta_a)_{j+1}=(\delta'_a)_j$;
and put
\begin{equation}\label{eq:elimination-new-terms}
P_a:=\lambda_a\cdot\exp\left(\frac{1}{2}\|d\delta_a\|^2+\langle\delta'_a,\delta_a\rangle\right)\cdot\exp(i\langle\delta'_a+d^*d\delta_a,\hat A\rangle)\cdot\mathfrak{V}_{\delta_a}\cdot\Psi_a ,
\end{equation}
so that, in the notation of Definition~\ref{def:shifted-elements},
\[
P_a=\exp\left(\frac{1}{2}\|d\delta_a\|^2\right)\cdot\varphi_{d\delta_a}\cdot\mathfrak{V}_{\delta_a}\cdot\left(\lambda_a\cdot\exp(i\langle\delta'_a,\hat A\rangle)\cdot\Psi_a\right)^{(\delta_a)} .
\]
Then every $P_a$ lies in $\mathcal{P}$, has lowest plane at least $j+1$ and highest plane at most the larger of the two numbers $j+2$ and the highest plane of $P$, and is read in $B^{+8}$ (Definition~\ref{def:enlarged-boxes});
the cochains $\delta_a$ and the elements $P_a$ are defined on $\Z^3$, and no torus and no factor enters their definition;
and
\[
Z_{\mathbb{T}}[\Phi\cdot P]=\sum_{a\in I}Z_{\mathbb{T}}\left[\Phi^{(\delta_a)}\cdot P_a\right]
\]
for every torus $\mathbb{T}$ in which $B^{+8}$ embeds and every entire factor $\Phi$ on $\mathbb{T}$.
For every $a\in I$ the cochain $\delta'':=\delta'_a+d^*d\delta_a$ has the components
\[
\begin{aligned}
\delta''_j&=0, & \delta''^{\perp}_j&=0,\\
\delta''_{j+1}&=(\delta'_a)_{j+1}+(2+d_{\mathrm{s}}^{*}d_{\mathrm{s}})(\delta'_a)_j, & \delta''^{\perp}_{j+1}&=(\delta'_a)^{\perp}_{j+1}+d_{\mathrm{s}}^{*}(\delta'_a)_j,\\
\delta''_{j+2}&=(\delta'_a)_{j+2}-(\delta'_a)_j ,
\end{aligned}
\]
all its other components coincide with those of $\delta'_a$, and
\[
\|d\delta_a\|^2=2\|(\delta'_a)_j\|^2+\|d_{\mathrm{s}}(\delta'_a)_j\|^2,
\qquad
\langle\delta'_a,\delta_a\rangle=\langle(\delta'_a)_{j+1},(\delta'_a)_j\rangle .
\]
\item\label{it:highest-plane}
\emph{Highest plane.}
Let $P\in\mathcal{P}$ be read in $B$, and let the highest plane of $P$ be at most $j$.
Put
\[
P^{\downarrow}:=\mathfrak{R}^{-}_{\mathbb{X}_{j,1}}\cdots\mathfrak{R}^{-}_{\mathbb{X}_{j,5}}\mathfrak{R}^{-}_{\mathbb{X}_{j,6}}P ,
\]
where $\mathfrak{R}^{-}_{\mathbb{X}_{j,6}}$ is applied first, and let $P^{\downarrow}=\sum_{a\in I}\lambda_a\cdot\exp(i\langle\delta'_a,\hat A\rangle)\cdot\Psi_a$ be its unique representation in the basis of $\mathcal{P}$, with a finite index set $I$.
For $a\in I$ let $\delta_a\in C^1_{\mathrm{c}}$ be the cochain that is supported on the links of the plane $\{x_1=j-1\}$ and has there the component $(\delta_a)_{j-1}=(\delta'_a)_j$, and define $P_a$ by \eqref{eq:elimination-new-terms}.
Then every $P_a$ lies in $\mathcal{P}$, has highest plane at most $j-1$ and lowest plane at least the smaller of the two numbers $j-2$ and the lowest plane of $P$, and is read in $B^{+8}$;
the cochains $\delta_a$ and the elements $P_a$ are defined on $\Z^3$, and no torus and no factor enters their definition;
and
\[
Z_{\mathbb{T}}[\Phi\cdot P]=\sum_{a\in I}Z_{\mathbb{T}}\left[\Phi^{(\delta_a)}\cdot P_a\right]
\]
for every torus $\mathbb{T}$ in which $B^{+8}$ embeds and every entire factor $\Phi$ on $\mathbb{T}$.
For every $a\in I$ the cochain $\delta'':=\delta'_a+d^*d\delta_a$ has the components
\[
\begin{aligned}
\delta''_j&=0, & \delta''^{\perp}_{j-1}&=0,\\
\delta''_{j-1}&=(\delta'_a)_{j-1}+(2+d_{\mathrm{s}}^{*}d_{\mathrm{s}})(\delta'_a)_j, & \delta''^{\perp}_{j-2}&=(\delta'_a)^{\perp}_{j-2}-d_{\mathrm{s}}^{*}(\delta'_a)_j,\\
\delta''_{j-2}&=(\delta'_a)_{j-2}-(\delta'_a)_j ,
\end{aligned}
\]
all its other components coincide with those of $\delta'_a$, and
\[
\|d\delta_a\|^2=2\|(\delta'_a)_j\|^2+\|d_{\mathrm{s}}(\delta'_a)_j\|^2,
\qquad
\langle\delta'_a,\delta_a\rangle=\langle(\delta'_a)_{j-1},(\delta'_a)_j\rangle .
\]
\item\label{it:reflected-elimination}
\emph{Reflected identity.}
Let $P$, $I$ and $(\delta_a,P_a)_{a\in I}$ be as in \eqref{it:lowest-plane} or as in \eqref{it:highest-plane}, let $\Pi$ be a mirror and $\theta:=\theta_\Pi$.
For every torus $\mathbb{T}$ in which $\theta(B^{+8})$ embeds and every entire factor $\Phi$ on $\mathbb{T}$,
\[
Z_{\mathbb{T}}[\Phi\cdot\Theta_\Pi P]=\sum_{a\in I}Z_{\mathbb{T}}\left[\Phi^{[-\theta^*d\delta_a]}\cdot\Theta_\Pi P_a\right],
\]
where $\Phi^{[-\theta^*d\delta_a]}$ is the shifted factor $F\mapsto\Phi(F+i\,\theta^*d\delta_a)$ of Definition~\ref{def:factors}.
\item\label{it:plane-estimate}
\emph{Estimate.}
Let $\mathfrak{n}=(n_i)$ be a sequence of sides that satisfies \textup{(Div)} (Definition~\ref{def:sides}),
let $\mathfrak{d}$ be data with the observable $O$ (Definition~\ref{def:data}),
let $h$ be the width of the cell structure $\mathcal{Q}_{\mathfrak{d}}$,
and let $C_{\mathcal{Q}_{\mathfrak{d}}}$ be formed with $\mathcal{Q}_{\mathfrak{d}}$ and $O$ as in Proposition~\ref{prop:contour-shift}\eqref{it:shift-arrays}.
Let $\mathcal{U}=(P;(\eta_a,P_a)_{a\in I};B')$ be any universal identity (Definition~\ref{def:universal-identity}), and put
\[
P^{\flat}:=\sum_{a\in I}P_a,
\qquad
C:=\max_{a\in I}C_{\mathcal{Q}_{\mathfrak{d}}}(\eta_a).
\]
For all $t\geq0$ and all $i$ so large that $D(n_i)$ carries $\mathcal{Q}_{\mathfrak{d}}$ and $(B')^{+h}$ embeds in $D(n_i)$,
\begin{equation}\label{eq:elimination-estimate}
\left|\mu^t_i(P)-\mu^t_i(P^{\flat})\right|\leq(\exp(Ct)-1)\cdot\sum_{a\in I}\sup_{\beta}\left|\mu^t_i(\varphi_\beta\cdot P_a)\right| ,
\end{equation}
the supremum being taken over all $\beta\in C^2_{\mathrm{c}}$ supported in $(B')^{+h}$.
In particular, let $P$, $I$ and $(\delta_a,P_a)_{a\in I}$ be as in \eqref{it:lowest-plane} or as in \eqref{it:highest-plane}, let $\Pi$ be a mirror with $\theta:=\theta_\Pi$, and put
\[
P^{\flat}:=\sum_{a\in I}P_a,
\qquad
C:=\max_{a\in I}C_{\mathcal{Q}_{\mathfrak{d}}}(d\delta_a),
\qquad
C':=\max_{a\in I}C_{\mathcal{Q}_{\mathfrak{d}}}(-\theta^*d\delta_a).
\]
For all $t\geq0$ and all $i$ so large that $D(n_i)$ carries $\mathcal{Q}_{\mathfrak{d}}$ and $(B^{+8})^{+h}$ embeds in $D(n_i)$,
\[
\left|\mu^t_i(P)-\mu^t_i(P^{\flat})\right|\leq(\exp(Ct)-1)\cdot\sum_{a\in I}\sup_{\beta}\left|\mu^t_i(\varphi_\beta\cdot P_a)\right| ,
\]
the supremum being taken over all $\beta\in C^2_{\mathrm{c}}$ supported in $(B^{+8})^{+h}$;
and for all $t\geq0$ and all $i$ so large that $D(n_i)$ carries $\mathcal{Q}_{\mathfrak{d}}$ and $(\theta(B^{+8}))^{+h}$ embeds in $D(n_i)$,
\[
\left|\mu^t_i(\Theta_\Pi P)-\mu^t_i(\Theta_\Pi P^{\flat})\right|\leq(\exp(C't)-1)\cdot\sum_{a\in I}\sup_{\beta}\left|\mu^t_i(\varphi_\beta\cdot\Theta_\Pi P_a)\right| ,
\]
the supremum being taken over all $\beta\in C^2_{\mathrm{c}}$ supported in $(\theta(B^{+8}))^{+h}$.
\end{enumerate}
\end{proposition}

\begin{proof}
\eqref{it:lowest-plane}
\emph{The fields of the thick layer $j$.}
As the lowest plane of $P$ is at least $j$, the element $P$ has no field in a thin layer below $\mathbb{X}_{j,1}$, and it reads no link that meets a plane $\{x_1=j'\}$ with $j'<j$.
Proposition~\ref{prop:schwinger-dyson}\eqref{it:replacement-fields}, applied to the six maps one after the other, shows by induction on $k=1,\dots,6$ that $\mathfrak{R}^{+}_{\mathbb{X}_{j,k}}\cdots\mathfrak{R}^{+}_{\mathbb{X}_{j,1}}P$ has no field in $\mathbb{X}_{j,k}$ or below it.
For $k=6$ this says that every field of $P^{\uparrow}$ lies in a thick layer $j'\geq j+1$.
By the same item, a field of $P^{\uparrow}$ that is not a field of $P$ lies in $\mathbb{X}_{j,k}^{+}$ or in $\mathbb{X}_{j,k}^{++}$ for some $k\in\{1,\dots,6\}$, hence in $\mathbb{X}_{j,6}^{++}=\mathbb{X}_{j+1,2}$ or below it, and in particular in a thick layer $j'\leq j+1$.
Six applications of Proposition~\ref{prop:schwinger-dyson}\eqref{it:replacement-identity} show that $P^{\uparrow}$ is read in $B^{+6}$ and that
\begin{equation}\label{eq:elimination-upward}
Z_{\mathbb{T}}[\Phi\cdot P]=Z_{\mathbb{T}}[\Phi\cdot P^{\uparrow}]
\end{equation}
for every torus $\mathbb{T}$ in which $B^{+6}$ embeds and every factor $\Phi$ on $\mathbb{T}$;
here the boxes $B^{+1},\dots,B^{+5}$ of the intermediate steps embed in $\mathbb{T}$ as well, since they are contained in $B^{+6}$.
Next, by Proposition~\ref{prop:schwinger-dyson}\eqref{it:replacement-links}, applied to the six maps one after the other,
every link that $P^{\uparrow}$ reads and $P$ does not read joins $\pi(x)$ to $\pi(x')$ for a bond $\{x,x'\}$ of type~$3$
with an end in one of the thin layers $\mathbb{X}_{j,1}^{+},\dots,\mathbb{X}_{j,6}^{+}$, that is, in $\mathbb{X}_{j,2},\dots,\mathbb{X}_{j,6}$ or $\mathbb{X}_{j+1,1}$.
By Lemma~\ref{lem:layer-bonds}\eqref{it:layer-types} and~\eqref{it:layer-phases} the bonds of type~$3$ with an end in $\mathbb{X}_{j,2},\dots,\mathbb{X}_{j,5}$ lie over links of the plane $\{x_1=j\}$, and those with an end in $\mathbb{X}_{j,6}$ or in $\mathbb{X}_{j+1,1}$ lie over links from the plane $\{x_1=j\}$ to the plane $\{x_1=j+1\}$.
Consequently $P^{\uparrow}$ reads no link that meets a plane $\{x_1=j'\}$ with $j'<j$.
Finally, the maps $\mathfrak{R}^{+}_{\mathbb{L}}$ are defined on $\Z^3$ (Proposition~\ref{prop:schwinger-dyson}\eqref{it:replacement-maps});
hence so are $P^{\uparrow}$, its representation, the cochains $\delta_a$ and, by \eqref{eq:elimination-new-terms}, the elements $P_a$, and none of them involves a torus or a factor.

\emph{Gauge invariance on the plane $\{x_1=j\}$.}
Fix $a\in I$ and abbreviate $\delta':=\delta'_a$, so that $(\delta'_a)_j=\delta'_j$.
By the first paragraph the monomial $\Psi_a$ has no field in the thick layer $j$, so that its charge density $\mathfrak{q}(\Psi_a)$ vanishes at the photon sites of the plane $\{x_1=j\}$.
The basis element $\exp(i\langle\delta',\hat A\rangle)\cdot\Psi_a$ of $\mathcal{P}$ is gauge invariant, $d^*\delta'=q\cdot\mathfrak{q}(\Psi_a)$ (Definition~\ref{def:polynomials}).
Together with Lemma~\ref{lem:slice-formulas}\eqref{it:slice-adjoint-one} this gives, for every site $y$ of the slice,
\[
0=(d^*\delta')(j,y)=(d_{\mathrm{s}}^{*}\delta'_j)(y)+\delta'^{\perp}_{j-1}(y)-\delta'^{\perp}_j(y).
\]
A link from $(j-1,y)$ to $(j,y)$ meets the plane $\{x_1=j-1\}$, so $P^{\uparrow}$ does not read it, and $\delta'^{\perp}_{j-1}=0$.
Therefore
\begin{equation}\label{eq:elimination-gauge-lowest}
\delta'^{\perp}_j=d_{\mathrm{s}}^{*}\delta'_j .
\end{equation}

\emph{The shift of the contour.}
For a link $b=(y,\nu)$ of the slice and $k\in\Z$ write $(k,b)$ for the link $((k,y),\nu)$ of the plane $\{x_1=k\}$, and let
\[
S:=\{(j+1,b):\ \delta'_j(b)\neq0\},
\]
a finite set of links of the plane $\{x_1=j+1\}$.
The cochain $\delta_a$ is supported on $S$ and is regarded as an element of $\R^S$.
Let $\mathbb{T}$ be a torus in which $B^{+8}$ embeds.
The torus $\mathbb{T}$, the set $S$, the box $B^{+8}$ and the element $\lambda_a\cdot\exp(i\langle\delta',\hat A\rangle)\cdot\Psi_a$ of $\mathcal{P}$ satisfy \eqref{eq:elimination-shift-box}.
Indeed, this element reads only links and sites that $P^{\uparrow}$ reads, so it is read in $B^{+6}$, and a fortiori in $B^{+8}$;
and if $(j+1,b)\in S$, then $P^{\uparrow}$ reads the link $(j,b)$, which is therefore contained in $B^{+6}$, whence
\[
(j+1,b)+[-1,1]^3=(j,b)+e_1+[-1,1]^3\subset(j,b)+[-2,2]^3\subset B^{+8}.
\]
By Definition~\ref{def:shifted-elements},
\[
\left(\lambda_a\cdot\exp(i\langle\delta',\hat A\rangle)\cdot\Psi_a\right)^{(\delta_a)}=\lambda_a\cdot\exp(\langle\delta',\delta_a\rangle)\cdot\exp(i\langle\delta',\hat A\rangle)\cdot\Psi_a ,
\]
and $\varphi_{d\delta_a}=\exp(i\langle d^*d\delta_a,\hat A\rangle)$ by Proposition~\ref{prop:contour-shift}\eqref{it:shift-identity};
comparison with \eqref{eq:elimination-new-terms} gives the second expression for $P_a$ in \eqref{it:lowest-plane}.
In view of that expression, Proposition~\ref{prop:contour-shift}\eqref{it:shift-identity}, applied to the element $\lambda_a\cdot\exp(i\langle\delta',\hat A\rangle)\cdot\Psi_a$ and the cochain $\delta_a$, shows that $P_a$ lies in $\mathcal{P}$, that it is read in $B^{+8}$, and that
\[
Z_{\mathbb{T}}\left[\Phi\cdot\lambda_a\cdot\exp(i\langle\delta',\hat A\rangle)\cdot\Psi_a\right]=Z_{\mathbb{T}}\left[\Phi^{(\delta_a)}\cdot P_a\right]
\]
for every entire factor $\Phi$ on $\mathbb{T}$.
An entire factor is a factor, and $B^{+6}$ embeds in $\mathbb{T}$ because it is contained in $B^{+8}$;
so \eqref{eq:elimination-upward} holds, and the sum over $a\in I$ gives
\[
Z_{\mathbb{T}}[\Phi\cdot P]=Z_{\mathbb{T}}[\Phi\cdot P^{\uparrow}]=\sum_{a\in I}Z_{\mathbb{T}}\left[\Phi^{(\delta_a)}\cdot P_a\right].
\]

\emph{The components of $\delta''$.}
The cochain $\delta_a$ is supported on the links of the plane $\{x_1=j+1\}$, where its component is $\delta'_j$.
Lemma~\ref{lem:slice-formulas}\eqref{it:slice-one-plane}, applied to $\delta_a$, yields $\|d\delta_a\|^2=2\|\delta'_j\|^2+\|d_{\mathrm{s}}\delta'_j\|^2$ and
\[
\begin{aligned}
(d^*d\delta_a)_j&=-\delta'_j, & (d^*d\delta_a)^{\perp}_j&=-d_{\mathrm{s}}^{*}\delta'_j,\\
(d^*d\delta_a)_{j+1}&=(2+d_{\mathrm{s}}^{*}d_{\mathrm{s}})\delta'_j, & (d^*d\delta_a)^{\perp}_{j+1}&=d_{\mathrm{s}}^{*}\delta'_j,\\
(d^*d\delta_a)_{j+2}&=-\delta'_j ,
\end{aligned}
\]
all other components of $d^*d\delta_a$ being zero.
Hence $\delta''_j=\delta'_j-\delta'_j=0$ and, by \eqref{eq:elimination-gauge-lowest}, $\delta''^{\perp}_j=d_{\mathrm{s}}^{*}\delta'_j-d_{\mathrm{s}}^{*}\delta'_j=0$;
the other formulas for the components of $\delta''$ in \eqref{it:lowest-plane} are read off from the list.
As $\delta_a$ vanishes off the plane $\{x_1=j+1\}$ and equals $\delta'_j$ on it, $\langle\delta',\delta_a\rangle=\langle\delta'_{j+1},\delta'_j\rangle$.

\emph{The planes read by $P_a$.}
The product of two basis elements of $\mathcal{P}$ is either zero or, up to sign, a basis element whose exponent is the sum of the two exponents and whose monomial carries the fields of both monomials;
so a product of elements of $\mathcal{P}$ reads only links and sites that one of its factors reads.
By \eqref{eq:elimination-new-terms}, $P_a$ is a scalar multiple of the product of $\exp(i\langle\delta'',\hat A\rangle)\cdot\Psi_a$, which lies in $\mathcal{P}$ because $d^*\delta''=d^*\delta'=q\cdot\mathfrak{q}(\Psi_a)$ (as $d^*d^*=0$), with the even element $\mathfrak{V}_{\delta_a}$ of $\mathcal{P}$.
By Proposition~\ref{prop:contour-shift}\eqref{it:shift-polynomial}, $\mathfrak{V}_{\delta_a}$ reads only links of $S$ and sites of $\mathbb{X}_{j+1,2},\dots,\mathbb{X}_{j+1,5}$.
Hence every link read by $P_a$ lies in the support of $\delta''$ or in $S$, and every site read by $P_a$ is the site of a field of $\Psi_a$ or lies in the thick layer $j+1$.

The cochain $\delta''$ vanishes on every link that meets a plane $\{x_1=j'\}$ with $j'\leq j$:
on the links of the plane $\{x_1=j\}$ and on the links from this plane to the plane $\{x_1=j+1\}$ by the two identities just proved, and on all other such links because both $\delta'$ (first paragraph) and $d^*d\delta_a$ vanish there.
The links of $S$ lie in the plane $\{x_1=j+1\}$, and the fields of $\Psi_a$ are fields of $P^{\uparrow}$, which lie in thick layers $j'\geq j+1$.
So the lowest plane of $P_a$ is at least $j+1$.

It remains to bound the highest plane.
A link read by $P_a$ is read by $P$,
or is one of the links that $P^{\uparrow}$ reads and $P$ does not read (these lie in the plane $\{x_1=j\}$ or join it to the plane $\{x_1=j+1\}$),
or belongs to the support of $d^*d\delta_a$, whose links meet only the planes $\{x_1=j\}$, $\{x_1=j+1\}$, $\{x_1=j+2\}$,
or belongs to $S$.
A site read by $P_a$ carries a field of $P$, or a field of $P^{\uparrow}$ that is not a field of $P$ and lies in a thick layer $j'\leq j+1$, or lies in the thick layer $j+1$.
Therefore the highest plane of $P_a$ is at most the larger of $j+2$ and the highest plane of $P$.

\eqref{it:highest-plane}
\emph{The downward maps.}
By Definition~\ref{def:two-layer}, the hypothesis that the highest plane of $P$ is at most $j$ means two things:
no field of $P$ lies in a thin layer above $\mathbb{X}_{j,6}$,
and no link read by $P$ meets a plane $\{x_1=j'\}$ with $j'>j$.
In $P^{\downarrow}$ the six maps act in the order $\mathfrak{R}^{-}_{\mathbb{X}_{j,6}},\mathfrak{R}^{-}_{\mathbb{X}_{j,5}},\dots,\mathfrak{R}^{-}_{\mathbb{X}_{j,1}}$.
For $k=6,5,\dots,1$ the element $\mathfrak{R}^{-}_{\mathbb{X}_{j,k}}\cdots\mathfrak{R}^{-}_{\mathbb{X}_{j,6}}P$ has no field in $\mathbb{X}_{j,k}$ and none above it.
Indeed, the element to which $\mathfrak{R}^{-}_{\mathbb{X}_{j,k}}$ is applied has no field above $\mathbb{X}_{j,k}$, by the hypothesis on $P$ if $k=6$ and by the preceding step if $k<6$;
so Proposition~\ref{prop:schwinger-dyson}\eqref{it:replacement-fields}, in its form for the downward maps, gives the assertion by descending induction on $k$.
For $k=1$ it says that every field of $P^{\downarrow}$ lies in a thick layer $j'\le j-1$.
By the same item, a field of $P^{\downarrow}$ that is not a field of $P$ lies in $\mathbb{X}_{j,k}^{-}$ or in $\mathbb{X}_{j,k}^{--}$ for some $k\in\{1,\dots,6\}$.
Each of these thin layers is $\mathbb{X}_{j,1}^{--}=\mathbb{X}_{j-1,5}$ or lies above it,
and so such a field lies in a thick layer $j'\ge j-1$.

Proposition~\ref{prop:schwinger-dyson}\eqref{it:replacement-identity}, again for the downward maps, shows that after $s$ of the six steps the element is read in $B^{+s}$, $s=1,\dots,6$,
and that the $s$-th step replaces an element $G$ by an element $G'$ with $Z_{\mathbb{T}}[\Phi\cdot G]=Z_{\mathbb{T}}[\Phi\cdot G']$ for every torus $\mathbb{T}$ in which $B^{+s}$ embeds and every factor $\Phi$ on $\mathbb{T}$.
All these boxes are contained in $B^{+6}$ and embed wherever $B^{+6}$ does.
Hence $P^{\downarrow}$ is read in $B^{+6}$, and
\begin{equation}\label{eq:elimination-downward}
Z_{\mathbb{T}}[\Phi\cdot P]=Z_{\mathbb{T}}[\Phi\cdot P^{\downarrow}]
\end{equation}
for every torus $\mathbb{T}$ in which $B^{+6}$ embeds and every factor $\Phi$ on $\mathbb{T}$.

Consider now a link that $P^{\downarrow}$ reads and $P$ does not read.
One of the six steps turns an element that does not read this link into one that reads it.
According to Proposition~\ref{prop:schwinger-dyson}\eqref{it:replacement-links} for the downward maps, applied to that step,
the link joins $\pi(x)$ to $\pi(x')$ for a bond $\{x,x'\}$ of type~$3$ with an end in one of the thin layers $\mathbb{X}_{j,6}^{-},\dots,\mathbb{X}_{j,1}^{-}$,
that is, in $\mathbb{X}_{j,5}$, $\mathbb{X}_{j,4}$, $\mathbb{X}_{j,3}$, $\mathbb{X}_{j,2}$, $\mathbb{X}_{j,1}$ or $\mathbb{X}_{j-1,6}$.
By Lemma~\ref{lem:layer-bonds}\eqref{it:layer-types} and~\eqref{it:layer-phases}, a bond of type~$3$ with an end in $\mathbb{X}_{j,2},\dots,\mathbb{X}_{j,5}$ lies over a link of the plane $\{x_1=j\}$,
while a bond of type~$3$ with an end in $\mathbb{X}_{j,1}$ or in $\mathbb{X}_{j-1,6}$ joins these two thin layers and lies over a link from the plane $\{x_1=j-1\}$ to the plane $\{x_1=j\}$.
So the new links meet no plane other than $\{x_1=j-1\}$ and $\{x_1=j\}$.
It follows that $P^{\downarrow}$ reads no link that meets a plane $\{x_1=j'\}$ with $j'>j$,
and, together with the description of the new fields, that the lowest plane of $P^{\downarrow}$ is not smaller than the smaller of the two numbers $j-1$ and the lowest plane of $P$.

The maps $\mathfrak{R}^{-}_{\mathbb{L}}$ are defined on $\Z^3$, and no torus and no factor enters their definition (Proposition~\ref{prop:schwinger-dyson}\eqref{it:replacement-maps}).
The same is therefore true of $P^{\downarrow}$, of its representation in the basis of $\mathcal{P}$, of the cochains $\delta_a$ and, by \eqref{eq:elimination-new-terms}, of the elements $P_a$.

\emph{Gauge invariance at the plane $\{x_1=j\}$.}
Fix $a\in I$ and put $\delta':=\delta'_a$; thus $(\delta'_a)_j=\delta'_j$.
The fields of $\Psi_a$ are fields of $P^{\downarrow}$, and none of them lies in the thick layer $j$, which consists of the sites $x$ with $\pi(x)$ in the plane $\{x_1=j\}$.
Hence the charge density $\mathfrak{q}(\Psi_a)$ vanishes at every photon site of this plane.
The basis element $\exp(i\langle\delta',\hat A\rangle)\cdot\Psi_a$ is gauge invariant, $d^*\delta'=q\cdot\mathfrak{q}(\Psi_a)$ (Definition~\ref{def:polynomials});
writing $d^*\delta'$ at the site $(j,y)$ in slice components by Lemma~\ref{lem:slice-formulas}\eqref{it:slice-adjoint-one}, one obtains
\[
0=(d^*\delta')(j,y)=(d_{\mathrm{s}}^{*}\delta'_j)(y)+\delta'^{\perp}_{j-1}(y)-\delta'^{\perp}_j(y)
\]
for every site $y$ of the slice.
The link from $(j,y)$ to $(j+1,y)$ meets the plane $\{x_1=j+1\}$;
by the preceding paragraphs $P^{\downarrow}$ does not read it, whence $\delta'^{\perp}_j=0$.
Consequently
\begin{equation}\label{eq:elimination-gauge-highest}
\delta'^{\perp}_{j-1}=-d_{\mathrm{s}}^{*}\delta'_j .
\end{equation}

\emph{The contour of the plane $\{x_1=j-1\}$.}
With the links $(k,b)$ of the planes $\{x_1=k\}$ named as above, let
\[
S:=\{(j-1,b):\ \delta'_j(b)\neq0\}.
\]
This is a finite set of links of the plane $\{x_1=j-1\}$, namely the support of $\delta_a$;
accordingly, $\delta_a$ is regarded as an element of $\R^S$.
Let $\mathbb{T}$ be a torus in which $B^{+8}$ embeds.
The torus $\mathbb{T}$, the set $S$, the box $B^{+8}$ and the element $\lambda_a\cdot\exp(i\langle\delta',\hat A\rangle)\cdot\Psi_a$ of $\mathcal{P}$ satisfy \eqref{eq:elimination-shift-box}, for two reasons.
First, every link and every site read by this element is read by $P^{\downarrow}$, because the element is one of the terms of the representation of $P^{\downarrow}$;
so it is read in $B^{+6}$, and a fortiori in $B^{+8}$.
Second, if $(j-1,b)\in S$, then $\delta'_j(b)\neq0$, so that $P^{\downarrow}$ reads the link $(j,b)$, which is then contained in $B^{+6}$;
therefore
\[
(j-1,b)+[-1,1]^3=(j,b)-e_1+[-1,1]^3\subset(j,b)+[-2,2]^3\subset B^{+8}.
\]
Proposition~\ref{prop:contour-shift}\eqref{it:shift-identity}, applied with the cochain $\delta_a$, now yields
\[
Z_{\mathbb{T}}\left[\Phi\cdot\lambda_a\cdot\exp(i\langle\delta',\hat A\rangle)\cdot\Psi_a\right]
=Z_{\mathbb{T}}\left[\Phi^{(\delta_a)}\cdot\exp\left(\frac12\|d\delta_a\|^2\right)\cdot\varphi_{d\delta_a}\cdot\mathfrak{V}_{\delta_a}\cdot\left(\lambda_a\cdot\exp(i\langle\delta',\hat A\rangle)\cdot\Psi_a\right)^{(\delta_a)}\right]
\]
for every entire factor $\Phi$ on $\mathbb{T}$,
where the product of the last four factors on the right lies in $\mathcal{P}$ and is read in $B^{+8}$.
That product is $P_a$.
Indeed, Definition~\ref{def:shifted-elements} gives
$\left(\lambda_a\cdot\exp(i\langle\delta',\hat A\rangle)\cdot\Psi_a\right)^{(\delta_a)}=\lambda_a\cdot\exp(\langle\delta',\delta_a\rangle)\cdot\exp(i\langle\delta',\hat A\rangle)\cdot\Psi_a$,
the character is $\varphi_{d\delta_a}=\exp(i\langle d^*d\delta_a,\hat A\rangle)$ by Proposition~\ref{prop:contour-shift}\eqref{it:shift-identity},
and coefficient functions commute with the elements of the Grassmann algebra;
collecting the scalar factors one arrives at \eqref{eq:elimination-new-terms}.
Thus $P_a\in\mathcal{P}$, $P_a$ is read in $B^{+8}$, and
\[
Z_{\mathbb{T}}\left[\Phi\cdot\lambda_a\cdot\exp(i\langle\delta',\hat A\rangle)\cdot\Psi_a\right]=Z_{\mathbb{T}}\left[\Phi^{(\delta_a)}\cdot P_a\right]
\]
for every entire factor $\Phi$ on $\mathbb{T}$.
An entire factor is a factor, and the box $B^{+6}$ embeds in $\mathbb{T}$ together with $B^{+8}$.
Summation over $a\in I$ and \eqref{eq:elimination-downward} give
\[
Z_{\mathbb{T}}[\Phi\cdot P]=Z_{\mathbb{T}}[\Phi\cdot P^{\downarrow}]=\sum_{a\in I}Z_{\mathbb{T}}\left[\Phi^{(\delta_a)}\cdot P_a\right]
\]
for every torus $\mathbb{T}$ in which $B^{+8}$ embeds and every entire factor $\Phi$ on $\mathbb{T}$.

\emph{The new exponent.}
The cochain $\delta_a$ vanishes off the links of the plane $\{x_1=j-1\}$, and its component there is $\delta'_j$.
Lemma~\ref{lem:slice-formulas}\eqref{it:slice-one-plane}, applied to this plane, gives $\|d\delta_a\|^2=2\|\delta'_j\|^2+\|d_{\mathrm{s}}\delta'_j\|^2$ and
\[
\begin{aligned}
(d^*d\delta_a)_{j-2}&=-\delta'_j, & (d^*d\delta_a)^{\perp}_{j-2}&=-d_{\mathrm{s}}^{*}\delta'_j,\\
(d^*d\delta_a)_{j-1}&=(2+d_{\mathrm{s}}^{*}d_{\mathrm{s}})\delta'_j, & (d^*d\delta_a)^{\perp}_{j-1}&=d_{\mathrm{s}}^{*}\delta'_j,\\
(d^*d\delta_a)_{j}&=-\delta'_j ,
\end{aligned}
\]
all other components of $d^*d\delta_a$ being zero.
For $\delta''=\delta'+d^*d\delta_a$ this gives $\delta''_j=\delta'_j-\delta'_j=0$,
and, by \eqref{eq:elimination-gauge-highest}, $\delta''^{\perp}_{j-1}=-d_{\mathrm{s}}^{*}\delta'_j+d_{\mathrm{s}}^{*}\delta'_j=0$.
Adding the list to the components of $\delta'$ gives the formulas for $\delta''_{j-1}$, $\delta''^{\perp}_{j-2}$ and $\delta''_{j-2}$ in \eqref{it:highest-plane},
and shows that all other components of $\delta''$ are those of $\delta'$.
Finally, $\delta_a$ coincides with $\delta'_j$ on the plane $\{x_1=j-1\}$ and vanishes elsewhere, so that $\langle\delta',\delta_a\rangle=\langle\delta'_{j-1},\delta'_j\rangle$.

\emph{The planes read by $P_a$.}
When two basis elements of $\mathcal{P}$ are multiplied, the exponents add, and the product of the two monomials is either zero or, up to sign, a monomial that carries the fields of both.
A product of elements of $\mathcal{P}$ therefore reads only links and sites that are read by one of its factors.
By \eqref{eq:elimination-new-terms}, $P_a$ is a constant multiple of the product of $\mathfrak{V}_{\delta_a}$ with $\exp(i\langle\delta'',\hat A\rangle)\cdot\Psi_a$,
and the latter element lies in $\mathcal{P}$, since $d^*\delta''=d^*\delta'$ by $d^*d^*=0$.
By Proposition~\ref{prop:contour-shift}\eqref{it:shift-polynomial}, the polynomial $\mathfrak{V}_{\delta_a}$ reads only links of $S$ and sites of $\mathbb{X}_{j-1,2},\dots,\mathbb{X}_{j-1,5}$.
Consequently every link read by $P_a$ belongs to the support of $\delta''$ or to $S$,
and every site read by $P_a$ carries a field of $\Psi_a$ or lies in the thick layer $j-1$.

Let $b$ be a link that meets a plane $\{x_1=k\}$ with $k\ge j$.
If $b$ lies in the plane $\{x_1=j\}$ or joins the plane $\{x_1=j-1\}$ to it, then $\delta''(b)=0$ because $\delta''_j=0$ and $\delta''^{\perp}_{j-1}=0$.
Otherwise $b$ meets a plane $\{x_1=k'\}$ with $k'>j$.
Then $P^{\downarrow}$ does not read $b$, so $\delta'(b)=0$;
and $(d^*d\delta_a)(b)=0$, as the components of $d^*d\delta_a$ that may differ from zero belong to the planes $\{x_1=j-2\}$, $\{x_1=j-1\}$, $\{x_1=j\}$ and to the links from $\{x_1=j-2\}$ to $\{x_1=j-1\}$ and from $\{x_1=j-1\}$ to $\{x_1=j\}$.
In either case $\delta''(b)=0$.
Moreover, the links of $S$ lie in the plane $\{x_1=j-1\}$, the fields of $\Psi_a$ lie in thick layers with index at most $j-1$, and the remaining sites read by $P_a$ lie in the thick layer $j-1$.
So the highest plane of $P_a$ is at most $j-1$.

For the lowest plane, recall that the support of $\delta'$ consists of links read by $P^{\downarrow}$, each of which is read by $P$ or meets only the planes $\{x_1=j-1\}$ and $\{x_1=j\}$,
that the links in the support of $d^*d\delta_a$ meet only the planes $\{x_1=j-2\}$, $\{x_1=j-1\}$ and $\{x_1=j\}$,
and that the links of $S$ meet only the plane $\{x_1=j-1\}$.
Hence a link read by $P_a$ is read by $P$, or every plane $\{x_1=k\}$ that it meets has $k\ge j-2$.
A site read by $P_a$ carries a field of $P$, or a field of $P^{\downarrow}$ that is not a field of $P$ and lies in a thick layer $j'\ge j-1$, or lies in the thick layer $j-1$.
The lowest plane of $P_a$ is thus not smaller than the smaller of the two numbers $j-2$ and the lowest plane of $P$.

\eqref{it:reflected-elimination}
Let $\mathbb{T}=(\Gamma,\chi)$ be a torus in which $\theta(B^{+8})$ embeds, let $\Phi$ be an entire factor on $\mathbb{T}$, and put $\eta_a:=d\delta_a$ for $a\in I$;
on entire factors, the shift by $\delta_a$ is the shift by $\eta_a$ (Definition~\ref{def:factors}).
The reflection $\theta$ is an involution of $\R^3$ whose linear part $\vec\theta$ maps the cube $[-1,1]^3$ onto itself.
It maps the sets $\theta(B^{+8})+[-1,1]^3+g$, $g\in\Gamma$, which are pairwise disjoint, onto the sets $B^{+8}+[-1,1]^3+\vec\theta g$.
Hence $B^{+8}$ embeds in the reflected torus $\theta\mathbb{T}=(\theta\Gamma,\chi\circ\vec\theta)$.
Moreover $\Phi^{\theta}$ is an entire factor on $\theta\mathbb{T}$, $\theta(\theta\mathbb{T})=\mathbb{T}$, and a factor reflected twice is the factor itself (Definition~\ref{def:factors}).
The element $P$ is read in $B$, hence in $B^{+8}$, and each $P_a$ is read in $B^{+8}$ by \eqref{it:lowest-plane}, resp. \eqref{it:highest-plane}.

Three identities are combined.
Lemma~\ref{lem:reflection-identities}\eqref{it:reflected-functional}, applied with the box $B^{+8}$ to the factor $\Phi$ and the element $P$, gives
\[
Z_{\mathbb{T}}[\Phi\cdot\Theta_\Pi P]=\overline{Z_{\theta\mathbb{T}}[\Phi^{\theta}\cdot P]} .
\]
The identity of \eqref{it:lowest-plane}, resp. \eqref{it:highest-plane}, on the torus $\theta\mathbb{T}$ and with the entire factor $\Phi^{\theta}$ gives
\[
Z_{\theta\mathbb{T}}[\Phi^{\theta}\cdot P]=\sum_{a\in I}Z_{\theta\mathbb{T}}\left[(\Phi^{\theta})^{[\eta_a]}\cdot P_a\right].
\]
For $a\in I$ the function $\Phi_a:=\left((\Phi^{\theta})^{[\eta_a]}\right)^{\theta}$ is an entire factor on $\theta(\theta\mathbb{T})=\mathbb{T}$ with $\Phi_a^{\theta}=(\Phi^{\theta})^{[\eta_a]}$;
Lemma~\ref{lem:reflection-identities}\eqref{it:reflected-functional}, applied to $\Phi_a$ and $P_a$, gives
\[
Z_{\mathbb{T}}[\Phi_a\cdot\Theta_\Pi P_a]=\overline{Z_{\theta\mathbb{T}}\left[(\Phi^{\theta})^{[\eta_a]}\cdot P_a\right]} .
\]
Taking complex conjugates in the second identity and inserting the result into the first, one finds
\[
Z_{\mathbb{T}}[\Phi\cdot\Theta_\Pi P]=\sum_{a\in I}\overline{Z_{\theta\mathbb{T}}\left[(\Phi^{\theta})^{[\eta_a]}\cdot P_a\right]}=\sum_{a\in I}Z_{\mathbb{T}}[\Phi_a\cdot\Theta_\Pi P_a].
\]
It remains to identify $\Phi_a$.
In the factor $(\Phi^{\theta})^{[\eta_a]}$ on $\theta\mathbb{T}$ the cochain $\eta_a$ stands for its periodisation over $\theta\Gamma$,
and $\theta^*$ maps this periodisation to the periodisation over $\Gamma$ of $\theta^*d\delta_a$.
By Lemma~\ref{lem:reflection-identities}\eqref{it:reflected-shift}, therefore, $\Phi_a=\Phi^{[-\theta^*d\delta_a]}$, that is, $\Phi_a(F)=\Phi(F+i\,\theta^*d\delta_a)$.

\eqref{it:plane-estimate}
Let $t\ge0$, and let $i$ be such that $D(n_i)$ carries $\mathcal{Q}_{\mathfrak{d}}$ and $(B')^{+h}$ embeds in $D(n_i)$.
Since $B'\subset(B')^{+h}$, the box $B'$ embeds in $D(n_i)$ as well (Definition~\ref{def:integrands}).
The torus $D(n_i)$ contains $n_i^3/|Q_0|$ cells of $\mathcal{Q}_{\mathfrak{d}}$;
on it the folded array $O_{\mathfrak{d}}$ is the sum of the mirror copies of $O$ in these cells, each of which is a trigonometric polynomial,
and $|O_{\mathfrak{d}}|\le|O|_1\cdot n_i^3/|Q_0|$ on real $F$ (Definition~\ref{def:arrays}).
Thus $\exp(tO_{\mathfrak{d}})$ is an entire factor on $D(n_i)$ (Definition~\ref{def:polynomials}),
and Lemma~\ref{lem:perturbed-denominators} together with Theorem~\ref{thm:partition-positivity} gives
\[
Z_{D(n_i)}\left[\exp(tO_{\mathfrak{d}})\right]\ge\exp\left(-t\cdot|O|_1\cdot\frac{n_i^3}{|Q_0|}\right)\cdot Z_{D(n_i)}[1]>0 .
\]
The elements $P$, $P_a$ and $\varphi_\beta\cdot P_a$, $\beta\in C^2_{\mathrm{c}}$, embed in $D(n_i)$, because $P$ and the $P_a$ are read in $B'$;
so their values under $\mu^t_i$ are defined (Definition~\ref{def:data}).

The identity \eqref{eq:elimination-universal} of $\mathcal{U}$, with the torus $D(n_i)$ and the entire factor $\exp(tO_{\mathfrak{d}})$, reads
\[
Z_{D(n_i)}\left[\exp(tO_{\mathfrak{d}})\cdot P\right]=\sum_{a\in I}Z_{D(n_i)}\left[\left(\exp(tO_{\mathfrak{d}})\right)^{[\eta_a]}\cdot P_a\right].
\]
Let $a\in I$.
The cochain $\eta_a$ is supported on plaquettes contained in $B'$ (Definition~\ref{def:universal-identity}).
Proposition~\ref{prop:contour-shift}\eqref{it:shift-arrays}, applied to $\mathcal{Q}_{\mathfrak{d}}$, $O$ and $\eta_a$, with $B_{\eta_a}:=B'$, provides complex numbers $b_{a,\alpha}(t)$ and cochains $\beta_{a,\alpha}\in C^2_{\mathrm{c}}$ supported in $(B')^{+h}$, indexed by a countable set, such that
\[
\left(\exp(tO_{\mathfrak{d}})\right)^{[\eta_a]}=\exp(tO_{\mathfrak{d}})\cdot\left(1+\sum_\alpha b_{a,\alpha}(t)\cdot\varphi_{\beta_{a,\alpha}}\right),
\qquad
\sum_\alpha\left|b_{a,\alpha}(t)\right|\le\exp\left(C_{\mathcal{Q}_{\mathfrak{d}}}(\eta_a)\cdot t\right)-1,
\]
the series converging uniformly on real $F$;
after multiplication by the bounded function $\exp(tO_{\mathfrak{d}})$ the convergence is still uniform.
As $P_a$ is read in $B'$, which embeds in $D(n_i)$, the map $\Phi\mapsto Z_{D(n_i)}[\Phi\cdot P_a]$ is continuous for the supremum norm on bounded Borel functions $\Phi$ (Lemma~\ref{lem:torus-elements}\eqref{it:torus-convergence}),
and the series may be integrated term by term.
Summation over $a\in I$ gives
\begin{equation}\label{eq:elimination-expansion}
Z_{D(n_i)}\left[\exp(tO_{\mathfrak{d}})\cdot P\right]
=\sum_{a\in I}\left(Z_{D(n_i)}\left[\exp(tO_{\mathfrak{d}})\cdot P_a\right]+\sum_\alpha b_{a,\alpha}(t)\cdot Z_{D(n_i)}\left[\exp(tO_{\mathfrak{d}})\cdot\varphi_{\beta_{a,\alpha}}\cdot P_a\right]\right).
\end{equation}
Divide \eqref{eq:elimination-expansion} by the positive number $Z_{D(n_i)}[\exp(tO_{\mathfrak{d}})]$ and recall $P^{\flat}=\sum_{a\in I}P_a$:
\[
\mu^t_i(P)-\mu^t_i(P^{\flat})=\sum_{a\in I}\sum_\alpha b_{a,\alpha}(t)\cdot\mu^t_i\left(\varphi_{\beta_{a,\alpha}}\cdot P_a\right).
\]
Every $\beta_{a,\alpha}$ is supported in $(B')^{+h}$, so that $\left|\mu^t_i(\varphi_{\beta_{a,\alpha}}\cdot P_a)\right|\le\sup_\beta\left|\mu^t_i(\varphi_\beta\cdot P_a)\right|$, the supremum being taken over all $\beta\in C^2_{\mathrm{c}}$ supported in $(B')^{+h}$.
Since $C_{\mathcal{Q}_{\mathfrak{d}}}(\eta_a)\le C$ and $t\ge0$,
\[
\left|\mu^t_i(P)-\mu^t_i(P^{\flat})\right|
\le\sum_{a\in I}\left(\exp\left(C_{\mathcal{Q}_{\mathfrak{d}}}(\eta_a)\cdot t\right)-1\right)\cdot\sup_\beta\left|\mu^t_i(\varphi_\beta\cdot P_a)\right|
\le(\exp(Ct)-1)\cdot\sum_{a\in I}\sup_\beta\left|\mu^t_i(\varphi_\beta\cdot P_a)\right| ,
\]
which is \eqref{eq:elimination-estimate}.

Now let $P$, $I$ and $(\delta_a,P_a)_{a\in I}$ be as in \eqref{it:lowest-plane} or as in \eqref{it:highest-plane}, and let $\Pi$ be a mirror with $\theta=\theta_\Pi$.
The triple $(P;(d\delta_a,P_a)_{a\in I};B^{+8})$ is a universal identity.
Indeed, $P$ is read in $B\subset B^{+8}$, and the $P_a$ belong to $\mathcal{P}$ and are read in $B^{+8}$.
Each $d\delta_a$ is a real element of $C^2_{\mathrm{c}}$ whose support consists of plaquettes containing a link of the support of $\delta_a$.
Such a link is $(j+1,b)=(j,b)+e_1$ in the case of \eqref{it:lowest-plane} and $(j-1,b)=(j,b)-e_1$ in the case of \eqref{it:highest-plane},
for a link $b$ of the slice with $(\delta'_a)_j(b)\neq0$.
Then $P^{\uparrow}$, resp. $P^{\downarrow}$, reads the link $(j,b)$;
as this element is read in $B^{+6}$ (Proposition~\ref{prop:schwinger-dyson}\eqref{it:replacement-identity}, applied six times), the link $(j,b)$ is contained in $B^{+6}$.
A plaquette containing $(j,b)\pm e_1$ therefore lies in $(j,b)\pm e_1+[-1,1]^3\subset(j,b)+[-2,2]^3\subset B^{+8}$.
Finally, \eqref{eq:elimination-universal} with $\eta_a=d\delta_a$ is the identity of \eqref{it:lowest-plane}, resp. \eqref{it:highest-plane}, since $\Phi^{(\delta_a)}=\Phi^{[d\delta_a]}$.
With $B'=B^{+8}$ and $C=\max_aC_{\mathcal{Q}_{\mathfrak{d}}}(d\delta_a)$, the condition on $i$ in \eqref{it:plane-estimate} is the one under which \eqref{eq:elimination-estimate} holds for this universal identity,
and \eqref{eq:elimination-estimate} becomes the first of the two inequalities in \eqref{it:plane-estimate}.

The triple $(\Theta_\Pi P;(-\theta^*d\delta_a,\Theta_\Pi P_a)_{a\in I};\theta(B^{+8}))$ is a universal identity as well.
The elements $\Theta_\Pi P$ and $\Theta_\Pi P_a$ belong to $\mathcal{P}$ and are read in $\theta(B^{+8})$ (Lemma~\ref{lem:reflection-identities}).
The cochain $-\theta^*d\delta_a\in C^2_{\mathrm{c}}$ is real and supported on $\theta(\supp d\delta_a)$ (Definition~\ref{def:factors}), a set of plaquettes contained in $\theta(B^{+8})$.
The identity \eqref{eq:elimination-universal} for this triple is \eqref{it:reflected-elimination}.
Since $\Theta_\Pi$ is additive, the sum of the elements on the right side of this universal identity is $\sum_{a\in I}\Theta_\Pi P_a=\Theta_\Pi P^{\flat}$.
So \eqref{eq:elimination-estimate} for it, with $B'=\theta(B^{+8})$ and $C'=\max_aC_{\mathcal{Q}_{\mathfrak{d}}}(-\theta^*d\delta_a)$ in the place of $C$,
holds for all $i$ such that $D(n_i)$ carries $\mathcal{Q}_{\mathfrak{d}}$ and $(\theta(B^{+8}))^{+h}$ embeds in $D(n_i)$, and it is the second inequality.
\end{proof}

\subsection{The class $\mathcal{P}$}\label{sec:polynomial-class}

Recall from Definition~\ref{def:covariant-classes} that each $P\in\mathcal{P}$ is, in exactly one way, a finite sum
\[
P=\sum_a\lambda_a\cdot\exp(i\langle\delta_a,\hat A\rangle)\cdot\Psi_a
\]
of basis elements with coefficients $\lambda_a\neq0$ and pairwise different pairs $(\delta_a,\Psi_a)$, where $d^{*}\delta_a=q\cdot\mathfrak{q}(\Psi_a)$;
this is the representation of $P$.
For a link $b$ of $\Z^3$ let $u_b:=\exp(iqA(b))$ be its phase;
if $\hat A$ is a local lift on a box that contains $b$, then $u_b=\exp(iq\hat A(b))$ (Definition~\ref{def:integrands}).

\begin{proposition}[Structure and stability]\label{prop:class-structure}
Let $\rho$ and $M$ be real, and let $w_0,\dots,w_3>0$ and $q>0$.
\begin{enumerate}
\item\label{it:class-span}
\emph{Structure.}
The space $\mathcal{P}$ is a commutative algebra with unit, and all its elements are even.
It is the linear span of the elements
\[
\varphi_\beta\cdot\exp\left(iq\langle\mathfrak{j},\hat A\rangle\right)\cdot\Psi
\qquad\text{with}\qquad
d^{*}\mathfrak{j}=\mathfrak{q}(\Psi),
\]
where $\beta\in C^2_{\mathrm{c}}$, $\mathfrak{j}\in C^1_{\mathrm{c}}(\Z)$ and $\Psi$ is a monomial.
Since $\exp(iq\langle\mathfrak{j},\hat A\rangle)=\prod_bu_b^{\mathfrak{j}(b)}$, the product being taken over the links $b$ of $\Z^3$,
the space $\mathcal{P}$ is the linear span of the products of a character of the curvature
with a gauge invariant polynomial in the link phases $u_b^{\pm1}=\exp(\pm iqA(b))$ and in the fields $\psi$, $\bar\psi$.
The exponent $\delta$ of every basis element of $\mathcal{P}$ lies in $d^{*}C^2_{\mathrm{c}}+q\cdot C^1_{\mathrm{c}}(\Z)$.
The coefficient function $\hat A\mapsto\exp(i\langle\delta,\hat A\rangle)$ of a basis element has modulus one for real $\hat A$;
it extends to an entire function of the complex variables $\hat A(b)$,
and with $\|\delta\|_1:=\sum_b|\delta(b)|$ one has, for every $r\geq0$ and every complex $\hat A$ with $|\Im\hat A(b)|\leq r$ for all links $b$,
\[
\left|\exp(i\langle\delta,\hat A\rangle)\right|\leq\exp\left(\|\delta\|_1\cdot r\right).
\]
\item\label{it:class-stable}
\emph{Stability.}
Each of the following maps sends $\mathcal{P}$ into $\mathcal{P}$:
the multiplication by an element of $\mathcal{P}$, in particular by a character $\varphi_\beta$, $\beta\in C^2_{\mathrm{c}}$,
and by the element $\mathfrak{V}_\delta$ of Definition~\ref{def:shifted-elements}
for every finite set $S$ of links contained in one coordinate plane and every $\delta\in\R^S$;
for every real $\delta\in C^1_{\mathrm{c}}$, the map $P\mapsto P^{(\delta)}$,
where $P^{(\delta)}$ arises from $P$ when $\hat A$ is replaced by $\hat A-i\delta$ in the coefficient functions;
if $\rho^2\neq1$, the maps $\mathfrak{R}^{+}_{\mathbb{L}}$ and $\mathfrak{R}^{-}_{\mathbb{L}}$ of Proposition~\ref{prop:schwinger-dyson}\eqref{it:replacement-maps} for all thin layers $\mathbb{L}$,
and the maps $\mathfrak{R}^{\mu,+}$ and $\mathfrak{R}^{\mu,-}$ of Definition~\ref{def:transported-objects} for the thin layers of the directions $e_2$ and $e_3$;
the map $\Theta_\Pi$ of Definition~\ref{def:grassmann-reflections} and \eqref{eq:setting-reflection-coefficients} for every mirror $\Pi$,
and hence the map $G\mapsto G^g$ of \eqref{eq:setting-group-action} for every $g\in W$.
Explicitly, let $\Pi=(\ell,c)$ be a mirror, $\theta:=\theta_\Pi$ and $n:=\hat n_\ell$,
let $\delta\in C^1_{\mathrm{c}}$ be real, and let $\Psi=\xi_1\cdots\xi_k$ be a monomial with the generators $\xi_1,\dots,\xi_k$.
Then
\[
\Theta_\Pi\left(\exp(i\langle\delta,\hat A\rangle)\cdot\Psi\right)=\exp\left(-i\langle\theta^{*}\delta,\hat A\rangle\right)\cdot\Theta_\Pi(\xi_k)\cdots\Theta_\Pi(\xi_1),
\]
where, for $x\in\mathbb{X}$ and $\alpha\in\{1,\dots,4\}$,
\[
\Theta_\Pi\psi_\alpha(x)=-\sum_{\alpha'=1}^{4}\bar\psi_{\alpha'}(\theta x)\,(\gamma_n)_{\alpha'\alpha},
\qquad
\Theta_\Pi\bar\psi_\alpha(x)=-\sum_{\alpha'=1}^{4}(\gamma_n)_{\alpha\alpha'}\,\psi_{\alpha'}(\theta x);
\]
in particular $\Theta_\Pi\varphi_\beta=\varphi_{-\theta^{*}\beta}$ for $\beta\in C^2_{\mathrm{c}}$.
If $P\in\mathcal{P}$ is read in a box $B$ (Definition~\ref{def:enlarged-boxes}),
then $\Theta_\Pi P$ is read in $\theta B$, and $P^g$ is read in $gB$ for every $g\in W$.

\emph{Even elements.}
Let $g=\theta_\Pi\theta_{\Pi'}$ for two mirrors $\Pi=(\ell,c)$ and $\Pi'=(\ell',c')$ with the upward unit normals $n=\hat n_\ell$ and $n'=\hat n_{\ell'}$;
the matrices $g_2$ and $g_3$ of Definition~\ref{def:other-directions} are of this form,
and every even element of $W$ is a product of elements of this form.
Then $G\mapsto G^g=\Theta_\Pi\Theta_{\Pi'}G$ is a linear automorphism of the algebra $\mathcal{P}$ with the following properties,
in which $g_{*}$ is the push-forward of cochains of Definition~\ref{def:factors}.
For every basis element $\exp(i\langle\delta,\hat A\rangle)\cdot\Psi$,
\[
\left(\exp(i\langle\delta,\hat A\rangle)\cdot\Psi\right)^g=\exp(i\langle g_{*}\delta,\hat A\rangle)\cdot\Psi^g,
\]
where $\Psi^g$ is the product, in the order of the factors of $\Psi$, of the images of the generators of $\Psi$ under the substitution
\[
\psi(x)\mapsto\gamma_{n'}\gamma_n\,\psi(gx),
\qquad
\bar\psi(x)\mapsto\bar\psi(gx)\,\gamma_n\gamma_{n'}
\]
of \eqref{eq:setting-two-reflections}.
For $\beta\in C^2_{\mathrm{c}}$ one has $(\varphi_\beta)^g=\varphi_{g_{*}\beta}$, and for every bond $\{x,x'\}$
\[
\begin{aligned}
&\left(\bar\psi(x)\cdot w(x,x')\cdot U_A(x,x')\cdot\psi(x')\right)^g\\
&\qquad=\bar\psi(gx)\cdot w(gx,gx')\cdot U_A(gx,gx')\cdot\psi(gx').
\end{aligned}
\]
Let $S$ be a finite set of links contained in one coordinate plane, and let $\delta\in\R^S$.
Then $gS$ is again a finite set of links contained in one coordinate plane, the cochain $g_{*}\delta$ is supported on $gS$, and
\[
(\mathfrak{Q}_\delta)^g=\mathfrak{Q}_{g_{*}\delta},
\qquad
(\mathfrak{V}_\delta)^g=\mathfrak{V}_{g_{*}\delta}.
\]
For $P\in\mathcal{P}$ and every real $\delta\in C^1_{\mathrm{c}}$,
\[
\left(P^{(\delta)}\right)^g=\left(P^g\right)^{(g_{*}\delta)}.
\]
Consequently, for $P\in\mathcal{P}$ and for $S$ and $\delta\in\R^S$ as above,
the image under $g$ of the element
\[
\exp\left(\frac{1}{2}\|d\delta\|^2\right)\cdot\varphi_{d\delta}\cdot\mathfrak{V}_\delta\cdot P^{(\delta)}
\]
that stands on the right side of the identity in Proposition~\ref{prop:contour-shift}\eqref{it:shift-identity}
is the element formed in the same way from $P^g$ and $g_{*}\delta$;
and for $\mu\in\{2,3\}$ the image under $g_\mu$ of an element $P_a$ of Proposition~\ref{prop:outer-plane}\eqref{it:lowest-plane} or~\eqref{it:highest-plane}
is given by \eqref{eq:elimination-new-terms} with $\delta'_a$, $\delta_a$ and $\Psi_a$ replaced by $g_{\mu*}\delta'_a$, $g_{\mu*}\delta_a$ and $\Psi_a^{g_\mu}$.
\end{enumerate}
\end{proposition}

\begin{proof}
\eqref{it:class-span}
\emph{Evenness, products and unit.}
For every basis element $\exp(i\langle\delta,\hat A\rangle)\cdot\Psi$ one has $\sum_{y\in\Z^3}\mathfrak{q}(\Psi)(y)=0$,
and every element of $\mathcal{P}$ is even (Subsection~\ref{sec:covariant-classes}).
Let $\exp(i\langle\delta,\hat A\rangle)\cdot\Psi$ and $\exp(i\langle\delta',\hat A\rangle)\cdot\Psi'$ be two basis elements.
The coefficient functions are scalar, so the product of the two elements is $\exp(i\langle\delta+\delta',\hat A\rangle)\cdot\Psi\Psi'$.
If $\Psi$ and $\Psi'$ have a generator in common, then $\Psi\Psi'=0$;
otherwise $\Psi\Psi'$ is $\pm$ the monomial whose generators are those of $\Psi$ together with those of $\Psi'$,
and its charge density is $\mathfrak{q}(\Psi)+\mathfrak{q}(\Psi')$, because the numbers counted in the definition of the charge density (Definition~\ref{def:chamber-operator}) add up.
As $d^{*}(\delta+\delta')=q\cdot\left(\mathfrak{q}(\Psi)+\mathfrak{q}(\Psi')\right)$, the product is $0$ or $\pm$ a basis element,
and by bilinearity $\mathcal{P}$ is closed under products.
The unit $1=\exp(i\langle0,\hat A\rangle)\cdot1$ is a basis element, the empty monomial having charge density zero.
Even elements of a Grassmann algebra commute with each other, and the coefficient functions are scalar;
so $\mathcal{P}$ is a commutative algebra with unit.

\emph{Characters and integer exponents.}
Let $\exp(i\langle\delta,\hat A\rangle)\cdot\Psi$ be a basis element.
If $\delta=0$, then $\mathfrak{q}(\Psi)=0$, and the element is $\varphi_0\cdot\exp(iq\langle0,\hat A\rangle)\cdot\Psi$.
Let $\delta\neq0$.
Let $\mathbb{B}$ be the smallest integer box (Definition~\ref{def:integer-box}) that contains the links of $\supp\delta$ and the sites of $\supp\mathfrak{q}(\Psi)$,
and let $V$ be the set of the sites of $\Z^3$ that lie in $\mathbb{B}$.
As $\mathbb{B}$ is a product of three integer intervals, the graph formed by $V$ and the links with both end points in $V$ is connected;
the integer function $\mathfrak{q}(\Psi)$ vanishes outside $V$, and its sum is zero (Subsection~\ref{sec:covariant-classes}).
Lemma~\ref{lem:lines-fillings}\eqref{it:integer-lines} provides $\mathfrak{j}\in C^1_{\mathrm{c}}(\Z)$,
supported on links with both end points in $V$, that is, on links contained in $\mathbb{B}$, with $d^{*}\mathfrak{j}=\mathfrak{q}(\Psi)$.
The real one-cochain $\delta-q\cdot\mathfrak{j}$ is supported on links contained in $\mathbb{B}$,
and $d^{*}(\delta-q\cdot\mathfrak{j})=q\cdot\mathfrak{q}(\Psi)-q\cdot\mathfrak{q}(\Psi)=0$;
so Lemma~\ref{lem:lines-fillings}\eqref{it:integer-filling} provides a real two-cochain $\beta$,
supported on plaquettes contained in $\mathbb{B}$, with $d^{*}\beta=\delta-q\cdot\mathfrak{j}$.
If $\hat A$ is a local lift on a box that contains $\mathbb{B}$, then $d\hat A=F$ on the plaquettes of $\mathbb{B}$ (Definition~\ref{def:integrands}),
hence $\langle d^{*}\beta,\hat A\rangle=\langle\beta,d\hat A\rangle=\langle\beta,F\rangle$, and
\[
\exp(i\langle\delta,\hat A\rangle)\cdot\Psi
=\exp(i\langle d^{*}\beta,\hat A\rangle)\cdot\exp(iq\langle\mathfrak{j},\hat A\rangle)\cdot\Psi
=\varphi_\beta\cdot\exp(iq\langle\mathfrak{j},\hat A\rangle)\cdot\Psi .
\]
Conversely, let $\beta$, $\mathfrak{j}$ and $\Psi$ be as in the statement.
Then $\varphi_\beta\cdot\exp(iq\langle\mathfrak{j},\hat A\rangle)\cdot\Psi=\exp(i\langle d^{*}\beta+q\cdot\mathfrak{j},\hat A\rangle)\cdot\Psi$,
the cochain $d^{*}\beta+q\cdot\mathfrak{j}\in C^1_{\mathrm{c}}$ is real,
and $d^{*}(d^{*}\beta+q\cdot\mathfrak{j})=q\cdot\mathfrak{q}(\Psi)$ because $d^{*}d^{*}=0$;
thus the element is a basis element of $\mathcal{P}$.
This proves that $\mathcal{P}$ is the stated linear span, and that every exponent lies in $d^{*}C^2_{\mathrm{c}}+q\cdot C^1_{\mathrm{c}}(\Z)$.

\emph{Link phases.}
The values $\mathfrak{j}(b)$ are integers, and $\exp(iq\hat A(b))=u_b$ on a local lift;
this gives the product formula, in which $u_b^{-1}=\exp(-iqA(b))$ occurs where $\mathfrak{j}(b)<0$.
The gauge transformation \eqref{eq:setting-gauge} by a real function $\lambda$ on $\Z^3$ replaces $u_b$ by $u_b\cdot\exp(-iq\cdot(d\lambda)(b))$
and $\Psi$ by $\exp(iq\langle\mathfrak{q}(\Psi),\lambda\rangle)\cdot\Psi$;
it therefore multiplies $\prod_bu_b^{\mathfrak{j}(b)}\cdot\Psi$ by
\[
\exp\left(-iq\langle\mathfrak{j},d\lambda\rangle\right)\cdot\exp\left(iq\langle\mathfrak{q}(\Psi),\lambda\rangle\right)
=\exp\left(iq\langle\mathfrak{q}(\Psi)-d^{*}\mathfrak{j},\lambda\rangle\right),
\]
a number that equals $1$ for every $\lambda$ if and only if $d^{*}\mathfrak{j}=\mathfrak{q}(\Psi)$.
The products $\prod_bu_b^{\mathfrak{j}(b)}\cdot\Psi$ with pairwise different pairs $(\mathfrak{j},\Psi)$ are linearly independent
(different monomials are linearly independent, and so are the functions $A\mapsto\exp(iq\langle\mathfrak{j},A\rangle)$ for different $\mathfrak{j}$),
and each gauge transformation multiplies each of them by a number;
so a polynomial in the link phases $u_b^{\pm1}$ and in the generators is gauge invariant
if and only if it is a linear combination of such products with $d^{*}\mathfrak{j}=\mathfrak{q}(\Psi)$.
Together with the description of $\mathcal{P}$ just proved, this gives its description by characters and gauge invariant polynomials.

\emph{Coefficient functions.}
For real $\hat A$ the number $\langle\delta,\hat A\rangle$ is real, and $|\exp(i\langle\delta,\hat A\rangle)|=1$.
The function $\exp\left(i\sum_b\delta(b)\cdot\hat A(b)\right)$ depends on the finitely many variables $\hat A(b)$, $b\in\supp\delta$, and is entire in them;
if $|\Im\hat A(b)|\leq r$ for every link $b$, then
\[
\left|\exp(i\langle\delta,\hat A\rangle)\right|=\exp\left(-\sum_b\delta(b)\cdot\Im\hat A(b)\right)\leq\exp\left(\|\delta\|_1\cdot r\right).
\]

\eqref{it:class-stable}
\emph{Products and shifts.}
By \eqref{it:class-span}, $\mathcal{P}$ is closed under products, and $\varphi_\beta$ is the element of \eqref{it:class-span} with $\mathfrak{j}=0$ and $\Psi=1$.
Let $S$ be a finite set of links contained in one coordinate plane, and let $\delta\in\R^S$.
The hypotheses of Proposition~\ref{prop:contour-shift}, among them \eqref{eq:elimination-shift-box}, are satisfied by this set $S$,
by the element $1$, which is read in every box, by a box $B$ that contains $b+[-1,1]^3$ for every $b\in S$,
and by a cubic torus of sufficiently large side, in which $B$ embeds.
By Proposition~\ref{prop:contour-shift}\eqref{it:shift-polynomial}, $\mathfrak{V}_\delta$ belongs to $\mathcal{P}$.
Now let $P=\sum_a\lambda_a\cdot\exp(i\langle\delta_a,\hat A\rangle)\cdot\Psi_a$ be the representation of $P\in\mathcal{P}$, and let $\delta\in C^1_{\mathrm{c}}$ be an arbitrary real cochain.
As $\exp(i\langle\delta_a,\hat A-i\delta\rangle)=\exp(\langle\delta_a,\delta\rangle)\cdot\exp(i\langle\delta_a,\hat A\rangle)$,
\[
P^{(\delta)}=\sum_a\lambda_a\cdot\exp(\langle\delta_a,\delta\rangle)\cdot\exp(i\langle\delta_a,\hat A\rangle)\cdot\Psi_a ,
\]
a linear combination of the same basis elements with other coefficients.

\emph{One mirror.}
Let $\Pi=(\ell,c)$, $\theta$, $n$, $\delta$ and $\Psi=\xi_1\cdots\xi_k$ be as in the statement.
By \eqref{eq:setting-reflection-coefficients}, $\Theta_\Pi$ replaces the coefficient function $\exp(i\langle\delta,\hat A\rangle)$ of $\Psi$
by the complex conjugate of $\exp(i\langle\delta,\theta^{*}\hat A\rangle)$, that is, for real $\hat A$, by
\[
\exp\left(-i\langle\delta,\theta^{*}\hat A\rangle\right)=\exp\left(-i\langle\theta^{*}\delta,\hat A\rangle\right);
\]
indeed $\theta^{*}$ is orthogonal and $(\theta^{*})^2=(\theta^2)^{*}=1$ (Definition~\ref{def:cochains}),
so that $\theta^{*}$ is its own adjoint for the pairing $\langle\cdot,\cdot\rangle$.
By \eqref{eq:setting-reflection-fields}, $\Theta_\Pi$ reverses the order of products and acts on the generators by the two formulas of the statement,
which are \eqref{eq:setting-reflection-fields} written in components.
This proves the formula for $\Theta_\Pi\left(\exp(i\langle\delta,\hat A\rangle)\cdot\Psi\right)$.

Let now $\exp(i\langle\delta,\hat A\rangle)\cdot\Psi$ be a basis element.
The image under $\Theta_\Pi$ of a generator at the site $x$ is a linear combination of generators of the other kind at the site $\theta x$,
and $\pi(\theta x)=\theta\pi(x)$ by Proposition~\ref{prop:geometry}\eqref{it:geometry-projection}.
So a factor $\psi_\alpha(x)$ of $\Psi$ contributes a generator $\bar\psi_{\alpha'}(\theta x)$ over the photon site $\theta\pi(x)$,
and a factor $\bar\psi_\alpha(x)$ contributes a generator $\psi_{\alpha'}(\theta x)$ over the same photon site.
Since $\theta$ is an involution, every non-zero product of generators in the expansion of $\Theta_\Pi(\xi_k)\cdots\Theta_\Pi(\xi_1)$
has the charge density $-\mathfrak{q}(\Psi)\circ\theta$.
The map $\theta^{*}$ commutes with $d$ (Definition~\ref{def:cochains});
taking adjoints in $\theta^{*}d=d\,\theta^{*}$ and using that $\theta^{*}$ is its own adjoint, one finds $d^{*}\theta^{*}=\theta^{*}d^{*}$.
On zero-cochains $\theta^{*}$ is the composition with $\theta$.
Therefore
\[
d^{*}\left(-\theta^{*}\delta\right)=-\theta^{*}d^{*}\delta=-q\cdot\theta^{*}\mathfrak{q}(\Psi)=q\cdot\left(-\mathfrak{q}(\Psi)\circ\theta\right),
\]
and $\Theta_\Pi\left(\exp(i\langle\delta,\hat A\rangle)\cdot\Psi\right)$ is a linear combination of basis elements of $\mathcal{P}$ with the exponent $-\theta^{*}\delta$.
Since $\Theta_\Pi$ is additive and antilinear, it maps $\mathcal{P}$ into $\mathcal{P}$.
For $\delta=d^{*}\beta$ and $\Psi=1$ the formula reads
$\Theta_\Pi\varphi_\beta=\exp(-i\langle\theta^{*}d^{*}\beta,\hat A\rangle)=\exp(i\langle d^{*}(-\theta^{*}\beta),\hat A\rangle)=\varphi_{-\theta^{*}\beta}$.
The group $W$ is generated by the reflections in the mirrors (Definition~\ref{def:mirrors}), each of which is its own inverse;
so every $g\in W$ is a word $g=\theta_{\Pi_1}\cdots\theta_{\Pi_r}$ in reflections.
As the elements of $\mathcal{P}$ are even, \eqref{eq:setting-group-action} gives $G^g=\Theta_{\Pi_1}\cdots\Theta_{\Pi_r}G$ for $G\in\mathcal{P}$,
and each of the maps $\Theta_{\Pi_1},\dots,\Theta_{\Pi_r}$ sends $\mathcal{P}$ into $\mathcal{P}$.

\emph{Reading sets.}
Let $P=\sum_a\lambda_a\cdot\exp(i\langle\delta_a,\hat A\rangle)\cdot\Psi_a$ be read in the box $B$:
the links read by $P$ are contained in $B$, and the sites $x$ read by $P$ lie in $\mathbb{X}(B)$, that is, their cubes $\pi(x)+\left[-\frac{1}{2},\frac{1}{2}\right]^3$ are contained in $B$.
The image of a box under $\theta$ is a box, because the linear part $\vec\theta\in\Ocub$ maps every cube centred at the origin onto itself.
The reflection $\theta$ maps the links contained in $B$ onto the links contained in $\theta B$,
and the cube $\pi(x)+\left[-\frac{1}{2},\frac{1}{2}\right]^3$ onto $\theta\pi(x)+\left[-\frac{1}{2},\frac{1}{2}\right]^3=\pi(\theta x)+\left[-\frac{1}{2},\frac{1}{2}\right]^3$;
hence $x\in\mathbb{X}(B)$ if and only if $\theta x\in\mathbb{X}(\theta B)$.
The support of $\theta^{*}\delta_a$ is $\theta(\supp\delta_a)$, because $(\theta^{*}\delta_a)(b)=\pm\delta_a(\theta b)$ for every link $b$.
By the formula for $\Theta_\Pi$, every basis element in the representation of $\Theta_\Pi P$ has an exponent $-\theta^{*}\delta_a$,
whose links are contained in $\theta B$, and a monomial whose generators sit at sites $\theta x$ with $x$ read by $P$, which lie in $\mathbb{X}(\theta B)$.
So $\Theta_\Pi P$ is read in $\theta B$.
Induction on the length $r$ of a word $g=\theta_{\Pi_1}\cdots\theta_{\Pi_r}$ shows that $P^g=\Theta_{\Pi_1}\cdots\Theta_{\Pi_r}P$ is read in $\theta_{\Pi_1}\cdots\theta_{\Pi_r}B=gB$.

\emph{Replacement maps.}
If $\rho^2\neq1$, the maps $\mathfrak{R}^{\pm}_{\mathbb{L}}$ are linear maps of $\mathcal{P}$ into $\mathcal{P}$ by Proposition~\ref{prop:schwinger-dyson}\eqref{it:replacement-maps}.
Each map $\mathfrak{R}^{\mu,\pm}$ of Definition~\ref{def:transported-objects} is the composition of $G\mapsto G^{g_\mu^{-1}}$,
of a map $\mathfrak{R}^{\pm}_{\mathbb{L}}$ and of $G\mapsto G^{g_\mu}$, each of which maps $\mathcal{P}$ into $\mathcal{P}$.

\emph{Even elements.}
Let $g=\theta_\Pi\theta_{\Pi'}$, and let $\vec\theta_\Pi$ and $\vec\theta_{\Pi'}$ be the linear parts of $\theta_\Pi$ and $\theta_{\Pi'}$,
the linear reflections with the normals $n$ and $n'$;
then $\vec g=\vec\theta_\Pi\vec\theta_{\Pi'}$ is the linear part of $g$.
By Subsection~\ref{sec:other-directions}, the matrices $g_2$ and $g_3$ are of this form,
the two mirrors being the diagonal mirrors $(x_1-x_2,0)$ and $(x_2-x_3,0)$.
The linear part of a reflection has determinant $-1$;
so an even element of $W$, written as a word in reflections, is a word of even length,
and it is the product of the elements formed by the consecutive pairs of its letters, each of which is a product of two reflections.

By \eqref{eq:setting-reflection-coefficients} and \eqref{eq:setting-reflection-fields}, the maps $\Theta_\Pi$ and $\Theta_{\Pi'}$ are antilinear
and reverse the order of products, also on elements with coefficient functions;
hence $G\mapsto G^g=\Theta_\Pi\Theta_{\Pi'}G$ is linear and multiplicative.
It maps $\mathcal{P}$ into $\mathcal{P}$, and so does $G\mapsto G^{g^{-1}}$, which is its inverse by \eqref{eq:setting-group-action};
so it is a linear automorphism of the algebra $\mathcal{P}$.

Applied twice, \eqref{eq:setting-reflection-coefficients} carries a coefficient function $a(\hat A)$ first to $\overline{a(\theta_{\Pi'}^{*}\hat A)}$ and then to
\[
\overline{\overline{a\left(\theta_{\Pi'}^{*}\theta_\Pi^{*}\hat A\right)}}=a\left((\theta_\Pi\theta_{\Pi'})^{*}\hat A\right)=a\left(g^{*}\hat A\right),
\]
since $(\theta_\Pi\theta_{\Pi'})^{*}=\theta_{\Pi'}^{*}\theta_\Pi^{*}$ by Definition~\ref{def:cochains}.
For $a(\hat A)=\exp(i\langle\delta,\hat A\rangle)$ this is $\exp(i\langle\delta,g^{*}\hat A\rangle)=\exp(i\langle g_{*}\delta,\hat A\rangle)$,
because $g^{*}$ is orthogonal with the inverse $g_{*}$ (Definition~\ref{def:factors}), so that $g_{*}$ is the adjoint of $g^{*}$.
On the generators the map acts by \eqref{eq:setting-two-reflections};
being multiplicative, it maps $\Psi$ to the product, in the same order, of the images of the generators of $\Psi$.
This proves the formula for basis elements.
For $\delta=d^{*}\beta$ and $\Psi=1$ it gives
$(\varphi_\beta)^g=\exp(i\langle g_{*}d^{*}\beta,\hat A\rangle)=\exp(i\langle d^{*}g_{*}\beta,\hat A\rangle)=\varphi_{g_{*}\beta}$,
since $g_{*}$ commutes with $d^{*}$ (Definition~\ref{def:factors}).

\emph{Hopping terms.}
By Definition~\ref{def:spin-matrices}, $\gamma_{n'}\gamma_v\gamma_{n'}=-\gamma_{\vec\theta_{\Pi'}v}$ and $\gamma_n\gamma_{\vec\theta_{\Pi'}v}\gamma_n=-\gamma_{\vec\theta_\Pi\vec\theta_{\Pi'}v}$ for $v\in\R^3$;
hence
\[
\gamma_n\gamma_{n'}\,\gamma_v\,\gamma_{n'}\gamma_n=\gamma_{\vec gv},\qquad v\in\R^3 .
\]
Let $\{x,x'\}$ be a bond of type $i$, and let $\hat d$ be the unit vector from $x$ to $x'$.
If $x=g'v_0$ with $g'\in W$, then $x'=g's_iv_0$ (Definition~\ref{def:fermion-sites});
thus $gx'=gg's_iv_0$ is the neighbour of type $i$ of $gx=gg'v_0$, and the unit vector from $gx$ to $gx'$ is $\vec g\hat d$.
As $\gamma_n^2=\gamma_{n'}^2=1$,
\[
\gamma_n\gamma_{n'}\cdot w(x,x')\cdot\gamma_{n'}\gamma_n=\frac{w_i}{2}\left(\rho-\gamma_{\vec g\hat d}\right)=w(gx,gx').
\]
The coefficient function of $\bar\psi(x)\cdot w(x,x')\cdot U_A(x,x')\cdot\psi(x')$ is $U_A(x,x')=\exp\left(iq\langle1_{(\pi(x),\pi(x'))},\hat A\rangle\right)$ by \eqref{eq:elimination-pair}.
By the rule for coefficient functions just proved, by $(g^{*}\hat A)(y,y')=\hat A(gy,gy')$ for oriented pairs (Definition~\ref{def:cochains})
and by $\pi(gx)=g\pi(x)$, $\pi(gx')=g\pi(x')$ (Proposition~\ref{prop:geometry}\eqref{it:geometry-projection}), it goes to
\[
\exp\left(iq\hat A(g\pi(x),g\pi(x'))\right)=\exp\left(iq\hat A(\pi(gx),\pi(gx'))\right)=U_A(gx,gx').
\]
By \eqref{eq:setting-two-reflections} and linearity, the Grassmann part
$\bar\psi(x)\cdot w(x,x')\cdot\psi(x')=\sum_{\alpha,\alpha'}w(x,x')_{\alpha\alpha'}\cdot\bar\psi_\alpha(x)\,\psi_{\alpha'}(x')$ goes to
\[
\bar\psi(gx)\cdot\gamma_n\gamma_{n'}\cdot w(x,x')\cdot\gamma_{n'}\gamma_n\cdot\psi(gx')=\bar\psi(gx)\cdot w(gx,gx')\cdot\psi(gx').
\]
The two computations together give the formula for the hopping terms.

\emph{The elements $\mathfrak{Q}_\delta$ and $\mathfrak{V}_\delta$.}
Let $S$ be a finite set of links contained in one coordinate plane, let $\delta\in\R^S$,
and read $\delta$ as an odd function of oriented pairs of neighbouring photon sites (Definition~\ref{def:cochains}).
By Proposition~\ref{prop:geometry}\eqref{it:geometry-projection}, the bonds of type~$3$ are the bonds whose two ends have different projections,
these projections are neighbouring photon sites,
and the bonds over a link $b$ are the eight bonds of type~$3$ whose ends project onto the two end points of $b$.
For such a bond, ordered as $(x,x')$ in the orientation of $b$, one has $\delta(\pi(x),\pi(x'))=\delta_b$ and $\delta(\pi(x'),\pi(x))=-\delta_b$.
So the two terms that the bond contributes to $\mathfrak{Q}_\delta$ in Definition~\ref{def:shifted-elements} are the terms of the ordered pairs $(x,x')$ and $(x',x)$ in
\[
\begin{aligned}
\mathfrak{Q}_\delta=\sum_{(x,x')}&\left(\exp\left(q\cdot\delta(\pi(x),\pi(x'))\right)-1\right)\\
&\cdot\bar\psi(x)\cdot w(x,x')\cdot U_A(x,x')\cdot\psi(x'),
\end{aligned}
\]
the sum being taken over all ordered pairs $(x,x')$ for which $\{x,x'\}$ is a bond of type~$3$;
a pair over a link outside $S$ contributes zero, since $\delta$ vanishes on that link.
The map $(x,x')\mapsto(gx,gx')$ is a bijection of this set of ordered pairs onto itself, because $g$ and $g^{-1}$ preserve the types of bonds.
Apply the linear map $G\mapsto G^g$ to the finitely many non-zero terms and use the formula for the hopping terms;
then re-index the sum by this bijection and use $\pi(g^{-1}x)=g^{-1}\pi(x)$:
\[
\begin{aligned}
(\mathfrak{Q}_\delta)^g=\sum_{(x,x')}&\left(\exp\left(q\cdot\delta\left(g^{-1}\pi(x),g^{-1}\pi(x')\right)\right)-1\right)\\
&\cdot\bar\psi(x)\cdot w(x,x')\cdot U_A(x,x')\cdot\psi(x').
\end{aligned}
\]
By Definition~\ref{def:cochains}, $\delta(g^{-1}y,g^{-1}y')=\left((g^{-1})^{*}\delta\right)(y,y')=(g_{*}\delta)(y,y')$.
The cochain $g_{*}\delta$ is supported on the image $gS$ of $S$ (Definition~\ref{def:factors}),
and $gS$ is a finite set of links in one coordinate plane:
the linear part of $g$ is a signed permutation matrix and its translation part lies in $\Z^3$,
so $g$ maps every plane $\{x_\nu=j\}$, $j\in\Z$, onto a plane of the same kind.
The same expression for $\mathfrak{Q}_{g_{*}\delta}$ therefore gives $(\mathfrak{Q}_\delta)^g=\mathfrak{Q}_{g_{*}\delta}$.
By Proposition~\ref{prop:contour-shift}\eqref{it:shift-polynomial}, applied as in the proof that $\mathfrak{V}_\delta\in\mathcal{P}$,
$\mathfrak{Q}_\delta$ is an even nilpotent element, and $\mathfrak{V}_\delta=\exp(\mathfrak{Q}_\delta)$ is a finite exponential series in it;
the same holds for the set $gS$ and the cochain $g_{*}\delta$.
The linear and multiplicative map $G\mapsto G^g$ carries this series term by term into the exponential series of $(\mathfrak{Q}_\delta)^g=\mathfrak{Q}_{g_{*}\delta}$, that is,
\[
(\mathfrak{V}_\delta)^g=\exp\left((\mathfrak{Q}_\delta)^g\right)=\exp\left(\mathfrak{Q}_{g_{*}\delta}\right)=\mathfrak{V}_{g_{*}\delta}.
\]

\emph{Shifted elements.}
Let $P=\sum_a\lambda_a\cdot\exp(i\langle\delta_a,\hat A\rangle)\cdot\Psi_a$ be the representation of $P\in\mathcal{P}$, and let $\delta\in C^1_{\mathrm{c}}$ be real.
The formula for $P^{(\delta)}$ and the formula for basis elements give
\[
\begin{aligned}
\left(P^{(\delta)}\right)^g&=\sum_a\lambda_a\cdot\exp(\langle\delta_a,\delta\rangle)\cdot\exp(i\langle g_{*}\delta_a,\hat A\rangle)\cdot\Psi_a^g,\\
P^g&=\sum_a\lambda_a\cdot\exp(i\langle g_{*}\delta_a,\hat A\rangle)\cdot\Psi_a^g .
\end{aligned}
\]
Each $\exp(i\langle g_{*}\delta_a,\hat A\rangle)\cdot\Psi_a^g$ lies in $\mathcal{P}$, being the image of a basis element,
and all basis elements in its representation have the exponent $g_{*}\delta_a$;
on them $G\mapsto G^{(g_{*}\delta)}$ is the multiplication by $\exp(\langle g_{*}\delta_a,g_{*}\delta\rangle)=\exp(\langle\delta_a,\delta\rangle)$, the map $g_{*}$ being orthogonal.
Hence $\left(P^{(\delta)}\right)^g=\left(P^g\right)^{(g_{*}\delta)}$.

\emph{The two consequences.}
The map $g_{*}$ is orthogonal and commutes with $d$ and with $d^{*}$ (Definition~\ref{def:factors});
so $g_{*}d\delta=d(g_{*}\delta)$, $\|d(g_{*}\delta)\|=\|d\delta\|$, $g_{*}d^{*}d\delta=d^{*}d(g_{*}\delta)$ and $\langle g_{*}\delta',g_{*}\delta\rangle=\langle\delta',\delta\rangle$
for all real $\delta,\delta'\in C^1_{\mathrm{c}}$.
As $G\mapsto G^g$ is multiplicative, the formulas above give,
for every finite set $S$ of links contained in one coordinate plane and every $\delta\in\R^S$,
\[
\begin{aligned}
&\left(\exp\left(\frac{1}{2}\|d\delta\|^2\right)\cdot\varphi_{d\delta}\cdot\mathfrak{V}_\delta\cdot P^{(\delta)}\right)^g\\
&\qquad=\exp\left(\frac{1}{2}\|d(g_{*}\delta)\|^2\right)\cdot\varphi_{d(g_{*}\delta)}\cdot\mathfrak{V}_{g_{*}\delta}\cdot\left(P^g\right)^{(g_{*}\delta)} .
\end{aligned}
\]
Now let $\mu\in\{2,3\}$, and let $P_a$ be given by \eqref{eq:elimination-new-terms}.
The cochain $\delta_a$ is supported on the links of the plane $\{x_1=j+1\}$, respectively $\{x_1=j-1\}$, where $j$ is the integer of Proposition~\ref{prop:outer-plane};
so the formulas above apply to $\delta_a$, with the finite set $\supp\delta_a$ of links of this plane in the place of $S$,
and $\exp(i\langle\delta'_a+d^{*}d\delta_a,\hat A\rangle)\cdot\Psi_a$ is a basis element of $\mathcal{P}$,
since $d^{*}(\delta'_a+d^{*}d\delta_a)=d^{*}\delta'_a=q\cdot\mathfrak{q}(\Psi_a)$.
As the coefficient functions are scalar and even elements of the Grassmann algebra commute,
\[
\begin{aligned}
P_a={}&\lambda_a\cdot\exp\left(\frac{1}{2}\|d\delta_a\|^2+\langle\delta'_a,\delta_a\rangle\right)\\
&\cdot\mathfrak{V}_{\delta_a}\cdot\left(\exp(i\langle\delta'_a+d^{*}d\delta_a,\hat A\rangle)\cdot\Psi_a\right).
\end{aligned}
\]
The matrix $g_\mu$ is of the form considered, and the formulas above, applied with $g=g_\mu$, yield
\[
\begin{aligned}
P_a^{g_\mu}={}&\lambda_a\cdot\exp\left(\frac{1}{2}\|d(g_{\mu*}\delta_a)\|^2+\langle g_{\mu*}\delta'_a,g_{\mu*}\delta_a\rangle\right)\\
&\cdot\exp\left(i\langle g_{\mu*}\delta'_a+d^{*}d(g_{\mu*}\delta_a),\hat A\rangle\right)\cdot\mathfrak{V}_{g_{\mu*}\delta_a}\cdot\Psi_a^{g_\mu},
\end{aligned}
\]
which is \eqref{eq:elimination-new-terms} with $\delta'_a$, $\delta_a$ and $\Psi_a$ replaced by $g_{\mu*}\delta'_a$, $g_{\mu*}\delta_a$ and $\Psi_a^{g_\mu}$.
\end{proof}

\subsection{The directions $e_2$ and $e_3$}\label{sec:other-directions}

The thin layers $\mathbb{X}_{j,k}$, the maps $\mathfrak{R}^{\pm}_{\mathbb{L}}$ and the lowest and highest planes of the elements of $\mathcal{P}$ belong to the direction $e_1$.
The corresponding objects of the directions $e_2$ and $e_3$ are their images under the cyclic permutation matrices of Definition~\ref{def:other-directions}.
Recall that $g_1=1$ and that, for $\mu=2,3$, $g_\mu\in\Ocub$ is the cyclic permutation matrix with $g_\mu e_1=e_\mu$.
Each of $g_2$ and $g_3$ is the product of the reflections in two diagonal mirrors through the origin:
\[
g_2=\theta_{(x_1-x_2,0)}\,\theta_{(x_2-x_3,0)},
\qquad
g_3=g_2^{-1}=\theta_{(x_2-x_3,0)}\,\theta_{(x_1-x_2,0)} .
\]
Indeed, by Definition~\ref{def:mirrors} the reflection $\theta_{(x_1-x_2,0)}$ exchanges the first two coordinates of a point of $\R^3$,
and $\theta_{(x_2-x_3,0)}$ exchanges the last two.
In the first product the right factor fixes $e_1$ and the left factor maps it to $e_2$;
the vector $e_2$ is mapped to $e_3$ by the right factor and then fixed;
and $e_3$ is mapped to $e_2$ and then to $e_1$.
Thus the first product maps $e_1\mapsto e_2\mapsto e_3\mapsto e_1$,
and its inverse, which is the second product, maps $e_1\mapsto e_3\mapsto e_2\mapsto e_1$.
In particular $g_2$ and $g_3$ are even elements of $W$.

Fix $\mu\in\{1,2,3\}$.
As $g_\mu$ is orthogonal and $g_\mu^{-1}e_\mu=e_1$, every $x\in\R^3$ satisfies
\[
(g_\mu x)_\mu=\langle g_\mu x,e_\mu\rangle=\langle x,e_1\rangle=x_1 ;
\]
so $g_\mu$ maps each plane on which $x_1$ is constant onto the plane on which $x_\mu$ takes the same constant value.
In particular, for $c\in\tfrac12\Z$ the mirror $(x_\mu,c)$ of Definition~\ref{def:mirrors} is the image of the mirror $(x_1,c)$,
and the two reflections are conjugate,
\[
\theta_{(x_\mu,c)}=g_\mu\,\theta_{(x_1,c)}\,g_\mu^{-1},
\]
both sides being the map $x\mapsto x-2(x_\mu-c)\,e_\mu$.
The map $\Theta_{\mu,c}=\Theta_{(x_\mu,c)}$ of Definition~\ref{def:other-directions} is the reflection of Definition~\ref{def:grassmann-reflections} at the mirror $(x_\mu,c)$.
For every even integrand $G$ the rule \eqref{eq:setting-group-action} for the action of $W$ gives
\begin{equation}\label{eq:elimination-conjugation}
\Theta_{\mu,c}\left(G^{g_\mu}\right)
=G^{\theta_{(x_\mu,c)}\,g_\mu}
=G^{g_\mu\,\theta_{(x_1,c)}}
=\left(\Theta_{(x_1,c)}G\right)^{g_\mu} .
\end{equation}

\begin{definition}\label{def:transported-objects}
Let $\mu\in\{1,2,3\}$.
The \emph{thin layers of the direction $e_\mu$} are the sets
\[
\mathbb{X}^{\mu}_{j,k}:=g_\mu\mathbb{X}_{j,k},\qquad j\in\Z,\ k=1,\dots,6 .
\]
They are ordered by the coordinate $x_\mu$, as the thin layers $\mathbb{X}_{j,k}$ are ordered by $x_1$.
For a thin layer $\mathbb{L}$ of the direction $e_1$, the successor and the predecessor of $g_\mu\mathbb{L}$ are
$(g_\mu\mathbb{L})^{+}:=g_\mu(\mathbb{L}^{+})$ and $(g_\mu\mathbb{L})^{-}:=g_\mu(\mathbb{L}^{-})$;
likewise $(g_\mu\mathbb{L})^{++}:=g_\mu(\mathbb{L}^{++})$ and $(g_\mu\mathbb{L})^{--}:=g_\mu(\mathbb{L}^{--})$.
The \emph{thick layer $j$ of the direction $e_\mu$} is the open slab
\[
\left\{|x_\mu-j|<\tfrac12\right\}=g_\mu\left\{|x_1-j|<\tfrac12\right\};
\]
it contains the links of the plane $\{x_\mu=j\}$ and the sites of $\mathbb{X}^{\mu}_{j,1},\dots,\mathbb{X}^{\mu}_{j,6}$.

Let $\rho^2\neq1$.
For a thin layer $\mathbb{L}$ of the direction $e_1$ and for $P\in\mathcal{P}$ put
\[
\mathfrak{R}^{\mu,+}_{g_\mu\mathbb{L}}(P):=\left(\mathfrak{R}^{+}_{\mathbb{L}}\left(P^{g_\mu^{-1}}\right)\right)^{g_\mu},
\qquad
\mathfrak{R}^{\mu,-}_{g_\mu\mathbb{L}}(P):=\left(\mathfrak{R}^{-}_{\mathbb{L}}\left(P^{g_\mu^{-1}}\right)\right)^{g_\mu}.
\]

The \emph{lowest} (\emph{highest}) \emph{plane of $P\in\mathcal{P}$ in the direction $e_\mu$} is the lowest (highest) plane of $P^{g_\mu^{-1}}$ in the sense of Definition~\ref{def:two-layer}.
\end{definition}

For $\mu=1$ these are the thin and thick layers, the maps $\mathfrak{R}^{\pm}_{\mathbb{L}}$ and the lowest and highest planes of the direction $e_1$.
The space of the two-layer elements at the link plane $(x_\mu,c)$, $c\in\Z+\tfrac12$, is $\mathcal{P}_1(\mu,c)=\left\{Y^{g_\mu}:Y\in\mathcal{P}_1(c)\right\}$ (Definition~\ref{def:other-directions}).
By \eqref{eq:setting-group-action} and Proposition~\ref{prop:class-structure}\eqref{it:class-stable},
the maps $G\mapsto G^{g_\mu}$ and $G\mapsto G^{g_\mu^{-1}}$ are linear maps of $\mathcal{P}$ into itself and are inverse to each other;
as $\mathfrak{R}^{+}_{\mathbb{L}}$ and $\mathfrak{R}^{-}_{\mathbb{L}}$ are linear maps of $\mathcal{P}$ into itself (Proposition~\ref{prop:schwinger-dyson}\eqref{it:replacement-maps}),
so are $\mathfrak{R}^{\mu,+}_{g_\mu\mathbb{L}}$ and $\mathfrak{R}^{\mu,-}_{g_\mu\mathbb{L}}$.

The thin layers of the direction $e_\mu$ can be read off from the coordinates of index $\mu$.
Let $x\in\mathbb{X}$.
The projection $\pi$ is equivariant under $W$ (Proposition~\ref{prop:geometry}\eqref{it:geometry-projection}), and $g_\mu$ is linear;
hence $\pi(g_\mu x)=g_\mu\pi(x)$ and $y(g_\mu x)=g_\mu x-g_\mu\pi(x)=g_\mu y(x)$,
so that the coordinates of index $\mu$ of $\pi(g_\mu x)$ and of $y(g_\mu x)$ are the first coordinates of $\pi(x)$ and of $y(x)$.
Since $g_\mu$ maps $\mathbb{X}=Wv_0$ onto itself,
$\mathbb{X}^{\mu}_{j,1},\dots,\mathbb{X}^{\mu}_{j,6}$ are the sets of the sites $x\in\mathbb{X}$ with $\pi(x)_\mu=j$ and
\[
y_\mu(x)=-v_1,\quad-v_2,\quad-v_3,\quad v_3,\quad v_2,\quad v_1,
\]
respectively.

\paragraph{Reading sets.}
Let $g\in W$ and $P\in\mathcal{P}$.
The links read by $P^g$ (Definition~\ref{def:two-layer}) are the images under $g$ of the links read by $P$,
and the sites read by $P^g$ are the images under $g$ of the sites read by $P$.
To prove this, let first $g=\theta_\Pi$ for a mirror $\Pi$, so that $P^g=\Theta_\Pi P$,
and let $P=\sum_a\lambda_a\cdot\exp(i\langle\delta_a,\hat A\rangle)\cdot\Psi_a$ be the representation of $P$ in the basis of $\mathcal{P}$.
By the formulas of Proposition~\ref{prop:class-structure}\eqref{it:class-stable},
$\Theta_\Pi$ maps $\exp(i\langle\delta_a,\hat A\rangle)\cdot\Psi_a$ to $\exp(-i\langle\theta_\Pi^{*}\delta_a,\hat A\rangle)$ times a product of linear combinations of generators at the sites $\theta_\Pi x$,
where $x$ ranges over the sites of the fields of $\Psi_a$;
and $\theta_\Pi^{*}\delta_a$ is supported on $\theta_\Pi(\supp\delta_a)$, since $\theta_\Pi^{*}$ is the push-forward by $\theta_\Pi$ (Definition~\ref{def:factors}).
Expanding these products into monomials and collecting equal basis elements gives the representation of $\Theta_\Pi P$.
Consequently every link and every site that $\Theta_\Pi P$ reads is the image under $\theta_\Pi$ of a link or a site that $P$ reads.
If $g=\theta_{\Pi_1}\cdots\theta_{\Pi_r}$ is a word in reflections, then $P^g=\Theta_{\Pi_1}\cdots\Theta_{\Pi_r}P$ by \eqref{eq:setting-group-action},
and induction on $r$ shows that every link and every site that $P^g$ reads is the image under $g$ of a link or a site that $P$ reads.
The same conclusion for the element $P^g$ and for $g^{-1}$, together with $\left(P^g\right)^{g^{-1}}=P$ \eqref{eq:setting-group-action}, gives the converse inclusion.
For the even elements $g_\mu$ the formulas of Proposition~\ref{prop:class-structure}\eqref{it:class-stable} take the form
\[
\left(\exp(i\langle\delta,\hat A\rangle)\cdot\Psi\right)^{g_\mu}=\exp(i\langle g_{\mu*}\delta,\hat A\rangle)\cdot\Psi^{g_\mu},
\]
with $g_{\mu*}=(g_\mu^{-1})^{*}$;
the factors of $\Psi^{g_\mu}$ are those of $\Psi$, in the same order,
each generator $\psi_\alpha(x)$ being replaced by a linear combination of $\psi_1(g_\mu x),\dots,\psi_4(g_\mu x)$
and each $\bar\psi_\alpha(x)$ by a linear combination of $\bar\psi_1(g_\mu x),\dots,\bar\psi_4(g_\mu x)$.

As a result, substituting $x_\mu$ for $x_1$ in the descriptions of Definition~\ref{def:two-layer} yields the objects of Definition~\ref{def:transported-objects}.
For instance, the lowest (highest) plane of $P\in\mathcal{P}$ in the direction $e_\mu$ is the least (largest) integer $j$ such that $P$ reads a site of the thick layer $j$ of the direction $e_\mu$ or a link that meets the plane $\{x_\mu=j\}$:
a site $g_\mu^{-1}x$ lies in $\left\{|x_1-j|<\tfrac12\right\}$ if and only if $x$ lies in $\left\{|x_\mu-j|<\tfrac12\right\}$,
and a link $g_\mu^{-1}b$ meets the plane $\{x_1=j\}$ if and only if $b$ meets the plane $\{x_\mu=j\}$.
Likewise, for $c=j_0+\tfrac12$ with $j_0\in\Z$, the space $\mathcal{P}_1(\mu,c)$ is the linear span of the elements $\exp(i\langle\delta,\hat A\rangle)\cdot\Psi$ of $\mathcal{P}$
in which $\delta$ is supported on the links of the planes $\{x_\mu=j_0\}$, $\{x_\mu=j_0+1\}$ and on the links between these two planes,
and $\Psi$ is a monomial in the fields of $\mathbb{X}^{\mu}_{j_0,6}\cup\mathbb{X}^{\mu}_{j_0+1,1}$.
Indeed, $g_\mu$ maps the links of the planes $\{x_1=j_0\}$, $\{x_1=j_0+1\}$ and the links between them onto the links just named,
and $\mathbb{X}_{j_0,6}\cup\mathbb{X}_{j_0+1,1}$ onto $\mathbb{X}^{\mu}_{j_0,6}\cup\mathbb{X}^{\mu}_{j_0+1,1}$.
If $Y$ is a two-layer monomial at $c$ (Definition~\ref{def:two-layer}), then $Y^{g_\mu}$ reads only links and sites of the kinds just named,
so that every term of its representation in the basis of $\mathcal{P}$ belongs to the span in question;
by linearity of $G\mapsto G^{g_\mu}$, the whole of $\mathcal{P}_1(\mu,c)$ lies in this span.
Conversely, let $P=\exp(i\langle\delta,\hat A\rangle)\cdot\Psi$ be an element of $\mathcal{P}$ of the kind described.
Then $P^{g_\mu^{-1}}$ reads only links of the planes $\{x_1=j_0\}$, $\{x_1=j_0+1\}$ and links between them, and only sites of $\mathbb{X}_{j_0,6}\cup\mathbb{X}_{j_0+1,1}$;
every term of its representation is therefore a two-layer monomial at $c$, so that $P^{g_\mu^{-1}}\in\mathcal{P}_1(c)$ and $P=\left(P^{g_\mu^{-1}}\right)^{g_\mu}\in\mathcal{P}_1(\mu,c)$.

These descriptions involve the coordinate $x_\mu$ alone.
The arguments above apply without change to any $g'\in\Ocub$ with $g'e_1=e_\mu$ in the place of $g_\mu$:
$g'$ is orthogonal, $\pi(g'x)=g'\pi(x)$ for $x\in\mathbb{X}$, and the links and sites read by $P^{g'}$ are the images under $g'$ of those read by $P$;
if $g'$ is odd, then $G\mapsto G^{g'}$ is antilinear \eqref{eq:setting-group-action}, and the set of the elements $Y^{g'}$, $Y\in\mathcal{P}_1(c)$, is still a complex linear space.
Hence the planes, the thin and the thick layers, the spaces $\mathcal{P}_1(\mu,c)$ and the lowest and highest planes of the direction $e_\mu$ remain the same when $g_\mu$ is replaced by such a matrix $g'$.
The maps $\mathfrak{R}^{\mu,\pm}_{g_\mu\mathbb{L}}$, and with them the element $Y^{\mu}_P$ of Theorem~\ref{thm:two-layer-reduction}, are formed with the matrix $g_\mu$ itself,
in the same way as the maps $\mathfrak{R}^{\pm}_{\mathbb{L}}$ and the element $Y_P$ of that theorem are formed with the total order of the generators fixed in Definition~\ref{def:covariant-classes}.

\begin{remark}\label{rem:universal-statements}
In the language of Definition~\ref{def:universal-identity}, the identities of Propositions~\ref{prop:schwinger-dyson}, \ref{prop:normal-form}, \ref{prop:contour-shift} and~\ref{prop:outer-plane}
say that certain triples are universal identities, or images of universal identities under one mirror.
In what follows $M\in\R$, $w_0,\dots,w_3>0$ and $q>0$, and $\Pi$ is a mirror.
\begin{enumerate}
\item[(i)]
Let $\rho^2\neq1$, let $\mathbb{L}$ be a thin layer, and let $P\in\mathcal{P}$ be read in a box $B$.
By Proposition~\ref{prop:schwinger-dyson}\eqref{it:replacement-identity}, each of the triples
\[
\left(P;\left(0,\mathfrak{R}^{+}_{\mathbb{L}}P\right);B^{+1}\right),
\qquad
\left(P;\left(0,\mathfrak{R}^{-}_{\mathbb{L}}P\right);B^{+1}\right),
\]
with an index set of one element, is a universal identity of the first kind.
Their images under the mirror $\Pi$ are the triples
$\left(\Theta_\Pi P;\left(0,\Theta_\Pi\mathfrak{R}^{\pm}_{\mathbb{L}}P\right);\theta_\Pi\left(B^{+1}\right)\right)$,
universal identities of the first kind by Lemma~\ref{lem:reflection-identities}\eqref{it:reflected-universal},
and Proposition~\ref{prop:schwinger-dyson}\eqref{it:mirrored-rule} is their identity \eqref{eq:elimination-universal}.
\item[(ii)]
Let $\rho^2\neq1$, $c\in\Z+\tfrac12$ and $P\in\mathcal{P}$, and let $I$, $\lambda_a$, $\beta_a$, $Y_a$ and $B'=B'(P,c)$ be as in Proposition~\ref{prop:normal-form}.
By Proposition~\ref{prop:normal-form}\eqref{it:normal-form-identity},
\[
\left(P;\left(0,\lambda_a\cdot\varphi_{\beta_a}\cdot Y_a\right)_{a\in I};B'\right)
\]
is a universal identity of the first kind.
Since $\Theta_\Pi$ is antilinear, its image under $\Pi$ is
$\left(\Theta_\Pi P;\left(0,\overline{\lambda_a}\cdot\Theta_\Pi\left(\varphi_{\beta_a}\cdot Y_a\right)\right)_{a\in I};\theta_\Pi B'\right)$,
and the identity \eqref{eq:elimination-universal} of this image is Proposition~\ref{prop:normal-form}\eqref{it:normal-form-mirror}.
\item[(iii)]
Let $\rho\in\R$, let $S$ be a finite set of links contained in one coordinate plane, let $\delta\in\R^S$,
and let $P\in\mathcal{P}$ be read in a box $B$ that contains $b+[-1,1]^3$ for every $b\in S$.
For every torus $\mathbb{T}$ in which $B$ embeds, $\mathbb{T}$, $S$, $B$ and $P$ satisfy \eqref{eq:elimination-shift-box};
moreover $d\delta$ is supported on plaquettes that contain a link of $S$, and these plaquettes lie in $B$.
So Proposition~\ref{prop:contour-shift}\eqref{it:shift-identity} states that
\[
\left(P;\left(d\delta,\ \exp\left(\frac12\|d\delta\|^2\right)\cdot\varphi_{d\delta}\cdot\mathfrak{V}_\delta\cdot P^{(\delta)}\right);B\right)
\]
is a universal identity.
\item[(iv)]
Let $\rho^2\neq1$ and $j\in\Z$, let $P\in\mathcal{P}$ be read in a box $B$,
and let $I$ and $(\delta_a,P_a)_{a\in I}$ be as in Proposition~\ref{prop:outer-plane}\eqref{it:lowest-plane} or as in Proposition~\ref{prop:outer-plane}\eqref{it:highest-plane}.
Then
\[
\left(P;\left(d\delta_a,P_a\right)_{a\in I};B^{+8}\right)
\]
is a universal identity.
For the condition on the supports, note that a plaquette of $\supp d\delta_a$ contains a link $b+e_1$, respectively $b-e_1$,
where $b$ is a link of the plane $\{x_1=j\}$ read by $P^{\uparrow}$, respectively by $P^{\downarrow}$;
this element is read in $B^{+6}$ by six applications of Proposition~\ref{prop:schwinger-dyson}\eqref{it:replacement-identity},
so that $b\subset B^{+6}$ and the plaquette lies in $b+[-2,2]^3\subset B^{+8}$.
The image of this universal identity under $\Pi$ is
$\left(\Theta_\Pi P;\left(-\theta_\Pi^{*}d\delta_a,\Theta_\Pi P_a\right)_{a\in I};\theta_\Pi\left(B^{+8}\right)\right)$,
and its identity \eqref{eq:elimination-universal} is Proposition~\ref{prop:outer-plane}\eqref{it:reflected-elimination}.
\end{enumerate}
\end{remark}

\begin{remark}\label{rem:transport}
Let $\mu\in\{2,3\}$, $M\in\R$, $w_0,\dots,w_3>0$ and $q>0$.
For the direction $e_\mu$ the statements of this section take the following form.
\begin{enumerate}
\item[(i)]
Let $A$ be a real or complex one-cochain on $\Z^3$.
Lemma~\ref{lem:layer-bonds} holds for the thin layers $\mathbb{X}^{\mu}_{j,k}$, ordered in the direction $e_\mu$, and for the phases $U_A$,
in the form obtained by applying it to the cochain $g_\mu^{*}A$ and carrying the result over by $g_\mu$.
Three facts underlie this.
First, if a site is written $x=g'v_0$ with $g'\in W$, then $g_\mu x^{(i)}=g_\mu g's_iv_0=(g_\mu x)^{(i)}$ for $i=0,\dots,3$ (Definition~\ref{def:fermion-sites});
so the bijection $g_\mu$ of $\mathbb{X}$ carries bonds to bonds and preserves their types.
Second, $\pi(g_\mu x)=g_\mu\pi(x)$, because $\pi$ is equivariant under $W$ (Proposition~\ref{prop:geometry}\eqref{it:geometry-projection}).
Third, for every bond $\{x,x'\}$ the second fact, the definition of the phases (Definition~\ref{def:chamber-operator}) and Definition~\ref{def:cochains} yield
\[
U_A(g_\mu x,g_\mu x')=\exp\left(iqA\left(g_\mu\pi(x),g_\mu\pi(x')\right)\right)=U_{g_\mu^{*}A}(x,x').
\]
In the lemma so obtained the planes $\{x_1=c\}$ are replaced by the planes $\{x_\mu=c\}$, the vectors $e_\nu$ by $g_\mu e_\nu$,
and the slice variables $a_j$ and $A_j$ of Definition~\ref{def:slices} by those of $g_\mu^{*}A$.
These are the values of $A$ on the images under $g_\mu$ of the links concerned, without any change of sign:
the matrix $g_\mu$ permutes $e_1$, $e_2$, $e_3$ and has no negative entry,
so it maps each link $(v,\nu)$, together with its orientation, onto the link $(g_\mu v,\nu')$ with $g_\mu e_\nu=e_{\nu'}$,
and Definition~\ref{def:cochains} gives $\left(g_\mu^{*}A\right)_\nu(v)=A_{\nu'}(g_\mu v)$.
\item[(ii)]
No direction is singled out in items \eqref{it:equation-of-motion} and \eqref{it:replacement-rules} of Proposition~\ref{prop:schwinger-dyson},
nor in Propositions~\ref{prop:a-priori}, \ref{prop:contour-shift} and~\ref{prop:class-structure}.
The equations of motion are formulated for an arbitrary site and the replacement rules for an arbitrary bond;
Proposition~\ref{prop:class-structure} concerns every mirror and every element of $W$;
the bounds of Proposition~\ref{prop:a-priori} apply to an arbitrary element of $\mathcal{P}$;
and in Proposition~\ref{prop:contour-shift} the links of $S$ may lie in a plane orthogonal to any one of $e_1$, $e_2$, $e_3$
(for the planes $\{x_\mu=j\}$ see \eqref{eq:elimination-other-planes}).
\item[(iii)]
Let $\rho^2\neq1$.
Remark~\ref{rem:universal-statements} exhibits the identities of Proposition~\ref{prop:schwinger-dyson}\eqref{it:replacement-identity},
of Proposition~\ref{prop:normal-form}\eqref{it:normal-form-identity}
and of Proposition~\ref{prop:outer-plane}\eqref{it:lowest-plane}, \eqref{it:highest-plane} as universal identities.
Let $\mathcal{U}$ be one of them and $\Pi$ a mirror.
Then $\mathcal{U}^{g_\mu}$ and $\mathcal{U}^{g_\mu\theta_\Pi}$ are universal identities as well (Lemma~\ref{lem:reflection-identities}\eqref{it:reflected-universal}), and
\[
g_\mu\,\theta_\Pi=\theta_{g_\mu\Pi}\,g_\mu ,
\]
where $g_\mu\Pi$ denotes the image of the mirror $\Pi=(\ell,c)$, that is, the plane $\left\{\ell\circ g_\mu^{-1}=c\right\}$.
This plane is again a mirror.
Indeed, since $g_\mu$ permutes the coordinates, $\ell\circ g_\mu^{-1}=\pm\ell'$ for one of the nine linear forms $\ell'$ of Definition~\ref{def:mirrors}, of the same kind as $\ell$;
the admissible positions of a form depend only on its kind and are symmetric under $c\mapsto-c$,
so that $\{\ell'=\pm c\}$ is a mirror, and its reflection is $\theta_{g_\mu\Pi}=g_\mu\,\theta_\Pi\,g_\mu^{-1}$.

Every element of $\mathcal{P}$ has the form $P^{g_\mu}$ with $P\in\mathcal{P}$.
By Definition~\ref{def:transported-objects},
$\mathfrak{R}^{\mu,\pm}_{g_\mu\mathbb{L}}\left(P^{g_\mu}\right)=\left(\mathfrak{R}^{\pm}_{\mathbb{L}}P\right)^{g_\mu}$ for every thin layer $\mathbb{L}$ of the direction $e_1$,
and the lowest and the highest plane of $P^{g_\mu}$ in the direction $e_\mu$ are those of $P$;
moreover $P^{g_\mu\theta_\Pi}=\Theta_{g_\mu\Pi}\left(P^{g_\mu}\right)$ by \eqref{eq:setting-group-action}.
Accordingly, $\mathcal{U}^{g_\mu}$ is the identity of the maps $\mathfrak{R}^{\mu,\pm}$, of the normal form at the plane $(x_\mu,c)$, $c\in\Z+\tfrac12$,
or of the elimination of the lowest or of the highest plane in the direction $e_\mu$, according to the case;
and $\mathcal{U}^{g_\mu\theta_\Pi}$ is the reflected form of this identity at the mirror $g_\mu\Pi$, which ranges over all mirrors as $\Pi$ does.
In this sense the identities of Proposition~\ref{prop:schwinger-dyson}\eqref{it:replacement-maps}, \eqref{it:mirrored-rule}, of Proposition~\ref{prop:normal-form}
and of Proposition~\ref{prop:outer-plane}\eqref{it:lowest-plane}--\eqref{it:reflected-elimination} hold in the direction $e_\mu$.
According to the formulas for even elements in Proposition~\ref{prop:class-structure}\eqref{it:class-stable},
the elements of $\mathcal{U}^{g_\mu}$ are computed from the formulas that define the elements of $\mathcal{U}$
by replacing every cochain by its image under $g_{\mu*}$ and every monomial $\Psi$ by $\Psi^{g_\mu}$; the complex coefficients do not change.
In the elimination of an outer plane, for instance, \eqref{eq:elimination-new-terms} holds with $g_{\mu*}\delta'_a$, $g_{\mu*}\delta_a$ and $\Psi_a^{g_\mu}$ in the places of $\delta'_a$, $\delta_a$ and $\Psi_a$;
the cochain $g_{\mu*}\delta_a$ arises from the restriction of $g_{\mu*}\delta'_a$ to the links of the plane $\{x_\mu=j\}$
by the translation by $e_\mu$ in the case of Proposition~\ref{prop:outer-plane}\eqref{it:lowest-plane} and by $-e_\mu$ in the case of Proposition~\ref{prop:outer-plane}\eqref{it:highest-plane}.
The operators $d$ and $d^{*}$, the hopping matrices $w(x,x')$, the phases $U_A$ and the polynomials $\mathfrak{V}$ enter unchanged;
whatever signs the orientations produce are absorbed in the push-forward $g_{\mu*}$.

With the thin layers and the order of the direction $e_\mu$, the maps $\mathfrak{R}^{\mu,\pm}$ also have the properties
\eqref{it:replacement-fields}, \eqref{it:replacement-links} and \eqref{it:replacement-characters} of Proposition~\ref{prop:schwinger-dyson}\eqref{it:replacement-maps}.
Let $\mathbb{L}$ be a thin layer of the direction $e_1$ and $P\in\mathcal{P}$.
As far as these properties speak of sites and links, they follow from the fact that the links and sites read by $P^{g_\mu}$ are the images under $g_\mu$ of those read by $P$,
together with $\pi(g_\mu x)=g_\mu\pi(x)$ and the fact that $g_\mu$ maps bonds of type~$3$ onto bonds of type~$3$;
property \eqref{it:replacement-characters} follows from the multiplicativity of $G\mapsto G^{g_\mu}$ and from $\left(\varphi_\beta\right)^{g_\mu}=\varphi_{g_{\mu*}\beta}$ (Proposition~\ref{prop:class-structure}\eqref{it:class-stable}).
The second assertion of \eqref{it:replacement-fields} speaks of single generators.
Let $x$ be a site outside $\mathbb{L}$, $\mathbb{L}^{+}$ and $\mathbb{L}^{++}$, and write $P=\sum_\Psi\Psi\cdot G_\Psi$,
where $\Psi$ ranges over the ordered monomials in the generators at $x$ and no $G_\Psi$ contains a generator at $x$;
in such an expansion of an element of $\mathcal{P}$ at a site, the fields at this site are the generators of the monomials with a non-zero coefficient.
With the replacement operators given by the formulas of Proposition~\ref{prop:schwinger-dyson}\eqref{it:replacement-rules},
the recursion \eqref{eq:elimination-recursion} makes sense for the product of every coefficient function $\exp(i\langle\delta,\hat A\rangle)$ with every monomial, gauge invariant or not,
and so extends $\mathfrak{R}^{+}_{\mathbb{L}}$ linearly to the span of these products.
It creates and removes no generator at $x$,
and the signs that arise when a monomial $\Psi$ in the generators at $x$ is moved across the replaced generator and across the generators that the replacement operator creates or differentiates cancel each other;
therefore $\mathfrak{R}^{+}_{\mathbb{L}}(P)=\sum_\Psi\Psi\cdot\mathfrak{R}^{+}_{\mathbb{L}}(G_\Psi)$.
Likewise, the substitution of generators that yields $\Psi^{g_\mu}$, together with $\delta\mapsto g_{\mu*}\delta$ in the coefficient functions,
extends $G\mapsto G^{g_\mu}$ to a linear and multiplicative bijection of the same span;
it sends each $\Psi$ to a linear combination of the ordered monomials in the generators at $g_\mu x$.
Collecting terms, one obtains for every such monomial $\Psi'$ a linear combination $G'$ of the elements $G_\Psi$
such that the coefficient of $\Psi'$ is $\left(G'\right)^{g_\mu}$ in $P^{g_\mu}$
and $\left(\mathfrak{R}^{+}_{\mathbb{L}}G'\right)^{g_\mu}$ in $\left(\mathfrak{R}^{+}_{\mathbb{L}}P\right)^{g_\mu}=\mathfrak{R}^{\mu,+}_{g_\mu\mathbb{L}}\left(P^{g_\mu}\right)$.
As the extended map is injective, the second coefficient vanishes whenever the first one does;
hence every field of $\mathfrak{R}^{\mu,+}_{g_\mu\mathbb{L}}\left(P^{g_\mu}\right)$ at the site $g_\mu x$ is a field of $P^{g_\mu}$.
The downward maps are treated alike, with $\mathbb{L}^{-}$ and $\mathbb{L}^{--}$ in the places of $\mathbb{L}^{+}$ and $\mathbb{L}^{++}$.
\item[(iv)]
The estimates, Proposition~\ref{prop:outer-plane}\eqref{it:plane-estimate} and Theorem~\ref{thm:two-layer-reduction}, cannot be transported in this way.
They concern the functional $\mu^t_i$, which is formed on the torus $D(n_i)$ with the array $O_{\mathfrak{d}}$ (Definition~\ref{def:data}).
The torus is not invariant under $g_\mu$:
the vector $\tfrac{n_i}{2}(e_1-e_2)$ is a period of $D(n_i)$, its image under $g_\mu$ has the third coordinate $\pm\tfrac{n_i}{2}$,
and the third coordinate of every period of $D(n_i)$ lies in $n_i\Z$.
In general the array is not invariant either:
its cells belong to the diagonal frame $(\ell_-,\ell_+,x_3)$ (Definition~\ref{def:cell-structures}), and $g_\mu$ maps this frame onto a different one.
Both estimates hold on $D(n_i)$ itself, in a form that covers the transported identities:
Proposition~\ref{prop:outer-plane}\eqref{it:plane-estimate} applies to every universal identity,
and the bound \eqref{eq:elimination-general} of Theorem~\ref{thm:two-layer-reduction} holds for the elements $P^g$ with an arbitrary $g\in W$.
\end{enumerate}
\end{remark}

\subsection{Two layers at a link plane}\label{sec:link-plane-reduction}

Elimination steps from below and from above, followed by the maps $\mathfrak{R}^{\pm}_{\mathbb{L}}$ in the two thick layers that remain,
reduce every element of $\mathcal{P}$ to a two-layer element at a link plane of the family $x_1$;
under the perturbed functionals the error is of order $t$.

\paragraph{The element $Y_P$.}
Let $\rho\in\R$ with $\rho^2\neq1$, let $M\in\R$, $w_0,\dots,w_3>0$ and $q>0$;
then the maps $\mathfrak{R}^{\pm}_{\mathbb{L}}$ of Proposition~\ref{prop:schwinger-dyson}\eqref{it:replacement-maps}
and the elimination steps of Proposition~\ref{prop:outer-plane} are defined.
Let $j_0\in\Z$ and $c:=j_0+\tfrac12$.
Let $P\in\mathcal{P}$, fix a box $B$ in which $P$ is read, and put
\[
B_s:=B^{+8s}\qquad(s=0,1,2,\dots)
\]
(Definition~\ref{def:enlarged-boxes}).
In two phases, elements $P_0:=P,P_1,\dots,P_r$ of $\mathcal{P}$ are defined, $P_s$ being read in $B_s$;
a third phase produces $Y_P$ from $P_r$.

\emph{First phase: elimination from below.}
Let $s\geq0$, suppose that $P_s\in\mathcal{P}$ has been defined and is read in $B_s$,
and let $j$ be the lowest plane of $P_s$ (Definition~\ref{def:two-layer}),
with $j=+\infty$ if $P_s$ is constant, according to the convention of Definition~\ref{def:covariant-classes}.
If $j\geq j_0$, the first phase ends, and $r_1:=s$.
Otherwise $j<j_0$, and Proposition~\ref{prop:outer-plane}\eqref{it:lowest-plane}, applied to $P_s$, the box $B_s$ and the plane $\{x_1=j\}$,
provides a finite index set $I_s$ and, for $a\in I_s$,
a cochain $\delta_{s,a}\in C^1_{\mathrm{c}}$ supported on the links of the plane $\{x_1=j+1\}$
and an element $P_{s,a}\in\mathcal{P}$ that is read in $B_s^{+8}=B_{s+1}$ and has lowest plane at least $j+1$.
Put
\[
P_{s+1}:=\sum_{a\in I_s}P_{s,a}.
\]
In the representation of Definition~\ref{def:covariant-classes}, the coefficient of a basis element in a sum is the sum of its coefficients in the summands;
so a sum of elements of $\mathcal{P}$ reads only links and sites that one of the summands reads.
Hence $P_{s+1}$ is read in $B_{s+1}$, and its lowest plane is at least $j+1$.
Each step thus raises the lowest plane by one or more;
as the first phase goes on only while this plane is below $j_0$, it ends after finitely many steps,
and the lowest plane of $P_{r_1}$ is at least $j_0$.

\emph{Second phase: elimination from above.}
Now let $s\geq r_1$, assume that $P_s$ is read in $B_s$ and has lowest plane at least $j_0$,
and let $j$ be the highest plane of $P_s$, with $j=-\infty$ if $P_s$ is constant (Definition~\ref{def:covariant-classes}).
If $j\leq j_0+1$, the second phase ends, and $r:=s$.
Otherwise $j\geq j_0+2$, and Proposition~\ref{prop:outer-plane}\eqref{it:highest-plane}, applied to $P_s$, the box $B_s$ and the plane $\{x_1=j\}$,
provides a finite index set $I_s$ and, for $a\in I_s$,
a cochain $\delta_{s,a}\in C^1_{\mathrm{c}}$ supported on the links of the plane $\{x_1=j-1\}$
and an element $P_{s,a}\in\mathcal{P}$ that is read in $B_{s+1}$ and has highest plane at most $j-1$;
its lowest plane is at least the smaller of the two numbers $j-2$ and the lowest plane of $P_s$,
hence at least $j_0$, because $j-2\geq j_0$.
Put again
\[
P_{s+1}:=\sum_{a\in I_s}P_{s,a};
\]
as before, $P_{s+1}$ is read in $B_{s+1}$, its highest plane is at most $j-1$, and its lowest plane is at least $j_0$,
so that $P_{s+1}$ satisfies the assumptions made on $P_s$.
The highest plane of $P_{s+1}$ is thus smaller than that of $P_s$ by one at least;
since the second phase goes on only while this plane exceeds $j_0+1$, it ends after finitely many steps, $r_2\geq0$ in number.
Then $r=r_1+r_2$, and $P_r$ has lowest plane at least $j_0$ and highest plane at most $j_0+1$.

\emph{Third phase: the two thick layers at $c$.}
The bounds $j_0$ and $j_0+1$ on the lowest and the highest plane of $P_r$ mean, by Definition~\ref{def:two-layer},
that each link read by $P_r$ lies in $\{x_1=j_0\}$ or in $\{x_1=j_0+1\}$ or joins these two planes,
and that each site read by $P_r$ belongs to the thick layer $j_0$ or to the thick layer $j_0+1$.
Put $P_{r,0}:=P_r$,
\[
P_{r,\ell}:=\mathfrak{R}^{+}_{\mathbb{X}_{j_0,\ell}}P_{r,\ell-1}\quad(\ell=1,\dots,5),
\qquad
P_{r,\ell}:=\mathfrak{R}^{-}_{\mathbb{X}_{j_0+1,12-\ell}}P_{r,\ell-1}\quad(\ell=6,\dots,10),
\]
and
\[
Y_P:=P_{r,10}
=\mathfrak{R}^{-}_{\mathbb{X}_{j_0+1,2}}\cdots\mathfrak{R}^{-}_{\mathbb{X}_{j_0+1,6}}\,
\mathfrak{R}^{+}_{\mathbb{X}_{j_0,5}}\cdots\mathfrak{R}^{+}_{\mathbb{X}_{j_0,1}}\,P_r .
\]
Each of the ten maps carries an element that is read in a box to an element that is read in this box enlarged by one
(Proposition~\ref{prop:schwinger-dyson}\eqref{it:replacement-identity});
so $P_{r,\ell}$ is read in $B_r^{+\ell}$, and $Y_P$ is read in $B^{+(8r+10)}$.

All maps and all elements of the three phases are defined on $\Z^3$:
no torus, no factor, no sequence of sides, no data, no index $i$ and no strength $t$ enters their definition.
Nor does the box $B$ enter:
in Proposition~\ref{prop:outer-plane}\eqref{it:lowest-plane} and~\eqref{it:highest-plane}
the index set, the cochains and the elements are determined by the given element and the plane alone,
and the maps $\mathfrak{R}^{\pm}_{\mathbb{L}}$ are maps of $\mathcal{P}$ into itself that involve no box
(Proposition~\ref{prop:schwinger-dyson}\eqref{it:replacement-maps}).
So the elements $P_s$, $P_{s,a}$, $P_{r,\ell}$ and $Y_P$ do not depend on the choice of $B$.

\begin{theorem}[Reduction to two layers]\label{thm:two-layer-reduction}
Let $0<\rho<1$, $M\in\R$, $w_0,\dots,w_3>0$ and $q>0$,
and let $\mathfrak{n}=(n_i)$ be a sequence of sides that satisfies \eqref{cond:divisibility} (Definition~\ref{def:sides}).
Let $c=j_0+\tfrac12$ with $j_0\in\Z$, let $P\in\mathcal{P}$, and let $Y_P$ be the element constructed above for $P$ and $c$.
Then $Y_P\in\mathcal{P}_1(c)$;
it is defined on $\Z^3$ and depends neither on $\mathfrak{n}$ nor on the data, the index $i$ or the strength $t$.
\begin{enumerate}
\item\label{it:reduction-link-plane}
For all data $\mathfrak{d}$ (Definition~\ref{def:data}) there are $C_{P,\mathfrak{d}}<\infty$ and $i_0$ such that
\[
\left|\mu^t_i(P)-\mu^t_i(Y_P)\right|\leq C_{P,\mathfrak{d}}\cdot t
\qquad\text{and}\qquad
\left|\mu^t_i(\Theta_c P)-\mu^t_i(\Theta_c Y_P)\right|\leq C_{P,\mathfrak{d}}\cdot t
\]
for $0\leq t\leq1$ and $i\geq i_0$, where $\Theta_c:=\Theta_{(x_1,c)}$.
\item\label{it:reduction-directions}
Let $\mu\in\{2,3\}$, let $g_\mu$, $\Theta_{\mu,c}$ and $\mathcal{P}_1(\mu,c)$ be as in Definition~\ref{def:other-directions}, and put
\[
\check P:=P^{g_\mu^{-1}},\qquad Y^{\mu}_P:=\left(Y_{\check P}\right)^{g_\mu},
\]
where $Y_{\check P}$ is the element constructed above for $\check P$ and $c$.
Then $Y^{\mu}_P\in\mathcal{P}_1(\mu,c)$;
it is defined on $\Z^3$ and depends neither on $\mathfrak{n}$ nor on the data, the index $i$ or the strength $t$.
For all data $\mathfrak{d}$ there are $C<\infty$ and $i_0$, depending on $P$, $\mu$, $c$ and $\mathfrak{d}$, such that
\[
\left|\mu^t_i(P)-\mu^t_i\left(Y^{\mu}_P\right)\right|\leq C\cdot t
\qquad\text{and}\qquad
\left|\mu^t_i\left(\Theta_{\mu,c} P\right)-\mu^t_i\left(\Theta_{\mu,c} Y^{\mu}_P\right)\right|\leq C\cdot t
\]
for $0\leq t\leq1$ and $i\geq i_0$.
\item\label{it:reduction-general}
For every $g\in W$ and all data $\mathfrak{d}$ there are $C_{P,g,\mathfrak{d}}<\infty$ and $i_0$ such that, for $0\leq t\leq1$ and $i\geq i_0$,
\begin{equation}\label{eq:elimination-general}
\left|\mu^t_i\left(P^g\right)-\mu^t_i\left((Y_P)^g\right)\right|\leq C_{P,g,\mathfrak{d}}\cdot t .
\end{equation}
\end{enumerate}
\end{theorem}

\begin{proof}
\emph{The element $Y_P$ lies in $\mathcal{P}_1(c)$.}
By the third phase of the construction, $P_r$ has fields only in the thick layers $j_0$ and $j_0+1$;
in particular it has no field below $\mathbb{X}_{j_0,1}$.
Proposition~\ref{prop:schwinger-dyson}\eqref{it:replacement-fields}, applied to the upward maps one after the other,
shows by induction on $\ell$ that $P_{r,\ell}$ has no field in or below $\mathbb{X}_{j_0,\ell}$ for $\ell=1,\dots,5$.
By the same item, every field of $P_{r,5}$ that is not a field of $P_r$ lies in $\mathbb{X}^{+}_{j_0,\ell}$ or in $\mathbb{X}^{++}_{j_0,\ell}$ for some $\ell\leq5$,
hence in $\mathbb{X}_{j_0+1,1}$ or below it.
So every field of $P_{r,5}$ lies in $\mathbb{X}_{j_0,6}$ or in the thick layer $j_0+1$,
and $P_{r,5}$ has no field above $\mathbb{X}_{j_0+1,6}$.
The same item for the downward maps shows by induction that $P_{r,\ell}$ has no field in or above $\mathbb{X}_{j_0+1,12-\ell}$ for $\ell=6,\dots,10$,
and that every field of $P_{r,10}$ that is not a field of $P_{r,5}$ lies in $\mathbb{X}^{-}_{j_0+1,k}$ or in $\mathbb{X}^{--}_{j_0+1,k}$ for some $k\in\{2,\dots,6\}$,
hence in $\mathbb{X}_{j_0,6}$ or above it.
Consequently every field of $Y_P=P_{r,10}$ lies in $\mathbb{X}_{j_0,6}\cup\mathbb{X}_{j_0+1,1}$,
the set of the ends of the bonds that cross the plane $\{x_1=c\}$.

Let $b$ be a link that $Y_P$ reads and $P_r$ does not read.
Then $b$ is added by one of the ten maps, and by Proposition~\ref{prop:schwinger-dyson}\eqref{it:replacement-links}
it joins $\pi(x)$ to $\pi(x')$ for a bond $\{x,x'\}$ of type~$3$ with an end in one of the thin layers
$\mathbb{X}^{+}_{j_0,\ell}=\mathbb{X}_{j_0,\ell+1}$, $1\leq\ell\leq5$, or $\mathbb{X}^{-}_{j_0+1,k}=\mathbb{X}_{j_0+1,k-1}$, $2\leq k\leq6$.
By Lemma~\ref{lem:layer-bonds}\eqref{it:layer-types} and~\eqref{it:layer-phases},
the bonds of type~$3$ with an end in $\mathbb{X}_{j_0,2},\dots,\mathbb{X}_{j_0,5}$ lie over links of the plane $\{x_1=j_0\}$,
those with an end in $\mathbb{X}_{j_0+1,2},\dots,\mathbb{X}_{j_0+1,5}$ lie over links of the plane $\{x_1=j_0+1\}$,
and those with an end in $\mathbb{X}_{j_0,6}$ or in $\mathbb{X}_{j_0+1,1}$ lie over links from the plane $\{x_1=j_0\}$ to the plane $\{x_1=j_0+1\}$.
Hence $Y_P$, like $P_r$, reads only links of these two planes and links between them.
Every basis element $\exp(i\langle\delta,\hat A\rangle)\cdot\Psi$ in the representation of $Y_P$ is gauge invariant,
its exponent $\delta$ is supported on links read by $Y_P$, and $\Psi$ is a monomial in fields at sites read by $Y_P$.
Each of these basis elements is therefore a two-layer monomial at $c$, and $Y_P\in\mathcal{P}_1(c)$ (Definition~\ref{def:two-layer}).
As no torus, factor, sequence of sides, data, index or strength enters the construction, $Y_P$ depends on none of them.

\emph{A chain of universal identities.}
Let $s<r$; the element $P_s$ is read in $B_s$, and $B_s^{+8}=B_{s+1}$.
By Proposition~\ref{prop:outer-plane}\eqref{it:lowest-plane} if $s<r_1$, and by Proposition~\ref{prop:outer-plane}\eqref{it:highest-plane} if $r_1\leq s<r$,
the triple
\[
\mathcal{U}_s:=\left(P_s;\ \left(d\delta_{s,a},P_{s,a}\right)_{a\in I_s};\ B_{s+1}\right)
\]
is a universal identity (Definition~\ref{def:universal-identity} and Remark~\ref{rem:universal-statements};
recall that $\Phi^{(\delta)}=\Phi^{[d\delta]}$ by Definition~\ref{def:factors}),
and the sum of the elements on its right side is $\sum_{a\in I_s}P_{s,a}=P_{s+1}$.
Let $\ell\in\{0,\dots,9\}$.
The elements $P_{r,\ell}$ and $P_{r,\ell+1}$ are read in $B_r^{+(\ell+1)}$,
and Proposition~\ref{prop:schwinger-dyson}\eqref{it:replacement-identity}, for the upward map if $\ell\leq4$ and for the downward map if $\ell\geq5$, gives
$Z_{\mathbb{T}}[\Phi\cdot P_{r,\ell}]=Z_{\mathbb{T}}[\Phi\cdot P_{r,\ell+1}]$
for every torus $\mathbb{T}$ in which $B_r^{+(\ell+1)}$ embeds and every factor $\Phi$ on $\mathbb{T}$.
Thus
\[
\mathcal{U}_{r,\ell}:=\left(P_{r,\ell};\ \left(0,P_{r,\ell+1}\right);\ B_r^{+(\ell+1)}\right),
\]
with an index set of one element, is a universal identity of the first kind.

\emph{Proof of \eqref{it:reduction-general}.}
Let $\mathfrak{d}$ be data with the observable $O$, let $h$ be the width of $\mathcal{Q}_{\mathfrak{d}}$,
let $C_{\mathcal{Q}_{\mathfrak{d}}}$ be formed with $\mathcal{Q}_{\mathfrak{d}}$ and $O$ as in Proposition~\ref{prop:contour-shift}\eqref{it:shift-arrays},
and let $g\in W$.
In general neither the torus $D(n_i)$ nor the array $O_{\mathfrak{d}}$ is invariant under $g$.
The identities of the construction are therefore carried over by $g$,
whereas all estimates are made on $D(n_i)$ itself;
there the identities are applied with the factor $\exp(t\cdot O_{\mathfrak{d}})$, whose shifts produce the errors in \eqref{eq:elimination-estimate}.
Write $P^g_s:=(P_s)^g$, $P^g_{s,a}:=(P_{s,a})^g$ and $P^g_{r,\ell}:=(P_{r,\ell})^g$.
By Lemma~\ref{lem:reflection-identities}\eqref{it:reflected-universal}, for every $s<r$ the image
\[
\mathcal{U}^g_s=\left(P^g_s;\ \left(\eta_{s,a},P^g_{s,a}\right)_{a\in I_s};\ gB_{s+1}\right),
\qquad
\eta_{s,a}:=\operatorname{sgn}(g)\cdot g_{*}d\delta_{s,a},
\]
of $\mathcal{U}_s$ is a universal identity,
and for $\ell=0,\dots,9$ the image $\mathcal{U}^g_{r,\ell}=\left(P^g_{r,\ell};\,\left(0,P^g_{r,\ell+1}\right);\,g\left(B_r^{+(\ell+1)}\right)\right)$ of $\mathcal{U}_{r,\ell}$
is a universal identity of the first kind (Definition~\ref{def:universal-identity}).
Applied to $P_{s+1}=\sum_{a\in I_s}P_{s,a}$, the additive map $G\mapsto G^g$ \eqref{eq:setting-group-action} gives $P^g_{s+1}=\sum_{a\in I_s}P^g_{s,a}$,
and each $P^g_{s,a}$ belongs to $\mathcal{P}$ by Proposition~\ref{prop:class-structure}\eqref{it:class-stable}.

For $s<r$ and $a\in I_s$ let $C_{s,a}:=C_{P^g_{s,a}}<\infty$ be the constant of the element $P^g_{s,a}$ (Definition~\ref{def:element-constant}).
Proposition~\ref{prop:a-priori}\eqref{it:a-priori-element}, applied to $P^g_{s,a}$ and to the box $(gB_{s+1})^{+h}$,
yields an integer $k_{s,a}\in\N$ and an index $i_{s,a}$ such that
\[
\left|\mu^t_i\left(\varphi_\beta\cdot P^g_{s,a}\right)\right|\leq C_{s,a}\cdot\gamma_{t,i}(k_{s,a})
\]
for all $i\geq i_{s,a}$, all $\beta\in C^2_{\mathrm{c}}$ with $\supp\beta\subset(gB_{s+1})^{+h}$ and all $t\geq0$.
Put
\[
C_s:=\max_{a\in I_s}C_{\mathcal{Q}_{\mathfrak{d}}}(\eta_{s,a}),
\]
and $C_s:=0$ if $I_s$ is empty.
Choose $i_0$ such that every $i\geq i_0$ has the following four properties.
The torus $D(n_i)$ carries $\mathcal{Q}_{\mathfrak{d}}$ (Definition~\ref{def:data}).
The box $\left(g\left(B^{+(8r+10)}\right)\right)^{+h}$ embeds in $D(n_i)$;
this holds for all large $i$,
because the period lattice of $D(n_i)=D(n_i/2,n_i,n_i)$ has the systole $n_i/2$ (Lemma~\ref{lem:systoles}) and $n_i\to\infty$.
One has $i\geq i_{s,a}$ for all $s<r$ and $a\in I_s$.
Finally $\Delta_i\leq1$, which holds for all large $i$ by Theorem~\ref{thm:four-tori}.
For $s<r$ and $\ell\leq10$ the boxes $B_{s+1}=B^{+8(s+1)}$ and $B_r^{+\ell}=B^{+(8r+\ell)}$ lie in $B^{+(8r+10)}$;
since $g$ and the enlargement by $h$ both preserve inclusions,
the box $\left(g\left(B^{+(8r+10)}\right)\right)^{+h}$ contains the boxes $(gB_{s+1})^{+h}$, $s<r$, and $g\left(B_r^{+\ell}\right)$, $\ell\leq10$,
and these boxes embed in $D(n_i)$ as well, for every $i\geq i_0$.

Now let $i\geq i_0$ and $0\leq t\leq1$.
For every $k\in\N$, Definition~\ref{def:amplification-factor} gives $\gamma_{t,i}(k)\leq\exp\left((2|O|_1+|Q_0|)\cdot k^3\right)$, since $t\leq1$ and $\Delta_i\leq1$;
and $\exp(C_s\cdot t)-1\leq(\exp(C_s)-1)\cdot t$, since the function $t\mapsto\exp(C_s\cdot t)-1$ is convex and vanishes at $t=0$.
For the universal identity $\mathcal{U}^g_s$ the sum of the elements on the right side is $P^g_{s+1}$,
and the box of the supremum in \eqref{eq:elimination-estimate} is $(gB_{s+1})^{+h}$.
Therefore \eqref{eq:elimination-estimate}, the bound provided by Proposition~\ref{prop:a-priori}\eqref{it:a-priori-element} and the two inequalities just noted give,
for $s=0,\dots,r-1$,
\begin{equation}\label{eq:elimination-step-bound}
\left|\mu^t_i\left(P^g_s\right)-\mu^t_i\left(P^g_{s+1}\right)\right|
\leq\left(\exp(C_s)-1\right)\cdot t\cdot\sum_{a\in I_s}C_{s,a}\cdot\exp\left(\left(2|O|_1+|Q_0|\right)\cdot k_{s,a}^3\right).
\end{equation}

The third phase contributes no error.
For $\ell=0,\dots,9$ the identity \eqref{eq:elimination-universal} of $\mathcal{U}^g_{r,\ell}$,
taken on the torus $D(n_i)$ with the factor $\exp(t\cdot O_{\mathfrak{d}})$, reads
\[
Z_{D(n_i)}\left[\exp(t\cdot O_{\mathfrak{d}})\cdot P^g_{r,\ell}\right]=Z_{D(n_i)}\left[\exp(t\cdot O_{\mathfrak{d}})\cdot P^g_{r,\ell+1}\right].
\]
On $D(n_i)$ one has $|O_{\mathfrak{d}}|\leq|O|_1\cdot n_i^3/|Q_0|$ (Definition~\ref{def:arrays}),
and Lemma~\ref{lem:perturbed-denominators} together with Theorem~\ref{thm:partition-positivity} gives
\[
Z_{D(n_i)}\left[\exp(t\cdot O_{\mathfrak{d}})\right]\geq\exp\left(-t\cdot|O|_1\cdot\frac{n_i^3}{|Q_0|}\right)\cdot Z_{D(n_i)}[1]>0 .
\]
Division by this number yields $\mu^t_i\left(P^g_{r,\ell}\right)=\mu^t_i\left(P^g_{r,\ell+1}\right)$ for $\ell=0,\dots,9$, and hence
\[
\mu^t_i\left(P^g_r\right)=\mu^t_i\left((Y_P)^g\right).
\]
As $P_0=P$, the triangle inequality and \eqref{eq:elimination-step-bound} give
\[
\left|\mu^t_i\left(P^g\right)-\mu^t_i\left((Y_P)^g\right)\right|
\leq\sum_{s=0}^{r-1}\left|\mu^t_i\left(P^g_s\right)-\mu^t_i\left(P^g_{s+1}\right)\right|
\leq C_{P,g,\mathfrak{d}}\cdot t
\]
for $i\geq i_0$ and $0\leq t\leq1$, with the finite constant
\[
C_{P,g,\mathfrak{d}}:=\sum_{s=0}^{r-1}\left(\exp(C_s)-1\right)\cdot\sum_{a\in I_s}C_{s,a}\cdot\exp\left(\left(2|O|_1+|Q_0|\right)\cdot k_{s,a}^3\right).
\]
This is \eqref{eq:elimination-general}.

\emph{Proof of \eqref{it:reduction-link-plane}.}
Let $\theta_c:=\theta_{(x_1,c)}$ be the reflection in the plane $\{x_1=c\}$.
For the identity element $1$ of $W$ one has $P^1=P$ and $(Y_P)^1=Y_P$;
for $g=\theta_c$, \eqref{eq:setting-group-action} gives $P^{\theta_c}=\Theta_c P$ and $(Y_P)^{\theta_c}=\Theta_c Y_P$.
So \eqref{eq:elimination-general} for these two elements of $W$ yields both inequalities of \eqref{it:reduction-link-plane}, with
\[
C_{P,\mathfrak{d}}:=\max\left(C_{P,1,\mathfrak{d}},\,C_{P,\theta_c,\mathfrak{d}}\right)
\]
and with the larger of the two indices $i_0$.
For $g=\theta_c$ one has $\operatorname{sgn}(\theta_c)=-1$ and $(\theta_c)_{*}=\theta_c^{*}$, so that $\eta_{s,a}=-\theta_c^{*}d\delta_{s,a}$.
In this case the identities $\mathcal{U}^{\theta_c}_s$ are the reflected identities of Proposition~\ref{prop:outer-plane}\eqref{it:reflected-elimination},
the estimate \eqref{eq:elimination-estimate} for them is the last inequality of Proposition~\ref{prop:outer-plane}\eqref{it:plane-estimate},
the identities $\mathcal{U}^{\theta_c}_{r,\ell}$ are those of Proposition~\ref{prop:schwinger-dyson}\eqref{it:mirrored-rule},
and, by Proposition~\ref{prop:a-priori}\eqref{it:a-priori-mirror}, the constants $C_{s,a}$ may be taken equal to $C_{P_{s,a}}$, the numbers that occur for $g=1$.

\emph{Proof of \eqref{it:reduction-directions}.}
Let $\mu\in\{2,3\}$.
The element $\check P=P^{g_\mu^{-1}}$ lies in $\mathcal{P}$ (Proposition~\ref{prop:class-structure}\eqref{it:class-stable}).
By the first part of this proof, applied to $\check P$, the element $Y_{\check P}$ lies in $\mathcal{P}_1(c)$,
so that $Y^{\mu}_P=(Y_{\check P})^{g_\mu}$ lies in $\mathcal{P}_1(\mu,c)$ (Definition~\ref{def:other-directions});
like $Y_{\check P}$, it is defined on $\Z^3$ and does not depend on $\mathfrak{n}$, $\mathfrak{d}$, $i$ or $t$.
Elements of $\mathcal{P}$ are even, and $(G^{g})^{g'}=G^{g'g}$ for even $G$ and $g,g'\in W$ \eqref{eq:setting-group-action}.
Hence $(\check P)^{g_\mu}=P^{g_\mu g_\mu^{-1}}=P$, while $(Y_{\check P})^{g_\mu}=Y^{\mu}_P$ by definition.
Moreover $G^{g_\mu\theta_c}=(G^{\theta_c})^{g_\mu}=(\Theta_c G)^{g_\mu}$ for $G\in\mathcal{P}$, and \eqref{eq:elimination-conjugation} turns this into
\[
(\check P)^{g_\mu\theta_c}=\Theta_{\mu,c}\left((\check P)^{g_\mu}\right)=\Theta_{\mu,c} P,
\qquad
(Y_{\check P})^{g_\mu\theta_c}=\Theta_{\mu,c}\left((Y_{\check P})^{g_\mu}\right)=\Theta_{\mu,c} Y^{\mu}_P .
\]
Now \eqref{eq:elimination-general}, applied to the element $\check P$ with $g=g_\mu$ and with $g=g_\mu\theta_c$,
gives the two inequalities of \eqref{it:reduction-directions}, with
\[
C:=\max\left(C_{\check P,g_\mu,\mathfrak{d}},\,C_{\check P,g_\mu\theta_c,\mathfrak{d}}\right)
\]
and with the larger of the two indices $i_0$.
\end{proof}

\begin{remark}[The Yang--Mills case]\label{rem:yang-mills-elimination}
Each step of this section,
namely the maps $\mathfrak{R}^{\pm}_{\mathbb{L}}$ of Proposition~\ref{prop:schwinger-dyson}\eqref{it:replacement-maps},
the normal form of Proposition~\ref{prop:normal-form},
the shift of the contour in Proposition~\ref{prop:contour-shift}\eqref{it:shift-identity}
and the elimination of an outer plane in Proposition~\ref{prop:outer-plane}\eqref{it:lowest-plane} and~\eqref{it:highest-plane},
sends elements of $\mathcal{P}$ to elements of $\mathcal{P}$ without any approximation,
and every term that it produces is again gauge invariant.
The reduction therefore involves neither a re-approximation by polynomials nor an average over the gauge group,
whereas the corresponding step of the argument for Yang--Mills theory in four dimensions contains both.
For the noncompact field both are out of reach.
In the Yang--Mills case the re-approximation is an application of the Stone--Weierstrass theorem on a compact space,
the product of copies of the compact structure group over a finite set of links,
while here the corresponding fibre is $\R$, which is not compact;
and the group $\R^{\Lambda}$ of the real gauge transformations on a finite set $\Lambda$ of sites carries no invariant probability measure.

The same difference appears in the eliminated observable.
In the Yang--Mills case it is obtained from the given observable by multiplication with a strictly positive weight,
by application of the conjugate $K^{-1}fK$ of a function $f$ of the eliminated layer by the crossing kernel $K$,
and by division by the weight.
Expanded in Peter--Weyl components, the conjugate $K^{-1}fK$ becomes an infinite series;
the observable that results, though smooth, is not a polynomial, and gauge invariance fails for its individual terms.
In the present case let $\Xi$ be a rectangular cross-section, with the slice torus $\Lambda_{\mathrm{s}}$ and the group $\Gamma_{\mathrm{s}}$ of Definition~\ref{def:slices}
and the operators of Definition~\ref{def:slice-operators},
and let $\delta\in\R^{\Lambda^{(1)}_{\mathrm{s}}}$ be a real one-cochain on the links of $\Lambda_{\mathrm{s}}$ such that $\langle\delta,\gamma\rangle\in2\pi\Z$ for every $\gamma\in\Gamma_{\mathrm{s}}$.
Then the function $A\mapsto\exp(i\langle\delta,A\rangle)$ is $\Gamma_{\mathrm{s}}$-periodic;
let $f:=\exp(i\langle\delta,A\rangle)$ be the unitary operator of multiplication by it,
with $E$, $\langle\delta,E\rangle$, $|E|^2$ and $|\delta|^2$ as in Remark~\ref{rem:imaginary-translation}.
By that remark, conjugation by the Gaussian factor $\exp\left(-\frac12|E|^2\right)$ of the crossing operator $K$ turns $f$ into
\[
f\cdot\exp\left(\tfrac12|\delta|^2\right)\cdot\exp(\langle\delta,E\rangle);
\]
the conjugate $\exp\left(\frac12|E|^2\right)\cdot f\cdot\exp\left(-\frac12|E|^2\right)$ and this operator are both defined exactly on the domain of $\exp(\langle\delta,E\rangle)$,
and they coincide there.
The imaginary translation $\exp(\langle\delta,E\rangle)$ shifts the argument $A$ to $A-i\delta$;
applied to an exponential, this operator returns a scalar multiple of an exponential.
As for the weight, let $\mathbb{T}$ be a torus, $S$ a finite set of links of one coordinate plane and $\delta\in\R^S$,
as in Proposition~\ref{prop:contour-shift}.
The substitution $A\mapsto A-i\delta$ multiplies the weight $\exp\left(-\frac12\|F\|^2\right)\cdot\exp\left(-S_{\mathbb{T}}(A)\right)$
in the definition of $Z_{\mathbb{T}}$ (Definition~\ref{def:functional}) by
\[
\exp\left(\tfrac12\|d\delta\|^2\right)\cdot\varphi_{d\delta}\cdot\mathfrak{V}_\delta ,
\]
which is an element of $\mathcal{P}$ (Proposition~\ref{prop:contour-shift}).
Hence the construction never divides by the weight.
\end{remark}

\section{No layered states}\label{sec:layered}

Throughout this section the Wilson parameter satisfies $0<\rho<1$,
while the weights $w_0,\dots,w_3>0$, the mass parameter $M\in\R$ and the lattice charge $q>0$ are arbitrary.

\subsection{Laws on $\Omega$}\label{sec:curvature-laws}

\paragraph{Laws.}
A \emph{law} is a Borel probability measure on $\Omega$.
For a law $\mu$ and a Borel function $G$ on $\Omega$ that is non-negative or $\mu$-integrable we write
\begin{equation}\label{eq:layers-law}
\mu(G):=\int\mathrm{d}\mu(F)\,G(F).
\end{equation}
The space $C^2$ of the real two-cochains is the product of countably many copies of $\R$, one for each plaquette;
with the product topology it is a Polish space.
Each value $(dF)(x)$ is a finite linear combination of coordinates of $F$, hence a continuous function of $F$.
So $\Omega$ is closed in $C^2$, and it is a Polish space (Definition~\ref{def:characters}).
Its Borel sets are the traces on $\Omega$ of the sets of the product $\sigma$-algebra of $C^2$.
Consequently the Borel $\sigma$-algebra of $\Omega$ is generated by the \emph{cylinder sets},
the preimages of the Borel subsets of $\R^{\Sigma_0}$ under the maps
\[
\Omega\to\R^{\Sigma_0},\qquad F\mapsto(F_\sigma)_{\sigma\in\Sigma_0},
\]
where $\Sigma_0$ ranges over the finite sets of plaquettes.
The cylinder sets include $\Omega$, and the intersection of two of them is again a cylinder set.

The weak topology on the set of laws is the coarsest topology for which $\mu\mapsto\mu(G)$ is continuous for every bounded continuous real function $G$ on $\Omega$;
thus laws $\mu_n$ converge weakly to a law $\mu$ if and only if $\mu_n(G)\to\mu(G)$ for every bounded continuous function $G$ on $\Omega$.
This topology is metrisable, by the Prokhorov metric \cite[Theorem~6.8]{Billingsley1999}.
Hence a set of laws is compact if and only if it is sequentially compact,
it is closed if and only if it contains the weak limit of each weakly convergent sequence of its elements,
and the weak limit points of a sequence of laws are exactly the weak limits of its subsequences.

\paragraph{The action of $W$.}
Let $g\in W$.
The pull-back $g^*$ of Definition~\ref{def:cochains} permutes the coordinates of $C^2$ up to signs and commutes with $d$;
it is therefore a homeomorphism of $\Omega$ onto itself, with inverse $(g^{-1})^*$.
Moreover $(gg')^*=g'^*g^*$ for $g'\in W$ (Definition~\ref{def:cochains}).
For a function $G$ on $\Omega$ put
\begin{equation}\label{eq:layers-action}
(gG)(F):=G(g^*F),
\end{equation}
as for trigonometric polynomials in Definition~\ref{def:characters}.
This defines an action of $W$: for $g,g'\in W$ one has $g(g'G)=(gg')G$, because
\[
(g(g'G))(F)=(g'G)(g^*F)=G(g'^*g^*F)=G((gg')^*F).
\]
The action commutes with products and with complex conjugation of functions.
A law $\mu$ is \emph{invariant under} $g$ if $\mu(gG)=\mu(G)$ for every bounded Borel function $G$ on $\Omega$,
that is, if the image $(g^*)_*\mu$ of $\mu$ under $g^*$ equals $\mu$.
A law is invariant at a mirror $(\ell,c)$ in the sense of Definition~\ref{def:reflection-positivity} if and only if it is invariant under $\theta_{\ell,c}$.

\paragraph{Local functions.}
For a subset $U$ of $\R^3$ let $\mathcal{O}(U)$ be the set of the functions $G\colon\Omega\to\C$ of the form
\[
G(F)=G((F_\sigma)_{\sigma\in\Sigma_0}),
\]
where $\Sigma_0$ is a finite set of plaquettes contained in $U$ and, by abuse of notation, $G$ also denotes a bounded Borel function on $\R^{\Sigma_0}$.
Let $\mathcal{O}_{\mathrm{cont}}(U)$ be the set of the elements of $\mathcal{O}(U)$ that are continuous on $\Omega$.

Fix a mirror $(\ell,c)$.
Every point of a plaquette is a convex combination of the four corners of the plaquette;
as $\ell$ is linear, the values of $\ell$ on a plaquette thus lie between the least and the greatest of its values at the corners.
Consequently a plaquette is contained in the open half-space $\{\ell>c\}$ if and only if $\ell>c$ at each of its corners,
that is, if and only if it belongs to the set $\Sigma^{+}(\ell,c)$ of Definition~\ref{def:reflection-positivity};
so the set $\mathcal{O}^{+}(\ell,c)$ of that definition coincides with $\mathcal{O}(\{\ell>c\})$.
In the notation~\eqref{eq:layers-action}, a law $\mu$ is thus reflection positive at $(\ell,c)$ if and only if
\[
\mu\left(\overline{\theta_{\ell,c}G}\cdot G\right)\ge0\qquad\text{for every }G\in\mathcal{O}(\{\ell>c\}).
\]

Now let $g\in W$.
Since $(g^*F)_\sigma=\pm F_{g\sigma}$ for every plaquette $\sigma$ (Definition~\ref{def:cochains}),
the map $G\mapsto gG$ carries a function of the $F_\sigma$, $\sigma\in\Sigma_0$, to a function of the $F_{g\sigma}$, $\sigma\in\Sigma_0$,
and $g\sigma\subset gU$ whenever $\sigma\subset U$;
it preserves continuity, $g^*$ being a homeomorphism.
Applied to $g$ and to $g^{-1}$, this gives
\begin{equation}\label{eq:layers-local}
g\,\mathcal{O}(U)=\mathcal{O}(gU),\qquad g\,\mathcal{O}_{\mathrm{cont}}(U)=\mathcal{O}_{\mathrm{cont}}(gU).
\end{equation}

\paragraph{Bosonic functions and the laws of tori.}
Let $\Pi$ be a mirror, and let $G$ be bosonic, that is, a function of $F$ (Definition~\ref{def:integrands}).
The only monomial of $G$ is the unit, with the coefficient $G(d\hat A)$, and $\Theta_\Pi(1)=1$.
By~\eqref{eq:setting-reflection-coefficients}, and because $\theta_\Pi^*$ commutes with $d$ (Definition~\ref{def:cochains}),
$\Theta_\Pi G$ is the complex conjugate of $\theta_\Pi G$:
\begin{equation}\label{eq:layers-bosonic}
(\Theta_\Pi G)(F)=\overline{G(\theta_\Pi^*F)}.
\end{equation}
Next let $\mathbb{T}=(\Gamma,\chi)$ be a torus and $G\in\mathcal{O}(\R^3)$.
As a product without boxes, $G$ belongs to $\mathcal{A}(\R^3)$ and embeds in $\mathbb{T}$;
no condition is imposed on the plaquettes read by a bounded factor, which is evaluated at the periodic extension $\widetilde F$ (Definition~\ref{def:integrands}).
The law $\mu_{\mathbb{T}}$ is the law of $\widetilde F$ under $Z_{\mathbb{T}}[\,\cdot\,]/Z_{\mathbb{T}}[1]$ (Definition~\ref{def:functional}); therefore
\[
\frac{Z_{\mathbb{T}}[G]}{Z_{\mathbb{T}}[1]}=\mu_{\mathbb{T}}(G).
\]
Finally let $\mathfrak{d}$ be data, $t\ge0$, and $i$ an index such that the torus $\mathbb{T}=(\Gamma,\chi):=D(n_i)$ carries $\mathcal{Q}_{\mathfrak{d}}$.
For the function $G$ above, Definition~\ref{def:data} relates the perturbed functional to the law $\mu^t_i$ on $\Omega$:
\[
\frac{Z_{\mathbb{T}}[\exp(t\cdot O_{\mathfrak{d}})\cdot G]}{Z_{\mathbb{T}}[\exp(t\cdot O_{\mathfrak{d}})]}=\mu^t_i(G),
\]
the right side being the integral of $G$ against this law in the notation~\eqref{eq:layers-law}.
The law $\mu^t_i$ is the image under $(A,z)\mapsto\widetilde F$ of the probability measure on $\Conf(\Gamma)$
whose density with respect to the measure of Definition~\ref{def:configurations} is the function
\[
\exp(t\cdot O_{\mathfrak{d}})\cdot\exp\left(-\frac12\|F\|^2\right)\cdot\det D^{\mathbb{T}}_A
\]
divided by its integral.
This function is non-negative, since $\det D^{\mathbb{T}}_A\ge0$ (Proposition~\ref{prop:accretivity});
its integral is $Z_{\mathbb{T}}[\exp(t\cdot O_{\mathfrak{d}})]$ (Definition~\ref{def:functional}),
and Lemma~\ref{lem:perturbed-denominators} and Theorem~\ref{thm:partition-positivity} bound this number from below:
\[
Z_{\mathbb{T}}[\exp(t\cdot O_{\mathfrak{d}})]\ge\exp\left(-t\cdot|O|_1\cdot\frac{n_i^3}{|Q_0|}\right)\cdot Z_{\mathbb{T}}[1]>0,
\]
the folded array being real and of modulus at most $|O|_1\cdot n_i^3/|Q_0|$ on $\mathbb{T}$ (Definition~\ref{def:arrays}).

\paragraph{Characters.}
Let $\beta\in C^2_{\mathrm{c}}$, let $\Pi$ be a mirror and $\theta:=\theta_\Pi$.
For $F\in\Omega$ and every plaquette $\sigma$ one has $(\theta^*F)_\sigma=\pm F_{\theta\sigma}$ and $(\theta^*\beta)(\theta\sigma)=\pm\beta(\sigma)$ with the same sign,
because $\theta^*\theta^*=(\theta\theta)^*$ is the identity (Definition~\ref{def:cochains}).
Hence $\beta(\sigma)\cdot(\theta^*F)_\sigma=(\theta^*\beta)(\theta\sigma)\cdot F_{\theta\sigma}$,
and summation over $\sigma$, with $\theta$ permuting the plaquettes, gives
\[
\langle\beta,\theta^*F\rangle=\langle\theta^*\beta,F\rangle .
\]
It follows that $\theta\varphi_\beta=\varphi_{\theta^*\beta}$, and by~\eqref{eq:layers-bosonic} that
\begin{equation}\label{eq:layers-characters}
\Theta_\Pi\varphi_\beta=\varphi_{-\theta_\Pi^*\beta},
\end{equation}
again a character, hence a bounded continuous function on $\Omega$.
If $\Pi$ is the link plane $(x_1,j_0+\tfrac12)$ with $j_0\in\Z$,
the pull-back $\theta_\Pi^*$ is written out in~\eqref{eq:setting-slice-reflections} of Definition~\ref{def:slices}: for every $j\in\Z$,
\[
F^{\mathrm{el}}_j(\theta_\Pi^*F)=-F^{\mathrm{el}}_{2j_0-j}(F),\qquad F^{\mathrm{mag}}_j(\theta_\Pi^*F)=F^{\mathrm{mag}}_{2j_0+1-j}(F).
\]

As an element of $\Gr(\mathbb{X})^{C^1}$ (Definition~\ref{def:integrands}), the character $\varphi_\beta$ is the map $\hat A\mapsto\varphi_\beta(d\hat A)$ on the real one-cochains $\hat A$.
For every $\hat A\in C^1$ one has $\langle\beta,d\hat A\rangle=\langle d^*\beta,\hat A\rangle$.
Indeed, let $\hat A'\in C^1_{\mathrm{c}}$ agree with $\hat A$ at the edges of the plaquettes $\sigma$ with $\beta(\sigma)\neq0$ and vanish at all other links.
Then $d\hat A'=d\hat A$ at these plaquettes.
Moreover, let $1_b$ denote the one-cochain with the value $1$ at a link $b$ and $0$ elsewhere;
if $b$ is not such an edge, then $d1_b$ vanishes at every plaquette $\sigma$ with $\beta(\sigma)\neq0$, so that $(d^*\beta)(b)=\langle\beta,d1_b\rangle=0$.
Therefore, by the definition of $d^*$ (Definition~\ref{def:cochains}),
\[
\langle\beta,d\hat A\rangle=\langle\beta,d\hat A'\rangle=\langle d^*\beta,\hat A'\rangle=\langle d^*\beta,\hat A\rangle .
\]
So $\varphi_\beta$ is the map $\hat A\mapsto\exp(i\langle\delta,\hat A\rangle)$ with the real cochain $\delta:=d^*\beta\in C^1_{\mathrm{c}}$,
which satisfies $d^*\delta=d^*d^*\beta=0$ because $d\circ d=0$;
it is therefore an element of $\mathcal{P}$ whose monomial is $\Psi=1$ (Definition~\ref{def:polynomials}).

\begin{lemma}[Characters determine a law]\label{lem:characters-determine}
\begin{enumerate}
\item[(a)] Let $\mu$ and $\mu'$ be laws on $\Omega$ with $\mu(O)=\mu'(O)$ for every $O\in\mathcal{T}$.
Then $\mu=\mu'$.
\item[(b)] Let $\mathcal{T}_{\Q}$ be the set of the finite sums $\sum_ka_k\cdot\varphi_{\beta_k}$ with $a_k\in\Q+i\Q$ and with $\beta_k\in C^2_{\mathrm{c}}$ taking rational values.
For every law $\mu$ on $\Omega$ the set $\mathcal{T}_{\Q}$ is dense in $L^2(\mu)$.
Moreover $\mathcal{T}_{\Q}$ is countable, and it is stable under every $g\in W$ and under complex conjugation.
\end{enumerate}
\end{lemma}

\begin{proof}
(a) Let $\beta\in C^2_{\mathrm{c}}$.
The functions
\[
\cos\langle\beta,F\rangle=\frac{\varphi_\beta(F)+\varphi_{-\beta}(F)}{2},\qquad
\sin\langle\beta,F\rangle=\frac{\varphi_\beta(F)-\varphi_{-\beta}(F)}{2i}
\]
are real-valued trigonometric polynomials, and $\varphi_\beta=\cos\langle\beta,\cdot\rangle+i\sin\langle\beta,\cdot\rangle$;
so the hypothesis gives $\mu(\varphi_\beta)=\mu'(\varphi_\beta)$.
Let $\Sigma_0$ be a finite set of plaquettes, and let $\mu_{\Sigma_0}$ and $\mu'_{\Sigma_0}$ be the images of $\mu$ and $\mu'$ under $F\mapsto(F_\sigma)_{\sigma\in\Sigma_0}$,
probability measures on $\R^{\Sigma_0}$.
Identify $\R^{\Sigma_0}$ with the space of the two-cochains that vanish outside $\Sigma_0$.
For such $\beta$ one has $\langle\beta,F\rangle=\sum_{\sigma\in\Sigma_0}\beta(\sigma)\cdot F_\sigma$;
so the characteristic function of $\mu_{\Sigma_0}$ takes the value $\mu(\varphi_\beta)$ at $\beta$,
and that of $\mu'_{\Sigma_0}$ the value $\mu'(\varphi_\beta)$.
By the uniqueness theorem for characteristic functions of probability measures on a Euclidean space \cite[Section~29]{Billingsley1995}, $\mu_{\Sigma_0}=\mu'_{\Sigma_0}$;
in other words, $\mu$ and $\mu'$ agree on the cylinder sets.
As noted at the beginning of this subsection, the cylinder sets generate the Borel $\sigma$-algebra of $\Omega$, and any two of them intersect in a cylinder set.
Two probability measures that agree on a family of sets closed under pairwise intersections agree on the $\sigma$-algebra generated by that family
(Dynkin's theorem, \cite[Theorem~3.3]{Billingsley1995});
hence $\mu=\mu'$.

(b) The elements of $\mathcal{T}_{\Q}$ are bounded and continuous, so $\mathcal{T}_{\Q}\subset L^2(\mu)$.
Let $H\in L^2(\mu)$ be orthogonal to $\mathcal{T}_{\Q}$.
In particular $\mu(\overline{H}\cdot\varphi_\beta)=0$ for every $\beta\in C^2_{\mathrm{c}}$ with rational values.
Fix a finite set $\Sigma_0$ of plaquettes and identify $\R^{\Sigma_0}$ with the two-cochains that vanish outside $\Sigma_0$, as in~(a).
The function $\beta\mapsto\mu(\overline{H}\cdot\varphi_\beta)$ is continuous on $\R^{\Sigma_0}$ by dominated convergence,
since $|\overline{H}\cdot\varphi_\beta|=|H|$ is $\mu$-integrable;
it vanishes at the rational points, which are dense, and hence everywhere.
As $\Sigma_0$ is arbitrary, $\mu(\overline{H}\cdot\varphi_\beta)=0$ for every $\beta\in C^2_{\mathrm{c}}$.
Let $\lambda$ be the complex Borel measure on $\Omega$ with the density $\overline{H}$ with respect to $\mu$,
which is finite because $H$ is $\mu$-integrable,
and let $\lambda_{\Sigma_0}$ be its image under $F\mapsto(F_\sigma)_{\sigma\in\Sigma_0}$.
The Fourier transform of $\lambda_{\Sigma_0}$ at $\beta\in\R^{\Sigma_0}$ is $\mu(\overline{H}\cdot\varphi_\beta)=0$.
The transform of the complex conjugate measure $\overline{\lambda_{\Sigma_0}}$ at $\beta$ is the complex conjugate of the transform of $\lambda_{\Sigma_0}$ at $-\beta$;
hence the real part and the imaginary part of $\lambda_{\Sigma_0}$ are finite signed measures with vanishing Fourier transform.
For each of them, the positive part and the negative part have the same Fourier transform, and in particular, at $\beta=0$, the same mass.
If this mass is zero, both parts vanish;
otherwise, divided by the mass, they are probability measures with the same characteristic function, and they coincide by the uniqueness theorem.
Therefore $\lambda_{\Sigma_0}=0$, that is, $\lambda$ vanishes on every cylinder set.
The Borel subsets of $\Omega$ on which $\lambda$ vanishes form a family that contains $\Omega$ and the cylinder sets
and is stable under proper differences and under increasing countable unions;
by Dynkin's $\pi$--$\lambda$ theorem \cite[Theorem~3.2]{Billingsley1995} this family contains every Borel set.
Hence $\overline{H}=0$ $\mu$-almost everywhere, and $H=0$ in $L^2(\mu)$.
Since $\mathcal{T}_{\Q}$ is a vector space over $\Q+i\Q$ and $\Q+i\Q$ is dense in $\C$, the closure of $\mathcal{T}_{\Q}$ in $L^2(\mu)$ is a closed complex subspace;
its orthogonal complement is $\{0\}$, so it is all of $L^2(\mu)$.

There are countably many finitely supported two-cochains with rational values and countably many elements of $\Q+i\Q$;
so $\mathcal{T}_{\Q}$ is countable.
Complex conjugation maps $\sum_ka_k\cdot\varphi_{\beta_k}$ to $\sum_k\overline{a_k}\cdot\varphi_{-\beta_k}$, an element of $\mathcal{T}_{\Q}$.
For $g\in W$ the pull-back $g^*$ is orthogonal, and $g\varphi_\beta=\varphi_{(g^{-1})^*\beta}$ (Definition~\ref{def:characters});
the cochain $(g^{-1})^*\beta$ takes rational values whenever $\beta$ does, so $g$ maps $\mathcal{T}_{\Q}$ into itself.
\end{proof}

\begin{lemma}[Gaussian domination: tails and compactness]\label{lem:domination-compactness}
For a law $\mu$ on $\Omega$ and $\beta\in C^2_{\mathrm{c}}$ consider the inequality
\begin{equation}\tag{G$_\beta$}\label{cond:domination-beta}
\int\mathrm{d}\mu(F)\,\exp(\langle\beta,F\rangle)\le\exp\left(\frac{1}{2}\,\|\beta\|^2\right).
\end{equation}
\begin{enumerate}
\item[(a)] Let $\sigma$ be a plaquette, $1_\sigma$ the two-cochain with the value $1$ at $\sigma$ and $0$ elsewhere, and $s>0$.
If $\mu$ satisfies~\eqref{cond:domination-beta} for $\beta=s\cdot1_\sigma$ and for $\beta=-s\cdot1_\sigma$, then
\[
\mu(\{F\in\Omega:\ |F_\sigma|\ge s\})\le2\exp\left(-\frac{s^2}{2}\right).
\]
\item[(b)] A set of laws on $\Omega$ each of which satisfies~\eqref{cond:domination-beta} for all $\beta=\pm s\cdot1_\sigma$,
$\sigma$ a plaquette and $s>0$, is uniformly tight.
\item[(c)] Let laws $\mu_n$ converge weakly to a law $\mu$, and let $\beta\in C^2_{\mathrm{c}}$.
If $\mu_n$ satisfies~\eqref{cond:domination-beta} for all large $n$, then $\mu$ satisfies~\eqref{cond:domination-beta}.
\item[(d)] The set $\mathcal{M}_{(\mathrm{G})}$ of the laws on $\Omega$ that satisfy~\eqref{cond:gaussian-domination} is convex and weakly compact.
\end{enumerate}
\end{lemma}

\begin{proof}
(a) For $\beta=s\cdot1_\sigma$ one has $\langle\beta,F\rangle=s\cdot F_\sigma$ and $\|\beta\|^2=s^2$.
On the set where $F_\sigma\ge s$ the function $\exp(s\cdot F_\sigma)$ is at least $\exp(s^2)$;
Chebyshev's inequality and~\eqref{cond:domination-beta} give
\[
\mu(\{F\in\Omega:\ F_\sigma\ge s\})\le\exp(-s^2)\cdot\int\mathrm{d}\mu(F)\,\exp(s\cdot F_\sigma)\le\exp(-s^2)\cdot\exp\left(\frac{s^2}{2}\right)=\exp\left(-\frac{s^2}{2}\right).
\]
The same argument with $\beta=-s\cdot1_\sigma$ bounds the measure of the set where $F_\sigma\le-s$, and the two bounds add up to~(a).

(b) Enumerate the plaquettes as $\sigma_1,\sigma_2,\dots$.
For $0<\varepsilon<1$ and $n\ge1$ put
\[
r_n:=\sqrt{2\log\frac{2^{n+1}}{\varepsilon}},\qquad
\Omega_\varepsilon:=\bigcap_{n\ge1}\{F\in\Omega:\ |F_{\sigma_n}|\le r_n\}.
\]
In the coordinates $(F_{\sigma_n})_{n\ge1}$ the set $\Omega_\varepsilon$ is the intersection of the closed set $\Omega$ with the product of the intervals $[-r_n,r_n]$,
which is compact by Tychonoff's theorem;
so $\Omega_\varepsilon$ is compact.
For every law $\mu$ of the set, (a) with $s=r_n>0$ gives
\[
\mu(\Omega\setminus\Omega_\varepsilon)\le\sum_{n\ge1}\mu(\{F\in\Omega:\ |F_{\sigma_n}|>r_n\})\le\sum_{n\ge1}2\exp\left(-\frac{r_n^2}{2}\right)=\sum_{n\ge1}\varepsilon\cdot2^{-n}=\varepsilon .
\]
The compact set $\Omega_\varepsilon$ does not depend on $\mu$; hence the set of laws is uniformly tight.

(c) Let $r>0$.
The function $F\mapsto\min(\exp(\langle\beta,F\rangle),r)$ is bounded and continuous on $\Omega$,
since $\langle\beta,F\rangle$ is a finite linear combination of coordinates of $F$.
Weak convergence therefore gives
\[
\int\mathrm{d}\mu(F)\,\min(\exp(\langle\beta,F\rangle),r)=\lim_{n\to\infty}\int\mathrm{d}\mu_n(F)\,\min(\exp(\langle\beta,F\rangle),r)\le\exp\left(\frac{1}{2}\,\|\beta\|^2\right),
\]
and monotone convergence as $r\to\infty$ yields~\eqref{cond:domination-beta} for $\mu$.

(d) By Definition~\ref{def:characters}, a law satisfies~\eqref{cond:gaussian-domination} if and only if it satisfies~\eqref{cond:domination-beta} for every $\beta\in C^2_{\mathrm{c}}$.
For fixed $\beta$ the left side of~\eqref{cond:domination-beta} is an affine function of $\mu$, with values in $[0,\infty]$;
so the laws that satisfy~\eqref{cond:domination-beta} form a convex set, and $\mathcal{M}_{(\mathrm{G})}$, the intersection of these sets, is convex.
Let laws $\mu_n\in\mathcal{M}_{(\mathrm{G})}$ converge weakly to a law $\mu$;
by~(c), applied to every $\beta\in C^2_{\mathrm{c}}$, the law $\mu$ lies in $\mathcal{M}_{(\mathrm{G})}$.
The weak topology being metrisable, $\mathcal{M}_{(\mathrm{G})}$ is closed.
By~(b) it is uniformly tight, hence relatively compact by Prokhorov's theorem \cite[Theorem~5.1]{Billingsley1999}.
A closed and relatively compact set is compact.
\end{proof}

\begin{lemma}[A mirror of a law]\label{lem:law-mirror}
Let $\Pi=(\ell,c)$ be a mirror and $\theta:=\theta_\Pi$.
\begin{enumerate}
\item[(a)] Let $\mu$ be a law on $\Omega$ that is invariant under $\theta$.
Then $G\mapsto\theta G$ is a unitary operator on $L^2(\mu)$, and
\[
\left|\mu\left(\overline{\theta G}\cdot H\right)\right|\le\|G\|_{L^2(\mu)}\cdot\|H\|_{L^2(\mu)}\qquad\text{for }G,H\in L^2(\mu).
\]
\item[(b)] Let $\mu$ be a law on $\Omega$ that is invariant under $\theta$,
and suppose that $\mu\left(\overline{\theta G}\cdot G\right)\ge0$ for every $G\in\mathcal{O}_{\mathrm{cont}}(\{\ell>c\})$.
Then $\mu$ is reflection positive at $\Pi$ in the sense of Definition~\ref{def:reflection-positivity},
that is, $\mu\left(\overline{\theta G}\cdot G\right)\ge0$ for every $G\in\mathcal{O}(\{\ell>c\})$;
moreover $\mu\left(\overline{\theta G}\cdot G\right)\ge0$ for every $G\in\mathcal{O}(\{\ell<c\})$.
\item[(c)] The set of the laws on $\Omega$ that are invariant and reflection positive at $\Pi$ is convex and weakly closed.
\end{enumerate}
\end{lemma}

\begin{proof}
(a) Since $(\theta^*)_*\mu=\mu$, one has $\int\mathrm{d}\mu(F)\,|G(\theta^*F)|^2=\int\mathrm{d}\mu(F)\,|G(F)|^2$ for every Borel function $G$ on $\Omega$.
So the linear map $G\mapsto\theta G$ carries $\mu$-null functions to $\mu$-null functions and preserves the norm of $L^2(\mu)$.
It is its own inverse, because $\theta(\theta G)=(\theta\theta)G=G$;
hence it is unitary.
The inequality is the Cauchy--Schwarz inequality in $L^2(\mu)$, combined with $\|\theta G\|_{L^2(\mu)}=\|G\|_{L^2(\mu)}$.

(b) Let $G\in\mathcal{O}(\{\ell>c\})$, and write $G(F)=G((F_\sigma)_{\sigma\in\Sigma_0})$
with a finite set $\Sigma_0$ of plaquettes contained in $\{\ell>c\}$ and a bounded Borel function $G$ on $\R^{\Sigma_0}$.
Let $\mu_{\Sigma_0}$ be the image of $\mu$ under $F\mapsto(F_\sigma)_{\sigma\in\Sigma_0}$, a finite Borel measure on $\R^{\Sigma_0}$.
The continuous functions with compact support are dense in $L^2(\mu_{\Sigma_0})$ \cite[Theorem~3.14]{Rudin1987};
choose such functions $G_n$ on $\R^{\Sigma_0}$ with $\|G-G_n\|_{L^2(\mu_{\Sigma_0})}\to0$.
Read as functions on $\Omega$, the $G_n$ belong to $\mathcal{O}_{\mathrm{cont}}(\{\ell>c\})$,
and $\|G-G_n\|_{L^2(\mu)}=\|G-G_n\|_{L^2(\mu_{\Sigma_0})}\to0$.
From the identity
\[
\overline{\theta G}\cdot G-\overline{\theta G_n}\cdot G_n=\overline{\theta(G-G_n)}\cdot G+\overline{\theta G_n}\cdot(G-G_n)
\]
and from~(a) one obtains
\[
\left|\mu\left(\overline{\theta G}\cdot G\right)-\mu\left(\overline{\theta G_n}\cdot G_n\right)\right|
\le\|G-G_n\|_{L^2(\mu)}\cdot\left(\|G\|_{L^2(\mu)}+\|G_n\|_{L^2(\mu)}\right)\longrightarrow0 .
\]
By hypothesis $\mu\left(\overline{\theta G_n}\cdot G_n\right)\ge0$ for every $n$; hence $\mu\left(\overline{\theta G}\cdot G\right)\ge0$.
Since $\mathcal{O}^{+}(\ell,c)=\mathcal{O}(\{\ell>c\})$,
this is reflection positivity at $\Pi$ in the sense of Definition~\ref{def:reflection-positivity}.

Now let $G\in\mathcal{O}(\{\ell<c\})$, and put $G':=\overline{\theta G}$.
Since $\ell\circ\theta=2c-\ell$ (Definition~\ref{def:mirrors}\eqref{it:mirror-height}), the reflection $\theta$ maps $\{\ell<c\}$ onto $\{\ell>c\}$,
and~\eqref{eq:layers-local} shows that $G'\in\mathcal{O}(\{\ell>c\})$.
Because $\theta\theta$ is the identity, $\overline{\theta G'}=G$, and therefore $\overline{\theta G}\cdot G=G'\cdot\overline{\theta G'}$.
The first part, applied to $G'$, gives
\[
\mu\left(\overline{\theta G}\cdot G\right)=\mu\left(\overline{\theta G'}\cdot G'\right)\ge0 .
\]

(c) For every bounded Borel function $G$ the map $\mu\mapsto\mu(G)$ is affine on the laws;
so invariance under $\theta$ and reflection positivity at $\Pi$ are preserved by convex combinations, and the set in question is convex.
As the weak topology is metrisable, closedness follows once the weak limit $\mu$ of a sequence of laws $\mu_n$ of the set is shown to belong to the set.
If $G$ is bounded and continuous, then so is $\theta G$, the map $\theta^*$ being a homeomorphism; hence
\[
\mu(\theta G)=\lim_{n\to\infty}\mu_n(\theta G)=\lim_{n\to\infty}\mu_n(G)=\mu(G).
\]
Thus $(\theta^*)_*\mu$ and $\mu$ have the same integrals of all bounded continuous functions,
and they coincide, a Borel probability measure on a metric space being determined by these integrals \cite[Theorem~1.2]{Billingsley1999}:
$\mu$ is invariant under $\theta$.
For $G\in\mathcal{O}_{\mathrm{cont}}(\{\ell>c\})$ the function $\overline{\theta G}\cdot G$ is bounded and continuous, so
\[
\mu\left(\overline{\theta G}\cdot G\right)=\lim_{n\to\infty}\mu_n\left(\overline{\theta G}\cdot G\right)\ge0 .
\]
By~(b), $\mu$ is reflection positive at $\Pi$.
\end{proof}

\begin{lemma}[Transport of a mirror]\label{lem:mirror-transport}
Let $\Pi=(\ell,c)$ and $\Pi'=(\ell',c')$ be mirrors, and let $g\in W$ satisfy $\ell'(gx)-c'=\ell(x)-c$ for all $x\in\R^3$.
Then $\theta_{\Pi'}=g\,\theta_\Pi\,g^{-1}$.
If a law $\mu$ on $\Omega$ is invariant under $g^{-1}$ and reflection positive at $\Pi$, then it is reflection positive at $\Pi'$.
\end{lemma}

\begin{proof}
The hypothesis on $g$ says that $g$ maps the plane $\{\ell=c\}$ onto $\{\ell'=c'\}$ and the half-space $\{\ell>c\}$ onto $\{\ell'>c'\}$.
The isometry $g\,\theta_\Pi\,g^{-1}$ of $\R^3$ has order two, fixes every point of $\{\ell'=c'\}$ and is not the identity.
An isometry of $\R^3$ that fixes a plane pointwise is either the identity or the reflection in this plane;
so $g\,\theta_\Pi\,g^{-1}=\theta_{\Pi'}$.

Let $G'\in\mathcal{O}(\{\ell'>c'\})$ and $G:=g^{-1}G'$.
By~\eqref{eq:layers-local}, $G\in\mathcal{O}(g^{-1}\{\ell'>c'\})=\mathcal{O}(\{\ell>c\})$.
The rule $g(g'G)=(gg')G$ that follows~\eqref{eq:layers-action} gives
\[
g^{-1}(\theta_{\Pi'}G')=(g^{-1}\theta_{\Pi'})G'=(\theta_\Pi g^{-1})G'=\theta_\Pi G .
\]
The action of $g^{-1}$ commutes with products and with complex conjugation.
By the invariance of $\mu$ under $g^{-1}$ and the reflection positivity of $\mu$ at $\Pi$,
\[
\mu\left(\overline{\theta_{\Pi'}G'}\cdot G'\right)=\mu\left(g^{-1}\left(\overline{\theta_{\Pi'}G'}\cdot G'\right)\right)=\mu\left(\overline{\theta_\Pi G}\cdot G\right)\ge0 .
\]
\end{proof}

\begin{lemma}[Laws with a diagonal mirror]\label{lem:diagonal-laws}
Let $c_0\in\Z$, and let $\mathcal{R}(c_0)$ be the set of the laws on $\Omega$ that satisfy~\eqref{cond:gaussian-domination}
and are invariant and reflection positive at the mirror $(\ell_-,c_0)$ (Definition~\ref{def:reflection-positivity}).
Then $\mathcal{R}(c_0)$ is convex and weakly compact.
\end{lemma}

\begin{proof}
The set $\mathcal{R}(c_0)$ is the intersection of $\mathcal{M}_{(\mathrm{G})}$, which is convex and weakly compact by Lemma~\ref{lem:domination-compactness}(d),
with the set of the laws that are invariant and reflection positive at $(\ell_-,c_0)$, which is convex and weakly closed by Lemma~\ref{lem:law-mirror}(c).
An intersection of two convex sets is convex, and a closed subset of a compact set is compact.
\end{proof}

\subsection{Data in other frames}\label{sec:transport-of-data}

The data of Definition~\ref{def:data}, the standard data, are built on the diagonal frame $\mathfrak{F}_{\mathrm{d}}$.
The elements of the cubic group $\Ocub$ carry them to data built on the images of this frame.

\begin{definition}[Image data]\label{def:image-data}
Let $g\in\Ocub$.
For a family $\ell\in\mathcal{L}$ the linear form $\ell\circ g^{-1}$ is $x\mapsto(gn_\ell)\cdot x$, since $g^{-1}$ is the transpose of $g$.
The signed permutation matrix $g$ maps each of the vectors $\pm e_\mu$ and $\pm e_\mu\pm e_\nu$, $\mu\neq\nu$, onto a vector of the same list,
and each vector of this list is $\pm n_{\ell'}$ for exactly one family $\ell'$ (Definition~\ref{def:mirrors}).
Hence there are exactly one family $\ell'$ and one sign such that
\[
\ell\circ g^{-1}=\pm\ell' .
\]
As $|n_{\ell'}|=|gn_\ell|=|n_\ell|$, the family $\ell'$ is an axis family if and only if $\ell$ is one,
so that $\Adm(\ell')=\Adm(\ell)$, a set that is symmetric under $c\mapsto-c$.
For $c\in\Adm(\ell)$ the image under $g$ of the plane $\{\ell=c\}$ is the plane $\left\{\ell\circ g^{-1}=c\right\}=\{\ell'=\pm c\}$, with the same sign.
The normals of two different families are not parallel, so that a mirror is determined by its plane.
The mirror
\[
g(\ell,c):=(\ell',\pm c),
\]
again with this sign, is the \emph{image} of the mirror $(\ell,c)$ under $g$;
for every mirror $\Pi$ and every $g'\in\Ocub$ one has $\left(gg'\right)\Pi=g\left(g'\Pi\right)$,
both sides being the mirror whose plane is the image under $gg'$ of the plane of $\Pi$.

Let $\mathfrak{F}=\left(\ell^1,\ell^2,\ell^3\right)$ be a frame,
and for $a=1,2,3$ let $\ell'^a$ be the family with $\ell^a\circ g^{-1}=\pm\ell'^a$.
Then $gn_{\ell^a}=\pm n_{\ell'^a}$ with the same sign, and $g$ is orthogonal;
so the normals of $\ell'^1$, $\ell'^2$, $\ell'^3$ are pairwise orthogonal, and the \emph{image} of $\mathfrak{F}$ under $g$,
\[
g\mathfrak{F}:=\left(\ell'^1,\ell'^2,\ell'^3\right),
\]
is a frame (Definition~\ref{def:cell-structures}), with $|n_{\ell'^a}|=|n_{\ell^a}|$ for $a=1,2,3$.
Let $\mathcal{Q}=(\mathfrak{F};o_1,o_2,o_3;h)$ be a cell structure.
The image of the mirror $\left(\ell^a,o_a\right)$ is a mirror $\left(\ell'^a,o'_a\right)$,
where $o'_a=o_a$ if $\ell^a\circ g^{-1}=\ell'^a$ and $o'_a=-o_a$ if $\ell^a\circ g^{-1}=-\ell'^a$.
The \emph{image} of $\mathcal{Q}$ under $g$ is the cell structure
\[
g\mathcal{Q}:=\left(g\mathfrak{F};o'_1,o'_2,o'_3;h\right),
\]
which has the width $h$ and the same cell volume as $\mathcal{Q}$.
The image of the wall $\left(\ell^a,o_a+j\cdot h\right)$, $j\in\Z$, is the wall $\left(\ell'^a,o'_a+j\cdot h\right)$ of $g\mathcal{Q}$ in the first case
and the wall $\left(\ell'^a,o'_a-j\cdot h\right)$ in the second;
so the walls of $g\mathcal{Q}$ are the images under $g$ of the walls of $\mathcal{Q}$.
An open cell $Q_\kappa$ of $\mathcal{Q}$ has the image
\[
gQ_\kappa=\left\{x\in\R^3:\ o_a+\kappa_a\cdot h<\ell^a\left(g^{-1}x\right)<o_a+(\kappa_a+1)\cdot h,\ a=1,2,3\right\}.
\]
In the first case the condition on the index $a$ reads $o'_a+\kappa_a\cdot h<\ell'^a(x)<o'_a+(\kappa_a+1)\cdot h$,
in the second case $o'_a-(\kappa_a+1)\cdot h<\ell'^a(x)<o'_a-\kappa_a\cdot h$.
Hence $gQ_\kappa$ is the open cell of $g\mathcal{Q}$ with the index $\kappa'$,
where $\kappa'_a=\kappa_a$ in the first case and $\kappa'_a=-\kappa_a-1$ in the second,
and the open cells of $g\mathcal{Q}$ are the sets $gQ_\kappa$, $\kappa\in\Z^3$.
The reflection in the image under $g$ of a wall is the conjugate by $g$ of the reflection in that wall; therefore
\[
W(g\mathcal{Q})=g\,W(\mathcal{Q})\,g^{-1}.
\]
Thus the walls, the open cells and the group of the walls of $g\mathcal{Q}$ are determined by those of $\mathcal{Q}$ and by $g$;
the two cases make a difference only for the positions and the indices.
For $g'\in\Ocub$ the rule $\left(gg'\right)\Pi=g\left(g'\Pi\right)$ for mirrors gives
$\left(gg'\right)\mathfrak{F}=g\left(g'\mathfrak{F}\right)$ and $\left(gg'\right)\mathcal{Q}=g\left(g'\mathcal{Q}\right)$.
Finally, let $\mathbb{T}(\mathfrak{F};p_1,p_2,p_3)=(\Gamma,\chi)$ be a rectangular torus (Definition~\ref{def:cell-structures}).
The conditions imposed there on $p_1,p_2,p_3$ hold for $g\mathfrak{F}$ as well, since $|n_{\ell'^a}|=|n_{\ell^a}|$ and $g\Z^3=\Z^3$.
The matrix $g$ maps the generator $\frac{p_a}{|n_{\ell^a}|^2}\cdot n_{\ell^a}$ of $\Gamma$ onto one of the vectors $\pm\frac{p_a}{|n_{\ell'^a}|^2}\cdot n_{\ell'^a}$,
and the character $\chi\circ g^{-1}$ of the image torus (Definition~\ref{def:tori}) takes the value $-1$ at both of them,
because $\chi$ takes the value $-1$ at both vectors $\pm\frac{p_a}{|n_{\ell^a}|^2}\cdot n_{\ell^a}$.
Therefore
\[
g\,\mathbb{T}(\mathfrak{F};p_1,p_2,p_3)=\mathbb{T}(g\mathfrak{F};p_1,p_2,p_3),
\]
and $g\mathbb{T}$ carries $g\mathcal{Q}$ whenever a torus $\mathbb{T}$ carries $\mathcal{Q}$.

Now let the sequence of sides $\mathfrak{n}=(n_i)$ satisfy \eqref{cond:divisibility},
and let $\mathfrak{d}=(c_0,c_+,c_3;h;O)$ be standard data,
with the cell structure $\mathcal{Q}_{\mathfrak{d}}$, its open cells $Q_\kappa$, $\kappa\in\Z^3$, and the folded array $O_{\mathfrak{d}}$ of Definition~\ref{def:data}.
The \emph{image data} $g\mathfrak{d}$ are given by the following objects.
\begin{enumerate}
\item[(a)]
The cell structure $g\mathcal{Q}_{\mathfrak{d}}$.
Its frame is $g\mathfrak{F}_{\mathrm{d}}$, its width is $h$,
and its positions are $c_0$, $c_+$, $c_3$ or their negatives,
according as $\ell_-\circ g^{-1}$, $\ell_+\circ g^{-1}$, $x_3\circ g^{-1}$ is a family or the negative of a family.
Its walls are the images under $g$ of the walls of $\mathcal{Q}_{\mathfrak{d}}$,
its open cells are the sets $gQ_\kappa$, $\kappa\in\Z^3$,
and $W\left(g\mathcal{Q}_{\mathfrak{d}}\right)=g\,W\left(\mathcal{Q}_{\mathfrak{d}}\right)g^{-1}$.
\item[(b)]
The observable $gO$, that is, the function $F\mapsto O\left(g^{*}F\right)$ of \eqref{eq:layers-action}.
If $O=\sum_ka_k\cdot\varphi_{\beta_k}$, then
\[
gO=\sum_ka_k\cdot\varphi_{\left(g^{-1}\right)^{*}\beta_k},
\]
because $g^{*}$ is orthogonal with the inverse $\left(g^{-1}\right)^{*}$ (Definition~\ref{def:cochains}),
so that $\left\langle\beta_k,g^{*}F\right\rangle=\left\langle\left(g^{-1}\right)^{*}\beta_k,F\right\rangle$.
Hence $gO$ is a real trigonometric polynomial with the same coefficients as $O$, and
\[
\supp(gO)=g\supp O\subset gQ_0,\qquad|gO|_1=|O|_1,\qquad r_{gO}=r_O;
\]
indeed, $\left(\left(g^{-1}\right)^{*}\beta\right)(\sigma)=\pm\beta\left(g^{-1}\sigma\right)$ for every plaquette $\sigma$,
and the signed permutation matrices $g$ and $g^{-1}$ map every cube $y+[-r,r]^3$ onto a cube of the same radius.
\item[(c)]
The torus $gD(n_i)=\left(g\Gamma,\chi\circ g^{-1}\right)$, where $D(n_i)=(\Gamma,\chi)$ (Definition~\ref{def:tori}).
\item[(d)]
The array
\[
O_{g\mathfrak{d}}:=\sum_{\kappa\in\Z^3}(gO)^{[g\kappa]},
\]
a formal sum with one term in each open cell, where $(gO)^{[g\kappa]}$ denotes the mirror copy of $gO$ in the cell $gQ_\kappa$ of $g\mathcal{Q}_{\mathfrak{d}}$
(Definitions~\ref{def:folds} and~\ref{def:arrays});
the mirror copies are defined, since $gO$ is bounded and depends only on finitely many $F_\sigma$, all with $\sigma\subset gQ_0$, so that $gO\in\mathcal{A}(gQ_0)$.
Let $\mathbb{T}'=(\Gamma',\chi')$ be a torus that carries $g\mathcal{Q}_{\mathfrak{d}}$.
On $\mathbb{T}'$ the array is the function
\[
\sum_{\kappa}(gO)^{[g\kappa]}\left(\widetilde F\right)
\]
on $\Conf(\Gamma')$, the sum extending over a set of indices $\kappa$ for which the cells $gQ_\kappa$ represent the cells of $\mathbb{T}'$.
This function does not depend on the choice of the representatives.
Indeed, for $v\in\Gamma'$ the translation $\tau_v$ belongs to $W\left(g\mathcal{Q}_{\mathfrak{d}}\right)$ (Definition~\ref{def:cell-structures}),
so that the mirror copy in the cell $\tau_vgQ_\kappa$ is the image under $\tau_v$ of the mirror copy in $gQ_\kappa$ (Definition~\ref{def:folds});
the two take the same value at $\widetilde F$, which is $\Gamma'$-periodic.
The sum has $|\mathbb{T}'|/|Q_0|$ terms, one for each cell of $\mathbb{T}'$.
Each of them is a mirror copy $(gO)^{g''}$ with $g''\in W$, which equals $g''(gO)$ by \eqref{eq:setting-group-action}, the function $gO$ being real;
it is therefore a real trigonometric polynomial with the coefficients of $O$ (Definition~\ref{def:characters}),
and its modulus is at most $|gO|_1=|O|_1$.
\item[(e)]
The perturbed functional.
Let $i$ be an index for which $D(n_i)$ carries $\mathcal{Q}_{\mathfrak{d}}$;
as $D(n_i)$ is the rectangular torus $\mathbb{T}(\mathfrak{F}_{\mathrm{d}};n_i,2n_i,n_i)$ (Definitions~\ref{def:cell-structures} and~\ref{def:sides}),
the torus $gD(n_i)$ carries $g\mathcal{Q}_{\mathfrak{d}}$.
Let $t\ge0$.
For $G\in\mathcal{A}(\R^3)$ that embeds in $gD(n_i)$,
\[
\mu^t_{i,g\mathfrak{d}}(G):=\frac{Z_{gD(n_i)}\left[\exp\left(t\cdot O_{g\mathfrak{d}}\right)\cdot G\right]}{Z_{gD(n_i)}\left[\exp\left(t\cdot O_{g\mathfrak{d}}\right)\right]}.
\]
As $g\Z^3=\Z^3$, the torus $gD(n_i)$ has as many photon sites as $D(n_i)$, namely $n_i^3$, and hence $n_i^3/|Q_0|$ cells.
By (d), the array $O_{g\mathfrak{d}}$ is on $gD(n_i)$ a real function of modulus at most $\frac{n_i^3}{|Q_0|}\cdot|O|_1$,
and Lemma~\ref{lem:perturbed-denominators} together with Theorem~\ref{thm:partition-positivity} gives
\[
Z_{gD(n_i)}\left[\exp\left(t\cdot O_{g\mathfrak{d}}\right)\right]\ge\exp\left(-t\cdot|O|_1\cdot\frac{n_i^3}{|Q_0|}\right)\cdot Z_{gD(n_i)}[1]>0.
\]
The determinant $\det D^{gD(n_i)}_A$ is non-negative (Proposition~\ref{prop:accretivity}).
Hence, as for standard data, on bosonic $G$ the functional $\mu^t_{i,g\mathfrak{d}}$ is the expectation of a probability measure on $\Conf(g\Gamma)$;
the law of $\widetilde F$ under this measure is a Borel probability measure on $\Omega$, again denoted $\mu^t_{i,g\mathfrak{d}}$,
and $\mu^t_{i,g\mathfrak{d}}(G)=\int\mathrm{d}\mu^t_{i,g\mathfrak{d}}(F)\,G(F)$ for every bosonic $G$ as above.
\end{enumerate}
The cubic torus is the same for $\mathfrak{d}$ and for $g\mathfrak{d}$: $gC(n_i)=C(n_i)$,
because the signed permutation matrices map the lattice $n_i\Z^3$ onto itself and permute the six vectors $\pm n_ie_\mu$,
at each of which the character of $C(n_i)$ takes the value $-1$.

Throughout this subsection, the word \emph{data} covers both standard data and image data.
Let $\mathfrak{d}=g'\mathfrak{d}'$ be image data, with standard data $\mathfrak{d}'$ and $g'\in\Ocub$, and let $g\in\Ocub$; one puts
\[
g\mathfrak{d}:=\left(gg'\right)\mathfrak{d}'.
\]
The cell structure, the observable and the torus of $g\mathfrak{d}$ arise from those of $\mathfrak{d}$ by an application of $g$:
$\left(gg'\right)\mathcal{Q}_{\mathfrak{d}'}=g\left(g'\mathcal{Q}_{\mathfrak{d}'}\right)$, as for every cell structure,
$\left(gg'\right)O'=g\left(g'O'\right)$ for the observable $O'$ of $\mathfrak{d}'$ by \eqref{eq:layers-action},
and $\left(gg'\right)D(n_i)=g\left(g'D(n_i)\right)$ by Definition~\ref{def:tori}.
The definitions (d) and (e) of the array and of the functional involve nothing but these three objects.
Two ways of writing $\mathfrak{d}$ as an image of standard data therefore lead to the same $g\mathfrak{d}$.
Moreover $1\mathfrak{d}=\mathfrak{d}$ and $g\left(g''\mathfrak{d}\right)=\left(gg''\right)\mathfrak{d}$ for all data $\mathfrak{d}$ and all $g,g''\in\Ocub$.
For data $\mathfrak{d}$ the symbols $\mathcal{Q}_{\mathfrak{d}}$, $\mathbb{T}_{i,\mathfrak{d}}$, $O_{\mathfrak{d}}$ and $\mu^t_{i,\mathfrak{d}}$ denote its cell structure, its torus, its array and its functional;
the observable of standard data $(c_0,c_+,c_3;h;O)$ is $O$, and the observable of $g\mathfrak{d}$ is the image under $g$ of the observable of $\mathfrak{d}$.
Thus $\mathbb{T}_{i,\mathfrak{d}}=D(n_i)$ for standard data, and $\mathbb{T}_{i,g\mathfrak{d}}=g\,\mathbb{T}_{i,\mathfrak{d}}$ for all data.
The open cells of $\mathcal{Q}_{\mathfrak{d}}$ are written $Q_\kappa$, $\kappa\in\Z^3$:
for standard data with the indices of Definition~\ref{def:cell-structures},
and for image data with the indices carried along by $g$, the cell of $g\mathcal{Q}_{\mathfrak{d}}$ written $Q_\kappa$ being $gQ_\kappa$,
whatever its index as an open cell of the cell structure $g\mathcal{Q}_{\mathfrak{d}}$.
So $Q_0$ always contains the support of the observable.
Comparison states for data $\mathfrak{d}$ are defined by Definition~\ref{def:comparison-states}, with the functionals $\mu^t_{i,\mathfrak{d}}$ in place of $\mu^t_i$.
\end{definition}

\emph{Three properties of a datum and a direction.}
Let $\mathfrak{d}'$ be data and $\mu\in\{1,2,3\}$,
and let $\Theta_{\mu,c}$ and $\mathcal{P}_1(\mu,c)$ be as in Definition~\ref{def:other-directions}.
The following properties of the pair $(\mathfrak{d}',\mu)$ concern the single datum $\mathfrak{d}'$
and all the coordinate link planes $(x_\mu,c)$, $c\in\Z+\tfrac12$.
The first is the limit relation of Theorem~\ref{thm:two-layer-principle}(b), formulated for the datum $\mathfrak{d}'$:
\begin{equation}\label{it:datum-hermitian}
\lim_{t\to0}\,\limsup_{i\to\infty}\,\left|\mu^t_{i,\mathfrak{d}'}\left(\Theta_{\mu,c}Y\right)-\overline{\mu^t_{i,\mathfrak{d}'}(Y)}\right|=0
\qquad\text{for all }c\in\Z+\tfrac12\text{ and all }Y\in\mathcal{P}_1(\mu,c).
\end{equation}
The second is the reduction of Theorem~\ref{thm:two-layer-reduction}, formulated for the datum $\mathfrak{d}'$ alone:
for every $c\in\Z+\tfrac12$ and every $P\in\mathcal{P}$ there are $Y_P\in\mathcal{P}_1(\mu,c)$, a constant $C<\infty$ and an index $i_0$,
none of which depends on $i$ or $t$, such that
\begin{equation}\label{it:datum-reduction}
\begin{aligned}
&\left|\mu^t_{i,\mathfrak{d}'}(P)-\mu^t_{i,\mathfrak{d}'}\left(Y_P\right)\right|\le C\cdot t\quad\text{and}\\
&\left|\mu^t_{i,\mathfrak{d}'}\left(\Theta_{\mu,c}P\right)-\mu^t_{i,\mathfrak{d}'}\left(\Theta_{\mu,c}Y_P\right)\right|\le C\cdot t
\qquad\text{for }0\le t\le1\text{ and }i\ge i_0.
\end{aligned}
\end{equation}
The third is the conclusion of Theorem~\ref{thm:no-layers} for the datum $\mathfrak{d}'$ and the direction $\mu$:
every comparison state $\nu_0$ for $\mathfrak{d}'$ is invariant under the reflections in the planes $(x_\mu,c)$, that is,
\begin{equation}\label{it:datum-invariance}
\left(\theta^{*}_{x_\mu,c}\right)_{*}\nu_0=\nu_0\qquad\text{for all }c\in\Z+\tfrac12.
\end{equation}

\begin{lemma}[Transport of data]\label{lem:data-transport}
Let $\mathfrak{d}$ be data, $g\in\Ocub$ and $t\ge0$,
put $\mathbb{T}_i:=\mathbb{T}_{i,\mathfrak{d}}$,
and let $i$ be so large that $\mathbb{T}_i$ is a torus that carries $\mathcal{Q}_{\mathfrak{d}}$.
For a mirror $\Pi$, $g\Pi$ denotes its image under $g$ (Definition~\ref{def:image-data}).
\begin{enumerate}
\item[(i)]
For every mirror $\Pi$ and every even integrand $G$,
\[
\left(\Theta_\Pi G\right)^g=\Theta_{g\Pi}\left(G^g\right).
\]
\item[(ii)]
$O_{g\mathfrak{d}}=\left(O_{\mathfrak{d}}\right)^g=g\,O_{\mathfrak{d}}$, term by term.
\item[(iii)]
Let $G\in\mathcal{A}(\R^3)$ embed in $\mathbb{T}_i$.
Then $G^g$ embeds in $g\mathbb{T}_i$, and
\[
\mu^t_{i,g\mathfrak{d}}\left(G^g\right)=
\begin{cases}
\mu^t_{i,\mathfrak{d}}(G)&\text{if }\det g=1,\\
\overline{\mu^t_{i,\mathfrak{d}}(G)}&\text{if }\det g=-1.
\end{cases}
\]
\item[(iv)]
The law $\mu^t_{i,\mathfrak{d}}$ is the image of the law $\mu^t_{i,g\mathfrak{d}}$ under $g^{*}$:
\[
\left(g^{*}\right)_{*}\mu^t_{i,g\mathfrak{d}}=\mu^t_{i,\mathfrak{d}}.
\]
A law $\nu$ on $\Omega$ is a comparison state for $g\mathfrak{d}$, obtained with a subsequence $(i_m)$ and numbers $t_k$,
if and only if $\left(g^{*}\right)_{*}\nu$ is a comparison state for $\mathfrak{d}$ obtained with the same $(i_m)$ and the same $t_k$.
\item[(v)]
Let $\mu,\mu'\in\{1,2,3\}$ be such that $ge_\mu=\pm e_{\mu'}$; below, $\pm$ is the sign in this relation.
\begin{enumerate}
\item[(a)]
For $c\in\Z+\tfrac12$ the element $g$ maps the mirror $(x_\mu,c)$ onto the mirror $(x_{\mu'},\pm c)$.
The map $G\mapsto G^g$ maps $\mathcal{P}$ onto $\mathcal{P}$ and $\mathcal{P}_1(\mu,c)$ onto $\mathcal{P}_1(\mu',\pm c)$.
In particular
\[
\mathcal{P}_1(\mu,c)=\left\{Y^{g''}:\ Y\in\mathcal{P}_1(c)\right\}\qquad\text{for every }g''\in\Ocub\text{ with }g''e_1=e_\mu .
\]
\item[(b)]
The pair $(\mathfrak{d},\mu)$ has the property \eqref{it:datum-hermitian} if and only if the pair $(g\mathfrak{d},\mu')$ has it;
the same holds for the property \eqref{it:datum-reduction} and for the property \eqref{it:datum-invariance}.
\end{enumerate}
\end{enumerate}
\end{lemma}

\begin{proof}
Item (i) rests on the composition rule for the action of $W$.
By \eqref{eq:setting-group-action}, $\Theta_\Pi G=G^{\theta_\Pi}$ for every even integrand $G$,
and $\left(G^{g'}\right)^{g''}=G^{g''g'}$ for all $g',g''\in W$.
The isometry $g\theta_\Pi g^{-1}$ has order two, fixes the plane of $g\Pi$ pointwise and is not the identity;
it is therefore the reflection $\theta_{g\Pi}$, and $g\theta_\Pi=\theta_{g\Pi}g$.
Hence
\[
\left(\Theta_\Pi G\right)^g=G^{g\theta_\Pi}=G^{\theta_{g\Pi}g}=\left(G^g\right)^{\theta_{g\Pi}}=\Theta_{g\Pi}\left(G^g\right).
\]

For (ii), let $O$ be the observable of $\mathfrak{d}$,
and for $\kappa\in\Z^3$ let $g_\kappa$ be the element of $W\left(\mathcal{Q}_{\mathfrak{d}}\right)$ with $g_\kappa Q_0=Q_\kappa$.
By Definition~\ref{def:folds} the term of $O_{\mathfrak{d}}$ in the cell $Q_\kappa$ is $O^{[\kappa]}=O^{g_\kappa}$,
and the composition rule gives
\[
\left(O^{g_\kappa}\right)^g=O^{gg_\kappa}=\left(O^g\right)^{gg_\kappa g^{-1}}.
\]
Here $gg_\kappa g^{-1}$ is the element of $W\left(g\mathcal{Q}_{\mathfrak{d}}\right)=g\,W\left(\mathcal{Q}_{\mathfrak{d}}\right)g^{-1}$ that maps $gQ_0$ onto $gQ_\kappa$;
the group identity holds for all data, the reflections in the walls of $g\mathcal{Q}_{\mathfrak{d}}$ being the conjugates by $g$ of the reflections in the walls of $\mathcal{Q}_{\mathfrak{d}}$.
The observable is real, so that $O^g=gO$ for either parity of $g$, by \eqref{eq:setting-group-action}.
Therefore $\left(O^{[\kappa]}\right)^g$ is the mirror copy of $gO$ in the cell $gQ_\kappa$, which is the term of $O_{g\mathfrak{d}}$ in this cell.
All terms are real functions of $F$,
and for a real function $\Phi$ of $F$ one has $\Phi^g=g\Phi$ by \eqref{eq:setting-group-action} and \eqref{eq:layers-action};
this gives the second equality, again term by term.

We turn to (iii).
Let $\Gamma$ be the period lattice of $\mathbb{T}_i$; then $g\Gamma$ is the period lattice of $g\mathbb{T}_i$.
The signed permutation matrix $g$ maps a box onto a box of the same radius, $g\left(y+[-r,r]^3\right)=gy+[-r,r]^3$,
and $g\left(B+[-1,1]^3+v\right)=gB+[-1,1]^3+gv$ for every box $B$ and every $v\in\Gamma$.
Hence boxes that embed disjointly in $\Gamma$ are mapped onto boxes that embed disjointly in $g\Gamma$ (Definition~\ref{def:integrands}).
By Lemma~\ref{lem:reflection-classes}, every $\Theta_\Pi$ maps $\mathcal{G}^{\mathrm{ev}}(B)$ onto $\mathcal{G}^{\mathrm{ev}}(\theta_\Pi B)$,
and a bounded Borel function of finitely many $F_\sigma$ onto a function of the same kind.
Since $g$ is a product of reflections and the factors of an admissible product are even, hence commute,
the map $G\mapsto G^g$ sends a product $\Phi\cdot G_1\cdots G_r$ with the boxes $B_1,\dots,B_r$
to the product $\Phi^g\cdot G_1^g\cdots G_r^g$ with the boxes $gB_1,\dots,gB_r$.
Consequently $G^g$ lies in $\mathcal{A}(\R^3)$ and embeds in $g\mathbb{T}_i$.

By Definition~\ref{def:arrays} for standard data, and by Definition~\ref{def:image-data}(d) for image data,
the factor $\exp\left(t\cdot O_{\mathfrak{d}}\right)$ is on $\mathbb{T}_i$ the function $\Phi_t(\widetilde F)$, where
\[
\Phi_t:=\exp\left(t\cdot\sum_{\kappa}O^{[\kappa]}\right)
\]
and $\kappa$ ranges over representatives of the cells of $\mathbb{T}_i$, one for each cell.
Only finitely many trigonometric polynomials enter this sum,
so that $\Phi_t$ is a bounded continuous function of finitely many plaquette variables.
Definition~\ref{def:integrands} admits such a function as the bounded factor of a product, whatever plaquettes it reads.
Multiplying the bounded factors of the products that make up $G$ by $\Phi_t$, one sees that $\exp\left(t\cdot O_{\mathfrak{d}}\right)\cdot G$ belongs to $\mathcal{A}(\R^3)$
and embeds in $\mathbb{T}_i$, with the boxes of $G$.
The map $G\mapsto G^g$ is a linear automorphism if $\det g=1$ and an antilinear anti-automorphism if $\det g=-1$ \eqref{eq:setting-group-action}.
In the first case $\left(\Phi_t\cdot G\right)^g=\Phi_t^g\cdot G^g$.
In the second case $\left(\Phi_t\cdot G\right)^g=G^g\cdot\Phi_t^g$, and the two factors commute, $\Phi_t^g$ being a scalar function of $F$.
As $\Phi_t$ is real, $\Phi_t^g=g\Phi_t$ in both cases, and (ii), applied term by term, gives
\[
g\Phi_t=\exp\left(t\cdot\sum_{\kappa}(gO)^{[g\kappa]}\right)
\]
with the same representatives $\kappa$.
The cells $gQ_\kappa$ represent the cells of $g\mathbb{T}_i$, so that, by Definition~\ref{def:image-data}(d), $g\Phi_t(\widetilde F)$ is the factor $\exp\left(t\cdot O_{g\mathfrak{d}}\right)$ on $g\mathbb{T}_i$.
In both cases, therefore,
\[
\left(\exp\left(t\cdot O_{\mathfrak{d}}\right)\cdot G\right)^g=\exp\left(t\cdot O_{g\mathfrak{d}}\right)\cdot G^g.
\]
Proposition~\ref{prop:symmetries}\eqref{it:functional-covariance} holds on every torus.
Applied on $\mathbb{T}_i$ to $\exp\left(t\cdot O_{\mathfrak{d}}\right)\cdot G$, it yields
\[
Z_{g\mathbb{T}_i}\left[\exp\left(t\cdot O_{g\mathfrak{d}}\right)\cdot G^g\right]=
\begin{cases}
Z_{\mathbb{T}_i}\left[\exp\left(t\cdot O_{\mathfrak{d}}\right)\cdot G\right]&\text{if }\det g=1,\\
\overline{Z_{\mathbb{T}_i}\left[\exp\left(t\cdot O_{\mathfrak{d}}\right)\cdot G\right]}&\text{if }\det g=-1.
\end{cases}
\]
For $G=1$ the right side is, in both cases, the denominator of $\mu^t_{i,\mathfrak{d}}$, a positive number;
it is therefore also the denominator of $\mu^t_{i,g\mathfrak{d}}$, and division gives (iii).

Consider now (iv), and let $\Phi\in\mathcal{O}(\R^3)$.
As a bounded factor without boxes, $\Phi$ belongs to $\mathcal{A}(\R^3)$ and embeds in every torus (Definition~\ref{def:integrands});
by \eqref{eq:setting-group-action}, $\Phi^g=g\Phi$ if $\det g=1$ and $\Phi^g=\overline{g\Phi}$ if $\det g=-1$.
On such functions the functionals $\mu^t_{i,\mathfrak{d}}$ and $\mu^t_{i,g\mathfrak{d}}$ are the integrals against the laws of the same names
(Definition~\ref{def:data}, Definition~\ref{def:image-data}(e), and \eqref{eq:layers-bosonic}).
If $\det g=1$, item (iii) reads $\mu^t_{i,g\mathfrak{d}}(g\Phi)=\mu^t_{i,\mathfrak{d}}(\Phi)$.
If $\det g=-1$, it reads $\mu^t_{i,g\mathfrak{d}}\left(\overline{g\Phi}\right)=\overline{\mu^t_{i,\mathfrak{d}}(\Phi)}$,
and complex conjugation, which commutes with the integral against a law, gives the same identity.
Hence
\[
\int\mathrm{d}\mu^t_{i,g\mathfrak{d}}(F)\,\Phi\left(g^{*}F\right)=\int\mathrm{d}\mu^t_{i,\mathfrak{d}}(F)\,\Phi(F)
\qquad\text{for all }\Phi\in\mathcal{O}(\R^3).
\]
The left side is the integral of $\Phi$ against $\left(g^{*}\right)_{*}\mu^t_{i,g\mathfrak{d}}$.
Every real trigonometric polynomial belongs to $\mathcal{O}(\R^3)$,
so the two laws coincide by Lemma~\ref{lem:characters-determine}(a).

The map $g^{*}$ is a homeomorphism of $\Omega$ with the inverse $\left(g^{-1}\right)^{*}$:
it permutes the coordinates $F_\sigma$ up to signs and commutes with $d$ (Definition~\ref{def:cochains}).
The composition of a bounded continuous function on $\Omega$ with $g^{*}$ or with $\left(g^{-1}\right)^{*}$ is again bounded and continuous;
hence a sequence of laws converges weakly to a law if and only if the images of its members under $g^{*}$ converge weakly to the image of that law.
Let $\nu$ be a comparison state for $g\mathfrak{d}$,
obtained with the subsequence $(i_m)$, the numbers $t_k$ and the laws $\nu_{t_k}$ of Definition~\ref{def:comparison-states}.
By the first assertion of (iv), $\mu^{t_k}_{i_m,\mathfrak{d}}=\left(g^{*}\right)_{*}\mu^{t_k}_{i_m,g\mathfrak{d}}$ converges weakly to $\left(g^{*}\right)_{*}\nu_{t_k}$ as $m\to\infty$, for every $k$,
and $\left(g^{*}\right)_{*}\nu_{t_k}$ converges weakly to $\left(g^{*}\right)_{*}\nu$ as $k\to\infty$.
So $\left(g^{*}\right)_{*}\nu$ is a comparison state for $\mathfrak{d}$, obtained with the same $(i_m)$ and $t_k$.
Conversely, let $\left(g^{*}\right)_{*}\nu$ be a comparison state for $\mathfrak{d}$, obtained with $(i_m)$, $t_k$ and laws $\nu'_{t_k}$.
The map $\left(\left(g^{-1}\right)^{*}\right)_{*}$ is inverse to $\left(g^{*}\right)_{*}$,
because $\left(g^{-1}\right)^{*}g^{*}=\left(gg^{-1}\right)^{*}$ is the identity (Definition~\ref{def:cochains}).
It carries $\mu^{t_k}_{i_m,\mathfrak{d}}$ to $\mu^{t_k}_{i_m,g\mathfrak{d}}$ and $\left(g^{*}\right)_{*}\nu$ to $\nu$, and it preserves weak convergence.
Hence $\mu^{t_k}_{i_m,g\mathfrak{d}}$ converges weakly to $\left(\left(g^{-1}\right)^{*}\right)_{*}\nu'_{t_k}$ as $m\to\infty$,
and these laws converge weakly to $\nu$ as $k\to\infty$:
$\nu$ is a comparison state for $g\mathfrak{d}$, obtained with the same $(i_m)$ and $t_k$.

It remains to prove (v).
Recall that $g_\mu$ is the cyclic permutation matrix with $g_\mu e_1=e_\mu$,
so that $\mathcal{P}_1(\mu,c)=\left\{Y^{g_\mu}:Y\in\mathcal{P}_1(c)\right\}$ (Definition~\ref{def:other-directions}).
Let $\pm$ be the sign in $ge_\mu=\pm e_{\mu'}$.
Since $g$ is orthogonal and $g^{-1}e_{\mu'}=\pm e_\mu$, one has $(gx)_{\mu'}=\pm x_\mu$ for $x\in\R^3$;
therefore $g\{x_\mu=c\}=\{x_{\mu'}=\pm c\}$, which is the first assertion of (a).
The class $\mathcal{P}$ is stable under every map $G\mapsto G^{g'}$, $g'\in W$ (Definition~\ref{def:polynomials}),
and $\left(G^g\right)^{g^{-1}}=G$; so $G\mapsto G^g$ maps $\mathcal{P}$ onto $\mathcal{P}$.

\emph{Images of generators.}
Let $\Pi$ be a mirror and $\exp\left(i\langle\delta,\hat A\rangle\right)\cdot\Psi$ a generator of $\mathcal{P}$.
The pull-back $\theta_\Pi^{*}$ is orthogonal and an involution on one-cochains,
so that $\left\langle\delta,\theta_\Pi^{*}\hat A\right\rangle=\left\langle\theta_\Pi^{*}\delta,\hat A\right\rangle$,
exactly as for two-cochains in \eqref{eq:layers-characters}.
With \eqref{eq:setting-reflection-coefficients} and \eqref{eq:setting-reflection-fields} it follows that
\[
\Theta_\Pi\left(\exp\left(i\langle\delta,\hat A\rangle\right)\cdot\Psi\right)=\exp\left(-i\left\langle\theta_\Pi^{*}\delta,\hat A\right\rangle\right)\cdot\Theta_\Pi(\Psi).
\]
On the right, the value $\left(\theta_\Pi^{*}\delta\right)(b)$ at a link $b$ is $\delta(\theta_\Pi b)$ up to an orientation sign,
so that $\supp\theta_\Pi^{*}\delta=\theta_\Pi(\supp\delta)$.
Moreover \eqref{eq:setting-reflection-fields} replaces a generator at a site $x$ by a combination of generators at the site $\theta_\Pi x$;
as $\Theta_\Pi$ is antilinear and reverses the order of the factors,
$\Theta_\Pi(\Psi)$ expands into monomials whose generators sit at the mirror images of the sites of $\Psi$.
Let $g'=\theta_1\cdots\theta_r\in W$ be written as a product of reflections.
Applying the displayed formula $r$ times, as \eqref{eq:setting-group-action} prescribes, term by term and with the antilinearity of the reflections,
and using $\theta_1^{*}\cdots\theta_r^{*}=\left(\theta_r\cdots\theta_1\right)^{*}=\left(g'^{-1}\right)^{*}$, one obtains
\[
\left(\exp\left(i\langle\delta,\hat A\rangle\right)\cdot\Psi\right)^{g'}=\exp\left(\pm i\left\langle\left(g'^{-1}\right)^{*}\delta,\hat A\right\rangle\right)\cdot\Psi^{g'};
\]
in this formula and up to the end of the paragraph, $\pm$ is $+$ if $g'$ is even and $-$ if $g'$ is odd.
The value of $\left(g'^{-1}\right)^{*}\delta$ at a link $b$ is $\delta\left(g'^{-1}b\right)$ up to an orientation sign,
so this exponent has the support $g'(\supp\delta)$;
and, the $r$ reflections moving every site in turn, the generators that occur in the expansion of $\Psi^{g'}$
sit at the images $\theta_1\cdots\theta_rx=g'x$ of the sites $x$ of $\Psi$.
Each monomial $\Psi'$ of this expansion, with the exponent $\pm\left(g'^{-1}\right)^{*}\delta$, is again a generator of $\mathcal{P}$.
Indeed, every reflection turns generators $\psi_\alpha$ into combinations of generators $\bar\psi_\beta$ and conversely \eqref{eq:setting-reflection-fields};
after the $r$ reflections the two kinds of generators are therefore preserved if $g'$ is even and interchanged if $g'$ is odd.
Together with the equivariance of $\pi$ under $W$ (Proposition~\ref{prop:geometry}\eqref{it:geometry-projection}) this gives the charge density
$\mathfrak{q}(\Psi')=\pm\mathfrak{q}(\Psi)\circ g'^{-1}$ (Definition~\ref{def:chamber-operator}).
The pull-back $\left(g'^{-1}\right)^{*}$ commutes with $d$ and is orthogonal (Definition~\ref{def:cochains}), hence commutes with $d^{*}$,
and the gauge invariance $d^{*}\delta=q\cdot\mathfrak{q}(\Psi)$ of the generator (Definition~\ref{def:polynomials} and \eqref{eq:setting-gauge}) gives
\[
d^{*}\left(\pm\left(g'^{-1}\right)^{*}\delta\right)=\pm q\cdot\mathfrak{q}(\Psi)\circ g'^{-1}=q\cdot\mathfrak{q}(\Psi').
\]

\emph{The direction $e_1$.}
Let $u\in\Ocub$ with $ue_1=\pm e_1$, where in this paragraph $\pm$ denotes the sign of this relation, and let $c=j_0+\tfrac12$.
If $ue_1=e_1$, then $u$ preserves the coordinate $x_1$;
it maps every plane $\{x_1=j\}$ onto itself, the links between the planes $\{x_1=j_0\}$ and $\{x_1=j_0+1\}$ onto such links,
and every thin layer onto itself, because the set $\mathbb{X}$ of the fermion sites is invariant under $W$ (Definition~\ref{def:fermion-sites}).
If $ue_1=-e_1$, then $u$ reverses the sign of the coordinate $x_1$.
Let $j_0':=-j_0-1$, so that $-c=j_0'+\tfrac12$.
Then $u$ maps the plane $\{x_1=j_0\}$ onto the plane $\{x_1=j_0'+1\}$ and the plane $\{x_1=j_0+1\}$ onto the plane $\{x_1=j_0'\}$,
the links between the first two onto the links between the last two, and, by Definition~\ref{def:two-layer},
\[
u\,\mathbb{X}_{j_0,6}=u\left(\mathbb{X}\cap\left\{x_1=j_0+v_1\right\}\right)=\mathbb{X}\cap\left\{x_1=j_0'+1-v_1\right\}=\mathbb{X}_{j_0'+1,1},
\qquad
u\,\mathbb{X}_{j_0+1,1}=\mathbb{X}_{j_0',6}.
\]
An element of $\mathcal{P}_1(c)$ is a linear combination of generators $\exp\left(i\langle\delta,\hat A\rangle\right)\cdot\Psi$ of $\mathcal{P}$
whose exponents $\delta$ are supported on the links of the planes $\{x_1=j_0\}$, $\{x_1=j_0+1\}$ and on the links between them,
and whose monomials $\Psi$ are monomials in the fields of $\mathbb{X}_{j_0,6}$ and $\mathbb{X}_{j_0+1,1}$ (Definition~\ref{def:two-layer}).
By the description of the images of generators, the image of such a generator under $G\mapsto G^u$ is a linear combination of generators of $\mathcal{P}$
whose exponents are supported on the images of these links and whose monomials are in the fields at the images of these sites,
that is, on the links of the two planes next to the plane $\{x_1=\pm c\}$ and on the links between them,
and in the fields of the two thin layers next to this plane.
As $G\mapsto G^u$ is linear or antilinear, $\left(\mathcal{P}_1(c)\right)^u\subset\mathcal{P}_1(\pm c)$.
The element $u^{-1}$ satisfies $u^{-1}e_1=\pm e_1$ with the same sign,
and the inclusion just proved, applied to $u^{-1}$ and to the plane $\pm c$, gives $\left(\mathcal{P}_1(\pm c)\right)^{u^{-1}}\subset\mathcal{P}_1(c)$.
Applying $G\mapsto G^u$ to this inclusion one obtains $\left(\mathcal{P}_1(c)\right)^u=\mathcal{P}_1(\pm c)$.

\emph{The general case.}
Let $\pm$ again be the sign in $ge_\mu=\pm e_{\mu'}$, and put $u:=g_{\mu'}^{-1}\,g\,g_\mu$.
Then $ue_1=g_{\mu'}^{-1}ge_\mu=\pm g_{\mu'}^{-1}e_{\mu'}=\pm e_1$,
so that the sign of $u$ in the preceding paragraph is the sign of $g$; moreover $gg_\mu=g_{\mu'}u$,
and the composition rule gives
\[
\left(\mathcal{P}_1(\mu,c)\right)^g=\left(\mathcal{P}_1(c)\right)^{gg_\mu}=\left(\mathcal{P}_1(c)\right)^{g_{\mu'}u}
=\left(\left(\mathcal{P}_1(c)\right)^u\right)^{g_{\mu'}}=\left(\mathcal{P}_1(\pm c)\right)^{g_{\mu'}}=\mathcal{P}_1(\mu',\pm c).
\]
Finally let $g''\in\Ocub$ with $g''e_1=e_\mu$.
The element $u:=g_\mu^{-1}g''$ fixes $e_1$, and
\[
\left(\mathcal{P}_1(c)\right)^{g''}=\left(\left(\mathcal{P}_1(c)\right)^u\right)^{g_\mu}=\left(\mathcal{P}_1(c)\right)^{g_\mu}=\mathcal{P}_1(\mu,c).
\]
This proves (a).

For (b), let $c\in\Z+\tfrac12$ and $\Pi:=(x_\mu,c)$; as in the general case, $\pm$ is the sign in $ge_\mu=\pm e_{\mu'}$.
By (a), $g\Pi=(x_{\mu'},\pm c)$, so that $\Theta_\Pi=\Theta_{\mu,c}$ and $\Theta_{g\Pi}=\Theta_{\mu',\pm c}$;
on even elements $\Theta_{g\Pi}$ does not depend on the orientation of the normal \eqref{eq:setting-reflection-fields}.
For a complex number $z$ put $z^g:=z$ if $\det g=1$ and $z^g:=\overline{z}$ if $\det g=-1$;
this map is additive, commutes with complex conjugation and preserves the modulus.
Let $G,P,Y\in\mathcal{P}$, and let $i$ be so large that these elements and their images under $\Theta_\Pi$ embed in $\mathbb{T}_i$.
Items (i) and (iii) give
\begin{align*}
\mu^t_{i,g\mathfrak{d}}\left(\Theta_{g\Pi}\left(G^g\right)\right)-\overline{\mu^t_{i,g\mathfrak{d}}\left(G^g\right)}
&=\mu^t_{i,g\mathfrak{d}}\left(\left(\Theta_\Pi G\right)^g\right)-\overline{\mu^t_{i,g\mathfrak{d}}\left(G^g\right)}\\
&=\left(\mu^t_{i,\mathfrak{d}}\left(\Theta_\Pi G\right)-\overline{\mu^t_{i,\mathfrak{d}}(G)}\right)^g,\\
\mu^t_{i,g\mathfrak{d}}\left(P^g\right)-\mu^t_{i,g\mathfrak{d}}\left(Y^g\right)
&=\left(\mu^t_{i,\mathfrak{d}}(P)-\mu^t_{i,\mathfrak{d}}(Y)\right)^g,\\
\mu^t_{i,g\mathfrak{d}}\left(\Theta_{g\Pi}\left(P^g\right)\right)-\mu^t_{i,g\mathfrak{d}}\left(\Theta_{g\Pi}\left(Y^g\right)\right)
&=\mu^t_{i,g\mathfrak{d}}\left(\left(\Theta_\Pi P\right)^g\right)-\mu^t_{i,g\mathfrak{d}}\left(\left(\Theta_\Pi Y\right)^g\right)\\
&=\left(\mu^t_{i,\mathfrak{d}}\left(\Theta_\Pi P\right)-\mu^t_{i,\mathfrak{d}}\left(\Theta_\Pi Y\right)\right)^g,
\end{align*}
and in each of these three identities the two sides have the same modulus.

\emph{The first property.}
By (a), $Y\mapsto Y^g$ is a bijection of $\mathcal{P}_1(\mu,c)$ onto $\mathcal{P}_1(\mu',\pm c)$,
and $c\mapsto\pm c$ is a bijection of $\Z+\tfrac12$ onto itself.
By the first identity, for every $c$ and every $Y\in\mathcal{P}_1(\mu,c)$ the modulus in \eqref{it:datum-hermitian} for $\mathfrak{d}$, the plane $(x_\mu,c)$ and $Y$
coincides, for all large $i$, with the modulus for $g\mathfrak{d}$, the plane $(x_{\mu'},\pm c)$ and $Y^g$.
Hence $(\mathfrak{d},\mu)$ has the property \eqref{it:datum-hermitian} if and only if $(g\mathfrak{d},\mu')$ has it.

\emph{The second property.}
Suppose that $(\mathfrak{d},\mu)$ has the property \eqref{it:datum-reduction}.
Let $c'\in\Z+\tfrac12$ and $P'\in\mathcal{P}$.
Put $c:=\pm c'$, so that $g\Pi=(x_{\mu'},c')$, and $P:=\left(P'\right)^{g^{-1}}$, an element of $\mathcal{P}$ with $P^g=P'$.
Let $Y_P\in\mathcal{P}_1(\mu,c)$, $C$ and $i_0$ be as in \eqref{it:datum-reduction} for $\mathfrak{d}$, the plane $(x_\mu,c)$ and $P$.
For $i\ge i_0$ the four numbers in \eqref{it:datum-reduction} are defined,
so that $P$, $Y_P$ and their images under $\Theta_\Pi$ embed in $\mathbb{T}_i$.
The element $Y_{P'}:=\left(Y_P\right)^g$ belongs to $\mathcal{P}_1(\mu',c')$ by (a),
and the second and third identities show that the two inequalities of \eqref{it:datum-reduction} hold for $g\mathfrak{d}$, the plane $(x_{\mu'},c')$ and $P'$,
with $Y_{P'}$ in place of $Y_P$ and with the same $C$ and $i_0$.
So $(g\mathfrak{d},\mu')$ has the property \eqref{it:datum-reduction}.
For the converse, apply what has just been shown to $g^{-1}$ and $g\mathfrak{d}$ in place of $g$ and $\mathfrak{d}$:
one has $g^{-1}(g\mathfrak{d})=\mathfrak{d}$ and $g^{-1}e_{\mu'}=\pm e_\mu$.

\emph{The third property.}
The map $\nu\mapsto\left(g^{*}\right)_{*}\nu$ carries the comparison states for $g\mathfrak{d}$ onto the comparison states for $\mathfrak{d}$.
By (iv) it maps the former into the latter, and every comparison state $\nu_0$ for $\mathfrak{d}$ is attained:
the law $\nu:=\left(\left(g^{-1}\right)^{*}\right)_{*}\nu_0$ satisfies $\left(g^{*}\right)_{*}\nu=\nu_0$, since $g^{*}\left(g^{-1}\right)^{*}=\left(g^{-1}g\right)^{*}$ is the identity,
and it is therefore, by (iv) again, a comparison state for $g\mathfrak{d}$.
For such $\nu_0$ and $\nu$, the rule $\left(g'g''\right)^{*}=g''^{*}g'^{*}$ for $g',g''\in W$ (Definition~\ref{def:cochains})
and the relation $g\theta_\Pi=\theta_{g\Pi}g$ from the proof of (i) give
\[
\left(\theta_\Pi^{*}\right)_{*}\nu_0=\left(\left(g\theta_\Pi\right)^{*}\right)_{*}\nu=\left(\left(\theta_{g\Pi}g\right)^{*}\right)_{*}\nu=\left(g^{*}\right)_{*}\left(\theta_{g\Pi}^{*}\right)_{*}\nu.
\]
Since $\left(g^{*}\right)_{*}$ is injective on laws, $\nu_0$ is invariant under $\theta_\Pi$ if and only if $\nu$ is invariant under $\theta_{g\Pi}$.
As $c$ ranges over $\Z+\tfrac12$, the mirror $g\Pi$ ranges over all the mirrors $(x_{\mu'},c')$, $c'\in\Z+\tfrac12$;
hence $(\mathfrak{d},\mu)$ has the property \eqref{it:datum-invariance} if and only if $(g\mathfrak{d},\mu')$ has it.
\end{proof}

\subsection{Invariance of comparison states}\label{sec:comparison-invariance}

\begin{theorem}[No layered states]\label{thm:no-layers}
Let $0<\rho<1$, $w_0,\dots,w_3>0$, $M\in\R$ and $q>0$, and let the sequence of sides $\mathfrak n=(n_i)$ satisfy~\eqref{cond:divisibility}.
Let $\mathfrak d$ be standard data (Definition~\ref{def:data}) or image data (Definition~\ref{def:image-data}), and let $\nu_0$ be a comparison state for $\mathfrak d$ (Definition~\ref{def:comparison-states}).
Then $\nu_0$ is invariant under the reflection $\theta_{x_\mu,c}$, that is,
\[
\left(\theta_{x_\mu,c}^{*}\right)_{*}\nu_0=\nu_0 ,
\]
for $\mu=1,2,3$ and every $c\in\Z+\tfrac12$.
Consequently $\nu_0$ is invariant under $\tau_v$ for every $v\in2\Z^3$.
\end{theorem}

\begin{proof}
\emph{Step~1: reduction to standard data.}
Let $\mathfrak d=g\mathfrak d'$ be image data, where $\mathfrak d'$ is standard data and $g\in\Ocub$; let $\mu\in\{1,2,3\}$, and let $\mu'$ be the index with $ge_{\mu'}=\pm e_\mu$.
Lemma~\ref{lem:data-transport}, applied to the standard data $\mathfrak d'$, to the element $g$ and to the statement~\eqref{it:datum-invariance}, shows that every comparison state for $\mathfrak d$ is invariant under $\theta_{x_\mu,c}$ for all $c\in\Z+\tfrac12$ if and only if every comparison state for $\mathfrak d'$ is invariant under $\theta_{x_{\mu'},c}$ for all $c\in\Z+\tfrac12$.
It is therefore enough to prove the invariance under the reflections for standard data, in each of the three directions.

Up to the end of Step~7, $\mathfrak d$ is standard data, the direction $\mu\in\{1,2,3\}$ is fixed, and $c=j_0+\tfrac12$ with $j_0\in\Z$.
Put $\theta_c:=\theta_{x_\mu,c}$ and $\Theta_c:=\Theta_{\mu,c}$.
The functional $\mu^t_i$ is the perturbed functional of $\mathfrak d$ on the adapted torus $D(n_i)$ (Definition~\ref{def:data}), and $O$ and $Q_0$ are the observable and the base cell of $\mathfrak d$.

\emph{Step~2: the character.}
Fix $\beta\in C^2_{\mathrm c}$ and put $P:=\varphi_\beta$.
By~\eqref{eq:layers-characters}, applied to the mirror $(x_\mu,c)$, the function $P$ lies in $\mathcal P$, and
\[
\Theta_cP=\varphi_{-\theta_c^*\beta} ;
\]
both $P$ and $\Theta_cP$ are continuous functions on $\Omega$ of modulus one that depend on finitely many plaquette variables.
Read on the torus $D(n_i)$, the element $P$ of $\mathcal P$ is the function $\varphi_\beta(\widetilde F)$ of the periodic extension.
Indeed, let $B$ be a box that contains every plaquette of $\supp\beta$ and embeds in $D(n_i)$, and let $\hat A$ be a local lift on $B$ (Definition~\ref{def:integrands}), extended by zero to the links outside $B$.
The links of $\supp d^*\beta$ lie in $B$, and $d\hat A=F$ on the plaquettes of $B$; hence
\[
\langle d^*\beta,\hat A\rangle=\langle\beta,d\hat A\rangle=\langle\beta,\widetilde F\rangle ,
\]
and the element $\exp\left(i\langle d^*\beta,\hat A\rangle\right)$ of $\mathcal P$ (Definition~\ref{def:polynomials}) is the function $\varphi_\beta(\widetilde F)$ (Lemma~\ref{lem:curvature-characters}).
The same holds for the character $\Theta_cP=\varphi_{-\theta_c^*\beta}$.

\emph{Step~3: reduction to the two layers at the plane.}
By Theorem~\ref{thm:two-layer-reduction}, in the form~\eqref{it:reduction-link-plane} if $\mu=1$ and in the form~\eqref{it:reduction-directions} if $\mu\in\{2,3\}$, there are an element $Y_P$ of $\mathcal P_1(\mu,c)$ (denoted $Y^\mu_P$ there if $\mu\in\{2,3\}$), a constant $C<\infty$ and an index $i_0$, all three depending only on $P$, $\mu$, $c$ and $\mathfrak d$, such that
\begin{equation}\label{eq:layers-elimination}
\left|\mu^t_i(P)-\mu^t_i(Y_P)\right|\le C\cdot t,\qquad\left|\mu^t_i(\Theta_cP)-\mu^t_i(\Theta_cY_P)\right|\le C\cdot t
\end{equation}
for $0\le t\le1$ and $i\ge i_0$.
By Definition~\ref{def:other-directions}, $Y_P=(Y')^{g_\mu}$ with $Y'\in\mathcal P_1(c)$, where $g_1$ is the identity.
Let
\[
Y'=\sum_\alpha a_\alpha\cdot Y'_\alpha ,\qquad Y'_\alpha=\exp\left(i\langle\delta_\alpha,\hat A\rangle\right)\cdot\Psi_\alpha ,
\]
be the representation of $Y'$ given in Definition~\ref{def:two-layer}: a finite sum with complex coefficients $a_\alpha$ and pairwise different monomials $Y'_\alpha$ of $\mathcal P_1(c)$.
Put $Y_\alpha:=(Y'_\alpha)^{g_\mu}$, a monomial of $\mathcal P_1(\mu,c)$; then $Y_\alpha=1$ if and only if $Y'_\alpha=1$.
The element $g_\mu$ of $W$ is even, so that $G\mapsto G^{g_\mu}$ is linear (Definition~\ref{def:other-directions}), and $\Theta_c$ is antilinear (see~\eqref{eq:setting-reflection-fields} and~\eqref{eq:setting-reflection-coefficients}); consequently
\[
Y_P=\sum_\alpha a_\alpha\cdot Y_\alpha ,\qquad\Theta_cY_P=\sum_\alpha\overline{a_\alpha}\cdot\Theta_cY_\alpha .
\]

\emph{Step~4: the two-layer principle.}
If $Y_\alpha=1$ for some $\alpha$, then $\Theta_cY_\alpha=1$ by~\eqref{eq:layers-bosonic}, and $\mu^t_i(\Theta_cY_\alpha)-\overline{\mu^t_i(Y_\alpha)}=0$ because $\mu^t_i(1)=1$.
For an index $\alpha$ with $Y_\alpha\ne1$, let $S_\alpha$, $\Lambda_{Y_\alpha}$, $X_{Y_\alpha,\varepsilon}$ and $\xi_{Y_\alpha}(\varepsilon)$ be the set of links, the set of sites, the element of the mid-plane algebra and the number that Definition~\ref{def:two-layer} attaches to the monomial $Y'_\alpha$ of $\mathcal P_1(c)$, as in Theorem~\ref{thm:two-layer-principle}(b).
As $Y'_\alpha\ne1$, its exponent or its Grassmann monomial is not trivial, and the set $\Lambda_{Y_\alpha}$ is not empty.
Let $r_{0,\alpha}$ be the radius of the box $B(c,S_\alpha,\Lambda_{Y_\alpha})$ of Definition~\ref{def:insertions}.
With the integers $k(\mathfrak d,r)$ and the thresholds $i_0(\mathfrak d,r)$ of Lemma~\ref{lem:coarse-structures}, put
\[
k_\alpha:=k(\mathfrak d,2r_{0,\alpha}),\qquad k_*:=\max_\alpha k_\alpha,\qquad i_1:=\max\left(i_0,\max_\alpha i_0(\mathfrak d,2r_{0,\alpha})\right),
\]
the maxima being taken over the $\alpha$ with $Y_\alpha\ne1$; if there is no such $\alpha$, put $k_*:=1$ and $i_1:=i_0$.
The index $i_1\ge i_0$ depends only on $P$, $\mu$, $c$ and $\mathfrak d$.
For $t\ge0$ the numbers $\gamma_{t,i}(k)$ and $\gamma_{t,i}(k)^2-1$ of Definition~\ref{def:amplification-factor} are non-decreasing in $k$, because $\Delta_i\ge0$ (Theorem~\ref{thm:four-tori}).
Theorem~\ref{thm:two-layer-principle}(b), applied to the monomial $Y'_\alpha$, gives the inequality~\eqref{eq:layers-monomial-bound} below with $k_\alpha$ in the place of $k_*$, for $\varepsilon>0$, $t\ge0$ and $i\ge i_0(\mathfrak d,2r_{0,\alpha})$.
Since $k_\alpha\le k_*$ and $i_1\ge i_0(\mathfrak d,2r_{0,\alpha})$, it follows that for every $\alpha$ with $Y_\alpha\ne1$ and for all $\varepsilon>0$, $t\ge0$ and $i\ge i_1$
\begin{multline}\label{eq:layers-monomial-bound}
\left|\mu^t_i(\Theta_cY_\alpha)-\overline{\mu^t_i(Y_\alpha)}\right|\le2\cdot2^{\frac{\#|S_\alpha|}{2}}\cdot\xi_{Y_\alpha}(\varepsilon)\cdot\gamma_{t,i}(k_*)\\
+4\left(\gamma_{t,i}(k_*)^2-1\right)^{\frac12}\cdot\|X_{Y_\alpha,\varepsilon}\|_\# .
\end{multline}

\emph{Step~5: the inequality at finite $i$, $t$ and $\varepsilon$.}
The functional $\mu^t_i$ is linear; together with the two identities at the end of Step~3 this gives
\begin{multline*}
\mu^t_i(\Theta_cP)-\overline{\mu^t_i(P)}=\left[\mu^t_i(\Theta_cP)-\mu^t_i(\Theta_cY_P)\right]\\
+\sum_\alpha\overline{a_\alpha}\cdot\left[\mu^t_i(\Theta_cY_\alpha)-\overline{\mu^t_i(Y_\alpha)}\right]+\overline{\mu^t_i(Y_P)-\mu^t_i(P)} ,
\end{multline*}
where the terms with $Y_\alpha=1$ vanish by Step~4.
By~\eqref{eq:layers-elimination} and~\eqref{eq:layers-monomial-bound}, for $0\le t\le1$, $\varepsilon>0$ and $i\ge i_1$,
\begin{multline}\label{eq:layers-finite}
\left|\mu^t_i(\Theta_cP)-\overline{\mu^t_i(P)}\right|\le2C\cdot t\\
+\sum_{\alpha:\,Y_\alpha\ne1}|a_\alpha|\cdot\left[2\cdot2^{\frac{\#|S_\alpha|}{2}}\cdot\xi_{Y_\alpha}(\varepsilon)\cdot\gamma_{t,i}(k_*)\right.\\
\left.{}+4\left(\gamma_{t,i}(k_*)^2-1\right)^{\frac12}\cdot\|X_{Y_\alpha,\varepsilon}\|_\#\right].
\end{multline}

\emph{Step~6: the limits $i\to\infty$, $t\to0$ and $\varepsilon\to0$, in this order.}
Let $(i_m)$ be a subsequence of the indices, $(t_k)$ a sequence of strengths $t_k>0$ with $t_k\to0$, and $(\nu_{t_k})$ a sequence of laws on $\Omega$, as in Definition~\ref{def:comparison-states} for the comparison state $\nu_0$: for every $k$ the laws $\mu^{t_k}_{i_m}$ converge weakly to $\nu_{t_k}$ as $m\to\infty$, and $\nu_{t_k}$ converges weakly to $\nu_0$ as $k\to\infty$.
The finitely many $k$ with $t_k>1$ are discarded.
Let $k$ and $\varepsilon>0$ be fixed.
By~\eqref{eq:layers-bosonic} and Step~2,
\[
\mu^{t_k}_{i_m}(P)=\int\mathrm{d}\mu^{t_k}_{i_m}(F)\,\varphi_\beta(F),
\]
and likewise for $\Theta_cP$; as $P$ and $\Theta_cP$ are bounded and continuous on $\Omega$, weak convergence gives $\mu^{t_k}_{i_m}(P)\to\nu_{t_k}(P)$ and $\mu^{t_k}_{i_m}(\Theta_cP)\to\nu_{t_k}(\Theta_cP)$ as $m\to\infty$.
On the right side of~\eqref{eq:layers-finite},
\[
\gamma_{t,i}(k_*)=\exp\left(2t\cdot|O|_1\cdot k_*^3\right)\cdot\exp\left(|Q_0|\cdot k_*^3\cdot\Delta_i\right),
\]
and $\Delta_i\to0$ as $i\to\infty$ (Theorem~\ref{thm:four-tori}); hence $\gamma_{t,i}(k_*)$ tends to $\exp\left(2t\cdot|O|_1\cdot k_*^3\right)$ as $i\to\infty$.
Taking $t=t_k$ and $i=i_m$ in~\eqref{eq:layers-finite}, for the $m$ with $i_m\ge i_1$, and letting $m\to\infty$, one obtains
\begin{multline}\label{eq:layers-volume-limit}
\left|\nu_{t_k}(\Theta_cP)-\overline{\nu_{t_k}(P)}\right|\le2C\cdot t_k\\
+\sum_{\alpha:\,Y_\alpha\ne1}|a_\alpha|\cdot\left[2\cdot2^{\frac{\#|S_\alpha|}{2}}\cdot\xi_{Y_\alpha}(\varepsilon)\cdot\exp\left(2t_k\cdot|O|_1\cdot k_*^3\right)\right.\\
\left.{}+4\left(\exp\left(4t_k\cdot|O|_1\cdot k_*^3\right)-1\right)^{\frac12}\cdot\|X_{Y_\alpha,\varepsilon}\|_\#\right].
\end{multline}
Next, $k$ tends to infinity while $\varepsilon$ stays fixed.
Then $t_k\to0$, so that $2C\cdot t_k\to0$, $\exp\left(2t_k\cdot|O|_1\cdot k_*^3\right)\to1$ and $\exp\left(4t_k\cdot|O|_1\cdot k_*^3\right)-1\to0$; moreover $\nu_{t_k}\to\nu_0$ weakly, and each $\|X_{Y_\alpha,\varepsilon}\|_\#$ is a finite number that does not depend on $k$ (Proposition~\ref{prop:two-layer-insertions}(c)).
So~\eqref{eq:layers-volume-limit} yields
\[
\left|\nu_0(\Theta_cP)-\overline{\nu_0(P)}\right|\le2\sum_{\alpha:\,Y_\alpha\ne1}|a_\alpha|\cdot2^{\frac{\#|S_\alpha|}{2}}\cdot\xi_{Y_\alpha}(\varepsilon)\qquad\text{for every }\varepsilon>0 .
\]
Since the left side is independent of $\varepsilon$ and each $\xi_{Y_\alpha}(\varepsilon)$ tends to $0$ with $\varepsilon$ (Proposition~\ref{prop:two-layer-insertions}(c)), one concludes that $\nu_0(\Theta_cP)=\overline{\nu_0(P)}$.

\emph{Step~7: from characters to the law.}
By Steps~2 and~6, for every $\beta\in C^2_{\mathrm c}$,
\[
\nu_0\left(\varphi_{-\theta_c^*\beta}\right)=\overline{\nu_0(\varphi_\beta)}=\nu_0(\varphi_{-\beta}) ;
\]
the second equality holds because $\nu_0$ is a positive measure and $\overline{\varphi_\beta}=\varphi_{-\beta}$.
With $-\beta$ in the place of $\beta$, and with $\theta_c\varphi_\beta=\varphi_{\theta_c^*\beta}$ from~\eqref{eq:layers-characters}, this reads $\nu_0(\theta_c\varphi_\beta)=\nu_0(\varphi_\beta)$ for every $\beta\in C^2_{\mathrm c}$.
By~\eqref{eq:layers-action},
\[
\nu_0(\theta_c\varphi_\beta)=\int\mathrm{d}\nu_0(F)\,\varphi_\beta(\theta_c^*F)=\int\mathrm{d}\left(\left(\theta_c^*\right)_*\nu_0\right)(F)\,\varphi_\beta(F) .
\]
Thus the two laws $\left(\theta_c^*\right)_*\nu_0$ and $\nu_0$ have the same integrals of all characters, and, by linearity, of all elements of $\mathcal T$; by Lemma~\ref{lem:characters-determine}(a) they coincide.

Two signs occur in this step, and they cancel.
The complex conjugation in the reflection of coefficient functions~\eqref{eq:setting-reflection-coefficients} produces the sign in $\Theta_c\varphi_\beta=\varphi_{-\theta_c^*\beta}$, and the positivity of $\nu_0$ the sign in $\overline{\nu_0(\varphi_\beta)}=\nu_0(\varphi_{-\beta})$.
What results is the invariance under $F\mapsto\theta_c^*F$, and not under $F\mapsto-\theta_c^*F$.
For a law that is invariant under $F\mapsto-F$, such as the law $\mu_{\mathbb T}$ of a torus (Proposition~\ref{prop:symmetries}\eqref{it:charge-conjugation}) or a weak limit of such laws, the two invariances are equivalent; the folded array $O_{\mathfrak d}$, however, is in general not an even function of $F$.

\emph{Step~8: translations.}
Let $\mathfrak d$ again be standard data or image data and $\nu_0$ a comparison state for $\mathfrak d$.
For $\mu\in\{1,2,3\}$ and $c\in\Z+\tfrac12$, the identity~\eqref{it:mirror-product} of Definition~\ref{def:mirrors}, for the family $\ell=x_\mu$, where $n_\ell=e_\mu$ and $|n_\ell|^2=1$, gives
\[
\theta_{x_\mu,c+1}\,\theta_{x_\mu,c}=\tau_{2e_\mu} .
\]
By~\eqref{eq:layers-action}, $\nu_0$ is invariant under an element of $W$ if its image under the pull-back by this element is $\nu_0$.
The elements of $W$ under which $\nu_0$ is invariant form a subgroup of $W$.
Indeed, the pull-back by a product of two elements of $W$ is the composition of their pull-backs in the reverse order (Definition~\ref{def:cochains}), so that invariance under two elements gives invariance under their product; and the pull-back by the inverse of an element is the inverse of its pull-back, so that invariance under an element gives invariance under its inverse.
By Steps~1 to~7 this subgroup contains $\theta_{x_\mu,c}$ and $\theta_{x_\mu,c+1}$, hence $\tau_{2e_\mu}$, for $\mu=1,2,3$.
It therefore contains $\tau_v$ for every $v$ in the group $2\Z^3$ generated by $2e_1$, $2e_2$ and $2e_3$.
\end{proof}

\subsection{Laws with nine mirror families}\label{sec:nine-mirrors}

Throughout this subsection, $\nu$ is a law on $\Omega$ that is invariant under $W$
and reflection positive at every mirror of the nine families
(Definitions~\ref{def:mirrors} and~\ref{def:reflection-positivity}).
The arguments of this subsection (reconstruction at a mirror, the cone, zero energy)
are those of the four-dimensional Yang--Mills case,
written for three dimensions, for the field $F$ on $\Omega$,
and for reflection positivity on the open side of a mirror.

Let $L^2:=L^2(\Omega,\nu)$, with the inner product and the norm
\[
\langle H,H'\rangle_2:=\nu\left(\overline{H}\cdot H'\right),
\qquad
\|H\|_2:=\langle H,H\rangle_2^{\frac12}.
\]
For $g\in W$ the map $H\mapsto gH$ of \eqref{eq:layers-action} is unitary on $L^2$:
it preserves the norm because $\nu$ is invariant under $g$,
and $H\mapsto g^{-1}H$ is its inverse.
The space $\mathcal{O}(\R^3)$ of the bounded Borel functions of finitely many coordinates $F_\sigma$
is dense in $L^2$:
\begin{equation}\label{eq:layers-density}
\inf_{G\in\mathcal{O}(\R^3)}\|H-G\|_2=0\qquad\text{for every }H\in L^2 .
\end{equation}
Indeed, consider the family of the Borel subsets of $\Omega$
whose indicators lie in the closure of $\mathcal{O}(\R^3)$ in $L^2$.
It contains $\Omega$ and the cylinder sets,
namely the sets of the $F\in\Omega$ for which $(F_\sigma)_{\sigma\in\Sigma_0}$ lies in a given Borel subset of $\R^{\Sigma_0}$,
with $\Sigma_0$ a finite set of plaquettes,
since the indicators of these sets belong to $\mathcal{O}(\R^3)$ itself.
It is stable under proper differences, the closure being a linear subspace,
and under increasing unions, since the indicators of an increasing sequence of sets
converge in $L^2$ to the indicator of their union, by dominated convergence.
The cylinder sets are stable under finite intersections
and generate the Borel $\sigma$-algebra of $\Omega$,
whose topology is induced by the countable product of the coordinates $F_\sigma$.
By Dynkin's $\pi$-$\lambda$ theorem the family consists of all Borel subsets of $\Omega$.
Consequently every simple function lies in the closure of $\mathcal{O}(\R^3)$,
and the simple functions are dense in $L^2$.

\begin{lemma}[Reconstruction at a lattice mirror]\label{lem:lattice-reconstruction}
Let $\nu$ be a law on $\Omega$ that is invariant under $W$
and reflection positive at every mirror.
Let $\Pi=(\ell,c)$ be a mirror and $\theta:=\theta_\Pi$,
let $v\in\Z^3$ be a positive multiple of $n_\ell$ such that $c':=c+\ell(v)/2$ lies in $\Adm(\ell)$,
and let $v'\in\Z^3$ satisfy $\ell(v')=0$.
For $G,\Phi\in\mathcal{O}^{+}(\ell,c)$ let
\[
\langle G,\Phi\rangle_\Pi:=\nu\left(\overline{\theta G}\cdot\Phi\right).
\]
\begin{enumerate}
\item[(a)] $\langle\cdot,\cdot\rangle_\Pi$ is a positive semidefinite Hermitian form on $\mathcal{O}^{+}(\ell,c)$,
and $\langle G,G\rangle_\Pi\le\|G\|_2^2\le\sup|G|^2$.
Let $\mathcal{H}_\Pi$ be the completion of the quotient of $\mathcal{O}^{+}(\ell,c)$
by the null vectors of this form, with the inner product again written $\langle\cdot,\cdot\rangle_\Pi$,
and let $[G]\in\mathcal{H}_\Pi$ be the class of $G$.
\item[(b)] There is exactly one bounded operator $T$ on $\mathcal{H}_\Pi$
with $T[G]=[\tau_vG]$ for all $G\in\mathcal{O}^{+}(\ell,c)$, and $0\le T\le1$.
\item[(c)] There is exactly one unitary operator $V$ on $\mathcal{H}_\Pi$
with $V[G]=[\tau_{v'}G]$ for all $G\in\mathcal{O}^{+}(\ell,c)$.
\item[(d)] $TV=VT$, and for $G,\Phi\in\mathcal{O}^{+}(\ell,c)$ and integers $n\ge0$, $j\in\Z$,
\[
\left\langle[G],T^nV^j[\Phi]\right\rangle_\Pi=\nu\left(\overline{\theta G}\cdot\tau_{nv+jv'}\Phi\right).
\]
\end{enumerate}
\end{lemma}

\begin{proof}
(a) The form is linear in $\Phi$ and conjugate linear in $G$.
On the diagonal, $\langle G,G\rangle_\Pi=\nu\left(\overline{\theta G}\cdot G\right)\ge0$
is the reflection positivity of $\nu$ at $\Pi$;
a sesquilinear form with real values on the diagonal is Hermitian, by the polarization identity.
Since $\nu$ is invariant under $\theta\in W$,
Lemma~\ref{lem:law-mirror}(a) yields $\langle G,G\rangle_\Pi\le\|G\|_2^2$,
and $\|G\|_2^2=\nu\left(|G|^2\right)\le\sup|G|^2$ because $\nu$ is a probability measure.
The Cauchy--Schwarz inequality for the form shows that its null vectors form a linear subspace,
so that the quotient and its completion $\mathcal{H}_\Pi$ are defined.

(b) Since $v$ is a positive multiple of $n_\ell$, one has $\ell(v)>0$
and $\tau_v\{\ell>c\}=\{\ell>c+\ell(v)\}$, and \eqref{eq:layers-local} gives
\[
\tau_v\,\mathcal{O}^{+}(\ell,c)=\mathcal{O}\left(\{\ell>c+\ell(v)\}\right)\subset\mathcal{O}^{+}(\ell,c).
\]
The linear part of $\theta$, the reflection in the plane orthogonal to $n_\ell$, maps $v$ to $-v$.
Hence $\theta\tau_v=\tau_{-v}\theta$ as maps of $\R^3$,
and $\theta(\tau_vG)=\tau_{-v}(\theta G)$ by \eqref{eq:layers-action}.

\emph{Symmetry.}
Let $G,\Phi\in\mathcal{O}^{+}(\ell,c)$.
The action of $W$ on functions is composition with $g^*$;
it is multiplicative and commutes with complex conjugation.
By the invariance of $\nu$ under $\tau_{-v}\in W$,
\[
\langle G,\tau_v\Phi\rangle_\Pi
=\nu\left(\tau_{-v}\left[\overline{\theta G}\cdot\tau_v\Phi\right]\right)
=\nu\left(\overline{\tau_{-v}\theta G}\cdot\Phi\right)
=\nu\left(\overline{\theta\tau_vG}\cdot\Phi\right)
=\langle\tau_vG,\Phi\rangle_\Pi .
\]

\emph{Positivity.}
One has
\[
v=\frac{\ell(v)}{|n_\ell|^2}\cdot n_\ell=\frac{2(c'-c)}{|n_\ell|^2}\cdot n_\ell ,
\]
so that $\theta_{\ell,c'}\theta_{\ell,c}=\tau_v$ by Definition~\ref{def:mirrors}\eqref{it:mirror-product},
and Definition~\ref{def:mirrors}\eqref{it:mirror-involution} turns this into $\theta=\theta_{\ell,c'}\tau_v$.
Thus $\theta G=\theta_{\ell,c'}(\tau_vG)$ for $G\in\mathcal{O}^{+}(\ell,c)$.
The function $\tau_vG$ lies in $\mathcal{O}\left(\{\ell>c+\ell(v)\}\right)$, and $c+\ell(v)>c'$;
so $\tau_vG\in\mathcal{O}^{+}(\ell,c')$,
and the reflection positivity of $\nu$ at the mirror $(\ell,c')$ yields
\[
\langle G,\tau_vG\rangle_\Pi=\nu\left(\overline{\theta_{\ell,c'}(\tau_vG)}\cdot\tau_vG\right)\ge0 .
\]

\emph{Contraction.}
Let $G\in\mathcal{O}^{+}(\ell,c)$, and for $n\ge0$ let
$a_n:=\left\langle\tau_v^nG,\tau_v^nG\right\rangle_\Pi^{\frac12}$.
Moving $\tau_v$ from the left to the right entry $n$ times by the symmetry,
and then applying the Cauchy--Schwarz inequality, one obtains
\[
a_n^2=\left\langle G,\tau_v^{2n}G\right\rangle_\Pi\le a_0\cdot a_{2n}.
\]
Taken for $n=1,2,4,\dots$, this yields, by induction on $k\ge1$,
\[
a_1\le a_0^{1-2^{-k}}\cdot a_{2^k}^{2^{-k}},
\]
while $a_{2^k}\le\sup\left|\tau_v^{2^k}G\right|\le\sup|G|$ by (a).
If $a_0>0$, then $\sup|G|\ge a_0>0$, and letting $k\to\infty$ one obtains $a_1\le a_0$;
if $a_0=0$, then $a_1^2\le a_0\cdot a_2=0$.
In both cases
\[
\langle\tau_vG,\tau_vG\rangle_\Pi\le\langle G,G\rangle_\Pi .
\]

\emph{The operator.}
By the contraction, $\tau_v$ maps null vectors to null vectors,
and $[G]\mapsto[\tau_vG]$ is a well defined linear contraction
on the subspace $\left\{[G]:G\in\mathcal{O}^{+}(\ell,c)\right\}$, which is dense in $\mathcal{H}_\Pi$.
It has exactly one continuous extension $T$ to $\mathcal{H}_\Pi$, and $\|T\|\le1$.
By the symmetry, $\langle[G],T[\Phi]\rangle_\Pi=\langle T[G],[\Phi]\rangle_\Pi$ on the dense subspace,
and by continuity $T$ is self-adjoint.
The positivity shows $\langle[G],T[G]\rangle_\Pi\ge0$ on the dense subspace, hence $T\ge0$;
together with $\|T\|\le1$ this is $0\le T\le1$.

(c) Since $\ell(v')=0$, the vector $v'$ is orthogonal to $n_\ell$ and is fixed by the linear part of $\theta$;
so $\theta\tau_{v'}=\tau_{v'}\theta$, and $\theta(\tau_{v'}G)=\tau_{v'}(\theta G)$ by \eqref{eq:layers-action}.
Moreover $\tau_{v'}\{\ell>c\}=\{\ell>c\}$,
whence $\tau_{v'}\mathcal{O}^{+}(\ell,c)=\mathcal{O}^{+}(\ell,c)$ by \eqref{eq:layers-local}.
Invariance of $\nu$ under $\tau_{v'}$ implies, for $G,\Phi\in\mathcal{O}^{+}(\ell,c)$,
\[
\langle\tau_{v'}G,\tau_{v'}\Phi\rangle_\Pi
=\nu\left(\tau_{v'}\left[\overline{\theta G}\cdot\Phi\right]\right)
=\langle G,\Phi\rangle_\Pi .
\]
Therefore $[G]\mapsto[\tau_{v'}G]$ is a well defined isometry of the dense subspace of the classes onto itself,
with inverse $[G]\mapsto[\tau_{-v'}G]$.
Its continuous extension $V$ is an isometry with closed and dense range, hence unitary,
and it is the only bounded operator with $V[G]=[\tau_{v'}G]$ for all $G$.

(d) The translations $\tau_v$ and $\tau_{v'}$ commute;
so $TV[G]=[\tau_v\tau_{v'}G]=[\tau_{v'}\tau_vG]=VT[G]$, and $TV=VT$ by density.
As $V^{-1}[\Phi]=[\tau_{-v'}\Phi]$, induction shows $T^nV^j[\Phi]=[\tau_{nv+jv'}\Phi]$
for $n\ge0$ and $j\in\Z$,
and the asserted identity is the definition of the form.
\end{proof}

\begin{lemma}[A two-constants estimate on the disc]\label{lem:two-constants}
Let $\overline{\mathbb{D}}$ be the closed and $\mathbb{D}$ the open unit disc of $\C$,
and $S^1$ the unit circle.
Let $\eta$ be a finite positive Borel measure on $\overline{\mathbb{D}}$,
and let $C\ge0$ be such that every polynomial $p$ in one complex variable satisfies
\begin{equation}\label{eq:layers-two-constants}
\int_{\overline{\mathbb{D}}}\mathrm{d}\eta(z)\,|p(z)|^2\le C\cdot\sup_{[0,1]}|p|\cdot\sup_{S^1}|p| .
\end{equation}
Then $\eta\left(\left\{z\in\overline{\mathbb{D}}:|1+z|^2<4|z|\right\}\right)=0$.
\end{lemma}

\begin{proof}
\emph{Step 1 (holomorphic functions).}
Let $p$ be holomorphic on an open disc centred at $0$ of radius larger than $1$.
Its Taylor polynomials converge to $p$ uniformly on $\overline{\mathbb{D}}$.
As $\eta$ is finite, both sides of \eqref{eq:layers-two-constants}, written for these polynomials,
converge to the corresponding sides for $p$;
hence \eqref{eq:layers-two-constants} holds for $p$.

\emph{Step 2 (test functions).}
For $\xi\in\mathbb{D}$ let
\[
w(\xi):=\frac{1+\xi}{1-\xi}=\frac{1-|\xi|^2+2i\cdot\Im\xi}{|1-\xi|^2},
\]
the second expression coming from $(1+\xi)\left(1-\overline{\xi}\right)=1-|\xi|^2+2i\cdot\Im\xi$.
Thus $\Re w>0$ on $\mathbb{D}$,
the principal logarithm $\operatorname{Log}w$ is holomorphic on $\mathbb{D}$ with $|\Im\operatorname{Log}w|<\pi/2$,
and it is odd, since $w(-\xi)=1/w(\xi)$.
Fix $s>0$.
The function $\xi\mapsto\cos\left(s\cdot\operatorname{Log}w(\xi)\right)$ is holomorphic and even on $\mathbb{D}$;
its Taylor series at $0$ converges on $\mathbb{D}$ and contains only even powers,
so that $\cos\left(s\cdot\operatorname{Log}w(\xi)\right)=g_s\left(\xi^2\right)$ with a function $g_s$ holomorphic on $\mathbb{D}$.
Put $f_s:=\left(\cosh(\pi s/2)\right)^{-1}\cdot g_s$.
For real $x$ and $y$,
\[
|\cos(x+iy)|^2=\cos^2x\cdot\cosh^2y+\sin^2x\cdot\sinh^2y=\cos^2x+\sinh^2y .
\]
Let $z\in\mathbb{D}$, let $\xi$ be a square root of $z$,
write $s\cdot\operatorname{Log}w(\xi)=x+iy$ with real $x,y$, and put
\[
\alpha(z):=|\arg w(\xi)|\in[0,\pi/2);
\]
this does not depend on the choice of $\xi$, because $\operatorname{Log}w$ is odd.
Then $|y|=s\cdot\alpha(z)<\pi s/2$, and from $|g_s(z)|^2=\cos^2x+\sinh^2y$ one obtains:
\begin{enumerate}
\item[(i)] $|f_s(z)|\le1$, because $\cos^2x+\sinh^2y\le\cosh^2y\le\cosh^2(\pi s/2)$;
\item[(ii)] $\cosh(\pi s/2)\cdot|f_s(z)|\le1$ if $z\in[0,1)$,
because $\xi$ may then be taken in $[0,1)$, so that $w(\xi)>0$ and $y=0$;
\item[(iii)] $\cosh(\pi s/2)\cdot|f_s(z)|\ge\sinh\left(s\cdot\alpha(z)\right)$,
because $\cos^2x+\sinh^2y\ge\sinh^2y$.
\end{enumerate}

\emph{Step 3 (the angle).}
By the second expression for $w$, $\tan\alpha(z)=2|\Im\xi|/\left(1-|\xi|^2\right)$ for $z=\xi^2\in\mathbb{D}$.
Since $|\Im\xi|^2=(|z|-\Re z)/2$, this reads
\[
\tan\alpha(z)=\frac{\sqrt{2\left(|z|-\Re z\right)}}{1-|z|}\qquad(z\in\mathbb{D}),
\]
and $\alpha$ is continuous on $\mathbb{D}$.
Expanding the squares one finds
\[
\left|1+\xi^2\right|^2-4|\xi|^2=1-2(\Re\xi)^2-6(\Im\xi)^2+|\xi|^4=\left(1-|\xi|^2\right)^2-4(\Im\xi)^2 ,
\]
so that for $z=\xi^2\in\mathbb{D}$
\[
\alpha(z)>\pi/4\iff2|\Im\xi|>1-|\xi|^2\iff|1+z|^2<4|z| .
\]
Fix once and for all a sequence $(r_n)_{n\ge1}$ in $(0,1)$ with $r_n\to1$.
Let $z\in S^1\setminus\{1\}$, and let $\xi_0$ be the square root of $z$ with $\Im\xi_0>0$.
For each $n$ the point $\xi_n:=\sqrt{r_n}\cdot\xi_0$ lies in $\mathbb{D}$ and is a square root of $r_nz\in\mathbb{D}$.
As $\xi_n\to\xi_0\neq1$ and the quotient $(1+\xi)/(1-\xi)$ is continuous at every $\xi\neq1$,
\[
w(\xi_n)=\frac{1+\xi_n}{1-\xi_n}\to\frac{1+\xi_0}{1-\xi_0}\qquad(n\to\infty).
\]
The identity $(1+\xi)\left(1-\overline{\xi}\right)=1-|\xi|^2+2i\cdot\Im\xi$ of Step~2 holds for every $\xi\in\C$;
at $\xi=\xi_0$, where $|\xi_0|=1$, it gives
\[
\frac{1+\xi_0}{1-\xi_0}=\frac{(1+\xi_0)\left(1-\overline{\xi_0}\right)}{|1-\xi_0|^2}=\frac{2i\cdot\Im\xi_0}{|1-\xi_0|^2},
\]
which lies on the positive imaginary axis because $\Im\xi_0>0$.
Hence $\alpha(r_nz)=|\arg w(\xi_n)|\to\pi/2$.
Let $\widetilde{\alpha}:=\alpha$ on $\mathbb{D}$ and $\widetilde{\alpha}:=\pi/2$ on $S^1\setminus\{1\}$.
This is a Borel function on $\overline{\mathbb{D}}\setminus\{1\}$,
$\alpha(r_nz)\to\widetilde{\alpha}(z)$ for every $z\in\overline{\mathbb{D}}\setminus\{1\}$
(for $z\in\mathbb{D}$ because $r_nz\to z$ and $\alpha$ is continuous on $\mathbb{D}$), and
\[
\left\{z\in\overline{\mathbb{D}}:|1+z|^2<4|z|\right\}
\subset\left\{z\in\overline{\mathbb{D}}\setminus\{1\}:\widetilde{\alpha}(z)>\pi/4\right\},
\]
since $\widetilde{\alpha}=\pi/2$ on $S^1\setminus\{1\}$, while $|1+z|^2=4=4|z|$ at $z=1$.

\emph{Step 4 (the estimate).}
For each $n$ the function $p(z):=f_s(r_nz)$ is holomorphic on $\{|z|<1/r_n\}$,
and (i), (ii) show that $\sup_{S^1}|p|\le1$ and $\cosh(\pi s/2)\cdot\sup_{[0,1]}|p|\le1$.
By Step 1,
\[
\cosh(\pi s/2)\cdot\int_{\overline{\mathbb{D}}}\mathrm{d}\eta(z)\,|f_s(r_nz)|^2\le C .
\]
By (iii), $\cosh^2(\pi s/2)\cdot|f_s(r_nz)|^2\ge\sinh^2\left(s\cdot\alpha(r_nz)\right)$,
and for $z\in\overline{\mathbb{D}}\setminus\{1\}$ the right side tends to
$\sinh^2\left(s\cdot\widetilde{\alpha}(z)\right)$ as $n\to\infty$.
Fatou's lemma therefore gives
\[
\int_{\overline{\mathbb{D}}\setminus\{1\}}\mathrm{d}\eta(z)\,\sinh^2\left(s\cdot\widetilde{\alpha}(z)\right)
\le\liminf_{n\to\infty}\cosh^2(\pi s/2)\cdot\int_{\overline{\mathbb{D}}}\mathrm{d}\eta(z)\,|f_s(r_nz)|^2
\le C\cdot\cosh(\pi s/2).
\]
For $\varepsilon>0$ let
$S_\varepsilon:=\left\{z\in\overline{\mathbb{D}}\setminus\{1\}:\widetilde{\alpha}(z)\ge\pi/4+\varepsilon\right\}$.
As $\sinh$ is increasing on $[0,\infty)$,
\[
\sinh^2\left(s(\pi/4+\varepsilon)\right)\cdot\eta(S_\varepsilon)\le C\cdot\cosh(\pi s/2).
\]
If $s(\pi/4+\varepsilon)\ge1$, then $\cosh(\pi s/2)\le\exp(\pi s/2)$
and $\sinh\left(s(\pi/4+\varepsilon)\right)\ge\frac14\exp\left(s(\pi/4+\varepsilon)\right)$, and therefore
\[
\eta(S_\varepsilon)\le16\,C\cdot\exp(-2s\varepsilon).
\]
Letting $s\to\infty$ one finds $\eta(S_\varepsilon)=0$ for every $\varepsilon>0$.
The set $\left\{\widetilde{\alpha}>\pi/4\right\}$ is the union of the sets $S_{1/k}$, $k=1,2,\dots$,
so it is an $\eta$-null set,
and with the inclusion of Step 3 it follows that $\eta\left(\left\{z\in\overline{\mathbb{D}}:|1+z|^2<4|z|\right\}\right)=0$.
\end{proof}

\begin{definition}[The axis picture]\label{def:axis-picture}
Let $\nu$ be a law on $\Omega$ that is invariant under $W$ and reflection positive (Definition~\ref{def:reflection-positivity}) at every mirror of the nine families (Definition~\ref{def:mirrors}),
and let $\mu,\mu'\in\{1,2,3\}$ with $\mu\neq\mu'$.
Let $\theta_\mu:=\theta_{x_\mu,-\frac{1}{2}}$, that is, $\theta_\mu x=x-(2x_\mu+1)e_\mu$, and
\[
\mathcal{O}^{+}_{\mu}:=\mathcal{O}^{+}\left(x_\mu,-\tfrac12\right)=\mathcal{O}\left(\left\{x_\mu>-\tfrac12\right\}\right);
\]
the plaquettes on which an element of $\mathcal{O}^{+}_{\mu}$ depends have all their corners in $\{x_\mu\geq0\}$.
Lemma~\ref{lem:lattice-reconstruction}, applied to the mirror $(x_\mu,-\frac12)$ with $v:=e_\mu$ and $v':=e_{\mu'}$ (so that $c'=0$, a site plane),
provides the reflected form
\[
\langle G,\Phi\rangle_\mu=\nu\left(\overline{\theta_\mu G}\cdot\Phi\right)\qquad\left(G,\Phi\in\mathcal{O}^{+}_{\mu}\right),
\]
the Hilbert space $\mathcal{H}_\mu$ with the class map $G\mapsto[G]_\mu$,
and two operators $T_\mu$ and $U_{\mu,\mu'}$ on $\mathcal{H}_\mu$, determined by
\[
T_\mu[G]_\mu=[\tau_{e_\mu}G]_\mu,\qquad U_{\mu,\mu'}[G]_\mu=[\tau_{e_{\mu'}}G]_\mu\qquad\left(G\in\mathcal{O}^{+}_{\mu}\right);
\]
one has $0\leq T_\mu\leq1$, the operator $U_{\mu,\mu'}$ is unitary, and $T_\mu U_{\mu,\mu'}=U_{\mu,\mu'}T_\mu$.
The \emph{axis picture} $\mu$, with second direction $\mu'$, consists of $\mathcal{H}_\mu$, $T_\mu$ and $U_{\mu,\mu'}$.
For fixed $\mu$ and $\mu'$ write $T:=T_\mu$ and $U:=U_{\mu,\mu'}$.

The operator $T$ is self-adjoint and commutes with $U$, hence also with $U^{*}=U^{-1}$;
so the $C^{*}$-algebra generated by $T$, $U$ and $U^{*}$ is commutative.
Every character of this algebra takes at $T$ a value in $[0,1]$ (as $0\leq T\leq1$) and at $U$ a value in $S^1$ (as $U^{*}U=1$), and it is determined by these two values;
in this way the maximal ideal space of the algebra is identified with a compact subset of $[0,1]\times S^1$.
The spectral theorem for commuting bounded normal operators \cite[Theorem~12.22]{Rudin1991}, applied to this algebra, gives a resolution of the identity on its maximal ideal space.
Transported to $[0,1]\times S^1$ by the identification just described, it is the \emph{joint spectral measure} $\mathfrak{E}_\mu$ of $(T,U)$:
a projection-valued measure on the Borel sets of $[0,1]\times S^1$ such that integration against $\mathfrak{E}_\mu$ is a $*$-homomorphism from the bounded Borel functions on $[0,1]\times S^1$ to the bounded operators on $\mathcal{H}_\mu$, with
\[
\int\mathrm{d}\mathfrak{E}_\mu(t,u)\,t=T,\qquad\int\mathrm{d}\mathfrak{E}_\mu(t,u)\,u=U.
\]
For $\psi\in\mathcal{H}_\mu$ the positive Borel measure $\mathfrak{e}_\psi:=\langle\psi,\mathfrak{E}_\mu(\cdot)\psi\rangle_\mu$ on $[0,1]\times S^1$ has the mass $\|\psi\|^2$, and
\[
\left\langle\psi,T^{n}U^{y}\psi\right\rangle_\mu=\int\mathrm{d}\mathfrak{e}_\psi(t,u)\,t^{n}\cdot u^{y}\qquad(n\geq0,\ y\in\Z).
\]
\end{definition}

Theorem~\ref{thm:cubic-limits} shows that, under its hypotheses on the parameters,
every weak limit point of the laws of the cubic tori $C(n_i)$, where $n_i\geq4$ are integers with $n_i\to\infty$,
is a law $\nu$ as in the following theorem; see Corollary~\ref{cor:all-cubic-tori}.

\begin{theorem}[The lattice cone]\label{thm:lattice-cone}
Let $\nu$ be a law on $\Omega$ that is invariant under $W$ and reflection positive (Definition~\ref{def:reflection-positivity}) at every mirror of the nine families,
and let $\mu,\mu'\in\{1,2,3\}$ with $\mu\neq\mu'$.
\begin{enumerate}
\item[(a)] In the axis picture $\mu$ with second direction $\mu'$ (Definition~\ref{def:axis-picture}),
the operator $T_\mu$ on the Hilbert space $\mathcal{H}_\mu$, induced by the translation by $e_\mu$, satisfies $0\leq T_\mu\leq1$,
the operator $U_{\mu,\mu'}$, induced by the translation by $e_{\mu'}$, is unitary,
the two operators commute, and
\[
(1-T_\mu)^2\geq T_\mu\cdot\left(2-U_{\mu,\mu'}-U_{\mu,\mu'}^{*}\right).
\]
\item[(b)] The joint spectral measure $\mathfrak{E}_\mu$ of $(T_\mu,U_{\mu,\mu'})$ is carried by the set
\[
\left\{(t,u)\in[0,1]\times S^1:\ (1-t)^2\geq t\cdot|1-u|^2\right\}.
\]
For $(t,u)$ with $t>0$ write $t=\exp(-\lambda)$, $\lambda\geq0$, and $u=\exp(ik)$, $k\in\R$;
then $\cosh\lambda+\cos k\geq2$ for $\mathfrak{E}_\mu$-almost every such $(t,u)$.
\item[(c)] (Lattice form of condition~\eqref{it:projections-coincide} of Definition~\ref{def:symmetry-conditions}.)
If $H\in L^2(\Omega,\nu)$ satisfies $\tau_{e_\mu}H=H$, then $\tau_vH=H$ for every $v\in\Z^3$.
\item[(d)] (Lattice form of condition~\eqref{it:components-cubic} of Definition~\ref{def:symmetry-conditions}.)
Let $\mathcal{B}^{\mathrm{inv}}$ be the $\sigma$-algebra of the Borel subsets of $\Omega$ that are invariant under $\tau_v^{*}$ for every $v\in\Z^3$,
and let $F\mapsto\nu_F$ be a regular conditional distribution of $\nu$ given $\mathcal{B}^{\mathrm{inv}}$.
For $\nu$-almost every $F$ the law $\nu_F$ is invariant under $\Z^3$, ergodic under $\Z^3$, and invariant under $W$.
\end{enumerate}
\end{theorem}

\begin{proof}[Proof of Theorem~\ref{thm:lattice-cone}(a) and (b)]
Apart from the inequality, the assertions of (a) are those of Lemma~\ref{lem:lattice-reconstruction}(b)--(d), as recorded in Definition~\ref{def:axis-picture}.
Fix $\mu\neq\mu'$; as in that definition, $T=T_\mu$ and $U=U_{\mu,\mu'}$.

\emph{Operator inequality and spectral support.}
Let first $f$ be an arbitrary real continuous function on $[0,1]\times S^1$.
For $\psi\in\mathcal{H}_\mu$ one has
\[
\left\langle\psi,\left(\int\mathrm{d}\mathfrak{E}_\mu(t,u)\,f(t,u)\right)\psi\right\rangle_\mu=\int\mathrm{d}\mathfrak{e}_\psi(t,u)\,f(t,u),
\qquad
\mathfrak{e}_\psi=\left\|\mathfrak{E}_\mu(\cdot)\psi\right\|^2 .
\]
If $\mathfrak{E}_\mu(\{f<0\})=0$, every $\mathfrak{e}_\psi$ is carried by $\{f\geq0\}$, and the operator $\int\mathrm{d}\mathfrak{E}_\mu(t,u)\,f(t,u)$ is non-negative.
Suppose, conversely, that $\mathfrak{E}_\mu(\{f<0\})\neq0$.
The set $\{f<0\}$ is the increasing union of the sets $\{f\leq-1/n\}$, $n\geq1$, and $\mathfrak{E}_\mu$ is countably additive in the strong operator topology;
hence $\mathfrak{E}_\mu(\{f\leq-\varepsilon\})\neq0$ for some $\varepsilon>0$.
For a unit vector $\psi$ in the range of $\mathfrak{E}_\mu(\{f\leq-\varepsilon\})$ the measure $\mathfrak{e}_\psi$ is carried by $\{f\leq-\varepsilon\}$, and therefore
\[
\left\langle\psi,\left(\int\mathrm{d}\mathfrak{E}_\mu(t,u)\,f(t,u)\right)\psi\right\rangle_\mu\leq-\varepsilon .
\]
Thus $\int\mathrm{d}\mathfrak{E}_\mu(t,u)\,f(t,u)\geq0$ if and only if $\mathfrak{E}_\mu(\{f<0\})=0$.

From now on $f$ denotes the function
\[
f(t,u):=(1-t)^2-t\cdot|1-u|^2\qquad\left((t,u)\in[0,1]\times S^1\right).
\]
For $|u|=1$ one has $|1-u|^2=2-u-\overline{u}$.
Integration against $\mathfrak{E}_\mu$ is a $*$-homomorphism that sends $t$ to $T$ and $u$ to $U$, hence $\overline{u}$ to $U^{*}$; consequently
\[
\int\mathrm{d}\mathfrak{E}_\mu(t,u)\,f(t,u)=(1-T)^2-T\cdot\left(2-U-U^{*}\right).
\]
The inequality of (a) is therefore equivalent to $\mathfrak{E}_\mu(\{f<0\})=0$, that is, to the first assertion of (b).
If $t=\exp(-\lambda)>0$ and $u=\exp(ik)$ with $k\in\R$, then
\[
\frac{f(t,u)}{t}=(2\cosh\lambda-2)-(2-2\cos k),
\]
so that on $\{t>0\}$ the inequality $f\geq0$ holds exactly where $\cosh\lambda+\cos k\geq2$.
The second assertion of (b) is thus a consequence of the first.

\emph{A second picture.}
Let $G\in\mathcal{O}^{+}_{\mu}$.
It depends on finitely many $F_\sigma$, and every corner of each of these plaquettes $\sigma$ has $x_\mu\geq0$;
so there is an integer $r\geq1$ with $G\in\mathcal{O}(\{0\leq x_\mu\leq r,\ |x_{\mu'}|\leq r\})$.
Let $\ell:=x_\mu+x_{\mu'}$, one of the nine families, let $\theta':=\theta_{\ell,r}$, and
\[
v_{+}:=e_\mu+e_{\mu'},\qquad v_{-}:=e_\mu-e_{\mu'}.
\]
Then $n_\ell=v_{+}$, $\ell(v_{+})=2$ and $\ell(v_{-})=0$, and $r+\frac12\ell(v_{+})=r+1$ lies in $\Adm(\ell)=\Z$.
Lemma~\ref{lem:lattice-reconstruction}, applied to the mirror $(\ell,r)$ with $v:=v_{+}$ and $v':=v_{-}$,
provides the reflected form $\langle\cdot,\cdot\rangle'$ on $\mathcal{O}^{+}(\ell,r)$, the Hilbert space $\mathcal{H}'$ with the class map $[\cdot]'$,
an operator $T_{+}$ on $\mathcal{H}'$ with $0\leq T_{+}\leq1$, induced by the translation by $v_{+}$,
and a unitary $V_{-}$ on $\mathcal{H}'$, induced by the translation by $v_{-}$.

By \eqref{eq:layers-local}, $\theta_\mu G$ lies in $\mathcal{O}(\{-r-1\leq x_\mu\leq-1,\ |x_{\mu'}|\leq r\})$, and $\ell\leq r-1$ on this set.
Since $\ell(\theta'x)=2r-\ell(x)$ (Definition~\ref{def:mirrors}\eqref{it:mirror-height}), the reflection $\theta'$ maps $\{\ell\leq r-1\}$ onto $\{\ell\geq r+1\}$;
by \eqref{eq:layers-local} again, the function
\[
\Phi:=\theta'(\theta_\mu G)
\]
lies in $\mathcal{O}(\{\ell\geq r+1\})\subset\mathcal{O}^{+}(\ell,r)$.
Moreover $\theta'\Phi=\theta_\mu G$, because $\theta'$ is an involution and \eqref{eq:layers-action} defines an action.
Likewise, the function $G_0:=\tau_{(2r+1)e_\mu}G$ lies in $\mathcal{O}(\{2r+1\leq x_\mu\leq3r+1,\ |x_{\mu'}|\leq r\})$, and $\ell\geq r+1$ on this set;
so $G_0\in\mathcal{O}^{+}(\ell,r)$ as well.
(With the shift by $2r$ of the corresponding argument for four-dimensional Yang--Mills theory,
the point with $x_\mu=2r$ and $x_{\mu'}=-r$, a corner of the region in which the shifted function is read, would lie on the mirror $\{\ell=r\}$,
whereas reflection positivity in the sense of Definition~\ref{def:reflection-positivity} concerns the open side $\{\ell>r\}$.)

For integers $k,j\geq0$,
\[
(2r+1+k+j)e_\mu+(k-j)e_{\mu'}=(2r+1)e_\mu+kv_{+}+jv_{-},
\]
so that $\tau_{kv_{+}+jv_{-}}G_0$ is the translate of $G$ by the vector on the left.
Lemma~\ref{lem:lattice-reconstruction}(d), applied in the axis picture with the exponents $2r+1+k+j$ and $k-j$
and in the second picture with the exponents $k$ and $j$, together with $\theta_\mu G=\theta'\Phi$, gives for all integers $k,j\geq0$
\begin{equation}\label{eq:layers-pictures}
\begin{aligned}
\left\langle[G]_\mu,T^{2r+1+k+j}U^{k-j}[G]_\mu\right\rangle_\mu
&=\nu\left(\overline{\theta_\mu G}\cdot\tau_{kv_{+}+jv_{-}}G_0\right)\\
&=\nu\left(\overline{\theta'\Phi}\cdot\tau_{kv_{+}+jv_{-}}G_0\right)
=\left\langle[\Phi]',T_{+}^{k}V_{-}^{j}[G_0]'\right\rangle'.
\end{aligned}
\end{equation}

\emph{An image measure on the closed disc.}
Let $G$, the integer $r$ and the objects of the second picture be fixed as above,
and let $\overline{\mathbb{D}}$ be the closed unit disc of $\C$.
The map $(t,u)\mapsto tu$ is continuous from $[0,1]\times S^1$ to $\overline{\mathbb{D}}$.
Let $\eta$ be the image under this map of the measure $t^{2r+1}\,\mathrm{d}\mathfrak{e}_{[G]_\mu}(t,u)$;
it is a positive Borel measure on $\overline{\mathbb{D}}$ of mass at most $\|[G]_\mu\|^2$,
and for every polynomial $p$ in one complex variable
\[
\int_{\overline{\mathbb{D}}}\mathrm{d}\eta(z)\,|p(z)|^2=\int\mathrm{d}\mathfrak{e}_{[G]_\mu}(t,u)\,t^{2r+1}\cdot|p(tu)|^2 .
\]
Write $p(z)=\sum_{k\geq0}a_k\cdot z^k$, with complex coefficients $a_k$ of which only finitely many differ from $0$,
and let
\[
p^{*}(z):=\sum_{k\geq0}\overline{a_k}\cdot z^k ,
\]
so that $p^{*}(z)=\overline{p\left(\overline{z}\right)}$.
For $t\in[0,1]$ and $u\in S^1$ one has $\overline{u}=u^{-1}$, and therefore
\[
|p(tu)|^2=p(tu)\cdot\overline{p(tu)}=\sum_{k,j\geq0}a_k\cdot\overline{a_j}\cdot t^{k+j}\cdot u^{k-j}.
\]
Integrating this identity against $t^{2r+1}\,\mathrm{d}\mathfrak{e}_{[G]_\mu}(t,u)$,
and using in turn the moment formula of Definition~\ref{def:axis-picture} for $\psi=[G]_\mu$ with the exponents $n=2r+1+k+j\geq0$ and $y=k-j$,
the identity~\eqref{eq:layers-pictures},
and the linearity of $\langle\cdot,\cdot\rangle'$ in its second argument,
one obtains
\[
\begin{aligned}
\int_{\overline{\mathbb{D}}}\mathrm{d}\eta(z)\,|p(z)|^2
&=\sum_{k,j\geq0}a_k\cdot\overline{a_j}\cdot\int\mathrm{d}\mathfrak{e}_{[G]_\mu}(t,u)\,t^{2r+1+k+j}\cdot u^{k-j}\\
&=\sum_{k,j\geq0}a_k\cdot\overline{a_j}\cdot\left\langle[G]_\mu,T^{2r+1+k+j}U^{k-j}[G]_\mu\right\rangle_\mu\\
&=\sum_{k,j\geq0}a_k\cdot\overline{a_j}\cdot\left\langle[\Phi]',T_{+}^{k}V_{-}^{j}[G_0]'\right\rangle'
=\left\langle[\Phi]',p(T_{+})p^{*}(V_{-})[G_0]'\right\rangle' .
\end{aligned}
\]
As $0\leq T_{+}\leq1$, the spectrum of $T_{+}$ lies in $[0,1]$; the spectrum of the unitary $V_{-}$ lies in $S^1$.
The norm of a polynomial in a normal operator equals the maximum of the modulus of that polynomial on the spectrum of the operator:
the polynomial in the operator is again normal, so its norm is its spectral radius,
and by the spectral mapping theorem its spectrum is the image of the spectrum of the operator under the polynomial.
Hence
\[
\left\|p(T_{+})\right\|\leq\sup_{[0,1]}|p|,
\qquad
\left\|p^{*}(V_{-})\right\|\leq\sup_{S^1}\left|p^{*}\right|=\sup_{S^1}|p| ,
\]
where the equality holds because $p^{*}(z)=\overline{p\left(\overline{z}\right)}$ and complex conjugation maps $S^1$ onto itself.
Let
\[
C:=\left\|[\Phi]'\right\|\cdot\left\|[G_0]'\right\| ,
\]
the norms being those of $\mathcal{H}'$.
The Cauchy--Schwarz inequality in $\mathcal{H}'$ now yields, for every polynomial $p$,
\[
\int_{\overline{\mathbb{D}}}\mathrm{d}\eta(z)\,|p(z)|^2
\leq\left\|[\Phi]'\right\|\cdot\left\|p(T_{+})\right\|\cdot\left\|p^{*}(V_{-})\right\|\cdot\left\|[G_0]'\right\|
\leq C\cdot\sup_{[0,1]}|p|\cdot\sup_{S^1}|p| .
\]
Thus~\eqref{eq:layers-two-constants} holds for the measure $\eta$ with this constant $C$ and every polynomial $p$,
and Lemma~\ref{lem:two-constants} gives
\[
\eta\left(\left\{z\in\overline{\mathbb{D}}:|1+z|^2<4|z|\right\}\right)=0 .
\]

\emph{Conclusion.}
Let $(t,u)\in[0,1]\times S^1$ with $t>0$.
Then $|tu|=t$, $|1+tu|^2=1+2t\cdot\Re u+t^2$ and $|1-u|^2=2-2\Re u$, so that
\[
|1+tu|^2-4|tu|=1-4t+t^2+2t\cdot\Re u=(1-t)^2-t\cdot|1-u|^2=f(t,u).
\]
Consider the set
\[
S_{-}:=\left\{(t,u)\in[0,1]\times S^1:\ t>0,\ f(t,u)<0\right\}.
\]
By the identity just displayed, the map $(t,u)\mapsto tu$ sends $S_{-}$ into $\left\{z\in\overline{\mathbb{D}}:|1+z|^2<4|z|\right\}$,
a null set of $\eta$; as $\eta$ is the image of $t^{2r+1}\,\mathrm{d}\mathfrak{e}_{[G]_\mu}(t,u)$ under this map,
\[
\int_{S_{-}}\mathrm{d}\mathfrak{e}_{[G]_\mu}(t,u)\,t^{2r+1}\leq\eta\left(\left\{z\in\overline{\mathbb{D}}:|1+z|^2<4|z|\right\}\right)=0 .
\]
The factor $t^{2r+1}$ is strictly positive on $S_{-}$, so $\mathfrak{e}_{[G]_\mu}(S_{-})=0$.
Since $\mathfrak{E}_\mu(S_{-})$ is an orthogonal projection,
\[
\left\|\mathfrak{E}_\mu(S_{-})[G]_\mu\right\|^2=\left\langle[G]_\mu,\mathfrak{E}_\mu(S_{-})[G]_\mu\right\rangle_\mu=\mathfrak{e}_{[G]_\mu}(S_{-})=0 .
\]
Thus $\mathfrak{E}_\mu(S_{-})[G]_\mu=0$, and this holds for every $G\in\mathcal{O}^{+}_{\mu}$.
By the construction of $\mathcal{H}_\mu$ as a completion, the classes $[G]_\mu$, $G\in\mathcal{O}^{+}_{\mu}$, form a dense subset of $\mathcal{H}_\mu$;
the bounded operator $\mathfrak{E}_\mu(S_{-})$ vanishes on this subset, and hence $\mathfrak{E}_\mu(S_{-})=0$.
At $t=0$ one has $f(0,u)=1$ for every $u\in S^1$, so the set $\{f<0\}$ has no point with $t=0$ and equals $S_{-}$.
Consequently $\mathfrak{E}_\mu(\{f<0\})=0$, which is the first assertion of (b).
The equivalence established at the beginning of the proof turns it into the inequality of (a),
and the second assertion of (b) was derived there from the first.
\end{proof}

\begin{proof}[Proof of Theorem~\ref{thm:lattice-cone}(c)]
Let $\mu'\in\{1,2,3\}$ be either of the two indices different from $\mu$.
Parts~(a) and~(b) hold for the pair $(\mu,\mu')$,
and $\mathcal{H}_\mu$, $[\cdot]_\mu$, $\langle\cdot,\cdot\rangle_\mu$, $T_\mu$, $U_{\mu,\mu'}$, $\mathfrak{E}_\mu$
and the scalar measures $\mathfrak{e}_\psi=\left\langle\psi,\mathfrak{E}_\mu(\cdot)\psi\right\rangle_\mu$, $\psi\in\mathcal{H}_\mu$,
denote the objects of the axis picture of this pair (Definition~\ref{def:axis-picture}).
Let $P_\mu$ be the orthogonal projection of $L^2$ onto the closed subspace
\[
\left\{H\in L^2:\ \tau_{e_\mu}H=H\right\},
\]
and let
\[
\widehat{P}_\mu:=\mathfrak{E}_\mu\left(\{1\}\times S^1\right)
\]
be the spectral projection of $T_\mu$ onto $\ker(1-T_\mu)$.

For every $\psi\in\mathcal{H}_\mu$ one has
\[
\left\|T_\mu^n\psi-\widehat{P}_\mu\psi\right\|^2
=\int\mathrm{d}\mathfrak{e}_\psi(t,u)\,\left|t^n-1_{\{1\}}(t)\right|^2
\longrightarrow0\qquad(n\to\infty)
\]
by dominated convergence:
the integrand is at most one and tends to zero at every point of $[0,1]\times S^1$,
and $\mathfrak{e}_\psi$ is a finite measure.
Thus $T_\mu^n\to\widehat{P}_\mu$ strongly.
By part~(b), $\mathfrak{E}_\mu$ is carried by the set of the $(t,u)$ with $(1-t)^2\ge t\cdot|1-u|^2$.
At $t=1$ this inequality reads $|1-u|^2\le0$,
so $\mathfrak{E}_\mu$ gives no mass to $\{1\}\times\left(S^1\setminus\{1\}\right)$,
and the functions $u\cdot1_{\{1\}}(t)$ and $1_{\{1\}}(t)$ agree $\mathfrak{E}_\mu$-almost everywhere.
Since $U_{\mu,\mu'}$ and $\widehat{P}_\mu$ are the integrals of $u$ and of $1_{\{1\}}(t)$ against $\mathfrak{E}_\mu$,
\[
U_{\mu,\mu'}\widehat{P}_\mu=\widehat{P}_\mu U_{\mu,\mu'}=\widehat{P}_\mu .
\]
In the variables of part~(b), with $t=\exp(-\lambda)$, zero energy $\lambda=0$ thus forces $u=1$, that is, zero momentum.

\emph{Step~1: the projection $\widehat{P}_\mu$ seen in $L^2$.}
For all $G,\Phi\in\mathcal{O}^{+}_{\mu}$,
\begin{equation}\label{eq:layers-zero-energy}
\left\langle[G]_\mu,\widehat{P}_\mu[\Phi]_\mu\right\rangle_\mu
=\left\langle\theta_\mu G,P_\mu\Phi\right\rangle_2 .
\end{equation}
Indeed, the law $\nu$ is invariant under $W$, so that $\tau_{e_\mu}$ is a unitary operator of $L^2$,
and von Neumann's mean ergodic theorem \cite[Theorem~II.11]{ReedSimon1980} gives
\[
\frac{1}{n}\sum_{k=0}^{n-1}\tau_{ke_\mu}\Phi\longrightarrow P_\mu\Phi
\qquad\text{in }L^2\ (n\to\infty).
\]
On the other hand,
Lemma~\ref{lem:lattice-reconstruction}(d), read in the axis picture with the exponent $0$ on $U_{\mu,\mu'}$,
states that for every integer $k\ge0$
\[
\left\langle[G]_\mu,T_\mu^k[\Phi]_\mu\right\rangle_\mu
=\nu\left(\overline{\theta_\mu G}\cdot\tau_{ke_\mu}\Phi\right)
=\left\langle\theta_\mu G,\tau_{ke_\mu}\Phi\right\rangle_2 .
\]
Average these identities over $0\le k\le n-1$ and let $n\to\infty$.
On the left, the means of the vectors $T_\mu^k[\Phi]_\mu$ converge in $\mathcal{H}_\mu$ to $\widehat{P}_\mu[\Phi]_\mu$,
because they are the Ces\`aro means of a convergent sequence.
On the right, Lemma~\ref{lem:law-mirror}(a) gives $\left|\langle\theta_\mu G,H'\rangle_2\right|\le\|G\|_2\cdot\|H'\|_2$ for every $H'\in L^2$;
applied to the difference $H'$ between the average $n^{-1}\sum_{k=0}^{n-1}\tau_{ke_\mu}\Phi$ and its limit $P_\mu\Phi$,
it shows that the averaged right sides converge to $\langle\theta_\mu G,P_\mu\Phi\rangle_2$.
This gives~\eqref{eq:layers-zero-energy}.

\emph{Step~2: $\theta_\mu$ fixes every vector of $\operatorname{ran}P_\mu$.}
The reflection $\theta_\mu$ replaces $x_\mu$ by $-x_\mu-1$ and keeps the other coordinates;
hence $\theta_\mu\tau_{e_\mu}\theta_\mu=\tau_{-e_\mu}$ as maps of $\R^3$.
By the composition rule of~\eqref{eq:layers-action} the same identity holds between the operators of $L^2$
that these maps define, all of which are unitary because $\nu$ is invariant under $W$;
as $\theta_\mu^2=1$ it may be written $\tau_{e_\mu}\theta_\mu=\theta_\mu\tau_{-e_\mu}$.
If $\tau_{e_\mu}H=H$, then $\tau_{-e_\mu}H=H$, and so $\tau_{e_\mu}\theta_\mu H=\theta_\mu H$.
Thus $\theta_\mu$ maps $\operatorname{ran}P_\mu$ into itself, and onto itself because $\theta_\mu^2=1$.
A unitary operator that maps a closed subspace onto itself maps the orthogonal complement onto itself as well;
consequently $\theta_\mu$ commutes with $P_\mu$,
and the restriction $\theta_\mu|_{\operatorname{ran}P_\mu}$ is a unitary involution.
It is its own inverse, hence its own adjoint: $\theta_\mu|_{\operatorname{ran}P_\mu}$ is self-adjoint.

Now let $H\in\operatorname{ran}P_\mu$.
By the density~\eqref{eq:layers-density} there are $G_k\in\mathcal{O}(\R^3)$ with $\|H-G_k\|_2\to0$.
Each $G_k$ depends on finitely many $F_\sigma$;
by~\eqref{eq:layers-local} there are integers $n_k\ge0$ such that
\[
G'_k:=\tau_{n_ke_\mu}G_k\in\mathcal{O}^{+}_{\mu}.
\]
As $\tau_{e_\mu}H=H$ and $\tau_{n_ke_\mu}$ is unitary,
\[
\left\|H-G'_k\right\|_2=\left\|\tau_{n_ke_\mu}\left(H-G_k\right)\right\|_2=\left\|H-G_k\right\|_2\longrightarrow0 .
\]
Since $G'_k\in\mathcal{O}^{+}_{\mu}$, the reflection positivity of $\nu$ at the mirror $(x_\mu,-\tfrac{1}{2})$
(Definition~\ref{def:reflection-positivity}) gives $\nu\left(\overline{\theta_\mu G'_k}\cdot G'_k\right)\ge0$ for every $k$.
Moreover
\[
\overline{\theta_\mu H}\cdot H-\overline{\theta_\mu G'_k}\cdot G'_k
=\overline{\theta_\mu\left(H-G'_k\right)}\cdot H+\overline{\theta_\mu G'_k}\cdot\left(H-G'_k\right),
\]
and Lemma~\ref{lem:law-mirror}(a) bounds the integrals of the two terms on the right by
$\|H-G'_k\|_2\cdot\|H\|_2$ and by $\|G'_k\|_2\cdot\|H-G'_k\|_2$.
Therefore
\[
\left\langle\theta_\mu H,H\right\rangle_2
=\nu\left(\overline{\theta_\mu H}\cdot H\right)
=\lim_{k\to\infty}\nu\left(\overline{\theta_\mu G'_k}\cdot G'_k\right)\ge0 .
\]
So the self-adjoint involution $\theta_\mu|_{\operatorname{ran}P_\mu}$ is non-negative.
Its spectrum lies in $\{1,-1\}\cap[0,\infty)=\{1\}$, and it is the identity.

\emph{Step~3: invariance under $\tau_{e_{\mu'}}$.}
Let $G,\Phi\in\mathcal{O}^{+}_{\mu}$.
From~\eqref{eq:layers-zero-energy}, the relations $P_\mu=P_\mu^*=P_\mu^2$ and Step~2,
\[
\left\langle[G]_\mu,\widehat{P}_\mu[\Phi]_\mu\right\rangle_\mu
=\left\langle P_\mu\theta_\mu G,P_\mu\Phi\right\rangle_2
=\left\langle\theta_\mu P_\mu G,P_\mu\Phi\right\rangle_2
=\left\langle P_\mu G,P_\mu\Phi\right\rangle_2 .
\]
For $G=\Phi$ the left side is the square of the norm of $\widehat{P}_\mu[\Phi]_\mu$ in $\mathcal{H}_\mu$,
since $\widehat{P}_\mu$ is an orthogonal projection; hence
\[
\left\|\widehat{P}_\mu[\Phi]_\mu\right\|=\left\|P_\mu\Phi\right\|_2
\qquad\left(\Phi\in\mathcal{O}^{+}_{\mu}\right).
\]
The maps $\Phi\mapsto[\Phi]_\mu$ and $\Phi\mapsto P_\mu\Phi$ are linear.
It follows that the assignment $\widehat{P}_\mu[\Phi]_\mu\mapsto P_\mu\Phi$ is well defined,
and that it is a linear isometry of the subspace
$\left\{\widehat{P}_\mu[\Phi]_\mu:\ \Phi\in\mathcal{O}^{+}_{\mu}\right\}$ of $\mathcal{H}_\mu$ into $L^2$.

The translation $\tau_{e_{\mu'}}$ is a unitary operator of $L^2$ that commutes with $\tau_{e_\mu}$.
It maps $\operatorname{ran}P_\mu$ and its orthogonal complement onto themselves, and so it commutes with $P_\mu$.
Let $\Phi\in\mathcal{O}^{+}_{\mu}$.
The half-space $\{x_\mu>-\tfrac{1}{2}\}$ is invariant under the translation by $e_{\mu'}$;
by~\eqref{eq:layers-local}, $\tau_{e_{\mu'}}\Phi\in\mathcal{O}^{+}_{\mu}$,
and $\left[\tau_{e_{\mu'}}\Phi\right]_\mu=U_{\mu,\mu'}[\Phi]_\mu$ by Lemma~\ref{lem:lattice-reconstruction}(c).
With $\widehat{P}_\mu U_{\mu,\mu'}=\widehat{P}_\mu$ one obtains
\[
\widehat{P}_\mu\left[\tau_{e_{\mu'}}\Phi\right]_\mu
=\widehat{P}_\mu U_{\mu,\mu'}[\Phi]_\mu
=\widehat{P}_\mu[\Phi]_\mu .
\]
The isometry maps the two sides to $P_\mu\tau_{e_{\mu'}}\Phi$ and to $P_\mu\Phi$, respectively; therefore
\[
\tau_{e_{\mu'}}P_\mu\Phi=P_\mu\tau_{e_{\mu'}}\Phi=P_\mu\Phi .
\]
The vectors $P_\mu\Phi$ with $\Phi\in\mathcal{O}^{+}_{\mu}$ are dense in $\operatorname{ran}P_\mu$:
for $H\in\operatorname{ran}P_\mu$ and the $G'_k$ of Step~2, $P_\mu G'_k\to P_\mu H=H$.
By continuity of $\tau_{e_{\mu'}}$, every $H\in\operatorname{ran}P_\mu$ satisfies $\tau_{e_{\mu'}}H=H$.

Finally, let $H\in L^2$ satisfy $\tau_{e_\mu}H=H$, that is, $H\in\operatorname{ran}P_\mu$.
Step~3 applies to each of the two indices $\mu'\neq\mu$.
Hence $H$ is invariant under $\tau_{e_1}$, $\tau_{e_2}$ and $\tau_{e_3}$,
and therefore under the group generated by them, which consists of the $\tau_v$ with $v\in\Z^3$.
\end{proof}

\begin{proof}[Proof of Theorem~\ref{thm:lattice-cone}(d)]
Denote by $P_{\Z^3}$ the orthogonal projection of $L^2$ onto the closed subspace of the vectors $H$
with $\tau_vH=H$ for all $v\in\Z^3$.

\emph{Invariant vectors.}
The range of $P_{\Z^3}$ consists of the classes of the $\mathcal{B}^{\mathrm{inv}}$-measurable functions in $L^2$,
and $P_{\Z^3}$ is the conditional expectation $\nu\left(\,\cdot\mid\mathcal{B}^{\mathrm{inv}}\right)$.
On the one hand, the level sets of a $\mathcal{B}^{\mathrm{inv}}$-measurable function $H$ are invariant under every $\tau_v^*$;
so $H(\tau_v^*F)=H(F)$ for all $F\in\Omega$ and $v\in\Z^3$, that is, $\tau_vH=H$.
On the other hand, let $H$ be a Borel function such that, for each $v\in\Z^3$, $\tau_vH=H$ holds $\nu$-almost everywhere, and put
\[
\Omega_H:=\left\{F\in\Omega:\ H\left(\tau_v^*F\right)=H(F)\ \text{for all }v\in\Z^3\right\}.
\]
Since $\Z^3$ is countable, $\Omega_H$ is a Borel set of full measure.
It is invariant: for $F\in\Omega_H$ and $v,v'\in\Z^3$,
\[
H\left(\tau_v^*\tau_{v'}^*F\right)=H\left(\tau_{v'+v}^*F\right)=H(F)=H\left(\tau_{v'}^*F\right),
\]
so that $\tau_{v'}^*F\in\Omega_H$;
with $-v'$ in the place of $v'$ this gives $\tau_{v'}^*\Omega_H=\Omega_H$.
Consequently the function $H\cdot1_{\Omega_H}$ is invariant under every $\tau_v^*$ at every point of $\Omega$.
It is therefore $\mathcal{B}^{\mathrm{inv}}$-measurable, and it equals $H$ almost everywhere.
Hence $\operatorname{ran}P_{\Z^3}$ is the space of the classes of the square integrable $\mathcal{B}^{\mathrm{inv}}$-measurable functions,
and the orthogonal projection onto this space is the conditional expectation given $\mathcal{B}^{\mathrm{inv}}$.

\emph{Regular conditional distributions.}
Because $\Omega$ is Polish (see~\eqref{eq:layers-law}),
regular conditional distributions of $\nu$ given $\mathcal{B}^{\mathrm{inv}}$ exist
\cite[Theorem~10.2.2]{Dudley2002}, \cite[Theorem~6.3]{Kallenberg2002}.
For the distribution $F\mapsto\nu_F$ of part~(d) this means the following:
$\nu_F$ is a law on $\Omega$ for every $F$;
$F\mapsto\nu_F(B)$ is $\mathcal{B}^{\mathrm{inv}}$-measurable for every Borel set $B$;
and for every bounded Borel function $G$ on $\Omega$
\[
\nu_F(G)=\left(P_{\Z^3}G\right)(F)\quad\text{for $\nu$-almost every }F,
\qquad
\nu(G)=\int\mathrm{d}\nu(F)\,\nu_F(G).
\]

\emph{Step~1: every reflection fixes $\operatorname{ran}P_{\Z^3}$ pointwise.}
Let $\Pi=(\ell,c)$ be a mirror.
For $v\in\Z^3$ the map $\theta_\Pi\tau_v\theta_\Pi$ is the translation by the image of $v$ under the linear part of $\theta_\Pi$,
which is again a vector of $\Z^3$.
If $H\in\operatorname{ran}P_{\Z^3}$, then for every $v\in\Z^3$
\[
\tau_v\theta_\Pi H=\theta_\Pi\left(\theta_\Pi\tau_v\theta_\Pi\right)H=\theta_\Pi H .
\]
So the unitary operator $\theta_\Pi$ of $L^2$ (Lemma~\ref{lem:law-mirror}(a)) maps $\operatorname{ran}P_{\Z^3}$ into itself,
and onto itself since $\theta_\Pi^2=1$.
Being unitary, $\theta_\Pi$ then preserves the orthogonal complement of this range too,
and consequently $\theta_\Pi P_{\Z^3}=P_{\Z^3}\theta_\Pi$.
On $\operatorname{ran}P_{\Z^3}$ the operator $\theta_\Pi$ acts as a unitary map that coincides with its own inverse;
as the inverse of a unitary map is its adjoint, the restriction of $\theta_\Pi$ to $\operatorname{ran}P_{\Z^3}$ is self-adjoint.

Let $H\in\operatorname{ran}P_{\Z^3}$.
Choose $G_k\in\mathcal{O}(\R^3)$ with $\|H-G_k\|_2\to0$, by~\eqref{eq:layers-density},
and vectors $v_k\in\Z^3$ with $\ell(v_k)$ so large that, by~\eqref{eq:layers-local},
\[
G'_k:=\tau_{v_k}G_k\in\mathcal{O}^{+}(\ell,c).
\]
Since $\tau_{v_k}H=H$,
\[
\left\|H-G'_k\right\|_2=\left\|\tau_{v_k}\left(H-G_k\right)\right\|_2=\left\|H-G_k\right\|_2\longrightarrow0 .
\]
Lemma~\ref{lem:law-mirror}(a), applied to the two terms on the right of
\[
\overline{\theta_\Pi H}\cdot H-\overline{\theta_\Pi G'_k}\cdot G'_k
=\overline{\theta_\Pi\left(H-G'_k\right)}\cdot H+\overline{\theta_\Pi G'_k}\cdot\left(H-G'_k\right),
\]
gives
\[
\left|\nu\left(\overline{\theta_\Pi H}\cdot H\right)-\nu\left(\overline{\theta_\Pi G'_k}\cdot G'_k\right)\right|
\le\left\|H-G'_k\right\|_2\cdot\left(\|H\|_2+\|G'_k\|_2\right)\longrightarrow0 ,
\]
while $\nu\left(\overline{\theta_\Pi G'_k}\cdot G'_k\right)\ge0$ for every $k$,
by the reflection positivity of $\nu$ at $\Pi$ (Definition~\ref{def:reflection-positivity}).
Thus
\[
\left\langle\theta_\Pi H,H\right\rangle_2=\nu\left(\overline{\theta_\Pi H}\cdot H\right)\ge0 ,
\]
and since $H\in\operatorname{ran}P_{\Z^3}$ was arbitrary, the restriction of $\theta_\Pi$ to $\operatorname{ran}P_{\Z^3}$ is a non-negative operator.
Its square is the identity, which confines its spectrum to the set $\{1,-1\}$,
and non-negativity excludes the point $-1$.
By the spectral theorem, a self-adjoint operator with spectrum contained in $\{1\}$ is the identity.
Hence $\theta_\Pi H=H$ for every $H\in\operatorname{ran}P_{\Z^3}$.

\emph{Step~2: invariance of the components under $W$.}
Let $g$ be one of the reflections $\theta_\Pi$, of which there are countably many,
and let $G$ belong to the countable set $\mathcal{T}_{\Q}$ of Lemma~\ref{lem:characters-determine}(b).
By Step~1, $g$ commutes with $P_{\Z^3}$ and acts as the identity on its range.
Therefore, for $\nu$-almost every $F$,
\[
\nu_F(gG)=\left(P_{\Z^3}gG\right)(F)=\left(gP_{\Z^3}G\right)(F)=\left(P_{\Z^3}G\right)(F)=\nu_F(G).
\]
There is one $\nu$-null set outside of which this identity holds for all pairs $(g,G)$ at once.
Let $F$ lie outside this null set.
For every $\beta\in C^2_{\mathrm{c}}$ with rational values the character $\varphi_\beta$ lies in $\mathcal{T}_{\Q}$,
so that $\nu_F(g\varphi_\beta)=\nu_F(\varphi_\beta)$.
An arbitrary $\beta\in C^2_{\mathrm{c}}$ is the pointwise limit of rational-valued $\beta_n\in C^2_{\mathrm{c}}$ supported in the support of $\beta$.
Then $\varphi_{\beta_n}\to\varphi_\beta$ and $g\varphi_{\beta_n}\to g\varphi_\beta$ at every point of $\Omega$,
all these functions have modulus one,
and dominated convergence yields $\nu_F(g\varphi_\beta)=\nu_F(\varphi_\beta)$ for all $\beta\in C^2_{\mathrm{c}}$.
Thus the image of $\nu_F$ under $g^*$ and $\nu_F$ agree on every character, hence on every trigonometric polynomial,
and Lemma~\ref{lem:characters-determine}(a) shows that they are equal: $\nu_F$ is invariant under $g$.
The group $W$ is generated by the reflections $\theta_\Pi$ (Definition~\ref{def:mirrors}),
and by the composition rule of~\eqref{eq:layers-action} a law that is invariant under $g$ and under $g'$ is invariant under $gg'$.
So for every $F$ outside this null set the law $\nu_F$ is invariant under $W$,
in particular under the translations $\tau_v$, $v\in\Z^3$.

\emph{Step~3: ergodicity of the components.}
For $r=1,2,\dots$ let
\[
\operatorname{Av}_r:=\frac{1}{r^3}\sum_{v\in\{0,\dots,r-1\}^3}\tau_v ;
\]
it is the product of the three one-dimensional means $r^{-1}\sum_{k=0}^{r-1}\tau_{ke_j}$, $j=1,2,3$.
Let $\mu$ be a law on $\Omega$ that is invariant under all $\tau_v$, $v\in\Z^3$.
On $L^2(\mu)$ the three one-dimensional means are contractions that commute with each other.
By the mean ergodic theorem \cite[Theorem~II.11]{ReedSimon1980} they converge strongly, as $r\to\infty$,
to the orthogonal projections onto the vectors that are invariant under $\tau_{e_1}$, $\tau_{e_2}$, $\tau_{e_3}$, respectively;
these projections commute, being strong limits of operators that commute.
A product of uniformly bounded, strongly convergent sequences of operators converges strongly to the product of the limits,
and a product of commuting orthogonal projections is the orthogonal projection onto the intersection of their ranges.
Hence $\operatorname{Av}_r$ converges strongly on $L^2(\mu)$ to the orthogonal projection $P^{\mu}_{\Z^3}$
onto the vectors that are invariant under all $\tau_v$, $v\in\Z^3$.

For $G\in\mathcal{T}_{\Q}$ let $\widetilde{G}$ be a bounded $\mathcal{B}^{\mathrm{inv}}$-measurable representative of $P_{\Z^3}G$
(for instance the function $F\mapsto\nu_F(G)$).
The disintegration $\nu(\cdot)=\int\mathrm{d}\nu(F)\,\nu_F(\cdot)$, applied to the bounded Borel function
$\left|\operatorname{Av}_rG-\widetilde{G}\right|^2$, and the preceding paragraph for the law $\nu$ give
\[
\int\mathrm{d}\nu(F)\,\left\|\operatorname{Av}_rG-\widetilde{G}\right\|^2_{L^2(\nu_F)}
=\left\|\operatorname{Av}_rG-\widetilde{G}\right\|_2^2\longrightarrow0
\qquad(r\to\infty).
\]
Enumerate $\mathcal{T}_{\Q}=\left\{G^{(1)},G^{(2)},\dots\right\}$,
let $\widetilde{G}^{(n)}$ be the representative chosen for $G^{(n)}$,
and choose integers $r_1<r_2<\cdots$ such that
\[
\left\|\operatorname{Av}_{r_k}G^{(n)}-\widetilde{G}^{(n)}\right\|_2^2\le2^{-k}
\qquad(1\le n\le k).
\]
For fixed $n$, the non-negative functions $F\mapsto\left\|\operatorname{Av}_{r_k}G^{(n)}-\widetilde{G}^{(n)}\right\|^2_{L^2(\nu_F)}$, $k\ge n$,
have $\nu$-integrals with a finite sum.
By monotone convergence their sum is finite for $\nu$-almost every $F$, so that they tend to zero there.
Taking the union of countably many null sets,
one finds a $\nu$-null set outside of which $\operatorname{Av}_{r_k}G\to\widetilde{G}$ in $L^2(\nu_F)$ for every $G\in\mathcal{T}_{\Q}$.

It remains to show that $\widetilde{G}$ is $\nu_F$-almost surely constant.
Let $a\in\Q+i\Q$ and $G\in\mathcal{T}_{\Q}$.
Since $\widetilde{G}$ is bounded and $\mathcal{B}^{\mathrm{inv}}$-measurable, so is $\left|\widetilde{G}-a\right|^2$;
in particular $P_{\Z^3}$ leaves this function unchanged, and the defining property of $F\mapsto\nu_F$ gives
\[
\nu_F\left(\left|\widetilde{G}-a\right|^2\right)=\left|\widetilde{G}(F)-a\right|^2
\]
for $\nu$-almost every $F$.
Only countably many pairs $(a,G)$ occur, and the union of the corresponding exceptional sets is again a $\nu$-null set.
Fix $F$ outside this union.
Expanding the square, one finds for every $a\in\C$
\[
\nu_F\left(\left|\widetilde{G}-a\right|^2\right)
=\nu_F\left(\left|\widetilde{G}\right|^2\right)-2\Re\left(\overline{a}\cdot\nu_F\left(\widetilde{G}\right)\right)+|a|^2 ;
\]
thus $\nu_F\left(\left|\widetilde{G}-a\right|^2\right)$, like $\left|\widetilde{G}(F)-a\right|^2$, depends continuously on $a\in\C$.
These two functions of $a$ coincide on the dense subset $\Q+i\Q$ of $\C$, hence on all of $\C$,
and the choice $a=\widetilde{G}(F)$ yields $\nu_F\left(\left|\widetilde{G}-\widetilde{G}(F)\right|^2\right)=0$:
the function $\widetilde{G}$ equals the constant $\widetilde{G}(F)$ $\nu_F$-almost surely.

Outside the union of the null sets found above, the law $\nu_F$ is invariant under the translations (Step~2),
and for every $G\in\mathcal{T}_{\Q}$
\[
P^{\nu_F}_{\Z^3}G=\lim_{k\to\infty}\operatorname{Av}_{r_k}G=\widetilde{G}=\widetilde{G}(F)
\qquad\text{in }L^2(\nu_F),
\]
a constant.
The set $\mathcal{T}_{\Q}$ is dense in $L^2(\nu_F)$ by Lemma~\ref{lem:characters-determine}(b),
and $P^{\nu_F}_{\Z^3}$ is continuous; so its range consists of constants.
Thus the vectors of $L^2(\nu_F)$ that are invariant under all $\tau_v$ are the constants.
In particular, the indicator of a Borel set that is invariant under all $\tau_v^*$ is $\nu_F$-almost surely constant,
so such a set has $\nu_F$-measure $0$ or $1$: the law $\nu_F$ is ergodic under $\Z^3$.
\end{proof}

\section{Arbitrary sides}\label{sec:sides}

This section contains the proof of Theorem~\ref{thm:cubic-limits}, for an arbitrary sequence of integers $n_i\ge 4$ with $n_i\to\infty$, even or odd.
Throughout, the hopping weights $w_0,\dots,w_3>0$, the mass parameter $M\in\R$ and the lattice charge $q>0$ are arbitrary,
and the Wilson parameter satisfies $-1\le\rho\le 1$ unless a statement names another range of $\rho$.

Condition~\eqref{cond:divisibility} is imposed on the sides of the tori that carry the comparison states,
since the cell structures of the data, and coarse structures whose widths grow without bound, have to fit these tori (Definition~\ref{def:data}, Lemma~\ref{lem:coarse-structures}\eqref{it:coarse-widths}).
The proof therefore forms the adapted tori, the data and the comparison states with a second sequence of sides, which satisfies~\eqref{cond:divisibility},
and the cubic tori $C(n_i)$ and their laws with the given sides.

In the argument that leads to Theorem~\ref{thm:cubic-limits}, the cubic torus $C(n)=C(n,n,n)$ occurs in two roles.
In the first role it is the torus whose law $\mu_{C(n)}$ is studied, and as such it enters at four places only:
\begin{enumerate}
\item\label{it:sides-jensen}
Jensen's inequality, applied to its law in Proposition~\ref{prop:fold-finite}\,(a);
\item\label{it:sides-trace}
the trace inequality of Theorem~\ref{thm:chessboard}\,(b), which carries the unfolded array from $C(n)$ to the adapted torus $D(n)$ in Proposition~\ref{prop:fold-finite}\,(b);
\item\label{it:sides-pressure}
the lower bound $f(n,n,n)\ge f_\infty$ for its pressure, the pressure of $C(n)$ entering Proposition~\ref{prop:fold-finite}\,(d) through Theorem~\ref{thm:four-tori};
\item\label{it:sides-symmetry}
the symmetries of $C(n)$ itself, Gaussian domination on $C(n)$ and reflection positivity at the coordinate mirrors of $C(n)$,
in the step~\ref{it:sides-limit-invariant} of the proof of Lemma~\ref{lem:diagonal-domination}
and in the steps~\ref{it:sides-tightness}--\ref{it:sides-axis} of the proof of Theorem~\ref{thm:cubic-limits}.
\end{enumerate}

In the second role the cubic torus is an auxiliary torus whose side is the side of the adapted torus.
For an even side $n\ge 8$, Theorem~\ref{thm:four-tori} gives $Z_{C(n)}[1]\le Z_{D(n)}[1]$, and therefore
\[
\frac{Z_{C(n/2,2n,n)}[1]}{Z_{D(n)}[1]}
\le\frac{Z_{C(n/2,2n,n)}[1]}{Z_{C(n)}[1]}
=\exp\left(n^3\cdot\left(f(n/2,2n,n)-f(n,n,n)\right)\right);
\]
for the sides $n_i$ of a sequence with \textup{(Div)} the right-hand side is $\exp\left(n_i^3\cdot\Delta_i\right)$, and this bound is step~\ref{it:amplification-ratio} of the proof of Theorem~\ref{thm:amplification}.
Through this step the auxiliary torus enters Propositions~\ref{prop:domination-tightness} and~\ref{prop:a-priori} as well as Theorems~\ref{thm:two-layer-principle}, \ref{thm:two-layer-reduction} and~\ref{thm:no-layers}.
Its side is then always an even number, namely the side of the adapted torus.
In the present section the adapted tori are formed with the second sequence $\mathfrak n'=(n'_i)$ (Definition~\ref{def:two-sequences}), and so is this auxiliary torus:
it is $C(n'_i)$, of the even side $n'_i$.

\paragraph{Every integer side.}
Any integer $n_i\ge 4$ may serve as a given side, with no condition on its parity or on its divisors.
The bound $n_i\ge 4$ is precisely the condition for $C(n_i)$ to be a torus:
the period lattice $n_i\Z^3$ has the systole $n_i$, and the systole of a period lattice is at least $4$ (Definitions~\ref{def:configurations} and~\ref{def:tori}).
At the places~\ref{it:sides-jensen}, \ref{it:sides-pressure} and~\ref{it:sides-symmetry} every integer side $n\ge 4$ is admissible.
For the places~\ref{it:sides-jensen} and~\ref{it:sides-symmetry} the reason is twofold.
First, Proposition~\ref{prop:symmetries}\eqref{it:charge-conjugation}, \eqref{it:functional-covariance}, Proposition~\ref{prop:accretivity} and Theorem~\ref{thm:coordinate-domination} are valid on every coordinate torus with three integer sides $\ge 4$.
Second, $C(n)=\mathbb{T}(\mathfrak{F}_{\mathrm{c}};n,n,n)$ is adapted to each axis family $x_\mu$, with the height period $n$ (Definitions~\ref{def:slabs} and~\ref{def:cell-structures}),
and for a position $c\in\tfrac12\Z=\Adm(x_\mu)$ (Definition~\ref{def:mirrors}) the position $c+n/2$ belongs to $\tfrac12\Z$ as well, whatever the parity of $n$;
if $n$ is odd, the mirrors $(x_\mu,c)$ and $(x_\mu,c+n/2)$ are of different kinds, a site plane and a link plane.
Corollary~\ref{cor:trace-forms}\eqref{it:trace-form} is therefore available on $C(n)$ for each such pair of mirrors.
At the place~\ref{it:sides-pressure}, Proposition~\ref{prop:coordinate-pressure} holds for all coordinate tori $C(a_1,a_2,a_3)$ with integers $a_\mu\ge 4$, odd sides included,
by virtue of Corollary~\ref{cor:trace-forms}\eqref{it:unit-step}:
for $c\in\tfrac12\Z$ and the half step $\mathfrak{S}_{1/2}$, which leads from the plane $\{x_1=c\}$ to the plane $\{x_1=c+\tfrac12\}$,
\[
Z_{C(a_1,a_2,a_3)}[1]=\Tr_{\mathcal{H}(c)}\left(\left(\mathfrak{S}_{1/2}^{*}\mathfrak{S}_{1/2}\right)^{a_1}\right).
\]
Both planes are mirrors of the family $x_1$, one of them a site plane and one a link plane.

The place~\ref{it:sides-trace} is of another kind.
In the earlier sections the parity of $n$ enters at this place through the straightened torus $\mathbb{T}^0(\ell_-)=D'(n/2,n,n)$ of $C(n)$ (Definition~\ref{def:slabs}, Proposition~\ref{prop:untwistings}),
whose lattice contains $(n/2)\,e_-$, and also through the tori $D(n/2,n,n)$ and $C(n/2,2n,n)$, which are built from that torus.
If $n$ is odd, then the height period $n$ of $C(n)$ for $\ell_-$ is not a multiple of $|n_{\ell_-}|^2=2$, and $C(n)$ has no straightened torus for $\ell_-$ (Lemma~\ref{lem:tori-over-section}\,(c)).

The trace inequality is derived from the H\"older inequality for slab operators in the Schatten classes (Proposition~\ref{prop:twisted-holder}),
and Proposition~\ref{prop:unequal-slabs} establishes that inequality for slabs of unequal heights.
There, for a period divided into $2k$ slabs, the slab of height $h_j$ is represented by a partition function on a torus $\mathbb{T}_j$ adapted to the same family,
with the same cross-section and with the height period $2kh_j$, raised to the power $\frac{1}{2k}$.
For the torus $\mathbb{T}$ whose functional is estimated and the family $\ell$ of the slabs, the tori $\mathbb{T}_j$ take the place that the straightened torus $\mathbb{T}^0(\ell)$ has in Proposition~\ref{prop:twisted-holder},
and they coincide with it when all heights are equal (Proposition~\ref{prop:unequal-slabs}\,(c)).
Applied to $\mathbb{T}=C(n)$ and $\ell=\ell_-$, the proposition involves the cross-section $\Xi(C(n),\ell_-)=(\Gamma_\parallel,\chi_\parallel)$, given by
\[
\Gamma_\parallel=\Z\,n e_+ +\Z\,n e_3,\qquad \chi_\parallel(n e_+)=1,\qquad \chi_\parallel(n e_3)=-1,
\]
the height period $n$ and the closing vector $n e_1$, at which it is attained;
all three are defined for every integer $n\ge 4$.
The tori $\mathbb{T}_j$ of the proposition then have the lattices $\Gamma_\parallel+\Z\,kh_j e_-$, the character $\chi_\parallel$ on $\Gamma_\parallel$ and the value $-1$ at $kh_j e_-$,
and their closing vectors $kh_j e_-$ are integer vectors for all integers $h_j$.

Let $\mathfrak d=(c_0,c_+,c_3;h;O)$ be data (Definition~\ref{def:data}), with the width $h$, the observable $O$ and the cell structure $\mathcal{Q}_{\mathfrak d}$,
formed with a side $n'\le n$ for which the adapted torus $D(n')=D(n'/2,n',n')$ carries $\mathcal{Q}_{\mathfrak d}$;
apart from this condition the side $n'$ may be chosen at will.
Lemma~\ref{lem:one-step} applies Proposition~\ref{prop:unequal-slabs} in one direction $\ell^a$ of a frame to a period of height $p$,
divided into $2k-1$ slabs of the height $h$ and one last slab, of the remaining height $p-(2k-1)h$, which receives no factor of the integrand.
Starting from $C(n)$, the lemma is applied once in each of the directions $\ell_-$, $\ell_+$ and $x_3$ of the diagonal frame.
The translates of $O$ that lie in cells outside the three empty slabs are thereby carried into cells of $D(n')$,
while the empty slabs are collected on three exceptional tori.
Since the share of these tori in the volume of $C(n)$ is at most a constant multiple of $(h+n-n')/n$,
their partition functions may be estimated by the bound $Z_{\mathbb{T}}[1]\le\exp\left(C\cdot|\mathbb{T}|\right)$ of Lemma~\ref{lem:uniform-ceiling},
where the constant $C$ depends only on $\rho$, $M$, $w_0,\dots,w_3$ and $q$.
In the second and in the third application the side $n$ appears only in the lattices of $D'(n'/2,n,n)$ and of $D(n'/2,n',n)$, as a coefficient of $e_+$ or of $e_3$.
The outcome is the fold inequality of Proposition~\ref{prop:fold-any-side}, which compares $C(n)$, for every integer side $n\ge n'$, with the adapted torus $D(n')$.

\paragraph{The second sequence.}
For a given sequence $\mathfrak n=(n_i)$, the adapted tori receive their sides from the reduced sequence of $\mathfrak n$ (Definition~\ref{def:reduced-sequence}).
For $n_i\ge 64$ let $k(n_i)$ be the largest integer $k\ge 1$ with $8\,k!\le n_i^{\frac12}$,
and let $n'_i:=\red(n_i)$ be obtained by rounding $n_i$ down to a multiple of $8\,k(n_i)!$;
the finitely many indices with $n_i<64$ receive $n'_i:=8$.
As $i\to\infty$ the integers $k(n_i)$ tend to infinity,
while for every index with $n_i\ge 64$ the difference $n_i-n'_i$, a remainder modulo $8\,k(n_i)!\le n_i^{\frac12}$, satisfies
\[
n_i-n'_i<n_i^{\frac12}\qquad(n_i\ge 64)
\]
(Lemma~\ref{lem:companions-exist}\,(a)).
Whatever the given sides, every $n'_i$ with $n_i\ge 64$ is thus divisible by $8\,k(n_i)!$.
The sequence $\mathfrak n'=(n'_i)$ satisfies \textup{(Div)} (Lemma~\ref{lem:companions-exist}\,(b));
consequently Proposition~\ref{prop:wall-positivity}, Theorem~\ref{thm:comparison-states}\,(a), (b), the statements of Sections~\ref{sec:amplification} and~\ref{sec:elimination} and Theorem~\ref{thm:no-layers} hold for $\mathfrak n'$ without any change.
The following convention keeps the two sequences apart.

\begin{definition}[Two sequences of sides]\label{def:two-sequences}
Throughout this section two sequences of positive integers with one and the same index set are given.
\begin{enumerate}
\item
The sequence $\mathfrak n=(n_i)$ of the given sides consists of integers $n_i\ge 4$, even or odd, with $n_i\to\infty$, and
\[
\mu_i:=\mu_{C(n_i)}
\]
is the law of the cubic torus $C(n_i)=C(n_i,n_i,n_i)$ (Definition~\ref{def:functional}).
Wherever Definition~\ref{def:sides}, formed with $\mathfrak n$, defines a law $\mu_i$, the two laws coincide;
here, however, $n_i$ is an arbitrary integer $\ge 4$.
\item
The sequence $\mathfrak n'=(n'_i)$ satisfies \textup{(Div)}.
In Theorem~\ref{thm:comparison-any-side}, in Lemma~\ref{lem:diagonal-domination} and in the proof of Theorem~\ref{thm:cubic-limits} it is a companion of $\mathfrak n$ (Definition~\ref{def:companions}), which exists by Lemma~\ref{lem:companions-exist}.
\item
The adapted tori and every object built on them belong to $\mathfrak n'$:
the tori $D'(n'_i)$, $D(n'_i)$, $C(n'_i/2,2n'_i,n'_i)$ and the numbers $\Delta_i$ of Definition~\ref{def:sides},
the data $\mathfrak d$ with their cell structure $\mathcal{Q}_{\mathfrak d}$, base cell $Q_0$ and lattice $\Lambda_h$,
the folded arrays $O_{\mathfrak d}$ and the functionals $\mu^t_i$ of Definitions~\ref{def:arrays} and~\ref{def:data},
the threshold $i_0(\mathfrak d)$ of Definition~\ref{def:threshold},
the numbers $k(\mathfrak d,r)$ and $i_0(\mathfrak d,r)$ of Lemma~\ref{lem:coarse-structures},
the amplification factors $\gamma_{t,i}(k)$ of Section~\ref{sec:setting},
the comparison states and the notion of subordination of Definition~\ref{def:comparison-states}
are those formed with the sequence $\mathfrak n'$.
Thus
\[
\begin{aligned}
&D'(n'_i)=D'(n'_i/2,n'_i,n'_i),\qquad D(n'_i)=D(n'_i/2,n'_i,n'_i),\\
&\Delta_i=f(n'_i/2,2n'_i,n'_i)-f(n'_i,n'_i,n'_i),\qquad 4h\in\mathcal{D}(\mathfrak n'),\\
&\mu^t_i(G)=\frac{Z_{D(n'_i)}\left[\exp(t\cdot O_{\mathfrak d})\cdot G\right]}{Z_{D(n'_i)}\left[\exp(t\cdot O_{\mathfrak d})\right]},
\end{aligned}
\]
and for $i\ge i_0(\mathfrak d)$ the number of the cells of $\mathcal{Q}_{\mathfrak d}$ in $D(n'_i)$ is $(n'_i)^3/|Q_0|$.
\item
A statement of an earlier section that is formulated for a sequence of sides is applied to $\mathfrak n'$ by substituting $\mathfrak n'$ for that sequence.
So read, it speaks of the cubic tori $C(n'_i)$, whose laws receive no name;
in particular, a cubic torus $C(n_i)$ of such a statement becomes $C(n'_i)$.
In this section the cubic torus of a given side $n_i$ is always written with that side, as $C(n_i)$ or $C(n_i,n_i,n_i)$.
\item
With the sequence $\mathfrak n$ itself only the tori $C(n_i)$, their laws $\mu_i$ and laws $\lambda_i$ on $\Omega$ with $\lambda_i\ll\mu_i$ are formed.
\item
Laws are understood as in~\eqref{eq:layers-law}:
the laws are the Borel probability measures on $\Omega$, and for a law $\mu$ one writes $\mu(G)=\int\mathrm{d}\mu(F)\,G(F)$ for every Borel function $G$ on $\Omega$ that is non-negative or $\mu$-integrable.
A sequence of laws converges weakly to a law $\mu$ if, for every bounded continuous function $G$ on $\Omega$, the integrals of $G$ with respect to its terms converge to $\mu(G)$,
and a law to which some subsequence of a given sequence of laws converges weakly is called a weak limit point of that sequence.
The letters $\lambda$, $\lambda_i$ stand for laws on $\Omega$, and for laws $\lambda\ll\mu$ the relative entropy of Definition~\ref{def:relative-entropy} is
\[
H(\lambda\,|\,\mu)=\lambda\left(\log\frac{\mathrm{d}\lambda}{\mathrm{d}\mu}\right)\in[0,\infty].
\]
\item
As in~\eqref{eq:fold-frame}, the families of the diagonal frame are $(\ell^1,\ell^2,\ell^3)=(\ell_-,\ell_+,x_3)$;
their normals $n_a:=n_{\ell^a}$ are $n_{\ell_-}=e_-$, $n_{\ell_+}=e_+$, $n_{x_3}=e_3$, so that $|n_a|^2=2,2,1$ for $a=1,2,3$;
for data $\mathfrak d=(c_0,c_+,c_3;h;O)$ the walls of $\mathcal{Q}_{\mathfrak d}$ lie at $\ell^a=o_a+jh$, $j\in\Z$, with $(o_1,o_2,o_3)=(c_0,c_+,c_3)$.
\end{enumerate}
\end{definition}

\subsection{Slabs of unequal heights}\label{sec:unequal-heights}

For a family $\ell\in\mathcal{L}$ the number $|n_\ell|^2/2$ equals $1/2$ if $\ell$ is an axis family and $1$ if $\ell$ is a diagonal family.
By Definition~\ref{def:mirrors},
\[
\Adm(\ell)=\frac{|n_\ell|^2}{2}\,\Z ;
\]
in particular $\Adm(\ell)$ is a subgroup of $\R$.

\begin{lemma}[Adapted tori over a cross-section]\label{lem:tori-over-section}
Let $\ell\in\mathcal{L}$, and let $\Xi=(\Gamma_\parallel,\chi_\parallel)$ be a cross-section for $\ell$ (Definition~\ref{def:slabs}), with the area
\[
|\Xi|:=\left[\Z^3\cap\{\ell=0\}:\Gamma_\parallel\right].
\]
Let $P$ be a positive element of $|n_\ell|^2\,\Z$, and put
\[
u:=\frac{P}{|n_\ell|^2}\,n_\ell .
\]
\begin{enumerate}
\item[\textup{(a)}]
The vector $u$ lies in $\Z^3$, and $\ell(u)=P$.
\item[\textup{(b)}]
Assume that $\sys(\Gamma_\parallel+\Z u)\ge4$, and put
\[
\mathbb{T}(\Xi;P):=\left(\Gamma_\parallel+\Z u,\ \chi'\right),
\]
where $\chi'$ is the character of $\Gamma_\parallel+\Z u$ with $\chi'=\chi_\parallel$ on $\Gamma_\parallel$ and $\chi'(u)=-1$.
Then $\mathbb{T}(\Xi;P)$ is a torus, and it is adapted to $\ell$.
Moreover
\[
\Xi\left(\mathbb{T}(\Xi;P),\ell\right)=\Xi,\qquad
p\left(\mathbb{T}(\Xi;P),\ell\right)=P,\qquad
\left|\mathbb{T}(\Xi;P)\right|=P\cdot|\Xi|,
\]
and $u$ is a closing vector of $\mathbb{T}(\Xi;P)$ for $\ell$.
\item[\textup{(c)}]
Let $\mathbb{T}$ be a torus with $\Xi(\mathbb{T},\ell)=\Xi$, and let $p:=p(\mathbb{T},\ell)$.
Then $|\mathbb{T}|=p\cdot|\Xi|$.
The straightened torus $\mathbb{T}^0(\ell)$ is defined if and only if
\[
p\in|n_\ell|^2\,\Z
\qquad\text{and}\qquad
\sys\left(\Gamma_\parallel+\Z\,\frac{p}{|n_\ell|^2}\,n_\ell\right)\ge4 ,
\]
and in that case $\mathbb{T}^0(\ell)=\mathbb{T}(\Xi;p)$.
\end{enumerate}
\end{lemma}

\begin{proof}
(a) Since $P\in|n_\ell|^2\Z$, the quotient $P/|n_\ell|^2$ is an integer:
for an axis family $|n_\ell|^2=1$,
and for a diagonal family $|n_\ell|^2=2$ and $P$ is even.
The vector $n_\ell$ has integer entries, so that $u\in\Z^3$;
and $\ell(n_\ell)=n_\ell\cdot n_\ell=|n_\ell|^2$ (Definition~\ref{def:mirrors}) gives $\ell(u)=P$.

(b) The form $\ell$ vanishes on $\Gamma_\parallel$ and takes the value $P\neq0$ at $u$.
Hence $\Gamma_\parallel\cap\Z u=\{0\}$, and the sum
\[
\Gamma':=\Gamma_\parallel+\Z u
\]
is direct.
As $\Gamma_\parallel$ has rank two (Definition~\ref{def:slabs}) and $u\in\Z^3$ by (a), $\Gamma'$ is a subgroup of $\Z^3$ of rank three.
A character of the direct sum $\Gamma'$ is determined by its restriction to $\Gamma_\parallel$ and by its value at $u$,
and both may be prescribed;
so $\chi'\colon\Gamma'\to\{1,-1\}$ is well defined.
By hypothesis $\sys(\Gamma')\ge4$.
Thus $\Gamma'$ is a period lattice (Definition~\ref{def:configurations}), and $\mathbb{T}(\Xi;P)$ is a torus (Definition~\ref{def:tori}).

An element $\gamma+i\cdot u$ of $\Gamma'$, with $\gamma\in\Gamma_\parallel$ and $i\in\Z$, has the height $\ell(\gamma+i\cdot u)=i\cdot P$.
Consequently $\Gamma'\cap\{\ell=0\}=\Gamma_\parallel$, and $\chi'$ restricts to $\chi_\parallel$ there;
that is, $\Xi(\mathbb{T}(\Xi;P),\ell)=\Xi$.
Moreover $\ell(\Gamma')=P\Z$, and $\Gamma'=(\Gamma'\cap\{\ell=0\})+\Z u$ with $\ell(u)=P>0$.
So $u$ is a closing vector, and $p(\mathbb{T}(\Xi;P),\ell)=\ell(u)=P$ (Definition~\ref{def:slabs}, Lemma~\ref{lem:closing-vectors}).
For this height period the vector $\frac{P}{|n_\ell|^2}\,n_\ell$ is $u$ itself, an integer vector,
and the lattice $(\Gamma'\cap\{\ell=0\})+\Z u=\Gamma'$ has systole at least $4$.
By Definition~\ref{def:slabs} the straightened torus of $\mathbb{T}(\Xi;P)$ for $\ell$ is therefore defined.
Its lattice is $\Gamma_\parallel+\Z u=\Gamma'$,
and its character agrees with $\chi_\parallel$ on $\Gamma_\parallel$ and equals $-1$ at $u$, as $\chi'$ does.
So the straightened torus is $\mathbb{T}(\Xi;P)$ itself, and $\mathbb{T}(\Xi;P)$ is adapted to $\ell$.
Finally, Lemma~\ref{lem:closing-vectors} gives
\[
\left|\mathbb{T}(\Xi;P)\right|=P\cdot\left[\Z^3\cap\{\ell=0\}:\Gamma_\parallel\right]=P\cdot|\Xi| .
\]

(c) Write $\mathbb{T}=(\Gamma,\chi)$.
As $\Xi(\mathbb{T},\ell)=\Xi$, one has $\Gamma\cap\{\ell=0\}=\Gamma_\parallel$, and $|\mathbb{T}|=p\cdot|\Xi|$ by Lemma~\ref{lem:closing-vectors}.
By Definition~\ref{def:slabs}, $\mathbb{T}^0(\ell)$ is defined if and only if $\frac{p}{|n_\ell|^2}\,n_\ell$ is an integer vector
and $\sys\left(\Gamma_\parallel+\Z\,\frac{p}{|n_\ell|^2}\,n_\ell\right)\ge4$;
it is then the torus with the lattice $\Gamma_\parallel+\Z\,\frac{p}{|n_\ell|^2}\,n_\ell$
and the character that agrees with $\chi_\parallel$ on $\Gamma_\parallel$ and takes the value $-1$ at $\frac{p}{|n_\ell|^2}\,n_\ell$.
One entry of $n_\ell$ equals $\pm1$.
Hence $\frac{p}{|n_\ell|^2}\,n_\ell$ is an integer vector if and only if $p/|n_\ell|^2$ is an integer,
that is, if and only if $p\in|n_\ell|^2\Z$.
In that case the torus just described is $\mathbb{T}(\Xi;p)$, by (b) applied with $P=p$.
\end{proof}

\begin{proposition}[H\"older's inequality for unequal slabs]\label{prop:unequal-slabs}
Let $-1\le\rho\le1$, $w_0,\dots,w_3>0$, $M\in\R$ and $q>0$.
Let $\mathbb{T}=(\Gamma,\chi)$ be a torus and $\ell\in\mathcal{L}$ a family, with the cross-section $\Xi=\Xi(\mathbb{T},\ell)=(\Gamma_\parallel,\chi_\parallel)$;
fix a closing vector $g$ of $\mathbb{T}$ for $\ell$, and put $p:=\ell(g)$.
For an integer $k\ge1$ let $2k+1$ positions
\[
c_0<c_1<\dots<c_{2k}=c_0+p
\]
in $\Adm(\ell)$ be given;
they bound $2k$ slabs, of the heights $h_j:=c_j-c_{j-1}$, $j=1,\dots,2k$.
Suppose that these heights satisfy
\begin{equation}\label{eq:sides-systole}
\sys\left(\Gamma_\parallel+\Z\,\frac{2k\cdot h_j}{|n_\ell|^2}\,n_\ell\right)\ge4\qquad(j=1,\dots,2k).
\end{equation}
\begin{enumerate}
\item[\textup{(a)}]
The heights $h_j$ lie in $\frac{|n_\ell|^2}{2}\,\Z$ and the numbers $2k\cdot h_j$ in $|n_\ell|^2\,\Z$;
moreover, for each $j$, the whole progression $c_{j-1}+i\cdot h_j$, $i\in\Z$, lies in $\Adm(\ell)$.
By these facts and by \eqref{eq:sides-systole}, Lemma~\ref{lem:tori-over-section} defines the tori
\[
\mathbb{T}_j:=\mathbb{T}(\Xi;2k\cdot h_j)\qquad(j=1,\dots,2k).
\]
Every $\mathbb{T}_j$ is a torus adapted to $\ell$ whose cross-section $\Xi(\mathbb{T}_j,\ell)$ is $\Xi$ and whose height period $p(\mathbb{T}_j,\ell)$ is $2k\cdot h_j$;
the numbers of photon sites are $|\mathbb{T}_j|=\frac{2k\cdot h_j}{p}\cdot|\mathbb{T}|$, and
\begin{equation}\label{eq:sides-volumes}
\sum_{j=1}^{2k}\frac{|\mathbb{T}_j|}{2k}=|\mathbb{T}| .
\end{equation}
\item[\textup{(b)}]
Let $G_j\in\mathcal{A}_\Xi[c_{j-1},c_j]$ for $j=1,\dots,2k$,
and let $\fold_\ell G_j$ be the fold of $G_j$ on the torus $\mathbb{T}_j$ (Definition~\ref{def:slab-copies})
formed with the planes $\{\ell=c_{j-1}+i\cdot h_j\}$, $i\in\Z$:
the base point is $c_{j-1}$, the height of the slabs is $h'=h_j$,
and $G_j$ lies in the first slab $\{c_{j-1}<\ell<c_j\}$.
Let $\mathfrak{S}$ be the map of Corollary~\ref{cor:slab-admissible} for the pair $(\ell,\Xi)$,
and $\|\cdot\|_{2k}$ the norm of Definition~\ref{def:even-schatten}.
Then
\begin{gather}
Z_{\mathbb{T}_j}\left[\fold_\ell G_j\right]=\left\|\mathfrak{S}[G_j]\right\|_{2k}^{2k}\ge0\qquad(j=1,\dots,2k),\label{eq:sides-slab-norm}\\
\left|Z_{\mathbb{T}}\left[G_1\cdots G_{2k}\right]\right|\le\prod_{j=1}^{2k}Z_{\mathbb{T}_j}\left[\fold_\ell G_j\right]^{\frac{1}{2k}} .\label{eq:sides-holder}
\end{gather}
On the left side of \eqref{eq:sides-holder} the product $G_1\cdots G_{2k}$ need not embed in $\mathbb{T}$;
by Corollary~\ref{cor:slab-admissible}(f) it lies in the space $\mathcal{A}(\mathbb{T})$ of Definition~\ref{def:torus-integrands}, so that it is a finite sum of products read on $\mathbb{T}$,
and $Z_{\mathbb{T}}$ stands there for the functional on such sums (Definition~\ref{def:functional}\eqref{it:read-on-torus}),
as in the trace formula of that corollary.
When every $G_j$ is a bounded Borel function of finitely many $F_\sigma$,
its copies are functions of the same kind (Lemma~\ref{lem:admissible-images}\eqref{it:images-boxes}),
and so is the argument of each functional in \eqref{eq:sides-slab-norm} and \eqref{eq:sides-holder}:
an element of $\mathcal{A}(\R^3)$ with no boxes, which embeds in every torus (Definition~\ref{def:integrands}).
In that case the values $Z_{\mathbb{T}}\left[G_1\cdots G_{2k}\right]$ and $Z_{\mathbb{T}_j}\left[\fold_\ell G_j\right]$ are given by the formula of Definition~\ref{def:functional}.
\item[\textup{(c)}]
If $h_1=\dots=h_{2k}=:h'$, then $\mathbb{T}^0(\ell)$ is defined and $\mathbb{T}_j=\mathbb{T}^0(\ell)$ for all $j$,
and \eqref{eq:sides-slab-norm}, \eqref{eq:sides-holder} are the identity \eqref{eq:tori-fold-norm} and the inequality \eqref{eq:tori-holder} of Proposition~\ref{prop:twisted-holder}.
\end{enumerate}
\end{proposition}

\begin{proof}
By Lemma~\ref{lem:closing-vectors}, $\Xi=\Xi(\mathbb{T},\ell)$ is a cross-section for $\ell$.

(a) The points $c_{j-1}$ and $c_j$ lie in the group $\Adm(\ell)=\frac{|n_\ell|^2}{2}\Z$,
and therefore so do $h_j=c_j-c_{j-1}$ and $c_{j-1}+i\cdot h_j$ for every $i\in\Z$.
The number $2h_j/|n_\ell|^2$ is thus an integer, and
\[
2k\cdot h_j=|n_\ell|^2\cdot k\cdot\frac{2h_j}{|n_\ell|^2}\in|n_\ell|^2\,\Z .
\]
As $h_j>0$, the number $P:=2k\cdot h_j$ is a positive element of $|n_\ell|^2\Z$.
For this $P$ the vector $u$ of Lemma~\ref{lem:tori-over-section} is $\frac{2k\cdot h_j}{|n_\ell|^2}\,n_\ell$,
and \eqref{eq:sides-systole} is the hypothesis of part (b) of that lemma.
Lemma~\ref{lem:tori-over-section}(b) shows that $\mathbb{T}_j$ is a torus adapted to $\ell$,
with $\Xi(\mathbb{T}_j,\ell)=\Xi$, $p(\mathbb{T}_j,\ell)=2k\cdot h_j$ and $|\mathbb{T}_j|=2k\cdot h_j\cdot|\Xi|$,
while Lemma~\ref{lem:tori-over-section}(c) gives $|\mathbb{T}|=p\cdot|\Xi|$.
Hence $|\mathbb{T}_j|=\frac{2k\cdot h_j}{p}\cdot|\mathbb{T}|$, and
\[
\sum_{j=1}^{2k}\frac{|\mathbb{T}_j|}{2k}
=|\Xi|\cdot\sum_{j=1}^{2k}h_j
=|\Xi|\cdot(c_{2k}-c_0)
=p\cdot|\Xi|
=|\mathbb{T}| ,
\]
which is \eqref{eq:sides-volumes}.

(b) Fix $j$.
For \eqref{eq:sides-slab-norm}, apply Proposition~\ref{prop:twisted-holder} to the torus $\mathbb{T}_j$ in the place of $\mathbb{T}$:
the family is $\ell$, the cross-section is $\Xi(\mathbb{T}_j,\ell)=\Xi$,
the closing vector is $\frac{2k\cdot h_j}{|n_\ell|^2}\,n_\ell$ (Lemma~\ref{lem:tori-over-section}(b)),
the integer is the same $k$, and $h':=h_j$, so that $p(\mathbb{T}_j,\ell)=2k\cdot h'$;
the planes are $c_{j-1}+i\cdot h_j$, $i\in\Z$, which lie in $\Adm(\ell)$ by (a),
and the element $G_j$ of $\mathcal{A}_\Xi[c_{j-1},c_{j-1}+h_j]$ occupies the first slab.
Since $\mathbb{T}_j$ is adapted to $\ell$, its straightened torus is defined and equals $\mathbb{T}_j$.
The identity \eqref{eq:tori-fold-norm} for $\mathbb{T}_j$ is therefore \eqref{eq:sides-slab-norm}.

For \eqref{eq:sides-holder}, note that by Corollary~\ref{cor:slab-admissible} the trace formula of Theorem~\ref{thm:slab-operators}\eqref{it:slab-trace}
applies to $\mathbb{T}$, the closing vector $g$, the $2k\ge2$ slabs between the planes $c_0<\dots<c_{2k}=c_0+p$,
and the elements $G_j\in\mathcal{A}_\Xi[c_{j-1},c_j]$.
It gives
\[
Z_{\mathbb{T}}\left[G_1\cdots G_{2k}\right]=\Tr_{\mathcal{H}(c_0)}\left(U\cdot\mathfrak{Y}_{2k}\cdots\mathfrak{Y}_1\right),
\]
where $\mathfrak{Y}_j:=\mathfrak{S}[G_j]\colon\mathcal{H}(c_{j-1})\to\mathcal{H}(c_j)$ is a Hilbert--Schmidt operator between separable Hilbert spaces, and
\[
U:=V_g^{-1}\cdot U_{\mathrm{f}}^{\frac{1+\chi(g)}{2}}\colon\mathcal{H}(c_0+p)\to\mathcal{H}(c_0)
\]
is a product of two unitaries.
In Corollary~\ref{cor:slab-admissible} the definitions of $\mathcal{H}(c)$, $V_v$, $U_{\mathrm{f}}$ and $\mathfrak{S}$ involve the parameters, $\ell$ and $\Xi$, but no torus;
since $\mathbb{T}_j$ has the same cross-section $\Xi$ for $\ell$ as $\mathbb{T}$,
the factor $\mathfrak{Y}_j$ of this trace is the operator whose norm appears in \eqref{eq:sides-slab-norm}.
Lemma~\ref{lem:holder-spaces}, applied to the spaces $\mathcal{H}(c_0),\dots,\mathcal{H}(c_{2k})$, the operators $\mathfrak{Y}_1,\dots,\mathfrak{Y}_{2k}$ and the unitary $U$, yields
\[
\left|\Tr_{\mathcal{H}(c_0)}\left(U\cdot\mathfrak{Y}_{2k}\cdots\mathfrak{Y}_1\right)\right|\le\prod_{j=1}^{2k}\left\|\mathfrak{Y}_j\right\|_{2k} .
\]
Together with \eqref{eq:sides-slab-norm} this gives
\[
\left|Z_{\mathbb{T}}\left[G_1\cdots G_{2k}\right]\right|
\le\prod_{j=1}^{2k}\left(\left\|\mathfrak{Y}_j\right\|_{2k}^{2k}\right)^{\frac{1}{2k}}
=\prod_{j=1}^{2k}Z_{\mathbb{T}_j}\left[\fold_\ell G_j\right]^{\frac{1}{2k}} .
\]

(c) Let all the heights be equal to $h'$.
Then $2k\cdot h'=\sum_jh_j=p$, so that $2k\cdot h_j=p$ and $\mathbb{T}_j=\mathbb{T}(\Xi;p)$ for every $j$.
By (a), $p=2k\cdot h_1$ lies in $|n_\ell|^2\Z$,
and \eqref{eq:sides-systole} is the condition on the systole in Lemma~\ref{lem:tori-over-section}(c);
by that lemma the straightened torus $\mathbb{T}^0=\mathbb{T}^0(\ell)=(\Gamma^0,\chi^0)$ is defined and equals $\mathbb{T}(\Xi;p)$.
The positions $c_0+i\cdot h'$, $i\in\Z$, lie in $\Adm(\ell)$ by (a), and $c_j=c_0+j\cdot h'$ for $0\le j\le2k$;
so $\mathbb{T}$, $\ell$, $g$, $k$, $h'$, $c_0$ and $G_1,\dots,G_{2k}$ satisfy the hypotheses of Proposition~\ref{prop:twisted-holder}.
It remains to compare the folds.
Fix $j$.
The planes $\{\ell=c_{j-1}+i\cdot h'\}$, $i\in\Z$, with which the fold of (b) is formed, are the planes $\{\ell=c_0+i\cdot h'\}$ of Proposition~\ref{prop:twisted-holder};
hence the reflections in them generate the same group,
and the copy of $G_j$ in a given slab is the same whether the slabs are numbered from the base point $c_{j-1}$ or from $c_0$ (Definition~\ref{def:slab-copies}).
In the numbering from $c_0$, the fold of (b) is the product of the copies of $G_j$ over the slabs $j,\dots,j+2k-1$,
while the fold of Proposition~\ref{prop:twisted-holder} is the product over the slabs $1,\dots,2k$;
for $j>1$ the two need not coincide as elements of $\mathcal{A}(\R^3)$.
By the last part of Definition~\ref{def:slab-copies}, applied to the adapted torus $\mathbb{T}^0$,
the products of the copies over any $2k$ consecutive slabs define one and the same function on $\Conf(\Gamma^0)$;
so the two folds have the same value under $Z_{\mathbb{T}^0}$.
With $\mathbb{T}_j=\mathbb{T}^0$ for every $j$, the identity \eqref{eq:sides-slab-norm} thus becomes \eqref{eq:tori-fold-norm},
and \eqref{eq:sides-holder} becomes \eqref{eq:tori-holder}.
\end{proof}

\begin{remark}\label{rem:unequal-slabs}
\begin{enumerate}
\item[(1)]
The straightened torus of $\mathbb{T}$ need not exist, and the hypotheses of Proposition~\ref{prop:unequal-slabs} do not require it.
Let $n\ge4$ be odd, and consider the cubic torus $C(n)$ and the family $\ell_-$.
As $\ell_-(n\Z^3)=n\Z$, the height period is $p(C(n),\ell_-)=n$;
it does not lie in $|n_{\ell_-}|^2\Z=2\Z$,
and $\frac{p}{|n_{\ell_-}|^2}\,n_{\ell_-}=\frac n2\,e_-$ is not an integer vector.
So the straightened torus of $C(n)$ for $\ell_-$ is not defined (Lemma~\ref{lem:tori-over-section}(c)).
For this reason the side $n$ is even in Proposition~\ref{prop:untwistings}(a),
where the straightened torus of $C(n)$ for $\ell_-$ is $D'(n/2,n,n)$.
In Proposition~\ref{prop:unequal-slabs}, on the other hand, the height periods $2k\cdot h_j$ of the tori $\mathbb{T}_j$ lie in $|n_\ell|^2\Z$ by part (a), whatever the parity of $p$;
so the proposition applies to $C(n)$ and $\ell_-$ for odd $n$ as well, whenever \eqref{eq:sides-systole} holds.
\item[(2)]
The number $2k$ of the slabs has to be even, as in Proposition~\ref{prop:twisted-holder} and for the reason given in Remark~\ref{rem:twist}.
Their heights, on the other hand, are arbitrary positive elements of $\frac{|n_\ell|^2}{2}\,\Z$,
subject only to $h_1+\dots+h_{2k}=p$ and to \eqref{eq:sides-systole}.
\item[(3)]
Let $C>0$, and multiply, for every torus $\mathbb{T}'$, the integrand in the formula of Definition~\ref{def:functional} by $C^{|\mathbb{T}'|}$,
that is, by one factor $C$ for each photon site of $\mathbb{T}'$.
Every functional $Z_{\mathbb{T}'}$ is then replaced by $C^{|\mathbb{T}'|}\cdot Z_{\mathbb{T}'}$.
The left side of \eqref{eq:sides-holder} acquires the factor $C^{|\mathbb{T}|}$, and the right side the factor
\[
\prod_{j=1}^{2k}C^{\frac{|\mathbb{T}_j|}{2k}}=C^{|\mathbb{T}|},
\]
the exponents adding up to $|\mathbb{T}|$ by \eqref{eq:sides-volumes}.
Hence both sides of \eqref{eq:sides-holder} are homogeneous of the same degree $|\mathbb{T}|$ in $C$.
\end{enumerate}
\end{remark}

\subsection{One wider, empty slab}\label{sec:empty-slab}

For real $\rho$ and $M$, hopping weights $w_0,\dots,w_3>0$ and a lattice charge $q>0$, let
\[
C_q=\sup_{y\in\R}\,\sum_{k\in\Z}\exp\left(-\frac{1}{2}\left(y+\frac{2\pi}{q}\,k\right)^{2}\right)
\]
be the constant of Lemma~\ref{lem:bounded-integrands}.
It is finite by that lemma, and it is positive, every term of the sums being positive.
Put
\begin{equation}\label{eq:sides-constants}
C_D:=|M|+\frac{1+|\rho|}{2}\sum_{i=0}^{3}w_i,
\qquad
\bar f:=192\log C_D+3\log\left(\frac{2\pi}{q}\cdot C_q\right).
\end{equation}
Since the weights are positive, $C_D>0$;
it bounds the operator norm of $D^{\mathbb{T}}_A$ for every torus $\mathbb{T}$ and every real $A$ (Proposition~\ref{prop:accretivity}\eqref{it:operator-norm}).
Hence $\bar f$ is a real number, and it depends only on $\rho$, $M$, $q$ and $w_0,\dots,w_3$.

\begin{lemma}[A bound on all partition functions]\label{lem:uniform-ceiling}
Let $\rho$ and $M$ be real, $w_0,\dots,w_3>0$ and $q>0$, and let $\bar f$ be given by \eqref{eq:sides-constants}.
Every torus $\mathbb{T}$ satisfies
\[
0\le Z_{\mathbb{T}}[1]\le\exp(\bar f\cdot|\mathbb{T}|).
\]
Consequently $f_\infty\le f(n_1,n_2,n_3)\le\bar f$ for all integers $n_1,n_2,n_3\ge4$.
\end{lemma}

\begin{proof}
Let $\Gamma$ be the period lattice of $\mathbb{T}$.
By Definitions~\ref{def:tori} and~\ref{def:functional}, $Z_{\mathbb{T}}[1]$ is the integral over $\Conf(\Gamma)$ of the function
\[
\exp\left(-\frac{1}{2}\|F\|^{2}\right)\cdot\det D^{\mathbb{T}}_A,
\]
which is non-negative by Proposition~\ref{prop:accretivity}\eqref{it:determinant-sign};
hence $Z_{\mathbb{T}}[1]\ge0$.
For the upper bound, the estimate \eqref{eq:positivity-ceiling}, which holds for every real $\rho$, reads in the notation \eqref{eq:sides-constants}
\[
\begin{aligned}
Z_{\mathbb{T}}[1]
&\le\left(\frac{2\pi}{q}\cdot C_q\right)^{3|\mathbb{T}|}\cdot C_D^{192|\mathbb{T}|}\\
&=\exp\left(|\mathbb{T}|\cdot\left(192\log C_D+3\log\left(\frac{2\pi}{q}\cdot C_q\right)\right)\right)
=\exp(\bar f\cdot|\mathbb{T}|).
\end{aligned}
\]

Now let $n_1,n_2,n_3\ge4$ be integers.
Each non-zero vector of $n_1\Z e_1+n_2\Z e_2+n_3\Z e_3$, the period lattice of $C(n_1,n_2,n_3)$, has a coordinate that is a non-zero multiple of some $n_\mu$;
so the systole of this lattice is at least $\min_\mu n_\mu\ge4$, and $C(n_1,n_2,n_3)$ is a torus with $|C(n_1,n_2,n_3)|=n_1n_2n_3$ (Definition~\ref{def:tori}).
Its partition function is positive by Theorem~\ref{thm:partition-positivity}.
The pressure
\[
f(n_1,n_2,n_3)=\frac{1}{n_1n_2n_3}\,\log Z_{C(n_1,n_2,n_3)}[1]
\]
of Definition~\ref{def:sides} is therefore a real number,
and the bound already proved, applied to $\mathbb{T}=C(n_1,n_2,n_3)$, yields $f(n_1,n_2,n_3)\le\bar f$.
Finally, $f_\infty$ is the infimum of $f$ over all triples of integers $\ge4$ (Proposition~\ref{prop:coordinate-pressure}(d)), so that $f_\infty\le f(n_1,n_2,n_3)$.
\end{proof}

Let $\mathcal{Q}=(\mathfrak{F};o_1,o_2,o_3;h)$ be a cell structure (Definition~\ref{def:cell-structures}).
Its frame $\mathfrak{F}=(\ell^1,\ell^2,\ell^3)$ consists of three families with pairwise orthogonal normals $n_1,n_2,n_3$,
its walls are the mirrors $(\ell^a,c)$ with $c\in o_a+h\Z$,
and $Q_\kappa$, $\kappa\in\Z^3$, are its open cells.
For $a\in\{1,2,3\}$ let $W_a$ be the group generated by the reflections in the walls $(\ell^a,c)$, $c\in o_a+h\Z$,
and for $i,j\in\Z$ let $g^a_{i\leftarrow j}$ be the element of $W_a$ that maps the layer $\{o_a+jh<\ell^a<o_a+(j+1)h\}$ onto the layer $\{o_a+ih<\ell^a<o_a+(i+1)h\}$ (Definition~\ref{def:cell-families}).
It maps every cell $Q_\kappa$ with $\kappa_a=j$ onto $Q_{\kappa^{a\to i}}$,
where $\kappa^{a\to i}\in\Z^3$ denotes $\kappa$ with its $a$-th entry replaced by $i$.

\begin{definition}[Finite fillings]\label{def:fillings}
A \emph{finite filling} of $\mathcal{Q}$ is a family $\mathfrak{G}=(G_\kappa)_{\kappa\in\mathcal{K}}$,
indexed by a finite subset $\mathcal{K}$ of $\Z^3$,
in which every $G_\kappa$ is a bounded complex Borel function of finitely many values $F_\sigma$ of the curvature at plaquettes $\sigma$ contained in the open cell $Q_\kappa$.
The product of $\mathfrak{G}$ is
\[
\prod\mathfrak{G}:=\prod_{\kappa\in\mathcal{K}}G_\kappa ,
\]
with $\prod\mathfrak{G}:=1$ if $\mathcal{K}$ is empty.
For $a\in\{1,2,3\}$, $j\in\Z$ and an integer $k\ge1$,
the \emph{fold of the layer $j$ of $\mathfrak{G}$ in the direction $a$ over $2k$ slabs}
is the family $\mathfrak{R}_a^{j,2k}\mathfrak{G}:=\left(G'_{\kappa'}\right)_{\kappa'\in\mathcal{K}'}$ given by
\[
\mathcal{K}':=\left\{\kappa^{a\to j'}:\ \kappa\in\mathcal{K},\ \kappa_a=j,\ j\le j'\le j+2k-1\right\},
\qquad
G'_{\kappa^{a\to j'}}:=\left(G_\kappa\right)^{g^a_{j'\leftarrow j}} .
\]
\end{definition}

Here $G^g$ is the image of $G$ under $g\in W$ \eqref{eq:setting-group-action};
for a bounded Borel function $G$ of the curvature,
$G^g(F)=G(g^*F)$ if $g$ is even and $G^g(F)=\overline{G(g^*F)}$ if $g$ is odd.
No torus enters the definition.
The product $\prod\mathfrak{G}$ is a bounded Borel function of finitely many $F_\sigma$;
it is therefore an element of $\mathcal{A}(\R^3)$, a product without integrands ($r=0$ in Definition~\ref{def:integrands}).
Such a product embeds in every torus,
and on a torus it is evaluated at the periodic extension $\widetilde F$ of the curvature.

The entries of $\mathfrak{R}_a^{j,2k}\mathfrak{G}$ are well defined:
the map $(\kappa,j')\mapsto\kappa^{a\to j'}$ is injective on $\{\kappa\in\mathcal{K}:\kappa_a=j\}\times\Z$,
since the $a$-th entry of $\kappa^{a\to j'}$ is $j'$, and the replacement of this entry by $j$ gives back $\kappa$.
As $g^a_{j'\leftarrow j}$ maps $Q_\kappa$ onto $Q_{\kappa^{a\to j'}}$,
Lemma~\ref{lem:admissible-images}\eqref{it:images-boxes} shows that $G'_{\kappa^{a\to j'}}$ is a bounded complex Borel function of finitely many $F_\sigma$ with $\sigma$ contained in $Q_{\kappa^{a\to j'}}$.
So $\mathfrak{R}_a^{j,2k}\mathfrak{G}$ is again a finite filling of $\mathcal{Q}$.
Its indices differ from those of $\mathfrak{G}$ in the $a$-th entry only:
in each direction other than $a$, the entries of the elements of $\mathcal{K}'$ are entries of elements $\kappa\in\mathcal{K}$ with $\kappa_a=j$.
For a family of cell contents on a torus, indexed by all of $\Z^3$ (Definition~\ref{def:cell-families}),
the same formula, read on the $2k$ layers $j\le\kappa'_a\le j+2k-1$,
gives the restriction of its fold $\mathfrak{R}_a^j$ to these layers.

\begin{lemma}[The fold of one layer]\label{lem:layer-fold}
Let $-1\le\rho\le1$.
Let $\mathcal{Q}=(\mathfrak{F};o_1,o_2,o_3;h)$ be a cell structure, $a\in\{1,2,3\}$, $\ell:=\ell^a$,
$\Xi=(\Gamma_\parallel,\chi_\parallel)$ a cross-section for $\ell$,
$\mathfrak{G}=(G_\kappa)_{\kappa\in\mathcal{K}}$ a finite filling of $\mathcal{Q}$ and $j\in\Z$.
Let
\[
\Phi:=\prod_{\kappa\in\mathcal{K},\ \kappa_a=j}G_\kappa
\]
be the content of the layer $j$, with $\Phi:=1$ if no element of $\mathcal{K}$ has the $a$-th entry $j$.
\begin{enumerate}
\item\label{it:sides-layer-slab}
The set $o_a+h\Z$ is contained in $\Adm(\ell)$,
and $\Phi$ lies in $\mathcal{A}_\Xi[o_a+jh,o_a+(j+1)h]$.
\item\label{it:sides-layer-fold}
Let $k\ge1$ be an integer with $\sys\left(\Gamma_\parallel+\Z\,\frac{2kh}{|n_a|^2}\,n_a\right)\ge4$,
and let $\mathbb{T}^{\mathrm{reg}}:=\mathbb{T}(\Xi;2kh)$ be the torus of Lemma~\ref{lem:tori-over-section}(b)
(the number $2kh$ lies in $2\Z$, which is contained in $|n_a|^2\Z$).
Let $\fold_\ell\Phi$ be the fold of $\Phi$ on $\mathbb{T}^{\mathrm{reg}}$ (Definition~\ref{def:slab-copies}) formed with the base point $o_a+jh$ and the height $h$;
its planes are $o_a+h\Z$, and by \eqref{it:sides-layer-slab} the element $\Phi$ belongs to its first slab $\{o_a+jh<\ell<o_a+(j+1)h\}$.
Then
\[
\fold_\ell\Phi=\prod\mathfrak{R}_a^{j,2k}\mathfrak{G},
\]
and
\begin{equation}\label{eq:sides-layer-norm}
Z_{\mathbb{T}^{\mathrm{reg}}}\left[\prod\mathfrak{R}_a^{j,2k}\mathfrak{G}\right]=\left\|\mathfrak{S}[\Phi]\right\|_{2k}^{2k}\ge0 ,
\end{equation}
where $\mathfrak{S}[\Phi]\colon\mathcal{H}(o_a+jh)\to\mathcal{H}(o_a+(j+1)h)$ is the slab operator of $\Phi$ for the pair $(\ell,\Xi)$ (Corollary~\ref{cor:slab-admissible})
and $\|\cdot\|_{2k}$ is the norm of Definition~\ref{def:even-schatten}.
\end{enumerate}
\end{lemma}

\begin{proof}
The position $o_a$ lies in $\Adm(\ell)$ (Definition~\ref{def:cell-structures}), $h$ is an integer,
and $\Adm(\ell)$, which is $\frac12\Z$ or $\Z$, is invariant under integer translations (Definition~\ref{def:mirrors});
hence $o_a+h\Z\subseteq\Adm(\ell)$.
A cell $Q_\kappa$ with $\kappa_a=j$ lies in the open slab $\{o_a+jh<\ell<o_a+(j+1)h\}$.
Consequently $\Phi$ is a bounded Borel function of finitely many $F_\sigma$, $\sigma$ contained in this slab:
an admissible product in $\mathcal{A}(\{o_a+jh<\ell<o_a+(j+1)h\})$ without factors of Gaussian growth,
which embeds in $\Gamma_\parallel$ (Definition~\ref{def:integrands}).
By Definition~\ref{def:slabs}, $\Phi$ belongs to $\mathcal{A}_\Xi[o_a+jh,o_a+(j+1)h]$;
the constant $1$ is covered as the empty product.
This proves \eqref{it:sides-layer-slab}.

For \eqref{it:sides-layer-fold}, take in Definition~\ref{def:slab-copies} the base point $c'_0:=o_a+jh$ and the height $h':=h$,
so that $c'_i=o_a+(j+i)h$ for $i\in\Z$.
These are the planes $o_a+h\Z$, which lie in $\Adm(\ell)$ by \eqref{it:sides-layer-slab};
the reflections $\theta_{\ell,c'_i}$, $i\in\Z$, are the reflections in the walls $(\ell^a,c)$, $c\in o_a+h\Z$,
so that the group generated by them is $W_a$.
The slab number $i$, namely $\{c'_{i-1}<\ell<c'_i\}$, is the layer of the cells with $\kappa_a=j+i-1$,
and $\Phi$ lies in the slab number $1$.
The element $g_{i\leftarrow1}$ of this group that maps the slab number $1$ onto the slab number $i$ is unique (Definition~\ref{def:slab-copies}),
and $g^a_{j+i-1\leftarrow j}$ has the same property;
so $g_{i\leftarrow1}=g^a_{j+i-1\leftarrow j}$.
By Lemma~\ref{lem:tori-over-section}(b), the torus $\mathbb{T}^{\mathrm{reg}}$ is adapted to $\ell$,
with $\Xi(\mathbb{T}^{\mathrm{reg}},\ell)=\Xi$ and $p(\mathbb{T}^{\mathrm{reg}},\ell)=2kh=2kh'$;
the fold of $\Phi$ on it is therefore defined.
With $j':=j+i-1$, Lemma~\ref{lem:admissible-images}\eqref{it:images-products} gives
\[
\fold_\ell\Phi=\prod_{i=1}^{2k}\Phi^{g_{i\leftarrow1}}
=\prod_{j'=j}^{j+2k-1}\ \prod_{\kappa\in\mathcal{K},\ \kappa_a=j}\left(G_\kappa\right)^{g^a_{j'\leftarrow j}}
=\prod\mathfrak{R}_a^{j,2k}\mathfrak{G},
\]
the last equality because $(\kappa,j')\mapsto\kappa^{a\to j'}$ maps
$\{\kappa\in\mathcal{K}:\kappa_a=j\}\times\{j,\dots,j+2k-1\}$ bijectively onto $\mathcal{K}'$.
Being adapted to $\ell$, the torus $\mathbb{T}^{\mathrm{reg}}$ is its own straightened torus.
The identity \eqref{eq:tori-fold-norm} of Proposition~\ref{prop:twisted-holder},
applied to $\mathbb{T}^{\mathrm{reg}}$ with the integer $k$, the height $h'=h$, the planes $c'_i$ and the element $\Phi$ of the first slab,
gives $Z_{\mathbb{T}^{\mathrm{reg}}}\left[\fold_\ell\Phi\right]=\left\|\mathfrak{S}[\Phi]\right\|_{2k}^{2k}\ge0$,
and the identity of the folds just obtained turns this into \eqref{eq:sides-layer-norm}.
\end{proof}

\begin{lemma}[One step with an empty slab]\label{lem:one-step}
Let $-1\le\rho\le1$.
Let $\mathcal{Q}=(\mathfrak{F};o_1,o_2,o_3;h)$ be a cell structure, $a\in\{1,2,3\}$, $\mathbb{T}=(\Gamma,\chi)$ a torus,
$\Xi:=\Xi(\mathbb{T},\ell^a)=(\Gamma_\parallel,\chi_\parallel)$ and $p:=p(\mathbb{T},\ell^a)$.
Let $k\ge1$ be an integer with $(2k-1)h<p$, and put
\[
h^{\mathrm{exc}}:=p-(2k-1)h>0,\qquad
\Gamma^{\mathrm{reg}}:=\Gamma_\parallel+\Z\,\frac{2kh}{|n_a|^2}\,n_a,\qquad
\Gamma^{\mathrm{exc}}:=\Gamma_\parallel+\Z\,\frac{2kh^{\mathrm{exc}}}{|n_a|^2}\,n_a .
\]
Assume that $\sys(\Gamma^{\mathrm{reg}})\ge4$ and $\sys(\Gamma^{\mathrm{exc}})\ge4$,
and let $\mathbb{T}^{\mathrm{reg}}:=\mathbb{T}(\Xi;2kh)$ and $\mathbb{T}^{\mathrm{exc}}:=\mathbb{T}(\Xi;2kh^{\mathrm{exc}})$
(the numbers $2kh$ and $2kh^{\mathrm{exc}}$ are even, hence multiples of $|n_a|^2$).
By Lemma~\ref{lem:tori-over-section}(b) these are tori adapted to $\ell^a$, with the lattices $\Gamma^{\mathrm{reg}}$ and $\Gamma^{\mathrm{exc}}$;
their characters are $\chi_\parallel$ on $\Gamma_\parallel$ and $-1$ at the generator $\frac{2kh}{|n_a|^2}\,n_a$,
respectively $\frac{2kh^{\mathrm{exc}}}{|n_a|^2}\,n_a$.
Let $\mathfrak{G}=(G_\kappa)_{\kappa\in\mathcal{K}}$ be a finite filling of $\mathcal{Q}$ with $0\le\kappa_a\le2k-2$ for all $\kappa\in\mathcal{K}$.
Then
\begin{gather}
Z_{\mathbb{T}^{\mathrm{reg}}}\left[\prod\mathfrak{R}_a^{j,2k}\mathfrak{G}\right]\ge0\quad(0\le j\le2k-2),
\qquad
Z_{\mathbb{T}^{\mathrm{exc}}}[1]\ge0,
\label{eq:sides-step-sign}\\
\left|Z_{\mathbb{T}}\left[\prod\mathfrak{G}\right]\right|
\le Z_{\mathbb{T}^{\mathrm{exc}}}[1]^{\frac{1}{2k}}\cdot\prod_{j=0}^{2k-2}Z_{\mathbb{T}^{\mathrm{reg}}}\left[\prod\mathfrak{R}_a^{j,2k}\mathfrak{G}\right]^{\frac{1}{2k}},
\label{eq:sides-step}\\
|\mathbb{T}|=\frac{2k-1}{2k}\,|\mathbb{T}^{\mathrm{reg}}|+\frac{1}{2k}\,|\mathbb{T}^{\mathrm{exc}}|.
\label{eq:sides-step-volumes}
\end{gather}
\end{lemma}

\begin{proof}
Write $\ell:=\ell^a$, fix a closing vector $g$ of $\mathbb{T}$ for $\ell$, so that $\ell(g)=p$ (Lemma~\ref{lem:closing-vectors}), and put
\[
c_j:=o_a+jh\quad(0\le j\le2k-1),\qquad c_{2k}:=o_a+p .
\]
By Lemma~\ref{lem:layer-fold}\eqref{it:sides-layer-slab} the points $c_0,\dots,c_{2k-1}$ lie in $\Adm(\ell)$;
so does $c_{2k}$, because $o_a\in\Adm(\ell)$ and $p$ is an integer.
Since $c_{2k}-c_{2k-1}=h^{\mathrm{exc}}>0$, they form a chain $c_0<c_1<\dots<c_{2k-1}<c_{2k}=c_0+p$,
whose heights $h_j:=c_j-c_{j-1}$ are $h_j=h$ for $1\le j\le2k-1$ and $h_{2k}=h^{\mathrm{exc}}$.
In Proposition~\ref{prop:unequal-slabs}, applied to $\mathbb{T}$, $\ell$, $g$, $k$ and this chain,
one therefore has $2kh_j=2kh$ and $\mathbb{T}_j=\mathbb{T}^{\mathrm{reg}}$ for $j\le2k-1$,
while $2kh_{2k}=2kh^{\mathrm{exc}}$ and $\mathbb{T}_{2k}=\mathbb{T}^{\mathrm{exc}}$.
The condition \eqref{eq:sides-systole} of that proposition is the assumption on $\sys(\Gamma^{\mathrm{reg}})$ and $\sys(\Gamma^{\mathrm{exc}})$,
and its identity \eqref{eq:sides-volumes} reads \eqref{eq:sides-step-volumes}.

It remains to identify the contents of the slabs and their folds.
For $1\le j\le2k-1$ let
\[
\Phi_j:=\prod_{\kappa\in\mathcal{K},\ \kappa_a=j-1}G_\kappa
\]
be the content of the layer $j-1$;
it lies in $\mathcal{A}_\Xi[c_{j-1},c_j]$ by Lemma~\ref{lem:layer-fold}\eqref{it:sides-layer-slab}.
The last slab receives $\Phi_{2k}:=1$, an element of $\mathcal{A}_\Xi[c_{2k-1},c_{2k}]$.
Every $\kappa\in\mathcal{K}$ satisfies $0\le\kappa_a\le2k-2$ and so lies in exactly one of the layers $0,\dots,2k-2$;
hence $\prod\mathfrak{G}=\Phi_1\cdots\Phi_{2k}$.
For $1\le j\le2k-1$ the fold of $\Phi_j$ in Proposition~\ref{prop:unequal-slabs}(b) is formed on $\mathbb{T}_j=\mathbb{T}^{\mathrm{reg}}$
with the base point $c_{j-1}=o_a+(j-1)h$ and the height $h_j=h$.
These are the torus, the base point and the height that Lemma~\ref{lem:layer-fold}\eqref{it:sides-layer-fold} prescribes for the layer $j-1$, whose content is $\Phi_j$,
and the condition on the systole in that lemma is the assumption $\sys(\Gamma^{\mathrm{reg}})\ge4$;
so the lemma gives $\fold_\ell\Phi_j=\prod\mathfrak{R}_a^{j-1,2k}\mathfrak{G}$.
The fold of $\Phi_{2k}=1$ on $\mathbb{T}^{\mathrm{exc}}$ is $1$, by Lemma~\ref{lem:admissible-images}\eqref{it:images-products}.
With these identifications, \eqref{eq:sides-slab-norm} yields \eqref{eq:sides-step-sign},
and \eqref{eq:sides-holder} yields \eqref{eq:sides-step}.
\end{proof}

\begin{remark}\label{rem:one-step}
\begin{enumerate}
\item
If $p=2kh$, then $h^{\mathrm{exc}}=h$ and $\mathbb{T}^{\mathrm{exc}}=\mathbb{T}^{\mathrm{reg}}=\mathbb{T}^0(\ell^a)$,
the straightened torus of $\mathbb{T}$ (Lemma~\ref{lem:tori-over-section}(c)).
In this case \eqref{eq:sides-step} is the inequality \eqref{eq:tori-holder} of Proposition~\ref{prop:twisted-holder}
for the planes $o_a+jh$, with the constant $1$ in the last of the $2k$ slabs.
\item
The exceptional slab $\{o_a+(2k-1)h<\ell^a<o_a+p\}$, of height $h^{\mathrm{exc}}$, carries no content of $\mathfrak{G}$.
The height period of $\mathbb{T}^{\mathrm{exc}}$ for $\ell^a$ is $2kh^{\mathrm{exc}}$ (Lemma~\ref{lem:tori-over-section}(b)),
so that the $2k$ slabs of a fold on this torus have the height $h^{\mathrm{exc}}$.
The copies of a content of the exceptional slab would therefore lie in cells whose width in the direction $a$ is $h^{\mathrm{exc}}$,
and not $h$ as for the cells of $\mathcal{Q}$.
As the slab is empty, its factor in \eqref{eq:sides-step} is $Z_{\mathbb{T}^{\mathrm{exc}}}[1]^{\frac{1}{2k}}$, in which no content occurs,
and every content of $\mathfrak{G}$ is carried into cells of $\mathcal{Q}$ on $\mathbb{T}^{\mathrm{reg}}$.
\item
Let $b\in\{1,2,3\}$ be a direction other than $a$.
Every fold $\mathfrak{R}_a^{j,2k}\mathfrak{G}$ is again a finite filling of $\mathcal{Q}$,
and the $b$-th entry of each of its indices is the $b$-th entry of an index of $\mathfrak{G}$.
If the indices of $\mathfrak{G}$ satisfy the condition that the lemma imposes on indices in the direction $b$, then so do the indices of all these folds,
and the lemma may be applied to each of them once more, in the direction $b$ and with $\mathbb{T}^{\mathrm{reg}}$ in the place of $\mathbb{T}$,
whenever its remaining hypotheses hold for this torus.
In the lemma and in its iterations, the frame $\mathfrak{F}$ of $\mathcal{Q}$ is arbitrary.
For the coordinate frame $\mathfrak{F}_{\mathrm{c}}$, three applications with the same $k$,
one in each of the directions $x_1$, $x_2$, $x_3$,
lead from a cubic torus $C(n)$ through $C(2kh,n,n)$ and $C(2kh,2kh,n)$ to the cubic torus $C(2kh)$,
provided that the hypotheses of the lemma hold at each step.
\end{enumerate}
\end{remark}

\subsection{Tori of different sides}\label{sec:different-sides}

\begin{proposition}[The fold inequality for arbitrary sides]\label{prop:fold-any-side}
Let $-1\le\rho\le1$, $w_0,\dots,w_3>0$, $M\in\R$ and $q>0$.
Let $\mathfrak{n}'=(n'_i)$ be a sequence of sides with \eqref{cond:divisibility};
data, their thresholds and the perturbed functionals $\mu^t_i$ are formed with $\mathfrak{n}'$,
as in Definition~\ref{def:two-sequences}.
Let $\mathfrak{d}=(c_0,c_+,c_3;h;O)$ be data, let $i\ge i_0(\mathfrak{d})$,
write $n':=n'_i$,
and let $n\ge n'$ be an arbitrary integer.
Define the three numbers
\[
k_1:=\frac{n'}{2h},\qquad k_2:=\frac{n'}{h},\qquad k_3:=\frac{n'}{2h},
\]
the set of the regular cells and its cardinality
\[
\mathcal{K}^{\mathrm{reg}}:=\left\{\kappa\in\Z^3:\ 0\le\kappa_a\le 2k_a-2,\ a=1,2,3\right\},\qquad
\#|\mathcal{K}^{\mathrm{reg}}|=(2k_1-1)(2k_2-1)(2k_3-1),
\]
the translation vectors and the partial unfolded array
\[
v(\kappa):=\frac{h}{2}\,(\kappa_1e_-+\kappa_2e_+)+h\kappa_3e_3\quad(\kappa\in\Z^3),\qquad
\widetilde O_{\mathcal{K}^{\mathrm{reg}}}:=\sum_{\kappa\in\mathcal{K}^{\mathrm{reg}}}\tau_{v(\kappa)}O,
\]
the regular and the exceptional volume
\[
V^{\mathrm{reg}}:=\#|\mathcal{K}^{\mathrm{reg}}|\cdot|Q_0|=\frac{(n'-h)^2(2n'-h)}{2},\qquad
V^{\mathrm{exc}}:=n^3-V^{\mathrm{reg}},
\]
and the error term
\[
\Delta'(n,n',h):=f\left(\frac{n'}{2},2n',n'\right)-f_\infty+\frac{V^{\mathrm{exc}}}{V^{\mathrm{reg}}}\left(\bar f-f_\infty\right),
\]
where $f_\infty=\inf f$ and $\bar f$ are the two bounds of the pressure in Lemma~\ref{lem:uniform-ceiling}.
Then $k_1$, $k_2$, $k_3$ are integers $\ge2$, $n'^3/|Q_0|=8k_1k_2k_3$,
$0<V^{\mathrm{reg}}<n'^3\le n^3$, and $\Delta'(n,n',h)\ge0$.
Moreover, for every $t>0$ the following four assertions hold,
in which $H(\lambda\,|\,\mu)\in[0,\infty]$ is the relative entropy of Definition~\ref{def:relative-entropy}.
\begin{enumerate}
\item[\textup{(a$'$)}]
For every law $\lambda$ on $\Omega$ with $\lambda\ll\mu_{C(n)}$ and $\lambda(\tau_vO)=\lambda(O)$ for all $v\in\Lambda_h$,
\[
\log\frac{Z_{C(n)}\left[\exp\left(t\widetilde O_{\mathcal{K}^{\mathrm{reg}}}\right)\right]}{Z_{C(n)}[1]}
\ \ge\ t\cdot\#|\mathcal{K}^{\mathrm{reg}}|\cdot\lambda(O)-H(\lambda\,|\,\mu_{C(n)}).
\]
The law $\mu_{C(n)}$ itself is such a law, and $H(\mu_{C(n)}\,|\,\mu_{C(n)})=0$.
\item[\textup{(b$'$)}]
One has
\[
Z_{C(n)}\left[\exp\left(t\widetilde O_{\mathcal{K}^{\mathrm{reg}}}\right)\right]
\ \le\ \exp\left(\bar f\cdot V^{\mathrm{exc}}\right)\cdot
Z_{D(n')}\left[\exp(tO_{\mathfrak{d}})\right]^{\frac{\#|\mathcal{K}^{\mathrm{reg}}|\cdot|Q_0|}{n'^3}}.
\]
\item[\textup{(d$'$)}]
One has
\[
\mu_{C(n)}(O)\ \le\ \mu^t_i(O)+\frac{|Q_0|\cdot\Delta'(n,n',h)}{t}.
\]
\item[\textup{(e$'$)}]
For every law $\lambda$ as in \textup{(a$'$)},
\[
\lambda(O)\ \le\ \mu^t_i(O)+\frac{|Q_0|}{t}\left(\Delta'(n,n',h)+\frac{H(\lambda\,|\,\mu_{C(n)})}{V^{\mathrm{reg}}}\right).
\]
\end{enumerate}
\end{proposition}

\begin{proof}
Throughout, $\ell^1=\ell_-$, $\ell^2=\ell_+$, $\ell^3=x_3$ are the families of the diagonal frame,
and $\operatorname{ht}=(\ell^1,\ell^2,\ell^3)$ is the map of height coordinates of Lemma~\ref{lem:heights}.
\begin{enumerate}
\item\label{it:sides-numbers}
\emph{The numbers.}
As $i\ge i_0(\mathfrak{d})$, the integer $4h$ divides $n'$ (Definition~\ref{def:threshold}).
Thus $n'/h$ is a multiple of $4$,
and $k_1=k_3=n'/(2h)$ and $k_2=n'/h$ are integers $\ge2$.
For data one has $h\ge18$ (Definition~\ref{def:threshold}), whence $n'\ge4h\ge72$.
By Definition~\ref{def:data}, $|Q_0|=h^3/2$,
and the number $n'^3/|Q_0|$ of the cells of $\mathcal{Q}_{\mathfrak{d}}$ in $D(n')$ satisfies
\[
\frac{n'^3}{|Q_0|}=\frac{2n'^3}{h^3}=2k_1\cdot2k_2\cdot2k_3=8k_1k_2k_3 .
\]
From $(2k_1-1)h=(2k_3-1)h=n'-h$ and $(2k_2-1)h=2n'-h$ one obtains
\[
V^{\mathrm{reg}}=\#|\mathcal{K}^{\mathrm{reg}}|\cdot\frac{h^3}{2}=\frac{(n'-h)^2(2n'-h)}{2}.
\]
Since $0<n'-h<n'$ and $0<2n'-h<2n'$,
it follows that
\[
0<V^{\mathrm{reg}}<\frac{n'^2\cdot2n'}{2}=n'^3\le n^3 ;
\]
in particular $V^{\mathrm{exc}}>0$.
The sides of $C(n'/2,2n',n')$ are at least $36$,
so $f(n'/2,2n',n')\ge f_\infty$ by Proposition~\ref{prop:coordinate-pressure}(d),
and $\bar f\ge f_\infty$ by Lemma~\ref{lem:uniform-ceiling}.
As $V^{\mathrm{exc}}/V^{\mathrm{reg}}>0$, it follows that $\Delta'(n,n',h)\ge0$.
For the rest of the proof let
\[
s:=n-n'\ge0,\qquad
\mathcal{K}_0:=\left\{\kappa\in\Z^3:\ 0\le\kappa_a\le 2k_a-1,\ a=1,2,3\right\},\qquad
\#|\mathcal{K}_0|=8k_1k_2k_3=\frac{n'^3}{|Q_0|},
\]
and let $\Gamma_i$ be the period lattice of $D(n')$.
Recall that $C(n)=C(n,n,n)$ is the cubic torus of side $n$ and $\mu_{C(n)}$ its law.

\item\label{it:sides-contents}
\emph{The contents.}
Fix $t>0$, and let $\kappa\in\Z^3$.
Because $h$ is even, the vectors $\frac{h}{2}e_-$, $\frac{h}{2}e_+$ and $he_3$ are integral,
and their height coordinates are
\[
\operatorname{ht}\left(\frac{h}{2}e_-\right)=(h,0,0),\qquad
\operatorname{ht}\left(\frac{h}{2}e_+\right)=(0,h,0),\qquad
\operatorname{ht}(he_3)=(0,0,h).
\]
By linearity $\operatorname{ht}(v(\kappa))=h\kappa$;
so Lemma~\ref{lem:wall-group}(d) gives $v(\kappa)\in\Lambda_h$ and $\tau_{v(\kappa)}Q_0=Q_\kappa$.
Put
\[
G_0:=\exp(tO),\qquad G_\kappa:=\exp\left(t\,\tau_{v(\kappa)}O\right)\quad(\kappa\in\Z^3).
\]
By Definition~\ref{def:characters}, $(\tau_vO)(F)=O(\tau_v^*F)$.
A translation is an even element of $W$, so the rule \eqref{eq:setting-group-action} gives
\[
(G_0)^{\tau_{v(\kappa)}}(F)=G_0\left(\tau_{v(\kappa)}^*F\right)=\exp\left(t\,O\left(\tau_{v(\kappa)}^*F\right)\right)=G_\kappa(F),
\]
that is, $G_\kappa=(G_0)^{\tau_{v(\kappa)}}$.
The polynomial $\tau_{v(\kappa)}O$ reads the plaquettes of $\supp O+v(\kappa)$ (Definition~\ref{def:characters}),
and these lie in $Q_0+v(\kappa)=Q_\kappa$, since $\supp O$ lies in $Q_0$ (Definition~\ref{def:data}).
Hence $G_\kappa$ is a positive, bounded and continuous function of finitely many $F_\sigma$, all $\sigma$ contained in $Q_\kappa$;
in particular $G_\kappa\in\mathcal{A}(Q_\kappa)$.
Consequently
\[
\mathfrak{G}:=(G_\kappa)_{\kappa\in\mathcal{K}^{\mathrm{reg}}}
\]
is a finite filling of $\mathcal{Q}_{\mathfrak{d}}$ (Definition~\ref{def:fillings}), and
\[
\prod\mathfrak{G}=\prod_{\kappa\in\mathcal{K}^{\mathrm{reg}}}G_\kappa=\exp\left(t\widetilde O_{\mathcal{K}^{\mathrm{reg}}}\right),\qquad
\left|\widetilde O_{\mathcal{K}^{\mathrm{reg}}}\right|\le\#|\mathcal{K}^{\mathrm{reg}}|\cdot|O|_1,
\]
the bound because each term $\tau_{v(\kappa)}O$ is bounded by $|O|_1$.

\item\label{it:sides-full-fold}
\emph{The full fold of one content.}
For every $j\in\Z^3$,
\[
Z_{D(n')}\left[\prod_{\kappa\in j+\mathcal{K}_0}(G_j)^{[\kappa]}\right]=Z_{D(n')}\left[\exp(tO_{\mathfrak{d}})\right],
\]
where $(G_j)^{[\kappa]}$ is the mirror copy in the cell $\kappa$ of $G_j\in\mathcal{A}(Q_j)$ (Definition~\ref{def:folds}).
To see this, fix $j$.
The torus $D(n')=\mathbb{T}(\mathfrak{F}_{\mathrm{d}};n',2n',n')$ carries $\mathcal{Q}_{\mathfrak{d}}$, because $n'$ and $2n'$ lie in $2h\Z$,
and its cells are the $\Gamma_i$-orbits of the cells $Q_\kappa$ (Definition~\ref{def:cell-structures}).
Since $n'$ is a multiple of $2h$,
Lemma~\ref{lem:wall-group}(f) shows that the generators $\frac{n'}{2}e_-$, $n'e_+$, $n'e_3$ of $\Gamma_i$ have the heights
$(n',0,0)$, $(0,2n',0)$, $(0,0,n')$,
and by Lemma~\ref{lem:wall-group}(d) the translations by them shift the indices of the cells by
$(2k_1,0,0)$, $(0,2k_2,0)$, $(0,0,2k_3)$.
Hence $j+\mathcal{K}_0$ is a set of representatives of the cells of $D(n')$.

Let $g\in\Gamma_i$ and $\kappa\in\Z^3$.
By Lemma~\ref{lem:wall-group}(f), $g$ lies in $2\Lambda_h$,
so that $\tau_g\in W(\mathcal{Q}_{\mathfrak{d}})$ by Lemma~\ref{lem:wall-group}(c);
moreover $\operatorname{ht}(g)=h\kappa_g$ with an index $\kappa_g\in\Z^3$,
and $\tau_gQ_\kappa=Q_{\kappa+\kappa_g}$, by Lemma~\ref{lem:wall-group}(d).
If $g'\in W(\mathcal{Q}_{\mathfrak{d}})$ maps $Q_j$ onto $Q_\kappa$,
then $\tau_gg'\in W(\mathcal{Q}_{\mathfrak{d}})$ maps $Q_j$ onto $Q_{\kappa+\kappa_g}$.
As $W(\mathcal{Q}_{\mathfrak{d}})$ acts simply transitively on the cells (Definition~\ref{def:cell-structures}),
Definition~\ref{def:folds} and the rule \eqref{eq:setting-group-action} give
\[
(G_j)^{[\kappa+\kappa_g]}=(G_j)^{\tau_gg'}=\left((G_j)^{g'}\right)^{\tau_g}=\left((G_j)^{[\kappa]}\right)^{\tau_g}.
\]
For a bounded Borel function $\Phi$ of finitely many $F_\sigma$ and a $\Gamma_i$-periodic $\widetilde F$
one has $\Phi^{\tau_g}\left(\widetilde F\right)=\Phi\left(\tau_g^*\widetilde F\right)$ by \eqref{eq:setting-group-action},
and $\Phi\left(\tau_g^*\widetilde F\right)=\Phi\left(\widetilde F\right)$ because $\tau_g^*\widetilde F=\widetilde F$.
So the function $\prod_{\kappa\in j+\mathcal{K}_0}(G_j)^{[\kappa]}$ on $\Conf(\Gamma_i)$ does not change
when $j+\mathcal{K}_0$ is replaced by any other set of representatives of the cells of $D(n')$:
it is the fold $\Fold_{D(n')}G_j$ of Definition~\ref{def:folds}, and
\[
Z_{D(n')}\left[\prod_{\kappa\in j+\mathcal{K}_0}(G_j)^{[\kappa]}\right]
=Z_{D(n')}\left[\Fold_{D(n')}G_j\right]
=\|G_j\|_{D(n')}^{\frac{n'^3}{|Q_0|}}.
\]
Now $G_j=(G_0)^{\tau_{v(j)}}$ by step~\ref{it:sides-contents},
the heights of $v(j)$ are $\ell^a(v(j))=hj_a\in h\Z$ for $a=1,2,3$,
and $\tau_{v(j)}D(n')=D(n')$ by Lemma~\ref{lem:wall-group}(f).
Lemma~\ref{lem:cell-symmetries} gives $\|G_j\|_{D(n')}=\|G_0\|_{D(n')}$,
and Proposition~\ref{prop:fold-finite}(b), read for the sequence $\mathfrak{n}'$ (Definition~\ref{def:two-sequences}), gives
\[
\|G_0\|_{D(n')}^{\frac{n'^3}{|Q_0|}}=Z_{D(n')}\left[\exp(tO_{\mathfrak{d}})\right].
\]
Together with the preceding display, this is the identity asserted at the beginning of the step.
\end{enumerate}

\emph{The assertion \textup{(a$'$)}.}
By step~\ref{it:sides-contents}, $\exp\left(t\widetilde O_{\mathcal{K}^{\mathrm{reg}}}\right)$ is a bounded continuous function of finitely many $F_\sigma$,
and Lemma~\ref{lem:bosonic-measure}(b) turns the quotient on the left of \textup{(a$'$)} into an expectation:
\[
\frac{Z_{C(n)}\left[\exp\left(t\widetilde O_{\mathcal{K}^{\mathrm{reg}}}\right)\right]}{Z_{C(n)}[1]}
=\mu_{C(n)}\left(\exp\left(t\widetilde O_{\mathcal{K}^{\mathrm{reg}}}\right)\right).
\]
Let $\lambda$ be a law on $\Omega$ with the two properties required in \textup{(a$'$)}.
Lemma~\ref{lem:jensen-entropy}, applied to the probability measures $\mu_{C(n)}$ and $\lambda\ll\mu_{C(n)}$ on $\Omega$
and to the bounded real function $t\widetilde O_{\mathcal{K}^{\mathrm{reg}}}$, gives
\[
\begin{aligned}
\log\mu_{C(n)}\left(\exp\left(t\widetilde O_{\mathcal{K}^{\mathrm{reg}}}\right)\right)
&\ge t\cdot\lambda\left(\widetilde O_{\mathcal{K}^{\mathrm{reg}}}\right)-H(\lambda\,|\,\mu_{C(n)})\\
&=t\sum_{\kappa\in\mathcal{K}^{\mathrm{reg}}}\lambda\left(\tau_{v(\kappa)}O\right)-H(\lambda\,|\,\mu_{C(n)})\\
&=t\cdot\#|\mathcal{K}^{\mathrm{reg}}|\cdot\lambda(O)-H(\lambda\,|\,\mu_{C(n)}),
\end{aligned}
\]
the last line because every $v(\kappa)$ lies in $\Lambda_h$ (step~\ref{it:sides-contents}),
where $\lambda(\tau_vO)=\lambda(O)$ by hypothesis.
Combined with the expression of the quotient as $\mu_{C(n)}\left(\exp\left(t\widetilde O_{\mathcal{K}^{\mathrm{reg}}}\right)\right)$,
this is the inequality of \textup{(a$'$)}.
As for $\mu_{C(n)}$ itself:
every translation by a vector of $\Z^3$ maps every torus onto itself (Lemma~\ref{lem:wall-group}(f)),
so $\tau_vC(n)=C(n)$, and $\mu_{C(n)}$ is invariant under $F\mapsto\tau_v^*F$ for every $v\in\Z^3$ (Lemma~\ref{lem:bosonic-measure}(c)).
Hence $\mu_{C(n)}(\tau_vO)=\mu_{C(n)}(O)$ for all $v\in\Z^3$, in particular for all $v\in\Lambda_h$.
Finally $\mu_{C(n)}\ll\mu_{C(n)}$ with the density $1$, so that $H(\mu_{C(n)}\,|\,\mu_{C(n)})=\mu_{C(n)}(\log1)=0$.

\emph{The assertion} \textup{(b$'$)}.
The inequality \textup{(b$'$)} results from three applications of Lemma~\ref{lem:one-step}
to the cell structure $\mathcal{Q}_{\mathfrak{d}}=(\mathfrak{F}_{\mathrm{d}};c_0,c_+,c_3;h)$ of the data,
successively in the directions $a=1,2,3$ of the diagonal frame~\eqref{eq:fold-frame}.
The torus of the first application is $\mathbb{T}_0:=C(n)$,
and each later application starts from the regular torus produced by the one before it.
In each of them $\Xi=(\Gamma_\parallel,\chi_\parallel)$ and $p$ are the cross-section and the height period of the torus for the family $\ell^a$ (Definition~\ref{def:slabs}).
The two tori of the lemma are $\mathbb{T}^{\mathrm{reg}}=\mathbb{T}(\Xi;2kh)$ and $\mathbb{T}^{\mathrm{exc}}=\mathbb{T}(\Xi;2kh^{\mathrm{exc}})$.
By Lemma~\ref{lem:tori-over-section}, for $P=2kh$ and for $P=2kh^{\mathrm{exc}}$ the torus $\mathbb{T}(\Xi;P)$ has the lattice $\Gamma_\parallel+\Z u$,
where $u$ is the integer vector normal to the planes of $\ell^a$ with $\ell^a(u)=P$,
and its character equals $\chi_\parallel$ on $\Gamma_\parallel$ and $-1$ at $u$.
The two conditions on systoles in Lemma~\ref{lem:one-step} are checked in the paragraph that follows the three steps.
\begin{enumerate}
\setcounter{enumi}{3}
\item\label{it:sides-first-step}
\emph{The direction $a=1$: the family $\ell_-$ on $\mathbb{T}_0=C(n)$, with $2k_1$ slabs.}
For integers $\alpha_1,\alpha_2,\alpha_3$ one has $\ell_-\left(n(\alpha_1,\alpha_2,\alpha_3)\right)=n(\alpha_1-\alpha_2)$.
The lattice of $\Xi(C(n),\ell_-)$ is therefore
\[
\Gamma_\parallel=\left\{n(\alpha_1,\alpha_1,\alpha_3):\ \alpha_1,\alpha_3\in\Z\right\}=\Z\,ne_++\Z\,ne_3,
\]
and by Definition~\ref{def:tori} its character takes the values
$\chi_\parallel(ne_+)=\chi(ne_1)\cdot\chi(ne_2)=(-1)(-1)=+1$ and $\chi_\parallel(ne_3)=-1$.
Moreover $\ell_-(n\Z^3)=n\Z$, so that $p=n$, for odd as well as for even $n$.
Since $(2k_1-1)h=n'-h<n$, the exceptional slab has the height
\[
h^{\mathrm{exc}}=n-(2k_1-1)h=h+s.
\]
The vector $u$ is $(2k_1h/2)\,e_-=(n'/2)\,e_-$ for $\mathbb{T}^{\mathrm{reg}}$, and $(2k_1h^{\mathrm{exc}}/2)\,e_-=k_1(h+s)\,e_-$ for $\mathbb{T}^{\mathrm{exc}}$.
Hence $\mathbb{T}^{\mathrm{reg}}$ has the lattice $\Z\,(n'/2)\,e_-+\Z\,ne_++\Z\,ne_3$,
with the character values $-1$, $+1$, $-1$ at these three generators,
and $\mathbb{T}^{\mathrm{exc}}$ arises from it when $n'/2$ is replaced by $k_1(h+s)$.
By Definition~\ref{def:tori},
\[
\mathbb{T}^{\mathrm{reg}}=D'(n'/2,n,n)=:\mathbb{T}_1,
\qquad
\mathbb{T}^{\mathrm{exc}}=D'\left(k_1(h+s),n,n\right)=:\mathbb{T}_1^{\mathrm{exc}}.
\]
\item\label{it:sides-second-step}
\emph{The direction $a=2$: the family $\ell_+$ on $\mathbb{T}_1$, with $2k_2$ slabs.}
Here $\ell_+(e_-)=0=\ell_+(e_3)$ and $\ell_+(e_+)=2$.
In the lattice $\Z\,(n'/2)\,e_-+\Z\,ne_++\Z\,ne_3$ of $\mathbb{T}_1$ the vectors on which $\ell_+$ vanishes form
\[
\Gamma_\parallel=\Z\,(n'/2)\,e_-+\Z\,ne_3,
\]
with $\chi_\parallel=-1$ at both generators,
and the height period is $p=\ell_+(ne_+)=2n$.
The closing vector $ne_+$ has $\chi(ne_+)=+1$,
so that the twist in the trace formula of Theorem~\ref{thm:slab-operators}\eqref{it:slab-trace} for $\mathbb{T}_1$ and $\ell_+$ contains the fermion parity $U_{\mathrm{f}}$.
Since $(2k_2-1)h=2n'-h<2n$,
\[
h^{\mathrm{exc}}=2n-(2k_2-1)h=h+2s,
\]
and the vector $u$ is $(2k_2h/2)\,e_+=n'e_+$ for $\mathbb{T}^{\mathrm{reg}}$ and $k_2(h+2s)\,e_+$ for $\mathbb{T}^{\mathrm{exc}}$.
Both tori take the value $-1$ at all three generators, and
\[
\mathbb{T}^{\mathrm{reg}}=D(n'/2,n',n)=:\mathbb{T}_2,
\qquad
\mathbb{T}^{\mathrm{exc}}=D\left(n'/2,k_2(h+2s),n\right)=:\mathbb{T}_2^{\mathrm{exc}}.
\]
\item\label{it:sides-third-step}
\emph{The direction $a=3$: the family $x_3$ on $\mathbb{T}_2$, with $2k_3$ slabs.}
The lattice of $\mathbb{T}_2$ is $\Z\,(n'/2)\,e_-+\Z\,n'e_++\Z\,ne_3$.
Its vectors with $x_3=0$ form $\Gamma_\parallel=\Z\,(n'/2)\,e_-+\Z\,n'e_+$, with $\chi_\parallel=-1$ at both generators,
and $p=x_3(ne_3)=n$.
Since $(2k_3-1)h=n'-h<n$,
\[
h^{\mathrm{exc}}=n-(2k_3-1)h=h+s,
\]
the vector $u$ is $2k_3h\,e_3=n'e_3$ for $\mathbb{T}^{\mathrm{reg}}$ and $2k_3(h+s)\,e_3$ for $\mathbb{T}^{\mathrm{exc}}$, and
\[
\mathbb{T}^{\mathrm{reg}}=D(n'/2,n',n')=D(n')=:\mathbb{T}_3,
\qquad
\mathbb{T}^{\mathrm{exc}}=D\left(n'/2,n',2k_3(h+s)\right)=:\mathbb{T}_3^{\mathrm{exc}}.
\]
\end{enumerate}

\emph{Systoles.}
By Lemma~\ref{lem:systoles}, for positive integers $a_1,a_2,a_3$ the lattice $\Z a_1\,e_-+\Z a_2\,e_++\Z a_3\,e_3$ has the systole $\min(a_1,a_2,a_3)$,
and the lattice $n\Z^3$ has the systole $n$.
The lattice of $\mathbb{T}_0=C(n)$ is $n\Z^3$,
and the lattices of the six tori $\mathbb{T}_1,\mathbb{T}_2,\mathbb{T}_3$ and $\mathbb{T}_1^{\mathrm{exc}},\mathbb{T}_2^{\mathrm{exc}},\mathbb{T}_3^{\mathrm{exc}}$
are of the form $\Z a_1\,e_-+\Z a_2\,e_++\Z a_3\,e_3$, with the triples $(a_1,a_2,a_3)$ found in the three steps.
As $k_1\ge2$, $k_2=2k_1\ge4$ and $k_3\ge2$,
each of these seven lattices has systole at least
\[
\min\left(\frac{n'}{2},\ 2(h+s),\ n\right)\ge36,
\]
because $n\ge n'\ge72$ and $h\ge18$ (step~\ref{it:sides-numbers}).
In particular $\mathbb{T}_1$ and $\mathbb{T}_2$ are tori,
and in each of the three steps the lattices of $\mathbb{T}^{\mathrm{reg}}$ and of $\mathbb{T}^{\mathrm{exc}}$ have systole $\ge4$, as Lemma~\ref{lem:one-step} requires.

\emph{The iterated folds.}
The filling $\mathfrak{G}$ satisfies the hypothesis of Lemma~\ref{lem:one-step} in step~\ref{it:sides-first-step},
since $0\le\kappa_1\le2k_1-2$ for every $\kappa\in\mathcal{K}^{\mathrm{reg}}$.
For $j=(j_1,j_2,j_3)\in\mathcal{K}^{\mathrm{reg}}$ define
\[
\mathfrak{G}^{j_1}:=\mathfrak{R}_1^{j_1,2k_1}\mathfrak{G},
\qquad
\mathfrak{G}^{j_1j_2}:=\mathfrak{R}_2^{j_2,2k_2}\mathfrak{G}^{j_1},
\qquad
\mathfrak{G}^{j_1j_2j_3}:=\mathfrak{R}_3^{j_3,2k_3}\mathfrak{G}^{j_1j_2};
\]
these are again finite fillings of $\mathcal{Q}_{\mathfrak{d}}$ (Remark~\ref{rem:one-step}(c)).
By Definition~\ref{def:fillings}, the filling $\mathfrak{G}^{j_1}$ is indexed by the $\kappa\in\Z^3$ with
\[
j_1\le\kappa_1\le j_1+2k_1-1,\qquad 0\le\kappa_2\le2k_2-2,\qquad 0\le\kappa_3\le2k_3-2,
\]
and in the cell $\kappa$ it has the content $\left(G_{\kappa^{1\to j_1}}\right)^{g^1_{\kappa_1\leftarrow j_1}}$.
All second entries of these indices lie between $0$ and $2k_2-2$, which is the hypothesis of Lemma~\ref{lem:one-step} in step~\ref{it:sides-second-step}.
The filling $\mathfrak{G}^{j_1j_2}$ is indexed by the $\kappa$ with $j_a\le\kappa_a\le j_a+2k_a-1$ for $a=1,2$ and $0\le\kappa_3\le2k_3-2$,
so that the hypothesis of step~\ref{it:sides-third-step} holds as well.
Finally, $\mathfrak{G}^{j_1j_2j_3}$ is indexed by $j+\mathcal{K}_0$,
and at $\kappa\in j+\mathcal{K}_0$ its content is the threefold image
\[
\left(\left(\left(G_j\right)^{g^1_{\kappa_1\leftarrow j_1}}\right)^{g^2_{\kappa_2\leftarrow j_2}}\right)^{g^3_{\kappa_3\leftarrow j_3}}
=\left(G_j\right)^{g^3_{\kappa_3\leftarrow j_3}\,g^2_{\kappa_2\leftarrow j_2}\,g^1_{\kappa_1\leftarrow j_1}}
=\left(G_j\right)^{[\kappa]}.
\]
The first equality is the composition rule $(G^g)^{g'}=G^{g'g}$ of~\eqref{eq:setting-group-action}.
The second holds because $g^3_{\kappa_3\leftarrow j_3}g^2_{\kappa_2\leftarrow j_2}g^1_{\kappa_1\leftarrow j_1}$
is the element of $W(\mathcal{Q}_{\mathfrak{d}})$ that maps $Q_j$ onto $Q_\kappa$ (Definition~\ref{def:cell-families}),
and $(G_j)^{[\kappa]}$ is the image of $G_j\in\mathcal{A}(Q_j)$ under this element (Definition~\ref{def:folds}).
So $\prod\mathfrak{G}^{j_1j_2j_3}=\prod_{\kappa\in j+\mathcal{K}_0}(G_j)^{[\kappa]}$,
and step~\ref{it:sides-full-fold} shows that
\[
Z_{D(n')}\left[\prod\mathfrak{G}^{j_1j_2j_3}\right]=Z_{D(n')}\left[\exp(tO_{\mathfrak{d}})\right]
\qquad\text{for every }j\in\mathcal{K}^{\mathrm{reg}}.
\]

\emph{The chain of inequalities.}
The number $Z_{\mathbb{T}_0}\left[\prod\mathfrak{G}\right]=Z_{C(n)}\left[\exp\left(t\widetilde O_{\mathcal{K}^{\mathrm{reg}}}\right)\right]$ is positive.
Indeed, by Lemma~\ref{lem:bosonic-measure}\eqref{it:bosonic-integral} it is the integral of the positive function $\exp\left(t\widetilde O_{\mathcal{K}^{\mathrm{reg}}}\right)$, evaluated at $\widetilde F$,
against the measure of Lemma~\ref{lem:bosonic-measure}\eqref{it:bosonic-mass},
and that measure is a finite positive measure with the positive total mass $Z_{C(n)}[1]$.
The first of the three lines below is therefore~\eqref{eq:sides-step} in step~\ref{it:sides-first-step}.
The second is~\eqref{eq:sides-step} in step~\ref{it:sides-second-step}, for the filling $\mathfrak{G}^{j_1}$;
its left side is non-negative by~\eqref{eq:sides-step-sign} in step~\ref{it:sides-first-step}.
The third is~\eqref{eq:sides-step} in step~\ref{it:sides-third-step}, for the filling $\mathfrak{G}^{j_1j_2}$,
whose left side is non-negative by~\eqref{eq:sides-step-sign} in step~\ref{it:sides-second-step}:
\[
\begin{aligned}
Z_{C(n)}\left[\exp\left(t\widetilde O_{\mathcal{K}^{\mathrm{reg}}}\right)\right]
&\le Z_{\mathbb{T}_1^{\mathrm{exc}}}[1]^{\frac{1}{2k_1}}\cdot\prod_{j_1=0}^{2k_1-2}Z_{\mathbb{T}_1}\left[\prod\mathfrak{G}^{j_1}\right]^{\frac{1}{2k_1}},\\
Z_{\mathbb{T}_1}\left[\prod\mathfrak{G}^{j_1}\right]
&\le Z_{\mathbb{T}_2^{\mathrm{exc}}}[1]^{\frac{1}{2k_2}}\cdot\prod_{j_2=0}^{2k_2-2}Z_{\mathbb{T}_2}\left[\prod\mathfrak{G}^{j_1j_2}\right]^{\frac{1}{2k_2}},\\
Z_{\mathbb{T}_2}\left[\prod\mathfrak{G}^{j_1j_2}\right]
&\le Z_{\mathbb{T}_3^{\mathrm{exc}}}[1]^{\frac{1}{2k_3}}\cdot\prod_{j_3=0}^{2k_3-2}Z_{\mathbb{T}_3}\left[\prod\mathfrak{G}^{j_1j_2j_3}\right]^{\frac{1}{2k_3}}.
\end{aligned}
\]
Every number in these lines is non-negative, by~\eqref{eq:sides-step-sign}, and every exponent is positive.
The bound of the third line may therefore be substituted into the second line, and the bound so obtained into the first.
This produces $\#|\mathcal{K}^{\mathrm{reg}}|$ regular factors, one for each $j\in\mathcal{K}^{\mathrm{reg}}$,
and every one of them is $Z_{D(n')}\left[\exp(tO_{\mathfrak{d}})\right]^{\frac{1}{8k_1k_2k_3}}$, where $8k_1k_2k_3=n'^3/|Q_0|$ (step~\ref{it:sides-numbers}).
The factor $Z_{\mathbb{T}_2^{\mathrm{exc}}}[1]$ occurs $2k_1-1$ times, with the exponent $\frac{1}{4k_1k_2}$ each time,
and $Z_{\mathbb{T}_3^{\mathrm{exc}}}[1]$ occurs $(2k_1-1)(2k_2-1)$ times, with the exponent $\frac{1}{8k_1k_2k_3}$.
Let
\[
\vartheta_1:=\frac{1}{2k_1},\qquad
\vartheta_2:=\frac{2k_1-1}{4k_1k_2},\qquad
\vartheta_3:=\frac{(2k_1-1)(2k_2-1)}{8k_1k_2k_3}.
\]
Then
\begin{equation}\label{eq:sides-three-steps}
\begin{split}
Z_{C(n)}\left[\exp\left(t\widetilde O_{\mathcal{K}^{\mathrm{reg}}}\right)\right]
\le{}&Z_{\mathbb{T}_1^{\mathrm{exc}}}[1]^{\vartheta_1}\cdot Z_{\mathbb{T}_2^{\mathrm{exc}}}[1]^{\vartheta_2}\cdot Z_{\mathbb{T}_3^{\mathrm{exc}}}[1]^{\vartheta_3}\\
&\cdot Z_{D(n')}\left[\exp(tO_{\mathfrak{d}})\right]^{\frac{\#|\mathcal{K}^{\mathrm{reg}}|\cdot|Q_0|}{n'^3}}.
\end{split}
\end{equation}

\emph{The exceptional tori.}
The identity~\eqref{eq:sides-step-volumes}, written once in each step, reads
\[
\begin{aligned}
n^3=|\mathbb{T}_0|&=\frac{|\mathbb{T}_1^{\mathrm{exc}}|}{2k_1}+\frac{2k_1-1}{2k_1}\cdot|\mathbb{T}_1|,\\
|\mathbb{T}_1|&=\frac{|\mathbb{T}_2^{\mathrm{exc}}|}{2k_2}+\frac{2k_2-1}{2k_2}\cdot|\mathbb{T}_2|,\\
|\mathbb{T}_2|&=\frac{|\mathbb{T}_3^{\mathrm{exc}}|}{2k_3}+\frac{2k_3-1}{2k_3}\cdot|\mathbb{T}_3|,
\end{aligned}
\]
and $|\mathbb{T}_3|=|D(n')|=n'^3$ by Definition~\ref{def:tori}.
Eliminating $|\mathbb{T}_2|$ and then $|\mathbb{T}_1|$ from these identities, one finds
\[
n^3=\vartheta_1\cdot|\mathbb{T}_1^{\mathrm{exc}}|+\vartheta_2\cdot|\mathbb{T}_2^{\mathrm{exc}}|+\vartheta_3\cdot|\mathbb{T}_3^{\mathrm{exc}}|
+\frac{\#|\mathcal{K}^{\mathrm{reg}}|}{8k_1k_2k_3}\cdot n'^3,
\]
and the last term equals $\#|\mathcal{K}^{\mathrm{reg}}|\cdot|Q_0|=V^{\mathrm{reg}}$.
Since $V^{\mathrm{exc}}=n^3-V^{\mathrm{reg}}$,
\begin{equation}\label{eq:sides-exceptional}
\vartheta_1\cdot|\mathbb{T}_1^{\mathrm{exc}}|+\vartheta_2\cdot|\mathbb{T}_2^{\mathrm{exc}}|+\vartheta_3\cdot|\mathbb{T}_3^{\mathrm{exc}}|=V^{\mathrm{exc}}.
\end{equation}
Lemma~\ref{lem:uniform-ceiling} gives $Z_{\mathbb{T}_a^{\mathrm{exc}}}[1]^{\vartheta_a}\le\exp\left(\bar f\cdot\vartheta_a\cdot|\mathbb{T}_a^{\mathrm{exc}}|\right)$ for $a=1,2,3$.
Inserting these three bounds into~\eqref{eq:sides-three-steps} and using~\eqref{eq:sides-exceptional}, one arrives at \textup{(b$'$)}.

\emph{The assertions} \textup{(d$'$)} \emph{and} \textup{(e$'$)}.
Every partition function that occurs in the estimate below is strictly positive.
For $Z_{C(n)}[1]$, $Z_{D(n')}[1]$ and $Z_{C(n'/2,2n',n')}[1]$ this is Theorem~\ref{thm:partition-positivity}.
Moreover $\left|\widetilde O_{\mathcal{K}^{\mathrm{reg}}}\right|\le\#|\mathcal{K}^{\mathrm{reg}}|\cdot|O|_1$ by step~\ref{it:sides-contents},
and $|O_{\mathfrak{d}}|\le\left(n'^3/|Q_0|\right)\cdot|O|_1$ on $D(n')$, which carries $\mathcal{Q}_{\mathfrak{d}}$ (Definition~\ref{def:arrays});
so Lemma~\ref{lem:perturbed-denominators} gives
\[
\begin{aligned}
Z_{C(n)}\left[\exp\left(t\widetilde O_{\mathcal{K}^{\mathrm{reg}}}\right)\right]
&\ge\exp\left(-t\cdot\#|\mathcal{K}^{\mathrm{reg}}|\cdot|O|_1\right)\cdot Z_{C(n)}[1]>0,\\
Z_{D(n')}\left[\exp(tO_{\mathfrak{d}})\right]
&\ge\exp\left(-t\cdot\frac{n'^3}{|Q_0|}\cdot|O|_1\right)\cdot Z_{D(n')}[1]>0.
\end{aligned}
\]
Let $\lambda$ be a law as in \textup{(a$'$)}.
If $H(\lambda\,|\,\mu_{C(n)})=\infty$, the right side of \textup{(e$'$)} is infinite and there is nothing to prove;
let therefore $H(\lambda\,|\,\mu_{C(n)})<\infty$.
Then
\[
\begin{aligned}
&t\cdot\#|\mathcal{K}^{\mathrm{reg}}|\cdot\lambda(O)-H(\lambda\,|\,\mu_{C(n)})\\
&\quad\le\log Z_{C(n)}\left[\exp\left(t\widetilde O_{\mathcal{K}^{\mathrm{reg}}}\right)\right]-\log Z_{C(n)}[1]\\
&\quad\le\bar f\cdot V^{\mathrm{exc}}+\frac{\#|\mathcal{K}^{\mathrm{reg}}|\cdot|Q_0|}{n'^3}\cdot\log Z_{D(n')}\left[\exp(tO_{\mathfrak{d}})\right]-n^3\cdot f(n,n,n)\\
&\quad\le\bar f\cdot V^{\mathrm{exc}}+\frac{\#|\mathcal{K}^{\mathrm{reg}}|\cdot|Q_0|}{n'^3}\cdot\left(t\cdot\frac{n'^3}{|Q_0|}\cdot\mu^t_i(O)+\log Z_{D(n')}[1]\right)-n^3\cdot f_\infty\\
&\quad\le\bar f\cdot V^{\mathrm{exc}}+t\cdot\#|\mathcal{K}^{\mathrm{reg}}|\cdot\mu^t_i(O)+V^{\mathrm{reg}}\cdot f\left(\frac{n'}{2},2n',n'\right)\\
&\qquad\quad-\left(V^{\mathrm{reg}}+V^{\mathrm{exc}}\right)\cdot f_\infty.
\end{aligned}
\]
The first inequality is \textup{(a$'$)}.
The second is \textup{(b$'$)}, combined with $\log Z_{C(n)}[1]=n^3\cdot f(n,n,n)$ (Definition~\ref{def:sides}).
The third rests on two facts:
Proposition~\ref{prop:fold-finite}(c), read for the sequence $\mathfrak{n}'$ (Definition~\ref{def:two-sequences}),
bounds $\log Z_{D(n')}\left[\exp(tO_{\mathfrak{d}})\right]$ by $t\cdot\left(n'^3/|Q_0|\right)\cdot\mu^t_i(O)+\log Z_{D(n')}[1]$;
and $f(n,n,n)\ge f_\infty$, because $f_\infty=\inf f$ (Proposition~\ref{prop:coordinate-pressure}(d)).
In the last inequality, Theorem~\ref{thm:four-tori}, applied to the even side $n'\ge8$, yields
\[
\log Z_{D(n')}[1]\le\log Z_{C(n'/2,2n',n')}[1]=n'^3\cdot f\left(\frac{n'}{2},2n',n'\right),
\]
since $|C(n'/2,2n',n')|=n'^3$;
moreover $\frac{\#|\mathcal{K}^{\mathrm{reg}}|\cdot|Q_0|}{n'^3}\cdot n'^3=V^{\mathrm{reg}}$ and $n^3=V^{\mathrm{reg}}+V^{\mathrm{exc}}$.
Comparing the two ends of the chain,
\[
\begin{aligned}
t\cdot\#|\mathcal{K}^{\mathrm{reg}}|\cdot\lambda(O)
\le{}&t\cdot\#|\mathcal{K}^{\mathrm{reg}}|\cdot\mu^t_i(O)
+V^{\mathrm{reg}}\cdot\left(f\left(\frac{n'}{2},2n',n'\right)-f_\infty\right)\\
&+V^{\mathrm{exc}}\cdot\left(\bar f-f_\infty\right)
+H(\lambda\,|\,\mu_{C(n)}).
\end{aligned}
\]
Division by the positive number $t\cdot\#|\mathcal{K}^{\mathrm{reg}}|=t\cdot V^{\mathrm{reg}}/|Q_0|$ gives \textup{(e$'$)}.
The law $\mu_{C(n)}$ is admissible in \textup{(a$'$)}, with $H(\mu_{C(n)}\,|\,\mu_{C(n)})=0$,
and for $\lambda=\mu_{C(n)}$ the inequality \textup{(e$'$)} becomes \textup{(d$'$)}.
\end{proof}

\begin{remark}\label{rem:fold-any-side}
In the notation of Proposition~\ref{prop:fold-any-side} let $s=n-n'$,
and let $\mathbb{T}_a^{\mathrm{exc}}$ and $\vartheta_a$, $a=1,2,3$, be the exceptional tori and the exponents of its proof.
\begin{enumerate}
\item
For $\lambda=\mu_{C(n)}$, the inequality \textup{(a$'$)} of that proposition is the inequality of Proposition~\ref{prop:fold-finite}(a),
written for the cubic torus $C(n)$ and the partial unfolded array $\widetilde O_{\mathcal{K}^{\mathrm{reg}}}$;
\textup{(b$'$)} and \textup{(d$'$)} are the analogues of Proposition~\ref{prop:fold-finite}(b), (d) for a cubic torus whose side $n\ge n'$ is arbitrary.
For $n=n'$, that is $s=0$, the partial unfolded array still omits the last layer of cells in each of the three directions,
and \textup{(d$'$)} is weaker than Proposition~\ref{prop:fold-finite}(d) for the sequence $\mathfrak{n}'$.
Indeed, the error term of the latter is $|Q_0|\cdot\Delta_i/t$ with $\Delta_i=f(n'/2,2n',n')-f(n',n',n')$, that of \textup{(d$'$)} is $|Q_0|\cdot\Delta'(n',n',h)/t$, and
\[
\Delta'(n',n',h)=\Delta_i+\left(f(n',n',n')-f_\infty\right)+\frac{V^{\mathrm{exc}}}{V^{\mathrm{reg}}}\cdot\left(\bar f-f_\infty\right),
\]
where $f(n',n',n')-f_\infty\ge0$.
The last term is at most $\frac{320}{63}\cdot\frac{h}{n'}\cdot\left(\bar f-f_\infty\right)$:
for $n=n'$ one has $V^{\mathrm{exc}}=\frac{5}{2}\,n'^2h-2n'h^2+\frac{1}{2}\,h^3\le\frac{5}{2}\,n'^2h$,
and $V^{\mathrm{reg}}=\frac{(n'-h)^2(2n'-h)}{2}\ge\frac{63}{128}\,n'^3$ because $n'=2k_1h\ge4h$ ($k_1\ge2$ by Proposition~\ref{prop:fold-any-side}).
\item
The inequality \textup{(d$'$)} is valid whenever $i\ge i_0(\mathfrak{d})$, $n\ge n'_i$ is an integer and $t>0$,
with no further relation between $i$, $n$ and $t$.
If $h$ and $t>0$ are kept fixed while $n'_i\to\infty$ and $n/n'_i\to1$,
its error term $|Q_0|\cdot\Delta'(n,n'_i,h)/t$ tends to $0$
(compare Theorem~\ref{thm:comparison-any-side}):
the first term $f(n'_i/2,2n'_i,n'_i)-f_\infty$ of $\Delta'(n,n'_i,h)$ tends to $0$ by Proposition~\ref{prop:coordinate-pressure}(d),
and in this regime the second term $\left(V^{\mathrm{exc}}/V^{\mathrm{reg}}\right)\cdot\left(\bar f-f_\infty\right)$ has the order of magnitude $(h+n-n'_i)/n$.
This convergence is not uniform in $t$.
\item
In~\eqref{eq:sides-three-steps} each exceptional torus contributes its partition function to the power $\vartheta_a$,
and Lemma~\ref{lem:uniform-ceiling} bounds these three factors one by one.
By~\eqref{eq:sides-exceptional} the three products $\vartheta_a\cdot|\mathbb{T}_a^{\mathrm{exc}}|$ have the sum $V^{\mathrm{exc}}$,
and $V^{\mathrm{exc}}=o(n^3)$ as $n'\to\infty$ and $n/n'\to1$ with $h$ fixed.
The volumes $|\mathbb{T}_a^{\mathrm{exc}}|$ themselves can exceed $n^3$ by a large factor;
for instance $|\mathbb{T}_1^{\mathrm{exc}}|=2k_1(h+s)\,n^2$ by Definition~\ref{def:tori}.
Only the products $\vartheta_a\cdot|\mathbb{T}_a^{\mathrm{exc}}|$ matter,
since the logarithm of the bound of Lemma~\ref{lem:uniform-ceiling} is $\bar f$ times the number $|\mathbb{T}|$ of photon sites of the torus,
with one and the same $\bar f$ for tori of every shape.
In closed form,
\[
\begin{aligned}
\vartheta_1\cdot|\mathbb{T}_1^{\mathrm{exc}}|&=(h+s)\,n^2,\\
\vartheta_2\cdot|\mathbb{T}_2^{\mathrm{exc}}|&=\frac{(n'-h)(h+2s)\,n}{2},\\
\vartheta_3\cdot|\mathbb{T}_3^{\mathrm{exc}}|&=\frac{(n'-h)(2n'-h)(h+s)}{2}.
\end{aligned}
\]
\end{enumerate}
\end{remark}

\subsection{Companion sequences}\label{sec:companion-sequences}

\begin{definition}[Companion sequences]\label{def:companions}
Let $\mathfrak{n}=(n_i)$ be a sequence of integers $n_i\ge4$ with $n_i\to\infty$.
A sequence $\mathfrak{n}'=(n'_i)$, indexed by the same set as $\mathfrak{n}$ and with all its terms positive integers,
is called a \emph{companion} of $\mathfrak{n}$ if it has the following three properties.
\begin{enumerate}
\item\label{it:companion-divisibility}
The sequence $\mathfrak{n}'$ satisfies condition \textup{(Div)} of Definition~\ref{def:sides}.
In particular $n'_i\to\infty$,
since each of the infinitely many elements of $\mathcal{D}(\mathfrak{n}')$ divides the positive integer $n'_i$ for all large $i$.
\item\label{it:companion-below}
One has $n'_i\le n_i$ for all but finitely many $i$.
\item\label{it:companion-ratio}
One has $n'_i/n_i\to1$.
\end{enumerate}
\end{definition}

\begin{definition}[The reduced sequence]\label{def:reduced-sequence}
For an integer $n\ge64$ put
\[
k(n):=\max\left\{k\in\Z:\ k\ge1,\ 8\,k!\le\sqrt{n}\right\},
\qquad
\red(n):=8\,k(n)!\cdot\left\lfloor\frac{n}{8\,k(n)!}\right\rfloor .
\]
Thus $\red(n)$ is the largest multiple of the modulus $8\,k(n)!$ that does not exceed $n$.
The set whose maximum defines $k(n)$ is non-empty, because $8\cdot1!=8=\sqrt{64}\le\sqrt{n}$,
and finite, because $k!\to\infty$.
The \emph{reduced sequence} of a sequence $\mathfrak{n}=(n_i)$ of integers $n_i\ge4$ with $n_i\to\infty$
is the sequence $\mathfrak{n}'=(n'_i)$ given by
\[
n'_i:=\red(n_i)\quad\text{if } n_i\ge64,
\qquad
n'_i:=8\quad\text{for the finitely many $i$ with } n_i<64 .
\]
\end{definition}

\begin{lemma}[Existence of companions]\label{lem:companions-exist}
\begin{enumerate}
\item\label{it:sides-modulus}
For every integer $n\ge64$ the modulus $8\,k(n)!$ divides $\red(n)$, and
\[
n-\sqrt{n}<\red(n)\le n,
\qquad
\red(n)>56 .
\]
The function $n\mapsto k(n)$ is non-decreasing,
and for all integers $k\ge1$ and $n\ge64$ one has $k(n)\ge k$ if and only if $n\ge64\,(k!)^2$.
\item\label{it:sides-reduced-companion}
Let $\mathfrak{n}=(n_i)$ be a sequence of integers $n_i\ge4$ with $n_i\to\infty$, and let $\mathfrak{n}'$ be its reduced sequence.
Then every positive integer belongs to $\mathcal{D}(\mathfrak{n}')$,
and $\mathfrak{n}'$ is a companion of $\mathfrak{n}$.
In particular every such sequence has a companion.
\item\label{it:sides-own-companion}
If a sequence $\mathfrak{n}=(n_i)$ of integers $n_i\ge4$ with $n_i\to\infty$ satisfies \textup{(Div)},
then $\mathfrak{n}$ is a companion of itself.
\end{enumerate}
\end{lemma}

\begin{proof}
\eqref{it:sides-modulus}
By definition $\red(n)$ is a multiple of $8\,k(n)!$,
and $n-\red(n)$ is the remainder of $n$ modulo $8\,k(n)!$.
Since $8\,k(n)!\le\sqrt{n}$,
\[
0\le n-\red(n)\le8\,k(n)!-1<\sqrt{n},
\]
which is the inequality $n-\sqrt{n}<\red(n)\le n$.
The map $n\mapsto n-\sqrt{n}$ is increasing on the integers $n\ge1$ and takes the value $64-8=56$ at $n=64$;
therefore $\red(n)>n-\sqrt{n}\ge56$.
For an integer $k\ge1$ the inequality $8\,k!\le\sqrt{n}$ is equivalent to $n\ge64\,(k!)^2$.
As $k\mapsto k!$ is non-decreasing, the integers $k\ge1$ with $8\,k!\le\sqrt{n}$ are exactly $1,\dots,k(n)$.
So $k(n)\ge k$ holds if and only if $8\,k!\le\sqrt{n}$, that is, if and only if $n\ge64\,(k!)^2$.
The last condition persists when $n$ increases;
hence $n\mapsto k(n)$ is non-decreasing.

\eqref{it:sides-reduced-companion}
Since $n_i\to\infty$, all but finitely many $i$ have $n_i\ge64$ and hence $n'_i=\red(n_i)$.
All terms $n'_i$ are positive integers:
$\red(n_i)>56$ by part~\eqref{it:sides-modulus}, and the remaining terms equal $8$.
Property~\eqref{it:companion-below} of Definition~\ref{def:companions} holds,
because $\red(n_i)\le n_i$ whenever $n_i\ge64$.
For property~\eqref{it:companion-ratio}, part~\eqref{it:sides-modulus} gives
\[
1-\frac{1}{\sqrt{n_i}}<\frac{n'_i}{n_i}\le1
\qquad\text{whenever } n_i\ge64 ,
\]
and the left-hand side tends to $1$;
consequently also $n'_i\to\infty$.
For property~\eqref{it:companion-divisibility}, let $k\ge1$ be an integer.
Then $n_i\ge64\,(k!)^2$ for all but finitely many $i$,
and for each such $i$ part~\eqref{it:sides-modulus} gives $k(n_i)\ge k$.
So $k$ divides $k(n_i)!$, which divides $8\,k(n_i)!$, which divides $n'_i$;
hence $k\in\mathcal{D}(\mathfrak{n}')$.
Every positive integer thus lies in $\mathcal{D}(\mathfrak{n}')$.
In particular $8\in\mathcal{D}(\mathfrak{n}')$ and $\mathcal{D}(\mathfrak{n}')$ is infinite,
which is condition \textup{(Div)} for $\mathfrak{n}'$.

\eqref{it:sides-own-companion}
Take $\mathfrak{n}'=\mathfrak{n}$.
Property~\eqref{it:companion-divisibility} of Definition~\ref{def:companions} is then the hypothesis,
and properties~\eqref{it:companion-below} and~\eqref{it:companion-ratio} hold trivially, since $n'_i=n_i$ for every $i$.
\end{proof}

\begin{remark}\label{rem:modulus}
\begin{enumerate}
\item
The argument for Lemma~\ref{lem:companions-exist}\eqref{it:sides-reduced-companion} produces a companion
also when the modulus $8\,k(n)!$ is replaced by any modulus that depends on $n$, tends to infinity and is $o(n)$,
and that is divisible by $8$ and by every integer of some fixed infinite set for all large $n$,
the threshold in $n$ depending on the integer.
A modulus independent of $n$ does not give a companion of every sequence,
even if it is chosen after the observable $O$, for instance as a multiple of $4h$ for an even width $h$:
if the $n_i$ are that modulus times pairwise distinct primes,
then the largest multiples of the modulus not exceeding the $n_i$ are the $n_i$ themselves,
so that the sequence they form is $\mathfrak{n}$,
and $\mathcal{D}(\mathfrak{n})$ is the set of the divisors of the modulus, a finite set;
hence $\mathfrak{n}$ does not satisfy \textup{(Div)}, which requires $\mathcal{D}(\mathfrak{n})$ to be infinite.
The inequalities $h''\ge2r+6h$ and $kh\ge16h''$ of Lemma~\ref{lem:coarse-structures}\eqref{it:coarse-widths}
force the widths $h''$ and $kh$ to tend to infinity with the radius $r$,
whereas for this sequence the conditions $4h''\in\mathcal{D}(\mathfrak{n})$ and $2kh\in\mathcal{D}(\mathfrak{n})$ of that lemma bound them by the modulus;
so the four conditions displayed there cannot hold together when $r$ is large.
\item
For the reduced sequence every positive integer lies in $\mathcal{D}(\mathfrak{n}')$,
by Lemma~\ref{lem:companions-exist}\eqref{it:sides-reduced-companion}.
The condition $4h\in\mathcal{D}(\mathfrak{n}')$ of Definition~\ref{def:data}
and the conditions $4h''\in\mathcal{D}(\mathfrak{n}')$ and $2kh\in\mathcal{D}(\mathfrak{n}')$ of Lemma~\ref{lem:coarse-structures}\eqref{it:coarse-widths}
therefore hold for all positive integers $h$, $h''$ and $k$;
in particular every even width $h$ is admitted.
\item
For the reduced sequence the threshold $i_0(\mathfrak{d})$ of Definition~\ref{def:threshold} can be large.
For $h=18$ the number $4h=72=8\cdot9$ divides $8\,k!$ if and only if $9$ divides $k!$,
that is, if and only if $k\ge6$.
By Lemma~\ref{lem:companions-exist}\eqref{it:sides-modulus},
$4h$ divides the modulus $8\,k(n)!$, and with it $\red(n)$, for every $n\ge64\,(6!)^2=33\,177\,600$,
whereas for $64\le n<64\,(6!)^2$ the modulus $8\,k(n)!$ is not divisible by $4h$.
\end{enumerate}
\end{remark}

If $\mathfrak{n}'$ is a companion of a sequence $\mathfrak{n}$ of integers $n_i\ge4$ with $n_i\to\infty$,
then $\mathfrak{n}'$ satisfies \textup{(Div)} by property~\eqref{it:companion-divisibility} of Definition~\ref{def:companions},
and Theorem~\ref{thm:comparison-states}(a), (b) hold for $\mathfrak{n}'$ and for data formed with it.

\begin{theorem}[Comparison states for arbitrary sides]\label{thm:comparison-any-side}
Let $-1\le\rho\le1$, $w_0,\dots,w_3>0$, $M\in\R$ and $q>0$.
Let $\mathfrak{n}=(n_i)$ be a sequence of integers $n_i\ge4$ with $n_i\to\infty$,
let $\mathfrak{n}'=(n'_i)$ be a companion of $\mathfrak{n}$,
let $\mu_i$ be the law of the cubic torus $C(n_i)$,
and let $\mathfrak{d}=(c_0,c_+,c_3;h;O)$ be data for the sequence $\mathfrak{n}'$ (Definition~\ref{def:data}),
with $\mu^t_i$, $i_0(\mathfrak{d})$ and the comparison states for $\mathfrak{d}$ formed with $\mathfrak{n}'$ as in Definition~\ref{def:two-sequences}.
\begin{enumerate}
\item\label{it:sides-error-term}
The number $\Delta'_i:=\Delta'(n_i,n'_i,h)$ of Proposition~\ref{prop:fold-any-side} is defined for all large $i$,
and $\Delta'_i\to0$.
\item\label{it:sides-cubic-law}
Let $\nu_0$ be a comparison state for $\mathfrak{d}$ subordinate to a subsequence $(i'_m)$ of the indices,
and suppose that the laws $\mu_{i'_m}$ converge weakly on $\Omega$ to a law $\nu$ as $m\to\infty$.
Then one has
\[
\nu(O)\le\nu_0(O).
\]
\item\label{it:sides-general-laws}
Suppose that to each index $i$ a law $\lambda_i$ on $\Omega$ with $\lambda_i\ll\mu_i$ is assigned,
that each $\lambda_i$ is invariant under $\tau_v$ for every $v\in\Z^3$,
and that $H(\lambda_i\,|\,\mu_i)=o(n_i^3)$ as $i\to\infty$, where $H$ is the relative entropy of Definition~\ref{def:relative-entropy}.
Then assertion~\eqref{it:sides-cubic-law} holds with $\lambda_{i'_m}$ in the place of $\mu_{i'_m}$.
\end{enumerate}
\end{theorem}

\begin{proof}
\eqref{it:sides-error-term}
By property~\eqref{it:companion-below} of Definition~\ref{def:companions} there is an index $i_1\ge i_0(\mathfrak{d})$ with $n'_i\le n_i$ for all $i\ge i_1$;
here $i_0(\mathfrak{d})$ is the threshold of Definition~\ref{def:threshold}, formed with $\mathfrak{n}'$.
For $i\ge i_1$ the hypotheses of Proposition~\ref{prop:fold-any-side} hold with the companion side $n'=n'_i$ and the side $n=n_i$.
So $\Delta'_i$ is defined for $i\ge i_1$, and
\[
\begin{aligned}
\Delta'_i&=f\left(\frac{n'_i}{2},2n'_i,n'_i\right)-f_\infty+\frac{V^{\mathrm{exc}}_i}{V^{\mathrm{reg}}_i}\cdot\left(\bar f-f_\infty\right),\\
V^{\mathrm{reg}}_i&=\frac{(n'_i-h)^2\cdot(2n'_i-h)}{2},
\qquad
V^{\mathrm{exc}}_i=n_i^3-V^{\mathrm{reg}}_i .
\end{aligned}
\]
Now let $i\to\infty$.
Then $n'_i\to\infty$ by property~\eqref{it:companion-divisibility} of Definition~\ref{def:companions},
so the smallest side $n'_i/2$ of $C(n'_i/2,2n'_i,n'_i)$ tends to infinity,
and $f(n'_i/2,2n'_i,n'_i)\to f_\infty$ by Proposition~\ref{prop:coordinate-pressure}(d).
Moreover
\begin{equation}\label{eq:sides-volume-ratio}
\frac{V^{\mathrm{reg}}_i}{n_i^3}
=\left(\frac{n'_i}{n_i}\right)^{3}\cdot\left(1-\frac{h}{n'_i}\right)^{2}\cdot\left(1-\frac{h}{2n'_i}\right)
\longrightarrow1
\end{equation}
by property~\eqref{it:companion-ratio} of Definition~\ref{def:companions} and $n'_i\to\infty$, the width $h$ being fixed.
Hence $V^{\mathrm{exc}}_i/V^{\mathrm{reg}}_i=n_i^3/V^{\mathrm{reg}}_i-1\to0$.
As $\bar f-f_\infty$ is a fixed number, $\Delta'_i\to0$.

\eqref{it:sides-general-laws}
Let the laws $\lambda_i$ be as in~\eqref{it:sides-general-laws},
let $\nu_0$ be a comparison state for $\mathfrak{d}$ subordinate to a subsequence $(i'_m)$ of the indices,
and let the laws $\lambda_{i'_m}$ converge weakly on $\Omega$ to a law $\nu$ as $m\to\infty$.
Choose $(i_m)$, $t_k>0$ and $\nu_{t_k}$ for $\nu_0$ as in Definition~\ref{def:comparison-states},
with $(i_m)$ a subsequence of $(i'_m)$.
The observable $O$ is a real trigonometric polynomial, hence bounded and continuous on $\Omega$.
Fix $k$.
Every $\lambda_i$ is invariant under all translations by vectors of $\Z^3$,
so $\lambda_i(\tau_vO)=\lambda_i(O)$ for every $v\in\Lambda_h$.
Proposition~\ref{prop:fold-any-side}\textup{(e$'$)}, applied at the index $i_m$ with $n=n_{i_m}$, the strength $t_k$ and the law $\lambda=\lambda_{i_m}$,
gives for every $m$ with $i_m\ge i_1$
\[
\lambda_{i_m}(O)\le\mu^{t_k}_{i_m}(O)
+\frac{|Q_0|}{t_k}\cdot\left(\Delta'_{i_m}+\frac{H(\lambda_{i_m}\,|\,\mu_{i_m})}{V^{\mathrm{reg}}_{i_m}}\right);
\]
for $\lambda_i=\mu_i$ this inequality, in which then $H(\mu_{i_m}\,|\,\mu_{i_m})=0$,
is Proposition~\ref{prop:fold-any-side}\textup{(d$'$)}.
Let $m\to\infty$ at the fixed strength $t_k>0$.
Then $\lambda_{i_m}(O)\to\nu(O)$, since $(i_m)$ is a subsequence of $(i'_m)$ and $O$ is bounded and continuous;
$\mu^{t_k}_{i_m}(O)\to\nu_{t_k}(O)$ by Definition~\ref{def:comparison-states};
$\Delta'_{i_m}\to0$ by~\eqref{it:sides-error-term};
and, by the hypothesis on the relative entropy and by~\eqref{eq:sides-volume-ratio},
\[
\frac{H(\lambda_i\,|\,\mu_i)}{V^{\mathrm{reg}}_i}
=\frac{H(\lambda_i\,|\,\mu_i)}{n_i^3}\cdot\frac{n_i^3}{V^{\mathrm{reg}}_i}\longrightarrow0 .
\]
Hence $\nu(O)\le\nu_{t_k}(O)$ for every $k$.
Finally let $k\to\infty$:
since $\nu_{t_k}(O)\to\nu_0(O)$ by Definition~\ref{def:comparison-states}, one obtains $\nu(O)\le\nu_0(O)$.

\eqref{it:sides-cubic-law}
The laws $\mu_i$ fulfil the conditions imposed on $\lambda_i$ in~\eqref{it:sides-general-laws}.
Indeed $\mu_i\ll\mu_i$ with the density $1$, so that $H(\mu_i\,|\,\mu_i)=0$;
and $\mu_i$ is invariant under $\tau_v$ for every $v\in\Z^3$,
since $\tau_v$ maps the torus $C(n_i)$ onto itself (Lemma~\ref{lem:wall-group}(f))
and Lemma~\ref{lem:bosonic-measure}(c) applies.
Assertion~\eqref{it:sides-cubic-law} is therefore the case $\lambda_i=\mu_i$ of~\eqref{it:sides-general-laws}.
\end{proof}

\begin{remark}\label{rem:order-volume-strength}
In the proof of Theorem~\ref{thm:comparison-any-side} the two limits are taken as Definition~\ref{def:comparison-states} prescribes:
first the volume grows ($m\to\infty$) while the strength $t_k$ is held fixed,
and then the strength tends to zero ($k\to\infty$).
The order is dictated by the error term $|Q_0|\cdot\Delta'_i/t$,
which by Theorem~\ref{thm:comparison-any-side}\eqref{it:sides-error-term} tends to $0$ as $i\to\infty$ for each fixed $t>0$, but carries the factor $1/t$;
in the proof of Theorem~\ref{thm:comparison-states}(c) the error term $|Q_0|\cdot\Delta_i/t$ plays the same part.
\end{remark}

\subsection{Limits along every sequence}\label{sec:every-sequence}

\begin{lemma}[Domination at a diagonal mirror]\label{lem:diagonal-domination}
Let $0<\rho<1$, $w_0,\dots,w_3>0$, $M\in\R$ and $q>0$.
Let $\mathfrak{n}=(n_i)$ be a sequence of integers $n_i\ge 4$ with $n_i\to\infty$,
and let $\mu_i$ be the law of the cubic torus $C(n_i)$.
For every $i$ let $\lambda_i$ be a law on $\Omega$
that is invariant under $\tau_v$ for every $v\in\Z^3$ and satisfies $\lambda_i\ll\mu_i$,
and assume that the relative entropies (Definition~\ref{def:relative-entropy}) satisfy
\[
H(\lambda_i\,|\,\mu_i)=o\left(n_i^3\right)\qquad(i\to\infty);
\]
for instance, $\lambda_i=\mu_i$.
Let $\nu$ be a weak limit point on $\Omega$ of the sequence $(\lambda_i)$, and let $c_0\in\Z$.
Then for every $O\in\mathcal{T}$ the set $\mathcal{R}(c_0)$ of Lemma~\ref{lem:diagonal-laws} contains a law $\nu_0$ with
\[
\nu(O)\le\nu_0(O).
\]
\end{lemma}

\begin{proof}
By Lemma~\ref{lem:companions-exist}\eqref{it:sides-reduced-companion} the sequence $\mathfrak{n}$ has a companion;
fix one, $\mathfrak{n}'=(n'_i)$ (Definition~\ref{def:companions}).
The data, the perturbed functionals $\mu^t_i$ and the comparison states that occur below are formed with $\mathfrak{n}'$,
in accordance with Definition~\ref{def:two-sequences},
and $\mathfrak{n}'$ satisfies \textup{(Div)} by Definition~\ref{def:companions}\eqref{it:companion-divisibility}.
Let $O\in\mathcal{T}$.
\begin{enumerate}
\item\label{it:sides-width}
\emph{Choice of the width.}
Put $r:=r_O\ge 1$,
and choose $y_0\in\R^3$ with $\supp O\subset y_0+[-r,r]^3$, as the definition of $r_O$ allows (Definition~\ref{def:characters}).
Since $\mathfrak{n}'$ satisfies \textup{(Div)},
Lemma~\ref{lem:coarse-structures}\eqref{it:coarse-widths}, applied to $\mathfrak{n}'$ with $h_0:=8r+9$,
provides an even integer $h$ with $4h\in\mathcal{D}(\mathfrak{n}')$ and
\[
h\ge 8r+9>8r+8.
\]
If $\mathfrak{n}'$ is the reduced sequence of $\mathfrak{n}$ (Definition~\ref{def:reduced-sequence}),
then $\mathcal{D}(\mathfrak{n}')$ is the set of all positive integers by Lemma~\ref{lem:companions-exist}\eqref{it:sides-reduced-companion},
and one may take $h:=8r+10$, an even integer with $h\ge 8r+9$.
\item\label{it:sides-position}
\emph{An even translation into an open cell.}
Choose an even integer $a$ in the open interval
\[
\left(c_0+2r-\ell_-(y_0),\ c_0+h-2r-\ell_-(y_0)\right).
\]
Such an $a$ exists:
the end points differ by $h-4r$, and $h-4r>2$ as $h>8r+8$ and $r\ge 1$,
while every open interval longer than $2$ contains an even integer.
The vector $v:=a\cdot e_1$ lies in $2\Z^3$, $a$ being even; let $y:=y_0+v$.
From $\ell_-(e_1)=1$ one obtains $\ell_-(y)=\ell_-(y_0)+a$, and the condition on $a$ becomes
\[
c_0<\ell_-(y)-2r,\qquad \ell_-(y)+2r<c_0+h.
\]
Next pick an integer $c_+$ and an element $c_3$ of $\Z+\tfrac12$ such that
\[
c_+<\ell_+(y)-2r,\qquad \ell_+(y)+2r<c_++h,\qquad c_3<y_3-r,\qquad y_3+r<c_3+h.
\]
The first two conditions ask that $c_+$ lie in $\left(\ell_+(y)+2r-h,\ \ell_+(y)-2r\right)$, an open interval of length $h-4r>1$,
and the last two ask that $c_3$ lie in $\left(y_3+r-h,\ y_3-r\right)$, an open interval of length $h-2r>1$;
any open interval longer than $1$ contains an integer as well as an element of $\Z+\tfrac12$.
For $x\in y+[-r,r]^3$ one has
\[
\left|\ell_\pm(x)-\ell_\pm(y)\right|\le 2r,\qquad |x_3-y_3|\le r,
\]
because $|n_{\ell_\pm}|_1=2$.
It follows that the cube $y+[-r,r]^3$ lies in the open cell
\[
Q_0=\left\{c_0<\ell_-<c_0+h,\ c_+<\ell_+<c_++h,\ c_3<x_3<c_3+h\right\}
\]
of the cell structure $\mathcal{Q}:=(\mathfrak{F}_{\mathrm{d}};c_0,c_+,c_3;h)$ (Definition~\ref{def:cell-structures}).
The translate $\tau_vO$ is again a real trigonometric polynomial, so $\tau_vO\in\mathcal{T}$;
it reads the plaquettes of $\supp O+v$ (Definition~\ref{def:characters}), and these lie in $y+[-r,r]^3\subset Q_0$.
Finally $r_{\tau_vO}=r$,
since the translation by $v$ maps the cubes that contain $\supp O$ onto the cubes that contain $\supp O+v$.
\item\label{it:sides-comparison}
\emph{A comparison state in $\mathcal{R}(c_0)$.}
Let $\mathfrak{d}':=(c_0,c_+,c_3;h;\tau_vO)$.
This tuple is data for the sequence $\mathfrak{n}'$ in the sense of Definition~\ref{def:data},
and its cell structure $\mathcal{Q}_{\mathfrak{d}'}$ is $\mathcal{Q}$.
Indeed, $c_0$ is an integer by hypothesis, and Step~\ref{it:sides-position} gives $c_+\in\Z$ and $c_3\in\Z+\tfrac12$;
by Step~\ref{it:sides-width}, $h$ is even and positive ($h\ge 8r+9$), with $4h\in\mathcal{D}(\mathfrak{n}')$;
by Step~\ref{it:sides-position}, $\tau_vO\in\mathcal{T}$ and $\supp\tau_vO\subset Q_0$, the open cell of $\mathcal{Q}$ with $\kappa=0$;
and $8r_{\tau_vO}+8=8r+8<h$.
Since $\nu$ is a weak limit point of $(\lambda_i)$,
there is a subsequence $(i'_m)$ of the indices with $\lambda_{i'_m}\to\nu$ weakly;
the indices of $\mathfrak{n}'$ are those of $\mathfrak{n}$ (Definition~\ref{def:companions}),
so $(i'_m)$ is a subsequence of the indices of $\mathfrak{n}'$ as well.

Theorem~\ref{thm:comparison-states}, read with $\mathfrak{n}'$ in the place of its sequence of sides (Definition~\ref{def:two-sequences}),
applies to $\mathfrak{n}'$ and $\mathfrak{d}'$.
Its hypotheses hold:
$-1\le\rho\le 1$, because $0<\rho<1$;
$w_0,\dots,w_3>0$, $M\in\R$ and $q>0$ by assumption;
the sequence $\mathfrak{n}'$ satisfies \textup{(Div)} by Definition~\ref{def:companions}\eqref{it:companion-divisibility};
and $\mathfrak{d}'$ is data for $\mathfrak{n}'$.
By Theorem~\ref{thm:comparison-states}\eqref{it:comparison-existence},
some comparison state $\nu_0$ for $\mathfrak{d}'$ is subordinate to $(i'_m)$.
By Theorem~\ref{thm:comparison-states}\eqref{it:comparison-properties},
$\nu_0$ satisfies \textup{(G)}, and at every wall of $\mathcal{Q}_{\mathfrak{d}'}=\mathcal{Q}$ it is invariant and reflection positive.
The diagonal mirror $(\ell_-,c_0)$ is a wall of $\mathcal{Q}$; consequently $\nu_0\in\mathcal{R}(c_0)$.

Theorem~\ref{thm:comparison-any-side} applies to $\mathfrak{n}$, its companion $\mathfrak{n}'$, the data $\mathfrak{d}'$ and the laws $\lambda_i$:
its condition $-1\le\rho\le 1$ holds,
and the conditions that it imposes on the laws $\lambda_i$ are hypotheses of the lemma.
As $\lambda_{i'_m}\to\nu$ weakly and $\nu_0$ is subordinate to $(i'_m)$, that theorem gives
\[
\nu(\tau_vO)\le\nu_0(\tau_vO).
\]
\item\label{it:sides-limit-invariant}
\emph{Invariance of $\nu$ under $\tau_v$.}
Every $\lambda_i$ is invariant under all translations by elements of $\Z^3$,
in particular under $\tau_v$
(for $\lambda_i=\mu_i$ this is Proposition~\ref{prop:symmetries}\eqref{it:functional-covariance}).
The functions $O$ and $\tau_vO$ are bounded and continuous on $\Omega$,
being finite linear combinations of characters $\varphi_\beta$ with $\beta\in C^2_{\mathrm{c}}$.
Weak convergence along $(i'_m)$ therefore yields
\[
\nu(\tau_vO)=\lim_{m\to\infty}\lambda_{i'_m}(\tau_vO)=\lim_{m\to\infty}\lambda_{i'_m}(O)=\nu(O).
\]
\item\label{it:sides-state-invariant}
\emph{Invariance of $\nu_0$ under $\tau_v$.}
By Step~\ref{it:sides-comparison}, the law $\nu_0$ is a comparison state for $\mathfrak{d}'$,
and $\mathfrak{d}'$ is standard data (Definition~\ref{def:data}) for $\mathfrak{n}'$, a sequence with \textup{(Div)}.
Hence Theorem~\ref{thm:no-layers}, in its case of standard data and read for $\mathfrak{n}'$,
applies to $\mathfrak{d}'$ and $\nu_0$;
its conditions on $\rho$, $w_0,\dots,w_3$, $M$ and $q$ are those of the lemma.
It shows that $\nu_0$ is invariant under the translations by elements of $2\Z^3$.
Since $a$ is even, $v=a\cdot e_1$ is such an element, whence
\[
\nu_0(\tau_vO)=\nu_0(O).
\]
\end{enumerate}
The comparison state $\nu_0$ of Step~\ref{it:sides-comparison} lies in $\mathcal{R}(c_0)$,
and Steps~\ref{it:sides-comparison}, \ref{it:sides-limit-invariant} and~\ref{it:sides-state-invariant} combine to
\[
\nu(O)=\nu(\tau_vO)\le\nu_0(\tau_vO)=\nu_0(O).
\]
\end{proof}

\begin{proof}[Proof of Theorem~\ref{thm:cubic-limits}]
Let $0<\rho<1$, $w_0,\dots,w_3>0$, $M\in\R$ and $q>0$,
let $\mathfrak{n}=(n_i)$ be a sequence of integers $n_i\ge 4$ with $n_i\to\infty$,
and let $\mu_i:=\mu_{C(n_i)}$ denote the distribution of $\widetilde F$
under the functional $\mathbb{E}_{C(n_i)}$ (Definition~\ref{def:functional}).
For every $i$ the period lattice of $C(n_i)=\mathbb{T}(\mathfrak{F}_{\mathrm{c}};n_i,n_i,n_i)$ is $n_i\Z^3$,
and $\sys(n_i\Z^3)=n_i\ge 4$;
so $C(n_i)$ is a torus (Definitions~\ref{def:configurations}, \ref{def:tori} and~\ref{def:cell-structures}).
Its partition function $Z_{C(n_i)}[1]$ is positive and finite by Theorem~\ref{thm:partition-positivity},
and $\mu_i$ is a law on $\Omega$ by Lemma~\ref{lem:bosonic-measure}.
\begin{enumerate}
\item\label{it:sides-tightness}
\emph{Tightness, Gaussian domination and closedness.}
Let $\beta\in C^2_{\mathrm{c}}$ have its support in a box $B$ that embeds in $\Gamma:=n_i\Z^3$ (Definition~\ref{def:integrands}).
The translates of $\supp\beta$ by the elements of $\Gamma$ are then pairwise disjoint,
so that the periodisation $\beta_\Gamma\in C^2(\Gamma)$ of $\beta$ satisfies
$\langle\beta_\Gamma,F\rangle=\langle\beta,\widetilde F\rangle$ and $\|\beta_\Gamma\|=\|\beta\|$.
Gaussian domination on $C(n_i)$, in the form of Theorem~\ref{thm:coordinate-domination}(d), gives
\[
\mu_i\left(\exp\left(\langle\beta,F\rangle\right)\right)\le\exp\left(\frac{\|\beta\|^{2}}{2}\right),
\]
that is, $\mu_i$ satisfies~\eqref{cond:domination-beta} for this $\beta$.
A plaquette $\sigma$ lies in a box of radius $1/2$;
enlarged by $[-1,1]^3$, this box becomes a cube of side $3<4\le n_i$,
and therefore it embeds in $n_i\Z^3$.
Hence every $\mu_i$ satisfies~\eqref{cond:domination-beta} for $\beta=\pm s\cdot 1_\sigma$,
for every plaquette $\sigma$ and every $s>0$,
and the sequence $(\mu_i)$ is tight by Lemma~\ref{lem:domination-compactness}(b).
By Prokhorov's theorem it has weak limit points.
Let $\nu$ be one of them, and let $(i'_m)$ be a subsequence of the indices with $\mu_{i'_m}\to\nu$ weakly.
Every $\beta\in C^2_{\mathrm{c}}$ has its support in a box of some radius $r'$,
and this box embeds in $n_i\Z^3$ as soon as $n_i>2r'+2$, that is, for all large $i$.
Lemma~\ref{lem:domination-compactness}(c) then shows that $\nu$ satisfies~\eqref{cond:domination-beta}
for every $\beta\in C^2_{\mathrm{c}}$, which is~\eqref{cond:gaussian-domination}.
Finally, $\nu$ is a law on $\Omega=\{dF=0\}$, so that $\nu(dF=0)=1$.
\item\label{it:sides-invariance}
\emph{Invariance.}
By Proposition~\ref{prop:symmetries}\eqref{it:functional-covariance} and~\eqref{it:charge-conjugation},
each $\mu_i$ is invariant under every $g\in W$ and under $F\mapsto -F$, whatever the side $n_i\ge 4$.
For $g\in W$ the map $g^{*}$ is a homeomorphism of $\Omega$, and so is $F\mapsto -F$.
Let $g\in W$, and let $G$ be a bounded continuous function on $\Omega$.
The function $gG$, given by $(gG)(F)=G(g^{*}F)$ as in~\eqref{eq:layers-action}, is again bounded and continuous, and
\[
\nu(gG)=\lim_{m\to\infty}\mu_{i'_m}(gG)=\lim_{m\to\infty}\mu_{i'_m}(G)=\nu(G).
\]
Since a law on a metric space is determined by the integrals of the bounded continuous functions,
$\nu$ is invariant under $g$.
The same argument, with $F\mapsto -F$ in the place of $g^{*}$, shows that $\nu$ is invariant under $F\mapsto -F$.
In particular $\nu$ is invariant at every mirror, the reflection in a mirror being an element of $W$.
\item\label{it:sides-axis}
\emph{The three axis families.}
Let $\Pi=(x_\mu,c)$ with $c\in\frac{1}{2}\Z$, put $\theta:=\theta_\Pi$,
and let $G\in\mathcal{O}_{\mathrm{cont}}(\{x_\mu>c\})$ (see~\eqref{eq:layers-local}):
a bounded continuous function of finitely many $F_\sigma$, each plaquette $\sigma$ contained in $\{x_\mu>c\}$.
Choose $h_G>0$ such that every corner of each of these plaquettes satisfies $c<x_\mu<c+h_G$.
Let $i$ be an index with $n_i>2h_G$; all large indices are such.
The torus $C(n_i)$ is adapted to $x_\mu$, with the height period $n_i$
(Definitions~\ref{def:slabs} and~\ref{def:cell-structures}).
As $c\in\frac{1}{2}\Z$ and $n_i$ is an integer,
$c+n_i/2$ lies in $\frac{1}{2}\Z=\Adm(x_\mu)$, for even and for odd $n_i$.
The plaquettes read by $G$ lie in the open slab $\{c<x_\mu<c+n_i/2\}$, since $h_G<n_i/2$.
Hence $G$ is an element of $\mathcal{A}\left(\{c<x_\mu<c+n_i/2\}\right)$ without factors of Gaussian growth;
such an element embeds in every lattice (Definition~\ref{def:integrands}),
and so $G$ lies in $\mathcal{A}_\Xi[c,c+n_i/2]$ for the cross-section $\Xi:=\Xi(C(n_i),x_\mu)$.
Corollary~\ref{cor:trace-forms}\eqref{it:trace-form} with $G'=G$,
together with~\eqref{eq:layers-bosonic} and Lemma~\ref{lem:bosonic-measure}(b), yields
\[
\mu_i\left(\overline{\theta G}\cdot G\right)
=\frac{Z_{C(n_i)}\left[\Theta_\Pi(G)\cdot G\right]}{Z_{C(n_i)}[1]}
=\frac{\Tr_{\mathcal{H}(c)}\left(\mathfrak{S}[G]^{*}\cdot\mathfrak{S}[G]\right)}{Z_{C(n_i)}[1]}
\ge 0 .
\]
The function $\overline{\theta G}\cdot G$ is bounded and continuous on $\Omega$; therefore
\[
\nu\left(\overline{\theta G}\cdot G\right)=\lim_{m\to\infty}\mu_{i'_m}\left(\overline{\theta G}\cdot G\right)\ge 0 .
\]
By step~\ref{it:sides-invariance}, $\nu$ is invariant under $\theta$,
and Lemma~\ref{lem:law-mirror}(b) shows that $\nu$ is reflection positive at $\Pi$.
\item\label{it:sides-wall}
\emph{Membership in $\mathcal{R}(c_0)$.}
Let $c_0\in\Z$.
The set $\mathcal{R}(c_0)$ is convex and weakly compact (Lemma~\ref{lem:diagonal-laws}).
Lemma~\ref{lem:diagonal-domination} applies with $\lambda_i=\mu_i$:
one has $\mu_i\ll\mu_i$ and $H(\mu_i\,|\,\mu_i)=0$,
and $\mu_i$ is invariant under $\tau_v$ for every $v\in\Z^3$
by Proposition~\ref{prop:symmetries}\eqref{it:functional-covariance}.
So for every $O\in\mathcal{T}$ there is a law $\nu_0\in\mathcal{R}(c_0)$ with $\nu(O)\le\nu_0(O)$.
By Lemma~\ref{lem:separation}, $\nu$ lies in $\mathcal{R}(c_0)$;
in particular, $\nu$ is reflection positive at $(\ell_-,c_0)$.
\item\label{it:sides-diagonal}
\emph{The six diagonal families.}
Let $\ell=x_\mu-x_{\mu'}$ or $\ell=x_\mu+x_{\mu'}$ with $\mu<\mu'$, let $c\in\Z$,
and let $\alpha$ be the third index, so that $\{\mu,\mu',\alpha\}=\{1,2,3\}$.
Let $g\in\Ocub$ be the matrix with $ge_1=e_\mu$ and $ge_3=e_\alpha$,
and with $ge_2=e_{\mu'}$ in the first case, $ge_2=-e_{\mu'}$ in the second.
Then $(gx)_\mu=x_1$, and $(gx)_{\mu'}=x_2$, resp.\ $(gx)_{\mu'}=-x_2$;
so $\ell(gx)=x_1-x_2=\ell_-(x)$ in both cases,
and $\ell(gx)-c=\ell_-(x)-c_0$ for all $x\in\R^3$, where $c_0:=c$.
The law $\nu$ is invariant under $W$ (step~\ref{it:sides-invariance})
and reflection positive at $(\ell_-,c_0)$ (step~\ref{it:sides-wall}).
By Lemma~\ref{lem:mirror-transport}, applied to the mirrors $(\ell_-,c_0)$ and $(\ell,c)$ and to $g$,
the law $\nu$ is reflection positive at $(\ell,c)$.
\end{enumerate}
These five steps contain every assertion of Theorem~\ref{thm:cubic-limits}:
tightness, Gaussian domination and closedness in step~\ref{it:sides-tightness},
the invariances in step~\ref{it:sides-invariance},
and reflection positivity at the mirrors of the three axis families and of the six diagonal families
in steps~\ref{it:sides-axis} and~\ref{it:sides-diagonal}.
\end{proof}

\begin{corollary}[All cubic tori]\label{cor:all-cubic-tori}
Let $0<\rho<1$, $w_0,\dots,w_3>0$, $M\in\R$ and $q>0$,
and for every integer $n\ge 4$ let $\mu_{C(n)}$ be the law of $\widetilde F$ under $\mathbb{E}_{C(n)}$
for the chamber operator with these parameters (Definitions~\ref{def:chamber-operator} and~\ref{def:functional}).
\begin{enumerate}
\item The family $\left(\mu_{C(n)}\right)_{n\ge 4}$ is tight on $\Omega$.
\item Let $\nu$ be a law on $\Omega$ that is the weak limit of $\left(\mu_{C(n_i)}\right)_i$
for some sequence of integers $n_i\ge 4$ with $n_i\to\infty$.
Then $\nu$ satisfies~\eqref{cond:gaussian-domination} and $\nu(dF=0)=1$,
it is invariant under $W$ and under $F\mapsto -F$,
and it is reflection positive (Definition~\ref{def:reflection-positivity})
at every mirror of the nine families (Definition~\ref{def:mirrors}).
\end{enumerate}
\end{corollary}

\begin{proof}
The argument of step~\ref{it:sides-tightness} in the proof of Theorem~\ref{thm:cubic-limits}
holds for every integer $n\ge 4$ in the place of $n_i$:
each law $\mu_{C(n)}$ satisfies~\eqref{cond:domination-beta} for $\beta=\pm s\cdot 1_\sigma$,
for every plaquette $\sigma$ and every $s>0$.
By Lemma~\ref{lem:domination-compactness}(b) the family $\left(\mu_{C(n)}\right)_{n\ge 4}$ is tight,
which is the first assertion.
A law $\nu$ as in the second assertion is a weak limit point of the sequence $\left(\mu_{C(n_i)}\right)_i$,
and Theorem~\ref{thm:cubic-limits}, applied to this sequence, gives the second assertion.
\end{proof}

Every law $\nu$ as in the second assertion of Corollary~\ref{cor:all-cubic-tori}
is invariant under $W$ and reflection positive at every mirror of the nine families,
and these are the hypotheses of Theorem~\ref{thm:lattice-cone}.
Assertions (a)--(d) of that theorem therefore hold for such a law $\nu$ and for all $\mu\neq\mu'$ in $\{1,2,3\}$:
the operator inequality of~(a),
the inequality $\cosh\lambda+\cos k\ge 2$ of~(b),
and the lattice forms~(c) and~(d) of the two symmetry conditions of Definition~\ref{def:symmetry-conditions}.

\begin{remark}\label{rem:two-sequences}
\begin{enumerate}
\item
A companion $\mathfrak{n}'$ of $\mathfrak{n}$ has the same index set as $\mathfrak{n}$ (Definition~\ref{def:companions}),
and the two sequences are coupled index by index.
For all large $i$, Proposition~\ref{prop:fold-any-side} with $n=n_i$ compares the law $\mu_i$ of the cubic torus $C(n_i)$
with the functional $\mu^t_i$ on $D(n'_i/2,n'_i,n'_i)$ of the same index $i$.
The limits are coupled in the same way:
a subsequence $(i'_m)$ of the indices along which the laws $\mu_i$ converge weakly is also a subsequence of the indices of $\mathfrak{n}'$,
and a comparison state for data formed with $\mathfrak{n}'$ is subordinate to $(i'_m)$
if it arises, as in Definition~\ref{def:comparison-states}, from the laws $\mu^{t_k}_{i_m}$ along a subsequence $(i_m)$ of $(i'_m)$.
\item
For each $O\in\mathcal{T}$, the law $\nu_0$ of Lemma~\ref{lem:diagonal-domination} is constructed in its proof
as a comparison state for data formed with a companion $\mathfrak{n}'$ of $\mathfrak{n}$,
and it depends on $\mathfrak{n}'$, on $O$ and on the subsequences of the indices.
All these laws lie in the set $\mathcal{R}(c_0)$ of the laws on $\Omega$ that satisfy~\eqref{cond:gaussian-domination}
and are invariant and reflection positive at the mirror $(\ell_-,c_0)$.
This set is convex and weakly compact (Lemma~\ref{lem:diagonal-laws}),
and Lemma~\ref{lem:separation} turns the inequalities $\nu(O)\le\nu_0(O)$, one for each $O\in\mathcal{T}$,
into the single inclusion $\nu\in\mathcal{R}(c_0)$ of step~\ref{it:sides-wall} in the proof of Theorem~\ref{thm:cubic-limits}.
\item
Let $\mathfrak{n}$ satisfy \textup{(Div)}.
Then $\mathfrak{n}$ is a companion of itself (Lemma~\ref{lem:companions-exist}(c)).
For $\mathfrak{n}'=\mathfrak{n}$ the data of Theorem~\ref{thm:comparison-any-side} are data for $\mathfrak{n}$,
and the threshold $i_0(\mathfrak{d})$, the functionals $\mu^t_i$ and the comparison states are formed with $\mathfrak{n}$ itself.
In finite volume, Proposition~\ref{prop:fold-any-side}\textup{(d$'$)} with $n=n'=n_i$ is then the inequality of Proposition~\ref{prop:fold-finite}\eqref{it:finite-fold},
with the error term $|Q_0|\cdot\Delta'(n_i,n_i,h)/t$ in the place of $|Q_0|\cdot\Delta_i/t$.
Both error terms tend to $0$ as $i\to\infty$ at fixed $t>0$
(Theorem~\ref{thm:four-tori} and Theorem~\ref{thm:comparison-any-side}\eqref{it:sides-error-term}),
so that the difference disappears in the limits of Definition~\ref{def:comparison-states}:
for $\lambda_i=\mu_i$ the inequality $\nu(O)\le\nu_0(O)$ of Theorem~\ref{thm:comparison-any-side}
is the inequality of Theorem~\ref{thm:comparison-states}\eqref{it:comparison-fold}.
\end{enumerate}
\end{remark}

\subsection{Real exponents}\label{sec:real-exponents}

Besides Proposition~\ref{prop:fold-any-side},
H\"older's inequality with real exponents gives a second fold inequality between tori of different sides:
Proposition~\ref{prop:comparable-even} for even and Proposition~\ref{prop:comparable-odd} for odd sides.

\emph{Norms with real exponents.}
Let $\mathcal{H}$ and $\mathcal{H}'$ be Hilbert spaces, and let $Y\colon\mathcal{H}\to\mathcal{H}'$ be a Hilbert--Schmidt operator.
For real $s\ge2$ put
\[
\|Y\|_s:=\left(\Tr_{\mathcal{H}}\left(Y^*\cdot Y\right)^{\frac{s}{2}}\right)^{\frac{1}{s}}.
\]
The trace is finite.
Indeed, $Y^*\cdot Y$ is a non-negative operator of trace class,
and as $s/2\ge1$, every eigenvalue $\xi$ of $Y^*\cdot Y$ satisfies $\xi^{\frac{s}{2}}\le\|Y^*\cdot Y\|^{\frac{s}{2}-1}\cdot\xi$.
For $s=2k$ with an integer $k\ge1$, the number $\|Y\|_{2k}$ is the norm of Definition~\ref{def:even-schatten}.

\emph{Pressures.}
For a torus $\mathbb{T}$ the number $Z_{\mathbb{T}}[1]$ is positive (Theorem~\ref{thm:partition-positivity}),
and the \emph{pressure} of $\mathbb{T}$ is
\[
f(\mathbb{T}):=\frac{1}{|\mathbb{T}|}\,\log Z_{\mathbb{T}}[1].
\]
On a coordinate torus it is the pressure of Definition~\ref{def:sides}: $f(C(n_1,n_2,n_3))=f(n_1,n_2,n_3)$.

\emph{Slab operators.}
Let $\ell\in\mathcal{L}$, and let $\Xi=(\Gamma_\parallel,\chi_\parallel)$ be a cross-section for $\ell$ (Definition~\ref{def:slabs}).
Recall that $\Adm(\ell)=\frac{|n_\ell|^2}{2}\,\Z$ and that the integer vector $n_\ell$ has the height $\ell(n_\ell)=|n_\ell|^2$ (Definition~\ref{def:mirrors}).
Let $\mathcal{H}(c)$, $c\in\Adm(\ell)$, be the mirror spaces of Definition~\ref{def:mirror-spaces} for $\ell$ and $\Xi$,
with the parity $U_{\mathrm{f}}$ and the operators $V_v\colon\mathcal{H}(c)\to\mathcal{H}(c+\ell(v))$, $v\in\Z^3$, defined there;
these spaces are separable, as shown after that definition.
By Theorem~\ref{thm:slab-operators}\eqref{it:slab-hilbert-schmidt}, $U_{\mathrm{f}}$ is a unitary involution of each $\mathcal{H}(c)$,
every $V_v$ is unitary, and $V_v\cdot V_{v'}=V_{v+v'}$ for all $v,v'\in\Z^3$.
For $c<c'$ in $\Adm(\ell)$, Corollary~\ref{cor:slab-admissible} gives a linear map $G\mapsto\mathfrak{S}[G]$ from $\mathcal{A}_\Xi[c,c']$ to the Hilbert--Schmidt operators from $\mathcal{H}(c)$ to $\mathcal{H}(c')$,
whose values are the slab operators.
The same corollary shows that the clauses \eqref{it:slab-composition} to \eqref{it:slab-trace} of Theorem~\ref{thm:slab-operators}, on composition, covariance, reflection and trace,
hold for elements of the classes $\mathcal{A}_\Xi[c,c']$,
with $G^{\tau_v}$ in the place of the translate of $G$ and $\Theta_{\ell,c'}G$ in the place of its reflection in the upper mirror.
For $c<c'$ in $\Adm(\ell)$ the constant function $1$ belongs to $\mathcal{A}_\Xi[c,c']$,
and its slab operator, regarded in this class, is denoted
\[
\mathfrak{S}[1]_{[c,c']}\colon\mathcal{H}(c)\to\mathcal{H}(c').
\]

\begin{lemma}[H\"older's inequality with real exponents]\label{lem:real-exponents}
Let $r\ge2$, let $\mathcal{H}_0,\dots,\mathcal{H}_r$ be separable Hilbert spaces,
let $Y_j\colon\mathcal{H}_{j-1}\to\mathcal{H}_j$ be Hilbert--Schmidt operators for $j=1,\dots,r$,
let $U\colon\mathcal{H}_r\to\mathcal{H}_0$ be unitary,
and let $s_1,\dots,s_r\in[2,\infty)$ be real numbers with $\sum_{j=1}^{r}\frac{1}{s_j}=1$.
Then
\[
\left|\Tr_{\mathcal{H}_0}\left(U\cdot Y_r\cdots Y_1\right)\right|\le\prod_{j=1}^{r}\|Y_j\|_{s_j}.
\]
Moreover, $\|Y\|_s\le\|Y\|_{s'}$ for every Hilbert--Schmidt operator $Y$ between two separable Hilbert spaces
and all real numbers $s\ge s'\ge2$.
\end{lemma}

\begin{proof}
The operator $U\cdot Y_r\cdots Y_1$ on $\mathcal{H}_0$ is the product of the two Hilbert--Schmidt operators $U\cdot Y_r$ and $Y_{r-1}\cdots Y_1$,
hence of trace class.
Let $\mathcal{H}:=\mathcal{H}_0\oplus\dots\oplus\mathcal{H}_{r-1}$, a separable Hilbert space,
let $\iota_j\colon\mathcal{H}_j\to\mathcal{H}$, $0\le j\le r-1$, be the canonical isometries, so that $\iota_j^*\cdot\iota_j=1$,
and let
\[
\hat Y_j:=\iota_j\cdot Y_j\cdot\iota_{j-1}^*\quad(1\le j\le r-1),\qquad
\hat Y_r:=\iota_0\cdot U\cdot Y_r\cdot\iota_{r-1}^*
\]
be the lifts~\eqref{eq:tori-direct-sum} of the operators $Y_j$ to $\mathcal{H}$.
From $\iota_j^*\cdot\iota_j=1$ and $U^*\cdot U=1$ one obtains
\[
\begin{aligned}
\hat Y_j^*\cdot\hat Y_j&=\iota_{j-1}\cdot Y_j^*\cdot Y_j\cdot\iota_{j-1}^*\qquad(1\le j\le r),\\
\hat Y_r\cdots\hat Y_1&=\iota_0\cdot\left(U\cdot Y_r\cdots Y_1\right)\cdot\iota_0^*.
\end{aligned}
\]
Each $\iota_j$ embeds $\mathcal{H}_j$ as a summand of $\mathcal{H}$.
If $\mathfrak{B}$ is a trace class operator on $\mathcal{H}_j$, $0\le j\le r-1$,
then $\iota_j\cdot\mathfrak{B}\cdot\iota_j^*$ acts as $\mathfrak{B}$ on this summand and annihilates the others;
its absolute value is $\iota_j\cdot|\mathfrak{B}|\cdot\iota_j^*$.
Computing traces in an orthonormal basis of $\mathcal{H}$ assembled from orthonormal bases of the summands,
one finds that $\iota_j\cdot\mathfrak{B}\cdot\iota_j^*$ is of trace class and that
\[
\Tr_{\mathcal{H}}\left(\iota_j\cdot\mathfrak{B}\cdot\iota_j^*\right)=\Tr_{\mathcal{H}_j}\mathfrak{B}.
\]

\emph{Real powers.}
Let $0\le j\le r-1$, let $\mathfrak{B}\ge0$ be a bounded operator on $\mathcal{H}_j$, and let $x>0$ be real.
Then
\[
\left(\iota_j\cdot\mathfrak{B}\cdot\iota_j^*\right)^x=\iota_j\cdot\mathfrak{B}^x\cdot\iota_j^*.
\]
For integer exponents this follows from $\iota_j^*\cdot\iota_j=1$.
Let $n\ge1$ be an integer.
The operator $\iota_j\cdot\mathfrak{B}^{\frac{1}{n}}\cdot\iota_j^*$ is non-negative, and its $n$-th power is $\iota_j\cdot\mathfrak{B}\cdot\iota_j^*$;
since a non-negative operator has exactly one non-negative $n$-th root,
$\left(\iota_j\cdot\mathfrak{B}\cdot\iota_j^*\right)^{\frac{1}{n}}=\iota_j\cdot\mathfrak{B}^{\frac{1}{n}}\cdot\iota_j^*$.
Integer powers of this identity give the assertion for every rational $x>0$.
The spectra of $\mathfrak{B}$ and of $\iota_j\cdot\mathfrak{B}\cdot\iota_j^*$ lie in $[0,\|\mathfrak{B}\|]$,
and on this interval the functions $\xi\mapsto\xi^x$ depend continuously on $x>0$ in the supremum norm;
by the continuous functional calculus both sides are therefore continuous in $x>0$ for the operator norm,
and the identity holds for every real $x>0$.

\emph{Conclusion.}
The separable space $\mathcal{H}$ carries the Schatten classes $\mathscr{J}_s(\mathcal{H})$ with real exponents $1\le s<\infty$ introduced in Section~\ref{sec:tori}:
$\mathscr{J}_s(\mathcal{H})$ consists of the compact operators $Y$ on $\mathcal{H}$ whose norm $\left(\Tr_{\mathcal{H}}|Y|^s\right)^{\frac{1}{s}}$, with $|Y|:=\left(Y^*\cdot Y\right)^{\frac12}$, is finite.
Now let $1\le j\le r$ and let $s\ge2$ be real.
With $\mathfrak{B}:=Y_j^*\cdot Y_j$, an operator on $\mathcal{H}_{j-1}$, and $x:=s/2$, the identity of real powers gives
$\left(\hat Y_j^*\cdot\hat Y_j\right)^{\frac{s}{2}}=\iota_{j-1}\cdot\left(Y_j^*\cdot Y_j\right)^{\frac{s}{2}}\cdot\iota_{j-1}^*$,
and hence
\[
\Tr_{\mathcal{H}}\left(\hat Y_j^*\cdot\hat Y_j\right)^{\frac{s}{2}}=\Tr_{\mathcal{H}_{j-1}}\left(Y_j^*\cdot Y_j\right)^{\frac{s}{2}},
\qquad\text{that is,}\qquad
\|\hat Y_j\|_s=\|Y_j\|_s .
\]
The case $s=2$ shows that $\hat Y_j$ is a Hilbert--Schmidt operator on $\mathcal{H}$.
Since $|\hat Y_j|^s=\left(\hat Y_j^*\cdot\hat Y_j\right)^{\frac{s}{2}}$, the norm of $\hat Y_j$ in $\mathscr{J}_s(\mathcal{H})$ is the number $\|\hat Y_j\|_s$ defined above;
moreover $\hat Y_j\in\mathscr{J}_2(\mathcal{H})\subset\mathscr{J}_{s_j}(\mathcal{H})$ by Lemma~\ref{lem:schatten-traces}\eqref{it:schatten-monotone}.
Lemma~\ref{lem:schatten-traces}\eqref{it:schatten-holder}, applied in $\mathcal{H}$ to the $r$ operators $\hat Y_j$ with the exponents $s_j$, yields
\[
\left|\Tr_{\mathcal{H}_0}\left(U\cdot Y_r\cdots Y_1\right)\right|
=\left|\Tr_{\mathcal{H}}\left(\hat Y_r\cdots\hat Y_1\right)\right|
\le\prod_{j=1}^{r}\|\hat Y_j\|_{s_j}
=\prod_{j=1}^{r}\|Y_j\|_{s_j}.
\]
For the second assertion let $Y\colon\mathcal{H}_0\to\mathcal{H}_1$ be a Hilbert--Schmidt operator between separable Hilbert spaces,
and form the lifts above with $r=2$, $Y_1:=Y$, $\mathcal{H}_2:=\mathcal{H}_0$, $Y_2:=Y^*$ and $U:=1$.
For real $s\ge s'\ge2$, Lemma~\ref{lem:schatten-traces}\eqref{it:schatten-monotone} in $\mathcal{H}=\mathcal{H}_0\oplus\mathcal{H}_1$ gives
$\|Y\|_s=\|\hat Y_1\|_s\le\|\hat Y_1\|_{s'}=\|Y\|_{s'}$.
\end{proof}

\begin{lemma}[Empty slabs]\label{lem:empty-powers}
Let $-1\le\rho\le1$, $w_0,\dots,w_3>0$, $M\in\R$ and $q>0$.
Let $\ell\in\mathcal{L}$, let $\Xi=(\Gamma_\parallel,\chi_\parallel)$ be a cross-section for $\ell$, and let $c\in\Adm(\ell)$.
With the slab operators of Corollary~\ref{cor:slab-admissible} for the pair $(\ell,\Xi)$, let
\[
\mathfrak{E}:=\mathfrak{S}[1]_{[c,\,c+|n_\ell|^2/2]}\colon\mathcal{H}(c)\to\mathcal{H}\left(c+\frac{|n_\ell|^2}{2}\right)
\]
be the \emph{elementary step} at $c$, and let $b\ge1$ be an integer.
\begin{enumerate}
\item[\textup{(a)}]
$\mathfrak{S}[1]_{[c,\,c+b\cdot|n_\ell|^2]}=V_{b\cdot n_\ell}\cdot\left(\mathfrak{E}^*\cdot\mathfrak{E}\right)^b$.
\item[\textup{(b)}]
For every real $s\ge2$,
$\left\|\mathfrak{S}[1]_{[c,\,c+b\cdot|n_\ell|^2]}\right\|_s^s=\Tr_{\mathcal{H}(c)}\left(\mathfrak{E}^*\cdot\mathfrak{E}\right)^{b\cdot s}$.
\item[\textup{(c)}]
Let $k\ge1$ be an integer with $\sys\left(\Gamma_\parallel+\Z\,k\cdot n_\ell\right)\ge4$,
and let $\mathbb{T}:=\mathbb{T}(\Xi;k\cdot|n_\ell|^2)$ be the torus of Lemma~\ref{lem:tori-over-section}.
Then
\[
Z_{\mathbb{T}}[1]=\Tr_{\mathcal{H}(c)}\left(\mathfrak{E}^*\cdot\mathfrak{E}\right)^k.
\]
\end{enumerate}
\end{lemma}

\begin{proof}
For $c'<c''$ in $\Adm(\ell)$ the class $\mathcal{A}_\Xi[c',c'']$ contains the constant $1$,
and $1^g=1$ for every $g\in W$ (Lemma~\ref{lem:admissible-images}\eqref{it:images-products}).
Thus $\Theta_{\ell,c'}1=1$ and $1^{\tau_v}=1$, as elements of the reflected and of the translated class.
All positions that occur below have the form $c+i\cdot\frac{|n_\ell|^2}{2}$ with $i\in\Z$, and so lie in $\Adm(\ell)$.

(a)
For a slab between positions $c'<c''$, the reflection rule of Theorem~\ref{thm:slab-operators}\eqref{it:slab-reflection} involves the vector $\frac{2(c''-c')}{|n_\ell|^2}\cdot n_\ell$;
for $c'=c$ and $c''=c+|n_\ell|^2/2$ this vector is $n_\ell$.
Applied to this slab and to the element $1$, the rule gives
\[
\mathfrak{S}[1]_{[c+|n_\ell|^2/2,\,c+|n_\ell|^2]}=V_{n_\ell}\cdot\mathfrak{E}^*,
\]
and the composition rule of Theorem~\ref{thm:slab-operators}\eqref{it:slab-composition} gives
\[
\mathfrak{S}[1]_{[c,\,c+|n_\ell|^2]}
=\mathfrak{S}[1]_{[c+|n_\ell|^2/2,\,c+|n_\ell|^2]}\cdot\mathfrak{S}[1]_{[c,\,c+|n_\ell|^2/2]}
=V_{n_\ell}\cdot\mathfrak{E}^*\cdot\mathfrak{E}.
\]
Let $i\ge0$ be an integer.
The vector $i\cdot n_\ell$ has the height $\ell(i\cdot n_\ell)=i\cdot|n_\ell|^2$,
and the covariance of Theorem~\ref{thm:slab-operators}\eqref{it:slab-covariance}, together with $V_v\cdot V_{v'}=V_{v+v'}$, yields
\[
\mathfrak{S}[1]_{[c+i\cdot|n_\ell|^2,\,c+(i+1)\cdot|n_\ell|^2]}
=V_{i\cdot n_\ell}\cdot V_{n_\ell}\cdot\mathfrak{E}^*\cdot\mathfrak{E}\cdot V_{i\cdot n_\ell}^{-1}
=V_{(i+1)\cdot n_\ell}\cdot\mathfrak{E}^*\cdot\mathfrak{E}\cdot V_{i\cdot n_\ell}^{-1}.
\]
After $b-1$ applications of the composition rule, $\mathfrak{S}[1]_{[c,\,c+b\cdot|n_\ell|^2]}$ is the product of these operators for $i=b-1,\dots,0$,
the factor with $i=b-1$ standing on the left:
\[
\begin{aligned}
\mathfrak{S}[1]_{[c,\,c+b\cdot|n_\ell|^2]}
&=\left(V_{b\cdot n_\ell}\cdot\mathfrak{E}^*\cdot\mathfrak{E}\cdot V_{(b-1)\cdot n_\ell}^{-1}\right)\cdot\left(V_{(b-1)\cdot n_\ell}\cdot\mathfrak{E}^*\cdot\mathfrak{E}\cdot V_{(b-2)\cdot n_\ell}^{-1}\right)\\
&\qquad\cdots\left(V_{n_\ell}\cdot\mathfrak{E}^*\cdot\mathfrak{E}\cdot V_0^{-1}\right)\\
&=V_{b\cdot n_\ell}\cdot\left(\mathfrak{E}^*\cdot\mathfrak{E}\right)^b.
\end{aligned}
\]
The product telescopes, and $V_0=1$, since $V_0\cdot V_0=V_0$ and $V_0$ is unitary.

(b)
Put $\mathfrak{S}^\circ:=\mathfrak{S}[1]_{[c,\,c+b\cdot|n_\ell|^2]}$.
As $V_{b\cdot n_\ell}$ is unitary and $\left(\mathfrak{E}^*\cdot\mathfrak{E}\right)^b$ is self-adjoint, (a) gives
\[
\left(\mathfrak{S}^\circ\right)^*\cdot\mathfrak{S}^\circ
=\left(\mathfrak{E}^*\cdot\mathfrak{E}\right)^b\cdot V_{b\cdot n_\ell}^*\cdot V_{b\cdot n_\ell}\cdot\left(\mathfrak{E}^*\cdot\mathfrak{E}\right)^b
=\left(\mathfrak{E}^*\cdot\mathfrak{E}\right)^{2b}.
\]
By the functional calculus of the non-negative operator $\mathfrak{E}^*\cdot\mathfrak{E}$,
$\left(\left(\mathfrak{S}^\circ\right)^*\cdot\mathfrak{S}^\circ\right)^{\frac{s}{2}}=\left(\mathfrak{E}^*\cdot\mathfrak{E}\right)^{b\cdot s}$,
and the trace of the left side is $\|\mathfrak{S}^\circ\|_s^s$.

(c)
The height $P:=k\cdot|n_\ell|^2$ is a positive element of $|n_\ell|^2\Z$, and $\frac{P}{|n_\ell|^2}\cdot n_\ell=k\cdot n_\ell$.
By Lemma~\ref{lem:tori-over-section}(b) and the hypothesis on the systole,
$\mathbb{T}$ is a torus adapted to $\ell$, that is, its own straightened torus;
moreover $\Xi(\mathbb{T},\ell)=\Xi$, $p(\mathbb{T},\ell)=P$, and $k\cdot n_\ell$ is a closing vector of $\mathbb{T}$ for $\ell$.
Put $h':=|n_\ell|^2/2$; then $P=2k\cdot h'$, and the positions $c+j\cdot h'$, $j\in\Z$, lie in $\Adm(\ell)$.
The identity~\eqref{eq:tori-fold-norm} of Proposition~\ref{prop:twisted-holder},
for the torus $\mathbb{T}$, the closing vector $k\cdot n_\ell$, the integer $k$, the height $h'$, the base point $c$
and the element $1$ of $\mathcal{A}_\Xi[c,c+h']$, reads
\[
Z_{\mathbb{T}}\left[\fold_\ell 1\right]=\|\mathfrak{E}\|_{2k}^{2k}=\Tr_{\mathcal{H}(c)}\left(\mathfrak{E}^*\cdot\mathfrak{E}\right)^k.
\]
The copies of $1$ in the slabs of Definition~\ref{def:slab-copies}, formed with $c_0:=c$ and $h'$, are images $1^g$ of $1$ under elements $g\in W$,
hence equal to $1$ (Lemma~\ref{lem:admissible-images}\eqref{it:images-products}).
So the fold of $1$ in the direction $\ell$, the product of these copies over the $2k$ slabs between the positions $c$ and $c+P$, is $\fold_\ell 1=1$,
and the left side of the last identity is $Z_{\mathbb{T}}[1]$.
\end{proof}

\begin{lemma}[A low empty slab]\label{lem:low-slab}
Let $-1\le\rho\le1$, $w_0,\dots,w_3>0$, $M\in\R$ and $q>0$.
Let $\mathcal{Q}=(\mathfrak{F};o_1,o_2,o_3;h)$ be a cell structure (Definition~\ref{def:cell-structures}),
let $a\in\{1,2,3\}$ and $\ell:=\ell^a$,
and let $\mathbb{T}=(\Gamma,\chi)$ be a torus, $\Xi:=\Xi(\mathbb{T},\ell)=(\Gamma_\parallel,\chi_\parallel)$ and $p:=p(\mathbb{T},\ell)$.
Let $k\ge1$ be an integer with $2k\cdot h\le p$, and put $h^\circ:=p-2k\cdot h$.
Assume that $h^\circ\le p/2$ and $h^\circ\in|n_\ell|^2\Z$;
that the straightened torus $\mathbb{T}^0:=\mathbb{T}^0(\ell)$ is defined,
so that $p\in|n_\ell|^2\Z$ and $\mathbb{T}^0=\mathbb{T}(\Xi;p)$ by Lemma~\ref{lem:tori-over-section}(c);
and that $\sys\left(\Gamma_\parallel+\Z\,\frac{2k\cdot h}{|n_a|^2}\cdot n_a\right)\ge4$.
Let $\mathbb{T}^{\mathrm{reg}}:=\mathbb{T}(\Xi;2k\cdot h)$.
Let $\mathfrak{G}=(G_\kappa)_{\kappa\in\mathcal{K}}$ be a finite filling of $\mathcal{Q}$ (Definition~\ref{def:fillings})
with $0\le\kappa_a\le2k-1$ for all $\kappa\in\mathcal{K}$.
Then $Z_{\mathbb{T}^{\mathrm{reg}}}\left[\prod\mathfrak{R}_a^{j,2k}\mathfrak{G}\right]\ge0$ for every $j\in\Z$,
$Z_{\mathbb{T}^0}[1]\ge0$,
and, with the convention $0^0:=1$,
\begin{equation}\label{eq:sides-low-slab}
\begin{aligned}
\left|Z_{\mathbb{T}}\left[\prod\mathfrak{G}\right]\right|
&\le Z_{\mathbb{T}^0}[1]^{\frac{h^\circ}{p}}\cdot\prod_{j=0}^{2k-1}Z_{\mathbb{T}^{\mathrm{reg}}}\left[\prod\mathfrak{R}_a^{j,2k}\mathfrak{G}\right]^{\frac{1}{2k}},\\
|\mathbb{T}|&=|\mathbb{T}^{\mathrm{reg}}|+\frac{h^\circ}{p}\cdot|\mathbb{T}^0|.
\end{aligned}
\end{equation}
\end{lemma}

\begin{proof}
\emph{Volumes.}
By Lemma~\ref{lem:tori-over-section}(c), $|\mathbb{T}|=p\cdot|\Xi|$.
Since $\ell=\ell^a$, one has $n_a=n_\ell$;
the heights $p$ and $2k\cdot h$ lie in $|n_\ell|^2\Z$, the second one because it is even.
So Lemma~\ref{lem:tori-over-section}(b) applies to $\mathbb{T}^0=\mathbb{T}(\Xi;p)$, whose lattice has systole at least $4$ by part~(c) of that lemma,
and to $\mathbb{T}^{\mathrm{reg}}$, by the hypothesis on the systole:
both are tori adapted to $\ell$ with the cross-section $\Xi$,
and $|\mathbb{T}^0|=p\cdot|\Xi|$, $|\mathbb{T}^{\mathrm{reg}}|=2k\cdot h\cdot|\Xi|$.
Therefore
\[
|\mathbb{T}^{\mathrm{reg}}|+\frac{h^\circ}{p}\cdot|\mathbb{T}^0|=\left(2k\cdot h+h^\circ\right)\cdot|\Xi|=p\cdot|\Xi|=|\mathbb{T}|,
\]
which is the second line of~\eqref{eq:sides-low-slab}.

\emph{The contents of the layers.}
Put $c_j:=o_a+j\cdot h$ for $0\le j\le2k$;
these positions lie in $\Adm(\ell)$ by Lemma~\ref{lem:layer-fold}\eqref{it:sides-layer-slab}.
For $1\le j\le2k$ let
\[
\Phi_j:=\prod_{\kappa\in\mathcal{K},\ \kappa_a=j-1}G_\kappa
\]
be the content of the layer $j-1$, with $\Phi_j:=1$ if no $\kappa\in\mathcal{K}$ has $\kappa_a=j-1$,
and let $Y_j:=\mathfrak{S}[\Phi_j]\colon\mathcal{H}(c_{j-1})\to\mathcal{H}(c_j)$.
By Lemma~\ref{lem:layer-fold}\eqref{it:sides-layer-slab}, $\Phi_j$ lies in $\mathcal{A}_\Xi[c_{j-1},c_j]$.
Lemma~\ref{lem:layer-fold}\eqref{it:sides-layer-fold}, applied to the layer $j-1$ under the hypothesis on the systole, gives by~\eqref{eq:sides-layer-norm}
\[
Z_{\mathbb{T}^{\mathrm{reg}}}\left[\prod\mathfrak{R}_a^{j-1,2k}\mathfrak{G}\right]=\|Y_j\|_{2k}^{2k}\ge0\qquad(1\le j\le2k);
\]
the same identity for an arbitrary layer $j\in\Z$ is the first assertion of the lemma.
Since $0\le\kappa_a\le2k-1$ for every $\kappa\in\mathcal{K}$,
each factor $G_\kappa$ of $\prod\mathfrak{G}$ occurs in exactly one of the products $\Phi_j$, and
\[
\prod\mathfrak{G}=\Phi_1\cdots\Phi_{2k}.
\]

\emph{The straightened torus.}
As $\mathbb{T}^0$ is defined, Lemma~\ref{lem:tori-over-section}(c) gives $p\in|n_\ell|^2\Z$
and $\sys\left(\Gamma_\parallel+\Z\,\frac{p}{|n_\ell|^2}\cdot n_\ell\right)\ge4$.
Hence $k_\circ:=p/|n_\ell|^2$ is a positive integer,
the condition on the systole in Lemma~\ref{lem:empty-powers}(c) holds for $k_\circ$,
and $\mathbb{T}(\Xi;k_\circ\cdot|n_\ell|^2)=\mathbb{T}(\Xi;p)=\mathbb{T}^0$.
Let $\mathfrak{E}$ be the elementary step at $c_{2k}$.
Lemma~\ref{lem:empty-powers}(c), with $k_\circ$ in the place of $k$, gives
\[
Z_{\mathbb{T}^0}[1]=\Tr_{\mathcal{H}(c_{2k})}\left(\mathfrak{E}^*\cdot\mathfrak{E}\right)^{k_\circ}\ge0,
\]
the trace of a non-negative operator of trace class;
this holds whether $h^\circ$ vanishes or not.

\emph{The closing vector.}
Let $g$ be a closing vector of $\mathbb{T}$ for $\ell$.
By Lemma~\ref{lem:closing-vectors}, $\ell(g)=p$, and the height period $p$ is a positive integer.
As $c_0=o_a$ lies in $\Adm(\ell)$ and $\Adm(\ell)+\Z=\Adm(\ell)$, the position $c_0+p$ lies in $\Adm(\ell)$, and
\[
U:=V_g^{-1}\cdot U_{\mathrm{f}}^{\frac{1+\chi(g)}{2}}\colon\mathcal{H}(c_0+p)\to\mathcal{H}(c_0)
\]
is a unitary operator, as a product of two unitaries.

\emph{The case $h^\circ=0$.}
Here $p=2k\cdot h$, $c_{2k}=c_0+p$ and $\mathbb{T}^{\mathrm{reg}}=\mathbb{T}(\Xi;p)=\mathbb{T}^0$.
The trace formula of Theorem~\ref{thm:slab-operators}\eqref{it:slab-trace}, applied to $\mathbb{T}$, the closing vector $g$,
the $2k\ge2$ slabs between the positions $c_0<c_1<\dots<c_{2k}=c_0+p$ and the elements $\Phi_1,\dots,\Phi_{2k}$,
gives $Z_{\mathbb{T}}\left[\prod\mathfrak{G}\right]=\Tr_{\mathcal{H}(c_0)}\left(U\cdot Y_{2k}\cdots Y_1\right)$.
Lemma~\ref{lem:holder-spaces}, applied to the $2k$ Hilbert--Schmidt operators $Y_j$ between the separable spaces $\mathcal{H}(c_j)$, now gives
\[
\left|Z_{\mathbb{T}}\left[\prod\mathfrak{G}\right]\right|
\le\prod_{j=1}^{2k}\|Y_j\|_{2k}
=\prod_{j=1}^{2k}\left(\|Y_j\|_{2k}^{2k}\right)^{\frac{1}{2k}}
=\prod_{j=0}^{2k-1}Z_{\mathbb{T}^{\mathrm{reg}}}\left[\prod\mathfrak{R}_a^{j,2k}\mathfrak{G}\right]^{\frac{1}{2k}}.
\]
As $Z_{\mathbb{T}^0}[1]^0=1$, this is the first line of~\eqref{eq:sides-low-slab}.

\emph{The case $h^\circ>0$.}
Put $c_{2k+1}:=o_a+p=c_0+p$, a position in $\Adm(\ell)$;
then $c_{2k}<c_{2k+1}$, as $c_{2k+1}-c_{2k}=h^\circ$.
Let
\[
\mathfrak{S}^\circ:=\mathfrak{S}[1]_{[c_{2k},\,c_{2k+1}]}\colon\mathcal{H}(c_{2k})\to\mathcal{H}(c_{2k+1})
\]
be the slab operator of the empty slab.
The trace formula of Theorem~\ref{thm:slab-operators}\eqref{it:slab-trace}, applied now to $\mathbb{T}$, the closing vector $g$,
the $2k+1\ge3$ slabs between the positions $c_0<\dots<c_{2k}<c_{2k+1}=c_0+p$ and the elements $\Phi_1,\dots,\Phi_{2k},1$, gives
\[
Z_{\mathbb{T}}\left[\prod\mathfrak{G}\right]=\Tr_{\mathcal{H}(c_0)}\left(U\cdot\mathfrak{S}^\circ\cdot Y_{2k}\cdots Y_1\right).
\]
Lemma~\ref{lem:real-exponents} applies with $r:=2k+1$, the operators $Y_1,\dots,Y_{2k}$ and $Y_{2k+1}:=\mathfrak{S}^\circ$, and the exponents
\[
s_j:=\frac{p}{h}\quad(1\le j\le2k),\qquad s_{2k+1}:=\frac{p}{h^\circ}:
\]
indeed $\sum_{j=1}^{2k+1}\frac{1}{s_j}=\frac{2k\cdot h+h^\circ}{p}=1$,
$p/h\ge2k\ge2$ because $2k\cdot h\le p$,
and $p/h^\circ\ge2$ because $h^\circ\le p/2$.
Together with the second assertion of Lemma~\ref{lem:real-exponents}, used for the exponents $p/h\ge2k\ge2$, it gives
\[
\left|Z_{\mathbb{T}}\left[\prod\mathfrak{G}\right]\right|
\le\|\mathfrak{S}^\circ\|_{p/h^\circ}\cdot\prod_{j=1}^{2k}\|Y_j\|_{p/h}
\le\|\mathfrak{S}^\circ\|_{p/h^\circ}\cdot\prod_{j=1}^{2k}\|Y_j\|_{2k}.
\]
The height of the empty slab is $h^\circ=b\cdot|n_\ell|^2$ with the integer $b:=h^\circ/|n_\ell|^2\ge1$,
so that $c_{2k+1}=c_{2k}+b\cdot|n_\ell|^2$.
Apply Lemma~\ref{lem:empty-powers}(b) at the position $c_{2k}$, with this $b$ and with $s:=p/h^\circ\ge2$.
Since $b\cdot s=p/|n_\ell|^2=k_\circ$, it gives, together with the identity for $Z_{\mathbb{T}^0}[1]$ obtained above,
\[
\|\mathfrak{S}^\circ\|_{p/h^\circ}^{\frac{p}{h^\circ}}=\Tr_{\mathcal{H}(c_{2k})}\left(\mathfrak{E}^*\cdot\mathfrak{E}\right)^{k_\circ}=Z_{\mathbb{T}^0}[1],
\qquad\text{that is,}\qquad
\|\mathfrak{S}^\circ\|_{p/h^\circ}=Z_{\mathbb{T}^0}[1]^{\frac{h^\circ}{p}}.
\]
Inserting this, and $\|Y_j\|_{2k}=Z_{\mathbb{T}^{\mathrm{reg}}}\left[\prod\mathfrak{R}_a^{j-1,2k}\mathfrak{G}\right]^{\frac{1}{2k}}$ for $1\le j\le2k$,
into the preceding inequality yields the first line of~\eqref{eq:sides-low-slab}.
\end{proof}

\subsection{Odd height periods}\label{sec:odd-periods}

\begin{remark}\label{rem:odd-sides}
For an odd side $n$, the comparison of the cubic torus $C(n)$ with the adapted torus $D(n')$ through low empty slabs
differs in two points from the case of an even side, Proposition~\ref{prop:comparable-even}.
First, the height period of $C(n)$ for the family $\ell_-$ is $n$.
Among the hypotheses of Lemma~\ref{lem:low-slab} is that the straightened torus $\mathbb{T}^0(\ell_-)$ of $C(n)$ be defined,
and this requires $n\in|n_{\ell_-}|^2\Z=2\Z$ (Lemma~\ref{lem:tori-over-section}(c)).
For odd $n$ the low empty slab has an odd height $h^\circ$.
Let $\mathfrak{E}$ be the elementary step of Lemma~\ref{lem:empty-powers} at the lower plane of the empty slab,
over the cross-section of $C(n)$ for $\ell_-$.
For the height period $p=n$, the identity~\eqref{eq:sides-empty-norm} shows that the slab operator of the empty slab
has the norm
\[
\left(\Tr\left(\mathfrak{E}^*\cdot\mathfrak{E}\right)^{\frac{n}{2}}\right)^{\frac{h^\circ}{n}}
\]
for the exponent $n/h^\circ$.
This is a power of a trace with the half-integer exponent $n/2$,
whereas Lemma~\ref{lem:empty-powers}(c) identifies $\Tr\left(\mathfrak{E}^*\cdot\mathfrak{E}\right)^{k}$ with a partition function for integers $k$ only.
Secondly, among the tori that enter the comparison there are tori whose period in the direction $e_3$ is the odd number $n$,
for instance $D(n'/2,n',n)$, and their pressures have to be compared with $f(n'/2,n'/2,n'/2)$;
in Theorem~\ref{thm:four-tori} the sides are even, and in Proposition~\ref{prop:periodic-perturbations} so are the entries of the shapes.
Two facts deal with both points:
the identification in Lemma~\ref{lem:empty-powers}(c),
and the monotonicity in Lemma~\ref{lem:schatten-traces}\eqref{it:schatten-monotone},
by which a Schatten norm does not increase when its exponent increases.
For $C(n)$ and $\ell_-$, Lemma~\ref{lem:odd-period} replaces the exponent $n/2$ by the integer $(n-1)/2$ with the help of the second fact,
and the first fact then identifies $\Tr\left(\mathfrak{E}^*\cdot\mathfrak{E}\right)^{\frac{n-1}{2}}$ with the partition function of $D'((n-1)/2,n,n)$;
in this way $\Tr\left(\mathfrak{E}^*\cdot\mathfrak{E}\right)^{\frac{n}{2}}$ is bounded by a power of that partition function.
For the pressures, the two facts lead to Lemma~\ref{lem:section-monotone}(c):
whatever the parities of the sides, the pressure of $D(a_1,a_2,a_3)$ does not increase when one of the sides grows,
all sides being integers $\ge4$.
With these two lemmas one arrives at Proposition~\ref{prop:comparable-odd}.
\end{remark}

\begin{lemma}[Monotonicity over one cross-section]\label{lem:section-monotone}
Let $-1\le\rho\le1$, $w_0,\dots,w_3>0$, $M\in\R$ and $q>0$.
Let $\ell\in\mathcal{L}$, let $\Xi=(\Gamma_\parallel,\chi_\parallel)$ be a cross-section for $\ell$
with the area $|\Xi|=[\Z^3\cap\{\ell=0\}:\Gamma_\parallel]$, let $c\in\Adm(\ell)$, and let
\[
\mathfrak{E}:=\mathfrak{S}[1]_{[c,c+|n_\ell|^2/2]}\colon\mathcal{H}(c)\to\mathcal{H}(c+|n_\ell|^2/2)
\]
be the elementary step, as in Lemma~\ref{lem:empty-powers}.
Since $\mathfrak{E}$ is a Hilbert--Schmidt operator (Corollary~\ref{cor:slab-admissible}),
$\mathfrak{E}^*\cdot\mathfrak{E}$ is a non-negative operator of trace class on $\mathcal{H}(c)$.
For real $s\ge1$, let $\mathscr{J}_s$ denote the Schatten class of $\mathcal{H}(c)$ with the exponent $s$ (Lemma~\ref{lem:schatten-traces}):
the compact operators on $\mathcal{H}(c)$ whose modulus has an $s$-th power of trace class,
normed by the $s$-th root of the trace of that power.
Put
\[
\|\mathfrak{E}^*\cdot\mathfrak{E}\|_s:=\left(\Tr\left(\mathfrak{E}^*\cdot\mathfrak{E}\right)^{s}\right)^{\frac{1}{s}};
\]
as the non-negative operator $\mathfrak{E}^*\cdot\mathfrak{E}$ coincides with its own modulus, this number is its norm in $\mathscr{J}_s$.
\begin{enumerate}
\item[(a)] For every integer $k\ge1$ with $\sys\left(\Gamma_\parallel+\Z\,k\cdot n_\ell\right)\ge4$,
\[
Z_{\mathbb{T}(\Xi;k\cdot|n_\ell|^2)}[1]=\Tr\left(\mathfrak{E}^*\cdot\mathfrak{E}\right)^{k}\ge0
\qquad\text{and}\qquad
|\mathbb{T}(\Xi;k\cdot|n_\ell|^2)|=k\cdot|n_\ell|^2\cdot|\Xi|.
\]
\item[(b)] For real $s\ge s'\ge1$ one has $\|\mathfrak{E}^*\cdot\mathfrak{E}\|_s\le\|\mathfrak{E}^*\cdot\mathfrak{E}\|_{s'}$.
Suppose that $Z_{\mathbb{T}(\Xi;k\cdot|n_\ell|^2)}[1]>0$ at some integer $k$ to which (a) applies.
Then this partition function is positive at every integer $k$ to which (a) applies,
the pressure $f(\mathbb{T})=|\mathbb{T}|^{-1}\log Z_{\mathbb{T}}[1]$ of $\mathbb{T}=\mathbb{T}(\Xi;k\cdot|n_\ell|^2)$ at each such $k$ is
\[
f\left(\mathbb{T}(\Xi;k\cdot|n_\ell|^2)\right)=\frac{\log\|\mathfrak{E}^*\cdot\mathfrak{E}\|_k}{|n_\ell|^2\cdot|\Xi|},
\]
and, as a function of $k$ on these integers, this pressure is non-increasing.
\item[(c)] Let $a_1,a_2,a_3\ge4$ be integers, even or odd.
If one of them is replaced by a larger integer while the other two are kept,
the pressure of the torus so obtained is at most $f(D(a_1,a_2,a_3))$.
\end{enumerate}
\end{lemma}

\begin{proof}
(a) Put $\mathbb{T}:=\mathbb{T}(\Xi;k\cdot|n_\ell|^2)$.
Lemma~\ref{lem:empty-powers}(c) gives $\Tr\left(\mathfrak{E}^*\cdot\mathfrak{E}\right)^{k}=Z_{\mathbb{T}}[1]$,
and Lemma~\ref{lem:tori-over-section}(b), with $P=k\cdot|n_\ell|^2$, gives $|\mathbb{T}|=k\cdot|n_\ell|^2\cdot|\Xi|$.
The operator $\left(\mathfrak{E}^*\cdot\mathfrak{E}\right)^k$ is non-negative, and hence so is its trace.

(b) Being of trace class, $\mathfrak{E}^*\cdot\mathfrak{E}$ lies in $\mathscr{J}_1$.
By Lemma~\ref{lem:schatten-traces}\eqref{it:schatten-monotone} it lies in $\mathscr{J}_s$ for every real $s\ge1$,
and its norms in these classes satisfy $\|\mathfrak{E}^*\cdot\mathfrak{E}\|_s\le\|\mathfrak{E}^*\cdot\mathfrak{E}\|_{s'}$ whenever $s\ge s'\ge1$.
These are the norms of the statement, since $\|\mathfrak{E}^*\cdot\mathfrak{E}\|_s^s=\Tr\left|\mathfrak{E}^*\cdot\mathfrak{E}\right|^s=\Tr\left(\mathfrak{E}^*\cdot\mathfrak{E}\right)^s$;
so the first assertion holds.
Let $\xi_1\ge\xi_2\ge\dots\ge0$ be the eigenvalues of $\mathfrak{E}^*\cdot\mathfrak{E}$, repeated according to their multiplicities.
Then
\[
\Tr\left(\mathfrak{E}^*\cdot\mathfrak{E}\right)^{k}=\sum_j\xi_j^k ,
\]
and this number is positive if and only if $\mathfrak{E}^*\cdot\mathfrak{E}\neq0$, a condition in which $k$ does not occur.
With (a) it follows that positivity of $Z_{\mathbb{T}(\Xi;k\cdot|n_\ell|^2)}[1]$ at some integer $k$ to which (a) applies forces $\mathfrak{E}^*\cdot\mathfrak{E}\neq0$,
and with it positivity at every such integer.
In that case (a) gives, for every $k$ to which (a) applies,
\[
f\left(\mathbb{T}(\Xi;k\cdot|n_\ell|^2)\right)
=\frac{\log\Tr\left(\mathfrak{E}^*\cdot\mathfrak{E}\right)^{k}}{k\cdot|n_\ell|^2\cdot|\Xi|}
=\frac{\log\|\mathfrak{E}^*\cdot\mathfrak{E}\|_k}{|n_\ell|^2\cdot|\Xi|}.
\]
Here $|n_\ell|^2\cdot|\Xi|>0$, and $k\mapsto\|\mathfrak{E}^*\cdot\mathfrak{E}\|_k$ is non-increasing by the first assertion;
so the pressure is non-increasing in $k$.

(c) Let $a_1,a_2,a_3\ge4$ be integers.
By Lemma~\ref{lem:systoles}, the period lattice $\Z a_1\,e_-+\Z a_2\,e_++\Z a_3\,e_3$ of $D(a_1,a_2,a_3)$ has the systole $\min(a_1,a_2,a_3)\ge4$.
So $D(a_1,a_2,a_3)$ is a torus (Definition~\ref{def:tori}), and $Z_{D(a_1,a_2,a_3)}[1]>0$ by Theorem~\ref{thm:partition-positivity}.
The formula for the maximum norm in Lemma~\ref{lem:systoles} shows in the same way that the three lattices
\[
\Z a_2\,e_++\Z a_3\,e_3,\qquad \Z a_1\,e_-+\Z a_3\,e_3,\qquad \Z a_1\,e_-+\Z a_2\,e_+
\]
have systole at least $4$.
Each of them is of rank two;
the first is contained in the plane $\{\ell_-=0\}$, the second in $\{\ell_+=0\}$, the third in $\{x_3=0\}$.
Let $\Xi_-$, $\Xi_+$ and $\Xi_3$ denote these lattices, each with the character that takes the value $-1$ on both generators;
they are cross-sections for $\ell_-$, $\ell_+$ and $x_3$ (Definition~\ref{def:slabs}).
The last generator of $\mathbb{T}(\Xi;P)$ is $\frac{P}{|n_\ell|^2}\cdot n_\ell$, with the character value $-1$ (Lemma~\ref{lem:tori-over-section}(b));
for $\ell_-$ and $P=2a_1$, for instance, it is $\frac{2a_1}{2}\cdot e_-=a_1\,e_-$.
With Definition~\ref{def:tori} this gives
\[
\begin{aligned}
D(a_1,a_2,a_3)&=\mathbb{T}(\Xi_-;2a_1)&&\text{for }\ell=\ell_-:\ |n_\ell|^2=2,\ n_\ell=e_-,\ k=a_1,\\
&=\mathbb{T}(\Xi_+;2a_2)&&\text{for }\ell=\ell_+:\ |n_\ell|^2=2,\ n_\ell=e_+,\ k=a_2,\\
&=\mathbb{T}(\Xi_3;a_3)&&\text{for }\ell=x_3:\ |n_\ell|^2=1,\ n_\ell=e_3,\ k=a_3,
\end{aligned}
\]
where in each line $k$ is the integer with $k\cdot|n_\ell|^2$ equal to the height period;
for the family $x_3$ the elementary step has the height $1/2$, so $k\cdot|n_\ell|^2=k=a_3$.
Let one of the three sides vary while the other two stay fixed,
and let $\ell$ and $\Xi$ be the family and the cross-section of that line of the preceding display in which the varying side is $k$;
this cross-section is formed from the two fixed sides alone.
For every integer $k\ge1$, the lattice $\Gamma_\parallel+\Z\,k\cdot n_\ell$ of (a) is generated by the three vectors $a_1\,e_-$, $a_2\,e_+$, $a_3\,e_3$
with the varying side replaced by $k$; for $\ell=\ell_-$, for instance, it is $\Z k\,e_-+\Z a_2\,e_++\Z a_3\,e_3$.
Lemma~\ref{lem:systoles}, which applies to arbitrary positive integers, gives this lattice the systole
equal to the smallest of the numbers $k$ and the two fixed sides.
Since the fixed sides are at least $4$, the condition of (a) holds precisely for the integers $k\ge4$,
and for each of them the torus $\mathbb{T}(\Xi;k\cdot|n_\ell|^2)$ is $D(a_1,a_2,a_3)$ with the varying side replaced by $k$,
by the same identification as in the display.
The partition functions of all these tori are positive, as shown at the start of the proof of (c) for arbitrary integer sides $\ge4$,
and (b) yields that the pressure of $D(a_1,a_2,a_3)$ is non-increasing in the varying side.
\end{proof}

\begin{remark}\label{rem:half-steps}
For $\ell=x_1$ and the cross-section $\Xi$ with the lattice $n_2\Z e_2+n_3\Z e_3$, $n_2,n_3\ge4$, and the character $-1$ on both generators,
one has $\mathbb{T}(\Xi;k)=C(k,n_2,n_3)$ for every integer $k\ge4$, and $\mathfrak{E}$ is the half step of Corollary~\ref{cor:trace-forms}\eqref{it:unit-step}.
In this case part (b) of Lemma~\ref{lem:section-monotone} is the monotonicity in $n_1$ of Proposition~\ref{prop:coordinate-pressure}(a),
and its proof is the argument given there:
a power of the partition function is a Schatten norm of a single non-negative operator of trace class,
and such a norm does not increase with the exponent.
For the integrand $1$, the monotonicity in Proposition~\ref{prop:periodic-perturbations}(a) rests on the same argument.
Lemma~\ref{lem:section-monotone} carries this argument over to any family $\ell$ and any cross-section $\Xi$ for it;
the half step is replaced there by the elementary step $\mathfrak{E}$, whose height is $|n_\ell|^2/2$.
In part (c), the height periods of $D(a_1,a_2,a_3)$ for $\ell_-$ and $\ell_+$ are $2a_1$ and $2a_2$,
which lie in $2\Z=|n_{\ell_\pm}|^2\Z$ for all integers $a_1$ and $a_2$.
For $x_3$ the elementary step is a half step, from a site plane to a link plane or from a link plane to a site plane,
and both kinds of planes are mirrors of the family $x_3$.
Hence the height period $a_3$ is unrestricted as well, and $a_1$, $a_2$, $a_3$ range over all integers $\ge4$, of either parity.
\end{remark}

\begin{lemma}[An odd height period]\label{lem:odd-period}
Let $-1\le\rho\le1$, $w_0,\dots,w_3>0$, $M\in\R$ and $q>0$.
Let $\mathcal{Q}=(\mathfrak{F};o_1,o_2,o_3;h)$ be a cell structure and $a\in\{1,2,3\}$ such that $\ell:=\ell^a$ is a diagonal family,
so that $|n_\ell|^2=2$ and $\Adm(\ell)=\Z$.
Let $\mathbb{T}=(\Gamma,\chi)$ be a torus, $\Xi:=\Xi(\mathbb{T},\ell)=(\Gamma_\parallel,\chi_\parallel)$,
and assume that the height period $p:=p(\mathbb{T},\ell)$ is odd.
Let $k\ge1$ be an integer with $2k\cdot h<p$, and put $h^\circ:=p-2k\cdot h$, an odd integer $\ge1$; in particular $p\ge3$.
Assume that $h^\circ\le p/2$ and that
\[
\sys\left(\Gamma_\parallel+\Z\,k\cdot h\cdot n_\ell\right)\ge4,\qquad
\sys\left(\Gamma_\parallel+\Z\,\frac{p-1}{2}\cdot n_\ell\right)\ge4.
\]
Put $\mathbb{T}^{\mathrm{reg}}:=\mathbb{T}(\Xi;2k\cdot h)$ and $\mathbb{T}^{\flat}:=\mathbb{T}(\Xi;p-1)$;
these are tori by Lemma~\ref{lem:tori-over-section}, since $2k\cdot h$ and $p-1$ lie in $2\Z=|n_\ell|^2\Z$.
Let $\mathfrak{G}=(G_\kappa)_{\kappa\in\mathcal{K}}$ be a finite filling of $\mathcal{Q}$,
with the product $\prod\mathfrak{G}=\prod_{\kappa\in\mathcal{K}}G_\kappa$ and the folds $\mathfrak{R}_a^{j,2k}\mathfrak{G}$ of its layers (Definition~\ref{def:fillings}),
such that $0\le\kappa_a\le2k-1$ for all $\kappa\in\mathcal{K}$.
Then $Z_{\mathbb{T}^{\mathrm{reg}}}\left[\prod\mathfrak{R}_a^{j,2k}\mathfrak{G}\right]\ge0$ for all $j\in\Z$, $Z_{\mathbb{T}^{\flat}}[1]\ge0$, and
\begin{equation}\label{eq:sides-odd-slab}
\begin{aligned}
\left|Z_{\mathbb{T}}\left[\prod\mathfrak{G}\right]\right|
&\le Z_{\mathbb{T}^{\flat}}[1]^{\frac{h^\circ}{p-1}}\cdot\prod_{j=0}^{2k-1}Z_{\mathbb{T}^{\mathrm{reg}}}\left[\prod\mathfrak{R}_a^{j,2k}\mathfrak{G}\right]^{\frac{1}{2k}},\\
|\mathbb{T}|&=|\mathbb{T}^{\mathrm{reg}}|+\frac{h^\circ}{p-1}\cdot|\mathbb{T}^{\flat}|.
\end{aligned}
\end{equation}
\end{lemma}

\begin{proof}
The entries of $n_\ell$ are $0$ and $\pm1$, and $p$ is odd; so $\frac{p}{2}\cdot n_\ell$ is not an integer vector,
and the straightened torus $\mathbb{T}^0(\ell)$ is not defined (Lemma~\ref{lem:tori-over-section}(c)).
In \eqref{eq:sides-odd-slab} the torus $\mathbb{T}^{\flat}$ stands where $\mathbb{T}^0(\ell)$ stands in \eqref{eq:sides-low-slab}.

\emph{Volumes.}
By Lemma~\ref{lem:tori-over-section}(b) and (c),
$|\mathbb{T}|=p\cdot|\Xi|$, $|\mathbb{T}^{\mathrm{reg}}|=2k\cdot h\cdot|\Xi|$ and $|\mathbb{T}^{\flat}|=(p-1)\cdot|\Xi|$.
Hence
\[
|\mathbb{T}^{\mathrm{reg}}|+\frac{h^\circ}{p-1}\cdot|\mathbb{T}^{\flat}|=\left(2k\cdot h+h^\circ\right)\cdot|\Xi|=p\cdot|\Xi|=|\mathbb{T}|.
\]

\emph{A trace formula.}
Put $c_j:=o_a+j\cdot h$ for $0\le j\le2k$, and $c_{2k+1}:=o_a+p$.
Since $o_a\in\Adm(\ell)=\Z$ and $h$, $p$ are integers, all these numbers lie in $\Adm(\ell)$,
and $c_0<c_1<\dots<c_{2k}<c_{2k+1}=c_0+p$ because $2k\cdot h<p$.
For $1\le j\le2k$ let
\[
\Phi_j:=\prod_{\kappa\in\mathcal{K},\ \kappa_a=j-1}G_\kappa
\]
be the content of the layer $j-1$ (with $\Phi_j:=1$ if no $\kappa\in\mathcal{K}$ has $\kappa_a=j-1$),
let $Y_j:=\mathfrak{S}[\Phi_j]\colon\mathcal{H}(c_{j-1})\to\mathcal{H}(c_j)$,
and let $\mathfrak{S}^{\circ}:=\mathfrak{S}[1]_{[c_{2k},c_{2k+1}]}$ be the slab operator of the empty slab;
the constant $1$ lies in $\mathcal{A}_\Xi[c_{2k},c_{2k+1}]$.
Lemma~\ref{lem:layer-fold}, applied to the layer $j-1$, shows that $\Phi_j$ lies in $\mathcal{A}_\Xi[c_{j-1},c_j]$,
and \eqref{eq:sides-layer-norm} reads
\[
Z_{\mathbb{T}^{\mathrm{reg}}}\left[\prod\mathfrak{R}_a^{j-1,2k}\mathfrak{G}\right]=\|Y_j\|_{2k}^{2k}\ge0\qquad(1\le j\le2k);
\]
applied to an arbitrary layer, the same lemma gives $Z_{\mathbb{T}^{\mathrm{reg}}}\left[\prod\mathfrak{R}_a^{j,2k}\mathfrak{G}\right]\ge0$ for every $j\in\Z$.
As $0\le\kappa_a\le2k-1$ on $\mathcal{K}$, each $G_\kappa$ is a factor of exactly one $\Phi_j$, and $\prod\mathfrak{G}=\Phi_1\cdots\Phi_{2k}$.
Let $g$ be a closing vector of $\mathbb{T}$ for $\ell$.
The trace formula, Theorem~\ref{thm:slab-operators}\eqref{it:slab-trace}, in the form of Corollary~\ref{cor:slab-admissible} for admissible products,
applied to $\mathbb{T}$, to the $2k+1\ge3$ slabs between the planes $c_0<\dots<c_{2k+1}$ and to the integrands $\Phi_1,\dots,\Phi_{2k},1$, gives
\[
Z_{\mathbb{T}}\left[\prod\mathfrak{G}\right]=\Tr_{\mathcal{H}(c_0)}\left(U\cdot\mathfrak{S}^{\circ}\cdot Y_{2k}\cdots Y_1\right),
\qquad U:=V_g^{-1}\cdot U_{\mathrm{f}}^{\frac{1+\chi(g)}{2}},
\]
where $U\colon\mathcal{H}(c_0+p)\to\mathcal{H}(c_0)$ is unitary.

\emph{H\"older's inequality.}
The exponents $s_j:=p/h$ for $1\le j\le2k$ and $s_{2k+1}:=p/h^\circ$ satisfy
\[
\sum_{j=1}^{2k+1}\frac{1}{s_j}=\frac{2k\cdot h+h^\circ}{p}=1,\qquad \frac{p}{h}>2k\ge2,\qquad \frac{p}{h^\circ}\ge2.
\]
Lemma~\ref{lem:real-exponents}, applied to the operators $Y_1,\dots,Y_{2k},\mathfrak{S}^{\circ}$ and to $U$,
and the monotonicity in its last assertion, applied with $p/h\ge2k\ge2$, give
\[
\left|\Tr\left(U\cdot\mathfrak{S}^{\circ}\cdot Y_{2k}\cdots Y_1\right)\right|
\le\|\mathfrak{S}^{\circ}\|_{p/h^\circ}\cdot\prod_{j=1}^{2k}\|Y_j\|_{p/h}
\le\|\mathfrak{S}^{\circ}\|_{p/h^\circ}\cdot\prod_{j=1}^{2k}\|Y_j\|_{2k}.
\]

\emph{The empty slab.}
Let $c:=c_{2k}$, let $\mathfrak{E}:=\mathfrak{S}[1]_{[c,c+1]}$ be the elementary step at $c$,
and let $b:=(h^\circ-1)/2$, an integer $\ge0$.
The operator of the empty slab factorises as
\[
\mathfrak{S}^{\circ}=V_{b\cdot n_\ell}\cdot\mathfrak{E}\cdot\left(\mathfrak{E}^*\cdot\mathfrak{E}\right)^b .
\]
For $b=0$ this reads $\mathfrak{S}^{\circ}=\mathfrak{S}[1]_{[c,c+1]}=\mathfrak{E}$, with $V_0=1$.
For $b\ge1$, Lemma~\ref{lem:empty-powers}(a) gives $\mathfrak{S}[1]_{[c,c+2b]}=V_{b\cdot n_\ell}\cdot\left(\mathfrak{E}^*\cdot\mathfrak{E}\right)^b$.
The vector $v:=b\cdot n_\ell$ has $\ell(v)=2b$, and the translate of the constant $1$ by $v$ is $1$;
so the covariance, Theorem~\ref{thm:slab-operators}\eqref{it:slab-covariance}, gives
$\mathfrak{S}[1]_{[c+2b,c+2b+1]}=V_{b\cdot n_\ell}\cdot\mathfrak{E}\cdot V_{b\cdot n_\ell}^{-1}$.
Since $c+2b+1=c+h^\circ=c_{2k+1}$, the composition rule, Theorem~\ref{thm:slab-operators}\eqref{it:slab-composition}, yields
\[
\begin{aligned}
\mathfrak{S}^{\circ}&=\mathfrak{S}[1]_{[c+2b,c+2b+1]}\cdot\mathfrak{S}[1]_{[c,c+2b]}
=V_{b\cdot n_\ell}\cdot\mathfrak{E}\cdot V_{b\cdot n_\ell}^{-1}\cdot V_{b\cdot n_\ell}\cdot\left(\mathfrak{E}^*\cdot\mathfrak{E}\right)^b\\
&=V_{b\cdot n_\ell}\cdot\mathfrak{E}\cdot\left(\mathfrak{E}^*\cdot\mathfrak{E}\right)^b .
\end{aligned}
\]
As $V_{b\cdot n_\ell}$ is unitary and $\mathfrak{E}^*\cdot\mathfrak{E}$ is self-adjoint,
\[
\left(\mathfrak{S}^{\circ}\right)^*\cdot\mathfrak{S}^{\circ}
=\left(\mathfrak{E}^*\cdot\mathfrak{E}\right)^b\cdot\mathfrak{E}^*\cdot\mathfrak{E}\cdot\left(\mathfrak{E}^*\cdot\mathfrak{E}\right)^b
=\left(\mathfrak{E}^*\cdot\mathfrak{E}\right)^{2b+1}
=\left(\mathfrak{E}^*\cdot\mathfrak{E}\right)^{h^\circ},
\]
and for real $s\ge2$ the functional calculus of the non-negative operator $\mathfrak{E}^*\cdot\mathfrak{E}$ gives
\[
\|\mathfrak{S}^{\circ}\|_s^s=\Tr\left(\left(\mathfrak{S}^{\circ}\right)^*\cdot\mathfrak{S}^{\circ}\right)^{\frac{s}{2}}=\Tr\left(\mathfrak{E}^*\cdot\mathfrak{E}\right)^{\frac{h^\circ\cdot s}{2}}.
\]
For $s=p/h^\circ$ this becomes
\begin{equation}\label{eq:sides-empty-norm}
\|\mathfrak{S}^{\circ}\|_{p/h^\circ}=\left(\Tr\left(\mathfrak{E}^*\cdot\mathfrak{E}\right)^{\frac{p}{2}}\right)^{\frac{h^\circ}{p}}=\|\mathfrak{E}^*\cdot\mathfrak{E}\|_{p/2}^{\frac{h^\circ}{2}}.
\end{equation}

\emph{Lowering the exponent.}
As $p\ge3$, one has $p/2\ge(p-1)/2\ge1$, and Lemma~\ref{lem:section-monotone}(b) gives
$\|\mathfrak{E}^*\cdot\mathfrak{E}\|_{p/2}\le\|\mathfrak{E}^*\cdot\mathfrak{E}\|_{(p-1)/2}$.
Therefore
\[
\|\mathfrak{S}^{\circ}\|_{p/h^\circ}\le\|\mathfrak{E}^*\cdot\mathfrak{E}\|_{(p-1)/2}^{\frac{h^\circ}{2}}
=\left(\Tr\left(\mathfrak{E}^*\cdot\mathfrak{E}\right)^{\frac{p-1}{2}}\right)^{\frac{h^\circ}{p-1}}.
\]
The integer $(p-1)/2\ge1$ meets the condition of Lemma~\ref{lem:section-monotone}(a):
at this value the lattice of that condition is $\Gamma_\parallel+\Z\,\frac{p-1}{2}\cdot n_\ell$, whose systole is at least $4$ by hypothesis,
and the torus of clause (a) has the height period $\frac{p-1}{2}\cdot|n_\ell|^2=p-1$, that is, it is $\mathbb{T}(\Xi;p-1)=\mathbb{T}^{\flat}$.
Clause (a), taken at the value $(p-1)/2$, therefore gives
\[
\Tr\left(\mathfrak{E}^*\cdot\mathfrak{E}\right)^{\frac{p-1}{2}}=Z_{\mathbb{T}^{\flat}}[1]\ge0 .
\]

\emph{Conclusion.}
The trace formula, H\"older's inequality and the bound on the empty slab combine to
\[
\begin{aligned}
\left|Z_{\mathbb{T}}\left[\prod\mathfrak{G}\right]\right|
&\le\left(\Tr\left(\mathfrak{E}^*\cdot\mathfrak{E}\right)^{\frac{p-1}{2}}\right)^{\frac{h^\circ}{p-1}}\cdot\prod_{j=1}^{2k}\left(\|Y_j\|_{2k}^{2k}\right)^{\frac{1}{2k}}\\
&=Z_{\mathbb{T}^{\flat}}[1]^{\frac{h^\circ}{p-1}}\cdot\prod_{j=1}^{2k}Z_{\mathbb{T}^{\mathrm{reg}}}\left[\prod\mathfrak{R}_a^{j-1,2k}\mathfrak{G}\right]^{\frac{1}{2k}},
\end{aligned}
\]
and the shift $j-1\mapsto j$ of the index in the last product turns this into the inequality \eqref{eq:sides-odd-slab}.
\end{proof}

\begin{remark}\label{rem:both-parities}
Lemmas~\ref{lem:low-slab} and~\ref{lem:odd-period} combine into one inequality for both parities of the height period.
Let $-1\le\rho\le1$, $w_0,\dots,w_3>0$, $M\in\R$ and $q>0$;
let $\mathcal{Q}=(\mathfrak{F};o_1,o_2,o_3;h)$ be a cell structure, $a\in\{1,2,3\}$, $\ell:=\ell^a$,
$\mathbb{T}$ a torus, $\Xi:=\Xi(\mathbb{T},\ell)=(\Gamma_\parallel,\chi_\parallel)$ and $p:=p(\mathbb{T},\ell)$;
let $k\ge1$ be an integer with $2k\cdot h\le p$, put $h^\circ:=p-2k\cdot h$,
and let $\mathfrak{G}$ be a finite filling of $\mathcal{Q}$ with $0\le\kappa_a\le2k-1$ on its index set.
Let $p'$ be the largest element of $|n_\ell|^2\Z$ with $p'\le p$.
For an axis family, $|n_\ell|^2\Z=\Z$ contains the integers $p$ and $h^\circ$;
for a diagonal family with $p$ even, $h^\circ$ is even as well, because $2k\cdot h$ is even.
In these two cases $h^\circ$ and $p$ lie in $|n_\ell|^2\Z$, as Lemma~\ref{lem:low-slab} requires, $p'=p$,
and $\mathbb{T}(\Xi;p)=\mathbb{T}^0(\ell)$ as soon as $\Gamma_\parallel+\Z\,\frac{p}{|n_\ell|^2}\cdot n_\ell$ has systole at least $4$ (Lemma~\ref{lem:tori-over-section}(c)).
Suppose finally that $\ell$ is a diagonal family and that $p$ is odd.
Then $2k\cdot h<p$, since the even number $2k\cdot h$ cannot equal $p$, and $h^\circ$ is odd;
this is the situation of Lemma~\ref{lem:odd-period}, in which $p'=p-1$ and $\mathbb{T}(\Xi;p')=\mathbb{T}^{\flat}$.
Assume now that $h^\circ\le p/2$ and that the two lattices
\[
\Gamma_\parallel+\Z\,\frac{2k\cdot h}{|n_\ell|^2}\cdot n_\ell,\qquad \Gamma_\parallel+\Z\,\frac{p'}{|n_\ell|^2}\cdot n_\ell
\]
have systole at least $4$; these are the remaining hypotheses of both lemmas.
Then, with $\mathbb{T}^{\mathrm{reg}}:=\mathbb{T}(\Xi;2k\cdot h)$, the empty slab contributes in every case the factor $Z_{\mathbb{T}(\Xi;p')}[1]^{\frac{h^\circ}{p'}}$:
\[
\begin{aligned}
\left|Z_{\mathbb{T}}\left[\prod\mathfrak{G}\right]\right|
&\le Z_{\mathbb{T}(\Xi;p')}[1]^{\frac{h^\circ}{p'}}\cdot\prod_{j=0}^{2k-1}Z_{\mathbb{T}^{\mathrm{reg}}}\left[\prod\mathfrak{R}_a^{j,2k}\mathfrak{G}\right]^{\frac{1}{2k}},\\
|\mathbb{T}|&=|\mathbb{T}^{\mathrm{reg}}|+\frac{h^\circ}{p'}\cdot|\mathbb{T}(\Xi;p')|.
\end{aligned}
\]
\end{remark}

\subsection{Comparable sides}\label{sec:comparable-sides}

\begin{proposition}[Comparable even sides]\label{prop:comparable-even}
Let $-1\le\rho\le1$, $w_0,\dots,w_3>0$, $M\in\R$ and $q>0$.
Let $\mathfrak n'=(n'_i)$ be a sequence of positive integers with $n'_i\to\infty$ that satisfies \textup{(Div)},
and let $\mathfrak d=(c_0,c_+,c_3;h;O)$ be data for $\mathfrak n'$.
The threshold $i_0(\mathfrak d)$, the adapted tori $D(n'_i)=D(n'_i/2,n'_i,n'_i)$ and the functionals
\[
\mu^t_i(G)=\frac{Z_{D(n'_i)}\left[\exp\left(tO_{\mathfrak d}\right)\cdot G\right]}{Z_{D(n'_i)}\left[\exp\left(tO_{\mathfrak d}\right)\right]}
\]
are those of Definitions~\ref{def:sides}, \ref{def:data} and~\ref{def:threshold}, formed with the sequence $\mathfrak n'$.
Let $i\ge i_0(\mathfrak d)$, put $n':=n'_i$, and let $n$ be an even integer with $n'\le n\le2n'$.
Then for every $t>0$ the law $\mu_{C(n)}$ of the cubic torus $C(n)=C(n,n,n)$ satisfies
\[
\mu_{C(n)}(O)\le\mu^t_i(O)+|Q_0|\cdot\left(\frac{n}{n'}\right)^{3}\cdot\frac{f(n'/2,n'/2,n'/2)-f_\infty}{t}.
\]
\end{proposition}

\begin{proof}
Fix $t>0$ and put $s:=n-n'$.
Since $i\ge i_0(\mathfrak d)$, the integer $4h$ divides $n'$ and $n'\ge72$ (Definition~\ref{def:threshold}).
In particular $n'$ is even, and so is $s$;
from $n'\le n\le2n'$ one has $0\le s\le n/2$.

\emph{Systoles.}
By Lemma~\ref{lem:systoles}, the period lattice of $C(n_1,n_2,n_3)$ has the systole $\min_\mu n_\mu$,
and the period lattice of $D(a,b,c)$, which is also that of $D'(a,b,c)$ (Definition~\ref{def:tori}), has the systole $\min(a,b,c)$.
Every side of the tori $C(n)$, $D'(n/2,n,n)$, $D'(n'/2,n,n)$, $D(n'/2,n,n)$, $D(n'/2,n',n)$, $D(n')$
and of the coordinate tori $C(n/2,2n,n)$, $C(n'/2,2n',n')$, $C(n'/2)$
is one of the integers $n'/2$, $n/2$, $n'$, $n$, $2n'$, $2n$,
and each of these is at least $n'/2\ge36$, because $n\ge n'\ge72$.
Hence the period lattices of all these tori have systole at least $n'/2\ge36$.

\emph{Cells, contents and the filling.}
As in step~\ref{it:sides-numbers} of the proof of Proposition~\ref{prop:fold-any-side}, let
\[
(2k_1,2k_2,2k_3):=\left(\frac{n'}{h},\ \frac{2n'}{h},\ \frac{n'}{h}\right),\qquad
\mathcal K_0:=\left\{\kappa\in\Z^3:\ 0\le\kappa_a\le2k_a-1,\ a=1,2,3\right\}.
\]
Since $4h$ divides $n'$, the numbers $2k_a$ are even integers $\ge4$,
and $\#|\mathcal K_0|=8k_1k_2k_3=2n'^3/h^3=n'^3/|Q_0|$.
For $\kappa\in\Z^3$ let
\[
v(\kappa):=\frac{h}{2}\left(\kappa_1e_-+\kappa_2e_+\right)+h\kappa_3e_3,\qquad
G_\kappa:=\exp\left(t\,\tau_{v(\kappa)}O\right).
\]
By step~\ref{it:sides-contents} of the proof of Proposition~\ref{prop:fold-any-side},
$G_\kappa$ is a positive bounded continuous function of finitely many $F_\sigma$ with $\sigma$ contained in the open cell $Q_\kappa$ of $\mathcal Q_{\mathfrak d}$.
Put
\[
\widetilde O_{\mathcal K_0}:=\sum_{\kappa\in\mathcal K_0}\tau_{v(\kappa)}O,\qquad
\mathfrak G:=\left(G_\kappa\right)_{\kappa\in\mathcal K_0}.
\]
Then $\mathfrak G$ is a finite filling of $\mathcal Q_{\mathfrak d}$ (Definition~\ref{def:fillings}),
and $\prod\mathfrak G=\exp\left(t\widetilde O_{\mathcal K_0}\right)$.

\emph{Jensen's inequality.}
Lemma~\ref{lem:bosonic-measure}(b), applied on $C(n)$ to the bounded function $\exp\left(t\widetilde O_{\mathcal K_0}\right)$ of finitely many $F_\sigma$,
and Jensen's inequality for the probability measure $\mu_{C(n)}$ and the convex function $\exp$ give
\begin{equation}\label{eq:sides-even-jensen}
\log\frac{Z_{C(n)}\left[\exp\left(t\widetilde O_{\mathcal K_0}\right)\right]}{Z_{C(n)}[1]}
=\log\mu_{C(n)}\left(\exp\left(t\widetilde O_{\mathcal K_0}\right)\right)
\ge t\cdot\mu_{C(n)}\left(\widetilde O_{\mathcal K_0}\right)
=t\cdot\frac{n'^3}{|Q_0|}\cdot\mu_{C(n)}(O).
\end{equation}
The last equality holds because $\tau_vC(n)=C(n)$ for every $v\in\Z^3$ (Lemma~\ref{lem:wall-group}(f)),
so that $\mu_{C(n)}$ is invariant under $F\mapsto\tau_v^*F$ (Lemma~\ref{lem:bosonic-measure}(c));
each of the $n'^3/|Q_0|$ terms $\mu_{C(n)}\left(\tau_{v(\kappa)}O\right)$, $\kappa\in\mathcal K_0$, therefore equals $\mu_{C(n)}(O)$.

\emph{Three applications of the lemma on a low empty slab.}
Lemma~\ref{lem:low-slab} is applied three times to the cell structure $\mathcal Q_{\mathfrak d}=(\mathfrak F_{\mathrm d};c_0,c_+,c_3;h)$.
Its families are $\ell^1=\ell_-$, $\ell^2=\ell_+$, $\ell^3=x_3$, with the normals $n_{\ell_-}=e_-$, $n_{\ell_+}=e_+$, $n_{x_3}=e_3$,
so that $|n_{\ell_-}|^2=|n_{\ell_+}|^2=2$ and $|n_{x_3}|^2=1$.
The cross-sections and the height periods of the three tori $C(n)$, $D'(n'/2,n,n)$, $D(n'/2,n',n)$ are those determined in
steps~\ref{it:sides-first-step}, \ref{it:sides-second-step} and~\ref{it:sides-third-step} of the proof of Proposition~\ref{prop:fold-any-side}.
The tori $\mathbb T(\Xi;P)$ are those of Lemma~\ref{lem:tori-over-section}:
the lattice $\Gamma_\parallel+\Z\,(P/|n_\ell|^2)\,n_\ell$, with the character $\chi_\parallel$ on $\Gamma_\parallel$ and the value $-1$ at the last generator.
Below, a torus or a cross-section is written as its lattice, followed by the values of its character on the generators in the order in which they are written.
The tori $\mathbb T(\Xi;p)$ and $\mathbb T(\Xi;2kh)$ found in the three steps are among those of the paragraph on systoles,
so their lattices have systole $\ge36$.
Hence in each step the straightened torus $\mathbb T^0=\mathbb T^0(\ell)$ is defined and equals $\mathbb T(\Xi;p)$ as soon as $p\in|n_\ell|^2\Z$ (Lemma~\ref{lem:tori-over-section}(c)),
and the lattice of $\mathbb T^{\mathrm{reg}}=\mathbb T(\Xi;2kh)$ has systole $\ge4$.

\emph{Step 1: the direction $\ell_-$ on $\mathbb T=C(n)$.}
The cross-section $\Xi=\Xi(C(n),\ell_-)$ has the lattice $\Gamma_\parallel=\Z\,ne_++\Z\,ne_3$
and the character $\chi_\parallel$ with $\chi_\parallel(ne_+)=+1$ and $\chi_\parallel(ne_3)=-1$;
the height period is $p=n$.
Take $2k:=2k_1$.
Then $2k_1h=n'\le n=p$ and $h^\circ=p-2k_1h=s$;
moreover $s\in2\Z=|n_{\ell_-}|^2\Z$, $s\le n/2=p/2$, and $p=n\in2\Z$ because $n$ is even.
So the hypotheses of Lemma~\ref{lem:low-slab} hold, and
\begin{align*}
\mathbb T^0&=\mathbb T(\Xi;n)=\left(\Z\,ne_++\Z\,ne_3+\Z\,\frac{n}{2}e_-;\ +1,-1,-1\right)=D'(n/2,n,n),\\
\mathbb T^{\mathrm{reg}}&=\mathbb T(\Xi;n')=D'(n'/2,n,n).
\end{align*}

\emph{Step 2: the direction $\ell_+$ on $\mathbb T=D'(n'/2,n,n)$.}
Here $\Xi=\left(\Z\,\frac{n'}{2}e_-+\Z\,ne_3;\ -1,-1\right)$ and $p=\ell_+(ne_+)=2n$.
Take $2k:=2k_2$.
Then $2k_2h=2n'\le2n=p$ and $h^\circ=2n-2n'=2s$;
moreover $2s\in2\Z=|n_{\ell_+}|^2\Z$, $2s\le n=p/2$, and $p=2n\in2\Z$.
So the hypotheses of Lemma~\ref{lem:low-slab} hold, and
\begin{align*}
\mathbb T^0&=\mathbb T(\Xi;2n)=\left(\Z\,\frac{n'}{2}e_-+\Z\,ne_3+\Z\,ne_+;\ -1,-1,-1\right)=D(n'/2,n,n),\\
\mathbb T^{\mathrm{reg}}&=\mathbb T(\Xi;2n')=D(n'/2,n',n).
\end{align*}

\emph{Step 3: the direction $x_3$ on $\mathbb T=D(n'/2,n',n)$.}
Here $\Xi=\left(\Z\,\frac{n'}{2}e_-+\Z\,n'e_+;\ -1,-1\right)$ and $p=n$.
Take $2k:=2k_3$.
Then $2k_3h=n'\le n=p$ and $h^\circ=s$;
moreover $s\in\Z=|n_{x_3}|^2\Z$, $s\le n/2=p/2$, and $p=n\in\Z$.
So the hypotheses of Lemma~\ref{lem:low-slab} hold, and
\begin{align*}
\mathbb T^0&=\mathbb T(\Xi;n)=\left(\Z\,\frac{n'}{2}e_-+\Z\,n'e_++\Z\,ne_3;\ -1,-1,-1\right)=D(n'/2,n',n)=\mathbb T,\\
\mathbb T^{\mathrm{reg}}&=\mathbb T(\Xi;n')=D(n').
\end{align*}

\emph{The iterated folds.}
For $j=(j_1,j_2,j_3)\in\mathcal K_0$ put
\[
\mathfrak G^{j_1}:=\mathfrak R_1^{j_1,2k_1}\mathfrak G,\qquad
\mathfrak G^{j_1j_2}:=\mathfrak R_2^{j_2,2k_2}\mathfrak G^{j_1},\qquad
\mathfrak G^{j_1j_2j_3}:=\mathfrak R_3^{j_3,2k_3}\mathfrak G^{j_1j_2}.
\]
By Definition~\ref{def:fillings} the fold of a layer in the direction $a$ changes only the $a$-th entries of the indices:
the other two entries of an index of the folded filling are those of an index of the filling that is folded.
Hence the first entries of the indices of $\mathfrak G$, the second entries of the indices of $\mathfrak G^{j_1}$
and the third entries of the indices of $\mathfrak G^{j_1j_2}$ are entries of elements of $\mathcal K_0$:
in each of the three steps the filling satisfies the hypothesis $0\le\kappa_a\le2k_a-1$ of Lemma~\ref{lem:low-slab}.
The index set of $\mathfrak G^{j_1j_2j_3}$ is $j+\mathcal K_0$, and its content in the cell $\kappa$ is
\[
\left(\left(G_j^{\,g^1_{\kappa_1\leftarrow j_1}}\right)^{g^2_{\kappa_2\leftarrow j_2}}\right)^{g^3_{\kappa_3\leftarrow j_3}}
=G_j^{\,g^3_{\kappa_3\leftarrow j_3}g^2_{\kappa_2\leftarrow j_2}g^1_{\kappa_1\leftarrow j_1}}
=G_j^{[\kappa]},
\]
by \eqref{eq:setting-group-action} and because $g^3_{\kappa_3\leftarrow j_3}g^2_{\kappa_2\leftarrow j_2}g^1_{\kappa_1\leftarrow j_1}$
is the element of $W(\mathcal Q_{\mathfrak d})$ that maps $Q_j$ onto $Q_\kappa$ (Definitions~\ref{def:cell-families} and~\ref{def:folds}).
Thus $\prod\mathfrak G^{j_1j_2j_3}=\prod_{\kappa\in j+\mathcal K_0}G_j^{[\kappa]}$,
and step~\ref{it:sides-full-fold} of the proof of Proposition~\ref{prop:fold-any-side} gives
\[
Z_{D(n')}\left[\prod\mathfrak G^{j_1j_2j_3}\right]=Z_{D(n')}\left[\exp\left(tO_{\mathfrak d}\right)\right]\qquad(j\in\mathcal K_0).
\]

\emph{The three inequalities.}
By Lemma~\ref{lem:bosonic-measure}(b), the number $Z_{C(n)}\left[\prod\mathfrak G\right]=Z_{C(n)}\left[\exp\left(t\widetilde O_{\mathcal K_0}\right)\right]$
equals $Z_{C(n)}[1]\cdot\mu_{C(n)}\left(\exp\left(t\widetilde O_{\mathcal K_0}\right)\right)$.
The first factor is positive by Lemma~\ref{lem:bosonic-measure}(a),
and the second, the mean of an everywhere positive function under a probability measure, is positive as well.
Hence $Z_{C(n)}\left[\prod\mathfrak G\right]>0$, and in Step~1 the left-hand side of \eqref{eq:sides-low-slab} may be written without the modulus.
Apply \eqref{eq:sides-low-slab} in Step~1 to the filling $\mathfrak G$;
in Step~2 to each filling $\mathfrak G^{j_1}$, for which $Z_{D'(n'/2,n,n)}\left[\prod\mathfrak G^{j_1}\right]\ge0$ by Lemma~\ref{lem:low-slab} in Step~1;
and in Step~3 to each filling $\mathfrak G^{j_1j_2}$, for which $Z_{D(n'/2,n',n)}\left[\prod\mathfrak G^{j_1j_2}\right]\ge0$ by Lemma~\ref{lem:low-slab} in Step~2.
With all products taken over $0\le j_a\le2k_a-1$, this gives
\begin{align*}
Z_{C(n)}\left[\exp\left(t\widetilde O_{\mathcal K_0}\right)\right]
&\le Z_{D'(n/2,n,n)}[1]^{\frac{s}{n}}\cdot\prod_{j_1}Z_{D'(n'/2,n,n)}\left[\prod\mathfrak G^{j_1}\right]^{\frac{1}{2k_1}},\\
Z_{D'(n'/2,n,n)}\left[\prod\mathfrak G^{j_1}\right]
&\le Z_{D(n'/2,n,n)}[1]^{\frac{2s}{2n}}\cdot\prod_{j_2}Z_{D(n'/2,n',n)}\left[\prod\mathfrak G^{j_1j_2}\right]^{\frac{1}{2k_2}},\\
Z_{D(n'/2,n',n)}\left[\prod\mathfrak G^{j_1j_2}\right]
&\le Z_{D(n'/2,n',n)}[1]^{\frac{s}{n}}\cdot\prod_{j_3}Z_{D(n')}\left[\exp\left(tO_{\mathfrak d}\right)\right]^{\frac{1}{2k_3}}.
\end{align*}
All numbers in these lines are $\ge0$, the exponents $1/(2k_a)$ are positive and $s/n\ge0$;
so the third line may be inserted in the second, and the result in the first.
In the combined inequality the factor $Z_{D(n'/2,n,n)}[1]$ carries the exponent $2k_1\cdot\frac{1}{2k_1}\cdot\frac{s}{n}=\frac{s}{n}$,
the factor $Z_{D(n'/2,n',n)}[1]$ the exponent $2k_1\cdot2k_2\cdot\frac{1}{2k_1\cdot2k_2}\cdot\frac{s}{n}=\frac{s}{n}$,
and there are $8k_1k_2k_3=n'^3/|Q_0|$ factors $Z_{D(n')}\left[\exp\left(tO_{\mathfrak d}\right)\right]^{\frac{1}{8k_1k_2k_3}}$.
Therefore
\begin{equation}\label{eq:sides-even-bound}
Z_{C(n)}\left[\exp\left(t\widetilde O_{\mathcal K_0}\right)\right]
\le\left(Z_{D'(n/2,n,n)}[1]\cdot Z_{D(n'/2,n,n)}[1]\cdot Z_{D(n'/2,n',n)}[1]\right)^{\frac{s}{n}}\cdot Z_{D(n')}\left[\exp\left(tO_{\mathfrak d}\right)\right].
\end{equation}

\emph{Pressures.}
By Theorem~\ref{thm:partition-positivity}, $Z_{\mathbb T}[1]>0$ for every torus $\mathbb T$;
so the pressures $f(\mathbb T)=|\mathbb T|^{-1}\log Z_{\mathbb T}[1]$ below are defined.
Put $f_{n'}:=f(n'/2,n'/2,n'/2)$, the pressure of the cubic torus $C(n'/2)$.
\begin{enumerate}
\item
The side $n$ is even and $n\ge8$.
By Theorem~\ref{thm:four-tori}, $Z_{D'(n/2,n,n)}[1]\le Z_{C(n/2,2n,n)}[1]$,
and $|D'(n/2,n,n)|=|C(n/2,2n,n)|=n^3$;
hence $f(D'(n/2,n,n))\le f(n/2,2n,n)$.
The pressure $f$ of coordinate tori is non-increasing in each side (Proposition~\ref{prop:coordinate-pressure}(a)),
and $n/2$, $2n$, $n$ are $\ge n'/2$;
hence $f(n/2,2n,n)\le f_{n'}$.
\item
The side $n'$ is even and $n'\ge8$.
By Theorem~\ref{thm:four-tori} with $n'$ in the place of $n$, $Z_{D(n')}[1]\le Z_{C(n'/2,2n',n')}[1]$,
and $|D(n')|=|C(n'/2,2n',n')|=n'^3$.
With Proposition~\ref{prop:coordinate-pressure}(a) this gives $f(D(n'))\le f(n'/2,2n',n')\le f_{n'}$.
\item\label{it:sides-pressures}
Every adapted torus $D(a,b,c)$ is the rectangular torus $\mathbb{T}(\mathfrak{F}_{\mathrm{d}};2a,2b,c)$ of the diagonal frame,
with the height periods $2a$, $2b$ and $c$ (Definition~\ref{def:cell-structures}).
Let $\mathcal{Q}^{(1)}:=(\mathfrak{F}_{\mathrm{d}};0,0,0;1)$ be the cell structure of this frame with all three positions $0$ and with the width $1$;
its positions are admissible, since $0$ lies in $\Adm(\ell_-)=\Adm(\ell_+)=\Z$ and in $\Adm(x_3)=\tfrac12\Z$ (Definition~\ref{def:mirrors}).
As the width is $1$, a shape $\mathfrak{m}=(m_1,m_2,m_3)$ for $\mathcal{Q}^{(1)}$ (Definition~\ref{def:shapes})
is a triple of even positive integers for which $\mathbb{T}(\mathfrak{F}_{\mathrm{d}};\mathfrak{m}):=\mathbb{T}(\mathfrak{F}_{\mathrm{d}};m_1,m_2,m_3)$ is defined and is a torus.
By Definition~\ref{def:cell-structures} this is the case exactly when $m_1/2$, $m_2/2$ and $m_3$ are at least $4$,
and then $\mathbb{T}(\mathfrak{F}_{\mathrm{d}};\mathfrak{m})=D(m_1/2,m_2/2,m_3)$.
Thus
\[
\mathcal{M}\left(\mathcal{Q}^{(1)}\right)=\left\{\mathfrak{m}\in(2\Z)^3:\ m_1\ge8,\ m_2\ge8,\ m_3\ge4\right\}.
\]
Take the constant $1$ as the content of the cell of $\mathcal{Q}^{(1)}$ with the index $0$;
it lies in $\mathcal{A}(U)$ for every open set $U$, being the product with $r=0$ and $\Phi=1$ in Definition~\ref{def:integrands}.
The action~\eqref{eq:setting-group-action} of $W$ fixes the constant $1$,
since on bosonic integrands $G^g(F)$ is $G(g^*F)$ or its complex conjugate.
Hence every mirror copy of $1$ is $1$ (Definition~\ref{def:folds}).
For a shape $\mathfrak{m}$ the torus $\mathbb{T}(\mathfrak{F}_{\mathrm{d}};\mathfrak{m})$ carries $\mathcal{Q}^{(1)}$, its height periods $m_1,m_2,m_3$ being even
(Definition~\ref{def:cell-structures}),
and $\Fold_{\mathbb{T}(\mathfrak{F}_{\mathrm{d}};\mathfrak{m})}1=1$:
the folded partition function of $1$ at the shape $\mathfrak{m}$ is $Z_{\mathbb{T}(\mathfrak{F}_{\mathrm{d}};\mathfrak{m})}[1]$,
which is positive by Theorem~\ref{thm:partition-positivity}.
Proposition~\ref{prop:periodic-perturbations}\eqref{it:shape-limit}, applied to $\mathcal{Q}^{(1)}$ and to the content $1$, therefore shows that the function
\[
\mathfrak{m}\longmapsto|\mathfrak{m}|^{-1}\cdot\log Z_{\mathbb{T}(\mathfrak{F}_{\mathrm{d}};\mathfrak{m})}[1],\qquad|\mathfrak{m}|=m_1m_2m_3,
\]
is non-increasing in each of the three entries of $\mathfrak{m}$ on $\mathcal{M}\left(\mathcal{Q}^{(1)}\right)$.
The torus $\mathbb{T}(\mathfrak{F}_{\mathrm{d}};\mathfrak{m})=D(m_1/2,m_2/2,m_3)$ has $2\cdot(m_1/2)\cdot(m_2/2)\cdot m_3=|\mathfrak{m}|/2$ sites (Definition~\ref{def:tori}),
so that its pressure is
\[
f\left(\mathbb{T}(\mathfrak{F}_{\mathrm{d}};\mathfrak{m})\right)=\frac{2}{|\mathfrak{m}|}\cdot\log Z_{\mathbb{T}(\mathfrak{F}_{\mathrm{d}};\mathfrak{m})}[1],
\]
and it is non-increasing in each entry of $\mathfrak{m}$ on $\mathcal{M}\left(\mathcal{Q}^{(1)}\right)$ as well.
Since $n'$ and $n$ are even, the three triples
\[
(n',2n,n),\qquad(n',2n',n),\qquad(n',2n',n')
\]
consist of even integers, all at least $n'\ge72$;
so they are shapes for $\mathcal{Q}^{(1)}$, and their rectangular tori are $D(n'/2,n,n)$, $D(n'/2,n',n)$ and $D(n')$.
The first two triples differ only in the second entry, where $2n\ge2n'$,
and the last two only in the third entry, where $n\ge n'$.
It follows that
\[
f(D(n'/2,n,n))\le f(D(n'/2,n',n))\le f(D(n')),
\]
and the last number is at most $f_{n'}$ by the preceding item.
\end{enumerate}
So $f(\mathbb T)\le f_{n'}$ for the four tori $\mathbb T=D'(n/2,n,n)$, $D(n'/2,n,n)$, $D(n'/2,n',n)$, $D(n')$.
Since $|D(a,b,c)|=|D'(a,b,c)|=2abc$ (Definition~\ref{def:tori}), their volumes are $n^3$, $n'\cdot n^2$, $n'^2\cdot n$ and $n'^3$, and
\[
s\cdot\left(n^2+n\cdot n'+n'^2\right)+n'^3=(n-n')\cdot\left(n^2+n\cdot n'+n'^2\right)+n'^3=n^3.
\]

\emph{Conclusion.}
The partition functions in \eqref{eq:sides-even-bound} are positive,
by Lemma~\ref{lem:bosonic-measure}(a), (ii) for the two perturbed ones and by Theorem~\ref{thm:partition-positivity} for the others;
so \eqref{eq:sides-even-bound} may be read on the logarithmic scale.
Proposition~\ref{prop:fold-finite}(c), for the sequence $\mathfrak n'$, gives
\[
\log Z_{D(n')}\left[\exp\left(tO_{\mathfrak d}\right)\right]\le t\cdot\frac{n'^3}{|Q_0|}\cdot\mu^t_i(O)+\log Z_{D(n')}[1].
\]
Moreover $\log Z_{\mathbb T}[1]=|\mathbb T|\cdot f(\mathbb T)\le|\mathbb T|\cdot f_{n'}$ for the four tori above,
$\log Z_{C(n)}[1]=n^3\cdot f(n,n,n)$, and $f(n,n,n)\ge f_\infty$, where $f_\infty=\inf f$ is finite (Proposition~\ref{prop:coordinate-pressure}(d)).
Together with \eqref{eq:sides-even-jensen} and \eqref{eq:sides-even-bound} this yields
\begin{align*}
t\cdot\frac{n'^3}{|Q_0|}\cdot\mu_{C(n)}(O)
&\le\log Z_{C(n)}\left[\exp\left(t\widetilde O_{\mathcal K_0}\right)\right]-\log Z_{C(n)}[1]\\
&\le\frac{s}{n}\cdot\left(n^3+n'\cdot n^2+n'^2\cdot n\right)\cdot f_{n'}+t\cdot\frac{n'^3}{|Q_0|}\cdot\mu^t_i(O)+n'^3\cdot f_{n'}-n^3\cdot f_\infty\\
&=t\cdot\frac{n'^3}{|Q_0|}\cdot\mu^t_i(O)+n^3\cdot\left(f_{n'}-f_\infty\right).
\end{align*}
Since $t\cdot n'^3/|Q_0|>0$ and $n^3\cdot|Q_0|/n'^3=|Q_0|\cdot(n/n')^3$, division by $t\cdot n'^3/|Q_0|$ leaves
\[
\mu_{C(n)}(O)\le\mu^t_i(O)+|Q_0|\cdot\left(\frac{n}{n'}\right)^{3}\cdot\frac{f_{n'}-f_\infty}{t}.
\]
\end{proof}

\begin{remark}\label{rem:pressure-comparison}
Let $i\ge i_0(\mathfrak{d})$ and $n'=n'_i$ be as in Proposition~\ref{prop:comparable-even};
thus $n'$ is even and $n'\ge72$ (Definition~\ref{def:threshold}).
Let $n\ge n'$ be an integer, odd or even.
The sides of the three tori $D(n'/2,n,n)$, $D(n'/2,n',n)$ and $D(n')$
are integers not smaller than $n'/2\ge36$.
The second of these tori arises from the first when the second side is lowered from $n$ to $n'$,
and the third arises from the second when the third side is lowered from $n$ to $n'$.
By Lemma~\ref{lem:section-monotone}(c), the pressure of $D(a_1,a_2,a_3)$ is non-increasing in each of the integers $a_1,a_2,a_3\ge4$, whatever their parity.
Consequently
\[
f(D(n'/2,n,n))\le f(D(n'/2,n',n))\le f(D(n')).
\]
\end{remark}

For an odd side $n$ the height period of $C(n)$ along $\ell_-$ is odd;
the slab left empty in this direction is then estimated by Lemma~\ref{lem:odd-period}.

\begin{proposition}[Comparable odd sides]\label{prop:comparable-odd}
Let $-1\le\rho\le 1$, $w_0,\dots,w_3>0$, $M\in\R$ and $q>0$.
Let $\mathfrak{n}'=(n'_i)$ be a sequence with \textup{(Div)},
let $\mathfrak{d}=(c_0,c_+,c_3;h;O)$ be data for $\mathfrak{n}'$ and $\mu^t_i$ the perturbed functionals on the tori $D(n'_i)$ (Definition~\ref{def:data}),
let $i\ge i_0(\mathfrak{d})$ (Definition~\ref{def:threshold}) and $n':=n'_i$,
and let $n$ be an odd integer with $n'<n\le 2n'$.
Then the law $\mu_{C(n)}$ of the cubic torus $C(n)=C(n,n,n)$ satisfies, for every $t>0$,
\[
\mu_{C(n)}(O)\le\mu^t_i(O)+|Q_0|\cdot\left(\frac{n}{n'}\right)^3\cdot\frac{f(n'/2,n'/2,n'/2)-f_\infty}{t}.
\]
\end{proposition}

\begin{proof}
Fix $t>0$.
The notation is that of Proposition~\ref{prop:fold-any-side}
and of steps~\ref{it:sides-numbers}, \ref{it:sides-contents} and~\ref{it:sides-full-fold} of its proof,
which are valid for every integer side $n\ge n'$:
the even integers $2k_1=n'/h$, $2k_2=2n'/h$ and $2k_3=n'/h$,
the translations $v(\kappa)$,
the contents $G_\kappa=\exp(t\,\tau_{v(\kappa)}O)$ for $\kappa\in\Z^3$, with $G_0=\exp(tO)$,
and the set
\[
\mathcal{K}_0=\{\kappa\in\Z^3:\ 0\le\kappa_a\le 2k_a-1,\ a=1,2,3\},
\qquad
\#|\mathcal{K}_0|=8k_1k_2k_3=\frac{n'^3}{|Q_0|}.
\]
Put $s:=n-n'$ and $a_1:=(n-1)/2$.
As $4h$ divides $n'$ (Definition~\ref{def:threshold}) and $n$ is odd, $s$ is an odd integer with $s\ge 1$,
and $s\le n/2$ because $n\le 2n'$.
Moreover $n-1\ge n'$, since $n>n'$, and $n'\ge 72$ by step~\ref{it:sides-numbers};
so $a_1$ is an integer with $a_1\ge n'/2\ge 36$.
Let
\[
\widetilde O_{\mathcal{K}_0}:=\sum_{\kappa\in\mathcal{K}_0}\tau_{v(\kappa)}O,
\qquad
\mathfrak{G}:=(G_\kappa)_{\kappa\in\mathcal{K}_0}.
\]
By step~\ref{it:sides-contents}, $\mathfrak{G}$ is a finite filling of $\mathcal{Q}_{\mathfrak{d}}$,
and $\prod\mathfrak{G}=\prod_{\kappa\in\mathcal{K}_0}G_\kappa=\exp(t\,\widetilde O_{\mathcal{K}_0})$.

\emph{Jensen's inequality.}
The function $\exp(t\,\widetilde O_{\mathcal{K}_0})$ is bounded, continuous and depends on finitely many $F_\sigma$.
By Lemma~\ref{lem:bosonic-measure}(b), the functional $Z_{C(n)}$ takes at this function the value $Z_{C(n)}[1]$ times its expectation under the probability measure $\mu_{C(n)}$.
Jensen's inequality for $\mu_{C(n)}$ and the exponential function therefore gives
\[
\log\frac{Z_{C(n)}\left[\exp\left(t\,\widetilde O_{\mathcal{K}_0}\right)\right]}{Z_{C(n)}[1]}
=\log\mu_{C(n)}\left(\exp\left(t\,\widetilde O_{\mathcal{K}_0}\right)\right)
\ge t\cdot\mu_{C(n)}\left(\widetilde O_{\mathcal{K}_0}\right)
=t\cdot\frac{n'^3}{|Q_0|}\cdot\mu_{C(n)}(O).
\]
In the last step, every translation $\tau_v$ with $v\in\Z^3$ maps $C(n)$ onto itself (Lemma~\ref{lem:wall-group}(f)),
so $\mu_{C(n)}$ is invariant under $\tau_v$ (Lemma~\ref{lem:bosonic-measure}(c)),
and each of the $n'^3/|Q_0|$ terms $\tau_{v(\kappa)}O$ has the expectation $\mu_{C(n)}(O)$.

\emph{Three steps.}
For an arbitrary integer $n\ge n'$, odd or even,
steps~\ref{it:sides-first-step}, \ref{it:sides-second-step} and~\ref{it:sides-third-step} of the proof of Proposition~\ref{prop:fold-any-side}(b$'$)
give the cross-section and the height period of $C(n)$ for $\ell_-$, of $D'(n'/2,n,n)$ for $\ell_+$, and of $D(n'/2,n',n)$ for $x_3$.
Below, a lattice is followed by the values of the character on the generators listed,
and $\mathbb{T}(\Xi;P)$ is the adapted torus of Lemma~\ref{lem:tori-over-section}.
By Lemma~\ref{lem:systoles}, for positive integers $a,b,c$ the period lattices of $D(a,b,c)$ and of $D'(a,b,c)$ have the systole $\min(a,b,c)$,
and that of $C(n)$ has the systole $n$.
In the three steps below, the tori $\mathbb{T}$, $\mathbb{T}^{\mathrm{reg}}$, $\mathbb{T}^{\flat}$ and $\mathbb{T}^0$ are of these kinds,
with sides among $n'/2$, $a_1$, $n'$ and $n$;
so the lattices of all of them have systole at least $\min(n'/2,a_1)=n'/2\ge 36$.
All three steps are taken with the cell structure $\mathcal{Q}_{\mathfrak{d}}$.

In the direction $\ell_-$, Lemma~\ref{lem:odd-period} is applied with $a=1$, $\mathbb{T}=C(n)$,
$\Xi=\Xi(C(n),\ell_-)=(\Z n e_++\Z n e_3;\ +1,-1)$, the odd height period $p=n$,
$2k=2k_1$, $2k_1h=n'<n$ and $h^\circ=n-n'=s\le n/2$.
Then
\[
\mathbb{T}^{\mathrm{reg}}=\mathbb{T}(\Xi;n')=D'(n'/2,n,n),
\qquad
\mathbb{T}^{\flat}=\mathbb{T}(\Xi;n-1)=\left(\Z n e_++\Z n e_3+\Z a_1e_-;\ +1,-1,-1\right)=D'(a_1,n,n),
\]
and the two lattices whose systole the lemma requires to be at least $4$ are the lattices of these two tori.
The filling $\mathfrak{G}$ satisfies $0\le\kappa_1\le 2k_1-1$ for all $\kappa\in\mathcal{K}_0$.

In the direction $\ell_+$, Lemma~\ref{lem:low-slab} is applied with $a=2$, $\mathbb{T}=D'(n'/2,n,n)$,
$\Xi=\Xi(\mathbb{T},\ell_+)=(\Z(n'/2)e_-+\Z n e_3;\ -1,-1)$, $p=2n$,
$2k=2k_2$, $2k_2h=2n'\le 2n$ and $h^\circ=2s$.
Its hypotheses hold: $2s\le n=p/2$;
$h^\circ$ and $p$ lie in $|n_{\ell_+}|^2\Z=2\Z$;
$\mathbb{T}^0=\mathbb{T}(\Xi;2n)=(\Z(n'/2)e_-+\Z n e_3+\Z n e_+;\ -1,-1,-1)=D(n'/2,n,n)$;
and $\mathbb{T}^{\mathrm{reg}}=\mathbb{T}(\Xi;2n')=D(n'/2,n',n)$.
Here the height period is $2n$, so the parity of $n$ plays no role.

In the direction $x_3$, Lemma~\ref{lem:low-slab} is applied with $a=3$, $\mathbb{T}=D(n'/2,n',n)$,
$\Xi=\Xi(\mathbb{T},x_3)=(\Z(n'/2)e_-+\Z n'e_+;\ -1,-1)$, $p=n$,
$2k=2k_3$, $2k_3h=n'\le n$ and $h^\circ=s$.
Its hypotheses hold: $s\le n/2$;
$s$ and $p=n$ lie in $|e_3|^2\Z=\Z$, since the elementary height along $x_3$ is $1/2$, so that the odd height period $n$ is admissible;
$\mathbb{T}^0=\mathbb{T}(\Xi;n)=(\Z(n'/2)e_-+\Z n'e_++\Z n e_3;\ -1,-1,-1)=\mathbb{T}$;
and $\mathbb{T}^{\mathrm{reg}}=\mathbb{T}(\Xi;n')=D(n'/2,n',n')=D(n'_i)$.

\emph{The fillings.}
For $j=(j_1,j_2,j_3)\in\mathcal{K}_0$ let
\[
\mathfrak{G}^{j_1}:=\mathfrak{R}_1^{j_1,2k_1}\mathfrak{G},
\qquad
\mathfrak{G}^{j_1j_2}:=\mathfrak{R}_2^{j_2,2k_2}\mathfrak{G}^{j_1},
\qquad
\mathfrak{G}^{j_1j_2j_3}:=\mathfrak{R}_3^{j_3,2k_3}\mathfrak{G}^{j_1j_2}.
\]
By Definition~\ref{def:fillings}, a fold $\mathfrak{R}_a^{j,2k}$ alters the $a$-th entries of the indices and no other entries:
in every direction $b\ne a$, the entries of the indices of the folded filling are entries of indices of the given filling.
The fold that turns $\mathfrak{G}$ into $\mathfrak{G}^{j_1}$ changes only first entries, and the fold that turns $\mathfrak{G}^{j_1}$ into $\mathfrak{G}^{j_1j_2}$ only second entries.
As the bounds $0\le\kappa_2\le 2k_2-1$ and $0\le\kappa_3\le 2k_3-1$ hold on $\mathcal{K}_0$,
the first of them holds for every index $\kappa$ of $\mathfrak{G}^{j_1}$, and the second for every index $\kappa$ of $\mathfrak{G}^{j_1j_2}$.
These are the bounds on the indices that Lemma~\ref{lem:low-slab} imposes in the directions $a=2$ and $a=3$.
The filling $\mathfrak{G}^{j_1j_2j_3}$ has the index set $j+\mathcal{K}_0$; at an index $\kappa$ of this set its content equals
\[
\left(\left(\left(G_j\right)^{g^1_{\kappa_1\leftarrow j_1}}\right)^{g^2_{\kappa_2\leftarrow j_2}}\right)^{g^3_{\kappa_3\leftarrow j_3}}
=\left(G_j\right)^{g^3_{\kappa_3\leftarrow j_3}g^2_{\kappa_2\leftarrow j_2}g^1_{\kappa_1\leftarrow j_1}}
=\left(G_j\right)^{[\kappa]}
\]
by \eqref{eq:setting-group-action}, Definition~\ref{def:cell-families} and Definition~\ref{def:folds}.
Therefore $\prod\mathfrak{G}^{j_1j_2j_3}=\prod_{\kappa\in j+\mathcal{K}_0}(G_j)^{[\kappa]}$,
and step~\ref{it:sides-full-fold} of the proof of Proposition~\ref{prop:fold-any-side} gives
\[
Z_{D(n'_i)}\left[{\textstyle\prod}\mathfrak{G}^{j_1j_2j_3}\right]=Z_{D(n'_i)}\left[\exp(tO_{\mathfrak{d}})\right]
\qquad(j\in\mathcal{K}_0).
\]

\emph{The inequalities.}
By Lemma~\ref{lem:bosonic-measure}(b), the number $Z_{C(n)}[\prod\mathfrak{G}]=Z_{C(n)}[\exp(t\,\widetilde O_{\mathcal{K}_0})]$
is the integral of the everywhere positive function $\exp(t\,\widetilde O_{\mathcal{K}_0})$ against the measure of Lemma~\ref{lem:bosonic-measure}(a),
a positive measure with the positive total mass $Z_{C(n)}[1]$;
hence $Z_{C(n)}[\prod\mathfrak{G}]>0$.
Apply \eqref{eq:sides-odd-slab} in the direction $\ell_-$;
then \eqref{eq:sides-low-slab} in the direction $\ell_+$ to each of the $2k_1$ numbers $Z_{D'(n'/2,n,n)}[\prod\mathfrak{G}^{j_1}]$,
which are non-negative by Lemma~\ref{lem:odd-period};
finally \eqref{eq:sides-low-slab} in the direction $x_3$ to each of the $2k_1\cdot 2k_2$ numbers $Z_{D(n'/2,n',n)}[\prod\mathfrak{G}^{j_1j_2}]$,
which are non-negative by Lemma~\ref{lem:low-slab}.
In the second step $h^\circ/p=2s/(2n)=s/n$.
With $0\le j_a\le 2k_a-1$ in all products, one obtains
\[
\begin{aligned}
Z_{C(n)}\left[\exp\left(t\,\widetilde O_{\mathcal{K}_0}\right)\right]
&\le Z_{D'(a_1,n,n)}[1]^{\frac{s}{n-1}}\cdot\prod_{j_1}Z_{D'(n'/2,n,n)}\left[{\textstyle\prod}\mathfrak{G}^{j_1}\right]^{\frac{1}{2k_1}},\\
Z_{D'(n'/2,n,n)}\left[{\textstyle\prod}\mathfrak{G}^{j_1}\right]
&\le Z_{D(n'/2,n,n)}[1]^{\frac{s}{n}}\cdot\prod_{j_2}Z_{D(n'/2,n',n)}\left[{\textstyle\prod}\mathfrak{G}^{j_1j_2}\right]^{\frac{1}{2k_2}},\\
Z_{D(n'/2,n',n)}\left[{\textstyle\prod}\mathfrak{G}^{j_1j_2}\right]
&\le Z_{D(n'/2,n',n)}[1]^{\frac{s}{n}}\cdot\prod_{j_3}Z_{D(n'_i)}\left[\exp(tO_{\mathfrak{d}})\right]^{\frac{1}{2k_3}}.
\end{aligned}
\]
Every number in these three lines is non-negative and every exponent is positive,
so the bounds can be chained:
the third line bounds each factor of the product in the second,
and the bound so obtained for the left-hand side of the second line bounds each factor of the product in the first.
After these substitutions, $Z_{D(n'/2,n,n)}[1]$ carries the exponent $2k_1\cdot\frac{1}{2k_1}\cdot\frac{s}{n}=\frac{s}{n}$,
$Z_{D(n'/2,n',n)}[1]$ carries the exponent $2k_1\cdot 2k_2\cdot\frac{1}{2k_1\cdot 2k_2}\cdot\frac{s}{n}=\frac{s}{n}$,
and the factor $Z_{D(n'_i)}[\exp(tO_{\mathfrak{d}})]$ occurs $8k_1k_2k_3$ times, each time with the exponent $\frac{1}{8k_1k_2k_3}$.
Thus
\begin{equation}\label{eq:sides-odd-bound}
Z_{C(n)}\left[\exp\left(t\,\widetilde O_{\mathcal{K}_0}\right)\right]
\le Z_{D'(a_1,n,n)}[1]^{\frac{s}{n-1}}\cdot\left(Z_{D(n'/2,n,n)}[1]\cdot Z_{D(n'/2,n',n)}[1]\right)^{\frac{s}{n}}\cdot Z_{D(n'_i)}\left[\exp(tO_{\mathfrak{d}})\right].
\end{equation}

\emph{Pressures.}
Every partition function $Z_{\mathbb{T}}[1]$ that occurs is positive by Theorem~\ref{thm:partition-positivity},
and $f(\mathbb{T})=|\mathbb{T}|^{-1}\log Z_{\mathbb{T}}[1]$ denotes the pressure of a torus $\mathbb{T}$.
Put $f_{n'}:=f(n'/2,n'/2,n'/2)$.

First, $Z_{D'(a_1,n,n)}[1]\le Z_{D(a_1,n,n)}[1]$.
This inequality is Lemma~\ref{lem:low-slab} for the empty filling of a cell structure other than $\mathcal{Q}_{\mathfrak{d}}$,
namely $(\mathfrak{F}_{\mathrm{d}};0,0,0;n)$, whose width is $n$;
the lemma is taken in the direction $a=2$, with the family $\ell_+$, the torus $\mathbb{T}=D'(a_1,n,n)$ and $k=1$, that is, with two slabs.
The cross-section is $\Xi(\mathbb{T},\ell_+)=(\Z a_1e_-+\Z ne_3;\ -1,-1)$, and the height period is $p=\ell_+(ne_+)=2n$.
Two slabs of the width $n$ make up the whole height period, $2\cdot n=2n=p$;
so the height of the empty slab is $h^\circ=p-2\cdot n=0$, and
\[
\mathbb{T}^0=\mathbb{T}^{\mathrm{reg}}=\mathbb{T}(\Xi;2n)=D(a_1,n,n),
\]
whose lattice has systole $a_1\ge36$.
The empty filling has the product $1$, and so have its folds, which are again empty fillings.
As $h^\circ=0$, the factor $Z_{\mathbb{T}^0}[1]^{\frac{h^\circ}{p}}$ is $1$, and \eqref{eq:sides-low-slab} becomes
\[
\left|Z_{D'(a_1,n,n)}[1]\right|\le\prod_{j=0}^{1}Z_{D(a_1,n,n)}[1]^{\frac{1}{2}}=Z_{D(a_1,n,n)}[1].
\]
As $|D'(a_1,n,n)|=|D(a_1,n,n)|=(n-1)\cdot n^2$, it follows that $f(D'(a_1,n,n))\le f(D(a_1,n,n))$.

Second, by Lemma~\ref{lem:section-monotone}(c), the pressure of a torus $D(a,b,c)$ with integer sides $a,b,c\ge4$ does not decrease
when one of these three sides is replaced by a smaller integer that is still at least $4$.
Since $a_1\ge n'/2$ and $n\ge n'$, the sides may be lowered one at a time:
the first from $a_1$ to $n'/2$, then the second from $n$ to $n'$, then the third from $n$ to $n'$,
and no side ever drops below $n'/2\ge36$.
This yields
\[
f(D(a_1,n,n))\le f(D(n'/2,n,n))\le f(D(n'/2,n',n))\le f(D(n'/2,n',n'))=f(D(n'_i)).
\]

Third, $n'$ is even and at least $8$, so Theorem~\ref{thm:four-tori} gives $Z_{D(n'_i)}[1]\le Z_{C(n'/2,2n',n')}[1]$.
Both tori have $n'^3$ sites; hence $f(D(n'_i))\le f(n'/2,2n',n')$.
Lowering the sides $2n'$ and $n'$ of this coordinate torus to $n'/2$ does not decrease its pressure (Proposition~\ref{prop:coordinate-pressure}(a)),
so that $f(n'/2,2n',n')\le f(n'/2,n'/2,n'/2)=f_{n'}$.

Consequently the pressures of the four tori $D'(a_1,n,n)$, $D(n'/2,n,n)$, $D(n'/2,n',n)$ and $D(n'_i)$ are at most $f_{n'}$.
Their volumes are $(n-1)\cdot n^2$, $n'\cdot n^2$, $n'^2\cdot n$ and $n'^3$ (Definition~\ref{def:tori}), and
\[
\frac{s}{n-1}\cdot(n-1)\cdot n^2+\frac{s}{n}\cdot\left(n'\cdot n^2+n'^2\cdot n\right)+n'^3
=s\cdot\left(n^2+n\cdot n'+n'^2\right)+n'^3
=(n-n')\cdot\left(n^2+n\cdot n'+n'^2\right)+n'^3
=n^3.
\]

\emph{Conclusion.}
All numbers in \eqref{eq:sides-odd-bound} are positive.
For the partition functions $Z_{\mathbb{T}}[1]$ this is Theorem~\ref{thm:partition-positivity}.
The other two, $Z_{C(n)}[\exp(t\,\widetilde O_{\mathcal{K}_0})]$ and $Z_{D(n'_i)}[\exp(tO_{\mathfrak{d}})]$,
are integrals of everywhere positive functions (Lemma~\ref{lem:bosonic-measure}(b))
against positive measures of positive total mass (Lemma~\ref{lem:bosonic-measure}(a)).
Take logarithms in \eqref{eq:sides-odd-bound}.
Proposition~\ref{prop:fold-finite}(c), read for the sequence $\mathfrak{n}'$, gives
\[
\log Z_{D(n'_i)}\left[\exp(tO_{\mathfrak{d}})\right]\le t\cdot\frac{n'^3}{|Q_0|}\cdot\mu^t_i(O)+\log Z_{D(n'_i)}[1].
\]
For each of the four tori, $\log Z_{\mathbb{T}}[1]=|\mathbb{T}|\cdot f(\mathbb{T})\le|\mathbb{T}|\cdot f_{n'}$,
and this bound enters multiplied by the positive number $\frac{s}{n-1}$, $\frac{s}{n}$, $\frac{s}{n}$ or $1$, respectively.
Finally $\log Z_{C(n)}[1]=n^3\cdot f(n,n,n)$, and $f(n,n,n)\ge f_\infty$ (Proposition~\ref{prop:coordinate-pressure}(d)).
Together with Jensen's inequality and the identity of the volumes, this yields
\[
\begin{aligned}
t\cdot\frac{n'^3}{|Q_0|}\cdot\mu_{C(n)}(O)
&\le s\cdot n^2\cdot f_{n'}+\frac{s}{n}\cdot\left(n'\cdot n^2+n'^2\cdot n\right)\cdot f_{n'}
+t\cdot\frac{n'^3}{|Q_0|}\cdot\mu^t_i(O)+n'^3\cdot f_{n'}-n^3\cdot f_\infty\\
&=t\cdot\frac{n'^3}{|Q_0|}\cdot\mu^t_i(O)+n^3\cdot\left(f_{n'}-f_\infty\right).
\end{aligned}
\]
Division by $t\cdot n'^3/|Q_0|$ turns this into the asserted inequality, since $n^3\cdot|Q_0|/n'^3=|Q_0|\cdot(n/n')^3$.
\end{proof}

Since $n'$ is even, an odd integer $n\ge n'$ is larger than $n'$.
The inequality of Propositions~\ref{prop:comparable-even} and~\ref{prop:comparable-odd} therefore holds for every integer $n$ with $n'\le n\le 2n'$:
for even $n$ by the first, for odd $n$ by the second.

\begin{remark}\label{rem:comparable-ratio}
Let $\mathfrak{n}=(n_i)$ be a sequence of integers $n_i\ge 4$ with $n_i\to\infty$, let $\mu_i$ be the law of $C(n_i)$,
and let $\mathfrak{n}'=(n'_i)$ be a sequence with \textup{(Div)} such that $n'_i\le n_i\le 2n'_i$ for all large $i$;
by the three conditions of Definition~\ref{def:companions}, every companion of $\mathfrak{n}$ is such a sequence.
Let $\mathfrak{d}$ be data for $\mathfrak{n}'$; the functionals $\mu^t_i$ and the comparison states for $\mathfrak{d}$ are formed with $\mathfrak{n}'$.
For all large $i$ and all $t>0$,
Proposition~\ref{prop:comparable-even} (if $n_i$ is even) and Proposition~\ref{prop:comparable-odd} (if $n_i$ is odd, hence larger than the even number $n'_i$) give
\[
\mu_i(O)\le\mu^t_i(O)+\frac{\varepsilon_i}{t},
\qquad
\varepsilon_i:=|Q_0|\cdot\left(\frac{n_i}{n'_i}\right)^3\cdot\left(f(n'_i/2,n'_i/2,n'_i/2)-f_\infty\right).
\]
The numbers $\varepsilon_i$ do not depend on $t$, and since $n'_i\to\infty$,
\[
0\le\varepsilon_i\le 8\,|Q_0|\cdot\left(f(n'_i/2,n'_i/2,n'_i/2)-f_\infty\right)\longrightarrow 0
\]
by Proposition~\ref{prop:coordinate-pressure}(d).
This yields the inequality $\nu(O)\le\nu_0(O)$ of Theorem~\ref{thm:comparison-any-side}(c) in the case $\lambda_i=\mu_i$.
Indeed, let $\mu_{i'_m}\to\nu$ weakly,
and let $\nu_0$ be a comparison state for $\mathfrak{d}$ subordinate to $(i'_m)$,
given by a subsequence $(i_m)$ of $(i'_m)$, strengths $t_k\to 0$ and laws $\nu_{t_k}$ as in Definition~\ref{def:comparison-states}.
For fixed $k$, the inequality at $t=t_k$ along $(i_m)$ gives $\nu(O)\le\nu_{t_k}(O)$ in the limit $m\to\infty$,
because $O$ is bounded and continuous on $\Omega$ and $\varepsilon_{i_m}\to 0$;
the limit $k\to\infty$ then gives $\nu(O)\le\nu_0(O)$.
For large $i$ the ratio $n'_i/n_i$ lies in $[1/2,1]$,
and it need not converge to $1$, whereas property~\eqref{it:companion-ratio} of Definition~\ref{def:companions} requires this of a companion.
Hence Theorem~\ref{thm:cubic-limits} can also be derived from Propositions~\ref{prop:comparable-even} and~\ref{prop:comparable-odd},
with any such sequence $\mathfrak{n}'$ in the place of a companion of $\mathfrak{n}$:
in the proof of Lemma~\ref{lem:diagonal-domination} for $\lambda_i=\mu_i$, the inequality just obtained takes the place of Theorem~\ref{thm:comparison-any-side}(c),
and every other step that involves $\mathfrak{n}'$ holds for all sequences with \textup{(Div)}.
\end{remark}

\subsection{The averaged law for arbitrary sides}\label{sec:averaged-arbitrary}

Conditioning the law of a cubic torus on the sector of flux zero gives the harmonic-averaged law.
For an integer $n\ge4$ let $\Phi_0$ be the indicator of the sector of flux zero on the cubic torus $C(n)=C(n,n,n)$,
a bounded Borel function of finitely many $F_\sigma$,
and let $\mu^{\mathrm{av}}_{C(n)}$ be the averaged law of Definition~\ref{def:averaged-law}:
the law of $\widetilde F$ under the sector functional $G\mapsto Z^{(0)}_{C(n)}[G]=Z_{C(n)}[\Phi_0\cdot G]$, normalised by $Z^{(0)}_{C(n)}[1]$.
Thus, for every bounded Borel function $G$ on $\Omega$, evaluated at $\widetilde F$,
\[
\mu^{\mathrm{av}}_{C(n)}(G)=\frac{\mathbb{E}_{C(n)}[\Phi_0\cdot G]}{\mathbb{E}_{C(n)}[\Phi_0]} ,
\]
where the weight $\mathbb{E}_{C(n)}[\Phi_0]=Z_{C(n)}[\Phi_0]/Z_{C(n)}[1]$ of the sector of flux zero is positive by Lemma~\ref{lem:sector-weights}(b).

\begin{theorem}[Limits of the averaged law]\label{thm:averaged-any-side}
Let the chamber operator have the Wilson parameter $0<\rho<1$, weights $w_0,\dots,w_3>0$ and a real mass parameter $M$, and let the lattice charge be $q>0$.
Let $\mathfrak{n}=(n_i)$ be a sequence of integers $n_i\ge4$ with $n_i\to\infty$, and put $\mu^{\mathrm{av}}_i:=\mu^{\mathrm{av}}_{C(n_i)}$.
Then the sequence $(\mu^{\mathrm{av}}_i)$ is tight on $\Omega$, and every weak limit point $\nu$ of it
satisfies~\eqref{cond:gaussian-domination} and $\nu(dF=0)=1$,
is invariant under $W$ and under $F\mapsto-F$,
and is reflection positive (Definition~\ref{def:reflection-positivity}) at every mirror of the six diagonal families (Definition~\ref{def:mirrors}).
If all $n_i$ are even, then $\nu$ is reflection positive at every mirror of the nine families.
\end{theorem}

\begin{proof}
\emph{Tightness, Gaussian domination, closedness.}
For an index $i$ let $\Gamma:=n_i\Z^3$ be the period lattice of $C(n_i)$.
Let $\beta\in C^2_{\mathrm{c}}$ be supported in a box that embeds in $\Gamma$ (Definition~\ref{def:integrands}),
and let $\beta_\Gamma\in C^2(C(n_i))$, $\beta_\Gamma(\sigma):=\sum_{g\in\Gamma}\beta(\sigma+g)$, be its periodisation.
One has $\langle\beta_\Gamma,F\rangle=\langle\beta,\widetilde F\rangle$ for every $F\in C^2(C(n_i))$, because $\widetilde F$ is $\Gamma$-periodic,
and $\|\beta_\Gamma\|=\|\beta\|$, because the translates of $\supp\beta$ by the elements of $\Gamma$ are pairwise disjoint.
Proposition~\ref{prop:averaged-domination}, applied on $C(n_i)$ to the cochain $\beta_\Gamma$, therefore yields
\[
\mu^{\mathrm{av}}_i\left(\exp\langle\beta,F\rangle\right)\le\exp\left(\frac{1}{2}\|\beta\|^2\right),
\]
that is, $\mu^{\mathrm{av}}_i$ satisfies~\eqref{cond:domination-beta}.
For a box $B$ of radius $r'$ the set $B+[-1,1]^3$ is a cube of side $2r'+2$;
its translates by $n_i\Z^3$ are pairwise disjoint as soon as $n_i>2r'+2$, and then $B$ embeds in $n_i\Z^3$.
Every plaquette $\sigma$ lies in a box of radius $\frac{1}{2}$, and such a box embeds in $n_i\Z^3$ because $n_i\ge4>3$.
Hence every $\mu^{\mathrm{av}}_i$ satisfies~\eqref{cond:domination-beta} for all $\beta=\pm s\cdot1_\sigma$, $\sigma$ a plaquette and $s>0$,
and Lemma~\ref{lem:domination-compactness}(b) shows that the sequence $(\mu^{\mathrm{av}}_i)$ is tight.

Let now $\nu$ be a weak limit point, and let $(i'_m)$ be a subsequence with $\mu^{\mathrm{av}}_{i'_m}\to\nu$ weakly.
Given $\beta\in C^2_{\mathrm{c}}$, choose a box of some radius $r'$ that contains the support of $\beta$;
since $n_i\to\infty$, this box embeds in $n_i\Z^3$ for all large $i$, and $\mu^{\mathrm{av}}_{i'_m}$ satisfies~\eqref{cond:domination-beta} for all large $m$.
Lemma~\ref{lem:domination-compactness}(c) gives~\eqref{cond:domination-beta} for $\nu$;
as $\beta$ was arbitrary, $\nu$ satisfies~\eqref{cond:gaussian-domination}.
Finally, $\nu$ is a law on $\Omega=\{F\in C^2: dF=0\}$, so that $\nu(dF=0)=1$.

\emph{Invariance.}
By Definition~\ref{def:averaged-law}, every $\mu^{\mathrm{av}}_i$ is invariant under $W$ and under $F\mapsto-F$.
For $g\in W$ the map $g^*$ is a homeomorphism of $\Omega$, and so is $F\mapsto-F$.
If $G$ is a bounded continuous function on $\Omega$, then so are $gG$, defined in~\eqref{eq:layers-action}, and $G(-\,\cdot\,)$, and
\[
\begin{aligned}
\nu(gG)&=\lim_{m\to\infty}\mu^{\mathrm{av}}_{i'_m}(gG)=\lim_{m\to\infty}\mu^{\mathrm{av}}_{i'_m}(G)=\nu(G),\\
\nu\left(G(-\,\cdot\,)\right)&=\lim_{m\to\infty}\mu^{\mathrm{av}}_{i'_m}\left(G(-\,\cdot\,)\right)=\lim_{m\to\infty}\mu^{\mathrm{av}}_{i'_m}(G)=\nu(G).
\end{aligned}
\]
A Borel probability measure on the metric space $\Omega$ is determined by its integrals of bounded continuous functions;
therefore $\nu$ is invariant under $W$ and under $F\mapsto-F$.

\emph{The six diagonal families.}
Let $\mu_i=\mu_{C(n_i)}$ be the law of the cubic torus of the given side (Definition~\ref{def:two-sequences}),
put $\lambda_i:=\mu^{\mathrm{av}}_i$,
and read the indicator $\Phi_0$ of the sector of flux zero on $C(n_i)$ as a function of $\widetilde F$.
For every bounded Borel function $G$ on $\Omega$,
\[
\lambda_i(G)=\frac{\mu_i(\Phi_0\cdot G)}{\mathbb{E}_{C(n_i)}[\Phi_0]} ,
\]
so that $\lambda_i\ll\mu_i$, with the density $\Phi_0/\mathbb{E}_{C(n_i)}[\Phi_0]$.
Since $\Phi_0=1$ holds $\lambda_i$-almost everywhere, the relative entropy of Definition~\ref{def:relative-entropy} is
\[
H(\lambda_i\,|\,\mu_i)=\lambda_i\left(\log\frac{\Phi_0}{\mathbb{E}_{C(n_i)}[\Phi_0]}\right)=-\log\mathbb{E}_{C(n_i)}[\Phi_0] .
\]
Lemma~\ref{lem:sector-weights}(c) provides an integer $n_0(q)\ge4$, depending on $q$ only, such that
\[
-\log\mathbb{E}_{C(n_i)}[\Phi_0]\le\log\left(2(2n_i+1)^3\right)
\]
for all $i$ with $n_i\ge n_0(q)$;
since $n_i\to\infty$, this holds for all large $i$,
and for each of the finitely many remaining $i$ the relative entropy is finite, because $\mathbb{E}_{C(n_i)}[\Phi_0]>0$ by Lemma~\ref{lem:sector-weights}(b).
As $\log\left(2(2n_i+1)^3\right)=o(n_i^3)$, it follows that $H(\lambda_i\,|\,\mu_i)=o(n_i^3)$.
In addition, each $\lambda_i$ is invariant under $\tau_v$ for all $v\in\Z^3$ (Definition~\ref{def:averaged-law}).

The laws $\lambda_i$ thus satisfy the hypotheses of Lemma~\ref{lem:diagonal-domination}, and $\nu$ is a weak limit point of $(\lambda_i)$.
Fix $c_0\in\Z$;
the set $\mathcal{R}(c_0)$ of Lemma~\ref{lem:diagonal-laws} consists of the laws on $\Omega$
that satisfy~\eqref{cond:gaussian-domination} and are invariant and reflection positive at $(\ell_-,c_0)$.
By Lemma~\ref{lem:diagonal-domination}, for every $O\in\mathcal{T}$ there is $\nu_0\in\mathcal{R}(c_0)$ with $\nu(O)\le\nu_0(O)$.
The set $\mathcal{R}(c_0)$ is convex and weakly compact by Lemma~\ref{lem:diagonal-laws};
Lemma~\ref{lem:separation} therefore places $\nu$ in $\mathcal{R}(c_0)$,
and in particular $\nu$ is reflection positive at $(\ell_-,c_0)$.

Consider now a mirror $\Pi=(\ell,c)$ of one of the six diagonal families, so that $c\in\Z$,
and let $g\in\Ocub$ be the matrix of step~\ref{it:sides-diagonal} in the proof of Theorem~\ref{thm:cubic-limits},
for which $\ell(gx)=\ell_-(x)$ for all $x\in\R^3$.
Then $\ell(gx)-c=\ell_-(x)-c$ for all $x\in\R^3$.
The law $\nu$ is invariant under $W$, hence under $g^{-1}$,
and it is reflection positive at $(\ell_-,c)$ by the preceding paragraph with $c_0:=c$.
Lemma~\ref{lem:mirror-transport} now shows that $\nu$ is reflection positive at $\Pi$.

\emph{The axis families, for even sides.}
Assume that every $n_i$ is even.
Let $\Pi=(x_\mu,c)$ be a mirror of an axis family, $c\in\frac{1}{2}\Z$, and put $\theta:=\theta_\Pi$.
Let $G\in\mathcal{O}_{\mathrm{cont}}(\{x_\mu>c\})$, a bounded continuous function on $\Omega$ of finitely many $F_\sigma$ with $\sigma\subset\{x_\mu>c\}$,
and choose $h_G>0$ such that all corners of the plaquettes read by $G$ satisfy $c<x_\mu<c+h_G$.
If $n_i>2h_G$, then $G$, evaluated at $\widetilde F$, is a bounded Borel function of the $F_\sigma$,
$\sigma$ a plaquette of $C(n_i)$ all of whose corners satisfy $c<x_\mu<c+n_i/2$,
and Proposition~\ref{prop:averaged-positivity}(a), in the direction $e_\mu$ and with the even side $n_i$, gives
\[
\mu^{\mathrm{av}}_i\left(\overline{\theta G}\cdot G\right)\ge0 .
\]
The function $\overline{\theta G}\cdot G$ is bounded and continuous on $\Omega$, and $n_{i'_m}>2h_G$ for all large $m$; consequently
\[
\nu\left(\overline{\theta G}\cdot G\right)=\lim_{m\to\infty}\mu^{\mathrm{av}}_{i'_m}\left(\overline{\theta G}\cdot G\right)\ge0 .
\]
This holds for every $G\in\mathcal{O}_{\mathrm{cont}}(\{x_\mu>c\})$, and $\nu$ is invariant under $\theta\in W$;
by Lemma~\ref{lem:law-mirror}(b), $\nu$ is reflection positive at $\Pi$.
\end{proof}

\begin{remark}\label{rem:averaged-any-side}
\begin{enumerate}
\item[(1)] Let Lemma~\ref{lem:diagonal-domination} be applied to the laws $\lambda_i=\mu^{\mathrm{av}}_i$ of the proof above.
In step~\ref{it:sides-comparison} of the proof of that lemma the law $\nu_0\in\mathcal{R}(c_0)$ is obtained
as a comparison state (Definition~\ref{def:comparison-states}) for the data $\mathfrak{d}'$ of that step,
and these are data for a companion $\mathfrak{n}'=(n'_i)$ of $\mathfrak{n}$ (Definitions~\ref{def:companions} and~\ref{def:two-sequences}).
Such a comparison state is formed with the all-flux functionals
\[
\mu^t_i(G)=\frac{Z_{D(n'_i)}\left[\exp(t\cdot O_{\mathfrak{d}'})\cdot G\right]}{Z_{D(n'_i)}\left[\exp(t\cdot O_{\mathfrak{d}'})\right]}
\]
on the adapted tori $D(n'_i)=D(n'_i/2,n'_i,n'_i)$,
in which all flux sectors are summed.
\item[(2)] Let $a_1,a_2,a_3\ge4$ be integers, $\mu\in\{1,2,3\}$ and $c\in\frac{1}{2}\Z$.
For the coordinate torus $C(a_1,a_2,a_3)$ and the reflection $\theta_{x_\mu,c}$,
Proposition~\ref{prop:averaged-positivity}(a) has the hypothesis that the side $a_\mu$ in the direction $e_\mu$ of the reflection is even;
in its proof the mirrors $(x_\mu,c)$ and $(x_\mu,c+a_\mu/2)$ are then of the same kind, both site planes or both link planes.
On the cubic torus $C(n_i)$ the side in each of the three directions is $n_i$.
For odd $n_i$ one has $n_i/2\in\Z+\frac{1}{2}$,
so that for every $c\in\frac{1}{2}\Z$ one of the mirrors $(x_\mu,c)$, $(x_\mu,c+n_i/2)$ is a site plane and the other one is a link plane;
the hypothesis of Proposition~\ref{prop:averaged-positivity}(a) then fails in each direction, and that proposition does not apply on $C(n_i)$.
\item[(3)] Lemma~\ref{lem:diagonal-domination} applies in the same way to laws conditioned on translation invariant events.
In the setting of Theorem~\ref{thm:averaged-any-side},
let $\mathcal{B}_i\subset\Omega$ be Borel sets that are invariant under $\tau_v$ for all $v\in\Z^3$
and satisfy $\mu_i(\mathcal{B}_i)\ge\exp\left(-o(n_i^3)\right)$,
and let $\lambda_i:=\mu_i(\,\cdot\,|\,\mathcal{B}_i)$.
Then $\lambda_i\ll\mu_i$ with the density $1_{\mathcal{B}_i}/\mu_i(\mathcal{B}_i)$,
one has $H(\lambda_i\,|\,\mu_i)=-\log\mu_i(\mathcal{B}_i)=o(n_i^3)$,
and $\lambda_i$ is invariant under $\tau_v$, $v\in\Z^3$,
since both $\mathcal{B}_i$ and $\mu_i$ are invariant under these translations
(for $\mu_i$ see step~\ref{it:sides-invariance} of the proof of Theorem~\ref{thm:cubic-limits}).
Consequently, for every weak limit point $\nu$ of $(\lambda_i)$, every $c_0\in\Z$ and every $O\in\mathcal{T}$
there is $\nu_0\in\mathcal{R}(c_0)$ with $\nu(O)\le\nu_0(O)$.
\end{enumerate}
\end{remark}

\section{The continuum limit}\label{sec:continuum}

\subsection{Units and isometries}\label{sec:units}

Throughout this section the parameters are those of Definition~\ref{def:cutoffs}: $0<\rho<1$, a fixed base point $v_0\in C_0$, the standard weights, $e>0$, an odd integer $L\ge3$, $0<\bar m\le1$, and as mass assignment an arbitrary real sequence $N\mapsto m_{0,N}$.

\begin{definition}[Averaging cubes]\label{def:averaging-cubes}
For a real number $\eta>0$ and a plaquette $\sigma=(x;\mu\nu)$ of $\Z^3$, the \emph{centre} of $\sigma$ and the \emph{averaging cube} of $\sigma$ at the spacing $\eta$ are
\[
y_\sigma:=x+\tfrac{1}{2}\left(e_\mu+e_\nu\right),\qquad Q_\eta(\sigma):=\eta\cdot\left(y_\sigma+\left[-\tfrac{1}{2},\tfrac{1}{2}\right)^3\right)\subset\R^3;
\]
the point $\eta\cdot y_\sigma$ is the centre of the plaquette $\eta\sigma$ of $\eta\Z^3$.
For fixed $\mu<\nu$ the centres $y_{(x;\mu\nu)}$, $x\in\Z^3$, form the translate $\frac{1}{2}(e_\mu+e_\nu)+\Z^3$ of $\Z^3$,
and the half-open unit cubes centred at the points of a translate of $\Z^3$ form a partition of $\R^3$;
hence, for every $\eta>0$ and each pair $\mu<\nu$, the cubes $Q_\eta(x;\mu\nu)$, $x\in\Z^3$, form a partition of $\R^3$.

Let $(N,T)$ be a cutoff, with the spacing $\eta=L^{-N}$, the side $n$, the number $R=n\eta$ and the torus $\Lambda_{\eta,R}=\eta\Z^3/R\Z^3$ of Definition~\ref{def:cutoffs}.
The exponent $\lfloor\log_L\bar m^{-1}\rfloor$ is a non-negative integer, because $0<\bar m\le1$;
so $n=8\cdot L^{\lfloor\log_L\bar m^{-1}\rfloor}\cdot L^{T+N}$ is even, $n/2=4\cdot L^{\lfloor\log_L\bar m^{-1}\rfloor}\cdot L^{T+N}$ is an integer $\ge4$, and $R=8\cdot L^{\lfloor\log_L\bar m^{-1}\rfloor}\cdot L^{T}$.
As in Definition~\ref{def:cutoffs}, a two-cochain of $\Lambda_{\eta,R}$ is read as an $n\Z^3$-periodic two-cochain of $\Z^3$, the plaquette $(x;\mu\nu)$ of $\Z^3$ standing for the plaquette $(\eta\cdot x;\mu\nu)$ of $\Lambda_{\eta,R}$,
and two such cochains are paired by
\[
\langle\beta,\beta'\rangle_\eta=\eta^3\sum_\sigma\beta(\sigma)\cdot\beta'(\sigma),\qquad|\beta|_\eta=\langle\beta,\beta\rangle_\eta^{\frac{1}{2}},
\]
the sum being taken over the plaquettes $\sigma$ of $\Lambda_{\eta,R}$, that is, over the plaquettes of the torus $\Z^3/n\Z^3$.
For a two-cochain $\beta$ of $\Lambda_{\eta,R}$, the two-cochain $\beta^{\mathrm u}:=\eta^{\frac{3}{2}}\cdot\beta$ is $\beta$ \emph{in lattice units}.
Since $\eta^3=\eta^{\frac{3}{2}}\cdot\eta^{\frac{3}{2}}$, all two-cochains $\beta$, $\beta'$ of $\Lambda_{\eta,R}$ satisfy
\[
\langle\beta,\beta'\rangle_\eta=\langle\beta^{\mathrm u},\beta'^{\mathrm u}\rangle,\qquad|\beta|_\eta=\|\beta^{\mathrm u}\|,
\]
where $\langle u,u'\rangle=\sum_\sigma u(\sigma)\cdot u'(\sigma)$, with $\sigma$ ranging over the plaquettes of $\Z^3/n\Z^3$, and $\|u\|=\langle u,u\rangle^{\frac{1}{2}}$ denote the pairing and the norm of $C^2(n\Z^3)$ in Definition~\ref{def:configurations}.
The physical curvature of Definition~\ref{def:cutoffs} is $F=\eta^{-\frac{3}{2}}\cdot F^{\mathrm u}$, with $F^{\mathrm u}$ the curvature in lattice units;
thus $F^{\mathrm u}$ arises from $F$ as $\beta^{\mathrm u}$ from $\beta$, and $\langle\beta,F\rangle_\eta=\langle\beta^{\mathrm u},F^{\mathrm u}\rangle$ \cite[Section~21.1.3]{NoncompactQED}.
At the cutoff $(N,T)$ the cube $Q_\eta(\sigma)$ is the cube $Q(\eta\sigma)$ of Definition~\ref{def:cutoffs}, formed with the representative $\eta\cdot x$ of the corner of $\eta\sigma$.

The same formulas define, for every real number $\eta>0$ and every integer $n\ge2$, with $R:=n\eta$, the torus $\Lambda_{\eta,R}=\eta\Z^3/R\Z^3$, the reading of its two-cochains as $n\Z^3$-periodic two-cochains of $\Z^3$, the pairing $\langle\cdot,\cdot\rangle_\eta$, the norm $|\cdot|_\eta$ and the two-cochains $\beta^{\mathrm u}$ in lattice units;
the identities $\langle\beta,\beta'\rangle_\eta=\langle\beta^{\mathrm u},\beta'^{\mathrm u}\rangle$ and $|\beta|_\eta=\|\beta^{\mathrm u}\|$ hold on these tori as well, with $\langle u,u'\rangle$ and $\|u\|$ formed by sums over the plaquettes of $\Z^3/n\Z^3$.
For all $\eta>0$ and $R\ge1$ with $R/\eta\in\Z$, the image of $Q_\eta(\sigma)$ in $\R^3/R\Z^3$ is the averaging cube of the plaquette of $\eta\Z^3/R\Z^3$ represented by $\eta\sigma$, in the sense of \cite[Sections~2.4 and~7.1.3]{NoncompactQED}.
\end{definition}

\begin{definition}[Actions of isometries]\label{def:isometry-actions}
Let $\mathfrak g$ be a Euclidean map, $\mathfrak g(y)=g\,y+a$ with $g\in\mathrm{O}(3)$ and $a\in\R^3$ (Definition~\ref{def:laws}); the matrix $g$ is the \emph{linear part} of $\mathfrak g$.
The value $f(y)$ of a two-form $f=(f_{\mu\nu})_{\mu<\nu}$ on $\R^3$, real or complex, is read as the element $\sum_{\mu<\nu}f_{\mu\nu}(y)\,e_\mu\wedge e_\nu$ of $\wedge^2\C^3$,
and $\wedge^2g$ is the linear map of $\wedge^2\C^3$ with $\left(\wedge^2g\right)(e_\mu\wedge e_\nu)=ge_\mu\wedge ge_\nu$.
With this reading the formula
\[
(\mathfrak g\cdot f)(y)=\left(\wedge^2g\right)f\left(\mathfrak g^{-1}y\right)\qquad(y\in\R^3)
\]
of Definition~\ref{def:laws} defines the two-form $\mathfrak g\cdot f$ for every two-form $f$ on $\R^3$;
on test forms it is the action of $\mathfrak g$ in Definition~\ref{def:laws} \cite[Section~8.1.2]{NoncompactQED}.
On $\mathcal S'$ and on functions on $\mathcal S'$ the Euclidean map $\mathfrak g$ acts as in Definition~\ref{def:laws};
so $\left(\mathfrak g^{-1}F\right)(f)=F(\mathfrak g\cdot f)$ for $F\in\mathcal S'$ and $f\in\mathcal S$, and $\hat{\mathfrak g}$ maps the function $F\mapsto F(f)$ onto the function $F\mapsto F(\mathfrak g\cdot f)$.
For $b\in\R^3$, $g\in\Ocub$ and $\nu\in\{1,2,3\}$, the test forms $\tau_bf$, $U_gf$ and $\mathsf R_\nu f$ of Definition~\ref{def:cutoffs}, where $U_g$ is the map $U_r$ for $r=g$, are $\tau_b\cdot f$, $g\cdot f$ and $\theta_\nu\cdot f$ (Definition~\ref{def:laws}; \cite[Section~2.5]{NoncompactQED}).

\emph{Isometries at a spacing.}
For $g\in W$ (Definition~\ref{def:mirrors}) and a real number $\eta>0$, the \emph{isometry of $g$ at the spacing $\eta$} is the Euclidean map $\mathfrak g:=\eta\,g\,\eta^{-1}$, that is, $\mathfrak g(y)=\eta\cdot g\left(\eta^{-1}y\right)$;
its linear part is the linear part $\vec g\in\Ocub$ of $g$.
A Euclidean map is the isometry at the spacing $\eta$ of an element of $W$ if and only if it maps $\eta\Z^3$ onto itself.
Indeed, every $g\in W$ has the form $g(x)=\vec g\,x+v$ with $v\in\Z^3$, as $W=\Z^3\rtimes\Ocub$ (Definition~\ref{def:mirrors});
its isometry at the spacing $\eta$ is $y\mapsto\vec g\,y+\eta\cdot v$, and it maps $\eta\Z^3$ onto itself, since $\vec g$ and $\vec g^{-1}$ have integer entries.
Conversely, let $\mathfrak g(y)=g\,y+a$ be a Euclidean map with $\mathfrak g(\eta\Z^3)=\eta\Z^3$.
Then $a=\mathfrak g(0)\in\eta\Z^3$, and the matrix $g$ maps $\eta\Z^3$, hence $\Z^3$, onto itself;
thus $g\in\mathrm{O}(3)\cap\mathrm{GL}(3,\Z)=\Ocub$, and $\mathfrak g$ is the isometry at the spacing $\eta$ of the element $x\mapsto g\,x+\eta^{-1}\cdot a$ of $W$.

\emph{Reflections in the planes of the nine families.}
For $\ell\in\mathcal L$ and $c\in\R$, the reflection in the plane $\{\ell=c\}$, which is $r_{\hat n_\ell,c/|n_\ell|}$ by Definition~\ref{def:laws}, is denoted $r_{\ell,c}$;
as $\ell(y)=n_\ell\cdot y$ and $\hat n_\ell=n_\ell/|n_\ell|$,
\[
r_{\ell,c}(y)=y-\frac{2}{|n_\ell|^2}\left(\ell(y)-c\right)\cdot n_\ell\qquad(y\in\R^3).
\]
For the axis family $x_\nu$ and $c=0$ this is the reflection $\theta_\nu$ in the plane $\{x_\nu=0\}$, which as a map of $\R^3$ is the element $\theta_{x_\nu,0}$ of $W$.
Let $\eta>0$.
If $c\in\eta\Adm(\ell)$, then $(\ell,c/\eta)$ is a mirror, and $r_{\ell,c}$ is the isometry at the spacing $\eta$ of the reflection $\theta_{\ell,c/\eta}\in W$ of Definition~\ref{def:mirrors}: by the linearity of $\ell$,
\[
\eta\cdot\theta_{\ell,c/\eta}\left(\eta^{-1}y\right)=y-\frac{2}{|n_\ell|^2}\left(\ell(y)-c\right)\cdot n_\ell=r_{\ell,c}(y)\qquad(y\in\R^3).
\]
If $c\notin\eta\Adm(\ell)$, then $r_{\ell,c}$ is not the isometry at the spacing $\eta$ of any element of $W$, because $r_{\ell,c}(0)=\frac{2c}{|n_\ell|^2}\cdot n_\ell$ does not lie in $\eta\Z^3$:
for an axis family $n_\ell$ is one of the vectors $e_\mu$ and $2c\notin\eta\Z$, while for a diagonal family $n_\ell=e_\mu\pm e_\nu$ with $\mu<\nu$, $|n_\ell|^2=2$ and $c\notin\eta\Z$.

\emph{Orientation signs and pull-backs.}
Let $g\in W$, and let $\sigma$ be a plaquette of $\Z^3$.
The sign in $(g^*u)(\sigma)=\pm u(g\sigma)$ (Definition~\ref{def:cochains}), which compares the image under $g$ of the orientation of $\sigma$ with the orientation of the plaquette $g\sigma$, is the \emph{orientation sign} of $g$ at $\sigma$.
For $g'\in W$, the identity $(gg')^*=g'^*g^*$ of Definition~\ref{def:cochains}, applied to the indicator cochain of the plaquette $gg'\sigma$ and evaluated at $\sigma$, shows that the orientation sign of $gg'$ at $\sigma$ is the product of the orientation sign of $g'$ at $\sigma$ and the orientation sign of $g$ at $g'\sigma$.
With $g^{-1}$ in the place of $g$ and $g$ in the place of $g'$, this product is the orientation sign of the identity map, which is $1$.
Hence $g^{-1}$ has at $g\sigma$ the same orientation sign as $g$ at $\sigma$: for every two-cochain $u$ the identities $\left(\left(g^{-1}\right)^*u\right)(g\sigma)=\pm u(\sigma)$ and $(g^*u)(\sigma)=\pm u(g\sigma)$ hold with the same sign.
Let $u,u'\in C^2$ be such that $\sum_\sigma|u(g\sigma)|\cdot|u'(\sigma)|<\infty$, the sum being taken over the plaquettes $\sigma$ of $\Z^3$.
Then the pairings $\langle g^*u,u'\rangle$ and $\left\langle u,\left(g^{-1}\right)^*u'\right\rangle$ converge absolutely, and both are equal to the sum, over the plaquettes $\sigma$, of the orientation sign of $g$ at $\sigma$ times $u(g\sigma)\cdot u'(\sigma)$;
for the second pairing one substitutes $g\sigma$ for the plaquette of summation, $\sigma\mapsto g\sigma$ being a bijection of the set of the plaquettes of $\Z^3$.
So the transpose of $g^*$ is $\left(g^{-1}\right)^*$, which is the inverse of $g^*$ by the identity $(gg')^*=g'^*g^*$, and $g^*$ is an orthogonal map of $C^2_{\mathrm c}$ onto itself.
If $\theta\in W$ is a reflection, then $\theta^{-1}=\theta$;
so $\theta^*$ is a symmetric involution of $C^2$, that is, $\theta^*\theta^*$ is the identity and $\langle\theta^*u,u'\rangle=\langle u,\theta^*u'\rangle$ whenever $\sum_\sigma|u(\theta\sigma)|\cdot|u'(\sigma)|<\infty$.
In particular $\theta^*$ is a symmetric orthogonal involution of $C^2_{\mathrm c}$.

\emph{Two-cochains of a torus.}
Let $\eta>0$, let $n\ge2$ be an integer and $R:=n\eta$, and let $\Lambda_{\eta,R}$ be the torus of Definition~\ref{def:averaging-cubes} formed with these numbers;
at a cutoff, $\eta$, $n$, $R$ and $\Lambda_{\eta,R}$ are those of Definition~\ref{def:cutoffs}.
Let $g\in W$, with the linear part $\vec g$.
For $m\in\Z^3$ one has $g\circ\tau_{nm}=\tau_{n\vec gm}\circ g$, and $\vec gm\in\Z^3$.
By the identity $(gg')^*=g'^*g^*$, therefore, $\tau_{nm}^*g^*u=g^*\tau_{n\vec gm}^*u=g^*u$ for every $n\Z^3$-periodic two-cochain $u$ of $\Z^3$;
so $g^*$ maps the two-cochains of $\Lambda_{\eta,R}$, read as $n\Z^3$-periodic two-cochains of $\Z^3$, into themselves.
Since $g(\sigma+nm)=g\sigma+n\vec gm$, the map $\sigma\mapsto g\sigma$ induces a permutation of the plaquettes of $\Z^3/n\Z^3$, and the computation of the transpose of $g^*$, with sums over the plaquettes of $\Z^3/n\Z^3$, gives
\[
\langle g^*\beta,\beta'\rangle_\eta=\left\langle\beta,\left(g^{-1}\right)^*\beta'\right\rangle_\eta
\]
for all two-cochains $\beta$, $\beta'$ of $\Lambda_{\eta,R}$.
Hence $g^*$ is orthogonal for $\langle\cdot,\cdot\rangle_\eta$, with the inverse $\left(g^{-1}\right)^*$, and for a reflection $\theta\in W$ the map $\theta^*$ is a symmetric orthogonal involution of the space of the two-cochains of $\Lambda_{\eta,R}$.

Let moreover $R\ge1$.
The maps $v_*$, $v\in\eta\Z^3$, and $r_*$, $r\in\Ocub$, of \cite[Lemma~7.6]{NoncompactQED} and the reflections $\hat{\mathsf R}_\nu$, $\nu\in\{1,2,3\}$, of \cite[Section~7.1.4]{NoncompactQED} are defined on the two-cochains of $\Lambda_{\eta,R}$ by the same formulas for every integer $n=R/\eta$, even or odd, and \cite[Lemma~7.6]{NoncompactQED} holds on $\Lambda_{\eta,R}$ \cite[Section~7.1.3]{NoncompactQED};
in particular they act on the two-cochains of the tori $\Lambda_{\eta,R}$ of the cutoffs, whose side $n$ is even.
Read in lattice units, they are pull-backs by elements of $W$: for $v\in\eta\Z^3$, $g\in\Ocub$ and $\nu\in\{1,2,3\}$,
\[
v_*=\tau_{-v/\eta}^{*},\qquad r_*=\left(g^{-1}\right)^{*}\quad(r=g),\qquad\hat{\mathsf R}_\nu=\theta_{x_\nu,0}^{*}.
\]
Indeed, let $\beta$ be a two-cochain of $\Lambda_{\eta,R}$ and $\sigma=(x;\mu\nu)$ a plaquette of $\Z^3$.
Both $v_*\beta$ and $\tau_{-v/\eta}^{*}\beta$ take the value $\beta(\sigma)$ at the plaquette $\sigma+\eta^{-1}v$, the translation $\tau_{-v/\eta}\in W$ preserving orientations (Definition~\ref{def:cochains}).
The cochain $r_*\beta$, $r=g$, takes at $g\sigma$ the value $\pm\beta(\sigma)$, where the sign compares $\left(\wedge^2g\right)(e_\mu\wedge e_\nu)$, the image under $g$ of the orientation of $\sigma$, with the orientation of $g\sigma$ \cite[Lemma~7.6]{NoncompactQED};
this is the orientation sign of $g$ at $\sigma$, hence that of $g^{-1}$ at $g\sigma$, and so $\left(\left(g^{-1}\right)^*\beta\right)(g\sigma)=\pm\beta(\sigma)$ with the same sign.
Finally, the linear part of $\theta_{x_\nu,0}$ changes the sign of $e_\nu$ and fixes the two other basis vectors;
hence $\theta_{x_\nu,0}$ maps every plaquette onto a plaquette with the same two directions, and its orientation sign is $-1$ at the plaquettes that contain the direction $\nu$ and $1$ at the others.
These are the signs in $\left(\hat{\mathsf R}_\nu\beta\right)(\sigma)=\pm\beta\left(\theta_{x_\nu,0}\,\sigma\right)$ \cite[Section~7.1.4]{NoncompactQED}.
\end{definition}

\begin{lemma}[Orientation signs]\label{lem:orientation-signs}
\begin{enumerate}
\item[(a)] Let $g\in\Ocub$.
Each of the vectors $ge_1$, $ge_2$, $ge_3$ is of the form $\pm e_k$, with different indices $k$ for different vectors.
For $\mu<\nu$ let $\mu'<\nu'$ be the indices such that $ge_\mu$ and $ge_\nu$ are multiples of $e_{\mu'}$ and $e_{\nu'}$, in some order.
Then
\begin{equation}\label{eq:continuum-signs}
\begin{aligned}
\left(\wedge^2g\right)(e_\mu\wedge e_\nu)&=\pm\,e_{\mu'}\wedge e_{\nu'},\\
(\mathfrak g\cdot f)_{\mu'\nu'}&=\pm\,f_{\mu\nu}\circ\mathfrak g^{-1},\\
\left((\tau_vg)^{*}u\right)(\sigma)&=\pm\,u\left((\tau_vg)\,\sigma\right),
\end{aligned}
\end{equation}
with one and the same sign in the three identities, for every Euclidean map $\mathfrak g$ with linear part $g$, every two-form $f$ on $\R^3$, every $v\in\Z^3$, every two-cochain $u$ and every plaquette $\sigma=(x;\mu\nu)$ of $\Z^3$.
This sign is the product of the signs of $ge_\mu$ and $ge_\nu$ if $ge_\mu$ is a multiple of $e_{\mu'}$, and the negative of this product if $ge_\mu$ is a multiple of $e_{\nu'}$.
The element $\tau_vg\colon y\mapsto g\,y+v$ of $W$ maps $\sigma$ onto the plaquette with the directions $\mu'\nu'$ and the centre $g\,y_\sigma+v$.
\item[(b)] Let $\ell\in\mathcal L$, and let $g:=1-2\,\hat n_\ell\hat n_\ell^{T}$, the common linear part of the reflections $\theta_{\ell,c}$, $c\in\Adm(\ell)$, and $r_{\ell,c}$, $c\in\R$.
Then $g\in\Ocub$, and the three vectors $\left(\wedge^2g\right)(e_1\wedge e_2)$, $\left(\wedge^2g\right)(e_1\wedge e_3)$, $\left(\wedge^2g\right)(e_2\wedge e_3)$ are, in this order, those of the row of $\ell$ in the list
\[
\begin{aligned}
&x_1\colon&&-e_1\wedge e_2,&&-e_1\wedge e_3,&&e_2\wedge e_3;\\
&x_2\colon&&-e_1\wedge e_2,&&e_1\wedge e_3,&&-e_2\wedge e_3;\\
&x_3\colon&&e_1\wedge e_2,&&-e_1\wedge e_3,&&-e_2\wedge e_3;\\
&x_1-x_2\colon&&-e_1\wedge e_2,&&e_2\wedge e_3,&&e_1\wedge e_3;\\
&x_1+x_2\colon&&-e_1\wedge e_2,&&-e_2\wedge e_3,&&-e_1\wedge e_3;\\
&x_1-x_3\colon&&-e_2\wedge e_3,&&-e_1\wedge e_3,&&-e_1\wedge e_2;\\
&x_1+x_3\colon&&e_2\wedge e_3,&&-e_1\wedge e_3,&&e_1\wedge e_2;\\
&x_2-x_3\colon&&e_1\wedge e_3,&&e_1\wedge e_2,&&-e_2\wedge e_3;\\
&x_2+x_3\colon&&-e_1\wedge e_3,&&-e_1\wedge e_2,&&-e_2\wedge e_3.
\end{aligned}
\]
If the row of $\ell$ gives $\left(\wedge^2g\right)(e_\mu\wedge e_\nu)=\pm e_{\mu'}\wedge e_{\nu'}$, then for every $c\in\Adm(\ell)$ the reflection $\theta_{\ell,c}$ maps each plaquette $(x;\mu\nu)$ onto a plaquette with the directions $\mu'\nu'$ and has the orientation sign $\pm1$ at $(x;\mu\nu)$, and $(r_{\ell,c}\cdot f)_{\mu'\nu'}=\pm f_{\mu\nu}\circ r_{\ell,c}$ for every $c\in\R$ and every two-form $f$.
\item[(c)] Let $\ast\colon\wedge^2\C^3\to\C^3$ be the linear map with $\ast(e_2\wedge e_3)=e_1$, $\ast(e_1\wedge e_3)=-e_2$, $\ast(e_1\wedge e_2)=e_3$, and for a two-form $f$ let $\ast f:=\ast\circ f=\left(f_{23},-f_{13},f_{12}\right)$ be its dual vector field.
For every Euclidean map $\mathfrak g$ with linear part $g\in\mathrm{O}(3)$ and every two-form $f$ one has, as in \cite[Section~8.9.4]{NoncompactQED},
\[
\ast(\mathfrak g\cdot f)=(\det g)\cdot g\,(\ast f)\circ\mathfrak g^{-1}.
\]
In particular, for a unit vector $n\in\R^3$, a number $c\in\R$ and the reflection $r_{n,c}$ in the plane $\{y\in\R^3:n\cdot y=c\}$ (Definition~\ref{def:laws}),
\[
\ast(r_{n,c}\cdot f)=\left(2\,n\,n^{T}-1\right)(\ast f)\circ r_{n,c};
\]
the component of $\ast f$ along $n$ is kept, and the components orthogonal to $n$ change sign.
\end{enumerate}
\end{lemma}

\begin{proof}
(a) The columns $ge_1$, $ge_2$, $ge_3$ of $g$ are orthonormal vectors with integer entries.
So each of them is $\pm e_k$ for some $k$, and their orthogonality forces different indices for different columns.
Therefore $\left(\wedge^2g\right)(e_\mu\wedge e_\nu)=ge_\mu\wedge ge_\nu$ is the product of the signs of $ge_\mu$ and $ge_\nu$ times $e_{\mu'}\wedge e_{\nu'}$ if $ge_\mu$ is a multiple of $e_{\mu'}$, and times $e_{\nu'}\wedge e_{\mu'}=-e_{\mu'}\wedge e_{\nu'}$ if $ge_\mu$ is a multiple of $e_{\nu'}$.
This is the first identity, with the sign described in the statement.
The map $\mu\nu\mapsto\mu'\nu'$ is a bijection of the set of the three pairs of directions, because $g$ permutes the three lines $\R e_1$, $\R e_2$, $\R e_3$.
Hence
\[
(\mathfrak g\cdot f)(y)=\sum_{\mu<\nu}f_{\mu\nu}\left(\mathfrak g^{-1}y\right)\left(\wedge^2g\right)(e_\mu\wedge e_\nu)=\sum_{\mu<\nu}\pm f_{\mu\nu}\left(\mathfrak g^{-1}y\right)e_{\mu'}\wedge e_{\nu'}
\]
is the expansion of $(\mathfrak g\cdot f)(y)$ in the basis of the vectors $e_{\mu'}\wedge e_{\nu'}$, $\mu'<\nu'$, and the second identity follows.
For the third identity let $\sigma=(x;\mu\nu)$ and $x':=g\,x+v$, a point of $\Z^3$.
The affine map $\tau_vg$ has the linear part $g$.
It maps the corners $x$, $x+e_\mu$, $x+e_\mu+e_\nu$, $x+e_\nu$ of $\sigma$, in this order, to $x'$, $x'+ge_\mu$, $x'+ge_\mu+ge_\nu$, $x'+ge_\nu$.
As $ge_\mu$ and $ge_\nu$ are $\pm e_{\mu'}$ and $\pm e_{\nu'}$ in some order, these four points are the corners of a plaquette with the directions $\mu'\nu'$.
Its centre, the mean of the four corners, is $x'+\frac{1}{2}\left(ge_\mu+ge_\nu\right)=g\,y_\sigma+v$.
By Definition~\ref{def:cochains} the plaquette $\sigma$ is oriented by $e_\mu\wedge e_\nu$, that is, by the ordered pair $(e_\mu,e_\nu)$ of its edge vectors at the corner $x$.
The image of this orientation under $\tau_vg$ is given by the pair $(ge_\mu,ge_\nu)$, that is, by $\left(\wedge^2g\right)(e_\mu\wedge e_\nu)=\pm e_{\mu'}\wedge e_{\nu'}$, whereas the plaquette $(\tau_vg)\sigma$ is oriented by $e_{\mu'}\wedge e_{\nu'}$.
So the orientation sign of $\tau_vg$ at $\sigma$ is the sign of the first identity, and the third identity is the definition of $(\tau_vg)^*$.

(b) The matrix $g$ maps $n_\ell$ to $-n_\ell$ and fixes every vector orthogonal to $n_\ell$.
For the axis family $\ell=x_i$ it therefore changes the sign of $e_i$ and fixes the two other basis vectors.
For a diagonal family one has $|n_\ell|^2=2$ and $g=1-n_\ell n_\ell^{T}$.
If $\ell=x_i-x_j$ with $i<j$, then $n_\ell=e_i-e_j$, $ge_i=e_i-n_\ell=e_j$ and $ge_j=e_j+n_\ell=e_i$.
If $\ell=x_i+x_j$ with $i<j$, then $n_\ell=e_i+e_j$, $ge_i=e_i-n_\ell=-e_j$ and $ge_j=e_j-n_\ell=-e_i$.
In both cases the third basis vector is orthogonal to $n_\ell$ and is fixed.
So in each of the nine cases $g$ permutes the basis vectors up to signs, and $g\in\Ocub$.
The rows of the list are the values $\left(\wedge^2g\right)(e_\mu\wedge e_\nu)=ge_\mu\wedge ge_\nu$, written in the basis $e_1\wedge e_2$, $e_1\wedge e_3$, $e_2\wedge e_3$.
For instance, for $\ell=x_1+x_3$ one has $ge_1=-e_3$, $ge_2=e_2$, $ge_3=-e_1$, and therefore
\[
\begin{aligned}
e_1\wedge e_2&\mapsto-e_3\wedge e_2=e_2\wedge e_3,\\
e_1\wedge e_3&\mapsto e_3\wedge e_1=-e_1\wedge e_3,\\
e_2\wedge e_3&\mapsto-e_2\wedge e_1=e_1\wedge e_2.
\end{aligned}
\]
The reflections $\theta_{\ell,c}$ and $r_{\ell,c}$ have the linear part $1-\frac{2}{|n_\ell|^2}\cdot n_\ell n_\ell^{T}=g$.
For $c\in\Adm(\ell)$ the reflection $\theta_{\ell,c}$ lies in $W$, so that $\theta_{\ell,c}=\tau_vg$ with $v=\theta_{\ell,c}(0)\in\Z^3$; moreover $r_{\ell,c}^{-1}=r_{\ell,c}$.
The last assertion of (b) is therefore (a), applied to $\tau_vg=\theta_{\ell,c}$ and to $\mathfrak g=r_{\ell,c}$.

(c) For $\xi,\xi',\xi''\in\R^3$ the scalar product $(\xi\times\xi')\cdot\xi''$ is the determinant $\det(\xi,\xi',\xi'')$ of the matrix with the columns $\xi$, $\xi'$, $\xi''$.
Both $\ast(\xi\wedge\xi')$ and $\xi\times\xi'$ are bilinear and alternating in $(\xi,\xi')$, and they agree for pairs of basis vectors; hence $\ast(\xi\wedge\xi')=\xi\times\xi'$.
Since $g$ preserves the scalar product,
\[
\begin{aligned}
(g\xi\times g\xi')\cdot g\xi''&=\det(g\xi,g\xi',g\xi'')=(\det g)\cdot\det(\xi,\xi',\xi'')\\
&=(\det g)\cdot(\xi\times\xi')\cdot\xi''=(\det g)\cdot\left(g(\xi\times\xi')\right)\cdot g\xi''.
\end{aligned}
\]
As $g\xi''$ ranges over $\R^3$ when $\xi''$ does, $g\xi\times g\xi'=(\det g)\cdot g(\xi\times\xi')$, that is, $\ast\left(\left(\wedge^2g\right)(\xi\wedge\xi')\right)=(\det g)\cdot g\,\ast(\xi\wedge\xi')$.
Both sides of the last identity are linear in $\xi\wedge\xi'$, and the elements $\xi\wedge\xi'$ span $\wedge^2\C^3$; so the identity holds with $\xi\wedge\xi'$ replaced by any element of $\wedge^2\C^3$.
For the element $f\left(\mathfrak g^{-1}y\right)$, $y\in\R^3$, it is the first formula at the point $y$.
The reflection $r_{n,c}\colon y\mapsto y-2\left(n\cdot y-c\right)\cdot n$ of Definition~\ref{def:laws} has the linear part $1-2\,n\,n^{T}$, whose determinant is $-1$, and $r_{n,c}^{-1}=r_{n,c}$; the first formula becomes the second.
Finally, $\left(2\,n\,n^{T}-1\right)n=n$, and $\left(2\,n\,n^{T}-1\right)\xi=-\xi$ for every $\xi$ orthogonal to $n$.
\end{proof}

By the rows of the axis families in Lemma~\ref{lem:orientation-signs}(b), for $\nu\in\{1,2,3\}$, $c\in\frac{1}{2}\Z$ and every two-cochain $u$ one has $\left(\theta_{x_\nu,c}^*u\right)(\sigma)=\pm u\left(\theta_{x_\nu,c}\,\sigma\right)$, with the sign $-1$ exactly when the plaquette $\sigma$ contains the direction $\nu$.
These are the signs of the electric and of the magnetic variables in \eqref{eq:setting-slice-reflections}, and for $c=0$ they are the signs of the reflection $\hat{\mathsf R}_\nu$ of \cite[Section~7.1.4]{NoncompactQED}.

\subsection{The test map}\label{sec:test-map}

For a real number $\eta>0$ and an integer $n\ge2$, even or odd, put $R:=n\cdot\eta$, and let $\Lambda_{\eta,R}=\eta\Z^3/R\Z^3$ be the torus of Definition~\ref{def:cutoffs}, formed with these numbers.
For a plaquette $\sigma=(x;\mu\nu)$ of $\Z^3$ let $x(\sigma):=x$ be its base point, and recall from Definition~\ref{def:averaging-cubes} its centre $y_\sigma=x+\frac12(e_\mu+e_\nu)$ and its averaging cube $Q_\eta(\sigma)=\eta\cdot\left(y_\sigma+\left[-\frac12,\frac12\right)^3\right)$, a half-open cube of side $\eta$.
For fixed $\mu<\nu$ the cubes $Q_\eta(x;\mu\nu)$, $x\in\Z^3$, are the translates of one of them by the vectors of $\eta\Z^3$; they form a partition of $\R^3$.

Let $[\sigma]$ denote the orbit of the plaquette $\sigma$ under the translations by the vectors of $n\Z^3$.
Two plaquettes $\eta\sigma$ and $\eta\sigma'$ of $\eta\Z^3$ have the same image in $\Lambda_{\eta,R}$ if and only if $[\sigma]=[\sigma']$, and every plaquette of $\Lambda_{\eta,R}$ is such an image;
the plaquette of $\Lambda_{\eta,R}$ obtained from $\eta\sigma$ is again written $[\sigma]$.
Accordingly, a two-cochain $u$ on $\Lambda_{\eta,R}$ is identified with its periodic extension $\sigma\mapsto u([\sigma])$, an $n\Z^3$-periodic two-cochain on $\Z^3$.
Since $Q_\eta(\sigma+n\cdot k)=Q_\eta(\sigma)+R\cdot k$ for $k\in\Z^3$, the cube $Q_\eta(\sigma)$ is determined by $[\sigma]$ up to a translation by a vector of $R\Z^3$;
for fixed $\mu<\nu$ the cubes $Q_\eta(x;\mu\nu)$ with $x\in\{0,\dots,n-1\}^3$ form a partition of a fundamental domain of $R\Z^3$.
Two-cochains $\beta$ and $F$ on $\Lambda_{\eta,R}$ are paired by
\[
\langle\beta,F\rangle_\eta=\eta^3\sum_{[\sigma]}\beta([\sigma])\cdot F([\sigma]),
\]
the sum being taken over the plaquettes of $\Lambda_{\eta,R}$, as in Definition~\ref{def:averaging-cubes};
complex cochains are paired by the same formula, without complex conjugation.
The pull-back $g^*$ of Definition~\ref{def:cochains}, by an isometry $g$ of $\R^3$ with $g\Z^3=\Z^3$, acts on complex cochains by the formula given there, $(g^*u)(\sigma)=\pm u(g\sigma)$ with the same signs; it is complex-linear.
The norm of a complex two-cochain $u$ on $\Z^3$ is given by $\|u\|^2=\sum_\sigma|u(\sigma)|^2$, and that of a complex two-cochain $u$ on $\Lambda_{\eta,R}$ by $|u|_\eta^2=\eta^3\sum_{[\sigma]}|u([\sigma])|^2$.

\begin{definition}[Test cochains]\label{def:test-cochains}
Let $f$ be a test form, real or complex, in the sense of Definition~\ref{def:cutoffs}: an element of $\mathcal S$, or $f=f'+if''$ with $f',f''\in\mathcal S$.
Its coefficients $f_{\mu\nu}$, $\mu<\nu$, are complex-valued Schwartz functions on $\R^3$, and its norm in $L^2(\R^3;\Lambda^2)$, written $\|f\|_2$ as there, is
\[
\|f\|_2=\left(\sum_{\mu<\nu}\int\mathrm{d}y\,|f_{\mu\nu}(y)|^2\right)^{\frac12}.
\]
Let $\eta>0$ and $R>0$ be real numbers.
The \emph{periodisation} of $f$ is the $R\Z^3$-periodic two-form
\[
f^R(y):=\sum_{k\in\Z^3}f(y+R\cdot k)\qquad(y\in\R^3),
\]
the series converging absolutely and uniformly on compact sets.
Put $Q_R:=[-R/2,R/2)^3$, write $\|f^R\|_{L^2(\R^3/R\Z^3)}$ for $\left(\sum_{\mu<\nu}\int_{Q_R}\mathrm{d}y\,|f^R_{\mu\nu}(y)|^2\right)^{\frac12}$, and let
\[
p_R(f):=\|f\|_2+26R^{-2}\cdot\left\|\,|y|^2f\,\right\|_{L^2(\R^3\setminus Q_R)},
\]
where $|y|$ is the Euclidean norm and $|y|^2f$ is the two-form $y\mapsto|y|^2\cdot f(y)$.

If $n:=R/\eta$ is an integer $\ge2$, the \emph{test cochain of $f$ on $\Lambda_{\eta,R}$} is the complex two-cochain $f_\eta$ on $\Lambda_{\eta,R}$ with
\[
f_\eta([\sigma]):=\eta^{-3}\int_{Q_\eta(\sigma)}\mathrm{d}y\,f^R_{\mu\nu}(y)\qquad\left(\sigma=(x;\mu\nu)\text{ a plaquette of }\Z^3\right);
\]
it does not depend on the representative $\sigma$ of $[\sigma]$, because $f^R$ is $R\Z^3$-periodic.
The map $f\mapsto f_\eta$ is the \emph{test map}; $f_\eta$ is real if $f$ is real.
For a two-cochain $F$ on $\Lambda_{\eta,R}$ the \emph{step two-form} $F^\sharp$ has the coefficients
\[
F^\sharp_{\mu\nu}(y):=F([\sigma])\qquad\text{for }y\in Q_\eta(\sigma)+R\Z^3,\ \sigma=(x;\mu\nu);
\]
they are well defined, because for fixed $\mu<\nu$ the cubes $Q_\eta(x;\mu\nu)$, $x\in\Z^3$, form a partition of $\R^3$, and they are bounded and $R\Z^3$-periodic.
The \emph{smeared field} is
\[
F^\sharp(f):=\int\mathrm{d}y\,\sum_{\mu<\nu}f_{\mu\nu}(y)\cdot F^\sharp_{\mu\nu}(y).
\]
For $R\ge1$ the periodisation, the test map and the step two-form are those of \cite[Section~7.1.3]{NoncompactQED}, and $p_R$ is the seminorm of \cite[Lemma~7.5]{NoncompactQED};
there the cube $Q_\eta(\sigma)$ is written $Q(\eta\sigma)$, the averaging cube of the plaquette $\eta\sigma$, and the number $\langle f_\eta,F\rangle_\eta$ is denoted $X_{\eta,R}(f)$.

Finally, the \emph{test cochain of $f$ on $\Z^3$} at the spacing $\eta$ is the complex two-cochain $f_{[\eta]}$ on $\Z^3$, in lattice units, with
\[
f_{[\eta]}(\sigma):=\eta^{-\frac32}\int_{Q_\eta(\sigma)}\mathrm{d}y\,f_{\mu\nu}(y)\qquad\left(\sigma=(x;\mu\nu)\text{ a plaquette of }\Z^3\right).
\]
\end{definition}

\begin{lemma}[Norms of test cochains]\label{lem:test-norms}
Let $\eta>0$ be a real number, $n\ge2$ an integer, even or odd, $R:=n\cdot\eta$, and $f$ a test form, real or complex.
\begin{enumerate}
\item\label{it:test-norms-sum}
One has
\[
\sum_\sigma|f_{[\eta]}(\sigma)|\cdot\left(1+|x(\sigma)|_\infty\right)\le\eta^{-\frac32}\sum_{\mu<\nu}\int\mathrm{d}y\,|f_{\mu\nu}(y)|\cdot\left(2+\frac{|y|_\infty}{\eta}\right)<\infty,
\]
the sum on the left being taken over all plaquettes $\sigma$ of $\Z^3$, and $\|f_{[\eta]}\|\le\|f\|_2$.
\item\label{it:test-norms-periodisation}
The cochain $\eta^{\frac32}\cdot f_\eta$ is the periodisation of $f_{[\eta]}$: for every plaquette $\sigma$ of $\Z^3$
\[
\eta^{\frac32}\cdot f_\eta([\sigma])=\sum_{k\in\Z^3}f_{[\eta]}(\sigma+n\cdot k),
\]
the series converging absolutely.
\item\label{it:test-norms-torus}
One has
\[
|f_\eta|_\eta\le\|f^R\|_{L^2(\R^3/R\Z^3)}\le p_R(f).
\]
If $\supp f\subset y_0+(-R/2,R/2)^3$ for some $y_0\in\R^3$, then $|f_\eta|_\eta\le\|f\|_2$.
If $R\ge1$, then $p_R(f)\le p_1(f)$; and $p_R(f)\to\|f\|_2$ as $R\to\infty$.
The map $p_1$ is a continuous seminorm on $\mathcal S$, and $p_1(h_\kappa)\le31\lambda_\kappa$ for the Hermite two-forms $h_\kappa$ and their eigenvalues $\lambda_\kappa$ of \cite[Section~7.1.5]{NoncompactQED}.
\end{enumerate}
\end{lemma}

\begin{proof}
Let $\sigma=(x;\mu\nu)$ be a plaquette of $\Z^3$ and $y\in Q_\eta(\sigma)$.
Then $y/\eta-x$ lies in $\frac12(e_\mu+e_\nu)+\left[-\frac12,\frac12\right)^3\subset\left[-\frac12,1\right)^3$, so that $|x(\sigma)|_\infty\le|y|_\infty/\eta+1$, that is, $1+|x(\sigma)|_\infty\le2+|y|_\infty/\eta$.
Hence
\[
|f_{[\eta]}(\sigma)|\cdot\left(1+|x(\sigma)|_\infty\right)\le\eta^{-\frac32}\int_{Q_\eta(\sigma)}\mathrm{d}y\,|f_{\mu\nu}(y)|\cdot\left(2+\frac{|y|_\infty}{\eta}\right).
\]
For fixed $\mu<\nu$ the cubes $Q_\eta(x;\mu\nu)$, $x\in\Z^3$, form a partition of $\R^3$; summation over $x$ and over the pairs $\mu<\nu$ gives the inequality of \eqref{it:test-norms-sum}, and its right side is finite because the $f_{\mu\nu}$ are Schwartz functions.
The Cauchy--Schwarz inequality on the cube $Q_\eta(\sigma)$, whose volume is $\eta^3$, gives
\[
|f_{[\eta]}(\sigma)|^2\le\eta^{-3}\cdot\eta^3\int_{Q_\eta(\sigma)}\mathrm{d}y\,|f_{\mu\nu}(y)|^2,
\]
and the same summation yields $\|f_{[\eta]}\|^2\le\|f\|_2^2$.

For \eqref{it:test-norms-periodisation} fix $\sigma=(x;\mu\nu)$.
The cubes $Q_\eta(\sigma+n\cdot k)=Q_\eta(\sigma)+R\cdot k$, $k\in\Z^3$, are pairwise disjoint, so that
\[
\sum_{k\in\Z^3}\int_{Q_\eta(\sigma)}\mathrm{d}y\,|f_{\mu\nu}(y+R\cdot k)|\le\int\mathrm{d}y\,|f_{\mu\nu}(y)|<\infty ,
\]
and summation and integration may be interchanged:
\[
\begin{aligned}
\eta^{\frac32}\cdot f_\eta([\sigma])&=\eta^{-\frac32}\sum_{k\in\Z^3}\int_{Q_\eta(\sigma)}\mathrm{d}y\,f_{\mu\nu}(y+R\cdot k)\\
&=\sum_{k\in\Z^3}\eta^{-\frac32}\int_{Q_\eta(\sigma+n\cdot k)}\mathrm{d}y\,f_{\mu\nu}(y)=\sum_{k\in\Z^3}f_{[\eta]}(\sigma+n\cdot k).
\end{aligned}
\]
The series converges absolutely by \eqref{it:test-norms-sum}.

The argument for \eqref{it:test-norms-torus} follows the proof of \cite[Lemma~7.5]{NoncompactQED}.
By the Cauchy--Schwarz inequality on $Q_\eta(\sigma)$,
\[
|f_\eta([\sigma])|^2\le\eta^{-3}\int_{Q_\eta(\sigma)}\mathrm{d}y\,|f^R_{\mu\nu}(y)|^2 .
\]
Multiply by $\eta^3$ and sum over the plaquettes of $\Lambda_{\eta,R}$, that is, over $x\in\{0,\dots,n-1\}^3$ and over the pairs $\mu<\nu$.
For fixed $\mu<\nu$ the cubes $Q_\eta(x;\mu\nu)$, $x\in\{0,\dots,n-1\}^3$, form a partition of a fundamental domain of $R\Z^3$, and $f^R$ is $R\Z^3$-periodic; hence
\[
|f_\eta|_\eta^2\le\sum_{\mu<\nu}\int_{Q_R}\mathrm{d}y\,|f^R_{\mu\nu}(y)|^2=\|f^R\|^2_{L^2(\R^3/R\Z^3)},
\]
which is the first inequality.

Suppose now that $\supp f\subset y_0+(-R/2,R/2)^3$.
If $y+R\cdot k$ and $y+R\cdot k'$ both lie in $\supp f$, then every coordinate of $R\cdot(k-k')$ is the difference of two points of an open interval of length $R$; thus $|R\cdot(k-k')|_\infty<R$, and $k=k'$.
So at every point $y\in\R^3$ at most one term of the series $f^R(y)$ differs from zero, and
\[
|f^R_{\mu\nu}(y)|^2=\sum_{k\in\Z^3}|f_{\mu\nu}(y+R\cdot k)|^2\qquad(\mu<\nu,\ y\in\R^3).
\]
Integration over $Q_R$ gives
\[
\|f^R\|^2_{L^2(\R^3/R\Z^3)}=\sum_{k\in\Z^3}\sum_{\mu<\nu}\int_{Q_R+R\cdot k}\mathrm{d}y\,|f_{\mu\nu}(y)|^2=\|f\|_2^2,
\]
and the first inequality yields $|f_\eta|_\eta\le\|f\|_2$.

For the second inequality let $a_k:=\|f\|_{L^2(Q_R+R\cdot k)}$ for $k\in\Z^3$; then $a_0\le\|f\|_2$.
On $Q_R$ the two-form $f^R$ is the sum of the two-forms $y\mapsto f(y+R\cdot k)$, whose norms in $L^2(Q_R)$ are the numbers $a_k$, and $Q_R$ is a fundamental domain of $R\Z^3$; by Minkowski's inequality
\[
\|f^R\|_{L^2(\R^3/R\Z^3)}=\|f^R\|_{L^2(Q_R)}\le\sum_{k\in\Z^3}a_k .
\]
Let $k\ne0$ and $y\in Q_R+R\cdot k$.
Then $|y|\ge|y|_\infty\ge R\cdot|k|_\infty-R/2\ge R\cdot|k|_\infty/2$, so that $1\le4R^{-2}\cdot|k|_\infty^{-2}\cdot|y|^2$ on $Q_R+R\cdot k$, and
\[
a_k\le4R^{-2}\cdot|k|_\infty^{-2}\cdot\left\|\,|y|^2f\,\right\|_{L^2(Q_R+R\cdot k)} .
\]
The cubes $Q_R+R\cdot k$, $k\ne0$, form a partition of $\R^3\setminus Q_R$, and the Cauchy--Schwarz inequality for the sum over $k\ne0$ gives
\[
\sum_{k\ne0}a_k\le4R^{-2}\cdot\left(\sum_{k\ne0}|k|_\infty^{-4}\right)^{\frac12}\cdot\left\|\,|y|^2f\,\right\|_{L^2(\R^3\setminus Q_R)} .
\]
For every integer $i\ge1$ one has $\#\left|\{k\in\Z^3:|k|_\infty=i\}\right|=(2i+1)^3-(2i-1)^3=24i^2+2$, whence
\[
\sum_{k\ne0}|k|_\infty^{-4}=\sum_{i\ge1}\left(24i^2+2\right)\cdot i^{-4}=4\pi^2+\frac{\pi^4}{45}<41.65 ,
\]
and $4\cdot41.65^{\frac12}<26$.
Together with $a_0\le\|f\|_2$ this proves $\|f^R\|_{L^2(\R^3/R\Z^3)}\le p_R(f)$.

If $R\ge1$, then $R^{-2}\le1$ and $Q_1\subset Q_R$, so that the second term of $p_R(f)$ is at most the second term of $p_1(f)$.
As $R\to\infty$ the sets $\R^3\setminus Q_R$ decrease to the empty set, and $y\mapsto|y|^4\cdot|f(y)|^2$ is integrable; by dominated convergence $\left\|\,|y|^2f\,\right\|_{L^2(\R^3\setminus Q_R)}\to0$, and therefore $p_R(f)\to\|f\|_2$.
Each of the two terms of $p_1$ is a seminorm on $\mathcal S$, and
\[
p_1(f)\le\|f\|_2+26\left\|\,|y|^2f\,\right\|_2 ,
\]
a sum of two continuous seminorms of $\mathcal S$; so $p_1$ is a continuous seminorm.
The inequality $p_1(h_\kappa)\le31\lambda_\kappa$ contains neither $\eta$ nor $R$; it is the last assertion of \cite[Lemma~7.5]{NoncompactQED}.
\end{proof}

\begin{lemma}[Equivariance of test cochains]\label{lem:test-equivariance}
Let $\eta>0$ be a real number and $g\in W$ (Definition~\ref{def:mirrors}), let $\vec g$ be the linear part of $g$, and let $\mathfrak g$ be the isometry $y\mapsto\eta\cdot g(y/\eta)$ of $g$ at the spacing $\eta$, which acts on two-forms by $(\mathfrak g\cdot f)(y)=\left(\wedge^2\vec g\right)f\left(\mathfrak g^{-1}y\right)$ (Definition~\ref{def:isometry-actions}).
Let $f$ be a test form, real or complex.
\begin{enumerate}
\item\label{it:test-equivariance-lattice}
One has $(\mathfrak g\cdot f)_{[\eta]}=(g^{-1})^*f_{[\eta]}$.
In particular $(r_{\ell,c}\cdot f)_{[\eta]}=\theta_{\ell,c/\eta}^*f_{[\eta]}$ for $\ell\in\mathcal L$ and $c\in\eta\Adm(\ell)$, where $r_{\ell,c}$ is the reflection of $\R^3$ in the plane $\{\ell=c\}$.
\item\label{it:test-equivariance-torus}
Let $n\ge2$ be an integer, even or odd, and $R:=n\cdot\eta$.
On the cubic torus $\Lambda_{\eta,R}$ one has $(\mathfrak g\cdot f)_\eta=(g^{-1})^*f_\eta$.
In particular $(\tau_vf)_\eta=\tau_{-v/\eta}^*f_\eta$ for $v\in\eta\Z^3$, $(U_gf)_\eta=(g^{-1})^*f_\eta$ if $g\in\Ocub$, and $(\mathsf R_\nu f)_\eta=\hat{\mathsf R}_\nu f_\eta$ for $\nu\in\{1,2,3\}$, where $\hat{\mathsf R}_\nu=\theta_{x_\nu,0}^*$.
\end{enumerate}
\end{lemma}

\begin{proof}
Fix a plaquette $\sigma=(x;\mu\nu)$ of $\Z^3$.
By Lemma~\ref{lem:orientation-signs}, $g\sigma$ is a plaquette $(x';\mu'\nu')$ with the centre $g(y_\sigma)$ (the affine map $g$ carries the square $\sigma$ onto a unit square of $\Z^3$ spanned by $\vec ge_\mu$ and $\vec ge_\nu$, and its centre onto the centre), and the sign rule \eqref{eq:continuum-signs}, applied to $\vec g$, gives a sign $\pm$ with
\[
\left(\wedge^2\vec g\right)(e_\mu\wedge e_\nu)=\pm e_{\mu'}\wedge e_{\nu'} .
\]
As $\wedge^2\vec g$ maps each of the vectors $e_1\wedge e_2$, $e_1\wedge e_3$, $e_2\wedge e_3$ to plus or minus one of them, the formula for $\mathfrak g\cdot f$ gives $(\mathfrak g\cdot f)_{\mu'\nu'}=\pm f_{\mu\nu}\circ\mathfrak g^{-1}$.
The matrix $\vec g$ permutes the coordinate axes up to signs, so the closed cube $\left[-\frac12,\frac12\right]^3$ is invariant under $\vec g$, and $\vec g\left[-\frac12,\frac12\right)^3$ differs from $\left[-\frac12,\frac12\right)^3$ by a null set.
Since $\mathfrak g\left(Q_\eta(\sigma)\right)=\eta\cdot\left(g(y_\sigma)+\vec g\left[-\frac12,\frac12\right)^3\right)$, the sets $\mathfrak gQ_\eta(\sigma)$ and $Q_\eta(g\sigma)$ differ by a null set.
The isometry $\mathfrak g$ preserves Lebesgue measure, and the substitution $y=\mathfrak gz$ gives
\[
\begin{aligned}
(\mathfrak g\cdot f)_{[\eta]}(g\sigma)&=\eta^{-\frac32}\int_{Q_\eta(g\sigma)}\mathrm{d}y\,(\mathfrak g\cdot f)_{\mu'\nu'}(y)\\
&=\pm\eta^{-\frac32}\int_{\mathfrak gQ_\eta(\sigma)}\mathrm{d}y\,f_{\mu\nu}\left(\mathfrak g^{-1}y\right)\\
&=\pm\eta^{-\frac32}\int_{Q_\eta(\sigma)}\mathrm{d}z\,f_{\mu\nu}(z)=\pm f_{[\eta]}(\sigma).
\end{aligned}
\]
On the other hand, $\wedge^2\vec g^{-1}$ is the inverse of $\wedge^2\vec g$, so that $\left(\wedge^2\vec g^{-1}\right)(e_{\mu'}\wedge e_{\nu'})=\pm e_\mu\wedge e_\nu$ with the same sign.
The sign rule \eqref{eq:continuum-signs} for $g^{-1}$ at the plaquette $g\sigma$, whose image under $g^{-1}$ is $\sigma$, therefore reads $\left((g^{-1})^*u\right)(g\sigma)=\pm u(\sigma)$ for every two-cochain $u$, real or complex, again with the same sign.
For $u=f_{[\eta]}$ the two identities give $(\mathfrak g\cdot f)_{[\eta]}(g\sigma)=\left((g^{-1})^*f_{[\eta]}\right)(g\sigma)$.
As $\sigma\mapsto g\sigma$ is a bijection of the set of plaquettes of $\Z^3$, the first assertion of \eqref{it:test-equivariance-lattice} follows.
If $\ell\in\mathcal L$ and $c\in\eta\Adm(\ell)$, then $c/\eta\in\Adm(\ell)$, the reflection $\theta:=\theta_{\ell,c/\eta}$ lies in $W$, $\theta^{-1}=\theta$, and its isometry at the spacing $\eta$ is
\[
\eta\cdot\theta(y/\eta)=y-\frac{2}{|n_\ell|^2}\left(\ell(y)-c\right)\cdot n_\ell=r_{\ell,c}(y)\qquad(y\in\R^3);
\]
this is the second assertion.

Turn to \eqref{it:test-equivariance-torus}.
The matrix $\vec g$ maps $n\Z^3$ onto itself, and $g(\sigma+n\cdot m)=g\sigma+n\cdot\vec gm$ for $m\in\Z^3$.
Hence $g$ and $g^{-1}$ map orbits $[\sigma]$ onto orbits; as the sign of \eqref{eq:continuum-signs} depends only on the directions of a plaquette, $(g^{-1})^*$ maps $n\Z^3$-periodic two-cochains, real or complex, to $n\Z^3$-periodic ones, that is, it acts on the two-cochains on $\Lambda_{\eta,R}$.
Let $\pm$ be the sign with $\left((g^{-1})^*u\right)(\sigma')=\pm u(g^{-1}\sigma')$ for every two-cochain $u$, real or complex, and every plaquette $\sigma'$ with the directions of $g\sigma$.
Lemma~\ref{lem:test-norms}\eqref{it:test-norms-periodisation}, applied to $\mathfrak g\cdot f$ and to $f$, the assertion \eqref{it:test-equivariance-lattice}, and the fact that $k\mapsto\vec g^{-1}k$ is a bijection of $\Z^3$ give
\[
\begin{aligned}
\eta^{\frac32}\cdot(\mathfrak g\cdot f)_\eta([g\sigma])&=\sum_{k\in\Z^3}(\mathfrak g\cdot f)_{[\eta]}(g\sigma+n\cdot k)=\sum_{k\in\Z^3}\left((g^{-1})^*f_{[\eta]}\right)\left(g(\sigma+n\cdot\vec g^{-1}k)\right)\\
&=\pm\sum_{m\in\Z^3}f_{[\eta]}(\sigma+n\cdot m)=\pm\eta^{\frac32}\cdot f_\eta([\sigma])=\eta^{\frac32}\cdot\left((g^{-1})^*f_\eta\right)([g\sigma]).
\end{aligned}
\]
Every plaquette of $\Lambda_{\eta,R}$ is of the form $[g\sigma]$, which proves the identity on $\Lambda_{\eta,R}$.
For $v\in\eta\Z^3$ and the translation $g=\tau_{v/\eta}$ the isometry $\mathfrak g$ is the translation by $v$, so that $\mathfrak g\cdot f=\tau_vf$ and $g^{-1}=\tau_{-v/\eta}$.
For $g\in\Ocub$ one has $\mathfrak g=g$ and $\mathfrak g\cdot f=U_gf$.
For $g=\theta_{x_\nu,0}$ one has $\mathfrak g=\theta_\nu$, $\mathfrak g\cdot f=\mathsf R_\nu f$ and $(g^{-1})^*=\theta_{x_\nu,0}^*=\hat{\mathsf R}_\nu$.
For $R\ge1$ these three identities are those of \cite[Lemma~7.6]{NoncompactQED}, where $\tau_{-v/\eta}^*$ and $(g^{-1})^*$, $g\in\Ocub$, are written $v_*$ and $r_*$.
\end{proof}

\begin{lemma}[Bianchi identity for test cochains]\label{lem:test-bianchi}
Let $\eta>0$ be a real number and $\xi$ a Schwartz function on $\R^3$, real or complex.
For $j\in\{1,2,3\}$ let $\partial^\eta_j\xi$ be the centred difference quotient
\[
\partial^\eta_j\xi(y):=\eta^{-1}\left(\xi\left(y+\frac{\eta}{2}\cdot e_j\right)-\xi\left(y-\frac{\eta}{2}\cdot e_j\right)\right)\qquad(y\in\R^3),
\]
and let $u^{(\eta)}$ and $d^*\xi$ be the test forms, real if $\xi$ is real, whose components, listed in the order $23,13,12$, are
\[
\begin{aligned}
\left(u^{(\eta)}_{23},u^{(\eta)}_{13},u^{(\eta)}_{12}\right)&=\left(-\partial^\eta_1\xi,\ \partial^\eta_2\xi,\ -\partial^\eta_3\xi\right),\\
\left((d^*\xi)_{23},(d^*\xi)_{13},(d^*\xi)_{12}\right)&=\left(-\partial_1\xi,\ \partial_2\xi,\ -\partial_3\xi\right).
\end{aligned}
\]
\begin{enumerate}
\item\label{it:test-bianchi-closed}
For every closed two-cochain $u$ on $\Z^3$ with
\[
\sup_\sigma\frac{|u(\sigma)|}{1+|x(\sigma)|_\infty}<\infty
\]
the series $\sum_\sigma u^{(\eta)}_{[\eta]}(\sigma)\cdot u(\sigma)$ over the plaquettes $\sigma$ of $\Z^3$ converges absolutely, and its sum is $0$.
\item\label{it:test-bianchi-estimate}
One has
\[
\left\|u^{(\eta)}-d^*\xi\right\|_2\le\frac{\eta^2}{24}\sum_{j=1}^3\left\|\partial_j^3\xi\right\|_2 ,
\]
the norms on the right being those of $L^2(\R^3)$.
\end{enumerate}
\end{lemma}

\begin{proof}
For \eqref{it:test-bianchi-closed}, attach to every cube $x+[0,1]^3$ of $\Z^3$, $x\in\Z^3$, the number
\[
\xi_{[\eta]}(x):=\eta^{-\frac52}\int_{\eta\cdot\left(x+[0,1)^3\right)}\mathrm{d}y\,\xi(y),
\]
an integral over the cube $\eta\cdot\left(x+\left(\frac12,\frac12,\frac12\right)+\left[-\frac12,\frac12\right)^3\right)=\eta\cdot\left(x+[0,1)^3\right)$.
The plaquette $(x;23)$ has the centre $x+\left(0,\frac12,\frac12\right)$; translated by $\frac{\eta}{2}\cdot e_1$, its averaging cube $Q_\eta(x;23)$ becomes $\eta\cdot\left(x+[0,1)^3\right)$, and translated by $-\frac{\eta}{2}\cdot e_1$ it becomes $\eta\cdot\left(x-e_1+[0,1)^3\right)$.
Therefore
\[
\begin{aligned}
u^{(\eta)}_{[\eta]}(x;23)&=-\eta^{-\frac52}\int_{Q_\eta(x;23)}\mathrm{d}y\,\left(\xi\left(y+\frac{\eta}{2}\cdot e_1\right)-\xi\left(y-\frac{\eta}{2}\cdot e_1\right)\right)\\
&=-\left(\xi_{[\eta]}(x)-\xi_{[\eta]}(x-e_1)\right),
\end{aligned}
\]
and in the same way
\[
u^{(\eta)}_{[\eta]}(x;13)=\xi_{[\eta]}(x)-\xi_{[\eta]}(x-e_2),\qquad u^{(\eta)}_{[\eta]}(x;12)=-\left(\xi_{[\eta]}(x)-\xi_{[\eta]}(x-e_3)\right).
\]
Write $u_{\mu\nu}(x):=u(x;\mu\nu)$.
By Definition~\ref{def:cochains} the coboundary of $u$ at the cube $x+[0,1]^3$ is
\[
(du)(x)=\left(u_{23}(x+e_1)-u_{23}(x)\right)-\left(u_{13}(x+e_2)-u_{13}(x)\right)+\left(u_{12}(x+e_3)-u_{12}(x)\right),
\]
so that the three numbers just computed are the coefficients of $u_{23}(x)$, $u_{13}(x)$ and $u_{12}(x)$ in $\sum_{x\in\Z^3}\xi_{[\eta]}(x)\cdot(du)(x)$.
All series converge absolutely.
Indeed, $|x|_\infty\le|y|_\infty/\eta+1$ for $y\in\eta\cdot\left(x+[0,1)^3\right)$, whence
\[
\sum_{x\in\Z^3}|\xi_{[\eta]}(x)|\cdot\left(2+|x|_\infty\right)\le\eta^{-\frac52}\int\mathrm{d}y\,|\xi(y)|\cdot\left(3+\frac{|y|_\infty}{\eta}\right)<\infty ;
\]
the six plaquettes that occur in $(du)(x)$ have the base points $x$, $x+e_1$, $x+e_2$, $x+e_3$, at which $|u|$ is at most $C\cdot\left(2+|x|_\infty\right)$, with $C$ the supremum in the hypothesis.
The series $\sum_\sigma u^{(\eta)}_{[\eta]}(\sigma)\cdot u(\sigma)$ converges absolutely by Lemma~\ref{lem:test-norms}\eqref{it:test-norms-sum}, applied to the test form $u^{(\eta)}$.
The absolutely convergent double series may be regrouped plaquette by plaquette, and as $u$ is closed,
\[
\sum_\sigma u^{(\eta)}_{[\eta]}(\sigma)\cdot u(\sigma)=\sum_{x\in\Z^3}\xi_{[\eta]}(x)\cdot(du)(x)=0 .
\]

The two-forms $u^{(\eta)}$ and $d^*\xi$ are those of \cite[Lemma~7.8]{NoncompactQED}, where $\partial^\eta_j$ is written $\delta_{\eta,j}$ and the function $\xi$ is written $g$.
That lemma contains the inequality of \eqref{it:test-bianchi-estimate}, in which neither $R$ nor $n$ occurs, with the proof that follows, and, in place of \eqref{it:test-bianchi-closed}, the corresponding identity for the closed two-cochains on the tori $\Lambda_{\eta,R}$ treated there.
To prove \eqref{it:test-bianchi-estimate}, write $\hat\xi(p)=\int\mathrm{d}y\,\exp(-ip\cdot y)\,\xi(y)$ for the Fourier transform of $\xi$, and use the same notation for the Schwartz functions $\partial^\eta_j\xi$, $\partial_j\xi$ and $\partial_j^3\xi$.
Then
\[
\widehat{\partial^\eta_j\xi}(p)=\frac{2i}{\eta}\sin\left(\frac{\eta\cdot p_j}{2}\right)\cdot\hat\xi(p),\qquad\widehat{\partial_j\xi}(p)=ip_j\cdot\hat\xi(p),\qquad\widehat{\partial_j^3\xi}(p)=-ip_j^3\cdot\hat\xi(p),
\]
so that $\partial^\eta_j-\partial_j$ acts as multiplication by $i\left(\frac{2}{\eta}\sin\left(\frac{\eta\cdot p_j}{2}\right)-p_j\right)$.
Since $|\sin s-s|\le\frac{|s|^3}{6}$ for real $s$, this multiplier has modulus at most $\frac{\eta^2}{24}\cdot|p_j|^3$, and Plancherel's theorem gives $\left\|(\partial^\eta_j-\partial_j)\xi\right\|_2\le\frac{\eta^2}{24}\left\|\partial_j^3\xi\right\|_2$.
The components of $u^{(\eta)}-d^*\xi$ are the functions $\pm(\partial^\eta_j-\partial_j)\xi$, $j=1,2,3$; hence
\[
\left\|u^{(\eta)}-d^*\xi\right\|_2=\left(\sum_{j=1}^3\left\|(\partial^\eta_j-\partial_j)\xi\right\|_2^2\right)^{\frac12}\le\sum_{j=1}^3\left\|(\partial^\eta_j-\partial_j)\xi\right\|_2\le\frac{\eta^2}{24}\sum_{j=1}^3\left\|\partial_j^3\xi\right\|_2 .
\]
\end{proof}

\begin{lemma}[Tests on the positive side]\label{lem:positive-tests}
Let $\eta>0$ be a real number, $n\ge2$ an even integer, $R:=n\cdot\eta$ and $\nu\in\{1,2,3\}$, and let $\Sigma^{++}_\nu$ be the set of the plaquettes of $\Lambda_{\eta,R}$ all of whose vertices have their $\nu$-th coordinate in $\{\eta,2\eta,\dots,(n/2-1)\cdot\eta\}$ modulo $R$.
Let $\delta>0$ and $r$ be real numbers, and let $f$ be a test form, real or complex, with
\[
\supp f\subset\{y\in\R^3:y_\nu\ge\delta\}\cap[-r,r]^3 .
\]
If $\eta\le\delta$ and $R\ge2r+3\eta$, in particular if $2\eta\le\delta$ and $R\ge2r+4\eta$, then $f_\eta$ vanishes outside $\Sigma^{++}_\nu$.
\end{lemma}

\begin{proof}
The conditions $2\eta\le\delta$ and $R\ge2r+4\eta$ imply $\eta\le\delta$ and $R\ge2r+3\eta$; assume the latter.
Since $Q_R$ is a fundamental domain of $R\Z^3$, every plaquette of $\Lambda_{\eta,R}$ is $[\sigma]$ for exactly one plaquette $\sigma$ of $\Z^3$ with $\eta\cdot y_\sigma\in Q_R$; fix this representative, and let $x:=x(\sigma)$ be its base point.
Every point $y$ of $Q_\eta(\sigma)=\eta\cdot y_\sigma+\left[-\frac{\eta}{2},\frac{\eta}{2}\right)^3$ satisfies $|y_j|\le\frac{R+\eta}{2}$ for $j=1,2,3$.
If $k\in\Z^3$ and $k_j\ne0$ for some $j$, then
\[
|y_j+R\cdot k_j|\ge R-|y_j|\ge\frac{R-\eta}{2}\ge r+\eta>r,
\]
so that $y+R\cdot k\notin\supp f$.
Hence $f^R=f$ on $Q_\eta(\sigma)$, and $f_\eta([\sigma])\ne0$ implies that $Q_\eta(\sigma)$ contains a point $y$ of $\supp f$.
For such a point $\delta\le y_\nu\le r$ and $\eta\cdot(y_\sigma)_\nu-\frac{\eta}{2}\le y_\nu<\eta\cdot(y_\sigma)_\nu+\frac{\eta}{2}$, whence
\[
\delta-\frac{\eta}{2}<\eta\cdot(y_\sigma)_\nu\le r+\frac{\eta}{2}.
\]
Suppose first that $\nu$ is not one of the two directions of $\sigma$.
Then $(y_\sigma)_\nu=x_\nu$ is an integer; $\eta\cdot x_\nu>\delta-\frac{\eta}{2}\ge\frac{\eta}{2}$ gives $x_\nu\ge1$, and $\eta\cdot x_\nu\le r+\frac{\eta}{2}\le\frac{R}{2}-\eta$ gives $x_\nu\le\frac{n}{2}-1$.
All four vertices of $\eta\sigma$ have the $\nu$-th coordinate $\eta\cdot x_\nu$.
Suppose next that $\nu$ is one of the two directions of $\sigma$.
Then $(y_\sigma)_\nu=x_\nu+\frac12$; $\eta\cdot x_\nu>\delta-\eta\ge0$ gives $x_\nu\ge1$, and $\eta\cdot(x_\nu+1)\le r+\eta\le\frac{R}{2}-\frac{\eta}{2}$ gives $x_\nu+1\le\frac{n}{2}-\frac12$.
As $n$ is even, $x_\nu+1$ and $n/2$ are integers, so that $x_\nu+1\le\frac{n}{2}-1$.
The vertices of $\eta\sigma$ have the $\nu$-th coordinates $\eta\cdot x_\nu$ and $\eta\cdot(x_\nu+1)$.
In both cases $[\sigma]\in\Sigma^{++}_\nu$.
\end{proof}

\begin{lemma}[Unfolding]\label{lem:test-unfolding}
Let $\eta>0$ be a real number, $n\ge2$ an integer and $R:=n\cdot\eta$.
Let $F^{\mathrm u}$ be a two-cochain on $\Z^3/n\Z^3$ with the periodic extension $\widetilde F^{\mathrm u}$ to $\Z^3$, and let $F:=\eta^{-\frac32}\cdot F^{\mathrm u}$, a two-cochain on $\Lambda_{\eta,R}$.
Then for every test form $f$, real or complex,
\[
F^\sharp(f)=\langle f_\eta,F\rangle_\eta=\sum_\sigma f_{[\eta]}(\sigma)\cdot\widetilde F^{\mathrm u}(\sigma),
\]
the sum being taken over all plaquettes $\sigma$ of $\Z^3$ and converging absolutely.
\end{lemma}

\begin{proof}
By the definition of the pairing,
\[
\langle f_\eta,F\rangle_\eta=\eta^3\sum_{[\sigma]}f_\eta([\sigma])\cdot\eta^{-\frac32}\cdot F^{\mathrm u}([\sigma])=\sum_{[\sigma]}\eta^{\frac32}\cdot f_\eta([\sigma])\cdot F^{\mathrm u}([\sigma]).
\]
Choose one representative $\sigma$ in each orbit; the plaquettes $\sigma+n\cdot k$, with $\sigma$ a representative and $k\in\Z^3$, are then all the plaquettes of $\Z^3$, each of them once.
Insert Lemma~\ref{lem:test-norms}\eqref{it:test-norms-periodisation} and use $\widetilde F^{\mathrm u}(\sigma+n\cdot k)=F^{\mathrm u}([\sigma])$:
\[
\langle f_\eta,F\rangle_\eta=\sum_{[\sigma]}\sum_{k\in\Z^3}f_{[\eta]}(\sigma+n\cdot k)\cdot\widetilde F^{\mathrm u}(\sigma+n\cdot k)=\sum_\sigma f_{[\eta]}(\sigma)\cdot\widetilde F^{\mathrm u}(\sigma).
\]
The regrouping is allowed: $\widetilde F^{\mathrm u}$ takes finitely many values, so that $\sum_\sigma|f_{[\eta]}(\sigma)|\cdot|\widetilde F^{\mathrm u}(\sigma)|\le\max|F^{\mathrm u}|\cdot\sum_\sigma|f_{[\eta]}(\sigma)|<\infty$ by Lemma~\ref{lem:test-norms}\eqref{it:test-norms-sum}.

For $R\ge1$ the identity $F^\sharp(f)=\langle f_\eta,F\rangle_\eta$ is the display (7.3) of \cite[Section~7.1.3]{NoncompactQED}; for every $R$ it follows from the definitions.
Fix $\mu<\nu$.
Since $f_{\mu\nu}$ is integrable and $F^\sharp_{\mu\nu}$ is bounded and $R\Z^3$-periodic, Fubini's theorem gives
\[
\int\mathrm{d}y\,f_{\mu\nu}(y)\cdot F^\sharp_{\mu\nu}(y)=\sum_{k\in\Z^3}\int_{Q_R}\mathrm{d}y\,f_{\mu\nu}(y+R\cdot k)\cdot F^\sharp_{\mu\nu}(y)=\int_{Q_R}\mathrm{d}y\,f^R_{\mu\nu}(y)\cdot F^\sharp_{\mu\nu}(y).
\]
The function $f^R_{\mu\nu}\cdot F^\sharp_{\mu\nu}$ is $R\Z^3$-periodic, the cubes $Q_\eta(x;\mu\nu)$, $x\in\{0,\dots,n-1\}^3$, form a partition of a fundamental domain of $R\Z^3$, and $F^\sharp_{\mu\nu}=F\left([(x;\mu\nu)]\right)$ on $Q_\eta(x;\mu\nu)$.
The last integral is therefore
\[
\sum_{x\in\{0,\dots,n-1\}^3}F\left([(x;\mu\nu)]\right)\cdot\int_{Q_\eta(x;\mu\nu)}\mathrm{d}y\,f^R_{\mu\nu}(y)=\eta^3\sum_{x\in\{0,\dots,n-1\}^3}f_\eta\left([(x;\mu\nu)]\right)\cdot F\left([(x;\mu\nu)]\right),
\]
and summation over the pairs $\mu<\nu$ gives $F^\sharp(f)=\langle f_\eta,F\rangle_\eta$.
\end{proof}

\subsection{The lattice law at a cutoff}\label{sec:cutoff-law}

For a cutoff $(N,T)$, with the spacing $\eta$, the side $n$ and the torus $\Lambda_{\eta,R}$ of Definition~\ref{def:cutoffs}, let $C^2(\Lambda_{\eta,R})$ denote the real vector space of the two-cochains of $\Lambda_{\eta,R}$.
Under the identification made in Definition~\ref{def:cutoffs} it is the space $C^2(n\Z^3)$ of Definition~\ref{def:configurations}, whose coboundary is denoted $d$;
the pairing $\langle\cdot,\cdot\rangle_\eta$ and the norm $|\cdot|_\eta$ on $C^2(\Lambda_{\eta,R})$ are those of Definition~\ref{def:cutoffs}.

\begin{proposition}[Symmetries, domination and positivity]\label{prop:cutoff-law}
Let $L\ge3$ be an odd integer, let $0<\bar m\le1$, $e>0$, $\rho\in(0,1)$ and $v_0\in C_0$, let the weights $w_0,\dots,w_3$ be standard (Definition~\ref{def:chamber-operator}), and let $(m_{0,N})_{N\ge0}$ be a mass assignment (Definition~\ref{def:cutoffs}).
Let $N,T\ge0$ be integers, and let
\[
\eta=L^{-N},\qquad n=8\cdot L^{\lfloor\log_L\bar m^{-1}\rfloor}\cdot L^{T+N},\qquad R=n\cdot\eta,\qquad\Lambda_{\eta,R}=\eta\Z^3/R\Z^3
\]
be the data of the cutoff $(N,T)$ (Definition~\ref{def:cutoffs}).
On the cubic torus $C(n)=C(n,n,n)$ let $\mathbb{E}_{C(n)}$ be the expectation of Definition~\ref{def:functional}, formed with the lattice charge $q=e\cdot\eta^{\frac12}$, the Wilson parameter $\rho$, the point $v_0$, the standard weights and the mass parameter $M=m_N+\frac{\rho}{2}\cdot\sum_{i=0}^{3}w_i$, where $m_N=\bar m\cdot\eta+m_{0,N}$ is the lattice mass.
Let $\mathbb{P}_{N,T}$ be the law under $\mathbb{E}_{C(n)}$ of the physical curvature $F=\eta^{-\frac32}\cdot F^{\mathrm u}\in C^2(\Lambda_{\eta,R})$, where $F^{\mathrm u}=dA+(2\pi/q)\cdot z$ is the curvature of the configuration $(A,z)\in\Conf(n\Z^3)$.
For a coordinate direction $\nu_0\in\{1,2,3\}$ let $\Sigma^{++}_{\nu_0}$ be the set of the plaquettes of $\Lambda_{\eta,R}$ all of whose vertices have their $\nu_0$-th coordinate in $\{\eta,2\eta,\dots,(n/2-1)\cdot\eta\}+R\Z$, as in Lemma~\ref{lem:positive-tests},
and let $\hat{\mathsf{R}}_{\nu_0}:=\theta_{x_{\nu_0},0}^{*}$ be the pull-back of two-cochains by the reflection $\theta_{x_{\nu_0},0}\in W$ in the plane $\{x_{\nu_0}=0\}$.
Then the following hold.
\begin{enumerate}
\item\label{it:cutoff-closed}
$\mathbb{P}_{N,T}$ is a Borel probability measure on $C^2(\Lambda_{\eta,R})$, and $\mathbb{P}_{N,T}(\ker d)=1$.
\item\label{it:cutoff-invariance}
$\mathbb{P}_{N,T}$ is invariant under $\tau_{-v/\eta}^{*}$ for every $v\in\eta\Z^3$ and under $(g^{-1})^{*}$ for every $g\in\Ocub$.
\item\label{it:cutoff-domination}
For every real $\beta\in C^2(\Lambda_{\eta,R})$,
\[
\int\mathrm{d}\mathbb{P}_{N,T}(F)\,\exp\left(\langle\beta,F\rangle_\eta\right)\le\exp\left(\frac12|\beta|_\eta^2\right).
\]
\item\label{it:cutoff-even}
$\mathbb{P}_{N,T}$ is invariant under $F\mapsto-F$.
\item\label{it:cutoff-positivity}
For $\nu_0=1,2,3$, for every finite family $(\beta_p)$ of real two-cochains of $\Lambda_{\eta,R}$ supported on $\Sigma^{++}_{\nu_0}$ and for every family $(a_p)$ of complex numbers with the same indices,
\begin{equation}\tag{Ad3$'$}\label{eq:continuum-admissible-even}
\sum_{p,p'}\overline{a_p}\cdot a_{p'}\cdot\int\mathrm{d}\mathbb{P}_{N,T}(F)\,\exp\left(i\left\langle\beta_{p'}-\hat{\mathsf{R}}_{\nu_0}\beta_p,F\right\rangle_\eta\right)\ge0.
\end{equation}
\item\label{it:cutoff-chessboard}
Let $\mathcal{Q}=(\mathfrak{F}_{\mathrm{c}};o_1,o_2,o_3;h)$ be a cell structure of the coordinate frame (Definition~\ref{def:cell-structures}; thus $o_1,o_2,o_3\in\frac12\Z$, and the width $h$ is a positive integer), and suppose that $n\in2h\Z$.
Then $C(n)$ carries $\mathcal{Q}$, and the number of its cells is $(n/h)^3$.
Fix a set of representatives of the cells of $C(n)$, and let $G_\kappa\in\mathcal{A}(Q_\kappa)$ for each $\kappa$ in this set.
With the folds of Definition~\ref{def:folds}, and with $\kappa$ ranging over the same set, one has
\[
\mathbb{E}_{C(n)}\left[\Fold_{C(n)}G_\kappa\right]\ge0
\qquad\text{and}\qquad
\left|\mathbb{E}_{C(n)}\left[\prod_\kappa G_\kappa\right]\right|\le\prod_\kappa\mathbb{E}_{C(n)}\left[\Fold_{C(n)}G_\kappa\right]^{\left(\frac{h}{n}\right)^3}.
\]
If $G_\kappa$ is bosonic, its mirror copy $G_\kappa^{[\kappa']}$ in the cell $Q_{\kappa'}$ is the function $F\mapsto G_\kappa(g^{*}F)$ when the element $g\in W(\mathcal{Q})$ with $gQ_\kappa=Q_{\kappa'}$ is even, and the function $F\mapsto\overline{G_\kappa(g^{*}F)}$ when $g$ is odd.
\end{enumerate}
\end{proposition}

\begin{proof}
The lattice $n\Z^3$ has systole $n\ge4$.
It is therefore a period lattice (Definition~\ref{def:configurations}), and $C(n)=(n\Z^3,\chi)$ is a torus with $\chi(n\cdot e_\mu)=-1$ for $\mu=1,2,3$ (Definition~\ref{def:tori}).
By Theorem~\ref{thm:partition-positivity}, $0<Z_{C(n)}[1]<\infty$.
Let $\Phi$ be a bounded Borel function on $C^2(n\Z^3)$.
For the bosonic integrand $\Phi(F^{\mathrm u})$ the integral of Definition~\ref{def:functional} reads
\[
Z_{C(n)}[\Phi]=\sum_{z}\int_{[0,2\pi/q)^{C(n)^{(1)}}}\mathrm{d}A\,\exp\left(-\frac12\|F^{\mathrm u}\|^2\right)\cdot\det D^{C(n)}_A\cdot\Phi(F^{\mathrm u}),
\]
the sum being taken over the closed $z\in C^2(n\Z^3;\Z)$ (Definition~\ref{def:configurations}).
It converges absolutely by Proposition~\ref{prop:berezin}\eqref{it:torus-integrability}, and its weight is non-negative, since $\det D^{C(n)}_A\ge0$ for real $A$ by Proposition~\ref{prop:accretivity}\eqref{it:determinant-sign}.
Hence the weight, divided by $Z_{C(n)}[1]$, is the density, with respect to the measure of Definition~\ref{def:configurations}, of a probability measure on $\Conf(n\Z^3)$, and $\Phi\mapsto\mathbb{E}_{C(n)}[\Phi]$ is the integral against the image of this probability measure under the Borel map $(A,z)\mapsto F^{\mathrm u}$.
The image of the same probability measure under $(A,z)\mapsto\eta^{-\frac32}\cdot F^{\mathrm u}$ is $\mathbb{P}_{N,T}$, which is thus a Borel probability measure on $C^2(\Lambda_{\eta,R})$.

\emph{Clause~\eqref{it:cutoff-closed}.}
For every configuration, $F^{\mathrm u}=dA+(2\pi/q)\cdot z$ with $dz=0$ (Definition~\ref{def:configurations}).
As $d\circ d=0$, it follows that $dF^{\mathrm u}=0$ and $dF=\eta^{-\frac32}\cdot dF^{\mathrm u}=0$; so $\mathbb{P}_{N,T}$ is carried by $\ker d$.

\emph{Clause~\eqref{it:cutoff-invariance}.}
Let $g\in W$.
By Proposition~\ref{prop:symmetries}\eqref{it:functional-covariance}, the law $\mu_{C(n)}$ of $\widetilde F^{\mathrm u}$ is invariant under $g^{*}$: the cochains $\widetilde F^{\mathrm u}$ and $g^{*}\widetilde F^{\mathrm u}$ have the same law under $\mathbb{E}_{C(n)}$.
The space $C^2(n\Z^3)$ is invariant under $g^{*}$ (Section~\ref{sec:units}); taking the periodic extension therefore commutes with $g^{*}$, and $g^{*}\widetilde F^{\mathrm u}$ is the periodic extension of $g^{*}F^{\mathrm u}$.
An element of $C^2(n\Z^3)$ is recovered from its periodic extension by restriction to the plaquettes of a fundamental domain of $n\Z^3$, and this restriction is a continuous map on $\Omega$.
Consequently $F^{\mathrm u}$ and $g^{*}F^{\mathrm u}$ have the same law, and $\mathbb{P}_{N,T}$, the law of $\eta^{-\frac32}\cdot F^{\mathrm u}$, is invariant under $g^{*}$.
For $v\in\eta\Z^3$ and $g\in\Ocub$, the translation $\tau_{-v/\eta}$ and the element $g^{-1}$ belong to $W=\Z^3\rtimes\Ocub$ (Definition~\ref{def:mirrors}), because $v/\eta\in\Z^3$; so the maps of clause~\eqref{it:cutoff-invariance} are pull-backs by elements of $W$.

\emph{Clause~\eqref{it:cutoff-domination}.}
Let $\beta\in C^2(\Lambda_{\eta,R})$ be real, and let $\beta^{\mathrm u}=\eta^{\frac32}\cdot\beta$ be $\beta$ in lattice units (Definition~\ref{def:averaging-cubes}).
By the identities $\langle\beta,F\rangle_\eta=\langle\beta^{\mathrm u},F^{\mathrm u}\rangle$ and $|\beta|_\eta=\|\beta^{\mathrm u}\|$ of Section~\ref{sec:units},
\[
\int\mathrm{d}\mathbb{P}_{N,T}(F)\,\exp\left(\langle\beta,F\rangle_\eta\right)=\mathbb{E}_{C(n)}\left[\exp\left(\langle\beta^{\mathrm u},F^{\mathrm u}\rangle\right)\right]\le\exp\left(\frac12\|\beta^{\mathrm u}\|^2\right)=\exp\left(\frac12|\beta|_\eta^2\right).
\]
The inequality is Theorem~\ref{thm:coordinate-domination}(a), applied to the real cochain $\beta^{\mathrm u}\in C^2(n\Z^3)$ on the coordinate torus $C(n)$, all of whose sides are $n\ge4$.

\emph{Clause~\eqref{it:cutoff-even}.}
By Proposition~\ref{prop:symmetries}\eqref{it:charge-conjugation}, the law $\mu_{C(n)}$ is invariant under multiplication by $-1$ on $\Omega$.
The argument given for clause~\eqref{it:cutoff-invariance}, with $-1$ in place of $g^{*}$, shows that $\mathbb{P}_{N,T}$ is invariant under $F\mapsto-F$.

\emph{Clause~\eqref{it:cutoff-positivity}.}
Fix $\nu_0$, the cochains $\beta_p$ and the numbers $a_p$; put $\ell=x_{\nu_0}$ and $\theta=\theta_{\ell,0}$.

\emph{The torus is adapted to $\ell$.}
The lattice $n\Z^3\cap\{\ell=0\}$ is spanned by the two vectors $n\cdot e_\mu$, $\mu\neq\nu_0$; it has systole $n\ge4$, and $\chi$ is $-1$ on both generators.
With this restriction of $\chi$ it is the cross-section $\Xi=\Xi(C(n),\ell)$.
The vector $n\cdot e_{\nu_0}$ lies in $n\Z^3$, one has $n\Z^3=(n\Z^3\cap\{\ell=0\})+\Z\cdot n\cdot e_{\nu_0}$, and $\ell(n\cdot e_{\nu_0})=n>0$.
Thus $n\cdot e_{\nu_0}$ is a closing vector, and the height period of $C(n)$ for $\ell$ is $n$ (Definition~\ref{def:slabs}).
Since $n_\ell=e_{\nu_0}$ has $|n_\ell|^2=1$, the vector $(n/|n_\ell|^2)\cdot n_\ell$ of Definition~\ref{def:slabs} is $n\cdot e_{\nu_0}$ itself; the lattice $(n\Z^3\cap\{\ell=0\})+\Z\cdot n\cdot e_{\nu_0}=n\Z^3$ has systole $n\ge4$, and $\chi(n\cdot e_{\nu_0})=-1$.
So the straightened torus $\mathbb{T}^0(\ell)$ of $\mathbb{T}=C(n)$ is $C(n)$: the torus $C(n)$ is adapted to $\ell$.

\emph{The two mirrors.}
As $n$ is even, the positions $-n/2$, $0$ and $n/2$ are integers and lie in $\Adm(\ell)=\frac12\Z$.
The mirrors $(\ell,0)$ and $(\ell,n/2)$ are therefore site planes, half a height period apart.

\emph{The trace form.}
Put $g:=n\cdot e_{\nu_0}$, and let $G,G'$ be arbitrary elements of $\mathcal{A}_\Xi[0,n/2]$.
By Lemma~\ref{lem:product-integrands}\eqref{it:sum-slab}, the reflection $\Theta_{\ell,0}G'$ belongs to $\mathcal{A}_\Xi[-n/2,0]$; reflected once more in $(\ell,0)$, the upper mirror of the slab between $-n/2$ and $0$, it gives $G'$ back, since $\Theta_{\ell,0}$ is an involution (Definition~\ref{def:grassmann-reflections} and \eqref{eq:setting-reflection-coefficients}).
By Corollary~\ref{cor:slab-admissible}, the reflection rule of Theorem~\ref{thm:slab-operators}\eqref{it:slab-reflection} holds for admissible products, with $\Theta_{\ell,c'}$ as the reflection in the upper mirror $(\ell,c')$ of a slab between positions $c<c'$.
For the slab between $-n/2$ and $0$ the vector of this rule is $v_{-n/2,0}=2\cdot\frac{n}{2}\cdot n_\ell/|n_\ell|^2=n\cdot e_{\nu_0}=g$,
so that the rule, applied to $\Theta_{\ell,0}G'$, gives $\mathfrak{S}[G']=V_g\cdot\mathfrak{S}[\Theta_{\ell,0}G']^{*}$; as $V_g$ is unitary, $\mathfrak{S}[\Theta_{\ell,0}G']=\mathfrak{S}[G']^{*}\cdot V_g$.
The trace formula of Theorem~\ref{thm:slab-operators}\eqref{it:slab-trace} likewise holds for admissible products (Corollary~\ref{cor:slab-admissible}).
Take in it the torus $C(n)$, whose cross-section for $\ell$ is $\Xi$, the closing vector $g$, for which $\ell(g)=n$, the number $r=2$, the positions $c_0=-n/2<c_1=0<c_2=n/2=c_0+\ell(g)$ and the elements $G_1=\Theta_{\ell,0}G'$, $G_2=G$;
after insertion of the last identity it reads
\[
Z_{C(n)}\left[\Theta_{\ell,0}(G')\cdot G\right]=\Tr_{\mathcal{H}(-n/2)}\left(V_g^{-1}\cdot U_{\mathrm{f}}^{\frac{1+\chi(g)}{2}}\cdot\mathfrak{S}[G]\cdot\mathfrak{S}[G']^{*}\cdot V_g\right).
\]
As $\chi(g)=-1$, the parity factor $U_{\mathrm{f}}^{\frac{1+\chi(g)}{2}}$ is the identity.
Conjugation by the unitary $V_g$ leaves the trace unchanged, and so does the exchange of two Hilbert--Schmidt factors, one from $\mathcal{H}(0)$ to $\mathcal{H}(n/2)$ and one back;
the right side is therefore $\Tr_{\mathcal{H}(0)}\left(\mathfrak{S}[G']^{*}\cdot\mathfrak{S}[G]\right)$ for all $G,G'\in\mathcal{A}_\Xi[0,n/2]$, and on $C(n)$ the trace form of Corollary~\ref{cor:trace-forms}\eqref{it:trace-form} contains no fermion parity $U_{\mathrm{f}}$.
The value $\chi(g)=-1$ decides this parity factor, but not the numbers $Z_{C(n)}[\Phi]$ at bounded Borel functions $\Phi$ of the curvature:
by Proposition~\ref{prop:symmetries}\eqref{it:sign-character}, these numbers and the law $\mu_{C(n)}$ are the same for every sign character of $n\Z^3$, and hence so are $\mathbb{P}_{N,T}$, the law $P^{\mathrm{ch}}_{N,T}$ of Definition~\ref{def:cutoffs} and the limit points of the family $\left(P^{\mathrm{ch}}_{N,T}\right)$ (Definition~\ref{def:limit-points}).

\emph{The integrand.}
Read in lattice units, a plaquette of $\Sigma^{++}_{\nu_0}$ is an $n\Z^3$-orbit of plaquettes of $\Z^3$.
Such an orbit has exactly one representative $\sigma$ whose vertices have their $\nu_0$-th coordinate in $\{1,\dots,n/2-1\}$ and whose base point $x(\sigma)$ has its two other coordinates in $\{0,\dots,n-1\}$.
As a closed subset of $\R^3$, this representative lies in the open slab $\{0<\ell<n/2\}$.
Since the $\beta_p$ are supported on $\Sigma^{++}_{\nu_0}$, the function
\[
G:=\sum_{p'}a_{p'}\cdot\exp\left(i\langle\beta_{p'},F\rangle_\eta\right)=\sum_{p'}a_{p'}\cdot\exp\left(i\sum_\sigma\beta^{\mathrm u}_{p'}(\sigma)\cdot\widetilde F^{\mathrm u}(\sigma)\right),
\]
where $\sigma$ ranges over these representatives, is a bounded continuous function of finitely many values $\widetilde F^{\mathrm u}(\sigma)$ at plaquettes contained in $\{0<\ell<n/2\}$.
It is thus an admissible product without factors of Gaussian growth, an element of $\mathcal{A}(\{0<\ell<n/2\})$ (Definition~\ref{def:integrands}).
A product without boxes embeds in every lattice (Definition~\ref{def:integrands}); hence $G\in\mathcal{A}_\Xi[0,n/2]$ (Definition~\ref{def:slabs}).

\emph{The reflected integrand.}
The function $G$ is bosonic, so that $(\Theta_{\ell,0}G)(F)=\overline{G(\theta^{*}F)}$ by the rules \eqref{eq:setting-reflection-coefficients} and \eqref{eq:setting-group-action}.
The pull-back $\theta^{*}$ is $\hat{\mathsf{R}}_{\nu_0}$, and, $\theta$ being its own inverse, $\langle\beta_p,\theta^{*}F\rangle_\eta=\langle\hat{\mathsf{R}}_{\nu_0}\beta_p,F\rangle_\eta$ (Section~\ref{sec:units}).
Consequently
\[
\overline{G(\theta^{*}F)}=\sum_p\overline{a_p}\cdot\exp\left(-i\langle\hat{\mathsf{R}}_{\nu_0}\beta_p,F\rangle_\eta\right)
\]
and
\[
\Theta_{\ell,0}(G)\cdot G=\sum_{p,p'}\overline{a_p}\cdot a_{p'}\cdot\exp\left(i\left\langle\beta_{p'}-\hat{\mathsf{R}}_{\nu_0}\beta_p,F\right\rangle_\eta\right).
\]

\emph{Positivity.}
Corollary~\ref{cor:trace-forms}\eqref{it:trace-form}, applied to the adapted torus $C(n)$, the mirrors $(\ell,0)$, $(\ell,n/2)$ and $G'=G\in\mathcal{A}_\Xi[0,n/2]$, gives
\[
Z_{C(n)}\left[\Theta_{\ell,0}(G)\cdot G\right]=\Tr_{\mathcal{H}(0)}\left(\mathfrak{S}[G]^{*}\cdot\mathfrak{S}[G]\right)\ge0,
\]
where $\mathfrak{S}[G]$ is the Hilbert--Schmidt slab operator of $G$.
The left side of \eqref{eq:continuum-admissible-even} equals $\mathbb{E}_{C(n)}[\Theta_{\ell,0}(G)\cdot G]=Z_{C(n)}[\Theta_{\ell,0}(G)\cdot G]/Z_{C(n)}[1]$ and is therefore non-negative, the denominator $Z_{C(n)}[1]$ being positive.

\emph{Clause~\eqref{it:cutoff-chessboard}.}
The torus $C(n)$ is the rectangular torus $\mathbb{T}(\mathfrak{F}_{\mathrm{c}};n,n,n)$, and its height periods $p_1=p_2=p_3=n$ lie in $2h\Z$; so $C(n)$ carries $\mathcal{Q}$ (Definition~\ref{def:cell-structures}).
For the coordinate frame $|Q|=h^3$, while $|C(n)|=n^3$.
Hence $C(n)$ has $|C(n)|/|Q|=(n/h)^3$ cells, and the chessboard seminorm of $G_\kappa$ on $C(n)$ (Definition~\ref{def:folds}) is $\|G_\kappa\|_{C(n)}=Z_{C(n)}\left[\Fold_{C(n)}G_\kappa\right]^{\left(\frac{h}{n}\right)^3}$.
Theorem~\ref{thm:chessboard}\eqref{it:chessboard-carried}, applied to $C(n)$ and $\mathcal{Q}$, yields
\[
Z_{C(n)}\left[\Fold_{C(n)}G_\kappa\right]\ge0,\qquad
\left|Z_{C(n)}\left[\prod_\kappa G_\kappa\right]\right|\le\prod_\kappa Z_{C(n)}\left[\Fold_{C(n)}G_\kappa\right]^{\left(\frac{h}{n}\right)^3}.
\]
The product over $\kappa$ has $(n/h)^3$ factors, so that $Z_{C(n)}[1]=\prod_\kappa Z_{C(n)}[1]^{\left(\frac{h}{n}\right)^3}$.
Dividing both inequalities by $Z_{C(n)}[1]>0$ gives the two inequalities of clause~\eqref{it:cutoff-chessboard}.
Finally, let $G_\kappa$ be bosonic, let $\kappa'\in\Z^3$, and let $g\in W(\mathcal{Q})$ be the element with $gQ_\kappa=Q_{\kappa'}$.
By Definition~\ref{def:folds}, $G_\kappa^{[\kappa']}=G_\kappa^{g}$, and the action \eqref{eq:setting-group-action} on bosonic integrands gives $G_\kappa^{g}(F)=G_\kappa(g^{*}F)$ if $g$ is even and $G_\kappa^{g}(F)=\overline{G_\kappa(g^{*}F)}$ if $g$ is odd.
\end{proof}

Clauses~\eqref{it:cutoff-invariance} and~\eqref{it:cutoff-domination} are the conditions (Ad1) and (Ad2) of \cite[Definition~7.1]{NoncompactQED} for the law $\mathbb{P}_{N,T}$;
the actions written $v_*$ and $r_*$ in \cite[Section~7.1.4]{NoncompactQED} are the maps $\tau_{-v/\eta}^{*}$ and $(g^{-1})^{*}$, with $g=r$ (Section~\ref{sec:units}).
The inequality (Ad3$'$) is the counterpart of (Ad3) for the even sides $n$ (Remark~\ref{rem:cutoff-law}).
As \cite[Section~7.1.4]{NoncompactQED} does for (Ad1)--(Ad3), the names (Ad1), (Ad2) and (Ad3$'$) are also given to the same properties of a single Borel probability measure on $C^2(\Lambda_{\eta,R})$.

\begin{remark}\label{rem:cutoff-law}
\begin{enumerate}
\item
In \cite[Section~7.1]{NoncompactQED} the sides $n$ are odd, and condition (Ad3) of \cite[Definition~7.1]{NoncompactQED} concerns the tests supported on the plaquettes whose vertices have their $\nu$-th coordinate in $\left\{0,\eta,\dots,\frac{n-1}{2}\cdot\eta\right\}$ and which are contained in neither of the planes $\{x_\nu=0\}$ and $\left\{x_\nu=\frac{n-1}{2}\cdot\eta\right\}$; such a plaquette may have an edge in one of these planes.
The plaquettes of $\Sigma^{++}_\nu$ are those contained in the open slab $\{0<x_\nu<R/2\}$, and none of them meets the planes $\{x_\nu=0\}$ and $\{x_\nu=R/2\}$.
Hence \eqref{eq:continuum-admissible-even} follows from the inequality into which (Ad3) turns when $\frac{n-1}{2}$ is replaced by $n/2$, and (Ad3$'$) is the weaker of these two conditions.
\item
Let $\mathcal{Q}$ be a cell structure as in Proposition~\ref{prop:cutoff-law}\eqref{it:cutoff-chessboard}.
For $a=1,2,3$, its walls $(x_a,o_a+j\cdot h)$, $j\in\Z$, are site planes if $o_a\in\Z$, and link planes if $o_a\in\Z+\frac12$.
Let $B$ be a box contained in an open cell $Q_\kappa$, and let $\Phi$ be a Borel function of the values $F_\sigma$, $\sigma\in B^{(2)}$, such that
\[
|\Phi(F)|\le C\cdot\exp\left(C\cdot\left(\sum_{\sigma\in B^{(2)}}F_\sigma^2\right)^{\frac12}\right)
\]
for some constant $C<\infty$, the growth allowed in the class of functions of \cite[Section~2.3]{NoncompactQED}.
The map $\hat A\mapsto\Phi(d\hat A)$ on $\R^{B^{(1)}}$ is a bosonic integrand that reads in $B$: it is Borel, and it is gauge invariant because it depends on $d\hat A$ only.
Since $C\cdot t\le C^2+\frac14\cdot t^2$ for $t\ge0$, it satisfies the growth bound of Definition~\ref{def:integrands} with the constant $C\cdot\exp(C^2)$ and with $\varepsilon_G=\frac14$, so that it lies in $\mathcal{G}(B)$;
as a product with one factor of Gaussian growth and the box $B\subset Q_\kappa$, it belongs to $\mathcal{A}(Q_\kappa)$.
\item
Besides Proposition~\ref{prop:cutoff-law}\eqref{it:cutoff-domination}, Theorem~\ref{thm:coordinate-domination} gives on $C(n)$ the bound for complex exponents in its clause~(a), the bounds on the moments and on the fourth cumulant in~(b), and the Peierls bound in~(c); these are the inequalities (a)--(c) of \cite[Corollary~10.8]{NoncompactQED}, which are stated there for the all-flux law of the companion paper, formed with the Wilson determinant.
Through the identities $\langle\beta,F\rangle_\eta=\langle\beta^{\mathrm u},F^{\mathrm u}\rangle$ and $|\beta|_\eta=\|\beta^{\mathrm u}\|$ of Section~\ref{sec:units} they hold for $\mathbb{P}_{N,T}$, with $\langle\beta,F\rangle_\eta$ and $|\beta|_\eta$ in place of $\langle\beta^{\mathrm u},F^{\mathrm u}\rangle$ and $\|\beta^{\mathrm u}\|$.
For example,
\[
\int\mathrm{d}\mathbb{P}_{N,T}(F)\,\langle\beta,F\rangle_\eta^{2k}\le\left(2k\cdot\exp(1)\cdot|\beta|_\eta^2\right)^k
\]
for every real $\beta\in C^2(\Lambda_{\eta,R})$ and every integer $k\ge1$.
\end{enumerate}
\end{remark}

\subsection{Limit points of admissible families}\label{sec:admissible-limits}

\begin{proposition}[Admissible families on even tori]\label{prop:admissible-families}
For every cutoff $(N,T)$ let $\eta=L^{-N}$, $R=8\cdot L^{\lfloor\log_L\bar m^{-1}\rfloor+T}$ and $\Lambda_{\eta,R}$ be as in Definition~\ref{def:cutoffs},
and let $\mathbb{P}_{N,T}$ be a Borel probability measure on the space $\ker d$ of the closed two-cochains on $\Lambda_{\eta,R}$ with the following three properties;
the first two are the conditions \textup{(Ad1)} and \textup{(Ad2)} of the definition of admissible families \cite[Definition~7.1]{NoncompactQED}, where the maps $\tau^{*}_{-v/\eta}$ and $(g^{-1})^{*}$ of \textup{(Ad1)} are written $v_*$ and $r_*$, with $r=g$.
\begin{enumerate}
\item[(Ad1)]
The measure $\mathbb{P}_{N,T}$ is invariant under the pull-backs $\tau^{*}_{-v/\eta}$, $v\in\eta\Z^3$, and $(g^{-1})^{*}$, $g\in\Ocub$, of the two-cochains on $\Lambda_{\eta,R}$ (Definition~\ref{def:cochains});
these maps commute with $d$ and therefore map $\ker d$ onto itself.
\item[(Ad2)]
For every $\beta\in C^2(\Lambda_{\eta,R})$,
\[
\int\mathrm{d}\mathbb{P}_{N,T}(F)\,\exp\langle\beta,F\rangle_\eta\le\exp\left(\frac12|\beta|_\eta^2\right).
\]
\item[(Ad3$'$)]
For every $\nu_0\in\{1,2,3\}$, every integer $s\ge1$, all $\beta_1,\dots,\beta_s\in C^2(\Lambda_{\eta,R})$ supported on the set $\Sigma^{++}_{\nu_0}$ of Lemma~\ref{lem:positive-tests} and all $a_1,\dots,a_s\in\C$,
with the pull-back $\hat{\mathsf{R}}_{\nu_0}=\theta^{*}_{x_{\nu_0},0}$ by the reflection in the plane $\{x_{\nu_0}=0\}$,
\[
\sum_{p,q=1}^{s}\overline{a_p}\cdot a_q\int\mathrm{d}\mathbb{P}_{N,T}(F)\,\exp\left(i\left\langle\beta_q-\hat{\mathsf{R}}_{\nu_0}\beta_p,F\right\rangle_\eta\right)\ge0 .
\]
\end{enumerate}
Let $P_{N,T}$ be the law of $F^{\sharp}$ under $\mathbb{P}_{N,T}$.
Then the following hold.
\begin{enumerate}
\item[(a)]
For every cutoff $(N,T)$ the assertions of \cite[Theorem~7.11]{NoncompactQED}, except those of its last paragraph, hold for the measure $\mathbb{P}_{N,T}$;
and the assertions of \cite[Theorem~7.12]{NoncompactQED} hold for the family $(P_{N,T})$.
\item[(b)]
Let $(N_k,T_k)_{k\in\N}$ be a sequence of cutoffs with $N_k\to\infty$ and $T_k\to\infty$,
and let $\mu$ be a Borel probability measure on $H_{-4}$ to which $P_{N_k,T_k}$ converges weakly on $H_{-4}$.
Then $\mu$ has the properties (1)--(8) of \cite[Theorem~7.15]{NoncompactQED}, all the properties of \cite[Theorem~7.17]{NoncompactQED} and all those of \cite[Corollary~7.18]{NoncompactQED}.
If moreover every measure $\mathbb{P}_{N_k,T_k}$ is invariant under $F\mapsto-F$, then $\mu$ is invariant under $F\mapsto-F$, and
\[
\int\mathrm{d}\mu(F)\,\prod_{i=1}^{n}F(f_i)=0
\]
for every odd $n\ge1$ and all real $f_1,\dots,f_n\in C^{\infty}_{\mathrm c}(\R^3;\Lambda^2)$.
\end{enumerate}
In (a) and (b) the statements of \cite{NoncompactQED} are read with the cutoffs of Definition~\ref{def:cutoffs} and the tori $\Lambda_{\eta,R}$ formed with them, and with the family $(\mathbb{P}_{N,T})$ in the place of an admissible family;
$\mathbb{P}_{N,T}$, $P_{N,T}$ and $F^{\sharp}(f)$ take the places of the measures $\mathrm{Pr}_{N,T}$, $\mathsf{P}_{N,T}$ and of the variables $X_{\eta,R}(f)$, $X_k(f)$ there;
and, in (b), the sequence $(N_k,T_k)$ and the expectation with respect to $\mathbb{P}_{N_k,T_k}$ take the places of the sequence of cutoffs and of the expectation $E_k$ there.
\end{proposition}

\begin{proof}
The argument follows the proofs of \cite[Theorems~7.11 and~7.12]{NoncompactQED}, of \cite[Theorem~7.15(2)]{NoncompactQED} and of \cite[Theorem~7.17(3)]{NoncompactQED},
together with those parts of the proofs of \cite[Theorem~7.15(1), (4)--(7) and~(9)]{NoncompactQED} that establish the conditions (L1)--(L4) of \cite[Section~7.5]{NoncompactQED} and, for lattice laws invariant under $F\mapsto-F$, the evenness of the characteristic functional of the limit.
In this argument Lemma~\ref{lem:test-norms} takes the place of \cite[Lemma~7.5]{NoncompactQED},
Lemma~\ref{lem:test-equivariance} that of \cite[Lemma~7.6]{NoncompactQED} and, together with (Ad1), that of \cite[Lemma~7.7(d)]{NoncompactQED},
Lemmas~\ref{lem:test-bianchi} and~\ref{lem:test-unfolding} that of \cite[Lemma~7.8]{NoncompactQED},
and Lemma~\ref{lem:positive-tests} that of \cite[Lemma~7.9]{NoncompactQED};
in the step on reflection positivity, (Ad3$'$) takes the place of the condition (Ad3) of \cite[Definition~7.1]{NoncompactQED}.

For every cutoff $(N,T)$ let $\mathbb{E}_{N,T}$ be the expectation with respect to $\mathbb{P}_{N,T}$, regarded as a measure on $C^2(\Lambda_{\eta,R})$ that is carried by $\ker d$.
For $F\in C^2(\Lambda_{\eta,R})$ and a test form $f$ one has $F^{\sharp}(f)=\langle f_\eta,F\rangle_\eta$ (Lemma~\ref{lem:test-unfolding});
this number depends complex-linearly on $f$, and it is real for real $f$.
Along a sequence $(N_k,T_k)$ of cutoffs write
\[
\eta_k=L^{-N_k},\qquad R_k=8\cdot L^{\lfloor\log_L\bar m^{-1}\rfloor+T_k},\qquad\mathbb{E}_k=\mathbb{E}_{N_k,T_k}.
\]
For a test form $f$, real or complex, $f_{\eta_k}$ denotes its test cochain on the torus $\Lambda_{\eta_k,R_k}$ of the cutoff $(N_k,T_k)$, that is, the cochain $f_\eta$ of Definition~\ref{def:test-cochains} formed with $\eta=\eta_k$ and $R=R_k$;
for an indexed test form $f_i$ it is written $(f_i)_{\eta_k}$.
The pairing $\langle\cdot,\cdot\rangle_{\eta_k}$ and the norm $|\cdot|_{\eta_k}$ are those of the two-cochains on $\Lambda_{\eta_k,R_k}$,
and under $\mathbb{E}_k$ the variable $F^{\sharp}(f)=\langle f_{\eta_k},F\rangle_{\eta_k}$ is the smeared field of the two-cochain $F$ on $\Lambda_{\eta_k,R_k}$, formed at the cutoff $(N_k,T_k)$.
\begin{enumerate}
\item[(1)]\itemmark{1}\label{it:admissible-bounds}
\emph{Bounds at a fixed cutoff.}
Fix a cutoff $(N,T)$.
By (Ad2), the hypothesis of \cite[Proposition~7.10]{NoncompactQED} is fulfilled by the probability space $\left(C^2(\Lambda_{\eta,R}),\mathbb{P}_{N,T}\right)$, the real vector space $C^2(\Lambda_{\eta,R})$ with the Hilbert norm $|\cdot|_\eta$, and the linear map that sends $\beta$ to the random variable $F\mapsto\langle\beta,F\rangle_\eta$.
The test map $f\mapsto f_\eta$ is linear and sends real test forms to real cochains, and $F^{\sharp}(f)=\langle f_\eta,F\rangle_\eta$.
The same hypothesis therefore holds for the real vector space of the real test forms, with the Hilbert seminorm $f\mapsto|f_\eta|_\eta$ and the map $f\mapsto F^{\sharp}(f)$; in this form it reads
\[
\mathbb{E}_{N,T}\exp\left(F^{\sharp}(f)\right)\le\exp\left(\frac12|f_\eta|_\eta^2\right)\qquad(f\in\mathcal S\text{ real}).
\]
As $R\ge8$, Lemma~\ref{lem:test-norms}\eqref{it:test-norms-torus} gives
\[
|f_\eta|_\eta\le p_R(f)\le p_1(f)
\]
for every test form $f$, real or complex, and $|f_\eta|_\eta\le\|f\|_2$ whenever $\supp f\subset y_0+\left(-\frac R2,\frac R2\right)^3$ for some $y_0\in\R^3$.
Consequently the assertions (1)--(7) of \cite[Proposition~7.10]{NoncompactQED} hold under $\mathbb{P}_{N,T}$ for the variables $\langle\beta,F\rangle_\eta$ with the norm $|\beta|_\eta$, and for the variables $F^{\sharp}(f)$ with the seminorm $|f|:=|f_\eta|_\eta$.
Together with the two bounds for $|f_\eta|_\eta$, these are the assertions of \cite[Theorem~7.11]{NoncompactQED} other than those of its last paragraph.
\item[(2)]\itemmark{2}\label{it:admissible-tightness}
\emph{Second moments and tightness.}
Let $h_\kappa$ be the Hermite two-forms of \cite[Section~7.1.5]{NoncompactQED}:
they form an orthonormal basis of $L^2(\R^3;\Lambda^2)$, they are real test forms, and $\mathsf{A}h_\kappa=\lambda_\kappa\cdot h_\kappa$ with $\mathsf{A}=1+|x|^2-\Delta$ and $\lambda_\kappa\ge4$;
moreover $\|F\|_{-3}^2=\sum_\kappa\lambda_\kappa^{-6}\cdot|F(h_\kappa)|^2$ for $F\in\mathcal S'$, and $\mathsf{Z}_a=\sum_\kappa\lambda_\kappa^{-a}$ for $a>3$.
For $F\in C^2(\Lambda_{\eta,R})$ and every index $\kappa$, the Cauchy--Schwarz inequality and Lemma~\ref{lem:test-norms}\eqref{it:test-norms-torus} give
\[
\left|F^{\sharp}(h_\kappa)\right|=\left|\langle(h_\kappa)_\eta,F\rangle_\eta\right|\le|(h_\kappa)_\eta|_\eta\cdot|F|_\eta\le p_1(h_\kappa)\cdot|F|_\eta\le31\lambda_\kappa\cdot|F|_\eta ,
\]
and hence
\[
\|F^{\sharp}\|_{-3}^2=\sum_\kappa\lambda_\kappa^{-6}\cdot\left|F^{\sharp}(h_\kappa)\right|^2\le961\cdot\mathsf{Z}_4\cdot|F|_\eta^2 .
\]
So $F\mapsto F^{\sharp}$ is a linear map of the finite-dimensional space $C^2(\Lambda_{\eta,R})$ into $H_{-3}$, hence continuous, and $P_{N,T}$ is a Borel probability measure on $H_{-3}$ and on $H_{-4}$.
Tonelli's theorem, step~\eqref{it:admissible-bounds} together with \cite[Proposition~7.10(2)]{NoncompactQED}, and Lemma~\ref{lem:test-norms}\eqref{it:test-norms-torus} give
\[
\int\mathrm{d}P_{N,T}(F)\,\|F\|_{-3}^2=\sum_\kappa\lambda_\kappa^{-6}\cdot\mathbb{E}_{N,T}\left[F^{\sharp}(h_\kappa)^2\right]\le\sum_\kappa\lambda_\kappa^{-6}\cdot p_1(h_\kappa)^2\le961\cdot\mathsf{Z}_4\le61 ,
\]
the last inequality because $\mathsf{Z}_4<\frac1{16}$ \cite[Section~7.1.5]{NoncompactQED}.

For $\varrho>0$ let $\mathcal{B}_\varrho=\left\{F\in H_{-3}:\|F\|_{-3}\le\varrho\right\}$, the set denoted $K_\varrho$ in \cite[Theorem~7.12(2)]{NoncompactQED}.
By the first assertion of \cite[Theorem~7.12(2)]{NoncompactQED}, which concerns $\mathcal S'$ alone, $\mathcal{B}_\varrho$ is a compact subset of $H_{-4}$.
For $F\in H_{-4}$ and $f\in\mathcal S$ one has $|F(f)|\le\|F\|_{-4}\cdot\|f\|_4$ \cite[Section~7.1.5]{NoncompactQED}, and $\|\cdot\|_4$ is bounded on every bounded subset of $\mathcal S$;
so the inclusion of $H_{-4}$ into $\mathcal S'$ is continuous for the strong topology of $\mathcal S'$, hence also for its weak-$*$ topology, and $\mathcal{B}_\varrho$ is compact in $\mathcal S'$ for both topologies.
By Chebyshev's inequality and the bound on the second moments,
\[
P_{N,T}\left(H_{-4}\setminus\mathcal{B}_\varrho\right)=P_{N,T}\left(\left\{F:\|F\|_{-3}>\varrho\right\}\right)\le61\cdot\varrho^{-2}
\]
for every cutoff $(N,T)$.
The family $(P_{N,T})$ is therefore tight on $H_{-4}$, and on $\mathcal S'$ with its strong and with its weak-$*$ topology.
Let $(N_k,T_k)$ be a sequence of cutoffs.
Since $H_{-4}$ is a separable Hilbert space, Prokhorov's theorem \cite[Theorem~5.1]{Billingsley1999} yields a subsequence along which $P_{N_k,T_k}$ converges weakly on $H_{-4}$ to a Borel probability measure $\mu$;
a bounded function on $\mathcal S'$ that is continuous for the strong or for the weak-$*$ topology restricts to a bounded continuous function on $H_{-4}$, so that the convergence holds on $\mathcal S'$ as well.
The set $\mathcal{B}_\varrho$ is closed in $H_{-4}$, and along the subsequence the portmanteau theorem \cite[Theorem~2.1]{Billingsley1999} gives
\[
\mu(\mathcal{B}_\varrho)\ge\limsup_{k\to\infty}P_{N_k,T_k}(\mathcal{B}_\varrho)\ge1-61\cdot\varrho^{-2}\qquad(\varrho>0);
\]
as $\mathcal{B}_\varrho\subset H_{-3}$ for every $\varrho$, it follows that $\mu(H_{-3})=1$.
These are the assertions of \cite[Theorem~7.12]{NoncompactQED}; together with step~\eqref{it:admissible-bounds} they prove assertion (a) of the proposition.
\end{enumerate}

For the remainder of the proof let $(N_k,T_k)$ and $\mu$ be as in assertion (b); then $\eta_k\to0$ and $R_k\to\infty$.
As in Definition~\ref{def:laws}, $\mu$ is also regarded as the probability measure $B\mapsto\mu(B\cap H_{-4})$ on $(\mathcal S',\Sigma)$;
its characteristic functional is
\[
\hat\mu(f)=\int\mathrm{d}\mu(F)\,\exp\left(iF(f)\right)\qquad(f\in\mathcal S\text{ real}).
\]
Let $f_1,\dots,f_s$ be test forms, real or complex.
The map $F\mapsto\left(F(f_1),\dots,F(f_s)\right)$ of $H_{-4}$ into $\C^s$ is continuous, by the inequality $|F(f)|\le\|F\|_{-4}\cdot\|f\|_4$,
and $P_{N_k,T_k}$ is the law of $F^{\sharp}$ under $\mathbb{P}_{N_k,T_k}$.
The mapping theorem \cite[Theorem~2.7]{Billingsley1999} therefore gives
\begin{equation}\label{eq:continuum-weak-limit}
\lim_{k\to\infty}\mathbb{E}_k\,u\left(F^{\sharp}(f_1),\dots,F^{\sharp}(f_s)\right)=\int\mathrm{d}\mu(F)\,u\left(F(f_1),\dots,F(f_s)\right)
\end{equation}
for every bounded continuous function $u\colon\C^s\to\C$.

\begin{enumerate}
\item[(3)]\itemmark{3}\label{it:admissible-domination}
\emph{Gaussian domination.}
Let $f\in\mathcal S$ be real and $\varepsilon>0$.
The function $F\mapsto\min\left(\exp F(f),\varepsilon^{-1}\right)$ is bounded and continuous on $H_{-4}$.
By step~\eqref{it:admissible-bounds} and Lemma~\ref{lem:test-norms}\eqref{it:test-norms-torus},
\[
\mathbb{E}_k\exp\left(F^{\sharp}(f)\right)\le\exp\left(\frac12|f_{\eta_k}|_{\eta_k}^2\right)\le\exp\left(\frac12p_{R_k}(f)^2\right),
\]
and $p_{R_k}(f)\to\|f\|_2$ as $k\to\infty$, because $R_k\to\infty$ (Lemma~\ref{lem:test-norms}\eqref{it:test-norms-torus}).
With \eqref{eq:continuum-weak-limit} it follows that
\[
\begin{aligned}
\int\mathrm{d}\mu(F)\,\min\left(\exp F(f),\varepsilon^{-1}\right)
&=\lim_{k\to\infty}\mathbb{E}_k\min\left(\exp F^{\sharp}(f),\varepsilon^{-1}\right)\\
&\le\liminf_{k\to\infty}\exp\left(\frac12p_{R_k}(f)^2\right)=\exp\left(\frac12\|f\|_2^2\right).
\end{aligned}
\]
As $\varepsilon\to0$, the left side increases to $\int\mathrm{d}\mu(F)\,\exp F(f)$ by monotone convergence; hence
\begin{equation}\tag{L1}\label{eq:continuum-limit-domination}
\int\mathrm{d}\mu(F)\,\exp F(f)\le\exp\left(\frac12\|f\|_2^2\right)\qquad(f\in\mathcal S\text{ real}).
\end{equation}
\item[(4)]\itemmark{4}\label{it:admissible-moments}
\emph{Generating functions, moments and cumulants.}
Let $f_1,\dots,f_s$ be test forms, real or complex, and for $z\in\C^s$ put
\[
\mathcal{Z}_k(z)=\mathbb{E}_k\exp\left(\sum_{i=1}^{s}z_i\cdot F^{\sharp}(f_i)\right),\qquad
\mathcal{Z}(z)=\int\mathrm{d}\mu(F)\,\exp\left(\sum_{i=1}^{s}z_i\cdot F(f_i)\right).
\]
For $z\in\C^s$ the real part $\Re\sum_iz_i\cdot f_i$ is a real test form, and the modulus of $\exp\left(\sum_iz_i\cdot F^{\sharp}(f_i)\right)$ is $\exp\left(F^{\sharp}\left(\Re\sum_iz_i\cdot f_i\right)\right)$.
Step~\eqref{it:admissible-bounds}, here (Ad2) for the cochain $\beta=2\cdot\left(\Re\sum_iz_i\cdot f_i\right)_{\eta_k}$, and Lemma~\ref{lem:test-norms}\eqref{it:test-norms-torus} give
\[
\begin{aligned}
\mathbb{E}_k\left|\exp\left(\sum_{i=1}^{s}z_i\cdot F^{\sharp}(f_i)\right)\right|^2
&=\mathbb{E}_k\exp\left(2F^{\sharp}\left(\Re\sum_{i=1}^{s}z_i\cdot f_i\right)\right)\\
&\le\exp\left(2\left|\left(\Re\sum_{i=1}^{s}z_i\cdot f_i\right)_{\eta_k}\right|_{\eta_k}^2\right)\le\exp\left(2p_1\left(\Re\sum_{i=1}^{s}z_i\cdot f_i\right)^2\right).
\end{aligned}
\]
The right side does not depend on $k$, and it is bounded for $z$ in a compact subset of $\C^s$, since $p_1$ is a seminorm and $\Re\sum_iz_i\cdot f_i=\sum_i\left(\Re z_i\cdot\Re f_i-\Im z_i\cdot\Im f_i\right)$.
Hence the real and the imaginary part of $\exp\left(\sum_iz_i\cdot F^{\sharp}(f_i)\right)$ are uniformly integrable with respect to the measures $\mathbb{P}_{N_k,T_k}$;
by \eqref{eq:continuum-weak-limit} their laws converge weakly to the laws of the real and the imaginary part of $F\mapsto\exp\left(\sum_iz_i\cdot F(f_i)\right)$ under $\mu$;
and the theorem on uniformly integrable sequences \cite[Theorem~3.5]{Billingsley1999} gives $\mathcal{Z}_k(z)\to\mathcal{Z}(z)$ for every $z\in\C^s$.

By step~\eqref{it:admissible-bounds} and \cite[Proposition~7.10(1)]{NoncompactQED}, each $\mathcal{Z}_k$ is an entire function on $\C^s$,
and $|\mathcal{Z}_k(z)|\le\exp\left(p_1\left(\Re\sum_iz_i\cdot f_i\right)^2\right)$ by the Cauchy--Schwarz inequality;
so the functions $\mathcal{Z}_k$ are bounded on every compact subset of $\C^s$, uniformly in $k$.
Cauchy's estimates make them equi-Lipschitz on every compact set.
The pointwise convergence is therefore locally uniform, $\mathcal{Z}$ is entire, and by Cauchy's integral formula every derivative of $\mathcal{Z}_k$ converges locally uniformly to the corresponding derivative of $\mathcal{Z}$.
The derivatives at $z=0$ are the moments, since the exponential moments of \cite[Proposition~7.10(1)]{NoncompactQED}, which hold under $\mathbb{P}_{N_k,T_k}$ by step~\eqref{it:admissible-bounds} and under $\mu$ by \eqref{eq:continuum-limit-domination}, allow differentiation under the integral.
Thus all moments converge:
\[
\mathbb{E}_k\prod_{i=1}^{s}F^{\sharp}(f_i)^{\alpha_i}\longrightarrow\int\mathrm{d}\mu(F)\,\prod_{i=1}^{s}F(f_i)^{\alpha_i}\qquad\left(k\to\infty,\ \alpha\in\{0,1,2,\dots\}^s\right).
\]

Fix a compact subset of the open set $\left\{z\in\C^s:\sum_i|z_i|\cdot\|f_i\|_2<\frac25\right\}$.
By Lemma~\ref{lem:test-norms}\eqref{it:test-norms-torus}, $\sum_i|z_i|\cdot|(f_i)_{\eta_k}|_{\eta_k}\le\sum_i|z_i|\cdot p_{R_k}(f_i)$, and $p_{R_k}(f_i)\to\|f_i\|_2$ for each $i$;
the compact set being contained in the open set, $\sum_i|z_i|\cdot|(f_i)_{\eta_k}|_{\eta_k}\le\frac25$ on it for all large $k$.
For such $k$ the triangle inequality gives $\left|\left(\sum_iz_i\cdot f_i\right)_{\eta_k}\right|_{\eta_k}\le\frac25$ on the compact set,
and \cite[Proposition~7.10(6)]{NoncompactQED}, which applies by step~\eqref{it:admissible-bounds}, gives $|\mathcal{Z}_k-1|\le\frac12$ there; hence also $|\mathcal{Z}-1|\le\frac12$ there.
The principal branch of the logarithm is uniformly continuous on $\left\{w\in\C:|w-1|\le\frac12\right\}$, so that $\log\mathcal{Z}_k\to\log\mathcal{Z}$ uniformly on the compact set.
Cauchy's integral formula on a closed polydisc centred at $0$ inside the open set above then shows that every derivative of $\log\mathcal{Z}_k$ at $0$, a joint cumulant of $F^{\sharp}(f_1),\dots,F^{\sharp}(f_s)$ under $\mathbb{P}_{N_k,T_k}$,
converges to the corresponding derivative of $\log\mathcal{Z}$, the joint cumulant of $F(f_1),\dots,F(f_s)$ under $\mu$.
This is assertion (2) of \cite[Theorem~7.15]{NoncompactQED} for the sequence $(N_k,T_k)$.
\item[(5)]\itemmark{5}\label{it:admissible-closed}
\emph{Closedness.}
Let $\xi\in C^{\infty}_{\mathrm c}(\R^3)$ be real, and let $u^{(\eta)}$ and $d^{*}\xi$ be the two-forms of Lemma~\ref{lem:test-bianchi}.
Let $F\in\ker d\subset C^2(\Lambda_{\eta_k,R_k})$, and let $\widetilde F^{\mathrm u}$ be the periodic extension to $\Z^3$ of the cochain $F^{\mathrm u}=\eta_k^{\frac32}\cdot F$ in lattice units;
it is a bounded closed two-cochain on $\Z^3$.
Lemma~\ref{lem:test-unfolding} and Lemma~\ref{lem:test-bianchi}\eqref{it:test-bianchi-closed} give
\[
F^{\sharp}\left(u^{(\eta_k)}\right)=\sum_\sigma u^{(\eta_k)}_{[\eta_k]}(\sigma)\cdot\widetilde F^{\mathrm u}(\sigma)=0 ,
\]
the sum being taken over the plaquettes $\sigma$ of $\Z^3$.
As $\mathbb{P}_{N_k,T_k}$ is carried by $\ker d$, the identity $F^{\sharp}\left(u^{(\eta_k)}\right)=0$ holds $\mathbb{P}_{N_k,T_k}$-almost surely.
The two-forms $d^{*}\xi$ and $u^{(\eta_k)}$ vanish at every point whose distance from $\supp\xi$ exceeds $\frac{\eta_k}2$;
since $R_k\to\infty$, their supports lie in $\left(-\frac{R_k}2,\frac{R_k}2\right)^3$ for all large $k$.
Put
\[
C_\xi=\left(\frac1{24}\sum_{j=1}^{3}\left\|\partial_j^3\xi\right\|_2\right)^2 .
\]
For all large $k$, step~\eqref{it:admissible-bounds} together with \cite[Proposition~7.10(2)]{NoncompactQED}, Lemma~\ref{lem:test-norms}\eqref{it:test-norms-torus} and Lemma~\ref{lem:test-bianchi}\eqref{it:test-bianchi-estimate} give
\[
\begin{aligned}
\mathbb{E}_k\left[F^{\sharp}(d^{*}\xi)^2\right]
&=\mathbb{E}_k\left[F^{\sharp}\left(d^{*}\xi-u^{(\eta_k)}\right)^2\right]\le\left|\left(d^{*}\xi-u^{(\eta_k)}\right)_{\eta_k}\right|_{\eta_k}^2\\
&\le\left\|d^{*}\xi-u^{(\eta_k)}\right\|_2^2\le C_\xi\cdot\eta_k^4 .
\end{aligned}
\]
Let $\varepsilon>0$.
The function $F\mapsto\min\left(F(d^{*}\xi)^2,\varepsilon^{-1}\right)$ is bounded and continuous on $H_{-4}$, and \eqref{eq:continuum-weak-limit} gives
\[
\int\mathrm{d}\mu(F)\,\min\left(F(d^{*}\xi)^2,\varepsilon^{-1}\right)=\lim_{k\to\infty}\mathbb{E}_k\min\left(F^{\sharp}(d^{*}\xi)^2,\varepsilon^{-1}\right)\le\limsup_{k\to\infty}C_\xi\cdot\eta_k^4=0 .
\]
Therefore $F(d^{*}\xi)=0$ for $\mu$-almost every $F$:
\begin{equation}\tag{L2}\label{eq:continuum-limit-closed}
\mu\left(\left\{F\in\mathcal S':F(d^{*}\xi)=0\right\}\right)=1\qquad\left(\xi\in C^{\infty}_{\mathrm c}(\R^3)\text{ real}\right).
\end{equation}
\item[(6)]\itemmark{6}\label{it:admissible-invariance}
\emph{Invariance.}
Let $f\in\mathcal S$ be real, let $g\in W$, and let $\mathfrak g\colon y\mapsto\eta_k\cdot g(y/\eta_k)$ be the isometry of $g$ at the spacing $\eta_k$.
By Lemma~\ref{lem:test-equivariance}\eqref{it:test-equivariance-torus}, $(\mathfrak g\cdot f)_{\eta_k}=(g^{-1})^{*}f_{\eta_k}$.
The pull-backs $g^{*}$ and $(g^{-1})^{*}$ map the two-cochains of $\Lambda_{\eta_k,R_k}$ into themselves,
and $\langle g^{*}\beta,\beta'\rangle_{\eta_k}=\langle\beta,(g^{-1})^{*}\beta'\rangle_{\eta_k}$ for all two-cochains $\beta$, $\beta'$ of $\Lambda_{\eta_k,R_k}$ (Definition~\ref{def:isometry-actions});
as $g^{*}(g^{-1})^{*}=(g^{-1}g)^{*}$ is the identity (Definition~\ref{def:cochains}), the map $(g^{-1})^{*}$ is orthogonal for $\langle\cdot,\cdot\rangle_{\eta_k}$. Hence
\[
\left\langle(\mathfrak g\cdot f)_{\eta_k},(g^{-1})^{*}F\right\rangle_{\eta_k}=\langle f_{\eta_k},F\rangle_{\eta_k}\qquad\left(F\in C^2(\Lambda_{\eta_k,R_k})\right).
\]
If $\mathbb{P}_{N_k,T_k}$ is invariant under $(g^{-1})^{*}$, it follows that $\mathbb{E}_k\exp\left(iF^{\sharp}(\mathfrak g\cdot f)\right)=\mathbb{E}_k\exp\left(iF^{\sharp}(f)\right)$.
By (Ad1) this is the case for every $g\in\Ocub$, for which $\mathfrak g=g$ and $\mathfrak g\cdot f=U_gf$,
and for the translation $g=\tau_{v/\eta_k}$ of $\Z^3$, $v\in\eta_k\Z^3$, for which $\mathfrak g\cdot f=\tau_vf$ and $(g^{-1})^{*}=\tau^{*}_{-v/\eta_k}$.
For $g\in\Ocub$ the identity holds for every $k$, and \eqref{eq:continuum-weak-limit} gives $\hat\mu(U_gf)=\hat\mu(f)$.

Now let $b\in\R^3$, and choose $b_k\in\eta_k\Z^3$ with $|b_k-b|_\infty\le\frac{\eta_k}2$.
Then $\mathbb{E}_k\exp\left(iF^{\sharp}(\tau_{b_k}f)\right)=\mathbb{E}_k\exp\left(iF^{\sharp}(f)\right)$, by what has just been shown.
Since $|\exp(ix)-\exp(iy)|\le|x-y|$ for real $x$ and $y$, the Cauchy--Schwarz inequality, step~\eqref{it:admissible-bounds} together with \cite[Proposition~7.10(2)]{NoncompactQED}, and Lemma~\ref{lem:test-norms}\eqref{it:test-norms-torus} give
\[
\begin{aligned}
\left|\mathbb{E}_k\exp\left(iF^{\sharp}(\tau_bf)\right)-\mathbb{E}_k\exp\left(iF^{\sharp}(\tau_{b_k}f)\right)\right|
&\le\left(\mathbb{E}_k\left[F^{\sharp}(\tau_bf-\tau_{b_k}f)^2\right]\right)^{\frac12}\\
&\le\left|(\tau_bf-\tau_{b_k}f)_{\eta_k}\right|_{\eta_k}\le p_1(\tau_bf-\tau_{b_k}f).
\end{aligned}
\]
The map $b'\mapsto\tau_{b'}f$ is continuous from $\R^3$ into $\mathcal S$, $p_1$ is a continuous seminorm on $\mathcal S$ (Lemma~\ref{lem:test-norms}\eqref{it:test-norms-torus}), and $b_k\to b$ because $\eta_k\to0$;
so the right side tends to $0$, and \eqref{eq:continuum-weak-limit} gives
\[
\hat\mu(\tau_bf)=\lim_{k\to\infty}\mathbb{E}_k\exp\left(iF^{\sharp}(\tau_bf)\right)=\lim_{k\to\infty}\mathbb{E}_k\exp\left(iF^{\sharp}(\tau_{b_k}f)\right)=\lim_{k\to\infty}\mathbb{E}_k\exp\left(iF^{\sharp}(f)\right)=\hat\mu(f).
\]
Altogether,
\begin{equation}\tag{L3}\label{eq:continuum-limit-invariance}
\hat\mu(\tau_bf)=\hat\mu(f)=\hat\mu(U_gf)\qquad\left(f\in\mathcal S\text{ real},\ b\in\R^3,\ g\in\Ocub\right).
\end{equation}
\item[(7)]\itemmark{7}\label{it:admissible-positivity}
\emph{Reflection positivity.}
Let $\nu\in\{1,2,3\}$, let $s\ge1$ be an integer and $a_1,\dots,a_s\in\C$,
and let $f_1,\dots,f_s\in C^{\infty}_{\mathrm c}(\R^3;\Lambda^2)$ be real with $\supp f_p\subset\{y\in\R^3:y_\nu>0\}$ for $p=1,\dots,s$.
The supports being compact, there are numbers $\delta>0$ and $r>0$ with
\[
\supp f_p\subset\{y\in\R^3:y_\nu\ge\delta\}\cap[-r,r]^3\qquad(p=1,\dots,s).
\]
For all large $k$ one has $2\eta_k\le\delta$ and $R_k\ge2r+4\eta_k$.
For such $k$, Lemma~\ref{lem:positive-tests}, applied on the torus $\Lambda_{\eta_k,R_k}$, whose side $R_k/\eta_k$ is even, shows that the real cochains $(f_p)_{\eta_k}$ are supported on $\Sigma^{++}_\nu$,
and Lemma~\ref{lem:test-equivariance}\eqref{it:test-equivariance-torus} gives $(\mathsf{R}_\nu f_p)_{\eta_k}=\hat{\mathsf{R}}_\nu(f_p)_{\eta_k}$; hence
\[
F^{\sharp}\left(f_q-\mathsf{R}_\nu f_p\right)=\left\langle(f_q)_{\eta_k}-\hat{\mathsf{R}}_\nu(f_p)_{\eta_k},F\right\rangle_{\eta_k}\qquad(p,q=1,\dots,s).
\]
Property (Ad3$'$) of $\mathbb{P}_{N_k,T_k}$, with $\nu_0=\nu$ and $\beta_p=(f_p)_{\eta_k}$, gives
\[
\sum_{p,q=1}^{s}\overline{a_p}\cdot a_q\cdot\mathbb{E}_k\exp\left(iF^{\sharp}\left(f_q-\mathsf{R}_\nu f_p\right)\right)\ge0 .
\]
Each $f_q-\mathsf{R}_\nu f_p$ is a real test form, and by \eqref{eq:continuum-weak-limit} the corresponding expectation converges to $\hat\mu\left(f_q-\mathsf{R}_\nu f_p\right)$ as $k\to\infty$.
In the limit,
\begin{equation}\tag{L4}\label{eq:continuum-limit-positivity}
\sum_{p,q=1}^{s}\overline{a_p}\cdot a_q\cdot\hat\mu\left(f_q-\mathsf{R}_\nu f_p\right)\ge0 .
\end{equation}
\item[(8)]\itemmark{8}\label{it:admissible-even}
\emph{Evenness.}
Suppose that every measure $\mathbb{P}_{N_k,T_k}$ is invariant under $F\mapsto-F$, and let $f\in\mathcal S$ be real.
As $F^{\sharp}(f)$ is linear in $F$,
\[
\mathbb{E}_k\exp\left(iF^{\sharp}(f)\right)=\mathbb{E}_k\exp\left(-iF^{\sharp}(f)\right)=\mathbb{E}_k\exp\left(iF^{\sharp}(-f)\right)
\]
for every $k$, and \eqref{eq:continuum-weak-limit} gives $\hat\mu(-f)=\hat\mu(f)$.
\item[(9)]
\emph{The remaining properties.}
By steps~\eqref{it:admissible-domination}, \eqref{it:admissible-closed}, \eqref{it:admissible-invariance} and~\eqref{it:admissible-positivity}, $\mu$ is a probability measure on $(\mathcal S',\Sigma)$ with the four properties \eqref{eq:continuum-limit-domination}--\eqref{eq:continuum-limit-positivity};
these are the conditions bearing the same names in \cite[Section~7.5]{NoncompactQED}, where the characteristic functional $\hat\mu$ is written $\Phi_\mu$.
By \cite[Proposition~7.19]{NoncompactQED}, $\mu$ therefore has the properties (1) and (3)--(8) of \cite[Theorem~7.15]{NoncompactQED}, the properties (1), (2), (4) and (5) of \cite[Theorem~7.17]{NoncompactQED}, and all the properties of \cite[Corollary~7.18]{NoncompactQED}.
If every $\mathbb{P}_{N_k,T_k}$ is invariant under $F\mapsto-F$, then $\hat\mu(-f)=\hat\mu(f)$ for every real $f\in\mathcal S$ by step~\eqref{it:admissible-even},
and the last assertion of \cite[Proposition~7.19]{NoncompactQED} shows that $\mu$ is invariant under $F\mapsto-F$ and that its moments of odd order vanish.
Property (2) of \cite[Theorem~7.15]{NoncompactQED} is step~\eqref{it:admissible-moments}.

There remains property (3) of \cite[Theorem~7.17]{NoncompactQED}.
Let $n\ge1$, and let $f\in\mathscr{S}_n$, the space of the functions on $\R^{3n}$ with values in the complexification of $\left(\Lambda^2\right)^{\otimes n}$ whose components are Schwartz functions \cite[Section~7.6]{NoncompactQED}.
For $\boldsymbol\kappa=(\kappa_1,\dots,\kappa_n)$ let $a_{\boldsymbol\kappa}(f)=\left\langle h_{\kappa_1}\otimes\dots\otimes h_{\kappa_n},f\right\rangle$, the scalar product being that of $L^2\left(\R^{3n};\left(\Lambda^2\right)^{\otimes n}\right)$, be the Hermite coefficients of $f$,
and for $F\in\mathcal S'$ let $F^{\otimes n}$ be the continuous linear functional on $\mathscr{S}_n$ with $F^{\otimes n}(f_1\otimes\dots\otimes f_n)=\prod_{i=1}^{n}F(f_i)$.
By assertion (2) of \cite[Theorem~7.17]{NoncompactQED}, which $\mu$ has by the preceding paragraph, the series
\[
F^{\otimes n}(f)=\sum_{\boldsymbol\kappa}a_{\boldsymbol\kappa}(f)\cdot\prod_{i=1}^{n}F(h_{\kappa_i})
\]
converges absolutely for every $F\in\mathcal S'$, in particular for $F^{\sharp}$ with $F\in C^2(\Lambda_{\eta_k,R_k})$;
moreover $\int\mathrm{d}\mu(F)\,F^{\otimes n}(f)$ is the value $S_n(f)$ of the Schwinger function of $\mu$ of order $n$.
The operator $\mathsf{A}$ is symmetric on $\mathcal S$, and $\mathsf{A}h_\kappa=\lambda_\kappa\cdot h_\kappa$;
so the Hermite coefficients of the tensor test $\mathsf{A}^{\otimes n}f$, obtained from $f$ by applying $\mathsf{A}$ in each of its $n$ variables, are the numbers $\prod_i\lambda_{\kappa_i}\cdot a_{\boldsymbol\kappa}(f)$.
Applied to $\mathsf{A}^{\otimes n}f$, \cite[Lemma~7.16]{NoncompactQED} gives
\[
\sum_{\boldsymbol\kappa}\left|a_{\boldsymbol\kappa}(f)\right|\cdot\prod_{i=1}^{n}\lambda_{\kappa_i}\le\mathsf{Z}_4^{\frac n2}\cdot\left\|\left(\mathsf{A}^3\right)^{\otimes n}f\right\|_2 .
\]
By step~\eqref{it:admissible-bounds}, \cite[Proposition~7.10(4)]{NoncompactQED} applies under $\mathbb{P}_{N_k,T_k}$ to the real test forms $h_{\kappa_1},\dots,h_{\kappa_n}$,
for which $|(h_{\kappa_i})_{\eta_k}|_{\eta_k}\le31\lambda_{\kappa_i}$ by Lemma~\ref{lem:test-norms}\eqref{it:test-norms-torus}.
For real test forms the constant of that assertion is the number $2\cdot2^{\frac n2}\cdot\Gamma\left(\frac n2+1\right)$, $\Gamma$ being the Gamma function;
it is denoted $\sigma_n$ in \cite[Proposition~7.10]{NoncompactQED}, which also shows that it does not exceed $2\left(2^n\cdot n!\right)^{\frac12}$.
Thus
\begin{equation}\label{eq:continuum-hermite-moments}
\mathbb{E}_k\prod_{i=1}^{n}\left|F^{\sharp}(h_{\kappa_i})\right|\le2\cdot2^{\frac n2}\cdot\Gamma\left(\frac n2+1\right)\cdot\prod_{i=1}^{n}31\lambda_{\kappa_i}
\end{equation}
for every $k$ and every multi-index $\boldsymbol\kappa$.
With Tonelli's theorem and the inequality for the Hermite coefficients it follows that
\[
\mathbb{E}_k\sum_{\boldsymbol\kappa}\left|a_{\boldsymbol\kappa}(f)\right|\cdot\prod_{i=1}^{n}\left|F^{\sharp}(h_{\kappa_i})\right|
\le2\cdot2^{\frac n2}\cdot\Gamma\left(\frac n2+1\right)\cdot31^n\cdot\mathsf{Z}_4^{\frac n2}\cdot\left\|\left(\mathsf{A}^3\right)^{\otimes n}f\right\|_2<\infty .
\]
Since the series for $\left(F^{\sharp}\right)^{\otimes n}(f)$ converges absolutely, the modulus of $\left(F^{\sharp}\right)^{\otimes n}(f)$ does not exceed the sum under the expectation on the left.
Therefore the function $F\mapsto\left(F^{\sharp}\right)^{\otimes n}(f)$ is integrable with respect to $\mathbb{P}_{N_k,T_k}$, its expectation satisfies
\[
\left|\mathbb{E}_k\left[\left(F^{\sharp}\right)^{\otimes n}(f)\right]\right|\le2\cdot2^{\frac n2}\cdot\Gamma\left(\frac n2+1\right)\cdot31^n\cdot\mathsf{Z}_4^{\frac n2}\cdot\left\|\left(\mathsf{A}^3\right)^{\otimes n}f\right\|_2 ,
\]
which is the bound of \cite[Theorem~7.17(3)]{NoncompactQED},
and Fubini's theorem gives
\[
\mathbb{E}_k\left[\left(F^{\sharp}\right)^{\otimes n}(f)\right]=\sum_{\boldsymbol\kappa}a_{\boldsymbol\kappa}(f)\cdot\mathbb{E}_k\prod_{i=1}^{n}F^{\sharp}(h_{\kappa_i}).
\]
For every $\boldsymbol\kappa$, step~\eqref{it:admissible-moments} gives $\mathbb{E}_k\prod_{i}F^{\sharp}(h_{\kappa_i})\to\int\mathrm{d}\mu(F)\,\prod_{i}F(h_{\kappa_i})$ as $k\to\infty$;
and by \eqref{eq:continuum-hermite-moments} the term of index $\boldsymbol\kappa$ of the last series has, for every $k$, modulus at most
\[
2\cdot2^{\frac n2}\cdot\Gamma\left(\frac n2+1\right)\cdot\left|a_{\boldsymbol\kappa}(f)\right|\cdot\prod_{i=1}^{n}31\lambda_{\kappa_i},
\]
a number that is summable over $\boldsymbol\kappa$ by the inequality for the Hermite coefficients.
Under $\mu$, \eqref{eq:continuum-limit-domination} is the hypothesis of \cite[Proposition~7.10]{NoncompactQED} for the real test forms with the norm $\|\cdot\|_2$ and the map $f\mapsto F(f)$;
as $\|h_\kappa\|_2=1$, assertion (4) of that proposition gives $\int\mathrm{d}\mu(F)\,\prod_{i=1}^{n}\left|F(h_{\kappa_i})\right|\le2\cdot2^{\frac n2}\cdot\Gamma\left(\frac n2+1\right)$ for every $\boldsymbol\kappa$.
Since $\lambda_\kappa\ge4$, the inequality for the Hermite coefficients gives $\sum_{\boldsymbol\kappa}\left|a_{\boldsymbol\kappa}(f)\right|<\infty$, and Tonelli's theorem yields
\[
\int\mathrm{d}\mu(F)\,\sum_{\boldsymbol\kappa}\left|a_{\boldsymbol\kappa}(f)\right|\cdot\prod_{i=1}^{n}\left|F(h_{\kappa_i})\right|\le2\cdot2^{\frac n2}\cdot\Gamma\left(\frac n2+1\right)\cdot\sum_{\boldsymbol\kappa}\left|a_{\boldsymbol\kappa}(f)\right|<\infty ;
\]
so Fubini's theorem applies under $\mu$ as well.
Dominated convergence for series now gives
\[
\lim_{k\to\infty}\mathbb{E}_k\left[\left(F^{\sharp}\right)^{\otimes n}(f)\right]=\sum_{\boldsymbol\kappa}a_{\boldsymbol\kappa}(f)\cdot\int\mathrm{d}\mu(F)\,\prod_{i=1}^{n}F(h_{\kappa_i})=\int\mathrm{d}\mu(F)\,F^{\otimes n}(f)=S_n(f).
\]
\end{enumerate}
\end{proof}

\begin{proposition}[Limit points of the chamber cutoff]\label{prop:chamber-limits}
Let $0<\rho<1$ and $v_0\in C_0$, let the weights be standard, let $L\ge3$ be an odd integer, and let $0<\bar m\le1$ and $e>0$.
Let $N\mapsto m_{0,N}$ be a mass assignment.
For every cutoff $(N,T)$, that is, for every pair of integers $N,T\ge0$,
let $P^{\mathrm{ch}}_{N,T}$ be the law on $\mathcal{S}'$ that Definition~\ref{def:cutoffs} forms with these data.
\begin{enumerate}
\item\label{it:chamber-tightness}
One has
\[
\sup_{N,T\ge0}\int\mathrm{d}P^{\mathrm{ch}}_{N,T}(F)\,\|F\|_{-3}^{2}\le61 .
\]
The family $\left(P^{\mathrm{ch}}_{N,T}\right)$ is tight as a family of laws on $H_{-4}$,
and as a family of laws on $\mathcal{S}'$ for each of the two topologies of $\mathcal{S}'$, the strong one and the weak-$*$ one;
every sequence $(N_i,T_i)_{i\ge1}$ of cutoffs has a subsequence along which $P^{\mathrm{ch}}_{N_i,T_i}$ converges weakly on $H_{-4}$.
\item\label{it:chamber-properties}
Every limit point $\mu$ of $\left(P^{\mathrm{ch}}_{N,T}\right)$ in the sense of Definition~\ref{def:limit-points} has the following properties,
in the form in which \cite[Theorem~3.1, Part~I\,(b)]{NoncompactQED} states them for the limit points of the families of laws considered there,
the family $\left(P^{\mathrm{ch}}_{N,T}\right)$ taking the place of those families.
\begin{itemize}
\item
The Laplace transform of $\mu$ is bounded by that of the free Maxwell law,
with the consequences stated there for the field smeared with elements of $L^2(\R^3;\Lambda^2)$ and for its generating functions,
the free Maxwell law being the centred Gaussian law on $\mathcal{S}'$ so named in \cite[Lemma~8.30]{NoncompactQED}.
\item
The moments and the joint cumulants of $\mu$ obey the bounds stated there,
and the moments obey the linear growth condition of Osterwalder and Schrader.
\item
The moments are symmetric in their arguments and real when all arguments are real, and $\mu$ is invariant under the map $F\mapsto-F$ of $\mathcal{S}'$.
\item
The law $\mu$ is determined by its moments;
one has $\mu\left(H_{-t}\right)=1$ for every $t>\frac{3}{2}$, and
\[
\mu\left(\left\{F\in\mathcal{S}':dF=0\right\}\right)=1 .
\]
\item
The law $\mu$ is invariant under the translations of $\R^3$, which act on $L^2(\mu)$ by a strongly continuous unitary group, and under $\Ocub$.
\item
The law $\mu$ is reflection positive at every plane $\{x_j=c\}$, $j\in\{1,2,3\}$, $c\in\R$.
\item
In each coordinate direction the reconstruction of Definition~\ref{def:laws} applies to $\mu$.
\item
Let $(N_k,T_k)_{k\ge1}$ be a sequence of cutoffs with $N_k\to\infty$, $T_k\to\infty$ and $P^{\mathrm{ch}}_{N_k,T_k}\to\mu$ weakly on $H_{-4}$,
and let $f_1,\dots,f_s$ be test forms, real or complex.
Then
\[
\int\mathrm{d}P^{\mathrm{ch}}_{N_k,T_k}(F)\,\exp\left(\sum_{p=1}^{s}z_p\cdot F(f_p)\right)
\longrightarrow
\int\mathrm{d}\mu(F)\,\exp\left(\sum_{p=1}^{s}z_p\cdot F(f_p)\right)
\qquad(k\to\infty),
\]
uniformly in $z=(z_1,\dots,z_s)$ on every compact subset of $\C^s$,
and every moment and every joint cumulant of $F(f_1),\dots,F(f_s)$ under $P^{\mathrm{ch}}_{N_k,T_k}$ converges to the corresponding moment or joint cumulant under $\mu$.
\end{itemize}
\item\label{it:chamber-components}
Every limit point $\mu$ satisfies the hypothesis~\eqref{cond:soft-hypothesis} of Definition~\ref{def:laws}.
Let $\mu=\int\mathrm{d}\nu(\xi)\,\mu_\xi$ be its ergodic decomposition under the translations of $\R^3$ \cite[Section~2.5]{NoncompactQED}.
Then for $\nu$-almost every $\xi$ the component $\mu_\xi$ has the properties that \cite[Theorem~3.1, Part~I\,(c)]{NoncompactQED} states for the components of a limit point;
among them, $\mu_\xi$ is invariant and ergodic under the translations of $\R^3$, satisfies \textup{(G)}, is carried by the closed distributions,
and is reflection positive at every plane $\{x_j=c\}$, $j\in\{1,2,3\}$, $c\in\R$.
\item\label{it:chamber-iterated}
Every volume-first limit point (Definition~\ref{def:limit-points}) is a limit point.
In the opposite order, where the spacing index tends to infinity first, let $(T_i)_{i\ge1}$ be a sequence of integers $T_i\ge0$ with $T_i\to\infty$,
suppose that for every $i$ the laws $P^{\mathrm{ch}}_{N,T_i}$ converge weakly on $H_{-4}$ along some sequence of spacing indices $N$ tending to infinity,
and suppose that the limits so obtained converge weakly on $H_{-4}$ to a law $\mu$ as $i\to\infty$;
then $\mu$ is a limit point.
\end{enumerate}
\end{proposition}

\begin{proof}
Fix the data of the statement.
For every cutoff $(N,T)$ let $\mathbb{P}_{N,T}$ be the law on $C^2(\Lambda_{\eta,R})$ that Proposition~\ref{prop:cutoff-law} forms with these data,
$\eta$ and $R$ being as in Definition~\ref{def:cutoffs},
and let $P_{N,T}$ be the law of $F^{\sharp}$ under $\mathbb{P}_{N,T}$, as in Proposition~\ref{prop:admissible-families}.
\begin{enumerate}
\item\label{it:chamber-admissible}
\emph{The family $\left(\mathbb{P}_{N,T}\right)$; assertion~\eqref{it:chamber-tightness}.}
By Proposition~\ref{prop:cutoff-law}(a), \eqref{it:cutoff-invariance}, \eqref{it:cutoff-domination} and~(e),
each $\mathbb{P}_{N,T}$ is a probability measure on $C^2(\Lambda_{\eta,R})$ with $\mathbb{P}_{N,T}(\ker d)=1$,
and it satisfies \textup{(Ad1)}, \textup{(Ad2)} and~\eqref{eq:continuum-admissible-even}.
These are the hypotheses of Proposition~\ref{prop:admissible-families}.
Each $\mathbb{P}_{N,T}$ is, moreover, invariant under $F\mapsto-F$, by Proposition~\ref{prop:cutoff-law}(d);
this is the hypothesis of step~\ref{it:admissible-even} of the proof of Proposition~\ref{prop:admissible-families}.
Definition~\ref{def:cutoffs} defines $P^{\mathrm{ch}}_{N,T}$ as the law of $F^{\sharp}$ under $\mathbb{P}_{N,T}$; hence
\[
P_{N,T}=P^{\mathrm{ch}}_{N,T}\qquad\text{for every cutoff }(N,T).
\]
In particular, the limit points of $\left(P^{\mathrm{ch}}_{N,T}\right)$ are the weak limits on $H_{-4}$ of $P_{N_k,T_k}$
along sequences of cutoffs with $N_k\to\infty$ and $T_k\to\infty$,
and the steps of the proof of Proposition~\ref{prop:admissible-families} apply to each of them.

Step~\ref{it:admissible-tightness} of that proof gives, at every cutoff,
\[
\int\mathrm{d}P_{N,T}(F)\,\|F\|_{-3}^{2}\le61
\qquad\text{and}\qquad
P_{N,T}\left(H_{-4}\setminus\mathcal{B}_\varrho\right)\le61\cdot\varrho^{-2}\quad(\varrho>0)
\]
for the balls $\mathcal{B}_\varrho=\left\{F\in H_{-3}:\|F\|_{-3}\le\varrho\right\}$,
which are compact in $H_{-4}$, and compact in $\mathcal{S}'$ both for the strong and for the weak-$*$ topology of $\mathcal{S}'$;
and it gives weakly convergent subsequences on $H_{-4}$.
As $P_{N,T}$ is carried by $H_{-3}$, the family $\left(P_{N,T}\right)$ is therefore tight on $H_{-4}$, and tight on $\mathcal{S}'$ for each of these two topologies.
With $P_{N,T}=P^{\mathrm{ch}}_{N,T}$ this is assertion~\eqref{it:chamber-tightness}.
\item\label{it:chamber-soft}
\emph{Limit points are soft; assertions~\eqref{it:chamber-properties} and~\eqref{it:chamber-components}.}
Let $\mu$ be a limit point of $\left(P^{\mathrm{ch}}_{N,T}\right)$,
and let $(N_k,T_k)_{k\ge1}$ be a sequence of cutoffs with $N_k\to\infty$, $T_k\to\infty$ and $P^{\mathrm{ch}}_{N_k,T_k}\to\mu$ weakly on $H_{-4}$.
Definition~\ref{def:limit-points} provides at least one such sequence, and what follows holds for each of them.
Let $\mathbb{E}_k$ be the expectation with respect to $\mathbb{P}_{N_k,T_k}$.
As in the proof of Proposition~\ref{prop:admissible-families}, $\mu$ is regarded also as a law on $\mathcal{S}'$ with its cylinder $\sigma$-algebra.

Steps~\ref{it:admissible-domination}, \ref{it:admissible-invariance} and~\ref{it:admissible-positivity} of that proof show that $\mu$ satisfies
\eqref{eq:continuum-limit-domination}, \eqref{eq:continuum-limit-invariance} and~\eqref{eq:continuum-limit-positivity}.
A law on $\mathcal{S}'$ satisfies these three conditions if and only if it satisfies the hypothesis~\eqref{cond:soft-hypothesis} \cite[Remark~7.20]{NoncompactQED};
so $\mu$ satisfies~\eqref{cond:soft-hypothesis}.
Step~\ref{it:admissible-closed} of the proof of Proposition~\ref{prop:admissible-families} shows that $\mu$ satisfies~\eqref{eq:continuum-limit-closed} as well.
For a law with all four conditions \eqref{eq:continuum-limit-domination}--\eqref{eq:continuum-limit-positivity},
\cite[Proposition~7.19(a)]{NoncompactQED} yields assertion~(4) of \cite[Theorem~7.15]{NoncompactQED};
this assertion contains the identity
\[
\mu\left(\left\{F\in\mathcal{S}':dF=0\right\}\right)=1 .
\]
By step~\ref{it:admissible-even} of the proof of Proposition~\ref{prop:admissible-families},
whose hypothesis was verified in step~\ref{it:chamber-admissible} of the present proof,
the characteristic functional
\[
\hat\mu(f)=\int\mathrm{d}\mu(F)\,\exp\left(iF(f)\right),
\]
written $\Phi_\mu$ in \cite[Section~7.5]{NoncompactQED}, satisfies $\hat\mu(-f)=\hat\mu(f)$ for every real $f\in\mathcal{S}$.
The last assertion of \cite[Proposition~7.19]{NoncompactQED} then shows that $\mu$ is invariant under the map $F\mapsto-F$ of $\mathcal{S}'$.

A law on $\mathcal{S}'$ is called \emph{soft} if it satisfies the hypothesis~\eqref{cond:soft-hypothesis},
if it is carried by the closed distributions, that is, if the set $\left\{F\in\mathcal{S}':dF=0\right\}$ has measure one,
and if it is invariant under $F\mapsto-F$ \cite[Definition~21.35]{NoncompactQED}.
Hence $\mu$ is soft.
Every soft law has all properties listed in~\eqref{it:chamber-properties} except the last one, which concerns sequences of cutoffs,
and it has the properties listed in~\eqref{it:chamber-components} \cite[Lemma~21.36]{NoncompactQED}.

It remains to establish the last property listed in~\eqref{it:chamber-properties}.
Let $f_1,\dots,f_s$ be test forms, real or complex.
By step~\ref{it:admissible-moments} of the proof of Proposition~\ref{prop:admissible-families}, the functions $z\mapsto\mathbb{E}_k\exp\left(\sum_{p=1}^{s}z_p\cdot F^{\sharp}(f_p)\right)$ converge,
uniformly on every compact subset of $\C^s$, to $z\mapsto\int\mathrm{d}\mu(F)\,\exp\left(\sum_{p=1}^{s}z_p\cdot F(f_p)\right)$,
and all moments and all joint cumulants of $F^{\sharp}(f_1),\dots,F^{\sharp}(f_s)$ under $\mathbb{P}_{N_k,T_k}$ converge to those of $F(f_1),\dots,F(f_s)$ under $\mu$.
The law of $F^{\sharp}$ under $\mathbb{P}_{N_k,T_k}$ is $P^{\mathrm{ch}}_{N_k,T_k}$ (step~\ref{it:chamber-admissible} of the present proof),
and for a complex test form $f=f'+if''$ the value $F(f)=F(f')+iF(f'')$ of Definition~\ref{def:cutoffs} is a $\Sigma$-measurable function of $F\in\mathcal{S}'$.
Consequently
\[
\mathbb{E}_k\exp\left(\sum_{p=1}^{s}z_p\cdot F^{\sharp}(f_p)\right)=\int\mathrm{d}P^{\mathrm{ch}}_{N_k,T_k}(F)\,\exp\left(\sum_{p=1}^{s}z_p\cdot F(f_p)\right)
\qquad\left(z\in\C^s\right),
\]
and the moments and joint cumulants of $F^{\sharp}(f_1),\dots,F^{\sharp}(f_s)$ under $\mathbb{P}_{N_k,T_k}$
coincide with those of $F(f_1),\dots,F(f_s)$ under $P^{\mathrm{ch}}_{N_k,T_k}$.
Since $(N_k,T_k)$ was an arbitrary sequence of cutoffs with $N_k\to\infty$, $T_k\to\infty$ and $P^{\mathrm{ch}}_{N_k,T_k}\to\mu$ weakly on $H_{-4}$,
this is the last property listed in~\eqref{it:chamber-properties},
and the assertions~\eqref{it:chamber-properties} and~\eqref{it:chamber-components} are proved.
\item\label{it:chamber-diagonal}
\emph{Iterated limits; assertion~\eqref{it:chamber-iterated}.}
The space $H_{-4}$ is a separable Hilbert space.
On a separable metric space, weak convergence of Borel probability measures is convergence with respect to the Prokhorov metric $d_{\mathrm{P}}$
\cite[Theorem~6.8]{Billingsley1999}.
Let $\mu$ be a volume-first limit point.
By Definition~\ref{def:limit-points} there are integers $N_i\ge0$, $i\ge1$, tending to infinity, and laws $\mu^{(N_i)}\in\mathcal{V}_{N_i}$
such that $\mu^{(N_i)}\to\mu$ weakly on $H_{-4}$;
and for every $i$ there is a sequence $(T_{i,j})_{j\ge1}$ of integers $T_{i,j}\ge0$ with $T_{i,j}\to\infty$
such that $P^{\mathrm{ch}}_{N_i,T_{i,j}}\to\mu^{(N_i)}$ weakly on $H_{-4}$ as $j\to\infty$.
For every $i$ choose, as in \cite[Remark~21.17]{NoncompactQED}, an index $j_i$ with
\[
T_{i,j_i}\ge i
\qquad\text{and}\qquad
d_{\mathrm{P}}\left(P^{\mathrm{ch}}_{N_i,T_{i,j_i}},\mu^{(N_i)}\right)\le\frac{1}{i};
\]
it exists, because $T_{i,j}\to\infty$ and $d_{\mathrm{P}}\left(P^{\mathrm{ch}}_{N_i,T_{i,j}},\mu^{(N_i)}\right)\to0$ as $j\to\infty$.
The triangle inequality gives
\[
d_{\mathrm{P}}\left(P^{\mathrm{ch}}_{N_i,T_{i,j_i}},\mu\right)\le\frac{1}{i}+d_{\mathrm{P}}\left(\mu^{(N_i)},\mu\right)\longrightarrow0
\qquad(i\to\infty).
\]
So $P^{\mathrm{ch}}_{N_i,T_{i,j_i}}\to\mu$ weakly on $H_{-4}$ along the cutoffs $\left(N_i,T_{i,j_i}\right)$, $i\ge1$, both of whose entries tend to infinity,
and $\mu$ is a limit point.

In the opposite order, let $(T_i)$ and $\mu$ be as in the last sentence of~\eqref{it:chamber-iterated},
and for every $i$ let $\mu'_i$ be the weak limit on $H_{-4}$ of $P^{\mathrm{ch}}_{N_{i,j},T_i}$ as $j\to\infty$,
taken along a sequence $(N_{i,j})_{j\ge1}$ of integers $N_{i,j}\ge0$ with $N_{i,j}\to\infty$.
For the same reason as above there is an index $j_i$ with $N_{i,j_i}\ge i$ and $d_{\mathrm{P}}\left(P^{\mathrm{ch}}_{N_{i,j_i},T_i},\mu'_i\right)\le\frac{1}{i}$,
and the triangle inequality, with $\mu'_i$ in the place of $\mu^{(N_i)}$, shows that $P^{\mathrm{ch}}_{N_{i,j_i},T_i}\to\mu$ weakly on $H_{-4}$.
Both entries of the cutoffs $\left(N_{i,j_i},T_i\right)$ tend to infinity; hence $\mu$ is a limit point in this case too.
\end{enumerate}
\end{proof}

\subsection{Volume limits at a fixed spacing}\label{sec:volume-limits}

In this subsection $L\ge3$ is the odd integer of Definition~\ref{def:cutoffs}, $N\ge0$ is a fixed integer, and $\eta=L^{-N}$.
The set of the plaquettes of $\Z^3$ is denoted by $(\Z^3)^{(2)}$.
The space $C^2$ of the real two-cochains on $\Z^3$, that is, of the real functions on $(\Z^3)^{(2)}$, carries the product topology;
the space $\Omega$ of the closed two-cochains (Definition~\ref{def:characters}), whose elements are written $F^{\mathrm u}$ in this subsection,
is a closed subspace of it, each of the conditions $dF^{\mathrm u}(x)=0$, $x\in\Z^3$, involving finitely many coordinates.
A plaquette $\sigma=(x;\mu\nu)$ of $\Z^3$ has the base point $x(\sigma)=x$, the centre $y_\sigma=x+\frac12\cdot(e_\mu+e_\nu)$
and the averaging cube $Q_\eta(\sigma)=\eta\cdot\left(y_\sigma+\left[-\frac12,\frac12\right)^3\right)$ of Definition~\ref{def:averaging-cubes};
for fixed $\mu<\nu$ the cubes $Q_\eta(x;\mu\nu)$, $x\in\Z^3$, form a partition of $\R^3$.
For a test form $f$ the cochain $f_{[\eta]}$ on $\Z^3$ is that of Definition~\ref{def:test-cochains}.

The Hermite two-forms $h_\kappa$ of \cite[Section~7.1.5]{NoncompactQED} are real and form an orthonormal basis of $L^2(\R^3;\Lambda^2)$;
each of them lies in $\mathcal S$, and $\mathsf A h_\kappa=\lambda_\kappa\cdot h_\kappa$ with $\lambda_\kappa\ge4$, where $\mathsf A=1+|y|^2-\Delta$ acts on each component.
For $F\in\mathcal S'$ and $s\in\R$ one has
\[
\|F\|_s^2=\sum_\kappa\lambda_\kappa^{2s}\cdot\left|F(h_\kappa)\right|^2\in[0,\infty],\qquad H_s=\left\{F\in\mathcal S':\|F\|_s<\infty\right\},
\]
and $|F(f)|\le\|F\|_{-4}\cdot\|f\|_4$ for $F\in H_{-4}$ and $f\in\mathcal S$.
The Borel $\sigma$-algebra of $H_{-4}$ is generated by the maps $F\mapsto F(f)$, $f\in\mathcal S$.
For $a>3$ put $\mathsf Z_a:=\sum_\kappa\lambda_\kappa^{-a}$; then $\mathsf Z_4<\frac{1}{16}$ \cite[Section~7.1.5, (7.4)]{NoncompactQED}.

\begin{definition}[The step map]\label{def:step-map}
Let $L\ge3$ be an odd integer, $N\ge0$ an integer, and $\eta=L^{-N}$.
\begin{enumerate}
\item[(a)]
For $F^{\mathrm u}\in\Omega$ let $\sharp_\eta F^{\mathrm u}$ be the two-form valued function on $\R^3$ whose component $\mu\nu$ is equal to $\eta^{-\frac32}\cdot F^{\mathrm u}(\sigma)$ on $Q_\eta(\sigma)$,
for every plaquette $\sigma=(x;\mu\nu)$ of $\Z^3$.
\item[(b)]
Let
\[
\Omega':=\left\{F^{\mathrm u}\in\Omega:\ \sup_{\sigma\in(\Z^3)^{(2)}}\frac{\left|F^{\mathrm u}(\sigma)\right|}{1+|x(\sigma)|_\infty}<\infty\right\},
\]
a Borel subset of $\Omega$.
Let $F^{\mathrm u}\in\Omega'$, and let $C$ be the supremum in this definition.
For $y\in Q_\eta(\sigma)$ one has $|x(\sigma)|_\infty\le|y|_\infty/\eta+1$;
so the components of $\sharp_\eta F^{\mathrm u}$ are measurable and bounded by $C\cdot\eta^{-\frac32}\cdot\left(2+|y|_\infty/\eta\right)$,
and $\sharp_\eta F^{\mathrm u}$ is identified with the element of $\mathcal S'$ whose value at $f\in\mathcal S$ is
\[
\sharp_\eta(F^{\mathrm u})(f):=\int\mathrm{d}y\,\sum_{\mu<\nu}f_{\mu\nu}(y)\cdot(\sharp_\eta F^{\mathrm u})_{\mu\nu}(y).
\]
\item[(c)]
Let
\[
\Omega^{(\eta)}:=\left\{F^{\mathrm u}\in\Omega':\ \sum_\kappa\lambda_\kappa^{-6}\cdot\left|\sharp_\eta(F^{\mathrm u})(h_\kappa)\right|^2<\infty\right\},
\]
again a Borel subset of $\Omega$.
For $F^{\mathrm u}\in\Omega^{(\eta)}$ the series in this definition is $\left\|\sharp_\eta F^{\mathrm u}\right\|_{-3}^2$, so that $\sharp_\eta F^{\mathrm u}$ lies in $H_{-3}$, a subset of $H_{-4}$.
As a map into $H_{-4}$, the step map $\sharp_\eta\colon\Omega\to H_{-4}$ sends $F^{\mathrm u}\in\Omega^{(\eta)}$ to $\sharp_\eta F^{\mathrm u}$
and every $F^{\mathrm u}\in\Omega\setminus\Omega^{(\eta)}$ to $0$; this map is Borel.
This is the meaning of $\sharp_\eta$ whenever its values are taken in $H_{-4}$, in particular in image measures $(\sharp_\eta)_*\nu$.
\item[(d)]
For every plaquette $\sigma=(x;\mu\nu)$ of $\Z^3$ fix a real form $u^\sigma\in C^\infty_{\mathrm c}(\R^3;\Lambda^2)$, the bump of $\sigma$,
whose only non-vanishing component is the component $\mu\nu$,
whose support lies in the interior of $Q_\eta(\sigma)$, and which satisfies
\[
\int\mathrm{d}y\,u^\sigma_{\mu\nu}(y)=\eta^{\frac32}.
\]
Then $(u^\sigma)_{[\eta]}=\mathbf 1_\sigma$, the indicator cochain of $\sigma$.
Define
\[
J_\eta\colon H_{-4}\to C^2,\qquad J_\eta(F)(\sigma):=F(u^\sigma)\qquad\left(\sigma\in(\Z^3)^{(2)}\right).
\]
The map $J_\eta$ is linear and continuous for the product topology of $C^2$, and
\[
J_\eta\left(\sharp_\eta F^{\mathrm u}\right)=F^{\mathrm u}\qquad\left(F^{\mathrm u}\in\Omega^{(\eta)}\right).
\]
\end{enumerate}
\end{definition}

The assertions made in Definition~\ref{def:step-map} are proved as follows.
The set $\Omega'$ is Borel,
since $F^{\mathrm u}\mapsto\sup_\sigma\left|F^{\mathrm u}(\sigma)\right|/(1+|x(\sigma)|_\infty)$ is the supremum of countably many continuous functions on $\Omega$.
Let $F^{\mathrm u}\in\Omega'$, and let $f$ be a test form, real or complex.
Then
\begin{equation}\label{eq:continuum-step-pairing}
\sharp_\eta(F^{\mathrm u})(f)=\int\mathrm{d}y\,\sum_{\mu<\nu}f_{\mu\nu}(y)\cdot(\sharp_\eta F^{\mathrm u})_{\mu\nu}(y)=\sum_{\sigma\in(\Z^3)^{(2)}}f_{[\eta]}(\sigma)\cdot F^{\mathrm u}(\sigma)=:\langle f_{[\eta]},F^{\mathrm u}\rangle,
\end{equation}
where the integral and the series converge absolutely.
Indeed, with $C$ as in Definition~\ref{def:step-map}(b), the integrand is bounded by
$C\cdot\eta^{-\frac32}\cdot\sum_{\mu<\nu}\left|f_{\mu\nu}(y)\right|\cdot\left(2+|y|_\infty/\eta\right)$, an integrable function, $f$ being a Schwartz form.
For fixed $\mu<\nu$ the cubes $Q_\eta(x;\mu\nu)$, $x\in\Z^3$, partition $\R^3$, and $(\sharp_\eta F^{\mathrm u})_{\mu\nu}$ is constant on each of them;
so the integral is the sum over the plaquettes $\sigma=(x;\mu\nu)$ of
\[
\eta^{-\frac32}\cdot F^{\mathrm u}(\sigma)\cdot\int_{Q_\eta(\sigma)}\mathrm{d}y\,f_{\mu\nu}(y)=f_{[\eta]}(\sigma)\cdot F^{\mathrm u}(\sigma),
\]
and $\sum_\sigma\left|f_{[\eta]}(\sigma)\right|\cdot\left|F^{\mathrm u}(\sigma)\right|\le C\cdot\sum_\sigma\left|f_{[\eta]}(\sigma)\right|\cdot(1+|x(\sigma)|_\infty)<\infty$
by Lemma~\ref{lem:test-norms}(a).

For $f\in\mathcal S$ the map $F^{\mathrm u}\mapsto\langle f_{[\eta]},F^{\mathrm u}\rangle$ is a Borel function on $\Omega'$:
on $\Omega'$ it is the pointwise limit, as $k\to\infty$, of the sums over the plaquettes $\sigma$ with $|x(\sigma)|_\infty\le k$,
each of which is a continuous function on $\Omega$ of finitely many coordinates.
Consequently $\Omega^{(\eta)}$ is a Borel subset of $\Omega$,
the series in its definition being a sum of countably many non-negative Borel functions on the Borel set $\Omega'$.
The map $\sharp_\eta\colon\Omega\to H_{-4}$ of Definition~\ref{def:step-map}(c) is Borel.
In fact, for $f\in\mathcal S$ the function $F^{\mathrm u}\mapsto(\sharp_\eta F^{\mathrm u})(f)$, with $\sharp_\eta F^{\mathrm u}\in H_{-4}$,
equals $\langle f_{[\eta]},F^{\mathrm u}\rangle$ on $\Omega^{(\eta)}$ and $0$ on $\Omega\setminus\Omega^{(\eta)}$, so that it is Borel on $\Omega$;
and the Borel $\sigma$-algebra of $H_{-4}$ is generated by the maps $F\mapsto F(f)$, $f\in\mathcal S$ \cite[Section~7.1.5]{NoncompactQED}.

Let $\sigma=(x;\mu\nu)$ and $\sigma'=(x';\mu'\nu')$ be plaquettes.
If $\mu'\nu'\ne\mu\nu$, the component $\mu'\nu'$ of $u^\sigma$ vanishes;
if $\mu'\nu'=\mu\nu$ and $x'\ne x$, the cube $Q_\eta(\sigma')$ is disjoint from $Q_\eta(\sigma)$, which contains the support of $u^\sigma$.
In both cases $(u^\sigma)_{[\eta]}(\sigma')=0$, while $(u^\sigma)_{[\eta]}(\sigma)=\eta^{-\frac32}\cdot\int\mathrm{d}y\,u^\sigma_{\mu\nu}(y)=1$;
hence $(u^\sigma)_{[\eta]}=\mathbf 1_\sigma$.
The map $J_\eta$ is linear, and it is continuous for the product topology of $C^2$,
each coordinate $F\mapsto F(u^\sigma)$ being a continuous linear form on $H_{-4}$.
For $F^{\mathrm u}\in\Omega^{(\eta)}$ and every plaquette $\sigma$, \eqref{eq:continuum-step-pairing} gives
$J_\eta(\sharp_\eta F^{\mathrm u})(\sigma)=\langle(u^\sigma)_{[\eta]},F^{\mathrm u}\rangle=\langle\mathbf 1_\sigma,F^{\mathrm u}\rangle=F^{\mathrm u}(\sigma)$,
so that $J_\eta(\sharp_\eta F^{\mathrm u})=F^{\mathrm u}$.

\begin{proposition}[Volume limits and lattice limits]\label{prop:volume-limits}
Let $L\ge3$ be an odd integer, $N\ge0$ an integer and $\eta=L^{-N}$,
and let $\sharp_\eta$, $\Omega'$, $\Omega^{(\eta)}$ and $J_\eta$ be as in Definition~\ref{def:step-map}.
\begin{enumerate}
\item[(a)]
Let $\nu$ be a Borel probability measure on $\Omega$ that satisfies \eqref{cond:gaussian-domination}.
Then $\nu(\Omega^{(\eta)})=1$; for every real test form $f\in\mathcal S$,
\[
\int\mathrm{d}\nu(F^{\mathrm u})\,\exp\langle f_{[\eta]},F^{\mathrm u}\rangle\le\exp\left(\frac{\|f_{[\eta]}\|^2}{2}\right)\le\exp\left(\frac{\|f\|_2^2}{2}\right);
\]
and
\[
\int\mathrm{d}\nu(F^{\mathrm u})\,\left\|\sharp_\eta F^{\mathrm u}\right\|_{-3}^2\le\mathsf Z_6<\frac{1}{16}<61,
\]
where $\mathsf Z_6=\sum_\kappa\lambda_\kappa^{-6}\le\frac{1}{16}\cdot\mathsf Z_4$, because $\lambda_\kappa\ge4$.
\item[(b)]
Let $0<\bar m\le1$, the Wilson parameter $\rho\in(0,1)$, the point $v_0\in C_0$, the standard weights, a charge $e>0$ and a mass assignment be given,
and let $P^{\mathrm{ch}}_{N,T}$ be the laws of Definition~\ref{def:cutoffs}.
For an integer $T\ge0$ let $n=8\cdot L^{\lfloor\log_L\bar m^{-1}\rfloor}\cdot L^{T+N}$,
and let $\mu_{C(n)}$ be the law on $\Omega$ of the cubic torus $C(n)=C(n,n,n)$ from which $P^{\mathrm{ch}}_{N,T}$ is formed.
Then
\[
P^{\mathrm{ch}}_{N,T}=(\sharp_\eta)_*\mu_{C(n)}\qquad\text{and}\qquad\mu_{C(n)}=(J_\eta)_*P^{\mathrm{ch}}_{N,T}.
\]
Let $(T_j)$ be a sequence of integers with $T_j\to\infty$ such that $P^{\mathrm{ch}}_{N,T_j}\to\mu^{(N)}$ weakly on $H_{-4}$,
so that $\mu^{(N)}\in\mathcal V_N$ (Definition~\ref{def:limit-points}).
Then the laws $\mu_{C(n_j)}$, $n_j=8\cdot L^{\lfloor\log_L\bar m^{-1}\rfloor}\cdot L^{T_j+N}$, converge weakly on $\Omega$ to
\[
\nu^{(N)}:=(J_\eta)_*\mu^{(N)};
\]
moreover $\mu^{(N)}=(\sharp_\eta)_*\nu^{(N)}$, and $\nu^{(N)}$ satisfies \eqref{cond:gaussian-domination}.
\item[(c)]
In the situation of (b), $\nu^{(N)}$ satisfies \eqref{cond:gaussian-domination} and $\nu^{(N)}\left(\left\{F^{\mathrm u}:dF^{\mathrm u}=0\right\}\right)=1$,
it is invariant under $W$ and under $F^{\mathrm u}\mapsto-F^{\mathrm u}$,
and it is reflection positive (Definition~\ref{def:reflection-positivity}) at every mirror of the nine families (Definition~\ref{def:mirrors}).
Consequently $\mu^{(N)}$ is invariant under the isometries $\mathfrak g=\eta g\eta^{-1}$ of $\R^3$, $g\in W$ (Definition~\ref{def:isometry-actions}),
that is, under $\eta\Z^3\rtimes\Ocub$, and under $F\mapsto-F$.
Moreover, let $\ell\in\mathcal L$, $c\in\eta\Adm(\ell)$ and $\delta>2\eta$,
let $r_{\ell,c}$ be the reflection of $\R^3$ in the plane $\{\ell=c\}$, and let $r_{\ell,c}\cdot f$ be the image of a test form $f$ under it (Definition~\ref{def:isometry-actions}).
Then for every integer $s\ge1$, all real $f_1,\dots,f_s\in C^\infty_{\mathrm c}(\R^3;\Lambda^2)$ with supports in $\{\ell\ge c+\delta\}$ and all $a_1,\dots,a_s\in\C$,
\[
\sum_{p,q=1}^{s}\overline{a_p}\cdot a_q\cdot\int\mathrm{d}\mu^{(N)}(F)\,\exp\left(i\left(F(f_q)-F(r_{\ell,c}\cdot f_p)\right)\right)\ge0.
\]
\end{enumerate}
\end{proposition}

\begin{proof}
(a) Let $\sigma$ be a plaquette of $\Z^3$.
For real $t$ the cochain $t\cdot\mathbf 1_\sigma$ is finitely supported, with $\|t\cdot\mathbf 1_\sigma\|=|t|$, and \eqref{cond:gaussian-domination} gives
\[
\int\mathrm{d}\nu(F^{\mathrm u})\,\exp\left(t\cdot F^{\mathrm u}(\sigma)\right)\le\exp\left(\frac{t^2}{2}\right).
\]
This is the hypothesis of \cite[Proposition~7.10]{NoncompactQED} for the probability space $(\Omega,\nu)$,
the line $\R\cdot\mathbf 1_\sigma$ with the Hilbert seminorm $|t\cdot\mathbf 1_\sigma|:=|t|$,
and the linear map that sends $t\cdot\mathbf 1_\sigma$ to the random variable $F^{\mathrm u}\mapsto t\cdot F^{\mathrm u}(\sigma)$.
Clause (3) of that proposition, the Chernoff bound, gives
\[
\nu\left(\left\{F^{\mathrm u}:\left|F^{\mathrm u}(\sigma)\right|\ge t\right\}\right)\le2\exp\left(-\frac{t^2}{2}\right)\qquad(t>0).
\]
For each of the three orientations and each integer $k\ge0$ there are at most $(2k+1)^3$ plaquettes $\sigma$ with $|x(\sigma)|_\infty=k$; hence
\[
\sum_{\sigma\in(\Z^3)^{(2)}}2\exp\left(-\frac{(1+|x(\sigma)|_\infty)^2}{2}\right)\le6\cdot\sum_{k\ge0}(2k+1)^3\cdot\exp\left(-\frac{(k+1)^2}{2}\right)<\infty.
\]
By the Borel--Cantelli lemma, for $\nu$-almost every $F^{\mathrm u}$ the inequality $\left|F^{\mathrm u}(\sigma)\right|<1+|x(\sigma)|_\infty$ holds for all but finitely many plaquettes $\sigma$;
the supremum in the definition of $\Omega'$ is then finite.
So $\nu(\Omega')=1$.

Let now $f$ be a real test form, and put $\beta:=f_{[\eta]}$, a real two-cochain on $\Z^3$, the coefficients of $f$ being real-valued.
For $r>0$ let $\beta_r\in C^2_{\mathrm c}$ be equal to $\beta$ on the finitely many plaquettes $\sigma$ with $|x(\sigma)|_\infty\le r$ and to $0$ on the others.
For $F^{\mathrm u}\in\Omega'$ the series $\langle\beta,F^{\mathrm u}\rangle$ converges absolutely by \eqref{eq:continuum-step-pairing},
so that $\langle\beta_r,F^{\mathrm u}\rangle\to\langle\beta,F^{\mathrm u}\rangle$ as $r\to\infty$.
Fatou's lemma and \eqref{cond:gaussian-domination} for the cochains $\beta_r$ give
\[
\int\mathrm{d}\nu(F^{\mathrm u})\,\exp\langle\beta,F^{\mathrm u}\rangle
\le\liminf_{r\to\infty}\int\mathrm{d}\nu(F^{\mathrm u})\,\exp\langle\beta_r,F^{\mathrm u}\rangle
\le\lim_{r\to\infty}\exp\left(\frac{\|\beta_r\|^2}{2}\right)=\exp\left(\frac{\|\beta\|^2}{2}\right),
\]
and $\|\beta\|=\|f_{[\eta]}\|\le\|f\|_2$ by Lemma~\ref{lem:test-norms}(a).
This proves the exponential bound of (a).

That bound is the hypothesis of \cite[Proposition~7.10]{NoncompactQED} for the probability space $(\Omega,\nu)$,
the real vector space of the real test forms with the Hilbert seminorm $f\mapsto\|f_{[\eta]}\|$,
which belongs to the positive semidefinite symmetric bilinear form $(f_1,f_2)\mapsto\langle(f_1)_{[\eta]},(f_2)_{[\eta]}\rangle$,
and the linear map that sends $f$ to the random variable equal to $\langle f_{[\eta]},F^{\mathrm u}\rangle$ on $\Omega'$ and to $0$ on $\Omega\setminus\Omega'$.
Clause (2) of that proposition gives
\[
\int\mathrm{d}\nu(F^{\mathrm u})\,\langle f_{[\eta]},F^{\mathrm u}\rangle^2\le\|f_{[\eta]}\|^2\le\|f\|_2^2.
\]
The Hermite two-forms are real elements of $\mathcal S$ with $\|h_\kappa\|_2=1$,
and $\sharp_\eta(F^{\mathrm u})(h_\kappa)=\langle(h_\kappa)_{[\eta]},F^{\mathrm u}\rangle$ for $F^{\mathrm u}\in\Omega'$, by \eqref{eq:continuum-step-pairing}.
By Tonelli's theorem, the integrand on the left being defined on the set $\Omega'$ of full measure,
\[
\int\mathrm{d}\nu(F^{\mathrm u})\,\sum_\kappa\lambda_\kappa^{-6}\cdot\left|\sharp_\eta(F^{\mathrm u})(h_\kappa)\right|^2
=\sum_\kappa\lambda_\kappa^{-6}\cdot\int\mathrm{d}\nu(F^{\mathrm u})\,\langle(h_\kappa)_{[\eta]},F^{\mathrm u}\rangle^2
\le\sum_\kappa\lambda_\kappa^{-6}=\mathsf Z_6 .
\]
This is the computation of the second moments in \cite[Theorem~7.12(1)]{NoncompactQED}, display (7.6) there;
on $\Z^3$ there is no periodisation, and the bound $\|(h_\kappa)_{[\eta]}\|\le\|h_\kappa\|_2=1$
takes the place of the bound $|(h_\kappa)_\eta|_\eta\le p_1(h_\kappa)$ of Lemma~\ref{lem:test-norms}(c).
Since $\lambda_\kappa^{-6}\le\frac{1}{16}\cdot\lambda_\kappa^{-4}$, one has $\mathsf Z_6\le\frac{1}{16}\cdot\mathsf Z_4<\frac{1}{16}$.
In particular the series $\sum_\kappa\lambda_\kappa^{-6}\cdot\left|\sharp_\eta(F^{\mathrm u})(h_\kappa)\right|^2$ is finite for $\nu$-almost every $F^{\mathrm u}$, that is, $\nu(\Omega^{(\eta)})=1$.
On $\Omega^{(\eta)}$ this series equals $\left\|\sharp_\eta F^{\mathrm u}\right\|_{-3}^2$, and $\Omega\setminus\Omega^{(\eta)}$ is a $\nu$-null set;
hence $\int\mathrm{d}\nu(F^{\mathrm u})\,\left\|\sharp_\eta F^{\mathrm u}\right\|_{-3}^2$ equals the left side of the last display, and it is at most $\mathsf Z_6$.

(b) The spacing $\eta=L^{-N}$ is fixed,
and the cubic torus at the cutoff $(N,T)$ has the side $n=8\cdot L^{\lfloor\log_L\bar m^{-1}\rfloor}\cdot L^{T+N}$ (Definition~\ref{def:cutoffs}).
\begin{enumerate}
\item[(1)]\itemmark{1}\label{it:volume-identities}
\emph{The two identities.}
Let $T\ge0$, let $F^{\mathrm u}$ be the curvature of a configuration on $C(n)$, a closed two-cochain on $\Z^3/n\Z^3$,
and let $F^{\sharp}$ be the step two-form of $\eta^{-\frac32}\cdot F^{\mathrm u}$ (Definition~\ref{def:cutoffs}).
Its periodic extension $\widetilde F^{\mathrm u}$ is closed and takes only finitely many values;
being bounded, it lies in $\Omega'$.
For every $f\in\mathcal S$, Lemma~\ref{lem:test-unfolding} and \eqref{eq:continuum-step-pairing} give
\[
F^{\sharp}(f)=\sum_{\sigma\in(\Z^3)^{(2)}}f_{[\eta]}(\sigma)\cdot\widetilde F^{\mathrm u}(\sigma)=\left(\sharp_\eta\widetilde F^{\mathrm u}\right)(f),
\]
so that $F^{\sharp}=\sharp_\eta\widetilde F^{\mathrm u}$ as elements of $\mathcal S'$.
The two-form $F^{\sharp}$ lies in $H_{-3}$ (Step~\ref{it:admissible-tightness} of the proof of Proposition~\ref{prop:admissible-families}), that is,
\[
\sum_{\kappa}\lambda_\kappa^{-6}\cdot\left|\left(\sharp_\eta\widetilde F^{\mathrm u}\right)(h_\kappa)\right|^2=\left\|F^{\sharp}\right\|_{-3}^2<\infty ,
\]
and therefore $\widetilde F^{\mathrm u}\in\Omega^{(\eta)}$.
Now $\mu_{C(n)}$ is the law of $\widetilde F^{\mathrm u}$ under $\mathbb E_{C(n)}$ (Definition~\ref{def:functional}),
and $P^{\mathrm{ch}}_{N,T}$ is the law of $F^{\sharp}=\sharp_\eta\widetilde F^{\mathrm u}$ (Definition~\ref{def:cutoffs}).
Hence
\[
P^{\mathrm{ch}}_{N,T}=\left(\sharp_\eta\right)_*\mu_{C(n)},
\qquad
\mu_{C(n)}\left(\Omega^{(\eta)}\right)=1 .
\]
As $J_\eta\circ\sharp_\eta$ is the identity of $\Omega^{(\eta)}$ (Definition~\ref{def:step-map}), the second identity follows:
\[
\left(J_\eta\right)_*P^{\mathrm{ch}}_{N,T}=\left(J_\eta\circ\sharp_\eta\right)_*\mu_{C(n)}=\mu_{C(n)} .
\]
\item[(2)]\itemmark{2}\label{it:volume-convergence}
\emph{Convergence on $\Omega$.}
Let $T_j\to\infty$ be a sequence along which $P^{\mathrm{ch}}_{N,T_j}\to\mu^{(N)}$ weakly on $H_{-4}$, and let $n_j$ be the side of the cubic torus at $(N,T_j)$.
The map $J_\eta\colon H_{-4}\to C^2$ is continuous for the product topology of $C^2$ (Definition~\ref{def:step-map}).
The first step and the mapping theorem \cite[Theorem~2.7]{Billingsley1999} give
\[
\mu_{C(n_j)}=\left(J_\eta\right)_*P^{\mathrm{ch}}_{N,T_j}\longrightarrow\left(J_\eta\right)_*\mu^{(N)}=\nu^{(N)}
\qquad\text{weakly on }C^2\text{ as }j\to\infty .
\]
The set $\Omega$ is closed in $C^2$ (Definition~\ref{def:characters}), and $\mu_{C(n_j)}(\Omega)=1$ for every $j$.
By the portmanteau theorem \cite[Theorem~2.1]{Billingsley1999},
$\nu^{(N)}(\Omega)\ge\limsup_{j}\mu_{C(n_j)}(\Omega)=1$, so $\nu^{(N)}$ is a law on $\Omega$.
Every closed subset of $\Omega$ is closed in $C^2$,
so its $\nu^{(N)}$-measure is at least the upper limit of its $\mu_{C(n_j)}$-measures;
the portmanteau theorem on $\Omega$ now shows that $\mu_{C(n_j)}\to\nu^{(N)}$ weakly on $\Omega$.
\item[(3)]\itemmark{3}\label{it:volume-domination}
\emph{Gaussian domination of $\nu^{(N)}$.}
Let $\beta\in C^2_{\mathrm c}$, and let $B=y+[-r,r]^3$, with $y\in\R^3$ and $r>0$, be a box that contains the support of $\beta$.
If $n_j>2r+2$, the cubes $B+[-1,1]^3+v$, $v\in n_j\Z^3$, of side $2r+2$ are pairwise disjoint,
so that $B$ embeds in $n_j\Z^3$ (Definition~\ref{def:integrands});
since $n_j\to\infty$, this is the case for all large $j$.
Theorem~\ref{thm:coordinate-domination}(d), which applies because $n_j\ge8$,
shows that $\mu_{C(n_j)}$ satisfies \eqref{cond:domination-beta} for all large $j$.
By the weak convergence on $\Omega$ proved in the previous step and Lemma~\ref{lem:domination-compactness}(c),
$\nu^{(N)}$ satisfies \eqref{cond:domination-beta}.
The cochain $\beta\in C^2_{\mathrm c}$ being arbitrary, $\nu^{(N)}$ satisfies \eqref{cond:gaussian-domination}.
\item[(4)]\itemmark{4}\label{it:volume-step-functions}
\emph{The law $\mu^{(N)}$ is carried by step functions.}
By the previous step, $\nu^{(N)}$ satisfies \eqref{cond:gaussian-domination},
and part (a), applied to $\nu=\nu^{(N)}$, gives $\nu^{(N)}\left(\Omega^{(\eta)}\right)=1$.
Since $\nu^{(N)}=\left(J_\eta\right)_*\mu^{(N)}$, this says that $J_\eta(F)\in\Omega^{(\eta)}$ for $\mu^{(N)}$-almost every $F\in H_{-4}$.
Let $\mathcal S_0$ be the countable set of \cite[Lemma~8.11\,(i)]{NoncompactQED}, denoted $\mathcal D_0$ there.
Its elements are real forms in $C^{\infty}_{\mathrm c}(\R^3;\Lambda^2)$,
and it is dense, for the topology of $\mathcal S$, in the space of the real test forms, which the companion paper denotes by $\mathcal S$.
For $f\in\mathcal S_0$ the cochain $f_{[\eta]}$ has finite support, $f$ having compact support, and
\[
f':=f-\sum_{\sigma\in(\Z^3)^{(2)}}f_{[\eta]}(\sigma)\cdot u^{\sigma}
\]
is a finite sum of test forms.
Since $\left(u^{\sigma}\right)_{[\eta]}=\mathbf 1_\sigma$ for every plaquette $\sigma$, one has $\left(f'\right)_{[\eta]}=0$.
For every $T$ and every two-cochain $F^{\mathrm u}$ on $\Z^3/n\Z^3$, Lemma~\ref{lem:test-unfolding} therefore gives
\[
F^{\sharp}(f')=\sum_{\sigma\in(\Z^3)^{(2)}}\left(f'\right)_{[\eta]}(\sigma)\cdot\widetilde F^{\mathrm u}(\sigma)=0 ,
\]
and hence $P^{\mathrm{ch}}_{N,T}\left\{F:F(f')=0\right\}=1$ for all $T$.
The set $\left\{F\in H_{-4}:F(f')=0\right\}$ is closed in $H_{-4}$, the map $F\mapsto F(f')$ being continuous on $H_{-4}$;
by the portmanteau theorem, $\mu^{(N)}\left\{F:F(f')=0\right\}=1$.
As $\mathcal S_0$ is countable, there is a set of full $\mu^{(N)}$-measure on which $J_\eta(F)\in\Omega^{(\eta)}$ and $F(f')=0$ for all $f\in\mathcal S_0$.
For $F$ in this set and $f\in\mathcal S_0$, linearity, the definition of $J_\eta$ and \eqref{eq:continuum-step-pairing} give
\[
F(f)=\sum_{\sigma\in(\Z^3)^{(2)}}f_{[\eta]}(\sigma)\cdot F\left(u^{\sigma}\right)
=\sum_{\sigma\in(\Z^3)^{(2)}}f_{[\eta]}(\sigma)\cdot J_\eta(F)(\sigma)
=\left(\sharp_\eta\left(J_\eta(F)\right)\right)(f) .
\]
Both $F$ and $\sharp_\eta\left(J_\eta(F)\right)$ lie in $\mathcal S'$,
so that their values at a real test form depend continuously on the form for the topology of $\mathcal S$;
since $\mathcal S_0$ is dense in the space of the real test forms, the two agree at every real test form.
An element of $\mathcal S'$ is determined by these values,
its value at a complex test form being its value at the real part plus $i$ times its value at the imaginary part (Definition~\ref{def:cutoffs}).
Hence $F=\sharp_\eta\left(J_\eta(F)\right)$ for $\mu^{(N)}$-almost every $F$, and
\[
\mu^{(N)}=\left(\sharp_\eta\circ J_\eta\right)_*\mu^{(N)}=\left(\sharp_\eta\right)_*\nu^{(N)} .
\]
\end{enumerate}

(c) Let $T_j$, $n_j$, $\mu^{(N)}$ and $\nu^{(N)}=\left(J_\eta\right)_*\mu^{(N)}$ be as in part (b).
At the fixed spacing $\eta$ the parameters of the laws $\mu_{C(n_j)}$ do not depend on $j$,
and they satisfy the hypotheses of Theorem~\ref{thm:cubic-limits}:
$0<\rho<1$;
the standard weights $w_i$ are positive (Definition~\ref{def:chamber-operator});
$M$ is real (Definition~\ref{def:cutoffs});
and $q=e\cdot\eta^{\frac12}>0$.
The sides $n_j=8\cdot L^{\lfloor\log_L\bar m^{-1}\rfloor}\cdot L^{T_j+N}$ satisfy $n_j\ge8$ and $n_j\to\infty$.
For each $j$ the law $\mu_{C(n_j)}$ is the law $\mu_{\mathbb T}$ of Definition~\ref{def:functional} for the torus $\mathbb T=C(n_j)$,
formed with the measure of Definition~\ref{def:configurations}, in which all closed integer two-cochains $z$ on the torus are summed;
so $\left(\mu_{C(n_j)}\right)_j$ is the sequence of the laws $\mu_i$ of Theorem~\ref{thm:cubic-limits} for these parameters and the sides $n_j$.
By part (b), $\nu^{(N)}$ is the weak limit on $\Omega$ of this sequence, in particular a weak limit point of it.
Theorem~\ref{thm:cubic-limits} therefore shows that $\nu^{(N)}$ satisfies \eqref{cond:gaussian-domination} and $\nu^{(N)}\left(dF^{\mathrm u}=0\right)=1$,
that it is invariant under $W$ and under $F^{\mathrm u}\mapsto-F^{\mathrm u}$,
and that it is reflection positive (Definition~\ref{def:reflection-positivity}) at every mirror of the nine families.
In particular $\nu^{(N)}$ is carried by $\Omega^{(\eta)}\subset\Omega'$, by part (a).
Since $\mu^{(N)}=\left(\sharp_\eta\right)_*\nu^{(N)}$, the pairing \eqref{eq:continuum-step-pairing} gives, for every real $f\in\mathcal S$,
\begin{equation}\label{eq:continuum-volume-functional}
\mu^{(N)}\left(\exp\left(i\,F(f)\right)\right)=\int\mathrm{d}\nu^{(N)}(F^{\mathrm u})\,\exp\left(i\left\langle f_{[\eta]},F^{\mathrm u}\right\rangle\right).
\end{equation}

\emph{Invariance.}
Let $g\in W$, let $\mathfrak g=\eta g\eta^{-1}$ be the isometry of $g$ at the spacing $\eta$ (Definition~\ref{def:isometry-actions}), and let $f\in\mathcal S$ be real.
The pull-back $g^{*}$ permutes the coordinates of $C^2$ up to signs, $\left(g^{*}\beta\right)(\sigma)=\pm\beta(g\sigma)$ for $\beta\in C^2$,
and it maps $\Omega$ onto itself, with inverse $\left(g^{-1}\right)^{*}$ (Definition~\ref{def:cochains}).
The linear part of $g$ is a signed permutation matrix, and $x(g\sigma)$ and $g\left(x(\sigma)\right)$ are corners of the same plaquette $g\sigma$; hence
\[
\left|x(g\sigma)\right|_\infty\le\left|x(\sigma)\right|_\infty+\left|g(0)\right|_\infty+1
\qquad\text{for every plaquette }\sigma .
\]
It follows that $g^{*}$ maps $\Omega'$ into itself, and, the same holding for $g^{-1}$, onto itself.
For $F^{\mathrm u}\in\Omega'$ the same bound and Lemma~\ref{lem:test-norms}(a) show that the series
$\sum_{\sigma}f_{[\eta]}(\sigma)\cdot\left(g^{*}F^{\mathrm u}\right)(\sigma)$ converges absolutely.
It may therefore be rearranged, and since the transpose of $g^{*}$ is $\left(g^{-1}\right)^{*}$ (Definition~\ref{def:isometry-actions}),
the first of the following equalities holds; the second is Lemma~\ref{lem:test-equivariance}(a):
\[
\left\langle f_{[\eta]},g^{*}F^{\mathrm u}\right\rangle
=\left\langle\left(g^{-1}\right)^{*}f_{[\eta]},F^{\mathrm u}\right\rangle
=\left\langle\left(\mathfrak g\cdot f\right)_{[\eta]},F^{\mathrm u}\right\rangle .
\]
With \eqref{eq:continuum-volume-functional} for the real test forms $\mathfrak g\cdot f$ and $f$, and the invariance of $\nu^{(N)}$ under $g$, one obtains
\[
\begin{aligned}
\mu^{(N)}\left(\exp\left(i\,F(\mathfrak g\cdot f)\right)\right)
&=\int\mathrm{d}\nu^{(N)}(F^{\mathrm u})\,\exp\left(i\left\langle f_{[\eta]},g^{*}F^{\mathrm u}\right\rangle\right)\\
&=\int\mathrm{d}\nu^{(N)}(F^{\mathrm u})\,\exp\left(i\left\langle f_{[\eta]},F^{\mathrm u}\right\rangle\right)
=\mu^{(N)}\left(\exp\left(i\,F(f)\right)\right).
\end{aligned}
\]
Since $\left(\mathfrak g^{-1}\omega\right)(f)=\omega(\mathfrak g\cdot f)$ for $\omega\in\mathcal S'$ (Definition~\ref{def:isometry-actions}),
the left side is the characteristic functional at $f$ of the image of $\mu^{(N)}$ under $\omega\mapsto\mathfrak g^{-1}\omega$.
Two laws on $(\mathcal S',\Sigma)$ whose characteristic functionals agree on $\mathcal S_0$ are equal \cite[Lemma~8.12\,(ii)]{NoncompactQED};
so this image is $\mu^{(N)}$ itself, and $\mu^{(N)}$ is invariant under $\mathfrak g$.
As $g\in W$ is arbitrary, $\mu^{(N)}$ is invariant under the isometries at the spacing $\eta$ of all elements of $W$, that is, under $\eta\Z^3\rtimes\Ocub$.
For the change of sign, $(-f)_{[\eta]}=-f_{[\eta]}$, and the invariance of $\nu^{(N)}$ under $F^{\mathrm u}\mapsto-F^{\mathrm u}$ gives in the same way
\[
\mu^{(N)}\left(\exp\left(i\,F(-f)\right)\right)
=\int\mathrm{d}\nu^{(N)}(F^{\mathrm u})\,\exp\left(i\left\langle f_{[\eta]},-F^{\mathrm u}\right\rangle\right)
=\mu^{(N)}\left(\exp\left(i\,F(f)\right)\right).
\]
The left side is the characteristic functional at $f$ of the image of $\mu^{(N)}$ under $\omega\mapsto-\omega$;
by the same lemma of the companion paper, $\mu^{(N)}$ is invariant under $F\mapsto-F$.

\emph{Positivity.}
Let $\ell\in\mathcal L$, $c\in\eta\Adm(\ell)$ and $\delta>2\eta$,
let $f_1,\dots,f_s\in C^{\infty}_{\mathrm c}(\R^3;\Lambda^2)$ be real with supports in $\{\ell\ge c+\delta\}$,
and let $a_1,\dots,a_s\in\C$.
As $c/\eta\in\Adm(\ell)$, the reflection $\theta:=\theta_{\ell,c/\eta}$ belongs to $W$,
and its isometry at the spacing $\eta$ is $r_{\ell,c}$ (Definition~\ref{def:isometry-actions}).
Let $p\in\{1,\dots,s\}$, and let $\sigma$ be a plaquette with $(f_p)_{[\eta]}(\sigma)\neq0$.
The integral that defines $(f_p)_{[\eta]}(\sigma)$ does not vanish,
so the averaging cube $Q_\eta(\sigma)$ contains a point $y$ of the support of $f_p$, and $\ell(y)\ge c+\delta$.
Let $x$ be a corner of $\sigma$.
The point $\eta\cdot x$ differs from the centre $\eta\cdot y_\sigma$ by $\eta/2$ in two coordinates and by $0$ in the third,
while $\left|y-\eta\cdot y_\sigma\right|_\infty\le\eta/2$;
hence $\left|\eta\cdot x-y\right|_\infty\le\eta$.
Let $|n_\ell|_1$ be the sum of the moduli of the coordinates of the vector $n_\ell$;
it equals $1$ for an axis family and $2$ for a diagonal family (Definition~\ref{def:mirrors}).
As $\ell(\eta\cdot x)-\ell(y)=n_\ell\cdot(\eta\cdot x-y)$, one obtains
\[
\left|\ell(\eta\cdot x)-\ell(y)\right|\le|n_\ell|_1\cdot\left|\eta\cdot x-y\right|_\infty\le|n_\ell|_1\cdot\eta\le2\eta .
\]
Therefore $\ell(\eta\cdot x)\ge c+\delta-2\eta>c$, that is, $\ell(x)>c/\eta$.
Every corner of $\sigma$ thus has height $\ell>c/\eta$, and $\sigma\in\Sigma^{+}(\ell,c/\eta)$.
Since $p$ is arbitrary and every $f_p$ has compact support,
each $(f_p)_{[\eta]}$ is a finitely supported cochain with support in $\Sigma^{+}(\ell,c/\eta)$.
Consequently
\[
\Phi(F^{\mathrm u}):=\sum_{q=1}^{s}a_q\cdot\exp\left(i\left\langle(f_q)_{[\eta]},F^{\mathrm u}\right\rangle\right)
\]
is a bounded continuous function of finitely many coordinates $F^{\mathrm u}(\sigma)$, $\sigma\in\Sigma^{+}(\ell,c/\eta)$,
and it belongs to $\mathcal O^{+}(\ell,c/\eta)$.
The map $\theta^{*}$ is a symmetric involution of $C^2$ (Definition~\ref{def:isometry-actions}),
and $\theta^{*}(f_p)_{[\eta]}=\left(r_{\ell,c}\cdot f_p\right)_{[\eta]}$ by Lemma~\ref{lem:test-equivariance}(a); hence
\[
\overline{\Phi\left(\theta^{*}F^{\mathrm u}\right)}
=\sum_{p=1}^{s}\overline{a_p}\cdot\exp\left(-i\left\langle\left(r_{\ell,c}\cdot f_p\right)_{[\eta]},F^{\mathrm u}\right\rangle\right).
\]
Reflection positivity of $\nu^{(N)}$ at the mirror $(\ell,c/\eta)$,
followed by \eqref{eq:continuum-volume-functional} for the real test forms $f_q-r_{\ell,c}\cdot f_p$, gives
\[
\begin{aligned}
0&\le\int\mathrm{d}\nu^{(N)}(F^{\mathrm u})\,\overline{\Phi\left(\theta^{*}F^{\mathrm u}\right)}\cdot\Phi(F^{\mathrm u})\\
&=\sum_{p,q=1}^{s}\overline{a_p}\cdot a_q\cdot\int\mathrm{d}\nu^{(N)}(F^{\mathrm u})\,\exp\left(i\left\langle\left(f_q-r_{\ell,c}\cdot f_p\right)_{[\eta]},F^{\mathrm u}\right\rangle\right)\\
&=\sum_{p,q=1}^{s}\overline{a_p}\cdot a_q\cdot\mu^{(N)}\left(\exp\left(i\left(F(f_q)-F\left(r_{\ell,c}\cdot f_p\right)\right)\right)\right).
\end{aligned}
\]
\end{proof}

\subsection{Volume-first limit points}\label{sec:volume-first}

\begin{proposition}[Nine families of mirrors]\label{prop:nine-families}
Let $L\ge3$ be odd, $0<\bar m\le1$, $e>0$, $0<\rho<1$ and $v_0\in C_0$,
let the weights $w_0,\dots,w_3$ be standard, and let $N\mapsto m_{0,N}$ be a mass assignment;
let $\left(P^{\mathrm{ch}}_{N,T}\right)$ be the family of laws of Definition~\ref{def:cutoffs} formed with these data.
Every volume-first limit point $\mu$ of $\left(P^{\mathrm{ch}}_{N,T}\right)$ (Definition~\ref{def:limit-points}) is a limit point of $\left(P^{\mathrm{ch}}_{N,T}\right)$;
in particular, $\mu$ has all the properties that Proposition~\ref{prop:chamber-limits} asserts of limit points.
Moreover, $\mu$ lies in $\mathcal{C}_9$ (Definition~\ref{def:laws}):
it is invariant under every translation of $\R^3$,
and for every plane $\{x_j=c\}$ with $1\le j\le3$ and every plane $\{x_j-x_k=c\}$ or $\{x_j+x_k=c\}$ with $1\le j<k\le3$, where $c\in\R$,
the reflection in the plane leaves $\mu$ invariant,
and $\mu$ is reflection positive at the plane in the sense of \cite[Section~8.1.3]{NoncompactQED}.
\end{proposition}

\begin{proof}
By Definition~\ref{def:limit-points} there are integers $N_i\ge0$ with $N_i\to\infty$ and volume limits $\mu^{(N_i)}\in\mathcal{V}_{N_i}$ such that $\mu^{(N_i)}\to\mu$ weakly on $H_{-4}$;
each $\mu^{(N_i)}$ is the weak limit on $H_{-4}$ of the laws $P^{\mathrm{ch}}_{N_i,T}$ along a sequence of volume indices $T$ tending to infinity.
Thus $\mu$ is an iterated limit, taken first in the volume index and then in the spacing index,
and Proposition~\ref{prop:chamber-limits}(d) shows that $\mu$ is a limit point of $\left(P^{\mathrm{ch}}_{N,T}\right)$.
This is the first assertion.
In what follows $\eta_i:=L^{-N_i}$ is the spacing of $\mu^{(N_i)}$,
and $\mu$ and the $\mu^{(N_i)}$ are regarded as laws on $\mathcal{S}'$, as in Definition~\ref{def:laws}.

\emph{Invariance.}
Being a limit point, $\mu$ has the properties of limit points listed in \cite[Theorem~3.1, Part~I\,(b)]{NoncompactQED} (Proposition~\ref{prop:chamber-limits}(b)).
Among them is the invariance of $\mu$ under the translations of $\R^3$ and under the cubic group:
$(\tau_b)_*\mu=\mu$ for every $b\in\R^3$, and $g_*\mu=\mu$ for every $g\in\Ocub$.
The planes named in the statement are the planes $\{\ell=c\}$ with $\ell\in\mathcal{L}$ and $c\in\R$ (Definition~\ref{def:laws}),
and the reflection in the plane $\{\ell=c\}$ is $r_{\ell,c}$ (Definition~\ref{def:isometry-actions}).
Fix $\ell\in\mathcal{L}$.
The reflection
\[
r_{\ell,0}\colon y\mapsto y-2|n_\ell|^{-2}\cdot\ell(y)\cdot n_\ell
\]
is linear and orthogonal, and it maps $\Z^3$ into $\Z^3$, because $n_\ell\in\Z^3$, $\ell(\Z^3)\subset\Z$ and $2|n_\ell|^{-2}\in\{1,2\}$.
Its matrix is therefore orthogonal with integer entries, and so is the inverse matrix, the transpose;
hence $r_{\ell,0}$ lies in $\mathrm{O}(3)\cap\mathrm{GL}(3,\Z)$, which is the cubic group $\Ocub$.
For $c\in\R$ put $b:=2c\cdot|n_\ell|^{-2}\cdot n_\ell$; then $r_{\ell,c}=\tau_b\circ r_{\ell,0}$.
The actions of Euclidean maps on $\mathcal{S}'$ and on laws (Definition~\ref{def:laws}) are compatible with composition, so that
\[
(r_{\ell,c})_*\mu=(\tau_b)_*\left((r_{\ell,0})_*\mu\right)=\mu .
\]

\emph{Reflection positivity at a plane.}
Let $c\in\R$.
By Definition~\ref{def:laws}, reflection positivity at the plane $\{\ell=c\}$ is reflection positivity at the plane $\{\hat n_\ell\cdot x=c/|n_\ell|\}$,
whose reflection is $r_{\hat n_\ell,c/|n_\ell|}=r_{\ell,c}$, the test forms being supported in $\{\ell>c\}$.
For the law $\mu$ it is the assertion that
\begin{equation}\label{eq:continuum-plane-positivity}
\sum_{p,q=1}^{s}\overline{a_p}\cdot a_q\cdot\int\mathrm{d}\mu(F)\,\exp\left(i\left(F(f_q)-F(r_{\ell,c}\cdot f_p)\right)\right)\ge0
\end{equation}
for every integer $s\ge1$, all $a_1,\dots,a_s\in\C$ and all real $f_1,\dots,f_s\in C^{\infty}_{\mathrm{c}}(\R^3;\Lambda^2)$
with $\supp f_p\subset\{\ell>c\}$ for $p=1,\dots,s$.

\emph{Positions in $\Z[L^{-1}]$.}
Let $c$ lie in $\Z[L^{-1}]$, the union over the integers $N\ge0$ of the sets $L^{-N}\Z$,
which increase with $N$ because $L$ is an integer.
Then $c\in L^{-N}\Z$ for all sufficiently large $N$;
since $N_i\to\infty$ and $\Z\subset\Adm(\ell)$, one has $c\in\eta_i\Z\subset\eta_i\Adm(\ell)$ for all large $i$.
Let $s$, $a_1,\dots,a_s$ and $f_1,\dots,f_s$ be as in~\eqref{eq:continuum-plane-positivity}.
The supports of the $f_p$ are compact subsets of the open half-space $\{\ell>c\}$;
hence there is $\delta>0$ with $\supp f_p\subset\{\ell\ge c+\delta\}$ for $p=1,\dots,s$.
For all large $i$ both $c\in\eta_i\Adm(\ell)$ and $2\eta_i<\delta$ hold,
and Proposition~\ref{prop:volume-limits}(c), applied at the spacing $\eta_i$ to the volume limit $\mu^{(N_i)}\in\mathcal{V}_{N_i}$,
gives~\eqref{eq:continuum-plane-positivity} with $\mu^{(N_i)}$ in the place of $\mu$.
For every $f\in\mathcal{S}$ the map $F\mapsto F(f)$ is a continuous linear functional on $H_{-4}$,
and $f_q$ and $r_{\ell,c}\cdot f_p$ belong to $\mathcal{S}$;
so for all $p,q$ the function $F\mapsto\exp\left(i\left(F(f_q)-F(r_{\ell,c}\cdot f_p)\right)\right)$ is bounded and continuous on $H_{-4}$,
and its integral against $\mu^{(N_i)}$ tends to its integral against $\mu$ as $i\to\infty$.
The left side of~\eqref{eq:continuum-plane-positivity} is thus the limit of the corresponding sums for the laws $\mu^{(N_i)}$,
which are $\ge0$ for all large $i$.
So $\mu$ is reflection positive at every plane $\{\ell=c\}$ with $c\in\Z[L^{-1}]$.

\emph{Arbitrary positions.}
Let now $c\in\R$ be arbitrary, let $c'\in\Z[L^{-1}]$, and put $v:=(c-c')\cdot|n_\ell|^{-2}\cdot n_\ell$, so that $\ell(v)=c-c'$.
For $y\in\R^3$ one has
\[
r_{\ell,c'}(y-v)+v=y-2|n_\ell|^{-2}\cdot\left(\ell(y)-\ell(v)-c'\right)\cdot n_\ell=r_{\ell,c}(y),
\]
that is, $r_{\ell,c}=\tau_v\circ r_{\ell,c'}\circ\tau_{-v}$.
Let $s\ge1$, $a_1,\dots,a_s\in\C$ and real $f_1,\dots,f_s\in C^{\infty}_{\mathrm{c}}(\R^3;\Lambda^2)$ with supports in $\{\ell>c\}$ be given,
and put $f'_p:=\tau_{-v}\cdot f_p$.
The support of $f'_p$ is $\supp f_p-v$, and $\ell(y-v)=\ell(y)-c+c'>c'$ for $y\in\supp f_p$;
so the $f'_p$ are real test forms with compact supports in $\{\ell>c'\}$.
The action on test forms is compatible with composition;
hence $\tau_v\cdot f'_q=f_q$ and $\tau_v\cdot\left(r_{\ell,c'}\cdot f'_p\right)=r_{\ell,c}\cdot f_p$.
The operator $U(v)=\hat\tau_v$ maps the function $F\mapsto F(f)$ onto the function $F\mapsto F(\tau_v\cdot f)$ (Definition~\ref{def:laws}),
and, being the composition with a map of $\mathcal{S}'$ onto itself, it commutes with the exponential;
therefore
\[
\exp\left(i\left(F(f_q)-F(r_{\ell,c}\cdot f_p)\right)\right)=\left(U(v)\Phi_{pq}\right)(F),\qquad
\Phi_{pq}(F):=\exp\left(i\left(F(f'_q)-F(r_{\ell,c'}\cdot f'_p)\right)\right).
\]
As $\mu$ is invariant under the translations of $\R^3$, one has $\mu(U(v)\Phi)=\mu(\Phi)$ for every bounded $\Sigma$-measurable function $\Phi$ on $\mathcal{S}'$.
Consequently the left side of~\eqref{eq:continuum-plane-positivity} formed with the plane $\{\ell=c\}$ and the forms $f_p$
equals the left side formed with the plane $\{\ell=c'\}$, the forms $f'_p$ and the same coefficients $a_p$,
and the latter is $\ge0$ by the case of the positions in $\Z[L^{-1}]$.
\end{proof}

\subsection{Two spectral lemmas}\label{sec:spectral-lemmas}

All laws in this subsection are laws on $\mathcal{S}'$ in the sense of Definition~\ref{def:laws}, that is, probability measures on $(\mathcal{S}',\Sigma)$.

\begin{definition}[Observables on the positive side]\label{def:positive-side}
Let $\mu$ be a law on $\mathcal{S}'$ and let $j\in\{1,2,3\}$.
The elements of $L^2(\mu)$, with the scalar product $\langle\cdot,\cdot\rangle_\mu$ of Definition~\ref{def:laws}, are called observables.
Put $\Sigma_j^+:=\Sigma_{\{x_j>0\}}$;
this is the $\sigma$-algebra on $\mathcal{S}'$ generated by the functions $F\mapsto F(f)$ with $f\in C_c^\infty(\R^3;\Lambda^2)$ real and $\supp f\subset\{x_j>0\}$.
The elements of the closed subspace $L^2(\Sigma_j^+,\mu)$ of $L^2(\mu)$, formed by the classes that contain a $\Sigma_j^+$-measurable function, are called the observables on the positive side.
\end{definition}

For $\mu$ and $j$ as in Definition~\ref{def:positive-side}, the space $L^2(\Sigma_j^+,\mu)$ is the space denoted $\mathcal{E}_+$ in \cite[Proposition~8.19]{NoncompactQED}.
Let now $\mu$ satisfy $(T_{123})$, $(R_j)$ and $(RP_j)$ (Definition~\ref{def:laws}).
The reflection $\Theta_j$ of Definition~\ref{def:laws}, given by $\Theta_j\Phi=\overline{\hat\theta_j\Phi}$, is antilinear, whereas the reflection that \cite[Section~8.1.2]{NoncompactQED} denotes $\Theta_j$ is the linear map $\hat\theta_j$;
accordingly, the form $(\Phi,\Phi')\mapsto\left\langle\hat\theta_j\Phi,\Phi'\right\rangle_\mu$ from which \cite[Proposition~8.19]{NoncompactQED} builds the reconstruction is the reflected form $(\Phi,\Phi')\mapsto\mu\left((\Theta_j\Phi)\cdot\Phi'\right)$ on $L^2(\Sigma_j^+,\mu)$.
Recall that the reconstruction $\left(\mathcal{H}_j,[\,\cdot\,]_j,\mathsf{H}_j,V_j,\Omega_j\right)$ of $\mu$ in the direction $j$ (Definition~\ref{def:laws}; \cite[Proposition~8.19\,(i)--(iii)]{NoncompactQED}) consists of a Hilbert space $\mathcal{H}_j$, with scalar product $\langle\cdot,\cdot\rangle$ and norm $\|\cdot\|$;
the canonical map $[\,\cdot\,]_j\colon L^2(\Sigma_j^+,\mu)\to\mathcal{H}_j$, a linear contraction with dense range such that
\[
\left\langle[\Phi]_j,[\Phi']_j\right\rangle=\mu\left((\Theta_j\Phi)\cdot\Phi'\right)\qquad\left(\Phi,\Phi'\in L^2(\Sigma_j^+,\mu)\right);
\]
the self-adjoint operator $\mathsf{H}_j\ge0$ in $\mathcal{H}_j$ with $\exp\left(-t\cdot\mathsf{H}_j\right)[\Phi]_j=\left[U(t\cdot e_j)\Phi\right]_j$ for $t\ge0$;
the unitary operators $V_j(y)$, $y\in e_j^{\perp}$, with $V_j(y)[\Phi]_j=\left[U(y)\Phi\right]_j$;
and the vector $\Omega_j=[1]_j$.

\begin{lemma}[Joint spectral measure]\label{lem:joint-spectrum}
Let $j\in\{1,2,3\}$, let $\mu$ be a law on $\mathcal{S}'$ that satisfies $(T_{123})$, $(R_j)$ and $(RP_j)$, and let $\left(\mathcal{H}_j,[\,\cdot\,]_j,\mathsf{H}_j,V_j,\Omega_j\right)$ be its reconstruction in the direction $j$.
The plane $e_j^{\perp}=\left\{y\in\R^3:y_j=0\right\}$ is identified with $\R^2$ through the coordinates $y=(y_i)_{i\neq j}$, and for $k=(k_i)_{i\neq j}\in\R^2$ one puts
\[
k\cdot y=\sum_{i\neq j}k_i\cdot y_i,\qquad|k|_\infty=\max_{i\neq j}|k_i|.
\]
\begin{enumerate}
\item[(a)] The Hilbert space $\mathcal{H}_j$ is separable.
The map $y\mapsto V_j(y)$ is a strongly continuous unitary representation of the additive group $e_j^{\perp}$ on $\mathcal{H}_j$, and $V_j(y)$ commutes with $\exp\left(-t\cdot\mathsf{H}_j\right)$ for all $t\ge0$ and $y\in e_j^{\perp}$.
\item[(b)] There is exactly one projection-valued measure $\mathfrak{E}_j$ on the Borel sets of the space $\R\times e_j^{\perp}$ of the pairs $(\lambda,k)$ of an energy $\lambda$ and a momentum $k$ such that
\[
\exp\left(i\,t\cdot\mathsf{H}_j\right)V_j(y)=\int\mathrm{d}\mathfrak{E}_j(\lambda,k)\,\exp\left(i\,t\cdot\lambda+i\,k\cdot y\right)\qquad\left(t\in\R,\ y\in e_j^{\perp}\right).
\]
The measure $\mathfrak{E}_j$ is carried by $[0,\infty)\times e_j^{\perp}$, and for every bounded Borel function $u$ on $[0,\infty)$ and every $y\in e_j^{\perp}$
\[
u(\mathsf{H}_j)V_j(y)=\int\mathrm{d}\mathfrak{E}_j(\lambda,k)\,u(\lambda)\cdot\exp\left(i\,k\cdot y\right);
\]
in particular
\[
\exp\left(-t\cdot\mathsf{H}_j\right)V_j(y)=\int\mathrm{d}\mathfrak{E}_j(\lambda,k)\,\exp\left(-t\cdot\lambda+i\,k\cdot y\right)\qquad\left(t\ge0,\ y\in e_j^{\perp}\right).
\]
\item[(c)] Let $i\neq j$.
The self-adjoint generator $\mathsf{P}_i$ of the unitary group $t\mapsto V_j(t\cdot e_i)$ of Definition~\ref{def:laws}, that is, the self-adjoint operator with $V_j(t\cdot e_i)=\exp\left(i\,t\cdot\mathsf{P}_i\right)$ for $t\in\R$, is
\[
\mathsf{P}_i=\int\mathrm{d}\mathfrak{E}_j(\lambda,k)\,k_i.
\]
The image of $\mathfrak{E}_j$ under the map $(\lambda,k)\mapsto(\lambda,k_i)$ is called the \emph{joint spectral measure} of $(\mathsf{H}_j,\mathsf{P}_i)$, and its support the \emph{joint spectrum} of $(\mathsf{H}_j,\mathsf{P}_i)$.
\item[(d)] Let $u$ be a real bounded Borel function on $[0,\infty)\times\R^2$.
Then $\int\mathrm{d}\mathfrak{E}_j(\lambda,k)\,u(\lambda,k)\ge0$ if and only if $\mathfrak{E}_j\left(\{u<0\}\right)=0$.
\end{enumerate}
For $\psi\in\mathcal{H}_j$ let $\varrho_\psi=\left\langle\psi,\mathfrak{E}_j(\cdot)\psi\right\rangle$;
it is a positive measure of total mass $\|\psi\|^2$.
\end{lemma}

\begin{proof}
(a) By \cite[Lemma~8.12\,(i)]{NoncompactQED} there is a countable set $\mathcal{X}_j$ of finite linear combinations of functions $F\mapsto\exp\left(i\,F(f)\right)$, with $f\in C_c^\infty(\R^3;\Lambda^2)$ real and $\supp f\subset\{x_j>0\}$, that is dense in $L^2(\Sigma_j^+,\mu)$;
it is the set denoted $\mathcal{A}^0(\{x_j>0\})$ there.
The map $[\,\cdot\,]_j$ is a contraction with dense range \cite[Proposition~8.19\,(i)]{NoncompactQED}.
Hence the countable set $\left\{[\Phi]_j:\Phi\in\mathcal{X}_j\right\}$ is dense in $\mathcal{H}_j$, and $\mathcal{H}_j$ is separable.

By $(T_{123})$ the law $\mu$ is invariant under every translation.
So, for $y\in e_j^{\perp}$, the operator $V_j(y)$ is unitary and commutes with every $\exp\left(-t\cdot\mathsf{H}_j\right)$, $t\ge0$ \cite[Proposition~8.19\,(iii)]{NoncompactQED}.
For $y,y'\in e_j^{\perp}$ and $\Phi\in L^2(\Sigma_j^+,\mu)$ one has $V_j(0)[\Phi]_j=[\Phi]_j$ and, since $U(y')\Phi$ lies again in $L^2(\Sigma_j^+,\mu)$ (the translation by $y'$ maps the half space $\{x_j>0\}$ onto itself),
\[
V_j(y)V_j(y')[\Phi]_j=\left[U(y)U(y')\Phi\right]_j=\left[U(y+y')\Phi\right]_j=V_j(y+y')[\Phi]_j;
\]
since the classes are dense, $V_j(0)=1$ and $V_j(y)V_j(y')=V_j(y+y')$.
The map $a\mapsto U(a)$ is strongly continuous on $L^2(\mu)$ \cite[Lemma~8.13\,(ii)]{NoncompactQED}, and $[\,\cdot\,]_j$ is a contraction;
therefore
\[
\left\|V_j(y)[\Phi]_j-[\Phi]_j\right\|=\left\|\left[U(y)\Phi-\Phi\right]_j\right\|\le\left\|U(y)\Phi-\Phi\right\|_{L^2(\mu)}\to0\qquad(y\to0).
\]
As the classes are dense and $\|V_j(y)\|=1$, it follows that $V_j(y)\psi\to\psi$ as $y\to0$ for every $\psi\in\mathcal{H}_j$;
by the group law, $y\mapsto V_j(y)\psi$ is continuous on all of $e_j^{\perp}$.

(b) Because $\mathsf{H}_j\ge0$, the operator $\exp\left(-\mathsf{H}_j\right)$ is bounded and self-adjoint with $0\le\exp\left(-\mathsf{H}_j\right)\le1$.
It is injective, since the function $\lambda\mapsto\exp(-\lambda)$ has no zero, and it commutes with every $V_j(y)$, by (a).
For $t\in\R$ let $u_t$ be the bounded Borel function on $[0,1]$ given by
\[
u_t(x)=\exp\left(-i\,t\cdot\log x\right)\quad(0<x\le1),\qquad u_t(0)=0.
\]
As $u_t\left(\exp(-\lambda)\right)=\exp\left(i\,t\cdot\lambda\right)$ for $\lambda\ge0$, the functional calculus gives $\exp\left(i\,t\cdot\mathsf{H}_j\right)=u_t\left(\exp\left(-\mathsf{H}_j\right)\right)$.
An operator that commutes with a bounded self-adjoint operator commutes with every bounded Borel function of it \cite[Theorem~VII.2\,(g)]{ReedSimon1980};
hence $V_j(y)$ commutes with $\exp\left(i\,t\cdot\mathsf{H}_j\right)$ for all $t\in\R$ and $y\in e_j^{\perp}$.
The unitary groups $t\mapsto\exp\left(i\,t\cdot\mathsf{H}_j\right)$ and $y\mapsto V_j(y)$ are strongly continuous, the first by the functional calculus of $\mathsf{H}_j$ and the second by (a), and they commute with each other.
Consequently
\[
(t,y)\mapsto\exp\left(i\,t\cdot\mathsf{H}_j\right)V_j(y)
\]
is a strongly continuous unitary representation of the group $\R\times e_j^{\perp}$, which is isomorphic to $\R^3$, on the separable Hilbert space $\mathcal{H}_j$.
A strongly continuous unitary representation of $\R^3$ on a separable Hilbert space is the Fourier transform of a projection-valued measure \cite[Theorem~VIII.12]{ReedSimon1980};
this yields $\mathfrak{E}_j$ with the identity stated in (b).
If $\mathfrak{E}'$ is a second projection-valued measure with the same identity, then for all $\psi,\psi'\in\mathcal{H}_j$ the finite complex measures $\left\langle\psi',\mathfrak{E}_j(\cdot)\psi\right\rangle$ and $\left\langle\psi',\mathfrak{E}'(\cdot)\psi\right\rangle$ have the same Fourier transform;
by the uniqueness theorem for Fourier transforms of finite complex measures \cite[Theorem~1.3.6]{Rudin1962} they coincide, and $\mathfrak{E}'=\mathfrak{E}_j$.

The image of $\mathfrak{E}_j$ under $(\lambda,k)\mapsto\lambda$ is a projection-valued measure on $\R$ whose Fourier transform is $t\mapsto\exp\left(i\,t\cdot\mathsf{H}_j\right)$, as the case $y=0$ of the identity shows.
By the spectral theorem the spectral measure of $\mathsf{H}_j$ has the same Fourier transform, and by the same uniqueness theorem the two projection-valued measures coincide.
Since $\mathsf{H}_j\ge0$, its spectral measure vanishes on $(-\infty,0)$;
so does the image of $\mathfrak{E}_j$, that is, $\mathfrak{E}_j\left((-\infty,0)\times e_j^{\perp}\right)=0$, and $\mathfrak{E}_j$ is carried by $[0,\infty)\times e_j^{\perp}$.
Expressed through this image measure, the functional calculus of $\mathsf{H}_j$ takes the form
\[
u(\mathsf{H}_j)=\int\mathrm{d}\mathfrak{E}_j(\lambda,k)\,u(\lambda)
\]
for every bounded Borel function $u$ on $[0,\infty)$.
For $t=0$ the identity defining $\mathfrak{E}_j$ reads
\begin{equation}\label{eq:continuum-translations}
V_j(y)=\int\mathrm{d}\mathfrak{E}_j(\lambda,k)\,\exp\left(i\,k\cdot y\right)\qquad\left(y\in e_j^{\perp}\right).
\end{equation}
The two operators $u(\mathsf{H}_j)$ and $V_j(y)$ are therefore the integrals, with respect to $\mathfrak{E}_j$, of the bounded Borel functions $(\lambda,k)\mapsto u(\lambda)$ and $(\lambda,k)\mapsto\exp\left(i\,k\cdot y\right)$.
Their product is the integral of the product of these two functions, $(\lambda,k)\mapsto u(\lambda)\cdot\exp\left(i\,k\cdot y\right)$, which is the formula for $u(\mathsf{H}_j)V_j(y)$ in (b);
the function $u(\lambda)=\exp(-t\cdot\lambda)$ with $t\ge0$ turns it into the last display of (b).

(c) The function $(\lambda,k)\mapsto k_i$ is real and Borel.
Therefore the operator $\int\mathrm{d}\mathfrak{E}_j(\lambda,k)\,k_i$ is self-adjoint, and its unitary group is
\[
\int\mathrm{d}\mathfrak{E}_j(\lambda,k)\,\exp\left(i\,t\cdot k_i\right)=V_j(t\cdot e_i)\qquad(t\in\R),
\]
where the equality is \eqref{eq:continuum-translations} at the point $y=t\cdot e_i$ of $e_j^{\perp}$, at which $k\cdot y=t\cdot k_i$.
By Stone's theorem \cite[Theorem~VIII.8]{ReedSimon1980} a strongly continuous unitary group has exactly one self-adjoint generator;
thus $\int\mathrm{d}\mathfrak{E}_j(\lambda,k)\,k_i$ is the generator $\mathsf{P}_i$ of $t\mapsto V_j(t\cdot e_i)$.

(d) For $\psi\in\mathcal{H}_j$ one has $\varrho_\psi=\left\|\mathfrak{E}_j(\cdot)\psi\right\|^2$ on the Borel sets, and
\[
\left\langle\psi,\left(\int\mathrm{d}\mathfrak{E}_j\,u\right)\psi\right\rangle=\int\mathrm{d}\varrho_\psi\,u.
\]
If $\mathfrak{E}_j\left(\{u<0\}\right)=0$, then $\varrho_\psi\left(\{u<0\}\right)=0$ for every $\psi$, the right side is non-negative, and $\int\mathrm{d}\mathfrak{E}_j\,u\ge0$.
Conversely, let $\int\mathrm{d}\mathfrak{E}_j\,u\ge0$.
For an integer $m\ge1$ put $S_m=\left\{u\le-\frac{1}{m}\right\}$, and let $\psi$ lie in the range of $\mathfrak{E}_j(S_m)$.
Then $\varrho_\psi$ is carried by $S_m$, and
\[
0\le\int\mathrm{d}\varrho_\psi\,u=\int_{S_m}\mathrm{d}\varrho_\psi\,u\le-\frac{1}{m}\cdot\varrho_\psi(S_m)=-\frac{1}{m}\cdot\|\psi\|^2,
\]
so that $\psi=0$.
Thus $\mathfrak{E}_j(S_m)=0$ for every $m\ge1$.
The set $\{u<0\}$ is the union of the increasing sequence $(S_m)_{m\ge1}$, and $\mathfrak{E}_j\left(\{u<0\}\right)=0$.
\end{proof}

\begin{lemma}[An abstract cone inequality]\label{lem:abstract-cone}
Let $\mathcal{H}$ and $\mathcal{H}'$ be Hilbert spaces with scalar products $\langle\cdot,\cdot\rangle$ and $\langle\cdot,\cdot\rangle'$, antilinear in the first entry, and norms $\|\cdot\|$.
Let $\mathfrak{T}$ be a self-adjoint operator on $\mathcal{H}$ with $0\le\mathfrak{T}\le1$, and $\mathfrak{U}$ a unitary operator on $\mathcal{H}$ with $\mathfrak{T}\mathfrak{U}=\mathfrak{U}\mathfrak{T}$.
Let $\mathfrak{T}_{+}$ be a self-adjoint operator on $\mathcal{H}'$ with $0\le\mathfrak{T}_{+}\le1$, and $\mathfrak{V}_{-}$ a unitary operator on $\mathcal{H}'$.
Let $\mathcal{H}_{0}\subset\mathcal{H}$ be a dense subset with the following property: to each $\psi\in\mathcal{H}_{0}$ there correspond an integer $r\ge0$ and vectors $\psi_{1},\psi_{2}\in\mathcal{H}'$ for which the identity
\begin{equation}\label{eq:continuum-two-pictures}
\left\langle\psi,\mathfrak{T}^{r+k+k'}\mathfrak{U}^{k-k'}\psi\right\rangle=\left\langle\psi_{1},\mathfrak{T}_{+}^{k}\mathfrak{V}_{-}^{k'}\psi_{2}\right\rangle'
\end{equation}
holds for all integers $k\ge0$ and $k'\ge0$.
Then
\[
\left(1-\mathfrak{T}\right)^2\ge\mathfrak{T}\left(2-\mathfrak{U}-\mathfrak{U}^*\right).
\]
\end{lemma}

\begin{proof}
From $\mathfrak{T}\mathfrak{U}=\mathfrak{U}\mathfrak{T}$ and $\mathfrak{U}^*=\mathfrak{U}^{-1}$ it follows that $\mathfrak{T}$ commutes with $\mathfrak{U}^*$ as well.
Let $\mathfrak{A}$ be the $C^*$-algebra generated by $1$, $\mathfrak{T}$ and $\mathfrak{U}$, that is, the norm closure of the set of polynomials in $\mathfrak{T}$, $\mathfrak{U}$ and $\mathfrak{U}^*$;
since these three operators commute with one another, $\mathfrak{A}$ is commutative.
By the Gelfand--Naimark theorem \cite[Theorem~11.18]{Rudin1991}, the Gelfand transform is an isometric $*$-isomorphism of $\mathfrak{A}$ onto the algebra of the continuous functions on the compact space of the characters of $\mathfrak{A}$.
A character $\chi$ of $\mathfrak{A}$ is a continuous $*$-homomorphism into $\C$, so it is determined by its values $\chi(\mathfrak{T})$ and $\chi(\mathfrak{U})$ on the generators.
The map $\chi\mapsto\left(\chi(\mathfrak{T}),\chi(\mathfrak{U})\right)$ is continuous on the space of characters, and injective by what has just been said;
it is therefore a homeomorphism of this compact space onto a compact set $\operatorname{sp}(\mathfrak{T},\mathfrak{U})\subset\C^2$, the joint spectrum of $\mathfrak{T}$ and $\mathfrak{U}$.
This set lies in $[0,1]\times S^1$.
Indeed, $\mathfrak{T}$ and $1-\mathfrak{T}$ are positive, so they are the squares of their positive square roots, which are self-adjoint elements of $\mathfrak{A}$;
as $\chi$ is real on self-adjoint elements, this gives $\chi(\mathfrak{T})\ge0$ and $1-\chi(\mathfrak{T})\ge0$.
Moreover $|\chi(\mathfrak{U})|^2=\chi\left(\mathfrak{U}^*\mathfrak{U}\right)=1$.
Consequently there is an isometric $*$-isomorphism $\vartheta\mapsto\vartheta(\mathfrak{T},\mathfrak{U})$ of the algebra $C\left(\operatorname{sp}(\mathfrak{T},\mathfrak{U})\right)$ of continuous functions onto $\mathfrak{A}$ that maps the coordinate functions $(t,\zeta)\mapsto t$ and $(t,\zeta)\mapsto\zeta$ to $\mathfrak{T}$ and $\mathfrak{U}$.

Let $\psi\in\mathcal{H}$.
The linear functional $\vartheta\mapsto\left\langle\psi,\vartheta(\mathfrak{T},\mathfrak{U})\psi\right\rangle$ on $C\left(\operatorname{sp}(\mathfrak{T},\mathfrak{U})\right)$ is positive, because $\vartheta(\mathfrak{T},\mathfrak{U})=\vartheta^{\frac{1}{2}}(\mathfrak{T},\mathfrak{U})^*\,\vartheta^{\frac{1}{2}}(\mathfrak{T},\mathfrak{U})$ for $\vartheta\ge0$.
By the Riesz representation theorem there is a finite positive Borel measure $\varrho_\psi$ on $\operatorname{sp}(\mathfrak{T},\mathfrak{U})$ such that
\[
\left\langle\psi,\vartheta(\mathfrak{T},\mathfrak{U})\psi\right\rangle=\int\mathrm{d}\varrho_\psi(t,\zeta)\,\vartheta(t,\zeta)\qquad\left(\vartheta\in C\left(\operatorname{sp}(\mathfrak{T},\mathfrak{U})\right)\right).
\]
Let $u$ be the real continuous function on $[0,1]\times S^1$ given by
\[
u(t,\zeta)=(1-t)^2-t\cdot|1-\zeta|^2=(1-t)^2-t\cdot\left(2-\zeta-\overline{\zeta}\right).
\]
Then $u(\mathfrak{T},\mathfrak{U})=\left(1-\mathfrak{T}\right)^2-\mathfrak{T}\left(2-\mathfrak{U}-\mathfrak{U}^*\right)$, and the assertion is $u(\mathfrak{T},\mathfrak{U})\ge0$.

Now let $\psi\in\mathcal{H}_{0}$, with $r$, $\psi_{1}$ and $\psi_{2}$ as in \eqref{eq:continuum-two-pictures}.
Let $\nu$ be the image, under the map $(t,\zeta)\mapsto t\cdot\zeta$ of $[0,1]\times S^1$ into the closed unit disc $\overline{\mathbb{D}}$, of the measure that has the density $(t,\zeta)\mapsto t^r$ with respect to $\varrho_\psi$ (the density $1$ if $r=0$).
It is a finite positive Borel measure on $\overline{\mathbb{D}}$.
Let $p(z)=\sum_k a_k\cdot z^k$ be a polynomial with complex coefficients, and let $p^*$ be the polynomial whose coefficients are their complex conjugates, $p^*(z)=\sum_k\overline{a_k}\cdot z^k$.
For $t\in[0,1]$ and $|\zeta|=1$,
\[
\left|p(t\cdot\zeta)\right|^2=\sum_{k,k'}a_k\cdot\overline{a_{k'}}\cdot t^{k+k'}\cdot\zeta^{k-k'},
\]
and the function $(t,\zeta)\mapsto t^{r+k+k'}\cdot\zeta^{k-k'}$ corresponds to the operator $\mathfrak{T}^{r+k+k'}\mathfrak{U}^{k-k'}$.
By the definition of $\nu$ and by \eqref{eq:continuum-two-pictures},
\[
\begin{aligned}
\int\mathrm{d}\nu(z)\,|p(z)|^2&=\int\mathrm{d}\varrho_\psi(t,\zeta)\,t^r\cdot\left|p(t\cdot\zeta)\right|^2=\sum_{k,k'}a_k\cdot\overline{a_{k'}}\cdot\left\langle\psi,\mathfrak{T}^{r+k+k'}\mathfrak{U}^{k-k'}\psi\right\rangle\\
&=\sum_{k,k'}a_k\cdot\overline{a_{k'}}\cdot\left\langle\psi_{1},\mathfrak{T}_{+}^{k}\mathfrak{V}_{-}^{k'}\psi_{2}\right\rangle'=\left\langle\psi_{1},p\left(\mathfrak{T}_{+}\right)p^*\left(\mathfrak{V}_{-}\right)\psi_{2}\right\rangle'.
\end{aligned}
\]
For a normal operator, the norm of a polynomial in the operator is the supremum of the modulus of the polynomial on the spectrum of the operator (by the spectral mapping theorem, since the norm of a normal operator equals its spectral radius).
Both $\mathfrak{T}_{+}$ and $\mathfrak{V}_{-}$ are normal;
the inequalities $0\le\mathfrak{T}_{+}\le1$ confine the spectrum of $\mathfrak{T}_{+}$ to $[0,1]$, and the spectrum of the unitary operator $\mathfrak{V}_{-}$ lies on $S^1$.
As moreover $p^*(z)=\overline{p(\overline{z})}$ and complex conjugation maps $S^1$ onto itself, one obtains
\[
\left\|p\left(\mathfrak{T}_{+}\right)\right\|\le\sup_{[0,1]}|p|,\qquad\left\|p^*\left(\mathfrak{V}_{-}\right)\right\|\le\sup_{S^1}|p^*|=\sup_{S^1}|p|.
\]
It follows that
\[
\int_{\overline{\mathbb{D}}}\mathrm{d}\nu(z)\,|p(z)|^2\le\|\psi_{1}\|\cdot\|\psi_{2}\|\cdot\sup_{[0,1]}|p|\cdot\sup_{S^1}|p|
\]
for every polynomial $p$.
This is the hypothesis \eqref{eq:layers-two-constants} of Lemma~\ref{lem:two-constants}, for the measure $\nu$ in the place of the measure there and with the constant $C=\|\psi_{1}\|\cdot\|\psi_{2}\|$.
That lemma gives
\[
\nu\left(\left\{z\in\overline{\mathbb{D}}:|1+z|^2<4|z|\right\}\right)=0.
\]

For $t>0$ and $|\zeta|=1$ one has $|t\cdot\zeta|=t$ and
\[
|1+t\cdot\zeta|^2-4t=1+2t\cdot\Re\zeta+t^2-4t=(1-t)^2-t\cdot\left(2-2\Re\zeta\right)=u(t,\zeta).
\]
So the map $(t,\zeta)\mapsto t\cdot\zeta$ carries the set $\left\{(t,\zeta)\in\operatorname{sp}(\mathfrak{T},\mathfrak{U}):t>0,\ u(t,\zeta)<0\right\}$ into the set $\left\{z\in\overline{\mathbb{D}}:|1+z|^2<4|z|\right\}$.
The integral of $t^r$ with respect to $\varrho_\psi$ over the former set is therefore at most the $\nu$-measure of the latter, which is $0$;
since $t^r>0$ on the former set, it is a $\varrho_\psi$-null set.
At the points with $t=0$ one has $u=1$.
Hence $u\ge0$ holds $\varrho_\psi$-almost everywhere, and
\[
\left\langle\psi,u(\mathfrak{T},\mathfrak{U})\psi\right\rangle=\int\mathrm{d}\varrho_\psi(t,\zeta)\,u(t,\zeta)\ge0.
\]
The function $\psi\mapsto\left\langle\psi,u(\mathfrak{T},\mathfrak{U})\psi\right\rangle$ is continuous on $\mathcal{H}$ and non-negative on the dense set $\mathcal{H}_{0}$, hence non-negative on $\mathcal{H}$:
$\left(1-\mathfrak{T}\right)^2-\mathfrak{T}\left(2-\mathfrak{U}-\mathfrak{U}^*\right)\ge0$.
\end{proof}

\begin{remark}\label{rem:abstract-cone}
\begin{enumerate}
\item[(1)] In Lemma~\ref{lem:abstract-cone} the operators $\mathfrak{T}$ and $\mathfrak{T}_{+}$ are positive, so that their spectra lie in $[0,1]$ and not merely in $[-1,1]$.
For a translation by one lattice step normal to a mirror, the positivity of the transfer operator of Lemma~\ref{lem:lattice-reconstruction}(b) is reflection positivity at the mirror halfway between that mirror and its translate, as in the lattice argument for four-dimensional Yang--Mills theory.
In the continuum, $\exp\left(-s\cdot\mathsf{H}_j\right)$ with $s>0$ is positive, being the square of the self-adjoint operator $\exp\left(-\frac{s}{2}\cdot\mathsf{H}_j\right)$.
\item[(2)] The identity \eqref{eq:continuum-two-pictures} relates two mirrors.
Suppose that, at a mirror normal to $e_j$, the operator $\mathfrak{T}$ is the transfer operator of a step of length $s$ along $e_j$, and $\mathfrak{U}$ is the translation by $s\cdot e_i$ for a direction $e_i$ with $i\neq j$.
Then $\mathfrak{T}^{r+k+k'}\mathfrak{U}^{k-k'}$ corresponds to the translation by
\[
s\cdot\left(r\cdot e_j+k\cdot(e_j+e_i)+k'\cdot(e_j-e_i)\right),
\]
which consists, after $r$ steps along $e_j$, of $k$ steps along $e_j+e_i$ and $k'$ steps along $e_j-e_i$.
At a second mirror, normal to $e_j+e_i$, the steps along $e_j+e_i$ are steps in time and the steps along $e_j-e_i$ are translations in space;
the right side of \eqref{eq:continuum-two-pictures} is the same matrix element, computed with the transfer operator $\mathfrak{T}_{+}$ and the translation $\mathfrak{V}_{-}$ of that second mirror.
\end{enumerate}
\end{remark}

\subsection{The energy-momentum cone}\label{sec:cone}

For an open set $U\subset\R^3$ let $C^\infty_{\mathrm c}(U;\Lambda^2)\subset\mathcal S$ be the space of the real smooth two-forms on $\R^3$ whose support is a compact subset of $U$;
in \cite[Section~8.1.1]{NoncompactQED} it is denoted $\mathcal D(U)$.
By Definition~\ref{def:laws}, $\Sigma_U$ is the $\sigma$-algebra generated by the functions $F\mapsto F(f)$ with $f\in C^\infty_{\mathrm c}(U;\Lambda^2)$.

\begin{theorem}[Cone theorem]\label{thm:cone}
Let $i\neq j$ be indices in $\{1,2,3\}$, and let $\mu$ be a law on $\mathcal S'$ that is invariant under the translations of $\R^3$, satisfies $(R_j)$ and $(RP_j)$, and is invariant under the reflection in the plane $\{x_j+x_i=0\}$ and reflection positive at this plane (Definition~\ref{def:laws}).
Let $\left(\mathcal H_j,[\,\cdot\,]_j,\mathsf H_j,V_j,\Omega_j\right)$ be the reconstruction of $\mu$ in the direction $j$, let $\mathsf P_i$ be the self-adjoint generator of the unitary group $t\mapsto V_j(t\cdot e_i)$, and let $\mathfrak E_j$ be the joint spectral measure of Lemma~\ref{lem:joint-spectrum}.
\begin{enumerate}
\item[(a)]
For every $s>0$ one has, as an inequality between bounded self-adjoint operators on $\mathcal H_j$,
\[
\left(1-\exp\left(-s\cdot\mathsf H_j\right)\right)^{2}\ge\exp\left(-s\cdot\mathsf H_j\right)\left(2-V_j(s\cdot e_i)-V_j(-s\cdot e_i)\right).
\]
\item[(b)]
For every $s>0$ the joint spectral measure of $(\mathsf H_j,\mathsf P_i)$ is carried by the set of the points $(\lambda,k_i)$ with $\cosh(s\cdot\lambda)+\cos(s\cdot k_i)\ge2$.
For fixed $s>0$ this property is equivalent to the inequality of~(a).
\item[(c)]
The joint spectral measure of $(\mathsf H_j,\mathsf P_i)$ is carried by the set of the points $(\lambda,k_i)$ with $\lambda\ge|k_i|$.
Consequently every vector $\psi$ in the domain of $\mathsf H_j$ lies in the domain of $\mathsf P_i$, and
\[
\|\mathsf P_i\psi\|\le\|\mathsf H_j\psi\|,\qquad\langle\psi,|\mathsf P_i|\psi\rangle\le\langle\psi,\mathsf H_j\psi\rangle;
\]
thus $\mathsf H_j\ge|\mathsf P_i|$ in the sense of quadratic forms, with the constant one.
\item[(d)]
(Five-point inequality.)
Let $\Phi\in L^2(\Sigma_j^+,\mu)$ (Definition~\ref{def:positive-side}), and put
\[
u(v)=\mu\left((\Theta_j\Phi)\cdot U(v)\Phi\right)\qquad(v\in\R^3),
\]
with the antilinear reflection $\Theta_j$ and the translations $U(v)$ of Definition~\ref{def:laws}.
Then, for every $t>0$,
\begin{equation}\label{eq:continuum-five-point}
u(0)-4u(t\cdot e_j)+u(2t\cdot e_j)+u(t\cdot e_j+t\cdot e_i)+u(t\cdot e_j-t\cdot e_i)\ge0.
\end{equation}
Conversely, let $t>0$;
if \eqref{eq:continuum-five-point} holds for all $\Phi$ in a subset of $L^2(\Sigma_j^+,\mu)$ whose classes form a dense subset of $\mathcal H_j$, then the inequality of~(a) holds with $s=t$.
\end{enumerate}
\end{theorem}

\begin{proof}
Euclidean maps act on test forms, on $\mathcal S'$, on functions on $\mathcal S'$ and on laws as in Definition~\ref{def:laws}.
For Euclidean maps $\mathfrak g$ and $\mathfrak g'$ one has $\hat{\mathfrak g}F(f)=F(\mathfrak g\cdot f)$ and $\widehat{\mathfrak g\mathfrak g'}=\hat{\mathfrak g}\,\hat{\mathfrak g}'$;
the map $\hat{\mathfrak g}$ is multiplicative, commutes with complex conjugation, and carries the $\Sigma_U$-measurable functions onto the $\Sigma_{\mathfrak gU}$-measurable functions, for every open set $U\subset\R^3$ \cite[Section~8.1.2]{NoncompactQED}.
Recall that $U(v)=\hat\tau_v$ and $\overline{\Theta_j\Phi}=\hat\theta_j\Phi$.
As $\mu$ satisfies $(T_{123})$, $(R_j)$ and $(RP_j)$, its reconstruction has the properties listed in Definition~\ref{def:laws}:
the class map $[\,\cdot\,]_j$ is a linear contraction of $L^2(\Sigma_j^+,\mu)$ into $\mathcal H_j$ with dense range, $\mathsf H_j\ge0$, every $V_j(b\cdot e_i)$ with $b\in\R$ is unitary and commutes with every $\exp\left(-a\cdot\mathsf H_j\right)$ with $a\ge0$, and
\begin{equation}\label{eq:continuum-reconstruction}
\left\langle[\Phi]_j,\exp\left(-a\cdot\mathsf H_j\right)V_j(b\cdot e_i)[\Phi']_j\right\rangle=\mu\left((\Theta_j\Phi)\cdot U(a\cdot e_j+b\cdot e_i)\Phi'\right)
\end{equation}
for $\Phi,\Phi'\in L^2(\Sigma_j^+,\mu)$, $a\ge0$ and $b\in\R$
(in \cite[Proposition~8.19]{NoncompactQED} the reflected form is written with the linear map $\Phi\mapsto\overline{\Theta_j\Phi}$).
\begin{enumerate}
\item[(1)]\itemmark{1}\label{it:cone-rotated-law}
\emph{The rotated law.}
Let $i'$ be the third index, so that $\{i,j,i'\}=\{1,2,3\}$, and put $n_\pm=\frac{1}{\sqrt2}\cdot\left(e_j\pm e_i\right)$.
The vectors $n_+$, $n_-$, $e_{i'}$ form an orthonormal basis of $\R^3$.
Let $g\in\mathrm O(3)$ be the map with $ge_j=n_+$, $ge_i=n_-$ and $ge_{i'}=e_{i'}$, and let $r=g\theta_jg^{-1}$.
As $\theta_j$ is the reflection in the plane orthogonal to $e_j$, the map $r$ is the reflection in the plane orthogonal to $n_+$, which is the plane $\{x_j+x_i=0\}$;
that is, $r=r_{n_+,0}$.
The map $g$ carries $\{x_j>0\}=\{e_j\cdot x>0\}$ onto $\{n_+\cdot x>0\}=\{x_j+x_i>0\}$, and $r$ exchanges the half spaces $\{x_j+x_i>0\}$ and $\{x_j+x_i<0\}$.

Let $\mu'=(g^{-1})_*\mu$.
By the definitions of an image law and of $\hat g$, one has $\mu'(\Phi)=\mu(\hat g\Phi)$ for every $\Sigma$-measurable function $\Phi$ that is non-negative or $\mu'$-integrable.
Hence $\hat g$ is a unitary operator from $L^2(\mu')$ onto $L^2(\mu)$, and $\hat g^{-1}=\widehat{g^{-1}}$ is a unitary operator from $L^2(\mu)$ onto $L^2(\mu')$;
the latter maps $L^2(\Sigma_{\{x_j+x_i>0\}},\mu)$ onto $L^2(\Sigma_j^+,\mu')$.
The law $\mu'$ satisfies $(T_{123})$, $(R_j)$ and $(RP_j)$.
Indeed, $\tau_vg^{-1}=g^{-1}\tau_{gv}$ for $v\in\R^3$, and $\theta_jg^{-1}=g^{-1}r$;
so the invariance of $\mu$ under the translations and under $r$ gives
\[
(\tau_v)_*\mu'=(g^{-1})_*(\tau_{gv})_*\mu=\mu',\qquad(\theta_j)_*\mu'=(g^{-1})_*r_*\mu=\mu'.
\]
Let $(f_p)$ be a finite family in $C^\infty_{\mathrm c}(\{x_j>0\};\Lambda^2)$ and $(a_p)$ a family of complex numbers with the same indices.
As $\hat gF(f)=F(g\cdot f)$ and $g\theta_j=rg$,
\begin{align*}
&\sum_{p,q}\overline{a_p}\cdot a_q\cdot\mu'\left(\exp\left(i\left(F(f_q)-F(\theta_j\cdot f_p)\right)\right)\right)\\
&\qquad=\sum_{p,q}\overline{a_p}\cdot a_q\cdot\mu\left(\exp\left(i\left(F(g\cdot f_q)-F\left(r\cdot(g\cdot f_p)\right)\right)\right)\right).
\end{align*}
The forms $g\cdot f_p$ lie in $C^\infty_{\mathrm c}(\{n_+\cdot x>0\};\Lambda^2)$.
The law $\mu$ is reflection positive at the plane $\{n_+\cdot x=0\}$ with the test forms supported in $\{n_+\cdot x>0\}$ (Definition~\ref{def:laws}; as $\mu$ is invariant under $r$, the two half spaces bounded by the plane give the same condition);
so the right side is non-negative.
This is $(RP_j)$ for $\mu'$.

Let $\left(\mathcal H',[\,\cdot\,]',\mathsf H',V'\right)$ be the space, the class map, the Hamiltonian and the spatial translations of the reconstruction of $\mu'$ in the direction $j$, and let $\langle\cdot,\cdot\rangle'$ be the scalar product of $\mathcal H'$.
The properties recalled above hold for them, with $\mu'$ in the place of $\mu$.
Let $\Phi_1,\Phi_2\in L^2(\Sigma_{\{x_j+x_i>0\}},\mu)$, $a\ge0$ and $b\in\R$.
Then $\hat g^{-1}\Phi_1$ and $\hat g^{-1}\Phi_2$ lie in $L^2(\Sigma_j^+,\mu')$, and \eqref{eq:continuum-reconstruction} for $\mu'$, together with $\mu'(\Phi)=\mu(\hat g\Phi)$ and the rules recalled above, gives
\begin{align*}
\left\langle\left[\hat g^{-1}\Phi_1\right]',\exp\left(-a\cdot\mathsf H'\right)V'(b\cdot e_i)\left[\hat g^{-1}\Phi_2\right]'\right\rangle'
&=\mu'\left(\left(\Theta_j\hat g^{-1}\Phi_1\right)\cdot U(a\cdot e_j+b\cdot e_i)\hat g^{-1}\Phi_2\right)\\
&=\mu\left(\overline{\hat g\hat\theta_j\hat g^{-1}\Phi_1}\cdot\hat gU(a\cdot e_j+b\cdot e_i)\hat g^{-1}\Phi_2\right).
\end{align*}
Since $g\theta_jg^{-1}=r$ and $g\tau_{a\cdot e_j+b\cdot e_i}g^{-1}=\tau_{a\cdot n_++b\cdot n_-}$, one has $\hat g\hat\theta_j\hat g^{-1}=\hat r$ and $\hat gU(a\cdot e_j+b\cdot e_i)\hat g^{-1}=U(a\cdot n_++b\cdot n_-)$.
Therefore
\begin{equation}\label{eq:continuum-rotated-picture}
\left\langle\left[\hat g^{-1}\Phi_1\right]',\exp\left(-a\cdot\mathsf H'\right)V'(b\cdot e_i)\left[\hat g^{-1}\Phi_2\right]'\right\rangle'=\left\langle\hat r\Phi_1,U(a\cdot n_++b\cdot n_-)\Phi_2\right\rangle_\mu .
\end{equation}

\item[(2)]\itemmark{2}\label{it:cone-dense-identity}
\emph{The identity of the two pictures.}
Fix $s>0$.
Let $\Phi=\sum_pa_p\cdot\exp\left(iF(f_p)\right)$ be a finite sum with $a_p\in\C$ and $f_p\in C^\infty_{\mathrm c}(\{x_j>0\};\Lambda^2)$.
The supports of the $f_p$ are compact subsets of $\{x_j>0\}$;
hence there is an integer $r_0\ge1$ such that, with $D=r_0\cdot s$, all of them lie in
\[
B=\left\{x\in\R^3:0<x_j<D,\ |x_i|<D\right\}.
\]
Let $y=\frac D2\cdot(e_j+e_i)$, so that $y_j+y_i=D$, and put
\[
\Phi_1=\hat rU(-y)\overline{\Theta_j\Phi},\qquad\Phi_2=U(2D\cdot e_j-y)\Phi.
\]
As $\overline{\Theta_j\Phi}=\hat\theta_j\Phi=\sum_pa_p\cdot\exp\left(iF(\theta_j\cdot f_p)\right)$, these functions are
\[
\Phi_1=\sum_pa_p\cdot\exp\left(iF\left(r\cdot\tau_{-y}\cdot\theta_j\cdot f_p\right)\right),\qquad\Phi_2=\sum_pa_p\cdot\exp\left(iF\left(\tau_{2D\cdot e_j-y}\cdot f_p\right)\right).
\]
The support of $\theta_j\cdot f_p$ is a compact subset of $\theta_jB=\left\{x\in\R^3:-D<x_j<0,\ |x_i|<D\right\}$, on which $x_j+x_i<D$;
the translation $\tau_{-y}$ lowers $x_j+x_i$ by $D$, and $r$ carries $\{x_j+x_i<0\}$ onto $\{x_j+x_i>0\}$.
On $B$ one has $x_j+x_i>-D$, and the translation $\tau_{2D\cdot e_j-y}$ raises $x_j+x_i$ by $D$.
So all the forms in the exponents have compact supports in $\{x_j+x_i>0\}$, and $\Phi_1$, $\Phi_2$ are bounded functions in $L^2(\Sigma_{\{x_j+x_i>0\}},\mu)$.

Let $k,k'\ge0$ be integers, and let
\[
v=(k+k')\cdot s\cdot e_j+(k-k')\cdot s\cdot e_i=\sqrt2\cdot s\cdot\left(k\cdot n_++k'\cdot n_-\right).
\]
By \eqref{eq:continuum-reconstruction}, by the invariance of $\mu$ under $\tau_{-y}$, which makes $U(-y)$ unitary on $L^2(\mu)$, and by $\hat r^2=1$, which holds because $r^2$ is the identity,
\begin{align*}
&\left\langle[\Phi]_j,\exp\left(-\left(2D+(k+k')\cdot s\right)\cdot\mathsf H_j\right)V_j\left((k-k')\cdot s\cdot e_i\right)[\Phi]_j\right\rangle\\
&\qquad=\mu\left((\Theta_j\Phi)\cdot U(2D\cdot e_j+v)\Phi\right)
=\left\langle\overline{\Theta_j\Phi},U(2D\cdot e_j+v)\Phi\right\rangle_\mu\\
&\qquad=\left\langle U(-y)\overline{\Theta_j\Phi},U(v)U(2D\cdot e_j-y)\Phi\right\rangle_\mu
=\left\langle\hat r\Phi_1,U(v)\Phi_2\right\rangle_\mu .
\end{align*}
Put
\[
\mathfrak T_+=\exp\left(-\sqrt2\cdot s\cdot\mathsf H'\right),\qquad\mathfrak V_-=V'\left(\sqrt2\cdot s\cdot e_i\right),\qquad\psi_1=\left[\hat g^{-1}\Phi_1\right]',\qquad\psi_2=\left[\hat g^{-1}\Phi_2\right]'.
\]
By \eqref{eq:continuum-rotated-picture} with $a=\sqrt2\cdot s\cdot k$ and $b=\sqrt2\cdot s\cdot k'$, and by the semigroup and group laws, the last member equals $\left\langle\psi_1,\mathfrak T_+^{k}\mathfrak V_-^{k'}\psi_2\right\rangle'$.
With
\[
\mathfrak T=\exp\left(-s\cdot\mathsf H_j\right),\qquad\mathfrak U=V_j(s\cdot e_i),
\]
the operator in the first member is $\mathfrak T^{2r_0+k+k'}\mathfrak U^{k-k'}$, because $D=r_0\cdot s$.
Hence the vector $\psi=[\Phi]_j$ satisfies the identity \eqref{eq:continuum-two-pictures} of Lemma~\ref{lem:abstract-cone}, with the integer $2r_0$ as its exponent and with the vectors $\psi_1$, $\psi_2$.

\item[(3)]
\emph{Proof of~(a).}
Let $\mathcal H_0$ be the set of the classes $[\Phi]_j$ of the finite sums $\Phi$ of step~\eqref{it:cone-dense-identity}.
These sums are dense in $L^2(\Sigma_j^+,\mu)$ \cite[Lemma~8.12\,(i)]{NoncompactQED}, and $[\,\cdot\,]_j$ is continuous with dense range;
so $\mathcal H_0$ is dense in $\mathcal H_j$.
The operators $\mathfrak T$ and $\mathfrak T_+$ are self-adjoint with $0\le\mathfrak T\le1$ and $0\le\mathfrak T_+\le1$, because $\mathsf H_j\ge0$ and $\mathsf H'\ge0$;
the operators $\mathfrak U$ and $\mathfrak V_-$ are unitary, and $\mathfrak T\mathfrak U=\mathfrak U\mathfrak T$.
By step~\eqref{it:cone-dense-identity}, every $\psi\in\mathcal H_0$ satisfies \eqref{eq:continuum-two-pictures} with suitable $\psi_1,\psi_2\in\mathcal H'$.
Lemma~\ref{lem:abstract-cone} yields $(1-\mathfrak T)^2\ge\mathfrak T\left(2-\mathfrak U-\mathfrak U^*\right)$, and $\mathfrak U^*=V_j(-s\cdot e_i)$;
this is~(a).

\item[(4)]\itemmark{4}\label{it:cone-spectral-support}
\emph{Proof of~(b) and~(c).}
By Lemma~\ref{lem:joint-spectrum}(b), the operators $\exp\left(-s\cdot\mathsf H_j\right)$ and $V_j(\pm s\cdot e_i)$ are the integrals of the functions $\exp(-s\cdot\lambda)$ and $\exp\left(\pm i\,s\cdot k_i\right)$ with respect to $\mathfrak E_j$;
and integration with respect to $\mathfrak E_j$ is a $*$-homomorphism on the bounded Borel functions.
Hence
\begin{align*}
&\left(1-\exp\left(-s\cdot\mathsf H_j\right)\right)^{2}-\exp\left(-s\cdot\mathsf H_j\right)\left(2-V_j(s\cdot e_i)-V_j(-s\cdot e_i)\right)\\
&\qquad=\int\mathrm d\mathfrak E_j(\lambda,k)\,\left(\left(1-\exp(-s\cdot\lambda)\right)^{2}-\exp(-s\cdot\lambda)\cdot\left|1-\exp\left(i\,s\cdot k_i\right)\right|^{2}\right).
\end{align*}
The bounded real function under the integral equals
\begin{align*}
&\exp(-s\cdot\lambda)\cdot\left(4\sinh^2\left(\frac{s\cdot\lambda}{2}\right)-4\sin^2\left(\frac{s\cdot k_i}{2}\right)\right)\\
&\qquad=2\exp(-s\cdot\lambda)\cdot\left(\cosh(s\cdot\lambda)+\cos(s\cdot k_i)-2\right),
\end{align*}
so it is negative exactly where $\cosh(s\cdot\lambda)+\cos(s\cdot k_i)<2$.
By Lemma~\ref{lem:joint-spectrum}(d), the operator in the first member is non-negative if and only if $\mathfrak E_j$ vanishes on this set, that is, if and only if the image of $\mathfrak E_j$ under $(\lambda,k)\mapsto(\lambda,k_i)$ is carried by the set of the $(\lambda,k_i)$ with $\cosh(s\cdot\lambda)+\cos(s\cdot k_i)\ge2$.
This image is the joint spectral measure of $(\mathsf H_j,\mathsf P_i)$ (Lemma~\ref{lem:joint-spectrum}(c)).
So the two properties in~(b) are equivalent for each $s>0$, and~(a) gives the first assertion of~(b).

For $\lambda\ge0$ the computation above shows that $\cosh(s\cdot\lambda)+\cos(s\cdot k_i)\ge2$ holds if and only if $\sinh\left(\frac{s\cdot\lambda}{2}\right)\ge\left|\sin\left(\frac{s\cdot k_i}{2}\right)\right|$.
The measure $\mathfrak E_j$ is carried by $[0,\infty)\times e_j^\perp$ (Lemma~\ref{lem:joint-spectrum}(b)).
By~(b) for the steps $s=1,\frac12,\frac13,\dots$, it is therefore carried by the intersection of the countably many sets
\[
\left\{(\lambda,k)\in[0,\infty)\times e_j^\perp:\ \frac2s\cdot\sinh\left(\frac{s\cdot\lambda}{2}\right)\ge\frac2s\cdot\left|\sin\left(\frac{s\cdot k_i}{2}\right)\right|\right\}.
\]
As $s\to0$, the left side of the inequality tends to $\lambda$ and the right side to $|k_i|$;
so $\lambda\ge|k_i|$ on the intersection.
This proves the first assertion of~(c).
Let $\psi\in\mathcal H_j$ and $\varrho_\psi=\langle\psi,\mathfrak E_j(\cdot)\psi\rangle$.
By Lemma~\ref{lem:joint-spectrum}(b) and~(c), $\mathsf H_j$ and $\mathsf P_i$ are the integrals of $\lambda$ and of $k_i$ with respect to $\mathfrak E_j$.
By the functional calculus, $\psi$ lies in the domain of $\mathsf H_j$ if and only if $\int\mathrm d\varrho_\psi(\lambda,k)\,\lambda^2<\infty$, and then
\[
\|\mathsf H_j\psi\|^2=\int\mathrm d\varrho_\psi(\lambda,k)\,\lambda^2,\qquad\langle\psi,\mathsf H_j\psi\rangle=\int\mathrm d\varrho_\psi(\lambda,k)\,\lambda;
\]
in the same way $\psi$ lies in the domain of $\mathsf P_i$ if and only if $\int\mathrm d\varrho_\psi(\lambda,k)\,k_i^2<\infty$, and then $\|\mathsf P_i\psi\|^2=\int\mathrm d\varrho_\psi(\lambda,k)\,k_i^2$ and $\langle\psi,|\mathsf P_i|\psi\rangle=\int\mathrm d\varrho_\psi(\lambda,k)\,|k_i|$.
Since $|k_i|\le\lambda$ holds $\varrho_\psi$-almost everywhere, the domain of $\mathsf H_j$ is contained in that of $\mathsf P_i$, and the two inequalities of~(c) follow.

\item[(5)]\itemmark{5}\label{it:cone-five-point}
\emph{Proof of~(d).}
Let $\Phi\in L^2(\Sigma_j^+,\mu)$ and $t>0$, and let now $\mathfrak T=\exp\left(-t\cdot\mathsf H_j\right)$ and $\mathfrak U=V_j(t\cdot e_i)$.
By \eqref{eq:continuum-reconstruction}, $\left\langle[\Phi]_j,\exp\left(-a\cdot\mathsf H_j\right)V_j(b\cdot e_i)[\Phi]_j\right\rangle=u(a\cdot e_j+b\cdot e_i)$ for $a\ge0$ and $b\in\R$;
in particular
\[
\begin{aligned}
u(0)&=\left\langle[\Phi]_j,[\Phi]_j\right\rangle,&\qquad u(t\cdot e_j)&=\left\langle[\Phi]_j,\mathfrak T[\Phi]_j\right\rangle,\\
u(2t\cdot e_j)&=\left\langle[\Phi]_j,\mathfrak T^2[\Phi]_j\right\rangle,&\qquad u(t\cdot e_j\pm t\cdot e_i)&=\left\langle[\Phi]_j,\mathfrak T\mathfrak U^{\pm1}[\Phi]_j\right\rangle.
\end{aligned}
\]
Since $\mathfrak U^*=\mathfrak U^{-1}$ and
\[
(1-\mathfrak T)^2-\mathfrak T\left(2-\mathfrak U-\mathfrak U^*\right)=1-4\mathfrak T+\mathfrak T^2+\mathfrak T\mathfrak U+\mathfrak T\mathfrak U^*,
\]
the left side of \eqref{eq:continuum-five-point} equals
\[
\left\langle[\Phi]_j,\left((1-\mathfrak T)^2-\mathfrak T\left(2-\mathfrak U-\mathfrak U^*\right)\right)[\Phi]_j\right\rangle,
\]
which is non-negative by~(a) with $s=t$.
Conversely, if this number is non-negative for all $\Phi$ in a set whose classes are dense in $\mathcal H_j$, then the quadratic form of the bounded self-adjoint operator $(1-\mathfrak T)^2-\mathfrak T\left(2-\mathfrak U-\mathfrak U^*\right)$, which is continuous on $\mathcal H_j$, is non-negative on all of $\mathcal H_j$;
that is,~(a) holds with $s=t$.
\end{enumerate}
\end{proof}

\paragraph{Tensor tests.}
Every $F\in\mathcal S'$ is extended to complex test forms by complex linearity: $F(f+if')=F(f)+iF(f')$ for real test forms $f$ and $f'$.
For an integer $n\ge1$ let $\mathscr S_n$ be the space of the Schwartz functions on $\R^{3n}=(\R^3)^n$ with values in the complexification of $(\Lambda^2)^{\otimes n}$, with the seminorms $q_m$, $m\ge0$, of \cite[Section~7.6]{NoncompactQED}, where this space is denoted $\mathcal S_n$;
let $\mathscr S_0=\C$.
A point of $\R^{3n}$ is written $x=(x_1,\dots,x_n)$ with $x_p\in\R^3$, and, for $j\in\{1,2,3\}$, $(x_p)_j$ is the $j$-th coordinate of $x_p$.
For test forms $f_1,\dots,f_n$, real or complex, the \emph{product} $f_1\otimes\dots\otimes f_n\in\mathscr S_n$ is the function $x\mapsto f_1(x_1)\otimes\dots\otimes f_n(x_n)$.
For $f\in\mathscr S_m$ and $f'\in\mathscr S_n$ the function $f\otimes f'\in\mathscr S_{m+n}$ is $(x,x')\mapsto f(x)\otimes f'(x')$, and $c\otimes f=f\otimes c=c\cdot f$ for $c\in\mathscr S_0$.
The derivatives of $f\otimes f'$ are the functions $\partial^\alpha f\otimes\partial^\beta f'$, and $1+|x|^2+|x'|^2\le\left(1+|x|^2\right)\cdot\left(1+|x'|^2\right)$;
hence $q_{m'}(f\otimes f')\le q_{m'}(f)\cdot q_{m'}(f')$ for every integer $m'\ge0$.
A Euclidean map $\mathfrak g\colon x\mapsto gx+v$ with $g\in\mathrm O(3)$ and $v\in\R^3$ acts on $\mathscr S_n$ by
\[
\left((\mathfrak g\,\cdot)^{\otimes n}f\right)(x)=\left(\wedge^2g\right)^{\otimes n}f\left(\mathfrak g^{-1}x_1,\dots,\mathfrak g^{-1}x_n\right)\qquad\left(x\in\R^{3n}\right);
\]
this is a continuous linear bijection of $\mathscr S_n$ onto itself, an affine change of the variables composed with a constant linear map of the values, and
\[
(\mathfrak g\,\cdot)^{\otimes n}(f_1\otimes\dots\otimes f_n)=(\mathfrak g\cdot f_1)\otimes\dots\otimes(\mathfrak g\cdot f_n),
\]
with the action on test forms of Definition~\ref{def:laws}.
The map $(\mathfrak g\,\cdot)^{\otimes0}$ is the identity of $\mathscr S_0$, and $\tau_v^{\otimes n}=(\tau_v\,\cdot)^{\otimes n}$ translates all $n$ arguments by $v$.
For $j\in\{1,2,3\}$ the \emph{reflections of tensor tests in the direction $j$} are the antilinear bijections $\Theta^{(n)}$ of $\mathscr S_n$ onto itself, $n\ge0$, given by
\[
\Theta^{(n)}f=\overline{(\theta_j\,\cdot)^{\otimes n}f}\quad\left(f\in\mathscr S_n,\ n\ge1\right),\qquad\Theta^{(0)}c=\overline c\quad\left(c\in\mathscr S_0\right).
\]

Let $h_\kappa$ be the Hermite two-forms, $\mathsf A$ the operator and $\lambda_\kappa$ its eigenvalues as in \cite[Section~7.1.5]{NoncompactQED}.
The $h_\kappa$ are real test forms and form an orthonormal basis of $L^2(\R^3;\Lambda^2)$;
the operator $\mathsf A$ is symmetric on the test forms, and $\mathsf Ah_\kappa=\lambda_\kappa\cdot h_\kappa$ with $\lambda_\kappa\ge4$;
every $F\in\mathcal S'$ lies in a space $H_{-s}$ with an integer $s\ge0$, and then $|F(h_\kappa)|\le\|F\|_{-s}\cdot\lambda_\kappa^{s}$ for every $\kappa$.
For a multi-index $\boldsymbol\kappa=(\kappa_1,\dots,\kappa_n)$ put $h_{\boldsymbol\kappa}=h_{\kappa_1}\otimes\dots\otimes h_{\kappa_n}$, and for $f\in\mathscr S_n$ put $a_{\boldsymbol\kappa}(f)=\langle h_{\boldsymbol\kappa},f\rangle_{L^2}$, which is written $c_{\boldsymbol\kappa}(f)$ in \cite[Section~7.6]{NoncompactQED}.
The $h_{\boldsymbol\kappa}$ form an orthonormal basis of $L^2\left(\R^{3n};(\Lambda^2)^{\otimes n}\right)$.
Let $J_1\subset J_2\subset\dots$ be finite sets of multi-indices whose union is the set of all multi-indices, and put $f_{J_k}=\sum_{\boldsymbol\kappa\in J_k}a_{\boldsymbol\kappa}(f)\cdot h_{\boldsymbol\kappa}$ for $f\in\mathscr S_n$;
then $f_{J_k}\to f$ in $\mathscr S_n$ as $k\to\infty$ \cite[Sections~7.1.5 and~7.6]{NoncompactQED}.
The $h_{\boldsymbol\kappa}$ are products of real test forms.
Hence the finite linear combinations of products of real test forms are dense in $\mathscr S_n$, and two continuous linear maps on $\mathscr S_n$ that agree on these products agree everywhere.

\begin{lemma}[Tensor tests]\label{lem:tensor-tests}
Let $\mu$ be a law on $\mathcal S'$ that satisfies \textup{(G)}, let $j\in\{1,2,3\}$, and let $\Theta^{(n)}$, $n\ge0$, be the reflections of tensor tests in the direction $j$.
\begin{enumerate}
\item[(a)]
Let $n\ge1$, and let $f_1,\dots,f_n\in L^2(\R^3;\Lambda^2)$ be real, $F(f_p)$ denoting the element of $L^2(\mu)$ given by \cite[Lemma~8.15\,(ii)]{NoncompactQED}.
Then $F(f_1)\cdots F(f_n)$ lies in $L^2(\mu)$, and
\[
\left\|F(f_1)\cdots F(f_n)\right\|_{L^2(\mu)}\le\left(2n\cdot2^{n}\cdot(n-1)!\right)^{\frac12}\cdot\prod_{p=1}^{n}\|f_p\|_2 .
\]
\item[(b)]
Let $n\ge1$ and $f\in\mathscr S_n$.
For every $F\in\mathcal S'$ the series
\begin{equation}\label{eq:continuum-tensor-series}
F^{(n)}(f)=\sum_{\boldsymbol\kappa}a_{\boldsymbol\kappa}(f)\cdot\prod_{p=1}^{n}F(h_{\kappa_p})
\end{equation}
converges absolutely;
for fixed $F$, the map $f\mapsto F^{(n)}(f)$ is a continuous linear functional on $\mathscr S_n$, and $F^{(n)}(f_1\otimes\dots\otimes f_n)=F(f_1)\cdots F(f_n)$ for all test forms $f_1,\dots,f_n$.
The function $F\mapsto F^{(n)}(f)$ on $\mathcal S'$ is $\Sigma$-measurable and lies in $L^2(\mu)$, and
\begin{align*}
\left\|F^{(n)}(f)\right\|_{L^2(\mu)}
&\le\left(2n\cdot2^{n}\cdot(n-1)!\right)^{\frac12}\cdot\sum_{\boldsymbol\kappa}\left|a_{\boldsymbol\kappa}(f)\right|\\
&\le\left(2n\cdot2^{n}\cdot(n-1)!\right)^{\frac12}\cdot4^{-n}\cdot\left(34\pi\sqrt3\right)^{n}\cdot q_{6n}(f).
\end{align*}
In particular $f\mapsto F^{(n)}(f)$ is the unique continuous linear map from $\mathscr S_n$ into $L^2(\mu)$ that takes the value $F(f_1)\cdots F(f_n)$ at $f_1\otimes\dots\otimes f_n$ for all real test forms $f_1,\dots,f_n$.
\item[(c)]
Let $F^{(0)}(c)=c$ for $c\in\mathscr S_0$.
For every $F\in\mathcal S'$, all integers $m,n\ge0$, $f\in\mathscr S_m$ and $f'\in\mathscr S_n$,
\[
F^{(m)}(f)\cdot F^{(n)}(f')=F^{(m+n)}(f\otimes f'),\qquad\overline{F^{(n)}(f')}=F^{(n)}\left(\overline{f'}\right).
\]
\item[(d)]
For every Euclidean map $\mathfrak g$, every integer $n\ge0$ and every $f\in\mathscr S_n$, one has $\hat{\mathfrak g}F^{(n)}(f)=F^{(n)}\left((\mathfrak g\,\cdot)^{\otimes n}f\right)$ at every point of $\mathcal S'$.
In particular $U(v)F^{(n)}(f)=F^{(n)}\left(\tau_v^{\otimes n}f\right)$ for $v\in\R^3$, and $\Theta_jF^{(n)}(f)=F^{(n)}\left(\Theta^{(n)}f\right)$.
\item[(e)]
Let $n\ge1$, and let $U\subset\R^3$ be open.
If the support of $f\in\mathscr S_n$ is a compact subset of $U^n$, there are finite linear combinations of products of elements of $C^\infty_{\mathrm c}(U;\Lambda^2)$, all supported in one compact subset of $U^n$, that converge to $f$ in $\mathscr S_n$, hence uniformly together with all their derivatives;
so these combinations are sequentially dense in the space of the smooth maps of $\R^{3n}$ into the complexification of $(\Lambda^2)^{\otimes n}$ whose supports are compact subsets of $U^n$, endowed with its usual topology.
If $f\in\mathscr S_n$ and $\supp f\subset U^n$, then $F^{(n)}(f)$ lies in $L^2(\Sigma_U,\mu)$.
\item[(f)]
Let $n\ge1$ and let $f\in\mathscr S_n$ be supported in $\left\{x\in\R^{3n}:(x_p)_j\ge0\text{ for }p=1,\dots,n\right\}$.
Then $F^{(n)}(f)$ lies in $L^2(\Sigma_j^+,\mu)$.
\end{enumerate}
\end{lemma}

\begin{proof}
(a)
Let $f\in L^2(\R^3;\Lambda^2)$ be real.
By \cite[Lemma~8.15\,(i), (ii)]{NoncompactQED}, where $F(f)$ is written $X(f)$, the moment bound for the exponent $2n$ reads
\[
\int\mathrm d\mu(F)\,|F(f)|^{2n}\le2n\cdot2^{n}\cdot(n-1)!\cdot\|f\|_2^{2n}.
\]
H\"older's inequality with $n$ exponents equal to $n$ yields
\[
\int\mathrm d\mu(F)\,\prod_{p=1}^{n}|F(f_p)|^2\le\prod_{p=1}^{n}\left(\int\mathrm d\mu(F)\,|F(f_p)|^{2n}\right)^{\frac1n}\le2n\cdot2^{n}\cdot(n-1)!\cdot\prod_{p=1}^{n}\|f_p\|_2^2 .
\]

(b)
Let $F\in\mathcal S'$, and let $s\ge0$ be an integer with $F\in H_{-s}$.
The operator $(\mathsf A^s)^{\otimes n}$, which applies $\mathsf A^s$ in each of the $n$ variables, maps $\mathscr S_n$ continuously into itself \cite[Section~7.6]{NoncompactQED};
as $\mathsf A$ is symmetric and $\mathsf Ah_\kappa=\lambda_\kappa\cdot h_\kappa$,
\[
a_{\boldsymbol\kappa}\left((\mathsf A^s)^{\otimes n}f\right)=\lambda_{\kappa_1}^{s}\cdots\lambda_{\kappa_n}^{s}\cdot a_{\boldsymbol\kappa}(f).
\]
Together with \cite[Lemma~7.16]{NoncompactQED}, applied to $(\mathsf A^s)^{\otimes n}f$, this gives
\[
\sum_{\boldsymbol\kappa}\left|a_{\boldsymbol\kappa}(f)\right|\cdot\prod_{p=1}^{n}\left|F(h_{\kappa_p})\right|
\le\|F\|_{-s}^{n}\cdot\sum_{\boldsymbol\kappa}\left|a_{\boldsymbol\kappa}\left((\mathsf A^s)^{\otimes n}f\right)\right|
\le\|F\|_{-s}^{n}\cdot\mathsf Z_4^{\frac n2}\cdot\left\|\left(\mathsf A^{s+2}\right)^{\otimes n}f\right\|_{L^2}.
\]
The last member is a continuous seminorm of $f$ on $\mathscr S_n$.
Hence the series \eqref{eq:continuum-tensor-series} converges absolutely, and $f\mapsto F^{(n)}(f)$ is a continuous linear functional on $\mathscr S_n$.
For test forms $f_1,\dots,f_n$, Fubini's theorem gives $a_{\boldsymbol\kappa}(f_1\otimes\dots\otimes f_n)=\prod_{p=1}^{n}\langle h_{\kappa_p},f_p\rangle_{L^2}$.
So the series for $f_1\otimes\dots\otimes f_n$ is the product of the $n$ series $\sum_\kappa\langle h_\kappa,f_p\rangle_{L^2}\cdot F(h_\kappa)$, which converge absolutely by the case $n=1$;
the $p$-th of them has the sum $F(f_p)$, because the Hermite expansion of $f_p$ converges to $f_p$ in the space of test forms and $F$ is continuous and complex linear.
This proves $F^{(n)}(f_1\otimes\dots\otimes f_n)=F(f_1)\cdots F(f_n)$.

The partial sums of \eqref{eq:continuum-tensor-series} over the sets $J_k$ are polynomials in the $\Sigma$-measurable functions $F\mapsto F(h_\kappa)$, and they converge at every point of $\mathcal S'$;
so $F\mapsto F^{(n)}(f)$ is $\Sigma$-measurable.
At every point $|F^{(n)}(f)|\le\sum_{\boldsymbol\kappa}|a_{\boldsymbol\kappa}(f)|\cdot\prod_p|F(h_{\kappa_p})|$.
Minkowski's inequality and monotone convergence, followed by~(a) for the real forms $h_{\kappa_1},\dots,h_{\kappa_n}$ of norm one, give the first bound of~(b);
\cite[Lemma~7.16]{NoncompactQED} gives the second.
So $f\mapsto F^{(n)}(f)$ is a continuous linear map from $\mathscr S_n$ into $L^2(\mu)$ with the asserted values on products;
it is the only one, by the density of the combinations of products in $\mathscr S_n$.

(c)
Fix $F\in\mathcal S'$.
If $m=0$ or $n=0$, the product rule is the convention $F^{(0)}(c)=c$ together with $c\otimes f=c\cdot f$.
Let $m,n\ge1$.
If $f$ and $f'$ are products of test forms, both sides of the product rule equal the product of the values of $F$ at the $m+n$ factors, by~(b).
Let $f'$ be such a product.
The maps $f\mapsto F^{(m)}(f)\cdot F^{(n)}(f')$ and $f\mapsto F^{(m+n)}(f\otimes f')$ are continuous linear functionals on $\mathscr S_m$, the second one because $q_{m'}(f\otimes f')\le q_{m'}(f)\cdot q_{m'}(f')$;
they agree on products, hence on $\mathscr S_m$.
For fixed $f\in\mathscr S_m$ the same argument in the variable $f'$ gives the product rule for every $f'\in\mathscr S_n$.
The $h_\kappa$ are real, and so is $F(h_\kappa)$;
hence $a_{\boldsymbol\kappa}\left(\overline{f'}\right)=\overline{a_{\boldsymbol\kappa}(f')}$, and \eqref{eq:continuum-tensor-series} gives $\overline{F^{(n)}(f')}=F^{(n)}\left(\overline{f'}\right)$;
for $n=0$ both sides equal the complex conjugate $\overline{f'}$ of the number $f'$, since $F^{(0)}(c)=c$.

(d)
For $n=0$ both sides equal $f$.
Let $n\ge1$ and $F\in\mathcal S'$.
By the definition of the action of $\mathfrak g^{-1}$ on $\mathcal S'$, $(\mathfrak g^{-1}F)(f')=F(\mathfrak g\cdot f')$ for every test form $f'$.
The maps $f\mapsto(\mathfrak g^{-1}F)^{(n)}(f)$ and $f\mapsto F^{(n)}\left((\mathfrak g\,\cdot)^{\otimes n}f\right)$ are continuous linear functionals on $\mathscr S_n$, and at a product $f_1\otimes\dots\otimes f_n$ both take the value $F(\mathfrak g\cdot f_1)\cdots F(\mathfrak g\cdot f_n)$;
so they coincide.
As $(\hat{\mathfrak g}\Phi)(F)=\Phi(\mathfrak g^{-1}F)$ for every function $\Phi$ on $\mathcal S'$, this is the first assertion.
The second is the case $\mathfrak g=\tau_v$, since $U(v)=\hat\tau_v$.
For the third, $\Theta_j\Phi=\overline{\hat\theta_j\Phi}$, and~(c) gives $\Theta_jF^{(n)}(f)=\overline{F^{(n)}\left((\theta_j\,\cdot)^{\otimes n}f\right)}=F^{(n)}\left(\Theta^{(n)}f\right)$.

(e)
Let first the support of $f$ be a compact subset of $U^n$.
The images of the compact set $\supp f\subset U^n$ under the $n$ projections $x\mapsto x_p$ are compact subsets of $U$.
Let $\chi$ be a real smooth function with compact support in $U$ that equals one on a neighbourhood of their union, and put $\chi^{\otimes n}(x)=\chi(x_1)\cdots\chi(x_n)$.
Then $\chi^{\otimes n}\cdot f=f$, and multiplication by $\chi^{\otimes n}$, a smooth function whose derivatives are all bounded, is continuous on $\mathscr S_n$.
As $f_{J_k}\to f$ in $\mathscr S_n$, also $\chi^{\otimes n}\cdot f_{J_k}\to f$ in $\mathscr S_n$, and
\[
\chi^{\otimes n}\cdot f_{J_k}=\sum_{\boldsymbol\kappa\in J_k}a_{\boldsymbol\kappa}(f)\cdot(\chi\cdot h_{\kappa_1})\otimes\dots\otimes(\chi\cdot h_{\kappa_n})
\]
is a finite linear combination of products of the forms $\chi\cdot h_\kappa\in C^\infty_{\mathrm c}(U;\Lambda^2)$, supported in the compact set $(\supp\chi)^n\subset U^n$.
Convergence in $\mathscr S_n$ implies uniform convergence of every derivative;
this gives the first assertion of~(e).
By~(b),
\[
F^{(n)}\left(\chi^{\otimes n}\cdot f_{J_k}\right)=\sum_{\boldsymbol\kappa\in J_k}a_{\boldsymbol\kappa}(f)\cdot\prod_{p=1}^{n}F(\chi\cdot h_{\kappa_p})
\]
is a polynomial in $\Sigma_U$-measurable functions;
it lies in $L^2(\mu)$, hence in $L^2(\Sigma_U,\mu)$, and it converges to $F^{(n)}(f)$ in $L^2(\mu)$ as $k\to\infty$, by~(b).
The subspace $L^2(\Sigma_U,\mu)$ is closed in $L^2(\mu)$;
so $F^{(n)}(f)$ lies in it.

Let now $f\in\mathscr S_n$ be any function with $\supp f\subset U^n$.
Let $\zeta$ be a real smooth function with compact support on $\R^3$ that equals one on the closed unit ball, and for $0<\delta\le1$ let $\zeta_\delta$ be the function $x\mapsto\zeta(\delta\cdot x_1)\cdots\zeta(\delta\cdot x_n)$ on $\R^{3n}$.
The support of $\zeta_\delta\cdot f$ lies in $\supp f\cap\supp\zeta_\delta$, a compact subset of $U^n$;
by the case just treated, $F^{(n)}(\zeta_\delta\cdot f)$ lies in $L^2(\Sigma_U,\mu)$.
For a multi-index $\alpha\neq0$ one has $\partial^\alpha\zeta_\delta(x)=\delta^{|\alpha|}\cdot\left(\partial^\alpha\zeta_1\right)(\delta\cdot x)$.
Hence, as $\delta\le1$, for every integer $m\ge0$ the derivatives of order at most $m$ of $1-\zeta_\delta$ are bounded by a number $C_m$ that does not depend on $\delta$;
they vanish on the open set of the points $x$ with $|x|<\frac1\delta$, on which $\zeta_\delta=1$.
At the points with $|x|\ge\frac1\delta$ one has $\left(1+|x|^2\right)^{\frac m2}\le\delta\cdot\left(1+|x|^2\right)^{\frac{m+1}{2}}$, and Leibniz's rule gives
\[
q_m\left((1-\zeta_\delta)\cdot f\right)\le2^m\cdot C_m\cdot\delta\cdot q_{m+1}(f).
\]
Thus $\zeta_\delta\cdot f\to f$ in $\mathscr S_n$ as $\delta\to0$, and $F^{(n)}(\zeta_\delta\cdot f)\to F^{(n)}(f)$ in $L^2(\mu)$ by~(b);
being closed, $L^2(\Sigma_U,\mu)$ contains the limit $F^{(n)}(f)$.

(f)
Let $U=\{x\in\R^3:x_j>0\}$, so that $L^2(\Sigma_U,\mu)=L^2(\Sigma_j^+,\mu)$ (Definition~\ref{def:positive-side}).
For $\varepsilon>0$ the function $f_\varepsilon=\tau_{\varepsilon\cdot e_j}^{\otimes n}f$ is supported in $\left\{x:(x_p)_j\ge\varepsilon\text{ for }p=1,\dots,n\right\}$, a subset of $U^n$;
by~(e), $F^{(n)}(f_\varepsilon)$ lies in $L^2(\Sigma_j^+,\mu)$.
For $v\in\R^3$ with $|v|\le1$, the mean value theorem and the inequality $1+|x|^2\le2\cdot\left(1+|x-y|^2\right)\cdot\left(1+|y|^2\right)$ for $x,y\in\R^{3n}$ give $q_m\left(\tau_v^{\otimes n}f-f\right)\le C\cdot|v|\cdot q_{m+1}(f)$ for every integer $m\ge0$, with a number $C$ that depends only on $m$ and $n$.
So $f_\varepsilon\to f$ in $\mathscr S_n$ as $\varepsilon\to0$, and $F^{(n)}(f_\varepsilon)\to F^{(n)}(f)$ in $L^2(\mu)$ by~(b).
The subspace $L^2(\Sigma_j^+,\mu)$ is closed;
so it contains $F^{(n)}(f)$.
\end{proof}

\paragraph{Schwinger functions.}
Let $\mu$ be a law on $\mathcal S'$ that satisfies \textup{(G)}.
For $n\ge1$ the \emph{Schwinger function} $S^\mu_n$ is the linear functional
\[
S^\mu_n(f)=\int\mathrm d\mu(F)\,F^{(n)}(f)\qquad\left(f\in\mathscr S_n\right),
\]
and $S^\mu_0(c)=c$ for $c\in\mathscr S_0$.
By Lemma~\ref{lem:tensor-tests}(b), $S^\mu_n$ is continuous on $\mathscr S_n$, and $S^\mu_n(f_1\otimes\dots\otimes f_n)=\int\mathrm d\mu(F)\,F(f_1)\cdots F(f_n)$ for all test forms $f_1,\dots,f_n$.

\begin{corollary}[Unbounded observables and Schwinger functions]\label{cor:unbounded-observables}
Let $i\neq j$ be indices in $\{1,2,3\}$, let $\mu$ be a law on $\mathcal S'$ that satisfies the hypotheses of Theorem~\ref{thm:cone} for these indices and \textup{(G)}, and let $\Theta^{(m)}$, $m\ge0$, be the reflections of tensor tests in the direction $j$.
Then the space $L^2(\Sigma_j^+,\mu)$ of Theorem~\ref{thm:cone}(d) contains
every complex polynomial in the functions $F(f)$ with real $f\in L^2(\R^3;\Lambda^2)$ that vanish almost everywhere on $\{x_j<0\}$;
every function $\exp\left(\sum_pz_p\cdot F(f_p)\right)$ with finitely many $z_p\in\C$ and such forms $f_p$;
and every function $F^{(n)}(f)$ with $n\ge1$ and $f\in\mathscr S_n$ supported in $\left\{x\in\R^{3n}:(x_p)_j\ge0\text{ for }p=1,\dots,n\right\}$.
Let $f_n\in\mathscr S_n$ for $n\ge0$, with $f_n=0$ for all but finitely many $n$ and with $f_n$ supported in $\left\{x\in\R^{3n}:(x_p)_j\ge0\text{ for }p=1,\dots,n\right\}$ for $n\ge1$, and put
\[
u(v)=\sum_{m,n\ge0}S^\mu_{m+n}\left(\Theta^{(m)}f_m\otimes\tau_v^{\otimes n}f_n\right)\qquad(v\in\R^3).
\]
Then $u$ satisfies the inequality \eqref{eq:continuum-five-point} for every $t>0$.
\end{corollary}

\begin{proof}
Let $f\in L^2(\R^3;\Lambda^2)$ be real and vanish almost everywhere on $\{x_j<0\}$.
As the plane $\{x_j=0\}$ is a null set, $f$ vanishes almost everywhere outside $\{x_j>0\}$;
so $F(f)$ lies in $L^2(\Sigma_j^+,\mu)$, and the bound of \textup{(G)} holds for $f$ \cite[Lemma~8.15\,(ii)]{NoncompactQED}.
Choose real $\Sigma_j^+$-measurable representatives of the functions $F(f)$ of this kind.
A finite product of them is $\Sigma_j^+$-measurable and square integrable, by Lemma~\ref{lem:tensor-tests}(a);
hence every complex polynomial in them lies in $L^2(\Sigma_j^+,\mu)$.
For finitely many such forms $f_p$ and numbers $z_p\in\C$, the function $\exp\left(\sum_pz_p\cdot F(f_p)\right)$ is $\Sigma_j^+$-measurable;
its squared modulus is $\exp F(2h)$ for the real form $h=\sum_p(\Re z_p)\cdot f_p$, which vanishes almost everywhere on $\{x_j<0\}$, and \textup{(G)} for $2h$ gives
\[
\int\mathrm d\mu(F)\,\exp F(2h)\le\exp\left(2\|h\|_2^2\right).
\]
The functions $F^{(n)}(f)$ of the statement lie in $L^2(\Sigma_j^+,\mu)$ by Lemma~\ref{lem:tensor-tests}(f).

Let now $(f_n)_{n\ge0}$ be as in the statement, and let $\Phi=\sum_{n\ge0}F^{(n)}(f_n)$, an element of $L^2(\Sigma_j^+,\mu)$.
Since $\Theta_j$ is additive and antilinear, Lemma~\ref{lem:tensor-tests}(c) and~(d) give, at every point of $\mathcal S'$ and for every $v\in\R^3$,
\[
(\Theta_j\Phi)\cdot U(v)\Phi=\sum_{m,n\ge0}F^{(m)}\left(\Theta^{(m)}f_m\right)\cdot F^{(n)}\left(\tau_v^{\otimes n}f_n\right)=\sum_{m,n\ge0}F^{(m+n)}\left(\Theta^{(m)}f_m\otimes\tau_v^{\otimes n}f_n\right).
\]
Each term on the right is $\mu$-integrable, by Lemma~\ref{lem:tensor-tests}(b), and its integral is $S^\mu_{m+n}\left(\Theta^{(m)}f_m\otimes\tau_v^{\otimes n}f_n\right)$.
Hence $\mu\left((\Theta_j\Phi)\cdot U(v)\Phi\right)=u(v)$ for every $v\in\R^3$.
Theorem~\ref{thm:cone}(d), applied to $\Phi$, now gives \eqref{eq:continuum-five-point} for every $t>0$.
\end{proof}

\subsubsection{The free Maxwell law}\label{sec:free-maxwell}

The space $L^2(\R^3;\Lambda^2)$ carries the scalar product
\[
\langle f,g\rangle:=\int\mathrm{d}x\,\sum_{\mu<\nu}\overline{f_{\mu\nu}(x)}\cdot g_{\mu\nu}(x)
\]
and the norm $\|f\|_2=\langle f,f\rangle^{\frac12}$.
To a two-form $f$ belongs its dual vector field $\ast f=(f_{23},-f_{13},f_{12})$, as in Lemma~\ref{lem:orientation-signs}(c);
since $|\ast f(x)|^2=\sum_{\mu<\nu}|f_{\mu\nu}(x)|^2$ at every point,
the map $f\mapsto\ast f$ is an isometry of $L^2(\R^3;\Lambda^2)$ onto the space of square integrable vector fields on $\R^3$,
whose scalar product is again written $\langle\cdot,\cdot\rangle$.
Fourier transforms of vector fields are formed componentwise, with the normalisation
\[
\hat b(p):=\int\mathrm{d}x\,\exp\left(-i\,p\cdot x\right)\cdot b(x)\qquad(p\in\R^3)
\]
for integrable $b$.
By Plancherel's theorem the transform extends to square integrable fields, and
$\langle b,b'\rangle=(2\pi)^{-3}\int\mathrm{d}^3p\,\overline{\hat b(p)}\cdot\hat b'(p)$.

\begin{definition}[Transverse projection and free Maxwell law]\label{def:free-maxwell}
For $p\in\R^3\setminus\{0\}$ let
\[
P^{\perp}(p):=1-\frac{p\,p^{T}}{|p|^{2}}
\]
be the orthogonal projection of $\C^3$ onto the orthogonal complement of $p$.
The \emph{transverse projection} $P^{\perp}$ is the operator on $L^2(\R^3;\Lambda^2)$ defined by
\[
\widehat{\ast\left(P^{\perp}f\right)}(p)=P^{\perp}(p)\,\widehat{\ast f}(p)\qquad\text{for almost every }p\in\R^3 .
\]
The \emph{free Maxwell law} is the law $\gamma$ on $(\mathcal{S}',\Sigma)$ with the characteristic functional
\[
\gamma\left(\exp\left(i\,F(f)\right)\right)=\exp\left(-\frac{1}{2}\cdot\left\|P^{\perp}f\right\|_2^{2}\right)\qquad(f\in\mathcal{S}).
\]
\end{definition}

Every matrix $P^{\perp}(p)$ is an orthogonal projection with real entries, and $P^{\perp}(-p)=P^{\perp}(p)$;
by Plancherel's theorem, $P^{\perp}$ is therefore an orthogonal projection of $L^2(\R^3;\Lambda^2)$ that maps real two-forms to real two-forms.
In \cite[Section~7.5]{NoncompactQED} this projection is written $\Pi_c$,
and \cite[Section~8.9.4]{NoncompactQED}, with the same normalisation of the Fourier transform, writes $\Pi$ for it as an operator on dual vector fields.
The law $\gamma$ is the free Maxwell law of \cite[Lemma~8.30]{NoncompactQED}.
Because $F(f)$ is linear in $f$, the characteristic functional shows that the variables $F(f)$, $f\in\mathcal{S}$, are jointly centred Gaussian under $\gamma$,
and polarisation of $\|P^{\perp}f\|_2^2$ gives their covariance:
\[
\gamma\left(F(f)\cdot F(g)\right)=\langle f,P^{\perp}g\rangle\qquad(f,g\in\mathcal{S}).
\]
By \cite[Lemma~8.30\,(i)]{NoncompactQED}, the law $\gamma$ satisfies \textup{(G)}, is invariant under all Euclidean maps of $\R^3$, and satisfies $(RP_j)$ for $j=1,2,3$;
in particular it has the properties \textup{(G)}, $(T_{123})$, $(R_j)$ and $(RP_j)$, $j=1,2,3$, of Definition~\ref{def:laws}.
Lemma~\ref{lem:joint-spectrum} therefore applies to $\gamma$ in each of the three directions.

\begin{lemma}[Uniqueness of Laplace--Fourier transforms]\label{lem:laplace-fourier}
Let $\nu_1$ and $\nu_2$ be finite positive Borel measures on $[0,\infty)\times\R^2$ such that
\[
\int\mathrm{d}\nu_1(\lambda,k)\,\exp\left(-n\cdot\lambda+i\,k\cdot y\right)=\int\mathrm{d}\nu_2(\lambda,k)\,\exp\left(-n\cdot\lambda+i\,k\cdot y\right)
\]
for every integer $n\ge0$ and every $y\in\R^2$.
Then $\nu_1=\nu_2$.
\end{lemma}

\begin{proof}
To begin with, fix an integer $n\ge0$.
For $a=1,2$ the assignment
\[
B\mapsto\int_{[0,\infty)\times B}\mathrm{d}\nu_a(\lambda,k)\,\exp(-n\cdot\lambda),\qquad B\subset\R^2\text{ a Borel set},
\]
is a finite positive Borel measure on $\R^2$,
and the hypothesis says that the two measures so obtained, for $a=1$ and for $a=2$, have the same Fourier transform.
By the uniqueness theorem for Fourier transforms of finite measures \cite[Theorem~1.3.6]{Rudin1962} they coincide; that is,
\[
\int_{[0,\infty)\times B}\mathrm{d}\nu_1(\lambda,k)\,\exp(-n\cdot\lambda)=\int_{[0,\infty)\times B}\mathrm{d}\nu_2(\lambda,k)\,\exp(-n\cdot\lambda)
\]
for every Borel set $B\subset\R^2$ and every integer $n\ge0$.

Now fix a Borel set $B\subset\R^2$.
For $a=1,2$ let $\nu_a^B$ be the image of the finite measure $S\mapsto\nu_a(S\times B)$ on $[0,\infty)$ under the map $\lambda\mapsto\exp(-\lambda)$ of $[0,\infty)$ into $(0,1]$;
it is a finite positive Borel measure on $[0,1]$ that gives no mass to the point $0$.
In terms of these measures the identity just proved reads
\[
\int\mathrm{d}\nu_1^B(t)\,t^{n}=\int\mathrm{d}\nu_2^B(t)\,t^{n}\qquad(n=0,1,2,\dots),
\]
so $\nu_1^B$ and $\nu_2^B$ give the same integral to every polynomial.
By the Weierstrass approximation theorem \cite[Theorem~7.26]{Rudin1976} every continuous function on $[0,1]$ is a uniform limit of polynomials;
as both measures are finite, they give the same integral to every continuous function on $[0,1]$.
Finite positive Borel measures on $[0,1]$ are regular, and the uniqueness part of the Riesz representation theorem \cite[Theorem~2.14]{Rudin1987} yields $\nu_1^B=\nu_2^B$.
Since $\lambda\mapsto\exp(-\lambda)$ maps $[0,\infty)$ homeomorphically onto $(0,1]$,
for every Borel set $S\subset[0,\infty)$ the set $\exp(-S):=\{\exp(-\lambda):\lambda\in S\}$ is a Borel subset of $(0,1]$ whose preimage is $S$, and
\[
\nu_1(S\times B)=\nu_1^B\left(\exp(-S)\right)=\nu_2^B\left(\exp(-S)\right)=\nu_2(S\times B).
\]
The sets $S\times B$, with $S\subset[0,\infty)$ and $B\subset\R^2$ Borel, are closed under finite intersections, the whole space $[0,\infty)\times\R^2$ is one of them,
and they generate the Borel $\sigma$-algebra of $[0,\infty)\times\R^2$.
Dynkin's $\pi$-$\lambda$ theorem therefore extends the equality of $\nu_1$ and $\nu_2$ on these sets to all Borel sets: $\nu_1=\nu_2$.
\end{proof}

\begin{lemma}[Linear observables of the free Maxwell law]\label{lem:maxwell-spectrum}
Let $j\in\{1,2,3\}$, let $\gamma$ be the free Maxwell law of Definition~\ref{def:free-maxwell},
and let $\left(\mathcal{H}_j,[\,\cdot\,]_j,\mathsf{H}_j,V_j,\Omega_j\right)$ be its reconstruction in the direction $j$,
with the joint spectral measure $\mathfrak{E}_j$ and the momenta $\mathsf{P}_i$, $i\ne j$, of Lemma~\ref{lem:joint-spectrum}.
Write the points of $\R^3$ as $x=s\cdot e_j+z$ with $s\in\R$ and $z\in e_j^{\perp}$,
where $e_j^{\perp}$ is identified with $\R^2$ as in Lemma~\ref{lem:joint-spectrum};
for $k=(k_i)_{i\ne j}\in\R^2$ put
\[
k\cdot z=\sum_{i\neq j}k_i\cdot z_i,\qquad|k|=|k|_2=\left(\sum_{i\neq j}k_i^{2}\right)^{\frac{1}{2}},
\]
so that $|k|$ is the Euclidean norm of $k$.
Let $f\in C^{\infty}_{\mathrm c}(\{x_j>0\};\Lambda^2)$ be a real two-form, put $b:=\ast f$ and $h:=\nabla\cdot b$, and for $s\in\R$ and $k\in\R^2$ let
\[
\tilde h(s,k):=\int_{\R^2}\mathrm{d}^2z\,\exp\left(-i\,k\cdot z\right)\cdot h(s\cdot e_j+z),\qquad
a_f(k):=\int_0^{\infty}\mathrm{d}s\,\exp(-|k|\cdot s)\cdot\tilde h(s,k).
\]
Let $\nu_f$ be the positive Borel measure on $\R^2$ given by
\[
\mathrm{d}\nu_f(k):=\frac{1}{(2\pi)^{2}}\cdot\frac{|a_f(k)|^{2}}{2|k|}\,\mathrm{d}^2k .
\]
\begin{enumerate}
\item[(a)]
For all $t\ge0$ and $y\in e_j^{\perp}$,
\[
\left\langle\theta_j\cdot f,P^{\perp}\tau_{t\cdot e_j+y}\cdot f\right\rangle=\int_{\R^2}\mathrm{d}\nu_f(k)\,\exp\left(-|k|\cdot t-i\,k\cdot y\right).
\]
The measure $\nu_f$ is finite, and its total mass is $\langle\theta_j\cdot f,P^{\perp}f\rangle\le\|f\|_2^{2}$.
\item[(b)]
The variable $F(f)$ lies in $L^2(\Sigma_j^+,\gamma)$,
and the spectral measure $\varrho_\psi$ of the vector $\psi:=[F(f)]_j$ of $\mathcal{H}_j$ is the image of $\nu_f$ under the map $k\mapsto(|k|,-k)$ of $\R^2$ into $[0,\infty)\times\R^2$.
In particular, $\varrho_\psi$ is carried by $\{(\lambda,k):\lambda=|k|_2\}$.
\item[(c)]
Suppose that $b=b_j\cdot e_j$, that is, the only component of $f$ that does not vanish identically is the one whose pair of indices does not contain $j$,
and that $b_j\in C^{\infty}_{\mathrm c}(\{x_j>0\})$ is non-negative and not identically zero.
Put
\[
\tilde b_j(s,k):=\int_{\R^2}\mathrm{d}^2z\,\exp\left(-i\,k\cdot z\right)\cdot b_j(s\cdot e_j+z),\qquad
\tilde a_f(k):=\int_0^{\infty}\mathrm{d}s\,\exp(-|k|\cdot s)\cdot\tilde b_j(s,k).
\]
Then $a_f(k)=|k|\cdot\tilde a_f(k)$ for $k\in\R^2$, the function $\tilde a_f$ is continuous, and $\tilde a_f(0)=\int\mathrm{d}x\,b_j(x)>0$.
Consequently there is $k_0>0$ such that, on $\{0<|k|<k_0\}$, the measure $\nu_f$ has with respect to Lebesgue measure the strictly positive density
\[
\frac{|k|\cdot|\tilde a_f(k)|^{2}}{8\pi^{2}} .
\]
\item[(d)]
Let $f$ be as in (c) and $\psi=[F(f)]_j$.
\begin{enumerate}
\item[(d1)]
$\varrho_\psi\left(\{0<\lambda<\varepsilon\}\right)>0$ for every $\varepsilon>0$;
hence $0$ is not an isolated point of the spectrum of $\mathsf{H}_j$.
\item[(d2)]
For $i\ne j$ and $0<\lambda_0<k_0$, the point $(\lambda,k_i)=(\lambda_0,\lambda_0)$ lies in the joint spectrum of $(\mathsf{H}_j,\mathsf{P}_i)$;
and for every $C>1$ there is a vector $\psi'$ in the domain of $\mathsf{H}_j$ with
$\langle\psi',\mathsf{H}_j\psi'\rangle<C\cdot\langle\psi',|\mathsf{P}_i|\psi'\rangle$.
\end{enumerate}
\end{enumerate}
\end{lemma}

\begin{proof}
\emph{(a)}
Let $v:=t\cdot e_j+y$, $b':=\ast(\theta_j\cdot f)$ and $b'':=\ast(\tau_v\cdot f)=b(\cdot-v)$.
Lemma~\ref{lem:orientation-signs}(c), applied to the reflection $\theta_j=r_{e_j,0}$, gives
\[
b'=\left(2e_je_j^{T}-1\right)\left(b\circ\theta_j\right),
\]
so that $b'_j=b_j\circ\theta_j$, while the two components of $b'$ transverse to $e_j$ are those of $b\circ\theta_j$ with the opposite sign.
Consequently $\nabla\cdot b'=-h\circ\theta_j$:
the derivative of $b_j\circ\theta_j$ in the direction $e_j$ acquires a minus sign through the chain rule, and the two transverse terms carry the minus sign of their components.
Furthermore $\nabla\cdot b''=h(\cdot-v)$.
Put $h_1:=\nabla\cdot b'$ and $h_2:=\nabla\cdot b''$;
let $\tilde h_1$ and $\tilde h_2$ be formed from $h_1$ and $h_2$ as $\tilde h$ is formed from $h$.
The support of $f$, which is also that of $b$, is a compact subset of the open half space $\{x_j>0\}$.
The reflection $\theta_j$ maps this half space onto $\{x_j<0\}$, and the translation by $v$ maps it onto $\{x_j>t\}$;
hence, by the formulas for $b'$ and $b''$, the support of $b'$ is a compact subset of $\{x_j<0\}$, and that of $b''$ a compact subset of $\{x_j>t\}$.
As $t\ge0$, these two open half spaces are disjoint, and so $\langle b',b''\rangle=0$.
The Fourier transform of a divergence is computed componentwise:
$\hat h_1(p)=i\,p\cdot\hat b'(p)$ and $\hat h_2(p)=i\,p\cdot\hat b''(p)$.
Since $f\mapsto\ast f$ is isometric, Plancherel's theorem and the definition of $P^{\perp}$ give
\begin{equation}\label{eq:continuum-maxwell-pairing}
\left\langle\theta_j\cdot f,P^{\perp}\tau_v\cdot f\right\rangle
=\langle b',b''\rangle-\int\frac{\mathrm{d}^3p}{(2\pi)^{3}}\,\frac{\overline{\left(p\cdot\hat b'(p)\right)}\left(p\cdot\hat b''(p)\right)}{|p|^{2}}
=-\int\frac{\mathrm{d}^3p}{(2\pi)^{3}}\,\frac{\overline{\hat h_1(p)}\cdot\hat h_2(p)}{|p|^{2}} .
\end{equation}
The integrand of the last member coincides with that of the integral in the middle member;
it is bounded in modulus by $|\hat b'(p)|\cdot|\hat b''(p)|$, which is integrable because $\hat b'$ and $\hat b''$ are square integrable.
By Fubini's theorem the last integral can therefore be computed as an iterated integral over $p=p_j\cdot e_j+k$, with $p_j\in\R$ and $k\in e_j^{\perp}=\R^2$, the integration in $p_j$ being performed first.
For the function $h$, and in the same way for $h_1$ and $h_2$,
\[
\hat h(p_j\cdot e_j+k)=\int\mathrm{d}s\,\exp\left(-i\,p_j\cdot s\right)\cdot\tilde h(s,k).
\]
Fix $k\ne0$ (the set $\{0\}$ is a null set of $\R^2$).
For this $k$ the slices $\tilde h_1(\cdot,k)$ and $\tilde h_2(\cdot,k)$ are bounded functions of compact support, while $\left(p_j^2+|k|^2\right)^{-1}$ is integrable in $p_j$.
The function $s\mapsto\exp(-|k|\cdot|s|)$ is continuous and integrable on $\R$, and its Fourier transform $p_j\mapsto2|k|/\left(p_j^2+|k|^2\right)$ is integrable as well;
so the function is the inverse Fourier transform of $p_j\mapsto2|k|/\left(p_j^2+|k|^2\right)$, and division by $2|k|$ gives
\[
\int\frac{\mathrm{d}p_j}{2\pi}\,\frac{\exp\left(i\,p_j\cdot(s-s')\right)}{p_j^{2}+|k|^{2}}=\frac{\exp\left(-|k|\cdot|s-s'|\right)}{2|k|}\qquad(s,s'\in\R).
\]
Together with Fubini's theorem in the variables $(p_j,s,s')$, this yields
\begin{equation}\label{eq:continuum-maxwell-slices}
\int\frac{\mathrm{d}p_j}{2\pi}\,\frac{\overline{\hat h_1(p_j\cdot e_j+k)}\cdot\hat h_2(p_j\cdot e_j+k)}{p_j^{2}+|k|^{2}}
=\int\mathrm{d}s\int\mathrm{d}s'\,\overline{\tilde h_1(s,k)}\cdot\tilde h_2(s',k)\cdot\frac{\exp\left(-|k|\cdot|s-s'|\right)}{2|k|} .
\end{equation}
As $\theta_j(s\cdot e_j+z)=-s\cdot e_j+z$, one has $\tilde h_1(s,k)=-\tilde h(-s,k)$;
moreover $\tilde h_2(s',k)=\exp\left(-i\,k\cdot y\right)\cdot\tilde h(s'-t,k)$.
Substitute $s=-s_1$ and $s'=s_2+t$.
Since $h$ vanishes on $\{x_j\le0\}$, only $s_1>0$ and $s_2>0$ contribute, and there $|s-s'|=s_1+s_2+t$.
After this substitution the right side of \eqref{eq:continuum-maxwell-slices} equals
\[
-\frac{\exp\left(-|k|\cdot t-i\,k\cdot y\right)}{2|k|}\cdot\overline{a_f(k)}\cdot a_f(k).
\]
As $(2\pi)^{-3}\,\mathrm{d}^3p=(2\pi)^{-2}\,\mathrm{d}^2k\cdot(2\pi)^{-1}\,\mathrm{d}p_j$ and $|p|^2=p_j^2+|k|^2$,
integration of this function over $k\in\R^2$ with the weight $(2\pi)^{-2}$, together with the minus sign in the last member of \eqref{eq:continuum-maxwell-pairing}, gives the identity of (a);
the integral over $k$ converges absolutely by Fubini's theorem.
For $t=0$ and $y=0$ its integrand is non-negative, and the identity reads
\[
\nu_f(\R^2)=\left\langle\theta_j\cdot f,P^{\perp}f\right\rangle\le\|\theta_j\cdot f\|_2\cdot\|f\|_2=\|f\|_2^{2},
\]
by the Cauchy--Schwarz inequality, because $P^{\perp}$ is an orthogonal projection and $\theta_j$ acts isometrically.
In particular $\nu_f$ is finite.

\emph{(b)}
The variable $F(f)$ is real, and it is $\Sigma_j^+$-measurable because $f$ has its support in $\{x_j>0\}$ (Definition~\ref{def:positive-side});
by \textup{(G)} and \cite[Lemma~8.15\,(i)]{NoncompactQED} it is square integrable under $\gamma$.
So $F(f)\in L^2(\Sigma_j^+,\gamma)$.
Let $t\ge0$ and $y\in e_j^{\perp}$.
By the properties of the reconstruction (Definition~\ref{def:laws}; \cite[Proposition~8.19\,(i)--(iii)]{NoncompactQED}),
and since $\Theta_jF(f)=F(\theta_j\cdot f)$ and $U(t\cdot e_j+y)F(f)=F(\tau_{t\cdot e_j+y}\cdot f)$ for the real form $f$,
\[
\left\langle\psi,\exp\left(-t\cdot\mathsf{H}_j\right)V_j(y)\psi\right\rangle
=\gamma\left(\left(\Theta_jF(f)\right)\cdot U(t\cdot e_j+y)F(f)\right)
=\gamma\left(F(\theta_j\cdot f)\cdot F(\tau_{t\cdot e_j+y}\cdot f)\right)
=\left\langle\theta_j\cdot f,P^{\perp}\tau_{t\cdot e_j+y}\cdot f\right\rangle .
\]
Lemma~\ref{lem:joint-spectrum}(b) expresses the left side as $\int\mathrm{d}\varrho_\psi(\lambda,k)\,\exp\left(-t\cdot\lambda+i\,k\cdot y\right)$,
and (a) expresses the right side as $\int\mathrm{d}\nu'(\lambda,k)\,\exp\left(-t\cdot\lambda+i\,k\cdot y\right)$,
where $\nu'$ denotes the image of $\nu_f$ under $k\mapsto(|k|,-k)$.
Both $\varrho_\psi$ and $\nu'$ are finite positive Borel measures on $[0,\infty)\times\R^2$, by Lemma~\ref{lem:joint-spectrum}(b) and by (a),
and their transforms agree, in particular for $t=0,1,2,\dots$;
Lemma~\ref{lem:laplace-fourier} now yields $\varrho_\psi=\nu'$.
The measure $\nu'$ is carried by the image of $k\mapsto(|k|,-k)$, which lies in $\{\lambda=|k|_2\}$.

\emph{(c)}
Here $h=\partial_jb_j$, so that $\tilde h(s,k)=\partial_s\tilde b_j(s,k)$.
For each $k$ the function $\tilde b_j(\cdot,k)$ is smooth and has compact support in $(0,\infty)$;
integrating by parts once, with vanishing boundary terms, one finds
\[
a_f(k)=\int_0^{\infty}\mathrm{d}s\,\exp(-|k|\cdot s)\cdot\partial_s\tilde b_j(s,k)=|k|\cdot\int_0^{\infty}\mathrm{d}s\,\exp(-|k|\cdot s)\cdot\tilde b_j(s,k)=|k|\cdot\tilde a_f(k).
\]
Written as an integral over $x=s\cdot e_j+z$, one has $\tilde a_f(k)=\int_{\{x_j>0\}}\mathrm{d}x\,\exp\left(-|k|\cdot x_j-i\,k\cdot z\right)\cdot b_j(x)$;
the integrand depends continuously on $k$ and is bounded in modulus by the integrable function $b_j$.
By dominated convergence $\tilde a_f$ is continuous, and $\tilde a_f(0)=\int\mathrm{d}x\,b_j(x)>0$ because $b_j$ is continuous, non-negative and not identically zero.
Choose $k_0>0$ with $\tilde a_f(k)\ne0$ for $|k|<k_0$.
On $\{0<|k|<k_0\}$ the density of $\nu_f$ with respect to Lebesgue measure is
\[
\frac{1}{(2\pi)^{2}}\cdot\frac{|k|^{2}\cdot|\tilde a_f(k)|^{2}}{2|k|}=\frac{|k|\cdot|\tilde a_f(k)|^{2}}{8\pi^{2}}>0 .
\]

\emph{(d1)}
By (b) and (c),
\[
\varrho_\psi\left(\{0<\lambda<\varepsilon\}\right)=\nu_f\left(\{0<|k|<\varepsilon\}\right)>0,
\]
since $\nu_f$ has a strictly positive density on the non-empty open set $\{0<|k|<\min(\varepsilon,k_0)\}$.
Thus $\mathfrak{E}_j\left(\{0<\lambda<\varepsilon\}\right)\ne0$ for every $\varepsilon>0$.
By Lemma~\ref{lem:joint-spectrum}(b) this projection is the spectral projection of $\mathsf{H}_j$ for the interval $(0,\varepsilon)$.
Every interval $(0,\varepsilon)$ therefore meets the spectrum of $\mathsf{H}_j$;
the spectrum is closed, so it contains $0$, and $0$ is not isolated in it.

\emph{(d2)}
Fix $i\ne j$ and $0<\lambda_0<k_0$, and for $0<\varepsilon<\lambda_0$ let
\[
S_\varepsilon:=\left\{(\lambda,k)\in[0,\infty)\times\R^2:\ |\lambda-\lambda_0|<\varepsilon,\ |k_i-\lambda_0|<\varepsilon\right\}.
\]
The preimage of $S_\varepsilon$ under $k\mapsto(|k|,-k)$ is an open subset of $\R^2$ that contains the point $k=-\lambda_0\cdot e_i$ of $e_j^{\perp}$, where $|k|=\lambda_0\in(0,k_0)$;
near this point $\nu_f$ has a continuous, strictly positive density by (c).
By (b), therefore,
\[
\varrho_\psi(S_\varepsilon)=\nu_f\left(\left\{k\in\R^2:(|k|,-k)\in S_\varepsilon\right\}\right)>0 .
\]
The set $S_\varepsilon$ is the preimage, under $(\lambda,k)\mapsto(\lambda,k_i)$, of the open square with centre $(\lambda_0,\lambda_0)$ and side $2\varepsilon$,
and these squares form a base of the neighbourhoods of $(\lambda_0,\lambda_0)$.
By Lemma~\ref{lem:joint-spectrum}(c) the joint spectral measure of $(\mathsf{H}_j,\mathsf{P}_i)$ is the image of $\mathfrak{E}_j$ under this map,
so its value on the square is $\mathfrak{E}_j(S_\varepsilon)$, which is not zero because $\|\mathfrak{E}_j(S_\varepsilon)\psi\|^2=\varrho_\psi(S_\varepsilon)>0$.
Hence the joint spectral measure vanishes on no neighbourhood of $(\lambda_0,\lambda_0)$, and this point lies in the joint spectrum.

Now let $C>1$.
Choose $\varepsilon\in(0,\lambda_0)$ with $\lambda_0+\varepsilon<C\cdot(\lambda_0-\varepsilon)$, and put $\psi':=\mathfrak{E}_j(S_\varepsilon)\psi$.
Then $\|\psi'\|^{2}=\varrho_\psi(S_\varepsilon)>0$, and the spectral measure of $\psi'$ is the restriction of $\varrho_\psi$ to $S_\varepsilon$.
At every point of $S_\varepsilon$ one has $0\le\lambda<\lambda_0+\varepsilon$ and $0<\lambda_0-\varepsilon<k_i<\lambda_0+\varepsilon$, hence $\lambda_0-\varepsilon<|k_i|<\lambda_0+\varepsilon$;
thus the two functions $(\lambda,k)\mapsto\lambda$ and $(\lambda,k)\mapsto|k_i|$ are bounded on $S_\varepsilon$, which carries the spectral measure of $\psi'$.
By Lemma~\ref{lem:joint-spectrum}(b), (c), the operators $\mathsf{H}_j$ and $|\mathsf{P}_i|$ are the integrals of these two functions with respect to $\mathfrak{E}_j$.
A vector belongs to the domain of such an integral when the square of the function is integrable with respect to the spectral measure of the vector,
and the scalar product of the vector with its image is then the integral of the function with respect to that measure.
Consequently $\psi'$ lies in the domain of $\mathsf{H}_j$ and in that of $|\mathsf{P}_i|$, and
\[
\langle\psi',\mathsf{H}_j\psi'\rangle\le(\lambda_0+\varepsilon)\cdot\|\psi'\|^{2}<C\cdot(\lambda_0-\varepsilon)\cdot\|\psi'\|^{2}\le C\cdot\langle\psi',|\mathsf{P}_i|\psi'\rangle .
\]
\end{proof}

\begin{remark}
For $j$, $f$ and $\psi=[F(f)]_j$ as in Lemma~\ref{lem:maxwell-spectrum}, the function $h$ is real, and therefore
\[
\tilde h(s,-k)=\overline{\tilde h(s,k)},\qquad a_f(-k)=\overline{a_f(k)}\qquad\left(s\in\R,\ k\in\R^2\right).
\]
In particular $\nu_f$ is invariant under $k\mapsto-k$.
Since the map $k\mapsto(|k|,k)$ is the map $k\mapsto(|k|,-k)$ preceded by $k\mapsto-k$, the spectral measure $\varrho_\psi$ is also the image of $\nu_f$ under $k\mapsto(|k|,k)$.

Let now $f'\in C^{\infty}_{\mathrm c}(\{x_j>0\};\Lambda^2)$ be a second real two-form;
form $h'$, $\tilde h'$, $a_{f'}$ and $\nu_{f'}$ from $f'$ as $h$, $\tilde h$, $a_f$ and $\nu_f$ are formed from $f$, and put $\psi'=[F(f')]_j$ and
\[
\mathrm{d}\nu_{f',f}(k):=\frac{1}{(2\pi)^{2}}\cdot\frac{\overline{a_{f'}(k)}\cdot a_f(k)}{2|k|}\,\mathrm{d}^2k .
\]
As $\nu_f$ and $\nu_{f'}$ are finite by Lemma~\ref{lem:maxwell-spectrum}(a), the Cauchy--Schwarz inequality shows that $\nu_{f',f}$ is a complex measure on $\R^2$ of total variation at most $\nu_{f'}(\R^2)^{\frac{1}{2}}\cdot\nu_f(\R^2)^{\frac{1}{2}}$; moreover $\nu_{f,f}=\nu_f$.
For $t\ge0$ and $y\in e_j^{\perp}$ one has
\begin{equation}\label{eq:continuum-cross-measures}
\int\mathrm{d}\left\langle\psi',\mathfrak{E}_j(\lambda,k)\psi\right\rangle\,\exp\left(-t\cdot\lambda+i\,k\cdot y\right)=\int_{\R^2}\mathrm{d}\nu_{f',f}(k)\,\exp\left(-|k|\cdot t-i\,k\cdot y\right).
\end{equation}
To see this, note first that by Lemma~\ref{lem:joint-spectrum}(b) the left side is $\left\langle\psi',\exp\left(-t\cdot\mathsf{H}_j\right)V_j(y)\psi\right\rangle$,
which the chain of equalities in the proof of Lemma~\ref{lem:maxwell-spectrum}(b), with $\Theta_jF(f')=F(\theta_j\cdot f')$ in the first factor, identifies with $\left\langle\theta_j\cdot f',P^{\perp}\tau_{t\cdot e_j+y}\cdot f\right\rangle$.
In the proof of Lemma~\ref{lem:maxwell-spectrum}(a), take now $b'=\ast(\theta_j\cdot f')$.
Its support is again a compact subset of $\{x_j<0\}$, while $h_1=-h'\circ\theta_j$ and $\tilde h_1(s,k)=-\tilde h'(-s,k)$.
As $h$ and $h'$ vanish on $\{x_j\le0\}$, the substitution $s=-s_1$, $s'=s_2+t$ turns the right side of \eqref{eq:continuum-maxwell-slices} into
\[
-\frac{\exp\left(-|k|\cdot t-i\,k\cdot y\right)}{2|k|}\cdot\overline{a_{f'}(k)}\cdot a_f(k),
\]
and integration over $k$ with the weight $(2\pi)^{-2}$, together with the minus sign in the last member of \eqref{eq:continuum-maxwell-pairing}, gives \eqref{eq:continuum-cross-measures}.

The complex measure $\left\langle\psi',\mathfrak{E}_j(\cdot)\psi\right\rangle$ on $[0,\infty)\times\R^2$ is the image of $\nu_{f',f}$ under the map $k\mapsto(|k|,-k)$.
For the proof, let $\nu_{f,f'}$ be formed as $\nu_{f',f}$, with the roles of $f$ and $f'$ exchanged.
As $f$ and $f'$ are arbitrary real two-forms in $C^{\infty}_{\mathrm c}(\{x_j>0\};\Lambda^2)$, the identity \eqref{eq:continuum-cross-measures} holds for every ordered pair of such forms;
in particular it remains true if $\psi'$, $\psi$ and $\nu_{f',f}$ are replaced by $\psi$, $\psi'$ and $\nu_{f,f'}$, by $\psi$, $\psi$ and $\nu_{f,f}=\nu_f$, or by $\psi'$, $\psi'$ and $\nu_{f',f'}=\nu_{f'}$.
Fix $n\in\{0,1,2,3\}$.
Since the scalar product of $\mathcal{H}_j$ is antilinear in the first entry and $\overline{i^n}=i^{-n}$, the vector $\psi+i^n\cdot\psi'$ has the spectral measure
\[
\varrho_{\psi+i^n\cdot\psi'}=\varrho_\psi+i^n\cdot\left\langle\psi,\mathfrak{E}_j(\cdot)\psi'\right\rangle+i^{-n}\cdot\left\langle\psi',\mathfrak{E}_j(\cdot)\psi\right\rangle+\varrho_{\psi'} ,
\]
a finite positive measure by Lemma~\ref{lem:joint-spectrum}.
For $t\ge0$ and $y\in e_j^{\perp}$, the four identities just listed turn the integral of $(\lambda,k)\mapsto\exp\left(-t\cdot\lambda+i\,k\cdot y\right)$ against the right side into the integral of $k\mapsto\exp\left(-|k|\cdot t-i\,k\cdot y\right)$ against the measure
\[
\nu_f+i^n\cdot\nu_{f,f'}+i^{-n}\cdot\nu_{f',f}+\nu_{f'}
\]
on $\R^2$, that is, into the integral of $(\lambda,k)\mapsto\exp\left(-t\cdot\lambda+i\,k\cdot y\right)$ against the image of this measure under $k\mapsto(|k|,-k)$.
This measure is a linear combination of finite measures, and its density with respect to Lebesgue measure is the product of $(2\pi)^{-2}\cdot(2|k|)^{-1}$ with
\[
|a_f(k)|^{2}+i^n\cdot\overline{a_f(k)}\cdot a_{f'}(k)+i^{-n}\cdot\overline{a_{f'}(k)}\cdot a_f(k)+|a_{f'}(k)|^{2}=\left|a_f(k)+i^n\cdot a_{f'}(k)\right|^{2}\ge0 ;
\]
it is therefore a finite positive measure.
So $\varrho_{\psi+i^n\cdot\psi'}$ and the image of this measure under $k\mapsto(|k|,-k)$ are finite positive Borel measures on $[0,\infty)\times\R^2$ whose integrals of $(\lambda,k)\mapsto\exp\left(-t\cdot\lambda+i\,k\cdot y\right)$ agree for all $t\ge0$ and $y\in e_j^{\perp}$, in particular for $t=0,1,2,\dots$;
Lemma~\ref{lem:laplace-fourier} shows that the two measures coincide.
Multiply now the identity for $\varrho_{\psi+i^n\cdot\psi'}$, and the measure $\nu_f+i^n\cdot\nu_{f,f'}+i^{-n}\cdot\nu_{f',f}+\nu_{f'}$, by $i^n$, add over $n=0,1,2,3$ and divide by $4$.
As $\sum_{n=0}^{3}i^n=\sum_{n=0}^{3}i^{2n}=0$ and $i^n\cdot i^{-n}=1$, this gives
\[
\begin{aligned}
\frac{1}{4}\sum_{n=0}^{3}i^n\cdot\varrho_{\psi+i^n\cdot\psi'}&=\left\langle\psi',\mathfrak{E}_j(\cdot)\psi\right\rangle,\\
\frac{1}{4}\sum_{n=0}^{3}i^n\cdot\left(\nu_f+i^n\cdot\nu_{f,f'}+i^{-n}\cdot\nu_{f',f}+\nu_{f'}\right)&=\nu_{f',f}.
\end{aligned}
\]
The image under $k\mapsto(|k|,-k)$ of a linear combination of finite measures on $\R^2$ is the same combination of their images,
and the image of $\nu_f+i^n\cdot\nu_{f,f'}+i^{-n}\cdot\nu_{f',f}+\nu_{f'}$ is $\varrho_{\psi+i^n\cdot\psi'}$ for each $n$.
By the second line of the last display, the image of $\nu_{f',f}$ is therefore $\frac{1}{4}\sum_{n=0}^{3}i^n\cdot\varrho_{\psi+i^n\cdot\psi'}$, which the first line identifies with $\left\langle\psi',\mathfrak{E}_j(\cdot)\psi\right\rangle$.

Since $f'$ is real as well, $a_{f'}(-k)=\overline{a_{f'}(k)}$;
hence the image of $\nu_{f',f}$ under $k\mapsto-k$ is the complex conjugate measure $\overline{\nu_{f',f}}$, which in general differs from $\nu_{f',f}$.
For instance, let $y_0\in e_j^{\perp}$ and $f'=\tau_{y_0}\cdot f$, again a real two-form with compact support in $\{x_j>0\}$, since the translation by $y_0$ maps this half space onto itself.
Then $h'=h(\cdot-y_0)$, as for $b''$ in the proof of Lemma~\ref{lem:maxwell-spectrum}(a), so that $\tilde h'(s,k)=\exp\left(-i\,k\cdot y_0\right)\cdot\tilde h(s,k)$ and
\[
\mathrm{d}\nu_{f',f}(k)=\exp\left(i\,k\cdot y_0\right)\cdot\mathrm{d}\nu_f(k).
\]
If moreover $f$ is as in Lemma~\ref{lem:maxwell-spectrum}(c) and $y_0\ne0$, then at $t=0$ and $y=y_0$ the right side of \eqref{eq:continuum-cross-measures} equals $\nu_f(\R^2)$, whereas
\[
\int_{\R^2}\mathrm{d}\nu_{f',f}(k)\,\exp\left(i\,k\cdot y_0\right)=\int_{\R^2}\mathrm{d}\nu_f(k)\,\cos\left(2\,k\cdot y_0\right)<\nu_f(\R^2);
\]
the equality holds because $\nu_f$ is invariant under $k\mapsto-k$,
and the inequality because the non-negative function $1-\cos\left(2\,k\cdot y_0\right)$ and the density of $\nu_f$ given in Lemma~\ref{lem:maxwell-spectrum}(c) are continuous and positive on a non-empty open subset of $\{0<|k|<k_0\}$.
The left member of the last display is the integral of the bounded continuous function $(\lambda,k)\mapsto\exp\left(i\,k\cdot y_0\right)$ against the image of $\nu_{f',f}$ under $k\mapsto(|k|,k)$,
while the integral of the same function against $\left\langle\psi',\mathfrak{E}_j(\cdot)\psi\right\rangle$ is the left side of \eqref{eq:continuum-cross-measures} at $t=0$ and $y=y_0$, that is, $\nu_f(\R^2)$.
Hence the complex measure $\left\langle\psi',\mathfrak{E}_j(\cdot)\psi\right\rangle$, which is the image of $\nu_{f',f}$ under $k\mapsto(|k|,-k)$, is in general not the image of $\nu_{f',f}$ under $k\mapsto(|k|,k)$.
\end{remark}

\begin{definition}[Scaled free Maxwell laws]\label{def:scaled-maxwell}
Let $\gamma$ be the free Maxwell law of Definition~\ref{def:free-maxwell}, and let $0<a\le1$.
The law $\gamma_a$ is the law of $a^{\frac12}F$ when $F$ has the law $\gamma$,
that is, the image of $\gamma$ under the map $F\mapsto a^{\frac12}F$ of $\mathcal{S}'$ onto itself;
this map is measurable for $\Sigma$, since $\left(a^{\frac12}F\right)(f)=a^{\frac12}\cdot F(f)$.
In particular $\gamma_1=\gamma$.
The characteristic functional of a law $\mu$ is the function $\hat\mu\colon\mathcal{S}\to\C$,
$\hat\mu(f):=\mu\left(\exp\left(iF(f)\right)\right)$.
By Definition~\ref{def:free-maxwell}, that of $\gamma_a$ is
\[
\hat\gamma_a(f)=\int\mathrm{d}\gamma_a(F)\,\exp\left(iF(f)\right)
=\hat\gamma\left(a^{\frac12}f\right)
=\exp\left(-\frac{a}{2}\cdot\|P^{\perp}f\|_2^2\right)\qquad(f\in\mathcal{S}).
\]
For finitely many $f_p\in\mathcal{S}$ and real numbers $t_p$, the linearity of $f\mapsto F(f)$ gives
\[
\gamma_a\left(\exp\left(i\sum_pt_p\cdot F(f_p)\right)\right)=\hat\gamma_a\left(\sum_pt_p\cdot f_p\right);
\]
by the formula for $\hat\gamma_a$, this joint characteristic function of the variables $F(f_p)$ under $\gamma_a$
is the exponential of a negative semidefinite quadratic form in the numbers $t_p$.
Hence the variables $F(f)$, $f\in\mathcal{S}$, are jointly centred Gaussian under $\gamma_a$,
and polarisation of the quadratic form $f\mapsto a\cdot\|P^{\perp}f\|_2^2$ gives their covariance
\[
\gamma_a\left(F(f)\cdot F(f')\right)=a\cdot\langle f,P^{\perp}f'\rangle\qquad(f,f'\in\mathcal{S}).
\]
Thus $\gamma_a$ is the centred Gaussian law with this covariance.
\end{definition}

\begin{lemma}\label{lem:scaled-maxwell}
Let $0<a\le1$.
The law $\gamma_a$ satisfies~\eqref{cond:soft-hypothesis},
is carried by the closed distributions, $\gamma_a\left(\left\{F\in\mathcal{S}':dF=0\right\}\right)=1$,
is invariant under $F\mapsto-F$,
and lies in $\mathcal{C}_9$.
For all real $f,f'\in C^{\infty}_{\mathrm{c}}(\R^3;\Lambda^2)$
\[
\lim_{|x|\to\infty}\gamma_a\left(\exp\left(iF(f)\right)\cdot\exp\left(iF(\tau_x\cdot f')\right)\right)
=\hat\gamma_a(f)\cdot\hat\gamma_a(f'),
\]
in whatever way $x\in\R^3$ tends to infinity;
and $\gamma_a$ is ergodic under the translations of $\R^3$, that is, $\mathcal{E}^{123}(\gamma_a)=\C\cdot1$.
A convex combination of finitely many laws on $\mathcal{S}'$,
each of which satisfies~\eqref{cond:soft-hypothesis}, is carried by the closed distributions,
is invariant under $F\mapsto-F$ and lies in $\mathcal{C}_9$,
has again each of these four properties.
\end{lemma}

\begin{proof}
By Definition~\ref{def:laws} and \cite[Section~8.1.3]{NoncompactQED},
a law $\mu$ is reflection positive at the plane $\{n\cdot x=c\}$,
where $n\in\R^3$ is a unit vector and $c\in\R$, if
\[
\sum_{p,q}\overline{z_p}\cdot z_q\cdot\hat\mu\left(f_q-r_{n,c}\cdot f_p\right)\ge0
\]
for all finite families of real forms $f_p\in C^{\infty}_{\mathrm{c}}(\{n\cdot x>c\};\Lambda^2)$ and numbers $z_p\in\C$;
indeed $\exp\left(i\left(F(f_q)-F(r_{n,c}\cdot f_p)\right)\right)=\exp\left(iF(f_q-r_{n,c}\cdot f_p)\right)$.

\emph{The law $\gamma$.}
By \cite[Lemma~8.30(i)]{NoncompactQED},
where $F(f)$ is written $X(f)$ and $P^{\perp}$ appears as the projection $\Pi$ acting on dual vector fields,
the law $\gamma$ satisfies \textup{(G)}, is invariant under every Euclidean map,
is carried by the closed distributions,
and satisfies $(RP_j)$ for $j=1,2,3$.
It is even reflection positive at every plane of $\R^3$.
To see this, let $n$ be a unit vector, let $c\in\R$, and let $g\in\mathrm{O}(3)$ satisfy $ge_1=n$;
as for the elements of $\Ocub$, $g$ also denotes the Euclidean map $x\mapsto gx$.
Since $e_1\cdot\left(g^{-1}x\right)=n\cdot x$ for $x\in\R^3$,
the map $g$ carries the half space $\{x_1>c\}$ onto $\{n\cdot x>c\}$, and
\[
r_{n,c}=g\circ r_{e_1,c}\circ g^{-1},
\]
where $r_{e_1,c}$ is the reflection in the plane $\{x_1=c\}$.
Let finitely many real $f_p\in C^{\infty}_{\mathrm{c}}(\{n\cdot x>c\};\Lambda^2)$ and $z_p\in\C$ be given,
and put $h_p:=g^{-1}\cdot f_p$.
Then $h_p\in C^{\infty}_{\mathrm{c}}(\{x_1>c\};\Lambda^2)$,
and $f_q-r_{n,c}\cdot f_p=g\cdot\left(h_q-r_{e_1,c}\cdot h_p\right)$.
The invariance of $\gamma$ under $g$ gives $\hat\gamma(g\cdot h)=\hat\gamma(h)$ for $h\in\mathcal{S}$, and therefore
\[
\sum_{p,q}\overline{z_p}\cdot z_q\cdot\hat\gamma\left(f_q-r_{n,c}\cdot f_p\right)
=\sum_{p,q}\overline{z_p}\cdot z_q\cdot\hat\gamma\left(h_q-r_{e_1,c}\cdot h_p\right).
\]
Being invariant under all Euclidean maps, $\gamma$ satisfies $(R_1)$ and $(T_1)$; it also satisfies $(RP_1)$.
By \cite[Lemma~8.16(5)]{NoncompactQED} it is therefore reflection positive at every plane $\{x_1=c\}$, $c\in\R$;
so the right side of the last identity is non-negative, and so is its left side.

\emph{Symmetries, positivity and domination of $\gamma_a$.}
Let $0<a\le1$.
For every Euclidean map $\mathfrak g$ and every $F\in\mathcal{S}'$
one has $\mathfrak g\left(a^{\frac12}F\right)=a^{\frac12}\cdot\mathfrak gF$ and $-\left(a^{\frac12}F\right)=a^{\frac12}\cdot(-F)$,
and $a^{\frac12}F$ is closed if and only if $F$ is closed.
As $\gamma_a$ is the image of $\gamma$ under $F\mapsto a^{\frac12}F$,
it follows that $\gamma_a$ is invariant under every Euclidean map,
that it is carried by the closed distributions,
and that it is invariant under $F\mapsto-F$ provided $\gamma$ is.
The image of $\gamma$ under $F\mapsto-F$ has the characteristic functional $f\mapsto\hat\gamma(-f)=\hat\gamma(f)$;
by the uniqueness of the law with this characteristic functional \cite[Lemma~8.30]{NoncompactQED}, this image is $\gamma$.

Let $\{n\cdot x=c\}$ be a plane, and let $f_p$, $z_p$ be as in the first paragraph.
The forms $a^{\frac12}f_p$ again lie in $C^{\infty}_{\mathrm{c}}(\{n\cdot x>c\};\Lambda^2)$,
and $r_{n,c}\cdot\left(a^{\frac12}f_p\right)=a^{\frac12}\cdot\left(r_{n,c}\cdot f_p\right)$; hence
\[
\sum_{p,q}\overline{z_p}\cdot z_q\cdot\hat\gamma_a\left(f_q-r_{n,c}\cdot f_p\right)
=\sum_{p,q}\overline{z_p}\cdot z_q\cdot\hat\gamma\left(a^{\frac12}f_q-r_{n,c}\cdot\left(a^{\frac12}f_p\right)\right)\ge0,
\]
by the reflection positivity of $\gamma$ at $\{n\cdot x=c\}$.
For $f\in\mathcal{S}$ the variable $F(f)$ is centred Gaussian under $\gamma$ with the variance $\|P^{\perp}f\|_2^2$, so that
\[
\int\mathrm{d}\gamma_a(F)\,\exp F(f)
=\int\mathrm{d}\gamma(F)\,\exp\left(a^{\frac12}\cdot F(f)\right)
=\exp\left(\frac{a}{2}\cdot\|P^{\perp}f\|_2^2\right)
\le\exp\left(\frac12\cdot\|f\|_2^2\right);
\]
here $a\le1$, and $\|P^{\perp}f\|_2\le\|f\|_2$ because $P^{\perp}$ is an orthogonal projection.
Thus $\gamma_a$ satisfies \textup{(G)}, it is invariant under the translations of $\R^3$ and under $\Ocub$,
and, by the case $n=e_j$, $c=0$, in which $r_{e_j,0}=\theta_j$, it satisfies $(RP_j)$ for $j=1,2,3$:
this is~\eqref{cond:soft-hypothesis}.
Moreover $\gamma_a$ is invariant under the reflection in every plane and reflection positive at every plane of $\R^3$,
in particular at the planes of the nine families;
together with its invariance under the translations, this shows $\gamma_a\in\mathcal{C}_9$.

\emph{Clustering.}
Let $f,f'\in C^{\infty}_{\mathrm{c}}(\R^3;\Lambda^2)$ be real, and let $x\in\R^3$.
One has $F(f)+F(\tau_x\cdot f')=F(f+\tau_x\cdot f')$, and, $P^{\perp}$ being an orthogonal projection,
\[
\|P^{\perp}(f+\tau_x\cdot f')\|_2^2
=\|P^{\perp}f\|_2^2+\|P^{\perp}(\tau_x\cdot f')\|_2^2+2\cdot\langle f,P^{\perp}(\tau_x\cdot f')\rangle .
\]
Moreover $\hat\gamma_a(\tau_x\cdot f')=\hat\gamma_a(f')$, since $\gamma_a$ is invariant under $\tau_x$.
Consequently
\[
\gamma_a\left(\exp\left(iF(f)\right)\cdot\exp\left(iF(\tau_x\cdot f')\right)\right)
=\hat\gamma_a(f+\tau_x\cdot f')
=\hat\gamma_a(f)\cdot\hat\gamma_a(f')\cdot\exp\left(-a\cdot\langle f,P^{\perp}(\tau_x\cdot f')\rangle\right).
\]
Choose $r>0$ such that the supports of $f$ and $f'$ lie in the cube $(-r,r)^3$, and let $|x|_\infty>2r$.
For $y\in\supp f$ and $y'\in\supp(\tau_x\cdot f')=x+\supp f'$
one has $|y-y'|\ge|y-y'|_\infty>|x|_\infty-2r$, where $|\cdot|$ is the Euclidean norm.
The dual vector field $\ast(\tau_x\cdot f')$ is the translate of $\ast f'$ by $x$,
and so has the same norm $\|\ast f'\|_{L^1}=\int\mathrm{d}y\,|(\ast f')(y)|$.
Applied to the vector fields $\ast f$ and $\ast(\tau_x\cdot f')$,
whose supports have the distance at least $|x|_\infty-2r$,
\cite[Lemma~8.30(ii)]{NoncompactQED} gives
\[
\left|\langle f,P^{\perp}(\tau_x\cdot f')\rangle\right|
\le\frac{\|\ast f\|_{L^1}\cdot\|\ast f'\|_{L^1}}{2\pi\cdot\left(|x|_\infty-2r\right)^3}
\qquad\left(|x|_\infty>2r\right).
\]
Since $|x|_\infty\ge3^{-\frac12}\cdot|x|$, the right side tends to zero as $|x|\to\infty$,
and the exponential factor above tends to one; this is the asserted limit.

\emph{Ergodicity.}
Let $\Phi$ and $\Phi'$ be finite linear combinations, with complex coefficients,
of functions $F\mapsto\exp\left(iF(f)\right)$ with real $f\in C^{\infty}_{\mathrm{c}}(\R^3;\Lambda^2)$.
The rules $\overline{\exp\left(iF(f)\right)}=\exp\left(iF(-f)\right)$,
$U(x)\exp\left(iF(f')\right)=\exp\left(iF(\tau_x\cdot f')\right)$
and $\hat\gamma_a(-f)=\overline{\hat\gamma_a(f)}$
turn the limit relation, applied term by term, into
\[
\lim_{|x|\to\infty}\gamma_a\left(\overline{\Phi}\cdot U(x)\Phi'\right)=\overline{\gamma_a(\Phi)}\cdot\gamma_a(\Phi').
\]
In particular this holds along every sequence of points of $\R^3$ that tends to infinity.
As $\gamma_a$ satisfies $(T_{123})$, \cite[Lemma~8.14(a)]{NoncompactQED}, applied with $J=\{1,2,3\}$,
shows that $\gamma_a$ is ergodic under the translations of $\R^3$.

\emph{Convex combinations.}
Let $\mu=\sum_kt_k\cdot\mu_k$ be a finite sum with $t_k\ge0$ and $\sum_kt_k=1$,
where every law $\mu_k$ satisfies~\eqref{cond:soft-hypothesis}, is carried by the closed distributions,
is invariant under $F\mapsto-F$ and lies in $\mathcal{C}_9$.
Each of these four properties is a conjunction of conditions,
each of which requires the integral of one fixed measurable function on $\mathcal{S}'$, bounded or with values in $[0,\infty]$,
to lie in one fixed convex set.
For the invariance under a translation, under an element of $\Ocub$,
under the reflection in a plane of the nine families or under $F\mapsto-F$,
the function is the difference of the indicator functions of a set $S\in\Sigma$ and of its preimage under the map,
and the set is $\{0\}$.
For \textup{(G)}, the function is $F\mapsto\exp F(f)$, $f\in\mathcal{S}$,
and the set is $\left[0,\exp\left(\frac12\cdot\|f\|_2^2\right)\right]$.
For reflection positivity at a plane $\{n\cdot x=c\}$, in particular for $(RP_j)$,
the function is $F\mapsto\sum_{p,q}\overline{z_p}\cdot z_q\cdot\exp\left(i\left(F(f_q)-F(r_{n,c}\cdot f_p)\right)\right)$,
with $f_p$, $z_p$ as in the first paragraph, and the set is $[0,\infty)$.
For being carried by the closed distributions,
the function is the indicator function of $\left\{F\in\mathcal{S}':dF=0\right\}$, and the set is $\{1\}$.
The integral of such a function with respect to $\mu$ is the sum of the $t_k$ times its integrals with respect to the $\mu_k$,
and a convex set contains every convex combination of its points;
so $\mu$ satisfies each of the conditions.
\end{proof}

\begin{remark}\label{rem:mirror-in-limit}
In the lattice argument for four-dimensional Yang--Mills theory,
the five-point inequality is first obtained on each lattice, for several steps together,
by means of the elementary inequalities
\[
\sinh(mt)\ge m\cdot\sinh t\quad(t\ge0),\qquad|\sin(mt)|\le m\cdot|\sin t|\quad(t\in\R),
\]
$m$ a positive integer;
the limit is then taken in this inequality,
and only for bounded lattice counterparts of products of fields whose supports are disjoint.
Here, by contrast, invariance under the reflections in the planes of the nine families
and reflection positivity at these planes are carried over to the limit:
by Proposition~\ref{prop:nine-families} every volume-first limit point lies in $\mathcal{C}_9$,
and Theorem~\ref{thm:cone}, applied to such a limit point, yields the cone.
A law of $\mathcal{C}_9$ satisfies the hypotheses of that theorem for every ordered pair $j\ne i$,
because the planes $\{x_j=0\}$ and $\{x_j+x_i=0\}$ belong to the nine families.
Consequently, for every $\mu\in\mathcal{C}_9$ and every ordered pair $j\ne i$,
the five-point inequality~\eqref{eq:continuum-five-point} holds for every $\Phi\in L^2(\Sigma_j^+,\mu)$, bounded or not;
a product of smeared fields enters with no condition on the supports of its test forms,
as soon as it lies in $L^2(\Sigma_j^+,\mu)$.
This applies to every ergodic component that lies in $\mathcal{C}_9$,
although such a component arises from the decomposition of a law and not as a limit of lattice laws.
Clause~(a) of Theorem~\ref{thm:cone} holds for every step $s>0$, so that all steps are covered at once.
Conversely, for every $s>0$ one has, on the set $\{\lambda\ge|k_i|\}$,
\[
\sinh\frac{s\lambda}{2}\ge\frac{s\lambda}{2}\ge\frac{s|k_i|}{2}\ge\left|\sin\frac{sk_i}{2}\right|,
\]
and hence
\[
\cosh(s\lambda)+\cos(sk_i)-2=2\cdot\left(\sinh^2\frac{s\lambda}{2}-\sin^2\frac{sk_i}{2}\right)\ge0;
\]
so clause~(c) of Theorem~\ref{thm:cone} implies clause~(b) for every $s>0$.
\end{remark}

\begin{remark}\label{rem:cone-optimal}
By Lemma~\ref{lem:scaled-maxwell} with $a=1$, the free Maxwell law $\gamma=\gamma_1$ lies in $\mathcal{C}_9$.
Let $j\ne i$, let $i'$ be the third index,
and let $\left(\mathcal{H}_j,[\,\cdot\,]_j,\mathsf{H}_j,V_j,\Omega_j\right)$ be the reconstruction of $\gamma$ in the direction $j$,
with the momentum $\mathsf{P}_i$, as in Lemma~\ref{lem:maxwell-spectrum}.
For every real $f\in C^{\infty}_{\mathrm{c}}(\{x_j>0\};\Lambda^2)$
the spectral measure $\varrho_\psi$ of the vector $\psi=[F(f)]_j$ of $\mathcal{H}_j$
is carried by the set $\{\lambda=|k|_2\}$ (Lemma~\ref{lem:maxwell-spectrum}(b)).
This set lies in the cone $\{\lambda\ge|k_i|\}$, because $|k|_2\ge|k_i|$,
and it meets the boundary $\{\lambda=|k_i|\}$ of that cone exactly in the two half-lines $\{k_{i'}=0,\ \lambda=|k_i|\}$.
Let $f$ be as in Lemma~\ref{lem:maxwell-spectrum}(c), and let $k_0>0$ be the number given there.
By Lemma~\ref{lem:maxwell-spectrum}(d2),
the joint spectrum of $(\mathsf{H}_j,\mathsf{P}_i)$ for $\gamma$ contains points with $\lambda=|k_i|>0$,
namely the points $(\lambda,k_i)=(\lambda_0,\lambda_0)$ with $0<\lambda_0<k_0$;
and for every $C>1$ there is a vector $\psi'$ in the domain of $\mathsf{H}_j$ with
\[
\langle\psi',\mathsf{H}_j\psi'\rangle<C\cdot\langle\psi',|\mathsf{P}_i|\psi'\rangle .
\]
So for no $C>1$ does the inequality $\mathsf{H}_j\ge C\cdot|\mathsf{P}_i|$ hold for all laws of $\mathcal{C}_9$:
the constant one in Theorem~\ref{thm:cone}(c) is optimal on $\mathcal{C}_9$.
\end{remark}

\subsection{Invariant vectors}\label{sec:invariant-vectors}

\begin{proposition}[Zero energy implies zero momentum]\label{prop:zero-momentum}
Let $\mu\in\mathcal{C}_9$ and $j\in\{1,2,3\}$.
Then $\mu$ belongs to the class $\mathcal{C}$;
let $(\mathcal{H}_j,[\,\cdot\,]_j,\mathsf{H}_j,V_j,\Omega_j)$ be its reconstruction in the direction $e_j$,
for $i\neq j$ let $\mathsf{P}_i$ be the self-adjoint generator of $u\mapsto V_j(u\cdot e_i)$,
and let $\mathfrak{E}_j$ be the joint spectral measure of Lemma~\ref{lem:joint-spectrum}, with $e_j^{\perp}$ identified with $\R^2$ as in that lemma.
With
\[
\mathcal{K}_j=\left\{(\lambda,k)\in[0,\infty)\times\R^2:\ \lambda\ge|k|_\infty\right\},
\]
the following hold.
\begin{enumerate}
\item[(a)] The projection-valued measure $\mathfrak{E}_j$ is carried by $\mathcal{K}_j$, and $\mathsf{H}_j\ge|\mathsf{P}_i|$ for both $i\neq j$.
\item[(b)] For both $i\neq j$ one has $\ker\mathsf{H}_j\subset\ker\mathsf{P}_i$.
The orthogonal projection $P_{\ker\mathsf{H}_j}$ of $\mathcal{H}_j$ onto $\ker\mathsf{H}_j$ equals $\mathfrak{E}_j\left(\{0\}\times\R^2\right)$, and
\[
V_j(y)\,P_{\ker\mathsf{H}_j}=P_{\ker\mathsf{H}_j}\,V_j(y)=P_{\ker\mathsf{H}_j}\qquad\left(y\in e_j^{\perp}\right).
\]
A vector of $\mathcal{H}_j$ lies in $\ker\mathsf{H}_j$ if and only if it is invariant under every $\exp(-t\cdot\mathsf{H}_j)$, $t\ge0$, and under every $V_j(y)$, $y\in e_j^{\perp}$.
\item[(c)] $\mathcal{E}^{j}(\mu)=\mathcal{E}^{123}(\mu)$.
Equivalently, $P^{j}_{\mu}=P^{123}_{\mu}$, which is the identity of Definition~\ref{def:symmetry-conditions}\eqref{it:projections-coincide} for the index $j$.
\end{enumerate}
\end{proposition}

\begin{proof}
Each coordinate plane $\{x_\nu=0\}$, $\nu\in\{1,2,3\}$, belongs to the nine families of planes, and so does the plane $\{x_j+x_i=0\}$ for $i\neq j$.
By the definition of $\mathcal{C}_9$ (Definition~\ref{def:laws}), the law $\mu$ is invariant under the translations of $\R^3$,
and for $\nu=1,2,3$ it is invariant under the reflection $\theta_\nu$ in $\{x_\nu=0\}$ and reflection positive at this plane.
These are the conditions $(T_{123})$, $(R_\nu)$ and $(RP_\nu)$, $\nu=1,2,3$, of \cite[Section~8.1.3]{NoncompactQED}; hence $\mu\in\mathcal{C}$.
In particular Lemma~\ref{lem:joint-spectrum} and \cite[Proposition~8.19]{NoncompactQED} apply to $\mu$ and $j$
(the companion paper writes $\rho$ for the law, $\mathcal{E}_+$ for $L^{2}(\Sigma_j^{+},\rho)$ and $H_j$ for $\mathsf{H}_j$).
For $i\neq j$ the law $\mu$ is, moreover, invariant under the reflection in $\{x_j+x_i=0\}$ and reflection positive at this plane;
so Theorem~\ref{thm:cone} applies to $\mu$ and to each of the two ordered pairs $(j,i)$, $i\neq j$.

\emph{(a)} Fix $i\neq j$.
By Lemma~\ref{lem:joint-spectrum}(c) the joint spectral measure of $(\mathsf{H}_j,\mathsf{P}_i)$ is the image of $\mathfrak{E}_j$ under $(\lambda,k)\mapsto(\lambda,k_i)$,
and by Theorem~\ref{thm:cone}(c) it is carried by the set of the points $(\lambda,k_i)$ with $\lambda\ge|k_i|$;
the same clause of the theorem gives $\mathsf{H}_j\ge|\mathsf{P}_i|$.
Consequently
\[
\mathfrak{E}_j\left(\left\{(\lambda,k):\ \lambda<|k_i|\right\}\right)=0\qquad(i\neq j).
\]
The number $|k|_\infty$ is the larger of the two numbers $|k_i|$, $i\neq j$;
so a point of $[0,\infty)\times\R^2$ that does not lie in $\mathcal{K}_j$ lies in one of the two sets just displayed.
As $\mathfrak{E}_j$ is carried by $[0,\infty)\times\R^2$ (Lemma~\ref{lem:joint-spectrum}(b)), it vanishes on the complement of $\mathcal{K}_j$.

\emph{(b)} Lemma~\ref{lem:joint-spectrum}(b), applied with $y=0$ to the indicator function of the set $\{0\}$,
shows that the spectral projection of $\mathsf{H}_j$ for $\{0\}$ equals $\mathfrak{E}_j\left(\{0\}\times\R^2\right)$;
and the spectral projection of a self-adjoint operator for $\{0\}$ is the orthogonal projection onto its kernel.
This identifies $P_{\ker\mathsf{H}_j}$.

Let $\psi\in\ker\mathsf{H}_j$, and let $\varrho_\psi=\langle\psi,\mathfrak{E}_j(\cdot)\psi\rangle$ be its spectral measure, as in Lemma~\ref{lem:joint-spectrum}.
Since $\psi=P_{\ker\mathsf{H}_j}\psi$, the measure $\varrho_\psi$ is carried by $\{0\}\times\R^2$;
by (a) it is also carried by $\mathcal{K}_j$.
It is therefore concentrated in the single point
\[
\left(\{0\}\times\R^2\right)\cap\mathcal{K}_j=\{(0,0)\}.
\]
For $y\in e_j^{\perp}$, Lemma~\ref{lem:joint-spectrum}(b) with $t=0$ represents $V_j(y)$ as the integral of $(\lambda,k)\mapsto\exp(ik\cdot y)$ against $\mathfrak{E}_j$; hence
\[
\left\|V_j(y)\psi-\psi\right\|^2=\int\mathrm{d}\varrho_\psi(\lambda,k)\,\left|\exp(ik\cdot y)-1\right|^2=0,
\]
because the integrand vanishes at $(0,0)$.
In the same manner, $\mathsf{P}_i$ is the integral of $(\lambda,k)\mapsto k_i$ against $\mathfrak{E}_j$ (Lemma~\ref{lem:joint-spectrum}(c)),
and $\int\mathrm{d}\varrho_\psi(\lambda,k)\,k_i^2=0$;
so $\psi$ lies in the domain of $\mathsf{P}_i$, and $\left\|\mathsf{P}_i\psi\right\|^2=\int\mathrm{d}\varrho_\psi(\lambda,k)\,k_i^2=0$.
One obtains $V_j(y)\psi=\psi$ for every $y\in e_j^{\perp}$, and $\ker\mathsf{H}_j\subset\ker\mathsf{P}_i$ for both $i\neq j$.

The range of $P_{\ker\mathsf{H}_j}$ is $\ker\mathsf{H}_j$, so the preceding paragraph yields $V_j(y)\,P_{\ker\mathsf{H}_j}=P_{\ker\mathsf{H}_j}$ for every $y\in e_j^{\perp}$.
Since $V_j(y)^{*}=V_j(-y)$, the adjoint of this identity reads $P_{\ker\mathsf{H}_j}\,V_j(-y)=P_{\ker\mathsf{H}_j}$;
with $-y$ in place of $y$ it becomes $P_{\ker\mathsf{H}_j}\,V_j(y)=P_{\ker\mathsf{H}_j}$.

Every vector of $\ker\mathsf{H}_j$ is invariant under $\exp(-t\cdot\mathsf{H}_j)$ for $t\ge0$, and, as shown above, under every $V_j(y)$.
Conversely, let $\psi\in\mathcal{H}_j$ satisfy $\exp(-t\cdot\mathsf{H}_j)\psi=\psi$ for all $t\ge0$.
Lemma~\ref{lem:joint-spectrum}(b) with $t=1$ and $y=0$ gives
\[
0=\left\|\exp(-\mathsf{H}_j)\psi-\psi\right\|^2=\int\mathrm{d}\varrho_\psi(\lambda,k)\,\left(\exp(-\lambda)-1\right)^2 .
\]
Thus $\varrho_\psi$ is carried by $\left\{(\lambda,k):\ \exp(-\lambda)=1\right\}=\{0\}\times\R^2$,
which means $\psi=\mathfrak{E}_j\left(\{0\}\times\R^2\right)\psi=P_{\ker\mathsf{H}_j}\psi$; so $\psi\in\ker\mathsf{H}_j$.

\emph{(c)} Let $\Phi\in\mathcal{E}^{j}(\mu)$.
By \cite[Lemma~8.18(i)]{NoncompactQED}, applied at the plane $\{x_j=0\}$, the vector $\Phi$ lies in $L^{2}(\Sigma_j^{+},\mu)$,
and by \cite[Proposition~8.19(iv)]{NoncompactQED} its class $[\Phi]_j$ lies in $\ker\mathsf{H}_j$.
Let $y\in e_j^{\perp}$.
The translation $\tau_y$ maps the half space $\{x_j>0\}$ onto itself; hence $U(y)$ maps $L^{2}(\Sigma_j^{+},\mu)$ into itself,
and $[U(y)\Phi']_j=V_j(y)[\Phi']_j$ for every $\Phi'$ in this space by \cite[Proposition~8.19(iii)]{NoncompactQED}.
Together with (b) this gives
\[
\left[U(y)\Phi-\Phi\right]_j=V_j(y)[\Phi]_j-[\Phi]_j=0 .
\]
Any two translations commute, so that $U(t\cdot e_j)\,U(y)\Phi=U(y)\,U(t\cdot e_j)\Phi=U(y)\Phi$ for $t\in\R$;
thus $U(y)\Phi-\Phi$ belongs to $\mathcal{E}^{j}(\mu)$.
On $\mathcal{E}^{j}(\mu)$ the class map $[\,\cdot\,]_j$ is an isometry \cite[Proposition~8.19(iv)]{NoncompactQED}, and therefore $\left\|U(y)\Phi-\Phi\right\|_{L^{2}(\mu)}=0$.
As $U(t\cdot e_j+y)=U(t\cdot e_j)\,U(y)$ for $t\in\R$ and $y\in e_j^{\perp}$, the vector $\Phi$ is invariant under $U(a)$ for every $a\in\R^3$, that is, $\Phi\in\mathcal{E}^{123}(\mu)$.
Conversely, a vector invariant under all translations is in particular invariant under those along $e_j$, whence $\mathcal{E}^{123}(\mu)\subset\mathcal{E}^{j}(\mu)$.
The two closed subspaces being equal, so are the orthogonal projections onto them.
\end{proof}

\subsection{Ergodic components}\label{sec:components}

\begin{proposition}[Properties of the ergodic components]\label{prop:ergodic-components}
Let $\mu$ be a law that satisfies \eqref{cond:soft-hypothesis} and lies in $\mathcal{C}_9$,
and let $\mu=\int\mathrm{d}\nu(\xi)\,\mu_\xi$ be its decomposition into $\R^3$-ergodic components
\textup{(}Definition~\ref{def:symmetry-conditions}\textup{)}.
For $\nu$-almost every $\xi$ the law $\mu_\xi$ has the following properties.
\begin{enumerate}
\item[\textup{(a)}] The law $\mu_\xi$ is invariant and ergodic under the translations of $\R^3$, and $\mu_\xi\in\mathcal{C}_9$.
\item[\textup{(b)}] The law $\mu_\xi$ satisfies \textup{(G)} with the same right side as $\mu$:
\[
\int\mathrm{d}\mu_\xi(F)\,\exp F(f)\le\exp\left(\frac{1}{2}\cdot\|f\|_2^2\right)\qquad\text{for every }f\in\mathcal{S}.
\]
\item[\textup{(c)}] The law $\mu_\xi$ is invariant under $\R^3\rtimes\Ocub$;
so $\mu$ satisfies condition~\eqref{it:components-cubic} of Definition~\ref{def:symmetry-conditions}.
\item[\textup{(d)}] For $j=1,2,3$ one has $\ker\mathsf{H}_j(\mu_\xi)=\C\cdot\Omega_j$,
where $\mathsf{H}_j(\mu_\xi)$ and $\Omega_j$ belong to the reconstruction of $\mu_\xi$ in the direction $j$
\textup{(}Definition~\ref{def:laws}\textup{)};
the joint spectral measure $\mathfrak{E}_j$ of this reconstruction \textup{(}Lemma~\ref{lem:joint-spectrum}\textup{)}
is carried by the cone
\[
\mathcal{K}_j=\left\{(\lambda,k)\in[0,\infty)\times\R^2:\ \lambda\ge|k|_\infty\right\},
\]
and $\mathsf{H}_j(\mu_\xi)\ge|\mathsf{P}_i|$ for both $i\neq j$.
\item[\textup{(e)}] For all $\Phi,\Phi'\in L^2(\mu_\xi)$,
\[
\mu_\xi\left(\overline{\Phi}\cdot U(a)\Phi'\right)\to\overline{\mu_\xi(\Phi)}\cdot\mu_\xi(\Phi')\qquad\text{as }|a|\to\infty,
\]
in whatever way $a\in\R^3$ tends to infinity.
\end{enumerate}
Moreover:
\begin{enumerate}
\item[\textup{(f)}] For each $S\in\Sigma$ with $\mu(S)=1$ one has $\mu_\xi(S)=1$ for $\nu$-almost every $\xi$,
the exceptional null set depending on $S$.
In particular, if $\mu(dF=0)=1$, then $\mu_\xi(dF=0)=1$ for $\nu$-almost every $\xi$.
\end{enumerate}
\end{proposition}

\begin{proof}
By Definition~\ref{def:laws}, the law $\mu$ lies in $\mathcal{C}_9\subset\mathcal{C}$:
it satisfies $(T_{123})$ and, for $j=1,2,3$, $(R_j)$ and $(RP_j)$,
and it is invariant under the reflection in each plane of the nine families and reflection positive at it.

\emph{(a)}
According to \cite[Theorem~8.1(a)]{NoncompactQED}, for $\nu$-almost every $\xi$ the law $\mu_\xi$ satisfies $(T_{123})$
and is ergodic under the translations of $\R^3$.
Let $j\in\{1,2,3\}$.
Since $\mu$ satisfies $(R_j)$ and $(RP_j)$, part~(b) of the same theorem shows that for $\nu$-almost every $\xi$
the law $\mu_\xi$ satisfies $(R_j)$ and is reflection positive at the plane $\{x_j=c\}$ for every $c\in\R$.
Next let $n$ be one of the six unit vectors
\[
\frac{1}{\sqrt{2}}\left(e_i+e_{i'}\right),\qquad\frac{1}{\sqrt{2}}\left(e_i-e_{i'}\right)\qquad(1\le i<i'\le 3).
\]
For $c\in\R$ one has $\{n\cdot x=c\}=\{x_i\pm x_{i'}=\sqrt{2}\cdot c\}$, where the sign is the one that occurs in $n$;
so the planes $\{n\cdot x=c\}$, $c\in\R$, form one of the six diagonal families,
and the plane $\{n\cdot x=0\}$ belongs to it.
Hence $\mu$ is invariant under the reflection $r_{n,0}$ in $\{n\cdot x=0\}$ and reflection positive at this plane.
These are the hypotheses of \cite[Theorem~8.8(v)]{NoncompactQED} for the law $\mu\in\mathcal{C}$ and the vector $n$,
and that theorem gives: for $\nu$-almost every $\xi$ the law $\mu_\xi$ is invariant under $r_{n,0}$
and reflection positive at the plane $\{n\cdot x=c\}$ for every $c\in\R$.
For a unit vector $n$ and $c\in\R$ the reflection in the plane $\{n\cdot x=c\}$ satisfies
\[
r_{n,c}(x)=x-2\cdot\left(n\cdot x-c\right)\cdot n=r_{n,0}(x)+2c\cdot n,
\qquad\text{that is,}\qquad r_{n,c}=\tau_{2c\cdot n}\circ r_{n,0};
\]
since $r_{e_j,0}=\theta_j$, the reflection in the plane $\{x_j=c\}$ is $\tau_{2c\cdot e_j}\circ\theta_j$.
A law that satisfies $(T_{123})$ and is invariant under $r_{n,0}$ is therefore invariant under $r_{n,c}$ for every $c\in\R$.
Each of the nine families thus contributes one assertion that holds outside a $\nu$-null set:
the three axis families by \cite[Theorem~8.1(b)]{NoncompactQED},
the six diagonal families by \cite[Theorem~8.8(v)]{NoncompactQED}.
Outside the union of these nine null sets and of the null set of \cite[Theorem~8.1(a)]{NoncompactQED},
the law $\mu_\xi$ is invariant and ergodic under the translations of $\R^3$
and is invariant and reflection positive at every plane of the nine families,
that is, $\mu_\xi\in\mathcal{C}_9$.

\emph{(b)}
The law $\mu$ satisfies \textup{(G)}.
By \cite[Theorem~8.1(c)]{NoncompactQED}, for $\nu$-almost every $\xi$ the law $\mu_\xi$ satisfies the bound of (b)
for all $f\in\mathcal{S}$ simultaneously, with the right side unchanged.

\emph{(c)}
Let $\xi$ be such that $\mu_\xi$ has the properties stated in (a).
For $n=(e_i-e_{i'})/\sqrt{2}$ one has
\[
r_{n,0}(x)=x-\left(x_i-x_{i'}\right)\cdot\left(e_i-e_{i'}\right),
\]
so $r_{n,0}$ exchanges the coordinates $x_i$ and $x_{i'}$ and fixes the third coordinate.
By (a), $\mu_\xi$ is invariant under these three transpositions of coordinates
and under the reflections $\theta_1,\theta_2,\theta_3$ in the planes $\{x_j=0\}$.
Every element of $\Ocub$ is the product of a permutation matrix and of a diagonal matrix with diagonal entries $\pm1$
\cite[Section~8.1.2]{NoncompactQED};
the permutation matrices are products of the matrices of transpositions,
and the diagonal matrices of $\Ocub$ are products of $\theta_1,\theta_2,\theta_3$.
So $\Ocub$ is the group generated by these six reflections, and $\mu_\xi$ is invariant under $\Ocub$.
Together with $(T_{123})$ this gives the invariance of $\mu_\xi$ under $\R^3\rtimes\Ocub$,
the group generated by the translations and by $\Ocub$.
As this holds for $\nu$-almost every $\xi$, the law $\mu$, which satisfies \eqref{cond:soft-hypothesis},
satisfies condition~\eqref{it:components-cubic} of Definition~\ref{def:symmetry-conditions}.

\emph{(d)}
Again let $\xi$ be such that $\mu_\xi$ has the properties stated in (a), and let $j\in\{1,2,3\}$.
As $\mu_\xi\in\mathcal{C}_9$, Proposition~\ref{prop:zero-momentum}(c), applied to the law $\mu_\xi$,
gives $P^j_{\mu_\xi}=P^{123}_{\mu_\xi}$.
The law $\mu_\xi$ is ergodic under the translations of $\R^3$, so $P^{123}_{\mu_\xi}$ is the orthogonal projection onto the constants:
\[
P^{j}_{\mu_\xi}\Phi=P^{123}_{\mu_\xi}\Phi=\mu_\xi(\Phi)\cdot 1\qquad\left(\Phi\in L^2(\mu_\xi),\ j=1,2,3\right).
\]
Thus $\mathcal{E}^j(\mu_\xi)=\C\cdot 1$:
the law $\mu_\xi$ is ergodic under the translations $\tau_{s\cdot e_j}$, $s\in\R$, in the sense of Definition~\ref{def:laws}.
By \cite[Theorem~8.3(ii)]{NoncompactQED}, for a law with $(T_j)$, $(R_j)$ and $(RP_j)$
that is ergodic under these translations,
the kernel of the Hamiltonian of the reconstruction in the direction $j$ is $\C\cdot\Omega_j$.
The law $\mu_\xi$ is such a law; hence $\ker\mathsf{H}_j(\mu_\xi)=\C\cdot\Omega_j$.
The assertion on the carrier of $\mathfrak{E}_j$ and the inequality $\mathsf{H}_j(\mu_\xi)\ge|\mathsf{P}_i|$, $i\neq j$,
are Proposition~\ref{prop:zero-momentum}(a) for the law $\mu_\xi\in\mathcal{C}_9$.

\emph{(e)}
Fix $\xi$ such that $\mu_\xi$ has the properties stated in (a), let $\Phi,\Phi'\in L^2(\mu_\xi)$,
and let $\left(a^{(k)}\right)_{k\ge1}$ be a sequence in $\R^3$ with $\left|a^{(k)}\right|\to\infty$.
It suffices to show that every subsequence has a further subsequence along which
$\mu_\xi\left(\overline{\Phi}\cdot U\left(a^{(k)}\right)\Phi'\right)$ tends to $\overline{\mu_\xi(\Phi)}\cdot\mu_\xi(\Phi')$.
Given a subsequence, the compactness of $[-\infty,\infty]$, used for the three coordinates one after the other,
yields a further subsequence and a set $J\subset\{1,2,3\}$ such that, along it,
$\left|a^{(k)}_i\right|\to\infty$ for $i\in J$, and $a^{(k)}_i$ converges to a real number $(a_\infty)_i$ for $i\notin J$.
The set $J$ is not empty, because $\left|a^{(k)}\right|\to\infty$.
Put $(a_\infty)_i:=0$ for $i\in J$ and $a_\infty:=\left((a_\infty)_1,(a_\infty)_2,(a_\infty)_3\right)$.
Since $\mu_\xi\in\mathcal{C}$, \cite[Theorem~8.3(i)]{NoncompactQED}, with $J$ in the place of the set called $K$ there,
gives along this subsequence
\[
\mu_\xi\left(\overline{\Phi}\cdot U\left(a^{(k)}\right)\Phi'\right)
=\left\langle\Phi,U\left(a^{(k)}\right)\Phi'\right\rangle_{\mu_\xi}
\to\left\langle\Phi,P^J_{\mu_\xi}U(a_\infty)\Phi'\right\rangle_{\mu_\xi},
\qquad P^J_{\mu_\xi}=\prod_{i\in J}P^i_{\mu_\xi}.
\]
By the proof of (d), every factor $P^i_{\mu_\xi}$ is the orthogonal projection onto the constants;
as $J$ is not empty, so is their product.
The limit is therefore
\[
\left\langle\Phi,\mu_\xi\left(U(a_\infty)\Phi'\right)\cdot 1\right\rangle_{\mu_\xi}
=\overline{\mu_\xi(\Phi)}\cdot\mu_\xi\left(U(a_\infty)\Phi'\right)
=\overline{\mu_\xi(\Phi)}\cdot\mu_\xi(\Phi'),
\]
where the last equality expresses the invariance of $\mu_\xi$ under the translation $\tau_{a_\infty}$.

\emph{(f)}
Every $S\in\Sigma$ with $\mu(S)=1$ satisfies $\mu_\xi(S)=1$ for $\nu$-almost every $\xi$,
by \cite[Theorem~8.1(d)]{NoncompactQED}.
The set $\{dF=0\}$ of the closed distributions belongs to $\Sigma$ \cite[Section~8.1.6]{NoncompactQED};
taking $S=\{dF=0\}$, one finds $\mu_\xi(dF=0)=1$ for $\nu$-almost every $\xi$ whenever $\mu(dF=0)=1$.
\end{proof}

\subsection{Decay of correlations}\label{sec:decay}

Let $\mu\in\mathcal{C}_9$ (Definition~\ref{def:laws}) and $j\in\{1,2,3\}$.
By Proposition~\ref{prop:zero-momentum} the law $\mu$ belongs to the class $\mathcal{C}$;
so it has the reconstruction $\left(\mathcal{H}_j(\mu),[\,\cdot\,]_j,\mathsf{H}_j,V_j,\Omega_j\right)$ in the direction $j$
and the joint spectral measure $\mathfrak{E}_j$ of Lemma~\ref{lem:joint-spectrum}.
The same proposition shows that $\mathfrak{E}_j$ is carried by the cone
\[
\mathcal{K}_j=\left\{(\lambda,k)\in[0,\infty)\times\R^2:\lambda\ge|k|_\infty\right\},
\]
that the orthogonal projection onto $\ker\mathsf{H}_j$ is $P_{\ker\mathsf{H}_j}=\mathfrak{E}_j(\{\lambda=0\})$,
that $V_j(y)P_{\ker\mathsf{H}_j}=P_{\ker\mathsf{H}_j}$ for every $y\in e_j^{\perp}$,
and that $P^j_\mu=P^{123}_\mu$.

The scalar product of $L^2(\mu)$ is $\langle\Phi,\Phi'\rangle_\mu=\mu\left(\overline{\Phi}\cdot\Phi'\right)$, antilinear in the first entry,
and its norm is written $\|\Phi\|_2:=\|\Phi\|_{L^2(\mu)}$.
The reflection $\theta_j$ acts on functions on $\mathcal{S}'$ by the linear map $\hat\theta_j$ of Definition~\ref{def:isometry-actions};
thus $\hat\theta_j\Phi=\overline{\Theta_j\Phi}$.
As a member of $\mathcal{C}_9$, the law $\mu$ is invariant under all translations of $\R^3$ and under $\theta_j$;
the second invariance is the condition $(R_j)$ \cite[Section~8.1.3]{NoncompactQED}.
Hence $\hat\theta_j$ and the operators $U(a)=\hat\tau_a$, $a\in\R^3$, are unitary on $L^2(\mu)$,
with $\hat\theta_j^2=1$ and $U(a)^{-1}=U(-a)$, and
\[
\hat\theta_j\,U(a)=U(\theta_ja)\,\hat\theta_j\qquad(a\in\R^3),
\]
because $\theta_j\circ\tau_a=\tau_{\theta_ja}\circ\theta_j$ and $\widehat{\mathfrak{g}_1\circ\mathfrak{g}_2}=\hat{\mathfrak{g}}_1\,\hat{\mathfrak{g}}_2$ for Euclidean maps $\mathfrak{g}_1$, $\mathfrak{g}_2$ \cite[Section~8.1.2]{NoncompactQED}.

Let $r>0$.
The open cube $(-r,r)^3$ satisfies $\theta_j(-r,r)^3=(-r,r)^3$,
and its translate $(-r,r)^3+r\cdot e_j$ lies in the half space $\{x_j>0\}$.
The operator $U(r\cdot e_j)=\hat\tau_{re_j}$ maps the $\Sigma_{(-r,r)^3}$-measurable functions onto the $\Sigma_{(-r,r)^3+re_j}$-measurable functions,
and $\hat\theta_j$ maps the $\Sigma_{(-r,r)^3}$-measurable functions onto themselves \cite[Section~8.1.2]{NoncompactQED}.
Every form in $C^{\infty}_{\mathrm c}\left((-r,r)^3+r\cdot e_j;\Lambda^2\right)$ belongs to $C^{\infty}_{\mathrm c}\left(\{x_j>0\};\Lambda^2\right)$,
so that $\Sigma_{(-r,r)^3+re_j}\subset\Sigma_j^{+}$.
Consequently, for $\Phi\in L^2(\Sigma_{(-r,r)^3},\mu)$ the function $U(r\cdot e_j)\Phi$ lies in $L^2(\Sigma_j^{+},\mu)$ and has the norm $\|\Phi\|_2$,
while $\hat\theta_j\Phi$ and $\Theta_j\Phi=\overline{\hat\theta_j\Phi}$ lie in $L^2(\Sigma_{(-r,r)^3},\mu)$ and have the norm $\|\Phi\|_2$.

\begin{definition}[Cube vectors]\label{def:cube-vectors}
Let $\mu\in\mathcal{C}_9$, $j\in\{1,2,3\}$, $r>0$ and $\Phi\in L^2(\Sigma_{(-r,r)^3},\mu)$.
The \emph{cube vector} of $\Phi$ in the direction $j$ is
\[
\psi^j_\Phi:=\left[U(r\cdot e_j)\Phi\right]_j\in\mathcal{H}_j(\mu).
\]
\end{definition}

The class map $[\,\cdot\,]_j$ is a linear contraction on $L^2(\Sigma_j^{+},\mu)$ \cite[Proposition~8.19\,(i)]{NoncompactQED},
and $U(r\cdot e_j)\Phi$ belongs to this space.
Hence the cube vector is well defined, it depends linearly on $\Phi$, and
\[
\left\|\psi^j_\Phi\right\|\le\left\|U(r\cdot e_j)\Phi\right\|_2=\|\Phi\|_2 .
\]

\begin{lemma}[A periodic kernel]\label{lem:periodic-kernel}
Let $(z_m)_{m\in\Z}$ be the even sequence of period four with $(z_0,z_1,z_2,z_3)=(0,1,-2,1)$,
let $B$ be the $2\pi$-periodic function on $\R$ with
\[
B(u)=8\left(1-\frac{4|u|}{\pi}\right)_{+}\qquad(|u|\le\pi),
\]
and put $b_m:=(2\pi)^{-1}\int_{-\pi}^{\pi}\mathrm{d}u\,B(u)\cdot\exp(-imu)$ for $m\in\Z$.
\begin{enumerate}
\item[(a)] One has $b_0=1$ and $b_m=\left(\frac{8\sin(m\pi/8)}{m\pi}\right)^{2}$ for $m\ne0$; $b_m=b_{-m}\ge0$;
$B(u)=\sum_{m\in\Z}b_m\cdot\exp(imu)$ for every $u\in\R$; and
\[
C_z:=\sum_{m\in\Z}|z_m|\cdot b_m=8.
\]
\item[(b)] The function $B_z(u):=\sum_{m\in\Z}z_m\cdot b_m\cdot\exp(imu)$ is real and continuous, $|B_z|\le8$,
and $B_z(u)=0$ whenever $\operatorname{dist}(u,2\pi\Z)\le\pi/4$.
\item[(c)] The function $\Psi(t):=2\sum_{m\ge1}z_m\cdot b_m\cdot t^{m-1}$, $t\in[0,1]$, is continuous, $|\Psi|\le8$, $\Psi(1)=0$, and
\[
\Psi\left(\exp(-x)\right)\ge\frac{\tanh x}{1+\sqrt{2}}\qquad(0<x<\infty).
\]
In particular $\Psi\ge0$ on $[0,1]$ and $\Psi>0$ on $[0,1)$.
\end{enumerate}
\end{lemma}

\begin{proof}
(a) The function $B$ is even, continuous and non-negative, and it vanishes for $\pi/4\le|u|\le\pi$.
With $a:=\pi/4$ one has $B(u)=8(1-|u|/a)$ for $|u|\le a$, and two integrations by parts give
\[
\int_{-a}^{a}\mathrm{d}u\,8\left(1-\frac{|u|}{a}\right)\cdot\cos(mu)=8a\cdot\left(\frac{2\sin(ma/2)}{ma}\right)^{2}\qquad(m\ne0),
\]
while the integral equals $8a$ for $m=0$.
Since $B$ is even, $b_m$ is $(2\pi)^{-1}$ times this integral;
as $8a/(2\pi)=1$ and $ma/2=m\pi/8$, the stated values of $b_m$ follow, and with them $b_m=b_{-m}\ge0$.
For $m\ne0$ one has $0\le b_m\le64/(\pi^2m^2)$, so the series $\sum_mb_m\cdot\exp(imu)$ converges absolutely and uniformly.
Its sum is a continuous $2\pi$-periodic function with the Fourier coefficients $b_m$,
hence equal to $B$ by the uniqueness theorem for Fourier coefficients of continuous periodic functions.

The identities
\[
z_m=\frac12\,i^m-(-1)^m+\frac12\,(-i)^m,\qquad|z_m|=1-\frac12\,i^m-\frac12\,(-i)^m\qquad(m\in\Z)
\]
hold, because all four members have period four in $m$ and the right sides take the values $(0,1,-2,1)$ and $(0,1,2,1)$ at $m=0,1,2,3$.
Since $i^m=\exp\left(\frac{im\pi}{2}\right)$, $(-i)^m=\exp\left(-\frac{im\pi}{2}\right)$ and $\sum_mb_m<\infty$, it follows that
\[
\begin{aligned}
C_z&=\sum_{m\in\Z}b_m-\frac12\sum_{m\in\Z}b_m\cdot\exp\left(\frac{im\pi}{2}\right)-\frac12\sum_{m\in\Z}b_m\cdot\exp\left(-\frac{im\pi}{2}\right)\\
&=B(0)-\frac12\,B\left(\frac{\pi}{2}\right)-\frac12\,B\left(-\frac{\pi}{2}\right)=8 .
\end{aligned}
\]

(b) The same identity for $z_m$ and (a) give
\begin{equation}\label{eq:continuum-kernel-shifts}
B_z(u)=\frac12\,B\left(u+\frac{\pi}{2}\right)-B(u+\pi)+\frac12\,B\left(u-\frac{\pi}{2}\right)\qquad(u\in\R);
\end{equation}
so $B_z$ is real and continuous.
As $B(v)\ne0$ only if $\operatorname{dist}(v,2\pi\Z)<\pi/4$,
the first term on the right of \eqref{eq:continuum-kernel-shifts} can be non-zero only where $\operatorname{dist}(u+\pi/2,2\pi\Z)<\pi/4$,
the second only where $\operatorname{dist}(u+\pi,2\pi\Z)<\pi/4$, and the third only where $\operatorname{dist}(u-\pi/2,2\pi\Z)<\pi/4$.
Any two of the four numbers $u$, $u+\pi/2$, $u+\pi$, $u-\pi/2$ differ by one of $\pm\pi/2$, $\pm\pi$, $\pm3\pi/2$, each of which lies at distance at least $\pi/2$ from $2\pi\Z$;
by the triangle inequality, the distances of the two numbers from $2\pi\Z$ add up to at least $\pi/2$.
So no two of the four numbers lie closer than $\pi/4$ to $2\pi\Z$,
and when $\operatorname{dist}(u,2\pi\Z)\le\pi/4$ none of $u+\pi/2$, $u+\pi$, $u-\pi/2$ does.
At every $u$ at most one of the three terms is therefore different from zero, and its modulus is at most $8$, since $0\le B\le8$;
if $\operatorname{dist}(u,2\pi\Z)\le\pi/4$, all three vanish.

(c) As $\sum_m|z_m|\cdot b_m=C_z<\infty$, the power series defining $\Psi$ converges absolutely and uniformly on $[0,1]$, and $\Psi$ is continuous.
Because $z_0=0$ and $m\mapsto z_m\cdot b_m$ is even, $|\Psi|\le2\sum_{m\ge1}|z_m|\cdot b_m=C_z=8$,
and $\Psi(1)=\sum_{m\in\Z}z_m\cdot b_m=B_z(0)=0$ by (b).

Let $x>0$, and let
\[
P_x(u):=\sum_{m\in\Z}\exp(-|m|x)\cdot\exp(imu)=\frac{\sinh x}{\cosh x-\cos u}\qquad(u\in\R)
\]
be the Poisson kernel; the closed form is the sum of two geometric series.
The series defining $B_z$ converges uniformly, so that $z_m\cdot b_m=(2\pi)^{-1}\int_{-\pi}^{\pi}\mathrm{d}u\,B_z(u)\cdot\exp(-imu)$;
the series defining $P_x$ converges uniformly as well, and $P_x$ is even.
Using once more that $z_0=0$ and that $m\mapsto z_m\cdot b_m$ is even, one obtains
\[
\begin{aligned}
\Psi\left(\exp(-x)\right)&=\exp(x)\cdot\sum_{m\in\Z}z_m\cdot b_m\cdot\exp(-|m|x)
=\frac{\exp(x)}{2\pi}\int_{-\pi}^{\pi}\mathrm{d}u\,B_z(u)\cdot P_x(u)\\
&=\frac{\exp(x)}{2\pi}\int_{-\pi}^{\pi}\mathrm{d}v\,B(v)\cdot\left(\frac12\,P_x\left(v-\frac{\pi}{2}\right)-P_x(v-\pi)+\frac12\,P_x\left(v+\frac{\pi}{2}\right)\right);
\end{aligned}
\]
the last step inserts \eqref{eq:continuum-kernel-shifts}, substitutes $v=u+\pi/2$, $v=u+\pi$, $v=u-\pi/2$ in the three terms,
and uses the $2\pi$-periodicity of $B$ and of $P_x$.
Since $\cos(v\mp\pi/2)=\pm\sin v$ and $\cos(v-\pi)=-\cos v$, the expression in the large parentheses equals
\[
\sinh x\cdot\left(\frac{\cosh x}{\cosh^2x-\sin^2v}-\frac{1}{\cosh x+\cos v}\right),
\]
and for $|v|\le\pi/4$, where $\sqrt{2}\cos v\ge1$, it is bounded from below by
\[
\sinh x\cdot\left(\frac{1}{\cosh x}-\frac{\sqrt{2}}{\sqrt{2}\cosh x+1}\right)=\frac{\tanh x}{\sqrt{2}\cosh x+1}.
\]
The function $B$ is non-negative, it vanishes for $\pi/4<|v|\le\pi$, and its mean $(2\pi)^{-1}\int_{-\pi}^{\pi}\mathrm{d}v\,B(v)$ is $b_0=1$.
Therefore
\[
\Psi\left(\exp(-x)\right)\ge\frac{\exp(x)\cdot\tanh x}{\sqrt{2}\cosh x+1}.
\]
For $x\ge0$,
\[
\frac{\exp(x)}{\sqrt{2}\cosh x+1}=\frac{2}{\sqrt{2}+\sqrt{2}\exp(-2x)+2\exp(-x)}\ge\frac{2}{2\sqrt{2}+2}=\frac{1}{1+\sqrt{2}},
\]
and the bound of (c) follows.
Finally $\Psi(0)=2z_1\cdot b_1=2b_1>0$.
The map $x\mapsto\exp(-x)$ sends $(0,\infty)$ onto $(0,1)$, and $\Psi(1)=0$; so $\Psi\ge0$ on $[0,1]$ and $\Psi>0$ on $[0,1)$.
\end{proof}

\begin{lemma}[Correlations of two cubes]\label{lem:two-cubes}
Let $\mu\in\mathcal{C}_9$, $j\in\{1,2,3\}$, $r>0$, $\Phi,\Phi'\in L^2(\Sigma_{(-r,r)^3},\mu)$ and $y\in\R^3$.
\begin{enumerate}
\item[(a)] If $y_j\ge2r$, then, with $y_\perp:=y-y_j\cdot e_j\in e_j^{\perp}$,
\[
\langle\Phi,U(y)\Phi'\rangle_\mu=\left\langle\psi^j_{\hat\theta_j\Phi},\exp\left(-(y_j-2r)\cdot\mathsf{H}_j\right)V_j(y_\perp)\,\psi^j_{\Phi'}\right\rangle .
\]
\item[(b)] One has $\langle\Phi,U(-y)\Phi'\rangle_\mu=\overline{\langle\Phi',U(y)\Phi\rangle_\mu}$.
\end{enumerate}
\end{lemma}

\begin{proof}
(a) For $\Phi_1,\Phi_2\in L^2(\Sigma_j^{+},\mu)$, $t\ge0$ and $y'\in e_j^{\perp}$, the reconstruction satisfies
\[
\begin{aligned}
\langle[\Phi_1]_j,[\Phi_2]_j\rangle&=\langle\hat\theta_j\Phi_1,\Phi_2\rangle_\mu,\\
\exp(-t\cdot\mathsf{H}_j)[\Phi_2]_j&=\left[U(t\cdot e_j)\Phi_2\right]_j,\\
V_j(y')[\Phi_2]_j&=\left[U(y')\Phi_2\right]_j
\end{aligned}
\]
\cite[Proposition~8.19\,(i)--(iii)]{NoncompactQED};
the companion writes $\Theta_j$ for the linear map $\hat\theta_j$ and $T(t)$ for $\exp(-t\cdot\mathsf{H}_j)$.
Here $U(y')$ maps $L^2(\Sigma_j^{+},\mu)$ onto itself, since $\tau_{y'}\{x_j>0\}=\{x_j>0\}$ \cite[Section~8.1.2]{NoncompactQED}.
Put $t:=y_j-2r\ge0$.
Since $t\cdot e_j+y_\perp+r\cdot e_j=y-r\cdot e_j$, the vector in the second entry of the right side of (a) is
\[
\exp(-t\cdot\mathsf{H}_j)V_j(y_\perp)\,\psi^j_{\Phi'}=\left[U(t\cdot e_j)U(y_\perp)U(r\cdot e_j)\Phi'\right]_j=\left[U(y-r\cdot e_j)\Phi'\right]_j .
\]
Moreover $\hat\theta_jU(r\cdot e_j)\hat\theta_j=U(-r\cdot e_j)$, by $\hat\theta_jU(a)=U(\theta_ja)\hat\theta_j$ and $\hat\theta_j^2=1$.
The right side of (a) is therefore
\[
\left\langle\hat\theta_jU(r\cdot e_j)\hat\theta_j\Phi,U(y-r\cdot e_j)\Phi'\right\rangle_\mu
=\left\langle U(-r\cdot e_j)\Phi,U(y-r\cdot e_j)\Phi'\right\rangle_\mu=\langle\Phi,U(y)\Phi'\rangle_\mu ,
\]
where the last step uses that $U(r\cdot e_j)$ is unitary with inverse $U(-r\cdot e_j)$.

(b) As $U(y)$ is unitary with inverse $U(-y)$,
$\langle\Phi,U(-y)\Phi'\rangle_\mu=\langle U(y)\Phi,\Phi'\rangle_\mu=\overline{\langle\Phi',U(y)\Phi\rangle_\mu}$.
\end{proof}

\begin{lemma}[Bounded momentum density]\label{lem:momentum-density}
Let $\mu\in\mathcal{C}_9$, $j\in\{1,2,3\}$ and $r>0$, and let $\Phi\in L^2(\Sigma_{(-r,r)^3},\mu)$ satisfy $\hat\theta_j\Phi=\pm\Phi$
for a fixed sign $\pm$ (equivalently, $\Theta_j\Phi=\pm\overline{\Phi}$).
Let $\psi:=\psi^j_\Phi$, and let $\varrho:=\varrho_\psi$ be its spectral measure (Lemma~\ref{lem:joint-spectrum}), with $e_j^{\perp}$ identified with $\R^2$ as there.
Let $\tilde\varrho$ be the image on $\R^2$ of the measure $\Psi\left(\exp(-2r\lambda)\right)\cdot\mathrm{d}\varrho(\lambda,k)$ on $[0,\infty)\times\R^2$
under the map $(\lambda,k)\mapsto k$, where $\Psi$ is the function of Lemma~\ref{lem:periodic-kernel}(c).
Then $\tilde\varrho$ is a finite positive measure, and $\tilde\varrho=h\,\mathrm{d}k$ for a continuous function $h$ on $\R^2$,
where $\mathrm{d}k$ denotes the Lebesgue measure; this function satisfies
\[
\begin{aligned}
&0\le h\le C_2\cdot r^2\cdot\|\Phi\|_2^2,\\
&C_2:=\frac{C'}{4\pi^2},\qquad C':=16\cdot C_z+64\left(\frac{16}{\pi}+\frac{64}{\pi^2}\right)=128+64\left(\frac{16}{\pi}+\frac{64}{\pi^2}\right),
\end{aligned}
\]
with the number $C_z=8$ of Lemma~\ref{lem:periodic-kernel}(a).
\end{lemma}

\begin{proof}
For $v\in\R^3$ put
\[
u(v):=\pm\langle\Phi,U(v)\Phi\rangle_\mu=\langle\hat\theta_j\Phi,U(v)\Phi\rangle_\mu ,
\]
with the sign of the hypothesis.
The Cauchy--Schwarz inequality and the unitarity of $U(v)$ give $|u(v)|\le\|\Phi\|_2^2$.
For $y\in e_j^{\perp}$ let
\[
K(y):=\sum_{m\in\Z}z_m\cdot b_m\cdot u(2rm\cdot e_j+y),
\]
with the numbers $z_m$, $b_m$ of Lemma~\ref{lem:periodic-kernel}.
By Lemma~\ref{lem:periodic-kernel}(a) the series converges absolutely, and $|K(y)|\le C_z\cdot\|\Phi\|_2^2$.
Since $z_0=0$, only the terms with $m\ne0$ contribute, that is, only the points $a\cdot e_j+y$ with $|a|\ge2r$.
For real $a$ and $y\in e_j^{\perp}$ the following hold.
\begin{enumerate}
\item\label{it:density-time}
If $a\ge2r$, then
\[
u(a\cdot e_j+y)=\left\langle\psi,\exp\left(-(a-2r)\cdot\mathsf{H}_j\right)V_j(y)\,\psi\right\rangle .
\]
Indeed, the point $a\cdot e_j+y$ has the $j$-th coordinate $a\ge2r$, and its component orthogonal to $e_j$ is $y$.
Lemma~\ref{lem:two-cubes}(a) with $\Phi'=\Phi$ expresses $\langle\Phi,U(a\cdot e_j+y)\Phi\rangle_\mu$
as the matrix element of $\exp\left(-(a-2r)\cdot\mathsf{H}_j\right)V_j(y)$ between $\psi^j_{\hat\theta_j\Phi}$ and $\psi$,
and $\psi^j_{\hat\theta_j\Phi}=\pm\psi$ by linearity; the sign is real and its square is $1$.
\item\label{it:density-reflection}
One has $u(-a\cdot e_j+y)=u(a\cdot e_j+y)$.
For the proof, recall that $\hat\theta_j$ is unitary, that $\hat\theta_jU(v)=U(\theta_jv)\hat\theta_j$, and that $\theta_j(-a\cdot e_j+y)=a\cdot e_j+y$.
Hence
\[
\begin{aligned}
\langle\Phi,U(-a\cdot e_j+y)\Phi\rangle_\mu
&=\langle\hat\theta_j\Phi,\hat\theta_jU(-a\cdot e_j+y)\Phi\rangle_\mu\\
&=\langle\hat\theta_j\Phi,U(a\cdot e_j+y)\hat\theta_j\Phi\rangle_\mu=\langle\Phi,U(a\cdot e_j+y)\Phi\rangle_\mu ,
\end{aligned}
\]
the last step because $\hat\theta_j\Phi=\pm\Phi$ with a real sign whose square is $1$.
\item\label{it:density-space}
Let $i\ne j$, let $i'$ be the third index, and let $y=y_i\cdot e_i+y_{i'}\cdot e_{i'}$ with $y_i\ge2r$.
Then, with the sign of the hypothesis and with the reconstruction and the cube vectors of $\mu$ in the direction $i$,
\[
u(a\cdot e_j+y)=\pm\left\langle\psi^i_{\hat\theta_i\Phi},\exp\left(-(y_i-2r)\cdot\mathsf{H}_i\right)V_i(a\cdot e_j)V_i(y_{i'}\cdot e_{i'})\,\psi^i_\Phi\right\rangle .
\]
This is Lemma~\ref{lem:two-cubes}(a) in the direction $i$, with $\Phi'=\Phi$, at the point $a\cdot e_j+y$:
its $i$-th coordinate is $y_i\ge2r$, its component orthogonal to $e_i$ is $a\cdot e_j+y_{i'}\cdot e_{i'}$,
and $V_i(a\cdot e_j+y_{i'}\cdot e_{i'})=V_i(a\cdot e_j)V_i(y_{i'}\cdot e_{i'})$ (Lemma~\ref{lem:joint-spectrum}(a)).
\item\label{it:density-conjugation}
One has $u(a\cdot e_j-y)=\overline{u(a\cdot e_j+y)}$.
Indeed, Lemma~\ref{lem:two-cubes}(b) with $\Phi'=\Phi$ gives $u(-a\cdot e_j-y)=\overline{u(a\cdot e_j+y)}$,
and $u(-a\cdot e_j-y)=u(a\cdot e_j-y)$ by \eqref{it:density-reflection}.
\item\label{it:density-fourier}
The kernel $K$ is the Fourier transform of $\tilde\varrho$:
\[
\begin{aligned}
K(y)&=2\sum_{m\ge1}z_m\cdot b_m\cdot\left\langle\psi,\exp\left(-2r(m-1)\cdot\mathsf{H}_j\right)V_j(y)\,\psi\right\rangle\\
&=\left\langle\psi,\Psi\left(\exp(-2r\cdot\mathsf{H}_j)\right)V_j(y)\,\psi\right\rangle\\
&=\int\mathrm{d}\varrho(\lambda,k)\,\Psi\left(\exp(-2r\lambda)\right)\cdot\exp(ik\cdot y)\\
&=\int_{\R^2}\mathrm{d}\tilde\varrho(k)\,\exp(ik\cdot y).
\end{aligned}
\]
Moreover $\tilde\varrho$ is a positive measure of mass at most $8\|\Phi\|_2^2$, and $K$ is continuous.
The first equality follows from $z_0=0$, from the evenness of $m\mapsto z_m\cdot b_m$, from \eqref{it:density-reflection},
and from \eqref{it:density-time}, which applies to $a=2rm\ge2r$ for $m\ge1$.
The functions $\lambda\mapsto2\sum_{m=1}^{n}z_m\cdot b_m\cdot\exp\left(-2r(m-1)\lambda\right)$
converge to $\lambda\mapsto\Psi\left(\exp(-2r\lambda)\right)$ uniformly on $[0,\infty)$ as $n\to\infty$, because $\sum_m|z_m|\cdot b_m<\infty$.
As $\mathfrak{E}_j$ is carried by $[0,\infty)\times\R^2$, Lemma~\ref{lem:joint-spectrum}(b) shows that the operator series
$2\sum_{m\ge1}z_m\cdot b_m\cdot\exp\left(-2r(m-1)\cdot\mathsf{H}_j\right)$ converges in norm to $\Psi\left(\exp(-2r\cdot\mathsf{H}_j)\right)$;
this is the second equality.
The third equality is Lemma~\ref{lem:joint-spectrum}(b) for the bounded continuous function $\lambda\mapsto\Psi\left(\exp(-2r\lambda)\right)$ on $[0,\infty)$,
together with $\varrho=\langle\psi,\mathfrak{E}_j(\cdot)\psi\rangle$; the fourth is the definition of $\tilde\varrho$.
By Lemma~\ref{lem:periodic-kernel}(c), $0\le\Psi\le8$, and $\varrho$ has the mass $\|\psi\|^2\le\|\Phi\|_2^2$;
this gives the positivity of $\tilde\varrho$ and the bound on its mass.
By dominated convergence, the Fourier transform $K$ of the finite measure $\tilde\varrho$ is continuous.
\item\label{it:density-decay}
If $|y|_\infty\ge2r$, then
\[
|K(y)|\le8\exp\left(-\frac{\pi\left(|y|_\infty-2r\right)}{8r}\right)\cdot\|\Phi\|_2^2 .
\]
To see this, choose $i\ne j$ with $|y_i|=|y|_\infty$, and let $i'$ be the third index.
By \eqref{it:density-conjugation} and since the numbers $z_m\cdot b_m$ are real, $K(-y)=\overline{K(y)}$;
replacing $y$ by $-y$ if necessary, one may assume $y_i\ge2r$.
Multiplying \eqref{it:density-space} with $a=2rm$ by $z_m\cdot b_m$ and summing over $m\in\Z$, one finds
\[
\begin{aligned}
K(y)&=\pm\left\langle\psi^i_{\hat\theta_i\Phi},\exp\left(-(y_i-2r)\cdot\mathsf{H}_i\right)B_{z,i}\,V_i(y_{i'}\cdot e_{i'})\,\psi^i_\Phi\right\rangle,\\
B_{z,i}&:=\sum_{m\in\Z}z_m\cdot b_m\cdot V_i(2rm\cdot e_j),
\end{aligned}
\]
where the operator series converges in norm.
Let $(\lambda,k)$ be the spectral variables of the reconstruction in the direction $i$;
the momentum $k$ has the two components $k_j$ and $k_{i'}$, and $k_j$ is the variable that belongs to $V_i(\cdot\,e_j)$.
Lemma~\ref{lem:joint-spectrum}(b) in the direction $i$ gives, for $t\ge0$,
\[
\exp(-t\cdot\mathsf{H}_i)B_{z,i}=\int\mathrm{d}\mathfrak{E}_i(\lambda,k)\,\exp(-t\cdot\lambda)\cdot B_z(2rk_j);
\]
for $t=0$ this is $B_{z,i}=\int\mathrm{d}\mathfrak{E}_i(\lambda,k)\,B_z(2rk_j)$.
If $B_z(2rk_j)\ne0$, then $|2rk_j|>\pi/4$ by Lemma~\ref{lem:periodic-kernel}(b), that is, $|k_j|>\pi/(8r)$;
on the cone $\mathcal{K}_i$, which carries $\mathfrak{E}_i$ (Proposition~\ref{prop:zero-momentum}(a)), this forces $\lambda\ge|k_j|>\pi/(8r)$.
Since $|B_z|\le8$, it follows that for every $t\ge0$
\[
\begin{aligned}
\left\|\exp(-t\cdot\mathsf{H}_i)B_{z,i}\right\|
&\le\sup_{(\lambda,k)\in\mathcal{K}_i}\left|\exp(-t\cdot\lambda)\cdot B_z(2rk_j)\right|\\
&\le8\exp\left(-\frac{\pi t}{8r}\right).
\end{aligned}
\]
The claim follows with $t=y_i-2r=|y|_\infty-2r$, since $\left\|\psi^i_{\hat\theta_i\Phi}\right\|\le\|\hat\theta_i\Phi\|_2=\|\Phi\|_2$
and $\left\|V_i(y_{i'}\cdot e_{i'})\,\psi^i_\Phi\right\|\le\|\Phi\|_2$.
\item\label{it:density-integrable}
The kernel is integrable:
\[
\int_{\R^2}\mathrm{d}y\,|K(y)|\le C'\cdot r^2\cdot\|\Phi\|_2^2 .
\]
The set $\{y\in\R^2:|y|_\infty\le s\}$ has the Lebesgue measure $4s^2$ for $s\ge0$;
so the image of the Lebesgue measure under $y\mapsto|y|_\infty$ has the density $8s$ on $(0,\infty)$.
On the square $\{|y|_\infty\le2r\}$, of measure $(4r)^2$, one has $|K|\le C_z\cdot\|\Phi\|_2^2$, and outside it \eqref{it:density-decay} applies.
Since
\[
\int_{2r}^{\infty}\mathrm{d}s\,s\cdot\exp\left(-\frac{\pi(s-2r)}{8r}\right)=2r\cdot\frac{8r}{\pi}+\frac{64r^2}{\pi^2},
\]
it follows that
\[
\begin{aligned}
\int_{\R^2}\mathrm{d}y\,|K(y)|
&\le\left(C_z\cdot(4r)^2+8\int_{2r}^{\infty}\mathrm{d}s\,8s\cdot\exp\left(-\frac{\pi(s-2r)}{8r}\right)\right)\cdot\|\Phi\|_2^2\\
&=\left(16\cdot C_z\cdot r^2+64\left(\frac{16r^2}{\pi}+\frac{64r^2}{\pi^2}\right)\right)\cdot\|\Phi\|_2^2=C'\cdot r^2\cdot\|\Phi\|_2^2 .
\end{aligned}
\]
\item\label{it:density-inversion}
The function
\[
h(k):=(2\pi)^{-2}\int_{\R^2}\mathrm{d}y\,K(y)\cdot\exp(-ik\cdot y)\qquad(k\in\R^2)
\]
is continuous, $|h|\le(2\pi)^{-2}\cdot C'\cdot r^2\cdot\|\Phi\|_2^2$, and $\tilde\varrho=h\,\mathrm{d}k$.
Continuity and the bound follow from \eqref{it:density-integrable}, by dominated convergence.
Let $\chi\in\mathcal{S}(\R^2)$ and $\hat\chi(y):=\int\mathrm{d}k\,\chi(k)\cdot\exp(-ik\cdot y)$;
by the Fourier inversion theorem, $\chi(k)=(2\pi)^{-2}\int\mathrm{d}y\,\hat\chi(y)\cdot\exp(ik\cdot y)$.
Fubini's theorem ($\hat\chi$ is integrable and $\tilde\varrho$ is finite), \eqref{it:density-fourier},
and Fubini's theorem once more ($\chi$ and $K$ are integrable) give
\[
\int\mathrm{d}\tilde\varrho(k)\,\chi(k)=(2\pi)^{-2}\int\mathrm{d}y\,\hat\chi(y)\cdot K(y)=\int\mathrm{d}k\,\chi(k)\cdot h(k).
\]
In particular $\tilde\varrho$ and $h\,\mathrm{d}k$ give the same integral to every $\chi\in C^{\infty}_{\mathrm c}(\R^2)\subset\mathcal{S}(\R^2)$.
Both are complex Borel measures on $\R^2$ that are finite on compact sets, the second because $h$ is bounded,
and such a measure is determined by its integrals against the functions of $C^{\infty}_{\mathrm c}(\R^2)$;
hence $\tilde\varrho=h\,\mathrm{d}k$.
\end{enumerate}
By \eqref{it:density-fourier} and \eqref{it:density-inversion}, $\tilde\varrho$ is a finite positive measure and $\tilde\varrho=h\,\mathrm{d}k$ with $h$ continuous.
The positivity of $\tilde\varrho$ gives $h\ge0$ almost everywhere, hence everywhere by continuity,
and \eqref{it:density-inversion} gives $h\le(2\pi)^{-2}\cdot C'\cdot r^2\cdot\|\Phi\|_2^2=C_2\cdot r^2\cdot\|\Phi\|_2^2$.
\end{proof}

The momentum density of Lemma~\ref{lem:momentum-density}, together with the cone $\mathcal{K}_j$ of Proposition~\ref{prop:zero-momentum},
bounds the spectral weight that a cube vector carries away from $\ker\mathsf{H}_j$.

\begin{lemma}[Decay on cube vectors]\label{lem:cube-decay}
Let $\mu\in\mathcal{C}_9$, $j\in\{1,2,3\}$, $r>0$ and $t>0$,
and let $C_2$ be the constant of Lemma~\ref{lem:momentum-density}.
Put
\[
C_3:=32\cdot(1+\sqrt{2})\cdot C_2 .
\]
Then for every $\Phi\in L^2(\Sigma_{(-r,r)^3},\mu)$ the cube vector $\psi^{j}_{\Phi}\in\mathcal{H}_j(\mu)$ of Definition~\ref{def:cube-vectors} satisfies
\[
0\le\langle\psi^{j}_{\Phi},\exp(-t\cdot\mathsf{H}_j)\,(1-P_{\ker\mathsf{H}_j})\,\psi^{j}_{\Phi}\rangle
\le C_3\cdot\|\Phi\|_2^2\cdot\left(\left(\frac{r}{t}\right)^{2}+\frac{r}{t}\right).
\]
\end{lemma}

\begin{proof}
Write $S:=\exp(-t\cdot\mathsf{H}_j)\,(1-P_{\ker\mathsf{H}_j})$.
By Proposition~\ref{prop:zero-momentum}(b) the projection $P_{\ker\mathsf{H}_j}$ equals $\mathfrak{E}_j(\{0\}\times\R^2)$,
and Lemma~\ref{lem:joint-spectrum}(b), taken at $y=0$, gives
\[
S=\int_{(0,\infty)\times\R^2}\mathrm{d}\mathfrak{E}_j(\lambda,k)\,\exp(-t\cdot\lambda),
\qquad
\langle\psi,S\psi\rangle=\int_{(0,\infty)\times\R^2}\mathrm{d}\varrho_\psi(\lambda,k)\,\exp(-t\cdot\lambda)\ge0
\]
for every $\psi\in\mathcal{H}_j$, where $\varrho_\psi=\langle\psi,\mathfrak{E}_j(\cdot)\psi\rangle$.
The left inequality follows.

\emph{Observables of definite parity.}
Suppose first that $\hat\theta_j\Phi=\pm\Phi$ for a fixed sign.
Let $\psi:=\psi^{j}_{\Phi}$,
and let $\varrho:=\varrho_\psi$, the measure $\tilde\varrho$ on $e_j^{\perp}=\R^2$ and its continuous density $h$ be those of Lemma~\ref{lem:momentum-density}.
Thus $\tilde\varrho$ is the image of the measure $\Psi(\exp(-2r\cdot\lambda))\,\mathrm{d}\varrho(\lambda,k)$ under $(\lambda,k)\mapsto k$,
and $0\le h\le C_2\cdot r^2\cdot\|\Phi\|_2^2$.
By Lemma~\ref{lem:periodic-kernel}(c) the function $\Psi$ is non-negative on $[0,1]$, and, with $x=2r\cdot\lambda$,
\[
\Psi(\exp(-2r\cdot\lambda))\ge\frac{\tanh(2r\cdot\lambda)}{1+\sqrt{2}}>0\qquad(\lambda>0).
\]
Since $\tilde\varrho$ has a density, it gives no mass to the point $0$:
\[
0=\tilde\varrho(\{0\})=\int_{[0,\infty)\times\{0\}}\mathrm{d}\varrho(\lambda,k)\,\Psi(\exp(-2r\cdot\lambda))
\ge\int_{(0,\infty)\times\{0\}}\mathrm{d}\varrho(\lambda,k)\,\Psi(\exp(-2r\cdot\lambda)) .
\]
The integrand of the last integral is strictly positive; hence $\varrho((0,\infty)\times\{0\})=0$.
By Proposition~\ref{prop:zero-momentum}(a) the measure $\varrho$ is carried by $\mathcal{K}_j$, on which $\lambda\ge|k|_\infty$;
so $\lambda\ge|k|_\infty>0$ at $\varrho$-almost every point with $k\neq0$.
The function $x\mapsto\exp(-t\cdot x)/\tanh(2r\cdot x)$ is decreasing on $(0,\infty)$.
Moreover $1/\tanh x\le1+1/x$ for $x>0$, since
\[
\coth x-1=\frac{2}{\exp(2x)-1}\le\frac{1}{x},
\]
and the image of the Lebesgue measure $\mathrm{d}k$ under $k\mapsto|k|_\infty$ has the density $8x$ on $(0,\infty)$,
the square $\{|k|_\infty\le x\}$ having the area $4x^2$.
Therefore
\[
\begin{aligned}
\langle\psi,S\psi\rangle
&=\int_{(0,\infty)\times(\R^2\setminus\{0\})}\mathrm{d}\varrho(\lambda,k)\,\exp(-t\cdot\lambda)\\
&\le(1+\sqrt{2})\cdot\int_{(0,\infty)\times(\R^2\setminus\{0\})}\mathrm{d}\varrho(\lambda,k)\,\frac{\exp(-t\cdot\lambda)}{\tanh(2r\cdot\lambda)}\cdot\Psi(\exp(-2r\cdot\lambda))\\
&\le(1+\sqrt{2})\cdot\int_{[0,\infty)\times(\R^2\setminus\{0\})}\mathrm{d}\varrho(\lambda,k)\,\frac{\exp(-t\cdot|k|_\infty)}{\tanh(2r\cdot|k|_\infty)}\cdot\Psi(\exp(-2r\cdot\lambda))\\
&=(1+\sqrt{2})\cdot\int_{\R^2\setminus\{0\}}\mathrm{d}k\,\frac{\exp(-t\cdot|k|_\infty)}{\tanh(2r\cdot|k|_\infty)}\cdot h(k)\\
&\le(1+\sqrt{2})\cdot C_2\cdot r^2\cdot\|\Phi\|_2^2\cdot\int_0^\infty\mathrm{d}x\,8x\cdot\exp(-t\cdot x)\cdot\left(1+\frac{1}{2r\cdot x}\right)\\
&=(1+\sqrt{2})\cdot C_2\cdot\|\Phi\|_2^2\cdot\left(8\cdot\left(\frac{r}{t}\right)^{2}+4\cdot\frac{r}{t}\right).
\end{aligned}
\]
The first step uses $\varrho((0,\infty)\times\{0\})=0$, and the second the lower bound for $\Psi$.
The third uses $\lambda\ge|k|_\infty>0$ and the monotonicity of $x\mapsto\exp(-t\cdot x)/\tanh(2r\cdot x)$;
the integrand being non-negative, the domain of integration may then be enlarged.
The fourth step is the transformation formula for an image measure:
the factor $\exp(-t\cdot|k|_\infty)/\tanh(2r\cdot|k|_\infty)$ is non-negative and depends on $k$ only,
and $\tilde\varrho=h\,\mathrm{d}k$ is the image of $\Psi(\exp(-2r\cdot\lambda))\,\mathrm{d}\varrho(\lambda,k)$ under $(\lambda,k)\mapsto k$.
The fifth uses the bound for $h$, the inequality $1/\tanh x\le1+1/x$ and the distribution of $|k|_\infty$.
As $8\cdot(r/t)^2+4\cdot r/t\le8\cdot((r/t)^2+r/t)$, every $\Phi\in L^2(\Sigma_{(-r,r)^3},\mu)$ with $\hat\theta_j\Phi=\pm\Phi$ satisfies
\begin{equation}\label{eq:continuum-decay-parity}
\langle\psi^{j}_{\Phi},S\,\psi^{j}_{\Phi}\rangle\le8\cdot(1+\sqrt{2})\cdot C_2\cdot\|\Phi\|_2^2\cdot\left(\left(\frac{r}{t}\right)^{2}+\frac{r}{t}\right).
\end{equation}

\emph{The general case.}
Since $\theta_j$ maps $(-r,r)^3$ onto itself and leaves $\mu$ invariant,
the operator $\hat\theta_j$ maps $L^2(\Sigma_{(-r,r)^3},\mu)$ isometrically onto itself, and $\hat\theta_j^2=1$.
For arbitrary $\Phi\in L^2(\Sigma_{(-r,r)^3},\mu)$ put
\[
\Phi_\pm:=\frac{1}{2}\cdot\left(\Phi\pm\hat\theta_j\Phi\right).
\]
These functions lie in $L^2(\Sigma_{(-r,r)^3},\mu)$, satisfy $\hat\theta_j\Phi_\pm=\pm\Phi_\pm$ and $\|\Phi_\pm\|_2\le\|\Phi\|_2$,
and $\psi^{j}_{\Phi}=\psi^{j}_{\Phi_+}+\psi^{j}_{\Phi_-}$, the cube vector being linear in the observable.
As $\psi\mapsto\langle\psi,S\psi\rangle^{\frac12}=\|S^{\frac12}\psi\|$ is a seminorm on $\mathcal{H}_j$,
\[
\langle\psi+\psi',S(\psi+\psi')\rangle\le2\cdot\langle\psi,S\psi\rangle+2\cdot\langle\psi',S\psi'\rangle\qquad(\psi,\psi'\in\mathcal{H}_j).
\]
With \eqref{eq:continuum-decay-parity} for $\Phi_+$ and for $\Phi_-$ this gives
\[
\begin{aligned}
\langle\psi^{j}_{\Phi},S\,\psi^{j}_{\Phi}\rangle
&\le16\cdot(1+\sqrt{2})\cdot C_2\cdot\left(\|\Phi_+\|_2^2+\|\Phi_-\|_2^2\right)\cdot\left(\left(\frac{r}{t}\right)^{2}+\frac{r}{t}\right)\\
&\le C_3\cdot\|\Phi\|_2^2\cdot\left(\left(\frac{r}{t}\right)^{2}+\frac{r}{t}\right),
\end{aligned}
\]
because $\|\Phi_+\|_2^2+\|\Phi_-\|_2^2\le2\cdot\|\Phi\|_2^2$ and $32\cdot(1+\sqrt{2})\cdot C_2=C_3$.
\end{proof}

For $\mu\in\mathcal{C}_9$, $r>0$ and $\Phi,\Phi'\in L^2(\Sigma_{(-r,r)^3},\mu)$,
Lemma~\ref{lem:cube-decay} leads to a bound of order $r/|a|_\infty$, as $|a|_\infty\to\infty$, for the difference between $\langle\Phi,U(a)\Phi'\rangle_\mu$ and $\langle\Phi,P^{123}_{\mu}\Phi'\rangle_\mu$.

\begin{theorem}[Clustering with a rate]\label{thm:clustering-rate}
Let $C_2$ and $C_3=32\cdot(1+\sqrt{2})\cdot C_2$ be the constants of Lemma~\ref{lem:momentum-density} and of Lemma~\ref{lem:cube-decay}, so that
\[
C_2=\frac{1}{4\pi^2}\cdot\left(128+64\cdot\left(\frac{16}{\pi}+\frac{64}{\pi^2}\right)\right),
\]
and put $C_1:=4$ and $C_0:=3\cdot C_3=96\cdot(1+\sqrt{2})\cdot C_2$.
Let $\mu\in\mathcal{C}_9$ and $r>0$, and let $\Phi,\Phi'\in L^2(\Sigma_{(-r,r)^3},\mu)$, bounded or not.
For every $a\in\R^3$ with $|a|_\infty>2r$, and with $t:=|a|_\infty-2r$,
\begin{equation}\label{eq:continuum-rate-bound}
|\langle\Phi,U(a)\Phi'\rangle_\mu-\langle\Phi,P^{123}_{\mu}\Phi'\rangle_\mu|
\le C_3\cdot\|\Phi\|_2\cdot\|\Phi'\|_2\cdot\left(\frac{r}{t}+\left(\frac{r}{t}\right)^{2}\right);
\end{equation}
in particular, for every $a\in\R^3$ with $|a|_\infty\ge C_1\cdot r$,
\begin{equation}\label{eq:continuum-rate-cube}
|\langle\Phi,U(a)\Phi'\rangle_\mu-\langle\Phi,P^{123}_{\mu}\Phi'\rangle_\mu|
\le C_0\cdot\|\Phi\|_2\cdot\|\Phi'\|_2\cdot\frac{r}{|a|_\infty}.
\end{equation}
If $\mu$ is ergodic under the translations of $\R^3$, then $\langle\Phi,P^{123}_{\mu}\Phi'\rangle_\mu=\overline{\mu(\Phi)}\cdot\mu(\Phi')$.
\end{theorem}

\begin{proof}
Choose $j\in\{1,2,3\}$ with $|a_j|=|a|_\infty$, and write $a_\perp:=a-a_j\cdot e_j\in e_j^{\perp}$.

\emph{The case $a_j>0$.}
Here $a_j=2r+t$.
Since $\theta_j$ maps $(-r,r)^3$ onto itself and leaves $\mu$ invariant,
the function $\hat\theta_j\Phi$ lies in $L^2(\Sigma_{(-r,r)^3},\mu)$ and has the norm $\|\Phi\|_2$.
Put $\psi_1:=\psi^{j}_{\hat\theta_j\Phi}$ and $\psi_2:=\psi^{j}_{\Phi'}$.
Lemma~\ref{lem:two-cubes}(a), with $y=a$, gives
\begin{equation}\label{eq:continuum-rate-matrix}
\langle\Phi,U(a)\Phi'\rangle_\mu=\langle\psi_1,\exp(-t\cdot\mathsf{H}_j)\,V_j(a_\perp)\,\psi_2\rangle .
\end{equation}
Split
\[
\exp(-t\cdot\mathsf{H}_j)=P_{\ker\mathsf{H}_j}+S,\qquad S:=\exp(-t\cdot\mathsf{H}_j)\,(1-P_{\ker\mathsf{H}_j})\ge0 .
\]

The first term gives the subtracted scalar product.
Indeed, $P_{\ker\mathsf{H}_j}V_j(a_\perp)=P_{\ker\mathsf{H}_j}$ by Proposition~\ref{prop:zero-momentum}(b),
and $P_{\ker\mathsf{H}_j}[G]_j=[P^{j}_{\mu}G]_j$ for $G\in L^2(\Sigma_j^{+},\mu)$ by \cite[Proposition~8.19(iv)]{NoncompactQED},
where $P_{\ker\mathsf{H}_j}$ is denoted $P_0$.
The projection $P^{j}_{\mu}$ onto the vectors invariant under the translations along $e_j$ satisfies
$P^{j}_{\mu}U(r\cdot e_j)=P^{j}_{\mu}=U(r\cdot e_j)P^{j}_{\mu}$.
The scalar product of $\mathcal{H}_j$ is $\langle[G']_j,[G]_j\rangle=\langle\hat\theta_jG',G\rangle_\mu$ for $G,G'\in L^2(\Sigma_j^{+},\mu)$
\cite[Proposition~8.19(i)]{NoncompactQED}, where $\hat\theta_j$ is written $\Theta_j$.
As $\hat\theta_jU(r\cdot e_j)\hat\theta_j=U(-r\cdot e_j)$, it follows that $\langle\psi_1,[G]_j\rangle=\langle U(-r\cdot e_j)\Phi,G\rangle_\mu$.
With Proposition~\ref{prop:zero-momentum}(c) in the last step, one obtains
\[
\begin{aligned}
\langle\psi_1,P_{\ker\mathsf{H}_j}V_j(a_\perp)\psi_2\rangle
&=\langle\psi_1,P_{\ker\mathsf{H}_j}\psi_2\rangle
=\langle\psi_1,[P^{j}_{\mu}U(r\cdot e_j)\Phi']_j\rangle
=\langle\psi_1,[P^{j}_{\mu}\Phi']_j\rangle\\
&=\langle U(-r\cdot e_j)\Phi,P^{j}_{\mu}\Phi'\rangle_\mu
=\langle\Phi,U(r\cdot e_j)P^{j}_{\mu}\Phi'\rangle_\mu\\
&=\langle\Phi,P^{j}_{\mu}\Phi'\rangle_\mu
=\langle\Phi,P^{123}_{\mu}\Phi'\rangle_\mu .
\end{aligned}
\]

For the second term, observe that $S$, $S^{\frac12}=\exp\left(-\frac{t}{2}\cdot\mathsf{H}_j\right)(1-P_{\ker\mathsf{H}_j})$ and the unitary $V_j(a_\perp)$
are integrals against the same projection-valued measure $\mathfrak{E}_j$
(Lemma~\ref{lem:joint-spectrum}(b) and Proposition~\ref{prop:zero-momentum}(b)), so that they commute.
The Cauchy--Schwarz inequality and Lemma~\ref{lem:cube-decay}, applied to $\hat\theta_j\Phi$ and to $\Phi'$, give
\[
\begin{aligned}
|\langle\psi_1,S\,V_j(a_\perp)\psi_2\rangle|
&=|\langle S^{\frac12}\psi_1,V_j(a_\perp)S^{\frac12}\psi_2\rangle|
\le\langle\psi_1,S\psi_1\rangle^{\frac12}\cdot\langle\psi_2,S\psi_2\rangle^{\frac12}\\
&\le C_3\cdot\|\Phi\|_2\cdot\|\Phi'\|_2\cdot\left(\left(\frac{r}{t}\right)^{2}+\frac{r}{t}\right).
\end{aligned}
\]
Inserted into \eqref{eq:continuum-rate-matrix}, the two terms give \eqref{eq:continuum-rate-bound} for $a_j>0$.

\emph{The case $a_j<0$.}
The vector $-a$ has $(-a)_j=|-a|_\infty>2r$, with the same $t$.
By the first case, applied to $-a$ and to the pair $(\Phi',\Phi)$,
the modulus of $\langle\Phi',U(-a)\Phi\rangle_\mu-\langle\Phi',P^{123}_{\mu}\Phi\rangle_\mu$ is at most the right side of \eqref{eq:continuum-rate-bound}.
By Lemma~\ref{lem:two-cubes}(b), $\langle\Phi,U(a)\Phi'\rangle_\mu=\overline{\langle\Phi',U(-a)\Phi\rangle_\mu}$,
and $\langle\Phi,P^{123}_{\mu}\Phi'\rangle_\mu=\overline{\langle\Phi',P^{123}_{\mu}\Phi\rangle_\mu}$ because $P^{123}_{\mu}$ is self-adjoint.
The two differences are complex conjugate to each other, and \eqref{eq:continuum-rate-bound} holds for $a$.

\emph{The bound in terms of $|a|_\infty$.}
Let $|a|_\infty\ge C_1\cdot r=4r$.
Then $t=|a|_\infty-2r\ge|a|_\infty/2$, hence $r/t\le2r/|a|_\infty\le1/2$, and
\[
\left(\frac{r}{t}\right)^{2}+\frac{r}{t}\le\frac{3}{2}\cdot\frac{r}{t}\le\frac{3r}{|a|_\infty} .
\]
With $C_0=3\cdot C_3$, the bound \eqref{eq:continuum-rate-bound} turns into \eqref{eq:continuum-rate-cube}.

\emph{Ergodic laws.}
If $\mu$ is ergodic under the translations of $\R^3$, then $\mathcal{E}^{123}(\mu)=\C\cdot1$ (Definition~\ref{def:laws}),
so that $P^{123}_{\mu}\Phi'=\mu(\Phi')\cdot1$ and
$\langle\Phi,P^{123}_{\mu}\Phi'\rangle_\mu=\mu(\Phi')\cdot\langle\Phi,1\rangle_\mu=\overline{\mu(\Phi)}\cdot\mu(\Phi')$.
\end{proof}

\begin{remark}\label{rem:rate-exponent}
\begin{enumerate}
\item[(1)]
The exponent one in \eqref{eq:continuum-rate-cube} reflects the dimension two of the momentum space $e_j^{\perp}$.
In dimension $d\ge3$ the same proof, with the momentum space $e_j^{\perp}\cong\R^{d-1}$, gives the bound
\[
C\cdot\|\Phi\|_2\cdot\|\Phi'\|_2\cdot\left(\left(\frac{r}{t}\right)^{d-2}+\left(\frac{r}{t}\right)^{d-1}\right),\qquad t=|a|_\infty-2r,
\]
with a constant $C$ that depends only on $d$:
the density of Lemma~\ref{lem:momentum-density} is then at most a constant multiple of $r^{d-1}\cdot\|\Phi\|_2^2$,
and the last integral in the proof of Lemma~\ref{lem:cube-decay} is at most a constant that depends only on $d$ times
\[
\int_0^\infty\mathrm{d}x\,x^{d-2}\cdot\exp(-t\cdot x)\cdot\left(1+\frac{1}{r\cdot x}\right).
\]
For $d=4$ the bound decays like the inverse square of the distance;
this is the rate of the corresponding clustering estimate for Yang--Mills theory in four dimensions.
The correlations of the massless free scalar field in dimension $d$ decay with the same exponent $d-2$.
Linear observables of the free Maxwell law $\gamma$ of Definition~\ref{def:free-maxwell},
which lies in $\mathcal{C}_9$ by Lemma~\ref{lem:scaled-maxwell} applied to $\gamma_1=\gamma$, decorrelate faster.
For real test forms $f$ and $f'$ with compact supports,
the covariance of $\gamma$, together with $U(a)F(f')=F(\tau_a\cdot f')$, gives the equality,
and \cite[Lemma~8.30(ii)]{NoncompactQED}, where $P^{\perp}$ appears as the projection $\Pi$ acting on dual vector fields, gives the bound in
\[
\langle F(f),U(a)F(f')\rangle_\gamma=\langle f,P^{\perp}(\tau_a\cdot f')\rangle=O\left(|a|^{-3}\right)\qquad(|a|\to\infty).
\]
\item[(2)]
The observables enter Lemma~\ref{lem:cube-decay} and Theorem~\ref{thm:clustering-rate} through their $L^2$ norms alone,
whereas the corresponding estimate for Yang--Mills theory in four dimensions works with supremum norms and with reflection sizes of the observables.
$L^2$ norms suffice here for two reasons:
only translations by at least $2r$ in the direction $e_j$ are expressed through the reconstruction
(in the proof of Lemma~\ref{lem:momentum-density} this is the vanishing of the coefficient $z_0$),
and the class map $[\,\cdot\,]_j$ is a contraction from $L^2(\Sigma_j^{+},\mu)$ into $\mathcal{H}_j$ \cite[Proposition~8.19(i)]{NoncompactQED},
so that every $\Phi\in L^2(\Sigma_{(-r,r)^3},\mu)$ satisfies
\[
\|\psi^{j}_{\Phi}\|^2=\langle\hat\theta_j\Phi,U(2r\cdot e_j)\Phi\rangle_\mu\le\|\Phi\|_2^2 .
\]
\end{enumerate}
\end{remark}

\subsection{Proof of Theorem~\ref{thm:continuum-limits}}\label{sec:continuum-proof}

\begin{proof}[Proof of Theorem~\ref{thm:continuum-limits}]
Throughout, the parameters are those of the theorem, the laws $P^{\mathrm{ch}}_{N,T}$ are those of Definition~\ref{def:cutoffs},
and limit points, the sets $\mathcal{V}_N$ and volume-first limit points are those of Definition~\ref{def:limit-points};
a Borel probability measure on $H_{-4}$ is identified with a law on $(\mathcal{S}',\Sigma)$ as in Definition~\ref{def:laws}.
For $\varrho>0$ let
\[
\mathcal{B}_{\varrho}:=\left\{F\in H_{-4}:\|F\|_{-3}\le\varrho\right\}
\]
be the closed ball of radius $\varrho$ of $H_{-3}$.
By \cite[Theorem~7.12\,(2)]{NoncompactQED}, where this ball is called $K_\varrho$, it is a compact subset of $H_{-4}$;
in particular it is closed in $H_{-4}$.

\emph{Clause~\eqref{it:continuum-soft}.}
This clause is the content of Proposition~\ref{prop:chamber-limits}.
Part~(a) of that proposition gives the bound $61$ for the second moments,
the tightness of the family $\left(P^{\mathrm{ch}}_{N,T}\right)$ on $H_{-4}$ and on $\mathcal{S}'$,
and the convergent subsequences.
Part~(b) gives every limit point the eight properties of \cite[Theorem~3.1, Part~I\,(b)]{NoncompactQED},
read with the family $\left(P^{\mathrm{ch}}_{N,T}\right)$ in place of the families of laws considered there.
Among them are the bound of the Laplace transform by that of the free Maxwell law, which implies \eqref{cond:gaussian-domination},
the closedness \eqref{cond:closedness},
the invariance under the translations of $\R^3$ and under $\Ocub$, hence under $\R^3\rtimes\Ocub$,
and the reflection positivity at the coordinate planes.
Part~(c) gives \eqref{cond:soft-hypothesis} for every limit point,
and for almost every ergodic component the properties listed in \cite[Theorem~3.1, Part~I\,(c)]{NoncompactQED}.
Part~(d) says that iterated limits, in either order, are limit points.

\emph{Clause~\eqref{it:continuum-volume-first}: existence.}
Fix an integer $N\ge0$ and a sequence of integers $T_j\to\infty$.
By Proposition~\ref{prop:chamber-limits}(a) the laws $P^{\mathrm{ch}}_{N,T_j}$, $j\ge1$, form a tight family on the separable Hilbert space $H_{-4}$;
so, by Prokhorov's theorem \cite[Theorem~5.1]{Billingsley1999}, a subsequence converges weakly on $H_{-4}$.
The indices $T_j$ of this subsequence still tend to infinity,
and its limit therefore lies in $\mathcal{V}_N$.
Hence $\mathcal{V}_N$ is not empty.

Next let $\mu^{(N)}\in\mathcal{V}_N$ be the weak limit on $H_{-4}$ of $P^{\mathrm{ch}}_{N,T_j}$ along a sequence $T_j\to\infty$, and let $\varrho>0$.
At every cutoff $(N,T)$ the bound $61$ and Chebyshev's inequality give
$P^{\mathrm{ch}}_{N,T}\left(H_{-4}\setminus\mathcal{B}_\varrho\right)\le61\varrho^{-2}$,
as in step~\eqref{it:admissible-tightness} of the proof of Proposition~\ref{prop:admissible-families}.
The set $\mathcal{B}_\varrho$ is closed, and the portmanteau theorem \cite[Theorem~2.1]{Billingsley1999} yields
\[
\mu^{(N)}\left(\mathcal{B}_\varrho\right)\ge\limsup_{j\to\infty}P^{\mathrm{ch}}_{N,T_j}\left(\mathcal{B}_\varrho\right)\ge1-61\varrho^{-2}.
\]
This lower bound holds for every $N$ and every element of $\mathcal{V}_N$, and $\mathcal{B}_\varrho$ is compact;
so the union of the sets $\mathcal{V}_N$ is a tight family of laws on $H_{-4}$.
By Prokhorov's theorem, every sequence $\mu^{(N_i)}\in\mathcal{V}_{N_i}$ with $N_i\to\infty$ has a subsequence that converges weakly on $H_{-4}$,
and by Definition~\ref{def:limit-points} the limit is a volume-first limit point.
Since no $\mathcal{V}_N$ is empty, one may choose $\mu^{(N)}\in\mathcal{V}_N$ for every $N$ and take $N_i:=i$;
so volume-first limit points exist.

From here on $\mu$ is a volume-first limit point.
By Proposition~\ref{prop:nine-families} it is a limit point,
so that it has every property that clause~\eqref{it:continuum-soft} asserts of a limit point;
and $\mu\in\mathcal{C}_9$.

\emph{Clause~\eqref{it:continuum-nine-planes}.}
By Proposition~\ref{prop:nine-families}, $\mu$ is invariant under the translations of $\R^3$,
and for every $c\in\R$ and every plane $\{\ell=c\}$ of the nine families of Definition~\ref{def:laws}
it is invariant under the reflection $r_{\ell,c}$ in this plane and reflection positive at it.
It remains to show the positivity for all observables in $L^2$ of the open side $\{\ell>c\}$.
Fix $\ell$ and $c$, and let $\hat r_{\ell,c}$ be the action of $r_{\ell,c}$ on functions on $\mathcal{S}'$ (Definition~\ref{def:isometry-actions}).
As $r_{\ell,c}$ is an involution, $\left(\hat r_{\ell,c}\Phi\right)(\omega)=\Phi(r_{\ell,c}\,\omega)$ for $\omega\in\mathcal{S}'$, and $\hat r_{\ell,c}F(f)=F(r_{\ell,c}\cdot f)$.
The invariance $(r_{\ell,c})_*\mu=\mu$ gives
\[
\int\mathrm{d}\mu\,\left|\hat r_{\ell,c}\Phi\right|^2=\int\mathrm{d}\mu\,\left|\Phi\right|^2\qquad\left(\Phi\in L^2(\mu)\right),
\]
and $\hat r_{\ell,c}^2=1$;
so $\hat r_{\ell,c}$ is a unitary operator of $L^2(\mu)$.
The real forms in $C^{\infty}_{\mathrm c}(\{\ell>c\};\Lambda^2)$ make up a real linear subspace of $\mathcal{S}$,
and the functions $F(f)$, $f$ in this subspace, generate the $\sigma$-algebra $\Sigma_{\{\ell>c\}}$.
By \cite[Lemma~7.14]{NoncompactQED}, applied to this subspace,
the complex linear span of the functions $\exp\left(iF(f)\right)$, $f\in C^{\infty}_{\mathrm c}(\{\ell>c\};\Lambda^2)$ real,
is dense in $L^2(\Sigma_{\{\ell>c\}},\mu)$.
Let
\[
\Phi=\sum_pa_p\cdot\exp\left(iF(f_p)\right)
\]
be a finite sum of this kind, with $a_p\in\C$ and real $f_p\in C^{\infty}_{\mathrm c}(\{\ell>c\};\Lambda^2)$.
Then $\hat r_{\ell,c}\Phi=\sum_pa_p\cdot\exp\left(iF(r_{\ell,c}\cdot f_p)\right)$,
and with the scalar product $\langle\Phi,\Phi'\rangle_\mu=\int\mathrm{d}\mu\,\overline{\Phi}\cdot\Phi'$ of $L^2(\mu)$ one finds
\[
\left\langle\hat r_{\ell,c}\Phi,\Phi\right\rangle_\mu
=\sum_{p,q}\overline{a_p}\cdot a_q\cdot\mu\left(\exp\left(i\left(F(f_q)-F(r_{\ell,c}\cdot f_p)\right)\right)\right)\ge0,
\]
by the reflection positivity of $\mu$ at $\{\ell=c\}$.
For $\Phi,\Phi'\in L^2(\mu)$ the unitarity of $\hat r_{\ell,c}$ gives
\[
\left|\left\langle\hat r_{\ell,c}\Phi,\Phi\right\rangle_\mu-\left\langle\hat r_{\ell,c}\Phi',\Phi'\right\rangle_\mu\right|
\le\|\Phi-\Phi'\|_2\cdot\left(\|\Phi\|_2+\|\Phi'\|_2\right),
\]
so the quadratic form $\Phi\mapsto\left\langle\hat r_{\ell,c}\Phi,\Phi\right\rangle_\mu$ is continuous on $L^2(\mu)$.
Its non-negativity on a dense subset of $L^2(\Sigma_{\{\ell>c\}},\mu)$ therefore extends to the whole space:
\begin{equation}\label{it:open-side-extension}
\left\langle\hat r_{\ell,c}\Phi,\Phi\right\rangle_\mu\ge0\qquad\text{for every }\Phi\in L^2(\Sigma_{\{\ell>c\}},\mu).
\end{equation}

\emph{Clause~\eqref{it:continuum-cone}.}
Let $j\neq i$.
The planes $\{x_j=0\}$ and $\{x_j+x_i=0\}$ belong to the nine families.
As $\mu\in\mathcal{C}_9$, the law $\mu$ is invariant under the translations of $\R^3$, satisfies \eqref{cond:reflection-invariance} and \eqref{cond:reflection-positivity},
and is invariant and reflection positive at the plane $\{x_j+x_i=0\}$:
these are the hypotheses of Theorem~\ref{thm:cone} for the pair $(j,i)$.
Part~(c) of that theorem says that the joint spectral measure of $(\mathsf{H}_j,\mathsf{P}_i)$ is carried by $\{(\lambda,k_i):\lambda\ge|k_i|\}$
and that $\mathsf{H}_j\ge|\mathsf{P}_i|$ with the constant one.
Part~(a), together with $V_j(s\cdot e_i)^*=V_j(-s\cdot e_i)$, gives
\[
\left(1-\exp\left(-s\cdot\mathsf{H}_j\right)\right)^2\ge\exp\left(-s\cdot\mathsf{H}_j\right)\cdot\left(2-V_j(s\cdot e_i)-V_j(-s\cdot e_i)\right)\qquad(s>0),
\]
and part~(d) is the five-point inequality for every $t>0$ and every $\Phi\in L^2(\Sigma_j^+,\mu)$.
The law $\mu$ satisfies \eqref{cond:gaussian-domination}, by the bound of the Laplace transform among the properties of clause~\eqref{it:continuum-soft}.
So Corollary~\ref{cor:unbounded-observables}, with the tensor tests of Lemma~\ref{lem:tensor-tests}, applies to $\mu$ as well:
the space $L^2(\Sigma_j^+,\mu)$ contains the polynomials and the exponentials in the $F(f)$,
$f\in L^2(\R^3;\Lambda^2)$ with $f=0$ almost everywhere on $\{x_j<0\}$,
and the five-point inequality holds in the form that the corollary gives for the Schwinger functions.

\emph{Clause~\eqref{it:continuum-vacuum}.}
Proposition~\ref{prop:zero-momentum} applies to $\mu\in\mathcal{C}_9$.
Its part~(b) gives $\ker\mathsf{H}_j\subset\ker\mathsf{P}_i$ for $i\neq j$,
and its part~(c) gives $P^j_\mu=P^{123}_\mu$ for $j=1,2,3$.
As $\mu$ satisfies \eqref{cond:soft-hypothesis}, the latter is the property~\eqref{it:projections-coincide} of Definition~\ref{def:symmetry-conditions}.

\emph{Clause~\eqref{it:continuum-rate}.}
Theorem~\ref{thm:clustering-rate}, applied to the law $\mu\in\mathcal{C}_9$, is this clause, with the absolute constants $C_0$ and $C_1$ of that theorem.

\emph{Clause~\eqref{it:continuum-components}.}
The law $\mu$ satisfies \eqref{cond:soft-hypothesis} by Proposition~\ref{prop:chamber-limits}(c), it lies in $\mathcal{C}_9$,
and $\mu(\{dF=0\})=1$ by the closedness among the properties of clause~\eqref{it:continuum-soft}.
Let $\mu=\int\mathrm{d}\nu(\xi)\,\mu_\xi$ be its ergodic decomposition under $\R^3$;
by \cite[Corollary~8.9]{NoncompactQED}, $\nu$ is the only probability measure carried by the laws ergodic under $\R^3$ that represents $\mu$ in this way.
Proposition~\ref{prop:ergodic-components} applies to $\mu$, and gives for $\nu$-almost every $\xi$:
the law $\mu_\xi$ is invariant and ergodic under $\R^3$ and lies in $\mathcal{C}_9$ (part~(a));
it satisfies \eqref{cond:gaussian-domination} with the same right side (part~(b));
it is invariant under $\R^3\rtimes\Ocub$ (part~(c)), so that $\mu$ has the property~\eqref{it:components-cubic} of Definition~\ref{def:symmetry-conditions};
one has $\ker\mathsf{H}_j(\mu_\xi)=\C\cdot\Omega_j$ and $\mathsf{H}_j(\mu_\xi)\ge|\mathsf{P}_i|$ for $j\neq i$ (part~(d));
$\mu_\xi\left(\overline{\Phi}\cdot U(a)\Phi'\right)\to\overline{\mu_\xi(\Phi)}\cdot\mu_\xi(\Phi')$ for all $\Phi,\Phi'\in L^2(\mu_\xi)$ as $|a|\to\infty$ in any manner (part~(e));
and $\mu_\xi(\{dF=0\})=1$ (part~(f)).
Let $\xi$ be such that these assertions hold.
Theorem~\ref{thm:clustering-rate} applies to the law $\mu_\xi\in\mathcal{C}_9$.
Since $\mu_\xi$ is ergodic under $\R^3$, the term $\left\langle\Phi,P^{123}_{\mu_\xi}\Phi'\right\rangle_{\mu_\xi}$ subtracted in that theorem equals $\overline{\mu_\xi(\Phi)}\cdot\mu_\xi(\Phi')$,
while $\langle\Phi,U(a)\Phi'\rangle_{\mu_\xi}=\mu_\xi\left(\overline{\Phi}\cdot U(a)\Phi'\right)$.
Hence, for every $r>0$, all $\Phi,\Phi'\in L^2(\Sigma_{(-r,r)^3},\mu_\xi)$ and all $a\in\R^3$ with $|a|_\infty\ge C_1\cdot r$,
\[
\left|\mu_\xi\left(\overline{\Phi}\cdot U(a)\Phi'\right)-\overline{\mu_\xi(\Phi)}\cdot\mu_\xi(\Phi')\right|
\le C_0\cdot\|\Phi\|_2\cdot\|\Phi'\|_2\cdot\frac{r}{|a|_\infty}.
\]
Finally, as $\mu$ satisfies \eqref{cond:soft-hypothesis}, \cite[Corollary~8.9]{NoncompactQED} applies to $\mu$,
and $\nu$-almost every $\mu_\xi$ has the properties~1--6 of that corollary.
The exceptional sets of $\xi$ that occur in this paragraph are finitely many $\nu$-null sets.
\end{proof}

\begin{remark}\label{rem:volume-first-moments}
Let $\mu$ be a volume-first limit point:
$\mu^{(N_i)}\to\mu$ weakly on $H_{-4}$ for some $\mu^{(N_i)}\in\mathcal{V}_{N_i}$, $N_i\to\infty$,
where $P^{\mathrm{ch}}_{N_i,T_j(N_i)}\to\mu^{(N_i)}$ weakly on $H_{-4}$ as $j\to\infty$, with $T_j(N_i)\to\infty$.
The last of the eight properties in \cite[Theorem~3.1, Part~I\,(b)]{NoncompactQED}, which clause~\eqref{it:continuum-soft} of Theorem~\ref{thm:continuum-limits} gives to $\mu$,
is the convergence of the generating functions, of the moments and of the cumulants to those of $\mu$
along every sequence of cutoffs, both coordinates of which tend to infinity, along which the laws $P^{\mathrm{ch}}_{N,T}$ converge to $\mu$ weakly on $H_{-4}$.
A diagonal choice provides such a sequence.
The Prokhorov metric $d_{\mathrm{P}}$ metrises weak convergence of laws on $H_{-4}$ \cite[Theorem~6.8]{Billingsley1999};
choose $j_i$ with $T_{j_i}(N_i)\ge i$ and $d_{\mathrm{P}}\left(P^{\mathrm{ch}}_{N_i,T_{j_i}(N_i)},\mu^{(N_i)}\right)\le1/i$.
Then $d_{\mathrm{P}}\left(P^{\mathrm{ch}}_{N_i,T_{j_i}(N_i)},\mu\right)\le1/i+d_{\mathrm{P}}\left(\mu^{(N_i)},\mu\right)\to0$,
so $P^{\mathrm{ch}}_{N_i,T_{j_i}(N_i)}\to\mu$ weakly on $H_{-4}$, and both coordinates of the cutoffs $\left(N_i,T_{j_i}(N_i)\right)$ tend to infinity.

The generating functions and the moments converge also along the two limits taken one after the other,
first $j\to\infty$ at a fixed $N_i$, then $i\to\infty$.
Let $f_1,\dots,f_s$ be complex test forms.
For $z\in\C^s$ put $f_z:=\Re\sum_pz_p\cdot f_p$, a real test form;
then
\[
\left|\exp\sum_{p=1}^sz_p\cdot F(f_p)\right|^2=\exp F(2f_z)\qquad\left(F\in H_{-4}\right).
\]
For every cutoff $(N,T)$, Proposition~\ref{prop:cutoff-law} shows that the law $\mathbb{P}_{N,T}$ defined there satisfies the hypotheses of Proposition~\ref{prop:admissible-families};
and $P^{\mathrm{ch}}_{N,T}$ is the law of $F^\sharp$ under $\mathbb{P}_{N,T}$ (Definition~\ref{def:cutoffs}).
Hence, at every cutoff $(N,T)$, step~\eqref{it:admissible-bounds} of the proof of Proposition~\ref{prop:admissible-families} gives
\[
\int\mathrm{d}P^{\mathrm{ch}}_{N,T}(F)\,\exp F(2f_z)\le\exp\left(2p_1(f_z)^2\right),
\]
where $p_1$ is the continuous seminorm on $\mathcal{S}$ of \cite[Lemma~7.5]{NoncompactQED};
the right side is bounded for $z$ in a compact subset of $\C^s$, uniformly in $(N,T)$.
For the volume limits, let $\mu^{(N)}\in\mathcal{V}_N$ and $\eta=L^{-N}$.
By Proposition~\ref{prop:volume-limits}(b), $\mu^{(N)}=(\sharp_\eta)_*\nu^{(N)}$,
where $\nu^{(N)}$ is the weak limit on $\Omega$ of the laws of the cubic tori,
and $\nu^{(N)}$ satisfies \eqref{cond:gaussian-domination}.
So Proposition~\ref{prop:volume-limits}(a) applies to $\nu^{(N)}$:
one has $\nu^{(N)}(\Omega^{(\eta)})=1$, and with \eqref{eq:continuum-step-pairing}, for every real $f\in\mathcal{S}$,
\[
\int\mathrm{d}\mu^{(N)}(F)\,\exp F(f)
=\int\mathrm{d}\nu^{(N)}(F^{\mathrm{u}})\,\exp\left\langle f_{[\eta]},F^{\mathrm{u}}\right\rangle
\le\exp\left(\frac{1}{2}\|f\|_2^2\right).
\]
Thus $\mu^{(N)}$ satisfies \eqref{cond:gaussian-domination}, and the choice $f=2f_z$ gives
\[
\int\mathrm{d}\mu^{(N)}(F)\,\exp F(2f_z)\le\exp\left(2\|f_z\|_2^2\right)\le\exp\left(2\left(\sum_{p=1}^s|z_p|\cdot\|f_p\|_2\right)^2\right),
\]
a bound that is uniform in $N$.
Each of the two sequences, $\left(P^{\mathrm{ch}}_{N_i,T_j(N_i)}\right)_{j\ge1}$ at a fixed $i$ and $\left(\mu^{(N_i)}\right)_{i\ge1}$,
converges weakly on $H_{-4}$, to $\mu^{(N_i)}$ and to $\mu$ respectively,
and along each of them the integrals of $\left|\exp\sum_pz_p\cdot F(f_p)\right|^2$ are bounded for $z$ in a compact set, uniformly along the sequence.
With these bounds in place of those of step~\eqref{it:admissible-bounds},
the argument of step~\eqref{it:admissible-moments} of the proof of Proposition~\ref{prop:admissible-families} applies to both sequences.
For fixed $z$ the real and imaginary parts of $F\mapsto\exp\sum_pz_p\cdot F(f_p)$ are continuous functions on $H_{-4}$.
Their laws converge weakly by the mapping theorem \cite[Theorem~2.7]{Billingsley1999};
they are uniformly integrable, because their second moments are bounded along the sequence;
so their integrals converge to those under the limit law \cite[Theorem~3.5]{Billingsley1999}.
The generating functions of the laws of either sequence are entire functions of $z$:
for the laws $P^{\mathrm{ch}}_{N_i,T_j(N_i)}$ by step~\eqref{it:admissible-bounds} of the proof of Proposition~\ref{prop:admissible-families},
and for $\mu^{(N_i)}$ because, by \eqref{cond:gaussian-domination}, every $F(f)$ with $f\in\mathcal{S}$ real has finite exponential moments of every order.
They are bounded on compact subsets of $\C^s$ uniformly along the sequence;
by Cauchy's estimates they are equi-Lipschitz on compact sets, so that their pointwise convergence is uniform on compact sets,
and by Cauchy's integral formula all their partial derivatives converge, again uniformly on compact sets.
The partial derivatives at $z=0$ are the moments of $F(f_1),\dots,F(f_s)$.
Therefore the generating functions converge locally uniformly in $z$, and all moments converge,
first from $P^{\mathrm{ch}}_{N_i,T_j(N_i)}$ to $\mu^{(N_i)}$ as $j\to\infty$, and then from $\mu^{(N_i)}$ to $\mu$ as $i\to\infty$.
\end{remark}

The last of the eight properties in \cite[Theorem~3.1, Part~I\,(b)]{NoncompactQED} concerns the laws $P^{\mathrm{ch}}_{N,T}$ along a sequence of cutoffs;
every other property that Theorem~\ref{thm:continuum-limits} asserts of a volume-first limit point is a property of the law alone.
A law $\mu$ on $\mathcal{S}'$ is said to have the properties $\mathscr{P}$ if the following holds.
\begin{enumerate}
\item[\textup{($\mathscr{P}$)}]\itemmark{$\mathscr{P}$}\label{cond:limit-properties}
The law $\mu$ has the properties of \cite[Theorem~3.1, Part~I\,(b)]{NoncompactQED} other than the last,
and the property~(c) of Part~I there;
and it has the properties \eqref{it:continuum-nine-planes}--\eqref{it:continuum-components} in clause~\eqref{it:continuum-volume-first} of Theorem~\ref{thm:continuum-limits},
with $\mu$ in place of the volume-first limit point.
\end{enumerate}

\begin{lemma}[Soft laws with nine families of mirrors]\label{lem:soft-laws}
Let $\mu$ be a law on $\mathcal{S}'$ that satisfies \eqref{cond:soft-hypothesis} and \eqref{cond:closedness},
is invariant under the map $F\mapsto-F$ of $\mathcal{S}'$, and lies in $\mathcal{C}_9$.
Then $\mu$ has the properties~\eqref{cond:limit-properties}.
\end{lemma}

\begin{proof}
The law $\mu$ is soft in the sense of \cite[Definition~21.35]{NoncompactQED}:
it satisfies \eqref{cond:soft-hypothesis}, it is carried by the closed distributions, and it is invariant under $F\mapsto-F$.
By \cite[Lemma~21.36]{NoncompactQED} it therefore has the properties of \cite[Theorem~3.1, Part~I\,(b)]{NoncompactQED} other than the last,
and the property~(c) of Part~I there.
The remaining properties are obtained for $\mu$ from the same statements as for a volume-first limit point in the proof of Theorem~\ref{thm:continuum-limits},
each of which applies to $\mu$.

For every plane $\{\ell=c\}$ of the nine families, $\mu$ is invariant under $r_{\ell,c}$ and reflection positive at $\{\ell=c\}$, since $\mu\in\mathcal{C}_9$.
With these two properties, the argument that proves \eqref{it:open-side-extension},
which takes the density of the exponentials from \cite[Lemma~7.14]{NoncompactQED},
gives $\left\langle\hat r_{\ell,c}\Phi,\Phi\right\rangle_\mu\ge0$ for every $\Phi\in L^2(\Sigma_{\{\ell>c\}},\mu)$;
together with $\mu\in\mathcal{C}_9$ this is the property~\eqref{it:continuum-nine-planes}.

Membership in $\mathcal{C}_9$ also gives the hypotheses of Theorem~\ref{thm:cone} for every pair $j\neq i$,
and $\mu$ satisfies \eqref{cond:gaussian-domination}, which is a part of \eqref{cond:soft-hypothesis}.
Hence Theorem~\ref{thm:cone} and Corollary~\ref{cor:unbounded-observables}, with the tensor tests of Lemma~\ref{lem:tensor-tests},
give the property~\eqref{it:continuum-cone}.
Proposition~\ref{prop:zero-momentum}(b), (c) gives the property~\eqref{it:continuum-vacuum},
and Theorem~\ref{thm:clustering-rate} gives the property~\eqref{it:continuum-rate}.

Let $\mu=\int\mathrm{d}\nu(\xi)\,\mu_\xi$ be the ergodic decomposition of $\mu$ under $\R^3$.
Proposition~\ref{prop:ergodic-components} applies to $\mu$, which satisfies \eqref{cond:soft-hypothesis} and lies in $\mathcal{C}_9$;
it gives the assertions of its parts~(a)--(e) for $\nu$-almost every $\mu_\xi$,
and its part~(f), with $\mu(\{dF=0\})=1$, gives $\mu_\xi(\{dF=0\})=1$ for $\nu$-almost every $\xi$.
For every $\xi$ outside these null sets the law $\mu_\xi$ lies in $\mathcal{C}_9$ and is ergodic under $\R^3$;
Theorem~\ref{thm:clustering-rate} applies to $\mu_\xi$, and its subtracted term $\left\langle\Phi,P^{123}_{\mu_\xi}\Phi'\right\rangle_{\mu_\xi}$ equals $\overline{\mu_\xi(\Phi)}\cdot\mu_\xi(\Phi')$,
which gives the bound with the rate $r/|a|_\infty$ for $\mu_\xi$.
Finally \cite[Corollary~8.9]{NoncompactQED}, applied to $\mu$, gives the properties~1--6 listed there for $\nu$-almost every $\mu_\xi$.
These are the assertions of the property~\eqref{it:continuum-components}.
\end{proof}

\begin{example}[The zero field]\label{ex:zero-field}
Let $\mu_0$ be the point mass at the distribution $F=0$, a law on $\mathcal{S}'$.
It is the law of $a^{\frac{1}{2}}\cdot F$ under the free Maxwell law for $a=0$,
a value of $a$ outside the range $0<a\le1$ of Definition~\ref{def:scaled-maxwell}.
For every real $f\in\mathcal{S}$,
\[
\mu_0\left(\exp\left(iF(f)\right)\right)=1,\qquad\mu_0\left(\exp F(f)\right)=1\le\exp\left(\frac{1}{2}\|f\|_2^2\right);
\]
so $\mu_0$ satisfies \eqref{cond:gaussian-domination}.
Every Euclidean map acts on $\mathcal{S}'$ by a linear bijection, which fixes the distribution $0$;
hence $\mu_0$ is invariant under every Euclidean map,
in particular under the translations of $\R^3$, under $\Ocub$ and under every reflection $r_{n,c}$ in a plane $\{n\cdot x=c\}$.
The map $F\mapsto-F$ fixes $0$ as well, and $dF=0$ at $F=0$, so that $\mu_0(\{dF=0\})=1$.
For every plane $\{n\cdot x=c\}$ and every finite family of coefficients $a_p\in\C$ and of real tests $f_p\in C^{\infty}_{\mathrm c}(\{n\cdot x>c\};\Lambda^2)$,
\[
\sum_{p,q}\overline{a_p}\cdot a_q\cdot\mu_0\left(\exp\left(i\left(F(f_q)-F(r_{n,c}\cdot f_p)\right)\right)\right)=\left|\sum_pa_p\right|^2\ge0,
\]
so $\mu_0$ is reflection positive at every plane.
It follows that $\mu_0$ satisfies \eqref{cond:soft-hypothesis} and \eqref{cond:closedness}, is invariant under $F\mapsto-F$, and lies in $\mathcal{C}_9$.
By Lemma~\ref{lem:soft-laws}, $\mu_0$ has the properties~\eqref{cond:limit-properties}.
Thus the properties~\eqref{cond:limit-properties}, which Theorem~\ref{thm:continuum-limits} asserts of every volume-first limit point, do not exclude the zero field.
\end{example}

\section{Averaged law and joint limits}\label{sec:averaged}

\subsection{Separation and relative entropy}\label{sec:separation-entropy}

\begin{definition}[Complex trigonometric polynomials]\label{def:complex-polynomials}
The complex linear span of the characters $\varphi_\beta$, $\beta\in C^2_{\mathrm{c}}$, of Definition~\ref{def:characters} is denoted by $\mathcal{T}_{\C}$.
Its elements are called complex trigonometric polynomials.
\end{definition}

By Definition~\ref{def:characters}, the space $\mathcal{T}$ of real trigonometric polynomials is the real-valued part of $\mathcal{T}_{\C}$.
For every $\beta\in C^2_{\mathrm{c}}$ the functions $F\mapsto\cos\langle\beta,F\rangle$ and $F\mapsto\sin\langle\beta,F\rangle$ on $\Omega$ belong to $\mathcal{T}$,
since
\[
\cos\langle\beta,F\rangle=\frac{\varphi_\beta(F)+\varphi_{-\beta}(F)}{2},
\qquad
\sin\langle\beta,F\rangle=\frac{\varphi_\beta(F)-\varphi_{-\beta}(F)}{2i}.
\]

\begin{lemma}[Separation]\label{lem:separation}
Let $\mathcal{R}$ be a convex set of laws on $\Omega$ that is compact for the weak topology, and let $\nu$ be a law on $\Omega$.
Suppose that for every $O\in\mathcal{T}$ there is $\nu_0\in\mathcal{R}$ with $\nu(O)\le\nu_0(O)$.
Then $\nu\in\mathcal{R}$.
\end{lemma}

\begin{proof}
Let $\mathcal{M}(\Omega)$ be the real vector space of the finite signed Borel measures on $\Omega$;
every law on $\Omega$ is an element of it.
Each $O\in\mathcal{T}$ is a bounded continuous function on $\Omega$.
Indeed, a character $\varphi_\beta$ has modulus one,
and $\langle\beta,F\rangle$ is a finite linear combination of the coordinates $F_\sigma$, $\sigma\in\supp\beta$,
hence a continuous function of $F$ for the product topology of $\Omega$.
For $O\in\mathcal{T}$ the integral $\lambda(O)=\int\mathrm{d}\lambda(F)\,O(F)$ therefore exists for every $\lambda\in\mathcal{M}(\Omega)$,
and $\lambda\mapsto\lambda(O)$ is a linear map from $\mathcal{M}(\Omega)$ to $\R$.
As $\mathcal{T}$ is a real vector space, these maps form a vector space of linear functionals on $\mathcal{M}(\Omega)$.
This space of functionals is identified with $\mathcal{T}$:
the point masses belong to $\mathcal{M}(\Omega)$, and the point mass at $F\in\Omega$ assigns to $O$ the value $O(F)$,
so that $O$ is determined by its functional.
The space $\mathcal{M}(\Omega)$ carries the topology $\sigma(\mathcal{M}(\Omega),\mathcal{T})$,
the coarsest topology for which all maps $\lambda\mapsto\lambda(O)$, $O\in\mathcal{T}$, are continuous.
It is defined by the seminorms $\lambda\mapsto|\lambda(O)|$, $O\in\mathcal{T}$, and is therefore locally convex.

\emph{The functionals separate the points of $\mathcal{M}(\Omega)$.}
Let $\lambda\in\mathcal{M}(\Omega)$ satisfy $\lambda(O)=0$ for every $O\in\mathcal{T}$,
and let $\lambda=\lambda^{+}-\lambda^{-}$ be its Jordan decomposition into two finite Borel measures.
For $\beta\in C^2_{\mathrm{c}}$ one has $\varphi_\beta=\cos\langle\beta,\cdot\,\rangle+i\sin\langle\beta,\cdot\,\rangle$,
where $\cos\langle\beta,\cdot\,\rangle$ and $\sin\langle\beta,\cdot\,\rangle$ belong to $\mathcal{T}$;
hence $\lambda^{+}(\varphi_\beta)=\lambda^{-}(\varphi_\beta)$.
Let $\Sigma_0$ be a finite set of plaquettes,
and let $\lambda^{+}_{\Sigma_0}$ and $\lambda^{-}_{\Sigma_0}$ be the images of $\lambda^{+}$ and $\lambda^{-}$ under the map $F\mapsto(F_\sigma)_{\sigma\in\Sigma_0}$ from $\Omega$ to $\R^{\Sigma_0}$.
If $\beta$ is supported in $\Sigma_0$, then
\[
\varphi_\beta(F)=\exp\left(i\sum_{\sigma\in\Sigma_0}\beta(\sigma)\cdot F_\sigma\right)
\]
depends on $(F_\sigma)_{\sigma\in\Sigma_0}$ only,
and $\lambda^{\pm}(\varphi_\beta)$ is the value of the Fourier transform of $\lambda^{\pm}_{\Sigma_0}$ at the point $(\beta(\sigma))_{\sigma\in\Sigma_0}$;
every point of $\R^{\Sigma_0}$ is of this form.
So the marginal $\lambda^{+}_{\Sigma_0}-\lambda^{-}_{\Sigma_0}$ of $\lambda$ on $\R^{\Sigma_0}$ has a vanishing Fourier transform:
the finite measures $\lambda^{+}_{\Sigma_0}$ and $\lambda^{-}_{\Sigma_0}$ have the same Fourier transform,
and they are equal by the uniqueness theorem for Fourier transforms of finite measures on $\R^{\Sigma_0}$.
Consequently $\lambda^{+}$ and $\lambda^{-}$ coincide on the cylinder sets,
the preimages of Borel subsets of $\R^{\Sigma_0}$ under $F\mapsto(F_\sigma)_{\sigma\in\Sigma_0}$, with $\Sigma_0$ finite.
The cylinder sets include $\Omega$ and are stable under finite intersections,
and the Borel $\sigma$-algebra of $\Omega$ is generated by them (see~\eqref{eq:layers-law}).
Two finite measures that agree on such a family agree on the $\sigma$-algebra generated by it (Dynkin's theorem);
hence $\lambda^{+}=\lambda^{-}$, that is, $\lambda=0$.

\emph{Dual space and compactness.}
The functionals $\lambda\mapsto\lambda(O)$, $O\in\mathcal{T}$, thus form a separating vector space of linear functionals on $\mathcal{M}(\Omega)$.
The topology defined by a separating vector space of linear functionals is a locally convex vector space topology,
and the continuous linear functionals for it are exactly the elements of that space \cite[Theorem~3.10]{Rudin1991}.
Hence $\mathcal{M}(\Omega)$ is a locally convex Hausdorff topological vector space,
and every continuous linear functional on it is of the form $\lambda\mapsto\lambda(O)$ with $O\in\mathcal{T}$.
The weak topology of laws on $\Omega$ (see~\eqref{eq:layers-law}) is the coarsest topology for which $\mu\mapsto\mu(G)$ is continuous for every bounded continuous real function $G$ on $\Omega$,
and the elements of $\mathcal{T}$ are such functions.
It follows that the inclusion of $\mathcal{R}$, with the weak topology, into $\mathcal{M}(\Omega)$ is continuous,
and $\mathcal{R}$ is a compact convex subset of $\mathcal{M}(\Omega)$.
It is not empty: the hypothesis, applied to $O=0$, provides an element $\nu_0$ of $\mathcal{R}$.

\emph{Separation.}
Suppose that $\nu\notin\mathcal{R}$.
Then $\mathcal{R}$ and $\{\nu\}$ are disjoint non-empty convex subsets of the locally convex space $\mathcal{M}(\Omega)$;
the first is compact, and the second is closed because the topology is Hausdorff.
By the separation theorem for a compact convex set and a disjoint closed convex set in a locally convex space \cite[Theorem~3.4\,(b)]{Rudin1991},
there is a continuous linear functional on $\mathcal{M}(\Omega)$ whose supremum over $\mathcal{R}$ is smaller than its value at $\nu$.
By the description of the dual space, this functional is $\lambda\mapsto\lambda(O)$ for some $O\in\mathcal{T}$;
as $\mathcal{R}$ is compact and the functional is continuous, the supremum is attained, and
\[
\max_{\lambda\in\mathcal{R}}\lambda(O)<\nu(O).
\]
For this $O$ the hypothesis yields $\nu_0\in\mathcal{R}$ with
\[
\nu(O)\le\nu_0(O)\le\max_{\lambda\in\mathcal{R}}\lambda(O),
\]
which is incompatible with the strict inequality above.
So $\nu$ lies in $\mathcal{R}$.
\end{proof}

\begin{definition}[Relative entropy]\label{def:relative-entropy}
Let $\mu$ and $\lambda$ be probability measures on a measurable space, for instance laws on $\Omega$, such that $\lambda\ll\mu$;
as for laws, $\mu(G)$ stands for the integral $\int\mathrm{d}\mu\,G$.
The relative entropy of $\lambda$ with respect to $\mu$ is
\[
H(\lambda\,|\,\mu):=\lambda\left(\log\frac{\mathrm{d}\lambda}{\mathrm{d}\mu}\right)\in[0,\infty].
\]
\end{definition}

The integral in this definition is well defined and non-negative.
The density $\mathrm{d}\lambda/\mathrm{d}\mu$ is positive $\lambda$-almost everywhere, and
\[
\lambda\left(\left(\frac{\mathrm{d}\lambda}{\mathrm{d}\mu}\right)^{-1}\right)
=\mu\left(\mathbf{1}\left[\frac{\mathrm{d}\lambda}{\mathrm{d}\mu}>0\right]\right)\le1.
\]
The logarithm of a positive number is smaller than the number;
so the negative part of $\log(\mathrm{d}\lambda/\mathrm{d}\mu)$ is bounded by $(\mathrm{d}\lambda/\mathrm{d}\mu)^{-1}$ and is $\lambda$-integrable,
and $H(\lambda\,|\,\mu)$ is defined as an element of $(-\infty,\infty]$.
If it is finite, Jensen's inequality for $\lambda$ and the concave function $\log$ gives
\[
-H(\lambda\,|\,\mu)
=\lambda\left(\log\left(\left(\frac{\mathrm{d}\lambda}{\mathrm{d}\mu}\right)^{-1}\right)\right)
\le\log\lambda\left(\left(\frac{\mathrm{d}\lambda}{\mathrm{d}\mu}\right)^{-1}\right)
\le0.
\]

\begin{lemma}[Jensen's inequality with relative entropy]\label{lem:jensen-entropy}
Let $\mu$ and $\lambda$ be probability measures on a measurable space with $\lambda\ll\mu$,
and let $\Phi$ be a bounded measurable real function on this space.
Then
\[
\log\mu(\exp\Phi)\ge\lambda(\Phi)-H(\lambda\,|\,\mu).
\]
\end{lemma}

\begin{proof}
If $H(\lambda\,|\,\mu)=\infty$, the right-hand side is $-\infty$, because $\lambda(\Phi)$ is finite, and there is nothing to prove.
Let $H(\lambda\,|\,\mu)<\infty$.
Then $\log(\mathrm{d}\lambda/\mathrm{d}\mu)$ is $\lambda$-integrable, and so is $\Phi-\log(\mathrm{d}\lambda/\mathrm{d}\mu)$.
The measure $\lambda$ has the density $\mathrm{d}\lambda/\mathrm{d}\mu$ with respect to $\mu$ and gives no mass to the set where this density vanishes.
Therefore
\[
\begin{aligned}
\mu(\exp\Phi)
&\ge\mu\left(\exp\Phi\cdot\mathbf{1}\left[\frac{\mathrm{d}\lambda}{\mathrm{d}\mu}>0\right]\right)
=\lambda\left(\exp\left(\Phi-\log\frac{\mathrm{d}\lambda}{\mathrm{d}\mu}\right)\right)\\
&\ge\exp\left(\lambda(\Phi)-\lambda\left(\log\frac{\mathrm{d}\lambda}{\mathrm{d}\mu}\right)\right)
=\exp\left(\lambda(\Phi)-H(\lambda\,|\,\mu)\right),
\end{aligned}
\]
where the second inequality is Jensen's inequality for the probability measure $\lambda$ and the convex function $\exp$.
Taking logarithms gives the assertion.
\end{proof}

\begin{example}
Let $\mu$ be a probability measure on a measurable space and $\mathcal{B}$ an event with $\mu(\mathcal{B})>0$.
The conditioned measure $\lambda:=\mu(\,\cdot\,|\,\mathcal{B})=\mu(\,\cdot\cap\mathcal{B})/\mu(\mathcal{B})$ satisfies $\lambda\ll\mu$, with the density
\[
\frac{\mathrm{d}\lambda}{\mathrm{d}\mu}=\mu(\mathcal{B})^{-1}\cdot\mathbf{1}_{\mathcal{B}}.
\]
Since $\lambda(\mathcal{B})=1$ and the density equals $\mu(\mathcal{B})^{-1}$ on $\mathcal{B}$,
\[
H(\lambda\,|\,\mu)=-\log\mu(\mathcal{B}).
\]
\end{example}

\subsection{Lattice laws with nine families of mirrors}\label{sec:nine-families}

\begin{proposition}[Tightness and soft limits]\label{prop:soft-limits}
Let $L\ge3$ be an odd integer (Definition~\ref{def:cutoffs}).
For an integer $N\ge0$ let $\eta=L^{-N}$,
and let $\sharp_\eta\colon\Omega\to H_{-4}$ be the step map of Definition~\ref{def:step-map}.
Let $\mathcal{N}$ be an infinite set of integers $N\ge0$.
For every $N\in\mathcal{N}$ let $\nu^{(N)}$ be a law on $\Omega$ (Definition~\ref{def:characters})
that satisfies \eqref{cond:gaussian-domination},
that is invariant under the maps $F^{\mathrm{u}}\mapsto g^{*}F^{\mathrm{u}}$, $g\in W$ (Definitions~\ref{def:cochains} and~\ref{def:mirrors}), and under $F^{\mathrm{u}}\mapsto-F^{\mathrm{u}}$,
and that is reflection positive, in the sense of Definition~\ref{def:reflection-positivity}, at every mirror $(\ell,c)$ with $\ell\in\mathcal{L}$ and $c\in\Adm(\ell)$.
Let
\[
\mu^{(N)}:=(\sharp_\eta)_{*}\nu^{(N)}\qquad\left(N\in\mathcal{N},\ \eta=L^{-N}\right)
\]
be the image of $\nu^{(N)}$ under $\sharp_\eta$.
Then $(\mu^{(N)})_{N\in\mathcal{N}}$ is a tight family of Borel probability measures on $H_{-4}$.
Moreover, let $(N_i)_{i\ge1}$ be a sequence in $\mathcal{N}$ with $N_i\to\infty$,
and let $\mu$ be a Borel probability measure on $H_{-4}$ such that $\mu^{(N_i)}\to\mu$ weakly on $H_{-4}$ as $i\to\infty$.
Then the following hold.
\begin{enumerate}
\item[\textup{(a)}]
The law $\mu$ satisfies \eqref{cond:soft-hypothesis} and \eqref{cond:closedness}, that is, $\mu(\{dF=0\})=1$,
and it is invariant under $F\mapsto-F$;
thus $\mu$ is soft in the sense of \cite[Definition~21.35]{NoncompactQED}.
Consequently $\mu$ has all the properties of a limit point listed in \cite[Theorem~3.1, Part~I\,(b)]{NoncompactQED} except the last one,
the convergence of moments and cumulants along a sequence of cutoffs;
and its ergodic decomposition under the translations of $\R^3$ has the properties listed in \cite[Theorem~3.1, Part~I\,(c)]{NoncompactQED}.
\item[\textup{(b)}]
The law $\mu$ belongs to the class $\mathcal{C}_9$ of Definition~\ref{def:laws}.
\end{enumerate}
\end{proposition}

\begin{proof}
\emph{Tightness.}
Fix $N\in\mathcal{N}$ and put $\eta=L^{-N}$.
Since $\nu^{(N)}$ satisfies \eqref{cond:gaussian-domination}, Proposition~\ref{prop:volume-limits}(a) applies to it:
the Borel set $\Omega^{(\eta)}\subset\Omega$ of Definition~\ref{def:step-map}, on which $\sharp_\eta$ takes its values in $H_{-3}$, satisfies $\nu^{(N)}(\Omega^{(\eta)})=1$, and
\[
\int\mathrm{d}\mu^{(N)}(F)\,\|F\|_{-3}^2
=\int\mathrm{d}\nu^{(N)}(F^{\mathrm{u}})\,\left\|\sharp_\eta F^{\mathrm{u}}\right\|_{-3}^2
\le\mathsf{Z}_6<\frac{1}{16},
\]
where $\|F\|_{-3}\in[0,\infty]$ is the norm of $H_{-3}$, equal to $\infty$ on $H_{-4}\setminus H_{-3}$.
The closed balls $\mathcal{B}_\varrho=\left\{F\in H_{-4}:\|F\|_{-3}\le\varrho\right\}$ of $H_{-3}$, $\varrho>0$, are compact in $H_{-4}$ \cite[Theorem~7.12\,(2)]{NoncompactQED},
and Chebyshev's inequality gives
\[
\mu^{(N)}\left(H_{-4}\setminus\mathcal{B}_\varrho\right)
\le\varrho^{-2}\cdot\int\mathrm{d}\mu^{(N)}(F)\,\|F\|_{-3}^2
<\frac{1}{16\varrho^2}
\qquad\left(N\in\mathcal{N},\ \varrho>0\right).
\]
Hence the family $(\mu^{(N)})_{N\in\mathcal{N}}$ is tight.

\emph{Test cochains.}
Again let $N\in\mathcal{N}$ and $\eta=L^{-N}$.
For a real $f\in\mathcal{S}$ the test cochain $f_{[\eta]}$ on $\Z^3$ at the spacing $\eta$ (Definition~\ref{def:test-cochains}) is the real two-cochain
whose value at a plaquette $\sigma$ is $\eta^{-\frac32}$ times the integral of the component of $f$ in the two directions of $\sigma$
over the averaging cube $Q_\eta(\sigma)=\eta\cdot\left(y_\sigma+\left[-\frac12,\frac12\right)^3\right)$, where $y_\sigma$ is the centre of $\sigma$ (Definition~\ref{def:averaging-cubes}).
If $f$ has compact support, only finitely many of these cubes meet $\supp f$, and $f_{[\eta]}$ has finite support.
By \eqref{eq:continuum-step-pairing},
\[
(\sharp_\eta F^{\mathrm{u}})(f)=\langle f_{[\eta]},F^{\mathrm{u}}\rangle\qquad\left(F^{\mathrm{u}}\in\Omega^{(\eta)}\right);
\]
thus $F(f)=\langle f_{[\eta]},F^{\mathrm{u}}\rangle$ for $F=\sharp_\eta F^{\mathrm{u}}$ and $\nu^{(N)}$-almost every $F^{\mathrm{u}}$.
The Cauchy--Schwarz inequality on each averaging cube, together with the fact that the averaging cubes of the plaquettes with the same two directions tile $\R^3$,
gives $\|f_{[\eta]}\|\le\|f\|_2$ (Lemma~\ref{lem:test-norms}(a)).

Consider next an element $g\in W$, with linear part $\vec g\in\Ocub$,
and let $g_\eta\colon\R^3\to\R^3$, $y\mapsto\eta\cdot g(y/\eta)$, be its isometry at the spacing $\eta$ (Definition~\ref{def:isometry-actions}).
It acts on two-forms by $(g_\eta\cdot f)(y)=(\wedge^2\vec g)\,f\left(g_\eta^{-1}y\right)$ \cite[Section~8.1.2]{NoncompactQED}.
The map $g$ sends every plaquette $\sigma$ to a plaquette $g\sigma$ and the centre $y_\sigma$ to $y_{g\sigma}$, and $\vec g$ preserves the cube $\left[-\frac12,\frac12\right]^3$;
hence $g_\eta$ maps $Q_\eta(\sigma)$ onto $Q_\eta(g\sigma)$ up to a null set.
Moreover, the component of $g_\eta\cdot f$ in the directions of $g\sigma$, at a point $y$, is $\pm$ the component of $f$ in the directions of $\sigma$ at $g_\eta^{-1}y$,
with the orientation sign of the pull-back $(g^{-1})^{*}$ at $g\sigma$ (Definition~\ref{def:cochains}).
A change of variables in the integrals over the averaging cubes therefore gives the identity
$(g_\eta\cdot f)_{[\eta]}=(g^{-1})^{*}f_{[\eta]}$ of Lemma~\ref{lem:test-equivariance}(a).
The pull-back $g^{*}$ is orthogonal, and $(g^{-1})^{*}=(g^{*})^{-1}$ (Definition~\ref{def:cochains});
so the transpose of $(g^{-1})^{*}$ is $g^{*}$, and for $f$ with compact support
\begin{equation}\label{eq:averaged-transposition}
\langle(g_\eta\cdot f)_{[\eta]},F^{\mathrm{u}}\rangle=\langle f_{[\eta]},g^{*}F^{\mathrm{u}}\rangle\qquad\left(F^{\mathrm{u}}\in\Omega\right),
\end{equation}
all sums being finite.

\emph{Gaussian domination of the limit.}
From now on $(N_i)$ and $\mu$ are as in the statement, and $\eta_i:=L^{-N_i}$.
For every $f\in\mathcal{S}$ the map $F\mapsto F(f)$ is a continuous linear functional on $H_{-4}$.
Let $f\in C^\infty_{\mathrm{c}}(\R^3;\Lambda^2)$ be real and $\varepsilon>0$.
The function $F\mapsto\min\left(\exp F(f),\varepsilon^{-1}\right)$ is bounded and continuous on $H_{-4}$,
and $\nu^{(N_i)}$ satisfies \eqref{cond:gaussian-domination} with the finitely supported cochain $f_{[\eta_i]}$.
With Lemma~\ref{lem:test-norms}(a) it follows that
\[
\begin{aligned}
\int\mathrm{d}\mu(F)\,\min\left(\exp F(f),\varepsilon^{-1}\right)
&=\lim_{i\to\infty}\int\mathrm{d}\mu^{(N_i)}(F)\,\min\left(\exp F(f),\varepsilon^{-1}\right)\\
&=\lim_{i\to\infty}\int\mathrm{d}\nu^{(N_i)}(F^{\mathrm{u}})\,\min\left(\exp\langle f_{[\eta_i]},F^{\mathrm{u}}\rangle,\varepsilon^{-1}\right)\\
&\le\sup_{i\ge1}\exp\left(\frac12\|f_{[\eta_i]}\|^2\right)
\le\exp\left(\frac12\|f\|_2^2\right).
\end{aligned}
\]
Letting $\varepsilon\to0$, the monotone convergence theorem gives $\int\mathrm{d}\mu(F)\,\exp F(f)\le\exp\left(\frac12\|f\|_2^2\right)$.
For a real $f\in\mathcal{S}$ choose real $f_j\in C^\infty_{\mathrm{c}}(\R^3;\Lambda^2)$ with $f_j\to f$ in $\mathcal{S}$.
Then $F(f_j)\to F(f)$ for every $F\in H_{-4}$ and $\|f_j\|_2\to\|f\|_2$, and Fatou's lemma yields
\[
\int\mathrm{d}\mu(F)\,\exp F(f)
\le\liminf_{j\to\infty}\int\mathrm{d}\mu(F)\,\exp F(f_j)
\le\exp\left(\frac12\|f\|_2^2\right).
\]
So $\mu$ satisfies \eqref{cond:gaussian-domination}.
This bound is the hypothesis of \cite[Proposition~7.10]{NoncompactQED} for the probability space $(H_{-4},\mu)$,
the real vector space of the real $f\in\mathcal{S}$ with the seminorm $\|f\|_2$, and the linear map $f\mapsto F(f)$, which is written $X$ there.
Assertion~(2) of that proposition gives
\begin{equation}\label{eq:averaged-second-moment}
\int\mathrm{d}\mu(F)\,F(f)^2\le\|f\|_2^2\qquad(f\in\mathcal{S}\ \text{real}).
\end{equation}
The same bound holds with $\mu^{(N)}$, $N\in\mathcal{N}$, in place of $\mu$.
Indeed, by Proposition~\ref{prop:volume-limits}(a) and \eqref{eq:continuum-step-pairing},
\[
\int\mathrm{d}\mu^{(N)}(F)\,\exp F(f)=\int\mathrm{d}\nu^{(N)}(F^{\mathrm{u}})\,\exp\langle f_{[\eta]},F^{\mathrm{u}}\rangle\le\exp\left(\frac12\|f\|_2^2\right)
\qquad(f\in\mathcal{S}\ \text{real},\ \eta=L^{-N}),
\]
which is the hypothesis of the same proposition for $(H_{-4},\mu^{(N)})$.

\emph{Invariance.}
Take $N\in\mathcal{N}$, $\eta=L^{-N}$, $g\in W$ and a real $f\in C^\infty_{\mathrm{c}}(\R^3;\Lambda^2)$;
then $g_\eta\cdot f$ is again real with compact support.
By \eqref{eq:continuum-step-pairing}, by \eqref{eq:averaged-transposition} and by the invariance of $\nu^{(N)}$ under $F^{\mathrm{u}}\mapsto g^{*}F^{\mathrm{u}}$,
\begin{equation}\label{eq:averaged-lattice-invariance}
\int\mathrm{d}\mu^{(N)}(F)\,\exp\left(iF(g_\eta\cdot f)\right)
=\int\mathrm{d}\nu^{(N)}(F^{\mathrm{u}})\,\exp\left(i\langle f_{[\eta]},g^{*}F^{\mathrm{u}}\rangle\right)
=\int\mathrm{d}\mu^{(N)}(F)\,\exp\left(iF(f)\right).
\end{equation}
An element $g\in\Ocub$ of $W$ is linear, so $g_\eta=g$ for every $\eta$.
For such $g$ the identity \eqref{eq:averaged-lattice-invariance} holds with the same test $g\cdot f$ for every $N=N_i$;
the integrands are bounded continuous functions on $H_{-4}$, and letting $i\to\infty$ gives
\[
\int\mathrm{d}\mu(F)\,\exp\left(iF(g\cdot f)\right)=\int\mathrm{d}\mu(F)\,\exp\left(iF(f)\right)\qquad(g\in\Ocub).
\]
Next let $b\in\R^3$, and choose $b_i\in\eta_i\Z^3$ with $|b-b_i|_\infty\le\eta_i$.
The translation of $\Z^3$ by $b_i/\eta_i$ belongs to $W$, and its isometry at the spacing $\eta_i$ is the translation $\tau_{b_i}$ of $\R^3$;
so \eqref{eq:averaged-lattice-invariance} gives
$\int\mathrm{d}\mu^{(N_i)}(F)\,\exp\left(iF(\tau_{b_i}\cdot f)\right)=\int\mathrm{d}\mu^{(N_i)}(F)\,\exp\left(iF(f)\right)$.
Since $|\exp(i\alpha)-\exp(i\alpha')|\le|\alpha-\alpha'|$ for real $\alpha,\alpha'$,
the Cauchy--Schwarz inequality and \eqref{eq:averaged-second-moment} for $\mu^{(N_i)}$ give
\[
\begin{aligned}
&\left|\int\mathrm{d}\mu^{(N_i)}(F)\,\left(\exp\left(iF(\tau_b\cdot f)\right)-\exp\left(iF(\tau_{b_i}\cdot f)\right)\right)\right|\\
&\qquad\le\left(\int\mathrm{d}\mu^{(N_i)}(F)\,F\left(\tau_b\cdot f-\tau_{b_i}\cdot f\right)^2\right)^{\frac12}
\le\left\|\tau_b\cdot f-\tau_{b_i}\cdot f\right\|_2 .
\end{aligned}
\]
The right side tends to $0$ as $i\to\infty$, because $b_i\to b$ and the translations act continuously on $L^2(\R^3;\Lambda^2)$.
Letting $i\to\infty$ yields $\int\mathrm{d}\mu(F)\,\exp\left(iF(\tau_b\cdot f)\right)=\int\mathrm{d}\mu(F)\,\exp\left(iF(f)\right)$.

The translations and $\Ocub$ generate the group $\R^3\rtimes\Ocub$ of the Euclidean maps whose linear part lies in $\Ocub$,
and this group maps the real elements of $C^\infty_{\mathrm{c}}(\R^3;\Lambda^2)$ onto themselves.
Hence the characteristic functional $f\mapsto\int\mathrm{d}\mu(F)\,\exp\left(iF(f)\right)$ is invariant under $f\mapsto\mathfrak{g}\cdot f$
for every $\mathfrak{g}\in\R^3\rtimes\Ocub$ and every real $f\in C^\infty_{\mathrm{c}}(\R^3;\Lambda^2)$.
It is invariant for every real $f\in\mathcal{S}$ as well:
choose real $f_j\in C^\infty_{\mathrm{c}}(\R^3;\Lambda^2)$ with $f_j\to f$ in $\mathcal{S}$;
then $\mathfrak{g}\cdot f_j\to\mathfrak{g}\cdot f$ in $\mathcal{S}$, the map $f\mapsto\mathfrak{g}\cdot f$ being a continuous linear bijection of $\mathcal{S}$ \cite[Section~8.1.2]{NoncompactQED},
so $F(f_j)\to F(f)$ and $F(\mathfrak{g}\cdot f_j)\to F(\mathfrak{g}\cdot f)$ for every $F\in H_{-4}$, and the dominated convergence theorem applies.
The image law $\mathfrak{g}_{*}\mu$ \cite[Section~8.1.2]{NoncompactQED} has the characteristic functional $f\mapsto\int\mathrm{d}\mu(F)\,\exp\left(iF(\mathfrak{g}^{-1}\cdot f)\right)$,
and a law on $\mathcal{S}'$ is determined by its characteristic functional on $\mathcal{S}$ \cite[Lemma~9.3\,(b)]{NoncompactQED}.
Therefore $\mathfrak{g}_{*}\mu=\mu$ for every $\mathfrak{g}\in\R^3\rtimes\Ocub$.
The invariance of $\nu^{(N)}$ under $F^{\mathrm{u}}\mapsto-F^{\mathrm{u}}$ gives, for real $f\in C^\infty_{\mathrm{c}}(\R^3;\Lambda^2)$,
\[
\int\mathrm{d}\mu^{(N)}(F)\,\exp\left(-iF(f)\right)
=\int\mathrm{d}\nu^{(N)}(F^{\mathrm{u}})\,\exp\left(-i\langle f_{[\eta]},F^{\mathrm{u}}\rangle\right)
=\int\mathrm{d}\mu^{(N)}(F)\,\exp\left(iF(f)\right).
\]
Letting $N=N_i$ with $i\to\infty$, extending the resulting identity to all real $f\in\mathcal{S}$ by dominated convergence,
and applying \cite[Lemma~9.3\,(b)]{NoncompactQED} once more, one finds that $\mu$ is invariant under $F\mapsto-F$.

\emph{Closedness.}
Let $\xi\in C^\infty_{\mathrm{c}}(\R^3)$ be real, let $N\in\mathcal{N}$ and $\eta=L^{-N}$,
and let $d^{*}\xi$ and $u^{(\eta)}$ be the two-forms of Lemma~\ref{lem:test-bianchi},
formed with the derivatives $\partial_j\xi$ and with the centred difference quotients $\partial^{\eta}_j\xi$ at the spacing $\eta$, respectively;
both are real and have compact support.
Every $F^{\mathrm{u}}\in\Omega^{(\eta)}$ is a closed two-cochain on $\Z^3$
for which $\left|F^{\mathrm{u}}(\sigma)\right|/\left(1+|x(\sigma)|_\infty\right)$ is bounded on the plaquettes $\sigma$ (Definition~\ref{def:step-map});
so \eqref{eq:continuum-step-pairing} and Lemma~\ref{lem:test-bianchi}(a) give
\[
(\sharp_\eta F^{\mathrm{u}})(u^{(\eta)})=\langle u^{(\eta)}_{[\eta]},F^{\mathrm{u}}\rangle=0\qquad\left(F^{\mathrm{u}}\in\Omega^{(\eta)}\right).
\]
As $\nu^{(N)}(\Omega^{(\eta)})=1$, one has $F(u^{(\eta)})=0$ for $\mu^{(N)}$-almost every $F$.
By \eqref{eq:averaged-second-moment} for $\mu^{(N)}$ and by Lemma~\ref{lem:test-bianchi}(b),
\[
\begin{aligned}
\int\mathrm{d}\mu^{(N)}(F)\,F(d^{*}\xi)^2
&=\int\mathrm{d}\mu^{(N)}(F)\,F\left(d^{*}\xi-u^{(\eta)}\right)^2
\le\left\|d^{*}\xi-u^{(\eta)}\right\|_2^2\\
&\le C_\xi\cdot\eta^4,
\qquad
C_\xi:=\frac{1}{576}\left(\sum_{j=1}^3\left\|\partial_j^3\xi\right\|_2\right)^2.
\end{aligned}
\]
This bound holds for every $N\in\mathcal{N}$, in particular for $N=N_i$ with $\eta=\eta_i$.
For $\varepsilon>0$ the function $F\mapsto\min\left(F(d^{*}\xi)^2,\varepsilon^{-1}\right)$ is bounded and continuous on $H_{-4}$; hence
\[
\int\mathrm{d}\mu(F)\,\min\left(F(d^{*}\xi)^2,\varepsilon^{-1}\right)
=\lim_{i\to\infty}\int\mathrm{d}\mu^{(N_i)}(F)\,\min\left(F(d^{*}\xi)^2,\varepsilon^{-1}\right)
\le\lim_{i\to\infty}C_\xi\cdot\eta_i^4=0,
\]
and $F(d^{*}\xi)=0$ for $\mu$-almost every $F$.

For $\xi$ in the space $\mathcal{S}(\R^3)$ of the real Schwartz functions on $\R^3$,
the components $(23)$, $(13)$, $(12)$ of the two-form $f=d^{*}\xi$ are $-\partial_1\xi$, $\partial_2\xi$, $-\partial_3\xi$,
so its dual vector field $\left(f_{23},-f_{13},f_{12}\right)$ is $-\nabla\xi$.
By the description of the closed distributions in \cite[Section~8.1.6]{NoncompactQED},
an element $F$ of $\mathcal{S}'$ is therefore closed if and only if $F(d^{*}\xi)=0$ for every $\xi\in\mathcal{S}(\R^3)$.
The real elements of $C^\infty_{\mathrm{c}}(\R^3)$ form a dense subset of $\mathcal{S}(\R^3)$, which is separable and metrizable;
so some countable set of real elements of $C^\infty_{\mathrm{c}}(\R^3)$ is dense in $\mathcal{S}(\R^3)$.
If $F(d^{*}\xi)=0$ for every $\xi$ in this countable set, then $F(d^{*}\xi)=0$ for every $\xi\in\mathcal{S}(\R^3)$,
because $\xi\mapsto d^{*}\xi$ is continuous from $\mathcal{S}(\R^3)$ to $\mathcal{S}$ and $F$ is continuous on $\mathcal{S}$.
Thus $\{dF=0\}$ is the intersection of countably many sets of $\mu$-measure one, and $\mu(\{dF=0\})=1$.

\emph{Mirrors.}
Fix $\ell\in\mathcal{L}$, an integer $N_0\ge0$ and $c\in L^{-N_0}\Z$,
and let $r_{\ell,c}$ be the reflection of $\R^3$ in the plane $\{\ell=c\}$ (Definition~\ref{def:isometry-actions}).
Let $k\ge1$, let $f_1,\dots,f_k\in C^\infty_{\mathrm{c}}(\R^3;\Lambda^2)$ be real with supports in the open half space $\{\ell>c\}$, and let $z_1,\dots,z_k\in\C$.
The supports are compact, so there is $\delta>0$ with $\supp f_a\subset\{\ell\ge c+\delta\}$ for $a=1,\dots,k$.
Choose $i_0$ such that $N_i\ge N_0$ and $2\eta_i<\delta$ for every $i\ge i_0$, and let $i\ge i_0$.
Then $c/\eta_i=c\cdot L^{N_i}$ is an integer, so $(\ell,c/\eta_i)$ is a mirror (Definition~\ref{def:mirrors});
put $\theta:=\theta_{\ell,c/\eta_i}$.

Suppose that $(f_a)_{[\eta_i]}(\sigma)\neq0$ for some $a\in\{1,\dots,k\}$ and some plaquette $\sigma$.
Then the averaging cube $Q_{\eta_i}(\sigma)$ meets $\supp f_a$, so it contains a point $y$ with $\ell(y)\ge c+\delta$.
Let $y'$ be $\eta_i$ times a corner of $\sigma$.
The point $y'$ differs from $\eta_i\cdot y_\sigma$ by $\frac{\eta_i}{2}$ in two coordinates and by $0$ in the third,
while $y$ differs from $\eta_i\cdot y_\sigma$ by at most $\frac{\eta_i}{2}$ in every coordinate; so $|y'-y|_\infty\le\eta_i$.
The vector $n_\ell$ has at most two non-zero coordinates, each equal to $\pm1$; hence $|\ell(y')-\ell(y)|\le2\eta_i$, and
\[
\ell(y')\ge c+\delta-2\eta_i>c .
\]
Every corner $y'/\eta_i$ of $\sigma$ therefore satisfies $\ell>c/\eta_i$, that is, $\sigma\in\Sigma^{+}(\ell,c/\eta_i)$.
So the finitely supported cochains $(f_a)_{[\eta_i]}$ are supported in $\Sigma^{+}(\ell,c/\eta_i)$, and the function
\[
\Phi(F^{\mathrm{u}}):=\sum_{a=1}^k z_a\cdot\exp\left(i\langle(f_a)_{[\eta_i]},F^{\mathrm{u}}\rangle\right)\qquad\left(F^{\mathrm{u}}\in\Omega\right)
\]
is a bounded continuous function of finitely many $F^{\mathrm{u}}_\sigma$, $\sigma\in\Sigma^{+}(\ell,c/\eta_i)$: it belongs to $\mathcal{O}^{+}(\ell,c/\eta_i)$.
The pull-back $\theta^{*}$ is an orthogonal involution (Definition~\ref{def:cochains}), hence symmetric,
and $\theta^{*}(f_a)_{[\eta_i]}=(r_{\ell,c}\cdot f_a)_{[\eta_i]}$ by Lemma~\ref{lem:test-equivariance}(a), as $c/\eta_i\in\Adm(\ell)$.
Consequently
\[
\overline{\Phi(\theta^{*}F^{\mathrm{u}})}\cdot\Phi(F^{\mathrm{u}})
=\sum_{a,a'=1}^k\overline{z_a}\cdot z_{a'}\cdot\exp\left(i\left(\langle(f_{a'})_{[\eta_i]},F^{\mathrm{u}}\rangle-\langle(r_{\ell,c}\cdot f_a)_{[\eta_i]},F^{\mathrm{u}}\rangle\right)\right).
\]
The tests $r_{\ell,c}\cdot f_a$ are real and have compact support.
The reflection positivity of $\nu^{(N_i)}$ at the mirror $(\ell,c/\eta_i)$ and \eqref{eq:continuum-step-pairing} give
\[
\begin{aligned}
&\sum_{a,a'=1}^k\overline{z_a}\cdot z_{a'}\cdot\int\mathrm{d}\mu^{(N_i)}(F)\,\exp\left(i\left(F(f_{a'})-F(r_{\ell,c}\cdot f_a)\right)\right)\\
&\qquad=\int\mathrm{d}\nu^{(N_i)}(F^{\mathrm{u}})\,\overline{\Phi(\theta^{*}F^{\mathrm{u}})}\cdot\Phi(F^{\mathrm{u}})
\ge0
\qquad(i\ge i_0).
\end{aligned}
\]
The integrands on the left are bounded continuous functions on $H_{-4}$; letting $i\to\infty$,
\begin{equation}\label{eq:averaged-plane-positivity}
\sum_{a,a'=1}^k\overline{z_a}\cdot z_{a'}\cdot\int\mathrm{d}\mu(F)\,\exp\left(i\left(F(f_{a'})-F(r_{\ell,c}\cdot f_a)\right)\right)\ge0 .
\end{equation}
As $k$, the real tests $f_a$ with supports in $\{\ell>c\}$ and the coefficients $z_a$ are arbitrary,
\eqref{eq:averaged-plane-positivity} states that $\mu$ is reflection positive at the plane $\{\ell=c\}$
in the sense of Definition~\ref{def:laws} and of \cite[Section~8.1.3]{NoncompactQED}.
Moreover $\theta\in W$ is the map $x\mapsto\vec\theta x+v$ with $\vec\theta\in\Ocub$ and $v\in\Z^3$ (Definition~\ref{def:mirrors}),
and by the linearity of $\ell$ its isometry $y\mapsto\eta_i\cdot\theta(y/\eta_i)=\vec\theta y+\eta_i\cdot v$ at the spacing $\eta_i$ is $r_{\ell,c}$.
Thus $r_{\ell,c}$ belongs to $\R^3\rtimes\Ocub$, and $\mu$ is invariant under it.

\emph{Arbitrary positions.}
Now let $\ell\in\mathcal{L}$ and $c'\in\R$ be arbitrary.
Choose $c\in L^{-N_0}\Z$ for some integer $N_0\ge0$, and $b\in\R^3$ with $\ell(b)=c'-c$.
The translation $\tau_b$ maps the plane $\{\ell=c\}$ onto $\{\ell=c'\}$ and the half space $\{\ell>c\}$ onto $\{\ell>c'\}$, and
\[
r_{\ell,c'}=\tau_b\circ r_{\ell,c}\circ\tau_{-b}.
\]
In particular $r_{\ell,c'}$ belongs to $\R^3\rtimes\Ocub$, and $\mu$ is invariant under it.
Let $f'_1,\dots,f'_k\in C^\infty_{\mathrm{c}}(\R^3;\Lambda^2)$ be real with supports in $\{\ell>c'\}$, let $z_1,\dots,z_k\in\C$, and put $f_a:=\tau_{-b}\cdot f'_a$.
The $f_a$ are real, with compact supports in $\{\ell>c\}$, and
\[
F(f'_{a'})-F(r_{\ell,c'}\cdot f'_a)=F\left(\tau_b\cdot\left(f_{a'}-r_{\ell,c}\cdot f_a\right)\right)\qquad\left(F\in H_{-4}\right).
\]
As $\mu$ is invariant under the translations, $\int\mathrm{d}\mu(F)\,\exp\left(iF(\tau_b\cdot f)\right)=\int\mathrm{d}\mu(F)\,\exp\left(iF(f)\right)$ for every real $f\in\mathcal{S}$.
So the sum in \eqref{eq:averaged-plane-positivity} formed with $c'$, the tests $f'_a$ and the coefficients $z_a$ equals the sum formed with $c$, the tests $f_a$ and the same coefficients, which is non-negative.
Hence, for every $\ell\in\mathcal{L}$ and every $c'\in\R$, the law $\mu$ is invariant under the reflection in the plane $\{\ell=c'\}$ and reflection positive at this plane.
Since $\mu$ is also invariant under the translations of $\R^3$, it belongs to $\mathcal{C}_9$, which is~(b).

\emph{Conclusion.}
For the axis family $\ell=x_j$ and $c'=0$ the plane $\{\ell=c'\}$ is $\{x_j=0\}$, the reflection in it is $\theta_j$, and the half space $\{\ell>c'\}$ is $\{x_j>0\}$;
so $\mu$ satisfies \eqref{cond:reflection-invariance} and \eqref{cond:reflection-positivity} for $j=1,2,3$.
Together with \eqref{cond:gaussian-domination} and the invariance under the translations of $\R^3$ and under $\Ocub$, these conditions make up \eqref{cond:soft-hypothesis}.
With \eqref{cond:closedness} and the invariance under $F\mapsto-F$, the law $\mu$ is soft,
and the remaining assertions of~(a) are those of \cite[Lemma~21.36]{NoncompactQED}.
\end{proof}

\subsection{The harmonic-averaged law}\label{sec:harmonic-average}

Throughout this subsection $-1\le\rho\le1$, the weights $w_i$ are positive, $M$ is real and $q>0$.
For a torus $\mathbb{T}=(\Gamma,\chi)$ the cochain spaces of its period lattice are also written
$C^k(\mathbb{T}):=C^k(\Gamma)$ and $C^k(\mathbb{T};\Z):=C^k(\Gamma;\Z)$,
with the pairing $\langle\cdot,\cdot\rangle$ and the norm $\|\cdot\|$ of Definition~\ref{def:configurations};
in particular, the curvature $F$ of a configuration is a closed element of $C^2(\mathbb{T})$.
The functional of a torus integrates over the configurations $(A,z)$ and sums over all closed integer two-cochains $z$ (Definition~\ref{def:configurations}).
On a coordinate torus the sums of $z$ over the coordinate two-tori sort these cochains into sectors.

\begin{definition}[Flux sectors and the averaged law]\label{def:averaged-law}
Let $-1\le\rho\le1$, $w_i>0$, $M\in\R$, $q>0$, let $\mathbb{T}=C(n_1,n_2,n_3)$ with $n_\mu\ge4$ (Definition~\ref{def:tori}), let $\Gamma$ be its period lattice, and $|n|=n_1n_2n_3$.
For $\mu<\nu$ let $\Sigma_{\mu\nu}$ be the set of the $n_\mu n_\nu$ plaquettes $(x;\mu\nu)$ of $\mathbb{T}$ that lie in one coordinate plane parallel to $e_\mu$ and $e_\nu$;
it is a coordinate two-torus of $\mathbb{T}$.
Let $(A,z)$ represent a point of $\Conf(\Gamma)$, and $F=dA+(2\pi/q)z$.
Each link of a coordinate two-torus is a face of exactly two of its plaquettes, with opposite incidence numbers;
so the sum of $dA$ over $\Sigma_{\mu\nu}$ vanishes, and
\[
\sum_{\sigma\in\Sigma_{\mu\nu}}F_\sigma=\frac{2\pi}{q}\,k_{\mu\nu},
\qquad
k_{\mu\nu}:=\sum_{\sigma\in\Sigma_{\mu\nu}}z_\sigma\in\Z .
\]
Since $z$ is closed, $k_{\mu\nu}$ does not depend on the choice of the plane \cite[(10.2)]{NoncompactQED}
(where the coordinate two-tori are written $\Pi_{\mu\nu}(t)$ and the flux $w=(w_{12},w_{13},w_{23})$).
The vector $k=(k_{23},k_{13},k_{12})\in\Z^3$ is the \emph{flux} of the configuration;
it is a function of $F$.
\begin{enumerate}
\item[(a)]\itemmark{a}\label{it:averaged-sectors}
For $k'\in\Z^3$ the indicator of the \emph{flux sector} $k'$ is the bounded Borel function of finitely many $F_\sigma$
\[
\Phi_{k'}(F):=\prod_{\mu<\nu}\mathbf{1}\left[\,\sum_{\sigma\in\Sigma_{\mu\nu}}F_\sigma=\frac{2\pi}{q}\,k'_{\mu\nu}\right],
\]
and, for every integrand $G$ for which $Z_{\mathbb{T}}[G]$ is defined (Definition~\ref{def:functional}),
\[
Z^{(k')}_{\mathbb{T}}[G]:=Z_{\mathbb{T}}[\Phi_{k'}G].
\]
The number $\mathbb{E}_{\mathbb{T}}[\Phi_{k'}]=Z^{(k')}_{\mathbb{T}}[1]/Z_{\mathbb{T}}[1]$ is the weight of the sector $k'$.
The functions $\Phi_k$, $k\in\Z^3$, take the values $0$ and $1$, and their sum is $1$.
For a bosonic integrand $G$ one has $|Z^{(k)}_{\mathbb{T}}[G]|\le Z_{\mathbb{T}}[\Phi_k\,|G|]$ by \eqref{eq:setting-modulus-bound};
on non-negative Borel functions of $F$ the functional is the integral against the non-negative measure $\exp(-\|F\|^2/2)\cdot\det D^{\mathbb{T}}_A$ on $\Conf(\Gamma)$ (Proposition~\ref{prop:accretivity}), and therefore
\[
\sum_{k\in\Z^3}\left|Z^{(k)}_{\mathbb{T}}[G]\right|\le\sum_{k\in\Z^3}Z_{\mathbb{T}}[\Phi_k\,|G|]=Z_{\mathbb{T}}[\,|G|\,].
\]
\item[(b)]\itemmark{b}\label{it:averaged-law-def}
The number $Z^{(0)}_{\mathbb{T}}[1]$ is the sum, over the closed $z\in C^2(\Gamma;\Z)$ of flux $0$, of the flux terms
\[
\int_{[0,2\pi/q)^{\mathbb{T}^{(1)}}}\mathrm{d}A\,\exp\left(-\frac{1}{2}\left\|dA+\frac{2\pi}{q}\,z\right\|^2\right)\cdot\det D^{\mathbb{T}}_A\;\ge0 ,
\]
and the term of $z=0$ is positive (Theorem~\ref{thm:floor}\,(d) if $\rho^2\ne1$; Lemma~\ref{lem:flux-terms} and Theorem~\ref{thm:closed-interval} if $\rho=\pm1$).
Hence $0<Z^{(0)}_{\mathbb{T}}[1]\le Z_{\mathbb{T}}[1]<\infty$ (Theorem~\ref{thm:partition-positivity}).
The \emph{harmonic-averaged law} $\mu^{\mathrm{av}}_{\mathbb{T}}$ is the law of $\widetilde F$ under the normalised functional $Z^{(0)}_{\mathbb{T}}[\,\cdot\,]/Z^{(0)}_{\mathbb{T}}[1]$.
It is the law $\mu_{\mathbb{T}}$ conditioned on the flux $0$:
$\mu^{\mathrm{av}}_{\mathbb{T}}$ has the density $\Phi_0/\mathbb{E}_{\mathbb{T}}[\Phi_0]$ with respect to $\mu_{\mathbb{T}}$.
\item[(c)]\itemmark{c}\label{it:averaged-pressure}
For sides $n_\mu\ge4$ the pressure of the sector of flux zero is $f^{(0)}(n_1,n_2,n_3):=|n|^{-1}\log Z^{(0)}_{C(n_1,n_2,n_3)}[1]$, and $f^{(0)}(n):=f^{(0)}(n,n,n)$.
\item[(d)]\itemmark{d}\label{it:averaged-harmonic}
For $k\in\Z^3$ let $\mathfrak{h}_k\in C^2(\mathbb{T})$ be the two-cochain whose value at every plaquette $(x;\mu\nu)$ of $\mathbb{T}$ is $(2\pi/q)\,k_{\mu\nu}/(n_\mu n_\nu)$;
it is the cochain $h_w$ of \cite[Lemma~10.2]{NoncompactQED}, with $k_{\mu\nu}$ in the place of $w_{\mu\nu}$.
As $\mathbb{T}$ has $|n|$ plaquettes of each of the three kinds,
\[
\|\mathfrak{h}_k\|^2=\left(\frac{2\pi}{q}\right)^2\sum_{\mu<\nu}\frac{|n|}{(n_\mu n_\nu)^2}\,k_{\mu\nu}^2 ,
\]
which equals $(2\pi/q)^2\,|k|_2^2/n$ on the cubic torus $C(n)$.
\item[(e)]\itemmark{e}\label{it:averaged-covariance}
Let $g\in W$.
Then $g\mathbb{T}$ is a coordinate torus, and $g$ maps every coordinate two-torus of $\mathbb{T}$ onto one of $g\mathbb{T}$, preserving the orientations of all its plaquettes or reversing all of them.
Hence for every $k\in\Z^3$ there is a vector $k^g\in\Z^3$, obtained from $k$ by a permutation of the components and changes of their signs, such that $\Phi_k(g^*F)=\Phi_{k^g}(F)$ for the closed cochains $F$ of $g\mathbb{T}$
($\Phi_k$ formed on $\mathbb{T}$, $\Phi_{k^g}$ on $g\mathbb{T}$);
in particular $0^g=0$.
For a scalar element $G$ of $\mathcal{I}(\mathbb{T})$, that is, one with values in the multiples of the unit of $\Gr_{\mathbb{T}}$ (Definition~\ref{def:torus-integrands}),
and a word $g=\theta_{\Pi^{(1)}}\cdots\theta_{\Pi^{(r)}}$ in the reflections in mirrors,
let $G^g:=\Theta_{\Pi^{(1)}}\cdots\Theta_{\Pi^{(r)}}G$ be formed with the maps of Lemma~\ref{lem:reflection-covariance};
by $r$ applications of that lemma, $G^g\in\mathcal{I}(g\mathbb{T})$.
The antilinear maps between Grassmann algebras that enter that lemma send the unit to the unit,
so that $(\Theta_\Pi G)(A',z')=\overline{G(\theta_\Pi^{*}A',\theta_\Pi^{*}z')}$ for a scalar $G$.
As $g^{*}=\theta^{*}_{\Pi^{(r)}}\cdots\theta^{*}_{\Pi^{(1)}}$ (Definition~\ref{def:cochains}) and $\det\vec g=(-1)^r$, it follows that, for $(A',z')\in\mathfrak{C}(g\Gamma)$,
\[
G^g(A',z')=G(g^{*}A',g^{*}z')\quad(g\ \text{even}),
\qquad
G^g(A',z')=\overline{G(g^{*}A',g^{*}z')}\quad(g\ \text{odd}).
\]
So $G^g$ does not depend on the word, and $G\mapsto G^g$ is multiplicative.
If $G$ is a function of the curvature $F$ of $\mathbb{T}$, then $G^g$ is the function $F'\mapsto G(g^{*}F')$, resp.\ $F'\mapsto\overline{G(g^{*}F')}$, of the curvature $F'$ of $g\mathbb{T}$, since $g^{*}$ commutes with $d$.
For a bosonic element $G$ of $\mathcal{A}(\R^3)$ that embeds in $\mathbb{T}$, the element $G^g$ of \eqref{eq:setting-group-action}, regarded as a function on $\Conf(g\Gamma)$, is given by the same formula, the periodic extension of $g^{*}F'$ being $g^{*}\widetilde{F'}$.
As $\Phi_k$ is real, $(\Phi_kG)^g=\Phi_{k^g}G^g$.
By Proposition~\ref{prop:symmetries}\eqref{it:functional-covariance}, applied to $\Phi_kG$, for $G\in\mathcal{A}(\R^3)$ that embeds in $\mathbb{T}$, and by Lemma~\ref{lem:reflection-covariance}, applied to $\Phi_kG$ and to the reflections of a word for $g$, for scalar $G\in\mathcal{I}(\mathbb{T})$,
\[
Z^{(k^g)}_{g\mathbb{T}}[G^g]=Z^{(k)}_{\mathbb{T}}[G]\quad(g\ \text{even}),
\qquad
Z^{(k^g)}_{g\mathbb{T}}[G^g]=\overline{Z^{(k)}_{\mathbb{T}}[G]}\quad(g\ \text{odd}).
\]
Since $\tau_v\mathbb{T}=\mathbb{T}$ for $v\in\Z^3$, and $gC(n)=C(n)$ for $g\in W$, the law $\mu^{\mathrm{av}}_{\mathbb{T}}$ is invariant under $\Z^3$, and $\mu^{\mathrm{av}}_{C(n)}$ is invariant under $W$.
Finally, $\mu^{\mathrm{av}}_{\mathbb{T}}$ is invariant under $F\mapsto-F$:
$\Phi_0(-F)=\Phi_0(F)$, $\det D^{\mathbb{T}}_{-A}=\det D^{\mathbb{T}}_A$ (Proposition~\ref{prop:symmetries}\eqref{it:charge-conjugation}), and $(A,z)\mapsto(-A,-z)$ preserves the measure of Definition~\ref{def:configurations}.
\end{enumerate}
\end{definition}

For a coordinate torus with time $e_1$ the word formula expresses $Z_{\mathbb{T}}$ through operators that are diagonal in the flux through the slices.
The two other components of the flux are reached by a Bloch angle.

\begin{lemma}[Bloch representation]\label{lem:bloch}
Let $-1\le\rho\le1$, $w_i>0$, $M\in\R$, $q>0$, and $\mathbb{T}=C(n_1,n_2,n_3)$ with $n_\mu\ge4$.
Take $e_1$ as time:
$\Lambda_{\mathrm{s}}$ is the slice torus, with the links $\Lambda^{(1)}_{\mathrm{s}}$ and the plaquettes $\Lambda^{(2)}_{\mathrm{s}}$,
$F^{\mathrm{el}}_j$ and $F^{\mathrm{mag}}_j$ are the electric and the magnetic part of $F$ (Definition~\ref{def:slices}),
and $\mathcal{H}^{\mathrm{l}}$, $\mathcal{H}^{\mathrm{s}}$, $\mathfrak{M}$, $K_\beta$ and $P_{\mathrm{G}}$ are those of Definition~\ref{def:slice-operators}, formed with a plaquette $\sigma_0$ of $\Lambda_{\mathrm{s}}$.
On the summand of flux index $k\in\Z$ of $\mathcal{H}^{\mathrm{l}}$ and of $\mathcal{H}^{\mathrm{s}}$, $F^{\mathrm{mag}}$ denotes the commuting multiplication operators $F^{\mathrm{mag}}(\sigma'):=(d_{\mathrm{s}}A)(\sigma')+(2\pi/q)k\,1_{\sigma_0}(\sigma')$, $\sigma'\in\Lambda^{(2)}_{\mathrm{s}}$;
for a bounded Borel function $u$ on $\R^{\Lambda^{(2)}_{\mathrm{s}}}$, $u(F^{\mathrm{mag}})$ is the multiplication by $u\bigl(d_{\mathrm{s}}A+(2\pi/q)k\,1_{\sigma_0}\bigr)$ on that summand, as in the word formula of Proposition~\ref{prop:word-formula}.
Fix a site $y^0=(y^0_2,y^0_3)$ of $\Lambda_{\mathrm{s}}$ and the two closed lines of slice links through it,
\[
S_2:=\{((y_2,y^0_3),2):y_2\in\Z/n_2\},\qquad S_3:=\{((y^0_2,y_3),3):y_3\in\Z/n_3\}.
\]
For $\vartheta=(\vartheta_2,\vartheta_3)\in[0,1]^2$ let
\[
\beta_\vartheta:=q\,(\vartheta_2\mathbf{1}_{S_2}+\vartheta_3\mathbf{1}_{S_3})\in\R^{\Lambda^{(1)}_{\mathrm{s}}},
\]
and for $k\in\Z$ let $P_k$ be the orthogonal projection of $\mathcal{H}^{\mathrm{l}}$, and of $\mathcal{H}^{\mathrm{s}}$, onto the summand of flux index $k$.
For $j\in\Z/n_1$ let $u_j$ be a bounded Borel function on $\R^{\Lambda^{(2)}_{\mathrm{s}}}$, and $v_j$ a Borel function on $\R^{\Lambda^{(1)}_{\mathrm{s}}}$ with $|v_j(y)|\le C\exp\bigl((1/2-\varepsilon_v)|y|^2\bigr)$ for some $C<\infty$ and $\varepsilon_v>0$.
Put
\[
G:=\prod_{j\in\Z/n_1}u_j(F^{\mathrm{mag}}_j)\cdot v_j(F^{\mathrm{el}}_j),
\]
and let $K^\vartheta_j$ be the operator on $\mathcal{H}^{\mathrm{l}}$ whose kernel is the periodisation over $\Gamma_{\mathrm{s}}$ of
\[
(A',A)\mapsto v_j(A'-A)\cdot\exp\bigl(i\langle\beta_\vartheta,A'-A\rangle\bigr)\cdot\prod_{b\in\Lambda^{(1)}_{\mathrm{s}}}g_1\bigl(A'(b)-A(b)\bigr),
\]
that is, the operator of the step $j$ of the word formula of Proposition~\ref{prop:word-formula}, formed with $v_j(y)\cdot\exp(i\langle\beta_\vartheta,y\rangle)$ in the place of $v_j(y)$.
\begin{enumerate}
\item[(a)]\itemmark{a}\label{it:bloch-fourier}
For every $k_{23}\in\Z$ and every $\vartheta\in[0,1]^2$ the series on the left converges absolutely, and
\begin{multline}\label{eq:averaged-fourier}
\sum_{k_{12},k_{13}\in\Z}\exp\bigl(2\pi i(\vartheta_2k_{12}+\vartheta_3k_{13})\bigr)\cdot Z^{(k_{23},k_{13},k_{12})}_{\mathbb{T}}[G]\\
=(2\pi)^{|n|}\cdot\Tr_{\mathcal{H}^{\mathrm{l}}}\left(P_{k_{23}}\prod_{j=n_1-1,\dots,0}\left((2\pi/q)^{|\Lambda_{\mathrm{s}}|}\,K^\vartheta_j\,P_{\mathrm{G}}\,\mathfrak{M}\,u_j(F^{\mathrm{mag}})\,\mathfrak{M}^{*}\right)\right).
\end{multline}
\item[(b)]\itemmark{b}\label{it:bloch-zero-flux}
One has
\begin{equation}\label{eq:averaged-bloch}
\begin{aligned}
Z^{(0)}_{\mathbb{T}}[G]&=(2\pi)^{|n|}\int_{[0,1]^2}\mathrm{d}\vartheta\,\Tr\left(P_0\prod_{j=n_1-1,\dots,0}\left((2\pi/q)^{|\Lambda_{\mathrm{s}}|}\,K^\vartheta_j\,P_{\mathrm{G}}\,\mathfrak{M}\,u_j(F^{\mathrm{mag}})\,\mathfrak{M}^{*}\right)\right),\\
Z^{(0)}_{\mathbb{T}}[1]&=(2\pi)^{|n|}\int_{[0,1]^2}\mathrm{d}\vartheta\,\Tr_{\mathcal{H}^{\mathrm{s}}}\left(P_0\,\bigl(\mathfrak{M}^{*}K_{\beta_\vartheta}\mathfrak{M}\bigr)^{n_1}\right),
\end{aligned}
\end{equation}
where $\mathfrak{M}^{*}K_{\beta_\vartheta}\mathfrak{M}$ is non-negative and of trace class.
\item[(c)]\itemmark{c}\label{it:bloch-directions}
Let $\mu\in\{2,3\}$, let $g_\mu$ be the cyclic permutation matrix with $g_\mu e_1=e_\mu$ of Definition~\ref{def:other-directions}, and let $\mathbb{T}':=g_\mu^{-1}\mathbb{T}$, a coordinate torus whose side in the direction $e_1$ is $n_\mu$.
Take $e_\mu$ as time, as in that definition:
the electric and the magnetic parts of $F$ are $F^{\mathrm{el}}_j(g_\mu^{*}F)$ and $F^{\mathrm{mag}}_j(g_\mu^{*}F)$, $j\in\Z/n_\mu$, formed on $\mathbb{T}'$,
the integrand $G$ is formed with them,
and the slice torus, the lines $S_2$, $S_3$, the Hilbert spaces and the operators are those of $\mathbb{T}'$ with time $e_1$.
Then (b) holds with $n_\mu$ in the place of $n_1$.
Moreover, for $\mu=2$ and every $k_{13}\in\Z$ the series
\[
\sum_{k_{23},k_{12}\in\Z}\exp\bigl(2\pi i(\vartheta_2k_{23}-\vartheta_3k_{12})\bigr)\cdot Z^{(k_{23},k_{13},k_{12})}_{\mathbb{T}}[G]
\]
converges absolutely, and its sum is the right side of \eqref{eq:averaged-fourier}, formed with time $e_2$ and with $P_{-k_{13}}$ in the place of $P_{k_{23}}$;
for $\mu=3$ and every $k_{12}\in\Z$ the series
\[
\sum_{k_{13},k_{23}\in\Z}\exp\bigl(-2\pi i(\vartheta_2k_{13}+\vartheta_3k_{23})\bigr)\cdot Z^{(k_{23},k_{13},k_{12})}_{\mathbb{T}}[G]
\]
converges absolutely, and its sum is the right side of \eqref{eq:averaged-fourier}, formed with time $e_3$ and with $P_{k_{12}}$ in the place of $P_{k_{23}}$.
\end{enumerate}
\end{lemma}

\begin{proof}
When the fermion weight is the determinant of the Wilson operator, the trace formula with Bloch angles is \cite[Proposition~11.5\,(iii)]{NoncompactQED};
here $\det D^{\mathbb{T}}_A$ takes its place, and the formula is derived from the word formula.

(a)
The coordinate two-torus $\Sigma_{12}$ in the plane $\{x_3=y^0_3\}$ consists of the plaquettes $((j,y_2,y^0_3);12)$, and $\Sigma_{13}$ in the plane $\{x_2=y^0_2\}$ of the plaquettes $((j,y^0_2,y_3);13)$, $j\in\Z/n_1$.
Since $F_{((j,y);1\nu)}=F^{\mathrm{el}}_j(y,\nu)$,
\[
\sum_{\sigma\in\Sigma_{12}}F_\sigma=\sum_{j\in\Z/n_1}\langle\mathbf{1}_{S_2},F^{\mathrm{el}}_j\rangle,
\qquad
\sum_{\sigma\in\Sigma_{13}}F_\sigma=\sum_{j\in\Z/n_1}\langle\mathbf{1}_{S_3},F^{\mathrm{el}}_j\rangle,
\]
and, as $q\cdot(2\pi/q)=2\pi$, the flux $k$ of $F$ satisfies
\[
\exp\bigl(2\pi i(\vartheta_2k_{12}+\vartheta_3k_{13})\bigr)=\prod_{j\in\Z/n_1}\exp\bigl(i\langle\beta_\vartheta,F^{\mathrm{el}}_j\rangle\bigr).
\]
Since $F^{\mathrm{mag}}_0(y)=F_{((0,y);23)}$, the coordinate two-torus $\Sigma_{23}$ in the plane $\{x_1=0\}$ gives $\sum_{\sigma\in\Sigma_{23}}F_\sigma=\sum_{\sigma'\in\Lambda^{(2)}_{\mathrm{s}}}F^{\mathrm{mag}}_0(\sigma')$.

Put $\varepsilon':=\min(\varepsilon_v,1/2)$.
The bound on $|v_j|$ holds with $\varepsilon'$ in the place of $\varepsilon_v$, since for $\varepsilon_v>1/2$ its right side is at most $C$ (Remark~\ref{rem:slab-growth});
as $1/2-\varepsilon'\ge0$ and $\sum_j|F^{\mathrm{el}}_j|^2\le\|F\|^2$, it follows that
\[
|G|\le C'\cdot\exp\bigl((1/2-\varepsilon')\|F\|^2\bigr),\qquad C':=C^{n_1}\cdot\prod_j\sup|u_j| .
\]
So $G$ and $|G|$ belong to $\mathcal{I}(\mathbb{T})$ (Definition~\ref{def:torus-integrands}), $Z_{\mathbb{T}}[\,|G|\,]<\infty$ (Proposition~\ref{prop:berezin}\eqref{it:torus-integrability}), and by Definition~\ref{def:averaged-law}\eqref{it:averaged-sectors} the series in \eqref{eq:averaged-fourier} converges absolutely.
For fixed $k_{23}$ the sum of the functions $\Phi_{(k_{23},k_{13},k_{12})}$ over $(k_{12},k_{13})\in\Z^2$ is the indicator of $\sum_{\sigma'}F^{\mathrm{mag}}_0(\sigma')=(2\pi/q)k_{23}$,
and on the set where $\Phi_k=1$ the flux of $F$ is $k$, so that $\exp(2\pi i(\vartheta_2k_{12}+\vartheta_3k_{13}))$ is there the product above.
By dominated convergence (the partial sums are bounded by $|G|$) the left side of \eqref{eq:averaged-fourier} equals
\[
\begin{aligned}
&Z_{\mathbb{T}}\left[\mathbf{1}\left[\,\sum_{\sigma'\in\Lambda^{(2)}_{\mathrm{s}}}F^{\mathrm{mag}}_0(\sigma')=\frac{2\pi}{q}\,k_{23}\right]\cdot\prod_{j}\exp\bigl(i\langle\beta_\vartheta,F^{\mathrm{el}}_j\rangle\bigr)\cdot G\right]\\
&\qquad=Z_{\mathbb{T}}\left[\,\prod_{j\in\Z/n_1}\tilde u_j(F^{\mathrm{mag}}_j)\cdot\tilde v_j(F^{\mathrm{el}}_j)\right],
\end{aligned}
\]
where $\tilde v_j(y):=v_j(y)\cdot\exp(i\langle\beta_\vartheta,y\rangle)$, $\tilde u_j:=u_j$ for $j\ne0$, and $\tilde u_0(\xi):=u_0(\xi)\cdot\mathbf{1}[\sum_{\sigma'}\xi(\sigma')=(2\pi/q)k_{23}]$ for $\xi\in\R^{\Lambda^{(2)}_{\mathrm{s}}}$.
These functions satisfy the hypotheses of Proposition~\ref{prop:word-formula} without local step data, since $|\tilde v_j|=|v_j|$;
in this case the plaquette $\sigma_0$ of the flux string is arbitrary (Proposition~\ref{prop:word-formula}\eqref{it:word-cases}).
The operator of the step $j$ formed with $\tilde v_j$ is $K^\vartheta_j$.
On the summand of flux index $k$ the operator $\tilde u_0(F^{\mathrm{mag}})$ is the multiplication by
\[
\begin{aligned}
&u_0\bigl(d_{\mathrm{s}}A+(2\pi/q)k\,1_{\sigma_0}\bigr)\cdot\mathbf{1}\left[\,\sum_{\sigma'\in\Lambda^{(2)}_{\mathrm{s}}}\bigl(d_{\mathrm{s}}A+(2\pi/q)k\,1_{\sigma_0}\bigr)(\sigma')=\frac{2\pi}{q}\,k_{23}\right]\\
&\qquad=u_0\bigl(d_{\mathrm{s}}A+(2\pi/q)k\,1_{\sigma_0}\bigr)\cdot\mathbf{1}[k=k_{23}],
\end{aligned}
\]
because the sum of $d_{\mathrm{s}}A$ over the plaquettes of the slice torus vanishes (each link of $\Lambda_{\mathrm{s}}$ is a face of exactly two plaquettes, with opposite incidence numbers).
Thus $\tilde u_0(F^{\mathrm{mag}})=u_0(F^{\mathrm{mag}})P_{k_{23}}$.
All operators of the word are diagonal in the flux index and commute with $P_{k_{23}}$, which can therefore be moved to the front of the trace;
the word formula, Proposition~\ref{prop:word-formula}\eqref{it:word-identity}, gives \eqref{eq:averaged-fourier}.

(b)
Let $k_{23}=0$.
By (a) the left side of \eqref{eq:averaged-fourier} is a Fourier series in $\vartheta$ with absolutely summable coefficients;
it converges uniformly, and its integral over $[0,1]^2$ is its coefficient at $k_{12}=k_{13}=0$, namely $Z^{(0)}_{\mathbb{T}}[G]$.
This is the first identity in \eqref{eq:averaged-bloch}.
For $u_j=v_j=1$ the kernel of $K^\vartheta_j$ is the periodisation of $\exp(i\langle\beta_\vartheta,A'-A\rangle)\cdot\prod_bg_1(A'(b)-A(b))$, and $(2\pi/q)^{|\Lambda_{\mathrm{s}}|}K^\vartheta_jP_{\mathrm{G}}=K_{\beta_\vartheta}$ (Proposition~\ref{prop:word-formula}\eqref{it:word-cases}).
The trace in the first identity is then $\Tr\bigl(P_0(K_{\beta_\vartheta}\mathfrak{M}\mathfrak{M}^{*})^{n_1}\bigr)$, and the cyclicity of the trace turns it into $\Tr\bigl(P_0(\mathfrak{M}^{*}K_{\beta_\vartheta}\mathfrak{M})^{n_1}\bigr)$, which gives the second identity;
the cyclicity applies, since $K_{\beta_\vartheta}\ge0$ (Lemma~\ref{lem:slice-kernels}), $K^{1/2}_{\beta_\vartheta}\mathfrak{M}$ is a Hilbert--Schmidt operator (Corollary~\ref{cor:trace-class}), and $P_0$ commutes with $K_{\beta_\vartheta}$ and with $\mathfrak{M}$.
Finally $\mathfrak{M}^{*}K_{\beta_\vartheta}\mathfrak{M}=(K^{1/2}_{\beta_\vartheta}\mathfrak{M})^{*}(K^{1/2}_{\beta_\vartheta}\mathfrak{M})$ is non-negative and of trace class.

(c)
The matrix $g_\mu$ is an even element of $W$ (Definition~\ref{def:other-directions}), and $\mathbb{T}'$ is the coordinate torus $C(n_2,n_3,n_1)$ for $\mu=2$, $C(n_3,n_1,n_2)$ for $\mu=3$.
By Definition~\ref{def:other-directions}, written out in Subsection~\ref{sec:time-directions}, the objects of time $e_\mu$ on $\mathbb{T}=g_\mu\mathbb{T}'$ are the images under $g_\mu$ of those of time $e_1$ on $\mathbb{T}'$:
the electric and the magnetic parts of $F$ in the direction $e_\mu$ are those of $g_\mu^{*}F$ in the direction $e_1$, and the slice torus, the Hilbert spaces and the operators are those of $\mathbb{T}'$, as in (c).
Let $G'$ be the integrand of time $e_1$ on $\mathbb{T}'$ formed with the same $u_j$ and $v_j$, so that $G(F)=G'(g_\mu^{*}F)$.
The bound on $|G'|$ obtained in the proof of (a), now on $\mathbb{T}'$, together with $\|g_\mu^{*}F\|=\|F\|$, shows that $G$, a function of $F$, is a scalar element of $\mathcal{I}(\mathbb{T})$.
Then $G^{g_\mu^{-1}}=G'$, since $G^{g_\mu^{-1}}(F')=G\bigl((g_\mu^{-1})^{*}F'\bigr)$ by Definition~\ref{def:averaged-law}\eqref{it:averaged-covariance} and $g_\mu^{*}(g_\mu^{-1})^{*}=(g_\mu^{-1}g_\mu)^{*}$ is the identity.
Let $\hat k:=k^{g_\mu^{-1}}$ be the vector of Definition~\ref{def:averaged-law}\eqref{it:averaged-covariance}:
the flux of $F$ on $\mathbb{T}$ is $k$ if and only if the flux of $g_\mu^{*}F$ on $\mathbb{T}'$ is $\hat k$.
It is read off from the action of $g_\mu$ on the plaquettes of $\mathbb{T}'$:
$g_2$ maps the plaquettes of the kinds $(12)$, $(13)$, $(23)$ onto plaquettes of the kinds $(23)$, $(12)$, $(13)$, preserving the orientation for the first kind and reversing it for the two others,
and $g_3$ maps them onto plaquettes of the kinds $(13)$, $(23)$, $(12)$, reversing the orientation for the first two kinds and preserving it for the third.
Hence
\[
\begin{aligned}
(\hat k_{12},\hat k_{13},\hat k_{23})&=(k_{23},-k_{12},-k_{13})&&\text{for }\mu=2,\\
(\hat k_{12},\hat k_{13},\hat k_{23})&=(-k_{13},-k_{23},k_{12})&&\text{for }\mu=3 .
\end{aligned}
\]
As $G$ is a scalar element of $\mathcal{I}(\mathbb{T})$ and $g_\mu^{-1}$ is even, Definition~\ref{def:averaged-law}\eqref{it:averaged-covariance} gives $Z^{(k)}_{\mathbb{T}}[G]=Z^{(\hat k)}_{\mathbb{T}'}[G']$ for every $k\in\Z^3$, and $\hat k=0$ for $k=0$.
So (b) on $\mathbb{T}'$ is (b) on $\mathbb{T}$ with time $e_\mu$.
The series of (a) on $\mathbb{T}'$, at fixed $\hat k_{23}$, is a sum over $(\hat k_{12},\hat k_{13})\in\Z^2$, that is, over $(k_{23},k_{12})$ at fixed $k_{13}=-\hat k_{23}$ for $\mu=2$, and over $(k_{13},k_{23})$ at fixed $k_{12}=\hat k_{23}$ for $\mu=3$;
with the values of $\hat k_{12}$, $\hat k_{13}$ above, (a) on $\mathbb{T}'$ becomes the two assertions of (c).
\end{proof}

\begin{proposition}[Reflection positivity and chessboard estimate]\label{prop:averaged-positivity}
Let $-1\le\rho\le1$, $w_i>0$, $M\in\R$, $q>0$, and $\mathbb{T}=C(n_1,n_2,n_3)$ with $n_\mu\ge4$.
\begin{enumerate}
\item[(a)]\itemmark{a}\label{it:averaged-reflection-positivity}
Let $n_1$ be even and $c\in\frac12\Z$.
For $\beta\in C^2(\mathbb{T})$ let $\varphi_\beta(F):=\exp(i\langle\beta,F\rangle)$ for the curvature $F\in C^2(\mathbb{T})$, the pairing being that of $C^2(\Gamma)$ for the period lattice $\Gamma$ of $\mathbb{T}$;
these functions are the characters of $\mathbb{T}$.
Let $G=\sum_\alpha a_\alpha\varphi_{\beta_\alpha}$ be a series of characters of $\mathbb{T}$ with complex coefficients, $\sum_\alpha|a_\alpha|<\infty$, every $\beta_\alpha\in C^2(\mathbb{T})$ being supported on the plaquettes of $\mathbb{T}$ all of whose corners satisfy $c<x_1<c+n_1/2$.
Then
\[
Z^{(0)}_{\mathbb{T}}\bigl[\Theta_{x_1,c}(G)\,G\bigr]\ge0 .
\]
Consequently $\mu^{\mathrm{av}}_{\mathbb{T}}\bigl(\overline{\Phi\circ\theta^{*}_{x_1,c}}\,\Phi\bigr)\ge0$ for every bounded Borel function $\Phi$ of finitely many $F_\sigma$, $\sigma$ a plaquette all of whose corners satisfy $c<x_1<c+n_1/2$.
The same holds for $\mu=2,3$, with $n_\mu$ and $x_\mu$ in the place of $n_1$ and $x_1$.
\item[(b)]\itemmark{b}\label{it:averaged-chessboard}
Let $\mathbb{T}$ carry a coordinate cell structure $\mathcal{Q}=(\mathfrak{F}_{\mathrm{c}};o_1,o_2,o_3;h)$, $o_\mu\in\frac12\Z$.
Then $Z^{(0)}_{\mathbb{T}}[\Fold_{\mathbb{T}}O]\ge0$ for every $O\in\mathcal{T}_{\C}$ (Definition~\ref{def:complex-polynomials}) whose support lies in an open cell, and for $O_\kappa\in\mathcal{T}_{\C}$ with $\supp O_\kappa\subset Q_\kappa$, one $\kappa$ for each cell of $\mathbb{T}$,
\[
\left|Z^{(0)}_{\mathbb{T}}\left[\,\prod_\kappa O_\kappa\right]\right|\le\prod_\kappa Z^{(0)}_{\mathbb{T}}\bigl[\Fold_{\mathbb{T}}O_\kappa\bigr]^{\frac{|Q|}{|\mathbb{T}|}} .
\]
\end{enumerate}
\end{proposition}

\begin{proof}
The proof follows the arguments of \cite[Corollary~21.21]{NoncompactQED} and \cite[Proposition~21.22]{NoncompactQED}, where the fermion weight is the determinant of the Wilson operator;
here $\det D^{\mathbb{T}}_A$ takes its place, and the operators are those of Lemma~\ref{lem:bloch}.
The geometry differs as well: there the width of the cells is half an odd integer and consecutive walls are alternately site planes and link planes, whereas on $\mathbb{T}$ the planes at which a word is cut are all of one kind and the width $h$ in (b) is an integer.
Time is $e_1$ except where the directions $e_2$, $e_3$ are treated, the notation is that of Lemma~\ref{lem:bloch}, and $\vartheta\in[0,1]^2$ is fixed until an integral over $\vartheta$ is taken.

\emph{Characters.}
Let $\Gamma$ be the period lattice of $\mathbb{T}$, and let $\varphi_\beta$, $\beta\in C^2(\mathbb{T})$, be the characters of (a).
For $\mu\in\{1,2,3\}$ and $c\in\frac12\Z$ the linear part of the reflection $\theta_{x_\mu,c}$ maps $\Gamma$ onto itself;
so $\theta_{x_\mu,c}$, like every translation $\tau_v$, $v\in\Z^3$, permutes the $\Gamma$-orbits of plaquettes,
and the pull-backs by these maps act on $C^2(\mathbb{T})$ and commute with the periodic extension $F\mapsto\widetilde F$.
On a function $G$ of the curvature of $\mathbb{T}$ the reflection acts by $(\Theta_{x_\mu,c}G)(F)=\overline{G(\theta^{*}_{x_\mu,c}F)}$, as in \eqref{eq:layers-bosonic};
since $\theta^{*}_{x_\mu,c}$ is an orthogonal involution of $C^2(\mathbb{T})$, one has $\langle\beta,\theta^{*}_{x_\mu,c}F\rangle=\langle\theta^{*}_{x_\mu,c}\beta,F\rangle$, and hence
\[
\Theta_{x_\mu,c}\,\varphi_\beta=\varphi_{-\theta^{*}_{x_\mu,c}\beta}\qquad\left(\beta\in C^2(\mathbb{T})\right).
\]
For $\beta\in C^2_{\mathrm{c}}$ let $\beta_\Gamma\in C^2(\mathbb{T})$ be its periodisation,
\[
\beta_\Gamma(\sigma):=\sum_{g\in\Gamma}\beta(\sigma+g),
\]
a finite sum that depends only on the plaquette of $\mathbb{T}$ represented by the plaquette $\sigma$ of $\Z^3$;
if $\supp\beta$ lies in a box that embeds in $\Gamma$, it is the cochain of Lemma~\ref{lem:curvature-characters}(b).
As $\widetilde F_{\sigma+g}=\widetilde F_\sigma$ for $g\in\Gamma$, one has $\langle\beta,\widetilde F\rangle=\langle\beta_\Gamma,F\rangle$, that is,
\[
\varphi_\beta(\widetilde F)=\varphi_{\beta_\Gamma}(F),
\]
with the function $\varphi_\beta$ of Definition~\ref{def:characters} on the left.
If $\beta$ is supported on plaquettes all of whose corners satisfy $c_1<x_\mu<c_2$, where $c_2-c_1\le n_\mu/2$, then so is $\beta_\Gamma$, as a cochain of $\mathbb{T}$, the periods changing $x_\mu$ by multiples of $n_\mu$.
Thus an element $O=\sum_\alpha a_\alpha\varphi_{\beta_\alpha}$ of $\mathcal{T}_{\C}$, read on $\mathbb{T}$ at $\widetilde F$, is the finite sum $\sum_\alpha a_\alpha\varphi_{(\beta_\alpha)_\Gamma}$ of characters of $\mathbb{T}$,
and this reading commutes with the action of the translations and of the reflections $\Theta_{x_\mu,c}$.
For $\beta\in C^2(\mathbb{T})$ let $\beta^{\mathrm{el}}_j(y,\nu):=\beta((j,y);1\nu)$ and $\beta^{\mathrm{mag}}_j(y):=\beta((j,y);23)$ be its electric part at the step $j$ and its magnetic part at the slice $j$, so that
\[
\varphi_\beta=\prod_{j\in\Z/n_1}\exp\bigl(i\langle\beta^{\mathrm{el}}_j,F^{\mathrm{el}}_j\rangle\bigr)\cdot\exp\bigl(i\langle\beta^{\mathrm{mag}}_j,F^{\mathrm{mag}}_j\rangle\bigr),
\]
and let $\mathfrak{N}[\beta^{\mathrm{mag}}_j]:=\mathfrak{M}\exp\bigl(i\langle\beta^{\mathrm{mag}}_j,F^{\mathrm{mag}}\rangle\bigr)\mathfrak{M}^{*}$, an operator on $\mathcal{H}^{\mathrm{l}}$.
As $\exp(i\langle\beta^{\mathrm{mag}}_j,F^{\mathrm{mag}}\rangle)$ is a unitary multiplication operator, $\mathfrak{N}[-\beta^{\mathrm{mag}}_j]=\mathfrak{N}[\beta^{\mathrm{mag}}_j]^{*}$.
With $u_j:=\exp(i\langle\beta^{\mathrm{mag}}_j,\cdot\,\rangle)$ and $v_j:=\exp(i\langle\beta^{\mathrm{el}}_j,\cdot\,\rangle)$ one has $(2\pi/q)^{|\Lambda_{\mathrm{s}}|}K^\vartheta_jP_{\mathrm{G}}=K_{\beta^{\mathrm{el}}_j+\beta_\vartheta}$ (Proposition~\ref{prop:word-formula}\eqref{it:word-cases}),
and Lemma~\ref{lem:bloch}\eqref{it:bloch-fourier} gives
\begin{multline}\label{eq:averaged-character-word}
\sum_{k_{12},k_{13}\in\Z}\exp\bigl(2\pi i(\vartheta_2k_{12}+\vartheta_3k_{13})\bigr)\cdot Z^{(0,k_{13},k_{12})}_{\mathbb{T}}[\varphi_\beta]\\
=(2\pi)^{|n|}\cdot\Tr\left(P_0\prod_{j=n_1-1,\dots,0}K_{\beta^{\mathrm{el}}_j+\beta_\vartheta}\,\mathfrak{N}[\beta^{\mathrm{mag}}_j]\right);
\end{multline}
by Lemma~\ref{lem:bloch}\eqref{it:bloch-zero-flux} the integral of the right side over $\vartheta\in[0,1]^2$ is $Z^{(0)}_{\mathbb{T}}[\varphi_\beta]$.
By linearity both statements extend to finite sums of characters.
The operators $K_{\beta'}$, $\beta'\in\R^{\Lambda^{(1)}_{\mathrm{s}}}$, are bounded, self-adjoint and non-negative, and $K^{1/2}_{\beta'}\mathfrak{M}$ is a Hilbert--Schmidt operator (Definition~\ref{def:slice-operators}, Lemma~\ref{lem:slice-kernels}, Corollary~\ref{cor:trace-class}).
All operators that occur are diagonal in the flux index and commute with $P_0$.

\emph{The part of a word between two planes.}
Let $c_1<c_2$ be points of $\frac12\Z$ of the same kind with $c_2-c_1\le n_1/2$, and let $\beta\in C^2(\mathbb{T})$ be supported on the plaquettes of $\mathbb{T}$ all of whose corners satisfy $c_1<x_1<c_2$;
indices of steps and of slices are read modulo $n_1$.
In the following ordered products later times stand to the left, and an empty product is the identity.
If $c_1=j_1+\frac12$ and $c_2=j_2+\frac12$ are link planes, then $\beta^{\mathrm{el}}_{j_1}=\beta^{\mathrm{el}}_{j_2}=0$, and
\[
\mathfrak{Y}[\beta]:=K^{1/2}_{\beta_\vartheta}\,\mathfrak{N}[\beta^{\mathrm{mag}}_{j_2}]\left(\prod_{j=j_2-1,\dots,j_1+1}K_{\beta^{\mathrm{el}}_j+\beta_\vartheta}\,\mathfrak{N}[\beta^{\mathrm{mag}}_j]\right)K^{1/2}_{\beta_\vartheta},
\]
an operator on $\mathcal{H}^{\mathrm{l}}$.
If $c_1=j_1$ and $c_2=j_2$ are site planes, then $\beta^{\mathrm{mag}}_{j_1}=\beta^{\mathrm{mag}}_{j_2}=0$ and $\beta^{\mathrm{el}}_{j_1}=\beta^{\mathrm{el}}_{j_2-1}=0$, and
\[
\mathfrak{Y}[\beta]:=\mathfrak{M}^{*}\left(\prod_{j=j_2-1,\dots,j_1+1}K_{\beta^{\mathrm{el}}_j+\beta_\vartheta}\,\mathfrak{N}[\beta^{\mathrm{mag}}_j]\right)K_{\beta_\vartheta}\,\mathfrak{M},
\]
an operator on $\mathcal{H}^{\mathrm{s}}$;
since $\beta^{\mathrm{el}}_{j_2-1}=0$, it begins with $\mathfrak{M}^{*}K_{\beta_\vartheta}$ in all cases.
So $\mathfrak{Y}[\beta]=\mathfrak{Y}^\vartheta_{c_1,c_2}[\beta]$ is the part of the product in \eqref{eq:averaged-character-word} between the two planes, the operator $K_{\beta_\vartheta}$ of a step that a link plane cuts being split as $K^{1/2}_{\beta_\vartheta}\cdot K^{1/2}_{\beta_\vartheta}$, and the operator $\mathfrak{M}\mathfrak{M}^{*}$ of a slice on a site plane as $\mathfrak{M}\cdot\mathfrak{M}^{*}$.
Put $\mathfrak{Y}[\varphi_\beta]:=\mathfrak{Y}[\beta]$, and extend $\mathfrak{Y}$ linearly to finite formal sums of such characters.
For an element $O=\sum_\alpha a_\alpha\varphi_{\beta_\alpha}$ of $\mathcal{T}_{\C}$ whose cochains $\beta_\alpha\in C^2_{\mathrm{c}}$ are supported on plaquettes all of whose corners satisfy $c_1<x_1<c_2$, the periodised cochains $(\beta_\alpha)_\Gamma$ are supported in the same way, and $\mathfrak{Y}[O]:=\sum_\alpha a_\alpha\,\mathfrak{Y}\bigl[(\beta_\alpha)_\Gamma\bigr]$ is formed with them.
It has three properties.

First, $\mathfrak{Y}[\beta]$ is a Hilbert--Schmidt operator:
it begins with $K^{1/2}_{\beta_\vartheta}\mathfrak{M}$, resp.\ with $\mathfrak{M}^{*}K^{1/2}_{\beta_\vartheta}=(K^{1/2}_{\beta_\vartheta}\mathfrak{M})^{*}$, and all its other factors are bounded.

Second, it does not change if $c_1$, $c_2$ and $\beta$ are translated by the same vector of $\Z e_1$, because the operators of a step and of a slice depend on its index only through the parts of $\beta$.

Third, a reflection in the upper plane gives the adjoint.
The reflection $\theta:=\theta_{x_1,c_2}$ maps the step $j$ onto the step $2c_2-1-j$ and the slice $j$ onto the slice $2c_2-j$, reversing the orientation of the electric plaquettes in the direction $e_1$;
hence
\[
F^{\mathrm{el}}_j(\theta^{*}F)=-F^{\mathrm{el}}_{2c_2-1-j}(F),\qquad F^{\mathrm{mag}}_j(\theta^{*}F)=F^{\mathrm{mag}}_{2c_2-j}(F),
\]
which for a link plane is \eqref{eq:setting-slice-reflections}.
By the rule $\Theta_{x_\mu,c}\varphi_\beta=\varphi_{-\theta^{*}_{x_\mu,c}\beta}$ for the characters of $\mathbb{T}$, one has $\Theta_{x_1,c_2}\varphi_\beta=\varphi_{\beta'}$ with $\beta':=-\theta^{*}\beta$, whose parts are $\beta'^{\mathrm{el}}_{2c_2-1-j}=\beta^{\mathrm{el}}_j$ and $\beta'^{\mathrm{mag}}_{2c_2-j}=-\beta^{\mathrm{mag}}_j$;
the cochain $\beta'$ is supported between $c_2$ and $2c_2-c_1$.
Written out, $\mathfrak{Y}^\vartheta_{c_2,2c_2-c_1}[\beta']$ consists of the factors of $\mathfrak{Y}^\vartheta_{c_1,c_2}[\beta]$ in the reverse order, each replaced by its adjoint:
the operators $K$ are self-adjoint, $\mathfrak{N}[-\beta^{\mathrm{mag}}_j]=\mathfrak{N}[\beta^{\mathrm{mag}}_j]^{*}$, the outer factors $\mathfrak{M}^{*}$ and $\mathfrak{M}$ of the site planes are each other's adjoints, and the shift $\beta_\vartheta$ is the same at every step.
Both sides being antilinear in $G$, it follows that
\begin{equation}\label{eq:averaged-adjoint-part}
\mathfrak{Y}^\vartheta_{c_2,2c_2-c_1}\bigl[\Theta_{x_1,c_2}G\bigr]=\mathfrak{Y}^\vartheta_{c_1,c_2}[G]^{*}
\end{equation}
for every finite formal sum $G$ of characters of $\mathbb{T}$ supported on the plaquettes all of whose corners satisfy $c_1<x_1<c_2$,
and hence, the periodisation commuting with $\theta^{*}$, for every element $G$ of $\mathcal{T}_{\C}$ whose cochains are supported in this way.

\emph{Proof of (a).}
Let $\beta_1,\beta_2\in C^2(\mathbb{T})$ be supported on the plaquettes all of whose corners satisfy $c<x_1<c+n_1/2$.
As $n_1$ is even, the planes $c$ and $c+n_1/2$ are of the same kind, and $\mathfrak{Y}:=\mathfrak{Y}^\vartheta_{c,c+n_1/2}$ is defined on these cochains.
On $\mathbb{T}$ the reflections $\theta_{x_1,c}$ and $\theta_{x_1,c+n_1/2}$ induce the same map, since their product is the translation by $n_1e_1$, a period of $\mathbb{T}$;
so $\Theta_{x_1,c}(\varphi_{\beta_1})=\Theta_{x_1,c+n_1/2}(\varphi_{\beta_1})$, a character supported between $c+n_1/2$ and $c+n_1$, whose part of the word between these two planes is $\mathfrak{Y}[\beta_1]^{*}$ by \eqref{eq:averaged-adjoint-part}.
Splitting the product in \eqref{eq:averaged-character-word} for the character $\Theta_{x_1,c}(\varphi_{\beta_1})\varphi_{\beta_2}$ at the two planes, and using the cyclicity of the trace (the two parts are Hilbert--Schmidt operators), one obtains
\begin{multline}\label{eq:averaged-two-halves}
\sum_{k_{12},k_{13}\in\Z}\exp\bigl(2\pi i(\vartheta_2k_{12}+\vartheta_3k_{13})\bigr)\cdot Z^{(0,k_{13},k_{12})}_{\mathbb{T}}\bigl[\Theta_{x_1,c}(\varphi_{\beta_1})\,\varphi_{\beta_2}\bigr]\\
=(2\pi)^{|n|}\cdot\Tr\bigl(P_0\,\mathfrak{Y}[\beta_1]^{*}\,\mathfrak{Y}[\beta_2]\bigr),
\end{multline}
and the integral of the right side over $\vartheta$ is $Z^{(0)}_{\mathbb{T}}[\Theta_{x_1,c}(\varphi_{\beta_1})\varphi_{\beta_2}]$.
For $\beta_1=\beta_2=\beta$ the right side of \eqref{eq:averaged-two-halves} is $(2\pi)^{|n|}\cdot\|P_0\,\mathfrak{Y}[\beta]\|^2_{\mathrm{HS}}$, because $P_0$ is an orthogonal projection that commutes with $\mathfrak{Y}[\beta]$;
the integrand $\Theta_{x_1,c}(\varphi_\beta)\varphi_\beta$ has modulus $1$, so that the left side has a modulus at most $Z_{\mathbb{T}}[1]$ by Definition~\ref{def:averaged-law}\eqref{it:averaged-sectors}.
Hence, for every such $\beta$ and every $\vartheta\in[0,1]^2$,
\begin{equation}\label{eq:averaged-halves-bound}
\|P_0\,\mathfrak{Y}[\beta]\|^2_{\mathrm{HS}}\le(2\pi)^{-|n|}\cdot Z_{\mathbb{T}}[1].
\end{equation}
Let now $G=\sum_\alpha a_\alpha\varphi_{\beta_\alpha}$ be as in (a), so that $\Theta_{x_1,c}(G)=\sum_\alpha\overline{a_\alpha}\,\Theta_{x_1,c}(\varphi_{\beta_\alpha})$.
By \eqref{eq:averaged-halves-bound} the series $P_0\,\mathfrak{Y}[G]:=\sum_\alpha a_\alpha P_0\,\mathfrak{Y}[\beta_\alpha]$ converges in the Hilbert--Schmidt norm, uniformly in $\vartheta$.
The double series $\Theta_{x_1,c}(G)\,G=\sum_{\alpha,\alpha'}\overline{a_\alpha}\cdot a_{\alpha'}\cdot\Theta_{x_1,c}(\varphi_{\beta_\alpha})\,\varphi_{\beta_{\alpha'}}$ converges uniformly,
and $|Z^{(0)}_{\mathbb{T}}[H]|\le Z_{\mathbb{T}}[1]\cdot\sup|H|$ for every bounded Borel function $H$ of $F$, by \eqref{eq:setting-modulus-bound}.
By \eqref{eq:averaged-two-halves}, integrated over $\vartheta$, and by dominated convergence (the summands are bounded by $|a_\alpha|\cdot|a_{\alpha'}|\cdot Z_{\mathbb{T}}[1]$, by \eqref{eq:averaged-halves-bound}),
\[
\begin{aligned}
Z^{(0)}_{\mathbb{T}}\bigl[\Theta_{x_1,c}(G)\,G\bigr]
&=(2\pi)^{|n|}\int_{[0,1]^2}\mathrm{d}\vartheta\sum_{\alpha,\alpha'}\overline{a_\alpha}\cdot a_{\alpha'}\cdot\Tr\bigl((P_0\mathfrak{Y}[\beta_\alpha])^{*}\,P_0\mathfrak{Y}[\beta_{\alpha'}]\bigr)\\
&=(2\pi)^{|n|}\int_{[0,1]^2}\mathrm{d}\vartheta\,\bigl\|P_0\,\mathfrak{Y}[G]\bigr\|^2_{\mathrm{HS}}\ge0 .
\end{aligned}
\]

Let $\mu\in\{2,3\}$, $n_\mu$ even, $c\in\frac12\Z$, and let $G=\sum_\alpha a_\alpha\varphi_{\beta_\alpha}$, $\sum_\alpha|a_\alpha|<\infty$, with every $\beta_\alpha$ supported on the plaquettes all of whose corners satisfy $c<x_\mu<c+n_\mu/2$.
Put $\mathbb{T}':=g_\mu^{-1}\mathbb{T}$, whose side in the direction $e_1$ is $n_\mu$, and $G':=G^{g_\mu^{-1}}$, formed as in Definition~\ref{def:averaged-law}\eqref{it:averaged-covariance} for the bounded scalar element $G$ of $\mathcal{I}(\mathbb{T})$, a function of $F$.
Since $g_\mu^{-1}$ maps $\{c<x_\mu<c+n_\mu/2\}$ onto $\{c<x_1<c+n_\mu/2\}$, $G'$ is a series of characters of $\mathbb{T}'$ as in the first part, with the same coefficients.
As $g_\mu^{-1}\theta_{x_\mu,c}\,g_\mu=\theta_{x_1,c}$, one has $(\Theta_{x_\mu,c}G)^{g_\mu^{-1}}=\Theta_{x_1,c}(G')$ by \eqref{eq:layers-bosonic} and Definition~\ref{def:averaged-law}\eqref{it:averaged-covariance}.
The element $g_\mu^{-1}$ is even and $0^{g_\mu^{-1}}=0$; by Definition~\ref{def:averaged-law}\eqref{it:averaged-covariance} and the first part, applied to $\mathbb{T}'$,
\[
Z^{(0)}_{\mathbb{T}}\bigl[\Theta_{x_\mu,c}(G)\,G\bigr]=Z^{(0)}_{\mathbb{T}'}\bigl[\Theta_{x_1,c}(G')\,G'\bigr]\ge0 .
\]

Let finally $\mu\in\{1,2,3\}$, $n_\mu$ even, $c\in\frac12\Z$, $\theta:=\theta_{x_\mu,c}$, and $\Phi(F)=u\bigl((F_\sigma)_{\sigma\in\Sigma_\Phi}\bigr)$ with a finite set $\Sigma_\Phi$ of plaquettes all of whose corners satisfy $c<x_\mu<c+n_\mu/2$ and a bounded Borel function $u$ on $\R^{\Sigma_\Phi}$.
Let $\lambda$ be the sum of the laws, under $\mu^{\mathrm{av}}_{\mathbb{T}}$, of the random vectors $(F_\sigma)_{\sigma\in\Sigma_\Phi}$ and $((\theta^{*}F)_\sigma)_{\sigma\in\Sigma_\Phi}$, a finite Borel measure on $\R^{\Sigma_\Phi}$.
The finite linear combinations of the functions $\xi\mapsto\exp(i\langle\beta,\xi\rangle)$, $\beta\in\R^{\Sigma_\Phi}$, are dense in $L^2(\lambda)$:
a function in $L^2(\lambda)$ that is orthogonal to all of them is the density of a finite complex measure whose Fourier transform vanishes, hence it vanishes $\lambda$-almost everywhere (uniqueness theorem for the Fourier transforms of finite measures).
Choose such combinations $u_m$ with $\|u_m-u\|_{L^2(\lambda)}\to0$, and put $\Phi'_m:=u_m\bigl((F_\sigma)_{\sigma\in\Sigma_\Phi}\bigr)$.
With $\beta\in\R^{\Sigma_\Phi}$ extended by zero to an element of $C^2_{\mathrm{c}}$, the function $\xi\mapsto\exp(i\langle\beta,\xi\rangle)$ at $\xi=(\widetilde F_\sigma)_{\sigma\in\Sigma_\Phi}$ is $\varphi_{\beta_\Gamma}(F)$, and $\beta_\Gamma$ is supported on plaquettes of $\mathbb{T}$ all of whose corners satisfy $c<x_\mu<c+n_\mu/2$;
so, as a function of the curvature of $\mathbb{T}$, $\Phi'_m$ is a finite sum of characters of the kind admitted above, and
by \eqref{eq:layers-bosonic} and the definition of $\mu^{\mathrm{av}}_{\mathbb{T}}$,
\[
\mu^{\mathrm{av}}_{\mathbb{T}}\bigl(\overline{\Phi'_m\circ\theta^{*}}\,\Phi'_m\bigr)=Z^{(0)}_{\mathbb{T}}\bigl[\Theta_{x_\mu,c}(\Phi'_m)\,\Phi'_m\bigr]\big/Z^{(0)}_{\mathbb{T}}[1]\ge0 .
\]
Since $\lambda$ dominates the laws of both random vectors, the Cauchy--Schwarz inequality gives
\[
\begin{aligned}
&\left|\mu^{\mathrm{av}}_{\mathbb{T}}\bigl(\overline{\Phi'_m\circ\theta^{*}}\,\Phi'_m\bigr)-\mu^{\mathrm{av}}_{\mathbb{T}}\bigl(\overline{\Phi\circ\theta^{*}}\,\Phi\bigr)\right|\\
&\qquad\le\mu^{\mathrm{av}}_{\mathbb{T}}\bigl(|\Phi'_m\circ\theta^{*}|\cdot|\Phi'_m-\Phi|\bigr)+\mu^{\mathrm{av}}_{\mathbb{T}}\bigl(|(\Phi'_m-\Phi)\circ\theta^{*}|\cdot|\Phi|\bigr)\\
&\qquad\le\|u_m-u\|_{L^2(\lambda)}\cdot\bigl(\|u_m\|_{L^2(\lambda)}+\|u\|_{L^2(\lambda)}\bigr),
\end{aligned}
\]
which tends to $0$.
So $\mu^{\mathrm{av}}_{\mathbb{T}}\bigl(\overline{\Phi\circ\theta^{*}}\,\Phi\bigr)\ge0$.

\emph{Proof of (b).}
As $\mathbb{T}$ carries $\mathcal{Q}$, the sides satisfy $n_a\in2h\Z$, $a=1,2,3$ (Definition~\ref{def:cell-structures}).

\emph{Step 1: one direction.}
Let $2r:=n_1/h$ and $c_j:=o_1+jh$ for $j\in\Z$;
these planes are of one kind, and $c_{j+1}-c_j=h\le n_1/2$.
For $j=0,\dots,2r-1$ let $G_j\in\mathcal{T}_{\C}$ have its support in the layer $\{c_j<x_1<c_{j+1}\}$, and let $\mathfrak{Y}_j:=\mathfrak{Y}^\vartheta_{c_j,c_{j+1}}[G_j]$, formed with the periodised cochains of $G_j$;
read on $\mathbb{T}$, $G_j$ is a finite sum of characters of $\mathbb{T}$ supported between $c_j$ and $c_{j+1}$, and $G_0\cdots G_{2r-1}$ a finite sum of products of such characters, one for each layer.
For a character of $\mathbb{T}$ supported in the union of these layers the product in \eqref{eq:averaged-character-word} splits at the planes $c_0,\dots,c_{2r}$;
by multilinearity and the cyclicity of the trace,
\[
Z^{(0)}_{\mathbb{T}}[G_0\cdots G_{2r-1}]=(2\pi)^{|n|}\int_{[0,1]^2}\mathrm{d}\vartheta\,\Tr\bigl(P_0\,\mathfrak{Y}_{2r-1}\cdots\mathfrak{Y}_0\bigr).
\]
The $\mathfrak{Y}_j$ act on one separable Hilbert space, $\mathcal{H}^{\mathrm{l}}$ if the $c_j$ are link planes and $\mathcal{H}^{\mathrm{s}}$ if they are site planes, and $P_0\,\mathfrak{Y}_j$ is a Hilbert--Schmidt operator, hence an element of $\mathscr{J}_{2r}$.
H\"older's inequality in the Schatten classes \cite[Lemma~10.6]{NoncompactQED} gives
\[
\begin{aligned}
\bigl|\Tr\bigl(P_0\,\mathfrak{Y}_{2r-1}\cdots\mathfrak{Y}_0\bigr)\bigr|
&=\bigl|\Tr\bigl((P_0\mathfrak{Y}_{2r-1})\cdots(P_0\mathfrak{Y}_0)\bigr)\bigr|\\
&\le\prod_{j=0}^{2r-1}\|P_0\,\mathfrak{Y}_j\|_{2r}
=\prod_{j=0}^{2r-1}\Tr\bigl(P_0\,(\mathfrak{Y}_j^{*}\mathfrak{Y}_j)^r\bigr)^{\frac{1}{2r}},
\end{aligned}
\]
and H\"older's inequality for the integral over $\vartheta$, with the $2r$ exponents $2r$, gives
\[
\bigl|Z^{(0)}_{\mathbb{T}}[G_0\cdots G_{2r-1}]\bigr|\le\prod_{j=0}^{2r-1}\left((2\pi)^{|n|}\int_{[0,1]^2}\mathrm{d}\vartheta\,\Tr\bigl(P_0\,(\mathfrak{Y}_j^{*}\mathfrak{Y}_j)^r\bigr)\right)^{\frac{1}{2r}} .
\]
Let $\fold_{x_1}G_j$ be the fold of $G_j$ in the direction $x_1$ of Definition~\ref{def:slab-copies}, formed with $\ell=x_1$, $c_0=o_1$, $h'=h$ and $k=r$, the layer $j$ being the slab $j+1$ of that definition.
It is defined: $\mathbb{T}$ is adapted to $x_1$, with the height period $n_1=2rh$ (Definition~\ref{def:cell-structures}),
and $G_j$, a finite sum of bounded functions of finitely many $F_\sigma$ with $\sigma$ contained in the open layer $j$, belongs to $\mathcal{A}_\Xi[c_j,c_{j+1}]$ for $\Xi:=\Xi(\mathbb{T},x_1)$ (Definition~\ref{def:slabs}).
For $j'=0,\dots,2r-1$ the element of the group $W(x_1;o_1,h)$ generated by the reflections $\theta_{x_1,c}$, $c\in o_1+h\Z$, that maps the layer $j$ onto the layer $j'$ is the translation by $(j'-j)h\,e_1$ if $j'-j$ is even, and the reflection $\theta_{x_1,c_{j+1}}$ followed by the translation by $(j'-j-1)h\,e_1$ if $j'-j$ is odd;
the image of $G_j$ under this element is the copy of $G_j$ in the layer $j'$, and $\fold_{x_1}G_j$ is the product of these $2r$ copies.
By the second and the third property of $\mathfrak{Y}$, the part of the word of $\fold_{x_1}G_j$ between $c_{j'}$ and $c_{j'+1}$ is $\mathfrak{Y}_j$ if $j'-j$ is even and $\mathfrak{Y}_j^{*}$ if $j'-j$ is odd.
By \eqref{eq:averaged-character-word} and the cyclicity of the trace, therefore,
\[
Z^{(0)}_{\mathbb{T}}[\fold_{x_1}G_j]=(2\pi)^{|n|}\int_{[0,1]^2}\mathrm{d}\vartheta\,\Tr\bigl(P_0\,(\mathfrak{Y}_j^{*}\mathfrak{Y}_j)^r\bigr)\ge0,
\]
since $\Tr\bigl(P_0(\mathfrak{Y}_j^{*}\mathfrak{Y}_j)^r\bigr)=\Tr\bigl((\mathfrak{Y}_j^{*}P_0\mathfrak{Y}_j)^r\bigr)\ge0$.
Hence
\begin{equation}\label{eq:averaged-slab-holder}
Z^{(0)}_{\mathbb{T}}[\fold_{x_1}G_j]\ge0,
\qquad
\bigl|Z^{(0)}_{\mathbb{T}}[G_0\cdots G_{2r-1}]\bigr|\le\prod_{j=0}^{2r-1}Z^{(0)}_{\mathbb{T}}[\fold_{x_1}G_j]^{\frac{1}{2r}} .
\end{equation}

\emph{Step 2: the directions $e_2$ and $e_3$.}
Let $\mu\in\{2,3\}$, $2r:=n_\mu/h$, $c_j:=o_\mu+jh$, and let $G_j\in\mathcal{T}_{\C}$ have its support in $\{c_j<x_\mu<c_{j+1}\}$, $j=0,\dots,2r-1$.
Then \eqref{eq:averaged-slab-holder} holds with $\fold_{x_\mu}$ in the place of $\fold_{x_1}$.
Indeed, let $\mathbb{T}':=g_\mu^{-1}\mathbb{T}$, whose side in the direction $e_1$ is $n_\mu$, and $G'_j:=G_j^{g_\mu^{-1}}$, an element of $\mathcal{T}_{\C}$ with support in $\{c_j<x_1<c_{j+1}\}$.
For an element $g$ of the group generated by the reflections $\theta_{x_\mu,c}$, $c\in o_\mu+h\Z$, one has $(G_j^{g})^{g_\mu^{-1}}=(G'_j)^{g_\mu^{-1}g\,g_\mu}$ by \eqref{eq:setting-group-action}, and $g_\mu^{-1}\theta_{x_\mu,c}\,g_\mu=\theta_{x_1,c}$;
as $G\mapsto G^{g_\mu^{-1}}$ is multiplicative, $(\fold_{x_\mu}G_j)^{g_\mu^{-1}}=\fold_{x_1}G'_j$, and $(G_0\cdots G_{2r-1})^{g_\mu^{-1}}=G'_0\cdots G'_{2r-1}$.
By Definition~\ref{def:averaged-law}\eqref{it:averaged-covariance} ($g_\mu^{-1}$ is even, $0^{g_\mu^{-1}}=0$) the numbers in question are the corresponding numbers on $\mathbb{T}'$, where Step 1 applies.

\emph{Step 3: positivity of the folds.}
In the rest of the proof $\kappa$, $\kappa'$ range over the $\kappa\in\Z^3$ with $0\le\kappa_a<n_a/h$, $a=1,2,3$.
The cells $Q_\kappa$ with these indices contain exactly one cell of each orbit of the period lattice $\bigoplus_a n_a\Z e_a$, because $\tau_{n_ae_a}Q_\kappa=Q_{\kappa+(n_a/h)e_a}$;
a function of $F$ that reads in a cell and its translates by the periods of $\mathbb{T}$ define the same function of the curvature of $\mathbb{T}$.
So we may and do take the indices of (b) in this range.
For $a\in\{1,2,3\}$ and $j,j'\in\Z$ let $g^{a}_{j'\leftarrow j}\in W(\mathcal{Q})$ be the element that maps the layer $\{o_a+jh<x_a<o_a+(j+1)h\}$ onto the layer $\{o_a+j'h<x_a<o_a+(j'+1)h\}$ and does not change the other two coordinates:
the translation by $(j'-j)h\,e_a$, which is the reflection in the wall $\{x_a=o_a\}$ followed by the reflection in the wall $\{x_a=o_a+(j'-j)h/2\}$, if $j'-j$ is even, and the reflection in the wall $\{x_a=o_a+(j+j'+1)h/2\}$ if $j'-j$ is odd.
Elements with different $a$ commute, and $g^3_{\kappa'_3\leftarrow\kappa_3}g^2_{\kappa'_2\leftarrow\kappa_2}g^1_{\kappa'_1\leftarrow\kappa_1}$ maps $Q_\kappa$ onto $Q_{\kappa'}$.
As $W(\mathcal{Q})$ acts simply transitively on the cells (Definition~\ref{def:cell-structures}), the mirror copies of an $O$ that reads in $Q_\kappa$ are (Definition~\ref{def:folds})
\[
O^{[\kappa']}=O^{\,g^3_{\kappa'_3\leftarrow\kappa_3}g^2_{\kappa'_2\leftarrow\kappa_2}g^1_{\kappa'_1\leftarrow\kappa_1}} .
\]
In Steps 1 and 2 the element of the group generated by the reflections $\theta_{x_a,c}$, $c\in o_a+h\Z$, that maps the layer $j$ onto the layer $j'$ is $g^a_{j'\leftarrow j}$;
so $\fold_{x_a}G_j=\prod_{j'=0}^{2r-1}G_j^{\,g^a_{j'\leftarrow j}}$ there.
Let $O\in\mathcal{T}_{\C}$ read in $Q_\kappa$, fix $j$, and let $G_j:=\prod_{\kappa':\kappa'_1=j}O^{[\kappa']}$ be the content of the layer $j$ of the direction $e_1$.
For $\kappa'_1=j$ one has $(O^{[\kappa']})^{g^1_{j'\leftarrow j}}=O^{[(j',\kappa'_2,\kappa'_3)]}$, by \eqref{eq:setting-group-action} and the uniqueness of the element of $W(\mathcal{Q})$ that maps $Q_\kappa$ onto a given cell.
Hence
\[
\fold_{x_1}G_j=\prod_{j'=0}^{2r-1}\ \prod_{\kappa':\kappa'_1=j}O^{[(j',\kappa'_2,\kappa'_3)]}=\prod_{\kappa'}O^{[\kappa']}=\Fold_{\mathbb{T}}O,
\]
and $Z^{(0)}_{\mathbb{T}}[\Fold_{\mathbb{T}}O]\ge0$ by \eqref{eq:averaged-slab-holder}.

\emph{Step 4: three directions.}
The estimate in one direction is applied in the directions $e_1$, $e_2$, $e_3$ in turn, in the scheme of Proposition~\ref{prop:three-directions}, the torus being the same at each step.
Let $O_\kappa$ be as in (b).
The content of the layer $j_1$ of the direction $e_1$ is $\prod_{\kappa:\kappa_1=j_1}O_\kappa$, and its fold is $\prod_\kappa(O_{(j_1,\kappa_2,\kappa_3)})^{g^1_{\kappa_1\leftarrow j_1}}$.
So Step 1 gives
\[
\left|Z^{(0)}_{\mathbb{T}}\left[\,\prod_\kappa O_\kappa\right]\right|\le\prod_{0\le j_1<n_1/h}Z^{(0)}_{\mathbb{T}}\left[\,\prod_\kappa\bigl(O_{(j_1,\kappa_2,\kappa_3)}\bigr)^{g^1_{\kappa_1\leftarrow j_1}}\right]^{\frac{h}{n_1}} .
\]
For fixed $j_1$ the factors $(O_{(j_1,\kappa_2,\kappa_3)})^{g^1_{\kappa_1\leftarrow j_1}}$ lie in $\mathcal{T}_{\C}$ and read in $Q_\kappa$.
Step 2 with $\mu=2$, applied to the contents of the layers of the direction $e_2$, and the rule $(G^g)^{g'}=G^{g'g}$ of \eqref{eq:setting-group-action} give
\[
Z^{(0)}_{\mathbb{T}}\left[\,\prod_\kappa\bigl(O_{(j_1,\kappa_2,\kappa_3)}\bigr)^{g^1_{\kappa_1\leftarrow j_1}}\right]\le\prod_{0\le j_2<n_2/h}Z^{(0)}_{\mathbb{T}}\left[\,\prod_\kappa\bigl(O_{(j_1,j_2,\kappa_3)}\bigr)^{g^2_{\kappa_2\leftarrow j_2}g^1_{\kappa_1\leftarrow j_1}}\right]^{\frac{h}{n_2}},
\]
and Step 2 with $\mu=3$, applied to the contents of the layers of the direction $e_3$, gives for fixed $j_1$, $j_2$
\[
\begin{aligned}
&Z^{(0)}_{\mathbb{T}}\left[\,\prod_\kappa\bigl(O_{(j_1,j_2,\kappa_3)}\bigr)^{g^2_{\kappa_2\leftarrow j_2}g^1_{\kappa_1\leftarrow j_1}}\right]\\
&\qquad\le\prod_{0\le j_3<n_3/h}Z^{(0)}_{\mathbb{T}}\left[\,\prod_\kappa\bigl(O_{(j_1,j_2,j_3)}\bigr)^{g^3_{\kappa_3\leftarrow j_3}g^2_{\kappa_2\leftarrow j_2}g^1_{\kappa_1\leftarrow j_1}}\right]^{\frac{h}{n_3}}\\
&\qquad=\prod_{0\le j_3<n_3/h}Z^{(0)}_{\mathbb{T}}\bigl[\Fold_{\mathbb{T}}O_{(j_1,j_2,j_3)}\bigr]^{\frac{h}{n_3}} .
\end{aligned}
\]
All numbers on the right are non-negative by Step 3.
The triples $(j_1,j_2,j_3)$ range over the same set as $\kappa$, and $\frac{h}{n_1}\cdot\frac{h}{n_2}\cdot\frac{h}{n_3}=\frac{h^3}{|n|}=\frac{|Q|}{|\mathbb{T}|}$;
combining the three inequalities gives the estimate of (b).
\end{proof}

\subsection{Domination and flux sectors}\label{sec:domination-sectors}

The Schatten classes $\mathscr{J}_p$, $1\le p<\infty$, of \cite[Section~10.3.1]{NoncompactQED} and the inequalities of \cite[Lemma~10.6]{NoncompactQED} concern the operators of one separable Hilbert space.
Operators between several spaces, such as the slice operator $\mathfrak{M}\colon\mathcal{H}^{\mathrm{s}}\to\mathcal{H}^{\mathrm{l}}$ of Definition~\ref{def:slice-operators}, are treated on the direct sum of the spaces.

\begin{lemma}[Operators between several Hilbert spaces]\label{lem:several-spaces}
Let $r\ge1$, let $\mathcal{H}_1,\dots,\mathcal{H}_r$ be separable Hilbert spaces, and let $\mathcal{H}:=\mathcal{H}_1\oplus\dots\oplus\mathcal{H}_r$ be their orthogonal direct sum.
A bounded operator $\mathfrak{Y}\colon\mathcal{H}_a\to\mathcal{H}_b$ is identified with the operator on $\mathcal{H}$ that agrees with $\mathfrak{Y}$ on $\mathcal{H}_a$ and vanishes on $\mathcal{H}_c$ for every $c\neq a$.
For such an operator and $1\le p<\infty$ put
\[
\|\mathfrak{Y}\|_p:=\left(\Tr_{\mathcal{H}_a}\left(\mathfrak{Y}^{*}\cdot\mathfrak{Y}\right)^{\frac{p}{2}}\right)^{\frac{1}{p}}\in[0,\infty];
\]
for even $p$ and a Hilbert--Schmidt operator $\mathfrak{Y}$ this is the norm of Definition~\ref{def:even-schatten}.
\begin{enumerate}
\item\label{it:several-spaces-norms}
The identification is compatible with products of composable operators, with adjoints and with operator norms.
The operator on $\mathcal{H}$ that corresponds to $\mathfrak{Y}$ has the same non-zero singular values as $\mathfrak{Y}$, with the same multiplicities; it belongs to $\mathscr{J}_p$ if and only if $\|\mathfrak{Y}\|_p<\infty$, and its Schatten norm is then $\|\mathfrak{Y}\|_p$.
\item\label{it:several-spaces-traces}
An operator of trace class on $\mathcal{H}_a$ corresponds to an operator of trace class on $\mathcal{H}$ with the same trace.
Let $\mathfrak{Y}\colon\mathcal{H}_a\to\mathcal{H}_b$ be bounded, and let $U$ and $U'$ be unitary operators on $\mathcal{H}_a$ and on $\mathcal{H}_b$.
For $1\le p'\le p<\infty$ one has $\|\mathfrak{Y}\|_p\le\|\mathfrak{Y}\|_{p'}$.
If $\|\mathfrak{Y}\|_p<\infty$, then $\|\mathfrak{Y}^{*}\|_p=\|U'\cdot\mathfrak{Y}\cdot U\|_p=\|\mathfrak{Y}\|_p$, and for even $p$
\[
\Tr_{\mathcal{H}_a}\left(\mathfrak{Y}^{*}\cdot\mathfrak{Y}\right)^{\frac{p}{2}}=\Tr_{\mathcal{H}_b}\left(\mathfrak{Y}\cdot\mathfrak{Y}^{*}\right)^{\frac{p}{2}}.
\]
\item\label{it:several-spaces-holder}
(H\"older's inequality for closed chains.)
Let $k\ge1$, let $a_0,a_1,\dots,a_k\in\{1,\dots,r\}$ with $a_k=a_0$, let $p_1,\dots,p_k\in[1,\infty)$ with $\sum_{i=1}^k\frac{1}{p_i}=1$, and for $i=1,\dots,k$ let $\mathfrak{Y}_i\colon\mathcal{H}_{a_{i-1}}\to\mathcal{H}_{a_i}$ be a bounded operator with $\|\mathfrak{Y}_i\|_{p_i}<\infty$.
Then $\mathfrak{Y}_k\cdots\mathfrak{Y}_1$ is of trace class on $\mathcal{H}_{a_0}$, and
\[
\left|\Tr_{\mathcal{H}_{a_0}}\left(\mathfrak{Y}_k\cdots\mathfrak{Y}_1\right)\right|\le\prod_{i=1}^k\|\mathfrak{Y}_i\|_{p_i}.
\]
\end{enumerate}
\end{lemma}

\begin{proof}
An element of $\mathcal{H}$ is a family $x=(x_c)_{c=1,\dots,r}$ with $x_c\in\mathcal{H}_c$.
The operator that corresponds to $\mathfrak{Y}\colon\mathcal{H}_a\to\mathcal{H}_b$ maps $x$ to the family whose component in $\mathcal{H}_b$ is $\mathfrak{Y}x_a$ and whose other components vanish.

\eqref{it:several-spaces-norms}
If $\mathfrak{Y}'\colon\mathcal{H}_b\to\mathcal{H}_c$ is a second bounded operator, the product of the two corresponding operators maps $x$ to the family with the component $\mathfrak{Y}'\mathfrak{Y}x_a$ in $\mathcal{H}_c$ and zeros elsewhere; it corresponds to $\mathfrak{Y}'\cdot\mathfrak{Y}$.
The adjoint of the operator that corresponds to $\mathfrak{Y}$ maps $\mathcal{H}_b$ into $\mathcal{H}_a$ by $\mathfrak{Y}^{*}$ and vanishes on the other summands, so it corresponds to $\mathfrak{Y}^{*}$.
The norm of the operator that corresponds to $\mathfrak{Y}$ is $\|\mathfrak{Y}\|$, because $\|\mathfrak{Y}x_a\|\le\|\mathfrak{Y}\|\cdot\|x\|$, with equality approached by vectors of $\mathcal{H}_a$.
Hence $\mathfrak{Y}^{*}\cdot\mathfrak{Y}$, taken on $\mathcal{H}$, is the non-negative operator that agrees with $\mathfrak{Y}^{*}\cdot\mathfrak{Y}$ on $\mathcal{H}_a$ and vanishes on the other summands; by the functional calculus the same holds for its powers with positive exponents, in particular for its square root.
So the moduli of $\mathfrak{Y}$ and of the corresponding operator have the same non-zero eigenvalues, with the same multiplicities, and the trace on $\mathcal{H}$ of the $p$-th power of the latter modulus is $\Tr_{\mathcal{H}_a}\left(\mathfrak{Y}^{*}\cdot\mathfrak{Y}\right)^{\frac{p}{2}}=\|\mathfrak{Y}\|_p^p$.
If this number is finite, then $\left(\mathfrak{Y}^{*}\cdot\mathfrak{Y}\right)^{\frac{p}{2}}$ is of trace class, its powers with positive exponents, among them $\mathfrak{Y}^{*}\cdot\mathfrak{Y}$, are compact, and so is $\mathfrak{Y}$; this gives the assertion on $\mathscr{J}_p$.

\eqref{it:several-spaces-traces}
Take an orthonormal basis of $\mathcal{H}$ that is the union of orthonormal bases of $\mathcal{H}_1,\dots,\mathcal{H}_r$.
Let $V$ be an operator of trace class on $\mathcal{H}_a$.
By \eqref{it:several-spaces-norms} the operator on $\mathcal{H}$ that corresponds to $V$ is of trace class, and its diagonal entries are those of $V$ at the basis vectors in $\mathcal{H}_a$ and $0$ at the others; so the two traces agree.
The operator that corresponds to $U'\cdot\mathfrak{Y}\cdot U$ is the product of the operator that corresponds to $\mathfrak{Y}$ with two unitary operators of $\mathcal{H}$: on the left the one that acts as $U'$ on $\mathcal{H}_b$ and as the identity on the other summands, on the right the one that acts as $U$ on $\mathcal{H}_a$ and as the identity on the other summands.
Now \eqref{it:several-spaces-norms} and \cite[Lemma~10.6\,(b)]{NoncompactQED}, applied on $\mathcal{H}$, give $\|\mathfrak{Y}\|_p\le\|\mathfrak{Y}\|_{p'}$ and, if $\|\mathfrak{Y}\|_p<\infty$, the equalities $\|\mathfrak{Y}^{*}\|_p=\|U'\cdot\mathfrak{Y}\cdot U\|_p=\|\mathfrak{Y}\|_p$.
By the definition of the norms, the equality $\|\mathfrak{Y}^{*}\|_p^p=\|\mathfrak{Y}\|_p^p$ reads $\Tr_{\mathcal{H}_b}\left(\mathfrak{Y}\cdot\mathfrak{Y}^{*}\right)^{\frac{p}{2}}=\Tr_{\mathcal{H}_a}\left(\mathfrak{Y}^{*}\cdot\mathfrak{Y}\right)^{\frac{p}{2}}$.

\eqref{it:several-spaces-holder}
By \eqref{it:several-spaces-norms} the operators that correspond to $\mathfrak{Y}_1,\dots,\mathfrak{Y}_k$ belong to $\mathscr{J}_{p_1},\dots,\mathscr{J}_{p_k}$, with the Schatten norms $\|\mathfrak{Y}_1\|_{p_1},\dots,\|\mathfrak{Y}_k\|_{p_k}$, and their product corresponds to the operator $\mathfrak{Y}_k\cdots\mathfrak{Y}_1$ on $\mathcal{H}_{a_0}$.
By \cite[Lemma~10.6\,(a)]{NoncompactQED} this product is of trace class on $\mathcal{H}$, and the modulus of its trace is at most $\prod_i\|\mathfrak{Y}_i\|_{p_i}$.
By \eqref{it:several-spaces-norms} the operator $\mathfrak{Y}_k\cdots\mathfrak{Y}_1$ has the same non-zero singular values and is therefore of trace class on $\mathcal{H}_{a_0}$, and by \eqref{it:several-spaces-traces} its trace equals the trace on $\mathcal{H}$ of the corresponding operator.
\end{proof}

\begin{proposition}[Gaussian domination for the averaged law]\label{prop:averaged-domination}
Let $-1\le\rho\le1$, $w_0,\dots,w_3>0$, $M\in\R$ and $q>0$, and let $\mathbb{T}=C(n_1,n_2,n_3)$ with $n_1,n_2,n_3\ge4$.
Then, for every $\beta\in C^2(\mathbb{T})$,
\[
0<Z^{(0)}_{\mathbb{T}}\left[\exp\left(-\langle\beta,F\rangle-\frac{1}{2}\|\beta\|^2\right)\right]\le Z^{(0)}_{\mathbb{T}}[1].
\]
Consequently $Z^{(0)}_{\mathbb{T}}\left[\exp\langle\beta,F\rangle\right]\le\exp\left(\frac{1}{2}\|\beta\|^2\right)\cdot Z^{(0)}_{\mathbb{T}}[1]$ for every $\beta\in C^2(\mathbb{T})$: under the averaged law $\mu^{\mathrm{av}}_{\mathbb{T}}$ the expectation of $\exp\langle\beta,F\rangle$ is at most $\exp\left(\frac{1}{2}\|\beta\|^2\right)$.
\end{proposition}

\begin{proof}
The proof follows that of \cite[Theorem~10.7]{NoncompactQED}, with the Bloch angle of Lemma~\ref{lem:bloch} and with a truncation of the magnetic factors.
For $\beta\in C^2(\mathbb{T})$ let $\mathcal{Z}(\beta)$ denote the middle term of the assertion.
Since $\Phi_0\le1$ and the integrand is positive, $\mathcal{Z}(\beta)\le Z_{\mathbb{T}}\left[\exp\left(-\langle\beta,F\rangle-\frac{1}{2}\|\beta\|^2\right)\right]$;
the right side is at most $Z_{\mathbb{T}}[1]$ by Theorem~\ref{thm:coordinate-domination}, and $Z_{\mathbb{T}}[1]<\infty$ by Theorem~\ref{thm:partition-positivity}.
So $\mathcal{Z}(\beta)$ is finite.

\emph{Positivity.}
As $\|F\|^2+2\langle\beta,F\rangle+\|\beta\|^2=\|F+\beta\|^2$, Definitions~\ref{def:configurations} and~\ref{def:functional} give
\[
\mathcal{Z}(\beta)=\sum_z\int_{[0,2\pi/q)^{\mathbb{T}^{(1)}}}\mathrm{d}A\,\exp\left(-\frac{1}{2}\|F+\beta\|^2\right)\cdot\det D^{\mathbb{T}}_A\cdot\Phi_0(F),\qquad F=dA+\frac{2\pi}{q}z,
\]
the sum being taken over the closed $z\in C^2(\mathbb{T};\Z)$.
The exponential factor is positive, and $\det D^{\mathbb{T}}_A\ge0$ by Proposition~\ref{prop:accretivity}.
For $\beta=0$ the right side is $Z^{(0)}_{\mathbb{T}}[1]$, the sum of the flux terms of $Z_{\mathbb{T}}[1]$ that belong to the closed $z$ of flux zero, $z=0$ among them; and every flux term of $Z_{\mathbb{T}}[1]$ is positive (Lemma~\ref{lem:flux-terms} and Theorem~\ref{thm:closed-interval}).
Hence $\det D^{\mathbb{T}}_A\cdot\Phi_0(F)$ is positive on a set of positive measure, and $\mathcal{Z}(\beta)>0$.

\emph{The inequality with $e_1$ as time.}
Let $\beta\in C^2(\mathbb{T})$.
For $j\in\Z/n_1\Z$ let $\beta^{\mathrm{el}}_j\in\R^{\Lambda^{(1)}_{\mathrm{s}}}$ and $\beta^{\mathrm{mag}}_j\in\R^{\Lambda^{(2)}_{\mathrm{s}}}$ be the electric part of $\beta$ at the step $j$ and its magnetic part at the slice $j$, in the notation of Definition~\ref{def:slices}: $\beta^{\mathrm{el}}_j(y,\nu):=\beta_{1\nu}(j,y)$ and $\beta^{\mathrm{mag}}_j(y):=\beta_{23}(j,y)$.
Every plaquette of $\mathbb{T}$ is either a $(1\nu)$-plaquette at a step or a $(23)$-plaquette at a slice; hence
\[
\langle\beta,F\rangle=\sum_j\left(\langle\beta^{\mathrm{el}}_j,F^{\mathrm{el}}_j\rangle+\langle\beta^{\mathrm{mag}}_j,F^{\mathrm{mag}}_j\rangle\right),\qquad\|\beta\|^2=\sum_j\left(|\beta^{\mathrm{el}}_j|^2+|\beta^{\mathrm{mag}}_j|^2\right).
\]
On $\R^{\Lambda^{(2)}_{\mathrm{s}}}$ let $u_j:=\exp\left(-\langle\beta^{\mathrm{mag}}_j,\,\cdot\,\rangle-\frac{1}{2}|\beta^{\mathrm{mag}}_j|^2\right)$ and, for $\varepsilon>0$, $u^{\varepsilon}_j:=\min\left(u_j,\frac{1}{\varepsilon}\right)$; on $\R^{\Lambda^{(1)}_{\mathrm{s}}}$ let $v_j(\xi):=\exp\left(-\langle\beta^{\mathrm{el}}_j,\xi\rangle-\frac{1}{2}|\beta^{\mathrm{el}}_j|^2\right)$.
Then
\[
\exp\left(-\langle\beta,F\rangle-\frac{1}{2}\|\beta\|^2\right)=\prod_ju_j(F^{\mathrm{mag}}_j)\cdot v_j(F^{\mathrm{el}}_j).
\]
The functions $u^{\varepsilon}_j$ are bounded and Borel, and $v_j(\xi)\le\exp\left(|\beta^{\mathrm{el}}_j|\cdot|\xi|\right)$, so that $v_j$ satisfies the growth condition of the word formula (Proposition~\ref{prop:word-formula}).
Hence the Bloch representation \eqref{eq:averaged-bloch} applies to the integrand with $u^{\varepsilon}_j$ in place of $u_j$:
\[
\begin{aligned}
&Z^{(0)}_{\mathbb{T}}\left[\prod_ju^{\varepsilon}_j(F^{\mathrm{mag}}_j)\cdot v_j(F^{\mathrm{el}}_j)\right]\\
&\qquad=(2\pi)^{|n|}\cdot\int_{[0,1]^2}\mathrm{d}\vartheta\,\Tr\left(P_0\cdot\prod_{j=n_1-1,\dots,0}\left((2\pi/q)^{|\Lambda_{\mathrm{s}}|}\cdot K^{\vartheta}_j\cdot P_{\mathrm{G}}\cdot\mathfrak{M}\cdot u^{\varepsilon}_j(F^{\mathrm{mag}})\cdot\mathfrak{M}^{*}\right)\right).
\end{aligned}
\]
Here $K^{\vartheta}_j$ is the operator of the step $j$ formed with $v_j(\xi)\cdot\exp\left(i\langle\beta_\vartheta,\xi\rangle\right)$; its kernel is the periodisation over $\Gamma_{\mathrm{s}}$ of the function $A'-A\mapsto v_j(A'-A)\cdot\exp\left(i\langle\beta_\vartheta,A'-A\rangle\right)\cdot\prod_bg_1\left(A'(b)-A(b)\right)$, the product being taken over the links $b$ of $\Lambda_{\mathrm{s}}$.
Since $\exp\left(-i\langle\xi,E\rangle\right)$ is the translation by $\xi$ (Definition~\ref{def:slice-operators}), an operator with such a periodised kernel is an integral of translations:
\[
K^{\vartheta}_j=\int_{\R^{\Lambda^{(1)}_{\mathrm{s}}}}\mathrm{d}\xi\,v_j(\xi)\cdot\exp\left(i\langle\beta_\vartheta,\xi\rangle\right)\cdot\prod_bg_1(\xi_b)\cdot\exp\left(-i\langle\xi,E\rangle\right).
\]
Moreover $v_j(\xi)\cdot\prod_bg_1(\xi_b)=\prod_bg_1\left(\xi_b+\beta^{\mathrm{el}}_j(b)\right)$.
With the complex number $z_j:=\exp\left(-i\langle\beta_\vartheta,\beta^{\mathrm{el}}_j\rangle\right)$ of modulus one it follows that
\[
\begin{aligned}
&v_j(\xi)\cdot\exp\left(i\langle\beta_\vartheta,\xi\rangle\right)\cdot\prod_bg_1(\xi_b)\\
&\qquad=z_j\cdot\exp\left(i\langle\beta_\vartheta,\xi+\beta^{\mathrm{el}}_j\rangle\right)\cdot\prod_bg_1\left(\xi_b+\beta^{\mathrm{el}}_j(b)\right).
\end{aligned}
\]
Let $U_j:=\exp\left(i\langle\beta^{\mathrm{el}}_j,E\rangle\right)$.
The substitution of $\xi-\beta^{\mathrm{el}}_j$ for $\xi$ in the integral, and the Fourier transform of the Gaussian, give
\[
\begin{aligned}
K^{\vartheta}_j&=z_j\cdot U_j\cdot\int_{\R^{\Lambda^{(1)}_{\mathrm{s}}}}\mathrm{d}\xi\,\exp\left(i\langle\beta_\vartheta,\xi\rangle\right)\cdot\prod_bg_1(\xi_b)\cdot\exp\left(-i\langle\xi,E\rangle\right)\\
&=z_j\cdot U_j\cdot\exp\left(-\frac{1}{2}|E-\beta_\vartheta|^2\right).
\end{aligned}
\]
The operator $U_j$ is unitary; as a function of $E$ it commutes with the translations of the slice variables and with constant matrices, hence with $P_{\mathrm{G}}$ and with $K_{\beta_\vartheta}^{\frac{1}{2}}$.
By Definition~\ref{def:slice-operators},
\[
(2\pi/q)^{|\Lambda_{\mathrm{s}}|}\cdot K^{\vartheta}_j\cdot P_{\mathrm{G}}=z_j\cdot U_j\cdot K_{\beta_\vartheta}=z_j\cdot K_{\beta_\vartheta}^{\frac{1}{2}}\cdot U_j\cdot K_{\beta_\vartheta}^{\frac{1}{2}}.
\]
The operators $U_j$, $K_{\beta_\vartheta}^{\frac{1}{2}}$, $\mathfrak{M}$ and the multiplication operators $u^{\varepsilon}_j(F^{\mathrm{mag}})$ are diagonal in the flux index (Definition~\ref{def:slice-operators}) and commute with $P_0$.
Put
\[
\mathfrak{Y}_j:=K_{\beta_\vartheta}^{\frac{1}{2}}\cdot\mathfrak{M}\cdot\left(u^{\varepsilon}_j\right)^{\frac{1}{2}}(F^{\mathrm{mag}})\colon\mathcal{H}^{\mathrm{s}}\to\mathcal{H}^{\mathrm{l}};
\]
it is a Hilbert--Schmidt operator, since $K_{\beta_\vartheta}^{\frac{1}{2}}\cdot\mathfrak{M}$ is one (Corollary~\ref{cor:trace-class}) and $u^{\varepsilon}_j$ is bounded.
When the last expression for the operators $(2\pi/q)^{|\Lambda_{\mathrm{s}}|}\cdot K^{\vartheta}_j\cdot P_{\mathrm{G}}$ is inserted into the trace and the first factor $K_{\beta_\vartheta}^{\frac{1}{2}}$ of the product under the trace is moved to the end by the cyclicity of the trace (the product of the remaining factors is of trace class), the trace in the integrand becomes, since $P_0=P_0^2$ commutes with all factors,
\[
\left(\prod_jz_j\right)\cdot\Tr\left(\prod_{j=n_1-1,\dots,0}\left(U_j\cdot P_0\cdot\mathfrak{Y}_j\cdot\left(P_0\cdot\mathfrak{Y}_j\right)^{*}\right)\right),
\]
the trace of a closed chain of the $2n_1$ operators $U_j\cdot P_0\cdot\mathfrak{Y}_j\colon\mathcal{H}^{\mathrm{s}}\to\mathcal{H}^{\mathrm{l}}$ and $\left(P_0\cdot\mathfrak{Y}_j\right)^{*}\colon\mathcal{H}^{\mathrm{l}}\to\mathcal{H}^{\mathrm{s}}$.
By Lemma~\ref{lem:several-spaces}\eqref{it:several-spaces-traces} each of these operators has the norm $\|P_0\cdot\mathfrak{Y}_j\|_{2n_1}\le\|P_0\cdot\mathfrak{Y}_j\|_2<\infty$, and Lemma~\ref{lem:several-spaces}\eqref{it:several-spaces-holder}, with the $2n_1$ exponents $2n_1$, bounds the modulus of the trace by
\[
\prod_j\|P_0\cdot\mathfrak{Y}_j\|_{2n_1}^2=\prod_j\left(\Tr_{\mathcal{H}^{\mathrm{s}}}\left(P_0\cdot\left(\mathfrak{Y}_j^{*}\cdot\mathfrak{Y}_j\right)^{n_1}\right)\right)^{\frac{1}{n_1}}.
\]
Now fix $j$.
In the Bloch representation \eqref{eq:averaged-bloch} of the integrand $\prod_{j'}u^{\varepsilon}_j(F^{\mathrm{mag}}_{j'})$, which carries the function $u^{\varepsilon}_j$ at every slice and the electric factors $1$, the operator $(2\pi/q)^{|\Lambda_{\mathrm{s}}|}\cdot K^{\vartheta}_{j'}\cdot P_{\mathrm{G}}$ equals $K_{\beta_\vartheta}$ for every $j'$, by the case $\beta^{\mathrm{el}}_j=0$ of the computation above.
Together with the cyclicity of the trace on $\mathcal{H}^{\mathrm{l}}$ and with Lemma~\ref{lem:several-spaces}\eqref{it:several-spaces-traces}, applied to $P_0\cdot\mathfrak{Y}_j$, it gives
\[
(2\pi)^{|n|}\cdot\int_{[0,1]^2}\mathrm{d}\vartheta\,\Tr_{\mathcal{H}^{\mathrm{s}}}\left(P_0\cdot\left(\mathfrak{Y}_j^{*}\cdot\mathfrak{Y}_j\right)^{n_1}\right)=Z^{(0)}_{\mathbb{T}}\left[\prod_{j'}u^{\varepsilon}_j(F^{\mathrm{mag}}_{j'})\right]\le\mathcal{Z}\left(\beta^{(j)}\right),
\]
where $\beta^{(j)}\in C^2(\mathbb{T})$ is the two-cochain with the electric parts $0$ and the magnetic part $\beta^{\mathrm{mag}}_j$ at every slice; the inequality holds because $u^{\varepsilon}_j\le u_j$ and $\exp\left(-\langle\beta^{(j)},F\rangle-\frac{1}{2}\|\beta^{(j)}\|^2\right)=\prod_{j'}u_j(F^{\mathrm{mag}}_{j'})$.
Since $|z_j|=1$, the Bloch representation of the integrand, the bound for the trace and H\"older's inequality for the integral over $\vartheta$, with $n_1$ exponents $n_1$, combine with the last display to
\[
\begin{aligned}
&Z^{(0)}_{\mathbb{T}}\left[\prod_ju^{\varepsilon}_j(F^{\mathrm{mag}}_j)\cdot v_j(F^{\mathrm{el}}_j)\right]\\
&\qquad\le\prod_j\left((2\pi)^{|n|}\cdot\int_{[0,1]^2}\mathrm{d}\vartheta\,\Tr_{\mathcal{H}^{\mathrm{s}}}\left(P_0\cdot\left(\mathfrak{Y}_j^{*}\cdot\mathfrak{Y}_j\right)^{n_1}\right)\right)^{\frac{1}{n_1}}\\
&\qquad\le\prod_j\mathcal{Z}\left(\beta^{(j)}\right)^{\frac{1}{n_1}}.
\end{aligned}
\]
When $\varepsilon$ decreases to $0$, $u^{\varepsilon}_j$ increases to $u_j$, and the monotone convergence theorem yields
\begin{equation}\label{eq:averaged-direction}
\mathcal{Z}(\beta)\le\prod_{j\in\Z/n_1\Z}\mathcal{Z}\left(\beta^{(j)}\right)^{\frac{1}{n_1}}.
\end{equation}

\emph{The inequality with $e_2$ as time.}
Lemma~\ref{lem:bloch} holds with $e_2$ as time, and with it the argument above; it proves \eqref{eq:averaged-direction} with $e_2$ in place of $e_1$, the electric parts of a two-cochain being then its $(12)$- and $(23)$-components and its magnetic parts its $(13)$-components.
The two-cochain $\beta^{(j)}$ has only $(23)$-components.
For the time $e_2$ its magnetic parts vanish, so that all the two-cochains formed from $\beta^{(j)}$ as $\beta^{(j)}$ was formed from $\beta$ are zero, and
\[
\mathcal{Z}\left(\beta^{(j)}\right)\le\prod_{j'\in\Z/n_2\Z}\mathcal{Z}(0)^{\frac{1}{n_2}}=\mathcal{Z}(0)=Z^{(0)}_{\mathbb{T}}[1].
\]
With \eqref{eq:averaged-direction} it follows that $\mathcal{Z}(\beta)\le Z^{(0)}_{\mathbb{T}}[1]$.

\emph{The averaged law.}
For $\beta\in C^2(\mathbb{T})$ one has $Z^{(0)}_{\mathbb{T}}\left[\exp\langle\beta,F\rangle\right]=\exp\left(\frac{1}{2}\|\beta\|^2\right)\cdot\mathcal{Z}(-\beta)$, which is at most $\exp\left(\frac{1}{2}\|\beta\|^2\right)\cdot Z^{(0)}_{\mathbb{T}}[1]$.
Division by $Z^{(0)}_{\mathbb{T}}[1]>0$ gives the bound for $\mu^{\mathrm{av}}_{\mathbb{T}}$, which is the law of $\widetilde F$ under $Z^{(0)}_{\mathbb{T}}[\,\cdot\,]/Z^{(0)}_{\mathbb{T}}[1]$ (Definition~\ref{def:averaged-law}).
\end{proof}

\begin{lemma}[Weights of the flux sectors]\label{lem:sector-weights}
Let $-1\le\rho\le1$, $w_0,\dots,w_3>0$, $M\in\R$ and $q>0$, and let $\mathbb{T}=C(n_1,n_2,n_3)$ with $n_1,n_2,n_3\ge4$.
\begin{enumerate}
\item\label{it:sector-weights-gaussian}
For every $k\in\Z^3$,
\[
\mathbb{E}_{\mathbb{T}}[\Phi_k]\le\exp\left(-\frac{1}{2}\|\mathfrak{h}_k\|^2\right).
\]
\item\label{it:sector-weights-zero}
$Z^{(k)}_{\mathbb{T}}[1]\le Z^{(0)}_{\mathbb{T}}[1]$ for every $k\in\Z^3$; consequently $Z^{(0)}_{\mathbb{T}}[1]>0$.
\item\label{it:sector-weights-cubic}
For an integer $n\ge4$ let $a:=\frac{2\pi^2}{q^2n}$.
Then, on the cubic torus $C(n)$ and for every integer $r\ge0$,
\[
\begin{aligned}
\mathbb{E}_{C(n)}[\Phi_0]&\ge\frac{1}{(2r+1)^3}\cdot\left(1-6\left(\sum_{j\in\Z}\exp(-aj^2)\right)^3\cdot\exp\left(-a(r+1)^2\right)\right),\\
\sum_{j\in\Z}\exp(-aj^2)&\le1+\left(\frac{\pi}{a}\right)^{\frac{1}{2}}.
\end{aligned}
\]
There is an integer $n_0(q)\ge4$, depending on $q$ only, such that
\[
-\log\mathbb{E}_{C(n)}[\Phi_0]\le\log\left(2(2n+1)^3\right)\qquad\text{for every integer }n\ge n_0(q).
\]
\end{enumerate}
\end{lemma}

\begin{proof}
Every configuration lies in exactly one flux sector (Definition~\ref{def:averaged-law}); so $Z_{\mathbb{T}}[G]=\sum_{k\in\Z^3}Z^{(k)}_{\mathbb{T}}[G]$ for every Borel function $G\ge0$ of $F$.

\eqref{it:sector-weights-gaussian}
The argument is that of \cite[Corollary~10.8\,(d)]{NoncompactQED}, where the fermion weight is the determinant of the Wilson operator; here $\det D^{\mathbb{T}}_A$ takes its place.
Let $z\in C^2(\mathbb{T};\Z)$ be closed with flux $k$, let $A\in C^1(\mathbb{T})$ and $F=dA+\frac{2\pi}{q}z$.
By \cite[Lemma~10.2\,(a)]{NoncompactQED} (there the cochain $z$, its flux and $\mathfrak{h}_k$ are written $\nu$, $w=(w_{12},w_{13},w_{23})$ and $h_w$, with $w_{\mu\nu}$ in the place of $k_{\mu\nu}$), the two-cochain $F-\mathfrak{h}_k$ is exact, and exact two-cochains are orthogonal to the constant ones \cite[(10.1)]{NoncompactQED}.
Hence $\langle\mathfrak{h},F\rangle=\langle\mathfrak{h},\mathfrak{h}_k\rangle$ on the sector $k$ for every constant two-cochain $\mathfrak{h}$.
Let $k'\in\Z^3$.
Theorem~\ref{thm:coordinate-domination}, applied with the two-cochain $-\mathfrak{h}_{k'}$, and the decomposition over the sectors give
\[
\begin{aligned}
Z_{\mathbb{T}}[1]&\ge Z_{\mathbb{T}}\left[\exp\left(\langle\mathfrak{h}_{k'},F\rangle-\frac{1}{2}\|\mathfrak{h}_{k'}\|^2\right)\right]=\sum_{k\in\Z^3}\exp\left(\langle\mathfrak{h}_{k'},\mathfrak{h}_k\rangle-\frac{1}{2}\|\mathfrak{h}_{k'}\|^2\right)\cdot Z^{(k)}_{\mathbb{T}}[1]\\
&\ge\exp\left(\frac{1}{2}\|\mathfrak{h}_{k'}\|^2\right)\cdot Z^{(k')}_{\mathbb{T}}[1],
\end{aligned}
\]
all terms of the sum being non-negative.
Division by $Z_{\mathbb{T}}[1]>0$ (Theorem~\ref{thm:partition-positivity}) gives \eqref{it:sector-weights-gaussian} for the flux $k'$.

\eqref{it:sector-weights-zero}
Here the argument of \cite[Proposition~21.23\,(ii)]{NoncompactQED} is carried out with the Bloch representation.
Fix $k_{23}\in\Z$, and for $\vartheta\in[0,1]^2$ let
\[
\widehat{Z}(\vartheta):=\sum_{k_{12},k_{13}\in\Z}\exp\left(2\pi i\left(\vartheta_2\cdot k_{12}+\vartheta_3\cdot k_{13}\right)\right)\cdot Z^{(k_{23},k_{13},k_{12})}_{\mathbb{T}}[1].
\]
By \eqref{eq:averaged-fourier} with $u_j=v_j=1$, in which $(2\pi/q)^{|\Lambda_{\mathrm{s}}|}\cdot K^{\vartheta}_j\cdot P_{\mathrm{G}}=K_{\beta_\vartheta}$ for every $j$, by the cyclicity of the trace and by Lemma~\ref{lem:several-spaces}\eqref{it:several-spaces-traces}, applied to $P_{k_{23}}\cdot K_{\beta_\vartheta}^{\frac{1}{2}}\cdot\mathfrak{M}$,
\[
\widehat{Z}(\vartheta)=(2\pi)^{|n|}\cdot\Tr\left(P_{k_{23}}\cdot\left(K_{\beta_\vartheta}\cdot\mathfrak{M}\cdot\mathfrak{M}^{*}\right)^{n_1}\right)=(2\pi)^{|n|}\cdot\Tr\left(P_{k_{23}}\cdot\left(\mathfrak{M}^{*}\cdot K_{\beta_\vartheta}\cdot\mathfrak{M}\right)^{n_1}\right)\ge0,
\]
the last operator being non-negative because $P_{k_{23}}$ commutes with $\mathfrak{M}^{*}\cdot K_{\beta_\vartheta}\cdot\mathfrak{M}\ge0$.
The series converges absolutely and uniformly, as $\sum_kZ^{(k)}_{\mathbb{T}}[1]=Z_{\mathbb{T}}[1]<\infty$ (Theorem~\ref{thm:partition-positivity}); so $\widehat{Z}$ is continuous, and its Fourier coefficients are the numbers $Z^{(k_{23},k_{13},k_{12})}_{\mathbb{T}}[1]\ge0$.
Each of them is obtained by integrating the series term by term against $\exp\left(-2\pi i\left(\vartheta_2\cdot k_{12}+\vartheta_3\cdot k_{13}\right)\right)$, a function of modulus one;
as $\widehat{Z}\ge0$, the modulus of that integral is at most the integral of $\widehat{Z}$ itself, which is the coefficient of index $(0,0)$:
\[
\begin{aligned}
Z^{(k_{23},k_{13},k_{12})}_{\mathbb{T}}[1]&=\int_{[0,1]^2}\mathrm{d}\vartheta\,\exp\left(-2\pi i\left(\vartheta_2\cdot k_{12}+\vartheta_3\cdot k_{13}\right)\right)\cdot\widehat{Z}(\vartheta)\\
&\le\int_{[0,1]^2}\mathrm{d}\vartheta\,\widehat{Z}(\vartheta)=Z^{(k_{23},0,0)}_{\mathbb{T}}[1].
\end{aligned}
\]
With $e_2$ as time (Lemma~\ref{lem:bloch}) the flux index of the summands is, up to sign, $k_{13}$, and the Fourier variables are, up to signs, $k_{12}$ and $k_{23}$; the same argument at the flux index $0$ gives $Z^{(k_{23},0,0)}_{\mathbb{T}}[1]\le Z^{(0,0,0)}_{\mathbb{T}}[1]$.
Finally $\sum_kZ^{(k)}_{\mathbb{T}}[1]=Z_{\mathbb{T}}[1]>0$ by Theorem~\ref{thm:partition-positivity}; so some $Z^{(k)}_{\mathbb{T}}[1]$ is positive, and $Z^{(0)}_{\mathbb{T}}[1]\ge Z^{(k)}_{\mathbb{T}}[1]>0$.

\eqref{it:sector-weights-cubic}
On $C(n)$ one has $\frac{1}{2}\|\mathfrak{h}_k\|^2=a|k|^2$ with $|k|^2=k_{23}^2+k_{13}^2+k_{12}^2$ (Definition~\ref{def:averaged-law}).
Let $|k|_\infty:=\max\left(|k_{23}|,|k_{13}|,|k_{12}|\right)$.
By \eqref{it:sector-weights-zero} for the $(2r+1)^3$ fluxes with $|k|_\infty\le r$ and by \eqref{it:sector-weights-gaussian} for the others,
\[
1=\sum_{k\in\Z^3}\mathbb{E}_{C(n)}[\Phi_k]\le(2r+1)^3\cdot\mathbb{E}_{C(n)}[\Phi_0]+\sum_{|k|_\infty>r}\exp\left(-a|k|^2\right).
\]
The set $\{|k|_\infty>r\}$ is the union of the three sets on which one component of $k$ has modulus larger than $r$.
Since $(r+1+j)^2\ge(r+1)^2+j^2$ for $j\ge0$,
\[
\sum_{|j|>r}\exp(-aj^2)=2\sum_{j\ge0}\exp\left(-a(r+1+j)^2\right)\le2\exp\left(-a(r+1)^2\right)\cdot\sum_{j\in\Z}\exp(-aj^2),
\]
and therefore
\[
\begin{aligned}
\sum_{|k|_\infty>r}\exp\left(-a|k|^2\right)&\le3\left(\sum_{|j|>r}\exp(-aj^2)\right)\cdot\left(\sum_{j\in\Z}\exp(-aj^2)\right)^2\\
&\le6\left(\sum_{j\in\Z}\exp(-aj^2)\right)^3\cdot\exp\left(-a(r+1)^2\right).
\end{aligned}
\]
Together with the first display of this paragraph this is the lower bound for $\mathbb{E}_{C(n)}[\Phi_0]$.
Comparison with an integral gives $\sum_{j\in\Z}\exp(-aj^2)\le1+2\int_0^\infty\mathrm{d}x\,\exp(-ax^2)=1+\left(\frac{\pi}{a}\right)^{\frac{1}{2}}$.
For $r=n$ one has $a(n+1)^2\ge an^2=\frac{2\pi^2n}{q^2}$ and $\left(\frac{\pi}{a}\right)^{\frac{1}{2}}=q\cdot\left(\frac{n}{2\pi}\right)^{\frac{1}{2}}$, so that
\[
6\left(\sum_{j\in\Z}\exp(-aj^2)\right)^3\cdot\exp\left(-a(n+1)^2\right)\le6\left(1+q\cdot\left(\frac{n}{2\pi}\right)^{\frac{1}{2}}\right)^3\cdot\exp\left(-\frac{2\pi^2n}{q^2}\right).
\]
The right side depends on $q$ and $n$ only and tends to $0$ as $n\to\infty$.
Let $n_0(q)\ge4$ be an integer such that the right side is at most $\frac{1}{2}$ for all $n\ge n_0(q)$.
For these $n$ the lower bound with $r=n$ gives $\mathbb{E}_{C(n)}[\Phi_0]\ge\frac{1}{2}(2n+1)^{-3}$.
\end{proof}

\begin{remark}
Under the hypotheses of Lemma~\ref{lem:sector-weights}, every $Z^{(k)}_{\mathbb{T}}[1]$, $k\in\Z^3$, is positive.
Indeed, let $z$ be the integer two-cochain of $\mathbb{T}$ with the components
\[
\begin{aligned}
z_{12}(x)&=k_{12}\cdot\mathbf{1}[x_1\equiv x_2\equiv0],\\
z_{13}(x)&=k_{13}\cdot\mathbf{1}[x_1\equiv x_3\equiv0],\\
z_{23}(x)&=k_{23}\cdot\mathbf{1}[x_2\equiv x_3\equiv0],
\end{aligned}
\]
the congruences being taken modulo the sides of $\mathbb{T}$.
No $z_{\mu\nu}$ depends on the third coordinate, so the three brackets in the formula for $dz$ of Definition~\ref{def:cochains} vanish and $z$ is closed; the sum of $z_{\mu\nu}$ over a coordinate two-torus $\Sigma_{\mu\nu}$ is $k_{\mu\nu}$, so $z$ has the flux $k$.
The term of $z$ in $Z^{(k)}_{\mathbb{T}}[1]$ is a flux term of $Z_{\mathbb{T}}[1]$, hence positive by Lemma~\ref{lem:flux-terms} and Theorem~\ref{thm:closed-interval}, and the other terms are non-negative.
\end{remark}

\begin{proposition}[Convexity without monotonicity]\label{prop:no-monotonicity}
Let $w_0,\dots,w_3>0$ and $q>0$.
\begin{enumerate}
\item\label{it:no-monotonicity-convex}
Let $-1\le\rho\le1$, $M\in\R$ and $n_2,n_3\ge4$.
The function $n_1\mapsto\log Z^{(0)}_{C(n_1,n_2,n_3)}[1]$ is convex on the integers $n_1\ge4$:
\[
2\log Z^{(0)}_{C(n_1,n_2,n_3)}[1]\le\log Z^{(0)}_{C(n_1-1,n_2,n_3)}[1]+\log Z^{(0)}_{C(n_1+1,n_2,n_3)}[1]\qquad(n_1\ge5).
\]
\item\label{it:no-monotonicity-fails}
Let $0\le\rho\le1$ and $M_0:=\frac{1+\rho}{2}\sum_{i=0}^3w_i$.
For an integer $n\ge2$ put
\[
\operatorname{Ld}(n):=\sum_{\xi\in(\Z/n\Z)^3,\ \xi\neq0}\log\left(\sum_{\mu=1}^3\left(2-2\cos\frac{2\pi\xi_\mu}{n}\right)\right).
\]
Let $n\ge4$ and $j\ge2$ be integers with
\[
q^2>2\pi\cdot\exp\left(\frac{\operatorname{Ld}(jn)-j^3\cdot\operatorname{Ld}(n)}{2(j^3-1)}\right);
\]
this holds if $q^2n>2\pi\cdot\exp\left(\frac{C_j}{j^3-1}\right)$, where $C_j:=(1+j+j^2)\cdot\left(\log j+j\cdot\left(1+\frac{j}{3.2}\right)\cdot\frac{2j^2}{2.56}\right)$.
Then there is $M_1<\infty$ such that, for every $M\ge M_1$,
\[
j^3\cdot\log Z^{(0)}_{C(n)}[1]-\log Z^{(0)}_{C(jn)}[1]<0,\qquad\text{that is,}\qquad f^{(0)}(jn,jn,jn)>f^{(0)}(n,n,n).
\]
Consequently, for these parameters, $f^{(0)}$ does not share the monotonicity of the pressure $f$, which is non-increasing in each side (Proposition~\ref{prop:coordinate-pressure});
and the inequality $\log Z^{(0)}_{C(jn)}[1]\le j^3\cdot\log Z^{(0)}_{C(n)}[1]$, which would express subadditivity under the decomposition of $C(jn)$ into $j^3$ cubes of side $n$, is false.
\item\label{it:no-monotonicity-norm}
Let $-1\le\rho\le1$, $M\in\R$ and $n_1,n_2,n_3\ge4$.
Take $e_1$ as time on $C(n_1,n_2,n_3)$, and let $P_0$, $\mathfrak{M}$ and $K_{\beta_\vartheta}$, $\vartheta\in[0,1]^2$, be the operators of Lemma~\ref{lem:bloch}.
They are operators of the cross-section with the sides $n_2$, $n_3$, determined by $\rho$, $M$, $w_0,\dots,w_3$, $q$ and by the plaquette $\sigma_0$ and the site $y^0$ fixed in that lemma; none of them depends on $n_1$.
Then
\[
Z^{(0)}_{C(n_1,n_2,n_3)}[1]^{\frac{1}{n_1}}=(2\pi)^{n_2n_3}\cdot\left(\int_{[0,1]^2}\mathrm{d}\vartheta\,\Tr\left(P_0\cdot\left(\mathfrak{M}^{*}\cdot K_{\beta_\vartheta}\cdot\mathfrak{M}\right)^{n_1}\right)\right)^{\frac{1}{n_1}},
\]
which is $(2\pi)^{n_2n_3}$ times the norm in $L^{n_1}$ of the eigenvalues of the non-negative operator $P_0\cdot\mathfrak{M}^{*}\cdot K_{\beta_\vartheta}\cdot\mathfrak{M}$, with respect to the product of Lebesgue measure on $[0,1]^2$ and counting measure.
\end{enumerate}
\end{proposition}

\begin{proof}
Let $-1\le\rho\le1$, $M\in\R$ and $n_2,n_3\ge4$ be as in \eqref{it:no-monotonicity-convex} and \eqref{it:no-monotonicity-norm}, and let $P_0$, $\mathfrak{M}$ and $K_{\beta_\vartheta}$ be the operators of~\eqref{it:no-monotonicity-norm}.
For $\vartheta\in[0,1]^2$ the operator $\mathfrak{M}^{*}\cdot K_{\beta_\vartheta}\cdot\mathfrak{M}=\left(K_{\beta_\vartheta}^{\frac{1}{2}}\cdot\mathfrak{M}\right)^{*}\cdot K_{\beta_\vartheta}^{\frac{1}{2}}\cdot\mathfrak{M}$ on $\mathcal{H}^{\mathrm{s}}$ is non-negative and of trace class (Corollary~\ref{cor:trace-class}), and it commutes with $P_0$, all these operators being diagonal in the flux index (Definition~\ref{def:slice-operators}).
Let $\lambda_1(\vartheta)\ge\lambda_2(\vartheta)\ge\dots\ge0$ be the eigenvalues of $P_0\cdot\mathfrak{M}^{*}\cdot K_{\beta_\vartheta}\cdot\mathfrak{M}$, repeated according to multiplicity and followed by zeros if there are finitely many.
Then, for every real $p\ge1$,
\[
\Tr\left(P_0\cdot\left(\mathfrak{M}^{*}\cdot K_{\beta_\vartheta}\cdot\mathfrak{M}\right)^p\right)=\sum_{i\ge1}\lambda_i(\vartheta)^p<\infty.
\]
The slice operators, the projection $P_0$ and the cochain $\beta_\vartheta$ belong to the cross-section with the sides $n_2$, $n_3$ and do not depend on $n_1$.
For every $n_1\ge4$, the formula for $Z^{(0)}_{C(n_1,n_2,n_3)}[1]$ in Lemma~\ref{lem:bloch}\eqref{it:bloch-zero-flux}, which is the case $u_j=v_j=1$ of that clause, reads in terms of these eigenvalues
\[
Z^{(0)}_{C(n_1,n_2,n_3)}[1]=(2\pi)^{n_1n_2n_3}\cdot\int_{[0,1]^2}\mathrm{d}\vartheta\,\sum_{i\ge1}\lambda_i(\vartheta)^{n_1}.
\]

\eqref{it:no-monotonicity-norm}
The identity follows from the last display by taking the $n_1$-th root.
Each $\lambda_i$ is a continuous function of $\vartheta$, so that the norm in $L^{n_1}$ is well defined.
Indeed, the map $\beta\mapsto K_\beta^{\frac{1}{2}}=(2\pi/q)^{\frac{|\Lambda_{\mathrm{s}}|}{2}}\cdot\exp\left(-\frac{1}{4}|E-\beta|^2\right)\cdot P_{\mathrm{G}}$ is continuous in the operator norm, the Gaussian being Lipschitz continuous;
hence $\vartheta\mapsto P_0\cdot\mathfrak{M}^{*}\cdot K_{\beta_\vartheta}\cdot\mathfrak{M}$ is continuous in the operator norm,
and by the min-max principle the $i$-th eigenvalues of two non-negative compact operators differ by at most the norm of their difference.

\eqref{it:no-monotonicity-convex}
Let $p,p'\ge1$.
The Cauchy--Schwarz inequality, first for the sum over $i$ and then for the integral over $\vartheta$, gives
\[
\int_{[0,1]^2}\mathrm{d}\vartheta\,\sum_i\lambda_i(\vartheta)^{\frac{p+p'}{2}}\le\left(\int_{[0,1]^2}\mathrm{d}\vartheta\,\sum_i\lambda_i(\vartheta)^p\right)^{\frac{1}{2}}\cdot\left(\int_{[0,1]^2}\mathrm{d}\vartheta\,\sum_i\lambda_i(\vartheta)^{p'}\right)^{\frac{1}{2}}.
\]
Take $p=n_1-1$ and $p'=n_1+1$.
The three partition functions are positive by Lemma~\ref{lem:sector-weights}\eqref{it:sector-weights-zero}, and $\log\left((2\pi)^{n_1n_2n_3}\right)$ is linear in $n_1$; so the identity for $Z^{(0)}_{C(n_1,n_2,n_3)}[1]$ turns the last inequality into the asserted one.

\eqref{it:no-monotonicity-fails}
Let $\mathbb{T}$ be one of the cubic tori $C(n)$ and $C(jn)$.
The hopping part $M-D^{\mathbb{T}}_A=\frac{\rho}{2}\cdot H_A-\frac{1}{2}\cdot J_A$ of the chamber operator (Definition~\ref{def:chamber-operator}), with the Hermitian and the anti-Hermitian hopping operators $H_A$ and $J_A$ of Proposition~\ref{prop:accretivity}, has norm at most $M_0$ for every $A\in C^1(\mathbb{T})$, by the norm bound of Proposition~\ref{prop:accretivity} with the mass parameter $0$.
For $M>M_0$ the singular values of $D^{\mathbb{T}}_A$ therefore lie in $[M-M_0,M+M_0]$.
The determinant $\det D^{\mathbb{T}}_A\ge0$ (Proposition~\ref{prop:accretivity}) is the product of these singular values, and the spinors of $\mathbb{T}$ form a space of dimension $4\cdot\#|\mathbb{X}_{\mathbb{T}}|=192\,|\mathbb{T}|$; hence
\[
(M-M_0)^{192|\mathbb{T}|}\le\det D^{\mathbb{T}}_A\le(M+M_0)^{192|\mathbb{T}|}.
\]
For $\mathbb{T}=C(n)$ let
\[
Z^{\mathrm{ph}}_n:=\sum_z\int_{[0,2\pi/q)^{\mathbb{T}^{(1)}}}\mathrm{d}A\,\exp\left(-\frac{1}{2}\left\|dA+\frac{2\pi}{q}z\right\|^2\right),
\]
the sum being taken over the closed $z\in C^2(\mathbb{T};\Z)$ of flux zero.
Integrating the bounds for the determinant against the Maxwell weight of the sector of flux zero, one obtains
\[
(M-M_0)^{192n^3}\cdot Z^{\mathrm{ph}}_n\le Z^{(0)}_{C(n)}[1]\le(M+M_0)^{192n^3}\cdot Z^{\mathrm{ph}}_n,
\]
and the same with $jn$ in place of $n$.
By \cite[(10.2)]{NoncompactQED} the closed integer two-cochains of flux zero are the coboundaries $dk$ of the integer one-cochains $k\in C^1(\mathbb{T};\Z)$, and $dk=dk'$ if and only if $k-k'$ is closed.
With $\Gamma_n:=\frac{2\pi}{q}\cdot\left\{k\in C^1(\mathbb{T};\Z):\ dk=0\right\}$, the translates of $[0,2\pi/q)^{\mathbb{T}^{(1)}}$ by $\frac{2\pi}{q}k$, where $k$ ranges over a set of representatives of $C^1(\mathbb{T};\Z)$ modulo the closed integer one-cochains, form a fundamental domain of $\Gamma_n$, and
\[
\begin{aligned}
Z^{\mathrm{ph}}_n&=\sum_k\int_{[0,2\pi/q)^{\mathbb{T}^{(1)}}}\mathrm{d}A\,\exp\left(-\frac{1}{2}\left\|d\left(A+\frac{2\pi}{q}k\right)\right\|^2\right)\\
&=\int_{C^1(\mathbb{T})/\Gamma_n}\mathrm{d}A\,\exp\left(-\frac{1}{2}\|dA\|^2\right).
\end{aligned}
\]
This is the number $Z^{\mathrm{ph}}_n$ of \cite[Lemma~21.32]{NoncompactQED}, which involves $q$ and $n$ only; by part~(a) of that lemma,
\[
Z^{\mathrm{ph}}_n=(2\pi/q)^{n^3+2}\cdot(2\pi)^{n^3-1}\cdot\exp\left(-\frac{1}{2}\operatorname{Ld}(n)\right).
\]
Since $j^3(n^3+2)-\left((jn)^3+2\right)=2(j^3-1)$ and $j^3(n^3-1)-\left((jn)^3-1\right)=-(j^3-1)$, it follows that
\[
j^3\cdot\log Z^{\mathrm{ph}}_n-\log Z^{\mathrm{ph}}_{jn}=2(j^3-1)\log\frac{2\pi}{q}-(j^3-1)\log(2\pi)+\frac{1}{2}\left(\operatorname{Ld}(jn)-j^3\cdot\operatorname{Ld}(n)\right)=D^{\mathrm{ph}}_j,
\]
where
\[
D^{\mathrm{ph}}_j:=(j^3-1)\log\frac{2\pi}{q^2}+\frac{1}{2}\left(\operatorname{Ld}(jn)-j^3\cdot\operatorname{Ld}(n)\right).
\]
The condition on $q$ is equivalent to $D^{\mathrm{ph}}_j<0$.
The number $C_j$ is the number denoted $K_j$ in \cite[Section~21.1.7]{NoncompactQED}, and \cite[Lemma~21.31]{NoncompactQED} gives $\frac{1}{2}\left(\operatorname{Ld}(jn)-j^3\cdot\operatorname{Ld}(n)\right)\le-(j^3-1)\log n+C_j$; so $q^2n>2\pi\cdot\exp\left(\frac{C_j}{j^3-1}\right)$ implies
\[
2\pi\cdot\exp\left(\frac{\operatorname{Ld}(jn)-j^3\cdot\operatorname{Ld}(n)}{2(j^3-1)}\right)\le\frac{2\pi}{n}\cdot\exp\left(\frac{C_j}{j^3-1}\right)<q^2.
\]
For $M>M_0$ the two-sided bounds give
\[
j^3\cdot\log Z^{(0)}_{C(n)}[1]-\log Z^{(0)}_{C(jn)}[1]\le D^{\mathrm{ph}}_j+192j^3n^3\cdot\log\frac{M+M_0}{M-M_0},
\]
the logarithms being defined by Lemma~\ref{lem:sector-weights}\eqref{it:sector-weights-zero}.
The last term is positive and decreases to $0$ as $M\to\infty$; since $D^{\mathrm{ph}}_j<0$, there is $M_1>M_0$ such that the right side is negative for every $M\ge M_1$.
Division by $j^3n^3=|C(jn)|$ gives $f^{(0)}(n,n,n)<f^{(0)}(jn,jn,jn)$.
If $f^{(0)}$ were non-increasing in each side, then $f^{(0)}(jn,jn,jn)\le f^{(0)}(n,jn,jn)\le f^{(0)}(n,n,jn)\le f^{(0)}(n,n,n)$.
\end{proof}

\begin{remark}
By Proposition~\ref{prop:no-monotonicity}\eqref{it:no-monotonicity-norm}, $Z^{(0)}_{C(n_1,n_2,n_3)}[1]^{\frac{1}{n_1}}$ is $(2\pi)^{n_2n_3}$ times the $L^{n_1}$-norm of a function that does not depend on $n_1$, with respect to the product of Lebesgue measure on $[0,1]^2$ and counting measure.
For every measure and every non-negative function the logarithm of the integral of the $p$-th power is convex in $p$, by H\"older's inequality; this is the convexity in Proposition~\ref{prop:no-monotonicity}\eqref{it:no-monotonicity-convex}.
For counting measure the $L^p$-norms are moreover non-increasing in $p$.
The functional of all flux sectors is of this kind: $Z_{C(n_1,n_2,n_3)}[1]^{\frac{1}{n_1}}=(2\pi)^{n_2n_3}\cdot\|\mathfrak{M}^{*}\cdot K\cdot\mathfrak{M}\|_{n_1}$ by Proposition~\ref{prop:word-formula}, and the pressure $f$ is non-increasing in $n_1$ (Proposition~\ref{prop:coordinate-pressure}).
For the product of Lebesgue measure on $[0,1]^2$ and counting measure the $L^p$-norms are in general not monotone in $p$, and Proposition~\ref{prop:no-monotonicity}\eqref{it:no-monotonicity-fails} gives parameters for which the pressure $f^{(0)}$ of the sector of flux zero is larger on $C(jn)$ than on $C(n)$.
\end{remark}

\subsection{Limits of the averaged law}\label{sec:averaged-limits}

Let $L$, $\bar m$, $v_0$ and the standard weights be as in Definition~\ref{def:cutoffs}.
For a cutoff $(N,T)$ let $\eta=L^{-N}$, $n=8\cdot L^{\lfloor\log_L\bar m^{-1}\rfloor}\cdot L^{T+N}$, $R=n\cdot\eta$ and $\Lambda_{\eta,R}$ be as there,
and let $\langle\cdot,\cdot\rangle_\eta$, $|\cdot|_\eta$ and the step two-form $F^\sharp$ of a two-cochain $F$ on $\Lambda_{\eta,R}$ be those of Definition~\ref{def:cutoffs};
through the identification made there, the two-cochains on $\Lambda_{\eta,R}$ are the two-cochains of the cubic torus $C(n)$.
Fix a Wilson parameter $-1\le\rho\le1$, a charge $e>0$ and a mass assignment $N\mapsto m_{0,N}$.
At the cutoff $(N,T)$ the chamber operator on $C(n)$ is formed with the parameter $\rho$, the point $v_0$, the standard weights,
the lattice charge $q=e\cdot\eta^{\frac12}$ and the mass parameter $M=m_N+\frac{\rho}{2}\cdot\sum_iw_i$, $m_N=\bar m\cdot\eta+m_{0,N}$;
for these parameters let $\mu^{\mathrm{av}}_{C(n)}$ be the averaged law of Definition~\ref{def:averaged-law} on $C(n)$, a law on $\Omega$.
Let $\mathbb{P}^{\mathrm{av}}_{N,T}$ be the law of $F=\eta^{-\frac32}\cdot F^{\mathrm u}$ on $C^2(\Lambda_{\eta,R})$
when the periodic extension $\widetilde F^{\mathrm u}$ of the curvature $F^{\mathrm u}$ of a configuration on $C(n)$ has the law $\mu^{\mathrm{av}}_{C(n)}$.

For $\nu_0\in\{1,2,3\}$ let $\Sigma^{++}_{\nu_0}$ be the set of the plaquettes of $\Lambda_{\eta,R}$
all of whose vertices have their $\nu_0$-coordinate in $\left\{\eta,2\eta,\dots,\left(\frac n2-1\right)\cdot\eta\right\}$ modulo $R$,
as in Lemma~\ref{lem:positive-tests},
and let $\hat{\mathsf R}_{\nu_0}=\theta^{*}_{x_{\nu_0},0}$ be the reflection of the two-cochains on $\Lambda_{\eta,R}$ in the site plane $\{x_{\nu_0}=0\}$.

\begin{corollary}[Admissibility and volume-first limits]\label{cor:averaged-admissible}
Let $L$, $\bar m$, $v_0$ and the weights be as in Definition~\ref{def:cutoffs}, let $-1\le\rho\le1$ and $e>0$,
and let $N\mapsto m_{0,N}$ be a mass assignment.
\begin{enumerate}
\item[(a)] For all integers $N,T\ge0$ the law $\mathbb{P}^{\mathrm{av}}_{N,T}$ is carried by the closed cochains,
$\mathbb{P}^{\mathrm{av}}_{N,T}(\ker d)=1$, and it is invariant under $F\mapsto-F$.
It is invariant under the maps $\tau_{-v/\eta}^{*}$, $v\in\eta\Z^3$, and $(g^{-1})^{*}$, $g\in\Ocub$, of $C^2(\Lambda_{\eta,R})$
(condition \textup{(Ad1)}).
For every real two-cochain $\beta$ on $\Lambda_{\eta,R}$
\[
\int\mathrm{d}\mathbb{P}^{\mathrm{av}}_{N,T}(F)\,\exp\langle\beta,F\rangle_\eta\le\exp\left(\frac12|\beta|_\eta^2\right)
\]
(condition \textup{(Ad2)}).
For $\nu_0\in\{1,2,3\}$, every integer $r\ge1$, all real two-cochains $\beta_1,\dots,\beta_r$ on $\Lambda_{\eta,R}$ supported on $\Sigma^{++}_{\nu_0}$
and all $a_1,\dots,a_r\in\C$
\[
\sum_{k,l=1}^{r}\overline{a_k}\cdot a_l\cdot\int\mathrm{d}\mathbb{P}^{\mathrm{av}}_{N,T}(F)\,\exp\left(i\left\langle\beta_l-\hat{\mathsf R}_{\nu_0}\beta_k,F\right\rangle_\eta\right)\ge0,
\]
that is, $\mathbb{P}^{\mathrm{av}}_{N,T}$ satisfies the condition~\eqref{eq:continuum-admissible-even} of Proposition~\ref{prop:cutoff-law}.
\item[(b)] Let in addition $0<\rho<1$.
For an integer $N\ge0$ and $\eta=L^{-N}$ let $\mathcal V^{\mathrm{av}}_N$ be the set of the laws $(\sharp_\eta)_{*}\nu$ on $H_{-4}$,
where $\sharp_\eta$ is the step map of Definition~\ref{def:step-map}
and $\nu$ is a weak limit point on $\Omega$ of a sequence $\left(\mu^{\mathrm{av}}_{C(n_j)}\right)_{j\ge1}$,
\[
n_j=8\cdot L^{\lfloor\log_L\bar m^{-1}\rfloor}\cdot L^{T_j+N},
\]
along a sequence of integers $T_j\to\infty$,
all the laws $\mu^{\mathrm{av}}_{C(n_j)}$ being formed with the lattice charge $q=e\cdot\eta^{\frac12}$ and the lattice mass $m_N$.
Then $\mathcal V^{\mathrm{av}}_N$ is not empty for every $N$, and the union of the sets $\mathcal V^{\mathrm{av}}_N$, $N\ge0$, is tight on $H_{-4}$.
Let $\mu$ be the weak limit on $H_{-4}$ of a sequence of laws $\mu^{(N_i)}\in\mathcal V^{\mathrm{av}}_{N_i}$ with $N_i\to\infty$.
Then $\mu$ has all the properties that \cite[Theorem~3.1, Part~I\,(b), (c)]{NoncompactQED} asserts of a limit point,
except the convergence of moments and cumulants along a defining sequence of cutoffs,
and $\mu$ lies in $\mathcal C_9$ (Definition~\ref{def:laws}).
For $i,j\in\{1,2,3\}$ with $j\ne i$ the reconstruction of $\mu$ in the direction $j$ satisfies $\mathsf H_j\ge|\mathsf P_i|$ with the constant one;
moreover, for every $\Phi\in L^2(\Sigma_j^+,\mu)$ (Definition~\ref{def:positive-side}) the function
$u(v):=\mu\left((\Theta_j\Phi)\cdot U(v)\Phi\right)$, $v\in\R^3$,
formed with the antilinear reflection $\Theta_j$ and the translations $U(v)$ of Definition~\ref{def:laws},
satisfies the five-point inequality~\eqref{eq:continuum-five-point} for every $t>0$.
The projections $P^j_\mu$, $j=1,2,3$, are equal to $P^{123}_\mu$;
this is the property~\eqref{it:projections-coincide} of Definition~\ref{def:symmetry-conditions}.
Almost every ergodic component $\mu_\xi$ of $\mu$ under the translations of $\R^3$ lies in $\mathcal C_9$,
is invariant under $\R^3\rtimes\Ocub$, satisfies~\eqref{cond:gaussian-domination} and $\mu_\xi(\{F\in\mathcal S':dF=0\})=1$,
has $\ker\mathsf H_j(\mu_\xi)=\C\cdot\Omega_j$ and $\mathsf H_j(\mu_\xi)\ge|\mathsf P_i|$ for $j\ne i$, and satisfies
\[
\left|\mu_\xi\left(\overline{\Phi}\cdot U(a)\Phi'\right)-\overline{\mu_\xi(\Phi)}\cdot\mu_\xi(\Phi')\right|\le C_0\cdot\|\Phi\|_2\cdot\|\Phi'\|_2\cdot\frac{r}{|a|_\infty}
\]
for every $r>0$, all $\Phi,\Phi'\in L^2(\Sigma_{(-r,r)^3},\mu_\xi)$ and all $a\in\R^3$ with $|a|_\infty\ge C_1\cdot r$,
where $C_0$ and $C_1$ are the absolute constants of Theorem~\ref{thm:clustering-rate}.
\end{enumerate}
\end{corollary}

\begin{proof}
Throughout the proof $N,T\ge0$ are integers, and $\eta$ and $n$ are as above.

\emph{(a).}
The curvature of a configuration $(A,z)$ on $C(n)$ is $F^{\mathrm u}=dA+\frac{2\pi}{q}\cdot z$ with $dz=0$.
Hence $dF^{\mathrm u}=0$, and $\mathbb{P}^{\mathrm{av}}_{N,T}(\ker d)=1$.
By Definition~\ref{def:averaged-law}\eqref{it:averaged-covariance} the law $\mu^{\mathrm{av}}_{C(n)}$ is invariant under $F\mapsto-F$ and under $W$.
The first of these invariances is at once that of $\mathbb{P}^{\mathrm{av}}_{N,T}$ under $F\mapsto-F$.
For the second, let $g\in W$.
The linear part of $g$ maps $n\Z^3$ onto itself;
so $g^*$ maps $n\Z^3$-periodic two-cochains on $\Z^3$ to $n\Z^3$-periodic ones,
and $\widetilde{g^*F^{\mathrm u}}=g^*\widetilde F^{\mathrm u}$ for every two-cochain $F^{\mathrm u}$ of $C(n)$.
The invariance of $\mu^{\mathrm{av}}_{C(n)}$ under $g$ is therefore the invariance of $\mathbb{P}^{\mathrm{av}}_{N,T}$ under the map $g^*$ of $C^2(\Lambda_{\eta,R})$.
The maps $\tau_{-v/\eta}^{*}$, $v\in\eta\Z^3$, and $(g^{-1})^{*}$, $g\in\Ocub$, are of this form, and \textup{(Ad1)} follows.

Next let $\beta$ be a real two-cochain on $\Lambda_{\eta,R}$, and let $\beta^{\mathrm u}=\eta^{\frac32}\cdot\beta$, regarded as a two-cochain of $C(n)$.
Then $\langle\beta,F\rangle_\eta=\langle\beta^{\mathrm u},F^{\mathrm u}\rangle$ and $|\beta|_\eta=\|\beta^{\mathrm u}\|$.
Proposition~\ref{prop:averaged-domination}, applied to $\beta^{\mathrm u}$, gives
\[
\int\mathrm{d}\mathbb{P}^{\mathrm{av}}_{N,T}(F)\,\exp\langle\beta,F\rangle_\eta
=\int\mathrm{d}\mu^{\mathrm{av}}_{C(n)}\left(\widetilde F^{\mathrm u}\right)\,\exp\langle\beta^{\mathrm u},F^{\mathrm u}\rangle
\le\exp\left(\frac12\|\beta^{\mathrm u}\|^2\right)
=\exp\left(\frac12|\beta|_\eta^2\right),
\]
which is \textup{(Ad2)}.

For the last inequality of~(a) let $\nu_0$, $r$, $\beta_1,\dots,\beta_r$ and $a_1,\dots,a_r$ be as there,
let $\beta^{\mathrm u}_k=\eta^{\frac32}\cdot\beta_k$, let $\theta=\theta_{x_{\nu_0},0}$, and put
\[
G(F^{\mathrm u})=\sum_{l=1}^{r}a_l\cdot\exp\left(i\langle\beta^{\mathrm u}_l,F^{\mathrm u}\rangle\right).
\]
A plaquette in the support of some $\beta^{\mathrm u}_l$ is a plaquette of $C(n)$ all of whose vertices satisfy $0<x_{\nu_0}<\frac n2$,
the coordinate $x_{\nu_0}$ being represented in $\{0,\dots,n-1\}$.
So $G$ is a bounded continuous function of finitely many $F^{\mathrm u}_\sigma$, each $\sigma$ being a plaquette of this kind.
The side $n$ is even, and $\{x_{\nu_0}=0\}$ and $\left\{x_{\nu_0}=\frac n2\right\}$ are two site planes.
Proposition~\ref{prop:averaged-positivity}(a), in the direction $\nu_0$ and with $c=0$, gives
\[
\int\mathrm{d}\mu^{\mathrm{av}}_{C(n)}\,\overline{G\circ\theta^*}\cdot G\ge0.
\]
On two-cochains $\theta^*$ acts as $\hat{\mathsf R}_{\nu_0}$, which is an orthogonal involution and hence symmetric;
therefore $\langle\beta^{\mathrm u}_k,\theta^*F^{\mathrm u}\rangle=\langle\hat{\mathsf R}_{\nu_0}\beta^{\mathrm u}_k,F^{\mathrm u}\rangle$, and
\[
\overline{G(\theta^*F^{\mathrm u})}\cdot G(F^{\mathrm u})=\sum_{k,l=1}^{r}\overline{a_k}\cdot a_l\cdot\exp\left(i\left\langle\beta^{\mathrm u}_l-\hat{\mathsf R}_{\nu_0}\beta^{\mathrm u}_k,F^{\mathrm u}\right\rangle\right).
\]
Since $\left\langle\beta^{\mathrm u}_l-\hat{\mathsf R}_{\nu_0}\beta^{\mathrm u}_k,F^{\mathrm u}\right\rangle=\left\langle\beta_l-\hat{\mathsf R}_{\nu_0}\beta_k,F\right\rangle_\eta$,
the last two displays together give the asserted inequality.

\emph{(b).}
Fix $N\ge0$, $\eta=L^{-N}$, and a sequence of integers $T_j\to\infty$.
The sides $n_j$ are even, $n_j\ge8$ and $n_j\to\infty$,
while the parameters $\rho\in(0,1)$, $w_i>0$, $M=m_N+\frac{\rho}{2}\cdot\sum_iw_i\in\R$ and $q=e\cdot\eta^{\frac12}>0$ do not depend on $j$.
Theorem~\ref{thm:averaged-any-side}, applied to the sequence of sides $(n_j)$, shows that $\left(\mu^{\mathrm{av}}_{C(n_j)}\right)_j$ is tight on $\Omega$,
and that every weak limit point $\nu$ of it satisfies~\textup{(G)}, is invariant under $W$ and under $F\mapsto-F$
and, all $n_j$ being even, is reflection positive at every mirror of the nine families.
By Prokhorov's theorem the sequence has a subsequence that converges weakly on the Polish space $\Omega$,
and the limit of such a subsequence is a weak limit point.
For $T_j=j$ this shows that $\mathcal V^{\mathrm{av}}_N$ is not empty.

Every law $\nu$ that enters the definition of $\mathcal V^{\mathrm{av}}_N$ thus satisfies \textup{(G)}.
By Proposition~\ref{prop:volume-limits}(a) one has $\nu(\Omega^{(\eta)})=1$ and
\[
\int\mathrm{d}\nu(F^{\mathrm u})\,\|\sharp_\eta F^{\mathrm u}\|_{-3}^2\le\mathsf Z_6<\frac{1}{16},
\]
so that $\int\mathrm{d}\mu(F)\,\|F\|_{-3}^2\le\mathsf Z_6$ for every $\mu$ in the union of the sets $\mathcal V^{\mathrm{av}}_N$, $N\ge0$;
the bound does not depend on $N$.
For $\varrho>0$ the closed ball $\mathcal B_\varrho=\{F\in H_{-4}:\|F\|_{-3}\le\varrho\}$ is compact in $H_{-4}$
(Step~\ref{it:admissible-tightness} of the proof of Proposition~\ref{prop:admissible-families}),
and Chebyshev's inequality gives
\[
\mu\left(H_{-4}\setminus\mathcal B_\varrho\right)\le\mathsf Z_6\cdot\varrho^{-2}
\]
for all these $\mu$.
Hence the union is tight on $H_{-4}$.

Now let $\mu^{(N_i)}\in\mathcal V^{\mathrm{av}}_{N_i}$, $N_i\to\infty$, converge weakly on $H_{-4}$ to $\mu$.
Since $N_i\to\infty$, some subsequence has strictly increasing indices, and it converges to $\mu$ as well;
so one may assume $N_1<N_2<\cdots$.
Let $\mathcal N=\{N_i:i\ge1\}$, and for $N\in\mathcal N$ let $\nu^{(N)}$ be a weak limit point on $\Omega$
as in the definition of $\mathcal V^{\mathrm{av}}_N$ with $\mu^{(N)}=(\sharp_{L^{-N}})_{*}\nu^{(N)}$.
By the first paragraph of the proof of~(b), every $\nu^{(N)}$ satisfies \textup{(G)}, is invariant under $W$ and under $F\mapsto-F$,
and is reflection positive at every mirror of the nine families.
Proposition~\ref{prop:soft-limits}, applied to the family $\left(\nu^{(N)}\right)_{N\in\mathcal N}$ and to the weak limit $\mu$ of $\left(\mu^{(N_i)}\right)_{i\ge1}$,
shows that $\mu$ is soft, that is, it satisfies~\eqref{cond:soft-hypothesis}, $\mu(\{F\in\mathcal S':dF=0\})=1$
and it is invariant under $F\mapsto-F$; and it shows that $\mu$ lies in $\mathcal C_9$.
By \cite[Lemma~21.36]{NoncompactQED} a soft law has all the properties that \cite[Theorem~3.1, Part~I\,(b), (c)]{NoncompactQED} asserts of a limit point
except the convergence along a defining sequence of cutoffs.

Since $\mu$ lies in $\mathcal C_9$, Theorem~\ref{thm:cone}(c), (d) gives, for $j\ne i$, $\mathsf H_j\ge|\mathsf P_i|$ with the constant one
and the five-point inequality for every $t>0$ and every $\Phi\in L^2(\Sigma_j^+,\mu)$,
and Proposition~\ref{prop:zero-momentum} gives $P^j_\mu=P^{123}_\mu$ for $j=1,2,3$.
As $\mu$ satisfies \textup{(H)} in addition, Proposition~\ref{prop:ergodic-components} applies to $\mu$:
almost every ergodic component $\mu_\xi$ lies in $\mathcal C_9$, satisfies \textup{(G)}, is invariant under $\R^3\rtimes\Ocub$,
has $\ker\mathsf H_j(\mu_\xi)=\C\cdot\Omega_j$ and $\mathsf H_j(\mu_\xi)\ge|\mathsf P_i|$,
and is carried by the closed forms, because $\mu$ is.
Every such $\mu_\xi$ is ergodic under the translations of $\R^3$ and lies in $\mathcal C_9$.
Theorem~\ref{thm:clustering-rate}, applied to $\mu_\xi$, gives the bound of the statement,
the subtracted term of an ergodic law being $\overline{\mu_\xi(\Phi)}\cdot\mu_\xi(\Phi')$.
\end{proof}

For the data of Corollary~\ref{cor:averaged-admissible} and integers $N,T\ge0$
let $P^{\mathrm{av}}_{N,T}$ be the law of the step two-form $F^\sharp$ when $F$ has the law $\mathbb{P}^{\mathrm{av}}_{N,T}$.
The map $F\mapsto F^\sharp$ of the finite-dimensional space $C^2(\Lambda_{\eta,R})$ into $H_{-3}$ is linear and continuous (Definition~\ref{def:cutoffs});
so $P^{\mathrm{av}}_{N,T}$ is a Borel probability measure on $H_{-4}$, and a law on $\mathcal S'$.
Limit points, volume limits and volume-first limit points of the family $\left(P^{\mathrm{av}}_{N,T}\right)$ are defined as in Definition~\ref{def:limit-points}.
Thus a limit point is a Borel probability measure $\mu$ on $H_{-4}$ such that $P^{\mathrm{av}}_{N_i,T_i}\to\mu$ weakly on $H_{-4}$
for some sequences $N_i\to\infty$ and $T_i\to\infty$;
a volume limit at the spacing $L^{-N}$ is a weak limit on $H_{-4}$ of a sequence $\left(P^{\mathrm{av}}_{N,T_j}\right)_{j\ge1}$ with $T_j\to\infty$;
and a volume-first limit point is a weak limit on $H_{-4}$ of a sequence $\mu^{(N_i)}$, $N_i\to\infty$,
in which $\mu^{(N_i)}$ is a volume limit at the spacing $L^{-N_i}$ for every $i$.

\begin{theorem}[Limit points of the averaged laws]\label{thm:averaged-limit-points}
Let $L$, $\bar m$, $v_0$ and the weights be as in Definition~\ref{def:cutoffs}, let $-1\le\rho\le1$ and $e>0$,
and let $N\mapsto m_{0,N}$ be a mass assignment.
\begin{enumerate}
\item[(i)] One has
\[
\sup_{N,T}\int\mathrm{d}P^{\mathrm{av}}_{N,T}(F)\,\|F\|_{-3}^2\le61;
\]
the family $\left(P^{\mathrm{av}}_{N,T}\right)$ is tight on $H_{-4}$ and on $\mathcal S'$,
and every sequence of cutoffs has a subsequence along which $P^{\mathrm{av}}_{N,T}$ converges weakly on $H_{-4}$.
Every limit point $\mu$ of the family has all the properties that \cite[Theorem~3.1, Part~I\,(b), (c)]{NoncompactQED} asserts of a limit point
other than the convergence of moments along a defining sequence of cutoffs,
and it has this convergence in the following form:
if $P^{\mathrm{av}}_{N_k,T_k}\to\mu$ weakly on $H_{-4}$ with $N_k\to\infty$ and $T_k\to\infty$,
then for all test forms $f_1,\dots,f_r$, real or complex,
\[
\int\mathrm{d}P^{\mathrm{av}}_{N_k,T_k}(F)\,\exp\left(\sum_{p=1}^{r}z_p\cdot F(f_p)\right)\longrightarrow\int\mathrm{d}\mu(F)\,\exp\left(\sum_{p=1}^{r}z_p\cdot F(f_p)\right)
\]
as $k\to\infty$, uniformly for $z=(z_1,\dots,z_r)$ in every compact subset of $\C^r$,
and all moments and all joint cumulants of $F(f_1),\dots,F(f_r)$ under $P^{\mathrm{av}}_{N_k,T_k}$ converge to those under $\mu$.
Iterated limits are limit points:
let $\mathcal N$ be an infinite set of integers,
for every $N\in\mathcal N$ let $\mu^{(N)}$ be the weak limit on $H_{-4}$ of $P^{\mathrm{av}}_{N,T_j(N)}$ as $j\to\infty$ for a sequence $T_j(N)\to\infty$,
and let $\mu$ be the weak limit on $H_{-4}$ of $\mu^{(N_i)}$ along a sequence $N_i\to\infty$ in $\mathcal N$;
then $\mu$ is a limit point of the family.
The same holds with the roles of $N$ and $T$ exchanged.
\item[(ii)] For all integers $N,T\ge0$, with $\eta$ and $n$ as above and with the maps $\sharp_\eta$ and $J_\eta$ of Definition~\ref{def:step-map},
\[
P^{\mathrm{av}}_{N,T}=(\sharp_\eta)_{*}\mu^{\mathrm{av}}_{C(n)},\qquad\mu^{\mathrm{av}}_{C(n)}=(J_\eta)_{*}P^{\mathrm{av}}_{N,T}.
\]
Let $N\ge0$, let $T_j\to\infty$ be a sequence of integers, and let $n_j=8\cdot L^{\lfloor\log_L\bar m^{-1}\rfloor}\cdot L^{T_j+N}$.
Then $\left(P^{\mathrm{av}}_{N,T_j}\right)_{j\ge1}$ converges weakly on $H_{-4}$
if and only if $\left(\mu^{\mathrm{av}}_{C(n_j)}\right)_{j\ge1}$ converges weakly on $\Omega$,
and the two limits $\mu$ and $\nu$ then satisfy $\mu=(\sharp_\eta)_{*}\nu$ and $\nu=(J_\eta)_{*}\mu$.
\item[(iii)] Let in addition $0<\rho<1$.
For every $N\ge0$ the set $\mathcal V^{\mathrm{av}}_N$ of Corollary~\ref{cor:averaged-admissible}(b)
is the set of the volume limits of the family $\left(P^{\mathrm{av}}_{N,T}\right)$ at the spacing $L^{-N}$.
Hence the volume-first limit points of the family are exactly the laws $\mu$ of Corollary~\ref{cor:averaged-admissible}(b).
Volume-first limit points exist:
$\mathcal V^{\mathrm{av}}_N$ is not empty for every $N$,
and every sequence $\mu^{(N_i)}\in\mathcal V^{\mathrm{av}}_{N_i}$ with $N_i\to\infty$ has a subsequence that converges weakly on $H_{-4}$.
Every volume-first limit point $\mu$ of the family is a limit point of it, so that $\mu$ has the properties stated in~(i);
and $\mu$ has the properties \eqref{it:continuum-nine-planes}--\eqref{it:continuum-components}
listed in item~\eqref{it:continuum-volume-first} of Theorem~\ref{thm:continuum-limits}.
\end{enumerate}
\end{theorem}

\begin{proof}
\emph{(i).}
By Corollary~\ref{cor:averaged-admissible}(a), for all integers $N,T\ge0$ the law $\mathbb{P}^{\mathrm{av}}_{N,T}$ is a probability measure on $\ker d\subset C^2(\Lambda_{\eta,R})$,
with $\eta=L^{-N}$ and $R=8\cdot L^{\lfloor\log_L\bar m^{-1}\rfloor}\cdot L^{T}$,
that satisfies \textup{(Ad1)}, \textup{(Ad2)} and~\eqref{eq:continuum-admissible-even},
and $P^{\mathrm{av}}_{N,T}$ is the law of $F^\sharp$ under it.
So $\left(\mathbb{P}^{\mathrm{av}}_{N,T}\right)$ is a family as in Proposition~\ref{prop:admissible-families},
and the steps of the proof of that proposition hold for it;
since its members are invariant under $F\mapsto-F$, so does Step~\ref{it:admissible-even}.
The bound on the second moments, the tightness on $H_{-4}$ and on $\mathcal S'$ and the weakly convergent subsequences
are Step~\ref{it:admissible-tightness}.

Let $\mu$ be a limit point, and let $P^{\mathrm{av}}_{N_k,T_k}\to\mu$ weakly on $H_{-4}$ with $N_k,T_k\to\infty$.
By Steps~\ref{it:admissible-domination}, \ref{it:admissible-invariance} and~\ref{it:admissible-positivity},
$\mu$ satisfies the conditions~\eqref{eq:continuum-limit-domination}, \eqref{eq:continuum-limit-invariance} and~\eqref{eq:continuum-limit-positivity}
of \cite[Section~7.5]{NoncompactQED},
and these three conditions together are the hypothesis~\eqref{cond:soft-hypothesis} \cite[Remark~7.20]{NoncompactQED}.
By Step~\ref{it:admissible-closed}, $\mu$ satisfies the remaining condition~\eqref{eq:continuum-limit-closed} of that section as well;
so \cite[Proposition~7.19\,(a)]{NoncompactQED} gives for $\mu$ the assertion~(4) of \cite[Theorem~7.15]{NoncompactQED},
namely $\mu(\{F\in\mathcal S':dF=0\})=1$.
By Step~\ref{it:admissible-even} the characteristic functional of $\mu$ is even,
and the last assertion of \cite[Proposition~7.19]{NoncompactQED} shows that $\mu$ is invariant under $F\mapsto-F$.
Thus $\mu$ is soft, and \cite[Lemma~21.36]{NoncompactQED} gives all the properties that \cite[Theorem~3.1, Part~I\,(b), (c)]{NoncompactQED} asserts of a limit point
other than the convergence along a defining sequence.
That convergence, in the form stated in~(i), is Step~\ref{it:admissible-moments}.

Let $d_{\mathrm P}$ denote the Prokhorov metric between Borel probability measures on $H_{-4}$;
as $H_{-4}$ is a separable Hilbert space, a sequence of such measures converges weakly if and only if it converges for $d_{\mathrm P}$ \cite[Theorem~6.8]{Billingsley1999}.
Let $\mathcal N$, $\mu^{(N)}$, $T_j(N)$, $N_i$ and $\mu$ be as in the assertion on iterated limits.
For every $i$ choose an index $j_i$ with $T_{j_i}(N_i)\ge i$ and $d_{\mathrm P}\left(P^{\mathrm{av}}_{N_i,T_{j_i}(N_i)},\mu^{(N_i)}\right)\le\frac1i$.
Then
\[
d_{\mathrm P}\left(P^{\mathrm{av}}_{N_i,T_{j_i}(N_i)},\mu\right)\le\frac1i+d_{\mathrm P}\left(\mu^{(N_i)},\mu\right)\longrightarrow0,
\]
while $N_i\to\infty$ and $T_{j_i}(N_i)\to\infty$; so $\mu$ is a limit point.
With the roles of $N$ and $T$ exchanged, $\mu^{(T)}$ is, for every $T$ of an infinite set,
the weak limit of $P^{\mathrm{av}}_{N_j(T),T}$ as $j\to\infty$ for a sequence $N_j(T)\to\infty$,
and $\mu$ is the weak limit of $\mu^{(T_i)}$ along a sequence $T_i\to\infty$ of elements of this set.
An index $j_i$ with $N_{j_i}(T_i)\ge i$ and $d_{\mathrm P}\left(P^{\mathrm{av}}_{N_{j_i}(T_i),T_i},\mu^{(T_i)}\right)\le\frac1i$
gives, by the triangle inequality for $d_{\mathrm P}$, $P^{\mathrm{av}}_{N_{j_i}(T_i),T_i}\to\mu$ weakly on $H_{-4}$,
with $N_{j_i}(T_i)\to\infty$ and $T_i\to\infty$.

\emph{(ii). The two identities.}
Consider the curvature $F^{\mathrm u}$ of a configuration on $C(n)$, and $F=\eta^{-\frac32}\cdot F^{\mathrm u}$.
The periodic extension $\widetilde F^{\mathrm u}$ is bounded, so it lies in the set $\Omega'$ of Definition~\ref{def:step-map},
and Lemma~\ref{lem:test-unfolding} together with~\eqref{eq:continuum-step-pairing} gives $F^\sharp(f)=\sharp_\eta\left(\widetilde F^{\mathrm u}\right)(f)$ for every $f\in\mathcal S$,
that is, $F^\sharp=\sharp_\eta\left(\widetilde F^{\mathrm u}\right)$ in $\mathcal S'$.
By Step~\ref{it:admissible-tightness} of the proof of Proposition~\ref{prop:admissible-families},
\[
\|F^\sharp\|_{-3}^2\le961\cdot\mathsf Z_4\cdot|F|_\eta^2<\infty
\]
for each single configuration.
Hence $\widetilde F^{\mathrm u}$ lies in $\Omega^{(\eta)}$, and $J_\eta\left(\sharp_\eta\left(\widetilde F^{\mathrm u}\right)\right)=\widetilde F^{\mathrm u}$ by Definition~\ref{def:step-map}.
Since $\mu^{\mathrm{av}}_{C(n)}$ is the law of $\widetilde F^{\mathrm u}$ and $P^{\mathrm{av}}_{N,T}$ is the law of $F^\sharp$, both identities follow.

\emph{Convergence on $H_{-4}$ implies convergence on $\Omega$.}
Suppose that $P^{\mathrm{av}}_{N,T_j}\to\mu$ weakly on $H_{-4}$.
The map $J_\eta$ is continuous from $H_{-4}$ into the space $C^2$ of all real two-cochains on $\Z^3$ with the product topology
(Definition~\ref{def:step-map}).
By the mapping theorem \cite[Theorem~2.7]{Billingsley1999} and the second identity,
\[
\mu^{\mathrm{av}}_{C(n_j)}=(J_\eta)_{*}P^{\mathrm{av}}_{N,T_j}\longrightarrow\nu:=(J_\eta)_{*}\mu
\]
weakly on $C^2$.
The set $\Omega$ is closed in $C^2$, and $\mu^{\mathrm{av}}_{C(n_j)}(\Omega)=1$ for every $j$;
by the portmanteau theorem \cite[Theorem~2.1]{Billingsley1999}, $\nu(\Omega)\ge\limsup_j\mu^{\mathrm{av}}_{C(n_j)}(\Omega)=1$.
Every closed subset of $\Omega$ is closed in $C^2$, so that the same theorem gives the weak convergence on $\Omega$.

\emph{The limit $\nu$ satisfies \textup{(G)}.}
Fix a real $\beta\in C^2_{\mathrm c}$, and let $r>0$ be such that the base points of the plaquettes in the support of $\beta$ lie in a cube of side $2r$.
If $n_j>2r$, two different plaquettes of the support of $\beta$ are not congruent modulo $n_j\Z^3$, and the periodisation
\[
\beta^{(n_j)}([\sigma])=\sum_{v\in\Z^3}\beta(\sigma+n_j\cdot v),
\]
a two-cochain of $C(n_j)$, has at most one non-zero term in each sum.
Therefore $\langle\beta,\widetilde F^{\mathrm u}\rangle=\langle\beta^{(n_j)},F^{\mathrm u}\rangle$ for every two-cochain $F^{\mathrm u}$ of $C(n_j)$,
and $\|\beta^{(n_j)}\|=\|\beta\|$.
Proposition~\ref{prop:averaged-domination}, applied to $\beta^{(n_j)}$, gives
\[
\int\mathrm{d}\mu^{\mathrm{av}}_{C(n_j)}(F)\,\exp\langle\beta,F\rangle\le\exp\left(\frac12\|\beta\|^2\right)\qquad(n_j>2r).
\]
For $\varepsilon>0$ the function $F\mapsto\min\left(\exp\langle\beta,F\rangle,\varepsilon^{-1}\right)$ depends on finitely many coordinates
and is bounded and continuous on $\Omega$.
The weak convergence on $\Omega$ yields
\[
\int\mathrm{d}\nu(F)\,\min\left(\exp\langle\beta,F\rangle,\varepsilon^{-1}\right)\le\exp\left(\frac12\|\beta\|^2\right),
\]
and monotone convergence as $\varepsilon\to0$ gives \textup{(G)} for $\nu$.

\emph{The law $\mu$ is the image of $\nu$ under the step map.}
By Proposition~\ref{prop:volume-limits}(a), applied to $\nu$, one has $\nu(\Omega^{(\eta)})=1$;
as $\nu=(J_\eta)_{*}\mu$, the cochain $J_\eta(F)$ lies in $\Omega^{(\eta)}$ for $\mu$-almost every $F$.
Let $\mathcal S_0$ be a countable subset of $C^\infty_{\mathrm c}(\R^3;\Lambda^2)\cap\mathcal S$ that is dense in $\mathcal S$;
it exists, since $\mathcal S$ is separable and $C^\infty_{\mathrm c}(\R^3;\Lambda^2)\cap\mathcal S$ is dense in it.
For $f\in\mathcal S_0$ let $f_{[\eta]}$ be the test cochain of $f$ on $\Z^3$ (Definition~\ref{def:test-cochains}).
Only finitely many of the cubes $Q_\eta(\sigma)$ meet the compact support of $f$;
so $f_{[\eta]}$ has finite support, and the test form
\[
f'=f-\sum_\sigma f_{[\eta]}(\sigma)\cdot u^\sigma
\]
is defined by a finite sum.
The test cochain depends linearly on the test form, and the bumps $u^\sigma$ of Definition~\ref{def:step-map}(d) satisfy $(u^\sigma)_{[\eta]}=\mathbf 1_\sigma$;
hence $(f')_{[\eta]}=0$.
By Lemma~\ref{lem:test-unfolding}, $F^\sharp(f')=0$ for every configuration on every torus $C(n)$; hence
\[
P^{\mathrm{av}}_{N,T}\left(\{F\in H_{-4}:F(f')=0\}\right)=1
\]
for every $T$.
This set is closed in $H_{-4}$, and the portmanteau theorem gives $\mu(\{F:F(f')=0\})=1$.
Consequently, for $\mu$-almost every $F$ and all $f\in\mathcal S_0$,
\[
F(f)=\sum_\sigma f_{[\eta]}(\sigma)\cdot F(u^\sigma)=\sum_\sigma f_{[\eta]}(\sigma)\cdot J_\eta(F)(\sigma)=\sharp_\eta\left(J_\eta(F)\right)(f),
\]
the last equality by~\eqref{eq:continuum-step-pairing}, since $J_\eta(F)$ lies in $\Omega^{(\eta)}$.
The two ends are continuous in $f\in\mathcal S$, and $\mathcal S_0$ is dense; so $F=\sharp_\eta(J_\eta(F))$ for $\mu$-almost every $F$, and
\[
\mu=(\sharp_\eta\circ J_\eta)_{*}\mu=(\sharp_\eta)_{*}\nu.
\]

\emph{Convergence on $\Omega$ implies convergence on $H_{-4}$.}
For the converse, let $\mu^{\mathrm{av}}_{C(n_j)}\to\nu$ weakly on $\Omega$, and let $(j_k)_{k\ge1}$ be a strictly increasing sequence of indices.
Applied to the sequence of cutoffs $(N,T_{j_k})$, $k\ge1$, assertion~(i) provides a subsequence $(j_{k_l})_{l\ge1}$ and a Borel probability measure $\mu'$ on $H_{-4}$
with $P^{\mathrm{av}}_{N,T_{j_{k_l}}}\to\mu'$ weakly on $H_{-4}$ as $l\to\infty$.
The indices $T_{j_{k_l}}$ tend to infinity;
so the three preceding paragraphs, with $(T_{j_{k_l}})_{l\ge1}$ in place of $(T_j)$ and $\mu'$ in place of $\mu$,
show that $\mu^{\mathrm{av}}_{C(n_{j_{k_l}})}\to(J_\eta)_{*}\mu'$ weakly on $\Omega$ and that $\mu'=(\sharp_\eta)_{*}(J_\eta)_{*}\mu'$.
The same laws $\mu^{\mathrm{av}}_{C(n_{j_{k_l}})}$ tend to $\nu$, as a subsequence of $\left(\mu^{\mathrm{av}}_{C(n_j)}\right)_{j\ge1}$,
and a weak limit on $\Omega$ is unique;
therefore $(J_\eta)_{*}\mu'=\nu$, and $\mu'=(\sharp_\eta)_{*}\nu$, whatever the sequence $(j_k)$ was.
Consequently, from each subsequence of $\left(P^{\mathrm{av}}_{N,T_j}\right)_{j\ge1}$ one can extract a subsequence
that tends to $(\sharp_\eta)_{*}\nu$ for the metric $d_{\mathrm P}$,
whose convergent sequences are exactly the weakly convergent ones \cite[Theorem~6.8]{Billingsley1999}.
In a metric space a sequence with this property converges to that point;
hence $P^{\mathrm{av}}_{N,T_j}\to(\sharp_\eta)_{*}\nu$ weakly on $H_{-4}$.

\emph{(iii).}
A weak limit point $\nu$ on $\Omega$ of $\left(\mu^{\mathrm{av}}_{C(n_j)}\right)_j$, for a sequence of integers $T_j\to\infty$,
is the weak limit along a subsequence $(j_k)_{k\ge1}$, and $(T_{j_k})_{k\ge1}$ is again a sequence of integers tending to infinity.
By~(ii), $(\sharp_\eta)_{*}\nu$ is the weak limit on $H_{-4}$ of $\left(P^{\mathrm{av}}_{N,T_{j_k}}\right)_{k\ge1}$, a volume limit at the spacing $L^{-N}$.
Conversely, let $P^{\mathrm{av}}_{N,T_j}\to\mu$ weakly on $H_{-4}$ for a sequence $T_j\to\infty$.
By~(ii), $\left(\mu^{\mathrm{av}}_{C(n_j)}\right)_j$ converges weakly on $\Omega$ to $\nu=(J_\eta)_{*}\mu$, which is a weak limit point,
and $\mu=(\sharp_\eta)_{*}\nu$ lies in $\mathcal V^{\mathrm{av}}_N$.
This proves the first assertion of~(iii), and the second follows from the definitions.
By Corollary~\ref{cor:averaged-admissible}(b) every $\mathcal V^{\mathrm{av}}_N$ is not empty and the union of the sets $\mathcal V^{\mathrm{av}}_N$ is tight on $H_{-4}$;
by Prokhorov's theorem every sequence $\mu^{(N_i)}\in\mathcal V^{\mathrm{av}}_{N_i}$ has a subsequence that converges weakly on $H_{-4}$.

Consider now a volume-first limit point $\mu$, the weak limit of volume limits $\mu^{(N_i)}$ at the spacings $L^{-N_i}$ with $N_i\to\infty$.
Some subsequence has strictly increasing indices and the same limit; so one may assume $N_1<N_2<\cdots$.
With $\mathcal N=\{N_i:i\ge1\}$, the assertion of~(i) on iterated limits shows that $\mu$ is a limit point of the family;
so $\mu$ has the properties stated in~(i).
In particular $\mu$ satisfies \textup{(H)}, hence \textup{(G)}, it satisfies $\mu(\{F\in\mathcal S':dF=0\})=1$,
and almost every ergodic component of $\mu$ has the properties that \cite[Theorem~3.1, Part~I\,(c)]{NoncompactQED} lists.
By Corollary~\ref{cor:averaged-admissible}(b), $\mu$ lies in $\mathcal C_9$.
It remains to verify the properties \eqref{it:continuum-nine-planes}--\eqref{it:continuum-components} of Theorem~\ref{thm:continuum-limits}\eqref{it:continuum-volume-first} for $\mu$.

\emph{The property~\eqref{it:continuum-nine-planes}.}
Fix a plane $\{\ell=c\}$ of the nine families, and let $r_{\ell,c}$ be the reflection of $\R^3$ in it.
As $\mu$ is invariant under $r_{\ell,c}$, the operator $\hat r_{\ell,c}$ is unitary on $L^2(\mu)$.
By \cite[Lemma~7.14]{NoncompactQED}, applied to the linear space of the real forms in $C^\infty_{\mathrm c}(\{\ell>c\};\Lambda^2)$,
whose $\sigma$-algebra there is $\Sigma_{\{\ell>c\}}$ here,
the linear span of the functions $\exp(iF(f))$, $f$ real in $C^\infty_{\mathrm c}(\{\ell>c\};\Lambda^2)$, is dense in $L^2(\Sigma_{\{\ell>c\}},\mu)$.
For $\Phi=\sum_ka_k\cdot\exp(iF(f_k))$ in this span
\[
\left\langle\hat r_{\ell,c}\Phi,\Phi\right\rangle_\mu=\sum_{k,l}\overline{a_k}\cdot a_l\cdot\mu\left(\exp\left(i\left(F(f_l)-F(r_{\ell,c}\cdot f_k)\right)\right)\right)\ge0,
\]
by the reflection positivity of $\mu$ at $\{\ell=c\}$.
The form $\Phi\mapsto\langle\hat r_{\ell,c}\Phi,\Phi\rangle_\mu$ is continuous on $L^2(\mu)$;
so the inequality holds for every $\Phi\in L^2(\Sigma_{\{\ell>c\}},\mu)$.

\emph{The properties~\eqref{it:continuum-cone}--\eqref{it:continuum-components}.}
The property~\eqref{it:continuum-cone} follows from Theorem~\ref{thm:cone}, applied to the law $\mu$ of $\mathcal C_9$,
and from Corollary~\ref{cor:unbounded-observables}, which applies because $\mu$ satisfies \textup{(G)}.
The property~\eqref{it:continuum-vacuum} is Proposition~\ref{prop:zero-momentum}, and the property~\eqref{it:continuum-rate} is Theorem~\ref{thm:clustering-rate},
both applied to $\mu$.
For the property~\eqref{it:continuum-components}: $\mu$ satisfies \textup{(H)}, lies in $\mathcal C_9$ and is carried by the closed forms,
so that Proposition~\ref{prop:ergodic-components} applies to it;
Theorem~\ref{thm:clustering-rate}, applied to the ergodic laws $\mu_\xi$ of $\mathcal C_9$, gives the clustering bound for almost every component;
finally, $\mu$ satisfies the hypothesis \textup{(H)} of \cite[Corollary~8.9]{NoncompactQED}, so that this corollary applies to $\mu$ itself:
by the uniqueness asserted there, the measure on the laws that it provides is the one of the ergodic decomposition of $\mu$ under $\R^3$,
and almost every component has the properties 1--6 listed in it.
\end{proof}

\subsection{Joint limits of spacing and volume}\label{sec:joint-limits}

Throughout this subsection $d_{\mathrm{P}}$ denotes the Prokhorov metric on the set of Borel probability measures on $H_{-4}$:
for two such measures $\mu$ and $\mu'$, $d_{\mathrm{P}}(\mu,\mu')$ is the infimum of the numbers $\varepsilon>0$ such that
\[
\mu(B)\le\mu'\left(B^{\varepsilon}\right)+\varepsilon
\qquad\text{and}\qquad
\mu'(B)\le\mu\left(B^{\varepsilon}\right)+\varepsilon
\]
for every Borel set $B\subset H_{-4}$,
where $B^{\varepsilon}$ denotes the set of the points of $H_{-4}$ at distance less than $\varepsilon$ from $B$.
Since $H_{-4}$ is a separable metric space, $d_{\mathrm{P}}$ metrises weak convergence:
a sequence $(\mu_k)_{k\ge1}$ of Borel probability measures on $H_{-4}$ converges weakly to $\mu$ if and only if $d_{\mathrm{P}}(\mu_k,\mu)\to0$
\cite[Theorem~6.8]{Billingsley1999}.
For a Borel probability measure $\mu$ on $H_{-4}$ and a non-empty set $\mathcal{V}$ of such measures put
\[
d_{\mathrm{P}}(\mu,\mathcal{V}):=\inf\left\{d_{\mathrm{P}}(\mu,\mu'):\ \mu'\in\mathcal{V}\right\}.
\]
As in Definition~\ref{def:limit-points}, the laws $P^{\mathrm{ch}}_{N,T}$ of Definition~\ref{def:cutoffs} are regarded as Borel probability measures on $H_{-4}$,
and weak convergence means weak convergence on $H_{-4}$.

\begin{theorem}[Joint limits]\label{thm:joint-limits}
Let $0<\rho<1$ and $v_0\in C_0$, let the weights $w_i$ be standard, let $L\ge3$ be an odd integer, let $0<\bar m\le1$ and $e>0$,
and let $N\mapsto m_{0,N}$ be a mass assignment.
For integers $N,T\ge0$ let $P^{\mathrm{ch}}_{N,T}$ be the law of Definition~\ref{def:cutoffs},
$N$ being the index of the spacing $L^{-N}$ and $T$ the index of the volume,
and let $\mathcal{V}_N$ be the set of the volume limits at the spacing $L^{-N}$ (Definition~\ref{def:limit-points}).
\begin{enumerate}
\item\label{it:joint-limits-distance}
For every $N$ the set $\mathcal{V}_N$ is not empty, and
\[
\lim_{T\to\infty}d_{\mathrm{P}}\left(P^{\mathrm{ch}}_{N,T},\mathcal{V}_N\right)=0 .
\]
Consequently, for every $N$ and every $\varepsilon>0$ there is an integer $T_0(N,\varepsilon)\ge0$ such that
$d_{\mathrm{P}}\left(P^{\mathrm{ch}}_{N,T},\mathcal{V}_N\right)\le\varepsilon$ for all integers $T\ge T_0(N,\varepsilon)$.
\item\label{it:joint-limits-first}
Let $(N_i)_{i\ge1}$ and $(T_i)_{i\ge1}$ be sequences of integers with $N_i\to\infty$, $T_i\to\infty$ and
\[
\lim_{i\to\infty}d_{\mathrm{P}}\left(P^{\mathrm{ch}}_{N_i,T_i},\mathcal{V}_{N_i}\right)=0 .
\]
Then every weak limit point of the sequence $\left(P^{\mathrm{ch}}_{N_i,T_i}\right)_{i\ge1}$ is a volume-first limit point;
in particular it has all the properties that Theorem~\ref{thm:continuum-limits}\eqref{it:continuum-volume-first} asserts of volume-first limit points.
The condition on the distances holds whenever $T_i\ge T_0(N_i,\varepsilon_i)$ for all $i$, for some sequence of numbers $\varepsilon_i>0$ with $\varepsilon_i\to0$;
the choice $T_i\ge\max\left(i,T_0(N_i,\varepsilon_i)\right)$ gives in addition $T_i\to\infty$.
Hence for every sequence of integers $N_i\to\infty$ there are integers $T_i\to\infty$ with
$d_{\mathrm{P}}\left(P^{\mathrm{ch}}_{N_i,T_i},\mathcal{V}_{N_i}\right)\to0$.
\end{enumerate}
\end{theorem}

\begin{proof}
Fix $N$.
The family $\left(P^{\mathrm{ch}}_{N,T}\right)_{N,T\ge0}$ is tight on $H_{-4}$ by Theorem~\ref{thm:continuum-limits}\eqref{it:continuum-soft},
and so is its subfamily $\left(P^{\mathrm{ch}}_{N,T}\right)_{T\ge0}$.
By Prokhorov's theorem, every sequence of integers $T_j\to\infty$ has a subsequence $(T_{j_k})_{k\ge1}$ along which $P^{\mathrm{ch}}_{N,T_{j_k}}$ converges weakly;
since $T_{j_k}\to\infty$, the limit belongs to $\mathcal{V}_N$ by Definition~\ref{def:limit-points}.
Applied to $T_j:=j$, this shows that $\mathcal{V}_N$ is not empty.
Suppose that $d_{\mathrm{P}}\left(P^{\mathrm{ch}}_{N,T},\mathcal{V}_N\right)$ does not tend to $0$ as $T\to\infty$.
Then there are $\delta>0$ and a sequence of integers $T_j\to\infty$ with $d_{\mathrm{P}}\left(P^{\mathrm{ch}}_{N,T_j},\mathcal{V}_N\right)\ge\delta$ for all $j$.
Let $(T_{j_k})_{k\ge1}$ be a subsequence along which $P^{\mathrm{ch}}_{N,T_{j_k}}$ converges weakly, to $\mu\in\mathcal{V}_N$ say.
As $d_{\mathrm{P}}$ metrises weak convergence,
\[
d_{\mathrm{P}}\left(P^{\mathrm{ch}}_{N,T_{j_k}},\mathcal{V}_N\right)\le d_{\mathrm{P}}\left(P^{\mathrm{ch}}_{N,T_{j_k}},\mu\right)\longrightarrow0\qquad(k\to\infty),
\]
which is incompatible with the lower bound $\delta$.
This proves the limit in~\eqref{it:joint-limits-distance}, and the existence of $T_0(N,\varepsilon)$ is a restatement of it.

For~\eqref{it:joint-limits-first}, let $\mu$ be a weak limit point of $\left(P^{\mathrm{ch}}_{N_i,T_i}\right)_{i\ge1}$, that is, the weak limit along a subsequence.
The three conditions $N_i\to\infty$, $T_i\to\infty$ and $d_{\mathrm{P}}\left(P^{\mathrm{ch}}_{N_i,T_i},\mathcal{V}_{N_i}\right)\to0$ hold along every subsequence.
After replacing $(N_i,T_i)_{i\ge1}$ by a suitable subsequence and renumbering, one may therefore assume that
$P^{\mathrm{ch}}_{N_i,T_i}\to\mu$ weakly and that $(N_i)_{i\ge1}$ is strictly increasing.
By~\eqref{it:joint-limits-distance} each set $\mathcal{V}_{N_i}$ is not empty, and the definition of the infimum provides $\mu^{(N_i)}\in\mathcal{V}_{N_i}$ with
\[
d_{\mathrm{P}}\left(P^{\mathrm{ch}}_{N_i,T_i},\mu^{(N_i)}\right)\le d_{\mathrm{P}}\left(P^{\mathrm{ch}}_{N_i,T_i},\mathcal{V}_{N_i}\right)+\frac{1}{i}.
\]
The triangle inequality for $d_{\mathrm{P}}$ gives
\[
\begin{aligned}
d_{\mathrm{P}}\left(\mu^{(N_i)},\mu\right)
&\le d_{\mathrm{P}}\left(\mu^{(N_i)},P^{\mathrm{ch}}_{N_i,T_i}\right)+d_{\mathrm{P}}\left(P^{\mathrm{ch}}_{N_i,T_i},\mu\right)\\
&\le d_{\mathrm{P}}\left(P^{\mathrm{ch}}_{N_i,T_i},\mathcal{V}_{N_i}\right)+\frac{1}{i}+d_{\mathrm{P}}\left(P^{\mathrm{ch}}_{N_i,T_i},\mu\right).
\end{aligned}
\]
On the right, the first term tends to $0$ by hypothesis, the second is $1/i$,
and the third tends to $0$ because $P^{\mathrm{ch}}_{N_i,T_i}\to\mu$ weakly and $d_{\mathrm{P}}$ metrises weak convergence.
Hence $\mu^{(N_i)}\to\mu$ weakly, with $\mu^{(N_i)}\in\mathcal{V}_{N_i}$ and $N_i\to\infty$:
by Definition~\ref{def:limit-points}, $\mu$ is a volume-first limit point,
and Theorem~\ref{thm:continuum-limits}\eqref{it:continuum-volume-first} applies to it.

It remains to justify the sufficient condition and the existence of the sequences $(T_i)$.
If $T_i\ge T_0(N_i,\varepsilon_i)$, then $d_{\mathrm{P}}\left(P^{\mathrm{ch}}_{N_i,T_i},\mathcal{V}_{N_i}\right)\le\varepsilon_i$ by~\eqref{it:joint-limits-distance},
and the right side tends to $0$ with $\varepsilon_i$.
Given integers $N_i\to\infty$, the integers $T_i:=\max\left(i,T_0(N_i,1/i)\right)$, $i\ge1$, satisfy $T_i\ge i$, so that $T_i\to\infty$,
and $d_{\mathrm{P}}\left(P^{\mathrm{ch}}_{N_i,T_i},\mathcal{V}_{N_i}\right)\le 1/i$.
\end{proof}

\begin{remark}\label{rem:joint-limits-tight}
Consider a family $(\lambda_{N,T})_{N,T\ge0}$ of Borel probability measures on $H_{-4}$, indexed by pairs of non-negative integers,
such that $(\lambda_{N,T})_{T\ge0}$ is tight for every $N$.
Let $\mathcal{V}_N(\lambda)$ be the set of the weak limits of the sequences $(\lambda_{N,T_j})_{j\ge1}$ with $T_j\to\infty$,
and call a weak limit of a sequence $\left(\mu^{(N_i)}\right)_{i\ge1}$ with $\mu^{(N_i)}\in\mathcal{V}_{N_i}(\lambda)$ and $N_i\to\infty$ a volume-first limit point of $(\lambda_{N,T})$;
for $\lambda_{N,T}=P^{\mathrm{ch}}_{N,T}$ these are the sets and the points of Definition~\ref{def:limit-points}.
The proof of Theorem~\ref{thm:joint-limits}, with the tightness of $(\lambda_{N,T})_{T\ge0}$ in place of Theorem~\ref{thm:continuum-limits}\eqref{it:continuum-soft}, establishes three facts.
For every $N$ the set $\mathcal{V}_N(\lambda)$ is not empty, and $d_{\mathrm{P}}\left(\lambda_{N,T},\mathcal{V}_N(\lambda)\right)\to0$ as $T\to\infty$.
If $N_i\to\infty$, $T_i\to\infty$ and $d_{\mathrm{P}}\left(\lambda_{N_i,T_i},\mathcal{V}_{N_i}(\lambda)\right)\to0$,
then every weak limit point of $(\lambda_{N_i,T_i})_{i\ge1}$ is a volume-first limit point of $(\lambda_{N,T})$.
For every sequence of integers $N_i\to\infty$ there are integers $T_i\to\infty$ with $d_{\mathrm{P}}\left(\lambda_{N_i,T_i},\mathcal{V}_{N_i}(\lambda)\right)\to0$.

The averaged laws form such a family.
Let $P^{\mathrm{av}}_{N,T}$ be the laws of Theorem~\ref{thm:averaged-limit-points} and $\mathcal{V}^{\mathrm{av}}_N$ the sets of Corollary~\ref{cor:averaged-admissible},
both formed with the parameters of Theorem~\ref{thm:joint-limits}.
By Theorem~\ref{thm:averaged-limit-points}, the family $\left(P^{\mathrm{av}}_{N,T}\right)_{T\ge0}$ is tight on $H_{-4}$ for every $N$;
the set $\mathcal{V}^{\mathrm{av}}_N$ consists of the weak limits of the sequences $\left(P^{\mathrm{av}}_{N,T_j}\right)_{j\ge1}$ with $T_j\to\infty$,
that is, $\mathcal{V}^{\mathrm{av}}_N=\mathcal{V}_N(\lambda)$ for $\lambda_{N,T}:=P^{\mathrm{av}}_{N,T}$;
and every volume-first limit point of $\left(P^{\mathrm{av}}_{N,T}\right)$ has all the properties that Theorem~\ref{thm:continuum-limits}\eqref{it:continuum-volume-first} asserts of volume-first limit points.
Consequently, for every $N$ the set $\mathcal{V}^{\mathrm{av}}_N$ is not empty and $d_{\mathrm{P}}\left(P^{\mathrm{av}}_{N,T},\mathcal{V}^{\mathrm{av}}_N\right)\to0$ as $T\to\infty$;
if $N_i\to\infty$, $T_i\to\infty$ and $d_{\mathrm{P}}\left(P^{\mathrm{av}}_{N_i,T_i},\mathcal{V}^{\mathrm{av}}_{N_i}\right)\to0$,
then every weak limit point of $\left(P^{\mathrm{av}}_{N_i,T_i}\right)_{i\ge1}$ is a volume-first limit point of $\left(P^{\mathrm{av}}_{N,T}\right)$ and has all these properties;
and for every sequence of integers $N_i\to\infty$ there are integers $T_i\to\infty$ with $d_{\mathrm{P}}\left(P^{\mathrm{av}}_{N_i,T_i},\mathcal{V}^{\mathrm{av}}_{N_i}\right)\to0$.
\end{remark}

\section{Rotations}\label{sec:rotations}

\subsection{The field and its symmetries}\label{sec:field-symmetries}

Throughout this section a law is a probability measure on the measurable space $(\mathcal S',\Sigma)$, where $\mathcal S'=\mathcal S'(\R^3;\Lambda^2)$,
and the notation of Definition~\ref{def:laws} is in force \cite[Section~8.1]{NoncompactQED}.

\begin{definition}[Field and Euclidean actions]\label{def:field-actions}
The dual vector field of a two-form $f$ on $\R^3$ is
\[
\ast f=\left(f_{23},-f_{13},f_{12}\right)
\]
\cite[Section~8.1.6]{NoncompactQED}, as in Definition~\ref{def:free-maxwell}.
Pointwise, $\ast$ is an isometry of $\Lambda^2$ onto $\R^3$,
and every $b\in\mathcal S(\R^3;\R^3)$ is of the form $b=\ast f$ for exactly one $f\in\mathcal S(\R^3;\Lambda^2)$.
For such $b$ the field $\mathbb B(b)$ is the function
\[
\mathbb B(b):=F(f)\qquad(b=\ast f)
\]
on $\mathcal S'$.
For a complex field $b=b'+ib''$ with $b',b''\in\mathcal S(\R^3;\R^3)$ one sets $\mathbb B(b):=\mathbb B(b')+i\mathbb B(b'')$;
thus $\overline{\mathbb B(b)}=\mathbb B(\overline{b})$.
The norm of $L^2(\R^3;\C^3)$ is denoted $\|\cdot\|$.
The unit vectors $e_1,e_2,e_3$ of $\R^3$ serve also as the unit vectors of the space $\C^3$ in which $b$ takes its values,
and for $a\in\{1,2,3\}$ one writes $\mathbb B_a(u):=\mathbb B(u\cdot e_a)$ for every scalar function $u$ on $\R^3$ for which the right side is defined.

A Euclidean map $\mathfrak g\colon x\mapsto gx+a$, with $g\in\mathrm O(3)$ and $a\in\R^3$,
acts on two-forms, on $\mathcal S'$, on functions on $\mathcal S'$ and on laws as in Definition~\ref{def:laws};
on vector fields $b\colon\R^3\to\C^3$ it acts by
\[
(\mathfrak g\cdot b)(x):=\det(g)\cdot g\,b\left(\mathfrak g^{-1}x\right).
\]
For $j\in\{1,2,3\}$ the reflection $\theta_j$ in the plane $\{x_j=0\}$ is also read as the diagonal matrix
whose entry in the place $j$ is $-1$ and whose two other diagonal entries are $1$;
its operator $\hat\theta_j$ on functions on $\mathcal S'$ is complex linear,
and $\Theta_j\Phi=\overline{\hat\theta_j\Phi}$ is the reflection of Definition~\ref{def:laws}.

For $j\neq i$ in $\{1,2,3\}$ and $t\in\R$ let $R^{ji}_t\in\mathrm{SO}(3)$ be the rotation that fixes the third unit vector and satisfies
\[
R^{ji}_te_j=\cos t\cdot e_j+\sin t\cdot e_i,\qquad R^{ji}_te_i=-\sin t\cdot e_j+\cos t\cdot e_i .
\]
At $t=\pi/2$ one obtains the \emph{quarter turn of the plane $(j,i)$}, which maps $e_j$ to $e_i$ and $e_i$ to $-e_j$;
as an orthogonal matrix with integer entries, $R^{ji}_{\pi/2}$ lies in $\Ocub$.
\end{definition}

These actions are compatible with the field.
For every two-form $f$ one has $\mathfrak g\cdot(\ast f)=\ast(\mathfrak g\cdot f)$
(Lemma~\ref{lem:orientation-signs}; \cite[Section~8.9.4]{NoncompactQED}),
and $\hat{\mathfrak g}$ maps $F\mapsto F(f)$ onto $F\mapsto F(\mathfrak g\cdot f)$;
hence $\hat{\mathfrak g}\,\mathbb B(b)=\mathbb B(\mathfrak g\cdot b)$ for every real field $b\in\mathcal S(\R^3;\R^3)$.
The map $b\mapsto\mathfrak g\cdot b$ is complex linear,
and it takes real fields to real fields because the matrix $\det(g)\cdot g$ is real:
for $b=b'+ib''$ with $b',b''\in\mathcal S(\R^3;\R^3)$ the fields $\mathfrak g\cdot b'$ and $\mathfrak g\cdot b''$ are real,
and $\mathfrak g\cdot b=\mathfrak g\cdot b'+i\,\mathfrak g\cdot b''$.
The operator $\hat{\mathfrak g}$, being a composition operator, is complex linear as well; therefore
\[
\begin{aligned}
\hat{\mathfrak g}\,\mathbb B(b)&=\hat{\mathfrak g}\,\mathbb B(b')+i\,\hat{\mathfrak g}\,\mathbb B(b'')\\
&=\mathbb B(\mathfrak g\cdot b')+i\,\mathbb B(\mathfrak g\cdot b'')=\mathbb B(\mathfrak g\cdot b).
\end{aligned}
\]
Altogether, for $b\in\mathcal S(\R^3;\C^3)$ and $a\in\R^3$, at every point of $\mathcal S'$,
\begin{equation}\label{eq:rotation-field-action}
\hat{\mathfrak g}\,\mathbb B(b)=\mathbb B(\mathfrak g\cdot b),\qquad U(a)\,\mathbb B(b)=\mathbb B(\tau_a\cdot b),\qquad(\tau_a\cdot b)(x)=b(x-a);
\end{equation}
the second identity is the first one for the translation $\tau_a\colon x\mapsto x+a$, whose operator is $U(a)=\hat\tau_a$.
For $\mathfrak g=\theta_j$ the action reads $(\theta_j\cdot b)(x)=-\theta_j\,b(\theta_jx)$,
because $\det\theta_j=-1$ and $\theta_j^{-1}=\theta_j$:
the component of $b$ normal to the plane $\{x_j=0\}$ is even under the reflection, and the two tangential components are odd.
If $u$ is a differentiable function on $\R^3$, the chain rule
$\nabla\left(u\circ\mathfrak g^{-1}\right)(x)=g\,(\nabla u)\left(\mathfrak g^{-1}x\right)$ yields
\begin{equation}\label{eq:rotation-gradient-action}
\mathfrak g\cdot\nabla u=\det(g)\cdot\nabla\left(u\circ\mathfrak g^{-1}\right).
\end{equation}

The reflection $\Theta_j$ is an antilinear, multiplicative involution of the functions on $\mathcal S'$.
When $\mu$ satisfies $(T_j)$, $(R_j)$ and $(RP_j)$, the scalar product of its reconstruction in the direction $e_j$ is
\[
\left\langle[\Phi]_j,[\Phi']_j\right\rangle=\left\langle\hat\theta_j\Phi,\Phi'\right\rangle_\mu=\mu\left((\Theta_j\Phi)\cdot\Phi'\right)
\qquad\left(\Phi,\Phi'\in L^2(\Sigma_j^+,\mu)\right)
\]
(Definition~\ref{def:laws}; here $\Sigma_j^+=\Sigma_{\{x_j>0\}}$, as in Definition~\ref{def:positive-side}).
For $b\in\mathcal S(\R^3;\C^3)$ and $a\in\R^3$,
\begin{equation}\label{eq:rotation-conjugation}
\Theta_j\mathbb B(b)=\mathbb B\left(\overline{\theta_j\cdot b}\right),\qquad
\Theta_jU(a)=U(\theta_ja)\,\Theta_j,\qquad
\Theta_j\exp\mathbb B(b)=\exp\mathbb B\left(\overline{\theta_j\cdot b}\right).
\end{equation}
Indeed, $\hat\theta_j$ is multiplicative and commutes with the composition with functions \cite[Section~8.1.2]{NoncompactQED};
so the first and the third identity follow from \eqref{eq:rotation-field-action} and from $\overline{\mathbb B(b)}=\mathbb B(\overline{b})$.
The second one follows from $\theta_j\tau_a=\tau_{\theta_ja}\theta_j$,
because the operator of a composition of Euclidean maps is the product of their operators \cite[Section~8.1.2]{NoncompactQED},
and because $U(\theta_ja)$ commutes with complex conjugation.

The hat carries two meanings, which are distinguished by what stands under it.
Over a Euclidean map, as in $\hat{\mathfrak g}$, in $\hat g$ for $g\in\mathrm O(3)$, in $\hat\theta_j$ or in $\hat r_{n,c}$,
it denotes the operator of that map on functions on $\mathcal S'$.
Over a function or a field it denotes the Fourier transform, normalised in every dimension as in Definition~\ref{def:free-maxwell}:
\[
\hat b(p):=\int\mathrm dx\,\exp(-ip\cdot x)\cdot b(x).
\]
The single exception is the transverse profile $\chi$ of a plate (Definition~\ref{def:labels-plates}), a real function of the coordinate $x_3$;
its transform
\[
\hat\chi(p_3):=\int\mathrm dx_3\,\exp\left(ip_3\cdot x_3\right)\cdot\chi(x_3)
\]
carries the opposite sign in the exponent:
the sign of the Laplace--Fourier transforms of Lemma~\ref{lem:spectral-matrix} and Definition~\ref{def:two-pictures},
which are integrated against the joint spectral measures $\mathfrak E_j$ of Lemma~\ref{lem:joint-spectrum}.
For a column $v\in\C^3$, $v^\dagger$ denotes its conjugate transpose.

\begin{definition}[Closedness, quarter turns and cones]\label{def:closedness-cones}
Let $\mu$ be a law, and let $j\neq i$ be elements of $\{1,2,3\}$.
The conditions \textup{(G)}, $(T_J)$, $(R_j)$, $(RP_j)$, reflection positivity at a plane,
the hypothesis \textup{(H)} and the classes $\mathcal C$, $\mathcal C_9$ are those of Definition~\ref{def:laws}
\cite[Sections~8.1.3 and~8.1.5]{NoncompactQED}.
Further conditions on $\mu$ are the following.
\begin{enumerate}
\item
The law $\mu$ satisfies \textup{(dF)} if it is carried by the closed distributions:
\begin{equation}\tag{dF}\label{cond:closedness}
\mu\left(\left\{F\in\mathcal S':dF=0\right\}\right)=1 .
\end{equation}
Here $F$ is closed, $dF=0$, if $F(f)=0$ for every $f\in\mathcal S(\R^3;\Lambda^2)$ such that $\ast f=\nabla u$ for some $u$
in the space $\mathcal S(\R^3)$ of real Schwartz functions on $\R^3$;
the closed distributions form a set in $\Sigma$ \cite[Section~8.1.6]{NoncompactQED}.
For real $u\in\mathcal S(\R^3)$ the field $\nabla u$ is $\ast f$ with $f\in\mathcal S(\R^3;\Lambda^2)$,
so that $\mathbb B(\nabla u)=F(f)$ vanishes at every closed $F$;
by linearity the same holds for complex $u$.
Hence a law with \textup{(dF)} satisfies
\begin{equation}\label{eq:rotation-gradient-fields}
\mu\left(\left\{\mathbb B(\nabla u)\neq0\right\}\right)=0\qquad\left(u\in\mathcal S(\R^3;\C)\right).
\end{equation}
\item
The law $\mu$ satisfies $(Q_{ji})$ if it is invariant under the quarter turn of the plane $(j,i)$:
\begin{equation}\tag{$Q_{ji}$}\label{cond:quarter-turn}
\left(R^{ji}_{\pi/2}\right)_*\mu=\mu .
\end{equation}
\item
Let $\mu$ satisfy $(T_{123})$, $(R_j)$ and $(RP_j)$,
let $\left(\mathcal H_j,[\,\cdot\,]_j,\mathsf H_j,V_j,\Omega_j\right)$ be its reconstruction in the direction $e_j$, formed on $L^2(\Sigma_j^+,\mu)$
(Definitions~\ref{def:laws} and~\ref{def:positive-side}; \cite[Proposition~8.19]{NoncompactQED}),
and let $\mathfrak E_j$ be its joint spectral measure on $[0,\infty)\times e_j^\perp$ (Lemma~\ref{lem:joint-spectrum}).
The variables of $\mathfrak E_j$ are written $(\omega,p_\perp)$, where $\omega$ is the energy and $p_\perp=(p_k)_{k\neq j}$ the momentum,
$e_j^\perp$ being identified with $\R^2$ and $p_\perp\cdot y=\sum_{k\neq j}p_k\cdot y_k$.
In these letters Lemma~\ref{lem:joint-spectrum}(b), (c) reads
\[
\exp\left(-t\cdot\mathsf H_j\right)V_j(y)=\int\mathrm d\mathfrak E_j(\omega,p_\perp)\,\exp\left(-t\cdot\omega+ip_\perp\cdot y\right)
\quad\left(t\ge0,\ y\in e_j^\perp\right),
\qquad
\mathsf P_i=\int\mathrm d\mathfrak E_j(\omega,p_\perp)\,p_i .
\]
The law $\mu$ satisfies $(C_{ji})$ if
\begin{equation}\tag{$C_{ji}$}\label{cond:cone}
\mathfrak E_j\left(\left\{(\omega,p_\perp):\omega<|p_i|\right\}\right)=0,
\end{equation}
that is, if the joint spectral measure of $(\mathsf H_j,\mathsf P_i)$, the image of $\mathfrak E_j$ under $(\omega,p_\perp)\mapsto(\omega,p_i)$,
is carried by $\{\omega\ge|p_i|\}$.
This is the inequality
\begin{equation}\label{eq:rotation-cone-operator}
\mathsf H_j\ge\left|\mathsf P_i\right|
\end{equation}
with the constant one, in the sense of Theorem~\ref{thm:cone}(c):
since $|p_i|\le\omega$ holds $\mathfrak E_j$-almost everywhere,
every $\psi$ in the domain of $\mathsf H_j$ lies in the domain of $\mathsf P_i$,
and $\|\mathsf P_i\psi\|\le\|\mathsf H_j\psi\|$, $\left\langle\psi,|\mathsf P_i|\psi\right\rangle\le\left\langle\psi,\mathsf H_j\psi\right\rangle$.
\item\label{it:one-field-space}
Let $\mu$ satisfy \textup{(G)}, $(T_{123})$, $(R_j)$ and $(RP_j)$,
and let $b\in L^2(\R^3;\C^3)$ vanish almost everywhere on $\{x_j<0\}$.
Then $\mathbb B(b)\in L^2(\Sigma_j^+,\mu)$ by Lemma~\ref{lem:domination-consequences}\eqref{it:field-extension},
so that the vector $[\mathbb B(b)]_j$ of $\mathcal H_j$ is defined.
The \emph{one-field space} $\mathcal H^1_j$ is the closed linear span in $\mathcal H_j$ of all these vectors.
For $t\ge0$ and $y\in e_j^\perp$ the translations $\tau_{te_j}$ and $\tau_y$ leave $\mu$ invariant;
by \cite[Proposition~8.19\,(ii), (iii)]{NoncompactQED} and Lemma~\ref{lem:domination-consequences}\eqref{it:field-extension},
\[
\exp\left(-t\cdot\mathsf H_j\right)\left[\mathbb B(b)\right]_j=\left[\mathbb B(\tau_{te_j}\cdot b)\right]_j,\qquad
V_j(y)\left[\mathbb B(b)\right]_j=\left[\mathbb B(\tau_y\cdot b)\right]_j ,
\]
and the fields $\tau_{te_j}\cdot b$ and $\tau_y\cdot b$ vanish almost everywhere on $\{x_j<0\}$ together with $b$.
Each of the bounded operators $\exp(-t\cdot\mathsf H_j)$, $t\ge0$, and $V_j(y)$, $y\in e_j^\perp$,
therefore maps the spanning vectors of $\mathcal H^1_j$ into $\mathcal H^1_j$, and with them the whole closed span.
The operators $\exp(-t\cdot\mathsf H_j)$ are self-adjoint, and the adjoint of $V_j(y)$ is $V_j(-y)$;
it follows that the orthogonal complement of $\mathcal H^1_j$ is mapped into itself as well,
and the orthogonal projection $Q$ of $\mathcal H_j$ onto $\mathcal H^1_j$ commutes with every $\exp(-t\cdot\mathsf H_j)$ and every $V_j(y)$.
By the functional calculus $Q$ then commutes with every bounded Borel function of $\exp(-\mathsf H_j)$,
in particular with $\exp(it\cdot\mathsf H_j)=h_t\left(\exp(-\mathsf H_j)\right)$, $t\in\R$,
where $h_t(s):=\exp(-it\cdot\log s)$ for $0<s\le1$ and $h_t(0):=0$;
consequently $Q$ commutes with $\exp(it\cdot\mathsf H_j)V_j(y)$.
The operator $2Q-1$ is unitary and self-adjoint,
so that $A\mapsto(2Q-1)\,\mathfrak E_j(A)\,(2Q-1)$ is a projection-valued measure on the Borel sets of $\R\times e_j^\perp$;
its integral against $\exp\left(it\cdot\omega+ip_\perp\cdot y\right)$ is $(2Q-1)\exp(it\cdot\mathsf H_j)V_j(y)\,(2Q-1)=\exp(it\cdot\mathsf H_j)V_j(y)$
for all $t\in\R$ and $y\in e_j^\perp$.
By the uniqueness in Lemma~\ref{lem:joint-spectrum}(b) this measure equals $\mathfrak E_j$.
Thus $Q$ commutes with $\mathfrak E_j(A)$ for every Borel set $A$,
and $\mathcal H^1_j$ is invariant under all spectral projections $\mathfrak E_j(A)$.
\item
Let $\mu$ be as in \eqref{it:one-field-space}.
It satisfies $(C^1_{ji})$, the inequality \eqref{eq:rotation-cone-operator} on the one-field space, if
\begin{equation}\tag{$C^1_{ji}$}\label{cond:one-field-cone}
\mathfrak E_j\left(\left\{(\omega,p_\perp):\omega<|p_i|\right\}\right)\psi=0\qquad\left(\psi\in\mathcal H^1_j\right).
\end{equation}
Condition $(C_{ji})$ implies $(C^1_{ji})$.
\item\label{it:one-plane-hypotheses}
The law $\mu$ satisfies the \emph{one-plane hypotheses} for the pair $(j,i)$
if it satisfies \textup{(G)}, \textup{(dF)}, $(T_{123})$, $(R_j)$, $(RP_j)$, $(Q_{ji})$ and $(C_{ji})$.
\end{enumerate}
\end{definition}

\begin{lemma}[Consequences of Gaussian domination]\label{lem:domination-consequences}
Let $\mu$ be a law that satisfies \textup{(G)}.
\begin{enumerate}
\item\label{it:field-extension}
The map $b\mapsto\mathbb B(b)$ has exactly one extension to a continuous complex linear map from $L^2(\R^3;\C^3)$ to $L^2(\mu)$,
again denoted $b\mapsto\mathbb B(b)$.
Its values lie in $\bigcap_{1\le p<\infty}L^p(\mu)$, and for every $b\in L^2(\R^3;\C^3)$ and $N=1,2,\dots$
\[
\begin{aligned}
&\overline{\mathbb B(b)}=\mathbb B(\overline{b}),\qquad\mu\left(\mathbb B(b)\right)=0,\qquad\mu\left(|\mathbb B(b)|^2\right)\le\|b\|^2,\qquad
\left|\exp\mathbb B(b)\right|=\exp\mathbb B(\Re b),\\
&\left\|\mathbb B(b)\right\|_{L^N(\mu)}\le c_N\cdot\|b\|,\qquad
c_N:=2^{\frac12}\cdot\left(N\cdot2^{\frac N2}\cdot\Gamma\left(\frac N2\right)\right)^{\frac1N};
\end{aligned}
\]
if $b$ is real, moreover $\mu\left(\exp\mathbb B(b)\right)\le\exp\left(\frac12\|b\|^2\right)$.
If $U$ is $\R^3$ or one of the half spaces $\{x_j>0\}$, $\{x_j<0\}$, $j\in\{1,2,3\}$,
and $b=0$ almost everywhere outside $U$, then $\mathbb B(b)\in L^2(\Sigma_U,\mu)$.
If $\mathfrak g$ is a Euclidean map with $\mathfrak g_*\mu=\mu$,
then $\hat{\mathfrak g}\,\mathbb B(b)=\mathbb B(\mathfrak g\cdot b)$ in $L^2(\mu)$ for every $b\in L^2(\R^3;\C^3)$.
\item\label{it:moment-forms}
For $N\ge1$ the moment
\[
S_N(b_1,\dots,b_N):=\mu\left(\mathbb B(b_1)\cdot\ldots\cdot\mathbb B(b_N)\right)
\]
is a symmetric $N$-linear form on $L^2(\R^3;\C^3)$ with $\left|S_N(b_1,\dots,b_N)\right|\le c_N^N\cdot\prod_{a=1}^N\|b_a\|$;
moreover $S_1=0$ and $\overline{S_N(b_1,\dots,b_N)}=S_N(\overline{b_1},\dots,\overline{b_N})$.
\item\label{it:moment-invariance}
A Euclidean map $\mathfrak g$ satisfies $\mathfrak g_*\mu=\mu$ if and only if
$S_N(\mathfrak g\cdot b_1,\dots,\mathfrak g\cdot b_N)=S_N(b_1,\dots,b_N)$ for every $N\ge1$ and all $b_1,\dots,b_N\in\mathcal S(\R^3;\R^3)$.
\end{enumerate}
\end{lemma}

\begin{proof}
Consider first \eqref{it:field-extension} for a real field $b\in L^2(\R^3;\R^3)$, and put $f:=\ast^{-1}b$, so that $\|f\|_2=\|b\|$.
By \cite[Lemma~8.15\,(i), (ii)]{NoncompactQED}, where $F(f)$ is written $X(f)$,
the map $\mathbb B$ extends from $\mathcal S(\R^3;\R^3)$ to a real linear map of norm at most one from $L^2(\R^3;\R^3)$ to $L^2(\mu)$;
the extension satisfies the inequality of \textup{(G)} for all real $b\in L^2(\R^3;\R^3)$, as well as
\[
\mu\left(\mathbb B(b)\right)=0,\qquad\mu\left(\mathbb B(b)^2\right)\le\|b\|^2,\qquad
\mu\left(|\mathbb B(b)|^N\right)\le N\cdot2^{\frac N2}\cdot\Gamma\left(\frac N2\right)\cdot\|b\|^N\quad(N\ge1),
\]
and $\mathbb B(b)\in L^2(\Sigma_U,\mu)$ whenever $b$ vanishes almost everywhere outside one of the sets $U$ of the statement.
For $b=b'+ib''$ with real $b',b''\in L^2(\R^3;\R^3)$ put $\mathbb B(b):=\mathbb B(b')+i\mathbb B(b'')$.
This map is complex linear and agrees with Definition~\ref{def:field-actions} on $\mathcal S(\R^3;\C^3)$;
as $\mathbb B(b')$ and $\mathbb B(b'')$ are real,
\[
\mu\left(|\mathbb B(b)|^2\right)=\mu\left(\mathbb B(b')^2\right)+\mu\left(\mathbb B(b'')^2\right)\le\|b'\|^2+\|b''\|^2=\|b\|^2 .
\]
So the map is continuous,
and it is the only continuous extension because $\mathcal S(\R^3;\C^3)$ is dense in $L^2(\R^3;\C^3)$.
The reality of $\mathbb B(b')$ and $\mathbb B(b'')$ also gives
$\overline{\mathbb B(b)}=\mathbb B(b')-i\mathbb B(b'')=\mathbb B(\overline{b})$ and $\left|\exp\mathbb B(b)\right|=\exp\mathbb B(b')=\exp\mathbb B(\Re b)$,
whereas linearity of the integral, together with the mean zero for real fields, yields
$\mu\left(\mathbb B(b)\right)=\mu\left(\mathbb B(b')\right)+i\mu\left(\mathbb B(b'')\right)=0$.
Minkowski's inequality and the moment bound for real fields give
\[
\left\|\mathbb B(b)\right\|_{L^N(\mu)}\le\left\|\mathbb B(b')\right\|_{L^N(\mu)}+\left\|\mathbb B(b'')\right\|_{L^N(\mu)}
\le\left(N\cdot2^{\frac N2}\cdot\Gamma\left(\frac N2\right)\right)^{\frac1N}\cdot\left(\|b'\|+\|b''\|\right),
\]
and $\|b'\|+\|b''\|\le2^{\frac12}\cdot\|b\|$.
This is the bound with the constant $c_N$;
as $\mu$ is a probability measure, it places $\mathbb B(b)$ in $L^p(\mu)$ for every $p<\infty$.
If $b$ vanishes almost everywhere outside $U$, so do $b'$ and $b''$, whence $\mathbb B(b)\in L^2(\Sigma_U,\mu)$.
Let finally $\mathfrak g_*\mu=\mu$.
For $b\in\mathcal S(\R^3;\C^3)$ the identity $\hat{\mathfrak g}\,\mathbb B(b)=\mathbb B(\mathfrak g\cdot b)$ holds at every point of $\mathcal S'$,
by \eqref{eq:rotation-field-action}.
The operator $\hat{\mathfrak g}$ is unitary on $L^2(\mu)$ \cite[Section~8.1.2]{NoncompactQED},
and $b\mapsto\mathfrak g\cdot b$ is an isometry of $L^2(\R^3;\C^3)$;
hence both sides of the identity are continuous functions of $b\in L^2(\R^3;\C^3)$ with values in $L^2(\mu)$,
and they agree on the dense subspace $\mathcal S(\R^3;\C^3)$.

In \eqref{it:moment-forms}, each $\mathbb B(b_a)$ lies in $L^N(\mu)$ with norm at most $c_N\cdot\|b_a\|$ by \eqref{it:field-extension},
and H\"older's inequality with $N$ exponents equal to $N$ gives the bound on $\left|S_N(b_1,\dots,b_N)\right|$.
The form is $N$-linear because $b\mapsto\mathbb B(b)$ is linear, and symmetric because the functions $\mathbb B(b_a)$ commute.
Further $S_1(b)=\mu\left(\mathbb B(b)\right)=0$,
and the rule $\overline{\mathbb B(b)}=\mathbb B(\overline{b})$ gives $\overline{S_N(b_1,\dots,b_N)}=S_N(\overline{b_1},\dots,\overline{b_N})$.

For \eqref{it:moment-invariance}, suppose first that $\mathfrak g_*\mu=\mu$,
so that $\mu(\hat{\mathfrak g}\Phi)=\mu(\Phi)$ for every $\Phi\in L^1(\mu)$ \cite[Section~8.1.2]{NoncompactQED}.
For $b_1,\dots,b_N\in\mathcal S(\R^3;\R^3)$ the product $\mathbb B(b_1)\cdot\ldots\cdot\mathbb B(b_N)$ is integrable by \eqref{it:moment-forms},
and the multiplicative operator $\hat{\mathfrak g}$ maps it to $\mathbb B(\mathfrak g\cdot b_1)\cdot\ldots\cdot\mathbb B(\mathfrak g\cdot b_N)$,
by \eqref{eq:rotation-field-action};
integration against $\mu$ gives the invariance of $S_N$.
Conversely, assume that all the moments are invariant, and let $b=\ast f$ with real $f\in\mathcal S(\R^3;\Lambda^2)$.
The partial sums of the series $\sum_{N\ge0}\frac{1}{N!}\cdot\left(i\mathbb B(b)\right)^N$ are bounded in modulus by
$\exp\left|\mathbb B(b)\right|\le\exp\mathbb B(b)+\exp\mathbb B(-b)$,
which is integrable by \textup{(G)} for $f$ and for $-f$.
By dominated convergence,
\[
\mu\left(\exp iF(f)\right)=1+\sum_{N\ge1}\frac{i^N}{N!}\cdot S_N(b,\dots,b),
\]
and the same identity holds with $\mathfrak g\cdot f$ in place of $f$,
because $\mathfrak g\cdot b=\ast(\mathfrak g\cdot f)$ again lies in $\mathcal S(\R^3;\R^3)$.
By hypothesis the two series agree term by term.
Since $\left(\mathfrak g^{-1}F\right)(f)=F(\mathfrak g\cdot f)$ for $F\in\mathcal S'$,
\[
\left(\left(\mathfrak g^{-1}\right)_*\mu\right)\left(\exp iF(f)\right)=\mu\left(\exp iF(\mathfrak g\cdot f)\right)=\mu\left(\exp iF(f)\right)
\]
for every real $f\in\mathcal S(\R^3;\Lambda^2)$.
By \cite[Lemma~8.12\,(ii)]{NoncompactQED}, two laws that assign the same integral to $F\mapsto\exp iF(f)$ for every real $f\in\mathcal S(\R^3;\Lambda^2)$ are equal.
Hence $\left(\mathfrak g^{-1}\right)_*\mu=\mu$, and applying $\mathfrak g_*$ to both sides gives $\mathfrak g_*\mu=\mu$.
\end{proof}

Let $\mu$ satisfy \textup{(G)} and $(R_j)$.
Then the identities \eqref{eq:rotation-conjugation} hold $\mu$-almost surely for every $b\in L^2(\R^3;\C^3)$.
To see this, apply Lemma~\ref{lem:domination-consequences}\eqref{it:field-extension} to $\mathfrak g=\theta_j$, which leaves $\mu$ invariant:
$\hat\theta_j\mathbb B(b)=\mathbb B(\theta_j\cdot b)$ in $L^2(\mu)$.
As $(\theta_j)_*\mu=\mu$, the operator $\hat\theta_j$ maps functions that agree $\mu$-almost surely to functions that agree $\mu$-almost surely;
it is multiplicative and commutes with the composition with functions \cite[Section~8.1.2]{NoncompactQED},
so that also $\hat\theta_j\exp\mathbb B(b)=\exp\mathbb B(\theta_j\cdot b)$ holds $\mu$-almost surely.
Complex conjugation, together with the rule $\overline{\mathbb B(b)}=\mathbb B(\overline{b})$ of Lemma~\ref{lem:domination-consequences}\eqref{it:field-extension},
turns these two identities into the first and the third identity of \eqref{eq:rotation-conjugation},
and the second identity does not involve $b$.

\begin{remark}[Domination by the free Maxwell field]\label{rem:maxwell-domination}
Let $P^\perp$ be the orthogonal projection of $L^2(\R^3;\C^3)$ with the Fourier multiplier
\[
P^\perp(p):=1-|p|^{-2}\cdot p\,p^T\qquad\left(p\in\R^3\setminus\{0\}\right);
\]
on real fields $b=\ast f$ it is the projection $P^\perp$ of Definition~\ref{def:free-maxwell}, read on the dual vector field.
A law $\mu$ satisfies \textup{(G)} and \textup{(dF)} if and only if
\[
\mu\left(\exp\mathbb B(b)\right)\le\exp\left(\frac12\left\|P^\perp b\right\|^2\right)\qquad\left(b\in\mathcal S(\R^3;\R^3)\right).
\]
This bound is the first inequality of the domination of the Laplace transform by that of the free Maxwell law
in \cite[Theorem~3.1, Part~I\,(b)]{NoncompactQED}, where $\left\|P^\perp(\ast f)\right\|^2$ is written $C_0(f,f)$.
For the proof of the equivalence, assume first \textup{(G)} and \textup{(dF)}.
The gradients $\nabla u$, $u\in\mathcal S(\R^3;\C)$, lie in the range of $1-P^\perp$, because $\widehat{\nabla u}(p)=i\hat u(p)\cdot p$ and $P^\perp(p)\,p=0$;
they are dense in it.
Indeed, let $v$ belong to this range and be orthogonal to all of them.
Plancherel's theorem gives
$\int\mathrm dp\,\overline{\hat u(p)}\cdot\left(p\cdot\hat v(p)\right)=0$ for every $u\in\mathcal S(\R^3;\C)$,
and $\hat u$ ranges over $\mathcal S(\R^3;\C)$ together with $u$;
so $p\cdot\hat v(p)=0$ for almost every $p$, which means $P^\perp v=v$.
As $v$ lies in the range of $1-P^\perp$, also $P^\perp v=0$, and therefore $v=0$.
By Lemma~\ref{lem:domination-consequences}\eqref{it:field-extension} the map $b\mapsto\mathbb B(b)$ is continuous from $L^2(\R^3;\C^3)$ to $L^2(\mu)$,
and by \eqref{eq:rotation-gradient-fields} it vanishes on gradients;
it therefore vanishes on the range of $1-P^\perp$, and $\mathbb B(b)=\mathbb B(P^\perp b)$ holds $\mu$-almost surely for every $b\in L^2(\R^3;\C^3)$.
For $b\in\mathcal S(\R^3;\R^3)$ the field $P^\perp b$ is real, the multiplier being real and even,
and Lemma~\ref{lem:domination-consequences}\eqref{it:field-extension} gives
\[
\mu\left(\exp\mathbb B(b)\right)=\mu\left(\exp\mathbb B(P^\perp b)\right)\le\exp\left(\frac12\left\|P^\perp b\right\|^2\right).
\]
Conversely, assume the bound.
Since $\left\|P^\perp b\right\|\le\|b\|$, it implies \textup{(G)}.
For real $u\in\mathcal S(\R^3)$ the field $b=\nabla u$ satisfies $P^\perp b=0$;
the bound for $t\cdot b$ and for $-t\cdot b$ gives $\mu\left(\cosh\left(t\cdot\mathbb B(b)\right)\right)\le1$ for every real $t$,
and as $\cosh\ge1$, with equality only at $0$, it follows that $\mathbb B(b)=0$ holds $\mu$-almost surely.
Thus $\mu\left(F(f)^2\right)=0$ for every $f\in\mathcal S(\R^3;\Lambda^2)$ with $\ast f=\nabla u$, $u\in\mathcal S(\R^3)$,
and $\mu$ is carried by the closed distributions \cite[Section~8.1.6]{NoncompactQED}, which is \textup{(dF)}.
\end{remark}

\subsection{Plates and the forward tube}\label{sec:plates-tube}

\begin{definition}[Labels and plates]\label{def:labels-plates}
Until the end of the proof of Proposition~\ref{prop:rotations-boosts}
the indices $j\neq i$ of Definition~\ref{def:closedness-cones} are fixed as $j=1$ and $i=2$.
The proof of Theorem~\ref{thm:one-plane} deduces the case of an arbitrary ordered pair from this one
by means of Lemma~\ref{lem:transport}\eqref{it:transport-pairs}.
For real $l$ and $t$ let
\[
U_l:=U(l\,e_1),\qquad\tau_l:=\tau_{l\,e_1},\qquad R_t:=R^{12}_t .
\]

\emph{Coordinates.}
A point of $\R^3$ is written $x=(y,x_3)$,
where $y=(y^0,y^1):=(x_1,x_2)\in\R^2$ is a point of the plane spanned by $e_1$ and $e_2$,
and $x_3$ is the transverse coordinate.
The coordinate $y^0$ is the time of the reconstruction at $\{x_1=0\}$.
In $U(y)$ and in $\tau_y$ a point $y\in\R^2$ stands for the vector $y^0e_1+y^1e_2$ of $\R^3$,
whereas the argument of $V_1$ is a vector of $e_1^\perp$.
For real $\upsilon$ let $n_\upsilon:=(\cos\upsilon,\sin\upsilon)\in\R^2$;
thus $n_0=(1,0)$ is the direction of time, and $n_{\pi/2}=(0,1)$.
The rotation $R_t$ maps the plane onto itself,
\[
R_t(y^0,y^1)=\left(y^0\cos t-y^1\sin t,\ y^0\sin t+y^1\cos t\right),
\]
and so do the reflections $\theta_1$ and $\theta_2$;
their restrictions $\theta_1(y^0,y^1)=(-y^0,y^1)$ and $\theta_2(y^0,y^1)=(y^0,-y^1)$ to the plane are denoted by the same letters.
The open half plane of positive times is $\mathbb{H}:=\{y\in\R^2:\ y^0>0\}$.

\emph{Eigenbasis, weights and signs.}
The vectors
\[
e^+:=e_1+i\,e_2,\qquad e^-:=e_1-i\,e_2,\qquad e^3:=e_3
\]
form a basis of $\C^3$, with $|e^\pm|=\sqrt2$ and $|e^3|=1$.
To $\sigma\in\{+,-,3\}$ are attached the weight $w(\sigma)$ and the sign $s(\sigma)$,
\[
w(+):=1,\quad w(-):=-1,\quad w(3):=0,\qquad s(+):=s(-):=1,\quad s(3):=-1,
\]
and $\sigma\mapsto\bar\sigma$ denotes the involution of $\{+,-,3\}$ that exchanges $+$ and $-$ and fixes $3$.

\emph{Labels and plates.}
A label is a pair $\lambda=(\sigma,\chi)$ of an element $\sigma$ of $\{+,-,3\}$ and a transverse profile $\chi\in L^2(\R;\R)$;
its conjugate label is $\bar\lambda:=(\bar\sigma,\chi)$.
For $\alpha\in L^2(\R^2;\C)$ the plate $\alpha\otimes\lambda\colon\R^3\to\C^3$ is the function
\[
(\alpha\otimes\lambda)(x):=\alpha(y)\cdot\chi(x_3)\cdot e^\sigma,\qquad x=(y,x_3),
\]
and the field smeared with it is
\[
\mathbb{B}_\lambda(\alpha):=\mathbb{B}(\alpha\otimes\lambda).
\]
For a tuple $\lambda=(\lambda_1,\dots,\lambda_N)$ of labels $\lambda_a=(\sigma_a,\chi_a)$ put
\[
w(\lambda):=\sum_{a=1}^{N}w(\sigma_a),\qquad s(\lambda):=\prod_{a=1}^{N}s(\sigma_a),\qquad\bar\lambda:=(\bar\lambda_1,\dots,\bar\lambda_N).
\]
\end{definition}

By Fubini's theorem a plate has the norm
\[
\|\alpha\otimes\lambda\|=|e^\sigma|\cdot\|\alpha\|_2\cdot\|\chi\|_2\le\sqrt2\cdot\|\alpha\|_2\cdot\|\chi\|_2 .
\]
So plates lie in $L^2(\R^3;\C^3)$,
and $\mathbb{B}_\lambda(\alpha)$ is an element of $L^2(\mu)$ for every law $\mu$ with \eqref{cond:gaussian-domination}
(Lemma~\ref{lem:domination-consequences}\eqref{it:field-extension});
if $\alpha$ and $\chi$ are Schwartz functions, $\mathbb{B}_\lambda(\alpha)$ is a function on $\mathcal{S}'$.
In terms of the components of the field:
with $u(x):=\alpha(y)\cdot\chi(x_3)$ one has $\alpha\otimes(\pm,\chi)=u\cdot e_1\pm i\,u\cdot e_2$ and $\alpha\otimes(3,\chi)=u\cdot e_3$,
hence, by complex linearity,
\[
\mathbb{B}_{(\pm,\chi)}(\alpha)=\mathbb{B}_1(u)\pm i\,\mathbb{B}_2(u),\qquad\mathbb{B}_{(3,\chi)}(\alpha)=\mathbb{B}_3(u),
\]
where $\mathbb{B}_3(u)$ is the component $F_{12}$ of the field in the plane, smeared with $u$.

\begin{lemma}[Plates under rotations and reflections]\label{lem:plate-identities}
Let $\lambda=(\sigma,\chi)$ be a label, $\alpha\in L^2(\R^2;\C)$ and $t\in\R$.
The reflections $\theta_1$ and $\theta_2$ act on $\C^3$ as the matrices $\diag(-1,1,1)$ and $\diag(1,-1,1)$,
on fields as Euclidean maps, and on the plane as in Definition~\ref{def:labels-plates}.
\begin{enumerate}
\item\label{it:plate-rotation}
One has $R_t\,e^\pm=\exp(\mp it)\cdot e^\pm$ and $R_t\,e^3=e^3$,
that is, $R_t\,e^\sigma=\exp\left(-i\,w(\sigma)\cdot t\right)\cdot e^\sigma$.
Hence
\[
R_t\cdot(\alpha\otimes\lambda)=\exp\left(-i\,w(\sigma)\cdot t\right)\cdot(\alpha\circ R_{-t})\otimes\lambda .
\]
\item\label{it:plate-reflection}
One has $(-\theta_1)\,e^\sigma=s(\sigma)\cdot\overline{e^\sigma}$.
Hence
\[
\overline{\theta_1\cdot(\alpha\otimes\lambda)}=s(\sigma)\cdot(\overline{\alpha}\circ\theta_1)\otimes\lambda,
\qquad
\Theta_1\mathbb{B}_\lambda(\alpha)=s(\sigma)\cdot\mathbb{B}_\lambda(\overline{\alpha}\circ\theta_1).
\]
\item\label{it:plate-second-reflection}
One has $(-\theta_2)\,e^\sigma=-\overline{e^\sigma}$ for each of the three values of $\sigma$.
Hence
\[
\Theta_2\mathbb{B}_\lambda(\alpha)=-\mathbb{B}_\lambda(\overline{\alpha}\circ\theta_2).
\]
\item\label{it:plate-conjugation}
One has $\overline{\mathbb{B}_\lambda(\alpha)}=\mathbb{B}_{\bar\lambda}(\overline{\alpha})$.
\item\label{it:plate-gradient}
Let $\partial_0$ and $\partial_1$ be the derivatives in $y^0$ and $y^1$, and
\[
\partial_z:=\frac12\left(\partial_0-i\,\partial_1\right),\qquad\partial_{\bar z}:=\frac12\left(\partial_0+i\,\partial_1\right).
\]
Let $\eta\in C^\infty_{\mathrm c}(\R^2;\C)$, let $\chi_0\in C^\infty_{\mathrm c}(\R;\R)$ have the derivative $\chi_0'$,
and let $u(x):=\eta(y)\cdot\chi_0(x_3)$.
Then
\[
\nabla u=(\partial_z\eta)\otimes(+,\chi_0)+(\partial_{\bar z}\eta)\otimes(-,\chi_0)+\eta\otimes(3,\chi_0').
\]
\end{enumerate}
The identities between functions of the field in \eqref{it:plate-reflection}, \eqref{it:plate-second-reflection} and \eqref{it:plate-conjugation}
hold at every point of $\mathcal{S}'$ if $\alpha$ and $\chi$ are Schwartz functions.
For a law $\mu$ with \eqref{cond:gaussian-domination} and arbitrary $\alpha$ and $\chi$ they hold $\mu$-almost surely:
the one in \eqref{it:plate-conjugation} always,
the one in \eqref{it:plate-reflection} if $\mu$ satisfies \eqref{cond:reflection-invariance} for $j=1$,
and the one in \eqref{it:plate-second-reflection} if $\mu$ satisfies \eqref{cond:reflection-invariance} for $j=2$.
\end{lemma}

\begin{proof}
For $g\in\mathrm{O}(3)$ the action on fields is $(g\cdot b)(x)=\det(g)\cdot g\,b(g^{-1}x)$ (Definition~\ref{def:field-actions}).

\eqref{it:plate-rotation}
Since $R_te_1=\cos t\cdot e_1+\sin t\cdot e_2$ and $R_te_2=-\sin t\cdot e_1+\cos t\cdot e_2$,
\[
R_t\,e^+=(\cos t-i\sin t)\cdot e_1+(\sin t+i\cos t)\cdot e_2=\exp(-it)\cdot e^+ .
\]
The matrix $R_t$ is real; complex conjugation gives $R_t\,e^-=\exp(it)\cdot e^-$, and $R_t\,e^3=e^3$ because $R_t$ fixes $e_3$.
As $\det R_t=1$ and $R_t^{-1}(y,x_3)=(R_{-t}y,x_3)$,
\[
\left(R_t\cdot(\alpha\otimes\lambda)\right)(x)=\alpha(R_{-t}y)\cdot\chi(x_3)\cdot R_t\,e^\sigma
=\exp\left(-i\,w(\sigma)\cdot t\right)\cdot\left((\alpha\circ R_{-t})\otimes\lambda\right)(x).
\]

\eqref{it:plate-reflection}
The matrix $-\theta_1=\diag(1,-1,-1)$ maps $e_1\pm i\,e_2$ to $e_1\mp i\,e_2=\overline{e_1\pm i\,e_2}$ and $e_3$ to $-e_3$;
this is the identity $(-\theta_1)\,e^\sigma=s(\sigma)\cdot\overline{e^\sigma}$.
Now $\det\theta_1=-1$, $\theta_1^{-1}=\theta_1$ and $\theta_1(y,x_3)=(\theta_1y,x_3)$, so that
\[
\left(\theta_1\cdot(\alpha\otimes\lambda)\right)(x)=\alpha(\theta_1y)\cdot\chi(x_3)\cdot(-\theta_1)\,e^\sigma
=s(\sigma)\cdot\alpha(\theta_1y)\cdot\chi(x_3)\cdot\overline{e^\sigma}.
\]
The numbers $s(\sigma)$ and $\chi(x_3)$ are real.
Complex conjugation therefore gives the first identity,
and \eqref{eq:rotation-conjugation} for $j=1$, that is, $\Theta_1\mathbb{B}(b)=\mathbb{B}\left(\overline{\theta_1\cdot b}\right)$, gives the second.
In contrast to $\theta_1$ itself, $\Theta_1$ preserves the label.
As $\overline{e^\pm}=e^\mp$ and $\overline{e^3}=e^3$, the last display says that $\theta_1\cdot(\alpha\otimes\lambda)=s(\sigma)\cdot(\alpha\circ\theta_1)\otimes\bar\lambda$:
the reflection, which inverts the orientation of the plane, replaces $\lambda$ by $\bar\lambda$,
and the complex conjugation in $\Theta_1$ restores $\lambda$.

\eqref{it:plate-second-reflection}
The matrix $-\theta_2=\diag(-1,1,-1)$ maps $e_1\pm i\,e_2$ to $-e_1\pm i\,e_2=-\overline{e_1\pm i\,e_2}$ and $e_3$ to $-e_3=-\overline{e_3}$.
With $\det\theta_2=-1$, $\theta_2^{-1}=\theta_2$ and $\theta_2(y,x_3)=(\theta_2y,x_3)$, the computation of \eqref{it:plate-reflection} gives
$\overline{\theta_2\cdot(\alpha\otimes\lambda)}=-(\overline{\alpha}\circ\theta_2)\otimes\lambda$,
and \eqref{eq:rotation-conjugation} for $j=2$ gives the claim.

\eqref{it:plate-conjugation}
One has $\overline{e^\pm}=e^\mp$ and $\overline{e^3}=e^3$, and $\chi$ is real; hence $\overline{\alpha\otimes\lambda}=\overline{\alpha}\otimes\bar\lambda$.
The claim follows from $\overline{\mathbb{B}(b)}=\mathbb{B}(\overline{b})$.

\eqref{it:plate-gradient}
Expanding the right side one checks that every $v=(v_1,v_2,v_3)\in\C^3$ satisfies
\[
v=\frac12\left(v_1-i\,v_2\right)\cdot e^++\frac12\left(v_1+i\,v_2\right)\cdot e^-+v_3\cdot e^3 .
\]
For $v=\nabla u(x)=\left(\chi_0(x_3)\cdot\partial_0\eta(y),\ \chi_0(x_3)\cdot\partial_1\eta(y),\ \eta(y)\cdot\chi_0'(x_3)\right)$
the three coefficients are $\chi_0(x_3)\cdot\partial_z\eta(y)$, $\chi_0(x_3)\cdot\partial_{\bar z}\eta(y)$ and $\eta(y)\cdot\chi_0'(x_3)$,
which is the asserted identity at the point $x$.

It remains to justify the use of \eqref{eq:rotation-conjugation} and of $\overline{\mathbb{B}(b)}=\mathbb{B}(\overline{b})$.
If $\alpha$ and $\chi$ are Schwartz functions, all fields that occur are Schwartz fields,
and both identities hold at every point of $\mathcal{S}'$ (Definition~\ref{def:field-actions}).
Let $\mu$ satisfy \eqref{cond:gaussian-domination}.
By Lemma~\ref{lem:domination-consequences}\eqref{it:field-extension}, $\overline{\mathbb{B}(b)}=\mathbb{B}(\overline{b})$ in $L^2(\mu)$ for every $b\in L^2(\R^3;\C^3)$;
if moreover $(\theta_j)_*\mu=\mu$, the same item gives $\hat\theta_j\mathbb{B}(b)=\mathbb{B}(\theta_j\cdot b)$ in $L^2(\mu)$,
and, after complex conjugation, $\Theta_j\mathbb{B}(b)=\overline{\hat\theta_j\mathbb{B}(b)}=\mathbb{B}\left(\overline{\theta_j\cdot b}\right)$ in $L^2(\mu)$ for every $b\in L^2(\R^3;\C^3)$.
\end{proof}

\begin{lemma}\label{lem:density-criterion}
Let $\mu$ be a law with \eqref{cond:gaussian-domination}, and let $S_N$ be its moments (Lemma~\ref{lem:domination-consequences}\eqref{it:moment-forms}).
Let $N\ge1$, let $\lambda=(\lambda_1,\dots,\lambda_N)$ be a tuple of labels, and let $t\in\R$.
Let $U\subset(\R^2)^N$ be an open set with $R_tU=U$, where $R_t$ acts on $(\R^2)^N$ entrywise,
and let $G\colon U\to\C$ be a continuous function such that
\[
S_N\left(\alpha_1\otimes\lambda_1,\dots,\alpha_N\otimes\lambda_N\right)
=\int_U\mathrm{d}Y\left(\prod_{a=1}^{N}\alpha_a(y_a)\right)G(Y),\qquad Y=(y_1,\dots,y_N),
\]
for all $\alpha_1,\dots,\alpha_N\in L^2(\R^2;\C)$ with compact supports and $\supp\alpha_1\times\dots\times\supp\alpha_N\subset U$.
Then the identity
\[
S_N\left(R_t\cdot(\alpha_1\otimes\lambda_1),\dots,R_t\cdot(\alpha_N\otimes\lambda_N)\right)=S_N\left(\alpha_1\otimes\lambda_1,\dots,\alpha_N\otimes\lambda_N\right)
\]
holds for all such $\alpha_1,\dots,\alpha_N$ if and only if
\begin{equation}\label{eq:rotation-density-criterion}
G(R_tY)=\exp\left(i\,w(\lambda)\cdot t\right)\cdot G(Y),\qquad Y\in U.
\end{equation}
\end{lemma}

\begin{proof}
All integrals below converge absolutely:
an element of $L^2(\R^2;\C)$ with compact support is integrable,
and $G$ is bounded on the compact set $\supp\alpha_1\times\dots\times\supp\alpha_N\subset U$.
Let $\alpha_1,\dots,\alpha_N$ be as in the hypothesis.
By Lemma~\ref{lem:plate-identities}\eqref{it:plate-rotation} and the multilinearity of $S_N$ (Lemma~\ref{lem:domination-consequences}\eqref{it:moment-forms}),
\[
\begin{aligned}
&S_N\left(R_t\cdot(\alpha_1\otimes\lambda_1),\dots,R_t\cdot(\alpha_N\otimes\lambda_N)\right)\\
&\qquad=\exp\left(-i\,w(\lambda)\cdot t\right)\cdot S_N\left((\alpha_1\circ R_{-t})\otimes\lambda_1,\dots,(\alpha_N\circ R_{-t})\otimes\lambda_N\right).
\end{aligned}
\]
The function $\alpha_a\circ R_{-t}$ has the compact support $R_t\supp\alpha_a$,
and the product of these supports is $R_t\left(\supp\alpha_1\times\dots\times\supp\alpha_N\right)\subset R_tU=U$;
so the hypothesis applies to the functions $\alpha_a\circ R_{-t}$.
The substitution $Y\mapsto R_tY$ preserves the Lebesgue measure of $(\R^2)^N$ and maps $U$ onto $U$.
Consequently
\[
\begin{aligned}
S_N\left((\alpha_1\circ R_{-t})\otimes\lambda_1,\dots,(\alpha_N\circ R_{-t})\otimes\lambda_N\right)
&=\int_U\mathrm{d}Y\left(\prod_{a=1}^{N}\alpha_a(R_{-t}y_a)\right)G(Y)\\
&=\int_U\mathrm{d}Y\left(\prod_{a=1}^{N}\alpha_a(y_a)\right)G(R_tY),
\end{aligned}
\]
and the left side minus the right side of the identity in question equals
\[
\int_U\mathrm{d}Y\left(\prod_{a=1}^{N}\alpha_a(y_a)\right)\left(\exp\left(-i\,w(\lambda)\cdot t\right)\cdot G(R_tY)-G(Y)\right).
\]
Under \eqref{eq:rotation-density-criterion} the integrand vanishes.
Conversely, suppose that the integral vanishes for all admissible $\alpha_1,\dots,\alpha_N$, in particular for the bounded ones.
The functions $Y\mapsto\exp\left(-i\,w(\lambda)\cdot t\right)\cdot G(R_tY)$ and $G$ are continuous on $U$,
so they coincide by Lemma~\ref{lem:analysis-facts}\eqref{it:equal-densities};
this is \eqref{eq:rotation-density-criterion}.
\end{proof}

\begin{lemma}[Plates are total]\label{lem:plates-total}
Let $A\subset\R^2$ be a measurable set.
\begin{enumerate}
\item\label{it:plates-total-square}
The finite sums of plates $\alpha\otimes\lambda$ with $\alpha=0$ outside $A$ are dense in the space of the fields $b\in L^2(\R^3;\C^3)$
with $b=0$ almost everywhere outside $A\times\R$.
\item\label{it:plates-total-smooth}
The same holds for the finite sums of such plates whose labels have profiles in $C^\infty_{\mathrm c}(\R;\R)$.
\end{enumerate}
\end{lemma}

\begin{proof}
A plate $\alpha\otimes\lambda$ with $\alpha=0$ outside $A$ vanishes outside $A\times\R$.

Consider first scalar functions.
Let $\alpha$ range over the elements of $L^2(\R^2;\C)$ that vanish outside $A$, and $\chi$ over $L^2(\R;\R)$.
The finite sums of the products $\alpha(y)\cdot\chi(x_3)$ form a complex vector space, although the profiles are real:
a complex multiple of $\alpha$ again vanishes outside $A$, so that complex coefficients can be absorbed into the first factor.
This space contains the indicator function of the product of a bounded measurable subset of $A$ with a bounded measurable subset of $\R$,
and it is dense in $L^2(A\times\R;\C)$:
a function in $L^2(A\times\R;\C)$ that is orthogonal to all of these indicator functions
has a vanishing integral over every bounded measurable subset of $A\times\R$ (Dynkin's $\pi$-$\lambda$ theorem),
and so it vanishes almost everywhere.

\eqref{it:plates-total-square}
Let $b\in L^2(\R^3;\C^3)$ vanish almost everywhere outside $A\times\R$, and let $\varepsilon>0$.
Since $e_1=\frac12\left(e^++e^-\right)$ and $e_2=\frac1{2i}\left(e^+-e^-\right)$,
\[
b=b_+\cdot e^++b_-\cdot e^-+b_3\cdot e^3,\qquad b_\pm:=\frac12\left(b_1\mp i\,b_2\right),
\]
where $b_1,b_2,b_3$ are the components of $b$; the three coefficients $b_\sigma$ lie in $L^2(A\times\R;\C)$.
By the preceding paragraph there are, for each $\sigma\in\{+,-,3\}$, finitely many $\alpha_{\sigma,k}\in L^2(\R^2;\C)$ that vanish outside $A$
and $\chi_{\sigma,k}\in L^2(\R;\R)$ such that $h_\sigma(x):=\sum_k\alpha_{\sigma,k}(y)\cdot\chi_{\sigma,k}(x_3)$ satisfies
\[
\|b_\sigma-h_\sigma\|_2\le\frac{\varepsilon}{3\sqrt2}.
\]
The finite sum of plates
\[
\sum_{\sigma\in\{+,-,3\}}\sum_k\alpha_{\sigma,k}\otimes(\sigma,\chi_{\sigma,k})=\sum_{\sigma\in\{+,-,3\}}h_\sigma\cdot e^\sigma
\]
has from $b$ the distance at most $\sum_\sigma|e^\sigma|\cdot\|b_\sigma-h_\sigma\|_2\le\varepsilon$, because $|e^\sigma|\le\sqrt2$ for each of the three values of $\sigma$.

\eqref{it:plates-total-smooth}
For fixed $\alpha$ and $\sigma$ the map $\chi\mapsto\alpha\otimes(\sigma,\chi)$ is linear,
and $\left\|\alpha\otimes(\sigma,\chi)\right\|=|e^\sigma|\cdot\|\alpha\|_2\cdot\|\chi\|_2$.
Since $C^\infty_{\mathrm c}(\R;\R)$ is dense in $L^2(\R;\R)$, replacing each profile in a finite sum of plates by a close profile in $C^\infty_{\mathrm c}(\R;\R)$
changes the sum by an arbitrarily small amount in norm, and the functions $\alpha$ are not changed.
So \eqref{it:plates-total-smooth} follows from \eqref{it:plates-total-square}.
\end{proof}

\paragraph{The complexified plane.}
For $\xi=(\xi^0,\xi^1)\in\C^2$ let
\[
\xi^\pm:=\xi^0\pm i\,\xi^1,\qquad d_{\mathcal{W}}(\xi):=\Re\xi^0-|\Im\xi^1|=\min\left(\Re\xi^+,\Re\xi^-\right),
\]
and $\Re\xi:=(\Re\xi^0,\Re\xi^1)\in\R^2$.
The linear forms $\xi\mapsto\xi^+$ and $\xi\mapsto\xi^-$ are coordinates on $\C^2$:
\[
\xi^0=\frac12\left(\xi^++\xi^-\right),\qquad\xi^1=\frac1{2i}\left(\xi^+-\xi^-\right).
\]
As $\Re\xi^\pm=\Re\xi^0\mp\Im\xi^1$ are real linear functions of $\xi$,
the function $d_{\mathcal{W}}\colon\C^2\to\R$ is positively homogeneous and superadditive,
\[
d_{\mathcal{W}}(\xi+\xi')\ge d_{\mathcal{W}}(\xi)+d_{\mathcal{W}}(\xi'),\qquad\xi,\xi'\in\C^2 .
\]
Therefore the forward tube
\[
\mathcal{W}:=\left\{\xi\in\C^2:\ d_{\mathcal{W}}(\xi)>0\right\}
\]
is an open convex cone, and its closure $\overline{\mathcal{W}}=\{\xi\in\C^2:\ d_{\mathcal{W}}(\xi)\ge0\}$ is a closed convex cone
(if $d_{\mathcal{W}}(\xi)\ge0$ and $\varepsilon>0$, then $d_{\mathcal{W}}(\xi+\varepsilon\,n_0)=d_{\mathcal{W}}(\xi)+\varepsilon>0$).
A real point $y\in\R^2$ has $d_{\mathcal{W}}(y)=y^0$, so that $\mathcal{W}\cap\R^2=\mathbb{H}$:
the real points of $\mathcal{W}$ are restricted in time only, and $y^1$ is arbitrary.
Further let
\[
\xi^*:=\left(\overline{\xi^0},-\overline{\xi^1}\right),\qquad\overline{\xi}:=\left(\overline{\xi^0},\overline{\xi^1}\right).
\]
For $\zeta\in\C$ let $R_\zeta$ be the complex linear map of $\C^2$ with
\[
(R_\zeta\xi)^\pm=\exp(\pm i\zeta)\cdot\xi^\pm .
\]
For real $\zeta=t$ it is the complex linear extension of the rotation $R_t$ of the plane:
for $y\in\R^2$ the formula for $R_t$ gives $(R_ty)^+=\exp(it)\cdot\left(y^0+i\,y^1\right)$, and $(R_ty)^-$ is the complex conjugate of this number.
For real $\beta$ the boost is $\Lambda_\beta:=R_{i\beta}$, and for real $\upsilon$ the rotated tube is $\mathcal{W}_\upsilon:=R_\upsilon\mathcal{W}$.
For a positive integer $N$ the elements $Z=(z_1,\dots,z_N)$ of $(\C^2)^N$ are called configurations.
Every map defined above on points of $\C^2$, or of $\R^2$, is applied to a configuration entry by entry;
for instance $\Lambda_\beta Z=(\Lambda_\beta z_1,\dots,\Lambda_\beta z_N)$.
The letters $X$ and $Y=(y_1,\dots,y_N)$ denote real configurations, that is, elements of $(\R^2)^N$,
and the real configurations whose $N$ entries are pairwise distinct form the subset $(\R^2)^N_{\neq}$.

\begin{lemma}[The forward tube]\label{lem:forward-tube}
Let $\xi\in\C^2$ and $\beta\in\R$.
\begin{enumerate}
\item\label{it:tube-adjoint}
One has $(\xi^*)^\pm=\overline{\xi^\pm}$;
hence $d_{\mathcal{W}}(\xi^*)=d_{\mathcal{W}}(\xi)$, and $\xi^*\in\mathcal{W}$ if and only if $\xi\in\mathcal{W}$.
One has $(\overline{\xi})^\pm=\overline{\xi^\mp}$;
hence $d_{\mathcal{W}}(\overline{\xi})=d_{\mathcal{W}}(\xi)$, and $\overline{\xi}\in\mathcal{W}$ if and only if $\xi\in\mathcal{W}$.
If $\xi$ is real, then $\xi^*=-\theta_1\xi$.
\item\label{it:tube-boosts}
One has $\Re(\Lambda_\beta\xi)^\pm=\exp(\mp\beta)\cdot\Re\xi^\pm$.
Hence $\Lambda_\beta\mathcal{W}=\mathcal{W}$, $(\Lambda_\beta\xi)^*=\Lambda_\beta(\xi^*)$,
$\overline{\Lambda_\beta y}=\Lambda_{-\beta}y$ for $y\in\R^2$, and
\[
d_{\mathcal{W}}(\Lambda_\beta\xi)\ge\exp(-|\beta|)\cdot d_{\mathcal{W}}(\xi)\quad\text{if }\xi\in\overline{\mathcal{W}},
\qquad
d_{\mathcal{W}}(\Lambda_\beta\xi)\ge\exp(|\beta|)\cdot d_{\mathcal{W}}(\xi)\quad\text{if }\xi\notin\overline{\mathcal{W}}.
\]
\item\label{it:boost-components}
One has
\[
(\Lambda_\beta\xi)^0=\xi^0\cdot\cosh\beta-i\,\xi^1\cdot\sinh\beta,\qquad(\Lambda_\beta\xi)^1=\xi^1\cdot\cosh\beta+i\,\xi^0\cdot\sinh\beta .
\]
\item\label{it:rotated-tubes}
Let $\upsilon\in\R$.
A real point $\xi$ lies in $\mathcal{W}_\upsilon$ if and only if $n_\upsilon\cdot\xi>0$.
If $\xi\in\mathcal{W}_\upsilon$, then $\Re\xi\in\mathcal{W}_\upsilon$.
Moreover $\Lambda_\beta\mathcal{W}_\upsilon=\mathcal{W}_\upsilon$.
\item\label{it:complex-rotation}
If $\xi$ is real, $\xi=|\xi|\,n_\vartheta$ with $\vartheta\in\R$, and $\zeta=t+i\beta$ with $t\in\R$, then
\[
\Re(R_\zeta\xi)^\pm=|\xi|\cdot\exp(\mp\beta)\cdot\cos(\vartheta+t).
\]
\end{enumerate}
\end{lemma}

\begin{proof}
\eqref{it:tube-adjoint}
One computes $(\xi^*)^\pm=\overline{\xi^0}\mp i\,\overline{\xi^1}=\overline{\xi^0\pm i\,\xi^1}=\overline{\xi^\pm}$
and $(\overline{\xi})^\pm=\overline{\xi^0}\pm i\,\overline{\xi^1}=\overline{\xi^0\mp i\,\xi^1}=\overline{\xi^\mp}$.
Complex conjugation does not change real parts, and $d_{\mathcal{W}}$ is the smaller of the real parts of $\xi^+$ and $\xi^-$;
so $d_{\mathcal{W}}(\xi^*)=d_{\mathcal{W}}(\overline{\xi})=d_{\mathcal{W}}(\xi)$, which gives the two equivalences.
For real $\xi$ one has $\xi^*=(\xi^0,-\xi^1)=-\theta_1\xi$.

\eqref{it:tube-boosts}
By definition $(\Lambda_\beta\xi)^\pm=\exp(\pm i\cdot i\beta)\cdot\xi^\pm=\exp(\mp\beta)\cdot\xi^\pm$, with a real factor;
this gives the identity for the real parts.
If $\xi\in\mathcal{W}$, both numbers $\Re(\Lambda_\beta\xi)^\pm$ are positive.
Hence $\Lambda_\beta\mathcal{W}\subset\mathcal{W}$ for every real $\beta$, and $\Lambda_{-\beta}=\Lambda_\beta^{-1}$ gives equality.
Next, by \eqref{it:tube-adjoint},
$\left((\Lambda_\beta\xi)^*\right)^\pm=\overline{\exp(\mp\beta)\cdot\xi^\pm}=\exp(\mp\beta)\cdot(\xi^*)^\pm=(\Lambda_\beta\xi^*)^\pm$;
and for real $y$, where $\overline{y^\mp}=y^\pm$,
$\left(\overline{\Lambda_\beta y}\right)^\pm=\overline{(\Lambda_\beta y)^\mp}=\exp(\pm\beta)\cdot y^\pm=(\Lambda_{-\beta}y)^\pm$.
For the inequalities note that $\exp(-|\beta|)\le\exp(\mp\beta)\le\exp(|\beta|)$.
If $\xi\in\overline{\mathcal{W}}$, then $\Re\xi^\pm\ge d_{\mathcal{W}}(\xi)\ge0$,
and $\Re(\Lambda_\beta\xi)^\pm\ge\exp(-|\beta|)\cdot\Re\xi^\pm\ge\exp(-|\beta|)\cdot d_{\mathcal{W}}(\xi)$ for both signs.
If $d_{\mathcal{W}}(\xi)<0$, then for each sign either $\Re\xi^\pm\ge0$, and $\Re(\Lambda_\beta\xi)^\pm\ge0>\exp(|\beta|)\cdot d_{\mathcal{W}}(\xi)$,
or $\Re\xi^\pm<0$, and $\Re(\Lambda_\beta\xi)^\pm\ge\exp(|\beta|)\cdot\Re\xi^\pm\ge\exp(|\beta|)\cdot d_{\mathcal{W}}(\xi)$.
The minimum over the two signs gives the two inequalities.

\eqref{it:boost-components}
Inserting $(\Lambda_\beta\xi)^\pm=\exp(\mp\beta)\cdot\xi^\pm$ and $\xi^\pm=\xi^0\pm i\,\xi^1$ into the expressions of the two coordinates through $\xi^+$ and $\xi^-$,
\[
\begin{aligned}
(\Lambda_\beta\xi)^0&=\frac12\left(\exp(-\beta)\cdot\xi^++\exp(\beta)\cdot\xi^-\right)=\xi^0\cdot\cosh\beta-i\,\xi^1\cdot\sinh\beta,\\
(\Lambda_\beta\xi)^1&=\frac1{2i}\left(\exp(-\beta)\cdot\xi^+-\exp(\beta)\cdot\xi^-\right)=\xi^1\cdot\cosh\beta+i\,\xi^0\cdot\sinh\beta .
\end{aligned}
\]

\eqref{it:rotated-tubes}
For real $\xi$ the point $R_{-\upsilon}\xi$ is real, and its first coordinate is $\xi^0\cos\upsilon+\xi^1\sin\upsilon=n_\upsilon\cdot\xi$;
a real point lies in $\mathcal{W}$ if and only if its first coordinate is positive.
Hence $\xi\in\mathcal{W}_\upsilon$, that is, $R_{-\upsilon}\xi\in\mathcal{W}$, holds if and only if $n_\upsilon\cdot\xi>0$.
If $\xi\in\mathcal{W}$, then $d_{\mathcal{W}}(\Re\xi)=\Re\xi^0>|\Im\xi^1|\ge0$, so that $\Re\xi\in\mathcal{W}$.
For general $\upsilon$ the matrix of $R_{-\upsilon}$ is real, so $\Re(R_{-\upsilon}\xi)=R_{-\upsilon}\Re\xi$;
if $\xi\in\mathcal{W}_\upsilon$, then $R_{-\upsilon}\xi\in\mathcal{W}$, hence $R_{-\upsilon}\Re\xi\in\mathcal{W}$, which means $\Re\xi\in\mathcal{W}_\upsilon$.
Finally, $\Lambda_\beta$ and $R_\upsilon$ act diagonally in the coordinates $(\xi^+,\xi^-)$ and therefore commute;
by \eqref{it:tube-boosts}, $\Lambda_\beta\mathcal{W}_\upsilon=R_\upsilon\Lambda_\beta\mathcal{W}=R_\upsilon\mathcal{W}=\mathcal{W}_\upsilon$.

\eqref{it:complex-rotation}
For $\xi=|\xi|\,n_\vartheta$ one has $\xi^\pm=|\xi|\cdot(\cos\vartheta\pm i\sin\vartheta)=|\xi|\cdot\exp(\pm i\vartheta)$, and therefore
\[
(R_\zeta\xi)^\pm=|\xi|\cdot\exp\left(\pm i(t+i\beta)\right)\cdot\exp(\pm i\vartheta)=|\xi|\cdot\exp(\mp\beta)\cdot\exp\left(\pm i(\vartheta+t)\right),
\]
whose real part is $|\xi|\cdot\exp(\mp\beta)\cdot\cos(\vartheta+t)$.
\end{proof}

\subsection{Analytic facts and transport}\label{sec:analytic-transport}

\begin{lemma}[Facts from analysis]\label{lem:analysis-facts}
Let $d$ and $N$ be positive integers.
\begin{enumerate}
\item\label{it:identity-theorem}
Let $D\subset\C^d$ be open and connected, and let $f$ be holomorphic on $D$.
If $f$ vanishes on a non-empty subset of $D\cap\R^d$ that is open in $\R^d$, then $f=0$ on $D$.
This statement, and for $d=1$ the classical fact that $f=0$ on $D$ as soon as the zeros of $f$ have a limit point in $D$,
are both referred to as the identity theorem.
\item\label{it:analytic-zeros}
Let $D\subset\R^d$ be open and connected, and let $f$ be real analytic on $D$.
If $f=0$ on a non-empty open subset of $D$, then $f=0$ on $D$.
\item\label{it:connected-complement}
Let $E_1,\dots,E_n$ be affine subspaces of $\R^d$ of codimension at least two.
Then $\R^d\setminus\left(E_1\cup\dots\cup E_n\right)$ is open, dense and connected.
\item\label{it:positive-kernels}
Let $X$ be a set.
A function $k\colon X\times X\to\C$ is a positive semidefinite kernel if $k(v,u)=\overline{k(u,v)}$ for all $u,v\in X$ and
\[
\sum_{p,q=1}^{n}\overline{c_p}\cdot c_q\cdot k(u_p,u_q)\ge0
\]
for every integer $n\ge1$, all $u_1,\dots,u_n\in X$ and all $c_1,\dots,c_n\in\C$.
Such a kernel satisfies $|k(u,v)|^2\le k(u,u)\cdot k(v,v)$ for all $u,v\in X$.
A pointwise limit of positive semidefinite kernels is a positive semidefinite kernel,
and so is $(u,v)\mapsto g(u)\cdot k(u,v)\cdot\overline{g(v)}$ for every positive semidefinite kernel $k$ and every function $g\colon X\to\C$.
\item\label{it:exponential-solutions}
Let $A$ be a finite set of real numbers, and let $f$ be a smooth function on $\R$ with
\[
\prod_{a\in A}\left(\frac{\mathrm{d}}{\mathrm{d}\beta}+a\right)f=0 .
\]
Then there are numbers $c_a$, $a\in A$, each of them a linear combination of the derivatives $f^{(l)}(0)$, $0\le l<\#|A|$,
with coefficients that depend on $A$ only, such that
\[
f(\beta)=\sum_{a\in A}c_a\cdot\exp(-a\cdot\beta)\qquad(\beta\in\R).
\]
The functions $\beta\mapsto\exp(-a\cdot\beta)$, $a\in A$, are linearly independent on every interval of positive length.
\item\label{it:polynomial-zeros}
The zero set of a real polynomial on $\R^d$ that is not the zero polynomial is closed and has empty interior.
\item\label{it:equal-densities}
Let $U\subset(\R^d)^N$ be open, write $Y=(y_1,\dots,y_N)$ for the points of $(\R^d)^N$,
and let $G_1$ and $G_2$ be continuous functions on $U$.
Suppose that
\[
\int\mathrm{d}Y\,\prod_{i=1}^{N}\alpha_i(y_i)\cdot\left(G_1-G_2\right)(Y)=0
\]
for all bounded measurable functions $\alpha_1,\dots,\alpha_N$ on $\R^d$ with compact supports such that $\supp\alpha_1\times\dots\times\supp\alpha_N\subset U$;
the integrand vanishes outside this compact set.
Then $G_1=G_2$.
The same conclusion holds if the identity is assumed only for such functions $\alpha_1,\dots,\alpha_N$ that belong to $C^\infty_{\mathrm c}(\R^d)$.
\item\label{it:mean-value}
Let $r>0$, and let $u$ be holomorphic on a neighbourhood of the closed polydisc $\{z\in\C^2:|z_1|\le r,\ |z_2|\le r\}$.
Then
\[
|u(0)|^2\le\left(\pi r^2\right)^{-2}\int_{\{|z_1|<r\}}\mathrm{d}^2z_1\int_{\{|z_2|<r\}}\mathrm{d}^2z_2\,\left|u(z_1,z_2)\right|^2,
\]
where $\mathrm{d}^2z$ denotes the Lebesgue measure of $\C$.
\item\label{it:normal-families}
Let $D\subset\C^d$ be open, and let $C\ge0$.
Every sequence of holomorphic functions on $D$ whose absolute values are bounded by $C$ on $D$ has a subsequence that converges uniformly on each compact subset of $D$;
thus the holomorphic functions $f$ on $D$ with $|f|\le C$ form a normal family (Montel's theorem).
\end{enumerate}
\end{lemma}

\begin{proof}
\eqref{it:identity-theorem}
Let $x_0$ be a point of a non-empty subset of $D\cap\R^d$, open in $\R^d$, on which $f$ vanishes.
Since $f$ is holomorphic, its complex partial derivatives of every order at $x_0$ are the partial derivatives at $x_0$ of the restriction of $f$ to $D\cap\R^d$,
and these vanish, because this restriction vanishes on a neighbourhood of $x_0$ in $\R^d$.
So all Taylor coefficients of $f$ at $x_0$ vanish, and $f=0$ on a polydisc about $x_0$.
The set of the points of $D$ near which $f$ vanishes identically is therefore not empty.
It is open, and it is closed in $D$:
at a point of $D$ that lies in its closure all derivatives of $f$ vanish by continuity, and then $f$ vanishes on a polydisc about that point.
As $D$ is connected, this set is $D$.

\eqref{it:analytic-zeros}
The set of the points at which all partial derivatives of the function vanish is closed by continuity,
open because near each point the function is the sum of its Taylor series,
and not empty by hypothesis.
By connectedness it is the whole domain.

\eqref{it:connected-complement}
Each $E_k$ is a closed affine subspace of dimension at most $d-2$,
and a finite union of proper affine subspaces is closed and has empty interior;
so the complement is open and dense.
Let $x$ and $x'$ be points of the complement.
If the segment from $x$ to a point $x''$ meets $E_k$, then $x''$ lies in the affine span of $\{x\}\cup E_k$,
whose dimension is at most $\dim E_k+1\le d-1$;
the same holds with $x'$ in place of $x$.
Choose $x''$ outside these $2n$ proper affine subspaces, which is possible because their union has empty interior;
then $x''$ lies in the complement as well, since each $E_k$ is contained in the span of $\{x\}\cup E_k$.
The segments from $x$ to $x''$ and from $x''$ to $x'$ lie in the complement, which is therefore connected.

\eqref{it:positive-kernels}
Let $u,v\in X$.
The case $n=2$, $u_1=u$, $u_2=v$ of the definition says that the Hermitian matrix with the rows $\left(k(u,u),k(u,v)\right)$ and $\left(k(v,u),k(v,v)\right)$ is positive semidefinite;
hence its determinant $k(u,u)\cdot k(v,v)-|k(u,v)|^2$ is non-negative.
The symmetry $k(v,u)=\overline{k(u,v)}$ and the inequalities of the definition are preserved under pointwise limits.
Finally, for $c_1,\dots,c_n\in\C$ and $c'_p:=c_p\cdot\overline{g(u_p)}$,
\[
\sum_{p,q=1}^{n}\overline{c_p}\cdot c_q\cdot g(u_p)\cdot k(u_p,u_q)\cdot\overline{g(u_q)}
=\sum_{p,q=1}^{n}\overline{c'_p}\cdot c'_q\cdot k(u_p,u_q)\ge0,
\]
and $g(v)\cdot k(v,u)\cdot\overline{g(u)}$ is the complex conjugate of $g(u)\cdot k(u,v)\cdot\overline{g(v)}$.

\eqref{it:exponential-solutions}
The operators $\frac{\mathrm{d}}{\mathrm{d}\beta}+a$, $a\in A$, commute, and $\left(\frac{\mathrm{d}}{\mathrm{d}\beta}+a\right)\exp(-a\cdot\beta)=0$.
So $\sum_{a\in A}c_a\cdot\exp(-a\cdot\beta)$ solves the equation for any numbers $c_a$,
and its derivative of order $l$ at $0$ is $\sum_{a\in A}c_a\cdot(-a)^l$.
The square matrix $\left((-a)^l\right)$, with rows indexed by $0\le l<\#|A|$ and columns indexed by $a\in A$,
is a Vandermonde matrix with the pairwise distinct nodes $-a$; hence it is invertible.
Its inverse, applied to the vector $\left(f^{(l)}(0)\right)_{0\le l<\#|A|}$, yields numbers $c_a$ of the asserted form
for which this function and $f$ have the same derivatives of order less than $\#|A|$ at $0$.
Their difference solves a linear differential equation of order $\#|A|$ with constant coefficients and has vanishing initial values at $0$;
by the uniqueness theorem for such equations it vanishes.
If $\sum_{a\in A}c_a\cdot\exp(-a\cdot\beta)=0$ on an interval of positive length, and $\beta_0$ is an interior point of this interval,
then the derivatives of order $l<\#|A|$ at $\beta_0$ give $\sum_{a\in A}(-a)^l\cdot c_a\cdot\exp(-a\cdot\beta_0)=0$,
and the invertibility of the same matrix gives $c_a\cdot\exp(-a\cdot\beta_0)=0$, that is, $c_a=0$, for every $a\in A$.

\eqref{it:polynomial-zeros}
The zero set is closed, the polynomial being continuous.
A polynomial is real analytic on the connected set $\R^d$.
If it vanished on a non-empty open set, it would vanish on $\R^d$ by \eqref{it:analytic-zeros},
and then all its coefficients, which are its partial derivatives at $0$ divided by factorials, would vanish.

\eqref{it:equal-densities}
Let $Y'=(y'_1,\dots,y'_N)\in U$, and let $\varepsilon>0$ be so small that the product of the closed balls $\overline{B(y'_i,\varepsilon)}$, $i=1,\dots,N$, lies in $U$.
For $\alpha_i$ take the indicator function of the open ball $B(y'_i,\varepsilon)$ divided by the volume of this ball,
or, in the second case, a non-negative function in $C^\infty_{\mathrm c}(\R^d)$ with support in $\overline{B(y'_i,\varepsilon)}$ and integral one.
Then $\prod_{i}\alpha_i(y_i)\,\mathrm{d}Y$ is a probability measure carried by the product of the closed balls,
and the hypothesis gives
\[
\begin{aligned}
\left|\left(G_1-G_2\right)(Y')\right|
&=\left|\int\mathrm{d}Y\,\prod_{i=1}^{N}\alpha_i(y_i)\cdot\left(\left(G_1-G_2\right)(Y')-\left(G_1-G_2\right)(Y)\right)\right|\\
&\le\sup_{Y}\left|\left(G_1-G_2\right)(Y')-\left(G_1-G_2\right)(Y)\right|,
\end{aligned}
\]
the supremum being taken over $Y$ in the product of the closed balls.
As $\varepsilon\to0$, the right side tends to $0$, because $G_1-G_2$ is continuous at $Y'$.
Hence $G_1(Y')=G_2(Y')$.

\eqref{it:mean-value}
Let $h$ be holomorphic on a neighbourhood of a closed disc $\{|\zeta|\le r\}$ of $\C$.
Its mean over each circle $\{|\zeta|=s\}$, $0<s\le r$, is $h(0)$, and integration over $s$ gives
\[
h(0)=\left(\pi r^2\right)^{-1}\int_{\{|\zeta|<r\}}\mathrm{d}^2\zeta\,h(\zeta).
\]
Apply this to $z_2\mapsto u(0,z_2)^2$, and to $z_1\mapsto u(z_1,z_2)^2$ for each $z_2$ with $|z_2|<r$.
By Fubini's theorem
\[
u(0)^2=\left(\pi r^2\right)^{-2}\int_{\{|z_1|<r\}}\mathrm{d}^2z_1\int_{\{|z_2|<r\}}\mathrm{d}^2z_2\,u(z_1,z_2)^2,
\]
and $|u(0)|^2=\left|u(0)^2\right|$ is at most the same expression with $|u|^2$ in place of $u^2$.

\eqref{it:normal-families}
This is Montel's theorem.
In several variables, as in one, the Cauchy estimates on polydiscs show that a uniformly bounded family of holomorphic functions is equicontinuous on every compact subset,
and the Arzel\`a--Ascoli theorem, together with a diagonal sequence over an exhaustion by compact sets, provides the subsequence.
\end{proof}

\begin{lemma}[Harmonics of a quasi-periodic entire function]\label{lem:harmonics}
Let $w$ be an integer, let $a\in\{\pi/2,\pi\}$ be the quasi-period, $n\ge0$ the exponent of growth and $C\ge0$ a constant,
and let $h$ be an entire function with
\[
h(\zeta+a)=\exp(iw\cdot a)\cdot h(\zeta)\quad(\zeta\in\C),\qquad
\left|h(t+i\beta)\right|\le C\cdot\exp(n\cdot|\beta|)\quad(t,\beta\in\R).
\]
Then
\[
h(\zeta)=\sum_{k}a_k\cdot\exp(ik\cdot\zeta)\qquad(\zeta\in\C),
\]
where $k$ ranges over the integers with $k-w\in(2\pi/a)\Z$ and $|k|\le n$; for each such $k$
\[
a_k=\frac{1}{a}\int_0^a\mathrm{d}t\,h(t+i\beta)\cdot\exp\left(-ik\cdot(t+i\beta)\right)\quad\text{for every }\beta\in\R,
\qquad\text{and}\qquad|a_k|\le C.
\]
\end{lemma}

\begin{proof}
By the first hypothesis, the entire function $\zeta\mapsto\exp(-iw\cdot\zeta)\cdot h(\zeta)$ has the period $a$.
The map $\zeta\mapsto\exp\left(\frac{2\pi i\cdot\zeta}{a}\right)$ takes $\C$ onto $\C\setminus\{0\}$;
two points have the same image exactly when they differ by an element of $a\Z$,
and the map has holomorphic local inverses, namely the branches of $z\mapsto\frac{a}{2\pi i}\cdot\log z$.
Hence there is a holomorphic function $H$ on $\C\setminus\{0\}$ with
\[
\exp(-iw\cdot\zeta)\cdot h(\zeta)=H\left(\exp\left(\frac{2\pi i\cdot\zeta}{a}\right)\right)\qquad(\zeta\in\C).
\]
The Laurent series of $H$ converges uniformly on every circle about $0$.
For $m\in\Z$ put $k_m:=w+\frac{2\pi\cdot m}{a}$, an integer since $\frac{2\pi}{a}\in\{2,4\}$,
and denote the coefficient of $z^m$ in this series by $a_{k_m}$:
\[
H(z)=\sum_{m\in\Z}a_{k_m}\cdot z^m\qquad\left(z\in\C\setminus\{0\}\right).
\]
The substitution $z=\exp\left(\frac{2\pi i\cdot\zeta}{a}\right)$ gives $h(\zeta)=\sum_{m\in\Z}a_{k_m}\cdot\exp(ik_m\cdot\zeta)$,
a sum over the integers $k\in w+(2\pi/a)\Z$.
Let $\beta\in\R$.
The circle of radius $\exp\left(-\frac{2\pi\cdot\beta}{a}\right)$ about $0$ is traversed once by $z=\exp\left(\frac{2\pi i\cdot(t+i\beta)}{a}\right)$, $0\le t\le a$,
and $\mathrm{d}z=\frac{2\pi i}{a}\cdot z\,\mathrm{d}t$ along it.
At this point $z$ one has $H(z)\cdot z^{-m}=h(t+i\beta)\cdot\exp\left(-ik_m\cdot(t+i\beta)\right)$,
so Cauchy's formula for the Laurent coefficients,
\[
a_{k_m}=\frac{1}{2\pi i}\oint\mathrm{d}z\,H(z)\cdot z^{-m-1},
\]
taken on this circle, becomes the asserted formula for $a_k$ with $k=k_m$.
Since $\left|\exp\left(-ik\cdot(t+i\beta)\right)\right|=\exp(k\cdot\beta)$, the second hypothesis gives
\[
|a_k|\le C\cdot\exp\left(n\cdot|\beta|+k\cdot\beta\right)\qquad(\beta\in\R).
\]
If $k>n$, the right side tends to $0$ as $\beta\to-\infty$; if $k<-n$, it tends to $0$ as $\beta\to+\infty$.
So $a_k=0$ unless $|k|\le n$, and the choice $\beta=0$ gives $|a_k|\le C$.
\end{proof}

\begin{lemma}[Transport by orthogonal maps]\label{lem:transport}
Let $g\in\mathrm{O}(3)$, let $\mu$ be a law on $\mathcal{S}'$, and put $\mu':=\left(g^{-1}\right)_*\mu$.
\begin{enumerate}
\item\label{it:transport-operators}
One has $\mu'(\Phi)=\mu(\hat g\Phi)$ for every measurable function $\Phi$ on $\mathcal{S}'$ for which one of the two sides is defined,
and $\hat g$ is a unitary operator from $L^2(\mu')$ onto $L^2(\mu)$
that maps $L^2(\Sigma_U,\mu')$ onto $L^2(\Sigma_{gU},\mu)$ for every open set $U\subset\R^3$.
For every $a\in\R^3$ and every Euclidean map $\mathfrak{g}$,
\[
\hat g\,U(a)\,\hat g^{-1}=U(ga),\qquad\hat g\,\hat{\mathfrak{g}}\,\hat g^{-1}=\widehat{g\mathfrak{g}g^{-1}},
\]
and $\hat g\,\mathbb{B}(b)=\mathbb{B}(g\cdot b)$ for every $b\in\mathcal{S}(\R^3;\C^3)$.
\item\label{it:transport-properties}
Each of the conditions \textup{(G)}, \textup{(dF)} and $(T_{123})$ holds for $\mu'$ exactly when it holds for $\mu$.
For every Euclidean map $\mathfrak{g}$ one has $\mathfrak{g}_*\mu'=\mu'$ if and only if $\left(g\mathfrak{g}g^{-1}\right)_*\mu=\mu$.
For every unit vector $n\in\R^3$ and every $c\in\R$, the law $\mu'$ is reflection positive at $\{n\cdot x=c\}$
if and only if $\mu$ is reflection positive at $\{(gn)\cdot x=c\}$.
\item\label{it:transport-all-planes}
Suppose that $\mu$ is invariant under the Euclidean group $\R^3\rtimes\mathrm{O}(3)$ and satisfies $(RP_1)$.
Then, for every unit vector $n\in\R^3$ and every $c\in\R$,
the law $\mu$ is invariant under $r_{n,c}$ and reflection positive at $\{n\cdot x=c\}$, and
\[
\left\langle\hat r_{n,c}\Phi,\Phi\right\rangle_\mu\ge0\qquad\text{for every }\Phi\in L^2\left(\Sigma_{\{n\cdot x>c\}},\mu\right).
\]
\item\label{it:transport-pairs}
Let $j\neq i$ be indices in $\{1,2,3\}$, let $g=g_{ji}$ be the permutation matrix with $g_{ji}e_1=e_j$ and $g_{ji}e_2=e_i$,
and suppose that $\mu$ satisfies $(T_{123})$, $(R_j)$ and $(RP_j)$.
Then $\mu'$ satisfies $(T_{123})$, $(R_1)$ and $(RP_1)$.
Let $\mathcal{H}_1(\mu')$, $[\,\cdot\,]_1$, $\mathsf{H}_1$ and $V_1$ be the space, the canonical map, the energy and the transverse translations of the reconstruction of $\mu'$ in the direction $1$,
$\mathsf{P}_2$ the generator of $s\mapsto V_1(s\cdot e_2)$ and $\mathfrak{E}_1$ the joint spectral measure of $\mu'$,
and let $\mathcal{H}_j(\mu)$, $[\,\cdot\,]_j$, $\mathsf{H}_j$, $V_j$, $\mathsf{P}_i$ and $\mathfrak{E}_j$ be the corresponding objects for $\mu$ in the direction $j$
(Definition~\ref{def:laws}, Lemma~\ref{lem:joint-spectrum});
when $j=1$, the symbols $[\,\cdot\,]_1$, $\mathsf{H}_1$, $V_1$ and $\mathfrak{E}_1$ still denote the objects of $\mu'$,
and the symbols $[\,\cdot\,]_j$, $\mathsf{H}_j$, $V_j$ and $\mathfrak{E}_j$ those of $\mu$, although the indices then agree.
The map $[\Phi]_1\mapsto[\hat g_{ji}\Phi]_j$, $\Phi\in L^2(\Sigma_1^+,\mu')$, extends to a unitary map $J$ of $\mathcal{H}_1(\mu')$ onto $\mathcal{H}_j(\mu)$ such that
\[
J\exp(-t\cdot\mathsf{H}_1)J^{-1}=\exp(-t\cdot\mathsf{H}_j)\quad(t\ge0),\qquad
JV_1(y)J^{-1}=V_j(g_{ji}y)\quad\left(y\in e_1^\perp\right),
\]
and such that, for every Borel set $A\subset[0,\infty)\times e_j^\perp$,
\[
\mathfrak{E}_j(A)=J\,\mathfrak{E}_1\left(\left\{(\omega,p_\perp):(\omega,g_{ji}p_\perp)\in A\right\}\right)J^{-1}.
\]
In particular $JV_1(s\cdot e_2)J^{-1}=V_j(s\cdot e_i)$ for real $s$,
and $J$ carries the joint spectral measure of $(\mathsf{H}_1,\mathsf{P}_2)$ to that of $(\mathsf{H}_j,\mathsf{P}_i)$.
If $\mu$ satisfies \textup{(G)}, then $J\mathcal{H}^1_1(\mu')=\mathcal{H}^1_j(\mu)$.
Consequently, $\mu'$ satisfies $(C_{12})$, resp.\ $(Q_{12})$, if and only if $\mu$ satisfies $(C_{ji})$, resp.\ $(Q_{ji})$;
if $\mu$ satisfies \textup{(G)}, then $\mu'$ satisfies $(C^1_{12})$ if and only if $\mu$ satisfies $(C^1_{ji})$;
$\mu'$ satisfies the hypotheses of Definition~\ref{def:closedness-cones}\eqref{it:one-plane-hypotheses} for the pair $(1,2)$,
namely \textup{(G)}, \textup{(dF)}, $(T_{123})$, $(R_1)$, $(RP_1)$, $(Q_{12})$ and $(C_{12})$,
if and only if $\mu$ satisfies them for the pair $(j,i)$;
and, for every real $t$, $\left(R^{12}_t\right)_*\mu'=\mu'$ if and only if $\left(R^{ji}_t\right)_*\mu=\mu$.
\end{enumerate}
\end{lemma}

\begin{proof}
\eqref{it:transport-operators}
By the definition of the image law,
$\mu'(\Phi)=\int\mathrm{d}\mu(F)\,\Phi\left(g^{-1}F\right)=\mu(\hat g\Phi)$ whenever one of the two sides is defined.
Applied to $|\Phi|^2$, this shows that $\hat g$ maps $L^2(\mu')$ isometrically into $L^2(\mu)$;
it is onto, with the inverse $\Phi\mapsto\Phi\circ g$.
The operator of a composition of two Euclidean maps is the composition of their operators,
the operators of Euclidean maps are multiplicative and commute with complex conjugation,
and $\hat g$ maps the $\Sigma_U$-measurable functions onto the $\Sigma_{gU}$-measurable functions,
because it maps $F\mapsto F(f)$ onto $F\mapsto F(g\cdot f)$ and $f\mapsto g\cdot f$ maps $C^\infty_{\mathrm c}(U;\Lambda^2)$ onto $C^\infty_{\mathrm c}(gU;\Lambda^2)$
\cite[Section~8.1.2]{NoncompactQED}.
This gives the assertion on $L^2(\Sigma_U,\mu')$, and, together with $g\tau_ag^{-1}=\tau_{ga}$ and $U(a)=\hat\tau_a$, the two displayed identities.
Finally, the identity $\hat g\,\mathbb{B}(b)=\mathbb{B}(g\cdot b)$ for $b\in\mathcal{S}(\R^3;\C^3)$
is the first identity of \eqref{eq:rotation-field-action} for the Euclidean map $x\mapsto gx$, whose operator is $\hat g$.

\eqref{it:transport-properties}
Let $f\in\mathcal{S}$.
By \eqref{it:transport-operators}, $\mu'\left(\exp F(f)\right)=\mu\left(\exp F(g\cdot f)\right)$,
and $\|g\cdot f\|_2=\|f\|_2$, since $\wedge^2g$ is orthogonal and $g$ preserves the Lebesgue measure.
So \textup{(G)} for $\mu$ implies \textup{(G)} for $\mu'$.
As $\mu=g_*\mu'$ arises from $\mu'$ in the same way, with $g^{-1}$ in place of $g$, the converse follows.

For $\mu$, and likewise for $\mu'$, condition \textup{(dF)} is equivalent to $\mu\left(\mathbb{B}(\nabla u)^2\right)=0$ for every real $u\in\mathcal{S}(\R^3)$:
the condition implies these identities by \eqref{eq:rotation-gradient-fields},
and they imply it by the criterion for closedness in \cite[Section~8.1.6]{NoncompactQED},
applied to the forms $f\in\mathcal{S}$ with $\ast f=\nabla u$, for which $F(f)=\mathbb{B}(\nabla u)$.
Since $\hat g$ is multiplicative, \eqref{it:transport-operators} and the chain rule \eqref{eq:rotation-gradient-action} give
\[
\mu'\left(\mathbb{B}(\nabla u)^2\right)=\mu\left(\mathbb{B}(g\cdot\nabla u)^2\right)=\mu\left(\mathbb{B}\left(\nabla\left(u\circ g^{-1}\right)\right)^2\right),
\]
because $g\cdot\nabla u=\det(g)\cdot\nabla\left(u\circ g^{-1}\right)$ and $\det(g)^2=1$.
As $u\mapsto u\circ g^{-1}$ is a bijection of the real Schwartz functions onto themselves,
$\mu'$ satisfies \textup{(dF)} if and only if $\mu$ does.

For a Euclidean map $\mathfrak{g}$,
\[
\mathfrak{g}_*\mu'=\left(\mathfrak{g}g^{-1}\right)_*\mu=\left(g^{-1}\right)_*\left(\left(g\mathfrak{g}g^{-1}\right)_*\mu\right),
\]
and $\left(g^{-1}\right)_*$ is a bijection of the set of laws onto itself, with the inverse $g_*$.
So $\mathfrak{g}_*\mu'=\mu'=\left(g^{-1}\right)_*\mu$ holds if and only if $\left(g\mathfrak{g}g^{-1}\right)_*\mu=\mu$.
For $\mathfrak{g}=\tau_a$ one has $g\tau_ag^{-1}=\tau_{ga}$, and $ga$ ranges over $\R^3$ together with $a$;
this gives the assertion on $(T_{123})$.

Let $n\in\R^3$ be a unit vector and $c\in\R$.
As $n\cdot g^{-1}x=(gn)\cdot x$ for $x\in\R^3$, the formula $r_{n,c}(x)=x-2\cdot(n\cdot x-c)\cdot n$ gives
$g\,r_{n,c}\,g^{-1}=r_{gn,c}$; moreover $g\{n\cdot x>c\}=\{(gn)\cdot x>c\}$.
So $f\mapsto g\cdot f$ maps the real forms in $C^\infty_{\mathrm c}(\{n\cdot x>c\};\Lambda^2)$ onto the real forms in $C^\infty_{\mathrm c}(\{(gn)\cdot x>c\};\Lambda^2)$.
For two forms $f$ and $f'$ of the first kind, \eqref{it:transport-operators} gives
\[
\mu'\left(\exp\left(i\left(F(f')-F(r_{n,c}\cdot f)\right)\right)\right)
=\mu\left(\exp\left(i\left(F(g\cdot f')-F\left(r_{gn,c}\cdot(g\cdot f)\right)\right)\right)\right),
\]
since $g\cdot\left(r_{n,c}\cdot f\right)=\left(g\,r_{n,c}\,g^{-1}\right)\cdot(g\cdot f)$.
Hence the sums that express the reflection positivity of $\mu'$ at $\{n\cdot x=c\}$, formed with $z_1,\dots,z_k\in\C$ and with forms $f_1,\dots,f_k$,
are the sums that express the reflection positivity of $\mu$ at $\{(gn)\cdot x=c\}$, formed with $z_1,\dots,z_k$ and with $g\cdot f_1,\dots,g\cdot f_k$;
and every family of real forms in $C^\infty_{\mathrm c}(\{(gn)\cdot x>c\};\Lambda^2)$ arises in this way.

\eqref{it:transport-all-planes}
Every reflection $r_{n,c}$ is a Euclidean map; so $(r_{n,c})_*\mu=\mu$.
Fix $n$ and $c$, and apply \eqref{it:transport-properties} to an element $g\in\mathrm{O}(3)$ with $ge_1=n$;
for this $g$ one has $\mu'=\left(g^{-1}\right)_*\mu=\mu$.
The law $\mu$ satisfies $(T_1)$, $(R_1)$ and $(RP_1)$, and so it is reflection positive at the plane $\{x_1=c\}$ \cite[Lemma~8.16\,(5)]{NoncompactQED};
by \eqref{it:transport-properties} it is reflection positive at $\{(ge_1)\cdot x=c\}=\{n\cdot x=c\}$.
Let $\mathfrak{g}:=\tau_{c\cdot n}\,g$, that is, $\mathfrak{g}(x)=gx+c\cdot n$.
Then $n\cdot\mathfrak{g}(x)=x_1+c$, and $\mathfrak{g}(\theta_1x)=gx-2x_1\cdot n+c\cdot n=r_{n,c}\left(\mathfrak{g}(x)\right)$ for $x\in\R^3$.
As the operator of a composition of Euclidean maps is the composition of their operators,
\begin{equation}\label{eq:rotation-plane-conjugation}
\mathfrak{g}\{x_1>0\}=\{n\cdot x>c\},\qquad
\mathfrak{g}\,\theta_1\,\mathfrak{g}^{-1}=r_{n,c},\qquad
\hat{\mathfrak{g}}\,\hat\theta_1\,\hat{\mathfrak{g}}^{-1}=\hat r_{n,c}.
\end{equation}
Since $\mathfrak{g}_*\mu=\mu$, the operator $\hat{\mathfrak{g}}$ is unitary on $L^2(\mu)$,
and by the first identity of \eqref{eq:rotation-plane-conjugation} it maps $L^2(\Sigma_1^+,\mu)$ onto $L^2\left(\Sigma_{\{n\cdot x>c\}},\mu\right)$ \cite[Section~8.1.2]{NoncompactQED}.
Let $\Phi\in L^2\left(\Sigma_{\{n\cdot x>c\}},\mu\right)$, and put $\Phi':=\hat{\mathfrak{g}}^{-1}\Phi\in L^2(\Sigma_1^+,\mu)$.
By the last identity of \eqref{eq:rotation-plane-conjugation},
\[
\left\langle\hat r_{n,c}\Phi,\Phi\right\rangle_\mu
=\left\langle\hat{\mathfrak{g}}\,\hat\theta_1\Phi',\hat{\mathfrak{g}}\Phi'\right\rangle_\mu
=\left\langle\hat\theta_1\Phi',\Phi'\right\rangle_\mu\ge0,
\]
where the inequality is \cite[Lemma~8.16\,(3)]{NoncompactQED}, in which $\hat\theta_1$ is written $\Theta_1$;
it applies because $\mu$ satisfies $(R_1)$ and $(RP_1)$.

\eqref{it:transport-pairs}
Write $g=g_{ji}$; it maps $e_3$ to the third unit vector.
By the computation in the proof of \eqref{it:transport-properties}, $g\theta_1g^{-1}=g\,r_{e_1,0}\,g^{-1}=r_{e_j,0}=\theta_j$, and $g\{x_1>0\}=\{x_j>0\}$.
So \eqref{it:transport-properties} shows that $\mu'$ satisfies $(T_{123})$,
and that $(R_1)$, resp.\ $(RP_1)$, holds for $\mu'$ if and only if $(R_j)$, resp.\ $(RP_j)$, holds for $\mu$.
In particular the reconstruction of $\mu'$ in the direction $1$ and its joint spectral measure are defined.

For $\Phi,\Phi'\in L^2(\Sigma_1^+,\mu')$ the functions $\hat g\Phi$ and $\hat g\Phi'$ lie in $L^2(\Sigma_j^+,\mu)$, by \eqref{it:transport-operators}.
From $\theta_jg=g\theta_1$ one gets $\hat\theta_j\hat g=\hat g\hat\theta_1$;
as $\hat g$ commutes with complex conjugation, $\Theta_j\hat g\Phi=\hat g\,\Theta_1\Phi$.
Since $\hat g$ is multiplicative, \eqref{it:transport-operators} gives
\begin{equation}\label{eq:rotation-transport-isometry}
\mu\left(\left(\Theta_j\hat g\Phi\right)\cdot\hat g\Phi'\right)
=\mu\left(\hat g\left(\left(\Theta_1\Phi\right)\cdot\Phi'\right)\right)
=\mu'\left(\left(\Theta_1\Phi\right)\cdot\Phi'\right).
\end{equation}
The outer members of \eqref{eq:rotation-transport-isometry} are the scalar products $\left\langle[\hat g\Phi]_j,[\hat g\Phi']_j\right\rangle$ in $\mathcal{H}_j(\mu)$
and $\left\langle[\Phi]_1,[\Phi']_1\right\rangle$ in $\mathcal{H}_1(\mu')$ (Definition~\ref{def:laws}).
Therefore $[\Phi]_1\mapsto[\hat g\Phi]_j$ is well defined and isometric on the classes, which are dense in $\mathcal{H}_1(\mu')$,
and it extends to an isometry $J$ of $\mathcal{H}_1(\mu')$ into $\mathcal{H}_j(\mu)$.
As $\hat g$ maps $L^2(\Sigma_1^+,\mu')$ onto $L^2(\Sigma_j^+,\mu)$, the range of $J$ contains every class $[\Phi'']_j$, $\Phi''\in L^2(\Sigma_j^+,\mu)$;
these classes are dense in $\mathcal{H}_j(\mu)$, and so $J$ is unitary.

Let $t\ge0$ and $y\in e_1^\perp$; then $gy\in e_j^\perp$.
The reconstruction of $\mu'$ satisfies $\exp(-t\cdot\mathsf{H}_1)[\Phi]_1=[U(t\cdot e_1)\Phi]_1$ and $V_1(y)[\Phi]_1=[U(y)\Phi]_1$ for $\Phi\in L^2(\Sigma_1^+,\mu')$,
and that of $\mu$ satisfies the corresponding relations in the direction $j$ (Definition~\ref{def:laws}).
With $\hat g\,U(a)=U(ga)\,\hat g$ from \eqref{it:transport-operators},
\[
\begin{aligned}
J\exp(-t\cdot\mathsf{H}_1)[\Phi]_1&=\left[U(t\cdot e_j)\,\hat g\Phi\right]_j=\exp(-t\cdot\mathsf{H}_j)\,J[\Phi]_1,\\
JV_1(y)[\Phi]_1&=\left[U(gy)\,\hat g\Phi\right]_j=V_j(gy)\,J[\Phi]_1 .
\end{aligned}
\]
The classes are dense and the operators are bounded;
hence $J\exp(-t\cdot\mathsf{H}_1)J^{-1}=\exp(-t\cdot\mathsf{H}_j)$ and $JV_1(y)J^{-1}=V_j(gy)$.
Since $\mathsf{H}_1=-\log\exp(-\mathsf{H}_1)$ and $\mathsf{H}_j=-\log\exp(-\mathsf{H}_j)$,
and since unitary equivalence is compatible with the functional calculus,
$J\mathsf{H}_1J^{-1}=\mathsf{H}_j$; so $J\exp(it\cdot\mathsf{H}_1)J^{-1}=\exp(it\cdot\mathsf{H}_j)$ for every real $t$.
For a Borel set $A\subset[0,\infty)\times e_j^\perp$ put
\[
\mathfrak{E}'_j(A):=J\,\mathfrak{E}_1\left(\left\{(\omega,p_\perp):(\omega,gp_\perp)\in A\right\}\right)J^{-1};
\]
this is a projection-valued measure on $[0,\infty)\times e_j^\perp$,
the image of $\mathfrak{E}_1$ under the homeomorphism $(\omega,p_\perp)\mapsto(\omega,gp_\perp)$ of $[0,\infty)\times e_1^\perp$ onto $[0,\infty)\times e_j^\perp$, conjugated by $J$.
As $(gp_\perp)\cdot y'=p_\perp\cdot\left(g^{-1}y'\right)$, Lemma~\ref{lem:joint-spectrum}\,(b) for $\mu'$ gives
\[
\begin{aligned}
\int\mathrm{d}\mathfrak{E}'_j(\omega,p_\perp)\,\exp\left(it\cdot\omega+ip_\perp\cdot y'\right)
&=J\exp(it\cdot\mathsf{H}_1)\,V_1\left(g^{-1}y'\right)J^{-1}\\
&=\exp(it\cdot\mathsf{H}_j)\,V_j(y')\qquad\left(t\in\R,\ y'\in e_j^\perp\right),
\end{aligned}
\]
and by Lemma~\ref{lem:joint-spectrum}\,(b) for $\mu$ the right side is also the integral of the same function with respect to $\mathfrak{E}_j$.
So, for $\psi,\psi'\in\mathcal{H}_j(\mu)$, the finite complex measures $\left\langle\psi,\mathfrak{E}'_j(\cdot)\psi'\right\rangle$ and $\left\langle\psi,\mathfrak{E}_j(\cdot)\psi'\right\rangle$
have the same Fourier transform on $\R\times e_j^\perp$;
by the uniqueness of the Fourier transform of a finite complex measure \cite[Theorem~1.3.6]{Rudin1962} they coincide.
Hence $\mathfrak{E}'_j=\mathfrak{E}_j$, which is the asserted formula for $\mathfrak{E}_j$.
For $y=s\cdot e_2$ one has $gy=s\cdot e_i$.
The component of $gp_\perp$ along $e_i$ is $p_\perp\cdot g^{-1}e_i=p_2$;
so the joint spectral measures of $(\mathsf{H}_j,\mathsf{P}_i)$ and of $(\mathsf{H}_1,\mathsf{P}_2)$,
which are the images of $\mathfrak{E}_j$ and of $\mathfrak{E}_1$ under $(\omega,p_\perp)\mapsto(\omega,p_i)$ and under $(\omega,p_\perp)\mapsto(\omega,p_2)$ (Lemma~\ref{lem:joint-spectrum}\,(c)),
correspond under $J$.

Let now $\mu$ satisfy \textup{(G)}.
Then $\mu'$ satisfies \textup{(G)} by \eqref{it:transport-properties},
and the identity $\hat g\,\mathbb{B}(b)=\mathbb{B}(g\cdot b)$ of \eqref{it:transport-operators} holds in $L^2(\mu)$ for every $b\in L^2(\R^3;\C^3)$:
both sides depend continuously on $b$, by Lemma~\ref{lem:domination-consequences}\eqref{it:field-extension}, by the unitarity of $\hat g$ and by $\|g\cdot b\|=\|b\|$,
and they agree on the dense subspace $\mathcal{S}(\R^3;\C^3)$.
Moreover $(g\cdot b)(x)=\det(g)\cdot g\,b\left(g^{-1}x\right)$, and $g^{-1}x$ lies in $\{x_1<0\}$ exactly when $x$ lies in $\{x_j<0\}$;
so $b\mapsto g\cdot b$ maps the fields that vanish almost everywhere on $\{x_1<0\}$ onto the fields that vanish almost everywhere on $\{x_j<0\}$.
Hence $J[\mathbb{B}(b)]_1=[\mathbb{B}(g\cdot b)]_j$ for every $b\in L^2(\R^3;\C^3)$ that vanishes almost everywhere on $\{x_1<0\}$;
for such $b$ one has $\mathbb{B}(b)\in L^2(\Sigma_1^+,\mu')$ and $\mathbb{B}(g\cdot b)\in L^2(\Sigma_j^+,\mu)$
(Definition~\ref{def:closedness-cones}\eqref{it:one-field-space}).
Consequently $J$ maps the vectors that span $\mathcal{H}^1_1(\mu')$ onto the vectors that span $\mathcal{H}^1_j(\mu)$,
and, being unitary, it carries the closed span of the former onto the closed span of the latter.

The preimage of $\{(\omega,p_\perp):\omega<|p_i|\}$ under $(\omega,p_\perp)\mapsto(\omega,gp_\perp)$ is $\{(\omega,p_\perp):\omega<|p_2|\}$, so that
\[
\mathfrak{E}_j\left(\left\{\omega<|p_i|\right\}\right)=J\,\mathfrak{E}_1\left(\left\{\omega<|p_2|\right\}\right)J^{-1}.
\]
This shows that $(C_{12})$ for $\mu'$ is equivalent to $(C_{ji})$ for $\mu$,
and, together with $J\mathcal{H}^1_1(\mu')=\mathcal{H}^1_j(\mu)$ under \textup{(G)}, that $(C^1_{12})$ for $\mu'$ is equivalent to $(C^1_{ji})$ for $\mu$.
The map $g$ takes $e_1$, $e_2$ and $e_3$ to $e_j$, $e_i$ and the third unit vector;
so $gR^{12}_tg^{-1}$ takes $e_j$ to $\cos t\cdot e_j+\sin t\cdot e_i$ and $e_i$ to $-\sin t\cdot e_j+\cos t\cdot e_i$, and it fixes the third unit vector.
Thus $gR^{12}_tg^{-1}=R^{ji}_t$, and \eqref{it:transport-properties} shows that $\left(R^{12}_t\right)_*\mu'=\mu'$ if and only if $\left(R^{ji}_t\right)_*\mu=\mu$;
for $t=\pi/2$ this is the equivalence of $(Q_{12})$ for $\mu'$ and $(Q_{ji})$ for $\mu$.
Together with \eqref{it:transport-properties} for \textup{(G)}, \textup{(dF)} and $(T_{123})$,
and with the equivalences for $(R_1)$ and $(RP_1)$ established above,
this proves the assertion on the hypotheses of Definition~\ref{def:closedness-cones}\eqref{it:one-plane-hypotheses}.
\end{proof}

\subsection{No contact terms}\label{sec:contact-terms}

The proposition of this subsection is a statement of functional analysis about bounded multilinear forms on $L^2$; no law enters it.
With its help, an identity that is known for entries whose supports lie over pairwise disjoint compact subsets of the plane extends to all entries in $L^2$, and hence to the law itself:
it is applied with $N=2$ in the proof of Theorem~\ref{thm:two-point},
and with arbitrary $N$ in the proof of Theorem~\ref{thm:one-plane}.
The argument for Yang--Mills theory in four dimensions has no counterpart of this step.

\begin{proposition}[No contact terms]\label{prop:no-contact-terms}
Let $U$ be an open subset of $\R^2$, and put
\[
L^2_U:=\{f\in L^2(\R^3;\C^3):\ f=0\text{ almost everywhere outside }U\times\R\}.
\]
Let $N\geq1$, and let $D\colon(L^2_U)^N\to\C$ be a bounded $N$-linear form with norm $\|D\|$, so that
\[
|D(f_1,\dots,f_N)|\leq\|D\|\cdot\prod_{i=1}^{N}\|f_i\|\qquad(f_1,\dots,f_N\in L^2_U),
\]
where $\|\cdot\|$ is the norm of $L^2(\R^3;\C^3)$.
Suppose that $D$ has the following two properties.
\begin{enumerate}
\item\label{it:contact-disjoint}
If $K_1,\dots,K_N$ are compact subsets of $U$ with $K_i\cap K_j=\emptyset$ for $i\neq j$,
and $f_1,\dots,f_N\in L^2_U$ satisfy $f_i=0$ almost everywhere outside $K_i\times\R$ for $i=1,\dots,N$,
then $D(f_1,\dots,f_N)=0$.
\item\label{it:contact-gradients}
If $f_1,\dots,f_N\in L^2_U$ and $f_\ell=\nabla u$ for some $\ell\in\{1,\dots,N\}$ and some $u\in C^{\infty}_{\mathrm c}(U\times\R;\C)$,
then $D(f_1,\dots,f_N)=0$.
\end{enumerate}
Then $D=0$.
\end{proposition}

\begin{proof}
The proof is by induction on $N$,
and at each stage the assertion is proved for all open subsets $U$ of $\R^2$ at once.
As in Definition~\ref{def:labels-plates}, a point of $\R^3$ is written $x=(y,x_3)$ with $y=(y^0,y^1)=(x_1,x_2)\in\R^2$.
A bounded Borel function $h$ on $\R^2$ acts on $L^2_U$ by multiplication, $(h\,f)(x):=h(y)\cdot f(x)$;
this operator has norm at most $\sup|h|$.
For a Borel subset $A$ of $\R^2$, $1_A$ stands for the indicator function of $A$ and for the multiplication by it,
which is an orthogonal projection of $L^2_U$.
The fields $f\in L^2_U$ that vanish outside $K\times\R$ for some compact subset $K$ of $U$ form a dense subspace of $L^2_U$,
since $U$ is the union of an increasing sequence of compact sets, along which $1_K\,f$ tends to $f$ in $L^2$ by dominated convergence.

\emph{One entry.}
For $N=1$, property~\eqref{it:contact-disjoint} says that $D$ vanishes on this dense subspace.
As $D$ is continuous, $D=0$.

\emph{Two entries.}
Let $N=2$.
For $v,v'\in\C^3$ let $v\cdot v':=v_1\cdot v'_1+v_2\cdot v'_2+v_3\cdot v'_3$ be the bilinear pairing, without complex conjugation.
The map $(f',f)\mapsto D(\overline{f'},f)$ is a bounded sesquilinear form on $L^2_U$.
By the Riesz representation theorem there is a bounded operator $T$ on $L^2_U$ such that
$D(\overline{f'},f)=\int\mathrm{d}x\,\overline{f'(x)}\cdot(T\,f)(x)$ for all $f,f'\in L^2_U$;
with $f':=\overline{f_1}$ and $f:=f_2$ this reads
\[
D(f_1,f_2)=\int\mathrm{d}x\,f_1(x)\cdot(T\,f_2)(x)\qquad(f_1,f_2\in L^2_U).
\]
Let $K$ and $K'$ be disjoint compact subsets of $U$.
For all $f_1,f_2\in L^2_U$ the fields $1_K\,f_1$ and $1_{K'}\,f_2$ vanish outside $K\times\R$ and $K'\times\R$ respectively,
so that property~\eqref{it:contact-disjoint} gives
\[
\int\mathrm{d}x\,f_1(x)\cdot(1_K\,T\,1_{K'}\,f_2)(x)=D(1_K\,f_1,1_{K'}\,f_2)=0;
\]
as $f_1$ is arbitrary, $1_K\,T\,1_{K'}=0$.
For Borel subsets $A$ and $A'$ of $U$ with $A\cap A'=\emptyset$ one has $1_A\,T\,1_{A'}=0$ as well.
The Lebesgue measure of $\R^2$ is inner regular on each of the bounded Borel sets $A\cap\{|y|<n\}$ and $A'\cap\{|y|<n\}$, $n\in\N$;
taking finite unions of compact subsets of these sets, one obtains increasing sequences of compact sets $K_n\subset A$ and $K'_n\subset A'$
for which $A\setminus\bigcup_nK_n$ and $A'\setminus\bigcup_nK'_n$ are null sets.
Dominated convergence gives $1_{K_n}\to1_A$ and $1_{K'_n}\to1_{A'}$ strongly,
and the operators $1_{K_n}\,T\,1_{K'_n}$, all of which vanish, converge strongly to $1_A\,T\,1_{A'}$.
For a Borel subset $A$ of $U$ this applies to the pairs $(U\setminus A,A)$ and $(A,U\setminus A)$:
\[
T\,1_A=1_A\,T\,1_A+1_{U\setminus A}\,T\,1_A=1_A\,T\,1_A,
\qquad
1_A\,T=1_A\,T\,1_A+1_A\,T\,1_{U\setminus A}=1_A\,T\,1_A.
\]
As $1_A=1_{A\cap U}$ on $L^2_U$ for every Borel subset $A$ of $\R^2$, it follows that
\begin{equation}\label{eq:rotation-commutation}
1_A\,T=T\,1_A\qquad\text{for every Borel subset $A$ of $\R^2$.}
\end{equation}
Every bounded Borel function on $\R^2$ is a uniform limit of simple Borel functions;
together with \eqref{eq:rotation-commutation}, this gives $h\,T=T\,h$ for every bounded Borel function $h$ on $\R^2$.
Consequently
\[
D(h\,f_1,f_2)=\int\mathrm{d}x\,f_1(x)\cdot(h\,T\,f_2)(x)=D(f_1,h\,f_2)\qquad(f_1,f_2\in L^2_U),
\]
and property~\eqref{it:contact-gradients} yields
\begin{equation}\label{eq:rotation-modulated-gradients}
D(h\,\nabla u,f_2)=D(\nabla u,h\,f_2)=0
\end{equation}
for every $u\in C^{\infty}_{\mathrm c}(U\times\R;\C)$, every bounded Borel function $h$ on $\R^2$ and every $f_2\in L^2_U$.

Let $A$ be a compact subset of $U$,
let $u_0\in C^{\infty}_{\mathrm c}(U)$ be equal to $1$ on a neighbourhood of $A$,
and let $\chi_0\in C^{\infty}_{\mathrm c}(\R)$.
A product such as $h\cdot1_A\cdot\chi_0'\cdot e_3$ denotes the vector field $x\mapsto h(y)\cdot1_A(y)\cdot\chi_0'(x_3)\cdot e_3$.
The three functions
\[
x\mapsto u_0(y)\cdot\chi_0(x_3),\qquad
x\mapsto u_0(y)\cdot y^0\cdot\chi_0(x_3),\qquad
x\mapsto u_0(y)\cdot y^1\cdot\chi_0(x_3)
\]
belong to $C^{\infty}_{\mathrm c}(U\times\R)$.
On $A\times\R$ the derivatives of $u_0$ vanish, and there the gradients of the three functions are
\[
\chi_0'\cdot e_3,\qquad\chi_0\cdot e_1+y^0\cdot\chi_0'\cdot e_3,\qquad\chi_0\cdot e_2+y^1\cdot\chi_0'\cdot e_3 .
\]
Apply \eqref{eq:rotation-modulated-gradients} with $h\cdot1_A$ in place of $h$.
The first function gives
\[
D(h\cdot1_A\cdot\chi_0'\cdot e_3,f_2)=0\qquad\text{for every bounded Borel function $h$ on $\R^2$.}
\]
The second function gives
\[
D(h\cdot1_A\cdot\chi_0\cdot e_1,f_2)=-D(h\cdot y^0\cdot1_A\cdot\chi_0'\cdot e_3,f_2)=0,
\]
where the last equality is the preceding identity for the function $y\mapsto h(y)\cdot y^0\cdot1_A(y)$,
which is bounded because $A$ is compact.
The third function gives
\[
D(h\cdot1_A\cdot\chi_0\cdot e_2,f_2)=-D(h\cdot y^1\cdot1_A\cdot\chi_0'\cdot e_3,f_2)=0,
\]
the last equality being the identity for the first function, now for the bounded function $y\mapsto h(y)\cdot y^1\cdot1_A(y)$.
As $A$ ranges over the compact subsets of $U$, $h$ over the bounded Borel functions on $\R^2$ and $\chi_0$ over $C^{\infty}_{\mathrm c}(\R)$,
the linear combinations of the functions $(y,x_3)\mapsto h(y)\cdot1_A(y)\cdot\chi_0(x_3)$ are dense in $L^2(U\times\R;\C)$,
and the same is true with $\chi_0'$ in place of $\chi_0$.
For the second family it suffices to know that the orthogonal complement of $\{\chi_0':\ \chi_0\in C^{\infty}_{\mathrm c}(\R)\}$ in $L^2(\R)$ is $\{0\}$.
A function in this complement has derivative zero in the sense of distributions;
it therefore agrees almost everywhere with a constant, and no constant other than $0$ is square integrable on $\R$.
Therefore the fields $h\cdot1_A\cdot\chi_0\cdot e_1$, $h\cdot1_A\cdot\chi_0\cdot e_2$ and $h\cdot1_A\cdot\chi_0'\cdot e_3$ form a total subset of $L^2_U$.
The continuous linear form $D(\cdot,f_2)$ vanishes on this subset, hence on $L^2_U$, for every $f_2\in L^2_U$;
that is, $D=0$.

\emph{Three or more entries.}
Let $N\geq3$, and assume that the proposition holds for every number of entries smaller than $N$ and every open subset of $\R^2$.
The induction hypothesis yields two claims.
\begin{enumerate}
\item\label{it:contact-one-entry}
Let $K$ be a compact subset of $U$ and $i\in\{1,\dots,N\}$.
If $f_i$ vanishes outside $K\times\R$ and every other entry vanishes outside $(U\setminus K)\times\R$, then $D(f_1,\dots,f_N)=0$.
\item\label{it:contact-two-clusters}
Let $I$ and $J$ be non-empty subsets of $\{1,\dots,N\}$ with $I\cap J=\emptyset$ and $I\cup J=\{1,\dots,N\}$, and let $K_J\subset U$ be compact.
If $f_\ell$ vanishes outside $K_J\times\R$ for every $\ell\in J$,
and $f_i$ vanishes outside $(U\setminus K_J)\times\R$ for every $i\in I$, then $D(f_1,\dots,f_N)=0$.
\end{enumerate}
To prove claim~\eqref{it:contact-one-entry}, keep $f_i$ fixed.
The set $U\setminus K$ is open, the other $N-1$ entries range over $L^2_{U\setminus K}$,
and $D$ becomes a bounded $(N-1)$-linear form on $L^2_{U\setminus K}$, of norm at most $\|D\|\cdot\|f_i\|$.
This form has property~\eqref{it:contact-disjoint} with $U\setminus K$ in place of $U$:
if each entry $f_\ell$, $\ell\neq i$, vanishes outside $K_\ell\times\R$, where the $K_\ell$ are pairwise disjoint compact subsets of $U\setminus K$,
then the sets $K_\ell$, $\ell\neq i$, together with $K$ are $N$ pairwise disjoint compact subsets of $U$,
and $f_i$ vanishes outside $K\times\R$, so that $D(f_1,\dots,f_N)=0$ by property~\eqref{it:contact-disjoint} of $D$.
Every function in $C^{\infty}_{\mathrm c}((U\setminus K)\times\R;\C)$, extended by $0$, lies in $C^{\infty}_{\mathrm c}(U\times\R;\C)$;
so the form inherits property~\eqref{it:contact-gradients} from $D$, again with $U\setminus K$ in place of $U$.
As $N-1<N$, the induction hypothesis applies to this form on the open set $U\setminus K$ and shows that it is identically zero.

For claim~\eqref{it:contact-two-clusters}, hold the entries $f_\ell$, $\ell\in J$, fixed;
the entries $f_i$, $i\in I$, then range over $L^2_{U\setminus K_J}$,
and $D$ becomes a bounded form in $\#|I|<N$ entries on this space.
It has property~\eqref{it:contact-gradients} with $U\setminus K_J$ in place of $U$,
as a function in $C^{\infty}_{\mathrm c}((U\setminus K_J)\times\R;\C)$ becomes an element of $C^{\infty}_{\mathrm c}(U\times\R;\C)$ when it is extended by $0$.
It also has property~\eqref{it:contact-disjoint} with $U\setminus K_J$ in place of $U$.
Indeed, suppose that each $f_i$, $i\in I$, vanishes outside $K_i\times\R$,
where the $K_i$ are pairwise disjoint compact subsets of $U\setminus K_J$, and fix $i'\in I$.
Every entry other than $f_{i'}$ vanishes outside $(U\setminus K_{i'})\times\R$:
for the indices in $I\setminus\{i'\}$ because the sets $K_i$ are pairwise disjoint,
and for the indices in $J$ because $K_J$ does not meet $K_{i'}$.
Claim~\eqref{it:contact-one-entry}, applied with the index $i'$ and the compact set $K_{i'}$, gives $D(f_1,\dots,f_N)=0$.
The induction hypothesis, applied with $\#|I|$ entries and the open set $U\setminus K_J$, shows that the form is zero.

It remains to show that $D(f_1,\dots,f_N)=0$ for arbitrary entries.
Since $D$ is continuous, since the fields that vanish outside $K\times\R$ for a compact subset $K$ of $U$ are dense in $L^2_U$,
and since a finite union of compact sets is compact,
one may assume that all entries vanish outside $K_0\times\R$ for one compact subset $K_0$ of $U$.
For $d>0$ let
\[
\mathcal{Q}_d:=\{d(k+[0,1)^2):\ k\in\Z^2\}
\]
be the partition of $\R^2$ into half-open squares of side $d$.
Only finitely many squares of $\mathcal{Q}_d$ meet $K_0$;
so $f_i=\sum_{Q\in\mathcal{Q}_d}1_Q\,f_i$ is a finite sum for every $i$, and by $N$-linearity
\[
D(f_1,\dots,f_N)=\sum_{Q_1,\dots,Q_N\in\mathcal{Q}_d}D(1_{Q_1}\,f_1,\dots,1_{Q_N}\,f_N).
\]
Fix squares $Q_1,\dots,Q_N\in\mathcal{Q}_d$ that do not all coincide,
and put $I:=\{i:Q_i=Q_1\}$ and $J:=\{1,\dots,N\}\setminus I$; then $J\neq\emptyset$.
For $0<\varepsilon<d/2$ and $Q\in\mathcal{Q}_d$ denote by $Q^\varepsilon$ the closed square of side $d-2\varepsilon$ concentric with $Q$.
Then $Q^\varepsilon\subset Q$, and $1_{Q^\varepsilon}\,f\to1_Q\,f$ in $L^2$ as $\varepsilon\to0$ for every $f\in L^2_U$,
because the union of the squares $Q^\varepsilon$ differs from $Q$ by a null set.
As distinct squares of $\mathcal{Q}_d$ are disjoint, the compact set
\[
K_J:=K_0\cap\bigcup_{\ell\in J}Q_\ell^\varepsilon
\]
does not meet the compact set $K_0\cap Q_1^\varepsilon$.
For $\ell\in J$ the field $1_{Q_\ell^\varepsilon}\,f_\ell$ vanishes outside $K_J\times\R$,
and for $i\in I$ the field $1_{Q_i^\varepsilon}\,f_i=1_{Q_1^\varepsilon}\,f_i$ vanishes outside $(K_0\cap Q_1^\varepsilon)\times\R$,
which is contained in $(U\setminus K_J)\times\R$.
Claim~\eqref{it:contact-two-clusters} therefore gives $D(1_{Q_1^\varepsilon}\,f_1,\dots,1_{Q_N^\varepsilon}\,f_N)=0$,
and as $\varepsilon\to0$ the continuity of $D$ shows that $D(1_{Q_1}\,f_1,\dots,1_{Q_N}\,f_N)=0$.
Only the terms with $Q_1=\dots=Q_N$ survive, and
\[
|D(f_1,\dots,f_N)|=\left|\sum_{Q\in\mathcal{Q}_d}D(1_Q\,f_1,\dots,1_Q\,f_N)\right|
\leq\|D\|\cdot\sum_{Q\in\mathcal{Q}_d}\prod_{i=1}^{N}\|1_Q\,f_i\|.
\]
Since $\sum_{Q\in\mathcal{Q}_d}\|1_Q\,f\|^2=\|f\|^2$ for $f\in L^2_U$, the Cauchy--Schwarz inequality gives
\[
\sum_{Q\in\mathcal{Q}_d}\|1_Q\,f_1\|\cdot\|1_Q\,f_2\|
\leq\left(\sum_{Q\in\mathcal{Q}_d}\|1_Q\,f_1\|^2\right)^{\frac{1}{2}}\cdot\left(\sum_{Q\in\mathcal{Q}_d}\|1_Q\,f_2\|^2\right)^{\frac{1}{2}}
=\|f_1\|\cdot\|f_2\|,
\]
and together with $\|1_Q\,f_i\|\leq\|f_i\|$ for $i\geq4$ this yields
\[
|D(f_1,\dots,f_N)|\leq\|D\|\cdot\max_{Q\in\mathcal{Q}_d}\|1_Q\,f_3\|\cdot\|f_1\|\cdot\|f_2\|\cdot\prod_{i=4}^{N}\|f_i\|.
\]
Here
\[
\|1_Q\,f_3\|^2=\int_Q\mathrm{d}y\int_{\R}\mathrm{d}x_3\,|f_3(y,x_3)|^2
\]
is the integral, over a set of measure $d^2$, of one function of $y$ that does not depend on $d$ and is integrable on $\R^2$ by Tonelli's theorem.
By the absolute continuity of the integral, $\max_{Q\in\mathcal{Q}_d}\|1_Q\,f_3\|$ tends to $0$ as $d\to0$.
The left-hand side of the last estimate does not depend on $d$, and so $D(f_1,\dots,f_N)=0$.
\end{proof}

\subsection{The insertion bound}\label{sec:insertion-bound}

\begin{definition}[Time slabs]\label{def:time-slabs}
For real numbers $c<c'$ let
\[
L^2_{[c,c']}:=\left\{b\in L^2(\R^3;\C^3):\ b=0\text{ almost everywhere outside the time slab }\{c\le x_1\le c'\}\right\}.
\]
Let $\mu$ be a law on $\mathcal{S}'$ that satisfies \eqref{cond:gaussian-domination}, $(T_1)$, $(R_1)$ and $(RP_1)$,
let $\Sigma_1^{+}:=\Sigma_{\{x_1>0\}}$,
and let
\[
L^{\infty-}_{+}:=\bigcap_{1\le p<\infty}L^p\left(\mathcal{S}',\Sigma_1^{+},\mu\right).
\]
\end{definition}

By H\"older's inequality, $L^{\infty-}_{+}$ is an algebra under pointwise multiplication.
It is a subspace of $L^2(\Sigma_1^{+},\mu)$, and it contains every bounded $\Sigma_1^{+}$-measurable function.
Let $b\in L^2(\R^3;\C^3)$ vanish almost everywhere on $\{x_1<0\}$.
By Lemma~\ref{lem:domination-consequences}(a) the field $\mathbb{B}(b)$ belongs to $L^2(\Sigma_1^{+},\mu)$ and to every $L^p(\mu)$, $p<\infty$;
moreover $|\exp(\mathbb{B}(b))|=\exp(\mathbb{B}(\Re b))$, and the same lemma, applied to the real field $p\cdot\Re b$, gives
\[
\mu\left(\exp\left(p\cdot\mathbb{B}(\Re b)\right)\right)\le\exp\left(\frac{1}{2}\cdot p^2\cdot\|\Re b\|^2\right)\qquad(1\le p<\infty),
\]
where $\|\cdot\|$ is the norm of $L^2(\R^3;\C^3)$.
Hence $L^{\infty-}_{+}$ contains the polynomials and the exponentials in such fields.
For $l\ge0$ the operator $U_l$ of Definition~\ref{def:labels-plates} maps $L^{\infty-}_{+}$ into itself:
it maps $\Sigma_1^{+}$-measurable functions to $\Sigma_{\{x_1>l\}}$-measurable ones, $\Sigma_{\{x_1>l\}}$ is contained in $\Sigma_1^{+}$,
and by $(T_1)$ it preserves the norm of every $L^p(\mu)$.

\begin{lemma}[Source operators]\label{lem:source-operators}
Let $\mu$ be a law on $\mathcal{S}'$ that satisfies \eqref{cond:gaussian-domination}, $(T_1)$, $(R_1)$ and $(RP_1)$.
Let $l>0$, and let $\Phi$ be a bounded $\Sigma_1^{+}$-measurable function on $\mathcal{S}'$.
For every $h\in L^2_{[0,l]}$ the function $\exp(\mathbb{B}(h))\cdot U_l\Phi$ belongs to $L^{\infty-}_{+}$, and
\[
\left\|\left[\exp(\mathbb{B}(h))\cdot U_l\Phi\right]_1\right\|\le\exp\left(\frac{1}{2}\cdot\|\Re h\|^2\right)\cdot\|[\Phi]_1\| .
\]
\end{lemma}

\begin{proof}
For $l>0$ and $h\in L^2_{[0,l]}$ the source operator $M_l(h)$ assigns to a bounded $\Sigma_1^{+}$-measurable function $\Phi$ the function
\[
M_l(h)\Phi:=\exp(\mathbb{B}(h))\cdot U_l\Phi .
\]
The field $h$ vanishes almost everywhere on $\{x_1<0\}$;
so $\mathbb{B}(h)$ is $\Sigma_1^{+}$-measurable, and $\exp(\mathbb{B}(h))$, whose modulus $\exp(\mathbb{B}(\Re h))$ lies in every $L^p(\mu)$, belongs to $L^{\infty-}_{+}$
(Lemma~\ref{lem:domination-consequences}(a)).
The bounded $\Sigma_1^{+}$-measurable function $U_l\Phi$ belongs to $L^{\infty-}_{+}$ as well,
and their product $M_l(h)\Phi$ lies in the algebra $L^{\infty-}_{+}$, which is the first assertion.
From now on $\Phi$ is fixed.
Recall that $\left\langle[\Phi']_1,[\Phi'']_1\right\rangle=\mu\left((\Theta_1\Phi')\cdot\Phi''\right)$ for $\Phi',\Phi''\in L^2(\Sigma_1^{+},\mu)$ (Definition~\ref{def:laws}).

\emph{Step 1: doubling.}
The doubled source of $h$ for the width $l$ is
\[
h^{\#}:=\tau_l\cdot\left(\overline{\theta_1\cdot h}+h\right).
\]
The field $\theta_1\cdot h$ vanishes almost everywhere outside $\{-l\le x_1\le0\}$.
Consequently $\tau_l\cdot\overline{\theta_1\cdot h}\in L^2_{[0,l]}$ and $\tau_l\cdot h\in L^2_{[l,2l]}$, so that $h^{\#}\in L^2_{[0,2l]}$.
The map $\Theta_1$ is multiplicative, and by \eqref{eq:rotation-conjugation}
$\Theta_1\exp(\mathbb{B}(h))=\exp\left(\mathbb{B}\left(\overline{\theta_1\cdot h}\right)\right)$ and $\Theta_1U_l=U_{-l}\Theta_1$.
The operator $U_l$ is multiplicative as well, $U_l\mathbb{B}(b)=\mathbb{B}(\tau_l\cdot b)$ for every $b\in L^2(\R^3;\C^3)$ (Lemma~\ref{lem:domination-consequences}(a)),
and $\mu(U_l\Phi')=\mu(\Phi')$ for every $\mu$-integrable function $\Phi'$ by $(T_1)$.
Applying $U_l$ to the integrand in the second of the following equalities, and using the linearity of $b\mapsto\mathbb{B}(b)$, one obtains
\begin{equation}\label{eq:rotation-doubling}
\begin{aligned}
\left\|\left[M_l(h)\Phi\right]_1\right\|^2
&=\mu\left(\exp\left(\mathbb{B}\left(\overline{\theta_1\cdot h}\right)\right)\cdot\left(U_{-l}\Theta_1\Phi\right)\cdot\exp(\mathbb{B}(h))\cdot U_l\Phi\right)\\
&=\mu\left(\left(\Theta_1\Phi\right)\cdot\exp\left(\mathbb{B}\left(h^{\#}\right)\right)\cdot U_{2l}\Phi\right)
=\left\langle[\Phi]_1,\left[M_{2l}\left(h^{\#}\right)\Phi\right]_1\right\rangle .
\end{aligned}
\end{equation}

\emph{Step 2: symmetric sources.}
Let $l'>0$.
An element $h'$ of $L^2_{[0,l']}$ is called \emph{symmetric} if
\[
\tau_{l'}\cdot\overline{\theta_1\cdot h'}=h' ,
\]
that is, if $h'$ is invariant under complex conjugation combined with the reflection in the middle plane $\{x_1=l'/2\}$ of the slab $\{0\le x_1\le l'\}$.
For real $s$ the reflection $\theta_1$ conjugates the translation by $s\,e_1$ into the translation by $-s\,e_1$, so that $\theta_1\cdot(\tau_s\cdot b)=\tau_{-s}\cdot(\theta_1\cdot b)$;
complex conjugation commutes with the actions of $\theta_1$ and of $\tau_s$;
and $b\mapsto\overline{\theta_1\cdot b}$ is an involution of $L^2(\R^3;\C^3)$.
Hence
\[
\tau_{2l}\cdot\overline{\theta_1\cdot h^{\#}}
=\tau_{2l}\cdot\tau_{-l}\cdot\left(\overline{\theta_1\cdot\overline{\theta_1\cdot h}}+\overline{\theta_1\cdot h}\right)
=\tau_l\cdot\left(h+\overline{\theta_1\cdot h}\right)
=h^{\#}:
\]
the doubled source $h^{\#}$ of Step~1 is symmetric, with $l'=2l$.

Let now $h'\in L^2_{[0,l']}$ be symmetric, and for integers $n\ge1$ put
\[
h'^{(n)}:=\sum_{k=0}^{n-1}\tau_{kl'}\cdot h'\in L^2_{[0,nl']} .
\]
Then
\[
\tau_{nl'}\cdot\overline{\theta_1\cdot h'^{(n)}}
=\sum_{k=0}^{n-1}\tau_{(n-1-k)l'}\cdot\left(\tau_{l'}\cdot\overline{\theta_1\cdot h'}\right)
=\sum_{k=0}^{n-1}\tau_{(n-1-k)l'}\cdot h'
=h'^{(n)},
\]
so that $h'^{(n)}$ is symmetric in the slab $\{0\le x_1\le nl'\}$.
Its doubled source for the width $nl'$ is
\[
\left(h'^{(n)}\right)^{\#}
=\tau_{nl'}\cdot\overline{\theta_1\cdot h'^{(n)}}+\tau_{nl'}\cdot h'^{(n)}
=h'^{(n)}+\sum_{k=n}^{2n-1}\tau_{kl'}\cdot h'
=h'^{(2n)} .
\]

\emph{Step 3: iteration.}
Let $h'\in L^2_{[0,l']}$ be symmetric, and for $a=0,1,2,\dots$ let
\[
m_a:=\left\|\left[M_{2^al'}\left(h'^{(2^a)}\right)\Phi\right]_1\right\| .
\]
By Step~2 the doubled source of $h'^{(2^a)}$ for the width $2^al'$ is $h'^{(2^{a+1})}$.
Therefore \eqref{eq:rotation-doubling}, applied to this width and this source, and the Cauchy--Schwarz inequality in $\mathcal{H}_1$ give
\[
m_a^2=\left\langle[\Phi]_1,\left[M_{2^{a+1}l'}\left(h'^{(2^{a+1})}\right)\Phi\right]_1\right\rangle\le\|[\Phi]_1\|\cdot m_{a+1},
\]
and induction on $a$ yields
\[
m_0\le\|[\Phi]_1\|^{1-2^{-a}}\cdot m_a^{2^{-a}}\qquad(a=1,2,3,\dots).
\]
The third member of \eqref{eq:rotation-doubling}, for the same width and source, expresses $m_a^2$ as the integral of $(\Theta_1\Phi)\cdot\exp\left(\mathbb{B}\left(h'^{(2^{a+1})}\right)\right)\cdot U_{2^{a+1}l'}\Phi$,
and so bounds $m_a$ in a second way.
One has $|\Theta_1\Phi|\le\sup|\Phi|$, $|U_{2^{a+1}l'}\Phi|\le\sup|\Phi|$ and $|\exp(\mathbb{B}(b))|=\exp(\mathbb{B}(\Re b))$.
The real part of $h'^{(2^{a+1})}$ is the sum of the $2^{a+1}$ fields $\tau_{kl'}\cdot\Re h'$, $0\le k<2^{a+1}$,
which vanish outside the slabs $\{kl'\le x_1\le(k+1)l'\}$ with pairwise disjoint interiors;
so its squared norm is $2^{a+1}\cdot\|\Re h'\|^2$.
Gaussian domination on $L^2$, that is, $\mu(\exp(\mathbb{B}(b)))\le\exp\left(\frac{1}{2}\cdot\|b\|^2\right)$ for real $b\in L^2(\R^3;\R^3)$ (Lemma~\ref{lem:domination-consequences}(a)), now gives
\[
\begin{aligned}
m_a^2&\le(\sup|\Phi|)^2\cdot\mu\left(\exp\left(\mathbb{B}\left(\Re h'^{(2^{a+1})}\right)\right)\right)
\le(\sup|\Phi|)^2\cdot\exp\left(\frac{1}{2}\cdot\left\|\Re h'^{(2^{a+1})}\right\|^2\right)\\
&=(\sup|\Phi|)^2\cdot\exp\left(2^a\cdot\|\Re h'\|^2\right).
\end{aligned}
\]
The power $2^{-a-1}$ of this inequality bounds $m_a^{2^{-a}}$, and inserted into the preceding estimate it gives
\[
m_0\le\|[\Phi]_1\|^{1-2^{-a}}\cdot(\sup|\Phi|)^{2^{-a}}\cdot\exp\left(\frac{1}{2}\cdot\|\Re h'\|^2\right)\qquad(a=1,2,3,\dots).
\]
If $\Phi\neq0$, then $(\sup|\Phi|)^{2^{-a}}\to1$ and $\|[\Phi]_1\|^{1-2^{-a}}\to\|[\Phi]_1\|$ as $a\to\infty$; if $\Phi=0$, then $m_0=0$.
As $h'^{(1)}=h'$, in either case
\[
\left\|\left[M_{l'}(h')\Phi\right]_1\right\|\le\exp\left(\frac{1}{2}\cdot\|\Re h'\|^2\right)\cdot\|[\Phi]_1\|\qquad\text{for every symmetric }h'\in L^2_{[0,l']} .
\]

\emph{Step 4: conclusion.}
By Step~2 the doubled source $h^{\#}\in L^2_{[0,2l]}$ of Step~1 is symmetric with $l'=2l$, and its real part
\[
\Re h^{\#}=\tau_l\cdot\left(\theta_1\cdot\Re h+\Re h\right)
\]
is the sum of two fields of norm $\|\Re h\|$ that vanish outside $\{0\le x_1\le l\}$ and outside $\{l\le x_1\le2l\}$, respectively.
Hence $\|\Re h^{\#}\|^2=2\cdot\|\Re h\|^2$, and Step~3 gives $\left\|\left[M_{2l}\left(h^{\#}\right)\Phi\right]_1\right\|\le\exp\left(\|\Re h\|^2\right)\cdot\|[\Phi]_1\|$.
Together with \eqref{eq:rotation-doubling} and the Cauchy--Schwarz inequality this yields
\[
\left\|\left[M_l(h)\Phi\right]_1\right\|^2\le\|[\Phi]_1\|\cdot\left\|\left[M_{2l}\left(h^{\#}\right)\Phi\right]_1\right\|\le\exp\left(\|\Re h\|^2\right)\cdot\|[\Phi]_1\|^2,
\]
which is the square of the asserted inequality.
\end{proof}

\begin{proposition}[The insertion bound]\label{prop:insertion-bound}
Let $\mu$ be a law on $\mathcal{S}'$ that satisfies \eqref{cond:gaussian-domination}, $(T_1)$, $(R_1)$ and $(RP_1)$.
Let $l>0$.
To each $f\in L^2_{[0,l]}$ there corresponds a unique bounded operator $\mathcal{B}_l(f)$ on $\mathcal{H}_1$ such that
\[
\mathcal{B}_l(f)[\Phi]_1=\left[\mathbb{B}(f)\cdot U_l\Phi\right]_1\qquad\left(\Phi\in L^{\infty-}_{+}\right).
\]
It depends complex linearly on $f$, and
\[
\left\|\mathcal{B}_l(f)\right\|\le k_1\cdot\|f\|,\qquad k_1:=\left(2\exp(1)\right)^{\frac{1}{2}};
\]
if $f$ is real, then even $\left\|\mathcal{B}_l(f)\right\|\le\exp\left(\frac{1}{2}\right)\cdot\|f\|$.
Moreover the following hold.
\begin{enumerate}
\item\label{it:insertion-localisation}
\emph{Localisation.}
If $0\le c<c'\le l$ and $f\in L^2_{[c,c']}$, then
\[
\mathcal{B}_l(f)=\exp\left(-c\cdot\mathsf{H}_1\right)\,\mathcal{B}_{c'-c}\left(\tau_{-c}\cdot f\right)\,\exp\left(-(l-c')\cdot\mathsf{H}_1\right).
\]
\item\label{it:insertion-symmetric}
\emph{Symmetric form.}
For $d>0$ and $f\in L^2_{[-d,d]}$ put $\mathcal{B}^{\mathrm{c}}_d(f):=\mathcal{B}_{2d}(\tau_d\cdot f)$.
Then
\[
\left\langle[\Phi']_1,\mathcal{B}^{\mathrm{c}}_d(f)[\Phi]_1\right\rangle=\mu\left(\left(\Theta_1U_d\Phi'\right)\cdot\mathbb{B}(f)\cdot U_d\Phi\right)\qquad\left(\Phi,\Phi'\in L^{\infty-}_{+}\right),
\]
and
\[
\mathcal{B}^{\mathrm{c}}_d(f)^{*}=\mathcal{B}^{\mathrm{c}}_d\left(\overline{\theta_1\cdot f}\right),\qquad\left\|\mathcal{B}^{\mathrm{c}}_d(f)\right\|\le k_1\cdot\|f\| .
\]
\item\label{it:insertion-covariance}
\emph{Covariance.}
Let $\mu$ satisfy moreover $(T_{123})$, let $d>0$ and $f\in L^2_{[-d,d]}$, and let $r$, $t$ and $a$ be real numbers with $r\ge|t|$.
Then
\[
\exp\left(-(r+t)\cdot\mathsf{H}_1+ia\cdot\mathsf{P}_2\right)\,\mathcal{B}^{\mathrm{c}}_d(f)\,\exp\left(-(r-t)\cdot\mathsf{H}_1-ia\cdot\mathsf{P}_2\right)
=\mathcal{B}^{\mathrm{c}}_{d+r}\left(\tau_{t e_1+a e_2}\cdot f\right),
\]
where, for $s\ge0$, $\exp\left(-s\cdot\mathsf{H}_1\pm ia\cdot\mathsf{P}_2\right)$ denotes the product of the commuting operators $\exp\left(-s\cdot\mathsf{H}_1\right)$ and $V_1(\pm a\cdot e_2)=\exp\left(\pm ia\cdot\mathsf{P}_2\right)$.
\end{enumerate}
In particular, for every $d>0$, every $f\in L^2_{[-d,d]}$ and all $\Phi_1,\Phi_2\in L^{\infty-}_{+}$ that are measurable with respect to $\Sigma_{\{x_1>d\}}$,
\[
\left|\mu\left(\left(\Theta_1\Phi_1\right)\cdot\mathbb{B}(f)\cdot\Phi_2\right)\right|\le k_1\cdot\|f\|\cdot\left\|\left[U_{-d}\Phi_1\right]_1\right\|\cdot\left\|\left[U_{-d}\Phi_2\right]_1\right\| .
\]
The same number $k_1$ serves for every $d$ and for all such $\Phi_1$ and $\Phi_2$;
when $\Phi_1$ and $\Phi_2$ are products of fields, it is thus independent of how many factors they have.
\end{proposition}

\begin{proof}
\emph{The estimate on bounded functions.}
Let $\Phi$ be a bounded $\Sigma_1^{+}$-measurable function, and let first $f\in L^2_{[0,l]}$ be real with $f\neq0$.
As $f$ vanishes almost everywhere on $\{x_1<0\}$, the powers $\mathbb{B}(f)^n$ belong to $L^{\infty-}_{+}$, and so do the products $\mathbb{B}(f)^n\cdot U_l\Phi$, $n=0,1,2,\dots$.
Let $z_0\in\C$.
For $|z|\le|z_0|$ the moduli of the terms of the series
\[
\exp\left(z\cdot\mathbb{B}(f)\right)\cdot U_l\Phi=\sum_{n=0}^{\infty}\frac{z^n}{n!}\cdot\mathbb{B}(f)^n\cdot U_l\Phi
\]
are at most the corresponding terms of the series of $\sup|\Phi|\cdot\exp\left(|z_0|\cdot|\mathbb{B}(f)|\right)$.
This function belongs to $L^2(\mu)$:
its square is at most $(\sup|\Phi|)^2\cdot\left(\exp\left(2|z_0|\cdot\mathbb{B}(f)\right)+\exp\left(-2|z_0|\cdot\mathbb{B}(f)\right)\right)$,
which is integrable by Lemma~\ref{lem:domination-consequences}(a).
By dominated convergence the series converges in $L^2(\mu)$, uniformly on the disc $\{|z|\le|z_0|\}$.
The canonical map $[\,\cdot\,]_1$ is a linear contraction of $L^2(\Sigma_1^{+},\mu)$ into $\mathcal{H}_1$ with dense range \cite[Proposition~8.19(i)]{NoncompactQED}
(the companion paper writes $\mathcal{E}_+$ for $L^2(\Sigma_1^{+},\mu)$, $\Theta_1$ for $\hat\theta_1$ and $T(s)$ for $\exp(-s\cdot\mathsf{H}_1)$).
Hence
\[
z\mapsto\left[\exp\left(z\cdot\mathbb{B}(f)\right)\cdot U_l\Phi\right]_1=\sum_{n=0}^{\infty}\frac{z^n}{n!}\cdot\left[\mathbb{B}(f)^n\cdot U_l\Phi\right]_1
\]
is an entire function with values in $\mathcal{H}_1$, and its coefficient of $z$ is $\left[\mathbb{B}(f)\cdot U_l\Phi\right]_1$.
As $\mathbb{B}$ is complex linear, $\exp(z\cdot\mathbb{B}(f))=\exp(\mathbb{B}(z\cdot f))$, where $z\cdot f\in L^2_{[0,l]}$ has the real part $(\Re z)\cdot f$.
By Lemma~\ref{lem:source-operators} the value of the function at $z$ has a norm at most $\exp\left(\frac{1}{2}\cdot(\Re z)^2\cdot\|f\|^2\right)\cdot\|[\Phi]_1\|$,
and on the circle $|z|=\|f\|^{-1}$ this is at most $\exp\left(\frac{1}{2}\right)\cdot\|[\Phi]_1\|$.
Cauchy's estimate for the coefficient of $z$ of an entire function with values in $\mathcal{H}_1$, taken on this circle, gives
\[
\left\|\left[\mathbb{B}(f)\cdot U_l\Phi\right]_1\right\|\le\exp\left(\frac{1}{2}\right)\cdot\|f\|\cdot\|[\Phi]_1\| ,
\]
and this holds trivially for $f=0$.
Now let $f\in L^2_{[0,l]}$ be arbitrary.
Its real and imaginary parts $\Re f$ and $\Im f$ lie in $L^2_{[0,l]}$, $\mathbb{B}(f)=\mathbb{B}(\Re f)+i\,\mathbb{B}(\Im f)$,
and $\|\Re f\|+\|\Im f\|\le2^{\frac{1}{2}}\cdot\|f\|$.
The case of real fields therefore gives
\[
\left\|\left[\mathbb{B}(f)\cdot U_l\Phi\right]_1\right\|\le\exp\left(\frac{1}{2}\right)\cdot\left(\|\Re f\|+\|\Im f\|\right)\cdot\|[\Phi]_1\|\le k_1\cdot\|f\|\cdot\|[\Phi]_1\| .
\]

\emph{The operator.}
The bounded $\Sigma_1^{+}$-measurable functions are dense in $L^2(\Sigma_1^{+},\mu)$, and $[\,\cdot\,]_1$ is a contraction with dense range;
so the classes of these functions are dense in $\mathcal{H}_1$.
Let $\Phi$ and $\Phi'$ be bounded $\Sigma_1^{+}$-measurable functions with $[\Phi]_1=[\Phi']_1$.
Since $[\,\cdot\,]_1$ is linear, the estimate just proved, applied to $\Phi-\Phi'$, gives
\[
\left\|\left[\mathbb{B}(f)\cdot U_l\Phi\right]_1-\left[\mathbb{B}(f)\cdot U_l\Phi'\right]_1\right\|\le k_1\cdot\|f\|\cdot\left\|[\Phi]_1-[\Phi']_1\right\|=0 .
\]
The rule $[\Phi]_1\mapsto\left[\mathbb{B}(f)\cdot U_l\Phi\right]_1$ therefore defines a map from the subspace of $\mathcal{H}_1$ formed by the classes of bounded $\Sigma_1^{+}$-measurable functions into $\mathcal{H}_1$;
this map is linear, and by the same estimate its norm does not exceed $k_1\cdot\|f\|$, nor $\exp\left(\frac{1}{2}\right)\cdot\|f\|$ when $f$ is real.
Let $\mathcal{B}_l(f)$ be its continuous extension to $\mathcal{H}_1$.
For $\Phi\in L^{\infty-}_{+}$ the truncations $\Phi\cdot1_{\{|\Phi|\le n\}}$, $n=1,2,\dots$, are bounded and $\Sigma_1^{+}$-measurable,
and they converge to $\Phi$ in $L^4(\mu)$ by dominated convergence; so their classes converge to $[\Phi]_1$.
Since $\mathbb{B}(f)\in L^4(\mu)$ and $U_l$ is isometric on $L^4(\mu)$, H\"older's inequality gives
\[
\left\|\mathbb{B}(f)\cdot U_l\left(\Phi\cdot1_{\{|\Phi|\le n\}}\right)-\mathbb{B}(f)\cdot U_l\Phi\right\|_{L^2(\mu)}
\le\|\mathbb{B}(f)\|_{L^4(\mu)}\cdot\left\|\Phi\cdot1_{\{|\Phi|\le n\}}-\Phi\right\|_{L^4(\mu)}\longrightarrow0 .
\]
Consequently
\[
\mathcal{B}_l(f)[\Phi]_1=\lim_{n\to\infty}\mathcal{B}_l(f)\left[\Phi\cdot1_{\{|\Phi|\le n\}}\right]_1=\lim_{n\to\infty}\left[\mathbb{B}(f)\cdot U_l\left(\Phi\cdot1_{\{|\Phi|\le n\}}\right)\right]_1=\left[\mathbb{B}(f)\cdot U_l\Phi\right]_1 .
\]
A bounded operator with this property agrees with $\mathcal{B}_l(f)$ on a dense set and therefore equals it.
For fixed $\Phi$ the right side depends complex linearly on $f$; by density, so does $\mathcal{B}_l(f)$.

\emph{Proof of \eqref{it:insertion-localisation}.}
Let $0\le c<c'\le l$ and $f\in L^2_{[c,c']}$.
Then $\tau_{-c}\cdot f\in L^2_{[0,c'-c]}$, and $\mathcal{B}_{c'-c}(\tau_{-c}\cdot f)$ is defined.
Let $\Phi\in L^{\infty-}_{+}$.
As $l-c'\ge0$, the function $U_{l-c'}\Phi$ lies in $L^{\infty-}_{+}$, and so does
\[
\Phi'':=\mathbb{B}\left(\tau_{-c}\cdot f\right)\cdot U_{c'-c}U_{l-c'}\Phi .
\]
The operator $U_c$ is multiplicative and $U_c\mathbb{B}(\tau_{-c}\cdot f)=\mathbb{B}(f)$; hence
\[
\mathbb{B}(f)\cdot U_l\Phi=U_c\Phi'' .
\]
For $s\ge0$ the operator $\exp(-s\cdot\mathsf{H}_1)$ maps the class of every element of $L^2(\Sigma_1^{+},\mu)$ to the class of its image under $U_s$ \cite[Proposition~8.19(ii)]{NoncompactQED}.
Therefore
\[
\begin{aligned}
\mathcal{B}_l(f)[\Phi]_1&=\left[U_c\Phi''\right]_1=\exp\left(-c\cdot\mathsf{H}_1\right)[\Phi'']_1
=\exp\left(-c\cdot\mathsf{H}_1\right)\,\mathcal{B}_{c'-c}\left(\tau_{-c}\cdot f\right)\left[U_{l-c'}\Phi\right]_1\\
&=\exp\left(-c\cdot\mathsf{H}_1\right)\,\mathcal{B}_{c'-c}\left(\tau_{-c}\cdot f\right)\,\exp\left(-(l-c')\cdot\mathsf{H}_1\right)[\Phi]_1 .
\end{aligned}
\]
The two sides of \eqref{it:insertion-localisation} are bounded operators that agree on the classes of $L^{\infty-}_{+}$, which are dense.

\emph{Proof of \eqref{it:insertion-symmetric}.}
Let $d>0$ and $f\in L^2_{[-d,d]}$.
Then $\tau_d\cdot f\in L^2_{[0,2d]}$, and $\|\tau_d\cdot f\|=\|f\|$ gives the bound on $\mathcal{B}^{\mathrm{c}}_d(f)$.
For $\Phi,\Phi'\in L^{\infty-}_{+}$ all integrands below are products of functions that lie in every $L^p(\mu)$, $p<\infty$, and are therefore integrable.
The definition of $\mathcal{B}_{2d}$ and the scalar product of classes give
\[
\left\langle[\Phi']_1,\mathcal{B}^{\mathrm{c}}_d(f)[\Phi]_1\right\rangle=\mu\left(\left(\Theta_1\Phi'\right)\cdot\mathbb{B}\left(\tau_d\cdot f\right)\cdot U_{2d}\Phi\right).
\]
By $(T_1)$ the integral does not change when $U_{-d}$ is applied to the integrand.
As $U_{-d}$ is multiplicative, $U_{-d}\mathbb{B}(\tau_d\cdot f)=\mathbb{B}(f)$, $U_{-d}U_{2d}=U_d$, and $U_{-d}\Theta_1=\Theta_1U_d$ by \eqref{eq:rotation-conjugation},
the image of the integrand is $\left(\Theta_1U_d\Phi'\right)\cdot\mathbb{B}(f)\cdot U_d\Phi$.
This is the formula for the matrix elements.

For the adjoint let $\Phi,\Phi'\in L^{\infty-}_{+}$.
The formula for the matrix elements, with the roles of $\Phi$ and $\Phi'$ exchanged, gives
\[
\begin{aligned}
\left\langle\mathcal{B}^{\mathrm{c}}_d(f)[\Phi']_1,[\Phi]_1\right\rangle
&=\overline{\mu\left(\left(\Theta_1U_d\Phi\right)\cdot\mathbb{B}(f)\cdot U_d\Phi'\right)}
=\mu\left(\left(\hat\theta_1U_d\Phi\right)\cdot\mathbb{B}\left(\overline{f}\right)\cdot\overline{U_d\Phi'}\right)\\
&=\mu\left(\left(U_d\Phi\right)\cdot\mathbb{B}\left(\theta_1\cdot\overline{f}\right)\cdot\left(\Theta_1U_d\Phi'\right)\right)
=\left\langle[\Phi']_1,\mathcal{B}^{\mathrm{c}}_d\left(\overline{\theta_1\cdot f}\right)[\Phi]_1\right\rangle .
\end{aligned}
\]
The second equality holds because complex conjugation turns $\Theta_1U_d\Phi$ into $\hat\theta_1U_d\Phi$ and $\mathbb{B}(f)$ into $\mathbb{B}\left(\overline{f}\right)$.
The third equality applies $\hat\theta_1$ to the integrand, which does not change the integral by $(R_1)$;
the operator $\hat\theta_1$ is multiplicative, $\hat\theta_1\hat\theta_1=1$,
$\hat\theta_1\mathbb{B}(b)=\mathbb{B}(\theta_1\cdot b)$ for $b\in L^2(\R^3;\C^3)$ (Lemma~\ref{lem:domination-consequences}(a)), and $\hat\theta_1\overline{U_d\Phi'}=\Theta_1U_d\Phi'$.
The fourth is the formula for the matrix elements with $\overline{\theta_1\cdot f}=\theta_1\cdot\overline{f}\in L^2_{[-d,d]}$ in the place of $f$.
Since the classes of $L^{\infty-}_{+}$ are dense, $\mathcal{B}^{\mathrm{c}}_d(f)^{*}=\mathcal{B}^{\mathrm{c}}_d\left(\overline{\theta_1\cdot f}\right)$.

\emph{Proof of \eqref{it:insertion-covariance}.}
Let $\mu$ satisfy moreover $(T_{123})$, so that $V_1$ and $\mathsf{P}_2$ are defined (Definition~\ref{def:laws}).
For $s\ge0$ the operators $\exp(-s\cdot\mathsf{H}_1)$ and $V_1(\pm a\cdot e_2)=\exp(\pm ia\cdot\mathsf{P}_2)$ commute,
and their product $\exp\left(-s\cdot\mathsf{H}_1\pm ia\cdot\mathsf{P}_2\right)$ maps the class of every element of $L^2(\Sigma_1^{+},\mu)$
to the class of its image under $U\left(s\cdot e_1\pm a\cdot e_2\right)$,
by \cite[Proposition~8.19(ii),~(iii)]{NoncompactQED} and Lemma~\ref{lem:joint-spectrum}(c).
Let $\Phi\in L^{\infty-}_{+}$, and put $\Phi'':=U\left((r-t)\cdot e_1-a\cdot e_2\right)\Phi$.
Since $r-t\ge0$, the function $\Phi''$ is measurable with respect to $\Sigma_{\{x_1>r-t\}}\subset\Sigma_1^{+}$,
and by $(T_{123})$ it lies in every $L^p(\mu)$;
so $\Phi''\in L^{\infty-}_{+}$ and $[\Phi'']_1=\exp\left(-(r-t)\cdot\mathsf{H}_1-ia\cdot\mathsf{P}_2\right)[\Phi]_1$.
Since moreover $r+t\ge0$ and $\mathbb{B}(\tau_d\cdot f)\cdot U_{2d}\Phi''\in L^{\infty-}_{+}$,
\[
\begin{aligned}
&\exp\left(-(r+t)\cdot\mathsf{H}_1+ia\cdot\mathsf{P}_2\right)\,\mathcal{B}^{\mathrm{c}}_d(f)\,\exp\left(-(r-t)\cdot\mathsf{H}_1-ia\cdot\mathsf{P}_2\right)[\Phi]_1\\
&\qquad=\left[U\left((r+t)\cdot e_1+a\cdot e_2\right)\left(\mathbb{B}\left(\tau_d\cdot f\right)\cdot U_{2d}\Phi''\right)\right]_1 .
\end{aligned}
\]
The translation operators are multiplicative, $U(x)\,U(x')=U(x+x')$, and $U(x)\mathbb{B}(b)=\mathbb{B}(\tau_x\cdot b)$ for $x,x'\in\R^3$ and $b\in L^2(\R^3;\C^3)$
(Lemma~\ref{lem:domination-consequences}(a) and $(T_{123})$).
As $(r+t)\cdot e_1+a\cdot e_2+d\cdot e_1=(d+r)\cdot e_1+(t\cdot e_1+a\cdot e_2)$,
and as the three translation vectors $(r+t)\cdot e_1+a\cdot e_2$, $2d\cdot e_1$ and $(r-t)\cdot e_1-a\cdot e_2$ add up to $2(d+r)\cdot e_1$,
the function in the last class is
\[
\mathbb{B}\left(\tau_{d+r}\cdot\left(\tau_{t e_1+a e_2}\cdot f\right)\right)\cdot U_{2(d+r)}\Phi .
\]
The field $\tau_{t e_1+a e_2}\cdot f$ vanishes almost everywhere outside $\{-d+t\le x_1\le d+t\}$, which lies in $\{-d-r\le x_1\le d+r\}$ because $|t|\le r$.
So it belongs to $L^2_{[-d-r,d+r]}$, and the class of the last function is $\mathcal{B}^{\mathrm{c}}_{d+r}\left(\tau_{t e_1+a e_2}\cdot f\right)[\Phi]_1$.
The two sides of \eqref{it:insertion-covariance} are bounded operators that agree on the dense set of the classes of $L^{\infty-}_{+}$.

\emph{Observables on $\{x_1>d\}$.}
Let $d>0$ and $f\in L^2_{[-d,d]}$, and let $\Phi_1,\Phi_2\in L^{\infty-}_{+}$ be measurable with respect to $\Sigma_{\{x_1>d\}}$.
The functions $U_{-d}\Phi_1$ and $U_{-d}\Phi_2$ are $\Sigma_1^{+}$-measurable and, by $(T_1)$, lie in every $L^p(\mu)$;
so they belong to $L^{\infty-}_{+}$.
The formula for the matrix elements in \eqref{it:insertion-symmetric}, with $U_{-d}\Phi_1$ in the place of $\Phi'$ and $U_{-d}\Phi_2$ in the place of $\Phi$, and the bound on $\mathcal{B}^{\mathrm{c}}_d(f)$ give
\[
\left|\mu\left(\left(\Theta_1\Phi_1\right)\cdot\mathbb{B}(f)\cdot\Phi_2\right)\right|
=\left|\left\langle\left[U_{-d}\Phi_1\right]_1,\mathcal{B}^{\mathrm{c}}_d(f)\left[U_{-d}\Phi_2\right]_1\right\rangle\right|
\le k_1\cdot\|f\|\cdot\left\|\left[U_{-d}\Phi_1\right]_1\right\|\cdot\left\|\left[U_{-d}\Phi_2\right]_1\right\| .
\]
\end{proof}

\subsection{Point plates and kernels}\label{sec:point-plates}

\subsubsection{The plate bound}\label{sec:plate-bound}

Let $\mu$ be a law on $\mathcal{S}'$ with $(T_{123})$, $(R_1)$ and $(RP_1)$,
and let $\mathfrak{E}_1$ be the joint spectral measure of its reconstruction in the direction $e_1$ (Lemma~\ref{lem:joint-spectrum}),
written in the variables $(\omega,p_\perp)$ with $p_\perp=(p_2,p_3)$.
It is carried by $[0,\infty)\times\R^2$ (Lemma~\ref{lem:joint-spectrum}(b)).
If $\mu$ satisfies the cone $(C_{12})$ as well,
the open set $\{\omega<|p_2|\}$ is an $\mathfrak{E}_1$-null set,
so that $\mathfrak{E}_1$ is carried by the closed set $\{\omega\ge|p_2|\}$.
At each point of this set and for each $\xi=(\xi^0,\xi^1)\in\C^2$,
\begin{equation}\label{eq:rotation-semigroup-bound}
\left|\exp\left(-\xi^0\omega+i\xi^1p_2\right)\right|
=\exp\left(-\omega\Re\xi^0-p_2\Im\xi^1\right)
\le\exp\left(-d_{\mathcal{W}}(\xi)\cdot\omega\right),
\end{equation}
because $|p_2\Im\xi^1|\le\omega\cdot|\Im\xi^1|$.
For $\xi$ in the closed cone $\overline{\mathcal{W}}$ the right side is at most one.
For such $\xi$ the continuous function $(\omega,p_\perp)\mapsto\exp\left(-\xi^0\omega+i\xi^1p_2\right)$
has therefore modulus at most one on the set $\{\omega\ge|p_2|\}$ that carries $\mathfrak{E}_1$,
and its integral against $\mathfrak{E}_1$, formally the operator $\exp\left(-\xi^0\mathsf{H}_1+i\xi^1\mathsf{P}_2\right)$, is a bounded operator.

\begin{lemma}[The semigroup of the cone]\label{lem:cone-semigroup}
Let $\mu$ be a law on $\mathcal{S}'$ that satisfies \textup{(G)}, $(T_{123})$, $(R_1)$, $(RP_1)$ and $(C_{12})$,
and for $\xi\in\overline{\mathcal{W}}$ let
\[
\mathsf{T}(\xi):=\int\mathrm{d}\mathfrak{E}_1(\omega,p_\perp)\,\exp\left(-\xi^0\omega+i\xi^1p_2\right).
\]
Each $\mathsf{T}(\xi)$ is a bounded operator on $\mathcal{H}_1$, and the following hold.
\begin{enumerate}
\item\label{it:semigroup-contraction}
For all $\xi,\xi'\in\overline{\mathcal{W}}$,
\[
\left\|\mathsf{T}(\xi)\right\|\le1,\qquad
\mathsf{T}(\xi)\,\mathsf{T}(\xi')=\mathsf{T}(\xi+\xi'),\qquad
\mathsf{T}(\xi)^*=\mathsf{T}(\xi^*),\qquad
\mathsf{T}(\xi)\,\Omega_1=\Omega_1 .
\]
\item\label{it:semigroup-holomorphic}
The map $\xi\mapsto\mathsf{T}(\xi)$ is strongly continuous on $\overline{\mathcal{W}}$ and holomorphic in the operator norm on $\mathcal{W}$.
\item\label{it:semigroup-translation}
Let $y=(y^0,y^1)\in\R^2$ with $y^0\ge0$, read as the vector $y^0e_1+y^1e_2$ of $\R^3$ in $U(y)$.
Then
\[
\mathsf{T}(y)\,[\Phi]_1=[U(y)\Phi]_1\qquad\text{for every }\Phi\in L^2(\Sigma_1^+,\mu).
\]
\end{enumerate}
\end{lemma}

\begin{proof}
By Lemma~\ref{lem:joint-spectrum}(b) and $(C_{12})$ the projection-valued measure $\mathfrak{E}_1$ is carried by $\{\omega\ge|p_2|\}$,
and on this set the integrand of $\mathsf{T}(\xi)$ is a continuous function of modulus at most one, by~\eqref{eq:rotation-semigroup-bound}.
Hence $\mathsf{T}(\xi)$ is a bounded operator with $\|\mathsf{T}(\xi)\|\le1$, and for every $\psi\in\mathcal{H}_1$
\[
\left\|\mathsf{T}(\xi)\psi\right\|^2=\int\mathrm{d}\varrho_\psi(\omega,p_\perp)\,\left|\exp\left(-\xi^0\omega+i\xi^1p_2\right)\right|^2 .
\]
Integration against $\mathfrak{E}_1$ is multiplicative on bounded Borel functions,
and it takes the complex conjugate of a function to the adjoint of its integral.
For $\xi,\xi'\in\overline{\mathcal{W}}$ the product of the integrands of $\mathsf{T}(\xi)$ and $\mathsf{T}(\xi')$ is the integrand of $\mathsf{T}(\xi+\xi')$,
and $\xi+\xi'$ lies in $\overline{\mathcal{W}}$, since the triangle inequality gives $d_{\mathcal{W}}(\xi+\xi')\ge d_{\mathcal{W}}(\xi)+d_{\mathcal{W}}(\xi')$.
The complex conjugate of the integrand of $\mathsf{T}(\xi)$ is $\exp\left(-\overline{\xi^0}\omega-i\overline{\xi^1}p_2\right)$,
which is the integrand at $\xi^*=\left(\overline{\xi^0},-\overline{\xi^1}\right)$,
and $\xi^*\in\overline{\mathcal{W}}$ because $d_{\mathcal{W}}(\xi^*)=d_{\mathcal{W}}(\xi)$ (Lemma~\ref{lem:forward-tube}\eqref{it:tube-adjoint}).
Thus $\mathsf{T}(\xi)\mathsf{T}(\xi')=\mathsf{T}(\xi+\xi')$ and $\mathsf{T}(\xi)^*=\mathsf{T}(\xi^*)$.

The vacuum $\Omega_1=[1]_1$ is a unit vector, as $\langle[1]_1,[1]_1\rangle=\mu(1)=1$;
the constant function $1$ is invariant under every $U(a)$,
and Definition~\ref{def:laws} gives $\exp\left(-t\mathsf{H}_1\right)V_1(y)\Omega_1=[U(te_1)U(y)1]_1=\Omega_1$ for all $t\ge0$ and $y\in e_1^\perp$.
By Lemma~\ref{lem:joint-spectrum}(b) the spectral measure of $\Omega_1$ therefore satisfies
\[
\int\mathrm{d}\varrho_{\Omega_1}(\omega,p_\perp)\,\exp\left(-t\omega+ip_\perp\cdot y\right)
=\left\langle\Omega_1,\exp\left(-t\mathsf{H}_1\right)V_1(y)\Omega_1\right\rangle=1
\qquad\left(t\ge0,\ y\in e_1^\perp\right),
\]
and the unit mass at the point $(\omega,p_\perp)=(0,0)$ has the same transform.
Both measures are finite and positive on $[0,\infty)\times\R^2$,
so Lemma~\ref{lem:laplace-fourier} identifies $\varrho_{\Omega_1}$ with this unit mass.
As the integrand of $\mathsf{T}(\xi)$ takes the value one at $(0,0)$,
\[
\left\|\mathsf{T}(\xi)\Omega_1-\Omega_1\right\|^2
=\int\mathrm{d}\varrho_{\Omega_1}(\omega,p_\perp)\,\left|\exp\left(-\xi^0\omega+i\xi^1p_2\right)-1\right|^2=0 ,
\]
and the last identity of~\eqref{it:semigroup-contraction} follows.

Turning to~\eqref{it:semigroup-holomorphic}, let $\psi\in\mathcal{H}_1$ and $\xi,\xi_0\in\overline{\mathcal{W}}$. Then
\[
\left\|\mathsf{T}(\xi)\psi-\mathsf{T}(\xi_0)\psi\right\|^2
=\int\mathrm{d}\varrho_\psi(\omega,p_\perp)\,\left|\exp\left(-\xi^0\omega+i\xi^1p_2\right)-\exp\left(-\xi_0^0\omega+i\xi_0^1p_2\right)\right|^2 .
\]
As $\xi\to\xi_0$ the integrand tends to zero at every point,
and on $\{\omega\ge|p_2|\}$, which carries the finite measure $\varrho_\psi$, it is bounded by $4$;
by dominated convergence, $\mathsf{T}$ is strongly continuous on $\overline{\mathcal{W}}$.
Now fix $\varepsilon>0$.
The function $d_{\mathcal{W}}$ is continuous and concave, so the set $\{d_{\mathcal{W}}>\varepsilon\}$ is open and convex.
Let $\xi$ and $\xi'$ lie in this set, put $\xi_s:=\xi+s(\xi'-\xi)$ for $0\le s\le1$ and $|\xi'-\xi|_1:=|\xi'^0-\xi^0|+|\xi'^1-\xi^1|$,
and let $\omega\ge|p_2|$.
The first and the second derivative of $s\mapsto\exp\left(-\xi_s^0\omega+i\xi_s^1p_2\right)$ are this function multiplied by
$-(\xi'^0-\xi^0)\omega+i(\xi'^1-\xi^1)p_2$ and by the square of this number.
Since $d_{\mathcal{W}}(\xi_s)>\varepsilon$ and $|p_2|\le\omega$, the bound~\eqref{eq:rotation-semigroup-bound} shows that their moduli are at most
\[
|\xi'-\xi|_1\cdot\omega\cdot\exp(-\varepsilon\omega)\le\frac{|\xi'-\xi|_1}{\exp(1)\cdot\varepsilon}
\qquad\text{and}\qquad
|\xi'-\xi|_1^2\cdot\omega^2\cdot\exp(-\varepsilon\omega)\le\frac{4\,|\xi'-\xi|_1^2}{\left(\exp(1)\cdot\varepsilon\right)^2},
\]
for the maxima of $\omega\exp(-\varepsilon\omega)$ and of $\omega^2\exp(-\varepsilon\omega)$ on $[0,\infty)$ are $\left(\exp(1)\cdot\varepsilon\right)^{-1}$ and $4\left(\exp(1)\cdot\varepsilon\right)^{-2}$.
Taylor's formula with integral remainder, applied at $s=1$, bounds the remainder by one half of the bound on the second derivative,
uniformly on $\{\omega\ge|p_2|\}$.
Integrating against $\mathfrak{E}_1$ one obtains
\[
\begin{aligned}
&\left\|\mathsf{T}(\xi')-\mathsf{T}(\xi)-\int\mathrm{d}\mathfrak{E}_1(\omega,p_\perp)\,\left(-(\xi'^0-\xi^0)\omega+i(\xi'^1-\xi^1)p_2\right)\exp\left(-\xi^0\omega+i\xi^1p_2\right)\right\|\\
&\qquad\le\frac{2\,|\xi'-\xi|_1^2}{\left(\exp(1)\cdot\varepsilon\right)^2}.
\end{aligned}
\]
The integral on the left is complex linear in $\xi'-\xi$, and by the bound on the first derivative it is an operator of norm at most $|\xi'-\xi|_1\cdot\left(\exp(1)\cdot\varepsilon\right)^{-1}$.
So $\mathsf{T}$ is complex differentiable in the operator norm at every point of $\{d_{\mathcal{W}}>\varepsilon\}$;
every point of $\mathcal{W}$ lies in such a set, and $\mathsf{T}$ is holomorphic in the operator norm on $\mathcal{W}$.

For~\eqref{it:semigroup-translation}, let $y=(y^0,y^1)$ be real with $y^0\ge0$.
As $p_\perp\cdot(y^1e_2)=y^1p_2$, Lemma~\ref{lem:joint-spectrum}(b), applied with $t=y^0$ and the vector $y^1e_2$ of $e_1^\perp$, gives
$\mathsf{T}(y)=\exp\left(-y^0\mathsf{H}_1\right)V_1(y^1e_2)$.
For $\Phi\in L^2(\Sigma_1^+,\mu)$, $t\ge0$ and $a\in e_1^\perp$ the reconstruction satisfies
$\exp\left(-t\mathsf{H}_1\right)[\Phi]_1=[U(te_1)\Phi]_1$ and $V_1(a)[\Phi]_1=[U(a)\Phi]_1$ (Definition~\ref{def:laws}),
and $U(y^1e_2)$ maps $L^2(\Sigma_1^+,\mu)$ onto itself.
Therefore $\mathsf{T}(y)[\Phi]_1=[U(y^0e_1)U(y^1e_2)\Phi]_1=[U(y)\Phi]_1$.
\end{proof}

Let $\delta_0$ be the Dirac delta at the origin of the plane.
For $\xi,\xi'\in\mathcal{W}$ and a label $\lambda$, the operator $\mathsf{T}(\xi)\,\mathbb{B}_\lambda(\delta_0)\,\mathsf{T}(\xi')$ has a formal meaning only:
the field is smeared with a plate concentrated at one point of the plane,
and placed between two operators of the semigroup.
The point plates $\mathcal{B}_\lambda(\xi,\xi')$ are bounded operators that take its place.

\begin{proposition}[Point plates]\label{prop:point-plates}
Let $\mu$ be a law on $\mathcal{S}'$ that satisfies \textup{(G)}, $(T_{123})$, $(R_1)$, $(RP_1)$ and $(C_{12})$, and put
\[
k_2:=\frac{8}{\pi}\exp\left(\frac12\right).
\]
There exist bounded operators $\mathcal{B}_\lambda(\xi,\xi')$ on $\mathcal{H}_1$,
one for each label $\lambda=(\sigma,\chi)$ and each pair $(\xi,\xi')\in\mathcal{W}\times\mathcal{W}$,
which depend holomorphically on $(\xi,\xi')$ in the operator norm and real linearly on $\chi$,
and which have the following properties.
\begin{enumerate}
\item\label{it:point-plate-bound}
For all $\xi,\xi'\in\mathcal{W}$,
\[
\left\|\mathcal{B}_\lambda(\xi,\xi')\right\|\le\frac{k_2\cdot\|\chi\|_2}{\min\left(d_{\mathcal{W}}(\xi),d_{\mathcal{W}}(\xi')\right)} .
\]
\item\label{it:point-plate-shift}
For all $\xi,\xi'\in\mathcal{W}$ and all $\kappa,\kappa'\in\overline{\mathcal{W}}$,
\[
\mathsf{T}(\kappa)\,\mathcal{B}_\lambda(\xi,\xi')\,\mathsf{T}(\kappa')=\mathcal{B}_\lambda(\kappa+\xi,\xi'+\kappa'),
\qquad
\mathcal{B}_\lambda(\xi,\xi')^*=s(\sigma)\cdot\mathcal{B}_\lambda(\xi'^*,\xi^*).
\]
\item\label{it:point-plate-integral}
For every $d>0$ and every $\alpha\in L^2(\R^2;\C)$ with compact support in $\{|y^0|<d\}$,
\[
\mathcal{B}^{\mathrm{c}}_d(\alpha\otimes\lambda)=\int\mathrm{d}y\,\alpha(y)\,\mathcal{B}_\lambda\left(d\,n_0+y,d\,n_0-y\right),
\]
a Bochner integral in the operator norm.
\end{enumerate}
Consequently, if $\alpha$ is as in~\eqref{it:point-plate-integral} and $\supp\alpha\subset\{|y^0|\le d'\}$ for some $d'<d$, then
\[
\left\|\mathcal{B}^{\mathrm{c}}_d(\alpha\otimes\lambda)\right\|\le\frac{k_2\cdot\|\chi\|_2\cdot\|\alpha\|_{L^1(\R^2)}}{d-d'} .
\]
\end{proposition}

\begin{proof}
Fix the label $\lambda=(\sigma,\chi)$.
For $\delta>0$ let
\[
j_\delta:=\left(\pi\delta^2\right)^{-1}1_{\{|y|<\delta\}}
\]
be the normalised indicator function of the disc of radius $\delta$ about the origin of the plane;
it is real, invariant under the reflection $\theta_1$ of the plane, and $\|j_\delta\|_{L^1(\R^2)}=1$.
The plate $j_\delta\otimes\lambda$ lies in $L^2_{[-\delta,\delta]}$. Put
\[
A_\delta:=\mathcal{B}^{\mathrm{c}}_\delta\left(j_\delta\otimes\lambda\right)
\]
and, for $\xi,\xi'\in\C^2$ with $d_{\mathcal{W}}(\xi)\ge\delta$ and $d_{\mathcal{W}}(\xi')\ge\delta$,
\[
\mathcal{B}^{\delta}_\lambda(\xi,\xi'):=\mathsf{T}(\xi-\delta n_0)\,A_\delta\,\mathsf{T}(\xi'-\delta n_0);
\]
here $\xi-\delta n_0$ and $\xi'-\delta n_0$ lie in $\overline{\mathcal{W}}$, since $d_{\mathcal{W}}(\xi-\delta n_0)=d_{\mathcal{W}}(\xi)-\delta$.

These regularised point plates have three properties.
First, for $\kappa,\kappa'\in\overline{\mathcal{W}}$ the semigroup law of Lemma~\ref{lem:cone-semigroup}\eqref{it:semigroup-contraction} gives
\[
\mathsf{T}(\kappa)\,\mathcal{B}^{\delta}_\lambda(\xi,\xi')\,\mathsf{T}(\kappa')=\mathcal{B}^{\delta}_\lambda(\kappa+\xi,\xi'+\kappa'),
\]
the right side being defined because $d_{\mathcal{W}}(\kappa+\xi)\ge d_{\mathcal{W}}(\kappa)+d_{\mathcal{W}}(\xi)\ge\delta$, and likewise $d_{\mathcal{W}}(\xi'+\kappa')\ge\delta$.
Second, Proposition~\ref{prop:insertion-bound}(b) and Lemma~\ref{lem:plate-identities}\eqref{it:plate-reflection} give
\[
A_\delta^*=\mathcal{B}^{\mathrm{c}}_\delta\left(\overline{\theta_1\cdot(j_\delta\otimes\lambda)}\right)
=s(\sigma)\cdot\mathcal{B}^{\mathrm{c}}_\delta\left((j_\delta\circ\theta_1)\otimes\lambda\right)
=s(\sigma)\cdot A_\delta ,
\]
because $j_\delta$ is real and $j_\delta\circ\theta_1=j_\delta$.
The map $\xi\mapsto\xi^*$ is real linear, $\delta$ is real,
and $n_0^*=n_0$, because $n_0=(1,0)$ is a real point with vanishing second coordinate;
hence $(\xi-\delta n_0)^*=\xi^*-\delta n_0$.
Together with $\mathsf{T}(\xi)^*=\mathsf{T}(\xi^*)$ and $d_{\mathcal{W}}(\xi^*)=d_{\mathcal{W}}(\xi)$ (Lemma~\ref{lem:forward-tube}\eqref{it:tube-adjoint}) this yields
\[
\mathcal{B}^{\delta}_\lambda(\xi,\xi')^*=s(\sigma)\cdot\mathcal{B}^{\delta}_\lambda(\xi'^*,\xi^*).
\]
So $\mathcal{B}^{\delta}_\lambda$ satisfies the two identities of~\eqref{it:point-plate-shift}.
Third, let $y=(y^0,y^1)$ be real and $r\ge|y^0|$.
The points $r\,n_0\pm y$ are real with first coordinate $r\pm y^0\ge0$,
and Lemma~\ref{lem:joint-spectrum}(b),~(c) identifies
\[
\mathsf{T}(r\,n_0\pm y)=\exp\left(-(r\pm y^0)\mathsf{H}_1\right)V_1\left(\pm y^1e_2\right)=\exp\left(-(r\pm y^0)\mathsf{H}_1\pm iy^1\mathsf{P}_2\right).
\]
Moreover $\tau_y\cdot(j_\delta\otimes\lambda)=(\tau_y\cdot j_\delta)\otimes\lambda$, where $(\tau_y\cdot j_\delta)(y')=j_\delta(y'-y)$.
Proposition~\ref{prop:insertion-bound}(c), applied to $j_\delta\otimes\lambda\in L^2_{[-\delta,\delta]}$ with the time shift $y^0$ and the spatial shift $y^1$
(admissible, as $r\ge|y^0|$), therefore gives
\begin{equation}\label{eq:rotation-plate-shift}
\mathcal{B}^{\delta}_\lambda\left((r+\delta)n_0+y,(r+\delta)n_0-y\right)
=\mathsf{T}(r\,n_0+y)\,A_\delta\,\mathsf{T}(r\,n_0-y)
=\mathcal{B}^{\mathrm{c}}_{\delta+r}\left((\tau_y\cdot j_\delta)\otimes\lambda\right)
\end{equation}
for every real $y$ and every $r\ge|y^0|$.

\emph{Step~1: a uniform bound.}
Fix $\xi,\xi'\in\mathcal{W}$, put $d_0:=\min\left(d_{\mathcal{W}}(\xi),d_{\mathcal{W}}(\xi')\right)$,
let $0<\delta<d_0$ and $r:=(d_0-\delta)/2$,
and let $\psi,\psi'$ be unit vectors of $\mathcal{H}_1$.
On the tube
\[
S_r:=\left\{z=(z^0,z^1)\in\C^2:\ |\Re z^0|<r,\ |\Im z^1|<r\right\}
\]
one has
\[
d_{\mathcal{W}}(\xi+z-\delta n_0)\ge d_{\mathcal{W}}(\xi)-\delta-|\Re z^0|-|\Im z^1|>d_0-\delta-2r=0,
\]
and in the same way $d_{\mathcal{W}}(\xi'-z-\delta n_0)>0$.
The function
\[
h(z):=\left\langle\psi,\mathcal{B}^{\delta}_\lambda(\xi+z,\xi'-z)\,\psi'\right\rangle,\qquad z\in S_r,
\]
is therefore defined, and it is holomorphic on the open set $S_r$ by Lemma~\ref{lem:cone-semigroup}\eqref{it:semigroup-holomorphic}.
For every $r'$ with $0<r'<r$ the tube $S_r$ contains the closed polydisc $\left\{|z^0|\le r',\ |z^1|\le r'\right\}$.

Write $z=y+i(t,t')$ with $y=(y^0,y^1)\in\R^2$ and $t,t'\in\R$;
then $z\in S_r$ means $|y^0|<r$ and $|t'|<r$, and
\[
\begin{aligned}
\xi+z-\delta n_0&=\left[\xi-(\delta+r)n_0+i(t,t')\right]+\left[r\,n_0+y\right],\\
\xi'-z-\delta n_0&=\left[r\,n_0-y\right]+\left[\xi'-(\delta+r)n_0-i(t,t')\right].
\end{aligned}
\]
The four brackets lie in $\overline{\mathcal{W}}$.
Indeed, $d_{\mathcal{W}}(\xi)-\delta\ge2r$, so that $d_{\mathcal{W}}$ of the first bracket is at least $d_{\mathcal{W}}(\xi)-\delta-r-|t'|\ge r-|t'|>0$;
the last bracket satisfies $d_{\mathcal{W}}\ge d_{\mathcal{W}}(\xi')-\delta-r-|t'|\ge r-|t'|>0$ for the same reason;
and $d_{\mathcal{W}}(r\,n_0\pm y)=r\pm y^0>0$.
The outer brackets can thus carry the imaginary part $t'$ of the spatial coordinate,
as $d_{\mathcal{W}}$ is at least $r>|t'|$ at $\xi-(\delta+r)n_0$ and at $\xi'-(\delta+r)n_0$:
a point of $\overline{\mathcal{W}}$ has an imaginary spatial part only together with a real time part at least as large.
By Lemma~\ref{lem:cone-semigroup}\eqref{it:semigroup-contraction} and~\eqref{eq:rotation-plate-shift},
\[
h\left(y+i(t,t')\right)=\left\langle\psi_1,\mathcal{B}^{\mathrm{c}}_{\delta+r}\left((\tau_y\cdot j_\delta)\otimes\lambda\right)\psi'_1\right\rangle,
\]
where the vectors
\[
\psi_1:=\mathsf{T}\left(\xi-(\delta+r)n_0+i(t,t')\right)^*\psi,
\qquad
\psi'_1:=\mathsf{T}\left(\xi'-(\delta+r)n_0-i(t,t')\right)\psi'
\]
have norm at most one and do not depend on $y$.

Now fix $t,t'$ with $|t|<r$ and $|t'|<r$, and let $\eta\in L^2\left((-r,r)^2\right)$, extended by zero to $\R^2$.
As translation acts continuously on $L^2(\R^2)$ and $\eta$ is integrable,
the $L^2(\R^2)$-valued function $y\mapsto\eta(y)\,\tau_y\cdot j_\delta$ is Bochner integrable,
and its integral is the convolution $\eta*j_\delta$.
This convolution vanishes outside $\{|y^0|<r+\delta\}$,
because $\eta$ vanishes outside $(-r,r)^2$ and $j_\delta$ outside the disc of radius $\delta$.
By Proposition~\ref{prop:insertion-bound}(b) the map $\alpha\mapsto\left\langle\psi_1,\mathcal{B}^{\mathrm{c}}_{\delta+r}(\alpha\otimes\lambda)\,\psi'_1\right\rangle$
is a bounded linear form on the functions $\alpha\in L^2(\R^2;\C)$ that vanish outside $\{|y^0|\le\delta+r\}$,
and as such it commutes with the Bochner integral.
With the bound of Proposition~\ref{prop:insertion-bound}(b),
with $\|\alpha\otimes\lambda\|\le2^{\frac12}\cdot\|\alpha\|_2\cdot\|\chi\|_2$ (as $|e^\sigma|\le2^{\frac12}$)
and with Young's inequality $\|\eta*j_\delta\|_2\le\|\eta\|_2\cdot\|j_\delta\|_{L^1(\R^2)}=\|\eta\|_2$,
one obtains
\[
\begin{aligned}
\left|\int\mathrm{d}y\,\eta(y)\,h\left(y+i(t,t')\right)\right|
&=\left|\left\langle\psi_1,\mathcal{B}^{\mathrm{c}}_{\delta+r}\left((\eta*j_\delta)\otimes\lambda\right)\psi'_1\right\rangle\right|\\
&\le2^{\frac12}\cdot k_1\cdot\|\eta*j_\delta\|_2\cdot\|\chi\|_2
\le2\exp\left(\frac12\right)\cdot\|\eta\|_2\cdot\|\chi\|_2 ,
\end{aligned}
\]
where $2^{\frac12}\cdot k_1=2\exp\left(\frac12\right)$.
The function $y\mapsto h\left(y+i(t,t')\right)$ is continuous on $(-r,r)^2$ and bounded there by $\|A_\delta\|$;
taking for $\eta$ its complex conjugate one finds
\[
\int_{(-r,r)^2}\mathrm{d}y\,\left|h\left(y+i(t,t')\right)\right|^2\le4\exp(1)\cdot\|\chi\|_2^2
\qquad\left(|t|<r,\ |t'|<r\right).
\]

Let $0<r'<r$.
The open polydisc $\left\{|z^0|<r',\ |z^1|<r'\right\}$ is contained in the box of the points $y+i(t,t')$ with $|y^0|,|y^1|,|t|,|t'|<r$,
so Lemma~\ref{lem:analysis-facts}\eqref{it:mean-value} and Tonelli's theorem give
\[
\begin{aligned}
|h(0)|^2
&\le\left(\pi r'^2\right)^{-2}\int_{-r}^{r}\mathrm{d}t\int_{-r}^{r}\mathrm{d}t'\int_{(-r,r)^2}\mathrm{d}y\,\left|h\left(y+i(t,t')\right)\right|^2\\
&\le\left(\pi r'^2\right)^{-2}\cdot(2r)^2\cdot4\exp(1)\cdot\|\chi\|_2^2 .
\end{aligned}
\]
Let $r'$ tend to $r$.
As $h(0)=\left\langle\psi,\mathcal{B}^{\delta}_\lambda(\xi,\xi')\,\psi'\right\rangle$ for arbitrary unit vectors $\psi,\psi'$,
and as $r=(d_0-\delta)/2$, it follows that
\begin{equation}\label{eq:rotation-plate-estimate}
\left\|\mathcal{B}^{\delta}_\lambda(\xi,\xi')\right\|\le\frac{4\exp\left(\frac12\right)\cdot\|\chi\|_2}{\pi r}=\frac{k_2\cdot\|\chi\|_2}{d_0-\delta},
\qquad 0<\delta<d_0=\min\left(d_{\mathcal{W}}(\xi),d_{\mathcal{W}}(\xi')\right).
\end{equation}

\emph{Step~2: the limit $\delta\to0$.}
Keep $\xi$, $\xi'$, $d_0$, $\psi$ and $\psi'$, put $r_0:=d_0/4$, and let
\[
D_{r_0}:=\left\{z\in\C^2:\ |z^0|<r_0,\ |z^1|<r_0\right\}.
\]
For $z\in D_{r_0}$ one has $d_{\mathcal{W}}(\xi+z)>d_0-2r_0=d_0/2$ and $d_{\mathcal{W}}(\xi'-z)>d_0/2$.
Hence, for $0<\delta<r_0$, the function
\[
h_\delta(z):=\left\langle\psi,\mathcal{B}^{\delta}_\lambda(\xi+z,\xi'-z)\,\psi'\right\rangle
\]
is defined and holomorphic on $D_{r_0}$,
and~\eqref{eq:rotation-plate-estimate} at the points $(\xi+z,\xi'-z)$ gives
\[
|h_\delta(z)|\le\frac{2k_2\cdot\|\chi\|_2}{d_0-2\delta}\le\frac{4k_2\cdot\|\chi\|_2}{d_0}\qquad(z\in D_{r_0}).
\]
Let $\eta\in L^2\left((-r_0,r_0)^2\right)$.
For real $y\in(-r_0,r_0)^2$,
\[
\begin{aligned}
\xi+y-\delta n_0&=\left[\xi-(\delta+r_0)n_0\right]+\left[r_0n_0+y\right],\\
\xi'-y-\delta n_0&=\left[r_0n_0-y\right]+\left[\xi'-(\delta+r_0)n_0\right],
\end{aligned}
\]
with $d_{\mathcal{W}}\left(\xi-(\delta+r_0)n_0\right)=d_{\mathcal{W}}(\xi)-\delta-r_0>d_0/2$, and likewise for $\xi'$.
Put $f_\delta:=(\eta*j_\delta)\otimes\lambda$.
Lemma~\ref{lem:cone-semigroup}\eqref{it:semigroup-contraction} and~\eqref{eq:rotation-plate-shift} with $r=r_0$,
followed by the integration against $\eta$ of Step~1, give
\[
\int\mathrm{d}y\,\eta(y)\,h_\delta(y)
=\left\langle\psi,\mathsf{T}\left(\xi-(\delta+r_0)n_0\right)\mathcal{B}^{\mathrm{c}}_{\delta+r_0}(f_\delta)\,\mathsf{T}\left(\xi'-(\delta+r_0)n_0\right)\psi'\right\rangle .
\]
Proposition~\ref{prop:insertion-bound}(c), with no shift and with $r_0-\delta$ in the place of $r$,
together with $\mathsf{T}\left((r_0-\delta)n_0\right)=\exp\left(-(r_0-\delta)\mathsf{H}_1\right)$ (the identification above with $y=0$), gives
\[
\mathcal{B}^{\mathrm{c}}_{2r_0}(f_\delta)=\mathsf{T}\left((r_0-\delta)n_0\right)\mathcal{B}^{\mathrm{c}}_{\delta+r_0}(f_\delta)\,\mathsf{T}\left((r_0-\delta)n_0\right),
\]
and $\mathsf{T}\left(\xi-(\delta+r_0)n_0\right)=\mathsf{T}\left(\xi-2r_0n_0\right)\mathsf{T}\left((r_0-\delta)n_0\right)$,
where $d_{\mathcal{W}}\left(\xi-2r_0n_0\right)\ge d_0/2$; the same holds for $\xi'$.
Consequently
\[
\int\mathrm{d}y\,\eta(y)\,h_\delta(y)
=\left\langle\psi,\mathsf{T}\left(\xi-2r_0n_0\right)\mathcal{B}^{\mathrm{c}}_{2r_0}\left((\eta*j_\delta)\otimes\lambda\right)\mathsf{T}\left(\xi'-2r_0n_0\right)\psi'\right\rangle,
\]
and as $\delta\to0$ this tends to
\[
\left\langle\psi,\mathsf{T}\left(\xi-2r_0n_0\right)\mathcal{B}^{\mathrm{c}}_{2r_0}\left(\eta\otimes\lambda\right)\mathsf{T}\left(\xi'-2r_0n_0\right)\psi'\right\rangle .
\]
To see the convergence, observe that $j_\delta\ge0$ has integral one and vanishes outside the disc of radius $\delta$;
by the continuity of translation on $L^2(\R^2)$ it follows that $\|\eta*j_\delta-\eta\|_2\to0$.
For $\delta<r_0$ the plates $(\eta*j_\delta)\otimes\lambda$ lie in $L^2_{[-2r_0,2r_0]}$, and they tend to $\eta\otimes\lambda$ in norm;
moreover $\mathcal{B}^{\mathrm{c}}_{2r_0}$ is a bounded linear map on this space (Proposition~\ref{prop:insertion-bound}(b)).

The functions $h_\delta$, $0<\delta<r_0$, are holomorphic and uniformly bounded on $D_{r_0}$,
so they form a normal family (Lemma~\ref{lem:analysis-facts}\eqref{it:normal-families}).
Hence, if $(\delta_n)$ is a sequence in $(0,r_0)$ with limit zero,
some subsequence of $(h_{\delta_n})$ converges uniformly on each compact subset of $D_{r_0}$.
Let $h^{(1)}$ and $h^{(2)}$ be two limits obtained in this way, possibly from different sequences.
Both are holomorphic on $D_{r_0}$ and bounded by $4k_2\cdot\|\chi\|_2/d_0$,
and by dominated convergence their integrals against every $\eta\in L^2\left((-r_0,r_0)^2\right)$ are equal to the limit just computed.
Being continuous, $h^{(1)}$ and $h^{(2)}$ agree on the real square $(-r_0,r_0)^2$,
hence on the connected open set $D_{r_0}$ by Lemma~\ref{lem:analysis-facts}\eqref{it:identity-theorem}.
So $h_\delta$ converges on $D_{r_0}$ as $\delta\to0$;
in particular $\left\langle\psi,\mathcal{B}^{\delta}_\lambda(\xi,\xi')\,\psi'\right\rangle=h_\delta(0)$ has a limit.
By homogeneity the limit exists for all $\psi,\psi'\in\mathcal{H}_1$.
It is a sesquilinear form in $(\psi,\psi')$, and by~\eqref{eq:rotation-plate-estimate} its modulus is at most $k_2\cdot\|\chi\|_2\cdot\|\psi\|\cdot\|\psi'\|/d_0$.
Let $\mathcal{B}_\lambda(\xi,\xi')$ be the bounded operator that represents this form (Riesz' theorem on bounded sesquilinear forms).
It is the weak limit of $\mathcal{B}^{\delta}_\lambda(\xi,\xi')$ as $\delta\to0$, and its norm is at most $k_2\cdot\|\chi\|_2/d_0$: this is~\eqref{it:point-plate-bound}.
The identities of~\eqref{it:point-plate-shift} hold for every $\mathcal{B}^{\delta}_\lambda$;
since $\mathsf{T}(\kappa)$ and $\mathsf{T}(\kappa')$ are bounded and the adjoint is continuous in the weak operator topology,
they hold for the weak limits as well.
Finally, $A_\delta$ is real linear in $\chi$, and therefore so are $\mathcal{B}^{\delta}_\lambda(\xi,\xi')$ and $\mathcal{B}_\lambda(\xi,\xi')$.

\emph{Step~3: holomorphy.}
Let $(\xi_0,\xi'_0)\in\mathcal{W}\times\mathcal{W}$ and $0<\varepsilon<\min\left(d_{\mathcal{W}}(\xi_0),d_{\mathcal{W}}(\xi'_0)\right)$.
The pairs $(\xi,\xi')$ with $d_{\mathcal{W}}(\xi-\xi_0+\varepsilon n_0)>0$ and $d_{\mathcal{W}}(\xi'-\xi'_0+\varepsilon n_0)>0$ form an open neighbourhood of $(\xi_0,\xi'_0)$,
and the points $\xi_0-\varepsilon n_0$ and $\xi'_0-\varepsilon n_0$ lie in $\mathcal{W}$.
On this neighbourhood~\eqref{it:point-plate-shift} gives
\[
\mathcal{B}_\lambda(\xi,\xi')=\mathsf{T}(\xi-\xi_0+\varepsilon n_0)\,\mathcal{B}_\lambda(\xi_0-\varepsilon n_0,\xi'_0-\varepsilon n_0)\,\mathsf{T}(\xi'-\xi'_0+\varepsilon n_0),
\]
and the right side is holomorphic in the operator norm by Lemma~\ref{lem:cone-semigroup}\eqref{it:semigroup-holomorphic}.

\emph{Step~4: the integral formula and the last bound.}
Let $d>0$, let $\alpha\in L^2(\R^2;\C)$ have compact support in $\{|y^0|<d\}$, and let $\psi,\psi'\in\mathcal{H}_1$.
By compactness there is $d'<d$ with $\supp\alpha\subset\{|y^0|\le d'\}$,
and $\alpha$ is integrable, being square integrable with compact support.
For $0<\delta<d-d'$ and $y\in\supp\alpha$, \eqref{eq:rotation-plate-shift} with $r=d-\delta>d'\ge|y^0|$ gives
$\mathcal{B}^{\delta}_\lambda\left(d\,n_0+y,d\,n_0-y\right)=\mathcal{B}^{\mathrm{c}}_d\left((\tau_y\cdot j_\delta)\otimes\lambda\right)$,
and the integration against $\alpha$ of Step~1 yields
\[
\int\mathrm{d}y\,\alpha(y)\left\langle\psi,\mathcal{B}^{\delta}_\lambda\left(d\,n_0+y,d\,n_0-y\right)\psi'\right\rangle
=\left\langle\psi,\mathcal{B}^{\mathrm{c}}_d\left((\alpha*j_\delta)\otimes\lambda\right)\psi'\right\rangle .
\]
As $\delta\to0$, $\alpha*j_\delta\to\alpha$ in $L^2(\R^2)$, and these functions vanish outside $\{|y^0|<d\}$;
so the right side tends to $\left\langle\psi,\mathcal{B}^{\mathrm{c}}_d(\alpha\otimes\lambda)\,\psi'\right\rangle$.
On the left, the integrand tends to $\alpha(y)\left\langle\psi,\mathcal{B}_\lambda\left(d\,n_0+y,d\,n_0-y\right)\psi'\right\rangle$ at every $y\in\supp\alpha$, by Step~2.
Moreover $\min\left(d_{\mathcal{W}}(d\,n_0+y),d_{\mathcal{W}}(d\,n_0-y)\right)=d-|y^0|\ge d-d'$ on $\supp\alpha$,
so that, for $0<\delta\le(d-d')/2$,~\eqref{eq:rotation-plate-estimate} bounds the modulus of the integrand by
\[
\frac{k_2\cdot\|\chi\|_2\cdot|\alpha(y)|}{d-d'-\delta}\cdot\|\psi\|\cdot\|\psi'\|
\le\frac{2k_2\cdot\|\chi\|_2\cdot|\alpha(y)|}{d-d'}\cdot\|\psi\|\cdot\|\psi'\|,
\]
an integrable function of $y$ that does not depend on $\delta$.
Dominated convergence therefore gives
\[
\int\mathrm{d}y\,\alpha(y)\left\langle\psi,\mathcal{B}_\lambda\left(d\,n_0+y,d\,n_0-y\right)\psi'\right\rangle
=\left\langle\psi,\mathcal{B}^{\mathrm{c}}_d(\alpha\otimes\lambda)\,\psi'\right\rangle .
\]
By Step~3 the map $y\mapsto\mathcal{B}_\lambda\left(d\,n_0+y,d\,n_0-y\right)$ is continuous in the operator norm on $\{|y^0|<d\}$,
and by~\eqref{it:point-plate-bound} its norm is at most $k_2\cdot\|\chi\|_2/(d-d')$ on $\supp\alpha$.
Hence $y\mapsto\alpha(y)\,\mathcal{B}_\lambda\left(d\,n_0+y,d\,n_0-y\right)$ is Bochner integrable in the operator norm,
and, by the last identity, its integral has the same matrix elements as $\mathcal{B}^{\mathrm{c}}_d(\alpha\otimes\lambda)$.
This proves~\eqref{it:point-plate-integral}.
Together with~\eqref{it:point-plate-bound} it gives
\[
\begin{aligned}
\left\|\mathcal{B}^{\mathrm{c}}_d(\alpha\otimes\lambda)\right\|
&\le\int\mathrm{d}y\,|\alpha(y)|\cdot\left\|\mathcal{B}_\lambda\left(d\,n_0+y,d\,n_0-y\right)\right\|\\
&\le k_2\cdot\|\chi\|_2\cdot\int\mathrm{d}y\,\frac{|\alpha(y)|}{d-|y^0|}
\le\frac{k_2\cdot\|\chi\|_2\cdot\|\alpha\|_{L^1(\R^2)}}{d-d'} .
\end{aligned}
\]
\end{proof}

\begin{remark}
Let $\alpha$, $d$ and $d'$ be as in the last assertion of Proposition~\ref{prop:point-plates}, and let $\Phi,\Phi'\in L^{\infty-}_{+}$.
Proposition~\ref{prop:insertion-bound}(b) expresses the matrix elements of $\mathcal{B}^{\mathrm{c}}_d(\alpha\otimes\lambda)$ through the law,
and the last assertion becomes
\[
\left|\mu\left(\Theta_1(U_d\Phi')\cdot\mathbb{B}_\lambda(\alpha)\cdot U_d\Phi\right)\right|
\le\frac{k_2\cdot\|\chi\|_2\cdot\|\alpha\|_{L^1(\R^2)}}{d-d'}\cdot\left\|[\Phi']_1\right\|\cdot\left\|[\Phi]_1\right\| .
\]
This is the plate bound with the exponent one:
the number $d-d'$, which does not exceed the distance from the support of $\alpha$ to the boundary $\{|y^0|=d\}$ of the strip $\{|y^0|<d\}$,
enters through its first negative power.
The bound holds at every scale $d>0$,
and its constant depends neither on $\Phi$ nor on $\Phi'$,
in particular not on the number of fields to the left and to the right of the plate.
It is the plate bound of the argument for four-dimensional Yang--Mills theory, here with the exponent one.
\end{remark}

\subsubsection{Kernels and plate vectors}\label{sec:plate-vectors}

In a product of two point plates, the semigroup $\mathsf{T}$ can be shifted from one factor to the other.

\begin{lemma}[Products of point plates]\label{lem:inner-arguments}
Let $\mu$ be a law on $\mathcal{S}'$ that satisfies \textup{(G)}, $(T_{123})$, $(R_1)$, $(RP_1)$ and the cone condition $(C_{12})$,
and let $\mathcal{B}_\lambda(\xi,\xi')$ denote its point plates (Proposition~\ref{prop:point-plates}).
Let $\lambda$ and $\lambda'$ be labels, let $\xi,\xi'\in\mathcal{W}$,
and let $\kappa_1,\kappa_2,\kappa_1',\kappa_2'\in\mathcal{W}$ satisfy $\kappa_1+\kappa_2=\kappa_1'+\kappa_2'$.
Then
\begin{equation}\label{eq:rotation-inner-arguments}
\mathcal{B}_\lambda(\xi,\kappa_1)\,\mathcal{B}_{\lambda'}(\kappa_2,\xi')=\mathcal{B}_\lambda(\xi,\kappa_1')\,\mathcal{B}_{\lambda'}(\kappa_2',\xi').
\end{equation}
\end{lemma}

\begin{proof}
Fix $\varepsilon>0$ with $\varepsilon\leq\min\{d_{\mathcal{W}}(\kappa_1),d_{\mathcal{W}}(\kappa_2),d_{\mathcal{W}}(\kappa_1'),d_{\mathcal{W}}(\kappa_2')\}$.
Since $d_{\mathcal{W}}(\kappa-\varepsilon n_0)=d_{\mathcal{W}}(\kappa)-\varepsilon$ for every $\kappa\in\C^2$,
the four points $\kappa_1-\varepsilon n_0$, $\kappa_2-\varepsilon n_0$, $\kappa_1'-\varepsilon n_0$, $\kappa_2'-\varepsilon n_0$ lie in $\overline{\mathcal{W}}$.
By Proposition~\ref{prop:point-plates}\eqref{it:point-plate-shift},
\[
\begin{aligned}
\mathcal{B}_\lambda(\xi,\kappa_1)&=\mathcal{B}_\lambda(\xi,\varepsilon n_0)\,\mathsf{T}(\kappa_1-\varepsilon n_0),\\
\mathcal{B}_{\lambda'}(\kappa_2,\xi')&=\mathsf{T}(\kappa_2-\varepsilon n_0)\,\mathcal{B}_{\lambda'}(\varepsilon n_0,\xi'),
\end{aligned}
\]
and the semigroup law of Lemma~\ref{lem:cone-semigroup}\eqref{it:semigroup-contraction} turns the left side of \eqref{eq:rotation-inner-arguments} into
\[
\mathcal{B}_\lambda(\xi,\varepsilon n_0)\,\mathsf{T}(\kappa_1+\kappa_2-2\varepsilon n_0)\,\mathcal{B}_{\lambda'}(\varepsilon n_0,\xi').
\]
The same two steps turn the right side into
$\mathcal{B}_\lambda(\xi,\varepsilon n_0)\,\mathsf{T}(\kappa_1'+\kappa_2'-2\varepsilon n_0)\,\mathcal{B}_{\lambda'}(\varepsilon n_0,\xi')$,
which is the same operator.
\end{proof}

\begin{definition}[Kernels and plate vectors]\label{def:plate-kernels}
Let $\mu$ be a law on $\mathcal{S}'$ that satisfies \textup{(G)}, $(T_{123})$, $(R_1)$, $(RP_1)$ and $(C_{12})$,
with the point plates $\mathcal{B}_\lambda(\xi,\xi')$ of Proposition~\ref{prop:point-plates},
and let $\kappa\in\mathcal{W}$.

Let $N\geq1$, let $\lambda=(\lambda_1,\dots,\lambda_N)$ be a tuple of labels, and let $\xi=(\xi_1,\dots,\xi_{N-1})\in\mathcal{W}^{N-1}$.
The \emph{kernel} $G_N(\xi;\lambda)=G_N(\xi_1,\dots,\xi_{N-1};\lambda_1,\dots,\lambda_N)$ is the number
\[
G_N(\xi;\lambda):=\left\langle\Omega_1,\,\mathcal{B}_{\lambda_1}(\kappa,\xi_1/2)\,\mathcal{B}_{\lambda_2}(\xi_1/2,\xi_2/2)\cdots\mathcal{B}_{\lambda_N}(\xi_{N-1}/2,\kappa)\,\Omega_1\right\rangle;
\]
for $N=1$ there is no variable $\xi$, and $G_1(\lambda_1)=\langle\Omega_1,\mathcal{B}_{\lambda_1}(\kappa,\kappa)\,\Omega_1\rangle$.

For $n\geq0$ let
\[
\mathcal{W}^{(n)}:=\left\{(x_1,\dots,x_n)\in(\C^2)^n:\ x_1\in\mathcal{W},\ x_{a+1}-x_a\in\mathcal{W}\ \text{for}\ 1\leq a<n\right\},
\]
so that $\mathcal{W}^{(0)}$ is the one-point set $(\C^2)^0$.
For $n\geq1$, a tuple of labels $\lambda=(\lambda_1,\dots,\lambda_n)$ and $x=(x_1,\dots,x_n)\in\mathcal{W}^{(n)}$,
the \emph{plate vector} $\Psi(x;\lambda)=\Psi(x;\lambda_1,\dots,\lambda_n)\in\mathcal{H}_1$ is
\[
\begin{aligned}
\Psi(x;\lambda):={}&\mathcal{B}_{\lambda_1}(x_1,(x_2-x_1)/2)\,\mathcal{B}_{\lambda_2}((x_2-x_1)/2,(x_3-x_2)/2)\\
&\cdots\mathcal{B}_{\lambda_n}((x_n-x_{n-1})/2,\kappa)\,\Omega_1,
\end{aligned}
\]
which is $\mathcal{B}_{\lambda_1}(x_1,\kappa)\,\Omega_1$ for $n=1$;
for $n=0$ the plate vector is $\Psi:=\Omega_1$.

Functions $\alpha_1,\dots,\alpha_N$ on $\R^2$ have \emph{time-ordered supports} if their supports are compact and
\[
\sup\left\{y^0:\ y\in\supp\alpha_a\right\}<\inf\left\{y^0:\ y\in\supp\alpha_{a+1}\right\}\qquad(1\leq a<N).
\]

All arguments of point plates in these formulas lie in $\mathcal{W}$, since $\mathcal{W}$ is a convex cone.
A real point $y$ of $\C^2$ lies in $\mathcal{W}$ if and only if $y^0>0$, because $d_{\mathcal{W}}(y)=y^0$;
hence the real points of $\mathcal{W}^{(n)}$ are the configurations $(y_1,\dots,y_n)$ of points of $\mathbb{H}$ with $y_1^0<y_2^0<\dots<y_n^0$.

The values of $G_N$ and $\Psi$ do not depend on $\kappa$.
More precisely, the two outer arguments of the product that defines $G_N$ may be replaced by any two points of $\mathcal{W}$, chosen independently of each other,
and the last argument of the product that defines $\Psi$ by any point of $\mathcal{W}$, without changing the values.
Indeed, let $\kappa'\in\mathcal{W}$ and $0<\varepsilon\leq d_{\mathcal{W}}(\kappa')$.
Then $\kappa'-\varepsilon n_0\in\overline{\mathcal{W}}$,
and Proposition~\ref{prop:point-plates}\eqref{it:point-plate-shift} gives, for every label $\lambda$ and every $\xi\in\mathcal{W}$,
\[
\begin{aligned}
\mathcal{B}_\lambda(\kappa',\xi)&=\mathsf{T}(\kappa'-\varepsilon n_0)\,\mathcal{B}_\lambda(\varepsilon n_0,\xi),\\
\mathcal{B}_\lambda(\xi,\kappa')&=\mathcal{B}_\lambda(\xi,\varepsilon n_0)\,\mathsf{T}(\kappa'-\varepsilon n_0).
\end{aligned}
\]
By Lemma~\ref{lem:cone-semigroup}\eqref{it:semigroup-contraction} and Lemma~\ref{lem:forward-tube}\eqref{it:tube-adjoint},
every $\kappa_0\in\overline{\mathcal{W}}$ satisfies $\mathsf{T}(\kappa_0)\,\Omega_1=\Omega_1$ and $\mathsf{T}(\kappa_0)^{*}\,\Omega_1=\mathsf{T}(\kappa_0^{*})\,\Omega_1=\Omega_1$.
Consequently
\[
\begin{aligned}
\langle\Omega_1,\mathcal{B}_\lambda(\kappa',\xi)\,\psi\rangle&=\langle\Omega_1,\mathcal{B}_\lambda(\varepsilon n_0,\xi)\,\psi\rangle\qquad(\psi\in\mathcal{H}_1),\\
\mathcal{B}_\lambda(\xi,\kappa')\,\Omega_1&=\mathcal{B}_\lambda(\xi,\varepsilon n_0)\,\Omega_1.
\end{aligned}
\]
Applied to the first and to the last factor of the products above,
with $\varepsilon$ not larger than $d_{\mathcal{W}}$ of any outer argument of two given choices,
these identities show that the two choices give the same values.
\end{definition}

\begin{proposition}[Bounds and Gram formula for the kernels]\label{prop:kernels}
Let $\mu$ be a law on $\mathcal{S}'$ that satisfies \textup{(G)}, $(T_{123})$, $(R_1)$, $(RP_1)$ and $(C_{12})$,
and let $G_N$ and $\Psi$ be its kernels and plate vectors (Definition~\ref{def:plate-kernels}).
Labels are written $\lambda_a=(\sigma_a,\chi_a)$.
\begin{enumerate}
\item\label{it:kernel-bound}
$G_N$ is holomorphic on $\mathcal{W}^{N-1}$ and real multilinear in the profiles $\chi_1,\dots,\chi_N$.
For $N\geq2$, $\xi\in\mathcal{W}^{N-1}$ and labels $\lambda=(\lambda_1,\dots,\lambda_N)$,
\begin{equation}\label{eq:rotation-kernel-bound}
\left|G_N(\xi;\lambda)\right|\leq\prod_{a=1}^{N}\left(2k_2\cdot\|\chi_a\|_2\right)\cdot\left(\min_{1\leq a\leq N-1}d_{\mathcal{W}}(\xi_a)\right)^{-N},
\end{equation}
so that every plate contributes one factor $2k_2\cdot\|\chi_a\|_2$ and one power of $\left(\min_a d_{\mathcal{W}}(\xi_a)\right)^{-1}$.
For $N=1$ one has $G_1(\lambda_1)=0$ for every label $\lambda_1$.
\item\label{it:kernel-density}
Let $\alpha_1,\dots,\alpha_N\in L^2(\R^2;\C)$ have time-ordered supports, let $\lambda=(\lambda_1,\dots,\lambda_N)$,
and write $Y=(y_1,\dots,y_N)$ for the points of $(\R^2)^N$.
Then the first of the two identities
\begin{equation}\label{eq:rotation-kernel-density}
\begin{aligned}
S_N(\alpha_1\otimes\lambda_1,\dots,\alpha_N\otimes\lambda_N)&=\int\mathrm{d}Y\,\prod_{a=1}^{N}\alpha_a(y_a)\cdot G_N(y_2-y_1,\dots,y_N-y_{N-1};\lambda),\\
\left[\prod_{a=1}^{N}\mathbb{B}_{\lambda_a}(\alpha_a)\right]_1&=\int\mathrm{d}Y\,\prod_{a=1}^{N}\alpha_a(y_a)\cdot\Psi(Y;\lambda)
\end{aligned}
\end{equation}
holds, and if the supports lie in $\mathbb{H}$, the second holds as well.
Both integrals extend over $\supp\alpha_1\times\dots\times\supp\alpha_N$;
there the arguments of $G_N$ lie in $\mathcal{W}^{N-1}$, and $Y$ lies in $\mathcal{W}^{(N)}$ if the supports lie in $\mathbb{H}$.
The second integral is a Bochner integral in $\mathcal{H}_1$.
\item\label{it:plate-vectors}
$\Psi$ is holomorphic on $\mathcal{W}^{(n)}$ with values in $\mathcal{H}_1$,
and $\mathsf{T}(\kappa')\,\Psi(x;\lambda)=\Psi(x_1+\kappa',\dots,x_n+\kappa';\lambda)$ for $x\in\mathcal{W}^{(n)}$ and $\kappa'\in\overline{\mathcal{W}}$.
Let $x\in\mathcal{W}^{(n)}$ and $z\in\mathcal{W}^{(n')}$,
and let $\lambda=(\lambda_1,\dots,\lambda_n)$ and $\lambda'=(\lambda_1',\dots,\lambda_{n'}')$ be tuples of labels.
If $n\geq1$ and $n'\geq1$, then
\begin{multline}\label{eq:rotation-gram}
\langle\Psi(x;\lambda),\Psi(z;\lambda')\rangle=s(\lambda)\cdot G_{n+n'}((x_n-x_{n-1})^{*},\dots,(x_2-x_1)^{*},\\
x_1^{*}+z_1,\,z_2-z_1,\dots,z_{n'}-z_{n'-1};\ \lambda_n,\dots,\lambda_1,\lambda_1',\dots,\lambda_{n'}').
\end{multline}
Moreover, $\langle\Omega_1,\Psi(z;\lambda')\rangle=G_{n'}(z_2-z_1,\dots,z_{n'}-z_{n'-1};\lambda')$ if $n'\geq1$;
\[
\langle\Psi(x;\lambda),\Omega_1\rangle=s(\lambda)\cdot G_n((x_n-x_{n-1})^{*},\dots,(x_2-x_1)^{*};\lambda_n,\dots,\lambda_1)
\]
if $n\geq1$, and for $n=1$ this number is $s(\lambda_1)\cdot G_1(\lambda_1)=0$;
finally $\langle\Omega_1,\Omega_1\rangle=1$.
\item\label{it:kernel-reality}
For $N\geq1$, $\xi\in\mathcal{W}^{N-1}$ and labels $\lambda=(\lambda_1,\dots,\lambda_N)$,
\begin{equation}\label{eq:rotation-kernel-conjugation}
\overline{G_N(\xi;\lambda)}=G_N(\overline{\xi};\bar\lambda),
\end{equation}
where $\overline{\xi}:=(\overline{\xi_1},\dots,\overline{\xi_{N-1}})$ and $\bar\lambda=(\bar\lambda_1,\dots,\bar\lambda_N)$.
\end{enumerate}
\end{proposition}

\begin{proof}
The vacuum vector $\Omega_1=[1]_1$ satisfies $\langle\Omega_1,[\Phi]_1\rangle=\mu(\Phi)$ for $\Phi\in L^2(\Sigma_1^+,\mu)$,
and in particular $\|\Omega_1\|=1$ \cite[Proposition~8.19\,(i)]{NoncompactQED} (the law is called $\rho$ there).

\emph{Holomorphy and the bound in \eqref{it:kernel-bound}.}
By Proposition~\ref{prop:point-plates}, each factor of the product that defines $G_N$ is holomorphic in the operator norm as a function of its two arguments,
which are affine functions of $\xi$, and it is real linear in the profile of its label.
Hence the product is holomorphic on $\mathcal{W}^{N-1}$ and real multilinear in $(\chi_1,\dots,\chi_N)$,
and so is its matrix element between $\Omega_1$ and $\Omega_1$.
Now let $N\geq2$ and $\xi\in\mathcal{W}^{N-1}$, put $r:=\min_a d_{\mathcal{W}}(\xi_a)>0$, and choose $\kappa:=tn_0$ with $t\geq r/2$.
Every argument of a point plate in the product is then either $tn_0$, with $d_{\mathcal{W}}(tn_0)=t$,
or one of the points $\xi_a/2$, with $d_{\mathcal{W}}(\xi_a/2)=d_{\mathcal{W}}(\xi_a)/2$;
in both cases $d_{\mathcal{W}}$ is at least $r/2$.
By Proposition~\ref{prop:point-plates}\eqref{it:point-plate-bound}, the factor with the label $\lambda_a$ has a norm at most $2k_2\cdot\|\chi_a\|_2\cdot r^{-1}$,
and \eqref{eq:rotation-kernel-bound} follows from $\|\Omega_1\|=1$.

\emph{The identities \eqref{eq:rotation-kernel-density} for supports in $\mathbb{H}$.}
Let the supports of $\alpha_1,\dots,\alpha_N$ be time-ordered and contained in $\mathbb{H}$.
If $\alpha_a=0$ for some $a$, then $\mathbb{B}_{\lambda_a}(\alpha_a)=\mathbb{B}(0)=0$, so that both left sides vanish,
and so do both integrals on the right, whose integrands contain the factor $\alpha_a(y_a)$.
Let therefore $\alpha_a\neq0$ for every $a$.
The supports are then non-empty compact subsets of $\mathbb{H}$,
and the time component of each point of $\supp\alpha_a$ is smaller than that of each point of $\supp\alpha_{a+1}$ ($1\leq a<N$).
Hence there are numbers $c_a\in\R$ and $d_a>0$ ($1\leq a\leq N$)
with $\supp\alpha_a\subset\{y\in\R^2:\ |y^0-c_a|<d_a\}$ and
\[
0\leq c_1-d_1,\qquad c_a+d_a\leq c_{a+1}-d_{a+1}\quad(1\leq a<N).
\]
It suffices to take for $c_1-d_1$ a non-negative number below the time components of all points of $\supp\alpha_1$,
for the common value of $c_a+d_a$ and $c_{a+1}-d_{a+1}$ a number above the time components on $\supp\alpha_a$ and below those on $\supp\alpha_{a+1}$ ($1\leq a<N$),
and for $c_N+d_N$ a number above the time components on $\supp\alpha_N$.
The plate $f_a:=\alpha_a\otimes\lambda_a\in L^2(\R^3;\C^3)$ vanishes outside the slab $\{x:\ |x_1-c_a|<d_a\}$.
The claim is that, for $1\leq a_0\leq N$,
\begin{multline}\label{eq:rotation-ordered-product}
\left[\prod_{a=a_0}^{N}\mathbb{B}(f_a)\right]_1=\exp\left(-(c_{a_0}-d_{a_0})\mathsf{H}_1\right)\mathcal{B}^{\mathrm{c}}_{d_{a_0}}(\tau_{-c_{a_0}}\cdot f_{a_0})\\
\exp\left(-(c_{a_0+1}-d_{a_0+1}-c_{a_0}-d_{a_0})\mathsf{H}_1\right)\mathcal{B}^{\mathrm{c}}_{d_{a_0+1}}(\tau_{-c_{a_0+1}}\cdot f_{a_0+1})\\
\cdots\exp\left(-(c_N-d_N-c_{N-1}-d_{N-1})\mathsf{H}_1\right)\mathcal{B}^{\mathrm{c}}_{d_N}(\tau_{-c_N}\cdot f_N)\,\Omega_1;
\end{multline}
the right side has one exponential and one operator $\mathcal{B}^{\mathrm{c}}_{d_a}(\tau_{-c_a}\cdot f_a)$ for each $a=a_0,\dots,N$,
and for $a_0=N$ it is $\exp\left(-(c_N-d_N)\mathsf{H}_1\right)\mathcal{B}^{\mathrm{c}}_{d_N}(\tau_{-c_N}\cdot f_N)\,\Omega_1$.
These operators are defined in Proposition~\ref{prop:insertion-bound}\eqref{it:insertion-symmetric},
since $\tau_{-c_a}\cdot f_a$ vanishes outside $\{|x_1|<d_a\}$ and so lies in $L^2_{[-d_a,d_a]}$.

The claim is proved by downward induction on $a_0$,
simultaneously for all fields $f_{a_0},\dots,f_N\in L^2(\R^3;\C^3)$ and all numbers $c_a\in\R$, $d_a>0$
such that $f_a$ vanishes outside $\{|x_1-c_a|<d_a\}$ ($a_0\leq a\leq N$), $c_{a_0}-d_{a_0}\geq0$, and $c_a+d_a\leq c_{a+1}-d_{a+1}$ ($a_0\leq a<N$).
For such data all $f_a$ vanish on $\{x_1<0\}$;
so every product of the fields $\mathbb{B}(f_a)$ lies in $L^{\infty-}_{+}$
(Lemma~\ref{lem:domination-consequences}\eqref{it:field-extension} and Definition~\ref{def:time-slabs}) and has a class in $\mathcal{H}_1$.
Let $a_0\leq N$, and if $a_0<N$, assume the claim for $a_0+1$ and all data of this kind.
Put $l:=c_{a_0}+d_{a_0}$, which is positive, and
\[
\Phi'':=\prod_{a=a_0+1}^{N}\mathbb{B}(f_a),\qquad\Phi':=\prod_{a=a_0+1}^{N}\mathbb{B}(\tau_{-l}\cdot f_a),
\]
both products being $1$ if $a_0=N$.
For $a>a_0$ the field $f_a$ vanishes on $\{x_1<c_a-d_a\}$, and this set contains $\{x_1<l\}$;
so $\tau_{-l}\cdot f_a$ vanishes on $\{x_1<0\}$, and $\Phi'\in L^{\infty-}_{+}$.
Moreover $U_l\Phi'=\Phi''$, since $U_l$ is multiplicative and $U_l\mathbb{B}(b)=\mathbb{B}(\tau_l\cdot b)$ for $b\in L^2(\R^3;\C^3)$
(Lemma~\ref{lem:domination-consequences}\eqref{it:field-extension}, the law being invariant under translations by $(T_{123})$).
As $c_{a_0}-d_{a_0}\geq0$, the field $f_{a_0}$ lies in $L^2_{[0,l]}$, and Proposition~\ref{prop:insertion-bound} gives
\[
\left[\mathbb{B}(f_{a_0})\,\Phi''\right]_1=\left[\mathbb{B}(f_{a_0})\,U_l\Phi'\right]_1=\mathcal{B}_l(f_{a_0})\,[\Phi']_1.
\]
Item~\eqref{it:insertion-localisation} of that proposition, applied to the slab $\{c_{a_0}-d_{a_0}\leq x_1\leq l\}$ outside of which $f_{a_0}$ vanishes,
gives the first of the following equalities;
the second is the definition of $\mathcal{B}^{\mathrm{c}}_{d_{a_0}}$ in Proposition~\ref{prop:insertion-bound}\eqref{it:insertion-symmetric},
since $\tau_{-(c_{a_0}-d_{a_0})}\cdot f_{a_0}=\tau_{d_{a_0}}\cdot\left(\tau_{-c_{a_0}}\cdot f_{a_0}\right)$:
\[
\begin{aligned}
\mathcal{B}_l(f_{a_0})&=\exp\left(-(c_{a_0}-d_{a_0})\mathsf{H}_1\right)\mathcal{B}_{2d_{a_0}}\left(\tau_{-(c_{a_0}-d_{a_0})}\cdot f_{a_0}\right)\\
&=\exp\left(-(c_{a_0}-d_{a_0})\mathsf{H}_1\right)\mathcal{B}^{\mathrm{c}}_{d_{a_0}}(\tau_{-c_{a_0}}\cdot f_{a_0}).
\end{aligned}
\]
For $a_0=N$ one has $[\Phi']_1=[1]_1=\Omega_1$, and \eqref{eq:rotation-ordered-product} holds.
For $a_0<N$, the fields $\tau_{-l}\cdot f_a$ with the numbers $c_a-l$ and $d_a$ ($a_0<a\leq N$) are data of the kind above for $a_0+1$,
because $(c_{a_0+1}-l)-d_{a_0+1}\geq0$.
The induction hypothesis expresses $[\Phi']_1$ through these data:
the first exponent on its right side is $-\left((c_{a_0+1}-l)-d_{a_0+1}\right)\mathsf{H}_1=-(c_{a_0+1}-d_{a_0+1}-c_{a_0}-d_{a_0})\mathsf{H}_1$,
the other exponents do not change under the common shift of the numbers $c_a$,
and the remaining factors are $\mathcal{B}^{\mathrm{c}}_{d_a}\left(\tau_{-(c_a-l)}\cdot\tau_{-l}\cdot f_a\right)=\mathcal{B}^{\mathrm{c}}_{d_a}(\tau_{-c_a}\cdot f_a)$, $a_0<a\leq N$.
Inserting this into the last two displays gives \eqref{eq:rotation-ordered-product} at $a_0$.

Let now $a_0=1$, and let $1\leq a\leq N$.
The function $\alpha_a(\cdot+c_an_0)$ has compact support in $\{|y^0|<d_a\}$, and $\tau_{-c_a}\cdot f_a=\alpha_a(\cdot+c_an_0)\otimes\lambda_a$.
Proposition~\ref{prop:point-plates}\eqref{it:point-plate-integral}, followed by the substitution $y=y_a-c_an_0$, gives
\[
\mathcal{B}^{\mathrm{c}}_{d_a}(\tau_{-c_a}\cdot f_a)=\int\mathrm{d}y_a\,\alpha_a(y_a)\,\mathcal{B}_{\lambda_a}((d_a-c_a)n_0+y_a,\,(d_a+c_a)n_0-y_a),
\]
a Bochner integral in the operator norm over $\supp\alpha_a$.
For $y_a$ in the open strip $\{|y_a^0-c_a|<d_a\}$, which contains $\supp\alpha_a$,
the two arguments of the point plate in this integral are real points with the positive time components $d_a-c_a+y_a^0$ and $d_a+c_a-y_a^0$, hence points of $\mathcal{W}$.
On this strip the point plate is therefore a continuous function of $y_a$ in the operator norm (Proposition~\ref{prop:point-plates}),
and its norm is bounded on the compact subset $\supp\alpha_a$;
the scalar factor $\alpha_a$, square integrable with compact support, is integrable.
When these $N$ integrals are inserted into \eqref{eq:rotation-ordered-product},
the bounded operators $\exp(-t\mathsf{H}_1)$ can be taken inside,
and the product of the integrals, applied to $\Omega_1$, becomes one Bochner integral in $\mathcal{H}_1$ over $\supp\alpha_1\times\dots\times\supp\alpha_N$;
its integrand is $\prod_a\alpha_a(y_a)$ times the product of the exponentials and of the point plates, applied to $\Omega_1$.
For $t\geq0$,
\[
\exp(-t\mathsf{H}_1)=\int\mathrm{d}\mathfrak{E}_1(\omega,p_\perp)\,\exp(-t\omega)=\mathsf{T}(tn_0)
\]
by Lemma~\ref{lem:joint-spectrum} and the definition of $\mathsf{T}$.
By Proposition~\ref{prop:point-plates}\eqref{it:point-plate-shift}, the factor $\mathsf{T}((c_1-d_1)n_0)$ in front of the first point plate is absorbed into its first argument,
which becomes $(c_1-d_1)n_0+(d_1-c_1)n_0+y_1=y_1$,
and, for $1\leq a<N$, the factor $\mathsf{T}((c_{a+1}-d_{a+1}-c_a-d_a)n_0)$ is absorbed into the second argument of the $a$-th point plate;
all these shifts lie in $\overline{\mathcal{W}}$.
At a point $Y$ of $\supp\alpha_1\times\dots\times\supp\alpha_N$,
the inner arguments between the $a$-th and the $(a+1)$-st point plate are now $(c_{a+1}-d_{a+1})n_0-y_a$ and $(d_{a+1}-c_{a+1})n_0+y_{a+1}$.
They are real points of $\mathcal{W}$, with the time components $c_{a+1}-d_{a+1}-y_a^0>0$ and $y_{a+1}^0-c_{a+1}+d_{a+1}>0$, and their sum is
\[
\begin{aligned}
&\left((d_a+c_a)n_0-y_a\right)+(c_{a+1}-d_{a+1}-c_a-d_a)n_0+\left((d_{a+1}-c_{a+1})n_0+y_{a+1}\right)\\
&\qquad=y_{a+1}-y_a,
\end{aligned}
\]
a real point of $\mathcal{W}$ as well.
By \eqref{eq:rotation-inner-arguments}, both inner arguments may be replaced by $(y_{a+1}-y_a)/2$.
The first argument $y_1$ lies in $\mathcal{W}$,
and the last argument $(d_N+c_N)n_0-y_N$ is a point of $\mathcal{W}$ that may be replaced by $\kappa$ (Definition~\ref{def:plate-kernels}).
So the integrand equals $\prod_a\alpha_a(y_a)\cdot\Psi(Y;\lambda)$,
and \eqref{eq:rotation-ordered-product} with $a_0=1$ becomes the second identity of \eqref{eq:rotation-kernel-density}.

The first identity follows by taking the scalar product with $\Omega_1$, which commutes with the Bochner integral.
On the left, $\langle\Omega_1,[\prod_a\mathbb{B}(f_a)]_1\rangle=\mu\left(\prod_a\mathbb{B}(f_a)\right)=S_N(f_1,\dots,f_N)$.
On the right, $\langle\Omega_1,\Psi(Y;\lambda)\rangle=G_N(y_2-y_1,\dots,y_N-y_{N-1};\lambda)$ by Definition~\ref{def:plate-kernels},
the outer arguments of the product being $y_1$ and $\kappa$.

\emph{Arbitrary time-ordered supports.}
Let $t\in\R$, and replace each $\alpha_a$ by $\alpha_a(\cdot-tn_0)$.
The left side of the first identity of \eqref{eq:rotation-kernel-density} does not change:
$\alpha_a(\cdot-tn_0)\otimes\lambda_a=\tau_{te_1}\cdot(\alpha_a\otimes\lambda_a)$,
$\mathbb{B}(\tau_{te_1}\cdot b)=U(te_1)\mathbb{B}(b)$ by Lemma~\ref{lem:domination-consequences}\eqref{it:field-extension},
$U(te_1)$ is multiplicative, and $\mu(U(te_1)\Phi)=\mu(\Phi)$ by $(T_{123})$.
Nor does the right side, by the substitution $y_a\mapsto y_a+tn_0$ ($1\leq a\leq N$), which leaves the differences $y_{a+1}-y_a$ unchanged.
For large $t$ the translated supports lie in $\mathbb{H}$, and the first identity holds for all time-ordered supports.

\emph{The kernel $G_1$.}
For $N=1$ and any $\alpha_1\in L^2(\R^2;\C)$ with compact support, the first identity of \eqref{eq:rotation-kernel-density} reads
\[
S_1(\alpha_1\otimes\lambda_1)=G_1(\lambda_1)\cdot\int\mathrm{d}y_1\,\alpha_1(y_1).
\]
Its left side vanishes, because $S_1=0$ (Lemma~\ref{lem:domination-consequences}\eqref{it:moment-forms});
a function $\alpha_1$ with $\int\mathrm{d}y_1\,\alpha_1(y_1)\neq0$ gives $G_1(\lambda_1)=0$.
This completes~\eqref{it:kernel-bound}.

\emph{Plate vectors.}
Let $n\geq1$.
For $x$ in the open set $\mathcal{W}^{(n)}$, the arguments of the point plates in the definition of $\Psi(x;\lambda)$ are affine functions of $x$ with values in $\mathcal{W}$;
by Proposition~\ref{prop:point-plates}, their product is holomorphic on $\mathcal{W}^{(n)}$ in the operator norm, and so is $x\mapsto\Psi(x;\lambda)$.
For $\kappa'\in\overline{\mathcal{W}}$, the configuration $(x_1+\kappa',\dots,x_n+\kappa')$ has the same differences as $x$,
and $d_{\mathcal{W}}(x_1+\kappa')\geq d_{\mathcal{W}}(x_1)+d_{\mathcal{W}}(\kappa')>0$; so it lies in $\mathcal{W}^{(n)}$.
Proposition~\ref{prop:point-plates}\eqref{it:point-plate-shift} gives $\mathsf{T}(\kappa')\,\mathcal{B}_{\lambda_1}(x_1,\kappa_0)=\mathcal{B}_{\lambda_1}(x_1+\kappa',\kappa_0)$ for $\kappa_0\in\mathcal{W}$,
and the other factors depend on $x$ only through the differences $x_{a+1}-x_a$;
this proves $\mathsf{T}(\kappa')\,\Psi(x;\lambda)=\Psi(x_1+\kappa',\dots,x_n+\kappa';\lambda)$.
For $n=0$ the identity reads $\mathsf{T}(\kappa')\,\Omega_1=\Omega_1$ (Lemma~\ref{lem:cone-semigroup}\eqref{it:semigroup-contraction}).

For the scalar products, let $n\geq1$, write $\xi_a:=x_{a+1}-x_a$ ($1\leq a<n$),
and let $A$ be the product of point plates in the definition of $\Psi(x;\lambda)$, so that $\Psi(x;\lambda)=A\,\Omega_1$.
Proposition~\ref{prop:point-plates}\eqref{it:point-plate-shift} gives $\mathcal{B}_{\lambda_a}(\kappa_1,\kappa_2)^{*}=s(\sigma_a)\cdot\mathcal{B}_{\lambda_a}(\kappa_2^{*},\kappa_1^{*})$ for $\kappa_1,\kappa_2\in\mathcal{W}$; hence
\[
\begin{aligned}
A^{*}&=s(\lambda)\cdot\mathcal{B}_{\lambda_n}(\kappa^{*},\xi_{n-1}^{*}/2)\,\mathcal{B}_{\lambda_{n-1}}(\xi_{n-1}^{*}/2,\xi_{n-2}^{*}/2)\\
&\qquad\cdots\mathcal{B}_{\lambda_1}(\xi_1^{*}/2,x_1^{*}),
\end{aligned}
\]
and $A^{*}=s(\lambda_1)\cdot\mathcal{B}_{\lambda_1}(\kappa^{*},x_1^{*})$ if $n=1$;
the points $\kappa^{*}$, $\xi_a^{*}/2$ and $x_1^{*}$ lie in $\mathcal{W}$ by Lemma~\ref{lem:forward-tube}\eqref{it:tube-adjoint}.
If $n'\geq1$, let $A'$ be the product of point plates in the definition of $\Psi(z;\lambda')$,
so that $\langle\Psi(x;\lambda),\Psi(z;\lambda')\rangle=\langle\Omega_1,A^{*}A'\,\Omega_1\rangle$.
In $A^{*}A'$ the last factor of $A^{*}$, with the second argument $x_1^{*}$, is followed by the first factor of $A'$, with the first argument $z_1$;
these inner arguments add up to $x_1^{*}+z_1\in\mathcal{W}$, and by \eqref{eq:rotation-inner-arguments} both may be replaced by $(x_1^{*}+z_1)/2$.
After this replacement, $A^{*}A'$ is $s(\lambda)$ times the product that defines $G_{n+n'}$ at the arguments and with the labels of \eqref{eq:rotation-gram},
except that its outer arguments are $\kappa^{*}$ and $\kappa$;
by Definition~\ref{def:plate-kernels} this does not change the matrix element, and \eqref{eq:rotation-gram} follows.
For $n'=0$ the same expression of $A^{*}$ gives
\[
\langle\Psi(x;\lambda),\Omega_1\rangle=\langle\Omega_1,A^{*}\,\Omega_1\rangle=s(\lambda)\cdot G_n(\xi_{n-1}^{*},\dots,\xi_1^{*};\lambda_n,\dots,\lambda_1),
\]
the outer arguments being $\kappa^{*}$ and $x_1^{*}$; for $n=1$ this is $s(\lambda_1)\cdot G_1(\lambda_1)=0$ by \eqref{it:kernel-bound}.
For $n=0$ and $n'\geq1$, the number $\langle\Omega_1,\Psi(z;\lambda')\rangle=\langle\Omega_1,A'\,\Omega_1\rangle$
is the matrix element that defines $G_{n'}(z_2-z_1,\dots,z_{n'}-z_{n'-1};\lambda')$, with the outer arguments $z_1$ and $\kappa$;
and $\langle\Omega_1,\Omega_1\rangle=\|\Omega_1\|^2=1$.

\emph{Conjugation.}
For $N=1$ both sides of \eqref{eq:rotation-kernel-conjugation} vanish by \eqref{it:kernel-bound}.
Let $N\geq2$, write $\xi(Y):=(y_2-y_1,\dots,y_N-y_{N-1})$ for $Y\in(\R^2)^N$,
and let $\alpha_1,\dots,\alpha_N\in L^2(\R^2;\C)$ have time-ordered supports;
then $\overline{\alpha_1},\dots,\overline{\alpha_N}$ lie in $L^2(\R^2;\C)$ and have time-ordered supports as well.
By Lemma~\ref{lem:domination-consequences}\eqref{it:moment-forms} and Lemma~\ref{lem:plate-identities}\eqref{it:plate-conjugation},
\[
\overline{S_N(\alpha_1\otimes\lambda_1,\dots,\alpha_N\otimes\lambda_N)}=S_N(\overline{\alpha_1}\otimes\bar\lambda_1,\dots,\overline{\alpha_N}\otimes\bar\lambda_N),
\]
and the first identity of \eqref{eq:rotation-kernel-density}, applied to both sides, turns this into
\[
\int\mathrm{d}Y\,\prod_{a=1}^{N}\overline{\alpha_a(y_a)}\cdot\left(\overline{G_N(\xi(Y);\lambda)}-G_N(\xi(Y);\bar\lambda)\right)=0.
\]
The set $U:=\{Y\in(\R^2)^N:\ y_1^0<y_2^0<\dots<y_N^0\}$ is open, and the two functions of $Y$ in the bracket are continuous on it.
Let $\eta_1,\dots,\eta_N$ be bounded measurable functions on $\R^2$ with compact supports and $\supp\eta_1\times\dots\times\supp\eta_N\subset U$;
they lie in $L^2(\R^2;\C)$.
If none of them vanishes, their supports are time-ordered,
because a point $Y$ of the product whose entries $y_a$ and $y_{a+1}$ maximise and minimise the time component on $\supp\eta_a$ and on $\supp\eta_{a+1}$ lies in $U$;
so the last identity holds for $\alpha_a:=\overline{\eta_a}$, that is, with $\eta_a$ in place of $\overline{\alpha_a}$,
and it holds trivially if one of them vanishes.
By Lemma~\ref{lem:analysis-facts}\eqref{it:equal-densities}, $\overline{G_N(\xi(Y);\lambda)}=G_N(\xi(Y);\bar\lambda)$ for $Y\in U$,
that is, at every real point of $\mathcal{W}^{N-1}$.
By Lemma~\ref{lem:forward-tube}\eqref{it:tube-adjoint}, complex conjugation maps $\mathcal{W}$ onto itself;
hence $\xi\mapsto\overline{G_N(\overline{\xi};\lambda)}$ is a holomorphic function on $\mathcal{W}^{N-1}$, and it agrees with $G_N(\cdot;\bar\lambda)$ at the real points.
These form a non-empty open subset of $(\R^2)^{N-1}$, and $\mathcal{W}^{N-1}$ is convex, hence connected;
by Lemma~\ref{lem:analysis-facts}\eqref{it:identity-theorem} the two functions agree on all of $\mathcal{W}^{N-1}$.
At the point $\overline{\xi}$ this equality is \eqref{eq:rotation-kernel-conjugation}.
\end{proof}

\subsection{The two-point function}\label{sec:two-point-function}

\subsubsection{Spectral measures of plates}\label{sec:plate-spectra}

The Cauchy density
\[
P_1(v):=\frac{1}{\pi\left(1+v^2\right)}\qquad(v\in\R)
\]
satisfies
\[
\int\mathrm{d}v\,P_1(v)\cdot\exp(ip\cdot v)=\exp(-|p|)\quad(p\in\R),
\qquad
\|P_1\|_1=1,
\qquad
\|P_1\|_2^2=\frac{1}{2\pi}.
\]

\begin{lemma}[The spectral matrix of one field]\label{lem:spectral-matrix}
Let $j\in\{1,2,3\}$, and let $\mu$ be a law that satisfies \textup{(G)}, $(T_{123})$, $(R_j)$ and $(RP_j)$.
Let $\mathcal{H}_j$, $[\,\cdot\,]_j$, $\mathsf{H}_j$ and $V_j$ be the Hilbert space, the class map, the Hamiltonian and the transverse translations of its reconstruction in the direction $e_j$,
and let $\mathfrak{E}_j$ be its joint spectral measure (Lemma~\ref{lem:joint-spectrum}),
written in the variables $(\omega,p_\perp)\in[0,\infty)\times e_j^\perp$,
where $e_j^\perp$ is identified with $\R^2$ through $p_\perp=(p_i)_{i\neq j}$.
For $u\in L^1\cap L^2(\R^3;\C)$ with $u=0$ almost everywhere on $\{x_j<0\}$ put
\[
\tilde u(\omega,p_\perp):=\int\mathrm{d}x\,u(x)\cdot\exp\left(-x_j\cdot\omega+ip_\perp\cdot x_\perp\right),
\qquad x_\perp:=x-x_je_j ;
\]
it is a continuous function on $[0,\infty)\times e_j^\perp$, and $|\tilde u|\le\|u\|_1$.
There are complex Radon measures $\varrho_{ab}=\varrho^{(j)}_{ab}$, $a,b\in\{1,2,3\}$, on $[0,\infty)\times\R^2$,
that is, set functions on the bounded Borel sets whose restriction to every bounded Borel set is a complex measure,
with the following properties.
\begin{enumerate}
\item\label{it:spectral-matrix-positive}
One has $\varrho_{ab}=r_{ab}\cdot\varrho_0$, $a,b\in\{1,2,3\}$, for some positive Radon measure $\varrho_0$ on $[0,\infty)\times\R^2$
and some Borel functions $r_{ab}\colon[0,\infty)\times\R^2\to\C$
such that, at each point $(\omega,p_\perp)$, the $3\times3$ matrix with the entries $r_{ab}(\omega,p_\perp)$ is Hermitian and positive semidefinite.
It follows that for bounded Borel functions $c_a,c'_a\colon[0,\infty)\times\R^2\to\C$, $a\in\{1,2,3\}$,
the measures $\varrho_{cc'}:=\sum_{a,b}\overline{c_a}\cdot c'_b\cdot\varrho_{ab}$ satisfy $\varrho_{cc}\ge0$ and
\[
\int\mathrm{d}\left|\varrho_{cc'}\right|\,|G_1|\cdot|G_2|
\le\left(\int\mathrm{d}\varrho_{cc}\,|G_1|^2\right)^{\frac12}\cdot\left(\int\mathrm{d}\varrho_{c'c'}\,|G_2|^2\right)^{\frac12}
\]
for all Borel functions $G_1,G_2\colon[0,\infty)\times\R^2\to\C$.
\item\label{it:spectral-matrix-formula}
For all $u,u'$ as above and all $a,b\in\{1,2,3\}$ one has
$\int\mathrm{d}\varrho_{aa}\,|\tilde u|^2=\left\|[\mathbb{B}_a(u)]_j\right\|^2<\infty$ and
\begin{equation}\label{eq:rotation-spectral-matrix}
\left\langle[\mathbb{B}_a(u)]_j,[\mathbb{B}_b(u')]_j\right\rangle=\int\mathrm{d}\varrho_{ab}\,\overline{\tilde u}\cdot\tilde u' .
\end{equation}
\item\label{it:spectral-matrix-support}
One has $|\varrho_{ab}|(A)=0$ for every Borel set $A$
such that $\mathfrak{E}_j(A)\psi=0$ for all vectors $\psi$ of the one-field space $\mathcal{H}^1_j$ (Definition~\ref{def:closedness-cones}).
\end{enumerate}
\end{lemma}

\begin{proof}
On the support of $u$ one has $x_j\ge0$ almost everywhere, and $\omega\ge0$;
so the integrand that defines $\tilde u$ has modulus at most $|u(x)|$.
This gives $|\tilde u|\le\|u\|_1$ and, by dominated convergence, the continuity of $\tilde u$.
For a bounded Borel function $h\colon[0,\infty)\times\R^2\to\C$ let
\[
h(\mathsf{H}_j,\mathsf{P}):=\int\mathrm{d}\mathfrak{E}_j\,h
\]
be the bounded operator on $\mathcal{H}_j$ defined by the joint spectral measure of $\mathsf{H}_j$ and of the generators $\mathsf{P}_i$, $i\neq j$.
It depends linearly and multiplicatively on $h$, its adjoint is $\overline{h}(\mathsf{H}_j,\mathsf{P})$,
one has $1_A(\mathsf{H}_j,\mathsf{P})=\mathfrak{E}_j(A)$ for every Borel set $A$,
and $\left\langle\psi',h(\mathsf{H}_j,\mathsf{P})\psi\right\rangle=\int\mathrm{d}\left\langle\psi',\mathfrak{E}_j\psi\right\rangle\,h$ for all $\psi,\psi'\in\mathcal{H}_j$.

\emph{Step 1: convolutions.}
Let $u$ be as in the statement, and let $\alpha\in L^2(\R^3;\C)$ vanish almost everywhere on $\{x_j<0\}$.
For $x\in\R^3$ let $\tau_x\alpha$ be the translate $x'\mapsto\alpha(x'-x)$;
it vanishes almost everywhere on $\{x'_j<x_j\}$, and hence on $\{x'_j<0\}$ when $x_j\ge0$.
Translations act isometrically and strongly continuously on $L^2(\R^3;\C)$,
and $u$ is integrable and vanishes almost everywhere on $\{x_j<0\}$.
The convolution of $u$ and $\alpha$ is therefore the Bochner integral
\[
u*\alpha=\int_{\{x_j\ge0\}}\mathrm{d}x\,u(x)\cdot\tau_x\alpha
\]
of a function with values in the closed subspace of $L^2(\R^3;\C)$ formed by the functions that vanish almost everywhere on $\{x'_j<0\}$;
in particular, $u*\alpha$ lies in this subspace.
If a field $b\in L^2(\R^3;\C^3)$ vanishes almost everywhere on $\{x_j<0\}$,
then $\mathbb{B}(b)$ lies in $L^2(\Sigma_j^+,\mu)$ and $\mu\left(|\mathbb{B}(b)|^2\right)\le\|b\|^2$
(Lemma~\ref{lem:domination-consequences}\eqref{it:field-extension});
moreover the class map $[\,\cdot\,]_j$ is a linear contraction of $L^2(\Sigma_j^+,\mu)$ into $\mathcal{H}_j$
(\cite[Proposition~8.19\,(i)]{NoncompactQED}, where $L^2(\Sigma_j^+,\mu)$ is denoted $\mathcal{E}_+$).
Hence $b\mapsto[\mathbb{B}(b)]_j$ is a bounded linear map on these fields, with $\left\|[\mathbb{B}(b)]_j\right\|\le\|b\|$.
It commutes with the Bochner integral:
\[
[\mathbb{B}_a(u*\alpha)]_j=\int_{\{x_j\ge0\}}\mathrm{d}x\,u(x)\cdot[\mathbb{B}_a(\tau_x\alpha)]_j .
\]
The invariance $(T_{123})$ and Lemma~\ref{lem:domination-consequences}\eqref{it:field-extension}
give $\mathbb{B}_a(\tau_x\alpha)=U(x)\mathbb{B}_a(\alpha)$ for every $x\in\R^3$,
and for $x_j\ge0$ \cite[Proposition~8.19\,(ii), (iii)]{NoncompactQED}, in which $\exp(-t\cdot\mathsf{H}_j)$ is denoted $T(t)$, gives
\[
[U(x)\mathbb{B}_a(\alpha)]_j=\exp(-x_j\cdot\mathsf{H}_j)\,V_j(x_\perp)\,[\mathbb{B}_a(\alpha)]_j .
\]
In the present variables, Lemma~\ref{lem:joint-spectrum}(b) reads
\[
\exp(-t\cdot\mathsf{H}_j)\,V_j(y)=\int\mathrm{d}\mathfrak{E}_j(\omega,p_\perp)\,\exp\left(-t\cdot\omega+ip_\perp\cdot y\right)
\qquad\left(t\ge0,\ y\in e_j^\perp\right).
\]
Let $\psi:=[\mathbb{B}_a(\alpha)]_j$ and $\psi'\in\mathcal{H}_j$.
The complex measure $\left\langle\psi',\mathfrak{E}_j\psi\right\rangle$ is finite,
and on the domain of integration the integrand of the following double integral is bounded in modulus by $|u(x)|$;
Fubini's theorem therefore gives
\[
\begin{aligned}
\left\langle\psi',[\mathbb{B}_a(u*\alpha)]_j\right\rangle
&=\int_{\{x_j\ge0\}}\mathrm{d}x\,u(x)\int\mathrm{d}\left\langle\psi',\mathfrak{E}_j\psi\right\rangle(\omega,p_\perp)\,\exp\left(-x_j\cdot\omega+ip_\perp\cdot x_\perp\right)\\
&=\left\langle\psi',\tilde u(\mathsf{H}_j,\mathsf{P})\psi\right\rangle .
\end{aligned}
\]
As $\psi'$ is arbitrary,
\begin{equation}\label{eq:rotation-field-convolution}
[\mathbb{B}_a(u*\alpha)]_j=\tilde u(\mathsf{H}_j,\mathsf{P})\,[\mathbb{B}_a(\alpha)]_j\qquad(a=1,2,3).
\end{equation}

\emph{Step 2: a reference field.}
The function
\[
\alpha_0(x):=1_{(0,1)}(x_j)\cdot\prod_{i\neq j}P_1(x_i)
\]
lies in $L^1\cap L^2(\R^3;\C)$ and vanishes on $\{x_j<0\}$.
By the transform of $P_1$ recalled before the lemma,
\[
\tilde\alpha_0(\omega,p_\perp)=\frac{1-\exp(-\omega)}{\omega}\cdot\exp\left(-|p_\perp|_1\right),
\qquad |p_\perp|_1:=\sum_{i\neq j}|p_i|,
\]
where the first factor is read as $1$ at $\omega=0$;
thus $\tilde\alpha_0$ is continuous and strictly positive.
Put $\psi_a:=[\mathbb{B}_a(\alpha_0)]_j$ for $a\in\{1,2,3\}$.
Since $u*\alpha_0=\alpha_0*u$, the identity \eqref{eq:rotation-field-convolution},
applied once to the pair $(u,\alpha_0)$ and once to the pair $(\alpha_0,u)$, yields
\begin{equation}\label{eq:rotation-reference-field}
\tilde\alpha_0(\mathsf{H}_j,\mathsf{P})\,[\mathbb{B}_a(u)]_j=\tilde u(\mathsf{H}_j,\mathsf{P})\,\psi_a .
\end{equation}

\emph{Step 3: the measures.}
Let $m_{ab}:=\left\langle\psi_a,\mathfrak{E}_j(\cdot)\psi_b\right\rangle$ for $a,b\in\{1,2,3\}$;
these are finite complex Borel measures on $[0,\infty)\times\R^2$.
For $z\in\C^3$ and $\psi_z:=\sum_az_a\psi_a$ one has
$\sum_{a,b}\overline{z_a}\cdot z_b\cdot m_{ab}=\left\langle\psi_z,\mathfrak{E}_j(\cdot)\psi_z\right\rangle\ge0$.
Let $m_0:=m_{11}+m_{22}+m_{33}$.
For every Borel set $A$ the Cauchy--Schwarz inequality gives
\[
\left|m_{ab}(A)\right|=\left|\left\langle\mathfrak{E}_j(A)\psi_a,\mathfrak{E}_j(A)\psi_b\right\rangle\right|
\le\left(m_{aa}(A)\cdot m_{bb}(A)\right)^{\frac12}\le m_0(A);
\]
hence every $m_{ab}$ is absolutely continuous with respect to $m_0$,
and the Radon--Nikodym theorem provides Borel functions $r_{ab}$ with $m_{ab}=r_{ab}\cdot m_0$.
For each $z\in(\mathbb{Q}+i\mathbb{Q})^3$ the function $\sum_{a,b}\overline{z_a}\cdot z_b\cdot r_{ab}$ is the density of a positive measure,
so it is non-negative $m_0$-almost everywhere.
Outside an $m_0$-null Borel set these countably many inequalities hold simultaneously;
by continuity in $z$, the quadratic form $z\mapsto\sum_{a,b}\overline{z_a}\cdot z_b\cdot r_{ab}$ is then non-negative on all of $\C^3$,
and a complex matrix whose quadratic form takes only real values is Hermitian.
Once the $r_{ab}$ are replaced by $\delta_{ab}$ on that null set,
the matrix $\left(r_{ab}\right)_{a,b}$ is Hermitian and positive semidefinite at every point, and still $m_{ab}=r_{ab}\cdot m_0$.
Being continuous and strictly positive, $\tilde\alpha_0$ has a positive lower bound on every bounded set.
Therefore
\[
\varrho_0:=|\tilde\alpha_0|^{-2}\cdot m_0,
\qquad
\varrho_{ab}:=|\tilde\alpha_0|^{-2}\cdot m_{ab}=r_{ab}\cdot\varrho_0
\]
are a positive Radon measure and complex Radon measures on $[0,\infty)\times\R^2$.

\emph{Proof of \eqref{it:spectral-matrix-positive}.}
For bounded Borel functions $c_a,c'_a$ the measure $\varrho_{cc'}$ has the density $\sum_{a,b}\overline{c_a}\cdot c'_b\cdot r_{ab}$ with respect to $\varrho_0$.
For $c'=c$ this density is non-negative, so that $\varrho_{cc}\ge0$.
At every point, the Cauchy--Schwarz inequality for the positive semidefinite Hermitian form $\left(r_{ab}\right)_{a,b}$ gives
\[
\left|\sum_{a,b}\overline{c_a}\cdot c'_b\cdot r_{ab}\right|
\le\left(\sum_{a,b}\overline{c_a}\cdot c_b\cdot r_{ab}\right)^{\frac12}\cdot\left(\sum_{a,b}\overline{c'_a}\cdot c'_b\cdot r_{ab}\right)^{\frac12},
\]
and $|\varrho_{cc'}|=\left|\sum_{a,b}\overline{c_a}\cdot c'_b\cdot r_{ab}\right|\cdot\varrho_0$.
The asserted inequality now follows from the Cauchy--Schwarz inequality in $L^2(\varrho_0)$.

\emph{Proof of \eqref{it:spectral-matrix-support}.}
Let $A$ be a Borel set with $\mathfrak{E}_j(A)\psi=0$ for all $\psi\in\mathcal{H}^1_j$.
Every $\psi_b$ lies in $\mathcal{H}^1_j$, since the field $\alpha_0e_b$ vanishes on $\{x_j<0\}$.
For every Borel set $A'\subset A$ it follows that
$m_{ab}(A')=\left\langle\psi_a,\mathfrak{E}_j(A')\,\mathfrak{E}_j(A)\psi_b\right\rangle=0$.
So $|m_{ab}|(A)=0$, and $|\varrho_{ab}|(A)=0$ because $|\varrho_{ab}|=|\tilde\alpha_0|^{-2}\cdot|m_{ab}|$.

\emph{Step 4: proof of \eqref{it:spectral-matrix-formula}.}
For $n\ge1$ let $A_n:=\{\tilde\alpha_0>1/n\}$.
From $\tilde\alpha_0(\omega,p_\perp)\le\min\{1,1/\omega\}\cdot\exp\left(-|p_\perp|_1\right)$
one sees that $A_n$ is contained in $\{\omega<n,\ |p_\perp|_1<\log n\}$, hence bounded.
As $\tilde\alpha_0>0$ everywhere, the sets $A_n$ increase to $[0,\infty)\times\R^2$, and $\mathfrak{E}_j(A_n)\to1$ strongly as $n\to\infty$.
The function $1_{A_n}/\tilde\alpha_0$ is bounded by $n$;
the multiplicativity of $h\mapsto h(\mathsf{H}_j,\mathsf{P})$ and \eqref{eq:rotation-reference-field} give
\[
\mathfrak{E}_j(A_n)\,[\mathbb{B}_a(u)]_j
=\left(\frac{1_{A_n}}{\tilde\alpha_0}\right)(\mathsf{H}_j,\mathsf{P})\,\tilde\alpha_0(\mathsf{H}_j,\mathsf{P})\,[\mathbb{B}_a(u)]_j
=\left(\frac{1_{A_n}\cdot\tilde u}{\tilde\alpha_0}\right)(\mathsf{H}_j,\mathsf{P})\,\psi_a .
\]
For $u,u'$ as in the statement it follows that
\[
\begin{aligned}
\left\langle\mathfrak{E}_j(A_n)[\mathbb{B}_a(u)]_j,[\mathbb{B}_b(u')]_j\right\rangle
&=\left\langle\mathfrak{E}_j(A_n)[\mathbb{B}_a(u)]_j,\mathfrak{E}_j(A_n)[\mathbb{B}_b(u')]_j\right\rangle\\
&=\int_{A_n}\mathrm{d}m_{ab}\,\frac{\overline{\tilde u}\cdot\tilde u'}{|\tilde\alpha_0|^2}
=\int_{A_n}\mathrm{d}\varrho_{ab}\,\overline{\tilde u}\cdot\tilde u' .
\end{aligned}
\]
Take first $b=a$ and $u'=u$, and let $n\to\infty$.
The left side tends to $\left\|[\mathbb{B}_a(u)]_j\right\|^2$, which is at most $\|u\|_2^2$ by Step~1,
and the right side tends to $\int\mathrm{d}\varrho_{aa}\,|\tilde u|^2$ by monotone convergence,
$\varrho_{aa}$ being positive by \eqref{it:spectral-matrix-positive}.
This is the first assertion of \eqref{it:spectral-matrix-formula}.
For general $a,b$, part \eqref{it:spectral-matrix-positive} with the constant functions $c_{a'}:=\delta_{aa'}$ and $c'_{b'}:=\delta_{bb'}$,
for which $\varrho_{cc'}=\varrho_{ab}$, gives
\[
\int\mathrm{d}\left|\varrho_{ab}\right|\,|\tilde u|\cdot|\tilde u'|
\le\left(\int\mathrm{d}\varrho_{aa}\,|\tilde u|^2\right)^{\frac12}\cdot\left(\int\mathrm{d}\varrho_{bb}\,|\tilde u'|^2\right)^{\frac12}<\infty .
\]
Dominated convergence therefore applies to the right side of the preceding identity,
while its left side tends to $\left\langle[\mathbb{B}_a(u)]_j,[\mathbb{B}_b(u')]_j\right\rangle$;
the limit is \eqref{eq:rotation-spectral-matrix}.
\end{proof}

In the remainder of Subsection~\ref{sec:two-point-function}, the profiles of all labels lie in $C^\infty_{\mathrm c}(\R;\R)$.

\begin{definition}[Two pictures of the plane]\label{def:two-pictures}
Let $j\in\{1,2\}$, and let $i$ be the other element of $\{1,2\}$.
A point $y=(y^0,y^1)$ of the plane of Definition~\ref{def:labels-plates} is read as the vector $y^0e_1+y^1e_2$;
in the \emph{picture} $e_j$, its time coordinate is $y\cdot e_j$ and its space coordinate is $y\cdot e_i$.
As in Definition~\ref{def:labels-plates}, $\theta_j$ also denotes the restriction of the reflection $\theta_j$ to the plane,
$s$ is the sign of Definition~\ref{def:labels-plates}, and $e^\sigma_a$ is the $a$-th component of $e^\sigma$.
Put
\[
s_1(\sigma):=s(\sigma),\qquad s_2(\sigma):=-1\qquad\left(\sigma\in\{+,-,3\}\right),
\qquad\upsilon_1:=0,\qquad\upsilon_2:=\frac{\pi}{2}.
\]
For $\alpha\in L^1\cap L^2(\R^2;\C)$ with $\alpha=0$ almost everywhere on $\{y\cdot e_j<0\}$,
and for $\chi\in C^\infty_{\mathrm c}(\R;\R)$, put
\[
\begin{aligned}
\alpha^{[j]}(\omega,p_i)&:=\int\mathrm{d}y\,\alpha(y)\cdot\exp\left(-(y\cdot e_j)\cdot\omega+i\,(y\cdot e_i)\cdot p_i\right),\\
\hat\chi(p_3)&:=\int\mathrm{d}x_3\,\chi(x_3)\cdot\exp\left(ip_3\cdot x_3\right),
\end{aligned}
\]
where $(\omega,p_i)\in[0,\infty)\times\R$ and $p_3\in\R$.
Let $\mu$ be a law that satisfies \textup{(G)}, $(T_{123})$, $(R_j)$ and $(RP_j)$,
and let $\varrho^{(j)}_{ab}$, $a,b\in\{1,2,3\}$, be measures as in Lemma~\ref{lem:spectral-matrix},
in the variables $(\omega,p_\perp)=(\omega,p_i,p_3)$.
For labels $\lambda=(\sigma,\chi)$ and $\lambda'=(\sigma',\chi')$ put
\[
\varrho^{[j]}_{\lambda\lambda'}:=\sum_{a,b=1}^{3}\overline{e^{\sigma}_a}\cdot e^{\sigma'}_b\cdot\overline{\hat\chi(p_3)}\cdot\widehat{\chi'}(p_3)\cdot\varrho^{(j)}_{ab},
\]
a complex Radon measure on $[0,\infty)\times\R^2$.
\end{definition}

Unlike the Fourier transform $\hat b(p)=\int\mathrm{d}x\,\exp(-ip\cdot x)\cdot b(x)$ of a field $b$,
the transform $\hat\chi$ has the sign $+$ in its exponent,
in accordance with the factor $\exp\left(ip_\perp\cdot x_\perp\right)$ in the transforms $\tilde u$ of Lemma~\ref{lem:spectral-matrix}.
For a real profile the transforms with the two signs are complex conjugate to each other, and so they have the same modulus.

In both pictures the reflection of plates takes one form:
Lemma~\ref{lem:plate-identities}\eqref{it:plate-reflection} for $j=1$ and Lemma~\ref{lem:plate-identities}\eqref{it:plate-second-reflection} for $j=2$ read
\begin{equation}\label{eq:rotation-picture-reflection}
\Theta_j\mathbb{B}_\lambda(\alpha)=s_j(\sigma)\cdot\mathbb{B}_\lambda\left(\overline{\alpha}\circ\theta_j\right)
\qquad\left(\alpha\in L^2(\R^2;\C),\ \lambda=(\sigma,\chi)\right).
\end{equation}
For $\zeta\in\C$ let $n_\zeta:=(\cos\zeta,\sin\zeta)\in\C^2$, which for real $\zeta$ is the unit vector of Definition~\ref{def:labels-plates},
and extend the coordinates $y\cdot e_j$ and $y\cdot e_i$ linearly to $\C^2$.
For $r>0$ one has
\begin{equation}\label{eq:rotation-picture-polar}
(rn_\zeta)\cdot e_j=r\cdot\cos(\zeta-\upsilon_j),
\qquad
(rn_\zeta)\cdot e_i=(-1)^{j-1}r\cdot\sin(\zeta-\upsilon_j).
\end{equation}

Let $\alpha$ and $\chi$ be as in Definition~\ref{def:two-pictures}.
The function $u:x\mapsto\alpha(y)\cdot\chi(x_3)$, $x=(y,x_3)$, lies in $L^1\cap L^2(\R^3;\C)$ and vanishes on $\{x_j<0\}$,
and its transform in Lemma~\ref{lem:spectral-matrix} is
\[
\tilde u(\omega,p_\perp)=\alpha^{[j]}(\omega,p_i)\cdot\hat\chi(p_3),\qquad p_\perp=(p_i,p_3).
\]
The measure $\varrho^{[j]}_{\lambda\lambda'}$ of Definition~\ref{def:two-pictures} is the measure $\varrho_{cc'}$ of Lemma~\ref{lem:spectral-matrix}\eqref{it:spectral-matrix-positive}
for the bounded continuous functions $c_a:=e^\sigma_a\cdot\hat\chi(p_3)$ and $c'_a:=e^{\sigma'}_a\cdot\widehat{\chi'}(p_3)$;
in particular $\varrho^{[j]}_{\lambda\lambda}\ge0$.
Let $\alpha'$ be a second function as in Definition~\ref{def:two-pictures},
and let $u'$ be formed from $\alpha'$ and $\chi'$ as $u$ is formed from $\alpha$ and $\chi$.
Then $\mathbb{B}_\lambda(\alpha)=\sum_ae^\sigma_a\cdot\mathbb{B}_a(u)$ and $\mathbb{B}_{\lambda'}(\alpha')=\sum_be^{\sigma'}_b\cdot\mathbb{B}_b(u')$.
Expanding the scalar product of their classes sesquilinearly and applying \eqref{eq:rotation-spectral-matrix} to each term,
one obtains a sum of integrals with respect to the $\varrho^{(j)}_{ab}$;
each integrand $\overline{\tilde u}\cdot\tilde u'$ is integrable for $\left|\varrho^{(j)}_{ab}\right|$
by Lemma~\ref{lem:spectral-matrix}\eqref{it:spectral-matrix-positive}, \eqref{it:spectral-matrix-formula},
so that the sum is one integral with respect to $\varrho^{[j]}_{\lambda\lambda'}$:
\begin{equation}\label{eq:rotation-picture-matrix}
\left\langle[\mathbb{B}_\lambda(\alpha)]_j,[\mathbb{B}_{\lambda'}(\alpha')]_j\right\rangle
=\int\mathrm{d}\varrho^{[j]}_{\lambda\lambda'}\,\overline{\alpha^{[j]}}\cdot{\alpha'}^{[j]},
\qquad
\varrho^{[j]}_{\lambda\lambda}\ge0 .
\end{equation}

\begin{lemma}[Plate kernels in one picture]\label{lem:picture-kernels}
Let $j\in\{1,2\}$, let $i$ be the other element of $\{1,2\}$,
and let $\mu$ be a law that satisfies \textup{(G)}, $(T_{123})$, $(R_j)$, $(RP_j)$ and $(C^1_{ji})$.
Let $\lambda=(\sigma,\chi)$ and $\lambda'=(\sigma',\chi')$ be labels with profiles in $C^\infty_{\mathrm c}(\R;\R)$,
and let $\varrho^{[j]}_{\lambda\lambda'}$, $s_j$ and $\upsilon_j$ be as in Definition~\ref{def:two-pictures}.
\begin{enumerate}
\item\label{it:picture-kernel-size}
For every $t>0$,
\[
I_\lambda(t):=\int\mathrm{d}\varrho^{[j]}_{\lambda\lambda}(\omega,p_\perp)\,\exp(-t\cdot\omega)\le\frac{36\,\|\chi\|_2^2}{\pi t^2}.
\]
\item\label{it:picture-kernel-density}
The function $G^{[j]}_{\lambda\lambda'}$ from the half plane $\{y\in\R^2:y\cdot e_j>0\}$ to $\C$ defined by
\[
G^{[j]}_{\lambda\lambda'}(y):=s_j(\sigma)\cdot\int\mathrm{d}\varrho^{[j]}_{\lambda\lambda'}(\omega,p_\perp)\,\exp\left(-(y\cdot e_j)\cdot\omega+i\,(y\cdot e_i)\cdot p_i\right)
\]
is smooth.
If $\alpha,\alpha'\in L^2(\R^2;\C)$ have compact supports and
$\sup\{y\cdot e_j:y\in\supp\alpha\}<\inf\{y\cdot e_j:y\in\supp\alpha'\}$, then
\[
S_2(\alpha\otimes\lambda,\alpha'\otimes\lambda')=\int\mathrm{d}y\int\mathrm{d}y'\,\alpha(y)\cdot\alpha'(y')\cdot G^{[j]}_{\lambda\lambda'}(y'-y).
\]
\item\label{it:picture-kernel-continuation}
Let $r>0$.
The function
\[
h^{[j]}_{\lambda\lambda'}(\zeta):=s_j(\sigma)\cdot\int\mathrm{d}\varrho^{[j]}_{\lambda\lambda'}(\omega,p_\perp)\,
\exp\left(-r\cdot\omega\cdot\cos(\zeta-\upsilon_j)+i(-1)^{j-1}r\cdot p_i\cdot\sin(\zeta-\upsilon_j)\right)
\]
is holomorphic on the strip $\left\{\zeta\in\C:|\Re\zeta-\upsilon_j|<\frac{\pi}{2}\right\}$,
it satisfies $h^{[j]}_{\lambda\lambda'}(\zeta)=G^{[j]}_{\lambda\lambda'}(rn_\zeta)$ for every real $\zeta$ in the strip, and
\[
\left|h^{[j]}_{\lambda\lambda'}(t+i\beta)\right|
\le\frac{36}{\pi}\cdot\|\chi\|_2\cdot\|\chi'\|_2\cdot\left(r\cdot\cos(t-\upsilon_j)\right)^{-2}\cdot\exp(2|\beta|)
\]
for all real $t,\beta$ with $|t-\upsilon_j|<\frac{\pi}{2}$.
\item\label{it:picture-kernel-positivity}
For every $r>0$ and every real $\beta$, the function $h^{[j]}_{\lambda\lambda}$ of \eqref{it:picture-kernel-continuation} satisfies
$s_j(\sigma)\cdot h^{[j]}_{\lambda\lambda}(\upsilon_j+i\beta)\ge0$.
\end{enumerate}
\end{lemma}

\begin{proof}
Let $A:=\{(\omega,p_\perp):\omega<|p_i|\}$.
The condition $(C^1_{ji})$ says that $\mathfrak{E}_j(A)\psi=0$ for every $\psi\in\mathcal{H}^1_j$,
and Lemma~\ref{lem:spectral-matrix}\eqref{it:spectral-matrix-support} gives $\left|\varrho^{(j)}_{ab}\right|(A)=0$ for all $a,b$.
As $\varrho^{[j]}_{\lambda\lambda'}$ is a sum of the measures $\varrho^{(j)}_{ab}$ multiplied by bounded functions,
also $\left|\varrho^{[j]}_{\lambda\lambda'}\right|(A)=0$:
every measure $\varrho^{[j]}_{\lambda\lambda'}$ is carried by the cone $\{\omega\ge|p_i|\}$.

\emph{Proof of \eqref{it:picture-kernel-size}.}
For $c>0$ let
\[
\alpha_c(y):=c^{-2}\cdot1_{(c,2c)}(y\cdot e_j)\cdot P_1\left(\frac{y\cdot e_i}{c}\right).
\]
This function lies in $L^1\cap L^2(\R^2;\C)$ and vanishes on $\{y\cdot e_j<0\}$; by the transform of $P_1$,
\[
\alpha_c^{[j]}(\omega,p_i)=\left(\frac{1}{c}\int_c^{2c}\mathrm{d}v\,\exp(-v\cdot\omega)\right)\cdot\exp(-c\cdot|p_i|)\ge\exp(-2c\cdot\omega-c\cdot|p_i|).
\]
On the cone $\{\omega\ge|p_i|\}$, which carries $\varrho^{[j]}_{\lambda\lambda}$, it follows that $\alpha_c^{[j]}\ge\exp(-3c\cdot\omega)$.
Moreover $\|\alpha_c\|_2^2=c^{-4}\cdot c\cdot c\cdot\|P_1\|_2^2=\frac{1}{2\pi c^2}$,
and the plate $\alpha_c\otimes\lambda$ has the norm $|e^\sigma|\cdot\|\alpha_c\|_2\cdot\|\chi\|_2$ with $|e^\sigma|^2\le2$ (Definition~\ref{def:labels-plates}).
By \eqref{eq:rotation-picture-matrix}, the contractivity of $[\,\cdot\,]_j$ \cite[Proposition~8.19\,(i)]{NoncompactQED}
and Lemma~\ref{lem:domination-consequences}\eqref{it:field-extension},
\[
\begin{aligned}
\int\mathrm{d}\varrho^{[j]}_{\lambda\lambda}\,\exp(-6c\cdot\omega)
&\le\int\mathrm{d}\varrho^{[j]}_{\lambda\lambda}\,\left|\alpha_c^{[j]}\right|^2
=\left\|[\mathbb{B}_\lambda(\alpha_c)]_j\right\|^2
\le\mu\left(\left|\mathbb{B}_\lambda(\alpha_c)\right|^2\right)\\
&\le\|\alpha_c\otimes\lambda\|^2
\le2\,\|\alpha_c\|_2^2\cdot\|\chi\|_2^2
=\frac{\|\chi\|_2^2}{\pi c^2}.
\end{aligned}
\]
The choice $c:=t/6$ gives \eqref{it:picture-kernel-size}.

\emph{Proof of \eqref{it:picture-kernel-density}.}
Let $I_{\lambda'}$ be formed from $\varrho^{[j]}_{\lambda'\lambda'}$ as $I_\lambda$ is formed from $\varrho^{[j]}_{\lambda\lambda}$.
Lemma~\ref{lem:spectral-matrix}\eqref{it:spectral-matrix-positive} with $G_1=G_2=\exp\left(-\frac{t\cdot\omega}{2}\right)$, together with \eqref{it:picture-kernel-size}, gives
\begin{equation}\label{eq:rotation-picture-variation}
\int\mathrm{d}\left|\varrho^{[j]}_{\lambda\lambda'}\right|\,\exp(-t\cdot\omega)\le\left(I_\lambda(t)\cdot I_{\lambda'}(t)\right)^{\frac12}<\infty\qquad(t>0).
\end{equation}
Let $t>0$ and $y\cdot e_j\ge t$.
A partial derivative of order $n$ in $y$ of the integrand $\exp\left(-(y\cdot e_j)\cdot\omega+i\,(y\cdot e_i)\cdot p_i\right)$
is this integrand multiplied by $n$ factors, each equal to $-\omega$ or to $ip_i$;
on the cone $\{\omega\ge|p_i|\}$ its modulus is therefore at most
\[
\omega^n\cdot\exp(-t\cdot\omega)\le n!\cdot\left(\frac{2}{t}\right)^n\cdot\exp\left(-\frac{t\cdot\omega}{2}\right),
\]
a function that is integrable for $\left|\varrho^{[j]}_{\lambda\lambda'}\right|$ by \eqref{eq:rotation-picture-variation}.
Differentiation under the integral sign is thus permitted to every order on $\{y\cdot e_j>t\}$,
and $G^{[j]}_{\lambda\lambda'}$ is smooth on $\{y\cdot e_j>0\}$.

Now let $\alpha,\alpha'$ be as in \eqref{it:picture-kernel-density}.
For $v\in\R^2$ let $\tau_v\alpha$ be the translate $y\mapsto\alpha(y-v)$;
with $v$ read as a vector of $\R^3$ as in Definition~\ref{def:labels-plates},
the plate $(\tau_v\alpha)\otimes\lambda$ is the field $\tau_v\cdot(\alpha\otimes\lambda)$.
Neither side of the asserted identity changes when $\alpha$ and $\alpha'$ are replaced by $\tau_v\alpha$ and $\tau_v\alpha'$:
the left side does not, because $S_2(\tau_v\cdot f,\tau_v\cdot f')=S_2(f,f')$ for all $f,f'\in L^2(\R^3;\C^3)$
by $(T_{123})$ and Lemma~\ref{lem:domination-consequences}\eqref{it:field-extension},
and the right side does not, by the substitution $(y,y')\mapsto(y+v,y'+v)$.
After a translation by a suitable multiple of $e_j$ one may therefore assume $\supp\alpha\subset\{y\cdot e_j<0\}$ and $\supp\alpha'\subset\{y\cdot e_j>0\}$.
Let
\[
\delta:=\inf\{y'\cdot e_j:y'\in\supp\alpha'\}-\sup\{y\cdot e_j:y\in\supp\alpha\}>0.
\]
The function $\alpha^\theta:=\overline{\alpha}\circ\theta_j$ vanishes outside $\{y\cdot e_j>0\}$.
Applied to $\alpha^\theta$, the identity \eqref{eq:rotation-picture-reflection} gives
$\Theta_j\mathbb{B}_\lambda(\alpha^\theta)=s_j(\sigma)\cdot\mathbb{B}_\lambda(\alpha)$, that is, $\mathbb{B}_\lambda(\alpha)=s_j(\sigma)\cdot\Theta_j\mathbb{B}_\lambda(\alpha^\theta)$, since $s_j(\sigma)^2=1$.
The functions $\alpha^\theta$ and $\alpha'$ have compact supports, so they lie in $L^1\cap L^2(\R^2;\C)$,
and $\mathbb{B}_\lambda(\alpha^\theta)$, $\mathbb{B}_{\lambda'}(\alpha')$ lie in $L^2(\Sigma_j^+,\mu)$
(Lemma~\ref{lem:domination-consequences}\eqref{it:field-extension}).
The scalar product of the reconstruction satisfies $\left\langle[\Phi]_j,[\Phi']_j\right\rangle=\mu\left((\Theta_j\Phi)\cdot\Phi'\right)$ for $\Phi,\Phi'\in L^2(\Sigma_j^+,\mu)$
(\cite[Proposition~8.19\,(i)]{NoncompactQED}, where the reflection is the linear map written $\hat\theta_j$ here, so that $\Theta_j\Phi=\overline{\hat\theta_j\Phi}$).
With $S_2(f,f')=\mu\left(\mathbb{B}(f)\cdot\mathbb{B}(f')\right)$ for $f,f'\in L^2(\R^3;\C^3)$ (Lemma~\ref{lem:domination-consequences}\eqref{it:moment-forms}) and \eqref{eq:rotation-picture-matrix} one obtains
\[
S_2(\alpha\otimes\lambda,\alpha'\otimes\lambda')
=s_j(\sigma)\cdot\mu\left(\left(\Theta_j\mathbb{B}_\lambda(\alpha^\theta)\right)\cdot\mathbb{B}_{\lambda'}(\alpha')\right)
=s_j(\sigma)\cdot\int\mathrm{d}\varrho^{[j]}_{\lambda\lambda'}\,\overline{(\alpha^\theta)^{[j]}}\cdot{\alpha'}^{[j]} .
\]
As $\theta_j$ preserves the Lebesgue measure of the plane, reverses $y\cdot e_j$ and fixes $y\cdot e_i$,
\[
\begin{aligned}
\overline{(\alpha^\theta)^{[j]}(\omega,p_i)}
&=\int\mathrm{d}y\,\alpha(\theta_jy)\cdot\exp\left(-(y\cdot e_j)\cdot\omega-i\,(y\cdot e_i)\cdot p_i\right)\\
&=\int\mathrm{d}y\,\alpha(y)\cdot\exp\left((y\cdot e_j)\cdot\omega-i\,(y\cdot e_i)\cdot p_i\right),
\end{aligned}
\]
and consequently
\[
\overline{(\alpha^\theta)^{[j]}}\cdot{\alpha'}^{[j]}
=\int\mathrm{d}y\int\mathrm{d}y'\,\alpha(y)\cdot\alpha'(y')\cdot\exp\left(-\left((y'-y)\cdot e_j\right)\cdot\omega+i\left((y'-y)\cdot e_i\right)\cdot p_i\right).
\]
On the supports $(y'-y)\cdot e_j\ge\delta$, so the last integrand has modulus at most $|\alpha(y)|\cdot|\alpha'(y')|\cdot\exp(-\delta\cdot\omega)$,
which is integrable with respect to $\mathrm{d}y\,\mathrm{d}y'\,\mathrm{d}\left|\varrho^{[j]}_{\lambda\lambda'}\right|$ by \eqref{eq:rotation-picture-variation}.
Combining the formula for $S_2(\alpha\otimes\lambda,\alpha'\otimes\lambda')$ as an integral against $\varrho^{[j]}_{\lambda\lambda'}$
with this formula for the product of the two transforms,
and interchanging the integrations by Fubini's theorem, one arrives at the asserted identity.

\emph{Proof of \eqref{it:picture-kernel-continuation}.}
Let $\zeta=t+i\beta$ with real $t,\beta$ and $|t-\upsilon_j|<\frac{\pi}{2}$.
From
\[
\begin{aligned}
\cos(t-\upsilon_j+i\beta)&=\cos(t-\upsilon_j)\cdot\cosh\beta-i\sin(t-\upsilon_j)\cdot\sinh\beta,\\
\sin(t-\upsilon_j+i\beta)&=\sin(t-\upsilon_j)\cdot\cosh\beta+i\cos(t-\upsilon_j)\cdot\sinh\beta
\end{aligned}
\]
one finds that the real part of the exponent in the definition of $h^{[j]}_{\lambda\lambda'}$ equals
\[
-r\cdot\cos(t-\upsilon_j)\cdot\left(\omega\cdot\cosh\beta+(-1)^{j-1}p_i\cdot\sinh\beta\right).
\]
At every point of the cone $\{\omega\ge|p_i|\}$ one has
\[
\omega\cdot\cosh\beta+(-1)^{j-1}p_i\cdot\sinh\beta\ge\omega\cdot\left(\cosh\beta-|\sinh\beta|\right)=\omega\cdot\exp(-|\beta|);
\]
as $\cos(t-\upsilon_j)>0$, it follows that this real part is at most $-r\cdot\cos(t-\upsilon_j)\cdot\exp(-|\beta|)\cdot\omega$ on the cone.
(For $0<c<1$, the expression $\omega\cdot\cosh\beta+(-1)^{j-1}p_i\cdot\sinh\beta$ is non-negative at every point of the set $\{\omega\ge c\cdot|p_i|\}$
if and only if $\tanh|\beta|\le c$.)
The integrand is an entire function of $\zeta$.
On a compact subset of the strip, $r\cdot\cos(\Re\zeta-\upsilon_j)\cdot\exp(-|\Im\zeta|)$ has a positive lower bound $\varepsilon$,
so that there the integrand is bounded on the cone by $\exp(-\varepsilon\cdot\omega)$,
a function integrable for $\left|\varrho^{[j]}_{\lambda\lambda'}\right|$ by \eqref{eq:rotation-picture-variation}.
By dominated convergence, Fubini's theorem and Morera's theorem, $h^{[j]}_{\lambda\lambda'}$ is holomorphic on the strip.
With $t_\beta:=r\cdot\cos(t-\upsilon_j)\cdot\exp(-|\beta|)>0$, the estimate \eqref{eq:rotation-picture-variation} and part \eqref{it:picture-kernel-size} give
\[
\left|h^{[j]}_{\lambda\lambda'}(t+i\beta)\right|
\le\int\mathrm{d}\left|\varrho^{[j]}_{\lambda\lambda'}\right|\,\exp(-t_\beta\cdot\omega)
\le\left(I_\lambda(t_\beta)\cdot I_{\lambda'}(t_\beta)\right)^{\frac12}
\le\frac{36}{\pi}\cdot\|\chi\|_2\cdot\|\chi'\|_2\cdot t_\beta^{-2},
\]
which is the asserted bound.
For real $\zeta$ in the strip, \eqref{eq:rotation-picture-polar} shows that the exponent equals
$-\left((rn_\zeta)\cdot e_j\right)\cdot\omega+i\left((rn_\zeta)\cdot e_i\right)\cdot p_i$, with $(rn_\zeta)\cdot e_j=r\cdot\cos(\zeta-\upsilon_j)>0$;
hence $h^{[j]}_{\lambda\lambda'}(\zeta)=G^{[j]}_{\lambda\lambda'}(rn_\zeta)$.

\emph{Proof of \eqref{it:picture-kernel-positivity}.}
At $\zeta=\upsilon_j+i\beta$ the exponent equals $-r\cdot\omega\cdot\cosh\beta-(-1)^{j-1}r\cdot p_i\cdot\sinh\beta$, a real number.
Hence
\[
s_j(\sigma)\cdot h^{[j]}_{\lambda\lambda}(\upsilon_j+i\beta)
=s_j(\sigma)^2\cdot\int\mathrm{d}\varrho^{[j]}_{\lambda\lambda}(\omega,p_\perp)\,\exp\left(-r\cdot\omega\cdot\cosh\beta-(-1)^{j-1}r\cdot p_i\cdot\sinh\beta\right),
\]
and the right side is non-negative, because $\varrho^{[j]}_{\lambda\lambda}\ge0$ by \eqref{eq:rotation-picture-matrix}.
\end{proof}

\subsubsection{Invariance in one plane}\label{sec:two-point-plane}

\begin{theorem}[The two-point function]\label{thm:two-point}
Let $\mu$ be a law on $\mathcal{S}'$ that satisfies \textup{(G)}, \textup{(dF)}, $(T_{123})$, $(R_1)$ and $(RP_1)$,
and let $S_2$ be its two-point function (Lemma~\ref{lem:domination-consequences}\eqref{it:moment-forms}).
\begin{enumerate}
\item\label{it:two-point-invariance}
\emph{One plane.}
Assume the one-field cone $(C^1_{12})$ of Definition~\ref{def:closedness-cones} and one of the following two conditions.
\begin{enumerate}
\item\label{it:quarter-turn-two-point}
The two-point function is invariant under the quarter turn:
\begin{equation}\label{eq:rotation-two-point-quarter-turn}
S_2(R_{\pi/2}\cdot b,R_{\pi/2}\cdot b')=S_2(b,b')\qquad(b,b'\in L^2(\R^3;\C^3)).
\end{equation}
This is the case if $\mu$ satisfies $(Q_{12})$, by Lemma~\ref{lem:domination-consequences}\eqref{it:field-extension}.
\item\label{it:second-mirror-two-point}
The law $\mu$ satisfies $(R_2)$, $(RP_2)$ and $(C^1_{21})$;
here invariance under the reflection $\theta_2$ in the plane $\{x_2=0\}$, reflection positivity at this plane
and the one-field cone of the picture $e_2$ take the place of the quarter turn.
\end{enumerate}
Then, for the rotations $R_t=R^{12}_t$ of the plane $(1,2)$,
\[
S_2(R_t\cdot b,R_t\cdot b')=S_2(b,b')
\]
for every real $t$ and all $b,b'\in L^2(\R^3;\C^3)$, also when the supports of $b$ and $b'$ meet.
\item\label{it:two-point-frequencies}
\emph{Frequencies.}
Assume the hypotheses of \eqref{it:two-point-invariance},
and let $\lambda=(\sigma,\chi)$ and $\lambda'=(\sigma',\chi')$ be labels (Definition~\ref{def:labels-plates})
with profiles $\chi,\chi'\in C^\infty_{\mathrm c}(\R;\R)$.
There are a smooth function $G_{\lambda\lambda'}\colon\R^2\setminus\{0\}\to\C$ and a function $a_{\lambda\lambda'}\colon(0,\infty)\to\C$ such that
\[
S_2(\alpha\otimes\lambda,\alpha'\otimes\lambda')=\int\mathrm{d}y\int\mathrm{d}y'\,\alpha(y)\cdot\alpha'(y')\cdot G_{\lambda\lambda'}(y'-y)
\]
for all $\alpha,\alpha'\in L^2(\R^2;\C)$ with disjoint compact supports,
and such that, for $r>0$ and real $\vartheta$,
\[
G_{\lambda\lambda'}(rn_\vartheta)=a_{\lambda\lambda'}(r)\cdot\exp(i(w(\sigma)+w(\sigma'))\vartheta),
\qquad
|a_{\lambda\lambda'}(r)|\le\frac{72}{\pi}\cdot\|\chi\|_2\cdot\|\chi'\|_2\cdot r^{-2}.
\]
\item\label{it:two-point-kallen-lehmann}
\emph{K\"all\'en--Lehmann form.}
Assume that $S_2(g\cdot b,g\cdot b')=S_2(b,b')$ for all $g\in\mathrm{O}(3)$ and all $b,b'\in L^2(\R^3;\C^3)$.
Then there is exactly one finite positive Borel measure $\mathfrak{K}$ on $[0,\infty)$ such that
\[
S_2(\overline{b},b')=(2\pi)^{-3}\int\mathrm{d}p\,\hat b(p)^{\dagger}P^{\perp}(p)\,\hat b'(p)\int_{[0,\infty)}\mathrm{d}\mathfrak{K}(m)\,\frac{|p|^2}{|p|^2+m^2}
\qquad(b,b'\in L^2(\R^3;\C^3)),
\]
and its mass satisfies $\mathfrak{K}([0,\infty))\le1$.
Here $\hat b(p)=\int\mathrm{d}x\,\exp(-ip\cdot x)\cdot b(x)$ is the Fourier transform in the convention that follows Definition~\ref{def:field-actions},
$\hat b(p)^{\dagger}$ is its conjugate transpose,
and $P^{\perp}(p)=1-pp^{T}/|p|^2$ is the multiplier of the transverse projection $P^{\perp}$ (Remark~\ref{rem:maxwell-domination}).
For $b,b'\in L^2(\R^3;\C^3)$ with disjoint compact supports, moreover,
\[
S_2(b,b')=\int\mathrm{d}x\int\mathrm{d}x'\,b(x)^{T}S(x'-x)\,b'(x'),
\]
where the matrix kernel $S=(S_{ab})_{a,b=1}^{3}$ is real analytic on $\R^3\setminus\{0\}$ and given by
\[
S_{ab}(x)=(-\delta_{ab}\Delta+\partial_a\partial_b)\int_{[0,\infty)}\mathrm{d}\mathfrak{K}(m)\,\frac{\exp(-m|x|)}{4\pi|x|}\qquad(x\neq0),
\]
with the Kronecker symbol $\delta_{ab}$, the Laplace operator $\Delta$ and the derivative $\partial_a$ in the direction $e_a$.
The free Maxwell law $\gamma$ of Definition~\ref{def:free-maxwell} has $\mathfrak{K}=\delta_0$, the unit mass at $m=0$;
the law $\mu_0$ concentrated on $F=0$ has $\mathfrak{K}=0$.
\end{enumerate}
\end{theorem}

\begin{proof}[Proof of \eqref{it:two-point-invariance} and \eqref{it:two-point-frequencies}]
Assume the hypotheses of \eqref{it:two-point-invariance}.
Let $\lambda=(\sigma,\chi)$ and $\lambda'=(\sigma',\chi')$ be labels with profiles in $C^\infty_{\mathrm c}(\R;\R)$, and put
\[
w:=w(\sigma)+w(\sigma'),\qquad c:=\frac{72}{\pi}\cdot\|\chi\|_2\cdot\|\chi'\|_2;
\]
the weight $w$ is an integer.
For a real angle $\upsilon$, a pair $(K,K')$ of compact subsets of the plane is said to be \emph{ordered along} $n_\upsilon$ if
\[
\sup\{n_\upsilon\cdot y:\ y\in K\}<\inf\{n_\upsilon\cdot y:\ y\in K'\},
\]
and a pair $(\alpha,\alpha')$ of functions on the plane with compact supports is ordered along $n_\upsilon$ if $(\supp\alpha,\supp\alpha')$ is.
For such a pair, $n_\upsilon\cdot(y'-y)>0$ whenever $y\in\supp\alpha$ and $y'\in\supp\alpha'$.

\emph{Step 1. One kernel on the punctured plane.}
For $\upsilon\in\{0,\pi/2,\pi,3\pi/2\}$ a function $G^{(\upsilon)}$ from the open half plane $\{y\in\R^2:\ n_\upsilon\cdot y>0\}$ to $\C$ is defined as follows:
in the case \eqref{it:quarter-turn-two-point} by
\[
G^{(\upsilon)}(y):=\exp(iw\upsilon)\cdot G^{[1]}_{\lambda\lambda'}(R_{-\upsilon}y),
\]
and in the case \eqref{it:second-mirror-two-point} by
\[
G^{(\upsilon)}(y):=
\begin{cases}
G^{[1]}_{\lambda\lambda'}(y), & \upsilon=0,\\
G^{[2]}_{\lambda\lambda'}(y), & \upsilon=\pi/2,\\
G^{[1]}_{\lambda'\lambda}(-y), & \upsilon=\pi,\\
G^{[2]}_{\lambda'\lambda}(-y), & \upsilon=3\pi/2.
\end{cases}
\]
Here $G^{[j]}_{\lambda\lambda'}$ is the smooth kernel of Lemma~\ref{lem:picture-kernels}\eqref{it:picture-kernel-density} in the picture $e_j$,
defined on $\{y\cdot e_j>0\}$.
For $j=2$ that lemma is available in the case \eqref{it:second-mirror-two-point},
since $\mu$ then satisfies \textup{(G)}, $(T_{123})$, $(R_2)$, $(RP_2)$ and $(C^1_{21})$.
The identities
\[
\begin{aligned}
&y\cdot e_1=n_0\cdot y,\qquad
y\cdot e_2=n_{\pi/2}\cdot y,\qquad
n_0\cdot R_{-\upsilon}y=n_\upsilon\cdot y,\\
&n_0\cdot(-y)=n_\pi\cdot y,\qquad
n_{\pi/2}\cdot(-y)=n_{3\pi/2}\cdot y
\end{aligned}
\]
show that each $G^{(\upsilon)}$ is defined, and smooth, on its half plane.

Next, integration against $G^{(\upsilon)}$ computes $S_2$ on plates whose factors in the plane are ordered:
the density formula
\[
S_2(\alpha\otimes\lambda,\alpha'\otimes\lambda')=\int\mathrm{d}y\int\mathrm{d}y'\,\alpha(y)\cdot\alpha'(y')\cdot G^{(\upsilon)}(y'-y)
\]
holds for all $\alpha,\alpha'\in L^2(\R^2;\C)$ with compact supports such that $(\alpha,\alpha')$ is ordered along $n_\upsilon$.
In the case \eqref{it:quarter-turn-two-point} this is seen as follows.
Applied to $R_{-\pi/2}\cdot b$ and $R_{-\pi/2}\cdot b'$, the identity \eqref{eq:rotation-two-point-quarter-turn} shows that $S_2$ is invariant under $R_{-\pi/2}$,
and therefore under $R_{-\upsilon}$.
By Lemma~\ref{lem:plate-identities}\eqref{it:plate-rotation},
$R_{-\upsilon}\cdot(\alpha\otimes\lambda)=\exp(iw(\sigma)\upsilon)\cdot(\alpha\circ R_\upsilon)\otimes\lambda$,
and likewise for $\alpha'\otimes\lambda'$; hence
\[
\begin{aligned}
S_2(\alpha\otimes\lambda,\alpha'\otimes\lambda')
&=S_2(R_{-\upsilon}\cdot(\alpha\otimes\lambda),R_{-\upsilon}\cdot(\alpha'\otimes\lambda'))\\
&=\exp(iw\upsilon)\cdot S_2((\alpha\circ R_\upsilon)\otimes\lambda,(\alpha'\circ R_\upsilon)\otimes\lambda').
\end{aligned}
\]
The supports of $\alpha\circ R_\upsilon$ and $\alpha'\circ R_\upsilon$ are $R_{-\upsilon}\supp\alpha$ and $R_{-\upsilon}\supp\alpha'$.
As $n_0\cdot R_{-\upsilon}y=n_\upsilon\cdot y$, the pair $(\alpha\circ R_\upsilon,\alpha'\circ R_\upsilon)$ is ordered along $n_0$,
which is the order required in Lemma~\ref{lem:picture-kernels}\eqref{it:picture-kernel-density} for $j=1$.
That lemma, followed by the substitution $(y,y')\mapsto(R_{-\upsilon}y,R_{-\upsilon}y')$, which preserves the Lebesgue measure, gives
\[
S_2((\alpha\circ R_\upsilon)\otimes\lambda,(\alpha'\circ R_\upsilon)\otimes\lambda')
=\int\mathrm{d}y\int\mathrm{d}y'\,\alpha(y)\cdot\alpha'(y')\cdot G^{[1]}_{\lambda\lambda'}(R_{-\upsilon}(y'-y)),
\]
and multiplication by $\exp(iw\upsilon)$ yields the density formula for $G^{(\upsilon)}$.

In the case \eqref{it:second-mirror-two-point}, the density formula for $\upsilon=0$ and for $\upsilon=\pi/2$
is Lemma~\ref{lem:picture-kernels}\eqref{it:picture-kernel-density} itself, with $j=1$ and $j=2$ respectively.
Let now $\upsilon\in\{\pi,3\pi/2\}$, and let $j:=1$ if $\upsilon=\pi$ and $j:=2$ if $\upsilon=3\pi/2$.
Since $n_{\upsilon-\pi}=-n_\upsilon$,
the pair $(\alpha,\alpha')$ is ordered along $n_\upsilon$ if and only if $(\alpha',\alpha)$ is ordered along $n_{\upsilon-\pi}$,
and $n_{\upsilon-\pi}\cdot y=y\cdot e_j$.
The symmetry of $S_2$ (Lemma~\ref{lem:domination-consequences}\eqref{it:moment-forms})
and Lemma~\ref{lem:picture-kernels}\eqref{it:picture-kernel-density}, applied to the labels $(\lambda',\lambda)$ in the picture $e_j$, give
\[
S_2(\alpha\otimes\lambda,\alpha'\otimes\lambda')=S_2(\alpha'\otimes\lambda',\alpha\otimes\lambda)
=\int\mathrm{d}y'\int\mathrm{d}y\,\alpha'(y')\cdot\alpha(y)\cdot G^{[j]}_{\lambda'\lambda}(y-y'),
\]
and $G^{[j]}_{\lambda'\lambda}(y-y')=G^{(\upsilon)}(y'-y)$ by definition.

If $y=(y^0,y^1)\neq0$, at least one of the four numbers
$n_0\cdot y=y^0$, $n_{\pi/2}\cdot y=y^1$, $n_\pi\cdot y=-y^0$ and $n_{3\pi/2}\cdot y=-y^1$ is positive;
hence the union of the four half planes $\{n_\upsilon\cdot y>0\}$ is $\R^2\setminus\{0\}$.
Let $\upsilon$ and $\upsilon'$ be two of the four angles, and let
\[
U:=\{(y,y')\in(\R^2)^2:\ n_\upsilon\cdot(y'-y)>0,\ n_{\upsilon'}\cdot(y'-y)>0\}.
\]
The functions $(y,y')\mapsto G^{(\upsilon)}(y'-y)$ and $(y,y')\mapsto G^{(\upsilon')}(y'-y)$ are continuous on the open set $U$.
Let $\alpha,\alpha'$ be bounded measurable functions on $\R^2$ with compact supports and $\supp\alpha\times\supp\alpha'\subset U$;
in particular $\alpha,\alpha'\in L^2(\R^2;\C)$.
The extreme values of $n_\upsilon\cdot y$ and of $n_{\upsilon'}\cdot y$ on the compact supports are attained;
so $(\alpha,\alpha')$ is ordered along $n_\upsilon$ and along $n_{\upsilon'}$, and both functions are densities of the same form for such pairs.
By Lemma~\ref{lem:analysis-facts}\eqref{it:equal-densities} they coincide on $U$,
which means that $G^{(\upsilon)}=G^{(\upsilon')}$ on the intersection of the two half planes.
Consequently there is exactly one function $G_{\lambda\lambda'}$ on $\R^2\setminus\{0\}$ with
\begin{equation}\label{eq:rotation-four-frames}
G_{\lambda\lambda'}(y)=G^{(\upsilon)}(y)\qquad(n_\upsilon\cdot y>0,\ \upsilon\in\{0,\pi/2,\pi,3\pi/2\}),
\end{equation}
and $G_{\lambda\lambda'}$ is smooth, because each $G^{(\upsilon)}$ is.

There remains the case of supports in arbitrary relative position.
Let $\alpha,\alpha'\in L^2(\R^2;\C)$ vanish outside two disjoint compact sets $K$ and $K'$;
it is to be shown that the density formula of \eqref{it:two-point-frequencies}, with the kernel $G_{\lambda\lambda'}$, holds for $\alpha$ and $\alpha'$.
The number $\min\{|y-y'|_\infty:\ y\in K,\ y'\in K'\}$ is positive;
choose $d>0$ such that $3d$ is smaller than it, and put $Q_k:=d(k+[0,1)^2)$ for $k=(k^0,k^1)\in\Z^2$.
If $Q_k$ meets $K$ and $Q_{k'}$ meets $K'$, then $|k^0-k'^0|\ge3$ or $|k^1-k'^1|\ge3$:
otherwise $|y-y'|_\infty<3d$ for all $y\in Q_k$ and $y'\in Q_{k'}$, which is impossible for $y\in Q_k\cap K$ and $y'\in Q_{k'}\cap K'$.
The closures of $Q_k$ and $Q_{k'}$ are then ordered along one of the four directions $n_\upsilon$:
if, for instance, $k'^0\ge k^0+3$, they lie in $\{y^0\le d(k^0+1)\}$ and in $\{y^0\ge d(k^0+3)\}$ and are ordered along $n_0$,
and the three other possibilities lead to $n_\pi$, $n_{\pi/2}$ and $n_{3\pi/2}$.
For such $k,k'$ the density formula for $G^{(\upsilon)}$, applied to $\alpha1_{Q_k}$ and $\alpha'1_{Q_{k'}}$,
together with \eqref{eq:rotation-four-frames}, gives
\[
S_2((\alpha1_{Q_k})\otimes\lambda,(\alpha'1_{Q_{k'}})\otimes\lambda')
=\int\mathrm{d}y\int\mathrm{d}y'\,(\alpha1_{Q_k})(y)\cdot(\alpha'1_{Q_{k'}})(y')\cdot G_{\lambda\lambda'}(y'-y).
\]
Only finitely many squares meet $K\cup K'$.
Almost everywhere, $\alpha$ is the sum of the functions $\alpha1_{Q_k}$ with $Q_k\cap K\neq\emptyset$,
and $\alpha'$ is the sum of the functions $\alpha'1_{Q_{k'}}$ with $Q_{k'}\cap K'\neq\emptyset$.
Summation over these pairs, by the bilinearity of $S_2$ and the linearity of $\alpha\mapsto\alpha\otimes\lambda$,
gives the density formula of \eqref{it:two-point-frequencies} for $\alpha$ and $\alpha'$.

\emph{Step 2. Continuation in the angle.}
Fix $r>0$.
There is an entire function $h$ such that
\begin{equation}\label{eq:rotation-entire-angle}
\begin{aligned}
h(\vartheta)&=G_{\lambda\lambda'}(rn_\vartheta)&&\quad(\vartheta\in\R),\\
|h(t+i\beta)|&\le c\cdot r^{-2}\cdot\exp(2|\beta|)&&\quad(t,\beta\in\R),\\
h(\zeta+\pi/2)&=i^{w}\cdot h(\zeta)&&\quad(\zeta\in\C)\ \text{in the case \eqref{it:quarter-turn-two-point}},\\
h(\zeta+\pi)&=(-1)^{w}\cdot h(\zeta)&&\quad(\zeta\in\C)\ \text{in the case \eqref{it:second-mirror-two-point}}.
\end{aligned}
\end{equation}
To construct it, let $S_\upsilon:=\{\zeta\in\C:\ |\Re\zeta-\upsilon|<\pi/2\}$ for $\upsilon\in(\pi/2)\Z$.
Let $h^{[j]}_{\lambda\lambda'}$ and $h^{[j]}_{\lambda'\lambda}$ be the functions of Lemma~\ref{lem:picture-kernels}\eqref{it:picture-kernel-continuation}
for the radius $r$, for $j=1$, and in the case \eqref{it:second-mirror-two-point} also for $j=2$.
They are holomorphic on the strip $S_{\upsilon_j}$, where $\upsilon_1=0$ and $\upsilon_2=\pi/2$ (Definition~\ref{def:two-pictures}),
and $h^{[j]}_{\lambda\lambda'}(\vartheta)=G^{[j]}_{\lambda\lambda'}(rn_\vartheta)$ at the real points $\vartheta$ of $S_{\upsilon_j}$;
the same holds with $\lambda$ and $\lambda'$ exchanged.
On $S_\upsilon$ define, in the case \eqref{it:quarter-turn-two-point},
\[
h^{(\upsilon)}(\zeta):=\exp(iw\upsilon)\cdot h^{[1]}_{\lambda\lambda'}(\zeta-\upsilon),
\]
and, in the case \eqref{it:second-mirror-two-point}, for $k\in\Z$,
\[
h^{(\upsilon)}(\zeta):=
\begin{cases}
h^{[1]}_{\lambda\lambda'}(\zeta-2\pi k), & \upsilon=2\pi k,\\
h^{[2]}_{\lambda\lambda'}(\zeta-2\pi k), & \upsilon=2\pi k+\pi/2,\\
h^{[1]}_{\lambda'\lambda}(\zeta-\pi-2\pi k), & \upsilon=2\pi k+\pi,\\
h^{[2]}_{\lambda'\lambda}(\zeta-\pi-2\pi k), & \upsilon=2\pi k+3\pi/2.
\end{cases}
\]
In every case the argument on the right lies in the strip $S_{\upsilon_j}$ of the function on the right,
and its real part exceeds $\upsilon_j$ by exactly $\Re\zeta-\upsilon$;
so $h^{(\upsilon)}$ is holomorphic on $S_\upsilon$.

Extend the definition of Step~1 to all $\upsilon\in(\pi/2)\Z$ by $G^{(\upsilon+2\pi k)}:=G^{(\upsilon)}$, $k\in\Z$.
In the case \eqref{it:quarter-turn-two-point} this agrees with the formula of Step~1,
which does not change under $\upsilon\mapsto\upsilon+2\pi$ because $w$ is an integer;
and \eqref{eq:rotation-four-frames} now holds for every $\upsilon\in(\pi/2)\Z$.
Let $\vartheta$ be a real point of $S_\upsilon$.
In the case \eqref{it:quarter-turn-two-point}, since $R_{-\upsilon}n_\vartheta=n_{\vartheta-\upsilon}$,
\[
h^{(\upsilon)}(\vartheta)=\exp(iw\upsilon)\cdot G^{[1]}_{\lambda\lambda'}(rn_{\vartheta-\upsilon})
=\exp(iw\upsilon)\cdot G^{[1]}_{\lambda\lambda'}(R_{-\upsilon}rn_\vartheta)
=G^{(\upsilon)}(rn_\vartheta).
\]
In the case \eqref{it:second-mirror-two-point}, since $n_{\vartheta-2\pi k}=n_\vartheta$ and $n_{\vartheta-\pi}=-n_\vartheta$,
the four lines of the definition of $h^{(\upsilon)}(\vartheta)$ give
$G^{[1]}_{\lambda\lambda'}(rn_\vartheta)$, $G^{[2]}_{\lambda\lambda'}(rn_\vartheta)$,
$G^{[1]}_{\lambda'\lambda}(-rn_\vartheta)$ and $G^{[2]}_{\lambda'\lambda}(-rn_\vartheta)$,
which is $G^{(\upsilon)}(rn_\vartheta)$ in each line.
Since $n_\upsilon\cdot rn_\vartheta=r\cos(\vartheta-\upsilon)>0$,
\eqref{eq:rotation-four-frames} gives $h^{(\upsilon)}(\vartheta)=G_{\lambda\lambda'}(rn_\vartheta)$ in both cases.

Two distinct strips $S_\upsilon$ meet only if their indices differ by $\pi/2$.
For $\upsilon\in(\pi/2)\Z$ and $\upsilon':=\upsilon+\pi/2$,
the intersection $S_\upsilon\cap S_{\upsilon'}=\{\upsilon<\Re\zeta<\upsilon'\}$ is connected and contains the real interval $(\upsilon,\upsilon')$,
on which $h^{(\upsilon)}$ and $h^{(\upsilon')}$ both equal $G_{\lambda\lambda'}(rn_\vartheta)$;
by the identity theorem they coincide on $S_\upsilon\cap S_{\upsilon'}$.
As the strips cover $\C$, the functions $h^{(\upsilon)}$ are the restrictions of one entire function $h$,
and $h(\vartheta)=G_{\lambda\lambda'}(rn_\vartheta)$ for real $\vartheta$.

For the bound, let $t$ and $\beta$ be real and choose $\upsilon\in(\pi/2)\Z$ with $|t-\upsilon|\le\pi/4$, so that $(\cos(t-\upsilon))^{-2}\le2$.
The argument of the function $h^{[j]}$ that defines $h^{(\upsilon)}(t+i\beta)$ has imaginary part $\beta$,
and its real part exceeds $\upsilon_j$ by $t-\upsilon$.
Since $|\exp(iw\upsilon)|=1$, the estimate of Lemma~\ref{lem:picture-kernels}\eqref{it:picture-kernel-continuation},
in which $\chi$ and $\chi'$ enter symmetrically, gives
\[
|h(t+i\beta)|=|h^{(\upsilon)}(t+i\beta)|
\le\frac{36}{\pi}\cdot\|\chi\|_2\cdot\|\chi'\|_2\cdot(r\cos(t-\upsilon))^{-2}\cdot\exp(2|\beta|)
\le c\cdot r^{-2}\cdot\exp(2|\beta|).
\]

In the case \eqref{it:quarter-turn-two-point}, let $\zeta\in S_\upsilon$ and $\upsilon':=\upsilon+\pi/2$; then $\zeta+\pi/2\in S_{\upsilon'}$,
and the definition of the functions $h^{(\upsilon)}$ gives
\[
h(\zeta+\pi/2)=h^{(\upsilon')}(\zeta+\pi/2)=\exp(iw\upsilon')\cdot h^{[1]}_{\lambda\lambda'}(\zeta-\upsilon)
=i^{w}\cdot h^{(\upsilon)}(\zeta)=i^{w}\cdot h(\zeta),
\]
because $\exp(iw\upsilon')=i^{w}\cdot\exp(iw\upsilon)$.
In the case \eqref{it:second-mirror-two-point}, the law is invariant under $\theta_1$ and under $\theta_2$, by $(R_1)$ and $(R_2)$,
hence under $\theta_1\theta_2$, which is the rotation $R_\pi$ of $\R^3$.
By Lemma~\ref{lem:domination-consequences}\eqref{it:field-extension}, $S_2(R_\pi\cdot b,R_\pi\cdot b')=S_2(b,b')$ for all $b,b'\in L^2(\R^3;\C^3)$.
Let now $U:=\{(y,y')\in(\R^2)^2:\ y\neq y'\}$, an open set that $R_\pi$ maps onto itself.
By Step~1, the continuous function $(y,y')\mapsto G_{\lambda\lambda'}(y'-y)$ on $U$ is a density of $S_2$
on the pairs of plates $\alpha\otimes\lambda$, $\alpha'\otimes\lambda'$
in which $\alpha$ and $\alpha'$ have compact supports with $\supp\alpha\times\supp\alpha'\subset U$,
that is, disjoint compact supports.
The criterion \eqref{eq:rotation-density-criterion}, together with Lemma~\ref{lem:analysis-facts}\eqref{it:equal-densities}, gives
$G_{\lambda\lambda'}(R_\pi(y'-y))=\exp(iw\pi)\cdot G_{\lambda\lambda'}(y'-y)$ on $U$.
As $R_\pi y=-y$ on the plane, this says $G_{\lambda\lambda'}(-y)=(-1)^{w}\cdot G_{\lambda\lambda'}(y)$ for all $y\neq0$.
With $-n_\vartheta=n_{\vartheta+\pi}$ it follows that $h(\vartheta+\pi)=(-1)^{w}\cdot h(\vartheta)$ for real $\vartheta$.
The functions $\zeta\mapsto h(\zeta+\pi)$ and $\zeta\mapsto(-1)^{w}\cdot h(\zeta)$ are entire and agree on $\R$;
by the identity theorem, $h(\zeta+\pi)=(-1)^{w}\cdot h(\zeta)$ for all $\zeta\in\C$.

\emph{Step 3. The window of frequencies.}
Fix $r>0$, and let $h$ be the entire function of Step~2.
By \eqref{eq:rotation-entire-angle} it satisfies $|h(t+i\beta)|\le c\cdot r^{-2}\cdot\exp(2|\beta|)$ for all real $t$ and $\beta$.
Since $i^w=\exp\left(\frac{i\pi w}{2}\right)$ and $(-1)^w=\exp(i\pi w)$,
Lemma~\ref{lem:harmonics} applies to $h$ with the quasi-period $\frac{\pi}{2}$ in the case \eqref{it:quarter-turn-two-point} and $\pi$ in the case \eqref{it:second-mirror-two-point},
with the exponent of growth $2$ and with the constant $c\cdot r^{-2}$.
Let $a^k_{\lambda\lambda'}(r)$ be the coefficient of $\exp(ik\zeta)$ in the finite expansion of $h(\zeta)$ that the lemma provides;
the lemma also specifies the integers $k$ that can occur and bounds the coefficients.
At the real points $\vartheta$ one has $h(\vartheta)=G_{\lambda\lambda'}(rn_\vartheta)$, and $r>0$ was arbitrary.
This gives the \emph{window of frequencies}
\begin{equation}\label{eq:rotation-window}
\begin{aligned}
G_{\lambda\lambda'}(rn_\vartheta)&=\sum_{\substack{|k|\le2\\ k-w\in4\Z}}a^k_{\lambda\lambda'}(r)\cdot\exp(ik\vartheta)\quad\text{in the case \eqref{it:quarter-turn-two-point}},\\
G_{\lambda\lambda'}(rn_\vartheta)&=\sum_{\substack{|k|\le2\\ k-w\in2\Z}}a^k_{\lambda\lambda'}(r)\cdot\exp(ik\vartheta)\quad\text{in the case \eqref{it:second-mirror-two-point}},\\
\left|a^k_{\lambda\lambda'}(r)\right|&\le c\cdot r^{-2},
\end{aligned}
\qquad(r>0,\ \vartheta\in\R).
\end{equation}
In particular, the bound on the coefficients holds for small and for large radii alike.

In each case the sum in \eqref{eq:rotation-window} is a finite trigonometric sum,
so its coefficients are Fourier coefficients:
\[
a^k_{\lambda\lambda'}(r)=\frac{1}{2\pi}\int_0^{2\pi}\mathrm{d}\vartheta\,G_{\lambda\lambda'}(rn_\vartheta)\cdot\exp(-ik\vartheta).
\]
As $G_{\lambda\lambda'}$ is smooth on $\R^2\setminus\{0\}$, each $a^k_{\lambda\lambda'}$ is a smooth function of $r>0$.
For fixed $\sigma$ and $\sigma'$ it is real bilinear in the pair of profiles $(\chi,\chi')$.
Indeed, $S_2$ is bilinear and the plate $\alpha\otimes(\sigma,\chi)$ is linear in $\chi$,
while a continuous density of $(\alpha,\alpha')\mapsto S_2(\alpha\otimes\lambda,\alpha'\otimes\lambda')$ on the pairs with disjoint compact supports is unique by Lemma~\ref{lem:analysis-facts}\eqref{it:equal-densities};
hence $G_{\lambda\lambda'}$, and with it every $a^k_{\lambda\lambda'}(r)$, depends bilinearly on $(\chi,\chi')$.
The bound of Step~2 is of the order of $\exp(2|\beta|)$, and not $o(\exp(2|\beta|))$;
for this reason Lemma~\ref{lem:harmonics} does not exclude the two frequencies $k=\pm2$ at the ends of the window.

The harmonics obey the symmetry
\begin{equation}\label{eq:rotation-harmonic-symmetry}
a^k_{\lambda\lambda'}=(-1)^k\cdot a^k_{\lambda'\lambda}.
\end{equation}
To see this, recall that $S_2$ is symmetric (Lemma~\ref{lem:domination-consequences}\eqref{it:moment-forms}).
So $S_2(\alpha\otimes\lambda,\alpha'\otimes\lambda')=S_2(\alpha'\otimes\lambda',\alpha\otimes\lambda)$,
and the continuous functions $(y,y')\mapsto G_{\lambda\lambda'}(y'-y)$ and $(y,y')\mapsto G_{\lambda'\lambda}(y-y')$ are densities of the same form
on the open set $\{(y,y')\in(\R^2)^2:\ y\neq y'\}$.
By Lemma~\ref{lem:analysis-facts}\eqref{it:equal-densities} they coincide, that is, $G_{\lambda\lambda'}(y)=G_{\lambda'\lambda}(-y)$ for $y\neq0$.
The pair $(\lambda',\lambda)$ has the same weight $w$, so its expansion \eqref{eq:rotation-window} has the same frequencies;
with $-n_\vartheta=n_{\vartheta+\pi}$ one obtains
\[
\sum_k a^k_{\lambda\lambda'}(r)\cdot\exp(ik\vartheta)=\sum_k a^k_{\lambda'\lambda}(r)\cdot\exp\left(ik(\vartheta+\pi)\right)=\sum_k(-1)^k\cdot a^k_{\lambda'\lambda}(r)\cdot\exp(ik\vartheta),
\]
and the comparison of coefficients gives \eqref{eq:rotation-harmonic-symmetry}.

\emph{Step 4. The Bianchi identity.}
Write the points of the plane in polar coordinates, $y^0+iy^1=r\exp(i\vartheta)$.
The Wirtinger derivatives of Lemma~\ref{lem:plate-identities}\eqref{it:plate-gradient} then take the form
\[
\partial_z=\frac{\exp(-i\vartheta)}{2}\cdot\left(\partial_r-\frac{i}{r}\,\partial_\vartheta\right),
\qquad
\partial_{\bar z}=\frac{\exp(i\vartheta)}{2}\cdot\left(\partial_r+\frac{i}{r}\,\partial_\vartheta\right),
\]
and for $f\in C^1((0,\infty);\C)$ and an integer $k$ one finds
\begin{equation}\label{eq:rotation-wirtinger}
\begin{aligned}
\partial_z\left(f(r)\cdot\exp(ik\vartheta)\right)&=\frac12\exp(i(k-1)\vartheta)\cdot\left(f'(r)+\frac{k}{r}\cdot f(r)\right),\\
\partial_{\bar z}\left(f(r)\cdot\exp(ik\vartheta)\right)&=\frac12\exp(i(k+1)\vartheta)\cdot\left(f'(r)-\frac{k}{r}\cdot f(r)\right).
\end{aligned}
\end{equation}

Let $\eta,\alpha'\in C^\infty_{\mathrm c}(\R^2;\C)$ have disjoint supports,
let $\chi_0\in C^\infty_{\mathrm c}(\R;\R)$ have the derivative $\chi_0'$,
and let $\lambda'=(\sigma',\chi')$ be a label with $\chi'\in C^\infty_{\mathrm c}(\R;\R)$.
The function $u(x):=\eta(y)\cdot\chi_0(x_3)$ lies in $C^\infty_{\mathrm c}(\R^3;\C)$,
and Lemma~\ref{lem:plate-identities}\eqref{it:plate-gradient} splits its gradient into three plates:
\[
\nabla u=(\partial_z\eta)\otimes(+,\chi_0)+(\partial_{\bar z}\eta)\otimes(-,\chi_0)+\eta\otimes(3,\chi_0').
\]
The field vanishes on gradients, \eqref{eq:rotation-gradient-fields};
hence $S_2(\nabla u,\alpha'\otimes\lambda')=\mu\left(\mathbb{B}(\nabla u)\cdot\mathbb{B}_{\lambda'}(\alpha')\right)=0$, and by bilinearity
\[
S_2\left((\partial_z\eta)\otimes(+,\chi_0),\alpha'\otimes\lambda'\right)
+S_2\left((\partial_{\bar z}\eta)\otimes(-,\chi_0),\alpha'\otimes\lambda'\right)
+S_2\left(\eta\otimes(3,\chi_0'),\alpha'\otimes\lambda'\right)=0.
\]
In each term the plane parts have disjoint compact supports,
so the density formula of Step~1 applies, with the kernels $G_{(+,\chi_0)\lambda'}$, $G_{(-,\chi_0)\lambda'}$ and $G_{(3,\chi_0')\lambda'}$ respectively.
The compact set $\{y'-y:\ y\in\supp\eta,\ y'\in\supp\alpha'\}$ does not contain $0$.
Hence, for fixed $y'\in\supp\alpha'$, the functions $y\mapsto G_{(\pm,\chi_0)\lambda'}(y'-y)$ are smooth on a neighbourhood of $\supp\eta$,
and an integration by parts in $y$ gives
\[
\begin{aligned}
\int\mathrm{d}y\,(\partial_z\eta)(y)\cdot G_{(+,\chi_0)\lambda'}(y'-y)&=\int\mathrm{d}y\,\eta(y)\cdot\left(\partial_zG_{(+,\chi_0)\lambda'}\right)(y'-y),\\
\int\mathrm{d}y\,(\partial_{\bar z}\eta)(y)\cdot G_{(-,\chi_0)\lambda'}(y'-y)&=\int\mathrm{d}y\,\eta(y)\cdot\left(\partial_{\bar z}G_{(-,\chi_0)\lambda'}\right)(y'-y):
\end{aligned}
\]
the integration by parts changes the sign once,
and the inner derivative of $y\mapsto y'-y$ changes it a second time.
Therefore
\[
\int\mathrm{d}y\int\mathrm{d}y'\,\eta(y)\cdot\alpha'(y')\cdot\left(\partial_zG_{(+,\chi_0)\lambda'}+\partial_{\bar z}G_{(-,\chi_0)\lambda'}+G_{(3,\chi_0')\lambda'}\right)(y'-y)=0
\]
for all $\eta,\alpha'\in C^\infty_{\mathrm c}(\R^2;\C)$ with disjoint supports.
The function in brackets, evaluated at $y'-y$, is continuous on the open set $\{(y,y')\in(\R^2)^2:\ y\neq y'\}$,
and the pairs $(\eta,\alpha')$ just described are exactly the pairs of test functions in $C^\infty_{\mathrm c}$ whose product of supports lies in this set.
Lemma~\ref{lem:analysis-facts}\eqref{it:equal-densities}, applied with test functions in $C^\infty_{\mathrm c}$, shows that the function vanishes:
\begin{equation}\label{eq:rotation-bianchi}
\partial_zG_{(+,\chi_0)\lambda'}+\partial_{\bar z}G_{(-,\chi_0)\lambda'}+G_{(3,\chi_0')\lambda'}=0\qquad\text{on }\R^2\setminus\{0\}.
\end{equation}

Insert \eqref{eq:rotation-window} and \eqref{eq:rotation-wirtinger} into \eqref{eq:rotation-bianchi}.
From now on $a^k_{\sigma\sigma'}$ stands for $a^k_{\lambda\lambda'}$ when the profiles need not be displayed.
In \eqref{eq:rotation-bianchi} the first labels carry the profiles $\chi_0$, $\chi_0$ and $\chi_0'$,
the second label carries the profile $\chi'$,
and $\chi_0$ and $\chi'$ are arbitrary in $C^\infty_{\mathrm c}(\R;\R)$.
Put $a^k_{\sigma\sigma'}:=0$ for the integers $k$ that \eqref{eq:rotation-window} does not admit.
For every $r>0$ the left side of \eqref{eq:rotation-bianchi} at the point $rn_\vartheta$ becomes a finite linear combination of the functions $\vartheta\mapsto\exp(ik\vartheta)$, $k\in\Z$,
which are linearly independent;
so the coefficient of each of them vanishes:
\begin{equation}\label{eq:rotation-coefficients}
\frac12\left((a^{k+1}_{+\sigma'})'+\frac{k+1}{r}\cdot a^{k+1}_{+\sigma'}\right)
+\frac12\left((a^{k-1}_{-\sigma'})'-\frac{k-1}{r}\cdot a^{k-1}_{-\sigma'}\right)
+a^{k}_{3\sigma'}=0
\qquad(k\in\Z,\ r>0).
\end{equation}

Let $\kappa\ge0$ and $C\ge0$, and let $f\in C^1((0,\infty);\C)$ satisfy
\begin{equation}\label{eq:rotation-radial-decay}
r\cdot f'(r)=\kappa\cdot f(r),\qquad|f(r)|\le C\cdot r^{-2}\qquad(r>0).
\end{equation}
Then $f=0$.
Indeed, the derivative of $r\mapsto r^{-\kappa}\cdot f(r)$ is $r^{-\kappa-1}\cdot\left(r\cdot f'(r)-\kappa\cdot f(r)\right)=0$,
so that $f(r)=f(1)\cdot r^{\kappa}$;
then $|f(1)|\le C\cdot r^{-2-\kappa}$ for every $r>0$, and $r\to\infty$ gives $f(1)=0$.
For $\kappa=0$ the equation is $f'=0$.
Every harmonic $a^k_{\sigma\sigma'}$ satisfies the inequality in \eqref{eq:rotation-radial-decay},
with $C$ the constant of \eqref{eq:rotation-window} for its pair of labels;
so a harmonic that satisfies the equation in \eqref{eq:rotation-radial-decay} for some $\kappa\ge0$ vanishes.

\emph{The case \eqref{it:quarter-turn-two-point}.}
Here \eqref{eq:rotation-window} admits the following frequencies:
\[
\begin{aligned}
&\{0\}&&\text{for }(\sigma,\sigma')\in\{(3,3),(+,-),(-,+)\},\\
&\{1\}&&\text{for }(\sigma,\sigma')\in\{(+,3),(3,+)\},\\
&\{-1\}&&\text{for }(\sigma,\sigma')\in\{(-,3),(3,-)\},\\
&\{2,-2\}&&\text{for }(\sigma,\sigma')\in\{(+,+),(-,-)\}.
\end{aligned}
\]
So $k=w$ for every pair, with the two exceptions $a^{-2}_{++}$ and $a^{2}_{--}$.
For $\sigma'=+$ the three terms of \eqref{eq:rotation-bianchi} carry the frequencies $\{1,-3\}$, $\{1\}$ and $\{1\}$.
The coefficient of $\exp(-3i\vartheta)$, that is, \eqref{eq:rotation-coefficients} with $k=-3$, reduces to
\[
\frac12\left((a^{-2}_{++})'-\frac{2}{r}\cdot a^{-2}_{++}\right)=0 .
\]
Thus $a^{-2}_{++}$ satisfies \eqref{eq:rotation-radial-decay} with $\kappa=2$, and $a^{-2}_{++}=0$.
For $\sigma'=-$ the three terms carry the frequencies $\{-1\}$, $\{-1,3\}$ and $\{-1\}$;
the coefficient of $\exp(3i\vartheta)$, \eqref{eq:rotation-coefficients} with $k=3$, gives $(a^{2}_{--})'=\frac{2}{r}\cdot a^{2}_{--}$,
and again \eqref{eq:rotation-radial-decay} with $\kappa=2$ yields $a^{2}_{--}=0$.

\emph{The case \eqref{it:second-mirror-two-point}.}
Now \eqref{eq:rotation-window} admits the frequencies $\{0,2,-2\}$ if $w=0$, the frequencies $\{1,-1\}$ if $w=\pm1$, and the frequencies $\{0,2,-2\}$ if $w=\pm2$.
\begin{enumerate}
\item[(i)]
\emph{The harmonics $a^{\pm2}_{33}$ vanish.}
Let first $\sigma=\sigma'=3$ and $\chi'=\chi$, so that $\lambda=\lambda'=(3,\chi)$;
then $w=0$, and the signs of Definition~\ref{def:two-pictures} are $s_1(3)=s(3)=-1$ and $s_2(3)=-1$.
By its construction in Step~2, the function $h$ coincides with $h^{[1]}_{\lambda\lambda}$ on the strip $S_0$ and with $h^{[2]}_{\lambda\lambda}$ on the strip $S_{\pi/2}$;
by Step~3, $h(\zeta)=a^0+a^2\cdot\exp(2i\zeta)+a^{-2}\cdot\exp(-2i\zeta)$ with $a^k:=a^k_{\lambda\lambda}(r)$.
Lemma~\ref{lem:picture-kernels}\eqref{it:picture-kernel-positivity}, for $j=1$ and for $j=2$
(the hypotheses for $j=2$ hold in the case \eqref{it:second-mirror-two-point}), gives for all real $\beta$
\[
\begin{aligned}
-h(i\beta)&=-\left(a^0+a^2\cdot\exp(-2\beta)+a^{-2}\cdot\exp(2\beta)\right)\ge0,\\
-h\left(\frac{\pi}{2}+i\beta\right)&=-\left(a^0-a^2\cdot\exp(-2\beta)-a^{-2}\cdot\exp(2\beta)\right)\ge0 .
\end{aligned}
\]
Divide both inequalities by $\exp(2|\beta|)$.
The limit $\beta\to+\infty$ gives $-a^{-2}\ge0$ and $a^{-2}\ge0$;
the limit $\beta\to-\infty$ gives $-a^{2}\ge0$ and $a^{2}\ge0$.
Consequently $a^{2}_{(3,\chi)(3,\chi)}(r)=a^{-2}_{(3,\chi)(3,\chi)}(r)=0$ for every profile $\chi$.
For $k=\pm2$ the factor $(-1)^k$ in \eqref{eq:rotation-harmonic-symmetry} equals $1$,
so that $(\chi,\chi')\mapsto a^{k}_{(3,\chi)(3,\chi')}(r)$ is a symmetric real bilinear form;
such a form is determined by its values at the pairs $(\chi,\chi)$ (polarisation), and these are $0$.
Hence $a^{\pm2}_{(3,\chi)(3,\chi')}(r)=0$ for all profiles $\chi$ and $\chi'$.
\item[(ii)]
Take $\sigma'=+$ in \eqref{eq:rotation-bianchi}.
The frequency $-3$ occurs only in the term that contains $a^{-2}_{++}$, and the frequency $3$ only in the term that contains $a^{2}_{-+}$.
By \eqref{eq:rotation-coefficients} with $k=-3$ and with $k=3$,
\[
(a^{-2}_{++})'=\frac{2}{r}\cdot a^{-2}_{++},\qquad(a^{2}_{-+})'=\frac{2}{r}\cdot a^{2}_{-+}.
\]
For $\sigma'=-$ the frequency $-3$ occurs only with $a^{-2}_{+-}$ and the frequency $3$ only with $a^{2}_{--}$,
and the same two values of $k$ give $(a^{-2}_{+-})'=\frac{2}{r}\cdot a^{-2}_{+-}$ and $(a^{2}_{--})'=\frac{2}{r}\cdot a^{2}_{--}$.
By \eqref{eq:rotation-radial-decay} with $\kappa=2$,
\[
a^{-2}_{++}=a^{2}_{-+}=a^{-2}_{+-}=a^{2}_{--}=0 .
\]
Applied with the two profiles exchanged, \eqref{eq:rotation-harmonic-symmetry} turns $a^{2}_{-+}=0$ into $a^{2}_{+-}=0$ and $a^{-2}_{+-}=0$ into $a^{-2}_{-+}=0$.
So the kernels $G_{(+,\chi)(-,\chi')}$ and $G_{(-,\chi)(+,\chi')}$ have the frequency $0$ only.
\item[(iii)]
Take $\sigma'=3$.
The three terms of \eqref{eq:rotation-bianchi} carry the frequencies $\{0,-2\}$, coming from $a^{1}_{+3}$ and $a^{-1}_{+3}$,
the frequencies $\{2,0\}$, coming from $a^{1}_{-3}$ and $a^{-1}_{-3}$,
and, by (i), the frequency $\{0\}$.
The coefficients of $\exp(-2i\vartheta)$ and of $\exp(2i\vartheta)$, \eqref{eq:rotation-coefficients} with $k=-2$ and $k=2$, give
\[
(a^{-1}_{+3})'=\frac{1}{r}\cdot a^{-1}_{+3},\qquad(a^{1}_{-3})'=\frac{1}{r}\cdot a^{1}_{-3}.
\]
By \eqref{eq:rotation-radial-decay} with $\kappa=1$ both functions vanish,
and \eqref{eq:rotation-harmonic-symmetry} gives $a^{-1}_{3+}=a^{1}_{3-}=0$ as well.
\item[(iv)]
Take $\sigma'=+$ once more.
Three terms contribute to the frequency $-1$:
the term $\frac12\exp(-i\vartheta)\cdot(a^0_{++})'$ that $\partial_z$ produces from $a^0_{++}$;
the term $\partial_{\bar z}\left(a^{-2}_{-+}\cdot\exp(-2i\vartheta)\right)$, which vanishes by (ii);
and the term $a^{-1}_{3+}\cdot\exp(-i\vartheta)$, which vanishes by (iii).
Hence $(a^0_{++})'=0$, and \eqref{eq:rotation-radial-decay} with $\kappa=0$ gives $a^0_{++}=0$.
For $\sigma'=-$ the frequency $1$ receives the term $\frac12\exp(i\vartheta)\cdot(a^0_{--})'$ from $\partial_{\bar z}a^0_{--}$,
together with $\partial_z\left(a^{2}_{+-}\cdot\exp(2i\vartheta)\right)=0$ and $a^{1}_{3-}\cdot\exp(i\vartheta)=0$;
so $(a^0_{--})'=0$ and $a^0_{--}=0$.
\end{enumerate}

In both cases only the frequency $k=w$ survives in \eqref{eq:rotation-window}:
\[
G_{\lambda\lambda'}(rn_\vartheta)=a^{w}_{\lambda\lambda'}(r)\cdot\exp(iw\vartheta),\qquad
\left|a^{w}_{\lambda\lambda'}(r)\right|\le\frac{72}{\pi}\cdot\|\chi\|_2\cdot\|\chi'\|_2\cdot r^{-2}\qquad(r>0,\ \vartheta\in\R).
\]
Together with the density formula of Step~1, this is \eqref{it:two-point-frequencies}, with $a_{\lambda\lambda'}:=a^{w}_{\lambda\lambda'}$.
As $R_t(rn_\vartheta)=rn_{\vartheta+t}$, it also gives
\begin{equation}\label{eq:rotation-phase}
G_{\lambda\lambda'}(R_ty)=\exp(iwt)\cdot G_{\lambda\lambda'}(y)\qquad(t\in\R,\ y\in\R^2\setminus\{0\}).
\end{equation}

\emph{Step 5. From plates to $L^2$.}
For a real number $t$ define, for $b,b'\in L^2(\R^3;\C^3)$,
\[
D(b,b'):=S_2(R_t\cdot b,R_t\cdot b')-S_2(b,b').
\]
The map $b\mapsto R_t\cdot b$ is unitary on $L^2(\R^3;\C^3)$, since it rotates both the argument and the values of $b$,
and $S_2$ is a bounded bilinear form (Lemma~\ref{lem:domination-consequences}\eqref{it:moment-forms});
so $D$ is a bounded bilinear form on $L^2(\R^3;\C^3)$.

Consider first two plates $\alpha\otimes\lambda$ and $\alpha'\otimes\lambda'$ with profiles in $C^\infty_{\mathrm c}(\R;\R)$
whose plane parts $\alpha,\alpha'\in L^2(\R^2;\C)$ have disjoint compact supports.
The open set $U:=\{(y,y')\in(\R^2)^2:\ y\neq y'\}$ is invariant under $R_t$.
By Step~1 the continuous function $(y,y')\mapsto G_{\lambda\lambda'}(y'-y)$ on $U$ is a density of $S_2$ on such pairs of plates,
and by \eqref{eq:rotation-phase} it acquires the factor $\exp(iwt)$ when $(y,y')$ is replaced by $(R_ty,R_ty')$,
$w=w(\sigma)+w(\sigma')$ being the weight of the pair of labels.
The criterion \eqref{eq:rotation-density-criterion} therefore gives $D(\alpha\otimes\lambda,\alpha'\otimes\lambda')=0$.

Let now $K$ and $K'$ be disjoint compact subsets of $\R^2$.
By Lemma~\ref{lem:plates-total}\eqref{it:plates-total-smooth}, the finite sums of plates with profiles in $C^\infty_{\mathrm c}(\R;\R)$ whose plane parts vanish outside $K$
are dense in the space of the fields $b\in L^2(\R^3;\C^3)$ that vanish almost everywhere outside $K\times\R$, and likewise for $K'$.
Plates of this kind with plane parts vanishing outside $K$ and outside $K'$ respectively have disjoint compact plane supports.
Since $D$ is bilinear and bounded, it follows that $D(b,b')=0$ whenever $b$ vanishes outside $K\times\R$ and $b'$ vanishes outside $K'\times\R$.
This is property \eqref{it:contact-disjoint} of Proposition~\ref{prop:no-contact-terms} for the open set $\R^2$ and $N=2$.

Property \eqref{it:contact-gradients} of that proposition holds as well.
Let $u\in C^\infty_{\mathrm c}(\R^3;\C)$.
Since $\det R_t=1$, \eqref{eq:rotation-gradient-action} gives $R_t\cdot\nabla u=\nabla(u\circ R_{-t})$,
and $u\circ R_{-t}$ again lies in $C^\infty_{\mathrm c}(\R^3;\C)$.
By \eqref{eq:rotation-gradient-fields} the field vanishes on the gradients of $u$ and of $u\circ R_{-t}$,
so that $S_2(\nabla u,b')=0$ and $S_2(R_t\cdot\nabla u,R_t\cdot b')=S_2(\nabla(u\circ R_{-t}),R_t\cdot b')=0$ for every $b'\in L^2(\R^3;\C^3)$;
by the symmetry of $S_2$ the same holds with $\nabla u$ in the second entry.
Thus $D$ vanishes as soon as one of its entries is $\nabla u$.

Proposition~\ref{prop:no-contact-terms} now yields $D=0$.
As $t$ was arbitrary, $S_2(R_t\cdot b,R_t\cdot b')=S_2(b,b')$ for every real $t$ and all $b,b'\in L^2(\R^3;\C^3)$,
whether or not the supports of $b$ and $b'$ meet;
this is \eqref{it:two-point-invariance}.
\end{proof}

\begin{remark}
Gaussian domination \textup{(G)}, closedness \textup{(dF)}, the one-field cone $(C^1_{12})$,
and the requirement that \eqref{it:quarter-turn-two-point} or \eqref{it:second-mirror-two-point} hold,
are each necessary in Theorem~\ref{thm:two-point}\eqref{it:two-point-invariance}; see Proposition~\ref{prop:sharpness}.
For Yang--Mills theory in four dimensions, the corresponding argument disposes of the harmonic $\cos4\vartheta$ of the two-point function by means of a pair of signs:
the coordinate mirror supplies one of them, the diagonal mirror the other.
In the proof above exactly one step is of that type, namely step~(i) in the case \eqref{it:second-mirror-two-point}.
It works with two perpendicular mirrors and eliminates the harmonic of order~$2$,
that is, the frequencies $\pm2$ for a pair of labels $(3,\chi)$, $(3,\chi')$.
Every other frequency $k\neq w$ that the window \eqref{eq:rotation-window} admits is eliminated through the Bianchi identity \eqref{eq:rotation-bianchi}.
\end{remark}

\subsubsection{The K\"all\'en--Lehmann form}\label{sec:kallen-lehmann}

\begin{lemma}[Gaussian laws]\label{lem:gaussian-covariance}
Let $C$ be a real symmetric bilinear form on $\mathcal{S}(\R^3;\R^3)$ that is continuous for the topology of $\mathcal{S}(\R^3;\R^3)$ and satisfies $C(b,b)\ge0$ for every $b\in\mathcal{S}(\R^3;\R^3)$.
\begin{enumerate}
\item\label{it:gaussian-laws-criteria}
There is exactly one law $\gamma_C$ on $(\mathcal{S}',\Sigma)$ with
\[
\gamma_C\left(\exp\left(i\,\mathbb{B}(b)\right)\right)=\exp\left(-\frac{C(b,b)}{2}\right)\qquad\left(b\in\mathcal{S}(\R^3;\R^3)\right),
\]
and under $\gamma_C$ the random variables $\mathbb{B}(b)$, $b\in\mathcal{S}(\R^3;\R^3)$, are jointly Gaussian and centred, with the covariance $C$.
The law $\gamma_C$ satisfies \textup{(G)} if and only if $C(b,b)\le\|b\|^2$ for every $b\in\mathcal{S}(\R^3;\R^3)$.
It satisfies \textup{(dF)} if and only if $C(\nabla u,\nabla u)=0$ for every $u\in\mathcal{S}(\R^3;\R)$.
For a Euclidean map $\mathfrak{g}$ of $\R^3$ one has $\mathfrak{g}_*\gamma_C=\gamma_C$ if and only if $C(\mathfrak{g}\cdot b,\mathfrak{g}\cdot b)=C(b,b)$ for every $b\in\mathcal{S}(\R^3;\R^3)$.
\item\label{it:gaussian-laws-positivity}
Let $j\in\{1,2,3\}$, and let $C$ be invariant under $\theta_j$, that is, $C(\theta_j\cdot b,\theta_j\cdot b')=C(b,b')$ for all $b,b'\in\mathcal{S}(\R^3;\R^3)$.
Then $\gamma_C$ satisfies $(RP_j)$ if and only if $C(\theta_j\cdot b,b)\ge0$ for every $b\in C^{\infty}_{\mathrm c}(\{x_j>0\};\R^3)$.
\end{enumerate}
\end{lemma}

\begin{proof}
\eqref{it:gaussian-laws-criteria}
The map $f\mapsto *f$ is a linear homeomorphism of $\mathcal{S}(\R^3;\Lambda^2)$ onto $\mathcal{S}(\R^3;\R^3)$.
Hence the function $f\mapsto\exp\left(-\frac12C(*f,*f)\right)$ on $\mathcal{S}(\R^3;\Lambda^2)$ is continuous, and its value at $f=0$ is $1$.
It is also positive definite.
To see this, let $f_1,\dots,f_n\in\mathcal{S}(\R^3;\Lambda^2)$, put $b_p:=*f_p$, and let $c\in\C^n$.
The real symmetric matrix $\left(C(b_p,b_q)\right)_{p,q}$ is positive semidefinite, because
\[
\sum_{p,q}t_p\cdot t_q\cdot C(b_p,b_q)=C\left(\sum_pt_p\cdot b_p,\sum_pt_p\cdot b_p\right)\ge0\qquad(t_1,\dots,t_n\in\R).
\]
By the Schur product theorem all its entrywise powers are positive semidefinite, the power $0$ being the matrix with all entries $1$;
summing the exponential series, one finds that the entrywise exponential $\left(\exp C(b_p,b_q)\right)_{p,q}$ is positive semidefinite as well.
As $C(b_q-b_p,b_q-b_p)=C(b_p,b_p)+C(b_q,b_q)-2\,C(b_p,b_q)$, it follows, with $c'_p:=c_p\cdot\exp\left(-\frac12C(b_p,b_p)\right)$, that
\[
\sum_{p,q}\overline{c_p}\cdot c_q\cdot\exp\left(-\frac{C(b_q-b_p,b_q-b_p)}{2}\right)=\sum_{p,q}\overline{c'_p}\cdot c'_q\cdot\exp\left(C(b_p,b_q)\right)\ge0.
\]
By Minlos' theorem \cite[Section~3.4]{GlimmJaffe}, a continuous positive definite function on $\mathcal{S}(\R^3;\Lambda^2)$ with the value $1$ at $0$ is the characteristic functional of a law on $(\mathcal{S}',\Sigma)$;
and a law on $(\mathcal{S}',\Sigma)$ is determined by its characteristic functional \cite[Lemma~8.12\,(ii)]{NoncompactQED}.
Hence there is exactly one law $\gamma_C$ with $\gamma_C\left(\exp(iF(f))\right)=\exp\left(-\frac12C(*f,*f)\right)$ for all $f\in\mathcal{S}(\R^3;\Lambda^2)$;
since $\mathbb{B}(b)=F(f)$ for $b=*f$, this is the displayed identity.

For real $b_1,\dots,b_n\in\mathcal{S}(\R^3;\R^3)$ and $t\in\R^n$, the linearity of $b\mapsto\mathbb{B}(b)$ gives
\[
\gamma_C\left(\exp\left(i\sum_pt_p\cdot\mathbb{B}(b_p)\right)\right)=\exp\left(-\frac12\sum_{p,q}t_p\cdot t_q\cdot C(b_p,b_q)\right),
\]
which is the characteristic function of a centred Gaussian vector with the covariance matrix $\left(C(b_p,b_q)\right)_{p,q}$.
In particular, for $b,b'\in\mathcal{S}(\R^3;\R^3)$,
\[
\gamma_C\left(\mathbb{B}(b)\cdot\mathbb{B}(b')\right)=C(b,b'),\qquad\gamma_C\left(\exp\mathbb{B}(b)\right)=\exp\left(\frac{C(b,b)}{2}\right).
\]
Since $*$ is an isometry, \textup{(G)} asks that $\gamma_C(\exp\mathbb{B}(b))\le\exp\left(\frac12\|b\|^2\right)$ for every $b\in\mathcal{S}(\R^3;\R^3)$;
by the second identity this holds if and only if $C(b,b)\le\|b\|^2$ for all such $b$.

If $\gamma_C$ satisfies \textup{(dF)}, then $\mathbb{B}(\nabla u)=0$ almost surely for $u\in\mathcal{S}(\R^3;\R)$, by \eqref{eq:rotation-gradient-fields};
so $C(\nabla u,\nabla u)=\gamma_C\left(\mathbb{B}(\nabla u)^2\right)=0$.
Conversely, let $C(\nabla u,\nabla u)=0$ for every $u\in\mathcal{S}(\R^3;\R)$.
Then $\gamma_C\left(F(f)^2\right)=0$ for every $f\in\mathcal{S}(\R^3;\Lambda^2)$ with $*f=\nabla u$, $u\in\mathcal{S}(\R^3;\R)$,
and a law with this property is carried by the closed distributions \cite[Section~8.1.6]{NoncompactQED}.

Let $\mathfrak{g}$ be a Euclidean map.
For $b\in\mathcal{S}(\R^3;\R^3)$ the field $\mathfrak{g}\cdot b$ lies in $\mathcal{S}(\R^3;\R^3)$, and $\hat{\mathfrak{g}}\,\mathbb{B}(b)=\mathbb{B}(\mathfrak{g}\cdot b)$ by \eqref{eq:rotation-field-action}.
Therefore
\[
\begin{aligned}
\left((\mathfrak{g}^{-1})_*\gamma_C\right)\left(\exp\left(i\,\mathbb{B}(b)\right)\right)
&=\gamma_C\left(\exp\left(i\,\mathbb{B}(\mathfrak{g}\cdot b)\right)\right)\\
&=\exp\left(-\frac{C(\mathfrak{g}\cdot b,\mathfrak{g}\cdot b)}{2}\right).
\end{aligned}
\]
A law is determined by its characteristic functional \cite[Lemma~8.12\,(ii)]{NoncompactQED}.
So $(\mathfrak{g}^{-1})_*\gamma_C=\gamma_C$, which is equivalent to $\mathfrak{g}_*\gamma_C=\gamma_C$, holds if and only if $C(\mathfrak{g}\cdot b,\mathfrak{g}\cdot b)=C(b,b)$ for every $b\in\mathcal{S}(\R^3;\R^3)$.

\eqref{it:gaussian-laws-positivity}
Applied to $\mathfrak{g}=\theta_j$, part~\eqref{it:gaussian-laws-criteria} shows that $\gamma_C$ satisfies $(R_j)$.
Let $b_1,\dots,b_n\in C^{\infty}_{\mathrm c}(\{x_j>0\};\R^3)$, and put $M_{pq}:=C(\theta_j\cdot b_p,b_q)$.
The invariance of $C$ under $\theta_j$ and the identity $\theta_j\cdot(\theta_j\cdot b)=b$ give $C(\theta_j\cdot b_p,\theta_j\cdot b_p)=C(b_p,b_p)$ and $M_{pq}=C(b_p,\theta_j\cdot b_q)=M_{qp}$.
Thus the matrix $M:=(M_{pq})_{p,q}$ is real and symmetric, and part~\eqref{it:gaussian-laws-criteria} gives
\[
\begin{aligned}
\gamma_C\left(\exp\left(i\,\mathbb{B}(b_q)-i\,\mathbb{B}(\theta_j\cdot b_p)\right)\right)
&=\exp\left(-\frac{C(b_q-\theta_j\cdot b_p,\,b_q-\theta_j\cdot b_p)}{2}\right)\\
&=\exp\left(-\frac{C(b_p,b_p)}{2}\right)\cdot\exp\left(-\frac{C(b_q,b_q)}{2}\right)\cdot\exp(M_{pq}).
\end{aligned}
\]

\emph{The condition is sufficient.}
Assume that $C(\theta_j\cdot b,b)\ge0$ for every $b\in C^{\infty}_{\mathrm c}(\{x_j>0\};\R^3)$.
For real $t_1,\dots,t_n$ the field $b:=\sum_pt_p\cdot b_p$ lies in $C^{\infty}_{\mathrm c}(\{x_j>0\};\R^3)$, and $\sum_{p,q}t_p\cdot t_q\cdot M_{pq}=C(\theta_j\cdot b,b)\ge0$.
So $M$ is positive semidefinite, and so is $\left(\exp M_{pq}\right)_{p,q}$, by the Schur product theorem applied to each term of the exponential series.
For $c\in\C^n$ and $c'_p:=c_p\cdot\exp\left(-\frac12C(b_p,b_p)\right)$ it follows that
\[
\begin{aligned}
\sum_{p,q}\overline{c_p}\cdot c_q\cdot\gamma_C\left(\exp\left(i\,\mathbb{B}(b_q)-i\,\mathbb{B}(\theta_j\cdot b_p)\right)\right)
&=\sum_{p,q}\overline{c'_p}\cdot c'_q\cdot\exp(M_{pq})\\
&\ge0.
\end{aligned}
\]
For $b_p=*f_p$ with $f_p\in C^{\infty}_{\mathrm c}(\{x_j>0\};\Lambda^2)$ one has $\mathbb{B}(b_q)-\mathbb{B}(\theta_j\cdot b_p)=F(f_q)-F(\theta_j\cdot f_p)$;
so these inequalities are $(RP_j)$.

\emph{The condition is necessary.}
Let $\gamma_C$ satisfy $(RP_j)$, and let $b\in C^{\infty}_{\mathrm c}(\{x_j>0\};\R^3)$.
The random variable $\mathbb{B}(b)$ is $\Sigma_j^+$-measurable and Gaussian, so it lies in $L^2(\Sigma_j^+,\gamma_C)$;
and $\Theta_j\mathbb{B}(b)=\mathbb{B}(\theta_j\cdot b)$ by \eqref{eq:rotation-conjugation}, since $b$ is real.
As $\gamma_C$ satisfies $(R_j)$ as well, \cite[Lemma~8.16\,(3)]{NoncompactQED} gives $\gamma_C\left((\Theta_j\Phi)\cdot\Phi\right)\ge0$ for every $\Phi\in L^2(\Sigma_j^+,\gamma_C)$
(the companion paper writes $\Theta_j$ for the complex linear operator $\hat\theta_j$, so that the scalar product in that lemma is $\langle\hat\theta_j\Phi,\Phi\rangle_{\gamma_C}$, which is the left side of this inequality).
For $\Phi=\mathbb{B}(b)$ this reads
\[
C(\theta_j\cdot b,b)=\gamma_C\left(\mathbb{B}(\theta_j\cdot b)\cdot\mathbb{B}(b)\right)\ge0.
\]
\end{proof}

\begin{proof}[Proof of Theorem~\ref{thm:two-point}\eqref{it:two-point-kallen-lehmann}]
Let $\mu$ satisfy \textup{(G)}, \textup{(dF)}, $(T_{123})$, $(R_1)$, $(RP_1)$, and let $S_2(g\cdot b,g\cdot b')=S_2(b,b')$ for all $g\in\mathrm{O}(3)$ and $b,b'\in L^2(\R^3;\C^3)$.
Its reconstruction in the direction $e_1$ is $(\mathcal{H}_1,[\,\cdot\,]_1,\mathsf{H}_1,V_1,\Omega_1)$, and $\mathfrak{E}_1$ is its joint spectral measure on $[0,\infty)\times\R^2$, with the variables $(\omega,p_\perp)$, $p_\perp=(p_2,p_3)$.
Let $\varrho_{ab}:=\varrho^{(1)}_{ab}$, $a,b\in\{1,2,3\}$, be the spectral matrix of $\mu$ of Lemma~\ref{lem:spectral-matrix} for $j=1$.

\emph{Step 1 (the one-field cones).}
Let $C$ be the restriction of $S_2$ to $\mathcal{S}(\R^3;\R^3)\times\mathcal{S}(\R^3;\R^3)$.
By Lemma~\ref{lem:domination-consequences}\eqref{it:moment-forms} it is a real symmetric bilinear form.
For real $b,b'\in L^2(\R^3;\R^3)$, Lemma~\ref{lem:domination-consequences}\eqref{it:field-extension} and the Cauchy--Schwarz inequality give
\[
|S_2(b,b')|\le\|b\|\cdot\|b'\|,\qquad0\le S_2(b,b)=\mu\left(\mathbb{B}(b)^2\right)\le\|b\|^2.
\]
So $C$ is continuous for the topology of $\mathcal{S}(\R^3;\R^3)$ and satisfies $0\le C(b,b)\le\|b\|^2$;
Lemma~\ref{lem:gaussian-covariance} applies to $C$, and $\gamma_{S_2}:=\gamma_C$ denotes the law that it provides.
A Euclidean map $\mathfrak{g}:x\mapsto gx+a$ acts on fields by $\mathfrak{g}\cdot b=\tau_a\cdot(g\cdot b)$,
and $S_2$ is invariant under the translations, by $(T_{123})$, and under $\mathrm{O}(3)$, by hypothesis.
By Lemma~\ref{lem:gaussian-covariance}\eqref{it:gaussian-laws-criteria}, the law $\gamma_{S_2}$ satisfies \textup{(G)} and is invariant under every Euclidean map.
The form $C$ is invariant under $\theta_1\in\mathrm{O}(3)$.
For $b\in C^{\infty}_{\mathrm c}(\{x_1>0\};\R^3)$ the variable $\mathbb{B}(b)$ lies in $L^2(\Sigma_1^+,\mu)$ (Lemma~\ref{lem:domination-consequences}\eqref{it:field-extension}), and \eqref{eq:rotation-conjugation} gives
\[
C(\theta_1\cdot b,b)=\mu\left((\Theta_1\mathbb{B}(b))\cdot\mathbb{B}(b)\right)=\left\|[\mathbb{B}(b)]_1\right\|^2\ge0,
\]
the norm being that of $\mathcal{H}_1$.
By Lemma~\ref{lem:gaussian-covariance}\eqref{it:gaussian-laws-positivity}, $\gamma_{S_2}$ satisfies $(RP_1)$.
Lemma~\ref{lem:transport}\eqref{it:transport-all-planes} now shows that $\gamma_{S_2}$ is invariant under the reflection in every plane of $\R^3$ and reflection positive at it;
being invariant under the translations, $\gamma_{S_2}$ lies in $\mathcal{C}_9$ (Definition~\ref{def:laws}).
Theorem~\ref{thm:cone}(c), applied to $\gamma_{S_2}$ with $(j,i)=(1,2)$ and with $(j,i)=(1,3)$, gives $(C_{12})$ and $(C_{13})$ for $\gamma_{S_2}$:
the joint spectral measure of $\gamma_{S_2}$ in the direction $e_1$ vanishes on $\{\omega<|p_2|\}$ and on $\{\omega<|p_3|\}$.

The objects that belong to $\gamma_{S_2}$ carry a prime: $(\mathcal{H}_1',[\,\cdot\,]_1',\mathsf{H}_1',V_1',\Omega_1')$ is its reconstruction in the direction $e_1$, and $\mathfrak{E}_1'$ is its joint spectral measure.
Both laws satisfy \textup{(G)}.
By Lemma~\ref{lem:domination-consequences}\eqref{it:moment-forms} their second moments are bounded bilinear forms on $L^2(\R^3;\C^3)$, and by Lemma~\ref{lem:gaussian-covariance}\eqref{it:gaussian-laws-criteria} they agree on $\mathcal{S}(\R^3;\R^3)\times\mathcal{S}(\R^3;\R^3)$.
As $\mathcal{S}(\R^3;\R^3)$ spans a dense subspace of $L^2(\R^3;\C^3)$ over $\C$, the two second moments agree on $L^2(\R^3;\C^3)\times L^2(\R^3;\C^3)$, and $S_2$ stands for both.
Let $b\in L^2(\R^3;\R^3)$ vanish almost everywhere on $\{x_1<0\}$, and put $\psi:=[\mathbb{B}(b)]_1$ and $\psi':=[\mathbb{B}(b)]_1'$.
For $t\ge0$ and $y\in e_1^\perp$, \cite[Proposition~8.19]{NoncompactQED}, Lemma~\ref{lem:domination-consequences}\eqref{it:field-extension} and \eqref{eq:rotation-conjugation} give
\[
\begin{aligned}
\left\langle\psi,\exp(-t\cdot\mathsf{H}_1)\,V_1(y)\,\psi\right\rangle
&=\mu\left((\Theta_1\mathbb{B}(b))\cdot U(t\cdot e_1+y)\,\mathbb{B}(b)\right)\\
&=S_2\left(\theta_1\cdot b,\tau_{t\cdot e_1+y}\cdot b\right),
\end{aligned}
\]
and the same computation for $\gamma_{S_2}$ gives $\langle\psi',\exp(-t\cdot\mathsf{H}_1')\,V_1'(y)\,\psi'\rangle=S_2(\theta_1\cdot b,\tau_{t\cdot e_1+y}\cdot b)$.
By Lemma~\ref{lem:joint-spectrum} the two left sides are
\[
\int\mathrm{d}\varrho_\psi(\omega,p_\perp)\,\exp(-t\omega+ip_\perp\cdot y)\qquad\text{and}\qquad\int\mathrm{d}\varrho'_{\psi'}(\omega,p_\perp)\,\exp(-t\omega+ip_\perp\cdot y),
\]
where $\varrho_\psi:=\langle\psi,\mathfrak{E}_1(\cdot)\,\psi\rangle$ and $\varrho'_{\psi'}:=\langle\psi',\mathfrak{E}_1'(\cdot)\,\psi'\rangle$.
These finite positive measures on $[0,\infty)\times\R^2$ have therefore the same Laplace--Fourier transform, and Lemma~\ref{lem:laplace-fourier} gives
\begin{equation}\label{eq:rotation-one-field-spectra}
\left\langle\psi,\mathfrak{E}_1(A)\,\psi\right\rangle=\left\langle\psi',\mathfrak{E}_1'(A)\,\psi'\right\rangle\qquad\text{for every Borel set }A\subset[0,\infty)\times\R^2.
\end{equation}
For $A=\{\omega<|p_i|\}$, $i=2,3$, the right side vanishes, by $(C_{12})$ and $(C_{13})$ for $\gamma_{S_2}$;
hence $\|\mathfrak{E}_1(\{\omega<|p_i|\})\,[\mathbb{B}(b)]_1\|^2=0$.
Since $[\mathbb{B}(b'+ib'')]_1=[\mathbb{B}(b')]_1+i\,[\mathbb{B}(b'')]_1$, the vectors $[\mathbb{B}(b)]_1$ with real $b$ as above span a dense subspace of the one-field space $\mathcal{H}^1_1$ of $\mu$ (Definition~\ref{def:closedness-cones}).
So $\mathfrak{E}_1(\{\omega<|p_i|\})$ vanishes on $\mathcal{H}^1_1$ for $i=2,3$: \emph{the law $\mu$ satisfies $(C^1_{12})$ and $(C^1_{13})$.}
As $\mathfrak{E}_1$ is a projection-valued measure, $\mathfrak{E}_1$ of the union $\{\omega<|p_2|\}\cup\{\omega<|p_3|\}$ vanishes on $\mathcal{H}^1_1$ as well, and this union is the complement of
\[
K:=\left\{(\omega,p_\perp)\in[0,\infty)\times\R^2:\ \omega\ge|p_2|,\ \omega\ge|p_3|\right\}.
\]
By Lemma~\ref{lem:spectral-matrix}\eqref{it:spectral-matrix-support} all measures $\varrho_{ab}$ are carried by $K$.

\emph{Step 2 (size).}
The Cauchy density $P_1(v)=\pi^{-1}(1+v^2)^{-1}$ satisfies $\int\mathrm{d}v\,P_1(v)\cdot\exp(ipv)=\exp(-|p|)$ for real $p$, and $\|P_1\|_2^2=(2\pi)^{-1}$.
For $c>0$ let
\[
u_c(x):=c^{-3}\cdot1_{(c,2c)}(x_1)\cdot P_1(x_2/c)\cdot P_1(x_3/c)\qquad(x\in\R^3).
\]
The function $u_c$ lies in $L^1\cap L^2(\R^3;\C)$ and vanishes on $\{x_1<0\}$.
With $x_\perp:=x-x_1\cdot e_1$, its Laplace--Fourier transform for $j=1$ is
\[
\begin{aligned}
\tilde u_c(\omega,p_\perp)
&=\int\mathrm{d}x\,u_c(x)\cdot\exp(-x_1\omega+ip_\perp\cdot x_\perp)\\
&=\left(\frac1c\cdot\int_c^{2c}\mathrm{d}s\,\exp(-s\omega)\right)\cdot\exp\left(-c|p_2|-c|p_3|\right);
\end{aligned}
\]
it is at least $\exp(-2c\omega-c|p_2|-c|p_3|)$, hence at least $\exp(-4c\omega)$ on $K$.
Moreover
\[
\|u_c\|^2=c^{-6}\cdot c\cdot\left(\frac{c}{2\pi}\right)^2=\frac{1}{4\pi^2c^3}.
\]
The measures $\varrho_{aa}$ are positive (Lemma~\ref{lem:spectral-matrix}\eqref{it:spectral-matrix-positive}) and carried by $K$.
So Lemma~\ref{lem:spectral-matrix}\eqref{it:spectral-matrix-formula}, the contractivity of $[\,\cdot\,]_1$ \cite[Proposition~8.19\,(i)]{NoncompactQED} and Lemma~\ref{lem:domination-consequences}\eqref{it:field-extension} give
\[
\begin{aligned}
\int\mathrm{d}\varrho_{aa}\,\exp(-8c\omega)
&\le\int\mathrm{d}\varrho_{aa}\,|\tilde u_c|^2=\left\|[\mathbb{B}_a(u_c)]_1\right\|^2\\
&\le\mu\left(|\mathbb{B}_a(u_c)|^2\right)\le\|u_c\|^2=\frac{1}{4\pi^2c^3}.
\end{aligned}
\]
With $t=8c$ this becomes
\begin{equation}\label{eq:rotation-size}
I(t):=\max_{a\in\{1,2,3\}}\int\mathrm{d}\varrho_{aa}(\omega,p_\perp)\,\exp(-t\omega)\le\frac{128}{\pi^2t^3}<13\,t^{-3}\qquad(t>0).
\end{equation}

\emph{Step 3 (the kernel).}
Read $\theta_1$ as the matrix $\diag(-1,1,1)$; the diagonal entries of $-\theta_1$ are $(-\theta_1)_{11}=1$ and $(-\theta_1)_{22}=(-\theta_1)_{33}=-1$.
For $x\in\R^3$ with $x_1>0$ put
\[
S_{ab}(x):=(-\theta_1)_{aa}\cdot\int\mathrm{d}\varrho_{ab}(\omega,p_\perp)\,\exp\left(-x_1\omega+i(p_2x_2+p_3x_3)\right)\qquad(a,b\in\{1,2,3\}).
\]
By Lemma~\ref{lem:spectral-matrix}\eqref{it:spectral-matrix-positive} and \eqref{eq:rotation-size},
\[
\begin{aligned}
|S_{ab}(x)|
&\le\int\mathrm{d}|\varrho_{ab}|\,\exp(-x_1\omega)\\
&\le\left(\int\mathrm{d}\varrho_{aa}\,\exp(-x_1\omega)\right)^{\frac12}\cdot\left(\int\mathrm{d}\varrho_{bb}\,\exp(-x_1\omega)\right)^{\frac12}\le I(x_1).
\end{aligned}
\]
The same integral defines $S_{ab}(z)$ for $z\in\C^3$ with $s:=\Re z_1-|\Im z_2|-|\Im z_3|>0$.
Indeed, $|p_2|\le\omega$ and $|p_3|\le\omega$ on $K$, so that there the modulus of the integrand is
\[
\exp\left(-\omega\Re z_1-p_2\Im z_2-p_3\Im z_3\right)\le\exp(-s\omega),
\]
while $\int\mathrm{d}|\varrho_{ab}|\,\exp(-s\omega)\le I(s)<\infty$.
The integrand is entire in $z$, and this domination is locally uniform on the open set $\{z\in\C^3:\Re z_1>|\Im z_2|+|\Im z_3|\}$;
by dominated convergence and by the theorems of Fubini and Morera in each variable, $S_{ab}$ is holomorphic there.
This open set contains the real points with $x_1>0$; hence $S_{ab}$ is real analytic on $\{x\in\R^3:x_1>0\}$.

Two functions or fields $b,b'$ with compact supports are \emph{ordered along} a unit vector $n$ if $\sup\{n\cdot x:x\in\supp b\}<\inf\{n\cdot x:x\in\supp b'\}$.
Let $\alpha,\alpha'\in L^2(\R^3;\C)$ have compact supports ordered along $e_1$.
Then
\[
S_2(\alpha\cdot e_a,\alpha'\cdot e_b)=\int\mathrm{d}x\int\mathrm{d}x'\,\alpha(x)\cdot\alpha'(x')\cdot S_{ab}(x'-x).
\]
Both sides are unchanged if $\alpha$ and $\alpha'$ are translated by the same vector, the left side by $(T_{123})$ and the right side by a change of variables.
So one may assume that $\supp\alpha\subset\{x_1<0\}$ and $\supp\alpha'\subset\{x_1>0\}$;
let $\delta>0$ be such that $x'_1-x_1\ge\delta$ for $x\in\supp\alpha$ and $x'\in\supp\alpha'$.
The function $\alpha^\theta:=\overline{\alpha}\circ\theta_1$ has its support in $\{x_1>0\}$; like $\alpha'$, it lies in $L^1\cap L^2(\R^3;\C)$ and vanishes on $\{x_1<0\}$.
Moreover $\theta_1\cdot(\alpha^\theta\cdot e_a)=(-\theta_1)_{aa}\cdot\overline{\alpha}\cdot e_a$, so \eqref{eq:rotation-conjugation} gives $\Theta_1\mathbb{B}_a(\alpha^\theta)=(-\theta_1)_{aa}\cdot\mathbb{B}_a(\alpha)$.
By \eqref{eq:rotation-spectral-matrix} it follows that
\[
\begin{aligned}
S_2(\alpha\cdot e_a,\alpha'\cdot e_b)
&=(-\theta_1)_{aa}\cdot\mu\left((\Theta_1\mathbb{B}_a(\alpha^\theta))\cdot\mathbb{B}_b(\alpha')\right)\\
&=(-\theta_1)_{aa}\cdot\left\langle[\mathbb{B}_a(\alpha^\theta)]_1,[\mathbb{B}_b(\alpha')]_1\right\rangle\\
&=(-\theta_1)_{aa}\cdot\int\mathrm{d}\varrho_{ab}\,\overline{\widetilde{\alpha^\theta}}\cdot\widetilde{\alpha'}.
\end{aligned}
\]
The substitution $x\mapsto\theta_1x$ gives
\[
\begin{aligned}
\overline{\widetilde{\alpha^\theta}(\omega,p_\perp)}
&=\int\mathrm{d}x\,\alpha(\theta_1x)\cdot\exp(-x_1\omega-ip_\perp\cdot x_\perp)\\
&=\int\mathrm{d}x\,\alpha(x)\cdot\exp(x_1\omega-ip_\perp\cdot x_\perp),
\end{aligned}
\]
and therefore
\[
\begin{aligned}
&\overline{\widetilde{\alpha^\theta}(\omega,p_\perp)}\cdot\widetilde{\alpha'}(\omega,p_\perp)\\
&\qquad=\int\mathrm{d}x\int\mathrm{d}x'\,\alpha(x)\cdot\alpha'(x')\cdot\exp\left(-(x'_1-x_1)\omega+ip_\perp\cdot(x'-x)_\perp\right).
\end{aligned}
\]
On $\supp\alpha\times\supp\alpha'$ the integrand is bounded by $|\alpha(x)|\cdot|\alpha'(x')|\cdot\exp(-\delta\omega)$.
This bound is integrable for $\mathrm{d}x\,\mathrm{d}x'\,\mathrm{d}|\varrho_{ab}|$:
$\alpha$ and $\alpha'$ are integrable, having compact supports, and $\int\mathrm{d}|\varrho_{ab}|\,\exp(-\delta\omega)\le I(\delta)$ as above.
Fubini's theorem and the definition of $S_{ab}$ give the claim.

Let $S(x):=(S_{kl}(x))_{k,l}$ for $x_1>0$.
For fields $b,b'\in L^2(\R^3;\C^3)$ with compact supports ordered along $e_1$, the components $b_k$, $b'_l$ have their supports in those of $b$, $b'$;
summing the identity just proved over $k$ and $l$, one obtains
\[
S_2(b,b')=\int\mathrm{d}x\int\mathrm{d}x'\,b(x)^{T}S(x'-x)\,b'(x').
\]
Let $g\in\mathrm{O}(3)$.
As $(g\cdot b)(u)=\det(g)\cdot g\,b(g^{-1}u)$, the supports of $g\cdot b$ and $g\cdot b'$ are $g\supp b$ and $g\supp b'$, and $(ge_1)\cdot(gx)=x_1$.
Hence the fields $g\cdot b$, $g\cdot b'$ are ordered along $ge_1$, and every pair of fields with compact supports ordered along $ge_1$ arises in this way.
The hypothesis on $S_2$, the identity $b(x)=\det(g)\cdot g^{T}(g\cdot b)(gx)$ with $\det(g)^2=1$, and the substitution $u=gx$, $u'=gx'$ give
\[
\begin{aligned}
S_2(g\cdot b,g\cdot b')&=S_2(b,b')\\
&=\int\mathrm{d}u\int\mathrm{d}u'\,(g\cdot b)(u)^{T}\,g\,S\left(g^{-1}(u'-u)\right)g^{T}\,(g\cdot b')(u').
\end{aligned}
\]
Since $(g^{-1}x)_1=(ge_1)\cdot x$, the function
\[
S^{(g)}(x):=g\,S(g^{-1}x)\,g^{T}
\]
is defined and real analytic on the half space $\{x:(ge_1)\cdot x>0\}$, and it is a density of $S_2$ on the pairs of fields with compact supports ordered along $ge_1$.

Let $g,g'\in\mathrm{O}(3)$, and let $U$ be the open set of the $(x,x')\in(\R^3)^2$ with $(ge_1)\cdot(x'-x)>0$ and $(g'e_1)\cdot(x'-x)>0$.
Let $\alpha,\alpha'$ be bounded measurable functions on $\R^3$ with compact supports such that $\supp\alpha\times\supp\alpha'\subset U$.
The minimum of $(ge_1)\cdot(x'-x)$ over the compact set $\supp\alpha\times\supp\alpha'$ is $\min_{\supp\alpha'}(ge_1)\cdot x'-\max_{\supp\alpha}(ge_1)\cdot x$, and it is positive; the same holds for $g'$.
So $\alpha\cdot e_k$ and $\alpha'\cdot e_l$ are ordered along $ge_1$ and along $g'e_1$, and
\[
\int\mathrm{d}x\int\mathrm{d}x'\,\alpha(x)\cdot\alpha'(x')\cdot\left(S^{(g)}_{kl}-S^{(g')}_{kl}\right)(x'-x)=0.
\]
By Lemma~\ref{lem:analysis-facts}\eqref{it:equal-densities}, applied on $U$, the functions $(x,x')\mapsto S^{(g)}_{kl}(x'-x)$ and $(x,x')\mapsto S^{(g')}_{kl}(x'-x)$ coincide on $U$;
taking $x=0$, one finds $S^{(g)}=S^{(g')}$ on $\{(ge_1)\cdot x>0\}\cap\{(g'e_1)\cdot x>0\}$.
The half spaces $\{(ge_1)\cdot x>0\}$, $g\in\mathrm{O}(3)$, cover $\R^3\setminus\{0\}$.
Hence there is exactly one function $S=(S_{kl})_{k,l}$ on $\R^3\setminus\{0\}$ with $S=S^{(g)}$ on $\{(ge_1)\cdot x>0\}$ for every $g\in\mathrm{O}(3)$.
It is real analytic, and for $g=1$ one sees that it coincides on $\{x_1>0\}$ with the function $S$ defined there.

The function $S$ is a density of $S_2$ on all pairs of fields with disjoint compact supports: for $b,b'\in L^2(\R^3;\C^3)$ with disjoint compact supports,
\begin{equation}\label{eq:rotation-disjoint-density}
S_2(b,b')=\int\mathrm{d}x\int\mathrm{d}x'\,b(x)^{T}S(x'-x)\,b'(x').
\end{equation}
To see this, let $\delta>0$ be the distance of $\supp b$ and $\supp b'$ in the norm $|\cdot|_\infty$, choose $d>0$ with $3d<\delta$, and let $Q_m:=d\cdot(m+[0,1)^3)$ for $m\in\Z^3$.
These cubes are pairwise disjoint and cover $\R^3$, and finitely many of them meet $\supp b\cup\supp b'$.
Let $Q_m$ meet $\supp b$ and $Q_{m'}$ meet $\supp b'$.
If $|m_l-m'_l|\le2$ for $l=1,2,3$, then $|v-v'|_\infty<3d<\delta$ for all $v\in Q_m$ and $v'\in Q_{m'}$, which is impossible for $v\in\supp b$ and $v'\in\supp b'$.
So $|m_l-m'_l|\ge3$ for some $l$.
If $m'_l\ge m_l+3$, then $v'_l-v_l\ge2d$ for all $v$ in the closure of $Q_m$ and $v'$ in the closure of $Q_{m'}$, and the fields $1_{Q_m}\cdot b$, $1_{Q_{m'}}\cdot b'$ are ordered along $e_l$;
if $m'_l\le m_l-3$, they are ordered along $-e_l$.
Both $e_l$ and $-e_l$ are of the form $ge_1$ with $g\in\mathrm{O}(3)$.
So \eqref{eq:rotation-disjoint-density} holds for each pair $(1_{Q_m}\cdot b,1_{Q_{m'}}\cdot b')$, and by bilinearity it holds for $(b,b')$.
The integral converges absolutely, because $b$ and $b'$ are integrable and $S$ is bounded on the compact subset $\{x'-x:x\in\supp b,\ x'\in\supp b'\}$ of $\R^3\setminus\{0\}$.

The function $S$ has the following properties, for $x\in\R^3\setminus\{0\}$ and $a,b\in\{1,2,3\}$:
\begin{equation}\label{eq:rotation-kernel-covariance}
\begin{aligned}
&S(gx)=g\,S(x)\,g^{T}\quad(g\in\mathrm{O}(3)),\qquad|S_{ab}(x)|\le3\,I(|x|)<39\,|x|^{-3},\\
&S(x)\in\R^{3\times3},\qquad\sum_a\partial_aS_{ab}(x)=0.
\end{aligned}
\end{equation}

\emph{Covariance.}
Let $x\neq0$ and $g_0\in\mathrm{O}(3)$, and choose $g\in\mathrm{O}(3)$ with $(ge_1)\cdot x>0$.
Then $(g_0ge_1)\cdot(g_0x)=(ge_1)\cdot x>0$, and
\[
S(g_0x)=S^{(g_0g)}(g_0x)=g_0g\,S(g^{-1}x)\,g^{T}g_0^{T}=g_0\,S(x)\,g_0^{T}.
\]

\emph{Bound.}
Let $x\neq0$, and let $g\in\mathrm{O}(3)$ with $x=|x|\cdot ge_1$.
By the covariance, $S_{ab}(x)=\sum_{k,l}g_{ak}\cdot g_{bl}\cdot S_{kl}(|x|\cdot e_1)$.
Here $|S_{kl}(|x|\cdot e_1)|\le I(|x|)$ by the bound on $\{x_1>0\}$, and $\sum_{k,l}|g_{ak}|\cdot|g_{bl}|\le3$ by the Cauchy--Schwarz inequality, the rows of $g$ being unit vectors.
With \eqref{eq:rotation-size},
\[
|S_{ab}(x)|\le3\,I(|x|)\le\frac{384}{\pi^2|x|^3}<39\,|x|^{-3}.
\]

\emph{Reality.}
By Lemma~\ref{lem:domination-consequences}\eqref{it:moment-forms}, $S_2(b,b')=\overline{S_2(\overline{b},\overline{b'})}$.
Applying \eqref{eq:rotation-disjoint-density} to $\overline{b}$, $\overline{b'}$ and conjugating, one finds that $x\mapsto\overline{S(x)}$ is a density of $S_2$ on the pairs with disjoint compact supports as well.
Taking the fields $\alpha\cdot e_a$ and $\alpha'\cdot e_b$, where $\alpha$, $\alpha'$ are bounded measurable functions on $\R^3$ with disjoint compact supports, Lemma~\ref{lem:analysis-facts}\eqref{it:equal-densities} on the open set $\{(x,x'):x\neq x'\}$ of $(\R^3)^2$ gives $\overline{S_{ab}}=S_{ab}$.

\emph{Divergence.}
Let $u,\alpha'\in C^{\infty}_{\mathrm c}(\R^3;\C)$ have disjoint supports, and let $b\in\{1,2,3\}$.
By \textup{(dF)} and \eqref{eq:rotation-gradient-fields}, $\mathbb{B}(\nabla u)=0$ almost surely; so $S_2(\nabla u,\alpha'\cdot e_b)=0$.
The field $\nabla u=\sum_a(\partial_au)\cdot e_a$ has its support in $\supp u$, and for fixed $x'\in\supp\alpha'$ the function $x\mapsto S_{ab}(x'-x)$ is smooth on a neighbourhood of $\supp u$.
By \eqref{eq:rotation-disjoint-density} and one integration by parts in $x$, in which the integration by parts and the chain rule each change the sign once,
\[
\begin{aligned}
0&=\sum_a\int\mathrm{d}x\int\mathrm{d}x'\,(\partial_au)(x)\cdot\alpha'(x')\cdot S_{ab}(x'-x)\\
&=\int\mathrm{d}x\int\mathrm{d}x'\,u(x)\cdot\alpha'(x')\cdot\sum_a(\partial_aS_{ab})(x'-x).
\end{aligned}
\]
The function $(x,x')\mapsto\sum_a(\partial_aS_{ab})(x'-x)$ is continuous on the open set $\{(x,x'):x\neq x'\}$ of $(\R^3)^2$, and the last identity holds for all $u,\alpha'\in C^{\infty}_{\mathrm c}(\R^3;\C)$ with $\supp u\times\supp\alpha'$ contained in this set.
By Lemma~\ref{lem:analysis-facts}\eqref{it:equal-densities}, in its form for test functions in $C^{\infty}_{\mathrm c}$, this function vanishes; that is, $\sum_a\partial_aS_{ab}=0$ on $\R^3\setminus\{0\}$.

\emph{Step 4 (the potential).}
Fix $r>0$.
If $g\in\mathrm{O}(3)$ satisfies $ge_1=e_1$, then \eqref{eq:rotation-kernel-covariance} gives $S(re_1)=g\,S(re_1)\,g^{T}$,
so the matrix $S(re_1)$ commutes with every element of the stabiliser $1\oplus\mathrm{O}(2)$ of $e_1$.
A matrix with this property is diagonal, with equal second and third diagonal entries.
Indeed, its two off-diagonal blocks are a column and a row in $\C^2$ that are fixed by every element of $\mathrm{O}(2)$, in particular by $-1$, so they vanish.
Its lower $2\times2$ block commutes with a quarter turn of the plane $e_1^{\perp}$, which has the two distinct eigenvalues $i$ and $-i$, with the eigenlines $\C(1,i)$ and $\C(1,-i)$;
hence the block maps each of these two lines into itself.
The block also commutes with the reflection $(x_2,x_3)\mapsto(x_2,-x_3)$, which exchanges the two lines;
so its eigenvalues on them coincide, and the block is a multiple of the identity.

Let $a_L(r):=S_{11}(re_1)$ and $a_T(r):=S_{22}(re_1)=S_{33}(re_1)$, so that $S(re_1)=\diag\left(a_L(r),a_T(r),a_T(r)\right)$.
Given $x\neq0$, put $r:=|x|$ and choose $g\in\mathrm{O}(3)$ with $ge_1=x/r$.
By \eqref{eq:rotation-kernel-covariance},
$S(x)=g\,S(re_1)\,g^{T}=a_T(r)\cdot1+\left(a_L(r)-a_T(r)\right)\cdot r^{-2}xx^{T}$, that is,
\begin{equation}\label{eq:rotation-isotropic-kernel}
S_{ab}(x)=a_T(r)\cdot\delta_{ab}+\left(a_L(r)-a_T(r)\right)\cdot\frac{x_ax_b}{r^2}\qquad(r=|x|>0).
\end{equation}
The functions $a_L$ and $a_T$ are real valued and real analytic on $(0,\infty)$, because $S$ is a real, real analytic matrix function on $\R^3\setminus\{0\}$;
by \eqref{eq:rotation-kernel-covariance}, $|a_L(r)|<39\,r^{-3}$ and $|a_T(r)|<39\,r^{-3}$.
For $t>0$ the point $te_1$ lies in $\{x_1>0\}$, where $S_{11}$ is given by the integral formula of Step~3 with the sign $(-\theta_1)_{11}=1$; thus
\[
a_L(t)=S_{11}(te_1)=\int\mathrm{d}\varrho_{11}(\omega,p_\perp)\,\exp(-t\omega)\ge0,
\]
since $\varrho_{11}$ is a positive measure by Lemma~\ref{lem:spectral-matrix}\eqref{it:spectral-matrix-positive}.

The divergence of $S$ can be read off from \eqref{eq:rotation-isotropic-kernel}.
The derivative of $x\mapsto a_T(|x|)$ in the direction $e_a$ is $a_T'(r)\cdot x_a/r$, and likewise for $a_L$;
moreover $\sum_a\partial_a\left(x_ax_b/r^2\right)=4x_b/r^2-2x_b/r^2=2x_b/r^2$.
Hence
\[
\sum_a\partial_aS_{ab}(x)=\frac{x_b}{r}\cdot\left(a_L'(r)+\frac{2\cdot\left(a_L(r)-a_T(r)\right)}{r}\right).
\]
The left side vanishes by \eqref{eq:rotation-kernel-covariance}.
Choosing $x$ with $x_b\neq0$, one finds
\[
a_L'(r)=-\frac{2\cdot\left(a_L(r)-a_T(r)\right)}{r}\qquad(r>0).
\]

The \emph{potential} is the function
\[
P(r):=\int_r^\infty\mathrm{d}s\,\frac{s}{2}\cdot a_L(s)\qquad(r>0).
\]
As $0\le a_L(s)<39\,s^{-3}$, the integral converges, and $0\le P(r)\le39/(2r)\le20/r$.
The function $P$ is twice differentiable, with
\[
P'=-\frac{r\cdot a_L}{2},\qquad P''=-\frac{a_L}{2}-\frac{r\cdot a_L'}{2}=\frac{a_L}{2}-a_T,
\]
where the last equality is the equation for $a_L'$; hence
\[
-P''-\frac{P'}{r}=a_T,\qquad P''-\frac{P'}{r}=a_L-a_T.
\]
For the radial function $x\mapsto P(|x|)$ on $\R^3\setminus\{0\}$ one has, with $r=|x|$,
\[
\begin{aligned}
\partial_a\partial_bP(|x|)&=P''(r)\cdot\frac{x_ax_b}{r^2}+\frac{P'(r)}{r}\cdot\left(\delta_{ab}-\frac{x_ax_b}{r^2}\right),\\
\Delta P(|x|)&=P''(r)+\frac{2\cdot P'(r)}{r},
\end{aligned}
\]
so that
\[
\begin{aligned}
\left(-\delta_{ab}\Delta+\partial_a\partial_b\right)P(|x|)
&=\left(-P''(r)-\frac{P'(r)}{r}\right)\cdot\delta_{ab}\\
&\qquad+\left(P''(r)-\frac{P'(r)}{r}\right)\cdot\frac{x_ax_b}{r^2}.
\end{aligned}
\]
Comparison with \eqref{eq:rotation-isotropic-kernel} yields
\begin{equation}\label{eq:rotation-potential}
S_{ab}(x)=\left(-\delta_{ab}\Delta+\partial_a\partial_b\right)P(|x|)\qquad(x\neq0).
\end{equation}

\emph{Step 5 (a completely monotone function).}
Let $c$ be the function on $(0,\infty)$ defined by
\[
c(t):=-\int_{\R^2}\mathrm{d}x_\perp\,(S_{22}+S_{33})(t,x_\perp),
\]
where $(t,x_\perp)$ stands for the point $te_1+x_\perp$ of $\R^3$.
By \eqref{eq:rotation-kernel-covariance} the modulus of the integrand is less than $78\,(t^2+|x_\perp|^2)^{-\frac{3}{2}}$, an integrable function on $\R^2$;
so the integral converges absolutely.

\emph{(i) The function $c$ is the Laplace transform of a positive measure on $[0,\infty)$.}
For $\varepsilon>0$ and $t>0$ put
\[
c_\varepsilon(t):=-\int_{\R^2}\mathrm{d}x_\perp\,\exp\left(-\frac{\varepsilon|x_\perp|^2}{2}\right)\cdot(S_{22}+S_{33})(t,x_\perp).
\]
Since $(-\theta_1)_{22}=(-\theta_1)_{33}=-1$, the integral formula of Step~3 reads, on the plane $\{x_1=t\}$,
\[
-(S_{22}+S_{33})(t,x_\perp)=\int\mathrm{d}(\varrho_{22}+\varrho_{33})(\omega,p_\perp)\,\exp\left(-t\omega+ip_\perp\cdot x_\perp\right).
\]
Here $\varrho_{22}+\varrho_{33}$ is a positive measure (Lemma~\ref{lem:spectral-matrix}\eqref{it:spectral-matrix-positive}),
and $\int\mathrm{d}(\varrho_{22}+\varrho_{33})(\omega,p_\perp)\,\exp(-t\omega)\le2\cdot I(t)<\infty$ by \eqref{eq:rotation-size}.
The double integral that results for $c_\varepsilon(t)$ is therefore absolutely convergent.
Fubini's theorem, together with
\[
\int_{\R^2}\mathrm{d}x_\perp\,\exp\left(-\frac{\varepsilon|x_\perp|^2}{2}+ip_\perp\cdot x_\perp\right)=\frac{2\pi}{\varepsilon}\cdot\exp\left(-\frac{|p_\perp|^2}{2\varepsilon}\right),
\]
yields
\[
c_\varepsilon(t)=\int\mathrm{d}(\varrho_{22}+\varrho_{33})(\omega,p_\perp)\,\exp(-t\omega)\cdot\frac{2\pi}{\varepsilon}\cdot\exp\left(-\frac{|p_\perp|^2}{2\varepsilon}\right).
\]
So $c_\varepsilon$ is the Laplace transform, finite at every $t>0$, of a positive measure on $[0,\infty)$:
the image under $(\omega,p_\perp)\mapsto\omega$ of the measure whose density with respect to $\varrho_{22}+\varrho_{33}$ is $(\omega,p_\perp)\mapsto\frac{2\pi}{\varepsilon}\cdot\exp\left(-\frac{|p_\perp|^2}{2\varepsilon}\right)$.
As $\varepsilon\to0$, dominated convergence, with the integrable majorant $|S_{22}+S_{33}|(t,\cdot)$, gives $c_\varepsilon(t)\to c(t)$ for every $t>0$.
By the extended continuity theorem for Laplace transforms \cite[Section~XIII.1, Theorem~2a]{Feller1971},
a limit, finite at every point of $(0,\infty)$, of Laplace transforms of positive measures on $[0,\infty)$ is itself the Laplace transform of a positive measure on $[0,\infty)$.
Applied along a sequence $\varepsilon_n\to0$, this theorem provides a positive Borel measure $\varrho_c$ on $[0,\infty)$ with
\begin{equation}\label{eq:rotation-monotone-laplace}
c(t)=\int_{[0,\infty)}\mathrm{d}\varrho_c(m)\,\exp(-mt)\qquad(t>0).
\end{equation}
In particular $c\ge0$.

\emph{(ii) The function $c$ is $-4\pi$ times the derivative of $t\mapsto t\cdot P(t)$.}
Fix $t>0$, and on the plane $\{x_1=t\}$ put $r:=\left(t^2+|x_\perp|^2\right)^{\frac{1}{2}}$.
By \eqref{eq:rotation-isotropic-kernel},
\[
\begin{aligned}
(S_{22}+S_{33})(t,x_\perp)&=2\cdot a_T(r)+\left(a_L(r)-a_T(r)\right)\cdot\frac{|x_\perp|^2}{r^2}\\
&=(a_T+a_L)(r)-(a_L-a_T)(r)\cdot\frac{t^2}{r^2},
\end{aligned}
\]
a function of $r$ alone.
For such a function, integration over $\R^2$ in polar coordinates followed by the substitution $r^2=t^2+|x_\perp|^2$ replaces $\int_{\R^2}\mathrm{d}x_\perp$ by $2\pi\int_t^\infty\mathrm{d}r\,r$.
Both terms are absolutely integrable in this sense, since $r\cdot|a_T+a_L|(r)<78\,r^{-2}$ and $r\cdot|a_L-a_T|(r)\cdot t^2r^{-2}<78\,t^2r^{-4}$.

In the first term, the identities $a_L=-2\cdot P'/r$ and $a_T=-P''-P'/r$ of Step~4 give $r\cdot(a_T+a_L)=-r\cdot P''-3\cdot P'=-(r\cdot P')'-2\cdot P'$.
As $r\to\infty$, both $r\cdot P'(r)=-r^2\cdot a_L(r)/2$ and $P(r)$ tend to $0$ (Step~4).
Therefore
\[
\int_t^\infty\mathrm{d}r\,r\cdot(a_T+a_L)(r)=\left[-r\cdot P'(r)-2\cdot P(r)\right]_t^\infty=t\cdot P'(t)+2\cdot P(t).
\]
In the second term, $(a_L-a_T)/r=\left(P''-P'/r\right)/r=(P'/r)'$, and $P'(r)/r=-a_L(r)/2\to0$ as $r\to\infty$; hence
\[
-t^2\int_t^\infty\mathrm{d}r\,r\cdot\frac{(a_L-a_T)(r)}{r^2}=-t^2\int_t^\infty\mathrm{d}r\,\left(\frac{P'}{r}\right)'(r)=t\cdot P'(t).
\]
The two contributions add up to $2t\cdot P'(t)+2\cdot P(t)=2\cdot\left(t\cdot P\right)'(t)$.
Multiplication by $-2\pi$ gives
\begin{equation}\label{eq:rotation-monotone-derivative}
c(t)=-4\pi\cdot\frac{\mathrm{d}}{\mathrm{d}t}\left(t\cdot P(t)\right)\qquad(t>0).
\end{equation}

\emph{Step 6 (the measure).}
Let $k(t):=t\cdot P(t)$ for $t>0$.
By Step~4, $0\le k\le20$;
by \eqref{eq:rotation-monotone-derivative} and \eqref{eq:rotation-monotone-laplace}, $-k'=c/(4\pi)\ge0$.
So $k$ is non-increasing, the limit $k(\infty):=\lim_{t\to\infty}k(t)\ge0$ exists, and
\[
\int_t^\infty\mathrm{d}s\,c(s)=4\pi\cdot\left(k(t)-k(\infty)\right)<\infty\qquad(t>0).
\]
According to \eqref{eq:rotation-monotone-laplace}, $c(s)\ge\varrho_c(\{0\})$ for every $s>0$;
the finiteness of the last integral therefore forces $\varrho_c(\{0\})=0$.
Tonelli's theorem now gives, for $t>0$,
\[
\begin{aligned}
4\pi\cdot k(t)&=4\pi\cdot k(\infty)+\int_t^\infty\mathrm{d}s\int\mathrm{d}\varrho_c(m)\,\exp(-ms)\\
&=4\pi\cdot k(\infty)+\int_{(0,\infty)}\mathrm{d}\varrho_c(m)\,\frac{\exp(-mt)}{m}.
\end{aligned}
\]
Let $\delta_0$ be the unit mass at $m=0$, and let $\mathfrak{K}$ be the positive Borel measure on $[0,\infty)$ given on Borel sets $A\subset[0,\infty)$ by
\[
\mathfrak{K}(A):=4\pi\cdot k(\infty)\cdot\delta_0(A)+\int_{A\setminus\{0\}}\mathrm{d}\varrho_c(m)\,\frac{1}{m}.
\]
Then $4\pi t\cdot P(t)=\int\mathrm{d}\mathfrak{K}(m)\,\exp(-mt)$ for every $t>0$, that is,
\begin{equation}\label{eq:rotation-potential-measure}
P(r)=\int_{[0,\infty)}\mathrm{d}\mathfrak{K}(m)\,\frac{\exp(-mr)}{4\pi r}\qquad(r>0),
\end{equation}
and monotone convergence as $t\downarrow0$ gives $\mathfrak{K}([0,\infty))=\lim_{t\to0}4\pi\cdot k(t)\le80\pi<\infty$.
Inserting \eqref{eq:rotation-potential-measure} in \eqref{eq:rotation-potential}, one obtains for the kernel $S_{ab}$ the formula of Theorem~\ref{thm:two-point}\eqref{it:two-point-kallen-lehmann};
by Step~3, $S$ is real analytic on $\R^3\setminus\{0\}$ and is the density of $S_2$ on all pairs of fields with disjoint compact supports.

\emph{Step 7 (the multiplier; no contact term).}
Consider the form $(b,b')\mapsto S_2(\overline b,b')$ on $L^2(\R^3;\C^3)$.
It is antilinear in $b$, linear in $b'$ and Hermitian, because $S_2$ is symmetric and $\overline{S_2(b,b')}=S_2(\overline b,\overline{b'})$ (Lemma~\ref{lem:domination-consequences}\eqref{it:moment-forms}).
It satisfies $0\le S_2(\overline b,b)=\mu\left(|\mathbb{B}(b)|^2\right)\le\|b\|^2$ (Lemma~\ref{lem:domination-consequences}\eqref{it:field-extension}),
and by $(T_{123})$ it is invariant under all translations of $\R^3$.
By the Riesz representation of bounded sesquilinear forms, it is therefore the form of a bounded operator on $L^2(\R^3;\C^3)$ that lies between $0$ and $1$ and commutes with all translations.
Each of the nine matrix entries of this operator is a bounded operator on $L^2(\R^3;\C)$ commuting with the translations,
hence a Fourier multiplier with a bounded measurable symbol \cite[Chapter~I, Theorem~3.18]{SteinWeiss1971}.
With Plancherel's theorem for the Fourier transform $\hat b(p)=\int\mathrm{d}x\,\exp(-ip\cdot x)\cdot b(x)$, one obtains a measurable function $M$ on $\R^3$ with values in the Hermitian $3\times3$ matrices such that
\begin{equation}\label{eq:rotation-lehmann-form}
S_2(\overline b,b')=(2\pi)^{-3}\int_{\R^3}\mathrm{d}p\,\hat b(p)^{\dagger}M(p)\,\hat b'(p)\qquad\left(b,b'\in L^2(\R^3;\C^3)\right)
\end{equation}
and $0\le M(p)\le1$ for almost every $p$;
the last bounds follow from those of the form, applied to $b=u\,v$ with $u\in L^2(\R^3;\C)$ and $v$ in a countable dense subset of $\C^3$.

By (dF), $\mathbb{B}(\nabla u)=0$ almost surely for every $u\in\mathcal{S}(\R^3;\C)$, see \eqref{eq:rotation-gradient-fields};
so $S_2(\overline b,\nabla u)=0$ for all $b\in L^2(\R^3;\C^3)$.
The Fourier transform of $\nabla u$ is $p\mapsto ip\cdot\hat u(p)$; hence \eqref{eq:rotation-lehmann-form} gives $M(p)\,p\cdot\hat u(p)=0$ for almost every $p$.
Taking for $u$ a Gaussian, whose Fourier transform has no zero, one gets
\[
M(p)\,p=0\qquad\text{for almost every }p\in\R^3.
\]

Next, by Step~4 the function $x\mapsto P(|x|)$ on $\R^3\setminus\{0\}$ is non-negative and bounded by $20/|x|$;
it is locally integrable and defines a tempered distribution on $\R^3$.
Its Fourier transform is the locally integrable function $p\mapsto\int\mathrm{d}\mathfrak{K}(m)\,(|p|^2+m^2)^{-1}$.
Indeed, the Fourier transform of $x\mapsto\exp(-m|x|)/(4\pi|x|)$ is $p\mapsto(|p|^2+m^2)^{-1}\le|p|^{-2}$, a locally integrable function;
since $\mathfrak{K}$ is finite, Fubini's theorem applies to \eqref{eq:rotation-potential-measure} paired with a test function in $\mathcal{S}(\R^3;\C)$.
Consequently the tempered distribution
\[
K^{\mathfrak{K}}_{ab}:=\left(-\delta_{ab}\Delta+\partial_a\partial_b\right)P(|\cdot|)
\]
on $\R^3$ has the Fourier transform $(M_{\mathfrak{K}})_{ab}$, where, for $p\neq0$,
\begin{equation}\label{eq:rotation-lehmann-multiplier}
\begin{aligned}
M_{\mathfrak{K}}(p)&:=\left(|p|^2\cdot1-pp^{T}\right)\cdot\int\mathrm{d}\mathfrak{K}(m)\,\frac{1}{|p|^2+m^2}\\
&=P^{\perp}(p)\cdot\int\mathrm{d}\mathfrak{K}(m)\,\frac{|p|^2}{|p|^2+m^2};
\end{aligned}
\end{equation}
this is a bounded, real and even function of $p$, whose entries are bounded by $\mathfrak{K}([0,\infty))$.
The definition of the Fourier transform of a tempered distribution and the inversion formula on $\mathcal{S}(\R^3;\C)$ (Parseval's formula) give
$\langle K^{\mathfrak{K}}_{ab},u\rangle=(2\pi)^{-3}\int\mathrm{d}p\,(M_{\mathfrak{K}})_{ab}(-p)\cdot\hat u(p)$ for $u\in\mathcal{S}(\R^3;\C)$;
as $M_{\mathfrak{K}}$ is even,
\[
\langle K^{\mathfrak{K}}_{ab},u\rangle=(2\pi)^{-3}\int\mathrm{d}p\,(M_{\mathfrak{K}})_{ab}(p)\cdot\hat u(p)\qquad\left(u\in\mathcal{S}(\R^3;\C)\right).
\]
Let $K_{ab}$ be the tempered distribution $u\mapsto(2\pi)^{-3}\int\mathrm{d}p\,M_{ab}(p)\cdot\hat u(p)$ on $\R^3$.

The distribution $K_{ab}$ reproduces $S_2$ in the following way.
For $\alpha\in\mathcal{S}(\R^3;\C)$ put $\check\alpha(x):=\alpha(-x)$.
For $\alpha,\alpha'\in\mathcal{S}(\R^3;\C)$ the Fourier transform of $\check\alpha*\alpha'$ is $p\mapsto\hat\alpha(-p)\cdot\hat\alpha'(p)$, and that of $\overline\alpha$ is $p\mapsto\overline{\hat\alpha(-p)}$.
With $b:=\overline\alpha\,e_a$ and $b':=\alpha'e_b$ in \eqref{eq:rotation-lehmann-form} it follows that
\[
S_2(\alpha e_a,\alpha'e_b)=(2\pi)^{-3}\int\mathrm{d}p\,\hat\alpha(-p)\cdot M_{ab}(p)\cdot\hat\alpha'(p)=\langle K_{ab},\check\alpha*\alpha'\rangle.
\]
Now let $\alpha,\alpha'\in C^\infty_{\mathrm c}(\R^3;\C)$ have disjoint supports, and put $u:=\check\alpha*\alpha'$, so that
\[
u(x)=\int\mathrm{d}x'\,\alpha(x')\cdot\alpha'(x'+x).
\]
This function vanishes whenever $|x|$ is smaller than the distance between the two supports, in particular near $0$.
On $\R^3\setminus\{0\}$ the function $x\mapsto P(|x|)$ is smooth, and there, by \eqref{eq:rotation-potential}, the distribution $K^{\mathfrak{K}}_{ab}$ is the function $S_{ab}$.
Since $S$ is the density of $S_2$ on pairs with disjoint compact supports (Step~3),
\[
\begin{aligned}
\langle K^{\mathfrak{K}}_{ab},u\rangle&=\int\mathrm{d}x\,S_{ab}(x)\cdot u(x)
=\int\mathrm{d}x\int\mathrm{d}x'\,\alpha(x)\cdot\alpha'(x')\cdot S_{ab}(x'-x)\\
&=S_2(\alpha e_a,\alpha'e_b)=\langle K_{ab},u\rangle.
\end{aligned}
\]

Put $K'_{ab}:=K_{ab}-K^{\mathfrak{K}}_{ab}$.
Let $u\in C^\infty_{\mathrm c}(\R^3\setminus\{0\};\C)$, and let $(\alpha_n)$ be an approximate identity in $C^\infty_{\mathrm c}(\R^3)$ whose supports lie in balls about $0$ of radius less than half the distance of $\supp u$ from $0$.
Then $\supp\alpha_n$ and $\supp u$ are disjoint, and the preceding identity, applied to $\alpha_n$ and $u$, gives $\langle K'_{ab},\check\alpha_n*u\rangle=0$;
moreover $\check\alpha_n*u\to u$ in $\mathcal{S}(\R^3;\C)$.
Hence $\langle K'_{ab},u\rangle=0$, and the support of $K'_{ab}$ is contained in $\{0\}$.
A distribution supported by a single point is a finite linear combination of derivatives of the Dirac measure at that point \cite[Theorem~6.25]{Rudin1991}.
By the inversion formula $u(0)=(2\pi)^{-3}\int\mathrm{d}p\,\hat u(p)$ and since $\partial_au$ has the Fourier transform $p\mapsto ip_a\cdot\hat u(p)$,
every such combination at $0$ acts as $u\mapsto(2\pi)^{-3}\int\mathrm{d}p\,Q(p)\cdot\hat u(p)$ with a polynomial $Q$.
So there are polynomials $Q_{ab}$ with $\langle K'_{ab},u\rangle=(2\pi)^{-3}\int\mathrm{d}p\,Q_{ab}(p)\cdot\hat u(p)$ for all $u\in\mathcal{S}(\R^3;\C)$;
as $\hat u$ ranges over all of $\mathcal{S}(\R^3;\C)$, it follows that $M_{ab}-(M_{\mathfrak{K}})_{ab}=Q_{ab}$ almost everywhere.
The left side is bounded, and a bounded polynomial is constant.
Therefore $M-M_{\mathfrak{K}}$ coincides almost everywhere with a constant matrix $K_0$.
Now $M(p)\,p=0$ for almost every $p$, and $M_{\mathfrak{K}}(p)\,p=0$ for $p\neq0$ because $P^{\perp}(p)\,p=0$;
thus $K_0\,p=0$ for almost every $p\in\R^3$, whence $K_0=0$:
\[
M=M_{\mathfrak{K}}\qquad\text{almost everywhere on }\R^3.
\]
Together with \eqref{eq:rotation-lehmann-form} and \eqref{eq:rotation-lehmann-multiplier}, this is the representation of $S_2$ in Theorem~\ref{thm:two-point}\eqref{it:two-point-kallen-lehmann}, with the measure $\mathfrak{K}$ of Step~6.

\emph{Step 8 (mass and uniqueness).}
For $s>0$ let
\[
h_{\mathfrak{K}}(s):=\int\mathrm{d}\mathfrak{K}(m)\,\frac{s^2}{s^2+m^2}.
\]
The integrand takes its values in $[0,1]$ and $\mathfrak{K}$ is finite;
so $h_{\mathfrak{K}}$ is continuous on $(0,\infty)$ by dominated convergence, and $M_{\mathfrak{K}}(p)=P^{\perp}(p)\cdot h_{\mathfrak{K}}(|p|)$ for $p\neq0$ by \eqref{eq:rotation-lehmann-multiplier}.
If $p\neq0$ and $v$ is a unit vector orthogonal to $p$, then $v^{\dagger}M_{\mathfrak{K}}(p)\,v=h_{\mathfrak{K}}(|p|)$.
As $M=M_{\mathfrak{K}}$ almost everywhere and $M\le1$, one gets $h_{\mathfrak{K}}(|p|)\le1$ for almost every $p\in\R^3$.
The set of all $s>0$ with $h_{\mathfrak{K}}(s)>1$ is open, and the $p$ with $|p|$ in this set form a null set;
so this set is empty, and $h_{\mathfrak{K}}\le1$ on $(0,\infty)$.
For every $m\ge0$ the quotient $s^2/(s^2+m^2)$ increases to $1$ as $s\to\infty$, and monotone convergence gives
\[
\mathfrak{K}([0,\infty))=\lim_{s\to\infty}h_{\mathfrak{K}}(s)\le1.
\]

Let $\mathfrak{K}'$ be a finite positive Borel measure on $[0,\infty)$ for which the representation of Theorem~\ref{thm:two-point}\eqref{it:two-point-kallen-lehmann} holds with $\mathfrak{K}'$ in place of $\mathfrak{K}$,
and form $M_{\mathfrak{K}'}$ and $h_{\mathfrak{K}'}$ in the same way.
Then $M_{\mathfrak{K}}$ and $M_{\mathfrak{K}'}$ are bounded multipliers that represent the same form $(b,b')\mapsto S_2(\overline b,b')$ in the manner of \eqref{eq:rotation-lehmann-form};
a bounded multiplier is determined almost everywhere by this form, so $M_{\mathfrak{K}}=M_{\mathfrak{K}'}$ almost everywhere.
Since $P^{\perp}(p)\neq0$ for $p\neq0$, the continuous functions $h_{\mathfrak{K}}$ and $h_{\mathfrak{K}'}$ agree at $|p|$ for almost every $p\in\R^3$,
and, by the argument just used, they agree everywhere on $(0,\infty)$.
Division by $s^2$ gives
\begin{equation}\label{eq:rotation-stieltjes-equality}
\int\mathrm{d}\mathfrak{K}(m)\,\frac{1}{s^2+m^2}=\int\mathrm{d}\mathfrak{K}'(m)\,\frac{1}{s^2+m^2}\qquad(s>0):
\end{equation}
the Stieltjes transforms of the images of $\mathfrak{K}$ and $\mathfrak{K}'$ under $m\mapsto m^2$ coincide at every point $s^2>0$.

These transforms are Laplace transforms of finite measures, up to a shift of the variable.
For $s\ge1$, Tonelli's theorem and $(s^2+m^2)^{-1}=\int_0^\infty\mathrm{d}t\,\exp\left(-(s^2+m^2)\cdot t\right)$ give
\[
\begin{aligned}
&\int\mathrm{d}\mathfrak{K}(m)\,\frac{1}{s^2+m^2}\\
&\qquad=\int_0^\infty\mathrm{d}t\,\exp\left(-(s^2-1)\cdot t\right)\cdot\exp(-t)\cdot\int\mathrm{d}\mathfrak{K}(m)\,\exp(-tm^2),
\end{aligned}
\]
and the same holds with $\mathfrak{K}'$.
The function $t\mapsto\exp(-t)\cdot\int\mathrm{d}\mathfrak{K}(m)\,\exp(-tm^2)$ is continuous on $[0,\infty)$, by dominated convergence, and bounded by $\exp(-t)\cdot\mathfrak{K}([0,\infty))$.
It is therefore the density, with respect to Lebesgue measure, of a finite positive measure on $[0,\infty)$,
whose Laplace transform at the point $s^2-1\ge0$ is the left side of the last display.
The same holds for $\mathfrak{K}'$, and by \eqref{eq:rotation-stieltjes-equality} the two finite measures so obtained have the same Laplace transform on $[0,\infty)$.
Two finite positive Borel measures on $[0,\infty)$ with the same Laplace transform on $[0,\infty)$ coincide:
the common value of the transforms at $0$ is their mass, and if it is positive, the two measures divided by it are probability distributions with the same Laplace transform,
which coincide by the uniqueness theorem for Laplace transforms \cite[Section~XIII.1, Theorem~1]{Feller1971}.
Hence the two finite measures above are equal;
their continuous densities agree almost everywhere, hence everywhere on $[0,\infty)$:
\[
\int\mathrm{d}\mathfrak{K}(m)\,\exp(-tm^2)=\int\mathrm{d}\mathfrak{K}'(m)\,\exp(-tm^2)\qquad(t\ge0).
\]
The two sides are the Laplace transforms of the images of $\mathfrak{K}$ and $\mathfrak{K}'$ under $m\mapsto m^2$, which are finite positive Borel measures on $[0,\infty)$.
By the uniqueness statement just proved, these images are equal;
since $m\mapsto m^2$ is a homeomorphism of $[0,\infty)$ onto itself, $\mathfrak{K}'=\mathfrak{K}$.

For the free Maxwell law $\gamma$ of Definition~\ref{def:free-maxwell} one has $\gamma\left(\exp(i\mathbb{B}(b))\right)=\exp\left(-\frac{1}{2}\cdot\left\|P^{\perp}b\right\|^2\right)$ for real $b\in\mathcal{S}(\R^3;\R^3)$.
So the variables $\mathbb{B}(b)$ are jointly centred Gaussian with $S_2(b,b')=\langle P^{\perp}b,P^{\perp}b'\rangle$ for real $b,b'$,
and by bilinearity and continuity
$S_2(\overline b,b')=\langle P^{\perp}b,P^{\perp}b'\rangle=(2\pi)^{-3}\int\mathrm{d}p\,\hat b(p)^{\dagger}P^{\perp}(p)\,\hat b'(p)$ for all $b,b'\in L^2(\R^3;\C^3)$.
This is the representation with $M_{\delta_0}=P^{\perp}$, and the uniqueness gives $\mathfrak{K}=\delta_0$ for $\gamma$.
For the law $\mu_0$ of the field $F=0$ one has $S_2=0$, hence $M=0$ and $\mathfrak{K}=0$.
\end{proof}

\subsection{The anomaly index}\label{sec:anomaly}

Throughout this subsection $\mu$ is a law on $\mathcal{S}'$ that satisfies \textup{(G)}, $(T_{123})$, $(R_1)$, $(RP_1)$, $(Q_{12})$ and $(C_{12})$,
and $N\ge2$ is an integer unless something else is said.
The argument follows the corresponding argument for four-dimensional Yang--Mills theory
(Remark~\ref{rem:yang-mills-comparison}),
carried out here for a field whose components carry the weights $0$ and $\pm1$, in three dimensions.
The tensor structure of the field enters at two places.
First, the harmonics of the kernels in the angle of rotation are indexed by their anomaly index $m=(k-w)/4$
(Definition~\ref{def:anomaly-index}),
where $k$ is the frequency and $w$ the sum of the weights of the labels,
and not by the frequency $k$ itself.
Second, the antilinear reflection $\Theta_1$ keeps the label of a plate
(Lemma~\ref{lem:plate-identities}\eqref{it:plate-reflection}),
so that the weights add up correctly when the Cauchy--Schwarz inequality is applied to plate vectors.

Recall that the maps of $\C^2$ defined before Lemma~\ref{lem:forward-tube} act entrywise on configurations $Z=(z_1,\dots,z_N)\in(\C^2)^N$,
and that $X$ and $Y$ stand for real configurations.
Whether a real configuration is called $X$ or $Y$, its entries are denoted $y_1,\dots,y_N\in\R^2$,
and the coordinates of the point $y_a$ are written $y_a=(y_a^0,y_a^1)$.
A \emph{frame} is a number $\upsilon\in(\pi/2)\Z$;
with it goes the tube $\mathcal{W}_\upsilon=R_\upsilon\mathcal{W}$ of Lemma~\ref{lem:forward-tube}\eqref{it:rotated-tubes}.
Let $\mathfrak{S}_N$ be the group of the permutations of $\{1,\dots,N\}$.
For $\mathfrak{s}\in\mathfrak{S}_N$, a configuration $Z$ and a tuple $\lambda=(\lambda_1,\dots,\lambda_N)$ of labels put
\[
\mathfrak{s}Z:=\left(z_{\mathfrak{s}(a)}\right)_{1\le a\le N},
\qquad
\mathfrak{s}\lambda:=\left(\lambda_{\mathfrak{s}(a)}\right)_{1\le a\le N}.
\]
Then $w(\mathfrak{s}\lambda)=w(\lambda)$ and $s(\mathfrak{s}\lambda)=s(\lambda)$,
and for two permutations one has $\mathfrak{s}(\mathfrak{s}'Z)=(\mathfrak{s}'\circ\mathfrak{s})Z$ and $\mathfrak{s}(\mathfrak{s}'\lambda)=(\mathfrak{s}'\circ\mathfrak{s})\lambda$.
For a real configuration $X$ let
$\xi(X):=\left(y_{a+1}-y_a\right)_{1\le a<N}$
be the tuple of its consecutive differences.
A \emph{labelled configuration} is a pair $(X;\lambda)$ of a configuration $X\in(\R^2)^N_{\neq}$ and a tuple $\lambda$ of $N$ labels.
When a tuple $\lambda$ of labels is fixed, $w$ denotes its weight $w(\lambda)$, an integer.

\subsubsection{One function on all tubes}\label{sec:tube-function}

\begin{definition}[Tubes of configurations]\label{def:configuration-tubes}
Let $\mu$ be a law on $\mathcal{S}'$ with \textup{(G)}, $(T_{123})$, $(R_1)$, $(RP_1)$, $(Q_{12})$ and $(C_{12})$,
let $N\ge2$, let $G_N$ be the kernels of $\mu$ (Definition~\ref{def:plate-kernels}),
and let $\lambda$ be a tuple of $N$ labels, $w=w(\lambda)$.
For a permutation $\mathfrak{s}\in\mathfrak{S}_N$ and a frame $\upsilon$ the \emph{tube} of $(\mathfrak{s},\upsilon)$ is
\[
\mathcal{Z}(\mathfrak{s},\upsilon):=\left\{Z\in(\C^2)^N:\ z_{\mathfrak{s}(a+1)}-z_{\mathfrak{s}(a)}\in\mathcal{W}_\upsilon\ \text{for}\ 1\le a<N\right\},
\]
and the function $G^{\mathfrak{s}}_\upsilon(\cdot\,;\lambda)\colon\mathcal{Z}(\mathfrak{s},\upsilon)\to\C$ is
\[
G^{\mathfrak{s}}_\upsilon(Z;\lambda):=\exp(iw\upsilon)\cdot G_N\left(\left(R_{-\upsilon}\left(z_{\mathfrak{s}(a+1)}-z_{\mathfrak{s}(a)}\right)\right)_{1\le a<N};\mathfrak{s}\lambda\right).
\]
The union of all tubes is $\mathcal{Z}_N:=\bigcup_{\mathfrak{s},\upsilon}\mathcal{Z}(\mathfrak{s},\upsilon)$,
and the set of its real points is $\mathcal{O}_N:=\mathcal{Z}_N\cap(\R^2)^N$.
\end{definition}

For $Z$ in the tube the differences $z_{\mathfrak{s}(a+1)}-z_{\mathfrak{s}(a)}$ lie in $R_\upsilon\mathcal{W}$,
so the arguments of $G_N$ in this formula lie in $\mathcal{W}^{N-1}$.
Since $G_N$ is holomorphic on $\mathcal{W}^{N-1}$ (Proposition~\ref{prop:kernels}\eqref{it:kernel-bound}),
the function $G^{\mathfrak{s}}_\upsilon(\cdot\,;\lambda)$ is holomorphic on $\mathcal{Z}(\mathfrak{s},\upsilon)$.
As $(R_\zeta\xi)^{\pm}=\exp(\pm i\zeta)\cdot\xi^{\pm}$, one has $R_{\upsilon+2\pi}=R_\upsilon$,
and $\exp\left(iw(\upsilon+2\pi)\right)=\exp(iw\upsilon)$ because $w$ is an integer.
Adding $2\pi$ to the frame therefore changes neither the tube nor the function,
and only finitely many tubes occur.

\begin{lemma}[One function on all tubes]\label{lem:one-function}
Let $\mu$ be a law on $\mathcal{S}'$ with \textup{(G)}, $(T_{123})$, $(R_1)$, $(RP_1)$, $(Q_{12})$ and $(C_{12})$,
let $N\ge2$, and let $\lambda$ be a tuple of $N$ labels, $w=w(\lambda)$.
\begin{enumerate}
\item\label{it:tube-cones}
Every tube $\mathcal{Z}(\mathfrak{s},\upsilon)$ is an open convex cone in $(\C^2)^N$.
It is invariant under the translations $Z\mapsto Z+(v,\dots,v)$, $v\in\C^2$,
under the boosts $\Lambda_\beta$, $\beta\in\R$,
under the conjugation $Z\mapsto\overline{Z}$ and under $Z\mapsto\Re Z$.
One has $R_{\pi/2}\mathcal{Z}(\mathfrak{s},\upsilon)=\mathcal{Z}(\mathfrak{s},\upsilon+\pi/2)$.
The real points of $\mathcal{Z}(\mathfrak{s},\upsilon)$ are the $X\in(\R^2)^N$ with
\[
n_\upsilon\cdot y_{\mathfrak{s}(1)}<n_\upsilon\cdot y_{\mathfrak{s}(2)}<\dots<n_\upsilon\cdot y_{\mathfrak{s}(N)};
\]
only the coordinate along $n_\upsilon$ is constrained, and the coordinate orthogonal to $n_\upsilon$ is free.
\item\label{it:tube-densities}
Let $\mathfrak{s}\in\mathfrak{S}_N$, let $\upsilon$ be a frame,
and let $\alpha_1,\dots,\alpha_N\in L^2(\R^2;\C)$ have compact supports with $\supp\alpha_1\times\dots\times\supp\alpha_N\subset\mathcal{Z}(\mathfrak{s},\upsilon)$.
Then
\[
S_N(\alpha_1\otimes\lambda_1,\dots,\alpha_N\otimes\lambda_N)=\int\mathrm{d}Y\,\prod_{a=1}^N\alpha_a(y_a)\cdot G^{\mathfrak{s}}_\upsilon(Y;\lambda).
\]
\item\label{it:tube-gluing}
For all $\mathfrak{s},\mathfrak{s}'\in\mathfrak{S}_N$ and all frames $\upsilon$, $\upsilon'$
one has $G^{\mathfrak{s}}_\upsilon(\cdot\,;\lambda)=G^{\mathfrak{s}'}_{\upsilon'}(\cdot\,;\lambda)$ on $\mathcal{Z}(\mathfrak{s},\upsilon)\cap\mathcal{Z}(\mathfrak{s}',\upsilon')$.
Consequently there is exactly one holomorphic function $\mathcal{G}_N(\cdot\,;\lambda)\colon\mathcal{Z}_N\to\C$
with $\mathcal{G}_N(\cdot\,;\lambda)=G^{\mathfrak{s}}_\upsilon(\cdot\,;\lambda)$ on every tube $\mathcal{Z}(\mathfrak{s},\upsilon)$.
\item\label{it:tube-symmetries}
Let $\mathfrak{s}'\in\mathfrak{S}_N$ and $v\in\C^2$.
The set $\mathcal{Z}_N$ is invariant under the maps $Z\mapsto\mathfrak{s}'Z$, $Z\mapsto Z+(v,\dots,v)$, $R_{\pi/2}$ and $Z\mapsto\overline{Z}$,
and $\Lambda_\beta\mathcal{Z}_N=\mathcal{Z}_N$ for every real $\beta$.
For $Z\in\mathcal{Z}_N$ one has
\[
\mathcal{G}_N(\mathfrak{s}'Z;\mathfrak{s}'\lambda)=\mathcal{G}_N(Z;\lambda),
\qquad
\mathcal{G}_N\left(Z+(v,\dots,v);\lambda\right)=\mathcal{G}_N(Z;\lambda),
\]
\[
\mathcal{G}_N(R_{\pi/2}Z;\lambda)=i^w\cdot\mathcal{G}_N(Z;\lambda),
\qquad
\overline{\mathcal{G}_N(Z;\lambda)}=\mathcal{G}_N(\overline{Z};\bar\lambda).
\]
\item\label{it:ordered-configurations}
A configuration $X\in(\R^2)^N$ belongs to $\mathcal{O}_N$ if and only if
at least one of the two maps $a\mapsto y_a^0$ and $a\mapsto y_a^1$ is injective on $\{1,\dots,N\}$.
The set $\mathcal{O}_N$ is open and dense in $(\R^2)^N$, and it is connected.
\end{enumerate}
\end{lemma}

\begin{proof}
\eqref{it:tube-cones}
The tube $\mathcal{W}$ is an open convex cone in $\C^2$, and $R_\upsilon$ is an invertible $\C$-linear map;
so $\mathcal{W}_\upsilon$ is an open convex cone.
It is invariant under the boosts and under $\xi\mapsto\Re\xi$ by Lemma~\ref{lem:forward-tube}\eqref{it:rotated-tubes}.
It is invariant under the conjugation:
$R_\upsilon$ is a real matrix, so $\overline{R_\upsilon\xi}=R_\upsilon\overline{\xi}$ for $\xi\in\C^2$,
and $\overline{\xi}\in\mathcal{W}$ if and only if $\xi\in\mathcal{W}$ (Lemma~\ref{lem:forward-tube}\eqref{it:tube-adjoint});
hence $\overline{R_\upsilon\xi}\in\mathcal{W}_\upsilon$ for every $\xi\in\mathcal{W}$.
The tube $\mathcal{Z}(\mathfrak{s},\upsilon)$ is the preimage of $(\mathcal{W}_\upsilon)^{N-1}$ under the $\C$-linear map
$Z\mapsto\left(z_{\mathfrak{s}(a+1)}-z_{\mathfrak{s}(a)}\right)_{1\le a<N}$.
This map commutes with the boosts, which act entrywise and linearly;
its coefficients are integers, so it commutes with the entrywise conjugation and with $\Re$;
and it vanishes on the configurations $(v,\dots,v)$.
Hence $\mathcal{Z}(\mathfrak{s},\upsilon)$ is an open convex cone with the four invariances.
The maps $R_\zeta$ are diagonal in the coordinates $\xi^{\pm}$;
so $R_{\pi/2}R_\upsilon=R_{\upsilon+\pi/2}$ and $R_{\pi/2}\mathcal{W}_\upsilon=\mathcal{W}_{\upsilon+\pi/2}$,
and since $R_{\pi/2}$ commutes with the map of differences,
$R_{\pi/2}\mathcal{Z}(\mathfrak{s},\upsilon)=\mathcal{Z}(\mathfrak{s},\upsilon+\pi/2)$.
Finally, a real $\xi$ lies in $\mathcal{W}_\upsilon$ if and only if $n_\upsilon\cdot\xi>0$
(Lemma~\ref{lem:forward-tube}\eqref{it:rotated-tubes});
applied to the differences $y_{\mathfrak{s}(a+1)}-y_{\mathfrak{s}(a)}$, this describes the real points.

\eqref{it:tube-densities}
If one of the $\alpha_a$ vanishes, both sides vanish; let all supports be non-empty.
The case of a general $\mathfrak{s}$ reduces to the identity.
Indeed, put $\alpha'_a:=\alpha_{\mathfrak{s}(a)}$ and $\lambda':=\mathfrak{s}\lambda$.
The form $S_N$ is symmetric (Lemma~\ref{lem:domination-consequences}\eqref{it:moment-forms}),
so the left side does not change when the plates $\alpha_a\otimes\lambda_a$ are replaced by the $\alpha'_a\otimes\lambda'_a$.
A real $Y$ lies in $\mathcal{Z}(\mathfrak{s},\upsilon)$ if and only if $\mathfrak{s}Y$ lies in $\mathcal{Z}(\mathrm{id},\upsilon)$;
so the product of the supports of the $\alpha'_a$ lies in $\mathcal{Z}(\mathrm{id},\upsilon)$.
Moreover $\prod_a\alpha_a(y_a)=\prod_a\alpha'_a\left((\mathfrak{s}Y)_a\right)$ and $G^{\mathfrak{s}}_\upsilon(Y;\lambda)=G^{\mathrm{id}}_\upsilon(\mathfrak{s}Y;\mathfrak{s}\lambda)$,
and the substitution $Y\mapsto\mathfrak{s}Y$ preserves the Lebesgue measure of $(\R^2)^N$.
So the formula for $(\mathfrak{s},\lambda)$ follows from the formula for $(\mathrm{id},\mathfrak{s}\lambda)$.

Let now $\mathfrak{s}=\mathrm{id}$ and $g:=R^{12}_{-\upsilon}$.
As $\upsilon$ is an integer multiple of $\pi/2$, the rotation $g$ is an integer power of the quarter turn $R^{12}_{\pi/2}$;
hence $(Q_{12})$ gives $g_*\mu=\mu$.
By Lemma~\ref{lem:domination-consequences}\eqref{it:field-extension}, $\hat g\,\mathbb{B}(b)=\mathbb{B}(g\cdot b)$ in $L^2(\mu)$ for every $b\in L^2(\R^3;\C^3)$,
where $\hat g\Phi=\Phi\circ g^{-1}$ is multiplicative and preserves $\mu$.
Hence $S_N(g\cdot b_1,\dots,g\cdot b_N)=S_N(b_1,\dots,b_N)$.
With Lemma~\ref{lem:plate-identities}\eqref{it:plate-rotation} at the angle $-\upsilon$,
which gives $g\cdot(\alpha_a\otimes\lambda_a)=\exp(iw(\sigma_a)\upsilon)\cdot(\alpha_a\circ R_\upsilon)\otimes\lambda_a$ for $\lambda_a=(\sigma_a,\chi_a)$,
and with the multilinearity of $S_N$, it follows that
\[
S_N(\alpha_1\otimes\lambda_1,\dots,\alpha_N\otimes\lambda_N)
=\exp(iw\upsilon)\cdot S_N\left((\alpha_1\circ R_\upsilon)\otimes\lambda_1,\dots,(\alpha_N\circ R_\upsilon)\otimes\lambda_N\right).
\]
By \eqref{it:tube-cones}, $n_\upsilon\cdot y_1<\dots<n_\upsilon\cdot y_N$ whenever $y_a\in\supp\alpha_a$ for all $a$;
as the supports are compact,
\[
\max\left\{n_\upsilon\cdot y:\ y\in\supp\alpha_a\right\}<\min\left\{n_\upsilon\cdot y:\ y\in\supp\alpha_{a+1}\right\}
\qquad(1\le a<N).
\]
The support of $\alpha_a\circ R_\upsilon$ is $R_{-\upsilon}\supp\alpha_a$,
and the time coordinate of $R_{-\upsilon}y$ is $n_0\cdot R_{-\upsilon}y=n_\upsilon\cdot y$.
So the functions $\alpha_a\circ R_\upsilon$ have time-ordered supports,
and Proposition~\ref{prop:kernels}\eqref{it:kernel-density} gives
\[
S_N\left((\alpha_1\circ R_\upsilon)\otimes\lambda_1,\dots,(\alpha_N\circ R_\upsilon)\otimes\lambda_N\right)
=\int\mathrm{d}Y\,\prod_{a=1}^N\alpha_a(R_\upsilon y_a)\cdot G_N\left(\left(y_{a+1}-y_a\right)_a;\lambda\right).
\]
The substitution $y_a\mapsto R_{-\upsilon}y_a$, which preserves the Lebesgue measure,
turns the right side into $\int\mathrm{d}Y\,\prod_a\alpha_a(y_a)\cdot G_N\left(\left(R_{-\upsilon}(y_{a+1}-y_a)\right)_a;\lambda\right)$;
multiplied by $\exp(iw\upsilon)$, this is the asserted formula.

\eqref{it:tube-gluing}
There is nothing to prove if $I:=\mathcal{Z}(\mathfrak{s},\upsilon)\cap\mathcal{Z}(\mathfrak{s}',\upsilon')$ is empty; let it be non-empty.
By \eqref{it:tube-cones}, $I$ is open and convex, and $\Re Z\in I$ for every $Z\in I$;
so $U:=I\cap(\R^2)^N$ is a non-empty open subset of $(\R^2)^N$.
For bounded $\alpha_a$ with compact supports and $\supp\alpha_1\times\dots\times\supp\alpha_N\subset U$,
part \eqref{it:tube-densities}, applied to both tubes, gives
\[
\int\mathrm{d}Y\,\prod_{a=1}^N\alpha_a(y_a)\cdot\left(G^{\mathfrak{s}}_\upsilon(Y;\lambda)-G^{\mathfrak{s}'}_{\upsilon'}(Y;\lambda)\right)=0.
\]
Both functions are continuous on $U$, so they coincide there by Lemma~\ref{lem:analysis-facts}\eqref{it:equal-densities}.
Their difference is holomorphic on the connected open set $I$ and vanishes on the non-empty open subset $U$ of its real points;
by Lemma~\ref{lem:analysis-facts}\eqref{it:identity-theorem} it vanishes on $I$.
Hence the functions $G^{\mathfrak{s}}_\upsilon(\cdot\,;\lambda)$ fit together to one function $\mathcal{G}_N(\cdot\,;\lambda)$ on $\mathcal{Z}_N$.
It is holomorphic, since it agrees with a holomorphic function on each of the open sets $\mathcal{Z}(\mathfrak{s},\upsilon)$,
and it is unique, since these sets cover $\mathcal{Z}_N$.

\eqref{it:tube-symmetries}
Let $\mathfrak{s},\mathfrak{s}'\in\mathfrak{S}_N$.
Since $(\mathfrak{s}'Z)_{\mathfrak{s}(a)}=z_{\mathfrak{s}'(\mathfrak{s}(a))}$ and $\mathfrak{s}(\mathfrak{s}'\lambda)=(\mathfrak{s}'\circ\mathfrak{s})\lambda$,
the configuration $\mathfrak{s}'Z$ lies in $\mathcal{Z}(\mathfrak{s},\upsilon)$ if and only if $Z$ lies in $\mathcal{Z}(\mathfrak{s}'\circ\mathfrak{s},\upsilon)$, and then
\[
G^{\mathfrak{s}}_\upsilon(\mathfrak{s}'Z;\mathfrak{s}'\lambda)=G^{\mathfrak{s}'\circ\mathfrak{s}}_\upsilon(Z;\lambda)
\]
(recall $w(\mathfrak{s}'\lambda)=w$).
As $\mathfrak{s}$ ranges over $\mathfrak{S}_N$, so does $\mathfrak{s}'\circ\mathfrak{s}$;
this gives $\mathfrak{s}'\mathcal{Z}_N=\mathcal{Z}_N$ and the first identity.
The functions $G^{\mathfrak{s}}_\upsilon$ depend on differences only, and the tubes are invariant under the translations by \eqref{it:tube-cones};
this gives the second identity.
For $Z\in\mathcal{Z}(\mathfrak{s},\upsilon)$ one has $R_{\pi/2}Z\in\mathcal{Z}(\mathfrak{s},\upsilon+\pi/2)$ by \eqref{it:tube-cones};
so $R_{\pi/2}\mathcal{Z}_N\subset\mathcal{Z}_N$.
Since moreover $R_{-\upsilon-\pi/2}R_{\pi/2}=R_{-\upsilon}$,
\[
G^{\mathfrak{s}}_{\upsilon+\pi/2}(R_{\pi/2}Z;\lambda)
=\exp\left(iw(\upsilon+\pi/2)\right)\cdot G_N\left(\left(R_{-\upsilon}\left(z_{\mathfrak{s}(a+1)}-z_{\mathfrak{s}(a)}\right)\right)_a;\mathfrak{s}\lambda\right)
=i^w\cdot G^{\mathfrak{s}}_\upsilon(Z;\lambda),
\]
which is the third identity.
The tubes are invariant under the boosts by \eqref{it:tube-cones};
so $\Lambda_\beta\mathcal{Z}_N\subset\mathcal{Z}_N$ for all real $\beta$, and equality follows from $\Lambda_\beta\Lambda_{-\beta}=\mathrm{id}$.
Finally, $w(\bar\lambda)=-w$, the matrix $R_{-\upsilon}$ is real,
and $\overline{G_N(\xi;\lambda')}=G_N(\overline{\xi};\bar\lambda')$ for every tuple $\lambda'$ by Proposition~\ref{prop:kernels}\eqref{it:kernel-reality};
therefore
\[
\overline{G^{\mathfrak{s}}_\upsilon(Z;\lambda)}
=\exp(-iw\upsilon)\cdot G_N\left(\left(R_{-\upsilon}\left(\overline{z}_{\mathfrak{s}(a+1)}-\overline{z}_{\mathfrak{s}(a)}\right)\right)_a;\mathfrak{s}\bar\lambda\right)
=G^{\mathfrak{s}}_\upsilon(\overline{Z};\bar\lambda),
\]
where $\overline{Z}$ lies in $\mathcal{Z}(\mathfrak{s},\upsilon)$ by \eqref{it:tube-cones}.

\eqref{it:ordered-configurations}
By \eqref{it:tube-cones}, a real $X$ lies in a tube with the frame $\upsilon$ if and only if the numbers $n_\upsilon\cdot y_1,\dots,n_\upsilon\cdot y_N$ are pairwise distinct:
if they are, the permutation $\mathfrak{s}$ that orders them increasingly gives $X\in\mathcal{Z}(\mathfrak{s},\upsilon)$.
For $\upsilon\in(\pi/2)\Z$ one has $n_\upsilon=\pm n_0$ or $n_\upsilon=\pm n_{\pi/2}$, and $n_0\cdot y=y^0$, $n_{\pi/2}\cdot y=y^1$;
this is the description of $\mathcal{O}_N$.
The complement of $\mathcal{O}_N$ in $(\R^2)^N$ is the union, over $a<b$ and $k<l$, of the linear subspaces
$\left\{y_a^0=y_b^0,\ y_k^1=y_l^1\right\}$, each of codimension two.
By Lemma~\ref{lem:analysis-facts}\eqref{it:connected-complement}, $\mathcal{O}_N$ is open, dense and connected.
\end{proof}

The cone $(C_{12})$ with the constant one makes the kernels $G_N$ holomorphic on the whole of $\mathcal{W}^{N-1}$.
The invariance of the tubes under all boosts, and the fact that their real points are subject to conditions on one coordinate only,
are properties of this domain.

\paragraph{Harmonics.}
The function $\mathcal{G}_N$ has only finitely many harmonics in the angle of rotation.
For a complex number $\beta$ put $\Lambda_\beta:=R_{i\beta}$, as for real $\beta$;
in the coordinates $\xi^{\pm}$ the map $\Lambda_\beta$ is the multiplication by $\exp(\mp\beta)$.

\begin{lemma}[Finitely many harmonics]\label{lem:finitely-many-harmonics}
Let $\mu$ be a law on $\mathcal{S}'$ with \textup{(G)}, $(T_{123})$, $(R_1)$, $(RP_1)$, $(Q_{12})$ and $(C_{12})$,
let $N\ge2$, and let $\lambda$ be a tuple of $N$ labels, $w=w(\lambda)$.
There is exactly one family of holomorphic functions $c_m(\cdot\,;\lambda)\colon\mathcal{Z}_N\to\C$, $m\in\Z$,
such that $c_m(\cdot\,;\lambda)$ vanishes identically for every $m$ with $|w+4m|>N$ and
\begin{equation}\label{eq:rotation-boost-expansion}
\mathcal{G}_N(\Lambda_\beta Z;\lambda)=\sum_{m\in\Z}c_m(Z;\lambda)\cdot\exp\left(-(w+4m)\beta\right)
\qquad(Z\in\mathcal{Z}_N,\ \beta\in\R).
\end{equation}
For $Z\in\mathcal{Z}_N$, $\mathfrak{s}\in\mathfrak{S}_N$, $v\in\C^2$ and real $\beta$ they satisfy
\[
\begin{aligned}
c_m(\mathfrak{s}Z;\mathfrak{s}\lambda)&=c_m(Z;\lambda),
&\qquad c_m\left(Z+(v,\dots,v);\lambda\right)&=c_m(Z;\lambda),\\
c_m(\Lambda_\beta Z;\lambda)&=\exp\left(-(w+4m)\beta\right)\cdot c_m(Z;\lambda),
&\qquad c_m(R_{\pi/2}Z;\lambda)&=i^w\cdot c_m(Z;\lambda),\\
\sum_{m\in\Z}c_m(Z;\lambda)&=\mathcal{G}_N(Z;\lambda),
\end{aligned}
\]
and $\overline{c_m(X;\lambda)}=c_{-m}(X;\bar\lambda)$ for every $X\in\mathcal{O}_N$.
\end{lemma}

\begin{proof}
\emph{Pencils.}
Fix $\vartheta_0\in(0,\pi/4)$.
A real configuration $X$ is a \emph{pencil} if its consecutive differences have the form
$\xi_a:=y_{a+1}-y_a=|\xi_a|\,n_{\vartheta_a}$ with $|\xi_a|>0$ and $\vartheta_a\in(-\vartheta_0,\vartheta_0)$ for $1\le a<N$;
for a pencil put $r_X:=\min_a|\xi_a|>0$.
Let $X$ be a pencil, $\zeta=t+i\beta$ with real $t$, $\beta$, and $\upsilon$ a frame.
By Lemma~\ref{lem:forward-tube}\eqref{it:complex-rotation},
\[
\Re\left(R_{\zeta-\upsilon}\xi_a\right)^{\pm}=|\xi_a|\cdot\exp(\mp\beta)\cdot\cos(\vartheta_a+t-\upsilon),
\]
which is positive if $|t-\upsilon|<\pi/2-\vartheta_0$.
In that case $R_{-\upsilon}R_\zeta\xi_a=R_{\zeta-\upsilon}\xi_a\in\mathcal{W}$ for all $a$, that is, $R_\zeta X\in\mathcal{Z}(\mathrm{id},\upsilon)$.
The strips $\{\zeta:\ |\Re\zeta-\upsilon|<\pi/2-\vartheta_0\}$, $\upsilon\in(\pi/2)\Z$, have the width $\pi-2\vartheta_0>\pi/2$ and cover $\C$.
The map $\zeta\mapsto R_\zeta X$ is entire; hence
\[
h_X(\zeta):=\mathcal{G}_N(R_\zeta X;\lambda)
\]
is an entire function of $\zeta$.
As $R_{\zeta+\pi/2}=R_{\pi/2}R_\zeta$, Lemma~\ref{lem:one-function}\eqref{it:tube-symmetries} gives $h_X(\zeta+\pi/2)=i^w\cdot h_X(\zeta)$.
Given $t$, choose a frame $\upsilon$ with $|t-\upsilon|\le\pi/4$;
as $\pi/4<\pi/2-\vartheta_0$, the configuration $R_\zeta X$ lies in $\mathcal{Z}(\mathrm{id},\upsilon)$.
Moreover $|\vartheta_a+t-\upsilon|<\vartheta_0+\pi/4<\pi/2$, so that
\[
d_{\mathcal{W}}\left(R_{\zeta-\upsilon}\xi_a\right)=\min_{\pm}\Re\left(R_{\zeta-\upsilon}\xi_a\right)^{\pm}\ge r_X\cdot\exp(-|\beta|)\cdot\cos(\vartheta_0+\pi/4).
\]
Since $h_X(\zeta)=G^{\mathrm{id}}_\upsilon(R_\zeta X;\lambda)=\exp(iw\upsilon)\cdot G_N\left(\left(R_{\zeta-\upsilon}\xi_a\right)_a;\lambda\right)$,
the bound \eqref{eq:rotation-kernel-bound} yields
\[
\left|h_X(t+i\beta)\right|\le\prod_{a=1}^N\left(2k_2\cdot\|\chi_a\|_2\right)\cdot\left(r_X\cdot\cos(\vartheta_0+\pi/4)\right)^{-N}\cdot\exp(N|\beta|)
\qquad(t,\beta\in\R),
\]
where $\lambda_a=(\sigma_a,\chi_a)$.
Lemma~\ref{lem:harmonics}, with the quasi-period $\pi/2$ and the exponent $N$,
shows that $h_X$ is a finite sum of exponentials $\exp(ik\zeta)$ with $k-w\in4\Z$ and $|k|\le N$.
Writing $k=w+4m$ and calling the coefficients $c_m(X;\lambda)$, one has
\begin{equation}\label{eq:rotation-pencil-expansion}
\mathcal{G}_N(R_\zeta X;\lambda)=\sum_{m\in\Z}c_m(X;\lambda)\cdot\exp\left(i(w+4m)\zeta\right)
\qquad(\zeta\in\C),
\end{equation}
with $c_m(X;\lambda)=0$ unless $|w+4m|\le N$.
At $\zeta=i\beta$, where $R_{i\beta}=\Lambda_\beta$, this is \eqref{eq:rotation-boost-expansion} for $Z=X$.

\emph{Tubes.}
Let $\mathcal{M}:=\{m\in\Z:\ |w+4m|\le N\}$, a finite set, and $P(u):=\prod_{m\in\mathcal{M}}(u+w+4m)$.
Fix $\mathfrak{s}\in\mathfrak{S}_N$ and a frame $\upsilon$.
The point $\Lambda_\beta Z$ depends holomorphically on $(\beta,Z)\in\C\times(\C^2)^N$,
and it lies in $\mathcal{Z}(\mathfrak{s},\upsilon)$ whenever $\beta$ is real and $Z\in\mathcal{Z}(\mathfrak{s},\upsilon)$ (Lemma~\ref{lem:one-function}\eqref{it:tube-cones}).
Since the tube is open, the pairs $(\beta,Z)$ with $\Lambda_\beta Z\in\mathcal{Z}(\mathfrak{s},\upsilon)$ form an open neighbourhood of $\R\times\mathcal{Z}(\mathfrak{s},\upsilon)$ in $\C\times(\C^2)^N$,
and on this neighbourhood $(\beta,Z)\mapsto G^{\mathfrak{s}}_\upsilon(\Lambda_\beta Z;\lambda)$ is holomorphic.
In particular it is smooth in $\beta\in\R$,
and for every $l\ge0$ and every real $\beta$ the derivative $\partial_\beta^l\,G^{\mathfrak{s}}_\upsilon(\Lambda_\beta Z;\lambda)$ is holomorphic in $Z\in\mathcal{Z}(\mathfrak{s},\upsilon)$.

Put
\[
U:=\left\{X\in(\R^2)^N:\ R_{-\upsilon}(\mathfrak{s}X)\ \text{is a pencil}\right\},
\]
and for $X\in U$ write $X':=R_{-\upsilon}(\mathfrak{s}X)$.
The set $U$ is open, and it is not empty (for a pencil $X'$ the configuration $X$ with $\mathfrak{s}X=R_\upsilon X'$ belongs to it).
It consists of real points of $\mathcal{Z}(\mathfrak{s},\upsilon)$:
the differences of $X'$ are the $R_{-\upsilon}\left(y_{\mathfrak{s}(a+1)}-y_{\mathfrak{s}(a)}\right)$,
and the time coordinates of the differences of a pencil are positive,
so that $n_\upsilon\cdot\left(y_{\mathfrak{s}(a+1)}-y_{\mathfrak{s}(a)}\right)>0$.
For $X\in U$, since $R_{-\upsilon}$ commutes with $\Lambda_\beta$ and $\Lambda_\beta X'\in\mathcal{Z}(\mathrm{id},0)$ by the first part of the proof,
\[
G^{\mathfrak{s}}_\upsilon(\Lambda_\beta X;\lambda)
=\exp(iw\upsilon)\cdot G_N\left(\Lambda_\beta\,\xi(X');\mathfrak{s}\lambda\right)
=\exp(iw\upsilon)\cdot\mathcal{G}_N(\Lambda_\beta X';\mathfrak{s}\lambda)
\qquad(\beta\in\R).
\]
By the first part of the proof, applied to the tuple $\mathfrak{s}\lambda$, which has the weight $w$,
the right side is a linear combination of the functions $\exp\left(-(w+4m)\beta\right)$, $m\in\mathcal{M}$.
Each of them is annihilated by $P(\mathrm{d}/\mathrm{d}\beta)$, because $P(-(w+4m))=0$ for $m\in\mathcal{M}$.
So $P(\partial_\beta)\,G^{\mathfrak{s}}_\upsilon(\Lambda_\beta X;\lambda)=0$ for $X\in U$ and real $\beta$.
Let $\beta$ be real.
The function $Z\mapsto P(\partial_\beta)\,G^{\mathfrak{s}}_\upsilon(\Lambda_\beta Z;\lambda)$ has the open convex, hence connected, domain $\mathcal{Z}(\mathfrak{s},\upsilon)$ and is holomorphic there;
its zeros include the set $U$, which is open in $(\R^2)^N$, not empty, and contained in $\mathcal{Z}(\mathfrak{s},\upsilon)\cap(\R^2)^N$.
By Lemma~\ref{lem:analysis-facts}\eqref{it:identity-theorem}, it is zero on the whole tube.
Lemma~\ref{lem:analysis-facts}\eqref{it:exponential-solutions}, applied to the finite set $\{w+4m:\ m\in\mathcal{M}\}$, now gives
\[
G^{\mathfrak{s}}_\upsilon(\Lambda_\beta Z;\lambda)=\sum_{m\in\mathcal{M}}c^{\mathfrak{s},\upsilon}_m(Z)\cdot\exp\left(-(w+4m)\beta\right)
\qquad\left(Z\in\mathcal{Z}(\mathfrak{s},\upsilon),\ \beta\in\R\right),
\]
where each $c^{\mathfrak{s},\upsilon}_m(Z)$ is a fixed linear combination of the derivatives
$\partial_\beta^l\,G^{\mathfrak{s}}_\upsilon(\Lambda_\beta Z;\lambda)$ at $\beta=0$, $l<\#|\mathcal{M}|$;
in particular $c^{\mathfrak{s},\upsilon}_m$ is holomorphic on the tube.

Let $Z$ lie in two tubes $\mathcal{Z}(\mathfrak{s},\upsilon)$ and $\mathcal{Z}(\mathfrak{s}',\upsilon')$.
Their intersection is invariant under the boosts, so it contains $\Lambda_\beta Z$ for all real $\beta$,
and $G^{\mathfrak{s}}_\upsilon(\Lambda_\beta Z;\lambda)=G^{\mathfrak{s}'}_{\upsilon'}(\Lambda_\beta Z;\lambda)$ by Lemma~\ref{lem:one-function}\eqref{it:tube-gluing}.
The exponentials being linearly independent (Lemma~\ref{lem:analysis-facts}\eqref{it:exponential-solutions}),
$c^{\mathfrak{s},\upsilon}_m(Z)=c^{\mathfrak{s}',\upsilon'}_m(Z)$ for all $m\in\mathcal{M}$.
Therefore $c_m(Z;\lambda):=c^{\mathfrak{s},\upsilon}_m(Z)$ for $Z\in\mathcal{Z}(\mathfrak{s},\upsilon)$ and $m\in\mathcal{M}$,
and $c_m(\cdot\,;\lambda):=0$ for $m\notin\mathcal{M}$,
defines holomorphic functions on $\mathcal{Z}_N$;
they satisfy \eqref{eq:rotation-boost-expansion},
because $\mathcal{G}_N(\Lambda_\beta Z;\lambda)=G^{\mathfrak{s}}_\upsilon(\Lambda_\beta Z;\lambda)$ for $Z\in\mathcal{Z}(\mathfrak{s},\upsilon)$.
At every $Z$, the right side of \eqref{eq:rotation-boost-expansion} is a finite sum of exponentials in $\beta$ with pairwise distinct exponents,
whose coefficients are determined by the left side;
this proves uniqueness.
In particular, at a pencil $X$ the values $c_m(X;\lambda)$ are the coefficients found in the first part of the proof,
and \eqref{eq:rotation-pencil-expansion} holds with them.

\emph{The listed properties.}
They follow by comparison of the coefficients of the exponentials in \eqref{eq:rotation-boost-expansion},
together with Lemma~\ref{lem:one-function}\eqref{it:tube-symmetries}.
The boost $\Lambda_\beta$ commutes with every permutation of the entries;
hence $\mathcal{G}_N(\Lambda_\beta\,\mathfrak{s}Z;\mathfrak{s}\lambda)=\mathcal{G}_N(\mathfrak{s}\,\Lambda_\beta Z;\mathfrak{s}\lambda)=\mathcal{G}_N(\Lambda_\beta Z;\lambda)$,
and both sides are expanded with the same weight $w$; this gives the first identity.
Next, $\Lambda_\beta\left(Z+(v,\dots,v)\right)=\Lambda_\beta Z+(\Lambda_\beta v,\dots,\Lambda_\beta v)$, which gives the second.
From $\Lambda_\beta\Lambda_{\beta'}=\Lambda_{\beta+\beta'}$ one obtains
\[
\begin{aligned}
&\sum_{m\in\Z}c_m(\Lambda_{\beta'}Z;\lambda)\cdot\exp\left(-(w+4m)\beta\right)
=\mathcal{G}_N(\Lambda_{\beta+\beta'}Z;\lambda)\\
&\qquad=\sum_{m\in\Z}c_m(Z;\lambda)\cdot\exp\left(-(w+4m)\beta'\right)\cdot\exp\left(-(w+4m)\beta\right),
\end{aligned}
\]
which gives the third.
The fourth follows from $\Lambda_\beta R_{\pi/2}=R_{\pi/2}\Lambda_\beta$, and the fifth is \eqref{eq:rotation-boost-expansion} at $\beta=0$.
For the last, let $X\in\mathcal{O}_N$.
For real $\beta$ one has $\overline{\Lambda_\beta X}=\Lambda_{-\beta}X$ (Lemma~\ref{lem:forward-tube}\eqref{it:tube-boosts}) and $w(\bar\lambda)=-w$;
so Lemma~\ref{lem:one-function}\eqref{it:tube-symmetries} gives
\[
\overline{\mathcal{G}_N(\Lambda_\beta X;\lambda)}
=\mathcal{G}_N(\Lambda_{-\beta}X;\bar\lambda)
=\sum_{m'\in\Z}c_{m'}(X;\bar\lambda)\cdot\exp\left((-w+4m')\beta\right),
\]
while the left side equals $\sum_{m}\overline{c_m(X;\lambda)}\cdot\exp\left(-(w+4m)\beta\right)$.
The exponents $-(w+4m)$ and $-w+4m'$ agree exactly when $m'=-m$;
hence $\overline{c_m(X;\lambda)}=c_{-m}(X;\bar\lambda)$.
\end{proof}

\begin{definition}[The anomaly index]\label{def:anomaly-index}
Let $\mu$ be a law on $\mathcal{S}'$ with \textup{(G)}, $(T_{123})$, $(R_1)$, $(RP_1)$, $(Q_{12})$ and $(C_{12})$,
and let $c_m$ be the functions of Lemma~\ref{lem:finitely-many-harmonics}.
For a tuple $\lambda$ of $N\ge2$ labels and $m\in\Z$, the function $c_m(\cdot\,;\lambda)$ is called the \emph{harmonic} of $\mathcal{G}_N(\cdot\,;\lambda)$
with the \emph{anomaly index} $m$;
by \eqref{eq:rotation-pencil-expansion} it carries the frequency $k=w(\lambda)+4m$ in the angle of rotation, so that $m=\left(k-w(\lambda)\right)/4$.
The \emph{anomaly index} of the $N$-point functions of $\mu$ is $\mathfrak{a}(0):=\mathfrak{a}(1):=0$ and, for $N\ge2$,
\[
\mathfrak{a}(N):=\max\left\{|m|:\ c_m(\cdot\,;\lambda)\ \text{is not identically zero on}\ \mathcal{Z}_N\ \text{for some tuple}\ \lambda\ \text{of}\ N\ \text{labels}\right\},
\]
with $\mathfrak{a}(N):=0$ if all the $c_m$ vanish.
It has the following properties.
\begin{enumerate}
\item\label{it:anomaly-linear-bound}
\emph{Linear bound.} One has $\mathfrak{a}(N)\le N/2$ for every $N$.
\item\label{it:anomaly-time-ordered}
Let $N\ge2$.
If $c_m(\cdot\,;\lambda)$ is not identically zero, there is a permutation $\mathfrak{s}\in\mathfrak{S}_N$
such that $c_m(\cdot\,;\mathfrak{s}\lambda)$ is not identically zero on the strictly time-ordered configurations,
that is, on the set $\{y_1^0<y_2^0<\dots<y_N^0\}$ of the real points of $\mathcal{Z}(\mathrm{id},0)$.
\item\label{it:anomaly-both-signs}
\emph{Both signs occur.} If $\mathfrak{a}(N)\ge1$, there are tuples $\lambda$, $\lambda'$ of $N$ labels
such that neither $c_{\mathfrak{a}(N)}(\cdot\,;\lambda)$ nor $c_{-\mathfrak{a}(N)}(\cdot\,;\lambda')$ is identically zero.
\end{enumerate}
\end{definition}

\begin{proof}[Proof of the three properties]
\eqref{it:anomaly-linear-bound}
There is nothing to prove for $N\le1$.
For $N\ge2$, each weight $w(\sigma_a)$ lies in $\{-1,0,1\}$, so $|w|\le N$;
and a harmonic $c_m(\cdot\,;\lambda)$ that is not identically zero has $|w+4m|\le N$ by Lemma~\ref{lem:finitely-many-harmonics}.
Hence $4|m|\le|w+4m|+|w|\le2N$.
The bound is linear in $N$ because in the kernel bound \eqref{eq:rotation-kernel-bound}
each of the $N$ plates contributes the inverse distance to the boundary of the tube with the exponent one.

\eqref{it:anomaly-time-ordered}
The function $c_m(\cdot\,;\lambda)$ does not vanish identically on some tube $\mathcal{Z}(\mathfrak{s},\upsilon)$, $\upsilon=l\pi/2$ with $l\in\Z$.
By Lemma~\ref{lem:one-function}\eqref{it:tube-cones}, $R_{\pi/2}^{-l}$ maps $\mathcal{Z}(\mathfrak{s},\upsilon)$ onto $\mathcal{Z}(\mathfrak{s},0)$,
and $c_m(R_{\pi/2}^{-l}Z;\lambda)=i^{-lw}\cdot c_m(Z;\lambda)$ by Lemma~\ref{lem:finitely-many-harmonics};
so $c_m(\cdot\,;\lambda)$ does not vanish identically on $\mathcal{Z}(\mathfrak{s},0)$.
A configuration $Z$ lies in $\mathcal{Z}(\mathfrak{s},0)$ if and only if $\mathfrak{s}Z$ lies in $\mathcal{Z}(\mathrm{id},0)$,
and $c_m(\mathfrak{s}Z;\mathfrak{s}\lambda)=c_m(Z;\lambda)$;
so $c_m(\cdot\,;\mathfrak{s}\lambda)$ does not vanish identically on $\mathcal{Z}(\mathrm{id},0)$.
This tube is convex, and its real points form a non-empty open subset of $(\R^2)^N$;
by Lemma~\ref{lem:analysis-facts}\eqref{it:identity-theorem}, $c_m(\cdot\,;\mathfrak{s}\lambda)$ does not vanish identically on them.

\eqref{it:anomaly-both-signs}
Choose a tuple $\lambda$ and an integer $m$ with $|m|=\mathfrak{a}(N)$ such that $c_m(\cdot\,;\lambda)$ is not identically zero.
By \eqref{it:anomaly-time-ordered} there are $\mathfrak{s}\in\mathfrak{S}_N$ and a real point $X$ of $\mathcal{Z}(\mathrm{id},0)$,
which lies in $\mathcal{O}_N$, with $c_m(X;\mathfrak{s}\lambda)\neq0$.
By Lemma~\ref{lem:finitely-many-harmonics}, $c_{-m}(X;\overline{\mathfrak{s}\lambda})=\overline{c_m(X;\mathfrak{s}\lambda)}\neq0$.
So the harmonics of the indices $m$ and $-m$ both occur, one of these indices being $\mathfrak{a}(N)$ and the other $-\mathfrak{a}(N)$.
\end{proof}

\begin{remark}
The harmonic $c_m$ carries the frequency $w+4m$.
Hence $\mathfrak{a}(N)=0$ says that, for every tuple $\lambda$ of $N$ labels, no frequency other than $w$ occurs in $\mathcal{G}_N(\cdot\,;\lambda)$,
that is, that all the functions $\mathcal{G}_N(\cdot\,;\lambda)$ are covariant under the rotations in the sense of \eqref{eq:rotation-density-criterion}.
Precisely, $\mathfrak{a}(N)=0$ holds if and only if
\[
\mathcal{G}_N(R_tY;\lambda)=\exp(iwt)\cdot\mathcal{G}_N(Y;\lambda)
\]
for every tuple $\lambda$ of $N$ labels, every $Y\in\mathcal{O}_N$ and every real $t$ such that $R_{t'}Y\in\mathcal{O}_N$ for all $t'$ between $0$ and $t$.
To see this, let $\lambda$ and $Y\in\mathcal{O}_N$ be given.
The numbers $t$ just described form an open interval that contains $0$,
and the $\zeta=t+i\beta$ with $t$ in this interval and real $\beta$ form an open strip, which contains the interval and the imaginary axis.
For such $\zeta$ one has $R_tY\in\mathcal{O}_N\subset\mathcal{Z}_N$,
hence $R_\zeta Y=\Lambda_\beta(R_tY)\in\Lambda_\beta\mathcal{Z}_N=\mathcal{Z}_N$ (Lemma~\ref{lem:one-function}\eqref{it:tube-symmetries});
as $\zeta\mapsto R_\zeta Y$ is entire, the function $\zeta\mapsto\exp(-iw\zeta)\cdot\mathcal{G}_N(R_\zeta Y;\lambda)$ is holomorphic on the strip.
By the identity theorem (Lemma~\ref{lem:analysis-facts}\eqref{it:identity-theorem}),
it is constant on the interval if and only if it is constant on the imaginary axis,
that is, if and only if $\mathcal{G}_N(\Lambda_\beta Y;\lambda)=\exp(-w\beta)\cdot\mathcal{G}_N(Y;\lambda)$ for all real $\beta$.
By \eqref{eq:rotation-boost-expansion} at $Z=Y$, by the identity $\sum_mc_m(Y;\lambda)=\mathcal{G}_N(Y;\lambda)$ of Lemma~\ref{lem:finitely-many-harmonics}
and by the linear independence of the exponentials (Lemma~\ref{lem:analysis-facts}\eqref{it:exponential-solutions}),
this last identity holds if and only if $c_m(Y;\lambda)=0$ for every $m\neq0$.
As $\mathfrak{a}(N)=0$ means that $c_m(\cdot\,;\lambda)$ vanishes on $\mathcal{Z}_N$ for every tuple $\lambda$ and every $m\neq0$,
it remains to note that a harmonic $c_m(\cdot\,;\lambda)$ which vanishes on $\mathcal{O}_N$ vanishes on $\mathcal{Z}_N$:
on each tube it is holomorphic, the tube is convex,
and the real points of the tube lie in $\mathcal{O}_N$ and form a non-empty open subset of $(\R^2)^N$ (Lemma~\ref{lem:one-function}\eqref{it:tube-cones}),
so that Lemma~\ref{lem:analysis-facts}\eqref{it:identity-theorem} applies.
The moments are tied to these functions by Lemma~\ref{lem:density-criterion}:
for a real $t$ and an open set $U\subset\mathcal{O}_N$ with $R_tU=U$ on which $\mathcal{G}_N(\cdot\,;\lambda)$ is a density of $S_N$ as in that lemma,
the identity \eqref{eq:rotation-density-criterion} for $G=\mathcal{G}_N(\cdot\,;\lambda)$ holds on $U$
if and only if $R_t$ leaves the moments $S_N$ of the plates $\alpha_a\otimes\lambda_a$ unchanged
for all $\alpha_a\in L^2(\R^2;\C)$ with compact supports and $\supp\alpha_1\times\dots\times\supp\alpha_N\subset U$.

The linear bound, convexity and divisibility are formulated for the anomaly index $m$, and not for the frequency $k$.
Indeed, for a law that is invariant under the rotations of the plane,
the frequency $k=w$ takes every integer value between $-N$ and $N$:
for the free Maxwell law of Definition~\ref{def:free-maxwell} and $N=2a$ with a positive integer $a$,
the $2a$-point function of plates whose labels all have the first entry $\sigma=+$ has the frequency $k=2a=N$.
So the covariant terms, the harmonics $c_0$ with the frequency $k=w$, attain the linear bound $|k|\le N$,
and convexity, the linear bound and divisibility, formulated for the top frequency in place of the anomaly index, would give no information.
\end{remark}

\subsubsection{Exact phases and positivity}\label{sec:exact-phases}

The coordinate $y^0$ of a point $y=(y^0,y^1)$ of the plane is called its \emph{time}, and $y^1$ its \emph{abscissa};
a \emph{chord} of a configuration $X=(y_1,\dots,y_N)$ is a difference $y_a-y_c$ of two of its points, $a\neq c$.

\begin{lemma}[Exact phases]\label{lem:exact-phases}
Let $\mu$ be a law on $\mathcal{S}'$ that satisfies \textup{(G)}, $(T_{123})$, $(R_1)$, $(RP_1)$, $(Q_{12})$ and $(C_{12})$.
Let $N\geq2$, let $\lambda$ be a tuple of $N$ labels, put $w:=w(\lambda)$,
and let $c_m=c_m(\cdot\,;\lambda)$, $m\in\Z$, be the holomorphic functions on $\mathcal{Z}_N$ of Lemma~\ref{lem:finitely-many-harmonics}.
\begin{enumerate}
\item\label{it:phase-interval}
Let $\mathfrak{s}$ be a permutation of $\{1,\dots,N\}$, let $\upsilon\in\frac{\pi}{2}\Z$, and let $X$ be a real point of $\mathcal{Z}(\mathfrak{s},\upsilon)$.
Then the set
\[
I_X:=\{t\in(-\pi,\pi):\ R_tX\in\mathcal{Z}(\mathfrak{s},\upsilon)\}
\]
is an open interval containing $0$, and
\[
c_m(R_tX;\lambda)=\exp(i(w+4m)t)\cdot c_m(X;\lambda)\qquad(t\in I_X,\ m\in\Z).
\]
\item\label{it:generic-configurations}
Let $\mathcal{O}^*_N$ be the set of the configurations $X=(y_1,\dots,y_N)\in(\R^2)^N_{\neq}$ such that
\[
(y_a-y_c)\cdot(y_k-y_l)\neq0
\]
for any two different two-element subsets $\{a,c\}$ and $\{k,l\}$ of $\{1,\dots,N\}$;
these are the configurations of distinct points no two chords of which are perpendicular.
The set $\mathcal{O}^*_N$ is open and dense in $(\R^2)^N$, invariant under every rotation $R_t$, and contained in $\mathcal{O}_N$.
\item\label{it:phase-ordered}
If $t$ is real and both $X$ and $R_tX$ lie in $\mathcal{O}_N$, then
$c_m(R_tX;\lambda)=\exp(i(w+4m)t)\cdot c_m(X;\lambda)$ for every $m\in\Z$.
\end{enumerate}
\end{lemma}

\begin{proof}
\eqref{it:phase-interval}
Put $\xi_a:=y_{\mathfrak{s}(a+1)}-y_{\mathfrak{s}(a)}$ for $1\leq a<N$.
By the description of the real points of the tube in Lemma~\ref{lem:one-function}\eqref{it:tube-cones}, $n_\upsilon\cdot\xi_a>0$.
For a real $\xi$ the point $R_{-\upsilon}\xi$ is real,
and by the formula for the rotations in Definition~\ref{def:labels-plates} its time is $\xi^0\cdot\cos\upsilon+\xi^1\cdot\sin\upsilon=n_\upsilon\cdot\xi$;
hence $R_{-\upsilon}\xi_a=|\xi_a|\,n_{\vartheta_a}$ with an angle $\vartheta_a\in\left(-\frac{\pi}{2},\frac{\pi}{2}\right)$,
and, rotations of the plane being commutative, for real $t$
\[
n_\upsilon\cdot R_t\xi_a=(R_{-\upsilon}R_t\xi_a)^0=(R_tR_{-\upsilon}\xi_a)^0=|\xi_a|\cdot\cos(\vartheta_a+t).
\]
By Lemma~\ref{lem:forward-tube}\eqref{it:rotated-tubes}, $R_t\xi_a$ lies in $\mathcal{W}_\upsilon$ if and only if $\cos(\vartheta_a+t)>0$.
For $t\in(-\pi,\pi)$ this condition means $t\in\left(-\frac{\pi}{2}-\vartheta_a,\frac{\pi}{2}-\vartheta_a\right)$, and this interval contains $0$.
The consecutive differences of $R_tX$ along $\mathfrak{s}$ are the vectors $R_t\xi_a$;
therefore $I_X$ is the intersection of these $N-1$ intervals, an open interval containing $0$.

Let $\zeta=t+i\beta$ with $t\in I_X$ and $\beta\in\R$.
The complex rotation $R_\zeta$ multiplies the coordinates $\xi^{\pm}$ of a point $\xi\in\C^2$ by $\exp(\pm i\zeta)$;
so the complex rotations form a group, and $R_\zeta=\Lambda_\beta R_t$.
The configuration $R_tX$ lies in $\mathcal{Z}(\mathfrak{s},\upsilon)$, and this tube is invariant under $\Lambda_\beta$ (Lemma~\ref{lem:one-function}\eqref{it:tube-cones});
hence $R_\zeta X$ lies in it as well.
The map $\zeta\mapsto R_\zeta X$ is entire, and $\mathcal{G}_N(\cdot\,;\lambda)$ is holomorphic on $\mathcal{Z}_N$;
consequently $\zeta\mapsto\mathcal{G}_N(R_\zeta X;\lambda)$ is holomorphic on the strip $I_X+i\R$.
On the imaginary axis one has $R_{i\beta}=\Lambda_\beta$, and the expansion \eqref{eq:rotation-boost-expansion} at $Z=X$ reads
\[
\mathcal{G}_N(R_{i\beta}X;\lambda)=\sum_mc_m(X;\lambda)\cdot\exp(i(w+4m)\cdot i\beta).
\]
The right side is the restriction to the imaginary axis of an entire function of $\zeta$, the sum being finite,
and the strip is connected and contains the imaginary axis.
The identity theorem (Lemma~\ref{lem:analysis-facts}\eqref{it:identity-theorem}) gives
\begin{equation}\label{eq:rotation-strip}
\mathcal{G}_N(R_\zeta X;\lambda)=\sum_mc_m(X;\lambda)\cdot\exp(i(w+4m)\zeta)\qquad(\zeta\in I_X+i\R).
\end{equation}
At $\zeta=t+i\beta$ the right side of \eqref{eq:rotation-strip} equals $\sum_m\exp(i(w+4m)t)\cdot c_m(X;\lambda)\cdot\exp(-(w+4m)\beta)$,
while its left side is $\mathcal{G}_N(\Lambda_\beta R_tX;\lambda)=\sum_mc_m(R_tX;\lambda)\cdot\exp(-(w+4m)\beta)$,
by \eqref{eq:rotation-boost-expansion} at $Z=R_tX$.
Both are finite combinations of the functions $\beta\mapsto\exp(-(w+4m)\beta)$ with pairwise different exponents,
which are linearly independent (Lemma~\ref{lem:analysis-facts}\eqref{it:exponential-solutions});
comparison of the coefficients gives the assertion.

\eqref{it:generic-configurations}
The set $\mathcal{O}^*_N$ is the complement in $(\R^2)^N$ of the zero set of the product of the polynomials $|y_a-y_c|^2$, $a<c$,
and $(y_a-y_c)\cdot(y_k-y_l)$, $\{a,c\}\neq\{k,l\}$.
None of these factors is the zero polynomial.
The polynomial $|y_a-y_c|^2$ takes the value $1$ at a configuration with $y_a=0$ and $y_c=n_0$.
If $\{a,c\}$ and $\{k,l\}$ are disjoint, $(y_a-y_c)\cdot(y_k-y_l)$ takes the value $1$ at a configuration with $y_a=y_k=0$ and $y_c=y_l=n_0$;
if the two pairs have one index in common, it takes the value $1$ or $-1$ at a configuration in which the common point is $0$ and the two other points are $n_0$.
So the product is a non-zero polynomial, whose zero set is closed with empty interior (Lemma~\ref{lem:analysis-facts}\eqref{it:polynomial-zeros}),
and $\mathcal{O}^*_N$ is open and dense.
Rotations preserve scalar products and keep distinct points distinct; hence $R_t\mathcal{O}^*_N=\mathcal{O}^*_N$.
Let now $X\in\mathcal{O}^*_N$, and suppose that $y_a^0=y_c^0$ and $y_k^1=y_l^1$ for some $a\neq c$ and $k\neq l$.
Then $y_a-y_c$ is a multiple of $n_{\pi/2}$ and $y_k-y_l$ is a multiple of $n_0$, so that $(y_a-y_c)\cdot(y_k-y_l)=0$.
If $\{a,c\}\neq\{k,l\}$, this contradicts the definition of $\mathcal{O}^*_N$;
if $\{a,c\}=\{k,l\}$, it forces $y_a=y_c$, which contradicts $X\in(\R^2)^N_{\neq}$.
Therefore the times of $X$ are pairwise distinct, or its abscissae are,
and $X\in\mathcal{O}_N$ by Lemma~\ref{lem:one-function}\eqref{it:ordered-configurations}.

\eqref{it:phase-ordered}
Suppose first that $X\in\mathcal{O}^*_N$.
By \eqref{it:generic-configurations}, $R_sX\in\mathcal{O}_N$ for every real $s$; thus $R_sX$ is a real point of some tube.
Part~\eqref{it:phase-interval}, applied to $R_sX$ and to that tube, provides an open interval $I$ containing $0$ such that
\[
c_m(R_{s+t'}X;\lambda)=\exp(i(w+4m)t')\cdot c_m(R_sX;\lambda)\qquad(t'\in I).
\]
So the function $s\mapsto\exp(-i(w+4m)s)\cdot c_m(R_sX;\lambda)$ is locally constant on $\R$, hence constant,
and the comparison of its values at $0$ and at $t$ is the assertion.
For an arbitrary $X$, consider the set of the configurations $Y$ such that $Y$ and $R_tY$ lie in $\mathcal{O}_N$.
This set is open and contains $X$, and the two functions $Y\mapsto c_m(R_tY;\lambda)$ and $Y\mapsto\exp(i(w+4m)t)\cdot c_m(Y;\lambda)$ are continuous on it,
since $c_m(\cdot\,;\lambda)$ is holomorphic on $\mathcal{Z}_N$.
They agree on the intersection of the set with $\mathcal{O}^*_N$, which is dense in it by \eqref{it:generic-configurations};
hence they agree at $X$.
\end{proof}

\begin{lemma}[Harmonics on all configurations]\label{lem:harmonic-extension}
Let $\mu$ be a law on $\mathcal{S}'$ that satisfies \textup{(G)}, $(T_{123})$, $(R_1)$, $(RP_1)$, $(Q_{12})$ and $(C_{12})$.
Let $N\geq2$, let $\lambda$ be a tuple of $N$ labels, put $w:=w(\lambda)$,
and let $c_m$, $m\in\Z$, be the functions of Lemma~\ref{lem:finitely-many-harmonics}.
For $X\in(\R^2)^N_{\neq}$ the set
\[
A(X):=\{t\in\R:\ R_tX\in\mathcal{O}_N\}
\]
contains every $t$ for which no chord of $R_tX$ is parallel to $n_{\pi/2}$, that is, every real $t$ outside a set that is finite modulo $\pi$;
and the number
\[
\tilde c_m(X;\lambda):=\exp(-i(w+4m)t)\cdot c_m(R_tX;\lambda),\qquad t\in A(X),
\]
does not depend on the choice of $t$ in $A(X)$.
Let $\tilde G_N(X;\lambda):=\sum_{m\in\Z}\tilde c_m(X;\lambda)$, and write $\tilde c^{(2a)}_m$ for $\tilde c_m$ when $N=2a$.
Then the following hold.
\begin{enumerate}
\item\label{it:harmonic-analytic}
The function $\tilde c_m(\cdot\,;\lambda)$ is real analytic on $(\R^2)^N_{\neq}$, which is open and connected,
and it coincides with $c_m(\cdot\,;\lambda)$ on $\mathcal{O}_N$.
\item\label{it:harmonic-symmetric}
$\tilde c_m(\mathfrak{s}X;\mathfrak{s}\lambda)=\tilde c_m(X;\lambda)$ for every permutation $\mathfrak{s}$ of $\{1,\dots,N\}$:
the function $\tilde c_m$ depends only on the labelled set $\{(y_1,\lambda_1),\dots,(y_N,\lambda_N)\}$.
\item\label{it:harmonic-translation}
$\tilde c_m(X+(v,\dots,v);\lambda)=\tilde c_m(X;\lambda)$ for every $v\in\R^2$.
\item\label{it:harmonic-phase}
$\tilde c_m(R_tX;\lambda)=\exp(i(w+4m)t)\cdot\tilde c_m(X;\lambda)$ for every real $t$.
Together with \eqref{it:harmonic-translation}, the same identity holds for the rotation
$X\mapsto(v,\dots,v)+R_t(X-(v,\dots,v))$ by the angle $t$ about an arbitrary point $v\in\R^2$.
\item\label{it:harmonic-sum}
The sum that defines $\tilde G_N(X;\lambda)$ has only finitely many non-zero terms.
The function $\tilde G_N(\cdot\,;\lambda)$ is real analytic on $(\R^2)^N_{\neq}$, it equals $\mathcal{G}_N(\cdot\,;\lambda)$ on $\mathcal{O}_N$, and
\[
\tilde G_N(R_tX;\lambda)=\sum_m\exp(i(w+4m)t)\cdot\tilde c_m(X;\lambda)\qquad(t\in\R).
\]
\item\label{it:harmonic-top}
$\tilde c_m(\cdot\,;\lambda)=0$ whenever $|m|>\mathfrak{a}(N)$;
and if $\mathfrak{a}(N)\geq1$, then $\tilde c_{-\mathfrak{a}(N)}(\cdot\,;\lambda')$ is not identically zero for some tuple $\lambda'$ of $N$ labels.
\end{enumerate}
\end{lemma}

\begin{proof}
Let $X\in(\R^2)^N_{\neq}$, and let $a\neq c$.
Write the chord $y_a-y_c$ of $X$ as $|y_a-y_c|\,n_\vartheta$;
the rotation $R_t$ carries it into the chord of $R_tX$ with the time component $|y_a-y_c|\cdot\cos(\vartheta+t)$,
which is parallel to $n_{\pi/2}$ exactly when $t\in\frac{\pi}{2}-\vartheta+\pi\Z$.
If no chord of $R_tX$ is parallel to $n_{\pi/2}$, the times of $R_tX$ are pairwise distinct,
and $R_tX\in\mathcal{O}_N$ by Lemma~\ref{lem:one-function}\eqref{it:ordered-configurations}.
This proves the assertion on $A(X)$.
For $t,t'\in A(X)$, Lemma~\ref{lem:exact-phases}\eqref{it:phase-ordered}, applied to the configuration $R_tX$ and the angle $t'-t$, gives
\[
c_m(R_{t'}X;\lambda)=\exp(i(w+4m)(t'-t))\cdot c_m(R_tX;\lambda),
\]
so that $\tilde c_m(X;\lambda)$ does not depend on $t$.

\eqref{it:harmonic-analytic}
Fix $t\in A(X)$.
As $\mathcal{O}_N$ is open, $t\in A(Y)$ for all $Y$ in a neighbourhood $U\subseteq(\R^2)^N_{\neq}$ of $X$,
and on $U$ the function $\tilde c_m(\cdot\,;\lambda)$ is $Y\mapsto\exp(-i(w+4m)t)\cdot c_m(R_tY;\lambda)$:
the linear map $Y\mapsto R_tY$, followed by the restriction of the holomorphic function $c_m(\cdot\,;\lambda)$ to the real points of $\mathcal{Z}_N$.
Hence $\tilde c_m(\cdot\,;\lambda)$ is real analytic near $X$.
If $X\in\mathcal{O}_N$, the choice $t=0$ gives $\tilde c_m(X;\lambda)=c_m(X;\lambda)$.
Finally, $(\R^2)^N_{\neq}$ is the complement in $(\R^2)^N$ of the subspaces $\{y_a=y_c\}$, $a<c$, each of codimension $2$;
it is open and connected by Lemma~\ref{lem:analysis-facts}\eqref{it:connected-complement}.

\eqref{it:harmonic-symmetric}
The set $\mathcal{O}_N$ is invariant under permutations of the points, so that $A(\mathfrak{s}X)=A(X)$.
Moreover $R_t\mathfrak{s}X=\mathfrak{s}R_tX$, $w(\mathfrak{s}\lambda)=w$, and $c_m(\mathfrak{s}Z;\mathfrak{s}\lambda)=c_m(Z;\lambda)$ by Lemma~\ref{lem:finitely-many-harmonics}.
For $t\in A(X)$ it follows that
\[
\tilde c_m(\mathfrak{s}X;\mathfrak{s}\lambda)=\exp(-i(w+4m)t)\cdot c_m(\mathfrak{s}R_tX;\mathfrak{s}\lambda)=\exp(-i(w+4m)t)\cdot c_m(R_tX;\lambda)=\tilde c_m(X;\lambda).
\]

\eqref{it:harmonic-translation}
One has $R_t(X+(v,\dots,v))=R_tX+(R_tv,\dots,R_tv)$, and $\mathcal{O}_N$ is invariant under translations, so that $A(X+(v,\dots,v))=A(X)$.
The functions $c_m$ are invariant under translations by Lemma~\ref{lem:finitely-many-harmonics};
hence, for $t\in A(X)$, both sides of the assertion equal $\exp(-i(w+4m)t)\cdot c_m(R_tX;\lambda)$.

\eqref{it:harmonic-phase}
A real number $t'$ lies in $A(R_tX)$ if and only if $t+t'$ lies in $A(X)$, and for such $t'$
\[
\tilde c_m(R_tX;\lambda)=\exp(-i(w+4m)t')\cdot c_m(R_{t+t'}X;\lambda)=\exp(i(w+4m)t)\cdot\tilde c_m(X;\lambda).
\]
The rotation about $v$ is the translation by $-v$, followed by $R_t$ and by the translation by $v$;
so its case follows from \eqref{it:harmonic-translation}, the identity just proved, and \eqref{it:harmonic-translation} once more.

\eqref{it:harmonic-sum}
By Lemma~\ref{lem:finitely-many-harmonics}, $c_m(\cdot\,;\lambda)=0$ whenever $|w+4m|>N$, and for such $m$ also $\tilde c_m(\cdot\,;\lambda)=0$, by its definition;
moreover $\sum_mc_m(\cdot\,;\lambda)=\mathcal{G}_N(\cdot\,;\lambda)$ on $\mathcal{Z}_N$.
The assertions follow from \eqref{it:harmonic-analytic} and \eqref{it:harmonic-phase}.

\eqref{it:harmonic-top}
By the definition of $\mathfrak{a}(N)$ (Definition~\ref{def:anomaly-index}), for $|m|>\mathfrak{a}(N)$ the function $c_m(\cdot\,;\lambda)$ vanishes identically on $\mathcal{Z}_N$, whatever the labels;
then $\tilde c_m(\cdot\,;\lambda)=0$ by its definition.
Let $\mathfrak{a}(N)\geq1$.
By Definition~\ref{def:anomaly-index}\eqref{it:anomaly-both-signs} there is a tuple of $N$ labels for which $c_{-\mathfrak{a}(N)}$ does not vanish identically,
and Definition~\ref{def:anomaly-index}\eqref{it:anomaly-time-ordered}, applied to it, provides a permutation $\lambda'$ of this tuple
such that $c_{-\mathfrak{a}(N)}(\cdot\,;\lambda')$ does not vanish identically on the real points of $\mathcal{Z}(\mathrm{id},0)$.
These points lie in $\mathcal{O}_N$, where $\tilde c_{-\mathfrak{a}(N)}(\cdot\,;\lambda')=c_{-\mathfrak{a}(N)}(\cdot\,;\lambda')$ by \eqref{it:harmonic-analytic};
so $\tilde c_{-\mathfrak{a}(N)}(\cdot\,;\lambda')$ is not identically zero.
\end{proof}

When a configuration $X\in\mathcal{O}^*_N$ is turned, the order of its points in a frame changes.
At an angle $s$ at which two times of $R_sX$ coincide, no two of its abscissae coincide,
since $R_sX\in\mathcal{O}^*_N\subseteq\mathcal{O}_N$ (Lemma~\ref{lem:exact-phases}\eqref{it:generic-configurations} and Lemma~\ref{lem:one-function}\eqref{it:ordered-configurations});
so $R_sX$ is a real point of a tube of the frame $\frac{\pi}{2}$,
and the kernel of that frame, which agrees with the kernels of the frame $0$ wherever both are defined (Lemma~\ref{lem:one-function}\eqref{it:tube-gluing}),
carries the harmonic across the coincidence.
Of all rotations of the plane, the hypotheses on $\mu$ contain the quarter turn alone;
the kernels of all frames are formed with the reconstruction in the direction $e_1$, and the other angles are reached by analytic continuation.
The rotations move the positions and leave the labels unchanged,
and the phase in Lemma~\ref{lem:harmonic-extension}\eqref{it:harmonic-phase} is exact because $e^{+}$, $e^{-}$, $e^{3}$ are eigenvectors of every rotation $R_t$
(Lemma~\ref{lem:plate-identities}\eqref{it:plate-rotation}).

\begin{lemma}[Convexity]\label{lem:anomaly-convexity}
Let $\mu$ be a law on $\mathcal{S}'$ that satisfies \textup{(G)}, $(T_{123})$, $(R_1)$, $(RP_1)$, $(Q_{12})$ and $(C_{12})$.
Then
\[
2\,\mathfrak{a}(a+b)\leq\mathfrak{a}(2a)+\mathfrak{a}(2b)
\]
for all integers $a,b\geq0$.
\end{lemma}

\begin{proof}
Let $N:=a+b$ and $M:=\mathfrak{a}(N)$.
The anomaly index is non-negative and $\mathfrak{a}(0)=\mathfrak{a}(1)=0$;
so there is nothing to prove if $N\leq1$ or $M=0$, and one may assume $N\geq2$ and $M\geq1$.
By Definition~\ref{def:anomaly-index}\eqref{it:anomaly-both-signs} and \eqref{it:anomaly-time-ordered}
there are a tuple $\lambda$ of $N$ labels and a real point $X=(y_1,\dots,y_N)$ of $\mathcal{Z}(\mathrm{id},0)$,
that is, a configuration with $y_1^0<\dots<y_N^0$, such that $c_{-M}(X;\lambda)\neq0$.
The functions $c_m$ are invariant under translations (Lemma~\ref{lem:finitely-many-harmonics}), and a translation in time keeps $X$ strictly time-ordered;
one may therefore assume $y_a^0<0<y_{a+1}^0$,
where only $0<y_1^0$ is required if $a=0$, and only $y_N^0<0$ if $b=0$.
Let
\[
x_k:=\theta_1y_{a+1-k},\quad\lambda^x_k:=\lambda_{a+1-k}\quad(1\leq k\leq a),\qquad
z_l:=y_{a+l},\quad\lambda^z_l:=\lambda_{a+l}\quad(1\leq l\leq b),
\]
and let $x:=(x_1,\dots,x_a)$, $z:=(z_1,\dots,z_b)$, $\lambda^x:=(\lambda^x_1,\dots,\lambda^x_a)$ and $\lambda^z:=(\lambda^z_1,\dots,\lambda^z_b)$.
The times $x_k^0=-y^0_{a+1-k}$ are positive and increasing in $k$, and so are the times $z_l^0=y^0_{a+l}$ in $l$;
hence $x$ and $z$ are real points of $\mathcal{W}^{(a)}$ and of $\mathcal{W}^{(b)}$ (Definition~\ref{def:plate-kernels}).
The weight splits as $w=w_x+w_z$, with $w_x:=w(\lambda^x)$ and $w_z:=w(\lambda^z)$;
an empty tuple of labels, such as $\lambda^x$ for $a=0$ and $\lambda^z$ for $b=0$, has the weight $0$ and the sign $1$.
For real $\beta$ one has $\Lambda_\beta\mathcal{W}^{(n)}=\mathcal{W}^{(n)}$ and $(\Lambda_\beta\xi)^*=\Lambda_\beta(\xi^*)$ by Lemma~\ref{lem:forward-tube}\eqref{it:tube-boosts},
and $\xi^*=-\theta_1\xi$ for real $\xi$ by Lemma~\ref{lem:forward-tube}\eqref{it:tube-adjoint}.
Therefore
\[
(x_{k+1}-x_k)^*=y_{a+1-k}-y_{a-k}\quad(1\leq k<a),\qquad x_1^*+z_1=y_{a+1}-y_a\quad\text{if }a\geq1\text{ and }b\geq1.
\]
For every real $\beta$, Proposition~\ref{prop:kernels}\eqref{it:plate-vectors}, applied to $\Lambda_\beta x\in\mathcal{W}^{(a)}$ and $\Lambda_\beta z\in\mathcal{W}^{(b)}$, gives
\begin{align}
\langle\Psi(\Lambda_\beta x;\lambda^x),\Psi(\Lambda_\beta z;\lambda^z)\rangle
&=s(\lambda^x)\cdot G_N(\Lambda_\beta(y_2-y_1),\dots,\Lambda_\beta(y_N-y_{N-1});\lambda_1,\dots,\lambda_N)\notag\\
&=s(\lambda^x)\cdot\mathcal{G}_N(\Lambda_\beta X;\lambda).\label{eq:rotation-split}
\end{align}
If $a\geq1$ and $b\geq1$, the first equality is the Gram formula \eqref{eq:rotation-gram} of that item:
by the identities for $(x_{k+1}-x_k)^*$ and $x_1^*+z_1$, by $z_{l+1}-z_l=y_{a+l+1}-y_{a+l}$ and by $(\Lambda_\beta\xi)^*=\Lambda_\beta(\xi^*)$, its arguments are $\Lambda_\beta(y_2-y_1),\dots,\Lambda_\beta(y_N-y_{N-1})$,
and its labels $\lambda^x_a,\dots,\lambda^x_1,\lambda^z_1,\dots,\lambda^z_b$ are $\lambda_1,\dots,\lambda_N$.
If $a=0$, then $x$ and $\lambda^x$ are empty, $\Psi(\cdot\,;\lambda^x)=\Omega_1$, $z=X$ and $\lambda^z=\lambda$,
and the first equality is the formula of the same item for $\langle\Omega_1,\Psi(\Lambda_\beta z;\lambda^z)\rangle$.
If $b=0$, then $N=a$ and $\Psi(\cdot\,;\lambda^z)=\Omega_1$,
and the first equality is the formula of the same item for $\langle\Psi(\Lambda_\beta x;\lambda^x),\Omega_1\rangle$:
its arguments $(\Lambda_\beta x_a-\Lambda_\beta x_{a-1})^*,\dots,(\Lambda_\beta x_2-\Lambda_\beta x_1)^*$ are, by the identity for $(x_{k+1}-x_k)^*$ and by $(\Lambda_\beta\xi)^*=\Lambda_\beta(\xi^*)$, the points $\Lambda_\beta(y_2-y_1),\dots,\Lambda_\beta(y_N-y_{N-1})$,
and its labels $\lambda^x_a,\dots,\lambda^x_1$ are $\lambda_1,\dots,\lambda_N$.
In the same way, the Gram formula \eqref{eq:rotation-gram} for the pair $(\Lambda_\beta x,\Lambda_\beta x)$ if $a\geq1$,
and for the pair $(\Lambda_\beta z,\Lambda_\beta z)$ if $b\geq1$, gives
\[
\|\Psi(\Lambda_\beta x;\lambda^x)\|^2=s(\lambda^x)\cdot\mathcal{G}_{2a}(\Lambda_\beta X';\lambda'),\qquad
\|\Psi(\Lambda_\beta z;\lambda^z)\|^2=s(\lambda^z)\cdot\mathcal{G}_{2b}(\Lambda_\beta X'';\lambda''),
\]
where
\begin{equation}\label{eq:rotation-doubled}
\begin{aligned}
X'&:=(\theta_1x_a,\dots,\theta_1x_1,x_1,\dots,x_a),&\lambda'&:=(\lambda_1,\dots,\lambda_a,\lambda_a,\dots,\lambda_1),\\
X''&:=(\theta_1z_b,\dots,\theta_1z_1,z_1,\dots,z_b),&\lambda''&:=(\lambda_N,\dots,\lambda_{a+1},\lambda_{a+1},\dots,\lambda_N);
\end{aligned}
\end{equation}
indeed, as $\xi^*=-\theta_1\xi$ for real $\xi$, the arguments of these two Gram formulas are the images under $\Lambda_\beta$ of the consecutive differences of $X'$ and of $X''$,
and their labels are $\lambda'$ and $\lambda''$.
Here $X$, $X'$ and $X''$ are strictly time-ordered, so that $\Lambda_\beta X$, $\Lambda_\beta X'$, $\Lambda_\beta X''$ lie in the tubes $\mathcal{Z}(\mathrm{id},0)$ (Lemma~\ref{lem:one-function}\eqref{it:tube-cones});
on these tubes the functions $\mathcal{G}_N$, $\mathcal{G}_{2a}$, $\mathcal{G}_{2b}$ are the kernels $G_N$, $G_{2a}$, $G_{2b}$ evaluated at the consecutive differences
(Definition~\ref{def:configuration-tubes} and Lemma~\ref{lem:one-function}\eqref{it:tube-gluing}).
The reflected points of $X'$ carry the labels of their mirror images, so that $\lambda'$ has the weight $2w_x$;
likewise $\lambda''$ has the weight $2w_z$.
If $a=0$, the first norm is $\|\Omega_1\|=1$, and if $b=0$, the second norm is $\|\Omega_1\|=1$,
because $\langle\Omega_1,\Omega_1\rangle=1$ by Proposition~\ref{prop:kernels}\eqref{it:plate-vectors}; moreover $\mathfrak{a}(0)=0$ by definition.

By the Cauchy--Schwarz inequality in $\mathcal{H}_1$, the squared modulus of \eqref{eq:rotation-split} is at most the product of the two squared norms.
The norms are non-negative, and $|s(\cdot)|=1$;
so the squared norms are the moduli of $\mathcal{G}_{2a}(\Lambda_\beta X';\lambda')$ and of $\mathcal{G}_{2b}(\Lambda_\beta X'';\lambda'')$.
Expand the three values of the functions $\mathcal{G}$ by \eqref{eq:rotation-boost-expansion}, with the weights $w$, $2w_x$, $2w_z$;
the coefficients $c_m(X';\lambda')$ vanish for $|m|>\mathfrak{a}(2a)$, and the coefficients $c_m(X'';\lambda'')$ for $|m|>\mathfrak{a}(2b)$ (Definition~\ref{def:anomaly-index}).
After multiplication by $\exp(2w\beta)=\exp(2w_x\beta)\cdot\exp(2w_z\beta)$ one obtains, for every real $\beta$,
\[
\begin{aligned}
&\left|\sum_mc_m(X;\lambda)\cdot\exp(-4m\beta)\right|^2\\
&\qquad\leq\left|\sum_{|m|\leq\mathfrak{a}(2a)}c_m(X';\lambda')\cdot\exp(-4m\beta)\right|\cdot\left|\sum_{|m|\leq\mathfrak{a}(2b)}c_m(X'';\lambda'')\cdot\exp(-4m\beta)\right|,
\end{aligned}
\]
where the first factor on the right is to be read as $1$ if $a=0$, and the second as $1$ if $b=0$.
Since $c_m(X;\lambda)=0$ for $|m|>M$ and $c_{-M}(X;\lambda)\neq0$,
the left side is at least $\frac{1}{2}\,|c_{-M}(X;\lambda)|^2\cdot\exp(8M\beta)$ for all sufficiently large $\beta$.
The right side is at most $C\cdot\exp(4(\mathfrak{a}(2a)+\mathfrak{a}(2b))\beta)$ for $\beta\geq0$, with a constant $C$ independent of $\beta$.
Letting $\beta\to+\infty$ gives $8M\leq4(\mathfrak{a}(2a)+\mathfrak{a}(2b))$.
\end{proof}

In \eqref{eq:rotation-doubled} each reflected point $\theta_1x_k$ carries the label of $x_k$, because $\Theta_1$ keeps the labels:
the reflection $\theta_1$ alone carries a plate into a multiple of a plate with the conjugate label, in which $+$ and $-$ are exchanged,
and the complex conjugation contained in $\Theta_1$ exchanges them back (Lemma~\ref{lem:plate-identities}\eqref{it:plate-reflection}).
As a consequence $\lambda'$ and $\lambda''$ have the weights $2w_x$ and $2w_z$,
and the single factor $\exp(2w\beta)$ removes the weights from both sides of the inequality at once;
the exponential rates $8M$ and $4(\mathfrak{a}(2a)+\mathfrak{a}(2b))$ compared in the limit $\beta\to+\infty$ contain the anomaly indices alone.

\begin{lemma}[Positivity of the top harmonic]\label{lem:top-harmonic}
Let $\mu$ be a law on $\mathcal{S}'$ that satisfies \textup{(G)}, $(T_{123})$, $(R_1)$, $(RP_1)$, $(Q_{12})$ and $(C_{12})$, let $a\geq1$, and put $M:=\mathfrak{a}(2a)$.
A labelled configuration of $a$ distinct points of $\mathbb{H}$ is a pair $(x;\lambda)$ of a tuple $x=(x_1,\dots,x_a)$ of pairwise distinct points of $\mathbb{H}$,
ties in time being allowed, and a tuple $\lambda=(\lambda_1,\dots,\lambda_a)$ of labels.
For two such pairs $(x;\lambda)$ and $(z;\lambda')$, the tuple $(\theta_1x,z):=(\theta_1x_1,\dots,\theta_1x_a,z_1,\dots,z_a)$ is a configuration of $2a$ distinct points,
$a$ of them in $-\mathbb{H}$ and $a$ of them in $\mathbb{H}$;
with the labels $(\lambda,\lambda'):=(\lambda_1,\dots,\lambda_a,\lambda'_1,\dots,\lambda'_a)$ put
\[
K_a((x;\lambda),(z;\lambda')):=s(\lambda)\cdot\tilde c^{(2a)}_{-M}(\theta_1x,z;\lambda,\lambda').
\]
\begin{enumerate}
\item\label{it:top-harmonic-kernel}
On the set of the labelled configurations of $a$ distinct points of $\mathbb{H}$, the kernel $K_a$ is Hermitian and positive semidefinite.
\item\label{it:top-harmonic-mirror}
Let $L$ be a line in the plane, $\theta_L$ the reflection in $L$, $L_+$ one of the two open half planes bounded by $L$, $v$ a point of $L$,
and $\upsilon$ a real number such that $n_\upsilon$ is the unit normal of $L$ that points into $L_+$.
Labelled configurations of $a$ distinct points of $L_+$ are defined as those of $\mathbb{H}$, with $L_+$ in place of $\mathbb{H}$.
For two of them, $(x;\lambda)$ and $(z;\lambda')$, the tuple $(\theta_Lx,z):=(\theta_Lx_1,\dots,\theta_Lx_a,z_1,\dots,z_a)$ is formed as $(\theta_1x,z)$, with $\theta_L$ in place of $\theta_1$;
it is a configuration of $2a$ distinct points, $a$ of them in $\theta_LL_+$ and $a$ of them in $L_+$,
because $\theta_L$ maps $L_+$ onto the other open half plane bounded by $L$, and it carries the labels $(\lambda,\lambda')$.
For such $(x;\lambda)$ and $(z;\lambda')$, let $x'$ and $z'$ be the tuples with the entries
$x'_k:=R_{-\upsilon}(x_k-v)$ and $z'_k:=R_{-\upsilon}(z_k-v)$;
they consist of pairwise distinct points of $\mathbb{H}$, and
\[
\tilde c^{(2a)}_{-M}(\theta_Lx,z;\lambda,\lambda')=\exp(i(w(\lambda)+w(\lambda')-4M)\upsilon)\cdot s(\lambda)\cdot K_a((x';\lambda),(z';\lambda')).
\]
Consequently
\[
\left|\tilde c^{(2a)}_{-M}(\theta_Lx,z;\lambda,\lambda')\right|^2
\leq\left|\tilde c^{(2a)}_{-M}(\theta_Lx,x;\lambda,\lambda)\right|\cdot\left|\tilde c^{(2a)}_{-M}(\theta_Lz,z;\lambda',\lambda')\right|;
\]
in particular, if $\tilde c^{(2a)}_{-M}(\theta_Lx,x;\lambda,\lambda)=0$, then $\tilde c^{(2a)}_{-M}(\theta_Lx,z;\lambda,\lambda')=0$
for every labelled configuration $(z;\lambda')$ of $a$ distinct points of $L_+$.
\end{enumerate}
\end{lemma}

\begin{proof}
\eqref{it:top-harmonic-kernel}
Consider first labelled configurations $(x;\lambda)$ and $(z;\lambda')$ whose times are pairwise distinct.
A permutation of the points of $x$ together with their labels changes neither $s(\lambda)$ nor, by Lemma~\ref{lem:harmonic-extension}\eqref{it:harmonic-symmetric}, the value of $\tilde c^{(2a)}_{-M}$;
the same holds for $z$.
One may therefore assume $0<x_1^0<\dots<x_a^0$ and $0<z_1^0<\dots<z_a^0$, so that $x$ and $z$ are real points of $\mathcal{W}^{(a)}$,
and the configuration $(\theta_1x_a,\dots,\theta_1x_1,z_1,\dots,z_a)$ is strictly time-ordered.
The computation that leads to \eqref{eq:rotation-split}, together with \eqref{eq:rotation-boost-expansion} and Lemma~\ref{lem:harmonic-extension}\eqref{it:harmonic-analytic}, \eqref{it:harmonic-symmetric},
gives for every real $\beta$
\[
\begin{aligned}
\langle\Psi(\Lambda_\beta x;\lambda),\Psi(\Lambda_\beta z;\lambda')\rangle
&=s(\lambda)\cdot\mathcal{G}_{2a}(\Lambda_\beta(\theta_1x_a,\dots,\theta_1x_1,z_1,\dots,z_a);\lambda_a,\dots,\lambda_1,\lambda'_1,\dots,\lambda'_a)\\
&=s(\lambda)\cdot\sum_{|m|\leq M}\tilde c^{(2a)}_m(\theta_1x,z;\lambda,\lambda')\cdot\exp(-(w(\lambda)+w(\lambda')+4m)\beta);
\end{aligned}
\]
the terms with $|m|>M$ are absent because of the definition of $M=\mathfrak{a}(2a)$.
As a function of the pair $((x;\lambda),(z;\lambda'))$, the left side is a Gram kernel, hence positive semidefinite.
Multiply it by $\exp(w(\lambda)\beta)\cdot\exp(w(\lambda')\beta)$, the product of a positive function of $(x;\lambda)$ and the same function of $(z;\lambda')$,
and by the positive number $\exp(-4M\beta)$;
by Lemma~\ref{lem:analysis-facts}\eqref{it:positive-kernels} the resulting kernels
\[
s(\lambda)\cdot\sum_{|m|\leq M}\tilde c^{(2a)}_m(\theta_1x,z;\lambda,\lambda')\cdot\exp(-4(m+M)\beta)
\]
are positive semidefinite.
As $\beta\to+\infty$ the terms with $m>-M$ tend to $0$, and the kernels converge pointwise to $K_a$.
By the same fact, $K_a$ is positive semidefinite on the labelled configurations whose times are pairwise distinct.

That $K_a$ is positive semidefinite on all labelled configurations means that, for every finite family of them,
the matrix of the values of $K_a$ at the pairs of members of the family is positive semidefinite.
Among the configurations of $a$ distinct points of $\mathbb{H}$, those with pairwise distinct times are dense,
and $\tilde c^{(2a)}_{-M}$ is continuous on $(\R^2)^{2a}_{\neq}$ by Lemma~\ref{lem:harmonic-extension}\eqref{it:harmonic-analytic}.
Approximating each member of a finite family by configurations with pairwise distinct times and the same labels,
one obtains the matrix of the family as a limit of positive semidefinite matrices, and it is positive semidefinite.
A positive semidefinite kernel is Hermitian.

\eqref{it:top-harmonic-mirror}
Let $\tau_v$ be the translation $y\mapsto y+v$ of the plane.
The map $R_{-\upsilon}\tau_{-v}$ carries $L$ onto the line $\{y^0=0\}$ and $L_+$ onto $\mathbb{H}$, because $R_{-\upsilon}n_\upsilon=n_0$;
hence $\theta_L=\tau_vR_\upsilon\theta_1R_{-\upsilon}\tau_{-v}$.
In particular the time of $x'_k$ is $n_\upsilon\cdot(x_k-v)>0$, and $x'$ and $z'$ consist of pairwise distinct points of $\mathbb{H}$.
Moreover $\theta_Lx_k=v+R_\upsilon\theta_1x'_k$ and $z_k=v+R_\upsilon z'_k$, that is,
\[
(\theta_Lx,z)=(v,\dots,v)+R_\upsilon(\theta_1x',z').
\]
The labels $(\lambda,\lambda')$ have the weight $w(\lambda)+w(\lambda')$,
and Lemma~\ref{lem:harmonic-extension}\eqref{it:harmonic-translation}, \eqref{it:harmonic-phase}, applied with $m=-M$, give
\[
\tilde c^{(2a)}_{-M}(\theta_Lx,z;\lambda,\lambda')=\exp(i(w(\lambda)+w(\lambda')-4M)\upsilon)\cdot\tilde c^{(2a)}_{-M}(\theta_1x',z';\lambda,\lambda').
\]
Since $s(\lambda)^2=1$, the last factor is $s(\lambda)\cdot K_a((x';\lambda),(z';\lambda'))$, which proves the identity.
The exponential factor has modulus $1$.
For $z=x$ and $\lambda'=\lambda$ the identity gives $\left|\tilde c^{(2a)}_{-M}(\theta_Lx,x;\lambda,\lambda)\right|=K_a((x';\lambda),(x';\lambda))$,
the right side being non-negative by \eqref{it:top-harmonic-kernel};
in the same way $\left|\tilde c^{(2a)}_{-M}(\theta_Lz,z;\lambda',\lambda')\right|=K_a((z';\lambda'),(z';\lambda'))$.
The asserted inequality is therefore the inequality of Lemma~\ref{lem:analysis-facts}\eqref{it:positive-kernels} for the kernel $K_a$
at the two labelled configurations $(x';\lambda)$ and $(z';\lambda')$, and the last assertion is a consequence of it.
\end{proof}

The hypotheses contain reflection positivity of $\mu$ at the plane $\{x_1=0\}$ only, and no positivity at an oblique line of the plane;
part~\eqref{it:top-harmonic-mirror} of Lemma~\ref{lem:top-harmonic} derives a positivity at every line, for the top harmonic alone.
The unimodular factor in that part depends on the weights of the labels of $x$ and of $z$;
if all labels have the weight $0$, as in the case of a scalar field, it is the constant $\exp(-4iM\upsilon)$.
The inequality of that part involves its modulus only.

\subsubsection{Polygons and arithmetic}\label{sec:polygons}

For an integer $n\ge0$, a \emph{labelled set of $n$ distinct points} is a set of $n$ pairs $(y,\lambda)$, each formed by a point $y$ of the plane and a label $\lambda$, whose first entries are pairwise distinct
(for $n=0$ it is the empty set);
its elements are its labelled points, and their first entries are its points.
Let $(X;\lambda)$ be a labelled configuration of $N$ distinct points, with $X=(y_1,\dots,y_N)$ and $\lambda=(\lambda_1,\dots,\lambda_N)$.
The labelled set of $(X;\lambda)$ is $\{(y_1,\lambda_1),\dots,(y_N,\lambda_N)\}$, a labelled set of $N$ distinct points.
By Lemma~\ref{lem:harmonic-extension}\eqref{it:harmonic-symmetric}, the number $\tilde c_m(X;\lambda)$ depends on the labelled set of $(X;\lambda)$ only;
accordingly, $\tilde c_m$ is evaluated at labelled sets of $N$ distinct points as well.
A labelled configuration contains a labelled set $S$ if $S$ is a subset of its labelled set.
For a line $L$ of the plane, $\theta_L$ is the reflection in $L$;
the two open half planes bounded by $L$ are the open sides of $L$, and when one of them is called $L_+$, the other is called $L_-$.
The mirror image of a labelled set $S$ in $L$, with the same labels, is
\[
\theta_LS:=\{(\theta_Ly,\lambda):(y,\lambda)\in S\}.
\]

\begin{lemma}[Selection and halving]\label{lem:selection-halving}
Let $\mu$ be a law on $\mathcal{S}'$ that satisfies \textup{(G)}, $(T_{123})$, $(R_1)$, $(RP_1)$, $(Q_{12})$ and $(C_{12})$;
the functions $\tilde c_m$ below are those of $\mu$.
\begin{enumerate}
\item\label{it:selection}
\emph{Selection.}
Let $N\ge2$, let $(X;\lambda)$ be a labelled configuration of $N$ distinct points, $X=(y_1,\dots,y_N)$, $\lambda=(\lambda_1,\dots,\lambda_N)$, and put $w:=w(\lambda)$.
Let $t$ be real, $v\in\R^2$, and let $\mathfrak{s}$ be a permutation of $\{1,\dots,N\}$ such that
\[
v+R_t(y_k-v)=y_{\mathfrak{s}(k)}
\quad\text{and}\quad
\lambda_{\mathfrak{s}(k)}=\lambda_k
\qquad(k=1,\dots,N),
\]
so that the rotation by the angle $t$ about $v$ permutes the labelled points.
Then
\[
(1-\exp(i(w+4m)t))\cdot\tilde c_m(X;\lambda)=0
\qquad(m\in\Z).
\]
In particular, let $a\ge1$, and let the $2a$ points of $X$ be the vertices of a regular $2a$-gon, all with one and the same label $(\sigma,\chi)$;
then $\tilde c^{(2a)}_m(X;\lambda)=0$ for every integer $m$ such that $a$ does not divide $2m$.
\item\label{it:halving}
\emph{Halving.}
Let $a\ge1$, $M:=\mathfrak{a}(2a)$ and $c_*:=\tilde c^{(2a)}_{-M}$.
For a labelled set $S$ of at most $2a$ distinct points let $\mathcal{N}(S)$ be the statement that $c_*$ vanishes at every labelled configuration of $2a$ distinct points that contains $S$.
Let $L$ be a line of the plane, $L_+$ one of its open sides, $S'$ a labelled set of $j'\le a$ distinct points of $L_+$, and $S$ a subset of $\theta_LS'\cup S'$.
Then $\mathcal{N}(S)$ implies $\mathcal{N}(S')$ and $\mathcal{N}(\theta_LS')$.
\end{enumerate}
\end{lemma}

\begin{proof}
\eqref{it:selection}
By Lemma~\ref{lem:harmonic-extension}\eqref{it:harmonic-translation} and~\eqref{it:harmonic-phase}, applied to the configuration $X-(v,\dots,v)$,
\[
\tilde c_m((v,\dots,v)+R_t(X-(v,\dots,v));\lambda)
=\exp(i(w+4m)t)\cdot\tilde c_m(X-(v,\dots,v);\lambda)
=\exp(i(w+4m)t)\cdot\tilde c_m(X;\lambda).
\]
The $k$-th point of the configuration on the left is $v+R_t(y_k-v)=y_{\mathfrak{s}(k)}$, so that this configuration is $\mathfrak{s}X$;
moreover $\mathfrak{s}\lambda=\lambda$.
By Lemma~\ref{lem:harmonic-extension}\eqref{it:harmonic-symmetric},
\[
\tilde c_m(\mathfrak{s}X;\lambda)=\tilde c_m(\mathfrak{s}X;\mathfrak{s}\lambda)=\tilde c_m(X;\lambda),
\]
and the two identities together give the assertion.

For the polygon, number the vertices counterclockwise.
The rotation by the angle $\pi/a$ about the centre maps each vertex to the next one, and all labels coincide;
so the first part applies with $N=2a$ and $t=\pi/a$.
Here $w=2a\cdot w(\sigma)$, and $w(\sigma)$ is an integer; hence
\[
\exp\left(\frac{i(w+4m)\pi}{a}\right)
=\exp(2\pi i\cdot w(\sigma))\cdot\exp\left(\frac{4\pi im}{a}\right)
=\exp\left(\frac{4\pi im}{a}\right),
\]
so that the weight disappears from the condition.
The last number equals $1$ if and only if $2m/a$ is an integer, that is, if and only if $a$ divides $2m$.

\eqref{it:halving}
Let $\lambda''$ be a tuple of $a-j'$ labels and $\lambda'$ a tuple of $a$ labels.
For $y=(y_1,\dots,y_{a-j'})\in(\R^2)^{a-j'}$ and $z=(z_1,\dots,z_a)\in(\R^2)^a$ with pairwise distinct entries,
$(y;\lambda'')$ and $(z;\lambda')$ denote the labelled sets $\{(y_k,\lambda''_k):1\le k\le a-j'\}$ and $\{(z_l,\lambda'_l):1\le l\le a\}$.
Let $D$ be the set of the pairs $(y,z)$ for which the $2a-j'$ points $y_1,\dots,y_{a-j'},z_1,\dots,z_a$ are pairwise distinct and distinct from the points of $\theta_LS'$.

\emph{Points on opposite sides.}
Let $(y,z)\in D$, with $y_1,\dots,y_{a-j'}$ in $L_-$ and $z_1,\dots,z_a$ in $L_+$.
The labelled set $R:=\{(\theta_Ly_k,\lambda''_k):1\le k\le a-j'\}$ consists of $a-j'$ distinct points of $L_+$, none of which is a point of $S'$;
so $x:=S'\cup R$ is a labelled set of $a$ distinct points of $L_+$.
The $2a$ points of $\theta_Lx\cup x$ are distinct, $a$ of them lying in $L_-$ and $a$ in $L_+$,
and this labelled set contains $\theta_LS'\cup S'$, which contains $S$.
By $\mathcal{N}(S)$, the function $c_*$ vanishes at $\theta_Lx\cup x$.
Lemma~\ref{lem:top-harmonic}(b), applied to $x$, with its points listed in any order, and to $(z;\lambda')$, shows that $c_*$ vanishes at $\theta_Lx\cup(z;\lambda')$.
As $\theta_LR=(y;\lambda'')$, this labelled set is $\theta_LS'\cup(y;\lambda'')\cup(z;\lambda')$.

\emph{All positions.}
The set $D$ is the complement in $(\R^2)^{2a-j'}$ of finitely many affine subspaces of codimension $2$,
namely those defined by an equation $y_k=y_{k'}$, $z_l=z_{l'}$ or $y_k=z_l$, or by the equality of one $y_k$ or one $z_l$ with a given point of $\theta_LS'$.
By Lemma~\ref{lem:analysis-facts}\eqref{it:connected-complement}, $D$ is open and connected.
The function
\[
(y,z)\mapsto c_*(\theta_LS'\cup(y;\lambda'')\cup(z;\lambda'))
\]
is the composition of $c_*$ with an affine map, and it is real analytic on $D$ by Lemma~\ref{lem:harmonic-extension}\eqref{it:harmonic-analytic}.
It vanishes on the open subset of $D$ on which all $y_k$ lie in $L_-$ and all $z_l$ in $L_+$, which is not empty.
By Lemma~\ref{lem:analysis-facts}\eqref{it:analytic-zeros}, it vanishes on $D$.
Every labelled configuration of $2a$ distinct points that contains $\theta_LS'$ has the labelled set $\theta_LS'\cup(y;\lambda'')\cup(z;\lambda')$ for suitable tuples $\lambda''$, $\lambda'$ and some $(y,z)\in D$;
since $\lambda''$ and $\lambda'$ were arbitrary, $\mathcal{N}(\theta_LS')$ holds.

\emph{The other side.}
The labelled set $S'':=\theta_LS'$ consists of $j'$ distinct points of the open side $L_-$, its mirror image is $\theta_LS''=S'$, and $S$ is a subset of $\theta_LS''\cup S''$.
The argument above, with $L_-$ in place of $L_+$ and $S''$ in place of $S'$, gives $\mathcal{N}(\theta_LS'')$, which is $\mathcal{N}(S')$.
\end{proof}

\begin{proposition}[Regular polygons]\label{prop:polygons}
Let $\mu$ be a law on $\mathcal{S}'$ that satisfies \textup{(G)}, $(T_{123})$, $(R_1)$, $(RP_1)$, $(Q_{12})$ and $(C_{12})$.
Then $a$ divides $2\,\mathfrak{a}(2a)$ for every integer $a\ge1$.
\end{proposition}

\begin{proof}
Suppose that for some integer $a\ge1$ the number $a$ does not divide $2M$, where $M:=\mathfrak{a}(2a)$,
and let $c_*$ and $\mathcal{N}$ be as in Lemma~\ref{lem:selection-halving}\eqref{it:halving}.
Then $M\ge1$, and $a\ge3$, because $1$ and $2$ divide $2M$.
Fix a label $(\sigma,\chi)$.
For a regular $2a$-gon of the plane, of arbitrary centre, size and orientation, let $P_k$, $k\in\Z/2a\Z$, be its vertices in cyclic order.
For $k_0\in\Z/2a\Z$ and $1\le j\le2a$, the labelled set
\[
\{(P_k,(\sigma,\chi)):k=k_0+1,\dots,k_0+j\}
\]
is an arc of length $j$.
For $1\le j\le2a$ let $(\mathcal{N}_j)$ be the statement that $\mathcal{N}(S)$ holds for every arc $S$ of length $j$ of every regular $2a$-gon.

\emph{The statement $(\mathcal{N}_{2a})$ holds.}
An arc of length $2a$ consists of all vertices of the polygon, each with the label $(\sigma,\chi)$,
and the only labelled set of $2a$ distinct points that contains it is the arc itself.
By Lemma~\ref{lem:selection-halving}\eqref{it:selection} with $m=-M$, the function $c_*$ vanishes there, since $a$ does not divide $2M$.

\emph{For $2\le j\le2a$, the statement $(\mathcal{N}_j)$ implies $(\mathcal{N}_{j'})$, where $j'$ is the least integer $\ge j/2$.}
One has $j'\le a$ and $j\le2j'\le2a$.
Let $S'=\{(P_k,(\sigma,\chi)):k=k_0+1,\dots,k_0+j'\}$ be an arc of length $j'$ of a regular $2a$-gon,
and let $L$ be the perpendicular bisector of the edge $[P_{k_0},P_{k_0+1}]$.
The line $L$ contains the centre, the reflection $\theta_L$ maps $P_{k_0+k}$ to $P_{k_0+1-k}$ for every $k$, and no vertex lies on $L$.
Measured from the half line that starts at the centre and meets the midpoint of the edge $[P_{k_0},P_{k_0+1}]$,
the vertex $P_{k_0+k}$ has the polar angle $(2k-1)\pi/(2a)$, which lies in $(0,\pi)$ for $k=1,\dots,a$;
so these $a$ vertices lie in one open side $L_+$ of $L$.
Hence $S'$ is a labelled set of $j'$ distinct points of $L_+$, and
\[
\theta_LS'\cup S'=\{(P_k,(\sigma,\chi)):k=k_0-j'+1,\dots,k_0+j'\}
\]
is an arc of length $2j'\ge j$.
Its first $j$ vertices $P_{k_0-j'+1},\dots,P_{k_0-j'+j}$, with the label $(\sigma,\chi)$, form an arc $S$ of length $j$ contained in $\theta_LS'\cup S'$,
and $\mathcal{N}(S)$ holds by $(\mathcal{N}_j)$.
Lemma~\ref{lem:selection-halving}\eqref{it:halving} gives $\mathcal{N}(S')$.

\emph{Conclusion.}
Starting from $j=2a$ and replacing $j$ by $j'$ repeatedly, one reaches $j=1$ after finitely many steps, because $j'<j$ whenever $j\ge2$;
hence $(\mathcal{N}_1)$ holds.
Since any point of the plane is a vertex of some regular $2a$-gon, and an arc of length $1$ consists of one vertex with the label $(\sigma,\chi)$,
the statement $(\mathcal{N}_1)$ says that $c_*$ vanishes at every labelled configuration of $2a$ distinct points one of whose labels is $(\sigma,\chi)$.
As the label $(\sigma,\chi)$ was arbitrary, $\tilde c^{(2a)}_{-M}(\cdot;\lambda)=0$ for every tuple $\lambda$ of $2a$ labels.
This contradicts Lemma~\ref{lem:harmonic-extension}\eqref{it:harmonic-top}, because $M\ge1$.
\end{proof}

Two facts underlie this argument.
The functions $\tilde c_m$ are exactly symmetric under permutations of the labelled points, because the random variables $\mathbb{B}(b)$ commute.
And the labels $+$, $-$, $3$ belong to the eigenvectors $e^{+}$, $e^{-}$, $e^{3}$ of the rotations;
this makes the property that all vertices carry the same label stable under all reflections of the plane,
so that an arc and its mirror image in the proof above form an arc again.
For the components $\mathbb{B}_1$, $\mathbb{B}_2$, a reflection in a line that is not parallel to a coordinate axis would mix the labels.

\begin{lemma}[Convex integer sequences]\label{lem:convex-sequences}
Let $\nu:\{0,1,2,\dots\}\to\{0,1,2,\dots\}$ satisfy $\nu(0)=0$ and, for some constant $C$,
\[
2\nu(a)\le\nu(a-1)+\nu(a+1)\quad(a\ge1),
\qquad
\nu(a)\le C\cdot a\quad(a\ge0),
\]
and suppose that $a$ divides $2\nu(a)$ for every $a\ge1$.
Then $\nu(a)=a\cdot\nu(1)$ for every $a\ge0$.
\end{lemma}

\begin{proof}
The increments $d_a:=\nu(a)-\nu(a-1)$, $a\ge1$, are integers, and by the first inequality they are non-decreasing in $a$.
If $d_{a'}>C$ for some $a'\ge1$, then for $a\ge a'$
\[
\nu(a)=\nu(a'-1)+\sum_{k=a'}^{a}d_k\ge(a-a'+1)\cdot d_{a'},
\]
and the right-hand side exceeds $C\cdot a$ for large $a$, against the second inequality.
So $d_1\le d_a\le C$ for every $a\ge1$.
Being bounded, non-decreasing and integer-valued, the sequence $(d_a)$ is eventually constant:
there are an integer $a_1\ge1$ and an integer $d$ with $d_a=d$ for all $a\ge a_1$.
The number
\[
c:=\sum_{a=1}^{a_1-1}(d-d_a)
\]
is a non-negative integer, since $d_a\le d$ for every $a$, and $\nu(a)=\sum_{k=1}^{a}d_k=d\cdot a-c$ for $a\ge a_1$.
Hence $a$ divides $2\nu(a)=2d\cdot a-2c$, and therefore $2c$, for every $a\ge a_1$;
this forces $c=0$.
It follows that $d_a=d$ for all $a\ge1$, that $d=d_1=\nu(1)$, and that $\nu(a)=a\cdot\nu(1)$.
\end{proof}

\begin{proposition}[The anomaly index vanishes]\label{prop:anomaly-vanishes}
Let $\mu$ be a law on $\mathcal{S}'$ that satisfies \textup{(G)}, \textup{(dF)}, $(T_{123})$, $(R_1)$, $(RP_1)$, $(Q_{12})$ and $(C_{12})$.
Then $\mathfrak{a}(N)=0$ for every $N$.
In other words, for every $N\ge2$ and every tuple $\lambda$ of $N$ labels, with $w=w(\lambda)$,
\[
\tilde G_N(R_tX;\lambda)=\exp(iwt)\cdot\tilde G_N(X;\lambda)
\qquad\text{for all }X\in(\R^2)^N_{\neq}\text{ and all real }t,
\]
\[
\mathcal{G}_N(\Lambda_\beta Z;\lambda)=\exp(-w\beta)\cdot\mathcal{G}_N(Z;\lambda)
\qquad\text{for all }Z\in\mathcal{Z}_N\text{ and all real }\beta.
\]
If $\mu$ satisfies \textup{(G)}, $(T_{123})$, $(R_1)$, $(RP_1)$, $(Q_{12})$ and $(C_{12})$, but not necessarily \textup{(dF)}, then
\[
\mathfrak{a}(2a)=a\cdot\mathfrak{a}(2)
\quad\text{and}\quad
\mathfrak{a}(N)\le\frac{N}{2}\cdot\mathfrak{a}(2)
\qquad\text{for all integers }a,N\ge0;
\]
so the two-point functions decide whether the anomaly index vanishes at every order.
\end{proposition}

\begin{proof}
\emph{Without closedness.}
Let $\mu$ satisfy the six conditions of the last assertion.
The sequence $\nu(a):=\mathfrak{a}(2a)$, $a\ge0$, takes its values in $\{0,1,2,\dots\}$ and satisfies the hypotheses of Lemma~\ref{lem:convex-sequences}:
$\nu(0)=\mathfrak{a}(0)=0$;
the inequality $2\nu(a)\le\nu(a-1)+\nu(a+1)$, $a\ge1$, is Lemma~\ref{lem:anomaly-convexity} for the pair of integers $a-1$, $a+1$;
the linear bound holds with $C=1$, by Definition~\ref{def:anomaly-index}\eqref{it:anomaly-linear-bound};
and $a$ divides $2\nu(a)$ by Proposition~\ref{prop:polygons}.
Hence $\mathfrak{a}(2a)=a\cdot\mathfrak{a}(2)$ for every $a\ge0$.
For $N\ge0$ put $a:=\lfloor N/2\rfloor$ and $b:=N-a$; Lemma~\ref{lem:anomaly-convexity} gives
\[
2\,\mathfrak{a}(N)\le\mathfrak{a}(2a)+\mathfrak{a}(2b)=(a+b)\cdot\mathfrak{a}(2)=N\cdot\mathfrak{a}(2).
\]

\emph{Two points.}
Let now $\mu$ satisfy \textup{(dF)} as well.
By Definition~\ref{def:closedness-cones}, $(C_{12})$ implies $(C^1_{12})$,
and by Theorem~\ref{thm:two-point}\eqref{it:quarter-turn-two-point}, $(Q_{12})$ gives the invariance of $S_2$ under the quarter turn.
So Theorem~\ref{thm:two-point}(a) applies:
\[
S_2(R_t\cdot b,R_t\cdot b')=S_2(b,b')
\qquad(t\in\R,\ b,b'\in L^2(\R^3;\C^3)).
\]
Let $\lambda=(\lambda_1,\lambda_2)$ be a pair of labels, $w:=w(\lambda)$,
and let $X_0=(x_1,x_2)$ be a real point of $\mathcal{Z}(\mathrm{id},0)$, so that $x_1^0<x_2^0$ by Lemma~\ref{lem:one-function}(a).
The time $(R_ty)^0$ depends continuously on $(t,y)$;
so there are $r>0$ and $t_0\in(0,\pi)$ such that $(R_ty)^0<(R_ty')^0$ whenever $|t|\le t_0$, $|y-x_1|\le r$ and $|y'-x_2|\le r$.
Let $U\subset(\R^2)^2$ be the product of the open discs of radius $r$ about $x_1$ and about $x_2$;
then $R_tX$ lies in $\mathcal{Z}(\mathrm{id},0)$ for all $X\in U$ and $|t|\le t_0$.
Let $\alpha_1,\alpha_2\in L^2(\R^2;\C)$ have compact supports in the first and in the second of these discs, and let $|t|\le t_0$.
The supports of $\alpha_1$, $\alpha_2$ are time-ordered, and so are the supports of $\alpha_1\circ R_{-t}$, $\alpha_2\circ R_{-t}$, which are their images under $R_t$.
Lemma~\ref{lem:one-function}(b) and~(c), for the identity permutation and the frame $0$, the invariance just stated,
and Lemma~\ref{lem:plate-identities}\eqref{it:plate-rotation} give, with $Y=(y_1,y_2)$,
\[
\begin{aligned}
\int\mathrm{d}Y\,\alpha_1(y_1)\cdot\alpha_2(y_2)\cdot\mathcal{G}_2(Y;\lambda)
&=S_2(\alpha_1\otimes\lambda_1,\alpha_2\otimes\lambda_2)
=S_2(R_t\cdot(\alpha_1\otimes\lambda_1),R_t\cdot(\alpha_2\otimes\lambda_2))\\
&=\exp(-iwt)\cdot S_2((\alpha_1\circ R_{-t})\otimes\lambda_1,(\alpha_2\circ R_{-t})\otimes\lambda_2)\\
&=\exp(-iwt)\cdot\int\mathrm{d}Y\,\alpha_1(R_{-t}y_1)\cdot\alpha_2(R_{-t}y_2)\cdot\mathcal{G}_2(Y;\lambda)\\
&=\exp(-iwt)\cdot\int\mathrm{d}Y\,\alpha_1(y_1)\cdot\alpha_2(y_2)\cdot\mathcal{G}_2(R_tY;\lambda),
\end{aligned}
\]
where the last step is the substitution $Y\mapsto R_tY$, which preserves the Lebesgue measure.
The functions $X\mapsto\mathcal{G}_2(X;\lambda)$ and $X\mapsto\exp(-iwt)\cdot\mathcal{G}_2(R_tX;\lambda)$ are continuous on $U$,
and the identity holds for all bounded $\alpha_1$, $\alpha_2$ with compact supports in the two discs.
By Lemma~\ref{lem:analysis-facts}\eqref{it:equal-densities},
\[
\mathcal{G}_2(R_tX;\lambda)=\exp(iwt)\cdot\mathcal{G}_2(X;\lambda)
\qquad(X\in U,\ |t|\le t_0).
\]
On the other hand, for $X\in U$ the interval $I_X$ of Lemma~\ref{lem:exact-phases}(a), formed with the identity permutation and the frame $0$, contains $[-t_0,t_0]$,
and $\mathcal{G}_2=\sum_mc_m$ on $\mathcal{Z}_2$ by Lemma~\ref{lem:finitely-many-harmonics}.
So Lemma~\ref{lem:exact-phases}(a) gives
\[
\mathcal{G}_2(R_tX;\lambda)=\sum_mc_m(R_tX;\lambda)=\sum_m\exp(i(w+4m)t)\cdot c_m(X;\lambda)
\qquad(X\in U,\ |t|\le t_0).
\]
For fixed $X\in U$ the two expressions show that the trigonometric polynomial $t\mapsto\sum_mc_m(X;\lambda)\cdot\exp(4imt)$,
a finite sum by Lemma~\ref{lem:finitely-many-harmonics}, is constant on $[-t_0,t_0]$.
By the identity theorem it is constant on $\R$,
and integration against $\exp(-4imt)$ over $[0,\pi/2]$ shows that $c_m(X;\lambda)=0$ for $m\ne0$.
The point $X_0$ and the pair $\lambda$ being arbitrary,
$c_m(\cdot;\lambda)$ vanishes on the real points of $\mathcal{Z}(\mathrm{id},0)$ for every $m\ne0$ and every pair of labels $\lambda$.
By Definition~\ref{def:anomaly-index}\eqref{it:anomaly-time-ordered}, $\mathfrak{a}(2)=0$.
With the first part, $\mathfrak{a}(N)\le(N/2)\cdot\mathfrak{a}(2)=0$ for every $N$.

\emph{The two identities.}
Let $N\ge2$, let $\lambda$ be a tuple of $N$ labels and $w=w(\lambda)$.
Since $\mathfrak{a}(N)=0$, Lemma~\ref{lem:harmonic-extension}\eqref{it:harmonic-top} gives $\tilde c_m(\cdot;\lambda)=0$ for $m\ne0$.
Hence $\tilde G_N(\cdot;\lambda)=\tilde c_0(\cdot;\lambda)$, and Lemma~\ref{lem:harmonic-extension}\eqref{it:harmonic-sum} gives
\[
\tilde G_N(R_tX;\lambda)=\exp(iwt)\cdot\tilde c_0(X;\lambda)=\exp(iwt)\cdot\tilde G_N(X;\lambda)
\qquad(X\in(\R^2)^N_{\neq},\ t\in\R).
\]
By the definition of $\mathfrak{a}(N)$, the functions $c_m(\cdot;\lambda)$ with $m\ne0$ vanish on $\mathcal{Z}_N$.
So \eqref{eq:rotation-boost-expansion} reads $\mathcal{G}_N(\Lambda_\beta Z;\lambda)=c_0(Z;\lambda)\cdot\exp(-w\beta)$,
and at $\beta=0$ it gives $c_0(\cdot;\lambda)=\mathcal{G}_N(\cdot;\lambda)$; this is the second identity.
Conversely, the second identity and the uniqueness of the coefficients in \eqref{eq:rotation-boost-expansion} give $c_m(\cdot;\lambda)=0$ for $m\ne0$, that is, $\mathfrak{a}(N)=0$.
\end{proof}

\begin{remark}[The role of the exponent one]\label{rem:exponent-one}
The window $|w+4m|\le N$ of Lemma~\ref{lem:finitely-many-harmonics} comes from the bound of Proposition~\ref{prop:kernels}\eqref{it:kernel-bound},
in which the distance to the boundary of the forward tube enters with the power $-N$, the exponent one for each of the $N$ plates.
\begin{enumerate}
\item[(1)]
\emph{Labels $3$ only.}
For $N\ge2$ let $\mathfrak{a}_3(N)$ be defined as $\mathfrak{a}(N)$, with the maximum taken only over the tuples of labels $(3,\chi_1),\dots,(3,\chi_N)$,
and let $\mathfrak{a}_3(0):=\mathfrak{a}_3(1):=0$;
these tuples describe the correlation functions of the component $\mathbb{B}_3=F_{12}$ alone.
The proof of Lemma~\ref{lem:anomaly-convexity} does not leave these tuples:
the labels of the doubled configurations in \eqref{eq:rotation-doubled} are entries of the given tuple, and permuting the labels or replacing $\sigma$ by $\bar\sigma$ keeps $\sigma=3$.
Read for these tuples, it gives $2\,\mathfrak{a}_3(a+b)\le\mathfrak{a}_3(2a)+\mathfrak{a}_3(2b)$ for all integers $a,b\ge0$.
For these tuples $w=0$, and the window gives $4\,\mathfrak{a}_3(N)\le N$.
So $\nu_3(a):=\mathfrak{a}_3(2a)$ is a sequence of non-negative integers with $2\nu_3(a)\le\nu_3(a-1)+\nu_3(a+1)$ for $a\ge1$,
with $\nu_3(a)\le a/2$, and with $\nu_3(0)=\nu_3(1)=0$.
If $\nu_3$ took a positive value, let $a_0$ be the least $a$ with $\nu_3(a)\ge1$;
then $\nu_3(a_0-1)=0$, by convexity the increments $\nu_3(a)-\nu_3(a-1)$ are at least $1$ for all $a\ge a_0$,
and $\nu_3(a)\ge a-a_0+1>a/2$ for large $a$, against the bound.
Hence $\nu_3=0$, and the inequality above with $a=\lfloor N/2\rfloor$, $b=N-a$ gives $\mathfrak{a}_3(N)=0$ for every $N$.
For the component $F_{12}$ alone the anomaly index therefore vanishes by the linear bound and convexity,
without regular polygons, without the Bianchi identity and without information on the two-point function.
By power counting, the field $\mathbb{B}$ has the scaling dimension $3/2$,
and the smearing in the transverse coordinate with a profile in $L^2$ lowers it by $1/2$, to the exponent one of each plate;
shared out plate by plate, the window thus suggests that a field of dimension $3/2$ has room for the frequency one only, and none for a harmonic of frequency $4$.
In this sector this heuristic is exact at all orders.
\item[(2)]
\emph{Labels $+$ and $-$.}
When the labels $+$ and $-$ occur, the window yields only $\nu(a)=\mathfrak{a}(2a)\le a$.
Each sequence $\nu(a)=(a-a_0)_+:=\max\{a-a_0,0\}$ with an integer $a_0\ge1$ is convex, has $\nu(0)=\nu(1)=0$ and satisfies $\nu(a)\le a$.
Of convexity, the linear bound and divisibility, only the divisibility of Proposition~\ref{prop:polygons} excludes it;
for $a_0=1$ divisibility fails first at $a=3$, as $3$ does not divide $4$.
So for the components of weight $\pm1$ the exponent one brings no simplification:
the tensor field leaves no slack here, just as in Yang--Mills theory.
\item[(3)]
\emph{Two points.}
For $N=2$ the window $|w+4m|\le2$ leaves $m\ne0$ only for the pairs of labels with $\sigma_1=\sigma_2=+$, where $m=-1$,
and with $\sigma_1=\sigma_2=-$, where $m=1$.
The diagonal Maxwell laws of Proposition~\ref{prop:sharpness}(d) realise these harmonics;
they satisfy all hypotheses of Proposition~\ref{prop:anomaly-vanishes} except \textup{(dF)},
and for them $\mathfrak{a}(2a)=a$ for every $a$: the linear bound is attained.
\end{enumerate}
\end{remark}

\subsection{No singular part}\label{sec:singular-part}

Let $\mu$ be a law as in Proposition~\ref{prop:no-singular-part} below,
let $N\geq2$, and let $\lambda=(\lambda_1,\dots,\lambda_N)$ be a tuple of labels.
The time of a point $y=(y^0,y^1)$ of the plane is $y^0$, its abscissa is $y^1$.
By Lemma~\ref{lem:one-function}(b), (c) and Lemma~\ref{lem:harmonic-extension}\eqref{it:harmonic-sum},
the plate function $\tilde G_N(\,\cdot\,;\lambda)$ is a density of the $N$-linear form
\[
(\alpha_1,\dots,\alpha_N)\longmapsto S_N(\alpha_1\otimes\lambda_1,\dots,\alpha_N\otimes\lambda_N)
\]
on the tuples whose supports lie in small discs about the points of a configuration in $\mathcal{O}_N$.
A configuration of $N$ distinct points that does not belong to $\mathcal{O}_N$
has two points with the same time and two points with the same abscissa (Lemma~\ref{lem:one-function}(e)).
At such configurations the form has no singular part (Proposition~\ref{prop:no-singular-part}).

Moments of two clusters of fields that are separated in time depend analytically on a displacement of one cluster along $e_2$.

\begin{lemma}[Cluster moves]\label{lem:cluster-moves}
Let $\mu$ be a law on $\mathcal{S}'$ that satisfies \textup{(G)}, $(T_{123})$, $(R_1)$, $(RP_1)$, $(Q_{12})$ and $(C_{12})$.
Let $c\in\R$ and $\delta>0$.
Let $\Phi'$ be a polynomial in fields $\mathbb{B}(b)$ whose arguments $b\in L^2(\R^3;\C^3)$ vanish almost everywhere outside the half space $\{x_1<c-\delta\}$,
and let $\Phi''$ be a polynomial in fields $\mathbb{B}(b)$ whose arguments $b\in L^2(\R^3;\C^3)$ vanish almost everywhere outside $\{x_1>c+\delta\}$.
Then the function
\[
\R\ni s\longmapsto\mu\left(\Phi'\cdot U(se_2)\Phi''\right)
\]
is the restriction to $\R$ of a function that is holomorphic on the strip $\{s\in\C:|\Im s|<2\delta\}$.
\end{lemma}

\begin{proof}
By $(T_{123})$ the law $\mu$ is invariant under the translation by $-ce_1$;
the operator $U(-ce_1)$ is multiplicative, commutes with $U(se_2)$, and maps $\mathbb{B}(b)$ to $\mathbb{B}(\tau_{-ce_1}\cdot b)$
(Lemma~\ref{lem:domination-consequences}(a)).
Hence
\[
\mu\left(\Phi'\cdot U(se_2)\Phi''\right)=\mu\left(U(-ce_1)\Phi'\cdot U(se_2)U(-ce_1)\Phi''\right),
\]
and the polynomials $U(-ce_1)\Phi'$, $U(-ce_1)\Phi''$ satisfy the hypotheses with $0$ in place of $c$.
So it suffices to treat the case $c=0$, which is assumed from now on.

Since $\Theta_1$ is antilinear and multiplicative, \eqref{eq:rotation-conjugation} shows that $\Theta_1\Phi'$ is a polynomial in the fields $\mathbb{B}(\overline{\theta_1\cdot b})$,
where $b$ ranges over the arguments of the fields in $\Phi'$;
these new arguments vanish outside $\{x_1>\delta\}$, and $\Theta_1\Theta_1\Phi'=\Phi'$.
Consequently $U_{-\delta}\Theta_1\Phi'$ and $U_{-\delta}\Phi''$, where $U_{-\delta}=U(-\delta e_1)$, are polynomials in fields whose arguments vanish on $\{x_1<0\}$;
they belong to $L^{\infty-}_{+}$ (Definition~\ref{def:time-slabs}),
and so does $U(se_2)\Phi''$ for every real $s$, its arguments vanishing outside $\{x_1>\delta\}$.
Put
\[
\psi':=[U_{-\delta}\Theta_1\Phi']_1,\qquad\psi'':=[U_{-\delta}\Phi'']_1 .
\]
The reconstruction in the direction $e_1$ \cite[Proposition~8.19\,(i)--(iii)]{NoncompactQED}
(where the semigroup $\exp(-t\mathsf{H}_1)$ is written $T(t)$ and the scalar product of classes is expressed through the linear map $\hat\theta_1$)
has the properties
\[
\left\langle[\Phi]_1,[\Phi_1]_1\right\rangle=\mu\left((\Theta_1\Phi)\cdot\Phi_1\right),\qquad
\exp(-t\mathsf{H}_1)\,[\Phi]_1=[U(te_1)\Phi]_1,\qquad
V_1(y)\,[\Phi]_1=[U(y)\Phi]_1
\]
for $\Phi,\Phi_1\in L^2(\Sigma_1^+,\mu)$, $t\geq0$ and $y\in e_1^\perp$;
the operators $\exp(-t\mathsf{H}_1)$ are self-adjoint and commute with the unitary operators $V_1(y)$.
Hence, for real $s$,
\[
[\Theta_1\Phi']_1=\exp(-\delta\mathsf{H}_1)\,\psi',\qquad[U(se_2)\Phi'']_1=V_1(se_2)\exp(-\delta\mathsf{H}_1)\,\psi'' .
\]
Moreover
\[
\exp(-2\delta\mathsf{H}_1)\,V_1(se_2)=\int\mathrm{d}\mathfrak{E}_1(\omega,p_\perp)\,\exp(-2\delta\omega+isp_2)
\]
by Lemma~\ref{lem:joint-spectrum}.
With $\xi_s:=(2\delta,s)\in\C^2$ and the definition of $\mathsf{T}$ in Lemma~\ref{lem:cone-semigroup} it follows that, for real $s$,
\begin{align*}
\mu\left(\Phi'\cdot U(se_2)\Phi''\right)
&=\left\langle[\Theta_1\Phi']_1,[U(se_2)\Phi'']_1\right\rangle
=\left\langle\psi',\exp(-\delta\mathsf{H}_1)\,V_1(se_2)\exp(-\delta\mathsf{H}_1)\,\psi''\right\rangle\\
&=\left\langle\psi',\mathsf{T}(\xi_s)\,\psi''\right\rangle .
\end{align*}
For complex $s$ one has $d_{\mathcal{W}}(\xi_s)=2\delta-|\Im s|$;
so $\xi_s$ lies in $\mathcal{W}$ when $|\Im s|<2\delta$.
On this strip the map $s\mapsto\mathsf{T}(\xi_s)$, the composition of $\mathsf{T}$ with an affine map of $\C$ into $\C^2$,
is holomorphic in the operator norm (Lemma~\ref{lem:cone-semigroup}\eqref{it:semigroup-holomorphic}).
Therefore $s\mapsto\left\langle\psi',\mathsf{T}(\xi_s)\,\psi''\right\rangle$ is holomorphic on the strip,
and on the real axis it coincides with $\mu\left(\Phi'\cdot U(se_2)\Phi''\right)$.
\end{proof}

\begin{proposition}[No singular part]\label{prop:no-singular-part}
Let $\mu$ be a law on $\mathcal{S}'$ that satisfies \textup{(G)}, $(T_{123})$, $(R_1)$, $(RP_1)$, $(Q_{12})$ and $(C_{12})$, and let $N\geq2$.
Let $\lambda=(\lambda_1,\dots,\lambda_N)$ be labels (Definition~\ref{def:labels-plates}),
and let $\alpha_1,\dots,\alpha_N\in L^2(\R^2;\C)$ have compact, pairwise disjoint supports.
Then
\[
S_N(\alpha_1\otimes\lambda_1,\dots,\alpha_N\otimes\lambda_N)=\int_{(\R^2)^N}\mathrm{d}Y\,\prod_{a=1}^{N}\alpha_a(y_a)\cdot\tilde G_N(Y;\lambda),
\]
where $Y=(y_1,\dots,y_N)$ and $\tilde G_N(\,\cdot\,;\lambda)$ is the plate function of Lemma~\ref{lem:harmonic-extension}\eqref{it:harmonic-sum},
a real analytic function on $(\R^2)^N_{\neq}$.
\end{proposition}

\begin{proof}
Put $K:=\prod_{a=1}^{N}\supp\alpha_a$.
It is a compact subset of $(\R^2)^N_{\neq}$, since the supports are compact and pairwise disjoint.
The integrand vanishes almost everywhere outside $K$;
on $K$ the function $\tilde G_N(\,\cdot\,;\lambda)$ is continuous, hence bounded;
and every $\alpha_a$ is integrable, its support being compact.
So the integral converges absolutely.

A tuple $(\alpha_1,\dots,\alpha_N)$ of functions in $L^2(\R^2;\C)$ is called \emph{admissible} if the supports of its entries are compact and pairwise disjoint.
For admissible tuples let $D(\alpha_1,\dots,\alpha_N)$ be the left side of the asserted identity minus its right side;
$D$ is linear in each entry.
For $y\in\R^2$ and $r>0$ let $B(y,r)$ be the open disc of radius $r$ about $y$.
Let $X=(y_1,\dots,y_N)\in(\R^2)^N_{\neq}$.
One says that $D$ \emph{vanishes near $X$ with the radius $r>0$} if $D(\alpha_1,\dots,\alpha_N)=0$ for every admissible tuple with $\supp\alpha_a\subset B(y_a,r)$ for $a=1,\dots,N$;
then $D$ vanishes near $X$ with every smaller radius as well.

\emph{Step 1: configurations in $\mathcal{O}_N$.}
Let $X\in\mathcal{O}_N$.
As $\mathcal{O}_N=\mathcal{Z}_N\cap(\R^2)^N$ and $\mathcal{Z}_N$ is the union of the tubes (Definition~\ref{def:configuration-tubes}),
$X$ is a real point of a tube $\mathcal{Z}(\mathfrak{s},\upsilon)$.
The tube is open (Lemma~\ref{lem:one-function}(a)); choose $r>0$ with $\prod_{a}B(y_a,r)\subset\mathcal{Z}(\mathfrak{s},\upsilon)$.
For an admissible tuple with $\supp\alpha_a\subset B(y_a,r)$ for all $a$, the product of the supports lies in the tube, and Lemma~\ref{lem:one-function}(b), (c) gives
\[
S_N(\alpha_1\otimes\lambda_1,\dots,\alpha_N\otimes\lambda_N)=\int_{(\R^2)^N}\mathrm{d}Y\,\prod_{a=1}^{N}\alpha_a(y_a)\cdot\mathcal{G}_N(Y;\lambda).
\]
The real points of the tube belong to $\mathcal{O}_N$, where $\mathcal{G}_N(\,\cdot\,;\lambda)=\tilde G_N(\,\cdot\,;\lambda)$ by Lemma~\ref{lem:harmonic-extension}\eqref{it:harmonic-sum}.
Hence $D(\alpha_1,\dots,\alpha_N)=0$, and $D$ vanishes near $X$ with the radius $r$.

\emph{Step 2: all configurations of distinct points.}
For $X=(y_1,\dots,y_N)\in(\R^2)^N_{\neq}$, with $y_a=(y_a^0,y_a^1)$, let
\[
k(X):=\#\left|\left\{\{a,a'\}:\ a\neq a',\ y_a^1=y_{a'}^1\right\}\right|
\]
be the number of pairs of points of $X$ with equal abscissae.
It is shown by induction on $k(X)$ that $D$ vanishes near every $X\in(\R^2)^N_{\neq}$ with some radius.
If $k(X)=0$, the abscissae of the points of $X$ are pairwise distinct, so $X\in\mathcal{O}_N$ by Lemma~\ref{lem:one-function}(e), and Step~1 applies.
Let $k(X)>0$, and assume that $D$ vanishes near every $X'\in(\R^2)^N_{\neq}$ with $k(X')<k(X)$.

Choose indices $l\neq l'$ with $y_l^1=y_{l'}^1$.
As $y_l\neq y_{l'}$, the two times differ; after interchanging $l$ and $l'$ if necessary, $y_l^0<y_{l'}^0$.
Choose a level $c\in(y_l^0,y_{l'}^0)$ that differs from all the times $y_a^0$, and put
\[
J:=\{a:\ y_a^0>c\},\qquad I:=\{a:\ y_a^0<c\},\qquad\delta:=\min_{a\in J}y_a^0-\max_{a\in I}y_a^0 .
\]
The index set $\{1,\dots,N\}$ is the disjoint union of $I$ and $J$, with $l\in I$ and $l'\in J$, and $\delta>0$.
For real $s$ let $X(s):=(y_1(s),\dots,y_N(s))$, where
\[
y_a(s):=y_a\quad(a\in I),\qquad y_a(s):=y_a+s\,n_{\pi/2}\quad(a\in J);
\]
so the points with indices in $J$ are moved by $s$ in the direction $n_{\pi/2}=(0,1)$ of the abscissa.
Inside $I$, and inside $J$, the points of $X(s)$ are the points of $X$ translated by one and the same vector,
and a point with index in $I$ and a point with index in $J$ have different times;
hence $X(s)\in(\R^2)^N_{\neq}$ for every real $s$.

Fix $s_0$ with
\[
0<s_0<\min\left\{|y_a^1-y_{a'}^1|:\ y_a^1\neq y_{a'}^1\right\},
\]
any $s_0>0$ being allowed if all the abscissae of the points of $X$ are equal.
For $a,a'$ both in $I$ or both in $J$ one has $y_a^1(s_0)-y_{a'}^1(s_0)=y_a^1-y_{a'}^1$.
For $a\in I$ and $a'\in J$ one has $y_a^1(s_0)-y_{a'}^1(s_0)=y_a^1-y_{a'}^1-s_0$;
this equals $-s_0\neq0$ if $y_a^1=y_{a'}^1$, and it does not vanish if $y_a^1\neq y_{a'}^1$ either, since then $|y_a^1-y_{a'}^1|>s_0$.
Thus the coincidences of abscissae inside $I$ and inside $J$ are the same in $X(s_0)$ as in $X$,
while $X(s_0)$ has no coincidence between $I$ and $J$, whereas $X$ has the coincidence of $l$ and $l'$.
Therefore $k(X(s_0))<k(X)$.

By the induction hypothesis, applied to $X(s_0)$, there is $r>0$ such that $D(\alpha_1,\dots,\alpha_N)=0$ for every admissible tuple with $\supp\alpha_a\subset B(y_a(s_0),2r)$ for all $a$.
Decreasing $r$, one may assume moreover
\[
4r<\delta,\qquad 4r<\min_{a\neq a'}|y_a-y_{a'}| .
\]
Let $(\alpha_1,\dots,\alpha_N)$ be an admissible tuple with $\supp\alpha_a\subset B(y_a,r)$ for all $a$.
For $v\in\R^2$ and a function $\alpha$ on the plane let $(\tau_v\alpha)(y):=\alpha(y-v)$.
For real $s$ put
\[
\alpha_a^{(s)}:=\alpha_a\quad(a\in I),\qquad\alpha_a^{(s)}:=\tau_{s\,n_{\pi/2}}\alpha_a\quad(a\in J),\qquad F(s):=D\left(\alpha_1^{(s)},\dots,\alpha_N^{(s)}\right).
\]
The tuple $(\alpha_1^{(s)},\dots,\alpha_N^{(s)})$ is admissible for every real $s$, so that $F$ is defined on $\R$.
Indeed, $\supp\alpha_a^{(s)}\subset B(y_a(s),r)$.
Inside $I$, and inside $J$, these discs are translates by one vector of the discs $B(y_a,r)$, which are pairwise disjoint because $2r<\min_{a\neq a'}|y_a-y_{a'}|$.
For $a\in I$ the points of $B(y_a(s),r)$ have times less than $\max_{b\in I}y_b^0+r$,
for $a'\in J$ the points of $B(y_{a'}(s),r)$ have times greater than $\min_{b\in J}y_b^0-r$,
and the second bound exceeds the first by $\delta-2r>0$.

The first term of $F(s)$ is a moment of two clusters of fields.
Let
\[
\Phi':=\prod_{a\in I}\mathbb{B}_{\lambda_a}(\alpha_a),\qquad\Phi'':=\prod_{a\in J}\mathbb{B}_{\lambda_a}(\alpha_a).
\]
For $a\in J$ the plate $(\tau_{s\,n_{\pi/2}}\alpha_a)\otimes\lambda_a$ is $\tau_{se_2}\cdot(\alpha_a\otimes\lambda_a)$;
so $\mathbb{B}_{\lambda_a}(\alpha_a^{(s)})=U(se_2)\,\mathbb{B}_{\lambda_a}(\alpha_a)$ by Lemma~\ref{lem:domination-consequences}(a) and $(T_{123})$.
As $U(se_2)$ is multiplicative,
\[
S_N\left(\alpha_1^{(s)}\otimes\lambda_1,\dots,\alpha_N^{(s)}\otimes\lambda_N\right)=\mu\left(\Phi'\cdot U(se_2)\Phi''\right).
\]
The arguments $\alpha_a\otimes\lambda_a$, $a\in I$, of the fields in $\Phi'$ vanish outside $B(y_a,r)\times\R\subset\{x_1<\max_{b\in I}y_b^0+r\}$,
and those of the fields in $\Phi''$ vanish outside $\{x_1>\min_{b\in J}y_b^0-r\}$.
Lemma~\ref{lem:cluster-moves}, applied with the level $\tfrac12\left(\max_{b\in I}y_b^0+\min_{b\in J}y_b^0\right)$ and the margin $\delta/2-r>0$
in the roles of the numbers $c$ and $\delta$ of that lemma,
shows that the first term of $F(s)$ is the restriction to $\R$ of a holomorphic function on the strip $\{|\Im s|<\delta-2r\}$.
In particular it is real analytic on $\R$.

The second term of $F(s)$ is an integral against $\tilde G_N$.
Let $v_J\in(\R^2)^N$ be the configuration whose $a$-th point is $n_{\pi/2}$ for $a\in J$ and $0$ for $a\in I$.
The substitution $y_a\mapsto y_a+s\,n_{\pi/2}$, $a\in J$, gives
\[
\int_{(\R^2)^N}\mathrm{d}Y\,\prod_{a=1}^{N}\alpha_a^{(s)}(y_a)\cdot\tilde G_N(Y;\lambda)=\int_{(\R^2)^N}\mathrm{d}Y\,\prod_{a=1}^{N}\alpha_a(y_a)\cdot\tilde G_N(Y+s\,v_J;\lambda).
\]
Let $s_1<s_2$.
The set $K':=\{Y+s\,v_J:\ Y\in\prod_a\supp\alpha_a,\ s\in[s_1,s_2]\}$ is compact,
and it lies in $(\R^2)^N_{\neq}$, because the $a$-th point of $Y+s\,v_J$ lies in $\supp\alpha_a^{(s)}$ and these supports are pairwise disjoint.
The function $\tilde G_N(\,\cdot\,;\lambda)$ is real analytic on $(\R^2)^N_{\neq}$
(Lemma~\ref{lem:harmonic-extension}\eqref{it:harmonic-analytic}, \eqref{it:harmonic-sum}).
Hence every point $Z_0$ of $K'$ is the centre of a polydisc in $(\C^2)^N$
to which $\tilde G_N(\,\cdot\,;\lambda)$ extends holomorphically, with a bounded extension;
a polydisc about $Z_0$ is the set of the $Z\in(\C^2)^N$ each of whose $2N$ complex coordinates
differs from the corresponding coordinate of $Z_0$ by less than a fixed positive number, the radius of the polydisc.
Finitely many of the polydiscs with the same centres and half the radii cover $K'$;
a polydisc of the least of these half radii about a point of $K'$ lies in one of the polydiscs of full radius.
So there are $\varepsilon>0$ and $C<\infty$ such that, for every $Z_0\in K'$,
the function $\tilde G_N(\,\cdot\,;\lambda)$ extends holomorphically to the polydisc of radius $\varepsilon$ about $Z_0$,
with modulus at most $C$ there.
Let $Y\in\prod_a\supp\alpha_a$ and $s'\in[s_1,s_2]$.
The coordinates of $Y+s\,v_J$ and of $Y+s'\,v_J$ differ by at most $|s-s'|$;
so, near $s'$, the function $s\mapsto\tilde G_N(Y+s\,v_J;\lambda)$ is the restriction of a holomorphic function
on the disc $\{|s-s'|<\varepsilon\}$ of modulus at most $C$,
and Cauchy's estimates bound its derivative of order $j$ at $s'$ by $C\cdot j!\cdot\varepsilon^{-j}$.
Differentiation under the integral sign shows that the second term of $F(s)$ is smooth on $(s_1,s_2)$,
with derivatives of order $j$ bounded by $C\cdot j!\cdot\varepsilon^{-j}\cdot\prod_{a}\|\alpha_a\|_{L^1}$;
by Taylor's formula it is real analytic on $(s_1,s_2)$.
As $s_1<s_2$ are arbitrary, $F$ is real analytic on $\R$.

For $|s-s_0|<r$ and every $a$ one has $\supp\alpha_a^{(s)}\subset B(y_a(s),r)\subset B(y_a(s_0),2r)$;
hence $F(s)=0$ by the choice of $r$.
By Lemma~\ref{lem:analysis-facts}\eqref{it:analytic-zeros}, $F$ vanishes on all of $\R$.
In particular $D(\alpha_1,\dots,\alpha_N)=F(0)=0$: the form $D$ vanishes near $X$ with the radius $r$.
This completes the induction.

\emph{Step 3: arbitrary supports.}
Let $(\alpha_1,\dots,\alpha_N)$ be admissible, and $K=\prod_a\supp\alpha_a$ as above.
By Step~2 every $X=(y_1,\dots,y_N)\in K$ has a radius $r_X>0$ with which $D$ vanishes near $X$.
The open sets $\prod_aB(y_a,r_X/2)$, $X\in K$, cover the compact set $K$.
Choose finitely many of them that cover $K$, belonging to configurations $X^{(n)}=(y_1^{(n)},\dots,y_N^{(n)})$ with radii $r_n:=r_{X^{(n)}}$.
Let $d>0$ with $2\sqrt{2}\,d<\min_nr_n$, and let $\mathcal{Q}_d$ be the set of the squares $d(k+[0,1)^2)$, $k\in\Z^2$, each of diameter $\sqrt{2}\,d$.
Every $\alpha_a$ is the sum of the functions $1_Q\alpha_a$, $Q\in\mathcal{Q}_d$, of which only finitely many are not zero.
By linearity in each entry,
\[
D(\alpha_1,\dots,\alpha_N)=\sum_{Q_1,\dots,Q_N\in\mathcal{Q}_d}D\left(1_{Q_1}\alpha_1,\dots,1_{Q_N}\alpha_N\right),
\]
a finite sum of values on admissible tuples, since $\supp(1_Q\alpha_a)\subset\overline{Q}\cap\supp\alpha_a$.
Consider one term.
If $1_{Q_a}\alpha_a=0$ for some $a$, the term vanishes.
Otherwise every $Q_a$ contains a point $y'_a$ of $\supp\alpha_a$, because $\alpha_a$ vanishes almost everywhere outside its support.
The configuration $(y'_1,\dots,y'_N)$ lies in $K$, hence in $\prod_aB(y_a^{(n)},r_n/2)$ for some $n$.
Every point of $\overline{Q_a}$ has a distance at most $\sqrt{2}\,d<r_n/2$ from $y'_a$, hence a distance less than $r_n$ from $y_a^{(n)}$.
So $\supp(1_{Q_a}\alpha_a)\subset B(y_a^{(n)},r_n)$ for every $a$,
and the term vanishes because $D$ vanishes near $X^{(n)}$ with the radius $r_n$.
Summing over all terms, $D(\alpha_1,\dots,\alpha_N)=0$.
\end{proof}

\subsection{Invariance under rotations}\label{sec:rotation-theorems}

\begin{theorem}[Rotations of one plane]\label{thm:one-plane}
Let $j\neq i$ be indices in $\{1,2,3\}$,
and let $\mu$ be a law on $(\mathcal{S}',\Sigma)$, $\mathcal{S}'=\mathcal{S}'(\R^3;\Lambda^2)$,
with the following seven properties.
\begin{itemize}
\item[\textup{(G)}] $\int\mathrm{d}\mu(F)\,\exp F(f)\le\exp\left(\frac{1}{2}\|f\|_2^2\right)$ for every $f\in\mathcal{S}(\R^3;\Lambda^2)$.
\item[\textup{(dF)}] $\mu(dF=0)=1$.
\item[$(T_{123})$] The law $\mu$ is invariant under the translations of $\R^3$.
\item[$(R_j)$] The law $\mu$ is invariant under the reflection $\theta_j$ in the plane $\{x_j=0\}$.
\item[$(RP_j)$] The law $\mu$ is reflection positive at the plane $\{x_j=0\}$.
\item[$(Q_{ji})$] The law $\mu$ is invariant under the quarter turn $R^{ji}_{\pi/2}$ of the plane $(j,i)$.
\item[$(C_{ji})$] The joint spectral measure $\mathfrak{E}_j$ of the reconstruction of $\mu$ in the direction $j$
(Definition~\ref{def:laws}, Lemma~\ref{lem:joint-spectrum}) satisfies
$\mathfrak{E}_j\left(\left\{(\omega,p_\perp):\omega<|p_i|\right\}\right)=0$;
that is, $\mathsf{H}_j\ge|\mathsf{P}_i|$, with the constant one, on $\mathcal{H}_j$.
\end{itemize}
Then, for every integer $N\ge1$, every real $t$ and all $b_1,\dots,b_N\in L^2(\R^3;\C^3)$,
whose supports are subject to no condition,
\[
S_N\left(R^{ji}_t\cdot b_1,\dots,R^{ji}_t\cdot b_N\right)=S_N(b_1,\dots,b_N),
\]
and $\left(R^{ji}_t\right)_*\mu=\mu$ for every real $t$.
Thus $\mu$ is invariant under all rotations of the plane $(j,i)$.
\end{theorem}

\begin{proof}
\emph{Reduction to the plane $(1,2)$.}
Let $g_{ji}$ be the permutation matrix with $g_{ji}e_1=e_j$ and $g_{ji}e_2=e_i$,
and put $\mu':=\left(g_{ji}^{-1}\right)_*\mu$.
By Lemma~\ref{lem:transport}(d), the law $\mu'$ has the seven properties of the theorem for the pair $(1,2)$,
and $\left(R^{12}_t\right)_*\mu'=\mu'$ holds if and only if $\left(R^{ji}_t\right)_*\mu=\mu$.
Hence the invariance of $\mu$ under all $R^{ji}_t$ follows from the invariance of $\mu'$ under all $R^{12}_t$.
The identity of the moments, in turn, is a consequence of the invariance of the law.
Indeed, let $\mathfrak{g}$ be a Euclidean map with $\mathfrak{g}_*\mu=\mu$.
Then $\mu\left(\hat{\mathfrak{g}}\Phi\right)=\mu(\Phi)$ for every $\Phi\in L^1(\mu)$,
and $\hat{\mathfrak{g}}$ is multiplicative \cite[Section~8.1.2]{NoncompactQED}
(where a Euclidean map is written $g$);
moreover $\hat{\mathfrak{g}}\,\mathbb{B}(b)=\mathbb{B}(\mathfrak{g}\cdot b)$ in $L^2(\mu)$ for every $b\in L^2(\R^3;\C^3)$,
by Lemma~\ref{lem:domination-consequences}(a).
Each $\mathbb{B}(b_a)$ lies in $L^N(\mu)$ by the same lemma,
so that the product $\mathbb{B}(b_1)\cdots\mathbb{B}(b_N)$ is integrable by H\"older's inequality, and
\begin{equation}\label{eq:rotation-invariant-moments}
S_N\left(\mathfrak{g}\cdot b_1,\dots,\mathfrak{g}\cdot b_N\right)
=\mu\left(\hat{\mathfrak{g}}\left(\mathbb{B}(b_1)\cdots\mathbb{B}(b_N)\right)\right)
=S_N(b_1,\dots,b_N)
\end{equation}
for every $N\ge1$ and all $b_1,\dots,b_N\in L^2(\R^3;\C^3)$.
For $\mathfrak{g}=R^{ji}_t$ this is the identity of the moments asserted in the theorem.
It therefore suffices to prove $\left(R^{12}_t\right)_*\mu=\mu$ for every real $t$ and every law $\mu$ with the seven properties for the pair $(1,2)$.

From now on $(j,i)=(1,2)$, and $R_t:=R^{12}_t$ as in Definition~\ref{def:labels-plates}.
Fix a real number $t$ and an integer $N\ge2$;
for $N=1$ both sides of the identity vanish, since $S_1=0$.
For $f_1,\dots,f_N\in L^2(\R^3;\C^3)$ put
\[
D(f_1,\dots,f_N):=S_N\left(R_t\cdot f_1,\dots,R_t\cdot f_N\right)-S_N(f_1,\dots,f_N).
\]
The map $b\mapsto R_t\cdot b$ is an isometry of $L^2(\R^3;\C^3)$.
By Lemma~\ref{lem:domination-consequences}(b), with the constant $c_N$ of that lemma,
$D$ is therefore an $N$-linear form on $L^2(\R^3;\C^3)$ with
\[
\left|D(f_1,\dots,f_N)\right|\le2\,c_N^N\cdot\prod_{a=1}^N\|f_a\|.
\]
The open set $U=\R^2$ of the plane gives $L^2_U=L^2(\R^3;\C^3)$,
and the two hypotheses of Proposition~\ref{prop:no-contact-terms} are now checked for $D$ and this $U$.

\emph{Gradients.}
Let one of the entries be $f_a=\nabla u$ with $u\in C^{\infty}_{\mathrm c}(\R^3;\C)$.
By \textup{(dF)}, $\mathbb{B}(\nabla u)=0$ holds $\mu$-almost surely, see~\eqref{eq:rotation-gradient-fields}.
Since $\det R_t=1$, the chain rule~\eqref{eq:rotation-gradient-action} gives $R_t\cdot\nabla u=\nabla(u\circ R_{-t})$,
and $u\circ R_{-t}$ again belongs to $C^{\infty}_{\mathrm c}(\R^3;\C)$;
so $\mathbb{B}(R_t\cdot f_a)=0$ holds $\mu$-almost surely as well.
Each of the two moments in $D(f_1,\dots,f_N)$ is the integral of a product with a factor that vanishes $\mu$-almost surely,
and $D(f_1,\dots,f_N)=0$.

\emph{Entries over disjoint compact sets.}
Take compact sets $K_1,\dots,K_N$ in the plane $\R^2$, no two of which meet.
Consider first plates $f_a=\alpha_a\otimes\lambda_a$, $a=1,\dots,N$,
with $\alpha_a\in L^2(\R^2;\C)$ vanishing outside $K_a$ and arbitrary labels $\lambda_a$;
put $\lambda:=(\lambda_1,\dots,\lambda_N)$ and $w:=w(\lambda)$.
The function $\alpha_a\circ R_{-t}$ vanishes outside the compact set $R_tK_a$, and these sets are pairwise disjoint.
Propositions~\ref{prop:no-singular-part} and~\ref{prop:anomaly-vanishes} apply to $\mu$ and to the integer $N\ge2$,
since $\mu$ satisfies \textup{(G)}, \textup{(dF)}, $(T_{123})$, $(R_1)$, $(RP_1)$, $(Q_{12})$ and $(C_{12})$.
Identity~\eqref{it:plate-rotation} of Lemma~\ref{lem:plate-identities} and the $N$-linearity of $S_N$,
Proposition~\ref{prop:no-singular-part} for the supports $R_tK_a$,
the substitution $Y\mapsto R_tY$, which preserves the Lebesgue measure of $(\R^2)^N$,
Proposition~\ref{prop:anomaly-vanishes}, applicable because $\prod_aK_a$ lies in $(\R^2)^N_{\neq}$,
and Proposition~\ref{prop:no-singular-part} for the supports $K_a$ give, in this order,
\begin{align*}
S_N\left(R_t\cdot(\alpha_1\otimes\lambda_1),\dots,R_t\cdot(\alpha_N\otimes\lambda_N)\right)
&=\exp(-iwt)\int\mathrm{d}Y\,\prod_{a=1}^N\alpha_a\left(R_{-t}y_a\right)\cdot\tilde G_N(Y;\lambda)\\
&=\exp(-iwt)\int\mathrm{d}Y\,\prod_{a=1}^N\alpha_a(y_a)\cdot\tilde G_N(R_tY;\lambda)\\
&=\int\mathrm{d}Y\,\prod_{a=1}^N\alpha_a(y_a)\cdot\tilde G_N(Y;\lambda)\\
&=S_N\left(\alpha_1\otimes\lambda_1,\dots,\alpha_N\otimes\lambda_N\right).
\end{align*}
Hence $D(f_1,\dots,f_N)=0$ for every choice of plates of this kind.
The $N$-linearity of $D$ extends this to the case where each $f_a$ is a finite sum of plates $\alpha\otimes\lambda$ with $\alpha=0$ outside $K_a$.
Such sums are dense in the space of the $b\in L^2(\R^3;\C^3)$ that vanish almost everywhere outside $K_a\times\R$,
by Lemma~\ref{lem:plates-total}\eqref{it:plates-total-square};
as $D$ is continuous, $D(f_1,\dots,f_N)=0$ whenever each $f_a$ vanishes almost everywhere outside $K_a\times\R$.

\emph{Conclusion for the plane $(1,2)$.}
Proposition~\ref{prop:no-contact-terms} now gives $D=0$:
\[
S_N\left(R_t\cdot b_1,\dots,R_t\cdot b_N\right)=S_N(b_1,\dots,b_N)
\qquad\left(b_1,\dots,b_N\in L^2(\R^3;\C^3)\right).
\]
This holds for every $N\ge1$ and every real $t$.
The real Schwartz fields belong to $L^2(\R^3;\C^3)$,
and Lemma~\ref{lem:domination-consequences}(c) yields $(R_t)_*\mu=\mu$.
\end{proof}

\begin{theorem}[Invariance under the orthogonal group]\label{thm:orthogonal-group}
Let $\mu$ be a law on $(\mathcal{S}',\Sigma)$ that satisfies \textup{(dF)} and \textup{(H)};
the latter means \textup{(G)}, invariance under $\R^3\rtimes\Ocub$, and $(RP_j)$ for $j=1,2,3$.
Suppose that $(C_{ji})$ holds for one ordered pair of indices $j\neq i$.
Then the following hold.
\begin{itemize}
\item The law $\mu$ is invariant under $\R^3\rtimes\mathrm{O}(3)$.
\item $S_N\left(g\cdot b_1,\dots,g\cdot b_N\right)=S_N(b_1,\dots,b_N)$ for every $N\ge1$, every $g\in\mathrm{O}(3)$ and all $b_1,\dots,b_N\in L^2(\R^3;\C^3)$.
\item For every unit vector $n\in\R^3$ and every $c\in\R$,
the law $\mu$ is invariant under the reflection $r_{n,c}$ in the plane $\{n\cdot x=c\}$,
it is reflection positive at this plane,
and $\left\langle\hat r_{n,c}\Phi,\Phi\right\rangle_\mu\ge0$ for every $\Phi\in L^2\left(\Sigma_{\{n\cdot x>c\}},\mu\right)$,
the space of the observables of the open side $\{n\cdot x>c\}$.
\item $(C_{j'i'})$ holds for every ordered pair of indices $j'\neq i'$.
\item The two-point function $S_2$ has the K\"all\'en--Lehmann form of Theorem~\ref{thm:two-point}(c),
with a positive Borel measure on $[0,\infty)$ of mass at most one.
\end{itemize}
More generally, let $\mu$ be a law on $(\mathcal{S}',\Sigma)$ with \textup{(G)}, \textup{(dF)} and $(T_{123})$,
and let $(j,i)$ and $(j',i')$ be ordered pairs of distinct indices with $\{j,i\}\neq\{j',i'\}$
such that $\mu$ satisfies $(R_j)$, $(RP_j)$, $(Q_{ji})$, $(C_{ji})$ and $(R_{j'})$, $(RP_{j'})$, $(Q_{j'i'})$, $(C_{j'i'})$.
Then $\mu$ is invariant under $\mathrm{SO}(3)$.
\end{theorem}

\begin{proof}
Under \textup{(H)} the law $\mu$ satisfies $(T_{123})$,
and, since $\theta_1,\theta_2,\theta_3$ belong to $\Ocub$, also $(R_l)$ and $(RP_l)$ for $l=1,2,3$;
so its reconstruction is defined in each of the three directions.

\emph{The cone for every ordered pair.}
Let $(j',i')$ be an ordered pair of distinct indices, and let $g\in\Ocub$ satisfy $ge_j=e_{j'}$ and $ge_i=e_{i'}$;
the permutation matrix that maps $e_j$, $e_i$ and the remaining unit vector to $e_{j'}$, $e_{i'}$ and the remaining unit vector is such an element.
By \textup{(H)}, $g_*\mu=\mu$.
As $g$ is orthogonal, $g\theta_jg^{-1}=\theta_{j'}$ and $g\{x_j>0\}=\{x_{j'}>0\}$.
Lemma~\ref{lem:transport}(a), applied to $g$ and to the law $\left(g^{-1}\right)_*\mu=\mu$, shows the following:
$\mu\left(\hat g\Phi\right)=\mu(\Phi)$ for $\Phi\in L^1(\mu)$;
$\hat g$ is a unitary operator of $L^2(\mu)$ that maps $L^2\left(\Sigma_j^+,\mu\right)$ onto $L^2\left(\Sigma_{j'}^+,\mu\right)$;
$\hat gU(a)=U(ga)\hat g$ for $a\in\R^3$;
and $\hat g\hat\theta_j=\hat\theta_{j'}\hat g$.
The operator $\hat g$ is multiplicative and commutes with complex conjugation \cite[Section~8.1.2]{NoncompactQED};
hence $\hat g\Theta_j=\Theta_{j'}\hat g$, and for $\Phi,\Phi'\in L^2\left(\Sigma_j^+,\mu\right)$
\[
\left\langle\left[\hat g\Phi\right]_{j'},\left[\hat g\Phi'\right]_{j'}\right\rangle
=\mu\left(\left(\Theta_{j'}\hat g\Phi\right)\cdot\left(\hat g\Phi'\right)\right)
=\mu\left(\hat g\left(\left(\Theta_j\Phi\right)\cdot\Phi'\right)\right)
=\left\langle[\Phi]_j,[\Phi']_j\right\rangle .
\]
The classes of the elements of $L^2\left(\Sigma_j^+,\mu\right)$ are dense in $\mathcal{H}_j$,
and those of $L^2\left(\Sigma_{j'}^+,\mu\right)$ in $\mathcal{H}_{j'}$ (Definition~\ref{def:laws}).
So $[\Phi]_j\mapsto\left[\hat g\Phi\right]_{j'}$ is well defined and isometric on these classes,
and it extends to a unitary operator $J_g$ of $\mathcal{H}_j$ onto $\mathcal{H}_{j'}$.
For $\Phi\in L^2\left(\Sigma_j^+,\mu\right)$, $t\ge0$ and $y\in e_j^\perp$ the reconstruction satisfies
$\exp(-t\mathsf{H}_j)V_j(y)[\Phi]_j=[U(te_j+y)\Phi]_j$;
moreover $gy\in e_{j'}^\perp$ and $\hat gU(te_j+y)=U(te_{j'}+gy)\hat g$.
Therefore
\[
J_g\exp(-t\mathsf{H}_j)V_j(y)J_g^{-1}=\exp(-t\mathsf{H}_{j'})V_{j'}(gy)\qquad\left(t\ge0,\ y\in e_j^\perp\right);
\]
for $t=0$ this reads $J_gV_j(y)J_g^{-1}=V_{j'}(gy)$, and for $y=0$, $t=1$ it reads $J_g\exp(-\mathsf{H}_j)J_g^{-1}=\exp(-\mathsf{H}_{j'})$.
For $\mathsf{H}=\mathsf{H}_j$ and for $\mathsf{H}=\mathsf{H}_{j'}$, and for real $t$,
the operator $\exp(it\mathsf{H})$ is the bounded Borel function $x\mapsto\exp(-it\log x)$, with the value $0$ at $x=0$, of the bounded self-adjoint operator $\exp(-\mathsf{H})$.
The bounded Borel functional calculus is compatible with unitary equivalence;
so $J_g\exp(it\mathsf{H}_j)J_g^{-1}=\exp(it\mathsf{H}_{j'})$, and
\[
J_g\exp(it\mathsf{H}_j)V_j(y)J_g^{-1}=\exp(it\mathsf{H}_{j'})V_{j'}(gy)\qquad\left(t\in\R,\ y\in e_j^\perp\right).
\]
For a Borel set $A$ of $\R\times e_{j'}^\perp$ let
\[
\mathfrak{E}'(A):=J_g\,\mathfrak{E}_j\left(\left\{(\omega,p_\perp)\in\R\times e_j^\perp:(\omega,gp_\perp)\in A\right\}\right)J_g^{-1}.
\]
This is a projection-valued measure on $\R\times e_{j'}^\perp$,
and, as $(gp_\perp)\cdot y'=p_\perp\cdot\left(g^{-1}y'\right)$ for $y'\in e_{j'}^\perp$,
\[
\int\mathrm{d}\mathfrak{E}'(\omega,p')\,\exp\left(it\omega+ip'\cdot y'\right)
=J_g\exp(it\mathsf{H}_j)V_j\left(g^{-1}y'\right)J_g^{-1}
=\exp(it\mathsf{H}_{j'})V_{j'}(y')
\]
for all real $t$ and all $y'\in e_{j'}^\perp$.
By the uniqueness in Lemma~\ref{lem:joint-spectrum}(b), $\mathfrak{E}'=\mathfrak{E}_{j'}$; that is,
\begin{equation}\label{eq:rotation-pair-transport}
J_g[\Phi]_j=\left[\hat g\Phi\right]_{j'},\qquad
J_g\,\mathfrak{E}_j(A)\,J_g^{-1}=\mathfrak{E}_{j'}\left(\left\{(\omega,gp_\perp):(\omega,p_\perp)\in A\right\}\right)
\end{equation}
for $\Phi\in L^2\left(\Sigma_j^+,\mu\right)$ and every Borel set $A$ of $\R\times e_j^\perp$.
The component of $gp_\perp$ along $e_{i'}$ is $p_\perp\cdot g^{-1}e_{i'}=p_i$.
Hence conjugation by $J_g$ takes the joint spectral measure of $\mathsf{H}_j$ and $\mathsf{P}_i$ into the joint spectral measure of $\mathsf{H}_{j'}$ and $\mathsf{P}_{i'}$;
in particular, \eqref{eq:rotation-pair-transport} with $A=\{(\omega,p_\perp):\omega<|p_i|\}$ turns $(C_{ji})$ into
$\mathfrak{E}_{j'}\left(\left\{(\omega,p'):\omega<|p'_{i'}|\right\}\right)=0$, which is $(C_{j'i'})$.

\emph{Rotations.}
The quarter turns of the three coordinate planes belong to $\Ocub$.
By \textup{(H)}, \textup{(dF)} and the first step, $\mu$ satisfies, for every ordered pair $(j',i')$ of distinct indices,
\textup{(G)}, \textup{(dF)}, $(T_{123})$, $(R_{j'})$, $(RP_{j'})$, $(Q_{j'i'})$ and $(C_{j'i'})$.
Theorem~\ref{thm:one-plane} therefore applies to every ordered pair,
and $\mu$ is invariant under all rotations of the three coordinate planes.
For every law on $(\mathcal{S}',\Sigma)$, the elements of $\mathrm{O}(3)$ that leave it invariant form a subgroup of $\mathrm{O}(3)$,
since $(g_1g_2)_*=(g_1)_*\circ(g_2)_*$ for $g_1,g_2\in\mathrm{O}(3)$.
Moreover, let $(j_1,i_1)$ and $(j_2,i_2)$ be any two ordered pairs of distinct indices with $\{j_1,i_1\}\neq\{j_2,i_2\}$,
and let $e_{l_1}$ and $e_{l_2}$, $l_1\neq l_2$, be the unit vectors orthogonal to the planes $(j_1,i_1)$ and $(j_2,i_2)$.
Every element of $\mathrm{SO}(3)$ is of the form $R^{12}_{t_1}R^{23}_{t_2}R^{12}_{t_3}$ with real $t_1,t_2,t_3$ (Euler angles);
here $R^{12}_{t}$ fixes $e_3$ and $R^{23}_{t}$ fixes $e_1$.
Conjugation by a permutation matrix that maps $e_3$ to $e_{l_1}$ and $e_1$ to $e_{l_2}$ maps $\mathrm{SO}(3)$ onto itself;
it maps the elements of $\mathrm{SO}(3)$ that fix $e_3$ onto those that fix $e_{l_1}$, which are the rotations of the plane $(j_1,i_1)$,
and the elements that fix $e_1$ onto those that fix $e_{l_2}$, which are the rotations of the plane $(j_2,i_2)$.
So every element of $\mathrm{SO}(3)$ is the product of a rotation of the plane $(j_1,i_1)$,
a rotation of the plane $(j_2,i_2)$ and a rotation of the plane $(j_1,i_1)$.
Applied to $\mu$, which is left invariant by the rotations of each coordinate plane,
the subgroup property and this decomposition give the invariance of $\mu$ under $\mathrm{SO}(3)$.
The element $-1=\theta_1\theta_2\theta_3$ belongs to $\Ocub$, and therefore $(-1)_*\mu=\mu$ by \textup{(H)}.
As $\mathrm{O}(3)=\mathrm{SO}(3)\cup(-\mathrm{SO}(3))$, the law $\mu$ is invariant under $\mathrm{O}(3)$ and, by $(T_{123})$, under $\R^3\rtimes\mathrm{O}(3)$.

\emph{Moments, mirrors and the two-point function.}
For $g\in\mathrm{O}(3)$, the identity of the moments is~\eqref{eq:rotation-invariant-moments} with $\mathfrak{g}=g$.
The law $\mu$ is invariant under $\R^3\rtimes\mathrm{O}(3)$ and satisfies $(RP_1)$;
Lemma~\ref{lem:transport}(c) gives the invariance under $r_{n,c}$, the reflection positivity at $\{n\cdot x=c\}$
and the positivity of $\left\langle\hat r_{n,c}\Phi,\Phi\right\rangle_\mu$ on $L^2\left(\Sigma_{\{n\cdot x>c\}},\mu\right)$, for every unit vector $n$ and every real $c$.
As to the two-point function, $\mu$ satisfies \textup{(G)}, \textup{(dF)}, $(T_{123})$, $(R_1)$, $(RP_1)$,
and $S_2(g\cdot b,g\cdot b')=S_2(b,b')$ for all $g\in\mathrm{O}(3)$ and $b,b'\in L^2(\R^3;\C^3)$ by~\eqref{eq:rotation-invariant-moments} with $N=2$.
Theorem~\ref{thm:two-point}(c) applies and gives the K\"all\'en--Lehmann form of $S_2$,
with a unique finite positive Borel measure on $[0,\infty)$ of mass at most one.

\emph{The last assertion.}
In this step $\mu$ is the law of the last assertion of the theorem:
it satisfies \textup{(G)}, \textup{(dF)}, $(T_{123})$, $(R_j)$, $(RP_j)$, $(Q_{ji})$, $(C_{ji})$ and $(R_{j'})$, $(RP_{j'})$, $(Q_{j'i'})$, $(C_{j'i'})$
for two ordered pairs $(j,i)$ and $(j',i')$ of distinct indices with $\{j,i\}\neq\{j',i'\}$.
For each of the two pairs, $\mu$ has the seven properties of Theorem~\ref{thm:one-plane};
by that theorem, $\left(R^{ji}_t\right)_*\mu=\mu$ and $\left(R^{j'i'}_t\right)_*\mu=\mu$ for every real $t$.
Neither of the two facts on $\mathrm{O}(3)$ from the step on rotations depends on the law:
the elements of $\mathrm{O}(3)$ that leave a law invariant form a subgroup of $\mathrm{O}(3)$ for every law on $(\mathcal{S}',\Sigma)$,
and the decomposition of $\mathrm{SO}(3)$ holds for any two different coordinate planes.
Taken with $(j_1,i_1)=(j,i)$ and $(j_2,i_2)=(j',i')$, they give the invariance of $\mu$ under $\mathrm{SO}(3)$.
\end{proof}

\subsection{Unitary rotations and boosts}\label{sec:unitary-boosts}

In this subsection the ordered pair is $(1,2)$, as in Definition~\ref{def:labels-plates},
and the reconstruction is the one in the direction $e_1$.
Its joint spectral measure $\mathfrak{E}_1$ is written in the variables $(\omega,p_\perp)$, $p_\perp=(p_2,p_3)\in e_1^\perp$,
and for $\psi\in\mathcal{H}_1$ the positive measure $\varrho_\psi:=\left\langle\psi,\mathfrak{E}_1(\cdot)\,\psi\right\rangle$ is the spectral measure of $\psi$
(Lemma~\ref{lem:joint-spectrum}).

\begin{proposition}[Unitary rotations, the round cone and boosts]\label{prop:rotations-boosts}
\begin{enumerate}
\item\label{it:unitary-rotations}
\emph{Spatial rotations.}
Let $\mu$ be a law on $\mathcal{S}'$ that satisfies $(T_{123})$, $(R_1)$ and $(RP_1)$ and is invariant under the group
\[
\mathrm{O}(2)_1:=\left\{g\in\mathrm{O}(3):\ ge_1=e_1\right\}.
\]
Then, for every $g\in\mathrm{O}(2)_1$, the assignment
\[
[\Phi]_1\longmapsto[\hat g\Phi]_1\qquad\left(\Phi\in L^2(\Sigma_1^+,\mu)\right)
\]
extends to a unitary operator $\mathcal{U}(g)$ of $\mathcal{H}_1$,
and $\mathcal{U}$ is a strongly continuous unitary representation of $\mathrm{O}(2)_1$ on $\mathcal{H}_1$.
For $g\in\mathrm{O}(2)_1$, $t\geq0$ and $y\in e_1^\perp$,
\begin{align*}
\mathcal{U}(g)\Omega_1&=\Omega_1,\\
\mathcal{U}(g)\exp(-t\cdot\mathsf{H}_1)&=\exp(-t\cdot\mathsf{H}_1)\,\mathcal{U}(g),\\
\mathcal{U}(g)V_1(y)\,\mathcal{U}(g)^{-1}&=V_1(gy),
\end{align*}
and for every Borel subset $A$ of $[0,\infty)\times e_1^\perp$
\[
\mathcal{U}(g)\,\mathfrak{E}_1(A)\,\mathcal{U}(g)^{-1}=\mathfrak{E}_1(gA),
\qquad
gA:=\left\{(\omega,gp_\perp):\ (\omega,p_\perp)\in A\right\}.
\]
\item\label{it:round-cone}
\emph{The round cone.}
If $\mu$ satisfies, besides the hypotheses of~\eqref{it:unitary-rotations}, the cone condition $(C_{12})$,
then $\mathfrak{E}_1$ is carried by the round forward cone
\[
\left\{(\omega,p_\perp):\ \omega\geq|p_\perp|\right\},
\qquad
|p_\perp|:=\sqrt{p_2^2+p_3^2}\, .
\]
Consequently, for every unit vector $n=n_2e_2+n_3e_3$ of $e_1^\perp$,
\[
\mathfrak{E}_1\left(\left\{(\omega,p_\perp):\ \omega<\left|n\cdot p_\perp\right|\right\}\right)=0,
\qquad
n\cdot p_\perp:=n_2p_2+n_3p_3,
\]
that is, $\mathsf{H}_1\geq|n\cdot\mathsf{P}|$ with the constant one,
where $n\cdot\mathsf{P}$ is the generator of the unitary group $u\mapsto V_1(u\cdot n)$;
and every $\psi$ in the domain of $\mathsf{H}_1$ lies in the domains of $\mathsf{P}_2$ and $\mathsf{P}_3$, with
\[
\left\|\mathsf{H}_1\psi\right\|^2\geq\left\|\mathsf{P}_2\psi\right\|^2+\left\|\mathsf{P}_3\psi\right\|^2 .
\]
In particular, if a law $\mu$ on $\mathcal{S}'$ satisfies \eqref{cond:soft-hypothesis}, \eqref{cond:closedness}
and $(C_{kl})$ for one ordered pair $k\neq l$ of indices in $\{1,2,3\}$,
then for each $j\in\{1,2,3\}$ the joint spectral measure $\mathfrak{E}_j$ of the reconstruction in the direction $e_j$
is carried by $\left\{\omega\geq|p_\perp|\right\}$, where now $p_\perp=(p_i)_{i\neq j}$ and $|p_\perp|$ is its Euclidean norm,
and the two consequences above hold with $\mathsf{H}_j$, the generators $\mathsf{P}_i$, $i\neq j$, and the unit vectors of $e_j^\perp$.
\item\label{it:unitary-boosts}
\emph{Boosts.}
Let $\mu$ satisfy \eqref{cond:gaussian-domination}, \eqref{cond:closedness}, $(T_{123})$, $(R_1)$, $(RP_1)$, $(Q_{12})$ and $(C_{12})$,
and let $\mathcal{H}^{\Psi}_{12}$ be the closed span in $\mathcal{H}_1$ of the plate vectors $\Psi(x;\lambda)$ of Definition~\ref{def:plate-kernels},
where $n\geq0$, $x\in\mathcal{W}^{(n)}$ and $\lambda$ is a tuple of $n$ labels.
Then $\mathcal{H}^{\Psi}_{12}$ contains $\Omega_1$ and the class $\left[\prod_{a=1}^N\mathbb{B}_{\lambda_a}(\alpha_a)\right]_1$
for every $N\geq1$, all labels $\lambda_1,\dots,\lambda_N$ and all $\alpha_1,\dots,\alpha_N\in L^2(\R^2;\C)$
with time-ordered supports in $\mathbb{H}$ (Definition~\ref{def:plate-kernels});
it reduces the operators $\mathsf{T}(\kappa)$, $\kappa\in\overline{\mathcal{W}}$;
and it is invariant under the point plates $\mathcal{B}_{\lambda_0}(\xi,\xi')$, where $\lambda_0$ is a label and $(\xi,\xi')\in\mathcal{W}\times\mathcal{W}$.
There is exactly one strongly continuous unitary group $\left(\mathfrak{L}(\beta)\right)_{\beta\in\R}$ on $\mathcal{H}^{\Psi}_{12}$ such that
\begin{equation}\label{eq:rotation-boosts}
\mathfrak{L}(\beta)\Psi(x;\lambda)=\exp\left(w(\lambda)\cdot\beta\right)\cdot\Psi\left(\Lambda_\beta x;\lambda\right)
\end{equation}
for all real $\beta$, all $n\geq0$, all $x\in\mathcal{W}^{(n)}$ and all tuples $\lambda$ of $n$ labels.
For real $\beta$, $\kappa\in\overline{\mathcal{W}}$, $(\xi,\xi')\in\mathcal{W}\times\mathcal{W}$ and every label $\lambda_0=(\sigma,\chi)$ one has, on $\mathcal{H}^{\Psi}_{12}$,
\begin{align*}
\mathfrak{L}(\beta)\Omega_1&=\Omega_1,\\
\mathfrak{L}(\beta)\mathsf{T}(\kappa)\mathfrak{L}(\beta)^{-1}&=\mathsf{T}(\Lambda_\beta\kappa),\\
\mathfrak{L}(\beta)\mathcal{B}_{\lambda_0}(\xi,\xi')\mathfrak{L}(\beta)^{-1}&=\exp\left(w(\sigma)\cdot\beta\right)\cdot\mathcal{B}_{\lambda_0}\left(\Lambda_\beta\xi,\Lambda_\beta\xi'\right),
\end{align*}
and, for $t\geq0$ and real $u$,
\begin{align*}
\mathfrak{L}(\beta)\exp(-t\cdot\mathsf{H}_1)\,\mathfrak{L}(\beta)^{-1}&=\exp\left(-t\cdot\left(\mathsf{H}_1\cdot\cosh\beta+\mathsf{P}_2\cdot\sinh\beta\right)\right),\\
\mathfrak{L}(\beta)\exp(i\,u\cdot\mathsf{P}_2)\,\mathfrak{L}(\beta)^{-1}&=\exp\left(i\,u\cdot\left(\mathsf{P}_2\cdot\cosh\beta+\mathsf{H}_1\cdot\sinh\beta\right)\right),
\end{align*}
where the right sides are the integrals against $\mathfrak{E}_1$ of the functions
$(\omega,p_\perp)\mapsto\exp\left(-t\cdot\left(\omega\cdot\cosh\beta+p_2\cdot\sinh\beta\right)\right)$
and $(\omega,p_\perp)\mapsto\exp\left(i\,u\cdot\left(p_2\cdot\cosh\beta+\omega\cdot\sinh\beta\right)\right)$.
\end{enumerate}
\end{proposition}

\begin{proof}
\emph{Part~\eqref{it:unitary-rotations}.}
Let $g\in\mathrm{O}(2)_1$.
The operator $\hat g$ acts on functions on $\mathcal{S}'$ by composition with a map of $\mathcal{S}'$;
so it is multiplicative and commutes with complex conjugation.
Since $g_*\mu=\mu$, it is unitary on $L^2(\mu)$, and $\mu(\hat g\Phi)=\mu(\Phi)$ for $\Phi\in L^1(\mu)$;
moreover $\widehat{gh}=\hat g\hat h$ for $h\in\mathrm{O}(3)$,
and $\hat g$ maps $\Sigma_U$-measurable functions onto $\Sigma_{gU}$-measurable functions for every open $U\subset\R^3$
\cite[Section~8.1.2]{NoncompactQED}.
As $g\{x_1>0\}=\{x_1>0\}$, the operator $\hat g$ maps $L^2(\Sigma_1^+,\mu)$ onto itself.
The matrix $g$ fixes $e_1$ and maps $e_1^\perp$ onto itself,
while $\theta_1$ is $-1$ on $\R e_1$ and the identity on $e_1^\perp$;
hence $g\theta_1=\theta_1g$, and $\hat g\hat\theta_1=\hat\theta_1\hat g$.
Since $\Theta_1\Phi$ is the complex conjugate of $\hat\theta_1\Phi$, also $\hat g\Theta_1=\Theta_1\hat g$.
For $\Phi,\Phi'\in L^2(\Sigma_1^+,\mu)$ it follows that
\[
\left\langle[\hat g\Phi]_1,[\hat g\Phi']_1\right\rangle
=\mu\left(\left(\Theta_1\hat g\Phi\right)\cdot\hat g\Phi'\right)
=\mu\left(\hat g\left(\left(\Theta_1\Phi\right)\cdot\Phi'\right)\right)
=\mu\left(\left(\Theta_1\Phi\right)\cdot\Phi'\right)
=\left\langle[\Phi]_1,[\Phi']_1\right\rangle .
\]
If $[\Phi]_1=[\Phi'']_1$, this identity, applied to $\Phi-\Phi''$, gives
$\left\|[\hat g\Phi]_1-[\hat g\Phi'']_1\right\|=\left\|[\Phi-\Phi'']_1\right\|=0$.
So the assignment $[\Phi]_1\mapsto[\hat g\Phi]_1$ is well defined, linear and isometric on the classes.
The canonical map $[\,\cdot\,]_1$ of $L^2(\Sigma_1^+,\mu)$ into $\mathcal{H}_1$ is a contraction with dense range, and $\Omega_1=[1]_1$
\cite[Proposition~8.19\,(i)]{NoncompactQED};
there $L^2(\Sigma_1^+,\mu)$ is denoted $\mathcal{E}_+$, and the semigroup $\exp(-t\cdot\mathsf{H}_1)$ is denoted $T(t)$.
Hence the assignment extends to an isometry $\mathcal{U}(g)$ of $\mathcal{H}_1$,
and $\mathcal{U}(g)\Omega_1=\Omega_1$ because $\hat g1=1$.
The rule $\widehat{gh}=\hat g\hat h$ gives $\mathcal{U}(gh)=\mathcal{U}(g)\,\mathcal{U}(h)$ for $g,h\in\mathrm{O}(2)_1$ on the classes, hence on $\mathcal{H}_1$, and $\mathcal{U}(1)=1$;
thus $\mathcal{U}(g^{-1})$ is an inverse of $\mathcal{U}(g)$, and $\mathcal{U}(g)$ is unitary.

\emph{The commutation relations.}
Lemma~\ref{lem:transport}\eqref{it:transport-operators}, applied to $g$ and to the law $\left(g^{-1}\right)_*\mu=\mu$,
gives $\hat gU(a)\hat g^{-1}=U(ga)$ for $a\in\R^3$.
For $\Phi\in L^2(\Sigma_1^+,\mu)$, $t\geq0$ and $y\in e_1^\perp$,
\[
\exp(-t\cdot\mathsf{H}_1)[\Phi]_1=\left[U(t\cdot e_1)\Phi\right]_1,
\qquad
V_1(y)[\Phi]_1=\left[U(y)\Phi\right]_1
\]
\cite[Proposition~8.19\,(ii),~(iii)]{NoncompactQED}.
As $ge_1=e_1$ and $gy\in e_1^\perp$,
\begin{align*}
\mathcal{U}(g)\exp(-t\cdot\mathsf{H}_1)[\Phi]_1&=\left[\hat gU(t\cdot e_1)\Phi\right]_1=\left[U(t\cdot e_1)\hat g\Phi\right]_1=\exp(-t\cdot\mathsf{H}_1)\,\mathcal{U}(g)[\Phi]_1,\\
\mathcal{U}(g)V_1(y)[\Phi]_1&=\left[\hat gU(y)\Phi\right]_1=\left[U(gy)\hat g\Phi\right]_1=V_1(gy)\,\mathcal{U}(g)[\Phi]_1,
\end{align*}
and both identities extend from the dense set of the classes to $\mathcal{H}_1$.

\emph{Continuity.}
The canonical map $[\,\cdot\,]_1$ being a contraction,
$\left\|\mathcal{U}(g)[\Phi]_1-[\Phi]_1\right\|\leq\left\|\hat g\Phi-\Phi\right\|_{L^2(\mu)}$ for $\Phi\in L^2(\Sigma_1^+,\mu)$.
Let $f\in C^{\infty}_{\mathrm c}(\{x_1>0\};\Lambda^2)$ and $\Phi:=\exp\left(iF(f)\right)$.
Then $\hat g\Phi=\exp\left(iF(g\cdot f)\right)$ \cite[Section~8.1.2]{NoncompactQED}, and, $f$ being real,
\[
\left\|\hat g\Phi-\Phi\right\|^2_{L^2(\mu)}=\int\mathrm{d}\mu(F)\,\left|\exp\left(iF(g\cdot f-f)\right)-1\right|^2 .
\]
Let $g_k\to1$ in $\mathrm{O}(2)_1$.
Then $g_k\cdot f\to f$ in $\mathcal{S}$, since all derivatives converge uniformly and the supports stay in a fixed ball;
so $F(g_k\cdot f-f)\to0$ for every $F\in\mathcal{S}'$.
The integrand is at most $4$, and dominated convergence gives $\left\|\widehat{g_k}\Phi-\Phi\right\|_{L^2(\mu)}\to0$.
By \cite[Lemma~8.12\,(i)]{NoncompactQED}, the finite linear combinations with coefficients in $\mathbb{Q}+i\mathbb{Q}$ of the functions $\exp\left(iF(f)\right)$,
with $f$ in a countable set of real forms of $C^{\infty}_{\mathrm c}(\{x_1>0\};\Lambda^2)$,
form a dense subset of $L^2(\Sigma_1^+,\mu)$; it is the set denoted $\mathcal{A}^0(\{x_1>0\})$ there.
A fortiori the linear combinations of the functions $\Phi$ above are dense in $L^2(\Sigma_1^+,\mu)$,
and their classes are dense in $\mathcal{H}_1$, the canonical map being continuous with dense range.
For $\psi\in\mathcal{H}_1$ and the class $\psi'$ of such a combination,
\[
\left\|\mathcal{U}(g_k)\psi-\psi\right\|\leq2\left\|\psi-\psi'\right\|+\left\|\mathcal{U}(g_k)\psi'-\psi'\right\|,
\]
because $\mathcal{U}(g_k)$ is unitary; hence $\mathcal{U}(g_k)\psi\to\psi$.
Strong continuity at an arbitrary $g_0\in\mathrm{O}(2)_1$ follows from
$\mathcal{U}(g)\psi-\mathcal{U}(g_0)\psi=\mathcal{U}(g_0)\left(\mathcal{U}\left(g_0^{-1}g\right)\psi-\psi\right)$.

\emph{The spectral measure.}
The second relation with $t=1$ says that $\mathcal{U}(g)$ commutes with the bounded self-adjoint operator $\exp(-\mathsf{H}_1)$.
A bounded operator that commutes with a bounded self-adjoint operator commutes with each bounded Borel function of it,
by the functional calculus for bounded Borel functions \cite[Theorem~VII.2\,(g)]{ReedSimon1980}.
For real $t$ let $h_t(x):=\exp(-i\,t\cdot\log x)$ for $0<x\leq1$ and $h_t(0):=0$.
Since $h_t\left(\exp(-\omega)\right)=\exp(i\,t\cdot\omega)$ for $\omega\geq0$, one has $h_t\left(\exp(-\mathsf{H}_1)\right)=\exp(i\,t\cdot\mathsf{H}_1)$;
so $\mathcal{U}(g)$ commutes with $\exp(i\,t\cdot\mathsf{H}_1)$ for every real $t$.
For a Borel set $A\subset\R\times e_1^\perp$ let $gA$ be formed as in~\eqref{it:unitary-rotations}, and put
\[
\mathfrak{E}^{g}(A):=\mathcal{U}(g)^{-1}\,\mathfrak{E}_1(gA)\,\mathcal{U}(g).
\]
As $gA$ is the image of $A$ under the homeomorphism $(\omega,p_\perp)\mapsto(\omega,gp_\perp)$ of $\R\times e_1^\perp$,
$\mathfrak{E}^{g}$ is a projection-valued measure,
and integrals against $\mathfrak{E}^{g}$ arise from integrals against $\mathfrak{E}_1$ by the substitution $p_\perp\mapsto g^{-1}p_\perp$, followed by conjugation with $\mathcal{U}(g)$.
By Lemma~\ref{lem:joint-spectrum}(b), $\mathfrak{E}_1$ is the only projection-valued measure on $\R\times e_1^\perp$ with
\[
\exp(i\,t\cdot\mathsf{H}_1)V_1(y)=\int\mathrm{d}\mathfrak{E}_1(\omega,p_\perp)\,\exp\left(i\,t\cdot\omega+i\,p_\perp\cdot y\right)
\qquad\left(t\in\R,\ y\in e_1^\perp\right).
\]
Because $\left(g^{-1}p_\perp\right)\cdot y=p_\perp\cdot(gy)$, it follows that
\begin{align*}
&\int\mathrm{d}\mathfrak{E}^{g}(\omega,p_\perp)\,\exp\left(i\,t\cdot\omega+i\,p_\perp\cdot y\right)\\
&\qquad=\mathcal{U}(g)^{-1}\left(\int\mathrm{d}\mathfrak{E}_1(\omega,p_\perp)\,\exp\left(i\,t\cdot\omega+i\,p_\perp\cdot(gy)\right)\right)\mathcal{U}(g)\\
&\qquad=\mathcal{U}(g)^{-1}\exp(i\,t\cdot\mathsf{H}_1)V_1(gy)\,\mathcal{U}(g)
=\exp(i\,t\cdot\mathsf{H}_1)V_1(y);
\end{align*}
in the last step $\mathcal{U}(g)$ commutes with $\exp(i\,t\cdot\mathsf{H}_1)$, and $\mathcal{U}(g)^{-1}V_1(gy)\,\mathcal{U}(g)=V_1(y)$.
The uniqueness gives $\mathfrak{E}^{g}=\mathfrak{E}_1$, which is the last relation of~\eqref{it:unitary-rotations}.

\emph{Part~\eqref{it:round-cone}.}
Let $n$ be a unit vector of $e_1^\perp$, and let $g\in\mathrm{O}(2)_1$ be a rotation about the axis $\R e_1$ with $ge_2=n$.
The cone condition $(C_{12})$ says that $A:=\left\{(\omega,p_\perp):\ \omega<|p_2|\right\}$ is an $\mathfrak{E}_1$-null set.
Since $e_2\cdot\left(g^{-1}p_\perp\right)=(ge_2)\cdot p_\perp=n\cdot p_\perp$,
the set $gA$ is $\left\{(\omega,p_\perp):\ \omega<|n\cdot p_\perp|\right\}$, and~\eqref{it:unitary-rotations} gives
\[
\mathfrak{E}_1\left(\left\{\omega<\left|n\cdot p_\perp\right|\right\}\right)=\mathcal{U}(g)\,\mathfrak{E}_1(A)\,\mathcal{U}(g)^{-1}=0 .
\]
By Lemma~\ref{lem:joint-spectrum}(b) with $t=0$,
$V_1(u\cdot n)=\int\mathrm{d}\mathfrak{E}_1(\omega,p_\perp)\,\exp\left(i\,u\cdot(n\cdot p_\perp)\right)$ for real $u$;
so $n\cdot\mathsf{P}$ is the integral of $n\cdot p_\perp$ against $\mathfrak{E}_1$,
and the identity just obtained is the inequality $\mathsf{H}_1\geq|n\cdot\mathsf{P}|$ in the sense of the statement.
Now let $(n_k)_{k\geq1}$ be a dense sequence in the unit circle of $e_1^\perp$.
If $\omega<|p_\perp|$, then $p_\perp\neq0$;
as $n\mapsto n\cdot p_\perp$ is continuous on the unit circle and takes the value $|p_\perp|$ at $p_\perp/|p_\perp|$,
there is $k$ with $\omega<|n_k\cdot p_\perp|$.
Hence $\left\{\omega<|p_\perp|\right\}$ is the union of the $\mathfrak{E}_1$-null sets $\left\{\omega<|n_k\cdot p_\perp|\right\}$, $k\geq1$,
and is itself $\mathfrak{E}_1$-null: $\mathfrak{E}_1$ is carried by $\left\{\omega\geq|p_\perp|\right\}$.
Conversely, the null sets $\left\{\omega<|n\cdot p_\perp|\right\}$ follow from this support property:
by the Cauchy--Schwarz inequality, $|n\cdot p_\perp|\leq|p_\perp|$ for every unit vector $n$ of $e_1^\perp$,
so that $\left\{\omega<|n\cdot p_\perp|\right\}\subset\left\{\omega<|p_\perp|\right\}$.

Let $\psi$ lie in the domain of $\mathsf{H}_1$.
By Lemma~\ref{lem:joint-spectrum}(b),~(c) and the spectral theorem,
$\mathsf{H}_1$ and $\mathsf{P}_i$ are the integrals of $\omega$ and of $p_i$ against $\mathfrak{E}_1$,
and $\left\|\mathsf{H}_1\psi\right\|^2=\int\mathrm{d}\varrho_\psi(\omega,p_\perp)\,\omega^2<\infty$.
As $p_2^2+p_3^2\leq\omega^2$ holds $\varrho_\psi$-almost everywhere, $p_2$ and $p_3$ are square integrable for $\varrho_\psi$;
so $\psi$ lies in the domains of $\mathsf{P}_2$ and $\mathsf{P}_3$, and
\[
\left\|\mathsf{P}_2\psi\right\|^2+\left\|\mathsf{P}_3\psi\right\|^2
=\int\mathrm{d}\varrho_\psi(\omega,p_\perp)\,\left(p_2^2+p_3^2\right)
\leq\int\mathrm{d}\varrho_\psi(\omega,p_\perp)\,\omega^2
=\left\|\mathsf{H}_1\psi\right\|^2 .
\]

Let now $\mu$ satisfy \eqref{cond:soft-hypothesis}, \eqref{cond:closedness} and $(C_{kl})$ for one ordered pair $k\neq l$.
By Theorem~\ref{thm:orthogonal-group}, $\mu$ is invariant under $\R^3\rtimes\mathrm{O}(3)$ and satisfies $(C_{ji})$ for all ordered pairs $j\neq i$.
Together with~\eqref{cond:soft-hypothesis} this gives $(T_{123})$, and $(R_j)$, $(RP_j)$ for $j=1,2,3$,
as well as the invariance under $\mathrm{O}(2)_1$;
so the case $j=1$ is the first part of~\eqref{it:round-cone}.
Let $j\in\{2,3\}$, choose $i\neq j$, and let $g_{ji}$ be the permutation matrix with $g_{ji}e_1=e_j$ and $g_{ji}e_2=e_i$.
As $\mu$ is invariant under $\mathrm{O}(3)$, $\left(g_{ji}^{-1}\right)_*\mu=\mu$,
and Lemma~\ref{lem:transport}\eqref{it:transport-pairs} provides a unitary operator $J$ of $\mathcal{H}_1$ onto $\mathcal{H}_j$ with
\[
\mathfrak{E}_j(A)=J\,\mathfrak{E}_1\left(\left\{(\omega,p_\perp):\ \left(\omega,g_{ji}p_\perp\right)\in A\right\}\right)J^{-1}
\]
for every Borel set $A\subset[0,\infty)\times e_j^\perp$.
Since $|g_{ji}p_\perp|=|p_\perp|$, the set $\left\{\omega<|p_\perp|\right\}$ of $[0,\infty)\times e_j^\perp$
corresponds to the $\mathfrak{E}_1$-null set $\left\{\omega<|p_\perp|\right\}$ of $[0,\infty)\times e_1^\perp$.
Thus $\mathfrak{E}_j$ is carried by $\left\{\omega\geq|p_\perp|\right\}$.
Let $n$ be a unit vector of $e_j^\perp$.
Again $|n\cdot p_\perp|\leq|p_\perp|$, so that $\left\{\omega<|n\cdot p_\perp|\right\}$ lies in the $\mathfrak{E}_j$-null set $\left\{\omega<|p_\perp|\right\}$ and is itself $\mathfrak{E}_j$-null.
Lemma~\ref{lem:joint-spectrum}(b) for the direction $e_j$, with $t=0$, gives
$V_j(u\cdot n)=\int\mathrm{d}\mathfrak{E}_j(\omega,p_\perp)\,\exp\left(i\,u\cdot(n\cdot p_\perp)\right)$ for real $u$;
so the generator $n\cdot\mathsf{P}$ of $u\mapsto V_j(u\cdot n)$ is the integral of $n\cdot p_\perp$ against $\mathfrak{E}_j$,
and $\mathsf{H}_j\geq|n\cdot\mathsf{P}|$ in the sense of the statement.
The norm inequality, together with the inclusion of the domains, follows from the support property of $\mathfrak{E}_j$
by the computation made above for $j=1$, with Lemma~\ref{lem:joint-spectrum}(b),~(c) for the direction $e_j$,
and with $\mathfrak{E}_j$, $\mathsf{H}_j$ and the generators of the translations of $e_j^\perp$ in place of $\mathfrak{E}_1$, $\mathsf{H}_1$, $\mathsf{P}_2$ and $\mathsf{P}_3$.

\emph{Part~\eqref{it:unitary-boosts}: the subspace $\mathcal{H}^{\Psi}_{12}$.}
The plate vector with $n=0$ is $\Omega_1$.
Let $\alpha_1,\dots,\alpha_N$ and $\lambda=(\lambda_1,\dots,\lambda_N)$ be as in~\eqref{it:unitary-boosts}.
By Proposition~\ref{prop:kernels}\eqref{it:kernel-density},
\[
\left[\prod_{a=1}^N\mathbb{B}_{\lambda_a}(\alpha_a)\right]_1=\int\mathrm{d}Y\,\prod_{a=1}^N\alpha_a(y_a)\cdot\Psi(Y;\lambda),
\]
a Bochner integral over the compact product of the supports;
this product consists of strictly time-ordered configurations in $\mathbb{H}$, so it lies in $\mathcal{W}^{(N)}$ (Definition~\ref{def:plate-kernels}).
The integrand takes its values in the closed subspace $\mathcal{H}^{\Psi}_{12}$, and so does the integral.

Let $\kappa\in\overline{\mathcal{W}}$, $n\geq1$ and $x\in\mathcal{W}^{(n)}$.
The configuration $x+(\kappa,\dots,\kappa)$ lies in $\mathcal{W}^{(n)}$,
because its differences are those of $x$ and $d_{\mathcal{W}}(x_1+\kappa)\geq d_{\mathcal{W}}(x_1)+d_{\mathcal{W}}(\kappa)>0$.
By Proposition~\ref{prop:kernels}\eqref{it:plate-vectors},
\[
\mathsf{T}(\kappa)\Psi(x;\lambda)=\Psi\left(x+(\kappa,\dots,\kappa);\lambda\right),
\]
and $\mathsf{T}(\kappa)\Omega_1=\Omega_1$ by Lemma~\ref{lem:cone-semigroup}\eqref{it:semigroup-contraction}.
So $\mathsf{T}(\kappa)$ maps $\mathcal{H}^{\Psi}_{12}$ into itself.
By the same lemma $\mathsf{T}(\kappa)^*=\mathsf{T}(\kappa^*)$,
and $\kappa^*\in\overline{\mathcal{W}}$ since $d_{\mathcal{W}}(\kappa^*)=d_{\mathcal{W}}(\kappa)$ (Lemma~\ref{lem:forward-tube}\eqref{it:tube-adjoint});
so $\mathsf{T}(\kappa)^*$ maps $\mathcal{H}^{\Psi}_{12}$ into itself as well, and $\mathcal{H}^{\Psi}_{12}$ reduces $\mathsf{T}(\kappa)$.

Let $\lambda_0$ be a label, $(\xi,\xi')\in\mathcal{W}\times\mathcal{W}$, $x\in\mathcal{W}^{(n)}$ and $\lambda=(\lambda_1,\dots,\lambda_n)$, and put
\[
\tilde x:=\left(\xi,\ \xi+\xi'+x_1,\ \dots,\ \xi+\xi'+x_n\right)\in(\C^2)^{n+1}.
\]
Since $\mathcal{W}$ is a convex cone, $\tilde x_1=\xi$ and $\tilde x_2-\tilde x_1=\xi'+x_1$ lie in $\mathcal{W}$,
and $\tilde x_{a+1}-\tilde x_a=x_a-x_{a-1}\in\mathcal{W}$ for $2\leq a\leq n$; thus $\tilde x\in\mathcal{W}^{(n+1)}$.
The identity
\begin{equation}\label{eq:rotation-plate-product}
\mathcal{B}_{\lambda_0}(\xi,\xi')\Psi(x;\lambda)=\Psi\left(\tilde x;\lambda_0,\lambda_1,\dots,\lambda_n\right)
\end{equation}
holds.
For $n=0$ both sides equal $\mathcal{B}_{\lambda_0}(\xi,\kappa)\Omega_1$ for every $\kappa\in\mathcal{W}$,
the last argument in the definition of the plate vectors being arbitrary (Definition~\ref{def:plate-kernels}).
For $n\geq1$ the first factor of $\Psi(x;\lambda)$ is $\mathcal{B}_{\lambda_1}(x_1,\xi'')$,
with $\xi''=\frac{1}{2}(x_2-x_1)$ if $n\geq2$ and $\xi''\in\mathcal{W}$ arbitrary if $n=1$.
The inner arguments $\xi'$, $x_1$ and $\frac{1}{2}(\xi'+x_1)$, $\frac{1}{2}(\xi'+x_1)$ lie in $\mathcal{W}$ and have the same sum;
so Lemma~\ref{lem:inner-arguments} gives
\begin{align*}
&\mathcal{B}_{\lambda_0}(\xi,\xi')\mathcal{B}_{\lambda_1}(x_1,\xi'')\\
&\qquad=\mathcal{B}_{\lambda_0}\left(\tilde x_1,\frac{1}{2}\left(\tilde x_2-\tilde x_1\right)\right)\mathcal{B}_{\lambda_1}\left(\frac{1}{2}\left(\tilde x_2-\tilde x_1\right),\xi''\right).
\end{align*}
As $\tilde x_{a+1}-\tilde x_a=x_a-x_{a-1}$ for $a\geq2$,
the product of this operator with the remaining factors of $\Psi(x;\lambda)$, applied to $\Omega_1$, is $\Psi(\tilde x;\lambda_0,\lambda_1,\dots,\lambda_n)$.
In particular $\mathcal{H}^{\Psi}_{12}$ is invariant under the bounded operator $\mathcal{B}_{\lambda_0}(\xi,\xi')$.

\emph{The group $\mathfrak{L}$.}
By Lemma~\ref{lem:forward-tube}\eqref{it:tube-boosts}, $\Lambda_\beta\mathcal{W}=\mathcal{W}$ and $\left(\Lambda_\beta\xi\right)^*=\Lambda_\beta\left(\xi^*\right)$;
as $\Lambda_\beta$ is linear and invertible, it maps $\overline{\mathcal{W}}$ onto itself, and, acting entrywise, $\mathcal{W}^{(n)}$ onto $\mathcal{W}^{(n)}$.
Let $x\in\mathcal{W}^{(n)}$ and $x'\in\mathcal{W}^{(n')}$, let $\lambda$ and $\lambda'$ be tuples of $n$ and of $n'$ labels,
and put $N:=n+n'$ and $\lambda'':=(\lambda_n,\dots,\lambda_1,\lambda'_1,\dots,\lambda'_{n'})$, so that $w(\lambda'')=w(\lambda)+w(\lambda')$;
for the empty tuple the weight is $0$ and the sign is $1$, as an empty sum and an empty product.
For $N\geq1$, Proposition~\ref{prop:kernels}\eqref{it:plate-vectors} writes $\left\langle\Psi(x;\lambda),\Psi(x';\lambda')\right\rangle$ as $s(\lambda)\cdot G_N(\xi;\lambda'')$:
if $n\geq1$ and $n'\geq1$, by the Gram formula~\eqref{eq:rotation-gram},
and if $n=0$ or $n'=0$, by the formula given there for $\left\langle\Omega_1,\Psi(x';\lambda')\right\rangle$ or for $\left\langle\Psi(x;\lambda),\Omega_1\right\rangle$.
In each case the entries of $\xi=(\xi_1,\dots,\xi_{N-1})$ are formed from the entries of $x$ and of $x'$ by differences, sums and the involution ${}^*$ of $\C^2$.
These operations commute with $\Lambda_\beta$; hence
\[
\left\langle\Psi(\Lambda_\beta x;\lambda),\Psi(\Lambda_\beta x';\lambda')\right\rangle
=s(\lambda)\cdot G_N\left(\Lambda_\beta\xi_1,\dots,\Lambda_\beta\xi_{N-1};\lambda''\right).
\]
Let $N\geq2$.
The entries $\xi_a$ lie in $\mathcal{W}$, because $\mathcal{W}^*=\mathcal{W}$ and $\mathcal{W}+\mathcal{W}\subset\mathcal{W}$.
Let $Z=(z_1,\dots,z_N)$ be given by $z_1:=0$ and $z_{a+1}:=z_a+\xi_a$.
Then $Z$ and $\Lambda_\beta Z$ lie in the tube $\mathcal{Z}(\mathrm{id},0)$ of the identity permutation and the frame $0$,
with the consecutive differences $\xi_a$ and $\Lambda_\beta\xi_a$, respectively.
On this tube $\mathcal{G}_N(\cdot;\lambda'')$ is the function $G^{\mathrm{id}}_0(\cdot;\lambda'')$, that is, $G_N$ evaluated at the consecutive differences
(Definition~\ref{def:configuration-tubes} and Lemma~\ref{lem:one-function}\eqref{it:tube-gluing}).
Proposition~\ref{prop:anomaly-vanishes} therefore gives
\begin{align*}
G_N\left(\Lambda_\beta\xi;\lambda''\right)
&=\mathcal{G}_N\left(\Lambda_\beta Z;\lambda''\right)
=\exp\left(-w(\lambda'')\cdot\beta\right)\cdot\mathcal{G}_N\left(Z;\lambda''\right)\\
&=\exp\left(-\left(w(\lambda)+w(\lambda')\right)\cdot\beta\right)\cdot G_N\left(\xi;\lambda''\right).
\end{align*}
As the factors $\exp(w(\lambda)\cdot\beta)$ and $\exp(w(\lambda')\cdot\beta)$ are real, it follows that
\begin{align*}
&\left\langle\exp\left(w(\lambda)\cdot\beta\right)\cdot\Psi(\Lambda_\beta x;\lambda),\ \exp\left(w(\lambda')\cdot\beta\right)\cdot\Psi(\Lambda_\beta x';\lambda')\right\rangle\\
&\qquad=\left\langle\Psi(x;\lambda),\Psi(x';\lambda')\right\rangle .
\end{align*}
The same identity holds for $N\leq1$:
for $N=0$ both sides equal $\left\langle\Omega_1,\Omega_1\right\rangle=1$;
for $N=1$ the formulas of Proposition~\ref{prop:kernels}\eqref{it:plate-vectors} quoted above make both sides multiples of $G_1(\lambda'')$,
and $G_1=0$ by Proposition~\ref{prop:kernels}\eqref{it:kernel-bound}.

Consequently the assignment
\[
\sum_ka_k\cdot\Psi\left(x^{(k)};\lambda^{(k)}\right)\longmapsto\sum_ka_k\cdot\exp\left(w\left(\lambda^{(k)}\right)\cdot\beta\right)\cdot\Psi\left(\Lambda_\beta x^{(k)};\lambda^{(k)}\right)
\]
on finite linear combinations of plate vectors, with complex coefficients $a_k$, preserves inner products.
It is therefore well defined, a vanishing combination being mapped to a vector of norm zero,
and it is a linear isometry of the span of the plate vectors into $\mathcal{H}^{\Psi}_{12}$;
its continuous extension is an isometry $\mathfrak{L}(\beta)$ of $\mathcal{H}^{\Psi}_{12}$ that satisfies~\eqref{eq:rotation-boosts}.
As $R_\zeta$ multiplies the coordinates $\xi^+$ and $\xi^-$ on $\C^2$ by $\exp(i\zeta)$ and $\exp(-i\zeta)$,
one has $R_\zeta R_{\zeta'}=R_{\zeta+\zeta'}$ for complex angles $\zeta,\zeta'$;
hence $\Lambda_\beta\Lambda_{\beta'}=\Lambda_{\beta+\beta'}$, and $\Lambda_0$ is the identity.
By~\eqref{eq:rotation-boosts}, $\mathfrak{L}(\beta)\mathfrak{L}(\beta')=\mathfrak{L}(\beta+\beta')$ and $\mathfrak{L}(0)=1$ on the plate vectors, hence on $\mathcal{H}^{\Psi}_{12}$;
so $\mathfrak{L}(-\beta)$ is an inverse of $\mathfrak{L}(\beta)$, and $\mathfrak{L}(\beta)$ is unitary.
For fixed $x$ and $\lambda$ the map $\beta\mapsto\exp(w(\lambda)\cdot\beta)\cdot\Psi(\Lambda_\beta x;\lambda)$ is continuous,
$\Psi(\cdot;\lambda)$ being holomorphic on $\mathcal{W}^{(n)}$ (Proposition~\ref{prop:kernels}\eqref{it:plate-vectors});
together with $\left\|\mathfrak{L}(\beta)\right\|=1$ this yields the strong continuity of $\mathfrak{L}$ on $\mathcal{H}^{\Psi}_{12}$.
A bounded operator on $\mathcal{H}^{\Psi}_{12}$ is determined by its values on the plate vectors, whose span is dense;
this gives the uniqueness.

\emph{The relations.}
The case $n=0$ of~\eqref{eq:rotation-boosts} is $\mathfrak{L}(\beta)\Omega_1=\Omega_1$.
Let $\kappa\in\overline{\mathcal{W}}$, $n\geq1$, $x\in\mathcal{W}^{(n)}$ and $\lambda$ a tuple of $n$ labels.
Since $\Lambda_\beta\left(x+(\kappa,\dots,\kappa)\right)=\Lambda_\beta x+(\Lambda_\beta\kappa,\dots,\Lambda_\beta\kappa)$ and $\Lambda_\beta\kappa\in\overline{\mathcal{W}}$,
\eqref{eq:rotation-boosts} and Proposition~\ref{prop:kernels}\eqref{it:plate-vectors} give
\begin{align*}
\mathfrak{L}(\beta)\mathsf{T}(\kappa)\Psi(x;\lambda)
&=\exp\left(w(\lambda)\cdot\beta\right)\cdot\Psi\left(\Lambda_\beta x+(\Lambda_\beta\kappa,\dots,\Lambda_\beta\kappa);\lambda\right)\\
&=\mathsf{T}(\Lambda_\beta\kappa)\mathfrak{L}(\beta)\Psi(x;\lambda);
\end{align*}
for $n=0$ both ends equal $\Omega_1$.
Let $\lambda_0=(\sigma,\chi)$ be a label and $(\xi,\xi')\in\mathcal{W}\times\mathcal{W}$.
The configuration that~\eqref{eq:rotation-plate-product} attaches to $\Lambda_\beta\xi$, $\Lambda_\beta\xi'$ and $\Lambda_\beta x$ is $\Lambda_\beta\tilde x$,
and the tuple $(\lambda_0,\lambda_1,\dots,\lambda_n)$ has the weight $w(\sigma)+w(\lambda)$;
so~\eqref{eq:rotation-plate-product} and~\eqref{eq:rotation-boosts} give
\begin{align*}
\mathfrak{L}(\beta)\mathcal{B}_{\lambda_0}(\xi,\xi')\Psi(x;\lambda)
&=\exp\left(\left(w(\sigma)+w(\lambda)\right)\cdot\beta\right)\cdot\Psi\left(\Lambda_\beta\tilde x;\lambda_0,\lambda_1,\dots,\lambda_n\right)\\
&=\exp\left(w(\sigma)\cdot\beta\right)\cdot\mathcal{B}_{\lambda_0}\left(\Lambda_\beta\xi,\Lambda_\beta\xi'\right)\mathfrak{L}(\beta)\Psi(x;\lambda).
\end{align*}
All operators in these two identities are bounded and leave $\mathcal{H}^{\Psi}_{12}$ invariant;
by linearity and continuity the identities hold on $\mathcal{H}^{\Psi}_{12}$, which proves the second and the third relation.
Finally, Lemma~\ref{lem:forward-tube}\eqref{it:boost-components} gives
\begin{align*}
-\left(\Lambda_\beta\kappa\right)^0\cdot\omega+i\,\left(\Lambda_\beta\kappa\right)^1\cdot p_2
&=-\kappa^0\cdot\left(\omega\cdot\cosh\beta+p_2\cdot\sinh\beta\right)\\
&\qquad+i\,\kappa^1\cdot\left(p_2\cdot\cosh\beta+\omega\cdot\sinh\beta\right),
\end{align*}
and by the definition of $\mathsf{T}$ (Lemma~\ref{lem:cone-semigroup}) the operator $\mathsf{T}(\Lambda_\beta\kappa)$
is the integral against $\mathfrak{E}_1$ of the exponential of this function of $(\omega,p_\perp)$.
The points $(t,0)$, $t\geq0$, and $(0,u)$, $u\in\R$, lie in $\overline{\mathcal{W}}$,
and $\mathsf{T}((t,0))=\exp(-t\cdot\mathsf{H}_1)$, $\mathsf{T}((0,u))=V_1(u\cdot e_2)=\exp(i\,u\cdot\mathsf{P}_2)$ by Lemma~\ref{lem:joint-spectrum}(b),~(c).
For these two choices of $\kappa$ the second relation becomes the last two identities of~\eqref{it:unitary-boosts}.
\end{proof}

The last two identities of Proposition~\ref{prop:rotations-boosts}\eqref{it:unitary-boosts} say that conjugation by $\mathfrak{L}(\beta)$
replaces the pair $(\mathsf{H}_1,\mathsf{P}_2)$ by $\left(\mathsf{H}_1\cdot\cosh\beta+\mathsf{P}_2\cdot\sinh\beta,\ \mathsf{P}_2\cdot\cosh\beta+\mathsf{H}_1\cdot\sinh\beta\right)$:
the pair transforms as an energy--momentum vector under a boost.
By~\eqref{eq:rotation-boosts} and the relation for the point plates,
the point plates $\mathcal{B}_{\lambda_0}$ of the labels $\lambda_0=(+,\chi)$, $(-,\chi)$, $(3,\chi)$,
which are formed with the components $\mathbb{B}_1+i\mathbb{B}_2$, $\mathbb{B}_1-i\mathbb{B}_2$, $\mathbb{B}_3$ of the field,
acquire under conjugation by $\mathfrak{L}(\beta)$ the factors $\exp(\beta)$, $\exp(-\beta)$, $1$,
as the corresponding components of a two-form do under a boost.

\emph{Plate vectors and boosts of the other planes.}
Let $(j,i)\neq(1,2)$ be an ordered pair of distinct indices,
let $\mu$ satisfy \eqref{cond:gaussian-domination}, \eqref{cond:closedness}, $(T_{123})$, $(R_j)$, $(RP_j)$, $(Q_{ji})$ and $(C_{ji})$,
and let $\mu':=\left(g_{ji}^{-1}\right)_*\mu$ and the unitary operator $J$ of $\mathcal{H}_1(\mu')$ onto $\mathcal{H}_j(\mu)$ be those of Lemma~\ref{lem:transport}\eqref{it:transport-pairs};
by that lemma, $\mu'$ satisfies the hypotheses of Proposition~\ref{prop:rotations-boosts}\eqref{it:unitary-boosts}.
The plate vectors of the plane $(j,i)$ are defined as the vectors $J\Psi(x;\lambda)$, where $\Psi(x;\lambda)$ ranges over the plate vectors of $\mu'$;
their closed span is $\mathcal{H}^{\Psi}_{ji}:=J\mathcal{H}^{\Psi}_{12}$, with $\mathcal{H}^{\Psi}_{12}$ formed for $\mu'$;
and the boosts of the plane $(j,i)$ are the unitary operators $\mathfrak{L}_{ji}(\beta):=J\mathfrak{L}(\beta)J^{-1}$ of $\mathcal{H}^{\Psi}_{ji}$, $\beta\in\R$,
where $\mathfrak{L}$ is the group of Proposition~\ref{prop:rotations-boosts}\eqref{it:unitary-boosts} for $\mu'$.
By Lemma~\ref{lem:transport}\eqref{it:transport-pairs}, $J\exp(-t\cdot\mathsf{H}_1)J^{-1}=\exp(-t\cdot\mathsf{H}_j)$ for $t\geq0$, $JV_1(u\cdot e_2)J^{-1}=V_j(u\cdot e_i)$ for real $u$,
and $J$ carries the joint spectral measure of $(\mathsf{H}_1,\mathsf{P}_2)$ to that of $(\mathsf{H}_j,\mathsf{P}_i)$.
Both joint spectral measures are carried by the cone $\left\{(\omega,p):\ \omega\geq|p|\right\}$, by $(C_{12})$ for $\mu'$ and $(C_{ji})$ for $\mu$;
so $J$ carries every Borel function of $(\mathsf{H}_1,\mathsf{P}_2)$ that is bounded on this cone to the same function of $(\mathsf{H}_j,\mathsf{P}_i)$.
Hence, on $\mathcal{H}^{\Psi}_{ji}$, for real $\beta$, $t\geq0$ and real $u$,
\begin{align*}
\mathfrak{L}_{ji}(\beta)\exp(-t\cdot\mathsf{H}_j)\,\mathfrak{L}_{ji}(\beta)^{-1}&=\exp\left(-t\cdot\left(\mathsf{H}_j\cdot\cosh\beta+\mathsf{P}_i\cdot\sinh\beta\right)\right),\\
\mathfrak{L}_{ji}(\beta)\exp(i\,u\cdot\mathsf{P}_i)\,\mathfrak{L}_{ji}(\beta)^{-1}&=\exp\left(i\,u\cdot\left(\mathsf{P}_i\cdot\cosh\beta+\mathsf{H}_j\cdot\sinh\beta\right)\right).
\end{align*}

\subsection{Diagonal mirrors and the cone}\label{sec:diagonal-mirrors}

\subsubsection{One diagonal mirror}\label{sec:one-diagonal-mirror}

\begin{proposition}[One diagonal mirror gives the cone]\label{prop:diagonal-mirror}
Let $j\neq i$ in $\{1,2,3\}$, and put $n_+:=(e_j+e_i)/\sqrt{2}$.
Let $\mu$ be a law on $\mathcal{S}'$ that is invariant under the translations of $\R^3$,
satisfies~\eqref{cond:reflection-invariance} and~\eqref{cond:reflection-positivity},
and is invariant under the reflection $r_{n_+,0}$ in the plane $\{x_j+x_i=0\}$ and reflection positive at this plane
(Definition~\ref{def:laws}).
Let $D>0$, let
\[
B:=\left\{x\in\R^3:\ 0<x_j<D,\ |x_i|<D\right\},
\]
and let $\Phi\in L^2(\mu)$ be measurable with respect to $\Sigma_B$
(so that $\Phi\in L^2(\Sigma_j^+,\mu)$, as $B\subset\{x_j>0\}$).
Then the spectral measure $\varrho_\psi$ of $\psi:=[\Phi]_j$ (Lemma~\ref{lem:joint-spectrum}) is carried by the set
\[
\left\{(\omega,p_\perp)\in[0,\infty)\times e_j^\perp:\ \omega\ge|p_i|\right\}.
\]
Consequently $\mu$ satisfies~\eqref{cond:cone}.
\end{proposition}

\begin{proof}
Write $\|\cdot\|_2$ for the norm of $L^2(\mu)$ and $\langle\cdot,\cdot\rangle_\mu$ for its scalar product, which is antilinear in the first entry.
For a Euclidean map $\mathfrak g$ the operator $\hat{\mathfrak g}$ on functions on $\mathcal{S}'$ is multiplicative and commutes with complex conjugation,
and it is unitary on $L^2(\mu)$ if $\mathfrak g_*\mu=\mu$ \cite[Section~8.1.2]{NoncompactQED}.
By the invariances of $\mu$, the translation operators $U(a)$, $a\in\R^3$, and the operators $\hat\theta_j$ and $\hat r_{n_+,0}$ are therefore unitary on $L^2(\mu)$;
moreover $\Theta_j\Phi'=\overline{\hat\theta_j\Phi'}$ for every $\Phi'\in L^2(\mu)$.
Let $\varrho$ be the image of $\varrho_\psi$ under the map $(\omega,p_\perp)\mapsto(\omega,p_i)$ from $[0,\infty)\times e_j^\perp$ to $[0,\infty)\times\R$;
it is a finite positive measure.
The set $[0,\infty)\times\R$ is identified with the closed half-plane $\{z\in\C:\ \Re z\ge0\}$ through $(\omega,p_i)\mapsto z:=\omega-ip_i$.
For $a,c\ge0$ put
\[
F(a,c):=\mu\left((\Theta_j\Phi)\cdot U\left(2D\cdot e_j+a\cdot(e_j+e_i)+c\cdot(e_j-e_i)\right)\Phi\right).
\]
The Cauchy--Schwarz inequality and the unitarity of $\hat\theta_j$ and of the translations give $|F(a,c)|\le\|\Phi\|_2^2$.

\emph{The first picture.}
The translation vector in $F(a,c)$ is $(2D+a+c)\cdot e_j+(a-c)\cdot e_i$:
its component along $e_j$ is $2D+a+c\ge0$, and $(a-c)\cdot e_i$ lies in $e_j^\perp$.
By \cite[Proposition~8.19\,(i)--(iii)]{NoncompactQED}
(where the linear operator $\hat\theta_j$ is written $\Theta_j$ and $\exp(-t\cdot\mathsf{H}_j)$ is written $T(t)$),
\[
\begin{aligned}
F(a,c)
&=\left\langle[\Phi]_j,\left[U\left((2D+a+c)\cdot e_j\right)U\left((a-c)\cdot e_i\right)\Phi\right]_j\right\rangle\\
&=\left\langle\psi,\exp\left(-(2D+a+c)\cdot\mathsf{H}_j\right)V_j\left((a-c)\cdot e_i\right)\psi\right\rangle.
\end{aligned}
\]
Lemma~\ref{lem:joint-spectrum}(b) writes the operator on the right as an integral over $\mathfrak{E}_j$;
since the exponent depends on $p_\perp$ only through $p_i$, one obtains, with $z=\omega-ip_i$,
\[
\begin{aligned}
F(a,c)
&=\int\mathrm{d}\varrho(\omega,p_i)\,\exp\left(-(2D+a+c)\cdot\omega+i(a-c)\cdot p_i\right)\\
&=\int\mathrm{d}\varrho(\omega,p_i)\,\exp(-2D\cdot\omega)\cdot\exp(-a\cdot z)\cdot\exp\left(-c\cdot\overline{z}\right).
\end{aligned}
\]

\emph{The second picture.}
Let $i'$ be the third index and $n_-:=(e_j-e_i)/\sqrt{2}$.
Let $g\in\mathrm{O}(3)$, the law $\mu'=(g^{-1})_*\mu$ and its reconstruction $(\mathcal{H}',[\,\cdot\,]',\mathsf{H}',V')$ along $e_j$
be as in step~\ref{it:cone-rotated-law} of the proof of Theorem~\ref{thm:cone}.
Thus $ge_j=n_+$, $ge_i=n_-$, $ge_{i'}=e_{i'}$ and $r_{n_+,0}=g\theta_jg^{-1}$;
the operator $\hat g^{-1}$ is unitary from $L^2(\mu)$ onto $L^2(\mu')$ and maps $L^2(\Sigma_{\{x_j+x_i>0\}},\mu)$ onto $L^2(\Sigma_j^+,\mu')$;
and $\mu'$ is invariant under the translations of $\R^3$ and satisfies~\eqref{cond:reflection-invariance} and~\eqref{cond:reflection-positivity}.
The scalar product and the norm of $\mathcal{H}'$ are written $\langle\cdot,\cdot\rangle'$ and $\|\cdot\|'$.
Put $y:=\frac{D}{2}\cdot(e_j+e_i)$ and
\[
\Phi_1:=\hat r_{n_+,0}U(-y)\hat\theta_j\Phi,\qquad\Phi_2:=U(2D\cdot e_j-y)\Phi.
\]
Both observables belong to $L^2(\Sigma_{\{x_j+x_i>0\}},\mu)$.
The operator $\hat{\mathfrak g}$ of a Euclidean map $\mathfrak g$ carries, for every open set $U\subset\R^3$,
the $\Sigma_U$-measurable functions onto the $\Sigma_{\mathfrak gU}$-measurable functions \cite[Section~8.1.2]{NoncompactQED},
and $\Sigma_U\subset\Sigma_{U'}$ whenever $U\subset U'$ are open sets.
Accordingly $\hat\theta_j\Phi$ is measurable with respect to $\Sigma_{\theta_jB}$, where $\theta_jB=\{x\in\R^3:\ -D<x_j<0,\ |x_i|<D\}$,
and $x_j+x_i<D$ on $\theta_jB$.
The linear function $x\mapsto x_j+x_i$ takes the value $D$ at $y$, so the translation $\tau_{-y}$ lowers it by $D$;
thus $U(-y)\hat\theta_j\Phi$ is measurable with respect to $\Sigma_{\theta_jB-y}$, and $\theta_jB-y\subset\{x_j+x_i<0\}$.
The reflection $r_{n_+,0}$ carries $\{x_j+x_i<0\}$ onto $\{x_j+x_i>0\}$,
so that $\Phi_1$ is measurable with respect to $\Sigma_{\{x_j+x_i>0\}}$.
On $B$ one has $x_j+x_i>-D$, and the translation $\tau_{2D\cdot e_j-y}$ raises $x_j+x_i$ by $D$;
so $\Phi_2$ is measurable with respect to $\Sigma_{B+2D\cdot e_j-y}$, and $B+2D\cdot e_j-y\subset\{x_j+x_i>0\}$.
As the operators $\hat\theta_j$, $U(-y)$, $\hat r_{n_+,0}$ and $U(2D\cdot e_j-y)$ are unitary, $\|\Phi_1\|_2=\|\Phi_2\|_2=\|\Phi\|_2$.

The invariance of $\mu$ under $\tau_{-y}$, applied to the product that defines $F(a,c)$, gives
\[
F(a,c)=\mu\left(\left(U(-y)\Theta_j\Phi\right)\cdot U\left(a\cdot(e_j+e_i)+c\cdot(e_j-e_i)\right)U(2D\cdot e_j-y)\Phi\right).
\]
As $\hat r_{n_+,0}^2=1$, the first factor equals $\overline{U(-y)\hat\theta_j\Phi}=\overline{\hat r_{n_+,0}\Phi_1}$;
moreover $a\cdot(e_j+e_i)+c\cdot(e_j-e_i)=\sqrt{2}\cdot a\cdot n_++\sqrt{2}\cdot c\cdot n_-$.
Formula~\eqref{eq:continuum-rotated-picture}, applied with $\sqrt{2}\cdot a\ge0$ and $\sqrt{2}\cdot c$ in the places of $a$ and $b$, therefore yields
\[
\begin{aligned}
F(a,c)
&=\left\langle\hat r_{n_+,0}\Phi_1,U\left(\sqrt{2}\cdot a\cdot n_++\sqrt{2}\cdot c\cdot n_-\right)\Phi_2\right\rangle_\mu\\
&=\left\langle\psi_1,\exp\left(-\sqrt{2}\cdot a\cdot\mathsf{H}'\right)V'\left(\sqrt{2}\cdot c\cdot e_i\right)\psi_2\right\rangle',
\end{aligned}
\]
where
\[
\psi_1:=\left[\hat g^{-1}\Phi_1\right]',\qquad\psi_2:=\left[\hat g^{-1}\Phi_2\right]'.
\]
The map $[\,\cdot\,]'$ is a contraction \cite[Proposition~8.19\,(i)]{NoncompactQED}, applied to $\mu'$, and $\hat g^{-1}$ is unitary;
so $\|\psi_1\|'\le\|\Phi\|_2$ and $\|\psi_2\|'\le\|\Phi\|_2$.

Lemma~\ref{lem:joint-spectrum}, applied to $\mu'$, provides the joint spectral measure $\mathfrak{E}'$ of $\mu'$ in the direction $e_j$:
a projection-valued measure on $[0,\infty)\times e_j^\perp$ with
\[
\exp\left(-\sqrt{2}\cdot a\cdot\mathsf{H}'\right)V'\left(\sqrt{2}\cdot c\cdot e_i\right)
=\int\mathrm{d}\mathfrak{E}'(\omega,p_\perp)\,\exp\left(-\sqrt{2}\cdot a\cdot\omega+i\sqrt{2}\cdot c\cdot p_i\right).
\]
Let $\varrho'$ be the image of the complex measure $\langle\psi_1,\mathfrak{E}'(\cdot)\psi_2\rangle'$
under the map $(\omega,p_\perp)\mapsto(\sqrt{2}\cdot\omega,\sqrt{2}\cdot p_i)$ from $[0,\infty)\times e_j^\perp$ to $[0,\infty)\times\R$,
and denote the variables of $[0,\infty)\times\R$ in this picture by $(\omega',p')$.
For every finite partition $(A_k)_k$ of $[0,\infty)\times e_j^\perp$ into Borel sets,
the projections $\mathfrak{E}'(A_k)$ are mutually orthogonal with sum $1$,
and the Cauchy--Schwarz inequality, first in $\mathcal{H}'$ and then for finite sums, gives
\[
\sum_k\left|\left\langle\psi_1,\mathfrak{E}'(A_k)\psi_2\right\rangle'\right|
\le\sum_k\left\|\mathfrak{E}'(A_k)\psi_1\right\|'\cdot\left\|\mathfrak{E}'(A_k)\psi_2\right\|'
\le\|\psi_1\|'\cdot\|\psi_2\|'.
\]
An image measure has at most the total variation of the original one.
Hence $\varrho'$ is a complex measure on $[0,\infty)\times\R$ of total variation at most $\|\Phi\|_2^2$, and
\[
F(a,c)=\int\mathrm{d}\varrho'(\omega',p')\,\exp\left(-a\cdot\omega'+ic\cdot p'\right).
\]

\emph{An inequality between Laplace transforms.}
For $h\in L^1(0,\infty)$ let
\[
\tilde h(z):=\int_0^\infty\mathrm{d}a\,h(a)\cdot\exp(-a\cdot z)\qquad(\Re z\ge0).
\]
Multiply $F(a,c)$ by $h(a)\cdot\overline{h(c)}$ and integrate over $a,c\ge0$.
In each of the two representations of $F$ the exponentials have modulus at most $1$, because $\omega\ge0$ and $\omega'\ge0$;
so the integrand is bounded in modulus by $|h(a)|\cdot|h(c)|$,
an integrable function of $(a,c)$, while $\varrho$ and $\varrho'$ have finite total variation;
by Fubini's theorem the integrations over $a$ and $c$ may be carried out first.
Since
\[
\int_0^\infty\mathrm{d}c\,\overline{h(c)}\cdot\exp\left(-c\cdot\overline{z}\right)=\overline{\tilde h(z)},
\qquad
\int_0^\infty\mathrm{d}c\,\overline{h(c)}\cdot\exp(ic\cdot p')=\overline{\tilde h(ip')},
\]
the two pictures give
\begin{equation}\label{eq:rotation-two-pictures}
\begin{split}
\int\mathrm{d}\varrho(\omega,p_i)\,\exp(-2D\cdot\omega)\cdot\left|\tilde h(\omega-ip_i)\right|^2
&=\int\mathrm{d}\varrho'(\omega',p')\,\tilde h(\omega')\cdot\overline{\tilde h(ip')}\\
&\le\|\Phi\|_2^2\cdot\sup_{\omega'\ge0}\left|\tilde h(\omega')\right|\cdot\sup_{p'\in\R}\left|\tilde h(ip')\right|.
\end{split}
\end{equation}

\emph{Test functions.}
For $z\in\C\setminus\{0\}$ let $\arg z\in(-\pi,\pi]$ be the principal argument and $\log z:=\log|z|+i\arg z$ the principal logarithm;
all complex powers below are principal, that is, a power of $z$ is $\exp$ of the exponent times $\log z$.
Fix a sign $s\in\{1,-1\}$.
For every integer $n\ge1$ define a function $\tilde h_n$ on the closed half-plane $\{\Re z\ge0\}$ by
\[
\tilde h_n(z):=\left(\exp(-1)\cdot z^{-\frac{4is}{\pi}}\right)^n\cdot z\cdot(1+z)^{-3}\quad(z\neq0),
\qquad
\tilde h_n(0):=0.
\]
For $z\neq0$ one has $\left|z^{-\frac{4is}{\pi}}\right|=\exp\left(\frac{4s}{\pi}\cdot\arg z\right)$.
If $\Re z\ge0$, then $|1+z|^2=1+2\Re z+|z|^2$ is at least $1$ and at least $|z|^2$, so that $|z|\cdot|1+z|^{-3}\le1$.
Hence
\[
\left|\tilde h_n(z)\right|\le\exp\left(-n+\frac{4s}{\pi}\cdot n\cdot\arg z\right)\qquad(\Re z\ge0,\ z\neq0).
\]
On the positive axis $\arg z=0$, and $|\tilde h_n|\le\exp(-n)$ there.
On the imaginary axis $|\arg z|=\pi/2$ away from the origin, and $|\tilde h_n|\le\exp(-n+2n)=\exp(n)$ there.
So the product of $\sup_{\omega'\ge0}|\tilde h_n(\omega')|$ and $\sup_{p'\in\R}|\tilde h_n(ip')|$ is at most $1$.

\emph{Each $\tilde h_n$ is the Laplace transform of an integrable function.}
Put $\nu:=1-\frac{4is}{\pi}\cdot n$;
then $\Re\nu=1$ and $\Im\nu\neq0$.
For $z\neq0$ the principal powers satisfy $\left(z^{-\frac{4is}{\pi}}\right)^n\cdot z=\exp(\nu\cdot\log z)=z^\nu$, so that
\[
\tilde h_n(z)=\exp(-n)\cdot z^\nu\cdot(1+z)^{-3}\qquad(\Re z\ge0,\ z\neq0).
\]
Let $\Re z\ge0$ and $z\neq0$.
The arguments of $z$ and of $1+z$ lie in $[-\pi/2,\pi/2]$, they do not have opposite signs, and $|\arg(1+z)|\le|\arg z|$.
Therefore $\arg z-\arg(1+z)$ lies in $[-\pi/2,\pi/2]$ and is the principal argument of $z/(1+z)$;
the principal logarithms satisfy $\log\frac{z}{1+z}=\log z-\log(1+z)$, and with $u:=(1+z)^{-1}$
\[
z^\nu\cdot(1+z)^{-3}
=\left(\frac{z}{1+z}\right)^\nu\cdot(1+z)^{-(3-\nu)}
=(1-u)^\nu\cdot(1+z)^{-(3-\nu)}.
\]

The second factor is a Laplace transform:
\[
\int_0^\infty\mathrm{d}a\,a^{2-\nu}\cdot\exp(-a)\cdot\exp(-a\cdot z)=\Gamma(3-\nu)\cdot(1+z)^{-(3-\nu)}\qquad(\Re z\ge0).
\]
Indeed, both sides are holomorphic functions of $z$ on $\{\Re z>-1\}$,
and for real $z>-1$ they agree by the substitution of $a/(1+z)$ for $a$.
Here $\left|a^{2-\nu}\right|=a$, so $a\mapsto a^{2-\nu}\cdot\exp(-a)/\Gamma(3-\nu)$ is integrable on $(0,\infty)$;
$\Gamma(3-\nu)\neq0$ because the Gamma function has no zeros.

The first factor is handled by the binomial series.
As $\nu$ is not a non-negative integer,
\[
\binom{\nu}{k}=\frac{\nu\cdot(\nu-1)\cdots(\nu-k+1)}{k!}=(-1)^k\cdot\frac{\Gamma(k-\nu)}{\Gamma(-\nu)\cdot\Gamma(k+1)}\qquad(k\ge0),
\]
and Stirling's formula gives $\left|\binom{\nu}{k}\right|=O\left(k^{-1-\Re\nu}\right)=O\left(k^{-2}\right)$ as $k\to\infty$;
thus $\sum_{k\ge0}\left|\binom{\nu}{k}\right|<\infty$.
For $|u|<1$ one has $(1-u)^\nu=\sum_{k\ge0}\binom{\nu}{k}\cdot(-u)^k$.
The series converges uniformly on the closed disc $\{|u|\le1\}$, so that its sum is continuous there.
The function $(1-u)^\nu$ is continuous on $\{|u|\le1,\ u\neq1\}$, where $\Re(1-u)\ge0$ and $1-u\neq0$,
and it tends to $0$ as $u\to1$, because $\left|(1-u)^\nu\right|\le|1-u|^{\Re\nu}\cdot\exp\left(\frac{\pi}{2}\cdot|\Im\nu|\right)$ there.
With the value $0$ at $u=1$ the identity therefore holds on the closed disc.
For $\Re z\ge0$ the point $u=(1+z)^{-1}$ lies in this disc, and $u=1$ only for $z=0$.

For $k\ge1$, the same Gamma integral with the exponent $k$ in the place of $3-\nu$ shows that $(1+z)^{-k}$ is the Laplace transform of the probability density $a^{k-1}\cdot\exp(-a)/(k-1)!$ on $(0,\infty)$;
the constant $1$ is the Laplace transform of the unit point mass at $0$.
Consider the finite complex measure on $[0,\infty)$ that gives the mass $1$ to $\{0\}$ and has the density
\[
\sum_{k\ge1}(-1)^k\cdot\binom{\nu}{k}\cdot\frac{a^{k-1}\cdot\exp(-a)}{(k-1)!}
\]
on $(0,\infty)$, a series that converges in $L^1(0,\infty)$ since $\sum_k\left|\binom{\nu}{k}\right|<\infty$.
Its Laplace transform at a point $z$ with $\Re z\ge0$ is
\[
\sum_{k\ge0}\binom{\nu}{k}\cdot\left(-(1+z)^{-1}\right)^k=\left(1-(1+z)^{-1}\right)^\nu,
\]
the exchange of sum and integral being justified by the absolute convergence.
Combining the two factors, $\exp(n)\cdot\tilde h_n(z)$ is, for every $z$ with $\Re z\ge0$,
the product of the Laplace transform of this complex measure and the Laplace transform of the integrable function $a\mapsto a^{2-\nu}\cdot\exp(-a)/\Gamma(3-\nu)$;
at $z=0$ both sides vanish.
The product is the Laplace transform of the convolution of the measure with the function,
and this convolution is integrable on $(0,\infty)$, by Fubini's theorem.
So $\tilde h_n$ is the Laplace transform of a function $h_n\in L^1(0,\infty)$,
and~\eqref{eq:rotation-two-pictures} applies to $h=h_n$.

\emph{The cone for $\varrho$.}
By~\eqref{eq:rotation-two-pictures} and the bound on the product of the two suprema,
\[
\int\mathrm{d}\varrho(\omega,p_i)\,\exp(-2D\cdot\omega)\cdot\left|\tilde h_n(\omega-ip_i)\right|^2\le\|\Phi\|_2^2
\]
for every $n\ge1$ and for both signs $s$.
Let $\delta>0$, and let $S_\delta$ be the sector of all $z\neq0$ with $\Re z\ge0$ and $s\cdot\arg z\ge\pi/4+\delta$;
it depends on $s$ and is a Borel subset of the half-plane $[0,\infty)\times\R$.
For $z\in S_\delta$,
\[
\left|\exp(-1)\cdot z^{-\frac{4is}{\pi}}\right|
=\exp\left(-1+\frac{4s}{\pi}\cdot\arg z\right)
\ge\exp\left(-1+\frac{4}{\pi}\cdot\left(\frac{\pi}{4}+\delta\right)\right)
=\exp\left(\frac{4\delta}{\pi}\right),
\]
so that $\left|\tilde h_n(z)\right|^2\ge\exp\left(\frac{8\delta}{\pi}\cdot n\right)\cdot|z|^2\cdot|1+z|^{-6}$, and
\[
\exp\left(\frac{8\delta}{\pi}\cdot n\right)\cdot\int_{S_\delta}\mathrm{d}\varrho(\omega,p_i)\,\exp(-2D\cdot\omega)\cdot|z|^2\cdot|1+z|^{-6}\le\|\Phi\|_2^2\qquad(n\ge1).
\]
Letting $n\to\infty$ shows that the integral vanishes.
Its integrand is strictly positive on $S_\delta$, hence $\varrho(S_\delta)=0$.
The union of the sectors $S_\delta$ over both signs and over all rational $\delta>0$ is the set of all $z\neq0$ with $\Re z\ge0$ and $|\arg z|>\pi/4$,
and $\varrho$ gives it no mass.
For $z=\omega-ip_i$ with $\omega\ge0$, the condition ``$z=0$ or $|\arg z|\le\pi/4$'' is equivalent to $|p_i|\le\omega$.
Therefore $\varrho(\{\omega<|p_i|\})=0$, and by the definition of $\varrho$ the measure $\varrho_\psi$ is carried by $\{\omega\ge|p_i|\}$.

\emph{The condition~\eqref{cond:cone}.}
Let $A:=\{(\omega,p_\perp)\in[0,\infty)\times e_j^\perp:\ \omega<|p_i|\}$.
Let $\Phi$ now be a finite linear combination of functions $\exp(iF(f))$ with $f\in C^\infty_{\mathrm c}(\{x_j>0\};\Lambda^2)$.
The finitely many supports are compact subsets of $\{x_j>0\}$;
they lie in $B$ as soon as $D$ is large enough, and then $\Phi$ is measurable with respect to $\Sigma_B$.
By the first assertion, $\left\|\mathfrak{E}_j(A)[\Phi]_j\right\|^2=\varrho_{[\Phi]_j}(A)=0$.
These functions $\Phi$ are dense in $L^2(\Sigma_j^+,\mu)$ by \cite[Lemma~8.12\,(i)]{NoncompactQED}
(where the field $F(f)$ is written $X(f)$),
and $[\,\cdot\,]_j$ is a contraction with dense range \cite[Proposition~8.19\,(i)]{NoncompactQED};
so their classes are dense in $\mathcal{H}_j$.
The bounded operator $\mathfrak{E}_j(A)$ vanishes on this dense set, hence $\mathfrak{E}_j(A)=0$,
which is condition~\eqref{cond:cone} of Definition~\ref{def:closedness-cones}.
\end{proof}

\subsubsection{Equivalent forms of the hypothesis}\label{sec:equivalent-forms}

\begin{theorem}[Equivalent forms of the cone]\label{thm:cone-equivalences}
Let $\mu$ be a law on $\mathcal{S}'$ that satisfies \textup{(H)} and \textup{(dF)}.
Then the following seven conditions are equivalent.
\begin{enumerate}
\item\label{it:equiv-one-cone}
For some ordered pair $j\neq i$ in $\{1,2,3\}$ the cone condition $(C_{ji})$ of Definition~\ref{def:closedness-cones} holds, which is the inequality $\mathsf{H}_j\ge|\mathsf{P}_i|$ with the constant one.
\item\label{it:equiv-all-cones}
The cone condition $(C_{ji})$ holds for all six ordered pairs $j\neq i$.
\item\label{it:equiv-five-point}
There are an ordered pair $j\neq i$, a set $I\subset(0,\infty)$ with $\inf I=0$,
and a set of elements of $L^2(\Sigma_j^+,\mu)$ whose classes are dense in $\mathcal{H}_j$,
such that for every $t\in I$ and every $\Phi$ in this set the function
\[
u(v):=\mu\left((\Theta_j\Phi)\cdot U(v)\Phi\right)\qquad(v\in\R^3)
\]
satisfies the five-point inequality
\begin{equation}\label{eq:rotation-five-point}
u(0)-4u(te_j)+u(2te_j)+u(te_j+te_i)+u(te_j-te_i)\ge0.
\end{equation}
\item\label{it:equiv-orthogonal}
The law $\mu$ is invariant under $\mathrm{O}(3)$.
\item\label{it:equiv-diagonal-plane}
The law $\mu$ is reflection positive at one of the six diagonal planes $\{x_i=x_l\}$ and $\{x_i=-x_l\}$, $1\le i<l\le3$.
\item\label{it:equiv-nine-families}
The law $\mu$ belongs to $\mathcal{C}_9$.
\item\label{it:equiv-all-planes}
The law $\mu$ is reflection positive at every plane of $\R^3$.
\end{enumerate}
Furthermore, the following four conditions on the one-field spaces of $\mu$ (Definition~\ref{def:closedness-cones})
and on its two-point function $S_2$ (Lemma~\ref{lem:domination-consequences}\eqref{it:moment-forms}) are equivalent.
\begin{enumerate}
\setcounter{enumi}{7}
\item\label{it:equiv-one-field-cone}
For some ordered pair $j\neq i$ the one-field cone condition $(C^1_{ji})$ holds.
\item\label{it:equiv-two-point-orthogonal}
For every $g\in\mathrm{O}(3)$ and all $b,b'\in L^2(\R^3;\C^3)$ one has $S_2(g\cdot b,g\cdot b')=S_2(b,b')$.
\item\label{it:equiv-two-point-diagonal}
For one of the six diagonal planes of \eqref{it:equiv-diagonal-plane} and a unit normal $n$ of it,
\[
S_2(r_{n,0}\cdot b,b)\ge0
\]
for every $b\in L^2(\R^3;\R^3)$ that vanishes almost everywhere on $\{n\cdot x<0\}$.
\item\label{it:equiv-two-point-kallen-lehmann}
The two-point function has the K\"all\'en--Lehmann form: there is a finite positive Borel measure $\mathfrak{K}$ on $[0,\infty)$ with $\mathfrak{K}([0,\infty))\le1$ such that
\begin{equation}\label{eq:rotation-kallen-lehmann}
S_2(\overline{b},b')=(2\pi)^{-3}\int\mathrm{d}p\,\hat b(p)^{\dagger}P^{\perp}(p)\,\hat b'(p)\cdot\int_{[0,\infty)}\mathrm{d}\mathfrak{K}(m)\,\frac{|p|^2}{|p|^2+m^2}
\end{equation}
for all $b,b'\in L^2(\R^3;\C^3)$;
here $\hat b(p)=\int\mathrm{d}x\,\exp(-ip\cdot x)\,b(x)$, $\hat b(p)^{\dagger}$ is its conjugate transpose, and $P^{\perp}(p)=1-pp^{T}/|p|^2$ for $p\neq0$.
\end{enumerate}
\end{theorem}

\begin{proof}
The hypothesis \textup{(H)} contains \textup{(G)}, $(T_{123})$ and $(RP_j)$ for $j=1,2,3$;
it also gives $(R_j)$, since the reflection $\theta_j$ belongs to $\Ocub$.
For each $j$ let $(\mathcal{H}_j,[\,\cdot\,]_j,\mathsf{H}_j,V_j,\Omega_j)$ be the reconstruction of $\mu$ in the direction $e_j$,
and let $\mathfrak{E}_j$ be its joint spectral measure on $[0,\infty)\times e_j^{\perp}$, written in the variables $(\omega,p_\perp)$, $p_\perp=(p_{i'})_{i'\neq j}$ (Lemma~\ref{lem:joint-spectrum});
for $\psi\in\mathcal{H}_j$ the spectral measure $\varrho_\psi:=\langle\psi,\mathfrak{E}_j(\cdot)\psi\rangle$ is a positive measure of mass $\|\psi\|^2$.

\emph{Cone, orthogonal group and planes of reflection.}
Theorem~\ref{thm:orthogonal-group} shows that \eqref{it:equiv-one-cone} implies \eqref{it:equiv-orthogonal} and \eqref{it:equiv-all-cones}.
Conversely, \eqref{it:equiv-all-cones} contains \eqref{it:equiv-one-cone}.

Under \eqref{it:equiv-orthogonal} the law $\mu$ is invariant under the translations and under $\mathrm{O}(3)$, hence under $\R^3\rtimes\mathrm{O}(3)$, and it satisfies $(RP_1)$.
By Lemma~\ref{lem:transport}\eqref{it:transport-all-planes} it is then invariant under the reflection in every plane of $\R^3$ and reflection positive at every such plane;
in particular \eqref{it:equiv-all-planes} holds.

Assume \eqref{it:equiv-all-planes}.
A plane of the nine families has the form $\{n\cdot x=c\}$ with $c\in\R$ and $n$ one of the vectors $e_i$, $(e_i+e_l)/\sqrt2$, $(e_i-e_l)/\sqrt2$;
its reflection is $r_{n,c}=\tau_{2cn}\,r_{n,0}$, because $r_{n,c}x=x-2(n\cdot x-c)\,n$.
For these $n$ the linear map $r_{n,0}$ is a signed permutation matrix and so belongs to $\Ocub$.
By \textup{(H)} the law $\mu$ is invariant under the translations and under $\Ocub$, hence under each $r_{n,c}$ of this kind.
Together with the reflection positivity given by \eqref{it:equiv-all-planes}, this means $\mu\in\mathcal{C}_9$, which is \eqref{it:equiv-nine-families}.

The six diagonal planes $\{x_i-x_l=0\}$ and $\{x_i+x_l=0\}$ are planes of the nine families of Definition~\ref{def:laws};
therefore \eqref{it:equiv-nine-families} implies \eqref{it:equiv-diagonal-plane}.

Now let $\mu$ be reflection positive at a diagonal plane.
Write this plane as $\{n\cdot x=0\}$, where the unit normal $n$ is chosen so that the reflection positivity refers to the open half space $\{n\cdot x>0\}$;
then $n$ is one of the twelve vectors $(\pm e_i\pm e_l)/\sqrt2$, $i<l$.
The signed permutation matrices act transitively on these twelve vectors, and one may choose $g\in\Ocub$ with $gn=(e_1+e_2)/\sqrt2$.
By \textup{(H)}, $(g^{-1})_*\mu=\mu$.
Lemma~\ref{lem:transport}\eqref{it:transport-properties}, applied with this $g$, shows that $\mu$ is reflection positive at the plane $\{(gn)\cdot x=0\}=\{x_1+x_2=0\}$, with respect to the half space $\{x_1+x_2>0\}$.
The reflection in this plane is the map $(x_1,x_2,x_3)\mapsto(-x_2,-x_1,x_3)$, an element of $\Ocub$, so $\mu$ is invariant under it.
Hence $\mu$ is invariant under the translations, satisfies $(R_1)$ and $(RP_1)$, and is invariant and reflection positive at $\{x_1+x_2=0\}$:
these are the hypotheses of Theorem~\ref{thm:cone} for the pair $(j,i)=(1,2)$.
Theorem~\ref{thm:cone}(c) gives $(C_{12})$, and Proposition~\ref{prop:diagonal-mirror} gives it as well.
This is \eqref{it:equiv-one-cone}.

\emph{The five-point inequalities.}
Fix an ordered pair $j\neq i$.
It suffices to show that $(C_{ji})$ holds if and only if \eqref{it:equiv-five-point} holds with this pair.
By \cite[Proposition~8.19\,(ii), (iii)]{NoncompactQED}, where the semigroup $\exp(-s\cdot\mathsf{H}_j)$ is written $T(s)$,
\[
\exp(-s\cdot\mathsf{H}_j)[\Phi']_j=[U(se_j)\Phi']_j,\qquad V_j(y)[\Phi']_j=[U(y)\Phi']_j
\]
for $\Phi'\in L^2(\Sigma_j^+,\mu)$, $s\ge0$ and $y\in e_j^{\perp}$.
Let $\Phi\in L^2(\Sigma_j^+,\mu)$, and let $u$ be the function attached to $\Phi$ in \eqref{it:equiv-five-point}.
Since $\langle[\Phi]_j,[\Phi']_j\rangle=\mu\left((\Theta_j\Phi)\cdot\Phi'\right)$ and $U(se_j)U(y)=U(se_j+y)$, one obtains
\[
u(se_j+y)=\left\langle[\Phi]_j,\exp(-s\cdot\mathsf{H}_j)V_j(y)[\Phi]_j\right\rangle\qquad\left(s\ge0,\ y\in e_j^{\perp}\right).
\]
Let $t>0$.
As $V_j(-te_i)=V_j(te_i)^{*}$, the left side of \eqref{eq:rotation-five-point} equals $\langle[\Phi]_j,Q_t[\Phi]_j\rangle$, where
\begin{equation}\label{eq:rotation-five-point-operator}
\begin{aligned}
Q_t&:=1-4\exp(-t\cdot\mathsf{H}_j)+\exp(-2t\cdot\mathsf{H}_j)+\exp(-t\cdot\mathsf{H}_j)\left(V_j(te_i)+V_j(te_i)^{*}\right)\\
&=\int\mathrm{d}\mathfrak{E}_j(\omega,p_\perp)\,2\exp(-t\omega)\cdot\left(\cosh(t\omega)+\cos(tp_i)-2\right).
\end{aligned}
\end{equation}
For the second line, Lemma~\ref{lem:joint-spectrum}(b) writes the first line as the integral of $1-4\exp(-t\omega)+\exp(-2t\omega)+2\exp(-t\omega)\cdot\cos(tp_i)$ against $\mathfrak{E}_j$,
and $2\exp(-t\omega)\cdot\cosh(t\omega)=1+\exp(-2t\omega)$.
So $Q_t$ is a bounded self-adjoint operator on $\mathcal{H}_j$, and its quadratic form $\psi\mapsto\langle\psi,Q_t\psi\rangle$ is continuous.
This form is non-negative on a dense set of vectors if and only if $Q_t\ge0$.
By \eqref{eq:rotation-five-point-operator} and Lemma~\ref{lem:joint-spectrum}(d), applied to the bounded real integrand, $Q_t\ge0$ holds if and only if $\mathfrak{E}_j$ is carried by
\[
A_t:=\left\{(\omega,p_\perp)\in[0,\infty)\times e_j^{\perp}:\ \cosh(t\omega)+\cos(tp_i)\ge2\right\}.
\]

Suppose that $(C_{ji})$ holds, so that $\mathfrak{E}_j$ is carried by $\{\omega\ge|p_i|\}$.
For real $s$ one has $\cosh s+\cos s=2\sum_{k\ge0}s^{4k}/(4k)!\ge2$.
Consequently $\omega\ge|p_i|$ implies $\cosh(t\omega)\ge\cosh(tp_i)\ge2-\cos(tp_i)$, and $\mathfrak{E}_j$ is carried by $A_t$ for every $t>0$.
Thus $Q_t\ge0$, and \eqref{eq:rotation-five-point} holds for every $t>0$ and every $\Phi\in L^2(\Sigma_j^+,\mu)$.
The canonical map $[\,\cdot\,]_j$ has dense range \cite[Proposition~8.19\,(i)]{NoncompactQED};
so \eqref{it:equiv-five-point} holds with $I=(0,\infty)$ and the whole space $L^2(\Sigma_j^+,\mu)$.

Conversely, let $I$ and a set of elements of $L^2(\Sigma_j^+,\mu)$ be as in \eqref{it:equiv-five-point} for the pair $(j,i)$, and let $(t_k)_{k\ge1}$ be a sequence in $I$ with $t_k\to0$.
For each $k$ the quadratic form of $Q_{t_k}$ is non-negative at the classes of the given elements, and these classes are dense in $\mathcal{H}_j$.
Hence $\mathfrak{E}_j$ is carried by every $A_{t_k}$, and therefore by the intersection $\bigcap_{k\ge1}A_{t_k}$.
On $A_t$ one has $\cosh(t\omega)-1\ge1-\cos(tp_i)$, that is, $\sinh(t\omega/2)\ge|\sin(tp_i/2)|$, as $\omega\ge0$.
At a point of the intersection, multiply this inequality for $t=t_k$ by $2/t_k$ and let $k\to\infty$; the limit is $\omega\ge|p_i|$.
Therefore $\mathfrak{E}_j(\{\omega<|p_i|\})=0$, and $(C_{ji})$ holds.

\emph{The two-point function.}
Let $\mathfrak{g}$ be a Euclidean map with $\mathfrak{g}_*\mu=\mu$.
Then $\mu(\hat{\mathfrak{g}}\Phi)=\mu(\Phi)$ for $\Phi\in L^1(\mu)$, the operator $\hat{\mathfrak{g}}$ is multiplicative,
and $\hat{\mathfrak{g}}\,\mathbb{B}(b)=\mathbb{B}(\mathfrak{g}\cdot b)$ in $L^2(\mu)$ for every $b\in L^2(\R^3;\C^3)$ (Lemma~\ref{lem:domination-consequences}\eqref{it:field-extension}).
Consequently
\[
S_2(\mathfrak{g}\cdot b,\mathfrak{g}\cdot b')=\mu\left(\hat{\mathfrak{g}}\left(\mathbb{B}(b)\cdot\mathbb{B}(b')\right)\right)=S_2(b,b')\qquad\left(b,b'\in L^2(\R^3;\C^3)\right).
\]
By \textup{(H)} this applies to all translations and to all elements of $\Ocub$.

Assume \eqref{it:equiv-one-field-cone}, for the pair $(j,i)$.
For an ordered pair $j'\neq i'$ let $g_{j'i'}$ be the permutation matrix with $g_{j'i'}e_1=e_{j'}$ and $g_{j'i'}e_2=e_{i'}$.
It belongs to $\Ocub$, so that $(g_{j'i'}^{-1})_*\mu=\mu$;
Lemma~\ref{lem:transport}\eqref{it:transport-pairs}, applied with $g_{j'i'}$ (compare \eqref{eq:rotation-pair-transport}), therefore shows that $\mu$ satisfies $(C^1_{j'i'})$ if and only if it satisfies $(C^1_{12})$.
Taking first $(j',i')=(j,i)$, one obtains $(C^1_{12})$, and then $(C^1_{j'i'})$ for all six ordered pairs.
The quarter turn $R^{12}_{\pi/2}$ lies in $\Ocub$, so $S_2$ satisfies the identity \eqref{eq:rotation-two-point-quarter-turn} of Theorem~\ref{thm:two-point}.
The law $\mu$ satisfies \textup{(G)}, \textup{(dF)}, $(T_{123})$, $(R_1)$, $(RP_1)$ and $(C^1_{12})$, and Theorem~\ref{thm:two-point}\eqref{it:two-point-invariance} gives
\[
S_2(R^{12}_t\cdot b,R^{12}_t\cdot b')=S_2(b,b')\qquad\left(t\in\R,\ b,b'\in L^2(\R^3;\C^3)\right).
\]
For an ordered pair $j'\neq i'$ the matrix $g_{j'i'}R^{12}_tg_{j'i'}^{-1}$ fixes the third unit vector and maps $e_{j'}$ to $\cos t\,e_{j'}+\sin t\,e_{i'}$ and $e_{i'}$ to $-\sin t\,e_{j'}+\cos t\,e_{i'}$;
it is $R^{j'i'}_t$.
Since $(gg')\cdot b=g\cdot (g'\cdot b)$ for $g,g'\in\mathrm{O}(3)$, the elements of $\mathrm{O}(3)$ that leave $S_2$ invariant form a group.
This group contains $g_{j'i'}$ and $R^{12}_t$, hence $R^{j'i'}_t$;
so it contains the rotations of the three coordinate planes, and also the matrix $-1$, which lies in $\Ocub$.
The rotations of the planes $(1,2)$ and $(2,3)$ already yield all of $\mathrm{SO}(3)$:
every element of $\mathrm{SO}(3)$ is a product $R^{12}_{t_1}R^{23}_{t_2}R^{12}_{t_3}$ with real $t_1,t_2,t_3$ (Euler angles).
As $\mathrm{O}(3)=\mathrm{SO}(3)\cup(-\mathrm{SO}(3))$, the group is all of $\mathrm{O}(3)$, and \eqref{it:equiv-two-point-orthogonal} holds.

The law $\mu$ satisfies \textup{(G)}, \textup{(dF)}, $(T_{123})$, $(R_1)$ and $(RP_1)$;
so Theorem~\ref{thm:two-point}\eqref{it:two-point-kallen-lehmann} shows that \eqref{it:equiv-two-point-orthogonal} implies \eqref{it:equiv-two-point-kallen-lehmann}.

Assume \eqref{it:equiv-two-point-kallen-lehmann}, and let $g\in\mathrm{O}(3)$.
For $b\in L^1\cap L^2(\R^3;\C^3)$ the substitution $x=gy$ gives
\[
\widehat{g\cdot b}(p)=\det(g)\cdot g\,\hat b(g^{-1}p),
\]
and by the continuity of the Fourier transform on $L^2$ this holds for every $b\in L^2(\R^3;\C^3)$.
Insert it, for $g\cdot b$ and $g\cdot b'$, into the right side of \eqref{eq:rotation-kallen-lehmann}, and substitute $p=gq$.
Since $\det(g)^2=1$, $|gq|=|q|$ and $P^{\perp}(gq)=g\,P^{\perp}(q)\,g^{T}$, the result is the right side of \eqref{eq:rotation-kallen-lehmann} for $b$ and $b'$.
Hence $S_2(\overline{g\cdot b},g\cdot b')=S_2(\overline{b},b')$ for all $b,b'\in L^2(\R^3;\C^3)$.
As $\overline{g\cdot b}=g\cdot \overline{b}$ and $b\mapsto\overline{b}$ is a bijection of $L^2(\R^3;\C^3)$, this is \eqref{it:equiv-two-point-orthogonal}.

Assume \eqref{it:equiv-two-point-orthogonal}.
Let $n$ be a unit normal of one of the six diagonal planes, and let $b\in L^2(\R^3;\R^3)$ vanish almost everywhere on $\{n\cdot x<0\}$.
Choose $g\in\mathrm{O}(3)$ with $ge_1=n$.
As $r_{n,0}x=x-2(n\cdot x)\,n$, one has $g^{-1}r_{n,0}\,g=\theta_1$;
and the real field $b_1:=g^{-1}\cdot b$ vanishes almost everywhere on $\{x_1<0\}$, because $(gx)\cdot n=x_1$.
By \eqref{it:equiv-two-point-orthogonal},
\[
S_2(r_{n,0}\cdot b,b)=S_2\left(g^{-1}\cdot (r_{n,0}\cdot b),g^{-1}\cdot b\right)=S_2(\theta_1\cdot b_1,b_1).
\]
The field $\mathbb{B}(b_1)$ lies in $L^2(\Sigma_1^+,\mu)$ (Lemma~\ref{lem:domination-consequences}\eqref{it:field-extension}),
and $\Theta_1\mathbb{B}(b_1)=\mathbb{B}(\theta_1\cdot b_1)$ by \eqref{eq:rotation-conjugation}, $b_1$ being real.
Therefore
\[
S_2(\theta_1\cdot b_1,b_1)=\mu\left((\Theta_1\mathbb{B}(b_1))\cdot\mathbb{B}(b_1)\right)=\left\|[\mathbb{B}(b_1)]_1\right\|^2\ge0.
\]
This proves \eqref{it:equiv-two-point-diagonal}, for each of the six diagonal planes and each of its two unit normals.

It remains to derive \eqref{it:equiv-one-field-cone} from \eqref{it:equiv-two-point-diagonal}; let $n$ be the unit normal given there.
Let $C$ be the restriction of $S_2$ to $\mathcal{S}(\R^3;\R^3)$.
It is a real symmetric bilinear form, continuous for the topology of $\mathcal{S}(\R^3;\R^3)$ by Lemma~\ref{lem:domination-consequences}\eqref{it:moment-forms},
and $0\le C(b,b)=\mu\left(\mathbb{B}(b)^2\right)\le\|b\|^2$ by Lemma~\ref{lem:domination-consequences}\eqref{it:field-extension}.
Let $\gamma_C$ be the centred Gaussian law with the covariance $C$ (Lemma~\ref{lem:gaussian-covariance}\eqref{it:gaussian-laws-criteria}).
It satisfies \textup{(H)} and \textup{(dF)}, by the criteria of Lemma~\ref{lem:gaussian-covariance}.
Indeed, the bound $C(b,b)\le\|b\|^2$ gives \textup{(G)};
for $u\in\mathcal{S}(\R^3;\R)$ one has $\mathbb{B}(\nabla u)=0$ $\mu$-almost surely by \eqref{eq:rotation-gradient-fields}, so $C(\nabla u,\nabla u)=0$, which gives \textup{(dF)};
the form $C$ is invariant under the translations and under $\Ocub$, as $S_2$ is, so $\gamma_C$ is invariant under them;
and for $j'\in\{1,2,3\}$ and $b\in C^{\infty}_{\mathrm c}(\{x_{j'}>0\};\R^3)$, for which $\mathbb{B}(b)\in L^2(\Sigma_{j'}^+,\mu)$, \eqref{eq:rotation-conjugation} gives
\[
C(\theta_{j'}\cdot b,b)=\mu\left((\Theta_{j'}\mathbb{B}(b))\cdot\mathbb{B}(b)\right)=\left\|[\mathbb{B}(b)]_{j'}\right\|^2\ge0,
\]
whence $(RP_{j'})$ by Lemma~\ref{lem:gaussian-covariance}\eqref{it:gaussian-laws-positivity}, $C$ being invariant under $\theta_{j'}\in\Ocub$.

Next, $\gamma_C$ is reflection positive at the plane $\{n\cdot x=0\}$.
To see this, choose $g\in\mathrm{O}(3)$ with $ge_1=n$.
By Lemma~\ref{lem:transport}\eqref{it:transport-operators}, for $b\in\mathcal{S}(\R^3;\R^3)$,
\[
\left((g^{-1})_*\gamma_C\right)\left(\exp\left(i\,\mathbb{B}(b)\right)\right)=\gamma_C\left(\exp\left(i\,\mathbb{B}(g\cdot b)\right)\right)=\exp\left(-\frac{C(g\cdot b,g\cdot b)}{2}\right);
\]
by the uniqueness in Lemma~\ref{lem:gaussian-covariance}\eqref{it:gaussian-laws-criteria}, $(g^{-1})_*\gamma_C$ is the Gaussian law $\gamma_{C'}$ with the covariance $C'(b,b'):=C(g\cdot b,g\cdot b')$.
For a diagonal normal $n$ the reflection $r_{n,0}$ is a signed permutation matrix, hence an element of $\Ocub$, under which $C$ is invariant;
since $g\theta_1=r_{n,0}\,g$, the form $C'$ is invariant under $\theta_1$.
For $b\in C^{\infty}_{\mathrm c}(\{x_1>0\};\R^3)$ the field $b_1:=g\cdot b$ is real and vanishes outside $\{n\cdot x>0\}$, so \eqref{it:equiv-two-point-diagonal} gives
\[
C'(\theta_1\cdot b,b)=C(r_{n,0}\cdot b_1,b_1)=S_2(r_{n,0}\cdot b_1,b_1)\ge0.
\]
By Lemma~\ref{lem:gaussian-covariance}\eqref{it:gaussian-laws-positivity} the law $\gamma_{C'}$ satisfies $(RP_1)$,
and by Lemma~\ref{lem:transport}\eqref{it:transport-properties} the law $\gamma_C$ is reflection positive at $\{(ge_1)\cdot x=0\}=\{n\cdot x=0\}$.
Thus $\gamma_C$ satisfies \textup{(H)}, \textup{(dF)} and \eqref{it:equiv-diagonal-plane};
by the equivalences already proved, it satisfies \eqref{it:equiv-all-cones}.

Finally, fix an ordered pair $j\neq i$ and a field $b\in L^2(\R^3;\R^3)$ that vanishes almost everywhere on $\{x_j<0\}$.
Each of the two laws $\mu$ and $\gamma_C$ satisfies \textup{(G)}, $(T_{123})$, $(R_j)$ and $(RP_j)$,
and $\mathbb{B}(b)$ belongs to the space of its observables on the positive side (Lemma~\ref{lem:domination-consequences}\eqref{it:field-extension}).
The two laws have the same two-point function:
under $\gamma_C$ the variables $\mathbb{B}(b')$, $b'\in\mathcal{S}(\R^3;\R^3)$, have the covariance $C$ (Lemma~\ref{lem:gaussian-covariance}\eqref{it:gaussian-laws-criteria}),
and both two-point functions are bounded bilinear forms on $L^2(\R^3;\C^3)$ (Lemma~\ref{lem:domination-consequences}\eqref{it:moment-forms}), in which $\mathcal{S}(\R^3;\R^3)$ spans a dense subspace.
For either law let $\psi$ be the class of $\mathbb{B}(b)$ in its reconstruction space in the direction $e_j$, and let $\varrho_\psi$ be the spectral measure of $\psi$ with respect to its joint spectral measure.
The reconstruction identities of \cite[Proposition~8.19\,(ii), (iii)]{NoncompactQED},
the identity $U(a)\mathbb{B}(b)=\mathbb{B}(\tau_a\cdot b)$ of \eqref{eq:rotation-field-action}, extended to $b\in L^2(\R^3;\C^3)$ by continuity,
the relation $\Theta_j\mathbb{B}(b)=\mathbb{B}(\theta_j\cdot b)$ from \eqref{eq:rotation-conjugation},
and Lemma~\ref{lem:joint-spectrum}(b) give, for $t\ge0$ and $y\in e_j^{\perp}$,
\[
\int\mathrm{d}\varrho_\psi(\omega,p_\perp)\,\exp(-t\omega+ip_\perp\cdot y)=\left\langle\psi,\exp(-t\cdot\mathsf{H}_j)V_j(y)\psi\right\rangle=S_2\left(\theta_j\cdot b,\tau_{te_j+y}\cdot b\right).
\]
The right side is the same for the two laws.
By Lemma~\ref{lem:laplace-fourier} the two spectral measures coincide (compare \eqref{eq:rotation-one-field-spectra}).
For $\gamma_C$ this measure gives no mass to $\{\omega<|p_i|\}$, by $(C_{ji})$ for $\gamma_C$;
for $\mu$ it follows that $\left\|\mathfrak{E}_j(\{\omega<|p_i|\})[\mathbb{B}(b)]_j\right\|^2=0$.
Since $[\mathbb{B}(b'+ib'')]_j=[\mathbb{B}(b')]_j+i\,[\mathbb{B}(b'')]_j$, the classes $[\mathbb{B}(b)]_j$ with real $b$ as above span a dense subspace of the one-field space $\mathcal{H}^1_j$ of $\mu$.
The projection $\mathfrak{E}_j(\{\omega<|p_i|\})$ thus vanishes on $\mathcal{H}^1_j$:
the law $\mu$ satisfies $(C^1_{ji})$ for every ordered pair $j\neq i$, and in particular \eqref{it:equiv-one-field-cone}.
\end{proof}

\subsection{Sharpness of the hypotheses}\label{sec:sharpness}

\subsubsection{Gaussian tools}\label{sec:gaussian-tools}

For a centred Gaussian law the cone condition needs to be verified only on the classes of the linear functionals of the field.

\begin{lemma}[The cone of a Gaussian law]\label{lem:gaussian-cone}
Let $\mu$ be a centred Gaussian law on $(\mathcal{S}',\Sigma)$:
the random variables $F(f)$, $f\in\mathcal{S}$ real, are jointly centred Gaussian under $\mu$.
Suppose that $\mu$ satisfies $(T_{123})$, $(R_1)$ and $(RP_1)$,
let $\mathfrak{E}_1$ be its joint spectral measure (Lemma~\ref{lem:joint-spectrum}),
and let $c>0$ and $i\in\{2,3\}$.
For real $f\in C^\infty_{\mathrm c}(\{x_1>0\};\Lambda^2)$ the random variable $F(f)$ is $\Sigma_1^+$-measurable, since $\supp f\subset\{x_1>0\}$,
and square integrable, being Gaussian;
so it lies in $L^2(\Sigma_1^+,\mu)$ and has a class $[F(f)]_1\in\mathcal{H}_1$.
If
\[
\varrho_\psi\left(\left\{(\omega,p_\perp):\ \omega<c|p_i|\right\}\right)=0
\qquad\text{for }\psi=[F(f)]_1\text{ and every real }f\in C^\infty_{\mathrm c}(\{x_1>0\};\Lambda^2),
\]
then $\mathfrak{E}_1\left(\left\{(\omega,p_\perp):\ \omega<c|p_i|\right\}\right)=0$.
\end{lemma}

\begin{proof}
Let $K$ be a subset of $[0,\infty)\times\R^2$ with
\begin{equation}\label{eq:rotation-cone-properties}
\overline{K}=K,\qquad(0,0)\in K,\qquad K+K\subset K,
\end{equation}
and let $\complement K:=\left([0,\infty)\times\R^2\right)\setminus K$ be its complement.
It will be shown that $\mathfrak{E}_1(\complement K)=0$ whenever $\varrho_\psi(\complement K)=0$ for $\psi=[F(f)]_1$ and every real $f\in C^\infty_{\mathrm c}(\{x_1>0\};\Lambda^2)$.
This contains the lemma.
Indeed, the set $K_c:=\left\{(\omega,p_\perp)\in[0,\infty)\times\R^2:\ \omega\ge c|p_i|\right\}$ is closed and contains the origin,
and $\omega\ge c|p_i|$, $\omega'\ge c|p'_i|$ imply $\omega+\omega'\ge c|p_i+p'_i|$;
so $K_c$ has the properties~\eqref{eq:rotation-cone-properties}, and $\complement K_c=\{\omega<c|p_i|\}$.
A finite complex Borel measure $m$ on $[0,\infty)\times\R^2$ is said to be carried by $K$ if $|m|(\complement K)=0$, where $|m|$ is the total variation measure of $m$;
its norm is $\|m\|:=|m|\left([0,\infty)\times\R^2\right)$.
The transform of $m$ is the function
\[
(t,y)\longmapsto\int\mathrm{d}m(\omega,p_\perp)\,\exp\left(-t\omega+ip_\perp\cdot y\right)\qquad\left(t\ge0,\ y\in\R^2\right).
\]

\emph{Exponentials of the field.}
For real $f,f'\in\mathcal{S}$ put $C(f,f'):=\mu\left(F(f)\cdot F(f')\right)$.
By $(T_{123})$ and $(R_1)$ the form $C$ is invariant under all translations and under $\theta_1$.
Let $f_1,\dots,f_n\in C^\infty_{\mathrm c}(\{x_1>0\};\Lambda^2)$ be real, let $z_1,\dots,z_n\in\C$, and put
\[
\Phi:=\sum_{k=1}^n z_k\exp\left(iF(f_k)\right),\qquad
z'_k:=z_k\cdot\exp\left(-\tfrac12C(f_k,f_k)\right).
\]
Let $t\ge0$, $y\in e_1^{\perp}$ and $v:=te_1+y$.
By \cite[Proposition~8.19(i)--(iii)]{NoncompactQED}
(where the law is called $\rho$, the space $L^2(\Sigma_1^+,\rho)$ is denoted $\mathcal{E}_+$, and $\exp(-t\mathsf{H}_1)$ is denoted $T(t)$)
one has
\[
\left\langle[\Phi']_1,\exp(-t\mathsf{H}_1)V_1(y)[\Phi'']_1\right\rangle=\mu\left((\Theta_1\Phi')\cdot U(v)\Phi''\right)
\qquad\left(\Phi',\Phi''\in L^2(\Sigma_1^+,\mu)\right).
\]
By~\eqref{eq:rotation-conjugation} and Definition~\ref{def:field-actions}, for real $f$,
\[
\begin{aligned}
\Theta_1F(f)&=F(\theta_1\cdot f),&\qquad\Theta_1\exp\left(iF(f)\right)&=\exp\left(-iF(\theta_1\cdot f)\right),\\
U(v)F(f)&=F(\tau_v\cdot f),&\qquad U(v)\exp\left(iF(f)\right)&=\exp\left(iF(\tau_v\cdot f)\right).
\end{aligned}
\]
The random variable $F(\tau_v\cdot f_l)-F(\theta_1\cdot f_k)$ is centred Gaussian, and by the invariance of $C$ its variance is $C(f_l,f_l)+C(f_k,f_k)-2C(\theta_1\cdot f_k,\tau_v\cdot f_l)$.
It follows that
\[
\begin{aligned}
\left\langle[\Phi]_1,\exp(-t\mathsf{H}_1)V_1(y)[\Phi]_1\right\rangle
&=\sum_{k,l=1}^n\overline{z_k}\,z_l\cdot\mu\left(\exp\left(i\left(F(\tau_v\cdot f_l)-F(\theta_1\cdot f_k)\right)\right)\right)\\
&=\sum_{k,l=1}^n\overline{z'_k}\,z'_l\cdot\exp\left(C(\theta_1\cdot f_k,\tau_v\cdot f_l)\right).
\end{aligned}
\]
Let $\psi_k:=[F(f_k)]_1$, and let $m_{kl}:=\langle\psi_k,\mathfrak{E}_1(\cdot)\,\psi_l\rangle$; these are finite complex measures on $[0,\infty)\times\R^2$.
The identities above, applied to $F(f_k)$ and $F(f_l)$, together with Lemma~\ref{lem:joint-spectrum}(b), give
\[
C(\theta_1\cdot f_k,\tau_v\cdot f_l)=\left\langle\psi_k,\exp(-t\mathsf{H}_1)V_1(y)\,\psi_l\right\rangle
=\int\mathrm{d}m_{kl}(\omega,p_\perp)\,\exp\left(-t\omega+ip_\perp\cdot y\right).
\]

\emph{The measures $m_{kl}$ are carried by $K$.}
Let $A\subset[0,\infty)\times\R^2$ be a Borel set.
For every finite partition of $A$ into Borel sets $A'$, the Cauchy--Schwarz inequality, applied twice, gives
\[
\sum_{A'}\left|m_{kl}(A')\right|
\le\sum_{A'}\left\|\mathfrak{E}_1(A')\psi_k\right\|\cdot\left\|\mathfrak{E}_1(A')\psi_l\right\|
\le\left(m_{kk}(A)\cdot m_{ll}(A)\right)^{\frac12},
\]
so that $|m_{kl}|(A)\le\left(m_{kk}(A)\cdot m_{ll}(A)\right)^{\frac12}$.
By hypothesis $m_{kk}=\varrho_{\psi_k}$ vanishes on $\complement K$; hence every $m_{kl}$ is carried by $K$.

\emph{Exponentials of measures.}
On the additive semigroup $[0,\infty)\times\R^2$ the convolution $m*m'$ of two finite complex measures is the image of $m\otimes m'$ under the addition map.
One has $\|m*m'\|\le\|m\|\cdot\|m'\|$, and the transform of $m*m'$ is the product of the transforms of $m$ and $m'$,
because $(\omega,p_\perp)\mapsto\exp(-t\omega+ip_\perp\cdot y)$ is multiplicative.
If $m$ and $m'$ are carried by $K$, so is $m*m'$, since $|m*m'|$ is at most the image of $|m|\otimes|m'|$ under addition and $K+K\subset K$.
Consequently the series
\[
\exp_*(m):=\sum_{r\ge0}\frac{m^{*r}}{r!},\qquad m^{*0}:=\delta_{(0,0)}\ \text{(the unit mass at the origin)},
\]
converges in total variation, as $\|m^{*r}\|\le\|m\|^r$.
For $t\ge0$ the transform at $(t,y)$ is a linear functional of norm at most one on the finite complex measures;
so the transform of $\exp_*(m)$ is the exponential of the transform of $m$.
Since $(0,0)\in K$, the measure $\exp_*(m)$ is carried by $K$ if $m$ is.

\emph{Uniqueness for complex measures.}
Let $m$ and $m'$ be finite complex Borel measures on $[0,\infty)\times\R^2$ whose transforms coincide at all points $(t,y)$ with $t\in\{0,1,2,\dots\}$ and $y\in\R^2$.
Then $m=m'$.
For positive measures this is Lemma~\ref{lem:laplace-fourier}; the argument for complex measures is as follows.
Fix $t\in\{0,1,2,\dots\}$.
The finite complex measures
\[
B\longmapsto\int_{[0,\infty)\times B}\mathrm{d}m(\omega,p_\perp)\,\exp(-t\omega),\qquad
B\longmapsto\int_{[0,\infty)\times B}\mathrm{d}m'(\omega,p_\perp)\,\exp(-t\omega)
\]
on $\R^2$ have the same Fourier transform; so they coincide, by the uniqueness theorem for Fourier transforms of finite complex measures \cite[Theorem~1.3.6]{Rudin1962}.
Now fix a Borel set $B\subset\R^2$, and let $m_B$, $m'_B$ be the images of the finite complex measures $A\mapsto m(A\times B)$, $A\mapsto m'(A\times B)$ on $[0,\infty)$ under $\omega\mapsto\exp(-\omega)$, regarded as measures on $[0,1]$.
What has just been shown says that $\int\mathrm{d}m_B(s)\,s^t=\int\mathrm{d}m'_B(s)\,s^t$ for $t=0,1,2,\dots$.
By Weierstrass's approximation theorem $m_B$ and $m'_B$ have the same integrals of all continuous functions on $[0,1]$,
and by the uniqueness in the Riesz representation theorem for complex measures $m_B=m'_B$.
As $\omega\mapsto\exp(-\omega)$ is a homeomorphism of $[0,\infty)$ onto $(0,1]$, one obtains $m(A\times B)=m'(A\times B)$ for all Borel sets $A\subset[0,\infty)$ and $B\subset\R^2$.
The Borel sets on which $m$ and $m'$ agree contain the whole space and are closed under complements and under countable disjoint unions;
they form a Dynkin system that contains the $\pi$-system of the rectangles $A\times B$, and Dynkin's $\pi$--$\lambda$ theorem gives $m=m'$.

\emph{Conclusion.}
By Lemma~\ref{lem:joint-spectrum}(b) and the first step, the transform of the finite positive measure $\varrho_{[\Phi]_1}$ is
\[
(t,y)\longmapsto\sum_{k,l=1}^n\overline{z'_k}\,z'_l\cdot\exp\left(\int\mathrm{d}m_{kl}(\omega,p_\perp)\,\exp\left(-t\omega+ip_\perp\cdot y\right)\right),
\]
which is the transform of the finite complex measure $\nu:=\sum_{k,l}\overline{z'_k}\,z'_l\cdot\exp_*(m_{kl})$.
Hence $\varrho_{[\Phi]_1}=\nu$.
The measure $\nu$ is carried by $K$, and therefore
\[
\left\|\mathfrak{E}_1(\complement K)[\Phi]_1\right\|^2=\varrho_{[\Phi]_1}(\complement K)=\nu(\complement K)=0.
\]
The finite linear combinations of the functions $\exp(iF(f))$, $f\in C^\infty_{\mathrm c}(\{x_1>0\};\Lambda^2)$ real, are dense in $L^2(\Sigma_1^+,\mu)$
by \cite[Lemma~8.12(i)]{NoncompactQED}, where the tests range even over a countable subset of $C^\infty_{\mathrm c}(\{x_1>0\};\Lambda^2)$;
and the class map $[\,\cdot\,]_1$ is a contraction with dense range \cite[Proposition~8.19(i)]{NoncompactQED}.
So the vectors $[\Phi]_1$ are dense in $\mathcal{H}_1$, and the bounded operator $\mathfrak{E}_1(\complement K)$ is zero.
\end{proof}

The covariance of a centred Gaussian law can be prescribed by a matrix valued Fourier multiplier.

\begin{lemma}[Gaussian laws with a multiplier]\label{lem:multiplier-laws}
Let $M$ be a measurable map of $\R^3$ into the real $3\times3$ matrices such that $M(p)$ is symmetric and positive semidefinite for every $p$.
Suppose that $M$ is even, $M(-p)=M(p)$, and polynomially bounded:
for some $k\in\N$ the function $p\mapsto\left(1+|p|^2\right)^{-k}\|M(p)\|$ is bounded.
For $b,b'\in\mathcal{S}(\R^3;\R^3)$ put
\[
C_M(b,b'):=(2\pi)^{-3}\int\mathrm{d}p\,\hat b(p)^{\dagger}M(p)\,\hat b'(p),
\qquad\text{where }\hat b(p)=\int\mathrm{d}x\,\exp(-ip\cdot x)\cdot b(x).
\]
Then $C_M$ is a real symmetric bilinear form on $\mathcal{S}(\R^3;\R^3)$, continuous for the topology of $\mathcal{S}$,
with $C_M(b,b)\ge0$ and $C_M(\tau_a\cdot b,\tau_a\cdot b')=C_M(b,b')$ for all $a\in\R^3$.
The centred Gaussian law $\gamma_M:=\gamma_{C_M}$ with the covariance $C_M$ (Lemma~\ref{lem:gaussian-covariance}(a)) satisfies $(T_{123})$, and:
\begin{itemize}
\item \textup{(G)} holds for $\gamma_M$ if and only if $M(p)\le1$ for almost every $p$;
\item \textup{(dF)} holds for $\gamma_M$ if and only if $M(p)\,p=0$ for almost every $p$;
\item for $g\in\mathrm{O}(3)$ one has $g_*\gamma_M=\gamma_M$ if and only if $M(gp)=g\,M(p)\,g^{T}$ for almost every $p$.
\end{itemize}
\end{lemma}

\begin{proof}
For real $b$ one has $\hat b(-p)=\overline{\hat b(p)}$.
As $M$ is real, symmetric and even, the substitution $p\mapsto-p$ gives $\overline{C_M(b,b')}=C_M(b,b')$;
and $C_M(b',b)=\overline{C_M(b,b')}$ because every $M(p)$ is Hermitian.
So $C_M$ is real and symmetric, and $C_M(b,b)\ge0$ because $M(p)\ge0$.
Since $\int\mathrm{d}p\,\left(1+|p|^2\right)^{-2}<\infty$, the number $|C_M(b,b')|$ is at most a constant times
$\sup_p\left(1+|p|^2\right)^{k+1}|\hat b(p)|\cdot\sup_p\left(1+|p|^2\right)|\hat b'(p)|$,
and the Fourier transform is continuous on $\mathcal{S}$; this is the continuity of $C_M$.
The transform of $\tau_a\cdot b$ is $p\mapsto\exp(-ip\cdot a)\,\hat b(p)$, and the factor $\exp(-ip\cdot a)$ cancels in $C_M$.
By Lemma~\ref{lem:gaussian-covariance}(a) the law $\gamma_M$ exists, and it is invariant under all translations because $C_M$ is.

\emph{Claim.}
Let $D$ be a measurable, even and polynomially bounded map of $\R^3$ into the real symmetric $3\times3$ matrices
such that $\int\mathrm{d}p\,\hat b(p)^{\dagger}D(p)\,\hat b(p)\ge0$ for every real $b\in\mathcal{S}(\R^3;\R^3)$.
Then $D(p)\ge0$ for almost every $p$.
Indeed, let $v\in\R^3$, and let $u\in\mathcal{S}(\R^3)$ be real and even.
The field $b$ with $\hat b(p)=u(p)\,v$ lies in $\mathcal{S}(\R^3;\R^3)$ and is real, because $\int\mathrm{d}p\,\exp(ip\cdot x)\cdot u(p)=\int\mathrm{d}p\,\cos(p\cdot x)\cdot u(p)$;
so $\int\mathrm{d}p\,u(p)^2\cdot v^{T}D(p)\,v\ge0$.
Let $h\in C^\infty_{\mathrm c}(\R^3)$ be even and non-negative.
For $\varepsilon>0$ the function $u_\varepsilon:=\left(h+\varepsilon\exp\left(-|p|^2\right)\right)^{\frac12}$ is smooth, real and even,
and it equals $\varepsilon^{\frac12}\exp\left(-\tfrac12|p|^2\right)$ outside the support of $h$; so $u_\varepsilon\in\mathcal{S}(\R^3)$.
The limit $\varepsilon\to0$ in $\int\mathrm{d}p\,u_\varepsilon(p)^2\cdot v^{T}D(p)\,v\ge0$ gives $\int\mathrm{d}p\,h(p)\cdot v^{T}D(p)\,v\ge0$.
As $p\mapsto v^{T}D(p)\,v$ is even, the same holds for every non-negative $h\in C^\infty_{\mathrm c}(\R^3)$ (replace $h$ by $p\mapsto\tfrac12\left(h(p)+h(-p)\right)$).
Hence $v^{T}D(p)\,v\ge0$ for $p$ outside a null set.
Taking all $v$ of a countable dense subset of $\R^3$ at once, one obtains $D(p)\ge0$ for almost every $p$.

\emph{Gaussian domination.}
By Lemma~\ref{lem:gaussian-covariance}(a), \textup{(G)} holds for $\gamma_M$ if and only if $C_M(b,b)\le\|b\|^2$ for all real $b\in\mathcal{S}(\R^3;\R^3)$.
By Plancherel's theorem $\|b\|^2=(2\pi)^{-3}\int\mathrm{d}p\,|\hat b(p)|^2$;
so \textup{(G)} holds if and only if $\int\mathrm{d}p\,\hat b(p)^{\dagger}\left(1-M(p)\right)\hat b(p)\ge0$ for these $b$.
This is the case if $M\le1$ almost everywhere, and conversely the claim with $D=1-M$ gives $M\le1$ almost everywhere.

\emph{Closedness.}
For real $u\in\mathcal{S}(\R^3)$ the transform of $\nabla u$ is $p\mapsto ip\,\hat u(p)$, so that
\[
C_M(\nabla u,\nabla u)=(2\pi)^{-3}\int\mathrm{d}p\,|\hat u(p)|^2\cdot p^{T}M(p)\,p.
\]
Since $M(p)\ge0$, one has $p^{T}M(p)\,p\ge0$, with equality if and only if $M(p)\,p=0$.
If $M(p)\,p=0$ for almost every $p$, then $C_M(\nabla u,\nabla u)=0$ for every $u$.
Conversely, if $C_M(\nabla u,\nabla u)=0$ for $u(x)=\exp\left(-|x|^2\right)$, whose transform has no zero, then $p^{T}M(p)\,p=0$, that is, $M(p)\,p=0$, for almost every $p$.
By Lemma~\ref{lem:gaussian-covariance}(a), \textup{(dF)} holds for $\gamma_M$ if and only if $C_M(\nabla u,\nabla u)=0$ for all real $u\in\mathcal{S}(\R^3)$.

\emph{Invariance.}
Let $g\in\mathrm{O}(3)$.
By Definition~\ref{def:field-actions}, $(g\cdot b)(x)=\det(g)\cdot g\,b(g^{-1}x)$; hence the transform of $g\cdot b$ is $p\mapsto\det(g)\cdot g\,\hat b(g^{-1}p)$,
and the change of variables $p\mapsto gp$, which preserves Lebesgue measure, gives
\[
C_M(g\cdot b,g\cdot b)=(2\pi)^{-3}\int\mathrm{d}p\,\hat b(p)^{\dagger}g^{T}M(gp)\,g\,\hat b(p).
\]
By Lemma~\ref{lem:gaussian-covariance}(a), $g_*\gamma_M=\gamma_M$ if and only if $C_M(g\cdot b,g\cdot b)=C_M(b,b)$ for all real $b\in\mathcal{S}(\R^3;\R^3)$.
This holds if $M(gp)=g\,M(p)\,g^{T}$ for almost every $p$.
Conversely, the function $D(p):=g^{T}M(gp)\,g-M(p)$ is measurable, even, polynomially bounded and real symmetric,
and the claim applied to $D$ and to $-D$ gives $D=0$ almost everywhere, that is, $M(gp)=g\,M(p)\,g^{T}$ for almost every $p$.
\end{proof}

Write $p=(p_1,p_\perp)$ and $x=(x_1,x_\perp)$ with $p_\perp=(p_2,p_3)$ and $x_\perp=(x_2,x_3)$.
The plane $e_1^{\perp}$ is identified with $\R^2$ through the coordinates $(x_2,x_3)$, so that $p_\perp\cdot y=p_2y_2+p_3y_3$ for $y\in e_1^{\perp}$.
As a matrix, $\theta_1=\diag(-1,1,1)$; so $\theta_1p=(-p_1,p_\perp)$ and $-\theta_1=\diag(1,-1,-1)$.
When the multiplier is a sum of partial fractions in $p_1$, the covariance between a field and its reflected translates is a sum of residues.

\begin{lemma}[Time kernels]\label{lem:time-kernels}
Let $M$ be a measurable, even and polynomially bounded function on $\R^3$ whose values are positive semidefinite real symmetric $3\times3$ matrices,
and suppose that $\theta_1M(-p_1,p_\perp)\,\theta_1=M(p_1,p_\perp)$ for all $p$.
Let $n\ge1$, and let $\omega_1,\dots,\omega_n$ be real measurable functions on $\R^2$ such that $\omega_l(p_\perp)>0$ for almost every $p_\perp\in\R^2$.
For almost every $p_\perp\in\R^2$ let polynomials $Q_{p_\perp}$ and $N_1(\cdot,p_\perp),\dots,N_n(\cdot,p_\perp)$ in one variable, with coefficients in the complex $3\times3$ matrices, be given,
in such a way that the coefficients of each $N_l(\cdot,p_\perp)$ are measurable functions of $p_\perp$ and that, for almost every $p_\perp$,
\begin{equation}\label{eq:rotation-partial-fractions}
M(p_1,p_\perp)=Q_{p_\perp}(p_1)+\sum_{l=1}^n\frac{N_l(p_1,p_\perp)}{p_1^2+\omega_l(p_\perp)^2}\qquad(p_1\in\R).
\end{equation}
Put $A_l(p_\perp):=-\theta_1N_l\left(-i\omega_l(p_\perp),p_\perp\right)$.
For a real $b\in C^\infty_{\mathrm c}(\{x_1>0\};\R^3)$ let
\[
\begin{aligned}
\hat b_\perp(x_1,p_\perp)&:=\int\mathrm{d}^2x_\perp\,\exp(-ip_\perp\cdot x_\perp)\cdot b(x_1,x_\perp),\\
L_l(p_\perp)&:=\int_0^\infty\mathrm{d}x_1\,\exp\left(-\omega_l(p_\perp)x_1\right)\cdot\hat b_\perp(x_1,p_\perp)\in\C^3.
\end{aligned}
\]
Then, with $C_M$ and $\gamma_M$ as in Lemma~\ref{lem:multiplier-laws}, for $t\ge0$ and $y\in e_1^{\perp}$
\begin{equation}\label{eq:rotation-time-kernel}
\begin{split}
&C_M(\theta_1\cdot b,\tau_{te_1+y}\cdot b)\\
&\quad=(2\pi)^{-2}\int\mathrm{d}^2p_\perp\,\exp(-ip_\perp\cdot y)\sum_{l=1}^n\frac{\exp\left(-\omega_l(p_\perp)t\right)}{2\omega_l(p_\perp)}\,L_l(p_\perp)^{\dagger}A_l(p_\perp)\,L_l(p_\perp).
\end{split}
\end{equation}
If $A_1(p_\perp),\dots,A_n(p_\perp)$ are Hermitian and positive semidefinite for almost every $p_\perp$,
then $\gamma_M$ satisfies $(R_1)$ and $(RP_1)$.
In that case the field $\mathbb{B}(b)$ lies in $L^2(\Sigma_1^+,\gamma_M)$,
being $\Sigma_1^+$-measurable (the support of $b$ is a compact subset of $\{x_1>0\}$) and, as a Gaussian random variable, square integrable;
and, for the joint spectral measure $\mathfrak{E}_1$ of $\gamma_M$ (Lemma~\ref{lem:joint-spectrum}),
the measure $\varrho_{[\mathbb{B}(b)]_1}$ coincides with the sum over $l=1,\dots,n$ of the images under the maps $p_\perp\mapsto(\omega,p'_\perp)=\left(\omega_l(p_\perp),-p_\perp\right)$ of the positive measures
\[
(2\pi)^{-2}\left(2\omega_l(p_\perp)\right)^{-1}\cdot L_l(p_\perp)^{\dagger}A_l(p_\perp)\,L_l(p_\perp)\,\mathrm{d}^2p_\perp .
\]
\end{lemma}

\begin{proof}
\emph{A rule for integrals over $p_1$.}
Let $h\in C^\infty_{\mathrm c}((0,\infty))$ with the transform $\hat h(p_1):=\int\mathrm{d}u\,\exp(-ip_1u)\cdot h(u)$, let $\omega>0$, and let $P$ be a polynomial.
Then
\begin{equation}\label{eq:rotation-residue-rule}
\begin{aligned}
\frac{1}{2\pi}\int\mathrm{d}p_1\,\hat h(p_1)\cdot P(p_1)&=0,\\
\frac{1}{2\pi}\int\mathrm{d}p_1\,\frac{\hat h(p_1)\cdot P(p_1)}{p_1^2+\omega^2}&=\frac{P(-i\omega)}{2\omega}\int_0^\infty\mathrm{d}u\,\exp(-\omega u)\cdot h(u).
\end{aligned}
\end{equation}
For the first formula note that $p_1^k\hat h$ is the transform of $\left(-i\,\mathrm{d}/\mathrm{d}u\right)^kh$;
by the Fourier inversion formula the integral equals $\left(P(-i\,\mathrm{d}/\mathrm{d}u)\,h\right)(0)$, which is zero because $h$ vanishes near $0$.
Next, $(2\pi)^{-1}\int\mathrm{d}p_1\,\exp(-ip_1u)\cdot\left(p_1^2+\omega^2\right)^{-1}=(2\omega)^{-1}\cdot\exp(-\omega|u|)$.
Writing $\hat h(p_1)$ as an integral over $u$ and exchanging the two integrations,
which is legitimate by Fubini's theorem since $h$ and $p_1\mapsto\left(p_1^2+\omega^2\right)^{-1}$ are integrable,
one obtains the case $P=1$ of the second formula, because $h$ vanishes outside $(0,\infty)$.
For $P(p_1)=p_1$, apply this to $-ih'$, whose transform is $p_1\hat h$, and integrate by parts:
\[
\frac{1}{2\pi}\int\mathrm{d}p_1\,\frac{p_1\,\hat h(p_1)}{p_1^2+\omega^2}
=\frac{-i}{2\omega}\int_0^\infty\mathrm{d}u\,\exp(-\omega u)\cdot h'(u)
=\frac{-i\omega}{2\omega}\int_0^\infty\mathrm{d}u\,\exp(-\omega u)\cdot h(u).
\]
In general $P(p_1)=R(p_1)\left(p_1^2+\omega^2\right)+c_0+c_1p_1$ with a polynomial $R$ and complex numbers $c_0,c_1$, and $c_0+c_1\cdot(-i\omega)=P(-i\omega)$;
the first formula for $R$ and the two special cases give the second formula.

\emph{The covariance as an integral over $p_\perp$.}
By Definition~\ref{def:field-actions}, the transforms of $\theta_1\cdot b$ and $\tau_{te_1+y}\cdot b$ are
\[
(p_1,p_\perp)\longmapsto-\theta_1\hat b(-p_1,p_\perp)\qquad\text{and}\qquad p\longmapsto\exp\left(-ip_1t-ip_\perp\cdot y\right)\cdot\hat b(p).
\]
Moreover $\hat b(p_1,p_\perp)=\int\mathrm{d}s\,\exp(-ip_1s)\cdot\hat b_\perp(s,p_\perp)$, and therefore, for the components $a=1,2,3$,
$\overline{\hat b_a(-p_1,p_\perp)}=\int\mathrm{d}s\,\exp(-ip_1s)\cdot\overline{\hat b_{\perp,a}(s,p_\perp)}$.
For $a,a'\in\{1,2,3\}$ let
\[
h_{aa'}(u;p_\perp):=\int\mathrm{d}s\,\overline{\hat b_{\perp,a}(s,p_\perp)}\cdot\hat b_{\perp,a'}(u-s-t,p_\perp).
\]
As $b$ has compact support in $\{x_1>0\}$ and $t\ge0$, the function $u\mapsto h_{aa'}(u;p_\perp)$ lies in $C^\infty_{\mathrm c}((0,\infty))$,
and its transform is $p_1\mapsto\overline{\hat b_a(-p_1,p_\perp)}\cdot\hat b_{a'}(p_1,p_\perp)\cdot\exp(-ip_1t)$.
Hence
\[
\begin{aligned}
&C_M(\theta_1\cdot b,\tau_{te_1+y}\cdot b)\\
&\qquad=(2\pi)^{-3}\int\mathrm{d}^2p_\perp\,\exp(-ip_\perp\cdot y)\int\mathrm{d}p_1\sum_{a,a'=1}^3(-\theta_1)_{aa}\,M_{aa'}(p_1,p_\perp)\cdot\hat h_{aa'}(p_1;p_\perp).
\end{aligned}
\]
The integrand is absolutely integrable over $\R^3$, since $\hat b$ decreases faster than every power of $|p|$ while $\|M(p)\|$ grows at most like a power of $|p|$.
By Fubini's theorem the integral over $p$ may therefore be computed first over $p_1$ and then over $p_\perp$, as written, and the inner integral over $p_1$ is an integrable function of $p_\perp$.
For almost every $p_\perp$ insert~\eqref{eq:rotation-partial-fractions} and apply~\eqref{eq:rotation-residue-rule} entrywise, with $h=h_{aa'}(\cdot\,;p_\perp)$.
The polynomial $Q_{p_\perp}$ contributes nothing.
For $\omega>0$ the substitution $u=s+s'+t$ gives
\[
\begin{aligned}
\int_0^\infty\mathrm{d}u\,\exp(-\omega u)\cdot h_{aa'}(u;p_\perp)
&=\exp(-\omega t)\cdot\overline{\int_0^\infty\mathrm{d}s\,\exp(-\omega s)\cdot\hat b_{\perp,a}(s,p_\perp)}\\
&\qquad\cdot\int_0^\infty\mathrm{d}s'\,\exp(-\omega s')\cdot\hat b_{\perp,a'}(s',p_\perp);
\end{aligned}
\]
so, with $L_{l,a}$ the components of $L_l$, the term of $N_l$ contributes
\[
\begin{aligned}
&\frac{\exp\left(-\omega_l(p_\perp)t\right)}{2\omega_l(p_\perp)}\sum_{a,a'=1}^3\overline{L_{l,a}(p_\perp)}\cdot(-\theta_1)_{aa}\left[N_l\left(-i\omega_l(p_\perp),p_\perp\right)\right]_{aa'}\cdot L_{l,a'}(p_\perp)\\
&\qquad=\frac{\exp\left(-\omega_l(p_\perp)t\right)}{2\omega_l(p_\perp)}\,L_l(p_\perp)^{\dagger}A_l(p_\perp)\,L_l(p_\perp).
\end{aligned}
\]
This proves~\eqref{eq:rotation-time-kernel}, and the integrand there is an integrable function of $p_\perp$.

\emph{Consequences.}
The hypothesis $\theta_1M(-p_1,p_\perp)\,\theta_1=M(p_1,p_\perp)$ says that $M(\theta_1p)=\theta_1M(p)\,\theta_1^{T}$ for all $p$;
by Lemma~\ref{lem:multiplier-laws}, $\gamma_M$ satisfies $(R_1)$.
So Lemma~\ref{lem:gaussian-covariance}(a) gives $C_M(\theta_1\cdot b',\theta_1\cdot b')=C_M(b',b')$ for every real $b'\in\mathcal{S}(\R^3;\R^3)$;
as $C_M$ is symmetric and bilinear and $b'\mapsto\theta_1\cdot b'$ is linear, polarisation yields $C_M(\theta_1\cdot b',\theta_1\cdot b'')=C_M(b',b'')$ for all real $b',b''\in\mathcal{S}(\R^3;\R^3)$.
Let now $A_1(p_\perp),\dots,A_n(p_\perp)$ be Hermitian and positive semidefinite for almost every $p_\perp$.
At $t=0$, $y=0$, \eqref{eq:rotation-time-kernel} gives $C_M(\theta_1\cdot b,b)\ge0$ for every real $b\in C^\infty_{\mathrm c}(\{x_1>0\};\R^3)$,
the integrand being non-negative almost everywhere;
and Lemma~\ref{lem:gaussian-covariance}(b), which applies since $C_M$ is invariant under $\theta_1$, gives $(RP_1)$.
At $t=0$, $y=0$ each of the $n$ summands of the integrand of~\eqref{eq:rotation-time-kernel} is a measurable function of $p_\perp$ that is non-negative almost everywhere
(the $L_l$ and the $A_l$ are measurable, because the $\omega_l$ and the coefficients of the $N_l(\cdot,p_\perp)$ are), and their sum is integrable;
so each summand is integrable, and the $n$ measures of the statement are finite.
The law $\gamma_M$ satisfies $(T_{123})$ (Lemma~\ref{lem:multiplier-laws}), $(R_1)$ and $(RP_1)$;
let $\mathfrak{E}_1$ be its joint spectral measure (Lemma~\ref{lem:joint-spectrum}), and let $\psi:=[\mathbb{B}(b)]_1$.
By Lemma~\ref{lem:gaussian-covariance}(a), $\gamma_M\left(\mathbb{B}(b')\cdot\mathbb{B}(b'')\right)=C_M(b',b'')$ for real $b',b''\in\mathcal{S}(\R^3;\R^3)$.
By \cite[Proposition~8.19(i)--(iii)]{NoncompactQED}, \eqref{eq:rotation-conjugation} and Definition~\ref{def:field-actions}, for $t\ge0$ and $y\in e_1^{\perp}$,
\[
\begin{aligned}
\left\langle\psi,\exp(-t\mathsf{H}_1)V_1(y)\,\psi\right\rangle
&=\gamma_M\left(\left(\Theta_1\mathbb{B}(b)\right)\cdot\mathbb{B}(\tau_{te_1+y}\cdot b)\right)\\
&=\gamma_M\left(\mathbb{B}(\theta_1\cdot b)\cdot\mathbb{B}(\tau_{te_1+y}\cdot b)\right)
=C_M(\theta_1\cdot b,\tau_{te_1+y}\cdot b).
\end{aligned}
\]
By Lemma~\ref{lem:joint-spectrum}(b) the left side is $\int\mathrm{d}\varrho_\psi(\omega,p'_\perp)\,\exp\left(-t\omega+ip'_\perp\cdot y\right)$.
By~\eqref{eq:rotation-time-kernel} the right side is the integral of the same function against the sum of the image measures of the statement,
because the image of $p_\perp$ under the $l$-th map is $\left(\omega_l(p_\perp),-p_\perp\right)$.
Both measures are finite and positive, and both are carried by $[0,\infty)\times\R^2$:
the first by Lemma~\ref{lem:joint-spectrum}(b), the second because $\omega_l(p_\perp)>0$ for almost every $p_\perp$.
Lemma~\ref{lem:laplace-fourier} shows that they coincide.
\end{proof}

The cone with an arbitrary constant $c>0$ is described by five-point inequalities with unequal steps, and it is preserved under finite mixtures.

\begin{lemma}[Cones with a constant; mixtures]\label{lem:cone-constant}
Let $j\neq i$ be elements of $\{1,2,3\}$, and let $c>0$.
For a law $\nu$ with $(T_{123})$, $(R_j)$ and $(RP_j)$ let $\mathfrak{E}_j(\nu)$ be its joint spectral measure.
\begin{enumerate}
\item\label{it:cone-constant-five-point}
Let $\mu$ satisfy $(T_{123})$, $(R_j)$ and $(RP_j)$, and for $\Phi\in L^2(\Sigma_j^+,\mu)$ let $u(v):=\mu\left((\Theta_j\Phi)\cdot U(v)\Phi\right)$, $v\in\R^3$.
Then $\mathfrak{E}_j(\mu)\left(\{\omega<c|p_i|\}\right)=0$ if and only if
\begin{equation}\label{eq:rotation-five-point-c}
u(0)-4u(te_j)+u(2te_j)+u(te_j+cte_i)+u(te_j-cte_i)\ge0
\end{equation}
for all $t>0$ and all $\Phi\in L^2(\Sigma_j^+,\mu)$;
and this is the case if and only if~\eqref{eq:rotation-five-point-c} holds for all $t>0$ and all $\Phi$ in one subset of $L^2(\Sigma_j^+,\mu)$ whose classes are dense in $\mathcal{H}_j$.
\item\label{it:cone-constant-mixtures}
Let $\mu$ be a convex combination of finitely many laws $\mu_k$ that satisfy $(T_{123})$, $(R_j)$ and $(RP_j)$.
Then $\mu$ satisfies $(T_{123})$, $(R_j)$ and $(RP_j)$;
and if $\mathfrak{E}_j(\mu_k)\left(\{\omega<c|p_i|\}\right)=0$ for every $k$, then $\mathfrak{E}_j(\mu)\left(\{\omega<c|p_i|\}\right)=0$.
\end{enumerate}
\end{lemma}

\begin{proof}
\emph{Part~\eqref{it:cone-constant-five-point}.}
Let $t>0$, $\Phi\in L^2(\Sigma_j^+,\mu)$, and $\psi:=[\Phi]_j$.
By \cite[Proposition~8.19(i)--(iii)]{NoncompactQED}, $u(ae_j+a'e_i)=\left\langle\psi,\exp(-a\mathsf{H}_j)V_j(a'e_i)\,\psi\right\rangle$ for $a\ge0$ and $a'\in\R$.
Since $V_j(-cte_i)=V_j(cte_i)^*$, the left side of~\eqref{eq:rotation-five-point-c} equals $\langle\psi,Q_t\psi\rangle$ with the bounded self-adjoint operator
\[
\begin{aligned}
Q_t&:=1-4\exp(-t\mathsf{H}_j)+\exp(-2t\mathsf{H}_j)+\exp(-t\mathsf{H}_j)\left(V_j(cte_i)+V_j(cte_i)^*\right)\\
&=\int\mathrm{d}\mathfrak{E}_j(\omega,p_\perp)\,2\exp(-t\omega)\cdot\left(\cosh(t\omega)+\cos(ctp_i)-2\right);
\end{aligned}
\]
the second expression follows from Lemma~\ref{lem:joint-spectrum}(b) and the identity
$1-4x+x^2+2x\cos(ctp_i)=2x\left(\tfrac12\left(x+x^{-1}\right)-2+\cos(ctp_i)\right)$ at $x=\exp(-t\omega)$.
Hence the left side of~\eqref{eq:rotation-five-point-c} is $\int\mathrm{d}\varrho_\psi(\omega,p_\perp)\,2\exp(-t\omega)\cdot\left(\cosh(t\omega)+\cos(ctp_i)-2\right)$.

Suppose that $\mathfrak{E}_j(\mu)\left(\{\omega<c|p_i|\}\right)=0$.
Where $\omega\ge c|p_i|$, one has $\cosh(t\omega)\ge\cosh(ctp_i)\ge2-\cos(ctp_i)$;
the second inequality holds because $\cosh s+\cos s=2\sum_{r\ge0}s^{4r}/(4r)!\ge2$ for real $s$.
So the integrand is non-negative on a set that carries $\varrho_\psi$, and~\eqref{eq:rotation-five-point-c} holds for all $t>0$ and all $\Phi$.

Suppose conversely that~\eqref{eq:rotation-five-point-c} holds for all $t>0$ and all $\Phi$ in a subset of $L^2(\Sigma_j^+,\mu)$ whose classes are dense in $\mathcal{H}_j$.
The quadratic form of $Q_t$ is continuous and non-negative on a dense set of vectors; so $Q_t\ge0$,
and Lemma~\ref{lem:joint-spectrum}(d), applied to the bounded real function $2\exp(-t\omega)\cdot\left(\cosh(t\omega)+\cos(ctp_i)-2\right)$, gives
$\mathfrak{E}_j(\mu)\left(\left\{\cosh(t\omega)+\cos(ctp_i)<2\right\}\right)=0$.
As $\cosh(t\omega)-1=2\sinh^2(t\omega/2)$ and $1-\cos(ctp_i)=2\sin^2(ctp_i/2)$, the complementary set is $\left\{\sinh(t\omega/2)\ge|\sin(ctp_i/2)|\right\}$.
Take $t=1/n$, $n=1,2,\dots$.
On the intersection of these sets, which carries $\mathfrak{E}_j(\mu)$, one has
\[
2n\sinh\left(\frac{\omega}{2n}\right)\ge2n\left|\sin\left(\frac{cp_i}{2n}\right)\right|\qquad(n=1,2,\dots),
\]
and the limit $n\to\infty$ gives $\omega\ge c|p_i|$.
So $\mathfrak{E}_j(\mu)\left(\{\omega<c|p_i|\}\right)=0$.
Since the inequalities for all $\Phi\in L^2(\Sigma_j^+,\mu)$ contain those for any subset, and the classes of all of $L^2(\Sigma_j^+,\mu)$ are dense in $\mathcal{H}_j$ \cite[Proposition~8.19(i)]{NoncompactQED}, the three conditions are equivalent.

\emph{Part~\eqref{it:cone-constant-mixtures}.}
Invariance of a law under a map and reflection positivity at a plane are expressed by equalities and inequalities that are linear in the law;
so $\mu$ satisfies $(T_{123})$, $(R_j)$ and $(RP_j)$, and $\mathfrak{E}_j(\mu)$ is defined.
Let $\Phi$ be a finite linear combination of functions $\exp(iF(f))$, $f\in C^\infty_{\mathrm c}(\{x_j>0\};\Lambda^2)$ real.
It is bounded and $\Sigma_j^+$-measurable, so it lies in $L^2(\Sigma_j^+,\mu)$ and in every $L^2(\Sigma_j^+,\mu_k)$.
For each $v$ the function $(\Theta_j\Phi)\cdot U(v)\Phi$ on $\mathcal{S}'$ is bounded and does not depend on the law;
so the function $u$ formed with $\mu$ is the same convex combination of the functions $u$ formed with the $\mu_k$.
By part~\eqref{it:cone-constant-five-point}, applied to each $\mu_k$, each of the latter satisfies~\eqref{eq:rotation-five-point-c} for every $t>0$, and so does their convex combination.
The classes of these $\Phi$ are dense in $\mathcal{H}_j(\mu)$, by \cite[Lemma~8.12(i)]{NoncompactQED} and \cite[Proposition~8.19(i)]{NoncompactQED}.
Part~\eqref{it:cone-constant-five-point}, applied to $\mu$, gives $\mathfrak{E}_j(\mu)\left(\{\omega<c|p_i|\}\right)=0$.
\end{proof}

For fields of the following shape the vectors $L_l$ of Lemma~\ref{lem:time-kernels} are explicit.

\begin{definition}[A test field]\label{def:test-field}
Let $h\in C^\infty_{\mathrm c}((0,\infty))$ be non-negative and not identically zero, and let $v\in C^\infty_{\mathrm c}(\R^2)$ be real and not identically zero.
The test field formed with $h$ and $v$ is
\[
b_*(x):=h(x_1)\cdot v(x_2,x_3)\cdot e_1\qquad(x\in\R^3).
\]
\end{definition}

The field $b_*$ is real and lies in $C^\infty_{\mathrm c}(\{x_1>0\};\R^3)$.
With $\hat v(p_\perp):=\int\mathrm{d}^2x_\perp\,\exp(-ip_\perp\cdot x_\perp)\cdot v(x_\perp)$ one has $\hat b_{*,\perp}(x_1,p_\perp)=h(x_1)\cdot\hat v(p_\perp)\cdot e_1$.
Hence, for every real measurable function $\omega_l$ on $\R^2$ that is positive almost everywhere, the vector $L_l$ of Lemma~\ref{lem:time-kernels} formed with $b=b_*$ is
\[
L_l(p_\perp)=\hat v(p_\perp)\cdot\left(\int_0^\infty\mathrm{d}x_1\,h(x_1)\cdot\exp\left(-\omega_l(p_\perp)x_1\right)\right)\cdot e_1.
\]
The integral is positive, since $h$ is continuous, non-negative and not identically zero.
As $v$ has compact support and is not identically zero, $\hat v$ is the restriction to $\R^2$ of an entire function on $\C^2$ and does not vanish identically;
so $\hat v$ is a real analytic function on $\R^2$ that is not identically zero, and its zero set is a Lebesgue null set.
Therefore $L_l(p_\perp)\in\C e_1$ for every $p_\perp$, and $L_l(p_\perp)\neq0$ outside a Lebesgue null set of $p_\perp$.

\subsubsection{Eight laws}\label{sec:eight-laws}

\begin{proposition}[Sharpness]\label{prop:sharpness}
\begin{enumerate}
\item\label{it:sharp-constant}
\emph{The constant one.}
For every $c\in(0,1)$ there is a law $\mu^{(c)}$ on $\mathcal{S}'$ with \textup{(H)} and \textup{(dF)} whose joint spectral measures satisfy
\[
\mathfrak{E}_j\left(\{\omega<c\cdot|p_i|\}\right)=0,
\qquad\text{that is,}\qquad
\mathsf{H}_j\ge c\cdot|\mathsf{P}_i|,
\]
for every ordered pair $(j,i)$ of distinct elements of $\{1,2,3\}$,
and which is not invariant under the rotation $R^{12}_{\pi/4}$;
in particular $\mu^{(c)}$ is not invariant under $\mathrm{O}(3)$.
\item\label{it:sharp-quarter-turn}
\emph{The quarter turn.}
There is a Gaussian law with \textup{(G)}, \textup{(dF)} and $(T_{123})$, with $(R_j)$ and $(RP_j)$ for $j=1,2,3$, and with $(C_{12})$ and $(C_{13})$,
which is not invariant under the quarter turn $R^{12}_{\pi/2}$ of the plane $(1,2)$.
Consequently the hypothesis $(Q_{ji})$ cannot be omitted from Theorem~\ref{thm:one-plane},
and the requirement that \eqref{it:quarter-turn-two-point} or \eqref{it:second-mirror-two-point} hold
cannot be omitted from Theorem~\ref{thm:two-point}\eqref{it:two-point-invariance}.
\item\label{it:sharp-domination}
\emph{Gaussian domination.}
For every $\varepsilon\in(0,1]$ there is a Gaussian law $\gamma^{\mathrm{f}}_\varepsilon$
which satisfies \textup{(dF)} and $(RP_j)$ for $j=1,2,3$, is invariant under $\R^3\rtimes\Ocub$,
and whose joint spectral measures lie in the round cone with a mass gap:
for $j=1,2,3$ the measure $\mathfrak{E}_j$ is carried by
\[
\{(0,0)\}\cup\left\{(\omega,p_\perp):\ \omega\ge\left(1+|p_\perp|^2\right)^{\frac12}\right\}.
\]
The law $\gamma^{\mathrm{f}}_\varepsilon$ is not invariant under $\mathrm{O}(3)$, and it does not satisfy \textup{(G)}.
\item\label{it:sharp-closedness-two-point}
\emph{Closedness, for the two-point function.}
For every $m\ge0$ there is a Gaussian law $\gamma^{\oplus}_m$ with \textup{(H)} and with $(C_{ji})$ for every ordered pair $(j,i)$ of distinct indices,
which does not satisfy \textup{(dF)} and is not invariant under $\mathrm{O}(3)$.
Let $S_2$ be its two-point function (Lemma~\ref{lem:domination-consequences}\eqref{it:moment-forms}),
let $\chi\in C^\infty_{\mathrm c}(\R;\R)$ be a profile with $\chi\neq0$, and put $\lambda:=(+,\chi)$.
There is a smooth function $G_{\lambda\lambda}$ on $\R^2\setminus\{0\}$ with
\[
S_2(\alpha\otimes\lambda,\alpha'\otimes\lambda)=\int\mathrm{d}y\int\mathrm{d}y'\,\alpha(y)\cdot\alpha'(y')\cdot G_{\lambda\lambda}(y'-y)
\]
for all $\alpha,\alpha'\in L^2(\R^2;\C)$ with disjoint compact supports.
The harmonics of this kernel,
\[
a^k_{\lambda\lambda}(r):=\frac{1}{2\pi}\int_0^{2\pi}\mathrm{d}\vartheta\,G_{\lambda\lambda}(rn_\vartheta)\cdot\exp(-ik\vartheta)\qquad(r>0,\ k\in\Z),
\]
satisfy $\left|a^{-2}_{\lambda\lambda}\right|=\left|a^{2}_{\lambda\lambda}\right|$ on $(0,\infty)$,
and the harmonic $a^{-2}_{\lambda\lambda}$ of the anomalous frequency $-2$ does not vanish identically.
The anomaly index of $\gamma^{\oplus}_m$ (Definition~\ref{def:anomaly-index}) satisfies $\mathfrak{a}(2)=1$, and $\mathfrak{a}(2n)=n$ for every integer $n\ge1$.
\item\label{it:sharp-contact-terms}
\emph{Closedness, for the contact terms.}
Let $D$ be a triple of pairwise distinct numbers in $(0,1]$,
and let $m_1$ be the law of \cite[Example~8.26]{NoncompactQED} formed with $D$, the mean of six Gaussian white noises whose variances are the permutations of $D$.
The law $m_1$ satisfies \textup{(H)},
and it satisfies $(C_{ji})$ for every ordered pair $(j,i)$, because $\mathsf{H}_j=0$ and $\mathsf{P}_i=0$ for it.
Its two-point function is invariant under $\mathrm{O}(3)$,
and $S_N(b_1,\dots,b_N)=0$ for every $N\ge1$ whenever $b_1,\dots,b_N\in C^\infty_{\mathrm c}(\R^3;\R^3)$ have pairwise disjoint supports.
The law $m_1$ is not invariant under $\mathrm{O}(3)$, and it does not satisfy \textup{(dF)}.
\item\label{it:sharp-soft-properties}
\emph{The cone is not a consequence of \textup{(H)}, \textup{(dF)} and evenness, even with saturation at short distances.}
For all $\kappa>0$ and $m>0$ there is a Gaussian law $\gamma^{\mathrm{d}}_{\kappa,m}$ with \textup{(H)} and \textup{(dF)}, invariant under $F\mapsto-F$,
which is the law $\gamma_M$ of Lemma~\ref{lem:multiplier-laws} for a multiplier $M$ with
\[
\lim_{s\to0}M(p/s)=P^\perp(p)\qquad\text{for every }p\in\R^3\text{ with }p_1p_2p_3\neq0,
\]
and for which $\mathfrak{E}_1\left(\{\omega<m,\ |p_2|>C\}\right)\neq0$ for every $C\in\R$;
consequently there are no $c>0$ and $C\in\R$ with $\mathsf{H}_1\ge c\cdot|\mathsf{P}_2|-C$.
The law $\gamma^{\mathrm{d}}_{\kappa,m}$ is not invariant under $\mathrm{O}(3)$.
\item\label{it:sharp-planes}
\emph{Rotations of one plane only; no plate bound without the cone.}
\begin{enumerate}
\item\label{it:sharp-stack}
There is a Gaussian law $\gamma^{\mathrm{s}}_3$ with \textup{(G)}, \textup{(dF)}, $(T_{123})$, $(R_1)$, $(RP_1)$, $(Q_{12})$ and $(C_{12})$,
the hypotheses of Theorem~\ref{thm:one-plane} for the pair $(1,2)$;
it is therefore invariant under all rotations $R^{12}_t$, $t\in\R$, as that theorem asserts.
It does not satisfy $(C_{13})$, and it is not invariant under $R^{13}_{\pi/2}$.
\item\label{it:sharp-three-stacks}
There is a Gaussian law $\bar\gamma^{\mathrm{s}}$ with \textup{(H)} and \textup{(dF)} for which the bound of Proposition~\ref{prop:point-plates} fails:
for every $\chi\in L^2(\R;\R)$ with $\chi\neq0$ and $\lambda:=(3,\chi)$,
\[
\sup\frac{\left|S_2(\alpha\otimes\lambda,\alpha'\otimes\lambda)\right|}{\|\alpha\|_{L^1(\R^2)}\cdot\|\alpha'\|_{L^1(\R^2)}}=\infty,
\]
where $S_2$ is the two-point function of $\bar\gamma^{\mathrm{s}}$,
and the supremum is taken over all non-zero $\alpha,\alpha'\in L^1\cap L^2(\R^2;\C)$
with $\supp\alpha\subset\{-2<y^0<-1\}$ and $\supp\alpha'\subset\{1<y^0<2\}$.
\end{enumerate}
\end{enumerate}
\end{proposition}

\paragraph{Stretched Maxwell laws}
Let $\gamma$ be the free Maxwell law of Definition~\ref{def:free-maxwell}.
It is the law on $\mathcal{S}'$ with
\[
\gamma\left(\exp\left(i\mathbb{B}(b)\right)\right)=\exp\left(-\frac12\left\|P^\perp b\right\|^2\right)\qquad\left(b\in\mathcal{S}(\R^3;\R^3)\right),
\]
so that under $\gamma$ every $\mathbb{B}(b)$, $b\in\mathcal{S}(\R^3;\R^3)$, is a centred normal random variable with the variance $\|P^\perp b\|^2$;
it satisfies \textup{(G)}, is invariant under all Euclidean maps, is carried by the closed distributions, and satisfies $(RP_j)$ for $j=1,2,3$
\cite[Lemma~8.30(i)]{NoncompactQED}
(the companion paper writes $\Pi$ for $P^\perp$ and $X(f)$ for $F(f)=\mathbb{B}(\ast f)$).

For a triple $a=(a_1,a_2,a_3)\in(0,\infty)^3$ let $A:=\diag(a_1,a_2,a_3)$,
and for a vector field $b$ on $\R^3$ let
\[
b^a(x):=A^{-1}b\left(A^{-1}x\right)\qquad\left(x\in\R^3\right).
\]
With $a':=\left(a_1^{-1},a_2^{-1},a_3^{-1}\right)$ one has $(b^a)^{a'}=b=(b^{a'})^a$.
The map $b\mapsto b^a$ is linear, and it is continuous on $\mathcal{S}(\R^3;\R^3)$, like its inverse $b\mapsto b^{a'}$,
since each of them is a linear change of variables followed by the multiplication with a constant matrix;
so $b\mapsto b^a$ is a linear homeomorphism of $\mathcal{S}(\R^3;\R^3)$ onto itself.
Because $\supp b^a=A\,\supp b$ and $A$ maps each half space $\{x_j>0\}$ onto itself,
the map $b\mapsto b^a$ also maps $C^\infty_{\mathrm c}(\{x_j>0\};\R^3)$ onto itself, $j=1,2,3$.
For $a\in[\tfrac14,\tfrac12]^3$ the field $b^a$ is the field denoted $\Phi_{u\#}b$ in \cite[(8.8)]{NoncompactQED}, for the constant triple $u=a$.

For $f\in\mathcal{S}(\R^3;\Lambda^2)$ let $f^a$ be the two-form with $\ast f^a=(\ast f)^a$,
and for $F\in\mathcal{S}'$ let
\[
(T_aF)(f):=F(f^a)\qquad\left(f\in\mathcal{S}(\R^3;\Lambda^2)\right).
\]
As $f\mapsto f^a$ is a continuous linear bijection of $\mathcal{S}(\R^3;\Lambda^2)$ with the inverse $f\mapsto f^{a'}$,
one has $T_aF\in\mathcal{S}'$, and $T_a$ is a bijection of $\mathcal{S}'$ onto itself with the inverse $T_{a'}$.
Both maps are measurable:
the $\sigma$-algebra $\Sigma$ of $\mathcal{S}'$ is generated by the maps $F\mapsto F(f)$, and $F\mapsto(T_aF)(f)=F(f^a)$ is one of them.
By construction $\mathbb{B}(b)\circ T_a=\mathbb{B}(b^a)$ for real $b\in\mathcal{S}(\R^3;\R^3)$, and by complex linearity also for complex $b$.
The \emph{stretched Maxwell law} with the triple $a$ is the image
\[
\gamma^a:=(T_a)_*\gamma .
\]
By the transformation rule for image measures,
$\gamma^a(\Phi)=\gamma(\Phi\circ T_a)$ for every $\Sigma$-measurable function $\Phi$ on $\mathcal{S}'$ that takes its values in $[0,\infty]$ or is $\gamma^a$-integrable;
in the second case $\Phi\circ T_a$ is $\gamma$-integrable.
In particular
\begin{equation}\label{eq:rotation-stretched-law}
\gamma^a\left(\exp\left(i\mathbb{B}(b)\right)\right)=\gamma\left(\exp\left(i\mathbb{B}(b^a)\right)\right)=\exp\left(-\frac12\left\|P^\perp b^a\right\|^2\right)
\qquad\left(b\in\mathcal{S}(\R^3;\R^3)\right).
\end{equation}
For a law $\nu$ with $(T_{123})$, $(R_j)$ and $(RP_j)$, the reconstruction space, the Hamiltonian and the joint spectral measure of $\nu$ in the direction $e_j$
are written $\mathcal{H}_j(\nu)$, $\mathsf{H}_j(\nu)$ and $\mathfrak{E}_j(\nu)$ when the law has to be named (Lemma~\ref{lem:joint-spectrum}).

\begin{lemma}[Stretched Maxwell laws]\label{lem:stretched-laws}
Let $a\in(0,\infty)^3$ and $A=\diag(a_1,a_2,a_3)$,
and let $M^a$ be the function on $\R^3$ with $M^a(0):=0$ and
\[
M^a(p):=\det(A)\cdot\left(A^{-2}-\frac{A^{-2}pp^{T}A^{-2}}{p^{T}A^{-2}p}\right)\qquad(p\neq0).
\]
\begin{enumerate}
\item\label{it:stretch-multiplier}
One has $\gamma^a=\gamma_{M^a}$.
\item\label{it:stretch-actions}
Let $b$ be a vector field and $u$ a differentiable function on $\R^3$, and let $v\in\R^3$. Then
\[
(\tau_v\cdot b)^a=\tau_{Av}\cdot b^a,\qquad(\nabla u)^a=\nabla\left(u\circ A^{-1}\right),
\]
and $(g\cdot b)^a=g\cdot b^a$ for every $g\in\mathrm{O}(3)$ with $gA=Ag$; in particular $(\theta_j\cdot b)^a=\theta_j\cdot b^a$ for $j=1,2,3$.
If $g\in\Ocub$ and a permutation $\mathfrak{s}\in\mathfrak{S}_3$ satisfy $ge_k=\pm e_{\mathfrak{s}(k)}$ for $k=1,2,3$, then
\[
(g\cdot b)^a=g\cdot b^{a\circ\mathfrak{s}},\qquad\text{where}\quad a\circ\mathfrak{s}:=\left(a_{\mathfrak{s}(1)},a_{\mathfrak{s}(2)},a_{\mathfrak{s}(3)}\right).
\]
\item\label{it:stretch-properties}
The law $\gamma^a$ satisfies $(T_{123})$, $(R_j)$ and $(RP_j)$ for $j=1,2,3$, and \textup{(dF)}, and it is invariant under $F\mapsto-F$.
If
\[
a_i\cdot a_k\le a_l\qquad\text{whenever }\{i,k,l\}=\{1,2,3\},
\]
then $\gamma^a$ satisfies \textup{(G)}.
\item\label{it:stretch-cone}
Let $j\neq i$.
The joint spectral measure $\mathfrak{E}_j(\gamma^a)$ is carried by $\{\omega\ge(a_j/a_i)\cdot|p_i|\}$,
and the constant $a_j/a_i$ is optimal:
$\mathfrak{E}_j(\gamma^a)\left(\{\omega<c\cdot|p_i|\}\right)\neq0$ for every $c>a_j/a_i$.
\item\label{it:stretch-permutation}
If $g\in\Ocub$ and $\mathfrak{s}\in\mathfrak{S}_3$ satisfy $ge_k=\pm e_{\mathfrak{s}(k)}$ for $k=1,2,3$, then $(g^{-1})_*\gamma^a=\gamma^{a\circ\mathfrak{s}}$.
\item\label{it:stretch-trace}
With $y_l:=a_l^{-2}$ for $l=1,2,3$,
\[
\operatorname{tr}M^a(p)=\det(A)\cdot\left(\sum_{l=1}^3y_l-\frac{\sum_{l=1}^3y_l^2\cdot p_l^2}{\sum_{l=1}^3y_l\cdot p_l^2}\right)\qquad\left(p\in\R^3\setminus\{0\}\right).
\]
\end{enumerate}
\end{lemma}

\begin{proof}
\eqref{it:stretch-multiplier}
For $p\neq0$ the matrix $P^\perp(A^{-1}p)=1-A^{-1}pp^{T}A^{-1}/(p^{T}A^{-2}p)$ is an orthogonal projection, and
\[
M^a(p)=\det(A)\cdot A^{-1}P^\perp\left(A^{-1}p\right)A^{-1}.
\]
Hence $M^a$ is a measurable even function with values in the real symmetric matrices, and $0\le M^a(p)\le\det(A)\cdot A^{-2}$;
so $M^a$ is a multiplier as in Lemma~\ref{lem:multiplier-laws}, and $\gamma_{M^a}$ is the centred Gaussian law with the covariance $C_{M^a}$.
Let $b\in\mathcal{S}(\R^3;\R^3)$.
The substitution $x=Az$ and the identity $p\cdot Az=(Ap)\cdot z$ give
\[
\widehat{b^a}(p)=\int\mathrm{d}x\,\exp(-ip\cdot x)\cdot A^{-1}b\left(A^{-1}x\right)=\det(A)\cdot A^{-1}\hat b(Ap)\qquad\left(p\in\R^3\right).
\]
Plancherel's theorem, the definition of $P^\perp$ and the substitution $p\mapsto A^{-1}p$, whose Jacobian determinant is $\det(A)^{-1}$, yield
\[
\begin{aligned}
\left\|P^\perp b^a\right\|^2
&=(2\pi)^{-3}\int\mathrm{d}p\,\det(A)^2\cdot\hat b(Ap)^{\dagger}A^{-1}P^\perp(p)A^{-1}\hat b(Ap)\\
&=(2\pi)^{-3}\int\mathrm{d}p\,\hat b(p)^{\dagger}\left(\det(A)\cdot A^{-1}P^\perp\left(A^{-1}p\right)A^{-1}\right)\hat b(p)
=C_{M^a}(b,b).
\end{aligned}
\]
By \eqref{eq:rotation-stretched-law} the laws $\gamma^a$ and $\gamma_{M^a}$ have the same characteristic functional
$b\mapsto\exp\left(-\frac12C_{M^a}(b,b)\right)$ on $\mathcal{S}(\R^3;\R^3)$.
A law on $\mathcal{S}'$ is determined by its characteristic functional \cite[Lemma~8.12(ii)]{NoncompactQED}; hence $\gamma^a=\gamma_{M^a}$.

\eqref{it:stretch-actions}
Let $x\in\R^3$.
For the translation,
\[
(\tau_v\cdot b)^a(x)=A^{-1}b\left(A^{-1}x-v\right)=A^{-1}b\left(A^{-1}(x-Av)\right)=b^a(x-Av)=(\tau_{Av}\cdot b^a)(x).
\]
The chain rule and the symmetry of $A$ give $\nabla(u\circ A^{-1})(x)=A^{-1}(\nabla u)(A^{-1}x)=(\nabla u)^a(x)$.
Let $g\in\mathrm{O}(3)$ commute with $A$; then $g$ commutes with $A^{-1}$,
and the action $(g\cdot b)(x)=\det(g)\cdot g\,b(g^{-1}x)$ of Definition~\ref{def:field-actions} gives
\[
(g\cdot b)^a(x)=\det(g)\cdot A^{-1}g\,b\left(g^{-1}A^{-1}x\right)=\det(g)\cdot g\,A^{-1}b\left(A^{-1}g^{-1}x\right)=(g\cdot b^a)(x).
\]
The reflections $\theta_j$ are diagonal matrices and commute with $A$.
Let finally $g\in\Ocub$ and $\mathfrak{s}$ be as in the statement, and put $A_{\mathfrak{s}}:=\diag(a\circ\mathfrak{s})$.
For $k=1,2,3$ one has $g^{-1}e_{\mathfrak{s}(k)}=\pm e_k$, with the sign of $ge_k=\pm e_{\mathfrak{s}(k)}$; therefore
\[
g^{-1}Ag\,e_k=\pm g^{-1}Ae_{\mathfrak{s}(k)}=\pm a_{\mathfrak{s}(k)}\cdot g^{-1}e_{\mathfrak{s}(k)}=a_{\mathfrak{s}(k)}\cdot e_k .
\]
So $g^{-1}Ag=A_{\mathfrak{s}}$, that is, $A^{-1}g=g\,A_{\mathfrak{s}}^{-1}$ and $g^{-1}A^{-1}=A_{\mathfrak{s}}^{-1}g^{-1}$, and
\[
\begin{aligned}
(g\cdot b)^a(x)&=\det(g)\cdot A^{-1}g\,b\left(g^{-1}A^{-1}x\right)\\
&=\det(g)\cdot g\,A_{\mathfrak{s}}^{-1}b\left(A_{\mathfrak{s}}^{-1}g^{-1}x\right)=\left(g\cdot b^{a\circ\mathfrak{s}}\right)(x).
\end{aligned}
\]

\eqref{it:stretch-properties}
Let $\mu$ be a law on $\mathcal{S}'$ and $\mathfrak{g}$ a Euclidean map.
The operator $\hat{\mathfrak{g}}$ is the composition with the map $F\mapsto\mathfrak{g}^{-1}F$ (Definition~\ref{def:laws});
so it commutes with the composition with functions, $\hat{\mathfrak{g}}\exp(i\Phi)=\exp\left(i\,\hat{\mathfrak{g}}\Phi\right)$ for every function $\Phi$ on $\mathcal{S}'$.
Together with $\hat{\mathfrak{g}}\,\mathbb{B}(b)=\mathbb{B}(\mathfrak{g}\cdot b)$ (Definition~\ref{def:field-actions}) this gives
\[
\left((\mathfrak{g}^{-1})_*\mu\right)\left(\exp\left(i\mathbb{B}(b)\right)\right)
=\mu\left(\hat{\mathfrak{g}}\exp\left(i\mathbb{B}(b)\right)\right)
=\mu\left(\exp\left(i\mathbb{B}(\mathfrak{g}\cdot b)\right)\right)
\]
for real $b\in\mathcal{S}(\R^3;\R^3)$.
By \cite[Lemma~8.12(ii)]{NoncompactQED}, $\mu$ is therefore invariant under $\mathfrak{g}$
as soon as $\mu(\exp(i\mathbb{B}(\mathfrak{g}\cdot b)))=\mu(\exp(i\mathbb{B}(b)))$ for every such $b$.
Now let $b\in\mathcal{S}(\R^3;\R^3)$, $v\in\R^3$ and $j\in\{1,2,3\}$.
By \eqref{eq:rotation-stretched-law}, by \eqref{it:stretch-actions} and by the invariance of $\gamma$ under $\tau_{Av}$ and $\theta_j$,
\[
\begin{aligned}
\gamma^a\left(\exp\left(i\mathbb{B}(\tau_v\cdot b)\right)\right)
&=\gamma\left(\exp\left(i\mathbb{B}(\tau_{Av}\cdot b^a)\right)\right)
=\gamma\left(\exp\left(i\mathbb{B}(b^a)\right)\right)
=\gamma^a\left(\exp\left(i\mathbb{B}(b)\right)\right),\\
\gamma^a\left(\exp\left(i\mathbb{B}(\theta_j\cdot b)\right)\right)
&=\gamma\left(\exp\left(i\mathbb{B}(\theta_j\cdot b^a)\right)\right)
=\gamma\left(\exp\left(i\mathbb{B}(b^a)\right)\right)
=\gamma^a\left(\exp\left(i\mathbb{B}(b)\right)\right);
\end{aligned}
\]
so $\gamma^a$ satisfies $(T_{123})$ and $(R_j)$.
For reflection positivity let $n\ge1$, $z_1,\dots,z_n\in\C$ and $b_1,\dots,b_n\in C^\infty_{\mathrm c}(\{x_j>0\};\R^3)$.
Since $\mathbb{B}(b)\circ T_a=\mathbb{B}(b^a)$ and $(\theta_j\cdot b_k)^a=\theta_j\cdot b_k^a$,
\[
\begin{aligned}
&\sum_{k,k'=1}^n\overline{z_k}\cdot z_{k'}\cdot\gamma^a\left(\exp\left(i\left(\mathbb{B}(b_{k'})-\mathbb{B}(\theta_j\cdot b_k)\right)\right)\right)\\
&\qquad=\sum_{k,k'=1}^n\overline{z_k}\cdot z_{k'}\cdot\gamma\left(\exp\left(i\left(\mathbb{B}(b_{k'}^a)-\mathbb{B}(\theta_j\cdot b_k^a)\right)\right)\right)\ge0,
\end{aligned}
\]
because $b_1^a,\dots,b_n^a$ lie in $C^\infty_{\mathrm c}(\{x_j>0\};\R^3)$ and $\gamma$ satisfies $(RP_j)$.
This is $(RP_j)$ for $\gamma^a$.

Let $u\in\mathcal{S}(\R^3)$ be real.
Then $u\circ A^{-1}\in\mathcal{S}(\R^3)$, and the Fourier transform of $\nabla(u\circ A^{-1})$ is the function $p\mapsto i\,\widehat{u\circ A^{-1}}(p)\cdot p$,
which is annihilated by $P^\perp(p)$; so $P^\perp\nabla(u\circ A^{-1})=0$, and by \eqref{it:stretch-actions}
\[
\gamma^a\left(\mathbb{B}(\nabla u)^2\right)=\gamma\left(\mathbb{B}\left(\nabla(u\circ A^{-1})\right)^2\right)=\left\|P^\perp\nabla\left(u\circ A^{-1}\right)\right\|^2=0.
\]
A law with this property for every real $u\in\mathcal{S}(\R^3)$ is carried by the closed distributions \cite[Section~8.1.6]{NoncompactQED};
thus $\gamma^a$ satisfies \textup{(dF)}.
The right side of \eqref{eq:rotation-stretched-law} does not change when $b$ is replaced by $-b$,
and the image of $\gamma^a$ under $F\mapsto-F$ has the characteristic functional $b\mapsto\gamma^a\left(\exp\left(i\mathbb{B}(-b)\right)\right)$;
by \cite[Lemma~8.12(ii)]{NoncompactQED} this image is $\gamma^a$.

For Gaussian domination let $b\in\mathcal{S}(\R^3;\R^3)$.
As $\mathbb{B}(b^a)$ is under $\gamma$ a centred normal random variable with the variance $\|P^\perp b^a\|^2\le\|b^a\|^2$,
\[
\gamma^a\left(\exp\mathbb{B}(b)\right)=\gamma\left(\exp\mathbb{B}(b^a)\right)
=\exp\left(\frac12\left\|P^\perp b^a\right\|^2\right)\le\exp\left(\frac12\left\|b^a\right\|^2\right).
\]
The substitution $x=Az$ gives
\[
\left\|b^a\right\|^2=\det(A)\cdot\int\mathrm{d}z\,\left|A^{-1}b(z)\right|^2=\sum_{l=1}^3\frac{a_1\cdot a_2\cdot a_3}{a_l^2}\cdot\int\mathrm{d}z\,\left(e_l\cdot b(z)\right)^2,
\]
and $a_1\cdot a_2\cdot a_3/a_l^2=a_i\cdot a_k/a_l$ when $\{i,k,l\}=\{1,2,3\}$.
Under the condition of the statement each of these three coefficients is at most one;
so $\|b^a\|^2\le\|b\|^2$ and $\gamma^a(\exp\mathbb{B}(b))\le\exp\left(\frac12\|b\|^2\right)$, which is \textup{(G)}.

\eqref{it:stretch-cone}
For a function $\Phi$ on $\mathcal{S}'$ put $J_a\Phi:=\Phi\circ T_a$.
As $\gamma(|J_a\Phi|^2)=\gamma^a(|\Phi|^2)$ and $\Phi''\mapsto\Phi''\circ T_{a'}$ is the inverse map,
$J_a$ is a unitary operator of $L^2(\gamma^a)$ onto $L^2(\gamma)$,
and $\gamma(J_a\Phi'')=\gamma^a(\Phi'')$ for $\Phi''\in L^1(\gamma^a)$.
Being a composition operator, $J_a$ is multiplicative and commutes with the composition with functions:
$J_a(\Phi\cdot\Phi')=(J_a\Phi)\cdot(J_a\Phi')$ and $J_a\exp(i\Phi)=\exp\left(i\,J_a\Phi\right)$ for functions $\Phi,\Phi'$ on $\mathcal{S}'$.
Moreover $J_a\mathbb{B}(b)=\mathbb{B}(b^a)$, hence $J_a\exp\left(i\mathbb{B}(b)\right)=\exp\left(i\mathbb{B}(b^a)\right)$.
For real $b\in\mathcal{S}(\R^3;\R^3)$ and $v\in\R^3$ one has
$U(v)\exp\left(i\mathbb{B}(b)\right)=\exp\left(i\mathbb{B}(\tau_v\cdot b)\right)$,
since $U(v)=\hat\tau_v$ commutes with the composition with functions and maps $\mathbb{B}(b)$ to $\mathbb{B}(\tau_v\cdot b)$ (Definition~\ref{def:field-actions}),
and $\Theta_j\exp\left(i\mathbb{B}(b)\right)=\exp\left(-i\mathbb{B}(\theta_j\cdot b)\right)$ by \eqref{eq:rotation-conjugation}.
Together with \eqref{it:stretch-actions} this gives
\[
\begin{aligned}
J_aU(v)\exp\left(i\mathbb{B}(b)\right)&=\exp\left(i\mathbb{B}(\tau_{Av}\cdot b^a)\right)=U(Av)J_a\exp\left(i\mathbb{B}(b)\right),\\
J_a\Theta_j\exp\left(i\mathbb{B}(b)\right)&=\exp\left(-i\mathbb{B}(\theta_j\cdot b^a)\right)=\Theta_jJ_a\exp\left(i\mathbb{B}(b)\right).
\end{aligned}
\]
The finite linear combinations of the functions $\exp(i\mathbb{B}(b))$, $b\in\mathcal{S}(\R^3;\R^3)$, are dense in $L^2(\gamma^a)$ \cite[Lemma~8.12(i)]{NoncompactQED},
the operators $U(v)$ are unitary and $\Theta_j$ is antilinear and isometric on $L^2(\gamma^a)$ and on $L^2(\gamma)$,
by \eqref{it:stretch-properties} and by the invariance of $\gamma$ under all Euclidean maps;
consequently
\[
J_aU(v)=U(Av)J_a,\qquad J_a\Theta_j=\Theta_jJ_a .
\]
The $\sigma$-algebra $\Sigma_j^+$ is generated by the functions $\mathbb{B}(b)$, $b\in C^\infty_{\mathrm c}(\{x_j>0\};\R^3)$,
and $b\mapsto b^a$, $b\mapsto b^{a'}$ map this space onto itself;
as $\mathbb{B}(b)\circ T_a=\mathbb{B}(b^a)$ and $\mathbb{B}(b)\circ T_{a'}=\mathbb{B}(b^{a'})$, both $T_a$ and $T_{a'}$ are measurable for $\Sigma_j^+$.
Hence $J_a$ maps $L^2(\Sigma_j^+,\gamma^a)$ onto $L^2(\Sigma_j^+,\gamma)$.

The scalar product of the reconstruction of a law $\nu$ in the direction $e_j$ is
$\langle[\Phi]_j,[\Phi']_j\rangle=\nu\left((\Theta_j\Phi)\cdot\Phi'\right)$ for $\Phi,\Phi'\in L^2(\Sigma_j^+,\nu)$,
and these classes are dense in $\mathcal{H}_j(\nu)$ \cite[Proposition~8.19(i)]{NoncompactQED}
(there the scalar product is written $\langle\hat\theta_j\Phi,\Phi'\rangle_\nu$ with the linear operator $\hat\theta_j$, and $\exp(-t\cdot\mathsf{H}_j)$ is written $T(t)$).
For $\Phi,\Phi'\in L^2(\Sigma_j^+,\gamma^a)$ the properties of $J_a$ give
\[
\begin{aligned}
\left\langle[J_a\Phi]_j,[J_a\Phi']_j\right\rangle
&=\gamma\left((\Theta_jJ_a\Phi)\cdot J_a\Phi'\right)
=\gamma\left(J_a\left((\Theta_j\Phi)\cdot\Phi'\right)\right)\\
&=\gamma^a\left((\Theta_j\Phi)\cdot\Phi'\right)
=\left\langle[\Phi]_j,[\Phi']_j\right\rangle .
\end{aligned}
\]
So $[\Phi]_j\mapsto[J_a\Phi]_j$ is well defined and isometric, and it maps a dense subspace of $\mathcal{H}_j(\gamma^a)$ onto a dense subspace of $\mathcal{H}_j(\gamma)$;
it extends to a unitary operator $\tilde J_a$ of $\mathcal{H}_j(\gamma^a)$ onto $\mathcal{H}_j(\gamma)$.
Let $t\ge0$ and $y\in e_j^\perp$.
By \cite[Proposition~8.19(ii),(iii)]{NoncompactQED},
$\exp\left(-t\cdot\mathsf{H}_j(\nu)\right)V_j(y)[\Phi]_j=[U(t\cdot e_j+y)\Phi]_j$ for $\Phi\in L^2(\Sigma_j^+,\nu)$, $\nu=\gamma^a,\gamma$.
As $A(t\cdot e_j+y)=a_j\cdot t\cdot e_j+Ay$ with $Ay\in e_j^\perp$, it follows that
\[
\begin{aligned}
\tilde J_a\exp\left(-t\cdot\mathsf{H}_j(\gamma^a)\right)V_j(y)[\Phi]_j
&=\left[J_aU(t\cdot e_j+y)\Phi\right]_j=\left[U(a_j\cdot t\cdot e_j+Ay)J_a\Phi\right]_j\\
&=\exp\left(-a_j\cdot t\cdot\mathsf{H}_j(\gamma)\right)V_j(Ay)\,\tilde J_a[\Phi]_j ,
\end{aligned}
\]
where $V_j(y)$ acts on $\mathcal{H}_j(\gamma^a)$ and $V_j(Ay)$ on $\mathcal{H}_j(\gamma)$.
By density,
$\tilde J_a\exp\left(-t\cdot\mathsf{H}_j(\gamma^a)\right)V_j(y)\,\tilde J_a^{-1}=\exp\left(-a_j\cdot t\cdot\mathsf{H}_j(\gamma)\right)V_j(Ay)$ for all $t\ge0$ and $y\in e_j^\perp$.
For real $s$ let $h_s$ be the function on $[0,1]$ with $h_s(0):=0$ and $h_s\left(\exp(-\omega)\right):=\exp(is\cdot\omega)$ for $\omega\ge0$
(every point of $(0,1]$ is $\exp(-\omega)$ for exactly one $\omega\ge0$); it is a bounded Borel function.
For a self-adjoint operator $\mathsf{H}\ge0$ the operator $\exp(-\mathsf{H})$ is injective, and $h_s\left(\exp(-\mathsf{H})\right)=\exp(is\cdot\mathsf{H})$.
Conjugation with a unitary operator commutes with the bounded Borel functional calculus;
applied to the case $t=1$, $y=0$, this gives $\tilde J_a\exp\left(is\cdot\mathsf{H}_j(\gamma^a)\right)\tilde J_a^{-1}=\exp\left(ia_j\cdot s\cdot\mathsf{H}_j(\gamma)\right)$,
and together with the case $t=0$,
\[
\begin{aligned}
\tilde J_a\exp\left(is\cdot\mathsf{H}_j(\gamma^a)\right)V_j(y)\,\tilde J_a^{-1}
&=\exp\left(ia_j\cdot s\cdot\mathsf{H}_j(\gamma)\right)V_j(Ay)\\
&=\int\mathrm{d}\mathfrak{E}_j(\gamma)(\omega,p_\perp)\,\exp\left(is\cdot a_j\cdot\omega+i(Ap_\perp)\cdot y\right)
\end{aligned}
\]
for all $s\in\R$ and $y\in e_j^\perp$, by Lemma~\ref{lem:joint-spectrum}(b) and the identity $p_\perp\cdot Ay=(Ap_\perp)\cdot y$.
So the projection-valued measure
$B\mapsto\tilde J_a^{-1}\,\mathfrak{E}_j(\gamma)\left(\left\{(\omega,p_\perp):\ (a_j\cdot\omega,Ap_\perp)\in B\right\}\right)\tilde J_a$
on the Borel sets of $\R\times e_j^\perp$ has the Fourier transform $(s,y)\mapsto\exp\left(is\cdot\mathsf{H}_j(\gamma^a)\right)V_j(y)$,
and the uniqueness in Lemma~\ref{lem:joint-spectrum}(b) gives
\begin{equation}\label{eq:rotation-stretched-spectrum}
\tilde J_a\,\mathfrak{E}_j(\gamma^a)(B)\,\tilde J_a^{-1}=\mathfrak{E}_j(\gamma)\left(\left\{(\omega,p_\perp):\ (a_j\cdot\omega,Ap_\perp)\in B\right\}\right)
\end{equation}
for every Borel set $B$ of $\R\times e_j^\perp$.

The free Maxwell law $\gamma$ is the law $\gamma_1$ of Definition~\ref{def:scaled-maxwell}, scaled with the factor one;
so Lemma~\ref{lem:scaled-maxwell}, applied with the factor one, shows that $\gamma$ lies in $\mathcal{C}_9$.
In particular $\gamma$ is invariant under $\R^3$ and under $\theta_j$, satisfies $(RP_j)$, and is invariant and reflection positive at the plane $\{x_j+x_i=0\}$.
By Theorem~\ref{thm:cone}(c) the joint spectral measure of $(\mathsf{H}_j(\gamma),\mathsf{P}_i)$ is carried by $\{\omega\ge|p_i|\}$,
that is, $\mathfrak{E}_j(\gamma)\left(\{\omega<|p_i|\}\right)=0$.
For $B:=\{\omega<(a_j/a_i)\cdot|p_i|\}$ the set on the right of \eqref{eq:rotation-stretched-spectrum} is
$\{(\omega,p_\perp):\ a_j\cdot\omega<a_j\cdot|p_i|\}=\{\omega<|p_i|\}$;
so $\mathfrak{E}_j(\gamma^a)(B)=0$.

For the optimality choose a two-form $f$ as in Lemma~\ref{lem:maxwell-spectrum}(c),
let $k_0>0$ be the number given there, and let $0<r<k_0$.
By Lemma~\ref{lem:maxwell-spectrum}(d2) the point $(r,r)$ lies in the joint spectrum of $(\mathsf{H}_j(\gamma),\mathsf{P}_i)$,
the support of the image of $\mathfrak{E}_j(\gamma)$ under $(\omega,p_\perp)\mapsto(\omega,p_i)$ (Lemma~\ref{lem:joint-spectrum}(c)).
By \eqref{eq:rotation-stretched-spectrum} the image of $\mathfrak{E}_j(\gamma^a)$ under $(\omega,p_\perp)\mapsto(\omega,p_i)$
is unitarily equivalent to the image of $\mathfrak{E}_j(\gamma)$ under $(\omega,p_\perp)\mapsto(a_j\cdot\omega,a_i\cdot p_i)$.
As $(\omega,p_i)\mapsto(a_j\cdot\omega,a_i\cdot p_i)$ is a homeomorphism of $[0,\infty)\times\R$ onto itself,
the point $(a_j\cdot r,a_i\cdot r)$ lies in the joint spectrum of $(\mathsf{H}_j(\gamma^a),\mathsf{P}_i)$.
Let $c>a_j/a_i$.
The open set $\{(\omega,p_i):\ \omega<c\cdot|p_i|\}$ contains $(a_j\cdot r,a_i\cdot r)$, since $a_j\cdot r<c\cdot a_i\cdot r$;
so the joint spectral measure of $(\mathsf{H}_j(\gamma^a),\mathsf{P}_i)$ does not vanish on it, and $\mathfrak{E}_j(\gamma^a)\left(\{\omega<c\cdot|p_i|\}\right)\neq0$.

\eqref{it:stretch-permutation}
Let $b\in\mathcal{S}(\R^3;\R^3)$.
The identity for $(\mathfrak{g}^{-1})_*\mu$ at the beginning of the proof of \eqref{it:stretch-properties},
\eqref{eq:rotation-stretched-law}, \eqref{it:stretch-actions} and the invariance of $\gamma$ under $g$ give
\[
\begin{aligned}
\left((g^{-1})_*\gamma^a\right)\left(\exp\left(i\mathbb{B}(b)\right)\right)
&=\gamma^a\left(\exp\left(i\mathbb{B}(g\cdot b)\right)\right)
=\gamma\left(\exp\left(i\mathbb{B}(g\cdot b^{a\circ\mathfrak{s}})\right)\right)\\
&=\gamma\left(\exp\left(i\mathbb{B}(b^{a\circ\mathfrak{s}})\right)\right)
=\gamma^{a\circ\mathfrak{s}}\left(\exp\left(i\mathbb{B}(b)\right)\right),
\end{aligned}
\]
and $(g^{-1})_*\gamma^a=\gamma^{a\circ\mathfrak{s}}$ by \cite[Lemma~8.12(ii)]{NoncompactQED}.

\eqref{it:stretch-trace}
For $p\neq0$ one has $\operatorname{tr}A^{-2}=\sum_ly_l$,
$\operatorname{tr}\left(A^{-2}pp^{T}A^{-2}\right)=\left|A^{-2}p\right|^2=\sum_ly_l^2\cdot p_l^2$ and $p^{T}A^{-2}p=\sum_ly_l\cdot p_l^2$;
inserted into the definition of $M^a$, these give the formula.
\end{proof}

\begin{proof}[Proof of Proposition~\ref{prop:sharpness}\eqref{it:sharp-constant}]
Fix $c\in(0,1)$ and put
\[
a:=\left(c^{2},c^{\frac32},c\right),\qquad
\mu^{(c)}:=\frac16\cdot\sum_{\mathfrak s\in\mathfrak S_3}\gamma^{a\circ\mathfrak s}.
\]
The three numbers
\[
\frac{a_2\cdot a_3}{a_1}=c^{\frac12},\qquad
\frac{a_1\cdot a_3}{a_2}=c^{\frac32},\qquad
\frac{a_1\cdot a_2}{a_3}=c^{\frac52}
\]
are smaller than $1$, so the triple $a$ satisfies the condition $a_i\cdot a_k\le a_l$, $\{i,k,l\}=\{1,2,3\}$, of Lemma~\ref{lem:stretched-laws}\eqref{it:stretch-properties}.
This condition does not change when the entries of the triple are permuted, and it therefore holds for each of the six triples $a\circ\mathfrak s$.
By the same item every law $\gamma^{a\circ\mathfrak s}$ satisfies \textup{(G)}, \textup{(dF)}, $(T_{123})$, and $(R_j)$, $(RP_j)$ for $j=1,2,3$.
Each of these properties is an equality or an inequality that is linear in the law;
so the convex combination $\mu^{(c)}$ has them as well.

For $g\in\Ocub$ let $\mathfrak s_g\in\mathfrak S_3$ be the permutation with $ge_k=\pm e_{\mathfrak s_g(k)}$ for $k=1,2,3$.
Lemma~\ref{lem:stretched-laws}\eqref{it:stretch-permutation}, applied to the triple $a\circ\mathfrak s$, gives
\[
\left(g^{-1}\right)_*\gamma^{a\circ\mathfrak s}=\gamma^{a\circ(\mathfrak s\circ\mathfrak s_g)}\qquad(\mathfrak s\in\mathfrak S_3),
\]
and $\mathfrak s\mapsto\mathfrak s\circ\mathfrak s_g$ is a bijection of $\mathfrak S_3$.
Hence $g^{-1}$ permutes the six laws $\gamma^{a\circ\mathfrak s}$ and leaves their mean $\mu^{(c)}$ unchanged.
Thus $\mu^{(c)}$ is invariant under $\Ocub$, and together with the properties already established this is \textup{(H)}.

For $i\ne j$ the quotients $a_j/a_i$ are $c^{\frac12}$, $c$, $c^{\frac12}$ and their inverses;
the least of them is $a_1/a_3=c$, and each triple $a\circ\mathfrak s$ has the same six quotients.
By Lemma~\ref{lem:stretched-laws}\eqref{it:stretch-cone}, for every ordered pair $(j,i)$ the joint spectral measure $\mathfrak E_j(\gamma^{a\circ\mathfrak s})$ is carried by the set
\[
\left\{(\omega,p_\perp):\ \omega\ge\frac{a_{\mathfrak s(j)}}{a_{\mathfrak s(i)}}\cdot|p_i|\right\},
\]
which is contained in $\{\omega\ge c\cdot|p_i|\}$.
Lemma~\ref{lem:cone-constant}\eqref{it:cone-constant-mixtures} carries this over to the mixture:
$\mathfrak E_j(\mu^{(c)})(\{\omega<c\cdot|p_i|\})=0$ for every ordered pair $(j,i)$, that is, $\mathsf H_j\ge c\cdot|\mathsf P_i|$.

It remains to show that $\mu^{(c)}$ is not invariant under the rotation $R^{12}_{\pi/4}$.
By Lemma~\ref{lem:stretched-laws}\eqref{it:stretch-multiplier} the law $\gamma^{a\circ\mathfrak s}$ is the Gaussian law $\gamma_{M^{a\circ\mathfrak s}}$, whose two-point function is the covariance $C_{M^{a\circ\mathfrak s}}$,
and the multiplier $M^{a\circ\mathfrak s}$ is continuous on $\R^3\setminus\{0\}$.
Averaging over $\mathfrak s$, the two-point function of $\mu^{(c)}$ is $C_{\bar M}$, with
\[
\bar M:=\frac16\cdot\sum_{\mathfrak s\in\mathfrak S_3}M^{a\circ\mathfrak s},
\]
a function that is again continuous on $\R^3\setminus\{0\}$.
Suppose that $g_*\mu^{(c)}=\mu^{(c)}$ for some $g\in\mathrm O(3)$.
Then $C_{\bar M}(g\cdot b,g\cdot b)=C_{\bar M}(b,b)$ for every real $b\in\mathcal S(\R^3;\R^3)$ (Lemma~\ref{lem:domination-consequences}\eqref{it:moment-invariance}),
so the Gaussian law $\gamma_{\bar M}$ is invariant under $g$ (Lemma~\ref{lem:gaussian-covariance}\eqref{it:gaussian-laws-criteria}),
and Lemma~\ref{lem:multiplier-laws} gives $\bar M(gk)=g\,\bar M(k)\,g^{T}$ for almost every $k$, hence, by continuity, for every $k\ne0$.
In particular $\operatorname{tr}\bar M(gk)=\operatorname{tr}\bar M(k)$ for every $k\ne0$.

Let $A:=\diag(a)$, $y_l:=a_l^{-2}$ for $l=1,2,3$, and $Y:=y_1+y_2+y_3$;
the numbers $y_1=c^{-4}$, $y_2=c^{-3}$, $y_3=c^{-2}$ are pairwise distinct.
The diagonal matrix of the triple $a\circ\mathfrak s$ has the determinant $\det A$,
and the triple has the numbers $y_{\mathfrak s(1)},y_{\mathfrak s(2)},y_{\mathfrak s(3)}$ in place of $y_1,y_2,y_3$.
Lemma~\ref{lem:stretched-laws}\eqref{it:stretch-trace} therefore gives
\[
\operatorname{tr}M^{a\circ\mathfrak s}(k)=\det A\cdot\left(Y-\frac{\sum_{l=1}^{3}y_{\mathfrak s(l)}^2\cdot k_l^2}{\sum_{l=1}^{3}y_{\mathfrak s(l)}\cdot k_l^2}\right)\qquad(k\ne0).
\]
At $k=e_1$ the quotient is $y_{\mathfrak s(1)}$, which takes each of the values $y_1,y_2,y_3$ for two of the six permutations; hence
\[
\operatorname{tr}\bar M(e_1)=\det A\cdot\left(Y-\frac{Y}{3}\right).
\]
At $k=n_+:=2^{-\frac12}\cdot(e_1+e_2)$ one has $k_1^2=k_2^2=\frac12$ and $k_3=0$,
and the quotient is $(y_{\mathfrak s(1)}^2+y_{\mathfrak s(2)}^2)/(y_{\mathfrak s(1)}+y_{\mathfrak s(2)})$.
Each two-element subset $\{y,y'\}$ of $\{y_1,y_2,y_3\}$ is the set $\{y_{\mathfrak s(1)},y_{\mathfrak s(2)}\}$ for two of the six permutations; therefore
\[
\operatorname{tr}\bar M(n_+)=\det A\cdot\left(Y-\frac13\cdot\sum_{\{y,y'\}}\frac{y^2+y'^2}{y+y'}\right),
\]
where the sum extends over the three two-element subsets of $\{y_1,y_2,y_3\}$.
For $y\ne y'$,
\[
\frac{y^2+y'^2}{y+y'}-\frac{y+y'}{2}=\frac{(y-y')^2}{2\cdot(y+y')}>0,
\]
while $\frac13\cdot\sum_{\{y,y'\}}\frac{y+y'}{2}=\frac{Y}{3}$, because each $y_l$ lies in two of the three subsets.
Consequently $\operatorname{tr}\bar M(n_+)<\operatorname{tr}\bar M(e_1)$.
As $R^{12}_{\pi/4}e_1=n_+$, the law $\mu^{(c)}$ is not invariant under $R^{12}_{\pi/4}$, and a fortiori not under $\mathrm O(3)$.
\end{proof}

The comparison of the two traces is the computation in the proof of \cite[Proposition~21.40(c)]{NoncompactQED}, where it is made for the triple $\left(\frac14,\frac13,\frac12\right)$.
For $a=\left(\frac14,\frac13,\frac12\right)$ the mean $\frac16\cdot\sum_{\mathfrak s}\gamma^{a\circ\mathfrak s}$ is the law $m_A$ of \cite[Example~8.33]{NoncompactQED},
where the triple is written $\lambda$ and the law $\gamma^{a\circ\mathfrak s}$ is written $\gamma^{(\mathfrak s)}$;
the least quotient $a_j/a_i$ of this triple is $\frac12$.

\begin{proof}[Proof of Proposition~\ref{prop:sharpness}\eqref{it:sharp-quarter-turn}]
Let $a:=\left(\frac12,\frac14,\frac14\right)$ and $A:=\diag(a)$.
Here
\[
\frac{a_2\cdot a_3}{a_1}=\frac18,\qquad\frac{a_1\cdot a_3}{a_2}=\frac{a_1\cdot a_2}{a_3}=\frac12,
\]
and Lemma~\ref{lem:stretched-laws}\eqref{it:stretch-properties} shows that the Gaussian law $\gamma^a$ satisfies \textup{(G)}, \textup{(dF)}, $(T_{123})$, and $(R_j)$, $(RP_j)$ for $j=1,2,3$.

The cones are read off from Lemma~\ref{lem:stretched-laws}\eqref{it:stretch-cone}.
In the picture $e_1$ the constants are $a_1/a_2=a_1/a_3=2$:
the joint spectral measure $\mathfrak E_1(\gamma^a)$ is carried by $\{\omega\ge2\cdot|p_2|\}\cap\{\omega\ge2\cdot|p_3|\}$, and $(C_{12})$, $(C_{13})$ hold
(in this picture the speed of light of $\gamma^a$ along $e_2$ and along $e_3$ equals $2$).
In the picture $e_2$ the constant for the direction $e_1$ is $a_2/a_1=\frac12$, and it is optimal;
so $\mathfrak E_2(\gamma^a)(\{\omega<|p_1|\})\ne0$, and $(C_{21})$ fails.
For the remaining ordered pairs the constants are $a_2/a_3=a_3/a_2=1$, so that $(C_{23})$ and $(C_{32})$ hold,
and $a_3/a_1=\frac12$, which is optimal, so that $(C_{31})$ fails.

Since $a_2=a_3$, every rotation $R^{23}_t$ commutes with $A$,
and Lemma~\ref{lem:stretched-laws}\eqref{it:stretch-actions} gives $(R^{23}_t\cdot b)^a=R^{23}_t\cdot b^a$.
Together with the invariance of the free Maxwell law $\gamma$ under all Euclidean maps this yields
\[
\begin{aligned}
\gamma^a\left(\exp\left(i\mathbb B(R^{23}_t\cdot b)\right)\right)
&=\gamma\left(\exp\left(i\mathbb B\left((R^{23}_t\cdot b)^a\right)\right)\right)
=\gamma\left(\exp\left(i\mathbb B(R^{23}_t\cdot b^a)\right)\right)\\
&=\gamma\left(\exp\left(i\mathbb B(b^a)\right)\right)
=\gamma^a\left(\exp\left(i\mathbb B(b)\right)\right)
\end{aligned}
\]
for every real $t$ and every real $b\in\mathcal S(\R^3;\R^3)$.
So $\gamma^a$ is invariant under all rotations $R^{23}_t$,
in accordance with Theorem~\ref{thm:one-plane} for the pair $(2,3)$, all of whose hypotheses $\gamma^a$ satisfies.

Finally, $\det A=\frac1{32}$ and $\left(a_1^{-2},a_2^{-2},a_3^{-2}\right)=(4,16,16)$, and Lemma~\ref{lem:stretched-laws}\eqref{it:stretch-trace} gives
\[
\operatorname{tr}M^a(e_1)=\frac{16+16}{32}=1,\qquad\operatorname{tr}M^a(e_2)=\frac{4+16}{32}=\frac58.
\]
The law $\gamma^a$ is $\gamma_{M^a}$ (Lemma~\ref{lem:stretched-laws}\eqref{it:stretch-multiplier}), and $M^a$ is continuous on $\R^3\setminus\{0\}$.
If the two-point function $C_{M^a}$ of $\gamma^a$ were invariant under the quarter turn $g:=R^{12}_{\pi/2}$,
then $\gamma^a$ would be invariant under $g$ (Lemma~\ref{lem:gaussian-covariance}\eqref{it:gaussian-laws-criteria}),
Lemma~\ref{lem:multiplier-laws} would give $M^a(gk)=g\,M^a(k)\,g^{T}$ for almost every $k$, hence for every $k\ne0$,
and the two traces would coincide, because $ge_1=e_2$.
Therefore neither $\gamma^a$ nor its two-point function is invariant under the quarter turn of the plane $(1,2)$.
As $\gamma^a$ satisfies \textup{(G)}, \textup{(dF)}, $(T_{123})$, $(R_1)$, $(RP_1)$ and $(C_{12})$, hence also $(C^1_{12})$,
it has all hypotheses of Theorem~\ref{thm:one-plane} for the pair $(1,2)$ except $(Q_{12})$,
and all hypotheses of Theorem~\ref{thm:two-point}\eqref{it:two-point-invariance} except the alternative between \eqref{it:quarter-turn-two-point} and \eqref{it:second-mirror-two-point};
the conclusions of both theorems fail for it at the angle $\pi/2$.
\end{proof}

\paragraph{The diagonal Maxwell law.}
For $m\ge0$ let $M$ be the function on $\R^3$ with values in the diagonal $3\times3$ matrices given by $M(0):=0$ and
\[
M_{ab}(p):=\delta_{ab}\cdot\frac{|p|^2-p_a^2}{|p|^2+m^2}\qquad\left(a,b\in\{1,2,3\},\ p\in\R^3\setminus\{0\}\right)
\]
(for $m>0$ the same formula gives the value $0$ at $p=0$),
and let $\gamma^{\oplus}_m:=\gamma_M$ be the corresponding Gaussian law of Lemma~\ref{lem:multiplier-laws}.
For $m=0$ and $p\ne0$ the diagonal entries of $M(p)$ are those of $P^{\perp}(p)$;
then $(\mathbb B_1,\mathbb B_2,\mathbb B_3)$ has under $\gamma^{\oplus}_0$ the law of $(\mathbb B^{(1)}_1,\mathbb B^{(2)}_2,\mathbb B^{(3)}_3)$,
where $\mathbb B^{(1)},\mathbb B^{(2)},\mathbb B^{(3)}$ are the fields of three independent samples of the free Maxwell law $\gamma$.

\begin{proof}[Proof of Proposition~\ref{prop:sharpness}\eqref{it:sharp-closedness-two-point}]
\emph{Domination and the cubic symmetries.}
The function $M$ is measurable, even and bounded, and its values are diagonal matrices with entries in $[0,1]$.
By Lemma~\ref{lem:multiplier-laws} the law $\gamma^{\oplus}_m$ satisfies $(T_{123})$ and, since $M\le1$, also \textup{(G)}.
Let $g\in\Ocub$, and let $\mathfrak s\in\mathfrak S_3$ and the signs $\varepsilon_a\in\{\pm1\}$ be such that $ge_a=\varepsilon_a\cdot e_{\mathfrak s(a)}$ for $a=1,2,3$.
For $p\ne0$ the matrix $g\,M(p)\,g^{T}$ is diagonal, with the entry $M_{aa}(p)$ at the place $(\mathfrak s(a),\mathfrak s(a))$,
and $|gp|=|p|$, $(gp)_{\mathfrak s(a)}=\varepsilon_a\cdot p_a$.
Hence
\[
\left(g\,M(p)\,g^{T}\right)_{\mathfrak s(a)\mathfrak s(a)}=\frac{|p|^2-p_a^2}{|p|^2+m^2}=\frac{|gp|^2-(gp)_{\mathfrak s(a)}^2}{|gp|^2+m^2}=M_{\mathfrak s(a)\mathfrak s(a)}(gp),
\]
that is, $M(gp)=g\,M(p)\,g^{T}$; at $p=0$ both sides vanish.
So $\gamma^{\oplus}_m$ is invariant under $\Ocub$ by Lemma~\ref{lem:multiplier-laws}.

\emph{Reflection positivity and the cones.}
Write $p=(p_1,p_\perp)$ with $p_\perp=(p_2,p_3)$.
Let $\omega$ be the function on $\R^2$ with $\omega(p_\perp):=\left(|p_\perp|^2+m^2\right)^{\frac12}$ wherever this number is positive,
and with $\omega(0):=1$ in the case $m=0$;
it is measurable and takes its values in $(0,\infty)$.
For $p_\perp\ne0$ one has $|p|^2+m^2=p_1^2+\omega(p_\perp)^2$, and therefore
\[
M_{11}(p)=\frac{|p_\perp|^2}{p_1^2+\omega(p_\perp)^2},\qquad
M_{aa}(p)=1-\frac{p_a^2+m^2}{p_1^2+\omega(p_\perp)^2}\quad(a=2,3).
\]
So, for every $p_\perp\ne0$, the multiplier has the form \eqref{eq:rotation-partial-fractions} with $n=1$, $\omega_1=\omega$, the constant polynomial part $Q_{p_\perp}=\diag(0,1,1)$ and the numerator
\[
N_1(p_1,p_\perp)=\diag\left(|p_\perp|^2,\ -\left(p_2^2+m^2\right),\ -\left(p_3^2+m^2\right)\right),
\]
which does not depend on $p_1$.
Moreover $\theta_1\,M(-p_1,p_\perp)\,\theta_1=M(p_1,p_\perp)$ for every $p$, since $M$ is diagonal and even in $p_1$.
The residue matrix
\[
A_1(p_\perp)=-\theta_1\,N_1\left(-i\omega(p_\perp),p_\perp\right)=\diag\left(|p_\perp|^2,\ p_2^2+m^2,\ p_3^2+m^2\right)
\]
is positive semidefinite.
(The component $\mathbb B_1$, normal to the mirror $\{x_1=0\}$ and even under $\theta_1$, has a positive time kernel;
the tangential components $\mathbb B_2$, $\mathbb B_3$, which are odd, have a negative time kernel together with a white noise in time, which comes from $Q_{p_\perp}$;
the factor $-\theta_1$ turns both kinds of residues into positive ones.)
By Lemma~\ref{lem:time-kernels} the law $\gamma^{\oplus}_m$ satisfies $(R_1)$ and $(RP_1)$;
moreover, for every real $b\in C^\infty_{\mathrm c}(\{x_1>0\};\R^3)$, the spectral measure of $[\mathbb B(b)]_1$ is the image,
under the map $p_\perp\mapsto\left(\omega(p_\perp),-p_\perp\right)$, of a positive measure on $\R^2$ that has a density with respect to the Lebesgue measure.
This map sends every $p_\perp\ne0$ into the mass shell
\[
\left\{(\omega,p_\perp):\ \omega=\left(|p_\perp|^2+m^2\right)^{\frac12}\right\},
\]
and the point $p_\perp=0$ forms a null set;
so the spectral measure of $[\mathbb B(b)]_1$ is carried by the mass shell,
a subset of the round cone $\{\omega\ge|p_\perp|\}$ and hence of $\{\omega\ge|p_i|\}$ for $i=2,3$.
For $j=2,3$ choose $g\in\Ocub$ with $ge_1=e_j$;
as $\gamma^{\oplus}_m$ is invariant under $g$, Lemma~\ref{lem:transport}\eqref{it:transport-properties} carries $(RP_1)$ over to $(RP_j)$.
So $\gamma^{\oplus}_m$ satisfies \textup{(H)}.
Lemma~\ref{lem:gaussian-cone}, with $c=1$ and $i=2,3$, gives $(C_{12})$ and $(C_{13})$.
For an arbitrary ordered pair $(j,i)$ the permutation matrix $g_{ji}$ lies in $\Ocub$ and leaves $\gamma^{\oplus}_m$ invariant,
and Lemma~\ref{lem:transport}\eqref{it:transport-pairs} turns $(C_{12})$ into $(C_{ji})$.
The proof of Lemma~\ref{lem:gaussian-cone} is valid for every closed subset $K$ of $[0,\infty)\times\R^2$ with $(0,0)\in K$ and $K+K\subset K$ \eqref{eq:rotation-cone-properties}:
if the spectral measures of all the vectors $[F(f)]_1$, $f\in C^\infty_{\mathrm c}(\{x_1>0\};\Lambda^2)$ real, are carried by $K$, then so is $\mathfrak E_1$.
The round cone is such a set, by the triangle inequality in $\R^2$;
therefore $\mathfrak E_1(\gamma^{\oplus}_m)$ is carried by $\{\omega\ge|p_\perp|\}$.

\emph{Closedness fails.}
For $p\ne0$,
\[
p^{T}M(p)\,p=\sum_{a=1}^{3}p_a^2\cdot\frac{|p|^2-p_a^2}{|p|^2+m^2},
\]
and the right side is positive unless $p$ is parallel to a coordinate axis:
if $p_a\ne0$ and $p_b\ne0$ for some $a\ne b$, the term of index $a$ is at least $p_a^2\cdot p_b^2/(|p|^2+m^2)$.
So $M(p)\,p\ne0$ for almost every $p$, and $\gamma^{\oplus}_m$ violates \textup{(dF)} by Lemma~\ref{lem:multiplier-laws}.

\emph{No invariance under rotations.}
The function $M$ is continuous on $\R^3\setminus\{0\}$, where the denominator $|p|^2+m^2$ is positive.
For $r>0$ and $n_+:=2^{-\frac12}\cdot(e_1+e_2)$,
\[
M(re_1)=\frac{r^2}{r^2+m^2}\cdot\diag(0,1,1),\qquad
M(rn_+)=\frac{r^2}{r^2+m^2}\cdot\diag\left(\frac12,\frac12,1\right),
\]
and these two matrices have different eigenvalues.
If $\gamma^{\oplus}_m$ were invariant under $g:=R^{12}_{\pi/4}$, Lemma~\ref{lem:multiplier-laws} would give $M(gp)=g\,M(p)\,g^{T}$ for almost every $p$, by continuity for every $p\ne0$;
with $p=re_1$ and $gp=rn_+$ the two matrices would be similar.
So $\gamma^{\oplus}_m$ is not invariant under $R^{12}_{\pi/4}$.

\emph{The anomalous harmonic.}
The law $\gamma^{\oplus}_m$ satisfies \textup{(G)}, $(T_{123})$, $(R_1)$, $(RP_1)$ and $(C^1_{12})$, the last one because of $(C_{12})$;
these are the hypotheses of Lemma~\ref{lem:picture-kernels} for $j=1$.
Its two-point function satisfies \eqref{eq:rotation-two-point-quarter-turn}, since the quarter turn $R^{12}_{\pi/2}$ lies in $\Ocub$.
In the proof of Theorem~\ref{thm:two-point}, in the case \eqref{it:quarter-turn-two-point},
the kernels $G_{\lambda\lambda'}$ of \eqref{eq:rotation-four-frames}, the entire functions of \eqref{eq:rotation-entire-angle} and the window \eqref{eq:rotation-window}
are derived from Lemma~\ref{lem:picture-kernels} for $j=1$, from \eqref{eq:rotation-two-point-quarter-turn}
and from the bilinearity of $S_2$ (Lemma~\ref{lem:domination-consequences}\eqref{it:moment-forms}),
with the help of Lemmas~\ref{lem:plate-identities}, \ref{lem:analysis-facts} and~\ref{lem:harmonics} and of the identity theorem for holomorphic functions;
that construction therefore applies to $\gamma^{\oplus}_m$.
Thus, for labels $\lambda$, $\lambda'$ with profiles in $C^\infty_{\mathrm c}(\R;\R)$, there is a smooth function $G_{\lambda\lambda'}$ on $\R^2\setminus\{0\}$ with
\[
S_2(\alpha\otimes\lambda,\alpha'\otimes\lambda')=\int\mathrm{d}y\int\mathrm{d}y'\,\alpha(y)\cdot\alpha'(y')\cdot G_{\lambda\lambda'}(y'-y)
\]
for all $\alpha,\alpha'\in L^2(\R^2;\C)$ with disjoint compact supports, and $G_{\lambda\lambda'}$ satisfies \eqref{eq:rotation-window}.

Let $\chi,\chi'\in C^\infty_{\mathrm c}(\R;\R)$, $\lambda:=(+,\chi)$ and $\lambda':=(+,\chi')$.
Since $w(+)+w(+)=2$, the window admits the frequencies $2$ and $-2$ only:
\[
G_{\lambda\lambda'}(rn_\vartheta)=a^{2}_{\lambda\lambda'}(r)\cdot\exp(2i\vartheta)+a^{-2}_{\lambda\lambda'}(r)\cdot\exp(-2i\vartheta)\qquad(r>0,\ \vartheta\in\R).
\]
The two-point function of $\gamma_M$ is given on all of $L^2(\R^3;\C^3)$ by
\[
S_2(b,b')=(2\pi)^{-3}\int\mathrm{d}p\,\hat b(-p)^{T}M(p)\,\hat b'(p):
\]
both sides are complex bilinear and bounded on $L^2(\R^3;\C^3)$ (Lemma~\ref{lem:domination-consequences}\eqref{it:moment-forms}, and $0\le M\le1$),
they agree for real $b,b'\in\mathcal S(\R^3;\R^3)$, where the right side is $C_M(b,b')$ because $\hat b(-p)=\overline{\hat b(p)}$,
and the real Schwartz fields span a dense subspace of $L^2(\R^3;\C^3)$.
For the plate $b=\alpha\otimes\lambda$, that is $b(x)=\alpha(y)\cdot\chi(x_3)\cdot e^+$, one has
\[
\hat b(p)=\hat\alpha(p_1,p_2)\cdot\overline{\hat\chi(p_3)}\cdot e^+,
\]
with the two-dimensional Fourier transform $\hat\alpha$ and the transform $\hat\chi(p_3)=\int\mathrm{d}x_3\,\exp(ip_3x_3)\cdot\chi(x_3)$ of the profile;
as $\chi$ is real, $\overline{\hat\chi(-p_3)}=\hat\chi(p_3)$.
Moreover $(e^+)^{T}M(p)\,e^+=M_{11}(p)-M_{22}(p)$, because $e^+=e_1+ie_2$ and $M$ is diagonal, and
\[
M_{11}(p)-M_{22}(p)=\frac{p_2^2-p_1^2}{|p|^2+m^2}=-\frac{(p_1+ip_2)^2+(p_1-ip_2)^2}{2\cdot\left(|p|^2+m^2\right)}.
\]
Hence
\[
S_2(\alpha\otimes\lambda,\alpha'\otimes\lambda')=(2\pi)^{-3}\int\mathrm{d}p\,\hat\alpha(-p_1,-p_2)\cdot\hat\alpha'(p_1,p_2)\cdot\hat\chi(p_3)\cdot\overline{\hat\chi'(p_3)}\cdot\frac{p_2^2-p_1^2}{|p|^2+m^2},
\]
and the last factor is even in $p_2$.
Replacing $\alpha$ and $\alpha'$ by $\alpha\circ\theta_2$ and $\alpha'\circ\theta_2$ replaces $\hat\alpha(p_1,p_2)$ by $\hat\alpha(p_1,-p_2)$, and likewise for $\alpha'$;
after the substitution $p_2\mapsto-p_2$ the integral is unchanged:
\[
S_2\left((\alpha\circ\theta_2)\otimes\lambda,(\alpha'\circ\theta_2)\otimes\lambda'\right)=S_2(\alpha\otimes\lambda,\alpha'\otimes\lambda').
\]
Let $\alpha$ and $\alpha'$ have disjoint compact supports; then so have $\alpha\circ\theta_2$ and $\alpha'\circ\theta_2$,
and the substitution $(y,y')\mapsto(\theta_2y,\theta_2y')$ in the density formula turns this identity into
\[
\int\mathrm{d}y\int\mathrm{d}y'\,\alpha(y)\cdot\alpha'(y')\cdot\left(G_{\lambda\lambda'}\left(\theta_2(y'-y)\right)-G_{\lambda\lambda'}(y'-y)\right)=0.
\]
Lemma~\ref{lem:analysis-facts}\eqref{it:equal-densities}, applied on the open set $\{(y,y')\in(\R^2)^2:\ y\ne y'\}$, gives $G_{\lambda\lambda'}(\theta_2z)=G_{\lambda\lambda'}(z)$ for every $z\ne0$.
As $\theta_2(rn_\vartheta)=rn_{-\vartheta}$, the comparison of the coefficients of $\exp(2i\vartheta)$ and $\exp(-2i\vartheta)$ shows
\[
a^{-2}_{\lambda\lambda'}=a^{2}_{\lambda\lambda'}.
\]

Let now $\chi'=\chi\ne0$, so that $\lambda'=\lambda$, and suppose that $G_{\lambda\lambda}$ vanishes identically.
Put
\[
h(p_1,p_2):=\int\mathrm{d}p_3\,|\hat\chi(p_3)|^2\cdot\frac{p_2^2-p_1^2}{|p|^2+m^2}\qquad\left(p=(p_1,p_2,p_3)\right).
\]
As $|p_2^2-p_1^2|\le|p|^2+m^2$ and $\hat\chi$ is square integrable, dominated convergence shows that $h$ is bounded and continuous on $\R^2$.
Let $T$ be the bounded operator on $L^2(\R^2;\C)$ with the Fourier multiplier $h$.
By the formula for $S_2$, Fubini's theorem and Parseval's formula,
\[
\begin{aligned}
S_2(\alpha\otimes\lambda,\alpha'\otimes\lambda)
&=(2\pi)^{-3}\int\mathrm{d}p_1\int\mathrm{d}p_2\,\hat\alpha(-p_1,-p_2)\cdot\hat\alpha'(p_1,p_2)\cdot h(p_1,p_2)\\
&=(2\pi)^{-1}\int\mathrm{d}y\,\alpha(y)\cdot(T\alpha')(y)
\end{aligned}
\]
for all $\alpha,\alpha'\in L^2(\R^2;\C)$.
The density formula and $G_{\lambda\lambda}=0$ give $1_K\,T\,1_{K'}=0$ for all disjoint compact subsets $K$, $K'$ of $\R^2$.
If $A$ and $A'$ are disjoint Borel subsets of $\R^2$, the inner regularity of the Lebesgue measure provides increasing sequences of compact sets $K_n\subset A$ and $K'_n\subset A'$
with $\left|A\setminus\bigcup_nK_n\right|=\left|A'\setminus\bigcup_nK'_n\right|=0$;
then $1_{K_n}\to1_A$ and $1_{K'_n}\to1_{A'}$ strongly, and $1_A\,T\,1_{A'}=0$.
With $A':=\R^2\setminus A$ this yields $T1_A=1_A\,T\,1_A=1_AT$, as in \eqref{eq:rotation-commutation}.
So $T$ commutes with the multiplication by every bounded Borel function, a uniform limit of finite linear combinations of indicator functions,
in particular with the multiplication by $y\mapsto\exp(iv\cdot y)$ for $v\in\R^2$, which acts on Fourier transforms as the translation by $v$.
It follows that $h(p'-v)=h(p')$ for almost every $p'\in\R^2$, by continuity for every $p'$;
as $v$ is arbitrary, $h$ is constant.
But $h(p_2,p_1)=-h(p_1,p_2)$, so that the constant is $0$, whereas
\[
h(p_1,0)=-\int\mathrm{d}p_3\,|\hat\chi(p_3)|^2\cdot\frac{p_1^2}{p_1^2+p_3^2+m^2}<0\qquad(p_1\ne0),
\]
because $\hat\chi$ is continuous and not identically zero.
This contradiction shows that $G_{\lambda\lambda}$ does not vanish identically.
As $G_{\lambda\lambda}(rn_\vartheta)=2a^{2}_{\lambda\lambda}(r)\cdot\cos(2\vartheta)$, there is $r_0>0$ with
\[
a^{-2}_{\lambda\lambda}(r_0)=a^{2}_{\lambda\lambda}(r_0)\ne0:
\]
the kernel $G_{\lambda\lambda}$ contains the anomalous frequency $-2$ with the same coefficient as the frequency $2$.

This harmonic is now read through the anomaly index.
The law $\gamma^{\oplus}_m$ satisfies \textup{(G)}, $(T_{123})$, $(R_1)$, $(RP_1)$, $(Q_{12})$ and $(C_{12})$;
so the function $\mathcal G_2$ on $\mathcal Z_2$ of Lemma~\ref{lem:one-function}, its harmonics $c_\ell$, $\ell\in\Z$, of Lemma~\ref{lem:finitely-many-harmonics}
and the anomaly index $\mathfrak a(N)$ of Definition~\ref{def:anomaly-index} are defined for it.
Let $\lambda_1$, $\lambda_2$ be labels with profiles in $C^\infty_{\mathrm c}(\R;\R)$, put $w:=w(\lambda_1)+w(\lambda_2)$,
and set $a^{k}_{\lambda_1\lambda_2}:=0$ for the integers $k$ that do not occur in \eqref{eq:rotation-window}.
The real points of $\mathcal Z(\mathrm{id},0)$ form the open set of the configurations $(y_1,y_2)$ with $y_1^0<y_2^0$.
On this set the functions $(y_1,y_2)\mapsto\mathcal G_2\left((y_1,y_2);\lambda_1,\lambda_2\right)$ and $(y_1,y_2)\mapsto G_{\lambda_1\lambda_2}(y_2-y_1)$ are continuous,
and both are densities of $S_2(\alpha_1\otimes\lambda_1,\alpha_2\otimes\lambda_2)$ for the $\alpha_1$, $\alpha_2$ with compact supports whose product lies in it,
by Lemma~\ref{lem:one-function}\eqref{it:tube-densities} and by the density formula above;
by Lemma~\ref{lem:analysis-facts}\eqref{it:equal-densities} they coincide.
Let $X=(y_1,y_2)$ be such a configuration, and write $y_2-y_1=rn_\vartheta$ with $r>0$ and $|\vartheta|<\pi/2$.
For $|\vartheta+t|<\pi/2$ the configuration $R_tX$ is again a real point of $\mathcal Z(\mathrm{id},0)$, and $R_t(y_2-y_1)=rn_{\vartheta+t}$;
these $t$ form the interval $I_X$ of Lemma~\ref{lem:exact-phases}\eqref{it:phase-interval}.
For $t\in I_X$, that lemma, the identity $\sum_\ell c_\ell=\mathcal G_2$ of Lemma~\ref{lem:finitely-many-harmonics} and \eqref{eq:rotation-window} give
\[
\sum_{\ell\in\Z}c_\ell(X;\lambda_1,\lambda_2)\cdot\exp\left(i(w+4\ell)t\right)=\mathcal G_2(R_tX;\lambda_1,\lambda_2)=\sum_{k\in\Z}a^{k}_{\lambda_1\lambda_2}(r)\cdot\exp\left(ik(\vartheta+t)\right).
\]
Both sides are trigonometric polynomials in $t$ with frequencies in $w+4\Z$, and they agree on an interval;
so their coefficients agree:
\[
c_\ell(X;\lambda_1,\lambda_2)=a^{w+4\ell}_{\lambda_1\lambda_2}(r)\cdot\exp\left(i(w+4\ell)\vartheta\right)\qquad(\ell\in\Z).
\]
For $\lambda_1=\lambda_2=\lambda=(+,\chi)$ with $\chi\ne0$ one has $w=2$,
and for $\ell=-1$ and $X=(0,r_0n_0)$ this gives $c_{-1}(X;\lambda,\lambda)=a^{-2}_{\lambda\lambda}(r_0)\ne0$.
So $c_{-1}(\cdot\,;\lambda,\lambda)$ is not identically zero, and $\mathfrak a(2)\ge1$;
on the other hand $\mathfrak a(2)\le1$ by Definition~\ref{def:anomaly-index}\eqref{it:anomaly-linear-bound}.
Hence $\mathfrak a(2)=1$.
Proposition~\ref{prop:anomaly-vanishes} asserts, for laws with \textup{(G)}, $(T_{123})$, $(R_1)$, $(RP_1)$, $(Q_{12})$ and $(C_{12})$ that need not satisfy \textup{(dF)},
that $\mathfrak a(2n)=n\cdot\mathfrak a(2)$ for every integer $n\ge0$;
for $\gamma^{\oplus}_m$ this gives $\mathfrak a(2n)=n$ for every integer $n\ge1$.
\end{proof}

The law $\gamma^{\oplus}_m$ has all hypotheses of Theorem~\ref{thm:one-plane} for the pair $(1,2)$ except \textup{(dF)}.
A field whose plate kernels obey the bound of \eqref{eq:rotation-window} thus has room for the anomalous frequency $-2$ of the pair of labels $(+,\chi)$, $(+,\chi)$;
in the proof of Theorem~\ref{thm:two-point} it is the Bianchi identity \eqref{eq:rotation-bianchi}, a consequence of \textup{(dF)}, that removes it.

\paragraph{White noises with permuted variances.}
Let $D=(d_{12},d_{13},d_{23})$ be a triple of pairwise distinct numbers in $(0,1]$, and put $d_{li}:=d_{il}$.
For a real test two-form $f$ with the components $f_{il}$ let
\[
q_D(f):=\sum_{i<l}d_{il}\cdot\|f_{il}\|^2,
\]
and let $\gamma_D$ be the centred Gaussian law with
\[
\gamma_D\left(\exp(iF(f))\right)=\exp\left(-\frac12\cdot q_D(f)\right).
\]
For $\mathfrak s\in\mathfrak S_3$ let $\mathfrak sD$ be the triple with $(\mathfrak sD)_{\mathfrak s(i)\mathfrak s(l)}=d_{il}$; its entries are those of $D$ in another order.
The law
\[
m_1:=\frac16\cdot\sum_{\mathfrak s\in\mathfrak S_3}\gamma_{\mathfrak sD}
\]
is the law $m_1$ of \cite[Example~8.26]{NoncompactQED}, where $F(f)$ is written $X(f)$.

\begin{proof}[Proof of Proposition~\ref{prop:sharpness}\eqref{it:sharp-contact-terms}]
\emph{The hypothesis \textup{(H)}.}
By \cite[Example~8.26\,(1)]{NoncompactQED} each law $\gamma_{\mathfrak sD}$ satisfies \textup{(G)}, $(T_{123})$ and $(RP_j)$ for $j=1,2,3$,
and it is invariant under the diagonal sign matrices, in particular under $\theta_1$, $\theta_2$, $\theta_3$.
By \cite[Example~8.26\,(3)]{NoncompactQED} every element of $\Ocub$ permutes the six laws $\gamma_{\mathfrak sD}$, and $m_1$ satisfies \textup{(H)}.

\emph{The reconstruction is trivial.}
Let $j\in\{1,2,3\}$, let $\Sigma_j^-$ be the $\sigma$-algebra generated by the $F(f)$ with real $f\in C^\infty_{\mathrm c}(\{x_j<0\};\Lambda^2)$,
and let $\gamma_D$ stand for any of the six laws $\gamma_{\mathfrak sD}$.
If $f$ and $f'$ are real test two-forms with disjoint supports, then $\int\mathrm{d}x\,f_{il}(x)\cdot f'_{il}(x)=0$, so that $q_D(f+f')=q_D(f)+q_D(f')$ and
\[
\gamma_D\left(\exp\left(iF(f)+iF(f')\right)\right)=\gamma_D\left(\exp(iF(f))\right)\cdot\gamma_D\left(\exp(iF(f'))\right).
\]
Applied to $f=\sum_at_a\cdot f_a$ and $f'=\sum_bt'_b\cdot f'_b$ with real numbers $t_a$, $t'_b$ and real $f_a\in C^\infty_{\mathrm c}(\{x_j>0\};\Lambda^2)$, $f'_b\in C^\infty_{\mathrm c}(\{x_j<0\};\Lambda^2)$,
this says that the joint characteristic function of the two families $(F(f_a))_a$ and $(F(f'_b))_b$ factorises.
By Dynkin's $\pi$--$\lambda$ theorem the $\sigma$-algebras $\Sigma_j^+$ and $\Sigma_j^-$ are independent under $\gamma_D$.
For $\Phi,\Phi'\in L^2(\Sigma_j^+,\gamma_D)$ the function $\Theta_j\Phi$ is $\Sigma_j^-$-measurable, and $\gamma_D(\Theta_j\Phi)=\overline{\gamma_D(\Phi)}$ by $(R_j)$; therefore
\[
\gamma_D\left((\Theta_j\Phi)\cdot\Phi'\right)=\overline{\gamma_D(\Phi)}\cdot\gamma_D(\Phi').
\]
As $m_1\ge\frac16\cdot\gamma_{\mathfrak sD}$, the space $L^2(\Sigma_j^+,m_1)$ is contained in each space $L^2(\Sigma_j^+,\gamma_{\mathfrak sD})$,
and in the reconstruction space $\mathcal H_j$ of $m_1$
\[
\left\langle[\Phi]_j,[\Phi']_j\right\rangle=m_1\left((\Theta_j\Phi)\cdot\Phi'\right)=\frac16\cdot\sum_{\mathfrak s\in\mathfrak S_3}\overline{\gamma_{\mathfrak sD}(\Phi)}\cdot\gamma_{\mathfrak sD}(\Phi')\qquad\left(\Phi,\Phi'\in L^2(\Sigma_j^+,m_1)\right).
\]
Consequently
\[
\left\|[\Phi]_j-[\Phi']_j\right\|^2=\frac16\cdot\sum_{\mathfrak s\in\mathfrak S_3}\left|\gamma_{\mathfrak sD}(\Phi)-\gamma_{\mathfrak sD}(\Phi')\right|^2:
\]
the class $[\Phi]_j$ is determined by the six numbers $\gamma_{\mathfrak sD}(\Phi)$, and these do not change under translations, by $(T_{123})$.
For $t\ge0$ and $y\in e_j^{\perp}$ the operator $U(te_j+y)$ maps $L^2(\Sigma_j^+,m_1)$ into itself, and by \cite[Proposition~8.19(ii), (iii)]{NoncompactQED}
\[
\exp(-t\mathsf H_j)\,V_j(y)\,[\Phi]_j=[U(te_j+y)\Phi]_j=[\Phi]_j.
\]
Hence $\exp(-t\mathsf H_j)=1$ and $V_j(y)=1$ for all $t\ge0$ and $y\in e_j^{\perp}$:
$\mathsf H_j=0$, $\mathsf P_i=0$ for $i\ne j$, and $\mathfrak E_j$ is concentrated at the point $(0,0)$.
This point does not lie in $\{\omega<|p_i|\}$, so every $(C_{ji})$ holds.

\emph{The two-point function.}
For a real test two-form $f$ and $b=\ast f$ the variance of $F(f)$ under $\gamma_{\mathfrak sD}$ is $q_{\mathfrak sD}(f)$, so that
\[
S_2(b,b)=m_1\left(F(f)^2\right)=\frac16\cdot\sum_{\mathfrak s\in\mathfrak S_3}\sum_{i<l}(\mathfrak sD)_{il}\cdot\|f_{il}\|^2.
\]
For each pair $i<l$ the number $(\mathfrak sD)_{il}$ takes each of the values $d_{12}$, $d_{13}$, $d_{23}$ for two of the six permutations; hence
\[
S_2(b,b)=\frac{d_{12}+d_{13}+d_{23}}{3}\cdot\|b\|^2.
\]
By polarisation $S_2(b,b')$ is the same multiple of the scalar product of $b$ and $b'$, and $S_2$ is invariant under $\mathrm O(3)$, which acts isometrically on $L^2(\R^3;\R^3)$.

\emph{Disjoint supports.}
Let $N\ge1$, and let $b_1,\dots,b_N\in C^\infty_{\mathrm c}(\R^3;\R^3)$ have pairwise disjoint supports, $b_a=\ast f_a$.
The covariance of $F(f_a)$ and $F(f_{a'})$ under $\gamma_{\mathfrak sD}$ is $\sum_{i<l}(\mathfrak sD)_{il}\cdot\int\mathrm{d}x\,f_{a,il}(x)\cdot f_{a',il}(x)$, which vanishes for $a\ne a'$.
So the centred jointly Gaussian variables $F(f_1),\dots,F(f_N)$ are pairwise uncorrelated, hence independent under $\gamma_{\mathfrak sD}$, and
\[
\gamma_{\mathfrak sD}\left(\prod_{a=1}^{N}\mathbb B(b_a)\right)=\prod_{a=1}^{N}\gamma_{\mathfrak sD}\left(\mathbb B(b_a)\right)=0.
\]
Averaging over $\mathfrak s$ gives $S_N(b_1,\dots,b_N)=0$: all moments of $m_1$ vanish on such tuples, and by multilinearity also on complex fields with pairwise disjoint supports.

\emph{Closedness fails.}
Let $u\in\mathcal S(\R^3)$ be real and not identically zero.
By the formula for $S_2$, $m_1\left(\mathbb B(\nabla u)^2\right)=\frac{d_{12}+d_{13}+d_{23}}{3}\cdot\|\nabla u\|^2>0$,
whereas \textup{(dF)} would give $\mathbb B(\nabla u)=0$ almost surely \eqref{eq:rotation-gradient-fields}.

\emph{No invariance under rotations.}
Let $u\in\mathcal S(\R^3)$ be real and not identically zero, and $b:=u\,e_1$, so that $b=\ast f$ with $f_{23}=u$ and $f_{13}=f_{12}=0$.
Then $q_{\mathfrak sD}(f)=(\mathfrak sD)_{23}\cdot\|u\|^2$, and $(\mathfrak sD)_{23}$ takes each value $d\in\{d_{12},d_{13},d_{23}\}$ for two permutations; so
\[
m_1\left(\exp(i\mathbb B(b))\right)=\frac13\cdot\sum_{d}\exp\left(-\frac{d\cdot\|u\|^2}{2}\right),
\]
the sum being taken over the three numbers $d\in\{d_{12},d_{13},d_{23}\}$.
Let $g:=R^{12}_{\pi/4}$.
As $\det g=1$ and $ge_1=n_+:=2^{-\frac12}\cdot(e_1+e_2)$, one has $(g\cdot b)(x)=u(g^{-1}x)\cdot n_+$;
the components $(g\cdot b)_1=(g\cdot f)_{23}$ and $(g\cdot b)_2=-(g\cdot f)_{13}$ have the squared norm $\frac12\cdot\|u\|^2$ each, and $(g\cdot f)_{12}=0$.
Hence $q_{\mathfrak sD}(g\cdot f)=\frac12\cdot\left((\mathfrak sD)_{23}+(\mathfrak sD)_{13}\right)\cdot\|u\|^2$.
The set $\{(\mathfrak sD)_{23},(\mathfrak sD)_{13}\}$ is a two-element subset of $\{d_{12},d_{13},d_{23}\}$, and each such subset occurs for two permutations; so
\[
\begin{aligned}
m_1\left(\exp(i\mathbb B(g\cdot b))\right)
&=\frac16\cdot\sum_{\mathfrak s\in\mathfrak S_3}\exp\left(-\frac{\left((\mathfrak sD)_{23}+(\mathfrak sD)_{13}\right)\cdot\|u\|^2}{4}\right)\\
&=\frac13\cdot\sum_{\{d,d'\}}\exp\left(-\frac{(d+d')\cdot\|u\|^2}{4}\right),
\end{aligned}
\]
the sum being taken over the three two-element subsets $\{d,d'\}$ of $\{d_{12},d_{13},d_{23}\}$.
By the strict convexity of the exponential function,
\[
\exp\left(-\frac{(d+d')\cdot\|u\|^2}{4}\right)<\frac12\cdot\left(\exp\left(-\frac{d\cdot\|u\|^2}{2}\right)+\exp\left(-\frac{d'\cdot\|u\|^2}{2}\right)\right)\qquad(d\ne d'),
\]
and the sum over the three subsets, in which each of the three numbers occurs twice, gives
\[
m_1\left(\exp(i\mathbb B(g\cdot b))\right)<m_1\left(\exp(i\mathbb B(b))\right).
\]
Since $\hat g\exp(i\mathbb B(b))=\exp(i\mathbb B(g\cdot b))$, the law $m_1$ is not invariant under $R^{12}_{\pi/4}$, nor under $\mathrm O(3)$.
\end{proof}

Invariance of all moments on test functions with pairwise disjoint supports, which is the form of invariance that the argument for four-dimensional Yang--Mills theory provides,
therefore does not imply the invariance of a law:
for $m_1$ the whole defect of invariance lies in the contact terms.
For a closed field this is excluded by Proposition~\ref{prop:no-contact-terms}, applied to the difference of the rotated and the unrotated moments, which vanishes on gradients by \textup{(dF)}.

\subsubsection{A form factor and a dielectric}\label{sec:form-factor-dielectric}

The laws $\gamma^{\mathrm{f}}_\varepsilon$ and $\gamma^{\mathrm{d}}_{\kappa,m}$
of Proposition~\ref{prop:sharpness}\eqref{it:sharp-domination} and~\eqref{it:sharp-soft-properties}
are both of the form $\gamma_M$, with a multiplier $M(p_1,p_\perp)$ that is,
for almost every $p_\perp$, a rational function of $p_1$ whose poles are simple and lie on the imaginary axis.

\begin{proof}[Proof of Proposition~\ref{prop:sharpness}\eqref{it:sharp-domination}]
Fix $0<\varepsilon\leq1$ and put
\[
\begin{aligned}
h_\varepsilon(p)&:=\frac{1+\varepsilon\left(p_1^4+p_2^4+p_3^4-\frac{3}{5}\right)}{|p|^2+1},\\
M_{ab}(p)&:=h_\varepsilon(p)\cdot\left(|p|^2\delta_{ab}-p_ap_b\right)\qquad(a,b=1,2,3),
\end{aligned}
\]
and $\gamma^{\mathrm{f}}_\varepsilon:=\gamma_M$.

\emph{Symmetries, closedness, and the failure of domination.}
The numerator of $h_\varepsilon$ is at least $1-\frac{3\varepsilon}{5}>0$, so that $h_\varepsilon>0$;
and for $p\neq0$ the matrix with the entries $|p|^2\delta_{ab}-p_ap_b$ is $|p|^2P^\perp(p)\geq0$, while it vanishes at $p=0$.
Hence $M$ is a continuous function on $\R^3$ with values in the real symmetric positive semidefinite matrices;
it is even and polynomially bounded, and $M(p)p=0$.
By Lemma~\ref{lem:multiplier-laws}, $\gamma^{\mathrm{f}}_\varepsilon$ is a centred Gaussian law with $(T_{123})$,
and it satisfies \textup{(dF)}.
Let $g\in\Ocub$.
The coordinates of $gp$ are those of $p$, permuted and with signs;
therefore $|gp|=|p|$ and $h_\varepsilon(gp)=h_\varepsilon(p)$,
and together with $P^\perp(gp)=g\,P^\perp(p)\,g^{T}$ this gives $M(gp)=g\,M(p)\,g^{T}$.
By Lemma~\ref{lem:multiplier-laws}, $\gamma^{\mathrm{f}}_\varepsilon$ is invariant under $\R^3\rtimes\Ocub$.
For $g=\theta_1$ the last identity becomes $\theta_1M(-p_1,p_\perp)\,\theta_1=M(p_1,p_\perp)$.

For $p\neq0$ the number $|p|^2\cdot h_\varepsilon(p)$ is an eigenvalue of $M(p)$, with the plane orthogonal to $p$ as eigenspace.
Since $p_1^4+p_2^4+p_3^4\geq|p|^4/3$, this eigenvalue tends to infinity with $|p|$.
So $M\leq1$ fails outside a ball, and $\gamma^{\mathrm{f}}_\varepsilon$ violates \textup{(G)} by Lemma~\ref{lem:multiplier-laws}.

The trace $\operatorname{tr}M(p)=2|p|^2\cdot h_\varepsilon(p)$ is not a function of $|p|$:
at the unit vectors $e_1$ and $2^{-\frac{1}{2}}(e_1+e_2)$ the numerators of $h_\varepsilon$ are $1+\frac{2\varepsilon}{5}$ and $1-\frac{\varepsilon}{10}$.
If $\gamma^{\mathrm{f}}_\varepsilon$ were invariant under $\mathrm{O}(3)$,
Lemma~\ref{lem:multiplier-laws} would give $M(gp)=g\,M(p)\,g^{T}$ for every $g\in\mathrm{O}(3)$ and almost every $p$,
hence for every $p$, as $M$ is continuous;
taking traces, $\operatorname{tr}M$ would be a function of $|p|$.
Consequently $\gamma^{\mathrm{f}}_\varepsilon$ is not invariant under $\mathrm{O}(3)$.

\emph{Partial fractions in $p_1$.}
Let $\omega(p_\perp):=\left(1+|p_\perp|^2\right)^{\frac{1}{2}}$, so that $|p|^2+1=p_1^2+\omega^2$.
The identity $p_1^4=(p_1^2+\omega^2)(p_1^2-\omega^2)+\omega^4$ gives
\[
h_\varepsilon(p)=\varepsilon\left(p_1^2-\omega^2\right)+\frac{r(p_\perp)}{p_1^2+\omega^2}
\]
with
\[
r(p_\perp):=1+\varepsilon\left(\omega^4+p_2^4+p_3^4-\frac{3}{5}\right)\geq1+\frac{2\varepsilon}{5}>0,
\]
where the inequality holds because $\omega\geq1$.
This is the decomposition \eqref{eq:rotation-partial-fractions} of Lemma~\ref{lem:time-kernels} with one pole, $\omega_1:=\omega$:
the polynomial part $Q_{p_\perp}(p_1)$ has the entries $\varepsilon\left(p_1^2-\omega^2\right)\cdot\left((p_1^2+|p_\perp|^2)\,\delta_{ab}-p_ap_b\right)$,
and the numerator $N_1(p_1,p_\perp)$ has the entries $r(p_\perp)\cdot\left((p_1^2+|p_\perp|^2)\,\delta_{ab}-p_ap_b\right)$;
the coefficients of these polynomials in $p_1$ are continuous functions of $p_\perp$.
At $p_1=-i\omega$ one has $p_1^2+|p_\perp|^2=|p_\perp|^2-\omega^2$,
and the matrix $A_1=-\theta_1N_1(-i\omega,p_\perp)$ of Lemma~\ref{lem:time-kernels} has the entries
\[
\begin{aligned}
&(A_1)_{11}=r\,|p_\perp|^2,\qquad(A_1)_{1a}=i\omega\,r\,p_a,\qquad(A_1)_{a1}=-i\omega\,r\,p_a,\\
&(A_1)_{ab}=r\left(\left(\omega^2-|p_\perp|^2\right)\delta_{ab}+p_ap_b\right)\qquad(a,b\in\{2,3\}),
\end{aligned}
\]
where $r=r(p_\perp)$ and $\omega=\omega(p_\perp)$;
it is a Hermitian matrix.
For $p_\perp\neq0$ the vectors $e_1$, $|p_\perp|^{-1}(0,p_2,p_3)$ and $|p_\perp|^{-1}(0,-p_3,p_2)$ form an orthonormal basis of $\R^3$.
If $c\in\C^3$ has the coordinates $c_1,c_2,c_3$ in this basis, then
\[
\frac{c^{\dagger}A_1(p_\perp)\,c}{r(p_\perp)}
=\left|\,|p_\perp|\,c_1+i\omega\,c_2\right|^2+\left(\omega^2-|p_\perp|^2\right)|c_3|^2,
\]
since the left-hand side equals
\[
|p_\perp|^2\cdot|c_1|^2+i\omega\,|p_\perp|\cdot\left(\overline{c_1}\,c_2-c_1\overline{c_2}\right)
+\left(\omega^2-|p_\perp|^2\right)\left(|c_2|^2+|c_3|^2\right)+|p_\perp|^2\cdot|c_2|^2 .
\]
Here $\omega^2-|p_\perp|^2=1$.
Hence $A_1(p_\perp)$ is positive semidefinite for $p_\perp\neq0$, and so is $A_1(0)=r(0)\cdot\diag(0,1,1)$.

\emph{Reflection positivity and the spectrum.}
By Lemma~\ref{lem:time-kernels}, $\gamma^{\mathrm{f}}_\varepsilon$ satisfies $(R_1)$ and $(RP_1)$,
and for every real $b\in C^\infty_{\mathrm c}(\{x_1>0\};\R^3)$ the spectral measure $\varrho_{[\mathbb{B}(b)]_1}$
is the image of a finite positive measure on $\R^2$ under $p_\perp\mapsto(\omega(p_\perp),-p_\perp)$;
it is therefore carried by the mass shell $\left\{(\omega,p_\perp):\ \omega=\left(1+|p_\perp|^2\right)^{\frac{1}{2}}\right\}$.
As $\mathbb{B}(b)=F(f)$ for $b=\ast f$, this holds for the spectral measures of the vectors $[F(f)]_1$, $f\in C^\infty_{\mathrm c}(\{x_1>0\};\Lambda^2)$ real.
The law is invariant under the permutation matrices,
so Lemma~\ref{lem:transport}\eqref{it:transport-properties} turns its reflection positivity at $\{x_1=0\}$ into reflection positivity at $\{x_j=0\}$:
$(RP_j)$ holds for $j=1,2,3$.

The set
\[
K':=\{(0,0)\}\cup\left\{(\omega,p_\perp)\in[0,\infty)\times\R^2:\ \omega\geq\left(1+|p_\perp|^2\right)^{\frac{1}{2}}\right\}
\]
is closed, contains $(0,0)$ and the mass shell, and satisfies $K'+K'\subset K'$.
For the last property it suffices to take two points $(\omega,p_\perp)$ and $(\omega',p'_\perp)$ of the second set;
the triangle inequality in $\R^3$ for the vectors $(1,p_\perp)$ and $(1,p'_\perp)$ gives
\[
\omega+\omega'\geq\left(1+|p_\perp|^2\right)^{\frac{1}{2}}+\left(1+|p'_\perp|^2\right)^{\frac{1}{2}}
\geq\left(4+|p_\perp+p'_\perp|^2\right)^{\frac{1}{2}}
\geq\left(1+|p_\perp+p'_\perp|^2\right)^{\frac{1}{2}}.
\]
So $K'$ has the properties~\eqref{eq:rotation-cone-properties}.
The proof of Lemma~\ref{lem:gaussian-cone} shows, for a centred Gaussian law with $(T_{123})$, $(R_1)$ and $(RP_1)$,
that the complement of a set with the properties~\eqref{eq:rotation-cone-properties} is a null set of $\mathfrak{E}_1$
as soon as it is a null set of the spectral measures of the vectors $[F(f)]_1$, $f\in C^\infty_{\mathrm c}(\{x_1>0\};\Lambda^2)$ real.
The law $\gamma^{\mathrm{f}}_\varepsilon$ has $(T_{123})$ by Lemma~\ref{lem:multiplier-laws};
applied to it and to $K'$, this shows that $\mathfrak{E}_1$ is carried by $K'$.

Now let $j\in\{2,3\}$, let $i\neq j$,
and let $g_{ji}$ be the permutation matrix of Lemma~\ref{lem:transport}\eqref{it:transport-pairs}, with $g_{ji}e_1=e_j$ and $g_{ji}e_2=e_i$.
Since $\left(g_{ji}^{-1}\right)_*\gamma^{\mathrm{f}}_\varepsilon=\gamma^{\mathrm{f}}_\varepsilon$,
that lemma provides a unitary map $J$ of $\mathcal{H}_1$ onto $\mathcal{H}_j$ such that
$J\exp(-t\mathsf{H}_1)J^{-1}=\exp(-t\mathsf{H}_j)$ for $t\geq0$
and $J\,V_1(y)\,J^{-1}=V_j(g_{ji}y)$ for all $y\in e_1^\perp$
(because $\hat g_{ji}\,U(a)=U(g_{ji}a)\,\hat g_{ji}$);
by the same lemma,
\[
\mathfrak{E}_j(A)=J\,\mathfrak{E}_1\left(\left\{(\omega,p_\perp):\ (\omega,g_{ji}p_\perp)\in A\right\}\right)J^{-1}
\qquad\text{for every Borel set }A\subset[0,\infty)\times e_j^\perp .
\]
Membership in $K'$ depends only on $\omega$ and $|p_\perp|$, and $|g_{ji}p_\perp|=|p_\perp|$;
with the complement of $K'$ as $A$, it follows that $\mathfrak{E}_j$ is carried by $K'$ as well.
Since $K'\subset\{\omega\geq|p_\perp|\}$, this gives, for $j=1,2,3$, both the round cone and the mass gap.
\end{proof}

\begin{proof}[Proof of Proposition~\ref{prop:sharpness}\eqref{it:sharp-soft-properties}]
Let $\kappa>0$ and $m>0$.
For $c=1,2,3$ and $p\in\R^3$ put
\[
d_c(p):=1+\kappa\sum_{a\neq c}\frac{1}{p_a^2+m^2},
\qquad
D(p):=\diag\left(d_1(p),d_2(p),d_3(p)\right),
\qquad
\nu(p):=p^{T}D(p)^{-1}p;
\]
for $p\neq0$ put
\[
M(p):=D^{-1}-\frac{D^{-1}p\,p^{T}D^{-1}}{\nu}\qquad\left(D=D(p),\ \nu=\nu(p)\right),
\]
let $M(0):=0$, and let $\gamma^{\mathrm{d}}_{\kappa,m}:=\gamma_M$.
(For $p\neq0$, $M(p)$ is the covariance matrix of a centred Gaussian vector of $\R^3$ with the covariance $D(p)^{-1}$,
conditioned on the plane orthogonal to $p$.)

\emph{Properties other than reflection positivity.}
Since $d_c\geq1$, one has $D\geq1$, and $\nu(p)>0$ for $p\neq0$;
as $D$ is continuous ($m>0$), the multiplier $M$ is continuous on $\R^3\setminus\{0\}$.
With $u:=D^{-\frac{1}{2}}p$,
\[
M=D^{-\frac{1}{2}}\left(1-\frac{uu^{T}}{|u|^2}\right)D^{-\frac{1}{2}}\geq0,
\qquad
Mp=D^{-1}p-D^{-1}p=0;
\]
so \textup{(dF)} holds by Lemma~\ref{lem:multiplier-laws}.
Moreover $MD=1-D^{-1}pp^{T}/\nu$, and
\[
MDM=\left(1-\frac{D^{-1}pp^{T}}{\nu}\right)\left(D^{-1}-\frac{D^{-1}pp^{T}D^{-1}}{\nu}\right)=M.
\]
As $D\geq1$, it follows that $M=MDM\geq M^2$, so that the eigenvalues of $M$ lie in $[0,1]$.
This gives \textup{(G)} by Lemma~\ref{lem:multiplier-laws};
and since $Mp=0$, also $M=P^\perp(p)\,M\,P^\perp(p)\leq P^\perp(p)$.
The multiplier is even.
Let $g\in\Ocub$, with $ge_a=\pm e_{\mathfrak{s}(a)}$ for a permutation $\mathfrak{s}$ of $\{1,2,3\}$.
Then $(gp)_{\mathfrak{s}(a)}^2=p_a^2$, hence
\[
d_{\mathfrak{s}(c)}(gp)=d_c(p),\qquad D(gp)=g\,D(p)\,g^{T},\qquad \nu(gp)=\nu(p),\qquad M(gp)=g\,M(p)\,g^{T}
\]
for all $p$.
By Lemma~\ref{lem:multiplier-laws}, $\gamma^{\mathrm{d}}_{\kappa,m}$ has $(T_{123})$ and is invariant under $\Ocub$, hence under $\R^3\rtimes\Ocub$.
For $g=\theta_1$ the identity becomes $\theta_1M(-p_1,p_\perp)\,\theta_1=M(p_1,p_\perp)$ for all $p$.
Read entry by entry, it says that $M_{11}$ and the four entries $M_{ab}$ with $a,b\in\{2,3\}$ are even functions of $p_1$,
whereas the four entries $M_{1a}=M_{a1}$ with $a\in\{2,3\}$ are odd functions of $p_1$.
Finally let $p_1p_2p_3\neq0$.
Then $d_c(p/s)\to1$ as $s\to0$, for $c=1,2,3$, and
\[
M(p/s)=D^{-1}-\frac{D^{-1}p\,p^{T}D^{-1}}{p^{T}D^{-1}p}\qquad\left(D=D(p/s)\right)
\]
tends to $1-pp^{T}/|p|^2=P^\perp(p)$.

\emph{The zeros of $\nu$ as a function of $p_1^2$.}
Fix $p_\perp$ with $p_2p_3\neq0$ and $p_2^2\neq p_3^2$; almost every $p_\perp$ is of this kind.
Write $z:=p_1^2$.
Here $d_1=1+\kappa\left(p_2^2+m^2\right)^{-1}+\kappa\left(p_3^2+m^2\right)^{-1}$ is determined by $p_\perp$ alone,
and for $\{j,j'\}=\{2,3\}$ one has $d_j=a_j+\kappa/(z+m^2)$ with $a_j:=1+\kappa/(p_{j'}^2+m^2)$, so that
\[
\frac{1}{d_j}=\frac{z+m^2}{a_j\,(z+M_j)},
\qquad
M_j:=m^2+\frac{\kappa}{a_j}>m^2,
\]
and $M_2\neq M_3$, because $a_2\neq a_3$.
Hence $\nu=\tilde\nu(z)$, where
\[
\begin{aligned}
\tilde\nu(z)&:=\frac{z}{d_1}+\sum_{j=2,3}\frac{p_j^2}{a_j}\left(1-\frac{M_j-m^2}{z+M_j}\right),\\
\tilde\nu'(z)&=\frac{1}{d_1}+\sum_{j=2,3}\frac{p_j^2\,(M_j-m^2)}{a_j\,(z+M_j)^2}>0\qquad\text{for real }z\notin\{-M_2,-M_3\}.
\end{aligned}
\]
Multiplied by $(z+M_2)(z+M_3)$, the function $\tilde\nu$ becomes a polynomial of degree three in $z$.
Moreover $\tilde\nu$ maps each of the intervals $(-\infty,-\max_jM_j)$, $(-\max_jM_j,-\min_jM_j)$ and $(-\min_jM_j,\infty)$
strictly increasingly onto $\R$:
the term $z/d_1$ dominates as $z\to\pm\infty$,
and as $z$ approaches $-M_j$ from the left or from the right, $\tilde\nu(z)$ tends to $+\infty$ or to $-\infty$, respectively.
This polynomial therefore has three real simple zeros $z_1<z_2<z_3$, one in each interval, and these are all the zeros of $\tilde\nu$ in $\C$.
The points $-m^2$ and $0$ lie in the last interval, and
\[
\tilde\nu(-m^2)=-\frac{m^2}{d_1}<0<\sum_{j=2,3}\frac{p_j^2\,m^2}{a_j\,M_j}=\tilde\nu(0);
\]
hence $-m^2<z_3<0$.
Put $\omega_l(p_\perp):=(-z_l)^{\frac{1}{2}}>0$ for $l=1,2,3$; in particular $\omega_3(p_\perp)<m$.
On the null set of the remaining $p_\perp$ put $\omega_l(p_\perp):=1$.

\emph{No pole where $z=-M_j$.}
The exchange of the second and the third coordinates reduces the case $j=3$ to the case $j=2$:
if $g\in\Ocub$ is the transposition of $e_2$ and $e_3$, then $M(p_1,p_3,p_2)=g\,M(p_1,p_2,p_3)\,g^{T}$ by the identity $M(gp)=g\,M(p)\,g^{T}$,
the pair $(p_3,p_2)$ satisfies the same conditions as $(p_2,p_3)$,
and $M_3$, formed at $(p_2,p_3)$, equals $M_2$, formed at $(p_3,p_2)$.
For $j=2$, the function $X:=1/d_2$ tends to infinity as $z\to-M_2$, since the numerator $z+m^2$ of $1/d_2$ tends to $m^2-M_2\neq0$.
Write
\[
D^{-1}p=X\,p_2\,e_2+v,
\qquad
\nu=p_2^2X+N_0,
\qquad
D^{-1}=X\,e_2e_2^{T}+K_0,
\]
with
\[
v:=\left(\frac{p_1}{d_1},0,\frac{p_3}{d_3}\right),
\qquad
N_0:=\frac{p_1^2}{d_1}+\frac{p_3^2}{d_3},
\qquad
K_0:=\diag\left(\frac{1}{d_1},0,\frac{1}{d_3}\right);
\]
these are holomorphic in $p_1$ near the points with $p_1^2=-M_2$, because $M_3\neq M_2$.
From $(p_2^2X+N_0)^{-1}=(p_2^2X)^{-1}\left(1-N_0/(p_2^2X)+O(X^{-2})\right)$ one obtains
\[
\frac{D^{-1}p\,p^{T}D^{-1}}{\nu}=X\,e_2e_2^{T}-\frac{N_0}{p_2^2}\,e_2e_2^{T}+\frac{e_2v^{T}+v\,e_2^{T}}{p_2}+O\left(X^{-1}\right).
\]
In the difference $M=D^{-1}-D^{-1}p\,p^{T}D^{-1}/\nu$ the two contributions of order $X$, both in the entry $(2,2)$, are equal and opposite, and
\[
M=K_0+\frac{N_0}{p_2^2}\,e_2e_2^{T}-\frac{e_2v^{T}+v\,e_2^{T}}{p_2}+O\left(X^{-1}\right).
\]
Consequently $M$ has a finite limit as $p_1$ tends to either of the two points with $p_1^2=-M_2$.

\emph{Residues.}
The function $p_1\mapsto M(p_1,p_\perp)$ is rational.
It is bounded on the real axis, where $0\leq M\leq1$, and therefore bounded at infinity.
The matrix $D^{-1}$ is holomorphic in $p_1$ except at the points with $p_1^2\in\{-M_2,-M_3\}$, where $M$ has no pole;
so the poles of $M$ lie among the zeros of $p_1\mapsto\tilde\nu(p_1^2)$,
which are the simple zeros $p_1=\pm i\omega_l$, $l=1,2,3$ (as $\tilde\nu'(z_l)\neq0$ and $z_l\neq0$).
The partial fractions of $M$ in $p_1$ therefore read
\[
\begin{aligned}
M&=C_{p_\perp}+\sum_{l=1}^{3}\left(\frac{r^+_l}{p_1-i\omega_l}+\frac{r^-_l}{p_1+i\omega_l}\right)
=C_{p_\perp}+\sum_{l=1}^{3}\frac{N_l(p_1,p_\perp)}{p_1^2+\omega_l^2},\\
N_l(p_1,p_\perp)&:=\left(r^+_l+r^-_l\right)p_1+i\omega_l\left(r^+_l-r^-_l\right),
\end{aligned}
\]
with matrices $C_{p_\perp}$ and $r^\pm_l$ that do not depend on $p_1$:
this is \eqref{eq:rotation-partial-fractions} with three poles.
Since $p_1^2+\omega_l^2=z-z_l$ and $D^{-1}$ is holomorphic at $z_l\notin\{-M_2,-M_3\}$,
\begin{equation}\label{eq:rotation-dielectric-residue}
\begin{aligned}
N_l(-i\omega_l,p_\perp)=-2i\omega_l\,r^-_l
&=\lim_{p_1\to-i\omega_l}\left(p_1^2+\omega_l^2\right)M(p_1,p_\perp)\\
&=-v^l(v^l)^{T}\lim_{z\to z_l}\frac{z-z_l}{\tilde\nu(z)}=-\frac{v^l(v^l)^{T}}{\tilde\nu'(z_l)},
\end{aligned}
\end{equation}
where $v^l$ is the value of $D^{-1}p$ at $p_1=-i\omega_l$, that is,
\[
v^l=\left(-\frac{i\omega_l}{d_1},\frac{p_2}{d_2},\frac{p_3}{d_3}\right)\qquad\text{with $d_2,d_3$ taken at $z=z_l$.}
\]
The first coordinate of $v^l$ is purely imaginary, and the other two are real; thus $\theta_1v^l=\overline{v^l}$,
and the matrix of Lemma~\ref{lem:time-kernels} is
\[
A_l=-\theta_1N_l(-i\omega_l,p_\perp)=\frac{\theta_1v^l(v^l)^{T}}{\tilde\nu'(z_l)}=\frac{\overline{v^l}\,\left(\overline{v^l}\right)^{\dagger}}{\tilde\nu'(z_l)},
\]
that is, $(A_l)_{ab}=\overline{v^l_a}\,v^l_b/\tilde\nu'(z_l)$:
a Hermitian positive semidefinite matrix of rank one.
The numerator of $\tilde\nu$ is a cubic polynomial whose coefficients depend continuously on $p_\perp$;
its roots $z_1<z_2<z_3$, ordered by size, are measurable functions of $p_\perp$, and so are the $\omega_l$.
As $M(\overline{p_1},p_\perp)=\overline{M(p_1,p_\perp)}$, the residues satisfy $r^+_l=\overline{r^-_l}$;
so $N_l(p_1,p_\perp)=2\operatorname{Re}(r^-_l)\,p_1+2\omega_l\operatorname{Im}(r^-_l)$,
and by \eqref{eq:rotation-dielectric-residue} the coefficients of $N_l$, like the matrices $A_l$, are measurable functions of $p_\perp$.
By Lemma~\ref{lem:time-kernels}, $\gamma^{\mathrm{d}}_{\kappa,m}$ satisfies $(R_1)$ and $(RP_1)$.
The law is invariant under the permutation matrices, and Lemma~\ref{lem:transport}\eqref{it:transport-properties} gives $(RP_j)$ for $j=1,2,3$.
Together with the properties established first, $\gamma^{\mathrm{d}}_{\kappa,m}$ satisfies \textup{(H)} and \textup{(dF)};
being a centred Gaussian law, it is invariant under $F\mapsto-F$.

\emph{No cone.}
Let $b_*$ be the test field of Definition~\ref{def:test-field}.
As shown after that definition, for almost every $p_\perp$ the vector $L_3(p_\perp)$ of Lemma~\ref{lem:time-kernels}, formed with $b=b_*$,
is a non-zero multiple of $e_1$, so that $L_3^{\dagger}A_3L_3=|v^3_1|^2\cdot|(L_3)_1|^2/\tilde\nu'(z_3)$, where $(L_3)_1$ is the first coordinate of $L_3$.
By that lemma, $\varrho_{[\mathbb{B}(b_*)]_1}$ is the sum of three positive measures,
one of which is the image under $p_\perp\mapsto(\omega_3(p_\perp),-p_\perp)$ of the measure
\[
\frac{|v^3_1|^2\cdot|(L_3)_1|^2}{(2\pi)^2\cdot2\omega_3\cdot\tilde\nu'(z_3)}\,\mathrm{d}^2p_\perp ;
\]
its density is positive almost everywhere, because $v^3_1=-i\omega_3/d_1\neq0$.
As $\omega_3<m$, this image gives positive mass to the set $\{(\omega,p_\perp):\ \omega<m,\ |p_2|>C\}$ for every real $C$,
and so $\mathfrak{E}_1(\{\omega<m,\ |p_2|>C\})\neq0$ for every real $C$.
Given $c>0$ and a real $C$, the set $\{\omega<m,\ |p_2|>(m+|C|)/c\}$ is contained in $\{\omega<c|p_2|-C\}$;
hence $\mathfrak{E}_1(\{\omega<c|p_2|-C\})\neq0$,
and by countable additivity $\mathfrak{E}_1(S)\neq0$ for $S:=\{\omega<c|p_2|-C,\ |p_2|<R\}$ and some $R>0$.
For a unit vector $\psi$ in the range of $\mathfrak{E}_1(S)$ the measure $\varrho_\psi$ is carried by $S$, on which $\omega$ and $|p_2|$ are bounded, and
\[
\left\langle\psi,\left(c|\mathsf{P}_2|-C-\mathsf{H}_1\right)\psi\right\rangle=\int\mathrm{d}\varrho_\psi(\omega,p_\perp)\,\left(c|p_2|-C-\omega\right)>0 :
\]
the inequality $\mathsf{H}_1\geq c|\mathsf{P}_2|-C$ fails.

If $\gamma^{\mathrm{d}}_{\kappa,m}$ were invariant under $\mathrm{O}(3)$,
then Theorem~\ref{thm:cone-equivalences}, \eqref{it:equiv-orthogonal}$\Rightarrow$\eqref{it:equiv-all-cones}, would give $(C_{12})$,
that is $\mathfrak{E}_1(\{\omega<|p_2|\})=0$;
but $\{\omega<m,\ |p_2|>m\}$ is a subset of $\{\omega<|p_2|\}$, and its $\mathfrak{E}_1$-measure is not $0$.
So $\gamma^{\mathrm{d}}_{\kappa,m}$ is not invariant under $\mathrm{O}(3)$.

The same conclusion can be read off from the multiplier.
Suppose that $\gamma^{\mathrm{d}}_{\kappa,m}$ is invariant under the rotation $R^{23}_{\pi/4}$.
Lemma~\ref{lem:multiplier-laws} gives $M\left(R^{23}_{\pi/4}p\right)=R^{23}_{\pi/4}\,M(p)\,\left(R^{23}_{\pi/4}\right)^{T}$ for almost every $p$,
hence for every $p\neq0$, since $M$ is continuous on $\R^3\setminus\{0\}$
($D$ is continuous because $m>0$, and $\nu>0$ off the origin).
As $R^{23}_{\pi/4}$ fixes $e_1$, it follows that $M_{11}\left(R^{23}_{\pi/4}p\right)=M_{11}(p)$ for all $p\neq0$.
For $p$ orthogonal to $e_1$ one has $M_{11}(p)=1/d_1(p)$.
At $p=te_2$ and at $R^{23}_{\pi/4}p=2^{-\frac{1}{2}}t\,(e_2+e_3)$, $t>0$, this is
\[
\left(1+\frac{\kappa}{t^2+m^2}+\frac{\kappa}{m^2}\right)^{-1}
\qquad\text{and}\qquad
\left(1+\frac{4\kappa}{t^2+2m^2}\right)^{-1},
\]
respectively; as $t\to\infty$ the first value tends to $\left(1+\kappa/m^2\right)^{-1}<1$ and the second to $1$,
which contradicts $M_{11}\left(R^{23}_{\pi/4}p\right)=M_{11}(p)$ for large $t$.
\end{proof}

\subsubsection{Stacks}\label{sec:stacks}

The laws of Proposition~\ref{prop:sharpness}\eqref{it:sharp-planes} are centred Gaussian laws
whose multipliers are built from projections of rank one.
For $k\in\{1,2,3\}$ and $p\in\R^3\setminus\R e_k$ let
\[
M^{(k)}(p):=\frac{u\,u^{T}}{|u|^{2}},\qquad u:=e_k\times p,
\]
and let $M^{(k)}(p):=0$ for $p\in\R e_k$.
Thus $M^{(k)}(p)$ is the orthogonal projection of $\R^3$ onto the line $\R\,(e_k\times p)$,
the direction orthogonal to $e_k$ and to $p$.
The function $M^{(k)}$ is measurable (it is continuous off $\R e_k$), even and bounded by $1$,
and its values are real symmetric positive semidefinite matrices;
so Lemma~\ref{lem:multiplier-laws} defines the centred Gaussian laws
\[
\gamma^{\mathrm s}_k:=\gamma_{M^{(k)}}\quad(k=1,2,3),
\qquad
\bar\gamma^{\mathrm s}:=\gamma_{\bar M},\quad
\bar M:=\tfrac13\left(M^{(1)}+M^{(2)}+M^{(3)}\right).
\]
The laws $\gamma^{\mathrm s}_k$ are the \emph{stacks},
and $\bar\gamma^{\mathrm s}$ is the law of the \emph{three stacks}.
(Formally, under $\gamma^{\mathrm s}_3$ the components $\mathbb{B}_1$ and $\mathbb{B}_2$ of the field
are the derivatives $\partial_2$ and $-\partial_1$ of one scalar field,
which is a massless free field in the variables $(x_1,x_2)$ and white noise in $x_3$,
while $\mathbb{B}_3=0$.)

\paragraph{Common properties.}
For $p\neq0$ the vector $e_k\times p$ is orthogonal to $p$; hence
\[
0\le M^{(k)}(p)\le P^{\perp}(p)\le 1,\qquad M^{(k)}(p)\,p=0,
\]
and the same holds for $\bar M$.
By Lemma~\ref{lem:multiplier-laws} the laws $\gamma^{\mathrm s}_1$, $\gamma^{\mathrm s}_2$, $\gamma^{\mathrm s}_3$ and $\bar\gamma^{\mathrm s}$
satisfy (G) and (dF);
their covariances are invariant under translations (same lemma),
so that they satisfy $(T_{123})$ as well
(Lemma~\ref{lem:gaussian-covariance}\eqref{it:gaussian-laws-criteria}).
Every $g\in\mathrm{O}(3)$ satisfies $g(x\times x')=\det(g)\cdot(gx)\times(gx')$ for $x,x'\in\R^3$.
If $g\in\mathrm{O}(3)$ and $ge_k=\pm e_{k'}$,
then $e_{k'}\times gp=\pm\det(g)\cdot g\,(e_k\times p)$, and consequently
\begin{equation}\label{eq:rotation-stack-covariance}
M^{(k')}(gp)=g\,M^{(k)}(p)\,g^{T}\qquad(p\in\R^3).
\end{equation}
Each $g\in\Ocub$ maps $e_k$ to $\pm e_{\mathfrak{s}(k)}$, $k=1,2,3$,
for a permutation $\mathfrak{s}$ of $\{1,2,3\}$;
by \eqref{eq:rotation-stack-covariance}, $M^{(\mathfrak{s}(k))}(gp)=g\,M^{(k)}(p)\,g^{T}$,
so the elements of $\Ocub$ permute the three multipliers.
If $ge_3=\pm e_3$, then \eqref{eq:rotation-stack-covariance} reads
$M^{(3)}(gp)=g\,M^{(3)}(p)\,g^{T}$,
and $\gamma^{\mathrm s}_3$ is invariant under $g$ by Lemma~\ref{lem:multiplier-laws}.
In particular $\gamma^{\mathrm s}_3$ is invariant under all rotations $R^{12}_t$,
under $\theta_1$, $\theta_2$, $\theta_3$,
and under the permutation matrix $g_{21}$, which exchanges $e_1$ and $e_2$ and fixes $e_3$.

\paragraph{Reflection positivity and cones for $\gamma^{\mathrm s}_3$.}
Write $p=(p_1,p_\perp)$ with $p_\perp=(p_2,p_3)$.
Since $e_3\times p=(-p_2,p_1,0)$,
the entries of $M^{(3)}(p)$ that are not identically zero are
\[
M^{(3)}_{11}(p)=\frac{p_2^2}{p_1^2+p_2^2},\qquad
M^{(3)}_{12}(p)=M^{(3)}_{21}(p)=-\frac{p_1p_2}{p_1^2+p_2^2},\qquad
M^{(3)}_{22}(p)=\frac{p_1^2}{p_1^2+p_2^2}=1-\frac{p_2^2}{p_1^2+p_2^2}.
\]
For $p_2\neq0$ this is the decomposition \eqref{eq:rotation-partial-fractions} with one pole:
\[
M^{(3)}(p_1,p_\perp)=Q_{p_\perp}(p_1)+\frac{N_1(p_1,p_\perp)}{p_1^2+\omega_1(p_\perp)^2},
\qquad
Q_{p_\perp}(p_1):=e_2e_2^{T},\quad
\omega_1(p_\perp):=|p_2|,
\]
with the numerator
\[
N_1(p_1,p_\perp):=p_2^2\left(e_1e_1^{T}-e_2e_2^{T}\right)-p_1p_2\left(e_1e_2^{T}+e_2e_1^{T}\right).
\]
The symmetry $\theta_1M^{(3)}(-p_1,p_\perp)\,\theta_1=M^{(3)}(p_1,p_\perp)$
that Lemma~\ref{lem:time-kernels} requires
is \eqref{eq:rotation-stack-covariance} for $g=\theta_1$.
At the pole $p_1=-i|p_2|$ the signed residue matrix is
\[
A_1(p_\perp)=-\theta_1N_1(-i|p_2|,p_\perp)
=p_2^2\left(e_1e_1^{T}+e_2e_2^{T}\right)+i\,|p_2|\,p_2\left(e_1e_2^{T}-e_2e_1^{T}\right).
\]
Its block in the coordinates $1,2$ is Hermitian, with determinant $0$ and trace $2p_2^2$,
and its other entries vanish;
so $A_1(p_\perp)$ is positive semidefinite, of rank one.
By Lemma~\ref{lem:time-kernels}, $\gamma^{\mathrm s}_3$ satisfies $(R_1)$ and $(RP_1)$,
and for every real $b\in C^{\infty}_{\mathrm c}(\{x_1>0\};\R^3)$
the spectral measure of $[\mathbb{B}(b)]_1$ is the image of the positive measure
\[
(2\pi)^{-2}\,(2|p_2|)^{-1}\,L_1(p_\perp)^{\dagger}A_1(p_\perp)\,L_1(p_\perp)\,\mathrm{d}^2p_\perp
\]
under $p_\perp\mapsto(|p_2|,-p_\perp)$, where $L_1$ is formed from $b$ as in that lemma.
These one-field spectral measures are carried by the shell $\{\omega=|p_2|\}$,
on which $p_3$ is free.
In particular they vanish on $\{\omega<|p_2|\}$,
and Lemma~\ref{lem:gaussian-cone}, with $c=1$ and $i=2$, gives $(C_{12})$.

In the picture $e_2$ the same properties follow by transport.
The law $\gamma^{\mathrm s}_3$ is invariant under $g:=R^{12}_{\pi/2}$, and $ge_1=e_2$;
Lemma~\ref{lem:transport}\eqref{it:transport-properties},
applied to $\gamma^{\mathrm s}_3=(g^{-1})_*\gamma^{\mathrm s}_3$,
turns $(RP_1)$ into $(RP_2)$.
So $\gamma^{\mathrm s}_3$ has $(T_{123})$, $(R_2)$ and $(RP_2)$, and it is invariant under $g_{21}$;
Lemma~\ref{lem:transport}\eqref{it:transport-pairs},
applied with the pair $(2,1)$ to $\gamma^{\mathrm s}_3=(g_{21}^{-1})_*\gamma^{\mathrm s}_3$,
shows that $(C_{12})$ for $\gamma^{\mathrm s}_3$ is $(C_{21})$ for $\gamma^{\mathrm s}_3$.
Hence $(C_{21})$ holds.

Reflection positivity at the plane $\{x_3=0\}$ rests on the fact that $M^{(3)}$ does not depend on $p_3$;
write $M^{(3)}(p_1,p_2)$ for its value.
For real $b,b'\in\mathcal{S}(\R^3;\R^3)$ let $\hat b_\parallel(p_1,p_2,x_3)$
denote the Fourier transform of $b$ in the variables $(x_1,x_2)$.
Plancherel's theorem in the variable $p_3$ gives
\[
C_{M^{(3)}}(b',b)=(2\pi)^{-2}\int\mathrm{d}p_1\,\mathrm{d}p_2\int\mathrm{d}x_3\;
\widehat{b'}_{\parallel}(p_1,p_2,x_3)^{\dagger}\,M^{(3)}(p_1,p_2)\,\hat b_\parallel(p_1,p_2,x_3).
\]
If $b'$ vanishes on $\{x_3>0\}$ and $b$ vanishes on $\{x_3<0\}$,
the integrand vanishes identically, and $C_{M^{(3)}}(b',b)=0$.
For $b\in C^{\infty}_{\mathrm c}(\{x_3>0\};\R^3)$ the field $\theta_3\cdot b$ vanishes on $\{x_3>0\}$;
thus $C_{M^{(3)}}(\theta_3\cdot b,b)=0\ge0$,
and Lemma~\ref{lem:gaussian-covariance}\eqref{it:gaussian-laws-positivity},
together with the invariance under $\theta_3$, gives $(RP_3)$.

\begin{proof}[Proof of Proposition~\ref{prop:sharpness}\eqref{it:sharp-stack}]
By the common properties and the verifications for $\gamma^{\mathrm s}_3$ above,
the stack $\gamma^{\mathrm s}_3$ satisfies (G), (dF), $(T_{123})$, $(R_1)$, $(RP_1)$ and $(C_{12})$,
and it is invariant under the quarter turn $R^{12}_{\pi/2}$, which is $(Q_{12})$.
These are the hypotheses of Theorem~\ref{thm:one-plane} for the pair $(1,2)$
(Definition~\ref{def:closedness-cones}\eqref{it:one-plane-hypotheses}).
In accordance with that theorem, $\gamma^{\mathrm s}_3$ is invariant under every rotation $R^{12}_t$,
as the common properties show directly.

To see that $(C_{13})$ fails, let $b_*$ be the test field of Definition~\ref{def:test-field},
and let $L_1$ be formed from $b_*$ as in Lemma~\ref{lem:time-kernels}, with the pole $\omega_1(p_\perp)=|p_2|$.
For almost every $p_\perp$ the vector $L_1(p_\perp)$ is a non-zero multiple of $e_1$;
hence $L_1^{\dagger}A_1L_1=(A_1)_{11}\cdot|L_1|^2=p_2^2\cdot|L_1|^2$.
By Lemma~\ref{lem:time-kernels}, applied as above,
the spectral measure $\varrho_\psi$ of $\psi:=[\mathbb{B}(b_*)]_1$
is the image under $p_\perp\mapsto(|p_2|,-p_\perp)$ of the measure with the density
\[
(2\pi)^{-2}\,(2|p_2|)^{-1}\,p_2^2\cdot|L_1(p_\perp)|^2,
\]
which is positive for almost every $p_\perp$.
The preimage of $\{\omega<|p_3|\}$ under this map is the open set $\{|p_2|<|p_3|\}$,
which has positive Lebesgue measure;
so $\varrho_\psi(\{\omega<|p_3|\})>0$.
Under $(C_{13})$ one would have
$\varrho_\psi(\{\omega<|p_3|\})=\langle\psi,\mathfrak{E}_1(\{\omega<|p_3|\})\,\psi\rangle=0$.
Therefore $\gamma^{\mathrm s}_3$ violates $(C_{13})$.

Suppose finally that $\gamma^{\mathrm s}_3$ were invariant under $g:=R^{13}_{\pi/2}$.
Then $M^{(3)}(gp)=g\,M^{(3)}(p)\,g^{T}$ for almost every $p$ (Lemma~\ref{lem:multiplier-laws}).
As $ge_1=e_3$, that is, $g^{T}e_3=e_1$,
the entries in the place $(3,3)$ give
\[
M^{(3)}_{33}(gp)=\left(g\,M^{(3)}(p)\,g^{T}\right)_{33}=M^{(3)}_{11}(p)\qquad\text{for almost every }p.
\]
This is impossible:
$M^{(3)}_{33}=0$, while $M^{(3)}_{11}(p)=p_2^2/(p_1^2+p_2^2)\neq0$ whenever $p_2\neq0$,
that is, for almost every $p$.
\end{proof}

\begin{remark}
The stack $\gamma^{\mathrm s}_3$ violates $(C_{23})$ as well,
and it is not invariant under the quarter turn $R^{23}_{\pi/2}$.
Indeed, the permutation matrix $g_{13}$ exchanges $e_2$ and $e_3$,
and $g_{23}$ maps $e_1$, $e_2$, $e_3$ to $e_2$, $e_3$, $e_1$;
put $\mu_{13}:=(g_{13}^{-1})_*\gamma^{\mathrm s}_3$ and $\mu_{23}:=(g_{23}^{-1})_*\gamma^{\mathrm s}_3$.
Since $\gamma^{\mathrm s}_3$ has $(T_{123})$, $(R_1)$, $(RP_1)$, $(R_2)$ and $(RP_2)$,
Lemma~\ref{lem:transport}\eqref{it:transport-pairs} shows that
$(C_{13})$ for $\gamma^{\mathrm s}_3$ is $(C_{12})$ for $\mu_{13}$,
and that $(C_{23})$ for $\gamma^{\mathrm s}_3$ is $(C_{12})$ for $\mu_{23}$.
The matrix $g_{23}\,g_{13}^{-1}$ maps $e_1$ to $e_2$, $e_2$ to $e_1$ and $e_3$ to $e_3$;
so $\gamma^{\mathrm s}_3$ is invariant under it, and $\mu_{23}=\mu_{13}$.
As $(C_{13})$ fails for $\gamma^{\mathrm s}_3$, so does $(C_{23})$.
Next, $g:=R^{23}_{\pi/2}$ maps $e_2$ to $e_3$;
invariance under $g$ would give $M^{(3)}_{33}(gp)=M^{(3)}_{22}(p)$ for almost every $p$,
which is impossible, since $M^{(3)}_{33}=0$
and $M^{(3)}_{22}(p)=p_1^2/(p_1^2+p_2^2)\neq0$ for almost every $p$.
So among the quarter turns of the three coordinate planes only those of the plane $(1,2)$
leave $\gamma^{\mathrm s}_3$ invariant,
and $\gamma^{\mathrm s}_3$ violates both $(C_{13})$ and $(C_{23})$.
\end{remark}

\begin{proof}[Proof of Proposition~\ref{prop:sharpness}\eqref{it:sharp-three-stacks}]
First, $\bar\gamma^{\mathrm s}$ satisfies (H) and (dF).
The common properties give (G), (dF) and $(T_{123})$;
and since the elements of $\Ocub$ permute the three multipliers,
$\bar M(gp)=g\,\bar M(p)\,g^{T}$ for $g\in\Ocub$,
so that $\bar\gamma^{\mathrm s}$ is invariant under $\Ocub$ (Lemma~\ref{lem:multiplier-laws}).
For reflection positivity fix $j\in\{1,2,3\}$.
Each stack is invariant under $\theta_j$,
by \eqref{eq:rotation-stack-covariance} with $g=\theta_j$, since $\theta_je_k=\pm e_k$ (Lemma~\ref{lem:multiplier-laws}).
The stack $\gamma^{\mathrm s}_3$ satisfies $(RP_1)$, $(RP_2)$ and $(RP_3)$, as shown above.
For $k\in\{1,2\}$ let $g\in\Ocub$ be the permutation matrix that exchanges $e_k$ and $e_3$.
By \eqref{eq:rotation-stack-covariance}, $M^{(3)}(gp)=g\,M^{(k)}(p)\,g^{T}$ for $p\in\R^3$.
Let $b\in\mathcal{S}(\R^3;\R^3)$.
The field $g\cdot b$ is $x\mapsto\det(g)\cdot g\,b(g^{-1}x)$ (Definition~\ref{def:field-actions});
as $g$ preserves Lebesgue measure and $p\cdot gx=(g^{-1}p)\cdot x$,
its Fourier transform is $p\mapsto\det(g)\cdot g\,\hat b(g^{-1}p)$.
Inserting this into the integral that defines $C_{M^{(3)}}$ (Lemma~\ref{lem:multiplier-laws})
and substituting $gp$ for $p$, one obtains, since $\det(g)^2=1$,
\[
\begin{aligned}
C_{M^{(3)}}(g\cdot b,g\cdot b)
&=(2\pi)^{-3}\int\mathrm{d}p\,\hat b(p)^{\dagger}g^{T}M^{(3)}(gp)\,g\,\hat b(p)\\
&=(2\pi)^{-3}\int\mathrm{d}p\,\hat b(p)^{\dagger}M^{(k)}(p)\,\hat b(p)=C_{M^{(k)}}(b,b),
\end{aligned}
\]
where the second equality holds because $M^{(3)}(gp)=g\,M^{(k)}(p)\,g^{T}$ and $g^{T}g=1$.
By Lemma~\ref{lem:transport}\eqref{it:transport-operators},
the law $(g^{-1})_*\gamma^{\mathrm s}_3$ assigns to $\exp(i\mathbb{B}(b))$ the value
\[
\exp\left(-\tfrac12\,C_{M^{(3)}}(g\cdot b,g\cdot b)\right)=\exp\left(-\tfrac12\,C_{M^{(k)}}(b,b)\right),
\]
so that $(g^{-1})_*\gamma^{\mathrm s}_3=\gamma^{\mathrm s}_k$
(Lemma~\ref{lem:gaussian-covariance}\eqref{it:gaussian-laws-criteria}).
Now $ge_j=e_{j'}$ for some $j'\in\{1,2,3\}$,
and Lemma~\ref{lem:transport}\eqref{it:transport-properties}
carries $(RP_{j'})$ for $\gamma^{\mathrm s}_3$ to $(RP_j)$ for $\gamma^{\mathrm s}_k$.
Thus each stack satisfies $(R_j)$ and $(RP_j)$,
and Lemma~\ref{lem:gaussian-covariance}\eqref{it:gaussian-laws-positivity}
gives $C_{M^{(k)}}(\theta_j\cdot b,b)\ge0$ for $k=1,2,3$
and every real $b\in C^{\infty}_{\mathrm c}(\{x_j>0\};\R^3)$.
The covariance $C_{\bar M}=\tfrac13\left(C_{M^{(1)}}+C_{M^{(2)}}+C_{M^{(3)}}\right)$
is therefore invariant under $\theta_j$ and satisfies $C_{\bar M}(\theta_j\cdot b,b)\ge0$ for the same $b$;
by Lemma~\ref{lem:gaussian-covariance}\eqref{it:gaussian-laws-positivity},
$\bar\gamma^{\mathrm s}$ satisfies $(RP_j)$.
With (G) and the invariance under translations and under $\Ocub$, this is (H).

Now let $\lambda=(3,\chi)$ with $\chi\in L^2(\R;\R)$, $\chi\neq0$.
Let $\eta,\eta'\in C^{\infty}_{\mathrm c}(\R)$ be non-negative and not identically zero,
with $\supp\eta\subset(-2,-1)$ and $\supp\eta'\subset(1,2)$,
let $h\in L^2(\R)$ be real with compact support, and put
\[
\alpha(y):=\eta(y^0)\cdot h(y^1),\qquad\alpha'(y):=\eta'(y^0)\cdot h(y^1)\qquad\left(y=(y^0,y^1)\in\R^2\right).
\]
The plates $\alpha\otimes\lambda$ and $\alpha'\otimes\lambda$ are the real fields
$x\mapsto\alpha(y)\cdot\chi(x_3)\cdot e_3$ and $x\mapsto\alpha'(y)\cdot\chi(x_3)\cdot e_3$,
and $\alpha$, $\alpha'$ have their supports in $\{-2<y^0<-1\}$ and in $\{1<y^0<2\}$.
Under the centred Gaussian law $\bar\gamma^{\mathrm s}$ one has $S_2(b,b')=C_{\bar M}(b,b')$
for real $b,b'\in\mathcal{S}(\R^3;\R^3)$
(Lemma~\ref{lem:gaussian-covariance}\eqref{it:gaussian-laws-criteria});
both sides are bounded bilinear forms on the real fields in $L^2(\R^3;\R^3)$
(Lemma~\ref{lem:domination-consequences}\eqref{it:moment-forms}, and $\bar M\le1$),
so the identity extends to these fields.
The Fourier transform of $\alpha\otimes\lambda$ is the product of $\hat\alpha(p_1,p_2)$,
of the transform of $\chi$ at $p_3$ and of $e_3$,
and $\hat\alpha(p_1,p_2)=\hat\eta(p_1)\cdot\hat h(p_2)$;
moreover $\overline{\hat\eta(p_1)}=\hat\eta(-p_1)$, as $\eta$ is real.
Only the entry $\bar M_{33}=\frac13\left(M^{(1)}_{33}+M^{(2)}_{33}+M^{(3)}_{33}\right)$ enters, and
\[
M^{(1)}_{33}(p)=\frac{p_2^2}{p_2^2+p_3^2},\qquad
M^{(2)}_{33}(p)=\frac{p_1^2}{p_1^2+p_3^2},\qquad
M^{(3)}_{33}(p)=0,
\]
because $e_1\times p=(0,-p_3,p_2)$, $e_2\times p=(p_3,0,-p_1)$ and $e_3\times p=(-p_2,p_1,0)$.
Hence
\[
S_2(\alpha\otimes\lambda,\alpha'\otimes\lambda)
=(2\pi)^{-3}\int\mathrm{d}p\;\hat\eta(-p_1)\cdot\hat\eta'(p_1)\cdot|\hat h(p_2)|^2\cdot|\hat\chi(p_3)|^2
\cdot\tfrac13\left(M^{(1)}_{33}(p)+M^{(2)}_{33}(p)\right).
\]
The integrand is integrable:
$\hat\eta$ and $\hat\eta'$ are Schwartz functions, $|\hat h|^2$ and $|\hat\chi|^2$ are integrable,
and $0\le M^{(k)}_{33}\le1$.
The term with $M^{(1)}_{33}$, which does not depend on $p_1$, vanishes, because by Parseval's formula
\[
\int\mathrm{d}p_1\,\hat\eta(-p_1)\cdot\hat\eta'(p_1)=2\pi\int\mathrm{d}t\,\eta(t)\cdot\eta'(t)=0,
\]
the supports of $\eta$ and $\eta'$ being disjoint.
In the term with $M^{(2)}_{33}=1-p_3^2/(p_1^2+p_3^2)$ the constant $1$ contributes $0$ for the same reason.
For $\omega>0$ and real $s$ one has
\[
(2\pi)^{-1}\int\mathrm{d}p_1\,\frac{\exp(-ip_1s)}{p_1^2+\omega^2}=\frac{\exp(-\omega|s|)}{2\omega},
\qquad
\hat\eta(-p_1)\cdot\hat\eta'(p_1)=\int\mathrm{d}t\int\mathrm{d}t'\,\eta(t)\cdot\eta'(t')\cdot\exp\left(-ip_1(t'-t)\right);
\]
Fubini's theorem gives, for $p_3\neq0$,
\[
\int\mathrm{d}p_1\,\hat\eta(-p_1)\cdot\hat\eta'(p_1)\cdot M^{(2)}_{33}(p)
=-2\pi\cdot\frac{|p_3|}{2}\int\mathrm{d}t\int\mathrm{d}t'\,\eta(t)\cdot\eta'(t')\cdot\exp\left(-|p_3|\,|t'-t|\right).
\]
Finally $\int\mathrm{d}p_2\,|\hat h(p_2)|^2=2\pi\|h\|_2^2$ by Plancherel's theorem.
Collecting the three integrations,
\begin{equation}\label{eq:rotation-stacks-moment}
S_2(\alpha\otimes\lambda,\alpha'\otimes\lambda)
=-\frac13\,\|h\|_2^2\cdot\int\mathrm{d}t\int\mathrm{d}t'\,\eta(t)\cdot\eta'(t')\cdot G(t'-t),
\qquad
G(t):=(2\pi)^{-1}\int\mathrm{d}p_3\,|\hat\chi(p_3)|^2\cdot\frac{|p_3|}{2}\cdot\exp\left(-|p_3|\,|t|\right).
\end{equation}
On the product of the supports of $\eta$ and $\eta'$ one has $2<t'-t<4$.
For $|t|\ge2$ the function $G$ is finite, since $|p_3|\exp(-|p_3|\,|t|)$ is bounded,
and positive, since $\hat\chi$ does not vanish almost everywhere ($\chi\neq0$);
so all the interchanges of integrations above concern absolutely convergent integrals,
and the double integral in \eqref{eq:rotation-stacks-moment} is positive,
$\eta$ and $\eta'$ being non-negative, continuous and not identically zero.

Choose $h:=h_n:=n\cdot1_{[0,1/n]}$, $n\ge1$,
and let $\alpha_n$, $\alpha'_n$ be the corresponding functions on the plane.
Then $\|\alpha_n\|_{L^1(\R^2)}=\|\eta\|_{L^1(\R)}$, $\|\alpha'_n\|_{L^1(\R^2)}=\|\eta'\|_{L^1(\R)}$ and $\|h_n\|_2^2=n$,
and \eqref{eq:rotation-stacks-moment} gives
\[
\frac{|S_2(\alpha_n\otimes\lambda,\alpha'_n\otimes\lambda)|}{\|\alpha_n\|_{L^1(\R^2)}\cdot\|\alpha'_n\|_{L^1(\R^2)}}
=\frac n3\cdot\frac{\int\mathrm{d}t\int\mathrm{d}t'\,\eta(t)\cdot\eta'(t')\cdot G(t'-t)}{\|\eta\|_{L^1(\R)}\cdot\|\eta'\|_{L^1(\R)}}
\longrightarrow\infty\qquad(n\to\infty).
\]
So the supremum in Proposition~\ref{prop:sharpness}\eqref{it:sharp-three-stacks} is infinite.
For a law with the hypotheses of Proposition~\ref{prop:point-plates}, on the other hand,
the bound \eqref{eq:rotation-kernel-bound} for $N=2$ and the formula \eqref{eq:rotation-kernel-density}
of Proposition~\ref{prop:kernels},
with $d_{\mathcal{W}}(y'-y)=y'^0-y^0>2$ for $y\in\supp\alpha$ and $y'\in\supp\alpha'$, give
\[
|S_2(\alpha\otimes\lambda,\alpha'\otimes\lambda)|\le4k_2^2\cdot\|\chi\|_2^2\cdot2^{-2}\cdot\|\alpha\|_{L^1(\R^2)}\cdot\|\alpha'\|_{L^1(\R^2)}
\]
for all such $\alpha$ and $\alpha'$;
this is the bound that fails for $\bar\gamma^{\mathrm s}$.
\end{proof}

\begin{remark}
The right side of \eqref{eq:rotation-stacks-moment} comes from the stack $\gamma^{\mathrm s}_2$ alone:
the term with $M^{(1)}_{33}$ vanishes for these plates, and $M^{(3)}_{33}=0$.
The multiplier $M^{(2)}$ does not depend on $p_2$, so $\gamma^{\mathrm s}_2$ is white noise in $x_2$.
Its entries in the coordinates $1,3$ are those of $M^{(3)}$ in the coordinates $1,2$,
with $p_3$ in the place of $p_2$, and its other entries vanish;
Lemma~\ref{lem:time-kernels}, applied as for $\gamma^{\mathrm s}_3$, shows that in the picture $e_1$
the one-field spectral measures of $\gamma^{\mathrm s}_2$ are carried by the shell $\{\omega=|p_3|\}$,
on which $p_2$ is free.
\end{remark}

\subsection{Euclidean invariance of the limit points}\label{sec:euclidean-limit-points}

\begin{proof}[Proof of Theorem~\ref{thm:euclidean}, Part~I]
Clause \eqref{it:limit-general-laws} contains (a)--(e) as the case of a volume-first limit point $\mu$ of the family $\left(P^{\mathrm{ch}}_{N,T}\right)$,
provided that $\mu$ satisfies (H) and (dF) and lies in $\mathcal{C}_9$.
So it suffices to verify these three properties and to prove \eqref{it:limit-general-laws}.
Let $\mu$ be a volume-first limit point of the family $\left(P^{\mathrm{ch}}_{N,T}\right)$ (Definition~\ref{def:limit-points}),
in the setting of Theorem~\ref{thm:continuum-limits}.
By clauses \eqref{it:continuum-soft} and \eqref{it:continuum-volume-first} of that theorem,
$\mu$ is a limit point of this family
and has the properties of \cite[Theorem~3.1, Part~I\,(b)]{NoncompactQED}.
Closedness is one of them, so $\mu$ satisfies (dF);
and $\mu$ satisfies (H) by Proposition~\ref{prop:chamber-limits}(c).
It lies in $\mathcal{C}_9$ by Theorem~\ref{thm:continuum-limits}\eqref{it:continuum-nine-planes}.

\emph{The cone for all pairs.}
From now on $\mu$ is a law on $(\mathcal{S}',\Sigma)$ with (H) and (dF) that lies in $\mathcal{C}_9$,
and $\mu=\int\mathrm{d}\nu(\xi)\,\mu_\xi$ is its ergodic decomposition under the translations of $\R^3$
\cite[Theorem~8.1]{NoncompactQED},
written with an index $\xi$ as in \cite[Definition~9.5]{NoncompactQED}.
Let $(j,i)$ be an ordered pair of distinct indices.
As a member of $\mathcal{C}_9$, the law $\mu$ is invariant under $\R^3$, satisfies $(R_j)$ and $(RP_j)$,
and is invariant and reflection positive at the plane $\{x_j+x_i=0\}$.
Theorem~\ref{thm:cone}(c), or Proposition~\ref{prop:diagonal-mirror}, therefore gives $(C_{ji})$;
as the pair is arbitrary, $(C_{ji})$ holds for every ordered pair $(j,i)$ (Proposition~\ref{prop:zero-momentum}(a)).
Since $\Ocub$ contains the reflections $\theta_j$ and the quarter turns,
$\mu$ satisfies, for every ordered pair $(j,i)$, all the conditions
(G), (dF), $(T_{123})$, $(R_j)$, $(RP_j)$, $(Q_{ji})$, $(C_{ji})$
of Definition~\ref{def:closedness-cones}\eqref{it:one-plane-hypotheses}.
In particular, for every ordered pair $(j,i)\neq(1,2)$ the plate vectors of the plane $(j,i)$,
their closed span $\mathcal{H}^{\Psi}_{ji}$ and the boosts $\mathfrak{L}_{ji}$
are defined for $\mu$, as at the end of Section~\ref{sec:unitary-boosts}.

\emph{Clauses (a) and (b).}
Theorem~\ref{thm:orthogonal-group}, applied with the ordered pair $(1,2)$,
shows that $\mu$ is invariant under $\R^3\rtimes\mathrm{O}(3)$,
and that $\mu$ is invariant under the reflection in every plane of $\R^3$ and reflection positive at it,
on all of $L^2$ of the open side.
The second assertion is (b).
Let $\mathfrak{g}$ be a Euclidean map.
Since $\mathfrak{g}_*\mu=\mu$, one has $\mu\left(\hat{\mathfrak{g}}\Phi\right)=\mu(\Phi)$ for $\Phi\in L^1(\mu)$;
the operator $\hat{\mathfrak{g}}$ is multiplicative,
and $\hat{\mathfrak{g}}\,\mathbb{B}(b)=\mathbb{B}(\mathfrak{g}\cdot b)$ in $L^2(\mu)$ for all $b\in L^2(\R^3;\C^3)$
by Lemma~\ref{lem:domination-consequences}(a).
For two-forms $f_1,\dots,f_n\in L^2(\R^3;\Lambda^2)$ put
\[
S_n(f_1,\dots,f_n):=S_n({\ast}f_1,\dots,{\ast}f_n)=\mu\left(F(f_1)\cdots F(f_n)\right).
\]
As $\mathfrak{g}\cdot({\ast}f)={\ast}(\mathfrak{g}\cdot f)$ (Definition~\ref{def:field-actions}), it follows that
\[
S_n(\mathfrak{g}\cdot f_1,\dots,\mathfrak{g}\cdot f_n)
=\mu\left(\hat{\mathfrak{g}}\left(\mathbb{B}({\ast}f_1)\cdots\mathbb{B}({\ast}f_n)\right)\right)
=S_n(f_1,\dots,f_n)
\]
for every $n$, with no condition on the supports of the $f_a$.
This is (a).

\emph{Clause (c): the picture $e_1$ and the round cone.}
By (a), $\mu$ is invariant under $\mathrm{O}(2)_1=\{g\in\mathrm{O}(3):ge_1=e_1\}$.
For $g\in\mathrm{O}(3)$ and a set $A$ of pairs $(\omega,p_\perp)$ write
$gA:=\{(\omega,gp_\perp):(\omega,p_\perp)\in A\}$, the momentum $p_\perp$ being read as a vector of $\R^3$.
Proposition~\ref{prop:rotations-boosts}(a) gives the strongly continuous unitary representation $\mathcal{U}$
of $\mathrm{O}(2)_1$ on $\mathcal{H}_1(\mu)$ with $\mathcal{U}(g)[\Phi]_1=\left[\hat g\Phi\right]_1$;
it fixes $\Omega_1$, commutes with $\exp(-t\cdot\mathsf{H}_1)$ for $t\ge0$,
and satisfies $\mathcal{U}(g)V_1(y)\mathcal{U}(g)^{-1}=V_1(gy)$
and $\mathcal{U}(g)\mathfrak{E}_1(A)\mathcal{U}(g)^{-1}=\mathfrak{E}_1(gA)$.
Proposition~\ref{prop:rotations-boosts}(c) provides the boosts $\mathfrak{L}$ on $\mathcal{H}^{\Psi}_{12}$.
As $\mu$ satisfies (H), (dF) and $(C_{12})$, the last assertion of Proposition~\ref{prop:rotations-boosts}(b) applies to it:
in each picture $e_j$, $j=1,2,3$, the joint spectral measure $\mathfrak{E}_j$ is carried by the round forward cone
$\{\omega\ge|p_\perp|\}$.

\emph{Transport to the picture $e_j$.}
Let $(j,i)$ be an ordered pair of distinct indices,
and let $g_{ji}$ be the permutation matrix with $g_{ji}e_1=e_j$ and $g_{ji}e_2=e_i$.
It lies in $\Ocub$, so that $\left(g_{ji}^{-1}\right)_*\mu=\mu$.
Lemma~\ref{lem:transport}\eqref{it:transport-pairs}, applied with $g=g_{ji}$, in which the transported law is $\mu$ itself,
shows that the map $[\Phi]_1\mapsto\left[\hat g_{ji}\Phi\right]_j$, $\Phi\in L^2(\Sigma_1^+,\mu)$,
extends to a unitary map $J$ of $\mathcal{H}_1(\mu)$ onto $\mathcal{H}_j(\mu)$ with
\begin{equation}\label{eq:rotation-transport-relations}
J\exp(-t\cdot\mathsf{H}_1)J^{-1}=\exp(-t\cdot\mathsf{H}_j),\qquad
J\,V_1(y)\,J^{-1}=V_j(g_{ji}y),\qquad
\mathfrak{E}_j(A)=J\,\mathfrak{E}_1\left(g_{ji}^{-1}A\right)J^{-1}
\end{equation}
for $t\ge0$, $y\in e_1^\perp$ and every Borel set $A\subset[0,\infty)\times e_j^\perp$;
here $g_{ji}^{-1}A$ is the set of the pairs $(\omega,p_\perp)$ with $(\omega,g_{ji}p_\perp)\in A$.
For $(j,i)=(1,2)$ both $g_{12}$ and $J$ are identities.

\emph{Rotations of $e_j^\perp$.}
Let $\mathrm{O}(2)_j:=\{g\in\mathrm{O}(3):ge_j=e_j\}$;
restricted to $e_j^\perp$, its elements are the rotations and reflections of $e_j^\perp$.
For $g\in\mathrm{O}(2)_j$ the matrix $g_{ji}^{-1}gg_{ji}$ lies in $\mathrm{O}(2)_1$; put
\[
\mathcal{U}_j(g):=J\,\mathcal{U}(g_{ji}^{-1}gg_{ji})\,J^{-1}.
\]
Let $\Phi\in L^2(\Sigma_j^+,\mu)$.
The function $\Phi_1:=\hat g_{ji}^{-1}\Phi$ lies in $L^2(\Sigma_1^+,\mu)$ and $J[\Phi_1]_1=[\Phi]_j$;
by Lemma~\ref{lem:transport}(a), the operator $\hat g_{ji}\,\widehat{g_{ji}^{-1}gg_{ji}}\,\hat g_{ji}^{-1}$
is the operator $\hat g$ of $g$.
Therefore
\[
\mathcal{U}_j(g)[\Phi]_j=J\left[\widehat{g_{ji}^{-1}gg_{ji}}\,\Phi_1\right]_1
=\left[\hat g\,\hat g_{ji}\Phi_1\right]_j=\left[\hat g\Phi\right]_j .
\]
In particular $\mathcal{U}_j$ does not depend on $i$, and $\mathcal{U}_1=\mathcal{U}$.
Conjugation by the unitary map $J$ carries the properties of $\mathcal{U}$ over to $\mathcal{U}_j$.
It is a strongly continuous unitary representation of $\mathrm{O}(2)_j$ on $\mathcal{H}_j(\mu)$,
because $g\mapsto g_{ji}^{-1}gg_{ji}$ is a continuous isomorphism of $\mathrm{O}(2)_j$ onto $\mathrm{O}(2)_1$.
It fixes $\Omega_j$, the class of the constant function $1$, which equals $J\Omega_1$.
It commutes with $\exp(-t\cdot\mathsf{H}_j)=J\exp(-t\cdot\mathsf{H}_1)J^{-1}$ for $t\ge0$, hence with $\mathsf{H}_j$.
Finally it rotates the momenta:
since $(g_{ji}^{-1}gg_{ji})\,g_{ji}^{-1}=g_{ji}^{-1}g$,
the relations of $\mathcal{U}$ together with \eqref{eq:rotation-transport-relations} give
\begin{align*}
\mathcal{U}_j(g)V_j(y)\mathcal{U}_j(g)^{-1}&=J\,V_1(g_{ji}^{-1}gy)\,J^{-1}=V_j(gy)\qquad(y\in e_j^\perp),\\
\mathcal{U}_j(g)\mathfrak{E}_j(A)\mathcal{U}_j(g)^{-1}&=J\,\mathfrak{E}_1(g_{ji}^{-1}gA)\,J^{-1}=\mathfrak{E}_j(gA).
\end{align*}

\emph{Boosts.}
For $(j,i)=(1,2)$ the boosts are those of Proposition~\ref{prop:rotations-boosts}(c); let $(j,i)\neq(1,2)$.
Since $\left(g_{ji}^{-1}\right)_*\mu=\mu$, the plate vectors of the plane $(j,i)$ are the vectors $J\Psi(x;\lambda)$,
formed with the map $J$ of the transport and the plate vectors $\Psi(x;\lambda)$ of $\mu$ itself (Definition~\ref{def:plate-kernels});
moreover $\mathcal{H}^{\Psi}_{ji}=J\mathcal{H}^{\Psi}_{12}$ and $\mathfrak{L}_{ji}(\beta)=J\,\mathfrak{L}(\beta)\,J^{-1}$ for real $\beta$,
with the space $\mathcal{H}^{\Psi}_{12}$ and the group $\mathfrak{L}$ of Proposition~\ref{prop:rotations-boosts}(c) for $\mu$.
On $\mathcal{H}^{\Psi}_{ji}$ the relations of that proposition for $\exp(-t\cdot\mathsf{H}_1)$ and $\exp(i\,u\cdot\mathsf{P}_2)$
hold for $\mathfrak{L}_{ji}$ with $(\mathsf{H}_j,\mathsf{P}_i)$ in the place of $(\mathsf{H}_1,\mathsf{P}_2)$,
as shown at the end of Section~\ref{sec:unitary-boosts}.
Conjugation by $J$ carries over the remaining conclusions of the proposition.
The map $J$ sends $\mathcal{H}^{\Psi}_{12}$ onto $\mathcal{H}^{\Psi}_{ji}$ and $\Omega_1$ to $\Omega_j$;
so $\mathcal{H}^{\Psi}_{ji}$ contains $\Omega_j$, the group $\mathfrak{L}_{ji}$ fixes $\Omega_j$,
and $\mathfrak{L}_{ji}$ is the unique strongly continuous unitary group on $\mathcal{H}^{\Psi}_{ji}$ with
\[
\mathfrak{L}_{ji}(\beta)\,J\Psi(x;\lambda)=\exp\left(w(\lambda)\cdot\beta\right)\cdot J\Psi(\Lambda_\beta x;\lambda)
\]
for all real $\beta$, all $n\ge0$, all $x\in\mathcal{W}^{(n)}$ and all tuples $\lambda$ of $n$ labels.
By the last relation in \eqref{eq:rotation-transport-relations},
and since $(g_{ji}p_\perp)\cdot e_i=p_\perp\cdot\left(g_{ji}^{-1}e_i\right)=p_\perp\cdot e_2=p_2$,
\[
J\mathsf{T}(\kappa)J^{-1}=\int\mathrm{d}\mathfrak{E}_j(\omega,p_\perp)\,\exp\left(-\kappa^0\cdot\omega+i\,\kappa^1\cdot p_i\right)
\qquad\left(\kappa\in\overline{\mathcal{W}}\right):
\]
this operator is formed from $(\mathsf{H}_j,\mathsf{P}_i)$ as $\mathsf{T}(\kappa)$ is formed from $(\mathsf{H}_1,\mathsf{P}_2)$.
The space $\mathcal{H}^{\Psi}_{ji}$ reduces the operators $J\mathsf{T}(\kappa)J^{-1}$, $\kappa\in\overline{\mathcal{W}}$,
and it is invariant under the operators $J\mathcal{B}_{\lambda_0}(\kappa,\kappa')J^{-1}$,
where $\lambda_0=(\sigma,\chi)$ is a label and $\kappa,\kappa'\in\mathcal{W}$.
On $\mathcal{H}^{\Psi}_{ji}$ one has, for real $\beta$,
with $\kappa\in\overline{\mathcal{W}}$ in the first identity and $\kappa,\kappa'\in\mathcal{W}$ in the second,
\begin{align*}
\mathfrak{L}_{ji}(\beta)\,J\mathsf{T}(\kappa)J^{-1}\,\mathfrak{L}_{ji}(\beta)^{-1}&=J\mathsf{T}(\Lambda_\beta\kappa)J^{-1},\\
\mathfrak{L}_{ji}(\beta)\,J\mathcal{B}_{\lambda_0}(\kappa,\kappa')J^{-1}\,\mathfrak{L}_{ji}(\beta)^{-1}
&=\exp\left(w(\sigma)\cdot\beta\right)\cdot J\mathcal{B}_{\lambda_0}(\Lambda_\beta\kappa,\Lambda_\beta\kappa')J^{-1}.
\end{align*}
Together with the representations $\mathcal{U}_j$ and the round cone in each picture, this is (c).

\emph{Clause (d).}
Every $b\in L^2(\R^3;\C^3)$ is ${\ast}f$ for exactly one $f\in L^2(\R^3;\Lambda^2)$, and $g\cdot({\ast}f)={\ast}(g\cdot f)$;
so (a) gives $S_2(g\cdot b,g\cdot b')=S_2(b,b')$ for all $g\in\mathrm{O}(3)$ and $b,b'\in L^2(\R^3;\C^3)$.
Moreover $\mu$ satisfies (G), (dF), $(T_{123})$, $(R_1)$, $(RP_1)$.
Theorem~\ref{thm:two-point}(c) therefore provides exactly one finite positive Borel measure $\mathfrak{K}_\mu$ on $[0,\infty)$
for which the K\"all\'en--Lehmann identity of that theorem holds for all $b,b'\in L^2(\R^3;\C^3)$,
and its mass $\mathfrak{K}_\mu([0,\infty))$ is at most $1$.
For real two-forms $f,f'\in L^2(\R^3;\Lambda^2)$ the fields $b={\ast}f$ and $b'={\ast}f'$ are real,
$S_2\left(\overline{b},b'\right)=S_2(f,f')$, and the identity becomes the formula of (d).
Conversely, let $\mathfrak{K}'$ be a positive Borel measure on $[0,\infty)$ of mass at most $1$
for which the formula of (d) holds for all real $f,f'\in L^2(\R^3;\Lambda^2)$.
The right side of the identity of Theorem~\ref{thm:two-point}(c), formed with $\mathfrak{K}'$ in the place of $\mathfrak{K}_\mu$,
is then defined for all $b,b'\in L^2(\R^3;\C^3)$:
its inner integral lies between $0$ and $\mathfrak{K}'([0,\infty))\le1$,
and $\left|\hat b(p)^{\dagger}P^{\perp}(p)\,\hat b'(p)\right|\le|\hat b(p)|\cdot|\hat b'(p)|$,
which is integrable by the Cauchy--Schwarz inequality and Plancherel's theorem.
Like the left side, this right side is conjugate linear in $b$ and linear in $b'$.
Each field in $L^2(\R^3;\C^3)$ is a real field plus $i$ times a real field,
and every real field is the dual ${\ast}f$ of a real two-form $f$;
so the identity holds with $\mathfrak{K}'$ for all $b,b'\in L^2(\R^3;\C^3)$,
and the uniqueness in Theorem~\ref{thm:two-point}(c) gives $\mathfrak{K}'=\mathfrak{K}_\mu$.

\emph{Clause (e).}
Proposition~\ref{prop:ergodic-components} applies to $\mu$ and to its ergodic decomposition.
Outside a $\nu$-null set of indices $\xi$, the component $\mu_\xi$ is invariant and ergodic under $\R^3$
and lies in $\mathcal{C}_9$ (Proposition~\ref{prop:ergodic-components}(a)),
satisfies (G) with the same right side (Proposition~\ref{prop:ergodic-components}(b)),
is invariant under $\R^3\rtimes\Ocub$ (Proposition~\ref{prop:ergodic-components}(c)),
has $\ker\mathsf{H}_j(\mu_\xi)=\C\,\Omega_j$ for $j=1,2,3$ (Proposition~\ref{prop:ergodic-components}(d)),
and satisfies $\mu_\xi(dF=0)=1$ (Proposition~\ref{prop:ergodic-components}(f), as $\mu(dF=0)=1$).
For such $\xi$ the law $\mu_\xi$ satisfies $(RP_1)$, $(RP_2)$, $(RP_3)$ as a member of $\mathcal{C}_9$;
so it satisfies (H) and (dF) and lies in $\mathcal{C}_9$,
and clauses (a)--(d), proved above for every such law, hold for $\mu_\xi$,
with a K\"all\'en--Lehmann measure $\mathfrak{K}_\xi$ in (d).
Since $\mu_\xi$ lies in $\mathcal{C}_9$ and is ergodic under $\R^3$,
Theorem~\ref{thm:clustering-rate}, in its form for ergodic laws, gives for every $r>0$,
all $\Phi,\Phi'\in L^2(\Sigma_{(-r,r)^3},\mu_\xi)$ and all $a\in\R^3$ with $|a|_\infty\ge C_1\cdot r$
\[
\left|\mu_\xi\left(\overline{\Phi}\cdot U(a)\Phi'\right)-\overline{\mu_\xi(\Phi)}\cdot\mu_\xi(\Phi')\right|
\le C_0\cdot\|\Phi\|_2\cdot\|\Phi'\|_2\cdot\frac{r}{|a|_\infty}.
\]
\end{proof}

\begin{remark}\label{rem:euclidean-soft-limits}
Let $\mu$ be a weak limit on $H_{-4}$ of the laws $\mu^{(N_i)}$ of Proposition~\ref{prop:soft-limits}
along a sequence $N_i\to\infty$.
By clause (a) of that proposition $\mu$ satisfies (H) and (dF), and by its clause (b) it lies in $\mathcal{C}_9$.
The laws $\mu$ of Corollary~\ref{cor:averaged-admissible}(b), formed with the harmonic-averaged law, are limits of this kind.
To see this, let $\mu^{(N_i)}\in\mathcal{V}^{\mathrm{av}}_{N_i}$ converge weakly on $H_{-4}$ to $\mu$, with $N_i\to\infty$;
some subsequence has strictly increasing indices and the same limit, so that one may assume $N_1<N_2<\cdots$.
By the definition of the sets $\mathcal{V}^{\mathrm{av}}_N$ in that corollary,
$\mu^{(N_i)}=\left(\sharp_{\eta_i}\right)_*\nu^{(N_i)}$, where $\eta_i:=L^{-N_i}$, $\sharp_{\eta_i}$ is the step map of Definition~\ref{def:step-map},
and $\nu^{(N_i)}$ is a weak limit point on $\Omega$ of a sequence of averaged laws on the cubic tori with the sides
$8\cdot L^{\lfloor\log_L\bar m^{-1}\rfloor}\cdot L^{T_j+N_i}$, $j\ge1$, where $T_j\to\infty$;
these sides are even, at least $8$, and tend to infinity.
As assumed in (b) of that corollary, $0<\rho<1$;
the standard weights are positive (Definition~\ref{def:chamber-operator}),
the mass parameter is real and the lattice charge $e\cdot\eta_i^{\frac12}$ is positive.
So Theorem~\ref{thm:averaged-any-side} applies to this sequence of sides and shows that $\nu^{(N_i)}$ satisfies (G),
is invariant under $W$ and under $F\mapsto-F$, and, the sides being even, is reflection positive at every mirror of the nine families.
These are the hypotheses of Proposition~\ref{prop:soft-limits} for the laws $\nu^{(N)}$, $N$ in the infinite set $\{N_i:i\ge1\}$,
and $\mu$ is the weak limit of their images $\mu^{(N_i)}$.
By Theorem~\ref{thm:euclidean}\eqref{it:limit-general-laws}, the assertions \eqref{it:limit-euclidean}--\eqref{it:limit-components} of that theorem
therefore hold for every such law $\mu$, with its own $\R^3$-ergodic decomposition.
\end{remark}

\subsection{The Osterwalder--Schrader axioms}\label{sec:os-axioms}

The axioms of Osterwalder and Schrader are conditions on moments taken on tensor test functions.

\begin{definition}[Schwinger functions on tensor tests]\label{def:tensor-schwinger}
The tensor tests are those introduced in Section~\ref{sec:cone} before Lemma~\ref{lem:tensor-tests}:
for an integer $n\ge1$, $\mathscr S_n$ is the space of the Schwartz functions on $\R^{3n}=(\R^3)^n$ with values in the complexification of $(\Lambda^2)^{\otimes n}$, and $\mathscr S_0=\C$.
A point of $\R^{3n}$ is also written $y=(y_1,\dots,y_n)$ with $y_i\in\R^3$.
The seminorms $q_m$ of $\mathscr S_n$ are those of \cite[Sections~2.5 and~7.6]{NoncompactQED}:
\[
q_m(f)=\max_{|\alpha|\le m}\ \sup_{y\in\R^{3n}}\left(1+|y|^2\right)^{\frac m2}\cdot\left|\partial^\alpha f(y)\right|\qquad(m=0,1,2,\dots),
\]
where $\alpha$ ranges over the multi-indices in $\{0,1,2,\dots\}^{3n}$ and $\left|\partial^\alpha f(y)\right|$ is the largest of the moduli of the components of $\partial^\alpha f(y)$.
The product $f_1\otimes\dots\otimes f_n\in\mathscr S_n$ of test forms $f_1,\dots,f_n$, real or complex,
the function $f\otimes f'\in\mathscr S_{k+l}$ for $f\in\mathscr S_k$ and $f'\in\mathscr S_l$,
the map $(\mathfrak g\,\cdot)^{\otimes n}$ of a Euclidean map $\mathfrak g$ of $\R^3$
and the translation $\tau_a^{\otimes n}=(\tau_a\,\cdot)^{\otimes n}$, $a\in\R^3$, are as defined there;
in particular $(\mathfrak g\,\cdot)^{\otimes n}(f_1\otimes\dots\otimes f_n)=(\mathfrak g\cdot f_1)\otimes\dots\otimes(\mathfrak g\cdot f_n)$, and $(\mathfrak g\,\cdot)^{\otimes0}$ is the identity of $\mathscr S_0$.
\begin{enumerate}
\item[(a)] Let $\mu$ be a law on $\mathcal S'$ that satisfies \eqref{cond:gaussian-domination}.
For $n\ge1$ and $f\in\mathscr S_n$, $F^{(n)}(f)$ is the function on $\mathcal S'$ of Lemma~\ref{lem:tensor-tests}(b);
it lies in $L^2(\mu)$, and $f\mapsto F^{(n)}(f)$ is the only continuous linear map from $\mathscr S_n$ into $L^2(\mu)$
with $F^{(n)}(f_1\otimes\dots\otimes f_n)=F(f_1)\cdots F(f_n)$ for all real test forms $f_1,\dots,f_n$.
For $c\in\mathscr S_0$, $F^{(0)}(c)=c$, as in Lemma~\ref{lem:tensor-tests}(c).
The \emph{Schwinger functions} of $\mu$ are the functionals $S_n:=S^\mu_n$ defined after Lemma~\ref{lem:tensor-tests}:
\[
S_n(f)=\int\mathrm{d}\mu\,F^{(n)}(f)\qquad\left(f\in\mathscr S_n,\ n\ge0\right).
\]
Thus $S_0(c)=c$, and $S_n$ is defined on all of $\mathscr S_n$, with no condition on $f$ at the points where two of the variables $y_1,\dots,y_n$ coincide.
\item[(b)] For a unit vector $n\in\R^3$, an integer $m\ge0$ and $f\in\mathscr S_m$ let
\[
\Theta^{n}_mf:=\overline{(r_n\,\cdot)^{\otimes m}f},
\]
where $r_n=r_{n,0}$ is the reflection in the plane $\{n\cdot x=0\}$ and the bar denotes complex conjugation of the components;
in particular $\Theta^n_0c=\overline c$ for $c\in\mathscr S_0$.
For $j\in\{1,2,3\}$ one has $r_{e_j}=\theta_j$, so that $\Theta^{e_j}_m$ is the map $\Theta^{(m)}$ of Section~\ref{sec:cone} formed with the direction $j$.
\item[(c)] For a permutation $\mathfrak s$ of $\{1,\dots,n\}$ and $f\in\mathscr S_n$ let $f^{\mathfrak s}\in\mathscr S_n$ be the function whose value at $y$ is the image of $f\left(y_{\mathfrak s^{-1}(1)},\dots,y_{\mathfrak s^{-1}(n)}\right)$
under the linear map of the complexification of $(\Lambda^2)^{\otimes n}$ that sends $v_1\otimes\dots\otimes v_n$ to $v_{\mathfrak s(1)}\otimes\dots\otimes v_{\mathfrak s(n)}$ for $v_1,\dots,v_n\in\Lambda^2$.
Then $(f_1\otimes\dots\otimes f_n)^{\mathfrak s}=f_{\mathfrak s(1)}\otimes\dots\otimes f_{\mathfrak s(n)}$ for test forms $f_1,\dots,f_n$, real or complex:
the arguments are permuted together with their tensor indices.
\end{enumerate}
\end{definition}

The conditions \eqref{cond:os-growth}, \eqref{cond:os-invariance}, \eqref{cond:os-positivity}, \eqref{cond:os-symmetry} and \eqref{cond:os-cluster} in Part~II of Theorem~\ref{thm:euclidean}
bear the numbers of the axioms of linear growth, Euclidean invariance, positivity, symmetry and clustering in \cite{OsterwalderSchrader1975}.

\begin{proof}[Proof of Theorem~\ref{thm:euclidean}, Part~II]
Every law named in Part~II satisfies \eqref{cond:soft-hypothesis} and \eqref{cond:closedness} and lies in $\mathcal C_9$;
so it suffices to prove Part~II for an arbitrary law $\mu$ on $\mathcal S'$ with these properties.
Indeed, a volume-first limit point is an iterated limit, hence a limit point (Theorem~\ref{thm:continuum-limits}\eqref{it:continuum-soft});
so it satisfies \eqref{cond:soft-hypothesis} by Proposition~\ref{prop:chamber-limits}(c) and \eqref{cond:closedness} by Theorem~\ref{thm:continuum-limits}\eqref{it:continuum-soft},
and it lies in $\mathcal C_9$ by Theorem~\ref{thm:continuum-limits}\eqref{it:continuum-nine-planes}.
Next, let $\mu$ satisfy \eqref{cond:soft-hypothesis} and \eqref{cond:closedness} and lie in $\mathcal C_9$, and let $\mu=\int\mathrm{d}\nu(\xi)\,\mu_\xi$ be its ergodic decomposition under the translations of $\R^3$.
By Proposition~\ref{prop:ergodic-components}(a), (b), (c) and~(f), for $\nu$-almost every $\xi$ the component $\mu_\xi$ is invariant and ergodic under the translations of $\R^3$ and lies in $\mathcal C_9$,
satisfies \eqref{cond:gaussian-domination}, is invariant under $\Ocub$, and satisfies \eqref{cond:closedness}.
Membership in $\mathcal C_9$ includes \eqref{cond:reflection-positivity} for $j=1,2,3$;
so such a component satisfies \eqref{cond:soft-hypothesis}, and, since it is ergodic, the assertions of Part~II on ergodic laws apply to it.

From now on $\mu$ is a law on $\mathcal S'$ with \eqref{cond:soft-hypothesis} and \eqref{cond:closedness} that lies in $\mathcal C_9$.
It satisfies the conditions \textup{(L1)}--\textup{(L4)} of \cite[Section~7.5]{NoncompactQED}.
Indeed, \textup{(L1)} is \eqref{cond:gaussian-domination}.
Next, \textup{(L2)} follows from \eqref{cond:closedness}, since a closed distribution $F$ satisfies $F(d^*u)=0$ for every $u\in\mathcal S(\R^3)$, where $d^*u\in\mathcal S(\R^3;\Lambda^2)$ denotes the two-form whose dual vector field is $-\nabla u$, as in \cite[Section~2.5]{NoncompactQED}.
The condition \textup{(L3)}, the invariance of the characteristic functional of $\mu$ under the translations and under $\Ocub$, follows from the invariance of $\mu$.
Finally, \textup{(L4)} is \eqref{cond:reflection-positivity} for $j=1,2,3$, written with the characteristic functional.
By \cite[Proposition~7.19]{NoncompactQED}, the assertions (1) and (3)--(8) of \cite[Theorem~7.15]{NoncompactQED} and the assertions (1), (2), (4) and (5) of \cite[Theorem~7.17]{NoncompactQED} hold with $\mu$ in the place of the limit point considered there.
In \cite[Section~7.6]{NoncompactQED} the space $\mathscr S_n$ is denoted $\mathcal S_n$;
for $n\ge1$ the Schwinger function $S_n$ formed there with $\mu$ is the unique continuous linear functional on this space
whose value at $f_1\otimes\dots\otimes f_n$ is $\int\mathrm{d}\mu\,F(f_1)\cdots F(f_n)$ for all test forms $f_1,\dots,f_n$ \cite[Theorem~7.17(2)]{NoncompactQED}.
By Lemma~\ref{lem:tensor-tests}(b), the functional $S_n$ of Definition~\ref{def:tensor-schwinger} is continuous and takes the same values at these products, for test forms real or complex;
so the two functionals coincide.

\emph{Linear growth.}
For $n\ge1$ and $f\in\mathscr S_n$, \cite[Theorem~7.17(2)]{NoncompactQED} gives
\[
\left|S_n(f)\right|\le2\cdot66^n\cdot(n!)^{\frac12}\cdot q_{6n}(f),
\]
and this is \eqref{cond:os-growth}.

\emph{Euclidean invariance.}
By Theorem~\ref{thm:euclidean}\eqref{it:limit-euclidean}, which holds for $\mu$ by \eqref{it:limit-general-laws}, $\mathfrak g_*\mu=\mu$ for every Euclidean map $\mathfrak g$ of $\R^3$, reflections included;
for such $\mathfrak g$ one has $\int\mathrm{d}\mu\,\hat{\mathfrak g}\Phi_1=\int\mathrm{d}\mu\,\Phi_1$ for $\Phi_1\in L^1(\mu)$ \cite[Section~8.1.2]{NoncompactQED}.
By Lemma~\ref{lem:tensor-tests}(d), $\hat{\mathfrak g}F^{(n)}(f)=F^{(n)}\left((\mathfrak g\,\cdot)^{\otimes n}f\right)$ at every point of $\mathcal S'$, for $n\ge0$ and $f\in\mathscr S_n$;
as $F^{(n)}(f)$ lies in $L^2(\mu)\subset L^1(\mu)$, it follows that
\[
S_n\left((\mathfrak g\,\cdot)^{\otimes n}f\right)=\int\mathrm{d}\mu\,\hat{\mathfrak g}F^{(n)}(f)=\int\mathrm{d}\mu\,F^{(n)}(f)=S_n(f)\qquad\left(f\in\mathscr S_n\right).
\]
This is \eqref{cond:os-invariance}.
For the maps of $\R^3\rtimes\Ocub$, which are the products of a translation and an element of $\Ocub$, the same identity is also contained in \cite[Theorem~7.17(4)]{NoncompactQED}.

\emph{Positivity.}
Let $n$ be a unit vector, $N\ge0$ an integer, $f_0\in\C$, and $f_k\in\mathscr S_k$ for $1\le k\le N$, with
\[
\supp f_k\subset\left\{y\in\R^{3k}:\ n\cdot y_i\ge0\text{ for }1\le i\le k\right\}.
\]
If $n=e_j$ for some $j\in\{1,2,3\}$, then
\[
\sum_{k,l=0}^NS_{k+l}\left(\Theta^n_kf_k\otimes f_l\right)\ge0
\]
by \cite[Theorem~7.17(5)]{NoncompactQED}:
the reflections of tensor tests used there, formed with the direction $j$, are the maps $\Theta^{e_j}_k$, because $r_{e_j}=\theta_j$, and the supports are restricted there to the same closed half spaces.
For an arbitrary unit vector $n$ choose $g\in\mathrm{O}(3)$ with $ge_1=n$, and put $f'_0:=f_0$ and $f'_k:=(g^{-1}\,\cdot)^{\otimes k}f_k$ for $1\le k\le N$.
Since $f'_k(y)=(\wedge^2g^{-1})^{\otimes k}f_k(gy_1,\dots,gy_k)$ and $n\cdot(gy_i)=(ge_1)\cdot(gy_i)=(y_i)_1$, the support of $f'_k$ lies in $\left\{y:(y_i)_1\ge0\text{ for }1\le i\le k\right\}$.
The reflection $g\theta_1g^{-1}$ fixes the plane $g\{x_1=0\}=\{n\cdot x=0\}$ and maps $n$ to $-n$;
so it is $r_n$.
The maps $(\mathfrak g\,\cdot)^{\otimes k}$ commute with complex conjugation, the maps $\wedge^2$ of real matrices being real, and they act factor by factor on tensor products;
moreover $\left((\mathfrak g_1\mathfrak g_2)\,\cdot\right)^{\otimes k}=(\mathfrak g_1\,\cdot)^{\otimes k}(\mathfrak g_2\,\cdot)^{\otimes k}$ for Euclidean maps $\mathfrak g_1$ and $\mathfrak g_2$,
because $(\mathfrak g_1\mathfrak g_2)\cdot h=\mathfrak g_1\cdot(\mathfrak g_2\cdot h)$ for two-forms $h$ \cite[Section~8.1.2]{NoncompactQED}.
Hence $(g\,\cdot)^{\otimes l}f'_l=f_l$,
\[
(g\,\cdot)^{\otimes k}\,\Theta^{e_1}_kf'_k=\overline{\left((g\theta_1g^{-1})\,\cdot\right)^{\otimes k}f_k}=\Theta^n_kf_k,
\qquad
(g\,\cdot)^{\otimes(k+l)}\left(\Theta^{e_1}_kf'_k\otimes f'_l\right)=\Theta^n_kf_k\otimes f_l,
\]
and \eqref{cond:os-invariance}, applied to the map $g$, yields
\[
\sum_{k,l=0}^NS_{k+l}\left(\Theta^n_kf_k\otimes f_l\right)=\sum_{k,l=0}^NS_{k+l}\left(\Theta^{e_1}_kf'_k\otimes f'_l\right)\ge0,
\]
the inequality being the case $n=e_1$, applied to the tests $f'_k$.
This is \eqref{cond:os-positivity}.

\emph{Symmetry.}
Let $n\ge1$, and let $\mathfrak s$ be a permutation of $\{1,\dots,n\}$.
For real test forms $f_1,\dots,f_n$ the factors $F(f_i)$ are functions on $\mathcal S'$ and commute, so that
\[
S_n\left((f_1\otimes\dots\otimes f_n)^{\mathfrak s}\right)=\int\mathrm{d}\mu\,F(f_{\mathfrak s(1)})\cdots F(f_{\mathfrak s(n)})=\int\mathrm{d}\mu\,F(f_1)\cdots F(f_n)=S_n(f_1\otimes\dots\otimes f_n);
\]
on these products, this is the symmetry of the Schwinger functions stated in \cite[Theorem~7.17(1)]{NoncompactQED}.
The map $f\mapsto f^{\mathfrak s}$ is a continuous linear bijection of $\mathscr S_n$ onto itself, a permutation of the variables followed by a constant linear map of the values that permutes the products of basis vectors;
and $S_n$ is continuous on $\mathscr S_n$.
So the continuous linear functionals $f\mapsto S_n(f^{\mathfrak s})$ and $S_n$ agree on the products of real test forms,
hence on the finite complex linear combinations of these products, which are dense in $\mathscr S_n$ (Section~\ref{sec:cone}, before Lemma~\ref{lem:tensor-tests});
they are equal, and \eqref{cond:os-symmetry} holds.

\emph{Clustering.}
Let $m,n\ge0$, $f\in\mathscr S_m$ and $f'\in\mathscr S_n$, and put
\[
\Phi:=\overline{F^{(m)}(f)},\qquad\Phi':=F^{(n)}(f');
\]
both lie in $L^2(\mu)$, whose scalar product is $\langle\Phi_1,\Phi_2\rangle_\mu=\int\mathrm{d}\mu\,\overline{\Phi_1}\cdot\Phi_2$.
Let $a\in\R^3$.
Since $\mu$ is invariant under the translations, $(\tau_a)_*\mu=\mu$, and the operator $U(a)=\hat\tau_a$ is unitary on $L^2(\mu)$ (Definition~\ref{def:laws}).
The product rule of Lemma~\ref{lem:tensor-tests}(c), followed by Lemma~\ref{lem:tensor-tests}(d), gives at every point of $\mathcal S'$
\[
F^{(m+n)}\left(f\otimes\tau_a^{\otimes n}f'\right)=F^{(m)}(f)\cdot F^{(n)}\left(\tau_a^{\otimes n}f'\right)=F^{(m)}(f)\cdot U(a)F^{(n)}(f'),
\]
and integration with respect to $\mu$ yields
\begin{equation}\label{eq:rotation-cluster-pairing}
S_{m+n}\left(f\otimes\tau_a^{\otimes n}f'\right)=\int\mathrm{d}\mu\,F^{(m)}(f)\cdot U(a)F^{(n)}(f')=\langle\Phi,U(a)\Phi'\rangle_\mu.
\end{equation}

Let $\mu$ be ergodic under the translations of $\R^3$.
The point mass at $\mu$ on the space of laws is carried by the laws that are ergodic under the translations, and it represents $\mu$ as a mixture of such laws;
by the uniqueness assertion of \cite[Theorem~8.1(a)]{NoncompactQED}, taken with $J=\{1,2,3\}$, it is the decomposition measure $\nu$ of $\mu$.
So a property of $\nu$-almost every component of $\mu$ is a property of $\mu$ itself, and Proposition~\ref{prop:ergodic-components}(e), applied to $\Phi$ and $\Phi'$, gives
\[
\langle\Phi,U(a)\Phi'\rangle_\mu\longrightarrow\overline{\int\mathrm{d}\mu\,\Phi}\cdot\int\mathrm{d}\mu\,\Phi'=S_m(f)\cdot S_n(f')
\]
as $|a|\to\infty$ in any manner.
By \eqref{eq:rotation-cluster-pairing} this is the first assertion of \eqref{cond:os-cluster}.

Let, moreover, $r>0$, and let $f$ and $f'$ have compact supports in the open sets $\left((-r,r)^3\right)^m$ and $\left((-r,r)^3\right)^n$, respectively;
for $m=0$ this is no condition on $f$, and for $n=0$ none on $f'$.
Then $\Phi$ and $\Phi'$ lie in $L^2(\Sigma_{(-r,r)^3},\mu)$.
Indeed, for $m\ge1$ Lemma~\ref{lem:tensor-tests}(e), taken with the open set $(-r,r)^3$, gives $F^{(m)}(f)\in L^2(\Sigma_{(-r,r)^3},\mu)$,
and complex conjugation preserves measurability for $\Sigma_{(-r,r)^3}$;
for $m=0$ the function $\Phi=\overline f$ is a constant, and a constant lies in $L^2(\Sigma_U,\mu)$ for every open set $U\subset\R^3$.
The same holds for $\Phi'$.
Theorem~\ref{thm:clustering-rate}, applied to $\Phi$ and $\Phi'$, gives, with its absolute constants $C_0$ and $C_1$,
\[
\left|\langle\Phi,U(a)\Phi'\rangle_\mu-\langle\Phi,P^{123}_\mu\Phi'\rangle_\mu\right|\le C_0\cdot\|\Phi\|_{L^2(\mu)}\cdot\|\Phi'\|_{L^2(\mu)}\cdot\frac{r}{|a|_\infty}\qquad\left(|a|_\infty\ge C_1\cdot r\right),
\]
and, $\mu$ being ergodic, the same theorem gives $\langle\Phi,P^{123}_\mu\Phi'\rangle_\mu=\overline{\int\mathrm{d}\mu\,\Phi}\cdot\int\mathrm{d}\mu\,\Phi'=S_m(f)\cdot S_n(f')$.
Since $\|\Phi\|_{L^2(\mu)}=\|F^{(m)}(f)\|_{L^2(\mu)}$, this inequality together with \eqref{eq:rotation-cluster-pairing} is the bound in \eqref{cond:os-cluster}.

Drop now the assumption of ergodicity.
For all $\Phi_1,\Phi_2\in L^2(\mu)$,
\[
\lim_{|a|_\infty\to\infty}\langle\Phi_1,U(a)\Phi_2\rangle_\mu=\langle\Phi_1,P^{123}_\mu\Phi_2\rangle_\mu.
\]
If $\Phi_1$ and $\Phi_2$ lie in $L^2(\Sigma_{(-r,r)^3},\mu)$ for some $r>0$, this follows from Theorem~\ref{thm:clustering-rate}, which holds for every law in $\mathcal C_9$:
the right side of its bound tends to $0$ as $|a|_\infty\to\infty$.
The union of these spaces over $r>0$ is dense in $L^2(\mu)$.
Indeed, by \cite[Lemma~8.12(i)]{NoncompactQED}, taken with the open set $\R^3$, the finite linear combinations of functions $\exp\left(iF(h)\right)$, with $h$ in a countable set of real elements of $C^{\infty}_{\mathrm c}(\R^3;\Lambda^2)$, are dense in $L^2(\Sigma_{\R^3},\mu)$,
which is all of $L^2(\mu)$, because $\Sigma_{\R^3}=\Sigma$ \cite[Lemma~8.11(ii)]{NoncompactQED};
and such a combination is measurable for $\Sigma_{(-r,r)^3}$ as soon as $(-r,r)^3$ contains the supports of the forms $h$ that occur in it.
Given $\varepsilon>0$, choose $r>0$ and $\Phi_{1,\varepsilon},\Phi_{2,\varepsilon}\in L^2(\Sigma_{(-r,r)^3},\mu)$ with $\|\Phi_1-\Phi_{1,\varepsilon}\|_{L^2(\mu)}\le\varepsilon$ and $\|\Phi_2-\Phi_{2,\varepsilon}\|_{L^2(\mu)}\le\varepsilon$.
The operator $U(a)$ is unitary, and $P^{123}_\mu$ is an orthogonal projection;
so both are contractions, and
\[
\left|\langle\Phi_1,U(a)\Phi_2\rangle_\mu-\langle\Phi_{1,\varepsilon},U(a)\Phi_{2,\varepsilon}\rangle_\mu\right|\le\varepsilon\cdot\|\Phi_2\|_{L^2(\mu)}+\left(\|\Phi_1\|_{L^2(\mu)}+\varepsilon\right)\cdot\varepsilon,
\]
together with the same estimate for $P^{123}_\mu$ in the place of $U(a)$.
It follows that
\[
\limsup_{|a|_\infty\to\infty}\left|\langle\Phi_1,U(a)\Phi_2\rangle_\mu-\langle\Phi_1,P^{123}_\mu\Phi_2\rangle_\mu\right|\le2\varepsilon\cdot\left(\|\Phi_1\|_{L^2(\mu)}+\|\Phi_2\|_{L^2(\mu)}+\varepsilon\right)
\]
for every $\varepsilon>0$, which proves the limit.
Taken with $\Phi_1=\Phi$ and $\Phi_2=\Phi'$ and combined with \eqref{eq:rotation-cluster-pairing}, it shows that, for all $f\in\mathscr S_m$ and $f'\in\mathscr S_n$,
the limit of $S_{m+n}\left(f\otimes\tau_a^{\otimes n}f'\right)$ as $|a|_\infty\to\infty$ is $\left\langle\overline{F^{(m)}(f)},P^{123}_\mu F^{(n)}(f')\right\rangle_\mu$:
the last assertion of \eqref{cond:os-cluster}.

\emph{Reality.}
By Lemma~\ref{lem:tensor-tests}(c), $\overline{F^{(n)}(f)}=F^{(n)}\left(\overline f\right)$ at every point of $\mathcal S'$, for $n\ge0$ and $f\in\mathscr S_n$;
integration with respect to $\mu$ gives $\overline{S_n(f)}=S_n(\overline f)$,
and this is \eqref{it:schwinger-reality}.

\emph{Bianchi identity.}
For $n\ge1$, $u\in\mathcal S(\R^3)$ and $f'\in\mathscr S_{n-1}$, with the two-form $d^*u$ introduced above,
\[
S_n\left(d^*u\otimes f'\right)=0
\]
by \cite[Theorem~7.17(4)]{NoncompactQED}, for complex $u$ by linearity.
This is \eqref{it:schwinger-bianchi}.

\emph{Determinacy.}
Let $\mu'$ be a law on $\mathcal S'$ such that, for every $n\ge1$ and all real $f_1,\dots,f_n\in C^{\infty}_{\mathrm c}(\R^3;\Lambda^2)$, the function $F(f_1)\cdots F(f_n)$ is $\mu'$-integrable and
\[
\int\mathrm{d}\mu'\,F(f_1)\cdots F(f_n)=S_n(f_1\otimes\dots\otimes f_n).
\]
The right side equals $\int\mathrm{d}\mu\,F(f_1)\cdots F(f_n)$, the moment of $\mu$ in the sense of \cite[Theorem~7.15(8)]{NoncompactQED};
by that assertion, $\mu'=\mu$.
In particular, no law other than $\mu$ has the Schwinger functions $S_n$:
this is \eqref{it:schwinger-determinacy}.

\emph{Spectrum and vacuum.}
By Proposition~\ref{prop:zero-momentum}(a) the law $\mu$ satisfies \eqref{cond:cone} for every ordered pair $(j,i)$;
so it satisfies the hypotheses of Theorem~\ref{thm:orthogonal-group}, and by Proposition~\ref{prop:rotations-boosts}(b), in each of the three pictures $e_j$, the joint spectral measure $\mathfrak E_j$ is carried by the round forward cone $\{\omega\ge|p_\perp|\}$.
This is the first clause of Theorem~\ref{thm:euclidean}\eqref{it:limit-spectrum} for the law $\mu$.
Let $\mu$ be ergodic under the translations of $\R^3$.
As shown above for \eqref{cond:os-cluster}, its decomposition measure is the point mass at $\mu$;
so Proposition~\ref{prop:ergodic-components}(d), applied to $\mu$ itself, gives $\ker\mathsf H_j=\C\,\Omega_j$ for $j=1,2,3$,
which is the first assertion of \eqref{it:schwinger-spectrum}.

\emph{The assertion \eqref{it:schwinger-spectrum} at an arbitrary plane.}
Here $\mu$ is a law on $\mathcal{S}'$ with \textup{(H)} and \textup{(dF)} that lies in $\mathcal{C}_9$,
so that the assertions (a) and (b) of Theorem~\ref{thm:euclidean} hold for $\mu$.
The operator $\hat{\mathfrak{g}}$ of a Euclidean map $\mathfrak{g}$ has three properties \cite[Section~8.1.2]{NoncompactQED}:
it is unitary on $L^2(\mu)$ if $\mathfrak{g}_*\mu=\mu$;
it maps $\Sigma_U$-measurable functions to $\Sigma_{\mathfrak{g}(U)}$-measurable ones, for every open set $U\subset\R^3$;
and $\widehat{\mathfrak{g}\circ\mathfrak{g}'}=\hat{\mathfrak{g}}\,\widehat{\mathfrak{g}'}$ for every further Euclidean map $\mathfrak{g}'$.
In particular, if $\mathfrak{g}_*\mu=\mu$, then $\hat{\mathfrak{g}}$ maps $L^2(\Sigma_U,\mu)$ onto $L^2(\Sigma_{\mathfrak{g}(U)},\mu)$, its inverse being the operator of $\mathfrak{g}^{-1}$.
Fix a unit vector $n\in\R^3$ and a number $c\in\R$.
The reflection in the plane $\{n\cdot x=c\}$ is $r_{n,c}\colon x\mapsto x-2\cdot(n\cdot x-c)\cdot n$;
the open half space $\{n\cdot x>c\}$ is one side of this plane.

By Theorem~\ref{thm:euclidean}(b), $(r_{n,c})_*\mu=\mu$, so that $\hat r_{n,c}$ is unitary on $L^2(\mu)$;
as $r_{n,c}\circ r_{n,c}$ is the identity, $\hat r_{n,c}\,\hat r_{n,c}=1$.
Thus $\hat r_{n,c}$ is a unitary involution, hence self-adjoint, and
\[
(\Phi,\Phi')\longmapsto\langle\hat r_{n,c}\Phi,\Phi'\rangle_\mu
\]
is a Hermitian form on $L^2(\Sigma_{\{n\cdot x>c\}},\mu)$;
it is positive semidefinite there, again by Theorem~\ref{thm:euclidean}(b).
The reconstruction at the plane $\{n\cdot x=c\}$ consists of the Hilbert space $\mathcal{H}_{n,c}$,
the completion of the quotient of $L^2(\Sigma_{\{n\cdot x>c\}},\mu)$ by the null vectors of this form,
together with the canonical map $[\,\cdot\,]_{n,c}\colon L^2(\Sigma_{\{n\cdot x>c\}},\mu)\to\mathcal{H}_{n,c}$;
thus
\[
\langle[\Phi]_{n,c},[\Phi']_{n,c}\rangle=\langle\hat r_{n,c}\Phi,\Phi'\rangle_\mu\qquad\left(\Phi,\Phi'\in L^2(\Sigma_{\{n\cdot x>c\}},\mu)\right).
\]
For $n=e_j$ and $c=0$ one has $r_{e_j,0}=\theta_j$, $\{e_j\cdot x>0\}=\{x_j>0\}$ and $\langle\hat\theta_j\Phi,\Phi'\rangle_\mu=\mu\left((\Theta_j\Phi)\cdot\Phi'\right)$;
so $\mathcal{H}_{e_j,0}$ is the space $\mathcal{H}_j$ of the reconstruction in the direction $e_j$, and $[\,\cdot\,]_{e_j,0}=[\,\cdot\,]_j$ \cite[Proposition~8.19\,(i)]{NoncompactQED}.

A Euclidean map carries the picture $e_1$ to the plane $\{n\cdot x=c\}$.
Choose $g\in\mathrm{O}(3)$ with $g\,e_1=n$, and let $\mathfrak{g}:=\tau_{c\cdot n}\circ g$ be the map $x\mapsto g\,x+c\cdot n$.
By Theorem~\ref{thm:euclidean}(a), $\mathfrak{g}_*\mu=\mu$, so that $\hat{\mathfrak{g}}$ is unitary on $L^2(\mu)$.
By \eqref{eq:rotation-plane-conjugation},
$\mathfrak{g}\{x_1>0\}=\{n\cdot x>c\}$, $\mathfrak{g}\,\theta_1\,\mathfrak{g}^{-1}=r_{n,c}$ and $\hat{\mathfrak{g}}\,\hat\theta_1\,\hat{\mathfrak{g}}^{-1}=\hat r_{n,c}$.
The first of these identities shows that $\hat{\mathfrak{g}}$ maps $L^2(\Sigma_1^+,\mu)$ onto $L^2(\Sigma_{\{n\cdot x>c\}},\mu)$.
Moreover $\mathfrak{g}(x+a)=\mathfrak{g}(x)+g\,a$ for $x,a\in\R^3$, that is, $\mathfrak{g}\,\tau_a\,\mathfrak{g}^{-1}=\tau_{g a}$;
as $U(a)=\hat\tau_a$, the identity $\widehat{\mathfrak{g}\circ\mathfrak{g}'}=\hat{\mathfrak{g}}\,\widehat{\mathfrak{g}'}$ gives
\[
\hat{\mathfrak{g}}\,U(a)\,\hat{\mathfrak{g}}^{-1}=U(g\,a)\qquad(a\in\R^3).
\]
Since $\hat r_{n,c}\,\hat{\mathfrak{g}}=\hat{\mathfrak{g}}\,\hat\theta_1$ and $\hat{\mathfrak{g}}$ is unitary, one has, for $\Phi,\Phi'\in L^2(\Sigma_1^+,\mu)$,
\[
\langle\hat r_{n,c}\,\hat{\mathfrak{g}}\Phi,\hat{\mathfrak{g}}\Phi'\rangle_\mu
=\langle\hat{\mathfrak{g}}\,\hat\theta_1\Phi,\hat{\mathfrak{g}}\Phi'\rangle_\mu
=\langle\hat\theta_1\Phi,\Phi'\rangle_\mu
=\mu\left((\Theta_1\Phi)\cdot\Phi'\right)
=\langle[\Phi]_1,[\Phi']_1\rangle.
\]
Consequently $[\Phi]_1\mapsto[\hat{\mathfrak{g}}\Phi]_{n,c}$ is a well defined isometry
from the classes of $L^2(\Sigma_1^+,\mu)$, which are dense in $\mathcal{H}_1$,
onto the classes of $L^2(\Sigma_{\{n\cdot x>c\}},\mu)$, which are dense in $\mathcal{H}_{n,c}$.
It extends to a unitary map $J$ of $\mathcal{H}_1$ onto $\mathcal{H}_{n,c}$.
Now let $t\ge0$ and $y\in e_1^\perp$.
The translation $U(t\cdot e_1+y)$ maps $L^2(\Sigma_1^+,\mu)$ into itself, and
$\exp(-t\cdot\mathsf{H}_1)\,V_1(y)\,[\Phi]_1=[U(t\cdot e_1+y)\Phi]_1$ for $\Phi\in L^2(\Sigma_1^+,\mu)$ \cite[Proposition~8.19\,(ii),\,(iii)]{NoncompactQED}.
Therefore
\begin{equation}\label{eq:rotation-transport}
J\exp(-t\cdot\mathsf{H}_1)\,V_1(y)\,[\Phi]_1=[\hat{\mathfrak{g}}\,U(t\cdot e_1+y)\Phi]_{n,c}=[U(t\cdot n+g\,y)\,\hat{\mathfrak{g}}\Phi]_{n,c}.
\end{equation}

The map $g$ carries $e_1^\perp$ onto $n^\perp$, and one defines
\[
\mathsf{H}_n:=J\,\mathsf{H}_1\,J^{-1},\qquad V_n(y'):=J\,V_1(g^{-1}y')\,J^{-1}\quad(y'\in n^\perp).
\]
The operator $\mathsf{H}_n$ is self-adjoint and non-negative,
$\exp(-t\cdot\mathsf{H}_n)=J\exp(-t\cdot\mathsf{H}_1)\,J^{-1}$, $t\ge0$, is a self-adjoint contraction semigroup on $\mathcal{H}_{n,c}$,
and $y'\mapsto V_n(y')$ is a unitary group on $\mathcal{H}_{n,c}$.
Every $\Phi''\in L^2(\Sigma_{\{n\cdot x>c\}},\mu)$ is of the form $\hat{\mathfrak{g}}\Phi$ with $\Phi\in L^2(\Sigma_1^+,\mu)$, and then $J[\Phi]_1=[\Phi'']_{n,c}$;
with $y:=g^{-1}y'$, the identity \eqref{eq:rotation-transport} becomes
\begin{equation}\label{eq:rotation-plane-translations}
\exp(-t\cdot\mathsf{H}_n)\,V_n(y')\,[\Phi'']_{n,c}=[U(t\cdot n+y')\Phi'']_{n,c}\qquad\left(t\ge0,\ y'\in n^\perp,\ \Phi''\in L^2(\Sigma_{\{n\cdot x>c\}},\mu)\right).
\end{equation}
On classes, $\exp(-t\cdot\mathsf{H}_n)$ and $V_n(y')$ thus act as the translations by $t\cdot n$ along $n$ and by $y'$ in $n^\perp$.

Let $\mathfrak{E}_1$ be the joint spectral measure of the picture $e_1$, a projection-valued measure on $[0,\infty)\times e_1^\perp$;
its variable $p_\perp$ is read as a vector of $e_1^\perp\subset\R^3$.
By Lemma~\ref{lem:joint-spectrum}(b),
\[
\exp(-t\cdot\mathsf{H}_1)\,V_1(y)=\int\mathrm{d}\mathfrak{E}_1(\omega,p_\perp)\,\exp(-t\cdot\omega+i\,p_\perp\cdot y)\qquad(t\ge0,\ y\in e_1^\perp).
\]
For a Borel set $A\subset[0,\infty)\times n^\perp$ put
\[
\mathfrak{E}_n(A):=J\,\mathfrak{E}_1\left(\left\{(\omega,p_\perp):\ (\omega,g\,p_\perp)\in A\right\}\right)J^{-1}.
\]
This is a projection-valued measure on $[0,\infty)\times n^\perp$.
Since $(g\,p_\perp)\cdot y'=p_\perp\cdot(g^{-1}y')$, the operators of \eqref{eq:rotation-plane-translations} have the spectral representation
\[
\exp(-t\cdot\mathsf{H}_n)\,V_n(y')=\int\mathrm{d}\mathfrak{E}_n(\omega,p')\,\exp(-t\cdot\omega+i\,p'\cdot y')\qquad(t\ge0,\ y'\in n^\perp).
\]
Likewise, the first identity of Lemma~\ref{lem:joint-spectrum}(b) and $\exp(i\,t\cdot\mathsf{H}_n)=J\exp(i\,t\cdot\mathsf{H}_1)\,J^{-1}$, $t\in\R$, give
\[
\exp(i\,t\cdot\mathsf{H}_n)\,V_n(y')=\int\mathrm{d}\mathfrak{E}_n(\omega,p')\,\exp(i\,t\cdot\omega+i\,p'\cdot y')\qquad(t\in\R,\ y'\in n^\perp);
\]
so $\mathfrak{E}_n$ is the joint spectral measure of $\mathsf{H}_n$ and $V_n$,
which this identity determines by the uniqueness of the Fourier transform of a finite measure.
The measure $\mathfrak{E}_1$ is carried by the round cone $\{\omega\ge|p_\perp|\}$ (Proposition~\ref{prop:rotations-boosts}(b)), and $|g\,p_\perp|=|p_\perp|$;
hence $\mathfrak{E}_n$ is carried by $\{\omega\ge|p'|\}$.

Finally, $\ker\mathsf{H}_n=J\ker\mathsf{H}_1$ by the definition of $\mathsf{H}_n$.
Let $\mu$ be ergodic under the translations.
Then $\ker\mathsf{H}_1=\C\,\Omega_1$, as was derived above from Proposition~\ref{prop:ergodic-components}(d), applied to $\mu$ itself,
and $\Omega_1=[1]_1$ \cite[Proposition~8.19\,(i)]{NoncompactQED}.
As $\hat{\mathfrak{g}}1=1$, the vector $J\Omega_1$ is $[1]_{n,c}$, and therefore $\ker\mathsf{H}_n=\C\,[1]_{n,c}$.
\end{proof}

\begin{example}[The zero field and the axioms]\label{ex:zero-field-axioms}
Let $\mu$ be a volume-first limit point and $\mu_\xi$ its ergodic components, as in Theorem~\ref{thm:euclidean}.
All the estimates that the theorem gives for these laws are upper bounds:
Gaussian domination \textup{(G)} bounds the Laplace transform of $\mu$, the round forward cone bounds the momenta by the energy,
the rate of clustering bounds the truncated correlations, and $\mathfrak{K}_\mu([0,\infty))\le1$ bounds the mass of the K\"all\'en--Lehmann measure.
Let $\mu_0$ be the point mass at the field $F=0$, as in Example~\ref{ex:zero-field}.
It has every property that Theorem~\ref{thm:euclidean} asserts of $\mu$ and of the $\mu_\xi$,
and its K\"all\'en--Lehmann measure is $\mathfrak{K}_{\mu_0}=0$.
Hence these properties imply neither $\mu\neq\mu_0$ nor $\mathfrak{K}_\mu\neq0$.
\end{example}

\begin{proof}
For every real test form $f$ one has $F(f)=0$ $\mu_0$-almost surely.
Hence
\[
\mu_0\left(\exp(F(f))\right)=1\le\exp\left(\frac{1}{2}\|f\|_2^2\right),
\]
which is \textup{(G)}, and $\mu_0(dF=0)=1$, which is \textup{(dF)}.
Every Euclidean map fixes the field $F=0$;
so $\mu_0$ is invariant under all Euclidean maps, reflections included.
For a plane $\{n\cdot x=c\}$, forms $f_1,\dots,f_k\in C^\infty_{\mathrm c}(\{n\cdot x>c\};\Lambda^2)$ and numbers $z_1,\dots,z_k\in\C$,
\[
\sum_{p,q=1}^{k}\overline{z_p}\cdot z_q\cdot\mu_0\left(\exp\left(i\,F(f_q)-i\,F(r_{n,c}\cdot f_p)\right)\right)=\left|\sum_{p=1}^{k}z_p\right|^2\ge0;
\]
thus $\mu_0$ is reflection positive at every plane.
In particular $\mu_0$ satisfies \textup{(H)} and \textup{(dF)} and lies in $\mathcal{C}_9$, so that the assertions of Theorem~\ref{thm:euclidean} on such laws hold for $\mu_0$.

Every element of $L^2(\mu_0)$ is almost surely constant: $L^2(\mu_0)=\C$.
Therefore $\mu_0$ is ergodic under the translations, and for $j=1,2,3$ one has $\mathcal{H}_j=\C\,\Omega_j$, $\mathsf{H}_j=0$, and $\mathsf{P}_i=0$ for $i\neq j$.
The joint spectral measure $\mathfrak{E}_j$ is then concentrated at the origin of $[0,\infty)\times e_j^\perp$, a point of the round cone;
$\ker\mathsf{H}_j=\C\,\Omega_j$;
the unitary representations in Theorem~\ref{thm:euclidean}(c) fix $\Omega_j$ and so act as the identity;
and the left side of the clustering bound vanishes, because $\mu_0\left(\overline{\Phi}\cdot U(a)\Phi'\right)=\overline{\mu_0(\Phi)}\cdot\mu_0(\Phi')$ for constants $\Phi$, $\Phi'$.
In the same way $\mathcal{H}_{n,c}=\C\,[1]_{n,c}$ at every plane $\{n\cdot x=c\}$.
Finally, every Schwinger function with at least one argument vanishes;
in particular $S_2=0$, and the measure of the K\"all\'en--Lehmann form in Theorem~\ref{thm:euclidean}(d) is $\mathfrak{K}_{\mu_0}=0$.
\end{proof}

\subsection{Comparison with four-dimensional Yang--Mills theory}\label{sec:yang-mills-rotations}

\begin{remark}[Comparison with four-dimensional Yang--Mills theory]\label{rem:yang-mills-comparison}
The argument for four-dimensional Yang--Mills theory concerns a scalar field, $\operatorname{tr}F^2$, on $\R^4$,
smeared with test functions whose supports are pairwise disjoint.
Its steps compare with those of this section as follows.
\begin{enumerate}
\item
\emph{Functionals and moments.}
In that argument the functional is defined on tuples of separated test functions of class $C^1$ with compact supports,
and it comes with constants that depend on the number of entries and on further parameters.
For a law with \textup{(G)}, the moments $S_N$ are bounded symmetric $N$-linear forms on all of $L^2(\R^3;\C^3)^N$,
with $|S_N(b_1,\dots,b_N)|\le c_N^N\cdot\prod_{i=1}^N\|b_i\|$ (Lemma~\ref{lem:domination-consequences}\eqref{it:moment-forms});
no condition is imposed on the supports.
The moments are centred: $S_1=0$.
A Euclidean map $\mathfrak g\colon x\mapsto gx+a$ acts on the fields $b$ as on pseudovectors,
$(\mathfrak g\cdot b)(x)=\det(g)\cdot g\,b(\mathfrak g^{-1}x)$ (Definition~\ref{def:field-actions}).

\item
\emph{The insertion bound.}
That argument assumes it, or derives it from the lattice.
Here it is Proposition~\ref{prop:insertion-bound},
which follows from Gaussian domination \textup{(G)} by multiple reflections (Lemma~\ref{lem:source-operators}).
The bound is in the norm of $L^2(\R^3;\C^3)$ and contains no negative power of the width $l$ of the slab:
$\|\mathcal B_l(f)\|\le k_1\cdot\|f\|$ for every $l>0$ and every $f\in L^2_{[0,l]}$, and the constant $k_1$ does not depend on $l$.

\item
\emph{The plate bound.}
It is Proposition~\ref{prop:point-plates}, and it holds with the exponent one:
the bound in the last assertion of that proposition contains $d-d'$ only through the factor $(d-d')^{-1}$.
The objects attached to points are the point plates $\mathcal B_\lambda(\xi,\xi')$,
bounded operators on $\mathcal H_1$ that depend holomorphically on $(\xi,\xi')\in\mathcal W\times\mathcal W$.
They take the place of the limits of smeared products that Montel's theorem provides in that argument.

\item
\emph{Kernels of the plane.}
These are the kernels $G_N(\xi;\lambda)$ of Definition~\ref{def:plate-kernels},
holomorphic on $\mathcal W^{N-1}$ by Proposition~\ref{prop:kernels}.
They carry labels $\lambda_a=(\sigma_a,\chi_a)$,
and for $N\ge2$ the bound of Proposition~\ref{prop:kernels}\eqref{it:kernel-bound},
with its factor $\left(\min_a d_{\mathcal W}(\xi_a)\right)^{-N}$, holds at all scales.
The reality of the kernels takes the form $\overline{G_N(\xi;\lambda)}=G_N(\overline{\xi};\bar\lambda)$ of Proposition~\ref{prop:kernels}\eqref{it:kernel-reality}.

\item
\emph{The Gram kernel at boosted points.}
It follows at once from the plate vectors $\Psi(\Lambda_\beta x;\lambda)$.
By Lemma~\ref{lem:forward-tube}\eqref{it:tube-boosts} the boost $\Lambda_\beta$ maps $\mathcal W$ onto itself and commutes with $\xi\mapsto\xi^*$,
so the Gram formula of Proposition~\ref{prop:kernels}\eqref{it:plate-vectors} applies to boosted configurations as it stands.
The antilinear reflection $\Theta_1$ maps the field smeared with a plate with the label $\lambda_a=(\sigma_a,\chi_a)$
to $s(\sigma_a)$ times the field smeared with a plate with the same label (Lemma~\ref{lem:plate-identities}\eqref{it:plate-reflection}).
In that formula, therefore, every label of $\Psi(x;\lambda)$ reappears unchanged on the right side,
and the only factor contributed by the labels is the sign $s(\lambda)$.

\item
\emph{Finitely many harmonics.}
Here the tubes come first.
The harmonics $c_m(\,\cdot\,;\lambda)$ of Lemma~\ref{lem:finitely-many-harmonics} are holomorphic functions on the union $\mathcal Z_N$ of the tubes,
on which the one function $\mathcal G_N$ of Lemma~\ref{lem:one-function} is defined,
and they are the coefficients of the expansion~\eqref{eq:rotation-boost-expansion} of $\mathcal G_N(\Lambda_\beta Z;\lambda)$ in the exponentials $\exp(-(w+4m)\beta)$, $w=w(\lambda)$.
The harmonic $c_m$ has the frequency $k=w+4m$, and $c_m=0$ unless $|k|\le N$.
A harmonic is indexed by its anomaly index $m$, not by its frequency,
and $\mathfrak a(N)\le N/2$ (Definition~\ref{def:anomaly-index}\eqref{it:anomaly-linear-bound}).

\item
\emph{Convexity.}
Lemma~\ref{lem:anomaly-convexity} gives $2\mathfrak a(a+b)\le\mathfrak a(2a)+\mathfrak a(2b)$ for all integers $a,b\ge0$.
A strictly time-ordered configuration of $a+b$ points is translated in time so that its first $a$ points lie in $\{y^0<0\}$ and its last $b$ points in $\{y^0>0\}$;
the first $a$ points, reflected in the line $\{y^0=0\}$, form $x$ with the labels $\lambda^x$,
the last $b$ points form $z$ with the labels $\lambda^z$,
and $w_x=w(\lambda^x)$, $w_z=w(\lambda^z)$.
The Cauchy--Schwarz inequality for the plate vectors $\Psi(\Lambda_\beta x;\lambda^x)$ and $\Psi(\Lambda_\beta z;\lambda^z)$
is multiplied by $\exp(2w\beta)=\exp(2w_x\beta)\cdot\exp(2w_z\beta)$.
The factors match:
the two doubled configurations on the right of the inequality have the weights $2w_x$ and $2w_z$, whose mean is the weight $w=w_x+w_z$,
because in each of them a point and its mirror image carry the same label (Lemma~\ref{lem:plate-identities}\eqref{it:plate-reflection}).
The inequality is obtained directly at strictly time-ordered configurations, without pencils,
that is, without real configurations whose consecutive differences have polar angles in an interval $(-\vartheta_0,\vartheta_0)$ with $0<\vartheta_0<\pi/4$.

\item
\emph{One function on all tubes.}
Lemma~\ref{lem:one-function} glues the functions $G^{\mathfrak s}_\upsilon$ on the tubes $\mathcal Z(\mathfrak s,\upsilon)$,
$\mathfrak s$ a permutation and $\upsilon\in(\pi/2)\Z$ a frame, into one function $\mathcal G_N$ on $\mathcal Z_N$.
These functions carry the phases $\exp(iw\upsilon)$,
and the quarter turn acts by the factor $i^w$:
$\mathcal G_N(R_{\pi/2}Z;\lambda)=i^w\cdot\mathcal G_N(Z;\lambda)$ (Lemma~\ref{lem:one-function}\eqref{it:tube-symmetries}).

\item
\emph{Harmonics on the tubes, their covariance, and their extension to all configurations.}
These are Lemmas~\ref{lem:finitely-many-harmonics} and~\ref{lem:exact-phases} and the properties of Lemma~\ref{lem:harmonic-extension}.
The harmonic of index $m$ turns with an exact phase:
$\tilde c_m(R_tX;\lambda)=\exp(i(w+4m)t)\cdot\tilde c_m(X;\lambda)$ for every real $t$ (Lemma~\ref{lem:harmonic-extension}\eqref{it:harmonic-phase}).
Among the properties of the harmonics in that argument, one bears on their reality and on the reflection;
its place is taken by the conjugation rule $\overline{c_m(X;\lambda)}=c_{-m}(X;\bar\lambda)$ at real configurations $X$ (Lemma~\ref{lem:finitely-many-harmonics}).

\item
\emph{Positivity of the top harmonic.}
The positivity is Lemma~\ref{lem:top-harmonic}.
For labelled configurations $(x;\lambda)$ and $(z;\lambda')$ of $a$ points of the half plane $\mathbb H$ with pairwise distinct times,
the Gram kernel of the boosted plate vectors is multiplied, before the limit $\beta\to+\infty$, by $\exp(w(\lambda)\beta)\cdot\exp(w(\lambda')\beta)$,
the product of a positive function of $(x;\lambda)$ and the same function of $(z;\lambda')$,
which keeps the kernel positive semidefinite.
Let now $L$ be a mirror line of the plane in an arbitrary direction, with the unit normal $n_\upsilon$ that points into one of its open sides,
$\upsilon$ being here an arbitrary real number, not necessarily a frame,
and let $(x;\lambda)$, $(z;\lambda')$ be labelled configurations of $a$ distinct points on that side.
The top harmonic at $(\theta_Lx,z)$ is then a value of the positive semidefinite kernel $K_a$,
multiplied by $s(\lambda)$ and by the unimodular factor $\exp(i(w(\lambda)+w(\lambda')-4\mathfrak a(2a))\upsilon)$ (Lemma~\ref{lem:top-harmonic}\eqref{it:top-harmonic-mirror}).
For the scalar field of that argument the corresponding factor is a constant;
here it depends on the weights $w(\lambda)$ and $w(\lambda')$ of the two labelled configurations.

\item
\emph{Selection, halving, regular polygons, arithmetic.}
These are Lemma~\ref{lem:selection-halving}, Proposition~\ref{prop:polygons} and Lemma~\ref{lem:convex-sequences}.
The polygons carry equal labels.
For a regular $2a$-gon whose vertices all carry the label $(\sigma,\chi)$ the weight is $w=2a\cdot w(\sigma)$;
the rotation by $\pi/a$ about the centre permutes the labelled vertices,
so the harmonic of index $m$ vanishes there unless $\exp\left(\frac{i(w+4m)\pi}{a}\right)=1$ (Lemma~\ref{lem:selection-halving}\eqref{it:selection}).
Since $\exp(2\pi i\cdot w(\sigma))=1$, one has
$\exp\left(\frac{i(w+4m)\pi}{a}\right)=\exp(2\pi i\cdot w(\sigma))\cdot\exp\left(\frac{4\pi im}{a}\right)=\exp\left(\frac{4\pi im}{a}\right)$,
so that the weight no longer enters the condition,
and the divisibility is the same as in that argument, namely $a$ divides $2\mathfrak a(2a)$ for every $a\ge1$.

\item
\emph{The input from the two-point function.}
That argument obtains it from the diagonal mirrors.
Here it is Theorem~\ref{thm:two-point}.
The frequencies of the plate kernels of two points lie in the window $|k|\le2$ of~\eqref{eq:rotation-window},
which comes from the bound that \textup{(G)} gives on $L^2$.
In the presence of the quarter turn, the only frequency of this window that breaks the rotations is $k=-w$, $w=w(\sigma)+w(\sigma')$,
for the pairs of labels with $\sigma=\sigma'=+$ and with $\sigma=\sigma'=-$.
The Bianchi identity~\eqref{eq:rotation-bianchi} removes this one anomalous term:
it forces its coefficient to be a constant multiple of $r^2$,
and the bound by a multiple of $r^{-2}$ for all $r>0$ makes the constant vanish.

\item
\emph{Cluster moves and the singular part.}
Lemma~\ref{lem:cluster-moves} and Proposition~\ref{prop:no-singular-part} carry this step.
Everything takes place in the picture of $e_1$, that is, in the reconstruction at $\{x_1=0\}$,
and the clusters are moved along $e_2$ only.

\item
\emph{Separated test functions and the full group.}
That argument proceeds to separated test functions in $\R^4$ and to the full group.
Here this step is replaced by Proposition~\ref{prop:no-contact-terms} and Theorems~\ref{thm:one-plane} and~\ref{thm:orthogonal-group}.
No cluster is moved in the picture of the axis transverse to the plane.
Coinciding points are included:
the difference of the rotated and the unrotated moments is a bounded multilinear form on $L^2(\R^3;\C^3)$
that vanishes on gradients by \textup{(dF)} and on entries over pairwise disjoint compact subsets of the plane,
hence vanishes by Proposition~\ref{prop:no-contact-terms}.
The invariance is obtained for the moments themselves, and no step leads from moments to cumulants.
\end{enumerate}
\end{remark}

\begin{thebibliography}{99}

\bibitem{Billingsley1995}
P.~Billingsley, \emph{Probability and measure}, third edition,
John Wiley \& Sons, New York, 1995.

\bibitem{Billingsley1999}
P.~Billingsley, \emph{Convergence of probability measures}, second edition,
John Wiley \& Sons, New York, 1999.

\bibitem{Bourbaki1968}
N.~Bourbaki, \emph{Groupes et alg\`ebres de Lie. Chapitres~IV, V et~VI},
Hermann, Paris, 1968.

\bibitem{Dudley2002}
R.~M.~Dudley, \emph{Real analysis and probability},
Cambridge University Press, Cambridge, 2002.

\bibitem{Feller1971}
W.~Feller, \emph{An introduction to probability theory and its applications. Vol.~II},
second edition, John Wiley \& Sons, New York, 1971.

\bibitem{GlimmJaffe}
J.~Glimm and A.~Jaffe, \emph{Quantum physics. A functional integral point of view},
second edition, Springer-Verlag, New York, 1987.

\bibitem{Kallenberg2002}
O.~Kallenberg, \emph{Foundations of modern probability}, second edition,
Springer-Verlag, New York, 2002.

\bibitem{NoncompactQED}
\emph{Noncompact lattice QED in three dimensions: limit points of the curvature field
and their fourth cumulant} (companion paper).

\bibitem{OsterwalderSchrader1975}
K.~Osterwalder and R.~Schrader, \emph{Axioms for Euclidean Green's functions~II},
Commun.\ Math.\ Phys.\ 42 (1975) 281--305.

\bibitem{ReedSimon1980}
M.~Reed and B.~Simon, \emph{Methods of modern mathematical physics.~I. Functional analysis},
second edition, Academic Press, New York, 1980.

\bibitem{Rudin1962}
W.~Rudin, \emph{Fourier analysis on groups},
Interscience Publishers, New York, 1962.

\bibitem{Rudin1976}
W.~Rudin, \emph{Principles of mathematical analysis}, third edition,
McGraw-Hill, New York, 1976.

\bibitem{Rudin1987}
W.~Rudin, \emph{Real and complex analysis}, third edition,
McGraw-Hill, New York, 1987.

\bibitem{Rudin1991}
W.~Rudin, \emph{Functional analysis}, second edition,
McGraw-Hill, New York, 1991.

\bibitem{Simon2005}
B.~Simon, \emph{Trace ideals and their applications}, second edition,
American Mathematical Society, Providence, RI, 2005.

\bibitem{SteinWeiss1971}
E.~M.~Stein and G.~Weiss, \emph{Introduction to Fourier analysis on Euclidean spaces},
Princeton University Press, Princeton, NJ, 1971.

\end{thebibliography}

\end{document}
